%% =====================================================================
%%  Transseries: the polynomial--logarithmic calculus and its inversions
%%  Single consolidated volume for drafts/series-and-transseries/.
%%  This file is now the canonical source: the five source groups it was
%%  assembled from were merged into it and deleted on 2026-09-04, and the
%%  assembler cannot be re-run.  Edit here.
%% =====================================================================
\documentclass[11pt,openany]{book}
\usepackage[a4paper,margin=26mm,headheight=15pt]{geometry}
\usepackage[T1]{fontenc}
\usepackage[utf8]{inputenc}
% ---------------------------------------------------------------------
%  Lean crosswalk apparatus.
%
%  \lean{...} names a Lean declaration or module.  It is \detokenize'd inside
%  math and \nolinkurl'd outside, so it is safe in BOTH modes and needs no
%  escaping of underscores.  Prefer it to \path for declaration names: \path is
%  verbatim-like and is fatal inside math mode ("Command $ invalid in math
%  mode").  Existing \path{...} rows in this file are all outside math and are
%  left alone.
%
%  Lean names carry Unicode subscripts (eval-sub-2, and so on).  Inside \lean
%  and \path they are verbatim, but a \texttt row typesets them normally and
%  inputenc then refuses an undeclared character, so declare them here.
% ---------------------------------------------------------------------
\newcommand{\lean}[1]{%
  \ifmmode\text{\texttt{\detokenize{#1}}}\else\nolinkurl{#1}\fi}
\DeclareUnicodeCharacter{2080}{\ensuremath{_0}}
\DeclareUnicodeCharacter{2081}{\ensuremath{_1}}
\DeclareUnicodeCharacter{2082}{\ensuremath{_2}}
\DeclareUnicodeCharacter{2083}{\ensuremath{_3}}
\DeclareUnicodeCharacter{2084}{\ensuremath{_4}}
\DeclareUnicodeCharacter{2099}{\ensuremath{_n}}
\IfFileExists{libertinus.sty}{\usepackage{libertinus}}{\usepackage{lmodern}}
\usepackage{microtype}
\usepackage{amsmath,amssymb,amsthm,mathtools,mathrsfs,bm}
% Canonical mathematical notation for active FabiusFunction LaTeX documents.
%
% This file is intentionally a plain TeX input rather than a LaTeX package.
% Every standalone document loads it by an explicit path relative to that
% document's directory; included fragments inherit the definitions from their
% root document.  Keep this file independent of document layout and metadata.
%
% Contract:
%   * one command has one mathematical meaning throughout the active corpus;
%   * names encode semantic roles rather than a document-local choice of letter;
%   * visually similar but differently normalized objects have different print;
%   * every command below is defined normatively in the notation catalogue.
%
% The active corpus excludes docs/archive/.  Historical sources may retain
% their original notation only inside an explicitly labelled quotation.

\makeatletter
\@ifundefined{mathbb}{%
  \PackageError{fabius-notation}{The amssymb package must be loaded first}%
    {Load amsmath and amssymb before inputting fabius-notation.tex.}%
}{}
\@ifundefined{operatorname}{%
  \PackageError{fabius-notation}{The amsmath package must be loaded first}%
    {Load amsmath before inputting fabius-notation.tex.}%
}{}
\makeatother

% Version stamp used by the catalogue and by static audits.
\newcommand{\FabiusNotationVersion}{2026-08-31}

% Number systems.  NaturalNumbers is explicitly the nonnegative convention.
\newcommand{\NaturalNumbers}{\mathbb{N}_{0}}
\newcommand{\PositiveIntegers}{\mathbb{N}_{>0}}
\newcommand{\IntegerNumbers}{\mathbb{Z}}
\newcommand{\RationalNumbers}{\mathbb{Q}}
\newcommand{\RealNumbers}{\mathbb{R}}
\newcommand{\ComplexNumbers}{\mathbb{C}}
\newcommand{\ExtendedRealNumbers}{\overline{\mathbb{R}}}

% The three principal Fabius--Rvachev functions.  Subscripts are deliberate:
% a bare F and a calligraphic F have several unrelated uses in the corpus.
\newcommand{\FabiusBounded}{F_{\mathrm{bd}}}
\newcommand{\FabiusGlobal}{F_{\mathrm{gl}}}
\newcommand{\FabiusQuantile}{Q_{F}}
\newcommand{\FabiusClampedQuantile}{Q_{F,\mathrm{cl}}}
\newcommand{\RvachevUp}{\operatorname{up}}

% Constants, elementary operators, and delimiters.
\newcommand{\EulerE}{\mathrm{e}}
\newcommand{\ImaginaryUnit}{\mathrm{i}}
\newcommand{\Differential}{\mathop{}\!\mathrm{d}}
\newcommand{\AbsoluteValue}[1]{\left\lvert #1\right\rvert}
\newcommand{\Norm}[1]{\left\lVert #1\right\rVert}
\newcommand{\Floor}[1]{\left\lfloor #1\right\rfloor}
\newcommand{\Ceiling}[1]{\left\lceil #1\right\rceil}
\newcommand{\FractionalPart}[1]{\left\{#1\right\}}
\newcommand{\PositivePart}[1]{\left(#1\right)_{+}}
\newcommand{\IndicatorOf}[1]{\mathbf{1}_{#1}}
\newcommand{\SetOf}[1]{\left\{#1\right\}}
\newcommand{\InnerProduct}[1]{\left\langle #1\right\rangle}
\newcommand{\CardinalityOf}[1]{\left\lvert #1\right\rvert}
\newcommand{\DefinitionEquals}{\mathrel{\coloneqq}}
\newcommand{\Sign}{\operatorname{sgn}}
\newcommand{\IdentityOperator}{\operatorname{id}}
\newcommand{\SpanOperator}{\operatorname{span}}
\newcommand{\TraceOperator}{\operatorname{tr}}
\newcommand{\RankOperator}{\operatorname{rank}}
\newcommand{\DiagonalOperator}{\operatorname{diag}}
\newcommand{\ResidueOperator}{\operatorname*{Res}}
\newcommand{\LeastCommonMultiple}{\operatorname{lcm}}
\newcommand{\OrderOperator}{\operatorname{ord}}
\newcommand{\PrincipalLogarithm}{\operatorname{Log}}
\newcommand{\FallingFactorial}[2]{\left(#1\right)^{\underline{#2}}}
\newcommand{\RisingFactorial}[2]{\left(#1\right)^{\overline{#2}}}

% Support, topology, convolution, and metric notation.
\newcommand{\SupportOperator}{\operatorname{supp}}
\newcommand{\SupportOf}[1]{\SupportOperator\!\left(#1\right)}
\newcommand{\TopologicalSupportOf}[1]{\operatorname{tsupp}\!\left(#1\right)}
\newcommand{\Convolution}{\mathbin{*}}
\newcommand{\MetricDistance}{\operatorname{dist}}
\newcommand{\TotalVariationDistance}{d_{\mathrm{TV}}}

% Probability.  Equality and convergence in law are not metric distance.
\newcommand{\Probability}{\mathbb{P}}
\newcommand{\Expectation}{\mathbb{E}}
\newcommand{\Variance}{\operatorname{Var}}
\newcommand{\Covariance}{\operatorname{Cov}}
\newcommand{\LawOperator}{\operatorname{Law}}
\newcommand{\LawOf}[1]{\LawOperator\!\left(#1\right)}
\newcommand{\EqualInLaw}{\mathrel{\overset{\mathrm{law}}{=}}}
\newcommand{\ConvergesInLaw}{\mathrel{\xrightarrow{\mathrm{law}}}}

% Fourier and integral transforms.  The printed labels expose normalization.
% FourierTwoPi uses exp(-2*pi*i*x*xi); FourierAngular uses exp(-i*x*omega).
\newcommand{\FourierTwoPi}[1]{\widehat{#1}_{\,2\pi}}
\newcommand{\FourierAngular}[1]{\widehat{#1}_{\,\mathrm{ang}}}
\newcommand{\LaplaceTransformSymbol}{\mathcal{L}}
\newcommand{\MellinTransformSymbol}{\mathcal{M}}
\newcommand{\LaplaceTransformOf}[1]{\LaplaceTransformSymbol\!\left[#1\right]}
\newcommand{\MellinTransformOf}[1]{\MellinTransformSymbol\!\left[#1\right]}
\newcommand{\MomentGeneratingFunction}[1]{M_{#1}}
\newcommand{\CharacteristicFunction}[1]{\varphi_{#1}}

% The two standard sinc functions must never share one symbol.
\newcommand{\SincRad}{\operatorname{sinc}_{\mathrm{rad}}}
\newcommand{\SincPi}{\operatorname{sinc}_{\pi}}

% Binary arithmetic and Thue--Morse notation.
\newcommand{\BinaryDigitSum}{\operatorname{wt}_{2}}
\newcommand{\ThueMorseSignSymbol}{\varepsilon}
\newcommand{\ThueMorseSign}[1]{\varepsilon_{#1}}
\newcommand{\PadicValuation}[1]{\nu_{#1}}
\newcommand{\TwoAdicValuation}{\nu_{2}}

% q-series notation.  Every base is an explicit argument.
\newcommand{\QPochhammer}[3]{\left(#1;#2\right)_{#3}}
\newcommand{\GaussianBinomial}[3]{%
  \genfrac{[}{]}{0pt}{}{#1}{#2}_{#3}%
}
\newcommand{\QInteger}[2]{\left[#1\right]_{#2}}

% Named special functions.
\newcommand{\Polylogarithm}{\operatorname{Li}}
\newcommand{\LambertW}{W}
\newcommand{\GammaFunction}{\Gamma}
\newcommand{\RiemannZeta}{\zeta}

% Asymptotics.  BigO and LittleO are operator tokens for existing formula
% grammar; prefer the At forms so that the limiting regime is printed.
\newcommand{\BigO}{\mathcal{O}}
\newcommand{\LittleO}{o}
\newcommand{\BigOAt}[2]{\mathcal{O}_{#2}\!\left(#1\right)}
\newcommand{\LittleOAt}[2]{o_{#2}\!\left(#1\right)}

\usepackage{booktabs,longtable,tabularx,array,multirow}
\usepackage{graphicx}
\usepackage[dvipsnames,svgnames]{xcolor}
\usepackage{enumitem}
\usepackage{fancyhdr}
\usepackage{float}
\usepackage{needspace}
\usepackage{xurl}
\usepackage{listings}
\usepackage[most]{tcolorbox}
\usepackage[hidelinks]{hyperref}
\usepackage[nameinlink,noabbrev,capitalise]{cleveref}

\setlength{\parindent}{1.1em}
\setlength{\parskip}{0.15em}
\setlist{itemsep=0.25em,topsep=0.4em}
\allowdisplaybreaks
\renewcommand{\headrulewidth}{0.4pt}
\pagestyle{fancy}
\fancyhf{}
\fancyhead[L]{\small\nouppercase{\leftmark}}
\fancyhead[R]{\small\thepage}

\theoremstyle{plain}
\newtheorem{theorem}{Theorem}[chapter]
\newtheorem{proposition}[theorem]{Proposition}
\newtheorem{lemma}[theorem]{Lemma}
\newtheorem{corollary}[theorem]{Corollary}
\newtheorem{conjecture}[theorem]{Conjecture}
\newtheorem{identity}[theorem]{Identity}
\theoremstyle{definition}
\newtheorem{definition}[theorem]{Definition}
\newtheorem{example}[theorem]{Example}
\newtheorem{algorithm}[theorem]{Algorithm}
\newtheorem{exercise}[theorem]{Exercise}
\newtheorem{problem}[theorem]{Problem}
\newtheorem{convention}[theorem]{Convention}
\theoremstyle{remark}
\newtheorem{remark}[theorem]{Remark}
\newtheorem{warning}[theorem]{Warning}
\newtheorem{historical}[theorem]{Historical note}

%% the r-Lambert function at parameter 1, used in the x + W(x) part
\newcommand{\rWone}{\mathrm{W}_{\!1}}

%% the book class has no abstract environment
\newenvironment{abstract}%
  {\cleardoublepage\thispagestyle{empty}\begin{center}%
   \bfseries\large Abstract\end{center}\begin{quotation}\noindent\ignorespaces}%
  {\end{quotation}\clearpage}


\definecolor{linkblue}{RGB}{24,72,124}
\definecolor{deepgreen}{RGB}{23,102,69}
\definecolor{deepred}{RGB}{140,42,42}
\definecolor{deepgold}{RGB}{122,91,19}
\definecolor{shadegray}{RGB}{246,247,249}
\definecolor{rulegray}{RGB}{195,201,208}
\definecolor{palered}{RGB}{251,236,236}
\definecolor{palegold}{RGB}{251,245,230}
\definecolor{palegreen}{RGB}{233,245,243}

\crefname{exercise}{Exercise}{Exercises}
\Crefname{exercise}{Exercise}{Exercises}
\crefname{algorithm}{Algorithm}{Algorithms}
\Crefname{algorithm}{Algorithm}{Algorithms}
\crefname{chapter}{Part}{Parts}

\newtcolorbox{warningbox}[1][]{%
  enhanced,breakable,colback=palered,colframe=deepred,
  boxrule=0.6pt,arc=1.2mm,left=1.8mm,right=1.8mm,top=1.2mm,bottom=1.2mm,
  title={A common trap},fonttitle=\bfseries,#1}
\newtcolorbox{summarybox}[1][]{%
  enhanced,breakable,colback=palegold,colframe=deepgold,
  boxrule=0.6pt,arc=1.2mm,left=1.8mm,right=1.8mm,top=1.2mm,bottom=1.2mm,
  title={Section summary},fonttitle=\bfseries,#1}
\newtcolorbox{checkpointbox}[1][]{%
  enhanced,breakable,colback=palegreen,colframe=deepgreen,
  boxrule=0.6pt,arc=1.2mm,left=1.8mm,right=1.8mm,top=1.2mm,bottom=1.2mm,
  title={Structural checkpoint},fonttitle=\bfseries,#1}
\newtcolorbox{sourcebox}{colback=black!3,colframe=black!35,boxrule=0.4pt,
  left=4pt,right=4pt,top=3pt,bottom=3pt,breakable}
\newtcolorbox{repairbox}{colback=Goldenrod!12,colframe=Goldenrod!60!black,
  boxrule=0.4pt,left=4pt,right=4pt,top=3pt,bottom=3pt,breakable}

\lstdefinestyle{pseudo}{%
  basicstyle=\ttfamily\small,columns=fullflexible,frame=single,
  rulecolor=\color{rulegray},backgroundcolor=\color{shadegray},
  showstringspaces=false,keepspaces=true,
  xleftmargin=0.4em,xrightmargin=0.4em,aboveskip=0.8em,belowskip=0.8em}
\lstset{style=pseudo}

\newcommand{\K}{\mathbb{K}}
\newcommand{\PLpow}{\PolynomialLogPowerSeriesRing{\K}}   % K[L][[t]]
\newcommand{\PLlau}{\PolynomialLogLaurentRing{\K}}       % K[L]((t))
\newcommand{\PLrat}{\RationalLogLaurentField{\K}}        % K(L)((t))
\newcommand{\PLit}{\IteratedLaurentLogField{\K}}         % K((L^{-1}))((t))
\newcommand{\Gnear}{\NearIdentityPolynomialLogGroup{\K}} % X + K[L][[t]]
\newcommand{\Gfilt}[1]{\mathcal{G}^{(#1)}_{\mathrm{poly},\K}}
\newcommand{\ordt}{\nu_{t}}                              % outer valuation
\DeclareMathOperator{\lm}{lm}
\DeclareMathOperator{\lc}{lc}
\DeclareMathOperator{\degL}{deg_{L}}
\newcommand{\Trunc}[2]{\operatorname{Trunc}_{<#2}#1}     % exclusive cutoff
\newcommand{\DX}{D_{X}}                                  % distinguished derivation
\newcommand{\EulerOp}{\vartheta}                         % Euler operator X D_X
\newcommand{\shiftop}[2]{\Delta_{#1,#2}}                 % repeated-derivative operator
\newcommand{\compose}{\mathbin{\circ}}
\newcommand{\rev}[1]{\CompositionalInverseOf{#1}}
\newcommand{\recip}[1]{\MultiplicativeInverseOf{#1}}
\newcommand{\stirlingone}[2]{\UnsignedStirlingFirstKind{#1}{#2}}
\newcommand{\stirlingsigned}[2]{\SignedStirlingFirstKind{#1}{#2}}
\newcommand{\LamP}[1]{\LambertPolynomial{#1}}
\newcommand{\LamPsc}[1]{\ScaledLambertPolynomial{#1}}
\newcommand{\RevErr}[2]{\operatorname{Err}_{\mathrm{rev}}\!\left(#1,#2\right)}
\newcommand{\dprofile}[1]{d_{#1}}                        % log-degree profile
\newcommand{\precdom}{\prec}
\newcommand{\succdom}{\succ}
\newcommand{\preceqdom}{\preccurlyeq}
\newcommand{\asymdom}{\asymp}
\newcommand{\Nat}{\mathbb{N}}
\newcommand{\Int}{\mathbb{Z}}
\newcommand{\Rat}{\mathbb{Q}}
\newcommand{\Real}{\mathbb{R}}
\newcommand{\Cplx}{\mathbb{C}}
\newcommand{\e}{\mathrm{e}}
\newcommand{\ii}{\mathrm{i}}
\newcommand{\dd}{\mathrm{d}}
\newcommand{\Coef}[1]{[#1]\,}
\newcommand{\bigO}{\mathcal{O}}
\newcommand{\littleo}{o}
\DeclareMathOperator{\Res}{Res}
\DeclareMathOperator{\sgn}{sgn}
\DeclareMathOperator{\Li}{Li}
\DeclareMathOperator*{\argmin}{arg\,min}
\newcommand{\OB}{\widehat{B}}
\newcommand{\EB}{B}
\newcommand{\Wlam}{W}
\newcommand{\Xcore}{X}
\newcommand{\stirlingfirst}[2]{\genfrac{[}{]}{0pt}{}{#1}{#2}}
\newcommand{\stirlingsecond}[2]{\genfrac{\{}{\}}{0pt}{}{#1}{#2}}

\begin{document}
\frontmatter

\title{\textbf{Transseries}\\[0.35em]
  \Large The polynomial--logarithmic calculus,\\
  series reversal at infinity,\\
  and the inversion of rapidly growing functions}
\author{Vladimir Reshetnikov\\
  \small ProveIt formal-development synthesis}
\date{4 September 2026}

\maketitle

\begin{abstract}
This volume is the single consolidated treatment of the
\texttt{series-and-transseries} corpus.  It develops one formal scale --- the
polynomial--logarithmic transseries, whose monomials are $X^{-\mu}(\log X)^{k}$
--- from its ring structure through arithmetic, differentiation, composition
and reversal at infinity; it then specializes that calculus to the Lambert
$\LambertW$ function and the Wright $\omega$ function, converts formal
identities into Poincar\'e asymptotics with certificates, and finally applies
the whole apparatus to invert a dozen rapidly growing functions and sequences
to all orders.

The organizing observation is that inversion is not a collection of tricks.
For every subject treated here the same three steps suffice: isolate a
\emph{dominant block} that can be inverted in closed form; check that its
normalized slope is a unit; and revert the remainder by one triangular
recurrence whose only division is by that slope.  What changes between
subjects is which block, and what the slope is.  For the polynomial--logarithmic
scale the block is $X+\text{(lower)}$ and the slope is $1$; for a phase
$aX+b\log X$ the block is inverted by $\LambertW$ and the slope is $a+b/X$; for
$a x^{p}(\log x-\beta)$, where the unknown multiplies the logarithm, the block
is again inverted by $\LambertW$ but by a different substitution and the slope
is $v+1$; for a pure exponential the block is a logarithm and the slope is
constant.

The first eight parts develop the calculus.  The remaining parts apply it: to the reversal of $x+\LambertW(x)$; to four
combinatorial sequences --- the rooted trees of \textsc{a000081}, the double
factorial, the partition numbers, the swing factorial; to four special
functions --- Euler's $\Gamma$, the Barnes $G$-function, the hyperfactorial
$K$-function, the subfactorial; and to a real-argument Fibonacci function
whose inverse, alone among them, converges.  A synthesis part compares all of
them and records what the common frame does not explain.

Every part of this volume is a consolidation of independently written
articles; the provenance appendix lists each source and the part that absorbed
it, and the repair appendix lists every correction made to a source.  Where
several sources proved the same thing, one presentation is kept and the others
cite it; where a source asserted a named result without proof, a proof is
supplied or the statement is demoted to a remark that says what is actually
established.
\end{abstract}

\tableofcontents

\mainmatter


\part{Orientation}

\chapter{What a transseries is, and why one is needed}
\label{q0:sec:why}

\begin{sourcebox}
This part replaces four independently written expository introductions:
\texttt{Transseries\_Tutorial-1}, \texttt{-2},
\texttt{Transseries\_Tutorial-3} (\emph{Transseries for Mere Mortals}) and
\texttt{-4}.  Between them they state $58$ named results, of which $21$ are
already theorems of the later parts of this volume and are cited here rather
than restated.  Of the remainder, most are definitions --- the vocabulary is
genuinely common to all four --- and those are given once.  Five of the named
theorems are labelled in their sources ``stated without proof'' or
``schematic''; \cref{q0:sec:quoted} says which, and why they stay quoted.
Everything else stated here as a theorem is proved here.
\end{sourcebox}

\section{An expansion needs a scale}
\label{q0:sec:scale}

A limit remembers too little.  Knowing that $f(x)\to0$ says nothing about
whether $f$ is $1/x$, $\e^{-x}$ or $1/\log x$, and those three behave so
differently that no single statement covers them.  An asymptotic expansion
records more: it compares $f$ against a fixed yardstick.

\begin{definition}[Asymptotic scale]\label{q0:def:scale}
Let $\mathcal{S}$ be a filter base on a set --- for us always $x\to+\infty$
along a real ray, or $x\to\infty$ in a sector.  A sequence
$(\varphi_n)_{n\ge0}$ of eventually nonvanishing functions is an
\emph{asymptotic scale} if
\begin{equation}\label{q0:eq:scale}
 \varphi_{n+1}=\littleo(\varphi_n)\qquad\text{for every }n\ge0 .
\end{equation}
\end{definition}

\begin{definition}[Poincar\'e asymptotic expansion]\label{q0:def:poincare}
Given a scale $(\varphi_n)$ and constants $a_n$, write
\begin{equation}\label{q0:eq:poincare}
 f\sim\sum_{n\ge0}a_n\varphi_n
 \quad\text{to mean}\quad
 f-\sum_{n<N}a_n\varphi_n=\bigO(\varphi_N)\ \ \text{for every fixed }N .
\end{equation}
The number of retained terms is fixed \emph{before} the limit is taken.  No
claim is made that the infinite series converges, and in the cases that matter
in this volume it does not.
\end{definition}

\begin{proposition}[Uniqueness relative to a scale]\label{q0:prop:uniqueness}
The coefficients in \eqref{q0:eq:poincare} are uniquely determined by $f$ and
the scale:
\begin{equation}\label{q0:eq:coefficients}
 a_N=\lim\frac{f-\sum_{n<N}a_n\varphi_n}{\varphi_N},
\end{equation}
the limit being along $\mathcal S$.  In particular two different sequences of
coefficients cannot represent the same $f$ on the same scale.
\end{proposition}

\begin{proof}
Induction on $N$.  For $N=0$, \eqref{q0:eq:poincare} at $N=1$ gives
$f-a_0\varphi_0=\bigO(\varphi_1)=\littleo(\varphi_0)$, so
$f/\varphi_0\to a_0$; this both proves the formula and shows $a_0$ is
determined.  Suppose $a_0,\dots,a_{N-1}$ are determined.  Applying
\eqref{q0:eq:poincare} at $N+1$ to $g:=f-\sum_{n<N}a_n\varphi_n$ gives
$g-a_N\varphi_N=\bigO(\varphi_{N+1})=\littleo(\varphi_N)$, whence
$g/\varphi_N\to a_N$, which is \eqref{q0:eq:coefficients}.
\end{proof}

\begin{remark}[Uniqueness is relative, and that is the whole point]
\label{q0:rem:relative}
\Cref{q0:prop:uniqueness} says the coefficients are unique \emph{given the
scale}.  It does not say the expansion determines $f$.  It cannot: the two
functions $f$ and $f+\e^{-x}$ have identical expansions on the scale
$(x^{-n})_{n\ge0}$, because $\e^{-x}=\bigO(x^{-N})$ for every $N$.  This is
the single fact that makes transseries necessary, and \cref{q0:sec:flat}
develops it.
\end{remark}

\section{Beyond all algebraic orders}
\label{q0:sec:flat}

\begin{definition}[Flatness]\label{q0:def:flat}
A function $\varepsilon$ is \emph{flat} relative to a scale $(\varphi_n)$ if
$\varepsilon=\bigO(\varphi_n)$ for every $n$.  Relative to the power scale
$(x^{-n})$ this is often written $\varepsilon=\bigO(x^{-\infty})$.
\end{definition}

\begin{proposition}[The invisible-function problem]\label{q0:prop:invisible}
$\e^{-x}$ is flat relative to $(x^{-n})_{n\ge0}$, and the set of functions
flat relative to a given scale is a vector space closed under multiplication
by anything of at most polynomial growth.  Consequently a Poincar\'e
expansion on the power scale determines $f$ only up to a flat function, and no
refinement of the coefficients can recover the difference.
\end{proposition}

\begin{proof}
For any $N$, $x^{N}\e^{-x}\to0$ as $x\to+\infty$, since $\e^{x}$ exceeds
$x^{N+1}/(N+1)!$; hence $\e^{-x}=\littleo(x^{-N})$ for every $N$.  Closure
under sums and under multiplication by $\bigO(x^{M})$ is immediate from the
definition.  The last sentence is \cref{q0:prop:uniqueness} applied to $f$ and
$f+\varepsilon$: both have the same coefficients, so the expansion cannot
separate them.
\end{proof}

\begin{remark}[This volume is a catalogue of the phenomenon]
\label{q0:rem:catalogue}
Almost every chapter of this volume is, at bottom, an instance of
\cref{q0:prop:invisible}.  In \cref{q1:sec:top} the entire block structure of
the inverse of $x+\LambertW(x)$ is flat relative to the inverse-logarithmic
scale, so an expansion in powers of $1/\log z$ --- correct to every order ---
sees none of it.  In \cref{p6:sec:top} the constant $\tfrac12$ in
$\Gamma_+^{-1}$ is flat relative to the outer expansion of the Lambert core.
In \cref{q2:sec:bell} the third family of sectors is flat relative to the
first two.  In \cref{p8:sec:top} the whole subfactorial correction is flat
relative to every fixed exponential.  The remedy is always the same: enlarge
the scale until the invisible term becomes a monomial of it.
\end{remark}

\section{Divergence is not failure}
\label{q0:sec:divergence}

\begin{remark}[The Euler example]\label{q0:rem:euler}
The classical illustration is the Euler equation $x^{2}y'+y=x$, whose formal
solution is $\sum_{n\ge0}(-1)^{n}n!\,x^{n+1}$.  The series diverges for every
$x\ne0$, and yet its truncations approximate a genuine solution extremely
well: the error after $N$ terms is about $N!\,|x|^{N+1}$, which for fixed $x$
first decreases and then grows.  Stopping at the smallest term --- optimal
truncation --- leaves an error of order $\e^{-1/|x|}$, exponentially small.
The rigorous version of this, for the Gevrey-$1$ remainder bounds this volume
actually needs, is \cref{p0:thm:optimal-truncation}; the point to carry
forward is that a divergent expansion can be more informative than a
convergent one, and that the size of its least term is itself a meaningful
quantity.
\end{remark}

\begin{remark}[Where the exponentially small term goes]\label{q0:rem:beyond}
The $\e^{-1/|x|}$ left over by optimal truncation is exactly a flat function
in the sense of \cref{q0:def:flat}.  It is not noise: it is a term of the
answer, on a scale the power series cannot express.  A transseries is the
bookkeeping that lets one write it down --- and then, if one wishes, expand
\emph{its} coefficient in powers again, producing a second series, and so on.
\end{remark}

\chapter{The algebra of monomials}
\label{q0:sec:algebra}

The vocabulary below is common to all four source tutorials and to the
calculus parts of this volume; it is fixed once here.

\section{Ordered monomials and well-based support}
\label{q0:sec:monomials}

\begin{definition}[Ordered monomial group]\label{q0:def:group}
A \emph{monomial group} is a totally ordered abelian group
$(\mathfrak M,\cdot,\prec)$: multiplication of monomials is the group law, and
$\mathfrak m\prec\mathfrak n$ means $\mathfrak m$ is asymptotically negligible
against $\mathfrak n$.  Compatibility means $\mathfrak m\prec\mathfrak n$
implies $\mathfrak{mp}\prec\mathfrak{np}$ for every $\mathfrak p$.
\end{definition}

\begin{remark}[Why the ordering cannot be Archimedean]\label{q0:rem:archimedean}
In $\mathfrak M$ the monomials $x^{-1}$ and $\e^{-x}$ satisfy
$\e^{-x}\prec x^{-n}$ for every $n$; so no multiple of $\e^{-x}$ ever reaches
$x^{-1}$, and the ordered group is not Archimedean.  This is the algebraic
shadow of \cref{q0:prop:invisible}, and it is why the value group of a
transseries field is large and the field is not a Hardy field of finite rank.
\end{remark}

\begin{definition}[Well-based support and Hahn series]\label{q0:def:hahn}
A subset $S\subseteq\mathfrak M$ is \emph{well-based} (equivalently
\emph{reverse well-ordered}) if every nonempty subset of $S$ has a
$\prec$-greatest element; equivalently, $S$ contains no strictly increasing
infinite sequence.  Given a coefficient field $\K$, the \emph{Hahn
series} field $\K[[\mathfrak M]]$ consists of the formal sums
$\sum_{\mathfrak m\in S}c_{\mathfrak m}\mathfrak m$ with $S$ well-based.  The
\emph{leading monomial} of a nonzero such series is the greatest element of
its support, and its coefficient the \emph{leading coefficient}.
\end{definition}

The next two lemmas are what make that definition usable: they are the reason
products and infinite sums of Hahn series exist.  Both are stated in the
sources; one of them is stated there without proof.

\begin{lemma}[Neumann's lemma]\label{q0:lem:neumann}
Let $A,B\subseteq\mathfrak M$ be well-based.  Then $A\cdot B$ is well-based,
and every $\mathfrak c\in A\cdot B$ has only finitely many factorizations
$\mathfrak c=\mathfrak{ab}$ with $\mathfrak a\in A$, $\mathfrak b\in B$.
\end{lemma}

\begin{proof}
Write the group additively for the proof, so $A+B$ and $\mathfrak c=a+b$.

\emph{Finiteness of factorizations.}  Suppose $c=a_i+b_i$ for infinitely many
distinct pairs $(a_i,b_i)$, $i\in\Nat$.  Every sequence in a totally ordered
set has a monotone subsequence, so we may assume $(a_i)$ is monotone.  It
cannot be strictly increasing, since $A$ is well-based; so $(a_i)$ is
non-increasing.  Then $b_i=c-a_i$ is non-decreasing, and because the pairs are
distinct and determined by $a_i$, infinitely many $b_i$ are distinct, giving a
strictly increasing sequence in $B$ --- contradicting that $B$ is well-based.

\emph{$A+B$ is well-based.}  Suppose $c_1\prec c_2\prec\cdots$ in $A+B$, with
$c_i=a_i+b_i$.  Passing to a subsequence, assume $(a_i)$ monotone; as above it
is non-increasing.  Then $b_i=c_i-a_i$ is strictly increasing, since $c_i$ is
strictly increasing and $-a_i$ is non-decreasing --- again contradicting that
$B$ is well-based.
\end{proof}

\begin{lemma}[Dickson's lemma]\label{q0:lem:dickson}
Every subset of $\Nat^{k}$ has finitely many minimal elements for the
componentwise order; equivalently, $\Nat^{k}$ contains no infinite antichain.
\end{lemma}

\begin{proof}
Let $(v_i)_{i\ge1}$ be any sequence in $\Nat^{k}$.  Its first coordinates form
a sequence in $\Nat$, which has a non-decreasing subsequence; within that
subsequence the second coordinates again have a non-decreasing subsequence;
after $k$ such extractions we obtain a subsequence non-decreasing in every
coordinate, so its first two terms satisfy $v\le v'$ componentwise.  Hence no
infinite antichain exists, and a set whose minimal elements were infinite
would contain one.
\end{proof}

\begin{remark}[What these two are for]\label{q0:rem:why-lemmas}
\Cref{q0:lem:neumann} is what makes the \emph{product} of two Hahn series
well defined: each coefficient of the product is a finite sum.  It is used in
this volume at \cref{p6:thm:triangular}, where it is exactly the hypothesis
that lets a triangular recurrence terminate at each exponent, and it is the
reason that theorem needs only well-basedness and not the stronger local
finiteness one source assumed.  \Cref{q0:lem:dickson} is what makes a
\emph{grid} --- a support contained in a finitely generated monoid --- behave
like a finite object, and it underlies the closure properties that the
calculus parts use without further comment.
\end{remark}

\begin{definition}[Summable family]\label{q0:def:summable}
A family $(f_j)_{j\in J}$ of Hahn series is \emph{summable} if the union of
their supports is well-based and each monomial lies in only finitely many
supports.  Its sum is then defined coefficientwise.
\end{definition}

\section{Height, depth, and how monomials compare}
\label{q0:sec:height}

\begin{definition}[Exponential height and logarithmic depth]
\label{q0:def:height}
The \emph{exponential height} of a transmonomial is the number of nested
exponentials it carries; the \emph{logarithmic depth} is the number of nested
logarithms.  Thus $x=\e^{\log x}$ has height $0$ and depth $0$ in the usual
normalization, $\e^{x}$ has height $1$, and $\log\log x$ has depth $2$.
\end{definition}

\begin{proposition}[Height wins]\label{q0:prop:height}
For all $N,M\ge0$,
\begin{equation}\label{q0:eq:height}
 x^{N}(\log x)^{M}=\littleo\bigl(\e^{\varepsilon x}\bigr)
 \qquad\text{for every }\varepsilon>0,
\end{equation}
and consequently no product of powers and logarithms can compare with a
monomial of greater exponential height.  Dually, $(\log x)^{M}=\littleo(x^{\delta})$
for every $\delta>0$: a greater logarithmic depth is always negligible against
any positive power.
\end{proposition}

\begin{proof}
Taking logarithms, $N\log x+M\log\log x-\varepsilon x\to-\infty$, since
$\log x=\littleo(x)$ and $\log\log x=\littleo(\log x)$; exponentiating gives
\eqref{q0:eq:height}.  The dual statement is the same computation with
$M\log\log x-\delta\log x\to-\infty$.
\end{proof}

\begin{remark}[The comparison algorithm]\label{q0:rem:compare}
\Cref{q0:prop:height} is the top of a decision procedure: to compare two
transmonomials, compare exponential heights first; on a tie, compare the
exponents recursively; on a further tie, compare powers; and only at the
bottom compare logarithmic depths.  That is exactly the order in which the
scale part of this volume organizes the polynomial--logarithmic monomials, and
it is why the degree profile there is a lexicographic invariant rather than a
single number.
\end{remark}

\section{How to read this volume}
\label{q0:sec:roadmap}

The parts that follow fall into three groups.

\emph{The calculus} develops one scale in operational depth: the
polynomial--logarithmic transseries, whose monomials are
$X^{-\mu}(\log X)^{k}$.  It builds the rings, then arithmetic and
differentiation, then composition, then reversal at infinity; then it
specializes to $\LambertW$ and the Wright $\omega$ function; then it converts
formal identities into Poincar\'e asymptotics with certificates; then
algorithms and diagnostics.  A reader who wants the machinery should read
these in order.

\emph{The inversion apparatus} is a single part standing between the calculus
and its applications.  It contains what the calculus does not: the
exponential--power model, the monomial $\alpha$-reduction, the axiomatized
dominant core, perturbed inversion around it, and --- the part with no analogue
in the calculus --- the theory of inverting a \emph{sequence}, where three
different inverse objects must be distinguished and a staircase theorem
relates them.

\emph{The applications} are independent of one another.  Each takes a
rapidly growing function or sequence, identifies its dominant block, inverts
that block in closed form, and reverts the rest.  A reader interested in one
subject can read its part alone, following the citations back into the
calculus as needed.

\begin{remark}[The one thing worth carrying through all of it]
\label{q0:rem:moral}
Every part of this volume repeats one manoeuvre: \emph{normalize until what
remains is an ordinary reversion}.  Whether the normalization is the
substitution $u=\e^{-\lambda x}$ of \cref{p9:thm:factorization}, the Lambert
coordinate of \cref{p6:sec:core}, the change $q=L/(L-v)-\e^{-v}$ of
\cref{q1:prop:B}, or the saddle $r=\Wlam_0(n)$ of \cref{q2:thm:bell}, the shape
is the same: find the block that can be inverted exactly, check that its
normalized slope is a unit, and revert the remainder by a triangular
recurrence whose only division is by that slope.  What differs between
chapters is which block and what the slope is.
\end{remark}

\section{What is quoted rather than proved}
\label{q0:sec:quoted}

Five statements appear in the source tutorials as named theorems that they do
not prove, and this volume does not prove them either.  They are recorded here,
demoted from theorems to attributed quotations, so that nothing in the volume
carries the authority of a proof it does not have.

\begin{enumerate}
\item \emph{Real closedness of the transseries field.}  The field of
      logarithmic-exponential transseries is real closed.  This is a theorem
      of the literature (van der Hoeven; Aschenbrenner--van den
      Dries--van der Hoeven), and neither the sources nor this volume proves
      it.  Nothing in the volume depends on it.
\item \emph{Closure under antidifferentiation.}  Every transseries has a
      transseries antiderivative.  Again a theorem of the literature; the
      calculus part of this volume proves the corresponding statement only for
      the polynomial--logarithmic scale, where it is elementary and where a
      genuine obstruction --- the appearance of a new logarithm --- is visible
      and is handled explicitly.
\item \emph{The formal Taylor principle} and \emph{the formal contraction
      principle}, stated schematically in two of the sources.  This volume
      proves the versions it uses --- the composition part for the first, and
      the reversal part for the second --- on the polynomial--logarithmic
      scale with explicit hypotheses.  The schematic general forms are not
      established here.
\item \emph{Generalized Borel reconstruction.}  That a resurgent germ is
      recovered from its Borel transform by a directional Laplace integral,
      with Stokes jumps across singular directions, is quoted schematically.
      This volume uses only the elementary Watson-type statements it proves,
      and every hyperasymptotic claim in the applications is labelled as
      quoted or as open --- see \cref{p6:rem:not-claimed}.
\end{enumerate}

\begin{remark}[Why this list exists]\label{q0:rem:honesty}
A tutorial may reasonably state a deep theorem without proof, to orient the
reader.  A reference volume may not do so silently, because a later chapter
may then appear to rest on it.  Listing the five here, and checking that no
result of this volume uses any of them, is cheaper than tracking the
dependency later.
\end{remark}

\clearpage
\chapter{Reader's guide and the structure of this volume}
\label{plt:sec:readers-guide}

\section{The eight parts}
\label{plt:sec:mot-parts}

% ed.: the draft called "Part V may be read directly after Part IV" an
% exception to the rule that each part depends only on those before it.  It is
% not an exception at all --- V does come after IV --- and it also contradicts
% reading path 1 below, which does draw two statements from Parts II and III.
% The intended and correct content is the lightness of that dependency.
The volume is arranged so that each part depends only on those before it.  One
of those dependencies is lighter than the ordering suggests: a reader who wants
only the Lambert expansion may pass from Part~I directly to Parts~IV and~V,
because Part~V uses Parts~II and~III only through the two statements quoted in
reading path~1 below.

\begin{center}
\begin{tabular}{@{}lp{0.16\textwidth}p{0.54\textwidth}@{}}
\hline
Part & Subject & What it establishes \\
\hline
I & The scale &
  The generators \(t=X^{-1}\), \(L=\log X\); dominance and the rank-two
  valuation; the rings \(\PLpow\), \(\PLlau\), \(\PLrat\), \(\PLit\);
  coefficient support, leading block, leading monomial and leading scalar;
  exclusive truncation \(\Trunc{F}{N}\); \(t\)-adic completeness and formal
  summability. \\
II & Arithmetic and differential calculus &
  Cauchy convolution block by block; the exact unit criterion and the
  reciprocal and quotient recurrences; the distinguished derivation
  \(\DX=t\partial_{L}-t^{2}\partial_{t}\) and the Euler operator \(\EulerOp\);
  repeated derivatives \(\shiftop{n}{k}\); antidifferentiation and the resonant
  block; powers, logarithms and exponentials of small series. \\
III & Composition &
  Admissible inner arguments; substitution on the generators; the exact formal
  Taylor formula and its formal summability; finite coefficient formulas;
  associativity and the identity; the near-identity group \(\Gnear\); the exact
  points at which composition leaves the class. \\
IV & Series reversal &
  Existence and uniqueness of \(\rev{F}\) on \(\Gnear\); four algorithms ---
  contractive fixed point, triangular coefficient matching, Newton iteration
  with precision doubling, and Lagrange--B\"urmann at infinity; residual
  certificates \(\RevErr{F}{H}\); reversion with a nontrivial leading
  monomial. \\
V & Wright omega, the Lambert polynomials, and Lambert \(\LambertW\) &
  The reversion of \(X+L\); the canonical zero-based family \(\LamP{n}\) and
  the scaled integer family \(\LamPsc{n}\); recurrence, generating equation,
  operator formula, closed form through unsigned Stirling cycle numbers, degree
  law, Laplace transform; all complex branches, the lower real branch
  \(\LambertW_{-1}\), and the square-root branch point. \\
VI & From formal transseries to analytic asymptotics &
  What a formal identity does and does not give; blockwise Poincar\'e
  expansions with the logarithmic degree printed; transport of asymptotics
  through the algebraic operations; from a formal residual to an analytic
  inverse error; convergent versus divergent expansions; branch consistency in
  sectors. \\
VII & Algorithms, certificates, diagnostics &
  Sparse data structures for \(n\mapsto p_{n}(L)\); truncated primitives and
  their cost; reversion strategies compared; verification invariants, residual
  certificates, and a decision procedure for the recurring failure modes. \\
VIII & Extensions and the wider landscape &
  Puiseux and Hahn exponent groups; rational-log and Laurent-log coefficients;
  iterated logarithms and multiseries; power-log and Dulac classes at zero; the
  relation to the logarithmic-exponential transseries field and to Hardy
  fields \cite{hardy-orders,edgar-beginners,vanderhoeven-book,adh-book}. \\
\hline
\end{tabular}
\end{center}

\section{Appendices}
\label{plt:sec:mot-appendices}

\begin{itemize}
\item \emph{Bell polynomials and Fa\`a di Bruno bookkeeping.}  The partial and
  complete exponential Bell polynomials
  \(\ExponentialPartialBellPolynomial{n}{k}\) and
  \(\ExponentialCompleteBellPolynomial{n}\) in exactly the form used by the
  composition and reversion formulas of Parts~III and~IV.
\item \emph{Closed-form coefficient extraction.}  Explicit formulas for
  \([t^{n}]F\) and \([t^{n}L^{j}]F\) after each operation, with the conventions
  of \cref{plt:rmk:mot-normalizations} in force.
\item \emph{Lambert coefficient tables.}  \(\LamP{n}\) and \(\LamPsc{n}\)
  through a stated order, with Stirling diagonals
  \(\stirlingone{n}{k}\) and a reproducible residual check for each row.
\item \emph{Operational formula sheet.}  One entry per operation: hypotheses,
  formula, cost, and the certificate that validates a computed truncation.
\item \emph{Exercises with solutions.}
\item \emph{Notation concordance.}  The map from each source draft's local
  letters to the normative catalogue commands.
\item \emph{Provenance.}  Which source contributed which result, and the
  editorial repairs applied to it.
\end{itemize}

\section{Reading paths}
\label{plt:sec:mot-paths}

\begin{enumerate}
\item \emph{Only the Lambert \(\LambertW\) expansion.}  Read
  Part~I \(\to\) Part~IV \(\to\) Part~V.  Part~I supplies the scale, the ring
  and the valuation; Part~IV supplies the one theorem this path needs, namely
  that every element of \(\Gnear\) has a compositional inverse in \(\Gnear\)
  computable to any order by a finite calculation; Part~V performs the
  reversion of \(X+L\) and reads off every branch.  The only imports from
  Parts~II and~III along this path are the block derivative
  \eqref{plt:eq:mot-block-derivative} and the formal Taylor formula, both
  quoted with full statements where used.
\item \emph{The full calculus.}  Read Parts~I--IV in order.  This is the
  self-contained algebraic core.
\item \emph{Analytic use.}  After Part~I, read Part~VI.  A reader who intends
  to convert a formal truncation into a numerical bound should read Part~VI
  before computing anything, because none of the formal machinery by itself
  licenses an analytic claim.
\item \emph{Implementation.}  Read Part~I, then the recurrences of
  Parts~II--IV, then Part~VII with the formula-sheet appendix.
\item \emph{Context.}  Read Part~VIII for the embedding of this class into
  Hahn series, Dulac classes, Hardy fields, and the full
  logarithmic-exponential transseries field.
\end{enumerate}

\begin{checkpointbox}
After Part~I a reader should be able to: write an element of \(\PLlau\) in the
normal form \eqref{plt:eq:mot-normal-form}; compute \(\ordt\), \(\lm\),
\(\lc\), \(\degL\) and the log-degree profile \(\dprofile{F}\) of a given
series; decide from the leading block alone whether the series is a unit; say
which of \(\PLpow\), \(\PLlau\), \(\PLrat\), \(\PLit\) a given expression
belongs to; and state, for a displayed truncation, whether its tail is a formal
valuation tail or an analytic remainder.
\end{checkpointbox}

\section{Standing conventions and normalizations}
\label{plt:sec:mot-conventions}

\begin{remark}[Normalizations adopted, and how they differ from the source
drafts]
\label{plt:rmk:mot-normalizations}
This volume merges six independently written treatments.  Where their
conventions collided, the notation catalogue of the corpus is the arbiter, and
the following choices are in force everywhere.
\begin{enumerate}
\item \emph{Regime and field.}  Unless a statement says otherwise,
  \(X\to+\infty\) along the positive real ray, \(\K\) is a field of
  characteristic zero, and \(t=X^{-1}\), \(L=\log X\).  Complex work names its
  sector and its branch of \(\log\).  Local letters used by individual sources
  --- \(\lambda\), \(\ell\), \(x\) --- are renamings of \(L\) or \(X\) and are
  not additional generators.
\item \emph{Truncation is exclusive.}
  \(\Trunc{F}{N}=\sum_{n<N}p_{n}(L)t^{n}\) keeps every logarithmic monomial in
  every retained block; ``through outer order \(t^{N-1}\)'' means the exclusive
  cutoff \(<N\).  One source uses the inclusive \([F]_{\le N}\); the conversion
  is \([F]_{\le N}=\Trunc{F}{N+1}\), a one-block difference that must be
  applied whenever a formula is transported from that source.
\item \emph{Coefficient extraction.}  \([t^{n}]F=p_{n}(L)\) and
  \([t^{n}L^{j}]F=a_{n,j}\), with the series written to the right of the
  bracket; an absent coefficient is \(0\).  The same bracket glyph also denotes
  unsigned Stirling cycle numbers \(\stirlingone{n}{k}\) with two integer
  indices; the roles are separated by argument type and never abbreviated to
  the bare glyph in prose.
\item \emph{Lambert polynomials.}  The canonical zero-based family
  \(\LamP{n}\) of \eqref{plt:eq:mot-lambert-first} is used throughout, with
  \(\LamP{0}(u)=-u\).  The degree law \(\deg\LamP{n}=n\) is asserted only for
  \(n\ge1\).  The scaled family is
  \(\LamPsc{n}=n!\,\LamP{n}\in\IntegerNumbers[u]\) for \(n\ge1\); it is a
  storage normalization, not a different coefficient sequence.  The bare symbol
  \(P_{n}\) is not reserved for this family and is not used for it.
\item \emph{The three \(O\)-roles} are kept apart exactly as in the warning box
  of \cref{plt:sec:mot-three-statuses}.
\item \emph{Two inverses.}  \(\recip{F}\) is the multiplicative reciprocal and
  \(\rev{F}\) the compositional inverse.  They are unrelated operations, and the
  ambiguous notation \(F^{-1}\) is used for neither.
\item \emph{Filtration, not ideal.}  In \(\PLlau\) the generator \(t\) is a
  unit, so the principal ideal generated by \(t^{N}\) is all of \(\PLlau\), and
  \(t^{N}\PLpow\) is an additive valuation-filtration tail rather than an
  ideal.  Agreement below outer order \(N\) is therefore phrased as
  \(\ordt(F-G)\ge N\), as \(\FormalBigOAt{t^{N}}{t}\), or as equality of
  exclusive truncations --- never as a congruence modulo an ideal.  Genuine
  quotient-ring language is legitimate only in \(\PLpow\), where \((t^{N})\) is
  a proper ideal.
  % ed.: source draft 2 tabulates the output class of differentiation as
  % "t K[L]((t)) up to the leading power".  That set equals K[L]((t)), because
  % t is a unit there, so the row asserts nothing.  The correct statement is the
  % valuation inequality nu_t(D_X F) >= nu_t(F) + 1, proved in Part II, which
  % can be strict only when a constant leading block is annihilated.
\item \emph{Coefficient blocks are polynomials.}  In \(\PLpow\) and
  \(\PLlau\) every \(p_{n}\) lies in \(\K[L]\), so logarithmic exponents are
  nonnegative and each block is finite.  Membership implies no uniform bound on
  \(\deg_{L}p_{n}\); any such bound is a hypothesis or a proved invariant, never
  part of the definition.
  % ed.: source draft 4 prints the coefficient condition of its defining display
  % as "p_n in K[X]", naming the large variable where the logarithmic one is
  % meant.  The intended and correct condition is p_n in K[L].
\end{enumerate}
\end{remark}

\clearpage
\part{The polynomial--logarithmic scale}

% ---- fragment 01-motivation ----
% ---------------------------------------------------------------------------
% Cluster: mot -- Part I opening.  Motivation and reader's guide.
% ---------------------------------------------------------------------------

\chapter{Why powers and logarithms belong in one expansion}
\label{plt:sec:motivation}

This volume is about a single asymptotic scale and the exact algebra it
supports.  The scale is generated by two symbols,
\begin{equation}\label{plt:eq:mot-generators}
  t\DefinitionEquals X^{-1},
  \qquad
  L\DefinitionEquals\log X,
  \qquad X\to+\infty,
\end{equation}
and its monomials are the products \(t^{n}L^{j}\).  A reader who has met only
power series may reasonably ask why two generators are needed at all and, if
two, why exactly these two.  This chapter answers both questions.  It exhibits
the smallest problem that forces the scale --- inverting \(y\mapsto y+\log y\)
--- and then proves three separate minimality statements: a scale of pure
powers stalls after one term, a scale of pure logarithms cannot even begin, and
the ring generated by \(t\) and \(L\) is the smallest subring of germs at
\(+\infty\) that contains a logarithm and is closed under differentiation.

Throughout this chapter \(X\) is a real variable tending to \(+\infty\), all
logarithms are the real logarithm on \((0,\infty)\), and \(f\sim g\) means
\(f/g\to1\).  Complex branches are introduced only in Part~V, where every
branch of \(\LambertW\) and every choice of \(\log\) is printed with the
statement it belongs to.

\section{The inversion problem that starts everything}
\label{plt:sec:mot-inversion}

The whole subject is visible in one implicit equation with one logarithmic
term.

\begin{proposition}[The basic inverse and its elementary envelope]
\label{plt:prop:mot-omega-basic}
The map \(\Phi(y)=y+\log y\) is a strictly increasing bijection from
\((0,\infty)\) onto \(\RealNumbers\).  Denote its inverse by \(\WrightOmega\),
so that
\begin{equation}\label{plt:eq:mot-omega-equation}
  \WrightOmega(X)+\log\WrightOmega(X)=X
  \qquad(X\in\RealNumbers).
\end{equation}
Then \(\WrightOmega\) is real analytic and strictly increasing,
\(\WrightOmega(X)\to+\infty\) as \(X\to+\infty\), and for every \(X\ge1\),
writing \(L=\log X\),
\begin{equation}\label{plt:eq:mot-envelope}
  X-L\;\le\;\WrightOmega(X)\;\le\;X .
\end{equation}
Moreover \(\LambertW_{0}(z)=\WrightOmega(\log z)\) for every \(z>0\), where
\(\LambertW_{0}\) is the principal real branch of Lambert's function, defined
by \(\LambertW_{0}(z)\EulerE^{\LambertW_{0}(z)}=z\).
\end{proposition}

\begin{proof}
On \((0,\infty)\) we have \(\Phi'(y)=1+1/y>0\), so \(\Phi\) is strictly
increasing and real analytic with nowhere vanishing derivative;
\(\Phi(y)\to-\infty\) as \(y\downarrow0\) and \(\Phi(y)\to+\infty\) as
\(y\to+\infty\), so \(\Phi\) is a bijection onto \(\RealNumbers\).  The inverse
function theorem gives that \(\WrightOmega\) is real analytic and strictly
increasing, and \(\WrightOmega(X)\to+\infty\) as \(X\to+\infty\).

Fix \(X\ge1\) and put \(y=\WrightOmega(X)\).  If \(y<1\) then \(\log y<0\) and
hence \(X=y+\log y<1\), a contradiction; so \(y\ge1\) and \(\log y\ge0\).  From
\(y=X-\log y\) we get \(y\le X\); then \(\log y\le\log X=L\), whence
\(y=X-\log y\ge X-L\).  This is \eqref{plt:eq:mot-envelope}.

For the last claim let \(z>0\) and put \(w=\WrightOmega(\log z)\).  Then
\(w>0\) and \(w+\log w=\log z\), so \(\log(w\EulerE^{w})=\log z\) and
therefore \(w\EulerE^{w}=z\).  Since \(v\mapsto v\EulerE^{v}\) is strictly
increasing on \((-1,\infty)\) and \(w>0>-1\), the value \(w\) is the principal
branch \(\LambertW_{0}(z)\).
\end{proof}

The envelope \eqref{plt:eq:mot-envelope} already contains the message of this
volume in miniature.  It says that \(\WrightOmega(X)\) equals \(X\) up to a
correction of size at most \(\log X\); so the answer to a question posed purely
in powers of \(X\) is not expressible purely in powers of \(X\).  The next
subsection makes the correction explicit to two further orders.

\begin{remark}[Names]
\label{plt:rmk:mot-names}
The function \(\WrightOmega\) is Wright's omega function
\cite{wright-omega,corless-et-al-1996}.  Working with \(\WrightOmega\) rather
than with \(\LambertW\) removes one layer of logarithms from every display: the
two nested logarithms \(\log z\) and \(\log\log z\) of the classical Lambert
expansion become the single large variable \(X\) and the single logarithmic
coefficient variable \(L\).  Every result about \(\LambertW\) in Part~V is
obtained from a result about \(\WrightOmega\) by substituting a branch
coordinate for \(X\).
\end{remark}

\section{Two orders by hand}
\label{plt:sec:mot-two-orders}

The following computation is the motivating one.  It is short enough to do on
paper, and it already displays every structural feature of the scale: the
corrections are inverse powers of \(X\); their coefficients are polynomials in
\(\log X\); and the logarithmic degree grows with the inverse power.

\begin{proposition}[Two orders of the basic inverse]
\label{plt:prop:mot-two-orders}
With \(L=\log X\) and \(\WrightOmega\) as in
\cref{plt:prop:mot-omega-basic}, as \(X\to+\infty\) along the real axis,
\begin{equation}\label{plt:eq:mot-two-orders}
  \WrightOmega(X)
  =X-L+\frac{L}{X}+\frac{1}{X^{2}}\left(\frac{L^{2}}{2}-L\right)
   +\BigOAt{\frac{L^{3}}{X^{3}}}{X\to+\infty}.
\end{equation}
Moreover there is \(X_{0}\ge4\) such that for \(X\ge X_{0}\)
\begin{equation}\label{plt:eq:mot-psi2-envelope}
  \frac{1}{2}\cdot\frac{L}{X}
  \;\le\;\WrightOmega(X)-X+L\;\le\;
  2\cdot\frac{L}{X}.
\end{equation}
In particular \(\WrightOmega(X)-X\sim-L\) and
\(\WrightOmega(X)-X+L\sim L/X\).
\end{proposition}

\begin{proof}
Choose \(X_{0}\ge4\) so large that \(L\ge1\) and
\(\varepsilon\DefinitionEquals L/X\le\tfrac12\) for \(X\ge X_{0}\), and assume
\(X\ge X_{0}\) throughout.  Write
\[
  r\DefinitionEquals\WrightOmega(X)-X+L,
  \qquad
  u\DefinitionEquals\frac{L-r}{X},
\]
so that \(\WrightOmega(X)=X(1-u)\).  By \eqref{plt:eq:mot-envelope} we have
\(0\le r\le L\), hence
\begin{equation}\label{plt:eq:mot-u-range}
  0\le u\le\varepsilon\le\tfrac12 .
\end{equation}
Since \(\WrightOmega(X)>0\), we may take logarithms in
\(\WrightOmega(X)=X(1-u)\):
\(\log\WrightOmega(X)=L+\log(1-u)\).  Substituting this into
\eqref{plt:eq:mot-omega-equation} and solving for \(\WrightOmega(X)\) gives
\(\WrightOmega(X)=X-L-\log(1-u)\), that is,
\begin{equation}\label{plt:eq:mot-r-fixed-point}
  r=-\log(1-u)=\sum_{k\ge1}\frac{u^{k}}{k},
\end{equation}
the series converging because \(0\le u\le\tfrac12\).

\emph{First estimate.}  For \(0\le u\le\tfrac12\) one has
\(u\le-\log(1-u)\le2u\): the quotient \(-\log(1-u)/u=\sum_{k\ge1}u^{k-1}/k\) is
nondecreasing on \([0,\tfrac12]\), equals \(1\) at \(u=0\) and equals
\(2\log2<2\) at \(u=\tfrac12\).  Hence by
\eqref{plt:eq:mot-r-fixed-point} and \eqref{plt:eq:mot-u-range},
\(u\le r\le2u\le2\varepsilon\).  Feeding \(r\le2\varepsilon\) back into the
definition of \(u\) gives
\(u=(L-r)/X\ge(L-2L/X)/X\ge\tfrac12\varepsilon\), because \(X\ge4\).  Therefore
\begin{equation}\label{plt:eq:mot-r-two-sided}
  \tfrac12\varepsilon\le u\le\varepsilon,
  \qquad
  \tfrac12\varepsilon\le r\le2\varepsilon,
\end{equation}
which is \eqref{plt:eq:mot-psi2-envelope}.

\emph{Second estimate.}  From \eqref{plt:eq:mot-r-fixed-point} and
\(0\le u\le\tfrac12\) we get the two tail bounds
\[
  \Bigl|r-u-\frac{u^{2}}{2}\Bigr|
  =\Bigl|\sum_{k\ge3}\frac{u^{k}}{k}\Bigr|
  \le\frac{u^{3}}{3}\sum_{k\ge0}u^{k}
  \le\frac{2}{3}u^{3}\le\frac{2}{3}\varepsilon^{3},
\]
\[
  \AbsoluteValue{r-u}
  \le\frac{u^{2}}{2}\cdot\frac{1}{1-u}\le u^{2}\le\varepsilon^{2}.
\]
Next, \(u=\varepsilon-r/X\) with \(0\le r\le2\varepsilon\), so
\(\AbsoluteValue{u-\varepsilon}\le2\varepsilon/X\le2\varepsilon^{2}\), where the
last step uses \(1/X\le L/X=\varepsilon\), valid because \(L\ge1\).  Combining
this with the second tail bound gives
\(r=\varepsilon+\BigOAt{\varepsilon^{2}}{X\to+\infty}\), whence
\(r/X=\varepsilon/X+\BigOAt{\varepsilon^{2}/X}{X\to+\infty}\) and consequently
\begin{equation}\label{plt:eq:mot-u-refined}
  u=\varepsilon-\frac{r}{X}
   =\varepsilon-\frac{\varepsilon}{X}
    +\BigOAt{\frac{\varepsilon^{2}}{X}}{X\to+\infty}.
\end{equation}
Squaring \eqref{plt:eq:mot-u-refined} and using \(r=\BigOAt{\varepsilon}{X\to+\infty}\),
\[
  u^{2}=\Bigl(\varepsilon-\frac{r}{X}\Bigr)^{2}
       =\varepsilon^{2}-\frac{2\varepsilon r}{X}+\frac{r^{2}}{X^{2}}
       =\varepsilon^{2}+\BigOAt{\frac{\varepsilon^{2}}{X}}{X\to+\infty}.
\]
Adding the three contributions,
\[
  r=u+\frac{u^{2}}{2}+\BigOAt{\varepsilon^{3}}{X\to+\infty}
   =\varepsilon-\frac{\varepsilon}{X}+\frac{\varepsilon^{2}}{2}
    +\BigOAt{\frac{\varepsilon^{2}}{X}}{X\to+\infty}
    +\BigOAt{\varepsilon^{3}}{X\to+\infty}.
\]
Substituting \(\varepsilon=L/X\), the three main terms are \(L/X\),
\(-L/X^{2}\) and \(L^{2}/(2X^{2})\), while the two error terms are
\(\BigOAt{L^{2}/X^{3}}{X\to+\infty}\) and
\(\BigOAt{L^{3}/X^{3}}{X\to+\infty}\); both are
\(\BigOAt{L^{3}/X^{3}}{X\to+\infty}\) because \(L\ge1\).  Since
\(\WrightOmega(X)=X-L+r\), this is exactly \eqref{plt:eq:mot-two-orders}.

The two asymptotic equivalences follow: \(\WrightOmega(X)-X=-L+r\) with
\(0\le r\le2L/X=\LittleOAt{L}{X\to+\infty}\), and
\(\WrightOmega(X)-X+L=r\sim\varepsilon=L/X\).
\end{proof}

\begin{remark}[What the two coefficients are]
\label{plt:rmk:mot-lambert-preview}
The three coefficient blocks produced by the computation are the members
\(n=0,1,2\) of the canonical zero-based Lambert family; the next two members,
recorded here for orientation but derived only in Part~V, continue the
pattern:
\begin{gather}
  \LamP{0}(u)=-u,\qquad
  \LamP{1}(u)=u,\qquad
  \LamP{2}(u)=\frac{u^{2}}{2}-u,
  \label{plt:eq:mot-lambert-first}\\
  \LamP{3}(u)=\frac{u^{3}}{3}-\frac{3}{2}u^{2}+u,\qquad
  \LamP{4}(u)=\frac{u^{4}}{4}-\frac{11}{6}u^{3}+3u^{2}-u,
  \notag
\end{gather}
defined and studied in Part~V.  In that normalization
\eqref{plt:eq:mot-two-orders} reads
\[
  \WrightOmega(X)=X+\sum_{n=0}^{2}\LamP{n}(L)\,t^{n}
   +\BigOAt{L^{3}t^{3}}{X\to+\infty},
\]
the block \(n=0\) being the term \(-L\).  The degree law \(\deg\LamP{n}=n\)
holds for \(n\ge1\) only: \(\LamP{0}(u)=-u\) has degree one, so a blanket
statement ``\(\LamP{n}\) has degree \(n\)'' over \(n\in\NaturalNumbers\) is
false and is not made anywhere in this volume.
\end{remark}

\begin{remark}[The same computation for \(\LambertW\)]
\label{plt:rmk:mot-lambert-transfer}
Substituting \(X=\log z\) into \eqref{plt:eq:mot-two-orders} and using
\(\LambertW_{0}(z)=\WrightOmega(\log z)\) from
\cref{plt:prop:mot-omega-basic} gives, with the iterated logarithms
\(L_{1}=\log z\) and \(L_{2}=\log L_{1}\),
\begin{equation}\label{plt:eq:mot-W-two-orders}
  \LambertW_{0}(z)=L_{1}-L_{2}+\frac{L_{2}}{L_{1}}
   +\frac{L_{2}^{2}/2-L_{2}}{L_{1}^{2}}
   +\BigOAt{\frac{L_{2}^{3}}{L_{1}^{3}}}{z\to+\infty},
\end{equation}
the classical large-argument expansion \cite{corless-et-al-1996,dlmf-W}.  This
is the sense in which the algebra of this volume sits one level below Lambert's
function: the two nested logarithms of \(z\) become the single pair
\((X,L)\).  Here \(L_{1}\) and \(L_{2}\) are iterated logarithms and never
Lambert polynomials; the Lambert family is always written \(\LamP{n}\).
\end{remark}

\section{Asymptotic scales, dominance, and blocks}
\label{plt:sec:mot-dominance}

Before showing that no one-generator scale works, we fix what a scale is and
prove the one comparison on which everything rests.

\begin{definition}[Asymptotic scale and expansion; working definition]
\label{plt:def:mot-scale}
An \emph{asymptotic scale} at \(+\infty\) is a set \(\mathfrak{M}\) of
functions, each positive on some half-line, such that for all distinct
\(m,m'\in\mathfrak{M}\) either \(m/m'\to0\) or \(m'/m\to0\) as
\(X\to+\infty\).  A function \(f\) defined near \(+\infty\) has the
\emph{asymptotic expansion} \(\sum_{i\ge1}c_{i}m_{i}\) with respect to
\(\mathfrak{M}\) if \(c_{i}\in\RealNumbers\setminus\SetOf{0}\), the
\(m_{i}\in\mathfrak{M}\) satisfy \(m_{i+1}/m_{i}\to0\), and for every index
\(N\) present in the (finite or infinite) list,
\[
  f-\sum_{i\le N}c_{i}m_{i}=\LittleOAt{m_{N}}{X\to+\infty}.
\]
A finite such expansion is \emph{maximal} when it admits no further term, that
is, when the final remainder is not asymptotically equivalent to a constant
multiple of any element of \(\mathfrak{M}\).
\end{definition}
\emph{Lean status.}  This working definition has a formal counterpart in the sequence-indexed
form the rest of the corpus uses: \lean{Fabius.IsAsymptoticScale} in
\lean{TransseriesScale.lean}, for an arbitrary filter on an arbitrary normed field, asks that
each \(\varphi_{n}\) be eventually nonvanishing and that \(\varphi_{n+1}=\littleo(\varphi_{n})\),
and \lean{Fabius.IsPoincareExpansion} is the expansion relation.  That the
power-logarithmic monomials form such a scale is \lean{Fabius.isAsymptoticScale_plMonomial} in
\lean{TransseriesPolyLogScale.lean}, for any lexicographically decreasing sequence of exponent
pairs, with the two ladders used most often supplied ready-made:
\lean{Fabius.isAsymptoticScale_plMonomial_pow} for \(X^{-n}\) and
\lean{Fabius.isAsymptoticScale_plMonomial_log} for \(X^{c}(\log X)^{-n}\).  The set-indexed
formulation above, quantifying over an unordered \(\mathfrak M\), is not formalized as such;
the sequence form is what the uniqueness theorem \cref{q0:prop:uniqueness} is stated against.


This is the coarse blockwise notion adequate for a motivating discussion.
Part~VI replaces it by the sharper statements needed for analytic work,
including remainder estimates with the logarithmic degree printed and the exact
sense in which a truncated polynomial-logarithmic series approximates a
function.

\begin{lemma}[Lexicographic dominance of power-logarithmic monomials]
\label{plt:lem:mot-dominance}
Let \((a,b),(a',b')\in\RealNumbers^{2}\).  As \(X\to+\infty\),
\[
  \frac{X^{a}(\log X)^{b}}{X^{a'}(\log X)^{b'}}
  \longrightarrow
  \begin{cases}
    0,&\text{if }a<a',\text{ or }a=a'\text{ and }b<b',\\
    1,&\text{if }(a,b)=(a',b'),\\
    +\infty,&\text{if }a>a',\text{ or }a=a'\text{ and }b>b'.
  \end{cases}
\]
Consequently \(\SetOf{X^{a}(\log X)^{b}:(a,b)\in\RealNumbers^{2}}\) is an
asymptotic scale in the sense of \cref{plt:def:mot-scale}, totally ordered by
the lexicographic order on \((a,b)\); and in the generators
\eqref{plt:eq:mot-generators}, for \(n,m,j,k\in\IntegerNumbers\),
\begin{equation}\label{plt:eq:mot-dominance}
  t^{n}L^{j}\succdom t^{m}L^{k}
  \quad\Longleftrightarrow\quad
  n<m,\ \text{ or }\ (n=m\ \text{ and }\ j>k).
\end{equation}
\end{lemma}

\begin{proof}
Put \(s=\log X\), so \(s\to+\infty\), and write the ratio as
\[
  X^{a-a'}(\log X)^{b-b'}=\EulerE^{(a-a')s}s^{\,b-b'}
  =\exp\bigl((a-a')s+(b-b')\log s\bigr).
\]
If \(a<a'\) then \((a-a')s+(b-b')\log s\to-\infty\), because
\(\log s=\LittleOAt{s}{s\to+\infty}\) and \(a-a'<0\); hence the ratio tends to
\(0\).  The case \(a>a'\) is the reciprocal statement.  If \(a=a'\) the ratio
is \(s^{\,b-b'}\), which tends to \(0\), equals \(1\), or tends to
\(+\infty\) according as \(b<b'\), \(b=b'\), or \(b>b'\).  Since the trichotomy
is exhaustive, distinct exponent pairs give ratios tending to \(0\) or to
\(+\infty\), which is exactly the scale condition of
\cref{plt:def:mot-scale}.

% ed.: source draft 5 justifies the lexicographic order only in the case
% alpha > alpha', leaving the tie-breaking case alpha = alpha' unproved.  The
% display above proves both cases.
For \eqref{plt:eq:mot-dominance}, apply the trichotomy with \((a,b)=(-n,j)\)
and \((a',b')=(-m,k)\): the monomial \(t^{n}L^{j}\) dominates \(t^{m}L^{k}\)
exactly when \(-n>-m\), i.e.\ \(n<m\), or when \(n=m\) and \(j>k\).
\end{proof}
\emph{Lean status.}  Formal in \lean{TransseriesScaleDominance.lean}, with real exponents in
both slots: \lean{Fabius.plMonomial} is \(X^{a}(\log X)^{b}\), and the three cases of the
displayed limit are \lean{Fabius.tendsto_plMonomial_div_atTop_zero},
\lean{Fabius.tendsto_plMonomial_div_atTop_one} and \lean{Fabius.tendsto_plMonomial_div_atTop};
the workhorse is \lean{Fabius.tendsto_plMonomial_atTop_zero}, that \(X^{c}(\log X)^{d}\to0\)
for every \(c<0\) and \emph{every} real \(d\).  The lexicographic law in the volume's own
generators \(t=1/X\), \(L=\log X\) --- \eqref{plt:eq:mot-dominance} --- is
\lean{Fabius.plMonomial_generators_dominance}, with positivity past \(X=1\) recorded as
\lean{Fabius.plMonomial_pos} and the quotient rewritten by
\lean{Fabius.plMonomial_div_eventuallyEq}.  The dominance in \(\littleo\) form is
\lean{Fabius.isLittleO_plMonomial} in \lean{TransseriesPolyLogScale.lean}.


The dominance order is not a convention chosen to make bookkeeping convenient.
The grouping of monomials into inverse-power blocks is intrinsic.

\begin{proposition}[Blocks are the logarithmic-comparability classes]
\label{plt:prop:mot-blocks}
Let \(\mu=t^{n}L^{j}\) and \(\nu=t^{m}L^{k}\) with
\(n,m,j,k\in\IntegerNumbers\).  The following are equivalent.
\begin{enumerate}
\item \(n=m\).
\item There is an integer \(N\ge0\) with
  \((\log X)^{-N}\le\mu/\nu\le(\log X)^{N}\) for all sufficiently large \(X\).
\end{enumerate}
Hence ``differ by at most a fixed power of \(\log X\)'' is an equivalence
relation on the group of monomials \(t^{n}L^{j}\), its classes are exactly the
sets \(\SetOf{t^{n}L^{j}:j\in\IntegerNumbers}\) for fixed \(n\), and
\(\succdom\) descends to the total order of \(-n\) on the classes.
\end{proposition}

\begin{proof}
(i)\(\Rightarrow\)(ii): if \(n=m\) then \(\mu/\nu=L^{\,j-k}\), so
\(N=\AbsoluteValue{j-k}\) works, with equality on one side.

(ii)\(\Rightarrow\)(i): suppose \(n\ne m\).  If \(n>m\) then, with
\(s=\log X\), we have \(\mu/\nu=\EulerE^{(m-n)s}s^{\,j-k}\) and, for every
\(N\),
\[
  \frac{\mu/\nu}{s^{-N}}=\EulerE^{(m-n)s}s^{\,j-k+N}\longrightarrow0
\]
by \cref{plt:lem:mot-dominance}, since \(m-n<0\); so the lower bound in (ii)
fails for every \(N\).  If \(n<m\), the same computation applied to
\(\nu/\mu\) shows that the upper bound fails.

For the last sentence: the relation is reflexive (\(N=0\)) and symmetric
(replace \(N\) by \(N\)), and transitive with \(N\) replaced by the sum of the
two constants; by the equivalence just proved, its classes are the fibres of
\(t^{n}L^{j}\mapsto n\).  That \(\succdom\) induces the order of \(-n\) on
classes is \eqref{plt:eq:mot-dominance}.
\end{proof}

\Cref{plt:prop:mot-blocks} is the structural justification for the normal form
used throughout this volume: a polynomial-logarithmic transseries is written
\begin{equation}\label{plt:eq:mot-normal-form}
  F=\sum_{n\ge n_{0}}p_{n}(L)\,t^{n},
  \qquad p_{n}\in\K[L],\quad n_{0}\in\IntegerNumbers,
\end{equation}
one finite polynomial per comparability class, the classes indexed by the
integer \(n\) and ordered by \(-n\).  The formal ring housing
\eqref{plt:eq:mot-normal-form} is \(\PLlau=\K[L]((t))\), with power-series
subring \(\PLpow=\K[L][[t]]\); their construction, valuation, leading data and
truncation topology occupy the rest of Part~I.

\begin{remark}[Where the class ends]
\label{plt:rmk:mot-boundary}
Three failures of closure mark the exact boundary of
\eqref{plt:eq:mot-normal-form}, and each names its own minimal enlargement.
\begin{itemize}
\item \(\MultiplicativeInverseOf{(L+1)}\) is not a polynomial in \(L\).
  Unrestricted division needs the coefficient field \(\K(L)\), giving
  \(\PLrat=\K(L)((t))\); expanding those rational coefficients at
  \(L=\infty\) gives instead \(\PLit=\K((L^{-1}))((t))\).  These are different
  objects: \(1/(L+1)\) is exact in the first and becomes
  \(L^{-1}-L^{-2}+L^{-3}-\cdots\) in the second.
\item \(\log L\) arises from composing with an inner argument whose leading
  behaviour is logarithmic, and requires a second logarithmic level.
% ed.: the draft said X^{1/2} "arises from inverting a non-integer leading
% power".  That is backwards, and contradicts the reading of the Fabius
% endpoint law in the remark on its status: a half-integer exponent is produced
% by inverting an *integer* leading power other than the first (X^2 inverts to
% X^{1/2}); an inner argument that is already Puiseux has left the class before
% the inversion begins.
\item \(X^{1/2}\) arises from inverting a leading behaviour that is not
  linear: a leading power \(X^{k}\) with \(k\ne1\) inverts to one with
  exponents in \(k^{-1}\IntegerNumbers\).  This requires a Puiseux or Hahn
  exponent group.
\end{itemize}
% ed.: source draft 6 prints the refined scale as
% "..., X^{-n}L^2, X^{-n}L, X^{-n}, X^{-n}L^{-1}, ...", which silently places
% negative powers of L inside the polynomial-logarithmic class.  They are not
% there: coefficient blocks of K[L]((t)) are polynomials, so L-exponents are
% nonnegative, and L^{-1} belongs to the enlargements named above.
Negative powers of \(L\) do not belong to \(\PLlau\).  A display that mixes
\(L^{-1}\) into a coefficient block has already moved to \(\PLrat\) or
\(\PLit\) and must say so.  These enlargements are the subject of Part~VIII.
\end{remark}

\section{Why one generator is never enough}
\label{plt:sec:mot-one-generator}

We now show that neither obvious one-generator scale can carry the expansion
\eqref{plt:eq:mot-two-orders}.  The two failures are of different kinds and
deserve different names: the power scale \emph{stalls}, the logarithmic scale
is \emph{blind}.

\begin{proposition}[Both one-generator scales fail]
\label{plt:prop:mot-one-generator-fails}
Set \(\psi_{1}(X)=\WrightOmega(X)-X\) and
\(\psi_{2}(X)=\WrightOmega(X)-X+\log X\), and let
\[
  \mathfrak{M}_{\mathrm{pow}}=\SetOf{X^{a}:a\in\RealNumbers},
  \qquad
  \mathfrak{M}_{\log}=\SetOf{(\log X)^{b}:b\in\RealNumbers},
\]
both asymptotic scales by \cref{plt:lem:mot-dominance}.
\begin{enumerate}
\item \emph{(The power scale stalls.)}  With respect to
  \(\mathfrak{M}_{\mathrm{pow}}\), the function \(\WrightOmega\) has the
  one-term expansion \(X\), and that expansion is maximal.  Indeed
  \(\psi_{1}\sim-\log X\), so
  \(X^{a}=\LittleOAt{\AbsoluteValue{\psi_{1}}}{X\to+\infty}\) for every
  \(a\le0\) while \(\psi_{1}=\LittleOAt{X^{a}}{X\to+\infty}\) for every
  \(a>0\); no element of \(\mathfrak{M}_{\mathrm{pow}}\) is asymptotically
  proportional to \(\psi_{1}\).
\item \emph{(The logarithmic scale is blind.)}  With respect to
  \(\mathfrak{M}_{\log}\), the function \(\WrightOmega\) has no expansion at
  all, because \(\WrightOmega(X)/(\log X)^{b}\to+\infty\) for every
  \(b\in\RealNumbers\).  Moreover the residual \(\psi_{2}\) is strictly
  positive for large \(X\) and satisfies
  \(\psi_{2}=\LittleOAt{(\log X)^{b}}{X\to+\infty}\) for \emph{every}
  \(b\in\RealNumbers\); so \(\mathfrak{M}_{\log}\) assigns \(\psi_{2}\) the
  empty expansion and cannot distinguish \(\psi_{2}\) from \(0\).
\item \emph{(One mixed monomial repairs both failures.)}
  \(\psi_{2}(X)\sim X^{-1}\log X\).
\end{enumerate}
\end{proposition}

\begin{proof}
(i) By \cref{plt:prop:mot-two-orders}, \(\psi_{1}\sim-\log X\).  For
\(a\le0\) we have \(X^{a}\le1\) for \(X\ge1\) while
\(\AbsoluteValue{\psi_{1}}\to+\infty\), so
\(X^{a}/\AbsoluteValue{\psi_{1}}\to0\); for \(a>0\),
\(\psi_{1}/X^{a}\to0\) by \cref{plt:lem:mot-dominance} applied to
\(\log X\) and \(X^{a}\).  Hence there are no \(a\in\RealNumbers\) and
\(c\ne0\) with \(\psi_{1}\sim cX^{a}\): for \(a>0\) the quotient
\(\psi_{1}/X^{a}\) tends to \(0\), and for \(a\le0\) it is unbounded.  Finally
\(\WrightOmega\sim X\) by \eqref{plt:eq:mot-envelope}, and
\(\WrightOmega-X=\psi_{1}=\LittleOAt{X}{X\to+\infty}\), so \(X\) is a
legitimate first --- and, by the above, last --- term.

(ii) By \eqref{plt:eq:mot-envelope}, \(\WrightOmega(X)\ge X-\log X\), and
\((X-\log X)/(\log X)^{b}=X(\log X)^{-b}-(\log X)^{1-b}\to+\infty\) by
\cref{plt:lem:mot-dominance}, since the first summand dominates the second.  So
no leading term exists.  For the second claim,
\eqref{plt:eq:mot-psi2-envelope} gives
\(\tfrac12\log X/X\le\psi_{2}(X)\le2\log X/X\) for \(X\ge X_{0}\); in
particular \(\psi_{2}>0\).  For any \(b\),
\(\psi_{2}(X)(\log X)^{-b}\le2(\log X)^{1-b}/X\to0\), again by
\cref{plt:lem:mot-dominance}.

(iii) This is the last assertion of \cref{plt:prop:mot-two-orders}.
\end{proof}

\begin{remark}[The two failures are not symmetric]
\label{plt:rmk:mot-asymmetry}
Part~(i) is a \emph{gap}: the true size of the correction, \(\log X\), lies
strictly between \(X^{0}\) and \(X^{a}\) for every \(a>0\), so the power scale
simply runs out of monomials at the place where one is needed.  Part~(ii) is a
\emph{loss of resolution}: the logarithmic scale has monomials on both sides of
\(\psi_{2}\), but every one of them is too coarse, so an expansion exists and
carries no information.  A treatment that reports only ``a pure power series is
impossible'' has stated the weaker of the two obstructions.
\end{remark}

\section{The smallest closed class}
\label{plt:sec:mot-smallest}

The negative results say what fails.  The positive statement is an exact
minimality theorem, and it is where the pair \((t,L)\) is singled out.  Let
\(\mathcal{G}\) denote the ring of germs at \(+\infty\) of real functions
smooth on some half-line, with the derivation
\(\DX=\Differential/\Differential X\).

\begin{theorem}[The two-generator ring is the smallest differentially closed
class containing a logarithm]
\label{plt:thm:mot-smallest-differential-algebra}
Inside \(\mathcal{G}\), with \(t=X^{-1}\) and \(L=\log X\):
\begin{enumerate}
\item The smallest \(\DX\)-stable subring of \(\mathcal{G}\) containing
  \(\RationalNumbers\) and \(X\) is \(\RationalNumbers[X]\).  Every unbounded
  element of it grows at least as fast as \(X\); in particular it contains no
  element asymptotic to a nonzero multiple of \(\log X\).
\item The smallest \(\DX\)-stable subring of \(\mathcal{G}\) containing
  \(\RationalNumbers\) and \(L\) is \(\RationalNumbers[t,L]\).
\item The smallest \(\DX\)-stable subring of \(\mathcal{G}\) containing
  \(\RationalNumbers\), \(X\) and \(L\) is
  \(\RationalNumbers[X,t,L]=\RationalNumbers[t,t^{-1}][L]\), the ring of
  Laurent polynomials in \(t\) with coefficients polynomial in \(L\).
\item \(t\) and \(L\) are algebraically independent over \(\RealNumbers\).
  Hence \(\RealNumbers[t,L]\) is a genuine two-variable polynomial ring, and
  each of its elements has a \emph{unique} representation
  \(\sum_{n\ge0}p_{n}(L)t^{n}\) with \(p_{n}\in\RealNumbers[L]\) and only
  finitely many \(p_{n}\ne0\).
\item For \(n\in\IntegerNumbers\) and \(p\in\RationalNumbers[L]\),
  \begin{equation}\label{plt:eq:mot-block-derivative}
    \DX\bigl(t^{n}p(L)\bigr)=t^{n+1}\bigl(p'(L)-np(L)\bigr).
  \end{equation}
\end{enumerate}
The same statements hold with \(\RationalNumbers\) replaced by any subfield of
\(\RealNumbers\).  Abstractly, for a field \(\K\) of characteristic zero, the
ring in (iii) is the Laurent polynomial ring \(\K[t,t^{-1}][L]\) carrying the
unique derivation with \(\DX t=-t^{2}\) and \(\DX L=t\).
\end{theorem}

\begin{proof}
We first verify (v).  As germs, \(\DX t=-X^{-2}=-t^{2}\) and
\(\DX L=X^{-1}=t\); by the product and chain rules,
\[
  \DX\bigl(t^{n}p(L)\bigr)
  =nt^{n-1}(-t^{2})p(L)+t^{n}p'(L)\,t
  =t^{n+1}\bigl(p'(L)-np(L)\bigr).
\]

(i) \(\RationalNumbers[X]\) contains \(\RationalNumbers\) and \(X\) and is
\(\DX\)-stable, since \(\DX X=1\).  Conversely any \(\DX\)-stable subring
containing \(\RationalNumbers\) and \(X\) contains all powers of \(X\) and all
rational combinations of them, hence contains \(\RationalNumbers[X]\).  A
nonzero element of \(\RationalNumbers[X]\) is a polynomial of some degree
\(d\); if \(d=0\) it is a constant, and if \(d\ge1\) it is asymptotic to a
nonzero multiple of \(X^{d}\).  Since \(\log X/X^{d}\to0\) for \(d\ge1\) and
\(\log X\to\infty\), no element is asymptotic to a nonzero multiple of
\(\log X\).

(ii) Put \(R=\RationalNumbers[t,L]\).  It contains \(\RationalNumbers\) and
\(L\).  It is \(\DX\)-stable: a derivation is determined on a finitely
generated ring by its values on generators, and \(\DX t=-t^{2}\in R\),
\(\DX L=t\in R\), so \(\DX R\subseteq R\).  Conversely let \(S\) be
\(\DX\)-stable with \(\RationalNumbers\cup\SetOf{L}\subseteq S\).  Then
\(t=\DX L\in S\), so \(S\) contains the subring generated by
\(\RationalNumbers\), \(t\) and \(L\), namely \(R\).

(iii) Put \(R'=\RationalNumbers[X,t,L]\).  Since \(Xt=1\), \(R'\) is the ring
of Laurent polynomials in \(t\) with coefficients in
\(\RationalNumbers[L]\).  It is \(\DX\)-stable by \(\DX X=1\),
\(\DX t=-t^{2}\), \(\DX L=t\).  Any \(\DX\)-stable subring containing
\(\RationalNumbers\), \(X\) and \(L\) contains \(\DX L=t\), hence contains
\(R'\).

(iv) Suppose \(P\in\RealNumbers[T,\Lambda]\) is nonzero with \(P(t,L)=0\) as a
germ, i.e.\ \(P(X^{-1},\log X)=0\) for all large \(X\).  Write
\(P=\sum_{n\ge0}T^{n}q_{n}(\Lambda)\) with
\(q_{n}\in\RealNumbers[\Lambda]\), and let \(n_{0}\) be minimal with
\(q_{n_{0}}\ne0\).  Multiplying the relation by \(X^{n_{0}}\) gives, for all
large \(X\),
\[
  q_{n_{0}}(\log X)=-\sum_{n>n_{0}}X^{-(n-n_{0})}q_{n}(\log X).
\]
Each summand on the right is \(X^{-k}\) with \(k\ge1\) times a polynomial in
\(\log X\), hence tends to \(0\) by \cref{plt:lem:mot-dominance}; the sum is
finite, so the right side tends to \(0\).  The left side is a nonzero
polynomial evaluated along \(\log X\to+\infty\), so it tends to \(\pm\infty\)
or is a nonzero constant.  This contradiction proves algebraic independence;
uniqueness of the representation follows, since two representations of the same
germ would differ by a nontrivial polynomial relation.

For the abstract version, \(\K[t,t^{-1}][L]\) is the localization at \(t\) of
the polynomial ring \(\K[t,L]\).  A derivation of a polynomial ring is
determined freely by its values on the variables, and it extends uniquely to
any localization by the quotient rule; taking \(\DX t=-t^{2}\), \(\DX L=t\)
gives the derivation in question.
\end{proof}

\begin{remark}[Reading the theorem]
\label{plt:rmk:mot-reading-minimality}
Part~(i) says that the pure power world is closed but expressively poor: a
logarithm can never be produced from \(X\) by ring operations and
differentiation, so if a problem's answer contains a logarithm, that logarithm
must enter through an operation --- inversion, antidifferentiation,
composition --- that leaves the class.  Part~(ii) says the converse asymmetry
does not exist: once \(\log X\) is admitted, \(X^{-1}\) is forced immediately,
and then all of \(\RationalNumbers[t,L]\).  There is no ``logarithms only''
differential algebra.  This is the precise sense in which powers and logarithms
belong in one expansion.
% ed.: the draft called PLlau "the t-adic completion of the Laurent ring in
% (iii)".  That is exactly the error item 7 of the normalizations remark warns
% against: t is a unit in K[t,t^{-1}][L], so every power of the ideal (t) is
% the whole ring, the t-adic filtration is constant, and the associated
% completion collapses to zero.  The correct construction completes the
% polynomial subring K[L][t] and only then inverts t.
The ambient object \(\PLlau=\K[L]((t))\) of Part~I is obtained from the
Laurent ring of~(iii) by completing its \emph{polynomial} subring
\(\K[L][t]\) in the \(t\)-adic topology, which gives \(\PLpow=\K[L][[t]]\), and
then inverting \(t\).  Equivalently, \(\PLlau\) is the completion of
\(\K[t,t^{-1}][L]\) for the valuation \(\ordt\).  One may not complete
\(\K[t,t^{-1}][L]\) \(t\)-adically: there \(t\) is a unit, so \((t^{N})\) is
the whole ring for every \(N\) and the \(t\)-adic completion is trivial.
\end{remark}

\begin{corollary}[Both exponent directions are forced by \(\WrightOmega\)]
\label{plt:cor:mot-both-generators-needed}
Let \(S\subseteq\RealNumbers^{2}\) and put
\(\mathfrak{M}_{S}=\SetOf{X^{a}(\log X)^{b}:(a,b)\in S}\), an asymptotic scale
by \cref{plt:lem:mot-dominance}.  If \(\WrightOmega\) has an asymptotic
expansion with respect to \(\mathfrak{M}_{S}\) with at least three terms, then
\[
  (1,0)\in S,\qquad (0,1)\in S,\qquad (-1,1)\in S,
\]
with coefficients \(1\), \(-1\), \(1\) respectively.  Consequently the subgroup
of \((\RealNumbers^{2},+)\) generated by \(S\) contains
\(\IntegerNumbers^{2}\); in particular \(S\) is contained neither in
\(\RealNumbers\times\SetOf{0}\) (a scale of pure powers) nor in
\(\SetOf{0}\times\RealNumbers\) (a scale of pure logarithms).
\end{corollary}

\begin{proof}
By \cref{plt:lem:mot-dominance}, a monomial \(X^{a}(\log X)^{b}\) can be
asymptotically equivalent to \(cX^{a'}(\log X)^{b'}\) with \(c\in(0,\infty)\)
only if \((a,b)=(a',b')\) and \(c=1\), since otherwise the ratio tends to
\(0\) or to \(+\infty\).

By \cref{plt:def:mot-scale}, the first term satisfies
\(\WrightOmega\sim c_{1}m_{1}\).  By \eqref{plt:eq:mot-envelope},
\(\WrightOmega\sim X=X^{1}(\log X)^{0}\), so \(m_{1}\sim c_{1}^{-1}X\); by the
previous paragraph \(m_{1}=X\) and \(c_{1}=1\), whence \((1,0)\in S\).  The
first partial sum is therefore exactly \(X\), and the second term satisfies
\(\WrightOmega-X\sim c_{2}m_{2}\).  By \cref{plt:prop:mot-two-orders},
\(\WrightOmega-X\sim-\log X\); as \(m_{2}>0\), this forces \(c_{2}<0\) and
\(m_{2}\sim\AbsoluteValue{c_{2}}^{-1}\log X\), hence \(m_{2}=\log X\) and
\(c_{2}=-1\), so \((0,1)\in S\).  The second partial sum is exactly
\(X-\log X\), and the third term satisfies
\(\WrightOmega-X+\log X\sim c_{3}m_{3}\); by
\cref{plt:prop:mot-two-orders} this residual is \(\sim X^{-1}\log X\), so
\(m_{3}=X^{-1}\log X\), \(c_{3}=1\) and \((-1,1)\in S\).  Finally \((1,0)\) and
\((0,1)\) generate \(\IntegerNumbers^{2}\).
\end{proof}

\begin{remark}[Scope of the corollary]
\label{plt:rmk:mot-corollary-scope}
The corollary is stated for monomial scales in the two generators, and the
restriction is not cosmetic.  Nothing above rules out an exotic scale built
from other generators --- for instance one containing \(X+X^{1/2}\) and
\(X^{1/2}\) --- in which the first terms of \(\WrightOmega\) look different.
What \emph{is} ruled out, by \cref{plt:prop:mot-one-generator-fails}, is any
scale of powers alone or of logarithms alone, whatever the exponents; and what
is asserted by \cref{plt:thm:mot-smallest-differential-algebra} is that among
differentially closed rings the two-generator ring is minimal.  Together these
are the honest content of the phrase ``smallest closed class''.  A stronger
uniqueness statement --- that \(\PLlau\) is the \emph{unique} minimal
differential ring in which the reversion of \(X+L\) exists --- is not proved in
this volume and should not be inferred from it.
\end{remark}

\section{Where a logarithm is created}
\label{plt:sec:mot-creation}

\Cref{plt:thm:mot-smallest-differential-algebra}(i) says a logarithm cannot be
manufactured from powers by differentiation.  It can be manufactured by
\emph{anti}differentiation, and by the discrete analogue of
antidifferentiation, summation.  These two mechanisms account for nearly every
occurrence of the scale in practice, so we isolate them.

\begin{lemma}[Block antidifferentiation and the resonant block]
\label{plt:lem:mot-block-antiderivative}
Let \(\K\) be a field of characteristic zero and work in the differential ring
\(\K[t,t^{-1}][L]\) of
\cref{plt:thm:mot-smallest-differential-algebra}(iii).  Let \(c\in\K\) and
consider the \(\K\)-linear operator \(\partial_{L}-c\) on \(\K[L]\).
\begin{enumerate}
\item If \(c\ne0\), then \(\partial_{L}-c\) is a degree-preserving bijection of
  \(\K[L]\) onto itself.  Consequently, for \(n\ne1\) and \(p\in\K[L]\), the
  element \(t^{n}p(L)\) has a unique antiderivative of the shape
  \(t^{n-1}q(L)\) with \(q\in\K[L]\), and \(\deg q=\deg p\).
\item If \(c=0\), then \(\partial_{L}\) is surjective on \(\K[L]\) with kernel
  \(\K\), and \(\deg q=\deg(\partial_{L}q)+1\) for nonconstant \(q\).
  Consequently \(t\,p(L)\) has antiderivative \(q(L)\) with \(q'=p\), unique up
  to an additive constant, and \(\deg q=\deg p+1\).
\end{enumerate}
In particular \(t=X^{-1}\) integrates to \(L=\log X\).  The outer level
\(n=1\) is the unique block at which antidifferentiation raises logarithmic
degree.
\end{lemma}

\begin{proof}
(i) Write \(q=\sum_{i=0}^{d}q_{i}L^{i}\) with \(q_{d}\ne0\).  Then
\((\partial_{L}-c)q=-cq_{d}L^{d}+(\text{terms of degree}<d)\), so the operator
preserves degree, hence is injective and maps the finite-dimensional space of
polynomials of degree at most \(d\) into itself injectively, therefore
bijectively; taking the union over \(d\) gives bijectivity on \(\K[L]\).  For
the antiderivative statement, apply \eqref{plt:eq:mot-block-derivative} with
\(n-1\) in place of \(n\):
\(\DX(t^{n-1}q)=t^{n}\bigl(q'-(n-1)q\bigr)=t^{n}(\partial_{L}-c)q\) with
\(c=n-1\ne0\); solve \((\partial_{L}-c)q=p\), which has a unique solution of
the same degree.

(ii) \(\partial_{L}\) is surjective on \(\K[L]\) because
\(\partial_{L}\bigl(L^{i+1}/(i+1)\bigr)=L^{i}\), which uses
\(\operatorname{char}\K=0\); its kernel is \(\K\); and it lowers degree by
exactly one on nonconstant polynomials.  Taking \(n=1\) in
\eqref{plt:eq:mot-block-derivative}, with \(n-1=0\), gives
\(\DX q(L)=t\,q'(L)\), so the antiderivatives of \(t\,p(L)\) are exactly the
\(q\) with \(q'=p\).  The case \(p=1\), \(q=L\) is
\(\int X^{-1}\Differential X=\log X\).
\end{proof}
\emph{Lean status.}  Formal, both clauses, in \lean{TransseriesBlockAntiderivative.lean}
(a Mathlib-only leaf).  The operator \(\partial_{L}-c\) is \lean{Fabius.blockOperator}.
For \(c\ne0\) the inverse is \emph{explicit} rather than abstract:
\lean{Fabius.blockAntiderivative} is the finite Neumann series
\(-\sum_{k\le\deg f}c^{-(k+1)}\partial^{k}f\), finite because \(\partial\) is nilpotent in
bounded degree, and \lean{Fabius.blockOperator_blockAntiderivative} makes it a right inverse;
with \lean{Fabius.blockOperator_injective} and \lean{Fabius.natDegree_blockOperator} this gives
the degree-preserving bijection \lean{Fabius.blockOperator_bijective}.  For the resonant block
\(c=0\), \lean{Fabius.resonantAntiderivative} is an explicit primitive, whence
\lean{Fabius.derivative_surjective} over any field of characteristic zero and
\lean{Fabius.natDegree_resonantAntiderivative}, the degree rising by exactly one.  The
telescoping engine is \lean{Fabius.sum_sub_sum_shift}, stated for an arbitrary additive
commutative group.


\begin{lemma}[The discrete counterpart: a harmonic increment makes a logarithm]
\label{plt:lem:mot-harmonic}
Let \((w_{n})_{n\ge n_{1}}\) be real with \(w_{n}\to+\infty\) and suppose that
for constants \(c_{0}>0\) and \(c_{1}\in\RealNumbers\),
\begin{equation}\label{plt:eq:mot-harmonic-hypothesis}
  w_{n+1}=w_{n}+c_{0}+\frac{c_{1}}{w_{n}}
   +\BigOAt{w_{n}^{-2}}{n\to\infty}.
\end{equation}
Then there is a constant \(C\) with
\begin{equation}\label{plt:eq:mot-harmonic-conclusion}
  w_{n}=c_{0}n+\frac{c_{1}}{c_{0}}\log n+C+\LittleOAt{1}{n\to\infty}.
\end{equation}
\end{lemma}

\begin{proof}
Since \(w_{n}\to+\infty\), the last two terms of
\eqref{plt:eq:mot-harmonic-hypothesis} tend to \(0\), so
\(w_{n+1}-w_{n}\to c_{0}\).  By the Stolz--Cesàro theorem applied to the
increments, \(w_{n}/n\to c_{0}\); in particular \(w_{n}\sim c_{0}n\) and
\(1/w_{n}=\BigOAt{n^{-1}}{n\to\infty}\).

Put \(v_{n}=w_{n}-c_{0}n\).  Then
\(v_{n+1}-v_{n}=c_{1}/w_{n}+\BigOAt{w_{n}^{-2}}{n\to\infty}
=\BigOAt{n^{-1}}{n\to\infty}\); summing, \(v_{n}=\BigOAt{\log n}{n\to\infty}\).

Now bootstrap.  Since \(v_{n}=\BigOAt{\log n}{n\to\infty}=\LittleOAt{n}{n\to\infty}\),
we have \(\AbsoluteValue{v_{n}/(c_{0}n)}\le\tfrac12\) for large \(n\), and
\[
  \frac{1}{w_{n}}
  =\frac{1}{c_{0}n}\cdot\frac{1}{1+v_{n}/(c_{0}n)}
  =\frac{1}{c_{0}n}-\frac{v_{n}}{c_{0}^{2}n^{2}}
   +\BigOAt{\frac{v_{n}^{2}}{n^{3}}}{n\to\infty}
  =\frac{1}{c_{0}n}+\BigOAt{\frac{\log n}{n^{2}}}{n\to\infty}.
\]
Hence
\(v_{n+1}-v_{n}=\dfrac{c_{1}}{c_{0}n}
 +\BigOAt{\dfrac{\log n}{n^{2}}}{n\to\infty}\).  Summing from \(n_{1}\) to
\(n-1\), using
\(\sum_{k=n_{1}}^{n-1}k^{-1}=\log n+\gamma-\sum_{k<n_{1}}k^{-1}
+\LittleOAt{1}{n\to\infty}\) and the absolute convergence of
\(\sum_{k}k^{-2}\log k\) (so that its tail is \(\LittleOAt{1}{n\to\infty}\)),
we obtain \eqref{plt:eq:mot-harmonic-conclusion} for a suitable constant
\(C\).
\end{proof}

\begin{example}[Iterating the sine]
\label{plt:ex:mot-sine}
Let \(x_{0}\in(0,\pi)\) and \(x_{n+1}=\sin x_{n}\).  Since
\(0<\sin x\le\min(x,1)\) for \(x\in(0,\pi)\), the sequence stays in
\((0,1]\subset(0,\pi)\) and decreases; its limit \(\ell\) satisfies
\(\ell=\sin\ell\), so \(x_{n}\downarrow0\).  Put \(u_{n}=x_{n}^{-2}\).  From
\(\sin x=x\bigl(1-x^{2}/6+x^{4}/120+\BigOAt{x^{6}}{x\to0}\bigr)\) and
\((1+a)^{-2}=1-2a+3a^{2}+\BigOAt{a^{3}}{a\to0}\) with
\(a=-x^{2}/6+x^{4}/120\),
\[
  \Bigl(\frac{\sin x}{x}\Bigr)^{-2}
  =1+\frac{x^{2}}{3}+\Bigl(\frac{3}{36}-\frac{1}{60}\Bigr)x^{4}
   +\BigOAt{x^{6}}{x\to0}
  =1+\frac{x^{2}}{3}+\frac{x^{4}}{15}+\BigOAt{x^{6}}{x\to0}.
\]
Multiplying by \(u_{n}\) and using \(u_{n}x_{n}^{2}=1\),
\[
  u_{n+1}=u_{n}+\frac13+\frac{1}{15u_{n}}
   +\BigOAt{u_{n}^{-2}}{n\to\infty}.
\]
\Cref{plt:lem:mot-harmonic} with \(c_{0}=1/3\), \(c_{1}=1/15\) gives
\(u_{n}=n/3+\tfrac15\log n+C+\LittleOAt{1}{n\to\infty}\), and therefore
\begin{equation}\label{plt:eq:mot-sine-asymptotic}
  x_{n}=\sqrt{\frac{3}{n}}
   \left(1-\frac{3}{10}\cdot\frac{\log n}{n}
    +\BigOAt{\frac1n}{n\to\infty}\right).
\end{equation}
The coefficient of \(n^{-1}\) is a polynomial of degree one in \(\log n\), and
the whole expression is a Puiseux-logarithmic monomial times a
polynomial-logarithmic factor: with \(X=n\), \(t=X^{-1}\), \(L=\log X\), the
right side of \eqref{plt:eq:mot-sine-asymptotic} is
\(\sqrt3\,t^{1/2}\bigl(1-\tfrac{3}{10}Lt+\cdots\bigr)\), an element of the
Puiseux enlargement \(\K[L]\bigl((t^{1/2})\bigr)\).  Continuing the same
bootstrap produces further blocks, each again a polynomial in \(\log n\).  This
is de~Bruijn's example \cite[Ch.~8]{debruijn}.
\end{example}

\begin{example}[Iterating a parabolic germ]
\label{plt:ex:mot-parabolic}
Let \(f(z)=z-z^{2}+az^{3}+\BigOAt{z^{4}}{z\to0}\) be an analytic germ at \(0\)
tangent to the identity, and let \(z_{n}=f^{\circ n}(z_{0})\) for \(z_{0}>0\)
small enough that \(z_{n}\downarrow0\).  With \(w_{n}=1/z_{n}\), the expansion
\(\bigl(1-z+az^{2}+\BigOAt{z^{3}}{z\to0}\bigr)^{-1}
 =1+z+(1-a)z^{2}+\BigOAt{z^{3}}{z\to0}\) gives
\[
  w_{n+1}=w_{n}\bigl(1+z_{n}+(1-a)z_{n}^{2}+\BigOAt{z_{n}^{3}}{n\to\infty}\bigr)
   =w_{n}+1+\frac{1-a}{w_{n}}+\BigOAt{w_{n}^{-2}}{n\to\infty},
\]
so \cref{plt:lem:mot-harmonic} yields
\(w_{n}=n+(1-a)\log n+C+\LittleOAt{1}{n\to\infty}\) and hence
\[
  z_{n}=\frac1n-\frac{(1-a)\log n}{n^{2}}
   +\BigOAt{\frac{1}{n^{2}}}{n\to\infty}.
\]
The same logarithm appears in the Fatou coordinate \(\Phi\) of \(f\), the
solution of \(\Phi\circ f=\Phi+1\), whose leading behaviour as \(z\downarrow0\)
is \(\Phi(z)=1/z-(1-a)\log(1/z)+\text{const}+\LittleOAt{1}{z\downarrow0}\);
one checks directly that this shape reproduces the increment \(1\) to the
displayed order.  Normal-form and conjugacy theory for such germs is one of the
classical homes of this scale.
\end{example}

\begin{summarybox}
Two mechanisms create the logarithm, and they are one mechanism in continuous
and in discrete clothing.  Antidifferentiation preserves logarithmic degree on
every outer block except \(t^{1}\), where it raises it by one
(\cref{plt:lem:mot-block-antiderivative}); summation of increments is bounded
except for the harmonic one, which produces \(\log n\)
(\cref{plt:lem:mot-harmonic}).  Once one logarithm is present,
\cref{plt:thm:mot-smallest-differential-algebra}(ii) forces the whole ring
\(\K[t,L]\).
\end{summarybox}

\section{Provenance: where these expansions arise}
\label{plt:sec:mot-provenance}

\subsection{Implicit equations carrying a logarithmic term}

The generic source is a relation in which the unknown occurs both linearly and
inside a logarithm,
\begin{equation}\label{plt:eq:mot-generic-implicit}
  y+\alpha\log y+\frac{p_{1}(\log y)}{y}+\frac{p_{2}(\log y)}{y^{2}}+\cdots=X,
  \qquad \alpha\in\K,\quad p_{i}\in\K[\,\cdot\,],
\end{equation}
% ed.: the draft attached the exponentiated form to the full display, where it
% is false --- exponentiating \eqref{plt:eq:mot-generic-implicit} exponentiates
% the higher blocks as well.  The equivalence holds for the two-term equation
% only, and is stated for it.
whose two-term case \(y+\alpha\log y=X\) is equivalent to the exponentiated
form \(y^{\alpha}\EulerE^{y}=\EulerE^{X}\).  Dominant
balance gives \(y\sim X\); substituting that into the logarithm creates
\(\log X\); iterating creates powers of \(X^{-1}\) carrying polynomials in
\(\log X\).  Part~IV turns this heuristic into a theorem: read as a transseries
in the variable \(y\), the left side of \eqref{plt:eq:mot-generic-implicit}
lies in the near-identity class \(\Gnear\), and \(\Gnear\) is a group under
composition, so the solution is again of the same shape.

A classical instance outside the Lambert family is the size of the \(n\)-th
prime.  Inverting \(\pi(p_{n})=n\) against the logarithmic-integral
approximation gives Cipolla's expansion
\begin{equation}\label{plt:eq:mot-cipolla}
  \frac{p_{n}}{n}=\log n+\log\log n-1
   +\frac{\log\log n-2}{\log n}
   +\BigOAt{\frac{(\log\log n)^{2}}{(\log n)^{2}}}{n\to\infty}.
\end{equation}
Read with \(X=\log n\) and \(L=\log\log n\), the right side of
\eqref{plt:eq:mot-cipolla} is a truncation of an element of
\(X+\K[L][[t]]=\Gnear\): the very class that houses \(\WrightOmega\).  Quantile
asymptotics for laws with exponential-type tails have the same shape, for the
same reason --- one inverts a relation in which the unknown appears both
linearly and logarithmically.

\subsection{Dulac maps and first-return maps of hyperbolic sectors}

Let \(v\) be a real analytic planar vector field with a hyperbolic saddle at
the origin, and let \(D\) be the transition (Dulac, or corner) map from a
transversal on one separatrix to a transversal on the other, written in
coordinates \(x\downarrow0\) on the transversals.  Then \(D\) has an asymptotic
expansion in the power-logarithmic scale at zero,
\begin{equation}\label{plt:eq:mot-dulac-series}
  D(x)\;\sim\;\sum_{i\ge0}x^{\lambda_{i}}P_{i}\bigl(\log(1/x)\bigr),
  \qquad
  0<\lambda_{0}<\lambda_{1}<\cdots\to+\infty,
  \quad P_{i}\in\RealNumbers[\,\cdot\,],
\end{equation}
where \(\lambda_{0}\) is the hyperbolicity ratio of the saddle.  The logarithms
are forced precisely by resonances among the exponents: for a nonresonant
saddle the \(P_{i}\) are constants.  Composing finitely many such maps with
regular transition maps produces the first-return map of a hyperbolic
polycycle, whose asymptotic expansion is again of the form
\eqref{plt:eq:mot-dulac-series}.  The finiteness of the number of limit cycles
accumulating on such a polycycle --- Dulac's problem, settled independently by
Écalle \cite{ecalle} and by Ilyashenko --- turns on how much of this formal
expansion can be promoted to an analytic statement.  Normal forms for the
associated germs are studied in \cite{mardesic-normalforms,resman-dulac}, the
\(\varepsilon\)-neighbourhood length of an orbit of a Dulac map is computed in
the same scale \cite{mardesic-length}, and closure properties of the ambient
logarithmic transseries algebras are in \cite{peran-logtransseries}.

The change of variable \(X=1/x\), \(L=\log X=\log(1/x)\) carries
\eqref{plt:eq:mot-dulac-series} to \(\sum_{i}t^{\lambda_{i}}P_{i}(L)\), which
is the normal form \eqref{plt:eq:mot-normal-form} with a real exponent set;
Part~VIII treats that Hahn enlargement.  The core of this volume keeps the
integer lattice, where the algorithms are exact and finite.

\begin{warningbox}
Part of the dynamical literature uses \(-1/\log x\) as its small logarithmic
generator, so that the generator tends to \(0\) with \(x\), rather than
\(\log(1/x)\), which tends to \(+\infty\).  The two conventions differ by an
inversion of the coefficient variable, and a formula transported between them
without applying that inversion will be wrong.  In this volume \(L=\log X\)
always tends to \(+\infty\).
\end{warningbox}

\subsection{Iteration, normal forms, and recursions}

\Cref{plt:ex:mot-sine,plt:ex:mot-parabolic} are the model cases: whenever the
increment of a slowly diverging quantity carries a term of size \(1/w\), the
partial sums manufacture a logarithm, and the subsequent expansion is organized
in inverse powers of \(n\) with polynomial coefficients in \(\log n\).  In the
analytic classification of parabolic germs the same logarithm is the
logarithmic term of the Fatou coordinate; in the formal classification of maps
tangent to the identity at infinity the conjugating series is an element of
\(\Gnear\).  Symbolic algorithms for computing such expansions --- and for the
broader exp-log asymptotic inversion problem --- are due to Salvy and Shackell
\cite{salvy-shackell-1998,salvy-shackell-1999}; Part~VII specializes their
framework to the present class.

\subsection{The Fabius--Rvachev endpoint, where this corpus meets the scale}

The corpus in which this volume sits studies the Fabius function
\(\FabiusBounded\) and the Rvachev function \(\RvachevUp\).  Their contact with
the polynomial-logarithmic scale is at the endpoints, and it is structural
rather than decorative: the endpoint problem is an inversion problem of exactly
the type analysed above, and its solution needs two generators.

The corpus proves, and has verified in Lean, the effective flatness family
\begin{gather}
  \FabiusBounded(x)\le2^{\binom{n+1}{2}}x^{n},
  \qquad
  \FabiusBounded(x)\le\frac{2^{\binom{n+1}{2}}}{n!}\,x^{n},
  \label{plt:eq:mot-flatness}\\
  \text{valid for all }n\in\NaturalNumbers,\ x\ge0
  \text{ with }2^{n}x\le1 .
  \notag
\end{gather}
Optimizing the free index \(n\) in the sharper of the two turns the family into
a quadratic-in-\(\log\) decay law.

\begin{proposition}[A two-term endpoint upper bound from the flatness family]
\label{plt:prop:mot-fabius-endpoint}
Write \(\ell=\log(1/x)\) and \(a=\log2\), and set \(C=\tfrac12+3/a\).  For all
\(0<x\le\tfrac14\),
\begin{equation}\label{plt:eq:mot-fabius-endpoint}
  \log\FabiusBounded(x)
  \le-\frac{\ell^{2}}{2a}-\frac{\ell\log\ell}{a}+C\ell .
\end{equation}
\end{proposition}

\begin{proof}
Put \(m=\ell/a\); since \(x\le\tfrac14\) we have \(\ell\ge2a\), hence
\(m\ge2\).  Let \(n=\Floor{m}\ge2\) and \(\theta=m-n\in[0,1)\).  The hypothesis
\(2^{n}x\le1\) is equivalent to \(na\le\ell\), i.e.\ to \(n\le m\), which
holds.  Taking logarithms in the second bound of \eqref{plt:eq:mot-flatness}
and using \(\log x=-\ell\),
% ed.: the draft named the two parts of this estimate A and B, colliding with
% the B = log(1/y) of the corollary and the status remark that quote this
% proposition.  The internal quantities are renamed A_1 and A_2; B is reserved
% for the inverse-endpoint variable throughout the subsection.
\[
  \log\FabiusBounded(x)\le A_{1}+A_{2},
  \qquad
  A_{1}\DefinitionEquals\frac{an(n+1)}{2}-n\ell,
  \qquad
  A_{2}\DefinitionEquals-\log n! .
\]

\emph{The quadratic part.}  Using \(\ell=am\) and \(n=m-\theta\),
\[
  A_{1}=\frac{an}{2}\bigl(n+1-2m\bigr)
   =-\frac{a}{2}(m-\theta)(m+\theta-1)
   =-\frac{a}{2}\bigl(m^{2}-m+\theta-\theta^{2}\bigr).
\]
Since \(\theta-\theta^{2}\ge0\) on \([0,1)\), and since
\(am^{2}/2=\ell^{2}/(2a)\) and \(am/2=\ell/2\),
\[
  A_{1}\le-\frac{\ell^{2}}{2a}+\frac{\ell}{2}.
\]

\emph{The factorial part.}  From \(\EulerE^{n}\ge n^{n}/n!\) we get
\(n!\ge n^{n}\EulerE^{-n}\), hence \(A_{2}\le-n\log n+n\).  As \(m\ge2\) and
\(0\le\theta<1\), we have \(0\le\theta/m\le\tfrac12\), so
\(\log(1-\theta/m)\ge-2\theta/m\ge-2/m\) and therefore
\[
  n\log n=(m-\theta)\bigl(\log m+\log(1-\theta/m)\bigr)
   \ge(m-\theta)\log m-2 .
\]
Using \(0\le\theta<1\), \(\log m\ge\log2>0\) and \(n\le m\), this gives
\(A_{2}\le-m\log m+\log m+2+m\).  Now
\[
  m\log m=\frac{\ell}{a}\bigl(\log\ell-\log a\bigr)
   =\frac{\ell\log\ell}{a}-\frac{\ell\log a}{a},
\]
and \(a=\log2<1\) gives \(\log a<0\), so
\(-m\log m=-\ell\log\ell/a+\ell\log a/a\le-\ell\log\ell/a\).  Finally
\(\log m\le m=\ell/a\) and \(2\le m=\ell/a\), so
\[
  A_{2}\le-\frac{\ell\log\ell}{a}+\frac{3\ell}{a}.
\]
Adding the two estimates gives \eqref{plt:eq:mot-fabius-endpoint} with
\(C=\tfrac12+3/a\).
\end{proof}

\begin{corollary}[The inverse endpoint lives on a Puiseux-logarithmic scale]
\label{plt:cor:mot-fabius-inverse}
Let \(\FabiusQuantile\) be the inverse of the strictly increasing restriction
\(\FabiusBounded|_{[0,1]}\), and put \(B=\log(1/y)\).  Then, as
\(y\downarrow0\),
\begin{equation}\label{plt:eq:mot-fabius-inverse}
  \log\frac{1}{\FabiusQuantile(y)}
  \le\sqrt{2B\log2}+\BigOAt{1}{y\downarrow0}.
\end{equation}
\end{corollary}

\begin{proof}
Let \(x=\FabiusQuantile(y)\), so that \(y=\FabiusBounded(x)\) and
\(x\downarrow0\) as \(y\downarrow0\); for \(y\) small, \(x\le\tfrac14\).  Put
\(\ell=\log(1/x)\) and \(a=\log2\).  Since \(\ell\ge2a>1\) we have
\(\log\ell>0\), so \cref{plt:prop:mot-fabius-endpoint} gives
\[
  B=\log\frac{1}{y}=-\log\FabiusBounded(x)
   \ge\frac{\ell^{2}}{2a}+\frac{\ell\log\ell}{a}-C\ell
   \ge\frac{\ell^{2}}{2a}-C\ell .
\]
Hence \(\ell^{2}-2aC\ell-2aB\le0\), so
\(\ell\le aC+\sqrt{a^{2}C^{2}+2aB}\le2aC+\sqrt{2aB}\), using
\(\sqrt{p+q}\le\sqrt p+\sqrt q\).  As \(\ell=\log(1/\FabiusQuantile(y))\), this
is \eqref{plt:eq:mot-fabius-inverse}.
\end{proof}

\begin{remark}[What the corpus adds, and what must not be inferred]
\label{plt:rmk:mot-fabius-status}
\Cref{plt:prop:mot-fabius-endpoint,plt:cor:mot-fabius-inverse} are one-sided:
they rest on a family of upper bounds, and an upper bound for
\(\FabiusBounded\) yields only an upper bound for
\(\log(1/\FabiusQuantile(y))\).  The corpus obtains the matching lower bounds
by a lattice computation along the dyadic scale ladder, giving the two-sided
laws
\[
  \log\FabiusBounded(x)=-\frac{\ell^{2}}{2\log2}
   -\frac{\ell\log\ell}{\log2}+\BigOAt{\ell}{x\downarrow0},
\]
\[
  \log\frac{1}{\FabiusQuantile(y)}=\sqrt{2B\log2}-\tfrac12\log B
   +\BigOAt{1}{y\downarrow0},
\]
with \(\ell=\log(1/x)\) and \(B=\log(1/y)\).  Those results are quoted here,
not proved; their proofs belong to the endpoint volumes of the corpus.

Two structural points connect them to the present volume.  First, the inverse
law is not a polynomial-logarithmic series in the core sense.  Read with
\(X=B\), \(t=X^{-1}\) and \(L=\log X\), its displayed terms are
\(\sqrt{2\log2}\,X^{1/2}-\tfrac12L+\BigOAt{1}{X\to+\infty}\), an element of the
Puiseux enlargement \(\K[L]\bigl((t^{1/2})\bigr)\) of Part~VIII, not of
\(\PLlau\).  The half-integer exponent is the signature of inverting a leading
behaviour that is quadratic rather than linear, a situation treated in Part~IV
under reversion beyond the near-identity class.  Second, the corpus records
that the refined dyadic endpoint expansions carry \emph{log-periodic}
corrections --- bounded one-periodic functions of \(\ell/\log2\) --- and that
the natural coordinate for the associated saddle equations is the lower real
branch \(\LambertW_{-1}\).  A bounded oscillating function of \(\ell/\log2\) is
not a monomial of any power-logarithmic scale: by
\cref{plt:lem:mot-dominance} it is not asymptotic to a constant multiple of any
\(t^{n}L^{j}\), so it cannot appear as a coefficient in
\eqref{plt:eq:mot-normal-form}.  Its presence is a genuine departure from the
class studied here, and the algebra of this volume applies to such an expansion
only after the periodic factor has been separated out.  The appearance of
\(\LambertW_{-1}\), by contrast, is squarely inside the volume: that branch is
expanded in Part~V.
\end{remark}

\section{Formal, asymptotic, convergent}
\label{plt:sec:mot-three-statuses}

The material above has used three logically independent kinds of statement, and
the rest of the volume depends on never confusing them.

\begin{warningbox}
A displayed infinite expansion can have any of three statuses.
\begin{enumerate}
\item A \emph{formal identity}: an equality of coefficient arrays in a series
  ring.  Written with the formal tail \(\FormalBigOAt{t^{N}}{t}\), which means
  \(\ordt(F-G)\ge N\) and asserts no numerical bound, no convergence, and no
  uniformity in \(L\).
\item An \emph{asymptotic expansion}: a hierarchy of remainder estimates for an
  actual function on a stated domain.  Written \(\BigOAt{A(X)}{X\to+\infty}\),
  or \(\BigOAt{A(X)}{X\to\infty,\ X\in S}\) in a sector, with the limit and the
  domain printed.
\item A \emph{convergent expansion}: an actual infinite sum, with its domain of
  convergence stated.
\end{enumerate}
A fourth use of the same glyph, algorithmic complexity \(\BigO(N^{2})\), is
none of the three.  No one of these four meanings may be inferred from another;
in particular a bare \(\BigO\) is never acceptable for an asymptotic claim in
this volume.
\end{warningbox}

For the Lambert expansion all three statuses are in fact available on suitable
domains, which is exceptional rather than typical: the large-argument Lambert
series is not only asymptotic but convergent in appropriate real and complex
regions, a fact requiring separate analytic work
\cite{corless-et-al-1996,dlmf-W}.
% ed.: source draft 1 asserts convergence "on the principal real branch for
% z >= e", supported only by a general citation.  That threshold is not
% established anywhere in the six drafts, and the assertion also conflates a
% convergence claim with the separate, elementary observation that the series
% is exact at z = e.  We keep the coefficient statement, proved below, and
% defer the convergence domain to Part VI with its own citation.
% ed.: the draft called the asymptotic display \eqref{plt:eq:mot-W-two-orders}
% "exact at z = e".  A display whose remainder is an analytic O as z -> +infty
% asserts nothing at a single point, so that formulation commits precisely the
% conflation the box above forbids.  The correct elementary statement is a
% coefficient identity for the underlying formal series, and it is proved here.
What is elementary, and worth isolating, is a statement about the
\emph{coefficients}: the formal series underlying
\eqref{plt:eq:mot-W-two-orders}, namely
\(L_{1}+\sum_{n\ge0}\LamP{n}(L_{2})L_{1}^{-n}\), collapses to its leading term
at \(z=\EulerE\), where \(L_{1}=1\) and \(L_{2}=0\).  Indeed
\(\LamP{n}(0)=0\) for every \(n\in\NaturalNumbers\): in the defining formula
\(\LamP{n}(u)=\bigl((-1)^{n+1}/(n+1)!\bigr)
 \bigl[\prod_{j=0}^{n-1}(\partial_{u}-j)\bigr]u^{n+1}\)
each of the \(n\) factors either differentiates or rescales, so every monomial
of the result has degree at least \((n+1)-n=1\) and \(\LamP{n}\) has no
constant term.  Hence every correction block vanishes at \(z=\EulerE\) and the
series returns the correct value \(\LambertW_{0}(\EulerE)=1\).  This is an
identity between coefficients, not a convergence statement: a convergence
threshold is a different assertion and is not proved in this volume; Part~VI
states what is proved and cites what is not.

\begin{summarybox}
The inverse of \(y\mapsto y+\log y\) cannot be expanded in powers of \(X\)
alone --- the correction \(-\log X\) falls in a gap of that scale --- nor in
powers of \(\log X\) alone, which cannot even see the correction
\(X^{-1}\log X\).  It can be expanded in the two-generator scale \(t^{n}L^{j}\),
and that scale is forced: the smallest differentially closed ring containing a
logarithm is \(\K[t,L]\).  Dominance is lexicographic, one inverse power of
\(X\) defeating every fixed power of \(\log X\), so a series is organized into
inverse-power blocks, each a polynomial in \(L\).  The same scale organizes
implicit equations with a logarithmic term, Dulac and first-return maps,
parabolic normal forms, recursions with a harmonic increment, and the endpoint
asymptotics of the Fabius--Rvachev system.
\end{summarybox}

% ---- fragment 02-rings ----
\chapter{The formal polynomial--logarithmic rings}
\label{plt:sec:rings}

\section{Generators, coefficients, and what ``formal'' means here}
\label{plt:sec:ring-generators}

Throughout this volume \(\K\) is a field of characteristic zero.  The default
analytic regime is the positive ray \(X\to+\infty\); for complex work a sector
and a branch of \(\log\) must be named before any analytic statement is made,
and the algebra developed in this part is branch-neutral until such a choice is
imposed.

\begin{definition}[Generators]
\label{plt:def:ring-generators}
Set
\begin{equation}\label{plt:eq:ring-generators}
  t\DefinitionEquals X^{-1},
  \qquad
  L\DefinitionEquals\log X .
\end{equation}
In the formal algebra \(t\) and \(L\) are two \emph{independent} generators:
the underlying additive structure of every class below is a space of families
of scalars indexed by exponent pairs, and no algebraic relation between \(t\)
and \(L\) is imposed.
\end{definition}

Independence is a statement about the ring operations only, and it is exactly
as strong as it needs to be.

\begin{remark}[The two-variable illusion]
\label{plt:rmk:ring-two-variable}
For addition and multiplication, treating \(t\) and \(L\) as unrelated
indeterminates is harmless and is what makes the coefficient bookkeeping
finite.  For differentiation and for composition it is false.  The analytic
relation \(L=\log X\), \(t=X^{-1}\) re-enters through the distinguished
derivation
\[
  \DX=t\,\partial_L-t^2\,\partial_t,
  \qquad
  \DX t=-t^2,\qquad \DX L=t,
\]
studied in Part~II, and through the substitution rule
\(t\mapsto G^{-1}\), \(L\mapsto\log G\) of Part~III.  Every ``shifted''
operator \(\partial_L-n\) that appears later is a trace of this coupling.  A
calculation that silently uses \(\partial_L L=1\) together with
\(\partial_t L=0\) inside a differentiation is using the coupling, and must
say so.
\end{remark}

\begin{warningbox}[title={Formal is not analytic}]
Every statement in this part is an identity between formal objects: two
elements are equal when all their coefficients agree.  No such statement
asserts that substituting a numerical value of \(X\) makes an infinite sum
converge, nor that any function has the displayed expansion.  The passage from
a formal identity to an asymptotic or convergent statement is a separate
theorem with separately printed hypotheses.
\end{warningbox}

\section{The four ambient classes}
\label{plt:sec:ring-four-classes}

\begin{definition}[The four classes]
\label{plt:def:ring-four-classes}
Define
\begin{align}
  \PLpow&\DefinitionEquals\K[L][[t]]
    =\SetOf{\textstyle\sum_{n\ge0}p_n(L)t^n:p_n\in\K[L]},
    \label{plt:eq:ring-PLpow}\\
  \PLlau&\DefinitionEquals\K[L]((t))
    =\SetOf{\textstyle\sum_{n\ge n_0}p_n(L)t^n:
      n_0\in\IntegerNumbers,\ p_n\in\K[L]},
    \label{plt:eq:ring-PLlau}\\
  \PLrat&\DefinitionEquals\K(L)((t)),
    \label{plt:eq:ring-PLrat}\\
  \PLit&\DefinitionEquals\K((L^{-1}))((t)).
    \label{plt:eq:ring-PLit}
\end{align}
In each case the outer construction is a formal Laurent (or power) series in
\(t\); the coefficient object at level \(n\) is called the \emph{block} at
\(t^n\).  Addition is blockwise and multiplication is the Cauchy convolution
\[
  \Bigl(\sum_i p_it^i\Bigr)\Bigl(\sum_k q_kt^k\Bigr)
  =\sum_n\Bigl(\sum_{i+k=n}p_iq_k\Bigr)t^n ,
\]
in which every inner sum is finite because the outer supports are bounded
below.
\end{definition}

Two facts are definitional but deserve to be said once, because almost every
later argument uses them.  First, each block is a \emph{single} finite
polynomial (resp.\ rational function, resp.\ Laurent series in \(L^{-1}\)); no
uniform bound on its degree is implied.  Second, only finitely many negative
powers of \(t\), that is finitely many positive powers of \(X\), may occur.

\begin{proposition}[Localization]
\label{plt:prop:ring-localization}
\(\PLlau=\PLpow[t^{-1}]\): the natural map from the localization of \(\PLpow\)
at the multiplicative set \(\SetOf{t^m:m\ge0}\) to \(\PLlau\) is a ring
isomorphism.  The same holds with \(\K[L]\) replaced by \(\K(L)\) or
\(\K((L^{-1}))\).
\end{proposition}

\begin{proof}
The map sends \(F/t^m\) to \(t^{-m}F\); it is well defined and a ring
homomorphism.  It is surjective because an element of \(\PLlau\) with
\(\ordt\ge -m\) is \(t^{-m}\) times an element of \(\PLpow\).  It is injective
because \(t^{-m}\) is invertible in \(\PLlau\), so \(t^{-m}F=0\) forces
\(F=0\), whence \(F/t^m=0\) in the localization.
\end{proof}

\begin{proposition}[Strict inclusions]
\label{plt:prop:ring-chain}
There are inclusions of \(\K\)-algebras
\begin{equation}\label{plt:eq:ring-chain}
  \PLpow\subsetneq\PLlau\subsetneq\PLrat\subsetneq\PLit,
\end{equation}
each given by the identity on coefficients except the last, which is the
coefficientwise Laurent expansion at \(L=\infty\).
% ed.: the chain \eqref{plt:eq:ring-chain} exhibits three inclusions, not four;
% the count was wrong as written.
All three inclusions are
strict, and:
\begin{enumerate}[label=\textup{(\roman*)}]
\item \(\PLpow\) and \(\PLlau\) are integral domains and are not fields;
\item \(\PLrat\) and \(\PLit\) are fields.
\end{enumerate}
\end{proposition}

\begin{proof}
% ed.: the first inclusion is not a change of coefficient ring (the printed
% chain \(\K[L]\subset\K[L]\subset\K(L)\) made it look like one); it is the
% inclusion of a power-series ring in the corresponding Laurent ring.
\emph{The first two inclusions.}  The first is the inclusion
\(E[[t]]\subset E((t))\) for \(E=\K[L]\), a subring inclusion by
\cref{plt:def:ring-four-classes}; the second is the blockwise extension of the
coefficient inclusion \(\K[L]\subset\K(L)\), hence an injective ring
homomorphism.  Strictness: \(t^{-1}\in\PLlau\setminus\PLpow\), and
\(L^{-1}\in\PLrat\setminus\PLlau\) because \(L^{-1}\notin\K[L]\).

\emph{The third inclusion.}  The polynomial ring \(\K[L]\) sits inside
\(\K((L^{-1}))\) (a polynomial has finitely many positive and no negative
powers of \(L\)).  Every nonzero \(p=aL^d+\cdots\in\K[L]\), \(a\in\K^\times\),
is a unit there: writing \(u=L^{-1}\) we have \(p=au^{-d}(1+w)\) with
% ed.: this appeal was unsupported and mildly circular: the geometric-series
% corollary was stated only for coefficient rings whose list already contained
% \(\K((L^{-1}))\), the very object being constructed here.
% \cref{plt:cor:ring-geometric} is now proved for an arbitrary commutative
% coefficient ring, so the instance \(E=\K\), generator \(u\), is legitimate.
\(w\in u\K[u]\), and \(1+w\) is invertible in \(\K[[u]]\) by the geometric
series \(\sum_{r\ge0}(-w)^r\), which is \(u\)-adically summable
(\cref{plt:cor:ring-geometric} with coefficient ring \(\K\) and \(u\) in place
of \(t\)).  By the
universal property of the field of fractions, the inclusion
\(\K[L]\hookrightarrow\K((L^{-1}))\) extends uniquely to a field embedding
\(\K(L)\hookrightarrow\K((L^{-1}))\).  Extending blockwise gives
\(\PLrat\hookrightarrow\PLit\).  Strictness is
\cref{plt:prop:ring-PLit-strict}.

\emph{(i).}  A Laurent series ring \(E((t))\) over an integral domain \(E\) is
an integral domain: if \(F,G\ne0\), then by
\cref{plt:thm:ring-valuation-laws}\,(i) the block of \(FG\) at level
\(\ordt(F)+\ordt(G)\) is the product of two nonzero elements of \(E\), hence
nonzero.  Neither \(\PLpow\) nor \(\PLlau\) is a field: if \(LG=1\) with
\(G=\sum_nq_n(L)t^n\), comparison of the blocks at \(t^0\) gives
\(Lq_0(L)=1\) in \(\K[L]\), which is impossible on degree grounds.

\emph{(ii).}  Let \(E\) be a field and \(F\in E((t))\) nonzero, with
\(\nu=\ordt(F)\) and block \(c=p_\nu\in E^\times\).  Then
\(F=ct^\nu(1+U)\) with \(U=\sum_{n\ge1}(p_{\nu+n}/c)t^n\in tE[[t]]\), and
\(\sum_{r\ge0}(-U)^r\) is \(t\)-adically summable
(\cref{plt:cor:ring-geometric}), so \(F\) is invertible with
\(F^{-1}=c^{-1}t^{-\nu}\sum_{r\ge0}(-U)^r\).  Apply this with
\(E=\K(L)\) and \(E=\K((L^{-1}))\).
\end{proof}

\begin{corollary}[Domains]
\label{plt:cor:ring-domains}
All four classes are integral domains, and \(\ordt\) and the rank-two
valuation of \cref{plt:def:ring-rank-two} are defined on each of them.
\end{corollary}

\begin{proof}
Immediate from \cref{plt:prop:ring-chain}, since a subring of a field is a
domain and \(\PLpow\subset\PLlau\subset\PLit\).
\end{proof}

\begin{proposition}[The iterated field is strictly larger than the rational one]
\label{plt:prop:ring-PLit-strict}
The element
\begin{equation}\label{plt:eq:ring-lacunary}
  \Phi\DefinitionEquals\sum_{k\ge0}L^{-k^2}\in\K((L^{-1}))\subset\PLit
\end{equation}
does not lie in the image of \(\K(L)\); consequently
\(\PLrat\subsetneq\PLit\).
\end{proposition}

\begin{proof}
Put \(u=L^{-1}\), so that \(\K((L^{-1}))=\K((u))\), the image of \(\K(L)\) is
the subfield \(\K(u)\), and \(\Phi=\sum_{k\ge0}c_ku^k\in\K[[u]]\) with
\(c_k=1\) if \(k\) is a perfect square and \(c_k=0\) otherwise.  Suppose
\(\Phi\in\K(u)\).  Write \(\Phi=A/B\) in lowest terms with
\(A,B\in\K[u]\).  If \(B(0)=0\), then \(u\mid B\); since
\(\Phi\in\K[[u]]\) we get \(A=B\Phi\in u\K[[u]]\cap\K[u]\), so \(u\mid A\),
contradicting coprimality.  Hence \(B(0)=b_0\ne0\).  Let
\(B=\sum_{i=0}^{m}b_iu^i\) and \(D=\deg A\).  Comparing coefficients of
\(u^k\) in \(B\Phi=A\) gives, for every \(k>D\),
\[
  b_0c_k=-\sum_{i=1}^{m}b_ic_{k-i}.
\]
Therefore, if \(k>D+m\) and \(c_{k-1}=\dots=c_{k-m}=0\), then \(c_k=0\), and
by induction \(c_l=0\) for all \(l\ge k\).  Consecutive squares differ by
\((k+1)^2-k^2=2k+1\), so the sequence \((c_k)\) has runs of zeros of every
length; choose one of length \(m\) beginning beyond \(D+m\).  Then \(c_l=0\)
for all large \(l\), contradicting \(c_{k^2}=1\) for arbitrarily large \(k\).
\end{proof}

\begin{proposition}[The coefficient extensions are not localizations]
\label{plt:prop:ring-not-localization}
Let \(S=\K[L]\setminus\SetOf{0}\).  Then the localization
\(S^{-1}\PLlau\) is a proper subring of \(\PLrat\).  Explicitly,
\begin{equation}\label{plt:eq:ring-no-common-denominator}
  \Phi_0\DefinitionEquals\sum_{n\ge0}\frac{t^n}{L+n}\in\PLrat
\end{equation}
admits no representation \(F/q\) with \(F\in\PLlau\) and
\(q\in S\).
\end{proposition}

\begin{proof}
If \(\Phi_0=F/q\) then \(q\Phi_0=F\in\PLlau\), so for every \(n\ge0\) the
block \(q/(L+n)\) lies in \(\K[L]\), that is, \(L+n\) divides \(q\).  Since
\(\operatorname{char}\K=0\), the polynomials \(L+n\), \(n\ge0\), are pairwise
non-associate irreducibles, so \(q\) would have infinitely many pairwise
non-associate irreducible factors; hence \(q=0\), a contradiction.
\end{proof}

\begin{remark}[What each step buys, and what it costs]
\label{plt:rmk:ring-tradeoffs}
\Cref{plt:tab:ring-tradeoffs} records the trade-offs.  The two extensions on
the right are not interchangeable: \(\PLrat\) keeps rational dependence on
\(L\) \emph{exact} but does not resolve the secondary asymptotic scale, while
\(\PLit\) resolves that scale but replaces an exact coefficient such as
\((L+1)^{-1}\) by the infinite tail
\(L^{-1}-L^{-2}+L^{-3}-\cdots\) and therefore needs a second, inner cutoff
before anything can be stored or printed.
\end{remark}

\begin{table}[htbp]
\centering
\small
\caption{The four ambient classes.  ``Finite record'' means the data a
computer must hold to determine the object up to a declared precision.}
\label{plt:tab:ring-tradeoffs}
\begin{tabular}{@{}lp{0.30\textwidth}p{0.34\textwidth}@{}}
\hline
Class & Buys & Costs \\
\hline
\(\PLpow=\K[L][[t]]\)
 & \(t^N\PLpow\) is a genuine ideal, so residues modulo \(t^N\) form a ring;
   complete and separated; local finiteness is automatic
 & does not contain \(X\); not stable under multiplication by \(t^{-1}\) \\
\(\PLlau=\K[L]((t))\)
 & smallest of the four containing \(X\); still exact polynomial blocks;
   finite record \(n\mapsto p_n\)
 & not a field; \(t\) is a unit, so \(t^N\PLpow\) is only a filtration tail,
   not an ideal \\
\(\PLrat=\K(L)((t))\)
 & a field; exact rational dependence on \(L\); unrestricted division of
   blocks
 & not a localization of \(\PLlau\)
   (\cref{plt:prop:ring-not-localization}); a block has no intrinsic
   asymptotic expansion until one is chosen \\
\(\PLit=\K((L^{-1}))((t))\)
 & a field; the Hahn field of the rank-two value group
   (\cref{plt:thm:ring-hahn}); the secondary scale in \(L^{-1}\) is resolved
 & blocks are infinite objects, so a finite record needs an inner
   \(L^{-1}\)-cutoff as well as the outer \(t\)-cutoff; the embedding of
   \(\PLrat\) destroys exactness \\
\hline
\end{tabular}
\end{table}

\begin{remark}[The classes are not local rings]
\label{plt:rmk:ring-not-local}
% ed.: S4 calls \(t\K[\ell][[t]]\) ``the strictly small ideal'' of the Laurent
% ring \(\K[\ell]((t))\); it is neither small in that ring nor an ideal of it.
The subset \(t\PLpow\) is a prime ideal \emph{of \(\PLpow\)}, with
\(\PLpow/t\PLpow\cong\K[L]\); since \(\K[L]\) is not a field, \(t\PLpow\) is
not maximal and \(\PLpow\) is not a local ring.  In \(\PLlau\) the set
\(t\PLpow\) is not an ideal at all
(\cref{plt:prop:trunc-not-ideal}).  By contrast \(\K(L)[[t]]\) and
\(\K((L^{-1}))[[t]]\) are discrete valuation rings with residue fields
\(\K(L)\) and \(\K((L^{-1}))\).
\end{remark}

\section{Coefficient support}
\label{plt:sec:ring-support}

Expanding every block into monomials gives the double form
\begin{equation}\label{plt:eq:ring-double-form}
  F=\sum_{n\ge n_0}p_n(L)t^n
   =\sum_{\substack{n\in\IntegerNumbers\\ j\in\NaturalNumbers}}a_{n,j}\,t^nL^j,
  \qquad
  p_n(L)=\sum_{j\ge0}a_{n,j}L^j .
\end{equation}

\begin{definition}[Coefficient support]
\label{plt:def:ring-support}
For \(F\in\PLlau\) written as in \eqref{plt:eq:ring-double-form},
\[
  \CoefficientSupportOf{F}
  \DefinitionEquals
  \SetOf{(n,j)\in\IntegerNumbers\times\NaturalNumbers:a_{n,j}\ne0}.
\]
The zero series has empty coefficient support.  This is support in an
exponent-index lattice; it is unrelated to the pointwise support of a
function.
\end{definition}

\begin{definition}[Dominance]
\label{plt:def:ring-dominance}
For \((n,j),(m,k)\in\IntegerNumbers\times\IntegerNumbers\),
\begin{equation}\label{plt:eq:ring-dominance}
  t^nL^j\succdom t^mL^k
  \quad\Longleftrightarrow\quad
  n<m,\ \text{ or }\ (n=m\ \text{and}\ j>k).
\end{equation}
We write \(\precdom\), \(\preceqdom\) and \(\asymdom\) accordingly.
\end{definition}

\begin{lemma}[Analytic justification of dominance]
\label{plt:lem:ring-dominance-analytic}
% ed.: the statement imposed the standing hypothesis \(\alpha>\alpha'\) and
% then appended the case \(\alpha=\alpha'\), which that hypothesis excludes;
% the two cases are now alternatives of one hypothesis.  The regime
% \(X>\EulerE\) is printed here rather than produced inside the proof.
Let \(\alpha,\alpha',\beta,\beta'\) be fixed reals with either
\(\alpha>\alpha'\), or \(\alpha=\alpha'\) and \(\beta>\beta'\).  Then, on the
part \(X>\EulerE\) of the positive ray, where \(\log X>1\) and the quotient
below is positive,
\[
  \frac{X^{\alpha}(\log X)^{\beta}}{X^{\alpha'}(\log X)^{\beta'}}
  \longrightarrow+\infty
  \qquad(X\to+\infty).
\]
Consequently \eqref{plt:eq:ring-dominance} orders the monomials
\(t^nL^j=X^{-n}(\log X)^j\) by decreasing eventual size.
\end{lemma}

\begin{proof}
For \(X>\EulerE\) we have \(\log X>1\), so the quotient is positive and
\[
  \log\frac{X^{\alpha}(\log X)^{\beta}}{X^{\alpha'}(\log X)^{\beta'}}
  =(\alpha-\alpha')\log X+(\beta-\beta')\log\log X .
\]
In the first case, since \(\log\log X=\LittleOAt{\log X}{X\to+\infty}\) and
\(\alpha-\alpha'>0\), the right-hand side tends to \(+\infty\).  In the second
case, \(\alpha=\alpha'\) and \(\beta>\beta'\), so the expression reduces to
\((\beta-\beta')\log\log X\to+\infty\).
% ed.: the sources state this without the hypothesis \(X>\EulerE\); for
% \(1<X<\EulerE\) one has \(0<\log X<1\) and high logarithmic powers are
% small, so the comparison is genuinely eventual.
\end{proof}

\begin{remark}[Fixed exponents only]
\label{plt:rmk:ring-fixed-exponents}
\Cref{plt:lem:ring-dominance-analytic} is the two-generator fragment of
Hardy's comparison of orders of infinity \cite{hardy-orders}: it compares two
monomials with \emph{fixed} exponents.  It says nothing about a family in which the
logarithmic exponent grows with the inverse-power level; see
\cref{plt:rmk:trunc-no-size} for what this costs.
\end{remark}

\begin{theorem}[Characterization of admissible supports]
\label{plt:thm:ring-support}
For a subset \(S\subseteq\IntegerNumbers\times\NaturalNumbers\) the following
are equivalent.
\begin{enumerate}[label=\textup{(\alph*)}]
\item \(S=\CoefficientSupportOf{F}\) for some \(F\in\PLlau\).
\item The projection \(\SetOf{n:(n,j)\in S\ \text{for some }j}\) is bounded
  below, and for each \(n\) the fiber \(S_n=\SetOf{j:(n,j)\in S}\) is finite.
\item Every nonempty subset of \(S\) has a \(\succdom\)-greatest element
  \textup{(}reverse well-ordering\textup{)}.
\item \(S\) can be enumerated as a strictly \(\succdom\)-decreasing sequence
  \(s_0\succdom s_1\succdom s_2\succdom\cdots\), finite or of order type
  \(\omega\), exhausting \(S\).
\end{enumerate}
\end{theorem}

\begin{proof}
\emph{(a)\(\Leftrightarrow\)(b).}  By
\cref{plt:def:ring-four-classes} an element of \(\PLlau\) is exactly a family
\((a_{n,j})\) whose blocks \(p_n=\sum_ja_{n,j}L^j\) are polynomials --- that is,
have finite support in \(j\) --- and which vanishes for \(n<n_0\).  Conversely
any \(S\) with property (b) is the support of
\(\sum_{(n,j)\in S}t^nL^j\).

\emph{(b)\(\Rightarrow\)(d).}  If \(S=\emptyset\) there is nothing to prove.
Otherwise let \(n_0=\min\SetOf{n:S_n\ne\emptyset}\), which exists because the
projection is a nonempty subset of \(\IntegerNumbers\) bounded below.  List
first the finitely many elements of \(\SetOf{n_0}\times S_{n_0}\) in order of
decreasing \(j\), then those of \(\SetOf{n_0+1}\times S_{n_0+1}\), and so on.
By \eqref{plt:eq:ring-dominance} the resulting list is strictly
\(\succdom\)-decreasing, every element of \(S\) occupies a finite position, and
every element occurs.

\emph{(d)\(\Rightarrow\)(c).}  Given \(\emptyset\ne T\subseteq S\), let
\(k=\min\SetOf{i:s_i\in T}\); then \(s_k\succdom s_i\) for every other
\(s_i\in T\), so \(s_k\) is the greatest element of \(T\).

\emph{(c)\(\Rightarrow\)(b).}  If the projection of \(S\) were unbounded below,
then \(S\) itself would have no greatest element, since any \((n,j)\in S\) is
dominated by some \((n',j')\in S\) with \(n'<n\).  If some fiber \(S_n\) were
infinite, then \(\SetOf{n}\times S_n\) would be a nonempty subset with no
greatest element, because an infinite subset of \(\NaturalNumbers\) is
unbounded above.
\end{proof}

\begin{remark}[Order type]
\label{plt:rmk:ring-order-type}
Condition (d) is the sharp form of the ``well-based support'' hypothesis of
Hahn's construction \cite{hahn,edgar-beginners} in the present rank-two,
integer-lattice case: the support of an element of \(\PLlau\) is not merely
reverse well-ordered, it has order type at most \(\omega\).  The same is
\emph{not} true in \(\PLit\), where fibers may be infinite downwards; see
\cref{plt:thm:ring-hahn}.
\end{remark}

\begin{example}[A support lattice]
\label{plt:ex:ring-support-lattice}
The table below marks with \(\bullet\) the pairs \((n,j)\) of a support whose
log-degree profile is \(d_n=1,2,3,3,4,5,5\) for \(n=0,\dots,6\).  Each column
is finite; the columns start at \(n=0\); the greatest monomial of any
nonempty sub-diagram is found by moving as far left as possible and then as
far up as possible.
\begin{center}
\begin{tabular}{r|ccccccc}
\(j=5\) & \(\cdot\) & \(\cdot\) & \(\cdot\) & \(\cdot\) & \(\cdot\)
        & \(\bullet\) & \(\bullet\)\\
\(j=4\) & \(\cdot\) & \(\cdot\) & \(\cdot\) & \(\cdot\) & \(\bullet\)
        & \(\bullet\) & \(\bullet\)\\
\(j=3\) & \(\cdot\) & \(\cdot\) & \(\bullet\) & \(\bullet\) & \(\bullet\)
        & \(\bullet\) & \(\bullet\)\\
\(j=2\) & \(\cdot\) & \(\bullet\) & \(\bullet\) & \(\bullet\) & \(\bullet\)
        & \(\bullet\) & \(\bullet\)\\
\(j=1\) & \(\bullet\) & \(\bullet\) & \(\bullet\) & \(\bullet\) & \(\bullet\)
        & \(\bullet\) & \(\bullet\)\\
\(j=0\) & \(\bullet\) & \(\bullet\) & \(\bullet\) & \(\bullet\) & \(\bullet\)
        & \(\bullet\) & \(\bullet\)\\
\hline
        & \(n=0\) & \(1\) & \(2\) & \(3\) & \(4\) & \(5\) & \(6\)
\end{tabular}
\end{center}
\end{example}

\section{The Hahn-series viewpoint}
\label{plt:sec:ring-hahn}

Let \(\Gamma=\IntegerNumbers\times\IntegerNumbers\) be the additive group of
exponent pairs of the monomials \(t^nL^j\), \(n,j\in\IntegerNumbers\), ordered
by the rank-two valuation
\begin{equation}\label{plt:eq:ring-value-group}
  v(t^nL^j)=(n,-j),
\end{equation}
with \(\Gamma\) lexicographically ordered.  Thus a smaller value means a more
dominant monomial, and \eqref{plt:eq:ring-value-group} is exactly
\cref{plt:def:ring-dominance} rewritten.  The general framework --- an ordered
monomial group, a coefficient field, and a support condition making
convolution locally finite --- is Hahn's \cite{hahn}; its transseries
specialization is developed in
\cite{edgar-beginners,vanderhoeven-book,adh-book}.  The present rank-two case
is the smallest nontrivial instance, and the whole of
\cref{plt:sec:ring-support} is the statement that its support condition can be
checked coordinatewise.

\begin{lemma}[Rank two]
\label{plt:lem:ring-rank-two-group}
The convex subgroups of the ordered abelian group \(\Gamma\) are precisely
\(0\subset\SetOf{0}\times\IntegerNumbers\subset\Gamma\).  Hence \(\Gamma\) has
rank two, and \(v\) is a valuation of rank two.
\end{lemma}

\begin{proof}
\(\SetOf{0}\times\IntegerNumbers\) is convex: if \((0,b)<(a,c)<(0,b')\) then
\(a=0\) by the lexicographic rule.  Let \(H\) be a convex subgroup containing
some \((a,b)\) with \(a\ne0\); replacing \((a,b)\) by its negative we may
assume \(a>0\).  For any \((c,d)\in\Gamma\) choose \(N\) with \(Na>|c|\);
then \(-N(a,b)<(c,d)<N(a,b)\) and both bounds lie in \(H\), so \((c,d)\in H\)
by convexity.  Thus \(H=\Gamma\).  A convex subgroup contained in
\(\SetOf{0}\times\IntegerNumbers\) and containing some \((0,b)\ne0\) is all of
\(\SetOf{0}\times\IntegerNumbers\) by the same argument applied inside that
subgroup.
\end{proof}

\begin{theorem}[The iterated field is the Hahn field of \(\Gamma\)]
\label{plt:thm:ring-hahn}
Let \(\K[[t^{\IntegerNumbers}L^{\IntegerNumbers}]]\) denote the Hahn field of
formal sums \(\sum_{(n,j)}a_{n,j}t^nL^j\) with coefficients in \(\K\) whose
support is reverse well ordered for \(\succdom\), with the convolution
product.  Then
\[
  \K[[t^{\IntegerNumbers}L^{\IntegerNumbers}]]=\PLit
\]
as \(\K\)-algebras: a family \((a_{n,j})_{(n,j)\in\Gamma}\) has reverse
well-ordered support if and only if its \(n\)-projection is bounded below and
every fiber \(S_n\) is bounded above in \(j\); such families are exactly the
elements of \(\K((L^{-1}))((t))\), and the two products agree.
\end{theorem}

\begin{proof}
\emph{Support condition.}  Suppose the support \(S\) is reverse well ordered.
If its \(n\)-projection were unbounded below, \(S\) would have no
\(\succdom\)-greatest element, as in the proof of
\cref{plt:thm:ring-support}.  If some fiber \(S_n\) were unbounded above in
\(j\), then \(\SetOf{n}\times S_n\) would have no greatest element.
Conversely, assume the two conditions and let \(\emptyset\ne T\subseteq S\).
The \(n\)-projection of \(T\) is a nonempty subset of \(\IntegerNumbers\)
bounded below, so has a least element \(n_1\); the fiber \(T_{n_1}\) is a
nonempty subset of \(\IntegerNumbers\) bounded above, so has a greatest
element \(j_1\); and \((n_1,j_1)\) is \(\succdom\)-greatest in \(T\).

\emph{Identification of the underlying sets.}  A family whose \(n\)-projection
is bounded below by \(n_0\) and whose fibers are bounded above is precisely a
family \((p_n)_{n\ge n_0}\) with \(p_n=\sum_ja_{n,j}L^j\in\K((L^{-1}))\), that
is, an element of \(\K((L^{-1}))((t))\); and conversely.

\emph{Products.}  In \(\PLit\) the coefficient of \(t^n\) in a product is the
finite sum \(\sum_{i+k=n}p_iq_k\) (finite because both outer supports are
bounded below), and the coefficient of \(L^j\) in \(p_iq_k\) is the finite sum
\(\sum_{j_1+j_2=j}a_{i,j_1}b_{k,j_2}\) (finite because both inner supports are
bounded above).  This is exactly the Hahn convolution
\(\sum_{(n_1,j_1)+(n_2,j_2)=(n,j)}a_{n_1,j_1}b_{n_2,j_2}\), whose local
finiteness has just been verified.
\end{proof}

\begin{corollary}[Where the smaller classes sit]
\label{plt:cor:ring-hahn-subrings}
Under the identification of \cref{plt:thm:ring-hahn}:
\(\PLlau\) is the subring of elements with \emph{finite} fibers contained in
\(\NaturalNumbers\); \(\PLpow\) is the further subring with
\(n\)-projection contained in \(\NaturalNumbers\); and \(\PLrat\) is the set of
elements of \(\PLit\) \emph{every} block of which lies in the subfield
\(\K(L)\subset\K((L^{-1}))\).
\end{corollary}

\begin{proof}
The first two statements are \cref{plt:thm:ring-support}\,(b) together with
\cref{plt:def:ring-four-classes}.  The third is the definition of
\(\K(L)((t))\) transported along the blockwise embedding of
\cref{plt:prop:ring-chain}.
\end{proof}

\begin{remark}
\label{plt:rmk:ring-not-generated}
The third description cannot be improved to ``the subring generated by
\(\PLlau\) and the inverses of the nonzero elements of \(\K[L]\)'': every
element of that subring is a finite sum of finite products and therefore has
a \emph{common} denominator across all blocks, whereas
\(\Phi_0\) of \eqref{plt:eq:ring-no-common-denominator} has none, by
\cref{plt:prop:ring-not-localization}.
\end{remark}

\section{Why the iteration order is part of the type}
\label{plt:sec:ring-iteration-order}

Both \(\K((L^{-1}))((t))\) and \(\K((t))((L^{-1}))\) are fields of formal sums
of the monomials \(t^nL^j\), and both are naturally sets of families
\((a_{n,j})\in\K^{\Gamma}\).  They impose \emph{different} support conditions,
and neither is contained in the other.

\begin{proposition}[Incomparability]
\label{plt:prop:ring-iteration-order}
Regard both fields as sets of families \((a_{n,j})_{(n,j)\in\Gamma}\).  Then
\begin{align}
  \PLit&=\SetOf{(a_{n,j}):\ n\ \text{bounded below};\ \text{for each }n,\
    j\ \text{bounded above}},\label{plt:eq:ring-support-PLit}\\
  \K((t))((L^{-1}))&=\SetOf{(a_{n,j}):\ j\ \text{bounded above};\
    \text{for each }j,\ n\ \text{bounded below}}.
    \label{plt:eq:ring-support-reversed}
\end{align}
Neither set contains the other.  Explicit witnesses:
\begin{equation}\label{plt:eq:ring-order-witnesses}
  A=\sum_{n\ge0}t^nL^n\in\PLpow\subset\PLit,
  \qquad
  B=\sum_{r\ge0}t^{-r}L^{-r}\in\K((t))((L^{-1})),
\end{equation}
with \(A\notin\K((t))((L^{-1}))\) and \(B\notin\PLit\).
\end{proposition}

\begin{proof}
The two displayed descriptions are the definitions of an outer Laurent series
whose coefficients are inner Laurent series, read in the two possible orders.
For \(A\): its support is \(\SetOf{(n,n):n\ge0}\), whose \(j\)-coordinate is
unbounded above, so \(A\) violates the first condition in
\eqref{plt:eq:ring-support-reversed}; but its \(n\)-coordinate is bounded
below and each fiber is a single point, so \(A\in\PLit\) --- indeed
\(A\in\PLpow\), since its block at \(t^n\) is the polynomial \(L^n\).
For \(B\): its support is \(\SetOf{(-r,-r):r\ge0}\), whose \(n\)-coordinate is
unbounded below, so \(B\notin\PLit\); but its \(j\)-coordinate is bounded above
by \(0\) and each of its fibers over a fixed \(j=-r\) is the single point
\(n=-r\), so \(B\) satisfies \eqref{plt:eq:ring-support-reversed}.
\end{proof}

\begin{remark}[The common part, and why we choose \(t\) outermost]
\label{plt:rmk:ring-common-part}
Intersecting \eqref{plt:eq:ring-support-PLit} and
\eqref{plt:eq:ring-support-reversed} gives the families supported in a
quadrant \(\SetOf{n\ge n_0}\times\SetOf{j\le j_0}\).  These form a subring of
both fields, but not a field: \(1+tL\) lies in it and its inverse
\(\sum_{r\ge0}(-1)^rt^rL^r\) does not.  We take \(t\) outermost because the
intended hierarchy is \(t\precdom L^{-m}\precdom1\precdom L^m\precdom t^{-1}\)
for every fixed \(m>0\); this is the formal counterpart of
\(X^{-1}=\LittleOAt{(\log X)^{-m}}{X\to+\infty}\), and it is exactly the
statement that the convex subgroup \(\SetOf{0}\times\IntegerNumbers\) of
\cref{plt:lem:ring-rank-two-group} is the \emph{smaller} one.  The same
phenomenon --- that an iterated series field remembers the order in which its
scales were adjoined --- is standard in the general theory
\cite{edgar-composition,vanderhoeven-book}; what
\cref{plt:prop:ring-iteration-order} adds is that here the two orders are
\emph{incomparable}, not merely distinct.
\end{remark}

\section{The Puiseux-log extension}
\label{plt:sec:ring-puiseux}

\begin{definition}[Puiseux-log ring]
\label{plt:def:ring-puiseux}
For \(s\in\PositiveIntegers\) let
\begin{equation}\label{plt:eq:ring-puiseux}
  \K[L]\bigl((t^{1/s})\bigr)
  =\SetOf{\textstyle\sum_{r\ge r_0}p_{r}(L)\,t^{r/s}:
    r_0\in\IntegerNumbers,\ p_r\in\K[L]},
\end{equation}
the polynomial-logarithmic Laurent ring in the generator \(t^{1/s}\).  Its
monomials are \(t^{\alpha}L^j\) with \(\alpha\in s^{-1}\IntegerNumbers\) and
\(j\in\NaturalNumbers\).
\end{definition}

\begin{proposition}[Ramified tower]
\label{plt:prop:ring-puiseux-tower}
For \(s\mid s'\) the substitution \(t^{1/s}\mapsto(t^{1/s'})^{s'/s}\) is an
injective ring homomorphism
\(\K[L]((t^{1/s}))\hookrightarrow\K[L]((t^{1/s'}))\), and these maps are
compatible, so
\[
  \K[L]\bigl((t^{1/\infty})\bigr)
  \DefinitionEquals
  \bigcup_{s\in\PositiveIntegers}\K[L]\bigl((t^{1/s})\bigr)
\]
is a ring, containing \(\PLlau=\K[L]((t^{1/1}))\).  Every element of the union
has all its exponents in a single \(s^{-1}\IntegerNumbers\); consequently the
union is \emph{strictly} smaller than the Hahn ring of well-ordered supports
in \(\RationalNumbers\times\NaturalNumbers\).
\end{proposition}

\begin{proof}
The substitution is injective because it is injective on exponents and the
identity on coefficients, and it is a ring homomorphism because it is induced
by a monoid homomorphism of monomials; compatibility is immediate.  A directed
union of rings along injective homomorphisms is a ring.  For strictness,
consider \(\sum_{k\ge1}t^{1-1/k}\), whose support
\(\SetOf{1-1/k:k\ge1}\subset\RationalNumbers\) is well ordered in the
direction of increasing smallness and has no bounded denominator; it therefore
lies in the Hahn ring but in no \(\K[L]((t^{1/s}))\).
\end{proof}

\begin{remark}[What ramification costs]
\label{plt:rmk:ring-puiseux-cost}
Formally, nothing in \cref{plt:sec:ring-support} changes: the support is a
subset of \(s^{-1}\IntegerNumbers\times\NaturalNumbers\) with finite fibers and
lower-bounded projection, and \cref{plt:thm:ring-support} applies verbatim
after rescaling.  Analytically a branch of \(X^{1/s}\) must be chosen, and
that choice is data, not a convention.  Fractional exponents arise in this
volume only where they are forced --- at a degenerate leading derivative, and
at the square-root branch point of \(\LambertW\).  The corresponding classes
near a finite point, where one writes \(\sum_{\lambda\in S}z^{\lambda}
p_\lambda(\log z)\) as \(z\to0^+\), are the power--log and Dulac series of
local dynamics
\cite{mardesic-normalforms,mardesic-length,resman-dulac,peran-logtransseries};
the change of variable relating them to the present classes is carried out in
Part~V.
\end{remark}

\chapter{Valuation, leading data, and the degree profile}
\label{plt:sec:leading-data}

Throughout this chapter \(R\) denotes any one of the four classes of
\cref{plt:def:ring-four-classes}, with coefficient ring
\(E\in\SetOf{\K[L],\K(L),\K((L^{-1}))}\), and \(F=\sum_{n\ge n_0}p_nt^n\)
denotes the block representation of an element of \(R\).

\section{The outer valuation}
\label{plt:sec:lead-outer-valuation}

\begin{definition}[Outer valuation]
\label{plt:def:lead-ordt}
For nonzero \(F\in R\) set
\begin{equation}\label{plt:eq:lead-ordt}
  \ordt(F)\DefinitionEquals\min\SetOf{n\in\IntegerNumbers:p_n\ne0},
  \qquad
  \ordt(0)\DefinitionEquals+\infty .
\end{equation}
The symbol \(+\infty\) obeys \(+\infty+m=+\infty\) and
\(+\infty>m\) for every \(m\in\IntegerNumbers\).
\end{definition}

The minimum exists because the set of occurring \(n\) is a nonempty subset of
\(\IntegerNumbers\) bounded below.  The convention \(\ordt(0)=+\infty\) is not
cosmetic: it is what makes the laws below hold without case distinctions, and
it is what makes \(\SetOf{F:\ordt(F)\ge N}\) an additive subgroup.

\begin{theorem}[Valuation laws]
\label{plt:thm:ring-valuation-laws}
Let \(F,G\in R\).  Then
\begin{enumerate}[label=\textup{(\roman*)}]
\item \(\ordt(FG)=\ordt(F)+\ordt(G)\);
\item \(\ordt(F+G)\ge\min\SetOf{\ordt(F),\ordt(G)}\);
\item \(\ordt(F)=+\infty\) if and only if \(F=0\), and
  \(\ordt(-F)=\ordt(F)\).
\end{enumerate}
Moreover the inequality in \textup{(ii)} is \emph{strict} if and only if
\begin{equation}\label{plt:eq:lead-strict-condition}
  \ordt(F)=\ordt(G)<+\infty
  \quad\text{and}\quad
  p_{\ordt(F)}+q_{\ordt(G)}=0,
\end{equation}
% ed.: ``\(\ordt\) is a discrete valuation on \(R\) with valuation ring
% \(E[[t]]\)'' overstates the case for two of the four classes: \(\PLpow\) and
% \(\PLlau\) are not fields, and on \(\PLpow\) the image of \(\ordt\) is
% \(\NaturalNumbers\), not \(\IntegerNumbers\).  The claim is split accordingly.
where \(q_n\) denotes the blocks of \(G\).  In particular \(\ordt\) satisfies
the axioms of a valuation with values in
\(\IntegerNumbers\cup\SetOf{+\infty}\) on each of the four classes, and
\(\SetOf{F\in R:\ordt(F)\ge0}=E[[t]]\).  When \(E\) is a field --- that is,
for \(R=\PLrat\) and \(R=\PLit\) --- \(R\) is a field
(\cref{plt:prop:ring-chain}\,(ii)), \(\ordt\) is surjective
onto \(\IntegerNumbers\), and \(\ordt\) is a discrete valuation of \(R\) in
the usual sense, with valuation ring \(E[[t]]\).  For \(R=\PLpow\) and
\(R=\PLlau\) the same laws hold, but \(R\) is not a field and the phrase
``valuation ring'' is not available.
\end{theorem}

\begin{proof}
(iii) is immediate from \eqref{plt:eq:lead-ordt}.

(i)  If \(F=0\) or \(G=0\) both sides are \(+\infty\).  Otherwise put
\(m=\ordt(F)\), \(r=\ordt(G)\).  For \(n<m+r\) the block of \(FG\) at \(t^n\)
is \(\sum_{i+k=n}p_iq_k\); in each summand either \(i<m\) or \(k<r\) (else
\(i+k\ge m+r>n\)), so every summand vanishes.  At \(n=m+r\) the only
surviving pair is \((m,r)\), giving the block \(p_mq_r\), which is nonzero
because \(E\) is an integral domain.  Hence \(\ordt(FG)=m+r\).

(ii) and the strictness criterion.  Write \(\mu=\min\SetOf{\ordt F,\ordt G}\).
For \(n<\mu\) both \(p_n\) and \(q_n\) vanish, so the block of \(F+G\) at
\(t^n\) vanishes; this is (ii).  For the criterion, suppose first
\(\ordt F\ne\ordt G\), say \(\ordt F<\ordt G\).  Then the block of \(F+G\) at
\(t^{\ordt F}\) is \(p_{\ordt F}+q_{\ordt F}=p_{\ordt F}\ne0\), so equality
holds.  If \(\ordt F=\ordt G=\mu<+\infty\), the block of \(F+G\) at \(t^\mu\)
is \(p_\mu+q_\mu\), and the inequality is strict exactly when this vanishes.
If \(\mu=+\infty\) then \(F=G=0\) and both sides are \(+\infty\).

Finally, (i)--(iii) are the axioms of a valuation with values in
\(\IntegerNumbers\cup\SetOf{+\infty}\); the set \(\SetOf{\ordt\ge0}\) is, by
\eqref{plt:eq:lead-ordt}, the set of series with no negative blocks, namely
\(E[[t]]\).  That \(E((t))\) is a field when \(E\) is one --- so that
\(\ordt\) is then a discrete valuation of a field with valuation ring
\(E[[t]]\) --- is \cref{plt:prop:ring-chain}\,(ii), whose proof uses only
\cref{plt:cor:ring-geometric} and not the present theorem.
The argument used about \(E\) only that it is
an integral domain, so the same laws hold in \(E[[t]]\) and \(E((t))\) for
every integral domain \(E\); this is the form in which they are applied in
\cref{plt:prop:ring-chain}.
% ed.: the domain hypothesis enters only through the reverse inequality in (i).
% For an arbitrary commutative ring \(E\) the subadditivity
% \(\ordt(FG)\ge\ordt F+\ordt G\), together with (ii) and (iii), still holds in
% \(E[[t]]\) and \(E((t))\), by the first half of the computation in (i); this
% weaker form is what \cref{plt:cor:ring-geometric} uses, and stating it here
% is what keeps that corollary free of a circular appeal.
% ed.: S1 and S5 record the strictness of (ii) as happening ``exactly when the
% leading blocks cancel'', omitting the hypothesis \(\ordt F=\ordt G\); with
% unequal valuations no cancellation is possible and equality always holds.
\end{proof}

\begin{example}[Strictness needs equal valuations and full cancellation]
\label{plt:ex:lead-strictness}
With \(F=(L^2+L)t^3+t^4\) and \(G=(-L^2-L)t^3+t^5\) we have
\(\ordt F=\ordt G=3\) but \(\ordt(F+G)=4\).  Cancelling only part of a block
is not enough: \(F'=(L^2+L)t^3\) and \(G'=(-L^2+2)t^3\) give
\(F'+G'=(L+2)t^3\), still of valuation \(3\).  And with \(F''=t^3\),
\(G''=-t^4\) the valuations differ, so \(\ordt(F''+G'')=3\) whatever the
blocks are.
\end{example}

\section{The rank-two valuation and the three leading objects}
\label{plt:sec:lead-rank-two}

\begin{definition}[Leading data]
\label{plt:def:lead-leading-data}
% ed.: the display ``\(p_\nu=aL^d+\text{lower powers of }L\)'' is literally a
% polynomial expansion, but \(R\) here ranges over all four classes, where a
% block need not be a polynomial.  The reading is fixed explicitly.
Let \(F\in R\) be nonzero and \(\nu=\ordt(F)\).  Write the block \(p_\nu\) in
its \(L^{-1}\)-expansion inside \(\K((L^{-1}))\) --- for \(E=\K[L]\) this is
the polynomial itself, and for \(E=\K(L)\) it is the image under the embedding
of \cref{plt:prop:ring-chain} --- so that
\(p_{\nu}=aL^{d}+(\text{strictly lower powers of }L)\) with \(a\in\K^\times\)
and \(d=\degL p_\nu\) as in \cref{plt:def:ring-rank-two}.  Then:
\begin{itemize}
\item the \emph{leading coefficient block} of \(F\) is the whole coefficient
  \(p_\nu\);
\item the \emph{leading monomial} is \(\lm(F)=t^{\nu}L^{d}\);
\item the \emph{scalar leading coefficient} is \(\lc(F)=a\).
\end{itemize}
For \(F=0\) none of the three is defined.
\end{definition}

\begin{remark}[Leading block, leading monomial, leading scalar]
\label{plt:rmk:lead-three-objects}
\(p_\nu\) is a polynomial (or rational function, or Laurent series in
\(L^{-1}\)); \(\lm(F)\) is a monomial; \(\lc(F)\) is a scalar.  The
distinction is not pedantic.  For \(F=t(L+1)\) we have \(\lc(F)=1\in\K^\times\)
and \(\lm(F)=tL\), yet \(t^{-1}F=L+1\) is not invertible in \(\PLlau\): what
governs invertibility is the leading \emph{block}, not the leading scalar.
Every ``the leading coefficient is \(1\), so we may divide'' argument must
name which of the three objects it means.
\end{remark}

\begin{definition}[Rank-two valuation]
\label{plt:def:ring-rank-two}
For nonzero \(F\in R\) set
\begin{equation}\label{plt:eq:lead-rank-two}
  v(F)\DefinitionEquals\bigl(\ordt(F),\,-\degL p_{\ordt(F)}\bigr)
  \in\Gamma=\IntegerNumbers\times\IntegerNumbers,
\end{equation}
lexicographically ordered, and \(v(0)\DefinitionEquals\infty_\Gamma\), a
symbol larger than every element of \(\Gamma\) and absorbing under addition.
Here \(\degL\) of a nonzero block means the largest \(j\) with a nonzero
coefficient of \(L^{j}\); for \(E=\K((L^{-1}))\) this exists because inner
supports are bounded above (\cref{plt:thm:ring-hahn}), and for
\(E=\K(L)\) it is \(\deg\) numerator minus \(\deg\) denominator, computed in
\(\K((L^{-1}))\) under the embedding of \cref{plt:prop:ring-chain}.
Thus \(v(t^nL^j)=(n,-j)\), matching \eqref{plt:eq:ring-value-group}, and
\(\ordt\) is the coarsening of \(v\) obtained by projecting to the first
coordinate.
\end{definition}

\begin{theorem}[Rank-two valuation laws and leading-term arithmetic]
\label{plt:thm:lead-rank-two-laws}
For \(F,G\in R\),
\begin{enumerate}[label=\textup{(\roman*)}]
% ed.: ``equivalently'' was wrong: additivity of \(v\) is equivalent to
% multiplicativity of \(\lm\) alone; multiplicativity of \(\lc\) is a further
% assertion (proved below), not a restatement.
\item \(v(FG)=v(F)+v(G)\) --- equivalently, \(\lm(FG)=\lm(F)\lm(G)\) for
  nonzero \(F,G\) --- and, in addition,
  \[
    \lc(FG)=\lc(F)\lc(G)\qquad(F,G\ne0);
  \]
\item \(v(F+G)\ge\min\SetOf{v(F),v(G)}\), with equality unless
  \(v(F)=v(G)<\infty_\Gamma\) and \(\lc(F)+\lc(G)=0\);
\item if \(\lm(F)\ne\lm(G)\), then \(\lm(F+G)\) is the \(\succdom\)-greater of
  \(\lm(F)\) and \(\lm(G)\), and \(\lc(F+G)\) is the corresponding leading
  scalar.
\end{enumerate}
% ed.: ``with value group \(\Gamma\)'' is false as stated for two of the four
% classes: on \(\PLpow\) and \(\PLlau\) a block is a polynomial, so
% \(\degL\ge0\) and the image of \(v\) never meets \(\IntegerNumbers\times
% \PositiveIntegers\).  The image and the group it generates are separated.
Consequently \(v\) is a valuation on \(R\) with values in the
lexicographically ordered group \(\Gamma\).  Its image is all of \(\Gamma\)
exactly when \(E\) is \(\K(L)\) or \(\K((L^{-1}))\); on \(\PLlau\) the image
is \(\SetOf{(n,-d):n\in\IntegerNumbers,\ d\in\NaturalNumbers}\) and on
\(\PLpow\) it is \(\SetOf{(n,-d):n,d\in\NaturalNumbers}\).  In every case the
image generates \(\Gamma\), so the associated value group is \(\Gamma\), which
has rank two by \cref{plt:lem:ring-rank-two-group}.
\end{theorem}

\begin{proof}
(i)  The cases involving \(0\) are handled by the absorbing convention.  Let
\(F,G\ne0\), \(m=\ordt F\), \(r=\ordt G\).  By
\cref{plt:thm:ring-valuation-laws}\,(i) the leading block of \(FG\) is
\(p_mq_r\).  In \(E\) --- a polynomial ring, a rational function field, or a
field of Laurent series in \(L^{-1}\) --- the top \(L\)-degree is additive and
the top coefficients multiply, because \(\K\) is a domain:
\(\degL(p_mq_r)=\degL p_m+\degL q_r\) and the coefficient of the top power is
the product of the two top coefficients.  Hence
\(v(FG)=(m+r,-\degL p_m-\degL q_r)=v(F)+v(G)\), which is the displayed
statement about \(\lm\) and \(\lc\).

(ii)--(iii)  If \(\ordt F<\ordt G\), then \(v(F)<v(G)\) lexicographically and
the block of \(F+G\) at \(t^{\ordt F}\) equals \(p_{\ordt F}\), so
\(v(F+G)=v(F)=\min\).  If \(\ordt F=\ordt G=\mu\) but
\(\degL p_\mu>\degL q_\mu\), then again \(v(F)<v(G)\), and the top \(L\)-term
of \(p_\mu+q_\mu\) is that of \(p_\mu\), so \(v(F+G)=v(F)\).  These two cases
cover \(\lm(F)\ne\lm(G)\) and give (iii).  In the remaining case
\(\lm(F)=\lm(G)=t^\mu L^{d}\), the coefficient of \(L^d\) in \(p_\mu+q_\mu\)
is \(\lc(F)+\lc(G)\); if this is nonzero then \(v(F+G)=v(F)=v(G)\), and if it
vanishes then \(v(F+G)>v(F)\).

\emph{The image of \(v\).}  \(v(t^nL^{d})=(n,-d)\), and \(t^nL^{d}\) lies in
\(\PLpow\) for \(n,d\ge0\), in \(\PLlau\) for \(n\in\IntegerNumbers\),
\(d\ge0\), and in \(\PLrat\subset\PLit\) for all \(n,d\in\IntegerNumbers\);
conversely \(\degL\) of a block of an element of \(\PLlau\) is a nonnegative
integer and \(\ordt\ge0\) on \(\PLpow\), which gives the reverse inclusions.
Since \((1,0)\) and \((0,-1)\) lie in every one of these images, each
generates \(\Gamma\).  Rank two is then
\cref{plt:lem:ring-rank-two-group}.
\end{proof}

\begin{remark}[What \(\ordt\) forgets]
\label{plt:rmk:lead-coarse}
\(\ordt\) is deliberately coarser than \(v\): it ignores the internal
\(L\)-ordering of a block.  That is precisely what makes it usable as a
computational filtration --- the set \(\SetOf{\ordt\ge N}\) is a
\emph{finitely describable} tail, whereas \(\SetOf{v\ge(N,-d)}\) cuts a block
in half and is not stable under the operations of Part~II.  All truncation,
completeness and summability statements in this volume are for \(\ordt\); the
finer \(v\) is used only to compare individual monomials.
\end{remark}

\section{Uniqueness of the block representation}
\label{plt:sec:lead-uniqueness}

The block form \eqref{plt:eq:ring-double-form} is unique by construction if one
\emph{defines} an element to be its family of coefficients.  What actually
needs proof is that an element presented as a formally summable sum of
monomials --- which is how every later formula presents its output --- has
uniquely determined blocks, independently of the order of summation.  The
proof uses the summability calculus of \cref{plt:sec:trunc-summability}; we
state it here because it belongs to the leading-data discussion.

\begin{theorem}[Uniqueness of the block representation]
\label{plt:thm:lead-uniqueness}
Let \((a_{n,j})\) be a family of scalars indexed by
\(\IntegerNumbers\times\NaturalNumbers\) whose index set satisfies
\cref{plt:thm:ring-support}\,(b).  Then the family of monomial terms
\(\bigl(a_{n,j}t^nL^j\bigr)\) is formally summable in \(\PLlau\), its sum
\(F\) is independent of the enumeration, and
\[
  [t^n]F=\sum_{j}a_{n,j}L^j,
  \qquad
  [t^nL^j]F=a_{n,j}
  \qquad\text{for all }(n,j).
\]
In particular, if \(\sum a_{n,j}t^nL^j=\sum b_{n,j}t^nL^j\) with both index
sets admissible, then \(a_{n,j}=b_{n,j}\) for all \((n,j)\): the monomials
\(t^nL^j\) are linearly independent over \(\K\) in the strong, summable sense.
\end{theorem}

\begin{proof}
Summability: for \(N\in\IntegerNumbers\), a term \(a_{n,j}t^nL^j\) has
\(\ordt=n\), and the terms with \(n<N\) are indexed by
\(\bigcup_{n_0\le n<N}\SetOf{n}\times S_n\), a finite union of finite sets by
\cref{plt:thm:ring-support}\,(b).  Hence only finitely many terms have
\(\ordt<N\), which is
\cref{plt:def:trunc-summable}.  By
\cref{plt:thm:trunc-summability}\,(ii)--(iii) the sum exists, is independent
of the enumeration, and satisfies
\([t^n]F=\sum_{(n',j)}[t^n]\bigl(a_{n',j}t^{n'}L^j\bigr)
=\sum_ja_{n,j}L^j\), a finite sum.  Extracting the coefficient of \(L^j\) from
this polynomial gives \([t^nL^j]F=a_{n,j}\).  The last assertion follows
because \(F\) determines all its coefficients.
\end{proof}

\begin{proposition}[Leading-block factorization]
\label{plt:prop:lead-factorization}
Let \(E\) be a \emph{field} and \(F\in E((t))\) nonzero.  Then there are
unique \(\nu\in\IntegerNumbers\), \(c\in E^\times\) and \(U\in tE[[t]]\) with
\begin{equation}\label{plt:eq:lead-factorization}
  F=c\,t^{\nu}\,(1+U),
\end{equation}
namely \(\nu=\ordt(F)\), \(c=p_\nu\), and \(U=\sum_{n\ge1}(p_{\nu+n}/c)t^n\).
The factorization is available in \(\PLrat\) and in \(\PLit\); over
\(E=\K[L]\) it is available with \(c=p_\nu\) and \(U\in t\K(L)[[t]]\), but in
general \emph{not} with \(U\in t\K[L][[t]]\).
\end{proposition}

\begin{proof}
Existence over a field is the displayed formula, which makes sense because
\(c=p_\nu\in E^\times\).  Uniqueness: applying \(\ordt\) to
\eqref{plt:eq:lead-factorization} and using
\cref{plt:thm:ring-valuation-laws}\,(i) with \(\ordt(1+U)=0\) gives
\(\nu=\ordt(F)\); comparing blocks at \(t^\nu\) gives \(c=p_\nu\); and then
\(U=c^{-1}t^{-\nu}F-1\) is determined.  For the last assertion take
\(F=L+t\in\PLlau\).  A factorization \(F=L(1+U)\) with
\(U=\sum_{n\ge1}u_n(L)t^n\in t\K[L][[t]]\) would force, at \(t^1\),
\(1=Lu_1(L)\) in \(\K[L]\), which is impossible; and \(c\) must equal \(L\) by
the uniqueness just proved.
% ed.: S2 states the prefactor form \(A=cX^{\rho}L^{\beta}(1+U)\) ``whenever
% available''; the example \(L+t\) shows that over \(\K[L]\) it is available
% only after passing to \(\K(L)\), so no argument may assume it inside
% \(\PLlau\).
\end{proof}

\section{The log-degree profile}
\label{plt:sec:lead-profile}

\begin{definition}[Log-degree profile]
\label{plt:def:lead-profile}
For \(F\in\PLlau\) define \(\dprofile{F}:\IntegerNumbers\to
\NaturalNumbers\cup\SetOf{-\infty}\) by
\begin{equation}\label{plt:eq:lead-profile}
  \dprofile{F}(n)=
  \begin{cases}
    \degL p_n,&p_n\ne0,\\
    -\infty,&p_n=0,
  \end{cases}
\end{equation}
with the conventions \(-\infty+m=-\infty\) and \(\max(-\infty,m)=m\).  Absent
blocks are extended by zero, so \(\dprofile{F}(n)=-\infty\) for
\(n<\ordt(F)\).
\end{definition}

The profile is genuinely extra data: \(\ordt(F)\) is the first \(n\) at which
\(\dprofile{F}(n)>-\infty\), and nothing more about \(\dprofile{F}\) can be
recovered from it.

\begin{proposition}[Profile under the ring operations]
\label{plt:prop:lead-profile-laws}
For \(F,G\in\PLlau\) and every \(n\in\IntegerNumbers\),
\begin{align}
  \dprofile{F+G}(n)&\le\max\SetOf{\dprofile{F}(n),\dprofile{G}(n)},
    \label{plt:eq:lead-profile-sum}\\
  \dprofile{FG}(n)&\le
    \max_{i+k=n}\bigl(\dprofile{F}(i)+\dprofile{G}(k)\bigr),
    \label{plt:eq:lead-profile-product}
\end{align}
where in \eqref{plt:eq:lead-profile-product} the maximum is over the finitely
many pairs with \(i\ge\ordt F\) and \(k\ge\ordt G\), and is \(-\infty\) if
there are none.  Equality holds in \eqref{plt:eq:lead-profile-sum} unless the
top \(L\)-terms cancel, and in \eqref{plt:eq:lead-profile-product} unless the
sum of the products of top coefficients over the maximizing pairs vanishes.
\end{proposition}

\begin{proof}
For a sum of polynomials the degree is at most the maximum of the degrees,
with equality unless the coefficients of that maximal power cancel; apply this
to \(p_n+q_n\) for \eqref{plt:eq:lead-profile-sum} and to
\(\sum_{i+k=n}p_iq_k\) for \eqref{plt:eq:lead-profile-product}, using
\(\degL(p_iq_k)=\degL p_i+\degL q_k\) (valid with the \(-\infty\) convention).
The index set is finite because \(i\ge\ordt F\) and \(k=n-i\ge\ordt G\) force
\(\ordt F\le i\le n-\ordt G\).
\end{proof}

\begin{definition}[Series of bounded slope]
\label{plt:def:lead-slope-class}
For \(\sigma\ge0\) let
\[
  \mathcal D_\sigma
  \DefinitionEquals
  \SetOf{F\in\PLlau:\ \exists C\in\RealNumbers\ \forall n,\
    \dprofile{F}(n)\le\sigma n+C}.
\]
\end{definition}

\begin{proposition}[Slope transport]
\label{plt:prop:lead-slope}
Let \(F,G\in\PLlau\) be nonzero with \(m=\ordt F\), \(r=\ordt G\), and suppose
\[
  \dprofile{F}(n)\le A+\alpha n\ (n\ge m),
  \qquad
  \dprofile{G}(n)\le B+\beta n\ (n\ge r),
\]
with \(\alpha\ge\beta\ge0\).  Then
\begin{equation}\label{plt:eq:lead-slope-bound}
  \dprofile{FG}(n)\le\bigl(A+B-(\alpha-\beta)r\bigr)+\alpha n
  \qquad(n\ge m+r).
\end{equation}
In particular each \(\mathcal D_\sigma\) is a subring of \(\PLlau\) containing
\(\K[L]\) and \(t^{\pm1}\), and
\(\mathcal D_\sigma\subseteq\mathcal D_{\sigma'}\) for \(\sigma\le\sigma'\).
If moreover \(r\ge0\) --- in particular if \(F,G\in\PLpow\) --- then
\eqref{plt:eq:lead-slope-bound} implies the simpler bound
\(\dprofile{FG}(n)\le A+B+\alpha n\).
\end{proposition}

\begin{proof}
By \eqref{plt:eq:lead-profile-product}, for \(i+k=n\) with \(i\ge m\),
\(k\ge r\),
\[
  \dprofile{F}(i)+\dprofile{G}(k)
  \le A+B+\alpha i+\beta k
  = A+B+\beta n+(\alpha-\beta)i
  \le A+B+\beta n+(\alpha-\beta)(n-r),
\]
using \(i=n-k\le n-r\) and \(\alpha-\beta\ge0\); this is
\eqref{plt:eq:lead-slope-bound}.  For the subring claim, closure under
subtraction follows from \eqref{plt:eq:lead-profile-sum}, and closure under
multiplication from \eqref{plt:eq:lead-slope-bound}, applied with
\(\alpha=\beta=\sigma\) (so that the correction term \(-(\alpha-\beta)r\)
vanishes); the constants \(A,B\) are fixed for given \(F,G\).
% ed.: the monotonicity clause \(\mathcal D_\sigma\subseteq\mathcal
% D_{\sigma'}\) was asserted without proof, and it is not the trivial
% inequality \(\sigma n\le\sigma'n\): that fails for \(n<0\).  The argument
% below uses that \(\dprofile{F}\) is \(-\infty\) below \(\ordt F\).
For the monotonicity, let \(\sigma\le\sigma'\) and
\(F\in\mathcal D_\sigma\), say \(\dprofile{F}(n)\le\sigma n+C\) for all \(n\),
and put \(m=\ordt F\).  For \(n<m\) the left-hand side is \(-\infty\) and
nothing is to be proved; for \(n\ge m\) one has
\(\sigma n+C=\sigma'n+C+(\sigma-\sigma')n\le\sigma'n+C+(\sigma-\sigma')m\),
because \(\sigma-\sigma'\le0\).  Hence \(F\in\mathcal D_{\sigma'}\) with the
constant \(C+(\sigma-\sigma')m\).
% ed.: S2 states \(\deg c_n\le A+B+\max(\alpha,\beta)n\) for the product.  That
% form is correct only when the factor of smaller slope has nonnegative outer
% order; the counterexample of \cref{plt:ex:lead-slope-fails} violates it, and
% \eqref{plt:eq:lead-slope-bound} is the corrected statement, attained there
% with equality.
\end{proof}

\begin{example}[The naive slope bound fails on the Laurent ring]
\label{plt:ex:lead-slope-fails}
Take \(F=\sum_{n\ge0}L^nt^n\) with \(A=0\), \(\alpha=1\), \(m=0\), and
\(G=t^{-2}\) with \(B=0\), \(\beta=0\), \(r=-2\).  Then
\(FG=\sum_{n\ge-2}L^{n+2}t^{n}\), so \(\dprofile{FG}(n)=n+2\).  The naive
bound \(A+B+\max(\alpha,\beta)n=n\) is violated for every \(n\), while
\eqref{plt:eq:lead-slope-bound} gives \(0+0-(1-0)(-2)+n=n+2\), attained.
\end{example}

\begin{remark}[Degree bounds are hypotheses, never definitions]
\label{plt:rmk:lead-no-degree-bound}
No class in \cref{plt:def:ring-four-classes} imposes any bound on
\(\dprofile{F}\).  A statement of the form \(\dprofile{F}(n)\le\sigma n+C\) is
either an explicit hypothesis, or an invariant proved for the particular
construction at hand --- as it is for the Lambert polynomials, where
\(\deg\LamP{n}=n\) for \(n\ge1\).  It is never available merely because the
object lies in \(\PLlau\).  Two consequences are worth isolating.  First,
membership in \(\mathcal D_\sigma\) must be re-proved after every operation
that is not covered by \cref{plt:prop:lead-slope}; the subring property gives
sums and products only.  Second, a formal tail \(\FormalBigOAt{t^N}{t}\)
becomes an analytic estimate only after a degree bound on the discarded
blocks is supplied; see \cref{plt:rmk:trunc-no-size}.
\end{remark}

\begin{example}[Leading data and profile of a typical element]
\label{plt:ex:lead-worked}
Let
\[
  F=X^2+3XL+7+(L^3-2L)t+\sum_{n\ge2}(L^n+1)t^n
   =t^{-2}+3Lt^{-1}+7+(L^3-2L)t+\sum_{n\ge2}(L^n+1)t^n .
\]
Then \(\ordt(F)=-2\); the leading block is the constant polynomial \(1\);
\(\lm(F)=t^{-2}=X^2\); \(\lc(F)=1\); and \(v(F)=(-2,0)\).  The profile is
\[
  \dprofile{F}(-2)=0,\quad
  \dprofile{F}(-1)=1,\quad
  \dprofile{F}(0)=0,\quad
  \dprofile{F}(1)=3,\quad
  \dprofile{F}(n)=n\ (n\ge2),
\]
and \(-\infty\) for \(n<-2\).  It is not monotone, and the smallest affine
bound of slope \(1\) is \(\dprofile{F}(n)\le n+2\); thus
\(F\in\mathcal D_1\).  Here the leading block happens to be a nonzero scalar,
so \(F\) is invertible in \(\PLlau\) by the unit criterion of Part~II.  Delete
the term \(X^2\): the leading block of \(F-t^{-2}\) becomes \(3L\), whose
scalar leading coefficient is still \(3\in\K^\times\) --- and yet \(3L\) is not
a unit of \(\K[L]\), so the reciprocal of \(F-t^{-2}\) leaves \(\PLlau\).
This is \cref{plt:rmk:lead-three-objects} in a single example.
\end{example}

\chapter{Truncation, formal tails, and summability}
\label{plt:sec:truncation}

\section{The exclusive truncation operator}
\label{plt:sec:trunc-operator}

\begin{definition}[Exclusive outer truncation]
\label{plt:def:trunc-exclusive}
For \(N\in\IntegerNumbers\) and \(F=\sum_{n\ge n_0}p_n(L)t^n\in\PLlau\) put
\begin{equation}\label{plt:eq:trunc-exclusive}
  \Trunc{F}{N}\DefinitionEquals\sum_{n<N}p_n(L)t^n .
\end{equation}
The sum is finite because \(n_0\le n<N\); it is empty, hence \(0\), when
\(N\le n_0\), and \(\Trunc{0}{N}=0\).  Every retained block is retained
\emph{whole}: all powers of \(L\) at a retained level survive.  The phrase
``through outer order \(t^{N-1}\)'' means the exclusive cutoff \(<N\).
\end{definition}

\begin{remark}[Reconciliation with the inclusive convention]
\label{plt:rmk:trunc-inclusive}
% ed.: S6 uses the inclusive marker throughout; the catalogue makes the
% exclusive operator normative, so every transferred formula is shifted by one
% block.
Source~S6 writes \([F]_{\le N}=\sum_{n\le N}p_n(L)t^n\), which is
\emph{inclusive}.  The translation is
\begin{equation}\label{plt:eq:trunc-inclusive-dictionary}
  [F]_{\le N}=\Trunc{F}{N+1},
  \qquad
  \Trunc{F}{N}=[F]_{\le N-1}.
\end{equation}
Any statement imported from S6 --- in particular its residual and stopping
criteria --- must absorb this one-block shift.  We use
\eqref{plt:eq:trunc-exclusive} exclusively.
\end{remark}

\begin{proposition}[Elementary properties]
\label{plt:prop:trunc-basic}
For \(F,G\in\PLlau\), \(c\in\K\) and \(M,N\in\IntegerNumbers\):
\begin{enumerate}[label=\textup{(\roman*)}]
\item \(\Trunc{\cdot}{N}\) is \(\K\)-linear and idempotent, and
  \(\Trunc{(\Trunc{F}{M})}{N}=\Trunc{F}{\min(M,N)}\);
\item \(\ordt\bigl(F-\Trunc{F}{N}\bigr)\ge N\), and
  \(\ordt(\Trunc{F}{N})\ge\ordt(F)\);
\item \(\Trunc{F}{N}=\Trunc{G}{N}\) if and only if \(\ordt(F-G)\ge N\);
\item \(F=G\) if and only if \(\Trunc{F}{N}=\Trunc{G}{N}\) for every
  \(N\in\IntegerNumbers\);
\item \(\Trunc{(cF+G)}{N}=c\Trunc{F}{N}+\Trunc{G}{N}\).
\end{enumerate}
\end{proposition}

\begin{proof}
All five are read off blockwise from \eqref{plt:eq:trunc-exclusive}.  For (iii),
\(F-G\) has all blocks below \(N\) equal to \(p_n-q_n\), which vanish for all
\(n<N\) exactly when the two truncations agree; and \(\ordt(F-G)\ge N\) says
precisely that those blocks vanish.  For (iv), if the truncations agree for
all \(N\) then \(\ordt(F-G)\ge N\) for all \(N\), so \(F-G=0\) by
\cref{plt:thm:ring-valuation-laws}\,(iii).
\end{proof}

\begin{proposition}[Truncated product]
\label{plt:prop:trunc-product}
Let \(F,G\in\PLlau\) be nonzero, \(m=\ordt F\), \(r=\ordt G\), and
\(N\in\IntegerNumbers\).  Then
\begin{equation}\label{plt:eq:trunc-product}
  \Trunc{(FG)}{N}
  =\Trunc{\bigl(\Trunc{F}{N-r}\cdot\Trunc{G}{N-m}\bigr)}{N}.
\end{equation}
If in addition \(m\ge0\) and \(r\ge0\) --- in particular if
\(F,G\in\PLpow\) --- then also
\begin{equation}\label{plt:eq:trunc-product-pow}
  \Trunc{(FG)}{N}
  =\Trunc{\bigl(\Trunc{F}{N}\cdot\Trunc{G}{N}\bigr)}{N}.
\end{equation}
\end{proposition}

\begin{proof}
Write \(F=\Trunc{F}{N-r}+F_1\) and \(G=\Trunc{G}{N-m}+G_1\), so that
\(\ordt F_1\ge N-r\) and \(\ordt G_1\ge N-m\) by
\cref{plt:prop:trunc-basic}\,(ii).  Then
\[
  FG-\Trunc{F}{N-r}\Trunc{G}{N-m}
  =F_1G+\Trunc{F}{N-r}\,G_1 .
\]
By \cref{plt:thm:ring-valuation-laws}, \(\ordt(F_1G)\ge(N-r)+r=N\) and
\(\ordt\bigl(\Trunc{F}{N-r}G_1\bigr)\ge m+(N-m)=N\), using
\(\ordt(\Trunc{F}{N-r})\ge m\).  Hence the difference has \(\ordt\ge N\), and
\eqref{plt:eq:trunc-product} follows from
\cref{plt:prop:trunc-basic}\,(iii).  When \(m,r\ge0\) the same computation
with \(N\) in place of \(N-r\) and \(N-m\) gives
\(\ordt(F_1G)\ge N+r\ge N\) and
\(\ordt(\Trunc{F}{N}G_1)\ge m+N\ge N\), which is
\eqref{plt:eq:trunc-product-pow}.
\end{proof}

\begin{example}[The symmetric law fails on the Laurent ring]
\label{plt:ex:trunc-product-fails}
% ed.: S5 states \eqref{plt:eq:trunc-product-pow} without the hypothesis
% \(\ordt\ge0\); on \(\PLlau\) it is false, and
% \eqref{plt:eq:trunc-product} is the correct general form.
Let \(F=t^{-1}+1\), \(G=t^{-1}\) and \(N=0\).  Then \(FG=t^{-2}+t^{-1}\), so
\(\Trunc{(FG)}{0}=t^{-2}+t^{-1}\).  On the other hand
\(\Trunc{F}{0}=t^{-1}\) and \(\Trunc{G}{0}=t^{-1}\), whose product truncates
to \(t^{-2}\).  The block \(t^{-1}\) has been lost, because the term \(1\) of
\(F\), discarded at cutoff \(0\), meets the pole \(t^{-1}\) of \(G\).
Formula \eqref{plt:eq:trunc-product} keeps it: \(m=-1\), \(r=-1\), so one must
retain \(F\) below \(N-r=1\) and \(G\) below \(N-m=1\), giving
\((t^{-1}+1)t^{-1}=t^{-2}+t^{-1}\).
\end{example}

\begin{remark}[Precision accounting]
\label{plt:rmk:trunc-precision}
\Cref{plt:prop:trunc-product} is the prototype of every precision contract in
this volume: an operation must declare how much input precision it consumes.
Multiplication consumes \(-\ordt\) of the \emph{other} factor; equivalently,
each negative power of \(t\) in one factor costs one block of precision in the
other.  Composition with an inner argument of negative outer order amplifies
the requested precision by a factor, and is analysed in Part~III.
\end{remark}

\section{The formal tail token}
\label{plt:sec:trunc-formal-o}

\begin{definition}[Formal big-O]
\label{plt:def:trunc-formal-o}
For \(F,G\in\PLlau\) and \(N\in\IntegerNumbers\), the statement
\begin{equation}\label{plt:eq:trunc-formal-o}
  F=G+\FormalBigOAt{t^N}{t}
\end{equation}
means any of the following four equivalent things:
\[
  \ordt(F-G)\ge N;
  \qquad
  F-G\in t^N\PLpow;
  \qquad
  \Trunc{F}{N}=\Trunc{G}{N};
  \qquad
  [t^n](F-G)=0\ \text{for all }n<N.
\]
The token \(\FormalBigO\) is the same operator with its arguments left
implicit.
\end{definition}

\begin{proposition}[Equivalence, and the calculus of the token]
\label{plt:prop:trunc-formal-o}
The four conditions in \cref{plt:def:trunc-formal-o} are equivalent.
% ed.: \cref{plt:def:trunc-formal-o} defines the token only for an integer
% cutoff, so (ii) and (iii) below were type-incorrect when \(H\) or \(G\)
% vanishes and the printed exponent becomes \(+\infty\).  The reading of that
% degenerate case is now fixed explicitly.
Moreover, for \(F=G+\FormalBigOAt{t^N}{t}\) and any \(H,H'\in\PLlau\) and
\(M\in\IntegerNumbers\) --- where, in \textup{(ii)} and \textup{(iii)}, a tail
exponent that evaluates to \(+\infty\) (which happens only through a vanishing
factor) is to be read as the assertion that the difference in question is
\(0\):
\begin{enumerate}[label=\textup{(\roman*)}]
\item \(F+H=G+H+\FormalBigOAt{t^N}{t}\);
\item \(HF=HG+\FormalBigOAt{t^{N+\ordt H}}{t}\); in particular multiplication
  by \(H\) with \(\ordt H<0\) \emph{loses} \(|\ordt H|\) blocks of precision;
\item if also \(H=H'+\FormalBigOAt{t^{M}}{t}\), then
  \(FH=GH'+\FormalBigOAt{t^{K}}{t}\) with
  \(K=\min\SetOf{N+\ordt H,\ M+\ordt G}\);
\item \(\FormalBigOAt{t^N}{t}+\FormalBigOAt{t^M}{t}
  =\FormalBigOAt{t^{\min(N,M)}}{t}\) as a statement about the set
  \(t^N\PLpow+t^M\PLpow=t^{\min(N,M)}\PLpow\).
\end{enumerate}
\end{proposition}

\begin{proof}
Equivalence: \(\ordt(F-G)\ge N\) says all blocks of \(F-G\) below \(N\)
vanish, which is the fourth condition and, by
\cref{plt:prop:trunc-basic}\,(iii), the third; and a series with
\(\ordt\ge N\) is \(t^N\) times a series with \(\ordt\ge0\), that is, an
element of \(t^N\PLpow\), and conversely.

(i) is \cref{plt:prop:trunc-basic}\,(v).  (ii): \(H(F-G)\) has
\(\ordt=\ordt H+\ordt(F-G)\ge\ordt H+N\) by
\cref{plt:thm:ring-valuation-laws}\,(i).  (iii): write
\(FH-GH'=H(F-G)+G(H-H')\) and apply (ii) twice together with
\cref{plt:thm:ring-valuation-laws}\,(ii).  (iv): the inclusion
\(t^N\PLpow+t^M\PLpow\subseteq t^{\min}\PLpow\) is clear, and the reverse
holds because one of the two summands already equals \(t^{\min}\PLpow\).
\end{proof}

\begin{remark}[What the token does not assert]
\label{plt:rmk:trunc-no-size}
\(F=G+\FormalBigOAt{t^N}{t}\) asserts \emph{coefficient agreement below outer
order \(N\)} and nothing else.  It does not assert a numerical bound, a rate,
convergence, uniformity in \(L\), or the existence of any function with the
displayed expansion.  Concretely, the element
\[
  T=\sum_{n\ge N}L^{n^2}t^n
\]
satisfies \(T=\FormalBigOAt{t^N}{t}\), yet at \(X=\EulerE^{10}\), so that
\(L=10\), its very first term is
\(\EulerE^{-10N}10^{N^2}\), which for \(N=5\) exceeds \(1.9\times10^{3}\) and
grows without bound in \(N\).  The formal tail is small in the filtration, not
in size.  An estimate of size requires, in addition, a bound on
\(\dprofile{T}\) on the discarded range --- see
\cref{plt:rmk:lead-no-degree-bound} --- and then a separate analytic theorem.
\end{remark}

\begin{remark}[Contrast with the blockwise analytic statement]
\label{plt:rmk:trunc-blockwise}
% ed.: the displayed analytic statement was printed with none of \(f\), \(S\)
% or \(n_0\) introduced, contrary to the standing requirement that a statement
% name its regime and domain.  They are named here.
The analytic counterpart, proved for actual functions in Part~V, reads as
follows.  Let \(S\) be an unbounded subset of the positive ray (or, for
complex work, a sector on which a branch of \(\log\) has been fixed), let
\(f:S\to\ComplexNumbers\), and let
\(\sum_{n\ge n_0}p_n(L)t^n\in\PLlau\) be a candidate expansion.  One says that
\(f\) has this blockwise asymptotic expansion on \(S\) if for every
\(N>n_0\) there is an exponent \(M_N\in\NaturalNumbers\) with
\[
  f(X)-\sum_{n_0\le n<N}p_n(\log X)X^{-n}
  =\BigOAt{X^{-N}(\log X)^{M_N}}{X\to+\infty,\ X\in S}.
\]
% ed.: S1 offers instead the minimal Poincar\'e requirement that the remainder
% be little-o of the last retained term.  That condition does not determine
% the blocks: with \(p_n=L^5\) and \(q_n=L^5+L^3\) the two candidate blocks
% differ by \(L^3=\LittleOAt{L^5}{X\to+\infty}\), so uniqueness fails.  The
% displayed form, and S6's \(\LittleOAt{X^{-N}}{X\to+\infty}\) form, both do
% determine them.
The exponent \(M_N\) is part of the statement; it is exactly the data that
\cref{plt:rmk:trunc-no-size} shows the formal token cannot supply.
\end{remark}

\section{Why the Laurent tail is not an ideal}
\label{plt:sec:trunc-not-ideal}

\begin{proposition}[The tail is a filtration, not an ideal]
\label{plt:prop:trunc-not-ideal}
Fix \(N\in\IntegerNumbers\) and set \(\mathfrak F_N=t^N\PLpow\subseteq\PLlau\).
\begin{enumerate}[label=\textup{(\roman*)}]
\item \(t\) is a unit of \(\PLlau\), and the ideal of \(\PLlau\) generated by
  \(t^N\) is all of \(\PLlau\):
  \(\bigl(t^N\bigr)=\PLlau\).
\item \(\mathfrak F_N\) is an additive subgroup and a \(\PLpow\)-submodule of
  \(\PLlau\), but for no \(N\) is it an ideal of \(\PLlau\); indeed
  \(t^{-1}\cdot t^N=t^{N-1}\notin\mathfrak F_N\).
\item \(\mathfrak F_N=\SetOf{F\in\PLlau:\ordt F\ge N}\), the chain
  \((\mathfrak F_N)_{N\in\IntegerNumbers}\) is decreasing, exhaustive
  \(\bigl(\bigcup_N\mathfrak F_N=\PLlau\bigr)\) and separated
  \(\bigl(\bigcap_N\mathfrak F_N=0\bigr)\).
\item In \(\PLpow\), by contrast, \(t^N\PLpow\) for \(N\ge0\) is a genuine
  ideal and
  \[
    \PLpow/t^N\PLpow\ \cong\ \K[L][t]/(t^N)
  \]
  as \(\K[L]\)-algebras.
\end{enumerate}
\end{proposition}

\begin{proof}
(i)  \(t\cdot t^{-1}=1\) in \(\PLlau\), so \(t\in\PLlau^\times\) and hence
\(t^N\in\PLlau^\times\); an ideal containing a unit is the whole ring.

(ii)  \(\mathfrak F_N\) is closed under addition and under multiplication by
elements of \(\PLpow\) because \(\PLpow\) is a ring and
\(t^N\PLpow\cdot\PLpow= t^N\PLpow\).  It is not an ideal of \(\PLlau\):
\(t^N\in\mathfrak F_N\), \(t^{-1}\in\PLlau\), and
\(\ordt(t^{N-1})=N-1<N\), so \(t^{N-1}\notin\mathfrak F_N\) by (iii).

(iii)  \(F\in t^N\PLpow\) means \(F=t^NG\) with \(\ordt G\ge0\), i.e.
\(\ordt F\ge N\) by \cref{plt:thm:ring-valuation-laws}\,(i); conversely
\(t^{-N}F\in\PLpow\) whenever \(\ordt F\ge N\).  The chain is decreasing since
\(\ordt\ge N+1\) implies \(\ordt\ge N\); exhaustive because every element has
finite \(\ordt\) or is \(0\); separated because \(\ordt F\ge N\) for all \(N\)
forces \(\ordt F=+\infty\), i.e.\ \(F=0\).

(iv)  For \(N\ge0\), \(t^N\PLpow\) is the image of \(\PLpow\) under
multiplication by \(t^N\), hence an ideal of \(\PLpow\).  The inclusion
\(\K[L][t]\hookrightarrow\PLpow\) induces a \(\K[L]\)-algebra homomorphism
\(\varphi:\K[L][t]/(t^N)\to\PLpow/t^N\PLpow\).  It is surjective because
\(F\equiv\Trunc{F}{N}\pmod{t^N\PLpow}\) and \(\Trunc{F}{N}\in\K[L][t]\) for
\(F\in\PLpow\).  It is injective because a polynomial of \(t\)-degree \(<N\)
lying in \(t^N\PLpow\) has all blocks below \(N\) zero, hence is \(0\).
\end{proof}

\begin{checkpointbox}[title={The three legitimate replacements}]
% ed.: the label kind "box" is not in the normative list of kinds; these
% boxed statements are remarks, so each label is renamed to the "rmk" kind.
\label{plt:rmk:trunc-replacements}
On \(\PLlau\), never write ``modulo \(t^N\)'' or use quotient-ring language.
Use instead, in decreasing order of explicitness:
\begin{enumerate}
\item the valuation statement \(\ordt(F-G)\ge N\);
\item the formal tail token \(F=G+\FormalBigOAt{t^N}{t}\) of
  \cref{plt:def:trunc-formal-o};
\item equality of exclusive truncations \(\Trunc{F}{N}=\Trunc{G}{N}\).
\end{enumerate}
These three are equivalent by \cref{plt:prop:trunc-formal-o}.  Quotient-algebra
language becomes legitimate only after the finite Laurent principal part has
been removed --- that is, after multiplying by \(t^{-\ordt F}\) and working in
\(\PLpow\), where \cref{plt:prop:trunc-not-ideal}\,(iv) applies.
\end{checkpointbox}

\section{Coefficient extraction}
\label{plt:sec:trunc-extraction}

\begin{definition}[Extraction brackets]
\label{plt:def:trunc-extraction}
For \(F\in\PLlau\) written as in \eqref{plt:eq:ring-double-form},
\begin{equation}\label{plt:eq:trunc-extraction}
  [t^n]F\DefinitionEquals p_n(L)\in\K[L],
  \qquad
  [t^nL^j]F\DefinitionEquals a_{n,j}\in\K .
\end{equation}
The monomial is written to the \emph{left} of the bracket and the series to
the right.  An absent coefficient is \(0\).
\end{definition}

\begin{proposition}[Extraction is linear and filtration-continuous]
\label{plt:prop:trunc-extraction}
Each map \([t^n]\) and \([t^nL^j]\) is \(\K\)-linear, and
\([t^n]F=0\) for every \(F\in\mathfrak F_N\) with \(N>n\).  Consequently, if
\(F=G+\FormalBigOAt{t^N}{t}\) then \([t^n]F=[t^n]G\) for all \(n<N\), and
\(F\) is determined by the family \(\bigl([t^nL^j]F\bigr)_{(n,j)}\).
\end{proposition}

\begin{proof}
Linearity is blockwise addition.  If \(\ordt F\ge N>n\) then \(p_n=0\).  The
remaining statements follow from
\cref{plt:def:trunc-formal-o} and
\cref{plt:thm:lead-uniqueness}.
\end{proof}

\begin{remark}[Bracket collisions]
\label{plt:rmk:trunc-bracket}
The glyph \([\;\cdot\;]\) carries four unrelated meanings in this corpus:
coefficient extraction \([t^n]F\); the inclusive truncation \([F]_{\le N}\) of
S6, banished here by \cref{plt:rmk:trunc-inclusive}; the unsigned Stirling
cycle number \(\stirlingone{n}{k}\); and the Gaussian binomial coefficient.
They are told apart by argument type, and we always print the semantic form.
\end{remark}

\section{The separate logarithmic cutoff}
\label{plt:sec:trunc-log-cutoff}

\begin{warningbox}[title={A log-degree cap is not a truncation in the dominance order}]
% ed.: the label kind "box" is not in the normative list of kinds; these
% boxed statements are remarks, so each label is renamed to the "rmk" kind.
\label{plt:rmk:trunc-log-cutoff}
Bounding \(\degL p_n\) is an additional, independent approximation.  It is not
part of the \(t\)-adic topology, and --- this is the essential point --- inside
a fixed block the discarded high powers of \(L\) are the \emph{dominant}
terms, not the negligible ones: by
\cref{plt:def:ring-dominance}, \(t^nL^{j}\succdom t^nL^{k}\) whenever
\(j>k\).  Discarding them therefore removes the largest part of the block.
\end{warningbox}

\begin{example}[A log cap destroys a retained block]
\label{plt:ex:trunc-log-cap}
Let \(F=G=1+L^3t\in\PLpow\).  Then \(FG=1+2L^3t+L^6t^2\), so
\(\Trunc{(FG)}{2}=1+2L^3t\).  Capping the logarithmic degree at \(2\)
\emph{before} multiplying replaces both factors by \(1\) and returns
\(\Trunc{(FG)}{2}=1\): an entire retained \(t\)-block has been lost, not
merely perturbed.  Capping \emph{after} multiplying returns \(1\) as well.  No
\(t\)-adic precision statement detects the error, because both answers agree
modulo \(\FormalBigOAt{t^{1}}{t}\) and disagree at the retained level \(t^1\).
\end{example}

\begin{remark}[When an inner cutoff is legitimate]
\label{plt:rmk:trunc-inner-cutoff}
In \(\PLit\) the situation reverses: a block is a Laurent series in \(L^{-1}\)
whose discarded terms \(L^{-j}\) with \(j\) large \emph{are} the small ones,
so an inner cutoff is a genuine approximation in the rank-two order of
\cref{plt:def:ring-rank-two}.  A finite record in \(\PLit\) therefore needs a
two-dimensional precision region: an outer bound \(n<N\) and a level-dependent
inner bound \(j>J(n)\), with the staircase \(J\) chosen stable under the
intended convolutions.  An outer cutoff alone never makes an infinite
Laurent-log block finite.  In \(\PLlau\), by contrast, a log-degree cap is
legitimate only when a proved bound on \(\dprofile{F}\) says that the capped
powers are absent.
\end{remark}

\section{Completeness and separatedness}
\label{plt:sec:trunc-completeness}

\begin{definition}[The \(t\)-adic topology]
\label{plt:def:trunc-topology}
Give \(\PLlau\) the topology defined by the filtration
\((\mathfrak F_N)_{N\in\IntegerNumbers}\) of
\cref{plt:prop:trunc-not-ideal}, equivalently by the ultrametric
\begin{equation}\label{plt:eq:trunc-metric}
  d(F,G)=2^{-\ordt(F-G)},
  \qquad 2^{-\infty}\DefinitionEquals0 .
\end{equation}
A sequence \((F_r)_{r\ge0}\) is \emph{Cauchy} if for every \(N\) there is
\(R\) with \(\ordt(F_r-F_s)\ge N\) for all \(r,s\ge R\); it \emph{converges}
to \(F\) if \(\ordt(F_r-F)\to+\infty\).  We say \((F_r)\) \emph{stabilizes
coefficientwise} if for every \(n\) the block \([t^n]F_r\) is eventually
independent of \(r\).
\end{definition}

That \eqref{plt:eq:trunc-metric} is an ultrametric --- in particular
\(d(F,H)\le\max\SetOf{d(F,G),d(G,H)}\) --- is
\cref{plt:thm:ring-valuation-laws}\,(ii) restated, and \(d(F,G)=0\) iff
\(F=G\) is \cref{plt:thm:ring-valuation-laws}\,(iii).

\begin{theorem}[Completeness and separatedness]
\label{plt:thm:trunc-completeness}
\begin{enumerate}[label=\textup{(\roman*)}]
\item \emph{Separatedness.}  \(\bigcap_{N\ge0}t^N\PLpow=0\) in \(\PLpow\), and
  \(\bigcap_{N\in\IntegerNumbers}\mathfrak F_N=0\) in \(\PLlau\).
\item \emph{Cauchy \(=\) stabilizing truncations.}  \((F_r)\) is Cauchy if and
  only if for every \(N\) the truncations \(\Trunc{F_r}{N}\) are eventually
  independent of \(r\).
\item \emph{Completeness.}  \(\PLpow\) and \(\PLlau\) are complete: every
  Cauchy sequence converges, to a unique limit.  The same holds for
  \(E[[t]]\) and \(E((t))\) for any coefficient ring \(E\); in particular for
  \(\PLrat\) and \(\PLit\).
\item \emph{Coefficientwise stabilization is weaker on the Laurent ring.}
  Every Cauchy sequence stabilizes coefficientwise.  On \(\PLpow\) the
  converse holds.  On \(\PLlau\) it fails; it is restored exactly under the
  additional hypothesis that \(\inf_r\ordt(F_r)>-\infty\).
\end{enumerate}
\end{theorem}

\begin{proof}
(i)  If \(F\in\mathfrak F_N\) for every \(N\) then \(\ordt F\ge N\) for every
\(N\) by \cref{plt:prop:trunc-not-ideal}\,(iii), so \(\ordt F=+\infty\) and
\(F=0\).

(ii)  Immediate from \cref{plt:prop:trunc-basic}\,(iii): the condition
\(\ordt(F_r-F_s)\ge N\) is exactly \(\Trunc{F_r}{N}=\Trunc{F_s}{N}\).

(iii)  Let \((F_r)\) be Cauchy.  Apply (ii) with \(N=0\): there is \(R_0\)
with \(\Trunc{F_r}{0}=\Trunc{F_{R_0}}{0}\) for \(r\ge R_0\).  Set
\[
  n_0=\min\SetOf{0,\ \ordt\bigl(\Trunc{F_{R_0}}{0}\bigr)}\in\IntegerNumbers ,
\]
a genuine integer (equal to \(0\) if the truncation vanishes).  For
\(r\ge R_0\) all blocks of \(F_r\) below \(n_0\) vanish, since they agree with
those of \(\Trunc{F_{R_0}}{0}\) below \(0\) and there are none below \(n_0\);
hence \(\ordt(F_r)\ge n_0\).  This is the step at which the Cauchy hypothesis,
and not merely coefficientwise stabilization, is used: it supplies a
\emph{common} lower support bound.

Next, for each \(n\ge n_0\), choosing \(N=n+1\) in (ii) gives \(R_n\) with
\([t^n]F_r=[t^n]F_{R_n}=:p_n\) for \(r\ge R_n\).  Since \(\ordt F_r\ge n_0\)
for \(r\ge R_0\), we have \(p_n=0\) for \(n<n_0\), so
\(F\DefinitionEquals\sum_{n\ge n_0}p_nt^n\) is a well-defined element of
\(\PLlau\) (and of \(\PLpow\) if \(n_0\ge0\), which is the case when all
\(F_r\in\PLpow\)).  Given \(N\ge n_0\), take \(R=\max\SetOf{R_0,R_{n_0},\dots,
R_{N-1}}\); then \(\Trunc{F_r}{N}=\Trunc{F}{N}\) for \(r\ge R\), i.e.\
\(\ordt(F_r-F)\ge N\).  Hence \(F_r\to F\).  Uniqueness of the limit is (i).
The argument used nothing about \(\K[L]\) beyond its being a ring, so it
applies verbatim to \(E[[t]]\) and \(E((t))\).
% ed.: S2's completeness proof asserts the common lower bound because ``the
% sequence is already Laurent-bounded at some initial precision''.  That is not
% a hypothesis of the statement; the bound must be extracted, as above, from
% the Cauchy condition at cutoff \(N=0\).

(iv)  If \(F_r\to F\) then \([t^n]F_r=[t^n]F\) for \(r\) large, by
\cref{plt:prop:trunc-extraction}; so Cauchy sequences stabilize
coefficientwise.  On \(\PLpow\) suppose conversely that \((F_r)\) stabilizes
coefficientwise.  Given \(N\ge0\), the finitely many blocks
\(n=0,1,\dots,N-1\) all stabilize, say beyond \(R\); then
\(\Trunc{F_r}{N}\) is constant for \(r\ge R\), and (ii) gives Cauchy.  The
same argument works on \(\PLlau\) as soon as \(\ordt(F_r)\ge n_1\) for all
\(r\), because then only the finitely many blocks \(n_1\le n<N\) matter.
Failure without that hypothesis is
\cref{plt:ex:trunc-not-cauchy}.
% ed.: the word ``exactly'' in (iv) asserts that the common lower bound is also
% necessary; the sources supply only sufficiency, so the converse is proved
% here.
That the hypothesis is not merely sufficient but necessary --- which is what
``exactly'' asserts --- follows from (iii): a Cauchy sequence satisfies
\(\ordt(F_r)\ge n_0\) for all \(r\ge R_0\), and the finitely many terms with
\(r<R_0\) each have \(\ordt(F_r)\in\IntegerNumbers\cup\SetOf{+\infty}\), so
\(\inf_r\ordt(F_r)>-\infty\).
\end{proof}

\begin{example}[Coefficientwise stabilization without convergence]
\label{plt:ex:trunc-not-cauchy}
% ed.: S3 asserts that \(t\)-adic convergence is ``equivalent'' to
% coefficientwise stabilization and that both the power-series and the Laurent
% ring are complete ``in this sense''.  The equivalence is false on the Laurent
% ring, as follows; the completeness assertion itself is correct, but needs the
% proof of \cref{plt:thm:trunc-completeness}\,(iii).
Let
\[
  F_r=\sum_{n=-r}^{-1}t^n\in\PLlau
  \qquad(r\ge0),
\]
so \(F_0=0\).  For each fixed \(n\) the block \([t^n]F_r\) equals \(1\) for
all \(r\ge-n\) when \(n<0\), and \(0\) for all \(r\) when \(n\ge0\): the
sequence stabilizes coefficientwise.  But
\(\ordt(F_{r+1}-F_r)=-(r+1)\to-\infty\), so \((F_r)\) is not Cauchy, and there
is no \(F\in\PLlau\) with \([t^n]F=1\) for every \(n<0\), since such an \(F\)
would have infinitely many negative blocks.  The obstruction is exactly the
missing common lower bound: \(\inf_r\ordt(F_r)=-\infty\).
\end{example}

\begin{remark}[Separatedness is what makes residuals certificates]
\label{plt:rmk:trunc-separated}
Separatedness says that an element with \(\ordt=+\infty\) is \(0\).  Every
verification scheme in this volume rests on it: to prove \(F=G\) it suffices
to prove \(\Trunc{(F-G)}{N}=0\) for every \(N\), which is a sequence of finite
computations.  Completeness says the converse construction is available: a
consistent prescription of blocks defines an element.  Together they are the
formal substitute for ``the iteration converges''.
\end{remark}

\section{Formal summability}
\label{plt:sec:trunc-summability}

\begin{definition}[Formally summable family]
\label{plt:def:trunc-summable}
Let \(I\) be an index set.  A family \((S_i)_{i\in I}\) in \(\PLlau\) is
\emph{formally summable} if
\begin{equation}\label{plt:eq:trunc-summable}
  \text{for every }N\in\IntegerNumbers,\quad
  \SetOf{i\in I:\ordt(S_i)<N}\ \text{is finite.}
\end{equation}
For a sequence indexed by \(\NaturalNumbers\) this is the condition
\(\ordt(S_r)\to+\infty\).  Its sum is defined blockwise by
\begin{equation}\label{plt:eq:trunc-sum}
  [t^n]\sum_{i\in I}S_i\DefinitionEquals\sum_{i\in I}[t^n]S_i ,
\end{equation}
in which all but finitely many terms vanish.
\end{definition}

\begin{theorem}[The summability calculus]
\label{plt:thm:trunc-summability}
Let \((S_i)_{i\in I}\) be formally summable in \(\PLlau\).  Then:
\begin{enumerate}[label=\textup{(\roman*)}]
\item \eqref{plt:eq:trunc-sum} defines an element
  \(S=\sum_{i\in I}S_i\in\PLlau\), with
  \(\ordt(S)\ge\inf_i\ordt(S_i)>-\infty\);
\item for every \(N\), \(\Trunc{S}{N}=\sum_{i\in I_N}\Trunc{S_i}{N}\) where
  \(I_N=\SetOf{i:\ordt(S_i)<N}\) is finite; equivalently \(S\) is the limit of
  the net of finite partial sums, and for \(I=\NaturalNumbers\),
  \(\sum_{r\le R}S_r\to S\) in the \(t\)-adic topology;
\item the sum does not depend on the enumeration of \(I\): for any bijection
  \(\pi:I\to I\), \(\sum_iS_{\pi(i)}=S\);
\item summable families form a \(\K\)-module and \(\sum\) is \(\K\)-linear on
  it;
\item for \(H\in\PLlau\), the family \((HS_i)\) is summable and
  \(\sum_i HS_i=H\sum_iS_i\);
\item if \((T_k)_{k\in J}\) is also summable, then \((S_iT_k)_{(i,k)\in I\times
  J}\) is summable and
  \(\sum_{(i,k)}S_iT_k=\bigl(\sum_iS_i\bigr)\bigl(\sum_kT_k\bigr)\).
\end{enumerate}
\end{theorem}

\begin{proof}
(i)  Fix \(n\).  Terms with \([t^n]S_i\ne0\) satisfy \(\ordt S_i\le n\), and
by \eqref{plt:eq:trunc-summable} with \(N=n+1\) there are finitely many such
\(i\); so \eqref{plt:eq:trunc-sum} is a finite sum in \(\K[L]\) and defines a
block \(P_n\).  Applying \eqref{plt:eq:trunc-summable} with \(N=0\) shows that
% ed.: the identification of \(\mu\) was printed as an equality.  It is only a
% lower bound: if every \(\ordt S_i>0\) the right-hand side is \(0\) while
% \(\mu>0\).  Finiteness is all the argument needs, and all it gives.
\(\SetOf{i:\ordt S_i<0}\) is finite, hence
\[
  \mu\DefinitionEquals\inf_i\ordt(S_i)
  \ \ge\ \min\SetOf{0,\ \min_{i:\ordt S_i<0}\ordt S_i}\ >\ -\infty ,
\]
with the convention that the inner minimum is \(0\) when its index set is
empty.  Thus \(\mu\in\IntegerNumbers\) unless every \(S_i=0\), in which case
\(\mu=+\infty\).  For \(n<\mu\) every
\([t^n]S_i\) vanishes, so \(P_n=0\); thus \(S=\sum_{n\ge\mu}P_nt^n\in\PLlau\)
and \(\ordt S\ge\mu\).

(ii)  For \(n<N\) and \(i\notin I_N\) we have \(\ordt S_i\ge N>n\), so
\([t^n]S_i=0\); hence \([t^n]S=\sum_{i\in I_N}[t^n]S_i\) for all \(n<N\),
which is the displayed identity.  Consequently, for the partial sums
\(\Sigma_R=\sum_{r\le R}S_r\) with \(I=\NaturalNumbers\), choosing \(R\) with
\(I_N\subseteq\SetOf{0,\dots,R}\) gives \(\Trunc{\Sigma_{R'}}{N}=\Trunc{S}{N}\)
for all \(R'\ge R\), i.e.\ \(\ordt(\Sigma_{R'}-S)\ge N\).

(iii)  The defining formula \eqref{plt:eq:trunc-sum} is a sum over a finite
subset of \(I\) determined by \(n\) alone, and a finite sum in \(\K[L]\) is
independent of order; and \(\pi\) preserves the summability condition.

(iv)  If \((S_i)\) and \((S'_i)\) are summable then
\(\ordt(cS_i+S'_i)\ge\min\SetOf{\ordt S_i,\ordt S'_i}\), so
\(\SetOf{i:\ordt(cS_i+S'_i)<N}\subseteq\SetOf{i:\ordt S_i<N}\cup
\SetOf{i:\ordt S'_i<N}\) is finite.  Linearity of \(\sum\) is linearity of
\eqref{plt:eq:trunc-sum} blockwise.

(v)  \(\ordt(HS_i)=\ordt H+\ordt S_i\), so
\(\SetOf{i:\ordt(HS_i)<N}=\SetOf{i:\ordt S_i<N-\ordt H}\) is finite (if
\(H=0\) the claim is trivial).  For the identity, fix \(n\) and compute
\[
  [t^n]\sum_iHS_i
  =\sum_i\sum_{a+b=n}[t^a]H\,[t^b]S_i
  =\sum_{a+b=n}[t^a]H\sum_i[t^b]S_i
  =[t^n]\Bigl(H\sum_iS_i\Bigr),
\]
where the interchange is legitimate because every sum involved is finite: the
outer index \(a\) ranges over \(\ordt H\le a\le n-\mu\), and for each \(b\)
only finitely many \(i\) contribute.

(vi)  \(\ordt(S_iT_k)=\ordt S_i+\ordt T_k\).  Let \(\mu_S=\inf_i\ordt S_i\)
and \(\mu_T=\inf_k\ordt T_k\), both integers by (i).  If
\(\ordt S_i+\ordt T_k<N\), then \(\ordt S_i<N-\mu_T\) and
\(\ordt T_k<N-\mu_S\); so
\[
  \SetOf{(i,k):\ordt(S_iT_k)<N}
  \subseteq
  \SetOf{i:\ordt S_i<N-\mu_T}\times\SetOf{k:\ordt T_k<N-\mu_S},
\]
a product of two finite sets.  For the value, apply (v) with \(H=S_i\) to get
\(\sum_kS_iT_k=S_i\sum_kT_k\) for each \(i\); then
\cref{plt:prop:trunc-fubini}, whose proof uses only (i) and (ii) above, and
(v) again with \(H=\sum_kT_k\), give
\(\sum_{(i,k)}S_iT_k=\sum_i\bigl(S_i\sum_kT_k\bigr)
=\bigl(\sum_iS_i\bigr)\bigl(\sum_kT_k\bigr)\).
\end{proof}

\begin{corollary}[Geometric summability]
\label{plt:cor:ring-geometric}
% ed.: the corollary was stated only for \(E\in\SetOf{\K[L],\K(L),\K((L^{-1}))}\),
% yet it is invoked inside \cref{plt:prop:ring-chain} with \(E=\K\) and
% generator \(L^{-1}\), precisely in order to prove that \(\K((L^{-1}))\) is a
% field --- the fact that puts \(\K((L^{-1}))\) on that list.  Stating it for an
% arbitrary commutative coefficient ring removes the circle at no cost: the
% proof never needs a domain hypothesis.
Let \(E\) be a commutative ring with \(1\) and let \(U\in tE[[t]]\), i.e.\
\(\ordt U\ge1\).  Then \((U^r)_{r\ge0}\) is formally summable in \(E((t))\),
\(\sum_{r\ge0}U^r\in E[[t]]\), and
\[
  (1-U)\sum_{r\ge0}U^r=1 .
\]
This applies in particular to \(E=\K\), \(\K[L]\), \(\K(L)\) and
\(\K((L^{-1}))\), and --- reading \(L^{-1}\) for \(t\) --- inside
\(\K[[L^{-1}]]\subset\K((L^{-1}))\).
\end{corollary}

\begin{proof}
By \cref{plt:rmk:trunc-summability-classes} the summability calculus of
\cref{plt:thm:trunc-summability} is available in \(E((t))\) for every
commutative ring \(E\).  The Cauchy product gives the subadditivity
\(\ordt(FG)\ge\ordt F+\ordt G\) there --- this is the first half of the
computation in the proof of \cref{plt:thm:ring-valuation-laws}\,(i), which
needs no domain hypothesis --- so
\(\ordt(U^r)\ge r\,\ordt U\ge r\to+\infty\) and
\eqref{plt:eq:trunc-summable} holds.  Write \(V=\sum_{r\ge0}U^r\).  By
\cref{plt:thm:trunc-summability}\,(iv),(v),
\(V-UV=\sum_{r\ge0}U^r-\sum_{r\ge0}U^{r+1}
=\sum_{r\ge0}\bigl(U^r-U^{r+1}\bigr)\).  Fix \(N\ge1\); for \(N\le0\) both
\(\Trunc{(V-UV)}{N}\) and \(\Trunc{1}{N}\) vanish, since \(V-UV\) and \(1\)
lie in \(E[[t]]\).  Only the indices \(r<N\) contribute to
\(\Trunc{\cdot}{N}\), because \(\ordt(U^r-U^{r+1})\ge r\)
(\cref{plt:thm:trunc-summability}\,(ii)), and the resulting finite sum
telescopes:
\[
  \Trunc{(V-UV)}{N}
  =\Trunc{\Bigl(\sum_{0\le r<N}\bigl(U^r-U^{r+1}\bigr)\Bigr)}{N}
  =\Trunc{(1-U^{N})}{N}
  =\Trunc{1}{N},
\]
the last step because \(\ordt(U^N)\ge N\).  Since this holds for every \(N\),
\cref{plt:prop:trunc-basic}\,(iv) gives \(V-UV=1\).
\end{proof}

\begin{theorem}[Three formulations of summability]
\label{plt:thm:trunc-summability-criterion}
Let \((S_i)_{i\in I}\) be a family in \(\PLlau\).  Consider:
\begin{enumerate}[label=\textup{(\alph*)}]
\item \emph{Valuation form.}  For every \(N\), \(\SetOf{i:\ordt S_i<N}\) is
  finite \textup{(}\cref{plt:def:trunc-summable}\textup{)}.
\item \emph{Hahn form.}  \(\bigcup_i\CoefficientSupportOf{S_i}\) satisfies
  \cref{plt:thm:ring-support}\,(b), and every monomial \(t^nL^j\) lies in
  \(\CoefficientSupportOf{S_i}\) for only finitely many \(i\).
\item \emph{Block-local finiteness with a common bound.}  For every \(n\) the
  set \(\SetOf{i:[t^n]S_i\ne0}\) is finite, and
  \(\inf_i\ordt(S_i)>-\infty\).
\end{enumerate}
Then \textup{(a)}, \textup{(b)} and \textup{(c)} are equivalent.  If the family
lies in \(\PLpow\), the lower-bound clause in \textup{(c)} is automatic and may
be dropped.  On \(\PLlau\) it may not: block-local finiteness alone does not
imply summability.
\end{theorem}

\begin{proof}
\emph{(a)\(\Rightarrow\)(c).}  \(\SetOf{i:[t^n]S_i\ne0}\subseteq
\SetOf{i:\ordt S_i\le n}\), finite by (a) with \(N=n+1\); and the lower bound
was produced in the proof of
\cref{plt:thm:trunc-summability}\,(i).

\emph{(c)\(\Rightarrow\)(a).}  Let \(\mu=\inf_i\ordt S_i\in\IntegerNumbers\).
For \(N\le\mu\) the set in (a) is empty.  For \(N>\mu\),
\[
  \SetOf{i:\ordt S_i<N}
  \subseteq\bigcup_{n=\mu}^{N-1}\SetOf{i:[t^n]S_i\ne0},
\]
a finite union of finite sets.

\emph{(a)\(\Rightarrow\)(b).}  The union of supports has \(n\)-projection
bounded below by \(\mu\).  Its fiber over \(n\) is the union of the fibers of
the finitely many \(S_i\) with \(\ordt S_i\le n\), each finite, hence finite.
The monomial condition follows from
\(\SetOf{i:(n,j)\in\CoefficientSupportOf{S_i}}\subseteq\SetOf{i:\ordt S_i\le
n}\).

\emph{(b)\(\Rightarrow\)(a).}  Let \(n_0\) bound the \(n\)-projection of the
union from below and let \(U_n\) be its (finite) fiber over \(n\).  Given
\(N\),
\[
  \SetOf{i:\ordt S_i<N}
  \subseteq\bigcup_{n=n_0}^{N-1}\ \bigcup_{j\in U_n}
    \SetOf{i:(n,j)\in\CoefficientSupportOf{S_i}},
\]
a finite union --- finitely many \(n\), finitely many \(j\) for each --- of
finite sets.

\emph{The \(\PLpow\) case.}  There \(\ordt S_i\ge0\) for all \(i\), so
\(\inf_i\ordt S_i\ge0\).
\end{proof}

\begin{example}[Two families that just fail]
\label{plt:ex:trunc-summability-fails}
\emph{Local finiteness without a common lower bound.}  Let \(S_r=t^{-r}\),
\(r\ge0\), in \(\PLlau\).  Every block index \(n\) is hit by at most one
member, so block-local finiteness holds; but \(\ordt S_r=-r\to-\infty\), the
family is not summable, and the would-be sum \(\sum_{n\le0}t^{n}\) has an
infinite principal part and lies in no Laurent class.  This is exactly the
extra hypothesis that \cref{plt:thm:trunc-summability-criterion}\,(c) adds on
\(\PLlau\).

\emph{Monomial finiteness without block finiteness.}  Let \(S_r=tL^r\),
\(r\ge0\), in \(\PLpow\).  Each monomial \(tL^r\) occurs in exactly one
member, yet \(\ordt S_r=1\) for every \(r\), the family is not summable, and
the would-be block \(\sum_{r\ge0}L^r\) is not a polynomial.  Note that the
union of supports, \(\SetOf{(1,r):r\ge0\,}\), has an infinite fiber and so
violates the Hahn clause of
\cref{plt:thm:trunc-summability-criterion}\,(b).  Requiring only that each
\emph{monomial} occur finitely often is therefore not enough, even on
\(\PLpow\); the well-based clause is doing real work.
\end{example}

\begin{proposition}[Fubini for summable families]
\label{plt:prop:trunc-fubini}
Let \((S_{i,k})_{(i,k)\in I\times J}\) be formally summable.  Then for each
\(i\) the family \((S_{i,k})_{k\in J}\) is summable with sum \(T_i\); the
family \((T_i)_{i\in I}\) is summable; and
\[
  \sum_{i\in I}T_i=\sum_{(i,k)\in I\times J}S_{i,k}.
\]
The same holds with the roles of \(I\) and \(J\) exchanged, so iterated
summation may be performed in either order.
\end{proposition}

\begin{proof}
Summability of each slice is inherited, since
\(\SetOf{k:\ordt S_{i,k}<N}\) injects into
\(\SetOf{(i',k):\ordt S_{i',k}<N}\).  If \(\ordt S_{i,k}\ge N\) for all
\(k\in J\), then every block of \(T_i\) below \(N\) vanishes by
\eqref{plt:eq:trunc-sum}, so \(\ordt T_i\ge N\); hence
\(\SetOf{i:\ordt T_i<N}\subseteq\SetOf{i:\exists k,\ \ordt S_{i,k}<N}\), which
is finite.  Finally, for each \(n\),
\[
  [t^n]\sum_iT_i=\sum_i[t^n]T_i=\sum_i\sum_k[t^n]S_{i,k}
  =\sum_{(i,k)}[t^n]S_{i,k}=[t^n]\sum_{(i,k)}S_{i,k},
\]
all sums being finite, so ordinary finite-sum associativity applies.
\end{proof}

\begin{remark}[Summability on the other three classes]
\label{plt:rmk:trunc-summability-classes}
% ed.: the transfer was claimed only for the four ambient classes, but
% \cref{plt:cor:ring-geometric} needs it for \(E((t))\) with \(E\) an arbitrary
% commutative ring (the case \(E=\K\), generator \(L^{-1}\), is used in
% \cref{plt:prop:ring-chain}).  The scope of the transfer is widened, and the
% two places where the domain hypothesis entered are named.
\Cref{plt:def:trunc-summable} and
\cref{plt:thm:trunc-summability} were stated for \(\PLlau\) but use only the
valuation laws and the fact that each block lies in a fixed additive group;
they hold verbatim in \(E[[t]]\) and \(E((t))\) for every commutative ring
\(E\) --- in particular in \(\PLpow\), \(\PLrat\) and \(\PLit\).  The domain
hypothesis on \(E\) was used only for the equalities
\(\ordt(HS_i)=\ordt H+\ordt S_i\) and
\(\ordt(S_iT_k)=\ordt S_i+\ordt T_k\) in parts (v) and (vi) of that theorem,
and there the inequality \(\ge\) --- valid over any commutative ring --- is
all that the finiteness arguments need, the displayed set identities becoming
inclusions.  In \(\PLit\),
however, \eqref{plt:eq:trunc-sum} requires an additional check: the finitely
many contributing blocks are Laurent series in \(L^{-1}\), and their sum is
again one, so nothing is lost.  What does \emph{not} transfer is the
\emph{monomialwise} version: in \(\PLit\) a single block already involves
infinitely many monomials, so summability must be tested on blocks, never on
monomials.
\end{remark}

\section{The three \texorpdfstring{\(O\)}{O}-roles}
\label{plt:sec:trunc-three-o}

\begin{warningbox}[title={Never infer one \(O\)-role from another}]
% ed.: the label kind "box" is not in the normative list of kinds; these
% boxed statements are remarks, so each label is renamed to the "rmk" kind.
\label{plt:rmk:trunc-three-o}
Three unrelated meanings share the letter \(O\) in this subject, and the
corpus keeps them typographically distinct.
\end{warningbox}

\begin{center}
\begin{tabular}{@{}lp{0.30\textwidth}p{0.34\textwidth}@{}}
\hline
Printed form & Means & Does \emph{not} mean \\
\hline
\(\FormalBigOAt{t^N}{t}\)
 & formal valuation tail: \(\ordt(F-G)\ge N\)
   (\cref{plt:def:trunc-formal-o})
 & any numerical bound, rate, uniformity in \(L\), convergence, or the
   existence of a function \\
% ed.: the volume regime is \(X\to+\infty\) on the positive ray; the generic
% row printed the weaker \(X\to\infty\).
\(\BigOAt{A(X)}{X\to+\infty,\ X\in S}\)
 & analytic estimate: \(\AbsoluteValue{R(X)}\le C\AbsoluteValue{A(X)}\) on a
   stated late part of a stated domain, with branches fixed
 & anything formal; and it is void unless the regime and domain are printed \\
\(\BigO(N^2)\)
 & algorithmic cost in a declared model (coefficient-ring operations, bit
   operations, \dots)
 & a statement about the size or accuracy of the computed object \\
\hline
\end{tabular}
\end{center}

\begin{remark}[The bridge, and its price]
\label{plt:rmk:trunc-bridge}
Passing from the first column to the second is a theorem, never a notation
change.  It needs, at minimum: a function \(f\) on a stated domain; a proof
that \(f\) has the blockwise expansion in the sense of
\cref{plt:rmk:trunc-blockwise}; and, for each cutoff \(N\), an exponent
\(M_N\) bounding the logarithmic degree of the discarded material.  A bare
\(\BigO\) with no printed regime is never acceptable for an asymptotic claim.
\end{remark}

\begin{summarybox}
The formal ambient is \(\PLlau=\K[L]((t))\): supports are lattice sets with
finite fibers and lower-bounded projection, hence reverse well ordered of
order type at most \(\omega\).  The outer valuation \(\ordt\) is the
computational filtration; the rank-two valuation
\(v(t^nL^j)=(n,-j)\) is the monomial comparison, and \(\PLit\) is exactly the
Hahn field of its value group.  \(\PLpow\) is complete and separated and
\(t^N\PLpow\) is a genuine ideal there; in \(\PLlau\) the same set is only a
filtration tail, because \(t\) is a unit.  All infinite formulas in the
volume are legitimated by one criterion: a family is summable when
\(\ordt(S_i)\to+\infty\), which on the Laurent ring is strictly stronger than
block-local finiteness.  Leading block, leading monomial and leading scalar
are three different objects; log-degree bounds are always hypotheses or proved
invariants; and a formal tail carries no size information.
\end{summarybox}

\clearpage
\part{Arithmetic and differential calculus}

% ---- fragment 03-multiplication ----
\chapter{Multiplication: convolution of logarithmic blocks}
\label{plt:sec:multiplication}

Multiplication is the operation on which every later construction rests:
reciprocals, powers, composition, and reversion are all built from it by
triangular recurrences.  Its whole content is that the Cauchy rule in the
outer variable~$t$ is \emph{coefficientwise finite}, so that each output
block is an exact finite expression in finitely many input blocks.  Nothing
here is an analytic rearrangement, and nothing here requires convergence.

Throughout this chapter $\K$ is a field of characteristic zero, $X$ is the
large variable, the regime is $X\to+\infty$ on the positive ray unless a
different domain is announced, and
\begin{equation}\label{plt:eq:mul-generators}
 t\DefinitionEquals X^{-1},
 \qquad
 L\DefinitionEquals\log X .
\end{equation}
In the formal algebra $t$ and $L$ are independent generators; their analytic
coupling is recorded only by the distinguished derivation, which plays no
role in this chapter.  The four ambient classes are

\begin{center}
\begin{tabular}{@{}lll@{}}
\hline
class & coefficient ring & elements \\
\hline
$\PLpow$ & $\K[L]$ & $\sum_{n\ge0}p_n(L)t^n$ \\
$\PLlau$ & $\K[L]$ & $\sum_{n\ge n_0}p_n(L)t^n$, $n_0\in\IntegerNumbers$ \\
$\PLrat$ & $\K(L)$ & $\sum_{n\ge n_0}p_n(L)t^n$, $p_n$ rational in $L$ \\
$\PLit$  & $\K((L^{-1}))$ & $\sum_{n\ge n_0}t^n\sum_{j\ge j_n}a_{n,j}L^{-j}$ \\
\hline
\end{tabular}
\end{center}

All truncations below are \emph{exclusive}: $\Trunc{F}{N}$ keeps every
logarithmic monomial in every block $t^n$ with $n<N$, and discards the rest.

\begin{remark}[Reconciliation of the source conventions]
\label{plt:rmk:mul-conventions}
Two of the merged reports use an inclusive outer cutoff $[F]_{\le N}$.  The
translation is $[F]_{\le N}=\Trunc{F}{N+1}$, and every dependency range,
algorithm bound and error exponent quoted from those reports has been shifted
by one block accordingly.  Likewise, one report writes the coefficient
bracket as $[F]_M$; here the series always stands to the right of the
bracket, $[t^n]F=p_n(L)$ and $[t^nL^j]F=a_{n,j}$.
\end{remark}

\section{The block Cauchy product}
\label{plt:sec:mul-cauchy}

It is economical to prove the algebra of all four classes at once.  What the
three coefficient rings $\K[L]$, $\K(L)$, $\K((L^{-1}))$ have in common is a
degree function in $L$ together with a coefficient-extraction map obeying the
usual convolution rule.

\begin{definition}[Logarithmically graded domain]
\label{plt:def:mul-graded-domain}
A \emph{logarithmically graded domain over $\K$} is a commutative
$\K$-algebra $R$ together with $\K$-linear maps
$\kappa_d\colon R\to\K$, one for each $d\in\IntegerNumbers$, such that
for every nonzero $f\in R$ the set $\SetOf{d:\kappa_d(f)\ne0}$ is nonempty
and bounded above, and:
\begin{enumerate}
\item[(G1)] $f=0$ if and only if $\kappa_d(f)=0$ for every $d$;
\item[(G2)] $\kappa_d(fg)=\sum_{i+j=d}\kappa_i(f)\kappa_j(g)$ for all
$f,g\in R$ and all $d\in\IntegerNumbers$, the sum having only finitely many
nonzero terms;
\item[(G3)] $\kappa_0(1)=1$ and $\kappa_d(1)=0$ for $d\ne0$.
\end{enumerate}
For $f\ne0$ put $\degL f\DefinitionEquals\max\SetOf{d:\kappa_d(f)\ne0}$ and
$\lc_L(f)\DefinitionEquals\kappa_{\degL f}(f)\in\K^\times$; set
$\degL 0\DefinitionEquals-\infty$ and $\lc_L(0)\DefinitionEquals0$.
\end{definition}

Here $\lc_L$ is the leading coefficient \emph{inside one block}; it is not
the same object as the leading scalar $\lc$ of a series, defined in
\cref{plt:def:mul-leading-data} below, although the two agree on the leading
block.

\begin{lemma}[Elementary consequences]
\label{plt:lem:mul-graded-basics}
Let $R$ be a logarithmically graded domain over $\K$.  For all $f,g\in R$:
\begin{equation}\label{plt:eq:mul-degL-laws}
 \degL(fg)=\degL f+\degL g,
 \qquad
 \lc_L(fg)=\lc_L(f)\lc_L(g),
 \qquad
 \degL(f+g)\le\max(\degL f,\degL g).
\end{equation}
Here $-\infty+c=-\infty$ and $\max(-\infty,c)=c$.  In particular $R$ is an
integral domain.  Moreover, if $\degL f\le d$ and $\degL g\le e$, then
$\kappa_{d+e}(fg)=\kappa_d(f)\kappa_e(g)$.
\end{lemma}

\begin{proof}
If $f=0$ or $g=0$ every assertion is trivial under the stated conventions.
The third assertion in \eqref{plt:eq:mul-degL-laws} is immediate from
$\K$-linearity of each $\kappa_d$.  For the first two, assume $f,g\ne0$ and
set $d=\degL f$, $e=\degL g$.  If $c>d+e$ and $i+j=c$, then $i>d$ or $j>e$, so
every term of (G2) vanishes and $\kappa_c(fg)=0$.  If $c=d+e$, the only
possibly nonvanishing term is $i=d$, $j=e$, whence
\[
 \kappa_{d+e}(fg)=\kappa_d(f)\kappa_e(g)=\lc_L(f)\lc_L(g)\ne0,
\]
because $\K$ is a field and both factors are nonzero.  This proves
$\degL(fg)=d+e$, $\lc_L(fg)=\lc_L(f)\lc_L(g)$, and, by (G1), $fg\ne0$; so $R$
is a domain.  The last sentence is the same computation with $d,e$ merely
upper bounds, all other terms of (G2) again vanishing.
\end{proof}

\begin{lemma}[The three coefficient rings]
\label{plt:lem:mul-three-rings}
Each of
\[
 \K[L],\qquad \K(L),\qquad \K((L^{-1}))
\]
is a logarithmically graded domain over $\K$, with $\degL$ equal
respectively to the polynomial degree, to $\degL(p/q)=\deg p-\deg q$, and to
the largest exponent occurring in the Laurent expansion.  There are injective
$\K$-algebra homomorphisms
\begin{equation}\label{plt:eq:mul-coeff-chain}
 \K[L]\hookrightarrow\K(L)\hookrightarrow\K((L^{-1}))
\end{equation}
that preserve $\degL$ and $\lc_L$.
\end{lemma}

\begin{proof}
For $\K[L]$ take $\kappa_d(f)$ to be the coefficient of $L^d$ in $f$ for
$d\ge0$ and $\kappa_d=0$ for $d<0$; (G1)--(G3) are the definition of
polynomial multiplication.  For $\K((L^{-1}))$, an element is
$f=\sum_{j\le d}a_jL^j$ with only finitely many $j>0$ allowed but arbitrarily
negative $j$ permitted; put $\kappa_j(f)=a_j$.  Then (G1) and (G3) are clear,
and in (G2) the constraint $i+j=c$ with $i\le\degL f$, $j\le\degL g$ forces
$i\ge c-\degL g$, so the sum is finite; it is the definition of the Laurent
product.  The support of $f$ is bounded above by construction.

The inclusion $\K[L]\subset\K((L^{-1}))$ is a $\K$-algebra map, and every
nonzero element of $\K((L^{-1}))$ is invertible: writing $f\ne0$ as
$f=a_dL^d(1+u)$ with $d=\degL f$ and $u\in L^{-1}\K[[L^{-1}]]$, the family
$\bigl((-u)^r\bigr)_{r\ge0}$ is $L^{-1}$-adically summable and its sum
inverts $1+u$, so $f^{-1}=a_d^{-1}L^{-d}\sum_{r\ge0}(-u)^r$.  Hence
every nonzero polynomial becomes invertible, and the universal property of the
field of fractions supplies a unique $\K$-algebra homomorphism
$\iota\colon\K(L)\to\K((L^{-1}))$ extending the inclusion; it is injective
because $\K(L)$ is a field and $\iota(1)=1$.  Pulling the maps $\kappa_d$ back
along $\iota$ makes $\K(L)$ a logarithmically graded domain, and
\cref{plt:lem:mul-graded-basics} applied to $\iota(p)=\iota(q)\iota(p/q)$
gives $\degL(p/q)=\deg p-\deg q$ and
$\lc_L(p/q)=\lc_L(p)\lc_L(q)^{-1}$.  Both maps in
\eqref{plt:eq:mul-coeff-chain} preserve $\kappa_d$ by construction, hence
preserve $\degL$ and $\lc_L$.
\end{proof}

\begin{definition}[The block Cauchy product]
\label{plt:def:mul-cauchy}
Let $R$ be a logarithmically graded domain over $\K$ and let
\begin{equation}\label{plt:eq:mul-FG-blocks}
 F=\sum_{n\ge a}p_n\,t^n,
 \qquad
 G=\sum_{m\ge b}q_m\,t^m,
 \qquad p_n,q_m\in R,
\end{equation}
be elements of $R((t))$, the ring of formal Laurent series in $t$ over $R$.
Their product is
\begin{equation}\label{plt:eq:mul-cauchy}
 FG\DefinitionEquals\sum_{N\ge a+b}C_N\,t^N,
 \qquad
 C_N\DefinitionEquals\sum_{\substack{n+m=N\\ n\ge a,\ m\ge b}}p_nq_m
 =\sum_{n=a}^{N-b}p_n\,q_{N-n}.
\end{equation}
\end{definition}

\begin{lemma}[Finiteness of every fiber]
\label{plt:lem:mul-finite-fiber}
For fixed $N$ the index set in \eqref{plt:eq:mul-cauchy} is
$\SetOf{n\in\IntegerNumbers:a\le n\le N-b}$, of cardinality
$\max(N-a-b+1,\,0)$.  In particular it is finite, and it is empty exactly
when $N<a+b$.
\end{lemma}

\begin{proof}
The constraints $n\ge a$ and $m=N-n\ge b$ are equivalent to
$a\le n\le N-b$.  An interval of integers $[a,N-b]$ has $N-b-a+1$ elements
when $N-b\ge a$ and is empty otherwise.
\end{proof}

This is the entire mechanism.  There is no limit, no rearrangement, and no
hypothesis of numerical convergence: each $C_N$ is a finite sum of products
in~$R$.

\begin{theorem}[Closure of the four ambient classes]
\label{plt:thm:mul-closure}
Let $R$ be a logarithmically graded domain over $\K$.  Then
\eqref{plt:eq:mul-cauchy} makes $R((t))$ a commutative $\K$-algebra with unit
$1$, containing $R[[t]]$ as a subalgebra.  Consequently each of
% ed.: the four classes were displayed on one line, which overfills the text
% ed.: block by about 45pt; set as a two-by-two array instead.
\begin{align*}
 \PLpow&=\K[L][[t]], & \PLlau&=\K[L]((t)),\\
 \PLrat&=\K(L)((t)), & \PLit&=\K((L^{-1}))((t))
\end{align*}
is closed under multiplication, and the chain of inclusions
\begin{equation}\label{plt:eq:mul-class-chain}
 \PLpow\subset\PLlau\subset\PLrat\subset\PLit
\end{equation}
consists of injective $\K$-algebra homomorphisms.
\end{theorem}

\begin{proof}
By \cref{plt:lem:mul-finite-fiber} each $C_N$ is a well-defined element
of~$R$, and $C_N=0$ for $N<a+b$; so $FG$ again has support bounded below and
lies in $R((t))$.  Commutativity is the symmetry of the index set under
$(n,m)\mapsto(m,n)$; bilinearity over $\K$ is clear; and $1\cdot F=F$ because
the only admissible pair for $1=1\cdot t^0$ at level $N$ is $(0,N)$.  For
associativity let $H=\sum_{k\ge c}h_kt^k$.  Then
\[
 [t^N]\bigl((FG)H\bigr)
 =\sum_{r+k=N}\Bigl(\sum_{n+m=r}p_nq_m\Bigr)h_k
 =\sum_{n+m+k=N}p_nq_mh_k,
\]
the last sum ranging over the finite set of triples with $n\ge a$, $m\ge b$,
$k\ge c$: indeed $n+m+k=N$ with the three lower bounds confines each index to
a finite interval, so the regrouping is a finite rearrangement.  The same
computation with the other bracketing gives the same finite sum, whence
$(FG)H=F(GH)$.  Distributivity is likewise a finite regrouping.

If $\ordt F\ge0$ and $\ordt G\ge0$, then $C_N=0$ for $N<0$, so $R[[t]]$ is
closed under the product; it is a subalgebra.  Taking $R=\K[L]$ gives the
first two classes, $R=\K(L)$ the third, $R=\K((L^{-1}))$ the fourth, using
\cref{plt:lem:mul-three-rings}.  Finally, a $\K$-algebra homomorphism
$\varphi\colon R\to R'$ induces a $\K$-algebra homomorphism
$R((t))\to R'((t))$ by $\sum p_nt^n\mapsto\sum\varphi(p_n)t^n$, because
$\varphi$ commutes with the finite sums \eqref{plt:eq:mul-cauchy}; injectivity
is inherited blockwise.  Applying this to
\eqref{plt:eq:mul-coeff-chain} yields \eqref{plt:eq:mul-class-chain}.
\end{proof}

\begin{proposition}[The monomial-level convolution and the support bound]
\label{plt:prop:mul-monomial}
Let $F=\sum_{n,j}a_{n,j}t^nL^j$ and $G=\sum_{m,k}b_{m,k}t^mL^k$ lie in
$\PLlau$ (so $j,k\in\NaturalNumbers$, and $\PLpow$ is the case
$\ordt\ge0$) or in $\PLit$ (so $j,k\in\IntegerNumbers$, bounded above at
each fixed outer level).  Then
\begin{equation}\label{plt:eq:mul-2d-convolution}
 [t^NL^d](FG)
 =\sum_{n+m=N}\ \sum_{i+j=d}a_{n,i}b_{m,j},
\end{equation}
a finite sum, and
\begin{equation}\label{plt:eq:mul-support}
 \CoefficientSupportOf{FG}
 \subseteq
 \CoefficientSupportOf{F}+\CoefficientSupportOf{G},
\end{equation}
the Minkowski sum in $\IntegerNumbers^2$.
\end{proposition}

\begin{proof}
Formula \eqref{plt:eq:mul-2d-convolution} is (G2) applied to
\eqref{plt:eq:mul-cauchy}.  Finiteness holds because the outer sum is finite
by \cref{plt:lem:mul-finite-fiber} and, for each of its terms, the inner sum
is finite by (G2).  For \eqref{plt:eq:mul-support}: if
$(N,d)\notin\CoefficientSupportOf{F}+\CoefficientSupportOf{G}$, then every
summand of \eqref{plt:eq:mul-2d-convolution} has $a_{n,i}=0$ or $b_{m,j}=0$,
so $[t^NL^d](FG)=0$.
\end{proof}

\begin{remark}[$\PLrat$ has no monomial support of its own]
\label{plt:rmk:mul-rat-no-support}
\cref{plt:prop:mul-monomial} is deliberately not stated for $\PLrat$.  An
element of $\K(L)((t))$ has coefficient blocks that are rational functions of
$L$, and a rational function is not a (finite or infinite) combination of
powers of $L$ until one chooses an expansion point.  The coefficient support
$\CoefficientSupportOf{\cdot}$ becomes defined only after applying the
embedding $\PLrat\hookrightarrow\PLit$ of \cref{plt:thm:mul-closure}, which
expands each block at $L=\infty$ and turns the exact coefficient $1/(L+1)$
into the infinite series $L^{-1}-L^{-2}+L^{-3}-\cdots$.  Exact rational
dependence on $L$ and a monomial support in $L$ are different pieces of data;
only the second is what \eqref{plt:eq:mul-support} constrains.
\end{remark}

\begin{remark}[Why multiplication is a two-level, not a flat, operation]
\label{plt:rmk:mul-two-level}
Although \eqref{plt:eq:mul-2d-convolution} is a genuine two-dimensional
convolution, it should not be implemented as one.  The outer index is the
asymptotic hierarchy and the inner index is not: at fixed $N$ the inner sum
is bounded, whereas a cutoff imposed jointly on $|N|+|d|$ does not respect
dominance.  See \cref{plt:rmk:mul-rectangular-cutoff}.
\end{remark}

\begin{remark}[Hahn supports]
\label{plt:rmk:mul-hahn}
Nothing in \cref{plt:lem:mul-finite-fiber} uses that the exponents are
integers, only that the two supports are well ordered in the direction of
increasing smallness.  If $S,T\subseteq\Gamma$ are reverse-well-ordered
subsets of an ordered abelian group, then $S+T$ is reverse-well-ordered and
each $\gamma\in S+T$ has only finitely many decompositions $\gamma=s+t$ with
$s\in S$, $t\in T$; this is the Hahn support theorem
\cite{hahn,edgar-beginners}.  Everything in this section that does not
mention the specific value group $\IntegerNumbers$ transfers verbatim to the
Hahn power-log extension $\sum_{\alpha\in S}p_\alpha(L)t^\alpha$ and, in
particular, to the Puiseux-log extensions $\K[L]((t^{1/s}))$ used later for
fractional powers.  We do not repeat the statements.
\end{remark}

\section{Finite determination and precision bookkeeping}
\label{plt:sec:mul-finite-det}

The content of \cref{plt:lem:mul-finite-fiber} deserves to be recorded as a
dependency statement, because that is the form in which it is used: it says
exactly how much of the input must be known in order to know a given part of
the output, and it says that no more is needed and no less will do.

\begin{theorem}[Finite determination, with exact dependency range]
\label{plt:thm:mul-finite-determination}
Let $R$ be a logarithmically graded domain over $\K$, let $a,b,N\in
\IntegerNumbers$, and let $F,G\in R((t))$ satisfy $\ordt F\ge a$ and
$\ordt G\ge b$.  Then
\begin{equation}\label{plt:eq:mul-determination}
 [t^N](FG)=\sum_{n=a}^{N-b}p_n\,q_{N-n},
\end{equation}
so that $[t^N](FG)$ is a $\K$-bilinear function of the two finite tuples
\[
 (p_a,p_{a+1},\ldots,p_{N-b})
 \qquad\text{and}\qquad
 (q_b,q_{b+1},\ldots,q_{N-a}),
\]
equivalently of $\Trunc{F}{N-b+1}$ and $\Trunc{G}{N-a+1}$.  Both ranges are
exact: for every $n_0$ with $a\le n_0\le N-b$ there are $F,F'\in R((t))$ with
$\ordt\ge a$ that differ only in the block of index $n_0$, and a $G$ with
$\ordt G\ge b$, such that $[t^N](FG)\ne[t^N](F'G)$; and symmetrically in the
second argument.
\end{theorem}

\begin{proof}
Equation \eqref{plt:eq:mul-determination} and the bilinearity are
\cref{plt:def:mul-cauchy} together with \cref{plt:lem:mul-finite-fiber}; the
blocks listed are precisely those occurring in the sum.  For exactness fix
$n_0\in[a,N-b]$ and put $m_0=N-n_0$, so that $m_0\ge b$.  Take
$G=t^{m_0}$, which has $\ordt G=m_0\ge b$.  If $n_0>a$ set $F=t^a+t^{n_0}$
and $F'=t^a$; if $n_0=a$ set $F=t^a$ and $F'=2t^a$ (here $2\ne0$ because
$\operatorname{char}\K=0$).  In both cases $\ordt F=\ordt F'=a$, the two
series agree outside index $n_0$, and $[t^N](FG)-[t^N](F'G)$ equals $1$
respectively $-1$.  The symmetric statement follows by commutativity.
\end{proof}

\begin{checkpointbox}
\textbf{The precision rule.}  If $F$ is known through $\Trunc{F}{P}$ and $G$
through $\Trunc{G}{Q}$, and $\ordt F=a$, $\ordt G=b$, then $FG$ is known
through
\[
 \Trunc{FG}{\min(P+b,\;Q+a)}
\]
and, in the worst case, no further.  The binding entry is contributed by
whichever factor is paired with the \emph{more singular} partner: a factor
known to $P$ blocks is degraded by the valuation $b$ of the other factor.  In
particular, multiplying by a series with a pole of order $r$, that is with
valuation $-r$, costs exactly $r$ blocks of precision.
\end{checkpointbox}

% ed.: the bracketed title duplicated that of \cref{plt:prop:trunc-product}.
% ed.: The two statements are not the same result: the ring-chapter proposition
% ed.: is the law on \(\PLlau\) at the canonical cutoffs, this one is the law
% ed.: over an arbitrary logarithmically graded coefficient ring at arbitrary
% ed.: sufficient cutoffs, together with optimality and sharpness.  The titles
% ed.: now say so; see \cref{plt:rmk:mul-leading-restated}.
\begin{corollary}[Truncated product over a general coefficient ring, at
arbitrary sufficient cutoffs]
\label{plt:cor:mul-truncated-product}
Let $F,G\in R((t))$ be nonzero with $a=\ordt F$, $b=\ordt G$, and let
$N\in\IntegerNumbers$.  Then for all integers $P\ge N-b$ and $Q\ge N-a$,
\begin{equation}\label{plt:eq:mul-truncated-product}
 \Trunc{FG}{N}
 =\Trunc{\bigl(\Trunc{F}{P}\bigr)\bigl(\Trunc{G}{Q}\bigr)}{N};
\end{equation}
% ed.: the optimality of $(N-b,N-a)$ was asserted unconditionally.  For
% ed.: $N\le a+b$ both sides of \eqref{plt:eq:mul-truncated-product} vanish
% ed.: whatever $P,Q$ are, so no pair is forced and there is nothing to be
% ed.: optimal about; the clause is now stated under $N>a+b$, which is what
% ed.: the argument actually establishes.
and if $N>a+b$ the pair $(P,Q)=(N-b,\,N-a)$ is optimal: neither entry can be
lowered by one, that is, there exist $F,G$ with these valuations
for which $P=N-b-1$, respectively $Q=N-a-1$, falsifies
\eqref{plt:eq:mul-truncated-product}.  If $F,G\in\PLpow$ and $N\ge0$, then
$(P,Q)=(N,N)$ is admissible, so the cruder law
\begin{equation}\label{plt:eq:mul-truncated-product-crude}
 \Trunc{FG}{N}=\Trunc{\bigl(\Trunc{F}{N}\bigr)\bigl(\Trunc{G}{N}\bigr)}{N}
 \qquad(N\ge0)
\end{equation}
holds there.
\end{corollary}

\begin{proof}
Write $F=\Trunc{F}{P}+E$ and $G=\Trunc{G}{Q}+H$, so that $\ordt E\ge P$ and
$\ordt H\ge Q$, while $\ordt\Trunc{F}{P}\ge a$ and $\ordt\Trunc{G}{Q}\ge b$.
Then
\[
 FG-\bigl(\Trunc{F}{P}\bigr)\bigl(\Trunc{G}{Q}\bigr)
 =E\,\Trunc{G}{Q}+\Trunc{F}{P}\,H+EH,
\]
% ed.: the source bounds the third term by $P+Q$ and asserts $P+Q\ge N$.  That
% ed.: inequality does not follow from $P\ge N-b$ and $Q\ge N-a$: for
% ed.: $a=b=0$ and $N=P=Q=-5$ one has $P+Q=-10<N$.  The repair is to bound
% ed.: $\ordt H$ by $b$ rather than by $Q$, which is legitimate because every
% ed.: block of $H$ is a block of $G$.
whose $\ordt$ is at least $P+b\ge N$, $a+Q\ge N$ and $P+b\ge N$
respectively; for the third term note that $H=G-\Trunc{G}{Q}$ has
$\ordt H\ge b$, every block of $H$ being a block of $G$, so
$\ordt(EH)\ge P+b$.  Hence the two products agree below $t^N$, which is
\eqref{plt:eq:mul-truncated-product}.

For the optimality and the sharpness, assume $N>a+b$, so that $N-1\ge a+b$.
By \cref{plt:thm:mul-finite-determination} the block $[t^{N-1}](FG)$ already
depends on $p_{N-1-b}$ and on $q_{N-1-a}$, so no pair below $(N-b,N-a)$ can
serve for all $F,G$; the following pair of series shows that neither entry
may be lowered by one.  If $N-b-1>a$ put $F=t^a+t^{N-b-1}$; if
$N-b-1=a$ put $F=t^a$.  In both cases take $G=t^b$.  Then
$[t^{N-1}](FG)=1$, whereas $\Trunc{F}{N-b-1}$ omits the block $t^{N-b-1}$,
so $\Trunc{(\Trunc{F}{N-b-1})G}{N}$ has zero block at $t^{N-1}$.  (In the
first case $N-1>a+b$, so the surviving term $t^{a+b}$ does not sit at
level $N-1$.)  The second cutoff is symmetric.

For \eqref{plt:eq:mul-truncated-product-crude}, note that $a,b\ge0$ gives
$N\ge N-b$ and $N\ge N-a$, so $(P,Q)=(N,N)$ is admissible.
\end{proof}

\begin{warningbox}
The crude law \eqref{plt:eq:mul-truncated-product-crude} is \emph{false} on
the Laurent ring.  Take $F=t^{-1}$, $G=1+t$, $N=1$.  Then $FG=t^{-1}+1$, so
$\Trunc{FG}{1}=t^{-1}+1$, whereas
$\Trunc{F}{1}\Trunc{G}{1}=t^{-1}\cdot1=t^{-1}$.  The correct cutoff for $G$
is $N-\ordt F=2$, and indeed $t^{-1}(1+t)=t^{-1}+1$.  Implementations that
``normalize every series to start at $t^0$'' silently commit this error
whenever a genuine principal part is present.
\end{warningbox}

% ed.: S6, Theorem "Closure and finite determination", eq. (mult-error-rule).
% ed.: The source states the conclusion as O(t^{min(p+m_0, q+n_0)}) without
% ed.: any hypothesis relating p,q to the valuations n_0,m_0, and its proof
% ed.: silently bounds ord(E*Gtilde) by p+m_0 and ord(Ftilde*H) by q+n_0,
% ed.: which requires ord(Gtilde) >= m_0 and ord(Ftilde) >= n_0.  Neither
% ed.: follows from the stated hypotheses.  The corrected statement carries
% ed.: the third entry p+q, and the hypotheses p >= n_0, q >= m_0 under
% ed.: which the source's form is recovered.  Note that these hypotheses are
% ed.: sufficient but not necessary: the third entry is redundant exactly
% ed.: when p >= a or q >= b.
\begin{corollary}[Error rule for approximate factors]
\label{plt:cor:mul-error-rule}
Let $F,G\in\PLlau$ with $\ordt F=a$, $\ordt G=b$, and suppose
$\widetilde F,\widetilde G\in\PLlau$ satisfy
\[
 F=\widetilde F+\FormalBigOAt{t^{p}}{t},
 \qquad
 G=\widetilde G+\FormalBigOAt{t^{q}}{t}.
\]
Then
\begin{equation}\label{plt:eq:mul-error-rule}
 FG=\widetilde F\widetilde G
 +\FormalBigOAt{t^{\min(p+b,\;q+a,\;p+q)}}{t}.
\end{equation}
If in addition $p\ge a$ and $q\ge b$ --- that is, if neither approximation is
coarser than the leading order of the quantity it approximates --- then the
third entry is redundant and
\begin{equation}\label{plt:eq:mul-error-rule-normalized}
 FG=\widetilde F\widetilde G
 +\FormalBigOAt{t^{\min(p+b,\;q+a)}}{t}.
\end{equation}
\end{corollary}

\begin{proof}
Write $E=F-\widetilde F$ and $H=G-\widetilde G$, so $\ordt E\ge p$ and
$\ordt H\ge q$.  Expanding,
\begin{equation}\label{plt:eq:mul-error-expansion}
 FG-\widetilde F\widetilde G
 =E\widetilde G+\widetilde FH+EH .
\end{equation}
Now $\widetilde G=G-H$ gives
$\ordt\widetilde G\ge\min(\ordt G,\ordt H)\ge\min(b,q)$, and likewise
$\ordt\widetilde F\ge\min(a,p)$.  Since $\ordt$ is additive on products
(\cref{plt:thm:mul-leading-law} below, or directly
\cref{plt:lem:mul-finite-fiber}) and satisfies
$\ordt(A+B)\ge\min(\ordt A,\ordt B)$,
\[
 \ordt\bigl(FG-\widetilde F\widetilde G\bigr)
 \ \ge\
 \min\bigl(p+\min(b,q),\ \min(a,p)+q,\ p+q\bigr)
 \ =\ \min(p+b,\ q+a,\ p+q),
\]
which is \eqref{plt:eq:mul-error-rule}.  If $p\ge a$ and $q\ge b$, then
$p+q\ge a+q$ and $p+q\ge p+b$, so $p+q\ge\min(p+b,q+a)$ and the third entry
may be dropped.
\end{proof}

\begin{remark}[What \eqref{plt:eq:mul-error-rule} does and does not say]
\label{plt:rmk:mul-not-an-ideal}
The symbol $\FormalBigOAt{t^{N}}{t}$ asserts coefficient agreement below
outer order $N$ and nothing else: no numerical bound, no convergence, no
uniformity in $L$.  It is also not a statement about a quotient ring.  In
$\PLlau$ the element $t$ is a unit, so the ideal generated by $t^N$ is the
whole ring; the subset $t^N\PLpow$ is an additive piece of the
valuation filtration, not an ideal.  Congruences of the form
\eqref{plt:eq:mul-error-rule} must therefore be read as valuation
inequalities $\ordt(F-G)\ge N$, or equivalently as equalities of exclusive
truncations.  Quotient-algebra language becomes legitimate only after the
finite principal part has been removed and one works inside $\PLpow$, where
$(t^N)$ is a genuine ideal.
\end{remark}

The corresponding algorithm is the schoolbook convolution restricted to the
window supplied by \cref{plt:thm:mul-finite-determination}.

\begin{lstlisting}[style=pseudo,caption={Truncated block product},label={plt:alg:mul-truncated}]
# F, G sparse maps  exponent -> coefficient in R;  a = ord(F), b = ord(G)
# returns the blocks of Trunc_{<N}(F*G)
function block_product(F, G, a, b, N):
    C := empty map
    for (n, Pn) in F with a <= n <= N-b-1:
        for (m, Qm) in G with b <= m <= N-n-1:
            C[n+m] := C.get(n+m, 0) + Pn*Qm
    normalize every coefficient of C          # collect in L before zero-testing
    delete the blocks of C that are exactly 0
    return C
\end{lstlisting}

\begin{remark}[Cost, and what actually dominates it]
\label{plt:rmk:mul-cost}
With dense consecutive support the loop performs
$\BigO\bigl((N-a-b)^2\bigr)$ coefficient products; with sparse supports of
sizes $s_F,s_G$ it performs at most $s_Fs_G$.  If the retained blocks have
logarithmic degree $\BigO(d)$, each coefficient product costs
$\BigO(d^2)$ scalar operations by schoolbook polynomial arithmetic, so the
total is $\BigO\bigl((N-a-b)^2d^2\bigr)$; fast convolution in $t$ and fast
polynomial multiplication in $L$ reduce both factors.  These are complexity
$\BigO$'s and carry no analytic meaning.  In practice, at the moderate orders
typical of hand asymptotics, exact coefficient normalization --- expansion,
collection, and gcd reduction if the coefficients are rational functions of
$L$ --- costs more than the convolution itself, and the decisive
implementation choice is to never form a product whose outer order exceeds
the requested cutoff.
\end{remark}

\begin{remark}[A rectangular cutoff is not a truncation]
\label{plt:rmk:mul-rectangular-cutoff}
Discarding all monomials $t^nL^j$ with $|n|+|j|$ or $|j|$ above a threshold
does not respect dominance.  With a cap $|j|\le10$ one deletes $t^2L^{100}$
while retaining $t^3L^{-100}$, although
$t^3L^{-100}\precdom t^2L^{100}$ by an enormous margin.  Bounding
$\degL p_n$ is a second, independent approximation; it is legitimate only
when accompanied by a proved degree bound such as
\cref{plt:prop:mul-affine-profile}.
\end{remark}

\section{Leading data and the leading-term law}
\label{plt:sec:mul-leading}

% ed.: retitled to distinguish it from \cref{plt:def:lead-leading-data}, which
% ed.: is the same data on the four ambient classes; only the coefficient ring
% ed.: differs.  See \cref{plt:rmk:mul-leading-restated}.
\begin{definition}[Leading data over a logarithmically graded domain]
\label{plt:def:mul-leading-data}
Let $R$ be a logarithmically graded domain over $\K$ and let
$F=\sum_{n\ge n_0}p_nt^n\in R((t))$ be nonzero with
$\ordt F=\min\SetOf{n:p_n\ne0}$; set $\ordt0=+\infty$.  The
\emph{leading block} of $F$ is the whole coefficient $p_{\ordt F}\in R$, the
\emph{leading monomial} is
\begin{equation}\label{plt:eq:mul-leading-monomial}
 \lm(F)\DefinitionEquals t^{\ordt F}L^{\degL p_{\ordt F}},
\end{equation}
and the \emph{leading scalar} is
$\lc(F)\DefinitionEquals\lc_L\bigl(p_{\ordt F}\bigr)\in\K^\times$.  The
\emph{rank-two valuation} of $F$ is
\begin{equation}\label{plt:eq:mul-rank-two}
 v(F)\DefinitionEquals\bigl(\ordt F,\ -\degL p_{\ordt F}\bigr)
 \in\IntegerNumbers^2,
\end{equation}
with $\IntegerNumbers^2$ lexicographically ordered and $v(0)=+\infty$.  These
are three different objects and none is determined by the coarse integer
$\ordt F$ alone.
\end{definition}

% ed.: the leading data, the rank-two valuation, the log-degree profile, the
% ed.: exclusive truncation law and the summability calculus are all
% ed.: introduced in the ring chapter.  Restating them here without saying so
% ed.: would leave two definitions of the same symbols standing in one volume.
% ed.: This remark records that the restatements are verbatim generalizations
% ed.: over the coefficient ring and agree with the originals.
\begin{remark}[These are the objects of the ring chapter, over a general
coefficient ring]
\label{plt:rmk:mul-leading-restated}
\cref{plt:def:mul-leading-data} introduces no new object.  It repeats
\cref{plt:def:lead-leading-data} and \cref{plt:def:ring-rank-two} with the
three coefficient rings $\K[L]$, $\K(L)$, $\K((L^{-1}))$ replaced by an
arbitrary logarithmically graded domain $R$, so that the leading-term law can
be proved once for all four ambient classes simultaneously; on those four
classes the two readings agree.  The same applies to
\cref{plt:def:mul-profile} against \cref{plt:def:lead-profile} --- whose
codomain is $\NaturalNumbers\cup\SetOf{-\infty}$ because there the blocks are
polynomials, whereas over $\K(L)$ and $\K((L^{-1}))$ the value $\degL p_n$ may
be negative --- and to \cref{plt:cor:mul-truncated-product} against
\cref{plt:prop:trunc-product}, which the corollary strengthens to arbitrary
sufficient cutoffs $(P,Q)$ and supplements with optimality and sharpness.
Nothing below re-proves a ring-chapter result for its own sake; the
restatements exist only to make the general coefficient ring available.
\end{remark}

\begin{lemma}[Dominance is growth]
\label{plt:lem:mul-dominance}
On the positive ray, for integers $n,m$ and $j,k$,
\[
 t^nL^j\succdom t^mL^k
 \quad\Longleftrightarrow\quad
 \lim_{X\to+\infty}\frac{X^{-n}(\log X)^j}{X^{-m}(\log X)^k}=+\infty
 \quad\Longleftrightarrow\quad
 n<m,\ \text{or}\ (n=m\ \text{and}\ j>k).
\]
Consequently, for $F$ in $\PLlau$ or in $\PLit$, $\lm(F)$ is the
$\succdom$-largest element of $\CoefficientSupportOf{F}$.  (For $\PLrat$ the
statement is meaningful only after the expansion of
\cref{plt:rmk:mul-rat-no-support}.)
\end{lemma}

\begin{proof}
% ed.: the first equivalence of the display was left unargued.  It is not a
% ed.: theorem of this section at all: $\succdom$ is defined by the
% ed.: right-hand condition in \cref{plt:def:ring-dominance}, so only the
% ed.: analytic equivalence needs an argument, and that argument restates
% ed.: \cref{plt:lem:ring-dominance-analytic} for integer exponents.
By \cref{plt:def:ring-dominance} the relation $\succdom$ is \emph{defined} by
the right-hand condition, so the first equivalence is a definition and only
the second requires an argument; it is
\cref{plt:lem:ring-dominance-analytic} specialized to integer exponents, and
we repeat the two lines.  The ratio equals $X^{m-n}(\log X)^{j-k}$.  If $m>n$
it tends to $+\infty$
because $\log X=\LittleOAt{X^{\varepsilon}}{X\to+\infty}$ for every
$\varepsilon>0$; if $m=n$ it is $(\log X)^{j-k}$, which tends to $+\infty$
exactly when $j>k$; if $m<n$ it tends to $0$ for the same reason as the first
case.  For the last sentence: a monomial $t^nL^j$ in the support has
$n\ge\ordt F$, with equality attained; among those with $n=\ordt F$ the
largest $j$ with $\kappa_j(p_{\ordt F})\ne0$ is by definition
$\degL p_{\ordt F}$.
\end{proof}

\begin{theorem}[Leading-term law]
\label{plt:thm:mul-leading-law}
Let $R$ be a logarithmically graded domain over $\K$ and let $F,G\in R((t))$
be nonzero.  Then $FG\ne0$ and
\begin{align}
 \ordt(FG)&=\ordt F+\ordt G,
 \label{plt:eq:mul-ord-additive}\\
 \lm(FG)&=\lm(F)\,\lm(G),
 \label{plt:eq:mul-lm-law}\\
 \lc(FG)&=\lc(F)\,\lc(G),
 \label{plt:eq:mul-lc-law}\\
 v(FG)&=v(F)+v(G).
 \label{plt:eq:mul-v-additive}
\end{align}
Moreover $\ordt(F+G)\ge\min(\ordt F,\ordt G)$ and
$v(F+G)\ge\min\bigl(v(F),v(G)\bigr)$ in the lexicographic order.  Each is an
equality unless cancellation occurs at the common level: whole leading blocks
cancel in the first case, top logarithmic coefficients in the second.
\end{theorem}

\begin{proof}
Put $a=\ordt F$, $b=\ordt G$.  By \cref{plt:lem:mul-finite-fiber} the
coefficient $C_N$ of $FG$ vanishes for $N<a+b$, and for $N=a+b$ the index set
$[a,N-b]=\SetOf{a}$ is a single point, so $C_{a+b}=p_aq_b$.  Since $R$ is a
domain (\cref{plt:lem:mul-graded-basics}) and $p_a,q_b\ne0$, we get
$C_{a+b}\ne0$; hence $FG\ne0$ and \eqref{plt:eq:mul-ord-additive} holds.
Applying \eqref{plt:eq:mul-degL-laws} to $C_{a+b}=p_aq_b$ gives
\[
 \degL C_{a+b}=\degL p_a+\degL q_b,
 \qquad
 \lc_L(C_{a+b})=\lc_L(p_a)\lc_L(q_b),
\]
which are exactly \eqref{plt:eq:mul-lm-law} and \eqref{plt:eq:mul-lc-law};
\eqref{plt:eq:mul-v-additive} is the same pair of identities rewritten in the
notation \eqref{plt:eq:mul-rank-two}.

For the sum, let $c=\min(a,b)$.  Every block of $F+G$ of index $<c$ vanishes,
so $\ordt(F+G)\ge c$.  If $a<b$ then $[t^c](F+G)=p_a\ne0$ and the leading
data of $F+G$ and $F$ coincide, so $v(F+G)=v(F)=\min(v(F),v(G))$ --- note
that $a<b$ makes $v(F)<v(G)$ lexicographically.  If $a=b=c$ and
$p_c+q_c\ne0$, then $\ordt(F+G)=c$ and
$\degL(p_c+q_c)\le\max(\degL p_c,\degL q_c)$ by
\eqref{plt:eq:mul-degL-laws}, i.e.
$-\degL(p_c+q_c)\ge\min(-\degL p_c,-\degL q_c)$, which is
$v(F+G)\ge\min(v(F),v(G))$.  If $p_c+q_c=0$, then $\ordt(F+G)>c$ and
$v(F+G)>\min(v(F),v(G))$ outright.
\end{proof}

\begin{corollary}[Integrality and fields]
\label{plt:cor:mul-domain}
Each of $\PLpow$, $\PLlau$, $\PLrat$, $\PLit$ is an integral domain.  The map
$v$ of \eqref{plt:eq:mul-rank-two} is a rank-two valuation on $\PLrat$ and on
$\PLit$, with value group $\IntegerNumbers^2$ lexicographically ordered; on
$\PLlau$ it takes values in the sub-semigroup
$\IntegerNumbers\times\IntegerNumbers_{\le0}$.  The classes $\PLrat$ and
$\PLit$ are fields; $\PLpow$ and $\PLlau$ are not.
\end{corollary}

\begin{proof}
Integrality is \cref{plt:thm:mul-leading-law}, and the valuation axioms are
its last two displays together with $v(F)=+\infty$ iff $F=0$.  Surjectivity
onto $\IntegerNumbers^2$ for $\PLrat$ and $\PLit$ follows from
$v(t^nL^{-j})=(n,j)$, legitimate there because $L^{-j}$ lies in $\K(L)$ and
in $\K((L^{-1}))$; the group $\IntegerNumbers^2$ with the lexicographic order
has exactly one proper nontrivial convex subgroup,
$\SetOf{0}\times\IntegerNumbers$, so the valuation has rank two.  On $\PLlau$
the second coordinate is $-\degL p_{\ordt F}\le0$ because polynomial degrees
are nonnegative.

A Laurent-series ring $R((t))$ over a \emph{field} $R$ is a field, which
covers $\PLrat$ and $\PLit$: given $F\ne0$ with $\nu=\ordt F$ and leading
block $p_\nu\in R^\times$, write $F=p_\nu t^{\nu}(1+U)$ with
$U=\sum_{n>\nu}(p_n/p_\nu)t^{n-\nu}\in tR[[t]]$; the family
$\bigl((-U)^r\bigr)_{r\ge0}$ is summable because $\ordt(-U)^r\ge r$, and by
\cref{plt:lem:mul-summable}(2) its sum $V$ satisfies
\[
 (1+U)V=\sum_{r\ge0}(-U)^r-\sum_{r\ge0}(-U)^{r+1}=(-U)^0=1,
\]
so $F^{-1}=p_\nu^{-1}t^{-\nu}V$.  (\cref{plt:lem:mul-summable} is proved in
\cref{plt:sec:mul-powers} from \cref{plt:thm:mul-leading-law} alone, so there
is no circularity.)  Finally $L\in\PLlau$ is not invertible in
$\PLlau$: if $LH=1$ then \eqref{plt:eq:mul-lm-law} gives
$L\cdot\lm(H)=1$, forcing $\degL$ of the leading block of $H$ to be $-1$,
impossible for a polynomial.
% ed.: the statement asserts that neither $\PLpow$ nor $\PLlau$ is a field,
% ed.: but the argument as given covers only $\PLlau$.  One line closes it.
Since $L\in\PLpow\subseteq\PLlau$, an inverse of $L$ inside $\PLpow$ would in
particular be an inverse inside $\PLlau$; so $L$ is not invertible in
$\PLpow$ either, and neither class is a field.
\end{proof}

\begin{corollary}[Necessity half of the unit criterion]
\label{plt:cor:mul-unit-necessary}
If $F\in\PLlau$ is invertible in $\PLlau$, then its leading block is a
nonzero scalar: $p_{\ordt F}\in\K^\times$.  The same statement holds in
$\PLpow$ with the additional conclusion $\ordt F=0$.
\end{corollary}

\begin{proof}
If $FH=1$ then \eqref{plt:eq:mul-ord-additive} gives
$\ordt F+\ordt H=0$ and \eqref{plt:eq:mul-lm-law} gives
$\degL p_{\ordt F}+\degL h_{\ordt H}=0$.  Both degrees are nonnegative
integers, hence both vanish, so $p_{\ordt F}\in\K^\times$.  In $\PLpow$ both
orders are $\ge0$ and sum to $0$, so both are $0$.
\end{proof}

The converse --- that a scalar leading block suffices --- is a division
statement and is proved with the reciprocal recurrence; it is not a
multiplication fact and is not repeated here.

\section{Degree and support bounds}
\label{plt:sec:mul-degree}

% ed.: retitled to distinguish it from \cref{plt:def:lead-profile}, which is
% ed.: the same function on \(\PLlau\), where the blocks are polynomials and
% ed.: the codomain is \(\NaturalNumbers\cup\SetOf{-\infty}\).  See
% ed.: \cref{plt:rmk:mul-leading-restated}.
\begin{definition}[Log-degree profile on the four classes]
\label{plt:def:mul-profile}
For $F=\sum_np_nt^n$ in one of the four classes, the \emph{log-degree
profile} is the function $\dprofile{F}\colon\IntegerNumbers\to
\IntegerNumbers\cup\SetOf{-\infty}$,
\begin{equation}\label{plt:eq:mul-profile}
 \dprofile{F}(n)\DefinitionEquals
 \begin{cases}
  \degL p_n,&p_n\ne0,\\
  -\infty,&p_n=0.
 \end{cases}
\end{equation}
The profile is data additional to $\ordt F$; it is not recoverable from it.
For two profiles define the \emph{max-plus convolution}
\begin{equation}\label{plt:eq:mul-maxplus}
 \bigl(\dprofile{F}\ast\dprofile{G}\bigr)(N)
 \DefinitionEquals
 \max_{n+m=N}\bigl(\dprofile{F}(n)+\dprofile{G}(m)\bigr),
\end{equation}
with the conventions $-\infty+c=-\infty$ and $\max\emptyset=-\infty$.  Only
the finitely many pairs of \cref{plt:lem:mul-finite-fiber} can contribute a
value other than $-\infty$, so the maximum is effectively over a finite set.
\end{definition}

\begin{theorem}[Degree bound and exact criterion for attainment]
\label{plt:thm:mul-degree-bound}
Let $F,G\in R((t))$ with $R$ a logarithmically graded domain over $\K$, let
$N\in\IntegerNumbers$, and set
$M\DefinitionEquals(\dprofile{F}\ast\dprofile{G})(N)$.  Let
\[
 P_N\DefinitionEquals
 \SetOf{(n,m):n+m=N,\ p_n\ne0,\ q_m\ne0,\
        \dprofile{F}(n)+\dprofile{G}(m)=M}
\]
be the (finite, possibly empty) set of maximizing pairs.  Then
\begin{equation}\label{plt:eq:mul-degree-bound}
 \dprofile{FG}(N)\le M,
\end{equation}
% ed.: the criterion was stated as an unconditional ``if and only if''.  It
% ed.: fails in the degenerate case $P_N=\emptyset$ (equivalently
% ed.: $M=-\infty$): there $[t^N](FG)=0$, so equality holds although the empty
% ed.: sum \eqref{plt:eq:mul-top-coefficient} is zero.  The criterion is now
% ed.: confined to $M$ finite and the degenerate case is recorded separately.
If $P_N=\emptyset$, equivalently $M=-\infty$, then $[t^N](FG)=0$ and
\eqref{plt:eq:mul-degree-bound} holds with equality.  If $P_N\ne\emptyset$,
so that $M\in\IntegerNumbers$, then the coefficient of $L^{M}$ in the block
$[t^N](FG)$ is exactly
\begin{equation}\label{plt:eq:mul-top-coefficient}
 \kappa_{M}\bigl([t^N](FG)\bigr)
 =\sum_{(n,m)\in P_N}\lc_L(p_n)\,\lc_L(q_m),
\end{equation}
and equality holds in \eqref{plt:eq:mul-degree-bound} if and only if the
sum \eqref{plt:eq:mul-top-coefficient} is nonzero.
\end{theorem}

\begin{proof}
Each summand of $C_N=\sum_{n+m=N}p_nq_m$ has
$\degL(p_nq_m)=\dprofile{F}(n)+\dprofile{G}(m)\le M$ by
\eqref{plt:eq:mul-degL-laws}, with the convention that a vanishing factor
gives $-\infty$.  A finite sum of elements of $\degL\le M$ has
$\degL\le M$, which is \eqref{plt:eq:mul-degree-bound}.  If $M=-\infty$ then
every summand $p_nq_m$ vanishes, so $[t^N](FG)=0$ and both sides of
\eqref{plt:eq:mul-degree-bound} equal $-\infty$; assume from now on
$P_N\ne\emptyset$, so that $M\in\IntegerNumbers$.  By $\K$-linearity of
$\kappa_M$,
\[
 \kappa_M(C_N)=\sum_{n+m=N}\kappa_M(p_nq_m).
\]
If $(n,m)\notin P_N$ then either a factor vanishes or
$\degL(p_nq_m)<M$, and in both cases $\kappa_M(p_nq_m)=0$.  If
$(n,m)\in P_N$ then $\degL(p_nq_m)=M$ and
$\kappa_M(p_nq_m)=\lc_L(p_nq_m)=\lc_L(p_n)\lc_L(q_m)$ by
\cref{plt:lem:mul-graded-basics}.  This proves
\eqref{plt:eq:mul-top-coefficient}, and equality in
\eqref{plt:eq:mul-degree-bound} means precisely $\kappa_M(C_N)\ne0$.
\end{proof}

% ed.: S6, Proposition "Degree and support bounds", asserts equality "whenever
% ed.: the maximizing blocks are nonzero" under the hypothesis that all
% ed.: coefficients are nonnegative over an ordered field.  That hypothesis is
% ed.: far stronger than necessary and the quoted sufficient condition is
% ed.: stated ambiguously (a zero block has degree -infinity and cannot
% ed.: maximize at all).  The criterion below is exact and the positivity
% ed.: hypothesis has been reduced to the leading coefficients of the
% ed.: maximizing blocks only.
\begin{corollary}[Three sufficient conditions for equality]
\label{plt:cor:mul-degree-equality}
In the situation of \cref{plt:thm:mul-degree-bound}, equality holds in
\eqref{plt:eq:mul-degree-bound} in each of the following cases.
\begin{enumerate}
\item $P_N$ is a singleton.
\item $N=\ordt F+\ordt G$ (both factors nonzero); then $P_N$ is the
singleton $\SetOf{(\ordt F,\ordt G)}$ and one recovers
\eqref{plt:eq:mul-lm-law}.
% ed.: condition (3) as first drafted omitted $P_N\ne\emptyset$; over an
% ed.: ordered field its hypotheses are satisfied vacuously when $P_N$ is
% ed.: empty, and the conclusion would then be the false assertion that the
% ed.: sum \eqref{plt:eq:mul-top-coefficient} is nonzero.
\item $P_N\ne\emptyset$, $\K$ is an ordered field, and $\lc_L(p_n)>0$,
$\lc_L(q_m)>0$ for every $(n,m)\in P_N$.
\end{enumerate}
\end{corollary}

\begin{proof}
(1) The sum \eqref{plt:eq:mul-top-coefficient} is a single product of two
elements of $\K^\times$, hence nonzero.  (2) By
\cref{plt:lem:mul-finite-fiber} the only admissible pair at
$N=\ordt F+\ordt G$ is $(\ordt F,\ordt G)$, and both blocks are nonzero, so
$P_N$ is that singleton; apply (1).  (3) The sum is a nonempty sum of
strictly positive elements of an ordered field.
\end{proof}

\begin{example}[The bound is attained, at every order and by every pair]
\label{plt:ex:mul-degree-sharp}
Over $\K=\RationalNumbers$ let
\[
 F=G=\sum_{n\ge0}L^nt^n\in\PLpow,
 \qquad\text{so}\qquad
 \dprofile{F}(n)=\dprofile{G}(n)=n .
\]
Then $(\dprofile{F}\ast\dprofile{G})(N)=\max_{n+m=N}(n+m)=N$, and
\emph{every} one of the $N+1$ admissible pairs is maximizing, so $P_N$ is as
large as it can be and \cref{plt:cor:mul-degree-equality}(1) does not apply.
Nevertheless
\[
 [t^N](FG)=\sum_{n=0}^{N}L^n\cdot L^{N-n}=(N+1)L^N,
\]
so $\dprofile{FG}(N)=N$ for every $N$: the bound
\eqref{plt:eq:mul-degree-bound} is attained at every order.  Formula
\eqref{plt:eq:mul-top-coefficient} predicts exactly the top coefficient
$N+1=\sum_{(n,m)\in P_N}1\cdot1$.
\end{example}

\begin{warningbox}
\cref{plt:ex:mul-degree-sharp} is the place where characteristic zero is
used.  Over a field of characteristic $p>0$ the same computation gives
$[t^{p-1}](FG)=p\,L^{p-1}=0$, so the block vanishes entirely and
$\dprofile{FG}(p-1)=-\infty$ although the bound is $p-1$.  Every statement in
this section is asserted over a field of characteristic zero, and this is not
a formality.
\end{warningbox}

\begin{example}[The bound can fail by an arbitrary amount]
\label{plt:ex:mul-degree-drop}
Let $F=1+L^{d}t$ and $G=1-L^{d}t$ in $\PLpow$, $d\ge1$.  Then
$(\dprofile{F}\ast\dprofile{G})(1)=\max(0+d,\ d+0)=d$, while
$[t^1](FG)=-L^{d}+L^{d}=0$, so $\dprofile{FG}(1)=-\infty$.  Here
$P_1=\SetOf{(0,1),(1,0)}$ and the sum
\eqref{plt:eq:mul-top-coefficient} is $1\cdot(-1)+1\cdot1=0$, as it must be.
A drop by exactly one occurs for $F=1+(L^2+L)t$, $G=1-(L^2-L)t$, where the
bound at $N=1$ is $2$ while $[t^1](FG)=2L$ has degree $1$: here the top
coefficients cancel but the next ones do not.
\end{example}

\begin{proposition}[Affine profiles and the profile subalgebra]
\label{plt:prop:mul-affine-profile}
Let $\alpha,\alpha'\ge0$ and $\beta,\gamma\in\RealNumbers$.  Let
$F,G\in\PLlau$ be nonzero with $a=\ordt F$, $b=\ordt G$ and
\[
 \dprofile{F}(n)\le\alpha n+\beta\ \ (n\ge a),
 \qquad
 \dprofile{G}(m)\le\alpha'm+\gamma\ \ (m\ge b).
\]
Put $\mu=\max(\alpha,\alpha')$.  Then for all $N$,
\begin{equation}\label{plt:eq:mul-affine-bound}
 \dprofile{FG}(N)\le\mu N+\beta+\gamma-(\mu-\alpha)a-(\mu-\alpha')b .
\end{equation}
In particular, if $a,b\ge0$ then $\dprofile{FG}(N)\le\mu N+\beta+\gamma$.
Consequently, for each $\alpha\ge0$ the set
\[
 \mathfrak A_\alpha
 \DefinitionEquals
 \SetOf{F\in\PLlau:\ \exists\beta\in\RealNumbers,\
        \dprofile{F}(n)\le\alpha n+\beta\ \text{for all }n}
\]
is a $\K$-subspace with
$\mathfrak A_\alpha\cdot\mathfrak A_{\alpha'}\subseteq
\mathfrak A_{\max(\alpha,\alpha')}$, and
$\mathfrak A\DefinitionEquals\bigcup_{\alpha\ge0}\mathfrak A_\alpha$ is a
$\K$-subalgebra of $\PLlau$ containing $\K[L][t,t^{-1}]$.
\end{proposition}

\begin{proof}
Fix $N$ and $n+m=N$ with $n\ge a$, $m\ge b$.  Then
\[
 \alpha n+\alpha'm
 =\mu(n+m)-(\mu-\alpha)n-(\mu-\alpha')m
 \le\mu N-(\mu-\alpha)a-(\mu-\alpha')b,
\]
because $\mu-\alpha\ge0$, $\mu-\alpha'\ge0$, $n\ge a$ and $m\ge b$.  Adding
$\beta+\gamma$ and taking the maximum over the finite index set gives
\eqref{plt:eq:mul-affine-bound} via \eqref{plt:eq:mul-degree-bound}.  If
$a,b\ge0$ then the two subtracted terms are $\ge0$, so they may be dropped.

For the algebra statements: $\mathfrak A_\alpha$ is a subspace because
$\dprofile{F+G}(n)\le\max(\dprofile{F}(n),\dprofile{G}(n))
\le\alpha n+\max(\beta,\gamma)$ and scalars do not change profiles.  The
product statement is \eqref{plt:eq:mul-affine-bound}, whose right-hand side is
$\mu N+\text{const}$.
% ed.: the inclusion $\mathfrak A_\alpha\subseteq\mathfrak A_{\alpha''}$ for
% ed.: $\alpha\le\alpha''$ was used without proof.  On the Laurent ring it is
% ed.: not immediate: for $n<0$ a larger slope gives a \emph{smaller} bound,
% ed.: so the constant has to be increased, and that is possible only because
% ed.: the support is bounded below.  The argument is supplied here.
For $\alpha\le\alpha''$ one has $\mathfrak A_\alpha\subseteq\mathfrak
A_{\alpha''}$.  Indeed, let $F\ne0$ satisfy $\dprofile{F}(n)\le\alpha n+\beta$
for all $n$, and put $a=\ordt F$.  For $n\ge0$ we have
$\alpha n+\beta\le\alpha''n+\beta$; for $a\le n<0$,
\[
 \alpha n+\beta=\alpha''n+(\alpha''-\alpha)(-n)+\beta
 \le\alpha''n+\beta+(\alpha''-\alpha)\max(0,-a);
\]
and $\dprofile{F}(n)=-\infty$ for $n<a$.  Hence
$\beta''=\beta+(\alpha''-\alpha)\max(0,-a)$ witnesses
$F\in\mathfrak A_{\alpha''}$, and $0\in\mathfrak A_{\alpha''}$ trivially.
Consequently the union $\mathfrak A$ is closed under
addition and multiplication and contains $1$; every Laurent polynomial
$\sum_{n=n_1}^{n_2}p_n(L)t^n$ lies in $\mathfrak A_0$ with
$\beta=\max_n\degL p_n$.
\end{proof}

\begin{remark}[$\mathfrak A$ is not $t$-adically closed]
\label{plt:rmk:mul-profile-not-closed}
The affine-profile condition is preserved by finite sums and products but not
by $t$-adic limits.  Take $F_r=t^rL^{r^2}$ for $r\ge1$; each $F_r$ lies in
$\mathfrak A_0$, and $\ordt F_r=r\to\infty$, so $(F_r)$ is summable, yet the
sum
$S=\sum_{r\ge1}t^rL^{r^2}$ has $\dprofile{S}(r)=r^2$, which is bounded by no
affine function.  Hence $S\notin\mathfrak A$.  A degree bound must therefore
be re-derived, not inherited, whenever an infinite construction is performed.
\end{remark}

\begin{proposition}[Profile of a power]
\label{plt:prop:mul-power-profile}
Let $F\in\PLpow$ with $\dprofile{F}(n)\le\alpha n+\beta$ for all $n\ge0$,
where $\alpha\ge0$ and $\beta\ge0$.  Then for every integer $k\ge0$,
\begin{equation}\label{plt:eq:mul-power-profile}
 \dprofile{F^k}(N)\le\alpha N+k\beta
 \qquad(N\ge0).
\end{equation}
\end{proposition}

\begin{proof}
Induction on $k$.  For $k=0$ we have $F^0=1$, so $\dprofile{F^0}(0)=0\le0$
and $\dprofile{F^0}(N)=-\infty$ for $N>0$.  Assume the bound for $k$.  By
\eqref{plt:eq:mul-degree-bound} applied to $F^{k+1}=F^k\cdot F$,
\[
 \dprofile{F^{k+1}}(N)
 \le\max_{n+m=N,\ n,m\ge0}\bigl((\alpha n+k\beta)+(\alpha m+\beta)\bigr)
 =\alpha N+(k+1)\beta .
\]
\end{proof}

\begin{remark}[The Lambert profile, corrected]
\label{plt:rmk:mul-lambert-profile}
Two of the sources record the degree law of the canonical Lambert family as
$\deg\LamP{n}=n$ without restriction, and use it to assert that products of
Lambert-type expansions have profile $\le N$.  The law
$\deg\LamP{n}=n$ holds only for $n\ge1$: the normalization
$\LamP{0}(u)=-u$ has degree $1$, not $0$.  The correct hypothesis for a
series $F=\sum_{n\ge0}\LamP{n}(L)t^n$ is therefore
$\dprofile{F}(n)\le n+1$ for all $n\ge0$, i.e. $\alpha=1$, $\beta=1$ in
\cref{plt:prop:mul-affine-profile}; a product of two such series has
$\dprofile{FG}(N)\le N+2$, and a $k$-th power has
$\dprofile{F^k}(N)\le N+k$ by \cref{plt:prop:mul-power-profile}.  The
unrestricted claim $\dprofile{FG}(N)\le N$ is false already at $N=0$, where
$[t^0](FG)=\LamP{0}(L)^2=L^2$ has degree $2$.
\end{remark}

\section{Powers and binomial powers}
\label{plt:sec:mul-powers}

Infinite sums enter through powers, so we first record the exact sense in
which they are legitimate.

% ed.: retitled, and its relation to \cref{plt:def:trunc-summable} made
% ed.: explicit: the condition is verbatim the one of the ring chapter, read
% ed.: over the coefficient ring $R$ instead of $\K[L]$, which
% ed.: \cref{plt:rmk:trunc-summability-classes} already licenses.  Only the
% ed.: two clauses of \cref{plt:lem:mul-summable} are used below, so they are
% ed.: the only ones recorded here.
\begin{definition}[Summable family over a logarithmically graded domain]
\label{plt:def:mul-summable}
Let $R$ be a logarithmically graded domain over $\K$.  A family
$(S_i)_{i\in I}$ in $R((t))$ is \emph{summable} if for every
$N\in\IntegerNumbers$ the set $\SetOf{i\in I:\ordt S_i<N}$ is finite.  Its
sum is defined blockwise:
$[t^n]\sum_{i}S_i\DefinitionEquals\sum_i[t^n]S_i$, a finite sum in the
coefficient ring.
\end{definition}

\begin{lemma}[Summable families are well behaved]
\label{plt:lem:mul-summable}
Let $(S_i)_{i\in I}$ and $(T_j)_{j\in J}$ be summable families in $R((t))$.
\begin{enumerate}
\item $\inf_i\ordt S_i>-\infty$, so $\sum_iS_i$ is a genuine element of
$R((t))$ (its support is bounded below).
\item The family $(S_iT_j)_{(i,j)\in I\times J}$ is summable and
$\bigl(\sum_iS_i\bigr)\bigl(\sum_jT_j\bigr)=\sum_{(i,j)}S_iT_j$.
\end{enumerate}
\end{lemma}

\begin{proof}
(1) Taking $N=0$ in \cref{plt:def:mul-summable}, only finitely many $i$ have
$\ordt S_i<0$, and each of those has $\ordt S_i\in\IntegerNumbers$ (a member
of that finite set is in particular nonzero).  Let $\lambda$ be the minimum of
that finite set of integers, or $\lambda=0$ if the set is empty.  Then
$\ordt S_i\ge\min(0,\lambda)$ for all $i$, and every block of
$\sum_iS_i$ below that bound vanishes.

% ed.: the source calls the two infima ``finite by (1)''; part~(1) gives only
% ed.: $>-\infty$, and the value $+\infty$ does occur, namely when every
% ed.: member of the family is zero.  That case is disposed of separately.
(2) Let $\alpha=\inf_i\ordt S_i$ and $\beta=\inf_j\ordt T_j$, both
$>-\infty$ by (1).  If $\alpha=+\infty$ or $\beta=+\infty$ then every $S_i$,
respectively every $T_j$, is zero, so every $S_iT_j$ is zero and both
assertions are trivial; assume therefore $\alpha,\beta\in\IntegerNumbers$.
Fix $N$.  If $\ordt(S_iT_j)=\ordt S_i+\ordt T_j<N$ then
$\ordt S_i<N-\beta$ and $\ordt T_j<N-\alpha$, and each of those conditions
holds for only finitely many indices; so the family is summable.  For the
identity, fix $n$ and compute in the coefficient ring, all sums having only
finitely many nonzero terms:
\[
 [t^n]\Bigl(\sum_iS_i\sum_jT_j\Bigr)
 =\sum_{p+q=n}\Bigl(\sum_i[t^p]S_i\Bigr)\Bigl(\sum_j[t^q]T_j\Bigr)
 =\sum_{i,j}\sum_{p+q=n}[t^p]S_i\,[t^q]T_j
 =\sum_{i,j}[t^n](S_iT_j).
\]
\end{proof}

\begin{remark}
\label{plt:rmk:mul-summable-vs-cauchy}
Part (1) is worth isolating.  For a \emph{sequence} in a Laurent ring,
$t$-adic Cauchyness alone does not guarantee a limit with a finite principal
part; a common eventual lower support bound must be imposed.  For a
\emph{summable family} no such extra hypothesis is needed, because the
finiteness condition at $N=0$ already forces a common lower bound.
\end{remark}

\begin{proposition}[Coefficients of an integer power]
\label{plt:prop:mul-integer-power}
Let $F=\sum_{n\ge0}p_nt^n\in\PLpow$ and $k\ge0$.  Then
\begin{equation}\label{plt:eq:mul-power-coefficient}
 [t^N]F^k
 =\sum_{\substack{n_1,\ldots,n_k\ge0\\ n_1+\cdots+n_k=N}}
   p_{n_1}\cdots p_{n_k}
 \DefinitionEquals C_{N,k}(p_0,p_1,\ldots),
\end{equation}
with $C_{0,0}=1$ and $C_{N,0}=0$ for $N>0$.  These \emph{convolution
polynomials} satisfy
\begin{equation}\label{plt:eq:mul-Cmr-recurrence}
 C_{N,k+1}=\sum_{n=0}^{N}p_n\,C_{N-n,k},
\end{equation}
which computes the whole triangular table in $\BigO(N^2k)$ coefficient
products.  If $p_0=0$, then $C_{N,k}=0$ for $N<k$, and, writing
$F=\sum_{n\ge1}p_nt^n$,
\begin{equation}\label{plt:eq:mul-bell}
 [t^N]F^k
 =\frac{k!}{N!}\,
 \ExponentialPartialBellPolynomial{N}{k}
   \bigl(1!\,p_1,\,2!\,p_2,\ldots,(N-k+1)!\,p_{N-k+1}\bigr).
\end{equation}
\end{proposition}

\begin{proof}
Formula \eqref{plt:eq:mul-power-coefficient} follows by induction on $k$ from
\cref{plt:def:mul-cauchy}: the case $k=0$ is the convention, and
\[
 [t^N]F^{k+1}=\sum_{n+m=N}p_n\,[t^m]F^k
 =\sum_{n=0}^{N}p_n\!\!\!\!\sum_{\substack{n_1,\ldots,n_k\ge0\\
   n_1+\cdots+n_k=N-n}}\!\!\!\!
   p_{n_1}\cdots p_{n_k},
\]
which is both \eqref{plt:eq:mul-Cmr-recurrence} and the composition sum for
$k+1$.  All sums are finite because every $n_i\ge0$ and they sum to $N$.  If
$p_0=0$, every $n_i$ in a nonvanishing term is $\ge1$, so $N\ge k$.

For \eqref{plt:eq:mul-bell}, recall the defining generating identity of the
exponential partial Bell polynomials \cite{comtet-book,comtet1970}: in
$R[[z]]$ over a commutative $\RationalNumbers$-algebra $R$,
\[
 \frac1{k!}\Bigl(\sum_{j\ge1}x_j\frac{z^j}{j!}\Bigr)^{k}
 =\sum_{N\ge k}\ExponentialPartialBellPolynomial{N}{k}(x_1,x_2,\ldots)
   \frac{z^N}{N!}.
\]
Apply it with $R=\K[L]$ (a $\RationalNumbers$-algebra, as
$\operatorname{char}\K=0$), $z=t$ and $x_j=j!\,p_j$, so that
$\sum_{j\ge1}x_jz^j/j!=F$.  Comparing coefficients of $t^N$ gives
$\frac1{k!}[t^N]F^k=\frac1{N!}
\ExponentialPartialBellPolynomial{N}{k}(1!p_1,2!p_2,\ldots)$, which is
\eqref{plt:eq:mul-bell}; only the arguments $x_1,\ldots,x_{N-k+1}$ occur
because a term of $\ExponentialPartialBellPolynomial{N}{k}$ is a product of
$k$ variables with index sum $N$, so no index can exceed $N-k+1$.
\end{proof}

\begin{theorem}[Binomial powers and their finite determination]
\label{plt:thm:mul-binomial}
Let $U\in\PLpow$ with $\ordt U\ge e$ for an integer $e\ge1$, write
$u_n=[t^n]U$, and let $\sigma\in\K$.  Put
$\binom{\sigma}{r}=\sigma(\sigma-1)\cdots(\sigma-r+1)/r!\in\K$, with
$\binom{\sigma}{0}=1$.  Then:
\begin{enumerate}
\item The family $\bigl(\binom{\sigma}{r}U^r\bigr)_{r\ge0}$ is summable, so
\begin{equation}\label{plt:eq:mul-binomial-def}
 (1+U)^\sigma\DefinitionEquals\sum_{r\ge0}\binom{\sigma}{r}U^r
\end{equation}
is a well-defined element of $\PLpow$ with constant block $1$.
\item \emph{(Finite determination.)}  For $N\ge0$,
\begin{equation}\label{plt:eq:mul-binomial-finite}
 [t^N](1+U)^\sigma
 =\sum_{r=0}^{\Floor{N/e}}\binom{\sigma}{r}\,[t^N]U^r,
\end{equation}
and this depends only on the blocks $u_e,u_{e+1},\ldots,u_N$, that is, only
on $\Trunc{U}{N+1}$.  The range of $r$ is exact: for $U=t^e$ with $e\mid N$
the term $r=N/e$ contributes $\binom{\sigma}{N/e}$, nonzero whenever
$\sigma\notin\SetOf{0,1,\ldots,N/e-1}$.  The top block enters exactly
linearly: for $N\ge1$,
\begin{equation}\label{plt:eq:mul-binomial-top}
 [t^N](1+U)^\sigma=\sigma u_N+\Psi_N(u_e,\ldots,u_{N-e}),
\end{equation}
where $\Psi_N$ is a universal polynomial with coefficients in
$\RationalNumbers[\sigma]$, with the convention that $\Psi_N=0$ when the
argument list is empty.  Hence for $\sigma\ne0$ the block $u_N$ cannot be
omitted from the dependency range.
\item \emph{(Exponent law.)}  For $\sigma,\tau\in\K$,
\begin{equation}\label{plt:eq:mul-binomial-group}
 (1+U)^\sigma(1+U)^\tau=(1+U)^{\sigma+\tau},
 \qquad (1+U)^0=1,
\end{equation}
and for $k\in\NaturalNumbers$ the series $(1+U)^k$ agrees with the $k$-fold
product.  Consequently $(1+U)^\sigma$ is a unit of $\PLpow$ with inverse
$(1+U)^{-\sigma}$, and $\sigma\mapsto(1+U)^\sigma$ is a homomorphism from
$(\K,+)$ to the units of $\PLpow$.
\item \emph{(Profile.)}  If $\dprofile{U}(n)\le\alpha n+\beta$ for all
$n\ge e$, with $\alpha\ge0$, then, writing $\beta^+=\max(\beta,0)$,
\begin{equation}\label{plt:eq:mul-binomial-profile}
 \dprofile{(1+U)^\sigma}(N)\le\Bigl(\alpha+\frac{\beta^+}{e}\Bigr)N
 \qquad(N\ge0).
\end{equation}
In particular $\mathfrak A\cap\PLpow$ is stable under
$U\mapsto(1+U)^\sigma$ whenever $\ordt U\ge1$.
\end{enumerate}
\end{theorem}

\begin{proof}
(1) By \eqref{plt:eq:mul-ord-additive}, $\ordt(U^r)\ge re$, which tends to
$+\infty$; so for each $N$ only the finitely many $r\le N/e$ can have
$\ordt(\binom{\sigma}{r}U^r)<N$.  \cref{plt:lem:mul-summable}(1) gives a
well-defined sum, and $\ordt U\ge e\ge1$ makes the $r=0$ term the only
contributor to $t^0$.

(2) Blockwise, $[t^N](1+U)^\sigma=\sum_{r\ge0}\binom{\sigma}{r}[t^N]U^r$, and
$[t^N]U^r=0$ once $re>N$; this is \eqref{plt:eq:mul-binomial-finite}.  By
\cref{plt:thm:mul-finite-determination} applied $r-1$ times, $[t^N]U^r$
depends only on the blocks $u_n$ with $e\le n\le N-(r-1)e$, and for $r\ge1$
this range is contained in $[e,N]$.  For the sharpness of the $r$-range take
$U=t^e$: then $U^r=t^{re}$ and the only contribution to $t^N$ comes from
$r=N/e$, with value $\binom{\sigma}{N/e}$, which vanishes exactly when
$\sigma$ is one of $0,1,\ldots,N/e-1$.  For
\eqref{plt:eq:mul-binomial-top}, the term $r=0$ contributes nothing to
$t^N$ for $N\ge1$, the term $r=1$ contributes $\sigma u_N$, and each term
with $r\ge2$ involves only $u_n$ with $n\le N-(r-1)e\le N-e$.

(3) By \cref{plt:lem:mul-summable}(2),
\[
 (1+U)^\sigma(1+U)^\tau
 =\sum_{r,s\ge0}\binom{\sigma}{r}\binom{\tau}{s}U^{r+s}
 =\sum_{m\ge0}\Bigl(\sum_{r+s=m}\binom{\sigma}{r}\binom{\tau}{s}\Bigr)U^m,
\]
the regrouping being legitimate because for each fixed $t$-block only
finitely many pairs $(r,s)$ contribute.  Vandermonde's convolution
\cite{graham-knuth-patashnik}
\[
 \sum_{r+s=m}\binom{\sigma}{r}\binom{\tau}{s}=\binom{\sigma+\tau}{m}
\]
holds for all $\sigma,\tau\in\K$: both sides are polynomial expressions in
$(\sigma,\tau)$ with coefficients in $\RationalNumbers$; they agree for all
pairs of nonnegative integers by comparing coefficients in
$(1+z)^{\sigma}(1+z)^{\tau}=(1+z)^{\sigma+\tau}$ in
$\RationalNumbers[z]$; a polynomial in two variables over the infinite domain
$\RationalNumbers$ vanishing on $\NaturalNumbers^2$ is the zero polynomial;
and $\RationalNumbers\subseteq\K$ because $\operatorname{char}\K=0$, so the
identity persists in $\K$.  This proves
\eqref{plt:eq:mul-binomial-group}; $(1+U)^0=1$ is immediate.  For
$k\in\NaturalNumbers$ we have $\binom{k}{r}=0$ for $r>k$, so
\eqref{plt:eq:mul-binomial-def} is the finite binomial expansion of the
$k$-fold product in the commutative ring $\PLpow$.  Invertibility follows by
taking $\tau=-\sigma$.

(4) By \eqref{plt:eq:mul-binomial-finite} it suffices to bound
$\degL[t^N]U^r$ for $1\le r\le\Floor{N/e}$ (the term $r=0$ contributes only
at $N=0$, where the bound reads $0\le0$).  By
\eqref{plt:eq:mul-power-coefficient} and \eqref{plt:eq:mul-degL-laws},
\[
 \degL[t^N]U^r
 \le\max_{\substack{n_1+\cdots+n_r=N\\ n_i\ge e}}
    \sum_{i=1}^{r}(\alpha n_i+\beta)
 =\alpha N+r\beta
 \le\alpha N+\Floor{N/e}\beta^+
 \le\Bigl(\alpha+\frac{\beta^+}{e}\Bigr)N,
\]
where the second inequality uses $r\le\Floor{N/e}$ when $\beta\ge0$ and
$r\beta\le0\le\Floor{N/e}\beta^+$ when $\beta<0$.  Taking the maximum over
$r$ gives \eqref{plt:eq:mul-binomial-profile}.
\end{proof}

\begin{corollary}[Formal roots]
\label{plt:cor:mul-roots}
With $U$ as above and $p,q\in\IntegerNumbers$, $q\ge1$,
\[
 \bigl((1+U)^{p/q}\bigr)^{q}=(1+U)^{p},
\]
the right-hand side being the ordinary $p$-th power (or the reciprocal of the
$|p|$-th power if $p<0$).  In particular $(1+U)^{1/q}$ is a $q$-th root of
$1+U$ in $\PLpow$, and it is the unique one with constant block $1$.
\end{corollary}

\begin{proof}
Iterating \eqref{plt:eq:mul-binomial-group} $q-1$ times gives
$((1+U)^{p/q})^q=(1+U)^{qp/q}=(1+U)^p$, and
\cref{plt:thm:mul-binomial}(3) identifies the right-hand side with the
ordinary power.  For uniqueness, suppose $V,W\in\PLpow$ have constant block
$1$ and $V^q=W^q=1+U$.  Then $(V-W)\sum_{i=0}^{q-1}V^iW^{q-1-i}=0$; the
second factor has constant block $q\ne0$ (characteristic zero), hence is
nonzero, and $\PLpow$ is a domain by \cref{plt:cor:mul-domain}, so $V=W$.
\end{proof}

\begin{remark}[Infinite products]
\label{plt:rmk:mul-infinite-products}
If $(U_r)_{r\ge1}$ lies in $\PLpow$ with $\ordt U_r\to+\infty$, the partial
products $P_R=\prod_{r\le R}(1+U_r)$ converge $t$-adically.  Indeed, given
$N$ choose $R_N$ with $\ordt U_r\ge N$ for $r>R_N$; for $R>R_N$,
\[
 P_R-P_{R_N}=P_{R_N}\Bigl(\prod_{R_N<r\le R}(1+U_r)-1\Bigr),
\]
and expanding the finite product shows that the bracket is a sum of products
of at least one $U_r$ with $r>R_N$, hence of $\ordt\ge N$; since
$\ordt P_{R_N}\ge0$ we get $\ordt(P_R-P_{R_N})\ge N$.  Thus $(P_R)$ is
$t$-adically Cauchy in the complete ring $\PLpow$ and
$\prod_{r\ge1}(1+U_r)$ is well defined.  On $\PLlau$ the same argument
applies once one notes that $\ordt U_r\to+\infty$ forces all but finitely
many factors to lie in $\PLpow$, so the partial products have a common lower
support bound.
\end{remark}

\section{Three worked computations}
\label{plt:sec:mul-worked}

\begin{example}[A Laurent product through four blocks, with degree
bookkeeping]
\label{plt:ex:mul-worked-laurent}
Over $\K=\RationalNumbers$ take
\begin{align*}
 F&=3t^{-2}+(L-1)t^{-1}+(L^2+2)+2L\,t+(L^2-L)t^2
   +\FormalBigOAt{t^{3}}{t},\\
 G&=t^{-1}+(2-L)+(L+1)t+L^2t^2+\FormalBigOAt{t^{3}}{t}.
\end{align*}
Here $\ordt F=-2$ and $\ordt G=-1$; $F$ is known through $\Trunc{F}{3}$ and
$G$ through $\Trunc{G}{3}$.  By the precision rule the product is known
through
\[
 \Trunc{FG}{\min(3+(-1),\,3+(-2))}=\Trunc{FG}{1},
\]
that is, through the block $t^{0}$ inclusive --- four blocks in all, starting
at $t^{-3}$.  The extra block $[t^2]F$ is wasted: it can influence only
$[t^{1}](FG)$ and beyond, for which $[t^{3}]G$ would also be needed.

Writing $A_n=[t^n]F$ and $B_m=[t^m]G$, the convolution
\eqref{plt:eq:mul-cauchy} gives
\begin{align*}
 C_{-3}&=A_{-2}B_{-1}=3,\\
 C_{-2}&=A_{-2}B_{0}+A_{-1}B_{-1}
        =3(2-L)+(L-1)=5-2L,\\
 C_{-1}&=A_{-2}B_{1}+A_{-1}B_{0}+A_{0}B_{-1}\\
       &=3(L+1)+(-L^2+3L-2)+(L^2+2)=6L+3,\\
 C_{0}&=A_{-2}B_{2}+A_{-1}B_{1}+A_{0}B_{0}+A_{1}B_{-1}\\
      &=3L^2+(L^2-1)+(-L^3+2L^2-2L+4)+2L
      =-L^3+6L^2+3 .
\end{align*}
Hence
\begin{equation}\label{plt:eq:mul-worked-laurent}
 FG=3t^{-3}+(5-2L)t^{-2}+(6L+3)t^{-1}+(-L^3+6L^2+3)
   +\FormalBigOAt{t^{1}}{t}.
\end{equation}

Leading data: $\lm(F)=t^{-2}$ with $\lc(F)=3$, $\lm(G)=t^{-1}$ with
$\lc(G)=1$; \cref{plt:thm:mul-leading-law} predicts $\lm(FG)=t^{-3}$ and
$\lc(FG)=3$, confirmed by $C_{-3}=3$.

The degree bookkeeping is instructive.  The profiles are
$\dprofile{F}(-2)=0$, $\dprofile{F}(-1)=1$, $\dprofile{F}(0)=2$,
$\dprofile{F}(1)=1$ and $\dprofile{G}(-1)=0$, $\dprofile{G}(0)=1$,
$\dprofile{G}(1)=1$, $\dprofile{G}(2)=2$.  Then:

\begin{center}
\begin{tabular}{@{}ccccc@{}}
\hline
$N$ & maximizing pairs $P_N$ & bound $M$ & top coefficient
 \eqref{plt:eq:mul-top-coefficient} & $\dprofile{FG}(N)$ \\
\hline
$-3$ & $(-2,-1)$ & $0$ & $3\cdot1=3$ & $0$ \\
$-2$ & $(-2,0),\ (-1,-1)$ & $1$ & $3(-1)+1\cdot1=-2$ & $1$ \\
$-1$ & $(-1,0),\ (0,-1)$ & $2$ & $1\cdot(-1)+1\cdot1=0$ & $1<2$ \\
$0$  & $(0,0)$ & $3$ & $1\cdot(-1)=-1$ & $3$ \\
\hline
\end{tabular}
\end{center}

At $N=-2$ two pairs compete and their contributions $-3$ and $+1$ do not
cancel, so the bound is attained, with top coefficient $-2$: indeed
$C_{-2}=-2L+5$.  At $N=-1$ two pairs compete and cancel exactly, so the
degree drops from the bound $2$ to the actual value $1$: indeed
$C_{-1}=6L+3$.  At $N=0$ the maximizing pair is unique, so
\cref{plt:cor:mul-degree-equality}(1) forces attainment, with top
coefficient $-1$: indeed $C_0=-L^3+\cdots$.  This is exactly the phenomenon
\cref{plt:thm:mul-degree-bound} isolates; no inspection of the finished
polynomial is needed to predict it.
\end{example}

\begin{example}[A square root through four blocks, verified by squaring]
\label{plt:ex:mul-worked-sqrt}
Let
\[
 U=2Lt+(L^2+2)t^2+4L\,t^3+\FormalBigOAt{t^{4}}{t}\in t\PLpow,
\]
so $e=1$ in \cref{plt:thm:mul-binomial}, and take $\sigma=\tfrac12$, with
\[
 \binom{1/2}{0}=1,\quad
 \binom{1/2}{1}=\tfrac12,\quad
 \binom{1/2}{2}=-\tfrac18,\quad
 \binom{1/2}{3}=\tfrac1{16}.
\]
By \eqref{plt:eq:mul-binomial-finite} only $r\le N$ contributes to $t^N$, so
through $t^3$ we need
\[
 U^2=4L^2t^2+(4L^3+8L)t^3+\FormalBigOAt{t^{4}}{t},
 \qquad
 U^3=8L^3t^3+\FormalBigOAt{t^{4}}{t}.
\]
Block by block:
\begin{align*}
 [t^1](1+U)^{1/2}&=\tfrac12(2L)=L,\\
 [t^2](1+U)^{1/2}&=\tfrac12(L^2+2)-\tfrac18(4L^2)
   =\tfrac{L^2}2+1-\tfrac{L^2}2=1,\\
 [t^3](1+U)^{1/2}&=\tfrac12(4L)-\tfrac18(4L^3+8L)+\tfrac1{16}(8L^3)
   =2L-\tfrac{L^3}2-L+\tfrac{L^3}2=L .
\end{align*}
Hence
\begin{equation}\label{plt:eq:mul-worked-sqrt}
 (1+U)^{1/2}=1+L\,t+t^2+L\,t^3+\FormalBigOAt{t^{4}}{t}.
\end{equation}
Two logarithmic cancellations occur, one at each of $t^2$ and $t^3$; the
degree bound \eqref{plt:eq:mul-binomial-profile} with $\alpha=1$,
$\beta^{+}=0$, $e=1$ predicts only $\dprofile{(1+U)^{1/2}}(N)\le N$, and the
actual
profile $0,1,0,1$ is strictly smaller at $N=2,3$.  A degree bound is an upper
bound; it is not a prediction.

\emph{Verification.}  \cref{plt:cor:mul-roots} says the square of
\eqref{plt:eq:mul-worked-sqrt} must reproduce $1+U$ through $t^3$.  Squaring
by \eqref{plt:eq:mul-cauchy}:
\[
 [t^0]=1,\quad
 [t^1]=2L,\quad
 [t^2]=2\cdot1+L^2=L^2+2,\quad
 [t^3]=2L+2L=4L,
\]
which are exactly the blocks $1,u_1,u_2,u_3$ of $1+U$.  Note the precision
accounting: squaring a series known through $\Trunc{\cdot}{4}$ with
$\ordt=0$ returns a result known through $\Trunc{\cdot}{4}$, so the check is
available at every block computed and at none beyond.
\end{example}

\begin{example}[A product in the iterated field]
\label{plt:ex:mul-worked-iterated}
In $\PLit$ the coefficient convolution is finite at the inner level too.
Take the two blocks
\[
 f=\frac1{L+1}=L^{-1}-L^{-2}+L^{-3}-\cdots,
 \qquad
 g=\frac1{L-1}=L^{-1}+L^{-2}+L^{-3}+\cdots
\]
in $\K((L^{-1}))$.  For $r\ge2$, the coefficient of $L^{-r}$ in $fg$ is a sum
over the finitely many decompositions $r=i+j$ with $i,j\ge1$:
\[
 \kappa_{-r}(fg)=\sum_{i=1}^{r-1}(-1)^{i+1}\cdot1
 =\begin{cases}1,&r\ \text{even},\\0,&r\ \text{odd},\end{cases}
\]
so $fg=L^{-2}+L^{-4}+L^{-6}+\cdots=(L^2-1)^{-1}$, as it must be.  The
degree law is visible: $\degL f=\degL g=-1$, $\degL(fg)=-2$, and
$\lc_L(f)=\lc_L(g)=\lc_L(fg)=1$.  Consequently, for
$F=f\,t^{-1}+\cdots$ and $G=g+\cdots$ in $\PLit$ one gets
$\lm(FG)=t^{-1}L^{-2}$ and $\lc(FG)=1$ by
\cref{plt:thm:mul-leading-law}, with no need to compute any further block.
Observe that $\lm$ here carries a \emph{negative} power of $L$: this is
legitimate in $\PLrat$ and $\PLit$ and impossible in $\PLlau$, which is
exactly the content of \cref{plt:cor:mul-domain}.
\end{example}

\section{The analytic product rule}
\label{plt:sec:mul-analytic}

Everything above is formal algebra.  We now record the exact circumstances
under which the block convolution computes the asymptotic expansion of an
actual product of functions.  Nothing in this subsection asserts
convergence: an asymptotic expansion is a statement about remainders, and the
formal product theorem is used as an input, not replaced.

\begin{definition}[Polynomial--logarithmic asymptotic expansion]
\label{plt:def:mul-asymptotic}
Let $\K\subseteq\ComplexNumbers$, let $X_0>1$, and let $S=[X_0,+\infty)$ with
the real branch $L=\log X$.  A function $f\colon S\to\ComplexNumbers$ is
\emph{PL-asymptotic} to $F\in\PLlau$, written $f\approx F$, if for every
$N\in\IntegerNumbers$ there exist $M_N\in\NaturalNumbers$, $C_N>0$ and
$X_N\ge X_0$ such that
\begin{equation}\label{plt:eq:mul-asymptotic}
 \AbsoluteValue{f(X)-\bigl(\Trunc{F}{N}\bigr)(X)}
 \le C_N\,X^{-N}(\log X)^{M_N}
 \qquad(X\ge X_N),
\end{equation}
where $\bigl(\Trunc{F}{N}\bigr)(X)=\sum_{\ordt F\le n<N}p_n(\log X)X^{-n}$ is
the actual finite sum obtained by substituting the generators.  Equivalently,
$f-\Trunc{F}{N}=\BigOAt{X^{-N}(\log X)^{M_N}}{X\to+\infty,\ X\in S}$.
\end{definition}

% ed.: retitled: ``Basic properties of a scale'' in the differential part
% ed.: (\cref{plt:lem:dif-scale-basic}) is about the coefficient field of a
% ed.: power--log scale, not about the relation $\approx$ between a function
% ed.: and a series.  The two lemmas share nothing but the word ``basic''.
\begin{lemma}[Basic properties of the expansion relation $\approx$]
\label{plt:lem:mul-asymptotic-basic}
Let $f,g\colon S\to\ComplexNumbers$ and $F,G\in\PLlau$.
\begin{enumerate}
\item If \eqref{plt:eq:mul-asymptotic} holds for arbitrarily large $N$, it
holds for every $N$; so it suffices to verify it along a cofinal set of
cutoffs.
\item $f\approx F$ and $f\approx F'$ imply $F=F'$.
\item If $f\approx F$ and $g\approx G$, then $f+g\approx F+G$ and
$\lambda f\approx\lambda F$ for $\lambda\in\K$.
\item $f\approx0$ if and only if $f=\BigOAt{X^{-N}}{X\to+\infty}$ for every
$N$.
\end{enumerate}
\end{lemma}

\begin{proof}
(1) Let $N'<N$ and suppose \eqref{plt:eq:mul-asymptotic} holds at $N$.  Then
\[
 f-\Trunc{F}{N'}=\bigl(f-\Trunc{F}{N}\bigr)
   +\sum_{N'\le n<N}p_n(\log X)X^{-n},
\]
a finite sum; the first term is $\BigOAt{X^{-N}(\log
X)^{M_N}}{X\to+\infty}\subseteq\BigOAt{X^{-N'}(\log X)^{M_N}}{X\to+\infty}$
since $N\ge N'$, and each term of the sum is
$\BigOAt{X^{-N'}(\log X)^{\degL p_n}}{X\to+\infty}$.

(2) Suppose $D=F-F'\ne0$, with $n_1=\ordt D$ and leading block $p$ of degree
$d$ and leading coefficient $c\ne0$.  Take $N=n_1+1$ and subtract the two
instances of \eqref{plt:eq:mul-asymptotic}:
\[
 \bigl(\Trunc{D}{n_1+1}\bigr)(X)=p(\log X)X^{-n_1}
 =\BigOAt{X^{-n_1-1}(\log X)^{M}}{X\to+\infty}
\]
for some $M$.  But $|p(\log X)|\sim|c|(\log X)^{d}$, so the left side
divided by $X^{-n_1-1}(\log X)^M$ is $\sim|c|X(\log X)^{d-M}\to+\infty$, a
contradiction.

(3) Immediate from
$\Trunc{F+G}{N}=\Trunc{F}{N}+\Trunc{G}{N}$ and the triangle inequality.

(4) If $f\approx0$ then $\Trunc{0}{N}=0$ and $f=\BigOAt{X^{-N}(\log
X)^{M_N}}{X\to+\infty}$ for every $N$; since $(\log
X)^{M_N}=\LittleOAt{X}{X\to+\infty}$, this gives
$f=\BigOAt{X^{-N+1}}{X\to+\infty}$ for every $N$, hence
$f=\BigOAt{X^{-N}}{X\to+\infty}$ for every $N$.  The converse is trivial.
\end{proof}

\begin{theorem}[Analytic product rule]
\label{plt:thm:mul-analytic-product}
Let $\K\subseteq\ComplexNumbers$, $S=[X_0,+\infty)$ with the real logarithm,
and let $f,g\colon S\to\ComplexNumbers$ with $f\approx F$ and $g\approx G$
for $F,G\in\PLlau$.  Then
\[
 fg\approx FG,
\]
where $FG$ is the formal Cauchy product \eqref{plt:eq:mul-cauchy}.
Consequently the set
$\mathcal A(S)=\SetOf{f:\exists F\in\PLlau,\ f\approx F}$ is a
$\K$-algebra, the map $f\mapsto F$ is a well-defined $\K$-algebra
homomorphism $\mathcal A(S)\to\PLlau$, and its kernel is the ideal of
functions that are $\BigOAt{X^{-N}}{X\to+\infty}$ for every $N$.
\end{theorem}

\begin{proof}
First dispose of the degenerate case.  Any $g$ with $g\approx G$ satisfies
$g=\BigOAt{X^{c}(\log X)^{m}}{X\to+\infty}$ for some $c\in\IntegerNumbers$
and $m\in\NaturalNumbers$: take $N=1$ in \eqref{plt:eq:mul-asymptotic}, note
that $\Trunc{G}{1}$ is a finite sum of monomials $t^nL^j$, and set
$c=\max(-\ordt G,-1)$ (with $c=-1$ if $G=0$).  Hence, if $F=0$, then
\cref{plt:lem:mul-asymptotic-basic}(4) gives
$f=\BigOAt{X^{-P}}{X\to+\infty}$ for every $P$, so
$fg=\BigOAt{X^{-P+c}(\log X)^{m}}{X\to+\infty}$ for every $P$, whence
$fg\approx0=FG$ by \cref{plt:lem:mul-asymptotic-basic}(4) again; the case
$G=0$ is symmetric.  Assume from now on $F,G\ne0$ and set $a=\ordt F$,
$b=\ordt G$.

Fix $N\in\IntegerNumbers$ and put
\[
 P=\max(N-b,\;a),
 \qquad
 Q=\max(N-a,\;b).
\]
Write $f_P=(\Trunc{F}{P})(X)$, $r=f-f_P$, $g_Q=(\Trunc{G}{Q})(X)$,
$s=g-g_Q$.  By hypothesis
$r=\BigOAt{X^{-P}(\log X)^{m_1}}{X\to+\infty}$ and
$s=\BigOAt{X^{-Q}(\log X)^{m_2}}{X\to+\infty}$ for suitable exponents.  Since
$\Trunc{F}{P}$ is a finite sum of blocks with outer exponents $\ge a$,
\[
 f_P=\BigOAt{X^{-a}(\log X)^{m_3}}{X\to+\infty},
 \qquad
 g_Q=\BigOAt{X^{-b}(\log X)^{m_4}}{X\to+\infty}
\]
with $m_3=\max_{a\le n<P}\degL p_n$ and similarly for $m_4$ (the estimates
being vacuously true if the sums are empty).  Therefore
\[
 fg-f_Pg_Q=f_Ps+rg_Q+rs
\]
is, term by term,
% ed.: the three-term estimate was set on one display line, overfilling the
% ed.: text block by about 53pt; broken across two lines.
\begin{align*}
 &\BigOAt{X^{-(a+Q)}(\log X)^{m_3+m_2}}{X\to+\infty}
 +\BigOAt{X^{-(P+b)}(\log X)^{m_1+m_4}}{X\to+\infty}\\
 &\qquad
 +\BigOAt{X^{-(P+Q)}(\log X)^{m_1+m_2}}{X\to+\infty}.
\end{align*}
Now $a+Q\ge a+(N-a)=N$, $P+b\ge(N-b)+b=N$, and $P+Q\ge a+(N-a)=N$; hence
\begin{equation}\label{plt:eq:mul-analytic-step1}
 fg-f_Pg_Q=\BigOAt{X^{-N}(\log X)^{M'}}{X\to+\infty}
\end{equation}
for some $M'$.

Next, $f_Pg_Q$ is the value at $X$ of the finite polynomial--logarithmic
expression $\bigl(\Trunc{F}{P}\bigr)\bigl(\Trunc{G}{Q}\bigr)\in\PLlau$,
because substitution of the generators $t=X^{-1}$, $L=\log X$ into a finite
sum of monomials is a ring homomorphism into functions on $S$.  Since
$P\ge N-b$ and $Q\ge N-a$, \cref{plt:cor:mul-truncated-product} gives
\[
 \Trunc{\bigl(\Trunc{F}{P}\bigr)\bigl(\Trunc{G}{Q}\bigr)}{N}
 =\Trunc{FG}{N}.
\]
Hence $\bigl(\Trunc{F}{P}\bigr)\bigl(\Trunc{G}{Q}\bigr)-\Trunc{FG}{N}$ is a
\emph{finite} sum of monomials $c\,t^nL^j$ with $n\ge N$, so its value at
$X$ is $\BigOAt{X^{-N}(\log X)^{M''}}{X\to+\infty}$ with $M''$ the largest
$j$ occurring.  Combining with \eqref{plt:eq:mul-analytic-step1},
\[
 fg-\bigl(\Trunc{FG}{N}\bigr)(X)
 =\BigOAt{X^{-N}(\log X)^{\max(M',M'')}}{X\to+\infty},
\]
which is \eqref{plt:eq:mul-asymptotic} for $fg$ and $FG$.  As $N$ was
arbitrary, $fg\approx FG$.

The algebra statements now follow: $\mathcal A(S)$ is closed under $+$,
scalars and $\cdot$ by \cref{plt:lem:mul-asymptotic-basic}(3) and the above;
the assignment $f\mapsto F$ is well defined by
\cref{plt:lem:mul-asymptotic-basic}(2) and is a homomorphism by the same two
statements; its kernel is described by
\cref{plt:lem:mul-asymptotic-basic}(4).
\end{proof}

\begin{remark}[Which hypotheses are actually used, and the sectorial version]
\label{plt:rmk:mul-analytic-hypotheses}
The proof uses only three things: that $f$ and $g$ have remainders bounded by
$X^{-N}$ times \emph{some} power of $\log X$; that finitely many blocks are
involved at each stage; and the formal finite-determination theorem
\cref{plt:thm:mul-finite-determination}, which is what allows the two
truncation cutoffs $P$ and $Q$ to be chosen independently of the shape of
$F$ and $G$.  No convergence and no uniformity in $L$ is used or obtained.

The same proof works verbatim on an unbounded set
$S\subseteq\SetOf{z:\AbsoluteValue{\arg z}\le\theta}$ with
$0<\theta<\pi$ and the principal logarithm, with $X^{-N}$ replaced by
$\AbsoluteValue{X}^{-N}$ and $(\log X)^{M}$ by
$(\log\AbsoluteValue{X})^{M}$; the only additional input is
$\AbsoluteValue{\log X}\le\log\AbsoluteValue{X}+\theta
\le2\log\AbsoluteValue{X}$ for $\AbsoluteValue{X}\ge\EulerE^{\theta}$, so
that a bound in $\AbsoluteValue{\log X}$ and a bound in
$\log\AbsoluteValue{X}$ differ only by a constant factor.  Uniformity claims
on a sector must be stated for a \emph{closed} subsector; the branch of
$\log$ is part of the statement and cannot be recovered from the formal
data.
\end{remark}

\begin{warningbox}
Three cautions, each corresponding to an overstatement found in the sources.
\begin{enumerate}
\item \emph{A formal identity is not an analytic theorem.}
\cref{plt:thm:mul-analytic-product} presupposes that $f$ and $g$
\emph{already have} expansions; it does not produce one.  The
homomorphism $f\mapsto F$ has a large kernel --- $\EulerE^{-X}\approx0$ ---
so a formal computation determines an actual function at best up to a flat
error.
\item \emph{The formal $\FormalBigOAt{t^N}{t}$ carries no uniformity in $L$.}
Asserting $F=G+\FormalBigOAt{t^N}{t}$ says only that the blocks agree below
outer order $N$.  In particular the constant $C_N$ and the exponent $M_N$ of
\eqref{plt:eq:mul-asymptotic} are not supplied by any formal statement and
must be produced separately.
\item \emph{The three $\BigO$-roles are distinct.}  The formal tail
$\FormalBigOAt{t^N}{t}$, the analytic estimate
$\BigOAt{X^{-N}(\log X)^{M}}{X\to+\infty}$ of
\eqref{plt:eq:mul-asymptotic}, and the complexity bound
$\BigO\bigl((N-a-b)^2d^2\bigr)$ of \cref{plt:rmk:mul-cost} are three
different assertions and none follows from the others.  A bare $\BigO$ for an
asymptotic claim, with the regime left to context, is never acceptable.
\end{enumerate}
\end{warningbox}

\begin{summarybox}
Multiplication is Cauchy convolution in $t=X^{-1}$ with exact coefficient
arithmetic in $L=\log X$.  Every output block below a cutoff is a finite
bilinear expression in finitely many input blocks, with the exact range
$a\le n\le N-b$, $b\le m\le N-a$; the resulting precision rule is
$\Trunc{FG}{\min(P+b,\,Q+a)}$ and it cannot be improved.  All four ambient
classes are closed and are integral domains, the leading monomial and leading
scalar are multiplicative, and the rank-two valuation
$v(F)=(\ordt F,-\degL p_{\ordt F})$ is additive.  The log-degree profile
obeys the max-plus bound $\dprofile{FG}\le\dprofile{F}\ast\dprofile{G}$,
with equality exactly when the maximizing pairs do not cancel; affine
profiles are stable under products, powers, and binomial powers, but not
under $t$-adic limits.  Binomial powers $(1+U)^\sigma$ are defined for every
$\sigma\in\K$ by a summable family, are finitely determined at every order,
and satisfy the exponent law.  Under printed remainder hypotheses on a stated
domain and branch, the formal product computes the asymptotic expansion of
the analytic product --- and nothing more.
\end{summarybox}

% ---- fragment 04-division ----
\chapter{Units, reciprocals, and division}
\label{plt:sec:division}

Addition and multiplication never leave \(\PLlau=\K[L]((t))\).  Division is the
first operation that does.  This chapter settles exactly when it does not --
the unit criterion -- and then follows the three escape routes that the
criterion forces: an accidental exact quotient inside \(\PLlau\), the
rational-log field \(\PLrat=\K(L)((t))\), and the iterated Laurent-log field
\(\PLit=\K((L^{-1}))((t))\).  Each of the three is a genuinely different
object, and conflating them is the most frequent error in informal
polynomial--logarithmic computation.

Throughout, \(\K\) is a field of characteristic zero, \(t\DefinitionEquals
X^{-1}\) and \(L\DefinitionEquals\log X\) are independent formal generators, and
a nonzero element of \(\PLlau\) is written
\begin{equation}\label{plt:eq:div-normal-form}
  F=\sum_{n\ge n_{0}}p_{n}(L)t^{n},
  \qquad p_{n}\in\K[L],\quad p_{n_{0}}\ne0,
\end{equation}
so that \(\ordt F=n_{0}\).  We call \(p_{\ordt F}\) the \emph{leading block} of
\(F\); it is a polynomial, not a scalar, and it must be distinguished from the
leading monomial \(\lm F=t^{n_{0}}L^{d}\) and from the scalar leading
coefficient \(\lc F\in\K^{\times}\), where \(d=\degL p_{n_{0}}\).  Truncation is
always the exclusive operator \(\Trunc{F}{N}=\sum_{n<N}p_{n}(L)t^{n}\), and
\(F=G+\FormalBigOAt{t^{N}}{t}\) abbreviates \(\ordt(F-G)\ge N\); it carries no
analytic content whatever.

\section{The exact unit criterion}
\label{plt:sec:div-units}

Two elementary facts do all the work.  Both are standard, but the second is the
step that three of the six sources compress into a single clause, and it is
precisely the step that decides the nonunit direction of the criterion.

\begin{lemma}[Units of the coefficient ring]\label{plt:lem:div-poly-units}
Let \(\K\) be a field.  Then \(\K[L]^{\times}=\K^{\times}\): a polynomial is
invertible in \(\K[L]\) if and only if it is a nonzero constant.
\end{lemma}

\begin{proof}
A nonzero constant \(c\) has inverse \(c^{-1}\).  Conversely, suppose
\(pq=1\) with \(p,q\in\K[L]\).  Since \(\K\) is a field, \(\K[L]\) is an
integral domain and \(\degL(pq)=\degL p+\degL q\).  Hence
\(\degL p+\degL q=\degL 1=0\), and as both degrees are nonnegative integers,
\(\degL p=\degL q=0\).  So \(p\in\K\), and \(p\ne0\) because \(pq=1\).
\end{proof}

\begin{lemma}[The Laurent ring is a domain and \(\ordt\) is additive]
\label{plt:lem:div-valuation}
Let \(F,G\in\PLlau\) be nonzero, with leading blocks \(p_{\ordt F}\) and
\(q_{\ordt G}\).  Then \(FG\ne0\),
\begin{equation}\label{plt:eq:div-valuation-additive}
  \ordt(FG)=\ordt F+\ordt G,
\end{equation}
and the leading block of \(FG\) is \(p_{\ordt F}\,q_{\ordt G}\).  In particular
\(\PLlau\) is an integral domain, and the same statements hold in \(\PLpow\),
in \(\K(L)((t))\), and in \(E((t))\) for any integral domain \(E\).
\end{lemma}

\begin{proof}
Write \(m=\ordt F\), \(r=\ordt G\).  The Cauchy convolution gives
\([t^{k}](FG)=\sum_{i+j=k}p_{i}q_{j}\), a finite sum because \(i\ge m\) and
\(j\ge r\) force \(m\le i\le k-r\).  For \(k<m+r\) every term has \(i<m\) or
\(j<r\) and therefore vanishes.  For \(k=m+r\) the only surviving pair is
\((i,j)=(m,r)\), so \([t^{m+r}](FG)=p_{m}q_{r}\), which is nonzero because
\(\K[L]\) is an integral domain.  Hence \(FG\ne0\) and both assertions hold.
\end{proof}

\begin{theorem}[Unit criterion]\label{plt:thm:div-unit-criterion}
Let \(F\in\PLlau\) be nonzero, written as in \cref{plt:eq:div-normal-form}.
The following are equivalent.
\begin{enumerate}
\item \(F\) is a unit of \(\PLlau\).
\item The leading block \(p_{n_{0}}\) is a nonzero constant, that is
      \(p_{n_{0}}\in\K^{\times}\).
\item \(F=c\,t^{n_{0}}(1+U)\) for some \(c\in\K^{\times}\) and some
      \(U\in t\PLpow\).
\end{enumerate}
For \(F\in\PLpow\) with \(n_{0}=0\), these are equivalent to \(p_{0}\in
\K^{\times}\), and \(F^{-1}\) then lies in \(\PLpow\).
\end{theorem}

\begin{proof}
\((2)\Leftrightarrow(3)\).  If \(p_{n_{0}}=c\in\K^{\times}\), put
\(U\DefinitionEquals c^{-1}t^{-n_{0}}F-1\); every block of \(c^{-1}t^{-n_{0}}F\)
of index \(\ge1\) is a polynomial and the index-\(0\) block is \(1\), so
\(U\in t\PLpow\).  Conversely, in \(F=ct^{n_{0}}(1+U)\) with \(U\in t\PLpow\)
the block of index \(n_{0}\) is \(c\).

\((3)\Rightarrow(1)\).  Set
\begin{equation}\label{plt:eq:div-geometric}
  G\DefinitionEquals c^{-1}t^{-n_{0}}\sum_{k\ge0}(-U)^{k}.
\end{equation}
Because \(\ordt(U^{k})\ge k\), for each \(N\) only the terms with \(k<N\)
contribute to blocks of index \(<N\); the family is therefore formally summable
and \(G\in\PLlau\), with polynomial blocks because \(\PLpow\) is closed under
multiplication.  For every \(N\ge0\),
\begin{align*}
  \Trunc{\Bigl((1+U)\sum_{k\ge0}(-U)^{k}\Bigr)}{N}
  &=\Trunc{\Bigl((1+U)\sum_{k<N}(-U)^{k}\Bigr)}{N}\\
  &=\Trunc{\bigl(1-(-U)^{N}\bigr)}{N}
   =\Trunc{1}{N},
\end{align*}
using the telescoping identity \((1+U)\sum_{k<N}(-U)^{k}=1-(-U)^{N}\) and
\(\ordt((-U)^{N})\ge N\).  Since two elements of \(\PLlau\) with equal exclusive
truncations at every order are equal, \((1+U)\sum_{k\ge0}(-U)^{k}=1\), whence
\(FG=1\).

% ed.: S5 and S6 assert the nonunit direction by citing the general
% power-series unit criterion; the additivity step that makes it valid on the
% Laurent ring is supplied here explicitly.
\((1)\Rightarrow(2)\).  Suppose \(FG=1\) with \(G\in\PLlau\).  Then \(G\ne0\),
and \cref{plt:lem:div-valuation} gives \(\ordt F+\ordt G=\ordt 1=0\) together
with \(p_{\ordt F}\,q_{\ordt G}=1\) in \(\K[L]\), where \(q_{\ordt G}\) is the
leading block of \(G\).  By \cref{plt:lem:div-poly-units},
\(p_{\ordt F}\in\K^{\times}\).

The statement for \(\PLpow\) is the case \(n_{0}=0\) of the above, together
with the observation that the \(G\) constructed in \cref{plt:eq:div-geometric}
lies in \(\PLpow\) when \(n_{0}=0\).
\end{proof}

\begin{corollary}[The unit group]\label{plt:cor:div-unit-group}
The group of units of \(\PLlau\) is
\begin{equation}\label{plt:eq:div-unit-group}
  \PLlau^{\times}
  =\SetOf{c\,t^{m}(1+U):m\in\IntegerNumbers,\ c\in\K^{\times},\
          U\in t\PLpow},
\end{equation}
and the factorisation is unique, so that
\(\PLlau^{\times}\cong\IntegerNumbers\times\K^{\times}\times(1+t\PLpow)\) as
abelian groups.  Moreover \(\PLpow^{\times}=\K^{\times}\cdot(1+t\PLpow)\).
\end{corollary}

\begin{proof}
The description of the set is \cref{plt:thm:div-unit-criterion}.  Uniqueness:
in \(F=ct^{m}(1+U)\) one has \(m=\ordt F\) by
\cref{plt:lem:div-valuation}, then \(c\) is the leading block, then
\(U=c^{-1}t^{-m}F-1\).
% ed.: the direct-product step was justified by "the three subgroups pairwise
% intersect trivially", which does not imply an internal direct product for
% three subgroups (the Klein four-group refutes it); the isomorphism is read
% off instead from surjectivity plus the uniqueness just proved.
The set \(1+t\PLpow\) is a subgroup of \(\PLlau^{\times}\): it is closed under
multiplication, and by \cref{plt:thm:div-unit-criterion} the inverse of
\(1+U\) is a unit of \(\PLpow\) with leading block \(1\), hence again of the
form \(1+U'\) with \(U'\in t\PLpow\).  Since \(\PLlau^{\times}\) is abelian,
\[
  \mu:\IntegerNumbers\times\K^{\times}\times(1+t\PLpow)
  \longrightarrow\PLlau^{\times},
  \qquad
  \mu\bigl(m,c,1+U\bigr)=c\,t^{m}(1+U),
\]
is a group homomorphism.  It is surjective by
\cref{plt:thm:div-unit-criterion} and injective because
\(\mu(m,c,1+U)=1=1\cdot t^{0}\cdot(1+0)\) forces \(m=0\), \(c=1\) and \(U=0\)
by the uniqueness just proved.  Hence \(\mu\) is an isomorphism.
\end{proof}

\begin{example}[A nonzero nonunit]\label{plt:ex:div-nonunit}
The element \(L+1\in\PLpow\) is nonzero, but it is not a unit of \(\PLlau\).
Indeed \(\ordt(L+1)=0\) and its leading block is \(L+1\), which is not
constant; \cref{plt:thm:div-unit-criterion} applies.  Unwinding the proof gives
the direct argument: if \((L+1)G=1\) with \(G=\sum_{n\ge n_{1}}q_{n}(L)t^{n}\)
nonzero, then \(\ordt G=0\) and \((L+1)q_{0}=1\) in \(\K[L]\), which is
impossible because \(\degL\bigl((L+1)q_{0}\bigr)=1+\degL q_{0}\ge1\).  The
same argument shows that \(L\), \(L+t\), \(L^{2}+t^{3}\) and
\(t^{-5}(L-7)\) are nonzero nonunits.  Their reciprocals exist in \(\PLrat\):
\begin{equation}\label{plt:eq:div-Lplust}
  \frac{1}{L+t}
  =\frac1L-\frac{t}{L^{2}}+\frac{t^{2}}{L^{3}}-\cdots
  \in\PLrat\setminus\PLlau .
\end{equation}
\end{example}

\begin{remark}[The dominant-monomial fallacy]\label{plt:rmk:div-dominant}
Factoring out the dominant monomial is \emph{not} enough to invert.  Writing
\(L+t=L(1+t/L)\) and expanding the second factor geometrically already commits
one to the coefficient ring \(\K(L)\) or \(\K((L^{-1}))\): the step that leaves
\(\PLlau\) is the inversion of the scalar-free factor \(L\), not the infinite
\(t\)-expansion.  \Cref{plt:eq:div-Lplust} makes the point in one line -- the
failure is visible already in the block of index \(0\).
\end{remark}

\section{Existence, uniqueness, and the reciprocal recurrence}
\label{plt:sec:div-recurrence}

\Cref{plt:eq:div-geometric} proves existence but is a poor algorithm.  The
practical construction is triangular.

\begin{theorem}[Reciprocal of a unit]\label{plt:thm:div-reciprocal}
Let \(B=\sum_{n\ge0}B_{n}(L)t^{n}\in\PLpow\) with \(B_{0}=c\in\K^{\times}\).
Then there is exactly one \(C\in\PLpow\) with \(BC=1\), namely
\(C=\recip{B}=\sum_{n\ge0}C_{n}(L)t^{n}\) with
\begin{align}
  C_{0}&=c^{-1},\label{plt:eq:div-recip-c0}\\
  C_{n}&=-c^{-1}\sum_{j=1}^{n}B_{j}\,C_{n-j}\qquad(n\ge1),
        \label{plt:eq:div-recip-cn}
\end{align}
and every \(C_{n}\) lies in \(\K[L]\).  Moreover \(C\) is the unique reciprocal
of \(B\) in \(\PLlau\), and for a Laurent unit \(F=t^{m}B\) one has
\(F^{-1}=t^{-m}C\).
\end{theorem}

\begin{proof}
\emph{Existence.}  Define \(C_{n}\in\K[L]\) by
\cref{plt:eq:div-recip-c0,plt:eq:div-recip-cn}; the right-hand side of
\cref{plt:eq:div-recip-cn} involves only \(C_{0},\dots,C_{n-1}\), so the
definition is legitimate by strong induction, and each \(C_{n}\) is a
\(\K\)-linear combination of products of polynomials, hence a polynomial.  Put
\(C\DefinitionEquals\sum_{n\ge0}C_{n}t^{n}\in\PLpow\).  For \(n\ge1\),
\[
  [t^{n}](BC)=\sum_{j=0}^{n}B_{j}C_{n-j}
  =cC_{n}+\sum_{j=1}^{n}B_{j}C_{n-j}=0
\]
by \cref{plt:eq:div-recip-cn}, and \([t^{0}](BC)=cC_{0}=1\).  Hence \(BC=1\).

\emph{Uniqueness.}  If \(BC=B\widetilde C=1\) with
\(C,\widetilde C\in\PLlau\), then
\(C=C(B\widetilde C)=(CB)\widetilde C=\widetilde C\).  This argument uses only
associativity and commutativity, so uniqueness holds in \(\PLlau\) -- and in
any commutative ring -- not merely in \(\PLpow\).

\emph{Laurent case.}  \(t^{m}\) is a unit with inverse \(t^{-m}\), so
\((t^{m}B)(t^{-m}C)=BC=1\).

% ed.: \cref{plt:cor:div-quotient} invokes this theorem "over the coefficient
% ring E", a generality the statement did not record; the proof uses only
% commutativity and invertibility of B_0, so the clause is supplied here.
\emph{Generality.}  Nothing above uses more about the coefficient ring than
that it is commutative and that \(B_{0}\) is invertible in it.  Hence for any
commutative ring \(R\) and any \(B\in R[[t]]\) with \(B_{0}\in R^{\times}\),
the same recurrence produces the unique reciprocal of \(B\) in \(R[[t]]\), and
that reciprocal remains the unique one in \(R((t))\).
\end{proof}

\begin{corollary}[Quotient recurrence]\label{plt:cor:div-quotient}
Let \(E\) be a field, \(A=\sum_{n\ge0}A_{n}t^{n}\) and
\(B=\sum_{n\ge0}B_{n}t^{n}\) in \(E[[t]]\) with \(B_{0}\ne0\).  Then
\(Q=A/B=\sum_{n\ge0}Q_{n}t^{n}\in E[[t]]\) is the unique solution of \(BQ=A\),
and
\begin{equation}\label{plt:eq:div-quotient-recurrence}
  Q_{n}=B_{0}^{-1}\Bigl(A_{n}-\sum_{j=1}^{n}B_{j}Q_{n-j}\Bigr),
  \qquad n\ge0,
\end{equation}
the sum being empty when \(n=0\).  When \(E=\K[L]\) is used instead of a field,
\cref{plt:eq:div-quotient-recurrence} still produces polynomial blocks provided
\(B_{0}\in\K^{\times}\), and then \(Q\in\PLpow\).
\end{corollary}

\begin{proof}
Existence and uniqueness follow from \cref{plt:thm:div-reciprocal} applied over
the coefficient ring \(E\) (or \(\K[L]\), where \(B_{0}\in\K^{\times}\) is
invertible): \(Q\DefinitionEquals A\,\recip{B}\) solves \(BQ=A\), and if
\(BQ=BQ'\) then \(Q=Q'\) after multiplying by \(\recip{B}\).  Comparing
coefficients of \(t^{n}\) in \(BQ=A\) gives
\(B_{0}Q_{n}+\sum_{j=1}^{n}B_{j}Q_{n-j}=A_{n}\), which is
\cref{plt:eq:div-quotient-recurrence}.  Since
\cref{plt:eq:div-quotient-recurrence} determines \(Q_{n}\) from
\(Q_{0},\dots,Q_{n-1}\), it also proves uniqueness directly.
\end{proof}

The next statement is the one an implementation actually needs: it says that a
finite computation with truncated inputs returns exactly the truncation of the
true quotient, and that the truncated quotient is characterised without
reference to the recurrence at all.

\begin{theorem}[Correctness of truncated division]
\label{plt:thm:div-truncated}
Let \(B\in\PLpow\) with \(B_{0}\in\K^{\times}\), let \(A\in\PLpow\), and let
\(N\ge0\).  Let \(Q_{0},\dots,Q_{N-1}\) be produced by
\cref{plt:eq:div-quotient-recurrence} and set
\(P\DefinitionEquals\sum_{n<N}Q_{n}t^{n}\).  Then
\begin{equation}\label{plt:eq:div-truncated-residual}
  BP=A+\FormalBigOAt{t^{N}}{t},
\end{equation}
and \(P\) is the \emph{only} element of \(\K[L][t]\) of \(t\)-degree \(<N\)
satisfying \cref{plt:eq:div-truncated-residual}.  Consequently
\(P=\Trunc{(A/B)}{N}\).
\end{theorem}

\begin{proof}
For \(n<N\) the coefficient \([t^{n}](BP)=\sum_{j=0}^{n}B_{j}Q_{n-j}\) involves
only \(Q_{0},\dots,Q_{n}\), all of which are among the computed blocks, and it
equals \(A_{n}\) by \cref{plt:eq:div-quotient-recurrence}.  This is
\cref{plt:eq:div-truncated-residual}.

Uniqueness: if \(P'\) is another such polynomial, put \(D=P-P'\), a polynomial
in \(t\) of degree \(<N\) with \(\ordt(BD)\ge N\).  If \(D\ne0\) then
\cref{plt:lem:div-valuation} gives \(\ordt(BD)=\ordt B+\ordt D=\ordt D<N\), a
contradiction; so \(D=0\).

Finally, \(B\bigl(A/B-P\bigr)=A-BP\) has \(\ordt\ge N\), and \(\ordt B=0\), so
\(\ordt(A/B-P)\ge N\) by \cref{plt:lem:div-valuation}.  As \(P\) has
\(t\)-degree \(<N\), this says exactly \(P=\Trunc{(A/B)}{N}\).
\end{proof}

% ed.: S6 states this correctness theorem with its local inclusive truncation
% [F]_{<=N}; the statement above is the exclusive normalisation demanded by the
% catalogue, and the two differ by exactly one block.
\begin{remark}[Truncation convention]\label{plt:rmk:div-truncation-convention}
Report S6 states \cref{plt:thm:div-truncated} in the inclusive form
\[
  B\sum_{n\le N}Q_{n}t^{n}=A+\FormalBigOAt{t^{N+1}}{t}.
\]
Because \([F]_{\le N}=\Trunc{F}{N+1}\), that statement is the case
\(N\mapsto N+1\) of the one above.  Anyone transferring a precision budget
between the two
conventions must account for the one-block shift; it is the commonest silent
off-by-one in this subject.
\end{remark}

\begin{proposition}[Exact input requirement: relative precision is preserved]
\label{plt:prop:div-relative-precision}
Let \(E\) be a field, and let \(A,B\in E((t))\) with \(A\ne0\ne B\) (the case
\(A=0\) being trivial).  Put \(a=\ordt A\) and \(b=\ordt B\), so that
\(\ordt(A/B)=a-b\).  For every \(M\ge0\), the first \(M\) blocks of \(A/B\),
namely the coefficients of \(t^{a-b},\dots,t^{a-b+M-1}\), are determined by the
first \(M\) blocks of \(A\) and the first \(M\) blocks of \(B\), and by nothing
else.  Equivalently, in absolute cutoffs,
\begin{equation}\label{plt:eq:div-guard}
  \Trunc{(A/B)}{N}
  \quad\text{is determined by}\quad
  \Trunc{A}{\,N+b\,}
  \ \text{and}\
  \Trunc{B}{\,N+2b-a\,}.
\end{equation}
\end{proposition}

\begin{proof}
Write \(\widetilde B=t^{-b}B\), so \(\ordt\widetilde B=0\) and
\(\widetilde B_{j}=B_{j+b}\).  Then \(A/B=t^{-b}\,(A/\widetilde B)\), and
\(G\DefinitionEquals A/\widetilde B\) has \(\ordt G=a\).  Applying
\cref{plt:eq:div-quotient-recurrence} to \(G\) after shifting indices by \(a\),
the block \(G_{a+i}\) is a function of \(A_{a},\dots,A_{a+i}\) and
\(\widetilde B_{0},\dots,\widetilde B_{i}\) alone.  Taking \(i<M\) gives the
first assertion, since \(\widetilde B_{0},\dots,\widetilde B_{M-1}\) are the
first \(M\) blocks of \(B\).  For \cref{plt:eq:div-guard}, note first that
\(\Trunc{(A/B)}{N}=0\) when \(N\le a-b\), so assume \(N>a-b\).  The requested
blocks of \(A/B\) are then those of index \(<N\), that is the first
\(M=N-(a-b)\) blocks; these need \(A\) through index \(a+M-1=N+b-1\) and \(B\)
through index \(b+M-1=N+2b-a-1\).
\end{proof}

% ed.: S6 states a "guard precision" principle in words ("compute the quotient
% beyond the final requested block") but never says by how much; the exact
% guard is the content of the proposition above.
\begin{remark}[Guard precision]\label{plt:rmk:div-guard}
\Cref{plt:prop:div-relative-precision} is the precise form of the informal rule
that a quotient must be computed ``beyond the requested block''.  The guard is
\(b=\ordt B\) blocks on the numerator and \(2b-a\) on the denominator; it is
zero exactly when \(a=b=0\), and it is what a symbolic system must propagate
through an expression tree instead of a single global cutoff.
% ed.: the negative-guard case was explained by "the denominator has a Laurent
% principal part", which accounts only for the numerator guard b; the
% denominator guard 2b-a is negative under the unrelated condition a > 2b.
A negative guard means that fewer input blocks suffice.  On the numerator this
happens exactly when \(b<0\), that is when the denominator has a Laurent
principal part; on the denominator it happens exactly when \(a>2b\), for
instance when a numerator of high \(t\)-order is divided by a unit.
\end{remark}

\begin{example}[A reciprocal computed to five blocks]
\label{plt:ex:div-reciprocal-worked}
Let
\begin{equation}\label{plt:eq:div-recip-example-B}
  B=1+(L+1)t+L^{2}t^{2}+(L-2)t^{3}\in\PLpow .
\end{equation}
Here \(B_{0}=1\in\K^{\times}\) and \(B_{n}=0\) for \(n\ge4\).
\Cref{plt:eq:div-recip-c0,plt:eq:div-recip-cn} give
\begin{align*}
  C_{0}&=1,\\
  C_{1}&=-(L+1),\\
  C_{2}&=-\bigl((L+1)C_{1}+L^{2}C_{0}\bigr)=2L+1,\\
  C_{3}&=-\bigl((L+1)C_{2}+L^{2}C_{1}+(L-2)C_{0}\bigr)
        =L^{3}-L^{2}-4L+1,\\
  C_{4}&=-\bigl((L+1)C_{3}+L^{2}C_{2}+(L-2)C_{1}\bigr)
        =-L^{4}-2L^{3}+5L^{2}+2L-3,
\end{align*}
so that
\begin{equation}\label{plt:eq:div-recip-example-C}
  \recip{B}=1-(L+1)t+(2L+1)t^{2}+(L^{3}-L^{2}-4L+1)t^{3}
  +(-L^{4}-2L^{3}+5L^{2}+2L-3)t^{4}+\FormalBigOAt{t^{5}}{t}.
\end{equation}
Multiplying \cref{plt:eq:div-recip-example-B} by
\cref{plt:eq:div-recip-example-C} and truncating at \(t^{5}\) returns \(1\);
by \cref{plt:thm:div-truncated} this residual check certifies the five blocks
independently of the method used to produce them.
\end{example}

\begin{example}[A quotient with polynomial blocks]
\label{plt:ex:div-quotient-worked}
With
\[
  A=1+(L+1)t+(2L^{2}-L)t^{2}+\FormalBigOAt{t^{3}}{t},
  \qquad
  B=1+(2-L)t+(L^{2}+3)t^{2}+\FormalBigOAt{t^{3}}{t},
\]
\cref{plt:eq:div-quotient-recurrence} gives
\(Q_{0}=1\), \(Q_{1}=(L+1)-(2-L)=2L-1\), and
\[
  Q_{2}=(2L^{2}-L)-(2-L)(2L-1)-(L^{2}+3)=3L^{2}-6L-1,
\]
so \(A/B=1+(2L-1)t+(3L^{2}-6L-1)t^{2}+\FormalBigOAt{t^{3}}{t}\).  The block
\(Q_{3}\) is \emph{not} determined by the displayed data:
\cref{plt:prop:div-relative-precision} with \(a=b=0\) says three input blocks
determine exactly three output blocks.  If one additionally posits
\(A_{3}=B_{3}=0\), then \(Q_{3}=L^{3}-11L^{2}+5L+5\).
\end{example}

% ed.: S1 prints the block Q_3 of this example immediately after inputs given
% only modulo t^3; the hypothesis that the omitted blocks vanish is a genuine
% extra assumption and is stated here.
\begin{warningbox}[title={Do not read a block that the data does not contain}]
An input written with a formal tail \(\FormalBigOAt{t^{3}}{t}\) determines the
quotient only through \(t^{2}\).  Continuing the recurrence past that point
silently substitutes the assumption ``the omitted blocks are zero''.  Both S1
and S6 print one extra block of a worked example on exactly this unstated
convention; when it is intended it must be said.
\end{warningbox}

\begin{proposition}[Degree profile of a reciprocal and of a quotient]
\label{plt:prop:div-degree-profile}
Let \(\alpha\ge0\) and \(\beta\ge0\) be real.  Suppose \(B\in\PLpow\) has
\(B_{0}\in\K^{\times}\) and \(\degL B_{n}\le\alpha n\) for all \(n\ge1\).  Then
the blocks of \(\recip{B}\) satisfy
\begin{equation}\label{plt:eq:div-degree-recip}
  \degL C_{n}\le\alpha n\qquad(n\ge0).
\end{equation}
If moreover \(A\in\PLpow\) has \(\degL A_{n}\le\alpha n+\beta\), then the
quotient blocks satisfy \(\degL Q_{n}\le\alpha n+\beta\).
\end{proposition}

\begin{proof}
Strong induction on \(n\).  For \(n=0\), \(C_{0}=B_{0}^{-1}\in\K\) has degree
\(0\le0\).  For \(n\ge1\), each summand of \cref{plt:eq:div-recip-cn} obeys
\(\degL(B_{j}C_{n-j})\le\alpha j+\alpha(n-j)=\alpha n\) by the inductive
hypothesis and the additivity of \(\degL\) on products in \(\K[L]\); a finite
sum cannot raise the bound, and multiplying by the scalar \(-B_{0}^{-1}\)
changes nothing.  The quotient case is identical, with the extra term \(A_{n}\)
of degree \(\le\alpha n+\beta\) and \(\degL(B_{j}Q_{n-j})\le\alpha
j+\alpha(n-j)+\beta\).
\end{proof}

\begin{remark}
The bound in \cref{plt:eq:div-degree-recip} may be strict: cancellation can
lower a block's degree, and it does so in
\cref{plt:ex:div-reciprocal-worked}, where \(\degL B_{2}=2\) and \(\degL
C_{2}=1\).  An affine profile \(\dprofile{F}(n)\le\alpha n+\beta\) is a
hypothesis or a proved invariant, never part of the definition of the ring; the
Lambert expansion is the archetype with \(\alpha=1\), \(\beta=0\).
\end{remark}

\section{Geometric inversion and two-level summability}
\label{plt:sec:div-geometric}

The geometric formula \cref{plt:eq:div-geometric} is the conceptual
construction, and in \(\PLlau\) its justification is purely \(t\)-adic:
\(\ordt(U^{k})\ge k\to\infty\).  Over a completed logarithmic coefficient field
this reason is no longer available, and a second level of the argument is
needed.  Several sources normalise a divisor in \(\PLit\) as
\(B=b\,t^{\nu}L^{\mu}(1+S)\) with \(S\precdom1\) and then assert summability of
\(\sum_{k}(-S)^{k}\) ``by the support conditions''.  That is true but not for
the \(t\)-adic reason: \(S\) may have \(\ordt S=0\), so \(\ordt(S^{k})=0\) for
every \(k\).  The correct statement is the following.

\begin{lemma}[Two-level summability]\label{plt:lem:div-two-level}
For \(s\in\K((L^{-1}))\) let \(\omega(s)\) denote the \(L^{-1}\)-adic order, so
that \(s\in L^{-\omega(s)}\K[[L^{-1}]]\) and \(\omega(0)=+\infty\); \(\omega\)
is additive on products (\cref{plt:lem:div-valuation} with \(L^{-1}\) in place
of \(t\) and \(E=\K\)) and satisfies
\(\omega(s+s')\ge\min\SetOf{\omega(s),\omega(s')}\).  Let
\(S=\sum_{n\ge0}s_{n}t^{n}\in\PLit\) satisfy
\begin{equation}\label{plt:eq:div-small-S}
  \ordt S\ge0
  \qquad\text{and}\qquad
  \omega(s_{0})\ge1 ,
\end{equation}
which is exactly the condition \(S\precdom1\) in the dominance order.  Fix
\(n\ge0\) and put
\[
  c_{n}\DefinitionEquals
  \max\SetOf{0,\,-\omega(s_{1}),\,\dots,\,-\omega(s_{n})}
  \in\NaturalNumbers .
\]
Then for every \(k\ge0\),
\begin{equation}\label{plt:eq:div-two-level-order}
  \omega\bigl([t^{n}]S^{k}\bigr)\ge k-n\,(1+c_{n}).
\end{equation}
Consequently, for each fixed monomial \(t^{n}L^{-j}\) only the finitely many
indices \(k\le j+n(1+c_{n})\) contribute to \(\sum_{k\ge0}(-S)^{k}\); the
family is summable in \(\PLit\), and \((1+S)\sum_{k\ge0}(-S)^{k}=1\).
\end{lemma}

\begin{proof}
Because \(\ordt S\ge0\),
\[
  [t^{n}]S^{k}=\sum_{\substack{(n_{1},\dots,n_{k})\in\NaturalNumbers^{k}\\
                               n_{1}+\cdots+n_{k}=n}}
    s_{n_{1}}\cdots s_{n_{k}},
\]
a finite sum in which every index satisfies \(n_{i}\le n\).  Fix one tuple and
let \(r\) be the number of indices with \(n_{i}\ge1\); then \(r\le n\), since
each such index contributes at least \(1\) to the total \(n\).  The \(k-r\)
factors with \(n_{i}=0\) have \(\omega(s_{0})\ge1\), and the \(r\) remaining
factors have \(\omega(s_{n_{i}})\ge-c_{n}\) by definition of \(c_{n}\).
Additivity of \(\omega\) on products gives
\[
  \omega\bigl(s_{n_{1}}\cdots s_{n_{k}}\bigr)
  \ge(k-r)\cdot1+r\cdot(-c_{n})
  \ge k-n-n c_{n},
\]
using \(r\le n\) and \(c_{n}\ge0\).  Taking the minimum over the finitely many
tuples yields \cref{plt:eq:div-two-level-order}.

A power \(L^{-j}\) can occur in an element of order \(\omega\) only when
\(j\ge\omega\); hence \([t^{n}L^{-j}]S^{k}=0\) as soon as
\(k-n(1+c_{n})>j\), which proves the finiteness assertion.

% ed.: pointwise finiteness alone does not place the sum in the field: one must
% also check that each block is a Laurent series in L^{-1}, i.e. that its
% support is bounded below.  The same bound gives this at once, and no source
% records the step.
Summability in \(\PLit\) needs one more line.  The outer support of every
\(S^{k}\) lies in \(n\ge0\) because \(\ordt S\ge0\), so the outer support of
the sum is bounded below.  For a fixed \(n\),
\cref{plt:eq:div-two-level-order} with \(k\ge0\) gives
\(\omega\bigl([t^{n}]S^{k}\bigr)\ge-n(1+c_{n})\) for \emph{every} \(k\), so
each contribution to the \(n\)-th block is supported in \(j\ge-n(1+c_{n})\);
that block is therefore a well-defined element of \(\K((L^{-1}))\), and the
sum lies in \(\PLit\).

For the identity, fix \((n,j)\) and choose \(K\) larger than
\(j+n(1+c_{n})\).  Telescoping gives
\((1+S)\sum_{k\le K}(-S)^{k}=1-(-S)^{K+1}\), and by the finiteness just proved
the coefficient of \(t^{n}L^{-j}\) in \((-S)^{K+1}\) vanishes and the
coefficient of \(t^{n}L^{-j}\) in \((1+S)\sum_{k\le K}(-S)^{k}\) already equals
that in \((1+S)\sum_{k\ge0}(-S)^{k}\).  Hence the two sides agree at every
monomial.
\end{proof}

\begin{remark}
\Cref{plt:lem:div-two-level} is why the catalogue insists that a finite
representation in \(\PLit\) needs an outer \(t\)-cutoff \emph{and} an inner
\(L^{-1}\)-cutoff, or one explicitly declared staircase.  A single outer cutoff
does not bound the work: the coefficient of \(t^{0}L^{-j}\) still receives
contributions from \(j+1\) powers of \(S\).
\end{remark}

\section{Why unrestricted division needs the rational-log field}
\label{plt:sec:div-rational-field}

\begin{theorem}[Laurent series over a field]\label{plt:thm:div-Ett-field}
Let \(E\) be a field.  Then \(E((t))\) is a field.  In particular
\(\PLrat=\K(L)((t))\) and \(\PLit=\K((L^{-1}))((t))\) are fields.
\end{theorem}

\begin{proof}
\(E((t))\) is a commutative ring with \(1\), and it is an integral domain by
\cref{plt:lem:div-valuation}.  Let \(F\ne0\) with \(m=\ordt F\) and leading
coefficient \(f_{m}\in E^{\times}\).  Put \(U\DefinitionEquals
f_{m}^{-1}t^{-m}F-1\in tE[[t]]\).  Exactly as in the proof of
\cref{plt:thm:div-unit-criterion}, \(\ordt(U^{k})\ge k\) makes
\(\sum_{k\ge0}(-U)^{k}\) formally summable in \(E[[t]]\), and telescoping gives
\((1+U)\sum_{k\ge0}(-U)^{k}=1\).  Hence
\(F\cdot f_{m}^{-1}t^{-m}\sum_{k\ge0}(-U)^{k}=1\).  Finally \(\K(L)\) and
\(\K((L^{-1}))\) are fields: the first is the fraction field of the integral
domain \(\K[L]\); the second is the Laurent series field over \(\K\) in the
variable \(L^{-1}\), a field by the case already proved.
\end{proof}

The next theorem says that \(\PLrat\) is not merely \emph{a} field containing
\(\PLlau\), but the smallest one of the shape ``Laurent series in \(t\) over a
coefficient ring between \(\K[L]\) and \(\K(L)\)''.  This is the precise sense
in which unrestricted division of \emph{coefficient blocks} forces \(\K(L)\).

\begin{theorem}[Minimality of the rational-log coefficient field]
\label{plt:thm:div-minimal-field}
Let \(E\) be a subring of \(\K(L)\) with \(\K[L]\subseteq E\).  Then \(E((t))\)
is a field if and only if \(E=\K(L)\).  Consequently \(\PLrat\) is the unique
smallest field of the form \(E((t))\) containing \(\PLlau\) with
\(E\subseteq\K(L)\).
\end{theorem}

\begin{proof}
If \(E=\K(L)\), apply \cref{plt:thm:div-Ett-field}.  Conversely suppose
\(E((t))\) is a field and let \(e\in E\), \(e\ne0\).  Its inverse
\(G=\sum_{n}G_{n}t^{n}\in E((t))\) satisfies \(eG=1\), i.e.\ \(eG_{0}=1\) and
\(eG_{n}=0\) for \(n\ne0\).  Since \(E\subseteq\K(L)\) is an integral domain
and \(e\ne0\), we get \(G_{n}=0\) for \(n\ne0\) and \(G_{0}=e^{-1}\in E\).
Hence \(E\) is a field.  A subfield of \(\K(L)\) containing \(\K[L]\) contains
\(p^{-1}\) for every nonzero \(p\in\K[L]\) and hence every quotient \(p/q\),
so \(E=\K(L)\).
\end{proof}

\section{An integral domain that is not a field, made precise}
\label{plt:sec:div-domain}

``\(\PLlau\) is an integral domain but not a field'' is the standard slogan.
It is worth saying exactly how far from a field it is, and -- a point none of
the six sources records -- that its field of fractions is \emph{strictly
smaller} than \(\PLrat\).  Passing to \(\PLrat\) is therefore not the same
operation as ``forming the fraction field''.

\begin{lemma}[Denominator lemma]\label{plt:lem:div-denominator}
Let \(A,B\in\PLlau\) with \(B\ne0\), and let \(b_{0}\in\K[L]\) be the leading
block of \(B\).  Then, computed in \(\PLrat\),
\begin{equation}\label{plt:eq:div-denominator}
  \frac AB\in\K[L][b_{0}^{-1}]\,((t)),
\end{equation}
where \(\K[L][b_{0}^{-1}]=\SetOf{p/b_{0}^{k}:p\in\K[L],\,k\ge0}\subseteq\K(L)\).
In particular every element of the fraction field of \(\PLlau\) has all of its
coefficient blocks in a \emph{single} localisation \(\K[L][q^{-1}]\) of
\(\K[L]\) at one nonzero polynomial \(q\).
\end{lemma}

\begin{proof}
Write \(m=\ordt B\) and \(S\DefinitionEquals\K[L][b_{0}^{-1}]\), a subring of
\(\K(L)\).  Then \(t^{-m}B=\sum_{k\ge0}b_{k}t^{k}\) with \(b_{k}\in\K[L]\), so
\(b_{0}^{-1}t^{-m}B=1+U\) with \(U\DefinitionEquals\sum_{k\ge1}(b_{k}/b_{0})
t^{k}\in tS[[t]]\).  By the proof of \cref{plt:thm:div-Ett-field} applied to
the ring \(S\) -- only summability and telescoping are used, not that \(S\) is
a field -- \((1+U)^{-1}=\sum_{k\ge0}(-U)^{k}\in S[[t]]\).  Hence
\(B^{-1}=t^{-m}b_{0}^{-1}(1+U)^{-1}\in S((t))\) and, since
\(A\in\PLlau\subseteq S((t))\), also \(A/B\in S((t))\).  The last sentence is
this statement with \(q=b_{0}\).
\end{proof}

\begin{theorem}[The fraction field of \(\PLlau\) is strictly smaller than
\(\PLrat\)]\label{plt:thm:div-frac-strict}
Let \(\K\) have characteristic zero.  Then
\begin{equation}\label{plt:eq:div-frac-strict}
  \PLlau\subsetneq\operatorname{Frac}(\PLlau)\subsetneq\PLrat .
\end{equation}
An explicit witness for the second strict inclusion is
\begin{equation}\label{plt:eq:div-frac-witness}
  F\DefinitionEquals\sum_{n\ge0}\frac{t^{n}}{L-n}\in\PLrat .
\end{equation}
\end{theorem}

\begin{proof}
The first inclusion is strict because \(L\) is a nonzero nonunit
(\cref{plt:ex:div-nonunit}).  For the second, \(\operatorname{Frac}(\PLlau)\)
embeds in the field \(\PLrat\) because \(\PLrat\) contains \(\PLlau\), so we
must only show \(F\notin\operatorname{Frac}(\PLlau)\).  Certainly
\(F\in\PLrat\): its blocks \((L-n)^{-1}\) lie in \(\K(L)\) and its support is
bounded below by \(0\).

Suppose \(F=A/B\) with \(A,B\in\PLlau\), \(B\ne0\).  By
\cref{plt:lem:div-denominator} there is a nonzero \(q\in\K[L]\) with every
block of \(F\) in \(\K[L][q^{-1}]\).  Fix \(n\ge0\).  Then
\((L-n)^{-1}=p/q^{k}\) for some \(p\in\K[L]\) and \(k\ge0\), i.e.\
\(q^{k}=(L-n)p\).  Since \(\K\) has characteristic zero, the elements
\(0,1,2,\dots\) of \(\K\) are pairwise distinct, so \(L-n\) is a nonassociate
irreducible of the unique factorisation domain \(\K[L]\) for each \(n\), and
\((L-n)\mid q^{k}\) forces \((L-n)\mid q\).  As this holds for every
\(n\ge0\), the polynomial \(q\) has infinitely many roots, hence \(q=0\), a
contradiction.
\end{proof}

\begin{corollary}[The precise structure]\label{plt:cor:div-domain-not-field}
\(\PLlau\) is an integral domain whose unit group is
\cref{plt:eq:div-unit-group}.  It is not a field, and its nonunits do not form
an ideal, so it is not a local ring: \(L\) and \(1-L\) are nonunits whose sum
is the unit \(1\).  The obstruction to division is therefore not concentrated
at one maximal ideal but is spread over the primes of \(\K[L]\): for each monic
irreducible \(\pi\in\K[L]\), the set \(\pi\PLlau\) is a proper nonzero ideal,
and \(\bigcap_{\pi}\pi\PLlau=0\).
\end{corollary}

\begin{proof}
The first two sentences have been proved.  For the last, \(\pi\PLlau\) is an
ideal; it is nonzero and proper because \(\pi\) is a nonzero nonunit
(\cref{plt:thm:div-unit-criterion}).  If \(F\ne0\) lay in \(\pi\PLlau\) for
every \(\pi\), its leading block \(p_{\ordt F}\) would be divisible in
\(\K[L]\) by every monic irreducible -- by \cref{plt:lem:div-valuation} the
leading block of \(\pi G\) is \(\pi\) times the leading block of \(G\) -- and
\(\K[L]\) has infinitely many pairwise nonassociate irreducibles (Euclid's
argument applies verbatim), forcing \(p_{\ordt F}=0\).
\end{proof}

\section{Conditional divisibility}
\label{plt:sec:div-conditional}

A nonunit divisor still divides some elements exactly.  The following
proposition makes the test explicit; it also shows that ``being divisible'' is
a genuinely global condition, so that no unconditional division theorem in
\(\PLlau\) is available.

\begin{proposition}[Coefficientwise divisibility test]
\label{plt:prop:div-conditional}
Let \(B\in\PLlau\) be nonzero with \(\ordt B=0\) and leading block \(B_{0}\in
\K[L]\), and let \(A\in\PLpow\).  Let \(Q=A/B\in\K(L)[[t]]\) have blocks
\(Q_{n}\) given by \cref{plt:eq:div-quotient-recurrence} over \(E=\K(L)\).
Then:
\begin{enumerate}
\item \(A/B\in\PLpow\) if and only if \(Q_{n}\in\K[L]\) for every \(n\ge0\);
\item given \(Q_{0},\dots,Q_{n-1}\in\K[L]\), the block \(Q_{n}\) lies in
      \(\K[L]\) if and only if \(B_{0}\) divides
      \(A_{n}-\sum_{j=1}^{n}B_{j}Q_{n-j}\) in \(\K[L]\);
\item the set \(\SetOf{A\in\PLlau:A/B\in\PLlau}\) is exactly the principal
      ideal \(B\PLlau\), which is nonzero, and proper as soon as \(B\) is not a
      unit.
\end{enumerate}
\end{proposition}

\begin{proof}
(1) is the definition of \(\PLpow\) as the set of series with polynomial
blocks, together with the uniqueness in \cref{plt:cor:div-quotient}: the blocks
of \(A/B\) computed over \(\K(L)\) are the only candidates.  (2) is immediate
from \cref{plt:eq:div-quotient-recurrence}, since \(Q_{n}=B_{0}^{-1}R_{n}\)
with \(R_{n}=A_{n}-\sum_{j=1}^{n}B_{j}Q_{n-j}\in\K[L]\) under the stated
hypothesis.  For (3): \(A/B\in\PLlau\) means \(A=B\cdot(A/B)\) with
\(A/B\in\PLlau\), i.e.\ \(A\in B\PLlau\); conversely \(A=BC\) with
\(C\in\PLlau\) gives \(A/B=C\).  The ideal \(B\PLlau\) contains \(B\ne0\), and
it equals \(\PLlau\) exactly when \(1\in B\PLlau\), i.e.\ when \(B\) is a unit.
\end{proof}

\begin{example}[Exact division by a nonunit, and a near miss]
\label{plt:ex:div-conditional}
Take \(B=L+t\), a nonzero nonunit.  Then
\begin{equation}\label{plt:eq:div-exact-nonunit}
  \frac{L^{2}-t^{2}}{L+t}=L-t\in\PLpow ,
\end{equation}
verified by the single multiplication \((L+t)(L-t)=L^{2}-t^{2}\).  So a nonunit
can divide exactly.  Change one sign and the divisibility fails at the third
step: with \(A=L^{2}+t^{2}\), the test of
\cref{plt:prop:div-conditional}\,(2) gives
\[
  Q_{0}=\frac{L^{2}}{L}=L,\qquad
  Q_{1}=\frac{0-1\cdot L}{L}=-1,\qquad
  Q_{2}=\frac{1-1\cdot(-1)}{L}=\frac2L\notin\K[L],
\]
consistently with the exact identity
\(\dfrac{L^{2}+t^{2}}{L+t}=L-t+\dfrac{2t^{2}}{L+t}\).  Two successful steps
prove nothing about the third.
\end{example}

\begin{remark}[The test is a semi-decision procedure]
\label{plt:rmk:div-semidecision}
\Cref{plt:prop:div-conditional}\,(2) refutes divisibility in finite time
whenever it fails at some finite order.  It never confirms divisibility in
finite time on its own, because the condition quantifies over all \(n\).
A positive answer needs a finite certificate: either a candidate quotient
\(C\in\PLlau\) with only finitely many blocks -- so that \(BC=A\) is a finite
identity, as in \cref{plt:eq:div-exact-nonunit} -- or an a priori bound
reducing the infinitely many conditions to finitely many, or a structural
argument.  Absent such a certificate, a computation that ``keeps coming out
polynomial'' is evidence, not proof.
\end{remark}

\begin{checkpointbox}
Three closure questions must be kept apart.
\begin{enumerate}
\item Does the quotient exist in \emph{some} field?  Always, in \(\PLrat\).
\item Does it have rational blocks and nothing worse?  Always, again in
      \(\PLrat\); this is \cref{plt:thm:div-Ett-field}.
\item Does it have polynomial blocks?  Only under the unit criterion, or by a
      certified accident as in \cref{plt:ex:div-conditional}.
\end{enumerate}
Answering (1) and reporting it as an answer to (3) is the standard failure
mode.
\end{checkpointbox}

\section{Newton iteration and the exact precision-doubling law}
\label{plt:sec:div-newton}

The reciprocal recurrence produces one block per step.  Newton's method for
the equation \(C^{-1}-B=0\) produces the map
\begin{equation}\label{plt:eq:div-newton-step}
  C\longmapsto C(2-BC),
\end{equation}
which doubles the number of correct blocks.  Every source states
\cref{plt:eq:div-newton-step} and the identity
\(1-B\,C(2-BC)=(1-BC)^{2}\); none
states which truncations of \(B\) and of \(C\) suffice to compute the doubled
output, which is precisely what makes the scheme an algorithm rather than an
identity.  Part (4) below supplies that statement.

\begin{theorem}[Precision doubling]\label{plt:thm:div-newton}
Let \(B\in\PLpow\) with \(B_{0}\in\K^{\times}\), let \(C\in\PLpow\), let
\(N\ge1\), and suppose
\begin{equation}\label{plt:eq:div-newton-invariant}
  \ordt(1-BC)\ge N .
\end{equation}
Put \(C'\DefinitionEquals C(2-BC)\).  Then:
\begin{enumerate}
\item \(1-BC'=(1-BC)^{2}\), hence \(\ordt(1-BC')\ge2N\);
\item \(\ordt\bigl(\recip{B}-C'\bigr)=\ordt(1-BC')\ge 2N\), and more
      precisely \(\ordt(\recip{B}-C)=\ordt(1-BC)\);
\item \(\ordt\bigl(1-B\Trunc{C'}{2N}\bigr)\ge2N\), so truncating the output at
      order \(2N\) preserves the invariant;
\item \(\Trunc{C'}{2N}\) depends on \(B\) and \(C\) only through
      \(\Trunc{B}{2N}\) and \(\Trunc{C}{N}\).  Explicitly, if
      \(\overline B\in\PLpow\) and \(\overline C\in\PLpow\) satisfy
      \(\ordt(B-\overline B)\ge2N\) and \(\ordt(C-\overline C)\ge N\), then
      \(\ordt\bigl(C'-\overline C(2-\overline B\,\overline C)\bigr)\ge2N\).
\end{enumerate}
\end{theorem}

\begin{proof}
(1)  \(BC'=BC(2-BC)=2BC-(BC)^{2}\), so
\(1-BC'=1-2BC+(BC)^{2}=(1-BC)^{2}\).  Valuations add
(\cref{plt:lem:div-valuation}), so \(\ordt(1-BC')\ge2N\).

(2)  \(B\) is a unit by \cref{plt:thm:div-unit-criterion} and
\(\ordt(\recip{B})=-\ordt B=0\).  From
\(\recip{B}-C'=\recip{B}\,(1-BC')\) and additivity of \(\ordt\) we get
\(\ordt(\recip{B}-C')=\ordt(1-BC')\).  The same computation with \(C\) in
place of \(C'\) gives the second identity.

(3)  \(C'-\Trunc{C'}{2N}\in t^{2N}\PLpow\), so
\(B(C'-\Trunc{C'}{2N})\in t^{2N}\PLpow\) because \(\ordt B=0\); adding this to
(1) gives \(\ordt(1-B\Trunc{C'}{2N})\ge2N\).

% ed.: the input-requirement statement (4) is absent from S2, S3, S5 and S6,
% each of which asserts doubling without saying which truncations suffice.
(4)  Put \(\delta=C-\overline C\) and \(\varepsilon=B-\overline B\), so
\(\ordt\delta\ge N\) and \(\ordt\varepsilon\ge2N\).  Since
\(C'=2C-BC^{2}\) and \(\overline C(2-\overline B\,\overline C)
=2\overline C-\overline B\,\overline C^{2}\),
\[
  C'-\overline C(2-\overline B\,\overline C)
  =2\delta-\bigl(BC^{2}-\overline B\,\overline C^{2}\bigr)
  =2\delta-\varepsilon C^{2}
    -\overline B\,(C+\overline C)\,\delta
  =\delta\bigl(2-\overline B(C+\overline C)\bigr)-\varepsilon C^{2},
\]
using \(BC^{2}-\overline B\,\overline C^{2}
=\varepsilon C^{2}+\overline B(C^{2}-\overline C^{2})\) and
\(C^{2}-\overline C^{2}=(C+\overline C)\delta\).  The second term has
\(\ordt\ge2N\).  For the first,
\[
  2-\overline B(C+\overline C)
  =\bigl(1-\overline BC\bigr)+\bigl(1-\overline B\,\overline C\bigr),
\]
and \(1-\overline BC=(1-BC)+\varepsilon C\) has \(\ordt\ge N\), while
\(1-\overline B\,\overline C=(1-BC)+\varepsilon C+B\delta-\varepsilon\delta\)
also has \(\ordt\ge N\).  Hence
\(\ordt\bigl(\delta(2-\overline B(C+\overline C))\bigr)\ge N+N=2N\), and the
whole difference has \(\ordt\ge2N\).
\end{proof}

\begin{algorithm}[Newton reciprocal to exclusive order \(N\)]
\label{plt:alg:div-newton}
% ed.: the precondition N >= 1 was not printed; with N = 0 the loop body never
% runs and the routine returns the scalar B_0^{-1} rather than
% Trunc_{<0}(1/B) = 0.
Given an integer \(N\ge1\) and \(B\in\PLpow\) with \(B_{0}\in\K^{\times}\),
known through order \(t^{N-1}\):
\begin{lstlisting}[style=pseudo]
function newton_reciprocal(B, N):
    require B[0] is a nonzero scalar in K
    C := 1/B[0]                    # exact through order t^0
    m := 1
    while m < N:
        m2 := min(2*m, N)
        E  := Trunc(1 - B*C, m2)   # valuation >= m by the invariant
        C  := Trunc(C*(1 + E), m2)
        m  := m2
    return C                       # equals Trunc_{<N} of the reciprocal
\end{lstlisting}
\end{algorithm}

\begin{corollary}[The iterates are exact truncations]
\label{plt:cor:div-newton-exact}
Let
\[
  C^{(0)}=B_{0}^{-1},
  \qquad
  C^{(k+1)}=\Trunc{\bigl(C^{(k)}(2-BC^{(k)})\bigr)}{2^{k+1}} .
\]
Then for every \(k\ge0\),
\begin{equation}\label{plt:eq:div-newton-exact}
  C^{(k)}=\Trunc{\recip{B}}{2^{k}} .
\end{equation}
In particular \cref{plt:alg:div-newton} returns \(\Trunc{\recip{B}}{N}\), and
the intermediate steps require no blocks of \(B\) beyond index \(2^{k+1}-1\)
at step \(k\).
\end{corollary}

\begin{proof}
Induction on \(k\).  For \(k=0\), \(1-BC^{(0)}=1-B_{0}^{-1}B\) has zero
constant block, so \(\ordt(1-BC^{(0)})\ge1\), and by
\cref{plt:thm:div-newton}\,(2) \(\ordt(\recip{B}-C^{(0)})\ge1\); as
\(C^{(0)}\) has \(t\)-degree \(0<1\), it equals \(\Trunc{\recip{B}}{1}\).
Assume \(C^{(k)}=\Trunc{\recip{B}}{2^{k}}\); then
\(\ordt(\recip{B}-C^{(k)})\ge2^{k}\), so \(\ordt(1-BC^{(k)})\ge2^{k}\) by
\cref{plt:thm:div-newton}\,(2).  Parts (1) and (3) give
\(\ordt(1-BC^{(k+1)})\ge2^{k+1}\), hence
\(\ordt(\recip{B}-C^{(k+1)})\ge2^{k+1}\) by (2) again; since \(C^{(k+1)}\) has
\(t\)-degree \(<2^{k+1}\), it is \(\Trunc{\recip{B}}{2^{k+1}}\).

% ed.: the induction covers only the exact powers of two, while
% \cref{plt:alg:div-newton} caps every step at N; the capped step needs its own
% line before the "in particular" clause is justified.
\emph{The capped step.}  In \cref{plt:alg:div-newton} the new precision is
\(m'=\min\SetOf{2m,N}\le2m\).  Suppose \(C=\Trunc{\recip{B}}{m}\) at the top of
the loop, so \(\ordt(\recip{B}-C)\ge m\) and hence \(\ordt(1-BC)\ge m\) by
\cref{plt:thm:div-newton}\,(2).  Part (1) gives \(\ordt(1-BC')\ge2m\ge m'\) for
\(C'=C(2-BC)\), and the computation in the proof of part (3), run with \(m'\)
in place of \(2N\), gives \(\ordt\bigl(1-B\Trunc{C'}{m'}\bigr)\ge m'\); by (2)
once more \(\ordt\bigl(\recip{B}-\Trunc{C'}{m'}\bigr)\ge m'\), and
\(\Trunc{C'}{m'}\) has \(t\)-degree \(<m'\), so it equals
\(\Trunc{\recip{B}}{m'}\).  The invariant holds initially with \(m=1\) and the
loop terminates with \(m=N\), so the returned value is
\(\Trunc{\recip{B}}{N}\) for every \(N\ge1\), not only for \(N\) a power of
two.  The statement about required blocks of \(B\) is
\cref{plt:thm:div-newton}\,(4).
\end{proof}

% ed.: S2 justifies the doubling by "the valuation is multiplicative"; the
% valuation is additive on products, which is what the argument uses.
\begin{remark}[A wording to avoid]\label{plt:rmk:div-valuation-additive}
Writing \(E=1-BC\), the step \(\ordt(E^{2})\ge2N\) holds because \(\ordt\) is
\emph{additive} on products (\cref{plt:lem:div-valuation}), not because it is
multiplicative;
\(\ordt\) is a homomorphism from the multiplicative monoid of nonzero elements
to \((\IntegerNumbers,+)\).
\end{remark}

\begin{remark}[Where the iteration lives]
\Cref{plt:thm:div-newton} used only that \(\ordt\) is additive and that
\(B\) has an invertible leading block; it therefore holds verbatim over the
coefficient rings \(\K(L)\) and \(\K((L^{-1}))\), and in \(\PLit\) after
normalising the leading Laurent-log coefficient.  It is an exact algebraic
identity, not an approximation: unlike numerical Newton iteration, there is no
convergence hypothesis and no loss of significance -- by
\cref{plt:cor:div-newton-exact} each iterate is literally a truncation of the
answer, so the logarithmic degrees appearing in the iterates are exactly those
of \(\recip{B}\) and cannot swell beyond them, provided one truncates at every
step.
\end{remark}

\section{Cost: the recurrence against the doubling scheme}
\label{plt:sec:div-cost}

% ed.: the convention "multiplications by the scalar B_0^{-1} are not counted"
% was stated loosely enough to exclude also the products B_j C_0 and C_0 E_j,
% under which reading the two schemes cost exactly the same and the N-1 gap
% claimed below disappears.  The counted set is pinned down here and the other
% reading is recorded at the end of \cref{plt:prop:div-cost-newton}.
We count multiplications in the coefficient ring \(\K[L]\): every product of
two coefficient blocks costs one \emph{block multiplication}, including those
products one of whose factors happens to be the scalar block
\(C_{0}=B_{0}^{-1}\).  Additions are free, and so is the single outer
normalisation by \(B_{0}^{-1}\) that
\cref{plt:eq:div-recip-cn,plt:eq:div-quotient-recurrence} apply once per step.
Let \(M(N)\) denote the number of block multiplications used by whatever
routine computes \(\Trunc{(PQ)}{N}\) from \(\Trunc{P}{N}\) and
\(\Trunc{Q}{N}\).

\begin{proposition}[Cost of the recurrence]\label{plt:prop:div-cost-naive}
Computing \(\Trunc{\recip{B}}{N}\) by
\cref{plt:eq:div-recip-c0,plt:eq:div-recip-cn} uses exactly
\begin{equation}\label{plt:eq:div-cost-naive}
  \binom{N}{2}=\frac{N(N-1)}{2}
\end{equation}
block multiplications, and computing \(\Trunc{(A/B)}{N}\) by
\cref{plt:eq:div-quotient-recurrence} uses the same number.
\end{proposition}

\begin{proof}
Step \(n\) forms the products \(B_{j}C_{n-j}\) for \(j=1,\dots,n\), that is
\(n\) block multiplications, for \(n=0,1,\dots,N-1\).  The total is
\(\sum_{n=0}^{N-1}n=N(N-1)/2\).
\end{proof}

\begin{proposition}[Cost of the doubling scheme]
\label{plt:prop:div-cost-newton}
Let \(N=2^{K}\).
\begin{enumerate}
\item If \(M\) is nondecreasing and satisfies \(2M(n)\le M(2n)\), then
      \cref{plt:alg:div-newton} uses at most \(4M(N)\) block multiplications.
\item With schoolbook truncated multiplication, and exploiting the invariant
      \(\ordt(1-BC)\ge m\) so that the first \(m\) blocks of \(BC\) are known
      in advance, \cref{plt:alg:div-newton} uses exactly
      \begin{equation}\label{plt:eq:div-cost-newton-exact}
        \frac{N^{2}+N-2}{2}
      \end{equation}
      block multiplications, which exceeds
      \cref{plt:eq:div-cost-naive} by exactly \(N-1\).
\end{enumerate}
\end{proposition}

\begin{proof}
(1)  The step from precision \(m\) to \(2m\) forms two truncated products at
order \(2m\), so costs at most \(2M(2m)\).  Summing over
\(m=1,2,4,\dots,2^{K-1}\) gives at most \(\sum_{j=1}^{K}2M(2^{j})\).  The
hypothesis \(2M(n)\le M(2n)\) gives \(M(2^{j})\le M(2^{K})/2^{K-j}\), so the
sum is at most \(2M(N)\sum_{j=1}^{K}2^{\,j-K}\le4M(N)\).

(2)  Consider the step from \(m\) to \(2m\), with \(C=\Trunc{\recip{B}}{m}\)
supported on indices \(0,\dots,m-1\).
\emph{Forming \(E=1-BC\) modulo \(t^{2m}\).}  By the invariant, the blocks of
\(BC\) of index \(<m\) are known to be \(1,0,\dots,0\); only indices
\(m,\dots,2m-1\) must be computed.  The needed products are \(B_{j}C_{i}\) with
\(0\le i<m\) and \(m\le i+j<2m\), and for each of the \(m\) values of \(i\)
there are exactly \(m\) admissible \(j\).  Cost: \(m^{2}\).
\emph{Forming \(C(1+E)\) modulo \(t^{2m}\).}  \(E\) is supported on indices
\(m,\dots,2m-1\), and \(C\) on \(0,\dots,m-1\); the products \(C_{i}E_{j}\)
with \(i+j<2m\) number \(\sum_{j=m}^{2m-1}(2m-j)=\sum_{s=1}^{m}s=m(m+1)/2\).
Total for the step: \(m^{2}+m(m+1)/2=(3m^{2}+m)/2\).  Summing over
\(m=2^{k}\), \(k=0,\dots,K-1\),
\[
  \sum_{k=0}^{K-1}\frac{3\cdot4^{k}+2^{k}}{2}
  =\frac{3}{2}\cdot\frac{4^{K}-1}{3}+\frac{2^{K}-1}{2}
  =\frac{N^{2}-1}{2}+\frac{N-1}{2}
  =\frac{N^{2}+N-2}{2}.
\]
Subtracting \cref{plt:eq:div-cost-naive} gives \((2N-2)/2=N-1\).

% ed.: the N-1 gap is fragile: it is created entirely by charging for products
% against the scalar block C_0.  Recording this keeps the comparison from being
% read as a robust separation, which it is not.
\emph{How fragile the gap is.}  Suppose instead that a product one of whose
factors is the scalar \(C_{0}=B_{0}^{-1}\) is charged nothing.  That removes
one product per step \(n=1,\dots,N-1\) from the recurrence, hence \(N-1\) in
all, and \(2m\) from the doubling step at precision \(m\) (the \(m\) products
\(B_{j}C_{0}\) and the \(m\) products \(C_{0}E_{j}\)), hence
\(\sum_{k=0}^{K-1}2\cdot2^{k}=2N-2\) in all.  Both counts then drop to
\(\binom{N-1}{2}\) and the two schemes cost \emph{exactly the same}.  The
\(N-1\) gap of part (2) is therefore an artefact of the accounting; what is
robust in the comparison is part (1).
\end{proof}

% ed.: S3 asserts that Newton inversion "is often faster than coefficient
% recursion for large N even with naive multiplication"; the exact counts above
% show the opposite at the level of coefficient multiplications, and the honest
% statement is the corollary.
% ed.: the corollary was headed by the iff "improves ... exactly when block
% multiplication is subquadratic".  Only the two implications below are proved
% (subquadratic wins asymptotically; schoolbook loses by N-1), and neither the
% converse nor the intermediate regime M(N) = cN^2 is covered; the claim is
% downgraded to what the two cost counts give.
\begin{corollary}[When doubling pays]\label{plt:cor:div-doubling-pays}
Which of the two schemes wins is decided by the growth of \(M\), and the two
counts settle the two extremes.  Precisely: if \(M(N)/N^{2}\to0\) as
\(N\to\infty\) (with \(M\) as in \cref{plt:prop:div-cost-newton}\,(1)), then
the doubling scheme uses \(o(N^{2})\) block multiplications against
\(\Theta(N^{2})\) for the recurrence; whereas with schoolbook block
multiplication it uses exactly \(N-1\) \emph{more} block multiplications than
the recurrence in the accounting fixed at the head of
\cref{plt:sec:div-cost}, and exactly as many under the alternative accounting
described at the end of the proof of \cref{plt:prop:div-cost-newton}.  Under
neither accounting is it faster.
\end{corollary}

\begin{proof}
Combine \cref{plt:prop:div-cost-naive,plt:prop:div-cost-newton}; the first
claim is part (1) with the stated growth hypothesis, the second is part (2).
\end{proof}

% ed.: S6 gives the cost as O(N^2 d^2) "if the typical coefficient degree is
% d"; under the affine log-degree profile that actually occurs there is no such
% constant d, and the honest count is quartic in N.
\begin{remark}[Scalar cost under an affine degree profile]
\label{plt:rmk:div-scalar-cost}
Block multiplications are not unit-cost.  Assume the hypothesis of
\cref{plt:prop:div-degree-profile} with \(\alpha>0\) an integer, so
\(\degL B_{n}\le\alpha n\) and \(\degL C_{n}\le\alpha n\), and count
\(\K\)-multiplications with schoolbook polynomial arithmetic in \(L\).  The
recurrence then costs at most
\[
  \sum_{n=1}^{N-1}\sum_{j=1}^{n}\bigl(\alpha j+1\bigr)
       \bigl(\alpha(n-j)+1\bigr)
  =\Theta\bigl(\alpha^{2}N^{4}\bigr),
  \qquad\text{with leading term }\frac{\alpha^{2}N^{4}}{24},
\]
because \(\sum_{j=1}^{n}j(n-j)=\binom{n+1}{3}\sim n^{3}/6\) and
\(\sum_{n<N}\binom{n+1}{3}=\binom{N+1}{4}\sim N^{4}/24\).  Quoting a bound
\(\BigO(N^{2}d^{2})\)
with a fixed \(d\) therefore understates the cost whenever the log-degree
profile grows; substituting the honest \(d\asymp\alpha N\) recovers
\(\alpha^{2}N^{4}\).  This quartic behaviour, not the quadratic block count, is
what makes truncating after \emph{every} multiplication the decisive
implementation discipline; see \cite{vanderhoeven-book} for the relaxed and
fast-multiplication variants.
\end{remark}

\section{The embedding \texorpdfstring{\(\PLrat\hookrightarrow\PLit\)}%
{K(L)((t)) into K((1/L))((t))}}
\label{plt:sec:div-embedding}

The third response to a nonconstant leading block is to expand each rational
coefficient at \(L=\infty\).  This resolves a finer asymptotic scale, and it is
the natural home of a multiseries implementation \cite{salvy-shackell-1998};
it is also the step most likely to be taken without noticing that the object
has changed.

\begin{definition}[Iterated Laurent-log field]\label{plt:def:div-iterated}
\(\K((L^{-1}))\) is the field of formal Laurent series in the variable
\(L^{-1}\): elements \(\sum_{j\ge j_{0}}a_{j}L^{-j}\) with
\(j_{0}\in\IntegerNumbers\) and \(a_{j}\in\K\), so that only finitely many
positive powers of \(L\) occur but arbitrarily many negative ones may.  Then
\(\PLit\DefinitionEquals\K((L^{-1}))((t))\), whose elements are
\(\sum_{n\ge n_{0}}t^{n}\sum_{j\ge j_{n}}a_{n,j}L^{-j}\).
\end{definition}

\begin{theorem}[The expansion map]\label{plt:thm:div-embedding}
There is exactly one \(\K\)-algebra homomorphism
\(\iota_{0}:\K(L)\to\K((L^{-1}))\) with \(\iota_{0}(L)=L\), namely Laurent
expansion at \(L=\infty\), and it is injective.  Applying it blockwise defines
an injective ring homomorphism \(\iota:\PLrat\to\PLit\) that restricts to the
identity on \(\PLlau\).  The map \(\iota\) is \emph{not} surjective; its image
is therefore a proper subfield of \(\PLit\).
\end{theorem}

\begin{proof}
\emph{Existence on coefficients.}  A polynomial \(\sum_{i\le d}c_{i}L^{i}\)
already lies in \(\K((L^{-1}))\) (as \(\sum_{j\ge-d}c_{-j}L^{-j}\)), giving an
injective ring map \(\K[L]\hookrightarrow\K((L^{-1}))\).  Since
\(\K((L^{-1}))\) is a field (\cref{plt:thm:div-Ett-field}), every nonzero
polynomial maps to a unit, so by the universal property of the localisation
\(\K(L)=\operatorname{Frac}\K[L]\) the map extends uniquely to a ring
homomorphism \(\iota_{0}:\K(L)\to\K((L^{-1}))\); a ring homomorphism out of a
field is injective.  Concretely, for \(q=q_{d}L^{d}(1+w)\) with
\(q_{d}\in\K^{\times}\) and \(w\in L^{-1}\K[L^{-1}]\),
\(\iota_{0}(1/q)=q_{d}^{-1}L^{-d}\sum_{k\ge0}(-w)^{k}\), the series being
summable because \(w\) has positive \(L^{-1}\)-order.

Uniqueness of \(\iota_{0}\) is the uniqueness clause in that universal
property: a \(\K\)-algebra map on \(\K(L)\) is determined by the image of
\(L\).

\emph{Extension to Laurent series.}  Define
\(\iota\bigl(\sum_{n\ge n_{0}}f_{n}t^{n}\bigr)
=\sum_{n\ge n_{0}}\iota_{0}(f_{n})t^{n}\).  The support stays bounded below, so
the image is in \(\PLit\).  Applying a ring homomorphism blockwise to Laurent
series is again a ring homomorphism, because each block of a sum or a Cauchy
product is a finite algebraic expression in the blocks of the arguments; and
\(\iota\) is injective because \(\iota_{0}\) is.  On \(\PLlau\) every block is
a polynomial and \(\iota_{0}\) is the inclusion, so \(\iota\) is the identity
there.

\emph{Not surjective.}  Since \(\operatorname{char}\K=0\), the binomial series
\(\sigma\DefinitionEquals\sum_{k\ge0}\binom{1/2}{k}L^{-k}\in\K[[L^{-1}]]\) is
defined and satisfies \(\sigma^{2}=1+L^{-1}\), the identity holding
coefficientwise because it holds in \(\RationalNumbers[[u]]\) for
\(u=L^{-1}\).  Suppose \(\sigma=\iota_{0}(f)\) for some \(f\in\K(L)\).  Then
\(f^{2}=(L+1)/L\), so \(\bigl(Lf\bigr)^{2}=L(L+1)\) with \(Lf\in\K(L)\).
Writing \(Lf=u/v\) in lowest terms gives \(u^{2}=L(L+1)v^{2}\); comparing
multiplicities of the irreducible \(L\) on the two sides of this equation in
the unique factorisation domain \(\K[L]\) yields an even number on the left and
an odd one on the right, a contradiction.  Hence \(\sigma\) is not in the image
of \(\iota_{0}\), and \(\sigma\) regarded as a \(t^{0}\)-block is not in the
image of \(\iota\).
\end{proof}

\begin{example}[The exact coefficient becomes an infinite series]
\label{plt:ex:div-embedding}
\begin{equation}\label{plt:eq:div-embedding-example}
  \iota\!\left(\frac1{L+1}\right)
  =L^{-1}-L^{-2}+L^{-3}-\cdots
  =\sum_{j\ge1}(-1)^{j-1}L^{-j},
  \qquad
  \iota\!\left(\frac1{L+t}\right)
  =\sum_{n\ge0}(-1)^{n}L^{-n-1}t^{n}.
\end{equation}
The left-hand sides are single finite data; the right-hand sides are infinite.
\end{example}

\begin{proposition}[The iteration order is part of the type]
\label{plt:prop:div-iteration-order}
Regard both \(\PLit=\K((L^{-1}))((t))\) and \(\K((t))((L^{-1}))\) as sets of
families \((a_{n,j})\) of scalars indexed by
\((n,j)\in\IntegerNumbers^{2}\), representing \(\sum a_{n,j}t^{n}L^{-j}\).
For the first, the support condition is: the set of \(n\) carrying a nonzero
entry is bounded below, and for each such \(n\) the set of \(j\) with
\(a_{n,j}\ne0\) is bounded below.  For the second, the roles of \(n\) and
\(j\) are exchanged.  Then neither set contains the other.  Witnesses:
\begin{equation}\label{plt:eq:div-iteration-order}
  \sum_{n\ge0}L^{n}t^{n}\in\K((L^{-1}))((t))\setminus\K((t))((L^{-1})),
  \qquad
  \sum_{j\ge0}L^{-j}t^{-j}\in\K((t))((L^{-1}))\setminus\K((L^{-1}))((t)).
\end{equation}
\end{proposition}

\begin{proof}
The first element has \(a_{n,-n}=1\) for \(n\ge0\) and all other entries zero.
Its outer index \(n\) is bounded below by \(0\), and for each \(n\) exactly one
\(j\) occurs, so the first support condition holds; indeed the element lies
already in \(\PLpow\).  But for the second condition one must read the family
by inner index: the set of \(j\) carrying a nonzero entry is
\(\SetOf{0,-1,-2,\dots}\), which is not bounded below, so the element is not in
\(\K((t))((L^{-1}))\).

The second element has \(a_{-j,j}=1\) for \(j\ge0\) and all other entries zero.
Its inner index \(j\) is bounded below by \(0\), and for each \(j\) exactly one
\(n\) occurs, so it lies in \(\K((t))((L^{-1}))\).  But the set of \(n\)
carrying a nonzero entry is \(\SetOf{0,-1,-2,\dots}\), so it has infinitely
many positive powers of \(X\) and is excluded from \(\K((L^{-1}))((t))\).
\end{proof}

% ed.: the catalogue asserts that the two iteration orders differ "in general";
% no source proves it, and the witnesses above supply the proof.
\begin{warningbox}[title={Expanding at \(L=\infty\) changes the object}]
\label{plt:rmk:div-embedding-warning}
Three distinct things go wrong if \(\iota\) is applied silently.
\begin{enumerate}
\item \emph{Finiteness.}  An exact block such as \(1/(L+1)\) becomes an
      infinite series.  A finite representation in \(\PLit\) therefore needs an
      outer \(t\)-cutoff \emph{and} an inner \(L^{-1}\)-cutoff, or one declared
      staircase in the lexicographic monomial order.  An outer cutoff alone
      does not make the data finite.
\item \emph{Algebra.}  \(\iota\) is injective but not surjective
      (\cref{plt:thm:div-embedding}), so \(\PLrat\) sits inside \(\PLit\) as a
      proper subfield: an identity established in \(\PLit\) need not descend to
      \(\PLrat\), and an element of \(\PLit\) need not be a rational function
      of \(L\) at all.  Moreover the leading data change meaning -- in
      \(\PLlau\) each block has a finite \(\degL\), while in \(\PLit\) the
      relevant invariant is the \(L^{-1}\)-adic order.
\item \emph{Asymptotics.}  Let \(\K\subseteq\ComplexNumbers\) and let
      \(f\in\K(L)\) have largest pole modulus \(\rho(f)\), with
      \(\rho(f)=0\) when \(f\) is a polynomial and \(\iota_{0}(f)\) a finite
      sum.  The numerical series \(\iota_{0}(f)\) converges for
      \(\AbsoluteValue{L}>\rho(f)\), that is for \(X>\EulerE^{\rho(f)}\) on the
      positive ray, and for nonpolynomial \(f\) diverges for
      \(\AbsoluteValue{L}<\rho(f)\).  There is in general \emph{no} threshold
      valid for all blocks at once: for
      \(F=\sum_{n\ge0}t^{n}/(L-n)\in\PLrat\) the \(n\)-th block requires
      \(X>\EulerE^{n}\).  Hence ``\(F\) equals its \(L^{-1}\)-expansion for
      large \(X\)'' is not a theorem about \(F\); it is a statement about one
      block at a time, and any uniform claim needs its own proof with the
      regime printed as \(\BigOAt{\cdot}{X\to+\infty}\).
\end{enumerate}
\end{warningbox}

\begin{remark}[When to expand and when not to]
Keeping \(1/(L+1)\) unexpanded is exact, finite, and preserves cancellation;
expanding it resolves the finer scale
\(t\precdom\cdots\precdom L^{-2}\precdom L^{-1}\precdom1\) and is required as
soon as one must compare two blocks of the same \(t\)-order.  The decision is a
modelling decision and should be recorded, not inferred: a system that expands
rational coefficients automatically will destroy exact cancellations and swell
expressions, while one that never expands cannot answer questions about the
inner scale.  The same tension appears in every multiseries implementation
\cite{salvy-shackell-1998,salvy-shackell-1999}.
\end{remark}

\begin{summarybox}
An element of \(\PLlau\) is invertible there exactly when its leading
coefficient block is a nonzero scalar; the leading \emph{block} is a
polynomial, and factoring out the dominant monomial does not by itself invert.
Reciprocals of units exist, are unique, and are computed by a triangular
recurrence whose truncated output \(P=\Trunc{(A/B)}{N}\) is certified by the
residual \(BP=A+\FormalBigOAt{t^{N}}{t}\) alone; relative precision is exactly
preserved, so \(M\) leading input blocks give \(M\) leading output blocks.
Newton's map \(C\mapsto C(2-BC)\) doubles the correct order and, with the input
requirement of \cref{plt:thm:div-newton}\,(4), yields exact truncations of the
% ed.: "beats the recurrence only when block multiplication is subquadratic"
% asserts the unproved converse; only the two extreme regimes are established.
reciprocal; it beats the recurrence asymptotically when block multiplication is
subquadratic, and does not beat it at all when multiplication is schoolbook.
Unrestricted division of coefficient blocks forces \(\PLrat\),
the smallest Laurent field over an intermediate coefficient ring, which
strictly contains the fraction field of \(\PLlau\).  Expanding each rational
coefficient at \(L=\infty\) embeds \(\PLrat\) into \(\PLit\) injectively but
not surjectively, and converts one exact datum into an infinite series with no
uniform analytic threshold.
\end{summarybox}

% ---- fragment 05-differential ----
\chapter{Powers, logarithms, exponentials, and the closure boundary}
\label{plt:sec:elementary}

Addition, multiplication and division were settled by finite coefficient
algebra.  The operations of this chapter -- powers with an arbitrary exponent,
logarithms, exponentials, and more generally the substitution of a series into
a scalar power series -- are the first for which the answer to ``does the
result stay in the class?'' is neither always yes nor always no.  The purpose
of the section is to replace the informal closure tables of the sources by
statements with hypotheses and proofs, and in particular to prove the two
negative results that fix the boundary of the polynomial--logarithmic scale:
a logarithm leaves the scale exactly when the leading coefficient block is
nonconstant, and an exponential leaves \emph{every} power--log scale exactly
when its argument is not of nonnegative outer order with an affine
\(t^{0}\)-block.

Throughout, \(\K\) is a field of characteristic zero, \(t\DefinitionEquals
X^{-1}\) and \(L\DefinitionEquals\log X\) are independent formal generators,
the default analytic regime is \(X\to+\infty\) on the positive ray, and a
nonzero \(F\in\PLlau=\K[L]((t))\) is written
\begin{equation}\label{plt:eq:dif-normal-form}
  F=\sum_{n\ge n_{0}}p_{n}(L)t^{n},
  \qquad p_{n}\in\K[L],\quad p_{n_{0}}\ne0,\quad n_{0}=\ordt F .
\end{equation}
We write \(p_{\ordt F}\) for the leading block, \(\degL p_{\ordt F}\) for its
degree, and \(t\PLpow\) for the elements of \(\PLpow=\K[L][[t]]\) of strictly
positive outer order.  Truncation is always the exclusive operator
\(\Trunc{F}{N}=\sum_{n<N}p_{n}(L)t^{n}\), and \(F=G+\FormalBigOAt{t^{N}}{t}\)
abbreviates \(\ordt(F-G)\ge N\), a statement about coefficients only.

Two facts proved elsewhere in the volume are used freely: the unit criterion
(\cref{plt:thm:div-unit-criterion}), namely that \(F\) is a unit of \(\PLlau\)
if and only if its leading block is a nonzero scalar; and the expansion
embedding \(\iota\colon\PLrat\hookrightarrow\PLit\)
(\cref{plt:thm:div-embedding}), Laurent expansion of each coefficient at
\(L=\infty\).  On \(\K((L^{-1}))\) we use the \(L^{-1}\)-adic order
\(\omega\) of \cref{plt:lem:div-two-level}: \(\omega(\sum_{j\ge j_{0}}a_{j}
L^{-j})=\min\SetOf{j:a_{j}\ne0}\), so that \(\omega(p)=-\degL p\) for a nonzero
polynomial \(p\), \(\omega\) is additive on products, and
\(\omega(h+h')\ge\min\SetOf{\omega(h),\omega(h')}\).

The distinguished derivation \(\DX=t\partial_{L}-t^{2}\partial_{t}\) appears in
the hypotheses of several statements below; it is constructed, and all of its
properties used here are proved, in \cref{plt:sec:differential}, which is
independent of the present chapter.  The only facts borrowed are the block
formula \(\DX(t^{n}p)=t^{n+1}(p'-np)\), the resulting bound
\(\ordt(\DX F)\ge\ordt F+1\), and \(\ker\DX=\K\).

\section{Substituting a small series into a scalar power series}
\label{plt:sec:dif-substitution}

Every construction in this section rests on one mechanism: a series of
positive outer order may be substituted into an arbitrary formal power series
over \(\K\), because the outer order of the \(r\)-th power grows with \(r\).
All six sources state this; we record it once, in the form that also delivers
the finite-determination statement that a computation actually needs.

% ed.: retitled.  Three ``Summable families'' statements stood in the volume
% ed.: with the same bracketed title.  This one is the blockwise assembly
% ed.: principle for a sequence in a Laurent ring over any of the coefficient
% ed.: rings used here; \cref{plt:lem:cmp-summable} is the corresponding
% ed.: statement for an arbitrary index set in a Hahn ring with real
% ed.: exponents, and \cref{plt:thm:trunc-summability} is the full calculus on
% ed.: \(\PLlau\).  On \(\PLlau\) the present lemma is the sequence case of
% ed.: \cref{plt:thm:trunc-summability}\,(i),(iv),(v); it is stated again
% ed.: because it is applied below in the \(L^{-1}\)-adic order and in a Hahn
% ed.: ambient, neither of which that theorem covers.  The hypothesis
% ed.: \(\inf_{r}\ordt F_{r}>-\infty\) is redundant for a sequence with
% ed.: \(\ordt F_{r}\to+\infty\) and is kept only because the proof uses it.
\begin{lemma}[Summable sequences and blockwise assembly]
\label{plt:lem:dif-summable}
Let \((F_{r})_{r\ge0}\) be a family in \(\PLlau\) with
\(\ordt F_{r}\to+\infty\) and \(\inf_{r}\ordt F_{r}>-\infty\).  Then for each
\(n\in\IntegerNumbers\) only finitely many \(r\) have \([t^{n}]F_{r}\ne0\), the
element
\[
  \sum_{r\ge0}F_{r}\DefinitionEquals
  \sum_{n\in\IntegerNumbers}\Bigl(\sum_{r\ge0}[t^{n}]F_{r}\Bigr)t^{n}
\]
is a well-defined member of \(\PLlau\), and the assignment is
\(\K\)-linear and compatible with multiplication by a fixed element:
\(G\sum_{r}F_{r}=\sum_{r}GF_{r}\).  The same statements hold in \(\PLpow\),
in \(\PLrat\) and in \(\PLit\), with the coefficient ring changed accordingly.
\end{lemma}

\begin{proof}
Fix \(n\).  Since \(\ordt F_{r}\to+\infty\), the set \(R_{n}=\SetOf{r:\ordt
F_{r}\le n}\) is finite, and \([t^{n}]F_{r}=0\) for \(r\notin R_{n}\); so the
inner sum is a finite sum in the coefficient ring.  The common lower bound
\(m=\inf_{r}\ordt F_{r}\) is finite and gives \([t^{n}]F_{r}=0\) for all
\(n<m\), so the support of the result is bounded below and the result lies in
\(\PLlau\).  Linearity is immediate.  For the last claim,
\([t^{n}](G\sum_{r}F_{r})=\sum_{i+j=n}[t^{i}]G\sum_{r}[t^{j}]F_{r}\), a finite
double sum that may be reordered.  No step used that the coefficients are
polynomials.
\end{proof}

% ed.: retitled.  This is substitution of a small series into a \emph{scalar}
% ed.: power series, \(\K[[z]]\to\PLpow\).  The composition part has a
% ed.: different theorem under the same name,
% ed.: \cref{plt:thm:cmp-substitution-hom}, which substitutes a
% ed.: monomial-leading transseries into a transseries; neither implies the
% ed.: other.
\begin{proposition}[Substitution of a small series into a scalar power series]
\label{plt:prop:dif-substitution}
Let \(U\in\PLpow\) with \(q\DefinitionEquals\ordt U\ge1\).  For
\(\Phi=\sum_{r\ge0}c_{r}z^{r}\in\K[[z]]\) put
\begin{equation}\label{plt:eq:dif-substitution}
  \Phi(U)\DefinitionEquals\sum_{r\ge0}c_{r}U^{r}.
\end{equation}
Then the sum is defined by \cref{plt:lem:dif-summable}, \(\Phi(U)\in\PLpow\)
with \(\ordt\Phi(U)\ge q\cdot\operatorname{ord}_{z}\Phi\), and
\[
  \Phi\longmapsto\Phi(U)
\]
is a \(\K\)-algebra homomorphism \(\K[[z]]\to\PLpow\) sending \(z\) to \(U\).
It is the unique such homomorphism that is continuous for the \(z\)-adic and
\(t\)-adic topologies.  The same holds with \(\PLpow\) replaced by
\(\K(L)[[t]]\) or \(\K((L^{-1}))[[t]]\).
\end{proposition}

\begin{proof}
\(\ordt(U^{r})\ge rq\to+\infty\) and all orders are \(\ge0\), so
\cref{plt:lem:dif-summable} applies and gives \(\Phi(U)\in\PLpow\); if
\(c_{r}=0\) for \(r<r_{0}\) then every surviving term has order \(\ge r_{0}q\).
Additivity is clear.  For multiplicativity, let \(\Psi=\sum_{s}d_{s}z^{s}\).
Then
\[
  \Phi(U)\Psi(U)
  =\sum_{r}c_{r}U^{r}\Psi(U)
  =\sum_{r}\sum_{s}c_{r}d_{s}U^{r+s},
\]
where the first step is the last clause of \cref{plt:lem:dif-summable} and the
second applies it again inside each term; the resulting double family is
summable because \(\ordt(U^{r+s})\ge(r+s)q\), so it may be summed by the
antidiagonals \(r+s=m\), giving \(\sum_{m}(\sum_{r+s=m}c_{r}d_{s})U^{m}
=(\Phi\Psi)(U)\).  Uniqueness: a continuous homomorphism is determined on the
dense subring \(\K[z]\), where it is determined by the image of \(z\).
\end{proof}

The next statement is what a computation uses.  The sources give only its
first half -- that few powers contribute -- and leave implicit which blocks of
\(U\) must be known; both halves are needed to certify a truncated result.

% ed.: retitled: the composition part states a different theorem under the
% ed.: same name (\cref{plt:thm:cmp-finite-determination}, finite
% ed.: determination of \(F\compose G\)).  The present one is about
% ed.: \(\Phi(U)\) with \(\Phi\) a scalar power series.
\begin{theorem}[Finite determination of a substituted scalar series]
\label{plt:thm:dif-finite-determination}
Let \(U\in\PLpow\) with \(\ordt U=q\ge1\) and \(\Phi=\sum_{r}c_{r}z^{r}\in
\K[[z]]\).  For \(N\le0\), \(\Trunc{\Phi(U)}{N}=0\).  For \(N\ge1\), put
\(m=\Floor{(N-1)/q}\).  Then
\begin{equation}\label{plt:eq:dif-finite-determination}
  \Trunc{\Phi(U)}{N}
  =\Trunc{\;\sum_{r=0}^{m}c_{r}\bigl(\Trunc{U}{N}\bigr)^{r}}{N}.
\end{equation}
In particular \(\Trunc{\Phi(U)}{N}\) depends only on the \(m+1\) scalars
\(c_{0},\dots,c_{m}\) and on the blocks \([t^{n}]U\) with \(n<N\), and
\([t^{N}]\Phi(U)\) receives contributions only from the powers \(U^{r}\) with
\(r\le\Floor{N/q}\).
\end{theorem}

\begin{proof}
Both sides lie in \(\PLpow\), so the case \(N\le0\) is trivial.  Let \(N\ge1\).
If \(r>m\) then \(r\ge m+1>(N-1)/q\), hence \(rq>N-1\) and \(rq\ge N\); since
\(\ordt(c_{r}U^{r})\ge rq\), such terms do not affect
\(\Trunc{\cdot}{N}\), which proves that only \(r\le m\) matter.  For the
replacement of \(U\) by \(\Trunc{U}{N}\), write \(E=U-\Trunc{U}{N}\), so
\(\ordt E\ge N\).  For \(r\ge1\),
\[
  U^{r}-\bigl(\Trunc{U}{N}\bigr)^{r}
  =E\sum_{i=0}^{r-1}U^{i}\bigl(\Trunc{U}{N}\bigr)^{r-1-i}
\]
has outer order \(\ge N\), because every factor in the sum has nonnegative
order.  Hence the two truncations agree.  The last clause is the same
computation with \(N\) replaced by \(N+1\): \(rq>N\) forces
\([t^{N}](c_{r}U^{r})=0\).
\end{proof}

\begin{corollary}[The elementary series]\label{plt:cor:dif-elementary}
Let \(U\in t\PLpow\) and \(\alpha\in\K\).  Then
\begin{align}
  (1+U)^{\alpha}&\DefinitionEquals\sum_{r\ge0}\binom{\alpha}{r}U^{r}
   \in1+t\PLpow,
   \label{plt:eq:dif-binomial}\\
  \log(1+U)&\DefinitionEquals\sum_{r\ge1}\frac{(-1)^{r+1}}{r}U^{r}
   \in t\PLpow,
   \label{plt:eq:dif-log1pU}\\
  \exp U&\DefinitionEquals\sum_{r\ge0}\frac{U^{r}}{r!}\in1+t\PLpow
   \label{plt:eq:dif-expU}
\end{align}
are well defined, and every identity between formal power series over \(\K\)
in one or two variables that involves only these operations transfers.  In
particular, for \(U,V\in t\PLpow\),
\begin{gather}
  (1+U)^{\alpha}(1+U)^{\beta}=(1+U)^{\alpha+\beta},\qquad
  \bigl((1+U)^{\alpha}\bigr)^{n}=(1+U)^{n\alpha}\ (n\in\IntegerNumbers),
  \label{plt:eq:dif-power-laws}\\
  \exp(U+V)=(\exp U)(\exp V),\qquad
  \exp\bigl(\log(1+U)\bigr)=1+U,\qquad
  \log(\exp U)=U,
  \label{plt:eq:dif-exp-log-laws}\\
  \log\bigl((1+U)(1+V)\bigr)=\log(1+U)+\log(1+V),\qquad
  (1+U)^{\alpha}=\exp\bigl(\alpha\log(1+U)\bigr).
  \label{plt:eq:dif-log-laws}
\end{gather}
\end{corollary}

\begin{proof}
Each of \cref{plt:eq:dif-binomial,plt:eq:dif-log1pU,plt:eq:dif-expU} is
\(\Phi(U)\) for a \(\Phi\in\K[[z]]\) (here \(\operatorname{char}\K=0\) is used
to divide by \(r\) and \(r!\)), so \cref{plt:prop:dif-substitution} applies;
the order statements come from \(\operatorname{ord}_{z}\Phi\).  For the
identities, note that a one-variable identity \(\Phi_{1}=\Phi_{2}\) in
\(\K[[z]]\) is carried to \(\Phi_{1}(U)=\Phi_{2}(U)\) by the homomorphism.  For
the two-variable identities, the same argument applies to the unique continuous
\(\K\)-algebra homomorphism \(\K[[z_{1},z_{2}]]\to\PLpow\) with
\(z_{i}\mapsto U,V\), which exists by the proof of
\cref{plt:prop:dif-substitution} verbatim; the required identities
\(\exp(z_{1}+z_{2})=\exp z_{1}\exp z_{2}\) and so on are classical in
\(\RationalNumbers[[z_{1},z_{2}]]\).  The two identities in
\cref{plt:eq:dif-exp-log-laws} that involve a composition use one further
step: if \(\Psi\in\K[[z]]\) has \(\Psi(0)=0\), then
\((\Phi\circ\Psi)(U)=\Phi(\Psi(U))\) for every \(\Phi\).  Indeed both
\(\Phi\mapsto(\Phi\circ\Psi)(U)\) and \(\Phi\mapsto\Phi(\Psi(U))\) are
continuous \(\K\)-algebra homomorphisms \(\K[[z]]\to\PLpow\) -- the first as
the composite of \(\Phi\mapsto\Phi\circ\Psi\), which is
\cref{plt:prop:dif-substitution} applied inside \(\K[[z]]\), with evaluation at
\(U\); the second directly, since \(\ordt\Psi(U)\ge1\) -- and both send
\(z\mapsto\Psi(U)\), so they agree by the uniqueness clause.  Applying this to
the classical identities \(\exp(\log(1+z))=1+z\) and \(\log(\exp z)=z\) in
\(\RationalNumbers[[z]]\) gives the two displayed identities.
\end{proof}

For explicit coefficients there are two standard devices, a recurrence and a
closed form; each is preferable in a different situation.

\begin{proposition}[Coefficients of a substituted series]
\label{plt:prop:dif-bell}
Let \(U=\sum_{n\ge1}u_{n}t^{n}\in t\PLpow\) and
\(\Phi=\sum_{r\ge0}c_{r}z^{r}\in\K[[z]]\).  Then \([t^{0}]\Phi(U)=c_{0}\) and,
for \(N\ge1\),
\begin{equation}\label{plt:eq:dif-bell}
  [t^{N}]\Phi(U)
  =\sum_{k=1}^{N}c_{k}\,\frac{k!}{N!}\,
   \ExponentialPartialBellPolynomial{N}{k}
   \bigl(1!\,u_{1},2!\,u_{2},\dots,(N-k+1)!\,u_{N-k+1}\bigr).
\end{equation}
In particular, with \(E=\exp U\) and \(M=\log(1+U)\),
\begin{align}
  [t^{N}]E&=\frac{1}{N!}\,
    \ExponentialCompleteBellPolynomial{N}\bigl(1!\,u_{1},\dots,N!\,u_{N}\bigr),
    \label{plt:eq:dif-exp-bell}\\
  [t^{N}]M&=\frac{1}{N!}\sum_{k=1}^{N}(-1)^{k-1}(k-1)!\,
    \ExponentialPartialBellPolynomial{N}{k}
    \bigl(1!\,u_{1},\dots,(N-k+1)!\,u_{N-k+1}\bigr).
    \label{plt:eq:dif-log-bell}
\end{align}
Moreover \(E\) and \(M\) satisfy the recurrences
\begin{equation}\label{plt:eq:dif-exp-recurrence}
  N\,[t^{N}]E=\sum_{m=1}^{N}m\,u_{m}\,[t^{N-m}]E,
  \qquad
  N\,[t^{N}]M=N\,u_{N}-\sum_{m=1}^{N-1}m\,[t^{m}]M\;u_{N-m},
\end{equation}
which compute all blocks through \(t^{N}\) with \(\BigO(N^{2})\) coefficient
multiplications.
\end{proposition}

\begin{proof}
The partial exponential Bell polynomials are defined \cite{comtet-book} by the
identity
\[
  \frac{1}{k!}\Bigl(\sum_{m\ge1}a_{m}\frac{z^{m}}{m!}\Bigr)^{k}
  =\sum_{n\ge k}\ExponentialPartialBellPolynomial{n}{k}(a_{1},a_{2},\dots)
   \frac{z^{n}}{n!}
\]
in \(\RationalNumbers[a_{1},a_{2},\dots][[z]]\); since it is an identity with
universal coefficients it holds after any specialisation into a commutative
\(\RationalNumbers\)-algebra.  Specialise \(a_{m}=m!\,u_{m}\in\K[L]\) and
\(z=t\); then \(\sum_{m}a_{m}z^{m}/m!=U\), and comparing blocks gives
\([t^{n}]U^{k}=(k!/n!)\ExponentialPartialBellPolynomial{n}{k}(1!u_{1},2!u_{2},
\dots)\), the arguments beyond index \(n-k+1\) being irrelevant.  Summing over
\(k\) and using \([t^{N}]U^{k}=0\) for \(k>N\) gives \cref{plt:eq:dif-bell}.
\Cref{plt:eq:dif-exp-bell,plt:eq:dif-log-bell} are the cases
\(c_{k}=1/k!\) and \(c_{k}=(-1)^{k-1}/k\), together with
\(\ExponentialCompleteBellPolynomial{N}=\sum_{k=1}^{N}
\ExponentialPartialBellPolynomial{N}{k}\) for \(N\ge1\).

For the recurrences, let \(\partial_{t}\) be the formal partial derivative
holding \(L\) fixed.  It is a derivation of \(\PLpow\) with
\(\ordt(\partial_{t}F)\ge\ordt F-1\), hence commutes with summable sums, so
applying it to \cref{plt:eq:dif-expU} termwise gives
\(\partial_{t}E=(\partial_{t}U)E\).  Comparing the coefficient of \(t^{N-1}\)
and using \([t^{N-1}]\partial_{t}E=N[t^{N}]E\) and
\([t^{m-1}]\partial_{t}U=mu_{m}\) yields the first recurrence.  Likewise
\(\partial_{t}M=(\partial_{t}U)(1+U)^{-1}\), that is
\((1+U)\partial_{t}M=\partial_{t}U\); comparing coefficients of \(t^{N-1}\)
gives \(N[t^{N}]M+\sum_{m=1}^{N-1}m[t^{m}]M\,u_{N-m}=Nu_{N}\).  Each recurrence
performs \(\BigO(N)\) coefficient multiplications per block.
\end{proof}

\begin{example}[A finite logarithm]\label{plt:ex:dif-log-finite}
Let \(U=(L+1)t+L^{2}t^{2}\), so \(q=1\).  By
\cref{plt:thm:dif-finite-determination} with \(N=4\), only \(r\le3\)
contribute below \(t^{4}\), and
\begin{align*}
  \log(1+U)
  &=U-\tfrac12U^{2}+\tfrac13U^{3}+\FormalBigOAt{t^{4}}{t}\\
  &=(L+1)t
   +\Bigl(L^{2}-\tfrac{(L+1)^{2}}{2}\Bigr)t^{2}
   +\Bigl(\tfrac{(L+1)^{3}}{3}-L^{2}(L+1)\Bigr)t^{3}
   +\FormalBigOAt{t^{4}}{t}.
\end{align*}
Every step is a finite polynomial computation, and the certificate
\(\exp(\log(1+U))-(1+U)=\FormalBigOAt{t^{4}}{t}\) checks it.
% ed.: S3 states the analogous expansion with U^4 = O(t^4) as the reason for
% stopping at r = 3; the correct reason is that r*q >= N kills the term, which
% is the general statement proved above and matters as soon as q > 1.
\end{example}

\section{The ambient exponential extension}
\label{plt:sec:dif-ambient}

The exponential of an element of \(\PLlau\) need not lie in \(\PLlau\), so a
statement such as ``\(\exp F\) leaves the class'' is meaningless until an
ambient object has been named in which \(\exp F\) exists.  We therefore fix one
once and for all.  Nothing below depends on which ambient is chosen: every
negative result is proved from the axioms alone, so it holds in all of them
simultaneously.

\begin{definition}[Exponential extension]\label{plt:def:dif-exp-extension}
An \emph{exponential extension} of \(\PLlau\) is a triple
\((\mathcal T,\partial,\exp)\) in which
\begin{enumerate}
\item \(\mathcal T\) is an integral domain containing \(\PLlau\) as a subring,
      and \(\partial\) is a derivation of \(\mathcal T\) restricting to
      \(\DX\) on \(\PLlau\);
\item \(\exp\colon(\mathcal T,+)\to(\mathcal T^{\times},\cdot)\) is a group
      homomorphism with \(\partial(\exp F)=(\partial F)\exp F\) for all
      \(F\in\mathcal T\);
\item (normalisation) \(\exp L=t^{-1}\) and \(\exp\bigl(t\PLpow\bigr)
      \subseteq1+t\PLpow\).
\end{enumerate}
We write \(\operatorname{Const}(\mathcal T)=\ker\partial\) for the constants,
say that \(G\in\mathcal T\) is \emph{a logarithm of} \(F\) when \(\exp G=F\),
and for \(\alpha\in\K\) and \(F\) admitting a logarithm \(G\) we set
\(F^{\alpha}\DefinitionEquals\exp(\alpha G)\), a value that depends on the
choice of \(G\).
\end{definition}

\begin{remark}[Realisations]\label{plt:rmk:dif-realisations}
For \(\K=\RealNumbers\) the field \(\mathbb{T}\) of logarithmic--exponential
transseries is such an extension, with \(\partial\) the transseries derivation
and \(\exp\) its exponential
\cite{adh-book,vanderhoeven-book,edgar-beginners}; so is any Hardy field
containing the germs at \(+\infty\) of \(X\mapsto X\) and \(X\mapsto\log X\)
and closed under \(\exp\), for instance the germs of the
\(\log\)--\(\exp\)-definable functions \cite{vdDMM-2001,hardy-orders}.
Condition (iii) is a normalisation, not a restriction: it merely says that the
ambient exponential extends the relation \(L=\log X\) and does not multiply
small exponentials by a stray constant.  \Cref{plt:lem:dif-series-exp} shows
that (iii) already forces \(\exp\) to agree on \(t\PLpow\) with the series of
\cref{plt:eq:dif-expU}.
\end{remark}

\begin{lemma}[The ambient exponential is the series exponential]
\label{plt:lem:dif-series-exp}
Let \((\mathcal T,\partial,\exp)\) be an exponential extension.  For every
\(U\in t\PLpow\),
\[
  \exp U=\sum_{r\ge0}\frac{U^{r}}{r!},
\]
the right-hand side being \cref{plt:eq:dif-expU}.  Consequently, for
\(V\in1+t\PLpow\) the element \(\log V\) of \cref{plt:eq:dif-log1pU} is a
logarithm of \(V\) in \(\mathcal T\).
\end{lemma}

\begin{proof}
Write \(S=\sum_{r\ge0}U^{r}/r!\in1+t\PLpow\), a unit of \(\PLpow\).  Since
\(\DX\) is a derivation with \(\ordt(\DX F)\ge\ordt F+1\), it commutes with
summable sums, so \(\DX S=\sum_{r\ge1}U^{r-1}(\DX U)/(r-1)!=(\DX U)S\).  Hence
\(Z\DefinitionEquals(\exp U)S^{-1}\) lies in \(\PLlau\) by (iii), and
\[
  \partial Z=(\partial\exp U)S^{-1}-(\exp U)S^{-2}\DX S
  =(\DX U)Z-(\DX U)Z=0 .
\]
Thus \(Z\in\ker\DX\cap\PLlau=\K\) (\cref{plt:cor:dif-kernel}), while
\(Z\in1+t\PLpow\) by (iii); the only scalar in \(1+t\PLpow\) is \(1\).  For the
last sentence, apply the first with \(U=\log V\) and use
\cref{plt:eq:dif-exp-log-laws}.
\end{proof}

\begin{lemma}[Logarithmic derivative]\label{plt:lem:dif-logderiv}
Let \((\mathcal T,\partial,\exp)\) be an exponential extension and let
\(F\in\mathcal T\) admit a logarithm \(G\).  Then \(F\in\mathcal T^{\times}\)
and
\begin{equation}\label{plt:eq:dif-logderiv}
  \partial G=\frac{\partial F}{F}.
\end{equation}
Conversely, if \(G_{1},G_{2}\) are logarithms of \(F\) then
\(G_{1}-G_{2}\in\operatorname{Const}(\mathcal T)\), and if \(H\in\mathcal T\)
satisfies \(\partial H=\partial F/F\) then \(H-G\in\operatorname{Const}
(\mathcal T)\).
\end{lemma}

\begin{proof}
\(F=\exp G\) is a unit by (ii).  Differentiating \(F=\exp G\) gives
\(\partial F=(\partial G)F\), which is \cref{plt:eq:dif-logderiv}.  If
\(\exp G_{1}=\exp G_{2}\) then \(\exp(G_{1}-G_{2})=1\), so
\(\partial(G_{1}-G_{2})=\partial(1)/1=0\).  The last claim is immediate from
\cref{plt:eq:dif-logderiv}.
\end{proof}

\section{Power--log scales}
\label{plt:sec:dif-scales}

To say that an operation leaves the polynomial--logarithmic class \emph{and
every reasonable enlargement of it} we must quantify over the enlargements.
The following definition collects them: fractional or real exponents of \(X\),
rational or Laurent coefficients in \(L\), and any enlargement of the scalars.

\begin{definition}[Power--log scale]\label{plt:def:dif-scale}
Let \(\K'\) be a field of characteristic zero containing \(\K\), let
\(\Gamma\) be an ordered additive subgroup of \(\RealNumbers\) with
\(\IntegerNumbers\subseteq\Gamma\), fixed together with an embedding
\(\Gamma\hookrightarrow\K'\) of additive groups extending
\(\IntegerNumbers\subseteq\K'\) (so that the products \(\gamma p_{\gamma}\)
below are defined), and let \(C\) be one of
\(\K'[L]\), \(\K'(L)\), \(\K'((L^{-1}))\), each equipped with the derivation
\(\partial_{L}\).  The associated \emph{power--log scale} is the Hahn series
ring
\[
  \mathcal H(C,\Gamma)
  \DefinitionEquals
  \Bigl\{\,F=\sum_{\gamma\in S}p_{\gamma}t^{\gamma}\ :\
    p_{\gamma}\in C,\ S\subseteq\Gamma\ \text{well ordered}\,\Bigr\},
\]
with the Hahn product \cite{hahn}, the valuation
\(\ordt F=\min\SetOf{\gamma:p_{\gamma}\ne0}\), and the operator
\begin{equation}\label{plt:eq:dif-scale-derivation}
  \DX F\DefinitionEquals
  \sum_{\gamma\in S}t^{\gamma+1}\bigl(\partial_{L}p_{\gamma}
  -\gamma p_{\gamma}\bigr).
\end{equation}
\end{definition}

Thus \(\PLlau=\mathcal H(\K[L],\IntegerNumbers)\),
\(\PLrat=\mathcal H(\K(L),\IntegerNumbers)\),
\(\PLit=\mathcal H(\K((L^{-1})),\IntegerNumbers)\), and the Puiseux-log ring
\(\K[L]((t^{1/s}))\) is \(\mathcal H(\K[L],s^{-1}\IntegerNumbers)\).

% ed.: retitled to separate it from
% ed.: \cref{plt:lem:mul-asymptotic-basic}, which collects the basic
% ed.: properties of the relation \(\approx\) between a function and a series.
\begin{lemma}[Basic properties of a power--log scale]
\label{plt:lem:dif-scale-basic}
Let \(\mathcal H=\mathcal H(C,\Gamma)\) be a power--log scale.
\begin{enumerate}
\item \(\mathcal H\) is an integral domain, \(\ordt\) is additive on nonzero
      products, and \(\DX\) is a derivation of \(\mathcal H\).
\item \(\omega(\partial_{L}h)\ge\omega(h)+1\) for every \(h\in C\), where
      \(\omega\) denotes the \(L^{-1}\)-adic order computed in
      \(\K'((L^{-1}))\) after the expansion embedding
      (\cref{plt:thm:div-embedding}).  Consequently, for \(h\ne0\),
      \(\omega(\partial_{L}h/h)\ge1\), and \(\partial_{L}h-\gamma h\ne0\) for
      every \(\gamma\in\Gamma\setminus\SetOf{0}\).
\item \(\ordt(\DX F)\ge\ordt F+1\) for every \(F\in\mathcal H\), and
      \(\ker\DX=\K'\).
\item If \(F\in\mathcal H\) is nonzero and \(F^{-1}\) exists in the fraction
      field of \(\mathcal H\), then \(\ordt(\DX F/F)\ge1\), and
\begin{equation}\label{plt:eq:dif-logderiv-block}
  [t^{1}]\Bigl(\frac{\DX F}{F}\Bigr)
  =\frac{\partial_{L}b}{b}-e,
  \qquad e=\ordt F,\quad b=[t^{e}]F .
\end{equation}
\end{enumerate}
\end{lemma}

\begin{proof}
(i)  Hahn multiplication is well defined and makes \(\mathcal H\) a domain when
\(C\) is a domain, and \(\ordt\) is additive, because the minimum of the
Minkowski sum of two well-ordered sets is the sum of the minima and the
corresponding coefficient is the product of the leading coefficients
\cite{hahn}.  That \(\DX\) is additive is clear; for the Leibniz rule it
suffices, by \cref{plt:lem:dif-summable} applied in the Hahn setting, to check
it on monomials:
\[
  \DX\bigl(t^{\alpha}p\cdot t^{\beta}q\bigr)
  =t^{\alpha+\beta+1}\bigl((pq)'-(\alpha+\beta)pq\bigr)
  =\bigl(\DX(t^{\alpha}p)\bigr)t^{\beta}q+t^{\alpha}p\,\DX(t^{\beta}q).
\]

(ii)  For \(h=\sum_{j\ge j_{0}}a_{j}L^{-j}\) with \(a_{j_{0}}\ne0\), we have
\(\partial_{L}h=\sum_{j}(-j)a_{j}L^{-j-1}\), whose terms all have order
\(\ge j_{0}+1\).  The case \(C=\K'(L)\) reduces to this one because the
expansion embedding \(\iota_{0}\) of \cref{plt:thm:div-embedding} commutes with
\(\partial_{L}\): it is a ring homomorphism fixing \(L\), hence fixing
polynomials and their derivatives, and applying it to
\(\partial_{L}(u/v)=(u'v-uv')/v^{2}\) for \(u,v\in\K'[L]\) gives
\(\iota_{0}(\partial_{L}(u/v))=\partial_{L}(\iota_{0}(u/v))\) by the quotient
rule in \(\K'((L^{-1}))\).  Then \(\omega(\partial_{L}h/h)
=\omega(\partial_{L}h)-\omega(h)\ge1\).  If \(\gamma\ne0\) then
\(\omega(\gamma h)=\omega(h)<\omega(\partial_{L}h)\), so no cancellation is
possible and \(\partial_{L}h-\gamma h\ne0\).

(iii)  The bound is immediate from \cref{plt:eq:dif-scale-derivation}.  If
\(\DX F=0\) then \(\partial_{L}p_{\gamma}=\gamma p_{\gamma}\) for every
\(\gamma\); by (ii) this forces \(p_{\gamma}=0\) for \(\gamma\ne0\).  For
\(\gamma=0\) we get \(\partial_{L}p_{0}=0\); writing
\(p_{0}=\sum_{j\ge j_{0}}a_{j}L^{-j}\) in \(\K'((L^{-1}))\), this says
\(-ja_{j}=0\) for all \(j\), hence \(a_{j}=0\) for \(j\ne0\) and
\(p_{0}=a_{0}\in\K'\).

(iv)  \(\ordt(\DX F/F)=\ordt(\DX F)-\ordt F\ge1\) by (iii) and additivity.
Write \(F=t^{e}b(1+S)\) with \(\ordt S>0\), which is possible in the fraction
field.  Then, \(\DX\) being a derivation,
\[
  \frac{\DX F}{F}
  =\frac{\DX(t^{e})}{t^{e}}+\frac{\DX b}{b}+\frac{\DX S}{1+S}
  =-et+t\frac{\partial_{L}b}{b}+\frac{\DX S}{1+S},
\]
using \(\DX(t^{e})=-et^{e+1}\) and \(\DX b=t\,\partial_{L}b\).  The last term
has order \(\ge\ordt S+1>1\), so it does not contribute to the block at
\(t^{1}\).
\end{proof}

\begin{remark}[A common ambient scale]\label{plt:rmk:dif-common-scale}
Given two power--log scales \(\mathcal H(C_{1},\Gamma_{1})\) and
\(\mathcal H(C_{2},\Gamma_{2})\) inside one exponential extension, both embed
\(\DX\)-compatibly in \(\mathcal H(\K'''((L^{-1})),\Gamma_{1}+\Gamma_{2})\),
where \(\K'''\) is the compositum of the two scalar fields: on coefficients this
is \cref{plt:thm:div-embedding}, and on exponents it is the inclusion of
well-ordered supports.  Hence in the theorems below one may always assume that
the scale containing the input also contains the output.
\end{remark}

\section{Exponentials: the exact boundary}
\label{plt:sec:dif-exponential}

The mechanism is the logarithmic derivative.  If \(E=\exp F\) lies in a
power--log scale, then \(\DX F=\DX E/E\) is a logarithmic derivative of an
element of that scale, and \cref{plt:lem:dif-scale-basic} constrains such an
object twice: its outer order is at least \(1\), and its block at \(t^{1}\) has
positive \(L^{-1}\)-order.  Those two constraints are the whole boundary.

\begin{proposition}[Necessary conditions]\label{plt:prop:dif-exp-necessary}
Let \((\mathcal T,\partial,\exp)\) be an exponential extension, let
\(\mathcal H=\mathcal H(C,\Gamma)\subseteq\mathcal T\) be a power--log scale,
and let \(F\in\mathcal H\) be such that \(E\DefinitionEquals\exp F\) lies in
\(\mathcal H\).  Put \(e=\ordt E\in\Gamma\) and \(b=[t^{e}]E\ne0\).  Then
\begin{equation}\label{plt:eq:dif-exp-necessity}
  \ordt F\ge0
  \qquad\text{and}\qquad
  \partial_{L}\bigl([t^{0}]F\bigr)+e=\frac{\partial_{L}b}{b} .
\end{equation}
In particular \(\omega\bigl(\partial_{L}([t^{0}]F)+e\bigr)\ge1\); and if
\([t^{0}]F\) is a polynomial then
\begin{equation}\label{plt:eq:dif-exp-affine}
  \partial_{L}\bigl([t^{0}]F\bigr)=-e,
  \qquad\text{that is}\qquad
  [t^{0}]F=a-eL\ \text{ with }a\in\K' .
\end{equation}
\end{proposition}

\begin{proof}
\(E\) is nonzero and invertible in the fraction field of \(\mathcal H\), to
which \(\ordt\) extends with additivity.  By \cref{plt:lem:dif-logderiv},
\(\DX F=\DX E/E\), so \cref{plt:lem:dif-scale-basic}(iv) gives
\(\ordt(\DX F)\ge1\) and
\[
  [t^{1}]\DX F=\frac{\partial_{L}b}{b}-e .
\]
Write \(F=\sum_{\gamma}p_{\gamma}t^{\gamma}\).  By
\cref{plt:eq:dif-scale-derivation}, \(\ordt(\DX F)\ge1\) forces
\(\partial_{L}p_{\gamma}-\gamma p_{\gamma}=0\) for every \(\gamma<0\), and
\cref{plt:lem:dif-scale-basic}(ii) then gives \(p_{\gamma}=0\) for
\(\gamma<0\); that is, \(\ordt F\ge0\).  Again by
\cref{plt:eq:dif-scale-derivation}, the block of \(\DX F\) at \(t^{1}\) comes
only from \(\gamma=0\) and equals \(\partial_{L}p_{0}\), which proves
\cref{plt:eq:dif-exp-necessity}.  Its right-hand side has \(\omega\ge1\) by
\cref{plt:lem:dif-scale-basic}(ii).  If \(p_{0}\) is a polynomial then
\(\partial_{L}p_{0}+e\) is a polynomial, and a nonzero polynomial \(h\) has
\(\omega(h)=-\degL h\le0\); so \(\partial_{L}p_{0}+e=0\), which is
\cref{plt:eq:dif-exp-affine}.
% ed.: S4 and S6 note only that a "large part" other than an integral multiple
% of L may require new monomials.  The decisive point, which no source states,
% is that a logarithmic derivative has positive L^{-1}-order, whereas the
% derivative of a nonaffine polynomial has order at most -1.
\end{proof}

\begin{corollary}[The three ways out]\label{plt:cor:dif-exp-three-ways}
Let \(F\in\PLlau\).
\begin{enumerate}
\item If \(\ordt F<0\) then \(\exp F\) lies in \emph{no} power--log scale
      whatsoever.
\item If \(\ordt F\ge0\) and \(\degL[t^{0}]F\ge2\) then \(\exp F\) lies in no
      power--log scale.
\item If \(\ordt F\ge0\) and \([t^{0}]F=a+mL\), then \(\exp F\) lies in no
      power--log scale whose exponent group omits \(m\); in particular it lies
      in no Puiseux-log ring at all when \(m\) is irrational.
\end{enumerate}
\end{corollary}

\begin{proof}
By \cref{plt:rmk:dif-common-scale} we may work inside a single scale
\(\mathcal H(C,\Gamma)\) containing \(F\), and \(\PLlau\) has polynomial
coefficients, so \cref{plt:eq:dif-exp-affine} applies.  (i) If \(\ordt F<0\)
this contradicts the first half of \cref{plt:eq:dif-exp-necessity}.  (ii) Here
\(\partial_{L}[t^{0}]F\) is a nonzero polynomial of degree
\(\degL[t^{0}]F-1\ge1\), so it is not the constant \(-e\), contradicting
\cref{plt:eq:dif-exp-affine}.  (iii) Now \(\partial_{L}[t^{0}]F=m\), so
\cref{plt:eq:dif-exp-affine} forces \(m=-e\in\Gamma\).
\end{proof}

\begin{theorem}[Exponential closure boundary]\label{plt:thm:dif-exp-boundary}
Let \((\mathcal T,\partial,\exp)\) be an exponential extension and
\(F\in\PLlau\).  Then \(\exp F\in\PLlau\) if and only if
\begin{equation}\label{plt:eq:dif-exp-criterion}
  \ordt F\ge0,
  \qquad
  [t^{0}]F=a+mL\quad\text{with }m\in\IntegerNumbers,\ a\in\K,\
  \exp a\in\K .
\end{equation}
In that case
\begin{equation}\label{plt:eq:dif-exp-value}
  \exp F=(\exp a)\,t^{-m}\,\exp\bigl(F-[t^{0}]F\bigr),
  \qquad
  \exp\bigl(F-[t^{0}]F\bigr)\in1+t\PLpow .
\end{equation}
\end{theorem}

\begin{proof}
\emph{Sufficiency.}  Put \(V=F-[t^{0}]F\in t\PLpow\).  Since \(\exp\) is a
homomorphism and \(\exp L=t^{-1}\), for \(m\in\IntegerNumbers\) we get
\(\exp(mL)=(\exp L)^{m}=t^{-m}\), so
\(\exp F=(\exp a)t^{-m}\exp V\); and \(\exp V\in1+t\PLpow\) by
\cref{plt:lem:dif-series-exp}.  All three factors lie in \(\PLlau\).

\emph{Necessity.}  \Cref{plt:prop:dif-exp-necessary} with
\(\mathcal H=\PLlau\), \(C=\K[L]\), \(\Gamma=\IntegerNumbers\) gives
\(\ordt F\ge0\) and \([t^{0}]F=a-eL\) with \(a\in\K\) and
\(e=\ordt(\exp F)\in\IntegerNumbers\); put \(m=-e\).  Then
\cref{plt:eq:dif-exp-value} holds by the sufficiency computation, and solving
it for \(\exp a\) exhibits \(\exp a\) as a product of elements of \(\PLlau\),
so \(\exp a\in\PLlau\).  Finally \(\partial(\exp a)=(\partial a)\exp a=0\), so
\(\exp a\in\ker\DX\cap\PLlau=\K\) by \cref{plt:cor:dif-kernel}.
\end{proof}

\begin{corollary}[Fractional and irrational slopes]
\label{plt:cor:dif-exp-hahn}
Let \(\mathcal H=\mathcal H(\K'[L],\Gamma)\) be a power--log scale with
polynomial coefficients on which the ambient exponential is normalised, in the
sense that \(\exp V\in1+\SetOf{W\in\mathcal H:\ordt W>0}\) whenever
\(V\in\mathcal H\) has \(\ordt V>0\).  For \(F\in\mathcal H\), \(\exp
F\in\mathcal H\) if and only if \(\ordt F\ge0\), \([t^{0}]F=a+mL\) with
\(m\in\Gamma\), and the constant \(\kappa\DefinitionEquals\exp(a+mL)\,t^{m}\)
lies in \(\K'\).  Thus \(\exp(\tfrac12L)=X^{1/2}\) lies in the Puiseux-log ring
\(\K[L]((t^{1/2}))\) but not in \(\PLlau\), and \(\exp(\pi L)=X^{\pi}\)
requires a Hahn scale with \(\pi\in\Gamma\).
\end{corollary}

\begin{proof}
That \(\kappa\) is a constant is the computation
\(\partial\kappa=mt\kappa-mt\kappa=0\), using \(\DX(t^{m})=-mt^{m+1}\).
Sufficiency: \(\exp F=\kappa\,t^{-m}\exp(F-[t^{0}]F)\) with all factors in
\(\mathcal H\).  Necessity: \cref{plt:eq:dif-exp-affine} gives \(\ordt F\ge0\)
and \([t^{0}]F=a+mL\) with \(m=-e\in\Gamma\), and then
\(\kappa=\exp(F)\exp(-(F-[t^{0}]F))t^{m}\in\mathcal H\cap
\operatorname{Const}(\mathcal T)=\K'\) by
\cref{plt:lem:dif-scale-basic}(iii).
\end{proof}

\begin{remark}[Laurent coefficients admit more exponentials]
\label{plt:rmk:dif-exp-laurent}
The affine criterion \cref{plt:eq:dif-exp-affine} is a statement about
\emph{polynomial} coefficients.  It genuinely fails for Laurent coefficients:
\(F=L^{-1}\) has \(\ordt F=0\) and a \(t^{0}\)-block that is not affine in
\(L\), yet
\(\exp(L^{-1})=\sum_{k\ge0}L^{-k}/k!\in\K((L^{-1}))\subseteq\PLit\).  There is
no contradiction: \cref{plt:eq:dif-exp-necessity} is satisfied with
\(e=0\) and \(b=\exp(L^{-1})\), because \(\partial_{L}b/b=-L^{-2}\) does have
positive \(L^{-1}\)-order.  The exact criterion in a scale with Laurent
coefficients is precisely \cref{plt:eq:dif-exp-necessity}: the \(t^{0}\)-block
of \(F\), shifted by \(-eL\), must be a logarithm inside \(C\) of an element of
\(C\).
\end{remark}

\begin{example}[The three witnesses]\label{plt:ex:dif-exp-witnesses}
\(\exp(t^{-1})=\EulerE^{X}\) and \(\exp(t^{-1}L)=X^{X}\) fail (i);
\(\exp(L^{2})=X^{\log X}\) fails (ii); \(\exp(\tfrac12L)=X^{1/2}\) fails (iii)
for \(\Gamma=\IntegerNumbers\) but succeeds in the Puiseux-log ring.  By
contrast \(\exp(3L+t)=X^{3}\EulerE^{1/X}\in\PLlau\), and
\(\exp(Lt)=1+Lt+\tfrac12L^{2}t^{2}+\cdots\in1+t\PLpow\): smallness is measured
by the dominance order, not by the size of the coefficient polynomials.
\end{example}

\begin{remark}[The exponential hierarchy]\label{plt:rmk:dif-hierarchy}
\Cref{plt:cor:dif-exp-three-ways}(i) is the structural reason for the
transseries hierarchy.  The elements of negative outer order are exactly the
\emph{purely large} elements of the scale, and their exponentials are new
transmonomials that no enlargement of the exponent group or of the coefficient
field can express: enlarging \(\Gamma\) adds monomials \(t^{\gamma}\), all of
which have logarithmic derivative \(-\gamma t\) of outer order \(1\), whereas
\(\EulerE^{X}\) has logarithmic derivative \(1\).  Closing under \(\exp\) and
\(\log\) therefore requires a monomial group built by transfinite iteration,
which is precisely the construction of the logarithmic--exponential
transseries \cite{adh-book,vanderhoeven-book,edgar-beginners,costin-book}.  The
polynomial--logarithmic scale is the fragment of that field met after a
dominant logarithmic change of variable, and its virtue is that all closure
questions have the finite answers proved here.
\end{remark}

\begin{corollary}[Which elements are exponentials]
\label{plt:cor:dif-exp-group}
Let \(\K=\RealNumbers\) and let \(\mathcal T\) be an exponential extension with
\(\exp(\RealNumbers)\subseteq\RealNumbers\) equal to the real exponential.
Then
\[
  \mathcal E\DefinitionEquals
  \SetOf{F\in\PLlau:\exp F\in\PLlau}
  =\RealNumbers\oplus\IntegerNumbers L\oplus t\PLpow,
\]
and \(\exp\) restricts to a group isomorphism from \((\mathcal E,+)\) onto the
subgroup of \(\PLlau^{\times}\) consisting of the units whose scalar leading
coefficient is positive.
\end{corollary}

\begin{proof}
The description of \(\mathcal E\) is \cref{plt:thm:dif-exp-boundary} with
\(\K=\RealNumbers\), since \(\exp a\in\RealNumbers\) for every
\(a\in\RealNumbers\).  \Cref{plt:eq:dif-exp-value} shows that \(\exp\) maps
\(\mathcal E\) into the units with positive leading coefficient
\(\EulerE^{a}\).  It is onto: a unit is \(c\,t^{-m}(1+S)\) with
\(c>0,m\in\IntegerNumbers,S\in t\PLpow\)
(\cref{plt:thm:div-unit-criterion}), and \(F=\log c+mL+\log(1+S)\in\mathcal E\)
has \(\exp F=c\,t^{-m}(1+S)\) by
\cref{plt:eq:dif-exp-log-laws,plt:lem:dif-series-exp}.  It is injective: if
\(\exp F=1\) with \(F\in\mathcal E\) then \(\DX F=0\) by
\cref{plt:lem:dif-logderiv}, so \(F\in\K=\RealNumbers\) by
\cref{plt:cor:dif-kernel}, and \(\EulerE^{F}=1\) gives \(F=0\).
\end{proof}

\section{Logarithms: the leading block decides}
\label{plt:sec:dif-logarithm}

We now prove the counterpart for logarithms.  One elementary observation about
Laurent coefficients does all the work in the negative direction, and it will
reappear in \cref{plt:sec:differential} as the antiderivative obstruction; we
state it here in the form needed.

\begin{definition}[Logarithmic residue at \(L=\infty\)]
\label{plt:def:dif-residue}
For \(h=\sum_{j\ge j_{0}}a_{j}L^{-j}\in\K'((L^{-1}))\) put
\(\operatorname{res}(h)\DefinitionEquals a_{1}\), the coefficient of
\(L^{-1}\).  For \(h\in\K'(L)\), \(\operatorname{res}(h)\) means
\(\operatorname{res}\) of its expansion at \(L=\infty\); for \(h\in\K'[L]\),
\(\operatorname{res}(h)=0\).  In the classical convention
\(\operatorname{Res}_{L=\infty}(h\,\Differential L)=-\operatorname{res}(h)\).
\end{definition}

\begin{lemma}[Derivatives have no residue]\label{plt:lem:dif-res-derivative}
\(\operatorname{res}(\partial_{L}h)=0\) for every \(h\in\K'((L^{-1}))\).
\end{lemma}

\begin{proof}
\(\partial_{L}\bigl(\sum_{j\ge j_{0}}a_{j}L^{-j}\bigr)
=\sum_{j\ge j_{0}}(-j)a_{j}L^{-j-1}\), and the exponent \(-j-1\) equals
\(-1\) only for \(j=0\), whose coefficient \(-0\cdot a_{0}\) is \(0\).
\end{proof}

% ed.: no source identifies the image of \partial_L on Laurent-log
% coefficients, although S5 and S6 both assert closure of the class under
% antiderivatives and then remark vaguely that "new logarithmic levels can be
% forced" once inverse powers of L are admitted.  The converse of the preceding
% lemma is what makes that remark precise, so it is proved here and used both
% for the logarithm boundary and for the antiderivative obstruction.
\begin{lemma}[Image and kernel of \(\partial_{L}\)]
\label{plt:lem:dif-image-laurent}
Let \(g\in\K'((L^{-1}))\).  There exists \(h\in\K'((L^{-1}))\) with
\(\partial_{L}h=g\) if and only if \(\operatorname{res}(g)=0\), and in that
case \(h\) is unique up to an additive constant.  Equivalently,
\[
  \ker\partial_{L}=\K',
  \qquad
  \operatorname{im}\partial_{L}
  =\SetOf{g\in\K'((L^{-1})):\operatorname{res}(g)=0},
\]
so \(\operatorname{res}\) induces an isomorphism
\(\K'((L^{-1}))/\operatorname{im}\partial_{L}\cong\K'\).
\end{lemma}

\begin{proof}
Write \(h=\sum_{j\ge j_{0}}a_{j}L^{-j}\) and
\(g=\sum_{k\ge k_{0}}b_{k}L^{-k}\).  Then
\(\partial_{L}h=\sum_{j\ge j_{0}}(-j)a_{j}L^{-j-1}\), so
\(\partial_{L}h=g\) holds if and only if
\begin{equation}\label{plt:eq:dif-image-coefficients}
  -j\,a_{j}=b_{j+1}
  \qquad\text{for every }j .
\end{equation}
For \(j\ne0\) this determines \(a_{j}=-b_{j+1}/j\) uniquely, and characteristic
zero is what makes the division legitimate.  For \(j=0\) the left side is
\(0\), so \eqref{plt:eq:dif-image-coefficients} is solvable exactly when
\(b_{1}=\operatorname{res}(g)=0\), and then \(a_{0}\) is unconstrained.  This
gives both the criterion and uniqueness up to the additive constant \(a_{0}\).

The candidate \(h\) so defined is a legitimate element of \(\K'((L^{-1}))\):
if \(b_{k}=0\) for \(k<k_{0}\) then \(a_{j}=0\) for \(j<k_{0}-1\), so the
support of \(h\) is bounded below by \(k_{0}-1\).  Finally
\(\ker\partial_{L}=\K'\) is the case \(g=0\) of the same computation.
\end{proof}

\begin{corollary}[Residue of a logarithmic derivative]
\label{plt:cor:dif-logderiv-residue}
For nonzero \(h\in\K'((L^{-1}))\), \(\operatorname{res}(\partial_{L}h/h)
=-\omega(h)\).  In particular, for a nonzero polynomial \(p\),
\(\operatorname{res}(\partial_{L}p/p)=\degL p\), and for
\(p=a\prod_{i}(L-\lambda_{i})^{m_{i}}\) over \(\overline{\K'}\) the residue of
\(\partial_{L}p/p\) at the finite point \(\lambda_{i}\) is \(m_{i}\ne0\).
\end{corollary}

\begin{proof}
Write \(h=a_{j_{0}}L^{-j_{0}}(1+w)\) with \(\omega(w)\ge1\) and
\(j_{0}=\omega(h)\).  Then \(\partial_{L}h/h=-j_{0}L^{-1}+\partial_{L}w/(1+w)\)
by the Leibniz rule, and \(\omega(\partial_{L}w/(1+w))\ge\omega(w)+1\ge2\) by
\cref{plt:lem:dif-scale-basic}(ii); so the coefficient of \(L^{-1}\) is
\(-j_{0}\).  For a polynomial, \(\omega(p)=-\degL p\).  The last claim is the
partial-fraction identity \(\partial_{L}p/p=\sum_{i}m_{i}/(L-\lambda_{i})\),
and \(m_{i}\ne0\) in \(\K'\) because \(\operatorname{char}\K'=0\).
\end{proof}

\begin{theorem}[Logarithm boundary]\label{plt:thm:dif-log-boundary}
Let \((\mathcal T,\partial,\exp)\) be an exponential extension, let
\(F\in\PLlau\) be nonzero with \(n_{0}=\ordt F\), leading block
\(p=p_{n_{0}}\), scalar leading coefficient \(a\in\K^{\times}\) and
\(d=\degL p\), and write \(F=t^{n_{0}}p\,(1+S)\) with \(S\in t\K(L)[[t]]\).
Assume \(F\) has a logarithm \(G\) in \(\mathcal T\) and that \(L\) has a
logarithm \(\Lambda\) in \(\mathcal T\).  Then:
\begin{enumerate}
\item \(G\) is determined modulo \(\operatorname{Const}(\mathcal T)\) by
\begin{equation}\label{plt:eq:dif-log-decomposition}
  G\equiv
  -n_{0}L
  +d\,\Lambda
  +\log\Bigl(\frac{p}{aL^{d}}\Bigr)
  +\log(1+S)
  \pmod{\operatorname{Const}(\mathcal T)},
\end{equation}
where \(\log(p/(aL^{d}))=\sum_{r\ge1}\frac{(-1)^{r+1}}{r}w^{r}\) with
\(w=p/(aL^{d})-1\in L^{-1}\K[L^{-1}]\), and \(\log(1+S)\) is the series
\cref{plt:eq:dif-log1pU} with coefficients in \(\K(L)\).  All the displayed
terms except \(d\Lambda\) lie in \(\PLit\).
\item \(G\) lies in some power--log scale if and only if \(d=0\), that is
      (\cref{plt:thm:div-unit-criterion}) if and only if \(F\) is a unit of
      \(\PLlau\).  When \(d=0\) one may take
      \(G=c-n_{0}L+\log(1+S)\in\K(c)[L]((t))\) for any \(c\in\mathcal T\) with
      \(\exp c=a\).
\end{enumerate}
\end{theorem}

\begin{proof}
(i)  Call the right-hand side of \cref{plt:eq:dif-log-decomposition} \(H\).  By
\cref{plt:lem:dif-logderiv} it suffices to prove \(\partial H=\DX F/F\), for
then \(G-H\) is a constant.  Now \(\partial\Lambda=\DX L/L=t/L\) by
\cref{plt:lem:dif-logderiv} applied to \(L\), and \(\DX\) applied to the two
logarithmic series termwise (legitimate by
\cref{plt:lem:dif-summable}, the families being summable in the
\(L^{-1}\)-adic and \(t\)-adic order respectively) gives
\(\DX\log(1+w)=\DX w/(1+w)=t\,\partial_{L}w/(1+w)\) and
\(\DX\log(1+S)=\DX S/(1+S)\).  Hence
\[
  \partial H
  =-n_{0}t+\frac{d\,t}{L}
   +t\,\frac{\partial_{L}w}{1+w}
   +\frac{\DX S}{1+S}
  =-n_{0}t+t\,\frac{\partial_{L}p}{p}+\frac{\DX S}{1+S},
\]
because \(p=aL^{d}(1+w)\) gives \(\partial_{L}p/p=d/L+\partial_{L}w/(1+w)\).
On the other side, \(F=t^{n_{0}}p(1+S)\) and the Leibniz rule give exactly the
same expression for \(\DX F/F\).  The membership claims are clear:
\(w\in L^{-1}\K[L^{-1}]\) makes \(\log(1+w)\in L^{-1}\K[[L^{-1}]]\), and
\(\log(1+S)\in t\K(L)[[t]]\subseteq\PLit\) after
\cref{plt:thm:div-embedding}.

(ii)  \emph{Sufficiency.}  If \(d=0\) then \(p=a\), \(w=0\), and
\cref{plt:eq:dif-log-decomposition} exhibits a logarithm in \(\K(c)[L]((t))\)
once the constant is fixed as \(c\); note \(S\in t\K[L][[t]]\) in this case,
since \(F=at^{n_{0}}(1+S)\) with \(F\in\PLlau\).

\emph{Necessity.}  Suppose \(G\in\mathcal H(C,\Gamma)\) and \(d\ge1\).  By
\cref{plt:lem:dif-logderiv} and the computation in (i),
\[
  \DX G=\frac{\DX F}{F},
  \qquad
  [t^{1}]\DX G=\frac{\partial_{L}p}{p}-n_{0}.
\]
On the other hand \cref{plt:eq:dif-scale-derivation} gives \([t^{1}]\DX G
=\partial_{L}g_{0}\) with \(g_{0}=[t^{0}]G\in C\).  Hence
\(\partial_{L}(g_{0}+n_{0}L)=\partial_{L}p/p\) with \(g_{0}+n_{0}L\in C\).
Apply \(\operatorname{res}\): the left side vanishes by
\cref{plt:lem:dif-res-derivative}, the right side equals \(d\ne0\) by
\cref{plt:cor:dif-logderiv-residue} -- a contradiction.  The argument covers
all three coefficient rings at once, since each embeds \(\partial_{L}\)-
compatibly in \(\K'((L^{-1}))\) (\cref{plt:lem:dif-scale-basic}(ii)); for
\(C=\K'[L]\) one may argue even more directly, since
\(\partial_{L}p/p\) has a pole when \(d\ge1\) and is therefore not a
polynomial.
\end{proof}

\begin{example}[A forced nested logarithm]\label{plt:ex:dif-log-nested}
Let \(F=L^{2}t^{-3}\bigl(1+(L+1)t\bigr)\), so \(n_{0}=-3\), \(p=L^{2}\),
\(a=1\), \(d=2\), \(w=0\).  Then
\[
  \log F\equiv 3L+2\Lambda+(L+1)t-\tfrac12(L+1)^{2}t^{2}
  +\tfrac13(L+1)^{3}t^{3}+\FormalBigOAt{t^{4}}{t}.
\]
The new generator \(\Lambda=\log\log X\) is forced by the factor \(L^{2}\), not
by the small correction, and it enters with coefficient exactly
\(d=2\).
% ed.: S6 states this example with a leading factor L^{-2}, which is not an
% element of K[L]((t)); the phenomenon is the same, and the general statement
% is that the coefficient of Lambda is -omega of the leading coefficient.
\end{example}

\begin{corollary}[Logarithms inside the class]\label{plt:cor:dif-log-unit}
Let \(F\in\PLlau\) be nonzero.  Then \(F\) has a logarithm in \(\K'[L]((t))\)
for some field \(\K'\supseteq\K\) if and only if \(F\) is a unit of
\(\PLlau\); one may take \(\K'=\K(c)\) with \(\exp c\) the scalar leading
coefficient of \(F\).  If \(F\) is not a unit, a logarithm exists in
\(\PLit\oplus\K\Lambda\oplus\operatorname{Const}(\mathcal T)\) and in no
power--log scale.
\end{corollary}

\begin{proof}
Immediate from \cref{plt:thm:dif-log-boundary}; the second sentence is
\cref{plt:eq:dif-log-decomposition}, whose terms other than \(d\Lambda\) lie in
\(\PLit\).
\end{proof}

\section{Powers}
\label{plt:sec:dif-powers}

\begin{proposition}[General powers]\label{plt:prop:dif-powers}
Let \(F\in\PLlau\) be a unit, \(F=a\,t^{n_{0}}(1+S)\) with \(a\in\K^{\times}\),
\(n_{0}\in\IntegerNumbers\), \(S\in t\PLpow\), let \(\alpha\in\K\), fix
\(c\in\mathcal T\) with \(\exp c=a\), and let
\(G=c-n_{0}L+\log(1+S)\) be the corresponding logarithm of \(F\)
(\cref{plt:thm:dif-log-boundary}).  Then
\begin{equation}\label{plt:eq:dif-power-value}
  F^{\alpha}=\exp(\alpha G)
  =\mu\;t^{\alpha n_{0}}\,(1+S)^{\alpha},
  \qquad
  (1+S)^{\alpha}=\sum_{r\ge0}\binom{\alpha}{r}S^{r}\in1+t\PLpow ,
\end{equation}
where \(\mu\DefinitionEquals\exp(\alpha c-\alpha n_{0}L)\,t^{\alpha n_{0}}\) is
a constant of \(\mathcal T\), equal to \(\exp(\alpha c)\) whenever
\(\alpha n_{0}\in\IntegerNumbers\).  Consequently:
\begin{enumerate}
\item if \(\alpha n_{0}\in\IntegerNumbers\) and \(\exp(\alpha c)\in\K'\), then
      \(F^{\alpha}\in\K'[L]((t))\);
\item if \(\alpha n_{0}\in s^{-1}\IntegerNumbers\setminus\IntegerNumbers\),
      then \(F^{\alpha}\) lies in the Puiseux-log ring \(\K'[L]((t^{1/s}))\)
      and in no scale with \(\alpha n_{0}\notin\Gamma\); for
      \(\alpha n_{0}\) irrational a Hahn scale is required;
\item if \(F\) is not a unit, with leading block \(p\) of degree \(d\ge1\),
      then \(F^{\alpha}\) lies in a power--log scale only if
      \(\alpha d\in\IntegerNumbers\).
\end{enumerate}
\end{proposition}

\begin{proof}
Since \(\exp\) is a homomorphism, \(\exp(\alpha G)=\exp(\alpha c-\alpha
n_{0}L)\exp(\alpha\log(1+S))\), and the second factor is
\((1+S)^{\alpha}\) by the last identity of \cref{plt:eq:dif-log-laws} together
with \cref{plt:lem:dif-series-exp}; inserting \(t^{\alpha n_{0}}
t^{-\alpha n_{0}}\) gives \cref{plt:eq:dif-power-value}.  That \(\mu\) is a
constant is the computation of \cref{plt:cor:dif-exp-hahn}, and for
\(\alpha n_{0}\in\IntegerNumbers\) the normalisation \(\exp L=t^{-1}\) gives
\(\exp(-\alpha n_{0}L)=t^{\alpha n_{0}}\), whence \(\mu=\exp(\alpha c)\).
Items (i) and (ii) follow, the negative half of (ii) from
\cref{plt:cor:dif-exp-three-ways}(iii) applied to \(\alpha G\).  For (iii),
\cref{plt:eq:dif-log-decomposition} shows that \(\alpha G\) contains the term
\(\alpha d\,\Lambda\); equivalently, \(\DX(F^{\alpha})/F^{\alpha}=\alpha\DX
F/F\) has \(\operatorname{res}=\alpha d\) by
\cref{plt:cor:dif-logderiv-residue}, while by that same corollary the
logarithmic derivative of a nonzero element of any power--log scale has
residue \(-\omega(b)\in\IntegerNumbers\) when the coefficients are polynomial,
and in general an integer whenever the leading coefficient lies in
\(\K'((L^{-1}))\).  Hence \(\alpha d\in\IntegerNumbers\).
\end{proof}

\begin{remark}[Branches]\label{plt:rmk:dif-branches}
Over \(\ComplexNumbers\) the scalars \(\exp c=a^{1}\) and \(\exp\alpha c\) both
depend on the choice of \(c\), and different choices differ by
\(\exp(2\pi\ImaginaryUnit k\alpha)\); over \(\RealNumbers\) with \(a>0\) the
principal choice \(c=\PrincipalLogarithm a\) is available.  A formula for
\(F^{\alpha}\) is incomplete until the branch is named, and the choice
propagates into every later block.  This is a genuine mathematical datum, not a
notational nuisance: the two square roots of \(X^{2}(1+t)\) are
\(\pm X(1+\tfrac12t-\tfrac18t^{2}+\cdots)\), and only the choice of sign makes
the expansion an asymptotic expansion of a given branch as \(X\to+\infty\).
\end{remark}

\section{The closure boundary}
\label{plt:sec:dif-boundary}

We can now state the closure boundary as a theorem rather than a table.  The
first part says that the polynomial--logarithmic class is generated from its
two generators by one construction, Laurent assembly; the second says that a
long list of further operations does not enlarge it; the third exhibits, for
each operation that does enlarge it, a witness and the theorem that proves
non-membership.

\begin{theorem}[The generated class and its boundary]
\label{plt:thm:dif-generated}
Let \((\mathcal T,\partial,\exp)\) be an exponential extension.
\begin{enumerate}
\item \(\PLlau\) is the smallest subring \(\mathcal A\subseteq\mathcal T\) that
      contains \(\K\cup\SetOf{t,t^{-1},L}\) and is closed under Laurent
      assembly: whenever \(p_{n}\in\K[L]\) for \(n\ge n_{0}\) and every
      \(p_{n}t^{n}\) lies in \(\mathcal A\), the element
      \(\sum_{n\ge n_{0}}p_{n}t^{n}\in\PLlau\) lies in \(\mathcal A\).
\item \(\PLlau\) is closed under: the ring operations; inversion of its units;
      \(\Phi(U)\) for \(\Phi\in\K[[z]]\) and \(\ordt U\ge1\); the derivation
      \(\DX\) and the Euler derivation \(\EulerOp\); antidifferentiation; the
      exponential of any element of \(\K_{0}\oplus\IntegerNumbers L\oplus
      t\PLpow\), where \(\K_{0}=\SetOf{a\in\K:\exp a\in\K}\); and the logarithm
      of any unit whose scalar leading coefficient has a logarithm in \(\K\).
\item Each of the following leaves \(\PLlau\), and none of the enlargements it
      forces can be avoided: reciprocal of a nonunit; a power with
      \(\alpha n_{0}\notin\IntegerNumbers\); the logarithm of a nonunit; the
      exponential of an element of negative outer order or with
      \(\degL[t^{0}]F\ge2\).  Witnesses and minimal enlargements are given in
      the proof and collected in \cref{plt:tab:dif-closure}.
\end{enumerate}
\end{theorem}

\begin{proof}
(i)  The intersection of any family of subrings closed under Laurent assembly
is again one, so a smallest \(\mathcal A\) exists.  \(\PLlau\) is a subring of
\(\mathcal T\) containing the generators and is closed under Laurent assembly
by definition, so \(\mathcal A\subseteq\PLlau\).  Conversely \(\mathcal A\)
contains \(\K[L][t,t^{-1}]\), hence every \(p_{n}t^{n}\), hence by assembly
every element of \(\PLlau\).

(ii)  Ring operations and inversion of units are
\cref{plt:thm:div-unit-criterion}; \(\Phi(U)\) is
\cref{plt:prop:dif-substitution}; \(\DX\) and \(\EulerOp\) are
\cref{plt:prop:dif-closure-classes}; antidifferentiation is
\cref{plt:thm:dif-antiderivative}; the exponential and logarithm clauses are
\cref{plt:thm:dif-exp-boundary} and \cref{plt:cor:dif-log-unit}.

(iii)  The witnesses and their proofs are: \(1/(L+1)\notin\PLlau\)
(\cref{plt:thm:div-frac-strict}), with minimal enlargement \(\PLrat\), or
\(\PLit\) after expansion; \(t^{1/2}=X^{-1/2}\) for \(F=t\), \(\alpha=1/2\)
(\cref{plt:prop:dif-powers}), with minimal enlargement a Puiseux-log or Hahn
scale; \(\log L=\Lambda\) for \(F=L\)
(\cref{plt:thm:dif-log-boundary}), with enlargement the adjunction of the new
generator \(\Lambda\); \(\EulerE^{X}\) and \(X^{\log X}\)
(\cref{plt:cor:dif-exp-three-ways}), with enlargement a genuinely
exponential transseries field.  In each case the cited statement proves
non-membership in every power--log scale of the relevant type, so no
enlargement of the coefficient field or of the exponent group alone suffices.
\end{proof}

\begin{table}[htbp]
\centering
\caption{Closure of the polynomial--logarithmic scale under the elementary
operations.  ``Scale'' means any power--log scale \(\mathcal H(C,\Gamma)\) in
the sense of the definition in this section.}
\label{plt:tab:dif-closure}
\small
\begin{tabular}{@{}p{0.24\textwidth}p{0.30\textwidth}p{0.36\textwidth}@{}}
\toprule
Operation & Hypothesis & Result \\
\midrule
\(F+G\), \(FG\) & none & \(\PLlau\) \\
\(1/F\) & leading block in \(\K^{\times}\) & \(\PLlau\) \\
\(1/F\) & leading block of degree \(\ge1\) & \(\PLrat\); after expansion
  \(\PLit\) \\
\(\Phi(U)\), \(\Phi\in\K[[z]]\) & \(\ordt U\ge1\) & \(\PLpow\)
  (\cref{plt:prop:dif-substitution}) \\
\(\DX F\), \(\EulerOp F\) & none & same class
  (\cref{plt:prop:dif-closure-classes}) \\
antiderivative & none & \(\PLlau\) (\cref{plt:thm:dif-antiderivative}) \\
antiderivative in \(\PLit\) & \(\operatorname{res}F=0\) & \(\PLit\); otherwise
  adjoin \(\Lambda\) (\cref{plt:thm:dif-anti-it}) \\
\(F^{\alpha}\) & \(F\) a unit, \(\alpha\ordt F\in\IntegerNumbers\),
  \(a^{\alpha}\in\K\) & \(\PLlau\) \\
\(F^{\alpha}\) & \(\alpha\ordt F\notin\IntegerNumbers\) & Puiseux-log or Hahn
  scale \\
\(\log F\) & \(F\) a unit, \(\log a\in\K\) & \(\PLlau\)
  (\cref{plt:cor:dif-log-unit}) \\
\(\log F\) & leading block of degree \(d\ge1\) & \(\PLit\oplus\K\Lambda\);
  no scale (\cref{plt:thm:dif-log-boundary}) \\
\(\exp F\) & \(\ordt F\ge0\), \([t^{0}]F=a+mL\), \(m\in\IntegerNumbers\),
  \(\EulerE^{a}\in\K\) & \(\PLlau\)
  (\cref{plt:thm:dif-exp-boundary}) \\
\(\exp F\) & \(\ordt F\ge0\), \([t^{0}]F=a+mL\), \(m\notin\IntegerNumbers\)
  & Hahn scale with \(m\in\Gamma\) \\
\(\exp F\) & \(\degL[t^{0}]F\ge2\), or \(\ordt F<0\) & no scale; exponential
  transseries (\cref{plt:cor:dif-exp-three-ways}) \\
\bottomrule
\end{tabular}
\end{table}

\begin{warningbox}
Three errors recur in the sources and in informal practice.
\begin{itemize}
\item \emph{``Small'' is not about the coefficients.}  \(Lt\) is small and
  \(L^{-1}\) is not, in the sense that decides summability of
  \(\sum_{r}c_{r}U^{r}\): the criterion is \(\ordt U\ge1\), not the numerical
  size of the blocks.  For \(\ordt U=0\) the series
  \(\log(1+U)\) need not be summable at all, and a normalisation must precede
  the expansion.
\item \emph{A closure table is not a theorem.}  ``\(\exp F\) may require new
  monomials'' is not a statement one can act on.  What is true is
  \cref{plt:cor:dif-exp-three-ways}: for \(\ordt F<0\) no enlargement of the
  coefficient field or of the exponent group ever suffices.
\item \emph{Formal identities are not analytic ones.}  Every statement in this
  section is an identity between formal objects.  That
  \(\exp(\log(1+U))=1+U\) holds coefficientwise says nothing about the
  convergence of either series, and an equality
  \(F=G+\FormalBigOAt{t^{N}}{t}\) carries no bound of the form
  \(\BigOAt{X^{-N}}{X\to+\infty}\) without a separate realisation theorem.
\end{itemize}
\end{warningbox}

\begin{summarybox}
Substitution of a series of positive outer order into a formal power series is
unconditional and finitely determined: \(\Trunc{\Phi(U)}{N}\) depends on
\(c_{0},\dots,c_{\Floor{(N-1)/q}}\) and on \(\Trunc{U}{N}\) alone.  Hence
\((1+U)^{\alpha}\), \(\log(1+U)\) and \(\exp U\) stay in \(\PLpow\).  Beyond
that, the leading data decide.  A logarithm stays in the scale precisely when
the leading block is a scalar; otherwise the generator \(\Lambda=\log L\)
appears with coefficient \(\degL p\), and the coefficients must be expanded at
\(L=\infty\).  An exponential stays in a power--log scale precisely when
\(\ordt F\ge0\) and the \(t^{0}\)-block is affine in \(L\) with slope in the
exponent group; otherwise the object is a new transmonomial and the
exponential hierarchy begins.
\end{summarybox}

\chapter{Differentiation and antiderivatives}
\label{plt:sec:differential}

Differentiation is the operation for which the polynomial--logarithmic scale
behaves best: it is unconditionally closed, it raises the outer order by
exactly one except in one identifiable case, and -- unlike almost every other
restricted asymptotic scale -- so does its inverse.  This chapter proves those
statements, in the sharp form, and then determines exactly which elements of
each ambient class have an antiderivative inside that class.  The answer is a
residue: for \(\PLlau\) there is no obstruction at all, for \(\PLit\) the
obstruction is a single linear functional, and for \(\PLrat\) there is one
linear condition at every pole of every coefficient block.

\section{The distinguished derivation}
\label{plt:sec:dif-derivation}

In the formal algebra \(t\) and \(L\) are independent generators.  The analytic
relation \(t=X^{-1}\), \(L=\log X\) enters the algebra in exactly one place: it
selects a derivation.

\begin{theorem}[Existence and uniqueness]\label{plt:thm:dif-derivation}
Let \(\partial_{L}\) and \(\partial_{t}\) be the formal partial derivatives of
\(\PLlau\), each holding the other generator fixed.
\begin{enumerate}
\item The operator
\begin{equation}\label{plt:eq:dif-DX-definition}
  \DX\DefinitionEquals t\,\partial_{L}-t^{2}\,\partial_{t}
\end{equation}
is a \(\K\)-derivation of \(\PLlau\) with \(\DX t=-t^{2}\) and \(\DX L=t\),
and it maps \(\PLpow\) into \(t\PLpow\).
\item \(\DX\) is the unique \(\K\)-derivation \(\delta\) of \(\PLlau\) with
      \(\delta t=-t^{2}\), \(\delta L=t\) and \(\delta(\PLpow)\subseteq
      \PLpow\).  The last hypothesis may be replaced by \(t\)-adic continuity
      of \(\delta\).
\item Without any such hypothesis, uniqueness still holds on the subring
      \(\K[L][t,t^{-1}]\).
\end{enumerate}
\end{theorem}

\begin{proof}
(i)  Both \(\partial_{L}\) and \(\partial_{t}\) are \(\K\)-derivations of
\(\PLlau\).  Indeed each is defined blockwise or termwise,
\(\partial_{L}\bigl(\sum p_{n}t^{n}\bigr)=\sum p_{n}'t^{n}\) and
\(\partial_{t}\bigl(\sum p_{n}t^{n}\bigr)=\sum np_{n}t^{n-1}\), both supports
being bounded below; and the Leibniz rule holds because every block of a
product is the finite sum \([t^{k}](FG)=\sum_{i+j=k}p_{i}q_{j}\), to which the
Leibniz rule for polynomials, respectively the classical computation
\(k\sum_{i+j=k}p_{i}q_{j}=\sum_{i+j=k}(ip_{i}q_{j}+jp_{i}q_{j})\), applies.  A
left \(\PLlau\)-linear combination of derivations is a derivation, so \(\DX\)
is one.  The values on the generators are immediate, and the block formula
\cref{plt:eq:dif-block} below gives \(\DX(\PLpow)\subseteq t\PLpow\).

(ii)  Let \(\delta\) be as stated and put \(\eta=\delta-\DX\), a
\(\K\)-derivation with \(\eta t=\eta L=0\), hence \(\eta(t^{-1})
=-t^{-2}\eta(t)=0\) and \(\eta=0\) on \(\K[L][t,t^{-1}]\).  Let
\(F\in\PLlau\) and \(N>\ordt F\).  Write \(F=\Trunc{F}{N}+t^{N}R\) with
\(\Trunc{F}{N}\in\K[L][t,t^{-1}]\) and \(R\in\PLpow\).  Then
\[
  \eta F=\eta\bigl(t^{N}R\bigr)
  =Nt^{N-1}\eta(t)R+t^{N}\eta R=t^{N}\eta R .
\]
By hypothesis \(\delta R\in\PLpow\) and \(\DX R\in t\PLpow\), so
\(\eta R\in\PLpow\) and \(\ordt(\eta F)\ge N\).  As \(N\) was arbitrary,
\(\eta F=0\).  If instead \(\eta\) is continuous, choose \(N_{0}\) with
\(\ordt(\eta G)\ge0\) whenever \(\ordt G\ge N_{0}\); applying the display with
\(N+N_{0}\) in place of \(N\) gives \(\eta F=t^{N}\eta(t^{N_{0}}R)\) of order
\(\ge N\), and the same conclusion.

(iii)  Contained in the first two sentences of (ii).
\end{proof}

\begin{remark}[The hypothesis is not decorative]
\label{plt:rmk:dif-uniqueness-hypothesis}
% ed.: the notation catalogue and four of the six sources assert that D_X is
% "the unique derivation" with the two generator values, with no hypothesis.
% The argument above needs a bound on delta on the power-series part; without
% one, the passage from the dense subring to the completion is not available.
The sources, and the notation catalogue, call \(\DX\) \emph{the} unique
derivation with \(\DX t=-t^{2}\) and \(\DX L=t\).  That is correct on the dense
subring \(\K[L][t,t^{-1}]\) and correct on \(\PLlau\) under the boundedness
hypothesis of (ii), which every derivation arising in practice satisfies; but
the passage to the completion genuinely uses a hypothesis, since a derivation
of \(\PLlau\) is not a priori bounded on \(\PLpow\).  We have stated the
hypothesis rather than suppressing it.
\end{remark}

\section{One block, and the Euler derivation}
\label{plt:sec:dif-block}

\begin{proposition}[Block formula]\label{plt:prop:dif-block}
For \(n\in\IntegerNumbers\) and \(p\in\K[L]\),
\begin{equation}\label{plt:eq:dif-block}
  \DX\bigl(t^{n}p(L)\bigr)=t^{n+1}\bigl(\partial_{L}-n\bigr)p(L),
\end{equation}
and hence, for \(F=\sum_{n\ge n_{0}}p_{n}t^{n}\in\PLlau\),
\begin{equation}\label{plt:eq:dif-series-derivative}
  \DX F=\sum_{n\ge n_{0}}t^{n+1}\bigl(p_{n}'-np_{n}\bigr),
  \qquad\text{that is}\qquad
  [t^{n+1}]\DX F=\bigl(\partial_{L}-n\bigr)[t^{n}]F .
\end{equation}
The same formula holds with \(p\) in \(\K(L)\) or \(\K((L^{-1}))\) and with
\(n\) replaced by any exponent \(\alpha\) of a Puiseux-log or Hahn scale:
\(\DX(t^{\alpha}p)=t^{\alpha+1}(p'-\alpha p)\).  Each output block receives a
contribution from exactly one input block.
\end{proposition}

\begin{proof}
\(\partial_{L}(t^{n}p)=t^{n}p'\) and \(\partial_{t}(t^{n}p)=nt^{n-1}p\), so
\(\DX(t^{n}p)=t^{n+1}p'-nt^{n+1}p\).  The series form follows because \(\DX\)
is additive and commutes with summable sums, being continuous by
\cref{plt:eq:dif-block}.  For the general scale this is
\cref{plt:eq:dif-scale-derivation}, which was verified to be a derivation in
\cref{plt:lem:dif-scale-basic}(i).
\end{proof}

\begin{proposition}[The Euler derivation]\label{plt:prop:dif-euler}
Put \(\EulerOp\DefinitionEquals t^{-1}\DX\), the formal counterpart of
\(X\Differential/\Differential X\).  Then \(\EulerOp\) is a
\(\K\)-derivation of \(\PLlau\) with \(\EulerOp t=-t\) and \(\EulerOp L=1\),
it preserves each outer block,
\begin{equation}\label{plt:eq:dif-euler-block}
  \EulerOp\bigl(t^{n}p(L)\bigr)=t^{n}\bigl(\partial_{L}-n\bigr)p(L),
\end{equation}
and therefore \(\ordt(\EulerOp F)\ge\ordt F\) for all \(F\).  On the block
\(t^{n}\K[L]\) the operator \(\EulerOp\) acts as \(\partial_{L}-n\), whose only
eigenvalue is \(-n\); it is invertible on that block if and only if
\(n\ne0\).
\end{proposition}

\begin{proof}
A derivation multiplied by a ring element is again a derivation, and
\(t^{-1}\in\PLlau\); the values on the generators and
\cref{plt:eq:dif-euler-block} follow from \cref{plt:eq:dif-block}.  On
\(t^{n}\K[L]\), the operator \(\partial_{L}-n\) is, in the basis
\(1,L,\dots,L^{D}\) of polynomials of degree \(\le D\), triangular with every
diagonal entry \(-n\), because \(\partial_{L}L^{i}=iL^{i-1}\) lowers the degree.
Hence its determinant on that space is \((-n)^{D+1}\), which vanishes exactly
when \(n=0\); and \(\partial_{L}\) itself is not injective on \(\K[L]\).
\end{proof}

\begin{remark}[Resonance is a zero eigenvalue]\label{plt:rmk:dif-resonance}
\Cref{plt:prop:dif-euler} is the conceptual source of everything that follows.
The scale decomposes into the \(\EulerOp\)-invariant blocks \(t^{n}\K[L]\); on
block \(n\) the Euler derivation is the scalar \(-n\) plus a nilpotent, so it
is invertible exactly when \(n\ne0\).  The block \(n=0\) is \emph{resonant}:
there \(\EulerOp\) is the nilpotent \(\partial_{L}\), with kernel \(\K\) and
cokernel \(0\) on \(\K[L]\) -- but with a one-dimensional cokernel on
\(\K((L^{-1}))\).  Every statement below about constants of integration,
logarithms and residues is a statement about that single block.
\end{remark}

\section{Repeated derivatives}
\label{plt:sec:dif-repeated}

\begin{definition}[Repeated-derivative operator]\label{plt:def:dif-shiftop}
For \(n\) in the exponent group and \(k\in\NaturalNumbers\) put
\begin{equation}\label{plt:eq:dif-shiftop}
  \shiftop{n}{k}\DefinitionEquals\prod_{j=0}^{k-1}
  \bigl(\partial_{L}-n-j\bigr),
  \qquad \shiftop{n}{0}\DefinitionEquals1 .
\end{equation}
The \(k=0\) product is empty and acts as the identity; the factors commute,
each being a polynomial in \(\partial_{L}\).
\end{definition}

% ed.: notation reconciliation.  Two sources write Delta_n for the single
% factor partial_L - n and never introduce the second index; the normative
% two-index operator is used here throughout.
Two of the sources abbreviate the single factor \(\partial_{L}-n\) by
\(\Delta_{n}\).  In the notation of this volume that operator is
\(\shiftop{n}{1}\): the second index counts derivatives and is never optional,
since \(\shiftop{n}{k}\) is \emph{not} the \(k\)-th power of
\(\shiftop{n}{1}\).

\begin{theorem}[Repeated derivative formula]\label{plt:thm:dif-repeated}
For every \(k\in\NaturalNumbers\), every \(n\), and every coefficient \(p\),
\begin{equation}\label{plt:eq:dif-repeated}
  \DX^{k}\bigl(t^{n}p(L)\bigr)=t^{n+k}\,\shiftop{n}{k}p(L),
\end{equation}
and \(\shiftop{n}{k+1}=\shiftop{n+1}{k}\circ(\partial_{L}-n)
=(\partial_{L}-n-k)\circ\shiftop{n}{k}\).
\end{theorem}

\begin{proof}
Induction on \(k\).  For \(k=0\) both sides are \(t^{n}p\), the empty product
being the identity -- the convention that the sources state and then, in
\cref{plt:eq:dif-repeated}'s applications, occasionally forget.  Assume
\cref{plt:eq:dif-repeated} for \(k\).  Applying \cref{plt:eq:dif-block} at
outer exponent \(n+k\) to the coefficient \(\shiftop{n}{k}p\) gives
\[
  \DX^{k+1}\bigl(t^{n}p\bigr)
  =t^{n+k+1}\bigl(\partial_{L}-n-k\bigr)\shiftop{n}{k}p
  =t^{n+k+1}\shiftop{n}{k+1}p ,
\]
the last step because the factors of \cref{plt:eq:dif-shiftop} commute.
\end{proof}

\begin{corollary}[Two specialisations]\label{plt:cor:dif-stirling}
For \(k\ge0\):
\begin{equation}\label{plt:eq:dif-power-of-t}
  \DX^{k}\bigl(t^{n}\bigr)
  =(-1)^{k}\RisingFactorial{n}{k}\,t^{n+k},
\end{equation}
with \(\RisingFactorial{n}{k}=n(n+1)\cdots(n+k-1)\) and
\(\RisingFactorial{n}{0}=1\); and, for \(p\in\K[L]\),
\begin{equation}\label{plt:eq:dif-stirling}
  \EulerOp^{k}p(L)=t^{-k}\DX^{k}p(L)
  =\shiftop{0}{k}p
  =\sum_{j=0}^{k}\stirlingsigned{k}{j}\,\partial_{L}^{\,j}p ,
\end{equation}
where \(\stirlingsigned{k}{j}\) are the signed Stirling numbers of the first
kind.  More generally,
\[
  \shiftop{n}{k}
  =\sum_{j=0}^{k}\stirlingsigned{k}{j}\,(\partial_{L}-n)^{j}.
\]
\end{corollary}

\begin{proof}
For \cref{plt:eq:dif-power-of-t}, \(\shiftop{n}{k}1=\prod_{j=0}^{k-1}(-n-j)
=(-1)^{k}\prod_{j=0}^{k-1}(n+j)\).  For \cref{plt:eq:dif-stirling}, the
generating identity of the signed Stirling numbers of the first kind is
\(\FallingFactorial{z}{k}=\prod_{j=0}^{k-1}(z-j)=\sum_{j}
\stirlingsigned{k}{j}z^{j}\) \cite{comtet-book,graham-knuth-patashnik};
substituting the commuting operator \(z=\partial_{L}\), respectively
\(z=\partial_{L}-n\), gives both displays, and \(\EulerOp^{k}=t^{-k}\DX^{k}\)
holds on the block \(n=0\) by \cref{plt:eq:dif-euler-block}.
% ed.: S1 prints the identity as X^k d^k/dX^k P(log X) = sum_j s(k,j) P^{(j)}
% with the sum starting at j = 1; that is correct only for k >= 1, since
% s(0,0) = 1.  The range j = 0..k with the empty-product convention is uniform.
\end{proof}

\begin{corollary}[Derivative budget]\label{plt:cor:dif-budget}
\(\ordt(\DX^{k}F)\ge\ordt F+k\) for every \(F\in\PLlau\) and every \(k\ge0\).
Consequently, in a Taylor-type family \(\sum_{k\ge0}H^{k}\DX^{k}F/k!\) with
\(\ordt H\ge r\ge0\), the \(k\)-th term has outer order at least
\(\ordt F+(r+1)k\); the family is summable as soon as \(r\ge0\), and the
\(k\)-th term is invisible below \(t^{N}\) once \((r+1)k>N-1-\ordt F\).
\end{corollary}

\begin{proof}
Iterate \cref{plt:lem:dif-scale-basic}(iii), or read off
\cref{plt:eq:dif-repeated}.
\end{proof}

\section{The valuation of a derivative, exactly}
\label{plt:sec:dif-valuation}

The sources state \(\ordt(\DX F)\ge\ordt F+1\) and then describe the failure of
equality by three cases.  The three cases are correct, but they leave the
actual value of \(\ordt(\DX F)\) undetermined in the third.  It is determined,
and the sharp statement is no harder.

\begin{theorem}[Sharp valuation]\label{plt:thm:dif-valuation}
Let \(F\in\PLlau\) and write \(a_{0,0}=[t^{0}L^{0}]F\in\K\) for its constant
term.
\begin{enumerate}
\item \(\ordt(\DX F)\ge\ordt F+1\) always.
\item If \(F\notin\K\), then
\begin{equation}\label{plt:eq:dif-sharp-valuation}
  \ordt\bigl(\DX F\bigr)=\ordt\bigl(F-a_{0,0}\bigr)+1 .
\end{equation}
\item Equality \(\ordt(\DX F)=\ordt F+1\) holds if and only if \(F\) is not of
      the form \(a_{0,0}+\widetilde F\) with \(a_{0,0}\ne0\) and
      \(\ordt\widetilde F\ge1\).  Explicitly, with \(n_{0}=\ordt F\),
      \(p=p_{n_{0}}\), \(d=\degL p\) and \(a\) the scalar leading coefficient:
      \begin{itemize}
      \item if \(n_{0}\ne0\), the leading term of \(\DX F\) is
        \(-n_{0}a\,t^{n_{0}+1}L^{d}\) and equality holds;
      \item if \(n_{0}=0\) and \(d>0\), the leading term of \(\DX F\) is
        \(da\,tL^{d-1}\) and equality holds;
      \item if \(n_{0}=d=0\), the leading block is annihilated and
        \(\ordt(\DX F)=\ordt(F-a)+1\ge2\).
      \end{itemize}
\end{enumerate}
The same statements hold in \(\PLrat\) and \(\PLit\) with \(d\) replaced by
\(-\omega(p_{n_{0}})\), and in any power--log scale, except that in the third
case the bound reads \(\ordt(\DX F)>1\) when the exponent group is not
\(\IntegerNumbers\).
\end{theorem}

\begin{proof}
(i) is \cref{plt:lem:dif-scale-basic}(iii).  For (iii), apply
\cref{plt:eq:dif-block} to the leading block: the coefficient of
\(t^{n_{0}+1}\) in \(\DX F\) is \((\partial_{L}-n_{0})p\).  If \(n_{0}\ne0\)
this is a polynomial of degree \(d\) with leading coefficient \(-n_{0}a\ne0\),
so it is nonzero.  If \(n_{0}=0\) it is \(p'\), of degree \(d-1\) with leading
coefficient \(da\), nonzero exactly when \(d>0\).  If \(n_{0}=d=0\) then
\(p=a\) and \(p'=0\), so the block vanishes.

For (ii), let \(\widetilde F=F-a_{0,0}\), which is nonzero since
\(F\notin\K\), and note \(\DX F=\DX\widetilde F\).  It suffices to show that
\(\widetilde F\) falls under one of the first two cases of (iii).  Let
\(m=\ordt\widetilde F\).  If \(m\ne0\) we are in the first case.  If \(m=0\)
then \([t^{0}]\widetilde F=[t^{0}]F-a_{0,0}\) is a nonzero polynomial with
vanishing constant term, hence of degree \(\ge1\), and we are in the second
case.  In both cases \(\ordt(\DX\widetilde F)=m+1\).  The version for the
larger classes is the same computation with \(\omega\) in place of
\(-\degL\), using \cref{plt:lem:dif-scale-basic}(ii).
\end{proof}

\begin{corollary}[Kernel]\label{plt:cor:dif-kernel}
\(\ker\DX=\K\) in \(\PLpow\), \(\PLlau\), \(\PLrat\), \(\PLit\), and
\(\ker\DX=\K'\) in every power--log scale \(\mathcal H(C,\Gamma)\).
Consequently \(\ker\EulerOp=\K\) as well.
\end{corollary}

\begin{proof}
\Cref{plt:lem:dif-scale-basic}(iii); alternatively, if \(F\notin\K\) then
\cref{plt:eq:dif-sharp-valuation} gives \(\ordt(\DX F)<+\infty\), so
\(\DX F\ne0\).  For \(\EulerOp=t^{-1}\DX\), multiplication by the unit
\(t^{-1}\) does not change the kernel.
\end{proof}

\section{Closure of the ambient classes}
\label{plt:sec:dif-closure}

\begin{proposition}[Differential closure]\label{plt:prop:dif-closure-classes}
\(\DX\) and \(\EulerOp\) map each of \(\PLpow\), \(\PLlau\), \(\PLrat\),
\(\PLit\), \(\K[L]((t^{1/s}))\) and every power--log scale into itself.
Moreover \(\DX(\PLpow)= t\PLpow\) and \(\DX(t\PLpow)=t^{2}\PLpow\), so
\(\PLpow\) and its powers \(t^{k}\PLpow\) form a filtration by
\(\DX\)-stable subgroups, and each of \(\PLrat\), \(\PLit\) is a differential
field.
\end{proposition}

\begin{proof}
Closure under \(\DX\) follows from \cref{plt:prop:dif-block}, because the
operator \(\partial_{L}-n\) maps each of the three coefficient rings into
itself.  For \(\EulerOp\) multiply by \(t^{-1}\), a unit of every class listed
except \(\PLpow\).  The two inclusions
\[
  \DX(\PLpow)\subseteq t\PLpow,
  \qquad
  \DX(t\PLpow)\subseteq t^{2}\PLpow
\]
are immediate from \cref{plt:eq:dif-block}.  Conversely, run the construction
of \cref{plt:thm:dif-antiderivative} below on an \(F\) with \(\ordt F\ge1\):
it produces blocks only at indices \(n\ge0\), so the antiderivative lies in
\(\PLpow\).  On an \(F\) with \(\ordt F\ge2\) it meets only nonresonant
indices \(n\ge1\), so the antiderivative lies in \(t\PLpow\).
\end{proof}

\begin{remark}[Iterated logarithms are not polynomially closed]
\label{plt:rmk:dif-tower}
Set \(L_{1}=L\) and \(L_{j+1}=\log L_{j}\).  Then, by
\cref{plt:lem:dif-logderiv},
\begin{equation}\label{plt:eq:dif-tower}
  \DX L_{j}=\frac{\DX L_{j-1}}{L_{j-1}},
  \qquad\text{hence}\qquad
  \DX L_{j}=\frac{t}{L_{1}L_{2}\cdots L_{j-1}}
  \quad (j\ge1),
\end{equation}
by induction from \(\DX L_{1}=t\).  Consequently a polynomial ring
\(\K[L_{1},\dots,L_{r}]((t))\) with \(r\ge2\) is \emph{not} closed under
\(\DX\): already \(\DX L_{2}=t/L_{1}\) leaves it.  The smallest natural
differential extension replaces the coefficient ring by
\(\K(L_{1},\dots,L_{r})\), or by iterated Laurent coefficients, and the chain
rule then reads
\[
  \DX\bigl(a(L_{1},\dots,L_{r})t^{n}\bigr)
  =t^{n+1}\Bigl(\partial_{L_{1}}a
   +\tfrac{1}{L_{1}}\partial_{L_{2}}a+\cdots
   +\tfrac{1}{L_{1}\cdots L_{r-1}}\partial_{L_{r}}a-na\Bigr).
\]
This is the exact sense in which one logarithmic level is special: it is the
only depth at which polynomial coefficients are differentially closed.
\end{remark}

\begin{warningbox}
The formal partial derivative \(\partial_{t}\), which extracts coefficients, is
not the derivative with respect to \(X\).  A system that treats \(L\) as a
symbol independent of \(X\) will compute
\(\DX(p(L)t^{n})=-np(L)t^{n+1}\), losing the term \(t^{n+1}p'(L)\) -- exactly
the term that carries all logarithmic information.  Conversely, differentiating
after the substitution \(L=-\log t\) is a different operation again.  Print the
differentiated variable.

A separate trap is analytic.  Termwise differentiation of an \emph{asymptotic}
expansion is not automatic: a remainder that is \(\LittleOAt{X^{-N}}{X\to
+\infty}\) may have a derivative that is not \(\LittleOAt{X^{-N-1}}{X\to
+\infty}\).  Formal closure of \(\PLlau\) under \(\DX\) is an algebraic fact
and licenses no analytic conclusion without a realisation theorem --
analyticity in a sector, or monotonicity plus a derivative estimate
\cite{olver-asymptotics,hardy-orders}.
\end{warningbox}

\section{The block equation and the resonant block}
\label{plt:sec:dif-block-equation}

Antidifferentiation is, by \cref{plt:eq:dif-series-derivative}, a
block-by-block problem: \(\DX G=F\) with \(G=\sum q_{n}t^{n}\) is the system
\begin{equation}\label{plt:eq:dif-block-system}
  \bigl(\partial_{L}-n\bigr)q_{n}=p_{n+1},
  \qquad n\in\IntegerNumbers,
\end{equation}
one independent equation per block, with the outer index shifted down by one.
Everything therefore reduces to the solvability of one linear ordinary
equation in one polynomial or coefficient unknown.

\begin{lemma}[Nonresonant blocks are triangular]\label{plt:lem:dif-triangular}
Let \(a\in\K\setminus\SetOf{0}\) and \(q\in\K[L]\) of degree \(D\).  Then
\(\bigl(\partial_{L}-a\bigr)p=q\) has the unique polynomial solution
\begin{equation}\label{plt:eq:dif-nonresonant-inverse}
  p=-\sum_{j=0}^{D}\frac{q^{(j)}}{a^{j+1}},
\end{equation}
a finite sum, and \(\degL p=D\) with leading coefficient \(-a^{-1}\) times that
of \(q\).  The same formula solves the equation uniquely in \(\K((L^{-1}))\)
for any \(q\in\K((L^{-1}))\), the sum then being convergent in the
\(L^{-1}\)-adic order; hence \(\partial_{L}-a\) is bijective on
\(\K[L]\) and on \(\K((L^{-1}))\).
\end{lemma}

\begin{proof}
In the basis \(1,L,\dots,L^{D}\), the operator \(\partial_{L}-a\) is triangular
with diagonal \(-a\ne0\), hence bijective on polynomials of degree \(\le D\);
this proves existence, uniqueness and the degree statement.  For the closed
form, apply \(\partial_{L}-a\) to the right-hand side of
\cref{plt:eq:dif-nonresonant-inverse}:
\[
  \bigl(\partial_{L}-a\bigr)
  \Bigl(-\sum_{j\ge0}\frac{q^{(j)}}{a^{j+1}}\Bigr)
  =-\sum_{j\ge0}\frac{q^{(j+1)}}{a^{j+1}}+\sum_{j\ge0}\frac{q^{(j)}}{a^{j}}
  =q,
\]
the sums telescoping and \(q^{(j)}=0\) for \(j>D\).  In \(\K((L^{-1}))\) the
same computation is valid because \(\omega(q^{(j)})\ge\omega(q)+j\to+\infty\)
by \cref{plt:lem:dif-scale-basic}(ii), so the family is summable and may be
differentiated termwise; injectivity is
\cref{plt:lem:dif-scale-basic}(ii).
\end{proof}

\begin{lemma}[The resonant block]\label{plt:lem:dif-resonant}
On \(\K[L]\), \(\partial_{L}\) is surjective with kernel \(\K\): every
\(q\in\K[L]\) of degree \(D\) has the antiderivatives
\(p=c+\sum_{i=0}^{D}\frac{q_{i}}{i+1}L^{i+1}\), of degree \(D+1\), one for each
\(c\in\K\).  On \(\K((L^{-1}))\), \(\partial_{L}\) has kernel \(\K\) and image
\(\SetOf{h:\operatorname{res}(h)=0}\), a subspace of codimension one.
\end{lemma}

\begin{proof}
The polynomial statement uses \(\operatorname{char}\K=0\) to divide by
\(i+1\); that \(\degL p=D+1\) is clear from the leading term, and \(p'=0\)
forces \(p\in\K\).  The statement for \(\K((L^{-1}))\) is
\cref{plt:lem:dif-image-laurent}, and \(\operatorname{res}\) is onto \(\K\)
because \(\operatorname{res}(L^{-1})=1\).
\end{proof}

% ed.: the motivation part proves the same dichotomy in the weaker form it
% ed.: needs there (\cref{plt:lem:mot-block-antiderivative}, on \(\K[L]\)
% ed.: only, with no closed form).  Rather than restate it, the two lemmas
% ed.: above are identified with it here, so that a reader meets the resonance
% ed.: once as a mechanism and once in the form the antiderivative theorem
% ed.: consumes.
\begin{remark}[The same dichotomy as in the motivation part]
\label{plt:rmk:dif-resonance-restated}
\Cref{plt:lem:dif-triangular} and \cref{plt:lem:dif-resonant} are the
nonresonant and resonant halves of
\cref{plt:lem:mot-block-antiderivative}, sharpened in the two ways the
antiderivative theorem needs: the nonresonant inverse is given in closed form
by \cref{plt:eq:dif-nonresonant-inverse}, and both halves are proved on
\(\K((L^{-1}))\) as well as on \(\K[L]\), where the resonant image acquires
the codimension-one residue obstruction.  Nothing about the mechanism is new:
the unique resonant outer level is still \(n=0\) for \(\partial_{L}-n\),
equivalently the level \(n=1\) of the series being antidifferentiated, and it
is still the only level at which the logarithmic degree can rise.
\end{remark}

\section{Antiderivatives in the polynomial--logarithmic class}
\label{plt:sec:dif-antiderivative}

\begin{theorem}[\(\DX\) is onto \(\PLlau\)]\label{plt:thm:dif-antiderivative}
The sequence
\begin{equation}\label{plt:eq:dif-exact-sequence}
  0\longrightarrow\K\longrightarrow\PLlau
  \stackrel{\DX}{\longrightarrow}\PLlau\longrightarrow0
\end{equation}
is exact.  That is: every \(F\in\PLlau\) has an antiderivative \(G\in\PLlau\),
and \(G\) is unique up to an additive constant of \(\K\).  Explicitly, with
\(F=\sum_{n\ge n_{0}}p_{n}t^{n}\), the antiderivatives are
\begin{equation}\label{plt:eq:dif-antiderivative-formula}
  G=c+\int p_{1}\,\Differential L
   \;+\;\sum_{\substack{n\ge n_{0}-1\\ n\ne0}}
   t^{n}\Bigl(-\sum_{j\ge0}\frac{p_{n+1}^{(j)}}{n^{j+1}}\Bigr),
   \qquad c\in\K,
\end{equation}
where \(\int p_{1}\Differential L\) is the polynomial antiderivative with zero
constant term.  The log-degree profile of \(G\) is
\begin{equation}\label{plt:eq:dif-profile}
  \dprofile{G}(n)=\dprofile{F}(n+1)\ (n\ne0),
  \qquad
  \dprofile{G}(0)=\dprofile{F}(1)+1
\end{equation}
for the choice \(c=0\), with the convention \(-\infty+1=-\infty\).  The same
holds in every Puiseux-log ring and in every power--log scale with polynomial
coefficients.
\end{theorem}

\begin{proof}
Exactness at the left and middle is \cref{plt:cor:dif-kernel}.  For
surjectivity, solve \cref{plt:eq:dif-block-system} block by block:
for \(n\ne0\) use \cref{plt:lem:dif-triangular} with \(a=n\), obtaining the
unique \(q_{n}\in\K[L]\) displayed in
\cref{plt:eq:dif-antiderivative-formula}; for \(n=0\) use
\cref{plt:lem:dif-resonant}, obtaining \(q_{0}=c+\int p_{1}\Differential L\).
The resulting family \((q_{n}t^{n})\) has support bounded below by
\(n_{0}-1\), so \cref{plt:lem:dif-summable} assembles it into an element
\(G\in\PLlau\), and \cref{plt:eq:dif-series-derivative} gives \(\DX G=F\).
Uniqueness up to \(\K\) is again \cref{plt:cor:dif-kernel}.  The profile
statement is the degree bookkeeping of
\cref{plt:lem:dif-triangular,plt:lem:dif-resonant}: the nonresonant solution
has the same degree as its right-hand side, the resonant one has degree one
higher.  Nothing in the argument used \(\Gamma=\IntegerNumbers\); for a general
scale with polynomial coefficients the only resonant exponent is still
\(\gamma=0\).
% ed.: S6 concludes only that "the polynomial-logarithmic ring is closed under
% termwise antiderivatives, up to an arbitrary constant", and S3 writes the
% weaker (and redundant) inclusion int P dX subset P + K.  The correct
% statement is surjectivity of D_X on K[L]((t)) with kernel K, which is what
% the block system actually proves and what is needed downstream.
\end{proof}

\begin{example}[Both blocks]\label{plt:ex:dif-antiderivative}
Take \(F=L^{2}t^{3}+L^{2}t\).  The first term is nonresonant with \(n=2\):
\cref{plt:eq:dif-nonresonant-inverse} with \(a=2\), \(q=L^{2}\) gives
\(q_{2}=-\tfrac12L^{2}-\tfrac12L-\tfrac14\).  The second is the resonant block
\(n=0\): \(q_{0}=\tfrac13L^{3}+c\).  Hence
\[
  \int\Bigl(\frac{(\log X)^{2}}{X^{3}}+\frac{(\log X)^{2}}{X}\Bigr)
  \Differential X
  =-\frac{2L^{2}+2L+1}{4X^{2}}+\frac{L^{3}}{3}+c .
\]
Differentiation checks both blocks.  Note the degree growth in the resonant
block only, as \cref{plt:eq:dif-profile} predicts.
\end{example}

\begin{remark}[Why the classical obstruction is invisible here]
\label{plt:rmk:dif-invisible}
\Cref{plt:thm:dif-antiderivative} may look too strong, since
\(\int\Differential X/(X\log X)=\log\log X\) is the standard example of an
integral that leaves an elementary scale.  There is no conflict: the integrand
is \(t/L\), which is not in \(\PLlau\) at all.  The obstruction is invisible in
\(\PLlau\) precisely because a polynomial has no poles and no residue; it
becomes visible, and is then exactly one linear condition, as soon as the
coefficient ring is enlarged.  That is the content of the next subsection.
\end{remark}

\section{The exact obstruction in the larger coefficient classes}
\label{plt:sec:dif-obstruction}

\begin{theorem}[Antiderivatives in \(\PLit\)]\label{plt:thm:dif-anti-it}
Define \(\operatorname{res}\colon\PLit\to\K\) by
\begin{equation}\label{plt:eq:dif-res-functional}
  \operatorname{res}(F)\DefinitionEquals[t^{1}L^{-1}]F
  =\operatorname{res}\bigl([t^{1}]F\bigr),
\end{equation}
the coefficient of the single monomial \(tL^{-1}\).  Then the sequence
\begin{equation}\label{plt:eq:dif-exact-it}
  0\longrightarrow\K\longrightarrow\PLit
  \stackrel{\DX}{\longrightarrow}\PLit
  \stackrel{\operatorname{res}}{\longrightarrow}\K\longrightarrow0
\end{equation}
is exact.  That is: \(F\in\PLit\) has an antiderivative in \(\PLit\) if and
only if \(\operatorname{res}(F)=0\); the antiderivative is then unique up to
\(\K\).  In general, if \(\Lambda\) is a logarithm of \(L\) in an exponential
extension, then
\begin{equation}\label{plt:eq:dif-anti-with-Lambda}
  F-\operatorname{res}(F)\,\frac{t}{L}\in\DX\bigl(\PLit\bigr),
\end{equation}
so that every \(F\in\PLit\) has an antiderivative in
\(\PLit\oplus\K\Lambda\).
\end{theorem}

\begin{proof}
Exactness at \(\K\) and at the first \(\PLit\) is \cref{plt:cor:dif-kernel}.
The system \cref{plt:eq:dif-block-system} has, by
\cref{plt:lem:dif-triangular}, a unique solution \(q_{n}\in\K((L^{-1}))\) for
every \(n\ne0\), and by \cref{plt:lem:dif-resonant} a solution \(q_{0}\) if and
only if \(\operatorname{res}(p_{1})=0\), in which case it is unique up to
\(\K\).  Since \(\operatorname{res}(p_{1})=\operatorname{res}(F)\) by
definition, and since the assembled family has support bounded below by
\(\ordt F-1\) (\cref{plt:lem:dif-summable}), this proves exactness at the
second \(\PLit\).  Surjectivity of \(\operatorname{res}\) follows from
\(\operatorname{res}(t/L)=1\).  For
\cref{plt:eq:dif-anti-with-Lambda}, subtracting
\(\operatorname{res}(F)t/L\) makes the residue vanish, and
\(\DX\Lambda=\DX L/L=t/L\) by \cref{plt:lem:dif-logderiv}.
\end{proof}

\begin{corollary}[Where the second logarithm comes from]
\label{plt:cor:dif-loglog}
\(F=t/L=1/(X\log X)\) has \(\operatorname{res}(F)=1\), so it has no
antiderivative in \(\PLit\), hence none in \(\PLrat\) or \(\PLlau\); it has the
antiderivative \(\Lambda=\log\log X\).  More generally, by
\cref{plt:cor:dif-logderiv-residue}, for nonzero \(F\) in \(\PLit\) with
leading coefficient \(b\),
\begin{equation}\label{plt:eq:dif-res-log-derivative}
  \operatorname{res}\Bigl(\frac{\DX F}{F}\Bigr)=-\omega(b)
  =\degL b\ \text{ when }b\text{ is a polynomial},
\end{equation}
so the logarithmic derivative of \(F\) has an antiderivative in \(\PLit\) if
and only if \(\omega(b)=0\).  This is \cref{plt:thm:dif-log-boundary}(ii)
again: \(\log F=\int\DX F/F\), and the two obstructions are the same functional.
\end{corollary}

\begin{proof}
Combine \cref{plt:thm:dif-anti-it} with
\cref{plt:lem:dif-scale-basic}(iv) and
\cref{plt:cor:dif-logderiv-residue}: \([t^{1}](\DX F/F)=\partial_{L}b/b-e\),
whose residue is \(-\omega(b)\), the constant \(e\) contributing nothing.
\end{proof}

The rational-coefficient field behaves differently.  The sources record, all of
them correctly, that \(\PLrat\) is closed under differentiation, and they prove
closure under antiderivatives for the polynomial class only; none of them
remarks that the antiderivative statement fails after the enlargement to
\(\K(L)\) -- an inference their closure tables invite.  It does fail, and not
only at the resonant block: the obstruction is a linear condition at every pole
of every coefficient block.

\begin{lemma}[Twisted residues]\label{plt:lem:dif-image-rational}
Let \(a\in\K\) and \(q\in\K(L)\).  For \(\lambda\in\overline{\K}\) define the
\emph{twisted residue}
\begin{equation}\label{plt:eq:dif-twisted-residue}
  \rho_{a,\lambda}(q)\DefinitionEquals
  \operatorname{Res}_{L=\lambda}
  \Bigl(\EulerE^{-a(L-\lambda)}q\Bigr)
  =\sum_{k\ge0}\frac{(-a)^{k}}{k!}\,
   \operatorname{Res}_{L=\lambda}\bigl((L-\lambda)^{k}q\bigr),
\end{equation}
a finite sum, and note \(\rho_{0,\lambda}=\operatorname{Res}_{\lambda}\).
Then
\begin{equation}\label{plt:eq:dif-rational-image}
  \bigl(\partial_{L}-a\bigr)\K(L)
  =\SetOf{q\in\K(L):\rho_{a,\lambda}(q)=0
    \text{ for every }\lambda\in\overline{\K}} .
\end{equation}
In particular \(\K[L]\subseteq(\partial_{L}-a)\K(L)\), while for \(a\ne0\)
neither \(1/L\) nor \(1/L^{2}\) lies in the image.
\end{lemma}

\begin{proof}
Work first over \(\overline{\K}\) and fix \(\lambda\); write \(u=L-\lambda\)
and expand in \(\overline{\K}((u))\).  The multiplier
\(\EulerE^{-au}=\sum_{k}(-au)^{k}/k!\) is a unit of
\(\overline{\K}[[u]]\), the displayed sum is finite because \(u^{k}q\) is
regular for \(k\) at least the pole order, and
\[
  \EulerE^{-au}\bigl(\partial_{L}-a\bigr)p
  =\partial_{L}\bigl(\EulerE^{-au}p\bigr)
\]
for every \(p\in\overline{\K}((u))\).  Since the residue of a derivative of a
Laurent series vanishes, \(\rho_{a,\lambda}\) kills the image; this proves the
inclusion \(\subseteq\), the value at a point where \(q\) is regular being
\(0\).

For \(\supseteq\), decompose \(q=Q+\sum_{\lambda}q_{\lambda}\) into a
polynomial and its polar parts over \(\overline{\K}\); \(\rho_{a,\lambda}(q)
=\rho_{a,\lambda}(q_{\lambda})\) because the other summands are regular at
\(\lambda\).  The polynomial part is in the image by
\cref{plt:lem:dif-triangular} (for \(a\ne0\)) or
\cref{plt:lem:dif-resonant} (for \(a=0\)).  Fix \(\lambda\) and let \(M\) be
the pole order of \(q_{\lambda}\).  If \(p\) has a pole of order \(m\ge1\) at
\(\lambda\), then \((\partial_{L}-a)p\) has a pole of order exactly \(m+1\)
there, since the top coefficient is \(-mc_{m}\ne0\).  Hence the linear map
sending polar parts of order \(\le M-1\) to polar parts of order \(\le M\) is
injective, so its image has codimension one in the \(M\)-dimensional space of
polar parts of order \(\le M\); the functional \(\rho_{a,\lambda}\) vanishes on
that image and is not identically zero on the target (it takes the value
\(d_{1}\) on \(d_{1}u^{-1}\)).  Therefore the image is exactly
\(\ker\rho_{a,\lambda}\), and \(q_{\lambda}\) lies in it.  Summing the local
solutions and the polynomial solution gives \(p\in\overline{\K}(L)\) with
\((\partial_{L}-a)p=q\).

Finally, descend to \(\K\).  If \(a\ne0\), \(p\) is unique by
\cref{plt:lem:dif-triangular}, hence fixed by every automorphism of
\(\overline{\K}\) over \(\K\) (which fixes \(q\), \(a\) and \(L\)), so
\(p\in\K(L)\).  If \(a=0\), \(p\) is unique up to a constant; choosing a finite
Galois extension \(E/\K\) with \(p\in E(L)\) and replacing \(p\) by
\([E:\K]^{-1}\sum_{\sigma\in\operatorname{Gal}(E/\K)}\sigma(p)\), which is
again an antiderivative because each \(\sigma(p)\) is one, gives an element of
\(\K(L)\).  The examples: \(1/L\) has \(\rho_{a,0}=1\); \(1/L^{2}\) has
\(\rho_{a,0}=-a\).
\end{proof}

\begin{theorem}[Antiderivatives in \(\PLrat\)]\label{plt:thm:dif-anti-rat}
Let \(F=\sum_{n\ge n_{0}}p_{n}t^{n}\in\PLrat\).  Then \(F\) has an
antiderivative in \(\PLrat\) if and only if
\begin{equation}\label{plt:eq:dif-rat-criterion}
  \rho_{\,n-1,\;\lambda}\bigl(p_{n}\bigr)=0
  \qquad\text{for every }n\in\IntegerNumbers
  \text{ and every }\lambda\in\overline{\K},
\end{equation}
and the antiderivative is then unique up to \(\K\).  The condition at \(n=1\)
is the vanishing of all residues of \([t^{1}]F\); the conditions at
\(n\ne1\) are the twisted conditions with weight
\(\EulerE^{-(n-1)(L-\lambda)}\).  Consequently:
\begin{enumerate}
\item every \(F\in\PLlau\) satisfies \cref{plt:eq:dif-rat-criterion}
      (polynomials have no poles), recovering
      \cref{plt:thm:dif-antiderivative};
\item \(t^{2}/L=1/(X^{2}\log X)\) has no antiderivative in \(\PLrat\); it does
      have one in \(\PLit\), namely
      \(t\sum_{j\ge0}(-1)^{j+1}j!\,L^{-j-1}\);
\item \(\PLrat\) is a differential field that is not closed under
      antidifferentiation, whereas \(\PLlau\) is a differential ring, not a
      field, that is closed under antidifferentiation.
\end{enumerate}
\end{theorem}

\begin{proof}
The system \cref{plt:eq:dif-block-system} is solvable in \(\K(L)\) block by
block if and only if \(p_{n+1}\in(\partial_{L}-n)\K(L)\) for every \(n\), which
by \cref{plt:lem:dif-image-rational} is \cref{plt:eq:dif-rat-criterion} after
the index shift \(n\mapsto n-1\).  Solutions assemble by
\cref{plt:lem:dif-summable} because the support is bounded below, and the
ambiguity is \(\ker\DX=\K\) by \cref{plt:cor:dif-kernel}.  (i) is immediate.
For (ii), \(n=2\) and \(\lambda=0\) give \(\rho_{1,0}(1/L)=1\ne0\), so the
criterion fails; the element displayed is the unique solution of
\((\partial_{L}-1)q_{1}=L^{-1}\) provided by
\cref{plt:lem:dif-triangular} in \(\K((L^{-1}))\), namely
\(q_{1}=-\sum_{j\ge0}(L^{-1})^{(j)}=-\sum_{j\ge0}(-1)^{j}j!\,L^{-j-1}\).
(iii) collects \cref{plt:prop:dif-closure-classes},
\cref{plt:thm:dif-antiderivative} and part (ii); that \(\PLlau\) is not a
field is \cref{plt:cor:div-domain-not-field}.
\end{proof}

\begin{remark}[A divergent formal antiderivative]\label{plt:rmk:dif-borel}
The series \(t\sum_{j\ge0}(-1)^{j+1}j!L^{-j-1}\) of
\cref{plt:thm:dif-anti-rat}(ii) is a perfectly good element of \(\PLit\): its
defining sum is \(L^{-1}\)-adically summable, and its blocks are determined by
finitely many operations each.  It is not, and does not claim to be, a
convergent series; the factorial growth of its coefficients is the algebraic
shadow of the fact that \(\int\Differential X/(X^{2}\log X)\) is not
elementary, and the analytic content of the expansion is a matter of
summation theory rather than of the formal algebra
\cite{ecalle,costin-book}.  Nothing in
\cref{plt:thm:dif-anti-it,plt:thm:dif-anti-rat} asserts convergence; the
statements are exact statements about coefficients.
\end{remark}

\begin{table}[htbp]
\centering
\caption{The differential calculus of the four ambient classes.  Here
\(\operatorname{res}F=[t^{1}L^{-1}]F\) and \(\rho_{a,\lambda}\) are the twisted
residues defined in this section.}
\label{plt:tab:dif-classes}
\small
\begin{tabular}{@{}p{0.20\textwidth}p{0.16\textwidth}p{0.18\textwidth}p{0.36\textwidth}@{}}
\toprule
Class & \(\DX\)-closed & \(\ker\DX\) & antiderivative exists in the class \\
\midrule
\(\PLpow=\K[L][[t]]\) & yes, into \(t\PLpow\) & \(\K\) & iff \(\ordt F\ge1\)
  (\cref{plt:prop:dif-closure-classes}) \\
\(\PLlau=\K[L]((t))\) & yes & \(\K\) & always
  (\cref{plt:thm:dif-antiderivative}) \\
\(\PLrat=\K(L)((t))\) & yes & \(\K\) & iff all \(\rho_{n-1,\lambda}([t^{n}]F)
  =0\) (\cref{plt:thm:dif-anti-rat}) \\
\(\PLit=\K((L^{-1}))((t))\) & yes & \(\K\) & iff \(\operatorname{res}F=0\)
  (\cref{plt:thm:dif-anti-it}) \\
\bottomrule
\end{tabular}
\end{table}

\begin{checkpointbox}
The whole of \cref{plt:sec:differential} turns on one line: \(\DX\) sends the
outer block \(n\) to the block \(n+1\) by the operator \(\partial_{L}-n\).
Everything else is a statement about that operator.  It is injective for
\(n\ne0\) and has kernel \(\K\) for \(n=0\); it is surjective on \(\K[L]\) for
every \(n\), surjective on \(\K((L^{-1}))\) for \(n\ne0\) and of corank one for
\(n=0\), and of infinite corank on \(\K(L)\) for every \(n\) as soon as poles
are admitted.  Reading the four ambient classes off from that single operator
is the fastest route through the material.
\end{checkpointbox}

\begin{summarybox}
The distinguished derivation \(\DX=t\partial_{L}-t^{2}\partial_{t}\) is the
unique bounded \(\K\)-derivation with \(\DX t=-t^{2}\) and \(\DX L=t\); on one
block it is \(\DX(t^{n}p)=t^{n+1}(\partial_{L}-n)p\), and \(k\) derivatives
give \(t^{n+k}\shiftop{n}{k}p\) with the empty-product convention.  The Euler
derivation \(\EulerOp=t^{-1}\DX\) preserves each block and acts there as
\(\partial_{L}-n\), whose invertibility fails exactly at the resonant block
\(n=0\).  Valuation rises by exactly one except when the leading block is a
nonzero constant, and then \(\ordt(\DX F)=\ordt(F-a_{0,0})+1\).  All four
ambient classes are closed under \(\DX\).  For antiderivatives the answer is a
residue: \(\DX\) is onto \(\PLlau\) with kernel \(\K\); on \(\PLit\) the single
functional \(\operatorname{res}F=[t^{1}L^{-1}]F\) is the entire obstruction, and
it is repaired by adjoining \(\Lambda=\log L\); on \(\PLrat\) there is one
twisted-residue condition at every pole of every block, so exact rational
coefficients are the one enlargement that costs closure under integration.
\end{summarybox}

\clearpage
\part{Composition}

% ---- fragment 06-composition ----
\chapter{Admissible composition and substitution}
\label{plt:sec:composition}

Throughout this part \(\K\) is a field of characteristic zero, \(X\) is the
large variable, and
\[
 t\DefinitionEquals X^{-1},
 \qquad
 L\DefinitionEquals\log X
\]
are the two generators.  Unless a sector is named, the analytic regime is the
positive real ray \(X\to+\infty\) and every logarithm is the real logarithm
there.  We write \(\PLpow\) for \(\K[L][[t]]\), \(\PLlau\) for \(\K[L]((t))\),
\(\ordt\) for the outer valuation, \(\Trunc{F}{N}\) for the \emph{exclusive}
truncation \(\sum_{n<N}p_n(L)t^n\), and \(\DX=t\partial_L-t^2\partial_t\) for
the distinguished derivation.  Recall the repeated-derivative operator
\begin{equation}\label{plt:eq:cmp-shiftop}
 \shiftop{n}{k}
 \DefinitionEquals\prod_{j=0}^{k-1}(\partial_L-n-j),
 \qquad
 \shiftop{n}{0}=1,
 \qquad
 \DX^k\bigl(t^np(L)\bigr)=t^{n+k}\shiftop{n}{k}p(L).
\end{equation}
The empty product at \(k=0\) is the identity operator; this convention is used
without further comment.

Composition is the first operation for which the analytic relation between the
two generators is essential.  Addition and multiplication hold \(L\) fixed and
merely recombine supports that are already present.  Substitution of
\(X\mapsto G(X)\) moves both generators at once, and it can move them out of
the ring: fractional powers, a second logarithmic level, or exponential
monomials may appear.  The whole content of this chapter is a set of
hypotheses under which the substitution is defined, closed, associative, and
computable block by block, together with a careful separation of that formal
statement from the analytic statement that it computes the asymptotics of an
actual composite function.

\section{What substitution has to mean}
\label{plt:sec:cmp-what-substitution-means}

For a polynomial in \(t\) and \(L\), substitution is finite algebra and needs
no theory.  For an infinite series
\(F=\sum_{n\ge n_0}p_n(L)t^n\) one wants to declare
\[
 F\compose G
 \DefinitionEquals
 \sum_{n\ge n_0}p_n(L\compose G)\,(t\compose G)^n,
\]
and three separate things must be checked.

\begin{enumerate}
\item \emph{The generators must land in a named ring.}  Both \(t\compose G\)
and \(L\compose G\) must be expansible in the chosen monomial scale.
\item \emph{The outer sum must be summable.}  For each output monomial, only
finitely many indices \(n\) may contribute; otherwise a coefficient is an
undefined infinite sum in \(\K\).
\item \emph{The result must have admissible support.}  In a Laurent setting
this means finitely many negative outer exponents; in a Hahn setting it means
a well-ordered exponent set.
\end{enumerate}

Each requirement fails for some perfectly meaningful \(G\).

\begin{example}[Four failures of naive substitution]
\label{plt:ex:cmp-naive-failures}
\begin{enumerate}
\item \(F=t\), \(G=X^{1/2}\): here \(t\compose G=t^{1/2}\notin\PLlau\), and a
Puiseux extension is forced.
\item \(F=t\), \(G=\log X\): here \(t\compose G=L^{-1}\notin\PLlau\), and
negative powers of \(L\) are forced, that is, a Laurent-log coefficient field.
\item \(F=L\), \(G=X\log X\): here \(L\compose G=L+\log L\), and a second
logarithmic level is forced.
\item \(F=L\), \(G=\EulerE^{X}\): here \(t\compose G=\EulerE^{-X}\), smaller
than every power of \(t\), and an exponential monomial is forced.
\end{enumerate}
Failure of summability is separate and more insidious.  Let
\(F=\sum_{n\ge0}t^n\) and let \(G=c\) be a nonzero constant, so that
\(t\compose G=c^{-1}\) has outer valuation \(0\).  Every summand
\(c^{-n}\) then lands in the block \(t^0\), and the coefficient of \(t^0\) in
the putative composite is \(\sum_{n\ge0}c^{-n}\), which is not an element of
\(\K\).  The same collapse occurs for any \(G\) whose leading exponent is
\(0\); this is why the hypothesis \(\rho>0\) below is not decorative.
\end{example}

\begin{remark}
Nothing in \cref{plt:ex:cmp-naive-failures} says that the composite function
does not exist.  Closure is a statement about a chosen representation class,
not about existence.  When a substitution leaves the class, the correct
response is to name the enlargement it demands, not to suppress the offending
symbol.
\end{remark}

We record the summability criterion once, in the generality needed later.
Let \(\Gamma\subseteq\RealNumbers\) be an ordered additive subgroup and let
\[
 \K[L][[t^\Gamma]]
 \DefinitionEquals
 \SetOf{\textstyle\sum_{\alpha\in S}p_\alpha(L)t^\alpha
   :\ p_\alpha\in\K[L],\ S\subseteq\Gamma\text{ well ordered}}
\]
be the associated polynomial-logarithmic Hahn ring.  For \(F\ne0\) write
\(E(F)\DefinitionEquals\SetOf{\alpha\in\Gamma:[t^\alpha]F\ne0}\) for its
exponent support, so that \(\ordt(F)=\min E(F)\) and
\(\CoefficientSupportOf{F}\) is the finer set of pairs \((\alpha,j)\) with
\(a_{\alpha,j}\ne0\).  For \(\Gamma=\IntegerNumbers\) this ring is exactly
\(\PLlau\), and for \(\Gamma=s^{-1}\IntegerNumbers\) it is the Puiseux-log ring
\(\K[L]((t^{1/s}))\).

% ed.: retitled.  This is the summability criterion in a Hahn ring with real
% ed.: exponents, for an arbitrary index set, together with the notion of a
% ed.: map compatible with summable families that the rest of the part uses.
% ed.: It is not \cref{plt:lem:dif-summable} (sequences in a Laurent ring)
% ed.: and not \cref{plt:thm:trunc-summability} (the calculus on \(\PLlau\));
% ed.: for \(\Gamma=\IntegerNumbers\) it restricts to the latter.
\begin{lemma}[Summable families in a polynomial--logarithmic Hahn ring]
\label{plt:lem:cmp-summable}
Let \(\mathcal A=\K[L][[t^\Gamma]]\) and let \((S_i)_{i\in I}\) be a family in
\(\mathcal A\) such that
\begin{equation}\label{plt:eq:cmp-summability-criterion}
 \SetOf{i\in I:\ \ordt(S_i)\le\beta}
 \quad\text{is finite for every }\beta\in\Gamma .
\end{equation}
Then \(\bigcup_iE(S_i)\) is a well-ordered
subset of \(\Gamma\), every exponent occurs in only finitely many \(S_i\), and
\[
 \sum_{i\in I}S_i
 \DefinitionEquals
 \sum_{\alpha}\Bigl(\sum_{i}[t^\alpha]S_i\Bigr)t^\alpha
\]
is a well-defined element of \(\mathcal A\).  The sum is unchanged by any
regrouping or reindexing of \(I\); in particular a doubly indexed summable
family may be summed in either order.  Finally, call a \(\K\)-linear map
\(\Psi:\mathcal A\to\mathcal B\) between two such rings \emph{compatible with
summable families} when there is a non-decreasing unbounded
\(\varphi\) with \(\ordt(\Psi a)\ge\varphi(\ordt a)\) for all \(a\) and
\(\Psi\bigl(\sum_iS_i\bigr)=\sum_i\Psi S_i\) for every summable family; the
first condition alone already forces \((\Psi S_i)\) to satisfy
\eqref{plt:eq:cmp-summability-criterion}, and a composite of two compatible
maps is compatible.
\end{lemma}

\begin{proof}
Let \(T\subseteq\bigcup_iE(S_i)\) be nonempty and pick \(\beta\in T\).  By
\eqref{plt:eq:cmp-summability-criterion} the index set
\(I_\beta=\SetOf{i:\ordt(S_i)\le\beta}\) is finite, and every element of
\(T\) that is \(\le\beta\) lies in \(\bigcup_{i\in I_\beta}E(S_i)\), a finite
union of well-ordered sets and hence well ordered.  Therefore
\(T\cap(-\infty,\beta]\) has a least element, which is the least element of
\(T\).  So the union is well ordered.  Taking \(\beta=\alpha\) shows that only
the finitely many \(i\in I_\alpha\) can contribute to the exponent
\(\alpha\), so each inner sum is finite and the displayed series has
well-ordered support.  Regrouping and reindexing change none of the finitely
many terms contributing to a fixed \(\alpha\).  For the last assertion, let
\(\varphi\) be non-decreasing and unbounded with
\(\ordt(\Psi a)\ge\varphi(\ordt a)\), and fix \(\beta\).  Choose
\(\beta_0\in\Gamma\) with \(\varphi(\gamma)>\beta\) for all
\(\gamma>\beta_0\).  Then
\(\SetOf{i:\ordt(\Psi S_i)\le\beta}\subseteq\SetOf{i:\ordt(S_i)\le\beta_0}\),
which is finite.
\end{proof}

% ed.: the catalogue warns that a t-adically Cauchy *sequence* in a Laurent
% ring may fail to converge without a common lower support bound.  For
% *families* satisfying (plt:eq:cmp-summability-criterion) the bound is
% automatic, as the proof shows; the two statements are not in conflict.
\begin{remark}\label{plt:rmk:cmp-lower-bound-automatic}
Condition \eqref{plt:eq:cmp-summability-criterion} at \(\beta=0\) already
forces a common lower support bound: only finitely many members have
valuation \(\le0\), and each of those has finite valuation.  A \(t\)-adically Cauchy
\emph{sequence} in a Laurent ring carries no such guarantee, which is why the
family formulation is the one used throughout.
\end{remark}

\section{Formal exponential and logarithm of a one-unit}
\label{plt:sec:cmp-exp-log}

Write \(\mathfrak m\DefinitionEquals t\PLpow\) for the maximal ideal of
\(\PLpow\); in a Hahn ring \(\mathfrak m\) denotes the ideal of elements of
strictly positive valuation.  For \(Z\in\mathfrak m\) put
\begin{equation}\label{plt:eq:cmp-exp-log-def}
 \EulerE^{Z}\DefinitionEquals\sum_{k\ge0}\frac{Z^k}{k!},
 \qquad
 \log(1+Z)\DefinitionEquals\sum_{r\ge1}\frac{(-1)^{r-1}}{r}Z^r .
\end{equation}
Both families satisfy \eqref{plt:eq:cmp-summability-criterion} because
\(\ordt(Z^k)\ge k\ordt(Z)\to+\infty\).

\begin{lemma}[Formal exponential and logarithm]\label{plt:lem:cmp-exp-log}
The maps \eqref{plt:eq:cmp-exp-log-def} are mutually inverse bijections
between \(\mathfrak m\) and \(1+\mathfrak m\), and
\[
 \EulerE^{Z_1+Z_2}=\EulerE^{Z_1}\EulerE^{Z_2},
 \qquad
 \log(P_1P_2)=\log P_1+\log P_2,
 \qquad
 \log(P^{\,n})=n\log P
\]
for \(Z_1,Z_2\in\mathfrak m\), \(P,P_1,P_2\in1+\mathfrak m\) and
\(n\in\IntegerNumbers\).  Moreover, if \(\Psi\) is a \(\K\)-algebra
homomorphism into another such ring, compatible with summable families and
carrying elements of positive valuation to elements of positive valuation,
then
\(\Psi(\EulerE^{Z})=\EulerE^{\Psi Z}\) and
\(\Psi(\log(1+Z))=\log(1+\Psi Z)\).
\end{lemma}

\begin{proof}
In \(\RationalNumbers[[z_1,z_2]]\) one has the classical identities
\(\exp(z_1+z_2)=\exp z_1\exp z_2\),
\(\log((1+z_1)(1+z_2))=\log(1+z_1)+\log(1+z_2)\),
\(\log\exp z_1=z_1\) and \(\exp\log(1+z_1)=1+z_1\).  Because
\(\ordt(Z_i)>0\), the monomial family
\(\bigl(a_{k_1k_2}Z_1^{k_1}Z_2^{k_2}\bigr)_{k_1,k_2}\) satisfies
\eqref{plt:eq:cmp-summability-criterion} for every choice of coefficients
\(a_{k_1k_2}\in\RationalNumbers\), so \(z_i\mapsto Z_i\) extends to a
\(\RationalNumbers\)-algebra homomorphism
\(\RationalNumbers[[z_1,z_2]]\to\mathcal A\) that is compatible with summable
families (\cref{plt:lem:cmp-summable}); applying it to the four identities gives the
first three assertions for \(n\ge0\), and \(\log(P^{-1})=-\log P\) follows
from \(\log(PP^{-1})=\log1=0\).  The last sentence is termwise application of
\(\Psi\) to \eqref{plt:eq:cmp-exp-log-def}.
\end{proof}

\section{Monomial-leading inner arguments and the substitution homomorphism}
\label{plt:sec:cmp-monomial-leading}

\begin{definition}[Admissible monomial-leading argument]
\label{plt:def:cmp-admissible-inner}
An \emph{admissible monomial-leading inner argument} over \(\K\) is a
quadruple \((c,\log c,\rho,U)\) consisting of
\[
 c\in\K^\times,
 \qquad
 \text{a chosen value }\log c\in\K,
 \qquad
 \rho\in\PositiveIntegers,
 \qquad
 U\in t\PLpow ,
\]
presented as the inner argument
\begin{equation}\label{plt:eq:cmp-G-monomial}
 G(X)=cX^{\rho}\bigl(1+U(X)\bigr)
 = c\,t^{-\rho}(1+U)\in\PLlau .
\end{equation}
We put
\begin{equation}\label{plt:eq:cmp-V-def}
 V\DefinitionEquals\log(1+U)\in t\PLpow .
\end{equation}
On the real ray we additionally require \(c>0\) and take \(\log c\) real, so
that \(G(X)\to+\infty\).
\end{definition}

The datum \(\log c\) is part of the inner argument, not a consequence of it.
If \(\K\) does not already contain a logarithm of \(c\), it must first be
enlarged; and if \(\K\subseteq\ComplexNumbers\) contains \(2\pi\ImaginaryUnit\),
then \(c\) admits several admissible logarithms, which give genuinely
different substitutions (\cref{plt:rmk:cmp-branch-matters}).

% ed.: retitled to separate it from \cref{plt:prop:dif-substitution}, which
% ed.: substitutes a small series into a scalar power series.  The present
% ed.: theorem substitutes a monomial-leading transseries into a transseries.
\begin{theorem}[Substitution homomorphism of a monomial-leading argument]
\label{plt:thm:cmp-substitution-hom}
Let \(G\) be as in \cref{plt:def:cmp-admissible-inner}.  For
\(F=\sum_{n\ge n_0}p_n(L)t^n\in\PLlau\) the family
\begin{equation}\label{plt:eq:cmp-master}
 \bigl(\,p_n(\rho L+\log c+V)\,c^{-n}t^{\rho n}(1+U)^{-n}\,\bigr)_{n\ge n_0}
\end{equation}
is summable in \(\PLlau\), so that
\begin{equation}\label{plt:eq:cmp-master-sum}
 F\compose G
 \DefinitionEquals
 \sum_{n\ge n_0}
 p_n\!\bigl(\rho L+\log c+V\bigr)\,
 c^{-n}\,t^{\rho n}\,(1+U)^{-n}
\end{equation}
is a well-defined element of \(\PLlau\).  The resulting map
\(\Phi_G:F\mapsto F\compose G\) is
\begin{enumerate}
\item a \(\K\)-algebra homomorphism \(\PLlau\to\PLlau\), compatible with
summable families;
\item determined by its values on the generators,
\begin{equation}\label{plt:eq:cmp-generator-images}
 \Phi_G(t)=c^{-1}t^{\rho}(1+U)^{-1},
 \qquad
 \Phi_G(L)=\rho L+\log c+V ;
\end{equation}
\item exactly valuation-multiplying:
\(\ordt(F\compose G)=\rho\,\ordt(F)\) for all \(F\).  In particular
\(\Phi_G\) is injective, and \(\Phi_G(\PLpow)\subseteq\PLpow\).
\end{enumerate}
\end{theorem}

\begin{proof}
Since \(\ordt(U)\ge1\), \cref{plt:lem:cmp-exp-log} gives \(V\in t\PLpow\), and
the binomial family \((\binom{-n}{r}U^r)_{r\ge0}\) is summable for every
\(n\in\IntegerNumbers\), so \((1+U)^{-n}\in1+t\PLpow\).  Because \(p_n\) is a
polynomial of some finite degree \(d_n\),
\begin{equation}\label{plt:eq:cmp-poly-shift}
 p_n(\rho L+\log c+V)
 =\sum_{j=0}^{d_n}\frac{p_n^{(j)}(\rho L+\log c)}{j!}\,V^{j}
\end{equation}
is a \emph{finite} sum, each summand lying in \(\PLpow\); here \(p_n^{(j)}\)
is the \(j\)-th derivative of \(p_n\) in its own variable and the polynomial
is evaluated at \(\rho L+\log c\).  Hence the \(n\)-th member of
\eqref{plt:eq:cmp-master} lies in \(t^{\rho n}\PLpow\), and its valuation is
at least \(\rho n\).  Since \(\rho\ge1\) and \(n\ge n_0\), condition
\eqref{plt:eq:cmp-summability-criterion} holds, and
\cref{plt:lem:cmp-summable} gives the sum
\eqref{plt:eq:cmp-master-sum}.

For the algebra structure, note that \eqref{plt:eq:cmp-master-sum} is exactly
"substitute \eqref{plt:eq:cmp-generator-images} into the expansion of \(F\)".
Concretely, let \(\Psi_0:\K[L][t,t^{-1}]\to\PLlau\) be the unique
\(\K\)-algebra homomorphism with \(\Psi_0(L)=\rho L+\log c+V\) and
\(\Psi_0(t)=c^{-1}t^{\rho}(1+U)^{-1}\); it exists because \(L\) and \(t\) are
algebraically independent over \(\K\) and \(\Psi_0(t)\) is a unit of
\(\PLlau\), the required inverse being \(c\,t^{-\rho}(1+U)\).  Then
\(\Phi_G(p_n(L)t^n)=\Psi_0(p_n(L)t^n)\) for every \(n\) and every \(p_n\), and
\(\Phi_G(F)=\sum_n\Psi_0(p_n(L)t^n)\).  Additivity of \(\Phi_G\) is clear.
For multiplicativity, let \(F=\sum_np_nt^n\) and \(F'=\sum_mq_mt^m\).  The
double family \(\bigl(\Psi_0(p_nt^n)\Psi_0(q_mt^m)\bigr)_{n,m}\) satisfies
\eqref{plt:eq:cmp-summability-criterion}, because its \((n,m)\) member has
valuation at least \(\rho(n+m)\); summing it in either order
(\cref{plt:lem:cmp-summable}) gives
\[
 \Phi_G(F)\Phi_G(F')
 =\sum_{n,m}\Psi_0\bigl(p_nq_mt^{n+m}\bigr)
 =\sum_{r}\Psi_0\Bigl(\sum_{n+m=r}p_nq_m\,t^{r}\Bigr)
 =\Phi_G(FF').
\]

% ed.: the draft deduced "compatible with summable families" from
% (plt:lem:cmp-summable) directly, but that lemma only supplies the valuation
% half of the definition; the interchange Phi_G(sum S_i) = sum Phi_G(S_i) needs
% the argument below.  It is used by (plt:lem:tay-determined) and by the
% associativity proof, so it is supplied here.
It remains to check compatibility with summable families in the sense of
\cref{plt:lem:cmp-summable}.  The valuation half holds with
\(\varphi(x)=\rho x\), which is non-decreasing and unbounded.  For the
interchange, let \((S_i)_{i\in I}\) be summable with sum \(S\) and write
\(S_i=\sum_np_{i,n}(L)t^n\).  By
\cref{plt:rmk:cmp-lower-bound-automatic} there is \(n_1\) with
\(\ordt(S_i)\ge n_1\) for all \(i\), so \(p_{i,n}=0\) for \(n<n_1\), and for
each fixed \(n\) only the finitely many \(i\) with \(\ordt(S_i)\le n\) can
have \(p_{i,n}\ne0\); consequently \([t^n]S=\sum_ip_{i,n}\) is a finite sum.
The doubly indexed family \(\bigl(\Psi_0(p_{i,n}t^n)\bigr)_{i,n}\) has
\((i,n)\) member of valuation at least \(\rho n\ge\rho n_1\) and therefore
satisfies \eqref{plt:eq:cmp-summability-criterion}.  Summing it with \(i\)
innermost gives \(\sum_n\Psi_0\bigl([t^n]S\cdot t^n\bigr)=\Phi_G(S)\), and
summing it with \(n\) innermost gives \(\sum_i\Phi_G(S_i)\); the two agree by
\cref{plt:lem:cmp-summable}.
This proves (i) and (ii).

For (iii), let \(n_0=\ordt(F)\).  Every summand with \(n>n_0\) has valuation
at least \(\rho n>\rho n_0\), and the summand with \(n=n_0\) equals
\(c^{-n_0}t^{\rho n_0}\bigl(p_{n_0}(\rho L+\log c)+W\bigr)\) with
\(\ordt(W)\ge1\), by \eqref{plt:eq:cmp-poly-shift} and
\((1+U)^{-n_0}\in1+t\PLpow\).  Since \(\rho\ne0\), the substitution
\(L\mapsto\rho L+\log c\) is a \(\K\)-algebra automorphism of \(\K[L]\), so
\(p_{n_0}(\rho L+\log c)\ne0\).  Hence
\([t^{\rho n_0}](F\compose G)=c^{-n_0}p_{n_0}(\rho L+\log c)\ne0\) and
\(\ordt(F\compose G)=\rho n_0\).  Injectivity and preservation of
\(\PLpow\) follow at once.
\end{proof}

\begin{corollary}[Scaling and pure power substitution]
\label{plt:cor:cmp-pure-power}
With \(U=0\),
\[
 F(cX^{\rho})
 =\sum_{n\ge n_0}c^{-n}t^{\rho n}\,p_n(\rho L+\log c).
\]
Thus a pure scaling acts on coefficient blocks by the affine change of
logarithmic variable \(L\mapsto\rho L+\log c\) and on outer exponents by
\(n\mapsto\rho n\); no new monomials are created.
\end{corollary}

\begin{proof}
Immediate from \eqref{plt:eq:cmp-master-sum} with \(U=V=0\).
\end{proof}

\begin{remark}[The choice of logarithm is part of the data]
\label{plt:rmk:cmp-branch-matters}
Take \(\K=\ComplexNumbers\), \(c=1\), \(\rho=1\), \(U=0\), so that
\(G=X\) is the identity map.  Choosing \(\log c=0\) gives
\(\Phi_G=\mathrm{id}\).  Choosing \(\log c=2\pi\ImaginaryUnit\) gives the
\emph{different} \(\K\)-algebra automorphism
\(t\mapsto t\), \(L\mapsto L+2\pi\ImaginaryUnit\), which is
\eqref{plt:eq:cmp-master-sum} for the very same \(G\).  Therefore
"\(F\compose G\)" is shorthand for \(\Phi_{(c,\log c,\rho,U)}(F)\), and a
symbolic implementation must carry \(\log c\) as data rather than replace it
by a principal value.  All statements below that compose two inner arguments
carry an explicit compatibility hypothesis on the chosen logarithms.
\end{remark}

\section{Finite determination}
\label{plt:sec:cmp-finite-determination}

The practical content of \cref{plt:thm:cmp-substitution-hom} is that each
output block is a finite algebraic expression in finitely many input blocks,
with an explicit and \emph{triangular} precision budget.

% ed.: retitled: \cref{plt:thm:dif-finite-determination} carries the same name
% ed.: for a different statement, the finite determination of \(\Phi(U)\) with
% ed.: \(\Phi\) a scalar power series.
\begin{theorem}[Finite determination of a composition]
\label{plt:thm:cmp-finite-determination}
Let \(G\) be as in \cref{plt:def:cmp-admissible-inner}, let
\(F=\sum_{n\ge a}p_n(L)t^n\in\PLlau\) with \(a=\ordt(F)\), and let
\(N\in\IntegerNumbers\).  Then
\begin{equation}\label{plt:eq:cmp-block-finite}
 [t^{N}](F\compose G)
 =\sum_{\substack{a\le n\\ \rho n\le N}}
   c^{-n}\!\!\!
   \sum_{\substack{j\ge0,\ r\ge0\\ j\le\deg p_n}}\!\!\!
   \frac{\binom{-n}{r}}{j!}\,
   [t^{\,N-\rho n}]\bigl(V^{j}U^{r}\bigr)\;
   p_n^{(j)}(\rho L+\log c),
\end{equation}
and in the inner sum only the pairs with \(j+r\le N-\rho n\) contribute.  Both
index ranges are finite, so every output block below any cutoff is a finite
sum.  Consequently
\(\Trunc{F\compose G}{N}\) is a finite algebraic expression in
\(c^{\pm1}\), \(\log c\), \(\rho\), the blocks \(p_n\) with \(a\le n<N/\rho\),
and the blocks of \(U\) below outer order \(N-\rho a\).
\end{theorem}

\begin{proof}
Insert \eqref{plt:eq:cmp-poly-shift} and the binomial expansion
\((1+U)^{-n}=\sum_{r\ge0}\binom{-n}{r}U^{r}\) into
\eqref{plt:eq:cmp-master-sum} and extract the coefficient of \(t^{N}\); the
triple family is summable because its \((n,j,r)\) member has valuation at
least \(\rho n+j+r\), so \cref{plt:lem:cmp-summable} licenses the
rearrangement.  The outer index is restricted by \(a\le n\) and
\(\rho n\le N\), that is \(n\le\Floor{N/\rho}\), and the inner indices by
\(j+r\le N-\rho n\), because \(\ordt(V^jU^r)\ge j+r\); both are finite ranges.
Finally, \(\Trunc{F\compose G}{N}\) collects
\eqref{plt:eq:cmp-block-finite} for all \(N'<N\); the constraint
\(\rho n<N\) is \(n<N/\rho\), and every occurrence of \(U\) or \(V\) sits
inside a factor \(t^{\rho n}\) with \(\rho n\ge\rho a\), so only the blocks of
\(U\) below outer order \(N-\rho a\) can be reached.
\end{proof}

The last clause is worth isolating, because it is the difference between a
usable and an unusable implementation: the required precision of the inner
data is not uniform across outer blocks.

\begin{corollary}[Triangular precision budget]
\label{plt:cor:cmp-precision-budget}
In \eqref{plt:eq:cmp-master-sum} the \(n\)-th summand needs \(V\) and
\((1+U)^{-n}\) only through outer order \(N-\rho n-1\).  Computing all of them
to the single global precision \(N-\rho a-1\) is correct but wasteful.
\end{corollary}

\begin{proof}
The \(n\)-th summand carries the prefactor \(t^{\rho n}\), so anything of
outer order \(\ge N-\rho n\) inside it contributes only to
\(\FormalBigOAt{t^{N}}{t}\).
\end{proof}

Finite determination also yields the two stability estimates that make
composition continuous in each argument.  They are used repeatedly in the
reversion theory.

\begin{proposition}[Stability in the outer and inner arguments]
\label{plt:prop:cmp-stability}
Let \(G,G_1,G_2\) be admissible monomial-leading arguments with the same
\(c\), \(\log c\), \(\rho\) and with \(U_1,U_2\in t\PLpow\).
\begin{enumerate}
\item If \(\ordt(F_1-F_2)\ge M\) then
\(\ordt(F_1\compose G-F_2\compose G)\ge\rho M\).
\item If \(\ordt(U_1-U_2)\ge M\ge1\) and \(a=\ordt(F)\), then
\(\ordt(F\compose G_1-F\compose G_2)\ge\rho a+M\).
\end{enumerate}
\end{proposition}

\begin{proof}
(i) is \cref{plt:thm:cmp-substitution-hom}(iii) applied to \(F_1-F_2\)
together with linearity of \(\Phi_G\).

For (ii) write \(V_i=\log(1+U_i)\) and \(D=U_1-U_2\).  From
\(U_1^{r}-U_2^{r}=D\sum_{i=0}^{r-1}U_1^{i}U_2^{r-1-i}\) and
\(\ordt(U_i)\ge1\) we get \(\ordt(U_1^r-U_2^r)\ge M+r-1\), whence
\(\ordt(V_1-V_2)\ge M\) by \eqref{plt:eq:cmp-exp-log-def} and
\cref{plt:lem:cmp-summable}.  The same computation applied to
\((1+U_i)^{-n}=\sum_r\binom{-n}{r}U_i^{r}\) gives
\(\ordt\bigl((1+U_1)^{-n}-(1+U_2)^{-n}\bigr)\ge M\).  For a polynomial
\(p\in\K[L]\) there is a two-variable polynomial \(q\) over \(\K\) with
\(p(y_1)-p(y_2)=(y_1-y_2)q(y_1,y_2)\); substituting
\(y_i=\rho L+\log c+V_i\), whose values have nonnegative valuation, gives
\(\ordt\bigl(p_n(\rho L+\log c+V_1)-p_n(\rho L+\log c+V_2)\bigr)\ge M\).
Writing the \(n\)-th summands of \eqref{plt:eq:cmp-master-sum} as
\(A_1B_1-A_2B_2=(A_1-A_2)B_1+A_2(B_1-B_2)\) with
\(A_i=p_n(\rho L+\log c+V_i)\) and \(B_i=c^{-n}t^{\rho n}(1+U_i)^{-n}\), all
factors having valuation at least \(0\) resp.\ \(\rho n\), we obtain a
valuation at least \(\rho n+M\ge\rho a+M\) for every \(n\).
\end{proof}

\section{The shifted-Taylor form of the composition formula}
\label{plt:sec:cmp-shifted-taylor}

Formula \eqref{plt:eq:cmp-master-sum} separates the logarithmic shift from the
power prefactor.  The two can be packaged into a single shifted operator,
which is both more compact and better adapted to coefficient recurrences.

\begin{lemma}[Shifted-Taylor identity]\label{plt:lem:cmp-shifted-taylor}
Let \(p\in\K[L]\), let \(n\in\IntegerNumbers\), and let \(W\) be an element of
positive valuation in a polynomial-logarithmic Hahn ring.  Then
\begin{equation}\label{plt:eq:cmp-shifted-taylor}
 \EulerE^{-nW}\,p(z+W)
 =\sum_{k\ge0}\frac{W^{k}}{k!}\,
   \bigl((\partial_z-n)^{k}p\bigr)(z),
\end{equation}
where \(z\) is any element of the ring, \(\partial_z\) denotes differentiation
of the polynomial \(p\) in its own variable, and the operator is applied to
\(p\) \emph{before} evaluation at \(z\).  The family on the right is summable;
for \(p\ne0\) it terminates exactly when \(n=0\), since for \(n\ne0\) the
polynomial \((\partial_z-n)^{k}p\) has the same degree as \(p\) and leading
coefficient \((-n)^{k}\) times that of \(p\).
\end{lemma}

\begin{proof}
Both sides are \(\K\)-linear in \(p\), so it suffices to expand.  On the left,
\[
 \EulerE^{-nW}p(z+W)
 =\Bigl(\sum_{i\ge0}\frac{(-n)^{i}W^{i}}{i!}\Bigr)
  \Bigl(\sum_{j\ge0}\frac{W^{j}}{j!}p^{(j)}(z)\Bigr),
\]
the second factor being the finite Taylor expansion of the polynomial \(p\).
Collecting the total power \(k=i+j\) gives the coefficient
\(\sum_{j=0}^{k}\frac{(-n)^{k-j}}{(k-j)!\,j!}p^{(j)}(z)\).  On the right, the
binomial theorem for the commuting operators \(\partial_z\) and \(-n\) gives
\[
 \frac{1}{k!}(\partial_z-n)^{k}p
 =\frac{1}{k!}\sum_{j=0}^{k}\binom{k}{j}(-n)^{k-j}p^{(j)}
 =\sum_{j=0}^{k}\frac{(-n)^{k-j}}{(k-j)!\,j!}p^{(j)} .
\]
The two coefficients agree.  Summability holds because
\(\ordt(W^{k})\ge k\ordt(W)\to+\infty\), and the rearrangement of the double
family is licensed by \cref{plt:lem:cmp-summable}.
\end{proof}

\begin{corollary}[Operator form of the composition formula]
\label{plt:cor:cmp-operator-form}
With \(G\) as in \cref{plt:def:cmp-admissible-inner},
\begin{equation}\label{plt:eq:cmp-operator-form}
 F\compose G
 =\sum_{n\ge n_0}c^{-n}t^{\rho n}\,E_n,
 \qquad
 E_n\DefinitionEquals
 \sum_{k\ge0}\frac{V^{k}}{k!}
   \bigl((\partial_z-n)^{k}p_n\bigr)(\rho L+\log c).
\end{equation}
\end{corollary}

\begin{proof}
By \cref{plt:def:cmp-admissible-inner} and \cref{plt:lem:cmp-exp-log},
\((1+U)^{-n}=\EulerE^{-nV}\).  Apply
\cref{plt:lem:cmp-shifted-taylor} with \(W=V\) and \(z=\rho L+\log c\), and
substitute into \eqref{plt:eq:cmp-master-sum}.
\end{proof}

% ed.: S2 writes the operator in (its) power-leading formula as ":e^{delta(D-n)}:"
% with D previously declared to be d/dL, while the identity requires
% differentiation in the *argument* of the coefficient polynomial.  We print
% partial_z for the dummy variable and evaluate afterwards.
% ed.: an earlier draft of this warningbox described the discrepancy as "a
% spurious factor rho^k".  That is false as soon as n != 0: the chain rule
% turns (d_z - n)^k into (rho d_z - n)^k, and the two are proportional only
% when n = 0.  Corrected below, with a numerical witness.
\begin{warningbox}
\textbf{Two ordering traps.}
First, \emph{normal ordering}.  The operator in
\eqref{plt:eq:cmp-operator-form} may be written suggestively as
\(\mathopen{:}\EulerE^{V(\partial_z-n)}\mathclose{:}\), the colons meaning
that every power of \(V\) stands to the left and is never differentiated.  It
is \emph{not} the operator exponential \(\exp\bigl(M_V(\partial_L-n)\bigr)\),
where \(M_V\) is multiplication by \(V\): iterating \(M_V\partial_L\) would let
\(\partial_L\) act on \(V\), which depends on \(L\).  The two agree only when
\(\partial_LV=0\).

Second, \emph{which variable is differentiated}.  In
\eqref{plt:eq:cmp-operator-form} the operator \(\partial_z\) acts on the
coefficient polynomial \(p_n\) in its own variable, and the result is
evaluated at \(\rho L+\log c\) afterwards.  Writing \(\partial_L\) instead is
harmless only when \(\rho=1\).  Indeed the chain rule gives
\[
 (\partial_L-n)^{k}\bigl[p_n(\rho L+\log c)\bigr]
 =\bigl((\rho\,\partial_z-n)^{k}p_n\bigr)(\rho L+\log c),
\]
so the misprint replaces the operator \((\partial_z-n)^{k}\) by
\((\rho\,\partial_z-n)^{k}\).  These two agree only for \(\rho=1\), and they
are not even proportional: the factor \(\rho^{k}\) relates them in the single
case \(n=0\).  For \(\rho=2\), \(n=1\), \(\log c=0\), \(p_n(z)=z\) and
\(k=1\), the correct coefficient of \(V\) is \(1-2L\) whereas the misprint
produces \(2-2L\).
\end{warningbox}

The operator form makes the block coefficients of a single outer level satisfy
an exact first-order recurrence.  We state it for arbitrary \(\rho\); the case
\(\rho=1\) is the form found in the sources.

\begin{proposition}[Exact block recurrence]\label{plt:prop:cmp-block-recurrence}
Assume \(\rho\in\K^\times\).  With \(E_n\) as in
\eqref{plt:eq:cmp-operator-form}, write
\(E_n=\sum_{r\ge0}E_{n,r}(L)t^{r}\) and
\begin{equation}\label{plt:eq:cmp-theta-def}
 \theta\DefinitionEquals
 \frac{\partial_tV}{\rho+\partial_LV}
 =\sum_{q\ge0}\theta_q(L)t^{q},
\end{equation}
the denominator being a unit of \(\PLpow\).  Then
\begin{equation}\label{plt:eq:cmp-E-pde}
 \partial_tE_n=\theta\,(\partial_L-n\rho)E_n,
\end{equation}
and therefore
\begin{equation}\label{plt:eq:cmp-E-recurrence}
 E_{n,0}=p_n(\rho L+\log c),
 \qquad
 r\,E_{n,r}
 =\sum_{q=0}^{r-1}\theta_q\,(\partial_L-n\rho)E_{n,r-1-q}
 \quad(r\ge1).
\end{equation}
\end{proposition}

\begin{proof}
By \cref{plt:cor:cmp-operator-form} and \cref{plt:lem:cmp-shifted-taylor},
\(E_n=\EulerE^{-nV}\,p_n(\rho L+\log c+V)\).  Write \(P=p_n\) and
\(w=\rho L+\log c+V\).  Since \(\partial_t\) and \(\partial_L\) are formal
partial derivatives and \(\partial_tw=\partial_tV\),
\(\partial_Lw=\rho+\partial_LV\),
\begin{align*}
 \partial_tE_n
 &=\EulerE^{-nV}\,\partial_tV\,\bigl(P'-nP\bigr)(w),\\
 \partial_LE_n
 &=\EulerE^{-nV}\Bigl[-n(\partial_LV)P(w)+(\rho+\partial_LV)P'(w)\Bigr].
\end{align*}
Subtracting \(n\rho E_n=\EulerE^{-nV}n\rho P(w)\) from the second line gives
\[
 (\partial_L-n\rho)E_n
 =\EulerE^{-nV}\,(\rho+\partial_LV)\,\bigl(P'-nP\bigr)(w),
\]
because \(-n(\partial_LV)-n\rho=-n(\rho+\partial_LV)\).  Dividing the two
displays yields \eqref{plt:eq:cmp-E-pde}; the denominator
\(\rho+\partial_LV\) is a unit because \(\rho\in\K^\times\) and
\(\partial_LV\in t\PLpow\).  Comparing coefficients of \(t^{r-1}\) in
\eqref{plt:eq:cmp-E-pde} gives \eqref{plt:eq:cmp-E-recurrence}; and
\(E_{n,0}=[t^{0}]E_n=p_n(\rho L+\log c)\) because \(\ordt(V^{k})\ge k\), so
only \(k=0\) contributes to the block \(t^{0}\) in
\eqref{plt:eq:cmp-operator-form}.
\end{proof}

\begin{remark}
The naive substitute \(\theta=\partial_tV\) is wrong.  It ignores that the
logarithmic shift itself depends on \(L\), and it already fails at \(r=2\)
whenever \(\partial_LV\ne0\).  The correction factor
\((\rho+\partial_LV)^{-1}\) is exactly the Jacobian of the change of
logarithmic coordinate.
\end{remark}

\section{Rational and real leading exponents}
\label{plt:sec:cmp-fractional}

Nothing in the proof of \cref{plt:thm:cmp-substitution-hom} used integrality of
\(\rho\) except to keep the output in \(\PLlau\).

\begin{proposition}[Puiseux-log and Hahn extensions]
\label{plt:prop:cmp-fractional-rho}
Let \(F=\sum_{n\ge a}p_n(L)t^n\in\PLlau\), let \(c\in\K^\times\) with a chosen
\(\log c\in\K\), let \(U\in t\PLpow\), and let \(\rho>0\) be real with
\(\rho\in\K\).  Define \(G=cX^{\rho}(1+U)\) and \(F\compose G\) by
\eqref{plt:eq:cmp-master-sum}.
\begin{enumerate}
\item If \(\rho=r/s\) with \(r,s\in\PositiveIntegers\), then
\(F\compose G\in\K[L]((t^{1/s}))\), the Puiseux-log ring obtained by adjoining
\(t^{1/s}\).
\item For arbitrary real \(\rho>0\), the family
\eqref{plt:eq:cmp-master} is summable in
\(\K[L][[t^{\Gamma}]]\) with
\(\Gamma=\IntegerNumbers+\rho\IntegerNumbers\), and
\[
 \CoefficientSupportOf{F\compose G}
 \subseteq
 \SetOf{(\rho n+m,\,j):\ n\ge a,\ m\in\NaturalNumbers,\ j\in\NaturalNumbers}.
\]
The exponent set \(\SetOf{\rho n+m:n\ge a,\ m\ge0}\) is locally finite: only
finitely many of its elements lie below any real bound.
\item In both cases the conclusions of
\cref{plt:thm:cmp-substitution-hom} and
\cref{plt:thm:cmp-finite-determination} hold with \(\PLlau\) replaced as
target by \(\K[L]((t^{1/s}))\) resp.\ \(\K[L][[t^{\Gamma}]]\), with \(N\)
ranging over \(\Gamma\) and \(\Floor{N/\rho}\) read as the largest admissible
\(n\).
\end{enumerate}
\end{proposition}

\begin{proof}
Only the support statement needs an argument.  Fix a real bound \(B\).  Since
\(\rho>0\), the inequality \(\rho n+m\le B\) with \(n\ge a\), \(m\ge0\) forces
\(a\le n\le B/\rho\), finitely many values, and for each such \(n\) it
forces \(0\le m\le B-\rho n\), again finitely many.  Hence the exponent set is
locally finite, in particular well ordered, and it is contained in the group
\(\Gamma\).  In particular the \(n\)-th member of \eqref{plt:eq:cmp-master}
has valuation at least \(\rho n\) and
\(\SetOf{n\ge a:\rho n\le\beta}\) is finite for every real \(\beta\), which is
\eqref{plt:eq:cmp-summability-criterion};
so \cref{plt:lem:cmp-summable} applies and every proof above goes through
word for word.  When \(\rho=r/s\) we have
\(\Gamma\subseteq s^{-1}\IntegerNumbers\) and local finiteness upgrades to
"finitely many negative exponents", which is the Puiseux-log condition.
\end{proof}

\begin{remark}
The union over all \(s\) of the rings \(\K[L]((t^{1/s}))\) admits rational
exponents but no single denominator bound; it is a field, but a fixed element
of it does not carry a uniform \(s\).  Real exponents genuinely require the
Hahn setting, where the support condition, not the exponent group, does the
work.  Complex exponents require in addition a declared tie-breaking order,
since ordering by real part alone is not a total order
\cite{hahn,vanderhoeven-book}.
\end{remark}

\begin{example}[Where the one-logarithm class ceases to be closed]
\label{plt:ex:cmp-closure-failures}
The following table is a diagnostic; each row is proved by writing down
\(t\compose G\) and \(L\compose G\).
\begin{center}
\begin{tabular}{@{}lll@{}}
\hline
Leading behavior of \(G\) & New object created & Minimal enlargement\\
\hline
\(cX^{\rho}\), \(\rho\in\PositiveIntegers\)
 & none & \(\PLlau\)\\
\(cX^{r/s}\)
 & \(t^{1/s}\) & \(\K[L]((t^{1/s}))\)\\
\(cX^{\rho}\), \(\rho\) real
 & \(t^{\rho}\) & \(\K[L][[t^{\Gamma}]]\)\\
\(cX^{\rho}L^{\sigma}\), \(\sigma\ne0\)
 & \(\log L\) & one further iterated-log level\\
\(cL^{\sigma}\), \(\sigma>0\)
 & \(L^{-\sigma}\), and \(\log L\) if some \(p_n\notin\K\)
 & completion order reversed\\
\(\EulerE^{X}\)
 & \(\EulerE^{-nX}\) & logarithmic-exponential transseries\\
\hline
\end{tabular}
\end{center}
For the fourth row, \(\log G=\rho L+\sigma\log L+\log c+\log(1+U)\), so a
single outer coefficient of positive degree in its argument already produces
\(\log\log X\).  The fifth row is the most drastic: \(t\compose G=c^{-1}L^{-\sigma}\)
has outer valuation \(0\), so the composite is not a series in \(t\) at all
and the roles of the two completions are exchanged; since
\(\PLit=\K((L^{-1}))((t))\ne\K((t))((L^{-1}))\), naming the target field is
part of the statement.  For the sixth,
\(t\compose\EulerE^{X}=\EulerE^{-X}\), which is smaller than every power of
\(t\); such height-changing substitutions belong to the field of
logarithmic-exponential transseries
\cite{edgar-beginners,vanderhoeven-book,adh-book}.
\end{example}

\begin{example}[A scale-changing substitution computed exactly]
\label{plt:ex:cmp-scale-change}
Let \(F(Y)=Y+\log Y+Y^{-1}\) and \(G(X)=3X^{2}L^{-1}\) on the ray
\(X\to+\infty\).  The inner argument is \emph{not} admissible in the sense of
\cref{plt:def:cmp-admissible-inner}, because of the factor \(L^{-1}\); this is
exactly why a new scale appears.  Substituting the generators,
\[
 F\compose G
 =3X^{2}L^{-1}+\bigl(\log3+2L-\log L\bigr)+\tfrac13Lt^{2}.
\]
The second logarithmic level appears exactly, not as a remainder; the result
lies in a two-logarithm class and in no smaller one.
\end{example}

\section{Generator substitution for a near-identity argument}
\label{plt:sec:cmp-generator-substitution}

The case \(c=1\), \(\log c=0\), \(\rho=1\) is the one that governs reversion.

\begin{definition}[Near-identity argument]\label{plt:def:cmp-near-identity}
For \(H\in\PLpow\) set
\[
 G=X+H=X(1+tH),
\]
so that \(U=tH\in t\PLpow\), \(c=1\), \(\log c=0\), \(\rho=1\).  We write
\(\Gnear=\SetOf{X+H:H\in\PLpow}\).
\end{definition}

\begin{proposition}[Exact generator substitution]
\label{plt:prop:cmp-generator-substitution}
For \(G=X+H\) with \(H\in\PLpow\),
\begin{equation}\label{plt:eq:cmp-generators-near-identity}
 t\compose G=\frac{t}{1+tH},
 \qquad
 L\compose G=L+\log(1+tH),
\end{equation}
and consequently, for \(F=\sum_{n\ge n_0}p_n(L)t^n\),
\begin{equation}\label{plt:eq:cmp-near-identity-master}
 F\compose(X+H)
 =\sum_{n\ge n_0}t^{n}(1+tH)^{-n}\,
   p_n\!\bigl(L+\log(1+tH)\bigr).
\end{equation}
Moreover \(\ordt(F\compose(X+H))=\ordt(F)\), so \(\Phi_{X+H}\) preserves the
valuation exactly and maps \(\PLpow\) injectively into \(\PLpow\).
\end{proposition}

\begin{proof}
Specialize \eqref{plt:eq:cmp-generator-images} and
\eqref{plt:eq:cmp-master-sum} to \(c=1\), \(\log c=0\), \(\rho=1\),
\(U=tH\); the valuation statement is
\cref{plt:thm:cmp-substitution-hom}(iii) with \(\rho=1\).
\end{proof}

% ed.: the sources state the rule in words; the numerical demonstration of the
% two standard mistakes is supplied here so that the size of the error is
% visible.
\begin{warningbox}
\textbf{Both generators move, and neither moves the way it looks.}
Let \(H\in\PLpow\) and \(G=X+H\).

\emph{Wrong rule 1: \(t\mapsto t+H\).}  The symbol \(t\) denotes \(X^{-1}\),
not \(X\).  Adding \(H\) to \(t\) models \(X^{-1}\mapsto X^{-1}+H\), which is
not \((X+H)^{-1}\).  The mistake is visible at once from valuations: for
\(H\notin t\PLpow\) one has \(\ordt(t+H)=0\), whereas \(\ordt(1/G)=1\) because
\(G\to\infty\).

\emph{Wrong rule 2: update \(L\), hold \(t\) fixed.}  Take
\(F=tL=\log X/X\) and \(H=L\), so \(G=X+\log X\).  The correct value, from
\eqref{plt:eq:cmp-near-identity-master}, is
\[
 F\compose G
 =\bigl(t-Lt^{2}+L^{2}t^{3}-\cdots\bigr)
  \bigl(L+Lt-\tfrac12L^{2}t^{2}+\cdots\bigr)
 =tL+(L-L^{2})t^{2}+\Bigl(L^{3}-\tfrac32L^{2}\Bigr)t^{3}
  +\FormalBigOAt{t^{4}}{t}.
\]
Substituting only \(L\mapsto L+\log(1+Lt)\) returns
\(tL+Lt^{2}-\tfrac12L^{2}t^{3}+\FormalBigOAt{t^{4}}{t}\), which misses the
term \(-L^{2}t^{2}\) entirely: the answer is already wrong in the first
corrected block.  Substituting only \(t\mapsto t/(1+Lt)\) returns
\(tL-L^{2}t^{2}+L^{3}t^{3}+\FormalBigOAt{t^{4}}{t}\) and misses \(+Lt^{2}\).
A correct implementation substitutes
the two generators \emph{simultaneously}: the replacement for \(L\) contains
\(t\), and the replacement for \(t\) contains \(L\), so no sequential
substitution is legitimate.
\end{warningbox}

\begin{proposition}[The near-identity monoid]\label{plt:prop:cmp-near-monoid}
For \(A,B\in\PLpow\),
\begin{equation}\label{plt:eq:cmp-near-group-law}
 (X+A)\compose(X+B)=X+B+A\compose(X+B),
\end{equation}
and \(A\compose(X+B)\in\PLpow\).  Hence \(\Gnear\) is a monoid under
composition, with identity \(X\), and the map
\(G\mapsto\Phi_G\) is an anti-homomorphism of monoids
(\cref{plt:thm:cmp-associativity}).
\end{proposition}

\begin{proof}
Apply \(\Phi_{X+B}\), a \(\K\)-algebra homomorphism, to \(X+A=t^{-1}+A\):
\(\Phi_{X+B}(t^{-1})=(\Phi_{X+B}t)^{-1}=X+B\), and
\(\Phi_{X+B}(A)=A\compose(X+B)\), which lies in \(\PLpow\) by
\cref{plt:prop:cmp-generator-substitution}.
\end{proof}

\begin{proposition}[Filtration continuity, with a one-block gain on the right]
\label{plt:prop:cmp-filtration-continuity}
Let \(F,F_1,F_2\in\PLlau\) and \(H,K\in\PLpow\).  Then
\begin{align}
 \ordt\bigl(F_1\compose(X+H)-F_2\compose(X+H)\bigr)
 &=\ordt(F_1-F_2),
 \label{plt:eq:cmp-continuity-left}\\
 \ordt\bigl(F\compose(X+H)-F\compose(X+K)\bigr)
 &\ge\ordt(F)+1+\ordt(H-K).
 \label{plt:eq:cmp-continuity-right}
\end{align}
\end{proposition}

\begin{proof}
\eqref{plt:eq:cmp-continuity-left} is
\cref{plt:prop:cmp-generator-substitution} applied to \(F_1-F_2\).  For
\eqref{plt:eq:cmp-continuity-right} use the Taylor formula
\eqref{plt:eq:cmp-taylor}, proved independently in
\cref{plt:sec:taylor}: the difference equals
\(\sum_{k\ge1}\frac{H^{k}-K^{k}}{k!}\DX^{k}F\).  Since
\(H^{k}-K^{k}=(H-K)\sum_{i=0}^{k-1}H^{i}K^{k-1-i}\) and all factors lie in
\(\PLpow\), \(\ordt(H^{k}-K^{k})\ge\ordt(H-K)\); and
\(\ordt(\DX^{k}F)\ge\ordt(F)+k\ge\ordt(F)+1\) for \(k\ge1\).
\end{proof}

\begin{example}[A worked near-identity composition]
\label{plt:ex:cmp-worked-near-identity}
Let
\[
 F=L^{2}+(L+1)t+(L^{2}-L)t^{2}+(2L+1)t^{3},
 \qquad
 G=X+L .
\]
Here \(U=Lt\) and, through outer order \(3\),
\[
 t\compose G=t-Lt^{2}+L^{2}t^{3}+\FormalBigOAt{t^{4}}{t},
 \qquad
 L\compose G=L+Lt-\tfrac12L^{2}t^{2}+\tfrac13L^{3}t^{3}
  +\FormalBigOAt{t^{4}}{t}.
\]
Simultaneous substitution and collection of blocks give
\begin{equation}\label{plt:eq:cmp-worked}
 F\compose G
 =L^{2}+(2L^{2}+L+1)t+(-L^{3}+L^{2}-L)t^{2}
  +\Bigl(\tfrac23L^{4}-2L^{3}+\tfrac72L^{2}+L+1\Bigr)t^{3}
  +\FormalBigOAt{t^{4}}{t}.
\end{equation}
\end{example}

\begin{example}[Increasing logarithmic degree is not an artifact]
\label{plt:ex:cmp-log-degree-growth}
Let \(F=L+L^{2}t\) and \(G=X+L\).  Then
\[
 F\compose G
 =L+(L^{2}+L)t+\Bigl(-L^{3}+\tfrac32L^{2}\Bigr)t^{2}
  +\Bigl(L^{4}-\tfrac83L^{3}+L^{2}\Bigr)t^{3}
  +\Bigl(-L^{5}+\tfrac{41}{12}L^{4}-2L^{3}\Bigr)t^{4}
  +\FormalBigOAt{t^{5}}{t}.
\]
The log-degree profile \(\dprofile{F\compose G}\) grows by one per block:
every extra power of \(\log(1+Lt)\) supplies another factor \(L\).  No
uniform bound on \(\degL\) survives composition, which is why the ring is
defined without one.
\end{example}

\section{The associativity domain}
\label{plt:sec:cmp-associativity}

Associativity is a theorem, not a formality: it holds on the nose only when
the composite inner argument is itself admissible \emph{and} the chosen
logarithms are compatible.

\begin{theorem}[Associativity]\label{plt:thm:cmp-associativity}
Let
\[
 G=c_{G}X^{\rho}(1+U),
 \qquad
 H=c_{H}X^{\sigma}(1+W)
\]
be admissible monomial-leading arguments with chosen logarithms
\(\log c_{G},\log c_{H}\in\K\), with
\(\rho,\sigma\in\PositiveIntegers\) and \(U,W\in t\PLpow\).  Put
\begin{equation}\label{plt:eq:cmp-composed-inner}
 c\DefinitionEquals c_{G}c_{H}^{\rho},
 \qquad
 \log c\DefinitionEquals\log c_{G}+\rho\log c_{H},
 \qquad
 1+U^{\ast}\DefinitionEquals(1+W)^{\rho}\bigl(1+\Phi_{H}U\bigr).
\end{equation}
Then \(U^{\ast}\in t\PLpow\), the quadruple
\((c,\log c,\rho\sigma,U^{\ast})\) is an admissible monomial-leading argument
representing \(G\compose H\), and for every \(F\in\PLlau\)
\begin{equation}\label{plt:eq:cmp-associativity}
 (F\compose G)\compose H=F\compose(G\compose H),
 \qquad\text{equivalently}\qquad
 \Phi_{H}\circ\Phi_{G}=\Phi_{G\compose H}.
\end{equation}
The identity \(X\) is a two-sided unit, and \(\Gnear\) is closed under
composition with the canonical choices \(c=1\), \(\log c=0\).
\end{theorem}

\begin{proof}
First, admissibility.  \(\Phi_{H}\) is a \(\K\)-algebra homomorphism
compatible with summable families
(\cref{plt:thm:cmp-substitution-hom}), and
\(\ordt(\Phi_{H}U)=\sigma\,\ordt(U)\ge\sigma\ge1\); also
\(\ordt\bigl((1+W)^{\rho}-1\bigr)\ge1\), since \(\rho\) is a positive integer
and \((1+W)^{\rho}-1\) is a sum of positive powers of \(W\).  Hence
\(U^{\ast}=(1+W)^{\rho}(1+\Phi_HU)-1\in t\PLpow\).  Furthermore
\[
 G\compose H
 =\Phi_{H}(G)
 =c_{G}\,(\Phi_{H}t)^{-\rho}\,(1+\Phi_{H}U)
 =c_{G}c_{H}^{\rho}\,t^{-\rho\sigma}(1+W)^{\rho}(1+\Phi_{H}U)
 =c\,t^{-\rho\sigma}(1+U^{\ast}),
\]
which is \eqref{plt:eq:cmp-composed-inner}.

Second, the identity.  Both \(\Phi_{H}\circ\Phi_{G}\) and
\(\Phi_{G\compose H}\) are \(\K\)-algebra homomorphisms
\(\PLlau\to\PLlau\) compatible with summable families, the first
because a composite of two such maps is one
(\cref{plt:lem:cmp-summable}).  By the definition
\eqref{plt:eq:cmp-master-sum} of \(\Phi\), a homomorphism with this
compatibility is determined by its values on \(t\) and \(L\): indeed it maps
\(p_n(L)t^n\) to \(p_n(\text{image of }L)\cdot(\text{image of }t)^{n}\) and
then sums.  So it suffices to compare the generators.

For \(t\):
\[
 \Phi_{H}\bigl(\Phi_{G}t\bigr)
 =\Phi_{H}\bigl(c_{G}^{-1}t^{\rho}(1+U)^{-1}\bigr)
 =c_{G}^{-1}\bigl(c_{H}^{-1}t^{\sigma}(1+W)^{-1}\bigr)^{\rho}
   \bigl(1+\Phi_{H}U\bigr)^{-1}
 =c^{-1}t^{\rho\sigma}(1+U^{\ast})^{-1}
 =\Phi_{G\compose H}(t).
\]
For \(L\), using \cref{plt:lem:cmp-exp-log} twice,
\begin{align*}
 \Phi_{H}\bigl(\Phi_{G}L\bigr)
 &=\Phi_{H}\bigl(\rho L+\log c_{G}+\log(1+U)\bigr)\\
 &=\rho\bigl(\sigma L+\log c_{H}+\log(1+W)\bigr)+\log c_{G}
   +\log\bigl(1+\Phi_{H}U\bigr)\\
 &=\rho\sigma L+\bigl(\log c_{G}+\rho\log c_{H}\bigr)
   +\log\Bigl((1+W)^{\rho}\bigl(1+\Phi_{H}U\bigr)\Bigr)\\
 &=\rho\sigma L+\log c+\log(1+U^{\ast})
 =\Phi_{G\compose H}(L).
\end{align*}
Hence the two homomorphisms agree.  Taking \(G=X\) with \(c=1\),
\(\log c=0\), \(\rho=1\), \(U=0\) gives \(\Phi_{X}=\mathrm{id}\), which is the
two-sided identity statement; closure of \(\Gnear\) is
\cref{plt:prop:cmp-near-monoid}.
\end{proof}

\begin{remark}[Right substitution reverses the apparent order]
\label{plt:rmk:cmp-order-reversal}
Equation \eqref{plt:eq:cmp-associativity} reads
\(\Phi_{H}\circ\Phi_{G}=\Phi_{G\compose H}\): the operator that acts
\emph{second} is the one indexed by the inner argument that acts
\emph{first}.  Substitution on the right is a contravariant operation, so
\(G\mapsto\Phi_{G}\) is an anti-homomorphism.  Implementations of nested
composition should test exactly this identity; reversing it silently produces
correct-looking output for commuting pairs and wrong output otherwise.
\end{remark}

\begin{remark}[The compatibility hypothesis on logarithms cannot be dropped]
Suppose \(\K=\ComplexNumbers\) and one chooses for \(G\compose H\) the
logarithm \(\log c+2\pi\ImaginaryUnit\) instead of
\(\log c_{G}+\rho\log c_{H}\).  Then \(\Phi_{G\compose H}(L)\) differs from
\(\Phi_{H}(\Phi_{G}L)\) by \(2\pi\ImaginaryUnit\), and
\eqref{plt:eq:cmp-associativity} fails for every \(F\) with a coefficient
block of positive degree.  Associativity is a statement about admissible
\emph{data}, not about the underlying maps alone.
\end{remark}

\begin{remark}[Rational and real exponents]
\Cref{plt:thm:cmp-associativity} was stated for
\(\rho,\sigma\in\PositiveIntegers\) so that every displayed object stays in
\(\PLlau\).  For \(\rho,\sigma>0\) lying in \(\K\), the same proof works
verbatim inside the Hahn ring \(\K[L][[t^{\Gamma}]]\) with
\(\Gamma=\IntegerNumbers+\rho\IntegerNumbers+\sigma\IntegerNumbers
 +\rho\sigma\IntegerNumbers\), provided one \emph{defines}
\((1+W)^{\rho}\DefinitionEquals\EulerE^{\rho\log(1+W)}\), which is legitimate
by \cref{plt:lem:cmp-exp-log} and agrees with the integer power when
\(\rho\in\IntegerNumbers\).  The exponent of the composite is \(\rho\sigma\),
so the class of admissible exponents must be closed under multiplication;
\(\PositiveIntegers\) and \(\RationalNumbers_{>0}\) are, a fixed denominator
lattice \(s^{-1}\IntegerNumbers\) is not.
\end{remark}

\begin{warningbox}
\textbf{Formal closure is not a license.}  \Cref{plt:thm:cmp-associativity}
asserts \((F\compose G)\compose H=F\compose(G\compose H)\) only when every
displayed composition exists in the selected support class.  A syntactically
meaningful nested expression in a larger Hahn field need not have summable
support; support theorems for general transseries composition are genuinely
harder and carry their own hypotheses \cite{edgar-composition,vanderhoeven-book}.
\end{warningbox}

\section{The chain rule}
\label{plt:sec:cmp-chain-rule}

\begin{theorem}[Chain rule]\label{plt:thm:cmp-chain-rule}
Let \(G\) be an admissible monomial-leading argument with \(\rho>0\).  Then
for every \(F\in\PLlau\),
\begin{equation}\label{plt:eq:cmp-chain-rule}
 \DX\bigl(F\compose G\bigr)=\bigl(\DX F\bigr)\compose G\cdot\DX G .
\end{equation}
\end{theorem}

\begin{proof}
Set \(\delta_1(F)=\DX(\Phi_{G}F)\) and
\(\delta_2(F)=\Phi_{G}(\DX F)\cdot\DX G\).  Both are \(\K\)-linear, both are
compatible with summable families (each is a composite of maps that transform
valuations by a non-decreasing unbounded function, so
\cref{plt:lem:cmp-summable} applies), and both satisfy the twisted Leibniz
rule
\(\delta(F_1F_2)=\delta(F_1)\Phi_{G}(F_2)+\Phi_{G}(F_1)\delta(F_2)\): for
\(\delta_1\) because \(\DX\) is a derivation and \(\Phi_G\) a homomorphism,
for \(\delta_2\) for the same reason.  It therefore suffices to check
\eqref{plt:eq:cmp-chain-rule} for \(F=t\) and \(F=L\); the twisted Leibniz
rule then propagates it to \(\K[L][t]\), the relation
\(\delta(t\cdot t^{-1})=0\) propagates it to \(t^{-1}\), and compatibility
with summable families finishes.

Write \(G=ct^{-\rho}(1+U)\).  From
\(\DX(t^{\alpha}p(L))=t^{\alpha+1}(p'-\alpha p)\) we get
\(\DX(t^{-\rho})=\rho t^{1-\rho}\), so
\begin{equation}\label{plt:eq:cmp-log-derivative}
 \frac{\DX G}{G}
 =\frac{c\rho t^{1-\rho}(1+U)+c\,t^{-\rho}\DX U}{c\,t^{-\rho}(1+U)}
 =\rho t+\frac{\DX U}{1+U}.
\end{equation}
For \(F=t\): \(\delta_1(t)=\DX(G^{-1})=-G^{-2}\DX G\), while
\(\delta_2(t)=\Phi_{G}(-t^{2})\DX G=-G^{-2}\DX G\).  For \(F=L\):
\(\delta_2(L)=\Phi_{G}(t)\DX G=\DX G/G\), while
\[
 \delta_1(L)=\DX\bigl(\rho L+\log c+\log(1+U)\bigr)
 =\rho t+\sum_{r\ge1}(-1)^{r-1}U^{r-1}\DX U
 =\rho t+\frac{\DX U}{1+U},
\]
termwise differentiation being legitimate by
\cref{plt:lem:cmp-summable}.  By \eqref{plt:eq:cmp-log-derivative} the two
agree.
\end{proof}

\section{Analytic admissibility}
\label{plt:sec:cmp-analytic}

Everything above is an identity between formal objects.  It becomes a
statement about functions only under additional hypotheses, which we now make
explicit.  Take \(\K\subseteq\ComplexNumbers\).

\begin{definition}[Asymptotic realization]\label{plt:def:cmp-realization}
Let \(S\) be a ray \([X_{0},\infty)\) or a closed subsector of a sector at
infinity on which a branch of \(\log\) has been fixed.  A function
\(f:S\to\ComplexNumbers\) \emph{realizes} \(F=\sum_{n\ge a}p_n(L)t^n\in\PLlau\)
on \(S\), written \(f\sim_{S}F\), if for every \(N\in\IntegerNumbers\)
\begin{equation}\label{plt:eq:cmp-realization}
 f(X)-\bigl(\Trunc{F}{N}\bigr)(X)
 =\LittleOAt{X^{1-N}}{X\to\infty,\ X\in S},
\end{equation}
where \(\bigl(\Trunc{F}{N}\bigr)(X)=\sum_{a\le n<N}p_n(\log X)X^{-n}\) is a
finite sum evaluated with the fixed branch.
\end{definition}

% ed.: the exclusive normalization is used throughout; a source using the
% inclusive cutoff states the same condition as
% f - sum_{n<=N} = o(X^{-N}), which is (plt:eq:cmp-realization) at N+1.
\begin{lemma}[Equivalent forms and uniqueness]\label{plt:lem:cmp-realization-basic}
\(f\sim_{S}F\) holds if and only if for every \(N\) and every
\(\varepsilon>0\),
\(f-\Trunc{F}{N}=\BigOAt{X^{-N+\varepsilon}}{X\to\infty,\ X\in S}\).  A
function realizes at most one \(F\).  The stronger relation
\(f-\Trunc{F}{N}=\BigOAt{X^{-N}}{X\to\infty,\ X\in S}\) \emph{fails} for a
realization whenever \(\degL p_{N}>0\), since the first omitted block is then
not \(\BigOAt{X^{-N}}{X\to\infty}\); the \(\varepsilon\) may not be dropped.
\end{lemma}

\begin{proof}
If \eqref{plt:eq:cmp-realization} holds for all \(N\), apply it at \(N+1\) and
add \(p_{N}(\log X)X^{-N}=\BigOAt{X^{-N+\varepsilon}}{X\to\infty}\).
Conversely the \(\varepsilon\)-form at \(N\) with \(\varepsilon<1\) implies
\eqref{plt:eq:cmp-realization}.  For uniqueness, if \(F\ne\tilde F\) first
differ in block \(M\) by a nonzero \(R\in\K[L]\), subtracting the two
relations \eqref{plt:eq:cmp-realization} at \(N=M+1\) gives
\(R(\log X)X^{-M}=\LittleOAt{X^{-M}}{X\to\infty}\), so
\(R(\log X)\to0\); a nonzero polynomial does not tend to \(0\) along
\(\log X\to\infty\).  The last sentence is the observation that
\(X^{-N}(\log X)^{d}\) is not \(\BigOAt{X^{-N}}{X\to\infty}\) for \(d>0\).
\end{proof}

\begin{lemma}[Realizations form a ring]\label{plt:lem:cmp-realization-ring}
If \(f\sim_{S}F\) and \(g\sim_{S}\tilde F\) then \(f+g\sim_{S}F+\tilde F\) and
\(fg\sim_{S}F\tilde F\).
\end{lemma}

\begin{proof}
Addition is clear.  For the product put \(a=\ordt F\),
\(\tilde a=\ordt\tilde F\) and fix \(N\) and \(\varepsilon>0\).
% ed.: the draft omitted this reduction and then asserted that the product of
% the two remainders is O(X^{-N+eps}); that product is only
% O(X^{-2N+a+atilde+eps/2}), which is the required bound precisely when
% N >= a + atilde.  The reduction below makes the omission harmless.
Since \(\Trunc{F\tilde F}{N'}-\Trunc{F\tilde F}{N}\) is, for \(N'\ge N\), a
finite sum of terms \(q(\log X)X^{-m}\) with \(m\ge N\) and hence
\(\BigOAt{X^{-N+\varepsilon}}{X\to\infty}\), the assertion at a cutoff \(N'\)
implies the assertion at every smaller cutoff; we may therefore assume
\(N\ge a+\tilde a\).  Write
\(f=\Trunc{F}{N-\tilde a}+R\), \(g=\Trunc{\tilde F}{N-a}+\tilde R\).  By
\cref{plt:lem:cmp-realization-basic},
\(R=\BigOAt{X^{-N+\tilde a+\varepsilon/4}}{X\to\infty}\) and
\(\tilde R=\BigOAt{X^{-N+a+\varepsilon/4}}{X\to\infty}\), while
\(\Trunc{F}{N-\tilde a}=\BigOAt{X^{-a+\varepsilon/4}}{X\to\infty}\) and
similarly for \(\tilde F\).  Hence each of the three cross terms is
\(\BigOAt{X^{-N+\varepsilon}}{X\to\infty}\); for the term \(R\tilde R\),
whose exponent is \(-2N+a+\tilde a+\varepsilon/2\), this is where
\(N\ge a+\tilde a\) is used.  The remaining product of two
finite sums is a finite sum of terms \(q(\log X)X^{-m}\); those with \(m<N\)
constitute \(\Trunc{F\tilde F}{N}\) by the definition of the Cauchy product,
and those with \(m\ge N\) are \(\BigOAt{X^{-N+\varepsilon}}{X\to\infty}\).
\end{proof}

\begin{theorem}[Analytic transfer under composition]
\label{plt:thm:cmp-analytic-transfer}
Let \(S'=[X_{1},\infty)\) and \(S=[Y_{0},\infty)\) be rays, let
\(f\sim_{S}F\) with \(F=\sum_{n\ge a}p_n(L)t^n\), and let
\(g:S'\to\RealNumbers\) satisfy the following three conditions.
\begin{enumerate}
\item[\textup{(A1)}] \emph{Regime.}  \(g(X)\to+\infty\) and
\(g(X)\in S\) for all sufficiently large \(X\in S'\).
\item[\textup{(A2)}] \emph{Monomial-leading realization.}  There are
\(c>0\), \(\rho>0\) real and \(U\in t\PLpow\) with real coefficients such
that \(u\DefinitionEquals c^{-1}X^{-\rho}g(X)-1\) satisfies
\(u\sim_{S'}U\) and \(u(X)\to0\).
\item[\textup{(A3)}] \emph{Branch compatibility.}
\(\log g(X)=\rho\log X+\log c+\log(1+u(X))\) for all large \(X\in S'\), with
the principal real logarithm on both sides.
\end{enumerate}
Then \(f\circ g\) realizes \(F\compose G\) on \(S'\), where
\(G=cX^{\rho}(1+U)\) and \(F\compose G\) is
\eqref{plt:eq:cmp-master-sum} in the Hahn ring of
\cref{plt:prop:cmp-fractional-rho}, whose support is locally finite so that
\(\Trunc{\cdot}{\beta}\) at a real cutoff \(\beta\) is again a finite sum;
that is, for every real \(\beta\),
\[
 f(g(X))-\bigl(\Trunc{F\compose G}{\beta}\bigr)(X)
 =\LittleOAt{X^{-\beta+1}}{X\to+\infty}.
\]
\end{theorem}

\begin{proof}
Fix \(\beta\) and \(\varepsilon\in(0,1)\).  Choose an integer
\(N>(\beta+1)/\rho\) with \(N>a\), and split
\(f=\Trunc{F}{N}+R_{N}\).  By \cref{plt:lem:cmp-realization-basic},
\(R_{N}(Y)=\BigOAt{Y^{-N+\varepsilon}}{Y\to+\infty}\).  By (A1) and (A2),
\(g(X)=cX^{\rho}(1+u(X))\) with \(u\to0\), so \(g(X)\asymp X^{\rho}\) and
\begin{equation}\label{plt:eq:cmp-remainder-transport}
 R_{N}\bigl(g(X)\bigr)
 =\BigOAt{X^{-\rho(N-\varepsilon)}}{X\to+\infty}
 =\LittleOAt{X^{-\beta}}{X\to+\infty}
\end{equation}
for \(\varepsilon\) small, since \(\rho N>\beta+1\).  This is the
\emph{remainder transport} step; note that it consumes \(\rho^{-1}\) times as
many outer orders of the outer expansion as it produces.

It remains to expand the finite sum \(\Trunc{F}{N}\circ g\).  By (A3) and
\cref{plt:lem:cmp-realization-basic}, \(\log(1+u(X))\) realizes
\(V=\log(1+U)\) on \(S'\): indeed for any \(M\),
\(\log(1+u)=\sum_{r=1}^{M}\frac{(-1)^{r-1}}{r}u^{r}
 +\BigOAt{\AbsoluteValue{u}^{M+1}}{X\to+\infty}\)
by Taylor's theorem with remainder for \(\log(1+\cdot)\) at \(0\), applicable
because \(u\to0\); each \(u^{r}\) realizes \(U^{r}\) by
\cref{plt:lem:cmp-realization-ring}, and
\(\AbsoluteValue{u}^{M+1}=\BigOAt{X^{-M-1+\varepsilon}}{X\to+\infty}\).
The same argument with the binomial series shows that
\((1+u(X))^{-n}\) realizes \((1+U)^{-n}\) for each \(n\).  Consequently, for
each \(n<N\),
\[
 p_n\bigl(\log g(X)\bigr)g(X)^{-n}
 =p_n\bigl(\rho\log X+\log c+\log(1+u(X))\bigr)\,
   c^{-n}X^{-\rho n}(1+u(X))^{-n},
\]
which by (A3), the finite expansion \eqref{plt:eq:cmp-poly-shift} and
\cref{plt:lem:cmp-realization-ring} realizes
\(p_n(\rho L+\log c+V)c^{-n}t^{\rho n}(1+U)^{-n}\) on \(S'\).  Summing the
finitely many \(n<N\) and using
\eqref{plt:eq:cmp-remainder-transport} gives the claim, because by
\cref{plt:thm:cmp-finite-determination} the blocks of \(F\compose G\) of
exponent \(<\beta\) come only from those \(n\) with \(\rho n<\beta\le\rho N\).
\end{proof}

\begin{example}[Each hypothesis is used]\label{plt:ex:cmp-analytic-failures}
\begin{enumerate}
\item \emph{(A1)/(A2) with \(\rho>0\) are not decorative.}  Let
\(f(Y)=\EulerE^{-Y}\).  Then \(f\sim_{S}0\) on any ray: the zero series is a
correct realization to all orders.  Substituting \(g(X)=\log X\), which has no
positive power of \(X\) as leading monomial, gives
\(f(g(X))=X^{-1}\), which realizes \(t\ne0\).  Substituting
\(g(X)=X^{-1}\to0\), which violates (A1), gives
\(f(g(X))\to1\), which realizes the series \(1\ne0\).  So composition is not
even well defined on realizations without (A1)--(A2); what makes the theorem
work is that \(\rho>0\) transports flatness to flatness.
\item \emph{(A3) is a hypothesis, not a consequence.}  On a complex sector
containing \(\arg X\) close to \(\pi\), the principal logarithm satisfies
\(\PrincipalLogarithm g(X)=\rho\PrincipalLogarithm X+\log c
 +\PrincipalLogarithm(1+u)\) only while
\(\AbsoluteValue{\rho\arg X+\arg(1+u)}<\pi\); crossing that threshold changes
the answer by \(2\pi\ImaginaryUnit\) in every coefficient block of positive
degree.  On the positive real ray with \(c>0\), (A3) is automatic.
\item \emph{Precision transport is lossy for \(\rho<1\).}  By
\eqref{plt:eq:cmp-remainder-transport}, an output valid through exponent
\(\beta\) requires the outer expansion through outer order \(N>\beta/\rho\).
For \(\rho=1/2\) one must double the outer precision; for \(\rho=2\) one gains.
\end{enumerate}
\end{example}

\begin{warningbox}
\textbf{Three separate obligations.}  A formal identity
\(F\compose G\) established by \cref{plt:thm:cmp-substitution-hom},
\cref{plt:thm:cmp-finite-determination} or
\cref{plt:thm:cmp-associativity} carries no
analytic content.  To claim an asymptotic statement one must additionally
exhibit (A1) a domain on which the inner map lands in the outer expansion's
regime, (A2) a monomial-leading realization of the inner map with a positive
exponent, and (A3) a branch identity for every logarithm and fractional power
used.  Differentiating a realization requires yet more: a Poincar\'e expansion
cannot in general be differentiated termwise without sectorial analyticity or
direct derivative estimates \cite{olver-asymptotics,debruijn}.  The formal
derivation \(\DX\) is always defined; its analytic interpretation is not.
\end{warningbox}


\chapter{The formal Taylor formula}
\label{plt:sec:taylor}

For near-identity inner arguments there is a second description of composition
that looks exactly like the classical Taylor expansion but is an exact
identity of formal series.  We prove it from scratch, because the summability
and the multiplicativity of the Taylor operator are both nontrivial and are
the two points at which the sources are thin.

\section{The normal-ordered Taylor operator}
\label{plt:sec:tay-operator}

\begin{definition}[Taylor operator]\label{plt:def:tay-operator}
For \(H\in\PLpow\) define, for \(F\in\PLlau\),
\begin{equation}\label{plt:eq:tay-operator}
 \mathcal T_{H}(F)\DefinitionEquals
 \sum_{k\ge0}\frac{H^{k}}{k!}\,\DX^{k}F .
\end{equation}
\end{definition}

\begin{lemma}[Summability, with an explicit derivative budget]
\label{plt:lem:tay-summability}
Let \(H\in\PLpow\) and \(F\in\PLlau\) with \(a=\ordt(F)\).  Then
\begin{equation}\label{plt:eq:tay-valuation-bound}
 \ordt\Bigl(\frac{H^{k}}{k!}\DX^{k}F\Bigr)
 \ge a+k+k\,\ordt(H)\ \ge\ a+k,
\end{equation}
so the family in \eqref{plt:eq:tay-operator} is summable and
\(\mathcal T_{H}(F)\in\PLlau\) with
\(\ordt(\mathcal T_{H}F)\ge a\).  Consequently, for every
\(N\in\IntegerNumbers\),
\begin{equation}\label{plt:eq:tay-truncation}
 \Trunc{\mathcal T_{H}(F)}{N}
 =\Trunc{\Bigl(\sum_{k=0}^{N-a-1}\frac{H^{k}}{k!}\DX^{k}F\Bigr)}{N},
\end{equation}
that is, only the indices \(k\le N-a-1\) can affect any block below \(t^{N}\),
and only \(k\le N-a\) can affect the block \(t^{N}\) itself.  If moreover
\(\ordt(H)\ge1\), the budget halves: only \(k\le(N-a)/2\) can affect
\(t^{N}\).
\end{lemma}

\begin{proof}
Write \(F=\sum_{n\ge a}p_n(L)t^n\).  By \eqref{plt:eq:cmp-shiftop} and
\cref{plt:lem:cmp-summable} applied to the continuous operator \(\DX\),
\(\DX^{k}F=\sum_{n\ge a}t^{n+k}\shiftop{n}{k}p_n\), so
\(\ordt(\DX^{k}F)\ge a+k\); the inequality may be strict, and is so exactly
when the operators \(\shiftop{n}{k}\) annihilate the leading blocks, but the
bound is all we need.  Multiplying by \(H^{k}\), whose valuation is
\(k\,\ordt(H)\ge0\), gives \eqref{plt:eq:tay-valuation-bound}.  The right-hand
side tends to \(+\infty\) with \(k\), which is
\eqref{plt:eq:cmp-summability-criterion}, so
\cref{plt:lem:cmp-summable} applies.  A summand contributes to the block
\(t^{N}\) only if \(a+k\le N\), and to a block strictly below \(t^{N}\) only
if \(a+k\le N-1\).  If \(\ordt(H)\ge1\), \eqref{plt:eq:tay-valuation-bound}
improves to \(a+2k\le N\).
\end{proof}

\begin{lemma}[The Taylor operator is a homomorphism]
\label{plt:lem:tay-homomorphism}
For every \(H\in\PLpow\), \(\mathcal T_{H}:\PLlau\to\PLlau\) is a
\(\K\)-algebra homomorphism compatible with summable families.
\end{lemma}

\begin{proof}
\(\K\)-linearity is clear from \eqref{plt:eq:tay-operator}, and
\(\mathcal T_{H}(1)=1\) because \(\DX1=0\).

% ed.: as in (plt:thm:cmp-substitution-hom), the earlier text checked the
% interchange only for the family of blocks of a single F; the definition of
% compatibility quantifies over all summable families, and that is what
% (plt:lem:tay-determined) consumes.  Argument supplied.
For compatibility with summable families, the valuation half holds with
\(\varphi=\mathrm{id}\), since
\(\ordt(\mathcal T_{H}F)\ge\ordt(F)\) by
\eqref{plt:eq:tay-valuation-bound}.  For the interchange, observe first that
\(\DX\) itself is compatible with summable families: it is \(\K\)-linear,
raises valuation by at least \(1\), and by \eqref{plt:eq:cmp-shiftop} the
block \([t^{n+1}]\DX F\) depends \(\K\)-linearly on the single block
\([t^{n}]F\), so \(\DX\bigl(\sum_iS_i\bigr)=\sum_i\DX S_i\) for every
summable family; the same then holds for each \(\DX^{k}\).  Now let
\((S_i)_{i\in I}\) be summable with sum \(S\).  By
\cref{plt:rmk:cmp-lower-bound-automatic} the valuations \(\ordt(S_i)\) are
bounded below, so the doubly indexed family
\(\bigl(\tfrac{H^{k}}{k!}\DX^{k}S_i\bigr)_{i,k}\), whose \((i,k)\) member has
valuation at least \(\ordt(S_i)+k\) by
\eqref{plt:eq:tay-valuation-bound}, satisfies
\eqref{plt:eq:cmp-summability-criterion}.  Summing it with \(i\) innermost
gives \(\mathcal T_{H}(S)\) and with \(k\) innermost gives
\(\sum_i\mathcal T_{H}(S_i)\).

For multiplicativity, first record the Leibniz rule for the derivation
\(\DX\): by induction on \(m\),
\(\DX^{m}(F_1F_2)=\sum_{i=0}^{m}\binom{m}{i}\DX^{i}F_1\,\DX^{m-i}F_2\).
Hence
\[
 \mathcal T_{H}(F_1F_2)
 =\sum_{m\ge0}\frac{H^{m}}{m!}
   \sum_{i=0}^{m}\binom{m}{i}\DX^{i}F_1\,\DX^{m-i}F_2
 =\sum_{m\ge0}\sum_{i=0}^{m}
   \frac{H^{i}\DX^{i}F_1}{i!}\cdot
   \frac{H^{m-i}\DX^{m-i}F_2}{(m-i)!} .
\]
The doubly indexed family
\(\bigl(\tfrac{H^{i}\DX^{i}F_1}{i!}\cdot
 \tfrac{H^{j}\DX^{j}F_2}{j!}\bigr)_{i,j\ge0}\)
has \(\ordt\ge\ordt(F_1)+\ordt(F_2)+i+j\) by
\eqref{plt:eq:tay-valuation-bound}, hence is summable; regrouping it by
\(i+j=m\) recovers the display, and regrouping it as a product of two
summable families gives \(\mathcal T_{H}(F_1)\mathcal T_{H}(F_2)\).  Both
regroupings are licensed by \cref{plt:lem:cmp-summable}.
\end{proof}

% ed.: this multiplicativity is exactly the step that S3 asserts ("both sides
% are continuous algebra homomorphisms") and that S6 replaces by "the equality
% then extends to polynomials, products, inverses"; it is not automatic,
% because T_H is a *normal-ordered* exponential and not exp(H D_X).
\begin{remark}
\Cref{plt:lem:tay-homomorphism} is not a formality.  The operator
\(\mathcal T_{H}\) is the normal-ordered exponential
\(\mathopen{:}\EulerE^{H\DX}\mathclose{:}\), in which every power of \(H\) is
placed to the left and is never differentiated.  It is \emph{not} the operator
exponential \(\exp(M_{H}\DX)\), since
\((M_{H}\DX)^{2}=M_{H}M_{\DX H}\DX+M_{H^{2}}\DX^{2}\).  The latter is also an
automorphism, but a different one.  What makes
\eqref{plt:eq:tay-operator} multiplicative is the Leibniz rule combined with
the elementary identity \(\binom{m}{i}/m!=1/(i!\,(m-i)!)\), not any operator
exponential formalism.
\end{remark}

\section{The formal Taylor formula}
\label{plt:sec:tay-formula}

\begin{lemma}[Generators determine such a homomorphism]
\label{plt:lem:tay-determined}
Let \(\mathcal A\) be a polynomial-logarithmic Hahn ring and let
\(\Psi_{1},\Psi_{2}:\PLlau\to\mathcal A\) be \(\K\)-algebra homomorphisms
compatible with summable families.  If \(\Psi_{1}(t)=\Psi_{2}(t)\) and
\(\Psi_{1}(L)=\Psi_{2}(L)\), then \(\Psi_{1}=\Psi_{2}\).
\end{lemma}

\begin{proof}
Both fix \(\K\) and agree on \(L\), hence agree on \(\K[L]\); they agree on
\(t\), hence on \(t^{n}\) for \(n\ge0\), and on \(t^{-1}\) because
\(\Psi_{i}(t)\Psi_{i}(t^{-1})=1\) forces
\(\Psi_{i}(t^{-1})=\Psi_{i}(t)^{-1}\).  So they agree on every
\(p_n(L)t^{n}\).  For \(F=\sum_np_nt^n\) the family
\((p_n(L)t^n)_n\) is summable, so
\(\Psi_{i}(F)=\sum_n\Psi_{i}(p_nt^n)\), and the two sums coincide.
\end{proof}

\begin{theorem}[Formal Taylor formula]\label{plt:thm:tay-formula}
For every \(F\in\PLlau\) and every \(H\in\PLpow\),
\begin{equation}\label{plt:eq:cmp-taylor}
 F\compose(X+H)
 =\sum_{k\ge0}\frac{H^{k}}{k!}\,\DX^{k}F .
\end{equation}
This is an exact identity in \(\PLlau\), not an approximation: both sides have
the same coefficient in every block.  If \(a=\ordt(F)\), the coefficient of
\(t^{N}\) on the right receives contributions only from \(k\le N-a\).
\end{theorem}

\begin{proof}
By \cref{plt:prop:cmp-generator-substitution} the map
\(\Phi_{X+H}\) is a \(\K\)-algebra homomorphism compatible with summable
families, and by \cref{plt:lem:tay-homomorphism} so is
\(\mathcal T_{H}\).  By \cref{plt:lem:tay-determined} it suffices to compare
them on \(t\) and \(L\).

For \(F=t\): from \(\DX t=-t^{2}\) and
\eqref{plt:eq:cmp-shiftop} with \(n=1\), \(p=1\), one gets
\(\DX^{k}t=t^{1+k}\prod_{j=0}^{k-1}(-1-j)=(-1)^{k}k!\,t^{k+1}\).  Hence
\[
 \mathcal T_{H}(t)
 =\sum_{k\ge0}\frac{H^{k}}{k!}(-1)^{k}k!\,t^{k+1}
 =t\sum_{k\ge0}(-tH)^{k}
 =\frac{t}{1+tH}
 =\Phi_{X+H}(t),
\]
the geometric family being summable because \(\ordt(tH)\ge1\).

For \(F=L\): \(\DX L=t\), so
\(\DX^{k}L=\DX^{k-1}t=(-1)^{k-1}(k-1)!\,t^{k}\) for \(k\ge1\).  Hence
\[
 \mathcal T_{H}(L)
 =L+\sum_{k\ge1}\frac{H^{k}}{k!}(-1)^{k-1}(k-1)!\,t^{k}
 =L+\sum_{k\ge1}\frac{(-1)^{k-1}}{k}(tH)^{k}
 =L+\log(1+tH)
 =\Phi_{X+H}(L),
\]
by \eqref{plt:eq:cmp-exp-log-def}.  The derivative budget is
\cref{plt:lem:tay-summability}.
\end{proof}

\begin{remark}[The small parameter is \(H/X\), not \(H\)]
\label{plt:rmk:tay-small-parameter}
No smallness of \(H\) is assumed, and none is available: \(H\) may be
\(L^{100}\), which tends to infinity.  What the proof uses is
\(H/X=tH\precdom1\), which is automatic for \(H\in\PLpow\) because every
polynomial in \(\log X\) is \(\LittleOAt{X}{X\to+\infty}\).  The mechanism is
that \(\DX\) raises the outer order by one while multiplication by \(H\) does
not lower it.  The base point is the asymptotic variable itself, and the
conclusion is an identity, not an estimate; there is no analytic
differentiability hypothesis anywhere.
\end{remark}

\begin{remark}[Where the mechanism fails]\label{plt:rmk:tay-failure}
In a larger transseries field, "\(\Delta\) small compared with \(X\)" is not
enough.  Take \(F(X)=\EulerE^{\EulerE^{X}}\) and
\(\Delta=\EulerE^{-X/2}\), so \(\Delta\precdom X\).  Then
\[
 \Delta\,\DX F
 =\EulerE^{-X/2}\EulerE^{X}\EulerE^{\EulerE^{X}}
 =\EulerE^{X/2}F\succdom F,
\]
so the \(k=1\) Taylor term dominates the \(k=0\) term and the family is not
summable.  The quantity that must be small is not \(\Delta\) but
\(\Delta\cdot\DX\mathfrak b/\mathfrak b\), where \(\mathfrak b\) is the
dominant monomial of \(F\): the increment multiplied by the logarithmic
derivative of the outer scale.  In the polynomial-logarithmic class this is
automatically small, because for \(\mathfrak b=t^{n}L^{j}\) one has
\(\DX\mathfrak b/\mathfrak b=t\bigl(jL^{-1}-n\bigr)\precdom1\), whereas for
\(\mathfrak b=\EulerE^{\EulerE^{X}}\) one has
\(\DX\mathfrak b/\mathfrak b=\EulerE^{X}\succdom1\).  Sharper Taylor
approximation theorems for larger classes require exactly such a
logarithmic-derivative hypothesis
\cite{edgar-composition,mantova-monotonicity}.
\end{remark}

% ed.: S4 and S5 prove the Taylor formula by "inserting a scalar parameter s,
% forming F(x+sA), and differentiating in s".  That argument is circular: the
% object F(x+sA) is defined only by the composition being constructed, and
% "formal differentiation in s" of it presupposes the chain rule for that
% composition.  The homomorphism argument above avoids the circle.
\begin{remark}[On a tempting but circular proof]
One sometimes sees the Taylor formula deduced by introducing a scalar
parameter \(s\), forming \(\Phi(s)=F(X+sH)\) "with \(H\) held constant", and
expanding \(\Phi\) in powers of \(s\) at \(s=0\).  This is not a proof: the
symbol \(F(X+sH)\) has no meaning prior to the composition one is trying to
compute, and the claim
\(\partial_{s}^{k}\Phi(0)=H^{k}\DX^{k}F\) is the chain rule for that same
composition.  One can rescue the argument by working in
\(\PLlau[[s]]\) and \emph{defining} \(\Phi(s)=\Phi_{X+sH}(F)\), but then one
must still prove that \(\partial_{s}\Phi=H\cdot\Phi_{X+sH}(\DX F)\), which is
\cref{plt:thm:cmp-chain-rule}.  The route through
\cref{plt:lem:tay-homomorphism} and \cref{plt:lem:tay-determined} is shorter
and self-contained.
\end{remark}

\section{The finite coefficient formula}
\label{plt:sec:tay-coefficients}

\begin{theorem}[Finite coefficient formula]\label{plt:thm:tay-coefficients}
Let \(F=\sum_{n\ge a}p_n(L)t^{n}\in\PLlau\) and \(H\in\PLpow\), and write
\begin{equation}\label{plt:eq:tay-Hk-blocks}
 H^{k}=\sum_{m\ge0}h_{k,m}(L)\,t^{m},
 \qquad
 h_{0,0}=1,\quad h_{0,m}=0\ (m>0).
\end{equation}
Then for every \(r\in\IntegerNumbers\)
\begin{equation}\label{plt:eq:tay-coefficient-formula}
 [t^{r}]\bigl(F\compose(X+H)\bigr)
 =\sum_{\substack{n\ge a,\ k\ge0,\ m\ge0\\ n+k+m=r}}
   \frac{h_{k,m}(L)}{k!}\,\shiftop{n}{k}p_n(L),
\end{equation}
a sum with exactly \(\binom{r-a+2}{2}\) index triples when \(r\ge a\), and
empty when \(r<a\).  If \(\ordt(H)\ge1\), the constraint \(m\ge k\) applies as
well, so \(n+2k\le r\) and roughly half the triples drop out.
\end{theorem}

\begin{proof}
By \cref{plt:thm:tay-formula} and \eqref{plt:eq:cmp-shiftop},
\[
 F\compose(X+H)
 =\sum_{k\ge0}\frac{H^{k}}{k!}\sum_{n\ge a}t^{n+k}\shiftop{n}{k}p_n
 =\sum_{k\ge0}\sum_{n\ge a}\sum_{m\ge0}
   \frac{h_{k,m}}{k!}\,\shiftop{n}{k}p_n\;t^{n+k+m},
\]
the triple family being summable because its \((n,k,m)\) member has
valuation at least \(n+k+m\).  Extracting \([t^{r}]\) gives
\eqref{plt:eq:tay-coefficient-formula}.  The affine constraint
\(n+k+m=r\) with \(n\ge a\) and \(k,m\ge0\) has exactly
\(\binom{r-a+2}{2}\) solutions.  If \(\ordt(H)\ge1\) then
\(\ordt(H^{k})\ge k\), so \(h_{k,m}=0\) for \(m<k\).
\end{proof}

\begin{remark}[Three separable pieces of work]
Formula \eqref{plt:eq:tay-coefficient-formula} factors the computation into
(a) the powers \(H^{k}\), computed once by truncated multiplication and reused
for every outer series; (b) the shifted differential operators
\(\shiftop{n}{k}\) acting on coefficient polynomials, which depend only on
\(F\); and (c) a finite index convolution.  It is therefore the formula of
choice when one inner correction \(H\) is composed with many outer series, and
the generator form \eqref{plt:eq:cmp-near-identity-master} is the formula of
choice when a single sparse \(F\) is composed once.
\end{remark}

\section{Convolution coefficients}
\label{plt:sec:tay-convolution}

The blocks \(h_{k,m}\) of \eqref{plt:eq:tay-Hk-blocks} deserve a name and a
recurrence.  Write
\begin{equation}\label{plt:eq:tay-H-blocks}
 H=\sum_{j\ge0}h_{j}(L)\,t^{j},
 \qquad h_j\in\K[L].
\end{equation}

\begin{definition}[Convolution coefficients]\label{plt:def:tay-convolution}
For \(m,r\in\NaturalNumbers\) put
\begin{equation}\label{plt:eq:tay-C-def}
 C_{m,r}
 \DefinitionEquals
 \sum_{\substack{j_{1},\ldots,j_{r}\ge0\\ j_{1}+\cdots+j_{r}=m}}
   h_{j_{1}}\cdots h_{j_{r}}\ \in\K[L].
\end{equation}
The value at \(r=0\) is an empty product summed over the empty composition, so
\begin{equation}\label{plt:eq:tay-C-boundary}
 C_{0,0}=1,
 \qquad
 C_{m,0}=0\quad(m>0).
\end{equation}
\end{definition}

% ed.: the two index orders in the corpus are transposes of each other; the
% collision is silent because both symbols are two-index arrays of polynomials.
\begin{warningbox}
\textbf{Index order.}  With \eqref{plt:eq:tay-Hk-blocks} and
\eqref{plt:eq:tay-C-def} one has
\[
 h_{k,m}=C_{m,k}:
\]
in \(h_{k,m}\) the \emph{first} index is the exponent of \(H\) and the second
is the outer order; in \(C_{m,r}\) it is the other way round.  Both notations
occur in the literature.  Any transfer between them must transpose.
\end{warningbox}

\begin{proposition}[Structure of the convolution coefficients]
\label{plt:prop:tay-convolution}
With the conventions \eqref{plt:eq:tay-C-boundary}:
\begin{enumerate}
\item \(H^{r}=\sum_{m\ge0}C_{m,r}\,t^{m}\) for every \(r\in\NaturalNumbers\).
\item \(C_{m,1}=h_{m}\), and for all \(m,r\ge0\)
\begin{equation}\label{plt:eq:tay-C-recurrence}
 C_{m,r+1}=\sum_{j=0}^{m}h_{j}\,C_{m-j,r}.
\end{equation}
\item If \(h_{0}=0\), then \(C_{m,r}=0\) whenever \(m<r\), and
\(C_{r,r}=h_{1}^{r}\).
\item If \(h_{0}=0\), then for \(m\ge r\ge1\)
\begin{equation}\label{plt:eq:tay-C-bell}
 C_{m,r}
 =\frac{r!}{m!}\,
  \ExponentialPartialBellPolynomial{m}{r}
   \bigl(1!\,h_{1},\,2!\,h_{2},\ldots,(m-r+1)!\,h_{m-r+1}\bigr).
\end{equation}
\end{enumerate}
\end{proposition}

\begin{proof}
(i) Expanding \(H^{r}=\bigl(\sum_{j}h_jt^{j}\bigr)^{r}\) and collecting the
total outer order \(m\) is exactly \eqref{plt:eq:tay-C-def}; for \(r=0\) the
statement is \(H^{0}=1\), which is
\eqref{plt:eq:tay-C-boundary}.  The rearrangement is legitimate because the
family indexed by \((j_{1},\ldots,j_{r})\) has valuation
\(j_{1}+\cdots+j_{r}\) and is summable.

(ii) Take \(r=1\) in (i).  For the recurrence, write
\(H^{r+1}=H\cdot H^{r}\) and read off the coefficient of \(t^{m}\) using (i);
the inner index is bounded by \(m\) because \(C_{m',r}\) is \(0\) for
\(m'<0\).  At \(r=0\) the recurrence returns
\(C_{m,1}=h_{m}C_{0,0}=h_{m}\), consistent with
\eqref{plt:eq:tay-C-boundary}.

(iii) If \(h_{0}=0\) then \(\ordt(H)\ge1\), so \(\ordt(H^{r})\ge r\) and (i)
forces \(C_{m,r}=0\) for \(m<r\).  For \(m=r\) the only composition of \(r\)
into \(r\) positive parts is \((1,\ldots,1)\), giving \(h_{1}^{r}\).

(iv) With \(h_{0}=0\), set \(\hat h_{j}=j!\,h_{j}\).  The partial exponential
Bell polynomials are defined by the generating identity
\[
 \sum_{m\ge r}\ExponentialPartialBellPolynomial{m}{r}
   (\hat h_{1},\hat h_{2},\ldots)\frac{z^{m}}{m!}
 =\frac{1}{r!}\Bigl(\sum_{j\ge1}\hat h_{j}\frac{z^{j}}{j!}\Bigr)^{r}
 =\frac{1}{r!}\Bigl(\sum_{j\ge1}h_{j}z^{j}\Bigr)^{r}
\]
\cite{comtet-book}.  Substituting \(z\mapsto t\), which is legitimate
coefficientwise since both sides are formal power series in one variable with
coefficients in \(\K[L]\), and comparing with (i) gives
\(\frac{C_{m,r}}{r!}=\frac{1}{m!}\ExponentialPartialBellPolynomial{m}{r}(\hat h)\),
which is \eqref{plt:eq:tay-C-bell}.
\end{proof}

\begin{example}[The first convolution coefficients]\label{plt:ex:tay-C-table}
For \(H=h_{0}+h_{1}t+h_{2}t^{2}+h_{3}t^{3}+\FormalBigOAt{t^{4}}{t}\):
\begin{center}
\begin{tabular}{@{}lllll@{}}
\hline
 & \(m=0\) & \(m=1\) & \(m=2\) & \(m=3\)\\
\hline
\(r=0\) & \(1\) & \(0\) & \(0\) & \(0\)\\
\(r=1\) & \(h_{0}\) & \(h_{1}\) & \(h_{2}\) & \(h_{3}\)\\
\(r=2\) & \(h_{0}^{2}\) & \(2h_{0}h_{1}\)
        & \(h_{1}^{2}+2h_{0}h_{2}\)
        & \(2h_{0}h_{3}+2h_{1}h_{2}\)\\
\(r=3\) & \(h_{0}^{3}\) & \(3h_{0}^{2}h_{1}\)
        & \(3h_{0}^{2}h_{2}+3h_{0}h_{1}^{2}\)
        & \(h_{1}^{3}+6h_{0}h_{1}h_{2}+3h_{0}^{2}h_{3}\)\\
\hline
\end{tabular}
\end{center}
Each entry is obtained from the row above by
\eqref{plt:eq:tay-C-recurrence}.  Setting \(h_{0}=0\) empties the lower-left
triangle, as \cref{plt:prop:tay-convolution}(iii) predicts.
\end{example}

\section{Low-order composition laws}
\label{plt:sec:tay-low-order}

We now write out \eqref{plt:eq:tay-coefficient-formula} for the case that
governs reversion: both maps tangent to the identity, through four outer
blocks.

\begin{proposition}[Composition through four blocks]
\label{plt:prop:tay-four-blocks}
Let
\begin{align}
 F&=X+f_{0}(L)+f_{1}(L)t+f_{2}(L)t^{2}+f_{3}(L)t^{3}
    +\FormalBigOAt{t^{4}}{t},
 \label{plt:eq:tay-F-blocks}\\
 G&=X+g_{0}(L)+g_{1}(L)t+g_{2}(L)t^{2}+g_{3}(L)t^{3}
    +\FormalBigOAt{t^{4}}{t},
 \label{plt:eq:tay-G-blocks}
\end{align}
with all \(f_{j},g_{j}\in\K[L]\), and write
\(F\compose G=X+h_{0}+h_{1}t+h_{2}t^{2}+h_{3}t^{3}+\FormalBigOAt{t^{4}}{t}\).
Primes denote \(\partial_L\).  Then
\begin{align}
 h_{0}={}&f_{0}+g_{0},\label{plt:eq:tay-h0}\\
 h_{1}={}&f_{1}+g_{1}+g_{0}f_{0}',\label{plt:eq:tay-h1}\\
 h_{2}={}&f_{2}+g_{2}+g_{1}f_{0}'+g_{0}\bigl(f_{1}'-f_{1}\bigr)
   +\frac{g_{0}^{2}}{2}\bigl(f_{0}''-f_{0}'\bigr),\label{plt:eq:tay-h2}\\
 h_{3}={}&f_{3}+g_{3}+g_{2}f_{0}'+g_{1}\bigl(f_{1}'-f_{1}\bigr)
   +g_{0}\bigl(f_{2}'-2f_{2}\bigr)
   +g_{0}g_{1}\bigl(f_{0}''-f_{0}'\bigr)\notag\\
 &{}+\frac{g_{0}^{2}}{2}\bigl(f_{1}''-3f_{1}'+2f_{1}\bigr)
   +\frac{g_{0}^{3}}{6}\bigl(f_{0}'''-3f_{0}''+2f_{0}'\bigr).
 \label{plt:eq:tay-h3}
\end{align}
\end{proposition}

\begin{proof}
Put \(H=G-X=g_{0}+g_{1}t+g_{2}t^{2}+g_{3}t^{3}+\FormalBigOAt{t^{4}}{t}\), so
\(H\in\PLpow\) and \(\ordt(F)=-1\).  By
\cref{plt:lem:tay-summability} only \(k\le 3-(-1)=4\) can reach the block
\(t^{3}\); but \(\DX^{k}X=0\) for \(k\ge2\) and
\(\ordt(\DX^{k}(F-X))\ge k\), so in fact only \(k\le3\) matter.  We need the
first three derivatives of \(F\), which \eqref{plt:eq:cmp-shiftop} supplies
blockwise with \(\shiftop{n}{k}=\prod_{j=0}^{k-1}(\partial_L-n-j)\):
\begin{align*}
 \DX F&=1+t f_{0}'+t^{2}\bigl(f_{1}'-f_{1}\bigr)
        +t^{3}\bigl(f_{2}'-2f_{2}\bigr)+\FormalBigOAt{t^{4}}{t},\\
 \DX^{2}F&=t^{2}\bigl(f_{0}''-f_{0}'\bigr)
        +t^{3}\bigl(f_{1}''-3f_{1}'+2f_{1}\bigr)+\FormalBigOAt{t^{4}}{t},\\
 \DX^{3}F&=t^{3}\bigl(f_{0}'''-3f_{0}''+2f_{0}'\bigr)
        +\FormalBigOAt{t^{4}}{t},
\end{align*}
using \(\shiftop{0}{1}=\partial_L\), \(\shiftop{1}{1}=\partial_L-1\),
\(\shiftop{2}{1}=\partial_L-2\), \(\shiftop{0}{2}=\partial_L(\partial_L-1)\),
\(\shiftop{1}{2}=(\partial_L-1)(\partial_L-2)\) and
\(\shiftop{0}{3}=\partial_L(\partial_L-1)(\partial_L-2)\).  Also
\(H^{2}=g_{0}^{2}+2g_{0}g_{1}t+\FormalBigOAt{t^{2}}{t}\) and
\(H^{3}=g_{0}^{3}+\FormalBigOAt{t}{t}\).  Substituting into
\eqref{plt:eq:cmp-taylor},
\begin{align*}
 F\compose G
 ={}&\bigl(X+f_{0}+f_{1}t+f_{2}t^{2}+f_{3}t^{3}\bigr)\\
 &+\bigl(g_{0}+g_{1}t+g_{2}t^{2}+g_{3}t^{3}\bigr)
   \bigl(1+tf_{0}'+t^{2}(f_{1}'-f_{1})+t^{3}(f_{2}'-2f_{2})\bigr)\\
 &+\tfrac12\bigl(g_{0}^{2}+2g_{0}g_{1}t\bigr)
   \bigl(t^{2}(f_{0}''-f_{0}')+t^{3}(f_{1}''-3f_{1}'+2f_{1})\bigr)\\
 &+\tfrac16 g_{0}^{3}\,t^{3}\bigl(f_{0}'''-3f_{0}''+2f_{0}'\bigr)
   +\FormalBigOAt{t^{4}}{t}.
\end{align*}
Collecting the blocks \(t^{0},\ldots,t^{3}\) gives
\eqref{plt:eq:tay-h0}--\eqref{plt:eq:tay-h3}.
\end{proof}

\begin{remark}[Which derivatives can occur, and with which coefficients]
The formulas are triangular: derivatives fall only on the \emph{outer}
coefficients \(f_{j}\), never on the inner \(g_{j}\), because it is the outer
logarithmic argument that is Taylor-expanded while the inner coefficients only
supply the increment.  At block \(t^{r}\) the highest derivative of \(f_{n}\)
that occurs is of order \(r-n\), and the operator applied to \(f_{n}\) at
Taylor order \(k\) is exactly \(\shiftop{n}{k}\).  The integers appearing in
\eqref{plt:eq:tay-h2}--\eqref{plt:eq:tay-h3} are therefore signed Stirling
numbers of the first kind: since
\(\shiftop{n}{k}=\prod_{j=0}^{k-1}\bigl((\partial_L-n)-j\bigr)
 =\FallingFactorial{(\partial_L-n)}{k}\),
\begin{equation}\label{plt:eq:tay-shiftop-stirling}
 \shiftop{n}{k}
 =\sum_{i=0}^{k}\stirlingsigned{k}{i}\,(\partial_L-n)^{i}.
\end{equation}
For \(n=0,k=2\) this gives \(\partial_L^{2}-\partial_L\); for \(n=1,k=2\),
\(\partial_L^{2}-3\partial_L+2\); for \(n=0,k=3\),
\(\partial_L^{3}-3\partial_L^{2}+2\partial_L\), which are precisely the three
brackets displayed above.
\end{remark}

\begin{example}[Consistency check: a constant shift]\label{plt:ex:tay-shift}
Take \(G=X+c\) with \(c\in\K\), so \(g_{0}=c\) and \(g_{1}=g_{2}=g_{3}=0\).
Then \eqref{plt:eq:tay-h0}--\eqref{plt:eq:tay-h2} reduce to
\[
 h_{0}=f_{0}+c,\qquad
 h_{1}=f_{1}+cf_{0}',\qquad
 h_{2}=f_{2}+c\bigl(f_{1}'-f_{1}\bigr)+\frac{c^{2}}{2}\bigl(f_{0}''-f_{0}'\bigr),
\]
which is \eqref{plt:eq:cmp-taylor} read off directly, since
\(\DX f_{0}(L)=tf_{0}'\) and \(\DX^{2}f_{0}(L)=t^{2}(f_{0}''-f_{0}')\).
\end{example}

\begin{corollary}[First commutator block]\label{plt:cor:tay-commutator}
Let \(F=X+f(L)\) and \(G=X+g(L)\) with \(f,g\in\K[L]\).  Then
\begin{equation}\label{plt:eq:tay-commutator}
 F\compose G-G\compose F
 =\bigl(gf'-fg'\bigr)\,t+\FormalBigOAt{t^{2}}{t}.
\end{equation}
In particular composition is associative
(\cref{plt:thm:cmp-associativity}) but not commutative, and the leading
obstruction is the Lie bracket of the one-dimensional vector fields
\(f\partial_L\) and \(g\partial_L\).
\end{corollary}

\begin{proof}
Apply \eqref{plt:eq:tay-h0}--\eqref{plt:eq:tay-h1} with
\(f_{0}=f\), \(g_{0}=g\) and all other blocks zero: \(h_{0}=f+g\) is
symmetric and \(h_{1}=gf'\); exchanging the roles of \(F\) and \(G\) gives
\(fg'\).
\end{proof}

\begin{example}[Explicit noncommutativity]\label{plt:ex:tay-noncommutative}
For \(F=X+L\) and \(G=X+L^{2}\),
\begin{align*}
 F\compose G&=X+L^{2}+L+L^{2}t-\tfrac12L^{4}t^{2}+\FormalBigOAt{t^{3}}{t},\\
 G\compose F&=X+L+L^{2}+2L^{2}t+\bigl(L^{2}-L^{3}\bigr)t^{2}
   +\FormalBigOAt{t^{3}}{t},
\end{align*}
so \(F\compose G-G\compose F=-L^{2}t+\FormalBigOAt{t^{2}}{t}\), in agreement
with \eqref{plt:eq:tay-commutator}: \(gf'-fg'=L^{2}-2L^{2}=-L^{2}\).
\end{example}

% ed.: the source that supplies this example specifies F only modulo Y^{-3} and
% G only modulo t^2, yet quotes h_3, which depends on the coefficient of Y^{-3}
% in F and on the coefficients of t^2, t^3 in G.  The data are restated here as
% exact finite series, which is what the quoted answer actually assumes.
\begin{example}[A four-block composition with nontrivial coefficients]
\label{plt:ex:tay-four-block}
Let \(F\) and \(G\) be the \emph{exact} finite series
\[
 F(Y)=Y+\bigl((\log Y)^{2}-1\bigr)+\frac{2\log Y}{Y}+\frac{1}{Y^{2}},
 \qquad
 G(X)=X+(1-L)+\frac{L^{2}}{X},
\]
so that \(f_{0}=L^{2}-1\), \(f_{1}=2L\), \(f_{2}=1\), \(f_{3}=0\), and
\(g_{0}=1-L\), \(g_{1}=L^{2}\), \(g_{2}=g_{3}=0\).  Then
\cref{plt:prop:tay-four-blocks} gives
\[
 h_{0}=L^{2}-L,
 \quad
 h_{1}=-L^{2}+4L,
 \quad
 h_{2}=L^{3}+5L^{2}-7L+4,
 \quad
 h_{3}=\tfrac43L^{4}-L^{3}-8L^{2}+\tfrac{41}{3}L-6,
\]
that is,
\[
 F\compose G
 =X+\bigl(L^{2}-L\bigr)+\frac{-L^{2}+4L}{X}
  +\frac{L^{3}+5L^{2}-7L+4}{X^{2}}
  +\frac{\tfrac43L^{4}-L^{3}-8L^{2}+\tfrac{41}{3}L-6}{X^{3}}
  +\FormalBigOAt{t^{4}}{t}.
\]
Had \(F\) and \(G\) been specified only modulo \(Y^{-3}\) and \(t^{2}\), the
block \(h_{3}\) would not have been determined, since it contains
\(f_{3}+g_{3}+g_{2}f_{0}'\).
\end{example}

\begin{example}[One step of a reversion]\label{plt:ex:tay-lambert-step}
Let \(F(Y)=Y+\log Y\) and \(G(X)=X-L+LX^{-1}\), so \(f_{0}=L\),
\(f_{1}=f_{2}=f_{3}=0\), \(g_{0}=-L\), \(g_{1}=L\), \(g_{2}=g_{3}=0\).  Then
\(h_{0}=-L+L=0\) and \(h_{1}=f_1+g_{1}+g_{0}f_{0}'=0+L-L=0\), while
\[
 h_{2}=f_2+g_{2}+g_{1}f_{0}'+g_{0}(f_{1}'-f_{1})
   +\frac{g_{0}^{2}}{2}(f_{0}''-f_{0}')
 =0+0+L+0+\frac{L^{2}}{2}(0-1)
 =L-\frac{L^{2}}{2}.
\]
Adding \(\LamP{2}(L)=\tfrac12L^{2}-L\) to the \(X^{-2}\) block of \(G\)
cancels \(h_{2}\), because that block enters \(h_{2}\) additively through
\(g_{2}\).  The canonical Lambert polynomial \(\LamP{2}\) is thus recovered as
the unique second-order correction, exactly as the reversion theory predicts.
\end{example}

\section{Composing with an outer analytic function}
\label{plt:sec:tay-faa-di-bruno}

Two further compositions occur constantly and are governed by Bell
polynomials: substituting a polynomial-logarithmic series of positive
valuation into a scalar power series, and differentiating a composite
repeatedly.

% ed.: this proposition restated three results of the elementary-series part
% ed.: under a different normalisation of the outer coefficients: the
% ed.: substitution homomorphism (\cref{plt:prop:dif-substitution}), its
% ed.: finite determination (\cref{plt:thm:dif-finite-determination}), and the
% ed.: Bell-polynomial coefficient formula \eqref{plt:eq:dif-bell} of
% ed.: \cref{plt:prop:dif-bell}, which the displayed formula reproduces
% ed.: verbatim once \(c_{k}=\varphi_{k}/k!\) is substituted.  The
% ed.: restatement and its proof are deleted in favour of the pointer below;
% ed.: the label is kept so that \cref{plt:ex:tay-standard-outer} and any
% ed.: citation from elsewhere in the volume continue to resolve.  Only the
% ed.: translation clause was genuinely new, and it is retained with its
% ed.: one-line proof, as is the second route to the coefficient formula.
\begin{remark}[Substitution into a scalar power series: where it is proved]
\label{plt:prop:tay-scalar-outer}
Let \(\varphi(z)=\sum_{k\ge0}\frac{\varphi_{k}}{k!}z^{k}\in\K[[z]]\) and let
\(W=\sum_{j\ge1}w_{j}(L)t^{j}\in t\PLpow\), so that \(\ordt W\ge1\).
Everything used below about
\(\varphi(W)\DefinitionEquals\sum_{k\ge0}\frac{\varphi_{k}}{k!}W^{k}\) is
proved in the elementary-series part, applied to
\(\Phi=\sum_{r\ge0}c_{r}z^{r}\) with \(c_{r}=\varphi_{r}/r!\) and \(U=W\).
The family is summable, \(\varphi(W)\in\PLpow\), and
\(\varphi\mapsto\varphi(W)\) is the unique continuous \(\K\)-algebra
homomorphism \(\K[[z]]\to\PLpow\) sending \(z\) to \(W\)
(\cref{plt:prop:dif-substitution}); only \(\varphi_{0},\ldots,\varphi_{r}\)
are needed for the block \(t^{r}\)
(\cref{plt:thm:dif-finite-determination}); and
\[
 [t^{0}]\varphi(W)=\varphi_{0},
 \qquad
 [t^{r}]\varphi(W)
 =\frac{1}{r!}\sum_{k=1}^{r}\varphi_{k}\,
   \ExponentialPartialBellPolynomial{r}{k}
   \bigl(1!\,w_{1},2!\,w_{2},\ldots,(r-k+1)!\,w_{r-k+1}\bigr)
 \quad(r\ge1),
\]
which is \eqref{plt:eq:dif-bell} of \cref{plt:prop:dif-bell} after that
normalisation.  Two additions to what is proved there.  First, if \(\varphi\)
is analytic at \(z_{0}\) with Taylor coefficients in \(\K\) and
\(W-z_{0}\in t\PLpow\), then the same three statements hold after
translation: apply them to
\(\psi(z)\DefinitionEquals\varphi(z_{0}+z)\in\K[[z]]\) and to \(W-z_{0}\).
Second, the coefficient formula also follows from the convolution structure
established in this part, without the generating identity for the partial Bell
polynomials: \cref{plt:prop:tay-convolution}(iv), applied to \(W\) in place of
\(H\) (whose block \(w_{0}\) vanishes), gives
\([t^{r}]W^{k}
=\frac{k!}{r!}\ExponentialPartialBellPolynomial{r}{k}(1!w_{1},\ldots)\) for
\(r\ge k\ge1\) and \(0\) for \(r<k\); multiplying by \(\varphi_{k}/k!\) and
summing over \(k\le r\) reproduces the display.
\end{remark}

\begin{example}\label{plt:ex:tay-standard-outer}
Taking \(\varphi(z)=(1+z)^{-n}\), \(\varphi(z)=\log(1+z)\) and
\(\varphi(z)=\EulerE^{z}\) in
\cref{plt:prop:tay-scalar-outer} recovers, respectively, the binomial factor
\((1+U)^{-n}\), the logarithmic increment \(V=\log(1+U)\), and
\(\EulerE^{V}=1+U\) used throughout \cref{plt:sec:composition}.  For the
exponential one may replace the partial by the complete Bell polynomial:
\(\EulerE^{W}=\sum_{r\ge0}\frac{1}{r!}
 \ExponentialCompleteBellPolynomial{r}(1!w_{1},\ldots,r!w_{r})\,t^{r}\),
with \(\ExponentialCompleteBellPolynomial{0}=1\).
\end{example}

For repeated differentiation of a composite the relevant identity is Fa\`a di
Bruno's.  We prove it in the exact generality needed, from a recurrence for
the partial Bell polynomials which we also prove.

\begin{lemma}[Shift recurrence for partial Bell polynomials]
\label{plt:lem:tay-bell-recurrence}
Let \(\mathcal D=\sum_{i\ge1}x_{i+1}\partial_{x_{i}}\) act on polynomials in
the indeterminates \(x_{1},x_{2},\ldots\) over \(\RationalNumbers\).  Then for
all \(n\ge0\) and all \(k\ge1\), with the convention
\(\ExponentialPartialBellPolynomial{n}{k}=0\) for \(k>n\),
\begin{equation}\label{plt:eq:tay-bell-recurrence}
 \ExponentialPartialBellPolynomial{n+1}{k}
 =\mathcal D\,\ExponentialPartialBellPolynomial{n}{k}
  +x_{1}\,\ExponentialPartialBellPolynomial{n}{k-1},
\end{equation}
with \(\ExponentialPartialBellPolynomial{n}{0}=0\) for \(n\ge1\) and
\(\ExponentialPartialBellPolynomial{0}{0}=1\).
\end{lemma}

\begin{proof}
Let \(\Xi(z)=\sum_{m\ge1}x_{m}z^{m}/m!\), so that
\(\sum_{n\ge k}\ExponentialPartialBellPolynomial{n}{k}z^{n}/n!=\Xi^{k}/k!\)
is the defining generating identity \cite{comtet-book}.  Applying
\(\mathcal D\) coefficientwise,
\(\mathcal D\Xi=\sum_{i\ge1}x_{i+1}z^{i}/i!=\Xi'(z)-x_{1}\), because
\(\Xi'(z)=\sum_{i\ge0}x_{i+1}z^{i}/i!\).  Since \(\mathcal D\) is a derivation
that does not act on \(z\),
\[
 \mathcal D\Bigl(\frac{\Xi^{k}}{k!}\Bigr)
 =\frac{\Xi^{k-1}}{(k-1)!}\,\mathcal D\Xi
 =\frac{\Xi^{k-1}}{(k-1)!}\Xi'
  -x_{1}\frac{\Xi^{k-1}}{(k-1)!}
 =\frac{\Differential}{\Differential z}\Bigl(\frac{\Xi^{k}}{k!}\Bigr)
  -x_{1}\frac{\Xi^{k-1}}{(k-1)!} .
\]
Comparing coefficients of \(z^{n}/n!\) and using
\(\frac{\Differential}{\Differential z}\sum_n a_n z^{n}/n!
 =\sum_n a_{n+1}z^{n}/n!\) gives
\(\mathcal D\ExponentialPartialBellPolynomial{n}{k}
 =\ExponentialPartialBellPolynomial{n+1}{k}
  -x_{1}\ExponentialPartialBellPolynomial{n}{k-1}\)
for all \(n\ge0\) and \(k\ge1\); when \(k>n+1\) both sides vanish, and when
\(k=n+1\) the identity reads \(0=x_{1}^{n+1}-x_{1}\cdot x_{1}^{n}\).
\end{proof}

\begin{proposition}[Fa\`a di Bruno for admissible composition]
\label{plt:prop:tay-faa-di-bruno}
Let \(G\) be an admissible monomial-leading argument with \(\rho>0\), and let
\(\gamma\DefinitionEquals\DX G\).  Then for every \(F\in\PLlau\) and every
\(n\ge1\),
\begin{equation}\label{plt:eq:tay-faa-di-bruno}
 \DX^{n}\bigl(F\compose G\bigr)
 =\sum_{k=1}^{n}\bigl(\DX^{k}F\bigr)\compose G\;
  \ExponentialPartialBellPolynomial{n}{k}
   \bigl(\gamma,\DX\gamma,\ldots,\DX^{\,n-k}\gamma\bigr).
\end{equation}
\end{proposition}

\begin{proof}
Induction on \(n\).  For \(n=1\) this is
\cref{plt:thm:cmp-chain-rule}, since
\(\ExponentialPartialBellPolynomial{1}{1}=x_{1}\).  Assume
\eqref{plt:eq:tay-faa-di-bruno} for \(n\).  Apply \(\DX\), using the product
rule, the chain rule \eqref{plt:eq:cmp-chain-rule} in the form
\(\DX\bigl((\DX^{k}F)\compose G\bigr)=(\DX^{k+1}F)\compose G\cdot\gamma\), and
the fact that
\(\DX\bigl[\ExponentialPartialBellPolynomial{n}{k}(\gamma,\DX\gamma,\ldots)\bigr]
 =\bigl(\mathcal D\ExponentialPartialBellPolynomial{n}{k}\bigr)
   (\gamma,\DX\gamma,\ldots)\),
which is the ordinary chain rule for a polynomial in the entries
\(x_{i}=\DX^{\,i-1}\gamma\).  This gives
\[
 \DX^{n+1}(F\compose G)
 =\sum_{k\ge1}\bigl(\DX^{k}F\bigr)\compose G\;
  \Bigl[\gamma\,\ExponentialPartialBellPolynomial{n}{k-1}
        +\mathcal D\,\ExponentialPartialBellPolynomial{n}{k}\Bigr]
   \bigl(\gamma,\DX\gamma,\ldots\bigr),
\]
after re-indexing the first sum.  By
\cref{plt:lem:tay-bell-recurrence} with \(x_{1}=\gamma\) the bracket equals
\(\ExponentialPartialBellPolynomial{n+1}{k}\), completing the induction.  All
rearrangements involve finitely many terms for each \(n\).
\end{proof}

\begin{remark}[Valuation bookkeeping]
For a near-identity \(G=X+H=t^{-1}+H\) one has
\(\DX(t^{-1})=t^{0}(0+1)=1\), hence \(\gamma=\DX G=1+\DX H\) with
\(\ordt(\DX H)\ge1\), and \(\DX^{i}\gamma=\DX^{i+1}H\) has
valuation at least \(i+1\) for \(i\ge1\).  Hence
\(\ExponentialPartialBellPolynomial{n}{k}(\gamma,\DX\gamma,\ldots)\) has
valuation at least \(n-k\), and the \(k\)-th term of
\eqref{plt:eq:tay-faa-di-bruno} has valuation at least
\(\ordt(F)+k+(n-k)=\ordt(F)+n\), consistent with
\(\ordt(\DX^{n}(F\compose G))\ge\ordt(F)+n\).  In this class the size of
\eqref{plt:eq:tay-faa-di-bruno} is a combinatorial, not an analytic,
difficulty.
\end{remark}

\section{Which engine to use}
\label{plt:sec:tay-comparison}

The generator form \eqref{plt:eq:cmp-near-identity-master} and the Taylor form
\eqref{plt:eq:cmp-taylor} have identical formal content, by
\cref{plt:thm:tay-formula}, and different computational profiles.

\begin{center}
\begin{tabular}{@{}lll@{}}
\hline
 & Generator substitution & Taylor composition\\
\hline
Core data & \((1+U)^{-1}\) and \(\log(1+U)\)
          & \(\DX^{k}F\) and \(H^{k}/k!\)\\
Reused across & many outer blocks of one \(F\)
              & many outer series with one \(H\)\\
Truncation index & powers of \(U\)
                 & derivative order \(k\le N-\ordt F\)\\
Valid for & any admissible \(\rho>0\) & \(\rho=1\), \(c=1\) only\\
Typical error & updating one generator only
              & applying it when \(H\notin\PLpow\)\\
\hline
\end{tabular}
\end{center}

\begin{checkpointbox}
Before composing, verify in order: (1) the inner argument has a leading
monomial \(cX^{\rho}\) with \(\rho>0\) and a chosen \(\log c\) in the
coefficient field; (2) the residual factor is \(1+U\) with
\(U\precdom1\); (3) the images \(1/G\) and \(\log G\) stay in the declared
exponent group and logarithmic depth; (4) both generators are updated
simultaneously; (5) the precision of \(U\) and \(V\) is set blockwise by
\cref{plt:cor:cmp-precision-budget}, not globally; (6) if an asymptotic
statement is intended, hypotheses (A1)--(A3) of
\cref{plt:thm:cmp-analytic-transfer} are separately verified.
\end{checkpointbox}

\begin{summarybox}
For an admissible monomial-leading \(G=cX^{\rho}(1+U)\) with a chosen
\(\log c\in\K\), composition is exact simultaneous substitution of the two
generators, \(t\mapsto c^{-1}t^{\rho}(1+U)^{-1}\) and
\(L\mapsto\rho L+\log c+\log(1+U)\).  It is a \(\K\)-algebra homomorphism,
multiplies valuations by \(\rho\), determines every output block by a finite
sum with a triangular precision budget, and is associative on the nose once
the chosen logarithms are compatible; right substitution reverses the apparent
operator order.  For \(\rho=1\), \(c=1\) it coincides with the exact formal
Taylor formula \(F\compose(X+H)=\sum_{k\ge0}H^{k}\DX^{k}F/k!\), whose
summability comes from the fact that \(\DX\) raises outer order while
multiplication by \(H\in\PLpow\) does not lower it, so that only
\(k\le N-\ordt(F)\) reach the block \(t^{N}\).  None of this is an analytic
statement; asymptotic validity requires the regime, branch and
remainder-transport hypotheses separately.
\end{summarybox}

% ---- fragment 07-group ----
\chapter{The near-identity composition group}
\label{plt:sec:near-identity-group}

Throughout this chapter \(\K\) is a field of characteristic zero, the large
variable is \(X\), and
\[
  t\DefinitionEquals X^{-1},
  \qquad
  L\DefinitionEquals\log X,
  \qquad
  \PLpow=\K[L][[t]],
  \qquad
  \PLlau=\K[L]((t)).
\]
The regime is \(X\to+\infty\) along the positive real ray; in a complex sector a
branch of \(\log\) must be fixed once and for all, and every statement below is
then read with that branch.  The outer valuation is \(\ordt\), the
distinguished derivation is \(\DX=t\partial_L-t^2\partial_t\), and
\(\shiftop{n}{k}=\prod_{j=0}^{k-1}(\partial_L-n-j)\) with
\(\shiftop{n}{0}=1\), so that
\begin{equation}\label{plt:eq:grp-block-derivative}
  \DX^k\bigl(t^np(L)\bigr)=t^{n+k}\,\shiftop{n}{k}p(L),
  \qquad n\in\IntegerNumbers,\ p\in\K[L].
\end{equation}
All truncations are the exclusive ones, \(\Trunc{F}{N}=\sum_{n<N}p_n(L)t^n\).

\section{The set, and why its correction term is unrestricted}

\begin{definition}[The near-identity class]\label{plt:def:grp-gnear}
\[
  \Gnear\DefinitionEquals\SetOf{X+A\ :\ A\in\PLpow}\subseteq\PLlau .
\]
\end{definition}

Membership imposes no smallness on \(A\) itself: the correction may be
\(L^{100}\), or any polynomial block of any degree.  What is small is the
\emph{relative} correction
\begin{equation}\label{plt:eq:grp-relative}
  \frac{X+A}{X}=1+tA\in1+t\PLpow,
\end{equation}
and every element of \(1+t\PLpow\) arises exactly once this way.  This is the
whole content of the phrase ``tangent to the identity'' in the present class:
\(tA\precdom1\) because \(t^nL^j\precdom1\) for every \(n\ge1\) and every
\(j\ge0\) in the dominance order of the positive ray.

% ed.: S6 writes $F/X = 1+tA = 1+o(1)$ "formally".  A little-o symbol with no
% regime attached is not a formal statement; the three O-roles must not be
% conflated.  The correct formal assertion is $tA \in t\PLpow$, equivalently
% $1+tA-1=\FormalBigOAt{t}{t}$, equivalently $tA\precdom1$.
\begin{warningbox}
\(1+tA=1+\FormalBigOAt{t}{t}\) is a statement about coefficient support only.
It asserts no numerical bound, no uniformity in \(L\), and no analytic
realization.  The corresponding analytic statement, \(G(X)/X\to1\) as
\(X\to+\infty\), is a separate hypothesis about a chosen realization of the
series.
\end{warningbox}

We shall use the formal exponential and logarithm on \(1+t\PLpow\) constantly.

% ed.: this lemma was stated and proved twice in the volume, here and as
% ed.: \cref{plt:lem:cmp-exp-log} in the composition part, with the same
% ed.: transport-of-identities proof.  The composition-part statement is the
% ed.: stronger of the two --- it is set in an arbitrary polynomial-logarithmic
% ed.: Hahn ring, and it adds the additivity of \(\exp\), the law
% ed.: \(\log(P^{n})=n\log P\) for all \(n\in\IntegerNumbers\), and the
% ed.: commutation with any homomorphism compatible with summable families,
% ed.: all of which are used below.  The restatement and its proof are deleted
% ed.: and replaced by this pointer; the label is kept so that the citations
% ed.: to it elsewhere in the volume continue to resolve.
\begin{remark}[Formal exponential and logarithm: where they are proved]
\label{plt:lem:grp-explog}
The assignments \(u\mapsto\exp u\) and \(1+u\mapsto\log(1+u)\), given by the
usual series, are mutually inverse bijections between \(t\PLpow\) and
\(1+t\PLpow\), and they carry products to sums.  This is
\cref{plt:lem:cmp-exp-log} applied to the Hahn ring
\(\mathcal A=\PLlau\) with \(\Gamma=\IntegerNumbers\) and
\(\mathfrak m=t\PLpow\); every reference to the present remark below is a
reference to that lemma, whose extra clauses --- additivity of \(\exp\),
\(\log(P^{\,n})=n\log P\) for \(n\in\IntegerNumbers\), and commutation with a
homomorphism compatible with summable families --- are used here as well.
\end{remark}

\section{Substitution, and the exact Taylor formula}

\begin{definition}[Substitution by a near-identity map]
\label{plt:def:grp-substitution}
Let \(G=X+A\in\Gnear\) and put
\[
  R\DefinitionEquals(1+tA)^{-1}\in1+t\PLpow,
  \qquad
  V\DefinitionEquals\log(1+tA)\in t\PLpow .
\]
For \(F=\sum_{n\ge n_0}p_n(L)t^n\in\PLlau\) define
\begin{equation}\label{plt:eq:grp-substitution}
  \Phi_G(F)=F\compose G
  \DefinitionEquals
  \sum_{n\ge n_0}t^nR^{\,n}\,p_n(L+V).
\end{equation}
\end{definition}

The two generator images encoded in \eqref{plt:eq:grp-substitution} are
\begin{equation}\label{plt:eq:grp-generators}
  t\compose G=\frac{t}{1+tA},
  \qquad
  L\compose G=L+\log(1+tA).
\end{equation}

\begin{warningbox}
The substitution is \(t\mapsto t(1+tA)^{-1}\), never \(t\mapsto t+A\): the
symbol \(t\) denotes \(X^{-1}\), not \(X\).  Updating \(L\) while leaving \(t\)
fixed, or conversely, is equally wrong; the two generators must be replaced
simultaneously.
\end{warningbox}

\begin{proposition}[Substitution operator]\label{plt:prop:grp-substitution}
Let \(G=X+A\in\Gnear\).  Then
\begin{enumerate}
\item the family in \eqref{plt:eq:grp-substitution} is formally summable, so
\(\Phi_G:\PLlau\to\PLlau\) is well defined, and
\(\Trunc{(F\compose G)}{N}\) depends only on
\(\Trunc{F}{N}\) and on \(\Trunc{A}{N-n_0}\);
\item \(\Phi_G\) is a \(\K\)-algebra endomorphism of \(\PLlau\);
\item \(\Phi_G(t)=t(1+tA)^{-1}\), \(\Phi_G(L)=L+\log(1+tA)\), and
\(\Phi_G(X)=G\);
\item \(\ordt\bigl(\Phi_G(F)\bigr)=\ordt(F)\) for every \(F\), and
\(\Phi_G(t^N\PLpow)\subseteq t^N\PLpow\) for every \(N\in\IntegerNumbers\);
in particular \(\Phi_G\) is injective and \(t\)-adically continuous;
\item \(\Phi_G(\PLpow)\subseteq\PLpow\) and
\(\Phi_G\bigl(\log(1+u)\bigr)=\log\bigl(1+\Phi_G(u)\bigr)\) for
\(u\in t\PLpow\).
\end{enumerate}
\end{proposition}

\begin{proof}
(1) Since \(R\in1+t\PLpow\) and \(V\in t\PLpow\), and since \(p_n\) is a
polynomial, \(p_n(L+V)=\sum_{j\le\degL p_n}\frac{p_n^{(j)}(L)}{j!}V^j\) is a
\emph{finite} sum lying in \(\PLpow\); likewise \(R^{\,n}\in\PLpow\) for
\(n\ge0\) and \(R^{\,n}=(1+tA)^{-n}\in\PLpow\) for \(n<0\).  The \(n\)-th
summand of \eqref{plt:eq:grp-substitution} therefore lies in \(t^n\PLpow\), and
\(\ordt\) of the \(n\)-th summand tends to \(+\infty\); the family is summable
and \(n_0\) bounds the support below.  Only the summands with \(n<N\) meet
outer order \(<N\), and within such a summand only \(R\) and \(V\) through
order \(t^{N-n-1}\) are needed, whence the stated dependency.

(2) Additivity is clear.  For multiplicativity write
\(F=\sum_np_nt^n\), \(\widetilde F=\sum_mq_mt^m\).  Evaluation
\(p\mapsto p(L+V)\) is a \(\K\)-algebra homomorphism \(\K[L]\to\PLpow\), so
\[
  \Phi_G(F\widetilde F)
  =\sum_{r}t^rR^{\,r}\Bigl(\sum_{n+m=r}p_nq_m\Bigr)(L+V)
  =\sum_{n,m}t^{n+m}R^{\,n+m}p_n(L+V)q_m(L+V)
  =\Phi_G(F)\Phi_G(\widetilde F),
\]
all rearrangements being legitimate because the families are summable.

(3) Take \(F=t\), \(F=L\), \(F=t^{-1}\) in \eqref{plt:eq:grp-substitution};
the last gives \(t^{-1}R^{-1}=t^{-1}(1+tA)=X+A=G\).

(4) The coefficient of \(t^{n_0}\) in \(\Phi_G(F)\) is \(p_{n_0}(L)\), because
\(R^{\,n_0}\in1+t\PLpow\) and \(p_{n_0}(L+V)\in p_{n_0}(L)+t\PLpow\); no other
summand reaches outer order \(n_0\).  Hence
\(\ordt(\Phi_GF)=n_0=\ordt(F)\), and the displayed inclusion is immediate from
\eqref{plt:eq:grp-substitution}.

(5) The inclusion is the case \(n_0\ge0\).  For the last identity,
\(\Phi_G\) is multiplicative and continuous, and
\(\Phi_G(u)\in t\PLpow\) by (4), so
\(\Phi_G(\log(1+u))=\sum_{k\ge1}\frac{(-1)^{k+1}}{k}\Phi_G(u)^k
 =\log(1+\Phi_G(u))\).
\end{proof}

% ed.: the exact formal Taylor formula was stated and proved twice, here and
% ed.: as \cref{plt:thm:tay-formula} in the composition part, where an entire
% ed.: subsection (the Taylor operator, its summability budget, its
% ed.: multiplicativity, and the generator-determination lemma) is built to
% ed.: prove it.  The statements are identical, so the restatement and its
% ed.: proof are deleted in favour of a pointer; the label is kept because the
% ed.: reversion part and the algorithms part cite it.  The repair recorded
% ed.: here --- S3 concludes from agreement on the generators without proving
% ed.: that the normally ordered operator is multiplicative --- is carried out
% ed.: at the surviving site, in \cref{plt:lem:tay-homomorphism}.
\begin{remark}[Exact formal Taylor formula]\label{plt:thm:grp-taylor}
For every \(F\in\PLlau\) and every \(A\in\PLpow\) the family
\((A^k\DX^kF/k!)_{k\ge0}\) is formally summable and
\(F\compose(X+A)=\sum_{k\ge0}A^k\DX^kF/k!\); if \(a=\ordt(F)\), only the
indices \(k\le N-a\) affect the coefficient of \(t^N\).  This is
\cref{plt:thm:tay-formula}, equation \eqref{plt:eq:cmp-taylor}, proved there;
it is an identity of formal series and carries no analytic content.  Every
citation of the present remark is a citation of that theorem.  Two features of
its proof are used repeatedly below: the operator
\(T_A=\sum_{k\ge0}\frac{A^k}{k!}\DX^k\) is a \emph{continuous} \(\K\)-algebra
endomorphism of \(\PLlau\) (\cref{plt:lem:tay-homomorphism}), and two such
endomorphisms agreeing on \(t\) and \(L\) are equal
(\cref{plt:lem:tay-determined}).
\end{remark}

The single most useful consequence is a gain of one outer block whenever a
near-identity argument is perturbed.

% ed.: the Lipschitz bound was printed as \(\ordt(B-C)+1\), dropping the
% ed.: summand \(\ordt(A)\) that the proof below already produces.  The
% ed.: sharper form is what the reversion part needs (it is what makes the
% ed.: Newton estimate of \cref{plt:thm:rev-newton-doubling} come out), and it
% ed.: was stated there a second time as a separate lemma.  The inequality is
% ed.: corrected here and the restatement in the reversion part is deleted;
% ed.: every use of the weaker form below is unaffected, since
% ed.: \(\ordt(A)\ge0\) on \(\PLpow\).
\begin{corollary}[Filtered gain, and the Lipschitz contraction]
\label{plt:cor:grp-gain}
Let \(A,B,C\in\PLpow\).  Then
\begin{align}
  \ordt\bigl(A\compose(X+B)-A\compose(X+C)\bigr)
  &\ \ge\ \ordt(A)+\ordt(B-C)+1,\label{plt:eq:grp-gain-lipschitz}\\
  \ordt\bigl(A\compose(X+B)-A\bigr)
  &\ \ge\ \ordt(A)+\ordt(B)+1,\label{plt:eq:grp-gain-absolute}
\end{align}
and, more precisely,
\begin{equation}\label{plt:eq:grp-gain-leading}
  A\compose(X+B)-A\ \equiv\ B\,\DX A
  \pmod{t^{\,\ordt(A)+2\ordt(B)+2}\PLpow}.
\end{equation}
\end{corollary}

\begin{proof}
By \cref{plt:thm:grp-taylor},
\(A\compose(X+B)-A\compose(X+C)=\sum_{k\ge1}\frac{B^k-C^k}{k!}\DX^kA\).
Since \(B^k-C^k=(B-C)\sum_{i=0}^{k-1}B^iC^{k-1-i}\) and all factors have
nonnegative valuation, \(\ordt(B^k-C^k)\ge\ordt(B-C)\); and
\(\ordt(\DX^kA)\ge\ordt(A)+k\).  Hence the \(k\)-th summand has
\(\ordt\ge\ordt(B-C)+\ordt(A)+k\), and the minimum over \(k\ge1\) is attained
at \(k=1\); this gives
\eqref{plt:eq:grp-gain-lipschitz}.  Taking \(C=0\) and using
\(\ordt(B^k)\ge k\,\ordt(B)\) gives, for the \(k\)-th term, the bound
\(\ordt(A)+k(1+\ordt(B))\), which is \(\ge\ordt(A)+\ordt(B)+1\) for \(k\ge1\)
and \(\ge\ordt(A)+2\ordt(B)+2\) for \(k\ge2\); this is
\eqref{plt:eq:grp-gain-absolute} and \eqref{plt:eq:grp-gain-leading}.
\end{proof}

% ed.: S4 proves the contraction step with a formal "mean value" identity
% A o (X+B) - A o (X+C) = (B-C)\int_0^1 A' o (X+C+s(B-C)) ds.  The integrand is
% a composition with a series carrying an auxiliary parameter s, and neither its
% existence nor the legitimacy of the term-by-term integration is established
% there.  Corollary \ref{plt:cor:grp-gain} replaces it by the exact Taylor
% difference, which needs nothing beyond Theorem \ref{plt:thm:grp-taylor}.

\section{The group theorem}

\begin{theorem}[The near-identity composition group]\label{plt:thm:grp-group}
\(\Gnear\) is a group under composition.  Explicitly:
\begin{enumerate}
\item \emph{(Closure)} for \(F=X+A\) and \(G=X+B\) in \(\Gnear\),
\begin{equation}\label{plt:eq:grp-closure}
  F\compose G=X+B+A\compose G\in\Gnear;
\end{equation}
\item \emph{(Associativity)} \(\Phi_H\circ\Phi_G=\Phi_{G\compose H}\) for
\(G,H\in\Gnear\); consequently
\((F\compose G)\compose H=F\compose(G\compose H)\) for every \(F\in\PLlau\) and
\(G,H\in\Gnear\);
\item \emph{(Identity)} \(X\compose G=G\) and \(F\compose X=F\);
\item \emph{(Inverses)} every \(F=X+A\in\Gnear\) has a unique two-sided
compositional inverse \(\CompositionalInverseOf{F}=X+B\in\Gnear\), and
\(\ordt(B)=\ordt(A)\).
\end{enumerate}
The group is not abelian.
\end{theorem}

\begin{proof}
\emph{(1)} \(\Phi_G(X+A)=\Phi_G(X)+\Phi_G(A)=G+A\compose G\) by
\cref{plt:prop:grp-substitution}(2),(3), and
\(A\compose G\in\PLpow\) by part (5) of the same proposition; so
\(F\compose G=X+\bigl(B+A\compose G\bigr)\) with correction in \(\PLpow\).

\emph{(2)} Both \(\Phi_H\circ\Phi_G\) and \(\Phi_{G\compose H}\) are
continuous \(\K\)-algebra endomorphisms of \(\PLlau\), so by
\cref{plt:lem:tay-determined} it suffices to check them on \(t\)
and \(L\).  Write \(G=X+A\), \(H=X+B\).  Since \(\Phi_H\) is a ring
homomorphism with \(\Phi_H(X)=H\) and \(\Phi_H(A)=A\compose H\),
\begin{equation}\label{plt:eq:grp-phi-of-G}
  \Phi_H(G)=H+A\compose H=G\compose H .
\end{equation}
Hence
\(\Phi_H(\Phi_G(t))=\Phi_H(G^{-1})=\Phi_H(G)^{-1}=(G\compose H)^{-1}
 =\Phi_{G\compose H}(t)\).
For \(L\), using \cref{plt:prop:grp-substitution}(5) and
\cref{plt:lem:grp-explog},
\begin{align*}
  \Phi_H\bigl(\Phi_G(L)\bigr)
  &=\Phi_H\bigl(L+\log(1+tA)\bigr)
   =L+\log(1+tB)+\log\Bigl(1+\tfrac{t}{1+tB}(A\compose H)\Bigr)\\
  &=L+\log\Bigl((1+tB)\Bigl(1+\tfrac{t(A\compose H)}{1+tB}\Bigr)\Bigr)
   =L+\log\bigl(1+t(B+A\compose H)\bigr)
   =\Phi_{G\compose H}(L),
\end{align*}
the last equality because the correction of \(G\compose H\) is
\(B+A\compose H\) by \eqref{plt:eq:grp-closure}.  Applying the resulting
operator identity to \(F\) gives
\((F\compose G)\compose H=F\compose(G\compose H)\).

% ed.: S5 proves associativity by "introducing finite truncations mod t^N", for
% which "ordinary substitution is associative".  Truncation is not a ring
% retraction compatible with substitution, and the claim that a truncated
% composite equals the truncation of the composite is precisely what needs
% proof.  S1 proposes the correct route (check on the generators) but does not
% carry out the verification for L; the display above is that verification.

\emph{(3)} \(\Phi_G(X)=G\) is \cref{plt:prop:grp-substitution}(3), and
\(\Phi_X=\mathrm{id}\) because \(A=0\) makes \(R=1\), \(V=0\) in
\eqref{plt:eq:grp-substitution}.

% ed.: the Picard construction of the inverse was carried out twice in the
% ed.: volume: here, and again --- with the sharp rate
% ed.: \(\ordt(B_{j+1}-B_{j})\ge(j+1)\alpha+j\), the first-order form of the
% ed.: inverse, and the uniqueness of the two-sided inverse --- as
% ed.: \cref{plt:thm:rev-existence}.  The two constructions are the same
% ed.: iteration; the reversion part proves strictly more.  The duplicate is
% ed.: deleted and item (4) is derived from the reversion theorem, which makes
% ed.: the dependency of the group's inverses on Part~IV explicit instead of
% ed.: hiding it behind two independent proofs.  There is no circularity:
% ed.: \cref{plt:thm:rev-existence} uses items (1)--(3) of the present theorem
% ed.: and the gain estimate \cref{plt:cor:grp-gain}, and nothing else from
% ed.: this section.
\emph{(4)}  By (1)--(3), \((\Gnear,\compose)\) is a monoid with identity
\(X\), and by \eqref{plt:eq:grp-closure} an element \(G=X+B\) of \(\Gnear\)
satisfies \(F\compose G=X\) if and only if
\begin{equation}\label{plt:eq:grp-inverse-fixed-point}
  B=-A\compose(X+B).
\end{equation}
\Cref{plt:thm:rev-existence} --- proved from items (1)--(3) above,
\cref{plt:cor:grp-gain}, and the \(t\)-adic completeness of \(\PLpow\), so that
the present item may cite it --- says that
\eqref{plt:eq:grp-inverse-fixed-point} has exactly one solution
\(B\in\PLpow\), that the resulting \(G=X+B\) is the unique element of
\(\Gnear\) with \(F\compose G=X\), that it satisfies \(G\compose F=X\) as well
and is the unique element of \(\Gnear\) doing so, and that
\(\ordt(B)=\ordt(A)\) with \([t^{\ordt(A)}]B=-[t^{\ordt(A)}]A\).  Writing
\(\CompositionalInverseOf{F}=\rev{F}=G\) gives (4).

\emph{Noncommutativity} is \cref{plt:ex:grp-noncommutative}.
\end{proof}

\begin{example}[Two maps that do not commute]\label{plt:ex:grp-noncommutative}
Let \(F=X+L\) and \(G=X+1\).  Then \(G\compose F=F+1=X+L+1\), whereas
\[
  F\compose G=G+L\compose G=X+1+\log(X+1)
  =X+1+L+\log(1+t)
  =X+1+L+t-\tfrac12t^2+\FormalBigOAt{t^3}{t}.
\]
Hence \(F\compose G-G\compose F=t-\frac12t^2+\FormalBigOAt{t^3}{t}\neq0\).
\end{example}

\begin{remark}[Direction of the dependency on Part~IV]
\label{plt:rmk:grp-dependency}
% ed.: this remark previously asserted that the section is independent of
% ed.: Part~IV and that inverses are established here.  That was true only
% ed.: because the same Picard construction was written out twice.  With the
% ed.: duplicate removed, the dependency is one-directional and is stated as
% ed.: such.
Items (1)--(3) of \cref{plt:thm:grp-group} are proved here from
\cref{plt:prop:grp-substitution} and \cref{plt:lem:tay-determined} alone, and
\cref{plt:cor:grp-gain} needs nothing beyond \cref{plt:thm:grp-taylor}.  Item
(4) is different: the existence, uniqueness, two-sidedness and leading data of
\(\CompositionalInverseOf{F}\) are proved once in the volume, in
\cref{plt:thm:rev-existence}, and item (4) cites that theorem.  The
dependency runs in one direction only.  Part~IV takes items (1)--(3) and
\cref{plt:cor:grp-gain} as given --- they are the imports recorded in
\cref{plt:rmk:rev-imports} --- and from them constructs the solution of
\eqref{plt:eq:grp-inverse-fixed-point} by Picard iteration with the sharp rate
\(\ordt(B_{j+1}-B_j)\ge(j+1)\alpha+j\), \(\alpha=\ordt(A)\); it adds the
first-order form
\(\CompositionalInverseOf{F}=X-A+\FormalBigOAt{t^{2\alpha+1}}{t}\)
(\cref{plt:cor:rev-first-order}), the residual certificates, the
Lagrange--B\"urmann coefficient formula and the algorithms that compute it.  A
reader who wants only the group law therefore needs, beyond this section, the
existence half of Part~IV and nothing else.
\end{remark}

\begin{remark}[Substitution reverses order]\label{plt:rmk:grp-anti}
By \cref{plt:thm:grp-group}(2), \(G\mapsto\Phi_G\) is an injective
\emph{anti}-homomorphism from \(\Gnear\) into the group of continuous
\(\K\)-algebra automorphisms of \(\PLlau\); injectivity holds because
\(\Phi_G(X)=G\).  Every sign discrepancy between group-theoretic and
vector-field formulas below can be traced to this order reversal.
\end{remark}

\section{What the group is not}

\begin{warningbox}
\(\Gnear\) is \emph{not} a multiplicative group of coefficient series.  Four
distinct statements are involved, and conflating them is the commonest error in
this subject.
\begin{itemize}
\item No element of \(\Gnear\) lies in \(\PLpow\): each has a
\(t^{-1}\) block.  The elements of \(\Gnear\) are maps, not coefficient
series.
\item The bijection \(X+A\mapsto1+tA\) of \eqref{plt:eq:grp-relative} onto the
multiplicative group \(1+t\PLpow\) is \emph{not} a group homomorphism.  Indeed
\(1+t\PLpow\) is abelian and, by \cref{plt:ex:grp-noncommutative}, \(\Gnear\)
is not.
\item \(\CompositionalInverseOf{F}\ne\MultiplicativeInverseOf{F}\).  For
\(F=X+L\), \(\MultiplicativeInverseOf{F}=t(1+tL)^{-1}=t-Lt^2+L^2t^3-\cdots\),
whereas \(\CompositionalInverseOf{F}=X-L+Lt+\FormalBigOAt{t^2}{t}\).
\item \(\Gnear\) is not the additive group of corrections either:
\(F\compose G=X+A+B+(A\compose G-A)\), and the last bracket vanishes only in
the associated graded object (\cref{plt:prop:grp-graded}).
\end{itemize}
Note also that \(t^N\PLpow\) is not an ideal of \(\PLlau\), because \(t\) is a
unit there; all filtration statements below are statements about the valuation
\(\ordt\), not about quotient rings.
\end{warningbox}

\section{The valuation filtration and its normal subgroups}

\begin{definition}[Contact filtration]\label{plt:def:grp-filtration}
For \(r\in\NaturalNumbers\) put
\[
  \Gfilt{r}\DefinitionEquals\SetOf{X+A\in\Gnear:\ \ordt(A)\ge r},
  \qquad
  \Gfilt{0}=\Gnear .
\]
Write \(F\equiv_rF'\) when the corrections of \(F,F'\in\Gnear\) differ by an
element of \(t^r\PLpow\), i.e.\ when \(F-F'=\FormalBigOAt{t^r}{t}\); one says
that \(F\) and \(F'\) have \emph{contact of order \(r\)}.
\end{definition}

\begin{proposition}[The filtration is by normal subgroups]
\label{plt:prop:grp-congruence}
For every \(r\ge0\), \(\equiv_r\) is a congruence for composition: if
\(F\equiv_rF'\) and \(G\equiv_rG'\) then \(F\compose G\equiv_rF'\compose G'\).
Consequently \(\Gnear/\!\equiv_r\) is a group, the quotient map is a
surjective homomorphism, and \(\Gfilt{r}\) is a normal subgroup of \(\Gnear\)
with \(\Gnear/\Gfilt{r}\cong\Gnear/\!\equiv_r\).
\end{proposition}

\begin{proof}
Write \(F=X+A\), \(F'=X+A'\), \(G=X+B\), \(G'=X+B'\).  Then
\[
  F\compose G-F'\compose G=(A-A')\compose G,
\]
which has \(\ordt=\ordt(A-A')\ge r\) by
\cref{plt:prop:grp-substitution}(4).  Next
\[
  F'\compose G-F'\compose G'=(B-B')+\bigl(A'\compose G-A'\compose G'\bigr),
\]
whose first term has \(\ordt\ge r\) and whose second has \(\ordt\ge r+1\) by
\eqref{plt:eq:grp-gain-lipschitz}.  Adding gives \(F\compose G\equiv_rF'\compose G'\).
Since \(\equiv_r\) is a congruence on a group, the quotient is a group and the
class of \(X\), which is exactly \(\Gfilt{r}\), is a normal subgroup.
\end{proof}

% ed.: S6 states the two continuity relations for the contact filtration but
% justifies the right-composition one by "H' = 1 + O(t)".  That is the reason
% for a gain in the *left* argument; the right-hand estimate is the exact
% isometry of Proposition \ref{plt:prop:grp-substitution}(4) and needs no
% derivative hypothesis at all.

\begin{proposition}[Separation and completeness]\label{plt:prop:grp-limit}
\(\bigcap_{r\ge0}\Gfilt{r}=\SetOf{X}\), and the natural map
\[
  \Gnear\longrightarrow\varprojlim_r\ \Gnear/\Gfilt{r}
\]
is an isomorphism of groups.
\end{proposition}

\begin{proof}
If \(X+A\in\Gfilt{r}\) for all \(r\) then \(\ordt(A)=+\infty\), so \(A=0\);
this is injectivity as well.  For surjectivity, let
\((F_r\Gfilt{r})_r\) be a compatible family.  Compatibility says
\(F_{r+1}\equiv_rF_r\), so the coefficient blocks of the corrections stabilise:
\([t^n]A_{F_r}\) is independent of \(r\) once \(r>n\).  Define \(A\) by
\([t^n]A\DefinitionEquals[t^n]A_{F_{n+1}}\); then \(A\in\PLpow\) and
\(X+A\equiv_rF_r\) for every \(r\).
\end{proof}

\begin{proposition}[Commutators drop one further block]
\label{plt:prop:grp-commutator}
For \(r,s\ge0\),
\[
  [\Gfilt{r},\Gfilt{s}]\subseteq\Gfilt{r+s+1},
  \qquad\text{where } [F,G]=\CompositionalInverseOf{F}\compose
  \CompositionalInverseOf{G}\compose F\compose G .
\]
More precisely, if \(F=X+A\in\Gfilt{r}\) and \(G=X+B\in\Gfilt{s}\) then the
correction of \([F,G]\) satisfies
\begin{equation}\label{plt:eq:grp-commutator-leading}
  [F,G]-X\ \equiv\ B\,\DX A-A\,\DX B
  \pmod{t^{\,r+s+2}\PLpow}.
\end{equation}
Consequently the lower central series obeys
\(\gamma_k(\Gnear)\subseteq\Gfilt{k-1}\), each quotient \(\Gnear/\Gfilt{r}\) is
nilpotent of class at most \(r\), and \(\Gnear\) is residually nilpotent.
\end{proposition}

\begin{proof}
Put \(P=F\compose G=X+A_P\) and \(Q=G\compose F=X+A_Q\).  By
\eqref{plt:eq:grp-closure} and \eqref{plt:eq:grp-gain-leading},
\[
  A_P=B+A\compose G\equiv A+B+B\,\DX A,
  \qquad
  A_Q=A+B\compose F\equiv A+B+A\,\DX B
\]
modulo \(t^{r+s+2}\PLpow\); note that the error terms in
\eqref{plt:eq:grp-gain-leading} are of order \(t^{r+2s+2}\) and
\(t^{s+2r+2}\), both at least \(t^{r+s+2}\).  Hence
\begin{equation}\label{plt:eq:grp-first-commutator}
  A_P-A_Q\equiv B\,\DX A-A\,\DX B \pmod{t^{\,r+s+2}\PLpow},
  \qquad
  \ordt(A_P-A_Q)\ge r+s+1 .
\end{equation}
Now \([F,G]=\CompositionalInverseOf{Q}\compose P\).  Writing
\(\CompositionalInverseOf{Q}=X+B_Q\) and using
\(X=\CompositionalInverseOf{Q}\compose Q=X+A_Q+B_Q\compose Q\), i.e.\
\(B_Q\compose Q=-A_Q\), we get
\[
  [F,G]=P+B_Q\compose P
       =X+\bigl(A_P-A_Q\bigr)+\bigl(B_Q\compose P-B_Q\compose Q\bigr).
\]
The last bracket has \(\ordt\ge\ordt(A_P-A_Q)+1\ge r+s+2\) by
\eqref{plt:eq:grp-gain-lipschitz}.  With
\eqref{plt:eq:grp-first-commutator} this proves both displayed claims.
The statement about \(\gamma_k\) follows by induction:
\(\gamma_1=\Gnear=\Gfilt{0}\) and
\(\gamma_{k+1}=[\gamma_k,\Gnear]\subseteq[\Gfilt{k-1},\Gfilt{0}]
 \subseteq\Gfilt{k}\).
\end{proof}

% ed.: S6 calls the group "pro-unipotent at finite truncation".  No algebraic
% group structure is available here (the graded pieces are the infinite
% dimensional space K[L]), and none is needed: Propositions
% \ref{plt:prop:grp-limit} and \ref{plt:prop:grp-commutator} give the exact
% provable statement, namely that Gnear is the inverse limit of nilpotent
% quotients.

\section{The associated graded object}

\begin{proposition}[Graded pieces]\label{plt:prop:grp-graded}
For every \(r\ge0\) the map
\[
  \Gfilt{r}/\Gfilt{r+1}\longrightarrow(\K[L],+),
  \qquad
  (X+A)\Gfilt{r+1}\longmapsto[t^r]A,
\]
is a group isomorphism.  In particular each graded piece is abelian and
uniquely divisible.
\end{proposition}

\begin{proof}
If \(F=X+A\) and \(G=X+B\) lie in \(\Gfilt{r}\), then by
\eqref{plt:eq:grp-closure} and \eqref{plt:eq:grp-gain-absolute}
\(F\compose G=X+A+B+(A\compose G-A)\) with
\(\ordt(A\compose G-A)\ge r+r+1\ge r+1\); hence
\([t^r](A_{F\compose G})=[t^r]A+[t^r]B\), so the map is a homomorphism.  It is
surjective (take \(A=p(L)t^r\)) and its kernel is \(\Gfilt{r+1}\).
\end{proof}

\begin{theorem}[The graded Lie algebra]\label{plt:thm:grp-gradedlie}
Let
\[
  \mathcal V\DefinitionEquals\K[L][t]\,\DX,
\]
the \(\K\)-span of the derivations \(a(L)t^{r}\DX\) with \(a\in\K[L]\) and
\(r\in\NaturalNumbers\), a Lie algebra under
\([\alpha\DX,\beta\DX]=(\alpha\DX\beta-\beta\DX\alpha)\DX\), graded by
\(\deg(a(L)t^r\DX)\DefinitionEquals r+1\); explicitly
\begin{equation}\label{plt:eq:grp-vf-bracket}
  \bigl[a\,t^r\DX,\ b\,t^s\DX\bigr]
  =t^{\,r+s+1}\Bigl(a\bigl(b'-sb\bigr)-b\bigl(a'-ra\bigr)\Bigr)\DX,
  \qquad{}'=\partial_L .
\end{equation}
Then the associated graded group
\(\operatorname{gr}\Gnear=\bigoplus_{r\ge0}\Gfilt{r}/\Gfilt{r+1}\), equipped
with the bracket induced by the group commutator, is a graded Lie algebra, and
\[
  \theta:\operatorname{gr}\Gnear\longrightarrow\mathcal V,
  \qquad
  \theta\bigl((X+A)\Gfilt{r+1}\bigr)=-\bigl([t^r]A\bigr)t^r\DX
\]
is an isomorphism of graded Lie algebras.
\end{theorem}

\begin{proof}
That \(\mathcal V\) is a Lie algebra needs no verification beyond closure:
for \(\alpha,\beta\in\K[L][t]\), \(\alpha\DX\beta-\beta\DX\alpha\in\K[L][t]\)
by \eqref{plt:eq:grp-block-derivative}, and the bracket of two derivations of
an associative algebra is a derivation, so the Jacobi identity is inherited.
Formula \eqref{plt:eq:grp-vf-bracket} is \eqref{plt:eq:grp-block-derivative}
applied twice, and it shows that the bracket adds the degrees:
\((r+1)+(s+1)=(r+s+1)+1\).

By \cref{plt:prop:grp-commutator} the group commutator induces a well-defined
biadditive map
\(\Gfilt{r}/\Gfilt{r+1}\times\Gfilt{s}/\Gfilt{s+1}\to
  \Gfilt{r+s+1}/\Gfilt{r+s+2}\):
replacing \(A\) by \(A+\widetilde A\) with \(\ordt\widetilde A\ge r+1\)
changes \(B\DX A-A\DX B\) by an element of \(t^{r+s+2}\PLpow\), and
biadditivity follows from \eqref{plt:eq:grp-commutator-leading} together with
\cref{plt:prop:grp-graded}.  Writing \(A=a(L)t^r+\cdots\) and
\(B=b(L)t^s+\cdots\), \eqref{plt:eq:grp-block-derivative} gives
\[
  [t^{\,r+s+1}]\bigl(B\,\DX A-A\,\DX B\bigr)
  =b\,(a'-ra)-a\,(b'-sb)
  =(ba'-ab')+(s-r)ab .
\]
On the other hand \eqref{plt:eq:grp-vf-bracket} gives
\[
  \bigl[-a\,t^r\DX,\ -b\,t^s\DX\bigr]
  =t^{\,r+s+1}\Bigl(a(b'-sb)-b(a'-ra)\Bigr)\DX
  =-t^{\,r+s+1}\Bigl((ba'-ab')+(s-r)ab\Bigr)\DX,
\]
which is exactly \(\theta\) applied to the previous display.  Thus \(\theta\)
intertwines the two brackets; it is bijective and degree preserving by
\cref{plt:prop:grp-graded}.  Since \(\mathcal V\) satisfies the Jacobi
identity, so does \(\operatorname{gr}\Gnear\).
\end{proof}

\begin{remark}
The sign in \(\theta\) is forced by \cref{plt:rmk:grp-anti}: composition and
operator composition run in opposite directions, so the induced bracket is the
opposite of the vector-field bracket.  For \(r=s=0\), that is for
\(F=X+f(L)\), \(G=X+g(L)\), \eqref{plt:eq:grp-commutator-leading} reads
\(F\compose G-G\compose F=t\,(gf'-fg')+\FormalBigOAt{t^2}{t}\), the classical
bracket of one-dimensional vector fields in the coordinate \(L=\log X\).
\end{remark}

\section{The exponential of a vector field, and iteration}

The Taylor operator \(T_A=\sum_k\frac{A^k}{k!}\DX^k\) of
\cref{plt:thm:grp-taylor} is \emph{normally ordered}: every power of \(A\)
stands to the left of every \(\DX\).  It must not be confused with the
operator exponential \(\exp(A\DX)\), in which \(\DX\) also differentiates
\(A\).  Both are automorphisms, and both are substitutions --- but by
different maps.

\begin{theorem}[Exponential correspondence]\label{plt:thm:grp-exp}
For \(A\in\PLpow\) the operator
\(\exp(A\DX)=\sum_{k\ge0}\frac1{k!}(A\DX)^k\) is a well-defined continuous
\(\K\)-algebra automorphism of \(\PLlau\), and
\begin{equation}\label{plt:eq:grp-exp-is-substitution}
  \exp(A\DX)=\Phi_{\mathcal E(A)},
  \qquad
  \mathcal E(A)\DefinitionEquals\exp(A\DX)(X)
   =X+A+\tfrac12A\,\DX A+\cdots\in\Gnear .
\end{equation}
The map \(\mathcal E:\PLpow\to\Gnear\) is a bijection, and
\(\mathcal E(A)\equiv X+A\pmod{t^{\,2\ordt(A)+1}\PLpow}\).
\end{theorem}

\begin{proof}
\emph{Well-definedness.}  \(\delta\DefinitionEquals A\DX\) is a derivation with
\(\delta(t^n\PLpow)\subseteq t^{n+1}\PLpow\), so for \(F\in t^{-m}\PLpow\) the
family \((\delta^kF/k!)_k\) is summable and \(\exp(\delta)F\in t^{-m}\PLpow\);
continuity and \(\K\)-linearity are clear.  Multiplicativity follows from the
Leibniz rule exactly as for \(T_A\), and \(\exp(-\delta)\) is a two-sided
inverse because the identity
\(\sum_{k}\frac{x^k}{k!}\sum_j\frac{(-x)^j}{j!}=1\) holds coefficientwise.

\emph{Identification with a substitution.}  Let \(\Psi_s=\exp(s\delta)\); for
each \(n\), \([t^n]\Psi_s(F)\) is a polynomial in \(s\) with coefficients in
\(\K[L]\), because \(\ordt(\delta^kF)\ge\ordt(F)+k\).  Put
\[
  u(s)=\Psi_s(L),
  \qquad
  v(s)=\log\bigl(\Psi_s(X)\bigr)
      \DefinitionEquals L+\log\bigl(1+t(\Psi_s(X)-X)\bigr),
\]
the latter legitimate because \(\Psi_s(X)-X\in\PLpow\), so that
\(t(\Psi_s(X)-X)\in t\PLpow\) and \cref{plt:lem:grp-explog} applies.
Differentiating in \(s\) coefficientwise, and using
\(\frac{\Differential}{\Differential s}\Psi_s=\Psi_s\circ\delta\),
\[
  u'(s)=\Psi_s(A\,\DX L)=\Psi_s(At),
  \qquad
  v'(s)=\frac{\Psi_s(A\,\DX X)}{\Psi_s(X)}
       =\Psi_s(A)\,\Psi_s(t)=\Psi_s(At),
\]
where \(\Psi_s(X)^{-1}=\Psi_s(t)\) because \(\Psi_s\) is a ring automorphism.
Since \(u(0)=v(0)=L\) and \(\K\) has characteristic zero, \(u=v\).  Setting
\(s=1\) and writing \(G=\mathcal E(A)=\Psi_1(X)\), we get
\(\Psi_1(L)=\log G=\Phi_G(L)\) and
\(\Psi_1(t)=\Psi_1(X)^{-1}=G^{-1}=\Phi_G(t)\); both maps are continuous algebra
homomorphisms, so \(\Psi_1=\Phi_G\).

\emph{Bijectivity.}  Write \(\mathcal E(A)=X+A+\mathcal R(A)\) with
\(\mathcal R(A)=\sum_{k\ge2}\frac1{k!}(A\DX)^{k-1}A\).  Each term of
\(\mathcal R(A)\) is a sum of products of \(k\) factors drawn from \(A\) and
its \(\DX\)-derivatives, carrying \(k-1\) derivatives in total, so
\(\ordt\ge k\,\ordt(A)+k-1\ge2\ordt(A)+1\); this proves the last assertion.
Replacing \(A\) by \(A'\) one factor at a time in each such product shows, for
\(A,A'\in\PLpow\),
\[
  \ordt\bigl(\mathcal R(A)-\mathcal R(A')\bigr)\ge\ordt(A-A')+1,
\]
because each single-factor difference contributes \(\ordt(A-A')\) and the
remaining \(k-1\ge1\) derivatives contribute at least \(1\).  Hence
\(\mathcal E(A)-\mathcal E(A')\equiv A-A'\) modulo
\(t^{\ordt(A-A')+1}\PLpow\).  Injectivity is immediate.  For surjectivity, let
\(G=X+A_\ast\in\Gnear\); construct \(A^{(r)}\) with
\(\mathcal E(A^{(r)})\equiv_rG\) by setting \(A^{(0)}=0\) and
\(A^{(r+1)}=A^{(r)}-e_rt^r\), where \(e_rt^r\) is the outer block of order
\(r\) of \(\mathcal E(A^{(r)})-G\).  The sequence is Cauchy and its limit
\(A\) satisfies \(\mathcal E(A)=G\) by continuity.
\end{proof}

\begin{corollary}[Unique divisibility and continuous iteration]
\label{plt:cor:grp-iteration}
For \(F\in\Gnear\) let \(A=\mathcal E^{-1}(F)\) and define
\[
  F^{[s]}\DefinitionEquals\exp(sA\DX)(X)=\mathcal E(sA),\qquad s\in\K .
\]
Then \(F^{[1]}=F\), \(F^{[0]}=X\),
\(F^{[s]}\compose F^{[s']}=F^{[s+s']}\), \(F^{[n]}\) is the \(n\)-fold
composite for \(n\in\PositiveIntegers\), and
\(F^{[-1]}=\CompositionalInverseOf{F}\).  In particular \(\Gnear\) is uniquely
divisible: for each \(n\in\PositiveIntegers\) there is exactly one
\(H\in\Gnear\) with \(H^{[n]}=F\), namely \(H=\mathcal E(A/n)\).  Moreover
\(F^{[s]}\equiv X+sA\pmod{t^{\,2\ordt(A)+1}\PLpow}\).
\end{corollary}

\begin{proof}
By \cref{plt:thm:grp-exp} and \cref{plt:thm:grp-group}(2),
\(\Phi_{F^{[s']}\compose F^{[s]}}=\Phi_{F^{[s]}}\Phi_{F^{[s']}}
 =\exp(sA\DX)\exp(s'A\DX)=\exp((s+s')A\DX)=\Phi_{F^{[s+s']}}\);
apply both sides to \(X\) and use injectivity of \(G\mapsto\Phi_G\).  The
remaining assertions are immediate, uniqueness of \(n\)-th roots because
\(H^{[n]}=\mathcal E(nB)\) for \(H=\mathcal E(B)\) and \(\mathcal E\) is
injective.
\end{proof}

% ed.: S6 introduces log(S_F) as a power series in S_F - I and writes the
% fractional iterate as "F^{[s]} = X S_F^s".  Three things are missing there:
% that S_F - I raises the filtration (true: (S_F-I)H = H o F - H has
% ordt >= ordt(H)+1 by Corollary \ref{plt:cor:grp-gain}); that S_F^s is again a
% *substitution* operator, i.e. that F^{[s]} exists in Gnear at all; and the
% notation, which should be S_F^s(X), not a product.  Theorem
% \ref{plt:thm:grp-exp} and Corollary \ref{plt:cor:grp-iteration} supply all
% three.

\begin{remark}[Normal ordering versus the flow]
\label{plt:rmk:grp-normal-ordering}
\(T_A\ne\exp(A\DX)\) in general.  \(T_A=\Phi_{X+A}\) substitutes the map
\(X+A\); \(\exp(A\DX)=\Phi_{\mathcal E(A)}\) is the time-one flow of the
vector field \(A\DX\), and
\(\mathcal E(A)=X+A+\frac12A\DX A+\cdots\).  For \(A=L\), say,
\(\mathcal E(L)=X+L+\frac12Lt+\FormalBigOAt{t^2}{t}\ne X+L\).  The two
operators agree only when \(A\DX A=0\) to all relevant orders, in particular
when \(A\in\K\).  Confusing them turns an exact composition formula into an
approximate one.
\end{remark}

\begin{summarybox}
\(\Gnear=X+\PLpow\) is a group under composition, with identity \(X\) and with
inverses produced by a fixed-point iteration that gains one outer block per
step.  Its contact filtration \((\Gfilt{r})\) consists of normal subgroups with
\([\Gfilt{r},\Gfilt{s}]\subseteq\Gfilt{r+s+1}\); the quotients are nilpotent,
the group is their inverse limit, and the associated graded object is the
graded Lie algebra \(\K[L][t]\DX\) of polynomial vector fields.  Exponentiating
those vector fields makes \(\Gnear\) uniquely divisible and gives a continuous
iteration \(F^{[s]}\).  None of this is a statement about multiplication of
coefficient series.
\end{summarybox}


\chapter{Scale-changing composition and deeper logarithms}
\label{plt:sec:scale-change}

\cref{plt:sec:near-identity-group} treated inner maps that perturb the scale.
This chapter treats inner maps that \emph{change} it.  The governing principle
is the same in both cases and is worth isolating before any case analysis.

\section{The substitution principle}

Every \(F\in\PLlau\) is built from the two generators \(t=X^{-1}\) and
\(L=\log X\) by sums, products and \(t\)-adic limits.  Therefore a substitution
\(X\mapsto G\) is completely determined by the two images
\begin{equation}\label{plt:eq:grp-two-images}
  t\compose G=\MultiplicativeInverseOf{G},
  \qquad
  L\compose G=\log G,
\end{equation}
and the only question is which ambient class contains them and is closed enough
to receive the resulting sums.  Answering that question \emph{first} is the
difference between a correct computation and a silently truncated one:
computing in a class that lacks a monomial forced by
\eqref{plt:eq:grp-two-images} does not make that monomial disappear, it makes
the retained terms wrong.

\begin{definition}[Nested logarithms, monomial groups, Hahn ambients]
\label{plt:def:grp-logtower}
Put
\[
  L_0\DefinitionEquals X,
  \qquad
  L_{j+1}\DefinitionEquals\log L_j,
\]
so that \(L_1=L=\log X\) and \(L_2=\log\log X\).  For \(r\ge1\) and an ordered
subgroup \(\Gamma\subseteq\RealNumbers\) let
\[
  \mathcal M_r(\Gamma)
  \DefinitionEquals
  \SetOf{L_0^{\alpha_0}L_1^{\alpha_1}\cdots L_r^{\alpha_r}
        :\ \alpha_0,\ldots,\alpha_r\in\Gamma},
\]
ordered by declaring one monomial dominant over another when, at the least
index \(j\) at which the exponents differ, the first exponent is the larger.
This is a totally ordered abelian group, and we write
\(\mathcal H_r(\K,\Gamma)\) for the Hahn field of series
\(\sum_{\mathfrak m\in S}c_{\mathfrak m}\mathfrak m\) with
\(c_{\mathfrak m}\in\K\) and \(S\subseteq\mathcal M_r(\Gamma)\) well ordered in
the direction of increasing smallness \cite{hahn,edgar-beginners}.  On the
positive ray this order is the asymptotic order as \(X\to+\infty\).  One has
\[
  \PLpow\subset\PLlau\subset\PLrat\subset\PLit\subset
  \mathcal H_1(\K,\IntegerNumbers),
  \qquad
  \PLlau\subset\K[L]\bigl((t^{1/s})\bigr)\subset
  \mathcal H_1\bigl(\K,\tfrac1s\IntegerNumbers\bigr).
\]
\end{definition}

\begin{warningbox}
\(\mathcal H_r(\K,\Gamma)\) is a bookkeeping device, not a licence.  Membership
is a statement about support; it says nothing about convergence, and for
non-integer exponents it presupposes a branch of \(X^{\alpha_0}\) and of every
logarithm involved.  All statements below are for \(X\to+\infty\) on the
positive real ray with the real positive branches; in a sector, each power and
each logarithm must be assigned a branch before the classification is applied.
\end{warningbox}

\section{Power--logarithm inner maps: the classification}

\begin{theorem}[Classification of power--log scale changes]
\label{plt:thm:grp-scale-classification}
Let \(\rho\in\RealNumbers_{>0}\), \(\sigma\in\RealNumbers\),
\(c\in\K^{\times}\) with a chosen \(\log c\in\K\), and \(U\in t\PLpow\).  Put
\[
  G=c\,X^{\rho}L^{\sigma}(1+U),
\]
and let \(\Gamma_0=\rho\IntegerNumbers+\IntegerNumbers\),
\(\Gamma_1=\sigma\IntegerNumbers+\IntegerNumbers\).  Then for every
\(F=\sum_{n\ge n_0}p_n(L)t^n\in\PLlau\) the family
\begin{equation}\label{plt:eq:grp-scale-formula}
  F\compose G
  =\sum_{n\ge n_0}
   c^{-n}\,t^{\rho n}L^{-\sigma n}(1+U)^{-n}\,
   p_n\bigl(\rho L+\sigma L_2+\log c+\log(1+U)\bigr)
\end{equation}
is well based and locally finite, so \(F\compose G\) is a well-defined element
of \(\mathcal H_2(\K,\Gamma_0+\Gamma_1)\), with \(L_0\)-exponents in
\(-\Gamma_0\), \(L_1\)-exponents in \(-\sigma\IntegerNumbers+\NaturalNumbers\)
and \(L_2\)-exponents in \(\NaturalNumbers\); and
\(\Phi_G:F\mapsto F\compose G\) is an injective \(\K\)-algebra homomorphism.
Moreover:
\begin{enumerate}
\item if \(\sigma=0\) and \(\rho\in\PositiveIntegers\), then
\(\Phi_G(\PLlau)\subseteq\PLlau\);
\item if \(\sigma=0\) and \(\rho=r/s\in\RationalNumbers_{>0}\) in lowest terms,
then \(\Phi_G(\PLlau)\subseteq\K[L]((t^{1/s}))\); for \(s>1\) this is sharp, in
that \(\Phi_G(t)\notin\PLlau\) and the subgroup of \(\RealNumbers\) generated
by the \(L_0\)-exponents occurring in \(\Phi_G(\PLlau)\) is exactly
\(\frac1s\IntegerNumbers\);
\item if \(\sigma=0\) and \(\rho\notin\RationalNumbers\), then no Puiseux-log
extension \(\K[L]((t^{1/s}))\) contains \(\Phi_G(t)\), and
\(\Gamma_0\) is dense in \(\RealNumbers\);
\item if \(\sigma\ne0\), then \(L_2\) occurs in \(F\compose G\) if and only if
\(p_n\notin\K\) for some \(n\); and if \(p_n\ne0\) with
\(\sigma n\notin-\NaturalNumbers\) then an \(L_1\)-exponent outside
\(\NaturalNumbers\) occurs.  In particular \(\Phi_G(\PLlau)\) is contained in
no depth-one class as soon as some coefficient block of some \(F\) is
nonconstant.
\end{enumerate}
\end{theorem}

\begin{proof}
Formula \eqref{plt:eq:grp-scale-formula} is \eqref{plt:eq:grp-two-images}
substituted into \(F=\sum p_n(L)t^n\), using
\[
  \MultiplicativeInverseOf{G}=c^{-1}t^{\rho}L^{-\sigma}(1+U)^{-1},
  \qquad
  \log G=\log c+\rho L_1+\sigma L_2+\log(1+U)
\]
with the chosen branches, together with the finite Taylor expansion
\(p_n(a+w)=\sum_{j\le\degL p_n}\frac{p_n^{(j)}(a)}{j!}w^j\) at
\(a=\rho L+\sigma L_2+\log c\), \(w=\log(1+U)\in t\PLpow\).

\emph{Support.}  The \(n\)-th summand has \(L_0\)-exponents in
\(-\rho n-\NaturalNumbers\), since \((1+U)^{-n}\) and \(\log(1+U)\) lie in
\(\PLpow\) and \(p_n\) is a polynomial.  Because \(\rho>0\), for each real
\(M\) only finitely many pairs
\((n,k)\in\IntegerNumbers_{\ge n_0}\times\NaturalNumbers\) satisfy
\(\rho n+k\le M\).  Hence the set of \(L_0\)-exponents occurring is well
ordered in the direction of increasing smallness, each is reached by only
finitely many pairs, and each contribution carries only finitely many monomials
in \(L_1,L_2\).  The family is therefore summable in
\(\mathcal H_2(\K,\Gamma_0+\Gamma_1)\), and the stated exponent ranges are read
off the display.

\emph{Homomorphism and injectivity.}  Multiplicativity is proved exactly as in
\cref{plt:prop:grp-substitution}(2), evaluation of polynomials being a ring
homomorphism.  For injectivity let \(F\ne0\) with \(\ordt F=n_0\).  Since
\(\rho\ne0\), the map \(Y\mapsto\rho L+\sigma L_2+\log c\) induces an injection
\(\K[Y]\hookrightarrow\K[L_1,L_2]\), so
\(p_{n_0}(\rho L+\sigma L_2+\log c)\ne0\); its dominant monomial times
\(c^{-n_0}X^{-\rho n_0}L^{-\sigma n_0}\) is the dominant monomial of the whole
sum, because every other summand and every \(U\)-correction has strictly
smaller \(L_0\)-exponent.  Hence \(F\compose G\ne0\).

\emph{(1)} With \(\sigma=0\) and \(\rho\in\PositiveIntegers\) all
\(L_0\)-exponents lie in \(\IntegerNumbers\), all \(L_1\)-exponents in
\(\NaturalNumbers\), no \(L_2\) occurs, and only finitely many
\(L_0\)-exponents exceed \(-\rho n_0\); this is exactly membership of
\(\PLlau\).

\emph{(2)} All exponents lie in \(\frac1s\IntegerNumbers\).  Taking \(U=0\) and
\(F=t^{k}\) gives \(\Phi_G(t^k)=c^{-k}t^{kr/s}\); as \(k\) ranges over
\(\IntegerNumbers\) these exponents generate
\(\frac{r}{s}\IntegerNumbers+\IntegerNumbers=\frac1s\IntegerNumbers\), because
\(\gcd(r,s)=1\).  For \(s>1\) we have \(\rho\notin\IntegerNumbers\), so
\(\Phi_G(t)\) has \(L_0\)-exponent \(-\rho\notin\IntegerNumbers\) and cannot
lie in \(\PLlau\), all of whose \(L_0\)-exponents are integers.

\emph{(3)} \(\Phi_G(t)\) has \(L_0\)-exponent \(-\rho\notin\RationalNumbers\),
while every element of \(\K[L]((t^{1/s}))\) has \(L_0\)-exponents in
\(\frac1s\IntegerNumbers\subseteq\RationalNumbers\).  Density of
\(\rho\IntegerNumbers+\IntegerNumbers\) for irrational \(\rho\) is the
classical equidistribution-free statement that a nondiscrete subgroup of
\(\RealNumbers\) is dense.

\emph{(4)} If every \(p_n\) lies in \(\K\), then
\eqref{plt:eq:grp-scale-formula} reduces to
\(\sum_nc^{-n}p_n\,t^{\rho n}L^{-\sigma n}(1+U)^{-n}\), in which \(L_2\) does
not occur.  Conversely, let \(n_1\) be least with \(\degL p_{n_1}\ge1\).
Summands with \(n<n_1\) have constant \(p_n\) and contribute no \(L_2\), so the
coefficient of \(X^{-\rho n_1}\) receives \(L_2\) only from \(n=n_1\), where it
equals \(c^{-n_1}L^{-\sigma n_1}p_{n_1}(\rho L+\sigma L_2+\log c)\); its
\(L_2\)-degree is \(\degL p_{n_1}\ge1\), with leading coefficient
\(\sigma^{\degL p_{n_1}}\) times that of \(p_{n_1}\), hence nonzero because
\(\sigma\ne0\).  The claim about \(L_1\)-exponents is read off the same
summand, whose \(L_1\)-exponents are \(-\sigma n+\NaturalNumbers\).
\end{proof}

\begin{corollary}[The three closure regimes for a pure power]
\label{plt:cor:grp-power-regimes}
For \(G=X^{\rho}\) with \(\rho>0\) one has \(t\compose G=t^{\rho}\),
\(L\compose G=\rho L\), and
\begin{equation}\label{plt:eq:grp-power-composition}
  \Bigl(\sum_{n\ge n_0}p_n(L)t^n\Bigr)\compose X^{\rho}
  =\sum_{n\ge n_0}p_n(\rho L)\,t^{\rho n}.
\end{equation}
The composite stays in \(\PLlau\) for \(\rho\in\PositiveIntegers\); requires
exactly the Puiseux-log extension \(\K[L]((t^{1/s}))\) for
\(\rho=r/s\in\RationalNumbers_{>0}\) in lowest terms; and requires a Hahn
exponent group containing \(\rho\IntegerNumbers+\IntegerNumbers\) for
irrational \(\rho\).  For \(\rho\in\PositiveIntegers\), \(\Phi_{X^{\rho}}\) is
an automorphism of \(\PLlau\) if and only if \(\rho=1\).
\end{corollary}

\begin{proof}
All but the last sentence is
\cref{plt:thm:grp-scale-classification}(1)--(3) with \(c=1\), \(\sigma=0\),
\(U=0\).  For \(\rho\ge2\), \eqref{plt:eq:grp-power-composition} shows every
\(L_0\)-exponent of the image lies in \(\rho\IntegerNumbers\), so \(t\) is not
in the image; for \(\rho=1\), \(\Phi_X\) is the identity.
\end{proof}

% ed.: S4's closure table records only the rational case ("c x^{r/s}:
% Puiseux-log series") and S1 says "arbitrary real beta require a Hahn support",
% neither separating the three regimes nor proving that each enlargement is
% necessary.  Corollary \ref{plt:cor:grp-power-regimes} does both.

\begin{proposition}[Dilations and the semidirect structure]
\label{plt:prop:grp-dilation}
Let \(C\subseteq\K^{\times}\) be a subgroup carrying a fixed group homomorphism
\(\mathrm{lg}:C\to(\K,+)\), a \emph{multiplicative branch choice}; an analytic
realisation additionally requires \(\EulerE^{\mathrm{lg}(c)}=c\).  For
\(c\in C\) write \(m_c(X)=cX\).  Then
\[
  t\compose m_c=c^{-1}t,
  \qquad
  L\compose m_c=L+\mathrm{lg}(c),
\]
\(\Phi_{m_c}\) is an automorphism of \(\PLlau\) preserving \(\PLpow\), the
dilations normalise \(\Gnear\) by
\[
  \CompositionalInverseOf{m_c}\compose(X+A)\compose m_c
  =X+c^{-1}\bigl(A\compose m_c\bigr)\in\Gnear,
\]
and \(\mathcal G_{\mathrm{lin}}=\SetOf{cX+A:\ c\in C,\ A\in\PLpow}\) is a group
under composition; it is the semidirect product of \(C\) acting on the normal
subgroup \(\Gnear\) by the displayed conjugation.
\end{proposition}

\begin{proof}
The generator images are immediate from \eqref{plt:eq:grp-two-images}, and they
show that \(\Phi_{m_c}\) maps \(\PLpow\) into \(\PLpow\) with two-sided inverse
\(\Phi_{m_{c^{-1}}}\), using \(\mathrm{lg}(c^{-1})=-\mathrm{lg}(c)\).  Next,
\((X+A)\compose m_c=cX+A\compose m_c\), and composing with \(m_{c^{-1}}\) on
the outside multiplies by \(c^{-1}\); since \(A\compose m_c\in\PLpow\), the
conjugate lies in \(\Gnear\).  Finally \(cX+A=m_c\compose(X+c^{-1}A)\) writes
\(\mathcal G_{\mathrm{lin}}\) as the product of the subgroup
\(\SetOf{m_c:c\in C}\cong C\) with the normal subgroup \(\Gnear\), the two
intersecting trivially; the group axioms follow from
\cref{plt:thm:grp-group} together with the displayed normalisation.
\end{proof}

\begin{remark}[Why the branch constant must be carried as data]
\label{plt:rmk:grp-branch}
For \(\K=\RealNumbers\) and \(C=\RealNumbers_{>0}\) the ordinary logarithm is a
homomorphism and \cref{plt:prop:grp-dilation} applies verbatim.  For
\(\K=\ComplexNumbers\) no multiplicative branch choice exists on all of
\(\ComplexNumbers^{\times}\): a homomorphism
\(\mathrm{lg}:\ComplexNumbers^{\times}\to(\ComplexNumbers,+)\) splitting
\(\exp\) would exhibit \(\ker\exp=2\pi\ImaginaryUnit\IntegerNumbers\) as a
direct summand of \((\ComplexNumbers,+)\); but the latter is a
\(\RationalNumbers\)-vector space, hence divisible, every direct summand of a
divisible group is divisible, and \(\IntegerNumbers\) is not.  Consequently a
symbolic implementation must carry the chosen value of \(\log c\) as data
attached to the composition rather than recompute it from a principal branch:
changing \(\PrincipalLogarithm c\) by \(2\pi\ImaginaryUnit j\) changes the
constant block of \(L\compose G\) and can move the result onto a different
branch of a multivalued inverse.
\end{remark}

\section{Two boundary cases: a logarithmic shift and a logarithmic factor}

The maps \(X\mapsto X+c\log X\) and \(X\mapsto X\log X\) look similar and
behave completely differently.

\begin{proposition}[A logarithmic shift stays inside, at a cost]
\label{plt:prop:grp-logdegree}
Let \(c\in\K^{\times}\) and \(G=X+cL\).  Then \(G\in\Gnear\); in particular
\(\Phi_G(\PLlau)\subseteq\PLlau\) and \(G\) is invertible in \(\Gnear\).  The
logarithmic degrees of the composite are nevertheless unbounded:
\[
  t\compose G=\frac{t}{1+cLt}=\sum_{n\ge1}(-c)^{n-1}L^{n-1}t^{n},
  \qquad
  L\compose G=L+\sum_{n\ge1}\frac{(-1)^{n+1}c^{n}}{n}L^{n}t^{n},
\]
so that, in the log-degree profile notation,
\(\dprofile{(t\compose G)}(n)=n-1\) and \(\dprofile{(L\compose G)}(n)=n\) for
\(n\ge1\).
\end{proposition}

\begin{proof}
\(cL\in\PLpow\), so \(G\in\Gnear\) and \cref{plt:prop:grp-substitution} and
\cref{plt:thm:grp-group} apply.  The two displays are the geometric and
logarithmic series for \(u=cLt\in t\PLpow\), legitimate by
\cref{plt:lem:grp-explog}; each coefficient block is a single monomial in
\(L\), so its degree is read off directly.
\end{proof}

\begin{remark}
\cref{plt:prop:grp-logdegree} is the reason why the definition of \(\PLpow\)
imposes no uniform bound on \(\degL p_n\): the smallest class closed under
composition by \(X+cL\) must already contain blocks whose degree grows linearly
in \(n\).  A linear bound \(\dprofile{F}(n)\le an+b\) is therefore a hypothesis
or a proved invariant, never part of the ring definition.  The compositional
inverse of \(X+cL\) is the Wright-omega object whose coefficient blocks are the
Lambert polynomials \(\LamP{n}\); see \cref{plt:thm:rev-existence} and
the surrounding material in Part~IV.
\end{remark}

\begin{proposition}[A logarithmic factor forces depth two]
\label{plt:prop:grp-xlogx}
Let \(G=XL=X\log X\).  Then \(t\compose G=tL^{-1}\), \(L\compose G=L+L_2\), and
for \(F=\sum_{n\ge n_0}p_n(L)t^n\in\PLlau\),
\begin{equation}\label{plt:eq:grp-xlogx}
  F\compose(XL)=\sum_{n\ge n_0}t^{n}L^{-n}\,p_n(L+L_2)
  \ \in\ \K[L_2]\bigl[L,L^{-1}\bigr]((t)).
\end{equation}
The symbol \(L_2\) occurs if and only if some \(p_n\) is nonconstant; a
negative power of \(L\) occurs if and only if \(p_n\ne0\) for some \(n>0\).
Hence
\[
  F\compose(XL)\in\PLlau
  \quad\Longleftrightarrow\quad
  F\in\K[X]=\K[t^{-1}].
\]
\end{proposition}

\begin{proof}
Apply \cref{plt:thm:grp-scale-classification} with \(c=1\), \(\rho=1\),
\(\sigma=1\), \(U=0\); \eqref{plt:eq:grp-scale-formula} becomes
\eqref{plt:eq:grp-xlogx}.  The \(L_2\) criterion is part~(4) of that theorem.
For the powers of \(L\): the coefficient of \(t^n\) is
\(L^{-n}p_n(L+L_2)\), whose \(L\)-exponents run over
\(\SetOf{-n,\ldots,-n+\degL p_n}\), so a negative exponent occurs precisely
when some \(n>0\) has \(p_n\ne0\) --- note that the top \(L_2\)-degree part of
\(p_n(L+L_2)\) has \(L\)-degree \(0\), so no cancellation can remove the
exponent \(-n\).  Both defects are excluded in \(\PLlau\), whose
\(L\)-exponents lie in \(\NaturalNumbers\) and which contains no \(L_2\).
Membership in \(\PLlau\) therefore forces every \(p_n\) to be a constant and
\(p_n=0\) for \(n>0\), that is, \(F\in\K[t^{-1}]=\K[X]\); conversely
\(X^{k}\compose(XL)=t^{-k}L^{k}\in\PLlau\) for \(k\ge0\).
\end{proof}

\begin{example}[Both enlargements at once]\label{plt:ex:grp-scale-change}
Let \(F(X)=X+\log X+X^{-1}\) and \(G(X)=3X^{2}L^{-1}\), with \(\log3\in\K\).
Then
\[
  F\compose G
  =3X^{2}L^{-1}+\log\bigl(3X^{2}L^{-1}\bigr)+\frac{L}{3X^{2}}
  =3X^{2}L^{-1}+2L-L_2+\log3+\tfrac13L\,t^{2}.
\]
Here \(\rho=2\) and \(\sigma=-1\): the leading block carries \(L^{-1}\) and the
second logarithm \(L_2\) appears \emph{exactly}, not as a remainder.  The
ambient predicted by \cref{plt:thm:grp-scale-classification} with
\(\Gamma_0=\Gamma_1=\IntegerNumbers\) is \(\K[L_2][L,L^{-1}]((t))\), which is
where the answer lies.
\end{example}


\section{The two extreme scale changes: logarithm and exponential}

\begin{proposition}[Composition with \(\log\) is a level shift]
\label{plt:prop:grp-shift}
Let \(\Sigma(F)\DefinitionEquals F\compose\log\).  Then
\[
  t\compose\log=L_1^{-1},
  \qquad
  L\compose\log=L_2,
  \qquad
  \Sigma\Bigl(\sum_{n\ge n_0}p_n(L)t^n\Bigr)
  =\sum_{n\ge n_0}p_n(L_2)\,L_1^{-n},
\]
and \(\Sigma\) is an isomorphism of ordered rings from
\(\PLlau=\K[L_1]((L_0^{-1}))\) onto \(\K[L_2]((L_1^{-1}))\), carrying the
valuation \(\ordt\) to the \(L_1^{-1}\)-adic valuation and the dominance order
to the dominance order.  It is a differential isomorphism for the derivation
\(D_{L_1}\) of the target, i.e.\ \(\Sigma\circ\DX=D_{L_1}\circ\Sigma\); it is
\emph{not} compatible with \(\DX\) on the target, where instead
\begin{equation}\label{plt:eq:grp-shift-derivation}
  \DX\circ\Sigma=t\,\bigl(\Sigma\circ\DX\bigr).
\end{equation}
In particular \(\K[L_2]((L_1^{-1}))\) is not stable under \(\DX\).
\end{proposition}

\begin{proof}
The generator images are \eqref{plt:eq:grp-two-images} with \(G=\log X\).
Substituting them into \(F=\sum p_n(L)t^n\) gives the displayed formula, which
is a relabelling of indeterminates: it sends the topological generators
\(t^{\pm1},L_1\) of the source bijectively to the topological generators
\(L_1^{\mp1},L_2\) of the target, preserving sums, products, limits, the
lower-bounded integer support, and the lexicographic dominance order.  Hence
\(\Sigma\) is an isomorphism of ordered rings, and
\(\Sigma\circ\DX=D_{L_1}\circ\Sigma\) because both sides are continuous
\(\K\)-linear maps satisfying the same Leibniz rule and agreeing on \(t\) and
\(L_1\).  Since
\(\DX=\frac{\Differential L_1}{\Differential X}\,D_{L_1}=t\,D_{L_1}\) on the
target, formula \eqref{plt:eq:grp-shift-derivation} follows.  Finally
\(\DX(L_1^{-1})=-tL_1^{-2}\) involves \(t\), which does not occur in
\(\K[L_2]((L_1^{-1}))\).
\end{proof}

\begin{proposition}[Composition with \(\exp\) leaves every logarithmic class]
\label{plt:prop:grp-exponential}
One has \(t\compose\EulerE^{X}=\EulerE^{-X}\) and
\(L\compose\EulerE^{X}=X\), so that
\begin{equation}\label{plt:eq:grp-exp-composition}
  \Bigl(\sum_{n\ge n_0}p_n(L)t^n\Bigr)\compose\EulerE^{X}
  =\sum_{n\ge n_0}p_n(X)\,\EulerE^{-nX}.
\end{equation}
The monomial \(\EulerE^{-X}\) belongs to no \(\mathcal M_r(\Gamma)\) of
\cref{plt:def:grp-logtower}: for
\(\mathfrak m=\prod_{j\le r}L_j^{\alpha_j}\) one has
\(\log\mathfrak m=\BigOAt{L_1}{X\to+\infty}\), whereas
\(\log\EulerE^{-X}=-X\) and \(X/L_1\to+\infty\).  Moreover the loss is total on
the analytic side: if \(f\) is a function on some \([X_1,+\infty)\) admitting
\(F\) as an asymptotic expansion in the scale \(\SetOf{X^{-n}(\log X)^j}\) with
\(a=\ordt F\ge1\) and \(d=\degL p_a\), then
\begin{equation}\label{plt:eq:grp-exp-collapse}
  f\bigl(\EulerE^{X}\bigr)
  =\BigOAt{\EulerE^{-aX}X^{d}}{X\to+\infty}
  =\LittleOAt{X^{-N}}{X\to+\infty}
  \quad\text{for every }N,
\end{equation}
so the polynomial--logarithmic expansion of \(X\mapsto f(\EulerE^{X})\) is
identically zero although the function itself need not be.
\end{proposition}

\begin{proof}
The generator images are \eqref{plt:eq:grp-two-images} with \(G=\EulerE^{X}\),
and \eqref{plt:eq:grp-exp-composition} is their substitution.  For the
monomials, \(\log\mathfrak m=\alpha_0L_1+\alpha_1L_2+\cdots+\alpha_{r-1}L_r
+\alpha_r\log L_r\), each summand being \(\BigOAt{L_1}{X\to+\infty}\).  For the
last claim, the hypothesis gives \(f(Y)=p_a(\log Y)Y^{-a}
 +\BigOAt{Y^{-a-1}(\log Y)^{d'}}{Y\to+\infty}\) for a suitable \(d'\), hence
\(\AbsoluteValue{f(Y)}\le CY^{-a}(\log Y)^{d}\) for \(Y\) large; put
\(Y=\EulerE^{X}\).  Since \(a\ge1\), \(\EulerE^{-aX}X^{d}\) is smaller than
every power \(X^{-N}\).
\end{proof}

\begin{remark}
\eqref{plt:eq:grp-exp-composition} is the reason a formal exponential change of
scale must be met by enlarging the \emph{monomial} group, not the coefficient
ring: the objects \(\EulerE^{-nX}X^{j}\) are new transmonomials, and the
resulting field is the logarithmic--exponential transseries field of positive
exponential height \cite{adh-book,edgar-beginners,vanderhoeven-book}.  The
polynomial--logarithmic class is exactly its height-zero, depth-one corner.
\end{remark}

\section{The table of scale changes}

\begin{table}[htbp]
\centering\small
\caption{Scale changes and the resulting ambient class.  Throughout
\(A\in\PLpow\), \(U\in t\PLpow\), \(c\in\K^{\times}\) with a chosen
\(\log c\in\K\), and the regime is \(X\to+\infty\) on the positive ray.}
\label{plt:tab:grp-scale-changes}
\begin{tabular}{|p{0.20\textwidth}|p{0.32\textwidth}|p{0.36\textwidth}|}
\hline
Inner map \(G\) & \(t\compose G\); \ \(L\compose G\) &
Smallest ambient class for \(F\compose G\)\\
\hline
\(X+A\) &
\(t(1+tA)^{-1}\); \ \(L+\log(1+tA)\) &
\(\PLlau\); a group (\cref{plt:thm:grp-group})\\
\hline
\(cX\) &
\(c^{-1}t\); \ \(L+\log c\) &
\(\PLlau\) (\cref{plt:prop:grp-dilation})\\
\hline
\(X+cL\) &
\(t(1+cLt)^{-1}\); \ \(L+\log(1+cLt)\) &
\(\PLlau\); \(\degL\) grows linearly
(\cref{plt:prop:grp-logdegree})\\
\hline
\(cX^{\rho}(1+U)\), \(\rho\in\PositiveIntegers\) &
\(c^{-1}t^{\rho}(1+U)^{-1}\); \ \(\rho L+\log c+\log(1+U)\) &
\(\PLlau\) (\cref{plt:thm:grp-scale-classification}(1))\\
\hline
\(X^{r/s}\), \(\gcd(r,s)=1\), \(s>1\) &
\(t^{r/s}\); \ \((r/s)L\) &
Puiseux-log \(\K[L]((t^{1/s}))\); not \(\PLlau\)\\
\hline
\(X^{\rho}\), \(\rho>0\) irrational &
\(t^{\rho}\); \ \(\rho L\) &
Hahn \(\mathcal H_1(\K,\rho\IntegerNumbers+\IntegerNumbers)\); no Puiseux
extension suffices\\
\hline
\(XL^{\sigma}\), \(\sigma\ne0\) &
\(tL^{-\sigma}\); \ \(L+\sigma L_2\) &
depth two: \(\mathcal H_2\); for
\(\sigma\in\IntegerNumbers\), \(\K[L_2][L,L^{-1}]((t))\)\\
\hline
\(cX^{\rho}L^{\sigma}(1+U)\), \(\rho>0\), \(\sigma\ne0\) &
\(c^{-1}t^{\rho}L^{-\sigma}(1+U)^{-1}\); \
\(\rho L+\sigma L_2+\log c+\log(1+U)\) &
\(\mathcal H_2(\K,\Gamma_0+\Gamma_1)\)
(\cref{plt:thm:grp-scale-classification})\\
\hline
\(\log X\) &
\(L^{-1}\); \ \(L_2\) &
\(\K[L_2]((L^{-1}))\): the level shift \(\Sigma\)
(\cref{plt:prop:grp-shift})\\
\hline
\(\EulerE^{X}\) &
\(\EulerE^{-X}\); \ \(X\) &
outside every \(\mathcal H_r\): logarithmic--exponential transseries
(\cref{plt:prop:grp-exponential})\\
\hline
\(cX^{\rho}\), \(\rho\le0\) &
--- &
inadmissible: \(G\) does not tend to \(+\infty\)\\
\hline
\(cX\), \(c<0\), \(\K\subseteq\RealNumbers\) &
--- &
inadmissible on the positive ray: \(\log c\notin\RealNumbers\); extend \(\K\)
and fix a sector\\
\hline
\(X=1/x\) (regime change) &
\(t\mapsto x\); \ \(L\mapsto-\log x\) &
\(\K[\Lambda]((x))\), \(\Lambda=-\log x\): an isomorphism, not an
endomorphism (\cref{plt:prop:grp-zero-dictionary})\\
\hline
\end{tabular}
\end{table}

\begin{warningbox}
The last three rows are not failures of a formula but failures of a
\emph{hypothesis}.  An expression may be perfectly meaningful while lying
outside \(\PLlau\); and conversely a formally computed composite is
meaningless if the inner map does not carry the outer expansion's regime.
Closure is a statement about a chosen representation class, never about
existence of the composed function.  In particular the whole table is
directional: that \(F\compose G\) is defined and lies in a given class says
nothing about \(G\compose F\).
\end{warningbox}

\section{The nested logarithmic scale}

\begin{proposition}[Derivation of the logarithmic tower]
\label{plt:prop:grp-tower-derivation}
With \(L_0=X\) and \(L_{j+1}=\log L_j\),
\begin{equation}\label{plt:eq:grp-tower-derivative}
  \DX L_j=\frac{t}{L_1L_2\cdots L_{j-1}}
  \qquad(j\ge1),
\end{equation}
the product being empty for \(j=1\).  Put
\[
  \mathcal D_r
  \DefinitionEquals
  \partial_{L_1}
  +\sum_{j=2}^{r}\frac{1}{L_1L_2\cdots L_{j-1}}\,\partial_{L_j},
  \qquad
  \mathcal L_r
  \DefinitionEquals
  \K\bigl[L_1^{\pm1},\ldots,L_{r-1}^{\pm1},L_r\bigr].
\]
Then \(\mathcal D_r(\mathcal L_r)\subseteq\mathcal L_r\), and for
\(n\in\IntegerNumbers\), \(p\in\mathcal L_r\),
\begin{equation}\label{plt:eq:grp-tower-block}
  \DX\bigl(t^np\bigr)=t^{n+1}\bigl(\mathcal D_r-n\bigr)p .
\end{equation}
By contrast, for \(r\ge2\) the polynomial ring
\(\K[L_1,\ldots,L_r]\) is \emph{not} stable under \(\mathcal D_r\):
\(\mathcal D_r L_2=L_1^{-1}\).
\end{proposition}

\begin{proof}
\eqref{plt:eq:grp-tower-derivative} holds for \(j=1\) since \(\DX L_1=t\), and
inductively \(\DX L_{j+1}=(\DX L_j)/L_j\).  Stability: \(\partial_{L_1}\)
maps \(\mathcal L_r\) into \(\mathcal L_r\) because the ring is Laurent in
\(L_1\); for \(2\le j\le r\), \(\partial_{L_j}\) lowers the \(L_j\)-degree and
the prefactor \((L_1\cdots L_{j-1})^{-1}\) lies in \(\mathcal L_r\) --- here it
matters that only \(L_1,\ldots,L_{r-1}\) need be inverted, since
\(\partial_{L_r}\) is multiplied by \((L_1\cdots L_{r-1})^{-1}\).  Formula
\eqref{plt:eq:grp-tower-block} follows from
\(\DX t^n=-nt^{n+1}\) and \eqref{plt:eq:grp-tower-derivative} by the Leibniz
rule.
\end{proof}

% ed.: S2 says only that "the rational field K(L_1,...,L_r), or a
% Laurent-monomial field, repairs this defect".  Proposition
% \ref{plt:prop:grp-tower-derivation} identifies the minimal repair inside the
% Laurent-monomial family and shows that L_r itself need not be inverted.

\begin{theorem}[The near-identity group at logarithmic depth \(r\)]
\label{plt:thm:grp-multilog-group}
Let \(r\ge1\) and
\[
  \mathcal N_r\DefinitionEquals\SetOf{X+A:\ A\in\mathcal L_r[[t]]},
  \qquad \mathcal N_1=\Gnear .
\]
Then \(\mathcal N_r\) is a group under composition, and every statement of
\cref{plt:sec:near-identity-group} holds verbatim with \(\K[L]\) replaced by
\(\mathcal L_r\), \(\partial_L\) by \(\mathcal D_r\), and \(\shiftop{n}{k}\) by
\(\prod_{j=0}^{k-1}(\mathcal D_r-n-j)\).  The corresponding statement for
\(\SetOf{X+A:A\in\K[L_1,\ldots,L_r][[t]]}\) is false for \(r\ge2\).
\end{theorem}

\begin{proof}
Fix \(G=X+A\) with \(A\in\mathcal L_r[[t]]\) and define \(\Phi_G\) on the
topological generators \(t^{\pm1},L_1^{\pm1},\ldots,L_{r-1}^{\pm1},L_r\) of
\(\mathcal L_r((t))\) by \(\Phi_G(t)=t(1+tA)^{-1}\),
\(\Phi_G(L_{j+1})=\log\bigl(\Phi_G(L_j)\bigr)\), starting from
\(\Phi_G(L_0)=X(1+tA)\).  Set \(w_0=tA\) and, for \(j\ge1\),
\begin{equation}\label{plt:eq:grp-tower-w}
  w_j\DefinitionEquals L_j^{-1}\log(1+w_{j-1}).
\end{equation}
An induction using \cref{plt:lem:grp-explog} gives
\(w_j\in t\mathcal L_r[[t]]\) and
\[
  \Phi_G(L_j)=L_j\bigl(1+w_j\bigr)
  \qquad(0\le j\le r-1),
  \qquad
  \Phi_G(L_r)=L_r+\log(1+w_{r-1}),
\]
since \(\log\bigl(L_j(1+w_j)\bigr)=L_{j+1}+\log(1+w_j)
 =L_{j+1}\bigl(1+L_{j+1}^{-1}\log(1+w_j)\bigr)\).  The construction inverts
exactly \(L_1,\ldots,L_{r-1}\) and never \(L_r\); this is what forces the
coefficient ring \(\mathcal L_r\) and no smaller one.  Extending \(\Phi_G\)
multiplicatively and \(t\)-adically reproduces
\cref{plt:prop:grp-substitution} word for word, and
\eqref{plt:eq:grp-tower-block} reproduces \(\DX(\mathcal L_r[[t]])\subseteq
t\mathcal L_r[[t]]\), which is the only property of \(\DX\) used in
\cref{plt:thm:grp-taylor} and \cref{plt:cor:grp-gain}.

Two points genuinely need rechecking, and both go through.  First, the
\emph{gain estimate}: if \(\ordt(B-C)=m\), then comparing
\eqref{plt:eq:grp-tower-w} for \(G=X+B\) and \(G'=X+C\) gives
\(\ordt(w_0-w_0')\ge m+1\) and, inductively,
\(\ordt(w_j-w_j')\ge m+1\), because
\(\log(1+w)-\log(1+w')=\log\bigl(1+\frac{w-w'}{1+w'}\bigr)\) and \(L_j^{-1}\)
is a unit of \(\mathcal L_r\) for \(j\le r-1\).  Hence every generator except
\(X\) itself has \(\Phi_G\)- and \(\Phi_{G'}\)-images differing by an element of
\(t^{m+1}\mathcal L_r[[t]]\), and expanding
\(\Phi_G(a_nt^n)-\Phi_{G'}(a_nt^n)\) by
\(PQ-P'Q'=(P-P')Q+P'(Q-Q')\) yields
\(\ordt\bigl(A\compose G-A\compose G'\bigr)\ge m+1\).  Second,
\emph{associativity}: \(\Phi_H(\Phi_G(L_0))=\Phi_H(G)=G\compose H
 =\Phi_{G\compose H}(L_0)\) as in \eqref{plt:eq:grp-phi-of-G}, and if the
identity \(\Phi_H\Phi_G(L_j)=\Phi_{G\compose H}(L_j)\) holds then applying the
continuous algebra homomorphism \(\Phi_H\) to
\(\Phi_G(L_{j+1})=L_{j+1}+\log(1+w_j)\) and using
\(\Phi_H(\log(1+u))=\log(1+\Phi_H(u))\) together with
\(\log(uv)=\log u+\log v\) gives the identity at level \(j+1\); the case
\(j+1=1\) is the computation displayed in the proof of
\cref{plt:thm:grp-group}(2).  With these, the fixed-point construction of
inverses (\cref{plt:thm:rev-existence}, whose proof uses only the gain
estimate \cref{plt:cor:grp-gain} and \(t\)-adic completeness, both available
verbatim in \(\mathcal L_r[[t]]\)) and the monoid argument are unchanged.

For the final assertion take \(r=2\), \(G=X+1\in\mathcal N_2\) and
\(F=L_2\in\K[L_1,L_2]\).  Then
\[
  F\compose G=\log\log(X+1)
  =L_2+\log\Bigl(1+L_1^{-1}\log(1+t)\Bigr)
  =L_2+L_1^{-1}t+\FormalBigOAt{t^2}{t},
\]
whose block of outer order \(1\) is \(L_1^{-1}\notin\K[L_1,L_2]\).  Hence
\(\SetOf{X+A:A\in\K[L_1,L_2][[t]]}\) is not closed under composition.
\end{proof}

\begin{proposition}[The unit criterion changes shape with depth]
\label{plt:prop:grp-multilog-units}
\(\mathcal L_r^{\times}
 =\SetOf{c\,L_1^{a_1}\cdots L_{r-1}^{a_{r-1}}:\ c\in\K^{\times},\
  a_i\in\IntegerNumbers}\).
Consequently a nonzero \(F=\sum_{n\ge n_0}p_nt^n\in\mathcal L_r((t))\) is a
unit if and only if its leading block \(p_{n_0}\) is a nonzero scalar times a
Laurent monomial in \(L_1,\ldots,L_{r-1}\).  For \(r=1\) this is the classical
criterion ``\(p_{n_0}\in\K^{\times}\)''; for \(r\ge2\) strictly more elements
are invertible, for instance \(L_1^{-1}+\FormalBigOAt{t}{t}\).
\end{proposition}

\begin{proof}
Write \(\mathcal L_r=R[L_r]\) with
\(R=\K[L_1^{\pm1},\ldots,L_{r-1}^{\pm1}]\).  Since \(R\) is an integral domain,
\((R[L_r])^{\times}=R^{\times}\); and by induction on the number of variables,
the units of a Laurent polynomial ring over a domain are the units of the base
times monomials.  The criterion for \(\mathcal L_r((t))\) is the standard one
for a Laurent series ring over a commutative ring: if \(p_{n_0}\) is a unit,
\(F=t^{n_0}p_{n_0}(1+v)\) with \(v\in t\mathcal L_r[[t]]\), and \((1+v)^{-1}\)
is given by the geometric series; conversely \(FF^{-1}=1\) forces
\(p_{n_0}\,[t^{-n_0}]F^{-1}=1\).
\end{proof}

\section{Lifting a one-logarithm result, and where the lift fails}

\begin{proposition}[Transport by the level shift]\label{plt:prop:grp-transport}
Let \(\mathcal S\) be any statement about the structure
\(\bigl(\PLlau;\ t,\ L,\ \ordt,\ \preceqdom,\ \DX\bigr)\) expressed purely in
that language.  Then \(\mathcal S\) holds if and only if the corresponding
statement holds for
\(\bigl(\K[L_2]((L_1^{-1}));\ L_1^{-1},\ L_2,\
 \nu_{L_1^{-1}},\ \preceqdom,\ D_{L_1}\bigr)\).  In particular
\[
  \SetOf{L_1+B:\ B\in\K[L_2][[L_1^{-1}]]}
\]
is a group under composition in the variable \(L_1\), with all the structure of
\cref{plt:thm:grp-group}, \cref{plt:prop:grp-congruence} and
\cref{plt:thm:grp-gradedlie}.
\end{proposition}

\begin{proof}
\cref{plt:prop:grp-shift} exhibits \(\Sigma\) as an isomorphism of the two
structures.
\end{proof}

\begin{remark}[The correct fast variable]
This is the precise content of the practical rule that a nested expansion
should be read in the right variable.  A series
\(\sum_{n\ge0}p_n(\log\log X)(\log X)^{-n}\), which is a depth-two object in
\(X\), is a \emph{one-logarithm} object in the fast variable
\(T=L_1=\log X\), with \(\Lambda=L_2=\log T\).  Choosing the wrong fast
variable makes such a problem look as though it fell outside the theory.
\end{remark}

\begin{warningbox}
Three things do \emph{not} lift.
\begin{itemize}
\item \emph{Statements about \(\DX\).}  By
\eqref{plt:eq:grp-shift-derivation}, \(\Sigma\) intertwines \(\DX\) with
\(D_{L_1}\), not with \(\DX\); and \(\K[L_2]((L_1^{-1}))\) is not even stable
under \(\DX\).  A result whose statement mentions the derivation of the
\emph{original} variable must be re-derived, not transported.
\item \emph{Genuinely mixed depth.}  \(\Sigma\) reaches only the shifted copy.
A series such as \(\sum_np_n(L_1,L_2)t^n\) with genuinely bivariate blocks lies
in neither \(\PLlau\) nor \(\Sigma(\PLlau)\), and no change of fast variable
makes it single-logarithmic; one must work in \(\mathcal L_r((t))\) and accept
\cref{plt:prop:grp-tower-derivation}.
\item \emph{The unit criterion.}  Its shape changes at depth \(\ge2\)
(\cref{plt:prop:grp-multilog-units}), so any argument that used
``leading block a nonzero scalar'' must be restated as ``leading block a unit
of the coefficient ring''.
\end{itemize}
\end{warningbox}

\section{From infinity to a finite point}

\begin{proposition}[The infinity--zero dictionary]
\label{plt:prop:grp-zero-dictionary}
Put \(x=X^{-1}\) and
\[
  \Lambda\DefinitionEquals\log\frac1x=-\log x=\log X=L .
\]
Then \(F\mapsto F(1/x)\) is an isomorphism of ordered rings
\[
  \iota:\ \K[L]((t))\longrightarrow\K[\Lambda]((x)),
  \qquad t\longmapsto x,\quad L\longmapsto\Lambda,
\]
matching the regime \(X\to+\infty\) with \(x\to0^{+}\) and the dominance order
with the order of decreasing size as \(x\to0^{+}\).  On the differential side
\begin{equation}\label{plt:eq:grp-zero-derivation}
  \frac{\Differential}{\Differential x}=-X^{2}\DX,
  \qquad
  \frac{\Differential}{\Differential x}
   =\partial_x-\frac1x\,\partial_\Lambda,
  \qquad
  \frac{\Differential}{\Differential x}\bigl(x^{n}P(\Lambda)\bigr)
   =x^{n-1}\bigl(nP-P'\bigr),
\end{equation}
the exact mirror of \eqref{plt:eq:grp-block-derivative}.  Under \(\iota\) the
group \(\Gnear\) corresponds to the group of germs
\[
  \SetOf{\,g(x)=\frac{x}{1+x\,a(x)}\ :\ a\in\K[\Lambda][[x]]\,}
\]
with the composition law of \(\Gnear\), the isomorphism being
\(G\mapsto j\compose G\compose j\) with \(j(x)=1/x\).
\end{proposition}

\begin{proof}
The map \(\iota\) is the relabelling \(t\mapsto x\), \(L\mapsto\Lambda\); it is
a ring isomorphism, and it preserves the order because
\(t^nL^j\succdom t^mL^k\) and \(x^n\Lambda^j\succdom x^m\Lambda^k\) are
governed by the same inequalities on \((n,j)\) and \((m,k)\).  From
\(x=X^{-1}\), \(\frac{\Differential x}{\Differential X}=-X^{-2}\), so
\(\DX=-x^{2}\frac{\Differential}{\Differential x}\), which is the first
identity; the second is the chain rule with \(\partial_x\Lambda=-1/x\); the
third follows:
\(\frac{\Differential}{\Differential x}(x^nP)=nx^{n-1}P-x^{n}\frac1xP'\).
Finally \(j\) is an involution, so \(G\mapsto j\compose G\compose j\) is a group
isomorphism, and for \(G=X+A\),
\(j\compose G\compose j=\MultiplicativeInverseOf{\bigl(x^{-1}+A(1/x)\bigr)}
 =x\bigl(1+x\,\iota(A)\bigr)^{-1}\).
\end{proof}

\begin{warningbox}
The dictionary is \(L=-\log x\), never \(L=\log x\).  Substituting the wrong
sign reverses \(\partial_L\) and turns the block operator
\(\shiftop{n}{k}=\prod_{j=0}^{k-1}(\partial_L-n-j)\) into a different operator;
all reversion coefficients then come out with wrong signs while remaining
formally consistent, which makes the error hard to detect.
\end{warningbox}

\begin{remark}[What is deferred]
\cref{plt:prop:grp-zero-dictionary} is only the integer-exponent dictionary.
Local dynamics at a finite point uses \emph{Dulac} classes: real or well-ordered
exponent sets \(x^{\alpha}\), the generator \(\ell=-1/\log x\) in place of
\(\Lambda\), and their iterated versions, for which composition and normal
forms require the machinery of
\cite{mardesic-normalforms,mardesic-length,resman-dulac,peran-logtransseries}.
That development, including the change between \(\Lambda\) and \(\ell\) and the
consequences of admitting real exponents, is carried out systematically in
\cref{plt:sec:dulac}.  Note also that the other natural passage to a finite
point, \(X=-\log x\), is the exponential change of scale of
\cref{plt:prop:grp-exponential} and leaves the class.
\end{remark}

\begin{summarybox}
A substitution is decided by the two generator images
\(\MultiplicativeInverseOf{G}\) and \(\log G\), and by nothing else.  For
\(G=cX^{\rho}L^{\sigma}(1+U)\) with \(\rho>0\): integral \(\rho\) and
\(\sigma=0\) keep the composite in \(\PLlau\); rational \(\rho=r/s\) forces
exactly the Puiseux-log extension \(\K[L]((t^{1/s}))\); irrational \(\rho\)
forces a Hahn exponent group; \(\sigma\ne0\) forces the second logarithm
\(L_2\) precisely when some coefficient block is nonconstant, together with
\(L\)-exponents outside \(\NaturalNumbers\).  Composition with \(\log\) is an
isomorphism onto the level-shifted copy --- the formal statement of the
``correct fast variable'' rule --- while composition with \(\exp\) leaves every
logarithmic class and annihilates the polynomial--logarithmic expansion
entirely.  At depth \(r\) the near-identity group survives, but only over the
coefficient ring \(\K[L_1^{\pm1},\ldots,L_{r-1}^{\pm1},L_r]\) that the
derivation and the generator substitutions force.
\end{summarybox}

\clearpage
\part{Series reversal at infinity}

% ---- fragment 08-reversion ----
\chapter{Existence and uniqueness of the near-identity inverse}
\label{plt:sec:reversion-existence}

\section{Conventions, and what is imported from the composition part}
\label{plt:sec:rev-conventions}

Throughout this part \(\K\) is a field of characteristic zero, the large
variable is \(X\) with
\[
 t\DefinitionEquals X^{-1},
 \qquad
 L\DefinitionEquals\log X,
\]
and the default analytic regime, whenever one is invoked at all, is
\(X\to+\infty\) on the positive real ray.  We work in
\[
 \PLpow=\K[L][[t]],
 \qquad
 \PLlau=\K[L]((t)),
 \qquad
 \PLrat=\K(L)((t)),
\]
with the outer valuation \(\ordt\) and the exclusive truncation
\(\Trunc{F}{N}=\sum_{n<N}p_n(L)t^n\).  The near-identity group is
\begin{equation}\label{plt:eq:rev-gnear}
 \Gnear\DefinitionEquals\SetOf{X+A:A\in\PLpow}
 =\SetOf{F\in\PLlau:\ \ordt(F)\ge-1\ \text{and}\ [t^{-1}]F=1}.
\end{equation}

\begin{definition}[Truncation and blocks of a near-identity map]
\label{plt:def:rev-trunc-gnear}
For \(G=X+B\in\Gnear\) with \(B=\sum_{n\ge0}b_n(L)t^n\) we call \(b_n\) the
\emph{\(n\)-th inverse block} of \(G\) and set
\[
 \Trunc{G}{N}\DefinitionEquals X+\Trunc{B}{N}=X+\sum_{n<N}b_n(L)t^n .
\]
For \(F\in\Gnear\), the phrase ``\(G\) is correct through outer block \(N\)''
means
\[
 \Trunc{G}{N+1}=\Trunc{\rev{F}}{N+1},
\]
that is, agreement in the \(N+1\) blocks \(b_0,\dots,b_N\).
\end{definition}

\begin{remark}[Inclusive versus exclusive cutoffs]
\label{plt:rmk:rev-cutoff-collision}
% ed.: report 6 states its reversion algorithms with the inclusive cutoff
% [F]_{\le N}; every cutoff index taken from that source is shifted by one here.
Report~6 writes its truncations inclusively as \([F]_{\le N}\).  The
translation is \([F]_{\le N}=\Trunc{F}{N+1}\), so every cutoff index borrowed
from that source has been raised by one below.  The place where this matters
most is the finite certificate (\cref{plt:cor:rev-finite-certificate}):
correctness \emph{through} block \(N\) is the exclusive cutoff \(<N+1\), and
the certifying residual is a \(\FormalBigOAt{t^{N+1}}{t}\), not a
\(\FormalBigOAt{t^{N}}{t}\).
\end{remark}

\begin{remark}[Standing facts imported from the composition part]
\label{plt:rmk:rev-imports}
Four facts about \(\PLlau\), \(\DX\) and \(\compose\) are used constantly.
They are proved elsewhere in this volume; we record them in the exact form
needed here and prove nothing about them except the completeness statement,
whose precise shape is used below.

\emph{(I1) Block derivative} (\cref{plt:thm:dif-repeated}).  For \(n\in\IntegerNumbers\) and \(p\in\K[L]\),
\begin{equation}\label{plt:eq:rev-block-derivative}
 \DX\bigl(t^np(L)\bigr)=t^{n+1}\bigl(\partial_L-n\bigr)p(L),
 \qquad
 \DX^k\bigl(t^np(L)\bigr)=t^{n+k}\shiftop{n}{k}p(L),
\end{equation}
where \(\shiftop{n}{k}=\prod_{j=0}^{k-1}(\partial_L-n-j)\) and
\(\shiftop{n}{0}=1\) is the empty product.

\emph{(I2) Exact formal Taylor formula} (\cref{plt:thm:tay-formula}).  For \(S\in\PLlau\) and
\(P\in\PLpow\) the family \(\bigl(P^k\DX^kS/k!\bigr)_{k\ge0}\) is formally
summable and
\begin{equation}\label{plt:eq:rev-taylor}
 S\compose(X+P)=\sum_{k\ge0}\frac{P^k}{k!}\,\DX^kS .
\end{equation}
This is an identity of formal series; it carries no analytic content.

% ed.: the closure of \(\Gnear\) under composition was stated and proved a
% ed.: second time in this part, as a separate lemma.  It is item (1) of the
% ed.: group theorem, proved there from the substitution proposition, so it is
% ed.: imported here with the other three standing facts.  Note that only
% ed.: items (1)--(3) of that theorem are imported: its item (4) cites
% ed.: \cref{plt:thm:rev-existence} below, and importing it here would be
% ed.: circular.
\emph{(I3) Closure, associativity and chain rule}
(\cref{plt:thm:grp-group}\,(1)--(3), \cref{plt:thm:cmp-chain-rule}).  For
\(F=X+A\) and \(G=X+B\) in \(\Gnear\),
\begin{equation}\label{plt:eq:rev-closure}
 F\compose G=X+B+A\compose G\in\Gnear ;
\end{equation}
\((S\compose H_1)\compose H_2
=S\compose(H_1\compose H_2)\) whenever all three displayed composites exist in
the ambient class; this covers near-identity inner arguments and monomial
inner arguments \(cX\) for which a value of \(\log c\) lies in \(\K\).  The
series \(X\) is a two-sided identity for \(\compose\), and
\begin{equation}\label{plt:eq:rev-chain-rule}
 \DX(S\compose H)=\bigl((\DX S)\compose H\bigr)\cdot\DX H .
\end{equation}

\emph{(I4) Completeness} (\cref{plt:thm:trunc-completeness}).  \(\PLpow\) is complete for the \(t\)-adic
filtration: if \((F_j)\) is Cauchy then for each \(n\) the block
\([t^n]F_j\in\K[L]\) is eventually constant, say \(p_n\); the series
\(F=\sum_{n\ge0}p_nt^n\) lies in \(\PLpow\) --- no uniform bound on
\(\deg_Lp_n\) is required --- and \(F_j\to F\).  On \(\PLlau\) the same
argument needs the extra hypothesis that the \(\ordt(F_j)\) are bounded below,
without which a Cauchy sequence need not define a Laurent series with a finite
principal part.
\end{remark}

\begin{lemma}[Valuation gain of the derivation]
\label{plt:lem:rev-deriv-gain}
For every \(S\in\PLlau\) and every \(k\in\NaturalNumbers\),
\(\ordt(\DX^kS)\ge\ordt(S)+k\).
\end{lemma}

\begin{proof}
By \eqref{plt:eq:rev-block-derivative}, \(\DX^k\) sends the block \(t^np\) to
\(t^{n+k}\shiftop{n}{k}p\), hence maps \(t^m\PLpow\) into \(t^{m+k}\PLpow\) for
every \(m\in\IntegerNumbers\).  Apply this with \(m=\ordt(S)\).  Equality may
fail, since \(\shiftop{n}{k}p\) can vanish; for instance
\(\DX(t^0\cdot1)=0\).
\end{proof}

\begin{lemma}[Right composition with a near-identity preserves leading data]
\label{plt:lem:rev-order-preserved}
Let \(S\in\PLlau\), \(P\in\PLpow\), and \(m=\ordt(S)\).  Then
\[
 \ordt\bigl(S\compose(X+P)\bigr)=m,
 \qquad
 [t^{m}]\bigl(S\compose(X+P)\bigr)=[t^{m}]S .
\]
In particular \(S\compose(X+P)=0\) if and only if \(S=0\), and right
composition by an element of \(\Gnear\) carries units of \(\PLpow\) to units of
\(\PLpow\).
\end{lemma}

\begin{proof}
In \eqref{plt:eq:rev-taylor} the term \(k=0\) is \(S\).  For \(k\ge1\),
\(\ordt\bigl(P^k\DX^kS/k!\bigr)\ge k\,\ordt(P)+m+k\ge m+k>m\), using
\cref{plt:lem:rev-deriv-gain} and \(\ordt(P)\ge0\).  So no earlier block is
created and the block at level \(m\) is unchanged.
\end{proof}

% ed.: the closure lemma stated here duplicated \cref{plt:thm:grp-group}\,(1),
% ed.: which is proved in the composition-group part with the same one-line
% ed.: argument.  It is deleted; closure is now imported as part of (I3) and
% ed.: appears as \eqref{plt:eq:rev-closure}.  For the record, the argument is
% ed.: that \(F\compose G=G+A\compose G\) and that
% ed.: \(A\compose G=\sum_{k\ge0}B^k\DX^kA/k!\) has all summands in
% ed.: \(\PLpow\) by \cref{plt:lem:rev-deriv-gain}.

The next lemma re-expands a composition around a shifted base point without
inverting anything.  It may therefore be used inside the existence proof, and
it is the engine of every estimate in this part.

\begin{lemma}[Two-point Taylor expansion]
\label{plt:lem:rev-shifted-taylor}
Let \(S\in\PLlau\) and \(P,Q\in\PLpow\).  Then the family
\(\bigl(Q^j(\DX^jS)\compose(X+P)/j!\bigr)_{j\ge0}\) is formally summable and
\begin{equation}\label{plt:eq:rev-shifted-taylor}
 S\compose\bigl(X+P+Q\bigr)
 =\sum_{j\ge0}\frac{Q^j}{j!}\,\bigl(\DX^jS\bigr)\compose(X+P).
\end{equation}
\end{lemma}

\begin{proof}
Summability: by \cref{plt:lem:rev-deriv-gain,plt:lem:rev-order-preserved},
\(\ordt\bigl(Q^j(\DX^jS)\compose(X+P)/j!\bigr)\ge j\,\ordt(Q)+\ordt(S)+j
\to\infty\).

Apply \eqref{plt:eq:rev-taylor} with \(P+Q\in\PLpow\) in place of \(P\):
\[
 S\compose(X+P+Q)
 =\sum_{k\ge0}\frac{(P+Q)^k}{k!}\DX^kS
 =\sum_{k\ge0}\ \sum_{i+j=k}\frac{P^i}{i!}\frac{Q^j}{j!}\DX^{i+j}S .
\]
Each term has \(\ordt\ge\ordt(S)+i+j\), so the double family is summable and
may be summed in either order.  Grouping by \(j\) and applying
\eqref{plt:eq:rev-taylor} once more, now to \(\DX^jS\),
\[
 \sum_{j\ge0}\frac{Q^j}{j!}\sum_{i\ge0}\frac{P^i}{i!}\DX^i\bigl(\DX^jS\bigr)
 =\sum_{j\ge0}\frac{Q^j}{j!}\bigl(\DX^jS\bigr)\compose(X+P),
\]
which is \eqref{plt:eq:rev-shifted-taylor}.
\end{proof}

\section{The inverse equation as a fixed point}
\label{plt:sec:rev-fixed-point}

\begin{proposition}[The reversion equation]
\label{plt:prop:rev-fixed-point-equation}
Let \(F=X+A\in\Gnear\) and \(G=X+B\in\Gnear\).  Then \(F\compose G=X\) if and
only if \(B=-A\compose(X+B)\).
\end{proposition}

\begin{proof}
By \eqref{plt:eq:rev-closure}, \(F\compose G=X+B+A\compose G\), which equals
\(X\) exactly when \(B+A\compose(X+B)=0\).
\end{proof}

Thus \(B\) is a fixed point of
\begin{equation}\label{plt:eq:rev-Phi}
 \Phi\colon\PLpow\to\PLpow,
 \qquad
 \Phi(B)\DefinitionEquals-A\compose(X+B).
\end{equation}
The map \(\Phi\) is neither linear nor defined on all of \(\PLlau\): its
argument must be an admissible near-identity displacement.

\begin{remark}[Variable names]
\label{plt:rmk:rev-variable-names}
Reports~3 and~6 write the reversion equation in a second letter,
\(B(Y)=-A(Y+B(Y))\), to separate visually the argument of the inverse from the
argument of \(F\).  This is a typographic device only: \(X\) is a free
generator of the formal algebra and the renaming is the identity map.  We use
\(X\) throughout, in agreement with the notation catalogue.
\end{remark}

% ed.: this estimate was proved twice in the volume: here from the two-point
% ed.: Taylor expansion, and in the composition-group part from the exact
% ed.: Taylor formula, where it is \eqref{plt:eq:grp-gain-lipschitz}.  The two
% ed.: arguments are the same computation, and the group-part statement was
% ed.: missing the summand \(\ordt(A)\) that its own proof produces; that
% ed.: omission is repaired at \cref{plt:cor:grp-gain} and the duplicate proof
% ed.: is deleted here.  The inequality is recorded below as a fifth standing
% ed.: import, in the instantiated form used throughout this part.
\begin{remark}[Filtered Lipschitz gain: the fifth standing import]
\label{plt:lem:rev-gain}
Let \(A\in\PLpow\) with \(\alpha\DefinitionEquals\ordt(A)\), and let
\(B,C\in\PLpow\).  Then
\begin{equation}\label{plt:eq:rev-gain}
 \ordt\bigl(A\compose(X+B)-A\compose(X+C)\bigr)
 \ \ge\ \alpha+\ordt(B-C)+1 ,
\end{equation}
which is \eqref{plt:eq:grp-gain-lipschitz} of \cref{plt:cor:grp-gain}, proved
there.  In particular, for the map \(\Phi\) of \eqref{plt:eq:rev-Phi},
\(\ordt\bigl(\Phi(B)-\Phi(C)\bigr)\ge\alpha+\ordt(B-C)+1\).  The reader who
prefers an argument internal to this part may instead put \(S=A\), \(P=C\),
\(Q=B-C\) in \cref{plt:lem:rev-shifted-taylor} and bound the \(j\)-th term by
\(\alpha+j(\ordt(Q)+1)\) using
\cref{plt:lem:rev-order-preserved,plt:lem:rev-deriv-gain}; the minimum over
\(j\ge1\) is at \(j=1\).
\end{remark}

\begin{remark}[What ``a gain of one block'' means, precisely]
\label{plt:rmk:rev-gain-meaning}
% ed.: all six sources call this a contraction; none says in what sense.  The
% statement is about nu_t, not about coefficient size, and the constant term
% alpha = nu_t(A) is absent from every source formulation.
All six reports describe \eqref{plt:eq:rev-gain} as a contraction.  Three
points must be kept apart.

\begin{enumerate}
\item \emph{It is a statement about \(\ordt\) alone.}  If \(B\) and \(C\) agree
in the blocks \(0,\dots,r-1\), then \(\Phi(B)\) and \(\Phi(C)\) agree in the
blocks \(0,\dots,r\): one further block of the output is pinned down.  Nothing
whatever is asserted about the size, the coefficients, or the logarithmic
degree of the blocks that still differ; \(\deg_L\) is not contracted and in
general grows.
\item \emph{It is a genuine metric contraction for the standard \(t\)-adic
absolute value.}  Setting \(\AbsoluteValue{F}_t=2^{-\ordt(F)}\) and
\(\AbsoluteValue{0}_t=0\), inequality \eqref{plt:eq:rev-gain} says exactly that
\(\Phi\) is Lipschitz with constant \(2^{-1-\alpha}\le\tfrac12\).  The space
\((\PLpow,\AbsoluteValue{\cdot}_t)\) is complete by (I4), so the Banach
fixed-point theorem applies verbatim.  We nevertheless give the
coefficient-stabilization proof below, because that is what an implementation
performs and because it yields the sharp iteration rate.
\item \emph{The gain is at least one, and in general exactly one.}  Take
\(A=L\), \(B=0\) and \(C=ct^r\) with \(c\in\K^\times\), \(r\ge0\).  Then
\(A\compose(X+C)-A\compose(X+B)=\log(1+ct^{r+1})
=ct^{r+1}+\FormalBigOAt{t^{r+2}}{t}\), and \eqref{plt:eq:rev-gain} is an
equality.  The summand \(\alpha=\ordt(A)\) in \eqref{plt:eq:rev-gain} does not
appear in any of the six sources; it is what makes the sharp Newton estimate of
\cref{plt:thm:rev-newton-doubling} come out.
\end{enumerate}
\end{remark}

\section{Existence, uniqueness, and the two-sided property}
\label{plt:sec:rev-existence}

\begin{theorem}[Near-identity inverse theorem]
\label{plt:thm:rev-existence}
Let \(\K\) have characteristic zero and let \(F=X+A\in\Gnear\) with
\(\alpha=\ordt(A)\).  Then:
\begin{enumerate}
\item the Picard iterates \(B_0=0\), \(B_{j+1}=\Phi(B_j)=-A\compose(X+B_j)\)
satisfy
\begin{equation}\label{plt:eq:rev-picard-rate}
 \ordt\bigl(B_{j+1}-B_j\bigr)\ \ge\ (j+1)\alpha+j
 \qquad(j\ge0),
\end{equation}
and converge \(t\)-adically to some \(B\in\PLpow\) with
\begin{equation}\label{plt:eq:rev-picard-accuracy}
 \ordt(B-B_j)\ \ge\ (j+1)\alpha+j ;
\end{equation}
\item \(B\) is the unique element of \(\PLpow\) with \(B=-A\compose(X+B)\), and
\(G\DefinitionEquals X+B\) is the unique element of \(\Gnear\) with
\(F\compose G=X\);
\item \(G\compose F=X\) as well, and \(G\) is also the unique element of
\(\Gnear\) with that property; we write \(G=\rev{F}\);
\item \(\ordt(B)=\alpha\) and \([t^{\alpha}]B=-[t^{\alpha}]A\).
\end{enumerate}
\end{theorem}

\begin{proof}
\emph{(1)}  For \(j=0\), \(B_1-B_0=-A\compose X=-A\) has \(\ordt\ge\alpha
=(0+1)\alpha+0\).  If \eqref{plt:eq:rev-picard-rate} holds for \(j-1\), then by
\cref{plt:lem:rev-gain}
\[
 \ordt(B_{j+1}-B_j)=\ordt\bigl(\Phi(B_j)-\Phi(B_{j-1})\bigr)
 \ge\alpha+\bigl(j\alpha+j-1\bigr)+1=(j+1)\alpha+j .
\]
The bound tends to \(+\infty\), so \((B_j)\) is Cauchy and converges to some
\(B\in\PLpow\) by (I4); moreover
\(\ordt(B-B_j)\ge\min_{i\ge j}\ordt(B_{i+1}-B_i)\ge(j+1)\alpha+j\), which is
\eqref{plt:eq:rev-picard-accuracy}.

That \(B\) solves \(B=\Phi(B)\) needs no separate continuity argument.  For
every \(j\), \(B-\Phi(B)=(B-B_{j+1})+\bigl(\Phi(B_j)-\Phi(B)\bigr)\).  The
first summand has \(\ordt\ge(j+2)\alpha+j+1\) by
\eqref{plt:eq:rev-picard-accuracy}; the second has
\(\ordt\ge\alpha+\ordt(B_j-B)+1\ge(j+2)\alpha+j+1\) by
\cref{plt:lem:rev-gain}.  Hence
\(\ordt\bigl(B-\Phi(B)\bigr)\ge(j+2)\alpha+j+1\) for every \(j\), which
exceeds every fixed integer as \(j\to\infty\); so \(B=\Phi(B)\).

\emph{(2)}  If \(B\) and \(\widetilde B\) are fixed points with
\(r=\ordt(B-\widetilde B)<\infty\), then \cref{plt:lem:rev-gain} gives
\(r=\ordt\bigl(\Phi(B)-\Phi(\widetilde B)\bigr)\ge\alpha+r+1>r\), a
contradiction; so \(B=\widetilde B\).  By
\cref{plt:prop:rev-fixed-point-equation}, \(G=X+B\) is then the unique right
inverse of \(F\) inside \(\Gnear\).

\emph{(3)}  Apply (1)--(2) to \(G\in\Gnear\): there is \(H\in\Gnear\) with
\(G\compose H=X\).  Using (I3),
\[
 F=F\compose X=F\compose(G\compose H)=(F\compose G)\compose H=X\compose H=H,
\]
so \(G\compose F=G\compose H=X\).  If \(G'\in\Gnear\) also satisfies
\(G'\compose F=X\), then
\(G'=G'\compose(F\compose G)=(G'\compose F)\compose G=X\compose G=G\).

\emph{(4)}  From \(B=-A\compose(X+B)\) and \cref{plt:lem:rev-order-preserved},
\(\ordt(B)=\ordt(A)=\alpha\) and \([t^\alpha]B=-[t^\alpha]A\).
\end{proof}

% ed.: this corollary restated \cref{plt:thm:grp-group} in full.  The group
% ed.: statement now stands once, in the composition-group part, and its item
% ed.: (4) cites \cref{plt:thm:rev-existence}; the restatement and its proof
% ed.: are replaced by the sentence below, which records what has just been
% ed.: supplied.
Together with closure and associativity, which are (I3), part (3) of
\cref{plt:thm:rev-existence} supplies the missing clause of
\cref{plt:thm:grp-group}: \((\Gnear,\compose)\) is a group with identity
\(X\), and \(F\mapsto\rev{F}\) is its inversion map.  This is the only place in
the volume where the inverse is constructed.

\begin{corollary}[First-order form of the inverse]
\label{plt:cor:rev-first-order}
For \(F=X+A\in\Gnear\) with \(\alpha=\ordt(A)\),
\(\rev{F}=X-A+\FormalBigOAt{t^{2\alpha+1}}{t}\).
\end{corollary}

\begin{proof}
Take \(j=1\) in \eqref{plt:eq:rev-picard-accuracy}: \(B_1=-A\) and
\(\ordt(B-B_1)\ge2\alpha+1\).
\end{proof}

\begin{remark}[A second proof that a one-sided inverse is two-sided]
\label{plt:rmk:rev-idempotent}
Report~1 argues differently, and the argument is worth keeping because it
localizes the only arithmetic input.  Let \(G\) be the right inverse and put
\(J=G\compose F=X+R\in\Gnear\).  By (I3),
\(J\compose J=G\compose(F\compose G)\compose F=G\compose F=J\), so \(J\) is a
composition idempotent.  On the other hand, by \eqref{plt:eq:rev-taylor},
\[
 J\compose J=(X+R)\compose(X+R)=X+R+R\compose(X+R)
 =X+2R+\sum_{k\ge1}\frac{R^k}{k!}\DX^kR ,
\]
and the last sum has \(\ordt\ge2\ordt(R)+1\) by
\cref{plt:lem:rev-deriv-gain}.  If \(R\ne0\) with \(\rho=\ordt(R)\), comparing
the \(t^\rho\) blocks in \(J\compose J=J\) gives \(2[t^\rho]R=[t^\rho]R\),
hence \([t^\rho]R=0\), a contradiction.  So \(R=0\).  This proof does not
construct a right inverse for \(G\), but it does compare two different integer
multiples of one block; the monoid argument in
\cref{plt:thm:rev-existence}(3) uses nothing beyond associativity.
\end{remark}

\begin{summarybox}
For \(F=X+A\in\Gnear\) the equation \(B=-A\compose(X+B)\) has exactly one
solution, because \(\DX A\) has strictly positive outer valuation: changing the
guess in block \(r\) changes \(-A\compose(X+\cdot)\) only from block
\(r+1+\ordt(A)\) onward.  Iteration from \(B_0=0\) therefore freezes at least
one further block per step, and \(t\)-adic completeness turns that gain into
existence.  Uniqueness, the group property of \(\Gnear\), and two-sidedness of
the inverse all follow.
\end{summarybox}

\section{The reversal residual and the finite certificate}
\label{plt:sec:rev-residual}

\begin{definition}[Reversal residual]
\label{plt:def:rev-residual}
For \(F,H\in\Gnear\) define
\begin{equation}\label{plt:eq:rev-residual-def}
 \RevErr{F}{H}\DefinitionEquals F\compose H-X\ \in\PLpow .
\end{equation}
We call \(\RevErr{F}{H}\) the \emph{right residual} of the candidate inverse
\(H\) for the map \(F\), and \(H\compose F-X\) the \emph{left residual}.
\end{definition}

\begin{warningbox}
The residual has \emph{two} arguments and is not determined by the candidate
\(H\) alone.  Several source reports abbreviate it by a one-argument symbol in
which the ambient map \(F\) is pushed into a subscript or into the surrounding
prose.  In a document that reverses more than one map --- for instance
\(X+\log X\) and its perturbation \(X+\log X+c/X\), both of which occur below
--- such a symbol is ambiguous, and the ambiguity is invisible in the typeset
formula.  Always print both arguments, as in \(\RevErr{F}{H}\).  The same
warning applies to the left residual, which is a different function of the
same pair.
\end{warningbox}

The next proposition is what makes residuals usable as certificates.  It is
stated in none of the six sources, although several of them silently pass from
``the residual has order \(m\)'' to ``\(m\) blocks are correct''.

\begin{proposition}[The residual measures the error exactly, on both sides]
\label{plt:prop:rev-residual-order}
Let \(F\in\Gnear\), \(G=\rev{F}\), \(H\in\Gnear\), and \(E=H-G\in\PLpow\).
Then
\begin{equation}\label{plt:eq:rev-residual-order}
 \ordt\bigl(\RevErr{F}{H}\bigr)
 =\ordt\bigl(H\compose F-X\bigr)
 =\ordt(E),
\end{equation}
and, if \(E\ne0\) with \(e=\ordt(E)\), all three series have the \emph{same}
leading coefficient block:
\[
 [t^{e}]\RevErr{F}{H}=[t^{e}]\bigl(H\compose F-X\bigr)=[t^{e}]E .
\]
\end{proposition}

\begin{proof}
\emph{Right residual.}  Write \(F=X+A\), \(G=X+B\), \(H=X+B+E\), and apply
\cref{plt:lem:rev-shifted-taylor} with \(S=F\), \(P=B\), \(Q=E\):
\[
 F\compose H
 =F\compose G+\bigl((\DX F)\compose G\bigr)E
 +\sum_{j\ge2}\frac{E^j}{j!}\bigl(\DX^jF\bigr)\compose G .
\]
Here \(F\compose G=X\).  Since \(\DX F=1+\DX A\) with \(\ordt(\DX A)\ge1\),
\cref{plt:lem:rev-order-preserved} gives, for
\(J\DefinitionEquals(\DX F)\compose G\), that \(\ordt(J)=0\) and
\([t^0]J=1\).  For \(j\ge2\), \(\DX^jF=\DX^jA\) has \(\ordt\ge j\), so the
\(j\)-th tail term has \(\ordt\ge je+j\ge2e+2\).  Therefore
\begin{equation}\label{plt:eq:rev-residual-expansion}
 \RevErr{F}{H}=J\,E+\FormalBigOAt{t^{2e+2}}{t},
 \qquad \ordt(J)=0,\quad[t^0]J=1 .
\end{equation}
Since \(e\ge0\) we have \(2e+2>e\), while \(\ordt(JE)=e\) and
\([t^{e}](JE)=[t^0]J\cdot[t^{e}]E=[t^{e}]E\).

\emph{Left residual.}  Right composition by \(F\) is substitution, hence a ring
homomorphism, so, using \(G\compose F=X\),
\[
 H\compose F-X=H\compose F-G\compose F=(H-G)\compose F=E\compose F,
\]
and \cref{plt:lem:rev-order-preserved} gives \(\ordt(E\compose F)=e\) with
\([t^{e}](E\compose F)=[t^{e}]E\).
\end{proof}

\begin{corollary}[Finite certificate]
\label{plt:cor:rev-finite-certificate}
Let \(F\in\Gnear\), \(N\ge0\), and
\(G_N=X+\sum_{j=0}^{N}b_j(L)t^j\) with \(b_j\in\K[L]\).  The following are
equivalent.
\begin{enumerate}
\item \(G_N=\Trunc{\rev{F}}{N+1}\), that is, \(G_N\) is correct through outer
block \(N\).
\item \(\RevErr{F}{G_N}=F\compose G_N-X=\FormalBigOAt{t^{N+1}}{t}\)
\emph{(right certificate)}.
\item \(G_N\compose F-X=\FormalBigOAt{t^{N+1}}{t}\) \emph{(left certificate)}.
\end{enumerate}
\end{corollary}

\begin{proof}
By \cref{plt:prop:rev-residual-order}, each of (2) and (3) is equivalent to
\(\ordt(G_N-\rev{F})\ge N+1\).  Since \(G_N\) has no block beyond \(t^N\), that
is exactly \(\Trunc{G_N}{N+1}=\Trunc{\rev{F}}{N+1}\), i.e.\ (1).
\end{proof}

\begin{proposition}[The residual reads off the next inverse block]
\label{plt:prop:rev-next-block}
Let \(F\in\Gnear\), write \(\rev{F}=X+\sum_{n\ge0}b_n(L)t^n\), and let
\(G_{N-1}=\Trunc{\rev{F}}{N}\) for \(N\ge0\).  Then
\begin{equation}\label{plt:eq:rev-next-block}
 \RevErr{F}{G_{N-1}}=-b_N(L)\,t^N+\FormalBigOAt{t^{N+1}}{t}.
\end{equation}
\end{proposition}

\begin{proof}
Here \(E=G_{N-1}-\rev{F}=-\sum_{n\ge N}b_nt^n\), so \(\ordt(E)\ge N\) and
\([t^N]E=-b_N\).  By \eqref{plt:eq:rev-residual-expansion} the residual has
\(\ordt\ge N\) and \([t^N]\RevErr{F}{G_{N-1}}=[t^N]E=-b_N\).
\end{proof}

\begin{remark}[Why both certificates are computed, and what they really test]
\label{plt:rmk:rev-two-certificates}
% ed.: report 1 asserts that at finite order the two residuals "need not have
% identical first omitted blocks"; for an exact G_N in the group that is false.
\cref{plt:prop:rev-residual-order} gives more than the equivalence of the two
certificates: for an exact \(G_N\in\Gnear\) the right and left residuals have
the same valuation \emph{and} the same leading coefficient block.  Report~1's
claim that at finite order the two ``need not have identical first omitted
blocks'' is therefore incorrect as stated.  What is true, and is the genuine
reason to compute both, is this.
\begin{itemize}
\item The two are computed along different code paths: \(F\compose G_N\)
substitutes the candidate into the map, \(G_N\compose F\) substitutes the map
into the candidate.  A transposed argument order, a wrong update of
\(L\compose G=L+\log\bigl(1+t(G-X)\bigr)\), or a routine that substitutes only
the \(t\)-part of the inner series will usually break exactly one of them.
\item As soon as the inputs are themselves truncated --- \(F\) known only to
some block, or \(G_N\) stored with different tails in the two calls --- the two
computed residuals are the residuals of \emph{different} exact elements, and
then their leading blocks may indeed differ.  That is a statement about
inconsistent precision management, not about the algebra of \(\Gnear\), and it
is quantified by \cref{plt:thm:rev-newton-guard} and
\cref{plt:lem:rev-precision-propagation}.
\end{itemize}
\end{remark}

\begin{proposition}[The derivative certificate is one level weaker]
\label{plt:prop:rev-derivative-certificate}
Let \(F,H\in\Gnear\), \(G=\rev{F}\), \(E=H-G\), \(e=\ordt(E)\).  Then
\begin{equation}\label{plt:eq:rev-derivative-certificate}
 \ordt\Bigl(\DX H-\MultiplicativeInverseOf{\bigl((\DX F)\compose H\bigr)}\Bigr)
 \ \ge\ e+1,
\end{equation}
with equality whenever \(e\ge1\).  For \(e=0\) the bound need not be attained:
if \(E=c\in\K^\times\) then \(\DX E=0\) and the left-hand side has valuation at
least \(2\).
\end{proposition}

\begin{proof}
Differentiating \(F\compose G=X\) by the chain rule
\eqref{plt:eq:rev-chain-rule} gives
\(\bigl((\DX F)\compose G\bigr)\DX G=1\), so
\(\DX G=\MultiplicativeInverseOf{\bigl((\DX F)\compose G\bigr)}\).  Hence
\[
 \DX H-\MultiplicativeInverseOf{\bigl((\DX F)\compose H\bigr)}
 =\DX E
 +\Bigl(\MultiplicativeInverseOf{\bigl((\DX F)\compose G\bigr)}
       -\MultiplicativeInverseOf{\bigl((\DX F)\compose H\bigr)}\Bigr).
\]
By \cref{plt:lem:rev-shifted-taylor} applied to \(S=\DX F\),
\((\DX F)\compose H-(\DX F)\compose G
=\sum_{j\ge1}\frac{E^j}{j!}(\DX^{j+1}A)\compose G\), of valuation
\(\ge e+2\) because \(\ordt(\DX^{j+1}A)\ge j+1\ge2\).  Both reciprocals have
valuation \(0\), so the bracket has valuation \(\ge e+2\).  Finally
\(\ordt(\DX E)\ge e+1\), with equality when \(e\ge1\): if
\([t^{e}]E=c\ne0\) then \([t^{e+1}]\DX E=(\partial_L-e)c\), and
\((\partial_L-e)c=0\) forces \(c'=ec\), impossible for a nonzero polynomial
when \(e\ne0\).  For \(e=0\) the operator \(\partial_L\) annihilates constants.
\end{proof}

\begin{remark}
\label{plt:rmk:rev-derivative-certificate}
% ed.: report 1 claims the derivative residual "often detects a wrong block one
% order earlier"; it detects it one level later, and is blind at level 0.
Report~1 recommends the derivative identity
\(\DX G=\MultiplicativeInverseOf{\bigl((\DX F)\compose G\bigr)}\) as a check
that ``often detects a wrong block one order earlier''.
\cref{plt:prop:rev-derivative-certificate} shows the opposite: the derivative
residual first becomes nonzero at least one level \emph{later} than
\(\RevErr{F}{H}\), and it under-reports by two levels an error whose leading
block is a nonzero constant at level \(0\).  It remains a useful
\emph{independent} check --- it exercises \(\DX\) and the reciprocal rather
than the composition engine --- but as a detector of the first wrong block it
is strictly weaker.
\end{remark}

\section{What the theorem does not say}
\label{plt:sec:rev-not-analytic}

\begin{warningbox}
\cref{plt:thm:rev-existence} is a theorem about formal series.  It asserts
nothing about the existence of an analytic inverse branch, about convergence of
\(\sum_nb_n(L)t^n\) for any value of \(X\), or about any constant uniform in
\(L\).  A formal certificate \(\RevErr{F}{G_N}=\FormalBigOAt{t^{N+1}}{t}\)
means coefficient agreement below outer order \(N+1\) and nothing more; in
particular it does not imply
\(f(g_N(X))-X=\BigOAt{X^{-N-1}}{X\to+\infty}\) for any realization
\(f,g_N\), since the discarded blocks carry unbounded powers of \(\log X\).
\end{warningbox}

\begin{proposition}[Analytic residual transfer]
\label{plt:prop:rev-analytic-transfer}
Let \(I\subseteq\RealNumbers\) be an interval, let \(f\colon I\to\RealNumbers\)
be continuously differentiable and injective, let \(y\in f(I)\), and let
\(x_*=\rev{f}(y)\).  Let \(x_N\in I\) and suppose
\(\AbsoluteValue{f'(x)}\ge m>0\) for all \(x\) between \(x_N\) and \(x_*\).
Then
\[
 \AbsoluteValue{x_N-x_*}\le\frac{\AbsoluteValue{f(x_N)-y}}{m}.
\]
\end{proposition}

\begin{proof}
By the mean value theorem there is \(\xi\) between \(x_N\) and \(x_*\) with
\(f(x_N)-y=f(x_N)-f(x_*)=f'(\xi)(x_N-x_*)\), whence
\(\AbsoluteValue{f(x_N)-y}\ge m\AbsoluteValue{x_N-x_*}\).
\end{proof}

\begin{remark}
\label{plt:rmk:rev-transfer-hypotheses}
Both hypotheses are essential and neither is supplied by the formal theory.
For \(f(x)=x+\log x\) on \((1,\infty)\) one may take \(m=1\), because
\(f'(x)=1+1/x>1\), so there a numerical residual does control the inverse
error --- but that is a statement about one realization on one ray, proved
analytically, and it is not a corollary of
\cref{plt:cor:rev-finite-certificate}.  Transferring a \emph{formal} residual
to an asymptotic estimate additionally requires a realization theorem for the
class, treated separately in this volume.
\end{remark}

\chapter{Triangular reversal, Newton iteration, and precision doubling}
\label{plt:sec:reversion-algorithms}

\section{The triangular correction}
\label{plt:sec:rev-triangular}

\begin{lemma}[Block perturbation]
\label{plt:lem:rev-block-perturbation}
Let \(F\in\Gnear\), \(G\in\Gnear\), \(N\in\NaturalNumbers\), \(c\in\K[L]\).
Then
\begin{equation}\label{plt:eq:rev-block-perturbation}
 F\compose\bigl(G+c(L)t^N\bigr)
 =F\compose G+c(L)t^N+\FormalBigOAt{t^{N+1}}{t}.
\end{equation}
\end{lemma}

\begin{proof}
Apply \cref{plt:lem:rev-shifted-taylor} with \(S=F\), \(P=G-X\),
\(Q=\delta\DefinitionEquals ct^N\):
\[
 F\compose(G+\delta)-F\compose G
 =\bigl((\DX F)\compose G\bigr)\delta
 +\sum_{j\ge2}\frac{\delta^j}{j!}\bigl(\DX^jF\bigr)\compose G .
\]
By \cref{plt:lem:rev-order-preserved}, \((\DX F)\compose G=1+\FormalBigOAt{t}{t}\),
so the linear term is \(ct^N+\FormalBigOAt{t^{N+1}}{t}\); and for \(j\ge2\) the
\(j\)-th term has \(\ordt\ge jN+j\ge2N+2>N\) since \(N\ge0\).
\end{proof}

\begin{theorem}[Triangular reversal]
\label{plt:thm:rev-triangular}
Let \(F\in\Gnear\).  Put \(G_{-1}=X\) and, for \(n=0,1,2,\dots\),
\begin{equation}\label{plt:eq:rev-triangular-step}
 r_n(L)\DefinitionEquals[t^n]\RevErr{F}{G_{n-1}},
 \qquad
 G_n\DefinitionEquals G_{n-1}-r_n(L)t^n .
\end{equation}
Then for every \(n\ge-1\),
\[
 \RevErr{F}{G_n}=\FormalBigOAt{t^{n+1}}{t},
 \qquad
 G_n=\Trunc{\rev{F}}{n+1},
\]
and \(r_n=-b_n\), where \(\rev{F}=X+\sum_{j\ge0}b_j(L)t^j\).
\end{theorem}

\begin{proof}
Induction on \(n\).  For \(n=-1\), \(\RevErr{F}{X}=F-X=A\) has \(\ordt\ge0\),
which is the assertion \(\FormalBigOAt{t^{0}}{t}\), and
\(\Trunc{\rev{F}}{0}=X\).  Assume
\(\RevErr{F}{G_{n-1}}=\FormalBigOAt{t^{n}}{t}\).  Then \(r_n\) is its
coefficient at level \(n\), and by \cref{plt:lem:rev-block-perturbation} with
\(c=-r_n\),
\[
 \RevErr{F}{G_n}=\RevErr{F}{G_{n-1}}-r_nt^n+\FormalBigOAt{t^{n+1}}{t}
 =\FormalBigOAt{t^{n+1}}{t},
\]
since the inherited residual and the correction contribute \(r_nt^n\) and
\(-r_nt^n\) at level \(n\) and nothing below.  As \(G_n\) has no block beyond
\(t^n\), \cref{plt:cor:rev-finite-certificate} identifies
\(G_n=\Trunc{\rev{F}}{n+1}\).  Finally \cref{plt:prop:rev-next-block}, applied
to \(G_{n-1}=\Trunc{\rev{F}}{n}\), gives \(r_n=-b_n\).
\end{proof}

\begin{algorithm}
\label{plt:alg:rev-triangular}
\textbf{Residual-cancellation reversal.}
\emph{Input:} \(F=X+A\in\Gnear\) with \(A\) known below \(t^{N+1}\); a target
block \(N\ge0\).  \emph{Output:} \(\Trunc{\rev{F}}{N+1}\).
\begin{enumerate}
\item \(G\leftarrow X\).
\item For \(n=0,1,\dots,N\): compute
 \(R\leftarrow\Trunc{\bigl(F\compose G-X\bigr)}{N+1}\); set
 \(r_n\leftarrow[t^n]R\); set \(G\leftarrow G-r_n(L)t^n\).
\item Return \(G\) together with the certificate
 \(\Trunc{\bigl(F\compose G-X\bigr)}{N+1}=0\).
\end{enumerate}
Every intermediate composition may be truncated below \(t^{N+1}\): by
\cref{plt:lem:rev-gain} a change of \(G\) at order \(\ge N+1\) changes
\(F\compose G\) only at order \(\ge N+1\), so a discarded block never returns.
\end{algorithm}

Step \eqref{plt:eq:rev-triangular-step} divides the residual block by the
leading block of \(\DX F\), which is \(1\) for \(F\in\Gnear\).  Outside
\(\Gnear\) this is exactly where the unit criterion of the coefficient ring
enters, and the sources treat it too briefly.

\begin{proposition}[Reversal and triangular correction at a general leading
slope]
\label{plt:prop:rev-general-slope}
Let \(a\in\K^\times\) and fix a value \(\log a\in\K\).  Let \(\theta\) be the
\(\K\)-algebra automorphism of \(\PLlau\) induced by the substitution
\(X\mapsto a^{-1}X\), that is \(\theta(t)=at\) and \(\theta(L)=L-\log a\);
equivalently \(\theta(S)=S\compose(a^{-1}X)\).  Let \(A\in\PLpow\),
\(F=aX+A\), \(G_0\DefinitionEquals X+a^{-1}A\in\Gnear\), and
\(\mathcal G_a\DefinitionEquals a^{-1}X+\PLpow=\theta(\Gnear)\).  Then:
\begin{enumerate}
\item \(\theta\) preserves \(\ordt\), maps \(\PLpow\) onto \(\PLpow\), and
satisfies \(\theta(S\compose H)=S\compose\theta(H)\) whenever the composites
exist;
\item \(F=a\,G_0\) (a product, not a composite), and \(F\) has exactly one
compositional inverse in \(\mathcal G_a\), namely
\(\rev{F}=\theta\bigl(\rev{G_0}\bigr)\), which is two-sided;
\item for every \(G\in\mathcal G_a\), \(c\in\K[L]\) and \(N\ge0\),
\begin{equation}\label{plt:eq:rev-general-slope}
 F\compose\bigl(G+c(L)t^N\bigr)=F\compose G+a\,c(L)\,t^N
 +\FormalBigOAt{t^{N+1}}{t},
\end{equation}
so a residual block \(r_N\) is cancelled by the correction \(-r_N/a\), which
stays in \(\K[L]\).
\end{enumerate}
\end{proposition}

\begin{proof}
\emph{(1)}  \(\theta(t^np(L))=a^nt^np(L-\log a)\), so \(\theta\) is a bijective
ring homomorphism preserving \(\ordt\) and \(\PLpow\); it is well defined
because \(\log a\in\K\).  The intertwining relation is (I3):
\(\theta(S\compose H)=(S\compose H)\compose(a^{-1}X)
=S\compose\bigl(H\compose(a^{-1}X)\bigr)=S\compose\theta(H)\).

\emph{(2)}  \(F=aX+A=a(X+a^{-1}A)=a\,G_0\).  Put \(H=\rev{G_0}\).  Then, by
(1), \(F\compose\theta(H)=\theta(F\compose H)=\theta\bigl(a(G_0\compose
H)\bigr)=\theta(aX)=a\cdot a^{-1}X=X\).  On the other side,
\(\theta(H)\compose F=H\compose(a^{-1}X)\compose F=H\compose(a^{-1}F)
=H\compose G_0=X\).  For uniqueness, let \(G'\in\mathcal G_a\) with
\(F\compose G'=X\) and write \(G'=\theta(\widetilde G)\) with
\(\widetilde G\in\Gnear\).  Then \(\theta(F\compose\widetilde G)=X\), so
\(F\compose\widetilde G=\theta^{-1}(X)=aX\), i.e.\
\(G_0\compose\widetilde G=X\), so \(\widetilde G=\rev{G_0}\) by
\cref{plt:thm:rev-existence}(2).

\emph{(3)}  Write \(G=\theta(\widetilde G)\) with \(\widetilde G\in\Gnear\) and
\(\delta=ct^N\), \(\eta=\theta^{-1}(\delta)=a^{-N}c(L+\log a)t^N\in\PLpow\).
Since \(\theta\) is additive, \(G+\delta=\theta(\widetilde G+\eta)\), so by (1)
and \(F=aG_0\),
\[
 F\compose(G+\delta)
 =\theta\bigl(F\compose(\widetilde G+\eta)\bigr)
 =\theta\bigl(a\,G_0\compose(\widetilde G+\eta)\bigr).
\]
By \cref{plt:lem:rev-block-perturbation},
\(G_0\compose(\widetilde G+\eta)=G_0\compose\widetilde G+\eta
+\FormalBigOAt{t^{N+1}}{t}\).  Applying \(a\theta(\cdot)\) and using
\(\theta(\eta)=\delta\) together with the fact that \(\theta\) preserves
\(\ordt\) gives \eqref{plt:eq:rev-general-slope}.
\end{proof}

\begin{remark}[Two repairs to the general-slope rule]
\label{plt:rmk:rev-slope-repairs}
% ed.: report 4 states the rule as "for a leading slope F = ax + o(x) the
% correction is divided by a", with no hypothesis on a and no discussion of a
% nonconstant leading block of D_X F.
\emph{(i) The scaling reduction needs \(\log a\in\K\).}  Report~3 performs the
reduction \(F=a\,G_0\), \(\rev{F}(X)=\rev{G_0}(X/a)\) and notes in passing that
the substitution ``shifts the logarithm by \(-\log c\)''; report~4 states the
division rule with no hypothesis at all.  As
\cref{plt:prop:rev-general-slope}(1) makes explicit, the substitution
\(X\mapsto a^{-1}X\) is an automorphism of \(\PLlau\) only if a value of
\(\log a\) lies in \(\K\); otherwise \(\K\) must first be extended, and the
inverse of a map with coefficients in \(\K\) acquires coefficients in
\(\K(\log a)\).  Composition with \(F\) itself is likewise undefined without
that value.

\emph{(ii) The divisor must be a unit of the coefficient ring, not merely
nonzero.}  What is really being inverted in
\eqref{plt:eq:rev-general-slope} is \(\lambda\DefinitionEquals[t^0](\DX F)\),
the leading block of the derivative.  For \(F=aX+A\) it equals \(a\in\K^\times\)
and all is well.  But \(\lambda\) can be a nonconstant polynomial: for
\(F=X\log X=t^{-1}L\) one has
\(\DX F=(\partial_L+1)L=L+1\), so \(\lambda=L+1\), which is not a unit of
\(\K[L]\).  Division by \(\lambda\) then leaves \(\PLpow\) for \(\PLrat\).  This
is not a removable inconvenience: \(X\log X\) has no compositional inverse in
\(\PLlau\) at all.  Indeed, for any \(G\in\Gnear\) the composite
\((X+B)\log(X+B)\) has \(\ordt=-1\) with \(t^{-1}\)-block \(L\ne1\), so it
cannot equal \(X\); and \(\log(X\log X)=L+\log L\) forces the second logarithm
into the scale.  A leading slope that is not a unit is a reliable symptom that
the inverse leaves the one-logarithm class.
\end{remark}

\section{Coefficient stabilization by functional iteration}
\label{plt:sec:rev-picard}

\begin{algorithm}
\label{plt:alg:rev-picard}
\textbf{Coefficient-stabilization reversal.}
\emph{Input:} \(F=X+A\in\Gnear\), target block \(N\ge0\).
\emph{Output:} \(\Trunc{\rev{F}}{N+1}\).
\begin{enumerate}
\item \(B\leftarrow0\).
\item Repeat \(N+1\) times:
 \(B\leftarrow\Trunc{\bigl(-A\compose(X+B)\bigr)}{N+1}\).
\item Return \(X+B\) with the certificate
 \(\Trunc{\bigl(F\compose(X+B)-X\bigr)}{N+1}=0\).
\end{enumerate}
One may stop as soon as two consecutive iterates agree below \(t^{N+1}\).  By
\cref{plt:thm:rev-existence}(1) that happens after at most \(N+1\) steps, and
after at most \(\Ceiling{(N+1-\alpha)/(\alpha+1)}\) steps when
\(\alpha=\ordt(A)\).
\end{algorithm}

\begin{proposition}[Truncation is safe]
\label{plt:prop:rev-truncation-safe}
Let \(A\in\PLpow\), \(M\ge0\), and \(B,\widetilde B\in\PLpow\) with
\(\ordt(B-\widetilde B)\ge M\).  Then
\(\Trunc{\Phi(B)}{M}=\Trunc{\Phi(\widetilde B)}{M}\).  Consequently, truncating
the iterate below \(t^{M}\) after every step of \cref{plt:alg:rev-picard}
alters no block below \(t^{M}\) of any later iterate, and the algorithm returns
the same answer as its untruncated form.
\end{proposition}

\begin{proof}
\cref{plt:lem:rev-gain} gives
\(\ordt\bigl(\Phi(B)-\Phi(\widetilde B)\bigr)\ge\ordt(A)+M+1>M\), which is the
first claim.  For the second, let \(B_j\) be the exact iterates and
\(\widehat B_j\) the truncated ones, and suppose
\(\ordt(B_j-\widehat B_j)\ge M\).  Then
\(\ordt\bigl(\Phi(B_j)-\Phi(\widehat B_j)\bigr)\ge M+1\), and truncating
\(\Phi(\widehat B_j)\) below \(t^M\) introduces an error of order \(\ge M\);
hence \(\ordt(B_{j+1}-\widehat B_{j+1})\ge M\), and the invariant persists.
\end{proof}

\begin{remark}[An algebraically equivalent iteration need not be
order-improving]
\label{plt:rmk:rev-bad-rearrangement}
An existence proof by iteration depends on the chosen normal form, not only on
the equation.  For \(F=X+\log X\) the inverse \(\WrightOmega=\rev{F}\)
satisfies \(\WrightOmega+\log\WrightOmega=X\).  The rearrangement
\begin{equation}\label{plt:eq:rev-good-rearrangement}
 \WrightOmega=X-\log\WrightOmega
\end{equation}
is order-improving: if \(\omega_1,\omega_2\in\Gnear\) and
\(\omega_1-\omega_2=\FormalBigOAt{t^{r}}{t}\), then
\(\omega_1\MultiplicativeInverseOf{\omega_2}=1+\FormalBigOAt{t^{r+1}}{t}\),
hence \(\log\omega_1-\log\omega_2=\FormalBigOAt{t^{r+1}}{t}\); this is
\cref{plt:lem:rev-gain} again.  The algebraically equivalent rearrangement
\begin{equation}\label{plt:eq:rev-bad-rearrangement}
 \WrightOmega=\EulerE^{\,X-\WrightOmega}
\end{equation}
fails for a sharper reason than slow convergence: \(\EulerE^{X}\) is not an
element of \(\PLlau\), so the right-hand side of
\eqref{plt:eq:rev-bad-rearrangement} does not define a map of the class and no
truncation of it is meaningful.  Report~6 makes the same point with an
exponential rearrangement of the Lambert equation, attributing the observation
to Edgar \cite{edgar-beginners}: a fixed-point scheme must increase asymptotic
precision, and algebraic equivalence is no evidence that it does.  The test is
supplied by \cref{plt:lem:rev-gain}: the iteration map must send a difference of
valuation \(r\) to one of valuation \(>r\), uniformly in \(r\).
\end{remark}

\section{Newton iteration}
\label{plt:sec:rev-newton}

\begin{lemma}[The Newton denominator is a unit]
\label{plt:lem:rev-newton-unit}
Let \(F=X+A\in\Gnear\) and \(G\in\Gnear\).  Then
\(J\DefinitionEquals(\DX F)\compose G\) satisfies \(J=1+\FormalBigOAt{t}{t}\)
and is a unit of \(\PLpow\), with
\(\MultiplicativeInverseOf{J}=\sum_{k\ge0}(1-J)^k\).
\end{lemma}

\begin{proof}
\(\DX F=1+\DX A\) and \(\ordt(\DX A)\ge\ordt(A)+1\ge1\) by
\cref{plt:lem:rev-deriv-gain}.  By \cref{plt:lem:rev-order-preserved},
\(\ordt(J)=0\) and \([t^0]J=1\), so \(W\DefinitionEquals1-J\) has
\(\ordt(W)\ge1\); the family \((W^k)_{k\ge0}\) is summable and
\(J\sum_{k\ge0}W^k=(1-W)\sum_{k\ge0}W^k=1\).  (This is the unit criterion in
\(\PLpow\): \(U\) is a unit exactly when \([t^0]U\in\K^\times\).)
\end{proof}

\begin{definition}[Newton reversion map]
\label{plt:def:rev-newton}
For \(F=X+A\in\Gnear\) and \(G\in\Gnear\) put
\begin{equation}\label{plt:eq:rev-newton}
 \mathcal N_F(G)\DefinitionEquals
 G-\frac{\RevErr{F}{G}}{(\DX F)\compose G}\ \in\Gnear .
\end{equation}
\end{definition}

\begin{theorem}[Precision doubling, exact arithmetic]
\label{plt:thm:rev-newton-doubling}
Let \(F=X+A\in\Gnear\), \(\alpha=\ordt(A)\), \(G\in\Gnear\), and
\(m=\ordt\bigl(\RevErr{F}{G}\bigr)\).  Then
\begin{equation}\label{plt:eq:rev-newton-doubling}
 \ordt\Bigl(\RevErr{F}{\mathcal N_F(G)}\Bigr)\ \ge\ 2m+2+\alpha .
\end{equation}
Equivalently, by \cref{plt:prop:rev-residual-order}: if \(G\) is correct through
outer block \(m-1\), then \(\mathcal N_F(G)\) is correct through outer block
\(2m+1+\alpha\); the number of correct blocks passes from \(m\) to
\(2m+2+\alpha\).
\end{theorem}

\begin{proof}
Put \(R=\RevErr{F}{G}\), \(J=(\DX F)\compose G\), and
\(C=-R\MultiplicativeInverseOf{J}\), so \(\mathcal N_F(G)=G+C\).  By
\cref{plt:lem:rev-newton-unit}, \(\ordt(C)=\ordt(R)=m\).  By
\cref{plt:lem:rev-shifted-taylor} with \(S=F\), \(P=G-X\), \(Q=C\),
\[
 \RevErr{F}{G+C}
 =\underbrace{R+JC}_{=0}
 +\sum_{j\ge2}\frac{C^j}{j!}\bigl(\DX^jF\bigr)\compose G .
\]
For \(j\ge2\), \(\DX^jF=\DX^jA\), and
\cref{plt:lem:rev-deriv-gain,plt:lem:rev-order-preserved} give
\(\ordt\bigl((\DX^jA)\compose G\bigr)\ge\alpha+j\).  Hence the \(j\)-th term
has \(\ordt\ge jm+\alpha+j=\alpha+j(m+1)\), minimal at \(j=2\) with value
\(2m+2+\alpha\).
\end{proof}

\begin{remark}[Diagnosis of the six sources]
\label{plt:rmk:rev-doubling-diagnosis}
% ed.: the sources give five mutually inconsistent exponents; one contradicts
% its own proof, and two slide from residual order to correct-block count.
The reports state the Newton estimate in five different ways.

\begin{itemize}
\item Reports~1, 5 and~6 assert the exponent \(2m+2\), which is
\eqref{plt:eq:rev-newton-doubling} with the term \(\alpha\) omitted; their
proofs are correct as far as they go.
\item Report~2 states \(2N\) under the extra hypothesis \(N\ge1\).  The
hypothesis is never used and is not needed; the case \(N=0\) is legitimate and
is precisely the situation at the first Newton step from \(G=X\), where the
residual is \(A\) itself.  The argument given there in fact yields \(2N+2\).
\item Report~3 states \(2m\) and adds only that ``in this particular scale the
order is often even higher''.  The improvement is not occasional but automatic
and equals \(2m+2+\ordt(A)\).
\item Report~4 is internally inconsistent: it announces \(t^{2N+1}\) in the
text and proves \(t^{2N+2}\) two lines later, bridging the gap with the phrase
``allowing coarse bounds and Laurent shifts''.  There are no Laurent shifts
here: every object in \eqref{plt:eq:rev-newton} lies in \(\PLpow\), and the
estimate is exact.
\end{itemize}

A second and subtler slide occurs in reports~1 and~5, which pass from ``the
residual has order \(2m+2\)'' to ``the number of correct blocks doubles''
without argument.  That step is legitimate, but only because of
\cref{plt:prop:rev-residual-order}, which identifies the residual valuation
with \(\ordt(G-\rev{F})\); no source proves or even states that identification.
Without it, a small residual is a priori compatible with a large coefficient
error --- exactly as it would be for a map whose derivative were not a unit.
\end{remark}

\begin{remark}[Sharpness]
\label{plt:rmk:rev-doubling-sharp}
The exponent in \eqref{plt:eq:rev-newton-doubling} cannot be improved in
general.  For \(F=X+\log X\) (so \(\alpha=0\)) and
\(G=X-L\MultiplicativeInverseOf{(1+t)}\) one has \(m=2\), and after one exact
Newton step the residual is \(-\tfrac18L^4t^6+\FormalBigOAt{t^{7}}{t}\), of
order precisely \(2m+2=6\).  See \cref{plt:ex:rev-guard}.
\end{remark}

\section{Precision doubling with truncated intermediates}
\label{plt:sec:rev-guard}

In an implementation none of \(R\), \(J\) or the quotient is known exactly, and
the exponent \(2m+2+\alpha\) is attained only if the truncations are deep
enough.  Reports~1, 2, 4, 5 and~6 omit the hypothesis; report~3 gestures at it
(``computed consistently to sufficient precision'').  The following two results
make it exact.

\begin{theorem}[Newton step with an inexact correction]
\label{plt:thm:rev-newton-guard}
Let \(F=X+A\in\Gnear\), \(\alpha=\ordt(A)\), \(G\in\Gnear\),
\(R=\RevErr{F}{G}\), \(m=\ordt(R)\), \(J=(\DX F)\compose G\), and let
\(C=-R\MultiplicativeInverseOf{J}\) be the exact Newton correction.  Let
\(\widehat C\in\PLpow\) satisfy \(\ordt(\widehat C-C)\ge\sigma\) with
\(\sigma\ge m\).  Then
\begin{equation}\label{plt:eq:rev-newton-guard}
 \ordt\Bigl(\RevErr{F}{G+\widehat C}\Bigr)
 \ \ge\ \min\bigl\{\,2m+2+\alpha,\ \sigma\,\bigr\}.
\end{equation}
\end{theorem}

\begin{proof}
Since \(\ordt(C)=m\) and \(\ordt(\widehat C-C)\ge\sigma\ge m\), we get
\(\ordt(\widehat C)\ge m\).  By \cref{plt:lem:rev-shifted-taylor} with
\(S=F\), \(P=G-X\), \(Q=\widehat C\),
\[
 \RevErr{F}{G+\widehat C}
 =R+J\widehat C
 +\sum_{j\ge2}\frac{\widehat C^{\,j}}{j!}\bigl(\DX^jA\bigr)\compose G .
\]
The tail has \(\ordt\ge2m+2+\alpha\) exactly as in
\cref{plt:thm:rev-newton-doubling}.  For the head, \(R+JC=0\) gives
\(R+J\widehat C=J(\widehat C-C)\), and \(\ordt(J)=0\), so
\(\ordt\bigl(R+J\widehat C\bigr)=\ordt(\widehat C-C)\ge\sigma\).
\end{proof}

\begin{lemma}[Precision propagation through the Newton step]
\label{plt:lem:rev-precision-propagation}
Keep the notation of \cref{plt:thm:rev-newton-guard} and let
\(\widehat R,\widehat K\in\PLpow\) and \(q\in\NaturalNumbers\) satisfy
\[
 \ordt\bigl(\widehat R-R\bigr)\ge\sigma_R\ge m,
 \qquad
 \ordt\bigl(\widehat K-\MultiplicativeInverseOf{J}\bigr)\ge\sigma_K\ge0 .
\]
Put \(\widehat C=-\Trunc{\bigl(\widehat R\,\widehat K\bigr)}{q}\).  Then
\(\ordt(\widehat C-C)\ge\min\{q,\ \sigma_R,\ m+\sigma_K\}\).  Consequently, to
certify \(\ordt\bigl(\RevErr{F}{G+\widehat C}\bigr)\ge q\) for a target
\(q\le2m+2+\alpha\), it suffices to take \(\sigma_R\ge q\) and
\(\sigma_K\ge q-m\).  Moreover:
\begin{enumerate}
\item \(\widehat K\) with \(\sigma_K=s\) is obtained from \(\Trunc{J}{s}\),
since for units \(J,\widehat J\) with \([t^0]J=[t^0]\widehat J=1\) one has
\(\MultiplicativeInverseOf{\widehat J}-\MultiplicativeInverseOf{J}
=(J-\widehat J)\MultiplicativeInverseOf{\widehat J}\MultiplicativeInverseOf{J}\);
\item \(\Trunc{J}{s}\) is determined by \(\Trunc{A}{s-1}\) and
\(\Trunc{G}{s-2}\), because changing \(G\) by \(\FormalBigOAt{t^{u}}{t}\)
changes \(J\) only by \(\FormalBigOAt{t^{u+2}}{t}\);
\item \(\Trunc{R}{q}\) is determined by \(\Trunc{A}{q}\) and
\(\Trunc{G}{q}\), and no coarser knowledge of \(G\) suffices: for
\(\epsilon\in\PLpow\),
\[
 \RevErr{F}{G+\epsilon}-\RevErr{F}{G}
 =J\epsilon+\FormalBigOAt{t^{2\ordt(\epsilon)+2}}{t},
\]
whose valuation is exactly \(\ordt(\epsilon)\).  Thus a block of \(G\) at
level \(u<q\) moves \(R\) at level \(u\), and if \(G\) is stored with fewer
than \(q\) blocks its zero padding is part of the definition of the element
being corrected and must be used consistently throughout the step.
\end{enumerate}
\end{lemma}

\begin{proof}
Write
\(\widehat C-C=-\bigl(\Trunc{(\widehat R\widehat K)}{q}-\widehat R\widehat K\bigr)
-\bigl(\widehat R\widehat K-R\MultiplicativeInverseOf{J}\bigr)\).  The first
bracket has \(\ordt\ge q\).  For the second,
\[
 \widehat R\widehat K-R\MultiplicativeInverseOf{J}
 =\bigl(\widehat R-R\bigr)\widehat K
 +R\bigl(\widehat K-\MultiplicativeInverseOf{J}\bigr),
\]
and \(\ordt(\widehat K)\ge\min\{0,\sigma_K\}=0\) because
\(\ordt(\MultiplicativeInverseOf{J})=0\); so its valuation is at least
\(\min\{\sigma_R,\ m+\sigma_K\}\).  With \(\sigma_R\ge q\), \(\sigma_K\ge q-m\)
the bound becomes \(q\), and \cref{plt:thm:rev-newton-guard} applies since
\(q\ge m\).

For (1), the displayed identity gives
\(\ordt\bigl(\MultiplicativeInverseOf{\widehat J}
-\MultiplicativeInverseOf{J}\bigr)=\ordt(J-\widehat J)\ge s\).  For (2),
\cref{plt:lem:rev-shifted-taylor} applied to \(S=\DX F\) gives
\((\DX F)\compose(G+\epsilon)-(\DX F)\compose G
=\sum_{j\ge1}\frac{\epsilon^j}{j!}(\DX^{j+1}A)\compose G\), whose valuation is
at least \(\ordt(\epsilon)+2\) because \(\ordt(\DX^{j+1}A)\ge j+1\ge2\).  For
(3), \cref{plt:lem:rev-shifted-taylor} with \(S=F\), \(P=G-X\),
\(Q=\epsilon\) gives \(\RevErr{F}{G+\epsilon}-\RevErr{F}{G}
=J\epsilon+\sum_{j\ge2}\frac{\epsilon^j}{j!}(\DX^jA)\compose G\), and the tail
has valuation at least \(2\ordt(\epsilon)+2>\ordt(\epsilon)\) while
\(\ordt(J\epsilon)=\ordt(\epsilon)\).
\end{proof}

\begin{corollary}[Doubling schedule]
\label{plt:cor:rev-newton-schedule}
Let \(F=X+A\in\Gnear\) and let \(N\ge0\) be a target block.  Put \(H_0=X\),
\(q_0=\ordt(A)\), and for \(i\ge0\)
\[
 q_{i+1}=\min\{N+1,\ 2q_i+2\},
 \qquad
 H_{i+1}=\Trunc{\Bigl(H_i
   -\Trunc{R_i}{q_{i+1}}\cdot
    \Trunc{\MultiplicativeInverseOf{J_i}}{\,q_{i+1}-q_i}\Bigr)}{q_{i+1}},
\]
where \(R_i=\RevErr{F}{H_i}\) and \(J_i=(\DX F)\compose H_i\).
Then \(\ordt\bigl(\RevErr{F}{H_i}\bigr)\ge q_i\) for every \(i\), and
\(H_i=\Trunc{\rev{F}}{N+1}\) as soon as \(q_i=N+1\).  The number of steps
is \(\BigO(\log N)\).
\end{corollary}

\begin{proof}
Induction on \(i\).  The case \(i=0\) is \(\RevErr{F}{X}=A\).  Assume
\(\ordt(R_i)\ge q_i\).  Apply
\cref{plt:lem:rev-precision-propagation} with \(m=q_i\), \(q=q_{i+1}\),
\(\sigma_R=q_{i+1}\) and \(\sigma_K=q_{i+1}-q_i\), which is \(\ge1\) because
\(q_{i+1}\ge\min\{N+1,2q_i+2\}>q_i\) while \(q_i<N+1\); the final outer
truncation \(\Trunc{\cdot}{q_{i+1}}\) contributes an error of valuation
\(\ge q_{i+1}\) and is absorbed in the same bound.  Then
\cref{plt:thm:rev-newton-guard} gives
\(\ordt(R_{i+1})\ge\min\{2q_i+2+\ordt(A),\ q_{i+1}\}=q_{i+1}\).  The
identification with \(\Trunc{\rev{F}}{N+1}\) is
\cref{plt:cor:rev-finite-certificate}.  Finally \(q_{i+1}+1\ge2(q_i+1)\) until
the cap is reached, so \(q_i+1\ge2^i\).
\end{proof}

\begin{warningbox}
A Newton step at target \(q\) treats the stored \(G\) as an exact element of
\(\Gnear\) --- zero-padded beyond its last stored block --- and must form
\(R=\RevErr{F}{G}\) below \(t^{q}\) from \emph{that same} element, even though
only \(m<q\) of its blocks are known to be correct.  The blocks between levels
\(m\) and \(q-1\) are the guard blocks: they are not required to be correct,
but they are required to be fixed and used identically in the residual, in the
Jacobian factor and in the update.  Computing \(J\) to the final cutoff while
computing \(R\) only to the current accuracy \(m\), or recomputing \(R\) from a
differently truncated copy of \(G\), destroys the doubling: by
\eqref{plt:eq:rev-newton-guard} the guaranteed order collapses to \(\sigma\),
and the collapse is invisible in the printed series.
\end{warningbox}

\begin{example}[Doubling, and its failure, for \(F=X+\log X\)]
\label{plt:ex:rev-guard}
Let \(F=X+\log X=X+L\), so \(A=L\), \(\alpha=0\), and write
\(\WrightOmega=\rev{F}\).  Here \(\DX F=1+t\), and for \(G=X+B\)
\[
 \RevErr{F}{G}=B+L+\log(1+Bt),
 \qquad
 J=(\DX F)\compose G=1+\MultiplicativeInverseOf{G}
  =1+\frac{t}{1+Bt}.
\]

\emph{First step.}  From \(H_0=X\): \(R_0=L\), \(m=0\),
\(J_0=1+t\), so
\[
 H_1=X-\frac{L}{1+t}=X-L+Lt-Lt^2+Lt^3-\cdots .
\]
\cref{plt:thm:rev-newton-doubling} predicts \(\ordt(R_1)\ge2\), and indeed
\[
 R_1=-\tfrac12L^2t^2
 +\bigl(L^2-\tfrac13L^3\bigr)t^3
 +\bigl(-\tfrac32L^2+L^3-\tfrac14L^4\bigr)t^4
 +\FormalBigOAt{t^{5}}{t},
\]
so \(m=2\) exactly.  By \cref{plt:prop:rev-residual-order} the error must have
the same leading block, and
\(H_1-\WrightOmega=-\tfrac12L^2t^2+\FormalBigOAt{t^{3}}{t}\).

\emph{Exact second step.}  With \(q=2m+2=6\) and all data carried below
\(t^{6}\),
\[
 \RevErr{F}{\mathcal N_F(H_1)}=-\tfrac18L^4t^6+\FormalBigOAt{t^{7}}{t},
\]
of order exactly \(6\): blocks \(0,\dots,5\) of \(\WrightOmega\) are correct,
six blocks in place of two.

\emph{Guard-starved second step.}  Suppose the residual is truncated one
doubling too early, \(\widehat R=\Trunc{R_1}{4}\), while \(J\) and the
division are still carried below \(t^{6}\).  Then \(\sigma_R=4\) and
\cref{plt:lem:rev-precision-propagation} guarantees only order \(4\) --- and
the bound is attained.  With
\(\rho_4\DefinitionEquals[t^4]R_1=-\tfrac14L^4+L^3-\tfrac32L^2\),
\[
 \widehat C-C=\bigl(R_1-\widehat R\bigr)\MultiplicativeInverseOf{J}
 =\rho_4t^4+\FormalBigOAt{t^{5}}{t},
\]
hence \(\RevErr{F}{H_1+\widehat C}=\rho_4t^4+\FormalBigOAt{t^{5}}{t}\) and,
by \cref{plt:prop:rev-residual-order}, the computed inverse first deviates from
\(\WrightOmega\) in block \(4\), by exactly \(\rho_4\).  Four correct blocks in
place of six: the loss equals the two blocks of residual that were discarded,
and no amount of extra precision in \(J\) or in the division repairs it.
\end{example}

\section{Four reversal strategies compared}
\label{plt:sec:rev-comparison}

Write \(C(n)\) for the cost of one composition \(F\compose G\) truncated below
\(t^{n}\), \(M(n)\) for a truncated multiplication, and \(I(n)\) for the
reciprocal of a unit of \(\PLpow\).  The counts below are counts of these
operations, not bit complexities: each block is a polynomial in \(L\) whose
degree grows with the block index --- linearly when \(\deg_L[t^0]A=1\), and
like \(dn\) when \(\deg_L[t^0]A=d\) --- so a further degree-dependent factor is
hidden inside \(C\), \(M\) and \(I\).

\begin{center}
\begin{tabular}{@{}p{0.14\textwidth}p{0.22\textwidth}p{0.23\textwidth}p{0.25\textwidth}@{}}
\hline
Strategy & Operation count to block \(N\) & Strength & Failure mode\\
\hline
Triangular,
 \cref{plt:thm:rev-triangular}
& \(N+1\) residual evaluations; naively \(\sum_{n\le N}C(n{+}1)\), less with
 incremental residual updating
& one exact new block per step, each of them printable; no division beyond the
 leading slope; the residual is itself the certificate
& fails when the leading block of \(\DX F\) is not a unit
 (\cref{plt:rmk:rev-slope-repairs}); strictly sequential, so it cannot exploit
 fast arithmetic\\
\hline
Functional iteration,
 \cref{plt:alg:rev-picard}
& \(N+1\) compositions below \(t^{N+1}\): \((N{+}1)\,C(N{+}1)\)
& shortest existence proof; needs composition only; insensitive to the initial
 guess
& gains only \(\ordt(A)+1\) blocks per step and recomputes low blocks each
 time; an algebraically equivalent rearrangement may gain nothing
 (\cref{plt:rmk:rev-bad-rearrangement})\\
\hline
Newton,
 \cref{plt:cor:rev-newton-schedule}
& \(\BigO(\log N)\) steps at precisions \(q_i\) with \(q_i+1\ge2^i\); total
 \(\sum_i\bigl(2C(q_i)+M(q_i)+I(q_i)\bigr)\)
& residual valuation more than doubles each step; fastest at high order
& needs \(\DX F\), a reciprocal and a composition per step; silently loses
 blocks when the guard precisions of
 \cref{plt:lem:rev-precision-propagation} are not respected\\
\hline
Lagrange--B\"urmann,
 \cref{plt:sec:lagrange-burmann}
& powers \(A^r\), \(r\le N+1\): \(N\) multiplications, plus \(\BigO(N^2)\)
 applications of \(\shiftop{n}{k}\)
& closed universal formula; no recursion; the natural route to coefficient
 identities and degree laws
& expression swell in \(A^r\) and in the repeated derivatives; logarithmic
 degrees reach \(\deg_L\approx d(N{+}1)\) although only \(N\) blocks are
 wanted\\
\hline
\end{tabular}
\end{center}

\begin{remark}[On the cost claims]
\label{plt:rmk:rev-cost-caveat}
% ed.: reports 1, 3, 4 and 6 assert Newton is "asymptotically superior" with no
% cost model stated; the claim is model-dependent and is weakened here.
The assertion that Newton is asymptotically superior, made in reports~1, 3, 4
and~6, is a statement about the cost model: it holds when \(C\), \(M\) and
\(I\) are superadditive, so that the geometric schedule \(q_0<q_1<\cdots<N+1\)
sums to a constant multiple of the cost at the final precision.  It is not a
theorem about the polynomial--logarithmic class, and it can fail if a
composition routine carries a large fixed overhead, or if the coefficient
degrees rather than the block count dominate the work.  We assert only what is
proved here: the number of Newton steps is \(\BigO(\log N)\), and every step is
subject to \cref{plt:lem:rev-precision-propagation}.  Practical schemes seed
with two or three triangular blocks, switch to Newton, and certify with an
independent Lagrange--B\"urmann coefficient; algorithms for the wider exp--log
setting are in
\cite{salvy-shackell-1998,salvy-shackell-1999,vanderhoeven-book}.
\end{remark}

\section{Two worked reversions with certificates}
\label{plt:sec:rev-worked}

\begin{example}[The inverse of \(X+\log X\), through block \(4\)]
\label{plt:ex:rev-omega}
Take \(F=X+\log X=X+L\in\Gnear\), so \(A=L\) and \(\alpha=0\).  For \(G=X+B\),
\begin{equation}\label{plt:eq:rev-omega-residual}
 \RevErr{F}{G}=B+L+\log\bigl(1+Bt\bigr),
\end{equation}
since \(F\compose G=G+\log G=X+B+L+\log(1+Bt)\).  Run
\cref{plt:thm:rev-triangular}.

\emph{Stage \(-1\).}  \(G_{-1}=X\), \(B=0\), so \(\RevErr{F}{G_{-1}}=L\);
\(r_0=L\) and \(b_0=-L\).

\emph{Stage \(0\).}  \(G_0=X-L\).  Now \(B+L=0\) and
\eqref{plt:eq:rev-omega-residual} collapses to a closed form,
\[
 \RevErr{F}{G_0}=\log(1-Lt)
 =-Lt-\tfrac12L^2t^2-\tfrac13L^3t^3-\tfrac14L^4t^4-\cdots,
\]
so \(r_1=-L\) and \(b_1=L\).

\emph{Stage \(1\).}  \(G_1=X-L+Lt\), and
\[
 \RevErr{F}{G_1}
 =\bigl(L-\tfrac12L^2\bigr)t^2
 +\bigl(L^2-\tfrac13L^3\bigr)t^3
 +\bigl(-\tfrac12L^2+L^3-\tfrac14L^4\bigr)t^4
 +\FormalBigOAt{t^{5}}{t},
\]
so \(r_2=L-\tfrac12L^2\) and \(b_2=\tfrac12L^2-L\).

\emph{Stage \(2\).}  \(G_2=G_1+\bigl(\tfrac12L^2-L\bigr)t^2\), and
\[
 \RevErr{F}{G_2}
 =\bigl(-\tfrac13L^3+\tfrac32L^2-L\bigr)t^3
 +\bigl(-\tfrac14L^4+\tfrac32L^3-\tfrac32L^2\bigr)t^4
 +\FormalBigOAt{t^{5}}{t},
\]
so \(b_3=\tfrac13L^3-\tfrac32L^2+L\).

\emph{Stage \(3\).}  \(G_3=G_2+\bigl(\tfrac13L^3-\tfrac32L^2+L\bigr)t^3\), and
\[
 \RevErr{F}{G_3}
 =\bigl(-\tfrac14L^4+\tfrac{11}6L^3-3L^2+L\bigr)t^4
 +\FormalBigOAt{t^{5}}{t},
\]
so \(b_4=\tfrac14L^4-\tfrac{11}6L^3+3L^2-L\).

\emph{Stage \(4\).}  \(G_4=G_3+b_4t^4\) and
\(\RevErr{F}{G_4}=\FormalBigOAt{t^{5}}{t}\), which is the certificate of
\cref{plt:cor:rev-finite-certificate} for \(N=4\).  Collecting,
\begin{equation}\label{plt:eq:rev-omega-expansion}
 \rev{(X+\log X)}
 =X-L+Lt+\Bigl(\frac{L^2}2-L\Bigr)t^2
 +\Bigl(\frac{L^3}3-\frac{3L^2}2+L\Bigr)t^3
 +\Bigl(\frac{L^4}4-\frac{11L^3}6+3L^2-L\Bigr)t^4
 +\FormalBigOAt{t^{5}}{t}.
\end{equation}
The blocks are the canonical zero-based Lambert polynomials,
\(b_n=\LamP{n}(L)\) for \(0\le n\le4\), with \(\LamP{0}(u)=-u\),
\(\LamP{1}(u)=u\), \(\LamP{2}(u)=\tfrac12u^2-u\),
\(\LamP{3}(u)=\tfrac13u^3-\tfrac32u^2+u\) and
\(\LamP{4}(u)=\tfrac14u^4-\tfrac{11}6u^3+3u^2-u\).  The map
\(\rev{(X+\log X)}\) is Wright's omega function \(\WrightOmega\)
\cite{corless-et-al-1996,wright-omega}; its polynomials, the degree law
\(\deg\LamP{n}=n\) for \(n\ge1\) --- which fails at \(n=0\), where
\(\deg\LamP{0}=1\) --- and the Stirling-number formula are developed in the
Lambert part of this volume.
\end{example}

\begin{remark}[The residual is minus the next Lambert polynomial]
\label{plt:rmk:rev-omega-certificate}
Every residual displayed in \cref{plt:ex:rev-omega} has leading block
\(-\LamP{n}(L)\) at level \(n\).  This is not an accident of the computation
but an instance of \cref{plt:prop:rev-next-block}: for
\(G_{n-1}=\Trunc{\WrightOmega}{n}\),
\[
 \RevErr{X+\log X}{\ \Trunc{\WrightOmega}{n}}
 =-\LamP{n}(L)\,t^{n}+\FormalBigOAt{t^{n+1}}{t}.
\]
This is a one-line certificate for any published table of Lambert polynomials,
and it is independent of whatever recurrence generated the table.
\end{remark}

\begin{example}[A two-parameter logarithmic-rational perturbation, through
block \(4\)]
\label{plt:ex:rev-ab}
Let \(a,b\in\K\) and
\[
 F(X)=X+a\log X+\frac{b}{X}=X+aL+bt\ \in\Gnear .
\]
For \(G=X+B\),
\begin{equation}\label{plt:eq:rev-ab-residual}
 \RevErr{F}{G}=B+aL+a\log(1+Bt)+\frac{bt}{1+Bt}.
\end{equation}
Running \cref{plt:thm:rev-triangular} on
\eqref{plt:eq:rev-ab-residual} gives \(r_n=-b_n\) with
\begin{align}
 b_0={}&-aL,\label{plt:eq:rev-ab-b0}\\
 b_1={}&a^2L-b,\label{plt:eq:rev-ab-b1}\\
 b_2={}&a^3\Bigl(\frac{L^2}2-L\Bigr)+ab\,(1-L),\label{plt:eq:rev-ab-b2}\\
 b_3={}&a^4\Bigl(\frac{L^3}3-\frac{3L^2}2+L\Bigr)
        +a^2b\,\bigl(-L^2+3L-1\bigr)-b^2,\label{plt:eq:rev-ab-b3}\\
 b_4={}&a^5\Bigl(\frac{L^4}4-\frac{11L^3}6+3L^2-L\Bigr)
        +a^3b\Bigl(-L^3+\frac{11L^2}2-6L+1\Bigr)
        +ab^2\Bigl(\frac52-3L\Bigr).\label{plt:eq:rev-ab-b4}
\end{align}
The certificates exhibited at the successive stages are
\begin{align*}
 \RevErr{F}{G_{-1}}&=aL+bt,\\
 \RevErr{F}{G_{0}}&=\bigl(b-a^2L\bigr)t+\FormalBigOAt{t^{2}}{t},\\
 \RevErr{F}{G_{1}}&=\Bigl[a^3\Bigl(L-\frac{L^2}2\Bigr)+ab\,(L-1)\Bigr]t^2
   +\FormalBigOAt{t^{3}}{t},\\
 \RevErr{F}{G_{2}}&=-b_3\,t^3+\FormalBigOAt{t^{4}}{t},\\
 \RevErr{F}{G_{3}}&=-b_4\,t^4+\FormalBigOAt{t^{5}}{t},\\
 \RevErr{F}{G_{4}}&=\FormalBigOAt{t^{5}}{t},
\end{align*}
their valuations being \(0,1,2,3,4\) and \(\ge5\) respectively (for
\(a\ne0\)).  Three checks come free.
\begin{itemize}
\item \emph{\(b=0\).}  Then \(b_n=a^{\,n+1}\LamP{n}(L)\) for \(0\le n\le4\); at
\(a=1\) this reproduces \eqref{plt:eq:rev-omega-expansion}.  The pattern
persists for every \(n\), since in the Lagrange--B\"urmann formula of
\cref{plt:sec:lagrange-burmann} only the index \(r=n+1\) contributes to block
\(n\) when \(A=aL\), giving
\(a^{n+1}\frac{(-1)^{n+1}}{(n+1)!}\shiftop{0}{n}L^{n+1}=a^{n+1}\LamP{n}(L)\).
\item \emph{\(a=0\).}  Then \(b_0=b_2=b_4=0\) and \(b_1=-b\), \(b_3=-b^2\).
The vanishing of the even blocks is forced by parity: \(F=X+bX^{-1}\) satisfies
\(F(-X)=-F(X)\), so by the uniqueness in \cref{plt:thm:rev-existence} its
inverse is odd as well, and \(\rev{F}(-X)=-\rev{F}(X)\) reads
\(b_n(-1)^n=-b_n\).  See \cref{plt:ex:rev-catalan}.
\item \emph{Degrees.}  For \(a\ne0\), \(\deg_Lb_n=n\) for \(1\le n\le4\) while
\(\deg_Lb_0=1\), matching the Lambert degree law and its exception at
\(n=0\).
\end{itemize}
Formulas \eqref{plt:eq:rev-ab-b0}--\eqref{plt:eq:rev-ab-b3} agree with the
specializations recorded in reports~1, 2, 3, 4, 5 and~6, which we have checked
term by term against each other and against the triangular recursion;
\eqref{plt:eq:rev-ab-b4} appears in none of them.
\end{example}

\begin{example}[An exact benchmark: the Catalan inverse]
\label{plt:ex:rev-catalan}
The case \(a=0\) of \cref{plt:ex:rev-ab} has a closed form and is therefore the
sharpest available test of signs and of the branch convention.  Let
\(c\in\K^\times\) and \(F=X+ct\).  Since \(\K\) has characteristic zero,
\[
 S\DefinitionEquals\bigl(1-4ct^2\bigr)^{1/2}
 =\sum_{k\ge0}\binom{1/2}{k}(-4c)^kt^{2k}\in\K[[t]]
\]
is well defined, and \(\binom{1/2}{k}(-4)^k=-\tfrac2k\binom{2k-2}{k-1}\) for
\(k\ge1\), i.e.\ \(-2\) times the Catalan number \(C_{k-1}\).  Put
\(G=\tfrac12X(1+S)\).  Then \(\ordt(G-X)\ge1\), so \(G\in\Gnear\), and
\[
 G^2-XG=\tfrac14X^2(1+S)\bigl[(1+S)-2\bigr]=\tfrac14X^2\bigl(S^2-1\bigr)
 =-cX^2t^2=-c,
\]
so \(G^2-XG+c=0\).  Dividing by \(G\), a unit of \(\PLlau\) because
\(\ordt(G)=-1\) with leading block \(1\), gives
\(G+c\MultiplicativeInverseOf{G}=X\), that is \(F\compose G=X\).  By
\cref{plt:thm:rev-existence} this \(G\) is \(\rev{F}\), and expanding,
\[
 \rev{\bigl(X+cX^{-1}\bigr)}
 =X-\sum_{k\ge1}C_{k-1}c^kt^{2k-1}
 =X-\frac cX-\frac{c^2}{X^3}-\frac{2c^3}{X^5}-\frac{5c^4}{X^7}-\cdots,
\]
with Catalan coefficients \(1,1,2,5,14,\dots\)
\cite{comtet-book,graham-knuth-patashnik}.  Blocks \(0,1\) and \(3\) reproduce
\eqref{plt:eq:rev-ab-b0}, \eqref{plt:eq:rev-ab-b1} and
\eqref{plt:eq:rev-ab-b3} at \(a=0\).  Note that the second root
\(\tfrac12X(1-S)\) has \(\ordt=1\) and does not lie in \(\Gnear\): the
near-identity normalization is exactly what selects the branch, and no
coefficient comparison is needed to discard the other one.
\end{example}

\begin{checkpointbox}
Before accepting a computed reversal, record five items: the coefficient field
\(\K\), including any logarithm of a scaling constant that had to be adjoined;
the exclusive cutoff \(N+1\) of the answer; the guard precisions actually used
for \(R\), for \(J\) and for the quotient; the valuation of the right residual;
and the valuation of the left residual.  A result whose residual valuation is
not \emph{strictly} greater than the advertised block index is not certified,
however plausible its coefficients look.
\end{checkpointbox}

\section{Exercises}
\label{plt:sec:rev-exercises}

\begin{exercise}
Let \(F=X+A\in\Gnear\) with \(\ordt(A)=\alpha\).  Show that
\cref{plt:alg:rev-picard} reaches block \(N\) after at most
\(\Ceiling{(N+1-\alpha)/(\alpha+1)}\) iterations, and exhibit an \(A\) with
\(\alpha=0\) for which \(N+1\) iterations are necessary.
\end{exercise}

\begin{exercise}
Carry out the guard-starved Newton step of \cref{plt:ex:rev-guard} in the other
direction: keep \(q=\sigma_R=6\) but under-resolve the Jacobian, taking
\(\sigma_K\in\SetOf{2,3,4}\).  Predict the resulting residual order from
\cref{plt:lem:rev-precision-propagation}, and verify that in each case the
bound \(\min\SetOf{q,\sigma_R,m+\sigma_K}\) is attained exactly.
\end{exercise}

\begin{exercise}
Let \(F=X\log X\).  Show that \(\ordt(\DX F)=0\) with leading block \(L+1\), and
that the first triangular correction already requires
\(\MultiplicativeInverseOf{(L+1)}\).  Show that \(F\) has no compositional
inverse in \(\PLlau\), identify the leading monomial of the true inverse, and
locate the first block in which \(\log\log X\) appears.
\end{exercise}

\begin{exercise}
Let \(F=X+A\in\Gnear\) and \(\alpha=\ordt(A)\).  Using
\cref{plt:prop:rev-residual-order} and \cref{plt:cor:rev-first-order}, show
that \(H\in\Gnear\) satisfies
\(\RevErr{F}{H}=\FormalBigOAt{t^{2\alpha+1}}{t}\) if and only if
\(H=X-A+\FormalBigOAt{t^{2\alpha+1}}{t}\).  Deduce that
\(X-\Trunc{A}{2\alpha+1}\) is the unique element of \(\Gnear\) whose
correction is supported below \(t^{2\alpha+1}\) and whose residual has
valuation at least \(2\alpha+1\).
\end{exercise}

% ---- fragment 09-lagrange ----
\chapter{The Lagrange--B\"urmann formula at infinity}
\label{plt:sec:lagrange-burmann}

Fixed-point and Newton reversion produce the inverse of a near-identity map one
filtration block at a time.  The Lagrange--B\"urmann formula produces all blocks
at once, as a single explicit family of differential monomials in the
perturbation.  This chapter states that formula, proves it \emph{formally} --
no analytic Lagrange inversion theorem is invoked and no convergence is used or
claimed -- derives the finite coefficient formula in the
polynomial--logarithmic basis together with a sharp determination of its index
range, proves the transport formula for a composed observable, and establishes
the degree law that governs the growth of logarithmic degrees under reversion.

Throughout, \(\K\) is a field of characteristic zero, \(X\) is the large
variable, and
\[
  t\DefinitionEquals X^{-1},
  \qquad
  L\DefinitionEquals\log X .
\]
The default analytic regime, whenever one is mentioned at all, is
\(X\to+\infty\) along the positive real ray.  Truncation is always
\emph{exclusive}: \(\Trunc{F}{N}=\sum_{n<N}p_n(L)t^n\).  Coefficient extraction
is written with the series to the right of the bracket, \([t^n]F=p_n(L)\) and
\([t^nL^j]F=a_{n,j}\).

\section{Reversion frames}
\label{plt:sec:lag-frames}

Every argument below uses only four features of \(\PLpow\): it is a
characteristic-zero algebra, it carries a valuation, it is complete for that
valuation, and its distinguished derivation strictly increases the valuation.
Isolating these makes the theorems of this section available verbatim in the
Puiseux--log, iterated Laurent-log and nested-logarithmic enlargements that
\cref{plt:sec:dominant-balance} is forced to introduce.  The cost of the
abstraction is one definition.

\begin{definition}[Reversion frame]
\label{plt:def:lag-frame}
Let \(\K\) be a field of characteristic zero and let \(\Gamma\subseteq\RealNumbers\)
be an ordered subgroup with \(1\in\Gamma\).  A \emph{reversion frame} over
\(\K\) with value group \(\Gamma\) is a quadruple
\((\mathfrak A,\nu,\partial,X)\) in which \(\mathfrak A\) is a commutative
\(\K\)-algebra, \(\partial\) is a \(\K\)-derivation of \(\mathfrak A\),
\(X\in\mathfrak A\), and
\(\nu\colon\mathfrak A\to\Gamma\cup\SetOf{+\infty}\) satisfies the following.
\begin{enumerate}
\item[(F1)] \(\nu(f)=+\infty\) exactly when \(f=0\);
      \(\nu(f+g)\ge\min\SetOf{\nu(f),\nu(g)}\); \(\nu(fg)=\nu(f)+\nu(g)\); and
      \(\nu(\lambda)=0\) for every \(\lambda\in\K^\times\).
\item[(F2)] Call a family \((f_i)_{i\in I}\) in \(\mathfrak A\)
      \emph{summable} when \(\SetOf{i\in I:\nu(f_i)<\gamma}\) is finite for
      every \(\gamma\in\Gamma\).  Every summable family has a sum
      \(\sum_{i}f_i\in\mathfrak A\); summation is \(\K\)-linear, satisfies
      \(\nu\bigl(\sum_if_i\bigr)\ge\inf_i\nu(f_i)\), is unchanged by
      regrouping the index set into disjoint blocks, and satisfies
      \(\bigl(\sum_if_i\bigr)\bigl(\sum_jg_j\bigr)=\sum_{i,j}f_ig_j\) for
      summable \((f_i)\) and \((g_j)\).
\item[(F3)] \(\nu(\partial f)\ge\nu(f)+1\) for every \(f\in\mathfrak A\).
\item[(F4)] \(\partial X=1\).
\end{enumerate}
We write
\[
  \mathfrak A_{>-1}\DefinitionEquals\SetOf{f\in\mathfrak A:\nu(f)>-1},
\]
the set of \emph{admissible perturbations}.  It is closed under addition and
under multiplication by \(\SetOf{\nu\ge0}\), and contains \(0\); it is not in
general a subring, since \(\nu(f^2)=2\nu(f)\) may drop below \(-1\).
\end{definition}

\begin{remark}
\label{plt:rmk:lag-why-minus-one}
The threshold is \(-1\), not \(0\), because the summability of every family
below is governed by the \emph{gain}
\begin{equation}
\label{plt:eq:lag-gain}
  \delta\DefinitionEquals1+\nu(A)>0
\end{equation}
of one derivative against one factor of the perturbation.  In the
polynomial--logarithmic frame the value group is \(\IntegerNumbers\), so
\(\nu_t(A)>-1\) is the same as \(A\in\PLpow\) and
\(\mathcal G(\PLlau)=\Gnear\); the distinction matters only in the
Puiseux--log frames of \cref{plt:sec:dominant-balance}, where a perturbation of
outer order \(-\tfrac12\) is admissible and is exactly what the reduction of a
ramified reversion produces (\cref{plt:ex:lag-puiseux}).
\end{remark}

Condition (F3) forces \(\partial\) to be \(\nu\)-continuous, so it commutes with
sums of summable families: if \((f_i)\) is summable then so is
\((\partial f_i)\), and \(\partial\sum_if_i=\sum_i\partial f_i\), because the two
sides agree modulo \(\SetOf{\nu\ge\gamma}\) for every \(\gamma\) after
discarding the finitely many indices with \(\nu(f_i)<\gamma\).

\begin{lemma}[The polynomial--logarithmic frame]
\label{plt:lem:lag-pl-frame}
\(\bigl(\PLlau,\ordt,\DX,\,t^{-1}\bigr)\) is a reversion frame over \(\K\)
with \(\Gamma=\IntegerNumbers\).
\end{lemma}

\begin{proof}
\(\PLlau=\K[L]((t))\) is a domain because \(\K[L]\) is, and (F1) is the
standard outer valuation together with multiplicativity of leading blocks.
For (F2), let \((f_i)\) be summable.  Taking \(\gamma=n+1\) shows that only
finitely many \(i\) have \(\ordt(f_i)\le n\), so
\([t^n]\sum_if_i\DefinitionEquals\sum_i[t^n]f_i\) is a finite sum of elements of
\(\K[L]\); taking \(\gamma=0\) shows that \(\inf_i\ordt(f_i)\) is finite, so the
resulting series has only finitely many negative powers of \(t\) and lies in
\(\PLlau\).  Linearity, the valuation bound, regrouping and the product rule
are then finite-sum statements at each fixed \(n\).  For (F3), the block
formula \(\DX(t^np(L))=t^{n+1}(\partial_L-n)p(L)\) gives
\(\ordt(\DX f)\ge\ordt(f)+1\).  Finally \(\DX(t^{-1})=\DX X=1\), which is (F4).
\end{proof}

\begin{remark}
\label{plt:rmk:lag-other-frames}
The same verification, word for word, applies to the Puiseux--log ring
\(\K[L]((t^{1/s}))\) with \(\Gamma=s^{-1}\IntegerNumbers\), to the iterated
Laurent-log field \(\PLit\) with \(\Gamma=\IntegerNumbers\), and to the
nested-logarithmic ring \(\K((L^{-1}))[\log L]((t))\), on which
\(\DX(t^np)=t^{n+1}\bigl(\partial_Lp+L^{-1}\partial_{\log L}p-np\bigr)\).  All four
are used in \cref{plt:sec:dominant-balance}.  What is \emph{not} covered is a
value group of rank two, such as the refinement of \(\ordt\) by logarithmic
degree; see \cref{plt:rmk:lag-rank-two}.
\end{remark}

\begin{lemma}[Substitution calculus]
\label{plt:lem:lag-substitution}
Let \((\mathfrak A,\nu,\partial,X)\) be a reversion frame and let
\(H\in\mathfrak A_{>-1}\).  Define
\begin{equation}
\label{plt:eq:lag-shift}
  T_H(f)\DefinitionEquals\sum_{k\ge0}\frac{H^k}{k!}\,\partial^kf,
  \qquad f\in\mathfrak A,
\end{equation}
and write \(f\compose(X+H)\DefinitionEquals T_H(f)\).  Then:
\begin{enumerate}
\item the family in \cref{plt:eq:lag-shift} is summable, its \(k\)-th member
      having \(\nu\ge\nu(f)+k\bigl(1+\nu(H)\bigr)\), and \(\nu(T_Hf)\ge\nu(f)\);
\item \(T_H\) is a \(\K\)-algebra endomorphism of \(\mathfrak A\), continuous
      for \(\nu\), with \(T_0=\mathrm{id}\) and \(T_H(X)=X+H\);
\item \(\partial\compose T_H=(1+\partial H)\cdot(T_H\compose\partial)\);
\item for \(H,H'\in\mathfrak A_{>-1}\),
      \(T_{H'}\compose T_H=T_{H''}\) with \(H''=H'+T_{H'}(H)\in\mathfrak A_{>-1}\);
\item \(\mathcal G(\mathfrak A)\DefinitionEquals\SetOf{X+H:H\in\mathfrak A_{>-1}}\)
      is a group under
      \((X+H)\compose(X+H')\DefinitionEquals T_{H'}(X+H)=X+H'+T_{H'}(H)\),
      and \(T\) is an anti-homomorphism from it to the monoid of endomorphisms
      of \(\mathfrak A\).  In the polynomial--logarithmic frame,
      \(\mathcal G(\PLlau)=\Gnear\).
\end{enumerate}
\end{lemma}

\begin{proof}
(1)  By (F1) and (F3),
\(\nu\bigl(H^k\partial^kf/k!\bigr)\ge k\nu(H)+\nu(f)+k
=\nu(f)+k\bigl(1+\nu(H)\bigr)\), which tends to \(+\infty\) because
\(1+\nu(H)>0\); summability and the valuation bound follow from (F2), the
minimum over \(k\ge0\) being attained at \(k=0\).

(2) Linearity is clear.  For multiplicativity, expand \(\partial^k(fg)\) by
Leibniz and regroup, which (F2) permits:
\[
  T_H(fg)=\sum_{k\ge0}\frac{H^k}{k!}\sum_{i+j=k}\binom{k}{i}\partial^if\,\partial^jg
        =\sum_{i,j\ge0}\frac{H^i\partial^if}{i!}\cdot\frac{H^j\partial^jg}{j!}
        =T_H(f)\,T_H(g).
\]
\(T_H(1)=1\) and \(T_H(\lambda)=\lambda\).  \(T_H(X)=X+H\) because
\(\partial X=1\) and \(\partial^kX=0\) for \(k\ge2\).  Continuity follows from
\(\nu(T_Hf)\ge\nu(f)\) and linearity.

(3)  Differentiate \cref{plt:eq:lag-shift} term by term, which is legitimate
because \(\partial\) commutes with summable sums:
\[
  \partial T_H(f)
  =\sum_{k\ge1}\frac{kH^{k-1}(\partial H)}{k!}\partial^kf
   +\sum_{k\ge0}\frac{H^k}{k!}\partial^{k+1}f
  =(\partial H)\,T_H(\partial f)+T_H(\partial f).
\]

(4)  Let \(\mathfrak A\langle z\rangle_a\) denote the set of formal series
\(P=\sum_{m\ge0}P_mz^m\) with \(P_m\in\mathfrak A\) and \(\nu(P_m)\ge m-a\),
and \(\mathfrak A\langle z\rangle=\bigcup_{a\in\Gamma}\mathfrak A\langle
z\rangle_a\).  By (F1) and (F2) this is a subring of \(\mathfrak A[[z]]\) with
\(\mathfrak A\langle z\rangle_a\mathfrak A\langle z\rangle_b\subseteq\mathfrak
A\langle z\rangle_{a+b}\), and for each \(g\in\mathfrak A_{>-1}\) the
evaluation \(\mathrm{ev}_g(P)=\sum_mP_mg^m\) is a well-defined \(\K\)-algebra
homomorphism \(\mathfrak A\langle z\rangle\to\mathfrak A\), because
\(\nu(P_mg^m)\ge m\bigl(1+\nu(g)\bigr)-a\to+\infty\) makes the family summable
and (F2) makes the product rule a rearrangement of a summable double family.
Put \(\tau(f)=\sum_k(\partial^kf)z^k/k!\); by the computation in (1),
\(\tau(f)\in\mathfrak A\langle z\rangle_{-\nu(f)}\), by the Leibniz argument in
(2) \(\tau\) is a ring homomorphism, and \(\partial_z\tau=\tau\partial\),
\(\mathrm{ev}_0\tau=\mathrm{id}\), \(T_g=\mathrm{ev}_g\compose\tau\).

Fix \(f\in\mathfrak A\) and consider the doubly indexed family
\[
  w_{j,k}\DefinitionEquals
  \frac{\partial^kf}{j!\,(k-j)!}\,\tau(H)^{\,j}\,z^{\,k-j}
  \in\mathfrak A\langle z\rangle,
  \qquad 0\le j\le k .
\]
Its coefficient of \(z^m\) is nonzero only when \(k-j\le m\), and
\(\nu\bigl([z^m]w_{j,k}\bigr)\ge(\nu(f)+k)+\bigl(j\nu(H)+m-k+j\bigr)
=\nu(f)+m+j\bigl(1+\nu(H)\bigr)\).  Since \(k\le m+j\) and
\(1+\nu(H)>0\), for each fixed \(m\) and each \(\gamma\) only finitely many
pairs \((j,k)\) give valuation below \(\gamma\); the family is therefore
summable in every \(z\)-degree, and may be regrouped freely.  Summing first
over \(k\) and then over \(j\) gives
\(\sum_{j}\tau(H)^j\partial_z^{\,j}\tau(f)/j!\), which by
\(\partial_z^{\,j}\tau=\tau\partial^{\,j}\) and multiplicativity of \(\tau\)
equals \(\tau(T_Hf)\); summing first over \(j\) and then over \(k\), and using
the binomial theorem, gives \(\sum_k\bigl(\partial^kf/k!\bigr)(z+\tau(H))^k\).
Hence
\[
  \tau\bigl(T_H f\bigr)
  =\sum_{k\ge0}\frac{\partial^kf}{k!}\bigl(z+\tau(H)\bigr)^k .
\]
Applying the ring homomorphism \(\mathrm{ev}_{H'}\), which likewise commutes
with the regrouping of this summable family, and using
\(\mathrm{ev}_{H'}\tau(H)=T_{H'}(H)\) and \(\mathrm{ev}_{H'}(z)=H'\), gives
\(T_{H'}(T_Hf)=\sum_k(\partial^kf)(H'+T_{H'}H)^k/k!=T_{H''}(f)\).  Finally
\(\nu(H'')\ge\min\SetOf{\nu(H'),\nu(T_{H'}H)}\ge
\min\SetOf{\nu(H'),\nu(H)}>-1\).

(5)  Part (4) says exactly that \(\compose\) is associative on
\(\mathcal G(\mathfrak A)\) and that \(T_{F\compose G}=T_G\compose T_F\); the
element \(X\) is a two-sided identity.  By
\cref{plt:prop:lag-fixed-point} below every element has a right inverse, and a
monoid in which every element has a right inverse is a group: if \(ab=X\) and
\(bc=X\) then \(ba=ba(bc)=b(ab)c=bc=X\).  There is no circularity, because the
proof of \cref{plt:prop:lag-fixed-point} uses only parts (1) and (2).
\end{proof}

\begin{proposition}[Right inverse: existence, uniqueness, characterisation]
\label{plt:prop:lag-fixed-point}
Let \((\mathfrak A,\nu,\partial,X)\) be a reversion frame,
\(A\in\mathfrak A_{>-1}\), and \(F=X+A\).  There is exactly one
\(B\in\mathfrak A_{>-1}\) with
\begin{equation}
\label{plt:eq:lag-fixedpoint}
  F\compose(X+B)=X,
  \qquad\text{equivalently}\qquad
  B=-T_B(A)=-A\compose(X+B).
\end{equation}
It satisfies \(\nu(B)\ge\nu(A)\).  We write \(\rev{F}=X+B\).
\end{proposition}

\begin{proof}
The two displayed conditions are the same because
\(F\compose(X+B)=T_B(X+A)=X+B+T_B(A)\).

\emph{A contraction estimate.}  Let \(B,B'\in\mathfrak A_{>-1}\) and put
\(m=\min\SetOf{\nu(B),\nu(B')}>-1\).  Then
\begin{equation}
\label{plt:eq:lag-contraction}
  \nu\bigl(T_{B'}(A)-T_B(A)\bigr)\ \ge\ \nu(B'-B)+\delta,
  \qquad \delta=1+\nu(A)>0 .
\end{equation}
Indeed \(T_{B'}(A)-T_B(A)=\sum_{k\ge1}\bigl(B'^{\,k}-B^k\bigr)\partial^kA/k!\),
the \(k=0\) terms cancelling, and for \(k\ge1\)
\(B'^{\,k}-B^k=(B'-B)\sum_{i=0}^{k-1}B'^{\,i}B^{k-1-i}\) has valuation at least
\(\nu(B'-B)+(k-1)m\), while \(\nu(\partial^kA)\ge\nu(A)+k\).  The \(k\)-th term
therefore has valuation at least
\(\nu(B'-B)+\nu(A)+k(1+m)-m\), which is increasing in \(k\) because
\(1+m>0\); its value at \(k=1\) is \(\nu(B'-B)+\nu(A)+1\).

\emph{Uniqueness.}  If \(B,B'\) both solve \cref{plt:eq:lag-fixedpoint}, then
\cref{plt:eq:lag-contraction} gives
\(\nu(B-B')\ge\nu(B'-B)+\delta\) with \(\delta>0\), forcing
\(\nu(B-B')=+\infty\).

\emph{Existence.}  Set \(B_0=0\) and \(B_{n+1}=-T_{B_n}(A)\); by
\cref{plt:lem:lag-substitution}(1), \(\nu(B_{n+1})\ge\nu(A)>-1\), so every
\(B_n\) lies in \(\mathfrak A_{>-1}\) and
\cref{plt:eq:lag-contraction} applies with
\(m\ge\min\SetOf{0,\nu(A)}\).  Hence
\(\nu(B_{n+1}-B_n)\ge\nu(B_n-B_{n-1})+\delta\), so
\(\nu(B_{n+1}-B_n)\ge\nu(A)+n\delta\).  The family \((B_{n+1}-B_n)_{n\ge0}\) is
summable, and \(B\DefinitionEquals\sum_{n\ge0}(B_{n+1}-B_n)\) satisfies
\(\nu(B-B_n)\ge\nu(A)+n\delta\).  Applying
\cref{plt:eq:lag-contraction} once more,
\(\nu\bigl(T_B(A)-T_{B_n}(A)\bigr)\ge\nu(A)+(n+1)\delta\), so
\(\nu\bigl(B+T_B(A)\bigr)\ge
\min\SetOf{\nu(B-B_{n+1}),\nu(T_{B_n}A-T_BA)}\to+\infty\)
and \(B=-T_B(A)\).  Finally \(\nu(B)=\nu(T_B(A))\ge\nu(A)\).
\end{proof}

\begin{remark}
The characterisation \cref{plt:eq:lag-fixedpoint} is the only property of
\(\rev{F}\) used below; the two-sidedness supplied by
\cref{plt:lem:lag-substitution}(5) is never needed in the proof of the
Lagrange--B\"urmann formula and is recorded only so that
``the'' inverse is unambiguous.
\end{remark}

\section{Formal summability}
\label{plt:sec:lag-summability}

\begin{lemma}[Order of the \(r\)-th Lagrange term]
\label{plt:lem:lag-order}
Let \((\mathfrak A,\nu,\partial,X)\) be a reversion frame and
\(A\in\mathfrak A\).  For every \(r\ge1\),
\begin{equation}
\label{plt:eq:lag-order}
  \nu\Bigl(\partial^{\,r-1}\bigl(A^r\bigr)\Bigr)\ \ge\ r\,\nu(A)+r-1
  \ =\ r\delta-1,
  \qquad \delta=1+\nu(A).
\end{equation}
If \(A\in\mathfrak A_{>-1}\), so that \(\delta>0\), the family
\(\bigl((-1)^r\partial^{\,r-1}(A^r)/r!\bigr)_{r\ge1}\) is summable, and its sum
\(S\) satisfies \(\nu(S)\ge\nu(A)>-1\), whence \(S\in\mathfrak A_{>-1}\).
\end{lemma}

\begin{proof}
\(\nu(A^r)=r\nu(A)\) by (F1) and \(r-1\) applications of (F3) add at least
\(r-1\).  If \(\delta>0\) the right-hand side of \cref{plt:eq:lag-order} tends
to \(+\infty\), so the family is summable by (F2); its minimum over \(r\ge1\) is
attained at \(r=1\) and equals \(\nu(A)\), which bounds \(\nu(S)\).
\end{proof}

The hypothesis \(\nu(A)>-1\) is not merely convenient: it is exactly the
summability condition, in the polynomial--logarithmic frame and in its
Puiseux--log enlargements alike.

\begin{proposition}[Sharp summability criterion]
\label{plt:prop:lag-summability-sharp}
Fix \(s\in\PositiveIntegers\) and let \(A\in\K[L]((t^{1/s}))\) be nonzero, with
\(\alpha\DefinitionEquals\ordt(A)\in s^{-1}\IntegerNumbers\).  The family
\(\bigl((-1)^r\DX^{\,r-1}(A^r)/r!\bigr)_{r\ge1}\) is summable if and only if
\(\alpha>-1\).  Indeed, if \(\alpha\le-1\) then
\begin{equation}
\label{plt:eq:lag-order-exact}
  \ordt\Bigl(\DX^{\,r-1}\bigl(A^r\bigr)\Bigr)=r(\alpha+1)-1
  \qquad\text{for every }r\ge1,
\end{equation}
which is nonincreasing in \(r\) and hence never tends to \(+\infty\).  For
\(s=1\), that is in \(\PLlau\), the value group is \(\IntegerNumbers\) and the
criterion reads \(\ordt(A)\ge0\), i.e.\ \(A\in\PLpow\).
\end{proposition}

\begin{proof}
Sufficiency is \cref{plt:lem:lag-order}.  For necessity assume
\(\alpha\le-1\) and write \(A=t^{\alpha}p(L)+A_1\) with \(p\ne0\) and
\(\ordt(A_1)>\alpha\).  Then
\(A^r=t^{r\alpha}p^r+(\text{terms of }\ordt>r\alpha)\).  Applying the
repeated-derivative formula
\(\DX^{\,r-1}\bigl(t^{n}p(L)\bigr)=t^{n+r-1}\shiftop{n}{r-1}p(L)\) with
\(n=r\alpha\),
\[
  \DX^{\,r-1}\bigl(A^r\bigr)
  =t^{r\alpha+r-1}\prod_{j=0}^{r-2}\bigl(\partial_L-r\alpha-j\bigr)p(L)^r
   +\bigl(\text{terms of }\ordt>r\alpha+r-1\bigr).
\]
For \(0\le j\le r-2\) one has \(-r\alpha-j\ge r-(r-2)=2>0\), so each factor
\(\partial_L-(r\alpha+j)\) multiplies the leading coefficient of a nonzero
polynomial by the nonzero scalar \(-r\alpha-j\) and in particular is injective
on \(\K[L]\setminus\SetOf{0}\).  Hence the displayed leading block is nonzero
and \cref{plt:eq:lag-order-exact} holds.  Since \(\alpha+1\le0\), the exponent
\(r(\alpha+1)-1\) is at most \(-1\) for all \(r\), so
\(\SetOf{r:\ordt<0}\) is infinite and the family is not summable.
\end{proof}

% ed.: S1, S2, S3, S4 and S5 all assert only the sufficiency direction ("the
% r-th term has t-order at least r-1"), and all state it with the hypothesis
% ord_t(A) >= 0.  The correct threshold is ord_t(A) > -1; the two agree on
% K[L]((t)), whose value group is Z, but they differ on the Puiseux--log
% enlargement, where the weaker hypothesis is exactly what the reduction of a
% ramified reversion needs (plt:ex:lag-puiseux).  None of the sources observes
% that the condition is also necessary, which is what makes the misuse in
% plt:warn:lag-scale-trap detectable.
\begin{remark}
\label{plt:rmk:lag-nearid-meaning}
\Cref{plt:prop:lag-summability-sharp} makes precise what ``near identity''
means here, and it is \emph{not} that \(A\) is bounded.  For
\(A=L=\log X\) one has \(A\to+\infty\) yet \(\ordt(A)=0>-1\), and the theory
applies; for \(A=X/L\) one has \(A=\LittleOAt{X}{X\to+\infty}\) yet
\(\ordt(A)=-1\), and it does not.  The decisive quantity is the outer order
measured against the threshold \(-1=\ordt(X)\), equivalently the dominance
comparison \(A\precdom X\) \emph{refined to the outer filtration}: logarithmic
gains are invisible to \(\ordt\), and a perturbation of the same outer order as
\(X\) itself is never admissible however many logarithms separate it from
\(X\).
\end{remark}

\section{The parameter identity}
\label{plt:sec:lag-parameter}

The classical Lagrange--B\"urmann theorem is a statement about a formal
parameter.  We prove the version we need -- the \emph{shifted}, or reversion,
form -- from scratch, over an arbitrary commutative \(\RationalNumbers\)-algebra
and with no invertibility hypothesis whatsoever.

\begin{lemma}[Lagrange reversion identity]
\label{plt:lem:lag-abstract}
Let \(C\) be a commutative \(\RationalNumbers\)-algebra and let
\(\psi\in C[[z]]\).  For \(f\in C[[z]]\) and
\(w\in\varepsilon\,C[[z]][[\varepsilon]]\) write
\(f_{[w]}\DefinitionEquals\sum_{k\ge0}(\partial_z^kf)w^k/k!\).  There is exactly one
\(u\in\varepsilon\,C[[z]][[\varepsilon]]\) with
\begin{equation}
\label{plt:eq:lag-param-fixedpoint}
  u=\varepsilon\,\psi_{[u]},
\end{equation}
and for every \(\gamma\in C[[z]]\) and every \(m\ge1\),
\begin{equation}
\label{plt:eq:lag-param}
  [\varepsilon^m]\,\gamma_{[u]}
  =\frac1{m!}\,\partial_z^{\,m-1}\bigl(\gamma'\psi^m\bigr),
  \qquad \gamma'\DefinitionEquals\partial_z\gamma .
\end{equation}
\end{lemma}

\begin{proof}
Existence and uniqueness of \(u\): comparing coefficients of \(\varepsilon^m\)
in \cref{plt:eq:lag-param-fixedpoint} determines \([\varepsilon^m]u\) from
\([\varepsilon^1]u,\ldots,[\varepsilon^{m-1}]u\), because \(w^k\) contributes to
\(\varepsilon^m\) only for \(k\le m\) and then only through those coefficients.

For fixed \(u\), the map \(f\mapsto f_{[u]}\) is a \(C\)-algebra homomorphism
of \(C[[z]]\) into \(C[[z]][[\varepsilon]]\) (Leibniz, exactly as in
\cref{plt:lem:lag-substitution}(2)), and differentiating its defining series
term by term gives the two chain rules
\begin{equation}
\label{plt:eq:lag-chain}
  \partial_\varepsilon\bigl(f_{[u]}\bigr)=f'_{[u]}\,\partial_\varepsilon u,
  \qquad
  \partial_z\bigl(f_{[u]}\bigr)=f'_{[u]}\bigl(1+\partial_zu\bigr).
\end{equation}
Applying \cref{plt:eq:lag-chain} to \cref{plt:eq:lag-param-fixedpoint},
\[
  \partial_\varepsilon u=\psi_{[u]}+\varepsilon\psi'_{[u]}\partial_\varepsilon u,
  \qquad
  \partial_zu=\varepsilon\psi'_{[u]}\bigl(1+\partial_zu\bigr).
\]
Since \(\varepsilon\psi'_{[u]}\in\varepsilon C[[z]][[\varepsilon]]\), the element
\(1-\varepsilon\psi'_{[u]}\) is a unit; solving the second relation gives
\(1+\partial_zu=\bigl(1-\varepsilon\psi'_{[u]}\bigr)^{-1}\), and substituting into
the first gives the key identity
\begin{equation}
\label{plt:eq:lag-key}
  \partial_\varepsilon u=\psi_{[u]}\bigl(1+\partial_zu\bigr).
\end{equation}
Now let \(\gamma\in C[[z]]\) be arbitrary and let \(\Gamma\in C[[z]]\) be the
formal antiderivative of \(\gamma'\psi\) with \(\Gamma(0)=0\), which exists
because \(C\) is a \(\RationalNumbers\)-algebra.  Using
\cref{plt:eq:lag-chain}, \cref{plt:eq:lag-key} and multiplicativity of
\(f\mapsto f_{[u]}\),
\[
  \partial_\varepsilon\bigl(\gamma_{[u]}\bigr)
  =\gamma'_{[u]}\psi_{[u]}\bigl(1+\partial_zu\bigr)
  =\bigl(\gamma'\psi\bigr)_{[u]}\bigl(1+\partial_zu\bigr)
  =\Gamma'_{[u]}\bigl(1+\partial_zu\bigr)
  =\partial_z\bigl(\Gamma_{[u]}\bigr).
\]
Extracting \([\varepsilon^{m-1}]\) from both ends and using
\([\varepsilon^{m-1}]\partial_\varepsilon=m\,[\varepsilon^m]\),
\begin{equation}
\label{plt:eq:lag-recursion}
  m\,[\varepsilon^m]\gamma_{[u]}
  =\partial_z\Bigl([\varepsilon^{m-1}]\Gamma_{[u]}\Bigr),
  \qquad m\ge1 .
\end{equation}
We prove \cref{plt:eq:lag-param} by induction on \(m\), the statement being
quantified over all \(\gamma\in C[[z]]\).  For \(m=1\),
\([\varepsilon^0]\Gamma_{[u]}=\Gamma\) because \(u\in\varepsilon
C[[z]][[\varepsilon]]\), so \cref{plt:eq:lag-recursion} gives
\([\varepsilon^1]\gamma_{[u]}=\Gamma'=\gamma'\psi\), which is
\cref{plt:eq:lag-param} with \(m=1\).  For \(m\ge2\), the inductive hypothesis
applied to \(\Gamma\) gives
\([\varepsilon^{m-1}]\Gamma_{[u]}=\partial_z^{\,m-2}\bigl(\Gamma'\psi^{m-1}\bigr)/(m-1)!
=\partial_z^{\,m-2}\bigl(\gamma'\psi^{m}\bigr)/(m-1)!\), and
\cref{plt:eq:lag-recursion} yields
\(m[\varepsilon^m]\gamma_{[u]}=\partial_z^{\,m-1}(\gamma'\psi^m)/(m-1)!\).
\end{proof}

% ed.: S3's appendix proves the unshifted formula by a formal residue change of
% variables, which is valid only when the constant term of the auxiliary series
% is invertible, and repairs this by a universal-coefficient localisation
% argument.  S6 replaces that argument by the phrase "an integration by parts
% reduces the resulting residue".  The proof above needs neither device.
\begin{remark}
\label{plt:rmk:lag-residue-route}
A second argument runs through formal residues.  Writing
\(\operatorname*{res}_\varepsilon\) for the coefficient of
\(\varepsilon^{-1}\) in a formal Laurent series, one extracts
\([\varepsilon^m]\) as a residue and substitutes
\(\varepsilon=z/\psi(z)\).  That change of variables is legitimate only when
\(\psi(0)\) is invertible, so the general case must be recovered by performing
the computation in the universal polynomial ring generated by the coefficients
of \(\psi\) with the constant coefficient inverted, observing that both sides
are polynomial in finitely many of those generators at each fixed \(m\), and
specialising.  The route is instructive, but the
\(\partial_\varepsilon\)--\(\partial_z\) exchange \cref{plt:eq:lag-key} used
above needs no localisation, no residue theory, and no hypothesis on
\(\psi(0)\).  Classical accounts of formal Lagrange inversion are in
\cite{comtet-book,graham-knuth-patashnik}.
\end{remark}

\section{Specialising the bookkeeping parameter}
\label{plt:sec:lag-specialisation}

Every source we merge performs the step ``now set \(\varepsilon=1\)'' and
justifies it by the sentence that only finitely many terms contribute below any
fixed outer order.  That sentence is true but it is not a proof: evaluation at
\(\varepsilon=1\) has to be a \emph{ring homomorphism} defined on the subalgebra
in which the computation takes place, and the fixed-point equation
\cref{plt:eq:lag-param-fixedpoint} contains an infinite sum that must be shown
to survive it.  The following lemma supplies the missing object.

\begin{lemma}[Weighted deformation algebra and evaluation at one]
\label{plt:lem:lag-eval}
Let \((\mathfrak A,\nu,\partial,X)\) be a reversion frame and fix a weight
\(\delta\in\Gamma_{>0}\).  For \(a\in\Gamma\) put
\[
  \mathfrak A\langle\varepsilon\rangle_a
  \DefinitionEquals\Bigl\{\,u=\sum_{m\ge0}u_m\varepsilon^m\in\mathfrak A[[\varepsilon]]
      \ :\ \nu(u_m)\ge m\delta-a\ \text{for all }m\,\Bigr\},
  \qquad
  \mathfrak A\langle\varepsilon\rangle
  =\bigcup_{a\in\Gamma}\mathfrak A\langle\varepsilon\rangle_a .
\]
Then:
\begin{enumerate}
\item \(\mathfrak A\langle\varepsilon\rangle\) is a subring of
      \(\mathfrak A[[\varepsilon]]\) containing \(\mathfrak A\) and
      \(\varepsilon\), with
      \(\mathfrak A\langle\varepsilon\rangle_a\,\mathfrak A\langle\varepsilon
      \rangle_b\subseteq\mathfrak A\langle\varepsilon\rangle_{a+b}\), and it is
      stable under the coefficientwise extension of \(\partial\);
\item \(\mathrm{ev}\colon\mathfrak A\langle\varepsilon\rangle\to\mathfrak A\),
      \(\mathrm{ev}(u)=\sum_{m\ge0}u_m\), is a well-defined \(\K\)-algebra
      homomorphism with \(\mathrm{ev}\compose\partial=\partial\compose\mathrm{ev}\),
      \(\mathrm{ev}(\varepsilon)=1\), and
      \(\nu(\mathrm{ev}(u))\ge-a\) for \(u\in\mathfrak A\langle\varepsilon
      \rangle_a\);
\item if \(\bigl(f^{(k)}\bigr)_{k\ge0}\) is a family in
      \(\mathfrak A\langle\varepsilon\rangle_a\) with
      \(f^{(k)}\in\varepsilon^k\mathfrak A[[\varepsilon]]\), then
      \(\sum_kf^{(k)}\in\mathfrak A\langle\varepsilon\rangle_a\), the family
      \(\bigl(\mathrm{ev}\,f^{(k)}\bigr)_k\) is summable in \(\mathfrak A\), and
      \(\mathrm{ev}\bigl(\sum_kf^{(k)}\bigr)=\sum_k\mathrm{ev}\bigl(f^{(k)}\bigr)\).
\end{enumerate}
\end{lemma}

\begin{proof}
(1)  If \(\nu(u_m)\ge m\delta-a\) and \(\nu(v_m)\ge m\delta-b\) then
\(\nu\bigl(\sum_{i+j=m}u_iv_j\bigr)\ge\min_{i+j=m}(i\delta-a+j\delta-b)
=m\delta-a-b\).  Stability under \(\partial\) follows from (F3), which gives
\(\nu(\partial u_m)\ge m\delta-(a-1)\).

(2)  For \(u\in\mathfrak A\langle\varepsilon\rangle_a\) the family
\((u_m)_{m\ge0}\) is summable, since \(\nu(u_m)\ge m\delta-a\to+\infty\); the
sum exists by (F2) and has \(\nu\ge\inf_m(m\delta-a)=-a\).  Additivity is
clear, and multiplicativity is the product rule of (F2):
\(\mathrm{ev}(uv)=\sum_m\sum_{i+j=m}u_iv_j=\sum_{i,j}u_iv_j
=\bigl(\sum_iu_i\bigr)\bigl(\sum_jv_j\bigr)\), the regrouping being licensed
because \(\nu(u_iv_j)\ge(i+j)\delta-a-b\) makes the double family summable.
Commutation with \(\partial\) is the continuity of \(\partial\).

(3)  Since \(f^{(k)}\in\varepsilon^k\mathfrak A[[\varepsilon]]\), for each \(m\)
only the indices \(k\le m\) contribute to \([\varepsilon^m]\sum_kf^{(k)}\), so
the sum is defined coefficientwise and lies in
\(\mathfrak A\langle\varepsilon\rangle_a\).  The doubly indexed family
\(\bigl([\varepsilon^m]f^{(k)}\bigr)_{k,m}\) is summable: it vanishes for
\(k>m\), and \(\nu\bigl([\varepsilon^m]f^{(k)}\bigr)\ge m\delta-a\), so
\(\SetOf{(k,m):\nu<\gamma}\subseteq
\SetOf{(k,m):m\delta<\gamma+a,\ k\le m}\) is finite.  Both iterated sums
therefore equal the unordered sum, by the regrouping clause of (F2).
\end{proof}

\section{The formula and the transport formula}
\label{plt:sec:lag-main}

\begin{theorem}[Lagrange--B\"urmann reversion at infinity]
\label{plt:thm:lag-burmann}
Let \((\mathfrak A,\nu,\partial,X)\) be a reversion frame and let
\(A\in\mathfrak A_{>-1}\), \(F=X+A\).  Then the family
\(\bigl((-1)^r\partial^{\,r-1}(A^r)/r!\bigr)_{r\ge1}\) is summable and
\begin{equation}
\label{plt:eq:lag-burmann}
  \rev{F}
  =X+\sum_{r\ge1}\frac{(-1)^r}{r!}\,\partial^{\,r-1}\bigl(A^r\bigr).
\end{equation}
In the polynomial--logarithmic frame this reads
\begin{equation}
\label{plt:eq:lag-burmann-pl}
  \rev{(X+A)}(X)
  =X+\sum_{r\ge1}\frac{(-1)^r}{r!}\,\DX^{\,r-1}\bigl(A(X)^r\bigr),
  \qquad A\in\PLpow ,
\end{equation}
and the \(r\)-th summand has \(\ordt\ge r-1\).
\end{theorem}

\begin{proof}
Summability is \cref{plt:lem:lag-order}; call the sum \(S\), so
\(\nu(S)\ge\nu(A)>-1\).  By \cref{plt:prop:lag-fixed-point} it suffices to
prove \(S=-T_S(A)\).  Write \(\delta=1+\nu(A)>0\) and use this \(\delta\) as
the weight in \cref{plt:lem:lag-eval}.

Take \(C=\mathfrak A\) in \cref{plt:lem:lag-abstract} and let
\(\tau\colon\mathfrak A\to\mathfrak A[[z]]\),
\(\tau(f)=\sum_k(\partial^kf)z^k/k!\), be the Taylor-shift homomorphism of
\cref{plt:lem:lag-substitution}(4); recall \(\partial_z\tau=\tau\partial\) and
\(\mathrm{ev}_{z=0}\tau=\mathrm{id}\).  Put
\[
  \psi\DefinitionEquals\tau(-A)\in\mathfrak A[[z]],
\]
and let \(u\in\varepsilon\mathfrak A[[z]][[\varepsilon]]\) be the unique solution
of \cref{plt:eq:lag-param-fixedpoint}.  Apply
\cref{plt:eq:lag-param} with \(\gamma=\tau(X)=X+z\), so that
\(\gamma'=1=\tau(1)\):
\[
  [\varepsilon^m]\,\gamma_{[u]}
  =\frac1{m!}\partial_z^{\,m-1}\bigl(\psi^m\bigr)
  =\frac{(-1)^m}{m!}\,\partial_z^{\,m-1}\tau\bigl(A^m\bigr)
  =\frac{(-1)^m}{m!}\,\tau\Bigl(\partial^{\,m-1}\bigl(A^m\bigr)\Bigr),
\]
using that \(\tau\) is a ring homomorphism intertwining \(\partial_z\) with
\(\partial\).  Let
\(\mathrm{ev}_{z=0}\colon\mathfrak A[[z]][[\varepsilon]]\to\mathfrak
A[[\varepsilon]]\) act coefficientwise in \(\varepsilon\); it is a ring
homomorphism.  Setting \(U\DefinitionEquals\mathrm{ev}_{z=0}(u)\in
\varepsilon\mathfrak A[[\varepsilon]]\) and using
\(\gamma_{[u]}=\sum_k(\partial_z^k\gamma)u^k/k!=(X+z)+u\), we obtain
\begin{equation}
\label{plt:eq:lag-U}
  U=\sum_{m\ge1}\frac{(-1)^m}{m!}\,\partial^{\,m-1}\bigl(A^m\bigr)\,\varepsilon^m .
\end{equation}
Likewise, applying \(\mathrm{ev}_{z=0}\) to
\cref{plt:eq:lag-param-fixedpoint} and using
\(\psi_{[u]}=\sum_k(\partial_z^k\psi)u^k/k!\) together with
\(\mathrm{ev}_{z=0}\partial_z^k\tau(-A)=\partial^k(-A)\),
\begin{equation}
\label{plt:eq:lag-Ufix}
  U=-\varepsilon\sum_{k\ge0}\frac{\partial^kA}{k!}\,U^k .
\end{equation}

Now specialise.  By \cref{plt:eq:lag-U} and \cref{plt:lem:lag-order},
\(\nu\bigl([\varepsilon^m]U\bigr)\ge m\delta-1\), so
\(U\in\mathfrak A\langle\varepsilon\rangle_1\) and
\(\mathrm{ev}(U)=S\) by construction.  Since \(U\) has zero constant term,
\([\varepsilon^m](U^k)=0\) for \(m<k\) and, splitting \(m=m_1+\cdots+m_k\) with
every \(m_i\ge1\),
\(\nu\bigl([\varepsilon^m](U^k)\bigr)\ge\sum_i(m_i\delta-1)=m\delta-k\); with
\(\nu(\partial^kA)\ge\nu(A)+k\ge k-1+\delta\ge k\cdot\min\SetOf{1,\delta}\) --
more simply, with \(\nu(\partial^kA)\ge\nu(A)+k\) -- this gives
\[
  \nu\Bigl([\varepsilon^m]\Bigl(\tfrac{\partial^kA}{k!}U^k\Bigr)\Bigr)
  \ \ge\ \bigl(\nu(A)+k\bigr)+\bigl(m\delta-k\bigr)=m\delta+\nu(A),
  \qquad
  \tfrac{\partial^kA}{k!}U^k\in\varepsilon^k\mathfrak A[[\varepsilon]] .
\]
So \cref{plt:lem:lag-eval}(3) applies to \(f^{(k)}=\partial^kA\,U^k/k!\) with
\(a=-\nu(A)\), and applying \(\mathrm{ev}\) to \cref{plt:eq:lag-Ufix} gives
\[
  S=\mathrm{ev}(U)=-\sum_{k\ge0}\frac{\partial^kA}{k!}\,\mathrm{ev}(U)^k
   =-\sum_{k\ge0}\frac{S^k}{k!}\partial^kA=-T_S(A).
\]
By \cref{plt:prop:lag-fixed-point}, \(X+S=\rev{F}\).
\end{proof}

The same construction, run with a general observable in place of \(X\),
produces the expansion of any transseries along the inverse map without first
building the inverse.

\begin{theorem}[Transport formula for a composed observable]
\label{plt:thm:lag-transport}
Let \((\mathfrak A,\nu,\partial,X)\) be a reversion frame, let
\(A\in\mathfrak A_{\ge0}\), \(F=X+A\), and let \(\varphi\in\mathfrak A\) be
arbitrary.  Then the family
\(\bigl((-1)^r\partial^{\,r-1}\bigl((\partial\varphi)A^r\bigr)/r!\bigr)_{r\ge1}\)
is summable, the \(r\)-th member having
\begin{equation}
\label{plt:eq:lag-transport-order}
  \nu\Bigl(\partial^{\,r-1}\bigl((\partial\varphi)A^r\bigr)\Bigr)
  \ \ge\ \nu(\varphi)+r\bigl(1+\nu(A)\bigr)\ \ge\ \nu(\varphi)+r,
\end{equation}
and
\begin{equation}
\label{plt:eq:lag-transport}
  \varphi\compose\rev{F}
  =\varphi+\sum_{r\ge1}\frac{(-1)^r}{r!}\,
      \partial^{\,r-1}\bigl((\partial\varphi)\,A^r\bigr).
\end{equation}
In particular \([t^N]\bigl(\varphi\compose\rev{F}\bigr)\) in the
polynomial--logarithmic frame receives contributions only from
\(r\le N-\ordt(\varphi)\).
\end{theorem}

% ed.: S3 states this "whenever both sides are interpreted to a fixed finite
% t-precision" and S6 "for any H for which the family is summable"; both
% hypotheses are vacuous or unverifiable.  The order estimate
% plt:eq:lag-transport-order shows that no hypothesis beyond membership in the
% frame is needed, since nu(phi) is finite for phi nonzero.
\begin{proof}
For \cref{plt:eq:lag-transport-order}: \(\nu(\partial\varphi)\ge\nu(\varphi)+1\)
by (F3), \(\nu(A^r)=r\nu(A)\) by (F1), and \(r-1\) further derivatives add at
least \(r-1\); the total is
\(\nu(\varphi)+1+r\nu(A)+r-1=\nu(\varphi)+r(1+\nu(A))\), which is at least
\(\nu(\varphi)+r\to+\infty\).  Summability follows from (F2).

Run the proof of \cref{plt:thm:lag-burmann} with \(\gamma\DefinitionEquals
\tau(\varphi)\) in \cref{plt:eq:lag-param}.  Since
\(\gamma'=\partial_z\tau(\varphi)=\tau(\partial\varphi)\) and \(\tau\) is a ring
homomorphism intertwining \(\partial_z\) with \(\partial\),
\[
  [\varepsilon^m]\gamma_{[u]}
  =\frac1{m!}\partial_z^{\,m-1}\bigl(\tau(\partial\varphi)\,\tau(-A)^m\bigr)
  =\frac{(-1)^m}{m!}\,
    \tau\Bigl(\partial^{\,m-1}\bigl((\partial\varphi)A^m\bigr)\Bigr).
\]
Applying \(\mathrm{ev}_{z=0}\) and writing
\(c_m\DefinitionEquals(-1)^m\partial^{\,m-1}\bigl((\partial\varphi)A^m\bigr)/m!\),
\begin{equation}
\label{plt:eq:lag-transport-eps}
  \sum_{k\ge0}\frac{\partial^k\varphi}{k!}\,U^k
  =\varphi+\sum_{m\ge1}c_m\,\varepsilon^m
  \qquad\text{in }\mathfrak A[[\varepsilon]],
\end{equation}
where \(U\) is as in \cref{plt:eq:lag-U}.  By
\cref{plt:eq:lag-transport-order}, \(\nu(c_m)\ge\nu(\varphi)+m\), so the
right-hand side of \cref{plt:eq:lag-transport-eps} lies in
\(\mathfrak A\langle\varepsilon\rangle_{-\nu(\varphi)}\).  On the left,
\(\nu\bigl([\varepsilon^m](\partial^k\varphi\,U^k/k!)\bigr)
\ge\bigl(\nu(\varphi)+k\bigr)+(m-k)=\nu(\varphi)+m\) and
\(\partial^k\varphi\,U^k/k!\in\varepsilon^k\mathfrak A[[\varepsilon]]\), so
\cref{plt:lem:lag-eval}(3) applies with \(a=-\nu(\varphi)\).  Applying
\(\mathrm{ev}\) to \cref{plt:eq:lag-transport-eps} and using
\(\mathrm{ev}(U)=S\) with \(X+S=\rev{F}\) gives
\(T_S(\varphi)=\varphi+\sum_{m\ge1}c_m\), which is
\cref{plt:eq:lag-transport}.
\end{proof}

\begin{corollary}
\label{plt:cor:lag-observables}
With \(F=X+A\), \(A\in\PLpow\):
\begin{align}
  \rev{F}&=X+\sum_{r\ge1}\frac{(-1)^r}{r!}\DX^{\,r-1}\bigl(A^r\bigr)
  &&(\varphi=X),\label{plt:eq:lag-obs-X}\\
  \log\rev{F}&=L+\sum_{r\ge1}\frac{(-1)^r}{r!}\DX^{\,r-1}\bigl(t\,A^r\bigr)
  &&(\varphi=L),\label{plt:eq:lag-obs-L}\\
  \bigl(\rev{F}\bigr)^{-q}&=t^q-q\sum_{r\ge1}\frac{(-1)^r}{r!}
      \DX^{\,r-1}\bigl(t^{q+1}A^r\bigr)
  &&(\varphi=t^q,\ q\in\IntegerNumbers).\label{plt:eq:lag-obs-power}
\end{align}
\end{corollary}

\begin{proof}
Substitute \(\varphi=X=t^{-1}\), \(\varphi=L\) and \(\varphi=t^q\) in
\cref{plt:eq:lag-transport}, using \(\DX X=1\), \(\DX L=t\) and
\(\DX t^q=-q\,t^{q+1}\).  In \cref{plt:eq:lag-obs-L} the left-hand side means
the observable \(L\) evaluated along \(\rev{F}\), that is
\(L\compose\rev{F}=\log\bigl(\rev{F}\bigr)\) in the sense of the generator
substitution \(L\compose(X+B)=L+\log(1+tB)\).
\end{proof}

\begin{example}
\label{plt:ex:lag-omega-log}
For \(A=L\), \cref{plt:eq:lag-obs-L} gives, with \(\WrightOmega=\rev{(X+L)}\),
\[
  \log\WrightOmega
  =L-tL+t^2\Bigl(L-\frac{L^2}{2}\Bigr)
   +t^3\Bigl(-\frac{L^3}{3}+\frac32L^2-L\Bigr)+\FormalBigOAt{t^4}{t},
\]
which is exactly \(X-\WrightOmega\) read off from the defining relation
\(\WrightOmega+\log\WrightOmega=X\), and therefore a nontrivial check of both
\cref{plt:thm:lag-transport} and the coefficient table for \(\LamP{n}\).
\end{example}

\section{The finite coefficient formula and its exact index range}
\label{plt:sec:lag-coefficients}

\begin{theorem}[Coefficient formula in the polynomial--logarithmic basis]
\label{plt:thm:lag-coefficients}
Let \(A\in\PLpow\), \(F=X+A\), and write
\[
  A=\sum_{m\ge0}a_m(L)\,t^m,
  \qquad
  A^r=\sum_{m\ge0}A_{r,m}(L)\,t^m,
  \qquad
  \rev{F}=X+\sum_{N\ge0}b_N(L)\,t^N .
\]
Then for every \(N\in\NaturalNumbers\),
\begin{equation}
\label{plt:eq:lag-coefficients}
  b_N(L)
  =\sum_{r=1}^{N+1}\frac{(-1)^r}{r!}
    \left[\prod_{s=N-r+1}^{N-1}\bigl(\partial_L-s\bigr)\right]
    A_{r,\,N-r+1}(L),
\end{equation}
where a product whose lower bound exceeds its upper bound is empty and acts as
the identity; this occurs exactly at \(r=1\).
\end{theorem}

\begin{proof}
By \cref{plt:thm:lag-burmann} and the summability of the family, the
coefficient of \(t^N\) may be computed term by term.  For fixed \(r\ge1\) and
\(m\ge0\), the repeated-derivative formula gives
\[
  \DX^{\,r-1}\bigl(A_{r,m}(L)t^m\bigr)
  =t^{m+r-1}\,\shiftop{m}{r-1}A_{r,m}(L),
  \qquad
  \shiftop{m}{r-1}=\prod_{j=0}^{r-2}\bigl(\partial_L-m-j\bigr).
\]
This contributes to \(t^N\) exactly when \(m=N-r+1\), and then the parameters
\(m+j\) with \(0\le j\le r-2\) run over the integers
\(N-r+1,N-r+2,\ldots,N-1\), so \(\shiftop{N-r+1}{r-1}=
\prod_{s=N-r+1}^{N-1}(\partial_L-s)\).  The constraint \(m\ge0\) is
\(r\le N+1\).  For \(r=1\) the operator is \(\shiftop{N}{0}=1\), the empty
product.
\end{proof}

\begin{proposition}[The range \(r\le N+1\) is exactly right]
\label{plt:prop:lag-range}
Fix \(N\in\NaturalNumbers\) and let \(A\in\PLpow\) be nonzero.
\begin{enumerate}
\item No summand with \(r>N+1\) contributes to \(b_N\): indeed
      \(\ordt\bigl(\DX^{\,r-1}(A^r)\bigr)\ge r-1>N\).
\item The bound cannot be lowered.  For \(N\ge1\), the summand \(r=N+1\)
      contributes a nonzero polynomial to \cref{plt:eq:lag-coefficients} if and
      only if \(\degL a_0\ge1\); for \(N=0\) the only summand is \(r=1\) and
      \(b_0=-a_0\ne0\).
\end{enumerate}
In particular, for \(A=L\) every \(N\ge0\) attains the bound, so no uniform
improvement of the range is possible.
\end{proposition}

\begin{proof}
(1) is \cref{plt:lem:lag-order}.  For (2), the \(r=N+1\) summand of
\cref{plt:eq:lag-coefficients} is
\(\frac{(-1)^{N+1}}{(N+1)!}\bigl[\prod_{s=0}^{N-1}(\partial_L-s)\bigr]a_0^{N+1}\),
because \(A_{N+1,0}=a_0^{N+1}\).  Write \(q=a_0^{N+1}\); since \(\K[L]\) is a
domain, \(q\ne0\) and \(\degL q=(N+1)\degL a_0\).  If \(\degL a_0=0\) then
\(q\) is a nonzero constant and \(\partial_Lq=0\), so the product vanishes.  If
\(\degL a_0\ge1\) then \(\partial_Lq\ne0\), and for \(s\ge1\) the operator
\(\partial_L-s\) multiplies the leading coefficient of any nonzero polynomial by
\(-s\ne0\), hence is injective; so the whole product is nonzero.  For
\(A=L\) one has \(a_0=L\) and \(\degL a_0=1\).
\end{proof}

\begin{remark}
\label{plt:rmk:lag-truncation-transfer}
\Cref{plt:prop:lag-range} is the exact stopping rule for computation:
\(\Trunc{\rev{F}}{N+1}\), that is, the inverse through outer order \(t^{N}\)
inclusive, needs \(A^r\) only for \(r\le N+1\) and only through outer order
\(t^{N-r+1}\).  In the inclusive notation of report~6 this is
\(\bigl[\rev{F}\bigr]_{\le N}=\Trunc{\rev{F}}{N+1}\); the one-block difference
between the two conventions is a frequent source of off-by-one errors when
tables are transferred between reports.
\end{remark}

\section{The degree law}
\label{plt:sec:lag-degree}

Reversion is finite at every outer block, but the \emph{logarithmic} degrees
inside those blocks grow.  Recall the log-degree profile
\(\dprofile{F}(n)=\degL p_n\), with \(\dprofile{F}(n)=-\infty\) when
\(p_n=0\).

\begin{theorem}[Degree law for linearly bounded profiles]
\label{plt:thm:lag-degree-bound}
Let \(A\in\PLpow\), \(F=X+A\), \(\rev{F}=X+B\).  Suppose there are real
\(\alpha,\beta\ge0\) with
\begin{equation}
\label{plt:eq:lag-profile-hyp}
  \dprofile{A}(n)\le\alpha n+\beta
  \qquad\text{for all }n\in\NaturalNumbers .
\end{equation}
Then
\begin{equation}
\label{plt:eq:lag-profile-concl}
  \dprofile{B}(N)\ \le\ N\max\SetOf{\alpha,\beta}+\beta
  \qquad\text{for all }N\in\NaturalNumbers .
\end{equation}
More precisely, \(\dprofile{B}(0)\le\beta\) and, for \(N\ge1\),
\begin{equation}
\label{plt:eq:lag-profile-refined}
  \dprofile{B}(N)\ \le\
  \max\Bigl\{\,(N+1)\beta-1,\
    \max_{1\le r\le N}\bigl[\alpha(N-r+1)+r\beta\bigr]\Bigr\}.
\end{equation}
In particular the class of perturbations with linearly bounded log-degree
profile is stable under reversion.
\end{theorem}

\begin{proof}
Every operator \(\partial_L-s\) satisfies \(\degL\bigl((\partial_L-s)p\bigr)\le\degL p\),
with the strict bound \(\degL(\partial_Lp)\le\degL p-1\) when \(s=0\).  In
\cref{plt:eq:lag-coefficients} the operator attached to index \(r\) is
\(\prod_{s=N-r+1}^{N-1}(\partial_L-s)\); the value \(s=0\) occurs in that range
precisely when \(N-r+1\le0\le N-1\), that is (given \(1\le r\le N+1\)) when
\(N\ge1\) and \(r=N+1\).

Next, \(A_{r,m}=\sum_{m_1+\cdots+m_r=m}a_{m_1}\cdots a_{m_r}\), so
\cref{plt:eq:lag-profile-hyp} gives
\(\degL A_{r,m}\le\sum_i(\alpha m_i+\beta)=\alpha m+r\beta\).  Hence the
\(r\)-th summand of \cref{plt:eq:lag-coefficients} has degree at most
\(\alpha(N-r+1)+r\beta\), improved by one when \(r=N+1\) and \(N\ge1\), where
\(m=0\) and the bound is \((N+1)\beta-1\).  Taking the maximum over
\(1\le r\le N+1\) gives \cref{plt:eq:lag-profile-refined}, and
\(\dprofile{B}(0)=\degL(-a_0)\le\beta\).

For \cref{plt:eq:lag-profile-concl}, note
\(\alpha(N-r+1)+r\beta=\alpha N+\alpha+r(\beta-\alpha)\) is monotone in \(r\),
so on \(1\le r\le N\) it is at most
\(\max\SetOf{\alpha N+\beta,\ \alpha+N\beta}\le N\max\SetOf{\alpha,\beta}+\beta\),
while \((N+1)\beta-1<N\beta+\beta\le N\max\SetOf{\alpha,\beta}+\beta\).
\end{proof}

The bound is attained, with an exact leading coefficient, in the situation the
sources treat.

\begin{theorem}[Exact degree and leading coefficient]
\label{plt:thm:lag-degree-sharp}
Let \(A=\sum_{n\ge0}a_n(L)t^n\in\PLpow\) with
\[
  \degL a_0=d\ge1,
  \qquad
  a_d\DefinitionEquals[L^d]a_0\ne0,
  \qquad
  \degL a_n\le d-1\ \text{ for all }n\ge1 .
\]
Then \(b_0=-a_0\) has degree \(d\), and for every \(N\ge1\)
\begin{equation}
\label{plt:eq:lag-degree-sharp}
  \degL b_N=d(N+1)-1,
  \qquad
  \bigl[L^{\,d(N+1)-1}\bigr]b_N=\frac{d}{N}\,a_d^{\,N+1}.
\end{equation}
Thus \(\dprofile{B}(N)=dN+(d-1)\) for \(N\ge1\): the logarithmic degree of the
leading block of the perturbation becomes the \emph{slope} of the inverse's
log-degree profile.
\end{theorem}

\begin{proof}
Consider the summands of \cref{plt:eq:lag-coefficients} for \(N\ge1\).

\emph{The summand \(r=N+1\).}  It equals
\(\frac{(-1)^{N+1}}{(N+1)!}\,\partial_L(\partial_L-1)\cdots(\partial_L-N+1)\,a_0^{N+1}\).
The polynomial \(a_0^{N+1}\) has degree \(d(N+1)\) and leading coefficient
\(a_d^{N+1}\).  The factor \(\partial_L\) lowers the degree by one and multiplies
the leading coefficient by \(d(N+1)\); each factor \(\partial_L-s\) with
\(1\le s\le N-1\) preserves the degree and multiplies the leading coefficient
by \(-s\).  Hence this summand has degree \(d(N+1)-1\) and leading coefficient
\[
  \frac{(-1)^{N+1}}{(N+1)!}\;d(N+1)\,a_d^{\,N+1}\,(-1)^{N-1}(N-1)!
  =\frac{d(N+1)(N-1)!}{(N+1)!}\,a_d^{\,N+1}
  =\frac{d}{N}\,a_d^{\,N+1},
\]
which is nonzero because \(\K\) has characteristic zero.

\emph{The summands \(1\le r\le N\).}  Here \(m=N-r+1\ge1\), so in every product
\(a_{m_1}\cdots a_{m_r}\) occurring in \(A_{r,m}\) at least one index
\(m_i\) is \(\ge1\) and contributes degree at most \(d-1\), the remaining
\(r-1\) factors contributing at most \(d\) each.  Therefore
\(\degL A_{r,N-r+1}\le rd-1\le Nd-1<d(N+1)-1\), and the attached operator does
not raise degrees.

Adding the summands, the term of degree \(d(N+1)-1\) survives with the stated
coefficient.
\end{proof}

\begin{corollary}[Pure polynomial--logarithmic perturbations]
\label{plt:cor:lag-poly-perturbation}
Let \(P\in\K[L]\) with \(\degL P=d\ge1\) and leading coefficient \(a_d\), and
let \(F=X+P(L)\).  Then only the summand \(r=N+1\) contributes to \(b_N\), so
\begin{equation}
\label{plt:eq:lag-poly-block}
  \rev{F}=X-P(L)+\sum_{N\ge1}b_N(L)t^N,
  \qquad
  b_N(L)=\frac{(-1)^{N+1}}{(N+1)!}
   \left[\prod_{s=0}^{N-1}\bigl(\partial_L-s\bigr)\right]P(L)^{N+1},
\end{equation}
and \(\degL b_N=d(N+1)-1\) with
\(\bigl[L^{d(N+1)-1}\bigr]b_N=(d/N)a_d^{\,N+1}\) for \(N\ge1\).  For
\(P=aL\) this specialises to \(b_N=a^{N+1}\LamP{N}(L)\), so
\begin{equation}
\label{plt:eq:lag-scaled-lambert}
  \rev{(X+a\log X)}(Y)=Y+\sum_{N\ge0}a^{\,N+1}\LamP{N}(\log Y)\,Y^{-N},
  \qquad a\in\K .
\end{equation}
\end{corollary}

\begin{proof}
For \(A=P(L)\) one has \(A_{r,m}=P^r\delta_{m,0}\), so
\cref{plt:eq:lag-coefficients} retains only \(r=N+1\), giving
\cref{plt:eq:lag-poly-block}; the degree statement is
\cref{plt:thm:lag-degree-sharp} with \(a_n=0\) for \(n\ge1\).  For \(P=aL\),
\cref{plt:eq:lag-poly-block} reads
\(b_N=a^{N+1}\frac{(-1)^{N+1}}{(N+1)!}\bigl[\prod_{s=0}^{N-1}(\partial_L-s)\bigr]L^{N+1}\),
which is the defining formula of the canonical zero-based Lambert polynomial
\(\LamP{N}\); the case \(N=0\) gives \(\LamP{0}(L)=-L\), so the displayed sum
may start at \(N=0\).
\end{proof}

\begin{remark}[Expression swell, stated correctly]
\label{plt:rmk:lag-swell}
\Cref{plt:thm:lag-degree-sharp} is a statement about \emph{sizes of exact
symbolic objects}, not about an algorithm.  It says that computing the inverse
through outer order \(t^N\) requires manipulating polynomials in \(L\) of degree
about \(dN\); it says nothing about the number of arithmetic operations, and
nothing whatever about the analytic size of the discarded tail.  The three
big-\(O\) roles must be kept apart: the tail statement
\(\rev{F}-\Trunc{\rev{F}}{N}=\FormalBigOAt{t^N}{t}\) is a formal valuation
statement with no analytic content.
\end{remark}

\section{Worked extraction of the first coefficients}
\label{plt:sec:lag-worked}

Writing a dot for \(\partial_L\), \cref{plt:eq:lag-coefficients} gives the first
four universal coefficients of the inverse of \(F=X+\sum_{n\ge0}a_n(L)t^n\)
directly.  Grouping by the Lagrange index \(r\):
\begin{align}
  b_0={}&-a_0,\label{plt:eq:lag-b0}\\[1mm]
  b_1={}&-a_1
        +a_0\dot a_0,\label{plt:eq:lag-b1}\\[1mm]
  b_2={}&-a_2
        +\bigl[\dot a_0a_1+a_0\dot a_1-a_0a_1\bigr]
        +\Bigl[-a_0\dot a_0^{\,2}+\tfrac12a_0^2\dot a_0
               -\tfrac12a_0^2\ddot a_0\Bigr],\label{plt:eq:lag-b2}\\[1mm]
  b_3={}&-a_3
        +\bigl[\dot a_0a_2+a_0\dot a_2+a_1\dot a_1-2a_0a_2-a_1^2\bigr]
        \notag\\
       &+\Bigl[3a_0\dot a_0a_1-a_0^2a_1-\dot a_0^{\,2}a_1-a_0\ddot a_0a_1
               -2a_0\dot a_0\dot a_1+\tfrac32a_0^2\dot a_1
               -\tfrac12a_0^2\ddot a_1\Bigr]\notag\\
       &+\Bigl[a_0\dot a_0^{\,3}-\tfrac32a_0^2\dot a_0^{\,2}
               +\tfrac13a_0^3\dot a_0
               +\tfrac32a_0^2\dot a_0\ddot a_0
               -\tfrac12a_0^3\ddot a_0
               +\tfrac16a_0^3\dddot a_0\Bigr].\label{plt:eq:lag-b3}
\end{align}
In each line the successive bracketed groups are the contributions of
\(r=2,3,4\); the leading \(-a_N\) is the contribution of \(r=1\).

\begin{proof}[Derivation]
Apply \cref{plt:eq:lag-coefficients}.  For \(N=0\) only \(r=1\) occurs and the
operator is empty.  For \(N=1\), \(r=1\) gives \(-A_{1,1}=-a_1\) and \(r=2\)
gives \(\tfrac12\partial_L(a_0^2)=a_0\dot a_0\).  For \(N=2\): \(r=1\) gives
\(-a_2\); \(r=2\) gives
\(\tfrac12(\partial_L-1)A_{2,1}=\tfrac12(\partial_L-1)(2a_0a_1)\), which is the
first bracket; \(r=3\) gives
\(-\tfrac16\partial_L(\partial_L-1)a_0^3\), and with
\(\partial_L(\partial_L-1)q=\ddot q-\dot q\) applied to \(q=a_0^3\),
\(\dot q=3a_0^2\dot a_0\) and
\(\ddot q=6a_0\dot a_0^{\,2}+3a_0^2\ddot a_0\), this is the second bracket.
For \(N=3\): \(r=1\) gives \(-a_3\); \(r=2\) gives
\(\tfrac12(\partial_L-2)A_{2,2}\) with \(A_{2,2}=2a_0a_2+a_1^2\); \(r=3\) gives
\(-\tfrac16(\partial_L-1)(\partial_L-2)A_{3,1}\) with \(A_{3,1}=3a_0^2a_1\) and
\((\partial_L-1)(\partial_L-2)q=\ddot q-3\dot q+2q\); \(r=4\) gives
\(\tfrac1{24}\partial_L(\partial_L-1)(\partial_L-2)a_0^4\) with
\(\partial_L(\partial_L-1)(\partial_L-2)q=\dddot q-3\ddot q+2\dot q\), and
\(\dot q=4a_0^3\dot a_0\),
\(\ddot q=12a_0^2\dot a_0^{\,2}+4a_0^3\ddot a_0\),
\(\dddot q=24a_0\dot a_0^{\,3}+36a_0^2\dot a_0\ddot a_0+4a_0^3\dddot a_0\).
Collecting gives the displayed brackets.
\end{proof}

% ed.: S4 gives only b_0, b_1, b_2 in this universal form; the b_3 line
% plt:eq:lag-b3 is derived here and verified against the two independent worked
% examples plt:ex:lag-mixed and plt:ex:lag-abc.
\begin{example}[A fully mixed perturbation]
\label{plt:ex:lag-mixed}
Let
\[
  F(X)=X+2\log X+1+\frac{3\log X-1}{X}+\frac{(\log X)^2+2}{X^2},
\]
so \(a_0=2L+1\), \(a_1=3L-1\), \(a_2=L^2+2\), \(a_n=0\) for \(n\ge3\).  Then
\cref{plt:eq:lag-b0}--\cref{plt:eq:lag-b3} give
\begin{align}
  \rev{F}(Y)={}&Y-2M-1+\frac{M+3}{Y}
    +\frac{-3M^2+7M-3}{Y^2}\notag\\
   &+\frac{-\tfrac{32}{3}M^3+25M^2-6M-\tfrac{59}{6}}{Y^3}
    +\FormalBigOAt{t^4}{t},
  \qquad M=\log Y.
  \label{plt:eq:lag-mixed-inverse}
\end{align}
Grouped by Lagrange index, the contributions to the block \(t^3\) are
\(-4M^3-5M^2+9M-4\) from \(r=2\), \(-12M^3+46M^2-11M-\tfrac{17}{2}\) from
\(r=3\), and \(\tfrac{16}{3}M^3-16M^2-4M+\tfrac83\) from \(r=4\); nothing comes
from \(r=1\).  A summand that is ``later'' in the Lagrange index therefore
still reaches the same outer block: the stopping rule is the filtration bound
of \cref{plt:prop:lag-range}, never the number of displayed summands.
\end{example}

\begin{example}[A logarithmic--rational family]
\label{plt:ex:lag-abc}
For \(F(X)=X+a\log X+bX^{-1}+cX^{-2}\), that is \(a_0=aL\), \(a_1=b\),
\(a_2=c\),
\begin{align}
  \rev{F}(Y)={}&Y-aM+\frac{a^2M-b}{Y}
   +\frac{a^3\bigl(\tfrac12M^2-M\bigr)-abM+ab-c}{Y^2}\notag\\
   &+\frac{a^4\bigl(\tfrac13M^3-\tfrac32M^2+M\bigr)
        +a^2b\bigl(-M^2+3M-1\bigr)-2acM+ac-b^2}{Y^3}
   +\FormalBigOAt{t^4}{t}.
   \label{plt:eq:lag-abc-inverse}
\end{align}
Setting \(b=c=0\) recovers \cref{plt:eq:lag-scaled-lambert}; setting
\(a=c=0\) recovers \cref{plt:ex:lag-catalan} with \(\kappa=b\).  The parameter
\(c\), which enters \(F\) at outer
order \(t^2\), first perturbs the inverse at the same order, and the product
\(b^2\) first appears at \(t^3\), where two rational perturbations can meet.
\end{example}

\begin{example}[An exact check: Catalan numbers]
\label{plt:ex:lag-catalan}
For \(F(X)=X+\kappa X^{-1}\), \cref{plt:eq:lag-burmann-pl} with \(A=\kappa t\)
and
\(\DX^{\,r-1}(t^r)=t^{2r-1}\prod_{j=0}^{r-2}(-r-j)
=(-1)^{r-1}\frac{(2r-2)!}{(r-1)!}t^{2r-1}\) gives
\[
  \rev{F}(Y)=Y-\sum_{r\ge1}\frac{(2r-2)!}{r!\,(r-1)!}\,\kappa^r\,Y^{-(2r-1)}
  =Y-\frac \kappa Y-\frac{\kappa^2}{Y^3}-\frac{2\kappa^3}{Y^5}
   -\frac{5\kappa^4}{Y^7}-\cdots,
\]
the coefficients \((2r-2)!/\bigl(r!\,(r-1)!\bigr)\) being the Catalan numbers.
This matches the branch \(\tfrac12\bigl(Y+\sqrt{Y^2-4\kappa}\bigr)\) of
\(Z+\kappa Z^{-1}=Y\) that is asymptotic to \(Y\), and simultaneously checks
the signs, the derivative order \(r-1\), and the branch selection implicit in
the near-identity normalisation.
\end{example}

% ed.: this proposition restated \cref{plt:cor:rev-finite-certificate}, and its
% ed.: proof re-derived \cref{plt:prop:rev-residual-order} along the way.  The
% ed.: statements are identical --- the hypothesis \(F=X+A\) with
% ed.: \(A\in\PLpow\) is membership in \(\Gnear\), and the two candidate
% ed.: truncations \(G_N\) are the same object.  The restatement and its proof
% ed.: are deleted in favour of the pointer below; every citation of the
% ed.: deleted label in this part now cites the reversion part.
The coefficients produced by \cref{plt:thm:lag-coefficients} are certified by
a finite computation, and the certificate is the one already established in
the reversion part: for \(F=X+A\in\Gnear\) and
\(G_N=X+\sum_{j=0}^{N}g_j(L)t^j\), the three conditions
\(G_N=\Trunc{\rev{F}}{N+1}\), \(\RevErr{F}{G_N}=\FormalBigOAt{t^{N+1}}{t}\)
and \(G_N\compose F-X=\FormalBigOAt{t^{N+1}}{t}\) are equivalent
(\cref{plt:cor:rev-finite-certificate}); moreover the order of either residual
equals \(\ordt(G_N-\rev{F})\) exactly, with the same leading block
(\cref{plt:prop:rev-residual-order}).  Nothing in that argument uses the
Lagrange--B\"urmann representation, so it certifies coefficients obtained by
any of the routes of this part.

\begin{example}
\label{plt:ex:lag-certificate-mixed}
Substituting \cref{plt:eq:lag-mixed-inverse} back into the \(F\) of
\cref{plt:ex:lag-mixed}, using the exact generator substitutions
\(t\compose G=t/(1+U)\) and \(L\compose G=L+\log(1+U)\) with \(U=t(G-Y)\),
gives
\[
  F\compose G-Y
  =\Bigl(28M^4-\tfrac{211}{3}M^3-15M^2+91M-\tfrac{59}{3}\Bigr)t^4
   +\FormalBigOAt{t^5}{t}.
\]
The vanishing through \(t^3\) certifies every coefficient claimed in
\cref{plt:eq:lag-mixed-inverse}; the nonzero \(t^4\) block is expected, being
the first block not represented in \(G\).  Computing one residual block beyond
the target is the cheapest available check, and it detects exactly the errors
that a second run of the same code would not.
\end{example}

\begin{summarybox}[title={What the formula does and does not say}]
For \(F=X+A\) with \(A\in\PLpow\),
\(\rev{F}=X+\sum_{r\ge1}\frac{(-1)^r}{r!}\DX^{\,r-1}(A^r)\), and the block
\(t^N\) is a finite sum over \(1\le r\le N+1\).  Summability holds precisely
because \(\ordt(A)\ge0\), and \emph{only} then
(\cref{plt:prop:lag-summability-sharp}).  The same construction transports any
observable \(\varphi\in\PLlau\) along the inverse without building it
(\cref{plt:thm:lag-transport}).  Logarithmic degrees grow linearly, with slope
\(\degL a_0\) (\cref{plt:thm:lag-degree-sharp}).  None of this asserts
convergence, an analytic inverse branch, or any bound uniform in \(L\).
\end{summarybox}


\chapter{Reversion beyond the identity scale}
\label{plt:sec:dominant-balance}

\Cref{plt:sec:lagrange-burmann} inverts \(F=X+A\) with \(\ordt(A)\ge0\).  Two
things can go wrong outside that class: the leading exponent of \(F\) need not
be \(1\), and the leading coefficient block need not be a scalar.  This chapter
shows that these are the \emph{only} two obstructions, proves a criterion that
decides in advance which enlargement of \(\PLlau\) the inverse will require,
turns dominant balance into an algorithm with a proof attached, and exhibits
the way a careless scale change silently corrupts an otherwise correct
near-identity computation.

Throughout this chapter, composition means composition of maps: whenever we
write \(F\compose G\) for germs at \(+\infty\), associativity and the
functorial identity
\begin{equation}
\label{plt:eq:lag-functorial}
  \rev{(F\compose G)}=\rev{G}\compose\rev{F}
\end{equation}
are the ordinary group-theoretic facts, valid in any group of invertible
maps, and are used without further comment for genuine germs.  For formal
series they hold in the near-identity group of any reversion frame by
\cref{plt:lem:lag-substitution}(5), and more generally whenever all displayed
composites are admissible in the chosen support class.

\section{Monomial-leading maps and the closure criterion}
\label{plt:sec:lag-criterion}

\begin{definition}[Monomial-leading map]
\label{plt:def:lag-monomial-leading}
Let \(\Gamma\subseteq\RationalNumbers\) be an exponent group.  A series
\(G\) is \emph{monomial-leading of exponent \(\mu\)} if
\begin{equation}
\label{plt:eq:lag-monomial-leading}
  G=dX^{\mu}\bigl(1+V\bigr),
  \qquad d\in\K^\times,\quad \mu\in\Gamma_{>0},\quad \ordt(V)>0,
\end{equation}
with \(\log d\in\K\).  This is the class for which the generator substitutions
\begin{equation}
\label{plt:eq:lag-generator-substitution}
  t\compose G=d^{-1}t^{\mu}(1+V)^{-1},
  \qquad
  L\compose G=\mu L+\log d+\log(1+V)
\end{equation}
stay inside the ambient coefficient ring, and hence the class for which
composition with an outer polynomial--logarithmic series is defined.
\end{definition}

% ed.: the requirement log d in K is stated in the catalogue and is omitted in
% S1, S2 and S6, which write log c freely over an unspecified coefficient
% field.  Over K = Q it fails already for d = 2; see
% plt:warn:lag-scale-trap.
\begin{remark}
\label{plt:rmk:lag-nonscalar-leading}
The scalar hypothesis \(d\in\K^\times\) cannot be relaxed to a polynomial
leading block.  If \(G=X^\mu q(L)(1+V)\) with \(q\in\K[L]\) of degree
\(e\ge1\) and leading coefficient \(q_e\), then
\(L\compose G=\mu L+\log q(L)+\log(1+V)\) and
\(\log q(L)=e\log L+\log q_e+\BigOAt{L^{-1}}{X\to+\infty}\), so a new generator
\(\log L\) appears.  That generator is genuinely new: if \(\log L=p(L)+R\) with
\(p\in\K[L]\) and \(R\to0\) as \(X\to+\infty\), then \(\log L-p(L)\to0\), which
is impossible, since for constant \(p\) the difference tends to \(+\infty\) and
for \(\degL p\ge1\) it tends to \(\pm\infty\).  This is the same argument used
in \cref{plt:prop:lag-loglog-forced}.
\end{remark}

\begin{theorem}[Closure criterion for reversion in the one-logarithm ring]
\label{plt:thm:lag-rigidity}
Let \(\Gamma\subseteq\RationalNumbers\) be an exponent group containing
\(\IntegerNumbers\), let \(\mathfrak R\) be the corresponding Puiseux--log ring
(so \(\mathfrak R=\PLlau\) when \(\Gamma=\IntegerNumbers\), and
\(\mathfrak R=\K[L]((t^{1/s}))\) when \(\Gamma=s^{-1}\IntegerNumbers\)), and let
\begin{equation}
\label{plt:eq:lag-F-monomial}
  F=cX^{\rho}L^{\sigma}\bigl(1+U\bigr),
  \qquad
  c\in\K^\times,\quad \rho\in\Gamma_{>0},\quad
  \sigma\in\NaturalNumbers,\quad \ordt(U)>0 .
\end{equation}
Then \(F\) admits a monomial-leading right inverse \(G\) in the sense of
\cref{plt:def:lag-monomial-leading}, over \(\mathfrak R\) or over any extension
of \(\mathfrak R\) obtained by enlarging the coefficient field alone, if and
only if
\begin{equation}
\label{plt:eq:lag-criterion}
  \rho^{-1}\in\Gamma
  \qquad\text{and}\qquad
  \sigma=0 .
\end{equation}
When \cref{plt:eq:lag-criterion} holds, the exponent of \(G\) is
\(\mu=\rho^{-1}\) and its leading scalar is any \(d\) with \(cd^{\rho}=1\), so
the construction requires a field \(\K'\supseteq\K\) containing such a \(d\)
together with \(\log d\).  In particular:
\begin{enumerate}
\item in \(\PLlau\) itself \(\bigl(\Gamma=\IntegerNumbers\bigr)\), the criterion
      is \(\rho=1\) and \(\sigma=0\), that is,
      \(c^{-1}F\in\Gnear(\PLpow)\);
\item enlarging the exponent group to \(\rho^{-1}\IntegerNumbers\) removes the
      exponent obstruction but never the obstruction \(\sigma\ge1\).
\end{enumerate}
\end{theorem}

\begin{proof}
Write \(F(Y)=\sum_{n\ge-\rho}p_n(\log Y)Y^{-n}\) with
\(p_{-\rho}(\lambda)=c\lambda^{\sigma}\) and \(\ordt(F-cX^\rho L^\sigma)>-\rho\),
and let \(G=dX^{\mu}(1+V)\) be monomial-leading.  By
\cref{plt:eq:lag-generator-substitution} and the monomial-leading composition
formula,
\[
  F\compose G
  =\sum_{n\ge-\rho}p_n\bigl(\mu L+\log d+\log(1+V)\bigr)
     d^{-n}t^{\mu n}(1+V)^{-n}.
\]
Since \(\ordt(V)>0\) and \(\ordt\log(1+V)>0\), the term \(n=-\rho\) contributes
\[
  c\,d^{\rho}\,\bigl(\mu L+\log d\bigr)^{\sigma}\,t^{-\mu\rho}
  +\bigl(\text{terms of }\ordt>-\mu\rho\bigr),
\]
and every term with \(n>-\rho\) has \(\ordt\ge-\mu n>-\mu\rho\).  Hence
\(\ordt(F\compose G)=-\mu\rho\) and the leading block of \(F\compose G\) is
\(cd^{\rho}(\mu L+\log d)^{\sigma}\).

Requiring \(F\compose G=X=t^{-1}\) forces \(\mu\rho=1\), so \(\mu=\rho^{-1}\)
must lie in \(\Gamma\), and forces
\begin{equation}
\label{plt:eq:lag-leading-block-equation}
  c\,d^{\rho}\,\bigl(\mu L+\log d\bigr)^{\sigma}=1 .
\end{equation}
If \(\sigma\ge1\), the left side of \cref{plt:eq:lag-leading-block-equation} is
a polynomial in \(L\) of degree \(\sigma\ge1\) with leading coefficient
\(cd^{\rho}\mu^{\sigma}\ne0\), so it cannot be the constant \(1\); hence
\(\sigma=0\), and then \cref{plt:eq:lag-leading-block-equation} reads
\(cd^{\rho}=1\).

Conversely, suppose \(\rho^{-1}\in\Gamma\) and \(\sigma=0\).  Choose \(d\) with
\(cd^{\rho}=1\) in a finite extension \(\K'\supseteq\K\) also containing
\(\log d\), and set \(F_0=cX^{\rho}\), \(\rev{F_0}=dX^{1/\rho}\).  Then
\(\Phi\DefinitionEquals F\compose\rev{F_0}\) satisfies
\(\Phi=cd^{\rho}X\bigl(1+U\compose\rev{F_0}\bigr)=X\bigl(1+\widetilde U\bigr)\)
with \(\ordt(\widetilde U)>0\), because
\(\ordt\bigl(U\compose\rev{F_0}\bigr)=\rho^{-1}\ordt(U)>0\).  Hence
\(\Phi\in\Gnear\) of the Puiseux--log frame
\(\K'[L]((t^{1/s}))\), \(s\) a common denominator of \(\rho^{-1}\) and the
exponents of \(\widetilde U\), and \cref{plt:thm:lag-burmann} produces
\(\rev{\Phi}\).  By \cref{plt:eq:lag-functorial},
\(\rev{F}=\rev{F_0}\compose\rev{\Phi}\), which is monomial-leading of exponent
\(\rho^{-1}\).
\end{proof}

\begin{proposition}[A leading logarithm forces a nested logarithm]
\label{plt:prop:lag-loglog-forced}
Let \(F=cX^\rho L^\sigma(1+U)\) as in \cref{plt:eq:lag-F-monomial} with
\(\sigma\ge1\), and let \(G\) be a germ at \(+\infty\) with \(F(G(Y))=Y\) for
all large \(Y\).  Put \(S=\log Y\) and \(u=\log G(Y)\).  Then
\begin{equation}
\label{plt:eq:lag-loglog}
  u=\frac1\rho\Bigl(S-\sigma\log S+\sigma\log\rho-\log c\Bigr)
    +\BigOAt{\tfrac{\log S}{S}}{S\to+\infty},
\end{equation}
so \(\log G\) contains the term \(-(\sigma/\rho)\log\log Y\) with nonzero
coefficient.  Consequently \(G\) lies in no Puiseux--log ring
\(\K'[S]((Y^{-1/s}))\) with \(\K'\supseteq\K\) and \(s\in\PositiveIntegers\):
a second logarithm is unavoidable.
\end{proposition}

\begin{proof}
Taking logarithms in \(F(G(Y))=Y\) and writing \(u=\log G(Y)\),
\begin{equation}
\label{plt:eq:lag-log-equation}
  S=\log c+\rho u+\sigma\log u+\log\bigl(1+U(G(Y))\bigr).
\end{equation}
Since \(\ordt(U)>0\) there is \(\delta>0\) with
\(U(X)=\BigOAt{X^{-\delta}}{X\to+\infty}\), so the last term of
\cref{plt:eq:lag-log-equation} is
\(\BigOAt{G(Y)^{-\delta}}{Y\to+\infty}\); since \(G(Y)\to+\infty\) like a
positive power of \(Y\), it is smaller than every power of \(1/S\).  From
\cref{plt:eq:lag-log-equation}, \(\rho u\sim S\), hence
\(\log u=\log S-\log\rho+\LittleOAt{1}{S\to+\infty}\); substituting back,
\(\rho u=S-\log c-\sigma\log S+\sigma\log\rho+\LittleOAt{1}{S\to+\infty}\),
which is \cref{plt:eq:lag-loglog}.

Suppose \(G\in\K'[S]((Y^{-1/s}))\).  Then
\(G=q(S)Y^{-\lambda}\bigl(1+\LittleOAt{1}{Y\to+\infty}\bigr)\) for some
\(\lambda\in s^{-1}\IntegerNumbers\) and some nonzero \(q\in\K'[S]\), so
\[
  u=\log G=-\lambda S+\log q(S)+\LittleOAt{1}{Y\to+\infty}
   =-\lambda S+(\deg q)\log S+\log(\mathrm{lc}\,q)
    +\LittleOAt{1}{Y\to+\infty}.
\]
Comparing with \cref{plt:eq:lag-loglog}, the coefficients of \(S\) give
\(\lambda=-1/\rho\) and those of \(\log S\) give \(\deg q=-\sigma/\rho<0\),
which is impossible for a polynomial.
\end{proof}

\begin{corollary}[Which enlargement is needed]
\label{plt:cor:lag-enlargement}
For \(F\) as in \cref{plt:eq:lag-F-monomial}:
\begin{center}
\begin{tabular}{@{}lll@{}}
\hline
leading data & obstruction & smallest ambient for \(\rev{F}\)\\
\hline
\(\rho=1,\ \sigma=0\) & none & \(\PLlau\) over \(\K(c^{-1},\log c)\)\\
\(\rho\in\RationalNumbers_{>0},\ \sigma=0\) & exponent &
   \(\K'[L]((t^{1/s}))\), \(s\) a denominator of \(\rho^{-1}\)\\
\(\rho\notin\RationalNumbers,\ \sigma=0\) & exponent &
   Hahn power-log with \(\rho^{-1}\IntegerNumbers\subseteq\Gamma\)\\
\(\sigma\ge1\) & logarithm & nested-logarithmic; \(\log\log X\) required\\
\hline
\end{tabular}
\end{center}
\end{corollary}

\begin{proof}
Rows one to three are \cref{plt:thm:lag-rigidity} together with the explicit
construction in its converse half; row four is
\cref{plt:prop:lag-loglog-forced} together with
\cref{plt:thm:lag-log-normal-form} below, which produces the inverse in the
nested-logarithmic frame.
\end{proof}

\section{Dominant balance as a theorem}
\label{plt:sec:lag-reduction}

The workflow ``guess the leading behaviour, normalise, iterate'' becomes a
theorem once one names the object that does the normalising.

\begin{theorem}[Reduction to the near-identity case]
\label{plt:thm:lag-reduction}
Let \(\mathfrak G\) be a group of invertible germs at \(+\infty\) (or of formal
maps in a fixed support class, with all displayed composites admissible), let
\(F\in\mathfrak G\), and let \(F_0\in\mathfrak G\) be any \emph{leading model},
that is, any element whose inverse is known.  Put
\begin{equation}
\label{plt:eq:lag-normalisation}
  \Phi\DefinitionEquals F\compose\rev{F_0}.
\end{equation}
Then \(F=\Phi\compose F_0\) and
\begin{equation}
\label{plt:eq:lag-reduction}
  \rev{F}=\rev{F_0}\compose\rev{\Phi}.
\end{equation}
If moreover \(\Phi=X+A\) with \(A\) in the nonnegative part of some reversion
frame, then \(\rev{\Phi}\) is given by \cref{plt:eq:lag-burmann}, and
\begin{equation}
\label{plt:eq:lag-reduction-explicit}
  \rev{F}
  =\rev{F_0}\compose\Bigl(X+\sum_{r\ge1}\frac{(-1)^r}{r!}
     \partial^{\,r-1}\bigl(A^r\bigr)\Bigr).
\end{equation}
The dual normalisation \(\Psi\DefinitionEquals\rev{F_0}\compose F\) gives
\(F=F_0\compose\Psi\) and \(\rev{F}=\rev{\Psi}\compose\rev{F_0}\).
\end{theorem}

\begin{proof}
\(\Phi\compose F_0=F\compose\rev{F_0}\compose F_0=F\).  Applying
\cref{plt:eq:lag-functorial} to \(F=\Phi\compose F_0\) gives
\(\rev{F}=\rev{F_0}\compose\rev{\Phi}\).  The last two statements are the same
computation with the factors exchanged, and
\cref{plt:eq:lag-reduction-explicit} is \cref{plt:thm:lag-burmann} substituted
into \cref{plt:eq:lag-reduction}.
\end{proof}

\Cref{plt:thm:lag-reduction} is content-free until a leading model is produced.
For a monomial-leading \(F\) the model is supplied by one canonical change of
scale, and this is where dominant balance actually happens.

\begin{theorem}[Logarithmic normal form of a power--log leading term]
\label{plt:thm:lag-log-normal-form}
Let \(c\in\K^\times\) with \(\log c\in\K\), let \(\rho\in\RationalNumbers_{>0}\)
with \(\log\rho\in\K\), let \(\sigma\in\NaturalNumbers\), and set
\(F_0(X)=cX^{\rho}(\log X)^{\sigma}\).  Write \(E(X)=\EulerE^{X}\) and
\(m_\rho(X)=\rho X\).  Then
\begin{equation}
\label{plt:eq:lag-log-conjugate}
  G\DefinitionEquals \rev{m_\rho}\compose\log\compose F_0\compose E
  \ =\ u+\frac{\sigma}{\rho}\log u+\frac{\log c}{\rho}
\end{equation}
as a map in the large variable \(u\); in particular \(G\) lies in the
near-identity group of the reversion frame
\(\K[\log u][[u^{-1}]]\) with perturbation
\(A(u)=(\sigma/\rho)\log u+(\log c)/\rho\), and
\begin{equation}
\label{plt:eq:lag-log-inverse}
  \rev{F_0}=E\compose\rev{G}\compose m_{1/\rho}\compose\log .
\end{equation}
Consequently, with \(S=\log Y\) and \(\Lambda=\log S\),
\begin{align}
  \log\rev{F_0}(Y)
  ={}&\frac1\rho\Bigl[\,S-\sigma\Lambda+\sigma\log\rho-\log c\notag\\
     &\qquad+\frac{\sigma^2\Lambda+\sigma\log c-\sigma^2\log\rho}{S}
      +\FormalBigOAt{S^{-2}}{S^{-1}}\Bigr],
  \label{plt:eq:lag-u-expansion}\\[1mm]
  \rev{F_0}(Y)
  ={}&c^{-1/\rho}\rho^{\sigma/\rho}\,Y^{1/\rho}S^{-\sigma/\rho}
      \Bigl[1+\frac{\sigma^2\Lambda+\sigma\log c-\sigma^2\log\rho}{\rho S}
      +\FormalBigOAt{S^{-2}}{S^{-1}}\Bigr].
  \label{plt:eq:lag-X-expansion}
\end{align}
\end{theorem}

\begin{proof}
For \(X=\EulerE^{u}\), \(\log F_0(\EulerE^{u})=\log c+\rho u+\sigma\log u\).
Post-composing with \(\rev{m_\rho}=m_{1/\rho}\) divides by \(\rho\) and gives
\cref{plt:eq:lag-log-conjugate}.  Since \(\log u\in\K[\log u][[u^{-1}]]\) has
zero outer order in the variable \(u\), the perturbation \(A\) has
\(\nu(A)=0\), so \(G\in\Gnear\) of that frame.  Reading
\cref{plt:eq:lag-log-conjugate} as
\(G=m_{1/\rho}\compose\log\compose F_0\compose E\) and inverting with
\cref{plt:eq:lag-functorial}, using \(\rev{E}=\log\) and
\(\rev{(\log)}=E\),
\[
  \rev{G}=\rev{E}\compose\rev{F_0}\compose\rev{(\log)}\compose\rev{(m_{1/\rho})}
        =\log\compose\rev{F_0}\compose E\compose m_{\rho},
\]
which rearranges to \cref{plt:eq:lag-log-inverse}.

For the expansions, apply \cref{plt:eq:lag-b0} and \cref{plt:eq:lag-b1} to
\(A(u)=a_0(\log u)\) with \(a_0(\lambda)=(\sigma/\rho)\lambda+(\log c)/\rho\)
and \(a_n=0\) for \(n\ge1\):
\[
  b_0=-a_0,\qquad
  b_1=a_0\dot a_0=\frac{\sigma}{\rho}\Bigl(\frac{\sigma}{\rho}\lambda
      +\frac{\log c}{\rho}\Bigr)
     =\frac{\sigma}{\rho^2}\bigl(\sigma\lambda+\log c\bigr).
\]
Hence, writing \(v\) for the large variable of \(\rev{G}\) and
\(\lambda=\log v\),
\[
  \rev{G}(v)=v-\frac{\sigma\lambda+\log c}{\rho}
   +\frac{\sigma(\sigma\lambda+\log c)}{\rho^2 v}
   +\FormalBigOAt{v^{-2}}{v^{-1}}.
\]
By \cref{plt:eq:lag-log-inverse}, \(\log\rev{F_0}(Y)=\rev{G}(S/\rho)\).
Substituting \(v=S/\rho\), so \(\lambda=\Lambda-\log\rho\) and
\(v^{-1}=\rho/S\),
\[
  \log\rev{F_0}(Y)
  =\frac{S}{\rho}-\frac{\sigma(\Lambda-\log\rho)+\log c}{\rho}
   +\frac{\sigma\bigl(\sigma(\Lambda-\log\rho)+\log c\bigr)}{\rho S}
   +\FormalBigOAt{S^{-2}}{S^{-1}},
\]
which is \cref{plt:eq:lag-u-expansion}.  Exponentiating and using
\(\EulerE^{S/\rho}=Y^{1/\rho}\),
\(\EulerE^{-\sigma\Lambda/\rho}=S^{-\sigma/\rho}\), and
\(\EulerE^{(\sigma\log\rho-\log c)/\rho}=\rho^{\sigma/\rho}c^{-1/\rho}\), gives
\cref{plt:eq:lag-X-expansion}, the exponential of the \(S^{-1}\) block being
expanded as \(1+w+\FormalBigOAt{S^{-2}}{S^{-1}}\).
\end{proof}

% ed.: S6 derives plt:eq:lag-u-expansion by inserting an undetermined ansatz
% u = rho^{-1}(S - sigma log S + C + (a log S + b)/S + ...) and matching three
% constants by hand.  The derivation above obtains the same coefficients from
% the near-identity coefficient formula, so the pattern continues to all orders
% with no further ansatz, and the branch and field hypotheses (log c, log rho
% in K) become visible.
\begin{remark}
\label{plt:rmk:lag-log-frame-scale}
\Cref{plt:eq:lag-log-conjugate} explains the ``one level down'' phenomenon:
after conjugation by \(\EulerE\), the expansion parameter is \(1/\log Y\) and
the coefficients are polynomials in \(\log\log Y\).  The polynomial--logarithmic
calculus is applied in the frame of the variable \(u=\log X\), not of \(X\).
Note also that \cref{plt:eq:lag-log-conjugate} is exact only for
\(U=0\); for \(F=F_0\cdot(1+U)\) with \(\ordt(U)>0\) the conjugated map carries
an extra term \(\rho^{-1}\log(1+U(\EulerE^{u}))\), which is smaller than every
power of \(u^{-1}\) and therefore invisible in the frame
\(\K[\log u][[u^{-1}]]\).  That term is not negligible after exponentiating
back, and it is precisely what the second stage
\cref{plt:eq:lag-normalisation} of \cref{plt:thm:lag-reduction} handles: with
\(\Phi=F\compose\rev{F_0}=X(1+U\compose\rev{F_0})\), the two scales
\(1/\log Y\) and \(Y^{-1/\rho}\) are separated rather than mixed.
\end{remark}

\begin{algorithm}[Reversion with a nonidentity leading term]
\label{plt:alg:lag-dominant-balance}
Input: a germ \(F\) at \(+\infty\) presented as
\(F=cX^{\rho}L^{\sigma}(1+U)\) with \(c\in\K^\times\),
\(\rho\in\RationalNumbers_{>0}\), \(\sigma\in\NaturalNumbers\),
\(\ordt(U)>0\); a target outer order \(N\).
\begin{enumerate}
\item \textbf{Extend the coefficient field.}  Replace \(\K\) by a field
      containing \(\log c\), \(\log\rho\) and \(c^{-1/\rho}\).  Record the
      branch of every root and logarithm chosen.
\item \textbf{Fix the ambient.}  If \(\sigma=0\) and \(\rho=1\), the ambient is
      \(\PLlau\).  If \(\sigma=0\) and \(\rho=p/q\) in lowest terms, the ambient
      is \(\K[L]((t^{1/p}))\).  If \(\sigma\ge1\), the ambient is the
      nested-logarithmic frame of \cref{plt:thm:lag-log-normal-form}.
      (\Cref{plt:thm:lag-rigidity} and \cref{plt:cor:lag-enlargement}.)
\item \textbf{Invert the leading model.}  Set \(F_0=cX^{\rho}L^{\sigma}\).  If
      \(\sigma=0\), \(\rev{F_0}=c^{-1/\rho}X^{1/\rho}\).  If \(\sigma\ge1\),
      obtain \(\rev{F_0}\) from
      \cref{plt:eq:lag-log-inverse}: revert
      \(G(u)=u+(\sigma/\rho)\log u+(\log c)/\rho\) by
      \cref{plt:thm:lag-burmann} in the \(u\)-frame to the required order, then
      evaluate at \(u=S/\rho\) and exponentiate.
\item \textbf{Normalise.}  Form \(\Phi=F\compose\rev{F_0}\).  By construction
      \(\Phi=X\bigl(1+U\compose\rev{F_0}\bigr)\), which is near-identity in the
      ambient of step~2.
\item \textbf{Revert the near-identity factor.}  Apply
      \cref{plt:thm:lag-burmann} or the coefficient formula
      \cref{plt:eq:lag-coefficients} to \(\Phi\), through the order fixed by
      \cref{plt:prop:lag-range}.
\item \textbf{Recompose.}  Return \(\rev{F}=\rev{F_0}\compose\rev{\Phi}\), as in
      \cref{plt:eq:lag-reduction}.  \emph{Both} generators must be substituted:
      \(t\) and \(L\) change together, by
      \cref{plt:eq:lag-generator-substitution}.
\item \textbf{Certify.}  Verify \(F\compose\rev{F}-X=\FormalBigOAt{t^{N+1}}{t}\)
      and, independently, \(\rev{F}\compose F-X=\FormalBigOAt{t^{N+1}}{t}\), as
      in \cref{plt:cor:rev-finite-certificate}.  The two residuals detect different
      substitution-order errors.
\end{enumerate}
\end{algorithm}

\section{Worked examples}
\label{plt:sec:lag-dominant-examples}

\begin{example}[Leading behaviour \(X+c\log X\): inside the theory]
\label{plt:ex:lag-affine-log}
The map \(F(X)=X+a\log X+b\) with \(a\ne0\) has \(A=aL+b\in\PLpow\), so
\(\ordt(A)=0\) and no enlargement is needed, even though
\(A\to+\infty\).  By \cref{plt:cor:lag-poly-perturbation} with
\(P=aL+b\), the inverse is
\begin{align}
  \rev{F}(Y)={}&Y-aM-b
    +\frac{a(aM+b)}{Y}
    +\frac{a(aM+b)(aM+b-2a)}{2Y^2}\notag\\
   &+\frac{a(aM+b)}{6Y^3}
     \bigl(2a^2M^2-9a^2M+4abM+6a^2-9ab+2b^2\bigr)
   +\FormalBigOAt{t^4}{t},
   \label{plt:eq:lag-affine-log-inverse}
\end{align}
with \(M=\log Y\), and \(\degL b_N=2(N+1)-1\) when \(\deg P=1\)\dots\ no: with
\(d=\degL P=1\), \cref{plt:eq:lag-degree-sharp} gives \(\degL b_N=N\) and
\(\bigl[M^{N}\bigr]b_N=a^{N+1}/N\) for \(N\ge1\).  The family also has an exact
closed form: \(Y=X+a\log X+b\) is equivalent to
\(\tfrac Xa\EulerE^{X/a}=\tfrac1a\EulerE^{(Y-b)/a}\), whence
\begin{equation}
\label{plt:eq:lag-affine-log-W}
  \rev{F}(Y)=a\,\LambertW_k\!\Bigl(\tfrac1a\EulerE^{(Y-b)/a}\Bigr),
\end{equation}
with the branch \(k\) selected by the intended real interval or sector.
\Cref{plt:eq:lag-affine-log-inverse} is the large-\(Y\) expansion of
\cref{plt:eq:lag-affine-log-W}; the formal calculation does not choose \(k\)
\cite{corless-et-al-1996,dlmf-W,wright-omega}.
\end{example}

\begin{example}[Leading behaviour \(cX^{\rho}\): Puiseux]
\label{plt:ex:lag-puiseux}
Let \(F(X)=X^2+X\).  Here \(\rho=2\), \(\sigma=0\), \(c=1\), so
\cref{plt:thm:lag-rigidity} forbids an inverse in \(\PLlau\) and
\cref{plt:cor:lag-enlargement} prescribes \(\K[L]((t^{1/2}))\).
\Cref{plt:alg:lag-dominant-balance} gives \(F_0=X^2\),
\(\rev{F_0}=X^{1/2}\), \(\Phi=F\compose\rev{F_0}=X+X^{1/2}\), so
\(A=X^{1/2}=t^{-1/2}\) has \(\ordt(A)=-\tfrac12<0\) and step~4 fails: the
normalisation must be redone with \(F_0=X^2+X\) treated exactly, or, more
simply, with the substitution \(X=Z-\tfrac12\).  Doing the latter,
\(F=Z^2-\tfrac14\), and
\[
  \rev{F}(Y)=\Bigl(Y+\tfrac14\Bigr)^{1/2}-\tfrac12
  =Y^{1/2}-\frac12+\frac18Y^{-1/2}-\frac1{128}Y^{-3/2}+\cdots,
\]
which is the branch \(\tfrac12\bigl(-1+\sqrt{1+4Y}\bigr)\) asymptotic to
\(Y^{1/2}\).  No integer-power ansatz can reproduce it.
\end{example}

% ed.: S4 asserts that after the substitution y = d x^{1/q}(1+V) "the formal
% implicit-function recursion is triangular" for any leading exponent q, which
% is true, but it does not record that the intermediate normalisation
% Phi = F o rev(F_0) can leave the near-identity class, as it does in
% plt:ex:lag-puiseux.  The correct general statement is the converse half of
% plt:thm:lag-rigidity, where Phi = X(1 + U o rev(F_0)) because F_0 is the
% exact leading monomial of F, not merely its leading power.
\begin{remark}
\label{plt:rmk:lag-normalisation-care}
\Cref{plt:ex:lag-puiseux} shows that the normalisation step of
\cref{plt:alg:lag-dominant-balance} must use the \emph{full} leading monomial
\(cX^\rho L^\sigma\) of \(F\), not merely a convenient power.  In
\cref{plt:thm:lag-rigidity} the identity
\(\Phi=X\bigl(1+U\compose\rev{F_0}\bigr)\) holds because \(F=F_0(1+U)\)
exactly; if \(F-F_0\) is not \(\LittleOAt{F_0}{X\to+\infty}\), the step
produces a \(\Phi\) outside \(\Gnear\) and
\cref{plt:prop:lag-summability-sharp} then says that
\cref{plt:eq:lag-burmann} is not available.
\end{remark}

\begin{example}[Leading behaviour \(X\log X\): nested logarithm]
\label{plt:ex:lag-xlogx}
Let \(F(X)=X\log X\), so \(\rho=1\), \(\sigma=1\), \(c=1\).  By
\cref{plt:thm:lag-rigidity} there is no monomial-leading inverse in
\(\PLlau\), and by \cref{plt:prop:lag-loglog-forced} none in any Puiseux--log
ring.  \Cref{plt:thm:lag-log-normal-form} gives the leading model in
logarithmic coordinates: \(G(u)=u+\log u\), whose inverse is Wright omega, so
\begin{equation}
\label{plt:eq:lag-xlogx-inverse}
  \rev{F}(Y)=\EulerE^{\WrightOmega(S)}=\frac{Y}{\LambertW(Y)},
  \qquad S=\log Y,
\end{equation}
the second equality because \(Z\log Z=Y\) with \(Z=\EulerE^{w}\) gives
\(w\EulerE^{w}=Y\).  Expanding \(\LambertW(Y)=S-\Lambda+\sum_{n\ge1}\LamP{n}(\Lambda)S^{-n}\)
with \(\Lambda=\log S\) and inverting the resulting series,
\begin{equation}
\label{plt:eq:lag-xlogx-expansion}
  \rev{F}(Y)=\frac{Y}{S}\left(1+\frac{\Lambda}{S}
   +\frac{\Lambda^2-\Lambda}{S^2}
   +\BigOAt{\tfrac{\Lambda^3}{S^3}}{Y\to+\infty}\right).
\end{equation}
Every coefficient is a polynomial in \(\Lambda=\log\log Y\); a one-logarithm
ansatz, however many terms it carries, cannot represent
\cref{plt:eq:lag-xlogx-expansion}.  The same mechanism produces depth two in
the inverses of Dulac transseries with a logarithm in the leading monomial
\cite{mardesic-length,resman-dulac}.
\end{example}

\begin{example}[General power--log monomial]
\label{plt:ex:lag-powerlog}
For \(Y=cX^{\rho}(\log X)^{\sigma}\) with \(c,\rho>0\),
\cref{plt:eq:lag-X-expansion} gives, in fully explicit form,
\[
  X=c^{-1/\rho}\rho^{\sigma/\rho}\,Y^{1/\rho}(\log Y)^{-\sigma/\rho}
   \left[1+\frac{\sigma^2\log\log Y+\sigma\log c-\sigma^2\log\rho}
                {\rho\log Y}
    +\BigOAt{\tfrac{(\log\log Y)^2}{(\log Y)^2}}{Y\to+\infty}\right].
\]
An exact Lambert form is also available when \(\sigma\ne0\): putting
\(s=\log X\), the equation becomes \((Y/c)^{1/\sigma}=s\EulerE^{(\rho/\sigma)s}\),
so
\begin{equation}
\label{plt:eq:lag-powerlog-W}
  X=\exp\!\left[\frac{\sigma}{\rho}
    \LambertW_k\!\left(\frac{\rho}{\sigma}\Bigl(\frac Yc\Bigr)^{1/\sigma}\right)
   \right].
\end{equation}
Because the argument of \(\LambertW\) contains a power of \(Y\), the expansion
of \cref{plt:eq:lag-powerlog-W} contains \(\log Y\) and \(\log\log Y\): a
second logarithm is forced whenever \(\sigma\ne0\), in agreement with
\cref{plt:prop:lag-loglog-forced}.
\end{example}

\begin{warningbox}[title={A scale change that silently corrupts a correct series}]
\label{plt:warn:lag-scale-trap}
Let \(F(X)=2X+\log X\).  This is not near-identity, but it is affine in the
leading scale, so \cref{plt:thm:lag-reduction} applies with the model
\(F_0=m_2\), \(m_2(Z)=2Z\).  Writing
\(G(X)=X+\tfrac12\log X\), one has \(F=m_2\compose G\) and therefore
\begin{equation}
\label{plt:eq:lag-scaled-correct}
  \rev{F}(Y)=\rev{G}\bigl(Y/2\bigr).
\end{equation}
By \cref{plt:eq:lag-scaled-lambert} with \(a=\tfrac12\),
\(\rev{G}(Z)=Z+\sum_{N\ge0}2^{-(N+1)}\LamP{N}(\log Z)Z^{-N}\).  Substituting
\(Z=Y/2\) changes \emph{both} generators:
\(Z^{-N}=2^{N}Y^{-N}\) \emph{and} \(\log Z=\log Y-\log2\).  Hence
\begin{equation}
\label{plt:eq:lag-scale-correct}
  \rev{F}(Y)=\frac Y2+\frac12\sum_{N\ge0}\LamP{N}\bigl(\log Y-\log2\bigr)Y^{-N}.
\end{equation}
The silent failure is to substitute only the power, leaving \(\log Z\) as
\(\log Y\).  The result,
\(\tfrac Y2+\tfrac12\sum_N\LamP{N}(\log Y)Y^{-N}\), is a perfectly well formed
element of \(\PLlau\); every block has the right shape, the right degree, and
the right sign pattern, and no formal check inside the near-identity
calculation detects anything.  It is wrong by a constant.  Numerically, with
\(Y=20+\log10=22.302585\ldots\), for which \(\rev{F}(Y)=10\) exactly, the
truncations through \(Y^{-3}\) give
\[
  \text{\cref{plt:eq:lag-scale-correct}}:\ 10.000004\ldots,
  \qquad
  \text{naive substitution}:\ 9.67026\ldots,
\]
the leading discrepancy being \(\tfrac12\log2=0.34657\ldots\).  Two further
observations belong with this example.
\begin{itemize}
\item The coefficient field must contain \(\log2\).  Over \(\K=\RationalNumbers\)
      the identity \cref{plt:eq:lag-scale-correct} is not even a statement about
      \(\PLlau\); the field has to be extended first
      (\cref{plt:alg:lag-dominant-balance}, step~1).
\item The residual certificate of \cref{plt:cor:rev-finite-certificate} does catch the
      error, because it is computed against \(F\), not against \(G\).  This is
      why step~7 of \cref{plt:alg:lag-dominant-balance} certifies against the
      original map and never against the normalised one.
\end{itemize}
\end{warningbox}

\begin{remark}[Sufficient, not necessary]
\label{plt:rmk:lag-rank-two}
\Cref{plt:prop:lag-summability-sharp} settles summability \emph{for the outer
valuation}.  It does not say that \cref{plt:eq:lag-burmann} fails whenever
\(\ordt(A)<0\); it says the theorem as proved does not apply.  Take
\(F(X)=X+X/L=X(1+1/L)\), so \(A=t^{-1}L^{-1}\in\PLit\) with \(\ordt(A)=-1\),
while \(A\precdom X\).  Here every Lagrange summand contributes to the single
outer block \(t^{-1}\):
\[
  -A=-XL^{-1},\quad
  \tfrac12\DX(A^2)=XL^{-2}-XL^{-3},\quad
  -\tfrac16\DX^{2}(A^3)=-XL^{-3}+\tfrac52XL^{-4}-2XL^{-5},
\]
and so on, so no finite number of summands determines any block.  The family is
nevertheless summable for the \(L^{-1}\)-adic topology of the coefficient field
\(\K((L^{-1}))\) at each fixed outer level, and the resulting series
\(Y\bigl(1-M^{-1}+M^{-2}-2M^{-3}+\cdots\bigr)\), \(M=\log Y\), does satisfy
\(F\compose\rev{F}-Y=\BigOAt{YM^{-4}}{Y\to+\infty}\), as direct substitution
confirms.  A theorem covering this case needs a rank-two value group and
Hahn-style summability, in which \(\SetOf{i:\nu(f_i)<\gamma}\) is replaced by a
well-ordered-support condition.  We have not carried that out, and we do not
claim it: the statement ``\cref{plt:eq:lag-burmann} holds whenever
\(A\precdom X\) in the rank-two order, with summation in the refined
filtration'' is left as an open problem.
\end{remark}

\begin{remark}[What does not belong here]
\label{plt:rmk:lag-branch-point}
Reversion at a critical point is a different problem in a different regime.  If
\(f(x)-f(x_0)=c_m(x-x_0)^m+\cdots\) with \(c_m\ne0\) and \(m\ge2\), the local
inverse begins \(x-x_0\sim c_m^{-1/m}(y-f(x_0))^{1/m}\) and has \(m\) branches;
Lambert \(\LambertW\) near \(-\EulerE^{-1}\) is the case \(m=2\)
\cite{corless-et-al-1996}.  This is a Puiseux expansion at a finite point, not
an expansion in the scale \(X\to+\infty\), and none of the machinery of this
section applies to it.  In the at-infinity setting the derivative never
vanishes on the relevant class: for \(F=X+A\) with \(A\in\PLpow\),
\(\DX F=1+\FormalBigOAt{t}{t}\) is a unit of \(\PLpow\).
\end{remark}

\begin{summarybox}[title={Beyond the identity scale}]
A monomial-leading \(F=cX^{\rho}L^{\sigma}(1+U)\) has a monomial-leading
inverse inside a Puiseux--log ring exactly when \(\rho^{-1}\) lies in the
exponent group and \(\sigma=0\)
(\cref{plt:thm:lag-rigidity}).  A noninteger \(\rho\) is repaired by enlarging
the exponent group; \(\sigma\ge1\) cannot be repaired that way and forces
\(\log\log X\) (\cref{plt:prop:lag-loglog-forced}).  In every case reversion
factors as \(\rev{F}=\rev{F_0}\compose\rev{\Phi}\) with
\(\Phi=F\compose\rev{F_0}\) near-identity
(\cref{plt:thm:lag-reduction}), and the leading model \(F_0\) is inverted by
conjugating with \(\EulerE\) and dividing by \(\rho\)
(\cref{plt:thm:lag-log-normal-form}).  Every scale change substitutes
\emph{both} generators; substituting only the power silently produces a
plausible and wrong series (\cref{plt:warn:lag-scale-trap}).
\end{summarybox}

\clearpage
\part{Wright omega, the Lambert polynomials, and Lambert \texorpdfstring{\(\LambertW\)}{W}}

% ---- fragment 10-omega ----
\chapter{Wright omega: the canonical logarithmic inverse}
\label{plt:sec:wright-omega}

Every result of Parts~I--IV was stated for a general polynomial--logarithmic
series.  This part specialises all of it to a single map,
\[
 F(Y)=Y+\log Y,
\]
and shows that this one reversion carries the whole Lambert circle of
problems.  The correction blocks of \(\rev{F}\) are the Lambert polynomials
\(\LamP{n}\); they are defined and studied in
\cref{plt:sec:lambert-polynomials}, and the present chapter explains where
they come from and exactly what is being asserted when one writes them down.

\section{Standing conventions for this part}
\label{plt:sec:om-conventions}

Throughout, \(\K\) is a field of characteristic zero, so that
\(\RationalNumbers\subseteq\K\).  On the formal side we work in the rings of
\cref{plt:def:ring-four-classes} with the large variable \(X\), the small generator
\(t=X^{-1}\), and the logarithmic generator \(L=\log X\); the near-identity
class is \(\Gnear=X+\K[L][[t]]\) of \cref{plt:def:grp-gnear}, the derivation is
\(\DX=t\,\partial_L-t^2\,\partial_t\), and truncation is always the
\emph{exclusive} \(\Trunc{F}{N}=\operatorname{Trunc}_{<N}F\).

On the analytic side the regime is \(x\to+\infty\) along the positive real
ray, with the real logarithm; every deviation from that regime --- a complex
sector, a nonprincipal branch, a limit at a finite point --- is announced
explicitly.

Two symbols carry a variable role and must be read locally.  In this section
\(L\) is the generator \(\log X\) of the coefficient ring.  In
\cref{plt:sec:lambert-polynomials} the Lambert polynomials are treated as
polynomials in an indeterminate \(u\), and \(s\) is the formal variable of
their generating series; in \cref{plt:thm:om-laplace} the letter \(s\) is,
instead, a complex parameter of an honest integral.  The two roles of \(s\)
are unrelated and never appear in the same statement.

\begin{warningbox}
Three inverses meet in this part and must not be conflated.  For a series
\(F\), the multiplicative inverse is \(\recip{F}\) (\cref{plt:thm:div-reciprocal})
and the compositional inverse is \(\rev{F}\) (\cref{plt:thm:rev-existence});
for the Lambert function, the subscript in \(\LambertW_{-1}\) is a
\emph{branch label} and is neither of the two.  The glyph \(-1\) is the same
in all three cases and the objects are different.
\end{warningbox}

\section{The real bijection and its inverse}
\label{plt:sec:om-real}

\begin{proposition}[The real logarithmic bijection]
\label{plt:prop:om-real-bijection}
The map
\[
 \phi\colon(0,\infty)\longrightarrow\RealNumbers,
 \qquad
 \phi(Y)=Y+\log Y,
\]
is a strictly increasing real-analytic bijection.  Its inverse
\(\WrightOmega\colon\RealNumbers\to(0,\infty)\) is strictly increasing and
real-analytic, and satisfies, for every \(x\in\RealNumbers\),
\begin{equation}\label{plt:eq:om-defining}
 \WrightOmega(x)+\log\WrightOmega(x)=x,
 \qquad
 \WrightOmega'(x)=\frac{\WrightOmega(x)}{1+\WrightOmega(x)},
 \qquad
 \WrightOmega(x)=\LambertW_0(\EulerE^{x}).
\end{equation}
Here \(\LambertW_0\) is the branch of the inverse of \(w\mapsto w\EulerE^{w}\)
taking values in \([-1,\infty)\).  Moreover \(\WrightOmega(1)=1\), and for
every \(x\) with \(\WrightOmega(x)>1\),
\begin{equation}\label{plt:eq:om-envelope}
 x-\log x<\WrightOmega(x)<x.
\end{equation}
\end{proposition}

\begin{proof}
On \((0,\infty)\) the function \(\phi\) is real-analytic with
\(\phi'(Y)=1+Y^{-1}>1>0\), so it is strictly increasing and injective.  As
\(Y\to0^{+}\) we have \(\log Y\to-\infty\) and \(Y\to0\), so
\(\phi(Y)\to-\infty\); as \(Y\to+\infty\), \(\phi(Y)\ge Y\to+\infty\).  By
the intermediate value theorem \(\phi\) is onto \(\RealNumbers\).  Hence
\(\phi\) is a bijection and \(\WrightOmega\DefinitionEquals\rev{\phi}\) is
defined on all of \(\RealNumbers\), strictly increasing, and real-analytic by
the inverse function theorem, \(\phi'\) being nowhere zero.

Differentiating \(\WrightOmega+\log\WrightOmega=x\) gives
\(\WrightOmega'(1+\WrightOmega^{-1})=1\), i.e.\ the displayed derivative
formula.  For the Lambert identity fix \(x\in\RealNumbers\) and put
\(w=\WrightOmega(x)>0\).  Then
\(\log(w\EulerE^{w})=\log w+w=x\), so \(w\EulerE^{w}=\EulerE^{x}\).  The map
\(w\mapsto w\EulerE^{w}\) is strictly increasing on \([-1,\infty)\) with value
\(-\EulerE^{-1}\) at \(w=-1\), hence injective there; since \(w>0>-1\), \(w\)
is the value of the principal branch, i.e.\ \(w=\LambertW_0(\EulerE^{x})\).

Finally \(\phi(1)=1+\log1=1\), so \(\WrightOmega(1)=1\).  If
\(Y=\WrightOmega(x)>1\) then \(\log Y>0\), so \(Y=x-\log Y<x\); and then
\(\log Y<\log x\), whence \(Y=x-\log Y>x-\log x\).
\end{proof}

\begin{remark}[What is not repeated here]
\label{plt:rmk:om-envelope}
The elementary two-sided envelope for \(\WrightOmega\) and the sharper
two-order statement were established in \cref{plt:prop:mot-omega-basic} and
\cref{plt:prop:mot-two-orders}; \cref{plt:eq:om-envelope} is recorded only
because the proof of \cref{plt:thm:om-analytic} uses the crude lower bound
\(\WrightOmega(x)>x-\log x\) directly.  In particular
\(\WrightOmega(x)\to+\infty\) and \(\WrightOmega(x)/x\to1\) as
\(x\to+\infty\).
\end{remark}

\begin{remark}[Totality versus asymptotics]
\label{plt:rmk:om-total}
\(\WrightOmega\) is a genuine function on all of \(\RealNumbers\); it is not
defined only near \(+\infty\).  The polynomial--logarithmic expansion studied
below describes it \emph{only} as \(x\to+\infty\), and says nothing about,
say, \(x\le0\), where \(\log x\) is not even real.  Several of the six source
drafts introduce \(\WrightOmega\) with the phrase ``for \(x\) sufficiently
large''; that restriction belongs to the expansion, not to the function.
\end{remark}

\section{Every Lambert problem is this one reversion}
\label{plt:sec:om-reduction}

The point of isolating \(\WrightOmega\) is that a one-parameter family of
logarithmic inversion problems collapses onto it exactly, with no new
coefficients.  We give the analytic reduction first and the formal one
second.

\begin{lemma}[Exact scaling reduction, analytic form]
\label{plt:lem:om-scaling-exact}
Let \(a>0\) be real.  Then \(y\mapsto y+a\log y\) is a strictly increasing
bijection of \((0,\infty)\) onto \(\RealNumbers\), and its inverse is
\begin{equation}\label{plt:eq:om-scaling-exact}
 y=a\,\WrightOmega\!\left(\frac{x}{a}-\log a\right),
 \qquad x\in\RealNumbers .
\end{equation}
Equivalently, for \(\alpha>0\) and \(X>0\), the unique \(y>0\) with
\(y^{\alpha}\EulerE^{y}=X\) is
\(y=\alpha\,\WrightOmega\!\left(\alpha^{-1}\log X-\log\alpha\right)\).
\end{lemma}

\begin{proof}
Write \(y=az\) with \(z>0\).  Then
\(y+a\log y=az+a\log a+a\log z=a\bigl(z+\log z\bigr)+a\log a\), so
\(y+a\log y=x\) is equivalent to \(z+\log z=x/a-\log a\), that is to
\(z=\WrightOmega(x/a-\log a)\) by \cref{plt:prop:om-real-bijection}.
Monotonicity and bijectivity are inherited from \(\phi\) through the same
substitution.  For the second form, take logarithms in
\(y^{\alpha}\EulerE^{y}=X\): since \(y>0\) and \(X>0\), this is equivalent to
\(y+\alpha\log y=\log X\), and \eqref{plt:eq:om-scaling-exact} applies with
\(a=\alpha\), \(x=\log X\).
\end{proof}

The formal counterpart is stronger: it holds over any characteristic-zero
\(\K\), for every scalar \(a\), and it exhibits the \emph{same} polynomials
with only a scalar reweighting.

\begin{theorem}[The one-parameter family of logarithmic reversions]
\label{plt:thm:om-alpha-family}
Let \(a\in\K\) and let \(F_a\DefinitionEquals X+aL\in\Gnear\).  Then
\(\rev{F_a}\in\Gnear\) exists, is unique, and equals
\begin{equation}\label{plt:eq:om-alpha-family}
 \rev{F_a}(X)
 =X+\sum_{n\ge0}a^{\,n+1}\LamP{n}(L)\,t^{n}
 =X-aL+\sum_{n\ge1}a^{\,n+1}\LamP{n}(L)\,t^{n},
\end{equation}
with \(\LamP{n}\) as in \cref{plt:def:lambert-polynomials}.  In particular
\(\rev{F_1}=\WrightOmega\) and
\begin{equation}\label{plt:eq:om-minus-log}
 \rev{(X-L)}(X)=X+L+\sum_{n\ge1}(-1)^{n+1}\LamP{n}(L)\,t^{n}.
\end{equation}
\end{theorem}

\begin{proof}
Existence and uniqueness are \cref{plt:thm:rev-existence}, since
\(aL\in\K[L][[t]]\).  For the formula, write
\(\mathcal P(t,L)\DefinitionEquals\sum_{n\ge0}\LamP{n}(L)t^{n}\in\K[L][[t]]\).
By \cref{plt:prop:om-implicit},
\begin{equation}\label{plt:eq:om-implicit-recall}
 \mathcal P+L+\log(1+t\mathcal P)=0 .
\end{equation}
Substituting \(t\mapsto at\) into \eqref{plt:eq:om-implicit-recall} --- a
legitimate \(\K\)-algebra endomorphism of \(\K[L][[t]]\), continuous for the
\(t\)-adic filtration, hence commuting with the formal logarithm --- gives
\[
 \mathcal P(at,L)+L+\log\bigl(1+at\,\mathcal P(at,L)\bigr)=0 .
\]
Put \(g\DefinitionEquals a\,\mathcal P(at,L)\in\K[L][[t]]\) and multiply by
\(a\):
\begin{equation}\label{plt:eq:om-alpha-implicit}
 g+aL+a\log(1+tg)=0 .
\end{equation}
Now let \(G=X+g\).  By generator substitution
(\cref{plt:prop:cmp-generator-substitution}), \(L\compose G=L+\log(1+tg)\),
so
\[
 F_a\compose G=G+a\,(L\compose G)=X+g+aL+a\log(1+tg)=X
\]
by \eqref{plt:eq:om-alpha-implicit}.  Thus \(G\) is a right inverse of
\(F_a\) in \(\Gnear\), hence the two-sided inverse by uniqueness
(\cref{plt:thm:rev-existence}(2),(3)).  Finally
\(g=a\sum_{n\ge0}\LamP{n}(L)a^{n}t^{n}=\sum_{n\ge0}a^{n+1}\LamP{n}(L)t^{n}\),
which is \eqref{plt:eq:om-alpha-family}; and
\(a^{n+1}\big|_{a=-1}=(-1)^{n+1}\), with
\((-1)^{1}\LamP{0}(L)=L\), gives \eqref{plt:eq:om-minus-log}.
\end{proof}

\begin{remark}[Why no ``second Lambert family'' exists]
\label{plt:rmk:om-second-family}
% ed.: S1's prose ("the same P_n occur, multiplied by (-1)^n") contradicts its
% own displayed expansion of the inverse of X - log X, where P_1 = L occurs
% with a plus sign.  The correct weight is (-1)^{n+1}; the weight (-1)^n
% belongs to the negated quantity -rev(X-L), which is what enters W_{-1}.
The inverse of \(X-\log X\) is often presented as a separate expansion with
its own alternating coefficients.  Equation \eqref{plt:eq:om-minus-log}
settles the matter: the blocks are \((-1)^{n+1}\LamP{n}\), the same
polynomials with a sign.  The frequently quoted weight \((-1)^{n}\) is
correct for \(-\rev{(X-L)}\), not for \(\rev{(X-L)}\) itself; that sign flip
is exactly the substitution \(w\mapsto-w\) used to reach the lower real
Lambert branch, and it is the source of a recurring off-by-one-sign error in
the source drafts.  Concretely, for \(x\to0^{-}\), putting
\(T=\log(1/(-x))\) and \(\ell=\log T\), the substitution
\(\LambertW_{-1}(x)=-v\) turns \(w\EulerE^{w}=x\) into \(v-\log v=T\), so
\eqref{plt:eq:om-minus-log} gives the block pattern
\[
 \LambertW_{-1}(x)\;\text{blocks}\;=\;-T-\ell+\sum_{n\ge1}(-1)^{n}\LamP{n}(\ell)T^{-n}.
\]
The formal identity is \eqref{plt:eq:om-minus-log}; the asymptotic assertion
about \(\LambertW_{-1}\) needs the analytic transfer of
\cref{plt:sec:om-analytic} applied on the branch in question and is not
proved here.
\end{remark}

\begin{remark}[Branch coordinates]
\label{plt:rmk:om-branches}
Taking a compatible logarithm in \(\LambertW(z)\EulerE^{\LambertW(z)}=z\)
gives \(\LambertW_k(z)+\log\LambertW_k(z)=\xi_k\) with
\(\xi_k=\PrincipalLogarithm z+2\pi\ImaginaryUnit k\).  Thus every large
branch of \(\LambertW\) is the \emph{same} reversion evaluated at a different
large variable, and the formal expansion of \(\LambertW_k\) is obtained from
\cref{plt:thm:om-expansion} by the substitutions \(X\mapsto\xi_k\),
\(L\mapsto\ell_k=\log_{\mathcal B}\xi_k\) for a declared branch
\(\mathcal B\).  The polynomials \(\LamP{n}\) are therefore
branch-independent; the coordinates \(\xi_k,\ell_k\), the branch
\(\mathcal B\), the sector, and the resulting analytic estimates are not.  We
do not develop the complex branch analysis here.
\end{remark}

\section{The formal expansion}
\label{plt:sec:om-expansion}

The reversion machinery of Part~IV applies verbatim, and in this special case
it degenerates in a useful way: each Lagrange--B\"urmann summand produces
exactly one block, so no cross-order collection is ever needed.  We first
record the operator identity that makes this visible.

\begin{lemma}[Repeated derivative of a coefficient block]
\label{plt:lem:om-repeated-derivative}
Let \(n\in\IntegerNumbers\), \(k\in\NaturalNumbers\) and \(p\in\K[L]\).  Then
\begin{equation}\label{plt:eq:om-repeated-derivative}
 \DX^{k}\bigl(t^{n}p(L)\bigr)=t^{\,n+k}\,\shiftop{n}{k}\,p(L),
 \qquad
 \shiftop{n}{k}=\prod_{j=0}^{k-1}\bigl(\partial_L-n-j\bigr),
\end{equation}
the empty product at \(k=0\) acting as the identity.
\end{lemma}

\begin{proof}
Induction on \(k\).  The case \(k=0\) is trivial.  Assume
\eqref{plt:eq:om-repeated-derivative} for \(k\) and write
\(q=\shiftop{n}{k}p\in\K[L]\).  Since \(\DX t=-t^{2}\) and \(\DX L=t\),
Leibniz and the chain rule for the derivation \(\DX\) give
\[
 \DX\bigl(t^{\,n+k}q(L)\bigr)
 =(n+k)t^{\,n+k-1}(-t^{2})q(L)+t^{\,n+k}q'(L)\,t
 =t^{\,n+k+1}\bigl(\partial_L-(n+k)\bigr)q(L).
\]
The factors of \(\shiftop{n}{k}\) are polynomials in \(\partial_L\) and
therefore commute, so
\(\bigl(\partial_L-(n+k)\bigr)\shiftop{n}{k}=\shiftop{n}{k+1}\), which is the
statement for \(k+1\).
\end{proof}

\begin{proposition}[The implicit equation for the correction]
\label{plt:prop:om-implicit}
Let \(\mathcal T\colon\K[L][[t]]\to\K[L][[t]]\) be
\(\mathcal T(U)=-L-\log(1+tU)\), the logarithm being the formal one of
\cref{plt:lem:cmp-exp-log}.  Then \(\mathcal T\) has a unique fixed point
\(\mathcal P\in\K[L][[t]]\); equivalently, the equation
\begin{equation}\label{plt:eq:om-implicit}
 \mathcal P+L+\log(1+t\mathcal P)=0
\end{equation}
has exactly one solution in \(\K[L][[t]]\).  That solution lies in
\(-L+t\K[L][[t]]\), satisfies \(\mathcal P(t,0)=0\), and
\begin{equation}\label{plt:eq:om-P-is-correction}
 \WrightOmega=\rev{(X+L)}=X+\mathcal P .
\end{equation}
\end{proposition}

\begin{proof}
For \(U\in\K[L][[t]]\) we have \(tU\in t\K[L][[t]]\), so the family
\(\bigl((-1)^{j-1}(tU)^{j}/j\bigr)_{j\ge1}\) is formally summable and
\(\log(1+tU)\in t\K[L][[t]]\); hence \(\mathcal T\) is well defined and maps
into \(-L+t\K[L][[t]]\).  If \(\ordt(U-V)\ge N\), then, \(1+tV\) being a unit
of \(\K[L][[t]]\),
\[
 \mathcal T(U)-\mathcal T(V)
 =-\log\frac{1+tU}{1+tV}
 =-\log\left(1+\frac{t(U-V)}{1+tV}\right)\in t^{\,N+1}\K[L][[t]],
\]
so \(\ordt\bigl(\mathcal T(U)-\mathcal T(V)\bigr)\ge\ordt(U-V)+1\).
% ed.: S5 states uniqueness only inside -L + t K[L][[t]]; the contraction
% estimate holds on all of K[L][[t]], so uniqueness is unconditional there and
% membership in -L + t K[L][[t]] is a conclusion, not a hypothesis.
Thus \(\mathcal T\) is a strict contraction for the \(t\)-adic filtration on
the whole of \(\K[L][[t]]\), which is separated and complete
(\cref{plt:prop:grp-limit}).  Two fixed points would satisfy
\(\ordt(U-V)\ge\ordt(U-V)+1\), forcing \(\ordt(U-V)=+\infty\) and \(U=V\);
and the iterates \(U_0=0\), \(U_{r+1}=\mathcal T(U_r)\) form a Cauchy
sequence whose limit is a fixed point.  Its membership in
\(-L+t\K[L][[t]]\) is the image statement above.

Setting \(L=0\) in \eqref{plt:eq:om-implicit} --- that is, applying the
\(\K\)-algebra map \(\K[L][[t]]\to\K[[t]]\), \(L\mapsto0\), which is
\(t\)-adically continuous and hence commutes with the formal logarithm ---
shows that \(V\DefinitionEquals\mathcal P(t,0)\) solves
\(V+\log(1+tV)=0\) in \(\K[[t]]\).  The same contraction argument applied to
\(V\mapsto-\log(1+tV)\) gives a unique solution, and \(V=0\) is one; hence
\(\mathcal P(t,0)=0\).

Finally put \(G=X+\mathcal P\in\Gnear\).  By
\cref{plt:prop:cmp-generator-substitution},
\((X+L)\compose G=G+L+\log(1+t\mathcal P)=X\) by
\eqref{plt:eq:om-implicit}, so \(G\) is a right inverse of \(X+L\), hence
\(G=\rev{(X+L)}\) by \cref{plt:thm:rev-existence}(2),(3).  That \(\rev{(X+L)}\) is what
the notation \(\WrightOmega\) denotes formally is
\cref{plt:prop:om-real-bijection} together with
\cref{plt:thm:om-analytic} below; at this point
\eqref{plt:eq:om-P-is-correction} is a definition of the symbol on the formal
side.
\end{proof}

\begin{theorem}[Polynomial--logarithmic expansion of \(\WrightOmega\)]
\label{plt:thm:om-expansion}
In \(\Gnear\),
\begin{equation}\label{plt:eq:om-expansion}
 \WrightOmega
 =X+\sum_{n\ge0}\LamP{n}(L)\,t^{n}
 =X-L+\sum_{n\ge1}\LamP{n}(L)\,t^{n},
\end{equation}
where \(\LamP{n}\) is the Lambert polynomial of
\cref{plt:def:lambert-polynomials}.  More precisely, the \(r\)-th
Lagrange--B\"urmann summand for \(F=X+L\) is a \emph{single} block:
\begin{equation}\label{plt:eq:om-one-block}
 \frac{(-1)^{r}}{r!}\,\DX^{\,r-1}\bigl(L^{r}\bigr)=\LamP{r-1}(L)\,t^{\,r-1},
 \qquad r\ge1 .
\end{equation}
\end{theorem}

\begin{proof}
% ed.: the citation named a label "plt:thm:lagrange-burmann" that is defined
% ed.: nowhere in the volume; the theorem meant is the one labelled
% ed.: plt:thm:lag-burmann in the Lagrange--Buermann part, whose
% ed.: polynomial--logarithmic form is equation plt:eq:lag-burmann-pl.
Apply the Lagrange--B\"urmann formula \cref{plt:thm:lag-burmann}, in the form
\eqref{plt:eq:lag-burmann-pl}, to
\(F=X+A\) with \(A=L\):
\[
 \rev{F}(X)=X+\sum_{r\ge1}\frac{(-1)^{r}}{r!}\DX^{\,r-1}\bigl(L^{r}\bigr),
\]
the family being formally summable because the \(r\)-th summand has
\(t\)-order at least \(r-1\).  By
\cref{plt:lem:om-repeated-derivative} with \(n=0\), \(k=r-1\), \(p=L^{r}\),
\[
 \DX^{\,r-1}\bigl(L^{r}\bigr)=t^{\,r-1}\shiftop{0}{r-1}L^{r}
 =t^{\,r-1}\left[\prod_{j=0}^{r-2}(\partial_L-j)\right]L^{r}.
\]
Setting \(n=r-1\), the scalar prefactor is
\(\frac{(-1)^{r}}{r!}=\frac{(-1)^{n+1}}{(n+1)!}\) and the operator is
\(\prod_{j=0}^{n-1}(\partial_L-j)\) applied to \(L^{n+1}\); by
\cref{plt:def:lambert-polynomials} the product of the two is exactly
\(\LamP{n}(L)\).  This is \eqref{plt:eq:om-one-block}.  Since the summand of
index \(r\) has \(t\)-order exactly \(r-1\) and no other summand meets that
order, summing over \(r\ge1\) gives \eqref{plt:eq:om-expansion} with no
collection of contributions; the \(r=1\) term is \(-L=\LamP{0}(L)\).
\end{proof}

\begin{remark}[Two normalisations of the same expansion]
\label{plt:rmk:om-normalisation}
The two displays in \eqref{plt:eq:om-expansion} are the same series.  The
zero-based one is canonical: it is the form in which
\eqref{plt:eq:om-one-block} says ``one Lagrange index per block'', and the
form in which \cref{plt:thm:om-alpha-family} needs no case distinction.  The
one-based form, which displays \(-L\) separately, is the classical way of
writing the Lambert expansion and is the one the source drafts mostly use.
The price of the zero-based normalisation is that
\(\deg\LamP{0}=1\ne0\): the degree law \(\deg\LamP{n}=n\) holds for
\(n\ge1\) only (\cref{plt:cor:lambert-coefficients}).
\end{remark}

\section{The truncation residual}
\label{plt:sec:om-residual}

\begin{definition}[Omega truncations]
\label{plt:def:om-truncation}
For \(N\in\NaturalNumbers\) set
\begin{equation}\label{plt:eq:om-truncation}
 \WrightOmega_N
 \DefinitionEquals X+\Trunc{\mathcal P}{N+1}
 =X-L+\sum_{n=1}^{N}\LamP{n}(L)\,t^{n}\in\Gnear .
\end{equation}
In the inclusive notation of source~6 this is \(X+[\mathcal P]_{\le N}\); the
one-block difference \([F]_{\le N}=\Trunc{F}{N+1}\) is the only discrepancy
between the two conventions and is a frequent source of apparent off-by-one
disagreements between the drafts.
\end{definition}

\begin{theorem}[Exact leading residual]
\label{plt:thm:om-residual}
For every \(N\in\NaturalNumbers\), in \(\K[L][[t]]\),
\begin{equation}\label{plt:eq:om-residual}
 \RevErr{X+L}{\WrightOmega_N}
 =\WrightOmega_N+\log\WrightOmega_N-X
 =-\LamP{N+1}(L)\,t^{\,N+1}+\FormalBigOAt{t^{\,N+2}}{t}.
\end{equation}
In particular \(\ordt\bigl(\RevErr{X+L}{\WrightOmega_N}\bigr)=N+1\) exactly,
for every \(N\ge0\).
\end{theorem}

\begin{proof}
Write \(\mathcal P_N=\Trunc{\mathcal P}{N+1}\) and
\(\Delta=\mathcal P-\mathcal P_N=\sum_{n\ge N+1}\LamP{n}(L)t^{n}\), so that
\(\ordt\Delta=N+1\) and \([t^{N+1}]\Delta=\LamP{N+1}(L)\).  By generator
substitution (\cref{plt:prop:cmp-generator-substitution}),
\[
 \RevErr{X+L}{\WrightOmega_N}
 =\mathcal P_N+L+\log(1+t\mathcal P_N).
\]
Using \(\mathcal P_N=\mathcal P-\Delta\), the identity
\(1+t\mathcal P_N=(1+t\mathcal P)\Bigl(1-\dfrac{t\Delta}{1+t\mathcal P}\Bigr)\),
and additivity of the formal logarithm on \(1+t\K[L][[t]]\)
(\cref{plt:lem:cmp-exp-log}), we get
\[
 \RevErr{X+L}{\WrightOmega_N}
 =\underbrace{\mathcal P+L+\log(1+t\mathcal P)}_{=0\ \text{by }\eqref{plt:eq:om-implicit}}
 \;-\;\Delta\;+\;\log\!\left(1-\frac{t\Delta}{1+t\mathcal P}\right).
\]
Since \(1+t\mathcal P\) is a unit of \(\K[L][[t]]\) and \(\ordt\Delta=N+1\),
the argument of the last logarithm lies in \(1+t^{\,N+2}\K[L][[t]]\), so that
logarithm lies in \(t^{\,N+2}\K[L][[t]]\).  Hence
\[
 \RevErr{X+L}{\WrightOmega_N}
 =-\Delta+\FormalBigOAt{t^{\,N+2}}{t}
 =-\LamP{N+1}(L)t^{\,N+1}+\FormalBigOAt{t^{\,N+2}}{t}.
\]
Exactness of the order follows because \(\LamP{N+1}\ne0\) for every
\(N\ge0\), by \cref{plt:cor:lambert-coefficients}.
\end{proof}

% ed.: retitled.  This is not the general finite certificate
% ed.: (\cref{plt:cor:rev-finite-certificate}, which it invokes) but its
% ed.: specialisation to the single map \(X+L\), where the criterion becomes a
% ed.: characterisation of the Lambert polynomials themselves.
\begin{corollary}[Finite certificate for the Lambert polynomials]
\label{plt:cor:om-certificate}
Let \(p_0,\ldots,p_N\in\K[L]\) and set
\(H=X+\sum_{n=0}^{N}p_n(L)t^{n}\).  Then \(p_n=\LamP{n}\) for
\(0\le n\le N\) if and only if
\(\RevErr{X+L}{H}=\FormalBigOAt{t^{\,N+1}}{t}\).
\end{corollary}

\begin{proof}
Necessity is \cref{plt:thm:om-residual}.  Conversely, if the residual has
\(t\)-order at least \(N+1\) then \(H\) is a right inverse of \(X+L\) through
block \(N\) in the sense of \cref{plt:cor:rev-finite-certificate}, so
\(\Trunc{H-X}{N+1}=\Trunc{\mathcal P}{N+1}\), i.e.\ \(p_n=\LamP{n}\) for
\(n\le N\).
\end{proof}

\begin{example}[The residual at \(N=2\)]
\label{plt:ex:om-residual-2}
With \(\mathcal P_2=-L+Lt+\bigl(\tfrac12L^{2}-L\bigr)t^{2}\),
\[
 \mathcal P_2+L+\log(1+t\mathcal P_2)
 =-\left(\tfrac13L^{3}-\tfrac32L^{2}+L\right)t^{3}+\FormalBigOAt{t^{4}}{t}
 =-\LamP{3}(L)t^{3}+\FormalBigOAt{t^{4}}{t},
\]
which one checks by expanding
\(\log(1+t\mathcal P_2)=t\mathcal P_2-\tfrac12t^{2}\mathcal P_2^{2}+\tfrac13t^{3}\mathcal P_2^{3}-\cdots\)
and collecting \(t^{0},t^{1},t^{2},t^{3}\).  This single computation
certifies \(\LamP{0},\LamP{1},\LamP{2}\) simultaneously and predicts
\(\LamP{3}\); it is the cheapest available unit test for an implementation of
the recurrence.
\end{example}

\begin{example}[Every truncation is exact at \(x=1\)]
\label{plt:ex:om-exact-at-one}
Because \(\LamP{n}(0)=0\) for all \(n\) (\cref{plt:thm:lambert-recurrence}),
substituting \(x=1\), where \(L=\log1=0\), gives
\(\WrightOmega_N(1)=1-0+0=1=\WrightOmega(1)\) for every \(N\).  The same
phenomenon in the Lambert coordinates is the exactness of the branch
expansion at \(z=\EulerE\), where \(\xi_0=1\) and \(\ell_0=0\).  It is a
consequence of the vanishing \(\LamP{n}(0)=0\), not of any convergence
property.
\end{example}

\section{Formal statement, analytic statement}
\label{plt:sec:om-analytic}

\Cref{plt:thm:om-residual} is an identity in \(\K[L][[t]]\).  It asserts
coefficient agreement below a cutoff and nothing else: no numerical bound, no
convergence, no uniformity in \(L\), no claim that \(\WrightOmega\) exists as
a function.  An analytic statement about the real function of
\cref{plt:prop:om-real-bijection} needs four further ingredients, none of
which is formal algebra:

\begin{enumerate}
\item a \emph{regime}: here the ray \(x\to+\infty\) with the real logarithm,
      so that \(L=\log x\) is a real quantity that \emph{grows with} \(x\);
\item an \emph{analytic remainder} for the truncated logarithm, replacing the
      formal \(\log(1+t\mathcal P_N)\) by a finite sum plus an estimated tail;
\item a \emph{lower bound on the derivative} of \(\phi\) between the true and
      the approximate root, to convert a residual bound into an error bound;
\item an \emph{admissibility check} that the approximation stays inside the
      domain of \(\phi\), here \(\WrightOmega_N(x)>0\).
\end{enumerate}

% ed.: S5 asserts that the residual-to-error principle "proves the asymptotic
% validity of every truncation directly from" the FORMAL residual identity,
% and S1 asserts an analytic O-bound for the branch residual "by
% construction".  Neither is a proof: a formal t-adic identity supplies no
% inequality.  The missing ingredient is item (2), supplied below.
The following theorem supplies all four.  Note that the estimate is
\emph{not} uniform in \(N\), and that the exponent of \(\log x\) in the
remainder is quadratic in \(N\), not linear; this is an honest feature of the
argument, not an artefact.

\begin{theorem}[Asymptotic validity on the positive real ray]
\label{plt:thm:om-analytic}
Fix \(N\in\NaturalNumbers\) and put \(K=(N+2)^{2}\).  Then, with
\(L=\log x\),
\begin{equation}\label{plt:eq:om-analytic-residual}
 \WrightOmega_N(x)+\log\WrightOmega_N(x)-x
 =-\frac{\LamP{N+1}(L)}{x^{\,N+1}}
  +\BigOAt{\frac{L^{K}}{x^{\,N+2}}}{x\to+\infty}
\end{equation}
and
\begin{equation}\label{plt:eq:om-analytic-error}
 \WrightOmega(x)-\WrightOmega_N(x)
 =\frac{\LamP{N+1}(L)}{x^{\,N+1}}
  +\BigOAt{\frac{L^{K}}{x^{\,N+2}}}{x\to+\infty}.
\end{equation}
Consequently
\begin{equation}\label{plt:eq:om-sharp-order}
 \WrightOmega(x)-\WrightOmega_N(x)
 \sim\frac{(\log x)^{N+1}}{(N+1)\,x^{\,N+1}},
 \qquad x\to+\infty,
\end{equation}
so \(\WrightOmega_N\) approximates \(\WrightOmega\) to order exactly
\(L^{N+1}x^{-(N+1)}\), and
\(X+\sum_{n\ge0}\LamP{n}(L)t^{n}\) is an asymptotic expansion of
\(\WrightOmega\) on the scale \(\SetOf{(\log x)^{m}x^{-n}}\) in the sense of
\cref{plt:def:mul-asymptotic}.
\end{theorem}

\begin{proof}
Write \(t=1/x\), \(L=\log x\), and
\(p=p(x)\DefinitionEquals-L+\sum_{n=1}^{N}\LamP{n}(L)x^{-n}\), so
\(\WrightOmega_N(x)=x+p=x(1+\varepsilon)\) with \(\varepsilon=p/x\).

\emph{Step 1: \(\varepsilon\) is small.}  For \(x\ge\EulerE\) we have
\(L\ge1\), and \(\AbsoluteValue{\LamP{n}(L)}\le c_nL^{n}\) with
\(c_n=\sum_{m}\AbsoluteValue{[u^{m}]\LamP{n}}\).  Hence, once \(x\ge L\),
\(\sum_{n=1}^{N}\AbsoluteValue{\LamP{n}(L)}x^{-n}\le\sum_{n=1}^{N}c_n(L/x)^{n}\le C_N\,L/x\).
Therefore \(\AbsoluteValue{p}\le L+C_NL/x\le2L\) and
\(\AbsoluteValue{\varepsilon}\le2L/x\) for all \(x\) beyond some \(x_N\),
which we enlarge so that \(\AbsoluteValue{\varepsilon}\le\tfrac12\).  In
particular \(\WrightOmega_N(x)\ge x/2>0\), settling item~(4) above, and
\(\log\WrightOmega_N(x)=L+\log(1+\varepsilon)\) with the real logarithm.

\emph{Step 2: an analytic remainder for the logarithm.}  For
\(\AbsoluteValue{\varepsilon}\le\tfrac12\),
\[
 \log(1+\varepsilon)
 =\sum_{j=1}^{N+1}\frac{(-1)^{j-1}}{j}\varepsilon^{j}+\theta,
 \qquad
 \AbsoluteValue{\theta}
 \le\sum_{j>N+1}\frac{\AbsoluteValue{\varepsilon}^{j}}{j}
 \le2\AbsoluteValue{\varepsilon}^{\,N+2}
 \le2\left(\frac{2L}{x}\right)^{N+2}.
\]
Consequently
\begin{equation}\label{plt:eq:om-Pi}
 \WrightOmega_N(x)+\log\WrightOmega_N(x)-x
 =\Pi(1/x,\log x)+\theta,
 \qquad
 \Pi\DefinitionEquals
 \mathcal P_N+L+\sum_{j=1}^{N+1}\frac{(-1)^{j-1}}{j}\bigl(t\mathcal P_N\bigr)^{j},
\end{equation}
where \(\Pi\in\K[L][t]\) is a \emph{polynomial} in \(t\) and \(L\), evaluated
at the real numbers \(1/x\) and \(\log x\).

\emph{Step 3: identifying \(\Pi\).}  In \(\K[L][[t]]\) we have
\(\ordt(t\mathcal P_N)\ge1\), so
\(\log(1+t\mathcal P_N)-\sum_{j\le N+1}\frac{(-1)^{j-1}}{j}(t\mathcal P_N)^{j}
\in t^{\,N+2}\K[L][[t]]\).  Combining this with
\cref{plt:thm:om-residual},
\[
 \Pi\equiv-\LamP{N+1}(L)\,t^{\,N+1}\pmod{t^{\,N+2}\K[L][t]} .
\]
Now \(\mathcal P_N\) has \(t\)-degree \(\le N\) and \(L\)-degree
\(\le\max(N,1)\), so \(\Pi\) has \(t\)-degree \(\le(N+1)^{2}\) and
\(L\)-degree \(\le(N+1)\max(N,1)\le K\).  Writing
\(\Pi=-\LamP{N+1}(L)t^{N+1}+\sum_{i=N+2}^{(N+1)^{2}}c_i(L)t^{i}\) with
\(\deg c_i\le K\), we obtain, for \(x\ge x_N\),
\[
 \Bigl|\Pi(1/x,\log x)+\LamP{N+1}(L)x^{-(N+1)}\Bigr|
 \le C_N'\,\frac{L^{K}}{x^{\,N+2}} .
\]
Since also \(\AbsoluteValue{\theta}\le2^{N+3}L^{N+2}x^{-(N+2)}\) and
\(N+2\le K\), \eqref{plt:eq:om-analytic-residual} follows from
\eqref{plt:eq:om-Pi}.

\emph{Step 4: residual to error.}  Denote the left side of
\eqref{plt:eq:om-analytic-residual} by \(R_N(x)\), so
\(\phi(\WrightOmega_N(x))-\phi(\WrightOmega(x))=R_N(x)\) with
\(\phi(Y)=Y+\log Y\).  Both \(\WrightOmega_N(x)\) and \(\WrightOmega(x)\) lie
in \((0,\infty)\), and by \cref{plt:eq:om-envelope} and Step~1 both exceed
\(x/2\) for \(x\ge x_N\) (enlarging \(x_N\) again).  By the mean value
theorem there is \(\zeta\) between them, hence \(\zeta>x/2\), with
\[
 \WrightOmega_N(x)-\WrightOmega(x)=\frac{R_N(x)}{\phi'(\zeta)}
 =\frac{R_N(x)}{1+\zeta^{-1}},
 \qquad
 \frac{1}{1+\zeta^{-1}}=1+\BigOAt{x^{-1}}{x\to+\infty}.
\]
This is the general residual-transfer principle of
\cref{plt:prop:rev-analytic-transfer} in the concrete case at hand, with
derivative lower bound \(m=1\).  Multiplying
\eqref{plt:eq:om-analytic-residual} by \(-(1+\zeta^{-1})^{-1}\) and absorbing
\(\AbsoluteValue{\LamP{N+1}(L)}x^{-(N+1)}\cdot\BigOAt{x^{-1}}{x\to+\infty}
=\BigOAt{L^{N+1}x^{-(N+2)}}{x\to+\infty}\) into the remainder gives
\eqref{plt:eq:om-analytic-error}.

\emph{Step 5.}  By \cref{plt:cor:lambert-coefficients},
\(\LamP{N+1}(L)=L^{N+1}/(N+1)+\BigOAt{L^{N}}{L\to+\infty}\), and
\(L^{K}x^{-(N+2)}=\LittleOAt{L^{N+1}x^{-(N+1)}}{x\to+\infty}\) because
\(L^{K-N-1}/x\to0\).  Hence \eqref{plt:eq:om-sharp-order}.  Since
\eqref{plt:eq:om-sharp-order} holds for every \(N\), the partial sums are an
asymptotic expansion on the stated scale.
\end{proof}

\begin{remark}[The rate in \eqref{plt:eq:om-sharp-order} is only logarithmic]
\label{plt:rmk:om-slow}
Statement \eqref{plt:eq:om-analytic-error} is the useful one:
the error of \(\WrightOmega_N\) equals the \emph{next block} with a relative
error \(\BigOAt{L^{K}/x}{x\to+\infty}\).  The simplified form
\eqref{plt:eq:om-sharp-order} discards the lower-degree part of
\(\LamP{N+1}\), whose largest term is \(-H_{N}L^{N}\); the relative error of
\eqref{plt:eq:om-sharp-order} is therefore only \(\BigOAt{1/L}{L\to\infty}\),
i.e.\ \(\BigOAt{1/\log x}{x\to+\infty}\).  For \(N=4\) and \(x=5\cdot10^{5}\)
the ratio of the true error to \(\LamP{5}(L)x^{-5}\) is
\(1.0000148\ldots\), while its ratio to \(L^{5}/(5x^{5})\) is still far from
\(1\).  Quoting \eqref{plt:eq:om-sharp-order} where
\eqref{plt:eq:om-analytic-error} is meant is a standard way to make a correct
theorem look numerically wrong.
\end{remark}

The next proposition isolates the one convergence statement that the formal
theory does yield, and makes clear why it is not the statement one wants.

\begin{proposition}[Convergence at frozen logarithm]
\label{plt:prop:om-frozen-convergence}
For each \(u_0\in\ComplexNumbers\) there are \(r>0\) and \(\rho>0\) such that
the series \(\sum_{n\ge1}\LamP{n}(u)s^{n}\) converges absolutely for
\(\AbsoluteValue{s}<r\), uniformly for
\(\AbsoluteValue{u-u_0}\le\rho\), and its sum \(Q(s,u)\) is the unique
holomorphic solution near \((0,u_0)\) of
\begin{equation}\label{plt:eq:om-frozen}
 Q=-\log(1-su+sQ),\qquad Q(0,u)=0 .
\end{equation}
\end{proposition}

\begin{proof}
Let \(\Phi(s,u,Q)=Q+\PrincipalLogarithm(1-su+sQ)\), holomorphic on a
neighbourhood of \(\{(0,u,0):u\in\ComplexNumbers\}\) because
\(1-su+sQ=1\) there.  We have \(\Phi(0,u,0)=0\) and
\(\partial_Q\Phi(0,u,0)=1+\frac{s}{1-su+sQ}\big|_{s=0}=1\ne0\).  The
holomorphic implicit function theorem gives \(r,\rho>0\) and a unique
holomorphic \(Q\) on \(\AbsoluteValue{s}<r\),
\(\AbsoluteValue{u-u_0}<\rho\), with \(Q(0,u)=0\) and \(\Phi=0\).  Its Taylor
coefficients in \(s\) satisfy the same triangular recursion as those of the
formal solution of \eqref{plt:eq:om-frozen}, which by
\cref{plt:thm:om-generating-equation} is
\(\sum_{n\ge1}\LamP{n}(u)s^{n}\); by uniqueness of Taylor coefficients the
two agree.
\end{proof}

\begin{remark}[What is \emph{not} proved here]
\label{plt:rmk:om-open}
\Cref{plt:prop:om-frozen-convergence} freezes \(u\) and lets \(s\to0\).  The
asymptotic problem does the opposite: it couples them by \(u=\log(1/s)\), so
that \(u\to\infty\) as \(s\to0\), and \cref{plt:prop:om-frozen-convergence}
says nothing about it, because \(r\) is not claimed to be bounded below as
\(\AbsoluteValue{u_0}\to\infty\).  In fact the implicit function theorem
fails on the locus \(\partial_Q\Phi=0\), which after eliminating \(Q\) from
\eqref{plt:eq:om-frozen} is \(u=s^{-1}+1-\log(-s)\); this suggests
\(r\asymp1/\AbsoluteValue{u}\) for large \(\AbsoluteValue{u}\), which would
be compatible with genuine convergence of
\(\sum_n\LamP{n}(\log x)x^{-n}\) for large \(x\).  We record that as a
heuristic only.  The literature asserts absolute convergence of the
associated double series on explicit regions
\cite{corless-et-al-1996,jeffrey-et-al-1995,dlmf-W}; no proof of any such
convergence statement is given here, and none follows from anything in
Parts~I--IV.  Likewise, uniformity of \eqref{plt:eq:om-analytic-error} in a
complex sector, in the branch index \(k\) of \cref{plt:rmk:om-branches}, or
in \(N\), is neither claimed nor proved.
\end{remark}


\chapter{The Lambert polynomials}
\label{plt:sec:lambert-polynomials}

This chapter is self-contained: the polynomials are defined by an operator
formula, and every property --- the recurrence, the closed Stirling form, the
coefficient structure, the integrality, the generating equation, the Laplace
transform --- is derived from that definition without using
\cref{plt:sec:wright-omega}.  The two sections meet in
\cref{plt:thm:om-generating-equation}, which reproves
\cref{plt:thm:om-expansion} without the reversion machinery.

Throughout, \(u\) is an indeterminate, \(\partial_u\) is
\(\Differential/\Differential u\), and \(\shiftop{0}{n}=\prod_{j=0}^{n-1}(\partial_u-j)\)
is the falling-factorial operator \(\partial_u^{\underline n}\), with the
empty product at \(n=0\) acting as the identity.

\section{Definition and first values}

\begin{definition}[Canonical zero-based Lambert polynomials]
\label{plt:def:lambert-polynomials}
For \(n\in\NaturalNumbers\) put
\begin{equation}\label{plt:eq:lambert-operator}
 \LamP{n}(u)
 \DefinitionEquals
 \frac{(-1)^{n+1}}{(n+1)!}
 \left[\prod_{j=0}^{n-1}\bigl(\partial_u-j\bigr)\right]u^{\,n+1}
 \;\in\;\RationalNumbers[u],
\end{equation}
the product being empty at \(n=0\).  Since \(\operatorname{char}\K=0\) we
regard \(\LamP{n}\in\K[u]\) whenever convenient.
\end{definition}

At \(n=0\) the operator is the identity and the prefactor is
\((-1)^{1}/1!=-1\), so \(\LamP{0}(u)=-u\).  The next values are
\begin{equation}\label{plt:eq:lambert-first-values}
 \LamP{1}(u)=u,\qquad
 \LamP{2}(u)=\frac{u^{2}}{2}-u,\qquad
 \LamP{3}(u)=\frac{u^{3}}{3}-\frac32u^{2}+u,
\end{equation}
computed as follows.  For \(n=1\),
\(\partial_uu^{2}=2u\) and \((-1)^{2}/2!=\tfrac12\), giving \(u\).  For
\(n=2\), \((\partial_u-1)\partial_uu^{3}=(\partial_u-1)(3u^{2})=6u-3u^{2}\)
and \((-1)^{3}/3!=-\tfrac16\), giving \(\tfrac12u^{2}-u\).  For \(n=3\),
\((\partial_u-2)(\partial_u-1)\partial_uu^{4}
=(\partial_u-2)(\partial_u-1)(4u^{3})
=(\partial_u-2)(12u^{2}-4u^{3})
=24u-12u^{2}-24u^{2}+8u^{3}\), and \((-1)^{4}/4!=\tfrac1{24}\), giving
\(\tfrac13u^{3}-\tfrac32u^{2}+u\).

\begin{warningbox}
The degree law \(\deg\LamP{n}=n\) is asserted for \(n\ge1\) only.  The
zero-index member \(\LamP{0}(u)=-u\) has degree \(1\).  Any statement
quantified over \(n\ge0\) that mentions the degree, the leading coefficient,
or a coefficient indexed relative to the top must exclude \(n=0\) explicitly.
\end{warningbox}

\section{The recurrence}

\begin{theorem}[Differential and integral recurrence]
\label{plt:thm:lambert-recurrence}
For every \(n\in\NaturalNumbers\),
\begin{equation}\label{plt:eq:lambert-boundary}
 \LamP{n}(0)=0,
\end{equation}
\begin{equation}\label{plt:eq:lambert-diff-rec}
 \frac{\Differential}{\Differential u}\LamP{n+1}(u)
 =n\,\LamP{n}(u)-\frac{\Differential}{\Differential u}\LamP{n}(u),
\end{equation}
\begin{equation}\label{plt:eq:lambert-int-rec}
 \LamP{n+1}(u)=-\LamP{n}(u)+n\int_{0}^{u}\LamP{n}(v)\,\Differential v .
\end{equation}
Moreover, \eqref{plt:eq:lambert-int-rec} for all \(n\ge0\) is equivalent to
the conjunction of \eqref{plt:eq:lambert-diff-rec} for all \(n\ge0\) with
\eqref{plt:eq:lambert-boundary} for all \(n\ge0\); and either version,
together with \(\LamP{0}(u)=-u\), determines the whole family uniquely.
\end{theorem}

\begin{proof}
\emph{Boundary value.}  Expanding the operator in
\eqref{plt:eq:lambert-operator} as a \(\RationalNumbers\)-linear combination
of \(\partial_u^{k}\), \(0\le k\le n\), and applying it to \(u^{n+1}\)
produces a linear combination of the monomials \(u^{\,n+1-k}\) with
\(n+1-k\ge1\).  Hence \(\LamP{n}\) has zero constant term, which is
\eqref{plt:eq:lambert-boundary}.

\emph{Differential recurrence.}  Write \(A_n=\shiftop{0}{n}\), so
\(A_{n+1}=(\partial_u-n)A_n\) and \(A_n\) commutes with \(\partial_u\), all
factors being polynomials in \(\partial_u\).  From
\eqref{plt:eq:lambert-operator},
\[
 \LamP{n+1}(u)=\frac{(-1)^{n}}{(n+2)!}\,(\partial_u-n)A_n\,u^{\,n+2}.
\]
Differentiate and commute \(\partial_u\) past \((\partial_u-n)A_n\):
\[
 \frac{\Differential}{\Differential u}\LamP{n+1}(u)
 =\frac{(-1)^{n}}{(n+2)!}(\partial_u-n)A_n\,\partial_uu^{\,n+2}
 =\frac{(-1)^{n}(n+2)}{(n+2)!}(\partial_u-n)A_n\,u^{\,n+1}
 =\frac{(-1)^{n}}{(n+1)!}(\partial_u-n)A_nu^{\,n+1}.
\]
By \eqref{plt:eq:lambert-operator} again,
\(A_nu^{\,n+1}=(-1)^{n+1}(n+1)!\,\LamP{n}(u)\), so
\[
 \frac{\Differential}{\Differential u}\LamP{n+1}(u)
 =\frac{(-1)^{n}(-1)^{n+1}(n+1)!}{(n+1)!}\,(\partial_u-n)\LamP{n}(u)
 =-(\partial_u-n)\LamP{n}(u),
\]
which is \eqref{plt:eq:lambert-diff-rec}.

\emph{Equivalence of the two forms.}  Given
\eqref{plt:eq:lambert-diff-rec} and \eqref{plt:eq:lambert-boundary},
integrate \eqref{plt:eq:lambert-diff-rec} from \(0\) to \(u\) and use
\(\LamP{n+1}(0)=\LamP{n}(0)=0\) to get
\eqref{plt:eq:lambert-int-rec}.  Conversely, differentiating
\eqref{plt:eq:lambert-int-rec} returns \eqref{plt:eq:lambert-diff-rec}, and
evaluating \eqref{plt:eq:lambert-int-rec} at \(u=0\) gives
\(\LamP{n+1}(0)=-\LamP{n}(0)\); since \(\LamP{0}(0)=0\), induction yields
\eqref{plt:eq:lambert-boundary} for all \(n\).

\emph{Uniqueness.}  \eqref{plt:eq:lambert-int-rec} computes \(\LamP{n+1}\)
from \(\LamP{n}\) by operations with no free constant, so the family is
determined by \(\LamP{0}=-u\).
\end{proof}

\begin{remark}[Reconciliation with the one-based drafts]
\label{plt:rmk:om-onebased}
% ed.: S3 and S6 start the recurrence at P_1 = u and impose P_{n}(0)=0 as a
% side condition of the theorem rather than proving it; S4 derives the linear
% recurrence by "equating coefficients" in a NONLINEAR generating ODE, which
% does not produce it.  The proof above is direct from the definition and
% proves the boundary condition instead of assuming it.
Sources~3 and~6 state the recurrence as starting from \(\LamP{1}(u)=u\) with
\(\LamP{n}(0)=0\) imposed.  The zero-based normalisation subsumes this: with
\(\LamP{0}=-u\), \eqref{plt:eq:lambert-int-rec} at \(n=0\) reads
\(\LamP{1}=-\LamP{0}+0=u\).  Source~4 derives
\eqref{plt:eq:lambert-diff-rec} by ``equating coefficients'' in the
first-order equation
\(T\mathcal L\,\partial_T\mathcal L+(1+T)\partial_T\mathcal L+\mathcal L=0\)
satisfied by the generating series; that equation is \emph{quadratic} in
\(\mathcal L\) and equating coefficients in it yields a convolution
recurrence, not the two-term linear one.  The correct generating-function
route runs through the \emph{linear} equation
\eqref{plt:eq:om-pde} below, and is given in
\cref{plt:thm:om-generating-equation}.
\end{remark}

\section{The closed Stirling form}

We use the signed and unsigned Stirling numbers of the first kind in the
catalogue normalisation:
\begin{equation}\label{plt:eq:om-stirling-def}
 \FallingFactorial{z}{n}=z(z-1)\cdots(z-n+1)=\sum_{k=0}^{n}\stirlingsigned{n}{k}z^{k},
 \qquad
 \stirlingone{n}{k}=(-1)^{n-k}\stirlingsigned{n}{k},
\end{equation}
with \(\stirlingsigned{0}{0}=1\) and \(\stirlingsigned{n}{k}=0\) for
\(k<0\) or \(k>n\).

\begin{theorem}[Stirling-number formula]
\label{plt:thm:lambert-stirling}
For every \(n\ge1\),
\begin{equation}\label{plt:eq:lambert-stirling}
 \LamP{n}(u)=\sum_{m=1}^{n}(-1)^{n-m}\stirlingone{n}{n-m+1}\frac{u^{m}}{m!}.
\end{equation}
Equivalently, writing \(\LamP{n}(u)=\sum_{m\ge0}\lambda_{n,m}u^{m}\),
\begin{equation}\label{plt:eq:lambert-lambda}
 \lambda_{n,m}=\frac{(-1)^{n-m}}{m!}\stirlingone{n}{n-m+1}
 =\frac{(-1)^{n+1}}{m!}\stirlingsigned{n}{n-m+1},
 \qquad 1\le m\le n,
\end{equation}
and \(\lambda_{n,m}=0\) for \(m=0\) and for \(m>n\).  If the upper limit in
\eqref{plt:eq:lambert-stirling} is raised from \(n\) to \(n+1\), the formula
holds for \(n=0\) as well and returns \(\LamP{0}(u)=-u\).
\end{theorem}

\begin{proof}
Substituting \(z\mapsto\partial_u\) in \eqref{plt:eq:om-stirling-def} --- a
substitution of a commuting operator into a polynomial identity, hence valid
--- gives
\(\shiftop{0}{n}=\sum_{k=0}^{n}\stirlingsigned{n}{k}\partial_u^{k}\).
Since \(\partial_u^{k}u^{\,n+1}=\frac{(n+1)!}{(n+1-k)!}u^{\,n+1-k}\) for
\(0\le k\le n+1\), \eqref{plt:eq:lambert-operator} gives
\[
 \LamP{n}(u)
 =\frac{(-1)^{n+1}}{(n+1)!}\sum_{k=0}^{n}\stirlingsigned{n}{k}
  \frac{(n+1)!}{(n+1-k)!}u^{\,n+1-k}
 =(-1)^{n+1}\sum_{k=0}^{n}\stirlingsigned{n}{k}\frac{u^{\,n+1-k}}{(n+1-k)!} .
\]
Put \(m=n+1-k\), so \(k=n-m+1\) and \(m\) runs from \(1\) to \(n+1\):
\[
 \LamP{n}(u)=(-1)^{n+1}\sum_{m=1}^{n+1}\stirlingsigned{n}{n-m+1}\frac{u^{m}}{m!}.
\]
This is the second equality of \eqref{plt:eq:lambert-lambda}, and the first
follows from
\(\stirlingsigned{n}{n-m+1}=(-1)^{n-(n-m+1)}\stirlingone{n}{n-m+1}
=(-1)^{m-1}\stirlingone{n}{n-m+1}\),
so that \((-1)^{n+1}(-1)^{m-1}=(-1)^{n+m}=(-1)^{n-m}\).  For \(n\ge1\), the
term \(m=n+1\) carries \(\stirlingsigned{n}{0}=0\) and may be dropped,
giving \eqref{plt:eq:lambert-stirling}.  For \(n=0\) the only term is
\(m=1\), with \(\stirlingsigned{0}{0}=1\) and prefactor \((-1)^{1}\), giving
\(-u\).
\end{proof}

\begin{remark}[A second proof, and where the Stirling numbers come from]
\label{plt:rmk:om-stirling-second}
Formula \eqref{plt:eq:lambert-stirling} can also be obtained from
\cref{plt:thm:lambert-recurrence} alone, and that route explains the
appearance of Stirling numbers structurally.  Put
\(\sigma_{n,m}\DefinitionEquals m!\,\lambda_{n,m}\).  The integral recurrence
\eqref{plt:eq:lambert-int-rec} reads, coefficientwise,
\(\lambda_{n+1,m}=-\lambda_{n,m}+\frac{n}{m}\lambda_{n,m-1}\), which in the
\(\sigma\) normalisation becomes
\begin{equation}\label{plt:eq:om-sigma-rec}
 \sigma_{n+1,m}=-\sigma_{n,m}+n\,\sigma_{n,m-1},
 \qquad \sigma_{0,1}=-1 .
\end{equation}
Setting \(\sigma_{n,m}=(-1)^{n-m}\stirlingone{n}{n-m+1}\) and \(k=n-m+2\)
turns \eqref{plt:eq:om-sigma-rec} into
\(\stirlingone{n+1}{k}=\stirlingone{n}{k-1}+n\stirlingone{n}{k}\), the
defining recurrence of the unsigned Stirling numbers.  Thus the Lambert
recurrence \emph{is} the Stirling recurrence, transported by the
normalisation \(u^{m}\mapsto u^{m}/m!\).  Source~4 uses this induction; it is
valid, but it presupposes the recurrence, whereas the proof above needs only
the definition.
\end{remark}

\section{Coefficient structure}

The four ``edge'' coefficients below all come from a single generating
polynomial, so we prove the Stirling evaluations rather than quote them.

\begin{lemma}[Four Stirling evaluations]
\label{plt:lem:om-stirling-values}
For \(n\ge1\) let \(H_r=\sum_{j=1}^{r}j^{-1}\) and
\(H_r^{(2)}=\sum_{j=1}^{r}j^{-2}\), with \(H_0=H_0^{(2)}=0\).  Then
\[
 \stirlingone{n}{1}=(n-1)!,\quad
 \stirlingone{n}{2}=(n-1)!\,H_{n-1},\quad
 \stirlingone{n}{3}=\frac{(n-1)!}{2}\Bigl(H_{n-1}^{2}-H_{n-1}^{(2)}\Bigr),
\]
and for \(n\ge2\), \(\stirlingone{n}{n-1}=\binom{n}{2}\), while
\(\stirlingone{n}{n}=1\).
\end{lemma}

\begin{proof}
Replacing \(z\) by \(-z\) in \eqref{plt:eq:om-stirling-def} gives the
unsigned form \(\RisingFactorial{z}{n}=\sum_{k}\stirlingone{n}{k}z^{k}\),
where \(\RisingFactorial{z}{n}=z(z+1)\cdots(z+n-1)\).  Since
\(\stirlingone{n}{0}=0\) for \(n\ge1\), divide by \(z\) and set
\[
 f_n(z)\DefinitionEquals\prod_{j=1}^{n-1}(z+j)=\sum_{k\ge1}\stirlingone{n}{k}z^{\,k-1}.
\]
Then \(\stirlingone{n}{1}=f_n(0)=(n-1)!\); \(f_n\) is monic of degree
\(n-1\), so \(\stirlingone{n}{n}=1\); and the coefficient of \(z^{\,n-2}\) in
\(f_n\) is \(\sum_{j=1}^{n-1}j=\binom{n}{2}\), i.e.\
\(\stirlingone{n}{n-1}=\binom{n}{2}\).  For the remaining two, logarithmic
differentiation gives
\(f_n'=f_n\sum_{j=1}^{n-1}(z+j)^{-1}\) and
\(f_n''=f_n\bigl[(\sum_j(z+j)^{-1})^{2}-\sum_j(z+j)^{-2}\bigr]\); evaluating
at \(z=0\),
\(\stirlingone{n}{2}=f_n'(0)=(n-1)!H_{n-1}\) and
\(\stirlingone{n}{3}=\tfrac12f_n''(0)=\tfrac12(n-1)!(H_{n-1}^{2}-H_{n-1}^{(2)})\).
\end{proof}

\begin{corollary}[Coefficient structure of \(\LamP{n}\)]
\label{plt:cor:lambert-coefficients}
Let \(n\ge1\).  Then:
\begin{enumerate}[label=\textup{(\roman*)}]
\item \(\deg\LamP{n}=n\), the constant term is \(0\), and \(u\mid\LamP{n}\);
\item all \(n\) coefficients \(\lambda_{n,1},\ldots,\lambda_{n,n}\) are
      nonzero and strictly alternate: \(\Sign\lambda_{n,m}=(-1)^{n-m}\);
\item \([u^{n}]\LamP{n}=1/n\);
\item \([u^{n-1}]\LamP{n}=-H_{n-1}\) for \(n\ge2\);
\item \([u]\LamP{n}=(-1)^{n-1}\), equivalently
      \({\LamP{n}}'(0)=(-1)^{n-1}\);
\item \([u^{n-2}]\LamP{n}=\dfrac{n-1}{2}\Bigl(H_{n-1}^{2}-H_{n-1}^{(2)}\Bigr)\)
      for \(n\ge3\);
\item \([u^{2}]\LamP{n}=(-1)^{n}\dfrac{n(n-1)}{4}\) for \(n\ge2\).
\end{enumerate}
In particular
\(\LamP{n}(u)=\dfrac{u^{n}}{n}-H_{n-1}u^{\,n-1}+\cdots+(-1)^{n-1}u\).
\end{corollary}

\begin{proof}
All statements read off \eqref{plt:eq:lambert-lambda} together with
\cref{plt:lem:om-stirling-values}.  (i) and (ii): for \(1\le m\le n\) the
index \(k=n-m+1\) satisfies \(1\le k\le n\), so \(\stirlingone{n}{k}>0\)
(it counts the permutations of \(n\) letters with \(k\) cycles, and there is
at least one); hence \(\lambda_{n,m}\ne0\) with sign \((-1)^{n-m}\), and
\(\lambda_{n,m}=0\) outside \(1\le m\le n\).  (iii): \(m=n\) gives
\(k=1\) and \(\lambda_{n,n}=\stirlingone{n}{1}/n!=(n-1)!/n!=1/n\).  (iv):
\(m=n-1\) gives \(k=2\) and
\(\lambda_{n,n-1}=-\stirlingone{n}{2}/(n-1)!=-H_{n-1}\).  (v): \(m=1\) gives
\(k=n\) and \(\lambda_{n,1}=(-1)^{n-1}\stirlingone{n}{n}=(-1)^{n-1}\).
(vi): \(m=n-2\) gives \(k=3\) and
\(\lambda_{n,n-2}=\stirlingone{n}{3}/(n-2)!
=\frac{(n-1)!}{2(n-2)!}(H_{n-1}^{2}-H_{n-1}^{(2)})\).
(vii): \(m=2\) gives \(k=n-1\) and
\(\lambda_{n,2}=(-1)^{n-2}\binom{n}{2}/2!=(-1)^{n}n(n-1)/4\).
\end{proof}

\begin{remark}[The formulas overlap consistently]
\label{plt:rmk:om-overlap}
At small \(n\) several of the seven statements describe the same
coefficient, and they agree.  At \(n=3\), (iv) and (vii) both give
\([u^{2}]\LamP{3}=-\tfrac32\), and (v) and (vi) both give
\([u]\LamP{3}=1\).  At \(n=4\), (vi) and (vii) both give
\([u^{2}]\LamP{4}=3\), since
\(\tfrac32(H_3^{2}-H_3^{(2)})=\tfrac32\bigl(\tfrac{121}{36}-\tfrac{49}{36}\bigr)=3\).
Formula (vi) also returns the correct value \(0\) at \(n=2\), where
\(H_1^{2}-H_1^{(2)}=0\) and \([u^{0}]\LamP{2}=0\); it is stated for
\(n\ge3\) only because \(u^{\,n-2}\) is not a monomial for \(n<2\).
\end{remark}

\section{Integrality and the scaled family}

\begin{corollary}[Integrality]
\label{plt:cor:lambert-integrality}
For \(n\ge1\) and \(0\le m\le n\),
\begin{equation}\label{plt:eq:om-sigma-integral}
 \sigma_{n,m}\DefinitionEquals m!\,[u^{m}]\LamP{n}(u)
 =(-1)^{n-m}\stirlingone{n}{n-m+1}\in\IntegerNumbers .
\end{equation}
Consequently the scaled family
\begin{equation}\label{plt:eq:om-scaled-def}
 \LamPsc{n}(u)\DefinitionEquals n!\,\LamP{n}(u)
 =\sum_{m=1}^{n}\frac{n!}{m!}\,\sigma_{n,m}u^{m}
\end{equation}
lies in \(\IntegerNumbers[u]\), is divisible by \(u\), has degree \(n\),
leading coefficient \((n-1)!\) and linear coefficient
\((-1)^{n-1}n!\), and satisfies
\begin{equation}\label{plt:eq:om-scaled-rec}
 \LamPsc{n+1}(u)=-(n+1)\LamPsc{n}(u)+n(n+1)\int_{0}^{u}\LamPsc{n}(v)\,\Differential v .
\end{equation}
The convention \(\LamPsc{0}(u)=0!\,\LamP{0}(u)=-u\) makes
\eqref{plt:eq:om-scaled-rec} valid at \(n=0\) as well.
\end{corollary}

\begin{proof}
\eqref{plt:eq:om-sigma-integral} is \eqref{plt:eq:lambert-lambda} rewritten,
and \(\stirlingone{n}{k}\in\IntegerNumbers\).  Then
\(n!\lambda_{n,m}=\frac{n!}{m!}\sigma_{n,m}\) is an integer because
\(m\le n\).  The listed data are \cref{plt:cor:lambert-coefficients}(i),
(iii), (v) multiplied by \(n!\).  Multiplying
\eqref{plt:eq:lambert-int-rec} by \((n+1)!\) and using
\((n+1)!=(n+1)\cdot n!\) gives \eqref{plt:eq:om-scaled-rec}; at \(n=0\) both
sides equal \(u\).
\end{proof}

\begin{warningbox}
\eqref{plt:eq:om-scaled-rec} is an identity in \(\RationalNumbers[u]\) and
does \emph{not} manifestly preserve \(\IntegerNumbers[u]\): the
antiderivative divides the coefficient of \(u^{m}\) by \(m+1\), and
\(n(n+1)/(m+1)\) need not be an integer.  Integrality of \(\LamPsc{n}\) is a
theorem (\cref{plt:cor:lambert-integrality}), proved through the Stirling
form, not a visible feature of the recurrence.  An implementation that wants
to stay inside \(\IntegerNumbers\) should iterate
\eqref{plt:eq:om-sigma-rec} on the integers \(\sigma_{n,m}\) --- equivalently,
build the Stirling triangle --- and divide by \(m!\) only at output time.
\end{warningbox}

The first scaled polynomials are
\begin{align}
 \LamPsc{0}(u)&=-u, & \LamPsc{1}(u)&=u, \notag\\
 \LamPsc{2}(u)&=u^{2}-2u, & \LamPsc{3}(u)&=2u^{3}-9u^{2}+6u, \notag\\
 \LamPsc{4}(u)&=6u^{4}-44u^{3}+72u^{2}-24u,
 & \LamPsc{5}(u)&=24u^{5}-250u^{4}+700u^{3}-600u^{2}+120u,
 \label{plt:eq:om-scaled-values}
\end{align}
\begin{equation}\label{plt:eq:om-scaled-six}
 \LamPsc{6}(u)=120u^{6}-1644u^{5}+6750u^{4}-10200u^{3}+5400u^{2}-720u .
\end{equation}
Further values, and the corresponding \(\LamP{n}\), are tabulated in
\cref{plt:sec:lambert-tables}.

\section{The generating equation}

\begin{lemma}[The linear generating equation]
\label{plt:lem:om-pde}
Let
\(\mathcal P(s,u)\DefinitionEquals\sum_{n\ge0}\LamP{n}(u)s^{n}\in\K[u][[s]]\).
Then
\begin{equation}\label{plt:eq:om-pde}
 s^{2}\,\partial_s\mathcal P=(1+s)\,\partial_u\mathcal P+1 .
\end{equation}
\end{lemma}

\begin{proof}
Multiply \eqref{plt:eq:lambert-diff-rec} by \(s^{\,n+1}\) and sum over
\(n\ge0\).  The right-hand side gives
\(\sum_{n\ge0}n\LamP{n}(u)s^{\,n+1}=s\cdot s\,\partial_s\mathcal P
=s^{2}\partial_s\mathcal P\).  The left-hand side gives
\[
 \sum_{n\ge0}\Bigl({\LamP{n+1}}'(u)+{\LamP{n}}'(u)\Bigr)s^{\,n+1}
 =\bigl(\partial_u\mathcal P-{\LamP{0}}'(u)\bigr)+s\,\partial_u\mathcal P
 =(1+s)\partial_u\mathcal P+1,
\]
using \({\LamP{0}}'(u)=-1\).
\end{proof}

\begin{theorem}[Implicit generating equation and formal uniqueness]
\label{plt:thm:om-generating-equation}
Let \(Q(s,u)\DefinitionEquals\sum_{n\ge1}\LamP{n}(u)s^{n}\in\K[u][[s]]\).
Then \(Q(0,u)=0\) and
\begin{equation}\label{plt:eq:om-generating}
 Q=-\log\bigl(1-su+sQ\bigr),
\end{equation}
the logarithm being the formal one on \(1+s\K[u][[s]]\).  Conversely,
\eqref{plt:eq:om-generating} together with \(Q(0,u)=0\) has exactly one
solution in \(\K[u][[s]]\); hence it determines the family
\((\LamP{n})_{n\ge1}\), and therefore also \(\LamP{0}=-u\) through
\eqref{plt:eq:lambert-int-rec}, uniquely.
\end{theorem}

\begin{proof}
Write \(\mathcal P=-u+Q\) as in \cref{plt:lem:om-pde}.  Since
\(s\mathcal P\in s\K[u][[s]]\), the element
\(E\DefinitionEquals\mathcal P+u+\log(1+s\mathcal P)\in\K[u][[s]]\) is
defined, and \eqref{plt:eq:om-generating} is exactly \(E=0\).

\emph{Step 1: \(E\) satisfies the homogeneous equation.}  Put
\(D=1+s\mathcal P\), a unit of \(\K[u][[s]]\).  Then
\[
 \partial_uE=\partial_u\mathcal P+1+\frac{s\,\partial_u\mathcal P}{D},
 \qquad
 \partial_sE=\partial_s\mathcal P+\frac{\mathcal P+s\,\partial_s\mathcal P}{D}.
\]
Using \(s^{2}\partial_s\mathcal P=(1+s)\partial_u\mathcal P+1\) from
\eqref{plt:eq:om-pde}, hence
\(s^{3}\partial_s\mathcal P=s(1+s)\partial_u\mathcal P+s\), we get
\[
 s^{2}\partial_sE-(1+s)\partial_uE
 =\bigl[(1+s)\partial_u\mathcal P+1\bigr]-(1+s)\partial_u\mathcal P-(1+s)
  +\frac{s^{2}\mathcal P+s(1+s)\partial_u\mathcal P+s
        -s(1+s)\partial_u\mathcal P}{D}.
\]
The fraction is \(\dfrac{s^{2}\mathcal P+s}{D}=\dfrac{s(1+s\mathcal P)}{D}=s\),
so the whole expression equals \(1-(1+s)+s=0\).  Thus
\begin{equation}\label{plt:eq:om-E-pde}
 s^{2}\partial_sE=(1+s)\partial_uE .
\end{equation}

\emph{Step 2: \(E\) vanishes on \(u=0\).}  By
\eqref{plt:eq:lambert-boundary}, \(\mathcal P(s,0)=\sum_n\LamP{n}(0)s^{n}=0\).
Applying the continuous \(\K\)-algebra map \(u\mapsto0\) to \(E\) gives
\(E(s,0)=0+0+\log(1+0)=0\).

\emph{Step 3.}  Write \(E=\sum_{n\ge0}e_n(u)s^{n}\).  Comparing coefficients
of \(s^{0}\) and \(s^{\,n+1}\) in \eqref{plt:eq:om-E-pde} gives
\(e_0'=0\) and \(e_{n+1}'=n\,e_n-e_n'\); Step~2 gives \(e_n(0)=0\) for all
\(n\).  Hence \(e_0\) is a constant with \(e_0(0)=0\), so \(e_0=0\); and
inductively, if \(e_n=0\) then \(e_{n+1}'=0\) and \(e_{n+1}(0)=0\), so
\(e_{n+1}=0\).  Thus \(E=0\), which is \eqref{plt:eq:om-generating}.

\emph{Uniqueness.}  Let \(\mathcal S(Q)=-\log(1-su+sQ)\).  For
\(Q\in\K[u][[s]]\) the argument lies in \(1+s\K[u][[s]]\), so
\(\mathcal S(Q)\in s\K[u][[s]]\) and in particular \(\mathcal S(Q)(0,u)=0\).
If \(Q_1-Q_2\in s^{N}\K[u][[s]]\) then, \(1-su+sQ_2\) being a unit,
\[
 \mathcal S(Q_1)-\mathcal S(Q_2)
 =-\log\left(1+\frac{s(Q_1-Q_2)}{1-su+sQ_2}\right)\in s^{\,N+1}\K[u][[s]].
\]
So \(\mathcal S\) is a strict contraction for the \(s\)-adic filtration on
\(\K[u][[s]]\), which is separated and complete; it therefore has exactly one
fixed point, and that fixed point automatically satisfies \(Q(0,u)=0\).
\end{proof}

% ed.: this corollary carried its own argument inside the statement, so it read
% as an assertion with no proof.  The claim and the argument are separated.
\begin{corollary}[Reversion-free proof of \cref{plt:thm:om-expansion}]
\label{plt:cor:om-reversion-free}
The generating equation \eqref{plt:eq:om-generating} determines the expansion
of \(\WrightOmega\) on its own:
\[
  \rev{(X+L)}=X+\mathcal P(t,L),
  \qquad
  \mathcal P(t,L)=\sum_{n\ge0}\LamP{n}(L)t^{n},
\]
so \cref{plt:thm:om-expansion} holds without appeal to Lagrange--B\"urmann and
without \cref{plt:lem:om-repeated-derivative}.
\end{corollary}

\begin{proof}
Substitute \(s\mapsto t\) and \(u\mapsto L\) in
\eqref{plt:eq:om-generating}, and add \(u=L\) to both sides of
\(\mathcal P=-u+Q\).  The result is exactly the assertion that
\(\mathcal P(t,L)\) satisfies the implicit equation
\eqref{plt:eq:om-implicit}.  By the uniqueness half of
\cref{plt:prop:om-implicit}, the solution of that equation in
\(\Gnear\) is unique, so \(\rev{(X+L)}=X+\mathcal P(t,L)\).

The two ingredients used here are the existence and uniqueness of the fixed
point of the generating equation and the uniqueness half of
\cref{plt:prop:om-implicit}; neither depends on
\cref{plt:thm:lag-burmann} or on the repeated-derivative computation, which is
what makes this route independent.
\end{proof}

\begin{remark}[Which route to use]
\label{plt:rmk:om-routes}
Three independent characterisations of the same family are now available: the
operator formula \eqref{plt:eq:lambert-operator}, the recurrence
\eqref{plt:eq:lambert-int-rec}, and the generating equation
\eqref{plt:eq:om-generating}.  The recurrence is the cheapest generator
(one exact polynomial integration and one subtraction per index); the
operator formula is the right tool for proving identities, since it makes the
falling factorial visible; the generating equation is the right tool for
uniqueness statements and for transporting the family to a new inversion
problem, as in \cref{plt:thm:om-alpha-family}.  Agreement of two of them on a
computed table is a genuine check only if the two were implemented
independently; \eqref{plt:eq:om-residual} of \cref{plt:thm:om-residual} is a
stronger test, because it exercises the composition code as well.
\end{remark}

\section{The Laplace transform and the exponential moments}

The remaining result is of a different type: it is an assertion about
convergent integrals of the real polynomial functions \(\LamP{n}\), not a
formal identity, and its hypotheses have real content.

\begin{theorem}[Laplace transform of \(\LamP{n}\)]
\label{plt:thm:om-laplace}
Let \(n\ge1\) and let \(s\in\ComplexNumbers\) with \(\Re s>0\).  Then the
integral below converges absolutely and
\begin{equation}\label{plt:eq:om-laplace}
 \int_{0}^{\infty}\EulerE^{-su}\,\LamP{n}(u)\,\Differential u
 =\frac{1}{s^{\,n+1}}\prod_{j=1}^{n-1}(j-s)
 =\frac{\RisingFactorial{1-s}{\,n-1}}{s^{\,n+1}},
\end{equation}
the product being empty (and equal to \(1\)) when \(n=1\).
\end{theorem}

\begin{proof}
Absolute convergence is clear: \(\LamP{n}\) is a polynomial and
\(\AbsoluteValue{\EulerE^{-su}}=\EulerE^{-u\Re s}\) with \(\Re s>0\), so the
integrand is \(\BigOAt{\EulerE^{-u\Re s/2}}{u\to+\infty}\).  Write
\(I_n(s)=\int_0^\infty\EulerE^{-su}\LamP{n}(u)\,\Differential u\).

\emph{The integration by parts.}  For \(0<b<\infty\),
\[
 \int_{0}^{b}\EulerE^{-su}{\LamP{n}}'(u)\,\Differential u
 =\Bigl[\EulerE^{-su}\LamP{n}(u)\Bigr]_{0}^{b}
  +s\int_{0}^{b}\EulerE^{-su}\LamP{n}(u)\,\Differential u .
\]
The boundary term at \(u=0\) vanishes because \(\LamP{n}(0)=0\)
(\cref{plt:thm:lambert-recurrence}); the boundary term at \(u=b\) tends to
\(0\) as \(b\to\infty\) because \(\Re s>0\) and \(\LamP{n}\) is a
polynomial.  Both hypotheses are used, and neither can be dropped: without
\(\Re s>0\) the integral does not converge, and without
\(\LamP{n}(0)=0\) an extra term \(-\LamP{n}(0)\) would survive.  Letting
\(b\to\infty\),
\begin{equation}\label{plt:eq:om-laplace-ibp}
 \int_{0}^{\infty}\EulerE^{-su}{\LamP{n}}'(u)\,\Differential u=s\,I_n(s).
\end{equation}

\emph{The recursion.}  Apply \eqref{plt:eq:om-laplace-ibp} to both sides of
\eqref{plt:eq:lambert-diff-rec}, which is an identity of polynomials and may
therefore be integrated against \(\EulerE^{-su}\):
\[
 s\,I_{n+1}(s)=n\,I_n(s)-s\,I_n(s),
 \qquad\text{i.e.}\qquad
 I_{n+1}(s)=\frac{n-s}{s}\,I_n(s)\quad(n\ge1).
\]
The base case is \(I_1(s)=\int_0^\infty\EulerE^{-su}u\,\Differential u=s^{-2}\).
Iterating,
\(I_n(s)=s^{-2}\prod_{j=1}^{n-1}\frac{j-s}{s}=s^{-(n+1)}\prod_{j=1}^{n-1}(j-s)\).
Finally \(\prod_{j=1}^{n-1}(j-s)=(1-s)(2-s)\cdots(n-1-s)
=\RisingFactorial{1-s}{\,n-1}\).
\end{proof}

\begin{corollary}[Exponential moment cancellations]
\label{plt:cor:om-moment-cancellation}
Let \(n\ge1\).  For every integer \(r\) with \(1\le r\le n-1\),
\begin{equation}\label{plt:eq:om-moment}
 \int_{0}^{\infty}\EulerE^{-ru}\,\LamP{n}(u)\,\Differential u=0 .
\end{equation}
The range is empty for \(n=1\), so the first genuine cancellation is
\(\int_0^\infty\EulerE^{-u}\LamP{2}(u)\,\Differential u=0\).
\end{corollary}

\begin{proof}
Such an \(r\) is a positive real number, so \(\Re r>0\) and
\cref{plt:thm:om-laplace} applies with \(s=r\).  The factor \(j=r\) occurs in
\(\prod_{j=1}^{n-1}(j-s)\) precisely because \(1\le r\le n-1\), and it
vanishes at \(s=r\).
\end{proof}

\begin{remark}[The transform is the falling factorial again]
\label{plt:rmk:om-laplace-stirling}
\Cref{plt:thm:om-laplace} and \cref{plt:thm:lambert-stirling} are equivalent,
and the equivalence explains the cancellations of
\cref{plt:cor:om-moment-cancellation}.  Indeed
\(\FallingFactorial{s}{n}=s\prod_{j=1}^{n-1}(s-j)
=(-1)^{n-1}s\prod_{j=1}^{n-1}(j-s)\), so by
\eqref{plt:eq:om-stirling-def}
\[
 \prod_{j=1}^{n-1}(j-s)
 =(-1)^{n-1}\sum_{k\ge1}\stirlingsigned{n}{k}s^{\,k-1}
 =\sum_{k\ge1}(-1)^{k-1}\stirlingone{n}{k}s^{\,k-1},
\]
whence
\(I_n(s)=\sum_{k=1}^{n}(-1)^{k-1}\stirlingone{n}{k}s^{-(n-k+2)}\).
Comparing with the termwise transform
\(I_n(s)=\sum_{m=1}^{n}\lambda_{n,m}\,m!\,s^{-(m+1)}\) and matching powers
via \(m=n-k+1\) returns exactly \eqref{plt:eq:lambert-lambda}.  Thus
\eqref{plt:eq:om-moment} says that the falling factorial
\(\FallingFactorial{s}{n}\) vanishes at \(s=1,\ldots,n-1\) --- the
cancellations are the roots of the very operator that defines \(\LamP{n}\) in
\eqref{plt:eq:lambert-operator}.
\end{remark}

\begin{warningbox}
The letter \(s\) carries two unrelated meanings in this section.  In
\cref{plt:lem:om-pde} and \cref{plt:thm:om-generating-equation} it is a
formal indeterminate and no analytic content is asserted.  In
\cref{plt:thm:om-laplace} it is a complex number constrained by
\(\Re s>0\), and the identity is an equality of convergent integrals with
analytic functions of \(s\) on the half-plane.  A ``Laplace transform of the
generating series'' mixing the two would be a different object and is not
claimed anywhere here.
\end{warningbox}

% ---- fragment 11-lambertw ----
\chapter{Lambert \texorpdfstring{\(\LambertW\)}{W}: branches and expansions}
\label{plt:sec:lambert-branches}

Every result of Parts~I--IV was formal: an identity between elements of
\(\PLpow\), \(\PLlau\), \(\PLrat\) or \(\PLit\), or between elements of the
near-identity group \(\Gnear\).  Lambert's \(\LambertW\) is where that formal
theory meets a genuine multivalued analytic function, and it is therefore the
place where the boundary between an algebraic identity and an analytic
theorem has to be drawn with care.  The six source reports draw it in six
different places, and none of them draws it entirely correctly; the present
section states exactly what is proved, proves it, and marks what is not.

The organising principle is a single substitution.  The equation
\(w\EulerE^{w}=z\) is not itself a polynomial--logarithmic reversion problem:
its natural variable \(z\) is exponentially larger than its solution.  After
one logarithm it becomes the Wright omega problem \(w+\log w=\xi\), which is
exactly the reversion of \(X+L\) in \(\Gnear\), already solved in
\cref{plt:sec:wright-omega}.  All of \(\LambertW\) --- every branch, the
lower real branch near zero, and the generalized Lambert families of
\cref{plt:sec:generalized-lambert} --- is that one reversion evaluated at
different arguments.  What changes from branch to branch is the argument, not
the reversion.

\section{The defining equation and the real branches}
\label{plt:sec:lw-defining}

\begin{definition}[Branches of \(\LambertW\)]
\label{plt:def:lw-branches}
For \(k\in\IntegerNumbers\), a \emph{branch} \(\LambertW_k\) of the Lambert
function is a determination of the inverse of
\[
 g\colon \ComplexNumbers\longrightarrow\ComplexNumbers,
 \qquad
 g(w)\DefinitionEquals w\EulerE^{w},
\]
so that
\begin{equation}\label{plt:eq:lw-defining}
 \LambertW_k(z)\EulerE^{\LambertW_k(z)}=z .
\end{equation}
The map \(g\) is entire and infinite-to-one, so \cref{plt:eq:lw-defining}
alone determines nothing: a branch is the pair consisting of a solution and a
rule for selecting it.  Following the notation catalogue, \(\LambertW_0\) is
the \emph{principal} branch, real and increasing on
\([-\EulerE^{-1},\infty)\) with \(\LambertW_0(0)=0\), analytic at \(0\), and
cut along \((-\infty,-\EulerE^{-1}]\); \(\LambertW_{-1}\) is the
\emph{lower real} branch, real on \([-\EulerE^{-1},0)\) with values in
\((-\infty,-1]\).  A branchless symbol \(\LambertW(z)\) is admissible only
where the domain forces the real principal branch.
\end{definition}

\begin{proposition}[The two real branches]
\label{plt:prop:lw-real-branches}
Let \(g(w)=w\EulerE^{w}\) on \(\RealNumbers\).  Then:
\begin{enumerate}
\item \(g\) is strictly decreasing on \((-\infty,-1]\), with image
      \([-\EulerE^{-1},0)\), and strictly increasing on \([-1,\infty)\),
      with image \([-\EulerE^{-1},\infty)\);
\item consequently \(g\) restricted to \([-1,\infty)\) has a continuous
      strictly increasing inverse \(\LambertW_0\colon
      [-\EulerE^{-1},\infty)\to[-1,\infty)\), and \(g\) restricted to
      \((-\infty,-1]\) has a continuous strictly decreasing inverse
      \(\LambertW_{-1}\colon[-\EulerE^{-1},0)\to(-\infty,-1]\);
\item for \(x\in(-\EulerE^{-1},0)\) the equation \(g(w)=x\) has exactly the
      two real solutions \(\LambertW_0(x)\in(-1,0)\) and
      \(\LambertW_{-1}(x)\in(-\infty,-1)\); for \(x=-\EulerE^{-1}\) it has
      the single solution \(-1\); for \(x\ge0\) it has exactly one; and for
      \(x<-\EulerE^{-1}\) none;
\item \(\LambertW_{-1}(x)\to-\infty\) as \(x\uparrow0\).
\end{enumerate}
\end{proposition}

\begin{proof}
We have \(g'(w)=(1+w)\EulerE^{w}\), which is negative for \(w<-1\), zero at
\(w=-1\) and positive for \(w>-1\); thus \(g\) is strictly decreasing on
\((-\infty,-1]\) and strictly increasing on \([-1,\infty)\), with the single
minimum \(g(-1)=-\EulerE^{-1}\).  On \((-\infty,-1]\) we have
\(g(w)<0\) for \(w<0\) and \(g(w)\to0^{-}\) as \(w\to-\infty\), because
\(\AbsoluteValue{w}\EulerE^{w}\to0\); hence \(g\) maps \((-\infty,-1]\)
bijectively onto \([-\EulerE^{-1},0)\), decreasingly.  On \([-1,\infty)\),
\(g(-1)=-\EulerE^{-1}\) and \(g(w)\to+\infty\), so \(g\) maps
\([-1,\infty)\) bijectively onto \([-\EulerE^{-1},\infty)\), increasingly.
Items (i)--(iii) follow, and (iv) is the statement that the decreasing
inverse of \(g|_{(-\infty,-1]}\) tends to \(-\infty\) as its argument tends
to the endpoint \(0^{-}\) of the image.
\end{proof}

\begin{lemma}[Derivative of a branch]
\label{plt:lem:lw-derivative}
Let \(w=\LambertW_k(z)\) be analytic near \(z_0\) with \(z_0\ne0\) and
\(w\ne-1\).  Then
\begin{equation}\label{plt:eq:lw-derivative}
 \frac{\Differential\LambertW_k}{\Differential z}(z)
 =\frac{\LambertW_k(z)}{z\bigl(1+\LambertW_k(z)\bigr)} .
\end{equation}
\end{lemma}

\begin{proof}
Differentiating \(w\EulerE^{w}=z\) gives
\((1+w)\EulerE^{w}w'=1\).  Since \(\EulerE^{w}=z/w\) whenever \(w\ne0\),
this is \((1+w)(z/w)w'=1\), which is \cref{plt:eq:lw-derivative}.  The
excluded cases \(z=0\) and \(w=-1\) are exactly the points where the
substitution \(\EulerE^{w}=z/w\) or the division by \(1+w\) is illegitimate;
\(w=-1\) is the branch point studied in
\cref{plt:sec:lw-branch-point}.
\end{proof}

\section{Branch coordinates and the substitution to the omega problem}
\label{plt:sec:lw-coordinates}

\begin{definition}[Branch coordinates]
\label{plt:def:lw-coordinates}
Let \(z\in\ComplexNumbers\setminus\SetOf{0}\) and \(k\in\IntegerNumbers\).
The \emph{outer branch coordinate} is
\begin{equation}\label{plt:eq:lw-xi}
 \xi_k\DefinitionEquals\PrincipalLogarithm z+2\pi\ImaginaryUnit k,
\end{equation}
with \(\PrincipalLogarithm\) the principal logarithm, so that
\(\EulerE^{\xi_k}=z\).  Given in addition \(\xi_k\ne0\), an \emph{inner
branch coordinate} is any \(\ell_k\in\ComplexNumbers\) with
\begin{equation}\label{plt:eq:lw-ell}
 \EulerE^{\ell_k}=\xi_k ,
\end{equation}
written \(\ell_k=\log_{\mathcal B_k}\xi_k\) when a particular determination
\(\mathcal B_k\) has been fixed.  The pair \((\xi_k,\ell_k)\) is the
\emph{branch datum}.  A \emph{coordinate sector} is a connected open
\(S_k\subseteq\ComplexNumbers\setminus\SetOf{0}\) on which
\(z\mapsto\xi_k\) and \(\xi_k\mapsto\ell_k\) admit single-valued analytic
determinations and on which \(\AbsoluteValue{\xi_k}\to\infty\) along the
limits considered.
\end{definition}

Three data must be recorded and are routinely conflated: the choice of
\(\PrincipalLogarithm z\) (equivalently, of the strip in which
\(\operatorname{Arg}z\) is taken), the integer \(k\), and the determination
\(\mathcal B_k\) of the second logarithm.  Reports~S1, S2 and S4 all write
\(\ell_k=\PrincipalLogarithm\xi_k\), i.e.\ they silently force
\(\mathcal B_k\) to be principal.  That is a legitimate normalisation and we
use it in all numerical statements, but it is a hypothesis, not a
convention-free fact: \cref{plt:prop:lw-branch-label} shows exactly where it
is used.

\begin{proposition}[Reduction to the omega problem]
\label{plt:prop:lw-log-reduction}
Let \(z\ne0\), \(k\in\IntegerNumbers\), and let \(w\in\ComplexNumbers
\setminus\SetOf{0}\).  The following are equivalent.
\begin{enumerate}
\item \(w\EulerE^{w}=z\) and there is a logarithm \(\lambda\) of \(w\)
      (that is, \(\EulerE^{\lambda}=w\)) with \(w+\lambda=\xi_k\).
\item \(w+\lambda=\xi_k\) for some \(\lambda\) with \(\EulerE^{\lambda}=w\).
\end{enumerate}
Moreover, if \(w\EulerE^{w}=z\) and \(\lambda\) is any logarithm of \(w\),
then \(w+\lambda=\xi_j\) for exactly one \(j\in\IntegerNumbers\), and every
\(j\) arises from a suitable choice of \(\lambda\).
\end{proposition}

\begin{proof}
Assume (ii).  Then
\(w\EulerE^{w}=\EulerE^{\lambda}\EulerE^{w}=\EulerE^{w+\lambda}
=\EulerE^{\xi_k}=z\), which is (i).  Conversely (i) contains (ii).

For the last statement, let \(w\EulerE^{w}=z\) and \(\EulerE^{\lambda}=w\).
Then \(\EulerE^{w+\lambda}=w\EulerE^{w}=z=\EulerE^{\PrincipalLogarithm z}\),
so \(w+\lambda-\PrincipalLogarithm z\in2\pi\ImaginaryUnit\IntegerNumbers\),
that is \(w+\lambda=\xi_j\) for a unique \(j\).  Replacing \(\lambda\) by
\(\lambda+2\pi\ImaginaryUnit m\) (still a logarithm of \(w\)) replaces
\(j\) by \(j+m\), so every integer occurs.
\end{proof}

\begin{remark}[What the substitution does and does not say]
\label{plt:rmk:lw-unwinding}
\Cref{plt:prop:lw-log-reduction} is the exact form of the informal step
``take a logarithm of \(w\EulerE^{w}=z\)''.  Two things must be noticed.
First, the branch index \(k\) is \emph{defined} by the relation
\(w+\lambda=\xi_k\); it is not an independent piece of data that can be
chosen after \(w\) and \(\lambda\).  Second, the equivalence is with the
\emph{additive} equation \(w+\lambda=\xi_k\), in which \(\lambda\) is a
logarithm of the unknown \(w\).  Whether \(\lambda\) is the principal
logarithm of \(w\) is a separate question, settled for large
\(\AbsoluteValue{\xi_k}\) by \cref{plt:prop:lw-branch-label}.  The
discrepancy \(\lambda-\PrincipalLogarithm w\), divided by
\(2\pi\ImaginaryUnit\), is the \emph{unwinding number}; the identity
\(\WrightOmega(s)=\LambertW(\EulerE^{s})\) asserted without qualification in
S1 and S3 holds only where that unwinding number vanishes.
\end{remark}

\section{The formal expansion and branch independence}
\label{plt:sec:lw-formal}

We first record the formal statement, which is an identity in the
near-identity group and involves no analysis whatsoever.  Throughout, \(\K\)
is a field of characteristic zero, \(X\) is the large variable,
\(t=X^{-1}\), \(L=\log X\), and \(\Gnear=X+\K[L][[t]]\).

\begin{theorem}[Formal affine-logarithmic reversion]
\label{plt:thm:lw-affine-log}
Let \(a\in\K\) and let \(F=X+aL\in\Gnear\).  Then the compositional inverse
of \(F\) in \(\Gnear\) is
\begin{equation}\label{plt:eq:lw-affine-formal}
 \rev{F}
 =X+a\sum_{n\ge0}a^{n}\LamP{n}(L)\,t^{n}
 =X-aL+\sum_{n\ge1}a^{n+1}\LamP{n}(L)\,t^{n}.
\end{equation}
In particular the coefficient of \(t^{n}\) in \(\rev{F}\) is the fixed
polynomial \(\LamP{n}\) multiplied by the scalar \(a^{n+1}\); no new
polynomial family occurs for any value of \(a\), and \(a=1\) recovers the
Wright omega expansion of \cref{plt:thm:om-expansion}.
\end{theorem}

\begin{proof}
Let \(\mathcal Q(s,u)=\sum_{n\ge1}\LamP{n}(u)s^{n}\in\K[u][[s]]\) be the
generating series of \cref{plt:thm:om-generating-equation}, characterised
there as the unique element of \(s\K[u][[s]]\) with
\begin{equation}\label{plt:eq:lw-generating}
 \mathcal Q=-\log\bigl(1-su+s\mathcal Q\bigr).
\end{equation}
Because \(at\) has positive \(t\)-order, the assignment
\(s\mapsto at\), \(u\mapsto L\) extends to a continuous \(\K\)-algebra
homomorphism \(\K[u][[s]]\to\K[L][[t]]\) for the respective adic
topologies.  Write
\[
 \mathcal Q_a\DefinitionEquals\mathcal Q(at,L)
 =\sum_{n\ge1}a^{n}\LamP{n}(L)\,t^{n}\in t\K[L][[t]] .
\]
Applying the homomorphism to \cref{plt:eq:lw-generating} gives, in
\(\K[L][[t]]\),
\begin{equation}\label{plt:eq:lw-generating-transported}
 \mathcal Q_a=-\log\bigl(1-atL+at\,\mathcal Q_a\bigr),
\end{equation}
the logarithm being the formal logarithm of
\cref{plt:lem:cmp-exp-log}, legitimate because the argument lies in
\(1+t\K[L][[t]]\).

Put \(G\DefinitionEquals X-aL+a\mathcal Q_a\).  Since
\(-aL+a\mathcal Q_a\in\K[L][[t]]\), we have \(G\in\Gnear\), and
\[
 G=X\Bigl(1-at\,(L-\mathcal Q_a)\Bigr),
\]
whose second factor lies in \(1+t\K[L][[t]]\).  By the exact generator
substitution rule of \cref{plt:prop:cmp-generator-substitution},
\[
 L\compose G=\log G=L+\log\bigl(1-atL+at\,\mathcal Q_a\bigr)
 =L-\mathcal Q_a,
\]
the last equality being \cref{plt:eq:lw-generating-transported}.  Hence
\[
 F\compose G=G+a\,(L\compose G)
 =X-aL+a\mathcal Q_a+aL-a\mathcal Q_a=X .
\]
So \(G\) is a right inverse of \(F\) inside \(\Gnear\).  By
\cref{plt:thm:rev-existence} the inverse in \(\Gnear\) exists and is unique
and two-sided, so \(G=\rev{F}\).  Finally
\(a\LamP{0}(L)=-aL\), which converts the second display of
\cref{plt:eq:lw-affine-formal} into the first.
\end{proof}

\begin{remark}[Two proofs in the sources, and why this one is preferred]
Reports~S2 and S4 obtain \cref{plt:eq:lw-affine-formal} by inspecting the
Lagrange--B\"urmann series and observing that its \(m\)-th term carries a
factor \(a^{m}\).  That argument is correct but needs the summability of the
Lagrange family (\cref{plt:prop:lag-summability-sharp}) and a separate
identification of the \(m\)-th term with \(\LamP{m}\).  The proof above
verifies a closed-form candidate against the defining equation and appeals
only to uniqueness in \(\Gnear\); it is shorter, it is the same proof for
every \(a\) including \(a=-1\), and it is the one that transports without
change to the analytic statement of \cref{plt:thm:lw-basepoint}.
\end{remark}

\begin{theorem}[Formal branch expansion]
\label{plt:thm:lw-formal-expansion}
Let \(k\in\IntegerNumbers\) and let \((\xi_k,\ell_k)\) be a branch datum in
the sense of \cref{plt:def:lw-coordinates}.  Substituting \(X\mapsto\xi_k\),
\(L\mapsto\ell_k\), \(t\mapsto\xi_k^{-1}\) in
\cref{plt:eq:lw-affine-formal} with \(a=1\) produces the formal series
\begin{equation}\label{plt:eq:lw-branch-expansion}
 \xi_k+\sum_{n\ge0}\frac{\LamP{n}(\ell_k)}{\xi_k^{\,n}}
 =\xi_k-\ell_k+\sum_{n\ge1}\frac{\LamP{n}(\ell_k)}{\xi_k^{\,n}} ,
\end{equation}
equivalently, by \cref{plt:thm:lambert-stirling},
\begin{equation}\label{plt:eq:lw-branch-stirling}
 \xi_k-\ell_k+\sum_{n\ge1}\frac{(-1)^{n}}{\xi_k^{\,n}}
 \sum_{m=1}^{n}\stirlingone{n}{n-m+1}\frac{(-\ell_k)^{m}}{m!} .
\end{equation}
The coefficient polynomials \(\LamP{n}\) do not depend on \(k\), on the
determination \(\mathcal B_k\), or on the sector.
\end{theorem}

\begin{proof}
The first display is \cref{plt:eq:lw-affine-formal} with \(a=1\) after the
stated substitution, which is a substitution of symbols and requires no
justification beyond \(\xi_k\ne0\).  For the second, insert the closed
coefficient formula
\(\LamP{n}(u)=\sum_{m=1}^{n}(-1)^{n-m}\stirlingone{n}{n-m+1}u^{m}/m!\)
of \cref{plt:thm:lambert-stirling} and use
\((-1)^{n-m}=(-1)^{n}(-1)^{m}\), which converts \(u^{m}\) into
\((-u)^{m}\) with an overall factor \((-1)^{n}\).  Branch independence is the
observation that \cref{plt:thm:lw-affine-log} was proved once, in
\(\K[L][[t]]\), with no branch data present.
\end{proof}

The next theorem is the precise form of the catalogue's assertion that
``the polynomial family is branch-independent; the coordinates
\(\xi_k,\ell_k\), their branches, and the analytic domain are not''.  Its
first half is the uniqueness of the coefficients, and its second half
exhibits the genuine \(k\)-dependence.

\begin{theorem}[Branch independence of the coefficients, dependence of the
coordinates]
\label{plt:thm:lw-branch-independence}
\hfill
\begin{enumerate}
\item \textup{(Universality.)}  There is exactly one sequence
      \((p_n)_{n\ge0}\) in \(\K[u]\) such that for every
      \(a\in\K\) the map \(X+aL\in\Gnear\) has inverse
      \(X+a\sum_{n\ge0}a^{n}p_n(L)t^{n}\); namely \(p_n=\LamP{n}\).
      A fortiori, for each fixed \(k\), the expansion
      \cref{plt:eq:lw-branch-expansion} determines the polynomials
      \(\LamP{n}\) uniquely, and the answer is the same for all \(k\).
\item \textup{(Uniqueness against the coordinates.)}  Fix \(k\) and a
      coordinate sector \(S_k\) on which
      \(\AbsoluteValue{\xi_k}\to\infty\) and
      \(\ell_k=\PrincipalLogarithm\xi_k\).  If \((q_n)_{n\ge0}\) in
      \(\ComplexNumbers[u]\) satisfies
      \(\sum_{n=0}^{N}q_n(\ell_k)\xi_k^{-n}
      =\BigOAt{\ell_k^{\,N+1}/\xi_k^{\,N+1}}
      {\AbsoluteValue{\xi_k}\to\infty,\ \xi_k\in S_k}\)
      for every \(N\), then \(q_n=0\) for all \(n\).
\item \textup{(The coordinates really differ.)}  For \(j\ne k\) and \(z\) in
      a common domain, \(\xi_j-\xi_k=2\pi\ImaginaryUnit(j-k)\ne0\), and the
      corresponding solutions \(w_j,w_k\) of \cref{plt:eq:lw-defining}
      supplied by \cref{plt:thm:lw-basepoint} satisfy
      \(w_j-w_k\to2\pi\ImaginaryUnit(j-k)\) as
      \(\AbsoluteValue{z}\to\infty\) with \(\operatorname{Arg}z\) bounded.
      In particular they are distinct functions.
\end{enumerate}
\end{theorem}

\begin{proof}
(i)  Existence is \cref{plt:thm:lw-affine-log}.  For uniqueness take
\(a=1\): if \(X+\sum_{n\ge0}p_n(L)t^{n}\) and
\(X+\sum_{n\ge0}\LamP{n}(L)t^{n}\) are both inverses of \(X+L\) in
\(\Gnear\), they are equal by the uniqueness clause of
\cref{plt:thm:rev-existence}, and equality in \(\K[L][[t]]\) is coefficientwise
by \cref{plt:thm:lead-uniqueness}.  (Uniqueness already follows from the
single value \(a=1\); no compatibility across different \(a\) is needed.)

(ii)  Induct on \(n\).  Suppose \(q_0=\dots=q_{n-1}=0\) (vacuous for
\(n=0\)).  Taking \(N=n\) in the hypothesis and multiplying by
\(\xi_k^{\,n}\) gives
\(q_n(\ell_k)=\BigOAt{\ell_k^{\,n+1}/\xi_k}
{\AbsoluteValue{\xi_k}\to\infty,\ \xi_k\in S_k}\).
Suppose \(q_n\ne0\), say of degree \(d\ge0\) with leading coefficient
\(c\ne0\).  Since \(\ell_k=\PrincipalLogarithm\xi_k\) we have
\(\AbsoluteValue{\ell_k}\le\log\AbsoluteValue{\xi_k}+\pi\) and
\(\AbsoluteValue{\ell_k}\ge\log\AbsoluteValue{\xi_k}\), so
\(\AbsoluteValue{\ell_k}\to\infty\) and
\(\AbsoluteValue{q_n(\ell_k)}\ge
\tfrac12\AbsoluteValue{c}\AbsoluteValue{\ell_k}^{d}\) for
\(\AbsoluteValue{\xi_k}\) large, while the right-hand side is at most
\(C(\log\AbsoluteValue{\xi_k}+\pi)^{n+1}/\AbsoluteValue{\xi_k}\to0\).  For
\(d\ge1\) the left side tends to \(\infty\) and for \(d=0\) it is the
nonzero constant \(\AbsoluteValue{c}\); either way we contradict the bound.
Hence \(q_n=0\).

(iii)  The identity \(\xi_j-\xi_k=2\pi\ImaginaryUnit(j-k)\) is immediate
from \cref{plt:eq:lw-xi}.
% ed.: the hypotheses of the two theorems invoked here were not printed;
% both hold once |z| is large, which is the regime of the claim.
Since \(\operatorname{Arg}z\) stays bounded, both
\(\AbsoluteValue{\xi_j}\) and \(\AbsoluteValue{\xi_k}\) tend to infinity
while \(\AbsoluteValue{\ell_j},\AbsoluteValue{\ell_k}\) grow only
logarithmically, so for \(\AbsoluteValue{z}\) large both branch data satisfy
\(\AbsoluteValue{\xi}\ge10(1+\AbsoluteValue{\ell})\); we restrict to that
range.  By \cref{plt:thm:lw-basepoint} applied with base
point \(\xi_j\), respectively \(\xi_k\), we have
\(w_j=\xi_j-\ell_j+\mathcal Y_j\) and \(w_k=\xi_k-\ell_k+\mathcal Y_k\),
where \cref{plt:thm:lw-uniform-remainder} with \(N=0\) bounds
\(\AbsoluteValue{\mathcal Y_j},\AbsoluteValue{\mathcal Y_k}\) by
\(10(1+\AbsoluteValue{\ell})/\AbsoluteValue{\xi}\), which tends to \(0\).
Hence
\[
 w_j-w_k-2\pi\ImaginaryUnit(j-k)
 =-(\ell_j-\ell_k)+\mathcal Y_j-\mathcal Y_k .
\]
Now \(\ell_j-\ell_k
=\PrincipalLogarithm(\xi_j)-\PrincipalLogarithm(\xi_k)\to0\), because
\(\xi_j/\xi_k=1+2\pi\ImaginaryUnit(j-k)/\xi_k\to1\) and both
\(\operatorname{Arg}\xi_j\) and \(\operatorname{Arg}\xi_k\) tend to \(0\)
when \(\operatorname{Arg}z\) stays bounded.  All three terms on the right
tend to \(0\), so \(w_j-w_k\to2\pi\ImaginaryUnit(j-k)\ne0\).
\end{proof}

\begin{warningbox}[title={Universality is a statement about one problem, not
about all of them}]
\Cref{plt:thm:lw-branch-independence}\,(i) says that the \emph{same}
polynomials serve every \(a\) and every branch.  It does not say that
\(\LamP{n}\) is the coefficient family of every logarithmic reversion.  Even
inside the affine-logarithmic family the scalar \(a^{n+1}\) is present, and
% ed.: the degree law and leading coefficient quoted here for X + aL^2 were
% printed for all n >= 0, but they fail at n = 0: the t^0 coefficient is
% -aL^2, of degree 2, and the formula 2a^{n+1}/n is undefined at n = 0.
% Restricted to n >= 1 and the n = 0 value recorded; verified by explicit
% reversion through n = 4 (c_1 = 2a^2L^3, c_2 = a^3L^5 - 5a^3L^4,
% c_3 = (2/3)a^4L^7 - 7a^4L^6 + 14a^4L^5).
% ed.: independently re-verified through n = 9 in exact rational arithmetic,
% with the residual of the reversion confirmed to vanish identically through
% t^9; that computation also produced the Catalan observation recorded below.
as soon as the perturbation ceases to be \(aL\) --- for instance for
\(X+aL^{2}\) with \(a\ne0\), whose \(t^{0}\)-coefficient is \(-aL^{2}\) and
whose \(t^{n}\)-coefficient for \(n\ge1\) has degree
\(2n+1\) in \(L\) and leading coefficient \(2a^{n+1}/n\) by
\cref{plt:thm:lag-degree-sharp} --- the coefficients are new polynomials.
The universality of \(\LamP{n}\) is exactly the universality of the single
map \(X+L\).

\begin{remark}[An unexplained Catalan pattern in the \(X+aL^{2}\) family]
\label{plt:rmk:lw-catalan}
Reversing \(X+aL^{2}\) in exact rational arithmetic through \(n=9\), the
\emph{lowest}-degree term of the \(t^{n}\)-coefficient is
\[
 (-1)^{n}\,\CatalanNumber{n+1}\;a^{n+1}L^{n+2},
 \qquad 1\le n\le9,
\]
where \(\CatalanNumber{m}=\binom{2m}{m}/(m+1)\): the coefficients are
\(2,-5,14,-42,132,-429,1430,-4862,16796\).  The top-degree end of the same
coefficient is governed by \cref{plt:thm:lag-degree-sharp}, which gives degree
\(2n+1\) and leading coefficient \(2a^{n+1}/n\); the bottom-degree end has no
such explanation here.

This is an observation, not a theorem.  Nine confirmations are evidence and
nothing more, and no proof is offered: the Lagrange--B\"urmann coefficient
formula of \cref{plt:thm:lag-coefficients} controls the top of each block,
because the highest power of \(L\) comes from the term of largest \(r\), while
the bottom of the block collects contributions from every \(r\) at once and is
not visibly a single Catalan-counted sum.  A proof would presumably identify
the \(L^{n+2}\) coefficient with a lattice-path or tree count already attached
to the quadratic-logarithm reversion; we have not found that identification.
It is recorded because a pattern this clean is either a theorem with a short
proof or a coincidence that ends at \(n=10\), and both outcomes are worth
knowing.
\end{remark}
\end{warningbox}

\section{The analytic status of the expansion}
\label{plt:sec:lw-analytic}

\Cref{plt:thm:lw-formal-expansion} is a theorem about symbols.  It asserts
nothing about the function \(\LambertW_k\).  The sources at this point
diverge sharply.  S1 and S3 assert Poincar\'e validity ``in an admissible
sector'' without defining the sector; S6 asserts convergence by an
implicit-function argument that quantifies incorrectly (see
\cref{plt:rmk:lw-s6-gap}); S1, S3 and S5 all repeat from the literature the
much stronger claim that for real \(z\) the series converges for
\(z\ge\EulerE\).  We prove here a self-contained analytic theorem which is
strong enough for every use made of the expansion in this volume --- it
gives absolute convergence, an explicit region, and an explicit uniform
remainder --- and we mark the literature claim as unverified.

The whole analytic content sits in one two-variable lemma.  Its point is
that the two small quantities in the problem are \(1/\xi_k\) and
\(\ell_k/\xi_k\), and that the solution is jointly analytic in them near the
origin.  Separating them is what makes the estimate uniform in \(\ell_k\),
which is precisely the step that no source performs.

\begin{lemma}[The two-variable solution]
\label{plt:lem:lw-two-variable}
Let
\[
 D\DefinitionEquals\SetOf{(\sigma,\tau)\in\ComplexNumbers^{2}:
 \AbsoluteValue{\sigma}+\AbsoluteValue{\tau}<\tfrac12}.
\]
For \((\sigma,\tau)\in D\) the equation
\begin{equation}\label{plt:eq:lw-two-var-fixed}
 y=-\PrincipalLogarithm\bigl(1-\tau+\sigma y\bigr)
\end{equation}
has a unique solution \(y=\mathcal Y(\sigma,\tau)\) in the closed unit disc
\(\AbsoluteValue{y}\le1\), and it satisfies
\(\AbsoluteValue{\mathcal Y(\sigma,\tau)}\le\log2\).  The function
\(\mathcal Y\) is holomorphic on \(D\); its Taylor series
\(\mathcal Y=\sum_{a,b\ge0}c_{ab}\sigma^{a}\tau^{b}\) converges absolutely
at every point of \(D\); the coefficients satisfy
\begin{equation}\label{plt:eq:lw-cauchy-bound}
 \AbsoluteValue{c_{ab}}\le5^{\,a+b}
 \qquad(a,b\ge0);
\end{equation}
and for every \(u\in\ComplexNumbers\) and every \(n\ge0\),
\begin{equation}\label{plt:eq:lw-c-to-P}
 \sum_{a+b=n}c_{ab}u^{b}
 =\begin{cases}
   0,&n=0,\\
   \LamP{n}(u),&n\ge1.
  \end{cases}
\end{equation}
Consequently \(c_{ab}=\lambda_{a+b,b}\) for \(a+b\ge1\), where
\(\LamP{n}(u)=\sum_{m}\lambda_{n,m}u^{m}\), and for all
\((\sigma,\tau)\in D\) and \(n\ge1\),
\begin{equation}\label{plt:eq:lw-block-bound}
 \Bigl\lvert\textstyle\sum_{a+b=n}c_{ab}\sigma^{a}\tau^{b}\Bigr\rvert
 \le\bigl(5(\AbsoluteValue{\sigma}+\AbsoluteValue{\tau})\bigr)^{n}.
\end{equation}
\end{lemma}

\begin{proof}
\emph{Existence, uniqueness and the bound.}  Fix \((\sigma,\tau)\in D\) and
put \(\theta=\AbsoluteValue{\sigma}+\AbsoluteValue{\tau}<\tfrac12\).  For
\(\AbsoluteValue{y}\le1\) we have
\(\AbsoluteValue{\sigma y-\tau}\le\AbsoluteValue{\sigma}
+\AbsoluteValue{\tau}=\theta<\tfrac12\), so
\(\zeta\DefinitionEquals1-\tau+\sigma y\) lies in the open disc of radius
\(\tfrac12\) about \(1\), where the principal logarithm is holomorphic and,
by the Maclaurin series \(\PrincipalLogarithm(1+\nu)
=\sum_{m\ge1}(-1)^{m+1}\nu^{m}/m\),
\begin{equation}\label{plt:eq:lw-log-bound}
 \AbsoluteValue{\PrincipalLogarithm(1+\nu)}
 \le\sum_{m\ge1}\frac{\AbsoluteValue{\nu}^{m}}{m}
 =-\log\bigl(1-\AbsoluteValue{\nu}\bigr)
 \qquad(\AbsoluteValue{\nu}<1).
\end{equation}
Hence the map \(\Phi(y)=-\PrincipalLogarithm(1-\tau+\sigma y)\) sends the
closed unit disc into the closed disc of radius
\(-\log(1-\theta)\le\log2<1\).  Moreover
\[
 \AbsoluteValue{\Phi'(y)}
 =\left\lvert\frac{\sigma}{1-\tau+\sigma y}\right\rvert
 \le\frac{\AbsoluteValue{\sigma}}{1-\theta}
 \le\frac{\theta}{1-\theta}<1 ,
\]
so \(\Phi\) is a contraction of the complete metric space
\(\SetOf{\AbsoluteValue{y}\le1}\).  Banach's fixed point theorem gives a
unique fixed point \(\mathcal Y(\sigma,\tau)\), and it lies in the image, so
\(\AbsoluteValue{\mathcal Y}\le-\log(1-\theta)\le\log2\).

\emph{Holomorphy.}  Let \(y_0\equiv0\) and \(y_{j+1}=\Phi(y_j)\).  Each
\(y_j\) is holomorphic on \(D\) as a function of \((\sigma,\tau)\), by
induction: \(1-\tau+\sigma y_j\) is holomorphic with values in
\(\SetOf{\AbsoluteValue{\zeta-1}<\tfrac12}\), on which
\(\PrincipalLogarithm\) is holomorphic.  On a compact
\(E\subset D\) we have
\(\sup_{E}\theta=\theta_E<\tfrac12\), hence
\(\AbsoluteValue{y_{j+1}-y_j}\le
\bigl(\theta_E/(1-\theta_E)\bigr)^{j}\AbsoluteValue{y_1-y_0}\)
with \(\theta_E/(1-\theta_E)<1\), so \((y_j)\) converges uniformly on \(E\)
to \(\mathcal Y\).  A locally uniform limit of holomorphic functions is
holomorphic.

\emph{Taylor expansion and Cauchy bounds.}  Let \((\sigma_0,\tau_0)\in D\)
and choose \(\rho_1>\AbsoluteValue{\sigma_0}\),
\(\rho_2>\AbsoluteValue{\tau_0}\) with \(\rho_1+\rho_2<\tfrac12\).  The
closed polydisc of polyradius \((\rho_1,\rho_2)\) is contained in \(D\)
and \(\AbsoluteValue{\mathcal Y}\le1\) there, so the Cauchy integral formula
on the distinguished boundary yields the absolutely and uniformly convergent
expansion \(\mathcal Y=\sum_{a,b}c_{ab}\sigma^{a}\tau^{b}\) on the open
polydisc, with \(\AbsoluteValue{c_{ab}}\le\rho_1^{-a}\rho_2^{-b}\).  Since
\((\sigma_0,\tau_0)\) was arbitrary in \(D\), the series converges
absolutely at every point of \(D\), and the coefficients \(c_{ab}\) do not
depend on the polydisc.  Taking \(\rho_1=\rho_2=\tfrac15\), which satisfies
\(\rho_1+\rho_2=\tfrac25<\tfrac12\), gives \cref{plt:eq:lw-cauchy-bound}.

\emph{Identification of the diagonal sums.}  Fix \(u\in\ComplexNumbers\) and
restrict to the line \(\tau=\sigma u\); the restriction lies in \(D\)
for \(\AbsoluteValue{\sigma}<\bigl(2(1+\AbsoluteValue{u})\bigr)^{-1}\).  On
that disc, \(s\mapsto\mathcal Y(s,su)\) is holomorphic, vanishes at
\(s=0\), and satisfies
\(\mathcal Y=-\PrincipalLogarithm(1-su+s\mathcal Y)\); hence its Taylor
coefficients form the unique element of \(s\ComplexNumbers[[s]]\)
satisfying \cref{plt:eq:lw-generating} at the given \(u\).  By the formal
uniqueness clause of \cref{plt:thm:om-generating-equation}, those
coefficients are \(\LamP{n}(u)\).  On the other hand, absolute convergence
permits the rearrangement
\(\mathcal Y(s,su)=\sum_{n\ge0}\bigl(\sum_{a+b=n}c_{ab}u^{b}\bigr)s^{n}\).
Comparing coefficients gives \cref{plt:eq:lw-c-to-P} for the given \(u\);
as \(u\) ranges over \(\ComplexNumbers\), two polynomials agreeing
everywhere are equal, and comparing coefficients of \(u^{b}\) gives
\(c_{ab}=\lambda_{a+b,b}\).

\emph{The block bound.}  For \(x,y\ge0\), \(\sum_{a+b=n}x^{a}y^{b}
\le\sum_{a+b=n}\binom{n}{a}x^{a}y^{b}=(x+y)^{n}\).  Combining this with
\cref{plt:eq:lw-cauchy-bound} gives \cref{plt:eq:lw-block-bound}.
\end{proof}

\begin{theorem}[Convergent branch expansion from an arbitrary base point]
\label{plt:thm:lw-basepoint}
Let \(\xi\in\ComplexNumbers\).  Let \(K\in\ComplexNumbers\setminus
\SetOf{0}\), let \(\log K\) be any logarithm of \(K\), and put
\begin{equation}\label{plt:eq:lw-M}
 M\DefinitionEquals K+\log K-\xi .
\end{equation}
Assume the \emph{seating condition}
\begin{equation}\label{plt:eq:lw-seating}
 2\bigl(1+\AbsoluteValue{M}\bigr)<\AbsoluteValue{K}.
\end{equation}
Then:
\begin{enumerate}
\item the series \(\sum_{n\ge0}\LamP{n}(M)K^{-n}\) converges absolutely;
\item the number
      \[
       w\DefinitionEquals K+\sum_{n\ge0}\frac{\LamP{n}(M)}{K^{n}}
       =K-M+\sum_{n\ge1}\frac{\LamP{n}(M)}{K^{n}}
      \]
      satisfies \(\AbsoluteValue{w-(K-M)}\le\log2\) and \(w\ne0\);
\item with \(\log_{*}w\DefinitionEquals\log K
      +\PrincipalLogarithm(w/K)\), which is a logarithm of \(w\),
      \[
       w+\log_{*}w=\xi,
       \qquad\text{hence}\qquad
       w\EulerE^{w}=\EulerE^{\xi}.
      \]
\end{enumerate}
\end{theorem}

\begin{proof}
Set \(\sigma=1/K\) and \(\tau=M/K\).  Then
\(\AbsoluteValue{\sigma}+\AbsoluteValue{\tau}
=(1+\AbsoluteValue{M})/\AbsoluteValue{K}<\tfrac12\) by
\cref{plt:eq:lw-seating}, so \((\sigma,\tau)\in D\).  Let
\(\mathcal Y=\mathcal Y(\sigma,\tau)\) be as in
\cref{plt:lem:lw-two-variable}.  Applying \cref{plt:eq:lw-c-to-P} with
\(u=M\) and using \(\tau=\sigma M\),
\[
 \mathcal Y=\sum_{n\ge1}\LamP{n}(M)\sigma^{n}
 =\sum_{n\ge1}\frac{\LamP{n}(M)}{K^{n}},
\]
the series converging absolutely by \cref{plt:lem:lw-two-variable}.  Since
\(\LamP{0}(M)=-M\), this is (i), and
\(w=K-M+\mathcal Y\) with
\(\AbsoluteValue{\mathcal Y}\le\log2\), which is the first half of (ii).
For \(w\ne0\): \(\AbsoluteValue{w}\ge\AbsoluteValue{K}-\AbsoluteValue{M}
-\log2>2+\AbsoluteValue{M}-\log2>1\), using
\cref{plt:eq:lw-seating}.

For (iii), compute
\[
 \frac wK=\frac{K-M+\mathcal Y}{K}
 =1-\tau+\sigma\mathcal Y ,
\]
which lies in the disc of radius \(\tfrac12\) about \(1\) by the proof of
\cref{plt:lem:lw-two-variable}; hence \(\PrincipalLogarithm(w/K)\) is
defined and, by \cref{plt:eq:lw-two-var-fixed},
\(\PrincipalLogarithm(w/K)=-\mathcal Y\).  Therefore
\(\log_{*}w=\log K-\mathcal Y\), and
\(\EulerE^{\log_{*}w}=K\cdot(w/K)=w\), so \(\log_{*}w\) is indeed a
logarithm of \(w\).  Finally
\[
 w+\log_{*}w=(K-M+\mathcal Y)+(\log K-\mathcal Y)
 =K+\log K-M=\xi
\]
by \cref{plt:eq:lw-M}, and exponentiating gives
\(w\EulerE^{w}=\EulerE^{w+\log_{*}w}=\EulerE^{\xi}\).
\end{proof}

\begin{theorem}[Explicit uniform remainder]
\label{plt:thm:lw-uniform-remainder}
In the situation of \cref{plt:thm:lw-basepoint}, put
\begin{equation}\label{plt:eq:lw-theta}
 \theta\DefinitionEquals\frac{5\bigl(1+\AbsoluteValue{M}\bigr)}
 {\AbsoluteValue{K}}
\end{equation}
and assume \(\theta\le\tfrac12\), that is
\(\AbsoluteValue{K}\ge10(1+\AbsoluteValue{M})\).  Then for every
\(N\in\NaturalNumbers\),
\begin{equation}\label{plt:eq:lw-remainder}
 \left\lvert
 w-K-\sum_{n=0}^{N}\frac{\LamP{n}(M)}{K^{n}}
 \right\rvert
 \le2\,\theta^{\,N+1}
 =2\left(\frac{5(1+\AbsoluteValue{M})}{\AbsoluteValue{K}}\right)^{N+1}.
\end{equation}
The bound is uniform: its constants \(5\) and \(2\) are absolute, and no
hypothesis beyond \(\AbsoluteValue{K}\ge10(1+\AbsoluteValue{M})\) is used.
\end{theorem}

\begin{proof}
With \(\sigma=1/K\), \(\tau=M/K\) as above,
\(5(\AbsoluteValue{\sigma}+\AbsoluteValue{\tau})=\theta\), so
\cref{plt:eq:lw-block-bound} reads
\(\AbsoluteValue{\LamP{n}(M)K^{-n}}\le\theta^{n}\) for \(n\ge1\).  Since
\(\theta\le\tfrac12\),
\[
 \left\lvert\sum_{n>N}\frac{\LamP{n}(M)}{K^{n}}\right\rvert
 \le\sum_{n>N}\theta^{n}
 =\frac{\theta^{N+1}}{1-\theta}\le2\theta^{N+1}. \qedhere
\]
\end{proof}

Specialising to the canonical seating \(K=\xi_k\) gives the statement the
notation catalogue asks for, now with every hypothesis printed and every
constant explicit.

\begin{corollary}[Canonical branch expansion]
\label{plt:cor:lw-canonical}
Let \(z\ne0\), \(k\in\IntegerNumbers\), and let \((\xi_k,\ell_k)\) be a
branch datum with \(\xi_k\ne0\).  Take \(K=\xi_k\) and \(\log K=\ell_k\) in
\cref{plt:thm:lw-basepoint}; then \(M=\ell_k\).  Consequently:
\begin{enumerate}
\item if \(2(1+\AbsoluteValue{\ell_k})<\AbsoluteValue{\xi_k}\), the series
      \[
       w=\xi_k-\ell_k+\sum_{n\ge1}\frac{\LamP{n}(\ell_k)}{\xi_k^{\,n}}
      \]
      converges absolutely and defines a solution of
      \(w\EulerE^{w}=z\) with \(\log_{*}w=\xi_k-w\);
\item if moreover \(\AbsoluteValue{\xi_k}\ge10(1+\AbsoluteValue{\ell_k})\),
      then for every \(N\),
      \[
       \left\lvert w-\xi_k+\ell_k
        -\sum_{n=1}^{N}\frac{\LamP{n}(\ell_k)}{\xi_k^{\,n}}\right\rvert
       \le2\left(\frac{5(1+\AbsoluteValue{\ell_k})}
        {\AbsoluteValue{\xi_k}}\right)^{N+1}
       =\BigOAt{\ell_k^{\,N+1}/\xi_k^{\,N+1}}
        {\AbsoluteValue{\xi_k}\to\infty,\ \AbsoluteValue{\ell_k}\ge1};
      \]
\item in particular, if \(\ell_k=\PrincipalLogarithm\xi_k\), both
      hypotheses hold as soon as \(\AbsoluteValue{\xi_k}\ge100\), and for
      real \(z\ge\EulerE^{6}\) with \(k=0\) hypothesis \textup{(i)} holds
      while for real \(z\ge\EulerE^{50}\) hypothesis \textup{(ii)} holds.
\end{enumerate}
\end{corollary}

\begin{proof}
(i) and (ii) are \cref{plt:thm:lw-basepoint,plt:thm:lw-uniform-remainder}
with \(M=\xi_k+\ell_k-\xi_k=\ell_k\); the \(\BigOAt{\cdot}{\cdot}\) form of
the bound in (ii) follows because
\((1+\AbsoluteValue{\ell_k})\le2\AbsoluteValue{\ell_k}\) once
\(\AbsoluteValue{\ell_k}\ge1\), so the displayed bound is at most
\(2\cdot10^{\,N+1}\AbsoluteValue{\ell_k/\xi_k}^{N+1}\).

(iii)  If \(\ell_k=\PrincipalLogarithm\xi_k\) then
\(\AbsoluteValue{\ell_k}\le\log\AbsoluteValue{\xi_k}+\pi\).  Put
\(R=\AbsoluteValue{\xi_k}\).  Hypothesis (ii) holds if
\(R\ge10(1+\pi+\log R)\); the function \(R\mapsto R-10(1+\pi+\log R)\) has
derivative \(1-10/R>0\) for \(R>10\) and is positive at \(R=100\)
(indeed \(10(1+\pi+\log100)<87.5\)), hence for all \(R\ge100\); and
hypothesis (i) is weaker.  For real \(z>1\) and \(k=0\), \(\xi_0=\log z>0\)
and \(\ell_0=\log\log z\) are real, so (i) requires
\(2(1+\AbsoluteValue{\log\xi_0})<\xi_0\); at \(\xi_0=6\) one has
\(2(1+\log6)<5.59<6\), and \(\xi_0\mapsto\xi_0-2-2\log\xi_0\) is increasing
for \(\xi_0>2\).  Similarly (ii) requires \(\xi_0\ge10(1+\log\xi_0)\), which
holds at \(\xi_0=50\) (\(10(1+\log50)<49.2\)) and hence for all
\(\xi_0\ge50\) since \(\xi_0\mapsto\xi_0-10-10\log\xi_0\) is increasing for
\(\xi_0>10\).
\end{proof}

\begin{proposition}[Identification of the branch label]
\label{plt:prop:lw-branch-label}
Let \(z\ne0\), \(k\in\IntegerNumbers\), and let \((\xi_k,\ell_k)\) be a
branch datum with \(2(1+\AbsoluteValue{\ell_k})<\AbsoluteValue{\xi_k}\).
Let \(w=w_k(z)\) be the solution produced by
\cref{plt:cor:lw-canonical}\,(i), and assume in addition
\begin{equation}\label{plt:eq:lw-cut-distance}
 \AbsoluteValue{\operatorname{Im}\ell_k}\le\pi-\tfrac34 ,
\end{equation}
that is, \(\xi_k\) is at angular distance at least \(\tfrac34\) from the cut
of the chosen determination of the second logarithm.  Then
\(\log_{*}w=\PrincipalLogarithm w\), so
\begin{equation}\label{plt:eq:lw-strip}
 w+\PrincipalLogarithm w=\xi_k,
 \qquad
 \operatorname{Im}w\in
 \bigl[\,2\pi k+\operatorname{Arg}z-\pi,\;
        2\pi k+\operatorname{Arg}z+\pi\,\bigr).
\end{equation}
In particular distinct \(k\) produce solutions in disjoint horizontal
strips, so \(k\mapsto w_k(z)\) is injective.  Independently of
\cref{plt:eq:lw-cut-distance}, \(w_k(z)\) is the unique solution of
\(w\EulerE^{w}=z\) in the closed disc
\(\AbsoluteValue{v-(\xi_k-\ell_k)}\le1\).
\end{proposition}

\begin{proof}
By the proof of \cref{plt:thm:lw-basepoint},
\(\log_{*}w=\ell_k-\mathcal Y\) with
\(\AbsoluteValue{\mathcal Y}\le\log2<\tfrac34\).  Hence
\(\AbsoluteValue{\operatorname{Im}\log_{*}w}
\le\AbsoluteValue{\operatorname{Im}\ell_k}+\log2
<(\pi-\tfrac34)+\tfrac34=\pi\).
A logarithm of \(w\) whose imaginary part lies in \((-\pi,\pi)\) is the
principal one, so \(\log_{*}w=\PrincipalLogarithm w\).  Then
\cref{plt:thm:lw-basepoint}\,(iii) reads
\(w+\PrincipalLogarithm w=\xi_k\); taking imaginary parts and using
\(\operatorname{Im}\PrincipalLogarithm w=\operatorname{Arg}w\in(-\pi,\pi]\)
and \(\operatorname{Im}\xi_k=\operatorname{Arg}z+2\pi k\) gives
\(\operatorname{Im}w=\operatorname{Arg}z+2\pi k-\operatorname{Arg}w\), which
lies in the stated interval.  Those intervals for distinct \(k\) are
disjoint half-open intervals of length \(2\pi\), which gives injectivity.

For the disc uniqueness, let \(v\EulerE^{v}=z\) with
\(\AbsoluteValue{v-(\xi_k-\ell_k)}\le1\), and put
\(\sigma=1/\xi_k\), \(\tau=\ell_k/\xi_k\), so
\((\sigma,\tau)\in D\) by the seating condition.  Then
\(\AbsoluteValue{v-\xi_k}\le\AbsoluteValue{\ell_k}+1
\le\tfrac12\AbsoluteValue{\xi_k}\), so
\(\AbsoluteValue{v/\xi_k-1}\le\tfrac12\) and
\(y\DefinitionEquals-\PrincipalLogarithm(v/\xi_k)\) is defined with
\(\AbsoluteValue{y}\le\log2\) by \cref{plt:eq:lw-log-bound}.  Now
\(\lambda\DefinitionEquals\ell_k-y\) satisfies
\(\EulerE^{\lambda}=\xi_k\cdot(v/\xi_k)=v\), so
\cref{plt:prop:lw-log-reduction} gives \(v+\lambda=\xi_k
+2\pi\ImaginaryUnit m\) for some \(m\in\IntegerNumbers\), that is
\(v=\xi_k-\ell_k+y+2\pi\ImaginaryUnit m\).  Since
\(\AbsoluteValue{v-(\xi_k-\ell_k)}\le1\) and
\(\AbsoluteValue{y}\le\log2\), we get
\(2\pi\AbsoluteValue{m}\le1+\log2<2\pi\), hence \(m=0\) and
\(v=\xi_k-\ell_k+y\).  Dividing by \(\xi_k\) gives
\(v/\xi_k=1-\tau+\sigma y\), so \(y=-\PrincipalLogarithm(1-\tau+\sigma y)\):
\(y\) is a fixed point of \cref{plt:eq:lw-two-var-fixed} in
\(\AbsoluteValue{y}\le1\).  By the uniqueness in
\cref{plt:lem:lw-two-variable}, \(y=\mathcal Y(\sigma,\tau)\), whence
\(v=w_k(z)\).
\end{proof}

% ed.: S1, S3, S4 and S5 all state the branch expansion "in an admissible
% sector" or "in a branch-compatible sector" without ever defining one, and
% none proves that the constructed solution is the branch conventionally
% labelled W_k.  Repaired: the estimate is proved on an explicit region with
% no sector at all (Theorem lw-uniform-remainder), and the labelling is a
% separate proposition with a printed distance-to-the-cut hypothesis
% (Proposition lw-branch-label).
\begin{remark}[No sector is needed for the estimate]
\label{plt:rmk:lw-no-sector}
The remainder bound \cref{plt:eq:lw-remainder} holds pointwise on the
explicitly described region
\(\SetOf{\AbsoluteValue{\xi_k}\ge10(1+\AbsoluteValue{\ell_k})}\), with no
sector hypothesis at all.  The sector enters for two other purposes, and it
is worth separating them:
\begin{enumerate}
\item to make \(z\mapsto\xi_k\) and \(\xi_k\mapsto\ell_k\) single-valued
      analytic functions, so that \(w_k\) is an analytic function of \(z\)
      rather than a value attached to a branch datum;
\item to identify \(w_k\) with the function that the standard references
      call \(\LambertW_k\); this is where the distance-to-the-cut hypothesis
      \cref{plt:eq:lw-cut-distance} is used.
\end{enumerate}
Sources that write ``in an admissible sector'' without saying which are
conflating these two roles with the estimate itself.  With
\cref{plt:cor:lw-canonical,plt:prop:lw-branch-label} one may take, for
instance,
% ed.: the exhibited S_k was written with non-strict inequalities, hence was
% not open, while a coordinate sector is required to be open by
% Definition lw-coordinates.  Made strict.
\[
 S_k=\SetOf{z:\AbsoluteValue{\operatorname{Arg}z}<\pi,\
 \AbsoluteValue{\xi_k}>100,\
 \AbsoluteValue{\operatorname{Arg}\xi_k}<\pi-\tfrac34},
\]
on which \(\xi_k\) and \(\ell_k=\PrincipalLogarithm\xi_k\) are analytic,
\(w_k\) is analytic, and all hypotheses hold.
\end{remark}

% ed.: S6 deduces convergence of the Lambert series from (a) convergence at
% frozen u and (b) |s u| -> 0 along the Lambert path.  That inference is
% invalid: the radius of convergence in s depends on u and shrinks as
% |u| -> infinity.  Repaired by proving the two-variable statement
% (Lemma lw-two-variable), which is what the argument actually needs.
\begin{remark}[The gap in the convergence argument of S6]
\label{plt:rmk:lw-s6-gap}
S6 argues: the implicit function theorem shows that for each fixed \(u\) the
series \(\sum_{n}\LamP{n}(u)s^{n}\) converges for
\(\AbsoluteValue{s}\) small; along the Lambert path \(s=1/\xi\),
\(u=\ell\), one has \(\AbsoluteValue{su}\to0\); hence the series converges.
The conclusion does not follow from the premises.  The radius of convergence
in \(s\) supplied by the implicit function theorem depends on \(u\) and
shrinks as \(\AbsoluteValue{u}\to\infty\), so a statement about
\(\AbsoluteValue{su}\) is not a statement about \(\AbsoluteValue{s}\) lying
inside that radius.  What repairs the argument is exactly the separation of
variables in \cref{plt:lem:lw-two-variable}: one must apply the implicit
function theorem in the two variables \((\sigma,\tau)=(1/\xi,\ell/\xi)\),
\emph{both} of which are small, rather than in \(s\) alone with \(u\)
frozen.  The frozen-\(u\) statement is true --- it is
\cref{plt:prop:om-frozen-convergence} --- but it is not sufficient.
\end{remark}

% ed.: S1, S3 and S5 each assert, as though it were established here, that
% the principal-branch series converges for every real z >= e.  Downgraded
% to an attributed literature report: nothing in this volume proves it, and
% S5's supporting observation (all terms vanish at z = e) is vacuous.
\begin{warningbox}[title={What is proved here, and what is only reported}]
\Cref{plt:cor:lw-canonical} proves absolute convergence of the canonical
series for real \(z\ge\EulerE^{6}\approx403\), and the explicit uniform
remainder for \(z\ge\EulerE^{50}\).  Reports S1, S3 and S5 all state, citing
\cite{dlmf-W,corless-et-al-1996}, that for the principal branch and real
\(z\) the series converges for every \(z\ge\EulerE\).  That claim is far
stronger than anything proved here and we have not verified it; it is
recorded as a literature report, not as a theorem of this volume.  Note that
the check offered by S5 --- ``at \(z=\EulerE\) one has \(\ell_0=0\), so
every term vanishes'' --- is no evidence for it at all: it shows only that
the series is trivially convergent at the single point where all its terms
are zero (\cref{plt:prop:lw-exact-at-e}).  Similarly, S5's report of the
convergence domains \(X>(\alpha\EulerE)^{\alpha}\) for \(\alpha\ge1\) and
\(X>\EulerE\) for \(\alpha<1\) in the generalized problem
\cite{jeffrey-et-al-1995} is repeated in
\cref{plt:sec:generalized-lambert} as attribution only.
\end{warningbox}

\section{The principal real specialization}
\label{plt:sec:lw-principal-real}

\begin{corollary}[Principal real branch at infinity]
\label{plt:cor:lw-principal-real}
Let \(z>1\) be real, and take \(k=0\), so that
\(\xi_0=\log z>0\) and \(\ell_0=\log\log z\) are real.  Write
\(L_1=\log z\) and \(L_2=\log\log z\).  Then for \(z\ge\EulerE^{50}\),
\begin{align}
 \LambertW_0(z)={}&L_1-L_2+\frac{L_2}{L_1}
 +\frac{L_2(L_2-2)}{2L_1^{2}}
 +\frac{L_2(2L_2^{2}-9L_2+6)}{6L_1^{3}}
 \notag\\
 &+\frac{L_2(3L_2^{3}-22L_2^{2}+36L_2-12)}{12L_1^{4}}
 \notag\\
 &+\frac{L_2(12L_2^{4}-125L_2^{3}+350L_2^{2}-300L_2+60)}{60L_1^{5}}
 +\varrho_5 ,
 \label{plt:eq:lw-principal-five}
\end{align}
with the explicit bound
\(\AbsoluteValue{\varrho_5}\le2\bigl(5(1+L_2)/L_1\bigr)^{6}\), and the
infinite series converges absolutely to \(\LambertW_0(z)\) already for
\(z\ge\EulerE^{6}\).
\end{corollary}

\begin{proof}
For real \(z>1\) the datum \((\xi_0,\ell_0)=(\log z,\log\log z)\) satisfies
\cref{plt:eq:lw-cut-distance} with \(\operatorname{Im}\ell_0=0\), so
\cref{plt:prop:lw-branch-label} applies and identifies the constructed
solution with the real principal branch: \(w_0(z)\) is real (the whole
construction is real for real data, because
\(\mathcal Y(\sigma,\tau)\) is real for real \((\sigma,\tau)\in D\) by
uniqueness of the fixed point in \(\SetOf{\AbsoluteValue{y}\le1}\)) and
\(w_0>1\), hence \(w_0=\LambertW_0(z)\) by
\cref{plt:prop:lw-real-branches}\,(iii).  The five displayed corrections are
\(\LamP{1},\dots,\LamP{5}\) evaluated at \(L_2\), in the factored forms of
\cref{plt:sec:lambert-tables}, and the bound on \(\varrho_5\) is
\cref{plt:cor:lw-canonical}\,(ii) with \(N=5\).  Convergence for
\(z\ge\EulerE^{6}\) is \cref{plt:cor:lw-canonical}\,(iii).
\end{proof}

\begin{proposition}[Exactness at \texorpdfstring{\(z=\EulerE\)}{z=e}, and
what it does not mean]
\label{plt:prop:lw-exact-at-e}
At \(z=\EulerE\) one has \(L_1=1\), \(L_2=0\), and every truncation
\(\xi_0-\ell_0+\sum_{n=1}^{N}\LamP{n}(\ell_0)\xi_0^{-n}\) equals \(1\),
which is the exact value \(\LambertW_0(\EulerE)=1\).  The reason is that
\(\LamP{n}(0)=0\) for every \(n\in\NaturalNumbers\).  This gives no
information whatever about convergence, because at that point every term of
the series is \(0\); indeed the seating condition
\cref{plt:eq:lw-seating} reads \(2<1\) there and fails.
\end{proposition}

\begin{proof}
\(\LambertW_0(\EulerE)=1\) because \(1\cdot\EulerE^{1}=\EulerE\) and
\(1>-1\), so \cref{plt:prop:lw-real-branches} identifies the branch.  For
\(n\ge1\), \cref{plt:cor:lambert-coefficients} expresses \(\LamP{n}\) as a
sum over \(1\le m\le n\) of multiples of \(u^{m}\), so \(\LamP{n}(0)=0\);
and \(\LamP{0}(u)=-u\) gives \(\LamP{0}(0)=0\).  Hence every truncation is
\(1-0+0=1\).  For the last sentence, \(\AbsoluteValue{K}=\xi_0=1\) and
\(\AbsoluteValue{M}=\AbsoluteValue{\ell_0}=0\), so
\cref{plt:eq:lw-seating} asserts \(2<1\).
\end{proof}

\section{The lower real branch near zero}
\label{plt:sec:lw-lower}

The catalogue reserves the symbol \(\varepsilon_{\mathrm{br}}\) for the
positive small argument of the lower branch; it is not an error tolerance
and not an inverse iterated logarithm.  We write \(\varepsilon\) for it in
this subsection, where no other \(\varepsilon\) occurs.

\begin{proposition}[The lower-branch equation]
\label{plt:prop:lw-lower-equation}
The map \(h(u)=u-\log u\) is a strictly increasing bijection from
\((1,\infty)\) onto \((1,\infty)\).  For
\(0<\varepsilon<\EulerE^{-1}\), put
\begin{equation}\label{plt:eq:lw-lower-coords}
 S\DefinitionEquals\log(1/\varepsilon)>1,
 \qquad
 L\DefinitionEquals\log S .
\end{equation}
Then \(u\DefinitionEquals-\LambertW_{-1}(-\varepsilon)\) is the unique
solution in \((1,\infty)\) of
\begin{equation}\label{plt:eq:lw-lower-fixed}
 u-\log u=S .
\end{equation}
\end{proposition}

\begin{proof}
\(h'(u)=1-1/u>0\) on \((1,\infty)\), \(h(1)=1\), and \(h(u)\to\infty\); so
\(h\) is a strictly increasing bijection onto \((1,\infty)\).  By
\cref{plt:prop:lw-real-branches}, \(\LambertW_{-1}(-\varepsilon)\in
(-\infty,-1)\) for \(0<\varepsilon<\EulerE^{-1}\), so \(u>1\).  Writing
\(w=-u\) in \(w\EulerE^{w}=-\varepsilon\) gives
\(u\EulerE^{-u}=\varepsilon\); taking the real logarithm,
\(\log u-u=\log\varepsilon=-S\), which is
\cref{plt:eq:lw-lower-fixed}.  Uniqueness is the injectivity of \(h\).
\end{proof}

\begin{theorem}[Expansion of the lower real branch]
\label{plt:thm:lw-lower-branch}
With \(S,L\) as in \cref{plt:eq:lw-lower-coords}, and provided
\(2(1+\AbsoluteValue{L})<S\), the series below converges absolutely and
\begin{equation}\label{plt:eq:lw-lower-expansion}
 \LambertW_{-1}(-\varepsilon)
 =-S-L+\sum_{n\ge1}(-1)^{n}\frac{\LamP{n}(L)}{S^{n}},
\end{equation}
equivalently, in the zero-based normalisation and with
\(\LamP{0}(L)=-L\),
\begin{equation}\label{plt:eq:lw-lower-expansion-clean}
 \LambertW_{-1}(-\varepsilon)
 =-S+\sum_{n\ge0}(-1)^{n}\frac{\LamP{n}(L)}{S^{n}} .
\end{equation}
If in addition \(S\ge10(1+\AbsoluteValue{L})\), then for every \(N\),
\begin{equation}\label{plt:eq:lw-lower-remainder}
 \left\lvert\LambertW_{-1}(-\varepsilon)+S+L
 -\sum_{n=1}^{N}(-1)^{n}\frac{\LamP{n}(L)}{S^{n}}\right\rvert
 \le2\left(\frac{5(1+\AbsoluteValue{L})}{S}\right)^{N+1}.
\end{equation}
Explicitly,
\begin{align}
 \LambertW_{-1}(-\varepsilon)={}&-S-L-\frac{L}{S}
 +\frac{L(L-2)}{2S^{2}}
 -\frac{L(2L^{2}-9L+6)}{6S^{3}}
 \notag\\
 &+\frac{L(3L^{3}-22L^{2}+36L-12)}{12S^{4}}+\cdots.
 \label{plt:eq:lw-lower-first-terms}
\end{align}
\end{theorem}

\begin{proof}
This is the analytic affine-logarithmic theorem with \(a=-1\).  Concretely,
apply \cref{plt:lem:lw-two-variable} with
\(\sigma=-1/S\) and \(\tau=-L/S\); then
\(\AbsoluteValue{\sigma}+\AbsoluteValue{\tau}
=(1+\AbsoluteValue{L})/S<\tfrac12\) by hypothesis, so
\((\sigma,\tau)\in D\), and \(\tau=\sigma L\), so
\cref{plt:eq:lw-c-to-P} gives
\[
 \mathcal Y=\sum_{n\ge1}\LamP{n}(L)\sigma^{n}
 =\sum_{n\ge1}(-1)^{n}\frac{\LamP{n}(L)}{S^{n}},
\]
absolutely convergent.  Put \(v\DefinitionEquals S+L-\mathcal Y\).  Then
\[
 \frac vS=1+\frac LS-\frac{\mathcal Y}{S}
 =1-\tau+\sigma\mathcal Y ,
\]
which lies in the disc of radius \(\tfrac12\) about \(1\); hence \(v>0\) and
\(\PrincipalLogarithm(v/S)=-\mathcal Y\) by
\cref{plt:eq:lw-two-var-fixed}.  All the data being real, \(\mathcal Y\) is
real and \(\log(v/S)=-\mathcal Y\) is the real logarithm.  Therefore
\[
 v-\log v=v-\log S-\log(v/S)=S+L-\mathcal Y-L+\mathcal Y=S .
\]
Also \(v>1\): the hypothesis \(2(1+\AbsoluteValue{L})<S\) gives
\(\AbsoluteValue{L}<S/2-1\) and \(S>2\), hence
\(v\ge S+L-\AbsoluteValue{\mathcal Y}\ge S-\AbsoluteValue{L}-\log2
>S/2+1-\log2>2-\log2>1\).  By the uniqueness in
\cref{plt:prop:lw-lower-equation}, \(v=u=-\LambertW_{-1}(-\varepsilon)\),
which gives \cref{plt:eq:lw-lower-expansion-clean} after using
\(\LamP{0}(L)=-L\).  The remainder bound is
\cref{plt:eq:lw-block-bound} exactly as in
\cref{plt:thm:lw-uniform-remainder}, with
\(5(\AbsoluteValue{\sigma}+\AbsoluteValue{\tau})
=5(1+\AbsoluteValue{L})/S\le\tfrac12\).  The displayed first terms are
\(\LamP{1},\dots,\LamP{4}\) with the alternating signs.
\end{proof}

% ed.: every source displays the lower-branch expansion with the factor
% (-1)^n and explains it only by analogy ("the same polynomials reappear
% with alternating block signs").  Repaired: the alternation is derived as
% the single value a = -1 of the scalar a^{n+1} in Theorem lw-affine-log,
% and the box states explicitly that no second polynomial family exists.
\begin{warningbox}[title={The alternating signs do not define a second
Lambert family}]
\label{plt:rmk:lw-no-second-family}
Every source displays \cref{plt:eq:lw-lower-first-terms} and remarks that
``the same polynomials reappear with alternating block signs''.  The
statement deserves to be made exactly, because the alternation is sometimes
absorbed into the polynomials, producing a spurious second family
\(\widetilde{\LamP{n}}=(-1)^{n}\LamP{n}\) with its own tables and its own
recurrence.  There is no second family.  By
\cref{plt:thm:lw-affine-log}, the coefficient of \(X^{-n}\) in the inverse
of \(X+aL\) is \(a^{n+1}\LamP{n}(L)\), a \emph{scalar} multiple of the one
canonical polynomial; the lower real branch is the single value \(a=-1\) of
that scalar, so the coefficient of \(S^{-n}\) is
\((-1)^{n+1}\LamP{n}(L)\) in the equation \(u-\log u=S\) and
\((-1)^{n}\LamP{n}(L)\) after the reflection \(w=-u\).  The sign is data
about the argument, not about the polynomial.  Concretely: the recurrence
\(\LamP{n+1}=-\LamP{n}+n\int_{0}^{u}\LamP{n}\) of
\cref{plt:thm:lambert-recurrence}, the Stirling formula of
\cref{plt:thm:lambert-stirling} and the tables of
\cref{plt:sec:lambert-tables} are used unchanged for the lower branch; only
the evaluation point and the sign of the outer variable change.
\end{warningbox}

\begin{remark}[Why the lower branch is not literally an \(\WrightOmega\) value]
Setting \(y=aY\) in \(y+a\log y=X\) turns it into
\(Y+\log Y=X/a-\log a\), so for \(a>0\) the affine-logarithmic problem is a
rescaled Wright omega evaluation.  For \(a=-1\) this substitution requires
\(\log(-1)\), and the resulting omega argument is
\(-X\mp\pi\ImaginaryUnit\): the real lower branch corresponds to a
\emph{complex} point of \(\WrightOmega\).  This is why
\cref{plt:thm:lw-lower-branch} is proved directly from
\cref{plt:lem:lw-two-variable} rather than by quoting the omega expansion,
and it is the precise content of the catalogue's warning that
\(\varepsilon_{\mathrm{br}}\) is a branch-point-independent parameter for a
problem of its own.
\end{remark}

\section{The square-root branch point}
\label{plt:sec:lw-branch-point}

\begin{lemma}[The critical point of \texorpdfstring{\(w\EulerE^{w}\)}{w exp w}]
\label{plt:lem:lw-critical-point}
Let \(g(w)=w\EulerE^{w}\) and \(w=-1+q\).  Then
\begin{equation}\label{plt:eq:lw-critical-expansion}
 1+\EulerE\,g(-1+q)
 =\sum_{n\ge2}\frac{(n-1)\,q^{n}}{n!}
 =\frac{q^{2}}{2}+\frac{q^{3}}{3}+\frac{q^{4}}{8}
  +\frac{q^{5}}{30}+\frac{q^{6}}{144}+\cdots,
\end{equation}
an entire function of \(q\).  In particular \(g(-1)=-\EulerE^{-1}\),
\(g'(-1)=0\) and \(g''(-1)=\EulerE^{-1}\ne0\), so \(-1\) is a critical point
of \(g\) of order exactly two.
\end{lemma}

\begin{proof}
\(\EulerE\,g(-1+q)=\EulerE(-1+q)\EulerE^{-1+q}=(q-1)\EulerE^{q}\).  Now
\[
 (q-1)\EulerE^{q}
 =\sum_{n\ge0}\frac{q^{n+1}}{n!}-\sum_{n\ge0}\frac{q^{n}}{n!}
 =-1+\sum_{n\ge1}q^{n}\left(\frac1{(n-1)!}-\frac1{n!}\right)
 =-1+\sum_{n\ge1}\frac{(n-1)q^{n}}{n!},
\]
and the \(n=1\) summand vanishes.  Adding \(1\) gives
\cref{plt:eq:lw-critical-expansion}.  Comparing coefficients,
\(\EulerE\,g'(-1)=0\) and
\(\EulerE\,g''(-1)/2!=\tfrac12\), i.e.\ \(g''(-1)=\EulerE^{-1}\).
\end{proof}

% ed.: S1 and S3 assert the Puiseux expansion at the branch point without
% constructing the local inverse; S1 says only "ordinary series reversion in
% p produces" the coefficients, and neither source proves that a
% power-logarithmic expansion is impossible there.  Repaired: the local
% biholomorphism is constructed, the coefficients are rederived (and carried
% one order further than any source, to p^6 = -221/8505, verified against a
% high-precision evaluation), and the impossibility is proved from the order
% of the local monodromy in Corollary lw-branch-point-monodromy.
\begin{theorem}[Puiseux expansion at the branch point]
\label{plt:thm:lw-branch-point}
Put
\begin{equation}\label{plt:eq:lw-p}
 p\DefinitionEquals\sqrt{2\bigl(1+\EulerE z\bigr)} .
\end{equation}
There are \(r>0\) and a function \(\varphi\) holomorphic on
\(\AbsoluteValue{p}<r\) with \(\varphi(0)=0\), \(\varphi'(0)=1\), such that
for \(\AbsoluteValue{1+\EulerE z}<r^{2}/2\) the two solutions of
\(w\EulerE^{w}=z\) near \(w=-1\) are exactly
\begin{equation}\label{plt:eq:lw-branchpoint-solutions}
 w=-1+\varphi(p)
 \qquad\text{and}\qquad
 w=-1+\varphi(-p),
\end{equation}
for either determination of the square root in \cref{plt:eq:lw-p}.  Its
expansion begins
\begin{equation}\label{plt:eq:lw-branchpoint-series}
 \varphi(p)=p-\frac{p^{2}}{3}+\frac{11p^{3}}{72}
 -\frac{43p^{4}}{540}+\frac{769p^{5}}{17280}
 -\frac{221p^{6}}{8505}+\cdots.
\end{equation}
In particular the leading behaviour is
\begin{equation}\label{plt:eq:lw-branchpoint-leading}
 w+1=\pm\sqrt{2(1+\EulerE z)}
 +\BigOAt{\AbsoluteValue{1+\EulerE z}}{z\to-\EulerE^{-1}} .
\end{equation}
For real \(z\) slightly larger than \(-\EulerE^{-1}\) and \(p>0\), the
branch \(\varphi(p)\) is \(\LambertW_0\) and \(\varphi(-p)\) is
\(\LambertW_{-1}\); the two are exchanged by \(p\mapsto-p\).
\end{theorem}

\begin{proof}
Write \cref{plt:eq:lw-critical-expansion} as
\(1+\EulerE z=\tfrac12q^{2}h(q)\) with
\[
 h(q)\DefinitionEquals\sum_{n\ge2}\frac{2(n-1)}{n!}q^{\,n-2}
 =1+\frac{2q}{3}+\frac{q^{2}}{4}+\frac{q^{3}}{15}
  +\frac{q^{4}}{72}+\cdots,
\]
entire with \(h(0)=1\).  Choose \(r_0>0\) with \(h\ne0\) on
\(\AbsoluteValue{q}<r_0\) and let \(h^{1/2}\) be the holomorphic square root
there with \(h^{1/2}(0)=1\).  Put \(\psi(q)=q\,h^{1/2}(q)\), so
\(\psi(0)=0\), \(\psi'(0)=1\) and
\begin{equation}\label{plt:eq:lw-psi}
 \psi(q)^{2}=q^{2}h(q)=2(1+\EulerE z)
 \qquad\text{whenever } z=g(-1+q).
\end{equation}
By the holomorphic inverse function theorem \(\psi\) is a biholomorphism of
a disc \(\AbsoluteValue{q}<r_1\le r_0\) onto a neighbourhood of \(0\)
containing a disc \(\AbsoluteValue{p}<r\); let \(\varphi=\psi^{-1}\) on that
disc, so \(\varphi(0)=0\) and \(\varphi'(0)=1/\psi'(0)=1\).

Now fix \(z\) with \(\AbsoluteValue{1+\EulerE z}<r^{2}/2\) and let \(p\) be
either square root of \(2(1+\EulerE z)\), so \(\AbsoluteValue{p}<r\).  If
\(w=-1+q\) with \(\AbsoluteValue{q}<r_1\) solves \(g(w)=z\), then
\cref{plt:eq:lw-psi} gives \(\psi(q)^{2}=p^{2}\), hence
\(\psi(q)=\pm p\) and \(q=\varphi(\pm p)\); conversely each of
\(q=\varphi(\pm p)\) satisfies \(\tfrac12q^{2}h(q)
=\tfrac12\psi(q)^{2}=\tfrac12p^{2}=1+\EulerE z\), i.e.\ \(g(-1+q)=z\) by
\cref{plt:eq:lw-critical-expansion}.  Shrinking \(r\) so that
\(\AbsoluteValue{\varphi}<r_1\) on \(\AbsoluteValue{p}<r\) makes this a
bijection between the two square roots and the two nearby solutions; they
are distinct for \(p\ne0\) because
\(\varphi(p)-\varphi(-p)=2p+\BigOAt{p^{3}}{p\to0}\).

For the coefficients, write \(\varphi(p)=\sum_{i\ge1}a_ip^{i}\) with
\(a_1=1\) and determine \(a_2,a_3,\dots\) by substituting into
\(\psi(\varphi(p))=p\), where
\(\psi(q)=q+\tfrac13q^{2}+\tfrac{5}{72}q^{3}+\tfrac{11}{1080}q^{4}
+\tfrac{59}{51840}q^{5}+\cdots\) is obtained by expanding
\(h^{1/2}\) from the displayed \(h\).  Comparing coefficients of
\(p^{2},\dots,p^{6}\) gives in turn
\(a_2=-\tfrac13\), \(a_3=\tfrac{11}{72}\), \(a_4=-\tfrac{43}{540}\),
\(a_5=\tfrac{769}{17280}\), \(a_6=-\tfrac{221}{8505}\), which is
\cref{plt:eq:lw-branchpoint-series}.  Since \(\varphi(p)=p+
\BigOAt{p^{2}}{p\to0}\) and \(p^{2}=2(1+\EulerE z)\), we get
\cref{plt:eq:lw-branchpoint-leading}.

Finally let \(z\in(-\EulerE^{-1},-\EulerE^{-1}+\delta)\) be real and take
\(p=+\sqrt{2(1+\EulerE z)}>0\) small.  Then \(\varphi(p)>0\) and
\(\varphi(-p)<0\), so \(-1+\varphi(p)>-1\) and \(-1+\varphi(-p)<-1\);
by \cref{plt:prop:lw-real-branches}\,(iii) these are \(\LambertW_0(z)\) and
\(\LambertW_{-1}(z)\) respectively.
\end{proof}

\begin{corollary}[No power--logarithmic expansion at the branch point]
\label{plt:cor:lw-branch-point-monodromy}
Let \(z_0=-\EulerE^{-1}\).
\begin{enumerate}
\item There is no single-valued holomorphic function \(F\) on a punctured
      disc \(0<\AbsoluteValue{z-z_0}<\rho\) with \(g(F(z))=z\) and
      \(F(z)\to-1\).
\item The local monodromy of the two-valued germ at \(z_0\) has order
      exactly two.
\item Consequently no determination of the inverse near \(-1\) has, on a slit
      disc, a representation
      \[
       F(z)=\sum_{j=0}^{d}A_j(z)\bigl(\log(z-z_0)\bigr)^{j}
      \]
      with \(d\ge1\), the \(A_j\) single-valued holomorphic on the punctured
      disc, and \(A_d\not\equiv0\).  That is, no logarithm of
      \(z-z_0\) occurs at any power; the correct local scale is the Puiseux
      scale \(\SetOf{(z-z_0)^{m/2}}\), on which
      \cref{plt:thm:lw-branch-point} is an expansion.
\end{enumerate}
\end{corollary}

\begin{proof}
(ii)  Continuing \(p=\sqrt{2\EulerE(z-z_0)}\) once counterclockwise around
\(z_0\) replaces \(p\) by \(-p\), hence exchanges the two determinations
\(-1+\varphi(\pm p)\) of \cref{plt:thm:lw-branch-point}; continuing twice
restores them.  The exchange is nontrivial because the two determinations
differ for \(p\ne0\).  So the monodromy permutation has order two.

(i)  A single-valued \(F\) would be invariant under one continuation, but
the only two candidate values are exchanged; contradiction.

(iii)  Write \(\Lambda=\log(z-z_0)\) and suppose
\(F=\sum_{j=0}^{d}A_j\Lambda^{j}\) with \(d\ge1\) and \(A_d\not\equiv0\).
Continuation \(m\) times around \(z_0\) replaces \(\Lambda\) by
\(\Lambda+2\pi\ImaginaryUnit m\) and hence \(F\) by
\(F_m=\sum_{j}A_j(\Lambda+2\pi\ImaginaryUnit m)^{j}\).  Expanding, \(F_m\)
is again of the same shape with top coefficient \(A_d\) and next
coefficient \(A_{d-1}+2\pi\ImaginaryUnit m\,d\,A_d\).  Since the
representation of a germ in the form \(\sum_{j\le d}A_j\Lambda^{j}\) with
single-valued \(A_j\) is unique --- the functions \(1,\Lambda,\dots,
\Lambda^{d}\) are linearly independent over the field of single-valued
meromorphic germs, because \(\Lambda\) is transcendental over it --- the
germs \(F_m\), \(m\in\IntegerNumbers\), are pairwise distinct.  So the germ
would have infinitely many determinations, contradicting (ii).  If instead
\(F\) were given by a convergent series in integer powers of \(z-z_0\), it
would be single-valued, contradicting (i).
\end{proof}

\begin{remark}[The general principle]
\Cref{plt:cor:lw-branch-point-monodromy} is the local instance of a rule
used repeatedly in \cref{plt:sec:generalized-lambert}: the fractional-power
scale of an inverse is dictated by the order of vanishing of the derivative
of the direct map.  A simple zero of \(g'\) --- equivalently, a critical
point of order two --- forces the exponent \(1/2\) and monodromy of order
two.  A zero of order \(m-1\) would force exponent \(1/m\); an essential
degeneracy would force something else again.  Logarithms in an inverse come
from a different mechanism entirely, namely a logarithm in the direct map
(\cref{plt:prop:lag-loglog-forced}), and the two mechanisms never substitute
for one another.
\end{remark}

\section{Residual correction and one Newton step in branch coordinates}
\label{plt:sec:lw-residual}

The formal theory of \cref{plt:sec:reversion-existence} measures the quality
of a truncated inverse by the reversal residual
\(\RevErr{F}{H}=F\compose H-X\), and
\cref{plt:prop:rev-residual-order} shows that its \(t\)-order is exactly the
number of correct blocks.  In branch coordinates the analytic counterpart is
the \emph{logarithmic residual}, and it is a much better numerical object
than the multiplicative residual \(v\EulerE^{v}-z\), which overflows as soon
as \(z\) is large.

\begin{definition}[Truncations and the logarithmic residual]
\label{plt:def:lw-residual}
Let \((\xi,\ell)\) be a branch datum with
\(\AbsoluteValue{\xi}\ge10(1+\AbsoluteValue{\ell})\).  For
\(N\in\NaturalNumbers\) put
\begin{equation}\label{plt:eq:lw-truncation}
 w_N\DefinitionEquals
 \xi-\ell+\sum_{n=1}^{N}\frac{\LamP{n}(\ell)}{\xi^{n}},
\end{equation}
so that \(w_N\) is the value of the exclusive truncation
\(\Trunc{\rev{(X+L)}}{N+1}\) at \((X,L)=(\xi,\ell)\).  On the closed disc
\(\mathcal D\DefinitionEquals
\SetOf{v:\AbsoluteValue{v-(\xi-\ell)}\le1}\) define
\(\log_{*}v\DefinitionEquals\ell+\PrincipalLogarithm(v/\xi)\) and
\(F(v)\DefinitionEquals v+\log_{*}v\); set
\begin{equation}\label{plt:eq:lw-residual}
 R_N\DefinitionEquals F(w_N)-\xi
 =w_N+\log_{*}w_N-\xi,
 \qquad
 \delta_N\DefinitionEquals-\frac{R_N}{1+1/w_N}.
\end{equation}
\end{definition}

\begin{proposition}[The disc is admissible]
\label{plt:prop:lw-disc}
Under the hypothesis of \cref{plt:def:lw-residual}:
\(\mathcal D\) is contained in the region where \(\log_{*}\) is defined and
holomorphic; for \(v\in\mathcal D\) one has
\(\AbsoluteValue{v}\ge\tfrac9{10}\AbsoluteValue{\xi}\ge9\),
\(\tfrac89\le\AbsoluteValue{F'(v)}\le\tfrac{10}9\), and
\(\AbsoluteValue{F''(v)}\le\tfrac54\AbsoluteValue{\xi}^{-2}\); and both
\(w_N\) (for every \(N\)) and the exact solution \(w\) of
\cref{plt:cor:lw-canonical} lie in \(\mathcal D\).
\end{proposition}

\begin{proof}
Write \(\theta=5(1+\AbsoluteValue{\ell})/\AbsoluteValue{\xi}\le\tfrac12\).
For \(v\in\mathcal D\),
\(\AbsoluteValue{v-\xi}\le\AbsoluteValue{\ell}+1
\le\AbsoluteValue{\xi}/10\), so
\(\AbsoluteValue{v/\xi-1}\le\tfrac1{10}<1\) and
\(\PrincipalLogarithm(v/\xi)\) is defined and holomorphic there; also
\(\AbsoluteValue{v}\ge\tfrac9{10}\AbsoluteValue{\xi}\), and
\(\AbsoluteValue{\xi}\ge10(1+\AbsoluteValue{\ell})\ge10\) gives
\(\AbsoluteValue{v}\ge9\).  Since
\(F'(v)=1+1/v\) and \(F''(v)=-1/v^{2}\), we get
\(\AbsoluteValue{F'(v)-1}\le\tfrac19\) and
\(\AbsoluteValue{F''(v)}\le(10/9)^{2}\AbsoluteValue{\xi}^{-2}
\le\tfrac54\AbsoluteValue{\xi}^{-2}\).  Finally
\(\AbsoluteValue{w_N-(\xi-\ell)}
\le\sum_{n\ge1}\theta^{n}=\theta/(1-\theta)\le1\) by
\cref{plt:eq:lw-block-bound}, and
\(\AbsoluteValue{w-(\xi-\ell)}=\AbsoluteValue{\mathcal Y}\le\log2\le1\).
\end{proof}

\begin{theorem}[Residual order, residual-to-error transfer, and one Newton
step]
\label{plt:thm:lw-newton-step}
In the situation of \cref{plt:def:lw-residual}, put
\(\theta=5(1+\AbsoluteValue{\ell})/\AbsoluteValue{\xi}\le\tfrac12\) and let
\(w\) be the exact solution of \cref{plt:cor:lw-canonical}.  Then, with
\(e\DefinitionEquals w-w_N\):
\begin{enumerate}
\item \(\AbsoluteValue{e}\le2\theta^{N+1}\) and
      \(\AbsoluteValue{R_N}\le\tfrac{20}9\theta^{N+1}\); in particular
      \[
       R_N=\BigOAt{\ell^{\,N+1}/\xi^{\,N+1}}
       {\AbsoluteValue{\xi}\to\infty,\ \AbsoluteValue{\ell}\ge1},
      \]
      the same order as the error itself;
\item \(\AbsoluteValue{e}\le\tfrac54\AbsoluteValue{R_N}\)
      \textup{(}residual-to-error transfer, with an explicit
      constant\textup{)};
\item \(\displaystyle
      \AbsoluteValue{e-\delta_N}
      \le\frac32\,\frac{\AbsoluteValue{R_N}^{2}}
      {\AbsoluteValue{\xi}^{2}}\)
      \textup{(}one Newton step is quadratically better\textup{)}.
\end{enumerate}
\end{theorem}

\begin{proof}
(i)  The bound on \(e\) is \cref{plt:thm:lw-uniform-remainder}.  The segment
\([w_N,w]\) lies in the convex set \(\mathcal D\) by
\cref{plt:prop:lw-disc}, and \(F(w)=\xi\) by
\cref{plt:thm:lw-basepoint}\,(iii) --- note that on \(\mathcal D\) the
function \(\log_{*}\) of \cref{plt:def:lw-residual} agrees with the
\(\log_{*}\) of \cref{plt:thm:lw-basepoint} --- so
\(\AbsoluteValue{R_N}=\AbsoluteValue{F(w_N)-F(w)}
\le\sup_{\mathcal D}\AbsoluteValue{F'}\cdot\AbsoluteValue{e}
\le\tfrac{10}9\cdot2\theta^{N+1}\).  The \(\BigOAt{\cdot}{\cdot}\) form
follows as in \cref{plt:cor:lw-canonical}\,(ii).

(ii) and (iii).  Taylor's formula with integral remainder along
\([w_N,w]\subset\mathcal D\) gives
\[
 0=F(w)-\xi=R_N+F'(w_N)e+E,
 \qquad
 E=\int_{0}^{1}(1-s)F''\bigl(w_N+se\bigr)e^{2}\,\Differential s,
\]
so \(\AbsoluteValue{E}\le\tfrac12\cdot
\tfrac54\AbsoluteValue{\xi}^{-2}\AbsoluteValue{e}^{2}
=\tfrac58\AbsoluteValue{e}^{2}\AbsoluteValue{\xi}^{-2}\).  Since both
points lie in \(\mathcal D\) we have \(\AbsoluteValue{e}\le2\), and
\(\AbsoluteValue{\xi}\ge10\), so
\(\AbsoluteValue{E}\le\tfrac58\cdot2\AbsoluteValue{e}/100
=\AbsoluteValue{e}/80\).  Using \(\AbsoluteValue{F'(w_N)}\ge\tfrac89\),
\[
 \tfrac89\AbsoluteValue{e}\le\AbsoluteValue{R_N}+\tfrac1{80}
 \AbsoluteValue{e},
 \qquad\text{hence}\qquad
 \AbsoluteValue{e}\le\frac{\AbsoluteValue{R_N}}{8/9-1/80}
 \le\tfrac54\AbsoluteValue{R_N},
\]
which is (ii).  Dividing the Taylor identity by \(F'(w_N)=1+1/w_N\) gives
\(e=\delta_N-E/F'(w_N)\), so
\[
 \AbsoluteValue{e-\delta_N}
 \le\frac98\AbsoluteValue{E}
 \le\frac98\cdot\frac58\cdot
 \frac{\AbsoluteValue{e}^{2}}{\AbsoluteValue{\xi}^{2}}
 \le\frac{45}{64}\cdot\frac{25}{16}\,
 \frac{\AbsoluteValue{R_N}^{2}}{\AbsoluteValue{\xi}^{2}}
 \le\frac32\,\frac{\AbsoluteValue{R_N}^{2}}{\AbsoluteValue{\xi}^{2}}. \qedhere
\]
\end{proof}

\begin{remark}[Where the transfer fails]
Every constant above degrades as \(\AbsoluteValue{v}\) approaches \(0\), and
the argument collapses entirely at \(v=-1\), where \(F'(v)=1+1/v=0\).  That
is exactly the branch point of \cref{plt:sec:lw-branch-point}: the
residual-to-error transfer is a statement about a nonvanishing derivative,
and at a critical point a small residual is compatible with an error of the
order of its square root.  This is the analytic shadow of the algebraic fact
recorded in \cref{plt:lem:rev-newton-unit}: Newton reversion needs the
denominator to be a unit.
\end{remark}

\begin{remark}[Exact identities as independent checks]
Two exact relations cost nothing and catch branch mistakes that coefficient
checks do not.  On any branch, \(\log_{*}\LambertW_k(z)=\xi_k
-\LambertW_k(z)\) and \(\EulerE^{\LambertW_k(z)}=z/\LambertW_k(z)\).  The
first lets one expand \(\log\LambertW_k\) without composing a logarithm with
a series, and the second lets one check
\(\LambertW_k\EulerE^{\LambertW_k}=z\) by a division rather than an
exponentiation.  A truncation that satisfies the coefficient recurrences but
violates the first identity has a branch mismatch, not a coefficient error.
\end{remark}

\section{Numerical behaviour of the truncations}
\label{plt:sec:lw-numerics}

All six sources tabulate truncation errors, and the tables disagree ---
not in their numbers but in what their column headings mean.  S1 and S4
tabulate against \(x=\log z\), S3 and S5 against \(z\).  Both families of
numbers are correct, and each is a subset of the other's grid; we tabulate
once, in the intrinsic variable \(\xi_0=\log z\), and record the
concordance.

\begin{table}[htbp]
\centering
\caption{Absolute error
\(\AbsoluteValue{w_N-\LambertW_0(\EulerE^{\xi_0})}\) of the truncation
\cref{plt:eq:lw-truncation} on the principal real branch,
\(\ell_0=\log\xi_0\).  Recomputed here to \(60\) significant digits and in
agreement with the tables of S1, S3, S4 and S5 wherever those overlap.}
\label{plt:tab:lw-errors}
\begin{tabular}{@{}rllll@{}}
\toprule
\(N\) & \(\xi_0=10\) & \(\xi_0=25\) & \(\xi_0=50\) & \(\xi_0=100\)\\
\midrule
\(0\)  & \(2.3201\cdot10^{-1}\) & \(1.3180\cdot10^{-1}\)
       & \(7.9742\cdot10^{-2}\) & \(4.6657\cdot10^{-2}\)\\
\(1\)  & \(1.7467\cdot10^{-3}\) & \(3.0442\cdot10^{-3}\)
       & \(1.5017\cdot10^{-3}\) & \(6.0513\cdot10^{-4}\)\\
\(2\)  & \(1.7370\cdot10^{-3}\) & \(9.4510\cdot10^{-5}\)
       & \(5.7352\cdot10^{-6}\) & \(5.2671\cdot10^{-6}\)\\
\(3\)  & \(1.5606\cdot10^{-4}\) & \(1.7340\cdot10^{-5}\)
       & \(1.5654\cdot10^{-6}\) & \(8.1538\cdot10^{-8}\)\\
\(5\)  & \(5.1336\cdot10^{-6}\) & \(2.7114\cdot10^{-8}\)
       & \(2.2625\cdot10^{-9}\) & \(1.5423\cdot10^{-10}\)\\
\(8\)  & \(1.4208\cdot10^{-8}\) & \(1.0628\cdot10^{-11}\)
       & \(3.9030\cdot10^{-13}\) & \(2.9383\cdot10^{-15}\)\\
\(12\) & \(1.9269\cdot10^{-11}\) & \(2.7265\cdot10^{-15}\)
       & \(1.8339\cdot10^{-18}\) & \(2.1347\cdot10^{-21}\)\\
\bottomrule
\end{tabular}
\end{table}

\begin{table}[htbp]
\centering
\caption{Absolute error of the truncation of
\cref{plt:eq:lw-lower-expansion} for
\(\LambertW_{-1}(-\varepsilon)\), with
\(S=\log(1/\varepsilon)\), \(L=\log S\).}
\label{plt:tab:lw-lower-errors}
\begin{tabular}{@{}rlll@{}}
\toprule
\(N\) & \(\varepsilon=10^{-3}\ (S=6.908)\)
      & \(\varepsilon=10^{-6}\ (S=13.816)\)
      & \(\varepsilon=10^{-20}\ (S=46.052)\)\\
\midrule
\(0\)  & \(2.7761\cdot10^{-1}\) & \(1.8521\cdot10^{-1}\)
       & \(8.1518\cdot10^{-2}\)\\
\(1\)  & \(2.1725\cdot10^{-3}\) & \(4.8547\cdot10^{-3}\)
       & \(1.6446\cdot10^{-3}\)\\
\(2\)  & \(3.5365\cdot10^{-3}\) & \(5.5019\cdot10^{-4}\)
       & \(7.5505\cdot10^{-6}\)\\
\(3\)  & \(2.9763\cdot10^{-4}\) & \(8.7527\cdot10^{-5}\)
       & \(1.8886\cdot10^{-6}\)\\
\(5\)  & \(1.8100\cdot10^{-5}\) & \(1.0098\cdot10^{-6}\)
       & \(3.0375\cdot10^{-9}\)\\
\(8\)  & \(1.8025\cdot10^{-8}\) & \(2.4269\cdot10^{-9}\)
       & \(6.2816\cdot10^{-13}\)\\
\(12\) & \(3.8085\cdot10^{-10}\) & \(2.5958\cdot10^{-13}\)
       & \(3.3774\cdot10^{-18}\)\\
\bottomrule
\end{tabular}
\end{table}

% ed.: the four source tables are individually correct but use two
% incompatible column conventions (S1 and S4 index by x = log z, S3 and S5
% by z), which makes them look contradictory when read side by side.  All
% four were recomputed to 60 digits and merged into one table indexed by the
% intrinsic variable xi_0 = log z; the concordance is recorded below,
% together with the observation that the S3/S5 columns at z = 10 and z = 100
% lie outside the region where any convergence statement is proved here.
\begin{remark}[Concordance of the source tables]
\label{plt:rmk:lw-table-concordance}
For the reader of a particular source: the S1 columns headed \(L=10,50,100\)
and the S4 columns headed \(x=10,25,100\) are columns of
\cref{plt:tab:lw-errors} (S4's index \(N\) counts from the zero-based
family, \(\WrightOmega_N(x)=x+\sum_{n=0}^{N}\LamP{n}(\log x)x^{-n}\), which
is our \(w_N\)).  The S3 and S5 columns headed \(z=10,100,10^{6}\)
correspond to \(\xi_0=2.3026,\,4.6052,\,13.8155\).  For the first two the
seating condition \cref{plt:eq:lw-seating} \emph{fails}
(\(2(1+0.8340)=3.67\not<2.30\) and \(2(1+1.5272)=5.05\not<4.61\)), so
nothing proved here applies to them; for the third it holds
(\(2(1+2.6259)=7.25<13.82\)), so the series converges there, but the
remainder hypothesis \(\AbsoluteValue{\xi_0}\ge10(1+
\AbsoluteValue{\ell_0})=36.3\) still fails and the explicit bound
\cref{plt:eq:lw-remainder} is unavailable.  Those columns are therefore
empirical illustrations, not instances of a theorem.  Every numerical value
in all four source tables was recomputed to \(60\) digits and found
correct.
\end{remark}

\begin{remark}[Non-monotonicity is not a contradiction]
At \(\xi_0=10\) the error at \(N=2\) is essentially the same as at \(N=1\),
and at \(\varepsilon=10^{-3}\) the error at \(N=2\) is \emph{larger} than at
\(N=1\).  Nothing is wrong: \cref{plt:thm:lw-uniform-remainder} bounds the
tail by \(2\theta^{N+1}\) with \(\theta=5(1+\AbsoluteValue{\ell})/
\AbsoluteValue{\xi}\), and at \(\xi_0=10\) one has \(\theta\approx1.65>1\),
so the theorem says nothing there at all.  Even inside the proved region the
bound is a geometric envelope, not a term-by-term ordering: individual
\(\LamP{n}(\ell)\) change sign (\cref{plt:cor:lambert-coefficients}), so two
consecutive corrections can nearly cancel.  The residual \(R_N\) of
\cref{plt:def:lw-residual}, not the size of the last retained block, is the
quantity that should decide where to truncate; by
\cref{plt:thm:lw-newton-step}\,(ii) it controls the error with an explicit
constant.
\end{remark}

\begin{algorithm}[Branch-aware evaluation]
\label{plt:alg:lw-evaluate}
\hfill
\begin{lstlisting}[style=pseudo]
input : z != 0, branch index k, order N, tolerance tol
1  xi := Log(z) + 2*pi*i*k                 # principal Log; record the choice
2  if xi = 0: fail (the coordinates are undefined)
3  ell := Log(xi)                          # record the determination B_k
4  if |xi| < 10*(1 + |ell|):
       warn: outside the region of Theorem lw-uniform-remainder
5  w := xi - ell
6  for n := 1 to N:  w := w + P_n(ell) / xi^n     # tables, or the recurrence
7  R := w + ell + Log(w/xi) - xi           # the same branch data as line 3
8  while |R| > tol:
       w := w - R / (1 + 1/w)              # Newton in branch coordinates
       R := w + ell + Log(w/xi) - xi
9  return w, R
\end{lstlisting}
Line~7 evaluates \(\log_{*}w\) as \(\ell+\PrincipalLogarithm(w/\xi)\), never
as \(\PrincipalLogarithm w\): by \cref{plt:prop:lw-branch-label} the two
agree only when the distance-to-the-cut hypothesis
\cref{plt:eq:lw-cut-distance} holds, and an implementation that silently
uses \(\PrincipalLogarithm w\) will report a spurious residual of size
\(2\pi\) near the cut.  The multiplicative residual
\(w\EulerE^{w}-z\) must not be substituted for line~7: it overflows for
\(\AbsoluteValue{z}\) large, whereas \(R\) never forms \(\EulerE^{w}\).
\end{algorithm}

\begin{summarybox}
Lambert \(\LambertW\) is one reversion evaluated at many arguments.  The
substitution \(w\EulerE^{w}=z\rightsquigarrow w+\log_{*}w=\xi_k\)
(\cref{plt:prop:lw-log-reduction}) turns each branch into the reversion of
\(X+L\), so the coefficient polynomials \(\LamP{n}\) are the same for every
branch (\cref{plt:thm:lw-branch-independence}) while the coordinates
\(\xi_k,\ell_k\) carry all the branch dependence.  The formal expansion is a
theorem in \(\Gnear\) with no hypotheses; its analytic validity is a
separate theorem with explicit ones
(\cref{plt:thm:lw-basepoint,plt:thm:lw-uniform-remainder}), and the
frequently repeated claim of convergence for real \(z\ge\EulerE\) is not
among them.  The lower real branch is the single scalar \(a=-1\) of the
affine-logarithmic family and introduces no new polynomials.  At
\(z=-\EulerE^{-1}\) the derivative of \(w\EulerE^{w}\) vanishes, the local
monodromy has order two, and the correct scale is a square root, not a
logarithm.
\end{summarybox}

\begin{exercise}
Show that the seating condition \cref{plt:eq:lw-seating} holds for
\(K=\xi_0=\log z\) and \(M=\ell_0\) whenever \(z\ge\EulerE^{6}\), but fails
for every \(z\le\EulerE^{5}\); and find the exact threshold to four decimal
places.
\end{exercise}

\begin{exercise}
Take \(K=\xi_k-\ell_k\) in \cref{plt:thm:lw-basepoint} and show that
\(M=\log(1-\ell_k/\xi_k)\), so that
\(\AbsoluteValue{M}=\BigOAt{\AbsoluteValue{\ell_k/\xi_k}}
{\AbsoluteValue{\xi_k}\to\infty}\).  Deduce that the seating condition for
this base point is satisfied on a strictly larger region than for
\(K=\xi_k\), and compare with \cref{plt:thm:lw-transverse}.
\end{exercise}

\begin{exercise}
Using \(\log_{*}\LambertW=\xi-\LambertW\) and the division rules of
\cref{plt:sec:division}, expand
\(\LambertW_k(z)/(1+\LambertW_k(z))\) through \(\xi^{-3}\) and verify
against \cref{plt:eq:lw-derivative}.
\end{exercise}

\begin{exercise}
Prove that the two determinations of \cref{plt:thm:lw-branch-point} satisfy
\(\varphi(p)+\varphi(-p)=-\tfrac23p^{2}
-\tfrac{43}{270}p^{4}+\BigOAt{p^{6}}{p\to0}\), and explain why this sum is
an even function of \(p\), hence a single-valued holomorphic function of
\(1+\EulerE z\) near \(0\).
\end{exercise}

\chapter{Generalized Lambert equations and reusable inverse templates}
\label{plt:sec:generalized-lambert}

\Cref{plt:sec:lambert-branches} solved one equation, \(w+\log w=\xi\), and
obtained every branch of \(\LambertW\) from it.  This chapter asks how far
that leverage extends.  The answer has three layers.  First, an explicitly
delimited family of equations reduces to \(w+\log w=\xi\) by an exact
algebraic change of variable, so that no new coefficient family is needed at
all: this is the class \(y^{\alpha}\EulerE^{\beta y}=x\) and its relatives
(\cref{plt:sec:lw-generalized-class}).  Second, a much larger class does not
reduce exactly but has the same \emph{shape} of dominant balance --- a
monomial times a power of a logarithm --- and for that class there is a
mechanical template producing the polynomial--logarithmic inverse, with the
Lambert polynomials appearing as the leading part and computable corrections
beyond (\cref{plt:sec:lw-template}).  Third, and equally important, there
are equations that look similar and are not in the class; the last worked
example is one of them, and the point of including it is that the reader
should be able to recognise the boundary before spending a page computing on
the wrong side of it.

\section{The affine-logarithmic core, analytically}
\label{plt:sec:lw-affine-analytic}

\Cref{plt:thm:lw-affine-log} is the formal statement.  Its analytic
counterpart is proved by exactly the argument of
\cref{plt:thm:lw-basepoint}, with the pair \((\sigma,\tau)\) chosen to carry
the parameter \(a\).

\begin{theorem}[Analytic affine-logarithmic reversion]
\label{plt:thm:lw-affine-analytic}
Let \(a\in\ComplexNumbers\setminus\SetOf{0}\),
\(X\in\ComplexNumbers\setminus\SetOf{0}\), and let \(\ell\) satisfy
\(\EulerE^{\ell}=X\).  Assume
\begin{equation}\label{plt:eq:lw-affine-seating}
 2\AbsoluteValue{a}\bigl(1+\AbsoluteValue{\ell}\bigr)
 <\AbsoluteValue{X}.
\end{equation}
Then the series below converges absolutely; the number
\begin{equation}\label{plt:eq:lw-affine-value}
 y\DefinitionEquals X+a\sum_{n\ge0}
 \left(\frac aX\right)^{n}\LamP{n}(\ell)
 =X-a\ell+\sum_{n\ge1}a^{n+1}\frac{\LamP{n}(\ell)}{X^{n}}
\end{equation}
satisfies \(\AbsoluteValue{y/X-1}\le\tfrac12\), hence \(y\ne0\); and with
\(\log_{*}y\DefinitionEquals\ell+\PrincipalLogarithm(y/X)\), which is a
logarithm of \(y\),
\begin{equation}\label{plt:eq:lw-affine-equation}
 y+a\log_{*}y=X .
\end{equation}
If moreover \(10\AbsoluteValue{a}(1+\AbsoluteValue{\ell})
\le\AbsoluteValue{X}\), then with
\(\theta_a\DefinitionEquals5\AbsoluteValue{a}(1+\AbsoluteValue{\ell})/
\AbsoluteValue{X}\le\tfrac12\) and every \(N\in\NaturalNumbers\),
\begin{equation}\label{plt:eq:lw-affine-remainder}
 \left\lvert y-X-a\sum_{n=0}^{N}
 \left(\frac aX\right)^{n}\LamP{n}(\ell)\right\rvert
 \le2\AbsoluteValue{a}\,\theta_a^{\,N+1}.
\end{equation}
All data being real, \(y\) is real.
\end{theorem}

\begin{proof}
Put \(\sigma=a/X\) and \(\tau=a\ell/X\), so that \(\tau=\sigma\ell\) and
\(\AbsoluteValue{\sigma}+\AbsoluteValue{\tau}
=\AbsoluteValue{a}(1+\AbsoluteValue{\ell})/\AbsoluteValue{X}<\tfrac12\) by
\cref{plt:eq:lw-affine-seating}; thus \((\sigma,\tau)\in D\).  Let
\(\mathcal Y=\mathcal Y(\sigma,\tau)\) be as in
\cref{plt:lem:lw-two-variable}.  By \cref{plt:eq:lw-c-to-P} with \(u=\ell\),
\[
 \mathcal Y=\sum_{n\ge1}\LamP{n}(\ell)\sigma^{n}
 =\sum_{n\ge1}a^{n}\frac{\LamP{n}(\ell)}{X^{n}},
\]
absolutely convergent, and since \(\LamP{0}(\ell)=-\ell\) the number
\(y\) of \cref{plt:eq:lw-affine-value} equals \(X-a\ell+a\mathcal Y\).
Then
\[
 \frac yX=1-\frac{a\ell}{X}+\frac{a\mathcal Y}{X}
 =1-\tau+\sigma\mathcal Y ,
\]
which lies in the disc of radius \(\tfrac12\) about \(1\), so
\(\AbsoluteValue{y/X-1}\le\tfrac12\), \(y\ne0\),
\(\PrincipalLogarithm(y/X)=-\mathcal Y\) by
\cref{plt:eq:lw-two-var-fixed}, and
\(\EulerE^{\log_{*}y}=X\cdot(y/X)=y\).  Therefore
\[
 y+a\log_{*}y=(X-a\ell+a\mathcal Y)+a(\ell-\mathcal Y)=X .
\]
The remainder bound is \cref{plt:eq:lw-block-bound} with
\(5(\AbsoluteValue{\sigma}+\AbsoluteValue{\tau})=\theta_a\), multiplied by
the outer factor \(\AbsoluteValue{a}\).  Reality follows from uniqueness of
the fixed point in \cref{plt:lem:lw-two-variable}: for real
\((\sigma,\tau)\) the contraction \(\Phi\) preserves the real interval
\([-1,1]\), so its fixed point is real.
\end{proof}

\section{The class \texorpdfstring{\(y^{\alpha}\EulerE^{\beta y}=x\)}
{y to the alpha times exp(beta y) = x}}
\label{plt:sec:lw-generalized-class}

\begin{proposition}[Exact reduction, with the branch bookkeeping made
explicit]
\label{plt:prop:lw-power-exp}
Let \(\alpha,\beta\in\ComplexNumbers\setminus\SetOf{0}\) and
\(x\in\ComplexNumbers\setminus\SetOf{0}\); fix a logarithm \(\xi\) of \(x\)
and a logarithm \(\gamma\) of \(\alpha/\beta\).  Then the map
\begin{equation}\label{plt:eq:lw-power-exp-map}
 (y,\lambda)\longmapsto(u,\mu)
 \DefinitionEquals\left(\frac{\beta}{\alpha}y,\ \lambda-\gamma\right)
\end{equation}
is a bijection between
\[
 \SetOf{(y,\lambda):\EulerE^{\lambda}=y,\ \beta y+\alpha\lambda=\xi}
 \quad\text{and}\quad
 \SetOf{(u,\mu):\EulerE^{\mu}=u,\ u+\mu=\zeta},
 \qquad
 \zeta\DefinitionEquals\frac{\xi}{\alpha}-\gamma .
\]
Consequently, writing \(x^{1/\alpha}\DefinitionEquals
\EulerE^{\xi/\alpha}\), every solution of \(y^{\alpha}\EulerE^{\beta y}=x\)
compatible with the chosen logarithms has the form
\begin{equation}\label{plt:eq:lw-power-exp-solution}
 y=\frac{\alpha}{\beta}\,u,
 \qquad
 u\EulerE^{u}=\EulerE^{\zeta}=\frac{\beta}{\alpha}\,x^{1/\alpha},
\end{equation}
that is, \(y=(\alpha/\beta)\LambertW_k\bigl((\beta/\alpha)
x^{1/\alpha}\bigr)\) for the branch label \(k\) determined by
\cref{plt:prop:lw-log-reduction}.
\end{proposition}

\begin{proof}
Let \(\EulerE^{\lambda}=y\) and \(\beta y+\alpha\lambda=\xi\); set
\(u=(\beta/\alpha)y\) and \(\mu=\lambda-\gamma\).  Then
\(\EulerE^{\mu}=\EulerE^{\lambda}\EulerE^{-\gamma}
=y\cdot(\beta/\alpha)=u\), and
\(u+\mu=(\beta/\alpha)y+\lambda-\gamma
=\alpha^{-1}(\beta y+\alpha\lambda)-\gamma=\xi/\alpha-\gamma=\zeta\).
The inverse map is \((u,\mu)\mapsto((\alpha/\beta)u,\mu+\gamma)\), so the
correspondence is bijective.  For the last statement,
\(\EulerE^{\mu}=u\) and \(u+\mu=\zeta\) give
\(u\EulerE^{u}=\EulerE^{u+\mu}=\EulerE^{\zeta}\) by
\cref{plt:prop:lw-log-reduction}, and
\(\EulerE^{\zeta}=\EulerE^{\xi/\alpha}\EulerE^{-\gamma}
=(\beta/\alpha)x^{1/\alpha}\).  Finally
\(y^{\alpha}\EulerE^{\beta y}\) means \(\EulerE^{\alpha\lambda+\beta y}\)
once the logarithm \(\lambda\) of \(y\) has been chosen, so the equation
\(y^{\alpha}\EulerE^{\beta y}=x\) is exactly \(\beta y+\alpha\lambda=\xi\)
modulo \(2\pi\ImaginaryUnit\IntegerNumbers\), i.e.\ exactly the left-hand
set after adjusting \(\xi\) by a multiple of \(2\pi\ImaginaryUnit\).
\end{proof}

% ed.: S3, S4, S5 and S6 all give the closed form
% y = (alpha/beta) W_k((beta/alpha) x^{1/alpha}) with the disclaimer
% "branches chosen consistently", which hides a genuine ambiguity: for
% non-integer alpha the symbol x^{1/alpha} is not determined by x.
% Repaired by Proposition lw-power-exp, which makes the three choices
% explicit and states the correspondence as a bijection.
\begin{warningbox}[title={``With branches chosen consistently'' is three
choices, not a disclaimer}]
Reports S3, S4, S5 and S6 all display
\(y=(\alpha/\beta)\LambertW_k((\beta/\alpha)x^{1/\alpha})\) followed by a
phrase such as ``with powers and branches chosen consistently''.
\Cref{plt:prop:lw-power-exp} says what the phrase must mean: one chooses
\emph{(a)} a logarithm \(\xi\) of \(x\), which fixes the meaning of
\(x^{1/\alpha}\) as \(\EulerE^{\xi/\alpha}\) and is \emph{not} determined by
\(x\) alone when \(\alpha\notin\IntegerNumbers\); \emph{(b)} a logarithm
\(\gamma\) of \(\alpha/\beta\); \emph{(c)} the branch label \(k\), which by
\cref{plt:prop:lw-log-reduction} is then determined by the solution, not
free.  Different choices at (a) or (b) give genuinely different solutions
\(y\) of the same equation, not different names for one.
\end{warningbox}

\begin{theorem}[Generalized Lambert expansion]
\label{plt:thm:lw-generalized-expansion}
In the situation of \cref{plt:prop:lw-power-exp}, put
\begin{equation}\label{plt:eq:lw-gen-coords}
 a\DefinitionEquals\frac{\alpha}{\beta},
 \qquad
 X\DefinitionEquals\frac{\xi}{\beta},
\end{equation}
and let \(\ell\) satisfy \(\EulerE^{\ell}=X\).  If
\(2\AbsoluteValue{a}(1+\AbsoluteValue{\ell})<\AbsoluteValue{X}\), then
\begin{equation}\label{plt:eq:lw-gen-expansion}
 y=X+a\sum_{n\ge0}\left(\frac aX\right)^{n}\LamP{n}(\ell)
 =\frac{\xi}{\beta}
 -\frac{\alpha}{\beta}\ell
 +\sum_{n\ge1}
 \frac{\alpha^{n+1}}{\beta}\,\frac{\LamP{n}(\ell)}{\xi^{n}}
\end{equation}
converges absolutely and solves \(y^{\alpha}\EulerE^{\beta y}=x\) with
\(\lambda=\ell+\PrincipalLogarithm(y/X)\) as the logarithm of \(y\); the
remainder after \(N\) terms is bounded by
\cref{plt:eq:lw-affine-remainder}.  In particular, for \(\beta=1\), real
\(x\to+\infty\), \(L_1=\log x\) and \(L_2=\log\log x\),
\begin{equation}\label{plt:eq:lw-gen-beta-one}
 y=L_1-\alpha L_2
 +\sum_{n\ge1}\alpha^{n+1}\frac{\LamP{n}(L_2)}{L_1^{n}} .
\end{equation}
\end{theorem}

\begin{proof}
By \cref{plt:prop:lw-power-exp} the equation is
\(\beta y+\alpha\lambda=\xi\) with \(\EulerE^{\lambda}=y\); dividing by
\(\beta\) gives \(y+a\lambda=X\).  \Cref{plt:thm:lw-affine-analytic} with
this \(a\) and \(X\) produces \(y\) as in \cref{plt:eq:lw-affine-value},
which is the first equality in \cref{plt:eq:lw-gen-expansion}; the second is
\(a(a/X)^{n}=\alpha^{n+1}\beta^{-n-1}\cdot\beta^{n}\xi^{-n}
=\alpha^{n+1}\beta^{-1}\xi^{-n}\).  For \(\beta=1\) we have \(X=\xi=L_1\),
\(\ell=L_2\), \(a=\alpha\), and \(a\cdot a^{n}\LamP{n}(L_2)L_1^{-n}
=\alpha^{n+1}\LamP{n}(L_2)L_1^{-n}\).
\end{proof}

\begin{remark}[Comparison with the sources, and with the literature]
\label{plt:rmk:lw-gen-sources}
\Cref{plt:eq:lw-gen-beta-one} is S5's generalized Lambert expansion and
S2's \(\alpha\)-expansion; both are correct and both are proved there only
formally.  What is added here is the analytic half: a printed seating
condition and an explicit remainder.  The inversion of
\(y^{\alpha}\EulerE^{y}\) is classical: Comtet treated it together with
\(y\log y\) \cite{comtet1970}, and Jeffrey, Corless, Hare and Knuth
expressed the coefficients through associated Stirling numbers and studied
the convergence of the resulting series \cite{jeffrey-et-al-1995}.  S5
reports their sufficient domains as \(x>(\alpha\EulerE)^{\alpha}\) for
\(\alpha\ge1\) and \(x>\EulerE\) for \(\alpha<1\).  Those domains are much
larger than the region \(2\AbsoluteValue{\alpha}(1+\AbsoluteValue{L_2})<L_1\)
proved here; they are recorded as attribution and are not theorems of this
volume.  See \cref{plt:sec:lw-associated} for the structural reason why a
different regrouping converges on a larger domain, and for what this volume
can prove about it.
\end{remark}

\begin{proposition}[The \texorpdfstring{\(r\)}{r}-Lambert equation: the same
expansion to all orders]
\label{plt:prop:lw-r-lambert}
Fix \(r\in\RealNumbers\).  There is \(z_r>0\) such that for real \(z>z_r\)
the equation
\begin{equation}\label{plt:eq:lw-r-lambert}
 w\EulerE^{w}+rw=z
\end{equation}
has a unique large positive solution \(w=\mathcal W_r(z)\), it tends
to \(+\infty\), and
\begin{equation}\label{plt:eq:lw-r-lambert-difference}
 \bigl\lvert\mathcal W_r(z)-\LambertW_0(z)\bigr\rvert
 \le\frac{3\AbsoluteValue{r}\log z}{z}
 \qquad(z\to+\infty).
\end{equation}
Consequently \(\mathcal W_r\) and \(\LambertW_0\) have the \emph{same}
polynomial--logarithmic asymptotic expansion to all orders in the scale
\(\SetOf{\ell_0^{\,j}\xi_0^{-n}}\): the parameter \(r\) is invisible beyond
all orders.
\end{proposition}

\begin{proof}
Let \(f(w)=w\EulerE^{w}+rw\).  Then \(f'(w)=(1+w)\EulerE^{w}+r\), and
\((1+w)\EulerE^{w}\to\infty\), so there is \(w_r\ge1\) with \(f'>0\) on
\([w_r,\infty)\); since \(f\to+\infty\), \(f\) is a strictly increasing
bijection from \([w_r,\infty)\) onto \([f(w_r),\infty)\).  Put
\(z_r=\max\SetOf{f(w_r),\EulerE^{2}}\); for \(z>z_r\) let \(w>w_r\) be the
unique preimage.  As \(z\to\infty\), \(w\to\infty\).
% ed.: the proof establishes the residual bound only "for z large", while
% the statement attaches it to the same z_r that governs existence.  The
% quantifier is repaired here by enlarging z_r once, which is legitimate
% because the statement is existential in z_r.
Each of the finitely many smallness conditions used below holds for all
sufficiently large \(z\); we enlarge \(z_r\) once, at the end of the proof,
so that all of them hold for \(z>z_r\).

From \cref{plt:eq:lw-r-lambert}, \(w\EulerE^{w}=z-rw\), so
\begin{equation}\label{plt:eq:lw-r-residual}
 w+\log w=\log(z-rw)=\xi_0+R,
 \qquad
 R\DefinitionEquals\log\left(1-\frac{rw}{z}\right),
 \qquad\xi_0=\log z,
\end{equation}
valid once \(z-rw>0\).  To see that this holds for \(z\) large, note first
that \(w\EulerE^{w}=z-rw\le z+\AbsoluteValue{r}w\); since \(w\to\infty\) we
may assume \(\EulerE^{w}\ge2\AbsoluteValue{r}\), and then
\(2\AbsoluteValue{r}w\le w\EulerE^{w}\le z+\AbsoluteValue{r}w\) forces
\(\AbsoluteValue{r}w\le z\) and hence \(w\EulerE^{w}\le2z\).  Since
\(\LambertW_0(y)\le\log y\) for \(y\ge\EulerE\), this gives
\(w\le\LambertW_0(2z)\le\log(2z)\le\tfrac32\log z\) for \(z\ge4\).  For
\(z\) large,
\(\AbsoluteValue{rw/z}\le\tfrac12\) and therefore
\(\AbsoluteValue{R}\le2\AbsoluteValue{r}w/z\le
3\AbsoluteValue{r}\log z/z\), using
\(\AbsoluteValue{\log(1-\nu)}\le2\AbsoluteValue{\nu}\) for
\(\AbsoluteValue{\nu}\le\tfrac12\) and \(w\le\log(2z)\le\tfrac32\log z\).

Let \(F(v)=v+\log v\) on \((0,\infty)\); then \(F'(v)=1+1/v>1\), so
\(\AbsoluteValue{F(v_1)-F(v_2)}\ge\AbsoluteValue{v_1-v_2}\).  The principal
branch satisfies \(F(\LambertW_0(z))=\xi_0\) (take the real logarithm of
\(\LambertW_0\EulerE^{\LambertW_0}=z\), legitimate since
\(\LambertW_0(z)>0\)), and \(F(w)=\xi_0+R\); hence
\(\AbsoluteValue{w-\LambertW_0(z)}\le\AbsoluteValue{R}\), which is
\cref{plt:eq:lw-r-lambert-difference}.

Finally \(\log z/z=\xi_0\EulerE^{-\xi_0}\), which is
\(\LittleOAt{\xi_0^{-n}}{\xi_0\to\infty}\) for every \(n\); so subtracting
\(\mathcal W_r\) from \(\LambertW_0\) changes no coefficient of the
expansion.
\end{proof}

\begin{remark}[Why this proposition is here]
None of the six sources treats the \(r\)-Lambert family, so nothing is
consolidated in \cref{plt:prop:lw-r-lambert}; it is included because the
family is the standard first generalisation of \(\LambertW\) in the
literature and because the answer is instructive.  It exhibits, in the
simplest possible setting, the limitation that every result of this volume
shares: a polynomial--logarithmic expansion sees a function only modulo
quantities that are smaller than every power of the large variable, and a
one-parameter family of genuinely different functions can share one
expansion.  The same phenomenon is revisited in
\cref{plt:rmk:lw-beyond-all-orders}.
\end{remark}

\section{The template}
\label{plt:sec:lw-template}

\begin{definition}[Monomial--logarithmic inverse datum]
\label{plt:def:lw-monomial-log-datum}
Let \(\K\) be a field of characteristic zero.  An inverse problem
\(\Phi(y)=x\) is \emph{of monomial--logarithmic type} with data
\((\rho,\mu,\beta,B)\) if, after a logarithm \(\xi\) of \(x\) and a
logarithm \(Z\) of \(y\) have been chosen, it is equivalent to
\begin{equation}\label{plt:eq:lw-template-equation}
 \rho Z+\mu\log Z+\beta+B(Z)=\xi ,
\end{equation}
where \(\rho\in\K^{\times}\), \(\mu,\beta\in\K\), and, writing
\(t=Z^{-1}\) and \(L=\log Z\) for the generators of the logarithmic chart,
\begin{equation}\label{plt:eq:lw-template-B}
 B=\sum_{m\ge1}b_m(L)\,t^{m}\in t\K[L][[t]],
 \qquad b_m\in\K[L].
\end{equation}
Exponentiating, this says
\(\Phi(y)=\EulerE^{\beta}y^{\rho}(\log y)^{\mu}
\EulerE^{B(\log y)}\), all powers being defined through the chosen
logarithms; the leading balance is the monomial
\(\EulerE^{\beta}y^{\rho}(\log y)^{\mu}\).
\end{definition}

The three hypotheses are worth naming, because each fails in a recognisable
way.
\begin{itemize}
\item[(H1)] \(\rho\ne0\): there is a genuine monomial in the balance.  If
      \(\rho=0\) the inverse changes exponential level and leaves every
      polynomial--logarithmic class
      (\cref{plt:prop:grp-exponential}); see
      \cref{plt:rmk:lw-rho-zero}.
\item[(H2)] \(\mu\) is a \emph{constant}, so that the logarithmic part of
      the balance is \(\mu\log Z\) and not something with a variable
      exponent.
\item[(H3)] \(B\) has strictly positive \(t\)-order and coefficients
      \emph{polynomial} in \(L\).  If \(B\) has \(t\)-order \(0\) or its
      coefficients are not polynomial, the reversion leaves
      \(\Gnear\); see \cref{plt:ex:lw-failure}.
\end{itemize}

\begin{theorem}[Reusable inverse template]
\label{plt:thm:lw-template}
Let \(\Phi(y)=x\) be of monomial--logarithmic type with data
\((\rho,\mu,\beta,B)\) as in \cref{plt:def:lw-monomial-log-datum}.  Put
\begin{equation}\label{plt:eq:lw-template-coords}
 a\DefinitionEquals\frac{\mu}{\rho},
 \qquad
 X\DefinitionEquals\frac{\xi-\beta}{\rho},
 \qquad
 t=X^{-1},
 \qquad
 \ell=\log X .
\end{equation}
Then:
\begin{enumerate}
\item \(F\DefinitionEquals X+aL+\rho^{-1}B\) lies in \(\Gnear\), and
      \cref{plt:eq:lw-template-equation} is equivalent to
      \(F(Z)=X\); hence there is a unique
      \(Z\in X+\K[\ell][[t]]\) solving it, namely \(Z=\rev{F}\) evaluated at
      \(X\);
\item writing \(\rev{F}=X+\sum_{n\ge0}q_n(L)t^{n}\), one has
      \(q_0=-aL\), and for every \(N\ge0\) the polynomial \(q_N\) is a
      polynomial function of \(a\) and of the coefficients of
      \(b_1,\dots,b_N\) only, given explicitly by the coefficient formula of
      \cref{plt:thm:lag-coefficients};
\item \(q_N=a^{N+1}\LamP{N}+E_N\), where \(E_N\) depends only on
      \(a,b_1,\dots,b_N\) and vanishes identically when
      \(b_1=\dots=b_N=0\); in particular
      \(q_N=a^{N+1}\LamP{N}\) for all \(N\) when \(B=0\);
\item re-exponentiating,
      \begin{equation}\label{plt:eq:lw-template-answer}
       y=\EulerE^{Z}
       =\left(\frac{x}{\EulerE^{\beta}}\right)^{1/\rho}
        X^{-\mu/\rho}
        \exp\!\left(\sum_{n\ge1}q_n(\ell)\,X^{-n}\right),
      \end{equation}
      with \(x^{1/\rho}\DefinitionEquals\EulerE^{\xi/\rho}\); the last
      factor lies in \(1+t\K[\ell][[t]]\).  In particular the leading
      behaviour is
      \begin{equation}\label{plt:eq:lw-template-leading}
       y\sim\EulerE^{-\beta/\rho}\,x^{1/\rho}
       \left(\frac{\xi}{\rho}\right)^{-\mu/\rho} .
      \end{equation}
\end{enumerate}
When \(1/\rho\) or \(\mu/\rho\) is not an integer, \cref{plt:eq:lw-template-answer}
lives in the ramified Puiseux--logarithmic tower of
\cref{plt:prop:ring-puiseux-tower}, not in \(\PLlau\).
\end{theorem}

\begin{proof}
(i)  Dividing \cref{plt:eq:lw-template-equation} by \(\rho\) and moving
\(\beta/\rho\) to the right gives
\(Z+a\log Z+\rho^{-1}B(Z)=X\).  The perturbation
\(A\DefinitionEquals aL+\rho^{-1}B\) lies in \(\K[L][[t]]\) by
\cref{plt:eq:lw-template-B}, so \(F=X+A\in\Gnear\)
(\cref{plt:def:grp-gnear}), and \cref{plt:thm:rev-existence} gives a unique
inverse in \(\Gnear\).

(ii)  Write \(A^{r}=\sum_{m\ge0}A_{r,m}(L)t^{m}\).  By
\cref{plt:thm:lag-coefficients},
\[
 q_N(L)=\sum_{r=1}^{N+1}\frac{(-1)^{r}}{r!}
 \left[\prod_{s=N-r+1}^{N-1}(\partial_L-s)\right]A_{r,N-r+1}(L),
\]
with the empty product acting as the identity.  Now \(A_{r,m}\) is a sum of
products of \(r\) blocks \(A_{j_1},\dots,A_{j_r}\) with
\(j_1+\dots+j_r=m\), so it depends only on \(A_0,\dots,A_m\); since
\(m=N-r+1\le N\), the whole formula involves only \(A_0=aL\) and
\(A_1=\rho^{-1}b_1,\dots,A_N=\rho^{-1}b_N\), polynomially.  The block
\(q_0\) is the \(r=1\), \(m=0\) term, namely \(-A_{1,0}=-A_0=-aL\).

(iii)  Setting \(b_1=\dots=b_N=0\) in the formula of (ii) leaves the
coefficients of \(\rev{(X+aL)}\), which are \(a^{N+1}\LamP{N}\) by
\cref{plt:thm:lw-affine-log}.  Define \(E_N\) as the difference; by (ii) it
is a polynomial function of \(a,b_1,\dots,b_N\) vanishing when the \(b\)'s
do.

(iv)  From \(Z=X+q_0(\ell)+\sum_{n\ge1}q_n(\ell)X^{-n}\) and
\(q_0(\ell)=-a\ell\),
\[
 \EulerE^{Z}
 =\EulerE^{X}\cdot\EulerE^{-a\ell}\cdot
 \exp\!\Bigl(\sum_{n\ge1}q_n(\ell)X^{-n}\Bigr).
\]
Here \(\EulerE^{-a\ell}=(\EulerE^{\ell})^{-a}=X^{-a}=X^{-\mu/\rho}\) and
\(\EulerE^{X}=\EulerE^{(\xi-\beta)/\rho}
=\EulerE^{\xi/\rho}\EulerE^{-\beta/\rho}
=x^{1/\rho}\EulerE^{-\beta/\rho}\).  The exponential of an element of
\(t\K[\ell][[t]]\) lies in \(1+t\K[\ell][[t]]\) by
\cref{plt:lem:cmp-exp-log}, which also gives
\cref{plt:eq:lw-template-leading} after replacing
\(X=(\xi-\beta)/\rho\) by \(\xi/\rho\), a substitution absorbed by the
\(1+t\K[\ell][[t]]\) factor.  The final sentence is
\cref{plt:prop:cmp-fractional-rho} together with
\cref{plt:prop:ring-puiseux-tower}.
\end{proof}

\begin{corollary}[Closed form when \texorpdfstring{\(B=0\)}{B = 0}]
\label{plt:cor:lw-template-closed}
If \(B=0\), then
\begin{equation}\label{plt:eq:lw-template-closed}
 Z=\frac{\xi-\beta}{\rho}
 +\frac{\mu}{\rho}\sum_{n\ge0}
 \left(\frac{\mu}{\xi-\beta}\right)^{n}\LamP{n}(\ell),
 \qquad
 \ell=\log\frac{\xi-\beta}{\rho},
\end{equation}
and the series converges absolutely, with the remainder bound
\cref{plt:eq:lw-affine-remainder}, as soon as
\(2\AbsoluteValue{\mu}(1+\AbsoluteValue{\ell})
<\AbsoluteValue{\xi-\beta}\).
\end{corollary}

\begin{proof}
\Cref{plt:thm:lw-affine-analytic} with \(a=\mu/\rho\) and
\(X=(\xi-\beta)/\rho\); note \(a/X=\mu/(\xi-\beta)\) and
\(a\cdot(a/X)^{n}\) has the prefactor \(\mu/\rho\) shown.  The seating
condition \(2\AbsoluteValue{a}(1+\AbsoluteValue{\ell})<\AbsoluteValue{X}\)
is \(2\AbsoluteValue{\mu}(1+\AbsoluteValue{\ell})
<\AbsoluteValue{\xi-\beta}\) after multiplying by
\(\AbsoluteValue{\rho}\).
\end{proof}

% ed.: the drafted sentence said rho appears "only in the leading term and
% inside ell", which the displayed formula contradicts: rho also sits in the
% overall prefactor mu/rho.  What is true is that rho is absent from the
% ratio a/X = mu/(xi - beta) that drives the series.
Note the shape of \cref{plt:eq:lw-template-closed}: the exponent \(\rho\) of
the monomial appears only in the leading term, in the overall prefactor
\(\mu/\rho\) and inside \(\ell\) --- never in the ratio
\(a/X=\mu/(\xi-\beta)\) that drives the series --- while
the exponent \(\mu\) of the logarithm is the expansion parameter.  This is
why the generalized Lambert family is a one-parameter family and not a
two-parameter one.

\begin{algorithm}[The template, mechanically]
\label{plt:alg:lw-template}
\hfill
\begin{lstlisting}[style=pseudo]
input : the direct map Phi, the target x, an order N
1  choose logarithms: xi with e^xi = x; Z will be a logarithm of y
2  rewrite Phi(y) = x as  rho*Z + mu*log Z + beta + B(Z) = xi
      # if this is impossible with constant rho != 0 and constant mu,
      # or with B of positive t-order and polynomial coefficients,
      # the template does not apply: see Example lw-failure
3  a := mu/rho ;  X := (xi - beta)/rho ;  ell := log X ;  t := 1/X
4  A := a*L + B/rho                       # an element of K[L][[t]]
5  Z := reverse(X + A) to order N         # triangular or Newton reversion
      # equivalently  q_n = a^{n+1} P_n(ell) + E_n(b_1..b_n)
6  y := exp(Z)                            # = (x e^-beta)^(1/rho) * X^(-a) * unit
7  residual := (rho*Z + mu*log Z + beta + B(Z) - xi)   # must be O(t^{N+1})
\end{lstlisting}
Step~5 is where the whole of \cref{plt:sec:reversion-algorithms} is reused
unchanged; step~7 is the finite certificate of
\cref{plt:cor:rev-finite-certificate}.  Step~6 is the only place where the
answer can leave \(\PLlau\), and it does so exactly when \(1/\rho\) or
\(\mu/\rho\) is not an integer.
\end{algorithm}

\section{Three worked examples}
\label{plt:sec:lw-examples}

\begin{example}[A fractional power forcing a Puiseux--logarithmic extension]
\label{plt:ex:lw-fractional}
Solve \(Y=X^{2}\log X\) for \(X\) as \(Y\to+\infty\).

Here \(\rho=2\), \(\mu=1\), \(\beta=0\), \(B=0\), and \(a=\mu/\rho=1/2\).
With \(\xi=\log Y\), the template coordinates are
\(X_{\ast}=\xi/2\) and \(\ell=\log(\xi/2)\), and
\(a/X_{\ast}=\mu/\xi=1/\xi\).  \Cref{plt:cor:lw-template-closed} gives
\begin{equation}\label{plt:eq:lw-ex1-Z}
 \log X=\frac{\xi}{2}
 +\frac12\sum_{n\ge0}\frac{\LamP{n}(\ell)}{\xi^{n}}
 =\frac{\xi}{2}-\frac{\ell}{2}
 +\frac{\ell}{2\xi}
 +\frac{\ell^{2}/2-\ell}{2\xi^{2}}+\cdots,
\end{equation}
convergent as soon as \(2(1+\AbsoluteValue{\ell})<\xi\).  Exponentiating,
\begin{equation}\label{plt:eq:lw-ex1-X}
 X=\sqrt{\frac{2Y}{\log Y}}
 \left(1+\frac{\ell}{2\xi}
 +\frac{3\ell^{2}/8-\ell/2}{\xi^{2}}
 +\BigOAt{\ell^{3}/\xi^{3}}{Y\to+\infty}\right),
 \quad
 \xi=\log Y,\ \ell=\log\frac{\log Y}{2}.
\end{equation}
Two features are structural rather than accidental.  First, the answer
contains \(Y^{1/2}\): by \cref{plt:thm:lw-template} the exponent is
\(1/\rho\), so any \(\rho\) with \(1/\rho\notin\IntegerNumbers\) takes the
answer out of \(\PLlau\) into the ramified tower of
\cref{plt:prop:ring-puiseux-tower}.  Second, the
answer contains \(\log\log Y\) through \(\ell\), although the problem
contained only one logarithm; this is \cref{plt:prop:lag-loglog-forced},
and it is why the coefficient ring has to be \(\K[\ell]\) and not \(\K\).
The relative errors of the three displayed truncations at \(Y=10^{20}\) are
\(3.4\cdot10^{-2}\), \(9.8\cdot10^{-4}\) and \(1.3\cdot10^{-5}\).

The general case \(Y=cX^{\rho}(\log X)^{\mu}\) is
\cref{plt:eq:lw-template-leading}:
\(X\sim c^{-1/\rho}\rho^{\mu/\rho}Y^{1/\rho}(\log Y)^{-\mu/\rho}\), which is
the leading form quoted in S3 and S6.
\end{example}

\begin{example}[An iterated logarithm, and where the coefficient field
breaks]
\label{plt:ex:lw-iterated-log}
Solve \(Y=F(X)=X+\alpha\log X+\gamma\log\log X\) for \(X\) as
\(Y\to+\infty\), with \(\alpha,\gamma\in\K\).

This is \emph{not} covered by \cref{plt:thm:lw-template} as stated.  In the
template's variables the equation is
\(Z+\alpha\log Z+\gamma\log\log Z=Y\), so besides the permitted block
\(\mu\log Z\) with \(\mu=\alpha\) there is a \emph{second} block of
\(t\)-order \(0\), namely \(\gamma\log\log Z\), whose coefficient
\(M=\log L\) is not in \(\K[L]\): hypothesis (H3) fails.  Equivalently, in
the reversion chart the perturbation is \(A=\alpha L+\gamma M\) with
\(L=\log Y\), \(M=\log L\), and \(A\notin\K[L][[t]]\).  The example is
covered by the depth-two version: by \cref{plt:thm:grp-multilog-group} the
near-identity group at logarithmic depth \(2\) is a group over the
coefficient ring generated by \(L\) and \(M\), and the proof of
\cref{plt:thm:lw-template} goes through verbatim once
\(\K[L]\) is replaced by that ring.  Writing
\begin{equation}\label{plt:eq:lw-ex2-data}
 A=\alpha L+\gamma M,
 \qquad
 \varpi\DefinitionEquals\alpha+\frac{\gamma}{L},
 \qquad
 \frac{\Differential A}{\Differential Y}=\frac{\varpi}{Y},
 \qquad
 \frac{\Differential^{2}A}{\Differential Y^{2}}
 =-\frac1{Y^{2}}\left(\alpha+\frac{\gamma}{L}
 +\frac{\gamma}{L^{2}}\right),
\end{equation}
the first three Lagrange--B\"urmann terms of
\cref{plt:thm:lag-burmann} give
\begin{align}
 \CompositionalInverseOf{F}(Y)
 ={}&Y-A+\frac{A\varpi}{Y}
 +\frac1{Y^{2}}\left[-A\varpi^{2}
 +\frac{A^{2}}{2}\left(\alpha+\frac{\gamma}{L}
 +\frac{\gamma}{L^{2}}\right)\right]
 \notag\\
 &+\FormalBigOAt{t^{3}}{t},
 \qquad t=Y^{-1}.
 \label{plt:eq:lw-ex2-inverse}
\end{align}
% ed.: the remainder in this display was printed as an analytic estimate
% BigO(L^4/Y^3) although the three blocks are computed by formal
% Lagrange--Burmann; per the volume's own rule the three O-roles must not be
% conflated.  Downgraded to the formal tail, with the analytic upgrade
% stated as a separate sentence carrying its hypothesis.
The remainder here is the \emph{formal} tail of \cref{plt:thm:lag-burmann},
not an analytic estimate: the three displayed blocks are produced by a
computation in the depth-two reversion chart and say nothing by themselves
about the function \(\CompositionalInverseOf{F}\).  They do become an
asymptotic statement, with remainder
\(\BigOAt{L^{4}/Y^{3}}{Y\to+\infty}\), once one adds that \(F\) realises its
expansion in the sense of \cref{plt:def:cmp-realization} --- which it does,
being a sum of three elementary functions --- and invokes the transfer of
\cref{plt:prop:rev-analytic-transfer}; that step is analysis, and it is not
supplied by the template.  The instructive point of the example is
elsewhere: it is where the coefficient ring first fails to be polynomial.
\end{example}

% ed.: S6 states that ``the coefficient of Y^{-2} is no longer a polynomial
% in L and M alone''.  That is true but not sharp: expanding the Y^{-1}
% block gives A*varpi = alpha^2 L + alpha*gamma M + alpha*gamma
% + gamma^2 M/L, so a negative power of L is already present one block
% earlier whenever gamma is nonzero.  The corrected statement and its proof
% are below.
\begin{proposition}[The first non-polynomial coefficient]
\label{plt:prop:lw-ex2-sharp}
In \cref{plt:eq:lw-ex2-inverse}, expand the block of \(Y^{-1}\):
\begin{equation}\label{plt:eq:lw-ex2-first-block}
 A\varpi=(\alpha L+\gamma M)\left(\alpha+\frac{\gamma}{L}\right)
 =\alpha^{2}L+\alpha\gamma M+\alpha\gamma
 +\frac{\gamma^{2}M}{L}.
\end{equation}
Hence for every \(\gamma\ne0\) the coefficient of \(Y^{-1}\) already fails
to lie in \(\K[L,M]\) --- one block earlier than S6 records --- and the
natural coefficient ring is \(\K((L^{-1}))[M]\).  For \(\gamma=0\) the
problem degenerates to the affine-logarithmic case and every block is in
\(\K[L]\).
\end{proposition}

\begin{proof}
The displayed product is immediate.  The functions \(L=\log Y\) and
\(M=\log\log Y\) are algebraically independent over \(\K\): if
\(P\in\K[U,V]\) were nonzero with \(P(L,M)\equiv0\) for large \(Y\), then
among the monomials \(L^{i}M^{j}\) occurring in \(P\) there would be a
unique dominant one by \cref{plt:lem:mot-dominance} (the order being
lexicographic in \((i,j)\), since \(M=\LittleOAt{L^{\epsilon}}{Y\to+\infty}\)
for every \(\epsilon>0\)), and a dominant term cannot cancel.  Hence
\(\K[L,M]\) is a polynomial ring in two indeterminates, in which \(L\) is
prime and does not divide \(M\); therefore \(M/L\notin\K[L,M]\).  Since
\(\alpha^{2}L+\alpha\gamma M+\alpha\gamma\in\K[L,M]\) and
\(\gamma^{2}\ne0\), the sum \cref{plt:eq:lw-ex2-first-block} lies outside
\(\K[L,M]\).  For \(\gamma=0\), \(A=\alpha L\) and
\cref{plt:thm:lw-affine-log} applies, giving
\(q_n=\alpha^{n+1}\LamP{n}(L)\in\K[L]\).
\end{proof}

\begin{remark}[Numerical check]
With \(\alpha=\gamma=1\) and \(Y=10^{10}\), the absolute errors of the
truncations of \cref{plt:eq:lw-ex2-inverse} at orders
\(Y^{0},Y^{-1},Y^{-2}\) are \(2.73\cdot10^{-9}\), \(3.29\cdot10^{-18}\) and
\(5.16\cdot10^{-27}\); the ratio of consecutive errors is about
\(Y^{-1}L\), as \cref{plt:eq:lw-ex2-inverse} predicts.
\end{remark}

\begin{example}[An equation that fails the hypotheses]
\label{plt:ex:lw-failure}
Solve
\begin{equation}\label{plt:eq:lw-ex3-map}
 Y=X+\frac{X}{\log X}
\end{equation}
for \(X\) as \(Y\to+\infty\).

The map looks like a perturbation of the identity, and it is one
analytically: \(X/\log X=\LittleOAt{X}{X\to+\infty}\).  It is not one
\emph{in the sense of \(\Gnear\)}.  Writing \(t=X^{-1}\) and \(L=\log X\),
the perturbation is \(L^{-1}t^{-1}\); read inside \(\PLit\) it has
\(\ordt=-1\), the same outer order as \(X\) itself, and its coefficient
\(L^{-1}\) is not in \(\K[L]\).  So hypothesis (H3) of
\cref{plt:def:lw-monomial-log-datum} fails twice over --- wrong outer order
and wrong coefficient ring --- and \cref{plt:thm:rev-existence} does not
apply.

The correct diagnosis is that the problem is a near-identity problem one
level down.  Put \(\lambda=\log Y\) and \(Z=\log X\); taking logarithms of
\cref{plt:eq:lw-ex3-map},
\begin{equation}\label{plt:eq:lw-ex3-log}
 \lambda=Z+\log\left(1+\frac1Z\right)
 =Z+\frac1Z-\frac1{2Z^{2}}+\frac1{3Z^{3}}-\cdots,
\end{equation}
which \emph{is} of monomial--logarithmic type with
\(\rho=1\), \(\mu=0\), \(\beta=0\) and
\(B(Z)=\log(1+1/Z)\in Z^{-1}\K[[Z^{-1}]]\).  \Cref{plt:thm:lw-template}
therefore applies in the logarithmic chart and gives, with all \(q_n\)
constants because \(a=\mu/\rho=0\) and every \(b_m\) is constant,
\begin{equation}\label{plt:eq:lw-ex3-Z}
 Z=\lambda-\frac1{\lambda}+\frac1{2\lambda^{2}}
 -\frac4{3\lambda^{3}}+\frac7{4\lambda^{4}}
 -\frac{121}{30\lambda^{5}}+\cdots.
\end{equation}
Re-exponentiating, or equivalently using
\(X=Y\bigl(1-1/(Z+1)\bigr)\),
\begin{equation}\label{plt:eq:lw-ex3-X}
 X=Y\left(1-\frac1{\lambda}+\frac1{\lambda^{2}}
 -\frac2{\lambda^{3}}+\frac7{2\lambda^{4}}
 -\frac{22}{3\lambda^{5}}
 +\frac{179}{12\lambda^{6}}-\cdots\right),
 \qquad\lambda=\log Y .
\end{equation}
The relative errors of the first four truncations at \(Y=10^{40}\) are
\(1.09\cdot10^{-2}\), \(1.17\cdot10^{-4}\), \(2.54\cdot10^{-6}\),
\(4.81\cdot10^{-8}\).
\end{example}

\begin{proposition}[The answer is not a polynomial--logarithmic series]
\label{plt:prop:lw-failure-not-pl}
Let \(X=X(Y)\) be the solution of \cref{plt:eq:lw-ex3-map} for
\(Y\to+\infty\).  There is no
\(G=\sum_{n\ge n_0}p_n(L)t^{n}\in\PLlau\) --- with \(t=Y^{-1}\),
\(L=\log Y\), coefficients \(p_n\in\K[L]\) --- that is an asymptotic
expansion of \(X\) in the sense of \cref{plt:def:mul-asymptotic}.  The
answer \cref{plt:eq:lw-ex3-X} does lie in \(\PLit=\K((L^{-1}))((t))\).
\end{proposition}

\begin{proof}
Suppose such a \(G\) existed.  Since \(X/Y\to1\), the outer order is
\(n_0=-1\), and \cref{plt:def:mul-asymptotic} implies in particular that,
for some \(d\in\NaturalNumbers\) and all large \(Y\),
\begin{equation}\label{plt:eq:lw-notpl}
 X=p_{-1}(L)\,Y+\BigOAt{L^{d}}{Y\to+\infty},
\end{equation}
the error absorbing the block \(p_0(L)\) together with the tail.
% ed.: the proof as drafted read the two leading blocks off
% Equation lw-ex3-X, which the template produces only FORMALLY; using it as
% an analytic asymptotic statement here would be circular.  Replaced by a
% direct three-line derivation from the defining equation, which needs
% nothing beyond X -> infinity.
The two leading terms of \(X/Y\) are obtained directly from
\cref{plt:eq:lw-ex3-map}, with no appeal to \cref{plt:eq:lw-ex3-X}: writing
\(Z=\log X\), the equation \(Y=X(1+1/Z)\) gives
\[
 \frac XY=\frac{Z}{Z+1}=1-\frac1{Z+1},
 \qquad
 L=\log Y=Z+\log\Bigl(1+\frac1Z\Bigr)=Z+\BigOAt{Z^{-1}}{Y\to+\infty},
\]
so \(Z=L+\BigOAt{L^{-1}}{Y\to+\infty}\) and hence
\[
 \frac XY=1-\frac1L+\BigOAt{L^{-2}}{Y\to+\infty},
\]
whereas \cref{plt:eq:lw-notpl} gives
\(X/Y=p_{-1}(L)+\BigOAt{L^{d}/Y}{Y\to+\infty}\).  Subtracting,
\[
 p_{-1}(L)-1+\frac1L
 =\BigOAt{L^{-2}}{Y\to+\infty}+\BigOAt{L^{d}/Y}{Y\to+\infty}.
\]
As \(Y\to+\infty\) both right-hand terms tend to \(0\) (the second because
\(L^{d}/Y\to0\) for every fixed \(d\)), while \(1/L\to0\); hence
\(p_{-1}(L)\to1\) along \(L\to+\infty\), which for a polynomial forces
\(p_{-1}\equiv1\).  But then the left side is exactly \(1/L\), which is not
\(\BigOAt{L^{-2}}{Y\to+\infty}+\BigOAt{L^{d}/Y}{Y\to+\infty}\), because
\(L\cdot(1/L)=1\not\to0\).  Contradiction.

For the positive half, \cref{plt:eq:lw-ex3-X} exhibits \(X\) as
\(c(L)\,t^{-1}\) with \(c(L)=1-L^{-1}+L^{-2}-2L^{-3}+\cdots\in
\K((L^{-1}))\), which is one block of an element of \(\PLit\)
(\cref{plt:def:ring-four-classes}); the inclusion
\(\PLlau\subsetneq\PLit\) and its strictness are
\cref{plt:prop:ring-PLit-strict}.
\end{proof}

\begin{remark}[The other way to fail: no monomial at all]
\label{plt:rmk:lw-rho-zero}
Hypothesis (H1) fails for \(Y=(\log X)^{2}\), where
\cref{plt:eq:lw-template-equation} reads \(0\cdot Z+2\log Z=\xi\).  The
solution is \(X=\exp\bigl(\EulerE^{\xi/2}\bigr)
=\exp\bigl(\sqrt{Y}\bigr)\) --- an exponential of a power of the large
variable, hence outside every polynomial--logarithmic class by
\cref{plt:prop:grp-exponential}.  No amount of coefficient-ring enlargement
recovers it: the failure is a change of exponential level, not of
coefficient field.  The rule of thumb is that (H1) controls the
\emph{level}, (H2) the shape of the leading logarithm, and (H3) the
coefficient field; a violation of each is repaired, if at all, in a
different direction.
\end{remark}

\begin{remark}[Beyond all orders]
\label{plt:rmk:lw-beyond-all-orders}
A last failure mode passes every hypothesis and still leaves the answer
undetermined.  For \(Y=X+\EulerE^{-X}\), the perturbation
\(\EulerE^{-X}\) is not in \(\K[L][[t]]\), but it is
\(\LittleOAt{t^{n}}{X\to+\infty}\) for every \(n\); the template applied to
the truncated data returns \(X=Y\), and every polynomial--logarithmic
truncation of the true inverse is indeed \(Y\).  The expansion is correct
and useless: it does not see \(\EulerE^{-Y}\).  The same phenomenon appears
in \cref{plt:prop:lw-r-lambert}, where an entire one-parameter family of
functions shares one expansion.  Recovering the invisible terms requires
transseries with exponential monomials \cite{ecalle,costin-book,adh-book},
which are outside the classes of \cref{plt:sec:rings}.
\end{remark}

\begin{remark}[Two classical instances, recorded for reference]
\label{plt:rmk:lw-classical-instances}
The template covers, with no new work, two examples that recur across the
sources.
\begin{enumerate}
\item \(Y=X\log X\).  Put \(u=\log X\); then \(u\EulerE^{u}=Y\), so
      \(X=\EulerE^{\LambertW(Y)}=Y/\LambertW(Y)\), and dividing by the
      branch expansion, using the reciprocal recurrence of
      \cref{plt:thm:div-reciprocal},
      \[
       X=\frac{Y}{L_1}\left(1+\frac{L_2}{L_1}
       +\frac{L_2^{2}-L_2}{L_1^{2}}
       +\frac{L_2^{3}-\tfrac52L_2^{2}+L_2}{L_1^{3}}
       +\BigOAt{L_2^{4}/L_1^{4}}{Y\to+\infty}\right),
      \]
      with \(L_1=\log Y\), \(L_2=\log\log Y\).  This is the form given in
      S3 and S6, and it is correct.
\item \(X=y/\log y\).  Setting \(v=y/X\) turns this into
      \(v-\log v=\log X\), which is the lower-branch equation of
      \cref{plt:prop:lw-lower-equation} with
      \(S=T\), \(T\DefinitionEquals\log X\); hence
      \(y=-X\LambertW_{-1}(-1/X)\) and, with \(\Lambda=\log T\),
      \[
       y=X\left(T+\Lambda
       +\sum_{n\ge1}(-1)^{n+1}\frac{\LamP{n}(\Lambda)}{T^{n}}\right),
      \]
      convergent whenever \(2(1+\AbsoluteValue{\Lambda})<T\) by
      \cref{plt:thm:lw-lower-branch}.  Note the sign: it is
      \((-1)^{n+1}\), not \((-1)^{n}\), because the reflection
      \(y=Xv\) has already been undone.
\end{enumerate}
\end{remark}

\begin{remark}[Reductions that are exact, and reductions that are not]
It is worth separating two things the sources sometimes present together.
\Cref{plt:prop:lw-power-exp} and \cref{plt:rmk:lw-classical-instances} are
\emph{exact} reductions: an algebraic change of variable turns the equation
into \(u+\log u=\zeta\) with no error term, so the whole of
\cref{plt:sec:lambert-branches} applies verbatim, convergence statements
included.  \Cref{plt:thm:lw-template} is not of that kind: it produces a
formal inverse in \(\Gnear\) whose analytic status must be argued
separately, either by \cref{plt:thm:cmp-analytic-transfer} (if the direct
map realises its expansion) or by an inverse-function estimate as in
\cref{plt:prop:rev-analytic-transfer}.  S6's example
\(Y=X+\alpha\log X+\beta+(\gamma\log X+\delta)/X+\cdots\), whose inverse is
\[
 X=Y-\alpha L-\beta
 +\frac{(\alpha^{2}-\gamma)L+\alpha\beta-\delta}{Y}
 +\BigOAt{L^{2}/Y^{2}}{Y\to+\infty},
 \qquad L=\log Y,
\]
is of the second kind.  Its displayed coefficients are correct, and they are
the case \(\rho=1\), \(\mu=\alpha\), \(b_1=\gamma\ell+\delta\) of
\cref{plt:thm:lw-template}: the template shifts the constant away, working
in \(X=Y-\beta\) and \(\ell=\log(Y-\beta)\), where
\(q_1(\ell)=(\alpha^{2}-\gamma)\ell-\delta\) with no constant term, and the
extra \(\alpha\beta\) in the display above is produced when \(\ell\) is
re-expanded as \(L-\beta/Y+\BigOAt{Y^{-2}}{Y\to+\infty}\).  The error term,
however, is legitimate only under a hypothesis on the direct map, not merely
on its formal expansion.
\end{remark}

\section{Reseating the expansion, and the associated Stirling numbers}
\label{plt:sec:lw-associated}

\Cref{plt:thm:lw-basepoint} was stated with an arbitrary base point \(K\)
for a reason: the choice \(K=\xi_k\) is convenient but not optimal, and the
literature on Lambert \(\LambertW\) contains a whole \emph{sequence} of
series obtained by other choices \cite{corless-jeffrey-knuth-1997}, as well
as rearrangements whose coefficients are associated Stirling numbers rather
than ordinary ones \cite{jeffrey-et-al-1995,comtet1970}.  This subsection
isolates the mechanism, which is entirely elementary once the two-variable
picture of \cref{plt:lem:lw-two-variable} is available: the coefficient
array \((c_{ab})\) can be summed along its diagonals \(a+b=n\), which gives
the Lambert polynomials, or along its rows \(a=\text{const}\), which gives
something else --- and the something else turns out to be the same family
again, evaluated elsewhere.

\begin{theorem}[Transverse resummation]
\label{plt:thm:lw-transverse}
In \(\RationalNumbers[[\sigma,\tau]]\), let \(\mathcal Y\) be the unique
solution with \(\mathcal Y(0,0)=0\) of
\(\mathcal Y=-\log(1-\tau+\sigma\mathcal Y)\).  Then
\begin{equation}\label{plt:eq:lw-transverse}
 \mathcal Y(\sigma,\tau)
 =\sum_{a\ge0}\LamP{a}\bigl(\log(1-\tau)\bigr)
 \left(\frac{\sigma}{1-\tau}\right)^{a},
\end{equation}
where \(\log(1-\tau)=-\sum_{m\ge1}\tau^{m}/m\in\tau\RationalNumbers[[\tau]]\).
Equivalently, writing
\(\LamP{n}(u)=\sum_{m}\lambda_{n,m}u^{m}\), for all \(a,b\ge0\) with
\(a+b\ge1\),
\begin{equation}\label{plt:eq:lw-transverse-coefficients}
 \lambda_{a+b,b}
 =[\tau^{b}]\left\{
 \frac{\LamP{a}\bigl(\log(1-\tau)\bigr)}{(1-\tau)^{a}}\right\}.
\end{equation}
\end{theorem}

\begin{proof}
Existence and uniqueness of \(\mathcal Y\) in
\(\RationalNumbers[[\sigma,\tau]]\) with \(\mathcal Y(0,0)=0\) follow from
the \((\sigma,\tau)\)-adic contraction argument, the map
\(y\mapsto-\log(1-\tau+\sigma y)\) raising the order by one on the maximal
ideal; the analytic realisation is
\cref{plt:lem:lw-two-variable}.  Put
\(\zeta=1-\tau\), \(\Lambda=-\log\zeta\in\tau\RationalNumbers[[\tau]]\) and
\(\widehat\sigma=\sigma/\zeta\).  Then
\[
 \mathcal Y=-\log(\zeta+\sigma\mathcal Y)
 =-\log\zeta-\log\bigl(1+\widehat\sigma\mathcal Y\bigr)
 =\Lambda-\log\bigl(1+\widehat\sigma\mathcal Y\bigr).
\]
Set \(\mathcal Q\DefinitionEquals\mathcal Y-\Lambda\).  Substituting
\(\mathcal Y=\Lambda+\mathcal Q\),
\[
 \mathcal Q=-\log\bigl(1+\widehat\sigma\Lambda
 +\widehat\sigma\mathcal Q\bigr)
 =-\log\bigl(1-\widehat\sigma\,(-\Lambda)
 +\widehat\sigma\mathcal Q\bigr),
\]
which is precisely the generating equation \cref{plt:eq:lw-generating} with
\(s=\widehat\sigma\) and \(u=-\Lambda\).  Moreover \(\mathcal Q\) has
positive order in \(\widehat\sigma\).  By the uniqueness clause of
\cref{plt:thm:om-generating-equation}, transported along the continuous
\(\RationalNumbers\)-algebra homomorphism
\(\RationalNumbers[u][[s]]\to\RationalNumbers[[\sigma,\tau]]\) sending
\(s\mapsto\widehat\sigma\) and \(u\mapsto-\Lambda\),
\[
 \mathcal Q=\sum_{a\ge1}\LamP{a}(-\Lambda)\,\widehat\sigma^{\,a}.
\]
Since \(\LamP{0}(-\Lambda)=\Lambda\), adding \(\Lambda\) gives
\(\mathcal Y=\sum_{a\ge0}\LamP{a}(-\Lambda)\widehat\sigma^{\,a}\), which is
\cref{plt:eq:lw-transverse} because \(-\Lambda=\log(1-\tau)\).

For \cref{plt:eq:lw-transverse-coefficients}, compare with the diagonal
grouping.  Restricting \cref{plt:eq:lw-c-to-P} to
\(\tau=\sigma u\) identified \(\sum_{a+b=n}c_{ab}u^{b}=\LamP{n}(u)\), i.e.\
\(c_{ab}=\lambda_{a+b,b}\) for \(a+b\ge1\).  Extracting
\(\sigma^{a}\tau^{b}\) from \cref{plt:eq:lw-transverse} gives
\(c_{ab}=[\tau^{b}]\{\LamP{a}(\log(1-\tau))(1-\tau)^{-a}\}\).
\end{proof}

\begin{corollary}[The edge coefficients of the Lambert polynomials, twice]
\label{plt:cor:lw-edge-coefficients}
Taking \(a=0\) and \(a=1\) in
\cref{plt:eq:lw-transverse-coefficients}:
\begin{align}
 \lambda_{b,b}&=[\tau^{b}]\bigl(-\log(1-\tau)\bigr)=\frac1b
 \qquad(b\ge1),
 \label{plt:eq:lw-edge-0}\\
 \lambda_{b+1,b}&=[\tau^{b}]\frac{\log(1-\tau)}{1-\tau}=-H_{b}
 \qquad(b\ge0),
 \label{plt:eq:lw-edge-1}
\end{align}
where \(H_b=\sum_{j=1}^{b}1/j\).  These recover the leading coefficient
\(1/n\) and the next coefficient \(-H_{n-1}\) of \(\LamP{n}\) recorded in
\cref{plt:cor:lambert-coefficients}, now as the Taylor coefficients of two
elementary functions.  Continuing to \(a=2\),
\[
 \frac{\LamP{2}\bigl(\log(1-\tau)\bigr)}{(1-\tau)^{2}}
 =\tau+3\tau^{2}+\tfrac{35}{6}\tau^{3}+\tfrac{75}{8}\tau^{4}
 +\tfrac{203}{15}\tau^{5}+\cdots,
\]
whose coefficients are
\(\lambda_{3,1},\lambda_{4,2},\lambda_{5,3},\lambda_{6,4},\lambda_{7,5}\).
\end{corollary}

\begin{proof}
For \(a=0\), \(\LamP{0}(u)=-u\) gives \(-\log(1-\tau)
=\sum_{b\ge1}\tau^{b}/b\).  For \(a=1\), \(\LamP{1}(u)=u\) gives
\(\log(1-\tau)/(1-\tau)\), and
\(\sum_{b\ge1}H_b\tau^{b}=-\log(1-\tau)/(1-\tau)\) is the classical
Cauchy-product identity
\(\bigl(\sum_{m\ge1}\tau^{m}/m\bigr)\bigl(\sum_{j\ge0}\tau^{j}\bigr)
=\sum_{b\ge1}H_b\tau^{b}\), with \(H_0=0\).  The listed values for
\(a=2\) are obtained by the same expansion.
\end{proof}

\begin{corollary}[Reseating at the two-term approximation]
\label{plt:cor:lw-reseated}
In \cref{plt:thm:lw-basepoint}, take the base point
\(\widetilde K\DefinitionEquals K-M\) with
\(\log\widetilde K\DefinitionEquals\log K
+\PrincipalLogarithm(1-M/K)\), assuming \(\AbsoluteValue{M/K}<1\).  Then the
associated parameter is
\begin{equation}\label{plt:eq:lw-reseated-M}
 \widetilde M=\PrincipalLogarithm\!\left(1-\frac MK\right),
 \qquad
 \AbsoluteValue{\widetilde M}
 \le-\log\!\left(1-\frac{\AbsoluteValue{M}}{\AbsoluteValue{K}}\right),
\end{equation}
so that \(\widetilde M\to0\) when \(M/K\to0\).  If
\(2(1+\AbsoluteValue{\widetilde M})<\AbsoluteValue{\widetilde K}\), then the
series
\(\widetilde w\DefinitionEquals\widetilde K
+\sum_{a\ge0}\LamP{a}(\widetilde M)\widetilde K^{-a}\)
converges absolutely, with the remainder bound
\cref{plt:eq:lw-remainder} in the reseated data, and
\(\widetilde w\EulerE^{\widetilde w}=\EulerE^{\xi}\).
% ed.: as drafted this corollary wrote the reseated sum as "w", silently
% asserting that the reseated base point returns the SAME solution as the
% canonical one, and its proof discharged the point in one clause.  That
% does not follow from the stated hypotheses: Theorem lw-basepoint attaches
% to each base point SOME solution of w e^w = e^xi, and the reseated one is
% only known to satisfy |w~ - (xi - ell)| <= |M~| + log 2, which may exceed
% the radius 1 of the uniqueness disc of Proposition lw-branch-label.
% Repaired: the identification is now a separate clause with the printed
% hypothesis 4(1 + |ell_k|) <= |xi_k|, under which the disc argument does
% apply, and the real case is noted to need nothing at all.
In the canonical case
\(K=\xi_k\), \(M=\ell_k\), assume in addition to the hypotheses of
\cref{plt:prop:lw-branch-label} that
\begin{equation}\label{plt:eq:lw-reseated-hyp}
 4\bigl(1+\AbsoluteValue{\ell_k}\bigr)\le\AbsoluteValue{\xi_k}.
\end{equation}
Then the reseated seating condition holds automatically,
\(\widetilde w=\LambertW_k(z)\), and
\begin{equation}\label{plt:eq:lw-reseated-canonical}
 \LambertW_k(z)
 =\bigl(\xi_k-\ell_k\bigr)
 +\sum_{a\ge0}\frac{\LamP{a}\bigl(
 \PrincipalLogarithm(1-\ell_k/\xi_k)\bigr)}{(\xi_k-\ell_k)^{a}} .
\end{equation}
For real \(z>1\) and \(k=0\) the identification
\(\widetilde w=\LambertW_0(z)\) needs neither
\cref{plt:eq:lw-reseated-hyp} nor \cref{plt:prop:lw-branch-label}: all the
reseated data are then real, hence so is \(\widetilde w\), and a positive
real argument has exactly one real preimage under \(w\EulerE^{w}\) by
\cref{plt:prop:lw-real-branches}\,\textup{(iii)}.
\end{corollary}

\begin{proof}
By \cref{plt:eq:lw-M}, the parameter attached to a base point
\(\widetilde K\) with a chosen logarithm is
\(\widetilde M=\widetilde K+\log\widetilde K-\xi\).  With
\(\widetilde K=K-M\) and the stated logarithm,
\[
 \widetilde M=(K-M)+\log K+\PrincipalLogarithm(1-M/K)-\xi
 =(K+\log K-\xi)-M+\PrincipalLogarithm(1-M/K),
\]
and the bracket is \(M\), leaving
\(\widetilde M=\PrincipalLogarithm(1-M/K)\).  The stated modulus bound is
\cref{plt:eq:lw-log-bound}.  Note \(\EulerE^{\log\widetilde K}
=K(1-M/K)=K-M=\widetilde K\), so the chosen logarithm is legitimate.  The
convergence statement, the remainder bound and
\(\widetilde w\EulerE^{\widetilde w}=\EulerE^{\xi}\) are then
\cref{plt:thm:lw-basepoint,plt:thm:lw-uniform-remainder} applied verbatim to
the reseated data \((\widetilde K,\widetilde M)\).

It remains to identify \(\widetilde w\) in the canonical case.  Assume
\cref{plt:eq:lw-reseated-hyp}, so \(\AbsoluteValue{\ell_k}
\le\tfrac14\AbsoluteValue{\xi_k}-1\) and in particular
\(\AbsoluteValue{\ell_k}/\AbsoluteValue{\xi_k}\le\tfrac14\) and
\(\AbsoluteValue{\xi_k}\ge4\).  By \cref{plt:eq:lw-reseated-M},
\(\AbsoluteValue{\widetilde M}\le-\log(1-\tfrac14)=\log\tfrac43<0.2877\),
while \(\AbsoluteValue{\widetilde K}\ge\AbsoluteValue{\xi_k}
-\AbsoluteValue{\ell_k}\ge\tfrac34\AbsoluteValue{\xi_k}+1\ge4\); since
\(2(1+\AbsoluteValue{\widetilde M})<2.576<4\), the reseated seating
condition holds automatically.  \Cref{plt:thm:lw-basepoint}\,(ii) applied to
the reseated data gives
\(\AbsoluteValue{\widetilde w-(\widetilde K-\widetilde M)}\le\log2\), hence
\[
 \bigl\lvert\widetilde w-(\xi_k-\ell_k)\bigr\rvert
 \le\AbsoluteValue{\widetilde M}+\log2<0.2877+0.6932<1 .
\]
So \(\widetilde w\) is a solution of \(w\EulerE^{w}=z\) lying in the closed
disc \(\AbsoluteValue{v-(\xi_k-\ell_k)}\le1\); by the uniqueness clause of
\cref{plt:prop:lw-branch-label} --- whose hypothesis
\(2(1+\AbsoluteValue{\ell_k})<\AbsoluteValue{\xi_k}\) follows from
\cref{plt:eq:lw-reseated-hyp} --- it equals \(w_k(z)\), which under the
distance-to-the-cut hypothesis \cref{plt:eq:lw-cut-distance} is
\(\LambertW_k(z)\).  This is \cref{plt:eq:lw-reseated-canonical}.

For the real case, let \(z>1\) be real and \(k=0\).  Then \(\xi_0=\log z>0\)
and \(\ell_0=\log\xi_0<\xi_0\), so \(\widetilde K=\xi_0-\ell_0>0\) and
\(\widetilde M=\log(1-\ell_0/\xi_0)\) are real; the fixed point
\(\mathcal Y(\sigma,\tau)\) of \cref{plt:lem:lw-two-variable} is real for
real \((\sigma,\tau)\in D\), because the contraction \(\Phi\) then maps
\([-1,1]\) into itself and its fixed point there must be the unique fixed
point in \(\AbsoluteValue{y}\le1\).  Hence \(\widetilde w\) is real and
satisfies \(\widetilde w\EulerE^{\widetilde w}=z>0\), and
\cref{plt:prop:lw-real-branches}\,(iii) leaves it no choice but
\(\LambertW_0(z)\).
\end{proof}

\begin{remark}[What this buys, quantitatively]
\label{plt:rmk:lw-reseated-numbers}
The reseated parameter \(\widetilde M=\PrincipalLogarithm(1-\ell_k/\xi_k)\)
is \(\BigOAt{\ell_k/\xi_k}{\AbsoluteValue{\xi_k}\to\infty}\), whereas
\(M=\ell_k\to\infty\).  The seating condition
\cref{plt:eq:lw-seating} therefore holds for the reseated series on a
% ed.: the canonical threshold was printed as 5.3569...; the root of
% xi = 2 + 2 log xi is 5.356693980..., i.e. 5.3567 to four decimals.
% Corrected, and the threshold of the reseated condition added.
strictly larger region: for real \(z\) it becomes
\(2(1+\AbsoluteValue{\log(1-\ell_0/\xi_0)})<\xi_0-\ell_0\), whose threshold
is \(\xi_0=4.2849\ldots\), so that it holds already at \(\xi_0=5\), whereas
the canonical condition \(2(1+\log\xi_0)<\xi_0\) needs
\(\xi_0>5.3567\ldots\) (the root of \(\xi=2+2\log\xi\), namely
\(5.356693980\ldots\)).  Numerically, with \(K=\xi_0\) and
\(\widetilde K=\xi_0-\ell_0\):
\begin{center}
\begin{tabular}{@{}rlllll@{}}
\toprule
\(\xi_0\) & base point & \(M\) & \(N=2\) & \(N=5\) & \(N=10\)\\
\midrule
\(5\)  & \(\xi_0\)             & \(+1.6094\)
       & \(6.4\cdot10^{-3}\) & \(3.0\cdot10^{-5}\) & \(1.4\cdot10^{-7}\)\\
       & \(\xi_0-\ell_0\)      & \(-0.3884\)
       & \(1.1\cdot10^{-2}\) & \(1.1\cdot10^{-3}\) & \(4.9\cdot10^{-5}\)\\
\(10\) & \(\xi_0\)             & \(+2.3026\)
       & \(1.7\cdot10^{-3}\) & \(5.1\cdot10^{-6}\) & \(6.4\cdot10^{-10}\)\\
       & \(\xi_0-\ell_0\)      & \(-0.2617\)
       & \(6.9\cdot10^{-4}\) & \(4.3\cdot10^{-6}\) & \(1.7\cdot10^{-9}\)\\
\(30\) & \(\xi_0\)             & \(+3.4012\)
       & \(4.1\cdot10^{-5}\) & \(1.5\cdot10^{-9}\) & \(8.2\cdot10^{-14}\)\\
       & \(\xi_0-\ell_0\)      & \(-0.1203\)
       & \(7.3\cdot10^{-6}\) & \(6.9\cdot10^{-10}\) & \(2.2\cdot10^{-16}\)\\
\bottomrule
\end{tabular}
\end{center}
% ed.: the drafted reading said reseating "is numerically better at xi_0 =
% 10 and 30", which its own table contradicts at xi_0 = 10, N = 10
% (6.4e-10 canonical against 1.7e-9 reseated).  Corrected to what the
% table shows.
The honest reading: reseating enlarges the region on which convergence is
\emph{proved}, and it is numerically better at \(\xi_0=30\) throughout the
displayed range, and at \(\xi_0=10\) at the two lower orders --- but not at
\(N=10\), where the canonical series is already ahead by a factor of about
three; at \(\xi_0=5\) the canonical series is better at every displayed
order, even though nothing here proves that it converges at all.  A larger
proved domain and a smaller error at a given point are different goods, and
the table shows that neither implies the other.
\end{remark}

\begin{remark}[Associated Stirling numbers]
\label{plt:rmk:lw-associated-stirling}
\Cref{plt:thm:lw-transverse} is the structural reason why a second family of
combinatorial numbers appears in this problem.  Summing the coefficient
array \((c_{ab})\) along diagonals \(a+b=n\) gives \(\LamP{n}\), whose
coefficients are ordinary Stirling cycle numbers
(\cref{plt:thm:lambert-stirling}).  Summing along rows \(a=\text{const}\)
gives the functions
\(\LamP{a}(\log(1-\tau))(1-\tau)^{-a}\), whose Taylor coefficients are a
different array.  Comtet inverted \(y^{\alpha}\EulerE^{y}\) and
\(y\log y\) \cite{comtet1970}, and Jeffrey, Corless, Hare and Knuth
identified the coefficients of the row-type rearrangements with
\emph{associated} Stirling numbers --- those counting permutations all of
whose cycles have length at least two --- and used them to obtain series
converging on larger domains \cite{jeffrey-et-al-1995}; Corless, Jeffrey and
Knuth developed the resulting sequence of series systematically
\cite{corless-jeffrey-knuth-1997}.  S5's observation that the two small
parameters are \(\sigma=\alpha/L_1\) and \(\tau=\alpha L_2/L_1\), and that
associated Stirling numbers arise from ``a different small-parameter
geometry'', is correct and is made precise by
\cref{plt:thm:lw-transverse,plt:cor:lw-reseated}: the different geometry is
the base-point change \(K\rightsquigarrow K-M\).  This volume does not prove
the identification of the row coefficients with associated Stirling numbers;
that identification is cited, not derived, and
\cref{plt:eq:lw-transverse-coefficients} is what is proved here.
\end{remark}

\begin{summarybox}
Three layers.  \emph{Exact reduction}: \(y^{\alpha}\EulerE^{\beta y}=x\) and
its relatives become \(u+\log u=\zeta\) by an algebraic substitution, so
\cref{plt:sec:lambert-branches} transfers verbatim, convergence included
(\cref{plt:prop:lw-power-exp,plt:thm:lw-generalized-expansion}); the only
new content is bookkeeping, and \(\LamP{n}\) reappears with the scalar
\(a^{n+1}\), \(a=\mu/\rho\).  \emph{Template}: any equation whose dominant
balance is a monomial times a power of a logarithm, perturbed by a series of
positive outer order with polynomial logarithmic coefficients, has a
mechanically computable polynomial--logarithmic inverse
(\cref{plt:thm:lw-template}), whose leading blocks are the Lambert
polynomials and whose later blocks differ from them by explicitly computable
corrections.  \emph{Boundary}: the hypotheses are not decorative.  A zero
monomial exponent changes exponential level
(\cref{plt:rmk:lw-rho-zero}); a perturbation of outer order zero forces the
iterated Laurent coefficient field \(\PLit\)
(\cref{plt:prop:lw-failure-not-pl}); an exponentially small perturbation is
invisible to the whole apparatus
(\cref{plt:rmk:lw-beyond-all-orders}).
\end{summarybox}

\begin{exercise}
Carry \cref{plt:eq:lw-gen-expansion} to \(n=3\) for
\(\alpha=\tfrac12,\beta=1\) and verify the seating condition at
\(x=\EulerE^{40}\).
\end{exercise}

\begin{exercise}
Apply \cref{plt:thm:lw-template} to
\(Y=X+\alpha\log X+\beta+(\gamma\log X+\delta)/X\), in the shifted chart
\(X_{\ast}=Y-\beta\), \(\ell=\log(Y-\beta)\), and check from
\cref{plt:thm:lag-coefficients} that
\(q_1(\ell)=(\alpha^{2}-\gamma)\ell-\delta\).  Identify the part
\(a^{2}\LamP{1}(\ell)=\alpha^{2}\ell\) and the correction
\(E_1=-\gamma\ell-\delta\).  Then re-expand \(\ell\) in
\(L=\log Y\) and recover the block
\(\bigl((\alpha^{2}-\gamma)L+\alpha\beta-\delta\bigr)/Y\) displayed above.
\end{exercise}

\begin{exercise}
Show that \(\LamP{3}(\log(1-\tau))(1-\tau)^{-3}\) has Taylor coefficients
\(0,-1,-5,-\tfrac{85}{6},-\tfrac{245}{8},\dots\), and identify them among
the coefficients of \(\LamP{3},\LamP{4},\LamP{5},\LamP{6},\LamP{7}\).
\end{exercise}

\begin{exercise}
Prove that the logarithmic integral satisfies
\(\operatorname{li}(x)\sim x/\log x\cdot\sum_{j\ge0}j!(\log x)^{-j}\) and
that formal inversion of this relation is an instance of
\cref{plt:thm:lw-template} with \(\rho=1\), \(\mu=-1\); deduce the shape of
the first three blocks of the inverse, and explain why an analytic theorem
about \(\operatorname{li}\) --- not the template --- is what is needed to
pass from the formal inverse to an asymptotic formula for the \(n\)th prime.
\end{exercise}

\clearpage
\part{From formal transseries to analytic asymptotics}

% ---- fragment 12-analytic ----
\chapter{Poincar\'e asymptotics on the polynomial--logarithmic scale}
\label{plt:sec:poincare}

Everything proved so far about \(\PLpow\), \(\PLlau\), \(\PLrat\) and
\(\PLit\) is algebra: statements about coefficient arrays, the valuation
\(\ordt\), and the \(t\)-adic filtration.  None of it mentions a function, a
limit, or an inequality.  This part supplies the bridge, and states its toll.
The toll is real: four of the six sources assert asymptotic validity,
uniformity in a parameter, or termwise differentiability on the strength of a
formal identity alone, and each such assertion is marked with a \verb|% ed.|
comment and repaired where it occurs below.

Two relations between a function and a series have already been used in this
volume: \(f\approx F\) of \cref{plt:def:mul-asymptotic}, introduced with a
\(\BigOAt{\cdot}{\cdot}\) remainder carrying an unspecified power of
\(\log X\), and \(f\sim_{S}F\) of \cref{plt:def:cmp-realization}, introduced
with a \(\LittleOAt{\cdot}{\cdot}\) remainder.  They were introduced
independently, in different parts, for different purposes.  The first order
of business is to prove that they are the same relation
(\cref{plt:thm:an-equivalence}), so that results proved about one may be
applied to the other without comment.  The second is to determine exactly how
much of a function the relation remembers (\cref{plt:thm:an-borel-ritt} and
\cref{plt:cor:an-fibres}); the answer is: everything except a flat
perturbation, and nothing less.

Throughout, \(\K\subseteq\ComplexNumbers\) is a field of characteristic zero,
the regime is \(X\to+\infty\) along the real ray \(S=[X_{0},+\infty)\) with
\(X_{0}>1\) and the real logarithm \(L=\log X\), and \(t=X^{-1}\).  Where a
statement holds verbatim on a closed subsector of a sector at infinity with a
fixed branch of the logarithm, this is said explicitly.

\section{What may be called a scale here}
\label{plt:sec:an-scale}

Poincar\'e's definition of an asymptotic expansion presupposes an
\emph{enumerated} scale: a sequence \(\phi_{0},\phi_{1},\ldots\) of positive
functions with \(\phi_{k+1}=\LittleOAt{\phi_{k}}{X\to+\infty}\).  Two of the
sources assert that the polynomial--logarithmic monomials
\(X^{-n}(\log X)^{j}\) form such a scale under the lexicographic dominance
order of \cref{plt:def:ring-dominance}.  They do not, and the reason is
structural rather than pedantic.

\begin{proposition}[The full monomial family admits no decreasing enumeration]
\label{plt:prop:an-no-enumeration}
Let
\[
 \Gamma=\SetOf{X^{-n}(\log X)^{j}\ :\ n\in\IntegerNumbers,\
 j\in\NaturalNumbers},
\]
ordered by dominance \(\precdom\) as in \cref{plt:def:ring-dominance}.  Then
for every \(n\in\IntegerNumbers\) the set
\(\SetOf{\mu\in\Gamma:\mu\precdom X^{-n}}\) has no \(\precdom\)-greatest
element.  Consequently there is no sequence
\((\phi_{k})_{k\in\NaturalNumbers}\) in \(\Gamma\) with
\(\phi_{k+1}\precdom\phi_{k}\) whose terms exhaust \(\Gamma\).
\end{proposition}

\begin{proof}
Fix \(n\).  A monomial \(X^{-m}(\log X)^{j}\) is \(\precdom X^{-n}=X^{-n}(\log
X)^{0}\) if and only if \(m>n\), since for \(m=n\) the only monomial below
\((\log X)^{0}\) in the block would need a negative logarithmic exponent, and
there is none in \(\Gamma\).  Given such a monomial, \(X^{-m}(\log
X)^{j+1}\) is again \(\precdom X^{-n}\) (still \(m>n\)) and strictly dominates
it.  So the set has no greatest element.

For the consequence, suppose \((\phi_{k})\) were a strictly decreasing
enumeration of \(\Gamma\).  The monomial \(X^{-n}\) occurs as some
\(\phi_{k}\), and then \(\phi_{k+1}\precdom X^{-n}\) while every element of
\(\Gamma\) below \(X^{-n}\) must occur later in the sequence, hence be
\(\preceqdom\phi_{k+1}\).  That makes \(\phi_{k+1}\) a greatest element of a
set which has none.
\end{proof}

% ed.: S2 states that "the monomials X^{-n}L^m form an asymptotic scale under
% ed.: the lexicographic order", and S6 opens its analytic chapter by letting
% ed.: (phi_nu) be "a totally ordered asymptotic scale, so that
% ed.: phi_{nu+1} = o(phi_nu) in the chosen enumeration" and then applying
% ed.: this to the polynomial-logarithmic monomials.  The pairwise dominance
% ed.: comparisons those drafts display are all correct; what fails is the
% ed.: existence of the enumeration, by the proposition above.  The repair is
% ed.: to fix a degree bound first, which is what the block formulation does
% ed.: automatically.
\begin{definition}[Admissible enumerated scale]
\label{plt:def:an-scale}
Let \(n_{0}\in\IntegerNumbers\) and let
\(\delta\colon\IntegerNumbers_{\ge n_{0}}\to\NaturalNumbers\) be any function.
The \emph{scale of profile \(\delta\)} is the family
\[
 \Sigma_{\delta}
 =\SetOf{X^{-n}(\log X)^{j}\ :\ n\ge n_{0},\ 0\le j\le\delta(n)},
\]
enumerated by increasing \(n\) and, inside each block, by decreasing \(j\).
\end{definition}

\begin{lemma}[A profile makes a genuine Poincar\'e scale]
\label{plt:lem:an-scale}
The enumeration of \cref{plt:def:an-scale} is strictly decreasing for
dominance, and consecutive members satisfy
\(\phi_{k+1}=\LittleOAt{\phi_{k}}{X\to+\infty}\).
\end{lemma}

\begin{proof}
Each block contributes the finitely many indices \(j=\delta(n),\ldots,0\), so
the enumeration is well defined and of order type \(\omega\).  Inside a block,
\(X^{-n}(\log X)^{j-1}/X^{-n}(\log X)^{j}=(\log X)^{-1}\to0\).  At a block
boundary, the successor of \(X^{-n}\) is \(X^{-n-1}(\log X)^{\delta(n+1)}\)
and the ratio is \((\log X)^{\delta(n+1)}/X\to0\).
\end{proof}

\section{The expansion relation, and the equivalent remainder conditions}
\label{plt:sec:an-definition}

\cref{plt:def:mul-asymptotic} defines \(f\approx F\) for \(F\in\PLlau\).  We
record the extension to rational logarithmic blocks, which is forced on us
later by division (\cref{plt:thm:an-quotient}) and is already needed to say
anything at all about so simple a function as \(1/\log X\)
(\cref{plt:ex:an-no-expansion}).

% ed.: this definition carried the same title as \cref{plt:def:mul-asymptotic},
% which it extends rather than repeats: the blocks here are rational in L and
% may have poles, so the evaluation needs the large-X caveat below.  Retitled so
% the two are distinguishable in the list of results.
\begin{definition}[Polynomial--logarithmic asymptotic expansion over the
  rational-log field]
\label{plt:def:an-expansion}
Let \(F=\sum_{n\ge n_{0}}p_{n}(L)t^{n}\) lie in \(\PLrat\), so that every
block \(p_{n}\) is a rational function of \(L\) over \(\K\).  Fix
\(N\in\IntegerNumbers\).  Only the finitely many blocks with
\(n_{0}\le n<N\) are involved below, and each has finitely many poles, so for
all sufficiently large \(X\) none of them has a pole at \(\log X\); for such
\(X\) put
\[
 \bigl(\Trunc{F}{N}\bigr)(X)=\sum_{n_{0}\le n<N}p_{n}(\log X)\,X^{-n},
\]
a finite sum of complex numbers.  A function \(f\colon S\to\ComplexNumbers\)
is \emph{PL-asymptotic} to \(F\), written \(f\approx F\), if for every
\(N\in\IntegerNumbers\) there are \(M_{N}\in\NaturalNumbers\), \(C_{N}>0\) and
\(X_{N}\ge X_{0}\) with
\begin{equation}\label{plt:eq:an-strong}
 \AbsoluteValue{f(X)-\bigl(\Trunc{F}{N}\bigr)(X)}
 \le C_{N}\,X^{-N}(\log X)^{M_{N}}
 \qquad(X\ge X_{N}).
\end{equation}
For \(F\in\PLlau\) this is exactly \cref{plt:def:mul-asymptotic}.
\end{definition}

\begin{remark}[Why the ambient is not \(\PLit\)]
\label{plt:rmk:an-not-it}
The definition does not extend to \(F\in\PLit\) as it stands: there a block
\(p_{n}\) lies in \(\K((L^{-1}))\), so it is an infinite series and not a
function of \(\log X\), and \(\bigl(\Trunc{F}{N}\bigr)(X)\) is not a
number.  The correct
object there is a \emph{doubly} indexed condition, in which one truncates
first in \(t\) and then, inside each retained block, in \(L^{-1}\); the two
limits do not commute, since \(t\) is smaller than every power of \(L^{-1}\)
(\cref{plt:prop:ring-iteration-order}).  Two of the sources use \(\PLit\)
coefficients and an undecorated one-parameter remainder in the same formula.
We do not develop the iterated notion here and no statement below uses it;
it is recorded as an obligation in \cref{plt:sec:repairs}.
\end{remark}

\begin{theorem}[Equivalent forms of the remainder condition]
\label{plt:thm:an-equivalence}
Let \(F=\sum_{n\ge n_{0}}p_{n}(L)t^{n}\in\PLrat\) and let \(f\colon
S\to\ComplexNumbers\).  Write \(k_{n}\in\IntegerNumbers\) for the degree of
\(p_{n}\) as a rational function, so that \(p_{n}(\ell)\sim c_{n}\ell^{k_{n}}\)
with \(c_{n}\ne0\) when \(p_{n}\ne0\).  The following are equivalent.
\begin{enumerate}
\item[\textup{(i)}] For every \(N\),
\(f-\Trunc{F}{N}=\LittleOAt{X^{1-N}}{X\to+\infty}\).
\item[\textup{(ii)}] For every \(N\) and every \(\varepsilon>0\),
\(f-\Trunc{F}{N}=\BigOAt{X^{-N+\varepsilon}}{X\to+\infty}\).
\item[\textup{(iii)}] \(f\approx F\), that is, \eqref{plt:eq:an-strong} holds
for every \(N\).
\item[\textup{(iv)}] For every \(N\),
\(f-\Trunc{F}{N}-p_{N}(\log X)X^{-N}=\LittleOAt{X^{-N}}{X\to+\infty}\).
\item[\textup{(v)}] There are \(M_{N}\in\NaturalNumbers\) with
\(f-\Trunc{F}{N+1}=\LittleOAt{X^{-N}(\log X)^{M_{N}}}{X\to+\infty}\) for
every \(N\).
\end{enumerate}
If moreover \(F\in\PLlau\), these are further equivalent to:
\begin{enumerate}
\item[\textup{(vi)}] for some \textup{(}equivalently, every\textup{)} profile
\(\delta\) with \(\delta(n)\ge\degL p_{n}\) for all \(n\), the function \(f\)
has the Poincar\'e expansion \(\sum_{k}c_{k}\phi_{k}\) relative to the scale
\(\Sigma_{\delta}\) of \cref{plt:def:an-scale}, where \(c_{k}\) is the
coefficient in \(F\) of the monomial \(\phi_{k}\); that is,
\(f-\sum_{i\le k}c_{i}\phi_{i}=\LittleOAt{\phi_{k}}{X\to+\infty}\) for every
\(k\).
\end{enumerate}
The equivalence is between the \emph{families} of conditions, and cannot be
localised to one cutoff.  Cutoff by cutoff only the implications
\textup{(iv)}\(\Rightarrow\)\textup{(iii)}\(\Rightarrow\)\textup{(ii)}
\(\Rightarrow\)\textup{(i)} hold; the reverse implications need the whole
family.  For instance, with \(f(X)=X^{-N+1/2}\) and \(F=0\), condition
\textup{(i)} holds at that \(N\) while \textup{(iii)} fails at that \(N\) for
every \(M_{N}\).
\end{theorem}

\begin{proof}
Throughout we use that for any fixed \(m\in\IntegerNumbers\), \(d\ge0\) and
\(\varepsilon>0\),
\begin{equation}\label{plt:eq:an-logs-lose}
 X^{-m}(\log X)^{d}=\LittleOAt{X^{-m+\varepsilon}}{X\to+\infty},
 \qquad
 X^{-m-1}(\log X)^{d}=\LittleOAt{X^{-m}}{X\to+\infty},
\end{equation}
both because \((\log X)^{d}/X^{\varepsilon}\to0\).  We also use that a nonzero
\(p\in\K(L)\) satisfies \(p(\log X)=\BigOAt{(\log X)^{\max(k,0)}}{X\to+\infty}\)
with \(k\) its degree, so each individual block of \(F\) obeys
\(p_{n}(\log X)X^{-n}=\BigOAt{X^{-n}(\log X)^{\max(k_{n},0)}}{X\to+\infty}\).

\emph{\textup{(iii)}\(\Rightarrow\)\textup{(ii)}\(\Rightarrow\)\textup{(i)}.}
The first is \eqref{plt:eq:an-logs-lose}; the second follows by taking
\(\varepsilon<1\).

\emph{\textup{(i)}\(\Rightarrow\)\textup{(iv)}.}  Condition (i) at the cutoff
\(N+1\) reads
\(f-\Trunc{F}{N+1}=\LittleOAt{X^{-N}}{X\to+\infty}\), and one has the
identity \(\Trunc{F}{N+1}=\Trunc{F}{N}+p_{N}(\log X)X^{-N}\).

\emph{\textup{(iv)}\(\Rightarrow\)\textup{(iii)}.}  Given \(N\), (iv) at that
\(N\) gives
\(f-\Trunc{F}{N}=p_{N}(\log X)X^{-N}+\LittleOAt{X^{-N}}{X\to+\infty}\), which
is \(\BigOAt{X^{-N}(\log X)^{M_{N}}}{X\to+\infty}\) with
\(M_{N}=\max(k_{N},0)\).

\emph{\textup{(iv)}\(\Rightarrow\)\textup{(v)}} with \(M_{N}=0\), since (iv)
at \(N\) is precisely
\(f-\Trunc{F}{N+1}=\LittleOAt{X^{-N}}{X\to+\infty}\).

\emph{\textup{(v)}\(\Rightarrow\)\textup{(i)}.}  Apply (v) at index \(N+1\):
\(f-\Trunc{F}{N+2}=\LittleOAt{X^{-N-1}(\log X)^{M_{N+1}}}{X\to+\infty}\),
which by \eqref{plt:eq:an-logs-lose} is
\(\LittleOAt{X^{-N}}{X\to+\infty}\).  Adding the single block
\(p_{N+1}(\log X)X^{-N-1}=\LittleOAt{X^{-N}}{X\to+\infty}\) gives
\(f-\Trunc{F}{N+1}=\LittleOAt{X^{-N}}{X\to+\infty}\), which is (i) at cutoff
\(N+1\).  As \(N\) is arbitrary, (i) holds at every cutoff.

\emph{\textup{(iii)}\(\Rightarrow\)\textup{(vi)}, for \(F\in\PLlau\).}  Fix
an admissible profile \(\delta\), and let \(\phi_{k}=X^{-n}(\log X)^{j}\).
The partial sum \(\sum_{i\le k}c_{i}\phi_{i}\) equals
\(\Trunc{F}{n}(X)+X^{-n}\sum_{i\ge j}a_{n,i}(\log X)^{i}\), where
\(p_{n}=\sum_{i}a_{n,i}L^{i}\).  Hence
\[
 f-\sum_{i\le k}c_{i}\phi_{i}
 =\Bigl(f-\Trunc{F}{n+1}\Bigr)+X^{-n}\sum_{i<j}a_{n,i}(\log X)^{i}.
\]
By (iii) at cutoff \(n+1\) and \eqref{plt:eq:an-logs-lose}, the first bracket
is \(\LittleOAt{X^{-n}}{X\to+\infty}\), hence
\(\LittleOAt{\phi_{k}}{X\to+\infty}\) because \(\phi_{k}\ge X^{-n}\); the
second is \(\BigOAt{X^{-n}(\log X)^{j-1}}{X\to+\infty}
=\LittleOAt{\phi_{k}}{X\to+\infty}\) (empty, hence zero, if \(j=0\)).

\emph{\textup{(vi)}\(\Rightarrow\)\textup{(i)}.}  Take \(\phi_{k}=X^{-n}\),
the last member of block \(n\); the corresponding partial sum is exactly
\(\Trunc{F}{n+1}(X)\), so (vi) gives
\(f-\Trunc{F}{n+1}=\LittleOAt{X^{-n}}{X\to+\infty}\) for every \(n\ge
n_{0}\), which is (i) at every cutoff \(>n_{0}\); at cutoffs
\(\le n_{0}\) the truncation is \(0\) and the estimate follows from the one at
\(n_{0}+1\).  That the choice of admissible \(\delta\) is immaterial is now
clear, since the conclusion (i) does not mention \(\delta\).

Finally, the last assertion: for \(f=X^{-N+1/2}\) and \(F=0\) one has
\(f=\LittleOAt{X^{1-N}}{X\to+\infty}\), while
\(X^{-N+1/2}/(X^{-N}(\log X)^{M})=X^{1/2}(\log X)^{-M}\to+\infty\) for every
\(M\).
\end{proof}

\begin{remark}[Concordance with the six sources and with this volume]
\label{plt:rmk:an-conventions}
\Cref{plt:thm:an-equivalence} reconciles every convention in play.  Form
(iii) is \cref{plt:def:mul-asymptotic}; form (i) is
\cref{plt:def:cmp-realization}; form (iv), written with the inclusive cutoff
\(\sum_{n\le N}\), is the definition used by S4 and S5; form (v), again
inclusive and with a logarithmic weight on the right, is the definition used
by S6.  All of them say the same thing, and consequently
\cref{plt:lem:cmp-realization-basic,plt:lem:cmp-realization-ring,plt:thm:cmp-analytic-transfer}
may be used interchangeably with
\cref{plt:lem:mul-asymptotic-basic,plt:thm:mul-analytic-product}.  What is
\emph{not} interchangeable is a single truncation order: the last clause of
the theorem shows that a one-cutoff statement in one convention is strictly
weaker or stronger than the corresponding one-cutoff statement in another,
which is why a printed asymptotic formula must always say which convention
and which cutoff it uses.
\end{remark}

% ed.: S6 writes the blockwise condition as
% ed.:   f - sum_{n<=N} p_n(log X) X^{-n} = o(X^{-N}(log X)^{d_N})
% ed.: and remarks that "one must specify that inner ordering" of the
% ed.: monomials inside the last block.  Form (v) of the theorem shows the
% ed.: logarithmic weight is harmless once all N are quantified, and form
% ed.: (vi) shows the inner ordering never has to be specified: the blockwise
% ed.: and monomialwise conditions coincide for every admissible profile.
\begin{proposition}[Uniqueness, with rational blocks]
\label{plt:prop:an-uniqueness}
A function \(f\colon S\to\ComplexNumbers\) is PL-asymptotic to at most one
\(F\in\PLrat\).  The same holds on a closed subsector
\(\SetOf{z:\AbsoluteValue{z}\ge R,\ \AbsoluteValue{\arg z-\theta_{0}}\le\theta}\)
of a sector at infinity carrying a fixed analytic branch of the logarithm.
\end{proposition}

\begin{proof}
Suppose \(f\approx F\) and \(f\approx\widetilde F\) with
\(D=F-\widetilde F\ne0\).  Put \(n_{1}=\ordt D\) and let \(q\in\K(L)\) be its
leading block, of degree \(k\in\IntegerNumbers\) and leading coefficient
\(c\ne0\), so \(q(\ell)/\ell^{k}\to c\) as \(\AbsoluteValue{\ell}\to\infty\).
Subtracting the two instances of \eqref{plt:eq:an-strong} at
\(N=n_{1}+1\) and using \(\Trunc{D}{n_{1}+1}=q(L)t^{n_{1}}\) gives
\[
 q(\log X)\,X^{-n_{1}}
 =\BigOAt{X^{-n_{1}-1}(\log X)^{M}}{X\to+\infty}
\]
for some \(M\in\NaturalNumbers\).  Dividing,
\(\AbsoluteValue{q(\log X)}\,X\,(\log X)^{-M}\) would have to stay bounded;
but it is asymptotic to \(\AbsoluteValue{c}X(\log X)^{k-M}\to+\infty\), a
contradiction.  On a subsector, \(\log z=\log\AbsoluteValue{z}+
\ImaginaryUnit\arg z\) has \(\AbsoluteValue{\arg z}\) bounded, so
\(\AbsoluteValue{\log z}\sim\log\AbsoluteValue{z}\) and the same computation
applies with \(\AbsoluteValue{z}\) in place of \(X\).
\end{proof}

\begin{remark}[Exactly what uniqueness gives, and what it does not]
\label{plt:rmk:an-uniqueness-scope}
\Cref{plt:prop:an-uniqueness} says the map \(f\mapsto F\) is well defined
where it is defined.  It does \emph{not} say:
\begin{enumerate}
\item that \(f\) is determined by \(F\).  It is not, by
\cref{plt:cor:an-fibres}; the fibre over \(F\) is an entire coset of the flat
ideal.
\item that \(F\) determines \(f\) even up to a null set, or on any finite
interval.  The relation is a statement about germs at \(+\infty\) only.
\item that \(F\) converges anywhere.  Convergence is a separate theorem about
one series and one domain; see \cref{plt:prop:om-frozen-convergence} and
\cref{plt:rmk:om-open} for the only case in this volume where the question is
settled at all, and there only in a form that does not answer the asymptotic
question.
\item that \(f\) is differentiable, or that \(f'\approx\DX F\) if it is
(\cref{plt:ex:an-diff-fails}).
\item that \(f\) exists.  That every \(F\) is realized is a theorem
(\cref{plt:thm:an-borel-ritt}), and its proof constructs \(f\) by hand; no
part of the formal algebra supplies it.
\item anything about the coefficients being computable from finitely many
values of \(f\).  Each \(p_{n}\) is a limit of quotients along the whole ray.
\end{enumerate}
\end{remark}

\begin{example}[A smooth positive function with no \(\PLlau\) expansion]
\label{plt:ex:an-no-expansion}
Let \(f(X)=1/\log X\) on \([\EulerE,\infty)\).  If \(f\approx F\) with
\(F\in\PLlau\), then \(\log X\cdot f(X)=1\) and \(\log X\approx L\), so
\cref{plt:thm:mul-analytic-product} gives \(LF=1\) in \(\PLlau\).  But \(L\)
is not a unit of \(\PLlau\) by \cref{plt:thm:div-unit-criterion}, since its
leading block \(L\) is not a unit of \(\K[L]\); contradiction.  On the other
hand \(f\approx L^{-1}\in\PLrat\) trivially, the remainder being zero at every
cutoff \(N\ge1\).  So the enlargement of \cref{plt:def:an-expansion} from
\(\PLlau\) to \(\PLrat\) is forced by an example this elementary, and is not a
convenience.
\end{example}

\section{The algebra of expansions and the flat ideal}
\label{plt:sec:an-algebra}

\begin{definition}[The expansion algebra and its flat ideal]
\label{plt:def:an-algebra}
Let \(\mathcal A\) be the set of germs at \(+\infty\) of functions
\(f\colon[X_{f},\infty)\to\ComplexNumbers\) for which there exists
\(F\in\PLlau\) with \(f\approx F\), and let
\[
 \mathcal N=\SetOf{f\in\mathcal A:f\approx0}
 =\SetOf{f:\ f=\BigOAt{X^{-N}}{X\to+\infty}\ \text{for every }N}.
\]
Members of \(\mathcal N\) are called \emph{flat}.  Write
\(\mathsf{exp}\colon\mathcal A\to\PLlau\), \(\mathsf{exp}(f)=F\), for the
expansion map, well defined by \cref{plt:prop:an-uniqueness}.
\end{definition}

The second description of \(\mathcal N\) is
\cref{plt:lem:mul-asymptotic-basic}(4), and that \(\mathcal N\) is an ideal of
the \(\K\)-algebra \(\mathcal A\) is part of
\cref{plt:thm:mul-analytic-product}: \(\mathcal N=\ker\mathsf{exp}\) and
\(\mathsf{exp}\) is a \(\K\)-algebra homomorphism.  The next proposition
records the three finer structural facts, which are not in the sources.

\begin{proposition}[Structure of the flat ideal]
\label{plt:prop:an-flat-ideal}
\(\mathcal N\) is a proper ideal of \(\mathcal A\), and
\begin{enumerate}
\item \(\mathcal N\) is prime but not maximal;
\item \(\mathcal N\) is not stable under differentiation, even on the
subalgebra of germs of entire functions;
\item \(\mathcal N\) contains no minimal nonzero scale: if \(f\in\mathcal N\)
is eventually nonvanishing then \(f^{2}\in\mathcal N\) and
\(f^{2}=\LittleOAt{f}{X\to+\infty}\).
\end{enumerate}
\end{proposition}

\begin{proof}
Properness: \(1\approx1\ne0\), so \(1\notin\mathcal N\).

(1)  By \cref{plt:thm:an-borel-ritt} below, \(\mathsf{exp}\) is onto, so
\(\mathcal A/\mathcal N\cong\PLlau\).  That ring is an integral domain
(\cref{plt:cor:mul-domain}) and is not a field
(\cref{plt:cor:div-domain-not-field}); hence \(\mathcal N\) is prime and not
maximal.  A witness for non-maximality is the germ of \(\log X\), whose class
is nonzero and not invertible: an inverse would force \(L\) to be a unit of
\(\PLlau\), refuted in \cref{plt:ex:an-no-expansion}.

(2)  See \cref{plt:ex:an-diff-fails}, where \(f\in\mathcal N\) is entire and
\(f'\notin\mathcal A\), a fortiori \(f'\notin\mathcal N\).

(3)  \(\mathcal N\) is an ideal and \(f\in\mathcal N\subseteq\mathcal A\), so
\(f^{2}\in\mathcal N\).  Taking \(N=1\) in the definition of \(\mathcal N\)
gives \(f(X)\to0\), so \(f^{2}/f=f\to0\) wherever \(f\ne0\), that is
\(f^{2}=\LittleOAt{f}{X\to+\infty}\).
\end{proof}

\begin{theorem}[Borel--Ritt on the polynomial--logarithmic scale]
\label{plt:thm:an-borel-ritt}
For every \(F\in\PLlau\) there is a \(C^{\infty}\) function
\(f\colon(0,\infty)\to\ComplexNumbers\), real valued if
\(\K\subseteq\RealNumbers\), with \(f\approx F\).  Consequently
\[
 0\longrightarrow\mathcal N\longrightarrow\mathcal A
 \stackrel{\mathsf{exp}}{\longrightarrow}\PLlau\longrightarrow0
\]
is an exact sequence of \(\K\)-algebras and \(\mathcal A/\mathcal N\cong\PLlau\).
\end{theorem}

\begin{proof}
Write \(F=\sum_{n\ge n_{0}}p_{n}(L)t^{n}\).  Fix once and for all a
\(C^{\infty}\) function \(\chi\colon\RealNumbers\to[0,1]\) with \(\chi=0\) on
\((-\infty,0]\) and \(\chi=1\) on \([1,\infty)\).  Choose real numbers
\(X_{n_{0}}<X_{n_{0}+1}<\cdots\) with \(X_{n}\ge\max(\EulerE,X_{n-1}+1)\),
\(X_{n}\to+\infty\), and
\begin{equation}\label{plt:eq:an-BR-choice}
 \AbsoluteValue{p_{n}(\log X)}\le 2^{-\AbsoluteValue{n}}X
 \qquad\text{for all }X\ge X_{n};
\end{equation}
this is possible because \(p_{n}(\log X)/X\to0\) for each fixed \(n\).  Set
\[
 f(X)=\sum_{n\ge n_{0}}\chi(X-X_{n})\,p_{n}(\log X)\,X^{-n}.
\]
On any bounded interval only the finitely many \(n\) with \(X_{n}<\sup\) of
that interval contribute, so the sum is locally finite and \(f\in
C^{\infty}((0,\infty))\).

By \cref{plt:lem:mul-asymptotic-basic}(1) it suffices to establish
\eqref{plt:eq:an-strong} for every \(N>n_{0}\).  Fix such an \(N\) and split
\[
 f-\Trunc{F}{N}
 =\underbrace{\sum_{n_{0}\le n<N}\bigl(\chi(X-X_{n})-1\bigr)p_{n}(\log
 X)X^{-n}}_{=:\,\Sigma_{1}}
 +\underbrace{\sum_{n\ge N}\chi(X-X_{n})p_{n}(\log X)X^{-n}}_{=:\,\Sigma_{2}}.
\]
\(\Sigma_{1}\) is a finite sum whose \(n\)-th term vanishes once
\(X\ge X_{n}+1\); hence \(\Sigma_{1}=0\) for \(X\ge X_{N-1}+1\).  In
\(\Sigma_{2}\), the term \(n=N\) is bounded by
\(\AbsoluteValue{p_{N}(\log X)}X^{-N}
\le C X^{-N}(\log X)^{M_{N}}\) with \(M_{N}=\max(\degL p_{N},0)\).  For
\(n\ge N+1\) the \(n\)-th term vanishes unless \(X\ge X_{n}\), and where it
does not vanish \eqref{plt:eq:an-BR-choice} and \(\chi\le1\) bound it by
\(2^{-\AbsoluteValue{n}}X^{1-n}\le2^{-\AbsoluteValue{n}}X^{-N}\) for
\(X\ge1\); summing \(\sum_{n\in\IntegerNumbers}2^{-\AbsoluteValue{n}}=3\),
\[
 \AbsoluteValue{\Sigma_{2}}\ \le\ CX^{-N}(\log X)^{M_{N}}+3X^{-N} .
\]
This is \eqref{plt:eq:an-strong} at the cutoff \(N\).

Exactness at \(\PLlau\) is the surjectivity just proved, exactness at
\(\mathcal A\) is \(\mathcal N=\ker\mathsf{exp}\)
(\cref{plt:thm:mul-analytic-product}), and \(\mathsf{exp}\) is a
\(\K\)-algebra map by the same theorem.
\end{proof}

\begin{corollary}[The converse fails, and by exactly one coset]
\label{plt:cor:an-fibres}
For \(F\in\PLlau\) the set of \(f\in\mathcal A\) with \(f\approx F\) is
nonempty, and equals \(f_{0}+\mathcal N\) for any one of its members
\(f_{0}\).  The
ideal \(\mathcal N\) is nonzero, since it contains \(\EulerE^{-X}\),
\(\EulerE^{-\sqrt X}\) and \(X^{-\log X}\); so infinitely many pairwise
distinct germs share every polynomial--logarithmic expansion, and no property
of a function that is not invariant under adding a flat germ can be read off
from its expansion.
\end{corollary}

\begin{proof}
Nonemptiness is \cref{plt:thm:an-borel-ritt}; the fibre description is the
statement that \(\mathsf{exp}\) is additive with kernel \(\mathcal N\).  For
membership in \(\mathcal N\): \(\EulerE^{-X}X^{N}\to0\),
\(\EulerE^{-\sqrt X}X^{N}\to0\) and \(X^{-\log X}X^{N}
=\EulerE^{-(\log X)^{2}+N\log X}\to0\) for every \(N\).
\end{proof}

\begin{example}[Flatness has no bottom]
\label{plt:ex:an-flat-tower}
Extending the dominance symbol of \cref{plt:def:ring-dominance} from
monomials to arbitrary positive germs by declaring
\(\phi\succdom\psi\) to mean \(\psi=\LittleOAt{\phi}{X\to+\infty}\), the germs
\[
 X^{-\log X}\ \succdom\ \EulerE^{-\sqrt X}\ \succdom\ \EulerE^{-X}
 \ \succdom\ \EulerE^{-\EulerE^{X}}
\]
all lie in \(\mathcal N\); iterating
\cref{plt:prop:an-flat-ideal}(3) produces from any one of them an infinite
strictly decreasing chain inside \(\mathcal N\).  The expansion map cannot
distinguish any of them from \(0\).  This is the precise sense in which a
polynomial--logarithmic expansion ``remembers the whole power--log sector and
nothing beyond it''; adjoining exponential monomials is the standard remedy
and takes one out of the class studied here
(\cref{plt:thm:dif-exp-boundary}).
\end{example}

\begin{theorem}[Quotients]
\label{plt:thm:an-quotient}
Let \(f\approx F\) and \(g\approx G\) with \(F,G\in\PLrat\) and \(G\ne0\).
Put \(m=\ordt G\) and let \(q\in\K(L)\) be the leading block of \(G\).  Then
\(g\) is eventually nonvanishing, and
\[
 \frac{f}{g}\approx\frac{F}{G}\ \in\PLrat .
\]
If \(F,G\in\PLlau\) and \(q\in\K^{\times}\), then \(F/G\in\PLlau\); if
\(q\notin\K\), then \(F/G\) need not lie in \(\PLlau\), the quotient
\(F=1\), \(G=L\) being the smallest example
(\cref{plt:ex:an-no-expansion}).
\end{theorem}

\begin{proof}
Write \(k\) for the degree of \(q\) and \(c\ne0\) for its leading
coefficient.  Applying \eqref{plt:eq:an-strong} to \(g\) at cutoff \(m+1\),
\[
 g(X)=X^{-m}\Bigl(q(\log X)+\BigOAt{X^{-1}(\log
 X)^{M}}{X\to+\infty}\Bigr),
\]
and \(\AbsoluteValue{q(\log X)}\sim\AbsoluteValue{c}(\log X)^{k}\) while
\(X^{-1}(\log X)^{M}=\LittleOAt{(\log X)^{k}}{X\to+\infty}\).  Hence for large
\(X\)
\begin{equation}\label{plt:eq:an-g-lower}
 \AbsoluteValue{g(X)}\ \ge\ \tfrac12\AbsoluteValue{c}\,X^{-m}(\log X)^{k}>0 ,
\end{equation}
so \(g\) is eventually nonvanishing and
\(\AbsoluteValue{1/g(X)}\le 2\AbsoluteValue{c}^{-1}X^{m}(\log X)^{-k}\).

Since \(\PLrat=\K(L)((t))\) is the field of Laurent series over the field
\(\K(L)\) (\cref{plt:thm:div-Ett-field}), \(H\DefinitionEquals F/G\in\PLrat\)
exists and \(F=GH\).  Fix \(N\).  The truncation \(\Trunc{H}{N}\) is a finite
polynomial--logarithmic expression, hence realizes itself:
\(\Trunc{H}{N}\approx\Trunc{H}{N}\), because for any cutoff \(N'\) the
difference \(\Trunc{H}{N}-\Trunc{\Trunc{H}{N}}{N'}\) is a finite sum of blocks
of index \(\ge N'\).  By \cref{plt:thm:mul-analytic-product} and linearity
(\cref{plt:lem:mul-asymptotic-basic}(3)),
\[
 \phi\DefinitionEquals f-g\cdot\Trunc{H}{N}
 \ \approx\ E\DefinitionEquals F-G\,\Trunc{H}{N}
 =G\bigl(H-\Trunc{H}{N}\bigr).
\]
By additivity of \(\ordt\) (\cref{plt:thm:ring-valuation-laws}),
\(\ordt E\ge m+N\), so \(\Trunc{E}{m+N}=0\) and \eqref{plt:eq:an-strong}
applied to \(\phi\) at cutoff \(m+N\) gives
\(\phi=\BigOAt{X^{-(m+N)}(\log X)^{M'}}{X\to+\infty}\).  Dividing by \(g\) and
using the lower bound \eqref{plt:eq:an-g-lower},
\[
 \frac{f}{g}-\Trunc{H}{N}
 =\frac{\phi}{g}
 =\BigOAt{X^{-N}(\log X)^{\max(M'-k,\,0)}}{X\to+\infty}.
\]
As \(N\) was arbitrary, \(f/g\approx H\).  For the last sentence: \(F/G\)
lies in \(\PLlau\) for every \(F\in\PLlau\) if and only if \(t^{-m}G\) is a
unit of \(\PLpow\), which by \cref{plt:thm:div-unit-criterion} happens if and
only if its leading block \(q\) lies in \(\K^{\times}\), and then \(1/G\) is
given by the recurrence of \cref{plt:thm:div-reciprocal}.
\end{proof}

% ed.: S5 states products and quotients together with a "proof sketch" and, in
% ed.: its part (ii), adds the hypothesis "and g is eventually nonzero".  That
% ed.: hypothesis is redundant: it is a consequence of the leading block being
% ed.: nonzero, as (plt:eq:an-g-lower) shows.  S5's part (iii) requires
% ed.: instead that "Q(log x) stays nonzero eventually", which is likewise
% ed.: automatic for a nonzero rational function.  The full proof above
% ed.: replaces the sketch, and in particular exhibits the logarithmic degree
% ed.: k of the leading block explicitly, which the sketch suppresses; it is
% ed.: what makes the remainder exponent M'-k, not M'.
\begin{proposition}[Substitution into a function holomorphic at the origin]
\label{plt:prop:an-holomorphic-substitution}
Let \(u\colon S\to\ComplexNumbers\) satisfy \(u\approx U\) with
\(U\in\PLrat\), \(\ordt U\ge1\).  Let \(h(z)=\sum_{k\ge0}h_{k}z^{k}\) be
holomorphic on \(\AbsoluteValue{z}<r\) with coefficients in \(\K\).  Then
\(u(X)\to0\), \(h\compose u\) is defined for large \(X\), and
\[
 h\compose u\ \approx\ \sum_{k\ge0}h_{k}U^{k},
\]
% ed.: the right-hand side was attributed to \cref{plt:prop:tay-scalar-outer},
% ed.: which is stated for arguments in t\PLpow.  Here U has rational blocks, so
% ed.: the ambient is \(\K(L)[[t]]\); that case is covered by the last sentence
% ed.: of \cref{plt:prop:dif-substitution}, which is the citation used now.
the right side being the formally summable substitution of
\cref{plt:prop:dif-substitution}, applied in \(\K(L)[[t]]\).
\end{proposition}

\begin{proof}
Since \(\ordt U\ge1\), \(\Trunc{U}{1}=0\), so \eqref{plt:eq:an-strong} at
cutoff \(1\) gives \(u=\BigOAt{X^{-1}(\log X)^{M_{1}}}{X\to+\infty}\), whence
\(u\to0\) and, for each \(k\ge0\),
\(u^{k}=\BigOAt{X^{-k}(\log X)^{kM_{1}}}{X\to+\infty}\).  In particular
\(\AbsoluteValue{u}\le r/2\) for large \(X\), so \(h\compose u\) is defined
there.

Fix \(N\) and choose an integer \(J\ge\max(N,0)\).  Write
\(h(z)=\sum_{k<J}h_{k}z^{k}+z^{J}\widetilde h(z)\) with \(\widetilde h\)
holomorphic on \(\AbsoluteValue{z}<r\); let
\(C_{J}=\max_{\AbsoluteValue{z}\le r/2}\AbsoluteValue{\widetilde h(z)}\).
Then
\[
 h(u)-\sum_{k<J}h_{k}u^{k}
 =u^{J}\widetilde h(u)
 =\BigOAt{X^{-J}(\log X)^{JM_{1}}}{X\to+\infty}
 =\BigOAt{X^{-N}(\log X)^{JM_{1}}}{X\to+\infty}.
\]
By \cref{plt:thm:mul-analytic-product} and linearity,
\(\sum_{k<J}h_{k}u^{k}\approx\sum_{k<J}h_{k}U^{k}\).  Finally
\(\ordt(U^{k})\ge k\ge J\ge N\) for \(k\ge J\), so
\(\Trunc{\sum_{k<J}h_{k}U^{k}}{N}=\Trunc{\sum_{k\ge0}h_{k}U^{k}}{N}\).
Combining the last three displays gives \eqref{plt:eq:an-strong} for
\(h\compose u\) at cutoff \(N\).
\end{proof}

\begin{corollary}[The three small-argument operations]
\label{plt:cor:an-log-power-exp}
Let \(u\approx U\) with \(\ordt U\ge1\), and let \(\alpha\in\K\).  Then, with
the real logarithm and the principal branch of the power for large real
\(X\),
\[
 \log(1+u)\approx\log(1+U),\qquad
 (1+u)^{\alpha}\approx(1+U)^{\alpha},\qquad
 \EulerE^{u}\approx\exp U,
\]
% ed.: the right-hand sides were identified outright with the formal operations
% ed.: of \cref{plt:lem:cmp-exp-log} and \cref{plt:thm:mul-binomial}, both of
% ed.: which are stated over \PLpow; for a genuinely rational \(U\) the formal
% ed.: object is the scalar substitution of \cref{plt:prop:dif-substitution} in
% ed.: \(\K(L)[[t]]\), and the identification with the named operations is
% ed.: asserted only where they are defined.  Restricted accordingly.
where the right-hand sides are the scalar substitutions of
\cref{plt:prop:dif-substitution} in \(\K(L)[[t]]\); when \(U\in\PLpow\) they
are the formal operations of \cref{plt:lem:cmp-exp-log} and
\cref{plt:thm:mul-binomial}.  In particular the
formal geometric, logarithmic and binomial series may be substituted
termwise, with no hypothesis beyond \(\ordt U\ge1\).
\end{corollary}

\begin{proof}
Apply \cref{plt:prop:an-holomorphic-substitution} to \(h(z)=\log(1+z)\) and
\(h(z)=(1+z)^{\alpha}\), both holomorphic on \(\AbsoluteValue{z}<1\), and to
\(h(z)=\EulerE^{z}\), entire; \(u\to0\) puts the argument in the disc for
large \(X\), and on the positive real ray the principal branches agree with
the power series.  That the formal sums are the substitutions named, and that
for \(U\in\PLpow\) they coincide with the operations of
\cref{plt:lem:cmp-exp-log} and \cref{plt:thm:mul-binomial}, is
\cref{plt:prop:dif-substitution} together with the uniqueness of the
continuous \(\K\)-algebra homomorphism sending \(z\) to \(U\).
\end{proof}

\begin{remark}[Composition: the hypotheses are geometric, not algebraic]
\label{plt:rmk:an-composition}
Composition with a general inner function is \emph{not} covered by
\cref{plt:prop:an-holomorphic-substitution} and is not free.  The statement to
use is \cref{plt:thm:cmp-analytic-transfer}, whose three hypotheses (A1)
regime, (A2) monomial-leading realization with a positive exponent, and (A3)
branch compatibility are each shown to be necessary in
\cref{plt:ex:cmp-analytic-failures}.  All five sources that discuss
composition list ``uniformity of the outer remainder over the range of the
inner map'' as a requirement; that is (A1) together with the remainder
transport step \eqref{plt:eq:cmp-remainder-transport}, and it costs
\(\rho^{-1}\) outer orders per output order.  None of the sources records that
cost.
\end{remark}

\section{Termwise differentiation is not automatic}
\label{plt:sec:an-differentiation}

\begin{example}[An entire flat function whose derivative has no expansion]
\label{plt:ex:an-diff-fails}
Let
\[
 f(X)=\EulerE^{-X}\sin\bigl(\EulerE^{X}\bigr),
 \qquad X\ge1 .
\]
Then \(f\) is the restriction to the ray of an entire function,
\(\AbsoluteValue{f(X)}\le\EulerE^{-X}\), so \(f\in\mathcal N\) and
\(f\approx0\); the formal derivative of \(0\) is \(0\).  But
\[
 f'(X)=\cos\bigl(\EulerE^{X}\bigr)-\EulerE^{-X}\sin\bigl(\EulerE^{X}\bigr),
\]
and \(f'\) is not PL-asymptotic to \emph{any} \(F\in\PLrat\).  Indeed
\(\AbsoluteValue{f'}\le1+\EulerE^{-1}\) is bounded, so if \(f'\approx F\) then
\(\ordt F\ge0\) --- otherwise \eqref{plt:eq:an-strong} at cutoff
\(\ordt F+1\) would force
\(\AbsoluteValue{f'}\sim\AbsoluteValue{p_{\ordt F}(\log X)}X^{-\ordt
F}\to\infty\) --- and then \eqref{plt:eq:an-strong} at cutoff \(1\) gives
\(f'(X)-p_{0}(\log X)\to0\).  Since
\(f'(X)=\cos(\EulerE^{X})+\LittleOAt{1}{X\to+\infty}\), this forces
\(p_{0}(\log X)\to1\) along \(X_{k}=\log(2k\pi)\) and \(p_{0}(\log X)\to-1\)
along \(Y_{k}=\log((2k+1)\pi)\).  Both sequences tend to \(+\infty\), and a
rational function has a single limit in
\(\K\cup\SetOf{\infty}\) at \(+\infty\); contradiction.
\end{example}

\begin{warningbox}[title={An expansion says nothing about the derivative}]
The example is entire, bounded, and flat.  Real-analyticity is therefore not a
sufficient hypothesis for termwise differentiation, nor is boundedness, nor is
flatness.  In particular the frequently seen inference
\begin{quote}
``\(f(X)=X+\BigOAt{\log X}{X\to+\infty}\), hence
\(f'(X)=1+\LittleOAt{1}{X\to+\infty}\)''
\end{quote}
is invalid: \(f(X)=X+\sin(\EulerE^{X})\) satisfies the hypothesis and has
\(f'(X)=1+\EulerE^{X}\cos(\EulerE^{X})\), which is unbounded.  Two
sufficient hypotheses that do work are given in
\cref{plt:thm:an-diff-rigid} and \cref{plt:thm:an-sectorial-diff}.
\end{warningbox}

% ed.: S6 writes, in the chapter that this part consolidates, "For
% ed.: f(X) = X + O(log X), one has f'(X) = 1 + o(1) under the relevant
% ed.: differentiability hypotheses".  No such hypotheses are named, and the
% ed.: implication is false without them, by the display in the box above.
% ed.: S6 uses the conclusion to justify choosing m close to 1 in its
% ed.: residual-transfer theorem.  The repair is (plt:prop:an-omega-computable)
% ed.: and (plt:thm:an-inverse-transfer), which take the derivative bound as a
% ed.: hypothesis and verify it directly for the map at hand instead of
% ed.: deducing it from an expansion.
\section{Termwise integration is unconditional}
\label{plt:sec:an-integration}

The asymmetry between the two operations is complete: differentiation needs an
extra hypothesis and integration needs none.  The reason is visible in
\cref{plt:ex:an-diff-fails}: differentiation amplifies fast oscillation, while
integration averages it away.

\begin{lemma}[The elementary tail integral]
\label{plt:lem:an-tail-integral}
Let \(N\ge2\) and \(M\in\NaturalNumbers\).  For \(X\ge\EulerE\),
\[
 \int_{X}^{\infty}s^{-N}(\log s)^{M}\,\Differential s
 \ \le\ C_{N,M}\,X^{1-N}(\log X)^{M},
 \qquad
 C_{N,M}=\frac{1}{N-1}\sum_{i=0}^{M}\frac{M!}{(M-i)!\,(N-1)^{i}} .
\]
\end{lemma}

\begin{proof}
Let \(I_{M}(X)\) denote the integral.  Integrating by parts with
\(u=(\log s)^{M}\) and \(\Differential v=s^{-N}\Differential s\), and using
\(s^{1-N}(\log s)^{M}\to0\),
\[
 I_{M}(X)=\frac{X^{1-N}(\log X)^{M}}{N-1}+\frac{M}{N-1}I_{M-1}(X),
 \qquad
 I_{0}(X)=\frac{X^{1-N}}{N-1}.
\]
Unwinding the recursion gives
\(I_{M}(X)=\frac{X^{1-N}}{N-1}\sum_{i=0}^{M}\frac{M!}{(M-i)!(N-1)^{i}}
(\log X)^{M-i}\), and \((\log X)^{M-i}\le(\log X)^{M}\) for \(X\ge\EulerE\).
\end{proof}

\begin{theorem}[Termwise integration]
\label{plt:thm:an-integration}
Let \(g\colon[X_{0},\infty)\to\ComplexNumbers\) be locally integrable with
\(g\approx G\in\PLlau\), and let \(A\in\PLlau\) be any antiderivative of
\(G\), that is \(\DX A=G\); such an \(A\) exists and is unique up to an
additive constant of \(\K\) by \cref{plt:thm:dif-antiderivative} and
\cref{plt:cor:dif-kernel}.  Then there is a constant \(c\in\ComplexNumbers\)
with
\[
 \Bigl(\int_{X_{0}}^{X}g(s)\,\Differential s\Bigr)-c\ \approx\ A ,
\]
equivalently \(\int_{X_{0}}^{X}g\approx A+c\) once the coefficient field is
enlarged from \(\K\) to \(\K(c)\).  That constant is the only ambiguity: no
hypothesis beyond \(g\approx G\) is needed, and none of \(g\) beyond local
integrability.
\end{theorem}

\begin{proof}
Write \(G=\sum_{n}g_{n}(L)t^{n}\) and \(A=\sum_{n}a_{n}(L)t^{n}\), and put
\(h(X)=\int_{X_{0}}^{X}g\).  We first record the exact formal identity
\begin{equation}\label{plt:eq:an-trunc-commutes}
 \frac{\Differential}{\Differential X}\Bigl[\bigl(\Trunc{A}{N-1}\bigr)(X)\Bigr]
 =\bigl(\Trunc{\DX A}{N}\bigr)(X)
 =\bigl(\Trunc{G}{N}\bigr)(X),
\end{equation}
valid for every \(N\).  Indeed by \cref{plt:prop:dif-block} the block of
\(\DX A\) at index \(m\) is \((\partial_{L}-(m-1))a_{m-1}\), so the blocks of
\(\DX A\) of index \(<N\) are exactly the images of the blocks of \(A\) of
index \(<N-1\); differentiating the finite sum
\(\sum_{n<N-1}a_{n}(\log X)X^{-n}\) term by term reproduces them.

Fix \(N\ge2\) and set \(\psi_{N}=h-\Trunc{A}{N-1}\).  By
\eqref{plt:eq:an-trunc-commutes}, \(\psi_{N}'=g-\Trunc{G}{N}=:r_{N}\), and by
hypothesis \(\AbsoluteValue{r_{N}(X)}\le C_{N}X^{-N}(\log X)^{M_{N}}\) for
\(X\ge X_{N}\).  By \cref{plt:lem:an-tail-integral} this is integrable on
\([X_{N},\infty)\), so
\(c_{N}\DefinitionEquals\lim_{X\to\infty}\psi_{N}(X)\) exists and
\begin{equation}\label{plt:eq:an-int-remainder}
 \AbsoluteValue{\psi_{N}(X)-c_{N}}
 =\Bigl|\int_{X}^{\infty}r_{N}\Bigr|
 \le C_{N}C_{N,M_{N}}\,X^{1-N}(\log X)^{M_{N}} .
\end{equation}
Moreover for \(N\ge2\),
\[
 c_{N+1}-c_{N}
 =\lim_{X\to\infty}\bigl(\Trunc{A}{N-1}-\Trunc{A}{N}\bigr)(X)
 =-\lim_{X\to\infty}a_{N-1}(\log X)\,X^{1-N}=0,
\]
since \(N-1\ge1\).  So all \(c_{N}\) with \(N\ge2\) coincide; call the common
value \(c\).

Then \eqref{plt:eq:an-int-remainder} says
\(h-c-\Trunc{A}{N-1}=\BigOAt{X^{-(N-1)}(\log X)^{M_{N}}}{X\to+\infty}\) for
every \(N\ge2\), that is, \eqref{plt:eq:an-strong} holds for \(h-c\) against
\(A\) at every cutoff \(\ge1\).  By
\cref{plt:lem:mul-asymptotic-basic}(1) a cofinal family of cutoffs suffices,
so \(h-c\approx A\).
\end{proof}

% ed.: the hypothesis "locally integrable" was missing here although the proof
% ed.: invokes \cref{plt:thm:an-integration}, which requires it; without it the
% ed.: integrals below need not exist.  Restored.
\begin{corollary}[The tail integral of an integrable expansion]
\label{plt:cor:an-tail}
If \(g\colon[X_{0},\infty)\to\ComplexNumbers\) is locally integrable and
\(g\approx G\) with \(\ordt G\ge2\), then \(g\) is integrable near
\(+\infty\) and
\[
 \int_{X}^{\infty}g(s)\,\Differential s\ \approx\ -A,
\]
where \(A\) is the unique antiderivative of \(G\) in \(\PLlau\) with
\(\ordt A\ge1\).
\end{corollary}

\begin{proof}
Existence and uniqueness of \(A\) with \(\ordt A\ge1\): writing
\(G=\sum_{n}g_{n}(L)t^{n}\) and \(A=\sum_{n}a_{n}(L)t^{n}\), by
\cref{plt:prop:dif-block} the block of \(A\) at index \(n\ge1\) must satisfy
\((\partial_{L}-n)a_{n}=g_{n+1}\), which for \(n\ne0\) has a unique solution
in \(\K[L]\) by \cref{plt:lem:dif-triangular}, and \(\ordt G\ge2\) means only
indices \(n\ge1\) occur; the resonant index \(n=0\), the one carrying the
constant of \cref{plt:cor:dif-kernel} and the extra logarithmic degree of
\eqref{plt:eq:dif-profile}, is not reached.  Now apply
\cref{plt:thm:an-integration}: with \(h(X)=\int_{X_{0}}^{X}g\) there is
\(c\in\ComplexNumbers\) with \(h-c\approx A\).  Since \(\ordt A\ge1\) we have
\(\bigl(\Trunc{A}{1}\bigr)(X)=0\), so \eqref{plt:eq:an-strong} at cutoff
\(1\) gives \(h(X)-c\to0\); hence \(\int_{X_{0}}^{\infty}g\) converges, with
value \(c\), and \(\int_{X}^{\infty}g=c-h(X)=-(h-c)\approx-A\).
\end{proof}

\begin{remark}[Where the resonant block bites]
\label{plt:rmk:an-resonance}
If \(\ordt G=1\), the antiderivative acquires a block at index \(0\) of
logarithmic degree one higher than that of \(g_{1}\)
(\eqref{plt:eq:dif-profile}): integrating \(p(\log s)s^{-1}\) produces a
polynomial in \(\log X\), not a constant.  This is the analytic shadow of the
resonance of \cref{plt:lem:dif-resonant}, and it is the reason
\cref{plt:cor:an-tail} needs \(\ordt G\ge2\) and not merely \(\ordt G\ge1\).
In \(\PLrat\) the same mechanism at the block \(L^{-1}t\) produces
\(\log\log X\) and leaves the class altogether
(\cref{plt:thm:dif-anti-rat,plt:cor:dif-loglog}).
\end{remark}

\begin{theorem}[Termwise differentiation is rigid where it is possible]
\label{plt:thm:an-diff-rigid}
Let \(f\) be differentiable on a ray with \(f\approx F\in\PLlau\).  If
\(f'\approx G\) for \emph{some} \(G\in\PLlau\), then \(G=\DX F\).
\end{theorem}

\begin{proof}
First, \(f'\) is locally integrable on a tail.  It is Borel measurable, being
a pointwise limit of continuous difference quotients, and \(f'\approx G\)
applied at the cutoff \(\ordt G+1\) exhibits \(f'\) on some \([X_{1},\infty)\)
as a continuous function plus a remainder bounded on compacts; hence \(f'\)
is bounded on compact subsets of \([X_{1},\infty)\).  A function that is
differentiable everywhere on a compact interval with Lebesgue integrable
derivative satisfies the fundamental theorem of calculus there, so
\(f(X)-f(X_{1})=\int_{X_{1}}^{X}f'\) for \(X\ge X_{1}\).

Now \cref{plt:thm:an-integration} applied to \(g=f'\) gives, for any
\(A\in\PLlau\) with \(\DX A=G\), a constant \(c\in\ComplexNumbers\) with
\(f-\gamma\approx A\), where \(\gamma=c+f(X_{1})\).  Also \(f\approx F\), so
\(f-\gamma\approx F-\gamma\).  Both \(A\) and \(F-\gamma\) lie in
\(\PLlau\) over the field \(\K(\gamma)\subseteq\ComplexNumbers\), so
\cref{plt:prop:an-uniqueness}, applied over that field, gives
\(F-\gamma=A\).  Applying \(\DX\) and using \(\ker\DX=\K(\gamma)\)
(\cref{plt:cor:dif-kernel}) gives \(\DX F=\DX A=G\).
\end{proof}

\begin{remark}[The exact shape of the extra hypothesis]
\label{plt:rmk:an-diff-hypothesis}
\Cref{plt:thm:an-diff-rigid} isolates the content of the folklore statement
``an asymptotic expansion may not be differentiated termwise''.  The formal
answer is never wrong; what can fail is that \(f'\) has \emph{no} expansion at
all.  So the extra hypothesis needed for
\(f\approx F\Rightarrow f'\approx\DX F\) is exactly:
\begin{center}
\emph{\(f'\) admits some polynomial--logarithmic asymptotic expansion.}
\end{center}
This is a hypothesis about \(f'\), not about \(F\), and no property of \(F\)
implies it.  \Cref{plt:thm:an-sectorial-diff} gives the standard checkable
sufficient condition; a Hardy-field or tame-germ hypothesis is another
\cite{vdDMM-1997,vdDMM-2001,adh-book}.
\end{remark}

\begin{theorem}[Sectorial expansions may be differentiated termwise]
\label{plt:thm:an-sectorial-diff}
Let \(0<\theta\le\pi\) and \(R>1\), let
\(S(\theta,R)=\SetOf{z\in\ComplexNumbers:\AbsoluteValue{z}\ge R,\
\AbsoluteValue{\arg z}<\theta}\) with the principal logarithm, and let \(f\)
be holomorphic on \(S(\theta,R)\).  Suppose \(F\in\PLrat\) and that for every
\(\theta'<\theta\) and every \(N\) there are \(C,M,R'\) with
\begin{equation}\label{plt:eq:an-sector-hyp}
 \AbsoluteValue{f(z)-\bigl(\Trunc{F}{N}\bigr)(z)}
 \le C\AbsoluteValue{z}^{-N}\bigl(\log\AbsoluteValue{z}\bigr)^{M}
 \qquad\bigl(\AbsoluteValue{z}\ge R',\ \AbsoluteValue{\arg z}\le\theta'\bigr).
\end{equation}
Then \(f'\) satisfies \eqref{plt:eq:an-sector-hyp} with \(F\) replaced by
\(\DX F\): the expansion may be differentiated termwise, on every closed
subsector.
\end{theorem}

\begin{proof}
Fix \(\theta'<\theta\) and \(N\).  Choose \(\theta''\) with
\(\theta'<\theta''<\theta\) and put \(\delta=\min\bigl(\tfrac12,
\sin(\theta''-\theta')\bigr)\).  For \(z\) with
\(\AbsoluteValue{\arg z}\le\theta'\), every \(w\) with
\(\AbsoluteValue{w-z}\le\delta\AbsoluteValue{z}\) satisfies
\(\AbsoluteValue{\arg w-\arg z}\le\arcsin\delta\le\theta''-\theta'\), hence
\(\AbsoluteValue{\arg w}\le\theta''\), and
\((1-\delta)\AbsoluteValue{z}\le\AbsoluteValue{w}\le(1+\delta)
\AbsoluteValue{z}\).

Apply \eqref{plt:eq:an-sector-hyp} on the subsector \(\theta''\) at cutoff
\(N+1\): setting \(G=f-\Trunc{F}{N+1}\), which is holomorphic there, there
are \(C,M,R''\) with
\[
 \AbsoluteValue{G(w)}\le C\AbsoluteValue{w}^{-N-1}
 \bigl(\log\AbsoluteValue{w}\bigr)^{M}
 \qquad\bigl(\AbsoluteValue{\arg w}\le\theta'',\ \AbsoluteValue{w}\ge
 R''\bigr).
\]
For \(\AbsoluteValue{z}\ge2R''\) with \(\AbsoluteValue{\arg z}\le\theta'\),
the circle \(\AbsoluteValue{w-z}=\delta\AbsoluteValue{z}\) lies in that
region, so Cauchy's estimate gives, with
% ed.: the displayed constant was written (1-\delta)^{-N-1}, which bounds
% ed.: |w|^{-N-1} by ((1-\delta)|z|)^{-N-1} only when N+1\ge0.  Cutoffs here run
% ed.: over all of Z (\ordt F may be negative), and for N\le-2 the maximum of
% ed.: |w|^{-N-1} on the circle is attained at the outer end instead.  Replaced
% ed.: by the two-sided constant \kappa, which covers both signs.
\(\kappa=\max\bigl((1-\delta)^{-N-1},(1+\delta)^{-N-1}\bigr)\),
\[
 \AbsoluteValue{G'(z)}
 \le\frac{1}{\delta\AbsoluteValue{z}}
   \max_{\AbsoluteValue{w-z}=\delta\AbsoluteValue{z}}\AbsoluteValue{G(w)}
 \le\frac{C\,\kappa\,\bigl(\log((1+\delta)\AbsoluteValue{z})\bigr)^{M}}
        {\delta\,\AbsoluteValue{z}^{\,N+2}}
 \le C'\AbsoluteValue{z}^{-(N+2)}\bigl(\log\AbsoluteValue{z}\bigr)^{M}
\]
for \(\AbsoluteValue{z}\) large.  Finally, exactly as in
\eqref{plt:eq:an-trunc-commutes}, the derivative of \(\Trunc{F}{N+1}\) with
respect to \(z\) equals \(\Trunc{\DX F}{N+2}\), so that
\[
 G'=f'-\Trunc{\DX F}{N+2}.
\]
The display above is therefore \eqref{plt:eq:an-sector-hyp} for \(f'\)
against \(\DX F\), at the cutoff \(N+2\).  Since \(N\) was arbitrary and
such cutoffs are cofinal, the conclusion holds at every cutoff by
\cref{plt:lem:mul-asymptotic-basic}(1).
\end{proof}

\begin{remark}[The hypothesis is sectorial, not real-analytic]
\label{plt:rmk:an-sector-needed}
\Cref{plt:ex:an-diff-fails} shows why \eqref{plt:eq:an-sector-hyp} must be
required on a genuine two-dimensional sector.  The function
\(\EulerE^{-z}\sin\EulerE^{z}\) is entire, and its restriction to the ray is
flat; but on any sector \(\AbsoluteValue{\arg z}\le\theta'\) with
\(\theta'>0\) it is unbounded, so it fails
\eqref{plt:eq:an-sector-hyp} against \(F=0\) at every cutoff.  ``Analytic''
alone is worth nothing here; ``analytic with a uniform remainder bound on a
sector of positive opening'' is worth Cauchy's estimate, which is the whole
of the proof.
\end{remark}

\begin{summarybox}
On the polynomial--logarithmic scale, the four elementary transfers behave as
follows.  \emph{Sums and products}: unconditional
(\cref{plt:lem:mul-asymptotic-basic}, \cref{plt:thm:mul-analytic-product}).
\emph{Quotients}: unconditional once the divisor's leading block is nonzero,
with the quotient living in \(\PLrat\) unless that block is a constant
(\cref{plt:thm:an-quotient}).  \emph{Substitution into a function holomorphic
at \(0\)}: unconditional (\cref{plt:prop:an-holomorphic-substitution}).
\emph{Composition with a general inner map}: three geometric hypotheses, each
necessary (\cref{plt:thm:cmp-analytic-transfer}).  \emph{Integration}:
unconditional (\cref{plt:thm:an-integration}).  \emph{Differentiation}: not
transferable at all without an extra hypothesis, but rigid once \(f'\) is
known to have any expansion (\cref{plt:thm:an-diff-rigid},
\cref{plt:thm:an-sectorial-diff}).  Every function with an expansion is a
member of a coset of the flat ideal, and every formal series is realized
(\cref{plt:thm:an-borel-ritt}, \cref{plt:cor:an-fibres}).
\end{summarybox}

\chapter{Residual-to-error transfer}
\label{plt:sec:residual-transfer}

A formal reversion produces a candidate \(G_{N}\) and a certificate: the
residual \(\RevErr{F}{G_{N}}\) has \(t\)-order at least \(N\)
(\cref{plt:cor:rev-finite-certificate}).  A user of the answer wants a
number: how far is \(G_{N}(x)\) from the true solution of \(F(y)=x\)?  The
passage from the first to the second is the subject of this chapter.  It has
three steps, and only the first is trivial:
\begin{enumerate}
\item convert a residual \emph{value} into an error bound (calculus);
\item convert the formal residual \emph{order} into an analytic residual
\emph{size} (\cref{plt:thm:an-residual-analytic}, which needs an analytic
hypothesis on the map being reversed);
\item verify that the true solution exists, is unique, and lies where the
derivative bound holds (\cref{plt:thm:an-enclosure}).
\end{enumerate}
Steps (2) and (3) are exactly what the sources omit.

\section{The transfer inequality and its enclosure form}
\label{plt:sec:an-transfer}

\begin{remark}[What \cref{plt:prop:rev-analytic-transfer} does and does not
assume]
\label{plt:rmk:an-transfer-recall}
% ed.: \cref{plt:prop:rev-analytic-transfer} is printed with the analytic
% ed.: inverse written \rev{f}.  In this volume that symbol is reserved for the
% ed.: FORMAL compositional inverse -- and the whole content of
% ed.: (plt:thm:an-inverse-transfer) is that the formal and the analytic
% ed.: inverse agree asymptotically, a statement that is vacuous if the two are
% ed.: written with one glyph.  This is the same collision the ed. note before
% ed.: (plt:sec:an-omega-constants) records against S2.  The recollection below
% ed.: therefore names x_* by the equation it solves.
\Cref{plt:prop:rev-analytic-transfer} states: if \(f\) is \(C^{1}\) and
injective on an interval \(I\), \(y\in f(I)\), \(x_{*}\) is the unique point of
\(I\) with \(f(x_{*})=y\),
\(x_{N}\in I\) and \(\AbsoluteValue{f'}\ge m>0\) between \(x_{N}\) and
\(x_{*}\), then
\begin{equation}\label{plt:eq:an-transfer}
 \AbsoluteValue{x_{N}-x_{*}}\le\frac{\AbsoluteValue{f(x_{N})-y}}{m}.
\end{equation}
Two observations about the hypotheses, neither of them in the sources.
First, injectivity is used only to \emph{name} \(x_{*}\); the mean-value
argument itself needs only differentiability on the segment joining
\(x_{N}\) to \(x_{*}\).  Second, and more seriously, the derivative bound is
required on an interval whose right endpoint is the unknown.  In an actual
computation one knows \(x_{N}\) and \(f(x_{N})-y\) and does not know
\(x_{*}\); a bound ``between \(x_{N}\) and \(x_{*}\)'' is therefore not
directly verifiable, and using \eqref{plt:eq:an-transfer} as stated is
circular unless a separate argument locates \(x_{*}\).
\Cref{plt:thm:an-enclosure} removes the circularity by hypothesising the
derivative bound on a \emph{chosen} interval around \(x_{N}\) and deducing
the existence of \(x_{*}\) inside it.
\end{remark}

\begin{theorem}[Residual enclosure: existence, uniqueness, and two-sided error]
\label{plt:thm:an-enclosure}
Let \(\rho>0\), let \(I=[x_{N}-\rho,x_{N}+\rho]\subseteq\RealNumbers\), and
let \(f\colon I\to\RealNumbers\) be continuously differentiable with
\[
 m\DefinitionEquals\min_{x\in I}\AbsoluteValue{f'(x)}>0,
 \qquad
 M\DefinitionEquals\max_{x\in I}\AbsoluteValue{f'(x)} .
\]
Let \(y\in\RealNumbers\) satisfy
\begin{equation}\label{plt:eq:an-enclosure-hyp}
 \AbsoluteValue{f(x_{N})-y}\ \le\ m\rho .
\end{equation}
Then \(f\) is strictly monotone on \(I\); there is exactly one \(x_{*}\in I\)
with \(f(x_{*})=y\); and
\begin{equation}\label{plt:eq:an-enclosure}
 \frac{\AbsoluteValue{f(x_{N})-y}}{M}
 \ \le\ \AbsoluteValue{x_{N}-x_{*}}\ \le\
 \frac{\AbsoluteValue{f(x_{N})-y}}{m}.
\end{equation}
No hypothesis of injectivity is imposed; it is a conclusion.
\end{theorem}

\begin{proof}
Since \(f'\) is continuous on the interval \(I\) and never zero, the
intermediate value theorem applied to \(f'\) shows that \(f'\) has constant
sign; replacing \(f\) by \(-f\) and \(y\) by \(-y\) if necessary, assume
\(f'\ge m>0\).  Then \(f\) is strictly increasing, hence injective.

By the mean value theorem,
\(f(x_{N}+\rho)-f(x_{N})\ge m\rho\ge y-f(x_{N})\) and
\(f(x_{N})-f(x_{N}-\rho)\ge m\rho\ge f(x_{N})-y\), using
\eqref{plt:eq:an-enclosure-hyp} in both.  Hence
\(f(x_{N}-\rho)\le y\le f(x_{N}+\rho)\), and the intermediate value theorem
applied to \(f\) produces \(x_{*}\in I\) with \(f(x_{*})=y\); it is unique by
injectivity.

Finally the mean value theorem gives \(\xi\) strictly between \(x_{N}\) and
\(x_{*}\) (or \(x_{N}=x_{*}\), in which case
\eqref{plt:eq:an-enclosure} is trivial) with
\(f(x_{N})-y=f(x_{N})-f(x_{*})=f'(\xi)(x_{N}-x_{*})\).  Since
\(m\le\AbsoluteValue{f'(\xi)}\le M\), both inequalities of
\eqref{plt:eq:an-enclosure} follow.
\end{proof}

\begin{theorem}[Complex enclosure by Rouch\'e]
\label{plt:thm:an-complex-enclosure}
Let \(f\) be holomorphic on a neighbourhood of the closed disc
\(\overline D=\SetOf{z:\AbsoluteValue{z-x_{N}}\le\rho}\subseteq\ComplexNumbers\),
put \(\mu=\AbsoluteValue{f'(x_{N})}\), assume \(\mu>0\), and suppose
\begin{equation}\label{plt:eq:an-complex-hyp}
 \AbsoluteValue{f'(z)-f'(x_{N})}\le\frac{\mu}{4}\ \ (z\in\overline D),
 \qquad
 \AbsoluteValue{f(x_{N})-y}<\frac{\mu\rho}{4}.
\end{equation}
Then \(f\) is injective on \(\overline D\), the equation \(f(z)=y\) has
exactly one solution \(x_{*}\) in \(D\), and
\begin{equation}\label{plt:eq:an-complex-bound}
 \frac{4\AbsoluteValue{f(x_{N})-y}}{5\mu}
 \ \le\ \AbsoluteValue{x_{N}-x_{*}}
 \ \le\ \frac{4\AbsoluteValue{f(x_{N})-y}}{3\mu}.
\end{equation}
\end{theorem}

\begin{proof}
Write \(g(z)=f(z)-y\) and let
\(h(z)=\bigl(f(x_{N})-y\bigr)+f'(x_{N})(z-x_{N})\) be its linearisation at
\(x_{N}\).  For \(z\in\overline D\), integrating \(f'-f'(x_{N})\) along the
segment from \(x_{N}\) to \(z\), which lies in the convex set \(\overline D\),
gives
\begin{equation}\label{plt:eq:an-linearisation}
 \AbsoluteValue{g(z)-h(z)}
 =\Bigl|\int_{0}^{1}\bigl(f'(x_{N}+s(z-x_{N}))-f'(x_{N})\bigr)(z-x_{N})
 \,\Differential s\Bigr|
 \le\frac{\mu}{4}\AbsoluteValue{z-x_{N}} .
\end{equation}
On the circle \(\AbsoluteValue{z-x_{N}}=\rho\) this is \(\le\mu\rho/4\), while
\(\AbsoluteValue{h(z)}\ge\mu\rho-\AbsoluteValue{f(x_{N})-y}
>\mu\rho-\mu\rho/4=3\mu\rho/4\).  Thus
\(\AbsoluteValue{g-h}<\AbsoluteValue{h}\) on the circle, and Rouch\'e's
theorem gives that \(g\) and \(h\) have equally many zeros in \(D\), counted
with multiplicity.  The only zero of \(h\) is
\(z_{0}=x_{N}-(f(x_{N})-y)/f'(x_{N})\), and
\(\AbsoluteValue{z_{0}-x_{N}}<(\mu\rho/4)/\mu=\rho/4<\rho\), so \(z_{0}\in
D\) and \(g\) has exactly one zero \(x_{*}\in D\).

For injectivity and the error bounds, note that for \(z_{1},z_{2}\in\overline
D\),
\[
 f(z_{1})-f(z_{2})=(z_{1}-z_{2})\int_{0}^{1}
 f'\bigl(z_{2}+s(z_{1}-z_{2})\bigr)\Differential s,
\]
and the integral differs from \(f'(x_{N})\) by at most \(\mu/4\) by
\eqref{plt:eq:an-complex-hyp}; hence its modulus lies in
\([\,3\mu/4,\;5\mu/4\,]\) because \(\AbsoluteValue{f'(x_{N})}=\mu\).
Injectivity follows from the lower end.  Taking \(z_{1}=x_{N}\) and
\(z_{2}=x_{*}\),
\[
 \tfrac{3\mu}{4}\AbsoluteValue{x_{N}-x_{*}}
 \ \le\ \AbsoluteValue{f(x_{N})-y}\ \le\
 \tfrac{5\mu}{4}\AbsoluteValue{x_{N}-x_{*}},
\]
which is \eqref{plt:eq:an-complex-bound}.
\end{proof}

\begin{remark}[Using a lower bound for \(\AbsoluteValue{f'(x_{N})}\)]
\label{plt:rmk:an-complex-mu}
If \(\mu\) is only known to satisfy \(0<\mu\le\AbsoluteValue{f'(x_{N})}\),
the proof of \cref{plt:thm:an-complex-enclosure} still gives injectivity,
existence and uniqueness of \(x_{*}\in D\), and the upper bound
\(\AbsoluteValue{x_{N}-x_{*}}\le\frac{4}{3\mu}\AbsoluteValue{f(x_{N})-y}\);
only the left-hand inequality of \eqref{plt:eq:an-complex-bound} uses
\(\AbsoluteValue{f'(x_{N})}=\mu\).  This is the asymmetry that produces the
constant \(5\) on the left and \(3\) on the right.
\end{remark}

\begin{warningbox}[title={A formal residual is not an inequality}]
The statement \(\RevErr{F}{G_{N}}=\FormalBigOAt{t^{N}}{t}\) is a membership in
a filtration step of \(\PLpow\) (\cref{plt:def:trunc-formal-o}); the
statement \(F(G_{N}(x))-x=\BigOAt{x^{-N}(\log x)^{M}}{x\to+\infty}\) is an
inequality between real numbers valid beyond some \(x_{N}\).  Neither implies
the other without a theorem.  The direction one needs is supplied by
\cref{plt:thm:an-residual-analytic}, and its hypothesis is that the map being
reversed \emph{has} a polynomial--logarithmic asymptotic expansion --- an
analytic fact about a function, of the kind that
\cref{plt:thm:an-borel-ritt} shows cannot be recovered from the coefficients.
\end{warningbox}

% ed.: S5 writes, immediately after its residual-to-error principle: "For
% ed.: f(x) = x + log x, this proves the asymptotic validity of every
% ed.: truncation omega_N directly from [the formal residual identity]."  It
% ed.: does not: the principle converts an analytic residual into an analytic
% ed.: error, and the formal identity provides no analytic residual.  The
% ed.: missing step is the one carried out in (plt:thm:om-analytic) for this
% ed.: map and in (plt:thm:an-residual-analytic) in general.  The same gap is
% ed.: present in S1 ("by construction") and is repaired in the same way.
\section{From the formal residual order to an analytic residual size}
\label{plt:sec:an-residual-size}

% ed.: the coefficient field was left at the ambient \(\K\subseteq\ComplexNumbers\)
% ed.: of this part, but the proof applies \cref{plt:thm:cmp-analytic-transfer},
% ed.: whose hypothesis (A2) asks for an inner realization with REAL
% ed.: coefficients, and both \(\varphi\) and \(h\) are real valued here.  The
% ed.: real coefficient field is now printed, as house rule 5 requires.
\begin{theorem}[Analytic size of a residual]
\label{plt:thm:an-residual-analytic}
Let \(\K\subseteq\RealNumbers\), let \(Y_{0}>1\), let \(\Phi=X+A\in\Gnear\), and let
\(\varphi\colon[Y_{0},\infty)\to\RealNumbers\) satisfy \(\varphi\approx\Phi\)
in the sense of \cref{plt:def:an-expansion}, with the real logarithm and the
regime \(Y\to+\infty\).  Let \(H=X+B\in\Gnear\) with \(B\in\K[L][t]\) a
\emph{finite} sum of blocks, and let
\(h(x)=x+B(\log x,1/x)\) be its evaluation for large real \(x\).  Put
\(r=\ordt\bigl(\RevErr{\Phi}{H}\bigr)\in\NaturalNumbers\cup\SetOf{+\infty}\).
Then \(h(x)\to+\infty\), \(h(x)/x\to1\), and for every finite
\(N\le r\) there are \(M,C,x_{1}\) with
\begin{equation}\label{plt:eq:an-residual-analytic}
 \AbsoluteValue{\varphi\bigl(h(x)\bigr)-x}
 \ \le\ C\,x^{-N}(\log x)^{M}
 \qquad(x\ge x_{1}).
\end{equation}
\end{theorem}

\begin{proof}
Since \(\ordt B\ge0\) and \(B\) is a finite sum of blocks,
\(B(\log x,1/x)=\BigOAt{(\log x)^{d}}{x\to+\infty}\) with
\(d=\max_{n}\degL b_{n}\); hence \(h(x)=x\bigl(1+u(x)\bigr)\) with
\(u(x)=B(\log x,1/x)/x\to0\), and \(h(x)\to+\infty\), \(h(x)/x\to1\).  In
particular \(h(x)\ge Y_{0}\) for large \(x\), so \(\varphi\compose h\) is
defined there.

We verify the three hypotheses of \cref{plt:thm:cmp-analytic-transfer} for
the outer function \(\varphi\) and the inner function \(h\).  (A1) is the
previous paragraph.  (A2) holds with \(c=1\), \(\rho=1\) and \(U=tB\in
t\PLpow\), whose coefficients are real because \(\K\subseteq\RealNumbers\), as
that hypothesis of \cref{plt:thm:cmp-analytic-transfer} requires: the function
\(u\) is the evaluation of the finite sum \(tB\), so
\(u\approx U\) exactly (its remainder vanishes at every cutoff beyond the
\(t\)-degree of \(tB\)), and \(u\to0\).  (A3) holds automatically on the
positive real ray with \(c=1>0\).  By \cref{plt:thm:an-equivalence} the
relation \(\approx\) of \cref{plt:def:an-expansion} and the relation
\(\sim_{S}\) used in \cref{plt:thm:cmp-analytic-transfer} coincide, so the
theorem applies and yields
\[
 \varphi\compose h\ \approx\ \Phi\compose H .
\]
The identity function realizes the series \(X\), so by
\cref{plt:lem:mul-asymptotic-basic}(3) the germ
\(x\mapsto\varphi(h(x))-x\) is PL-asymptotic to
\(\Phi\compose H-X=\RevErr{\Phi}{H}\).  Let \(N\le r\) be finite.  Since
\(\ordt\RevErr{\Phi}{H}=r\ge N\), we have
\(\Trunc{\RevErr{\Phi}{H}}{N}=0\), and \eqref{plt:eq:an-strong} at cutoff
\(N\) is exactly \eqref{plt:eq:an-residual-analytic}.
\end{proof}

\begin{theorem}[Transfer of a formal inverse expansion to the analytic
inverse]
\label{plt:thm:an-inverse-transfer}
% ed.: coefficient field printed: \cref{plt:thm:an-residual-analytic} is applied
% ed.: below, and it now carries \(\K\subseteq\RealNumbers\).
Let \(\K\subseteq\RealNumbers\) and assume:
\begin{enumerate}
\item[\textup{(H1)}] \emph{Regime and derivative bound.}
\(\varphi\colon[Y_{0},\infty)\to\RealNumbers\) is \(C^{1}\) with
\(\varphi'(Y)\ge m_{0}>0\) for all \(Y\ge Y_{0}\), and
\(\varphi(Y)\to+\infty\).  \textup{(}Then \(\varphi\) is a strictly
increasing bijection onto \([\varphi(Y_{0}),\infty)\); write \(\psi\) for its
inverse.\textup{)}
\item[\textup{(H2)}] \emph{Analytic expansion of the map.}
\(\varphi\approx\Phi\) for some \(\Phi=X+A\in\Gnear\), with the real
logarithm as \(Y\to+\infty\).
\end{enumerate}
Let \(G=\rev{\Phi}=X+B\in\Gnear\) be the formal compositional inverse
supplied by \cref{plt:thm:rev-existence}, and put
\(g_{N}(x)=x+\bigl(\Trunc{B}{N}\bigr)(x)\).
Then for every \(N\in\NaturalNumbers\) there are
\(M_{N}\in\NaturalNumbers\), \(C_{N}>0\) and a threshold \(\tau_{N}\) with
\begin{equation}\label{plt:eq:an-inverse-transfer}
 \AbsoluteValue{\psi(x)-g_{N}(x)}\ \le\ C_{N}\,x^{-N}(\log x)^{M_{N}}
 \qquad(x\ge\tau_{N}),
\end{equation}
and consequently \(\psi\approx G\): the formal compositional inverse is the
asymptotic expansion of the analytic inverse.
\end{theorem}

\begin{proof}
\emph{(H1) makes \(\psi\) well defined.}  \(\varphi'\ge m_{0}>0\) forces
\(\varphi\) strictly increasing, hence injective, and
\(\varphi(Y)\to+\infty\) with the intermediate value theorem makes it onto
\([\varphi(Y_{0}),\infty)\); \(\psi\) is its inverse, \(\psi(x)\ge Y_{0}\),
and \(\psi(x)\to+\infty\).

\emph{Formal order of the residual.}  Fix \(N\) and put
\(G_{N}=X+\Trunc{B}{N}\in\Gnear\), a finite sum of blocks.  Then
\(G-G_{N}=B-\Trunc{B}{N}\) has \(\ordt\ge N\), so by
\cref{plt:prop:rev-residual-order}
\[
 \ordt\bigl(\RevErr{\Phi}{G_{N}}\bigr)=\ordt(G_{N}-G)\ \ge\ N .
\]

\emph{Analytic size of the residual.}  By (H2) and
\cref{plt:thm:an-residual-analytic} applied with \(H=G_{N}\), there are
\(M_{N},C,x_{1}\) with
\(\AbsoluteValue{\varphi(g_{N}(x))-x}\le Cx^{-N}(\log x)^{M_{N}}\) for
\(x\ge x_{1}\).

\emph{Transfer.}  For \(x\) large, both \(g_{N}(x)\) and \(\psi(x)\) lie in
the interval \([Y_{0},\infty)\), on all of which
\(\AbsoluteValue{\varphi'}\ge m_{0}\).  Applying
\cref{plt:prop:rev-analytic-transfer} on that interval with \(f=\varphi\),
argument value \(y=x\), approximation \(g_{N}(x)\), exact solution
\(x_{*}=\psi(x)\) and \(m=m_{0}\) gives
\[
 \AbsoluteValue{\psi(x)-g_{N}(x)}
 \le\frac{\AbsoluteValue{\varphi(g_{N}(x))-x}}{m_{0}}
 \le\frac{C}{m_{0}}\,x^{-N}(\log x)^{M_{N}},
\]
which is \eqref{plt:eq:an-inverse-transfer}.  Since
\(g_{N}(x)=\bigl(\Trunc{G}{N}\bigr)(x)\) for every \(N\ge0\), this is
\eqref{plt:eq:an-strong} at every cutoff \(N\ge0\), and cutoffs \(N<0\)
follow from \cref{plt:lem:mul-asymptotic-basic}(1).  Hence \(\psi\approx G\).
\end{proof}

\begin{remark}[Which hypothesis does which job]
\label{plt:rmk:an-hypothesis-roles}
(H1) is the only source of the constant \(m_{0}\), and it must be verified on
the function, not deduced from an expansion: by
\cref{plt:ex:an-diff-fails} the expansion \(\Phi=X+A\) does not entail
\(\varphi'=1+\LittleOAt{1}{Y\to+\infty}\), or indeed any statement about
\(\varphi'\) whatever.  (H2) is what turns a \(t\)-adic order into an
inequality; it is an analytic hypothesis and no algebraic theorem supplies
it.  A third condition, the admissibility check \(g_{N}(x)\ge Y_{0}\), is
listed as a hypothesis in several of the sources and in the four-item list
preceding \cref{plt:thm:om-analytic}; it is \emph{not} a hypothesis of
\cref{plt:thm:an-inverse-transfer}, because \(g_{N}(x)=x+\BigOAt{(\log
x)^{d}}{x\to+\infty}\to+\infty\) makes it automatic in the limit.  It is
nonetheless indispensable at any \emph{specific} \(x\) at which one wants to
use the resulting inequality: a truncated formal inverse can take values
outside \([Y_{0},\infty)\) for moderate \(x\), and then
\(\varphi(g_{N}(x))\) is not even defined, so the numerical bound is vacuous
while the asymptotic theorem remains true.
% ed.: this sentence read "Item (2) of the list preceding
% ed.: \cref{plt:thm:om-analytic} is the specialisation of (H2) to
% ed.: \varphi(Y)=Y+\log Y", which is wrong in an instructive way: (H2)
% ed.: specialised to that map is the trivial statement \varphi\approx X+L (the
% ed.: remainder vanishes identically), whereas item (2) asks for an analytic
% ed.: remainder for log applied to the TRUNCATION.  Item (2) is the conclusion
% ed.: of \cref{plt:thm:an-residual-analytic}, not its hypothesis.  Corrected.
Item (2) of the list preceding \cref{plt:thm:om-analytic} --- an analytic
remainder for the truncated logarithm --- is not (H2) but what
\cref{plt:thm:an-residual-analytic} \emph{deduces} from (H2); for
\(\varphi(Y)=Y+\log Y\) the hypothesis (H2) itself is trivial, since
\(\Phi=X+L\) is realized by \(\varphi\) with zero remainder at every cutoff.
\end{remark}

% ed.: S2's transfer theorem lists the hypotheses "F(y) = y + o(y)" and
% ed.: "F'(y) = 1 + o(1)" and then assumes an ANALYTIC residual bound
% ed.: F(G_N(X)) - X = O(X^{-N-1}L^M) outright.  Assuming it is legitimate;
% ed.: the surrounding text, however, obtains it from the formal reversion, so
% ed.: the theorem as used is circular.  The first hypothesis is never used in
% ed.: S2's proof.  (H2) above is the honest replacement, and
% ed.: (plt:thm:an-residual-analytic) discharges it.  S2 also writes the
% ed.: analytic inverse as \rev{F}; in this volume that symbol denotes the
% ed.: FORMAL compositional inverse, and the two must not share a glyph in a
% ed.: statement whose whole content is that they agree.
\section{Explicit constants for the Wright omega truncations}
\label{plt:sec:an-omega-constants}

For the archetype \(\varphi(Y)=Y+\log Y\) the constants of
\cref{plt:thm:an-enclosure} can be written down exactly, with no asymptotic
hypothesis at all.  This is the strongest form in which a residual bound is
available anywhere in this volume, and it is available because
\(\varphi'(Y)=1+Y^{-1}\) is monotone and explicitly bounded on
\((0,\infty)\).

\begin{proposition}[A computable two-sided error bound, valid for every \(x>1\)]
\label{plt:prop:an-omega-computable}
Let \(x>1\), let \(N\in\NaturalNumbers\), let
\(\WrightOmega_{N}(x)=x-\log x+\sum_{n=1}^{N}\LamP{n}(\log x)x^{-n}\) be the
evaluation of the truncation of \cref{plt:def:om-truncation}, and suppose
\(\WrightOmega_{N}(x)>0\).  Put
\[
 R_{N}(x)=\WrightOmega_{N}(x)+\log\WrightOmega_{N}(x)-x,
 \qquad
 \mu_{N}(x)=\min\SetOf{\WrightOmega_{N}(x),\ x-\log x} .
\]
Then \(\mu_{N}(x)>0\) and
\begin{equation}\label{plt:eq:an-omega-bracket}
 \frac{\AbsoluteValue{R_{N}(x)}}{1+\mu_{N}(x)^{-1}}
 \ \le\ \AbsoluteValue{\WrightOmega(x)-\WrightOmega_{N}(x)}
 \ \le\ \AbsoluteValue{R_{N}(x)} ,
\end{equation}
and \(\WrightOmega(x)-\WrightOmega_{N}(x)\) has the sign of \(-R_{N}(x)\).
Both bounds are computable from \(x\) and \(N\) alone; neither involves an
implied constant, an unspecified threshold, or a limit.
\end{proposition}

\begin{proof}
By \cref{plt:prop:om-real-bijection}, \(\varphi(Y)=Y+\log Y\) is a strictly
increasing bijection \((0,\infty)\to\RealNumbers\) with
\(\varphi'(Y)=1+Y^{-1}\); the same proposition gives \(\WrightOmega(1)=1\) and
% ed.: the step "for x>1 the envelope gives omega(x) > x - log x" skipped the
% ed.: verification of that envelope's own hypothesis, which is omega(x)>1 and
% ed.: not x>1.  Supplied: omega is strictly increasing with omega(1)=1.
\(\WrightOmega\) strictly increasing, so \(\WrightOmega(x)>1\) for \(x>1\),
which is the hypothesis under which the envelope
\eqref{plt:eq:om-envelope} gives \(\WrightOmega(x)>x-\log x\).  Also
\(x-\log x\ge1>0\) for every \(x>0\), because \(x\mapsto x-\log x\) has its
minimum at \(x=1\).  Hence both \(\WrightOmega(x)\) and
\(\WrightOmega_{N}(x)\) lie in \([\mu_{N}(x),\infty)\), and so does the whole
segment between them.

On \([\mu_{N}(x),\infty)\) the derivative satisfies
\(1<\varphi'(Y)\le1+\mu_{N}(x)^{-1}\).  Since
\(\varphi(\WrightOmega(x))=x\), the mean value theorem gives \(\zeta\) in
that segment with
\[
 R_{N}(x)=\varphi\bigl(\WrightOmega_{N}(x)\bigr)
 -\varphi\bigl(\WrightOmega(x)\bigr)
 =\varphi'(\zeta)\bigl(\WrightOmega_{N}(x)-\WrightOmega(x)\bigr),
\]
whence the sign statement and, dividing by \(\varphi'(\zeta)\in
(1,1+\mu_{N}(x)^{-1}]\), the two inequalities
\eqref{plt:eq:an-omega-bracket}.
\end{proof}

\begin{proposition}[The two-term approximation, in closed form]
\label{plt:prop:an-omega-explicit0}
For every \(x\ge1\), with \(L=\log x\),
\begin{equation}\label{plt:eq:an-R0-closed}
 R_{0}(x)=(x-L)+\log(x-L)-x=\log\Bigl(1-\frac{L}{x}\Bigr)\ \le\ 0,
\end{equation}
and
\begin{equation}\label{plt:eq:an-R0-bracket}
 \frac{L}{x}\ \le\ \AbsoluteValue{R_{0}(x)}\ \le\ \frac{L}{x}+\frac{L^{2}}{x^{2}} .
\end{equation}
Consequently \(x-\log x\) always underestimates \(\WrightOmega(x)\), and
\begin{equation}\label{plt:eq:an-omega0-bracket}
 \frac{L}{x}\cdot\frac{x-L}{x-L+1}
 \ \le\ \WrightOmega(x)-(x-\log x)
 \ \le\ \frac{L}{x}+\frac{L^{2}}{x^{2}} .
\end{equation}
\end{proposition}

\begin{proof}
For \(x\ge1\) we have \(L\ge0\) and \(x-L\ge1>0\), so
\(\log(x-L)\) is defined and
\((x-L)+\log(x-L)-x=\log(x-L)-L=\log\frac{x-L}{x}=\log(1-L/x)\), which is
\(\le0\).  Set \(u=L/x\).  The function \(x\mapsto(\log x)/x\) attains its
maximum \(\EulerE^{-1}\) at \(x=\EulerE\), so \(0\le u\le\EulerE^{-1}<\tfrac12\).
Then
\[
 -\log(1-u)=\sum_{j\ge1}\frac{u^{j}}{j}
 \ \ge\ u,
 \qquad
 -\log(1-u)\le u+\frac12\sum_{j\ge2}u^{j}
 =u+\frac{u^{2}}{2(1-u)}\le u+u^{2},
\]
the last step because \(u\le\tfrac12\).  This is
\eqref{plt:eq:an-R0-bracket}.
% ed.: the proof invoked \cref{plt:prop:an-omega-computable}, which is stated
% ed.: for x>1, while this proposition is stated for x\ge1.  The endpoint is
% ed.: dispatched separately rather than left to a proposition that does not
% ed.: cover it.
At \(x=1\) all three quantities in
\eqref{plt:eq:an-omega0-bracket} vanish, since \(L=0\) and
\(\WrightOmega(1)=1=x-\log x\), so assume \(x>1\).  Now apply
\cref{plt:prop:an-omega-computable} with \(N=0\), where
\(\WrightOmega_{0}(x)=x-L\) and hence \(\mu_{0}(x)=x-L\): the sign of
\(-R_{0}\ge0\) gives \(\WrightOmega(x)\ge x-L\), the upper bound of
\eqref{plt:eq:an-omega-bracket} combined with the right half of
\eqref{plt:eq:an-R0-bracket} gives the right half of
\eqref{plt:eq:an-omega0-bracket}, and the lower bound of
\eqref{plt:eq:an-omega-bracket} combined with the left half of
\eqref{plt:eq:an-R0-bracket} gives
\(\AbsoluteValue{R_{0}}/(1+(x-L)^{-1})\ge (L/x)(x-L)/(x-L+1)\).
\end{proof}

\begin{remark}[What is \emph{not} obtained]
\label{plt:rmk:an-no-uniform-constant}
\Cref{plt:prop:an-omega-computable} is uniform in \(N\) only because it does
not evaluate \(R_{N}\); the closed form
\eqref{plt:eq:an-R0-closed} is available at \(N=0\) and, by the same
computation, at \(N=1\), but a constant \(C\) with
\(\AbsoluteValue{R_{N}(x)}\le C\AbsoluteValue{\LamP{N+1}(\log x)}x^{-N-1}\)
uniformly in \(N\) is \emph{not} obtained here and is not obtained in
\cref{plt:thm:om-analytic} either: the constant there depends on \(N\), and
the logarithmic exponent \(K=(N+2)^{2}\) in its remainder is quadratic in
\(N\).  Any statement in the literature of the form ``the error is bounded by
the first omitted term'' for this expansion, with a constant independent of
\(N\), is stronger than anything proved here.  We do not assert it.
\end{remark}

\section{A worked bound, to digits}
\label{plt:sec:an-worked}

\begin{example}[The enclosure at \(x=100\), \(N=3\)]
\label{plt:ex:an-worked}
Take \(x=100\), so that
\[
 L=\log 100=4.605170185988091368\ldots,
\]
and use the Lambert polynomials of \cref{plt:def:lambert-polynomials},
\[
 \LamP{1}(u)=u,\qquad
 \LamP{2}(u)=\tfrac12u^{2}-u,\qquad
 \LamP{3}(u)=\tfrac13u^{3}-\tfrac32u^{2}+u .
\]
The three correction terms are
\begin{align*}
 \LamP{1}(L)\,x^{-1}&=4.60517018598809137\times10^{-2},\\
 \LamP{2}(L)\,x^{-2}&=5.99862603496870465\times10^{-4},\\
 \LamP{3}(L)\,x^{-3}&=5.34863899981332245\times10^{-6},
\end{align*}
so that
\[
 \WrightOmega_{3}(100)=95.4414867271142862294\ldots
\]
Substituting into \(\varphi(Y)=Y+\log Y\) gives the residual
\[
 R_{3}(100)=\WrightOmega_{3}(100)+\log\WrightOmega_{3}(100)-100
 =8.23927836460666826\times10^{-8}.
\]
Here \(x-\log x=95.39482981401190863\ldots\) and
\(\WrightOmega_{3}(100)>x-\log x\), so
\(\mu_{3}=95.394829814\ldots\) and \(1+\mu_{3}^{-1}=1.010482748404\ldots\).
\Cref{plt:prop:an-omega-computable} therefore yields the rigorous enclosure
\[
 8.153804088\times10^{-8}
 \ \le\ \AbsoluteValue{\WrightOmega(100)-\WrightOmega_{3}(100)}
 \ \le\ 8.239278365\times10^{-8},
\]
with \(\WrightOmega(100)<\WrightOmega_{3}(100)\) because \(R_{3}>0\).  The
true value, computed independently to \(40\) digits from
\(\WrightOmega(100)=\LambertW_{0}(\EulerE^{100})
=95.44148664557583184017\ldots\), is
\[
 \AbsoluteValue{\WrightOmega(100)-\WrightOmega_{3}(100)}
 =8.153845439\times10^{-8},
\]
inside the enclosure and within \(5\times10^{-6}\) relative of its lower end.
The lower bound is close to sharp because \(\varphi'(\zeta)=1+\zeta^{-1}\) is
close to \(1\); the whole enclosure has relative width
% ed.: the relative width was quoted as \mu_3^{-1}.  Width divided by the upper
% ed.: end is 1-(1+\mu^{-1})^{-1}=(1+\mu)^{-1}, which is slightly smaller;
% ed.: recomputed, it is 1.0374 per cent against \mu_3^{-1}=1.0483 per cent.
\((1+\mu_{3})^{-1}=1.037\%\), and in particular at most
\(\mu_{3}^{-1}\approx1.05\%\).

For comparison, \cref{plt:thm:om-analytic} predicts the error
\(\LamP{4}(L)x^{-4}=-7.593614686\times10^{-8}\) up to a relative correction
\(\BigOAt{L^{K}/x}{x\to+\infty}\); the predicted and true errors agree to
about \(6.9\%\), which is the size of the ratio
\(\LamP{5}(L)x^{-5}/\LamP{4}(L)x^{-4}=7.2\times10^{-2}\).  The enclosure of
\cref{plt:prop:an-omega-computable} has width \(8.55\times10^{-10}\),
about seven times smaller than the \(5.60\times10^{-9}\) by which the
asymptotic prediction misses; and unlike the asymptotic statement it is a
theorem about this one value of \(x\), with no implied constant and no
threshold.
\end{example}

\begin{table}[htbp]
\centering
\caption{Rigorous enclosure of \(\AbsoluteValue{\WrightOmega(100)-\WrightOmega_{N}(100)}\)
by \cref{plt:prop:an-omega-computable}.  All entries computed at \(40\)
significant digits and rounded; \(\mu_{N}=x-\log x=95.394829814\ldots\) in
every row.}
\label{plt:tab:an-worked}
\footnotesize
\begin{tabular}{@{}rllll@{}}
\toprule
\(N\) & \(\AbsoluteValue{R_{N}(100)}\) & lower bound & true error & upper bound\\
\midrule
0 & \(4.714580382\times10^{-2}\)  & \(4.665671324\times10^{-2}\)  & \(4.665683156\times10^{-2}\)  & \(4.714580382\times10^{-2}\)\\
1 & \(6.114700456\times10^{-4}\)  & \(6.051266551\times10^{-4}\)  & \(6.051297040\times10^{-4}\)  & \(6.114700456\times10^{-4}\)\\
2 & \(5.322287245\times10^{-6}\)  & \(5.267073836\times10^{-6}\)  & \(5.267100545\times10^{-6}\)  & \(5.322287245\times10^{-6}\)\\
3 & \(8.239278365\times10^{-8}\)  & \(8.153804088\times10^{-8}\)  & \(8.153845439\times10^{-8}\)  & \(8.239278365\times10^{-8}\)\\
4 & \(5.661006395\times10^{-9}\)  & \(5.602279113\times10^{-9}\)  & \(5.602307524\times10^{-9}\)  & \(5.661006395\times10^{-9}\)\\
5 & \(1.558418770\times10^{-10}\) & \(1.542251733\times10^{-10}\) & \(1.542259554\times10^{-10}\) & \(1.558418770\times10^{-10}\)\\
\bottomrule
\end{tabular}
\end{table}

\begin{remark}[Cancellation in the residual, with the digit count]
\label{plt:rmk:an-cancellation}
\Cref{plt:tab:an-worked} is a table of \emph{differences of numbers of size
\(95\)}.  At \(N=5\) the residual is \(1.56\times10^{-10}\), so forming it
loses about
\(\log_{10}(95.44/1.56\times10^{-10})\approx11.8\) decimal digits.  A
double-precision evaluation of \(R_{N}(x)\) carries roughly \(16\) digits and
is therefore already worthless at \(N=6\), where \(R_{6}\approx2.05\times
10^{-12}\) and the loss is about \(13.7\) digits.  The bound of
\cref{plt:prop:an-omega-computable} is exact mathematics and useless
numerics unless the residual is formed in extended precision.  This is a
property of the bound, not of the expansion: the expansion itself is
evaluated by a stable Horner recurrence.  One source advises ``sufficient
guard precision'' when evaluating the coordinates; none quantifies the
cancellation in the residual itself, although it is that cancellation, and
not the coordinates, that decides how many digits are needed.
\end{remark}

\begin{remark}[The error is not monotone in \(N\)]
\label{plt:rmk:an-nonmonotone}
At \(x=100\) the errors in \cref{plt:tab:an-worked} decrease at every step,
but this is a feature of that value of \(x\), not a theorem.  At \(x=10\),
where \(\WrightOmega(10)=7.929420095019697348562\ldots\), the successive
errors \(\WrightOmega(10)-\WrightOmega_{N}(10)\) are
\[
 \begin{array}{llll}
 N=4:&\ 1.9026\times10^{-5}, & N=5:&\ 5.1336\times10^{-6},\\
 N=6:&\ 3.8679\times10^{-8}, & N=7:&\ -1.2836\times10^{-7},\\
 N=11:&\ 4.8696\times10^{-12}, & N=12:&\ -1.9269\times10^{-11},
 \end{array}
\]
so the absolute error increases from \(N=6\) to \(N=7\) and again from
\(N=11\) to \(N=12\).  This is the reason to truncate by residual size rather
than by formal order when a numerical answer is wanted.
\end{remark}

% ed.: S1 asserts, of its own error table at x = 10, that "the N = 2
% ed.: truncation is slightly worse than N = 1".  Its table gives
% ed.: 1.7466787e-3 at N = 1 and 1.7369609e-3 at N = 2, which is smaller, and
% ed.: recomputation at 40 digits confirms both table entries.  The prose
% ed.: contradicts the table it is describing.  What is true at x = 10 is that
% ed.: the N = 2 correction almost exactly cancels the N = 1 error and flips
% ed.: its sign, improving it by only 0.6 per cent, and that genuine
% ed.: non-monotonicity occurs later, at N = 6 -> 7 and N = 11 -> 12, as
% ed.: recorded above.  S3's table at log z = 2.3026 makes the analogous claim
% ed.: correctly, for the pair N = 2, 3.
\section{An example where the hypotheses fail}
\label{plt:sec:an-failure}

\begin{example}[The Lambert branch point: a residual of \(10^{-8}\) with an
error of \(2.3\times10^{-4}\)]
\label{plt:ex:an-branchpoint}
Let \(\varphi(Y)=Y\EulerE^{Y}\) on \(\RealNumbers\), so that
\(\varphi'(Y)=(1+Y)\EulerE^{Y}\) and \(\varphi'(-1)=0\); the two real
solutions of \(\varphi(Y)=z\) for \(-\EulerE^{-1}<z<0\) are
\(\LambertW_{0}(z)>-1\) and \(\LambertW_{-1}(z)<-1\)
\cite{corless-et-al-1996,dlmf-W}.  Take
\[
 \eta=10^{-8},
 \qquad
 z=-\EulerE^{-1}+\eta=-0.36787943117144232159552\ldots,
 \qquad
 Y_{\mathrm{app}}=-1 .
\]
The residual is exactly
\[
 \varphi(Y_{\mathrm{app}})-z=-\EulerE^{-1}-z=-\eta=-10^{-8},
\]
as small as one could wish.  The two true solutions are
\[
 \LambertW_{0}(z)=-0.99976685372178274818\ldots,
 \qquad
 \LambertW_{-1}(z)=-1.00023318252197543531\ldots,
\]
so that the actual error of \(Y_{\mathrm{app}}\) is
\(2.3314627822\times10^{-4}\) or \(2.3318252198\times10^{-4}\) according to
which root is meant --- larger than the residual by a factor of about
\(2.33\times10^{4}\).

Every hypothesis fails, and each failure is instructive.
\begin{enumerate}
\item \emph{No admissible \(m\).}  On the interval joining
\(Y_{\mathrm{app}}=-1\) to either root, \(\inf\AbsoluteValue{\varphi'}=0\),
attained at the endpoint \(-1\).  So there is no \(m>0\) satisfying the
hypothesis of \cref{plt:prop:rev-analytic-transfer}, and
\eqref{plt:eq:an-transfer} is not merely weak, it is unavailable.
\item \emph{No enclosure.}  In \cref{plt:thm:an-enclosure} one needs
\(m=\min_{I}\AbsoluteValue{\varphi'}>0\) on an interval \(I\) around
\(Y_{\mathrm{app}}\); any such \(I\) contains \(-1\), where \(\varphi'\)
vanishes, so \(m=0\) and \eqref{plt:eq:an-enclosure-hyp} cannot be satisfied
for any \(\rho\).
\item \emph{No uniqueness.}  The residual cannot distinguish the two roots:
both lie within \(2.34\times10^{-4}\) of \(Y_{\mathrm{app}}\), and the single
number \(-10^{-8}\) certifies neither.  A residual bound is a statement about
\emph{a} root, and at a critical point ``a root'' is not ``the root''.
\item \emph{The true law is a square root.}  Since \(\varphi(-1)=-\EulerE^{-1}\),
\(\varphi'(-1)=0\) and \(\varphi''(-1)=\EulerE^{-1}\), Taylor's theorem gives
\(\varphi(Y)+\EulerE^{-1}=\tfrac12\EulerE^{-1}(Y+1)^{2}
+\BigOAt{(Y+1)^{3}}{Y\to-1}\), so the error obeys
\(\AbsoluteValue{Y+1}\approx\sqrt{2\EulerE\,\AbsoluteValue{R}}\).  With
\(\AbsoluteValue{R}=\eta=10^{-8}\) this predicts
\(\sqrt{2\EulerE\eta}=2.3316439816\times10^{-4}\), which matches both true
errors to four significant figures.  A linear residual bound with any finite
constant is false here for all small \(\eta\), since
\(\sqrt{\eta}/\eta\to\infty\).
\end{enumerate}
The repair is a change of chart, not a better constant: in the variable
\(p=\sqrt{2\EulerE(z+\EulerE^{-1})}\) the two branches are analytic and the
transfer machinery applies again.  That is the analytic reason the branch
point carries a Puiseux, not a polynomial--logarithmic, expansion.
\end{example}

\begin{remark}[The three ways the transfer fails, in general]
\label{plt:rmk:an-failure-modes}
\Cref{plt:ex:an-branchpoint} exhibits the first and worst of three failure
modes.
\begin{enumerate}
\item \emph{A critical point of the map.}  \(\varphi'\) vanishes between the
approximation and the root.  The residual then controls only a fractional
power of the error, with the exponent \(1/k\) at a zero of \(\varphi'\) of
order \(k-1\).  No choice of constants rescues the linear bound.
\item \emph{A small derivative without a critical point.}  \(\varphi'\) is
bounded below by \(m>0\), but \(m\) is tiny.  Then
\eqref{plt:eq:an-enclosure} is true and useless: the error bound is inflated
by \(m^{-1}\), and the enclosure hypothesis
\eqref{plt:eq:an-enclosure-hyp} demands a residual smaller than \(m\rho\).
This is the situation for \(\LambertW_{-1}\) at moderate arguments, where
\(\varphi'\) is small but nonzero.
\item \emph{An unverifiable analytic residual.}  (H2) of
\cref{plt:thm:an-inverse-transfer} fails or is not checked: the formal
residual has high \(t\)-order and nothing is known about the size of
\(\varphi(g_{N}(x))-x\).  This mode is silent --- no computation goes wrong,
no inequality is violated, and the resulting statement is simply unproved.
It is the mode that occurs in the sources.
\end{enumerate}
\end{remark}

\begin{summarybox}
A residual is converted into an error bound by the mean value theorem, and
that step is elementary.  Everything else is not.  The derivative lower bound
must hold on an interval that contains the unknown root, which is circular
unless one instead \emph{encloses} the root
(\cref{plt:thm:an-enclosure}, and \cref{plt:thm:an-complex-enclosure} by
Rouch\'e in the complex case).  The analytic size of the residual does not
follow from its formal \(t\)-order; it follows from the formal order
\emph{together with} an analytic expansion of the map being reversed
(\cref{plt:thm:an-residual-analytic}), and with that in hand the formal
inverse is the asymptotic expansion of the analytic inverse
(\cref{plt:thm:an-inverse-transfer}).  For \(\varphi(Y)=Y+\log Y\) all
constants are explicit and the enclosure
% ed.: "with relative width 1/(x-\log x)": the exact relative width is
% ed.: (1+x-\log x)^{-1}; the quoted quantity is an upper bound for it, so the
% ed.: sentence is corrected to say so.
\eqref{plt:eq:an-omega-bracket} holds at every single \(x>1\), with relative
width at most \(1/(x-\log x)\); at \(x=100\), \(N=3\) it brackets the true error
\(8.153845\times10^{-8}\) between \(8.153804\times10^{-8}\) and
\(8.239278\times10^{-8}\).  Where the derivative vanishes --- the Lambert
branch point --- a residual of \(10^{-8}\) coexists with an error of
\(2.3\times10^{-4}\), and no constant repairs the bound.
\end{summarybox}

% ---- fragment 13-dulac ----
\chapter{Power--log series near a finite point and the Dulac conventions}
\label{plt:sec:dulac}

Every result proved so far has been stated for \(X\to+\infty\).  The literature
that motivates the class --- local dynamics of planar vector fields, normal
forms of parabolic and hyperbolic germs, Dulac's non-accumulation problem ---
works instead at a finite point, almost always \(x\to0^{+}\), and with
generators, signs, orderings and truncation conventions that differ from ours
in every one of those four respects.  The purpose of this chapter is to make
the passage exact.

\Cref{plt:prop:grp-zero-dictionary} already recorded the integer-exponent
dictionary \(x=X^{-1}\), \(\Lambda=\log(1/x)=L\), and
\cref{plt:prop:grp-shift} recorded the level shift \(F\mapsto F\compose\log\).
Those two propositions are the input here; the remark closing
\cref{plt:sec:scale-change} deferred to the present chapter three things they
do not supply:

\begin{enumerate}[label=\textup{(\alph*)}]
\item a statement of \emph{which} class each of the two natural substitutions
  produces, and a proof that the two are not interchangeable;
\item the Dulac conventions themselves --- what a Dulac series is, what a Dulac
  germ is, how the support condition compares with the one that defines
  \(\PLlau\), and where the two genuinely differ --- together with the reason
  the first-return map of a hyperbolic polycycle lands in the class;
\item the enlargements a finite point forces --- real exponents in \(x\), real
  exponents in the logarithm, and the two distinct mechanisms by which
  finiteness of the logarithmic degree fails.
\end{enumerate}

The contribution of this chapter is precision rather than machinery.  The
sources are thin on proofs here and thick on convention: S2 devotes a chapter
to the small-variable dictionary and states its results without proof; S1 and
S6 each give half a page.  Worse, S1, S2 and S6 do not use the same letter for
the same object, and the discrepancy is invisible in any single formula.  We
therefore state definitions, prove transport theorems, and record as
\emph{checkable} propositions the eight ways a transported statement goes
wrong.

\section{Zero-side generators, and the two substitutions}
\label{plt:sec:du-generators}

\begin{definition}[Local coordinate and zero-side generators]
\label{plt:def:du-generators}
Let \(a\in\RealNumbers\).  A \emph{local coordinate at \(a\) from the right} is
\(x=z-a\); from the left it is \(x=a-z\).  In either case the regime is
\(x\to0^{+}\), and we work on an interval \(0<x<x_{0}\) with
\(x_{0}<\EulerE^{-1}\), so that the two zero-side generators
\begin{equation}\label{plt:eq:du-generators}
  x,
  \qquad
  \Lambda\DefinitionEquals\log\frac1x=-\log x
\end{equation}
satisfy \(0<x<\EulerE^{-1}<1<\Lambda\) and \(\Lambda\to+\infty\).  We write
\(\Lambda_{2}\DefinitionEquals\log\Lambda\) for the second zero-side logarithm
and
\begin{equation}\label{plt:eq:du-small-log}
  \ell\DefinitionEquals\Lambda^{-1}=-\frac1{\log x}\in(0,1),
  \qquad \ell\to0^{+},
\end{equation}
for the \emph{small} logarithmic generator used in the normal-form literature.
In the formal algebra \(x\) and \(\Lambda\) are independent indeterminates, and
the four zero-side classes are
\[
  \K[\Lambda][[x]],
  \qquad
  \K[\Lambda]((x)),
  \qquad
  \K(\Lambda)((x)),
  \qquad
  \K((\Lambda^{-1}))((x))=\K((\ell))((x)),
\]
carrying the \(x\)-adic valuation \(\nu_{x}\), the exclusive truncation
\(\Trunc{f}{N}=\sum_{n<N}p_{n}(\Lambda)x^{n}\), and the derivation determined
by
\begin{equation}\label{plt:eq:du-derivation}
  \frac{\Differential x}{\Differential x}=1,
  \qquad
  \frac{\Differential\Lambda}{\Differential x}=-\frac1x .
\end{equation}
% ed.: this sentence originally attributed the derivation to all four
% classes.  Unlike its mirror \PLpow at infinity, the power-series class
% K[Lambda][[x]] is NOT stable under d/dx: d(Lambda)/dx = -1/x already leaves
% it.  This is trap (T6) in its sharpest form, and it is why the near-identity
% class of plt:def:du-near-identity is described inside K[Lambda]((x)).
The last three of these classes are stable under
\(\Differential/\Differential x\); the first is not, since
\(\Differential\Lambda/\Differential x=-x^{-1}\notin\K[\Lambda][[x]]\).  This
is the first asymmetry with the situation at infinity, where \(\PLpow\) is
\(\DX\)-stable, and it recurs as trap~\textup{(T6)} of
\cref{plt:prop:du-traps}.
For a complex germ one replaces the interval by a sector
\(\SetOf{0<\AbsoluteValue{x}<x_{0},\ \AbsoluteValue{\arg x}<\theta}\) and
fixes a branch of \(\log\); every statement below that mentions positivity or a
real branch is then a hypothesis, not a consequence.
\end{definition}

\begin{warningbox}[title={The letter \(\ell\) means two reciprocal things}]
\label{plt:warn:du-ell}
% ed.: S2 sets ell = log(1/x); S1 and S6, following the normal-form
% literature, set ell = -1/log x.  The two are reciprocal, and no single
% displayed formula distinguishes them.  We fix Lambda for the large log and
% never write a bare ell for it.
Source~S2 writes \(\ell=\log(1/x)\); sources~S1 and~S6, following
\cite{mardesic-normalforms,mardesic-length,resman-dulac,peran-logtransseries},
write \(\ell=-1/\log x\).  These are reciprocal:
\(\ell_{\mathrm{S2}}=\Lambda\to+\infty\) while
\(\ell_{\mathrm{S1,S6}}=\Lambda^{-1}\to0^{+}\).  A monomial \(x^{\alpha}\ell^{m}\)
therefore means \(x^{\alpha}\Lambda^{m}\) in one convention and
\(x^{\alpha}\Lambda^{-m}\) in the other; the exponent of the logarithm changes
sign, and with it the direction in which the logarithmic part of the dominance
order runs.  Throughout this volume the large zero-side logarithm is
\(\Lambda\) and the small one is written \(\Lambda^{-1}\) or, where the
comparison with the normal-form literature is the point, \(\ell\) with
\eqref{plt:eq:du-small-log} recalled on the spot.
\end{warningbox}

\begin{definition}[The reciprocal and the exponential substitution]
\label{plt:def:du-substitutions}
Let \(\mathcal M_{2}(\K,\RealNumbers)\) and
\(\mathcal H_{2}(\K,\RealNumbers)\) be the depth-two monomial group and Hahn
field of \cref{plt:def:grp-logtower}.  Write
\(\mathcal Z\DefinitionEquals\iota\bigl(\mathcal H_{2}(\K,\RealNumbers)\bigr)\)
for its zero-side copy under the relabelling
\begin{equation}\label{plt:eq:du-iota}
  \iota:\quad
  L_{0}=X\longmapsto x^{-1},
  \qquad
  L_{1}\longmapsto\Lambda,
  \qquad
  L_{2}\longmapsto\Lambda_{2},
\end{equation}
so that a monomial of \(\mathcal Z\) is \(x^{a}\Lambda^{b}\Lambda_{2}^{c}\)
with \((a,b,c)\in\RealNumbers^{3}\), and
\begin{equation}\label{plt:eq:du-zero-order}
  x^{a}\Lambda^{b}\Lambda_{2}^{c}
  \succdom
  x^{a'}\Lambda^{b'}\Lambda_{2}^{c'}
  \iff
  a<a',
  \text{ or }(a=a',\ b>b'),
  \text{ or }(a=a',\ b=b',\ c>c').
\end{equation}
This is \eqref{plt:eq:du-iota} applied to the order of
\cref{plt:def:grp-logtower}, since the exponent of \(L_{0}=x^{-1}\) is
\(-a\).  The two
substitutions are the \(\K\)-algebra embeddings of \(\PLlau\) into
\(\mathcal Z\) determined on generators by
\begin{align}
  \label{plt:eq:du-sigma}
  \sigma:&\qquad t\longmapsto x, &&L\longmapsto\Lambda
  &&\text{(\emph{reciprocal}: \(X=1/x\)),}\\
  \label{plt:eq:du-tau}
  \tau:&\qquad t\longmapsto\Lambda^{-1}, &&L\longmapsto\Lambda_{2}
  &&\text{(\emph{exponential}: \(X=\log(1/x)\), i.e.\ \(x=\EulerE^{-X}\)).}
\end{align}
On germs, \(\sigma\) is \(F\mapsto F(1/x)\), which is \(\iota\) restricted to
\(\PLlau\); and \(\tau\) is \(F\mapsto F(\log(1/x))\), which is
\(\iota\compose\Sigma\) with \(\Sigma=(\,\cdot\,)\compose\log\) the level shift
of \cref{plt:prop:grp-shift}.  Throughout this section \(\sigma\) and \(\tau\)
denote these two substitutions and nothing else; elsewhere in the volume
\(\sigma\) is a logarithmic exponent
(\cref{plt:thm:grp-scale-classification}) and \(\tau\) is the Taylor-shift
homomorphism used in the proof of \cref{plt:thm:lag-burmann}, whose role here
is played by \(\theta\).
\end{definition}

\begin{theorem}[What each substitution does to each class]
\label{plt:thm:du-two-substitutions}
Let \(\K\) have characteristic zero.
\begin{enumerate}[label=\textup{(\roman*)}]
\item \(\sigma\) is an isomorphism of ordered rings onto its image and
  restricts to isomorphisms
  \[
    \PLpow\to\K[\Lambda][[x]],
    \qquad
    \PLlau\to\K[\Lambda]((x)),
  \]
  \[
    \PLrat\to\K(\Lambda)((x)),
    \qquad
    \PLit\to\K((\Lambda^{-1}))((x)),
  \]
  and, for every \(s\in\PositiveIntegers\) and every ordered subgroup
  \(\Gamma\subseteq\RealNumbers\),
  \(\K[L]((t^{1/s}))\to\K[\Lambda]((x^{1/s}))\) and
  \(\mathcal H_{1}(\K,\Gamma)\to\iota\bigl(\mathcal H_{1}(\K,\Gamma)\bigr)\).
  It carries \(\ordt\) to \(\nu_{x}\), \(\Trunc{\cdot}{N}\) to
  \(\Trunc{\cdot}{N}\), \(\degL\) to \(\deg_{\Lambda}\), and the dominance
  order of \cref{plt:def:ring-dominance} to the order of decreasing eventual
  size as \(x\to0^{+}\).  It does \emph{not} carry \(\DX\) to
  \(\Differential/\Differential x\); instead
  \begin{equation}\label{plt:eq:du-sigma-derivation}
    \sigma\compose\DX
    =\Bigl(-x^{2}\frac{\Differential}{\Differential x}\Bigr)\compose\sigma .
  \end{equation}
\item \(\tau\) is an isomorphism of ordered rings onto
  \(\K[\Lambda_{2}]((\Lambda^{-1}))\), carrying \(\ordt\) to the
  \(\Lambda^{-1}\)-adic valuation and dominance to dominance.  It intertwines
  \(\DX\) with the derivation \(D_{\Lambda}\) of the target determined by
  \(D_{\Lambda}\Lambda=1\), \(D_{\Lambda}\Lambda_{2}=\Lambda^{-1}\); it does
  \emph{not} intertwine \(\DX\) with \(\Differential/\Differential x\), and
  \(\K[\Lambda_{2}]((\Lambda^{-1}))\) is not stable under
  \(\Differential/\Differential x\), because
  \begin{equation}\label{plt:eq:du-tau-derivation}
    \frac{\Differential}{\Differential x}=-\frac1x\,D_{\Lambda},
    \qquad
    \frac{\Differential}{\Differential x}\Lambda^{-1}=\frac1{x\Lambda^{2}} .
  \end{equation}
\item Inside \(\mathcal Z\),
  \begin{equation}\label{plt:eq:du-intersection}
    \sigma(\PLlau)\cap\tau(\PLlau)=\K[\Lambda] .
  \end{equation}
  More precisely, in the exponent coordinates \((a,b,c)\) of
  \eqref{plt:eq:du-zero-order}, \(\sigma(\PLlau)\) consists of the elements of
  \(\mathcal Z\) whose support lies in
  \(\IntegerNumbers\times\NaturalNumbers\times\SetOf0\) with bounded-below
  first coordinate, while \(\tau(\PLlau)\) consists of those whose support lies
  in \(\SetOf0\times\IntegerNumbers\times\NaturalNumbers\) with bounded-above
  second coordinate.
\end{enumerate}
\end{theorem}

\begin{proof}
(i)  The map \(\sigma\) is the relabelling \(t\mapsto x\), \(L\mapsto\Lambda\)
of two algebraically independent generators, extended continuously; it is
therefore a ring isomorphism onto the correspondingly relabelled ring, and
each of the six displayed restrictions is the same relabelling applied to the
same support condition.  Since \(\sigma\) does not move the exponent pair
\((n,j)\) of a monomial, it commutes with the operators
\(\ordt\mapsto\nu_{x}\), \(\Trunc{\cdot}{N}\), \(\degL\mapsto\deg_{\Lambda}\),
which are defined from that pair alone; and it preserves the order because
\eqref{plt:eq:ring-dominance} and its zero-side counterpart are governed by the
same inequalities on \((n,j)\).  That the zero-side order is the order of
decreasing eventual size is \cref{plt:lem:ring-dominance-analytic} composed
with \(X=1/x\).  Finally \eqref{plt:eq:du-sigma-derivation} is
\eqref{plt:eq:grp-zero-derivation}, restated as an identity of operators.

(ii)  Composing \cref{plt:prop:grp-shift} with the relabelling
\eqref{plt:eq:du-iota} gives at once that \(\tau\) is an isomorphism of ordered
rings of \(\PLlau=\K[L_{1}]((L_{0}^{-1}))\) onto
\(\K[L_{2}]((L_{1}^{-1}))=\K[\Lambda_{2}]((\Lambda^{-1}))\) with the stated
transport of valuation and order, and \(\Sigma\compose\DX=D_{L_{1}}\compose
\Sigma\) becomes \(\tau\compose\DX=D_{\Lambda}\compose\tau\).  For the negative
statement, the chain rule and \eqref{plt:eq:du-derivation} give
\(\Differential/\Differential x=(\Differential\Lambda/\Differential x)
D_{\Lambda}=-x^{-1}D_{\Lambda}\), whence
\(\Differential\Lambda^{-1}/\Differential x
=-x^{-1}\cdot(-\Lambda^{-2})=(x\Lambda^{2})^{-1}\), which involves \(x\) and
therefore does not lie in \(\K[\Lambda_{2}]((\Lambda^{-1}))\); this is
\eqref{plt:eq:grp-shift-derivation} written at zero.

(iii)  Both images lie in \(\mathcal Z\) and Hahn representations are unique,
so it suffices to identify supports, which we record in the coordinates
\((a,b,c)\) of \eqref{plt:eq:du-zero-order}.  An element of \(\PLlau\) is
\(\sum_{n\ge n_{0}}p_{n}(L)t^{n}\); applying \(\sigma\) gives support
\(\SetOf{(n,j,0)}\) with \(n\ge n_{0}\) and \(0\le j\le\deg p_{n}\), while
applying \(\tau\) gives support \(\SetOf{(0,-n,j)}\) with the same ranges.  An
element of the intersection has support in
\(\SetOf0\times\IntegerNumbers\times\SetOf0\), i.e.\ it is a series in
\(\Lambda\) alone; lying in \(\sigma(\PLlau)\) makes it a \emph{polynomial} in
\(\Lambda\), and every polynomial in \(\Lambda\) does lie in
\(\tau(\PLlau)=\K[\Lambda_{2}]((\Lambda^{-1}))\), since that ring admits
finitely many negative powers of its outer generator \(\Lambda^{-1}\).  Hence
the intersection is exactly \(\K[\Lambda]\).
\end{proof}

\begin{remark}[The transport theorem is deliberately trivial as algebra]
\label{plt:rmk:du-trivial}
Item~(i) is a relabelling and nothing more; this is the point.  All of the
difficulty in passing between infinity and a finite point is concentrated in
the two places where the relabelling is \emph{not} available: the derivation,
by \eqref{plt:eq:du-sigma-derivation}, and the analytic meaning of a monomial,
by \cref{plt:def:mul-asymptotic} versus \cref{plt:def:du-asymptotic} below.  A
transported statement that mentions neither is automatically correct; a
transported statement that mentions either must be re-derived.  The eight traps
of \cref{plt:prop:du-traps} are all instances of this one sentence.
\end{remark}

The two substitutions are not alternatives between which one may choose
freely; they produce classes that are, apart from the polynomials in
\(\Lambda\), disjoint.  The next proposition is the precise obstruction, and it
also supplies the practical criterion.

\begin{proposition}[The two scales at a finite point are incomparable]
\label{plt:prop:du-scale-separation}
Work on \(0<x<\EulerE^{-1}\) with the real branches, so that
\(\Lambda>1\) and \(\Lambda_{2}>0\).
\begin{enumerate}[label=\textup{(\roman*)}]
% ed.: this item said "every nonzero f in tau(PLlau), evaluated at the
% generators".  A general element of K[Lambda_2]((Lambda^{-1})) is a formal,
% typically divergent, series and has no evaluation; and the proof uses only
% asymptotic equivalence to the leading monomial.  Restated for functions.
\item Let \(F\in\PLlau\) be nonzero, with leading monomial \(c\,t^{n}L^{j}\),
  and let \(f\) be a function on \((0,x_{0})\) that realizes \(\tau(F)\) at
  leading order, that is, \(f\sim c\,\Lambda^{-n}\Lambda_{2}^{\,j}\) as
  \(x\to0^{+}\); this holds in particular whenever \(f\approx\tau(F)\) in the
  sense of \cref{plt:def:mul-asymptotic} read in the variable \(X=\Lambda\).
  Then
  \[
    \log\AbsoluteValue{f(x)}=\BigOAt{\Lambda_{2}}{x\to0^{+}}
    =\LittleOAt{\Lambda}{x\to0^{+}} ;
  \]
  in particular, for every \(\varepsilon>0\),
  \(x^{\varepsilon}\precdom f\precdom x^{-\varepsilon}\).
\item \(\Lambda_{2}=\log\log(1/x)\) is not asymptotically equivalent to
  \(c\,x^{\alpha}\Lambda^{\beta}\) for any \(c\ne0\) and any
  \(\alpha,\beta\in\RealNumbers\).  Consequently \(\Lambda_{2}\) admits no
  asymptotic expansion in the scale
  \(\SetOf{x^{\alpha}\Lambda^{\beta}:\alpha,\beta\in\RealNumbers}\), and in
  particular lies in none of the four \(\sigma\)-images, nor in
  \(\sigma\bigl(\mathcal H_{1}(\K,\RealNumbers)\bigr)\) realized on that scale.
\item Symmetrically, \(x\) is not asymptotically equivalent to
  \(c\,\Lambda^{\beta}\Lambda_{2}^{\gamma}\) for any \(c\ne0\) and any
  \(\beta,\gamma\in\RealNumbers\); so \(x\) has no expansion in the scale
  produced by \(\tau\).
\end{enumerate}
\end{proposition}

\begin{proof}
(i)  By hypothesis \(f/(c\,\Lambda^{-n}\Lambda_{2}^{\,j})\to1\).  That
\(f\approx\tau(F)\) implies this is the estimate
\eqref{plt:eq:mul-asymptotic} at the cutoff \(N=n+1\), which reads
\(\AbsoluteValue{f-p_{n}(\Lambda_{2})\Lambda^{-n}}\le
C\Lambda^{-(n+1)}\Lambda_{2}^{M}\), combined with
\(p_{n}(\Lambda_{2})\sim c\,\Lambda_{2}^{\,j}\) and
\(\Lambda^{-1}\Lambda_{2}^{M-j}\to0\).  Hence
\[
  \log\AbsoluteValue f
  =-n\Lambda_{2}+j\log\Lambda_{2}+\BigOAt{1}{x\to0^{+}}
  =\BigOAt{\Lambda_{2}}{x\to0^{+}},
\]
and \(\Lambda_{2}=\LittleOAt{\Lambda}{x\to0^{+}}\) with \(\log x=-\Lambda\)
gives \(\log\AbsoluteValue f/\log x\to0\), which is the second assertion.

(ii)  Suppose \(\Lambda_{2}\sim c\,x^{\alpha}\Lambda^{\beta}\).  Taking
logarithms,
\[
  \log\Lambda_{2}
  =\log\AbsoluteValue c-\alpha\Lambda+\beta\Lambda_{2}
   +\LittleOAt{1}{x\to0^{+}} .
\]
The left side is \(\LittleOAt{\Lambda_{2}}{x\to0^{+}}\), hence
\(\LittleOAt{\Lambda}{x\to0^{+}}\).  If \(\alpha\ne0\) the right side is
\(\asymp\Lambda\), a contradiction; so \(\alpha=0\).  Then the right side is
\(\beta\Lambda_{2}+\BigOAt{1}{x\to0^{+}}\) while the left side is
\(\LittleOAt{\Lambda_{2}}{x\to0^{+}}\), forcing \(\beta=0\); but then the right
side is bounded and the left side tends to \(+\infty\).  For the consequence:
if \(\Lambda_{2}\) had an expansion in that scale with a nonzero leading term,
that term would be such a \(c\,x^{\alpha}\Lambda^{\beta}\), and if the
expansion were \(0\) then \(\Lambda_{2}=\BigOAt{x^{N}}{x\to0^{+}}\) for some
\(N>0\), which is false.

(iii)  Suppose \(x\sim c\,\Lambda^{\beta}\Lambda_{2}^{\gamma}\).  Taking
logarithms, \(-\Lambda=\log\AbsoluteValue c+\beta\Lambda_{2}
+\gamma\log\Lambda_{2}+\LittleOAt{1}{x\to0^{+}}\).  The right side is
\(\BigOAt{\Lambda_{2}}{x\to0^{+}}=\LittleOAt{\Lambda}{x\to0^{+}}\) while the
left side is \(-\Lambda\), a contradiction.
\end{proof}

\begin{remark}[Which substitution to use]
\label{plt:rmk:du-choice}
\Cref{plt:prop:du-scale-separation} turns the choice into a measurement.  Given
a germ \(f\) at \(0^{+}\), form \(\log\AbsoluteValue{f(x)}/\log x\).
\begin{itemize}
\item If it tends to a finite limit \(\alpha\) and the corrections to
  \(f/x^{\alpha}\) are powers of \(x\) times polynomials in \(\Lambda\), the
  reciprocal substitution \(\sigma\) is the right one: \(f\) is a power--log
  germ and \(f(1/X)\) is a candidate element of \(\PLlau\) or one of its
  enlargements.
\item If it tends to \(0\) while \(f\) itself is unbounded or tends to \(0\),
  the germ lives on the logarithmic scale and \(\sigma\) is useless: by
  \cref{plt:prop:du-scale-separation}(ii) even the depth-one Hahn class in
  \(x\) and \(\Lambda\) cannot express \(\Lambda_{2}\).  Use \(\tau\), i.e.\
  take \(X=\log(1/x)\) as the large variable.
\item If it tends to \(\pm\infty\), neither substitution applies: the germ is
  beyond every logarithmic class, by \cref{plt:prop:grp-exponential} read
  through \(\iota\).
\end{itemize}
\Cref{plt:ex:du-lower-branch} is the standard germ for which the second case
occurs, and for which the first case fails in an instructive way.
\end{remark}

\begin{example}[The lower Lambert branch at a finite point]
\label{plt:ex:du-lower-branch}
For \(0<x<\EulerE^{-1}\) the branch \(\LambertW_{-1}(-x)\) is the unique
solution \(w<-1\) of \(w\EulerE^{w}=-x\).  Writing \(w=-U\) with \(U>1\), the
equation becomes \(U\EulerE^{-U}=x\), i.e.
\begin{equation}\label{plt:eq:du-lower-branch-equation}
  U-\log U=\Lambda .
\end{equation}
The map \(g(u)=u-\log u\) has \(g'(u)=1-1/u>0\) on \((1,+\infty)\), with
\(g(1)=1\) and \(g(u)\to+\infty\); it is therefore an increasing bijection
\((1,+\infty)\to(1,+\infty)\), and \(U=\CompositionalInverseOf{g}(\Lambda)\) is
well defined for \(\Lambda>1\), i.e.\ for \(0<x<\EulerE^{-1}\).

Now apply \(\tau\).  In the variables \(X=\Lambda\), \(L=\Lambda_{2}\), the
equation \eqref{plt:eq:du-lower-branch-equation} is exactly the reversion of
\(X\mapsto X-L\), which is the member \(a=-1\) of the family of
\cref{plt:thm:om-alpha-family}.  Substituting \eqref{plt:eq:om-minus-log} gives
the expansion --- with \(\approx\) the relation of
\cref{plt:def:mul-asymptotic} read in the large variable \(X=\Lambda\), not the
zero-side relation \(\approx_{0}\) of \cref{plt:def:du-asymptotic}, whose scale
is powers of \(x\) ---
\begin{equation}\label{plt:eq:du-lower-branch}
  \LambertW_{-1}(-x)
  \;\approx\;
  -\Lambda-\Lambda_{2}
   -\sum_{n\ge1}(-1)^{n+1}\frac{\LamP{n}(\Lambda_{2})}{\Lambda^{n}}
  =-\Lambda-\Lambda_{2}-\frac{\Lambda_{2}}{\Lambda}
   +\frac{\Lambda_{2}^{2}/2-\Lambda_{2}}{\Lambda^{2}}-\cdots,
\end{equation}
whose analytic validity is \cref{plt:prop:du-lower-branch-analytic}.  Two
points deserve emphasis.  First, the expansion uses \emph{no} power of \(x\):
it lies entirely in \(\tau(\PLlau)\), consistently with
\(\log\AbsoluteValue{\LambertW_{-1}(-x)}/\log x\to0\).  Second, the reciprocal
substitution genuinely fails here: by
\cref{plt:prop:du-scale-separation}(ii) the term \(\Lambda_{2}\) already lies
outside every \(\sigma\)-image, so no amount of enlarging the coefficient field
in \(\Lambda\) will produce \eqref{plt:eq:du-lower-branch}.  A second logarithm
at infinity would be needed --- and that is precisely the level shift that
\(\tau\) performs once and for all.
\end{example}

\begin{proposition}[Asymptotic validity of the lower-branch expansion]
\label{plt:prop:du-lower-branch-analytic}
% ed.: S6 displays the expansion with a bare remainder O(M^3/S^3) and no
% argument.  The estimate below is proved; the sharp form of S6's remainder is
% recovered in the last sentence.
For \(N\in\NaturalNumbers\) put
\[
  V_{N}(\Lambda)
  \DefinitionEquals
  \Lambda+\Lambda_{2}
   +\sum_{n=1}^{N}(-1)^{n+1}\frac{\LamP{n}(\Lambda_{2})}{\Lambda^{n}} .
\]
Then there are \(K_{N}\in\NaturalNumbers\), \(C_{N}>0\) and
\(\Lambda_{N}>1\) such that, with \(U=-\LambertW_{-1}(-x)\) and
\(\Lambda=\log(1/x)\),
\begin{equation}\label{plt:eq:du-lower-branch-error}
  \AbsoluteValue{U-V_{N}(\Lambda)}
  \le C_{N}\,\frac{\Lambda_{2}^{K_{N}}}{\Lambda^{N+1}}
  \qquad(\Lambda\ge\Lambda_{N}).
\end{equation}
% ed.: "asymptotic expansion on the scale of all monomials Lambda^{-n}
% Lambda_2^j" would presuppose that this doubly indexed family is a Poincare
% scale, which it is not (plt:prop:an-no-enumeration).  The precise assertion
% is the relation of plt:def:mul-asymptotic read in the variable X=Lambda,
% which is exactly what plt:eq:du-lower-branch-error states.
In particular \eqref{plt:eq:du-lower-branch} holds in the sense of
\cref{plt:def:mul-asymptotic} read in the large variable \(X=\Lambda\), with
\(L=\Lambda_{2}\): the estimate \eqref{plt:eq:du-lower-branch-error} is
\eqref{plt:eq:mul-asymptotic} at the cutoff \(N+1\), and the cutoffs \(N+1\),
\(N\in\NaturalNumbers\), are cofinal, so
\cref{plt:lem:mul-asymptotic-basic}(1) supplies the remaining ones.
Concretely,
\[
  \LambertW_{-1}(-x)
  =-\Lambda-\Lambda_{2}-\frac{\Lambda_{2}}{\Lambda}
   +\frac{\Lambda_{2}^{2}/2-\Lambda_{2}}{\Lambda^{2}}
   +\BigOAt{\frac{\Lambda_{2}^{3}}{\Lambda^{3}}}{x\to0^{+}} .
\]
\end{proposition}

\begin{proof}
Throughout, \(g(u)=u-\log u\) and \(U=\CompositionalInverseOf{g}(\Lambda)\) as
in \cref{plt:ex:du-lower-branch}.

\emph{Step 1: an a priori bound.}  From \(U=\Lambda+\log U\) and \(U>1\) we get
\(U>\Lambda\).  For \(u\ge8\) one has \(\log u\le u/2\), hence
\(g(u)\ge u/2\); so if \(\Lambda\ge4\) and \(u>2\Lambda\ge8\) then
\(g(u)>\Lambda\), and monotonicity of \(g\) gives \(U\le2\Lambda\).  Therefore
\(\Lambda<U\le2\Lambda\) and \(\Lambda_{2}<\log U\le\Lambda_{2}+\log2\) for
\(\Lambda\ge4\).

\emph{Step 2: the formal identity.}  Work in \(\K[L][[t]]\) and put
\(\Psi\DefinitionEquals L+\sum_{n\ge1}(-1)^{n+1}\LamP{n}(L)t^{n}\), so that
\eqref{plt:eq:om-minus-log} reads \(\rev{(X-L)}=X+\Psi\).  Applying
\(X-L\) to \(X+\Psi\) and using generator substitution
(\cref{plt:prop:cmp-generator-substitution}),
\(L\compose(X+\Psi)=L+\log(1+t\Psi)\), the defining relation becomes
\begin{equation}\label{plt:eq:du-psi-fixedpoint}
  \Psi=L+\log(1+t\Psi).
\end{equation}
Let \(\Psi_{N}\DefinitionEquals\Trunc{\Psi}{N+1}\).  Since
\(\ordt(\Psi-\Psi_{N})\ge N+1\) and \(\log(1+t\,\cdot\,)\) raises the outer
order by at least one, \(\ordt\bigl(\log(1+t\Psi)-\log(1+t\Psi_{N})\bigr)\ge
N+2\); subtracting \eqref{plt:eq:du-psi-fixedpoint} therefore gives
\begin{equation}\label{plt:eq:du-formal-residual}
  \ordt\Bigl(\Psi_{N}-L-\log(1+t\Psi_{N})\Bigr)\ge N+1 .
\end{equation}

\emph{Step 3: the analytic residual.}  Put
\(\rho_{N}(\Lambda)=\sum_{n=1}^{N}(-1)^{n+1}\LamP{n}(\Lambda_{2})\Lambda^{-n}\),
so \(V_{N}=\Lambda+\Lambda_{2}+\rho_{N}=\Lambda(1+w)\) with
\(w=(\Lambda_{2}+\rho_{N})/\Lambda\).  For \(\Lambda\) large,
\(\AbsoluteValue w\le2\Lambda_{2}/\Lambda\le\tfrac12\), so \(V_{N}>\Lambda/2>1\)
and
\[
  R_{N}\DefinitionEquals g(V_{N})-\Lambda
  =V_{N}-\log V_{N}-\Lambda
  =\rho_{N}-\log(1+w).
\]
Expand \(\log(1+w)=\sum_{k=1}^{N+1}\frac{(-1)^{k-1}}{k}w^{k}+\theta\) with
\(\AbsoluteValue\theta\le2\AbsoluteValue w^{N+2}\le
2^{N+3}\Lambda_{2}^{N+2}\Lambda^{-(N+2)}\), the bound on the tail being
\(\sum_{k\ge N+2}\AbsoluteValue w^{k}/k\le
\AbsoluteValue w^{N+2}/\bigl((N+2)(1-\AbsoluteValue w)\bigr)\).  What remains,
\[
  P\DefinitionEquals
  \rho_{N}-\sum_{k=1}^{N+1}\frac{(-1)^{k-1}}{k}w^{k},
\]
is the value at \(t=\Lambda^{-1}\), \(L=\Lambda_{2}\) of the
\emph{polynomial}
\(\mathcal P=(\Psi_{N}-L)-\sum_{k=1}^{N+1}\frac{(-1)^{k-1}}{k}(t\Psi_{N})^{k}
\in\K[L][t]\), because both are the same finite algebraic combination of
\(t\), \(L\) and \(\Psi_{N}\).  Now
\[
  \mathcal P
  =\underbrace{\Bigl(\Psi_{N}-L-\log(1+t\Psi_{N})\Bigr)}_{\ordt\ge N+1
    \text{ by \eqref{plt:eq:du-formal-residual}}}
  +\underbrace{\sum_{k\ge N+2}\frac{(-1)^{k-1}}{k}(t\Psi_{N})^{k}}
    _{\ordt\ge N+2},
\]
so \(\ordt\mathcal P\ge N+1\), i.e.\ \(\mathcal P=t^{N+1}Q\) with
\(Q\in\K[L][t]\).  Writing \(K_{N}=\degL Q\) and evaluating at
\(t=\Lambda^{-1}\le1\), \(L=\Lambda_{2}\ge1\) gives
\(\AbsoluteValue P\le C\Lambda_{2}^{K_{N}}\Lambda^{-(N+1)}\); enlarging
\(K_{N}\) to absorb \(\theta\),
\(\AbsoluteValue{R_{N}}\le C_{N}'\Lambda_{2}^{K_{N}}\Lambda^{-(N+1)}\).

\emph{Step 4: inversion.}  We have \(g(U)=\Lambda\) and
\(g(V_{N})=\Lambda+R_{N}\).  Both \(U\) and \(V_{N}\) exceed \(\Lambda/2\ge2\)
for \(\Lambda\) large, so on the interval between them \(g'(u)=1-1/u\ge1/2\).
The mean value theorem gives
\(\AbsoluteValue{U-V_{N}}\le2\AbsoluteValue{R_{N}}\), which is
\eqref{plt:eq:du-lower-branch-error}.

For the last display, apply \eqref{plt:eq:du-lower-branch-error} with \(N=3\)
and add back the \(n=3\) term: \(U-V_{2}=\LamP{3}(\Lambda_{2})\Lambda^{-3}
+\BigOAt{\Lambda_{2}^{K_{3}}\Lambda^{-4}}{x\to0^{+}}
=\BigOAt{\Lambda_{2}^{3}\Lambda^{-3}}{x\to0^{+}}\), since
\(\deg\LamP{3}=3\) and \(\Lambda_{2}^{K}\Lambda^{-4}=
\LittleOAt{\Lambda_{2}^{3}\Lambda^{-3}}{x\to0^{+}}\) for every fixed \(K\).
\end{proof}

\section{The transported dictionary, and the repeated derivative}
\label{plt:sec:du-dictionary}

The purely algebraic half of the dictionary is \(\sigma\); the analytic half
needs its own definition, because \cref{plt:def:mul-asymptotic} names the
regime \(X\to+\infty\) explicitly.

\begin{definition}[Zero-side asymptotic expansion]
\label{plt:def:du-asymptotic}
Let \(\K\subseteq\ComplexNumbers\), let \(0<x_{0}<\EulerE^{-1}\), and let
\(f\colon(0,x_{0})\to\ComplexNumbers\).  Let \(\hat f\) be a formal series in
the zero-side generators with left-finite \(x\)-support (in the sense of
\cref{plt:def:ext-support-classes}(b)) and polynomial coefficient blocks.  We
write \(f\approx_{0}\hat f\) if for every \(N\in\RealNumbers\) there are
\(M_{N}\in\NaturalNumbers\), \(C_{N}>0\) and \(x_{N}\in(0,x_{0}]\) with
\begin{equation}\label{plt:eq:du-asymptotic}
  \AbsoluteValue{f(x)-\bigl(\Trunc{\hat f}{N}\bigr)(x)}
  \le C_{N}\,x^{N}\Lambda^{M_{N}}
  \qquad(0<x<x_{N}),
\end{equation}
where \(\bigl(\Trunc{\hat f}{N}\bigr)(x)\) is the finite sum obtained by
substituting the generators.
\end{definition}

\begin{proposition}[The two asymptotic relations correspond]
\label{plt:prop:du-asymptotic-transport}
Let \(F\in\PLlau\) and let \(f\) be a function on \((0,x_{0})\).  Put
\(\tilde f(X)\DefinitionEquals f(1/X)\) on \((1/x_{0},+\infty)\).  Then
\(f\approx_{0}\sigma(F)\) if and only if \(\tilde f\approx F\) in the sense of
\cref{plt:def:mul-asymptotic}.  The same equivalence holds for the Puiseux-log
and Hahn enlargements, with \(N\) ranging over the relevant exponent group.
\end{proposition}

\begin{proof}
Substituting \(X=1/x\) turns \(\bigl(\Trunc{F}{N}\bigr)(X)\) into
\(\bigl(\Trunc{\sigma F}{N}\bigr)(x)\) monomial by monomial, because
\(\sigma\) commutes with \(\Trunc{\cdot}{N}\)
(\cref{plt:thm:du-two-substitutions}(i)); it turns \(\log X\) into \(\Lambda\)
and \(X^{-N}(\log X)^{M}\) into \(x^{N}\Lambda^{M}\); and it is a bijection
between \((1/x_{0},+\infty)\) and \((0,x_{0})\) matching \(X\ge X_{N}\) with
\(0<x\le1/X_{N}\).  Thus \eqref{plt:eq:mul-asymptotic} and
\eqref{plt:eq:du-asymptotic} are the same inequality read in two coordinates.
\end{proof}

We can now transport a genuinely differential statement, and we do it in both
directions, on the theorem whose zero-side form the Dulac literature uses most:
the repeated-derivative formula.  Recall
\(\shiftop{n}{k}=\prod_{j=0}^{k-1}(\partial_{L}-n-j)\), with the empty product
equal to the identity.

\begin{theorem}[Repeated derivative at a finite point]
\label{plt:thm:du-repeated-derivative}
Let \(p\in\K[\Lambda]\), \(n\in\IntegerNumbers\) and \(k\in\NaturalNumbers\).
Then, in \(\K[\Lambda]((x))\),
\begin{equation}\label{plt:eq:du-repeated}
  \frac{\Differential^{k}}{\Differential x^{k}}
  \bigl(x^{n}p(\Lambda)\bigr)
  =(-1)^{k}\,x^{n-k}\,\shiftop{n-k+1}{k}\,p(\Lambda)
  =x^{n-k}\prod_{j=0}^{k-1}\bigl(n-j-\partial_{\Lambda}\bigr)\,p(\Lambda),
\end{equation}
the factors of the product commuting.  In particular
\(\Differential(x^{n}p)/\Differential x=x^{n-1}(np-p')\), and
\(\nu_{x}\bigl(\Differential f/\Differential x\bigr)\ge\nu_{x}(f)-1\) for every
\(f\in\K[\Lambda]((x))\), with equality unless \(\nu_{x}(f)=0\) and the block
\([x^{0}]f\) is a nonzero constant.
\end{theorem}

\begin{proof}
By \eqref{plt:eq:du-sigma-derivation},
\(\Differential/\Differential x=-x^{-2}\,\sigma\DX\sigma^{-1}\) as operators on
\(\K[\Lambda]((x))\).  Iterating,
\[
  \frac{\Differential^{k}}{\Differential x^{k}}
  =\bigl(-x^{-2}\,\sigma\DX\sigma^{-1}\bigr)^{k},
\]
which is not a power of \(\sigma\DX\sigma^{-1}\) because \(x^{-2}\) and
\(\sigma\DX\sigma^{-1}\) do not commute; so we argue directly by induction on
\(k\), using the case \(k=1\) supplied by
\eqref{plt:eq:grp-zero-derivation}, namely
\(\Differential(x^{m}q)/\Differential x=x^{m-1}(mq-q')\).

For \(k=0\) both sides of \eqref{plt:eq:du-repeated} are \(x^{n}p\), the empty
product being the identity.  Assume \eqref{plt:eq:du-repeated} for \(k\), in
the second form.  Applying the case \(k=1\) at outer exponent \(n-k\) to the
coefficient \(q=\prod_{j=0}^{k-1}(n-j-\partial_{\Lambda})p\) gives
\[
  \frac{\Differential^{k+1}}{\Differential x^{k+1}}\bigl(x^{n}p\bigr)
  =x^{n-k-1}\bigl((n-k)q-q'\bigr)
  =x^{n-(k+1)}\bigl(n-k-\partial_{\Lambda}\bigr)q
  =x^{n-(k+1)}\prod_{j=0}^{k}\bigl(n-j-\partial_{\Lambda}\bigr)p,
\]
which is the case \(k+1\); the factors commute because they are polynomials in
the single operator \(\partial_{\Lambda}\).  Finally
\(\prod_{j=0}^{k-1}(n-j-\partial_{\Lambda})
=(-1)^{k}\prod_{j=0}^{k-1}(\partial_{\Lambda}-n+j)\), and as \(j\) runs from
\(0\) to \(k-1\) the shifts \(n-j\) run over \(n,n-1,\ldots,n-k+1\), the same
set as the shifts \((n-k+1)+i\) for \(0\le i\le k-1\); hence the product is
\((-1)^{k}\shiftop{n-k+1}{k}\), which is the first form.

For the valuation statement, let \(n_{0}=\nu_{x}(f)\) and \(p=[x^{n_{0}}]f\).
The coefficient of \(x^{n_{0}-1}\) in \(\Differential f/\Differential x\) is
\(n_{0}p-p'\).  If \(n_{0}\ne0\) this has the same \(\Lambda\)-degree as
\(p\) and leading coefficient \(n_{0}\) times that of \(p\), hence is nonzero;
if \(n_{0}=0\) it is \(-p'\), nonzero exactly when \(\deg_{\Lambda}p>0\).
\end{proof}

\begin{warningbox}[title={The shift index moves, and only by \(k\)}]
\label{plt:warn:du-shift-index}
% ed.: S2 gives the product form of (13.7) but never names the operator; the
% shift index n-k+1 is the trap, since the naive mirror of
% plt:eq:dif-repeated would be Delta_{n,k}, which is wrong for k >= 2.
The infinity-side formula \eqref{plt:eq:dif-repeated} reads
\(\DX^{k}(t^{n}p)=t^{n+k}\shiftop{n}{k}p\): the outer exponent moves
\emph{up} by \(k\) and the shift index stays at \(n\).  The finite-point
formula \eqref{plt:eq:du-repeated} reads
\(\Differential^{k}(x^{n}p)/\Differential x^{k}
=(-1)^{k}x^{n-k}\shiftop{n-k+1}{k}p\): the outer exponent moves \emph{down} by
\(k\), a sign \((-1)^{k}\) appears, and the shift index moves to \(n-k+1\).
Writing \(\shiftop{n}{k}\) on the zero side is correct for \(k\le1\) and wrong
for every \(k\ge2\).
\end{warningbox}

Transporting in the other direction is not a formality either, because the
statement the sources make at a finite point is stated for a \emph{real}
exponent, as Dulac exponents require, whereas \cref{plt:thm:dif-repeated} is
stated over \(\IntegerNumbers\).  Carrying it back therefore yields a
strengthening.

\begin{proposition}[The return trip, and a real-exponent extension]
\label{plt:prop:du-repeated-back}
Let \(\Gamma\subseteq\RealNumbers\) be an ordered subgroup and let
\(\alpha\in\Gamma\), \(p\in\K[L]\), \(k\in\NaturalNumbers\).  In the Hahn
power--log ring \(\K[L][[t^{\Gamma}]]\), with the derivation determined by
\(\DX t^{\alpha}=-\alpha t^{\alpha+1}\) and \(\DX L=t\),
\begin{equation}\label{plt:eq:du-repeated-back}
  \DX^{k}\bigl(t^{\alpha}p(L)\bigr)=t^{\alpha+k}\,\shiftop{\alpha}{k}\,p(L),
  \qquad
  \shiftop{\alpha}{k}=\prod_{j=0}^{k-1}(\partial_{L}-\alpha-j),
\end{equation}
and applying \(\sigma\) returns \eqref{plt:eq:du-repeated} with \(n\) replaced
by \(\alpha\).  In particular the two transports are mutually inverse: starting
from \eqref{plt:eq:dif-repeated}, carrying it to the finite point by
\cref{plt:thm:du-repeated-derivative} and carrying the result back by
\eqref{plt:eq:du-sigma-derivation} reproduces \eqref{plt:eq:dif-repeated}
verbatim, with the index substitutions \(n\mapsto n\) and
\(n\mapsto n-k+1\) composing to the identity on the pair
\((\text{outer exponent},\text{shift index})\).
\end{proposition}

\begin{proof}
The block rule \(\DX(t^{\alpha}p)=t^{\alpha+1}(p'-\alpha p)\) holds for real
\(\alpha\) by the Leibniz rule applied to \(\DX t^{\alpha}=-\alpha t^{\alpha+1}\)
and \(\DX L=t\); the induction proving
\cref{plt:thm:dif-repeated} uses nothing about \(n\) beyond this rule and the
commutation of the factors, so it goes through verbatim with \(n\) replaced by
\(\alpha\).  This is \eqref{plt:eq:du-repeated-back}.

Applying \(\sigma\) to \eqref{plt:eq:du-repeated-back} and using
\eqref{plt:eq:du-sigma-derivation} \(k\) times gives
\(\bigl(-x^{2}\Differential/\Differential x\bigr)^{k}(x^{\alpha}p(\Lambda))
=x^{\alpha+k}\shiftop{\alpha}{k}p(\Lambda)\); and the induction of
\cref{plt:thm:du-repeated-derivative}, which again uses only the case
\(k=1\), gives
\(\Differential^{k}(x^{\alpha}p)/\Differential x^{k}
=(-1)^{k}x^{\alpha-k}\shiftop{\alpha-k+1}{k}p\).  These two are consistent.
Indeed, one application of \(-x^{2}\Differential/\Differential x\) sends
\(x^{m}q\) to
\[
  -x^{2}\cdot x^{m-1}\bigl(mq-q'\bigr)
  =x^{m+1}\bigl(\partial_{\Lambda}-m\bigr)q,
\]
so starting from \(x^{\alpha}p\) the \((j+1)\)-st application acts at outer
exponent \(\alpha+j\) and contributes the factor
\(\partial_{\Lambda}-\alpha-j\); after \(k\) steps the result is
\(x^{\alpha+k}\prod_{j=0}^{k-1}(\partial_{\Lambda}-\alpha-j)p
=x^{\alpha+k}\shiftop{\alpha}{k}p\), which is exactly the \(\sigma\)-image of
\eqref{plt:eq:du-repeated-back}.  Reading the same computation backwards
recovers \eqref{plt:eq:dif-repeated}, the shift index of the composite being
\(\alpha\) again.
\end{proof}

\begin{remark}[What the sources say]
\label{plt:rmk:du-sources-repeated}
% ed.: S2's eq:Dulac-repeated-derivative is correct as printed; the repair is
% that it is stated without proof, without the empty-product convention, and
% without identifying the operator, so its k >= 2 index shift is invisible.
S2 displays \eqref{plt:eq:du-repeated} in the second, product form, for real
\(\lambda\) and with no proof and no empty-product convention.  The formula is
correct; what is missing is the identification of the operator, which is where
the \(k\ge2\) index shift of \cref{plt:warn:du-shift-index} becomes visible,
and the observation that the real-exponent case is a strengthening of
\cref{plt:thm:dif-repeated} rather than an instance of it.  S6 gives only the
case \(k=1\).
\end{remark}

\section{Eight traps, each as a checkable statement}
\label{plt:sec:du-traps}

\begin{proposition}[Transport traps]
\label{plt:prop:du-traps}
Let \(F=\sum_{n\ge n_{0}}p_{n}(L)t^{n}\in\PLlau\) and \(f=\sigma F\).  Each of
the following is an identity or an equivalence, and each is the correct form
of a step that is routinely taken in the reverse or the wrong direction.
\begin{enumerate}[label=\textup{(T\arabic*)}]
\item \emph{Index sign.}  Writing \(F\) in powers of the variable,
  \(F=\sum_{m\le-n_{0}}q_{m}(L)X^{m}\) with \(q_{-n}=p_{n}\), one has
  \(f=\sum_{m\le-n_{0}}q_{m}(\Lambda)x^{-m}\).  Thus
  \(\deg_{X}F=-\ordt F=-\nu_{x}f\): the outer index used at infinity and the
  one used at zero are negatives of each other, while \(\ordt\) and
  \(\nu_{x}\) agree.
\item \emph{Dominance against the variable.}
  \(X^{a}\succdom X^{b}\) if and only if \(a>b\), whereas
  \(x^{a}\succdom x^{b}\) if and only if \(a<b\).  The dominance order read on
  the exponent of the variable therefore reverses; read on the exponent of
  \(t\), respectively \(x\), it is preserved.
\item \emph{Leading term.}  The leading monomial of \(F\) is the one of largest
  \(X\)-exponent and that of \(f\) the one of smallest \(x\)-exponent, and
  \(\sigma\) matches them.  Consequently a series written in increasing powers
  of \(X\) has its leading term last, while a series written in increasing
  powers of \(x\) has it first.
\item \emph{Truncation bound.}  \(\Trunc{F}{N}\) retains exactly the monomials
  with \(X\)-exponent \(>-N\) --- a \emph{lower} bound --- and
  \(\Trunc{f}{N}\) retains exactly those with \(x\)-exponent \(<N\) --- an
  \emph{upper} bound.  The remainder \(F-\Trunc{F}{N}\) consists of the
  monomials of \emph{small} \(X\)-exponent and of \emph{large} \(x\)-exponent.
\item \emph{The derivation is not transported.}
  \(\sigma\DX\sigma^{-1}=-x^{2}\Differential/\Differential x\), not
  \(\Differential/\Differential x\).  On blocks, \(\partial_{L}-n\) becomes
  \(n-\partial_{\Lambda}\), and for \(k\) derivatives the shift index moves as
  in \cref{plt:warn:du-shift-index}.
\item \emph{Differentiation changes the remainder in opposite directions.}
  \(\ordt(\DX F)\ge\ordt F+1\) (\cref{plt:thm:dif-valuation}), whereas
  \(\nu_{x}(\Differential f/\Differential x)\ge\nu_{x}(f)-1\)
  (\cref{plt:thm:du-repeated-derivative}).  Hence
  \(\FormalBigOAt{t^{N}}{t}\) differentiates to \(\FormalBigOAt{t^{N+1}}{t}\)
  at infinity and \(\FormalBigOAt{x^{N}}{x}\) differentiates to
  \(\FormalBigOAt{x^{N-1}}{x}\) at a finite point: at infinity a derivative
  improves a formal tail, at a finite point it costs one order.
\item \emph{Sign of the logarithm.}  \(\Lambda=-\log x\).  A block written as a
  polynomial \(P\) in \(\log x\) equals \(p(\Lambda)\) with \(p(u)=P(-u)\);
  degrees agree, and every coefficient of odd degree changes sign.  Writing
  \(\log x\) for \(\Lambda\) therefore leaves every formula formally
  consistent while reversing the sign of half the coefficients.
\item \emph{Reciprocal logarithmic generator.}  With \(\ell=\Lambda^{-1}\), a
  monomial \(x^{\alpha}\ell^{\beta}\) is \(x^{\alpha}\Lambda^{-\beta}\), i.e.\
  \(X^{-\alpha}L^{-\beta}\).  The exponent of the logarithm changes sign; a
  family listed by increasing powers of \(\ell\) is listed by
  \emph{decreasing} size in the logarithmic direction, the opposite of a family
  listed by increasing powers of \(\Lambda\).
\end{enumerate}
\end{proposition}

\begin{proof}
(T1)  \(X^{m}=t^{-m}\) and \(\sigma(t^{-m})=x^{-m}\); the rest is the
definition of \(\ordt\) and \(\nu_{x}\).

(T2)  For \(a>b\), \(X^{a}/X^{b}=X^{a-b}\to+\infty\) as \(X\to+\infty\), so
\(X^{a}\succdom X^{b}\); and \(x^{a}/x^{b}=x^{a-b}\to0\) as \(x\to0^{+}\), so
\(x^{a}\precdom x^{b}\), i.e.\ \(x^{a}\succdom x^{b}\) exactly when \(a<b\).
On the exponent of \(t\), respectively \(x\),
\eqref{plt:eq:ring-dominance} is literally the same inequality on both sides,
which is \cref{plt:thm:du-two-substitutions}(i).

(T3)  Immediate from (T1) and (T2).

(T4)  By \cref{plt:def:trunc-exclusive}, \(\Trunc{F}{N}=\sum_{n<N}p_{n}t^{n}\);
the monomial \(t^{n}\) is \(X^{-n}\), and \(n<N\) is \(-n>-N\).  On the zero
side \(t^{n}\) becomes \(x^{n}\) and the condition stays \(n<N\).  The
statement about the remainder is the negation of the two conditions.

(T5)  \eqref{plt:eq:du-sigma-derivation} and
\cref{plt:thm:du-repeated-derivative}.  The two operators are genuinely
different: \(\sigma\DX\sigma^{-1}\) raises \(\nu_{x}\) by one and
\(\Differential/\Differential x\) lowers it by one.

(T6)  The two valuation statements are \cref{plt:thm:dif-valuation}(i) and the
last clause of \cref{plt:thm:du-repeated-derivative}.  The consequences for
formal tails follow because \(\FormalBigOAt{t^{N}}{t}\) means
\(\ordt\ge N\) (\cref{plt:def:trunc-formal-o}) and likewise at zero.

(T7)  \(P(\log x)=P(-\Lambda)=p(\Lambda)\) with \(p(u)=P(-u)\); the coefficient
of \(u^{j}\) in \(p\) is \((-1)^{j}\) times that of \(P\).

(T8)  \(\ell^{\beta}=\Lambda^{-\beta}\) by \eqref{plt:eq:du-small-log}, and
\(\Lambda=L\), \(x=t\) by \eqref{plt:eq:du-sigma}.  Since \(\ell\to0^{+}\) and
\(\Lambda\to+\infty\), \(\ell^{\beta}\precdom\ell^{\beta'}\) exactly when
\(\beta>\beta'\), the reverse of \eqref{plt:eq:ring-dominance} read in
\(\Lambda\).
\end{proof}

\begin{warningbox}[title={Three ways a transported theorem is wrong while looking right}]
\label{plt:warn:du-transport}
The dangerous traps are (T5)--(T8), because each of them leaves a formula
\emph{formally consistent}.  A calculation performed at a finite point with
\(\shiftop{n}{k}\) in place of \(\shiftop{n-k+1}{k}\), or with \(\log x\) in
place of \(\Lambda\), or with the reciprocal \(\ell\) in place of \(\Lambda\),
produces a self-consistent array of coefficients that satisfies every internal
check and is simply not the expansion of the intended function.  The cheapest
detection is a two-sided residual: substitute the computed inverse back and
verify that the residual has the predicted order \emph{and the predicted
sign}, as in \cref{plt:prop:rev-residual-order}.  The second cheapest is a
numerical spot check at one value of \(x\), which \(\Lambda\leftrightarrow\ell\)
confusion fails immediately.
\end{warningbox}

\section{A statement that looks transportable and is not: reversion at a
finite point}
\label{plt:sec:du-lagrange}

All three of S1, S2 and S6 assert that near-identity reversion and the
Lagrange--B\"urmann formula ``have exactly the same form'' at a finite point.
The formulas are indeed the same; the reason is not that they are transports of
one another, and asserting that they are is a mistake with consequences,
because the summability mechanism runs in the opposite direction.  The
following proposition is the obstruction, stated so that it can be checked.

\begin{proposition}[The reversion-frame axioms are not \(\sigma\)-invariant]
\label{plt:prop:du-no-frame}
% ed.: S2 (eq:small-Lagrange) states the near-zero Lagrange-Burmann formula as
% having "exactly the same form" and gives no proof; S1 and S6 do the same for
% the fixed-point construction.  The formula is true (Theorem
% plt:thm:du-lagrange-zero) but it is not the transport of
% plt:thm:lag-burmann: the transported frame has a different derivation and a
% different distinguished element, and no rescaling of nu_x repairs (F3).
Let \(\mathfrak B\DefinitionEquals\K[\Lambda]((x))\) and
\(\partial\DefinitionEquals\Differential/\Differential x\).
\begin{enumerate}[label=\textup{(\roman*)}]
\item \(\sigma\) carries the polynomial--logarithmic reversion frame
  \(\bigl(\PLlau,\ordt,\DX,t^{-1}\bigr)\) of \cref{plt:lem:lag-pl-frame} to the
  frame \(\bigl(\mathfrak B,\nu_{x},-x^{2}\partial,\,x^{-1}\bigr)\).  Its
  derivation is \(-x^{2}\partial\), not \(\partial\), and its distinguished
  element is \(x^{-1}\), not \(x\).
\item The quadruple \(\bigl(\mathfrak B,\nu_{x},\partial,x\bigr)\) --- the one
  that governs reversion of germs at \(0^{+}\) --- satisfies \textup{(F1)},
  \textup{(F2)} and \textup{(F4)} of \cref{plt:def:lag-frame} with
  \(\Gamma=\IntegerNumbers\), but \emph{not} \textup{(F3)}: instead
  \begin{equation}\label{plt:eq:du-F3-minus}
    \nu_{x}(\partial f)\ge\nu_{x}(f)-1
    \qquad(f\in\mathfrak B),
  \end{equation}
  and \textup{(F3)} fails already at \(f=x\), where
  \(\nu_{x}(\partial x)=0<\nu_{x}(x)+1=2\).
\item No rescaling repairs this: for \(c\in\RealNumbers^{\times}\) the map
  \(c\,\nu_{x}\) satisfies \textup{(F3)} for \(\partial\) only if \(c\le-1\),
  and for \(c<0\) it violates \textup{(F1)}, as witnessed by \(f=x\),
  \(g=-x+x^{5}\), for which \(c\nu_{x}(f+g)=5c<c=\min\SetOf{c\nu_{x}f,c\nu_{x}g}\).
\end{enumerate}
Consequently \cref{plt:thm:lag-burmann} does not apply to
\(\bigl(\mathfrak B,\nu_{x},\partial,x\bigr)\), and the near-zero
Lagrange--B\"urmann formula is a second theorem, with its own proof.
\end{proposition}

\begin{proof}
(i)  \(\sigma\) is a ring isomorphism carrying \(\ordt\) to \(\nu_{x}\)
(\cref{plt:thm:du-two-substitutions}(i)), \(\DX\) to \(-x^{2}\partial\)
(\eqref{plt:eq:du-sigma-derivation}) and \(X=t^{-1}\) to \(x^{-1}\).  Since the
frame axioms are formulated in terms of these four data, the image is a frame.

(ii)  (F1): \(\K[\Lambda]\) is a domain, so \(\nu_{x}\) is a valuation with
\(\nu_{x}(\lambda)=0\) for \(\lambda\in\K^{\times}\).  (F2): a family with
\(\SetOf{i:\nu_{x}(f_{i})<\gamma}\) finite for all \(\gamma\) has
coefficientwise finite sums at each level and only finitely many levels below
any bound, exactly as in the proof of \cref{plt:lem:lag-pl-frame}.  (F4):
\(\partial x=1\).  Inequality \eqref{plt:eq:du-F3-minus} is the last clause of
\cref{plt:thm:du-repeated-derivative}, and \(\partial x=1\) gives the stated
failure of (F3).

(iii)  If \(\nu=c\nu_{x}\) satisfies (F3), then applying it to \(f=x^{n}\)
with \(\partial^{n}x^{n}=n!\) gives \(0=\nu(n!)\ge\nu(x^{n})+n=n(c+1)\) for
every \(n\ge1\), so \(c\le-1\); in particular \(c<0\).  The displayed witness
then contradicts the ultrametric inequality in (F1).
\end{proof}

The correct zero-side objects are these.

\begin{definition}[The near-identity class at a finite point]
\label{plt:def:du-near-identity}
Put
\[
  \mathfrak m_{0}\DefinitionEquals x^{2}\K[\Lambda][[x]]
  =\SetOf{h\in\K[\Lambda]((x)):\nu_{x}(h)\ge2},
  \qquad
  \mathcal G_{0}\DefinitionEquals\SetOf{x+h:h\in\mathfrak m_{0}}.
\]
For \(B\in\mathfrak m_{0}\) define the \emph{shift operator}
\begin{equation}\label{plt:eq:du-shift}
  T_{B}(f)\DefinitionEquals
  \sum_{k\ge0}\frac{B^{k}}{k!}\,
  \frac{\Differential^{k}f}{\Differential x^{k}},
  \qquad f\in\K[\Lambda]((x)),
\end{equation}
and write \(f\compose(x+B)\DefinitionEquals T_{B}(f)\).
\end{definition}

\begin{lemma}[The shift operator is substitution of \(x+B\)]
\label{plt:lem:du-taylor-zero}
Let \(B\in\mathfrak m_{0}\).  Then:
\begin{enumerate}[label=\textup{(\roman*)}]
\item the family in \eqref{plt:eq:du-shift} is summable, its \(k\)-th member
  having \(\nu_{x}\ge\nu_{x}(f)+k\); hence \(\nu_{x}(T_{B}f)\ge\nu_{x}(f)\);
\item \(T_{B}\) is a \(\K\)-algebra endomorphism of \(\K[\Lambda]((x))\),
  continuous for \(\nu_{x}\), with \(T_{0}=\mathrm{id}\);
\item on generators,
  \begin{equation}\label{plt:eq:du-generator-substitution}
    T_{B}(x)=x+B,
    \qquad
    T_{B}(\Lambda)=\Lambda-\log\Bigl(1+\frac Bx\Bigr)=\log\frac1{x+B};
  \end{equation}
\item \(T_{B'}\compose T_{B}=T_{B''}\) with
  \(B''=B'+T_{B'}(B)\in\mathfrak m_{0}\); consequently \(\mathcal G_{0}\) is a
  monoid under \((x+B)\compose(x+B')\DefinitionEquals T_{B'}(x+B)=x+B''\), with
  neutral element \(x\), and \(T\) is an anti-homomorphism from it into the
  monoid of endomorphisms of \(\K[\Lambda]((x))\).
\end{enumerate}
\end{lemma}

\begin{proof}
(i)  \(\nu_{x}(B^{k})\ge2k\) and, by \eqref{plt:eq:du-F3-minus},
\(\nu_{x}(\Differential^{k}f/\Differential x^{k})\ge\nu_{x}(f)-k\); the product
has \(\nu_{x}\ge\nu_{x}(f)+k\to+\infty\), so (F2) for
\(\bigl(\mathfrak B,\nu_{x}\bigr)\) --- established in
\cref{plt:prop:du-no-frame}(ii) --- gives the sum, and the bound is the
minimum at \(k=0\).

(ii)  Linearity is clear.  For multiplicativity, expand
\(\Differential^{k}(fg)/\Differential x^{k}\) by Leibniz and regroup; the
double family \(\bigl(B^{i}\partial^{i}f/i!\bigr)\bigl(B^{j}\partial^{j}g/j!\bigr)\)
has \(\nu_{x}\ge\nu_{x}(f)+\nu_{x}(g)+i+j\), so it is summable and the
regrouping clause of (F2) applies:
\[
  T_{B}(fg)=\sum_{k}\frac{B^{k}}{k!}\sum_{i+j=k}\binom ki
  \partial^{i}f\,\partial^{j}g
  =\sum_{i,j}\frac{B^{i}\partial^{i}f}{i!}\cdot\frac{B^{j}\partial^{j}g}{j!}
  =T_{B}(f)\,T_{B}(g).
\]
Continuity follows from (i).

(iii)  \(\partial x=1\) and \(\partial^{k}x=0\) for \(k\ge2\), giving
\(T_{B}(x)=x+B\).  For \(\Lambda\): by \eqref{plt:eq:du-derivation},
\(\partial\Lambda=-x^{-1}\) and hence
\(\partial^{k}\Lambda=(-1)^{k}(k-1)!\,x^{-k}\) for \(k\ge1\), so
\[
  T_{B}(\Lambda)
  =\Lambda+\sum_{k\ge1}\frac{B^{k}}{k!}(-1)^{k}(k-1)!\,x^{-k}
  =\Lambda+\sum_{k\ge1}\frac{(-1)^{k}}{k}\Bigl(\frac Bx\Bigr)^{k}
  =\Lambda-\log\Bigl(1+\frac Bx\Bigr),
\]
the formal logarithm being defined because \(\nu_{x}(B/x)\ge1\).  Since
\(\Lambda=\log(1/x)\) and \(\log\bigl(1+B/x\bigr)=\log\bigl((x+B)/x\bigr)\),
this is \(\log\bigl(1/(x+B)\bigr)\).

(iv)  Both \(T_{B'}\compose T_{B}\) and \(T_{B''}\) are continuous
\(\K\)-algebra endomorphisms, and \(\K[\Lambda]((x))\) is topologically
generated by \(x^{\pm1}\) and \(\Lambda\); so it suffices to compare them
there.  On \(x\): \(T_{B'}(T_{B}x)=T_{B'}(x+B)=x+B'+T_{B'}(B)=x+B''\), and
\(T_{B''}(x)=x+B''\).  On \(x^{-1}\): a ring homomorphism sends the inverse of
a unit to the inverse of its image, and \(x+B'\) is a unit of
\(\K[\Lambda]((x))\), its leading block being the scalar \(1\)
(\cref{plt:thm:div-unit-criterion}, transported by
\cref{plt:thm:du-two-substitutions}(i)); so the values on \(x^{-1}\) are
determined by those on \(x\).  On \(\Lambda\): \(T_{B'}\) is a continuous ring
homomorphism, hence \(T_{B'}\bigl(\log(1+w)\bigr)=\log\bigl(1+T_{B'}(w)\bigr)\)
whenever \(\nu_{x}(w)>0\), and \(T_{B'}(B/x)=T_{B'}(B)/(x+B')\).  Therefore, by
(iii),
\begin{align*}
  T_{B'}\bigl(T_{B}\Lambda\bigr)
  &=\Lambda-\log\Bigl(1+\frac{B'}x\Bigr)
        -\log\Bigl(1+\frac{T_{B'}(B)}{x+B'}\Bigr)\\
  &=\Lambda-\log\Bigl(\frac{x+B'}{x}\cdot
        \frac{x+B'+T_{B'}(B)}{x+B'}\Bigr)
  =\Lambda-\log\Bigl(1+\frac{B''}x\Bigr),
\end{align*}
which is \(T_{B''}(\Lambda)\).  Finally
\(\nu_{x}(B'')\ge\min\SetOf{\nu_{x}(B'),\nu_{x}(B)}\ge2\)
by (i), so \(B''\in\mathfrak m_{0}\) and \(\mathcal G_{0}\) is closed; neutrality
of \(x\) is \(T_{0}=\mathrm{id}\), and associativity is (iv) itself.
\end{proof}

\begin{theorem}[Lagrange--B\"urmann reversion at a finite point]
\label{plt:thm:du-lagrange-zero}
Let \(h\in\mathfrak m_{0}\), i.e.\ \(h\in\K[\Lambda][[x]]\) with
\(\nu_{x}(h)\ge2\), and put \(\delta\DefinitionEquals\nu_{x}(h)-1\ge1\).  Then
the family
\[
  \Bigl(B_{r}\Bigr)_{r\ge1},
  \qquad
  B_{r}\DefinitionEquals\frac{(-1)^{r}}{r!}
  \frac{\Differential^{\,r-1}}{\Differential x^{\,r-1}}\bigl(h^{r}\bigr),
\]
is summable, with \(\nu_{x}(B_{r})\ge r\delta+1\); its sum \(B\) lies in
\(\mathfrak m_{0}\); and
\begin{equation}\label{plt:eq:du-lagrange-zero}
  \rev{(x+h)}
  =x+\sum_{r\ge1}\frac{(-1)^{r}}{r!}
  \frac{\Differential^{\,r-1}}{\Differential x^{\,r-1}}\bigl(h^{r}\bigr)
\end{equation}
is the two-sided compositional inverse of \(x+h\) in \(\mathcal G_{0}\).  In
particular \(\mathcal G_{0}\) is a group.  The same statement holds in
\(\K[\Lambda]((x^{1/s}))\) and in the Hahn enlargements, with \(\mathfrak m_{0}\)
replaced by \(\SetOf{h:\nu_{x}(h)>1}\); the threshold \(\nu_{x}(h)>1\) is
exactly the mirror of \(\ordt(A)>-1\) of \cref{plt:def:lag-frame}, and
\(1=\nu_{x}(x)\) is the mirror of \(-1=\ordt(X)\).
\end{theorem}

\begin{proof}
Write \(\partial=\Differential/\Differential x\) and
\(\mathfrak B=\K[\Lambda]((x))\).

\emph{Step 1: summability.}  \(\nu_{x}(h^{r})\ge r\nu_{x}(h)\) and
\(\partial^{r-1}\) costs at most \(r-1\) by \eqref{plt:eq:du-F3-minus}, so
\(\nu_{x}(B_{r})\ge r\nu_{x}(h)-(r-1)=r\delta+1\), which tends to \(+\infty\)
since \(\delta>0\).  Hence \((B_{r})\) is summable and
\(\nu_{x}(B)\ge\delta+1=\nu_{x}(h)\ge2\); also \(B\in\K[\Lambda][[x]]\), so
\(B\in\mathfrak m_{0}\) and \(x+B\in\mathcal G_{0}\).

\emph{Step 2: what has to be proved.}  By
\cref{plt:lem:du-taylor-zero}(iv),
\((x+h)\compose(x+B)=T_{B}(x+h)=x+B+T_{B}(h)\).  So \(x+B\) is a right inverse
of \(x+h\) in \(\mathcal G_{0}\) if and only if
\begin{equation}\label{plt:eq:du-lagrange-fixedpoint}
  B=-T_{B}(h).
\end{equation}
Granting \eqref{plt:eq:du-lagrange-fixedpoint}, every element of the monoid
\(\mathcal G_{0}\) has a right inverse; and in a monoid in which every element
has a right inverse, every element is invertible: if \(ab=x\) and \(bc=x\) then
\(a=a\compose(b\compose c)=(a\compose b)\compose c=c\), whence \(ba=bc=x\).  So
\(\mathcal G_{0}\) is a group and \(x+B=\rev{(x+h)}\) is the two-sided inverse.
It remains to prove \eqref{plt:eq:du-lagrange-fixedpoint}.

\emph{Step 3: the weighted deformation algebra.}  Define
\(\mathfrak B\langle\varepsilon\rangle_{a}\) and
\(\mathfrak B\langle\varepsilon\rangle\) exactly as in
\cref{plt:lem:lag-eval}, with the present \(\nu_{x}\), the present \(\partial\)
and the weight \(\delta\).  The proof of \cref{plt:lem:lag-eval} uses
\textup{(F3)} in one place only, to see that \(\partial\) maps
\(\mathfrak B\langle\varepsilon\rangle_{a}\) into some
\(\mathfrak B\langle\varepsilon\rangle_{a'}\); by \eqref{plt:eq:du-F3-minus}
that holds here with \(a'=a+1\) in place of \(a-1\).  Everything else in that
proof uses only \textup{(F1)} and \textup{(F2)}, which hold by
\cref{plt:prop:du-no-frame}(ii).  Thus all three conclusions of
\cref{plt:lem:lag-eval} are available: \(\mathfrak B\langle\varepsilon\rangle\)
is a \(\partial\)-stable subring of \(\mathfrak B[[\varepsilon]]\); evaluation
\(\mathrm{ev}(u)=\sum_{m}u_{m}\) is a \(\K\)-algebra homomorphism with
\(\mathrm{ev}(\varepsilon)=1\); and \(\mathrm{ev}\) commutes with sums of
families \(\bigl(f^{(k)}\bigr)\) lying in a common
\(\mathfrak B\langle\varepsilon\rangle_{a}\) with
\(f^{(k)}\in\varepsilon^{k}\mathfrak B[[\varepsilon]]\).

\emph{Step 4: the parameter identity.}  \Cref{plt:lem:lag-abstract} is stated
for an arbitrary commutative \(\RationalNumbers\)-algebra and uses no frame;
take \(C=\mathfrak B\).  Let
\[
  \theta\colon\mathfrak B\to\mathfrak B[[z]],
  \qquad
  \theta(f)=\sum_{k\ge0}\bigl(\partial^{k}f\bigr)\frac{z^{k}}{k!} .
\]
This is a \(\K\)-algebra homomorphism (Leibniz, as in
\cref{plt:lem:du-taylor-zero}(ii), with no summability needed because
\(\mathfrak B[[z]]\) is a formal power-series ring), it satisfies
\(\partial_{z}\theta=\theta\partial\) by a term shift, and
\(\mathrm{ev}_{z=0}\theta=\mathrm{id}\).  Put \(\psi=\theta(-h)\) and
\(\gamma=\theta(x)=x+z\), so \(\partial_{z}\gamma=1\).  Let
\(u\in\varepsilon\mathfrak B[[z]][[\varepsilon]]\) be the unique solution of
\(u=\varepsilon\psi_{[u]}\).  Since \(\gamma_{[u]}=x+z+u\), formula
\eqref{plt:eq:lag-param} gives, for \(m\ge1\),
\begin{equation}\label{plt:eq:du-lag-coefficients}
  [\varepsilon^{m}]u
  =\frac1{m!}\partial_{z}^{\,m-1}\bigl(\psi^{m}\bigr)
  =\frac{(-1)^{m}}{m!}\partial_{z}^{\,m-1}\theta\bigl(h^{m}\bigr)
  =\theta\bigl(B_{m}\bigr).
\end{equation}

\emph{Step 5: evaluation at \(z=0\).}  The map \(\mathrm{ev}_{z=0}\), applied
coefficientwise in \(\varepsilon\), is a \(\K\)-algebra homomorphism
\(\mathfrak B[[z]][[\varepsilon]]\to\mathfrak B[[\varepsilon]]\); each
\(\varepsilon\)-coefficient of \(\psi_{[u]}=\sum_{k}(\partial_{z}^{k}\psi)u^{k}/k!\)
is a finite sum because \(u\in\varepsilon(\cdot)\), so applying it to
\(u=\varepsilon\psi_{[u]}\) is legitimate.  Write
\(\tilde u\DefinitionEquals\mathrm{ev}_{z=0}(u)=\sum_{m\ge1}B_{m}\varepsilon^{m}\)
by \eqref{plt:eq:du-lag-coefficients}, and note
\(\mathrm{ev}_{z=0}\bigl(\partial_{z}^{k}\psi\bigr)
=\mathrm{ev}_{z=0}\theta\bigl(\partial^{k}(-h)\bigr)=-\partial^{k}h\).  Hence
\begin{equation}\label{plt:eq:du-lag-eps}
  \tilde u=\varepsilon\sum_{k\ge0}\frac{-\partial^{k}h}{k!}\,\tilde u^{\,k}
  \quad\text{in }\mathfrak B[[\varepsilon]].
\end{equation}
By Step~1, \(\nu_{x}(B_{m})\ge m\delta+1\), so
\(\tilde u\in\mathfrak B\langle\varepsilon\rangle_{0}\) and
\(\mathrm{ev}(\tilde u)=\sum_{m}B_{m}=B\).

\emph{Step 6: evaluation at \(\varepsilon=1\).}  Put
\(f^{(k)}\DefinitionEquals\varepsilon\,\frac{-\partial^{k}h}{k!}\tilde u^{\,k}
\in\varepsilon^{k+1}\mathfrak B[[\varepsilon]]\).  Its \(\varepsilon^{m}\)
coefficient is a finite sum of terms
\(\frac{-\partial^{k}h}{k!}B_{m_{1}}\cdots B_{m_{k}}\) with
\(m_{1}+\cdots+m_{k}=m-1\) and every \(m_{i}\ge1\), so
\[
  \nu_{x}\bigl([\varepsilon^{m}]f^{(k)}\bigr)
  \ge\bigl(\nu_{x}(h)-k\bigr)+\sum_{i=1}^{k}\bigl(m_{i}\delta+1\bigr)
  =\nu_{x}(h)+\delta(m-1)
  =\delta m+1
  \ge\delta m ,
\]
the bound being uniform in \(k\).  Thus every \(f^{(k)}\) lies in the
\emph{same} \(\mathfrak B\langle\varepsilon\rangle_{0}\), and
\cref{plt:lem:lag-eval}(3) applies to \eqref{plt:eq:du-lag-eps}:
\[
  B=\mathrm{ev}(\tilde u)
  =\sum_{k\ge0}\mathrm{ev}\bigl(f^{(k)}\bigr)
  =\sum_{k\ge0}\frac{-\partial^{k}h}{k!}\,\mathrm{ev}(\tilde u)^{k}
  =-\sum_{k\ge0}\frac{B^{k}}{k!}\,\partial^{k}h
  =-T_{B}(h),
\]
using \(\mathrm{ev}(\varepsilon)=1\) and multiplicativity.  This is
\eqref{plt:eq:du-lagrange-fixedpoint}.

% ed.: the appeal to the enlarged exponent group also needs
% plt:lem:du-taylor-zero there, which was stated for m_0; its proof uses
% nu_x(B) >= 2 only through nu_x(B) - 1 > 0, so it generalises as stated.
For the Puiseux-log and Hahn versions, every step used only
\(\delta=\nu_{x}(h)-1>0\) and the properties (F1), (F2) of \(\nu_{x}\), which
hold verbatim over the enlarged exponent group; and
\cref{plt:lem:du-taylor-zero}, whose proof uses \(\nu_{x}(B)\ge2\) only
through the estimate \(\nu_{x}\bigl(B^{k}\partial^{k}f/k!\bigr)\ge\nu_{x}(f)
+k\bigl(\nu_{x}(B)-1\bigr)\) and the membership
\(\nu_{x}(B/x)>0\), holds there with \(\mathfrak m_{0}\) replaced by
\(\SetOf{B:\nu_{x}(B)>1}\).
\end{proof}

\begin{proposition}[Logarithmic degree of a finite-point reversion]
\label{plt:prop:du-degree-zero}
Let \(h=\sum_{n\ge2}c_{n}(\Lambda)x^{n}\in\mathfrak m_{0}\) with
\(\deg_{\Lambda}c_{n}\le d\) for all \(n\).  Then, for every \(n\ge2\),
\begin{equation}\label{plt:eq:du-degree-zero}
  \deg_{\Lambda}\bigl[x^{n}\bigr]\rev{(x+h)}\le(n-1)\,d,
\end{equation}
and the bound is attained for \(h=ax^{2}\Lambda\) with \(a\ne0\), \(d=1\).
\end{proposition}

\begin{proof}
Differentiation does not raise the logarithmic degree: by
\cref{plt:thm:du-repeated-derivative} one has
\[
  \frac{\Differential}{\Differential x}\bigl(x^{m}p\bigr)=x^{m-1}(mp-p'),
  \qquad
  \deg_{\Lambda}(mp-p')\le\deg_{\Lambda}p .
\]
Multiplication adds degrees, so
every block of \(h^{r}\) has \(\deg_{\Lambda}\le rd\), and hence so does every
block of \(\partial^{r-1}(h^{r})\).  By Step~1 of
\cref{plt:thm:du-lagrange-zero} with \(\delta\ge1\),
\(\nu_{x}(B_{r})\ge r+1\), so only the terms with \(r\le n-1\) contribute to
\([x^{n}]\rev{(x+h)}\); the largest degree among them is at most \((n-1)d\).
Attainment is \cref{plt:ex:du-lagrange} below, where the \(x^{n}\) block has
degree exactly \(n-1\) for \(n=2,3,4\); in general the \(r=n-1\) term of
\(h=ax^{2}\Lambda\) contributes
\(\frac{(-1)^{n-1}}{(n-1)!}\partial^{n-2}\bigl(a^{n-1}x^{2n-2}\Lambda^{n-1}\bigr)\),
whose \(\Lambda^{n-1}\) coefficient is
\(\frac{(-1)^{n-1}a^{n-1}}{(n-1)!}\prod_{j=0}^{n-3}(2n-2-j)\ne0\) by
\eqref{plt:eq:du-repeated} --- the pure top-degree part of
\(\prod_{j}(m-j-\partial_{\Lambda})\Lambda^{n-1}\) being
\(\prod_{j}(m-j)\Lambda^{n-1}\) with \(m=2n-2\) --- and no other \(r\)
contributes to \(\Lambda^{n-1}\) at level \(x^{n}\).
\end{proof}

\begin{example}[A Dulac reversion computed at zero and checked at infinity]
\label{plt:ex:du-lagrange}
Let \(a\in\K^{\times}\) and \(f(x)=x+a x^{2}\Lambda\), a tangent-to-identity
Dulac germ.  Here \(h=ax^{2}\Lambda\), \(\nu_{x}(h)=2\), \(\delta=1\).  Using
\eqref{plt:eq:du-repeated} at each step:
\begin{align*}
  B_{1}&=-h=-ax^{2}\Lambda,\\
  B_{2}&=\tfrac12\partial\bigl(a^{2}x^{4}\Lambda^{2}\bigr)
        =\tfrac{a^{2}}2x^{3}\bigl(4\Lambda^{2}-2\Lambda\bigr)
        =a^{2}x^{3}\bigl(2\Lambda^{2}-\Lambda\bigr),\\
  B_{3}&=-\tfrac16\partial^{2}\bigl(a^{3}x^{6}\Lambda^{3}\bigr)
        =-\tfrac{a^{3}}6x^{4}(\partial_{\Lambda}-5)(\partial_{\Lambda}-6)\Lambda^{3}
        =-\tfrac{a^{3}}6x^{4}\bigl(30\Lambda^{3}-33\Lambda^{2}+6\Lambda\bigr),
\end{align*}
and \(\nu_{x}(B_{4})\ge5\).  Hence
\begin{equation}\label{plt:eq:du-lagrange-example}
  \rev{f}(x)
  =x-ax^{2}\Lambda
   +a^{2}x^{3}\bigl(2\Lambda^{2}-\Lambda\bigr)
   +a^{3}x^{4}\Bigl(-5\Lambda^{3}+\tfrac{11}2\Lambda^{2}-\Lambda\Bigr)
   +\FormalBigOAt{x^{5}}{x}.
\end{equation}
\emph{Direct verification.}  Write \(g=x+B\) with \(B=-a\Lambda x^{2}
+a^{2}(2\Lambda^{2}-\Lambda)x^{3}+\beta_{4}x^{4}\), \(w=B/x\), and use
\(\Lambda\compose g=\Lambda-\log(1+w)\) from
\eqref{plt:eq:du-generator-substitution}.  Since \(g^{2}\) has
\(\nu_{x}=2\), only two terms of \(\log(1+w)\) are needed:
\begin{align*}
  \log(1+w)&=-a\Lambda x+a^{2}\bigl(\tfrac32\Lambda^{2}-\Lambda\bigr)x^{2}
  +\FormalBigOAt{x^{3}}{x},\\
  g^{2}&=x^{2}-2a\Lambda x^{3}+a^{2}(5\Lambda^{2}-2\Lambda)x^{4}
  +\FormalBigOAt{x^{5}}{x}.
\end{align*}
Multiplying,
\[
  a\,g^{2}\,(\Lambda\compose g)
  =a\Lambda x^{2}+a^{2}\bigl(\Lambda-2\Lambda^{2}\bigr)x^{3}
   +a^{3}\bigl(5\Lambda^{3}-\tfrac{11}2\Lambda^{2}+\Lambda\bigr)x^{4}
   +\FormalBigOAt{x^{5}}{x},
\]
so that \(f\compose g=g+a g^{2}(\Lambda\compose g)=x\) through order \(x^{4}\)
exactly when \(\beta_{4}=a^{3}\bigl(-5\Lambda^{3}+\tfrac{11}2\Lambda^{2}
-\Lambda\bigr)\), which is \eqref{plt:eq:du-lagrange-example}.

\emph{The same germ at infinity.}  Conjugating by \(j(x)=1/x\) as in
\cref{plt:prop:grp-zero-dictionary} and writing \(L=\Lambda\), the germ
\(F\DefinitionEquals j\compose f\compose j\) is \(F(X)=1/f(1/X)\):
\[
  F(X)=\frac{1}{X^{-1}+aX^{-2}L}=\frac{X}{1+atL}
  =X-aL+a^{2}L^{2}t-a^{3}L^{3}t^{2}+\cdots\in\Gnear .
\]
Thus the single zero-side block \(ax^{2}\Lambda\) becomes an \emph{infinite}
perturbation \(A=-aL+a^{2}L^{2}t-\cdots\in\PLpow\) at infinity: the map
conjugation \(f\mapsto j\compose f\compose j\) is not the coefficientwise
dictionary \(\sigma\), even though \(\sigma\) is a term-by-term relabelling on
series.  Reversion at infinity now proceeds by
\eqref{plt:eq:lag-burmann-pl}, with \(\DX\) and the shift operators
\(\shiftop{n}{k}\) --- a different computation with a different summability
count, whose answer is \(j\compose\rev{f}\compose j\).  This is the concrete
form of \cref{plt:prop:du-no-frame}.
\end{example}

\begin{remark}[The source formula for tangent-to-identity Dulac germs]
\label{plt:rmk:du-s1-dulac}
% ed.: S1's eq:dulac-inverse-first is correct, but it is written in the
% reciprocal convention ell = -1/log x, so its exponents of the logarithm are
% the negatives of ours; transporting it without noting this produces the
% wrong sign on every logarithmic exponent.
S1 computes, for \(f(x)=x+a x^{\alpha}\ell^{m}\) with
\(\ell=-1/\log x=\Lambda^{-1}\) and \(\alpha>1\),
\[
  \rev{f}(x)=x-a x^{\alpha}\ell^{m}
  +a^{2}x^{2\alpha-1}\bigl(\alpha\ell^{2m}+m\ell^{2m+1}\bigr)+\cdots.
\]
Here \(h=ax^{\alpha}\ell^{m}\) has \(\nu_{x}(h)=\alpha>1\) but is not in
\(\K[\Lambda][[x]]\) when \(m>0\); the applicable form of
\cref{plt:thm:du-lagrange-zero} is therefore its last clause, in the Hahn
enlargement with real \(x\)-exponents and Laurent \(\Lambda\)-coefficients.
Granting that, the formula agrees with \eqref{plt:eq:du-lagrange-zero} at
\(r\le2\): from
\(\Differential\ell/\Differential x=\ell^{2}/x\) one gets
\(h'=ax^{\alpha-1}(\alpha\ell^{m}+m\ell^{m+1})\) and
\(\tfrac12\partial(h^{2})=hh'\).  In our generator the same statement reads
\(\rev{f}=x-ax^{\alpha}\Lambda^{-m}
+a^{2}x^{2\alpha-1}\bigl(\alpha\Lambda^{-2m}+m\Lambda^{-2m-1}\bigr)+\cdots\):
the exponents of the logarithm are the \emph{negatives} of S1's.  Note also
that the dominant logarithmic term of the \(x^{2\alpha-1}\) block is
\(\alpha\ell^{2m}\) when \(m>0\), since \(\ell\to0\); in \(\Lambda\) it is
\(\alpha\Lambda^{-2m}\), again the dominant one, so the ordering statement
survives the translation even though every exponent changes sign.  This is
\cref{plt:prop:du-traps}\,(T8) in action.
\end{remark}

\section{Dulac series, Dulac germs, and how the support condition differs}
\label{plt:sec:du-dulac}

Two distinct conventions travel under the name ``power--log series near
zero''.  The older one, of Dulac and of the finiteness literature, keeps
\emph{polynomial} coefficients in the large logarithm and enlarges the
exponents of \(x\) to a set of reals increasing to \(+\infty\).  The newer one,
of the normal-form literature
\cite{mardesic-normalforms,mardesic-length,resman-dulac,peran-logtransseries},
keeps \(x\)-exponents real and takes real exponents of the \emph{small}
logarithm \(\ell=\Lambda^{-1}\).  They are not the same class, they are not
nested, and their intersection is small.  We define both and prove the
comparison.

\begin{definition}[Dulac series]
\label{plt:def:du-dulac-series}
A set \(S\subseteq\RealNumbers\) is \emph{left-finite} if \(S\cap(-\infty,M)\)
is finite for every \(M\in\RealNumbers\); equivalently
(\cref{plt:prop:ext-leftfinite}(i)), \(S\) is finite or can be enumerated as a
strictly increasing sequence \(\alpha_{0}<\alpha_{1}<\cdots\to+\infty\).  A
\emph{Dulac series} over \(\K\) is a formal expression
\begin{equation}\label{plt:eq:du-dulac-series}
  \hat f=\sum_{\alpha\in S}x^{\alpha}p_{\alpha}(\Lambda),
  \qquad
  S\subseteq\RealNumbers\ \text{left-finite},
  \quad
  p_{\alpha}\in\K[\Lambda]\setminus\SetOf0,
\end{equation}
and \(\mathfrak D(\K)\) denotes the set of these.  For \(\hat f\ne0\) and
\(N\in\RealNumbers\) we write
\[
  \nu_{x}(\hat f)=\min S,
  \qquad
  \Trunc{\hat f}{N}=\sum_{\alpha<N}x^{\alpha}p_{\alpha}(\Lambda);
\]
both are meaningful because \(S\) is left-finite.  For a
sub-semigroup \(\Gamma\subseteq(\RealNumbers,+)\) and \(\alpha_{0}\in
\RealNumbers\) we say that \(\hat f\) is \emph{carried by the grid}
\(\alpha_{0}+\Gamma\) if \(S\subseteq\alpha_{0}+\Gamma\); this is
\cref{plt:def:ext-support-classes}(c).  Finally
\(\mathfrak D_{>0}(\K)\DefinitionEquals\SetOf{\hat f\in\mathfrak D(\K):
\nu_{x}(\hat f)>0}\).
\end{definition}

\begin{lemma}[Finitely generated positive semigroups are left-finite]
\label{plt:lem:du-fg-semigroup}
Let \(g_{1},\dots,g_{k}\in\RealNumbers_{>0}\) and let
\(\Gamma=\SetOf{n_{1}g_{1}+\cdots+n_{k}g_{k}:n_{i}\in\NaturalNumbers}\).  Then
\(\Gamma\) is left-finite, and every \(\gamma\in\Gamma\) has only finitely many
representations.  Consequently every grid \(\alpha_{0}+\Gamma\) is left-finite,
and a sum of two such grids is again one.
\end{lemma}

\begin{proof}
Put \(g_{\min}=\min_{i}g_{i}>0\).  If \(\sum n_{i}g_{i}\le M\) then
\(g_{\min}\sum n_{i}\le M\), so \(\sum n_{i}\le M/g_{\min}\); there are only
finitely many tuples \((n_{1},\dots,n_{k})\in\NaturalNumbers^{k}\) with bounded
sum, which proves both assertions at once.  Translation preserves
left-finiteness, and
\((\alpha_{0}+\Gamma)+(\alpha_{0}'+\Gamma')
=(\alpha_{0}+\alpha_{0}')+\Gamma''\) with \(\Gamma''\) generated by the union
of the two generating sets.
\end{proof}

\begin{proposition}[\(\mathfrak D(\K)\) is a complete valued domain]
\label{plt:prop:du-dulac-ring}
\(\mathfrak D(\K)\) is a commutative \(\K\)-algebra under the operations
\[
  \hat f+\hat g=\sum_{\alpha}x^{\alpha}\bigl(p_{\alpha}+q_{\alpha}\bigr),
  \qquad
  \hat f\hat g=\sum_{\gamma}x^{\gamma}
  \sum_{\alpha+\beta=\gamma}p_{\alpha}q_{\beta},
\]
both well defined; it is an integral domain; \(\nu_{x}\) is a valuation on it
with value group \(\RealNumbers\); it is complete and separated for the
filtration by \(\SetOf{\nu_{x}\ge N}\); and
\begin{equation}\label{plt:eq:du-units}
  \hat f\in\mathfrak D(\K)^{\times}
  \iff
  p_{\nu_{x}(\hat f)}\in\K^{\times},
\end{equation}
i.e.\ exactly when the leading coefficient block is a nonzero scalar --- the
verbatim mirror of \cref{plt:thm:div-unit-criterion}.
\end{proposition}

\begin{proof}
Well-definedness of the sum is clear (a union of two left-finite sets is
left-finite).  For the product, \cref{plt:prop:ext-leftfinite}(iii) gives that
\(S+T\) is left-finite and that every fiber
\(\SetOf{(\alpha,\beta)\in S\times T:\alpha+\beta=\gamma}\) is finite, so each
inner sum is a finite sum in \(\K[\Lambda]\).  Associativity, distributivity
and the \(\K\)-algebra axioms are then verified fiberwise.

Leading blocks multiply: \([x^{\nu_{x}\hat f+\nu_{x}\hat g}](\hat f\hat g)
=p_{\nu_{x}\hat f}q_{\nu_{x}\hat g}\ne0\) because \(\K[\Lambda]\) is a domain.
This gives both integrality and \(\nu_{x}(\hat f\hat g)=\nu_{x}\hat f
+\nu_{x}\hat g\); the ultrametric inequality is immediate.

Separatedness is clear.  For completeness, let \((\hat f_{k})\) be Cauchy:
for every \(M\) there is \(k_{M}\) with
\(\nu_{x}(\hat f_{k}-\hat f_{l})\ge M\) for \(k,l\ge k_{M}\).  Then for each
\(M\) the truncations \(\Trunc{\hat f_{k}}{M}\) are eventually constant; define
\(\hat f\) by these stable truncations.  Its support below any \(M\) equals the
support of \(\Trunc{\hat f_{k_{M}}}{M}\), which is finite, so the support of
\(\hat f\) is left-finite and \(\hat f\in\mathfrak D(\K)\), with
\(\hat f_{k}\to\hat f\).

For \eqref{plt:eq:du-units}: if \(\hat f\hat g=1\) then the leading blocks
multiply to \(1\) in \(\K[\Lambda]\), whose units are \(\K^{\times}\).
Conversely if \(p_{\alpha_{0}}=c\in\K^{\times}\) with \(\alpha_{0}=\nu_{x}\hat f\),
write \(\hat f=cx^{\alpha_{0}}(1+u)\) with \(\nu_{x}(u)>0\); the support of
\(u\) is left-finite and contained in \(\RealNumbers_{>0}\), so by
\cref{plt:prop:ext-leftfinite}(iv) the semigroup it generates is left-finite,
the family \((-u)^{k}\) is summable, and
\(\hat f^{-1}=c^{-1}x^{-\alpha_{0}}\sum_{k\ge0}(-u)^{k}\in\mathfrak D(\K)\).
\end{proof}

\begin{theorem}[Dulac series against the classes of this volume]
\label{plt:thm:du-dulac-vs-class}
Inside \(\mathfrak D(\K)\), with \(\sigma\) as in \eqref{plt:eq:du-sigma}:
\begin{enumerate}[label=\textup{(\roman*)}]
\item \(\sigma(\PLlau)=\SetOf{\hat f\in\mathfrak D(\K):S\subseteq\IntegerNumbers}\)
  and \(\sigma\bigl(\K[L]((t^{1/s}))\bigr)
  =\SetOf{\hat f:S\subseteq s^{-1}\IntegerNumbers}\); the inclusions
  \[
    \sigma(\PLlau)
    \subsetneq
    \bigcup_{s\ge1}\sigma\bigl(\K[L]((t^{1/s}))\bigr)
    \subsetneq
    \mathfrak D(\K)
  \]
  are strict, the second witnessed by \(x^{\sqrt2}\).
\item \(\mathfrak D(\K)\) is \emph{strictly} contained in the set of elements of
  \(\sigma\bigl(\mathcal H_{1}(\K,\RealNumbers)\bigr)\) all of whose
  coefficient blocks are polynomials in \(\Lambda\): the Hahn condition on such
  an element is well-basedness of its \(x\)-support, which is strictly weaker
  than left-finiteness once the exponent group is dense.  The series
  \(\sum_{k\ge1}x^{1-1/k}\) has well-based \(x\)-support of order type
  \(\omega\) and is not a Dulac series.
\item For \(S\subseteq\IntegerNumbers\) --- and more generally for
  \(S\subseteq\Gamma\) with \(\Gamma\subseteq\RealNumbers\) a discrete
  subgroup --- the three conditions ``bounded below'', ``left-finite'' and
  ``well-based'' coincide, which is \cref{plt:thm:ring-support}.  This is the
  exact sense in which the support condition defining \(\PLlau\) and the one
  defining \(\mathfrak D(\K)\) agree, and \textup{(ii)} is the exact sense in
  which they differ.
\end{enumerate}
\end{theorem}

\begin{proof}
(i)  By \cref{plt:thm:du-two-substitutions}(i), \(\sigma\) sends
\(\sum_{n\ge n_{0}}p_{n}(L)t^{n}\) to \(\sum_{n\ge n_{0}}p_{n}(\Lambda)x^{n}\),
whose support is a subset of \(\IntegerNumbers\) bounded below, hence
left-finite; conversely a left-finite \(S\subseteq\IntegerNumbers\) is bounded
below and \(\sum_{\alpha\in S}x^{\alpha}p_{\alpha}\) is the \(\sigma\)-image of
the corresponding element of \(\PLlau\).  The same argument with
\(s^{-1}\IntegerNumbers\) gives the Puiseux case.  Strictness of the first
inclusion: \(x^{1/2}\).  Of the second: \(x^{\sqrt2}\) lies in
\(\mathfrak D(\K)\) but its exponent lies in no \(s^{-1}\IntegerNumbers\).

(ii)  The displayed series has support
\(\SetOf{1-1/k:k\ge1}\), which is well ordered of order type \(\omega\) and
therefore admissible in a Hahn ring, but its intersection with
\((-\infty,1)\) is infinite, so it is not left-finite; this is
\cref{plt:prop:ext-leftfinite}(ii).  Every left-finite set is well based,
since a strictly decreasing sequence would have infinitely many terms below its
first, so the inclusion holds and is strict.

(iii)  A subset of a discrete subgroup \(\Gamma\subseteq\RealNumbers\) is
bounded below if and only if it is left-finite, because \(\Gamma\) itself is
left-finite above any bound; and bounded below is equivalent to well based in a
discrete group for the same reason.
\end{proof}

\begin{warningbox}[title={``Left-finite'' is the Dulac condition, ``well-based'' is not}]
\label{plt:warn:du-support}
The condition \(\alpha_{j}\to+\infty\) that every source imposes on a Dulac
series is left-finiteness, not well-basedness.  The distinction is invisible in
\(\PLlau\), where the two coincide by
\cref{plt:thm:du-dulac-vs-class}(iii), and is essential the moment real
% ed.: this read "has no finite truncation Trunc_{<1} other than 0", which is
% doubly wrong: Trunc_{<N} of that series IS finite (and in general nonzero)
% for every N < 1, and Trunc_{<1} is the whole infinite series, not 0.  The
% correct statement is that truncation stops being finite at the cutoff 1.
exponents enter: a Hahn series such as \(\sum_{k\ge1}x^{1-1/k}\) has a
perfectly legitimate support, a leading term, and a well-defined product, but
its truncation \(\Trunc{\cdot}{N}\) is an infinite sum for every \(N\ge1\), and
already \(\Trunc{\cdot}{1}\) is the whole series --- so every algorithm
of \cref{plt:sec:algorithms}, every finite-determination theorem, and the very
notion of ``computing the expansion to order \(N\)'' fails for it.  Dulac
theory needs left-finiteness precisely because it needs those.
\end{warningbox}

\begin{definition}[Dulac germ]
\label{plt:def:du-dulac-germ}
Let \(\K=\RealNumbers\).  A germ \(f\colon(0,x_{0})\to\RealNumbers\) is a
\emph{Dulac germ} if there is \(\hat f\in\mathfrak D_{>0}(\RealNumbers)\) with
\(f\approx_{0}\hat f\) in the sense of \cref{plt:def:du-asymptotic}.  We call
\(\hat f\)
the \emph{Dulac series of \(f\)}; it is unique when it exists, and we write
\(\nu_{x}(f)\DefinitionEquals\nu_{x}(\hat f)>0\) for its leading exponent.
\end{definition}

\begin{proposition}[Uniqueness, and the flat kernel]
\label{plt:prop:du-flat-kernel}
Let \(f,g\) be germs at \(0^{+}\) and
\(\hat f,\hat g\in\mathfrak D(\RealNumbers)\).
\begin{enumerate}[label=\textup{(\roman*)}]
\item If \(f\approx_{0}\hat f\) and \(f\approx_{0}\hat g\) then
  \(\hat f=\hat g\).
\item \(f\approx_{0}0\) if and only if
  \(f(x)=\BigOAt{x^{N}}{x\to0^{+}}\) for every \(N\).  Such germs exist and are
  nonzero: \(\EulerE^{-1/x}\approx_{0}0\).
% ed.: this item asserted that f |-> f-hat is a ring homomorphism ONTO
% D(R).  Surjectivity is a Borel--Ritt statement for left-finite real
% exponent sets; it is not proved here and is not needed for anything below,
% so the claim is downgraded to "into" and the gap is named.
\item If \(f\approx_{0}\hat f\) and \(g\approx_{0}\hat g\) then
  \(f+g\approx_{0}\hat f+\hat g\) and \(fg\approx_{0}\hat f\hat g\).
  Consequently the germs admitting a Dulac series form a ring, the map
  \(f\mapsto\hat f\) is a ring homomorphism \emph{into}
  \(\mathfrak D(\RealNumbers)\), and by \textup{(ii)} its kernel --- the flat
  germs --- is a nonzero ideal.  The map is therefore \emph{not} injective.
  Whether it is surjective --- a Borel--Ritt theorem for left-finite real
  exponent sets, of which \cref{plt:thm:an-borel-ritt} is the integer-exponent
  case at infinity --- is not decided here and is used nowhere below.
\end{enumerate}
\end{proposition}

\begin{proof}
(i)  Let \(\hat h=\hat f-\hat g\ne0\), with \(\alpha_{0}=\nu_{x}(\hat h)\) and
leading block \(p\ne0\).  Because the support of \(\hat h\) is left-finite,
there is \(\epsilon>0\) with no support element in
\((\alpha_{0},\alpha_{0}+\epsilon)\); put \(N=\alpha_{0}+\epsilon\), so that
\(\Trunc{\hat h}{N}=x^{\alpha_{0}}p(\Lambda)\).  Subtracting the two estimates
\eqref{plt:eq:du-asymptotic} at this \(N\) and using linearity of
\(\Trunc{\cdot}{N}\) gives
\(\AbsoluteValue{\Trunc{\hat h}{N}(x)}\le Cx^{N}\Lambda^{M}\).  But
\(\AbsoluteValue{x^{\alpha_{0}}p(\Lambda)}\ge c\,x^{\alpha_{0}}\Lambda^{\deg p}\)
for small \(x\), with \(c>0\) half the absolute value of the leading
coefficient of \(p\), and
\(x^{\alpha_{0}}\Lambda^{\deg p}/(x^{\alpha_{0}+\epsilon}\Lambda^{M})
=x^{-\epsilon}\Lambda^{\deg p-M}\to+\infty\).  Contradiction.

(ii)  If \(\hat f=0\) then \(\Trunc{\hat f}{N}=0\) and
\eqref{plt:eq:du-asymptotic} is exactly
\(\AbsoluteValue{f}\le C_{N}x^{N}\Lambda^{M_{N}}\), which for every \(N'\) is
\(\BigOAt{x^{N'}}{x\to0^{+}}\) upon taking \(N=N'+1\); the converse is
immediate.  Finally \(\EulerE^{-1/x}x^{-N}\to0\) for every \(N\).

(iii)  Additivity is the triangle inequality applied to
\eqref{plt:eq:du-asymptotic}, since \(\Trunc{\cdot}{N}\) is linear: take
\(C_{N}\) the sum and \(M_{N}\) the maximum of the two data.

For multiplicativity, first note that if \(\hat h\in\mathfrak D(\RealNumbers)\)
is nonzero with \(b=\nu_{x}(\hat h)\) and \(h\approx_{0}\hat h\), then
\(\Trunc{\hat h}{N'}\) is, for each \(N'\), a finite sum
\(\sum_{b\le\gamma<N'}x^{\gamma}q_{\gamma}(\Lambda)\), whence
\begin{equation}\label{plt:eq:du-crude-bound}
  \AbsoluteValue{\bigl(\Trunc{\hat h}{N'}\bigr)(x)}\le C\,x^{b}\Lambda^{M},
  \qquad
  \AbsoluteValue{h(x)}\le C'\,x^{b}\Lambda^{M'}
\end{equation}
for small \(x\), the second by adding the remainder at \(N'=b+1\).  Now let
\(a=\nu_{x}(\hat f)\), \(b=\nu_{x}(\hat g)\), both finite, fix
\(N\in\RealNumbers\) and put \(N_{1}=N-b\), \(N_{2}=N-a\).  Then
\[
  fg-\bigl(\Trunc{\hat f}{N_{1}}\bigr)\bigl(\Trunc{\hat g}{N_{2}}\bigr)
  =\bigl(f-\Trunc{\hat f}{N_{1}}\bigr)g
   +\bigl(\Trunc{\hat f}{N_{1}}\bigr)\bigl(g-\Trunc{\hat g}{N_{2}}\bigr),
\]
and by \eqref{plt:eq:du-asymptotic} together with
\eqref{plt:eq:du-crude-bound} the two summands are
\(\BigOAt{x^{N_{1}+b}\Lambda^{\ast}}{x\to0^{+}}\) and
\(\BigOAt{x^{a+N_{2}}\Lambda^{\ast}}{x\to0^{+}}\), both
\(\BigOAt{x^{N}\Lambda^{\ast}}{x\to0^{+}}\).  Finally
\(\bigl(\Trunc{\hat f}{N_{1}}\bigr)\bigl(\Trunc{\hat g}{N_{2}}\bigr)
-\Trunc{\hat f\hat g}{N}\) is a finite sum of monomials \(x^{\alpha+\beta}\)
with \(\alpha+\beta\ge N\): a pair \((\alpha,\beta)\) with \(\alpha+\beta<N\),
\(\alpha\ge a\), \(\beta\ge b\) automatically has \(\alpha<N-b=N_{1}\) and
\(\beta<N-a=N_{2}\), so no term of \(\Trunc{\hat f\hat g}{N}\) is missing from
the product of truncations.  Hence this difference is
\(\BigOAt{x^{N}\Lambda^{\ast}}{x\to0^{+}}\) as well, and \(fg\approx_{0}\hat f\hat g\).
If \(\hat f=0\), then \(f=\BigOAt{x^{N}}{x\to0^{+}}\) for every \(N\) by (ii)
and \(\AbsoluteValue g\le C'x^{b}\Lambda^{M'}\) by
\eqref{plt:eq:du-crude-bound}, so \(fg\approx_{0}0=\hat f\hat g\).
\end{proof}

\begin{warningbox}[title={A Dulac series is not a Dulac germ}]
\label{plt:warn:du-quasianalytic}
In \(\PLlau\) the series \emph{is} the object.  At a finite point this is
false in both directions.  A Dulac series is typically divergent, so it defines
no germ; and by \cref{plt:prop:du-flat-kernel}(ii) a germ is not determined by
its Dulac series, since exponentially flat germs are invisible to it.  That the
germs \emph{arising from analytic vector fields} are nevertheless determined by
their Dulac series --- quasi-analyticity of the corresponding class --- is a
deep theorem, equivalent in strength to the non-accumulation of limit cycles on
a polycycle, and is the content of the resolution of Dulac's problem by
\'Ecalle \cite{ecalle} and, independently, by Il'yashenko; see
\cite{resman-dulac,mardesic-length} for the modern account and for the
associated \(\varepsilon\)-neighbourhood invariants.  \emph{No} formal
statement in this section, and no amount of algebra in \(\mathfrak D(\K)\),
implies it.
\end{warningbox}

\begin{theorem}[Composition of Dulac series]
\label{plt:thm:du-composition-closure}
Let \(\K=\RealNumbers\).  Let
\(\hat f=\sum_{\alpha\in S}x^{\alpha}p_{\alpha}(\Lambda)\in\mathfrak D(\K)\) and
let
\[
  \hat g=c\,x^{\beta}\bigl(1+u\bigr),
  \qquad c>0,\quad\beta>0,\quad u\in\mathfrak D(\K),\quad\nu_{x}(u)>0 .
\]
Define
\begin{equation}\label{plt:eq:du-composition}
  \hat f\compose\hat g
  \DefinitionEquals
  \sum_{\alpha\in S}c^{\alpha}x^{\alpha\beta}(1+u)^{\alpha}\,
  p_{\alpha}\bigl(\beta\Lambda-\log c-\log(1+u)\bigr),
\end{equation}
where \((1+u)^{\alpha}=\sum_{k\ge0}\binom{\alpha}{k}u^{k}\) and
\(\log(1+u)=\sum_{k\ge1}\frac{(-1)^{k-1}}{k}u^{k}\).  Then:
\begin{enumerate}[label=\textup{(\roman*)}]
\item the family in \eqref{plt:eq:du-composition} is summable and
  \(\hat f\compose\hat g\in\mathfrak D(\K)\), the set of \(x\)-exponents
  occurring in it satisfying
  \begin{equation}\label{plt:eq:du-composition-support}
    \SetOf{\text{\(x\)-exponents of }\hat f\compose\hat g}
    \subseteq\beta S+\Gamma_{u},
  \end{equation}
  \(\Gamma_{u}\) denoting the sub-semigroup of \(\RealNumbers_{\ge0}\)
  generated by the \(x\)-support of \(u\);
\item if \(S\subseteq\alpha_{0}+\Gamma_{1}\) and the support of \(u\) lies in a
  finitely generated \(\Gamma_{2}\subseteq\RealNumbers_{>0}\), with
  \(\Gamma_{1}\) finitely generated, then \(\hat f\compose\hat g\) is carried
  by the grid \(\alpha_{0}\beta+\Gamma_{3}\), where \(\Gamma_{3}\) is generated
  by \(\beta\cdot(\text{generators of }\Gamma_{1})\) together with the
  generators of \(\Gamma_{2}\); in particular \(\Gamma_{3}\) is finitely
  generated;
\item \(\nu_{x}(\hat f\compose\hat g)=\beta\,\nu_{x}(\hat f)\), the leading
  block is \(c^{\alpha_{0}}p_{\alpha_{0}}(\beta\Lambda-\log c)\) with
  \(\alpha_{0}=\nu_{x}(\hat f)\), and
  \(\deg_{\Lambda}\) of the leading block equals \(\deg p_{\alpha_{0}}\);
\item \(\hat f\mapsto\hat f\compose\hat g\) is an injective \(\K\)-algebra
  homomorphism of \(\mathfrak D(\K)\) into itself, continuous for the
  filtration.
\end{enumerate}
\end{theorem}

\begin{proof}
Write \(T\) for the support of \(u\); it is left-finite and contained in
\(\RealNumbers_{>0}\), so by \cref{plt:prop:ext-leftfinite}(iv) the semigroup
\(\Gamma_{u}=T^{\ast}\) is left-finite and every element of it has only
finitely many representations as a sum of elements of \(T\).

\emph{The individual summands.}  Since \(\nu_{x}(u)>0\), the families
\((\binom\alpha k u^{k})_{k\ge0}\) and
\((\frac{(-1)^{k-1}}ku^{k})_{k\ge1}\) are summable in \(\mathfrak D(\K)\) ---
\(\nu_{x}(u^{k})\ge k\nu_{x}(u)\to+\infty\) --- so \((1+u)^{\alpha}\) and
\(\log(1+u)\) are defined, with supports contained in \(\Gamma_{u}\).
Substituting the element \(\beta\Lambda-\log c-\log(1+u)\) into the polynomial
\(p_{\alpha}\) is a finite algebraic operation, so
\(p_{\alpha}(\beta\Lambda-\log c-\log(1+u))\) lies in \(\mathfrak D(\K)\) with
support in \(\Gamma_{u}\) and with \(\Lambda\)-degree \(\deg p_{\alpha}\) at
the exponent \(0\).  Hence the \(\alpha\)-th summand of
\eqref{plt:eq:du-composition} has support in \(\alpha\beta+\Gamma_{u}\).

\emph{Summability.}  Fix \(M\in\RealNumbers\).  A summand can contribute below
\(M\) only if \(\alpha\beta\le M\, \), i.e.\ \(\alpha\le M/\beta\); since \(S\)
is left-finite and \(\beta>0\), only finitely many \(\alpha\) qualify.  Each of
those contributes a left-finite set of exponents.  Therefore the total support
meets \((-\infty,M)\) in a finite set, the sum is defined coefficientwise, and
\eqref{plt:eq:du-composition-support} holds; the result lies in
\(\mathfrak D(\K)\).

(ii)  \(\beta S+\Gamma_{u}\subseteq\beta\alpha_{0}+\beta\Gamma_{1}+\Gamma_{2}\),
and \(\beta\Gamma_{1}+\Gamma_{2}=\Gamma_{3}\) is generated by the displayed
finite set; \cref{plt:lem:du-fg-semigroup} makes it left-finite.

(iii)  The exponent \(\alpha_{0}\beta\) is attained only by the summand
\(\alpha=\alpha_{0}\) and, within it, only by the \(\Gamma_{u}\)-degree-zero
part, which is \(c^{\alpha_{0}}p_{\alpha_{0}}(\beta\Lambda-\log c)\); this is
nonzero of the same \(\Lambda\)-degree as \(p_{\alpha_{0}}\) because
\(\beta\ne0\).

(iv)  Additivity is clear from \eqref{plt:eq:du-composition}.  For
multiplicativity, write \(\hat g^{\alpha}\DefinitionEquals
c^{\alpha}x^{\alpha\beta}(1+u)^{\alpha}\) and
\(\Lambda\compose\hat g\DefinitionEquals\beta\Lambda-\log c-\log(1+u)\), so
that \eqref{plt:eq:du-composition} is the \(\K\)-linear extension of
\[
  x^{\alpha}\Lambda^{j}\longmapsto
  \hat g^{\alpha}\,\bigl(\Lambda\compose\hat g\bigr)^{j}
\]
over the (summable) monomial decomposition of \(\hat f\).  This assignment is
multiplicative on monomials: \(c^{\alpha+\alpha'}=c^{\alpha}c^{\alpha'}\), and
\((1+u)^{\alpha}(1+u)^{\alpha'}=(1+u)^{\alpha+\alpha'}\) because the
coefficient of \(u^{k}\) in the product is
\(\sum_{i+j=k}\binom\alpha i\binom{\alpha'}j=\binom{\alpha+\alpha'}k\) by the
Chu--Vandermonde identity, which is a polynomial identity in
\(\alpha,\alpha'\).  Multiplicativity on all of \(\mathfrak D(\K)\) then follows
by expanding both sides as summable double families and regrouping, which the
finiteness established above licenses.  Continuity and injectivity follow from
(iii), which gives \(\nu_{x}(\hat f\compose\hat g)=\beta\nu_{x}(\hat f)\).
\end{proof}

\begin{corollary}[The class does not depend on the analytic chart]
\label{plt:cor:du-chart-independence}
Let \(\varphi\) be a germ at \(0\) of a real-analytic diffeomorphism with
\(\varphi(0)=0\) and \(\varphi'(0)>0\).  Then
\(\hat f\mapsto\hat f\compose\hat\varphi\) maps \(\mathfrak D(\RealNumbers)\)
into itself, preserves \(\nu_{x}\) and the leading \(\Lambda\)-degree, and
carries a grid \(\alpha_{0}+\Gamma\) to the grid
\(\alpha_{0}+\Gamma'\) with \(\Gamma'\) generated by the generators of
\(\Gamma\) together with \(1\).  Consequently the notions ``Dulac series'',
``leading exponent'' and ``leading logarithmic degree'' do not depend on the
analytic parametrisation of the transversal used to define them; the
corresponding statement for Dulac \emph{germs} follows from the analytic
composition clause quoted in \cref{plt:thm:du-dulac-transition}.
\end{corollary}

\begin{proof}
Write \(\varphi(x)=x\psi(x)\) with \(\psi\) analytic and
\(c\DefinitionEquals\psi(0)=\varphi'(0)>0\).  Then
\(\hat\varphi=c\,x^{1}\bigl(1+u\bigr)\) with
\(u=\psi(x)/c-1\in x\RealNumbers[[x]]\), whose support lies in
\(\PositiveIntegers\); so \(\beta=1\) and \(\Gamma_{2}\) is the semigroup
generated by \(1\), and
\cref{plt:thm:du-composition-closure} applies, giving all the assertions.  The
last sentence follows because two analytic parametrisations of the same
transversal differ by such a \(\varphi\), possibly after replacing
\(\varphi\) by \(\CompositionalInverseOf{\varphi}\), which is again analytic
with positive derivative.
\end{proof}

\begin{proposition}[The two logarithmic conventions]
\label{plt:prop:du-two-log-conventions}
Recall \(\ell=\Lambda^{-1}\).  Then, inside the Hahn ambient \(\mathcal Z\):
\begin{enumerate}[label=\textup{(\roman*)}]
\item \(\K[\Lambda]\cap\K[[\ell]]=\K\); the Dulac coefficient ring
  \(\K[\Lambda]\) and the normal-form coefficient ring \(\K[[\ell]]\) meet only
  in the constants.
\item The smallest of the four classes of \cref{plt:def:du-generators}
  containing both is \(\K((\ell))((x))=\sigma(\PLit)\), the transported
  iterated Laurent-log field: its coefficient ring \(\K((\Lambda^{-1}))\)
  contains \(\K[\Lambda]\) and \(\K[[\ell]]\), and any ring containing both
  contains their sums.
\item Admitting real exponents of the logarithm, i.e.\ monomials
  \(x^{\alpha}\ell^{\gamma}=x^{\alpha}\Lambda^{-\gamma}\) with
  \(\gamma\in\RealNumbers\) and well-based \(\gamma\)-support, produces the
  rank-two Hahn ring \(\sigma\bigl(\mathcal H_{1}(\K,\RealNumbers)\bigr)\)
  refined by real logarithmic exponents; both conventions embed in it.
\item In the dominance order \eqref{plt:eq:du-zero-order},
  \(\ell^{\gamma}\precdom\ell^{\gamma'}\) if and only if \(\gamma>\gamma'\).
  A list of monomials \(x^{\alpha}\ell^{m}\) written by increasing \(m\) is
  therefore written by \emph{decreasing} size, while a list of
  \(x^{\alpha}\Lambda^{j}\) written by increasing \(j\) is written by
  increasing size.
\end{enumerate}
\end{proposition}

\begin{proof}
(i)  In \(\mathcal Z\) a polynomial in \(\Lambda\) has support in
\(\SetOf0\times\NaturalNumbers\times\SetOf0\) and an element of \(\K[[\ell]]\)
has support in \(\SetOf0\times(-\NaturalNumbers)\times\SetOf0\); the
intersection of the supports is \(\SetOf{(0,0,0)}\), and uniqueness of the
Hahn representation gives the claim.

(ii)  \(\K((\Lambda^{-1}))\) consists of the series
\(\sum_{j\ge j_{0}}a_{j}\Lambda^{-j}\) with \(j_{0}\in\IntegerNumbers\); taking
\(j_{0}\le0\) gives every polynomial in \(\Lambda\), and taking \(j_{0}=0\)
gives \(\K[[\ell]]\).  Among \(\K[\Lambda]\), \(\K(\Lambda)\) and
\(\K((\Lambda^{-1}))\), the first does not contain \(\ell\), and the second
does not contain \(\sum_{j\ge0}j!\,\ell^{j}\).  Indeed, the Laurent expansion
at \(\Lambda=\infty\) of a rational function \(P/Q\) has coefficients
satisfying the linear recurrence with constant coefficients determined by
\(Q\), whereas \((j!)\) satisfies none: if
\(\sum_{i=0}^{d}c_{i}(j+i)!=0\) for all large \(j\), division by \(j!\) gives
\(\sum_{i=0}^{d}c_{i}(j+1)(j+2)\cdots(j+i)=0\) for all large \(j\), a
polynomial identity in \(j\) whose leading coefficient is \(c_{d}\), so
\(c_{d}=0\) and, inductively, every \(c_{i}=0\).  So \(\K((\Lambda^{-1}))\) is
the smallest of the three.

(iii)  Immediate from \eqref{plt:eq:du-small-log}.

(iv)  \(\ell^{\gamma}=\Lambda^{-\gamma}\), and by \eqref{plt:eq:du-zero-order}
\(\Lambda^{b}\succdom\Lambda^{b'}\) iff \(b>b'\); put \(b=-\gamma\).
\end{proof}

\section{The transition map of a hyperbolic sector}
\label{plt:sec:du-saddle}

The reason the class of \cref{plt:def:du-dulac-series} is the right one is a
theorem about planar vector fields.  We state it precisely, we prove the model
case, and we mark clearly which part is quoted.

\begin{example}[The linear saddle: an exact power law, and where the logarithm
hides]
\label{plt:ex:du-linear-saddle}
Consider \(\dot u=\lambda_{u}u\), \(\dot v=-\lambda_{s}v\) on
\(\RealNumbers^{2}\) with \(\lambda_{u},\lambda_{s}>0\), and put
\[
  r\DefinitionEquals\frac{\lambda_{s}}{\lambda_{u}}>0
\]
(the \emph{hyperbolicity ratio}).  Fix \(\varepsilon>0\) and take the
transversals \(\Sigma^{\mathrm{in}}=\SetOf{v=\varepsilon,\ 0<u<\varepsilon}\)
with coordinate \(x=u\) and
\(\Sigma^{\mathrm{out}}=\SetOf{u=\varepsilon,\ 0<v<\varepsilon}\) with
coordinate \(y=v\).  The orbit through \((x,\varepsilon)\) is
\(u(\theta)=x\EulerE^{\lambda_{u}\theta}\),
\(v(\theta)=\varepsilon\EulerE^{-\lambda_{s}\theta}\); it meets
\(\Sigma^{\mathrm{out}}\) at the transit time
\begin{equation}\label{plt:eq:du-transit-time}
  T(x)=\frac1{\lambda_{u}}\log\frac{\varepsilon}{x}
      =\frac{1}{\lambda_{u}}\bigl(\Lambda+\log\varepsilon\bigr),
\end{equation}
and the transition map is
\begin{equation}\label{plt:eq:du-linear-transition}
  D(x)=\varepsilon\EulerE^{-\lambda_{s}T(x)}
      =\varepsilon^{1-r}x^{r},
\end{equation}
an exact power law with no logarithm.  The logarithm is nevertheless present:
by \eqref{plt:eq:du-transit-time} the time spent near the saddle is exactly
\(\Lambda/\lambda_{u}\) up to a constant, and it is the exact exponential
cancellation in \eqref{plt:eq:du-linear-transition} that removes it.  For an
analytic saddle that is not linearisable the cancellation is only approximate,
and the surviving discrepancy is what produces the polynomials \(p_{j}\) in
\cref{plt:thm:du-dulac-transition}.
\end{example}

\begin{theorem}[Dulac; quoted, not proved here]
\label{plt:thm:du-dulac-transition}
Let \(\mathcal X\) be a real-analytic vector field on a neighbourhood of
\(0\in\RealNumbers^{2}\) with a hyperbolic saddle at \(0\), of eigenvalues
\(\lambda_{u}>0>-\lambda_{s}\) and hyperbolicity ratio
\(r=\lambda_{s}/\lambda_{u}>0\).  Let \(\Sigma^{\mathrm{in}}\) and
\(\Sigma^{\mathrm{out}}\) be analytic arcs transverse to the stable and to the
unstable separatrix, analytically parametrised by \(x\ge0\) and \(y\ge0\) with
the separatrix at the parameter value \(0\).  Then the transition map
\(D\colon(0,\varepsilon)\to(0,\infty)\) obtained by following the flow from
\(\Sigma^{\mathrm{in}}\) to \(\Sigma^{\mathrm{out}}\) is a Dulac germ in the
sense of \cref{plt:def:du-dulac-germ}: there are \(c>0\), a left-finite set
\(\SetOf{r=\alpha_{0}<\alpha_{1}<\cdots}\subseteq\RealNumbers_{>0}\) contained
in a finitely generated sub-semigroup of \(\RealNumbers_{>0}\), and
polynomials \(p_{j}\in\RealNumbers[\Lambda]\) with \(p_{0}=c\) constant, such
that
\begin{equation}\label{plt:eq:du-dulac-transition}
  D\approx_{0}\hat D,
  \qquad
  \hat D(x)=c\,x^{r}+\sum_{j\ge1}x^{\alpha_{j}}p_{j}(\Lambda).
\end{equation}
Moreover the composition of two Dulac germs whose leading exponents are
positive is again a Dulac germ, with Dulac series the formal composite
\eqref{plt:eq:du-composition}.
\end{theorem}

\begin{remark}[Provenance, and what we have and have not verified]
\label{plt:rmk:du-quoted}
\Cref{plt:thm:du-dulac-transition} is classical: it originates with Dulac and
is proved in the modern finiteness literature; the accounts we have used are
\cite{resman-dulac,mardesic-length}, with the normal-form machinery that makes
the exponent grid explicit in
\cite{mardesic-normalforms,peran-logtransseries}.  Three points of hygiene.
\begin{enumerate}[label=\textup{(\alph*)}]
\item We have \emph{proved} the model case
  (\cref{plt:ex:du-linear-saddle}), the closure of the formal class under the
  composition \eqref{plt:eq:du-composition}
  (\cref{plt:thm:du-composition-closure}), and the chart-independence
  (\cref{plt:cor:du-chart-independence}).  What is quoted is the existence of
  the asymptotic expansion for an analytic saddle and the analytic half of the
  composition statement, which requires uniform remainder control of the kind
  \cref{plt:thm:cmp-analytic-transfer} supplies at infinity and which we have
  not carried out at a finite point.
\item The sources record further structure that we quote but do not use: that
  the exponents may be taken in the sub-semigroup of \(\RealNumbers_{>0}\)
  generated by \(1\) and \(r\), and that in the non-resonant case
  \(r\notin\RationalNumbers\) every \(p_{j}\) is constant, so that the
  expansion has no logarithms at all.  We do
  not assert these; only the left-finiteness and the finite generation of the
  grid are used below, and both are consequences of the displayed statement.
\item The polynomials are usually printed as \(P_{j}(\ln x)\) in that
  literature.  The translation to our normalisation is
  \(p_{j}(u)=P_{j}(-u)\), \cref{plt:prop:du-traps}\,(T7): the degrees agree and
  every odd-degree coefficient changes sign.
\end{enumerate}
\end{remark}

\begin{corollary}[The first-return map of a hyperbolic polycycle]
\label{plt:cor:du-return-map}
Let \(\gamma\) be a hyperbolic polycycle of a real-analytic planar vector
field: a cyclically ordered finite family of hyperbolic saddles
\(s_{1},\dots,s_{k}\), with ratios \(r_{1},\dots,r_{k}>0\), joined by regular
orbit arcs.  Let \(\Sigma\) be an analytic transversal at a regular point of
\(\gamma\), analytically parametrised by \(x\ge0\), and let \(P\) be the
first-return map on the side from which the orbits approach \(\gamma\).  Then
\(P\) is a Dulac germ.  Its Dulac series \(\hat P\) is the formal composite
\[
  \hat P=\hat R_{k}\compose\hat D_{k}\compose\cdots\compose
         \hat R_{1}\compose\hat D_{1},
\]
carried by a finitely generated grid, with
\begin{equation}\label{plt:eq:du-return-leading}
  \nu_{x}(\hat P)=r\DefinitionEquals\prod_{i=1}^{k}r_{i},
  \qquad
  \bigl[x^{r}\bigr]\hat P\in\RealNumbers_{>0}
  \ \text{a constant},
\end{equation}
% ed.: "P-hat is independent of the parametrisation" would be false as
% literally read -- a change of chart conjugates P-hat.  What is invariant is
% the membership and the leading data, which is what cor:du-chart-independence
% supplies.
and a change of the analytic parametrisation of \(\Sigma\) conjugates
\(\hat P\) by an analytic chart change \(\varphi\) as in
\cref{plt:cor:du-chart-independence}, which changes neither membership in
\(\mathfrak D(\RealNumbers)\), nor \(\nu_{x}(\hat P)=r\), nor the leading
logarithmic degree, nor the finite generation of the grid.  In particular, if
\(r\ne1\) then \(\hat P(x)-x\) has leading term
\(\bigl([x^{r}]\hat P\bigr)x^{r}\) when \(r<1\) and \(-x\) when \(r>1\); if
\(r=1\) the leading block of \(\hat P-x\) is not determined by the ratios
alone and may carry a logarithm.
\end{corollary}

\begin{proof}
Each corner map \(D_{i}\) is a Dulac germ with leading exponent \(r_{i}>0\), by
\cref{plt:thm:du-dulac-transition}; each regular transition \(R_{i}\) is a germ
of an analytic diffeomorphism fixing \(0\) with positive derivative, hence a
Dulac germ with \(\hat R_{i}\in x\RealNumbers[[x]]\), leading exponent \(1\)
and support in \(\PositiveIntegers\) (a convergent Taylor series is
\(\approx_{0}\) itself).  That \(P\) is the displayed composite of germs is the
definition of the return map.  By the last clause of
\cref{plt:thm:du-dulac-transition} the composite is a Dulac germ with Dulac
series the formal composite; by \cref{plt:thm:du-composition-closure}(i),(ii)
the formal composite lies in \(\mathfrak D(\RealNumbers)\) and is carried by a
finitely generated grid, since each factor is; by
\cref{plt:thm:du-composition-closure}(iii) applied \(2k\) times the leading
exponent multiplies, giving \(\prod_{i}r_{i}\cdot 1^{k}=r\).  The leading block
is a positive constant by the same induction: at each stage the outer leading
block is a positive constant \(\kappa\) and the inner leading coefficient a
positive constant \(c\), and \cref{plt:thm:du-composition-closure}(iii) gives
the composite leading block \(c^{\alpha_{0}}\kappa>0\).  Independence of the
chart is
\cref{plt:cor:du-chart-independence}.  The final sentence compares
\(\nu_{x}(\hat P)=r\) with \(\nu_{x}(x)=1\) in the dominance order
\eqref{plt:eq:du-zero-order}.
\end{proof}

\begin{remark}[Dulac's problem, and the exact place formal algebra stops]
\label{plt:rmk:du-dulac-problem}
Limit cycles accumulating on \(\gamma\) correspond to zeros of \(P(x)-x\)
accumulating at \(0\).  If \(\hat P(x)\ne x\), the formal algebra of
\cref{plt:prop:du-dulac-ring} shows that \(\hat P(x)-x\) has a leading block
\(x^{\alpha}p(\Lambda)\) with \(p\ne0\), and hence --- \emph{if} \(P\) were
determined by \(\hat P\) --- that \(P(x)-x\) has finitely many zeros near
\(0\).  The gap is exactly \cref{plt:prop:du-flat-kernel}: a germ whose Dulac
series is \(x\) may still oscillate, and nothing in
\(\mathfrak D(\RealNumbers)\) forbids it.  Closing that gap for germs of
analytic origin is Dulac's problem, solved by \'Ecalle \cite{ecalle} and by
Il'yashenko.  The algebraic content of this section --- closure of the class,
the grid, the leading data --- is what makes the statement of the problem
precise; it is not, and cannot be, an ingredient of its solution.  Compare the
o-minimal route: the unary germs definable in
\(\RealNumbers_{\mathrm{an},\exp}\) form a Hardy field whose asymptotic
expansions live in the field of logarithmic--exponential series
\cite{vdDMM-1997,vdDMM-2001}, so for a germ known to be definable the
quasi-analyticity is automatic; we make no claim about the definability of a
general Dulac germ.
\end{remark}

\section{Enlargements forced at a finite point}
\label{plt:sec:du-enlargements}

At infinity the enlargements of \cref{plt:sec:scale-change} were forced by
composition.  At a finite point three further mechanisms operate, and the third
is the one the sources treat least carefully.

\begin{proposition}[Real and ramified exponents]
\label{plt:prop:du-puiseux-zero}
Let \(r>0\) be the hyperbolicity ratio of \cref{plt:ex:du-linear-saddle} and
let \(D(x)=cx^{r}\) be the model transition map.
\begin{enumerate}[label=\textup{(\roman*)}]
\item If \(r\in\PositiveIntegers\) then \(\hat D\in\sigma(\PLlau)\).
\item If \(r=p/q\in\RationalNumbers_{>0}\) in lowest terms, then
  \(\hat D\in\sigma\bigl(\K[L]((t^{1/q}))\bigr)\) and in no
  \(\sigma\bigl(\K[L]((t^{1/s}))\bigr)\) with \(q\nmid s\); the Puiseux-log
  ramification index is exactly \(q\).
% ed.: "is the semigroup generated by r and 1" asserts an equality that the
% grid bookkeeping of plt:thm:du-composition-closure(ii) does not give; it
% gives a containment, which is all that is used.
\item If \(r\notin\RationalNumbers\) then \(\hat D\) lies in no Puiseux-log
  ring at all.  The grid produced by composing \(\hat D\) with analytic charts
  on either side is contained, by
  \cref{plt:cor:du-chart-independence} and
  \cref{plt:thm:du-composition-closure}(ii), in the semigroup generated by
  \(r\) and \(1\); it is left-finite by \cref{plt:lem:du-fg-semigroup}, but the group
  \(\IntegerNumbers+r\IntegerNumbers\) it generates is dense in
  \(\RealNumbers\), so no ramification index bounds it.
\end{enumerate}
\end{proposition}

\begin{proof}
(i) and the membership in (ii) are
\cref{plt:thm:du-dulac-vs-class}(i).  For the sharpness in (ii): the exponent
\(p/q\) lies in \(s^{-1}\IntegerNumbers\) if and only if \(sp/q\in
\IntegerNumbers\), i.e.\ if and only if \(q\mid sp\), i.e.\ --- since
\(\gcd(p,q)=1\) --- if and only if \(q\mid s\).  For (iii), an irrational
\(r\) lies in no \(s^{-1}\IntegerNumbers\); and
\(\IntegerNumbers+r\IntegerNumbers\) is a non-cyclic subgroup of
\(\RealNumbers\), hence dense.  This is
\cref{plt:thm:grp-scale-classification}(2),(3) read through \(\sigma\).
\end{proof}

\begin{theorem}[Two ways the logarithmic degree stops being finite]
\label{plt:thm:du-log-degree-failure}
Work in a Hahn ambient admitting monomials \(x^{a}\Lambda^{b}\) with
\(a,b\in\RealNumbers\) and well-based supports.
\begin{enumerate}[label=\textup{(\roman*)}]
\item \textup{(Division.)}  \(\bigl(x(\Lambda+1)\bigr)^{-1}
  =x^{-1}\sum_{k\ge0}(-1)^{k}\Lambda^{-k-1}\).  Its coefficient block is an
  infinite series in \(\Lambda^{-1}\); it lies in \(\sigma(\PLit)\) and in
  \(\sigma(\PLrat)\), but in no class with polynomial blocks.  The
  \emph{block} has ceased to be finite while every exponent is still an
  integer.
\item \textup{(Non-integer powers.)}  Let
  \(\rho\in\RealNumbers\setminus\NaturalNumbers\).  Then
  \[
    (\Lambda+1)^{\rho}
    =\Lambda^{\rho}\sum_{k\ge0}\binom{\rho}{k}\Lambda^{-k},
    \qquad
    \binom\rho k\ne0\ \text{for every }k\ge0 .
  \]
% ed.: the statement here read "it lies in none of sigma(PLlau),
% sigma(PLrat), sigma(PLit), nor in any Puiseux-log or Dulac class".  That is
% false for rho a negative integer, where item (i) exhibits the element inside
% sigma(PLrat) and sigma(PLit); the exclusion from those two classes needs
% rho not an integer, and only the exclusion from the polynomial-block classes
% holds for every rho outside the natural numbers.  Split accordingly.
  Consequently \(\bigl(x(\Lambda+1)\bigr)^{\rho}
  =x^{\rho}\Lambda^{\rho}\sum_{k}\binom\rho k\Lambda^{-k}\) has infinitely many
  nonzero logarithmic coefficients, hence lies in no class with polynomial
  coefficient blocks: not in \(\sigma(\PLlau)\), not in a Puiseux-log ring
  \(\sigma\bigl(\K[L]((t^{1/s}))\bigr)\), and not in \(\mathfrak D(\K)\).  If
  moreover \(\rho\notin\IntegerNumbers\), its leading logarithmic exponent
  \(\rho\) is not an integer, and it then lies in neither \(\sigma(\PLrat)\)
  nor \(\sigma(\PLit)\), every monomial of those two classes having
  logarithmic exponent in \(\IntegerNumbers\); for
  \(\rho\in\IntegerNumbers_{<0}\), by contrast, it does lie in both, as
  \textup{(i)} shows for \(\rho=-1\).  In the non-integer case the
  \emph{degree} has ceased to be a natural number.
\item Neither failure is exotic.  Failure \textup{(i)} is forced by a single
  division, since \(\Lambda+1\) is not a unit of \(\K[\Lambda]\)
  (\cref{plt:thm:div-unit-criterion} transported).  Failure \textup{(ii)} is
  forced by a single square root: the element \(x(\Lambda+1)\), a member of
  \(\sigma(\PLlau)\), has by \textup{(ii)} with \(\rho=1/2\) the square root
  \(x^{1/2}\Lambda^{1/2}\sum_{k\ge0}\binom{1/2}{k}\Lambda^{-k}\), which lies in
  no Dulac, Puiseux-log or iterated Laurent-log class.
\item Within \(\sigma(\PLlau)\) itself, where every block \emph{is} a
  polynomial, the degrees need not be bounded: \(\sum_{n\ge1}x^{n}\Lambda^{n}\)
  has \(\deg_{\Lambda}[x^{n}]=n\).  No statement in this volume may assume a
  uniform logarithmic degree.
\end{enumerate}
\end{theorem}

\begin{proof}
(i)  \((\Lambda+1)^{-1}=\Lambda^{-1}(1+\Lambda^{-1})^{-1}
=\sum_{k\ge0}(-1)^{k}\Lambda^{-k-1}\), the geometric series being summable in
the \(\Lambda^{-1}\)-adic topology.  All coefficients are \(\pm1\ne0\), so the
support is infinite; membership in \(\sigma(\PLit)\) is the definition of
\(\K((\Lambda^{-1}))\), and membership in \(\sigma(\PLrat)\) holds because
\((\Lambda+1)^{-1}\in\K(\Lambda)\).

(ii)  Write \(\Lambda+1=\Lambda(1+\Lambda^{-1})\); the binomial series
converges \(\Lambda^{-1}\)-adically and
\(\binom\rho k=\rho(\rho-1)\cdots(\rho-k+1)/k!\) vanishes only if
\(\rho\in\SetOf{0,1,\dots,k-1}\), which is excluded by
\(\rho\notin\NaturalNumbers\).  Since all these coefficients are nonzero, the
single \(x\)-block at the exponent \(\rho\) is an infinite series in
\(\Lambda\), so the element is not in any class whose blocks are polynomials
in \(\Lambda\), that is, not in \(\sigma(\PLlau)\), not in
\(\sigma\bigl(\K[L]((t^{1/s}))\bigr)\) and not in \(\mathfrak D(\K)\); Hahn
representations are unique, so no rewriting can repair this.  The leading
monomial is \(x^{\rho}\Lambda^{\rho}\); if \(\rho\notin\IntegerNumbers\) its
logarithmic exponent is not an integer, whereas every monomial of
\(\sigma(\PLrat)\) and of \(\sigma(\PLit)\) has logarithmic exponent in
\(\IntegerNumbers\), which excludes those two classes as well.  For
\(\rho\in\IntegerNumbers_{<0}\) no such exclusion is available, and indeed
\textup{(i)} places \(\rho=-1\) inside both.

(iii)  The reciprocal of \(x(\Lambda+1)\) is not in \(\sigma(\PLlau)\) because
\(\Lambda+1\) is not a unit of \(\K[\Lambda]\); this is
\cref{plt:thm:div-unit-criterion}.  The square-root assertion is
\textup{(ii)} with \(\rho=1/2\): the displayed element has leading monomial
\(x^{1/2}\Lambda^{1/2}\), whose logarithmic exponent \(1/2\) is not an integer,
and its logarithmic support is infinite because \(\binom{1/2}{k}\ne0\) for all
\(k\).

(iv)  Immediate; the series lies in \(\K[\Lambda][[x]]\) because each block is
a single monomial.
\end{proof}

\begin{proposition}[A logarithm in the inner monomial forces \(\Lambda_{2}\)]
\label{plt:prop:du-iterated-log-forced}
Let \(\hat g=c\,x^{\beta}\Lambda^{\gamma}(1+u)\) with \(c>0\), \(\beta>0\),
\(\gamma\in\RealNumbers\), \(\nu_{x}(u)>0\).  Then
\begin{equation}\label{plt:eq:du-inner-log}
  \Lambda\compose\hat g
  =\log\frac1{\hat g}
  =\beta\Lambda-\gamma\Lambda_{2}-\log c-\log(1+u),
\end{equation}
so that for \(\hat f=\sum_{\alpha}x^{\alpha}p_{\alpha}(\Lambda)\):
\begin{enumerate}[label=\textup{(\roman*)}]
\item if \(\gamma=0\), \cref{plt:thm:du-composition-closure} applies and
  \(\hat f\compose\hat g\) stays in \(\mathfrak D(\K)\);
% ed.: the second clause read "for every alpha with p_alpha nonzero and
% gamma*alpha not a natural number".  Only the summand of the SMALLEST alpha
% is protected from interference: for alpha' < alpha the summand of alpha'
% also contributes at the x-exponent alpha*beta (through its u-tail), with
% Lambda-exponents in gamma*alpha' + N, and these can meet gamma*alpha, so no
% cancellation argument is available for a general alpha.  Restricted to
% alpha_0 = nu_x(f-hat), where the leading block stands alone.
\item if \(\gamma\ne0\), the second logarithm \(\Lambda_{2}\) occurs in
  \(\hat f\compose\hat g\) if and only if \(p_{\alpha}\notin\K\) for some
  \(\alpha\); and if \(\alpha_{0}=\nu_{x}(\hat f)=\min S\) satisfies
  \(\gamma\alpha_{0}\notin\NaturalNumbers\), then a \(\Lambda\)-exponent
  outside \(\NaturalNumbers\) occurs.  In the first case
  \(\hat f\compose\hat g\) lies in no class of depth one at a finite point,
  that is, in no class whose monomials are \(x^{a}\Lambda^{b}\); in the second
  it lies in no depth-one class whose logarithmic exponents are natural
  numbers --- \(\K[\Lambda][[x]]\), \(\K[\Lambda]((x))\), the Puiseux-log
  rings \(\K[\Lambda]((x^{1/s}))\) and \(\mathfrak D(\K)\) --- and, if moreover
  \(\gamma\alpha_{0}\notin\IntegerNumbers\), in none of the four classes of
  \cref{plt:def:du-generators}.
\end{enumerate}
\end{proposition}

\begin{proof}
Formula \eqref{plt:eq:du-inner-log} is \(-\log\) of the product, using
\(\log x=-\Lambda\) and \(\log\Lambda=\Lambda_{2}\).  Item (i) is
\cref{plt:thm:du-composition-closure}.

For (ii), substitute the generator images into
\(\hat f=\sum_{\alpha}x^{\alpha}p_{\alpha}(\Lambda)\):
\[
  \hat f\compose\hat g
  =\sum_{\alpha\in S}c^{\alpha}x^{\alpha\beta}\Lambda^{\gamma\alpha}
   (1+u)^{\alpha}\,
   p_{\alpha}\bigl(\beta\Lambda-\gamma\Lambda_{2}-\log c-\log(1+u)\bigr),
\]
which is summable by the argument of
\cref{plt:thm:du-composition-closure} verbatim, since the extra factor
\(\Lambda^{\gamma\alpha}\) affects no \(x\)-exponent.
% ed.: the argument here claimed that injectivity of alpha |-> alpha*beta
% makes the summands non-interfering.  It does not: the summand of alpha'
% has x-support in alpha'*beta + Gamma_u, so for alpha' < alpha it also
% reaches the x-exponent alpha*beta.  What is true, and all that is needed,
% is that only summands with alpha' <= alpha reach it.
The \(\alpha\)-th summand has \(x\)-support in
\(\alpha\beta+\Gamma_{u}\subseteq[\alpha\beta,+\infty)\), so the \(x\)-exponent
\(\alpha\beta\) receives a contribution only from summands with
\(\alpha'\le\alpha\), and from the summand \(\alpha\) itself only through the
\(\Gamma_{u}\)-degree-zero part, in which the factors \((1+u)^{\alpha}\) and
\(\log(1+u)\) contribute only their constant terms \(1\) and \(0\).  That
part is
\begin{equation}\label{plt:eq:du-inner-log-block}
  c^{\alpha}\Lambda^{\gamma\alpha}\,
  p_{\alpha}\bigl(\beta\Lambda-\gamma\Lambda_{2}-\log c\bigr).
\end{equation}
Writing \(d=\deg p_{\alpha}\) and \(a_{d}\ne0\) for its leading coefficient,
the coefficient of \(\Lambda_{2}^{\,d}\) in \eqref{plt:eq:du-inner-log-block}
is \(c^{\alpha}\Lambda^{\gamma\alpha}a_{d}(-\gamma)^{d}\), which is nonzero
when \(d\ge1\) because \(\gamma\ne0\); and it is \(0\) for every \(\alpha\)
when every \(p_{\alpha}\) is a constant, in which case no summand contains
\(\Lambda_{2}\) at all, since \((1+u)^{\alpha}\) and \(\log(1+u)\) do not.
That settles the ``only if''.  For the ``if'', let
\(\alpha_{\ast}=\min\SetOf{\alpha\in S:p_{\alpha}\notin\K}\), which exists
because \(S\) is left-finite.  Every summand \(\alpha'<\alpha_{\ast}\) has
\(p_{\alpha'}\in\K\) and therefore contributes no \(\Lambda_{2}\) anywhere,
while the summand \(\alpha_{\ast}\) contributes at the \(x\)-exponent
\(\alpha_{\ast}\beta\) the nonzero \(\Lambda_{2}^{\,d}\)-coefficient just
computed; nothing can cancel it, so \(\Lambda_{2}\) occurs.

For the second assertion take \(\alpha=\alpha_{0}=\min S\).  No summand
\(\alpha'<\alpha_{0}\) exists, so the block of \(\hat f\compose\hat g\) at the
\(x\)-exponent \(\alpha_{0}\beta\) is exactly
\eqref{plt:eq:du-inner-log-block}.  Its \(\Lambda\)-exponents form the set
\(\gamma\alpha_{0}+\SetOf{0,1,\dots,\deg p_{\alpha_{0}}}\), and the exponent
\(\gamma\alpha_{0}\) itself occurs with coefficient
\(c^{\alpha_{0}}p_{\alpha_{0}}(-\gamma\Lambda_{2}-\log c)\), a nonzero
polynomial in \(\Lambda_{2}\) because \(\gamma\ne0\) and
\(p_{\alpha_{0}}\ne0\); so if \(\gamma\alpha_{0}\notin\NaturalNumbers\) a
\(\Lambda\)-exponent outside \(\NaturalNumbers\) occurs, and if
\(\gamma\alpha_{0}\notin\IntegerNumbers\) one outside \(\IntegerNumbers\)
occurs.  Finally, depth one at a finite point means
support in \(\SetOf{(a,b,0)}\) in the coordinates of
\eqref{plt:eq:du-zero-order}, which a nonzero \(\Lambda_{2}\)-exponent
contradicts; and a non-natural \(\Lambda\)-exponent excludes every class with
\(b\in\NaturalNumbers\).  Read through \(\sigma^{-1}\), this is
\cref{plt:thm:grp-scale-classification}(4).
\end{proof}

\begin{table}[htbp]
\centering
\caption{The enlargement ladder at a finite point.  Every row names the
smallest class containing the object produced, and the statement that forces
it.}
\label{plt:tab:du-enlargements}
\begin{tabular}{@{}lll@{}}
\toprule
what forces it & smallest zero-side class & reference\\
\midrule
nothing (integer exponents)
  & \(\K[\Lambda]((x))=\sigma(\PLlau)\)
  & \cref{plt:thm:du-two-substitutions}\\
division by a nonconstant block
  & \(\K(\Lambda)((x))\), then \(\K((\ell))((x))\)
  & \cref{plt:thm:du-log-degree-failure}(i)\\
ratio \(r=p/q\in\RationalNumbers\)
  & \(\K[\Lambda]((x^{1/q}))\)
  & \cref{plt:prop:du-puiseux-zero}(ii)\\
ratio \(r\notin\RationalNumbers\)
  & \(\mathfrak D(\K)\), grid in \(\NaturalNumbers+r\NaturalNumbers\)
  & \cref{plt:prop:du-puiseux-zero}(iii)\\
non-integer power of a block
  & real \(\Lambda\)-exponents, rank-two Hahn
  & \cref{plt:thm:du-log-degree-failure}(ii)\\
inner monomial with \(\Lambda^{\gamma}\), \(\gamma\ne0\)
  & depth two: \(\Lambda_{2}\) appears
  & \cref{plt:prop:du-iterated-log-forced}\\
a logarithmic leading scale
  & the \(\tau\)-image \(\K[\Lambda_{2}]((\ell))\)
  & \cref{plt:thm:du-two-substitutions}(ii)\\
an exponentially flat germ
  & none: the expansion is \(0\)
  & \cref{plt:prop:du-flat-kernel}(ii)\\
\bottomrule
\end{tabular}
\end{table}

\begin{remark}[What is not claimed]
\label{plt:rmk:du-not-claimed}
Three questions are left open here, and none of them is answered by anything
above.
\begin{enumerate}[label=\textup{(\alph*)}]
\item Whether the logarithmic degrees \(\deg p_{j}\) of the Dulac series of a
  hyperbolic-polycycle return map grow at most linearly in \(j\).
  \Cref{plt:prop:du-degree-zero} gives a linear bound for reversion of a
  perturbation of uniformly bounded logarithmic degree, and
  \cref{plt:thm:lag-degree-bound} gives the corresponding statement at
  infinity; neither applies to a composite of transition maps, whose
  logarithmic degrees are not supplied by \cref{plt:thm:du-dulac-transition} in
  any form we have verified.
\item Whether the analytic half of \cref{plt:thm:du-dulac-transition} --- that
  a composition of Dulac germs is a Dulac germ with the formally composed
  series --- can be proved by transporting \cref{plt:thm:cmp-analytic-transfer}
  through \cref{plt:prop:du-asymptotic-transport}.  The obstruction is that the
  latter is stated for integer exponents, and the uniform control needed for
  the outer expansion at real exponents has not been established here.
\item Whether every Dulac germ arising from an analytic saddle is definable in
  an o-minimal expansion of the real field; see
  \cref{plt:rmk:du-dulac-problem}.
\end{enumerate}
\end{remark}

\begin{summarybox}
At a finite point there are two substitutions, not one.  The reciprocal
substitution \(\sigma\colon X=1/x\) is a relabelling of \(\PLlau\) onto
\(\K[\Lambda]((x))\) with \(\Lambda=\log(1/x)\); it preserves valuation,
truncation, leading data and dominance, and carries \(\DX\) to
\(-x^{2}\Differential/\Differential x\), not to
\(\Differential/\Differential x\).  The exponential substitution
\(\tau\colon X=\log(1/x)\) lands on the disjoint logarithmic scale
\(\K[\Lambda_{2}]((\ell))\) and is the only one that reaches germs such as
\(\LambertW_{-1}(-x)\).  Every trap in transporting a theorem is an instance of
one sentence: a statement mentioning neither the derivation nor the analytic
meaning of a monomial transports automatically, and one mentioning either must
be re-derived --- as the near-zero Lagrange--B\"urmann formula must, since the
reversion-frame axiom \(\nu(\partial f)\ge\nu(f)+1\) becomes
\(\nu_{x}(\partial f)\ge\nu_{x}(f)-1\).  A Dulac series is an element of the
Hahn ring in \(x\) and \(\Lambda\) whose \(x\)-support is \emph{left-finite},
strictly more than the well-basedness a Hahn ring requires and exactly what
finite truncation needs; over \(\IntegerNumbers\) the two conditions coincide,
which is why the distinction is invisible in \(\PLlau\).  The Dulac class is
closed under composition with a positive-leading-exponent argument and under
analytic changes of chart, which is what puts the first-return map of a
hyperbolic polycycle in it; but the passage from a germ to its Dulac series has
a nonzero kernel, and closing that gap is Dulac's problem, not algebra.
\end{summarybox}

\clearpage
\part{Algorithms, certificates, and diagnostics}

% ---- fragment 14-algorithms ----
\chapter{Data structures and finite-order algorithms}
\label{plt:sec:algorithms}

Parts~I--IV established what the polynomial--logarithmic calculus \emph{is}:
which sums, products, quotients, compositions and inverses exist, and which
finitely many input blocks determine a given output block.  This chapter turns
those finiteness theorems into routines.  Nothing here is a new closure result;
everything here is a contract.  The mathematical content is of three kinds:
the exact statement of what a truncated object denotes, the exact statement of
how much precision each primitive delivers, and the exact statement of what
each primitive costs.

Two conventions are in force throughout, and both are consequences of the
normative catalogue rather than choices of this chapter.

\begin{itemize}
\item Truncation is \emph{exclusive}: \(\Trunc{F}{N}=\sum_{n<N}p_n(L)t^n\).
      A source that writes \([F]_{\le N}\) is using the inclusive marker, and
      \([F]_{\le N}=\Trunc{F}{N+1}\).  Every index bound below has been
      converted; the one-block discrepancy this conversion repairs is the
      subject of \cref{plt:sec:failure-modes}.
\item The three \(O\)-roles are never conflated.  In this chapter
      \(\FormalBigOAt{t^N}{t}\) is a statement about coefficients,
      \(\BigOAt{\cdot}{X\to+\infty}\) is a statement about functions, and
      \(\BigO(\cdot)\) counts machine operations.  A complexity bound is not an
      error bound and an error bound is not a complexity bound.
\end{itemize}

\begin{warningbox}
The tail \(t^N\PLpow\) is an additive subgroup of \(\PLlau\) and a filtration
step, \emph{not} an ideal of \(\PLlau\): in \(\PLlau=\K[L]((t))\) the element
\(t\) is a unit, so the ideal generated by \(t^N\) is the whole ring.  Every
statement below about ``working modulo \(t^N\)'' is therefore a statement about
cosets of an additive subgroup, and quotient-ring language is illegitimate on
the Laurent ring.  Reduction modulo \((t^N)\) is a genuine ring quotient only
after one restricts to \(\PLpow\).
\end{warningbox}

\section{What a computation stores}
\label{plt:sec:alg-storage}

\begin{definition}[Sparse block representation]\label{plt:def:alg-representation}
Let \(R\) be a commutative \(\K\)-algebra with a decidable zero test and a
distinguished canonical form; the cases of interest are
\(R=\K[L]\), \(R=\K(L)\), and \(R=\K((L^{-1}))\) truncated at a declared inner
cutoff, together with the tower of \cref{plt:def:div-iterated}.  A
\emph{truncated representative} over \(R\) is a pair
\[
 \mathcal T=\bigl(v,N\bigr),
 \qquad
 N\in\IntegerNumbers\cup\SetOf{+\infty},
\]
in which \(v\) is a finitely supported map
\[
 v\colon \IntegerNumbers\longrightarrow R,
 \qquad
 v(n)=0\ \text{for}\ n\ge N,
\]
stored as the strictly increasing list of the pairs \((n,v(n))\) with
\(v(n)\ne0\).  The outer index \(n\) is the primary key and the internal
representation of \(v(n)\) --- for \(R=\K[L]\), a polynomial stored by
logarithmic degree --- is subordinate to it.  The representative additionally
carries the name of \(R\) and, when the ground field contains chosen values of
logarithms or fractional powers, the branch data selecting them.
\end{definition}

The two-level key is forced by the dominance order of
\cref{plt:lem:mot-dominance}: \(t^nL^j\succdom t^mL^k\) compares \(n\) first
and \(j\) only within a fixed \(n\).  Any storage scheme that flattens
\((n,j)\) into a single total degree destroys the hierarchy, because
\(L^{D}\succdom1\) while \(L^D\) has the larger total degree; truncating by
total degree therefore discards \emph{dominant} monomials.  Blocks are finite
(\cref{plt:prop:mot-blocks}) but their degrees are unbounded across blocks
(\cref{plt:def:mul-profile}), so the second level must be a variable-length
object, never a fixed second array dimension.

\begin{definition}[Denotation]\label{plt:def:alg-denotation}
For a truncated representative \(\mathcal T=(v,N)\) over \(R=\K[L]\) put
\[
 \operatorname{Den}(\mathcal T)
 \DefinitionEquals
 \SetOf{F\in\PLlau:\Trunc{F}{N}=v}.
\]
\end{definition}

\begin{lemma}[The denotation is a coset of a filtration step]
\label{plt:lem:alg-denotation-coset}
Let \(\mathcal T=(v,N)\) with \(N\) finite.  Then
\[
 \operatorname{Den}(\mathcal T)=v+t^N\PLpow,
\]
and every \(F\in\operatorname{Den}(\mathcal T)\) satisfies \(\ordt(F)\ge m(\mathcal T)\),
where
\[
 m(\mathcal T)
 \DefinitionEquals
 \begin{cases}
  \ordt(v),&v\ne0,\\
  N,&v=0.
 \end{cases}
\]
The bound \(m(\mathcal T)\) is attained: it equals
\(\min\SetOf{\ordt(F):F\in\operatorname{Den}(\mathcal T)}\).
\end{lemma}

\begin{proof}
If \(\Trunc{F}{N}=v\) then every block of \(F-v\) with index \(<N\) vanishes,
so \(F-v\in t^N\PLpow\); conversely if \(F=v+w\) with \(w\in t^N\PLpow\) then
\(\Trunc{F}{N}=\Trunc{v}{N}=v\) because \(v\) is supported below \(N\).  This
proves the displayed equality.  For the valuation, \(\ordt(F)\ge
\min\SetOf{\ordt(v),N}\), and \(\ordt(v)\le N-1<N\) when \(v\ne0\), which gives
\(\ordt(F)\ge\ordt(v)=m(\mathcal T)\) with equality for \(F=v\).  When \(v=0\)
the set is \(t^N\PLpow\), whose minimum valuation is \(N\), attained at
\(t^N\).
\end{proof}

% ed.: S4 stores the least allowed exponent n_min as an independent field and
% then warns that a stale zero at that exponent corrupts order and unit tests.
% Lemma plt:lem:alg-denotation-coset shows the field is redundant: the order
% bound is a function of the canonical form, so the failure is eliminated by
% deriving it rather than by documenting it.
\begin{remark}[The order bound is derived, not stored]
\label{plt:rmk:alg-derived-order}
\Cref{plt:lem:alg-denotation-coset} shows that the valuation bound
\(m(\mathcal T)\) is determined by \((v,N)\).  An implementation must therefore
recompute it from the canonical form and must not carry it as an independent
mutable field: a retained zero block at the nominal least exponent makes a
stored bound stale, and a stale bound falsifies exactly the precision law of
\cref{plt:thm:alg-precision-calculus}, where \(m\) occurs on the right-hand
side.
\end{remark}

\begin{definition}[Relative precision]\label{plt:def:alg-relative-precision}
The \emph{absolute precision} of \(\mathcal T=(v,N)\) is \(N\); its
\emph{relative precision} is
\[
 \rho(\mathcal T)\DefinitionEquals N-m(\mathcal T)\in\NaturalNumbers\cup\SetOf{+\infty}.
\]
Thus \(\rho\) counts the number of outer blocks known from the leading one
onwards, and \(\rho(\mathcal T)=0\) exactly when \(v=0\).
\end{definition}

\begin{definition}[Soundness and sharpness]\label{plt:def:alg-sound}
Let \(\varphi\colon\mathcal D\to\PLlau\) be a partial operation on
\(k\)-tuples.  A routine \(\mathsf R\) that maps truncated representatives
\(\mathcal T_1,\dots,\mathcal T_k\) to a truncated representative
\(\mathsf R(\mathcal T_1,\dots,\mathcal T_k)=(w,N_{\mathrm{out}})\) is
\emph{sound} for \(\varphi\) if
\[
 \varphi\bigl(\operatorname{Den}(\mathcal T_1)\times\cdots\times\operatorname{Den}(\mathcal T_k)\bigr)
 \subseteq
 \operatorname{Den}\bigl(\mathsf R(\mathcal T_1,\dots,\mathcal T_k)\bigr)
\]
whenever every tuple in the product lies in \(\mathcal D\).  It is
\emph{sharp} if in addition no larger value of \(N_{\mathrm{out}}\) has this
property, that is, if there exist two admissible tuples whose \(\varphi\)-images
differ in the block of index \(N_{\mathrm{out}}\).
\end{definition}

Soundness is the invariant that every routine must preserve.  It is worth
isolating the reason that it is enough to check it one primitive at a time.

\begin{proposition}[Soundness composes]\label{plt:prop:alg-soundness-composes}
Let \(\mathsf R_1\) be sound for \(\varphi_1\) and \(\mathsf R_2\) sound for
\(\varphi_2\), and suppose the declared output type of \(\mathsf R_1\) is an
admissible input type of \(\mathsf R_2\).  Then
\(\mathsf R_2\circ\mathsf R_1\) is sound for \(\varphi_2\circ\varphi_1\) on the
locus where the latter is defined.
\end{proposition}

\begin{proof}
Let \(F\in\operatorname{Den}(\mathcal T)\) with \(\varphi_2(\varphi_1(F))\) defined.
Soundness of \(\mathsf R_1\) gives
\(\varphi_1(F)\in\operatorname{Den}(\mathsf R_1(\mathcal T))\).  Applying soundness of
\(\mathsf R_2\) to the representative \(\mathsf R_1(\mathcal T)\) gives
\(\varphi_2(\varphi_1(F))\in\operatorname{Den}(\mathsf R_2(\mathsf R_1(\mathcal T)))\), which
is the assertion.  The argument uses nothing about \(\mathsf R_1\) except that
its output denotation contains \(\varphi_1(F)\); in particular it is
indifferent to how many blocks \(\mathsf R_1\) discarded.
\end{proof}

\begin{corollary}[Truncating early is free]\label{plt:cor:alg-truncate-early}
If every primitive in a straight-line computation is sound, then truncating
each intermediate result to its own declared precision changes no block of the
final result that the final declared precision advertises.
\end{corollary}

\begin{proof}
Immediate from \cref{plt:prop:alg-soundness-composes}: the truncation of an
intermediate is the passage from an element of a denotation to the
representative of that denotation, and soundness of the next routine is stated
for the whole denotation.
\end{proof}

\Cref{plt:cor:alg-truncate-early} is the theorem behind the universal advice of
all six sources to truncate after every operation.  The advice is not a
heuristic trade of accuracy against speed: at the declared precision the answer
is bit-for-bit the same.

\begin{definition}[The routine invariant]\label{plt:def:alg-invariant}
A routine is \emph{conforming} when it guarantees all of:
\begin{enumerate}
\item[(P1)] \emph{Soundness} in the sense of \cref{plt:def:alg-sound}, with the
      declared \(N_{\mathrm{out}}\).
\item[(P2)] \emph{Canonicity}: the returned map has strictly increasing keys, no
      key \(\ge N_{\mathrm{out}}\), no zero value, and every value in the
      canonical form of \(R\).
\item[(P3)] \emph{Order honesty}: the order bound is recomputed from the
      returned canonical form (\cref{plt:rmk:alg-derived-order}).
\item[(P4)] \emph{Ring honesty}: if the computation left the declared \(R\) --- for
      instance by dividing by a nonunit block, \cref{plt:thm:div-unit-criterion}
      --- the enlarged \(R\) is part of the returned type and not a silent
      simplification.
\end{enumerate}
\end{definition}

\section{The precision calculus}
\label{plt:sec:alg-precision}

The following theorem is the arithmetic that (P1) requires.  All six sources
state fragments of it; none states it sharply, and the missing sharpness is the
source of the guard-block failures catalogued in
\cref{plt:sec:failure-modes}.

\begin{theorem}[Precision calculus, with sharpness]
\label{plt:thm:alg-precision-calculus}
Let \(\mathcal T_i=(v_i,N_i)\) be truncated representatives with finite \(N_i\),
order bounds \(m_i=m(\mathcal T_i)\) and relative precisions
\(\rho_i=N_i-m_i\).  Then the following output precisions are sound, and each
is sharp in the sense of \cref{plt:def:alg-sound} --- for \emph{every} instance
in cases \ref{plt:thm:alg-precision-calculus:sum}--\ref{plt:thm:alg-precision-calculus:deriv},
and for every instance satisfying the nondegeneracy hypothesis stated there in
case \ref{plt:thm:alg-precision-calculus:comp}.
\begin{enumerate}
\item\label{plt:thm:alg-precision-calculus:sum}
      \emph{Sum.}  \(N_{\mathrm{out}}=\min\SetOf{N_1,N_2}\).
\item\label{plt:thm:alg-precision-calculus:prod}
      \emph{Product.}  \(N_{\mathrm{out}}=\min\SetOf{N_1+m_2,\ N_2+m_1}\), with
      output order bound \(m_1+m_2\).
\item\label{plt:thm:alg-precision-calculus:recip}
      \emph{Reciprocal.}  If every \(F\in\operatorname{Den}(\mathcal T_1)\) is a unit of
      \(\PLlau\), then \(N_{\mathrm{out}}=N_1-2m_1\), with output order bound
      \(-m_1\).
\item\label{plt:thm:alg-precision-calculus:quot}
      \emph{Quotient.}  If every \(G\in\operatorname{Den}(\mathcal T_2)\) is a unit, then
      \(N_{\mathrm{out}}=\min\SetOf{N_1-m_2,\ N_2+m_1-2m_2}\), with output order
      bound \(m_1-m_2\).
% ed.: the previous version of (v) asserted that sharpness fails at N_1=0
% ed.: because DX annihilates the unknown constant block.  That is an artifact
% ed.: of scalar power series and is false over R=K[L]: the unknown block of
% ed.: index N_1 is an arbitrary polynomial, and a=t^{N_1}L already has
% ed.: DX a = t^{N_1+1}(1-N_1 L) \ne 0 for every N_1, including N_1=0.  What is
% ed.: true, and is what the sources were reaching for, is that the output
% ed.: ORDER bound can fail to be attained -- exactly when the leading block is
% ed.: a nonzero constant sitting at index 0 -- and that this is a loss of
% ed.: relative precision, not of absolute precision.
\item\label{plt:thm:alg-precision-calculus:deriv}
      \emph{Derivation.}  \(N_{\mathrm{out}}=N_1+1\) for \(\DX\), and this is
      sharp without exception.  The output order bound \(m_1+1\) is always
      sound; for \(v_1\ne0\) it is attained if and only if
      \((\partial_L-m_1)v_1(m_1)\ne0\), that is, unless
% ed.: the criterion was stated without the hypothesis v_1 != 0, at which it
% ed.: fails: there v_1(m_1)=0, yet the returned representative is (0,N_1+1),
% ed.: whose order bound is N_1+1=m_1+1, so the bound is attained.
      \(m_1=0\) and \(v_1(0)\in\K^{\times}\), while for \(v_1=0\) it is
      attained trivially, both sides being \(N_1+1\); in the exceptional case the true
      order is strictly larger and must be recomputed from the returned
      canonical form (\cref{plt:rmk:alg-derived-order}).
\item\label{plt:thm:alg-precision-calculus:comp}
      \emph{Near-identity composition.}  Let \(\mathcal T_1\) represent the
      outer series \(F\in\PLlau\) and let \(\mathcal T_2=(u,N_2)\) represent
      \(H\in\PLpow\), so that the inner map is \(X+H\in\Gnear\).  Then
      \[
        N_{\mathrm{out}}=\min\SetOf{N_1,\ N_2+m_1+1},
      \]
      with output order bound \(m_1\).  Sharpness holds for every instance in
      which \(N_1\le N_2+m_1+1\), and, when \(N_2+m_1+1<N_1\), for every
      instance with \(\ordt(\DX v_1)=m_1+1\).  The excluded case is real: for
      \(v_1=c\in\K^{\times}\) every admissible pair gives
      \(F\compose(X+H)=c+a\compose(X+H)\) with \(\ordt(a)\ge N_1\), so the
      composite is in fact determined to the larger cutoff \(N_1\), and the
      declared \(N_{\mathrm{out}}=N_2+1\) is sound but pessimistic.
\end{enumerate}
\end{theorem}

\begin{proof}
Write \(F=v_1+a\) and \(G=v_2+b\) with \(\ordt(a)\ge N_1\), \(\ordt(b)\ge N_2\);
by \cref{plt:lem:alg-denotation-coset} this parametrises the denotations, and
\(\ordt(v_i)\ge m_i\), \(N_i\ge m_i\).

% ed.: the sharpness witnesses below were originally exhibited only at the zero
% ed.: representative (sum) or for one of the two competing bounds (product,
% ed.: quotient), which proves sharpness of the LAW but not of each instance.
% ed.: Since Definition plt:def:alg-sound quantifies sharpness per instance,
% ed.: the witnesses are now chosen so that the block that differs is exactly
% ed.: the declared cutoff, for every instance.
\emph{\ref{plt:thm:alg-precision-calculus:sum}.}
\((F+G)-(v_1+v_2)=a+b\) has valuation at least \(\min\SetOf{N_1,N_2}\).
For sharpness assume \(N_1\le N_2\) (otherwise exchange the two arguments) and
compare the admissible pairs \((v_1,v_2)\) and \((v_1+t^{N_1},v_2)\): their
sums differ by \(t^{N_1}\), a nonzero block of index exactly
\(N_1=N_{\mathrm{out}}\).  This witness exists for every instance, not only for
\(v_1=v_2=0\).

\emph{\ref{plt:thm:alg-precision-calculus:prod}.}
\(FG-v_1v_2=v_1b+av_2+ab\).  The three summands have valuations at least
\(m_1+N_2\), \(N_1+m_2\) and \(N_1+N_2\); since \(N_i\ge m_i\) the third is not
smaller than the first two, so \(\ordt(FG-v_1v_2)\ge N_{\mathrm{out}}\).
Additivity of \(\ordt\) (\cref{plt:lem:div-valuation}) gives the order bound.
% ed.: the witness that perturbs the first argument produces a difference in
% ed.: the block N_1+m_2, which is the declared cutoff only when that bound is
% ed.: the smaller of the two.  The previous version used it unconditionally
% ed.: and so proved sharpness of the law, not of each instance as
% ed.: Definition plt:def:alg-sound requires.  The witness must be selected by
% ed.: which bound realises the minimum.
For sharpness one must exhibit, for the given instance, two admissible pairs
whose products differ in the block of index \emph{exactly}
\(N_{\mathrm{out}}\).
\emph{Case \(N_1+m_2\le N_2+m_1\) with \(v_2\ne0\).}  The admissible pairs
\((v_1,v_2)\) and \((v_1+t^{N_1},v_2)\) have products differing by
\(t^{N_1}v_2\), of valuation exactly \(N_1+\ordt(v_2)=N_1+m_2=N_{\mathrm{out}}\).
\emph{Case \(N_2+m_1\le N_1+m_2\) with \(v_1\ne0\).}  Symmetrically the pairs
\((v_1,v_2)\) and \((v_1,v_2+t^{N_2})\) have products differing by
\(t^{N_2}v_1\), of valuation exactly \(N_2+m_1=N_{\mathrm{out}}\).
\emph{Remaining cases.}  If \(v_2=0\) and \(v_1\ne0\) then \(m_2=N_2\) while
\(m_1\le N_1-1\), so \(N_2+m_1<N_1+N_2=N_1+m_2\) and the second case applies;
symmetrically if \(v_1=0\) and \(v_2\ne0\).  If \(v_1=v_2=0\) then
\(N_{\mathrm{out}}=N_1+N_2\) and the pairs \((0,0)\), \((t^{N_1},t^{N_2})\)
have products \(0\) and \(t^{N_1+N_2}\).  Every instance is covered.

\emph{\ref{plt:thm:alg-precision-calculus:recip}.}
Here \(v_1\ne0\) and \(m_1=\ordt(v_1)=\ordt(F)\).  From
\(\MultiplicativeInverseOf{F}-\MultiplicativeInverseOf{v_1}
=(v_1-F)/(Fv_1)=-a/(Fv_1)\) and additivity of \(\ordt\),
\[
 \ordt\bigl(\MultiplicativeInverseOf{F}-\MultiplicativeInverseOf{v_1}\bigr)
 =\ordt(a)-2m_1\ge N_1-2m_1 .
\]
Equality holds for \(a=t^{N_1}\), which proves sharpness.

\emph{\ref{plt:thm:alg-precision-calculus:quot}.}
With \(A=v_1+a\) and \(B=v_2+b\),
\[
 \frac AB-\frac{v_1}{v_2}
 =\frac{av_2-bv_1}{Bv_2},
\]
whose valuation is at least
\(\min\SetOf{N_1+m_2,N_2+m_1}-2m_2\), which is the assertion.
% ed.: "follows from the two witnesses used in (ii)" was too quick: as in (ii)
% ed.: the witness must be the one realising the minimum.  Here, unlike (ii),
% ed.: both witnesses exist for every instance, because the numerator is
% ed.: allowed to be zero and may then be replaced by t^{N_1}.
For sharpness note that \(\ordt(v_2)=m_2\) and \(N_2>m_2\), so
\(\ordt\bigl(v_2(v_2+t^{N_2})\bigr)=2m_2\) exactly.  Perturbing the numerator
from \(v_1\) to \(v_1+t^{N_1}\) at fixed denominator \(v_2\) changes the
quotient by \(t^{N_1}/v_2\), of valuation exactly \(N_1-m_2\).  Perturbing the
denominator from \(v_2\) to \(v_2+t^{N_2}\) at a fixed admissible numerator
\(A\) changes it by \(-At^{N_2}/\bigl(v_2(v_2+t^{N_2})\bigr)\), of valuation
exactly \(\ordt(A)+N_2-2m_2\); take \(A=v_1\) if \(v_1\ne0\) and \(A=t^{N_1}\)
if \(v_1=0\), so that \(\ordt(A)=m_1\) in both cases and the valuation is
\(m_1+N_2-2m_2\).  Whichever of the two competing bounds is the smaller, the
corresponding witness differs in the block of index exactly
\(N_{\mathrm{out}}\).

\emph{\ref{plt:thm:alg-precision-calculus:deriv}.}
\(\DX(F-v_1)=\DX a\) and \(\ordt(\DX a)\ge\ordt(a)+1\ge N_1+1\)
(\cref{plt:lem:rev-deriv-gain}).
% ed.: the proof still carried the discarded claim that sharpness fails at
% ed.: N_1=0, contradicting the corrected statement above.  The witness
% ed.: a=t^{N_1} is the wrong one: it is annihilated exactly at N_1=0.  Over
% ed.: R=K[L] the undetermined block is an arbitrary polynomial, and
% ed.: a=t^{N_1}L works for every N_1.
For sharpness take \(a=t^{N_1}L\), which is admissible because
\(\ordt(a)=N_1\).  Then
\(\DX a=t^{N_1+1}(\partial_L-N_1)L=t^{N_1+1}(1-N_1L)\ne0\) is a nonzero block
of index exactly \(N_1+1=N_{\mathrm{out}}\), for every
\(N_1\in\IntegerNumbers\), including \(N_1=0\).
For the order clause, the block of index \(m_1+1\) of \(\DX v_1\) is
\((\partial_L-m_1)v_1(m_1)\) by the block formula; the operator
\(\partial_L-m_1\) is injective on \(\K[L]\) for \(m_1\ne0\) and has kernel
exactly \(\K\) for \(m_1=0\), while \(v_1(m_1)\ne0\) whenever \(v_1\ne0\).
Hence that block vanishes precisely when \(m_1=0\) and
\(v_1(0)\in\K^{\times}\).  This --- and not any failure of the cutoff law ---
is the algorithmic shadow of \(\ker\DX=\K\) (\cref{plt:cor:dif-kernel}).

\emph{\ref{plt:thm:alg-precision-calculus:comp}.}
Write \(G=X+H\).  Perturbing the outer argument by \(\Phi\) with
\(\ordt(\Phi)\ge N_1\) changes the value by \(\Phi\compose G\), and
\begin{equation}\label{plt:eq:alg-right-composition-isometry}
 \ordt(\Phi\compose G)=\ordt(\Phi)
 \qquad(\Phi\in\PLlau,\ G\in\Gnear),
\end{equation}
because the exact Taylor formula \cref{plt:thm:grp-taylor} gives
\(\Phi\compose G=\sum_{k\ge0}H^k\DX^k\Phi/k!\), whose \(k=0\) term is \(\Phi\)
and whose \(k\ge1\) terms have valuation at least
\(\ordt(\Phi)+k>\ordt(\Phi)\); this is
\cref{plt:lem:rev-order-preserved}.  Hence the outer argument contributes
precision \(N_1\), sharply, by taking \(\Phi=t^{N_1}\).
Perturbing the inner correction by \(\eta\) with \(\ordt(\eta)\ge N_2\) changes
the value by
\[
 F\compose(X+H+\eta)-F\compose(X+H)
 =\sum_{k\ge1}\frac{(H+\eta)^k-H^k}{k!}\,\DX^kF .
\]
Since \(\ordt(H)\ge0\) we have \(\ordt((H+\eta)^k-H^k)\ge\ordt(\eta)\ge N_2\),
while \(\ordt(\DX^kF)\ge m_1+k\); so the \(k\)-th summand has valuation at
least \(N_2+m_1+k\), and the minimum over \(k\ge1\) is \(N_2+m_1+1\).
% ed.: the witness was a single pair (F=t^{m_1} with m_1 nonzero), which proves
% ed.: sharpness of the law but not the per-instance claim actually stated.
% ed.: Under the printed hypothesis ordt(DX v_1)=m_1+1 the generic witness
% ed.: works, and the hypothesis is exactly what makes the k=1 term dominate.
For sharpness in the regime \(N_2+m_1+1<N_1\), assume
\(\ordt(\DX v_1)=m_1+1\) and take \(F=v_1\), \(\eta=t^{N_2}\).  The \(k=1\)
summand of the displayed difference is exactly \(\eta\,\DX F=t^{N_2}\DX v_1\),
of valuation exactly \(N_2+m_1+1\); for \(k\ge2\) one has
\(\ordt\bigl((H+\eta)^k-H^k\bigr)\ge N_2\) and
\(\ordt(\DX^kF)\ge\ordt(\DX F)+(k-1)=m_1+k\ge m_1+2\), so all later summands
have valuation at least \(N_2+m_1+2\).  The admissible inner data \(H\) and
\(H+\eta\) therefore give composites differing in the block of index exactly
\(N_{\mathrm{out}}\).  Finally the output
order bound is \(m_1\) by \cref{plt:eq:alg-right-composition-isometry}.
\end{proof}

\begin{corollary}[Relative precision is the conserved quantity]
\label{plt:cor:alg-relative-precision}
In the notation of \cref{plt:thm:alg-precision-calculus}, and writing
\(h=m_2\ge0\) for the inner order in
part~\ref{plt:thm:alg-precision-calculus:comp},
\begin{align*}
 \rho(FG)&=\min\SetOf{\rho_1,\rho_2}, &
 \rho\bigl(\MultiplicativeInverseOf{F}\bigr)&=\rho_1, \\
 \rho(A/B)&=\min\SetOf{\rho_1,\rho_2}, &
 \rho\bigl(F\compose(X+H)\bigr)&=\min\SetOf{\rho_1,\ \rho_2+h+1}.
\end{align*}
% ed.: the inequality was stated as rho(DX F) >= rho_1.  It runs the other way:
% ed.: the absolute cutoff gains one block, but in the exceptional case of
% ed.: part (v) the output ORDER gains more than one, so the difference --
% ed.: which is rho -- can only fall.  The extreme case v_1 = c in K^x returns
% ed.: the zero representative, whose relative precision is 0.
% ed.: two further repairs to this clause.  (a) The equality criterion for the
% ed.: derivation fails at the zero representative: there \(v_1(m_1)=0\) while
% ed.: \(m_{\mathrm{out}}=N_1+1=m_1+1\), so equality holds although the
% ed.: criterion does not; the clause is restricted to \(v_1\ne0\).  (b) "no
% ed.: bound in either direction" was an overclaim for the sum: an upper bound
% ed.: does hold, \(\rho\le\max\{\rho_1,\rho_2\}\) (proved below), so what the
% ed.: witness shows is that no LOWER bound holds -- and \(\rho\) can drop to
% ed.: \(0\), not merely to \(1\).
For the derivation, \(\rho(\DX F)\le\rho_1\); for \(v_1\ne0\), equality holds
if and only if \((\partial_L-m_1)v_1(m_1)\ne0\), and in the exceptional case
the drop can be total, as it is for \(v_1=c\in\K^{\times}\); for \(v_1=0\)
both sides are \(0\).  For the sum,
\(\rho\le\max\SetOf{\rho_1,\rho_2}\) always, but no lower bound in terms of
\(\rho_1,\rho_2\) holds: \(\rho\) can drop to \(0\).
\end{corollary}

\begin{proof}
Each identity is \(N_{\mathrm{out}}-m_{\mathrm{out}}\) computed from
\cref{plt:thm:alg-precision-calculus}; for the product, the reciprocal, the
quotient and the composition respectively,
\begin{align*}
 \min\SetOf{N_1+m_2,N_2+m_1}-(m_1+m_2)
   &=\min\SetOf{N_1-m_1,\ N_2-m_2},\\
 (N_1-2m_1)-(-m_1)&=N_1-m_1,\\
 \min\SetOf{N_1-m_2,N_2+m_1-2m_2}-(m_1-m_2)
   &=\min\SetOf{N_1-m_1,\ N_2-m_2},\\
 \min\SetOf{N_1,N_2+m_1+1}-m_1
   &=\min\SetOf{\rho_1,\ (N_2-h)+h+1}.
\end{align*}
For the derivation, the output cutoff is
\(N_1+1\) and the output order is
\(m_{\mathrm{out}}=\ordt(\DX v_1)\ge m_1+1\), so
\(\rho(\DX F)=N_1+1-m_{\mathrm{out}}\le N_1-m_1=\rho_1\), with equality
exactly when \(m_{\mathrm{out}}=m_1+1\), which for \(v_1\ne0\) is by
\cref{plt:thm:alg-precision-calculus}\ref{plt:thm:alg-precision-calculus:deriv}
the condition \((\partial_L-m_1)v_1(m_1)\ne0\); for \(v_1=c\in\K^{\times}\)
the returned representative is \((0,N_1+1)\), of relative precision \(0\); and
for \(v_1=0\) one has \(m_1=N_1\) and \(m_{\mathrm{out}}=N_1+1=m_1+1\), so
\(\rho=0=\rho_1\).
% ed.: the upper bound for the sum was asserted to be absent; it is not.
For the sum, the returned representative is
\(\bigl(v_1+v_2,\min\SetOf{N_1,N_2}\bigr)\).  If \(v_1+v_2=0\) then
\(\rho=0\).  Otherwise \(m_{\mathrm{out}}=\ordt(v_1+v_2)\ge
\min\SetOf{m_1,m_2}\), where the right-hand side is read as \(m_2\) when
\(v_1=0\) and as \(m_1\) when \(v_2=0\); hence
\(\rho\le\min\SetOf{N_1,N_2}-\min\SetOf{m_1,m_2}\), and assuming without loss
of generality \(N_1\le N_2\) this last quantity equals \(N_1-m_1=\rho_1\) when
\(m_1\le m_2\) and is at most \(N_2-m_2=\rho_2\) when \(m_2<m_1\).  In either
case \(\rho\le\max\SetOf{\rho_1,\rho_2}\).  That no lower bound holds is shown
by
\[
 F=\sum_{n=-1}^{N-1}t^n,
 \qquad
 G=-\sum_{n=-1}^{N-2}t^n,
\]
both with \(m=-1\) and \(N_i=N\), hence \(\rho_i=N+1\); their sum is
\(t^{N-1}\), with order bound \(N-1\) and absolute precision \(N\), so
\(\rho=1\); and by \(G=-F\), for which the sum is the zero representative
\((0,N)\) and \(\rho=0\).
\end{proof}

\begin{remark}[Addition, not multiplication, is the dangerous primitive]
\label{plt:rmk:alg-cancellation}
\Cref{plt:cor:alg-relative-precision} inverts the usual folklore.  Products,
quotients and compositions transport relative precision unchanged or improved;
% ed.: "only addition can destroy it" overlooked the annihilating case of the
% ed.: derivation, which the corrected corollary now records.
the two primitives that can destroy it are addition, which destroys it exactly
when leading blocks cancel --- the formal counterpart of catastrophic
cancellation --- and, in the single exceptional case
\((\partial_L-m_1)v_1(m_1)=0\), the derivation, which destroys it for the same
reason one block further in.  A
routine that reports only absolute cutoffs cannot see this loss; a routine that
reports \(\rho\) sees it immediately, because \(\rho\) drops at the subtraction
that caused it.
\end{remark}

\begin{example}[The one-block guard]\label{plt:ex:alg-guard-block}
Let \(F=X=t^{-1}\), known exactly, and let \(G\in\PLpow\) be known to exclusive
precision \(N\) with \(m_2=0\).  Then
\cref{plt:thm:alg-precision-calculus}\ref{plt:thm:alg-precision-calculus:prod}
gives
\[
 N_{\mathrm{out}}=\min\SetOf{\infty+0,\ N+(-1)}=N-1 :
\]
the product is determined only one block less far than the factor.  The two admissible
completions \(G_1=1\) and \(G_2=1+t^{N}\) satisfy
\[
 \Trunc{G_1}{N}=\Trunc{G_2}{N}=1,
 \qquad
 FG_1=t^{-1},
 \qquad
 FG_2=t^{-1}+t^{N-1},
\]
so the two products differ in the block of index \(N-1\).
% ed.: the source sentence asked for G "through block N-1", which is the datum
% ed.: already assumed and yields the product only through block N-2.  Since
% ed.: multiplication by t^{-1} shifts indices down by one, the block of index
% ed.: n of the product is read off the block of index n+1 of G.
Conversely, to obtain the product through block \(N-1\) inclusively one must
request \(G\) through block \(N\), that is, one guard block beyond the naive
expectation, because the block of index \(n\) of \(FG\) is the block of index
\(n+1\) of \(G\).
\end{example}

% ed.: S2 states only that inputs known modulo t^N with nonnegative valuation
% give a sum and product known modulo t^N.  That is sound but not sharp: by
% part (ii) the product is known modulo t^{N+min(m_1,m_2)}, which is strictly
% better whenever a factor has positive order, and strictly worse -- the case
% S2's formulation hides -- as soon as a factor is Laurent.
\begin{remark}[Why the sharp law matters]
The sound-but-blunt rule ``both inputs modulo \(t^N\) give the output modulo
\(t^N\)'' is correct only when both order bounds are \(\ge0\).  On the Laurent
ring it is false, by \cref{plt:ex:alg-guard-block}, and the failure is silent:
the routine returns a block it has no right to return.  This is failure mode
\ref{plt:sec:fail-undetermined-block} of \cref{plt:sec:failure-modes}.
\end{remark}

\section{Cost models}
\label{plt:sec:alg-cost}

\begin{definition}[Two cost models]\label{plt:def:alg-cost-models}
In the \emph{coefficient-operation model} one counts multiplications,
additions, and inversions of units in the coefficient ring \(R\), each at unit
cost.  In the \emph{scalar model} one counts multiplications and additions in
\(\K\).  Complexity statements below use \(\BigO(\cdot)\) in the sense of
operation counting only; the notation carries no analytic content.
\end{definition}

The two models are not interchangeable, and the discrepancy between them is
governed by the log-degree profile \(\dprofile{F}\) of
\cref{plt:def:mul-profile}, which is data attached to a series and is
\emph{not} bounded by any function of the outer cutoff.

\begin{lemma}[Scalar cost of a block-bilinear routine]
\label{plt:lem:alg-scalar-cost}
Let a routine over \(R=\K[L]\) perform a multiset \(\mathcal M\) of coefficient
multiplications \(p\cdot q\), each carried out by schoolbook polynomial
multiplication.  Then it performs exactly
\[
 \sum_{(p,q)\in\mathcal M}(\deg p+1)(\deg q+1)
\]
scalar multiplications in \(\K\), with the convention that a zero polynomial
contributes the factor \(0\).
\end{lemma}

\begin{proof}
Schoolbook multiplication of polynomials of degrees \(d\) and \(e\) forms
exactly the \((d+1)(e+1)\) products of pairs of coefficients.  Summing over
\(\mathcal M\) gives the count.
\end{proof}

\begin{proposition}[Exact scalar cost of truncated block convolution]
\label{plt:prop:alg-convolution-cost}
Let \(F,G\in\PLpow\) and \(N\ge1\).  The schoolbook algorithm computing
\(\Trunc{FG}{N}\) as \(\sum_{i+j<N}p_iq_j\) performs exactly
\[
 \sum_{\substack{i,j\ge0\\ i+j<N}}
 \bigl(\dprofile{F}(i)+1\bigr)\bigl(\dprofile{G}(j)+1\bigr)
\]
scalar multiplications, where \(\dprofile{F}(n)+1\) is read as \(0\) when the
block \(p_n\) vanishes.  In the coefficient-operation model the count is at
most \(\binom{N+1}{2}\) multiplications in \(\K[L]\).
\end{proposition}

\begin{proof}
Immediate from \cref{plt:lem:alg-scalar-cost} together with
\cref{plt:def:mul-cauchy}: the pairs contributing to blocks below \(N\) are
exactly those with \(i,j\ge0\) and \(i+j<N\), of which there are
\(\binom{N+1}{2}\).
\end{proof}

\begin{corollary}[The scalar cost is not \(\BigO(N^2)\)]
\label{plt:cor:alg-cost-not-quadratic}
\begin{enumerate}
\item If \(\dprofile{F}(n),\dprofile{G}(n)\le D\) for all \(n<N\), the scalar
      cost is at most \((D+1)^2\binom{N+1}{2}\).
\item If \(\dprofile{F}(n)=\dprofile{G}(n)=n\) for all \(n<N\), the scalar cost
      is exactly \(\binom{N+3}{4}\).
\item No bound of the form \(\BigO(N^c)\), with an absolute implied constant,
      holds uniformly over \(\PLpow\): already for \(N=1\) the cost of
      multiplying \(F=G=L^{M}\) is \((M+1)^2\), unbounded in \(M\).
\end{enumerate}
\end{corollary}

\begin{proof}
Part~(i) is immediate.  For part~(ii),
\[
 \sum_{i+j=r}(i+1)(j+1)
 =\sum_{k=1}^{r+1}k\,(r+2-k)
 =(r+2)\frac{(r+1)(r+2)}2-\frac{(r+1)(r+2)(2r+3)}6
 =\binom{r+3}{3},
\]
using \(\sum_{k=1}^{M}k=M(M+1)/2\) and \(\sum_{k=1}^{M}k^2=M(M+1)(2M+1)/6\)
with \(M=r+1\).  Summing over \(0\le r\le N-1\) and applying the hockey-stick
identity \(\sum_{j=3}^{N+2}\binom j3=\binom{N+3}4\) gives the stated closed
form.  Part~(iii) is the displayed example.
\end{proof}

% ed.: "four times the naive quadratic guess in the exponent" is not a
% ed.: statement: the exponent is four rather than two.  Also, the profile of
% ed.: the inverse of X+L is d(0)=1, d(n)=n for n>=1, which differs from
% ed.: case (ii) in the block of index 0; the exact count is corrected below.
\begin{warningbox}
% ed.: "the profile of the compositional inverse of X+L" names the wrong
% ed.: object: that inverse has a block of index -1, namely X itself.  The
% ed.: profile in question is the profile of its CORRECTION.
The profile of case~(ii) is not artificial: up to the single block of index
\(0\) it is the profile of the correction \(\rev{(X+L)}-X\), by
\cref{plt:prop:fail-reversion-degree} below, for which the exact scalar count
is \(\binom{N+3}{4}+N^{2}+N+1\).  So the natural output of the reversion
algorithms already costs \(N^4/24+\BigO(N^3)\) scalar multiplications per
truncated product: the exponent is four, not the naively expected two.  Every
complexity statement in this section is therefore stated in the
coefficient-operation model and must be multiplied by the profile factor
before it describes a machine.
\end{warningbox}


\section{Multiplication}
\label{plt:sec:alg-multiplication}

\begin{algorithm}[Truncated block product]\label{plt:alg:multiply}
\textbf{Input.}  Truncated representatives \(\mathcal T_1=(v_1,N_1)\),
\(\mathcal T_2=(v_2,N_2)\) over \(R\), with derived order bounds \(m_1,m_2\).\\
\textbf{Precondition.}  None beyond \cref{plt:def:alg-invariant}.\\
\textbf{Output.}  \((w,N)\) with \(N=\min\SetOf{N_1+m_2,N_2+m_1}\) and
\(w=\Trunc{v_1v_2}{N}\).\\
\textbf{Postcondition.}  \(FG\in\operatorname{Den}(w,N)\) for every
\(F\in\operatorname{Den}(\mathcal T_1)\), \(G\in\operatorname{Den}(\mathcal T_2)\),
and \(N\) is the largest cutoff with this property.\\
\textbf{Steps.}
\begin{enumerate}
\item Compute \(m_1,m_2\) from the canonical forms and set
      \(N\leftarrow\min\SetOf{N_1+m_2,N_2+m_1}\).  If \(N\le m_1+m_2\) return
      the zero representative \((0,N)\).
\item Initialise an empty ordered map \(w\).
\item For each stored pair \((i,v_1(i))\) and each stored pair
      \((j,v_2(j))\) with \(i+j<N\), accumulate
      \(w(i+j)\leftarrow w(i+j)+v_1(i)v_2(j)\).
\item Normalise every touched value in \(R\) and delete the zero values.
\end{enumerate}
\textbf{Cost.}  At most \(\binom{N-m_1-m_2+1}{2}\) multiplications in \(R\),
and exactly \(s_1s_2\) minus the number of skipped pairs when the supports have
\(s_1,s_2\) stored blocks.  In the scalar model the count is that of
\cref{plt:prop:alg-convolution-cost}.
\end{algorithm}

\begin{proof}[Correctness]
Soundness and sharpness of the declared \(N\) are
\cref{plt:thm:alg-precision-calculus}\ref{plt:thm:alg-precision-calculus:prod}.
That the loop computes \(\Trunc{v_1v_2}{N}\) is
\cref{plt:def:mul-cauchy} together with the finiteness of every fiber
(\cref{plt:lem:mul-finite-fiber}); no pair with \(i+j\ge N\) can contribute to
a retained block, so skipping such pairs is exact rather than approximate.
\end{proof}

\begin{lstlisting}[style=pseudo]
function multiply(T1, T2):
    m1 := order(T1);  m2 := order(T2)
    N  := min(T1.cutoff + m2, T2.cutoff + m1)
    w  := empty ordered map
    for (i, a) in T1.blocks:
        if i >= N - m2: break          # keys are increasing
        for (j, b) in T2.blocks:
            if i + j >= N: break
            w[i+j] := normalize(w[i+j] + a*b)
    return canonicalize(w, N)
\end{lstlisting}

The two \texttt{break} statements are legitimate exactly because the keys are
stored in increasing order; with a hash map they would be wrong, and the
routine would also lose the determinism required in
\cref{plt:sec:cert-determinism}.

\section{Reciprocal, division, and powers}
\label{plt:sec:alg-division}

\begin{algorithm}[Reciprocal by the linear recurrence]
\label{plt:alg:reciprocal-linear}
\textbf{Input.}  \(\mathcal T=(v,N_1)\) over \(R\), with \(m_1=m(\mathcal T)\).\\
\textbf{Precondition.}  \(v\ne0\) and the leading block \(v(m_1)\) is a unit of
\(R\).  For \(R=\K[L]\) this means \(v(m_1)\in\K^{\times}\), by
\cref{plt:lem:div-poly-units,plt:thm:div-unit-criterion}; if the test fails the
routine must not proceed but must report the obstruction, per (P4) of
\cref{plt:def:alg-invariant}.\\
\textbf{Output.}  \((w,N_1-2m_1)\) with \(w=\Trunc{\MultiplicativeInverseOf{v}}{N_1-2m_1}\).\\
\textbf{Steps.}  Factor \(v=t^{m_1}\widetilde v\) with
\(\widetilde v=b_0+b_1t+\cdots\), \(b_0\) a unit; set \(c_0\leftarrow
\MultiplicativeInverseOf{b_0}\) and, for \(n=1,\dots,\rho-1\) where
\(\rho=N_1-m_1\),
\[
 c_n\leftarrow-\MultiplicativeInverseOf{b_0}\sum_{j=1}^{n}b_jc_{n-j};
\]
return \(t^{-m_1}\sum_{n<\rho}c_nt^n\).\\
\textbf{Cost.}  Exactly \(\binom{\rho}{2}\) products \(b_jc_{n-j}\) in \(R\),
plus \(\rho-1\) multiplications by the precomputed
\(\MultiplicativeInverseOf{b_0}\) and one inversion.  In the scalar model,
\(\sum_{n=1}^{\rho-1}\sum_{j=1}^{n}(\deg b_j+1)(\deg c_{n-j}+1)\).
\end{algorithm}

Correctness is \cref{plt:thm:div-reciprocal}; the precision statement is
\cref{plt:thm:alg-precision-calculus}\ref{plt:thm:alg-precision-calculus:recip},
and the coefficient-model count restates \cref{plt:prop:div-cost-naive} in the
present bookkeeping.

\begin{algorithm}[Reciprocal by Newton doubling]\label{plt:alg:reciprocal-newton}
\textbf{Input.}  As in \cref{plt:alg:reciprocal-linear}.\\
\textbf{Precondition.}  As in \cref{plt:alg:reciprocal-linear}.\\
\textbf{Output.}  The same representative as \cref{plt:alg:reciprocal-linear}.\\
\textbf{Steps.}  With \(\widetilde v\) as above, set
\(C\leftarrow\MultiplicativeInverseOf{b_0}\), \(\mu\leftarrow1\); while
\(\mu<\rho\), set \(\mu'\leftarrow\min\SetOf{2\mu,\rho}\) and
\[
 C\leftarrow\Trunc{C\,(2-\widetilde vC)}{\mu'},
 \qquad \mu\leftarrow\mu' .
\]
Return \(t^{-m_1}C\).\\
\textbf{Cost.}  \(\Ceiling{\log_2\rho}\) iterations, the \(i\)-th using two
truncated products at cutoff \(2^i\).
\end{algorithm}

\begin{proof}[Correctness]
The update squares the residual: with \(E=1-\widetilde vC\),
\[
 1-\widetilde v\,C(2-\widetilde vC)
 =1-2\widetilde vC+(\widetilde vC)^2=E^2 .
\]
Hence \(\ordt(E)\ge\mu\) implies \(\ordt(E^2)\ge2\mu\), and by
\cref{plt:thm:alg-precision-calculus}\ref{plt:thm:alg-precision-calculus:recip}
the reciprocal is then known to precision \(2\mu\).  That the iterates are
exactly the truncations of the reciprocal, so that no spurious blocks are
produced, is \cref{plt:cor:div-newton-exact}; the precision schedule and the
crossover with \cref{plt:alg:reciprocal-linear} are
\cref{plt:thm:div-newton,plt:prop:div-cost-newton,plt:cor:div-doubling-pays}.
\end{proof}

\begin{warningbox}
Newton doubling is stated for the \emph{reciprocal}, an operation on the
multiplicative structure.  It has nothing to do with the compositional inverse
beyond a formal analogy, and the two must not share a symbol: this volume
writes \(\MultiplicativeInverseOf{F}\) and \(\rev{F}\) for the two operations
and never a bare superscript \(-1\).
\end{warningbox}

\begin{algorithm}[Truncated division]\label{plt:alg:divide}
\textbf{Input.}  \(\mathcal T_1=(v_1,N_1)\) (numerator) and
\(\mathcal T_2=(v_2,N_2)\) (denominator).\\
\textbf{Precondition.}  The leading block of \(v_2\) is a unit of \(R\).\\
\textbf{Output.}  \((w,N)\) with
\(N=\min\SetOf{N_1-m_2,\ N_2+m_1-2m_2}\), equivalently the representative of
relative precision \(\min\SetOf{\rho_1,\rho_2}\) and order \(m_1-m_2\).\\
\textbf{Steps.}  Either compose \cref{plt:alg:reciprocal-linear} or
\cref{plt:alg:reciprocal-newton} with \cref{plt:alg:multiply}, or run the direct
recurrence of \cref{plt:cor:div-quotient} on the normalised operands, which
saves one temporary and \(\rho\) coefficient products.\\
\textbf{Cost.}  \(\binom{\rho}{2}+\BigO(\rho)\) coefficient multiplications for
the direct recurrence, where \(\rho=\min\SetOf{\rho_1,\rho_2}\).
\end{algorithm}

Correctness of the truncated recurrence is \cref{plt:thm:div-truncated}; that
the declared relative precision is exactly right, and in particular that the
numerator need not be known to a larger absolute cutoff than the denominator,
is \cref{plt:prop:div-relative-precision} together with
\cref{plt:cor:alg-relative-precision}.  The profile of the answer is
\cref{plt:prop:div-degree-profile}; it is in general \emph{not} bounded by the
profiles of the inputs, which is why the scalar cost of the recurrence cannot
be read off from the inputs alone.

For powers, two routines are available and they are not interchangeable.

\begin{proposition}[Power recurrence]\label{plt:prop:alg-power-recurrence}
Let \(\K\) have characteristic zero, let \(B=\sum_{n\ge0}b_nt^n\in\PLpow\) with
\(b_0\) a unit of \(R\), and let \(\alpha\in\K\).  Define \(C=B^{\alpha}\)
by \cref{plt:thm:mul-binomial}, that is,
\(C=b_0^{\alpha}(1+U)^{\alpha}\) with \(U=\MultiplicativeInverseOf{b_0}B-1\) of
positive order and a chosen value \(b_0^{\alpha}\in R\).  Then
\(C=\sum_{n\ge0}c_nt^n\) satisfies \(c_0=b_0^{\alpha}\) and, for \(n\ge1\),
\begin{equation}\label{plt:eq:alg-power-recurrence}
 n\,b_0c_n=\sum_{j=1}^{n}\bigl((\alpha+1)j-n\bigr)b_jc_{n-j}.
\end{equation}
\end{proposition}

\begin{proof}
Let \(\partial_t\) be the formal partial derivative in \(t\) with \(L\) held
fixed; it is an \(R\)-linear derivation of \(\PLpow\), and it is continuous for
the \(t\)-adic filtration because \(\partial_t(t^N\PLpow)\subseteq
t^{N-1}\PLpow\).  Write \(C=b_0^{\alpha}(1+U)^{\alpha}
=b_0^{\alpha}\sum_{k\ge0}\binom{\alpha}{k}U^{k}\), a summable family since
\(\ordt(U^k)\ge k\) (\cref{plt:lem:mul-summable}).  Applying \(\partial_t\)
termwise, which is legitimate for a summable family by continuity, and using
\(k\binom{\alpha}{k}=\alpha\binom{\alpha-1}{k-1}\),
\[
 \partial_tC
 =b_0^{\alpha}\sum_{k\ge1}k\binom{\alpha}{k}U^{k-1}\partial_tU
 =\alpha\,b_0^{\alpha}(1+U)^{\alpha-1}\partial_tU .
\]
Multiplying by \(B=b_0(1+U)\) and using \(\partial_tB=b_0\partial_tU\) gives
the differential relation
\begin{equation}\label{plt:eq:alg-power-ode}
 B\,\partial_tC=\alpha\,C\,\partial_tB .
\end{equation}
Extracting the coefficient of \(t^{n-1}\) from
\cref{plt:eq:alg-power-ode} yields
\[
 \sum_{j=0}^{n}(n-j)\,b_jc_{n-j}=\alpha\sum_{j=1}^{n}j\,b_jc_{n-j},
\]
that is, \(n b_0c_n+\sum_{j=1}^{n}(n-j)b_jc_{n-j}
=\alpha\sum_{j=1}^{n}jb_jc_{n-j}\), which rearranges to
\cref{plt:eq:alg-power-recurrence}.  Division by \(n\) uses characteristic
zero.
\end{proof}

\begin{algorithm}[Powers]\label{plt:alg:power}
\textbf{Input.}  \(\mathcal T=(v,N_1)\), an exponent \(\alpha\).\\
\textbf{Precondition.}  For \(\alpha\in\NaturalNumbers\): none.  For general
\(\alpha\in\K\): \(m_1=0\), \(b_0\) a unit of \(R\), and a chosen value
\(b_0^{\alpha}\in R\) recorded as branch data.\\
\textbf{Output.}  Relative precision \(\rho_1\), order \(\alpha m_1\)
(an integer when \(\alpha\in\IntegerNumbers\); otherwise the computation lives
in a Puiseux extension and must be typed accordingly).\\
\textbf{Steps.}  For \(\alpha\in\NaturalNumbers\) either binary powering by
\cref{plt:alg:multiply} or \cref{plt:eq:alg-power-recurrence}.  For general
\(\alpha\), \cref{plt:eq:alg-power-recurrence} only.\\
\textbf{Cost.}  Binary powering: \(\BigO(\log\alpha)\) truncated products,
hence \(\BigO(\rho^2\log\alpha)\) coefficient multiplications.
Recurrence~\cref{plt:eq:alg-power-recurrence}: exactly \(\binom{\rho}{2}\)
products \(b_jc_{n-j}\), independent of \(\alpha\).
\end{algorithm}

\begin{remark}
Binary powering is preferable only for very small \(\alpha\); the recurrence is
independent of \(\alpha\) altogether and is the \emph{only} option for
\(\alpha\notin\NaturalNumbers\), including the formal roots of
\cref{plt:cor:mul-roots}.  The profile of the answer is
\cref{plt:prop:mul-power-profile}.
\end{remark}

\section{Logarithm and exponential of a positive-order argument}
\label{plt:sec:alg-explog}

\begin{proposition}[Coefficient recurrences for \(\log\) and \(\exp\)]
\label{plt:prop:alg-explog-recurrences}
Let \(\K\) have characteristic zero and let
\(U=\sum_{n\ge1}u_nt^n\in t\PLpow\).
\begin{enumerate}
\item If \(\Lambda=\log(1+U)=\sum_{n\ge1}\lambda_nt^n\), then for \(n\ge1\)
\begin{equation}\label{plt:eq:alg-log-recurrence}
 \lambda_n=u_n-\frac1n\sum_{j=1}^{n-1}j\,\lambda_j\,u_{n-j}.
\end{equation}
\item If \(E=\exp(U)=\sum_{n\ge0}e_nt^n\), then \(e_0=1\) and for \(n\ge1\)
\begin{equation}\label{plt:eq:alg-exp-recurrence}
 e_n=\frac1n\sum_{j=1}^{n}j\,u_j\,e_{n-j}.
\end{equation}
\end{enumerate}
Both series lie in \(\PLpow\), with \(\ordt(\Lambda)=\ordt(U)\) and
\(E\in1+t\PLpow\), by \cref{plt:lem:cmp-exp-log}.
\end{proposition}

\begin{proof}
Both families \(\sum_{m\ge1}(-1)^{m+1}U^m/m\) and \(\sum_{k\ge0}U^k/k!\) are
summable, since \(\ordt(U^m)\ge m\) (\cref{plt:lem:mul-summable}), so their sums
are defined and \(\partial_t\) may be applied termwise by the continuity noted
in the proof of \cref{plt:prop:alg-power-recurrence}.  Termwise differentiation
gives
\[
 \partial_t\Lambda=\Bigl(\sum_{m\ge1}(-1)^{m+1}U^{m-1}\Bigr)\partial_tU
 =\frac{\partial_tU}{1+U},
 \qquad
 \partial_tE=E\,\partial_tU,
\]
the first because \(\sum_{m\ge1}(-1)^{m+1}U^{m-1}\) is the reciprocal of
\(1+U\) (\cref{plt:thm:div-reciprocal}).  Rewrite the first as
\((1+U)\partial_t\Lambda=\partial_tU\) and extract the coefficient of
\(t^{n-1}\):
\[
 n\lambda_n+\sum_{\substack{i\ge1,\ j\ge1\\ i+j=n}}u_i\,j\lambda_j
 =n\,u_n,
\]
which is \cref{plt:eq:alg-log-recurrence} after writing \(i=n-j\).  From the
second, the coefficient of \(t^{n-1}\) is
\(ne_n=\sum_{j=1}^{n}j\,u_j\,e_{n-j}\), which is
\cref{plt:eq:alg-exp-recurrence}.  Division by \(n\) uses characteristic zero.
\end{proof}

\begin{algorithm}[Logarithm and exponential]\label{plt:alg:explog}
\textbf{Input.}  \(\mathcal T=(u,N)\) representing \(U\).\\
\textbf{Precondition.}  \(m(\mathcal T)\ge1\).  The routine must \emph{refuse}
an argument of order \(0\): \(\exp\) of a nonconstant block leaves the ring
(\(\exp(L)=X\)), and \(\log\) of a unit \(b_0(1+U)\) with \(b_0\notin\K\), or of
an element of nonzero order, produces \(\log b_0\) or a multiple of \(L\) and
must be handled by the generator formulas of
\cref{plt:prop:cmp-generator-substitution}, not by these recurrences.\\
\textbf{Output.}  Precision \(N\) in both cases; order \(\ge1\) for \(\log\),
order \(0\) for \(\exp\).\\
\textbf{Steps.}  \cref{plt:eq:alg-log-recurrence} or
\cref{plt:eq:alg-exp-recurrence} for \(n=1,\dots,N-1\).\\
% ed.: the two recurrences do not cost the same: the logarithm sums j<=n-1 and
% ed.: the exponential sums j<=n, so the exact counts differ by N-1.  The
% ed.: single figure binom(N,2) was right only for the exponential.
\textbf{Cost.}  Exactly \(\binom{N-1}{2}\) coefficient multiplications
\(\lambda_ju_{n-j}\) for \cref{plt:eq:alg-log-recurrence} and exactly
\(\binom{N}{2}\) coefficient multiplications \(u_je_{n-j}\) for
\cref{plt:eq:alg-exp-recurrence}, against \(\BigO(N)\) truncated products,
hence \(\BigO(N^3)\) coefficient multiplications, for the naive summation of
the defining series.
\end{algorithm}

\begin{proof}[Precision]
Perturb \(U\) by \(\eta\) with \(\ordt(\eta)\ge N\).  Then
\(\log(1+U+\eta)-\log(1+U)=\log\bigl(1+\eta\MultiplicativeInverseOf{(1+U)}\bigr)\)
has order \(\ordt(\eta)\ge N\), and
\(\exp(U+\eta)-\exp(U)=\exp(U)\bigl(\exp(\eta)-1\bigr)\) likewise has order
\(\ordt(\eta)\ge N\).  Both are sharp, as \(\eta=t^N\) shows.
\end{proof}

\section{The derivation and its iterates}
\label{plt:sec:alg-derivation}

\begin{algorithm}[The distinguished derivation]\label{plt:alg:derivative}
\textbf{Input.}  \(\mathcal T=(v,N)\) over \(R=\K[L]\).\\
\textbf{Precondition.}  \(R\) supports \(\partial_L\).\\
\textbf{Output.}  \((w,N+1)\) with
\(w(n+1)=(\partial_L-n)v(n)\) for every stored key \(n\), and order bound
\(m+1\).\\
\textbf{Cost.}  One \(\partial_L\) and one scalar multiple per stored block:
\(\BigO(N)\) coefficient operations, \(\BigO\bigl(\sum_n(\dprofile{F}(n)+1)\bigr)\)
scalar operations.  No coefficient \emph{multiplication} is required.
\end{algorithm}

Correctness is the block formula \(\DX(t^np)=t^{n+1}(\partial_L-n)p\); the
precision gain and its single exception are
\cref{plt:thm:alg-precision-calculus}\ref{plt:thm:alg-precision-calculus:deriv}.
Iterates use the shifted operator \(\shiftop{n}{k}=\prod_{j=0}^{k-1}(\partial_L-n-j)\),
with \(\shiftop{n}{0}\) the empty product and hence the identity:
\[
 \DX^k\bigl(t^np(L)\bigr)=t^{n+k}\shiftop{n}{k}p(L).
\]

\begin{proposition}[Cost of iterated differentiation]
\label{plt:prop:alg-iterated-derivative}
Computing \(\DX^k F\) from \(F\) by \(k\) applications of
\cref{plt:alg:derivative} costs \(\BigO(kN)\) coefficient operations, whereas
applying \(\shiftop{n}{k}\) directly to a \emph{single} block costs
\(\BigO(k)\).  Moreover, for \(p\in\K[L]\) and \(k\ge1\),
\[
 \shiftop{n}{k}p=0
 \iff
 p=0,\quad\text{or}\quad p\in\K\ \text{and}\ n\le0\le n+k-1,
\]
and
\(\deg\shiftop{n}{k}p=\deg p\) whenever \(n,n+1,\dots,n+k-1\) are all nonzero,
while \(\deg\shiftop{n}{k}p=\deg p-1\) when one of them vanishes and
\(\deg p\ge1\).
\end{proposition}

% ed.: the claim was "shiftop{n}{k}p = 0 as soon as k > deg p and n = 0",
% ed.: presumably by analogy with iterated partial_L.  It is false: only the
% ed.: single factor with index j = -n is partial_L, all the others are
% ed.: bijective, so the operator drops the degree by exactly one and no
% ed.: further.  Already shiftop{0}{2}L = (partial_L - 1)(1) = -1, nonzero
% ed.: with k = 2 > 1 = deg L.  The corrected kernel statement is what the
% ed.: annihilation clause of the precision calculus needs.
\begin{proof}
The operation counts are immediate.  Each factor \(\partial_L-s\) with
\(s\ne0\) maps a polynomial of degree \(d\ge0\) to one of degree \(d\), its
leading coefficient multiplied by \(-s\); in particular it is injective, and
being degree-preserving and \(\K\)-linear on each finite-dimensional space of
polynomials of degree \(\le d\) it is bijective on \(\K[L]\).  The factor
\(\partial_L\) lowers the degree by one and has kernel \(\K\).  The factors
commute, each being a polynomial in \(\partial_L\).  If none of
\(n,\dots,n+k-1\) vanishes, \(\shiftop{n}{k}\) is a composite of bijections,
whence injective and degree-preserving.  If exactly one of them vanishes ---
they are distinct consecutive integers, so at most one can --- write
\(\shiftop{n}{k}=\partial_L\circ Q\) with \(Q\) the product of the remaining
factors, a degree-preserving bijection mapping \(\K\) onto \(\K\).  Then
\(\shiftop{n}{k}p=0\) if and only if \(Qp\in\K\), that is, if and only if
\(p\in\K\); and for \(\deg p\ge1\) the degree of \(\shiftop{n}{k}p\) is
\(\deg Qp-1=\deg p-1\).
\end{proof}

Antidifferentiation is the inverse problem and is needed both by
integration-based algorithms and by the diagnosis of
\cref{plt:sec:fail-resonance}.

\begin{corollary}[Antiderivative]\label{plt:cor:alg-antiderivative}
The map \(\DX\colon\PLpow\to t\PLpow\) is surjective and \(\K\)-linear with
kernel \(\K\).  Concretely, given \(G=\sum_{m\ge1}g_mt^m\), the solutions
\(Y=\sum_{n\ge0}y_nt^n\) of \(\DX Y=G\) are obtained blockwise from
\((\partial_L-n)y_n=g_{n+1}\); for \(n\ge1\) the solution is unique, and for
\(n=0\) it is \(y_0=\int_0^Lg_1+c\) with \(c\in\K\) arbitrary.
\end{corollary}

\begin{proof}
The blockwise reduction is the block formula for \(\DX\).  For \(n\ne0\) the
operator \(\partial_L-n\) is bijective on \(\K[L]\): it preserves the filtration
by degree and acts on the \((d+1)\)-dimensional space of polynomials of degree
\(\le d\) by a triangular matrix with all diagonal entries \(-n\ne0\).  For
\(n=0\) it is \(\partial_L\), which is surjective on \(\K[L]\) in characteristic
zero with kernel \(\K\).  Assembling the blocks gives surjectivity onto
\(t\PLpow\) and kernel \(\K\).  This is the ring-level form of
\cref{plt:lem:mot-block-antiderivative}.
\end{proof}

\begin{warningbox}
\Cref{plt:cor:alg-antiderivative} says that an antiderivative routine is
\emph{not} a function of its input.  A routine that returns one antiderivative
without recording how the resonant constant was chosen is nondeterministic in
the sense of \cref{plt:def:cert-determinism}, and two correct runs may disagree
by an element of \(\K\), which is asymptotically larger than every block of
positive order.
\end{warningbox}

\section{Composition}
\label{plt:sec:alg-composition}

Two engines are available.  The Taylor engine of \cref{plt:thm:grp-taylor}
applies to a near-identity inner map; the generator engine of
\cref{plt:prop:cmp-generator-substitution} applies to every admissible
monomial-leading inner argument of \cref{plt:def:cmp-admissible-inner}.  They
compute the same thing (\cref{plt:thm:cmp-substitution-hom}) and are the
canonical pair of independent implementations for the differential test of
\cref{plt:sec:cert-tests}.

\begin{proposition}[Sharp Taylor budget, exclusive convention]
\label{plt:prop:alg-taylor-budget}
Let \(F\in\PLlau\) with \(\ordt(F)\ge a\), let \(H\in\PLpow\) with
\(\ordt(H)\ge h\ge0\), and let \(N\in\IntegerNumbers\).  Then
\[
 \Trunc{F\compose(X+H)}{N}
 =\Trunc{\ \sum_{k=0}^{K-1}\frac{H^k}{k!}\DX^kF\ }{N},
 \qquad
 K=\Ceiling{\frac{N-a}{h+1}},
\]
and no smaller \(K\) works for every such pair \(F\) and \(H\).
Equivalently, the last retained index is \(k=\Floor{(N-1-a)/(h+1)}\).
\end{proposition}

\begin{proof}
By \cref{plt:thm:grp-taylor} the identity
\(F\compose(X+H)=\sum_{k\ge0}H^k\DX^kF/k!\) is exact, and the family is
summable because \(\ordt(H^k\DX^kF)\ge kh+a+k=a+k(h+1)\) by
\cref{plt:lem:rev-deriv-gain}.  A summand can influence a block of index
\(<N\) only if \(a+k(h+1)<N\), that is, \(k<(N-a)/(h+1)\).
% ed.: the count was quoted without the hypothesis N>a under which it is
% ed.: valid; for N<=a the ceiling is nonpositive and counts nothing, which is
% ed.: the right answer only because both sides of the displayed identity then
% ed.: vanish.  The case is separated.
For \(N>a\) the number of nonnegative integers \(k\) satisfying this is
\(\Ceiling{(N-a)/(h+1)}\) and the largest is \(\Floor{(N-1-a)/(h+1)}\).  For
\(N\le a\) there is no such \(k\), the stated \(K\) is nonpositive, the sum is
empty, and both sides of the displayed identity vanish because
\(\ordt\bigl(F\compose(X+H)\bigr)=\ordt(F)\ge a\ge N\) by
\cref{plt:eq:alg-right-composition-isometry}.
% ed.: the witness F=t^a fails whenever a<0 and K-1>|a|, because then
% ed.: shiftop{a}{K-1} annihilates the constant 1 -- e.g. a=-1, h=0, N=10
% ed.: needs k=9 while DX^k t^{-1}=0 already for k=2.  Taking the block to be L
% ed.: instead of 1 removes the exception, by the kernel description in
% ed.: plt:prop:alg-iterated-derivative.
Sharpness: take \(F=t^{a}L\) and \(H=t^{h}\), so that \(\ordt F=a\) and
\(\ordt H=h\).  Then \(\DX^kF=t^{a+k}\shiftop{a}{k}L\), and
\(\shiftop{a}{k}L\ne0\) for every \(k\ge0\) because \(L\notin\K\)
(\cref{plt:prop:alg-iterated-derivative}).  Hence \(H^k\DX^kF/k!\) is a
nonzero multiple of \(t^{a+k(h+1)}\) for every \(k\), and the summand with
\(k=K-1\) contributes a nonzero block of index \(a+(K-1)(h+1)<N\), while all
summands with \(k\ge K\) have valuation at least \(N\).  Deleting the term
\(k=K-1\) therefore changes the truncation.
\end{proof}

% ed.: S1 records the budget as k <= floor((N-r)/(h+1)).  Under the exclusive
% cutoff adopted here that admits one Taylor term too many; the sharp exclusive
% bound is k <= floor((N-1-a)/(h+1)).  The discrepancy is exactly the
% inclusive/exclusive one-block conversion of the catalogue and is harmless
% for correctness but not for cost.

\begin{algorithm}[Near-identity composition, Taylor engine]
\label{plt:alg:compose-taylor}
\textbf{Input.}  \(\mathcal T_F=(v_F,N_F)\), \(\mathcal T_H=(v_H,N_H)\) with
\(m_H\ge0\).\\
\textbf{Precondition.}  \(m_H\ge0\), so that \(X+H\in\Gnear\).  An inner
argument of negative order is \emph{not} near-identity and must be routed to
\cref{plt:alg:compose-generators} after the dominant monomial has been
normalised.\\
\textbf{Output.}  Precision \(N=\min\SetOf{N_F,\ N_H+m_F+1}\), order \(m_F\).\\
% ed.: the running power was truncated at the output cutoff N.  That is wrong
% ed.: as soon as m_F<0: the term P D contributes to a block of index below N
% ed.: through every block of P of index below N-ordt(D), and ordt(D)>=m_F+k
% ed.: can be negative.  Truncating P at N-ordt(D), with the CURRENT D,
% ed.: repairs it and never demands a block of H beyond N_H, since
% ed.: N-ordt(D)<=N-m_F-k<=N_H+1-k.
\textbf{Steps.}  Set \(S\leftarrow0\), \(P\leftarrow1\), \(D\leftarrow v_F\),
\(k\leftarrow0\).  While \(\ordt(P)+\ordt(D)<N\): set
\(S\leftarrow\Trunc{S+PD/k!}{N}\); then \(D\leftarrow\DX D\),
\(P\leftarrow\Trunc{PH}{\,N-\ordt(D)}\), \(k\leftarrow k+1\).  Return
\(S\).\\
\textbf{Cost.}  \(K=\Ceiling{(N-m_F)/(m_H+1)}\) iterations, each one truncated
product and one derivative:
\(\BigO\bigl(K\rho^2\bigr)=\BigO\bigl(\rho^3/(m_H+1)\bigr)\) coefficient
multiplications, with \(\rho=N-m_F\).
\end{algorithm}

Precision and sharpness are
\cref{plt:thm:alg-precision-calculus}\ref{plt:thm:alg-precision-calculus:comp};
the stopping test is \cref{plt:prop:alg-taylor-budget}, implemented as the
dynamic test \(\ordt(P)+\ordt(D)\ge N\), which is never weaker than the static
bound and is often much stronger because \(\DX\) may annihilate blocks.  The
cutoff \(N-\ordt(D)\) carried by the running power is the exact one demanded
by
\cref{plt:thm:alg-precision-calculus}\ref{plt:thm:alg-precision-calculus:prod}:
the discarded part of \(P\) has valuation at least \(N-\ordt(D)\), so its
product with \(D\) has valuation at least \(N\) and cannot reach a retained
block.  For \(m_F\ge0\) this cutoff is at most \(N\) and the naive choice
\(N\) is merely wasteful; for \(m_F<0\), which is the normal case when the
outer series is Laurent, it is strictly larger than \(N\) and the naive
choice is wrong.

For the generator engine one needs the elementary but useful observation that
substituting into the logarithmic generator is a \emph{finite} Taylor shift.

\begin{proposition}[Polynomial shift]\label{plt:prop:alg-shift}
Let \(p\in\K[L]\) of degree \(d\) and let \(V\in\PLpow\).  Then, in \(\PLpow\),
\begin{equation}\label{plt:eq:alg-shift}
 p(L+V)=\sum_{k=0}^{d}\frac{(\partial_L^kp)(L)}{k!}\,V^k,
\end{equation}
a sum of \(d+1\) terms.
\end{proposition}

\begin{proof}
Both sides are \(\K\)-linear in \(p\), so it suffices to take \(p=L^{e}\) with
\(e\le d\).  The binomial theorem gives
\(p(L+V)=\sum_{k=0}^{e}\binom ek L^{e-k}V^k\), and
\(\partial_L^kL^{e}/k!=\binom ekL^{e-k}\).  Terms with \(k>e\) vanish on both
sides.  Characteristic zero is used only to divide by \(k!\).
\end{proof}

\begin{algorithm}[Composition by generator images]
\label{plt:alg:compose-generators}
\textbf{Input.}  \(\mathcal T_F=(v_F,N_F)\); an inner argument
\(G=cX^{\varrho}(1+U)\) with \(c\in\K^{\times}\),
\(\varrho\in\PositiveIntegers\), \(U\in t\PLpow\) given by
\(\mathcal T_U=(v_U,N_U)\); a chosen value \(\log c\in\K\).\\
\textbf{Precondition.}  \(G\) admissible in the sense of
\cref{plt:def:cmp-admissible-inner}; \(\log c\in\K\), failing which \(\K\) must
first be enlarged; \(\varrho\in\PositiveIntegers\), failing which the output
lives in a Puiseux or Hahn extension
(\cref{plt:prop:cmp-fractional-rho}) and the type must say so.\\
% ed.: the output line used a cutoff N that no line of the algorithm defined,
% ed.: and attributed it to plt:cor:cmp-precision-budget, which furnishes the
% ed.: INTERNAL budget for the auxiliary powers, not the output cutoff.  The
% ed.: output cutoff must be a function of the input cutoffs -- item (2) of
% ed.: plt:cor:cert-reproducible -- and the one the stability estimates
% ed.: deliver is min{varrho N_F, varrho m_F + N_U}.
\textbf{Output.}  \(\Trunc{F\compose G}{N}\) with
\[
 N=\min\SetOf{\varrho N_F,\ \varrho m_F+N_U},
\]
which is sound by \cref{plt:prop:cmp-stability}: perturbing the outer argument
below its cutoff moves the value by at least \(\varrho N_F\), and perturbing
\(U\) below \(N_U\) moves it by at least \(\varrho m_F+N_U\).  The internal
budget for the auxiliary powers is \cref{plt:cor:cmp-precision-budget}.\\
\textbf{Steps.}
\begin{enumerate}
\item \(V\leftarrow\log(1+U)\) by \cref{plt:alg:explog};
      \(R\leftarrow\MultiplicativeInverseOf{(1+U)}\) by
      \cref{plt:alg:reciprocal-linear} or \cref{plt:alg:reciprocal-newton}.
% ed.: ALGORITHM BUG, the generator-engine twin of the one repaired in
% ed.: plt:alg:compose-taylor.  The cap on the powers of V was min{D,N-1},
% ed.: justified by ordt(V^k)>=k alone.  That ignores the factor t^{varrho n}
% ed.: multiplying block n: the summand for the outer index n contributes to a
% ed.: block of index below N whenever varrho n + k < N, so for a Laurent
% ed.: outer series (m_F<0) powers with k>=N are still needed.  Concretely,
% ed.: F=t^{-1}L^5, varrho=1, N=1 has cap min{5,0}=0 and returns a wrong block
% ed.: of index -1, although V^1 is required.  The correct cap is
% ed.: min{D, N-varrho*m_F-1}, which reduces to at most N-1 exactly when
% ed.: m_F>=0.
\item Let \(D=\max_n\deg v_F(n)\) over retained blocks and let \(m_F\) be the
      order bound of \(\mathcal T_F\).  Precompute
      \(V^k/k!\) for \(1\le k\le\min\SetOf{D,\ N-\varrho m_F-1}\).  The second
      entry of the minimum is exactly the budget of
      \cref{plt:cor:cmp-precision-budget}, and it cannot be lowered: the
      summand carrying the outer index
      \(n\) is \(c^{-n}t^{\varrho n}R^{n}\,(\partial_L^kp_n)(\ldots)V^k/k!\),
      of valuation at least \(\varrho n+k\), so it can reach a block of index
      \(<N\) only for \(k<N-\varrho n\), and \(n\ge m_F\) with
      \(\varrho\ge1\) makes \(N-\varrho m_F-1\) the largest useful \(k\).  For
      \(m_F\ge0\) this cap is at most \(N-1\); for \(m_F<0\) it is strictly
      larger, and the value \(N-1\) is \emph{wrong}, not merely wasteful.
\item Precompute \(R^{n}\) incrementally for the retained outer indices \(n\)
      (for negative \(n\), \(R^{n}=(1+U)^{|n|}\), an ordinary positive power).
\item Accumulate
      \(\sum_n c^{-n}t^{\varrho n}R^{n}\,
        p_n(\varrho L+\log c+V)\)
      with the inner evaluation performed by \cref{plt:eq:alg-shift} applied to
      \(p_n(\varrho L+\log c+\,\cdot\,)\), truncating after every product.
\end{enumerate}
% ed.: the two step counts inherited the wrong cap and the wrong index range;
% ed.: the retained outer indices run from m_F to N-1, of which there are
% ed.: N-m_F, not N.
\textbf{Cost.}  At most \(\min\SetOf{D,\ N-\varrho m_F-1}\) truncated products
for the powers of \(V\), at most \(N-m_F\) for the powers of \(R\), and
\(\sum_n(\deg v_F(n)+1)\) polynomial-by-series products for the shifts: in the
coefficient-operation model \(\BigO\bigl((N+D)N^2\bigr)\).
\end{algorithm}

\begin{remark}[Why the powers of \(V\) are shared]
\Cref{plt:prop:alg-shift} converts every one of the \(\BigO(N)\) polynomial
evaluations into a fixed linear combination of the \emph{same} powers
\(V^k/k!\).  This is the precise content of the caching advice given by all six
sources.  Horner evaluation in \(L+V\), by contrast, performs
\(\deg p_n\) truncated products \emph{per block} and does not share work
between blocks; it is asymptotically worse by a factor \(N\).
\end{remark}

\begin{warningbox}
Composition changes \emph{both} generators.  For \(G=X+H=X(1+tH)\),
\[
 t\compose G=\frac{t}{1+tH},
 \qquad
 L\compose G=L+\log(1+tH),
\]
by \cref{plt:prop:cmp-generator-substitution}.  Substituting \(t\mapsto t+H\),
or updating one generator and not the other, is not an approximation but a
different and wrong operation; see \cref{plt:sec:fail-generator}.
\end{warningbox}

\section{Reversion by the four strategies}
\label{plt:sec:alg-reversion}

Throughout this subsection \(F=X+A\in\Gnear\) with \(A\in\PLpow\),
\(h=\ordt(A)\ge0\), and \(G=\rev{F}=X+B\).  All four strategies compute the
same \(B\), by uniqueness (\cref{plt:thm:rev-existence}).  All four are
precision-neutral in the following sense, so none of them needs guard blocks in
the outer argument.

\begin{proposition}[Reversion neither gains nor loses precision]
\label{plt:prop:alg-reversion-precision}
If \(F_1=X+A_1\) and \(F_2=X+A_2\) lie in \(\Gnear\) with
\(\ordt(A_1-A_2)\ge N\), then \(\ordt(\rev{F_1}-\rev{F_2})\ge N\).
Consequently a routine that receives \(A\) at precision \(N\) may return
\(B\) at precision \(N\), and no larger cutoff is sound.
\end{proposition}

\begin{proof}
Put \(G_2=\rev{F_2}\).  Then
\(F_1\compose G_2-X=(F_2\compose G_2-X)+(A_1-A_2)\compose G_2
=(A_1-A_2)\compose G_2\), whose order equals \(\ordt(A_1-A_2)\ge N\) by
\cref{plt:eq:alg-right-composition-isometry}.  By
\cref{plt:thm:cert-two-sided} below, applied to \(F_1\) and the candidate
\(G_2\),
\(\ordt(G_2-\rev{F_1})=\ordt(F_1\compose G_2-X)\ge N\).  Sharpness: take
\(A_1=A_2+t^{N}\); the same computation gives
\(\ordt(G_2-\rev{F_1})=N\) exactly.  The statement is the block form of
\cref{plt:prop:rev-truncation-safe}.
\end{proof}

\begin{algorithm}[Strategy 1: Picard iteration]\label{plt:alg:reverse-picard}
\textbf{Input.}  \(A\) at precision \(N\).\\
\textbf{Precondition.}  \(\ordt(A)\ge0\).\\
\textbf{Output.}  \(B\) at precision \(N\).\\
\textbf{Steps.}  \(B_0\leftarrow0\); \(B_{r+1}\leftarrow
\Trunc{-\,A\compose(X+B_r)}{r+1}\) for \(r=0,\dots,N-1\); return \(B_N\).\\
\textbf{Cost.}  \(N\) compositions, the \(r\)-th at cutoff \(r+1\).
\end{algorithm}

Correctness is the fixed-point equation \(B=-A\compose(X+B)\) of
\cref{plt:prop:rev-fixed-point-equation} together with the one-block gain of
\cref{plt:lem:rev-gain}, which gives \(\ordt(B-B_r)\ge r\); the limit exists by
completeness (\cref{plt:prop:grp-limit}).  Picard is the cheapest strategy to
implement and the most expensive to run; its virtue is that it needs nothing
but composition.

\begin{algorithm}[Strategy 2: triangular block solve]
\label{plt:alg:reverse-triangular}
\textbf{Input.}  \(A\) at precision \(N\).\\
\textbf{Precondition.}  \(\ordt(A)\ge0\); the leading slope of \(F\) is \(1\).
For leading slope \(c\in\K^{\times}\) divide each correction by \(c\)
(\cref{plt:prop:rev-general-slope}).\\
\textbf{Output.}  \(B\) at precision \(N\).\\
\textbf{Steps.}  \(G\leftarrow X\); for \(n=0,\dots,N-1\): compute the block of
index \(n\) of \(\RevErr{F}{G}\), call it \(r_n(L)\), and set
\(G\leftarrow G-r_n(L)t^n\).\\
\textbf{Cost.}  \(N\) residual evaluations; with the block-perturbation update
of \cref{plt:lem:rev-block-perturbation} each step touches only the blocks that
the new correction can reach, and the total is one composition-sized
computation rather than \(N\).
\end{algorithm}

Correctness is \cref{plt:thm:rev-triangular}; that the residual block reads off
the next inverse block is \cref{plt:prop:rev-next-block}.  This is the strategy
of choice when only a few blocks are wanted, when the answer is required in
expanded differential-polynomial form, or when a lazy interface must produce
blocks on demand.

\begin{algorithm}[Strategy 3: Newton reversion with doubling]
\label{plt:alg:reverse-newton}
\textbf{Input.}  \(A\) at precision \(N\); an initial \(G\) correct to
precision \(\mu_0\ge1\).\\
\textbf{Precondition.}  \(\ordt(A)\ge0\); \(\DX F\compose G\) is a unit, which
holds automatically since it equals \(1+\FormalBigOAt{t}{t}\)
(\cref{plt:lem:rev-newton-unit}).\\
\textbf{Output.}  \(B\) at precision \(N\).\\
% ed.: the loop tested a variable mu that the step list never initialised.
\textbf{Steps.}  Set \(\mu\leftarrow\mu_0\).  While \(\mu<N\): set
\(\mu'\leftarrow\min\SetOf{2\mu,N}\) and
\[
 G\leftarrow\Trunc{\,G-\frac{F\compose G-X}{(\DX F)\compose G}\,}{\mu'},
 \qquad \mu\leftarrow\mu'.
\]
Recompute \(\ordt(\RevErr{F}{G})\) after each step and abort if it is smaller
than \(\mu\).\\
\textbf{Cost.}  \(\Ceiling{\log_2(N/\mu_0)}\) stages; each stage one
composition, one derivative-composition and one division at the stage cutoff.
\end{algorithm}

Correctness in exact arithmetic is \cref{plt:thm:rev-newton-doubling}; the
behaviour with inexact data, and hence the guard precision each stage must
request, is \cref{plt:thm:rev-newton-guard}, and the schedule is
\cref{plt:cor:rev-newton-schedule}.

% ed.: S4's driver hard-codes the schedule as q := min(N, max(1, 2p+2)).  The
% additive +2 is a fudge for an unstated guard; the guard actually required is
% the one computed by plt:thm:rev-newton-guard together with the composition
% law of plt:thm:alg-precision-calculus(vi), and the abort test above makes any
% residual shortfall visible instead of silently absorbing it.

\begin{algorithm}[Strategy 4: Lagrange--B\"urmann streaming]
\label{plt:alg:reverse-lagrange}
\textbf{Input.}  \(A\) at precision \(N\).\\
\textbf{Precondition.}  \(\ordt(A)\ge h\ge0\).\\
\textbf{Output.}  \(B\) at precision \(N\).\\
% ed.: the cutoff carried by the running power was left as an ellipsis, which
% ed.: is exactly the datum a precision contract may not omit.  It is N-r+1:
% ed.: the block of index n of DX^{r-1}P is built from the block of index
% ed.: n-(r-1) of P, so the blocks of P of index < N-r+1 are the ones the
% ed.: accumulator can see, and no earlier truncation is thereby too coarse,
% ed.: since the cutoff demanded of P decreases by one at each step while the
% ed.: product law (plt:thm:alg-precision-calculus(ii)) delivers
% ed.: min{(N-r+2)+h, N+(r-1)h} >= N-r+2 from the previous power and A.
\textbf{Steps.}  \(B\leftarrow0\), \(P\leftarrow1\); for \(r=1,2,\dots\): set
\(P\leftarrow\Trunc{PA}{\,N-r+1\,}\), form \(Q\leftarrow\DX^{r-1}P\) by
\(r-1\) applications of \cref{plt:alg:derivative}, and accumulate
\(B\leftarrow\Trunc{B+(-1)^rQ/r!}{N}\).  Stop when
\(r(h+1)-1\ge N\), that is, after
\(R=\Ceiling{(N+1)/(h+1)}-1\) terms.\\
\textbf{Cost.}  \(R\) truncated products and \(\BigO(R^2)\) derivative passes;
memory \(\BigO(N)\), since only the current power is retained.
\end{algorithm}

\begin{proof}[Stopping rule]
By \cref{plt:thm:lag-burmann}, \(B=\sum_{r\ge1}((-1)^r/r!)\DX^{r-1}(A^r)\).
Since \(\ordt(A^r)\ge rh\) and each \(\DX\) raises the order by at least one
(\cref{plt:lem:rev-deriv-gain}), \(\ordt(\DX^{r-1}(A^r))\ge rh+r-1=r(h+1)-1\).
A summand can meet a block of index \(<N\) only if \(r(h+1)-1<N\), i.e.
\(r<(N+1)/(h+1)\), which is the stated bound.  For \(h=0\) it reads \(r\le N\);
for the inclusive statement ``through block \(N\)'' it reads \(r\le N+1\), which
is the source of the apparent one-block discrepancy discussed in
\cref{plt:sec:fail-offbyone}.
\end{proof}

\begin{algorithm}[Hybrid driver]\label{plt:alg:reverse-hybrid}
\textbf{Input.}  A map \(F\) whose dominant part need not be \(X\); a target
\(N\).\\
\textbf{Steps.}
\begin{enumerate}
\item Normalise: factor the dominant monomial of \(F\) and invert it exactly.
      If the leading exponent is not a positive integer, or if the leading
      derivative vanishes, stop: the problem is not near-identity reversion and
      the class must be changed (\cref{plt:sec:alg-decision}).
\item Run \cref{plt:alg:reverse-triangular} for two or three blocks.  This
      exposes the coefficient ring actually required and the correct leading
      slope before any expensive stage.
\item Run \cref{plt:alg:reverse-newton} with the doubling schedule to the
      target cutoff.
\item Compute \emph{both} residuals of \cref{plt:thm:cert-two-sided} and
      compare their orders.
\item If either order falls short of the target, enlarge the working cutoff and
      repeat from step~3; do not adjust the claimed precision instead.
\end{enumerate}
\end{algorithm}

\begin{proposition}[Comparison of the four strategies]
\label{plt:prop:alg-reversion-comparison}
Let \(C(\nu)\) denote the cost of one composition at cutoff \(\nu\) in the
coefficient-operation model, and assume \(C\) is nondecreasing and satisfies
\(C(2\nu)\ge2C(\nu)\).  Then, to reach precision \(N\) from \(A\) of order
\(h\ge0\):
\begin{enumerate}
\item Picard costs \(\sum_{r\le N}C(r)\) and uses only composition;
\item the triangular solve costs \(\BigO(C(N))\) when the block-perturbation
      update is used, and \(\sum_{n\le N}C(n)\) when it is not;
\item Newton costs \(\sum_{i\le\log_2N}\bigl(2C(2^i)+\BigO(4^i)\bigr)
      \le4C(N)+\BigO(N^2)\);
% ed.: the product count was quoted as the ceiling itself; the stopping rule
% ed.: of plt:alg:reverse-lagrange runs r=1..R with R = ceil((N+1)/(h+1)) - 1.
\item Lagrange--B\"urmann costs
      \(R=\Ceiling{(N+1)/(h+1)}-1\) truncated products and
      \(\BigO\bigl((N/(h+1))^2\bigr)\) derivative passes, and uses no
      composition at all.
\end{enumerate}
\end{proposition}

\begin{proof}
Parts (i)--(iii) are the step counts stated with the algorithms; the estimate
in (iii) uses
\[
 \sum_{i\le\log_2N}C(2^{i})\le2C(N),
\]
which follows by summing the geometric majorisation implied by
\(C(2\nu)\ge2C(\nu)\), and the division at
stage \(i\) costs \(\BigO(4^i)\) by \cref{plt:alg:divide}.  Part (iv) is the
stopping rule of \cref{plt:alg:reverse-lagrange}: the \(r\)-th term needs one
product to update \(A^r\) and \(r-1\) derivative passes, and
\(\sum_{r\le R}(r-1)=\BigO(R^2)\).
\end{proof}

\begin{remark}[Which to use]
Lagrange--B\"urmann is the only strategy that never composes, and it is
therefore the strategy of choice for theoretical derivations and for maps whose
powers and derivatives simplify --- the Lambert polynomials being the model
case, as \cref{plt:prop:fail-reversion-degree} shows.  Its weakness is that
its cost is insensitive to cancellation: it forms \(A^r\) for every \(r\) up to
the stopping index even when most terms cancel.  Newton dominates at high
\(N\), provided composition is itself efficient; Picard and the triangular
solve dominate at low \(N\) and are the only ones that give blocks lazily.
\end{remark}

\begin{summarybox}
Every routine in this section returns a coset, not an element.  Its contract is
the triple (order bound, absolute cutoff, coefficient ring), its correctness
obligation is soundness at the declared cutoff, and its cost must be quoted in
a named model.  Relative precision, not absolute cutoff, is the quantity that
% ed.: "only addition destroys it, and only differentiation improves it" is
% ed.: wrong twice: differentiation improves the ABSOLUTE cutoff by one block
% ed.: while leaving the relative precision unchanged, and in the annihilating
% ed.: case it destroys the relative precision outright.
products, quotients and compositions conserve.  Addition destroys it whenever
leading blocks cancel; differentiation gains one block of absolute cutoff and
conserves the relative precision, except in the single annihilating case
\((\partial_L-m_1)v_1(m_1)=0\), where it destroys that too.
\end{summarybox}

\chapter{Residual certificates and reproducible verification}
\label{plt:sec:certificates}

\Cref{plt:sec:algorithms} says what each routine must deliver.  This chapter
says how one finds out whether it did.  The mechanism is uniform: every
operation of this volume is characterised by an equation that its output
satisfies exactly, so a computed candidate can be substituted back into that
equation and the valuation of the result compared with the declared cutoff.
The mathematics needed to make this a proof rather than a habit is small but
it is not empty, and it is exactly the following three points, each of which
is made precise below.

\begin{itemize}
\item A residual certifies a \emph{coset}, and only when the residual equation
      is \emph{faithful} --- when the valuation of the residual determines the
% ed.: "every primitive has a faithful equation, with exactly one exception"
% ed.: overstated what plt:tab:cert-gauges contains: the table covers the
% ed.: primitives that invert a cheaper one, and lists no equation at all for
% ed.: the sum, the product and the derivation, which invert nothing.  For
% ed.: those the natural test is the Leibniz identity, and it is unfaithful in
% ed.: precisely the way plt:rmk:cert-derivative-defect records.  The claim is
% ed.: restricted to what is proved.
      valuation of the error.  Every primitive of
      \cref{plt:sec:algorithms} that inverts a cheaper one --- reciprocal,
      quotient, root, logarithm, exponential, composition, reversion --- has a
      faithful equation with an explicit offset
      (\cref{plt:tab:cert-gauges}), with exactly one exception, the
      antiderivative, and the exception is the resonant constant
      (\cref{plt:prop:cert-antiderivative}).  The three primitives that invert
      nothing --- addition, multiplication, the derivation --- have no entry
      in that table: they are certified by independent recomputation, or by
      the unfaithful derivative tests of
      \cref{plt:rmk:cert-derivative-defect}.
\item A residual test is a statement about coefficients in \(R\), decided by
      comparison of canonical forms.  Substituting numbers for \(L\) does not
      decide it (\cref{plt:prop:cert-exact-zero}).
\item A residual computed by the same code path that produced the answer, and
      in particular the residual that \emph{drove} an iteration, certifies
      nothing at all (\cref{plt:prop:cert-driven-residual}).  This is the
      precise reason for computing both residuals of a reversion, and the
      reason is not that one of them is more accurate: at finite precision
      they certify exactly as far as each other
      (\cref{plt:thm:cert-two-sided}\ref{plt:thm:cert-two-sided:precision}).
\end{itemize}

\section{Faithful defining equations}
\label{plt:sec:cert-discipline}

\begin{definition}[Defining equation, faithfulness, gauge]
\label{plt:def:cert-certificate}
Let \(\varphi\) be a partial operation on the ambient class, defined on a set
\(\mathcal D\) of input tuples with values in a declared solution class
\(\mathcal Y\).  A \emph{defining equation} for \(\varphi\) is a map \(E\)
that assigns to an input tuple \(z\in\mathcal D\) and a candidate
\(Y\in\mathcal Y\) an element \(E(Y;z)\) of the ambient class, computable from
\(Y\) and \(z\) by finitely many conforming routines, such that
\[
 E(Y;z)=0\iff Y=\varphi(z).
\]
The equation is \emph{faithful with gauge \(\kappa\in\IntegerNumbers\)} if
moreover
\begin{equation}\label{plt:eq:cert-gauge}
 \ordt E(Y;z)=\ordt\bigl(Y-\varphi(z)\bigr)+\kappa
 \qquad\text{for all }z\in\mathcal D,\ Y\in\mathcal Y,
\end{equation}
with the convention \(\ordt0=+\infty\).  A \emph{residual certificate} for a
computed \(Y\) at cutoff \(M\) is the pair consisting of the computed residual
\(E(Y;z)\) and the verified assertion \(\ordt E(Y;z)\ge M+\kappa\).
\end{definition}

\begin{proposition}[What a faithful certificate certifies]
\label{plt:prop:cert-certifies}
Let \(E\) be faithful with gauge \(\kappa\) and let \(M\in\IntegerNumbers\).
Then, for every \(Y\in\mathcal Y\),
\[
 \ordt E(Y;z)\ge M+\kappa
 \iff
 \ordt\bigl(Y-\varphi(z)\bigr)\ge M
 \iff
 \Trunc{Y}{M}=\Trunc{\varphi(z)}{M}.
\]
In particular the certificate is an equivalence, not an implication: a
candidate that passes is correct below \(M\), and a candidate that is correct
below \(M\) passes.
\end{proposition}

\begin{proof}
The first equivalence is \eqref{plt:eq:cert-gauge}.  The second is
\cref{plt:lem:alg-denotation-coset}: \(\ordt(Y-\varphi(z))\ge M\) says exactly
that \(Y-\varphi(z)\in t^{M}\PLpow\), which says that \(Y\) and
\(\varphi(z)\) have the same blocks of index \(<M\).
\end{proof}

The gauge is not a decoration: it is the number of blocks by which the
residual test must be run \emph{beyond} the cutoff one wants to certify, and
it is the commonest source of an apparently failing test on a correct answer.
For the reciprocal of a series of order \(m\) it equals \(m\), which is
negative exactly when the series has a pole.

\begin{proposition}[Exact zero testing]
\label{plt:prop:cert-exact-zero}
Let \(\K\) be infinite and let \(R=\K[L]\).  Let \(s\ge1\) and let
\(\lambda_1,\dots,\lambda_s\in\K\) be any finite list of sample points.  Then
there is a nonzero \(p\in R\) of degree \(s\) with \(p(\lambda_i)=0\) for all
\(i\), namely \(p=\prod_{i=1}^{s}(L-\lambda_i)\).  Consequently no test that
evaluates residual blocks at finitely many values of \(L\) can decide whether
a residual block vanishes, and a residual test is sound only when it compares
canonical forms in \(R\).
\end{proposition}

\begin{proof}
The displayed polynomial is monic of degree \(s\), hence nonzero, and vanishes
at each \(\lambda_i\).  A test that inspects only the values
\(p(\lambda_1),\dots,p(\lambda_s)\) cannot distinguish it from \(0\).
\end{proof}

\begin{remark}[The three coefficient rings test differently]
\label{plt:rmk:cert-zero-test-rings}
For \(R=\K[L]\) the canonical form is the coefficient vector and the test is a
comparison with the zero vector.  For \(R=\K(L)\) one must normalise the
fraction and test the numerator, since an unnormalised numerator and
denominator may share a factor; testing a rational function at sample points
is subject to \cref{plt:prop:cert-exact-zero} twice over, once for the
numerator and once for a spurious pole.  For \(R=\K((L^{-1}))\), a block is
stored only to a declared \emph{inner} cutoff, so the test decides vanishing
only below that cutoff, and the inner cutoff is a second precision parameter
that the type must carry and the certificate must quote
(\cref{plt:prop:ring-iteration-order}).  A test that silently compares a
polynomial block with a truncated Laurent block is comparing objects of two
different types.
\end{remark}

\begin{warningbox}
A certificate is a statement about the pair (candidate, defining equation).
It is not a statement about the routine that produced the candidate, and it is
not a statement about the routine that computes the residual.  If the residual
is computed by the very operation whose correctness is in question, a passing
test is consistent with that operation being arbitrarily wrong; see
\cref{plt:prop:cert-driven-residual} and the engine of
\cref{plt:ex:cert-wrong-engine}, which passes its own residual test
identically while getting every block after the first one wrong.
\end{warningbox}

\section{The two residuals of a reversion}
\label{plt:sec:cert-two-sided}

Throughout this subsection \(F=X+A\in\Gnear\) with \(A\in\PLpow\), and \(G\in
\Gnear\) is a candidate for \(\rev{F}\).  The two residuals are the right
residual \(\RevErr{F}{G}=F\compose G-X\) and the left residual
\(G\compose F-X\) of \cref{plt:def:rev-residual}.

\begin{theorem}[Two-sided residual certificate]
\label{plt:thm:cert-two-sided}
Let \(F,G\in\Gnear\) and put \(\Delta=G-\rev{F}\in\PLpow\); the letter \(E\) is
reserved for defining equations, as in \cref{plt:def:cert-certificate}.
\begin{enumerate}
\item\label{plt:thm:cert-two-sided:orders}
      \emph{Both equations are faithful with gauge \(0\):}
      \[
       \ordt\bigl(\RevErr{F}{G}\bigr)
       =\ordt\bigl(G\compose F-X\bigr)
       =\ordt(\Delta),
      \]
      and when \(\Delta\ne0\) all three series have the same leading block.
\item\label{plt:thm:cert-two-sided:cert}
      \emph{Certification.}  For every \(M\in\NaturalNumbers\), each of
      \(\ordt(\RevErr{F}{G})\ge M\) and \(\ordt(G\compose F-X)\ge M\) is
      equivalent to \(\Trunc{G-X}{M}=\Trunc{\rev{F}-X}{M}\).
\item\label{plt:thm:cert-two-sided:precision}
      \emph{Attainable certification cutoff.}  Suppose \(F\) and \(G\) are
      given by conforming representatives of the corrections \(A\) and
      \(G-X\), with finite cutoffs \(N_A\) and \(N_B\).  Then each of the two
      residuals is determined by that data exactly to the cutoff
      \(\min\SetOf{N_A,N_B}\), and to no larger cutoff.  Hence neither
      residual can certify \(G\) beyond block \(\min\SetOf{N_A,N_B}-1\), and
      the two residuals certify exactly as far as each other.
\end{enumerate}
\end{theorem}

\begin{proof}
Part~\ref{plt:thm:cert-two-sided:orders} is
\cref{plt:prop:rev-residual-order}, recorded here in the form the certificate
discipline uses; in the language of
\cref{plt:def:cert-certificate} it says that
\(E(G;F)=\RevErr{F}{G}\) and \(E'(G;F)=G\compose F-X\) are defining equations
for \(\varphi=\rev{\cdot}\) on \(\mathcal Y=\Gnear\), faithful with gauge
\(\kappa=0\).  Part~\ref{plt:thm:cert-two-sided:cert} is then
\cref{plt:prop:cert-certifies}; it contains
\cref{plt:cor:rev-finite-certificate}, which is the special case in which
\(G-X\) is a polynomial in \(t\) of degree \(M-1\).

For part~\ref{plt:thm:cert-two-sided:precision}, apply
\cref{plt:thm:alg-precision-calculus}\ref{plt:thm:alg-precision-calculus:comp}
to \(F\compose G\): the outer argument is \(F\), whose representative has
order bound \(m_F=-1\) --- the block \(X=t^{-1}\) --- and cutoff \(N_A\); the
inner correction is \(G-X\), with cutoff \(N_B\).  The declared output cutoff
is therefore
\[
 \min\SetOf{N_A,\ N_B+m_F+1}=\min\SetOf{N_A,N_B},
\]
and subtracting the exactly known \(X\) changes neither cutoff nor
determinacy.  Sharpness in the outer-limited regime is automatic; in the
inner-limited regime it needs \(\ordt(\DX v_F)=m_F+1=0\), which holds for
every admissible \(F\) because \(\DX X=1\) and every other block of \(v_F\)
has index \(\ge0\), so \(\ordt(\DX v_F)=0\).  The same computation with the
roles of \(F\) and \(G\) exchanged gives \(\min\SetOf{N_B,N_A}\) for the left
residual.  Since certification below \(M\) requires the residual to be known
to be zero on all blocks of index \(<M\), and blocks of index
\(\ge\min\SetOf{N_A,N_B}\) are not determined by the data, no certificate can
reach beyond \(\min\SetOf{N_A,N_B}-1\).
\end{proof}

\begin{remark}[Why compute both, then]
\label{plt:rmk:cert-both}
\Cref{plt:thm:cert-two-sided} removes two folklore reasons for computing both
residuals and leaves one.  It is \emph{not} true that one residual is sharper
than the other: by \ref{plt:thm:cert-two-sided:orders} they have the same
order and even the same leading block, and by
\ref{plt:thm:cert-two-sided:precision} they are determined to the same cutoff
by the same data.  It is \emph{not} true that the left residual detects an
error the right one cannot: by \ref{plt:thm:cert-two-sided:orders} either one
detects an error exactly when there is one.  What remains, and what is
decisive in practice, is independence: the two residuals evaluate the
composition routine at arguments in the opposite order, so a composition
implemented correctly in one order and incorrectly in the other is caught, and
--- far more importantly --- an iteration driven by one of them cannot be
certified by that one (\cref{plt:prop:cert-driven-residual}).
\end{remark}

\begin{proposition}[A driven residual certifies nothing]
\label{plt:prop:cert-driven-residual}
Let \(\odot\) be \emph{any} binary operation on \(\Gnear\)-representatives
whatsoever, and suppose the triangular loop of
\cref{plt:alg:reverse-triangular} is run with \(\odot\) in place of
\(\compose\), the block \(r_n\) at each step being read from
\(F\odot G-X\).  If the loop terminates having driven the blocks of
\(F\odot G-X\) of index \(<N\) to zero, then the assertion
``\(\ordt(F\odot G-X)\ge N\)'' holds by construction and is independent of
whether \(\odot=\compose\).  Consequently it carries no information about
\(G\), and in particular it does not certify \(\Trunc{G-X}{N}=\Trunc{\rev
F-X}{N}\).
\end{proposition}

\begin{proof}
The loop's exit condition is precisely the assertion in question: the loop
terminates only when the blocks of \(F\odot G-X\) of index \(<N\) have been
made to vanish, and each step changes \(G\) so as to make one more of them
vanish.  Nothing in the argument used any property of \(\odot\).  Hence the
assertion is a theorem about the loop, not about \(G\), and it holds equally
for an operation \(\odot\ne\compose\), for which the conclusion
\(\Trunc{G-X}{N}=\Trunc{\rev F-X}{N}\) is false; an explicit such \(\odot\) is
exhibited in \cref{plt:ex:cert-wrong-engine}.
\end{proof}

\begin{example}[An engine that passes its own test and is wrong from the
second block on]
\label{plt:ex:cert-wrong-engine}
Define \(\Psi_G\colon\PLlau\to\PLlau\) as the \(\K[L]\)-algebra homomorphism
that substitutes the outer generator and leaves the logarithmic generator
alone:
\[
 \Psi_G\Bigl(\sum_np_n(L)t^n\Bigr)
 \DefinitionEquals\sum_np_n(L)\bigl(\MultiplicativeInverseOf{G}\bigr)^{n},
 \qquad G=X+B\in\Gnear .
\]
This is well defined: \(\MultiplicativeInverseOf{G}=t/(1+tB)\) has
\(\ordt=1\), so the family is summable
(\cref{plt:lem:mul-summable}).  It is the ``substitute \(t\), forget \(L\)''
engine of \cref{plt:sec:fail-generator}.  Take \(F=X+L\), the map whose
inverse is \(\WrightOmega\) (\cref{plt:thm:om-expansion}).  Then
\[
 F\odot G\DefinitionEquals\Psi_G(F)=\Psi_G(t^{-1})+\Psi_G(L)=G+L,
\]
so the triangular loop driven by \(F\odot G-X=B+L\) sets \(b_0=-L\) and every
later block to \(0\), terminating with
\[
 \widetilde G=X-L
 \qquad\text{and}\qquad
 F\odot\widetilde G-X=0\ \text{identically.}
\]
The true inverse is \(\WrightOmega=X-L+Lt+(\tfrac12L^2-L)t^2+\cdots\), so
\(\widetilde G\) is correct in the block of index \(0\) and wrong in every
later block.  The true right residual detects this at once:
\[
 \RevErr{F}{\widetilde G}
 =(X-L)+\log(X-L)-X
 =\log(1-Lt)
 =-Lt-\tfrac12L^2t^2-\cdots,
\]
of order \(1\), and the true left residual is
\[
 \widetilde G\compose F-X=(X+L)-\log(X+L)-X=-\log(1+Lt)=-Lt+\tfrac12L^2t^2-\cdots,
\]
also of order \(1\), with the same leading block \(-L\) as predicted by
\cref{plt:thm:cert-two-sided}\ref{plt:thm:cert-two-sided:orders}.
\end{example}

The engine of \cref{plt:ex:cert-wrong-engine} is worth one more moment,
because it shows which structural tests can and cannot see a wrong generator
image.

\begin{proposition}[The ring tests are blind to the generator; the chain rule
is not]
\label{plt:prop:cert-chain-rule-sharp}
Let \(G=X+B\in\Gnear\) with \(B\ne0\), \(h=\ordt(B)\ge0\), and let
\(\Psi_G\) be as in \cref{plt:ex:cert-wrong-engine}.
\begin{enumerate}
\item \(\Psi_G\) is a \(\K[L]\)-algebra homomorphism of \(\PLlau\), continuous
      for the \(t\)-adic filtration.  Hence it satisfies every identity built
      from addition, multiplication, reciprocals and \(\K[L]\)-linearity ---
      in particular
      \(\Psi_G(F_1F_2)=\Psi_G(F_1)\Psi_G(F_2)\) and
      \(\Psi_G(\MultiplicativeInverseOf{F})=\MultiplicativeInverseOf{\Psi_G(F)}\)
      --- so no test built from those operations distinguishes it from
      composition.
\item \(\Psi_G\) fails the chain rule already on the single generator \(L\):
      \[
       \DX\bigl(\Psi_G(L)\bigr)-\Psi_G(\DX L)\,\DX G
       =t-\MultiplicativeInverseOf{G}\,\DX G,
      \]
      and this element has valuation exactly \(h+2\), with block
      \(-(\partial_L-h-1)b_h\) there, where \(b_h\) is the leading block of
      \(B\).
\end{enumerate}
\end{proposition}

\begin{proof}
(i)  \(\Psi_G\) is defined on generators by \(t\mapsto
\MultiplicativeInverseOf{G}\), \(L\mapsto L\), and extended by the summable
family in the definition; summability, since
\(\ordt\bigl(\bigl(\MultiplicativeInverseOf{G}\bigr)^{n}\bigr)=n\), makes the
extension a continuous
\(\K[L]\)-algebra map, exactly as in
\cref{plt:thm:cmp-substitution-hom} with the logarithmic generator held
fixed.  Multiplicativity and compatibility with reciprocals are formal
consequences.

(ii)  \(\Psi_G(L)=L\) and \(\DX L=t\), so the left side is
\(t-\Psi_G(t)\DX G=t-\MultiplicativeInverseOf{G}\DX G\).  Now
\(\MultiplicativeInverseOf{G}=t\MultiplicativeInverseOf{(1+tB)}\) and
\(\DX G=1+\DX B\), so
\[
 \MultiplicativeInverseOf{G}\,\DX G
 =t\,(1+\DX B)\bigl(1-tB+t^2B^2-\cdots\bigr)
 =t\Bigl[1+(\DX B-tB)+W\Bigr],
\]
where \(W\) collects the products of at least two of the small terms
\(\DX B\), \(tB\), each of valuation \(\ge h+1\); hence
\(\ordt W\ge2h+2\ge h+2\).  By the block formula the block of index \(h+1\) of
\(\DX B\) is \((\partial_L-h)b_h\) and that of \(tB\) is \(b_h\), so
\(\DX B-tB\) has block \((\partial_L-h-1)b_h\) at index \(h+1\).  Since
\(h+1\ne0\), the operator \(\partial_L-h-1\) is injective on \(\K[L]\)
(\cref{plt:prop:alg-iterated-derivative}), so that block is nonzero and
\(\ordt(\DX B-tB)=h+1<h+2\le\ordt W\).  Therefore
\(t-\MultiplicativeInverseOf{G}\DX G=-t(\DX B-tB)-tW\) has valuation exactly
\(h+2\) and block \(-(\partial_L-h-1)b_h\) there.
\end{proof}

\begin{remark}
\Cref{plt:prop:cert-chain-rule-sharp} explains why the chain-rule test is
listed as mandatory in \cref{plt:sec:cert-tests} even though it looks
redundant next to the composition residual: the composition residual computed
with a wrong engine is the object of
\cref{plt:prop:cert-driven-residual}, whereas the chain rule couples the
engine to the derivation, which is an independent primitive
(\cref{plt:alg:derivative}) with an independent certificate.
\end{remark}

\section{Certificates for the algebraic primitives}
\label{plt:sec:cert-primitives}

\begin{proposition}[Quotient and reciprocal certificates]
\label{plt:prop:cert-reciprocal}
Let \(B\in\PLlau\) be a unit, \(m=\ordt(B)\), and \(A\in\PLlau\).  For every
candidate \(Q\in\PLlau\),
\[
 \ordt(BQ-A)=\ordt\Bigl(Q-\frac AB\Bigr)+m .
\]
Thus \(E(Q)=BQ-A\) is a defining equation for the quotient, faithful with
gauge \(m\); the case \(A=1\) is the reciprocal, with the same gauge.
\end{proposition}

\begin{proof}
\(BQ-A=B\bigl(Q-A/B\bigr)\), and \(\ordt\) is additive on \(\PLlau\)
(\cref{plt:lem:div-valuation}), with \(\ordt(B)=m\) finite because \(B\ne0\).
\end{proof}

% ed.: S2 (eq. recip-residual-test) and S3 (eq. divisionresidual) both state
% the residual test as "AB_N - 1 = O(t^N)" or "BC_{<N} - A = O(t^N)" and read
% off correctness through block N-1.  That reading is the gauge-zero one and
% is correct only for a denominator of order 0.  With a denominator of order m
% the same residual certifies through block N-m-1, which is fewer blocks when
% m>0 and more when the denominator has a pole.
\begin{remark}[The gauge is where the sources' division test goes wrong]
\label{plt:rmk:cert-division-gauge}
Two of the six sources state the division certificate in the form ``residual
\(\FormalBigOAt{t^{N}}{t}\) certifies the quotient through block \(N-1\)''.
By \cref{plt:prop:cert-reciprocal,plt:prop:cert-certifies} the correct
statement is that such a residual certifies the quotient through block
\(N-m-1\), where \(m=\ordt(B)\).  The two agree exactly when the denominator
is normalised to order \(0\), which is the case the sources have in mind and
the case an implementation should enforce before testing --- but the
enforcement is then part of the contract and must be checked, not assumed.
\end{remark}

\begin{proposition}[Logarithm and exponential certificates]
\label{plt:prop:cert-explog}
Let \(U\in t\PLpow\) and let \(\partial_t\) be the formal derivative in \(t\)
at fixed \(L\).
\begin{enumerate}
\item For every candidate \(\Lambda\in t\PLpow\),
      \(\ordt\bigl(\exp\Lambda-(1+U)\bigr)=\ordt\bigl(\Lambda-\log(1+U)\bigr)\):
      gauge \(0\).
\item For every candidate \(\Lambda\in t\PLpow\),
      \(\ordt\bigl((1+U)\partial_t\Lambda-\partial_tU\bigr)
      =\ordt\bigl(\Lambda-\log(1+U)\bigr)-1\): gauge \(-1\).
\item For every candidate \(\mathcal E\in1+t\PLpow\),
      \(\ordt\bigl(\log\mathcal E-U\bigr)=\ordt\bigl(\mathcal E-\exp U\bigr)\):
      gauge \(0\).
\end{enumerate}
\end{proposition}

\begin{proof}
Write \(\Delta=\Lambda-\log(1+U)\in t\PLpow\).

(i) By \cref{plt:lem:cmp-exp-log},
\(\exp\Lambda=\exp\bigl(\log(1+U)\bigr)\exp\Delta=(1+U)\exp\Delta\), so
\(\exp\Lambda-(1+U)=(1+U)(\exp\Delta-1)\).  Now
\(\exp\Delta-1=\sum_{k\ge1}\Delta^k/k!\) and \(\ordt(\Delta^k)\ge k\ordt\Delta
>\ordt\Delta\) for \(k\ge2\) because \(\ordt\Delta\ge1\); hence
\(\ordt(\exp\Delta-1)=\ordt\Delta\).  Multiplying by the unit \(1+U\) of order
\(0\) preserves the order.

(ii) Termwise differentiation of the summable family defining the logarithm
gives \((1+U)\partial_t\log(1+U)=\partial_tU\) (proof of
\cref{plt:prop:alg-explog-recurrences}), so
\((1+U)\partial_t\Lambda-\partial_tU=(1+U)\partial_t\Delta\).  If
\(\Delta\ne0\) has \(\ordt\Delta=n\ge1\) and leading block \(\delta_n\ne0\),
then \(\partial_t\Delta\) has block \(n\delta_n\ne0\) at index \(n-1\), using
characteristic zero, and no block below; so
\(\ordt(\partial_t\Delta)=n-1\).  Multiplying by \(1+U\) preserves it.

(iii) With \(\Theta=\mathcal E-\exp U\), one has
\(\log\mathcal E-U=\log\bigl(\mathcal E\MultiplicativeInverseOf{(\exp
U)}\bigr) =\log\bigl(1+\Theta\MultiplicativeInverseOf{(\exp U)}\bigr)\), and
\(\exp U\) is a unit of order \(0\), so the argument of the logarithm lies in
\(1+t\PLpow\) and \(\ordt\log(1+\Xi)=\ordt\Xi\) for \(\Xi\in t\PLpow\) by
\cref{plt:lem:cmp-exp-log}.  Hence the order equals
\(\ordt\Theta\).
\end{proof}

\begin{remark}[The differential certificate is cheaper, not stronger]
\label{plt:rmk:cert-differential-tradeoff}
Part~(ii) of \cref{plt:prop:cert-explog} detects an error one block earlier
than part~(i), and it costs one truncated product and two derivatives instead
of an exponential.  It does not, however, certify further.  If the data are
known to cutoff \(N\), the residual of (i) is known to cutoff \(N\) and
certifies \(\ordt\Delta\ge M\) for \(M\le N\); the residual of (ii) is known
only to cutoff \(N-1\), because \(\partial_t\) shifts the filtration by one,
and its threshold is \(M-1\), so it too certifies exactly \(M\le N\).  The two
reach the same block for opposite reasons.
\end{remark}

\begin{proposition}[Power and root certificates]
\label{plt:prop:cert-power}
Let \(B\in\PLpow\) whose block of index \(0\) is a scalar
\(b_0\in\K^{\times}\) --- equivalently, by
\cref{plt:thm:div-unit-criterion}, a unit of \(\PLpow\) --- let
\(p\in\IntegerNumbers\) and
\(q\in\PositiveIntegers\), and let \(C=B^{p/q}\) be the power determined by a
chosen value \(b_0^{p/q}\in\K^{\times}\) satisfying
\(\bigl(b_0^{p/q}\bigr)^{q}=b_0^{p}\), the powers being those of
% ed.: the hypothesis was "whose leading scalar equals b_0^{p/q}", but the
% ed.: leading scalar of D is by plt:def:lead-leading-data the coefficient of
% ed.: its leading MONOMIAL t^0 L^j, which for a block of positive L-degree is
% ed.: not the block; the proof uses the block.  The hypothesis is restated in
% ed.: the form the proof needs and the counterexample clause requires.
\cref{plt:thm:mul-binomial}.  Then \(C^{q}=B^{p}\), and for every candidate
\(D\in\PLpow\) whose block of index \(0\) is the scalar \(b_0^{p/q}\),
\[
 \ordt\bigl(D^{q}-B^{p}\bigr)=\ordt(D-C):
\]
gauge \(0\).  The restriction on the block of index \(0\) of \(D\) cannot be
dropped, and the equation cannot replace it: if \(\zeta\in\K^{\times}\)
satisfies \(\zeta^{q}=1\) and \(\zeta\ne1\), then \(D=\zeta C\) satisfies
\(D^{q}-B^{p}=0\) identically while \(D\ne C\).  The branch datum is part of
the specification and is certified only by comparing the blocks of index
\(0\).
\end{proposition}

\begin{proof}
The exponent law \((1+U)^{a}(1+U)^{b}=(1+U)^{a+b}\) for binomial powers is the
Vandermonde convolution
\(\sum_{k}\binom ak\binom b{n-k}=\binom{a+b}n\)
\cite{graham-knuth-patashnik}, valid in any \(\RationalNumbers\)-algebra;
iterating it \(q\) times with \(a=b=\cdots=p/q\) gives
\((1+U)^{p}\) where \(U=\MultiplicativeInverseOf{b_0}B-1\).  Combined with the
hypothesis \(\bigl(b_0^{p/q}\bigr)^{q}=b_0^{p}\) this yields \(C^{q}=B^{p}\).
Next,
\[
 D^{q}-C^{q}=(D-C)\sum_{i=0}^{q-1}D^{i}C^{\,q-1-i},
\]
and each summand has order \(0\) with block of index \(0\) equal to
\(d_0^{i}c_0^{\,q-1-i}\), where \(c_0=b_0^{p/q}\) and \(d_0\) is the block of
index \(0\) of \(D\).  Under the hypothesis \(d_0=c_0\) the sum has block
\(qc_0^{\,q-1}\ne0\) there, because \(\K\) has characteristic zero and \(c_0\ne0\);
hence the sum has order exactly \(0\) and the displayed identity follows from
additivity of \(\ordt\).  For the final assertion, \((\zeta C)^{q}=\zeta^{q}
C^{q}=C^{q}=B^{p}\) while \(\zeta C-C=(\zeta-1)C\ne0\).
\end{proof}

\begin{proposition}[A faithful composition certificate]
\label{plt:prop:cert-composition}
Let \(F\in\PLlau\), \(G\in\Gnear\), and let \(H\in\PLlau\) be a candidate for
\(F\compose G\).  Then
\[
 \ordt\bigl(H\compose\rev{G}-F\bigr)=\ordt\bigl(H-F\compose G\bigr):
\]
gauge \(0\).
\end{proposition}

\begin{proof}
Right composition with \(\rev{G}\in\Gnear\) is a ring homomorphism of
\(\PLlau\) (\cref{plt:thm:cmp-substitution-hom}) and satisfies
\((F\compose G)\compose\rev{G}=F\compose(G\compose\rev{G})=F\) by
associativity on the admissible domain
(\cref{plt:thm:cmp-associativity}) and \(G\compose\rev G=X\)
(\cref{plt:thm:grp-group}).  Hence
\(H\compose\rev G-F=(H-F\compose G)\compose\rev G\), whose order equals
\(\ordt(H-F\compose G)\) by
\cref{plt:eq:alg-right-composition-isometry}.
\end{proof}

\begin{remark}[Composition has one faithful certificate and several cheap
unfaithful ones]
\label{plt:rmk:cert-composition}
\Cref{plt:prop:cert-composition} is the only single-equation certificate for
composition in this volume, and it is not free: it requires the inverse of the
inner map, hence an independent reversion.  The cheap alternatives --- the
chain rule (\cref{plt:thm:cmp-chain-rule}), the multiplicativity
\((F_1F_2)\compose G=(F_1\compose G)(F_2\compose G)\), associativity, and the
agreement of the two engines of \cref{plt:sec:alg-composition} --- are
\emph{consistency} tests: each of them is an identity that a wrong engine may
or may not satisfy.  \Cref{plt:prop:cert-chain-rule-sharp} shows that the
multiplicative tests are satisfied by a wrong engine, and the chain rule is
not, so among the cheap tests the chain rule is the one to run.
\end{remark}

\begin{proposition}[The antiderivative equation is the one unfaithful
defining equation]
\label{plt:prop:cert-antiderivative}
Let \(G\in t\PLpow\).
\begin{enumerate}
\item The equation \(E(Y)=\DX Y-G\) does not define a single-valued operation
      on \(\PLpow\): its solution set is a coset of \(\K\)
      (\cref{plt:cor:alg-antiderivative}), and \(E\) is not faithful with any
      gauge, since \(E(Y+c)=E(Y)\) while \(\ordt\bigl((Y+c)-Y\bigr)=0\) for
      \(c\in\K^{\times}\).
\item Fix \(c\in\K\) and let
      \(\mathcal Y_c=\SetOf{Y\in\PLpow:\ (\text{block of index }0\text{ of }Y)
      \text{ takes the value }c\text{ at }L=0}\).  Then \(E\) restricted to
      \(\mathcal Y_c\) is a defining equation for a single-valued operation,
      and it is faithful with gauge \(1\).
\end{enumerate}
\end{proposition}

\begin{proof}
(i) is \cref{plt:cor:alg-antiderivative} together with the displayed
observation.  For (ii), let \(Y,Y^{*}\in\mathcal Y_c\) with
\(\DX Y^{*}=G\) and put \(\Delta=Y-Y^{*}\), so that \(E(Y)=\DX\Delta\).  The
block of index \(0\) of \(\Delta\) vanishes at \(L=0\), so it is not a nonzero
constant.  If \(\ordt\Delta=0\), its leading block \(\Delta_0\) is a nonzero
polynomial with \(\Delta_0(0)=0\), hence \(\deg\Delta_0\ge1\) and
\(\partial_L\Delta_0\ne0\) in characteristic zero; if \(\ordt\Delta=n\ge1\),
then \((\partial_L-n)\Delta_n\ne0\) because \(\partial_L-n\) is injective on
\(\K[L]\).  In both cases the block formula gives
\(\ordt(\DX\Delta)=\ordt(\Delta)+1\).  Uniqueness of the solution in
\(\mathcal Y_c\) is the case \(\Delta=0\).
\end{proof}

\begin{remark}[Every derivative-based certificate inherits the same defect]
\label{plt:rmk:cert-derivative-defect}
The Leibniz test \(\DX H-(F\,\DX G+G\,\DX F)\) for a candidate product \(H\),
and the chain-rule test for a candidate composition, are residuals of the form
\(\DX(\text{error})\).  By
\cref{plt:thm:alg-precision-calculus}\ref{plt:thm:alg-precision-calculus:deriv}
they have gauge \(1\) \emph{except} when the leading block of the error is a
constant sitting at index \(0\), which they annihilate.  A derivative-based
test therefore cannot see an error in \(\K\), which is exactly the error an
antiderivative can commit.  Pair them with one gauge-zero test.
\end{remark}

\begin{longtable}{@{}p{0.25\textwidth} p{0.34\textwidth} p{0.13\textwidth} p{0.20\textwidth}@{}}
\caption{Defining equations, their gauges, and where they are proved.  The
gauge is the number of blocks by which the residual test must be run beyond
the cutoff to be certified (\cref{plt:prop:cert-certifies}).}
\label{plt:tab:cert-gauges}\\
\toprule
Operation and candidate & Defining equation \(E\) & Gauge & Reference\\
\midrule
\endfirsthead
\toprule
Operation and candidate & Defining equation \(E\) & Gauge & Reference\\
\midrule
\endhead
Quotient \(Q\) of \(A/B\) & \(BQ-A\) & \(\ordt B\) &
  \cref{plt:prop:cert-reciprocal}\\
Reciprocal \(C\) of \(\MultiplicativeInverseOf{B}\) & \(BC-1\) & \(\ordt B\) &
  \cref{plt:prop:cert-reciprocal}\\
Rational power \(D\) of \(B^{p/q}\), \(\ordt B=0\) & \(D^{q}-B^{p}\) & \(0\) &
  \cref{plt:prop:cert-power}\\
Logarithm \(\Lambda\) of \(\log(1+U)\) & \(\exp\Lambda-(1+U)\) & \(0\) &
  \cref{plt:prop:cert-explog}(i)\\
Logarithm \(\Lambda\), differential form &
  \((1+U)\partial_t\Lambda-\partial_tU\) & \(-1\) &
  \cref{plt:prop:cert-explog}(ii)\\
Exponential \(\mathcal E\) of \(\exp U\) & \(\log\mathcal E-U\) & \(0\) &
  \cref{plt:prop:cert-explog}(iii)\\
Composition \(H\) of \(F\compose G\), \(G\in\Gnear\) &
  \(H\compose\rev{G}-F\) & \(0\) & \cref{plt:prop:cert-composition}\\
Reversion \(G\) of \(\rev{F}\), right & \(F\compose G-X\) & \(0\) &
  \cref{plt:thm:cert-two-sided}\\
Reversion \(G\) of \(\rev{F}\), left & \(G\compose F-X\) & \(0\) &
  \cref{plt:thm:cert-two-sided}\\
Omega truncation \(\WrightOmega_N\) &
  \(\WrightOmega_N+\log\WrightOmega_N-X\) & \(0\) &
  \cref{plt:thm:om-residual}\\
Antiderivative \(Y\) with declared constant &
  \(\DX Y-G\) & \(1\) on \(\mathcal Y_c\) &
  \cref{plt:prop:cert-antiderivative}\\
Antiderivative \(Y\), constant undeclared & \(\DX Y-G\) & none &
  \cref{plt:prop:cert-antiderivative}(i)\\
\bottomrule
\end{longtable}

\section{A specification for a self-checking implementation}
\label{plt:sec:cert-tests}

The tests below are chosen so that each one fails on a specific error that
this volume documents, and so that no test is checked by the code path it
tests.  They are stated as a specification: an implementation conforms when it
runs all of them at every build, over the declared coefficient ring, with
exact arithmetic.

Two of them need an object to test against.  The Wright omega expansion is the
canonical choice, because it is the reversion of the single map \(X+L\), its
% ed.: "four mutually independent identities" is not what is meant and is not
% ed.: true: all four characterise the same family, so they are logically
% ed.: equivalent, not independent.  What is claimed, and what matters, is
% ed.: that they share no code path.
blocks are the Lambert polynomials, and it satisfies four identities that
share no code path: the reversion residual (\cref{plt:thm:om-residual}), the integral
recurrence \eqref{plt:eq:lambert-int-rec}, the Stirling formula
(\cref{plt:thm:lambert-stirling}), and the block recurrence of the
differentiated equation, which we now record because it is the only one of the
four that exercises \(\DX\) and the truncated product rather than
\(\partial_L\) alone.

\begin{lemma}[The block recurrence of the omega equation, and the constant it
cannot see]
\label{plt:lem:cert-omega-ode}
Let \(\mathcal Q=\sum_{n\ge0}q_n(L)t^{n}\in\PLpow\) and
\(\Omega=X+\mathcal Q\in\Gnear\).
\begin{enumerate}
\item If \(\Omega+\log\Omega=X+c\) for some \(c\in\K\), then
\begin{equation}\label{plt:eq:cert-omega-ode}
 (1+t+t\mathcal Q)\,\DX\mathcal Q=-t .
\end{equation}
\item Conversely, \eqref{plt:eq:cert-omega-ode} holds if and only if
      \(\partial_Lq_0=-1\) and, for every \(n\ge1\),
\begin{equation}\label{plt:eq:cert-omega-recurrence}
 (\partial_L-n)q_n
 =-(\partial_L-n+1)q_{n-1}
 -\sum_{j=1}^{n}q_{n-j}\,(\partial_L-j+1)q_{j-1}.
\end{equation}
      In particular \eqref{plt:eq:cert-omega-ode} determines \(q_n\) uniquely
      from \(q_0,\dots,q_{n-1}\) for every \(n\ge1\), and determines \(q_0\)
      only up to an additive constant.
\item The solutions of \eqref{plt:eq:cert-omega-ode} in \(\PLpow\) are exactly
      the corrections \(\rev{(X+L-c)}-X\) for \(c\in\K\), one for each value
      of \(q_0(0)=c\); only \(c=0\) solves \(\Omega+\log\Omega=X\), and the
      residual of \cref{plt:thm:om-residual} detects \(c\ne0\) already in the
      block of index \(0\).
\end{enumerate}
\end{lemma}

\begin{proof}
(i)  Apply \(\DX\) to \(\Omega+\log\Omega=X+c\).  By
\cref{plt:lem:dif-logderiv}, \(\DX\log\Omega=(\DX\Omega)\MultiplicativeInverseOf{\Omega}\),
and \(\DX(X+c)=1\), so \((\DX\Omega)(\Omega+1)=\Omega\).  Substituting
\(\Omega=t^{-1}(1+t\mathcal Q)\), \(\Omega+1=t^{-1}(1+t+t\mathcal Q)\) and
\(\DX\Omega=1+\DX\mathcal Q\), and multiplying by \(t\), gives
\((1+\DX\mathcal Q)(1+t+t\mathcal Q)=1+t\mathcal Q\), which rearranges to
\eqref{plt:eq:cert-omega-ode}.

(ii)  By the block formula, \(\DX\mathcal Q=\sum_{m\ge1}d_mt^{m}\) with
\(d_m=(\partial_L-m+1)q_{m-1}\), and \(d_0=0\).  Extracting the coefficient of
\(t^{m}\) in \eqref{plt:eq:cert-omega-ode} gives
\[
 d_m+d_{m-1}+\sum_{j=1}^{m-1}q_{m-1-j}d_j
 =-\,[m=1]
 \qquad(m\ge1),
\]
where \([m=1]\) is the Iverson bracket and the sum is empty for \(m=1\).  The case \(m=1\) is \(d_1=-1\), that is
\(\partial_Lq_0=-1\); the case \(m=n+1\ge2\) is
\eqref{plt:eq:cert-omega-recurrence} after substituting the values of the
\(d_j\).  Since \(\partial_L-n\) is bijective on \(\K[L]\) for \(n\ne0\)
(\cref{plt:cor:alg-antiderivative}), \(q_n\) is determined; and
\(\partial_Lq_0=-1\) has solution set \(-L+\K\) in characteristic zero.

(iii)  For \(c\in\K\) the map \(F_c=X+L-c\) lies in \(\Gnear\), so
\(\Omega_c=\rev{F_c}\) exists and is unique
(\cref{plt:thm:rev-existence}), and \(F_c\compose\Omega_c=X\) reads
\(\Omega_c+\log\Omega_c=X+c\); by (i) its correction solves
\eqref{plt:eq:cert-omega-ode}, and its block of index \(0\) satisfies
\(q_0+L=c\), so \(q_0(0)=c\).  Conversely a solution of
\eqref{plt:eq:cert-omega-ode} has \(q_0=-L+c\) for some \(c\in\K\) by (ii),
and then agrees with the correction of \(\Omega_c\) block by block, again by
the uniqueness in (ii).  Finally \(\Omega_c+\log\Omega_c-X=c\), which is the
block of index \(0\) of the residual.
\end{proof}

\begin{remark}
\Cref{plt:thm:lambert-recurrence} records the recurrence for the Lambert
polynomials in its intrinsic two-term form, which is the right form for
proving things about them.  \Cref{plt:lem:cert-omega-ode} is the form an
implementation can run without knowing the family in advance: it is a
recurrence for the \emph{blocks of a series}, driven by \(\DX\) and one
truncated product, and it therefore tests exactly the primitives of
\cref{plt:sec:alg-derivation,plt:sec:alg-multiplication}.  Its resonance at
\(n=0\) is not an artefact of the derivation: it is
\cref{plt:cor:alg-antiderivative} in the smallest instance anybody actually
computes.
\end{remark}

\begin{example}[Hand-checkable test vectors]
\label{plt:ex:cert-vectors}
Every quantity below can be verified with pencil and paper, which is the point
of a test vector.
\begin{enumerate}
\item \emph{Lambert blocks.}
      \(\LamP{0}=-L\), \(\LamP{1}=L\), \(\LamP{2}=\tfrac12L^{2}-L\),
      \(\LamP{3}=\tfrac13L^{3}-\tfrac32L^{2}+L\).  At \(n=3\) the recurrence
      \eqref{plt:eq:cert-omega-recurrence} reads
      \[
       (\partial_L-3)\LamP{3}=-L^{3}+\tfrac{11}2L^{2}-6L+1,
      \]
      and its right-hand side,
      \[
       -(\partial_L-2)\LamP{2}
        -\bigl[\LamP{2}\,\partial_L\LamP{0}
              +\LamP{1}(\partial_L-1)\LamP{1}
              +\LamP{0}(\partial_L-2)\LamP{2}\bigr],
      \]
      is the same polynomial; both are computed from the displayed
      \(\LamP{n}\) in four lines.
\item \emph{Composition through the logarithmic generator.}  With \(H=1\), the
      Taylor engine gives \(\DX L=t\), \(\DX^{2}L=-t^{2}\),
      \(\DX^{3}L=2t^{3}\), hence
      \[
       L\compose(X+1)=L+t-\tfrac12t^{2}+\tfrac13t^{3}+\FormalBigOAt{t^{4}}{t}
       =L+\log(1+t)+\FormalBigOAt{t^{4}}{t},
      \]
      in agreement with the generator formula
      \cref{plt:prop:cmp-generator-substitution}.  An engine that returns
      \(L\) has substituted into the wrong generator.
\item \emph{Guard block.}  \(F=X\) exactly, and \(G\) known to cutoff \(1\),
      its only stored block being the constant \(1\) at index \(0\): the sound
      output cutoff for \(FG\) is
      \(\min\SetOf{+\infty,\,1-1}=0\), so a routine that returns the block of
      index \(0\) of \(FG\) has returned an undetermined block
      (\cref{plt:ex:alg-guard-block}).
\item \emph{Nonunit denominator.}
      \(\MultiplicativeInverseOf{(L+t)}=\sum_{n\ge0}(-1)^{n}t^{n}L^{-n-1}\),
      an element of \(\PLrat\) and not of \(\PLlau\); the identity is checked
      by multiplying out and telescoping.  A routine that returns a
      polynomial answer here has violated (P4) of
      \cref{plt:def:alg-invariant}.
\item \emph{Resonance.}  \(\DX Y=t\) has solution set \(L+\K\); a routine that
      returns \(L\) without recording the choice is nondeterministic in the
      sense of \cref{plt:def:cert-determinism}.
\item \emph{Sampling is not a zero test.}  The residual block
      \(\prod_{i=1}^{s}(L-\lambda_i)\) vanishes at every sample point
      \(\lambda_1,\dots,\lambda_s\) and is nonzero
      (\cref{plt:prop:cert-exact-zero}).
\end{enumerate}
\end{example}

\begin{longtable}{@{}p{0.05\textwidth} p{0.42\textwidth} p{0.45\textwidth}@{}}
\caption{The conformance suite.  Each test is a checked identity; the third
column names the error class it is there to catch, with the diagnosis in
\cref{plt:sec:failure-modes}.}
\label{plt:tab:cert-suite}\\
\toprule
& Checked identity & Catches\\
\midrule
\endfirsthead
\toprule
& Checked identity & Catches\\
\midrule
\endhead
T1 & Declared cutoffs of \(+,\times,/,\compose,\DX\) equal those of
     \cref{plt:thm:alg-precision-calculus} on random inputs, including inputs
     of negative order &
     the ``modulo \(t^{N}\) in, modulo \(t^{N}\) out'' fallacy;
     \cref{plt:sec:fail-undetermined-block}\\
T2 & \(\Trunc{F}{N}\) retains the block of index \(N-1\) and discards the
     block of index \(N\) &
     inclusive/exclusive confusion; \cref{plt:sec:fail-offbyone}\\
% ed.: T3 had the gauge subtracted where it must be added.  By
% ed.: plt:prop:cert-reciprocal, ordt(BQ-A)=ordt(Q-A/B)+ordt(B); a quotient
% ed.: sound at cutoff N has ordt(Q-A/B)>=N, so its residual vanishes below
% ed.: N+ordt(B).  The printed N-ordt(B) is the threshold a WRONG answer
% ed.: already passes when ordt(B)>0, so the test as printed was vacuous in
% ed.: exactly the case the remark plt:rmk:cert-division-gauge is about.
T3 & \(BQ-A\) vanishes below \(N+\ordt(B)\) after a quotient at cutoff \(N\)
     &
     a gauge dropped from the division test;
     \cref{plt:rmk:cert-division-gauge}\\
T4 & \(\exp\log(1+U)=1+U\) and \(\log\exp U=U\) to the declared cutoff &
     off-by-one in the two recurrences of
     \cref{plt:prop:alg-explog-recurrences}\\
T5 & \(\DX(F\compose G)=((\DX F)\compose G)\,\DX G\) on the single input
     \(F=L\) &
     substitution into the wrong generator;
     \cref{plt:prop:cert-chain-rule-sharp}\\
T6 & Taylor engine and generator engine agree to the declared cutoff &
     an error in either engine, but not in a wrong convention shared by both\\
T7 & \(H\compose\rev{G}=F\) to the declared cutoff &
     composition errors that the consistency tests miss;
     \cref{plt:prop:cert-composition}\\
T8 & \emph{Both} \(F\compose G-X\) and \(G\compose F-X\) vanish below
     \(\min\SetOf{N_A,N_B}\), and their leading blocks agree &
     a reversion certified by the residual that drove it;
     \cref{plt:prop:cert-driven-residual}\\
T9 & The four reversion strategies of \cref{plt:sec:alg-reversion} agree to a
     common cutoff &
     strategy-specific indexing errors, in particular the Lagrange stopping
     rule; \cref{plt:sec:fail-offbyone}\\
T10& \(\WrightOmega_N+\log\WrightOmega_N-X\) has order exactly \(N+1\) and
     leading block \(-\LamP{N+1}\) &
     any error in the Lambert blocks, and a truncation off by one, which shows
     as order \(N\) or \(N+2\); \cref{plt:thm:om-residual}\\
T11& The blocks produced by \eqref{plt:eq:cert-omega-recurrence}, by
     \eqref{plt:eq:lambert-int-rec} and by
     \cref{plt:thm:lambert-stirling} coincide &
     an error in \(\DX\), in the truncated product, or in the Stirling
     coefficients; \cref{plt:lem:cert-omega-ode}\\
T12& \(\DX Y-G=0\) \emph{and} the recorded resonant constant is reproduced &
     a lost constant of integration; \cref{plt:sec:fail-resonance}\\
T13& Residual blocks are tested by canonical form, and a planted residual
     \(\prod_i(L-\lambda_i)\) is rejected &
     numerical sampling used as a zero test;
     \cref{plt:prop:cert-exact-zero}\\
T14& Two runs on the same input, and runs with permuted input orderings,
     produce byte-identical canonical forms &
     nondeterminism; \cref{plt:sec:cert-determinism}\\
\bottomrule
\end{longtable}

\begin{remark}[What a passing suite does not establish]
\label{plt:rmk:cert-suite-limits}
Every test above is a formal identity in the coefficient ring.  A complete
pass establishes that the computed truncations satisfy the defining equations
below their declared cutoffs.  It establishes nothing about any function:
the passage to an analytic statement needs a realization theorem
(\cref{plt:def:cmp-realization,plt:thm:cmp-analytic-transfer}), a regime, and
a derivative bound
(\cref{plt:prop:rev-analytic-transfer,plt:thm:om-analytic}).  Nor does a pass
establish that the answer is what the user wanted: the class, the branch and
the coefficient ring are inputs, and a computation performed in the wrong
class can be internally impeccable.  \Cref{plt:sec:alg-decision} exists for
that reason.
\end{remark}

\section{Determinism and reproducibility}
\label{plt:sec:cert-determinism}

\begin{definition}[Deterministic routine, reproducible computation]
\label{plt:def:cert-determinism}
A routine is \emph{deterministic} when its returned representative is a
function of its input representatives together with the declared branch and
normalisation data --- and, in particular, is independent of the order in
which stored blocks are visited, of any hashing, and of any internally chosen
constant.  A computation is \emph{reproducible} when its record
\[
 \bigl(\text{series},\ R,\ \text{cutoff},\ \text{order bound},\
 \text{residual and its order},\ \text{branch and normalisation data}\bigr)
\]
is returned in full, so that a second run, or a different implementation, can
be compared with it item by item.
\end{definition}

\begin{proposition}[Soundness pins the answer]
\label{plt:prop:cert-uniqueness}
Let \(\varphi\) be single-valued on denotations, and let \(\mathsf R\) be
sound for \(\varphi\) with declared cutoff \(N_{\mathrm{out}}\), returning
\(w\).  Then for every admissible input tuple \((F_1,\dots,F_k)\) in the
denotations,
\[
 w=\Trunc{\varphi(F_1,\dots,F_k)}{N_{\mathrm{out}}} .
\]
Consequently any two sound routines for \(\varphi\) that declare the same
cutoff return the same value, and if in addition both are canonical in the
sense of (P2) of \cref{plt:def:alg-invariant}, they return the same stored
object.
\end{proposition}

\begin{proof}
Write \(z=(F_1,\dots,F_k)\).  Soundness says
\(\varphi(z)\in\operatorname{Den}(w,N_{\mathrm{out}})\), which by
\cref{plt:lem:alg-denotation-coset} says
\(\Trunc{\varphi(z)}{N_{\mathrm{out}}}=w\).  The right-hand side
does not depend on the routine, and by hypothesis it does not depend on which
element of the denotations was chosen either.  Canonicity makes the stored
representation of \(w\) unique.
\end{proof}

\begin{corollary}[Reproducibility is a theorem, and its failure is a
diagnosis]
\label{plt:cor:cert-reproducible}
A conforming implementation --- sound, canonical, with a declared cutoff that
is a function of the declared input data --- is deterministic on every
operation whose \(\varphi\) is single-valued on denotations.  Consequently, if
two runs of a conforming implementation disagree, then at least one of the
following holds, and these are the only possibilities:
\begin{enumerate}
\item the operation is not single-valued --- an antiderivative
      (\cref{plt:cor:alg-antiderivative}), a root or a fractional power
      (\cref{plt:prop:cert-power}), a logarithm of a constant, or a branch
      choice --- and the choice was made internally instead of being an input;
\item the declared cutoff is not a function of the declared input data, for
      instance because it was taken from a stale stored order bound
      (\cref{plt:rmk:alg-derived-order});
\item the returned object is not in canonical form, for instance because a
      zero block was retained or a rational coefficient was left
      unnormalised;
\item the implementation is unsound.
\end{enumerate}
\end{corollary}

\begin{proof}
The positive statement is \cref{plt:prop:cert-uniqueness}: with (P1), (P2) and
a cutoff determined by the inputs, the returned object is
\(\Trunc{\varphi(\cdot)}{N_{\mathrm{out}}}\) in canonical form, hence the same
in both runs.  For the list, negate the hypotheses one at a time: the proof
used exactly single-valuedness of \(\varphi\) on the denotations, the fact
that \(N_{\mathrm{out}}\) is the same in both runs, canonicity, and
soundness.  A disagreement therefore falsifies one of them, and each of the
four items is the failure of one.
\end{proof}

\begin{warningbox}
Three habits break \cref{plt:cor:cert-reproducible} in practice and are worth
naming.  Iterating over a hash map makes the accumulation order depend on the
hash seed; with exact arithmetic the mathematical value is unaffected but the
canonical form is reached by a different route and any timing-dependent
truncation heuristic becomes visible, which is why the ordered map of
\cref{plt:def:alg-representation} is part of the contract and not an
optimisation.  Floating-point coefficients break the exact zero test of
\cref{plt:prop:cert-exact-zero}, hence break canonicity, hence break
determinism.  And a global ``simplification'' pass applied to coefficients
outside the declared \(R\) can silently change the type, which is failure
(P4).
\end{warningbox}

\begin{summarybox}
A certificate is a defining equation plus a checked valuation, and it
certifies a coset up to the equation's gauge.  Every primitive of
% ed.: same overclaim as in the opening list: the table has no entry for the
% ed.: sum, the product or the derivation.
\cref{plt:sec:algorithms} that inverts a cheaper primitive has a faithful
equation with a known gauge, except the antiderivative, whose equation is
blind to a constant --- which is exactly the resonance; addition,
multiplication and the derivation have none, and are certified by
recomputation.  Both reversion residuals certify the same blocks; the reason
to compute both is that a residual which drove the iteration certifies
nothing.  Zero tests are comparisons of canonical forms, never numerical
samples.  Reproducibility is then not a matter of discipline but a corollary
of soundness, canonicity, and a cutoff that is a function of the inputs.
\end{summarybox}

\chapter{Failure modes, diagnostics, and a decision procedure}
\label{plt:sec:failure-modes}

Polynomial--logarithmic computations are forgiving at leading order and
unforgiving after it.  Each of the eight failure modes below is one that the
six source drafts either commit, warn against without a diagnosis, or warn
against with the wrong diagnosis.  Each is presented in the same shape:
\emph{symptom}, the observable behaviour; \emph{diagnosis}, the theorem that
was violated; \emph{minimal reproducing example}, the smallest instance in
which the error is visible, chosen so that it can be checked by hand and
installed as a regression test; and \emph{repair}, what the contract must say
so that the error cannot recur silently.  The section closes with a decision
procedure: given a problem, which class, which algorithm, which certificate.

\section{Truncation-convention off-by-one}
\label{plt:sec:fail-offbyone}

\emph{Symptom.}  The last requested block is missing, or is present but wrong;
two implementations agree in every block but the last; a residual has order one
less than predicted, or one more; a loop that should run to \(R\) runs to
\(R-1\) and the answer is correct except at the most expensive block, the one
nobody checks by hand.

\emph{Diagnosis.}  Two conventions are in circulation and they differ by one
block:
\[
 \Trunc{F}{N}=\sum_{n<N}p_n(L)t^{n}
 \qquad\text{against}\qquad
 [F]_{\le N}=\Trunc{F}{N+1},
\]
the first adopted here and the second used by source~6.  The conversion is
harmless in a statement and fatal in a loop bound, because a loop bound is
where the two conventions meet an index.  The three places in this volume
where the discrepancy is real are the Taylor budget
(\cref{plt:prop:alg-taylor-budget}), the Lagrange stopping rule
(\cref{plt:alg:reverse-lagrange}), and the certifying residual, which for
correctness \emph{through} block \(N\) must be a
\(\FormalBigOAt{t^{N+1}}{t}\) and not a \(\FormalBigOAt{t^{N}}{t}\)
(\cref{plt:cor:rev-finite-certificate}).

\emph{Minimal reproducing example.}  Reverting \(F=X+L\), whose correction has
\(h=\ordt(L)=0\), by Lagrange--B\"urmann
(\cref{plt:alg:reverse-lagrange}).  Each summand contributes exactly one
block: for \(r\ge1\),
\[
 \frac{(-1)^{r}}{r!}\DX^{\,r-1}(L^{r})=\LamP{r-1}(L)\,t^{\,r-1},
\]
which is \cref{plt:eq:om-one-block}.  Asking for the blocks of index \(0,1,2\) ---
that is, for the exclusive cutoff \(N=3\) --- therefore requires \(r=1,2,3\),
which is \(R=\Ceiling{(N+1)/(h+1)}-1=3\).  A loop written to the inclusive
reading of ``through block \(N\)'' but run against an exclusive cutoff stops at
\(r=2\) and returns
\[
 -L+Lt+0\cdot t^{2}
 \qquad\text{instead of}\qquad
 -L+Lt+\Bigl(\tfrac12L^{2}-L\Bigr)t^{2},
\]
losing exactly \(\LamP{2}\).  The right residual then has order \(2\) instead
of the predicted \(3\), which is the diagnosis printed in one line.

\emph{Repair.}  Carry the cutoff in the type, never in the prose; convert once,
at the boundary, with \([F]_{\le N}=\Trunc{F}{N+1}\); and let the residual
order be the arbiter, since by
\cref{plt:thm:cert-two-sided}\ref{plt:thm:cert-two-sided:orders} it is exactly
the first index at which the answer is wrong.

\section{Reading a block the input data does not determine}
\label{plt:sec:fail-undetermined-block}

\emph{Symptom.}  A routine returns a block that changes when an input block
\emph{beyond} the requested cutoff is changed; two runs with different but
equally admissible completions of the same input disagree; a computation is
reproducible on every test case in which the inputs happen to have order
\(\ge0\) and wrong on the first Laurent input.

\emph{Diagnosis.}  The folklore rule ``both inputs modulo \(t^{N}\), output
modulo \(t^{N}\)'' is the product law
\cref{plt:thm:alg-precision-calculus}\ref{plt:thm:alg-precision-calculus:prod}
specialised to \(m_1=m_2=0\).  The law is
\(N_{\mathrm{out}}=\min\SetOf{N_1+m_2,N_2+m_1}\), which is \emph{better} than
the folklore rule when a factor has positive order and \emph{worse} as soon as
a factor has a pole.  Worse means: the routine returns a block it has no right
to return, and the failure is silent, because the returned block is a
perfectly good function of the stored data --- it is just not a function of
the mathematical object.

\emph{Minimal reproducing example.}  \(F=X=t^{-1}\), known exactly; \(G\)
known to exclusive cutoff \(1\), its only stored block being \(1\) at index
\(0\).  The sound cutoff for \(FG\) is
\(\min\SetOf{+\infty,\,1+(-1)}=0\): only the block of index \(-1\) is
determined.  Indeed \(G_1=1\) and \(G_2=1+t\) have the same representative and
\(FG_1=t^{-1}\), \(FG_2=t^{-1}+1\) differ at index \(0\).  A routine that
declares cutoff \(1\) here has invented the block \(0\).

\emph{Repair.}  Compute \(N_{\mathrm{out}}\) from the derived order bounds
(\cref{plt:rmk:alg-derived-order}), never from the input cutoffs alone;
report the relative precision \(\rho\) alongside the cutoff, since by
\cref{plt:cor:alg-relative-precision} it is \(\rho\), not \(N\), that the
multiplicative operations conserve; and run test~T1 of
\cref{plt:tab:cert-suite} with at least one input of negative order, which is
the only kind of input that exhibits the bug.

\section{Substituting into the wrong generator}
\label{plt:sec:fail-generator}

\emph{Symptom.}  The leading block is right and every later block is wrong, so
the answer looks nearly correct and passes any eyeball test; the residual has
order exactly one more than the leading block; ring-based consistency tests
all pass.

\emph{Diagnosis.}  Composition moves \emph{both} generators
(\cref{plt:prop:cmp-generator-substitution}):
\[
 t\compose G=\MultiplicativeInverseOf{G},
 \qquad
 L\compose G=\log G .
\]
Only the second is easy to forget, and forgetting it is not an approximation
but a different operation --- one that is still a \(\K[L]\)-algebra
homomorphism, hence still passes every multiplicative test
(\cref{plt:prop:cert-chain-rule-sharp}(i)).  A related error is to
``substitute \(t\mapsto t+H\)'', which is not an operation on \(\PLlau\) at
all: it is the substitution appropriate to a series in a variable that is not
the reciprocal of the argument.

\emph{Minimal reproducing example.}  Two, at different depths.
\begin{enumerate}
\item One generator, one block.  \(L\compose(X+1)=L+\log(1+t)
      =L+t-\tfrac12t^{2}+\tfrac13t^{3}+\FormalBigOAt{t^{4}}{t}\)
      (\cref{plt:ex:cert-vectors}(ii)).  An engine that returns \(L\) has
      forgotten the logarithmic generator.  The chain-rule test on this single
      input detects it (\cref{plt:prop:cert-chain-rule-sharp}(ii)).
\item A whole reversion.  The engine \(\Psi_G\) of
      \cref{plt:ex:cert-wrong-engine} reverts \(X+L\) to \(X-L\), which is the
      correct block of index \(0\) of \(\WrightOmega\) and nothing else, while
      its own residual vanishes identically.
\end{enumerate}

\emph{Repair.}  Implement composition only through the generator images
(\cref{plt:alg:compose-generators}) or the Taylor engine
(\cref{plt:alg:compose-taylor}), never by textual substitution; run the
chain-rule test T5 and the faithful certificate T7 of
\cref{plt:tab:cert-suite}; and never certify a reversion with the residual
that drove it (\cref{plt:prop:cert-driven-residual}).

\section{Treating a formal tail as an analytic bound}
\label{plt:sec:fail-formal-tail}

\emph{Symptom.}  A numerical error bound of the form \(CX^{-N}\) that
experiment violates; a claim that a truncation ``converges''; a remainder
quoted as \(\BigO(X^{-3})\) that in fact contains \(X^{-3}(\log X)^{100}\); an
identity between formal objects presented as an identity between functions.

\emph{Diagnosis.}  \(\FormalBigOAt{t^{N}}{t}\) is a statement about
coefficient support and nothing else
(\cref{plt:def:trunc-formal-o,plt:prop:trunc-formal-o}).  It becomes a
statement about a function only through a realization
(\cref{plt:def:cmp-realization}) together with a theorem that transports it
(\cref{plt:thm:cmp-analytic-transfer,plt:prop:rev-analytic-transfer,plt:thm:om-analytic}).
Two further confusions travel with this one: formal summability
(\cref{plt:def:trunc-summable}) is a statement about valuations and implies
nothing about numerical convergence; and a bare \(\BigO\) with no printed
regime silently invites the reader to supply the wrong one.

\begin{proposition}[Two ways in which the formal tail says nothing analytic]
\label{plt:prop:fail-formal-tail}
\begin{enumerate}
\item \emph{Realizations are not unique.}  Let \(f\) be a function on a ray
      \((X_0,+\infty)\) with \(f\approx F\) in the sense of
      \cref{plt:def:cmp-realization}, and let \(g(X)=f(X)+\EulerE^{-X}\).
      Then \(g\approx F\) as well, and \(f\ne g\) everywhere.  Consequently no
      formal statement about \(F\) determines any value of a realization, and
      the difference of two realizations is invisible to every block.
\item \emph{Formal summability is not convergence.}  The family
      \(\bigl(n!\,t^{n}\bigr)_{n\ge0}\) is formally summable, with sum
      \(F=\sum_{n\ge0}n!\,t^{n}\in\PLpow\), yet for every real \(X>0\) the
      numerical series \(\sum_{n\ge0}n!\,X^{-n}\) diverges.
\end{enumerate}
\end{proposition}

\begin{proof}
(i)  For every fixed \(N\), \(\EulerE^{-X}=\LittleOAt{X^{-N}}{X\to+\infty}\),
since \(X^{N}\EulerE^{-X}\to0\); hence \(g\) has the same asymptotic
remainders as \(f\) at every order, which is the definition of
\(g\approx F\).  And \(\EulerE^{-X}>0\), so \(f\ne g\) pointwise.

(ii)  Formal summability requires \(\ordt(n!\,t^{n})=n\to\infty\)
(\cref{plt:def:trunc-summable}), which holds, so the sum exists in
\(\PLpow\); each block is a single scalar and no coefficient receives
infinitely many contributions.  For the divergence, fix \(X>0\); then
\(n!\,X^{-n}\to\infty\) as \(n\to\infty\), because
\(\bigl((n+1)!X^{-n-1}\bigr)/\bigl(n!X^{-n}\bigr)=(n+1)/X\ge2\) for
\(n\ge2X\).  A series whose terms do not tend to zero diverges.
\end{proof}

\emph{Minimal reproducing example.}  \Cref{plt:prop:fail-formal-tail}(ii) with
\(N=3\): the ``remainder after three blocks'' is formally
\(\FormalBigOAt{t^{3}}{t}\), while numerically, at \(X=10\), the terms
\(n!\,10^{-n}\) decrease only until \(n=9\) and \(n=10\), where both equal
\(9!/10^{9}=3.63\times10^{-4}\), and grow without bound afterwards.  The
formal tail marker predicts none of this and forbids none of it: it is not a
statement about numbers.

\emph{Repair.}  Never write a bare \(\BigO\) for an asymptotic claim.  Print
the regime, as in
\[
 \BigOAt{(\log X)^{d}X^{-3}}{X\to+\infty};
\]
state which realization theorem is being used, and if none is, say that the
statement is formal.  \Cref{plt:sec:om-analytic} is the model of the four ingredients an
analytic statement needs.

\section{Dividing by a nonunit}
\label{plt:sec:fail-nonunit}

\emph{Symptom.}  Coefficients acquire denominators in \(L\); block degrees go
negative; a routine that promised polynomial coefficients returns something
else, or --- worse --- silently truncates an infinite series in \(L^{-1}\) and
returns a polynomial that is wrong in every block; a residual test fails at
the leading block for no visible reason.

\emph{Diagnosis.}  The units of \(\K[L]\) are \(\K^{\times}\)
(\cref{plt:lem:div-poly-units}); hence an element of \(\PLlau\) is a unit if
and only if its leading block is a nonzero \emph{scalar}
(\cref{plt:thm:div-unit-criterion}).  The reciprocal recurrence divides by the
leading block at every step, so a leading block of positive \(L\)-degree takes
the computation out of \(\PLlau\) and into \(\PLrat\), or after expansion into
\(\PLit\).  This is not a defect of the algorithm but information about the
problem, and (P4) of \cref{plt:def:alg-invariant} requires it to appear in the
returned type.

\begin{proposition}[The smallest nonunit, and the two rings it forces]
\label{plt:prop:fail-nonunit}
\begin{enumerate}
\item \(B=L+t\) is not a unit of \(\PLlau\).
\item In \(\PLrat=\K(L)((t))\) it is a unit, with
      \(\MultiplicativeInverseOf{B}=\sum_{n\ge0}(-1)^{n}L^{-n-1}t^{n}\).
\item For \(B'=L+1+t\) the reciprocal is
      \(\sum_{n\ge0}(-1)^{n}(L+1)^{-n-1}t^{n}\), and each of its blocks, when
      expanded in \(\K((L^{-1}))\), is an infinite series; so the passage from
      \(\PLrat\) to \(\PLit\) introduces a second, inner cutoff, which the
      type must carry and the certificate must quote.
\end{enumerate}
\end{proposition}

\begin{proof}
(i)  The leading block of \(B\) is \(L\), which is not a unit of \(\K[L]\)
because a product \(Lq(L)\) has degree \(\ge1\); now apply
\cref{plt:thm:div-unit-criterion}.

(ii)  Multiply out and telescope:
\begin{align*}
 (L+t)\sum_{n\ge0}(-1)^{n}L^{-n-1}t^{n}
 &=\sum_{n\ge0}(-1)^{n}L^{-n}t^{n}
  +\sum_{n\ge0}(-1)^{n}L^{-n-1}t^{n+1}\\
 &=\sum_{n\ge0}(-1)^{n}L^{-n}t^{n}-\sum_{m\ge1}(-1)^{m}L^{-m}t^{m}=1 .
\end{align*}
Every block lies in \(\K(L)\) and the family is summable, the \(n\)-th term
having \(\ordt=n\).

(iii)  The same telescoping with \(L\) replaced by \(L+1\) gives the displayed
reciprocal.  Expanding \((L+1)^{-n-1}=L^{-n-1}(1+L^{-1})^{-n-1}
=\sum_{k\ge0}\binom{-n-1}{k}L^{-n-1-k}\) exhibits infinitely many nonzero
inner coefficients, since \(\binom{-n-1}{k}=(-1)^{k}\binom{n+k}{k}\ne0\) for
every \(k\).
\end{proof}

\emph{Minimal reproducing example.}  \Cref{plt:prop:fail-nonunit} with the
quotient \(1/(L+t)\) requested in \(\PLlau\): a conforming routine refuses and
reports the obstruction; a nonconforming one returns
\(L^{-1}-L^{-2}t+\cdots\) typed as a polynomial-coefficient object, and every
downstream degree computation is then meaningless.

\emph{Repair.}  Test the leading block for unit-hood before dividing, taking
logarithms, or forming a general power; enlarge \(R\) explicitly and record it
in the returned type; and remember that in \(\PLit\) the log-degree profile of
\cref{plt:def:mul-profile} is no longer the relevant invariant --- the inner
valuation is.

\section{Losing the constant of integration at the resonant block}
\label{plt:sec:fail-resonance}

\emph{Symptom.}  Two runs differ by an additive constant; an identity that
held before an integration fails after it; a derived expansion is a shift of
the intended one, so that every block of positive index is right and the
answer is still wrong; or, in the rational class, a routine returns nonsense
where the true answer needs \(\log L\).

\emph{Diagnosis.}  The block equation for \(\DX Y=G\) is
\((\partial_L-n)y_n=g_{n+1}\).  The operator \(\partial_L-n\) is bijective on
\(\K[L]\) for \(n\ne0\) and equals \(\partial_L\) for \(n=0\), whose kernel is
\(\K\) (\cref{plt:cor:alg-antiderivative}).  So the block of index \(0\) --- the
resonant block --- is determined only up to a constant, and the defining
equation \(\DX Y-G\) is blind to that constant
(\cref{plt:prop:cert-antiderivative}).  In \(\PLrat\) the same block carries an
obstruction as well as a freedom: the equation \(\partial_Ly_0=g_1\) is
solvable in \(\K(L)\) only when \(g_1\) has zero residue
(\cref{plt:lem:dif-res-derivative,plt:thm:dif-anti-rat}).

\emph{Minimal reproducing example.}  Two, one for the freedom and one for the
obstruction.
\begin{enumerate}
\item \emph{Freedom.}  Recover the \(\WrightOmega\) expansion from the
      differentiated equation \eqref{plt:eq:cert-omega-ode} instead of from
      \(\Omega+\log\Omega=X\).  By
      \cref{plt:lem:cert-omega-ode}(ii)--(iii) the differentiated equation
      determines every block of index \(\ge1\) uniquely and determines the
      block of index \(0\) only up to a constant \(c\), and each \(c\) gives
      the genuine reversion of the genuinely different map \(X+L-c\).  Every
      block of positive index is then wrong as well, since it is a block of
      the wrong map: solving
      \eqref{plt:eq:cert-omega-recurrence} at \(n=1\) with \(q_0=c-L\) gives
      \((\partial_L-1)q_1=1+c-L\), hence \(q_1=L-c\).  For \(c=1\) the
      computed answer begins \(1-L+(L-1)t\) where the intended one begins
      \(-L+Lt\); no block is correct, and the residual
      \(\Omega+\log\Omega-X=c\) exposes it in the block of index \(0\).
\item \emph{Obstruction.}  Antidifferentiate \(G=t/L\in\PLrat\).  The block
      equation at \(n=0\) is \(\partial_Ly_0=1/L\), which has no solution in
      \(\K(L)\), because the derivative of a rational function has zero
      residue at \(0\) (\cref{plt:lem:dif-res-derivative}) while \(1/L\) has
      residue \(1\).  The class must be enlarged to admit \(\log L\), which is
      \cref{plt:cor:dif-loglog}.
\end{enumerate}

\emph{Repair.}  Treat the antiderivative as an operation with an input, not as
a function: the resonant constant is part of the specification, it is returned
in the record of \cref{plt:def:cert-determinism}, and the certificate is the
pair (T12 of \cref{plt:tab:cert-suite}).  Before integrating in \(\PLrat\),
test the residue.

\section{Assuming a bound uniform in \texorpdfstring{\(L\)}{L}}
\label{plt:sec:fail-uniform-L}

\emph{Symptom.}  A truncation that is provably correct as \(X\to+\infty\)
behaves badly at every argument one can actually compute with; the ``next
correction'' is larger than the term it corrects; an error constant that was
quoted once turns out to depend on the number of blocks retained, or on the
log-degrees of the discarded ones.

\emph{Diagnosis.}  The dominance \(t^{n+1}L^{D}\precdom t^{n}\) is a statement
about \(X\to+\infty\) at \emph{fixed} \(D\)
(\cref{plt:lem:mot-dominance,plt:lem:ring-dominance-analytic}).  Since
% ed.: the crossover was quoted without the hypothesis D>=3 under which the
% ed.: root exists, and the asymptotic x^{*}(D) ~ D\log D was asserted; a
% ed.: one-line proof is supplied, since the volume states nothing else about
% ed.: this root.
\(L=\log X\) grows, the crossover beyond which the monomial \(L^{D}X^{-1}\)
stays small is, for \(D\ge3\), at \(\log X>x^{*}(D)\), where \(x^{*}(D)\) is
the largest root of \(x=D\log x\) (\cref{plt:prop:fail-no-uniform-constant}(iv);
for \(D\le2\) there is no crossover).  Moreover
\(x^{*}(D)\sim D\log D\) as \(D\to\infty\).  Indeed, writing
\(y=\log x^{*}\), the equation \(x^{*}=D\log x^{*}\) reads
\(y-\log y=\log D\); the left side is strictly increasing to \(+\infty\) on
\(y>1\), and \(y>\log D\ge\log3>1\), so \(y\to\infty\) with \(D\), whence
\(\log y=\LittleOAt{y}{D\to\infty}\), \(y\sim\log D\), and
\(x^{*}=Dy\sim D\log D\).
A profile that grows with the block index therefore pushes the crossover out
with it, and no constant covers all blocks at once.

\begin{proposition}[No constant is uniform in the number of blocks]
\label{plt:prop:fail-no-uniform-constant}
Fix an integer \(D\ge1\), take \(\K=\RealNumbers\), and let
\(F=\sum_{n\ge0}L^{Dn}t^{n}\in\PLpow\), so that
\(\dprofile{F}(n)=Dn\).  All statements below about numbers are on the
positive real ray with the real logarithm.
\begin{enumerate}
\item On the ray \(X>1\) with \(L=\log X\), the numerical series converges
      exactly when \(L^{D}<X\), with sum \(f(X)=X/(X-L^{D})\), and
      \(f\approx F\) as \(X\to+\infty\).
\item For each fixed \(N\ge1\), the remainder is
      \(f(X)-\sum_{n<N}L^{Dn}X^{-n}
      =\bigl(L^{D}X^{-1}\bigr)^{N}\bigl(1-L^{D}X^{-1}\bigr)^{-1}\),
      which is \(\LittleOAt{X^{-N+\varepsilon}}{X\to+\infty}\) for every
      \(\varepsilon>0\) but is \emph{not} \(\BigOAt{X^{-N}}{X\to+\infty}\).
\item There is no pair \((C,X_0)\) with
      \(\AbsoluteValue{f(X)-\sum_{n<N}L^{Dn}X^{-n}}\le CX^{-N}\) for all
      \(N\ge0\) and all \(X\ge X_0\).
% ed.: three repairs.  (a) The clause was stated for every D>=1, but
% ed.: x=D\log x has NO real root for D=1,2 (the function x-D\log x attains
% ed.: its minimum D(1-\log D)>0 there), so x^{*}(D) does not exist and the
% ed.: "if and only if" cannot be read; for those D the blocks decrease at
% ed.: every X>1.  (b) Even for D>=3 the equivalence as printed was false near
% ed.: X=1: x-D\log x has TWO roots, and the blocks also decrease below the
% ed.: smaller one -- at X=2 and D=10, L^{10}=0.025<2.  What x^{*}(D) is, is
% ed.: the threshold beyond which they decrease for good.  (c) The printed
% ed.: factor 2.42... is not the truncated decimal expansion of
% ed.: L^{10}/X=2.41584..., only its rounding.
\item Suppose \(D\ge3\).  Then \(x=D\log x\) has exactly two roots
      \(x_{-}(D)<D<x^{*}(D)\), and at a given \(X>1\) the blocks decrease in
      size if and only if \(\log X<x_{-}(D)\) or \(\log X>x^{*}(D)\); thus
      \(x^{*}(D)\) is the threshold beyond which they decrease for good.  For
      \(D=10\), \(x_{-}=1.1183\ldots\) and \(x^{*}=35.7715\ldots\); so at
      \(X=10^{15}\) the block of index \(1\) exceeds the block of index \(0\)
      by the factor \(L^{10}/X=2.4158\ldots\), and the expansion becomes
      permanently decreasing only for \(X>3.4\times10^{15}\).  For \(D\le2\)
      there is no threshold at all: the blocks decrease at every \(X>1\).
\end{enumerate}
\end{proposition}

\begin{proof}
(i)  The series is geometric with ratio \(L^{D}X^{-1}\), whence the
convergence criterion and the sum \(1/(1-L^{D}X^{-1})=X/(X-L^{D})\).  For the
asymptotic statement, fix \(N\ge0\); by (ii) with \(N+1\) in place of \(N\),
\[
 f(X)-\sum_{n\le N}L^{Dn}X^{-n}
 =\frac{L^{D(N+1)}X^{-(N+1)}}{1-L^{D}X^{-1}}
 =\LittleOAt{X^{-N}}{X\to+\infty},
\]
because \(L^{D(N+1)}X^{-1}=(\log X)^{D(N+1)}/X\to0\) and the denominator tends
to \(1\).  This is the level-wise condition of
\cref{plt:def:cmp-realization}, so \(f\approx F\).

(ii)  Sum the geometric tail.  The remainder equals
\(L^{DN}X^{-N}\bigl(1-L^{D}X^{-1}\bigr)^{-1}\ge L^{DN}X^{-N}\) for \(X\)
large, and \(L^{DN}\to\infty\), so the remainder divided by \(X^{-N}\) is
unbounded as \(X\to+\infty\); hence it is not
\(\BigOAt{X^{-N}}{X\to+\infty}\).  The little-o statement is the estimate just
made.

% ed.: the proof chose only X >= max{X_0, e^2}, at which the series can still
% ed.: diverge (D=10, X=e^2 has L^{D}=1024>X) and f(X) is then undefined; the
% ed.: choice must also put X in the convergence range, which is possible
% ed.: because L^{D}/X -> 0.
(iii)  Since \(L^{D}X^{-1}=(\log X)^{D}/X\to0\), one may fix an
\(X\ge\max\SetOf{X_0,\EulerE^{2}}\) with \(L^{D}<X\), so that \(f(X)\) is
defined by (i); then \(L=\log X\ge2\) and in
particular \(L^{D}\ge2\).  The remainder at cutoff \(N\) is at least
\(L^{DN}X^{-N}\), so the asserted bound would force \(C\ge L^{DN}\ge2^{DN}\)
for every \(N\), which is impossible for a finite \(C\).

(iv)  The block of index \(n\) evaluates to \(\bigl(L^{D}X^{-1}\bigr)^{n}\),
so the blocks decrease if and only if \(L^{D}<X\), that is
\(D\log L<L\) with \(L=\log X\), that is \(\psi(L)>0\) for
\(\psi(x)=x-D\log x\).  On \((0,\infty)\) the function \(\psi\) is smooth with
\(\psi'(x)=1-D/x\), so its only critical point is \(x=D\), a minimum, with
\(\psi(D)=D(1-\log D)\), and \(\psi\to+\infty\) at both ends.
% ed.: the existence of the threshold was asserted for every D and the second
% ed.: root was overlooked; both are repaired here.
Hence: for \(D\le2\) one has \(\log D<1\), so \(\psi(D)=D(1-\log D)>0\) and
therefore \(\psi>0\) everywhere, and the blocks decrease at every \(X>1\); for
\(D\ge3\) one has \(\psi(D)<0\), so \(\psi\) has exactly one root
\(x_{-}(D)\in(0,D)\) and exactly one root \(x^{*}(D)\in(D,\infty)\), and
\(\psi(L)>0\) is equivalent to \(L<x_{-}(D)\) or \(L>x^{*}(D)\).  For
\(D=10\), \(x-10\log x\) is \(+0.147\) at \(x=1.1\) and \(-0.092\) at
\(x=1.13\), and \(-0.553\) at \(x=35\) and \(+0.165\) at \(x=36\), so
\(x_{-}=1.1183\ldots\), \(x^{*}=35.7715\ldots\) and
\(\EulerE^{x^{*}}=3.43\times10^{15}\).  At
\(X=10^{15}\) one has \(L=34.5388\ldots\) and \(L^{10}/X=2.4158\ldots\).
\end{proof}

\emph{Minimal reproducing example.}
\Cref{plt:prop:fail-no-uniform-constant}(iv) with \(D=10\) at
\(X=10^{15}\): a series that is a perfectly valid asymptotic expansion, whose
first correction is \(2.4\) times its leading term.

\emph{Repair.}  Quote the log-degrees with the error term, as in
\(\BigOAt{(\log X)^{d}X^{-N}}{X\to+\infty}\); state the constant's dependence
on \(N\) and on the profile; and, for numerical work, compare the actual
magnitude of the first discarded block at the actual argument rather than
trusting the ordering of the monomials.

\section{Using a degree bound that was a hypothesis}
\label{plt:sec:fail-degree}

\emph{Symptom.}  A fixed-width coefficient array overflows, or silently
truncates the top logarithmic coefficients; a cost estimate is off by a factor
that grows with \(N\); a proof of a degree law is quoted without its
hypothesis.

\emph{Diagnosis.}  Two distinct claims are easily confused.  The degree law
for a reversion (\cref{plt:thm:lag-degree-bound}) is a theorem \emph{under the
hypothesis} that the profile of the input is linearly bounded, and it bounds
the output profile in terms of that hypothesis.  What is not true, under any
hypothesis, is that the output profile is bounded by the input profile: a
single reciprocal already destroys that, and reversion destroys it
spectacularly.  The log-degree profile is output data
(\cref{plt:def:mul-profile,plt:prop:lead-profile-laws,plt:prop:div-degree-profile}),
to be measured and reported, never an input constant.

\begin{proposition}[Bounded input, unbounded output; and the exact scalar cost]
\label{plt:prop:fail-reversion-degree}
\begin{enumerate}
\item \emph{One reciprocal suffices.}  \(B=1+Lt\) has
      \(\dprofile{B}(0)=0\), \(\dprofile{B}(1)=1\) and no other nonzero block,
      while
      \(\MultiplicativeInverseOf{B}=\sum_{n\ge0}(-1)^{n}L^{n}t^{n}\) has
      \(\dprofile{\MultiplicativeInverseOf{B}}(n)=n\) for every \(n\).
\item \emph{Reversion.}  \(F=X+L\) has a single nonzero correction block, of
      degree \(1\), and \(\rev{F}=\WrightOmega\) has correction blocks
      \(\LamP{n}(L)\) with
      \[
       \dprofile{\rev{F}-X}(0)=1,
       \qquad
       \dprofile{\rev{F}-X}(n)=n\quad(n\ge1).
      \]
\item \emph{Cost.}  For \(N\ge1\), the schoolbook truncated product of two
      series with the profile of (ii) costs exactly
      \[
       \binom{N+3}{4}+N^{2}+N+1
      \]
      scalar multiplications, which exceeds the naive
      ``\(\binom{N+1}{2}\) multiplications'' count of
      \cref{plt:prop:alg-convolution-cost} by a factor
      \(\Theta(N^{2})\).
\end{enumerate}
\end{proposition}

\begin{proof}
(i)  The geometric series identity \((1+Lt)\sum_{n\ge0}(-1)^{n}L^{n}t^{n}=1\)
is immediate by telescoping, and \(\deg\bigl((-1)^{n}L^{n}\bigr)=n\).

(ii)  \(\rev{(X+L)}=\WrightOmega=X+\sum_{n\ge0}\LamP{n}(L)t^{n}\) is
\cref{plt:thm:om-expansion}.  For the degrees, \(\LamP{0}(u)=-u\) has degree
\(1\), and for \(n\ge1\),
\(\LamP{n}=\bigl((-1)^{n+1}/(n+1)!\bigr)\shiftop{0}{n}u^{n+1}\)
by \cref{plt:def:lambert-polynomials}, where
\(\shiftop{0}{n}=\prod_{j=0}^{n-1}(\partial_u-j)\) contains the factor
\(\partial_u\) exactly once, at \(j=0\); by
\cref{plt:prop:alg-iterated-derivative} the degree therefore drops by exactly
one, giving \(\deg\LamP{n}=(n+1)-1=n\).  This is the degree law of
\cref{plt:cor:lambert-coefficients}.

(iii)  Write \(\mathcal P=\rev{F}-X\).  By
\cref{plt:prop:alg-convolution-cost} the scalar count is
\(\sum_{i+j<N}w(i)w(j)\) with \(w(n)=\dprofile{\mathcal P}(n)+1\), so
\(w(0)=2\) and \(w(n)=n+1\) for \(n\ge1\); no block vanishes, since
\(\LamP{n}\ne0\) for all \(n\).  Using the Iverson bracket, which is \(1\)
when the enclosed statement holds and \(0\) otherwise, we have
\(w(i)=(i+1)+[i=0]\) and hence
\[
 w(i)w(j)=(i+1)(j+1)+(i+1)[j=0]+(j+1)[i=0]+[i=0][j=0].
\]
Summing over \(i,j\ge0\) with \(i+j<N\): the first term gives
\(\binom{N+3}{4}\) by
\cref{plt:cor:alg-cost-not-quadratic}(ii); the second gives
\(\sum_{i=0}^{N-1}(i+1)=N(N+1)/2\), and the third the same by symmetry; the
fourth gives \(1\).  The total is
\(\binom{N+3}{4}+N(N+1)+1=\binom{N+3}{4}+N^{2}+N+1\).  Since
\(\binom{N+3}{4}=N^{4}/24+\BigO(N^{3})\) while
\(\binom{N+1}{2}=N^{2}/2+\BigO(N)\), the ratio is \(\Theta(N^{2})\).
\end{proof}

\emph{Minimal reproducing example.}
\Cref{plt:prop:fail-reversion-degree}(i): an implementation that stores
coefficient polynomials in an array of width \(\max_n\dprofile{B}(n)+1=2\),
inferred from the input, loses every logarithmic coefficient above degree
\(1\) from the third block of the reciprocal on.

\emph{Repair.}  Store blocks as variable-length objects
(\cref{plt:def:alg-representation}); measure the profile of the output and
return it; quote costs in the coefficient-operation model and multiply by the
measured profile factor, as
\cref{plt:cor:alg-cost-not-quadratic} requires; and when invoking a degree law,
print its hypothesis.

\section{A decision procedure}
\label{plt:sec:alg-decision}

The eight failure modes have a common shape: a decision was made implicitly.
The procedure below makes each of them explicit and in the order in which the
information becomes available.  Its first three steps decide the class, its
next two the algorithm and the precision budget, its last three the
certificate and the status of the answer.  Nothing in it is optional: skipping
step~3 is failure mode \cref{plt:sec:fail-nonunit}, skipping step~4 is
\cref{plt:sec:fail-undetermined-block}, and skipping step~7 is
\cref{plt:sec:fail-formal-tail}.

\begin{algorithm}[Decision procedure: class, algorithm, certificate]
\label{plt:alg:decision}
\textbf{Input.}  A problem: an expansion, an implicit equation, an inverse, or
a composition; and a target amount of information.\\
\textbf{Steps.}
\begin{enumerate}
\item \emph{Name the regime and the branch.}  Which variable tends where, on
      which ray or sector, with which logarithm branch.  For a problem at
      \(x\to0^{+}\), transport it to \(X\to+\infty\) by the dictionary of
      \cref{plt:prop:grp-zero-dictionary} before anything else; for a problem
      at a finite critical point, stop --- the scale is ramified, not
      polynomial--logarithmic.
\item \emph{Solve the leading balance and normalise.}  Factor the dominant
      monomial.  If the map is monomial-leading,
      \(cX^{\varrho}(1+\ldots)\) with \(\varrho>0\)
      (\cref{plt:def:lag-monomial-leading}), reduce to the near-identity case
      by \cref{plt:thm:lag-reduction}; a leading logarithmic factor is first
      put in normal form by \cref{plt:thm:lag-log-normal-form}.  If the
      leading derivative degenerates, the problem is not near-identity
      reversion at all and needs a ramified scale.
\item \emph{Compute the generator images before anything else.}  Form
      \(t\compose G=\MultiplicativeInverseOf{G}\) and \(L\compose G=\log G\)
      and read the smallest closed class off
      \cref{plt:tab:decide-class}.  Adopt that class; if it is larger than the
      one the problem was posed in, say so in the output type.  This is the
      generator-first rule, and it is much safer than expanding first and
      inferring the class from the mess.
\item \emph{Fix the precision budget backwards.}  From the target cutoff,
      propagate the required input cutoffs through
      \cref{plt:thm:alg-precision-calculus}, remembering that the guard blocks
      are not a safety margin but part of the law
      (\cref{plt:ex:alg-guard-block}), and that relative precision, not
      absolute cutoff, is what the multiplicative operations conserve
      (\cref{plt:cor:alg-relative-precision}).
\item \emph{Choose the algorithm} from \cref{plt:tab:decide-algorithm},
      by how many blocks are wanted, whether they are wanted lazily, and
      whether composition is available and cheap.
\item \emph{Choose the certificate} from \cref{plt:tab:cert-gauges}, compute
      the residual, and compare its order with the declared cutoff
      \emph{plus the gauge}.  If the answer was produced by an iteration, the
      certificate must not be the residual that drove it
      (\cref{plt:prop:cert-driven-residual}).
\item \emph{Declare the status of the answer.}  Formal --- an identity of
      coefficients, and nothing more; asymptotic --- name the realization and
      the transfer theorem
      (\cref{plt:def:cmp-realization,plt:thm:cmp-analytic-transfer,plt:prop:rev-analytic-transfer,plt:thm:om-analytic});
      or convergent --- name the convergence theorem, and note that this
      volume proves none for the \(\WrightOmega\) series
      (\cref{plt:rmk:om-open}).
\item \emph{Return the record} of \cref{plt:def:cert-determinism}: series,
      ring, cutoff, order bound, residual and its order, branch and
      normalisation data.
\end{enumerate}
\end{algorithm}

\begin{longtable}{@{}p{0.42\textwidth} p{0.28\textwidth} p{0.24\textwidth}@{}}
\caption{Step~3 of \cref{plt:alg:decision}: what the generator images force.
The class named is the smallest one in this volume that is closed for the
operation; a smaller one can be used only by refusing the operation.}
\label{plt:tab:decide-class}\\
\toprule
Observed data & Smallest closed class & Authority\\
\midrule
\endfirsthead
\toprule
Observed data & Smallest closed class & Authority\\
\midrule
\endhead
Sums, products, derivatives of blocks with polynomial coefficients &
  \(\PLlau=\K[L]((t))\) &
  \cref{plt:thm:mul-closure,plt:thm:dif-derivation}\\
Inner map \(X+H\), \(H\in\PLpow\); reversion of it &
  \(\PLlau\), group \(\Gnear\) &
  \cref{plt:thm:grp-group,plt:thm:rev-existence}\\
Inner map \(cX^{\varrho}(1+U)\), \(\varrho\in\PositiveIntegers\),
  \(\log c\in\K\) & \(\PLlau\) &
  \cref{plt:prop:cmp-generator-substitution}\\
Division, logarithm or general power whose leading block is a nonconstant
  polynomial & \(\PLrat=\K(L)((t))\) &
  \cref{plt:thm:div-unit-criterion,plt:thm:div-minimal-field}\\
A rational coefficient expanded for large \(L\); a second, inner cutoff
  appears & \(\PLit=\K((L^{-1}))((t))\) &
  \cref{plt:thm:div-embedding,plt:prop:fail-nonunit}\\
\(\varrho\in\RationalNumbers_{>0}\setminus\IntegerNumbers\), or a formal root &
  Puiseux--log tower &
  \cref{plt:prop:ring-puiseux-tower,plt:prop:cmp-fractional-rho,plt:cor:mul-roots}\\
\(\varrho\) irrational, or a non-cyclic exponent group & Hahn field of
  \(\Gamma\) & \cref{plt:thm:ring-hahn}\\
Leading term \(cX^{a}L^{b}\) with \(b\ne0\); or an inner map with a
  logarithmic factor & depth two: \(\log\log X\) &
  \cref{plt:prop:lag-loglog-forced,plt:cor:lag-enlargement,plt:prop:grp-xlogx}\\
Composition with \(\log\), or with a map of logarithmic depth \(r\) &
  depth-\(r\) tower &
  \cref{plt:prop:grp-shift,plt:thm:grp-multilog-group}\\
An exponential of a block that is not a constant; composition with \(\exp\) &
  outside every logarithmic class &
  \cref{plt:thm:dif-exp-boundary,plt:prop:grp-exponential}\\
Antiderivative in \(\PLrat\) of a block with nonzero residue & adjoin
  \(\log L\) &
  \cref{plt:thm:dif-anti-rat,plt:cor:dif-loglog}\\
\bottomrule
\end{longtable}

\begin{longtable}{@{}p{0.30\textwidth} p{0.34\textwidth} p{0.30\textwidth}@{}}
\caption{Steps~5 and~6 of \cref{plt:alg:decision}: which algorithm, and which
certificate goes with it.  Costs are in the coefficient-operation model of
\cref{plt:def:alg-cost-models} and must be multiplied by the profile factor of
\cref{plt:cor:alg-cost-not-quadratic} to describe a machine.}
\label{plt:tab:decide-algorithm}\\
\toprule
Goal & Algorithm & Certificate\\
\midrule
\endfirsthead
\toprule
Goal & Algorithm & Certificate\\
\midrule
\endhead
Product at cutoff \(N\) & \cref{plt:alg:multiply} &
  Leibniz test, plus T1 of \cref{plt:tab:cert-suite}\\
Reciprocal, small \(N\) & \cref{plt:alg:reciprocal-linear},
  \(\binom{\rho}{2}\) products & \(BC-1\), gauge \(\ordt B\)\\
Reciprocal, large \(N\) & \cref{plt:alg:reciprocal-newton},
  \(\Ceiling{\log_2\rho}\) stages & same, recomputed at each stage\\
Quotient & \cref{plt:alg:divide} & \(BQ-A\), gauge \(\ordt B\)\\
Power, \(\alpha\) a small nonnegative integer & binary powering,
  \cref{plt:alg:power} & \(D^{q}-B^{p}\) with \(q=1\)\\
Power, general \(\alpha\); formal roots &
  recurrence \cref{plt:eq:alg-power-recurrence} &
  \(D^{q}-B^{p}\) and a check of the block of index \(0\)
  (\cref{plt:prop:cert-power})\\
\(\log(1+U)\), \(\exp U\), \(\ordt U\ge1\) & \cref{plt:alg:explog} &
  \cref{plt:prop:cert-explog}, either form\\
Composition, inner near-identity & \cref{plt:alg:compose-taylor} &
  T5, T6, and \cref{plt:prop:cert-composition}\\
Composition, inner monomial-leading & \cref{plt:alg:compose-generators} &
  the same three\\
Reversion, a few blocks, or lazily & \cref{plt:alg:reverse-triangular} &
  \emph{left} residual, never the driving one\\
Reversion, large \(N\), composition cheap & \cref{plt:alg:reverse-newton}
  with \cref{plt:cor:rev-newton-schedule} &
  both residuals, plus the abort test\\
Reversion, closed form or a map whose powers simplify &
  \cref{plt:alg:reverse-lagrange} & both residuals; for \(X+L\), T10\\
Reversion with only a composition primitive & \cref{plt:alg:reverse-picard} &
  both residuals\\
Antiderivative & \cref{plt:cor:alg-antiderivative}, blockwise &
  \(\DX Y-G\) with the constant declared (T12)\\
Mixed or unknown leading data & \cref{plt:alg:reverse-hybrid} &
  as at each stage\\
\bottomrule
\end{longtable}

\begin{longtable}{@{}p{0.56\textwidth} p{0.38\textwidth}@{}}
\caption{Diagnostic index: an observed symptom and the section that explains
it.}
\label{plt:tab:diagnostic-index}\\
\toprule
Symptom & Diagnosis\\
\midrule
\endfirsthead
\toprule
Symptom & Diagnosis\\
\midrule
\endhead
Exactly one block, the last, disagrees between two implementations &
  \cref{plt:sec:fail-offbyone}\\
A residual has order one less, or one more, than predicted &
  \cref{plt:sec:fail-offbyone}\\
An output block changes when an input block beyond the cutoff changes &
  \cref{plt:sec:fail-undetermined-block}\\
Everything is correct until the first input with a pole &
  \cref{plt:sec:fail-undetermined-block}\\
The leading block is right and every later block is wrong &
  \cref{plt:sec:fail-generator}\\
Multiplicative tests pass and the chain-rule test fails &
  \cref{plt:sec:fail-generator}\\
A numerical error bound is violated in experiment &
  \cref{plt:sec:fail-formal-tail} or \cref{plt:sec:fail-uniform-L}\\
Denominators in \(L\) appear where polynomials were promised &
  \cref{plt:sec:fail-nonunit}\\
Two runs differ by an additive constant &
  \cref{plt:sec:fail-resonance}\\
Every block of positive index is a block of a shifted problem &
  \cref{plt:sec:fail-resonance}\\
The first correction is larger than the term it corrects &
  \cref{plt:sec:fail-uniform-L}\\
A coefficient array overflows, or a cost estimate is off by a growing factor &
  \cref{plt:sec:fail-degree}\\
Two runs of the same input differ at all &
  \cref{plt:cor:cert-reproducible}\\
\bottomrule
\end{longtable}

\begin{summarybox}
The eight failure modes are eight implicit decisions.  Truncation convention
and guard blocks are decisions about indices; the generator images and the
% ed.: the list ran to nine items against eight failure modes, the branch
% ed.: being the extra one; it is folded into the item it travels with.
unit test are decisions about the class; the resonant constant --- and with it
the branch datum of \cref{plt:prop:cert-power} --- is a decision about which of
several correct answers is wanted; the formal
tail and the uniformity in \(L\) are decisions about what kind of statement is
being made; and the degree bound is a decision about which hypotheses are in
force.  \Cref{plt:alg:decision} orders them so that each is made when its
information is available, and \cref{plt:tab:cert-gauges} says what to compute
afterwards to find out whether the decisions were kept.
\end{summarybox}

\clearpage
\part{Extensions and the wider landscape}

% ---- fragment 15-extensions ----
\chapter{Generalized power--log classes, nested logarithms, and multiseries}
\label{plt:sec:extensions}

Parts~I--IV worked inside a deliberately small ambient: integer powers of
\(t=X^{-1}\), polynomial (or rational, or Laurent) coefficients in the single
logarithm \(L=\log X\), and no exponentials.  Every closure failure recorded
there --- the reciprocal of a nonunit
(\cref{plt:thm:div-unit-criterion}), a fractional leading exponent
(\cref{plt:thm:lag-rigidity}), a logarithm inside the leading monomial
(\cref{plt:prop:lag-loglog-forced}), the exponential of an element of negative
outer order (\cref{plt:thm:dif-exp-boundary}) --- names a specific enlargement.
This chapter makes each enlargement precise and, for each, says exactly which
theorems of Parts~I--IV survive verbatim, which survive after a hypothesis is
changed, and which become false.  The failures are proved, because they are the
statements a computation can actually violate; a class that is described only
by what it contains is a class one will silently misuse.

\section{Five independent axes}
\label{plt:sec:ext-axes}

It is useful to separate the enlargements into five axes that can be turned
independently.  Nothing in this subsection is a theorem; it is the map of the
territory, and each entry is made precise below or in the place indicated.

\begin{table}[htbp]
\centering
\caption{The five axes of enlargement.  The class of Parts~I--IV is the corner
\(\bigl(\Gamma=\IntegerNumbers,\ C=\K[L],\ r=1,\ \text{height }0,\
\text{well-based}\bigr)\).}
\label{plt:tab:ext-axes}
\small
\begin{tabular}{@{}p{0.17\textwidth}p{0.26\textwidth}p{0.26\textwidth}p{0.23\textwidth}@{}}
\toprule
Axis & What is enlarged & Forced by & Analysed in \\
\midrule
Exponent group in \(t\) &
\(\IntegerNumbers\rightsquigarrow s^{-1}\IntegerNumbers
\rightsquigarrow\Gamma\subseteq\RealNumbers\) &
fractional or irrational leading exponent
(\cref{plt:thm:lag-rigidity}) &
\cref{plt:sec:ext-hahn}, \cref{plt:sec:ext-puiseux} \\
\addlinespace
Coefficient ring in \(L\) &
\(\K[L]\rightsquigarrow\K(L)\rightsquigarrow\K((L^{-1}))\) &
division by a nonunit (\cref{plt:thm:div-unit-criterion}); the axis stops
short of \(\sqrt L\) (\cref{plt:prop:ext-sqrtL}) &
\cref{plt:sec:rings}, \cref{plt:sec:ext-puiseux} \\
\addlinespace
Logarithmic depth \(r\) &
one logarithm \(\rightsquigarrow L_1,\ldots,L_r\) &
a logarithm in the leading monomial
(\cref{plt:prop:lag-loglog-forced}); an unmatched residue
(\cref{plt:thm:ext-tower-strict}) &
\cref{plt:sec:scale-change}, \cref{plt:sec:ext-depth} \\
\addlinespace
Exponential height &
height \(0\rightsquigarrow\) transmonomials \(\EulerE^{\mathfrak p}\) &
\(\exp\) of an element of negative outer order
(\cref{plt:thm:dif-exp-boundary}, \cref{plt:thm:ext-logderiv}) &
\cref{plt:sec:full-transseries} \\
\addlinespace
Support class &
grid-based \(\rightsquigarrow\) left-finite \(\rightsquigarrow\) well-based &
infinite summation; completion
(\cref{plt:thm:ext-lf-completion}) &
\cref{plt:sec:ext-supports} \\
\bottomrule
\end{tabular}
\end{table}

\begin{warningbox}[title={The axes are not interchangeable}]
An enlargement along one axis never repairs an obstruction on another.
Ramifying the exponent group of \(t\) does not produce \(\sqrt{L}\)
(\cref{plt:prop:ext-sqrtL}); adding logarithmic depth does not produce
\(\EulerE^{-X}\) (\cref{plt:cor:ext-no-exponentials}); enlarging the
coefficient ring does not remove the obstruction \(\sigma\ge1\) of
\cref{plt:thm:lag-rigidity}; and no amount of exponential height supplies the
antiderivative of \(1/(L_0L_1\cdots L_r)\)
(\cref{plt:thm:ext-tower-strict}).  A closure test must therefore be run once
per axis, and a diagnosis that does not name the axis is not a diagnosis.
\end{warningbox}

\section{Supports: what well-ordering does, and what it does not do}
\label{plt:sec:ext-supports}

Throughout this subsection \(\Gamma\) is a totally ordered abelian group,
written additively, and a subset \(S\subseteq\Gamma\) is thought of as a set of
exponents of \(t\): a \emph{larger} exponent means a \emph{smaller} monomial,
so the least element of a support is its dominant monomial, and the
\(\succdom\)-greatest element of \cref{plt:thm:ring-support}\,(c) is the least
exponent.  We keep the exponent formulation because it makes the group
structure visible; the reader who prefers monomials should read every
``least'' as ``dominant''.

\begin{definition}[Three support classes]
\label{plt:def:ext-support-classes}
Let \(S\subseteq\Gamma\).
\begin{enumerate}[label=\textup{(\alph*)}]
\item \(S\) is \emph{well based} \textup{(}equivalently, well ordered\textup{)}
  if every nonempty subset of \(S\) has a least element.
\item \(S\) is \emph{left-finite} if \(S\cap(-\infty,M)\) is finite for every
  \(M\in\Gamma\).
\item \(S\) is \emph{grid-based} if there are \(\mathfrak n\in\Gamma\) and
  finitely many \(\alpha_1,\ldots,\alpha_k\in\Gamma_{>0}\) with
  \[
    S\subseteq
    \mathfrak n+\SetOf{n_1\alpha_1+\cdots+n_k\alpha_k:\
      n_1,\ldots,n_k\in\NaturalNumbers}.
  \]
\end{enumerate}
Each of the three classes is closed under passing to subsets, under finite
unions, and under translation.
\end{definition}

\begin{proof}[Verification of the closure claims]
Subsets and translates are immediate in all three cases; for (c) a translate
changes only \(\mathfrak n\).  Finite unions: for (a) and (b) this is clear.
For (c), let \(S\subseteq\mathfrak n+\sum\NaturalNumbers\alpha_i\) and
\(T\subseteq\mathfrak m+\sum\NaturalNumbers\beta_j\), put
\(\mathfrak p=\min(\mathfrak n,\mathfrak m)\), and take as generating set the
\(\alpha_i\), the \(\beta_j\), and whichever of
\(\mathfrak n-\mathfrak p,\ \mathfrak m-\mathfrak p\) is nonzero; both of the
latter are \(\ge0\), so all generators used are \(>0\), and
\(S\cup T\subseteq\mathfrak p+\sum\NaturalNumbers(\cdots)\).
\end{proof}

The whole content of \cref{plt:thm:ring-support} is that for
\(\Gamma=\IntegerNumbers\) the three classes coincide with ``bounded below''.
They separate as soon as \(\Gamma\) is dense, or as soon as \(\Gamma\) fails to
be archimedean; \cref{plt:thm:ext-support-hierarchy} says exactly how.

\subsection{Well-ordering and local finiteness}

The next lemma is the reason the catalogue calls the support condition ``what
makes multiplication locally finite''.  We give the proof in full because the
sources assert it and because part~(iii), which is what actually makes
reciprocals and geometric series legal, is not the same statement as
parts~(i)--(ii).

\begin{lemma}[Convolution of well-based supports]
\label{plt:lem:ext-neumann}
Let \(\Gamma\) be a totally ordered abelian group.
\begin{enumerate}[label=\textup{(\roman*)}]
\item If \(S,T\subseteq\Gamma\) are well based then so is
  \(S+T=\SetOf{\sigma+\tau:\sigma\in S,\ \tau\in T}\).
\item If \(S,T\) are well based then for every \(\gamma\in\Gamma\) the fiber
  \(R_{\gamma}=\SetOf{(\sigma,\tau)\in S\times T:\sigma+\tau=\gamma}\) is
  finite.
\item Let \(S\subseteq\Gamma_{>0}\) be well based and, for \(k\ge0\), let
  \(S^{(k)}=S+\cdots+S\) \textup{(}\(k\) summands; \(S^{(0)}=\SetOf{0}\)\textup{)}.
  Then every \(S^{(k)}\) is well based, and for every \(k\) and every
  \(\gamma\) only finitely many multisets of \(k\) elements of \(S\) have sum
  \(\gamma\).  If moreover \(\Gamma\) is archimedean, then
  \(S^{\ast}=\bigcup_{k\ge0}S^{(k)}\) is well based and every
  \(\gamma\in S^{\ast}\) has, in total, only finitely many representations as a
  sum of elements of \(S\).
\end{enumerate}
\end{lemma}

\begin{proof}
We use twice the following remark: \emph{every sequence in a well-based set has
a non-decreasing subsequence}.  Indeed, given \((\sigma_k)_{k\ge1}\) in the
well-based set \(S\), the set \(\SetOf{\sigma_k:k\ge1}\) has a least element;
let \(k_1\) attain it.  Having chosen \(k_1<\cdots<k_i\), let \(k_{i+1}>k_i\)
attain the least element of the nonempty set \(\SetOf{\sigma_k:k>k_i}\).  Since
\(\SetOf{\sigma_k:k>k_i}\subseteq\SetOf{\sigma_k:k>k_{i-1}}\) while
\(\sigma_{k_i}\) is least in the larger set, \(\sigma_{k_i}\le\sigma_{k_{i+1}}\).

(ii)  Suppose \(R_{\gamma}\) is infinite.  As \(\tau=\gamma-\sigma\) is
determined by \(\sigma\), infinitely many distinct \(\sigma\) occur; choose
pairwise distinct \(\sigma_k\in S\) with \(\gamma-\sigma_k\in T\).  Passing to a
subsequence we may assume \(\sigma_{k_1}\le\sigma_{k_2}\le\cdots\), and since
they are pairwise distinct the inequalities are strict.  Then
\(\gamma-\sigma_{k_i}\) is a strictly decreasing sequence in \(T\), so
\(\SetOf{\gamma-\sigma_{k_i}:i\ge1}\) is a nonempty subset of \(T\) with no
least element, contradicting well-basedness of \(T\).

(i)  Suppose \(\gamma_1>\gamma_2>\cdots\) is strictly decreasing in \(S+T\),
say \(\gamma_k=\sigma_k+\tau_k\).  Pass to a subsequence along which
\((\sigma_k)\) is non-decreasing; along it \(\tau_k=\gamma_k-\sigma_k\) is
strictly decreasing in \(T\), a contradiction as before.  Hence \(S+T\) has no
infinite strictly decreasing sequence; a totally ordered set with that property
is well based, since a nonempty subset without a least element would produce
one by dependent choice.

(iii)  That \(S^{(k)}\) is well based follows from (i) by induction on \(k\).
For the representation count, induct on \(k\); the case \(k\le1\) is trivial
(the empty multiset for \(\gamma=0\), and at most one singleton otherwise).
Let \(k\ge2\) and suppose infinitely many multisets
\(\SetOf{s_1\le s_2\le\cdots\le s_k}\subseteq S\) have sum \(\gamma\).  If the
value \(s_1\) took only finitely many values, then for one of them infinitely
many multisets of size \(k-1\) would sum to \(\gamma-s_1\), contradicting the
inductive hypothesis.  So infinitely many distinct values of \(s_1\) occur; by
the remark, choose a strictly increasing sequence
\(s_1^{(1)}<s_1^{(2)}<\cdots\) of such values.  Then
\(\gamma-s_1^{(i)}\) is a strictly decreasing sequence in the well-based set
\(S^{(k-1)}\), a contradiction.

Now assume \(\Gamma\) archimedean.  Since \(S\) is well based and
\(S\subseteq\Gamma_{>0}\), it has a least element \(\delta>0\) (we may assume
\(S\ne\emptyset\), else there is nothing to prove).  A sum of \(k\) elements of
\(S\) is \(\ge k\delta\); by archimedeanity, for each \(M\) only finitely many
\(k\) satisfy \(k\delta<M\).  Hence for \(\emptyset\ne T\subseteq S^{\ast}\)
and \(\tau\in T\) we have
\(T\cap(-\infty,\tau]\subseteq\bigcup_{k\le K}S^{(k)}\) for some finite \(K\),
a finite union of well-based sets and therefore well based; its least element
is least in \(T\).  The total representation count of \(\gamma\) is the sum,
over the finitely many admissible \(k\), of the counts already bounded.
\end{proof}

\begin{remark}[Neumann's lemma without archimedeanity]
\label{plt:rmk:ext-neumann-general}
The last clause of \cref{plt:lem:ext-neumann}\,(iii) is true for an arbitrary
totally ordered abelian group --- this is Neumann's lemma, and it is what makes
a Hahn series ring over any such group a field when the coefficients form a
field \cite{hahn,edgar-beginners,vanderhoeven-book}.  The proof above uses
archimedeanity only in the last paragraph, and \(\Gamma\subseteq\RealNumbers\)
is archimedean, so every use of the clause in this volume with a
\emph{one-dimensional} exponent group is covered by the proof given.  The
monomial groups \(\mathcal M_r(\Gamma)\) of \cref{plt:def:grp-logtower} are
\emph{not} archimedean for \(r\ge1\), and there we quote Neumann's lemma rather
% ed.: the draft claimed that no argument of \cref{plt:sec:ext-depth} needs
% more than parts (i)--(ii).  That is not accurate: the statements of
% \cref{plt:thm:ext-tower-strict} and \cref{plt:cor:ext-depth-unbounded} call
% H_r a differential FIELD, and the field property of a Hahn series ring over a
% non-archimedean monomial group is exactly Neumann's lemma.  What is true is
% that every argument carried out here needs only (i)--(ii); the field property
% is imported, not reproved.
than reprove it.  Every \emph{argument} carried out in
\cref{plt:sec:ext-depth} uses only \cref{plt:lem:ext-neumann}\,(i)--(ii); the
one further input, used only through the word ``field'', is that
\(\mathcal H_r(\K,\Gamma)\) is a field, which is
\cref{plt:def:grp-logtower} together with Neumann's lemma and is quoted, not
reproved.
\end{remark}

\begin{proposition}[Well-basedness is the largest admissible support class]
\label{plt:prop:ext-wo-maximal}
Let \(\Gamma\) be a totally ordered abelian group and let \(T\subseteq\Gamma\)
fail to be well based.  Then there is a well-based \(S\subseteq\Gamma\), of
order type \(\omega\), and a \(\gamma\in\Gamma\) whose fiber
\(R_{\gamma}\subseteq S\times T\) is infinite.  Consequently, if \(\mathcal F\)
is a family of subsets of \(\Gamma\) that contains every well-based set and for
which every convolution fiber \(R_{\gamma}\subseteq S\times T\) with
\(S,T\in\mathcal F\) is finite, then \(\mathcal F\) is exactly the family of
well-based subsets of \(\Gamma\).
\end{proposition}

\begin{proof}
If \(T\) is not well based it contains a strictly decreasing sequence
\(\tau_1>\tau_2>\cdots\).  Put \(S=\SetOf{-\tau_k:k\ge1}\); then
\(-\tau_1<-\tau_2<\cdots\), so \(S\) is well based of order type \(\omega\).
With \(\gamma=0\), every pair \((-\tau_k,\tau_k)\) lies in \(R_0\), which is
therefore infinite.  For the second assertion, a member \(T\in\mathcal F\)
failing to be well based would, together with the well-based
\(S\in\mathcal F\) just produced, violate the finiteness hypothesis.
\end{proof}

\begin{remark}[What the sharp statement is, and what it is not]
\label{plt:rmk:ext-lf-not-necessary}
% ed.: S1 ("What changes is the proof that every target coefficient receives
% only finitely many contributions ... generalized supports make it a central
% axiom") and S3 ("The well-order condition guarantees that each coefficient
% receives only an admissible family of contributions") both present
% well-ordering as if it were equivalent to local finiteness of a single
% product.  It is sufficient, and it is necessary for the class, but it is not
% necessary for an individual product; the witness below shows this.
\Cref{plt:lem:ext-neumann} says well-basedness suffices, and
\cref{plt:prop:ext-wo-maximal} says it cannot be relaxed \emph{as a condition
on a class of supports}.  It is not necessary for a single product.  Take
\(\Gamma=\RationalNumbers\) and
\[
  T=\SetOf{1+\tfrac1k:\ k\in\PositiveIntegers}
   =\SetOf{2,\tfrac32,\tfrac43,\tfrac54,\ldots},
\]
which is strictly decreasing and therefore not well based.  Every element of
\(T+T\) lies in \((2,4]\), and a representation
\(\gamma=(1+\frac1k)+(1+\frac1l)\) with \(k\le l\) forces
\(\frac1k+\frac1l=\gamma-2\), hence \(\frac1k\ge\frac{\gamma-2}{2}\) and
\(\frac1k<\gamma-2\), so
% ed.: the naive bound "k,l <= 2/(gamma-2)" is false -- for gamma-2 = 1/2 the
% pair (3,6) is a representation with l = 6 > 4.  Only the smaller index is
% bounded above; the larger is then determined.
\[
  (\gamma-2)^{-1}<k\le 2(\gamma-2)^{-1},
\]
which leaves finitely many \(k\), and each \(k\) determines
\(l=\bigl(\gamma-2-\frac1k\bigr)^{-1}\) uniquely.  Thus \(T\) is a
non-well-based support for which the square of the corresponding formal sum has
finite coefficient fibers.  What fails is \emph{closure}: by
\cref{plt:prop:ext-wo-maximal} the same \(T\) has an infinite fiber against the
perfectly legitimate support \(\SetOf{-1-\frac1k:k\ge1}\), whose elements
increase to \(-1\).  A support condition is a condition on a class, not on an
element, and this is why the axiom is stated as well-basedness.
\end{remark}

\subsection{Left-finite supports}

Left-finiteness is the exact condition under which the finite bookkeeping of
Parts~II and~VI --- truncation, block recursion, sparse representation ---
% ed.: the draft said left-finiteness "is literally \cref{plt:thm:ring-support}(d)
% with Z replaced by Gamma".  That is false and contradicts
% \cref{plt:prop:ext-hahn-trunc-fails}(b) below, which observes that clause (d)
% read verbatim over a general Gamma is strictly weaker: it asks only for an
% exhausting enumeration of order type at most omega and does not demand that
% the enumeration be cofinal.  Over Z cofinality is automatic, which is why the
% two conditions are indistinguishable in Parts I--IV.
survives a change of exponent group.  It is
\cref{plt:thm:ring-support}\,(d) with \(\IntegerNumbers\) replaced by
\(\Gamma\), strengthened by the cofinality clause of
\cref{plt:prop:ext-leftfinite}\,(i), which is automatic over
\(\IntegerNumbers\) and is not automatic over a dense \(\Gamma\).

\begin{proposition}[Left-finite supports]
\label{plt:prop:ext-leftfinite}
Let \(\Gamma\) be a totally ordered abelian group with \(\Gamma\ne\SetOf{0}\)
and let \(S\subseteq\Gamma\).
\begin{enumerate}[label=\textup{(\roman*)}]
\item \(S\) is left-finite if and only if it can be enumerated as a strictly
  increasing sequence \(\gamma_0<\gamma_1<\cdots\), finite or of order type
  \(\omega\), which in the infinite case is cofinal in \(\Gamma\): for every
  \(M\in\Gamma\) some \(\gamma_k>M\).
\item Left-finite \(\Rightarrow\) well based, and the converse fails whenever
  \(\Gamma\) is densely ordered: for \(\Gamma=\RationalNumbers\) the set
  \(S=\SetOf{1-\frac1k:k\in\PositiveIntegers}\) is well based of order type
  \(\omega\) and is not left-finite.
\item If \(S,T\) are left-finite then so are \(S\cup T\), \(S+T\) and every
  translate of \(S\), and all convolution fibers are finite.
\item Assume in addition that \(\Gamma\) is archimedean.  If
  \(S\subseteq\Gamma_{>0}\) is left-finite then \(S^{\ast}\) is left-finite,
  and every \(\gamma\in S^{\ast}\) arises from only finitely many multisets of
  elements of \(S\).
\end{enumerate}
\end{proposition}

\begin{proof}
(i)  Let \(S\) be left-finite and nonempty.  For \(\mu\in S\) the set
\(S\cap(-\infty,\mu]\) is finite and nonempty, so \(S\) has a minimum
\(\gamma_0\); removing it leaves a left-finite set, and iterating produces a
strictly increasing enumeration in which \(\mu\in S\) is reached after at most
\(\#\bigl(S\cap(-\infty,\mu]\bigr)\) steps, so every element occurs.  If the
enumeration is infinite and \(\gamma_k\le M\) for all \(k\), pick
\(g\in\Gamma_{>0}\); then \(S\cap(-\infty,M+g)=S\) is infinite, contradicting
left-finiteness.  Conversely, given such an enumeration and \(M\in\Gamma\),
choose \(K\) with \(\gamma_K>M\); then
\(S\cap(-\infty,M)\subseteq\SetOf{\gamma_0,\ldots,\gamma_{K-1}}\).

(ii)  If \(\emptyset\ne T\subseteq S\) with \(S\) left-finite, pick
\(\tau\in T\); then \(T\cap(-\infty,\tau]\) is finite and nonempty, so it has a
least element, which is least in \(T\).  For the counterexample, the displayed
\(S\) is enumerated by the strictly increasing sequence \(1-\frac1k\), so every
nonempty subset has a least element; but \(S\subseteq(-\infty,1)\), so
\(S\cap(-\infty,1)=S\) is infinite.  It is exactly the cofinality clause of (i)
that fails.  In a general densely ordered \(\Gamma\) the same construction
works: pick \(a<b\) and, recursively, \(\gamma_{k+1}\) strictly between
\(\gamma_{k}\) and \(b\) starting from \(\gamma_{1}\in(a,b)\); then
\(\SetOf{\gamma_{k}}\) is well based of order type \(\omega\), and it is not
left-finite because it is infinite and contained in \((-\infty,b)\).

(iii)  Unions and translates are immediate.  A nonempty left-finite set has a
minimum by (i); write \(\sigma_0=\min S\), \(\tau_0=\min T\).  If
\(\sigma+\tau<M\) with \(\sigma\in S\), \(\tau\in T\), then
\(\sigma<M-\tau_0\) and \(\tau<M-\sigma_0\), so \(\sigma,\tau\) range over
finite sets; hence \((S+T)\cap(-\infty,M)\) is finite and every fiber is
finite.

(iv)  Put \(\delta=\min S>0\).  A sum of \(k\) elements of \(S\) is at least
\(k\delta\), so a sum \(<M\) forces \(k\delta<M\), and archimedeanity leaves
finitely many such \(k\).  Each summand of such a sum is
\(<M-(k-1)\delta\le M\), so lies in the finite set \(S\cap(-\infty,M)\).  There
are therefore finitely many admissible multisets, which proves both assertions.
\end{proof}

\begin{remark}[Archimedean is not decoration]
\label{plt:rmk:ext-archimedean}
Part~(iv) fails for non-archimedean \(\Gamma\).  Take
\(\Gamma=\IntegerNumbers\times\IntegerNumbers\) ordered lexicographically and
\(S=\SetOf{(0,1)}\), which is left-finite, indeed finite.  Then
\(S^{\ast}=\SetOf{(0,k):k\in\NaturalNumbers}\) is infinite and contained in
\((-\infty,(1,0))\), so it is not left-finite.  The failure is precisely the
failure of \(\IntegerNumbers\delta\) to be cofinal, and it is the same
phenomenon that would break separatedness of a filtration indexed by
\(\Gamma\).  This is why \cref{plt:def:dif-scale} fixes
\(\Gamma\subseteq\RealNumbers\); the restriction cannot be dropped without
redoing every summability argument of Parts~II--IV.  It also means that in the
nested-logarithmic setting of \cref{plt:sec:ext-depth}, where the monomial
group is non-archimedean by construction, ``left-finite'' is the wrong notion
and only well-basedness survives.
\end{remark}

\begin{theorem}[The support hierarchy, and where it collapses]
\label{plt:thm:ext-support-hierarchy}
Let \(\Gamma\) be a totally ordered abelian group.
\begin{enumerate}[label=\textup{(\roman*)}]
\item Grid-based \(\Rightarrow\) well based, always.
\item Left-finite \(\Rightarrow\) well based, always.
\item If \(\Gamma\) is archimedean, grid-based \(\Rightarrow\) left-finite.
\item If \(\Gamma\) is not archimedean, grid-based \(\not\Rightarrow\)
  left-finite.
\item If \(\Gamma\) is archimedean and not densely ordered --- equivalently, by
  H\"older's theorem, \(\Gamma\) is cyclic --- then all three classes coincide
  with ``bounded below''.
\item For \(\Gamma=\RationalNumbers\) both inclusions in
  \(\text{grid-based}\subsetneq\text{left-finite}\subsetneq\text{well based}\)
  are strict.
\end{enumerate}
\end{theorem}

\begin{proof}
(i)  Let \(S\subseteq\mathfrak n+M\) with
\(M=\SetOf{\sum_in_i\alpha_i:n_i\in\NaturalNumbers}\) and
\(\alpha_i\in\Gamma_{>0}\).  It suffices to show \(M\) is well based.  Suppose
\(v^{(1)},v^{(2)},\ldots\in\NaturalNumbers^{k}\) give a strictly decreasing
sequence \(\sum_iv^{(m)}_i\alpha_i\).  By Dickson's lemma the componentwise
order on \(\NaturalNumbers^{k}\) is a well-quasi-order, so there are
\(m<m'\) with \(v^{(m)}\le v^{(m')}\) componentwise; as all \(\alpha_i>0\)
this gives \(\sum_iv^{(m)}_i\alpha_i\le\sum_iv^{(m')}_i\alpha_i\), contradicting
strict decrease.  So \(M\) has no infinite strictly decreasing sequence and is
well based; well-basedness passes to subsets and translates.

(ii)  \Cref{plt:prop:ext-leftfinite}\,(ii).

(iii)  With \(S\) and \(M\) as in (i) and \(M'\in\Gamma\) given, an element
\(\mathfrak n+\sum_in_i\alpha_i<M'\) has \(n_i\alpha_i<M'-\mathfrak n\) for each
\(i\), because the omitted summands are \(\ge0\); archimedeanity leaves finitely
many \(n_i\) for each \(i\), hence finitely many tuples and finitely many
values.

(iv)  In \(\IntegerNumbers\times\IntegerNumbers\) ordered lexicographically the
set \(M=\SetOf{(0,n):n\in\NaturalNumbers}\) is grid-based with
\(\mathfrak n=(0,0)\) and the single generator \(\alpha_1=(0,1)>0\), and it is
not left-finite, as in \cref{plt:rmk:ext-archimedean}.

(v)  An archimedean group is order-isomorphic to a subgroup of
\(\RealNumbers\) (H\"older), and a subgroup of \(\RealNumbers\) is either dense
or cyclic.  If \(\Gamma=\IntegerNumbers\delta\) is cyclic, a set bounded below
is finite below every bound, hence left-finite; it is then well based by (ii);
and it is grid-based with generator \(\delta\) and base point any lower bound.
Conversely each class consists of sets bounded below: a well-based set has a
least element if nonempty.  This is \cref{plt:thm:ring-support} in the case
\(\Gamma=\IntegerNumbers\).

(vi)  Left-finite \(\subsetneq\) well based is
\cref{plt:prop:ext-leftfinite}\,(ii).  For grid-based \(\subsetneq\)
left-finite, take
\[
  S=\SetOf{k+\tfrac1k\ :\ k\in\PositiveIntegers}
   \subseteq\RationalNumbers .
\]
It is enumerated by a strictly increasing sequence tending to \(+\infty\),
hence left-finite by \cref{plt:prop:ext-leftfinite}\,(i).  If it were
grid-based,
it would lie in the subgroup of \(\RationalNumbers\) generated by
\(\mathfrak n,\alpha_1,\ldots,\alpha_k\); a finitely generated subgroup of
\(\RationalNumbers\) is cyclic, say \(N^{-1}\IntegerNumbers\), and
\(k+\frac1k\in N^{-1}\IntegerNumbers\) forces \(N/k\in\IntegerNumbers\), that
is \(k\mid N\), which holds for only finitely many \(k\).
\end{proof}

\begin{checkpointbox}
Three different things are being asked of a support, and they are not the same
thing.  \emph{Well based} is what makes the product defined at all
(\cref{plt:lem:ext-neumann}), and by \cref{plt:prop:ext-wo-maximal} it is the
weakest condition that can be imposed on a whole class.  \emph{Left-finite} is
what makes truncation produce finite data
(\cref{plt:thm:ext-lf-completion}\,(ii)).  \emph{Grid-based} is what makes the
exponent arithmetic itself finitary: exponents are integer vectors on a fixed
finite basis.  In the integer-exponent world of Parts~I--IV all three collapse
into ``bounded below'', which is exactly why the distinction is invisible there
and unavoidable here.
\end{checkpointbox}

\section{Real exponents: Hahn power--log scales}
\label{plt:sec:ext-hahn}

\Cref{plt:def:dif-scale} already introduced the power--log scale
\(\mathcal H(C,\Gamma)\): well-based supports in an ordered group
\(\IntegerNumbers\subseteq\Gamma\subseteq\RealNumbers\), blocks in
\(C\in\SetOf{\K'[L],\K'(L),\K'((L^{-1}))}\), and the derivation
\eqref{plt:eq:dif-scale-derivation}.  We now isolate inside it the subring on
which the finite bookkeeping of Parts~II and~VI is still meaningful, and prove
that outside that subring it is not.

\begin{definition}[The left-finite power--log scale]
\label{plt:def:ext-lf-scale}
Let \(\mathcal H=\mathcal H(C,\Gamma)\) be a power--log scale in the sense of
\cref{plt:def:dif-scale}.  For \(F=\sum_{\gamma\in S}p_{\gamma}t^{\gamma}\) put
\(\operatorname{Exp}(F)=\SetOf{\gamma\in\Gamma:p_{\gamma}\ne0}\), the
\emph{exponent set} of \(F\) --- the projection to \(\Gamma\) of
\(\CoefficientSupportOf{F}\).  Define
\[
  \mathcal H_{\mathrm{lf}}(C,\Gamma)
  \DefinitionEquals
  \SetOf{F\in\mathcal H(C,\Gamma):\ \operatorname{Exp}(F)\ \text{is
    left-finite}},
\]
and write \(C[t^{\Gamma}]\) for the subring of elements with finite exponent
set.
\end{definition}

\begin{theorem}[Truncation, completeness, and what \(\mathcal H_{\mathrm{lf}}\) is]
\label{plt:thm:ext-lf-completion}
Let \(\mathcal H=\mathcal H(C,\Gamma)\) be a power--log scale.
\begin{enumerate}[label=\textup{(\roman*)}]
\item \(\mathcal H_{\mathrm{lf}}(C,\Gamma)\) is a subring of \(\mathcal H\)
  containing \(C[t^{\Gamma}]\), stable under \(\DX\), and stable under the
  inversion of those of its elements that are units of \(\mathcal H\).
\item For \(F\in\mathcal H\): every truncation \(\Trunc{F}{N}\),
  \(N\in\Gamma\), has finite exponent set if and only if
  \(F\in\mathcal H_{\mathrm{lf}}(C,\Gamma)\).
\item Both \(\mathcal H\) and \(\mathcal H_{\mathrm{lf}}(C,\Gamma)\) are
  separated and complete for the filtration
  \(\mathfrak F_{N}=\SetOf{F\in\mathcal H:\ordt F\ge N}\), \(N\in\Gamma\),
  whose topology is metrizable because \(\IntegerNumbers\) is cofinal in
  \(\Gamma\).
\item \(\mathcal H_{\mathrm{lf}}(C,\Gamma)\) is the closure of
  \(C[t^{\Gamma}]\) in \(\mathcal H\); equivalently it \emph{is} the
  completion of \(C[t^{\Gamma}]\) for the induced filtration.  If \(\Gamma\) is
  dense in \(\RealNumbers\) then
  \(\mathcal H_{\mathrm{lf}}(C,\Gamma)\subsetneq\mathcal H(C,\Gamma)\).
\end{enumerate}
\end{theorem}

\begin{proof}
(i)  Closure under sums and products is
\cref{plt:prop:ext-leftfinite}\,(iii), the block of a product at \(t^{\gamma}\)
being the finite sum over the fiber, which is finite by
\cref{plt:lem:ext-neumann}\,(ii).  Stability under \(\DX\) holds because
\eqref{plt:eq:dif-scale-derivation} translates the exponent set by
\(1\in\Gamma\).  For inverses, let \(F\in\mathcal H_{\mathrm{lf}}\) be a unit of
\(\mathcal H\) and write \(e=\ordt F\), \(p_{e}=[t^{e}]F\).  If \(FG=1\) then
\(\ordt G=-e\) and, by the leading-block law
(\cref{plt:lem:dif-scale-basic}(i) and its proof, or directly: the coefficient
of \(t^{0}\) in \(FG\) receives exactly one contribution, from the two leading
blocks), \(p_{e}\,[t^{-e}]G=1\), so \(p_{e}\in C^{\times}\).  Hence
\(F=t^{e}p_{e}(1+v)\) with \(v=t^{-e}p_{e}^{-1}F-1\in\mathcal H_{\mathrm{lf}}\)
and \(\ordt v>0\); then
\(F^{-1}=t^{-e}p_{e}^{-1}\sum_{k\ge0}(-v)^{k}\).  Put
\(S=\operatorname{Exp}(v)\subseteq\Gamma_{>0}\), which is left-finite.  By
\cref{plt:prop:ext-leftfinite}\,(iv), \(S^{\ast}\) is left-finite and each of
its elements has finitely many representations, so the family
\(\bigl((-v)^{k}\bigr)_{k\ge0}\) is summable and its sum has exponent set
inside \(S^{\ast}\); translating by \(-e\) preserves left-finiteness.

(ii)  If \(\operatorname{Exp}(F)\) is left-finite then
\(\operatorname{Exp}(F)\cap(-\infty,N)\), which is
\(\operatorname{Exp}(\Trunc{F}{N})\), is finite.  Conversely, if every
\(\Trunc{F}{N}\) has finite exponent set then
\(\operatorname{Exp}(F)\cap(-\infty,N)\) is finite for every \(N\in\Gamma\).

(iii)  \emph{Metrizability and separatedness.}  Since \(1\in\Gamma\subseteq
\RealNumbers\) and \(\Gamma\) is archimedean, no element of \(\Gamma\) exceeds
every integer, so \(\IntegerNumbers\) is cofinal in \(\Gamma\) and the
countable subfamily \((\mathfrak F_{N})_{N\in\IntegerNumbers}\) is a
neighbourhood basis of \(0\).  If \(\ordt F\ge N\) for every
\(N\in\IntegerNumbers\) then \(F=0\), for otherwise \(\ordt F\in\Gamma\) would
exceed every integer.
\emph{Completeness.}  Let \((F_{r})\) be Cauchy: for each
\(N\in\IntegerNumbers\) there is \(r(N)\) with
\(\Trunc{F_{r}}{N}=\Trunc{F_{r(N)}}{N}\) for all \(r\ge r(N)\), and we may take
\(r(N)\) non-decreasing in \(N\).  Define \(F\) blockwise by
\([t^{\gamma}]F=[t^{\gamma}]F_{r(N)}\) for any integer \(N>\gamma\); the
definition is consistent.  Then
\[
  \operatorname{Exp}(F)
  =\bigcup_{N\in\IntegerNumbers}
   \Bigl(\operatorname{Exp}(F_{r(N)})\cap(-\infty,N)\Bigr),
\]
which is well based: for \(\emptyset\ne T\subseteq\operatorname{Exp}(F)\) pick
\(\tau\in T\) and an integer \(N>\tau\); then \(T\cap(-\infty,N)\) is a
nonempty subset of the well-based set \(\operatorname{Exp}(F_{r(N)})\) and its
least element is least in \(T\).  So \(F\in\mathcal H\) and
\(\ordt(F-F_{r})\ge N\) for \(r\ge r(N)\).  If all \(F_{r}\) lie in
\(\mathcal H_{\mathrm{lf}}\) then so does \(F\), because
\(\operatorname{Exp}(F)\cap(-\infty,N)=\operatorname{Exp}(F_{r(N)})\cap
(-\infty,N)\) is finite for every integer \(N\), and every \(M\in\Gamma\) is
below some integer.

(iv)  \(\mathcal H_{\mathrm{lf}}\) is closed in \(\mathcal H\) by the last
sentence of (iii); it contains \(C[t^{\Gamma}]\); and every
\(F\in\mathcal H_{\mathrm{lf}}\) is the limit of the finitely supported
truncations \(\Trunc{F}{N}\), \(N\in\IntegerNumbers\), by (ii).  Conversely a
limit of finitely supported elements is left-finite, since
\(\Trunc{F}{N}=\Trunc{G}{N}\) with \(G\) finitely supported makes
\(\Trunc{F}{N}\) finitely supported.  Hence \(\mathcal H_{\mathrm{lf}}\) is the
closure of \(C[t^{\Gamma}]\), and a closed subring containing a subring is its
completion when the ambient is complete and separated.  Finally, if \(\Gamma\)
is dense in \(\RealNumbers\) then \(\Gamma\cap(0,1)\) is infinite; an infinite
bounded subset of \(\RealNumbers\) contains a strictly monotone sequence, and
by passing to a subsequence we obtain \(\gamma_{1}<\gamma_{2}<\cdots\) in
\(\Gamma\cap(0,1)\).  Put \(F=1+t+\sum_{k\ge1}t^{\gamma_{k}}\).  Its exponent
set \(\SetOf{0}\cup\SetOf{\gamma_{k}:k\ge1}\cup\SetOf{1}\) is well based --- a
nonempty subset containing \(0\) has least element \(0\); one meeting
\(\SetOf{\gamma_{k}}\) has as least element the \(\gamma_{k}\) of least index;
and the only remaining subset is \(\SetOf{1}\) --- so \(F\in\mathcal H\); it is
not left-finite, since \(\operatorname{Exp}(F)\cap(-\infty,1)\) is infinite, so
\(F\notin\mathcal H_{\mathrm{lf}}\).
\end{proof}

\begin{corollary}[The left-finite scale carries the near-identity theory]
\label{plt:cor:ext-lf-group}
Let \(\mathcal H=\mathcal H(C,\Gamma)\) and let
\(\mathcal H_{\ge0}=\SetOf{A\in\mathcal H:\ordt A\ge0}\).
\begin{enumerate}[label=\textup{(\roman*)}]
\item Every statement of \cref{plt:sec:near-identity-group} and of
  \cref{plt:sec:reversion-existence} whose proof uses only
  \textup{(a)} that \(\DX\) raises \(\ordt\) by at least \(1\),
  \textup{(b)} that a family with \(\ordt\to\infty\) may be summed termwise,
  \textup{(c)} that a continuous \(\K\)-algebra endomorphism is determined by
  the image of the monomial \(t\) and of \(L\) together with the binomial rule
  for \((1+u)^{\gamma}\), \(\gamma\in\Gamma\), \(\ordt u>0\), and
  \textup{(d)} that \(\ordt\) is additive,
  holds in \(\mathcal H\) with \(\PLpow\) replaced by \(\mathcal H_{\ge0}\) and
  the integer filtration by the \(\Gamma\)-filtration \(\mathfrak F_{N}\).
  This covers the substitution operator
  \textup{(\cref{plt:prop:grp-substitution})}, the Taylor formula
% ed.: the draft cited a nonexistent label plt:thm:grp-taylor; the formal
% Taylor formula of this volume is \cref{plt:thm:tay-formula}.
  \textup{(\cref{plt:thm:tay-formula})}, the filtered gain
  \textup{(\cref{plt:cor:grp-gain})}, the group law
  \textup{(\cref{plt:thm:grp-group})}, separation and completeness of the
  contact filtration \textup{(\cref{plt:prop:grp-limit})}, and existence and
  uniqueness of the near-identity inverse
  \textup{(\cref{plt:thm:rev-existence})}.
\item The subset
  \[
    \SetOf{X+A\ :\ A\in\mathcal H_{\mathrm{lf}}(C,\Gamma),\ \ordt A\ge0}
  \]
  is a subgroup of the group of \textup{(i)}: it is closed under composition
  and under compositional inversion.
\end{enumerate}
\end{corollary}

\begin{proof}
(i)  Inspect the four inputs.  \(\DX\) raises \(\ordt\) by \(1\) on monomials
by \eqref{plt:eq:dif-scale-derivation} and hence on all of \(\mathcal H\);
summability of \((A^{k}\DX^{k}F/k!)_{k\ge0}\) follows from
\(\ordt(A^{k}\DX^{k}F)\ge\ordt F+k\), and \(\IntegerNumbers\) is cofinal in
\(\Gamma\), so only finitely many \(k\) affect any given block; the
substitution operator is fixed on a general monomial by
\(t^{\gamma}\compose(X+A)=t^{\gamma}(1+tA)^{-\gamma}\), the binomial series
being summable since \(\ordt(tA)\ge1\)
(\cref{plt:thm:mul-binomial}, \cref{plt:prop:dif-powers}), and on \(L\) by
\(L\compose(X+A)=L+\log(1+tA)\); additivity of \(\ordt\) is
\cref{plt:lem:dif-scale-basic}(i).  Each of the six listed statements uses only
these.

(ii)  Let \(A,B\in\mathcal H_{\mathrm{lf}}\) with
\(\ordt A,\ordt B\ge0\), and put \(S_{A}=\operatorname{Exp}(A)\),
\(S_{B}=\operatorname{Exp}(B)\), \(S_{B}'=S_{B}\cap\Gamma_{>0}\).  By the
Taylor formula, \(A\compose(X+B)=\sum_{k\ge0}B^{k}\DX^{k}A/k!\) and
\[
  \operatorname{Exp}\bigl(B^{k}\DX^{k}A\bigr)
  \ \subseteq\ (S_{A}+k)+(S_{B}')^{\ast}
  \ \subseteq\ (S_{A}+k)+\Gamma_{\ge0},
\]
because \(k\) copies of \(S_{B}\) contribute a sum of elements of
\(S_{B}'\) together with some copies of \(0\).  Fix \(M\in\Gamma\).  A
contribution to \((-\infty,M)\) requires \(\min S_{A}+k<M\), so only finitely
many \(k\) contribute, since \(\IntegerNumbers\) is cofinal in the archimedean
\(\Gamma\).  For each such \(k\), the set \((S_{A}+k)+(S_{B}')^{\ast}\) is
left-finite by \cref{plt:prop:ext-leftfinite}\,(iii)--(iv).  A finite union of
left-finite sets is left-finite, so \(A\compose(X+B)\in
\mathcal H_{\mathrm{lf}}\).  Closure under inversion follows: the inverse of
\(X+A\) is \(X+B\) with \(B\) the \(\mathfrak F\)-limit of
\(B_{0}=0\), \(B_{n+1}=-A\compose(X+B_{n})\)
(\cref{plt:thm:rev-existence}), every \(B_{n}\) lies in
\(\mathcal H_{\mathrm{lf}}\) by what has just been proved and induction, and
\(\mathcal H_{\mathrm{lf}}\) is closed in \(\mathcal H\) by
\cref{plt:thm:ext-lf-completion}\,(iv).
\end{proof}

\begin{warningbox}[title={What clause (i) deliberately omits}]
The statements about the \emph{associated graded} of the contact filtration ---
\cref{plt:prop:grp-commutator}, \cref{plt:prop:grp-graded},
\cref{plt:thm:grp-gradedlie}, \cref{plt:thm:grp-exp} --- are formulated using
that \(r\mapsto r+1\) is the successor in the index group.  Write
\(\mathcal G_{\gamma}\) for the step of the contact filtration of
\cref{plt:def:grp-filtration} indexed by \(\gamma\in\Gamma\).  For dense
% ed.: the draft used an undefined macro \Gfilt here; the volume preamble
% provides no such macro, so the filtration steps are written out.
\(\Gamma\) the quotients \(\mathcal G_{\gamma}/\mathcal G_{\gamma+1}\) remain
defined,
because the gain is the element \(1\in\Gamma\), but they are no longer single
steps of the filtration, and the identification of the graded object with a Lie
algebra is not the same statement.  We do not claim it here.  The same caution
applies to any argument that enumerates blocks one at a time; see the entry for
\cref{plt:thm:rev-triangular} in \cref{plt:tab:ext-survival}.
\end{warningbox}

\begin{proposition}[Truncation, and therefore all of Part~VI, fails in the
general Hahn scale]
\label{plt:prop:ext-hahn-trunc-fails}
Let \(\Gamma\subseteq\RealNumbers\) be dense with
\(\IntegerNumbers\subseteq\Gamma\), and let
\(\mathcal H=\mathcal H(\K[L],\Gamma)\).  There is
\(F\in\mathcal H\) with \(\ordt F=0\) and every block equal to \(1\) such that:
\begin{enumerate}[label=\textup{(\alph*)}]
\item \(\Trunc{F}{N}\) has infinite exponent set for every \(N\ge1\) in
  \(\Gamma\);
\item \(\CoefficientSupportOf{F}\) admits no enumeration as a strictly
  \(\succdom\)-decreasing sequence of order type at most \(\omega\) exhausting
  it, so \cref{plt:thm:ring-support}\,(d) is false in \(\mathcal H\), although
  \cref{plt:thm:ring-support}\,(c) remains true by definition of
  \(\mathcal H\);
\item \(F\) has no sparse block representation in the sense of
  \cref{plt:def:alg-representation} at any precision \(N\ge1\); consequently no
  routine of \cref{plt:sec:algorithms} accepts \(F\) as input, and the
  precision calculus of \cref{plt:thm:alg-precision-calculus} and the cost
  model of \cref{plt:prop:alg-convolution-cost} have no meaning for it.
\end{enumerate}
\end{proposition}

\begin{proof}
Take \(F=1+t+\sum_{k\ge1}t^{\gamma_{k}}\) with
\(\gamma_{k}\in\Gamma\cap(0,1)\) strictly increasing, as constructed in the
proof of \cref{plt:thm:ext-lf-completion}\,(iv); then \(\ordt F=0\) and every
block is \(1\).
(a)~All the \(\gamma_{k}\) and \(0\) are \(<1\le N\), so
\(\operatorname{Exp}(\Trunc{F}{N})\supseteq\SetOf{\gamma_{k}:k\ge1}\) is
infinite.
(b)~An enumeration \(s_{0}\succdom s_{1}\succdom\cdots\) of order type at most
\(\omega\) exhausting \(\operatorname{Exp}(F)\) would be an increasing
enumeration of a set whose order type is at most \(\omega\).  But
\(\operatorname{Exp}(F)\) has order type \(\omega+1\): the element \(1\) is
strictly greater than each of the infinitely many elements
\(0<\gamma_{1}<\gamma_{2}<\cdots\), so in any exhausting increasing enumeration
it would have to be preceded by infinitely many terms and could occupy no
finite position.  Hence \cref{plt:thm:ring-support}\,(d) fails.  By
\cref{plt:prop:ext-leftfinite}\,(i) the correct replacement for (d) in a
general \(\Gamma\) is not ``order type at most \(\omega\)'' but
left-finiteness, which additionally demands cofinality of the enumeration.
(c)~A sparse block representation is a \emph{finitely supported} map
\(v\) together with a cutoff \(N\), and by
\cref{plt:lem:alg-denotation-coset} its denotation is the coset
\(\sum_{\gamma<N}v(\gamma)t^{\gamma}+\mathfrak F_{N}\).  If this coset
contained \(F\) then \(\Trunc{F}{N}\) would equal the finitely supported
\(\sum_{\gamma<N}v(\gamma)t^{\gamma}\), contradicting (a).
\end{proof}

% ed.: S1 ("The arithmetic formulas from Part II therefore survive almost
% unchanged"), S3 ("The differential rule is again ...  Near-identity
% composition remains controlled") and S5 ("Near-identity composition and
% reversion continue to work when the support conditions ensure coefficientwise
% finiteness") all extend Parts I-IV to Hahn supports without recording that
% the finite bookkeeping of Part VI -- truncation, sparse representation,
% precision calculus -- requires strictly more than well-basedness.
% \cref{plt:prop:ext-hahn-trunc-fails} isolates what is lost, and
% \cref{plt:thm:ext-lf-completion} isolates the class where nothing is.

\begin{table}[htbp]
\centering
\caption{Transfer of Parts~I--IV and~VI to a Hahn power--log scale
\(\mathcal H(C,\Gamma)\) with \(\Gamma\subseteq\RealNumbers\) dense.
``Verbatim'' means the statement and its proof go through with
\(n\in\IntegerNumbers\) read as \(\gamma\in\Gamma\).}
\label{plt:tab:ext-survival}
\small
\begin{tabular}{@{}p{0.42\textwidth}p{0.13\textwidth}p{0.38\textwidth}@{}}
\toprule
Statement & Status & What changes \\
\midrule
Domain, valuation laws
(\cref{plt:thm:ring-valuation-laws}) & verbatim & --- \\
Unit criterion (\cref{plt:thm:div-unit-criterion}) & verbatim & --- \\
Reciprocal of a unit (\cref{plt:thm:div-reciprocal}) & modified & the
  block-by-block recurrence presupposes a successor exponent; replace it by the
  geometric series, summable by \cref{plt:lem:ext-neumann}\,(iii) \\
\(t\)-adic completeness (\cref{plt:thm:trunc-completeness}) & verbatim &
  \cref{plt:thm:ext-lf-completion}\,(iii) \\
Support characterisation, clause (d)
  (\cref{plt:thm:ring-support}) & \textbf{false} &
  \cref{plt:prop:ext-hahn-trunc-fails}\,(b) \\
Block and repeated-derivative formulas
  (\cref{plt:prop:dif-block}, \cref{plt:thm:dif-repeated}) & verbatim &
  the shift is \(\partial_{L}-\gamma\) with \(\gamma\in\Gamma\); the successive
  shifts \(\gamma,\gamma+1,\ldots\) stay integrally spaced \\
Sharp valuation of \(\DX\) (\cref{plt:thm:dif-valuation}), kernel & verbatim
  & --- \\
Antiderivatives (\cref{plt:thm:dif-antiderivative}) & verbatim & the resonant
  exponent is still the single value \(\gamma=0\) \\
Exponential boundary (\cref{plt:thm:dif-exp-boundary}) & modified & \(m\) may
  be any element of \(\Gamma\) (\cref{plt:cor:dif-exp-hahn}) \\
Taylor formula, near-identity group, reversion
  (\cref{plt:thm:grp-group}, \cref{plt:thm:rev-existence}) & verbatim &
  \cref{plt:cor:ext-lf-group} \\
Triangular reversal (\cref{plt:thm:rev-triangular}) & modified & ``the next
  block'' is undefined when \(\Gamma\) is dense; one must fix a cofinal
  sequence in \(\Gamma\) and cancel residuals along it \\
Lagrange--B\"urmann index range \(r\le N+1\)
  (\cref{plt:prop:lag-range}) & \textbf{open} & the bound is an integrality
  statement; see \cref{plt:rmk:ext-open-burmann} \\
Sparse representation, precision calculus, cost model
  (\cref{plt:def:alg-representation},
   \cref{plt:thm:alg-precision-calculus},
   \cref{plt:prop:alg-convolution-cost}) & \textbf{false} & restored verbatim
  on \(\mathcal H_{\mathrm{lf}}(C,\Gamma)\) \\
\bottomrule
\end{tabular}
\end{table}

\section{Puiseux-log classes, and the union over denominators}
\label{plt:sec:ext-puiseux}

The Puiseux-log ring \(\K[L]((t^{1/s}))\) of \cref{plt:def:ring-puiseux} and
the ramified tower of \cref{plt:prop:ring-puiseux-tower} are the minimal
enlargement forced by a rational leading exponent
(\cref{plt:cor:lag-enlargement}).  At a fixed denominator nothing is lost at
all, and the reason is worth stating as a theorem rather than as the folklore
``the derivative formula generalizes without change'' of S1, S3 and S5: the
ring is not merely similar to \(\PLlau\), it \emph{is} \(\PLlau\) in a
different large variable.

\begin{proposition}[A fixed denominator is a change of variable, not an
enlargement]
\label{plt:prop:ext-puiseux-transport}
Fix \(s\in\PositiveIntegers\) and set
\[
  Y\DefinitionEquals X^{1/s},
  \qquad
  \tau\DefinitionEquals Y^{-1}=t^{1/s},
  \qquad
  \Lambda\DefinitionEquals\log Y=s^{-1}L .
\]
Then the \(\K\)-linear map sending \(t^{r/s}L^{j}\) to \(s^{j}\tau^{r}
\Lambda^{j}\) is an isomorphism of ordered \(\K\)-algebras
\[
  \Phi_{s}:\ \K[L]\bigl((t^{1/s})\bigr)\ \longrightarrow\
  \K[\Lambda]((\tau)),
\]
onto the polynomial--logarithmic Laurent ring of \cref{plt:def:ring-four-classes}
built from the generators \(\tau,\Lambda\), and
\begin{equation}\label{plt:eq:ext-puiseux-derivations}
  \Phi_{s}\compose\DX\compose\Phi_{s}^{-1}
  =s^{-1}\tau^{\,s-1}\,D_{Y},
\end{equation}
where \(D_{Y}=\tau\partial_{\Lambda}-\tau^{2}\partial_{\tau}\) is the
distinguished derivation of the target.  Consequently every statement of
Parts~I--IV that mentions only the ring structure, the order, the valuation,
the truncation calculus and the derivation up to a unit factor holds verbatim
in \(\K[L]((t^{1/s}))\).  In particular the unit criterion, the reciprocal
recurrence, the completeness theorem, the sparse representation and the
precision calculus transfer with no change whatever, and
\begin{equation}\label{plt:eq:ext-puiseux-block}
  \DX\bigl(t^{\alpha}p(L)\bigr)
  =t^{\alpha+1}\bigl(\partial_{L}p-\alpha p\bigr),
  \qquad \alpha\in s^{-1}\IntegerNumbers .
\end{equation}
\end{proposition}

\begin{proof}
\(\Phi_{s}\) is the relabelling of generators \(t^{1/s}\mapsto\tau\),
\(L\mapsto s\Lambda\); it is bijective on monomials, hence a \(\K\)-algebra
isomorphism, and it preserves the dominance order because it preserves the
exponent pair \((r,j)\) and the order is lexicographic in that pair.  For
\eqref{plt:eq:ext-puiseux-derivations} it suffices to compare on the two
generators.  On the one hand \(\DX\tau=\DX t^{1/s}=-s^{-1}t^{1/s+1}
=-s^{-1}\tau^{\,s+1}\), and \(s^{-1}\tau^{\,s-1}D_{Y}\tau
=-s^{-1}\tau^{\,s-1}\tau^{2}=-s^{-1}\tau^{\,s+1}\).  On the other hand
\(\DX\Lambda=s^{-1}\DX L=s^{-1}t=s^{-1}\tau^{\,s}\), and
\(s^{-1}\tau^{\,s-1}D_{Y}\Lambda=s^{-1}\tau^{\,s-1}\tau=s^{-1}\tau^{\,s}\).
Both sides of \eqref{plt:eq:ext-puiseux-derivations} are derivations of
\(\K[\Lambda]((\tau))\) that raise the \(\tau\)-order, hence are continuous for
the \(\tau\)-adic filtration; their difference is a continuous derivation
vanishing on \(\tau\) and \(\Lambda\), so by the Leibniz rule it vanishes on
\(\K[\Lambda][\tau,\tau^{-1}]\) and by continuity and
\cref{plt:thm:trunc-completeness} on all of \(\K[\Lambda]((\tau))\).  Formula
\eqref{plt:eq:ext-puiseux-block} is
\eqref{plt:eq:dif-scale-derivation} for
\(\Gamma=s^{-1}\IntegerNumbers\), or equivalently
\eqref{plt:eq:ext-puiseux-derivations} read backwards.
\end{proof}

\begin{remark}[The one thing the change of variable does \emph{not} transport]
\label{plt:rmk:ext-puiseux-not-near-identity}
\(\Phi_{s}\) is an isomorphism of ordered differential rings up to a unit, but
it does not carry near-identity maps in \(X\) to near-identity maps in \(Y\)
identically: if \(F=X+A\) with \(A\in\K[L][[t^{1/s}]]\), then in the variable
\(Y\) the same map reads
\[
  Y\longmapsto\bigl(Y^{s}+A\bigr)^{1/s}
  =Y\bigl(1+\tau^{s}A\bigr)^{1/s}
  =Y+s^{-1}\tau^{\,s-1}A+\FormalBigOAt{\tau^{2s-1}}{\tau},
\]
which is again near-identity but with a perturbation of outer order \(s-1\)
higher.  Reversion therefore transports through the conjugation
\(F\mapsto j_{s}\compose F\compose\rev{j_{s}}\) with \(j_{s}(X)=X^{1/s}\), not
through \(\Phi_{s}\) alone.  This is the formal counterpart of the fact that
tangency to the identity is a property of a map together with a coordinate.
\end{remark}

The union over denominators behaves quite differently, and here the sources are
silent about the essential point.

\begin{theorem}[The Puiseux-log union is separated but not complete]
\label{plt:thm:ext-puiseux-incomplete}
Let
\[
  \mathfrak P\DefinitionEquals
  \K[L]\bigl((t^{1/\infty})\bigr)
  =\bigcup_{s\in\PositiveIntegers}\K[L]\bigl((t^{1/s})\bigr)
\]
be the ramified union of \cref{plt:prop:ring-puiseux-tower}, viewed inside the
power--log scale \(\mathcal H(\K[L],\RationalNumbers)\).  Then:
\begin{enumerate}[label=\textup{(\roman*)}]
\item \(\mathfrak P\subseteq
  \mathcal H_{\mathrm{lf}}(\K[L],\RationalNumbers)\), and \(\mathfrak P\)
  contains every element with finite exponent set;
\item \(\mathfrak P\) is separated for the \(\RationalNumbers\)-indexed
  filtration, but it is \emph{not} complete: the sequence of partial sums of
  \begin{equation}\label{plt:eq:ext-puiseux-witness}
    F\DefinitionEquals\sum_{k\ge2}t^{\,k-1/k}
  \end{equation}
  is Cauchy in \(\mathfrak P\) and has no limit in \(\mathfrak P\).  Hence
  \cref{plt:thm:trunc-completeness} is false for \(\mathfrak P\);
\item the closure of \(\mathfrak P\) inside
  \(\mathcal H(\K[L],\RationalNumbers)\) is exactly
  \(\mathcal H_{\mathrm{lf}}(\K[L],\RationalNumbers)\); that is, the
  left-finite Hahn ring over \(\RationalNumbers\) is the completion of the
  Puiseux-log union.
\end{enumerate}
\end{theorem}

\begin{proof}
(i)  An element of \(\K[L]((t^{1/s}))\) has exponent set contained in
\(s^{-1}\IntegerNumbers\) and bounded below, hence left-finite by
\cref{plt:thm:ext-support-hierarchy}\,(v) applied to the cyclic group
\(s^{-1}\IntegerNumbers\), or directly.  Conversely a finite set of rationals
has a common denominator \(s\), so any element with finite exponent set lies in
\(\K[L]((t^{1/s}))\).

(ii)  Separatedness is inherited from
\cref{plt:thm:ext-lf-completion}\,(iii).  Put \(F_{K}=\sum_{k=2}^{K}
t^{\,k-1/k}\); then \(F_{K}\in\K[L]((t^{1/K!}))\subseteq\mathfrak P\), and for
\(K'>K\) we have \(\ordt(F_{K'}-F_{K})=(K+1)-\frac1{K+1}>K\), so
\((F_{K})\) is Cauchy.  Its limit in the complete ring
\(\mathcal H(\K[L],\RationalNumbers)\) is \(F\) of
\eqref{plt:eq:ext-puiseux-witness}, whose exponent set
\(\SetOf{k-\frac1k:k\ge2}\) is strictly increasing and unbounded, hence
left-finite, so \(F\in\mathcal H_{\mathrm{lf}}\).  But \(F\notin\mathfrak P\):
if \(F\in\K[L]((t^{1/s}))\) then \(k-\frac1k\in s^{-1}\IntegerNumbers\) for all
\(k\ge2\), that is \(s/k\in\IntegerNumbers\), that is \(k\mid s\), which fails
for all \(k>s\).  By separatedness a limit is unique, so \((F_{K})\) has no
limit in \(\mathfrak P\).

(iii)  By (i), \(\K[L][t^{\RationalNumbers}]\subseteq\mathfrak P\subseteq
\mathcal H_{\mathrm{lf}}\).  Taking closures in
\(\mathcal H(\K[L],\RationalNumbers)\) and using
\cref{plt:thm:ext-lf-completion}\,(iv), which identifies the closure of
\(\K[L][t^{\RationalNumbers}]\) with \(\mathcal H_{\mathrm{lf}}\), and the fact
that \(\mathcal H_{\mathrm{lf}}\) is closed, all three closures coincide.
\end{proof}

\begin{remark}[Why this matters operationally]
\label{plt:rmk:ext-denominator-growth}
Every individual computation in the Puiseux tower happens at some finite
denominator, and \cref{plt:prop:ext-puiseux-transport} says that at a fixed
denominator the whole apparatus of Parts~I--IV applies unchanged.  What
\cref{plt:thm:ext-puiseux-incomplete} adds is that a \emph{limit} of such
computations may leave the tower.  An implementation that stores a single
denominator per object is therefore sound for any finite sequence of
operations and unsound as a model of the completed class; the completed class
is \(\mathcal H_{\mathrm{lf}}\), on which
\cref{plt:cor:ext-lf-group} guarantees closure under near-identity composition
and reversion but no denominator bound.  Bounding the denominator growth of an
explicit iteration is the content of
\cref{plt:rmk:ext-open-denominators}.
\end{remark}

\subsection{An obstruction no exponent group repairs}

\begin{proposition}[No power--log scale contains a square root of \(L\)]
\label{plt:prop:ext-sqrtL}
Let \(\mathcal H=\mathcal H(C,\Gamma)\) be a power--log scale in the sense of
\cref{plt:def:dif-scale}, so that
\(C\in\SetOf{\K'[L],\K'(L),\K'((L^{-1}))}\).  Then there is no
\(F\in\mathcal H\) with \(F^{2}=L\).  This holds for every exponent group
\(\Gamma\), in particular for \(\Gamma=\RealNumbers\).
\end{proposition}

\begin{proof}
Suppose \(F^{2}=L\).  Since \(\mathcal H\) is a domain and \(\ordt\) is
additive (\cref{plt:lem:dif-scale-basic}(i)), \(2\ordt F=\ordt L=0\), so
\(\ordt F=0\); as \(F\ne0\) and its exponent set is well based, the block
\(p_{0}=[t^{0}]F\) is defined and nonzero, and comparing blocks at \(t^{0}\)
gives \(p_{0}^{2}=L\) in \(C\).  It therefore suffices to show that \(L\) is
not a square in any of the three coefficient rings.

For \(C=\K'(L)\) (and hence for the subring \(\K'[L]\)): write
\(p_{0}=u/v\) with \(u,v\in\K'[L]\) coprime and \(v\ne0\); then
\(u^{2}=Lv^{2}\).  The polynomial \(L\) is irreducible in the unique
factorisation domain \(\K'[L]\), and its multiplicity in \(u^{2}\) is even
while its multiplicity in \(Lv^{2}\) is odd, a contradiction.

For \(C=\K'((L^{-1}))\): write \(p_{0}=\sum_{j\le d}c_{j}L^{j}\) with
\(c_{d}\ne0\).  Then \(p_{0}^{2}=c_{d}^{2}L^{2d}+(\text{lower powers})\), so
comparing the top exponent with that of \(L\) gives \(2d=1\), impossible for
\(d\in\IntegerNumbers\).
\end{proof}

\begin{remark}[Which axis repairs it]
\label{plt:rmk:ext-sqrtL-axis}
\(\sqrt L\) does exist in the rank-two Hahn field
\(\mathcal H_{1}(\K,\Gamma)\) of \cref{plt:def:grp-logtower} whenever
\(\tfrac12\in\Gamma\), since there \(L_{1}^{1/2}\) is a monomial.  The point of
\cref{plt:prop:ext-sqrtL} is that the four classes of Part~I and all their
Puiseux and Hahn refinements share one structural feature --- the exponent of
\(L\) inside a block runs over \(\IntegerNumbers\), never over a larger group
--- and that no amount of refinement along the \(t\)-exponent axis touches it.
The enlargement required is a ramification of the \emph{monomial group} in
which \(L\) itself carries fractional exponents; algebraically this is the step
from a scale with a coefficient \emph{ring} to a genuine rank-two Hahn field.
Concretely, \(\sqrt{\log X}\) arises the moment one inverts a map with a
half-integral logarithmic exponent.  For \(Y=X(\log X)^{1/2}\) with
\(X\to+\infty\), taking logarithms gives
\(\log Y=\log X+\tfrac12\log\log X\); dividing that relation by \(\log Y\) and
taking logarithms gives
\(\log\log X-\log\log Y
 =\log\bigl(1-\tfrac12\log\log X/\log Y\bigr)
 =\BigOAt{\log\log Y/\log Y}{Y\to+\infty}\), whence
% ed.: the draft recorded only "+ o(1)" here, which is too weak to produce the
% error term displayed below: an additive o(1) inside \log X contributes a
% relative error o(1/\log Y), larger than (\log\log Y)^2(\log Y)^{-2}.  The
% sharper form is what the two-line computation actually gives, and it is what
% the displayed remainder needs.
\[
  \log X=\log Y-\tfrac12\log\log Y
    +\BigOAt{\log\log Y/\log Y}{Y\to+\infty},
\]
and hence
\[
  X=Y\,(\log X)^{-1/2}
   =Y\,(\log Y)^{-1/2}
     \Bigl(1+\tfrac14\frac{\log\log Y}{\log Y}
     +\BigOAt{(\log\log Y)^{2}(\log Y)^{-2}}{Y\to+\infty}\Bigr),
\]
whose leading monomial \(Y(\log Y)^{-1/2}\) has a half-integral exponent on the
first logarithm; the expansion lives in
\(\mathcal H_{2}\bigl(\K,\tfrac12\IntegerNumbers\bigr)\) and in no power--log
scale.  The computation is the shape of \cref{plt:thm:lag-log-normal-form} at
% ed.: the draft concluded that the display "is exactly the sigma = 1/2
% instance" of a theorem stated only for sigma in N.  That is a claim about a
% proof belonging to another part, and it is not made here: the display above
% is a self-contained computation, and the agreement with the theorem is
% recorded as a check, not used as an input.
\(\rho=1\), \(\sigma=\tfrac12\); that theorem is stated for
\(\sigma\in\NaturalNumbers\), and we do not widen its hypothesis.  The display
above is proved where it stands, and it agrees with
\eqref{plt:eq:lag-X-expansion} at \(c=1\), \(\rho=1\), \(\sigma=\tfrac12\).
\end{remark}

% ed.: \cref{plt:thm:lag-log-normal-form} restricts to sigma in N although its
% proof uses only that sigma is a scalar of K; the half-integral instance is
% needed here and is verified directly above rather than by widening the
% hypothesis of a theorem belonging to another part.  The displayed
% coefficient 1/4 agrees with \eqref{plt:eq:lag-X-expansion} at c=1, rho=1,
% sigma=1/2.

\section{Nested logarithms, and the residue that forces each new level}
\label{plt:sec:ext-depth}

The nested-logarithm tower \(L_0=X\), \(L_{j+1}=\log L_j\), its monomial groups
\(\mathcal M_r(\Gamma)\) and its Hahn ambients
\(\mathcal H_r(\K,\Gamma)\) were introduced in
\cref{plt:def:grp-logtower}; the tower derivation \(\mathcal D_r\), the minimal
Laurent coefficient ring \(\mathcal L_r\), the near-identity group at depth
\(r\) and the depth-dependent unit criterion are
\cref{plt:prop:grp-tower-derivation}, \cref{plt:thm:grp-multilog-group} and
\cref{plt:prop:grp-multilog-units}.  Two facts remain to be proved, and they
are the ones that make the tower a genuine hierarchy rather than a notational
convenience: the derivation extends to the full Hahn ambient, and each level is
forced by an obstruction of residue type that no enlargement along any other
axis removes.

Throughout, \(\IntegerNumbers\subseteq\Gamma\subseteq\RealNumbers\) is an
ordered group equipped, as in \cref{plt:def:dif-scale}, with an embedding
\(\Gamma\hookrightarrow\K\) of additive groups.  We extend
\cref{plt:def:grp-logtower} to \(r=0\) by the same formula, so that
\(\mathcal M_0(\Gamma)=\SetOf{X^{\alpha_0}:\alpha_0\in\Gamma}\) and
\(\mathcal H_0(\K,\Gamma)\) is the logarithm-free Hahn field of real powers of
\(X\).  We write a monomial of \(\mathcal M_r(\Gamma)\) as
\begin{equation}\label{plt:eq:ext-monomial}
  \mathfrak m=\prod_{j=0}^{r}L_j^{\alpha_j(\mathfrak m)},
  \qquad \alpha_j(\mathfrak m)\in\Gamma,
  \qquad
  \mathfrak n_j\DefinitionEquals L_0L_1\cdots L_j\quad(0\le j\le r).
\end{equation}
Each \(\alpha_j:\mathcal M_r(\Gamma)\to\Gamma\) is a group homomorphism, and
\(\mathfrak n_j\succdom1\), so \(\mathfrak n_j^{-1}\precdom1\); moreover
\(\mathfrak n_0^{-1}=t\succdom\mathfrak n_1^{-1}\succdom\cdots\succdom
\mathfrak n_r^{-1}\).

\begin{lemma}[The tower derivation on the Hahn ambient]
\label{plt:lem:ext-tower-derivation}
For \(0\le j\le r\) let \(\vartheta_j\) act on
\(F=\sum_{\mathfrak m\in S}c_{\mathfrak m}\mathfrak m\in\mathcal H_r(\K,\Gamma)\)
by \(\vartheta_jF=\sum_{\mathfrak m\in S}\alpha_j(\mathfrak m)
c_{\mathfrak m}\mathfrak m\), and set
\begin{equation}\label{plt:eq:ext-tower-derivation}
  \partial_r\DefinitionEquals\sum_{j=0}^{r}\mathfrak n_j^{-1}\vartheta_j .
\end{equation}
Then:
\begin{enumerate}[label=\textup{(\roman*)}]
\item each \(\vartheta_j\) is a well-defined \(\K\)-linear derivation of
  \(\mathcal H_r(\K,\Gamma)\), and so is \(\partial_r\);
\item \(\partial_r L_j=\mathfrak n_{j-1}^{-1}\) for \(1\le j\le r\) and
% ed.: the draft pointed at an equation label plt:eq:grp-tower-derivative that
% the volume does not carry; the statement it needs is the enclosing
% proposition, which does exist.
  \(\partial_rL_0=1\), in agreement with
  \cref{plt:prop:grp-tower-derivation}; and for
  \(\mathfrak m\in\mathcal M_r(\Gamma)\),
  \begin{equation}\label{plt:eq:ext-log-derivative}
    \frac{\partial_r\mathfrak m}{\mathfrak m}
    =\sum_{j=0}^{r}\alpha_j(\mathfrak m)\,\mathfrak n_j^{-1};
  \end{equation}
\item on \(\PLlau\subseteq\mathcal H_1(\K,\IntegerNumbers)\), \(\partial_1\)
  is \(\DX\); on a power--log scale \(\mathcal H(\K[L],\Gamma)\) it is
  \eqref{plt:eq:dif-scale-derivation}; and on \(\mathcal L_r((t))\) it is
  \eqref{plt:eq:grp-tower-block}.
\end{enumerate}
\end{lemma}

\begin{proof}
(i)  \(\vartheta_jF\) has support inside that of \(F\), so it is well defined.
It is a derivation: the coefficient of \(\mathfrak p\) in \(\vartheta_j(FG)\)
is \(\alpha_j(\mathfrak p)\sum_{\mathfrak m\mathfrak n=\mathfrak p}
c_{\mathfrak m}d_{\mathfrak n}\), while the coefficient of \(\mathfrak p\) in
\((\vartheta_jF)G+F(\vartheta_jG)\) is
\(\sum_{\mathfrak m\mathfrak n=\mathfrak p}\bigl(\alpha_j(\mathfrak m)
+\alpha_j(\mathfrak n)\bigr)c_{\mathfrak m}d_{\mathfrak n}\); these agree
because \(\alpha_j\) is a homomorphism, and both sums are finite by
\cref{plt:lem:ext-neumann}\,(ii).  Multiplying a derivation by a fixed ring
element again gives a derivation, and a finite sum of derivations is a
derivation, so \(\partial_r\) is one.  Its support lies in
\(\bigcup_{j}\mathfrak n_j^{-1}S\), a finite union of translates of a well-based
set, and each monomial receives at most \(r+1\) contributions, one per \(j\),
because \(\mathfrak m\mapsto\mathfrak m\mathfrak n_j^{-1}\) is injective.

(ii)  For \(\mathfrak m=L_j\) one has \(\alpha_i(L_j)=\delta_{ij}\), so
\(\partial_rL_j=\mathfrak n_j^{-1}L_j=(L_0\cdots L_{j-1})^{-1}
=\mathfrak n_{j-1}^{-1}\) for \(j\ge1\), and \(\partial_rL_0=L_0^{-1}L_0=1\).
Formula \eqref{plt:eq:ext-log-derivative} is
\eqref{plt:eq:ext-tower-derivation} applied to the single monomial
\(\mathfrak m\).

(iii)  For \(\mathfrak m=t^{\gamma}L^{k}=L_0^{-\gamma}L_1^{k}\),
\eqref{plt:eq:ext-log-derivative} gives
\(\partial_1\mathfrak m=\mathfrak m\bigl(-\gamma L_0^{-1}+k(L_0L_1)^{-1}\bigr)
=t^{\gamma+1}\bigl(kL^{k-1}-\gamma L^{k}\bigr)
=t^{\gamma+1}(\partial_L-\gamma)L^{k}\), which is
\eqref{plt:eq:dif-scale-derivation}, and for \(\gamma=n\in\IntegerNumbers\)
this is \(\DX\).  The identification with
\eqref{plt:eq:grp-tower-block} is the same computation with
\(\mathfrak m=t^{n}L_1^{a_1}\cdots L_r^{a_r}\).
\end{proof}

The next theorem is the exact reason the tower cannot be truncated.  It is the
depth-\(r\) form of \cref{plt:cor:dif-loglog}, which is the case \(r=1\): there
the obstruction was the residue \([t^{1}L^{-1}]\), and
\(t^{1}L^{-1}=\mathfrak n_1^{-1}\).

\begin{theorem}[Each level of the tower is forced by an unmatched residue]
\label{plt:thm:ext-tower-strict}
Let \(r\ge0\) and let \(\mathcal H_r=\mathcal H_r(\K,\Gamma)\) with
\(\IntegerNumbers\subseteq\Gamma\subseteq\RealNumbers\).  Then
\begin{equation}\label{plt:eq:ext-residue}
  \bigl[\mathfrak n_r^{-1}\bigr]\,\partial_rF=0
  \qquad\text{for every }F\in\mathcal H_r,
\end{equation}
where \([\mathfrak n]G\) denotes the coefficient of the monomial
\(\mathfrak n\) in \(G\).  Consequently:
\begin{enumerate}[label=\textup{(\roman*)}]
\item \(\mathfrak n_r^{-1}=\dfrac1{L_0L_1\cdots L_r}\) has no antiderivative in
  \(\mathcal H_r\), so \(\partial_r\) is not surjective and
  \cref{plt:thm:dif-antiderivative} fails at every depth once the ambient is
  the full Hahn field;
\item if \((\mathcal T,\partial)\) is any differential field extension of
  \(\mathcal H_r\) whose derivation restricts to \(\partial_r\), and
  \(\Lambda\in\mathcal T\) satisfies \(\partial\Lambda=\mathfrak n_r^{-1}\)
  --- as \(L_{r+1}=\log L_r\) does --- then \(\Lambda\notin\mathcal H_r\).  In
  particular \(L_{r+1}\notin\mathcal H_r(\K,\Gamma)\) for every \(\Gamma\), and
  \[
    \mathcal H_0(\K,\Gamma)\subsetneq\mathcal H_1(\K,\Gamma)\subsetneq
    \mathcal H_2(\K,\Gamma)\subsetneq\cdots
  \]
  is a strictly increasing chain of differential subfields;
\item no \(\mathcal H_r(\K,\Gamma)\) is closed under the logarithm of its
  positive elements: \(L_r\in\mathcal H_r\) but \(\log L_r\notin\mathcal H_r\).
\end{enumerate}
\end{theorem}

\begin{proof}
Fix \(F=\sum_{\mathfrak m\in S}c_{\mathfrak m}\mathfrak m\in\mathcal H_r\).  By
\eqref{plt:eq:ext-tower-derivation}, the monomial \(\mathfrak n_r^{-1}\)
receives, for each index \(j\in\SetOf{0,\ldots,r\,}\), exactly one possible
contribution, namely from the unique monomial \(\mathfrak m_j\) with
\(\mathfrak m_j\mathfrak n_j^{-1}=\mathfrak n_r^{-1}\), that is
\[
  \mathfrak m_j=\mathfrak n_j\mathfrak n_r^{-1}
  =\frac{L_0L_1\cdots L_j}{L_0L_1\cdots L_r}
  =\prod_{i=j+1}^{r}L_i^{-1},
\]
the empty product for \(j=r\) being \(1\).  Hence
\[
  \bigl[\mathfrak n_r^{-1}\bigr]\partial_rF
  =\sum_{j=0}^{r}\alpha_j(\mathfrak m_j)\,c_{\mathfrak m_j},
\]
a finite sum, in which \(c_{\mathfrak m_j}\) is read as \(0\) when
\(\mathfrak m_j\notin S\).  But \(\mathfrak m_j\) involves only the letters
\(L_{j+1},\ldots,L_r\), so its \(L_j\)-exponent is
\(\alpha_j(\mathfrak m_j)=0\) for every \(j\), and every term vanishes.  This
proves \eqref{plt:eq:ext-residue}.

(i)  If \(\partial_rF=\mathfrak n_r^{-1}\) then the left side has
\(\mathfrak n_r^{-1}\)-coefficient \(0\) and the right side has
\(\mathfrak n_r^{-1}\)-coefficient \(1\).

(ii)  If \(\Lambda\in\mathcal H_r\) then
\(\partial\Lambda=\partial_r\Lambda\) has vanishing
\(\mathfrak n_r^{-1}\)-coefficient by \eqref{plt:eq:ext-residue}, contradicting
\(\partial\Lambda=\mathfrak n_r^{-1}\).  That \(L_{r+1}\) has this derivative is
\cref{plt:prop:grp-tower-derivation} at index \(r+1\), rewritten with
\(\mathfrak n_r=L_0L_1\cdots L_r\).  Inclusion
\(\mathcal H_r\subseteq\mathcal H_{r+1}\) holds because
\(\mathcal M_r(\Gamma)\) is the subgroup of \(\mathcal M_{r+1}(\Gamma)\) with
\(\alpha_{r+1}=0\), the induced order agrees, and well-basedness and the
convolution product are inherited; strictness is the membership
\(L_{r+1}\in\mathcal H_{r+1}\setminus\mathcal H_r\).  Each inclusion commutes
with the derivations, since \(\partial_{r+1}\) restricted to monomials with
\(\alpha_{r+1}=0\) is \(\partial_r\).

(iii)  Immediate from (ii) with \(\Lambda=L_{r+1}\).
\end{proof}

\begin{corollary}[Depth is not bounded, and the union is the smallest
logarithmically closed ambient]
\label{plt:cor:ext-depth-unbounded}
Put \(\mathcal H_{\infty}(\K,\Gamma)=\bigcup_{r\ge0}\mathcal H_r(\K,\Gamma)\),
a directed union of differential fields, hence a differential field.  Then
\(\mathcal H_{\infty}\) contains \(L_j\) for every \(j\), no proper member
\(\mathcal H_r\) of the chain does, and any subfield of an ambient differential
field that contains \(X\) and is closed under \(\log\) on the elements
\(L_0,L_1,L_2,\ldots\) contains all of them.  In particular a bound on
logarithmic depth is never a theorem about a class; it is a hypothesis about a
particular computation.
\end{corollary}

\begin{proof}
A directed union of subfields of a fixed field is a subfield, and the
derivations agree on overlaps by
\cref{plt:thm:ext-tower-strict}\,(ii); \(L_j\in\mathcal H_j\subseteq
\mathcal H_{\infty}\) while \(L_{r+1}\notin\mathcal H_r\) by the same theorem.
The last assertion is an induction on \(j\) starting from \(L_0=X\).
\end{proof}

\begin{warningbox}[title={A depth bound is a hypothesis, not a conclusion}]
\Cref{plt:prop:lag-loglog-forced} shows one logarithm in a leading monomial
forcing one new level; \cref{plt:thm:ext-tower-strict} shows that no level is
ever redundant.  Neither result bounds the depth needed by an arbitrary
composite of admissible operations, and this volume proves no such bound.  What
is proved is that the classification of scale changes
(\cref{plt:thm:grp-scale-classification}) raises the depth of a single
power--logarithm substitution by at most one; iterating it \(k\) times can
therefore be guaranteed only within depth \(r+k\).  Whether that is sharp for
every \(k\) is \cref{plt:rmk:ext-open-depth}.
\end{warningbox}

\section{Multiseries: nesting, the fast variable, and the completion order}
\label{plt:sec:ext-multiseries}

A \emph{multiseries} is an expansion whose coefficients are themselves
expansions, in a strictly slower scale, iterated to a finite depth.  Formally
this is the iterated Laurent construction of \cref{plt:def:div-iterated},
carried up the logarithmic tower.

\begin{definition}[Finite-depth multiseries field]
\label{plt:def:ext-multiseries}
For \(r\ge0\) put
\begin{equation}\label{plt:eq:ext-multiseries}
  \mathfrak J_r
  \DefinitionEquals
  \K\bigl((L_r^{-1})\bigr)\bigl((L_{r-1}^{-1})\bigr)\cdots
  \bigl((L_1^{-1})\bigr)\bigl((L_0^{-1})\bigr),
\end{equation}
the iterated field of formal Laurent series in \(L_0^{-1}=t\) whose
coefficients are formal Laurent series in \(L_1^{-1}\), whose coefficients are
formal Laurent series in \(L_2^{-1}\), and so on to depth \(r\); the outermost
construction carries the \emph{most dominant} scale.  Thus
\(\mathfrak J_0=\K((t))\) and \(\mathfrak J_1=\PLit\).  An element of
\(\mathfrak J_r\) is a family \((c_{m})_{m\in\IntegerNumbers^{r+1}}\) of
scalars, \(m=(m_0,\ldots,m_r)\) representing the monomial
\(L_0^{-m_0}L_1^{-m_1}\cdots L_r^{-m_r}\).
\end{definition}

\begin{theorem}[The iterated field is the Hahn field, and only in the matching
order]
\label{plt:thm:ext-iterated-hahn}
Let \(r\ge0\).
\begin{enumerate}[label=\textup{(\roman*)}]
\item \(\mathfrak J_r=\mathcal H_r(\K,\IntegerNumbers)\) as \(\K\)-algebras,
  the identification being the identity on monomials.  Explicitly, a family
  \((c_m)_{m\in\IntegerNumbers^{r+1}}\) lies in \(\mathfrak J_r\) if and only
  if its support is well ordered in the direction of increasing smallness for
  the dominance order of \cref{plt:def:grp-logtower}.
\item Let \(\pi\) be a permutation of \(\SetOf{0,1,\ldots,r}\) and let
  \(\mathfrak J_{\pi}\) be built as in \eqref{plt:eq:ext-multiseries} with the
  levels nested in the order \(\pi(0)\) outermost, \ldots, \(\pi(r)\)
  innermost.  If \(\pi\ne\operatorname{id}\) then neither of
  \(\mathfrak J_{\pi}\), \(\mathfrak J_r\) contains the other, as subsets of
  the space of families \((c_m)\).  More generally, if \(\pi\ne\pi'\) then
  neither of \(\mathfrak J_{\pi}\), \(\mathfrak J_{\pi'}\) contains the other;
  so the \((r+1)!\) nestings give \((r+1)!\) pairwise incomparable, and in
  particular pairwise distinct, subsets.
\end{enumerate}
\end{theorem}

\begin{proof}
(i)  Write \(\mathfrak m(m)=L_0^{-m_0}\cdots L_r^{-m_r}\).  By
\cref{plt:def:grp-logtower}, \(\mathfrak m(m)\succdom\mathfrak m(m')\) exactly
when \(m<m'\) lexicographically, so ``the support is well ordered in the
direction of increasing smallness'' means: every nonempty subset of the support,
read as a set of vectors \(m\in\IntegerNumbers^{r+1}\), has a lexicographically
least element.  Membership in \(\mathfrak J_r\), unfolded, says: the set of
\(m_0\) occurring is bounded below; for each such \(m_0\) the set of \(m_1\)
occurring with that \(m_0\) is bounded below; and so on through \(m_r\).  These
two conditions are equivalent.

\(\Leftarrow\): given a nonempty subset \(T\), let \(\mu_0\) be the least
\(m_0\) occurring in \(T\), which exists because the set is a nonempty subset
of \(\IntegerNumbers\) bounded below; restrict \(T\) to that value and let
\(\mu_1\) be the least \(m_1\) occurring, and so on.  The resulting vector
\((\mu_0,\ldots,\mu_r)\) lies in \(T\) and is lexicographically least.

\(\Rightarrow\): if for some \(k\) and some fixed prefix
\((\mu_0,\ldots,\mu_{k-1})\) occurring in the support the set of admissible
\(m_k\) were unbounded below, choose \(m^{(1)},m^{(2)},\ldots\) in the support
with that prefix and with \(m^{(i)}_k\) strictly decreasing; then
\(m^{(1)}>m^{(2)}>\cdots\) lexicographically, so
\(\SetOf{m^{(i)}:i\ge1}\) has no lexicographically least element.

(ii)  Since \(\pi\ne\operatorname{id}\) there are \(a<b\) in
\(\SetOf{0,\ldots,r}\) with \(\pi^{-1}(b)<\pi^{-1}(a)\): level \(b\) is nested
\emph{outside} level \(a\) in \(\mathfrak J_{\pi}\), while in
\(\mathfrak J_r\) level \(a\) is outside level \(b\).  Put
\[
  A=\sum_{n\ge0}L_a^{-n}L_b^{\,n},
  \qquad
  B=\sum_{n\ge0}L_a^{\,n}L_b^{-n},
\]
so that \(A\) has \(m_a=n\), \(m_b=-n\) and all other coordinates \(0\), and
\(B\) has \(m_a=-n\), \(m_b=n\).  Read in the order
\(0,1,\ldots,r\): for \(A\), the coordinates \(m_0,\ldots,m_{a-1}\) are \(0\),
the set of \(m_a\) is \(\NaturalNumbers\), bounded below, and once \(m_a=n\) is
fixed all later coordinates are determined; hence \(A\in\mathfrak J_r\)
(indeed \(A=\sum_n(L_b/L_a)^{n}\)).  For \(B\), the set of \(m_a\) occurring is
\(\SetOf{0,-1,-2,\ldots}\), unbounded below, so \(B\notin\mathfrak J_r\).  Read
in the order \(\pi\), the roles of \(a\) and \(b\) are exchanged and the same
computation gives \(B\in\mathfrak J_{\pi}\) and \(A\notin\mathfrak J_{\pi}\).

% ed.: the draft proved incomparability only against the identity nesting,
% while \cref{plt:rmk:ext-completion-order} asserted (r+1)! pairwise distinct
% iterated fields.  The argument already gives the general statement, because
% it uses nothing about the identity beyond the relative order of the two
% chosen levels; the last paragraph records this.
The argument used only that the two nestings order the pair \(a,b\)
oppositely, and \(\mathfrak J_r=\mathfrak J_{\operatorname{id}}\).  Two
distinct total orders on the finite set \(\SetOf{0,\ldots,r}\) must order some
pair oppositely, so for \(\pi\ne\pi'\) the same witnesses --- built from the
two levels on which \(\pi\) and \(\pi'\) disagree --- give an element of
\(\mathfrak J_{\pi}\setminus\mathfrak J_{\pi'}\) and an element of
\(\mathfrak J_{\pi'}\setminus\mathfrak J_{\pi}\).
\end{proof}

\begin{remark}[The completion order, in one sentence]
\label{plt:rmk:ext-completion-order}
\Cref{plt:thm:ext-iterated-hahn} says that iterating the Laurent construction
reproduces the Hahn field \emph{exactly when the nesting order is the reverse
of the dominance order} --- outermost variable most dominant --- and produces
an incomparable object otherwise.  The case \(r=1\) is
\cref{plt:prop:ring-iteration-order} and
\cref{plt:prop:div-iteration-order}; the general case shows that the phenomenon
is not an artefact of the rank-two situation but is present at every depth, and
that there are \((r+1)!\) pairwise distinct iterated fields on the same
monomial group, exactly one of which is the asymptotically meaningful one.  The
catalogue records the same warning as ``the iteration order is part of the
type''.
\end{remark}

\begin{remark}[Nested truncation: what a finite datum is]
\label{plt:rmk:ext-nested-truncation}
Membership in \(\mathfrak J_r\) does not make an element finitely
representable.  A single block of \(\mathfrak J_1=\PLit\) is already an
infinite Laurent series in \(L^{-1}\), so a finite datum needs \(r+1\)
cutoffs, or one declared staircase in the lexicographic monomial order; an
% ed.: the draft cited a label plt:rmk:div-embedding-warning that the volume
% does not carry; the statement being invoked is the expansion map itself.
outer cutoff alone bounds nothing.  This is the depth-\(r\) form of the caution
attached to \cref{plt:thm:div-embedding}, and it is the reason the
sparse representation of \cref{plt:def:alg-representation} fixes an inner
cutoff for the coefficient ring \(\K((L^{-1}))\).  A workable contract at depth
\(r\) is a monotone staircase: a bound \(m_0<N_0\), then
\(m_1<N_1(m_0)\), then \(m_2<N_2(m_0,m_1)\), and so on, chosen so that the
region is stable under the convolutions to be performed.  The choice is part of
the specification of the computation, not a detail of its implementation.
\end{remark}

\begin{warningbox}[title={Do not expand a rational coefficient before you must}]
The exact block \((L^{2}+1)/(L-1)\in\K(L)\) and its expansion
\(L+1+2L^{-1}+2L^{-2}+\cdots\in\K((L^{-1}))\) are different data about the same
object: the first is finite and exact, the second resolves a finer asymptotic
% ed.: the draft wrote "an identity proved after expansion need not descend".
% That contradicts the injectivity it cites in the same sentence: if
% \iota(a)=\iota(b) then a=b.  What non-surjectivity actually costs is that a
% computation carried out after expansion may leave the image, and that an
% exact cancellation in K(L) is, after expansion, visible only in the limit and
% at no finite truncation.  Both claims are restated correctly.
scale and is infinite.  By \cref{plt:thm:div-embedding} the expansion map
\(\iota\) is injective, so an identity proved after expansion \emph{does}
descend; but it is not surjective, so a computation carried out after expansion
may leave the image and return no exact rational block, and a cancellation that
is exact and finite in \(\K(L)\) survives expansion only in the limit --- at
every finite truncation it is invisible.  Expand when the requested scale
requires it, and record that you did.
\end{warningbox}

\subsection{Changing the fast variable}

The most useful practical observation about multiseries is that a nested
expansion is often an ordinary polynomial--logarithmic expansion in the right
variable.  The mechanism is already a theorem in this volume ---
\cref{plt:prop:grp-transport}, the level shift --- and the following example
shows it operating on the classical case.

\begin{example}[A logarithmic prefactor, and where the second logarithm sits]
\label{plt:ex:ext-fast-variable}
% ed.: the draft obtained the display by asserting that the proof of
% \cref{plt:thm:lag-log-normal-form}, which is stated only for sigma in N,
% works verbatim for every scalar beta.  That is a claim about a proof
% belonging to another part and is not made here; the two-line derivation is
% given instead, and the agreement with that theorem is recorded as a check.
Let \(\beta\in\K\) and invert \(Y=Z(\log Z)^{\beta}\) for \(Z\to+\infty\).
Write \(S=\log Y\), \(\Lambda=\log S\) and \(z=\log Z\), so that the relation
reads \(S=z+\beta\log z\), that is \(z=S-\beta\log z\).  Substituting
\(z=S-\beta\Lambda+w\) and expanding
\(\log z=\Lambda+\log\bigl(1-\beta\Lambda S^{-1}+wS^{-1}\bigr)\) forces
\(w=\beta^{2}\Lambda S^{-1}+\FormalBigOAt{S^{-2}}{S^{-1}}\), the reversion
being the near-identity one of \(\K[\Lambda][[S^{-1}]]\).  Since
\(\EulerE^{S}=Y\) and \(\EulerE^{-\beta\Lambda}=S^{-\beta}\), exponentiating
\(z\) gives
\begin{equation}\label{plt:eq:ext-logprefactor}
  Z=Y\,S^{-\beta}
    \Bigl(1+\beta^{2}\frac{\Lambda}{S}
      +\FormalBigOAt{S^{-2}}{S^{-1}}\Bigr).
\end{equation}
This is the shape of \cref{plt:thm:lag-log-normal-form} at \(c=1\),
\(\rho=1\), \(\sigma=\beta\), whose hypothesis \(\sigma\in\NaturalNumbers\) we
do not widen; the display is proved where it stands.
The right-hand side is \emph{not} an element of any power--log scale in the
variable \(Y\), because the second logarithm \(\Lambda\) occurs.  But the
bracket \emph{is} an element of \(\K[\Lambda][[S^{-1}]]\), that is, of the
polynomial--logarithmic power-series ring built on the fast variable
\(S=\log Y\) and its logarithm \(\Lambda=\log S\).  Under
\cref{plt:prop:grp-transport} every statement of Parts~I--IV therefore applies
to the bracket verbatim, with \(t\) read as \(S^{-1}\) and \(L\) read as
\(\Lambda\).  The only genuinely new object in
\eqref{plt:eq:ext-logprefactor} is the monomial prefactor \(YS^{-\beta}\),
which is where the depth increase is confined.
\end{example}

\begin{remark}[The pattern, and the honest limit of the pattern]
\label{plt:rmk:ext-fast-variable}
\Cref{plt:ex:ext-fast-variable} is a special case of the normalisation theorem
\cref{plt:thm:lag-reduction}: factor out a leading model whose inverse is
known, and the remainder is a near-identity problem in a scale that the leading
model itself selects.  What the pattern does \emph{not} do is remove the depth
increase; it localises it in the prefactor.  A claim of the form ``after the
right change of variable the problem is one-logarithmic'' is true of the
correction and false of the answer, and
\cref{plt:thm:ext-tower-strict} is the reason it must be false of the answer
whenever \(\beta\ne0\).
\end{remark}

\chapter{Relation with the full transseries field}
\label{plt:sec:full-transseries}

Everything in this volume happens at \emph{exponential height zero}: no
monomial of any class considered so far is an exponential of a large series.
This chapter places the polynomial--logarithmic class inside the ambient in
which exponentials are allowed, says precisely what that ambient adds, and is
careful about the division of labour.  Two things are proved here, because they
are what a user of the smaller class needs and because they sharpen results
already in the volume: an exact criterion, valid at every logarithmic depth,
for an element with a prescribed logarithmic derivative to lie in the class
(\cref{plt:thm:ext-logderiv}); and the observation that the filtered gain which
makes all of Parts~III--IV work is a property of the logarithmic scale and is
destroyed by a single exponential (\cref{plt:prop:ext-gain-fails}).  Everything
else --- the construction of the transseries field, its model theory, its
composition and summability theory --- is quoted from the literature and
explicitly \emph{not} reproved.

\section{The exponential boundary at every logarithmic depth}
\label{plt:sec:ext-exponential-boundary}

\Cref{plt:thm:dif-exp-boundary} settled which elements of \(\PLlau\) have
exponentials inside \(\PLlau\), and \cref{plt:cor:dif-exp-hahn} extended this
to a power--log scale.  Both are statements at logarithmic depth one.  The
following theorem is the depth-\(r\) statement, and it is what makes the
boundary structural rather than a feature of the particular class: it says that
in \(\mathcal H_r(\K,\Gamma)\) the logarithmic derivative of a nonzero element
is never larger than \(t\), and that its \(t\)-coefficient is forced to lie in
the exponent group.

\begin{theorem}[Logarithmic derivatives in a nested-logarithmic Hahn field]
\label{plt:thm:ext-logderiv}
Let \(r\ge0\), \(\IntegerNumbers\subseteq\Gamma\subseteq\RealNumbers\), and let
\(\mathcal H_r=\mathcal H_r(\K,\Gamma)\) with the derivation \(\partial_r\) of
\cref{plt:lem:ext-tower-derivation}.  Let \(F\in\mathcal H_r\) be nonzero with
dominant monomial \(\mathfrak m=\lm(F)\) and leading scalar \(c=\lc(F)\).
Then
\begin{equation}\label{plt:eq:ext-logderiv}
  \lm\bigl(\partial_rF\bigr)\ \preceqdom\ t\,\mathfrak m,
  \qquad
  \bigl[t\,\mathfrak m\bigr]\,\partial_rF=\alpha_0(\mathfrak m)\,c .
\end{equation}
Consequently, if \((\mathcal T,\partial)\) is a differential ring extension of
\(\mathcal H_r\) whose derivation restricts to \(\partial_r\), if
\(E\in\mathcal H_r\) is nonzero, and if \(u\in\mathcal H_r\) is nonzero with
\(\partial E=uE\), then
\begin{equation}\label{plt:eq:ext-logderiv-criterion}
  \lm(u)\preceqdom t,
  \qquad\text{and if }\lm(u)=t\ \text{then}\ \lc(u)=\alpha_0(\lm E)\in\Gamma .
\end{equation}
\end{theorem}

\begin{proof}
Write \(F=\sum_{\mathfrak n\in S}c_{\mathfrak n}\mathfrak n\) with
\(\mathfrak m=\max_{\succdom}S\).  By
\eqref{plt:eq:ext-tower-derivation}, every monomial occurring in
\(\partial_rF\) is of the form \(\mathfrak n\mathfrak n_j^{-1}\) with
\(\mathfrak n\in S\) and \(0\le j\le r\).  Since
\(\mathfrak n\preceqdom\mathfrak m\) and
\(\mathfrak n_j^{-1}\preceqdom\mathfrak n_0^{-1}=t\), each such monomial is
\(\preceqdom t\mathfrak m\), which is the first half of
\eqref{plt:eq:ext-logderiv}.  For the second half, the contributions to
\(t\mathfrak m=\mathfrak m\mathfrak n_0^{-1}\) come from the pairs
\((\mathfrak n,j)\) with \(\mathfrak n=\mathfrak m\mathfrak n_0^{-1}
\mathfrak n_j=\mathfrak m\,L_1L_2\cdots L_j\).  For \(j=0\) this is
\(\mathfrak n=\mathfrak m\), contributing \(\alpha_0(\mathfrak m)c\).  For
\(j\ge1\) the monomial \(\mathfrak m L_1\cdots L_j\) has the same
\(L_0\)-exponent as \(\mathfrak m\) and a strictly larger \(L_1\)-exponent,
hence \(\mathfrak mL_1\cdots L_j\succdom\mathfrak m\), so it does not lie in
\(S\) and contributes nothing.

For the consequence, note first that \(\mathcal H_r\) is an integral domain
with \(\lm(uE)=\lm(u)\lm(E)\) and \(\lc(uE)=\lc(u)\lc(E)\): the coefficient of
\(\lm(u)\lm(E)\) in \(uE\) is the finite sum
(\cref{plt:lem:ext-neumann}\,(ii)) over pairs
\(\mathfrak p\mathfrak q=\lm(u)\lm(E)\) with \(\mathfrak p\in
\CoefficientSupportOf{u}\), \(\mathfrak q\in\CoefficientSupportOf{E}\), and any
such pair other than \((\lm u,\lm E)\) would have
\(\mathfrak p\precdom\lm(u)\) or \(\mathfrak q\precdom\lm(E)\) and hence
\(\mathfrak p\mathfrak q\precdom\lm(u)\lm(E)\).  Now put
\(\mathfrak m=\lm(E)\).  From \(\partial E=\partial_rE=uE\) and
\eqref{plt:eq:ext-logderiv},
\(\lm(u)\mathfrak m=\lm(uE)=\lm(\partial_rE)\preceqdom t\mathfrak m\), so
\(\lm(u)\preceqdom t\).  If \(\lm(u)=t\), comparing the coefficients of
\(t\mathfrak m\) gives \(\lc(u)\lc(E)=\alpha_0(\mathfrak m)\lc(E)\), and
\(\lc(E)\ne0\).
\end{proof}

The criterion is applied by enlarging the exponent group so that the
logarithmic derivative itself becomes an element of the ambient; this costs
nothing, because \(\mathcal H_r(\K,\Gamma)\subseteq\mathcal H_r(\K,\Gamma')\)
for \(\Gamma\subseteq\Gamma'\), while the conclusion
\(\lc(u)=\alpha_0(\lm E)\) still reads the exponent of an element of
\(\mathcal M_r(\Gamma)\).

\begin{corollary}[No finite logarithmic depth admits an exponential]
\label{plt:cor:ext-no-exponentials}
Let \(\IntegerNumbers\subseteq\Gamma\subseteq\Gamma'\subseteq\RealNumbers\),
let \((\mathcal T,\partial)\) be a differential ring extension of
\(\mathcal H_r(\K,\Gamma')\) whose derivation restricts to \(\partial_r\), and
let \(E\in\mathcal T\) be nonzero with \(\partial E=uE\) for some nonzero
\(u\in\mathcal H_r(\K,\Gamma')\).  If \(E\in\mathcal H_r(\K,\Gamma)\) then
\(\lm(u)\preceqdom t\), and if \(\lm(u)=t\) then \(\lc(u)\in\Gamma\).
Consequently, for every \(r\ge0\) and every
\(\Gamma\subseteq\RealNumbers\):
\begin{enumerate}[label=\textup{(\roman*)}]
\item \(\EulerE^{cX}\notin\mathcal H_r(\K,\Gamma)\) for \(c\in\K^{\times}\);
\item \(\EulerE^{\sqrt X}\notin\mathcal H_r(\K,\Gamma)\) and
  \(X^{\log X}=\EulerE^{L^{2}}\notin\mathcal H_r(\K,\Gamma)\);
% ed.: "a \notin \Gamma" is a statement about the image of \Gamma in \K under
% the embedding of \cref{plt:def:dif-scale}; the exponent a is read as a scalar
% of \K, which is what the proof needs and what makes the enlargement of the
% exponent group in the draft's proof unnecessary.
\item no nonzero \(E\) in an ambient as above satisfies
  \(\partial E=aX^{-1}E\) with \(E\in\mathcal H_r(\K,\Gamma)\), when
  \(a\in\K\) does not lie in the image of \(\Gamma\); informally,
  \(X^{a}\notin\mathcal H_r(\K,\Gamma)\) whenever \(a\notin\Gamma\);
\item by contrast \(L\in\mathcal H_1\) and \(\EulerE^{1/X}\in\PLpow\), whose
  logarithmic derivatives are \(\mathfrak n_1^{-1}=t/L\) and \(-t^{2}\): the
  criterion is not vacuous in either direction.
\end{enumerate}
\end{corollary}

\begin{proof}
The displayed criterion is \eqref{plt:eq:ext-logderiv-criterion} applied inside
\(\mathcal H_r(\K,\Gamma')\), which contains \(E\) and \(u\); the refinement is
that \(\lm(E)\in\mathcal M_r(\Gamma)\), so \(\alpha_0(\lm E)\in\Gamma\).

(i)  Take \(\Gamma'=\Gamma\) and \(u=c\); then \(\lm(u)=1\succdom t\).
% ed.: the draft passed to \Gamma' = \Gamma + (1/2)Z without checking that
% \Gamma' is admissible, i.e. that the embedding \Gamma \hookrightarrow \K of
% \cref{plt:def:dif-scale} extends to it injectively.  It does, and the
% verification is one line, given here.
(ii)  For \(E=\EulerE^{\sqrt X}\) take \(\Gamma'=\Gamma+\tfrac12
\IntegerNumbers\), which is admissible: since \(2\cdot\tfrac12=1\in\Gamma\) and
\(\K\) has characteristic zero, sending \(\tfrac12\) to half the image of
\(1\) extends the embedding \(\Gamma\hookrightarrow\K\) to \(\Gamma'\), and the
extension is injective because \(\gamma+\tfrac12\in\Gamma\) would force
\(\tfrac12\in\Gamma\), in which case \(\Gamma'=\Gamma\).  Take
\(u=\tfrac12X^{-1/2}\), so
\(\lm(u)=L_0^{-1/2}\succdom L_0^{-1}=t\) because \(-\tfrac12>-1\).  For
\(E=\EulerE^{L^{2}}\) assume first \(r\ge1\); take \(\Gamma'=\Gamma\) and
\(u=2LX^{-1}\), so
\(\lm(u)=L_0^{-1}L_1\succdom L_0^{-1}\): the two monomials agree in the
\(L_0\)-exponent and the first has the larger \(L_1\)-exponent.  Both
contradict the criterion.  For \(r=0\) the second case follows from the case
\(r=1\), because \(\mathcal H_0(\K,\Gamma)\subseteq\mathcal H_1(\K,\Gamma)\)
by \cref{plt:thm:ext-tower-strict}\,(ii).
% ed.: the draft took \Gamma' = \Gamma + aZ here.  That enlargement is neither
% needed nor in general admissible -- for K = Q, \Gamma = Z and a irrational
% the group Z + aZ has rank two and admits no injective homomorphism into Q.
% No enlargement is required: a is a coefficient, not an exponent, and
% u = at already lies in H_r(K,\Gamma).
(iii)  Take \(\Gamma'=\Gamma\) and \(u=aX^{-1}=a\,t\), which lies in
\(\mathcal H_r(\K,\Gamma)\) because \(t\in\mathcal M_r(\Gamma)\) and \(a\) is a
scalar; then \(\lm(u)=t\) and \(\lc(u)=a\), which is not in the image of
\(\Gamma\).
(iv)  \(\lm(t/L)=L_0^{-1}L_1^{-1}\precdom L_0^{-1}=t\) and
\(\lm(-t^{2})=L_0^{-2}\precdom L_0^{-1}\); both are consistent with the
criterion, and both elements do lie in the stated classes.
\end{proof}

\begin{remark}[Relation to the depth-one boundary]
\label{plt:rmk:ext-boundary-comparison}
\Cref{plt:cor:ext-no-exponentials} does not supersede
\cref{plt:thm:dif-exp-boundary}: the latter is an exact
\emph{iff}, listing the elements whose exponential stays in \(\PLlau\), while
the former is a necessary condition of a different shape --- a constraint on
the logarithmic derivative --- whose merit is that it is uniform in the
logarithmic depth and therefore rules out an element from \emph{every} member
of the tower at once.  Read together they say: within the logarithmic world an
exponential can only be a power of \(X\) times something already there
% ed.: the draft pointed at an equation label plt:eq:dif-exp-value that the
% volume does not carry; the statement meant is the classification of which
% elements are exponentials.
(\cref{plt:cor:dif-exp-group}), and enlarging the logarithmic depth never
changes that.  The first genuinely new monomial is \(\EulerE^{X}\), and it is
new at every depth.
\end{remark}

\section{Where the class sits: the logarithmic--exponential field}
\label{plt:sec:ext-le-field}

The field \(\mathbb T\) of logarithmic--exponential transseries is built by
iterating two constructions on an ordered monomial group: adjoining
exponentials \(\EulerE^{\mathfrak p}\) of \emph{purely large} series
\(\mathfrak p\), and adjoining logarithms; and then closing under well-based
sums.  Two standard variants must be distinguished, and they correspond exactly
to the support classes of \cref{plt:def:ext-support-classes}: the
\emph{grid-based} transseries of \cite{vanderhoeven-book}, whose supports lie
on a finitely generated multiplicative grid, and the \emph{well-based} field of
\cite{adh-book,edgar-beginners}, whose supports need only be well ordered.
\Cref{plt:thm:ext-support-hierarchy} is the elementary shadow of that
distinction.

\begin{summarybox}
Quoted, not proved.  The construction of \(\mathbb T\) as an ordered,
exponentially closed, differential field with well-based supports and finite
exponential height is \cite{vanderhoeven-book,adh-book,edgar-beginners}.  Its
monomial group contains every \(\prod_{j=0}^{r}L_j^{\alpha_j}\) with
\(\alpha_j\in\RealNumbers\), and \(\mathbb T\) is closed under well-based sums
of its own monomials; consequently, for \(\K=\RealNumbers\) and every
\(\Gamma\subseteq\RealNumbers\) and \(r\ge0\), the field
\(\mathcal H_r(\RealNumbers,\Gamma)\) of \cref{plt:def:grp-logtower} is an
ordered differential subfield of \(\mathbb T\).  The four classes of
\cref{plt:def:ring-four-classes} therefore sit inside \(\mathbb T\) as the
subrings of \(\mathcal H_1(\K,\IntegerNumbers)\) described in
\cref{plt:cor:ring-hahn-subrings}.  We use this placement only for orientation;
no theorem of this volume depends on it.
\end{summarybox}

What the exponential hierarchy adds is exactly one thing, and
\cref{plt:cor:ext-no-exponentials} says it cannot be simulated: monomials whose
logarithmic derivative is not \(\preceqdom t\).  Everything else --- real
exponents, arbitrary logarithmic depth, well-based rather than left-finite
supports --- lives already at height zero and was treated in
\cref{plt:sec:extensions}.  The cost of the addition is severe, and it is worth
stating precisely, because it explains why Parts~III and~IV could be elementary
and the corresponding theory in \(\mathbb T\) cannot.

\begin{proposition}[The filtered gain is a logarithmic-scale phenomenon]
\label{plt:prop:ext-gain-fails}
% ed.: the draft assumed only "E nonzero and c a nonzero constant".  In a ring
% with zero divisors that does not force cE \ne 0, and the proof needs
% \partial E = cE \ne 0 in order to apply the hypothesis on \nu at G = E.  The
% hypothesis is stated in the form the proof uses; it is satisfied whenever
% \mathcal T is a domain.
Let \((\mathcal T,\partial)\) be a differential ring containing a constant
\(c\) and an element \(E\) with \(\partial E=cE\ne0\).  Then there is no
map \(\nu\) from \(\mathcal T\setminus\SetOf{0}\) to a partially ordered set
that is constant on nonzero constant multiples and satisfies
\(\nu(\partial G)>\nu(G)\) for every \(G\) with \(G,\partial G\ne0\).  In
particular the one-block gain \(\ordt(\DX A)\ge\ordt(A)+1\), which is the
decisive hypothesis behind
\cref{plt:cor:grp-gain}, \cref{plt:lem:rev-deriv-gain},
\cref{plt:thm:rev-existence} and the termination of every algorithm of
\cref{plt:sec:algorithms}, cannot survive into any ring containing
\(\EulerE^{X}\).
\end{proposition}

\begin{proof}
Apply the hypothesis to \(G=E\): \(\nu(\partial E)=\nu(cE)=\nu(E)\), while the
assumed property gives \(\nu(\partial E)>\nu(E)\).  For the second sentence,
take \(c=1\), \(E=\EulerE^{X}\), and \(\nu=\ordt\), which is constant on
nonzero scalar multiples.  By contrast, inside
\(\mathcal H_r(\K,\Gamma)\) the gain does hold:
\cref{plt:thm:ext-logderiv} gives
\(\lm(\partial_rG)\preceqdom t\,\lm(G)\precdom\lm(G)\) for every nonzero
\(G\).
\end{proof}

\begin{remark}[What replaces it]
\label{plt:rmk:ext-adh-replacement}
In \(\mathbb T\) the derivation is not valuation-increasing, and the
contraction arguments of Parts~III--IV have no direct analogue.  Their
replacement is the asymptotic-differential-algebra machinery of
\cite{adh-book}: the notion of an \(H\)-field, the asymptotic couple attached
to the valuation and the derivation, the Newton degree of a differential
polynomial, and the axioms \(\omega\)-freeness and newtonianity, which are what
make the analogue of ``the residual determines the next block'' true again ---
now for differential equations rather than for a single composition.  It is
this apparatus, and not any refinement of the elementary filtration, that
carries the theory across the exponential boundary.
\end{remark}

\section{The Hardy-field perspective}
\label{plt:sec:ext-hardy}

Everything above is formal.  The classical counterpart is Hardy's calculus of
\emph{orders of infinity} \cite{hardy-orders}: the logarithmico-exponential
functions --- germs at \(+\infty\) built from the variable and real constants
by the field operations, \(\exp\) and \(\log\) --- are eventually of constant
sign, eventually monotone, and pairwise comparable, so their germs form a
totally ordered differential field, a \emph{Hardy field}.  Three points deserve
to be kept apart, and the sources of this consolidation do not always keep
them apart.

\begin{enumerate}
\item A \emph{Hardy field} consists of germs of actual real functions; its
  order is eventual dominance of real values.
\item A \emph{transseries field} consists of formal objects; its order is the
  combinatorial dominance order on supports
  (\cref{plt:def:ring-dominance}).
\item An \emph{expansion map} relates the two.  Its existence for a given germ
  is a theorem, never a definition, and its injectivity fails in general:
  \(\EulerE^{-X}\) and \(0\) have the same polynomial--logarithmic expansion.
\end{enumerate}

Within this volume, item~(2) is what
\cref{plt:lem:ring-dominance-analytic} justifies at the level of monomials, and
item~(3) is what \cref{plt:def:cmp-realization} and
\cref{plt:thm:cmp-analytic-transfer} implement, with hypotheses printed.  The
non-injectivity in item~(3) is the exact reason the class is blind to
resurgence; see \cref{plt:sec:ext-resurgence}.

\section{Model theory of the ambient, quoted}
\label{plt:sec:ext-model-theory}

\begin{summarybox}
Quoted, not proved.  The following are stated as citations and used nowhere in
this volume.
\begin{itemize}
% ed.: the draft opened this list with a Dahn--G\"oring item citing the key
% dahn-goering-1987, which is not in the merged bibliography of this volume.
% Rather than cite a key that does not resolve, the item is dropped; it carried
% no content used anywhere else.
\item Van den Dries, Macintyre and Marker \cite{vdDMM-1997} construct the field
  of logarithmic--exponential power series, a real closed field closed under
  \(\exp\) and \(\log\), and identify it as a model of the theory of the real
  field with restricted analytic functions and exponentiation; in
  \cite{vdDMM-2001} they develop its ordered differential and compositional
  structure and its relation with Hardy fields of
  logarithmico-exponential functions.
\item Aschenbrenner, van den Dries and van der Hoeven \cite{adh-book} prove
  that \(\mathbb T\), as an ordered valued differential field, is an
  \(\omega\)-free newtonian Liouville closed \(H\)-field, that this
  axiomatisation is complete and model complete in their language, and develop
  the accompanying theory of asymptotic differential algebra.
\end{itemize}
These results are what license the phrase ``the ambient field is well
behaved''.  None of them is needed to justify a single computation in
Parts~I--VI, and none of them is reproved here.
\end{summarybox}

\section{Composition, recursion, and convergence in the larger setting}
\label{plt:sec:ext-composition-large}

Composition is the operation that degrades worst under enlargement, and the
reason is visible already in \cref{plt:sec:extensions}: substitution moves
supports, and the image of a well-based support need not be well based unless
the inner argument is controlled.  Edgar's treatment \cite{edgar-composition}
isolates the admissibility conditions --- in particular that the inner argument
be positive and infinitely large --- under which composition of transseries is
defined, associative, and compatible with differentiation; the corresponding
statements in this volume are
\cref{plt:def:cmp-admissible-inner},
\cref{plt:thm:cmp-substitution-hom} and
\cref{plt:thm:cmp-associativity}, proved there in the narrow setting where
admissibility is automatic.  Fractional iteration and the iterative logarithm
--- the formal vector field \(V\) with \(\exp(V)X=F\) --- are developed for
transseries in \cite{edgar-fractional}; in the present class they are
\cref{plt:thm:grp-exp} and \cref{plt:cor:grp-iteration}, where the graded Lie
algebra of the contact filtration makes the correspondence exact and the
divisibility unique.

\begin{warningbox}[title={Two hazards that only appear after enlargement}]
Both are invisible in \(\PLlau\) and must be checked in any larger class.
\begin{enumerate}
\item \emph{Loss of well-basedness under substitution.}  Composition
  substitutes an inner series into infinitely many outer monomials; the union
  of the resulting supports is well based only under a hypothesis, and
  \cref{plt:prop:ext-wo-maximal} shows the hypothesis cannot simply be dropped
  in favour of a weaker support class.
\item \emph{Failure of Taylor approximation.}  The formal Taylor formula
  \cref{plt:thm:tay-formula} is summable because \(\DX\) raises the outer
  order; \cref{plt:prop:ext-gain-fails} shows this fails in an exponentially
  closed field, so the Taylor expansion of a transseries around a displaced
  argument requires a separate theorem with its own hypotheses.  Such theorems
  exist \cite{mantova-monotonicity}, and they are not the trivial extension of
  \cref{plt:thm:tay-formula}.
\end{enumerate}
\end{warningbox}

\section{Borel summability and resurgence: what this class cannot see}
\label{plt:sec:ext-resurgence}

The formal theory of Parts~I--IV is indifferent to convergence, and this is not
an oversight: it is a division of labour.  A divergent polynomial--logarithmic
series may be attached to an actual function by Borel--Laplace summation
\cite{costin-book}, and the ambiguity in that attachment --- the Stokes
phenomenon, and the exponentially small terms that measure it --- is the
subject of resurgence \cite{ecalle}.  The following is the honest statement of
where the present class stands.

\begin{proposition}[The class is blind to exponentially small corrections]
\label{plt:prop:ext-no-resurgence}
\begin{enumerate}[label=\textup{(\roman*)}]
\item For every \(r\ge0\), every \(\Gamma\subseteq\RealNumbers\), every
  \(\mu\in\K^{\times}\) and every ambient as in
  \cref{plt:cor:ext-no-exponentials}, the element \(\EulerE^{-\mu X}\) does not
  lie in \(\mathcal H_r(\K,\Gamma)\).
\item Let \(S=[X_0,\infty)\) be a real ray, let \(\mu>0\) and
  \(\lambda\in\ComplexNumbers\), and let \(f:S\to\ComplexNumbers\) realize
  \(F\in\PLlau\) in the sense of \cref{plt:def:cmp-realization}.  Then the
  function \(g(X)=f(X)+\lambda\EulerE^{-\mu X}\) realizes the \emph{same}
  \(F\), and by \cref{plt:lem:cmp-realization-basic} no other element of
  \(\PLlau\).
\end{enumerate}
Consequently no truncation, no residual certificate
\textup{(\cref{plt:def:rev-residual})} and no error estimate of Parts~I--VI
distinguishes \(f\) from \(g\).
\end{proposition}

\begin{proof}
(i)  This is \cref{plt:cor:ext-no-exponentials}\,(i) with \(c=-\mu\).

(ii)  For every \(N\in\IntegerNumbers\) and every \(\mu>0\) one has
\(\EulerE^{-\mu X}=\LittleOAt{X^{1-N}}{X\to+\infty}\), since
\(X^{N-1}\EulerE^{-\mu X}\to0\).  Hence
\[
  g(X)-\bigl(\Trunc{F}{N}\bigr)(X)
  =\Bigl(f(X)-\bigl(\Trunc{F}{N}\bigr)(X)\Bigr)
   +\lambda\EulerE^{-\mu X}
  =\LittleOAt{X^{1-N}}{X\to+\infty},
\]
% ed.: the draft referred to an equation label plt:eq:cmp-realization that the
% volume does not carry; the condition invoked is the defining one of the
% realization relation.
which is the defining remainder condition of \cref{plt:def:cmp-realization} for
\(g\).  Uniqueness of the realized
series is \cref{plt:lem:cmp-realization-basic}.
\end{proof}

\begin{remark}[Why the interesting resurgence is not here]
\label{plt:rmk:ext-resurgence-honest}
Resurgence is a statement about the \emph{relations} between the exponentially
small sectors of a transseries and the singularities of its Borel transform.
\Cref{plt:prop:ext-no-resurgence} says the polynomial--logarithmic class
contains no such sector at all, at any logarithmic depth; the class is
therefore too small for the phenomena to be visible, not merely inconvenient
for them.  What the class \emph{can} contribute is the height-zero datum on
which those phenomena act: the formal power--log series that a Borel--Laplace
summation is asked to sum, together with the exact residual certificates of
\cref{plt:def:rev-residual} and \cref{plt:prop:rev-residual-order}, which
control the algebraic part of the error without any convergence hypothesis.
Where a convergence statement is actually available in this volume it is proved
and its scope is printed --- \cref{plt:prop:om-frozen-convergence} for the
% ed.: three appeals to a label plt:rmk:om-open, which the volume does not
% carry, are replaced by the statement they were standing in for; see also
% \cref{plt:rmk:ext-open-convergence} and \cref{plt:rmk:ext-open-uniformity}.
Lambert generating function at frozen logarithm --- and where it is not, it is
recorded as unproved in \cref{plt:sec:ext-open}.  No claim of Borel summability
is made anywhere in this volume.
\end{remark}

\section{Symbolic asymptotics for larger classes}
\label{plt:sec:ext-symbolic}

The algorithmic counterpart of this section is the multiseries technology of
Salvy and Shackell.  Their programme computes asymptotic expansions of
exp--log functions by discovering the scale rather than fixing it in advance:
% ed.: the draft also cited salvy-shackell-1992, a key that is not in the
% merged bibliography; the item it carried (functional inverses) is covered by
% salvy-shackell-1999, which is.
implicit functions and
functions of two variables \cite{salvy-shackell-1998}, and multiseries of
inverse functions \cite{salvy-shackell-1999}.  The relationship with
\cref{plt:sec:algorithms} is the following, and it is a genuine division of
labour rather than a comparison.

\begin{itemize}
\item Their algorithms decide, at each step, \emph{which} scale the next term
  lives in --- an inverse power of \(X\), an inverse power of \(\log X\), a new
  iterated logarithm, a root, or an exponential monomial --- and then recurse.
  \Cref{plt:cor:lag-enlargement}, \cref{plt:thm:lag-rigidity} and
  \cref{plt:thm:ext-tower-strict} are the corresponding decision content in the
  present class, stated as theorems about which enlargement is forced.
\item Once the scale is fixed, the coefficient arithmetic is exactly the block
  calculus of Parts~I--IV, and the cost model of
  \cref{plt:def:alg-cost-models} applies.  The sharp scalar cost of truncated
  block convolution (\cref{plt:prop:alg-convolution-cost}) and the precision
  calculus (\cref{plt:thm:alg-precision-calculus}) are statements that a
  general multiseries implementation does not have available, because they use
  the finiteness of each block.
\item The practical rule that follows is the one this volume has argued for
  throughout: begin in the smallest class that contains the problem, run the
  closure test for the next operation, and promote explicitly.  Promotion
  hidden inside a simplifier is how a computation acquires an unproved
  convergence claim.
\end{itemize}

\section{Open problems}
\label{plt:sec:ext-open}

The following are stated as open, not as theorems and not as conjectures with
suppressed hypotheses.  Each is precise enough to be attacked, and each is
raised, and left unsettled, by material in this volume.

\begin{remark}[Open problem: convergence of the Lambert expansion with coupled
logarithm]
\label{plt:rmk:ext-open-convergence}
Does there exist \(x_0>0\) such that
\(\sum_{n\ge1}\AbsoluteValue{\LamP{n}(\log x)}\,x^{-n}<\infty\) for every real
\(x>x_0\)?  Equivalently, is the radius of convergence \(r(u)\) of
\(s\mapsto\sum_{n\ge1}\LamP{n}(u)s^{n}\) bounded below by \(\kappa/u\) for some
\(\kappa>0\) and all large real \(u\)?  \Cref{plt:prop:om-frozen-convergence}
proves convergence for each \emph{frozen} \(u\) with no lower bound on
% ed.: the draft attributed the singularity locus to a label plt:rmk:om-open
% that the volume does not carry.  The locus is recomputed here from the
% generating equation, and it does come out as the draft stated it.
\(r(u)\).  Writing \(\WrightOmega(X)=X+h\) with
\(h=\sum_{n\ge0}\LamP{n}(u)s^{n}\), \(u=L\), \(s=t\), the defining relation
\(\WrightOmega+\log\WrightOmega=X\) becomes \(h+u+\log(1+sh)=0\) --- this is
the implicit generating equation of \cref{plt:thm:om-generating-equation}, and
the substitution just made rederives it, the value \(h(u,0)=\LamP{0}(u)=-u\)
fixing the branch.  Its \(h\)-derivative
\(1+s/(1+sh)\) vanishes exactly at \(1+sh=-s\), that is at
\(h=-1-s^{-1}\), and eliminating \(h\) puts the singularities of
\(s\mapsto h(u,s)\) on the curve
\(u=s^{-1}+1-\log(-s)\).  This suggests \(r(u)\asymp1/u\); nothing in
Parts~I--IV proves it.  A positive answer would upgrade
\cref{plt:thm:om-expansion} from an asymptotic expansion to a convergent one on
a half-line.
\end{remark}

\begin{remark}[Open problem: sharp logarithmic depth of a composite]
\label{plt:rmk:ext-open-depth}
Fix \(r\ge1\) and let \(G_1,\ldots,G_k\) be monomial-leading maps, each of
logarithmic depth at most \(r\) in the sense that its leading monomial lies in
\(\mathcal M_r(\Gamma)\).  % ed.: the draft extended the depth bound r+k to the compositional inverse as
% well.  Nothing in this volume proves a depth bound for a reversion: the
% classification of scale changes bounds the depth of a substitution, not of an
% inverse, and \cref{plt:prop:lag-loglog-forced} bounds it only from below.
% The claim is weakened to what is proved, and the missing half is folded into
% the open question.
\Cref{plt:thm:grp-scale-classification} guarantees
that a single such substitution raises depth by at most one, so
\(G_1\compose\cdots\compose G_k\) lies at depth at most \(r+k\); this volume
proves no corresponding bound for its compositional inverse.  Is \(r+k\)
attained by the composite, and is the inverse subject to any depth bound at
all?  More precisely: give an algorithm that, from
the leading data of \(G_1,\ldots,G_k\), computes the least \(r'\) with
\(\rev{(G_1\compose\cdots\compose G_k)}\in\mathcal H_{r'}\), together with a
proof of minimality.  The only lower bound proved here is the single step of
\cref{plt:prop:lag-loglog-forced}, and the only obstruction proved here is the
residue of \cref{plt:thm:ext-tower-strict}; neither is known to be complete.
\end{remark}

\begin{remark}[Open problem: a Lagrange--B\"urmann formula at dense exponents]
\label{plt:rmk:ext-open-burmann}
\Cref{plt:thm:lag-burmann} expresses the reversion of a near-identity map as
\(\sum_{r\ge1}\frac{(-1)^{r}}{r!}\partial^{\,r-1}(A^{r})\), and
\cref{plt:prop:lag-range} shows that the coefficient of \(t^{N}\) is
determined by the indices \(r\le N+1\) and by no fewer.  Both statements use
that the exponent group is \(\IntegerNumbers\).  In a Hahn power--log scale
\(\mathcal H(\K[L],\Gamma)\) with \(\Gamma\) dense, is there a summation index
set --- necessarily not \(\PositiveIntegers\) with the range \(r\le N+1\) ---
and an operator identity reducing to \cref{plt:thm:lag-burmann} at
\(\Gamma=\IntegerNumbers\)?  The near-identity inverse exists there by
\cref{plt:cor:ext-lf-group}; what is missing is a closed formula for it.
\end{remark}

\begin{remark}[Open problem: denominator growth in the Puiseux tower]
\label{plt:rmk:ext-open-denominators}
By \cref{plt:thm:ext-puiseux-incomplete} the Puiseux-log union
\(\mathfrak P\) is not complete, so an infinite iteration may leave it.  Give
an explicit bound, or prove that none exists, for the following: if
\(F=X+A\) with \(A\in\K[L][[t^{1/s}]]\), the Newton scheme of
\cref{plt:def:rev-newton} started at \(X\) produces iterates
\(H_0,H_1,\ldots\); is there a function \(s\mapsto\sigma(s)\), independent of
\(A\), with \(H_k\in\K[L]((t^{1/\sigma(s)}))\) for all \(k\)?  For the
triangular scheme of \cref{plt:thm:rev-triangular} the answer is yes with
\(\sigma(s)=s\), since every block of the inverse has exponent in
\(s^{-1}\IntegerNumbers\); the question is whether the quadratically
convergent scheme, which divides by \(F'\compose H_k\), preserves the
denominator.
\end{remark}

\begin{remark}[Open problem: uniformity of the analytic transfer]
\label{plt:rmk:ext-open-uniformity}
\Cref{plt:thm:om-analytic} and \cref{plt:thm:cmp-analytic-transfer} give, for
each fixed truncation order \(N\), an asymptotic error estimate on the positive
ray with a constant depending on \(N\).  Is there an estimate uniform in
\(N\) --- for instance a bound of the form
\(\AbsoluteValue{\WrightOmega(X)-\Trunc{\WrightOmega}{N}(X)}\le
C^{N}N!\,X^{-N}(\log X)^{N}\) valid for all \(N\) and all \(X\ge X_0\), with
\(C,X_0\) absolute?  Is there a version uniform in a complex sector and in the
branch index \(k\) of \(\LambertW_k\)?  None of these is claimed or proved
anywhere in this volume; the literature asserts
related convergence statements \cite{corless-et-al-1996,jeffrey-et-al-1995,
dlmf-W} whose proofs the present consolidation has not verified.
\end{remark}

\begin{remark}[Open problem: realisation of left-finite Hahn elements by germs]
\label{plt:rmk:ext-open-realization}
\Cref{plt:thm:ext-lf-completion} identifies
\(\mathcal H_{\mathrm{lf}}(\K[L],\Gamma)\) as the completion of the finitely
supported power--log series, and \cref{plt:cor:ext-lf-group} shows it carries
the whole near-identity theory.  Which of its elements are asymptotic
expansions of germs at \(+\infty\) lying in a common Hardy field, in the sense
of \cref{plt:def:cmp-realization}?  For \(\Gamma=\IntegerNumbers\) the question
is the classical one about power--log series; for
\(\Gamma=\RationalNumbers\), where by
\cref{plt:thm:ext-puiseux-incomplete} the class strictly contains the Puiseux
tower, we do not know a characterisation, nor an example of an element of
\(\mathcal H_{\mathrm{lf}}(\K[L],\RationalNumbers)\setminus\mathfrak P\)
that is provably realised.
\end{remark}

\begin{remark}[Open problem: a resurgent completion with block arithmetic]
\label{plt:rmk:ext-open-resurgence}
\Cref{plt:prop:ext-no-resurgence} shows the class cannot see exponentially
small corrections.  Is there a minimal enlargement --- adjoining a single
transmonomial \(\EulerE^{-\mu X}\), \(\mu>0\), with coefficients in
\(\PLlau\), that is, the ring
\(\PLlau[\![\EulerE^{-\mu X}]\!]\) with its natural derivation --- in which
the block arithmetic of Parts~I--IV survives verbatim, the unit criterion and
the reversion theorem included?  \Cref{plt:prop:ext-gain-fails} shows the
valuation gain of \(\DX\) fails on the new generator, so a positive answer
would need a two-index filtration (order in \(t\), order in
\(\EulerE^{-\mu X}\)) and a proof that the composition operator is continuous
for it.  We state this as open because the naive filtration does not work and
we have not verified any substitute.
\end{remark}

\begin{remark}[Open problem: which of these statements are machine-checkable
today]
\label{plt:rmk:ext-open-formalization}
The formalisation register of this volume records, for each stated result, its
status in a proof assistant.  The support-theoretic results of
\cref{plt:sec:ext-supports} are finitary in character --- they are statements
about ordered abelian groups and well-quasi-orders --- and should be reachable.
\Cref{plt:thm:ext-tower-strict} is a finite computation with a monomial
bookkeeping argument.  \Cref{plt:thm:ext-lf-completion} needs a formal Hahn
series library with a \(\Gamma\)-indexed filtration.  Which of these is
actually within reach of a current library, and at what cost, is not known to
us, and the register should be read as a plan rather than as a claim.
\end{remark}

\begin{summarybox}
The polynomial--logarithmic class is a proper, well-delimited fragment of the
logarithmic--exponential transseries field.  What separates it from the ambient
is a single structural fact, proved here at every logarithmic depth: inside the
class the logarithmic derivative of a nonzero element is never larger than
\(X^{-1}\) (\cref{plt:thm:ext-logderiv}).  Everything the class does easily ---
triangular reversion, terminating algorithms, exact residual certificates ---
follows from the valuation gain that this fact expresses, and everything it
cannot do --- exponentially small terms, Stokes data, resurgence ---
follows from the same fact read as a limitation
(\cref{plt:prop:ext-gain-fails}, \cref{plt:prop:ext-no-resurgence}).  The right
attitude is neither to inflate the class nor to abandon it, but to promote
explicitly, one axis at a time, with the closure test named.
\end{summarybox}

\clearpage
\appendix
\part*{Appendices}
\addcontentsline{toc}{part}{Appendices}

% ---- fragment 16-bell ----
\chapter{Bell polynomials and Fa\`a di Bruno bookkeeping}
\label{plt:sec:bell}

The exponential partial Bell polynomials appear at four separate places in the
body of this volume: in the coefficients of an integer power
(\cref{plt:prop:mul-integer-power}), in the coefficients of a series
substituted into a scalar power series
(\cref{plt:prop:dif-bell,plt:prop:tay-scalar-outer}), in the shift recurrence
that drives Fa\`a di Bruno for admissible composition
(\cref{plt:lem:tay-bell-recurrence,plt:prop:tay-faa-di-bruno}), and in the
closed form of the convolution polynomials
(\cref{plt:prop:tay-convolution}\,(iv)).
% ed.: the second group is a restatement inside the volume, not an error, but
% the draft called its two members "the same statement" and declared them
% interchangeable.  They are not: plt:prop:tay-scalar-outer restates
% plt:prop:dif-bell and then adds a clause about expanding the outer function
% about a base point z_0 \ne 0.  Nothing below leans on that clause -- see the
% editorial note at plt:lem:bell-chain, where the base point is made part of
% the datum instead -- but the two statements are not interchangeable.
(The second member of that group restates the first in the part on
composition, and adds a clause about a base point \(z_{0}\ne0\); only the
first is proved from scratch.)  In
each of those places the defining
generating identity is quoted from \cite{comtet-book}.  This appendix pays that
debt: it defines the family, proves the generating identity and the
recurrences from the definition, proves the two Stirling evaluations --- one of
which is exactly the arithmetic content of the Lambert polynomials of
\cref{plt:sec:lambert-polynomials} --- and proves the form of Fa\`a di Bruno
that the volume needs but does not state, namely the one whose outer function
is a scalar-analytic function of a polynomial--logarithmic series.

Throughout, \(\K\) is a field of characteristic zero, so that
\(\RationalNumbers\subseteq\K\) and every factorial occurring below is
invertible.  The polynomial identities of
\cref{plt:sec:bell-definition,plt:sec:bell-stirling} are identities in
\(\RationalNumbers[x_{1},x_{2},\ldots]\) and are therefore available over every
commutative \(\RationalNumbers\)-algebra, in particular over \(\K[L]\),
\(\K(L)\), \(\PLpow\) and \(\PLlau\).

\section{Definition, generating identity, and the two normalizations}
\label{plt:sec:bell-definition}

\begin{definition}[Exponential partial and complete Bell polynomials]
\label{plt:def:bell-partial}
Let \(x_{1},x_{2},\ldots\) be independent indeterminates.  For
\(n,k\in\NaturalNumbers\) put
\begin{equation}\label{plt:eq:bell-partial-def}
 \ExponentialPartialBellPolynomial{n}{k}(x_{1},x_{2},\ldots)
 \DefinitionEquals
 \sum_{\substack{m_{1},m_{2},\ldots\ge0\\
                 \sum_{j\ge1}m_{j}=k\\
                 \sum_{j\ge1}jm_{j}=n}}
 \frac{n!}{\prod_{j\ge1}m_{j}!\,(j!)^{m_{j}}}\;
 \prod_{j\ge1}x_{j}^{m_{j}}
 \ \in\ \RationalNumbers[x_{1},x_{2},\ldots],
\end{equation}
the sum ranging over all sequences \((m_{j})_{j\ge1}\) of nonnegative integers
with finite support satisfying the two displayed constraints, and the empty
product being \(1\).  The \emph{complete} polynomial is
\begin{equation}\label{plt:eq:bell-complete-def}
 \ExponentialCompleteBellPolynomial{n}(x_{1},\ldots,x_{n})
 \DefinitionEquals
 \sum_{k=0}^{n}
 \ExponentialPartialBellPolynomial{n}{k}(x_{1},\ldots,x_{n-k+1}).
\end{equation}
\end{definition}

The index set in \eqref{plt:eq:bell-partial-def} is finite: the constraint
\(\sum_{j}jm_{j}=n\) forces \(m_{j}=0\) for \(j>n\) and \(m_{j}\le n/j\).  It
is empty --- so that
\(\ExponentialPartialBellPolynomial{n}{k}=0\) --- whenever no such sequence
exists.  Two degenerate cases matter and are fixed once and for all by
\eqref{plt:eq:bell-partial-def} itself: for \(n=k=0\) the only sequence is
\(m\equiv0\), giving
\(\ExponentialPartialBellPolynomial{0}{0}=1\); and for \(n>0\), \(k=0\) the
constraints \(\sum m_{j}=0\) and \(\sum jm_{j}=n>0\) are incompatible, so
\(\ExponentialPartialBellPolynomial{n}{0}=0\).  Consequently
\(\ExponentialCompleteBellPolynomial{0}=1\).

\begin{lemma}[Support, degree, weight, and boundary values]
\label{plt:lem:bell-elementary}
Fix \(n,k\in\NaturalNumbers\) and write
\(B\DefinitionEquals\ExponentialPartialBellPolynomial{n}{k}\).
\begin{enumerate}[label=\textup{(\roman*)}]
\item \(B=0\) unless \(0\le k\le n\); moreover \(B=1\) when \(n=k=0\) and
  \(B=0\) when \(n\ge1\) and \(k=0\).  When \(1\le k\le n\), only the
  variables \(x_{1},\ldots,x_{n-k+1}\) occur in \(B\).
\item \(B\) is homogeneous of degree \(k\) in the variables and isobaric of
  weight \(n\) for the weighting \(\operatorname{wt}(x_{j})=j\).
  Equivalently, for indeterminates \(a,b\),
  \begin{equation}\label{plt:eq:bell-bihomogeneous}
   \ExponentialPartialBellPolynomial{n}{k}
     \bigl(abx_{1},ab^{2}x_{2},ab^{3}x_{3},\ldots\bigr)
   =a^{k}b^{n}\,
   \ExponentialPartialBellPolynomial{n}{k}(x_{1},x_{2},\ldots).
  \end{equation}
\item The boundary rows and columns are
  \begin{equation}\label{plt:eq:bell-boundary}
   \ExponentialPartialBellPolynomial{n}{n}=x_{1}^{n},
   \qquad
   \ExponentialPartialBellPolynomial{n}{1}=x_{n}\ (n\ge1),
   \qquad
   \ExponentialPartialBellPolynomial{n}{n-1}=\binom{n}{2}x_{1}^{n-2}x_{2}
   \ (n\ge2).
  \end{equation}
\end{enumerate}
\end{lemma}

\begin{proof}
(i)  A contributing sequence has \(\sum_{j}m_{j}=k\) and
\(\sum_{j}jm_{j}=n\); since every \(j\ge1\), the second sum is at least the
first, so \(n\ge k\).  If \(k=0\) the first sum forces \(m\equiv0\) and then
\(n=0\).  Suppose \(1\le k\le n\) and \(m_{i}\ge1\) for some \(i\).  The
remaining \(k-1\) parts contribute at least \(k-1\) to \(\sum_{j}jm_{j}\),
whence \(i\le n-(k-1)=n-k+1\).

(ii)  Each monomial \(\prod_{j}x_{j}^{m_{j}}\) occurring in
\eqref{plt:eq:bell-partial-def} has total degree \(\sum_{j}m_{j}=k\) and
weight \(\sum_{j}jm_{j}=n\).  Substituting \(x_{j}\mapsto ab^{j}x_{j}\)
multiplies that monomial by
\(a^{\sum_{j}m_{j}}b^{\sum_{j}jm_{j}}=a^{k}b^{n}\), which is
\eqref{plt:eq:bell-bihomogeneous}.

(iii)  For \(k=n\ge1\), \(\sum_{j}m_{j}=\sum_{j}jm_{j}\) forces \(m_{j}=0\) for
\(j\ge2\) and hence \(m_{1}=n\); the coefficient is
\(n!/(n!\,(1!)^{n})=1\).  For \(k=1\) the only sequence is \(m_{n}=1\) with all
other entries zero; the coefficient is \(n!/(1!\,(n!)^{1})=1\).  For \(k=n-1\)
with \(n\ge2\) the only sequence is \(m_{1}=n-2\), \(m_{2}=1\); its coefficient
is \(n!/\bigl((n-2)!\,1!\,(1!)^{n-2}(2!)^{1}\bigr)=n(n-1)/2=\binom{n}{2}\).
\end{proof}

\begin{proposition}[The defining generating identity]
\label{plt:prop:bell-gf}
Let \(R\) be a commutative \(\RationalNumbers\)-algebra, let
\(\xi_{1},\xi_{2},\ldots\in R\), and put
\(\Xi(z)\DefinitionEquals\sum_{j\ge1}\xi_{j}z^{j}/j!\in zR[[z]]\).  Then, in
\(R[[z]]\),
\begin{equation}\label{plt:eq:bell-egf}
 \frac{\Xi(z)^{k}}{k!}
 =\sum_{n\ge k}
   \ExponentialPartialBellPolynomial{n}{k}(\xi_{1},\xi_{2},\ldots)
   \frac{z^{n}}{n!}
 \qquad(k\in\NaturalNumbers),
\end{equation}
and the family \(\bigl(\Xi^{k}/k!\bigr)_{k\ge0}\) is summable in \(R[[z]]\)
with
\begin{equation}\label{plt:eq:bell-complete-egf}
 \exp\Xi(z)
 \DefinitionEquals\sum_{k\ge0}\frac{\Xi(z)^{k}}{k!}
 =\sum_{n\ge0}
   \ExponentialCompleteBellPolynomial{n}(\xi_{1},\ldots,\xi_{n})
   \frac{z^{n}}{n!} .
\end{equation}
Conversely, \eqref{plt:eq:bell-egf} taken over
\(R=\RationalNumbers[x_{1},x_{2},\ldots]\) with \(\xi_{j}=x_{j}\) determines
\eqref{plt:eq:bell-partial-def} uniquely.
\end{proposition}

\begin{proof}
Fix \(k\ge1\) (the case \(k=0\) is \(1=1\)).  By the multinomial theorem in
\(R[[z]]\), applied to the \(k\)-th power of a series with zero constant term,
\[
 \Xi(z)^{k}
 =\sum_{(j_{1},\ldots,j_{k})\in\PositiveIntegers^{k}}
   \prod_{i=1}^{k}\frac{\xi_{j_{i}}z^{j_{i}}}{j_{i}!}
 =\sum_{n\ge k}z^{n}
   \sum_{\substack{j_{1},\ldots,j_{k}\ge1\\ j_{1}+\cdots+j_{k}=n}}
   \prod_{i=1}^{k}\frac{\xi_{j_{i}}}{j_{i}!},
\]
the rearrangement being legitimate because for fixed \(n\) only the finitely
many tuples with \(j_{1}+\cdots+j_{k}=n\) contribute, and there are none when
\(n<k\).  Now group the tuples by their multiplicity vector: a tuple
\((j_{1},\ldots,j_{k})\) with \(m_{j}\DefinitionEquals\#\SetOf{i:j_{i}=j}\)
satisfies \(\sum_{j}m_{j}=k\) and \(\sum_{j}jm_{j}=n\), and the number of
tuples with a prescribed multiplicity vector is the multinomial coefficient
\(k!/\prod_{j}m_{j}!\).  Every tuple in one such class contributes the same
term \(\prod_{j}(\xi_{j}/j!)^{m_{j}}\).  Hence
\[
 [z^{n}]\,\Xi^{k}
 =\sum_{\substack{\sum_{j}m_{j}=k\\ \sum_{j}jm_{j}=n}}
   \frac{k!}{\prod_{j}m_{j}!}\prod_{j}\Bigl(\frac{\xi_{j}}{j!}\Bigr)^{m_{j}}
 =\frac{k!}{n!}\,
   \ExponentialPartialBellPolynomial{n}{k}(\xi_{1},\xi_{2},\ldots),
\]
by \eqref{plt:eq:bell-partial-def}.  Dividing by \(k!\) gives
\eqref{plt:eq:bell-egf}.

For \eqref{plt:eq:bell-complete-egf}: \(\Xi^{k}/k!\in z^{k}R[[z]]\), so the
family is summable for the \(z\)-adic filtration and its sum is computed
blockwise; by \eqref{plt:eq:bell-egf} and
\cref{plt:lem:bell-elementary}\,(i) the coefficient of \(z^{n}/n!\) in that
sum is \(\sum_{k=0}^{n}\ExponentialPartialBellPolynomial{n}{k}\), which is
\eqref{plt:eq:bell-complete-def}.  Finally, in the universal case
\(R=\RationalNumbers[x_{1},x_{2},\ldots]\) the coefficients of a formal power
series are uniquely determined, so \eqref{plt:eq:bell-egf} recovers
\eqref{plt:eq:bell-partial-def}.
\end{proof}

\Cref{plt:prop:bell-gf} is precisely the identity quoted without proof in the
proofs of \cref{plt:prop:mul-integer-power},
\cref{plt:prop:tay-convolution}\,(iv) and
\cref{plt:lem:tay-bell-recurrence}; those citations may now be read as
references to \eqref{plt:eq:bell-egf}.

The volume uses two different normalizations for the same combinatorial data,
and the conversion between them is a constant source of factorial errors.  The
next lemma states it once.

\begin{lemma}[Ordinary versus exponential normalization]
\label{plt:lem:bell-normalizations}
Let \(R\) be a commutative \(\RationalNumbers\)-algebra and
\(\xi_{1},\xi_{2},\ldots\in R\).  Define the \emph{ordinary} partial Bell
polynomial by
\begin{equation}\label{plt:eq:bell-ord-def}
 \OrdinaryPartialBellPolynomial{n}{k}(\xi_{1},\xi_{2},\ldots)
 \DefinitionEquals
 \sum_{\substack{m_{1},m_{2},\ldots\ge0\\ \sum_{j}m_{j}=k\\
                 \sum_{j}jm_{j}=n}}
 \frac{k!}{\prod_{j\ge1}m_{j}!}\prod_{j\ge1}\xi_{j}^{m_{j}},
\end{equation}
so that \(\bigl(\sum_{j\ge1}\xi_{j}z^{j}\bigr)^{k}
=\sum_{n\ge k}\OrdinaryPartialBellPolynomial{n}{k}(\xi_{1},\ldots)z^{n}\).
Then for all \(n,k\in\NaturalNumbers\)
\begin{equation}\label{plt:eq:bell-ord-exp}
 \OrdinaryPartialBellPolynomial{n}{k}(\xi_{1},\xi_{2},\ldots)
 =\frac{k!}{n!}\,
  \ExponentialPartialBellPolynomial{n}{k}
   \bigl(1!\,\xi_{1},2!\,\xi_{2},3!\,\xi_{3},\ldots\bigr).
\end{equation}
In particular, for \(H=\sum_{j\ge1}h_{j}(L)t^{j}\in t\PLpow\) the convolution
polynomials of \cref{plt:def:tay-convolution} satisfy
\(C_{n,k}=\OrdinaryPartialBellPolynomial{n}{k}(h_{1},h_{2},\ldots)\), which is
\cref{plt:prop:tay-convolution}\,(iv).
\end{lemma}

\begin{proof}
Apply \eqref{plt:eq:bell-bihomogeneous} in the substituted form: putting
\(\hat\xi_{j}=j!\,\xi_{j}\) in \eqref{plt:eq:bell-partial-def} gives
\[
 \ExponentialPartialBellPolynomial{n}{k}(\hat\xi_{1},\hat\xi_{2},\ldots)
 =\sum_{\substack{\sum_{j}m_{j}=k\\ \sum_{j}jm_{j}=n}}
   \frac{n!}{\prod_{j}m_{j}!\,(j!)^{m_{j}}}
   \prod_{j}(j!)^{m_{j}}\xi_{j}^{m_{j}}
 =\frac{n!}{k!}
   \sum_{\substack{\sum_{j}m_{j}=k\\ \sum_{j}jm_{j}=n}}
   \frac{k!}{\prod_{j}m_{j}!}\prod_{j}\xi_{j}^{m_{j}},
\]
which is \eqref{plt:eq:bell-ord-exp}.  The generating identity for
\eqref{plt:eq:bell-ord-def} follows by the same multiplicity grouping as in
the proof of \cref{plt:prop:bell-gf}, now without the factors \(1/j!\).  The
last sentence is \cref{plt:prop:tay-convolution}\,(i) combined with that
generating identity.
\end{proof}

\begin{warningbox}[title={Three two-index arrays, three index orders}]
The volume carries three arrays whose entries are indexed by a pair of
integers and which are easy to confuse.
\begin{itemize}
\item \(h_{k,m}=[t^{m}]H^{k}\) of \eqref{plt:eq:tay-Hk-blocks}: first index the
  exponent, second the outer order.
\item \(C_{m,r}\) of \cref{plt:def:tay-convolution}: first index the outer
  order, second the exponent; \(h_{k,m}=C_{m,k}\).
\item \(\ExponentialPartialBellPolynomial{n}{k}\) of
  \cref{plt:def:bell-partial}: first index the \emph{weight} \(n\), second the
  \emph{degree} \(k\); it is an ordinary polynomial in abstract
  indeterminates, and only after the substitution
  \(x_{j}\mapsto j!\,h_{j}\) and the factor \(k!/n!\) of
  \eqref{plt:eq:bell-ord-exp} does it become \(C_{n,k}\).
\end{itemize}
Transferring a formula between any two of them requires transposing the
indices, and between the last and either of the first two also requires the
factorial conversion.  A missing \(k!/n!\) is the single commonest error in
this circle of identities.
\end{warningbox}

\begin{proposition}[Two computational recurrences]
\label{plt:prop:bell-recurrences}
For all \(n\ge1\) and \(k\ge1\),
\begin{align}
 \ExponentialPartialBellPolynomial{n}{k}
 &=\sum_{j=1}^{n-k+1}\binom{n-1}{j-1}\,x_{j}\,
   \ExponentialPartialBellPolynomial{n-j}{k-1},
 \label{plt:eq:bell-comtet-rec}\\
 k\,\ExponentialPartialBellPolynomial{n}{k}
 &=\sum_{j=1}^{n-k+1}\binom{n}{j}\,x_{j}\,
   \ExponentialPartialBellPolynomial{n-j}{k-1}.
 \label{plt:eq:bell-column-rec}
\end{align}
Either recurrence, together with the boundary values
\(\ExponentialPartialBellPolynomial{0}{0}=1\),
\(\ExponentialPartialBellPolynomial{n}{0}=0\) for \(n\ge1\) and
\(\ExponentialPartialBellPolynomial{0}{k}=0\) for \(k\ge1\), determines the
whole array.
\end{proposition}

\begin{proof}
Work in \(R[[z]]\) with \(R=\RationalNumbers[x_{1},x_{2},\ldots]\) and
\(\Xi(z)=\sum_{j\ge1}x_{j}z^{j}/j!\).

For \eqref{plt:eq:bell-column-rec}, write \(\Xi^{k}/k!
=\frac1k\,\Xi\cdot\Xi^{k-1}/(k-1)!\) and extract \([z^{n}/n!]\).  For two
series \(P=\sum_{n}p_{n}z^{n}/n!\) and \(Q=\sum_{n}q_{n}z^{n}/n!\) one has
\([z^{n}/n!](PQ)=\sum_{j=0}^{n}\binom{n}{j}p_{j}q_{n-j}\); applying this with
\(P=\Xi\) (so \(p_{0}=0\), \(p_{j}=x_{j}\)) and
\(Q=\Xi^{k-1}/(k-1)!\) and using \eqref{plt:eq:bell-egf} gives
\(\frac1k\sum_{j\ge1}\binom{n}{j}x_{j}
\ExponentialPartialBellPolynomial{n-j}{k-1}\).  The upper limit is
\(n-k+1\) because \(\ExponentialPartialBellPolynomial{n-j}{k-1}=0\) unless
\(n-j\ge k-1\) (\cref{plt:lem:bell-elementary}\,(i)).

For \eqref{plt:eq:bell-comtet-rec}, differentiate
\eqref{plt:eq:bell-egf} with respect to \(z\):
\[
 \frac{\Differential}{\Differential z}\frac{\Xi^{k}}{k!}
 =\Xi'\,\frac{\Xi^{k-1}}{(k-1)!},
 \qquad
 \Xi'(z)=\sum_{i\ge0}x_{i+1}\frac{z^{i}}{i!} .
\]
The left side is \(\sum_{n\ge0}
\ExponentialPartialBellPolynomial{n+1}{k}z^{n}/n!\).  Applying the product
rule for exponential generating functions to the right side with
\(p_{i}=x_{i+1}\) gives
\(\sum_{i=0}^{n}\binom{n}{i}x_{i+1}
\ExponentialPartialBellPolynomial{n-i}{k-1}\).  Replacing \(n\) by \(n-1\) and
setting \(j=i+1\) turns this into
\(\ExponentialPartialBellPolynomial{n}{k}
=\sum_{j\ge1}\binom{n-1}{j-1}x_{j}
\ExponentialPartialBellPolynomial{n-j}{k-1}\), which is
\eqref{plt:eq:bell-comtet-rec}; the upper limit is again forced by
\cref{plt:lem:bell-elementary}\,(i).

Both recurrences express \(\ExponentialPartialBellPolynomial{n}{k}\) in terms
of entries \(\ExponentialPartialBellPolynomial{n-j}{k-1}\) with \(j\ge1\),
hence in terms of strictly earlier rows; with the three stated boundary
values --- which supply every entry of row \(0\) and every entry of column
\(0\) --- they determine the array by induction on \(n\).
\end{proof}

\begin{remark}[A third recurrence, and where each is used]
\label{plt:rmk:bell-third-recurrence}
\Cref{plt:lem:tay-bell-recurrence} proves a third recurrence,
\(\ExponentialPartialBellPolynomial{n+1}{k}
=\mathcal D\,\ExponentialPartialBellPolynomial{n}{k}
+x_{1}\ExponentialPartialBellPolynomial{n}{k-1}\) with
\(\mathcal D=\sum_{i\ge1}x_{i+1}\partial_{x_{i}}\), which is the one adapted to
\emph{differentiating} a Bell polynomial evaluated at a tower of derivatives;
it is the engine of every Fa\`a di Bruno induction, including
\cref{plt:prop:tay-faa-di-bruno} and
\cref{plt:lem:bell-abstract-faa} below.  By contrast
\eqref{plt:eq:bell-comtet-rec} and \eqref{plt:eq:bell-column-rec} are the ones
adapted to \emph{tabulating} the array: each computes one entry from
\(n-k+1\) entries of the previous column, so the whole triangle up to weight
% ed.: the draft wrote "$\BigO(N^{3})$ coefficient multiplications" without
% saying what a coefficient is here.  The count $\BigO(N^{3})$ is correct only
% for the operation "multiply one indeterminate into one stored entry"; the
% entries are not scalars, and the scalar count is genuinely superpolynomial.
% Left unqualified this is exactly the conflation that
% plt:cor:alg-cost-not-quadratic warns against, so the two levels are now
% separated.
\(N\) costs \(\BigO(N^{3})\) multiplications of one indeterminate \(x_{j}\)
into one previously computed entry.  That is a count of \emph{entry}
operations, not of scalar ones: the entry
\(\ExponentialPartialBellPolynomial{n}{k}\) has exactly one monomial for each
partition of \(n\) into \(k\) parts, so row \(n\) carries as many monomials as
there are partitions of \(n\), and the number of scalar coefficients written
while tabulating up to weight \(N\) therefore exceeds every polynomial in
\(N\).  The distinction between the two levels of counting is the one drawn
in \cref{plt:cor:alg-cost-not-quadratic}.
\end{remark}

\begin{example}[The array through weight six]
\label{plt:ex:bell-table}
\Cref{plt:tab:bell-partial} lists
\(\ExponentialPartialBellPolynomial{n}{k}\) for \(1\le k\le n\le6\); the
entries were generated from \eqref{plt:eq:bell-comtet-rec} and checked against
\eqref{plt:eq:bell-partial-def}.  The first column is
\(\ExponentialPartialBellPolynomial{n}{1}=x_{n}\), the diagonal is
\(x_{1}^{n}\) and the subdiagonal is \(\binom{n}{2}x_{1}^{n-2}x_{2}\), as
\eqref{plt:eq:bell-boundary} predicts.  Each entry is homogeneous of degree
\(k\) and isobaric of weight \(n\), which is the cheapest available check on a
hand computation: in
\(\ExponentialPartialBellPolynomial{6}{3}
=15x_{1}^{2}x_{4}+60x_{1}x_{2}x_{3}+15x_{2}^{3}\) every monomial has three
factors and subscript sum six.  A second check is that setting all
\(x_{j}=1\) must return the Stirling number
\(\StirlingSecondKind{n}{k}\) of \cref{plt:thm:bell-stirling}\,(ii); for
\(n=6\), \(k=3\) this reads \(15+60+15=90\).

\begin{table}[htbp]\centering\small
\caption{The exponential partial Bell polynomials
 $\ExponentialPartialBellPolynomial{n}{k}(x_{1},x_{2},\ldots)$
 for $1\le k\le n\le6$.}\label{plt:tab:bell-partial}
\begin{tabular}{@{}r l l@{}}
\toprule
$n$ & $k$ & $\ExponentialPartialBellPolynomial{n}{k}$\\
\midrule
$1$ & $1$ & $x_{1}$\\
\midrule
$2$ & $1$ & $x_{2}$\\
    & $2$ & $x_{1}^{2}$\\
\midrule
$3$ & $1$ & $x_{3}$\\
    & $2$ & $3x_{1}x_{2}$\\
    & $3$ & $x_{1}^{3}$\\
\midrule
$4$ & $1$ & $x_{4}$\\
    & $2$ & $4x_{1}x_{3}+3x_{2}^{2}$\\
    & $3$ & $6x_{1}^{2}x_{2}$\\
    & $4$ & $x_{1}^{4}$\\
\midrule
$5$ & $1$ & $x_{5}$\\
    & $2$ & $5x_{1}x_{4}+10x_{2}x_{3}$\\
    & $3$ & $10x_{1}^{2}x_{3}+15x_{1}x_{2}^{2}$\\
    & $4$ & $10x_{1}^{3}x_{2}$\\
    & $5$ & $x_{1}^{5}$\\
\midrule
$6$ & $1$ & $x_{6}$\\
    & $2$ & $6x_{1}x_{5}+15x_{2}x_{4}+10x_{3}^{2}$\\
    & $3$ & $15x_{1}^{2}x_{4}+60x_{1}x_{2}x_{3}+15x_{2}^{3}$\\
    & $4$ & $20x_{1}^{3}x_{3}+45x_{1}^{2}x_{2}^{2}$\\
    & $5$ & $15x_{1}^{4}x_{2}$\\
    & $6$ & $x_{1}^{6}$\\
\bottomrule
\end{tabular}
\end{table}
\end{example}

\section{Stirling evaluations, and the link to the Lambert polynomials}
\label{plt:sec:bell-stirling}

The unsigned Stirling numbers of the first kind were fixed in
\eqref{plt:eq:om-stirling-def} by
\(\FallingFactorial{z}{n}=\sum_{k}\stirlingsigned{n}{k}z^{k}\) and
\(\stirlingone{n}{k}=(-1)^{n-k}\stirlingsigned{n}{k}\).  We first record the
exponential generating function that converts that definition into a Bell
evaluation.

\begin{lemma}[Rising-factorial and Stirling generating identities]
\label{plt:lem:bell-rising}
In \(\RationalNumbers[z][[u]]\),
\begin{equation}\label{plt:eq:bell-rising-egf}
 \sum_{n\ge0}\RisingFactorial{z}{n}\,\frac{u^{n}}{n!}
 =\exp\Bigl(z\log\frac{1}{1-u}\Bigr),
 \qquad
 \RisingFactorial{z}{n}=z(z+1)\cdots(z+n-1),
\end{equation}
where \(\log\frac{1}{1-u}\DefinitionEquals\sum_{i\ge1}u^{i}/i\) and
\(\exp\) is the formal exponential of a series with zero constant term.
Consequently \(\RisingFactorial{z}{n}=\sum_{k=0}^{n}\stirlingone{n}{k}z^{k}\)
and, for every \(k\in\NaturalNumbers\),
\begin{equation}\label{plt:eq:bell-stirling-egf}
 \frac{1}{k!}\Bigl(\log\frac{1}{1-u}\Bigr)^{k}
 =\sum_{n\ge k}\stirlingone{n}{k}\,\frac{u^{n}}{n!} .
\end{equation}
\end{lemma}

\begin{proof}
The unsigned form of \eqref{plt:eq:om-stirling-def} is obtained by replacing
\(z\) with \(-z\): since
\(\FallingFactorial{(-z)}{n}=(-z)(-z-1)\cdots(-z-n+1)
=(-1)^{n}\RisingFactorial{z}{n}\), we get
\(\RisingFactorial{z}{n}=(-1)^{n}\sum_{k}\stirlingsigned{n}{k}(-z)^{k}
=\sum_{k}(-1)^{n-k}\stirlingsigned{n}{k}z^{k}
=\sum_{k}\stirlingone{n}{k}z^{k}\).

For \eqref{plt:eq:bell-rising-egf} we compare two solutions of one linear
differential equation in \(\RationalNumbers[z][[u]]\).  A series
\(W=\sum_{n\ge0}w_{n}u^{n}/n!\) satisfies
\begin{equation}\label{plt:eq:bell-rising-ode}
 (1-u)\,\partial_{u}W=zW,
 \qquad w_{0}=1,
\end{equation}
if and only if \(w_{n+1}=(z+n)w_{n}\) for every \(n\ge0\); indeed
\(\partial_{u}W=\sum_{n\ge0}w_{n+1}u^{n}/n!\) and
\(u\,\partial_{u}W=\sum_{n\ge1}n\,w_{n}u^{n}/n!\), so comparing coefficients
of \(u^{n}/n!\) in \eqref{plt:eq:bell-rising-ode} gives
\(w_{n+1}-nw_{n}=zw_{n}\).  In particular
\eqref{plt:eq:bell-rising-ode} has exactly one solution, and
\(w_{n}=\RisingFactorial{z}{n}\) is it, because
\(\RisingFactorial{z}{0}=1\) and
\(\RisingFactorial{z}{n+1}=(z+n)\RisingFactorial{z}{n}\).

Now put \(\ell(u)\DefinitionEquals\log\frac{1}{1-u}=\sum_{i\ge1}u^{i}/i\), so
that \(\partial_{u}\ell=\sum_{i\ge1}u^{i-1}=(1-u)^{-1}\), and
\(Z\DefinitionEquals\exp\bigl(z\ell(u)\bigr)
=\sum_{k\ge0}z^{k}\ell(u)^{k}/k!\), a summable family in
\(\RationalNumbers[z][[u]]\) because \(\ell^{k}\in u^{k}\RationalNumbers[[u]]\).
Termwise differentiation of that summable family gives
\(\partial_{u}Z=z(\partial_{u}\ell)Z=zZ/(1-u)\), that is
\((1-u)\partial_{u}Z=zZ\), and \(Z\) has constant term \(1\).  By the
uniqueness just proved, \(Z=\sum_{n\ge0}\RisingFactorial{z}{n}u^{n}/n!\),
which is \eqref{plt:eq:bell-rising-egf}.

Finally, compare the coefficient of \(z^{k}\) on the two sides of
\eqref{plt:eq:bell-rising-egf}: on the left it is
\(\sum_{n\ge0}\stirlingone{n}{k}u^{n}/n!\) by the expansion
\(\RisingFactorial{z}{n}=\sum_{k}\stirlingone{n}{k}z^{k}\) just established,
and on the right it is \(\ell(u)^{k}/k!\).  This is
\eqref{plt:eq:bell-stirling-egf}; the sum there may be started at \(n=k\)
because \(\stirlingone{n}{k}=0\) for \(k>n\).
\end{proof}

\begin{theorem}[Stirling evaluations of the Bell polynomials]
\label{plt:thm:bell-stirling}
For all \(n,k\in\NaturalNumbers\),
\begin{enumerate}[label=\textup{(\roman*)}]
\item \(\ExponentialPartialBellPolynomial{n}{k}
  \bigl(0!,1!,2!,\ldots\bigr)=\stirlingone{n}{k}\), the unsigned Stirling
  number of the first kind; equivalently
  \(\ExponentialPartialBellPolynomial{n}{k}\bigl((j-1)!\bigr)_{j\ge1}
  =(-1)^{n-k}\stirlingsigned{n}{k}\);
\item \(\ExponentialPartialBellPolynomial{n}{k}(1,1,1,\ldots)
  =\StirlingSecondKind{n}{k}\), the Stirling number of the second kind;
\item \(\ExponentialCompleteBellPolynomial{n}(0!,1!,2!,\ldots)=n!\).
\end{enumerate}
\end{theorem}

\begin{proof}
(i)  Apply \eqref{plt:eq:bell-egf} over
\(R=\RationalNumbers\) with \(\xi_{j}=(j-1)!\).  Then
\(\Xi(u)=\sum_{j\ge1}(j-1)!\,u^{j}/j!=\sum_{j\ge1}u^{j}/j
=\log\frac{1}{1-u}\), so \eqref{plt:eq:bell-egf} reads
\[
 \frac1{k!}\Bigl(\log\frac{1}{1-u}\Bigr)^{k}
 =\sum_{n\ge k}\ExponentialPartialBellPolynomial{n}{k}
   \bigl(0!,1!,2!,\ldots\bigr)\frac{u^{n}}{n!} .
\]
Comparing with \eqref{plt:eq:bell-stirling-egf} and using uniqueness of
coefficients in \(\RationalNumbers[[u]]\) gives (i).  The second form of (i)
is then the sign convention \eqref{plt:eq:om-stirling-def}.

(ii)  With \(\xi_{j}=1\) for all \(j\) one gets
\(\Xi(u)=\sum_{j\ge1}u^{j}/j!=\EulerE^{u}-1\), and
\eqref{plt:eq:bell-egf} becomes
\(\frac1{k!}(\EulerE^{u}-1)^{k}
=\sum_{n\ge k}\ExponentialPartialBellPolynomial{n}{k}(1,1,\ldots)u^{n}/n!\).
The left-hand side is the classical exponential generating function of the
Stirling numbers of the second kind, which we verify from the definition
\(\StirlingSecondKind{n}{k}
=\frac1{k!}\sum_{i=0}^{k}(-1)^{k-i}\binom{k}{i}i^{n}\).  Expanding
\((\EulerE^{u}-1)^{k}\) by the binomial theorem gives
\[
 \bigl(\EulerE^{u}-1\bigr)^{k}
 =\sum_{i=0}^{k}(-1)^{k-i}\binom{k}{i}\EulerE^{iu}
 =\sum_{n\ge0}\Bigl(\sum_{i=0}^{k}(-1)^{k-i}\binom{k}{i}i^{n}\Bigr)
   \frac{u^{n}}{n!},
\]
so \([u^{n}/n!]\frac{1}{k!}(\EulerE^{u}-1)^{k}=\StirlingSecondKind{n}{k}\).

(iii)  Sum (i) over \(0\le k\le n\) and use
\(\sum_{k=0}^{n}\stirlingone{n}{k}=\RisingFactorial{1}{n}=1\cdot2\cdots n=n!\),
which is the polynomial identity
\(\RisingFactorial{z}{n}=\sum_{k}\stirlingone{n}{k}z^{k}\) of
\cref{plt:lem:bell-rising} evaluated at \(z=1\).
\end{proof}

\begin{corollary}[The Lambert triangle is a slice of the Bell array]
\label{plt:cor:bell-lambert}
For \(n\ge1\) and \(1\le m\le n\), the coefficients
\(\lambda_{n,m}=[u^{m}]\LamP{n}(u)\) of the Lambert polynomials
(\cref{plt:def:lambert-polynomials,plt:thm:lambert-stirling}) are
\begin{equation}\label{plt:eq:bell-lambert}
\begin{split}
 \lambda_{n,m}
 &=\frac{(-1)^{n-m}}{m!}\,
  \ExponentialPartialBellPolynomial{n}{\,n-m+1}\bigl(0!,1!,2!,\ldots\bigr),\\
 \LamP{n}(u)
 &=\sum_{m=1}^{n}\frac{(-1)^{n-m}}{m!}\,
  \ExponentialPartialBellPolynomial{n}{\,n-m+1}
   \bigl(0!,1!,2!,\ldots\bigr)\,u^{m}.
\end{split}
\end{equation}
Reading the array by antidiagonals, the \(n\)-th Lambert polynomial consumes
exactly the \(n\)-th row of the Bell array evaluated at the factorial
arguments, in reversed column order and with alternating signs.
\end{corollary}

\begin{proof}
Substitute \cref{plt:thm:bell-stirling}\,(i) with \(k=n-m+1\) into
\eqref{plt:eq:lambert-lambda}.
\end{proof}

\begin{remark}[Why the factorial arguments are the right ones]
\label{plt:rmk:bell-why-factorials}
The evaluation \(x_{j}\mapsto(j-1)!\) looks arbitrary until one notices what
it encodes.  By the proof of \cref{plt:thm:bell-stirling}\,(i) it is exactly
the evaluation for which the generating series \(\Xi\) of
\cref{plt:prop:bell-gf} becomes a logarithm.  The Lambert polynomials arise
from reverting \(F=X+L\), that is from a perturbation which \emph{is} the
logarithm (\cref{plt:thm:om-expansion}); the appearance of cycle numbers in
\eqref{plt:eq:bell-lambert} is therefore not a coincidence of two unrelated
triangles but the same statement twice.  The table of
\(\stirlingone{n}{n-m+1}\) in \cref{plt:tab:lambert-triangle} is
simultaneously a table of Bell polynomials at factorial arguments.  Note that
the correspondence is between the \emph{unsigned} cycle numbers and the Bell
array: the signs in \eqref{plt:eq:bell-lambert} come from the alternating
prefactor of \eqref{plt:eq:lambert-operator} and not from the Bell
polynomials, which have nonnegative integer coefficients by
\eqref{plt:eq:bell-partial-def}.
\end{remark}

\section{Fa\`a di Bruno with a scalar-analytic outer function}
\label{plt:sec:bell-faa}

\Cref{plt:prop:tay-faa-di-bruno} computes \(\DX^{n}(F\compose G)\) when
\(F\) and \(G\) are both polynomial--logarithmic and \(\compose\) is the
composition of \cref{plt:sec:composition}.  The volume also needs the
complementary case, in which the outer object is not an element of
\(\PLlau\) at all but a scalar power series \(\varphi\in\K[[z]]\) applied to a
polynomial--logarithmic argument through the functional calculus of
\cref{plt:prop:tay-scalar-outer}.  The two statements have the same proof, and
we isolate that proof once.

\begin{lemma}[Abstract Fa\`a di Bruno for a Leibniz tower]
\label{plt:lem:bell-abstract-faa}
Let \(R\) be a commutative \(\RationalNumbers\)-algebra, \(\partial\) a
derivation of \(R\), \(\gamma\in R\), and let \((u_{k})_{k\ge0}\) be a sequence
in \(R\) satisfying
\begin{equation}\label{plt:eq:bell-leibniz-tower}
 \partial u_{k}=u_{k+1}\,\gamma\qquad(k\ge0).
\end{equation}
Then for every \(n\ge1\),
\begin{equation}\label{plt:eq:bell-abstract-faa}
 \partial^{\,n}u_{0}
 =\sum_{k=1}^{n}u_{k}\;
   \ExponentialPartialBellPolynomial{n}{k}
   \bigl(\gamma,\partial\gamma,\ldots,\partial^{\,n-k}\gamma\bigr).
\end{equation}
\end{lemma}

\begin{proof}
Induction on \(n\).  For \(n=1\), \eqref{plt:eq:bell-abstract-faa} reads
\(\partial u_{0}=u_{1}\ExponentialPartialBellPolynomial{1}{1}(\gamma)
=u_{1}\gamma\), which is \eqref{plt:eq:bell-leibniz-tower} at \(k=0\).

Assume \eqref{plt:eq:bell-abstract-faa} for some \(n\ge1\) and abbreviate
\(\beta_{n,k}\DefinitionEquals
\ExponentialPartialBellPolynomial{n}{k}(\gamma,\partial\gamma,\ldots)\), the
evaluation of the polynomial
\(\ExponentialPartialBellPolynomial{n}{k}\) at
\(x_{i}=\partial^{\,i-1}\gamma\).  Since \(\partial\) is a derivation and
\(\ExponentialPartialBellPolynomial{n}{k}\) is a polynomial in finitely many
of the \(x_{i}\), the ordinary chain rule for polynomials gives
\[
 \partial\beta_{n,k}
 =\sum_{i\ge1}
   \bigl(\partial_{x_{i}}\ExponentialPartialBellPolynomial{n}{k}\bigr)
   (\gamma,\partial\gamma,\ldots)\cdot\partial^{\,i}\gamma
 =\bigl(\mathcal D\ExponentialPartialBellPolynomial{n}{k}\bigr)
   (\gamma,\partial\gamma,\ldots),
 \qquad
 \mathcal D=\sum_{i\ge1}x_{i+1}\partial_{x_{i}},
\]
because \(\partial^{\,i}\gamma\) is the value of \(x_{i+1}\).  Hence, using
the Leibniz rule and \eqref{plt:eq:bell-leibniz-tower},
\[
 \partial^{\,n+1}u_{0}
 =\sum_{k\ge1}\Bigl(u_{k+1}\gamma\,\beta_{n,k}
   +u_{k}\bigl(\mathcal D\ExponentialPartialBellPolynomial{n}{k}\bigr)
   (\gamma,\partial\gamma,\ldots)\Bigr)
 =\sum_{k\ge1}u_{k}\Bigl[
   \gamma\,\ExponentialPartialBellPolynomial{n}{k-1}
   +\mathcal D\,\ExponentialPartialBellPolynomial{n}{k}\Bigr]
   (\gamma,\partial\gamma,\ldots),
\]
after shifting the index in the first sum (the term \(k=1\) of the shifted sum
carries \(\ExponentialPartialBellPolynomial{n}{0}=0\) for \(n\ge1\), so
nothing is lost).  By \cref{plt:lem:tay-bell-recurrence} evaluated at
\(x_{1}=\gamma\), the bracket is
\(\ExponentialPartialBellPolynomial{n+1}{k}\).  The sum truncates at
\(k=n+1\) by \cref{plt:lem:bell-elementary}\,(i), and the arguments needed are
\(x_{1},\ldots,x_{n+1-k+1}\), that is
\(\gamma,\ldots,\partial^{\,n+1-k}\gamma\).  All sums are finite.
\end{proof}

\begin{remark}[The two Fa\`a di Bruno statements are two towers]
\label{plt:rmk:bell-two-faa}
\Cref{plt:prop:tay-faa-di-bruno} is
\cref{plt:lem:bell-abstract-faa} in \(R=\PLlau\) with
\(\partial=\DX\), \(\gamma=\DX G\) and
\(u_{k}=\bigl(\DX^{k}F\bigr)\compose G\): the tower condition
\eqref{plt:eq:bell-leibniz-tower} is then exactly the chain rule
\cref{plt:thm:cmp-chain-rule}.  \Cref{plt:thm:bell-faa} below is the same
lemma with \(\gamma=\DX F\) and \(u_{k}=\varphi^{(k)}(F)\).  Nothing else
distinguishes the two; in particular neither is a corollary of the other,
since the towers are different.
\end{remark}

Before applying \cref{plt:lem:bell-abstract-faa} to a scalar-analytic outer
function we must know that the required tower exists, that is, that the chain
rule holds for the functional calculus \(\varphi\mapsto\varphi(W)\) of
\cref{plt:prop:tay-scalar-outer}.  For a \emph{polynomial} \(\varphi\) this is
the Leibniz rule; for a power series it needs the continuity of \(\DX\).

% ed.: the draft wrote \varphi=\sum(\varphi_k/k!)z^k -- coefficients at the
% origin -- and then set \varphi(F)=\sum(\varphi_k/k!)W^k with W=F-z_0.  That
% is \varphi evaluated at F only when z_0=0; in general the coefficients of
% the k-th power of W are \varphi^{(k)}(z_0)/k!, not \varphi_k/k!.  (Take
% \varphi(z)=z, z_0=1, F=1+t: the displayed sum returns t, not F.)  The
% justification offered, "apply the translate z\mapsto\varphi(z_0+z)", is also
% unavailable: re-expanding a formal element of \K[[z]] about a centre
% z_0\ne0 requires an infinite sum in \K for every coefficient, and \K carries
% no topology here.  Both defects are repaired at once by making the centre
% part of the datum -- the outer function IS its expansion at z_0.  Nothing
% downstream changes: the proof only ever uses c_k=\varphi_k/k! and the
% relation \varphi'(F)=\sum_{k\ge1}kc_kW^{k-1}.
\begin{lemma}[Chain rule for a scalar-analytic substitution]
\label{plt:lem:bell-chain}
Let \(z_{0}\in\K\) and let the outer function be given \emph{by its expansion
about \(z_{0}\)}: that is, let \(\varphi_{0},\varphi_{1},\ldots\in\K\) be
arbitrary scalars and write, formally,
\(\varphi(z)=\sum_{k\ge0}\frac{\varphi_{k}}{k!}(z-z_{0})^{k}\), the datum being
the pair \(\bigl(z_{0},(\varphi_{k})_{k\ge0}\bigr)\).  Let \(F\in\PLpow\)
satisfy \(\ordt(F-z_{0})\ge1\), write \(W\DefinitionEquals F-z_{0}\in
t\PLpow\), and put
\begin{equation}\label{plt:eq:bell-outer-eval}
 \varphi(F)\DefinitionEquals\sum_{k\ge0}\frac{\varphi_{k}}{k!}\,W^{k},
 \qquad
 \varphi^{(k)}(F)\DefinitionEquals
   \sum_{i\ge0}\frac{\varphi_{k+i}}{i!}\,W^{i}
 \qquad(k\ge0),
\end{equation}
each well defined by \cref{plt:prop:tay-scalar-outer} applied to the series
\(\sum_{k\ge0}(\varphi_{k}/k!)z^{k}\in\K[[z]]\) and to \(W\).  When
\(z_{0}=0\) this is the ordinary reading \(\varphi(F)=\varphi(W)\); when
\(z_{0}\ne0\) the centre must be part of the datum, since an element of
\(\K[[z]]\) presented at the origin cannot be re-expanded about \(z_{0}\)
inside \(\K[[z]]\).  Then
\begin{equation}\label{plt:eq:bell-chain}
 \DX\bigl(\varphi(F)\bigr)=\varphi'(F)\,\DX F .
\end{equation}
If \(\varphi\in\K[z]\) is a polynomial, \eqref{plt:eq:bell-chain} holds for
every \(F\in\PLlau\), with \(\varphi(F)\) the ordinary polynomial evaluation
and no valuation hypothesis.  The two readings of \(\varphi(F)\) agree when
both apply, since the Taylor expansion of a polynomial about \(z_{0}\) is
exact and finite.
\end{lemma}

\begin{proof}
Each \(\varphi^{(k)}(F)\) is defined by \eqref{plt:eq:bell-outer-eval}, the
family there being summable because \(\ordt(W^{i})\ge i\).  Write
\(c_{k}=\varphi_{k}/k!\in\K\), so that
\(\varphi(F)=\sum_{k\ge0}c_{k}W^{k}\) and
\(\varphi'(F)=\sum_{k\ge1}kc_{k}W^{k-1}\), both families being summable
because \(\ordt(W^{k})\ge k\) (\cref{plt:thm:ring-valuation-laws},
\cref{plt:lem:dif-summable}).  Note \(\DX W=\DX F\), because
\(\DX z_{0}=0\) for a scalar \(z_{0}\) (\cref{plt:cor:dif-kernel}).

The identity \eqref{plt:eq:bell-chain} is checked one block at a time.  Fix
\(N\ge0\) and split \(\varphi(F)=P_{N}+T_{N}\) with
\(P_{N}=\sum_{k\le N+1}c_{k}W^{k}\) and \(\ordt(T_{N})\ge N+2\).  Since
\(\DX\) is a \(\K\)-derivation of \(\PLlau\) with
\(\ordt(\DX H)\ge\ordt H+1\) (\cref{plt:thm:dif-derivation},
\cref{plt:thm:dif-valuation}\,(1)), we get \(\ordt(\DX T_{N})\ge N+3\), so
\[
 [t^{N}]\,\DX\bigl(\varphi(F)\bigr)
 =[t^{N}]\,\DX P_{N}
 =\sum_{k=1}^{N+1}kc_{k}\,[t^{N}]\bigl(W^{k-1}\DX W\bigr),
\]
the last step being the Leibniz rule applied to the \emph{finite} sum
\(P_{N}\).  On the other side, \(\ordt(W^{k-1}\DX W)\ge(k-1)+2=k+1\), so the
terms with \(k\ge N+2\) do not reach the block \(t^{N}\) and
\[
 [t^{N}]\bigl(\varphi'(F)\,\DX F\bigr)
 =\sum_{k\ge1}kc_{k}\,[t^{N}]\bigl(W^{k-1}\DX W\bigr)
 =\sum_{k=1}^{N+1}kc_{k}\,[t^{N}]\bigl(W^{k-1}\DX W\bigr).
\]
The two agree for every \(N\), and an element of \(\PLpow\) is determined by
its blocks (\cref{plt:prop:trunc-extraction}), which proves
\eqref{plt:eq:bell-chain}.  For a polynomial \(\varphi\) the sums are finite
from the start and only the Leibniz rule is used, so no valuation hypothesis
is needed and the identity holds in \(\PLlau\).
\end{proof}

\begin{theorem}[Fa\`a di Bruno for a scalar-analytic outer function]
\label{plt:thm:bell-faa}
Let \(z_{0}\in\K\), let \(\varphi\) be given by its expansion about \(z_{0}\)
as in \cref{plt:lem:bell-chain}, and let \(F\in\PLpow\) satisfy
\(\ordt(F-z_{0})\ge1\); or let \(\varphi\in\K[z]\) be a polynomial and
\(F\in\PLlau\) arbitrary.  Then for every \(n\ge1\)
\begin{equation}\label{plt:eq:bell-faa}
 \DX^{\,n}\bigl(\varphi(F)\bigr)
 =\sum_{k=1}^{n}\varphi^{(k)}(F)\;
   \ExponentialPartialBellPolynomial{n}{k}
   \bigl(\DX F,\DX^{2}F,\ldots,\DX^{\,n-k+1}F\bigr),
\end{equation}
an identity in \(\PLpow\) in the first case and in \(\PLlau\) in the second.
\end{theorem}

\begin{proof}
Apply \cref{plt:lem:bell-abstract-faa} in \(R=\PLlau\) with
\(\partial=\DX\), \(\gamma=\DX F\) and \(u_{k}=\varphi^{(k)}(F)\).  The tower
condition \eqref{plt:eq:bell-leibniz-tower} is
\(\DX\bigl(\varphi^{(k)}(F)\bigr)=\varphi^{(k+1)}(F)\,\DX F\), which is
\cref{plt:lem:bell-chain} applied to \(\varphi^{(k)}\) in place of
\(\varphi\).  The arguments of the Bell polynomial are
\(x_{i}=\partial^{\,i-1}\gamma=\DX^{\,i}F\), \(1\le i\le n-k+1\).  Membership
in \(\PLpow\) in the first case follows because \(\PLpow\) is closed under
\(\DX\) (\cref{plt:thm:dif-derivation}\,(i)) and under the substitution
(\cref{plt:prop:tay-scalar-outer}).
\end{proof}

\begin{corollary}[Derivatives of a power, and the Lagrange terms]
\label{plt:cor:bell-power-derivative}
For every \(A\in\PLlau\), every \(r\in\NaturalNumbers\) and every \(n\ge1\),
\begin{equation}\label{plt:eq:bell-power-derivative}
 \DX^{\,n}\bigl(A^{r}\bigr)
 =\sum_{k=1}^{\min\SetOf{n,r}}
   \FallingFactorial{r}{k}\;A^{\,r-k}\;
   \ExponentialPartialBellPolynomial{n}{k}
   \bigl(\DX A,\DX^{2}A,\ldots,\DX^{\,n-k+1}A\bigr),
\end{equation}
where \(\FallingFactorial{r}{k}=r(r-1)\cdots(r-k+1)\) is the falling
factorial, and where the terms with \(k>r\) are absent because
\(\FallingFactorial{r}{k}=0\) there.
\end{corollary}

\begin{proof}
Take \(\varphi(z)=z^{r}\) in the polynomial case of
\cref{plt:thm:bell-faa}; then
\(\varphi^{(k)}(z)=\FallingFactorial{r}{k}z^{\,r-k}\) for \(k\le r\) and
\(\varphi^{(k)}=0\) for \(k>r\).
\end{proof}

\begin{remark}[Valuation bookkeeping]
\label{plt:rmk:bell-valuation}
In the setting of \cref{plt:thm:bell-faa} with \(\ordt(F-z_{0})\ge1\), one has
\(\ordt(\DX^{\,i}F)\ge i+1\) by \cref{plt:thm:dif-valuation}, so a monomial of
\(\ExponentialPartialBellPolynomial{n}{k}\) with multiplicity vector
\((m_{j})\) evaluates to something of order at least
\(\sum_{j}m_{j}(j+1)=n+k\).  The \(k\)-th summand of
\eqref{plt:eq:bell-faa} therefore has \(\ordt\ge n+k\ge n+1\), matching the
bound \(\ordt\bigl(\DX^{\,n}(\varphi(F))\bigr)\ge n+1\) that
\cref{plt:thm:dif-valuation}\,(2) gives directly, since
\(\varphi(F)-\varphi(z_{0})\in t\PLpow\).  Consequently only
\(k\le N-n\) can affect the block \(t^{N}\): the combinatorial size of
\eqref{plt:eq:bell-faa} at a fixed target block is bounded independently of
any analytic consideration.
\end{remark}

\chapter{Closed-form coefficient extraction}
\label{plt:sec:coefficient-extraction}

Parts~II--IV compute coefficients by \emph{recurrences}: the reciprocal by
\eqref{plt:eq:div-recip-cn}, the quotient by
\eqref{plt:eq:div-quotient-recurrence}, the convolution polynomials by
\eqref{plt:eq:tay-C-recurrence}, the inverse by
\cref{plt:thm:rev-triangular}.  A recurrence is the right tool when all
coefficients up to a cutoff are wanted, and it is what an implementation
should run (\cref{plt:sec:algorithms}).  It is the wrong tool when a
\emph{single} block is wanted, when the block is wanted as an explicit
function of the input blocks, or when one has to prove something about the
block --- a degree bound, a sign, an integrality statement --- for all \(N\)
at once.

This chapter derives the non-recursive counterparts.  By a \emph{closed form}
we mean a formula for \([t^{N}]\) or \([t^{N}L^{d}]\) of the result as a
finite sum, indexed by an explicitly bounded set, of products of the input
blocks and of universal combinatorial constants, with no reference to
previously computed output blocks.  Every such formula below is derived from
the body of the volume and from \cref{plt:sec:bell}; none is quoted.

\section{The multi-index bookkeeping}
\label{plt:sec:cext-bookkeeping}

Two separate finiteness mechanisms are in play, and conflating them is the
source of the ``infinite sum'' complaints that the sources occasionally raise
against these formulas.

\begin{lemma}[The two finiteness mechanisms]
\label{plt:lem:cext-finiteness}
Let \(F=\sum_{n\ge n_{0}}p_{n}(L)t^{n}\in\PLlau\) and fix
\(N\in\IntegerNumbers\), \(d\in\NaturalNumbers\).
\begin{enumerate}[label=\textup{(\roman*)}]
\item \textup{(Outer, by valuation.)}  If \((S_{i})_{i\in I}\) is a summable
  family in \(\PLlau\), then \(\SetOf{i\in I:[t^{N}]S_{i}\ne0}\) is finite and
  \([t^{N}]\sum_{i}S_{i}=\sum_{i}[t^{N}]S_{i}\).
\item \textup{(Inner, by degree.)}  Each block \(p_{n}\) lies in \(\K[L]\),
  hence has finite \(\degL\), so \([t^{N}L^{d}]F=0\) for
  \(d>\dprofile{F}(N)\) and the block \(p_{N}\) is a finite sum of monomials.
\item The two are independent: a summable family may have blocks of unbounded
  \(\degL\) at a fixed level of the partial sums --- for instance
  \(S_{i}=t^{i}L^{i}\) is summable although
  \(\dprofile{\sum_{i\le M}S_{i}}\) is unbounded in \(M\) --- and conversely a
  family of bounded \(\degL\) need not be summable.  Only \textup{(i)} is
  part of the \(t\)-adic topology; a bound on \(\degL\) is an additional
  approximation and must be stated separately.
\end{enumerate}
\end{lemma}

\begin{proof}
(i)  is \cref{plt:thm:trunc-summability} together with
\cref{plt:prop:trunc-extraction}.  (ii)  is
\cref{plt:def:lead-profile} and the definition of \(\K[L]\).  For the first
half of (iii), \(\ordt(t^{i}L^{i})=i\to\infty\), so the family is summable,
while \(\dprofile{\ }\) of the \(M\)-th partial sum takes the value \(M\) at
level \(M\); for the second half, the constant family \(S_{i}=1\) has
\(\degL S_{i}=0\) and is not summable.
\end{proof}

\begin{warningbox}[title={The formal tail does not bound a logarithmic degree}]
The statement \(F=G+\FormalBigOAt{t^{N}}{t}\) says nothing whatever about
\(\degL\) of the discarded blocks (\cref{plt:def:trunc-formal-o}); truncation
is exclusive in \(t\) and \emph{whole} in \(L\)
(\cref{plt:def:trunc-exclusive}).  Every closed formula below therefore
returns a complete polynomial block, never a polynomial block truncated in
\(L\).  Discarding a high power of \(L\) inside a retained block is a second,
independent approximation and requires its own error statement.
\end{warningbox}

\section{Products}
\label{plt:sec:cext-product}

\begin{proposition}[Scalar coefficients of a product]
\label{plt:prop:cext-product}
Let \(F,G\in\PLlau\) be nonzero with \(a=\ordt F\), \(b=\ordt G\), and write
\(a_{n,e}=[t^{n}L^{e}]F\), \(b_{m,f}=[t^{m}L^{f}]G\) in the notation of
\cref{plt:def:trunc-extraction}.  Then for all \(N\in\IntegerNumbers\) and
\(d\in\NaturalNumbers\),
\begin{equation}\label{plt:eq:cext-product}
 [t^{N}L^{d}](FG)
 =\sum_{n=a}^{N-b}\;\sum_{e=0}^{d}\,a_{n,e}\,b_{N-n,\,d-e},
\end{equation}
a sum of at most \((N-a-b+1)(d+1)\) products, empty when \(N<a+b\).  The
nonzero terms satisfy the sharper constraints
\(e\le\dprofile{F}(n)\) and \(d-e\le\dprofile{G}(N-n)\); in particular
\([t^{N}L^{d}](FG)=0\) whenever
\(d>(\dprofile{F}\ast\dprofile{G})(N)
=\max_{n+m=N}\bigl(\dprofile{F}(n)+\dprofile{G}(m)\bigr)\).
\end{proposition}

\begin{proof}
By \cref{plt:def:mul-cauchy} and \cref{plt:lem:mul-finite-fiber},
\([t^{N}](FG)=\sum_{n=a}^{N-b}p_{n}q_{N-n}\), a finite sum in \(\K[L]\).
Extracting \([L^{d}]\) of a product of two polynomials is the Cauchy product
in \(L\), which gives \eqref{plt:eq:cext-product}; the index \(e\) is confined
to \([0,d]\) because both factors are polynomials, hence have no negative
powers of \(L\).  The sharper constraints and the vanishing statement are the
degree bound \eqref{plt:eq:mul-degree-bound} of
\cref{plt:thm:mul-degree-bound}, whose \(\ast\) is the \((\max,+)\)-convolution
displayed above.
\end{proof}

Formula \eqref{plt:eq:cext-product} is a genuine two-level convolution: outer
in \(t\) by valuation, inner in \(L\) by degree.  Its cost, and the fact that
the resulting scalar cost is \emph{not} \(\BigO(N^{2})\) once the inner level
is counted, are analysed in \cref{plt:prop:alg-convolution-cost} and
\cref{plt:cor:alg-cost-not-quadratic}.

\section{Powers, reciprocals, and the convolution polynomials}
\label{plt:sec:cext-power}

The convolution polynomials \(C_{m,r}\) of \cref{plt:def:tay-convolution}
--- defined by \eqref{plt:eq:tay-C-def} with the empty-product conventions
\(C_{0,0}=1\), \(C_{m,0}=0\) for \(m>0\), and computed by the recurrence
\eqref{plt:eq:tay-C-recurrence} --- are the blocks of the powers of a series.
\Cref{plt:prop:tay-convolution}\,(iv) gives their closed form only under the
hypothesis \(h_{0}=0\).  The general case is needed whenever the inner series
has a nonzero level-zero block, which is the normal situation for a
near-identity map.

\begin{proposition}[Closed form of the convolution polynomials, general case]
\label{plt:prop:cext-convolution-closed}
Let \(H=\sum_{j\ge0}h_{j}(L)t^{j}\in\PLpow\) and let \(C_{m,r}\) be as in
\eqref{plt:eq:tay-C-def}.  Then for all \(m,r\in\NaturalNumbers\),
\begin{equation}\label{plt:eq:cext-Cmr-closed}
 C_{m,r}
 =\sum_{k=0}^{\min\SetOf{m,r}}
   \binom{r}{k}\,h_{0}^{\,r-k}\,\frac{k!}{m!}\;
   \ExponentialPartialBellPolynomial{m}{k}
    \bigl(1!\,h_{1},2!\,h_{2},\ldots,(m-k+1)!\,h_{m-k+1}\bigr),
\end{equation}
where the term \(k=0\) carries
\(\ExponentialPartialBellPolynomial{m}{0}\) and is therefore \(h_{0}^{\,r}\)
% ed.: the draft said that for h_0=0 the formula "reduces to
% plt:eq:tay-C-bell", which is stated only for m>=r>=1; at r=0 it reduces
% instead to the boundary convention C_{m,0}=\delta_{m,0}.  Range added.
for \(m=0\) and \(0\) for \(m>0\).  When \(h_{0}=0\) only \(k=r\) survives:
for \(r\ge1\) formula \eqref{plt:eq:cext-Cmr-closed} then reduces to
\eqref{plt:eq:tay-C-bell}, and for \(r=0\) it returns the boundary values
\(C_{m,0}=\delta_{m,0}\) of \eqref{plt:eq:tay-C-def}.
\end{proposition}

\begin{proof}
Split \(H=h_{0}+H_{+}\) with \(H_{+}=\sum_{j\ge1}h_{j}t^{j}\in t\PLpow\).
Since \(h_{0}\) and \(H_{+}\) commute, the binomial theorem gives
\(H^{r}=\sum_{k=0}^{r}\binom{r}{k}h_{0}^{\,r-k}H_{+}^{k}\), a finite sum.  By
\cref{plt:prop:tay-convolution}\,(i) applied to \(H\) and by
\cref{plt:prop:mul-integer-power} applied to \(H_{+}\) (whose level-zero block
vanishes),
\[
 C_{m,r}=[t^{m}]H^{r}
 =\sum_{k=0}^{r}\binom{r}{k}h_{0}^{\,r-k}\,[t^{m}]H_{+}^{k},
 \qquad
 [t^{m}]H_{+}^{k}
 =\frac{k!}{m!}\,\ExponentialPartialBellPolynomial{m}{k}
   \bigl(1!h_{1},\ldots\bigr)\ \ (k\ge1),
\]
with \([t^{m}]H_{+}^{k}=0\) for \(m<k\) and \([t^{m}]H_{+}^{0}=\delta_{m,0}\).
Truncating the sum at \(k=\min\SetOf{m,r}\) is
\cref{plt:lem:bell-elementary}\,(i).  If \(h_{0}=0\) then, with the convention
\(0^{0}=1\), \(h_{0}^{\,r-k}=0\) unless \(k=r\); the surviving term is
\eqref{plt:eq:tay-C-bell} when \(r\ge1\), and is
\((0!/m!)\ExponentialPartialBellPolynomial{m}{0}=\delta_{m,0}\) when
\(r=0\).
\end{proof}

\begin{theorem}[Closed form for an arbitrary power of a unit]
\label{plt:thm:cext-power}
Let \(W=\sum_{j\ge1}w_{j}(L)t^{j}\in t\PLpow\) and \(\sigma\in\K\), and let
\((1+W)^{\sigma}\in\PLpow\) be the binomial power of
\cref{plt:thm:mul-binomial}.  Then \([t^{0}](1+W)^{\sigma}=1\) and, for
\(r\ge1\),
\begin{equation}\label{plt:eq:cext-power}
 [t^{r}](1+W)^{\sigma}
 =\frac{1}{r!}\sum_{k=1}^{r}
   \FallingFactorial{\sigma}{k}\;
   \ExponentialPartialBellPolynomial{r}{k}
    \bigl(1!\,w_{1},2!\,w_{2},\ldots,(r-k+1)!\,w_{r-k+1}\bigr),
\end{equation}
where \(\FallingFactorial{\sigma}{k}=\sigma(\sigma-1)\cdots(\sigma-k+1)\).
The right-hand side depends only on \(w_{1},\ldots,w_{r}\), and it is
\(\K\)-linear in \(w_{r}\) with coefficient \(\sigma\).

If \(F\in\PLlau\) is a unit, so that \(F=ct^{n_{0}}(1+W)\) with
\(c\in\K^{\times}\), \(n_{0}=\ordt F\) and \(W\in t\PLpow\) by
% ed.: the draft asserted "then F^\sigma = c^\sigma t^{\sigma n_0}(1+W)^\sigma"
% as though it were an identity.  For non-integer \sigma the symbol F^\sigma
% has no prior meaning in the volume, and the right-hand side depends on the
% chosen value of c^\sigma; it is a definition, and is now printed as one.
\cref{plt:thm:div-unit-criterion}, then, once a value \(c^{\sigma}\in\K\) is
chosen, one \emph{defines}
\(F^{\sigma}\DefinitionEquals c^{\sigma}t^{\sigma n_{0}}(1+W)^{\sigma}\); this
lies in \(\PLlau\) when \(\sigma n_{0}\in\IntegerNumbers\), and otherwise in
the Puiseux-log or Hahn extension of \cref{plt:prop:ring-puiseux-tower}.  For
\(\sigma\in\IntegerNumbers\) the definition returns the ring power of \(F\).
\end{theorem}

\begin{proof}
By \cref{plt:thm:mul-binomial}\,(i)--(ii),
\((1+W)^{\sigma}=\sum_{k\ge0}\binom{\sigma}{k}W^{k}\) with the family
summable, and
\([t^{r}](1+W)^{\sigma}=\sum_{k=0}^{r}\binom{\sigma}{k}[t^{r}]W^{k}\) since
\(\ordt(W^{k})\ge k\).  For \(k\ge1\),
\cref{plt:prop:mul-integer-power} (formula \eqref{plt:eq:mul-bell}, whose
hypothesis \(w_{0}=0\) holds) gives
\([t^{r}]W^{k}=\frac{k!}{r!}
\ExponentialPartialBellPolynomial{r}{k}(1!w_{1},2!w_{2},\ldots)\); for
\(k=0\), \([t^{r}]W^{0}=\delta_{r,0}\), which contributes nothing when
\(r\ge1\).  Since
\(\binom{\sigma}{k}k!=\FallingFactorial{\sigma}{k}\), summing gives
\eqref{plt:eq:cext-power}.

Dependence only on \(w_{1},\ldots,w_{r}\) is
\cref{plt:lem:bell-elementary}\,(i).  For the linear dependence on \(w_{r}\):
the variable \(x_{r}\) occurs in
\(\ExponentialPartialBellPolynomial{r}{k}\) only for \(k=1\) --- by
\cref{plt:lem:bell-elementary}\,(i) it can occur only when \(r\le r-k+1\),
i.e.\ \(k\le1\) --- and
\(\ExponentialPartialBellPolynomial{r}{1}=x_{r}\) by
\eqref{plt:eq:bell-boundary}, so the coefficient of \(w_{r}\) is
\(\frac{1}{r!}\FallingFactorial{\sigma}{1}\cdot r!\,=\sigma\).  This is
\eqref{plt:eq:mul-binomial-top} made explicit.

The factorization statement is \cref{plt:thm:div-unit-criterion}\,(3); the
exponent \(t^{\sigma n_{0}}\) lies in the integer-exponent scale exactly when
\(\sigma n_{0}\in\IntegerNumbers\), and otherwise requires the ramified tower
of \cref{plt:prop:ring-puiseux-tower}.  For \(\sigma\in\IntegerNumbers\) the
displayed definition returns the ring power, because \(c^{\sigma}\) and
\(t^{\sigma n_{0}}\) are then the ordinary powers and \((1+W)^{\sigma}\)
agrees with the \(\sigma\)-fold product by the exponent law of
\cref{plt:thm:mul-binomial}.
\end{proof}

\begin{corollary}[Closed form for a reciprocal]
\label{plt:cor:cext-reciprocal}
With \(W\) as in \cref{plt:thm:cext-power},
\begin{equation}\label{plt:eq:cext-reciprocal}
 [t^{r}]\recip{(1+W)}
 =\frac{1}{r!}\sum_{k=1}^{r}(-1)^{k}\,k!\;
   \ExponentialPartialBellPolynomial{r}{k}
    \bigl(1!\,w_{1},\ldots,(r-k+1)!\,w_{r-k+1}\bigr)
 \qquad(r\ge1).
\end{equation}
Consequently, for a unit \(F=ct^{n_{0}}(1+W)\) of \(\PLlau\),
\(\recip{F}=c^{-1}t^{-n_{0}}\bigl(1+\sum_{r\ge1}
[t^{r}]\recip{(1+W)}\,t^{r}\bigr)\).  In particular
\[
 [t^{1}]\recip{(1+W)}=-w_{1},
 \quad
 [t^{2}]\recip{(1+W)}=w_{1}^{2}-w_{2},
 \quad
 [t^{3}]\recip{(1+W)}=-w_{1}^{3}+2w_{1}w_{2}-w_{3}.
\]
\end{corollary}

\begin{proof}
Put \(\sigma=-1\) in \eqref{plt:eq:cext-power} and use
\(\FallingFactorial{(-1)}{k}=(-1)(-2)\cdots(-k)=(-1)^{k}k!\).  The
factorization is \cref{plt:thm:div-reciprocal}.  The three displayed values
follow from \cref{plt:tab:bell-partial}: for \(r=1\),
\(-1!\,x_{1}/1!\) at \(x_{1}=w_{1}\); for \(r=2\),
\(\bigl(-1!\,x_{2}+2!\,x_{1}^{2}\bigr)/2!\) at
\((x_{1},x_{2})=(w_{1},2w_{2})\), giving \(-w_{2}+w_{1}^{2}\); for \(r=3\),
\(\bigl(-x_{3}+2\cdot3x_{1}x_{2}-6x_{1}^{3}\bigr)/6\) at
\((x_{1},x_{2},x_{3})=(w_{1},2w_{2},6w_{3})\), giving
\(-w_{3}+2w_{1}w_{2}-w_{1}^{3}\).
\end{proof}

The Bell form \eqref{plt:eq:cext-reciprocal} is closed but its term count
grows like the partition function.  A second closed form, with a term count
that is a single determinant and which is the one to use for structural
arguments about signs and integrality, is the following.

\begin{proposition}[Toeplitz--Hessenberg determinant for the reciprocal]
\label{plt:prop:cext-hessenberg}
Let \(B=\sum_{n\ge0}B_{n}(L)t^{n}\in\PLpow\) with \(B_{0}=c\in\K^{\times}\),
and let \(\recip{B}=\sum_{r\ge0}C_{r}(L)t^{r}\) be its reciprocal
(\cref{plt:thm:div-reciprocal}).  For \(r\ge1\) let \(M_{r}\) be the
\(r\times r\) matrix over \(\K[L]\) with entries
\begin{equation}\label{plt:eq:cext-hessenberg-matrix}
 (M_{r})_{ij}\DefinitionEquals B_{\,i-j+1},
 \qquad 1\le i,j\le r,
 \qquad B_{m}\DefinitionEquals0\ \text{for }m<0 .
\end{equation}
Thus \(M_{r}\) is lower Hessenberg, with \(c\) on the superdiagonal, \(B_{1}\)
on the diagonal, and \(B_{r}\) in the lower-left corner.  Then
\begin{equation}\label{plt:eq:cext-hessenberg}
 C_{r}=\frac{(-1)^{r}}{c^{\,r+1}}\det M_{r}
 \qquad(r\ge1),
 \qquad C_{0}=c^{-1}.
\end{equation}
\end{proposition}

\begin{proof}
Put \(D_{0}\DefinitionEquals1\) and \(D_{r}\DefinitionEquals\det M_{r}\).  We
first show
\begin{equation}\label{plt:eq:cext-hessenberg-rec}
 D_{r}=-\sum_{m=1}^{r}(-1)^{m}B_{m}\,c^{\,m-1}D_{r-m}
 \qquad(r\ge1).
\end{equation}
Expand \(\det M_{r}\) along its last row, whose \(j\)-th entry is
\(B_{r-j+1}\):
\(D_{r}=\sum_{j=1}^{r}(-1)^{r+j}B_{r-j+1}\,\mu_{j}\), where \(\mu_{j}\) is the
minor obtained by deleting row \(r\) and column \(j\).  Its rows are indexed
by \(i\in\SetOf{1,\ldots,r-1}\) and its columns by
\(q\in\SetOf{1,\ldots,r}\setminus\SetOf{j}\), with entry \(B_{i-q+1}\).  If
\(i\le j-1\) and \(q\ge j+1\), then \(i-q+1\le j-1-j-1+1<0\), so that block
vanishes; the minor is therefore block lower-triangular with diagonal blocks
\begin{itemize}
\item rows \(1,\ldots,j-1\) against columns \(1,\ldots,j-1\), which is
  \(M_{j-1}\), of determinant \(D_{j-1}\); and
\item rows \(j,\ldots,r-1\) against columns \(j+1,\ldots,r\).  In the local
  indices \(a=i-j+1\) and \(b=q-j\), both running over
  \(1,\ldots,r-j\), its entry is \(B_{i-q+1}=B_{a-b}\); it therefore vanishes
  strictly above the diagonal and carries \(B_{0}=c\) on the diagonal, so it
  is lower triangular with determinant \(c^{\,r-j}\).
\end{itemize}
Hence \(\mu_{j}=D_{j-1}c^{\,r-j}\) and
\(D_{r}=\sum_{j=1}^{r}(-1)^{r+j}B_{r-j+1}c^{\,r-j}D_{j-1}\).  Substituting
\(m=r-j+1\), so that \(j-1=r-m\), \(r-j=m-1\) and
\((-1)^{r+j}=(-1)^{2r-m+1}=(-1)^{m+1}\), gives
\eqref{plt:eq:cext-hessenberg-rec}.

Now define \(\widetilde C_{0}=c^{-1}\) and
\(\widetilde C_{r}=(-1)^{r}c^{-(r+1)}D_{r}\) for \(r\ge1\).  Multiplying
\eqref{plt:eq:cext-hessenberg-rec} by \((-1)^{r}c^{-(r+1)}\) turns it into
\begin{align*}
 \widetilde C_{r}
 &=-\sum_{m=1}^{r}(-1)^{r+m}c^{\,m-1-(r+1)}B_{m}D_{r-m}
  =-c^{-1}\sum_{m=1}^{r}B_{m}\,(-1)^{r-m}c^{-(r-m+1)}D_{r-m}\\
 &=-c^{-1}\sum_{m=1}^{r}B_{m}\widetilde C_{r-m},
\end{align*}
using \((-1)^{r+m}=(-1)^{r-m}\) and \(D_{0}=1\) for the term \(m=r\).  This is
exactly the recurrence \eqref{plt:eq:div-recip-cn} characterizing the blocks
\(C_{r}\), and \(\widetilde C_{0}=C_{0}\) by
\eqref{plt:eq:div-recip-c0}; by the uniqueness part of
\cref{plt:thm:div-reciprocal}, \(\widetilde C_{r}=C_{r}\) for all \(r\).
\end{proof}

\begin{remark}[What the determinant form is good for]
\label{plt:rmk:cext-hessenberg-use}
\Cref{plt:prop:cext-hessenberg} makes two structural facts immediate that are
awkward to see from either \eqref{plt:eq:div-recip-cn} or
\eqref{plt:eq:cext-reciprocal}.  First, \(c^{\,r+1}C_{r}\) is a polynomial
with integer coefficients in \(B_{0},\ldots,B_{r}\), so the only denominators
in a reciprocal are powers of the leading scalar: reciprocation introduces no
new primes beyond those dividing \(c\).  Second, for \(r\ge1\),
\[
 \degL C_{r}
 \le\max\Bigl\{\textstyle\sum_{i=1}^{s}\dprofile{B}(m_{i})\ :\
   s\ge1,\ m_{i}\ge1,\ m_{1}+\cdots+m_{s}=r\Bigr\},
\]
because a signed product in the expansion of \(\det M_{r}\) is
\(\prod_{i=1}^{r}B_{\,i-\sigma(i)+1}\) for a permutation \(\sigma\), whose
subscripts sum to \(\sum_{i}(i-\sigma(i)+1)=r\); the factors with subscript
\(0\) are the scalar \(c\) and contribute nothing to \(\degL\), and the
remaining factors are indexed by a composition of \(r\) into positive parts.
Setting \(\dprofile{B}(m)\le\alpha m\) makes the right-hand side at most
\(\alpha r\), which is the linear bound of
\cref{plt:prop:div-degree-profile}, now obtained without a recursion.
Neither observation requires expanding the determinant.
\end{remark}

\section{Composition}
\label{plt:sec:cext-composition}

\Cref{plt:thm:tay-coefficients} already gives a closed formula, in terms of
the convolution polynomials.  Substituting
\cref{plt:prop:cext-convolution-closed} makes it closed in terms of the
\emph{input blocks} alone.

\begin{proposition}[Composition coefficients in Bell form]
\label{plt:prop:cext-composition-bell}
Let \(F=\sum_{n\ge n_{0}}p_{n}(L)t^{n}\in\PLlau\) and
\(H=\sum_{j\ge1}h_{j}(L)t^{j}\in t\PLpow\) (so \(h_{0}=0\)).  Then for every
\(N\in\IntegerNumbers\),
\begin{equation}\label{plt:eq:cext-composition-bell}
 [t^{N}]\bigl(F\compose(X+H)\bigr)
 =\sum_{n=n_{0}}^{N}\;\sum_{r=0}^{\Floor{(N-n)/2}}
   \frac{1}{(N-n-r)!}\,
   \ExponentialPartialBellPolynomial{\,N-n-r}{\,r}
    \bigl(1!\,h_{1},2!\,h_{2},\ldots\bigr)\;
   \shiftop{n}{r}\,p_{n},
\end{equation}
where \(\shiftop{n}{r}=\prod_{j=0}^{r-1}(\partial_{L}-n-j)\) is the operator of
\cref{plt:def:dif-shiftop}.  The sum is finite, with at most
\(\bigl\lfloor(N-n_{0}+2)^{2}/4\bigr\rfloor\) terms, and it involves only
\(p_{n_{0}},\ldots,p_{N}\) and \(h_{1},\ldots,h_{N-n_{0}}\).
\end{proposition}

\begin{proof}
Start from \eqref{plt:eq:tay-coefficient-formula}, in which the summand
indexed by \((n,k,m)\) with \(n+k+m=N\) is
\(\frac{h_{k,m}}{k!}\shiftop{n}{k}p_{n}\) and \(h_{k,m}=C_{m,k}\) by the index
dictionary of \cref{plt:sec:tay-convolution}.  Rename \(k\) as \(r\) and use
\(m=N-n-r\).  Since \(h_{0}=0\), \eqref{plt:eq:tay-C-bell} gives
\(C_{m,r}=\frac{r!}{m!}\ExponentialPartialBellPolynomial{m}{r}(1!h_{1},\ldots)\)
for \(m\ge r\ge1\), and \(C_{m,r}=0\) for \(m<r\); also
\(C_{0,0}=1=\frac{0!}{0!}\ExponentialPartialBellPolynomial{0}{0}\) and
\(C_{m,0}=0=\ExponentialPartialBellPolynomial{m}{0}\) for \(m>0\).  Hence in
every case \(C_{m,r}/r!=\ExponentialPartialBellPolynomial{m}{r}(\cdot)/m!\),
which is the displayed summand.  The constraint \(r\le m=N-n-r\), i.e.
\(2r\le N-n\), is \cref{plt:lem:bell-elementary}\,(i).  Counting the pairs
\((n,r)\) with \(n_{0}\le n\le N\) and \(0\le r\le\Floor{(N-n)/2}\) and
writing \(s=N-n\) gives \(\sum_{s=0}^{N-n_{0}}(\Floor{s/2}+1)
=\bigl\lfloor(N-n_{0}+2)^{2}/4\bigr\rfloor\).  Dependence on the listed blocks
only is visible from the displayed ranges: \(n\le N\), and the arguments of
\(\ExponentialPartialBellPolynomial{N-n-r}{r}\) reach at most
\(x_{(N-n-r)-r+1}\), so \(h_{j}\) occurs only for \(j\le N-n\le N-n_{0}\).
This agrees with, and refines for this special \(G\), the general finite
determination of \cref{plt:thm:cmp-finite-determination}.
\end{proof}

\begin{corollary}[Composition with a pure level-zero shift]
\label{plt:cor:cext-log-shift}
Let \(F=\sum_{n\ge n_{0}}p_{n}(L)t^{n}\in\PLlau\) and let
\(H=h_{0}(L)\in\K[L]\) have no positive \(t\)-levels, so that
\(X+H=X+h_{0}(L)\).  Then for every \(N\),
\begin{equation}\label{plt:eq:cext-log-shift}
 [t^{N}]\bigl(F\compose(X+h_{0})\bigr)
 =\sum_{n=n_{0}}^{N}\frac{h_{0}^{\,N-n}}{(N-n)!}\;\shiftop{n}{\,N-n}\,p_{n}.
\end{equation}
In particular, for \(F=X+A\) and \(h_{0}=aL+b\) with \(a,b\in\K\), the whole
composite is computed by a single sum with \(N-n_{0}+1\) terms and no
convolution at all.
\end{corollary}

\begin{proof}
Here \(H^{k}=h_{0}^{k}\) is concentrated at level \(0\), so
\(h_{k,m}=h_{0}^{k}\delta_{m,0}\) in
\eqref{plt:eq:tay-Hk-blocks}.  Substituting into
\eqref{plt:eq:tay-coefficient-formula} leaves only the triples
\((n,k,m)=(n,N-n,0)\), which is \eqref{plt:eq:cext-log-shift}.
\end{proof}

\begin{example}[Composition with \(X+L\)]
\label{plt:ex:cext-log-shift}
Take \(F=t=X^{-1}\), so \(n_{0}=1\), \(p_{1}=1\) and \(p_{n}=0\) otherwise.
Then \eqref{plt:eq:cext-log-shift} with \(h_{0}=L\) gives
\([t^{N}](t\compose(X+L))=\frac{L^{N-1}}{(N-1)!}\shiftop{1}{N-1}1\).  Since
\(\shiftop{1}{k}1=\prod_{j=0}^{k-1}(-1-j)=(-1)^{k}k!\), this is
\((-1)^{N-1}L^{N-1}\), so
\(t\compose(X+L)=\sum_{N\ge1}(-1)^{N-1}L^{N-1}t^{N}=t/(1+tL)\), in agreement
with the generator rule \eqref{plt:eq:cmp-generators-near-identity}.
\end{example}

\section{Reversion}
\label{plt:sec:cext-reversion}

Two closed forms are available for the inverse, and they answer different
questions.  The first is closed in the blocks of the perturbation.

\begin{corollary}[Reversion coefficients, closed in the input blocks]
\label{plt:cor:cext-reversion}
Let \(A=\sum_{m\ge0}a_{m}(L)t^{m}\in\PLpow\), \(F=X+A\), and
\(\rev{F}=X+\sum_{N\ge0}b_{N}(L)t^{N}\).  Then for every
\(N\in\NaturalNumbers\),
\begin{equation}\label{plt:eq:cext-reversion}
 b_{N}
 =\sum_{r=1}^{N+1}\frac{(-1)^{r}}{r!}
   \left[\prod_{s=N-r+1}^{N-1}(\partial_{L}-s)\right]
   C_{\,N-r+1,\;r}(a_{0},a_{1},\ldots),
\end{equation}
where \(C_{m,r}\) is the convolution polynomial of
\cref{plt:def:tay-convolution} formed from the blocks of \(A\), given in
closed form by \eqref{plt:eq:cext-Cmr-closed}, and where a product whose lower
bound exceeds its upper bound is empty.  Only
\(a_{0},\ldots,a_{N}\) occur, and the range \(r\le N+1\) cannot be shortened.
\end{corollary}

\begin{proof}
\Cref{plt:thm:lag-coefficients} is \eqref{plt:eq:cext-reversion} with
\(A_{r,N-r+1}\) in place of \(C_{N-r+1,r}\), and
\cref{plt:prop:tay-convolution}\,(i) identifies
\(A_{r,m}=[t^{m}]A^{r}=C_{m,r}\).  Since \(C_{m,r}\) involves only
\(a_{0},\ldots,a_{m}\) and \(m=N-r+1\le N\) here, only
\(a_{0},\ldots,a_{N}\) occur.  Sharpness of the range is
\cref{plt:prop:lag-range}.
\end{proof}

The special case in which the perturbation is a single level-zero block
degenerates spectacularly: exactly one Lagrange summand survives per output
block, and the whole inverse is given by one operator applied to one power.

\begin{corollary}[Reversion of a pure level-zero perturbation]
\label{plt:cor:cext-level-zero-reversion}
Let \(P\in\K[L]\) and \(F=X+P(L)\in\Gnear\).  Then
\begin{equation}\label{plt:eq:cext-level-zero-reversion}
 \rev{F}
 =X+\sum_{N\ge0}
   \frac{(-1)^{N+1}}{(N+1)!}\,
   \Bigl[\prod_{j=0}^{N-1}\bigl(\partial_{L}-j\bigr)\Bigr]P^{\,N+1}\;t^{N}
 =X+\sum_{N\ge0}\frac{(-1)^{N+1}}{(N+1)!}\,
   \shiftop{0}{N}\bigl(P^{\,N+1}\bigr)\,t^{N},
\end{equation}
the product being empty at \(N=0\), where the block is \(-P\).  For
\(P(L)=aL\) with \(a\in\K\) this is
\(\rev{F}=X+\sum_{N\ge0}a^{\,N+1}\LamP{N}(L)t^{N}\), which is
\cref{plt:thm:om-alpha-family}; for \(a=1\) it is the expansion of
\(\WrightOmega\) (\cref{plt:thm:om-expansion}).  For a constant \(P=c\in\K\)
it collapses to \(\rev{(X+c)}=X-c\).
\end{corollary}

\begin{proof}
Here \(a_{0}=P\) and \(a_{m}=0\) for \(m\ge1\), so by
\eqref{plt:eq:tay-C-def} the convolution polynomials are \(C_{0,r}=P^{r}\) and
\(C_{m,r}=0\) for \(m\ge1\).  In \eqref{plt:eq:cext-reversion} the index
\(m=N-r+1\) is therefore forced to \(0\), that is \(r=N+1\), and the operator
becomes \(\prod_{s=0}^{N-1}(\partial_{L}-s)=\shiftop{0}{N}\).  Substituting
gives \eqref{plt:eq:cext-level-zero-reversion}.

For \(P=aL\), \(\shiftop{0}{N}\bigl((aL)^{N+1}\bigr)
=a^{\,N+1}\shiftop{0}{N}L^{\,N+1}\), and comparison with
\eqref{plt:eq:lambert-operator} identifies the block as
\(a^{\,N+1}\LamP{N}(L)\).  For \(P=c\) constant and \(N\ge1\), the factor
\(\partial_{L}\) at \(s=0\) annihilates the constant \(c^{\,N+1}\), so every
block beyond \(N=0\) vanishes.
\end{proof}

% ed.: S3's formula sheet displays this identity as
% [t^n](rev F - X + P) = ((-1)^{n+1}/(n+1)!) D(D-1)...(D-n+1) P^{n+1}.
% The added "+P" is a slip: it is correct for n >= 1, where P (a level-zero
% block) contributes nothing, but false at n = 0, where the left side is
% -P + P = 0 while the right side is -P.  Removing the "+P", as in
% plt:eq:cext-level-zero-reversion, makes the formula correct for every
% n >= 0 and makes the Lambert case a specialization rather than an exception.
\begin{remark}[A repair to the level-zero reversion formula]
\label{plt:rmk:cext-level-zero-repair}
Source~S3 records \eqref{plt:eq:cext-level-zero-reversion} in the equivalent
but shifted form ``\([t^{n}](\rev{F}-X+P)=\cdots\)''.  That form is correct
for \(n\ge1\), because \(P\) sits at level zero and contributes nothing to
higher blocks, but it fails at \(n=0\), where the left-hand side is
\(-P+P=0\) while the right-hand side is \(-P\).  The correction is simply to
delete the \(+P\); \eqref{plt:eq:cext-level-zero-reversion} is then valid for
every \(N\ge0\), and the Lambert normalization \(\LamP{0}(u)=-u\) of
\cref{plt:def:lambert-polynomials} becomes the \(N=0\) instance rather than a
separately stipulated boundary value.
\end{remark}

The second closed form is closed in \(A\) and its \(\DX\)-derivatives.  It is
what one wants when \(A\) is not given blockwise --- for instance when it is a
symbolic expression, or when one wants a formula valid for a whole family of
perturbations at once.

\begin{theorem}[The Lagrange terms as differential polynomials]
\label{plt:thm:cext-universal}
Let \(A\in\PLpow\) and \(F=X+A\).  Then for every \(r\ge1\)
\begin{equation}\label{plt:eq:cext-lagrange-bell}
 \frac{(-1)^{r}}{r!}\,\DX^{\,r-1}\bigl(A^{r}\bigr)
 =(-1)^{r}\sum_{j=1}^{r}\frac{A^{j}}{j!}\;
   \ExponentialPartialBellPolynomial{\,r-1}{\,r-j}
    \bigl(\DX A,\DX^{2}A,\ldots,\DX^{\,j}A\bigr),
\end{equation}
and consequently
\begin{equation}\label{plt:eq:cext-universal-bell}
 \rev{F}
 =X+\sum_{r\ge1}(-1)^{r}\sum_{j=1}^{r}\frac{A^{j}}{j!}\;
   \ExponentialPartialBellPolynomial{\,r-1}{\,r-j}
    \bigl(\DX A,\ldots,\DX^{\,j}A\bigr),
\end{equation}
the outer family being summable.  Writing \(A'=\DX A\), \(A''=\DX^{2}A\) and
so on, the first five Lagrange terms are
\begin{equation}\label{plt:eq:cext-universal-terms}
\begin{split}
 \rev{F}-X
 ={}&-A\\
   &+AA'\\
   &-A(A')^{2}-\tfrac12A^{2}A''\\
   &+A(A')^{3}+\tfrac32A^{2}A'A''+\tfrac16A^{3}A'''\\
   &-A(A')^{4}-3A^{2}(A')^{2}A''-\tfrac12A^{3}(A'')^{2}
     -\tfrac23A^{3}A'A'''-\tfrac1{24}A^{4}A''''\\
   &+\FormalBigOAt{t^{5}}{t},
\end{split}
\end{equation}
each line being one value of \(r\), from \(r=1\) to \(r=5\).
\end{theorem}

\begin{proof}
Apply \cref{plt:cor:bell-power-derivative} with \(n=r-1\).  For \(r=1\) both
sides of \eqref{plt:eq:cext-lagrange-bell} are \(-A\) (the left side is
\(-\DX^{0}A\); the right side has only \(j=1\), with
\(\ExponentialPartialBellPolynomial{0}{0}=1\)).  For \(r\ge2\),
\eqref{plt:eq:bell-power-derivative} gives
\[
 \DX^{\,r-1}\bigl(A^{r}\bigr)
 =\sum_{k=1}^{r-1}\FallingFactorial{r}{k}\,A^{\,r-k}\,
   \ExponentialPartialBellPolynomial{\,r-1}{\,k}
   \bigl(\DX A,\ldots,\DX^{\,r-k}A\bigr),
\]
the index \(k\) stopping at \(r-1\) because
\(\ExponentialPartialBellPolynomial{r-1}{k}=0\) for \(k>r-1\).  Divide by
\(r!\) and use \(\FallingFactorial{r}{k}/r!=1/(r-k)!\); putting \(j=r-k\),
whose range is \(1\le j\le r-1\), turns the sum into
\(\sum_{j}\frac{A^{j}}{j!}
\ExponentialPartialBellPolynomial{r-1}{r-j}(\DX A,\ldots,\DX^{\,j}A)\).  The
missing term \(j=r\) carries
\(\ExponentialPartialBellPolynomial{r-1}{0}\), which vanishes for
\(r\ge2\), so it may be added without changing the value; this gives
\eqref{plt:eq:cext-lagrange-bell} in the stated uniform range
\(1\le j\le r\).  The argument list is
\(x_{1},\ldots,x_{(r-1)-(r-j)+1}=x_{1},\ldots,x_{j}\), that is
\(\DX A,\ldots,\DX^{\,j}A\).

Formula \eqref{plt:eq:cext-universal-bell} is now
\eqref{plt:eq:lag-burmann-pl} of \cref{plt:thm:lag-burmann} term by term, and
the summability of the outer family is the same theorem.

For \eqref{plt:eq:cext-universal-terms}, evaluate
\eqref{plt:eq:cext-lagrange-bell} for \(r\le5\) using
\cref{plt:tab:bell-partial}.  At \(r=2\): \(j=1\) gives
\(A\,\ExponentialPartialBellPolynomial{1}{1}(A')=AA'\) and \(j=2\) gives
\(\tfrac12A^{2}\ExponentialPartialBellPolynomial{1}{0}=0\); the sign is
\((-1)^{2}\).  At \(r=3\): \(j=1\) gives
\(A\,\ExponentialPartialBellPolynomial{2}{2}(A')=A(A')^{2}\), \(j=2\) gives
\(\tfrac12A^{2}\ExponentialPartialBellPolynomial{2}{1}(A',A'')
=\tfrac12A^{2}A''\), and \(j=3\) gives \(0\); the sign is \(-1\).  At
\(r=4\): \(\ExponentialPartialBellPolynomial{3}{3}=(A')^{3}\),
\(\ExponentialPartialBellPolynomial{3}{2}=3A'A''\),
\(\ExponentialPartialBellPolynomial{3}{1}=A'''\), giving
\(A(A')^{3}+\tfrac12A^{2}\cdot3A'A''+\tfrac16A^{3}A'''\).  At \(r=5\):
\(\ExponentialPartialBellPolynomial{4}{4}=(A')^{4}\),
\(\ExponentialPartialBellPolynomial{4}{3}=6(A')^{2}A''\),
\(\ExponentialPartialBellPolynomial{4}{2}=4A'A'''+3(A'')^{2}\),
\(\ExponentialPartialBellPolynomial{4}{1}=A''''\), giving
\(A(A')^{4}+\tfrac12A^{2}\cdot6(A')^{2}A''
+\tfrac16A^{3}\bigl(4A'A'''+3(A'')^{2}\bigr)+\tfrac1{24}A^{4}A''''\), which is
the last line with the sign \((-1)^{5}\).  Finally,
\cref{plt:lem:lag-order} gives \(\ordt\) at least \(r-1\ge5\) for every
omitted term, which is the recorded formal tail.
\end{proof}

\begin{remark}[Which reversion formula to run]
\label{plt:rmk:cext-which-reversion}
\Cref{plt:cor:cext-reversion} and \cref{plt:thm:cext-universal} are the same
series regrouped.  \Cref{plt:cor:cext-reversion} groups by output block and is
the formula to use when a single \(b_{N}\) is wanted, since it needs only
\(A^{r}\) for \(r\le N+1\) and one application of a differential operator per
term.  \Cref{plt:thm:cext-universal} groups by Lagrange index and is the
formula to use for structural statements, since the \(r\)-th line of
\eqref{plt:eq:cext-universal-terms} is a universal differential polynomial
with rational coefficients, independent of \(A\), of weight \(r\) in the
grading \(\operatorname{wt}(\DX^{i}A)=i+1\).  Neither is competitive as an
\emph{algorithm}: for computing all blocks up to a cutoff, triangular reversal
(\cref{plt:thm:rev-triangular}) or Newton reversion
(\cref{plt:thm:rev-newton-doubling}) is cheaper, as
\cref{plt:prop:alg-reversion-comparison} records.
\end{remark}

\section{What differs from the classical Lagrange inversion formula}
\label{plt:sec:cext-classical}

The classical formula, in the form standard in enumerative combinatorics
\cite{gessel-lagrange,flajolet-sedgewick}, reads as follows.  Let \(C\) be a
commutative \(\RationalNumbers\)-algebra, let \(\psi\in C[[z]]\), and let
\(u\in\varepsilon C[[\varepsilon]]\) be the unique solution of
\(u=\varepsilon\psi(u)\).  Then for every \(\eta\in C[[z]]\) and every
\(m\ge1\),
\begin{equation}\label{plt:eq:cext-classical-lb}
 [\varepsilon^{m}]\,\eta(u)
 =\frac1m\,[z^{m-1}]\bigl(\eta'(z)\psi(z)^{m}\bigr).
\end{equation}
\Cref{plt:thm:lag-burmann} is not this statement, and the differences are
worth naming precisely, because each of them is a place where a transplanted
classical proof breaks.

\begin{enumerate}
\item \emph{The output of one term.}  In
  \eqref{plt:eq:cext-classical-lb} the right-hand side is a single scalar,
  obtained by a coefficient extraction that evaluates a derivative at the
  origin.  In \eqref{plt:eq:lag-burmann-pl} the \(r\)-th summand
  \(\frac{(-1)^{r}}{r!}\DX^{\,r-1}(A^{r})\) is an entire element of
  \(\PLlau\), contributing to every block \(t^{N}\) with \(N\ge r-1\); the
  passage to one block is a second step, performed by
  \cref{plt:thm:lag-coefficients}.  Accordingly the volume's abstract
  ingredient \cref{plt:lem:lag-abstract} states
  \([\varepsilon^{m}]\gamma_{[u]}=\frac1{m!}\partial_{z}^{\,m-1}(\gamma'\psi^{m})\)
  as an identity of \emph{series} in \(z\), and specializes \(z\) only
  afterwards.  Evaluating at \(z=0\) recovers
  \eqref{plt:eq:cext-classical-lb}, because
  \(\frac1{m!}\bigl(\partial_{z}^{\,m-1}(\gamma'\psi^{m})\bigr)\big|_{z=0}
  =\frac1m\,[z^{m-1}](\gamma'\psi^{m})\).
\item \emph{The derivation does not annihilate the coefficients.}  In
  \eqref{plt:eq:cext-classical-lb} the operator \(\partial_{z}\) kills \(C\).
  Here \(\DX L=t\ne0\) (\cref{plt:thm:dif-derivation}), so differentiating a
  Lagrange term differentiates its coefficient polynomials; this is what
  produces the operator \(\shiftop{n}{k}\) of \cref{plt:def:dif-shiftop} in
  \eqref{plt:eq:cext-reversion} and, ultimately, the Lambert polynomials
  themselves.  Replacing \(\DX^{\,r-1}\) by ``extract the coefficient of
  \(t^{\,r-1}\)'' is therefore wrong by an operator, not by a constant.
\item \emph{What makes the family summable.}  Classically the grading is by
  the single variable \(\varepsilon\), and the solution is small because
  \(u\in\varepsilon C[[\varepsilon]]\).  Here \(A\) need not be small at all:
  \(\ordt(A)=0\) is the normal case, since a near-identity map may have a full
  level-zero block \(a_{0}(L)\).  Summability comes instead from the pairing
  of one derivative against one extra factor of \(A\), that is from the axiom
  \(\nu(\partial f)\ge\nu(f)+1\) of \cref{plt:def:lag-frame}; the resulting
  admissibility threshold is \(\nu(A)>-1\), not \(\nu(A)>0\)
  (\cref{plt:rmk:lag-why-minus-one}, \cref{plt:lem:lag-order}).
\item \emph{No formal residue is available.}  The usual proof of
  \eqref{plt:eq:cext-classical-lb} runs through the formal residue on
  \(C((z))\) and the identity \(\operatorname{Res}_{z}\partial_{z}(\cdot)=0\).
  \Cref{plt:prop:cext-no-residue} shows that no such functional exists on
  \(\PLlau\).  The volume's proof of \cref{plt:thm:lag-burmann} therefore
  goes through a weighted deformation algebra and an evaluation homomorphism
  (\cref{plt:lem:lag-eval}) instead of through residues.
\item \emph{The contributing range is finite and known.}  Classically every
  \(m\ge1\) contributes to its own coefficient.  Here only \(r\le N+1\)
  contributes to the block \(N\), and that bound is attained
  (\cref{plt:prop:lag-range}); this is what makes
  \eqref{plt:eq:cext-reversion} a finite formula rather than a limit.
\item \emph{No analytic content.}  \Cref{plt:thm:lag-burmann} is an identity
  of formal objects.  The transfer to a genuine inverse function of a real or
  complex variable is a separate theorem with separate hypotheses
  (\cref{plt:prop:rev-analytic-transfer}), exactly as the catalogue's
  distinction between \(\FormalBigOAt{t^{N}}{t}\) and
  \(\BigOAt{\cdot}{X\to+\infty}\) requires.
\end{enumerate}

\begin{proposition}[No formal residue on the polynomial--logarithmic Laurent
 ring]
\label{plt:prop:cext-no-residue}
There is no nonzero \(\K\)-linear functional
\(\rho\colon\PLlau\to\K\) with \(\rho\bigl(\DX F\bigr)=0\) for every
\(F\in\PLlau\).  The same holds for any \(\K\)-linear map from \(\PLlau\) to
any \(\K\)-vector space.
\end{proposition}

\begin{proof}
By \cref{plt:thm:dif-antiderivative} the map \(\DX\colon\PLlau\to\PLlau\) is
surjective.  If \(\rho\circ\DX=0\) then \(\rho\) vanishes on
\(\DX(\PLlau)=\PLlau\), so \(\rho=0\).  The same argument applies verbatim to
a linear map with any target.
\end{proof}

\begin{remark}[The residue that does exist is a residue in \(L\)]
\label{plt:rmk:cext-residue-in-L}
\Cref{plt:prop:cext-no-residue} does not contradict
\cref{plt:def:dif-residue}.  That definition attaches a residue to a
\emph{coefficient} \(h\in\K((L^{-1}))\), namely the coefficient of
\(L^{-1}\), and \cref{plt:lem:dif-res-derivative} shows that it annihilates
\(\partial_{L}\)-images.  It is a residue in the inner variable, defined only
after the coefficient ring has been enlarged to \(\K((L^{-1}))\) --- where
\(\partial_{L}\) is \emph{not} surjective (\cref{plt:lem:dif-image-laurent})
--- and it does not implement extraction of an outer block.  The obstruction
that a classical residue proof would need lives in the outer variable \(t\),
and there is none, because \(t\) is a unit of \(\PLlau\)
(\cref{plt:prop:trunc-not-ideal}\,(i)) and \(\DX\) is onto.
\end{remark}

\section{Which formula to use}
\label{plt:sec:cext-summary}

\begin{table}[htbp]\centering\small
\caption{Closed forms and their recursive counterparts.  ``Closed'' means a
 finite sum over an explicitly bounded index set in the input blocks alone;
 ``recursive'' means a triangular recurrence in the output blocks.}
\label{plt:tab:cext-summary}
\begin{tabular}{@{}p{0.20\textwidth}p{0.30\textwidth}p{0.42\textwidth}@{}}
\toprule
Task & Closed form & Recursive alternative, and when to prefer it\\
\midrule
Block of a product &
 \eqref{plt:eq:cext-product} (\cref{plt:prop:cext-product}) &
 none needed; the Cauchy product \eqref{plt:eq:mul-cauchy} is already finite.
 Use \cref{plt:prop:trunc-product} to truncate inputs safely.\\
\addlinespace
Block of an integer power &
 \eqref{plt:eq:mul-power-coefficient}, or
 \eqref{plt:eq:cext-Cmr-closed} in Bell form &
 the triangular recurrence \eqref{plt:eq:mul-Cmr-recurrence}; preferable when
 all powers \(r\le R\) are wanted, since it costs \(\BigO(N^{2}R)\) block
 products.\\
\addlinespace
Block of \((1+W)^{\sigma}\) &
 \eqref{plt:eq:cext-power} (\cref{plt:thm:cext-power}) &
 the power recurrence of \cref{plt:prop:alg-power-recurrence}; preferable for
 a full table.\\
\addlinespace
Block of a reciprocal &
 \eqref{plt:eq:cext-reciprocal}, or the determinant
 \eqref{plt:eq:cext-hessenberg} &
 \eqref{plt:eq:div-recip-cn} (\cref{plt:thm:div-reciprocal}) for a table;
 Newton doubling \cref{plt:thm:div-newton} when \(N\) is large.  Use the
 determinant for structural claims (\cref{plt:rmk:cext-hessenberg-use}).\\
\addlinespace
Block of a quotient &
 \(\recip{B}\) then \eqref{plt:eq:cext-product} &
 \eqref{plt:eq:div-quotient-recurrence} (\cref{plt:cor:div-quotient}); almost
 always preferable.\\
\addlinespace
Block of a composition &
 \eqref{plt:eq:tay-coefficient-formula}, or
 \eqref{plt:eq:cext-composition-bell} in Bell form &
 the block recurrence \cref{plt:prop:cmp-block-recurrence}; preferable when
 one inner map is composed with many outer series.\\
\addlinespace
Composition with a level-zero shift &
 \eqref{plt:eq:cext-log-shift} (\cref{plt:cor:cext-log-shift}) &
 none needed; the closed form is already a single sum.\\
\addlinespace
Block of a reversion &
 \eqref{plt:eq:cext-reversion} (\cref{plt:cor:cext-reversion}) &
 triangular reversal \cref{plt:thm:rev-triangular} or Newton
 \cref{plt:thm:rev-newton-doubling}; both cheaper for a full table
 (\cref{plt:prop:alg-reversion-comparison}).\\
\addlinespace
Reversion of a level-zero perturbation \(X+P(L)\) &
 \eqref{plt:eq:cext-level-zero-reversion}
 (\cref{plt:cor:cext-level-zero-reversion}) &
 none needed; one operator applied to one power gives every block.  Contains
 the Lambert family as the case \(P=aL\).\\
\addlinespace
Reversion as a differential polynomial &
 \eqref{plt:eq:cext-universal-bell},
 \eqref{plt:eq:cext-universal-terms} (\cref{plt:thm:cext-universal}) &
 no recursive counterpart; this is the form for symbolic \(A\) and for
 weight-graded structural arguments.\\
\bottomrule
\end{tabular}
\end{table}

\chapter{Operational formula sheet}
\label{plt:sec:formula-sheet}

This chapter is a navigation aid, not an exposition.  Every line states a
result of this volume and cites where that result is proved; nothing is
asserted here that is not proved there, and no line carries a hypothesis that
its citation does not carry.  Read a line together with its reference before
using it: the abbreviations below suppress hypotheses that are essential ---
which coefficient class an element lives in, whether a leading block is a
unit, whether an inner argument is admissible, whether a claim is formal or
analytic.

Standing conventions for the whole sheet.  \(\K\) is a field of characteristic
zero; \(X\) is the large variable and the default regime is \(X\to+\infty\)
along the positive real ray; \(t=X^{-1}\) and \(L=\log X\); \(F,G,H\) denote
series, \(p_{n},q_{n},a_{n},b_{n},h_{n}\) their blocks in \(\K[L]\), and
\(N,n,m,r,k\) integers.  Truncation is \emph{exclusive}
(\cref{plt:def:trunc-exclusive}).  Three symbols that look alike are kept
apart throughout: \(\FormalBigOAt{t^{N}}{t}\) is a formal valuation tail with
no analytic content, \(\BigOAt{A(X)}{X\to+\infty}\) is an analytic estimate on
the printed regime, and \(\BigO(N^{2})\) is an operation count.

\section{The four ambient classes}
\label{plt:sec:sheet-classes}

\begin{longtable}{@{}p{0.56\textwidth}p{0.38\textwidth}@{}}
\caption{Generators, classes, and their algebraic type.}
\label{plt:tab:sheet-classes}\\
\toprule
Statement & Proved at\\
\midrule
\endfirsthead
\toprule
Statement & Proved at\\
\midrule
\endhead
\(t\DefinitionEquals X^{-1}\), \(L\DefinitionEquals\log X\); independent
 generators in the formal algebra & \cref{plt:def:ring-generators}\\
\(\PLpow=\K[L][[t]]\), \(\PLlau=\K[L]((t))\), \(\PLrat=\K(L)((t))\),
 \(\PLit=\K((L^{-1}))((t))\) & \cref{plt:def:ring-four-classes}\\
\(\PLpow\subsetneq\PLlau\subsetneq\PLrat\subsetneq\PLit\); the first two are
 domains and not fields, the last two are fields &
 \cref{plt:prop:ring-chain}, \cref{plt:cor:ring-domains},
 \cref{plt:thm:div-Ett-field}\\
\(\PLit\) is the Hahn field of the lexicographic rank-two value group &
 \cref{plt:thm:ring-hahn}, \cref{plt:cor:ring-hahn-subrings}\\
The coefficient extensions \(\K[L]\rightsquigarrow\K(L)\rightsquigarrow
 \K((L^{-1}))\) are not localizations of \(\PLlau\) &
 \cref{plt:prop:ring-not-localization}\\
Fractional outer exponents require the Puiseux-log tower; a general real
 exponent requires a Hahn exponent group &
 \cref{plt:def:ring-puiseux}, \cref{plt:prop:ring-puiseux-tower}\\
\(\Gnear=\SetOf{X+A:A\in\PLpow}\) is a group under \(\compose\) &
 \cref{plt:def:grp-gnear}, \cref{plt:thm:grp-group}\\
Every class is \(t\)-adically complete and separated &
 \cref{plt:thm:trunc-completeness}\\
\bottomrule
\end{longtable}

\section{Valuation, leading data, and the dominance order}
\label{plt:sec:sheet-leading}

\begin{longtable}{@{}p{0.56\textwidth}p{0.38\textwidth}@{}}
\caption{Outer valuation, rank-two valuation, leading data, dominance,
 log-degree profile.}\label{plt:tab:sheet-leading}\\
\toprule
Statement & Proved at\\
\midrule
\endfirsthead
\toprule
Statement & Proved at\\
\midrule
\endhead
\(\ordt(F)=\min\SetOf{n:p_{n}\ne0}\), \(\ordt(0)=+\infty\) &
 \cref{plt:def:lead-ordt}\\
\(\ordt(FG)=\ordt F+\ordt G\);
 \(\ordt(F+G)\ge\min\SetOf{\ordt F,\ordt G}\), strict exactly when the
 leading blocks cancel & \cref{plt:thm:ring-valuation-laws}\\
Leading block \(p_{\ordt F}\), leading monomial
 \(\lm(F)=t^{\ordt F}L^{\degL p_{\ordt F}}\), leading scalar \(\lc(F)\) ---
 three different objects & \cref{plt:def:lead-leading-data},
 \cref{plt:rmk:lead-three-objects}\\
\(\lm(FG)=\lm(F)\lm(G)\) and \(\lc(FG)=\lc(F)\lc(G)\) &
 \cref{plt:thm:lead-rank-two-laws}, \cref{plt:thm:mul-leading-law}\\
The block representation of a nonzero element is unique &
 \cref{plt:thm:lead-uniqueness}\\
\(t^{n}L^{j}\succdom t^{m}L^{k}\) iff \(n<m\), or \(n=m\) and \(j>k\) &
 \cref{plt:def:ring-dominance}\\
Dominance is genuine growth as \(X\to+\infty\) &
 \cref{plt:lem:ring-dominance-analytic}\\
\(\dprofile{F}(n)=\degL p_{n}\), with \(-\infty\) at absent blocks &
 \cref{plt:def:lead-profile}\\
\(\dprofile{F+G}\le\max\), \(\dprofile{FG}\le\) the
 \((\max,+)\)-convolution, with an exact attainment criterion &
 \cref{plt:prop:lead-profile-laws}, \cref{plt:thm:mul-degree-bound}\\
\bottomrule
\end{longtable}

\section{Truncation, the formal tail, and summability}
\label{plt:sec:sheet-truncation}

\begin{longtable}{@{}p{0.56\textwidth}p{0.38\textwidth}@{}}
\caption{Exclusive truncation, the formal big-O token, summability.}
\label{plt:tab:sheet-truncation}\\
\toprule
Statement & Proved at\\
\midrule
\endfirsthead
\toprule
Statement & Proved at\\
\midrule
\endhead
\(\Trunc{F}{N}=\sum_{n<N}p_{n}(L)t^{n}\); every retained block is retained
 whole & \cref{plt:def:trunc-exclusive}\\
Inclusive-to-exclusive dictionary \([F]_{\le N}=\Trunc{F}{N+1}\) &
 \cref{plt:rmk:trunc-inclusive}, \eqref{plt:eq:trunc-inclusive-dictionary}\\
\(\Trunc{(FG)}{N}
 =\Trunc{\bigl(\Trunc{F}{N-\ordt G}\cdot\Trunc{G}{N-\ordt F}\bigr)}{N}\) &
 \cref{plt:prop:trunc-product}\\
\(F=G+\FormalBigOAt{t^{N}}{t}\) means \(\ordt(F-G)\ge N\), equivalently
 \(\Trunc{F}{N}=\Trunc{G}{N}\) & \cref{plt:def:trunc-formal-o},
 \cref{plt:prop:trunc-formal-o}\\
\(t^{N}\PLpow\) is a filtration step of \(\PLlau\), never an ideal, because
 \(t\) is a unit there & \cref{plt:prop:trunc-not-ideal}\\
\([t^{n}]F=p_{n}(L)\in\K[L]\) and \([t^{n}L^{j}]F=a_{n,j}\in\K\); extraction
 is linear and filtration-continuous & \cref{plt:def:trunc-extraction},
 \cref{plt:prop:trunc-extraction}\\
A family is summable iff \(\ordt\to+\infty\); sums may then be regrouped and
 multiplied termwise & \cref{plt:def:trunc-summable},
 \cref{plt:thm:trunc-summability},
 \cref{plt:thm:trunc-summability-criterion},
 \cref{plt:prop:trunc-fubini}\\
A bound on \(\degL\) is a second, independent approximation, not part of the
 \(t\)-adic topology & \cref{plt:lem:cext-finiteness}\\
\bottomrule
\end{longtable}

\section{The derivation, its powers, and antiderivatives}
\label{plt:sec:sheet-differential}

\begin{longtable}{@{}p{0.56\textwidth}p{0.38\textwidth}@{}}
\caption{The distinguished derivation and its calculus.}
\label{plt:tab:sheet-differential}\\
\toprule
Statement & Proved at\\
\midrule
\endfirsthead
\toprule
Statement & Proved at\\
\midrule
\endhead
\(\DX=t\,\partial_{L}-t^{2}\partial_{t}\); the unique \(\K\)-derivation with
 \(\DX t=-t^{2}\), \(\DX L=t\) and \(\DX(\PLpow)\subseteq\PLpow\) &
 \cref{plt:thm:dif-derivation}\\
Block formula
 \(\DX\bigl(t^{n}p(L)\bigr)=t^{n+1}\bigl(\partial_{L}-n\bigr)p(L)\) &
 \cref{plt:prop:dif-block}\\
Euler form \(\EulerOp=t^{-1}\DX\) preserves each block and acts as
 \(\partial_{L}-n\) on \(t^{n}\K[L]\) & \cref{plt:prop:dif-euler}\\
\(\shiftop{n}{k}=\prod_{j=0}^{k-1}\bigl(\partial_{L}-n-j\bigr)\),
 \(\shiftop{n}{0}=1\); the second index is never optional &
 \cref{plt:def:dif-shiftop}\\
\(\DX^{k}\bigl(t^{n}p(L)\bigr)=t^{n+k}\,\shiftop{n}{k}p(L)\) &
 \cref{plt:thm:dif-repeated}\\
\(\shiftop{0}{k}=\sum_{j}\stirlingsigned{k}{j}\partial_{L}^{\,j}\), and
 \(\DX^{k}t^{n}=(-1)^{k}\RisingFactorial{n}{k}\,t^{n+k}\) &
 \cref{plt:cor:dif-stirling}\\
\(\ordt(\DX F)\ge\ordt F+1\), with the sharp value and the three cases &
 \cref{plt:thm:dif-valuation}\\
\(\ker\DX\cap\PLlau=\K\) & \cref{plt:cor:dif-kernel}\\
\(\DX\) is onto \(\PLlau\); antiderivatives exist and are unique up to a
 constant of \(\K\) & \cref{plt:thm:dif-antiderivative}\\
Resonance: \((\partial_{L}-n)p=q\) is triangular for \(n\ne0\); the block
 \(n=0\) integrates and adds a constant &
 \cref{plt:lem:dif-triangular}, \cref{plt:lem:dif-resonant}\\
In \(\PLit\) and \(\PLrat\) antidifferentiation is obstructed by a residue at
 \(L=\infty\); this is where \(\log L\) is forced &
 \cref{plt:thm:dif-anti-it}, \cref{plt:thm:dif-anti-rat},
 \cref{plt:cor:dif-loglog}\\
\bottomrule
\end{longtable}

\section{Product, reciprocal, and division}
\label{plt:sec:sheet-arithmetic}

\begin{longtable}{@{}p{0.56\textwidth}p{0.38\textwidth}@{}}
\caption{Multiplicative calculus.}\label{plt:tab:sheet-arithmetic}\\
\toprule
Statement & Proved at\\
\midrule
\endfirsthead
\toprule
Statement & Proved at\\
\midrule
\endhead
\([t^{N}](FG)=\sum_{n+m=N}p_{n}q_{m}\), a finite sum for each \(N\); the four
 classes are closed under it & \cref{plt:def:mul-cauchy},
 \cref{plt:lem:mul-finite-fiber}, \cref{plt:thm:mul-closure}\\
\([t^{N}L^{d}](FG)=\sum_{n+m=N}\sum_{e+f=d}a_{n,e}b_{m,f}\) &
 \cref{plt:prop:cext-product}\\
\(F\) is a unit of \(\PLlau\) iff its leading block is a nonzero
 \emph{scalar}, iff \(F=c\,t^{n_{0}}(1+U)\) with \(c\in\K^{\times}\),
 \(U\in t\PLpow\) & \cref{plt:thm:div-unit-criterion},
 \cref{plt:cor:div-unit-group}\\
\(\recip{B}=\sum_{n}C_{n}t^{n}\) with \(C_{0}=c^{-1}\) and
 \(C_{n}=-c^{-1}\sum_{j=1}^{n}B_{j}C_{n-j}\), where \(B_{0}=c\in\K^{\times}\)
 & \cref{plt:thm:div-reciprocal}\\
% ed.: this row printed the division recurrence with no hypothesis on B_0,
% which repeats the defect of S1's sheet ("B_0 \ne 0").  In K[L]((t)) a
% nonzero polynomial block is not enough: B_0 must be a unit of K[L], i.e. a
% nonzero scalar, or the quotient blocks leave K[L].  Hypothesis printed.
\(Q_{n}=B_{0}^{-1}\bigl(A_{n}-\sum_{j=1}^{n}B_{j}Q_{n-j}\bigr)\), valid over
 \(\K[L]\) only when \(B_{0}\in\K^{\times}\) --- a nonzero \emph{scalar}, not
 merely a nonzero polynomial &
 \cref{plt:cor:div-quotient}, \cref{plt:thm:div-unit-criterion}\\
Newton reciprocal step \(C\mapsto C(2-BC)\) squares the residual
 \(1-BC\) and may be truncated at the doubled order &
 \cref{plt:thm:div-newton}, \cref{plt:cor:div-newton-exact}\\
Truncated division: relative precision is exactly preserved &
 \cref{plt:thm:div-truncated}, \cref{plt:prop:div-relative-precision}\\
\((1+U)^{\sigma}=\sum_{r\ge0}\binom{\sigma}{r}U^{r}\) for \(\ordt U\ge1\),
 with finite determination at each block & \cref{plt:thm:mul-binomial}\\
Closed forms:
 \(\displaystyle[t^{r}](1+W)^{\sigma}
 =\frac1{r!}\sum_{k=1}^{r}\FallingFactorial{\sigma}{k}\,
 \ExponentialPartialBellPolynomial{r}{k}(1!w_{1},2!w_{2},\ldots)\);
 \(\sigma=-1\) gives the reciprocal &
 \cref{plt:thm:cext-power}, \cref{plt:cor:cext-reciprocal}\\
\(c^{\,r+1}[t^{r}]\recip{B}\) is a Toeplitz--Hessenberg determinant in
 \(B_{0},\ldots,B_{r}\) & \cref{plt:prop:cext-hessenberg}\\
% ed.: this row read "\PLrat is the smallest field of the tower with
% polynomial-log coefficients", which inverts the content of
% plt:thm:div-minimal-field: \PLrat has RATIONAL-log coefficients, and the
% theorem says E((t)) is a field iff E = K(L).  Restated as proved.
\(\PLlau\) is a domain but not a field; among the rings \(E((t))\) with
 \(\K[L]\subseteq E\subseteq\K(L)\), exactly one is a field, namely
 \(\PLrat\) &
 \cref{plt:cor:mul-domain}, \cref{plt:thm:div-minimal-field},
 \cref{plt:cor:div-domain-not-field}\\
\bottomrule
\end{longtable}

\section{Composition and the exact Taylor formula}
\label{plt:sec:sheet-composition}

\begin{longtable}{@{}p{0.56\textwidth}p{0.38\textwidth}@{}}
\caption{Substitution, the Taylor formula, and the coefficient formula.}
\label{plt:tab:sheet-composition}\\
\toprule
Statement & Proved at\\
\midrule
\endfirsthead
\toprule
Statement & Proved at\\
\midrule
\endhead
Admissible inner argument \(G=cX^{\rho}(1+U)\), \(c\in\K^{\times}\),
 \(\rho>0\), \(U\in t\PLpow\) & \cref{plt:def:cmp-admissible-inner}\\
Generator images \(\Phi_{G}(t)=c^{-1}t^{\rho}(1+U)^{-1}\),
 \(\Phi_{G}(L)=\rho L+\log c+\log(1+U)\); \(\Phi_{G}\) is a \(\K\)-algebra
 homomorphism compatible with summable families &
 \cref{plt:thm:cmp-substitution-hom}\\
Near identity: \(t\compose(X+H)=t/(1+tH)\),
 \(L\compose(X+H)=L+\log(1+tH)\).  Substituting \(t\mapsto t+H\), or updating
 \(L\) alone, is wrong & \cref{plt:prop:cmp-generator-substitution}\\
\(\displaystyle F\compose(X+H)=\sum_{k\ge0}\frac{H^{k}}{k!}\,\DX^{k}F\),
 an exact identity, not an approximation & \cref{plt:thm:tay-formula}\\
\(\displaystyle[t^{r}]\bigl(F\compose(X+H)\bigr)
 =\sum_{n+k+m=r}\frac{h_{k,m}}{k!}\,\shiftop{n}{k}p_{n}\), with
 \(h_{k,m}=[t^{m}]H^{k}=C_{m,k}\) & \cref{plt:thm:tay-coefficients},
 \cref{plt:def:tay-convolution}\\
\(C_{0,0}=1\), \(C_{m,0}=0\) for \(m>0\),
 \(C_{m,r+1}=\sum_{j=0}^{m}h_{j}C_{m-j,r}\); closed form in
 \eqref{plt:eq:cext-Cmr-closed} & \cref{plt:prop:tay-convolution},
 \cref{plt:prop:cext-convolution-closed}\\
Bell form when \(h_{0}=0\), and the single-sum formula for a level-zero shift
 \(H=h_{0}(L)\) & \cref{plt:prop:cext-composition-bell},
 \cref{plt:cor:cext-log-shift}\\
Substituting into a scalar power series:
 \(\varphi(W)\) for \(W\in t\PLpow\), with Bell coefficients &
 \cref{plt:prop:dif-bell}, \cref{plt:prop:tay-scalar-outer}\\
\(\displaystyle\DX^{n}\bigl(\varphi(F)\bigr)
 =\sum_{k=1}^{n}\varphi^{(k)}(F)\,
 \ExponentialPartialBellPolynomial{n}{k}\bigl(\DX F,\ldots,
 \DX^{\,n-k+1}F\bigr)\) & \cref{plt:thm:bell-faa}\\
\(\displaystyle\DX^{n}(F\compose G)
 =\sum_{k=1}^{n}\bigl(\DX^{k}F\bigr)\compose G\;
 \ExponentialPartialBellPolynomial{n}{k}\bigl(\gamma,\ldots,
 \DX^{\,n-k}\gamma\bigr)\), \(\gamma=\DX G\) &
 \cref{plt:prop:tay-faa-di-bruno}, \cref{plt:lem:bell-abstract-faa}\\
Chain rule \(\DX(F\compose G)=\bigl((\DX F)\compose G\bigr)\DX G\);
 associativity holds inside the admissible class &
 \cref{plt:thm:cmp-chain-rule}, \cref{plt:thm:cmp-associativity}\\
Finite determination and the triangular precision budget &
 \cref{plt:thm:cmp-finite-determination},
 \cref{plt:cor:cmp-precision-budget}\\
Analytic transfer under composition is a separate theorem with its own
 hypotheses & \cref{plt:thm:cmp-analytic-transfer}\\
\bottomrule
\end{longtable}

\section{Reversion}
\label{plt:sec:sheet-reversion}

\begin{longtable}{@{}p{0.56\textwidth}p{0.38\textwidth}@{}}
\caption{The compositional inverse: fixed point, algorithms, closed
 formulas.}\label{plt:tab:sheet-reversion}\\
\toprule
Statement & Proved at\\
\midrule
\endfirsthead
\toprule
Statement & Proved at\\
\midrule
\endhead
For \(F=X+A\in\Gnear\) and \(G=X+B\): \(F\compose G=X\) iff
 \(B=-A\compose(X+B)\) & \cref{plt:prop:rev-fixed-point-equation}\\
The fixed point exists, is unique, and is a two-sided inverse; the Picard
 iterates gain one block per step & \cref{plt:thm:rev-existence},
 \cref{plt:lem:rev-gain}, \cref{plt:thm:grp-group}\\
\(\RevErr{F}{H}=F\compose H-X\); its order equals
 \(\ordt(H-\rev{F})\) exactly, on both sides, with the same leading block &
 \cref{plt:def:rev-residual}, \cref{plt:prop:rev-residual-order},
 \cref{plt:cor:rev-finite-certificate}\\
Triangular reversal: \(r_{N}=[t^{N}]\RevErr{F}{G_{N-1}}\),
 \(G_{N}=G_{N-1}-r_{N}t^{N}\), and \(G_{N}=\Trunc{\rev{F}}{N+1}\) &
 \cref{plt:thm:rev-triangular}, \cref{plt:prop:rev-next-block}\\
At leading slope \(a\in\K^{\times}\) the correction is \(-r_{N}/a\), and the
 whole problem conjugates to the near-identity one &
 \cref{plt:prop:rev-general-slope}\\
Truncating the inputs of a reversal step is safe &
 \cref{plt:prop:rev-truncation-safe}, \cref{plt:cor:alg-truncate-early}\\
Newton map \(\displaystyle\mathcal N_{F}(G)=G-\frac{\RevErr{F}{G}}
 {(\DX F)\compose G}\); the denominator is a unit &
 \cref{plt:def:rev-newton}, \cref{plt:lem:rev-newton-unit}\\
Exact arithmetic: \(\ordt\RevErr{F}{\mathcal N_{F}(G)}\ge2m+2+\alpha\) with
 \(m=\ordt\RevErr{F}{G}\), \(\alpha=\ordt A\) &
 \cref{plt:thm:rev-newton-doubling}\\
\textbf{Guard requirement.}  With an inexact correction accurate only to
 order \(\sigma\), the new residual order is
 \(\min\SetOf{2m+2+\alpha,\ \sigma\,}\); doubling is lost unless
 \(\sigma\ge2m+2+\alpha\) & \cref{plt:thm:rev-newton-guard},
 \cref{plt:lem:rev-precision-propagation},
 \cref{plt:cor:rev-newton-schedule}\\
\(\displaystyle\rev{(X+A)}
 =X+\sum_{r\ge1}\frac{(-1)^{r}}{r!}\,\DX^{\,r-1}\bigl(A^{r}\bigr)\),
 the \(r\)-th summand of order \(\ge r-1\) & \cref{plt:thm:lag-burmann}\\
\(\displaystyle b_{N}=\sum_{r=1}^{N+1}\frac{(-1)^{r}}{r!}
 \Bigl[\prod_{s=N-r+1}^{N-1}(\partial_{L}-s)\Bigr]A_{r,N-r+1}\), the empty
 product occurring exactly at \(r=1\) & \cref{plt:thm:lag-coefficients},
 \cref{plt:cor:cext-reversion}\\
The range \(r\le N+1\) is exactly right and is attained &
 \cref{plt:prop:lag-range}\\
Level-zero perturbation:
 \(\displaystyle\rev{(X+P)}=X+\sum_{N\ge0}\frac{(-1)^{N+1}}{(N+1)!}
 \shiftop{0}{N}\bigl(P^{\,N+1}\bigr)t^{N}\) for \(P\in\K[L]\) &
 \cref{plt:cor:cext-level-zero-reversion},
 \cref{plt:rmk:cext-level-zero-repair}\\
Transport of an observable:
 \(\varphi\compose\rev{F}=\varphi+\sum_{r\ge1}\frac{(-1)^{r}}{r!}
 \partial^{\,r-1}\bigl((\partial\varphi)A^{r}\bigr)\) &
 \cref{plt:thm:lag-transport}, \cref{plt:cor:lag-observables}\\
Universal differential-polynomial form, weight-graded by
 \(\operatorname{wt}(\DX^{i}A)=i+1\) & \cref{plt:thm:cext-universal}\\
Beyond the identity scale: reduce by the leading monomial; a leading
 logarithm forces a nested logarithm &
 \cref{plt:thm:lag-reduction}, \cref{plt:thm:lag-rigidity},
 \cref{plt:prop:lag-loglog-forced}\\
Analytic residual transfer needs injectivity and a derivative lower bound &
 \cref{plt:prop:rev-analytic-transfer}\\
\bottomrule
\end{longtable}

\section{Wright omega, the Lambert polynomials, and Lambert
 \texorpdfstring{\(\LambertW\)}{W}}
\label{plt:sec:sheet-omega}

\begin{longtable}{@{}p{0.56\textwidth}p{0.38\textwidth}@{}}
\caption{The canonical logarithmic reversion and its coefficient family.}
\label{plt:tab:sheet-omega}\\
\toprule
Statement & Proved at\\
\midrule
\endfirsthead
\toprule
Statement & Proved at\\
\midrule
\endhead
\(y\mapsto y+\log y\) is a bijection \((0,\infty)\to\RealNumbers\); its
 inverse is \(\WrightOmega\), with
 \(\WrightOmega(x)=\LambertW_{0}(\EulerE^{x})\) &
 \cref{plt:prop:om-real-bijection}\\
\(\displaystyle\WrightOmega=X+\sum_{n\ge0}\LamP{n}(L)t^{n}
 =X-L+\sum_{n\ge1}\LamP{n}(L)t^{n}\) in \(\Gnear\); each
 Lagrange--B\"urmann summand is a single block &
 \cref{plt:thm:om-expansion}\\
\(\displaystyle\LamP{n}(u)=\frac{(-1)^{n+1}}{(n+1)!}
 \Bigl[\prod_{j=0}^{n-1}(\partial_{u}-j)\Bigr]u^{\,n+1}\);
 \(\LamP{0}(u)=-u\), \(\LamP{1}(u)=u\),
 \(\LamP{2}(u)=\frac{u^{2}}{2}-u\) & \cref{plt:def:lambert-polynomials}\\
\(\LamP{n}(0)=0\);
 \(\LamP{n+1}=-\LamP{n}+n\int_{0}^{u}\LamP{n}\); equivalently
 \({\LamP{n+1}}'=n\LamP{n}-{\LamP{n}}'\) &
 \cref{plt:thm:lambert-recurrence}\\
\(\displaystyle\LamP{n}(u)=\sum_{m=1}^{n}(-1)^{n-m}
 \stirlingone{n}{n-m+1}\frac{u^{m}}{m!}\) \ \((n\ge1)\) &
 \cref{plt:thm:lambert-stirling}\\
\(\deg\LamP{n}=n\) for \(n\ge1\) only; \([u^{n}]\LamP{n}=1/n\);
 \(\lambda_{n,1}=(-1)^{n-1}\); \(\LamPsc{n}=n!\LamP{n}\in
 \IntegerNumbers[u]\) & \cref{plt:cor:lambert-coefficients},
 \cref{plt:cor:lambert-integrality}\\
\(\stirlingone{n}{n-m+1}
 =\ExponentialPartialBellPolynomial{n}{\,n-m+1}(0!,1!,2!,\ldots)\): the
 Lambert triangle is a slice of the Bell array &
 \cref{plt:thm:bell-stirling}, \cref{plt:cor:bell-lambert}\\
\(\rev{(X+aL)}=X+\sum_{n\ge0}a^{\,n+1}\LamP{n}(L)t^{n}\); the case
 \(a=-1\) gives \((-1)^{n+1}\LamP{n}\), \emph{not} \((-1)^{n}\LamP{n}\) &
 \cref{plt:thm:om-alpha-family}, \cref{plt:rmk:om-second-family}\\
Exact residual of the truncation:
 \(\WrightOmega_{N}+\log\WrightOmega_{N}-X
 =-\LamP{N+1}(L)t^{\,N+1}+\FormalBigOAt{t^{\,N+2}}{t}\) &
 \cref{plt:thm:om-residual}, \cref{plt:cor:om-certificate}\\
Analytic form, \(x\to+\infty\) on the positive ray, \(K=(N+2)^{2}\):
 \(\WrightOmega(x)-\WrightOmega_{N}(x)
 =\LamP{N+1}(L)x^{-(N+1)}
 +\BigOAt{L^{K}x^{-(N+2)}}{x\to+\infty}\) &
 \cref{plt:thm:om-analytic}\\
Branch coordinates: \(\LambertW_{k}(z)+\log\LambertW_{k}(z)=\xi_{k}\) with
 \(\xi_{k}=\PrincipalLogarithm z+2\pi\ImaginaryUnit k\); the formal expansion
 of \(\LambertW_{k}\) is \cref{plt:thm:om-expansion} under
 \(X\mapsto\xi_{k}\), \(L\mapsto\ell_{k}=\log_{\mathcal B}\xi_{k}\).  The
 polynomials are branch-independent; the coordinates, sector and analytic
 estimates are not, and the complex analysis is not developed &
 \cref{plt:rmk:om-branches}\\
Numerical tables of \(\LamP{n}\), of \(\stirlingone{n}{n-m+1}\), and of
 \(\LamP{n}(1)\) & \cref{plt:sec:lambert-tables},
 \cref{plt:tab:lambert-polynomials}, \cref{plt:tab:lambert-triangle},
 \cref{plt:tab:lambert-at-one}\\
\bottomrule
\end{longtable}

\section{Bell polynomials and closed coefficient extraction}
\label{plt:sec:sheet-bell}

\begin{longtable}{@{}p{0.56\textwidth}p{0.38\textwidth}@{}}
\caption{The combinatorial layer of this appendix.}
\label{plt:tab:sheet-bell}\\
\toprule
Statement & Proved at\\
\midrule
\endfirsthead
\toprule
Statement & Proved at\\
\midrule
\endhead
\(\ExponentialPartialBellPolynomial{n}{k}\) by multiplicity sum;
 \(\ExponentialPartialBellPolynomial{0}{0}=1\),
 \(\ExponentialPartialBellPolynomial{n}{0}=0\) for \(n\ge1\) &
 \cref{plt:def:bell-partial}\\
\(\displaystyle\frac{\Xi^{k}}{k!}=\sum_{n\ge k}
 \ExponentialPartialBellPolynomial{n}{k}\frac{z^{n}}{n!}\) for
 \(\Xi=\sum_{j\ge1}\xi_{j}z^{j}/j!\); and the complete version with
 \(\exp\Xi\) & \cref{plt:prop:bell-gf}\\
Degree \(k\), weight \(n\), bihomogeneity, boundary rows &
 \cref{plt:lem:bell-elementary}\\
\(\OrdinaryPartialBellPolynomial{n}{k}(\xi)=\frac{k!}{n!}
 \ExponentialPartialBellPolynomial{n}{k}(1!\xi_{1},2!\xi_{2},\ldots)\); this
 is the factor that converts \(C_{n,k}\) to the exponential normalization &
 \cref{plt:lem:bell-normalizations}\\
Tabulating recurrences
 \(\ExponentialPartialBellPolynomial{n}{k}
 =\sum_{j}\binom{n-1}{j-1}x_{j}
 \ExponentialPartialBellPolynomial{n-j}{k-1}\) and its \(\binom{n}{j}\)
 variant & \cref{plt:prop:bell-recurrences}\\
Differentiating recurrence
 \(\ExponentialPartialBellPolynomial{n+1}{k}
 =\mathcal D\ExponentialPartialBellPolynomial{n}{k}
 +x_{1}\ExponentialPartialBellPolynomial{n}{k-1}\) &
 \cref{plt:lem:tay-bell-recurrence}\\
\(\ExponentialPartialBellPolynomial{n}{k}(0!,1!,\ldots)=\stirlingone{n}{k}\);
 \(\ExponentialPartialBellPolynomial{n}{k}(1,1,\ldots)
 =\StirlingSecondKind{n}{k}\) & \cref{plt:thm:bell-stirling}\\
Abstract Fa\`a di Bruno for a Leibniz tower \(\partial u_{k}=u_{k+1}\gamma\) &
 \cref{plt:lem:bell-abstract-faa}\\
\(\displaystyle\DX^{n}(A^{r})=\sum_{k}\FallingFactorial{r}{k}A^{\,r-k}
 \ExponentialPartialBellPolynomial{n}{k}(\DX A,\ldots)\) &
 \cref{plt:cor:bell-power-derivative}\\
No nonzero linear functional on \(\PLlau\) annihilates all \(\DX\)-images:
 there is no formal residue in \(t\) & \cref{plt:prop:cext-no-residue}\\
Six named differences from the classical Lagrange inversion formula &
 \cref{plt:sec:cext-classical}\\
Which closed form to use for which task & \cref{plt:tab:cext-summary}\\
\bottomrule
\end{longtable}

\section{What this sheet does not assert}
\label{plt:sec:sheet-caveats}

\begin{warningbox}[title={Four things the sheet is silent about}]
\begin{enumerate}
\item \textbf{Analytic validity.}  Every entry above whose citation is a
  formal statement is an identity of formal series.  It becomes an assertion
  about a function only through an explicit transfer theorem
  (\cref{plt:thm:mul-analytic-product}, \cref{plt:thm:cmp-analytic-transfer},
  \cref{plt:prop:rev-analytic-transfer}, \cref{plt:thm:om-analytic}), each
  with its own hypotheses and its own printed regime.  A bare
  \(\FormalBigO\) never implies a bare \(\BigOAt{\cdot}{\cdot}\)
  (\cref{plt:def:trunc-formal-o}).
\item \textbf{Uniformity in \(L\).}  A block-level identity says nothing about
  the size of the coefficient polynomial at a given \(x\).  The analytic
  entries above therefore carry an explicit power of \(L\) in the error term
  --- for instance \(K=(N+2)^{2}\) in \cref{plt:thm:om-analytic} --- and none
  of them is uniform in \(N\).
\item \textbf{Degree truncation.}  Discarding high powers of \(L\) inside a
  retained block is not part of the \(t\)-adic calculus and none of the
  precision statements above covers it
  (\cref{plt:lem:cext-finiteness}, \cref{plt:def:trunc-exclusive}).
\item \textbf{Closure.}  Membership in a class is a hypothesis, never a
  conclusion of a formula.  Exponentials, general powers, logarithms of
  non-units and reversions at a non-identity leading slope leave \(\PLlau\)
  under conditions recorded in \cref{plt:thm:dif-exp-boundary},
  \cref{plt:thm:dif-log-boundary}, \cref{plt:thm:dif-generated} and
  \cref{plt:thm:lag-rigidity}, and only there.
\end{enumerate}
\end{warningbox}

% ---- fragment 17-tables ----
\chapter{Tables of the Lambert polynomials}
\label{plt:sec:lambert-tables}

All entries below were produced from the operator definition of
\cref{plt:def:lambert-polynomials} in exact rational arithmetic and then
verified three independent ways: against the Stirling closed form of
\cref{plt:thm:lambert-stirling}, against the differential recurrence of
\cref{plt:thm:lambert-recurrence}, and --- the only check that does not
presuppose any of the derived identities --- by substituting the resulting
series into the defining equation
\(\WrightOmega+\log\WrightOmega=X\) and confirming that every coefficient block
through outer order \(t^{8}\) vanishes identically.

\section{The scaled integer family}

By \cref{plt:cor:lambert-integrality}, \(n!\LamP{n}(u)\in\IntegerNumbers[u]\)
for \(n\ge1\); the table lists that integer polynomial together with the
factorial that divides it, which is the storage normalization
\(\LamPsc{n}=n!\LamP{n}\).  Read the row as
\(\LamP{n}(u)=\bigl(\text{entry}\bigr)/n!\).  The sources tabulate these
polynomials only through \(n=8\); rows \(9\)--\(12\) are new here.

\begin{longtable}{@{}r r p{0.70\textwidth}@{}}
\caption{The Lambert polynomials $\LamP{n}(u)=\LamPsc{n}(u)/n!$.}\label{plt:tab:lambert-polynomials}\\
\toprule
$n$ & $n!$ & \multicolumn{1}{c}{$\LamPsc{n}(u)=n!\,\LamP{n}(u)$}\\
\midrule
\endfirsthead
\toprule
$n$ & $n!$ & \multicolumn{1}{c}{$\LamPsc{n}(u)=n!\,\LamP{n}(u)$}\\
\midrule
\endhead
$0$ & $1$ & $-u$\\[0.25em]
$1$ & $1$ & $u$\\[0.25em]
$2$ & $2$ & $u^{2}-2u$\\[0.25em]
$3$ & $6$ & $2u^{3}-9u^{2}+6u$\\[0.25em]
$4$ & $24$ & $6u^{4}-44u^{3}+72u^{2}-24u$\\[0.25em]
$5$ & $120$ & $24u^{5}-250u^{4}+700u^{3}-600u^{2}+120u$\\[0.25em]
$6$ & $720$ & $120u^{6}-1{,}644u^{5}+6{,}750u^{4}-10{,}200u^{3}+5{,}400u^{2}-720u$\\[0.25em]
$7$ & $5{,}040$ & $720u^{7}-12{,}348u^{6}+68{,}208u^{5}-154{,}350u^{4}+147{,}000u^{3}-52{,}920u^{2}+5{,}040u$\\[0.25em]
$8$ & $40{,}320$ & $5{,}040u^{8}-104{,}544u^{7}+735{,}392u^{6}-2{,}274{,}384u^{5}+3{,}292{,}800u^{4}-2{,}163{,}840u^{3}+564{,}480u^{2}-40{,}320u$\\[0.25em]
$9$ & $362{,}880$ & $40{,}320u^{9}-986{,}256u^{8}+8{,}504{,}928u^{7}-33{,}911{,}136u^{6}+67{,}885{,}776u^{5}-68{,}584{,}320u^{4}+33{,}022{,}080u^{3}-6{,}531{,}840u^{2}+362{,}880u$\\[0.25em]
$10$ & $3{,}628{,}800$ & $362{,}880u^{10}-10{,}265{,}760u^{9}+105{,}543{,}000u^{8}-521{,}049{,}600u^{7}+1{,}357{,}398{,}000u^{6}-1{,}913{,}375{,}520u^{5}+1{,}428{,}840{,}000u^{4}-526{,}176{,}000u^{3}+81{,}648{,}000u^{2}-3{,}628{,}800u$\\[0.25em]
$11$ & $39{,}916{,}800$ & $3{,}628{,}800u^{11}-116{,}915{,}040u^{10}+1{,}402{,}893{,}360u^{9}-8{,}325{,}405{,}000u^{8}+27{,}062{,}085{,}600u^{7}-50{,}009{,}929{,}200u^{6}+52{,}481{,}610{,}720u^{5}-30{,}187{,}080{,}000u^{4}+8{,}781{,}696{,}000u^{3}-1{,}097{,}712{,}000u^{2}+39{,}916{,}800u$\\[0.25em]
$12$ & $479{,}001{,}600$ & $39{,}916{,}800u^{12}-1{,}446{,}526{,}080u^{11}+19{,}921{,}172{,}832u^{10}-138{,}940{,}660{,}320u^{9}+546{,}429{,}272{,}400u^{8}-1{,}267{,}789{,}406{,}400u^{7}+1{,}754{,}714{,}586{,}240u^{6}-1{,}426{,}718{,}240{,}640u^{5}+652{,}040{,}928{,}000u^{4}-153{,}679{,}680{,}000u^{3}+15{,}807{,}052{,}800u^{2}-479{,}001{,}600u$\\[0.25em]
\bottomrule
\end{longtable}

\begin{remark}[The $n=0$ row is not an exception]
\label{plt:rmk:lambert-zero-row}
\(\LamP{0}(u)=-u\) has degree~\(1\), not~\(0\).  The degree law
\(\deg\LamP{n}=n\) of \cref{plt:cor:lambert-coefficients} is therefore stated
only for \(n\ge1\), and the integrality normalization
\(\LamPsc{n}=n!\LamP{n}\) is likewise a statement about \(n\ge1\); the table's
\(n=0\) row records the value, not an instance of either law.
\end{remark}

\section{The coefficient triangle}

Write \(\LamP{n}(u)=\sum_{m=1}^{n}\lambda_{n,m}u^{m}\) for \(n\ge1\).  By
\cref{plt:thm:lambert-stirling},
\[
 \lambda_{n,m}=\frac{(-1)^{n-m}}{m!}\stirlingone{n}{n-m+1},
\]
so the signs alternate along each row and the magnitudes are unsigned Stirling
cycle numbers divided by a factorial.  The next table lists the integers
\(m!\,(-1)^{n-m}\lambda_{n,m}=\stirlingone{n}{n-m+1}\), which is the row of the
Stirling triangle read from the right.

\begin{table}[htbp]\centering\small
\caption{$\stirlingone{n}{n-m+1}=m!\,\AbsoluteValue{\lambda_{n,m}}$, the unsigned coefficient data of $\LamP{n}$.}\label{plt:tab:lambert-triangle}
\begin{tabular}{@{}rrrrrrrrrr@{}}
\toprule
$n\backslash m$ & $1$ & $2$ & $3$ & $4$ & $5$ & $6$ & $7$ & $8$ & $9$\\
\midrule
$1$ & $1$ &  &  &  &  &  &  &  & \\
$2$ & $1$ & $1$ &  &  &  &  &  &  & \\
$3$ & $1$ & $3$ & $2$ &  &  &  &  &  & \\
$4$ & $1$ & $6$ & $11$ & $6$ &  &  &  &  & \\
$5$ & $1$ & $10$ & $35$ & $50$ & $24$ &  &  &  & \\
$6$ & $1$ & $15$ & $85$ & $225$ & $274$ & $120$ &  &  & \\
$7$ & $1$ & $21$ & $175$ & $735$ & $1{,}624$ & $1{,}764$ & $720$ &  & \\
$8$ & $1$ & $28$ & $322$ & $1{,}960$ & $6{,}769$ & $13{,}132$ & $13{,}068$ & $5{,}040$ & \\
$9$ & $1$ & $36$ & $546$ & $4{,}536$ & $22{,}449$ & $67{,}284$ & $118{,}124$ & $109{,}584$ & $40{,}320$\\
\bottomrule
\end{tabular}
\end{table}

The first column is \(\stirlingone{n}{n}=1\), which is
\cref{plt:cor:lambert-coefficients}(iv) in the form
\(\lambda_{n,1}=(-1)^{n-1}\).  The last entry of row~\(n\) is
\(\stirlingone{n}{1}=(n-1)!\), giving the leading coefficient
\(\lambda_{n,n}=1/n\).  The penultimate entry is
\(\stirlingone{n}{2}=(n-1)!H_{n-1}\), giving
\(\lambda_{n,n-1}=-H_{n-1}\).

\section{Numerical check values}

The following exact values are convenient unit tests for an
implementation of the recurrence; each is \(\LamP{n}(1)\), which by
\cref{plt:thm:lambert-stirling} equals
\(\sum_{m=1}^{n}(-1)^{n-m}\stirlingone{n}{n-m+1}/m!\).

\begin{table}[htbp]\centering\small
\caption{Exact values $\LamP{n}(1)$.}\label{plt:tab:lambert-at-one}
\begin{tabular}{@{}r l r l@{}}
\toprule
$n$ & $\LamP{n}(1)$ & $n$ & $\LamP{n}(1)$\\
\midrule
$0$ & $-1$ & $7$ & $\frac{15}{56}$\\
$1$ & $1$ & $8$ & $\frac{457}{1260}$\\
$2$ & $-\frac{1}{2}$ & $9$ & $-\frac{49}{90}$\\
$3$ & $-\frac{1}{6}$ & $10$ & $-\frac{391}{2016}$\\
$4$ & $\frac{5}{12}$ & $11$ & $\frac{32213}{36960}$\\
$5$ & $-\frac{1}{20}$ & $12$ & $-\frac{299363}{1425600}$\\
$6$ & $-\frac{49}{120}$ &  & \\
\bottomrule
\end{tabular}
\end{table}


\section{An end-to-end numerical validation}
\label{plt:sec:numeric-validation}

The tables above and the identities of \cref{plt:sec:lambert-polynomials} are
algebraic.  A coefficient table can be internally consistent and still be
wrong in sign, because a sign error propagates through every identity that
derives one normalization from another.  The following check is independent of
all of them: it solves the defining equation numerically and compares.

Let \(\WrightOmega_{N}(x)=x+\sum_{n=0}^{N}\LamP{n}(L)x^{-n}\) with
\(L=\log x\), the truncation of \cref{plt:thm:om-expansion}.  Solving
\(y+\log y=x\) to fifty significant digits by Newton's method and comparing
gives the following.  The residual column is
\(\WrightOmega_{N}+\log\WrightOmega_{N}-x\) computed numerically; the predicted
column is \(-\LamP{N+1}(L)x^{-(N+1)}\), which is what
\cref{plt:eq:om-residual} asserts the residual must be to leading order.

\begin{table}[htbp]\centering\small
\caption{Observed versus predicted residual of the truncated Wright-omega
expansion.  The ratio tends to \(1\) as \(x\) grows, which is the content of
\cref{plt:eq:om-residual}; the departure at \(x=50\) is the omitted
\(\FormalBigOAt{t^{N+2}}{t}\) block making itself felt.}
\label{plt:tab:omega-residual}
\begin{tabular}{@{}r r l l r@{}}
\toprule
\(x\) & \(N\) & observed residual & predicted & ratio\\
\midrule
$50$   & $1$ & $-1.534239\times10^{-3}$  & $-1.495976\times10^{-3}$  & $1.0256$\\
$50$   & $2$ & $-5.859441\times10^{-6}$  & $-7.300570\times10^{-6}$  & $0.8026$\\
$50$   & $3$ & $+1.599261\times10^{-6}$  & $+1.473299\times10^{-6}$  & $1.0855$\\
$50$   & $4$ & $+9.405048\times10^{-8}$  & $+8.979402\times10^{-8}$  & $1.0474$\\
\addlinespace
$200$  & $1$ & $-2.211720\times10^{-4}$  & $-2.184442\times10^{-4}$  & $1.0125$\\
$200$  & $2$ & $-1.606059\times10^{-6}$  & $-1.596060\times10^{-6}$  & $1.0063$\\
$200$  & $3$ & $-1.802138\times10^{-9}$  & $-2.030332\times10^{-9}$  & $0.8876$\\
$200$  & $4$ & $+2.386204\times10^{-10}$ & $+2.316867\times10^{-10}$ & $1.0299$\\
\addlinespace
$1000$ & $1$ & $-1.701321\times10^{-5}$  & $-1.695079\times10^{-5}$  & $1.0037$\\
$1000$ & $2$ & $-4.535166\times10^{-8}$  & $-4.520477\times10^{-8}$  & $1.0032$\\
$1000$ & $3$ & $-1.013683\times10^{-10}$ & $-1.011740\times10^{-10}$ & $1.0019$\\
$1000$ & $4$ & $-9.250010\times10^{-14}$ & $-9.319150\times10^{-14}$ & $0.9926$\\
\bottomrule
\end{tabular}
\end{table}

Two features of the table are worth reading carefully, because both are
predictions of the theory rather than accidents.

First, the residual \emph{changes sign} between \(N=2\) and \(N=3\) at
\(x=50\) and between \(N=3\) and \(N=4\) at \(x=200\).  That is not
instability: by \cref{plt:cor:lambert-coefficients} the coefficients of
\(\LamP{n}\) alternate, so the leading residual
\(-\LamP{N+1}(L)x^{-(N+1)}\) changes sign with \(N\) whenever \(L\) is small
enough that the top coefficient does not yet dominate.  A sign error anywhere
in the chain would break this pattern rather than reproduce it.

Second, at \(x=1000\) the error falls by roughly three orders of magnitude per
added block:

\begin{center}\small
\begin{tabular}{@{}r l r l@{}}
\toprule
\(N\) & \(\AbsoluteValue{\WrightOmega_{N}-\WrightOmega}\) &
\(N\) & \(\AbsoluteValue{\WrightOmega_{N}-\WrightOmega}\)\\
\midrule
$0$ & $6.924751\times10^{-3}$  & $4$ & $9.240705\times10^{-14}$\\
$1$ & $1.699609\times10^{-5}$  & $5$ & $7.844485\times10^{-16}$\\
$2$ & $4.530604\times10^{-8}$  & $6$ & $6.759229\times10^{-18}$\\
$3$ & $1.012664\times10^{-10}$ & & \\
\bottomrule
\end{tabular}
\end{center}

\begin{warningbox}[title={What this table does \emph{not} show}]
\label{plt:rmk:numeric-not-convergence}
Improvement through \(N=6\) at one value of \(x\) is not convergence.  The
series of \cref{plt:thm:om-expansion} diverges for every fixed \(x\): by
\cref{plt:cor:lambert-coefficients} the leading coefficient of \(\LamP{n}\) is
\(1/n\) while its degree is \(n\), so for fixed \(x\) the terms eventually
grow.  What the table exhibits is asymptotic behaviour --- for each fixed
\(N\), the error is \(\BigOAt{L^{N+2}/x^{N+2}}{x\to\infty}\) --- and the
correct reading of the second table is a statement about \(N\) small relative
to \(\log x\), not a limit in \(N\).  Section~\ref{plt:sec:poincare} makes the
distinction precise.
\end{warningbox}

% ---- fragment 18-exercises ----
\chapter{Exercises with solutions}
\label{plt:sec:exercises}

The six source reports carry \(83\) exercises between them, distributed very
unevenly: report~1 sets \(58\), three at the end of each chapter, and supplies
no solutions; report~3 sets \(13\) and solves all of them; report~2 sets \(12\)
and solves all of them; reports~4, 5 and~6 set none but contain worked examples
that make good problems.  The overlap is heavy --- five of the six ask for the
reversion of \(X+aL\) in one guise or another --- and the quality is uneven,
ranging from genuine tests of the formalism down to arithmetic drill.

What follows is a single graded set of \(56\) exercises, harvested from all six
sources, deduplicated, corrected, and organised by the part of the volume they
exercise.  Every exercise carries a complete solution, not a hint.  Where a
source supplied a solution it was recomputed independently; the corrections are
marked in the \LaTeX{} source with \texttt{\%~ed.:} comments and are also
summarised in \cref{plt:sec:repairs}.  Where an important technique of the body
had no exercise anywhere in the corpus --- the sharp summability threshold, the
guard-precision rule, the derivative certificate, the analytic transfer --- one
has been added.

\section{How to use this appendix}
\label{plt:sec:exercises-howto}

Three conventions are worth restating, because most of the mistakes recorded
below come from breaking one of them.

\begin{enumerate}
\item \emph{Truncation is exclusive.}  \(\Trunc{F}{N}\) keeps the blocks of
outer order \(n<N\) (\cref{plt:def:trunc-exclusive}).  Sources~1 and~6 use the
inclusive cutoff \([F]_{\le N}=\Trunc{F}{N+1}\); every index taken from them
has been shifted by one block.  Exercises whose title contains the word
\emph{trap} exist precisely because that shift is invisible until it is wrong.
\item \emph{The three \(\BigOAt{\cdot}{\cdot}\)-roles are distinct.}
\(\FormalBigOAt{t^{N}}{t}\) is a statement about a formal valuation and carries
no analytic content whatsoever;
\(\BigOAt{A(X)}{X\to+\infty}\) is an inequality between functions on a stated
ray; a complexity bound is neither.  No exercise below asks you to pass from
the first to the second without an analytic hypothesis, and three of them ask
you to find where somebody did.
\item \emph{Every statement names its coefficient field and its regime.}  The
default field is a field \(\K\) of characteristic zero; the default regime is
\(X\to+\infty\) along the positive real ray with the real logarithm.
\end{enumerate}

Exercises marked \emph{(hard)} need either a genuinely new idea or a long exact
computation.  Two special kinds recur.

\begin{itemize}
\item \emph{Counterexample exercises} ask you to refute a plausible
strengthening of a theorem of the body, usually one obtained by deleting a
hypothesis.  These are
\cref{plt:exer:exr-block-local-not-summable,plt:exer:exr-tail-not-ideal,%
plt:exer:exr-no-prefactor,plt:exer:exr-unit-nonzero-block,%
plt:exer:exr-antiderivative-obstruction,plt:exer:exr-commutator,%
plt:exer:exr-constant-inner,plt:exer:exr-derivative-certificate,%
plt:exer:exr-summability-threshold,plt:exer:exr-degree-law-zero,%
plt:exer:exr-poincare}.  The hypotheses the sources kept dropping are exactly
the ones these exercises defend.
\item \emph{Traps} hand you a plausible-looking computation containing one
documented failure mode --- a truncation convention read in the wrong
direction, a formal tail treated as an analytic estimate, a division by a
nonunit, a scale change applied to one generator only --- and ask you to locate
the error and repair it.  These are
\cref{plt:exer:exr-trap-symmetric-truncation,plt:exer:exr-trap-nonunit,%
plt:exer:exr-trap-inner-precision,plt:exer:exr-trap-universal-formula,%
plt:exer:exr-trap-scale,plt:exer:exr-trap-omega-index,%
plt:exer:exr-trap-formal-to-analytic,plt:exer:exr-trap-product-guard}.
\end{itemize}

\section{The scale and its algebra}
\label{plt:sec:exercises-scale}

\begin{exercise}[Valuation of a sum, and how far it can jump]
\label{plt:exer:exr-valuation-jump}
Work in \(\PLlau\), using only \cref{plt:thm:ring-valuation-laws}.
\begin{enumerate}
\item Prove \(\ordt(F+G)\ge\min\SetOf{\ordt F,\ordt G}\), with equality
whenever \(\ordt F\ne\ordt G\).
\item For every \(k\ge1\) and every \(d\ge0\) exhibit \(F,G\in\PLpow\) with
\(\ordt F=\ordt G=0\), \(\degL\bigl([t^{0}]F\bigr)=d\), and
\(\ordt(F+G)=k\).
\item Show that strict inequality in (1) forces \(\ordt F=\ordt G=:n\) and
\([t^{n}]F=-[t^{n}]G\).
\end{enumerate}
\end{exercise}

\begin{proof}[Solution]
(1)  Let \(n=\min\SetOf{\ordt F,\ordt G}\).  For \(m<n\) both \([t^{m}]F\) and
\([t^{m}]G\) vanish, hence so does \([t^{m}](F+G)\); this is
\(\ordt(F+G)\ge n\).  If \(\ordt F<\ordt G\) then \(n=\ordt F\) and
\([t^{n}](F+G)=[t^{n}]F\ne0\), so \(\ordt(F+G)=n\); the case
\(\ordt G<\ordt F\) is symmetric.

(2)  Put
\[
 F=L^{d}+t^{k},\qquad G=-L^{d}.
\]
Both lie in \(\PLpow\), both have outer order \(0\), the \(t^{0}\)-block of
\(F\) has \(L\)-degree \(d\), and \(F+G=t^{k}\) has outer order exactly \(k\).

(3)  If \(\ordt F<\ordt G\) or \(\ordt G<\ordt F\) then (1) gives equality, so
strictness forces \(\ordt F=\ordt G=n\).  Then
\(0=[t^{n}](F+G)=[t^{n}]F+[t^{n}]G\), which is the assertion.  Note that the
converse fails: cancelling leading blocks guarantee only
\(\ordt(F+G)\ge n+1\) and say nothing about how much further the order jumps,
which is what (2) exploits.
\end{proof}

% ed.: source 1 asks only for a jump "by at least two", i.e. the case k = 2 of
% ed.: part (2), and does not ask for the converse.  The stronger form is
% ed.: proved above because part (3) is what the sum rule of
% ed.: \cref{plt:thm:ring-valuation-laws} actually asserts.

\begin{exercise}[Completeness of the power-series ring, block by block]
\label{plt:exer:exr-completeness}
Let \((F_r)_{r\ge0}\) be a \(t\)-adically Cauchy sequence in \(\PLpow\).
\begin{enumerate}
\item Show that for each \(n\) the block \([t^{n}]F_r\) is eventually
independent of \(r\); call its value \(p_n\in\K[L]\).
\item Show \(F\DefinitionEquals\sum_{n\ge0}p_nt^{n}\) lies in \(\PLpow\) and
\(F_r\to F\).
\item Locate the exact step at which \(\PLpow\), rather than \(\PLlau\), was
used, and show that a Cauchy sequence in \(\PLlau\) nevertheless always
satisfies \(\inf_r\ordt(F_r)>-\infty\), so that the same three steps do prove
completeness of \(\PLlau\) as well.
\end{enumerate}
\end{exercise}

\begin{proof}[Solution]
(1)  Cauchy means: for every \(N\) there is \(r_N\) with
\(\ordt(F_r-F_s)\ge N\) whenever \(r,s\ge r_N\).  Applying this with \(N=n+1\)
gives \([t^{n}]F_r=[t^{n}]F_s\) for all \(r,s\ge r_{n+1}\), which is the
assertion; write \(p_n\) for the common value.

(2)  Each \(p_n\) lies in \(\K[L]\) and the index set is
\(\NaturalNumbers\), so the support of \(F\) is bounded below by \(0\) and
\(F\in\PLpow\).  Fix \(N\) and let \(r\ge\max\SetOf{r_1,\dots,r_{N}}\).  For
\(n<N\) we have \([t^{n}]F_r=p_n=[t^{n}]F\), so \(\ordt(F-F_r)\ge N\).  As
\(N\) was arbitrary, \(F_r\to F\).  Uniqueness of the limit is separatedness,
\cref{plt:thm:trunc-completeness}\,(i).

(3)  The only step that used \(\PLpow\) is the claim in (2) that
\(\sum_np_nt^{n}\) is an element of the ring: on \(\PLlau\) the blockwise
limits \(p_n\) could be nonzero for infinitely many \(n<0\), and then the sum
is not an element of \(\PLlau\).  This is not idle: the sequence of
\cref{plt:ex:trunc-not-cauchy} stabilises coefficientwise to exactly such a
symbol.  But that sequence is not Cauchy, and no Cauchy sequence can behave
that way.  Indeed apply the Cauchy condition with \(N=0\): there is \(r_0\)
with \(\ordt(F_r-F_{r_0})\ge0\) for all \(r\ge r_0\), whence by part~(1) of
\cref{plt:exer:exr-valuation-jump}
\[
 \ordt(F_r)\ \ge\ \min\SetOf{\ordt(F_{r_0}),\,0}
 \qquad(r\ge r_0),
\]
and the finitely many indices \(r<r_0\) contribute a finite minimum.  Hence
\(\mu\DefinitionEquals\inf_r\ordt(F_r)>-\infty\), every \(F_r\) lies in
\(t^{\mu}\PLpow\), and steps (1) and (2) may be run verbatim inside that
subring.  This is \cref{plt:thm:trunc-completeness}\,(iii),(iv).
\end{proof}

% ed.: source 2 proves the common lower bound by asserting that "the sequence
% ed.: is already Laurent-bounded at some initial precision".  That is not a
% ed.: hypothesis of completeness; the bound has to be extracted from the
% ed.: Cauchy condition at cutoff N = 0, as in part (3).

\begin{exercise}[Summable powers]
\label{plt:exer:exr-summable-powers}
Let \(U\in\PLlau\).  Prove that the family \((U^{n})_{n\ge0}\) is formally
summable in the sense of \cref{plt:def:trunc-summable} if and only if
\(U=0\) or \(\ordt(U)\ge1\), and that in that case
\(\sum_{n\ge0}U^{n}=\recip{(1-U)}\).
\end{exercise}

\begin{proof}[Solution]
Assume \(U\ne0\) and put \(q=\ordt(U)\).  Since \(\PLlau\) is a domain and
\(\ordt\) is additive on products (\cref{plt:thm:ring-valuation-laws}),
\(\ordt(U^{n})=nq\) for every \(n\ge0\), with \(U^{0}=1\) of order \(0\).

If \(q\ge1\) then \(nq\to+\infty\), so for each \(N\) the set
\(\SetOf{n:\ordt(U^{n})<N}\) is contained in \(\SetOf{n:n<N}\) and is finite:
the family is summable.  If \(q\le0\) then \(nq\le0\) for all \(n\), so
\(\SetOf{n:\ordt(U^{n})<1}=\NaturalNumbers\) is infinite and the family is not
summable.

For the value, let \(S=\sum_{n\ge0}U^{n}\), which exists by the summability
just proved.  Multiplication is continuous for the \(t\)-adic filtration on a
summable family (\cref{plt:thm:trunc-summability}), so
\[
 (1-U)S=\sum_{n\ge0}U^{n}-\sum_{n\ge0}U^{n+1}
 =U^{0}=1,
\]
the middle step being a telescoping of two summable families.  Since
\(\ordt(U)\ge1\) the element \(1-U\) is a unit of \(\PLpow\) by
\cref{plt:thm:div-unit-criterion}, and its reciprocal is unique, so
\(S=\recip{(1-U)}\).  This is \cref{plt:cor:ring-geometric}.
\end{proof}

\begin{exercise}[Counterexample: block-local finiteness is not summability]
\label{plt:exer:exr-block-local-not-summable}
\Cref{plt:thm:trunc-summability-criterion}\,(c) characterises summability by
two conditions: every block receives finitely many contributions, \emph{and}
\(\inf_i\ordt(S_i)>-\infty\).  Show that the second clause cannot be deleted.
\begin{enumerate}
\item Exhibit a family in \(\PLlau\) that is block-locally finite but not
summable.
\item Exhibit one that is block-locally finite, not summable, and whose
blockwise sum \emph{is} a legitimate element of \(\PLlau\).  Conclude that
``the blockwise sum exists in the ring'' is not a substitute for summability
either.
\item Show that on \(\PLpow\) the second clause is automatic.
\end{enumerate}
\end{exercise}

\begin{proof}[Solution]
(1)  Take \(S_i=t^{-i}\) for \(i\ge1\).  The block \(t^{n}\) receives a
contribution from \(S_i\) only for \(i=-n\), so at most one contribution.  But
\(\ordt(S_i)=-i\), so \(\SetOf{i:\ordt(S_i)<0}\) is infinite and the family
fails \cref{plt:def:trunc-summable}.  Concretely, the partial sums
\(T_m=\sum_{i=1}^{m}t^{-i}\) satisfy \(\ordt(T_{m+1}-T_m)=-(m+1)\), so they are
not Cauchy, and the blockwise limits would give a symbol with infinitely many
negative blocks, which is not an element of \(\PLlau\)
(\cref{plt:thm:ring-support}\,(b)).

(2)  Take
\[
 S_i=t^{-i}-t^{-i-1}\qquad(i\ge1).
\]
The block \(t^{n}\) receives contributions only from \(i=-n\) and \(i=-n-1\),
so at most two; the family is block-locally finite.  Blockwise,
\([t^{-1}]\sum_iS_i=1\) and for \(n\le-2\) the two contributions are \(+1\)
(from \(i=-n\)) and \(-1\) (from \(i=-n-1\)), summing to \(0\).  The blockwise
sum is therefore \(t^{-1}\in\PLlau\).  Nevertheless
\(\ordt(S_i)=-i-1\to-\infty\), so the family is not summable, and indeed the
partial sums
\(\sum_{i=1}^{m}S_i=t^{-1}-t^{-m-1}\)
% ed.: the sentence printed here was garbled ("ordt(...)-differences are ...").
% ed.: What is meant is the outer order of the difference of two consecutive
% ed.: partial sums, which is what Cauchyness is about.
do not converge: consecutive partial sums differ by
\(S_{m+1}=t^{-m-1}-t^{-m-2}\), of outer order \(-(m+2)\to-\infty\), so the
sequence is not Cauchy.  So a blockwise sum can exist in the ring and still fail
to be the limit of the partial sums in any sense the filtration recognises.
Summability is a statement about the family, not about the array of
coefficients it produces.

(3)  If every \(S_i\in\PLpow\) then \(\ordt(S_i)\ge0\), so
\(\inf_i\ordt(S_i)\ge0>-\infty\).
\end{proof}

\begin{exercise}[Counterexample: the tail is a filtration, not an ideal]
\label{plt:exer:exr-tail-not-ideal}
A plausible reading of ``compute modulo \(t^{N}\)'' is that
\(\mathfrak F_N=t^{N}\PLpow\) is an ideal of \(\PLlau\) and that
\(\Trunc{\cdot}{N}\) is the quotient homomorphism.  Refute both halves.
\begin{enumerate}
\item Show \(\mathfrak F_N\) is not an ideal of \(\PLlau\) for any \(N\), and
that the ideal it generates is all of \(\PLlau\).
\item Show \(\Trunc{\cdot}{N}\) is not multiplicative on \(\PLlau\), by
exhibiting \(F,G\) with
\(\Trunc{(FG)}{N}\ne\Trunc{\bigl(\Trunc{F}{N}\Trunc{G}{N}\bigr)}{N}\).
\item State what does survive, and prove it.
\end{enumerate}
\end{exercise}

\begin{proof}[Solution]
(1)  \(t\) is a unit of \(\PLlau\), with inverse \(t^{-1}\).  Hence
\(t^{-1}\cdot t^{N}=t^{N-1}\), which has outer order \(N-1<N\) and so does not
lie in \(\mathfrak F_N\): \(\mathfrak F_N\) is not closed under multiplication
by ring elements.  Moreover \(t^{-N}\cdot t^{N}=1\), so the ideal generated by
\(t^{N}\) contains \(1\) and is all of \(\PLlau\).  In particular
\(\PLlau/(t^{N})=0\) and there is no quotient ring to map onto.  This is
\cref{plt:prop:trunc-not-ideal}\,(i),(ii).

(2)  Take
\[
 F=G=\frac{t^{-2}}{1-t}=\sum_{n\ge-2}t^{n},
 \qquad N=1 .
\]
Then \(FG=t^{-4}(1-t)^{-2}=\sum_{j\ge-4}(j+5)t^{j}\), so
\[
 \Trunc{(FG)}{1}=t^{-4}+2t^{-3}+3t^{-2}+4t^{-1}+5 .
\]
On the other hand \(\Trunc{F}{1}=\Trunc{G}{1}=t^{-2}(1+t+t^{2})\), and
\[
 \Trunc{F}{1}\,\Trunc{G}{1}
 =t^{-4}\bigl(1+t+t^{2}\bigr)^{2}
 =t^{-4}+2t^{-3}+3t^{-2}+2t^{-1}+1,
\]
which is already of \(t\)-degree \(0\), so truncating again changes nothing.
The blocks at \(t^{-1}\) and \(t^{0}\) are wrong, by \(2\) and by \(4\).

(3)  Three statements survive.  First, \(\mathfrak F_N\) is an additive
subgroup and a \(\PLpow\)-submodule of \(\PLlau\), so
\(\Trunc{\cdot}{N}\) is \(\K[L]\)-linear and
\(\Trunc{(F+G)}{N}=\Trunc{F}{N}+\Trunc{G}{N}\).  Second, on \(\PLpow\) the set
\(t^{N}\PLpow\) \emph{is} an ideal for \(N\ge0\) and
\(\PLpow/t^{N}\PLpow\cong\K[L][t]/(t^{N})\); this is where the quotient
language is legitimate.  Third, the correct multiplicative law on \(\PLlau\) is
the asymmetric one of \cref{plt:prop:trunc-product}: with \(m=\ordt F\) and
\(r=\ordt G\),
\[
 \Trunc{(FG)}{N}
 =\Trunc{\bigl(\Trunc{F}{N-r}\cdot\Trunc{G}{N-m}\bigr)}{N}.
\]
In the example \(m=r=-2\) and \(N=1\), so the required inputs are
\(\Trunc{F}{3}\) and \(\Trunc{G}{3}\), two blocks more than were used.  A
direct check: \(\Trunc{F}{3}=t^{-2}(1+t+t^{2}+t^{3}+t^{4})\), whose square is
\(t^{-4}(1+2t+3t^{2}+4t^{3}+5t^{4}+\cdots)\), and truncating at \(<1\)
reproduces the correct answer.
\end{proof}

\begin{exercise}[Trap: a symmetric truncation on the Laurent ring]
\label{plt:exer:exr-trap-symmetric-truncation}
A routine is asked for \(\Trunc{(FG)}{0}\) where
\[
 F=X^{2}+X+\frac{L}{X}+\FormalBigOAt{t^{2}}{t},
 \qquad
 G=X^{3}-X^{2}L+\FormalBigOAt{t^{-1}}{t}.
\]
It reasons: ``both factors are known through outer order \(t^{1}\)
respectively \(t^{-2}\); truncating each at \(<0\) and multiplying,
\(FG=(X^{2}+X)(X^{3}-X^{2}L)+\cdots
=X^{5}-X^{4}L+X^{4}-X^{3}L+\cdots\), so
\(\Trunc{(FG)}{0}=X^{5}+(1-L)X^{4}-LX^{3}\).''
Find the error, say which blocks of the answer are nevertheless correct, and
state exactly what precision in \(F\) and \(G\) the requested output needs.
\end{exercise}

\begin{proof}[Solution]
The error is that \(\Trunc{\cdot}{N}\) is not multiplicative on \(\PLlau\)
(\cref{plt:exer:exr-tail-not-ideal}), and the symmetric rule
\eqref{plt:eq:trunc-product-pow} is valid only when both factors have
nonnegative outer order.  Here \(\ordt F=-2\) and \(\ordt G=-3\), so
\cref{plt:prop:trunc-product} demands
\begin{align*}
 \Trunc{(FG)}{0}
 &=\Trunc{\bigl(\Trunc{F}{0-(-3)}\cdot\Trunc{G}{0-(-2)}\bigr)}{0}\\
 &=\Trunc{\bigl(\Trunc{F}{3}\cdot\Trunc{G}{2}\bigr)}{0}.
\end{align*}
That is, \(F\) is needed through outer order \(t^{2}\) and \(G\) through
\(t^{1}\); the routine used \(F\) through \(t^{-1}\) and \(G\) through
\(t^{-1}\).

Which blocks survive?  Write \(F_n=[t^{n}]F\) and \(G_n=[t^{n}]G\), so that
\(F_{-2}=1\), \(F_{-1}=1\), \(F_{0}=0\), \(F_{1}=L\) are known and
\(F_{n}\) is unknown for \(n\ge2\); while \(G_{-3}=1\), \(G_{-2}=-L\) are
known and \(G_{n}\) is unknown for \(n\ge-1\).  Convolving,
\begin{align*}
 [t^{-5}](FG)&=F_{-2}G_{-3}=1,\\
 [t^{-4}](FG)&=F_{-2}G_{-2}+F_{-1}G_{-3}=1-L,\\
 [t^{-3}](FG)&=F_{-2}G_{-1}+F_{-1}G_{-2}+F_{0}G_{-3}=G_{-1}-L .
\end{align*}
The first two are determined by the data and agree with the routine's
\(X^{5}\) and \((1-L)X^{4}\).  The third is not determined: it needs
\(G_{-1}\), which the routine both discarded and never knew.  The printed
\(-LX^{3}\) is what one gets by setting \(G_{-1}=0\), an assumption nothing
supports.  The blocks at \(t^{-2},t^{-1},t^{0}\) are likewise undetermined,
and the routine did not print them at all.

So the data are \emph{insufficient} for the requested output.  Indeed \(F\)
is known only modulo \(\FormalBigOAt{t^{2}}{t}\), which is exactly
\(\Trunc{F}{2}\), one block short of the \(\Trunc{F}{3}\) required; and
\(G\) is known only modulo \(\FormalBigOAt{t^{-1}}{t}\), three blocks short
of \(\Trunc{G}{2}\).  The honest output is
\[
 \Trunc{(FG)}{-3}=X^{5}+(1-L)X^{4},
\]
with the remaining blocks reported as unknown.  This is the error rule of \cref{plt:cor:mul-error-rule} read
backwards.  Two lessons: the symmetric rule
\eqref{plt:eq:trunc-product-pow} is a theorem about \(\PLpow\), not about
\(\PLlau\); and a Laurent factor makes a routine need \emph{more} input
precision than the output it is asked for, never less.
\end{proof}

\begin{exercise}[Counterexample: no monomial prefactor inside the polynomial ring]
\label{plt:exer:exr-no-prefactor}
Several sources write a nonzero \(A\in\PLlau\) in the prefactor normal form
\(A=cX^{\rho}L^{\beta}(1+U)\) with \(c\in\K^{\times}\),
\(\rho\in\IntegerNumbers\), \(\beta\in\NaturalNumbers\) and \(U\in t\PLpow\),
``whenever available''.  Show that for \(A=L+t\) it is not available, and that
it becomes available exactly after passing to \(\PLrat\).  Deduce that no
argument inside \(\PLlau\) may assume the normal form.
\end{exercise}

\begin{proof}[Solution]
Suppose \(L+t=cX^{\rho}L^{\beta}(1+U)=c\,t^{-\rho}L^{\beta}(1+U)\) with
\(U\in t\PLpow\).  Comparing outer orders, \(\ordt(L+t)=0\) while
\(\ordt\bigl(ct^{-\rho}L^{\beta}(1+U)\bigr)=-\rho\), so \(\rho=0\).  Comparing
the \(t^{0}\)-blocks then gives \(cL^{\beta}=L\), so \(c=1\) and \(\beta=1\).
Comparing the \(t^{1}\)-blocks gives
\[
 1=[t^{1}]\bigl(L(1+U)\bigr)=L\cdot u_1,
 \qquad u_1\DefinitionEquals[t^{1}]U\in\K[L],
\]
whence \(u_1=L^{-1}\), which is not a polynomial.  Contradiction.

Over \(\K(L)\) the factorisation exists and is the obvious one,
\(L+t=L\bigl(1+L^{-1}t\bigr)\) with \(L^{-1}t\in t\PLrat\).  So the obstruction
is not the shape of the leading monomial but the failure of \(\K[L]\) to be a
field: the general leading-block factorisation
\cref{plt:prop:lead-factorization} produces \(A=t^{n_0}p_{n_0}(L)(1+U)\) with
\(U\) in \(t\PLrat\), and the coefficient \(p_{n_0}\) can be inverted only in
\(\K(L)\).  Consequently any argument that silently divides by the leading
block --- for instance one that defines \(\dprofile{A}\) through the prefactor,
or that ``normalises \(A\) to be monic'' --- has left \(\PLlau\).
\end{proof}

% ed.: source 2 states the prefactor form "whenever available" without an
% ed.: instance of unavailability; L + t is one, and it is the reason the
% ed.: log-degree profile of \cref{plt:def:lead-profile} is defined blockwise
% ed.: rather than through a normal form.

\begin{exercise}[A prescribed drop in the log-degree profile]
\label{plt:exer:exr-profile-drop}
Let \(d\ge1\).  For each \(k\) with \(1\le k\le d\) construct a pair
\(F,G\in\PLpow\) with \(\ordt F=\ordt G=0\) and
\(\dprofile{F}(1)=\dprofile{G}(1)=d\) such that \(\dprofile{FG}(1)=d-k\);
construct one more pair with \(\dprofile{FG}(1)=-\infty\); and exhibit a pair
with no drop at all.  Then say exactly when equality holds in the bound
\(\dprofile{FG}(N)\le(\dprofile{F}\ast\dprofile{G})(N)\) of
\cref{plt:thm:mul-degree-bound}.
\end{exercise}

\begin{proof}[Solution]
For \(1\le k\le d\) put
\[
 F=1+\bigl(L^{d}+L^{d-k}\bigr)t,
 \qquad
 G=1-L^{d}t .
\]
Both have \(\ordt=0\), \(\dprofile{F}(0)=\dprofile{G}(0)=0\) and
\(\dprofile{F}(1)=\dprofile{G}(1)=d\), the term \(L^{d-k}\) being of strictly
lower degree than \(L^{d}\) because \(k\ge1\).  Now
\[
 [t^{1}](FG)=\bigl(L^{d}+L^{d-k}\bigr)\cdot1+1\cdot\bigl(-L^{d}\bigr)
 =L^{d-k},
\]
so \(\dprofile{FG}(1)=d-k\), while the convolution bound is
\((\dprofile{F}\ast\dprofile{G})(1)=\max\SetOf{0+d,\ d+0}=d\).  For the
degenerate case take \(F=1+L^{d}t\), \(G=1-L^{d}t\), where
\([t^{1}](FG)=0\) and \(\dprofile{FG}(1)=-\infty\).  For no drop take
\(F=G=1+L^{d}t\), where \([t^{1}](FG)=2L^{d}\) has degree \(d\); this is the
generic situation, and it is what makes the drop worth an exercise.

Equality is decided exactly by \cref{plt:thm:mul-degree-bound}.  Let
\(M=(\dprofile{F}\ast\dprofile{G})(N)\) and let \(P_N\) be the set of pairs
\((n,m)\) with \(n+m=N\), \([t^{n}]F\ne0\ne[t^{m}]G\) and
\(\dprofile{F}(n)+\dprofile{G}(m)=M\).  If \(P_N=\emptyset\) then
\([t^{N}](FG)=0\) and equality holds trivially with \(M=-\infty\).  If
\(P_N\ne\emptyset\) then the coefficient of \(L^{M}\) in \([t^{N}](FG)\) equals
\(\sum_{(n,m)\in P_N}\lc\bigl([t^{n}]F\bigr)\lc\bigl([t^{m}]G\bigr)\), and
equality holds if and only if that sum is nonzero.  In the dropping family
\(P_1=\SetOf{(0,1),(1,0)}\) and the sum is \(1\cdot(-1)+1\cdot1=0\), which is
why the degree drops; how far it drops is then decided by the next
coefficients, and the parameter \(k\) tunes exactly that.  In the last pair
the same sum is \(1\cdot1+1\cdot1=2\ne0\) and there is no drop.  Note that
condition (3) of \cref{plt:cor:mul-degree-equality} --- positive leading
coefficients over an ordered field --- excludes the dropping family and
covers the last one, as it must.
\end{proof}

\begin{exercise}[Which enlargement does each expression force?]
\label{plt:exer:exr-scale-promotion}
For each expression name the smallest class in which it can be represented,
choosing among \(\PLlau\), \(\PLrat\), \(\PLit\), a Puiseux--log ring, a ring
with a second logarithmic level, and a logarithmic-exponential transseries
field; and justify both membership and minimality.
\[
 \recip{(L+t)},\qquad
 \sum_{n\ge0}n!\,L^{-n},\qquad
 X^{1/2},\qquad
 \log(X\log X),\qquad
 \EulerE^{-X}.
\]
\end{exercise}

\begin{proof}[Solution]
\emph{\(\recip{(L+t)}\).}  Geometric expansion in \(t\) gives
\begin{equation}
\label{plt:eq:exr-recip-L-plus-t}
 \frac{1}{L+t}=\frac{1}{L}\cdot\frac{1}{1+L^{-1}t}
 =\sum_{n\ge0}(-1)^{n}L^{-n-1}t^{n},
\end{equation}
every coefficient of which lies in \(\K[L,L^{-1}]\subseteq\K(L)\).  So the
element lies in \(\PLrat\).  It does not lie in \(\PLlau\), because the
reciprocal in \(\PLrat\) is unique and \eqref{plt:eq:exr-recip-L-plus-t} has
non-polynomial coefficients; equivalently, \(L+t\) is not a unit of \(\PLlau\)
by \cref{plt:thm:div-unit-criterion}, its leading block \(L\) not being a
nonzero constant.  Minimality among the classes of the volume:
\cref{plt:thm:div-minimal-field} says \(\PLrat\) is the smallest field
\(E((t))\) with \(\K[L]\subseteq E\subseteq\K(L)\) containing \(\PLlau\), and
this example is one of the witnesses that forces the enlargement.

% ed.: source 3 answers this item with "a natural home is K((L^{-1}))((t))",
% ed.: i.e. \PLit.  That overshoots by a whole class: the coefficients in
% ed.: \eqref{plt:eq:exr-recip-L-plus-t} are Laurent *monomials* in L, so even
% ed.: the localisation K[L,L^{-1}]((t)) suffices, and among the four ambient
% ed.: classes the answer is \PLrat.  \PLit is strictly larger
% ed.: (\cref{plt:prop:ring-PLit-strict}), and the next item is what actually
% ed.: needs it.
\emph{\(\sum_{n\ge0}n!L^{-n}\).}  This is a legitimate element of
\(\K((L^{-1}))\), hence of \(\PLit\) (as a \(t^{0}\)-block).  It is not in
\(\K(L)\): a rational function of \(L\) has an eventually linearly recurrent
sequence of Laurent coefficients at \(L=\infty\), and \((n!)_{n\ge0}\) satisfies
no linear recurrence with constant coefficients, since \(n!/(n-1)!=n\) is
unbounded.  So this element separates \(\PLit\) from \(\PLrat\), which is
\cref{plt:prop:ring-PLit-strict}.

\emph{\(X^{1/2}=t^{-1/2}\).}  A Puiseux--log ring with exponent group
\(\tfrac12\IntegerNumbers\) contains it; no class with exponent group
\(\IntegerNumbers\) does, since the outer exponents of every element of
\(\PLlau,\PLrat,\PLit\) are integers.  Enlarging the coefficient field cannot
help: the exponent group is not touched by that.  See
\cref{plt:prop:ring-puiseux-tower}.

\emph{\(\log(X\log X)=L+\log L\).}  The symbol \(\log L\) is not in
\(\K((L^{-1}))\), let alone \(\K(L)\): if it were, it would be a Laurent series
in \(L^{-1}\), so \(\partial_L\log L=L^{-1}\) would lie in the image of
\(\partial_L\) on \(\K((L^{-1}))\), contradicting
\cref{plt:lem:dif-image-laurent} (the residue at \(L=\infty\) of \(L^{-1}\)
is~\(1\)).  A second logarithmic level is forced; this is item~(3) of
\cref{plt:ex:cmp-naive-failures}.

\emph{\(\EulerE^{-X}\).}  By \cref{plt:thm:dif-exp-boundary}, \(\exp F\) lies
in \(\PLlau\) only if \(\ordt F\ge0\); here \(F=-X\) has \(\ordt=-1\).  In
fact \(\EulerE^{-X}\precdom t^{N}\) for every \(N\), so it is not comparable
with any monomial of any power--log scale
(\cref{plt:cor:dif-exp-three-ways}); an exponential transmonomial, hence a
logarithmic-exponential transseries field, is required.
\end{proof}

\begin{exercise}[Admissible supports, and a support that is not one]
\label{plt:exer:exr-supports}
Decide, with proof, which of the following subsets of
\(\IntegerNumbers\times\NaturalNumbers\) are the coefficient support
\(\CoefficientSupportOf{F}\) of some \(F\in\PLlau\), using
\cref{plt:thm:ring-support}:
\begin{align*}
 S_1&=\SetOf{(n,n):n\ge0}, &
 S_2&=\SetOf{(n,j):n\ge0,\ 0\le j\le n},\\
 S_3&=\SetOf{(-n,0):n\ge0}, &
 S_4&=\SetOf{(n,0):n\ge0}\cup\SetOf{(0,j):j\ge0}.
\end{align*}
For each admissible one, exhibit an \(F\) realising it.
\end{exercise}

\begin{proof}[Solution]
Criterion (b) of \cref{plt:thm:ring-support}: the set of outer exponents
occurring must be bounded below, and each fibre \(S_n=\SetOf{j:(n,j)\in S}\)
must be finite.

\(S_1\): outer exponents \(\SetOf{n\ge0}\) bounded below by \(0\); fibre
\(S_n=\SetOf{n}\) finite.  Admissible; realised by
\(F=\sum_{n\ge0}L^{n}t^{n}\).

\(S_2\): exponents bounded below by \(0\); fibre \(S_n=\SetOf{0,\dots,n}\)
finite.  Admissible; realised by
\(F=\sum_{n\ge0}\bigl(\sum_{j=0}^{n}L^{j}\bigr)t^{n}\).

\(S_3\): exponents \(\SetOf{-n:n\ge0}\) are not bounded below.  Not
admissible.  This is the support the divergent family of
\cref{plt:exer:exr-block-local-not-summable}\,(1) would have produced.

\(S_4\): exponents bounded below by \(0\), but the fibre at \(n=0\) is
\(\NaturalNumbers\), infinite: a block would have to be a power series in
\(L\), not a polynomial.  Not admissible in \(\PLlau\).  Nor is it rescued by
passing to the Hahn ambient of \cref{plt:thm:ring-hahn}: that ambient still
demands a reverse well-ordered support, which is criterion~(c), and \(S_4\)
fails (c), since \(\SetOf{(0,j):j\ge0}\) has no \(\succdom\)-greatest element ---
larger \(j\) at fixed \(n\) means larger monomial \(t^{n}L^{j}\).  So \(S_4\)
is not admissible in any of the four classes.
\end{proof}

\section{Arithmetic: product, reciprocal, derivative}
\label{plt:sec:exercises-arithmetic}

\begin{exercise}[The block Cauchy product, and what it consumed]
\label{plt:exer:exr-product}
Let
\[
 A=1+(L+1)t+(L^{2}-1)t^{2},
 \qquad
 B=1+(2-L)t+Lt^{2},
\]
with all undisplayed blocks zero.  Compute \(AB\) through outer order
\(t^{3}\).  Then, without recomputing, say which blocks of \(A\) and of \(B\)
the block \([t^{3}](AB)\) actually depends on, and check your answer against
\cref{plt:thm:mul-finite-determination}.
\end{exercise}

\begin{proof}[Solution]
The blocks convolve:
\begin{align*}
 [t^{0}](AB)&=1,\\
 [t^{1}](AB)&=(L+1)+(2-L)=3,\\
 [t^{2}](AB)&=(L^{2}-1)+(L+1)(2-L)+L
  =L^{2}-1+(-L^{2}+L+2)+L=2L+1,\\
 [t^{3}](AB)&=(L+1)\cdot L+(L^{2}-1)(2-L)
  =L^{2}+L-L^{3}+2L^{2}+L-2\\
 &=-L^{3}+3L^{2}+2L-2 .
\end{align*}
Hence
\[
 AB=1+3t+(2L+1)t^{2}+\bigl(-L^{3}+3L^{2}+2L-2\bigr)t^{3}
   +\FormalBigOAt{t^{4}}{t},
\]
and in fact the product is exact through \(t^{4}\) as well, both inputs being
polynomials in \(t\).

Dependency: since \(\ordt A=\ordt B=0\), \cref{plt:thm:mul-finite-determination}
says \([t^{N}](AB)\) depends exactly on \([t^{n}]A\) and \([t^{m}]B\) for
\(0\le n,m\le N\).  For \(N=3\) that is \(A_0,\dots,A_3\) and
\(B_0,\dots,B_3\); as \(A_3=B_3=0\) by hypothesis, four blocks of each are
consulted and only three of each are nonzero.  Note that the hypothesis
``undisplayed blocks vanish'' is not decoration: without it \(A_3\) and
\(B_3\) are unknown and \([t^{3}](AB)\) is undetermined.
\end{proof}

% ed.: source 3 poses this product only through t^2 and source 1 poses a
% ed.: different pair only through t^3 without saying that the undisplayed
% ed.: blocks must be assumed zero.  The two are merged here and the
% ed.: hypothesis is printed, since \cref{plt:thm:mul-finite-determination} is
% ed.: the point of the exercise.

\begin{exercise}[Reciprocal of a unit, and where the unit hypothesis enters]
\label{plt:exer:exr-reciprocal}
Find \(C\in\PLpow\) with \(\bigl(1+(L-1)t+Lt^{2}\bigr)C=1\), through outer
order \(t^{4}\).  Then explain why \(C\) has polynomial blocks, and give the
one-line reason that the same recurrence applied to \(L+(L-1)t\) does not.
\end{exercise}

\begin{proof}[Solution]
Write \(B=1+B_1t+B_2t^{2}\) with \(B_1=L-1\), \(B_2=L\), and
\(C=\sum_{n\ge0}C_nt^{n}\).  By \cref{plt:thm:div-reciprocal} with \(c=1\),
\(C_0=1\) and \(C_n=-\sum_{j=1}^{n}B_jC_{n-j}\).  Hence
\begin{align*}
 C_1&=-B_1=1-L,\\
 C_2&=-(B_1C_1+B_2C_0)=(L-1)^{2}-L=L^{2}-3L+1,\\
 C_3&=-(B_1C_2+B_2C_1)
   =-(L-1)(L^{2}-3L+1)-L(1-L)\\
   &=-\bigl(L^{3}-4L^{2}+4L-1\bigr)-\bigl(L-L^{2}\bigr)
   =-L^{3}+5L^{2}-5L+1,\\
 C_4&=-(B_1C_3+B_2C_2)\\
   &=-(L-1)\bigl(-L^{3}+5L^{2}-5L+1\bigr)-L\bigl(L^{2}-3L+1\bigr)\\
   &=-\bigl(-L^{4}+6L^{3}-10L^{2}+6L-1\bigr)-\bigl(L^{3}-3L^{2}+L\bigr)\\
   &=L^{4}-7L^{3}+13L^{2}-7L+1 .
\end{align*}
So
\[
 C=1+(1-L)t+(L^{2}-3L+1)t^{2}+(-L^{3}+5L^{2}-5L+1)t^{3}
  +(L^{4}-7L^{3}+13L^{2}-7L+1)t^{4}+\FormalBigOAt{t^{5}}{t}.
\]

Every \(C_n\) is a polynomial because the recurrence divides only by
\(c=B_0=1\), a unit of \(\K\); the whole content of
\cref{plt:thm:div-unit-criterion} is that this is the only case in which the
division stays inside \(\K[L]\).  For \(B=L+(L-1)t\) the recurrence starts
\(C_0=L^{-1}\), which already leaves \(\K[L]\); \(B\) is not a unit of
\(\PLlau\) because its leading block \(L\) is not a nonzero constant.
\end{proof}

\begin{exercise}[A quotient, and the exact guard precision it needs]
\label{plt:exer:exr-quotient-guard}
Let
\[
 A=1+(2L-1)t+(L^{2}+3)t^{2}+(L^{3}-L)t^{3}+\FormalBigOAt{t^{4}}{t},
 \qquad
 B=1+Lt+(L^{2}-1)t^{2}+\FormalBigOAt{t^{3}}{t}.
\]
Compute \(\Trunc{(A/B)}{4}\), or explain why the data do not determine it, and
in either case state the sharp input requirement given by
\cref{plt:prop:div-relative-precision}.
\end{exercise}

\begin{proof}[Solution]
Here \(a=\ordt A=0\) and \(b=\ordt B=0\), so
\eqref{plt:eq:div-guard} reads: \(\Trunc{(A/B)}{N}\) is determined by
\(\Trunc{A}{N}\) and \(\Trunc{B}{N}\).  With \(N=4\) we need \(A\) and \(B\)
each through outer order \(t^{3}\).  \(A\) is given to that precision; \(B\) is
given only modulo \(\FormalBigOAt{t^{3}}{t}\), i.e.\ only \(\Trunc{B}{3}\).
So the block \([t^{3}](A/B)\) is \emph{not} determined: changing \(B_3\) by
\(\delta\) changes the quotient's third block by \(-\delta\).

What is determined is \(\Trunc{(A/B)}{3}\).  Using
\eqref{plt:eq:div-quotient-recurrence}, i.e.\ \(Q_n=A_n-\sum_{j\ge1}B_jQ_{n-j}\)
since \(B_0=1\):
\begin{align*}
 Q_0&=1,\\
 Q_1&=A_1-B_1Q_0=(2L-1)-L=L-1,\\
 Q_2&=A_2-B_1Q_1-B_2Q_0=(L^{2}+3)-L(L-1)-(L^{2}-1)=-L^{2}+L+4 .
\end{align*}
If in addition \(B_3=0\) is assumed --- which is what the source that supplies
this example intends --- then
\[
 Q_3=A_3-B_1Q_2-B_2Q_1-B_3Q_0
  =(L^{3}-L)-L(-L^{2}+L+4)-(L^{2}-1)(L-1)
  =L^{3}-4L-1,
\]
giving
\[
 \frac AB=1+(L-1)t+(-L^{2}+L+4)t^{2}+(L^{3}-4L-1)t^{3}
 +\FormalBigOAt{t^{4}}{t}.
\]
A useful check, and the one the same source recommends, is to multiply the
computed quotient by \(B\) and compare exclusive truncations at the requested
cutoff; by \cref{plt:thm:div-truncated} the truncated quotient is the unique
polynomial of \(t\)-degree \(<4\) passing that test.
\end{proof}

% ed.: source 5 prints the block Q_3 of this example immediately after inputs
% ed.: given only modulo t^3.  The hypothesis B_3 = 0 is a genuine extra
% ed.: assumption, and \cref{plt:prop:div-relative-precision} shows exactly
% ed.: how much is missing; it is stated here rather than assumed silently.

\begin{exercise}[Integer and generalised binomial powers]
\label{plt:exer:exr-binomial}
Let \(F=1+Lt\in\PLpow\).
\begin{enumerate}
\item Compute \([t^{n}]F^{k}\) for every integer \(k\ge0\).
\item Compute \([t^{n}]F^{\kappa}\) for arbitrary \(\kappa\in\K\), where
\(F^{\kappa}\DefinitionEquals\exp(\kappa\log F)\) as in
\cref{plt:cor:mul-roots}, and compare with (1).
\item Compute \(\dprofile{F^{\kappa}}\) and compare with the bound of
\cref{plt:prop:mul-power-profile}.
\item Expand \(\bigl(1+Lt+(L^{2}-1)t^{2}\bigr)^{\kappa}\) through \(t^{2}\).
\end{enumerate}
\end{exercise}

\begin{proof}[Solution]
(1)  \(F^{k}=(1+Lt)^{k}=\sum_{n=0}^{k}\binom{k}{n}L^{n}t^{n}\), so
\([t^{n}]F^{k}=\binom{k}{n}L^{n}\), zero for \(n>k\).

(2)  \(\log F=\log(1+Lt)=\sum_{m\ge1}\frac{(-1)^{m+1}}{m}L^{m}t^{m}\), and
exponentiating gives the generalised binomial series
\(F^{\kappa}=\sum_{n\ge0}\binom{\kappa}{n}L^{n}t^{n}\) with
\(\binom{\kappa}{n}=\FallingFactorial{\kappa}{n}/n!\).  So the coefficient
formula is \emph{the same}, \([t^{n}]F^{\kappa}=\binom{\kappa}{n}L^{n}\); the
only difference is that for \(\kappa=k\in\NaturalNumbers\) the falling
factorial \(\FallingFactorial{k}{n}\) vanishes for \(n>k\) and the series
terminates, while for \(\kappa\notin\NaturalNumbers\) every coefficient is
nonzero and the series is genuinely infinite.  Both are legitimate elements of
\(\PLpow\); the support condition of \cref{plt:thm:ring-support} does not care.

(3)  \(\dprofile{F^{\kappa}}(n)=n\) for all \(n\ge0\) when
\(\kappa\notin\NaturalNumbers\), and for \(0\le n\le k\) when \(\kappa=k\).
The hypothesis of \cref{plt:prop:mul-power-profile} holds with
\(\alpha=1,\beta=0\), and its conclusion \(\dprofile{F^{k}}(N)\le N\) is
attained: the bound is sharp here.

(4)  Put \(U=Lt+(L^{2}-1)t^{2}\), so \(U^{2}=L^{2}t^{2}+\FormalBigOAt{t^{3}}{t}\)
and
\[
 (1+U)^{\kappa}=1+\kappa U+\binom{\kappa}{2}U^{2}+\FormalBigOAt{t^{3}}{t}
 =1+\kappa Lt
  +\Bigl(\kappa(L^{2}-1)+\tfrac{\kappa(\kappa-1)}{2}L^{2}\Bigr)t^{2}
  +\FormalBigOAt{t^{3}}{t},
\]
that is, the \(t^{2}\)-block is
\(\tfrac12\kappa\bigl((\kappa+1)L^{2}-2\bigr)\).
\end{proof}

\begin{exercise}[Counterexample: a nonzero leading block is not enough]
\label{plt:exer:exr-unit-nonzero-block}
A plausible strengthening of \cref{plt:thm:div-unit-criterion} is: ``\(F\) is a
unit of \(\PLlau\) if and only if its leading block is nonzero''.  Refute it,
locate \(\recip{(L+t)}\) exactly, and deduce that \(\PLlau\) is an integral
domain that is not a field --- so that the four classes of
\cref{plt:def:ring-four-classes} really are four different objects.
\end{exercise}

\begin{proof}[Solution]
Every nonzero \(F\in\PLlau\) has, by definition, a nonzero leading block; if
the proposed criterion were correct then every nonzero element would be a
unit and \(\PLlau\) would be a field.  It is not.  Take \(F=L+t\).  Its leading
block is \(L\ne0\).  If \(F\) were a unit of \(\PLlau\), with reciprocal
\(C=\sum_{n\ge0}C_nt^{n}\in\PLpow\), then comparing \(t^{0}\)-blocks in
\(FC=1\) gives \(LC_0=1\), impossible in \(\K[L]\) since \(\degL(LC_0)\ge1\)
for \(C_0\ne0\).  The reciprocal does exist one class up, in \(\PLrat\), and is
given by \eqref{plt:eq:exr-recip-L-plus-t}.

That \(\PLlau\) is a domain is \cref{plt:cor:mul-domain}: outer orders add and
leading blocks multiply in \(\K[L]\), which is a domain, so a product of
nonzero elements is nonzero.  A domain that is not a field is exactly what the
unit criterion describes: the units are the elements \(ct^{n}(1+U)\) with
\(c\in\K^{\times}\), \(n\in\IntegerNumbers\), \(U\in t\PLpow\), a proper
subgroup of the nonzero elements.  The strengthening fails because the
condition is not that the leading block be nonzero but that it be a nonzero
\emph{constant}.
\end{proof}

\begin{exercise}[Trap: a division that leaves the ring]
\label{plt:exer:exr-trap-nonunit}
A student computes \(\recip{(L+t)}\) inside \(\PLlau\) as follows.
``By \cref{plt:thm:div-reciprocal} the reciprocal of
\(B=B_0+B_1t+\cdots\) has \(C_0=B_0^{-1}\) and
\(C_n=-B_0^{-1}\sum_{j=1}^{n}B_jC_{n-j}\).  Here \(B_0=L\), \(B_1=1\), so
\(C_0=L^{-1}\), \(C_1=-L^{-2}\), \(C_2=L^{-3}\), and in general
\(C_n=(-1)^{n}L^{-n-1}\).  Since every \(C_n\) is a Laurent polynomial in
\(L\) and Laurent polynomials are polynomials in \(L\) and \(L^{-1}\), the
answer \(\sum_{n\ge0}(-1)^{n}L^{-n-1}t^{n}\) lies in \(\PLlau\).''
Find every error, and state precisely what is true.
\end{exercise}

\begin{proof}[Solution]
There are two errors, one of hypothesis and one of vocabulary.

\emph{Hypothesis.}  \Cref{plt:thm:div-reciprocal} requires
\(B_0=c\in\K^{\times}\); the whole proof consists in dividing by \(c\), and
\(c^{-1}\) has to exist \emph{in the coefficient ring}.  Here \(B_0=L\) is not
invertible in \(\K[L]\), so the theorem does not apply, and the recurrence as
written is being run over the larger ring \(\K(L)\).  That is legitimate --- it
computes the reciprocal in \(\PLrat\) --- but it is no longer the theorem
quoted.

\emph{Vocabulary.}  The blocks of an element of \(\PLlau\) are
\emph{polynomials} in \(L\), not Laurent polynomials.  \(\K[L]\) is the
coefficient ring; \(\K[L,L^{-1}]\) is a strictly larger ring, and
\(L^{-1}\notin\K[L]\).  So the concluding sentence is false at its last step.

What is true: the computed series is correct, and it is the reciprocal of
\(L+t\) in \(\PLrat\); all its blocks in fact lie in the localisation
\(\K[L,L^{-1}]\), so it lies in the intermediate ring
\(\K[L,L^{-1}]((t))\subsetneq\PLrat\).  It does not lie in \(\PLlau\), by
\cref{plt:exer:exr-unit-nonzero-block}.  The general statement replacing the
misquoted theorem is \cref{plt:prop:div-conditional}: with \(\ordt B=0\) and
\(A\in\PLpow\), the quotient \(A/B\) lies in \(\PLpow\) if and only if every
block of the \(\K(L)\)-recurrence happens to be a polynomial, which happens if
and only if \(A\in B\PLlau\).
\end{proof}

\begin{exercise}[Exact divisibility conditions, block by block]
\label{plt:exer:exr-divisibility}
Let \(B=(L+1)+t\) and let \(A=A_0(L)+A_1(L)t+\FormalBigOAt{t^{2}}{t}\) with
\(A_0,A_1\in\K[L]\).  Find necessary and sufficient conditions on \(A_0,A_1\)
for \(\Trunc{(A/B)}{2}\) to have polynomial blocks, and express them as two
evaluations at a single point of \(\K\).
\end{exercise}

\begin{proof}[Solution]
Work over \(E=\K(L)\), where \(B\) is a unit, and let \(Q=A/B\) have blocks
\(Q_n\).  From \(BQ=A\), comparing blocks,
\[
 (L+1)Q_0=A_0,
 \qquad
 (L+1)Q_1+Q_0=A_1 .
\]
Thus \(Q_0\in\K[L]\) if and only if \((L+1)\mid A_0\); and given that,
\(Q_1\in\K[L]\) if and only if \((L+1)\mid(A_1-Q_0)\).  This is
\cref{plt:prop:div-conditional}\,(2) for the present \(B\).

To turn the two conditions into evaluations, note that for \(P\in\K[L]\) one
has \((L+1)\mid P\) if and only if \(P(-1)=0\).  The first condition is
therefore \(A_0(-1)=0\).  Writing \(A_0=(L+1)Q_0\) and differentiating,
\(A_0'=Q_0+(L+1)Q_0'\), so \(Q_0(-1)=A_0'(-1)\).  The second condition,
\(A_1(-1)=Q_0(-1)\), becomes
\[
 \boxed{\quad A_0(-1)=0
 \qquad\text{and}\qquad
 A_1(-1)=A_0'(-1).\quad}
\]
Both are conditions at the single point \(L=-1\), the root of the leading
block.  Two remarks.  First, the conditions are not vacuous: for
\(A=1\) the first already fails, which is another proof that \(B\) is not a
unit.  Second, the pattern continues: \(Q_n\in\K[L]\) requires
\((L+1)\mid(A_n-Q_{n-1})\), so the whole quotient lies in \(\PLpow\) exactly
when \(A\in B\PLpow\), which is
\cref{plt:prop:div-conditional}\,(3).
\end{proof}

\begin{exercise}[Repeated derivatives of one block, three ways]
\label{plt:exer:exr-block-derivative}
\begin{enumerate}
\item Differentiate \(A=X^{2}+L^{3}+(L^{2}-1)X^{-2}\) once, using the block
formula of \cref{plt:prop:dif-block}.
\item Compute \(\DX^{k}\bigl(L^{3}\bigr)\) for \(k=1,2,3,4\) using
\cref{plt:thm:dif-repeated}.
\item Recompute \(\DX^{4}\bigl(L^{3}\bigr)\) from the Stirling expansion
\(\shiftop{0}{k}=\sum_{j=0}^{k}\stirlingsigned{k}{j}\partial_L^{\,j}\)
of \cref{plt:cor:dif-stirling}, and check the two answers agree.
\end{enumerate}
\end{exercise}

\begin{proof}[Solution]
(1)  \Cref{plt:prop:dif-block} is \(\DX\bigl(t^{n}p(L)\bigr)=t^{n+1}(p'-np)\).
Apply it to the three blocks, with \(n=-2,0,2\):
\begin{align*}
 \DX\bigl(t^{-2}\bigr)&=t^{-1}(0+2)=2X,\\
 \DX\bigl(L^{3}\bigr)&=t^{1}\bigl(3L^{2}-0\bigr)=3L^{2}t,\\
 \DX\bigl(t^{2}(L^{2}-1)\bigr)&=t^{3}\bigl(2L-2(L^{2}-1)\bigr)
   =\bigl(-2L^{2}+2L+2\bigr)t^{3}.
\end{align*}
Hence \(\DX A=2X+3L^{2}X^{-1}+(-2L^{2}+2L+2)X^{-3}\).

(2)  For a \(t^{0}\)-block, \(\DX^{k}p=t^{k}\shiftop{0}{k}p\) with
\(\shiftop{0}{k}=(\partial_L-k+1)\cdots(\partial_L-1)\partial_L\).  With
\(p=L^{3}\):
\begin{align*}
 \shiftop{0}{1}L^{3}&=3L^{2},\\
 \shiftop{0}{2}L^{3}&=(\partial_L-1)(3L^{2})=6L-3L^{2},\\
 \shiftop{0}{3}L^{3}&=(\partial_L-2)(6L-3L^{2})=6-6L-12L+6L^{2}
  =6L^{2}-18L+6,\\
 \shiftop{0}{4}L^{3}&=(\partial_L-3)(6L^{2}-18L+6)
  =12L-18-18L^{2}+54L-18=-18L^{2}+66L-36 .
\end{align*}
So \(\DX L^{3}=3L^{2}t\), \(\DX^{2}L^{3}=(6L-3L^{2})t^{2}\),
\(\DX^{3}L^{3}=(6L^{2}-18L+6)t^{3}\),
\(\DX^{4}L^{3}=(-18L^{2}+66L-36)t^{4}\).

(3)  The signed Stirling numbers of the first kind in row \(4\) are
\(\stirlingsigned{4}{0}=0\), \(\stirlingsigned{4}{1}=-6\),
\(\stirlingsigned{4}{2}=11\), \(\stirlingsigned{4}{3}=-6\),
\(\stirlingsigned{4}{4}=1\).  Applying
\(\sum_j\stirlingsigned{4}{j}\partial_L^{\,j}\) to \(L^{3}\), whose
derivatives are \(3L^{2},6L,6,0\),
\[
 -6\cdot3L^{2}+11\cdot6L-6\cdot6+1\cdot0=-18L^{2}+66L-36,
\]
which agrees with (2).  Equivalently
\(X^{4}\dfrac{\Differential^{4}}{\Differential X^{4}}(\log X)^{3}
 =\sum_{j=0}^{4}\stirlingsigned{4}{j}\bigl((\log X)^{3}\bigr)^{(j)}\);
note that the sum must start at \(j=0\), because
\(\stirlingsigned{k}{0}=0\) only for \(k\ge1\) and equals \(1\) at \(k=0\).
\end{proof}

% ed.: source 1 prints the Stirling identity with the sum starting at j = 1.
% ed.: That is correct only for k >= 1; the uniform range j = 0..k with the
% ed.: empty-product convention covers k = 0 as well, and is used above.

\begin{exercise}[Antiderivatives of a single block, resonant and not]
\label{plt:exer:exr-antiderivative}
Find polynomial-logarithmic antiderivatives of
\[
 F_1=t^{2}(L+1)
 \qquad\text{and}\qquad
 F_2=t(L+1),
\]
and verify each by differentiating.  Explain why the two computations look
different, and identify the block index at which the difference occurs.
\end{exercise}

\begin{proof}[Solution]
By \cref{plt:prop:dif-block}, \(\DX\bigl(t^{n}Q\bigr)=t^{n+1}(Q'-nQ)\), so an
antiderivative of \(t^{n+1}P\) is \(t^{n}Q\) with \(Q'-nQ=P\).

\emph{\(F_1=t^{2}(L+1)\), so \(n=1\).}  Solve \(Q'-Q=L+1\) in \(\K[L]\).  With
\(Q=\alpha L+\beta\): \(\alpha-\alpha L-\beta=L+1\), so \(\alpha=-1\) and
\(\alpha-\beta=1\), i.e.\ \(\beta=-2\).  Thus \(Q=-L-2\) and
\[
 \int F_1\,\Differential X=-t(L+2)+c,\qquad c\in\K .
\]
Check: \(\DX\bigl(-t(L+2)\bigr)=t^{2}\bigl(-1-(-(L+2))\bigr)=t^{2}(L+1)\).

\emph{\(F_2=t(L+1)\), so \(n=0\).}  Now the equation is \(Q'=L+1\), with
solution \(Q=\tfrac12L^{2}+L\) up to a constant, and
\[
 \int F_2\,\Differential X=\tfrac12L^{2}+L+c .
\]
Check: \(\DX\bigl(\tfrac12L^{2}+L\bigr)=t\bigl(L+1-0\bigr)=t(L+1)\).

The difference is the operator being inverted.  For \(n\ne0\) the map
\(Q\mapsto Q'-nQ\) is bijective on \(\K[L]\): it is triangular with diagonal
\(-n\ne0\) in the monomial basis, and the explicit inverse is the terminating
sum \(Q=-\sum_{j\ge0}P^{(j)}/n^{j+1}\) of
\eqref{plt:eq:dif-antiderivative-formula}.  This is the nonresonant case of
\cref{plt:lem:dif-triangular}, and it \emph{lowers} the \(L\)-degree, or at
worst preserves it.  For \(n=0\) the map is \(Q\mapsto Q'\), which is
surjective on \(\K[L]\) but has kernel \(\K\) and \emph{raises} the degree by
one; this is the resonant block of \cref{plt:lem:dif-resonant}, and it is the
only place where the log-degree profile of an antiderivative exceeds that of
its derivative, as recorded in \eqref{plt:eq:dif-profile}.  The resonant index
is the one whose outer exponent in \(F\) is \(1\), i.e.\ the block \(t^{1}\) of
the integrand.
\end{proof}

\begin{exercise}[Counterexample: surjectivity of \(\DX\) is not inherited]
\label{plt:exer:exr-antiderivative-obstruction}
By \cref{plt:thm:dif-antiderivative} the derivation \(\DX\) maps \(\PLlau\)
onto \(\PLlau\), with kernel \(\K\).  A plausible strengthening is that it
maps each of the four classes onto itself.  Refute this.  Show that the
element
\[
 F=\frac tL=\frac1{X\log X},
\]
which lies in \(\PLrat\) and in \(\PLit\), has no antiderivative in
\(\PLit\), and identify the antiderivative it does have.
\end{exercise}

\begin{proof}[Solution]
Suppose \(G=\sum_{n}g_n(L)t^{n}\in\PLit\), so each
\(g_n\in\K((L^{-1}))\), satisfies \(\DX G=t/L\).  By
\cref{plt:prop:dif-block}, extended over the coefficient ring
\(\K((L^{-1}))\), the block of \(\DX G\) at \(t^{n+1}\) is
\(g_n'-ng_n\), where \({}'=\partial_L\).  Comparing at \(t^{1}\), i.e.\
\(n=0\), gives
\[
 g_0'=\frac1L .
\]
Now \(\partial_L\) maps \(\K((L^{-1}))\) into itself, and for
\(g_0=\sum_{k\ge k_0}c_kL^{-k}\) one has
\(g_0'=\sum_{k\ge k_0}(-k)c_kL^{-k-1}\).  The coefficient of \(L^{-1}\) in
\(g_0'\) is the one with \(-k-1=-1\), i.e.\ \(k=0\), and it equals
\(-0\cdot c_0=0\).  So \(L^{-1}\) is never in the image: the residue at
\(L=\infty\),
\(\operatorname{res}(g)\DefinitionEquals[L^{-1}]g\), vanishes identically on
derivatives (\cref{plt:lem:dif-res-derivative}), and
\(\operatorname{res}(t/L)=1\ne0\).  Hence no such \(g_0\) exists, and
\(F\) has no antiderivative in \(\PLit\); a fortiori none in \(\PLrat\)
(\cref{plt:lem:dif-image-laurent}).

The obstruction is exactly one function, the residue, and it is met by
adjoining one new symbol: \(\Lambda=\log\log X\) satisfies
\(\DX\Lambda=t/L\), because
\(\frac{\Differential}{\Differential X}\log\log X=\frac{1}{X\log X}\).  This is
\cref{plt:cor:dif-loglog}, and it is where the second logarithmic level of the
whole theory comes from: not from composition, but from integration.  Note
that no contradiction with \cref{plt:thm:dif-antiderivative} arises, since
\(t/L\notin\PLlau\).
\end{proof}

\begin{exercise}[Which formal exponentials stay in the class?]
\label{plt:exer:exr-exponentials}
Decide, for each \(F\) below, whether \(\exp F\in\PLlau\), and if not, name
the smallest enlargement in which it lives:
\[
 L,\qquad 2L+t,\qquad t,\qquad \tfrac12L,\qquad L^{2},\qquad t^{-1}.
\]
State the criterion you used and check it in each case.
\end{exercise}

\begin{proof}[Solution]
The criterion is \cref{plt:thm:dif-exp-boundary}: \(\exp F\in\PLlau\) if and
only if \(\ordt F\ge0\) and \([t^{0}]F=a+mL\) with \(m\in\IntegerNumbers\),
\(a\in\K\) and \(\exp a\in\K\); and then
\(\exp F=(\exp a)t^{-m}\exp\bigl(F-[t^{0}]F\bigr)\) with the last factor in
\(1+t\PLpow\).

\(F=L\): \(\ordt=0\), \([t^{0}]F=0+1\cdot L\), \(m=1\), \(a=0\).  Yes:
\(\exp L=t^{-1}=X\).

\(F=2L+t\): \(\ordt=0\), \([t^{0}]F=2L\), \(m=2\), \(a=0\).  Yes:
\(\exp F=t^{-2}\exp(t)=X^{2}\sum_{k\ge0}t^{k}/k!\in\PLlau\).

\(F=t\): \(\ordt=1\ge0\), \([t^{0}]F=0\), \(m=0\), \(a=0\).  Yes, and
\(\exp t\in1+t\PLpow\).

\(F=\tfrac12L\): \(m=\tfrac12\notin\IntegerNumbers\).  No; \(\exp F=X^{1/2}\)
needs the Puiseux--log ring with exponent group \(\tfrac12\IntegerNumbers\)
(\cref{plt:cor:dif-exp-hahn}).

\(F=L^{2}\): \([t^{0}]F=L^{2}\) is not affine in \(L\).  No.  Here the failure
is not repaired by any coefficient-field or exponent-group enlargement of a
power--log scale: \(\exp(L^{2})=X^{\log X}\) grows faster than every power of
\(X\) and slower than \(\exp(cX)\), so it is a new transmonomial
(\cref{plt:cor:dif-exp-three-ways}).

\(F=t^{-1}=X\): \(\ordt=-1<0\).  No; \(\EulerE^{X}\) is an exponential
transmonomial, and the smallest home is a logarithmic-exponential transseries
field.

A remark worth extracting: the three failures are of three different kinds ---
a fractional slope, a nonaffine \(t^{0}\)-block, a negative outer order ---
and \cref{plt:cor:dif-exp-three-ways} says these are the only three.
\end{proof}

\section{Composition}
\label{plt:sec:exercises-composition}

\begin{exercise}[Substituting a constant shift into one monomial]
\label{plt:exer:exr-generator-substitution}
Let \(a\in\K\) and \(G=X+a\).  Compute \((tL)\compose G\) through outer order
\(t^{3}\), and then \((tL)\compose(cX)\) exactly, for \(c\in\K^{\times}\) with
\(\log c\in\K\).
\end{exercise}

\begin{proof}[Solution]
By \cref{plt:prop:cmp-generator-substitution}, with \(H=a\),
\[
 t\compose G=\frac{t}{1+at}=t-at^{2}+a^{2}t^{3}+\FormalBigOAt{t^{4}}{t},
 \qquad
 L\compose G=L+\log(1+at)=L+at-\frac{a^{2}}{2}t^{2}+\FormalBigOAt{t^{3}}{t}.
\]
Composition is a ring homomorphism (\cref{plt:thm:cmp-substitution-hom}), so
\((tL)\compose G=(t\compose G)(L\compose G)\), and multiplying,
\begin{align*}
 (tL)\compose(X+a)
 &=\bigl(t-at^{2}+a^{2}t^{3}\bigr)
   \Bigl(L+at-\tfrac{a^{2}}{2}t^{2}\Bigr)+\FormalBigOAt{t^{4}}{t}\\
 &=tL+a(1-L)t^{2}+a^{2}\Bigl(L-\tfrac32\Bigr)t^{3}
   +\FormalBigOAt{t^{4}}{t}.
\end{align*}

For the dilation \(G=cX\) the substitution is exact and finite:
\(t\compose G=c^{-1}t\) and \(L\compose G=L+\log c\), so
\[
 (tL)\compose(cX)=c^{-1}t\,(L+\log c),
\]
and more generally
\(F\compose(cX)=\sum_{n}c^{-n}p_n(L+\log c)t^{n}\) for
\(F=\sum_np_n(L)t^{n}\).  Two things are visible here that matter later.
First, a dilation changes \emph{both} generators, the power by \(c^{-n}\) and
the logarithm by the shift \(\log c\); forgetting the second is the failure
mode of \cref{plt:exer:exr-trap-scale}.  Second, the identity is a statement
about \(\PLlau\) only if \(\log c\in\K\); over \(\K=\RationalNumbers\) and
\(c=2\) the right-hand side does not live in \(\PLlau\) at all, and the field
must be enlarged first.
\end{proof}

\begin{exercise}[A composition through two blocks]
\label{plt:exer:exr-compose-two-blocks}
Let \(F(Y)=Y+(\log Y)^{2}\) and \(G(X)=X-L\).  Compute \(F\compose G\) through
outer order \(t^{2}\).
\end{exercise}

\begin{proof}[Solution]
Write \(F=t^{-1}+L^{2}\) in its own generators and \(G=X+H\) with \(H=-L\), so
\(H\in\PLpow\) and \(G\in\Gnear\).  By
\cref{plt:prop:cmp-generator-substitution},
\[
 L\compose G=L+\log(1-Lt)
 =L-Lt-\tfrac12L^{2}t^{2}+\FormalBigOAt{t^{3}}{t}.
\]
Therefore
\begin{align*}
 (L\compose G)^{2}
 &=\Bigl(L-Lt-\tfrac12L^{2}t^{2}\Bigr)^{2}+\FormalBigOAt{t^{3}}{t}\\
 &=L^{2}-2L^{2}t+\bigl(L^{2}-L^{3}\bigr)t^{2}+\FormalBigOAt{t^{3}}{t},
\end{align*}
the \(t^{2}\)-block being \(2\cdot L\cdot(-\tfrac12L^{2})+(-L)^{2}\).  Adding
\(G=X-L\),
\[
 F\compose G=X+\bigl(L^{2}-L\bigr)-\frac{2L^{2}}{X}
  +\frac{L^{2}-L^{3}}{X^{2}}+\FormalBigOAt{t^{3}}{t}.
\]
A check on the shape: \(\ordt(F\compose G)=\ordt F=-1\), as
\cref{plt:prop:cmp-generator-substitution} requires, and no block acquired a
degree in \(L\) larger than \(2+\) (the number of substituted logarithms),
consistent with \cref{plt:ex:cmp-log-degree-growth}.
\end{proof}

\begin{exercise}[Composition is multiplicative, from the generators]
\label{plt:exer:exr-compose-multiplicative}
Let \(G=X+H\) with \(H\in\PLpow\), so \(G\in\Gnear\).  Prove that
\[
 (F\Phi)\compose G=(F\compose G)(\Phi\compose G)
 \qquad\text{for all }F,\Phi\in\PLlau,
\]
\emph{directly} from the exact substitution formula
\eqref{plt:eq:cmp-near-identity-master}, without invoking
\cref{plt:thm:cmp-substitution-hom}.
\end{exercise}

\begin{proof}[Solution]
Let \(F=\sum_{n\ge n_0}p_n(L)t^{n}\) and
\(\Phi=\sum_{m\ge m_0}q_m(L)t^{m}\), and let \(G=X+H\in\Gnear\).  Write
\[
 \tau\DefinitionEquals t\compose G=\frac{t}{1+tH},
 \qquad
 \lambda\DefinitionEquals L\compose G=L+\log(1+tH),
\]
so that \eqref{plt:eq:cmp-near-identity-master} reads
\(F\compose G=\sum_n\tau^{n}p_n(\lambda)\).  Two facts are needed.

\emph{(a) The families are summable.}  Since \(\ordt(\tau)=1\) and
\(\ordt(\lambda-L)\ge1\), the summand \(\tau^{n}p_n(\lambda)\) has outer order
exactly \(n\) when \(p_n\ne0\), and is zero otherwise
(\cref{plt:prop:cmp-generator-substitution}); in either case its order is at
least \(n\), so
\(\SetOf{n:\ordt<N}\) is finite for every \(N\) and
\cref{plt:def:trunc-summable} is satisfied.  The same holds for \(\Phi\) and
for the Cauchy product.

\emph{(b) Multiplication distributes over summable families.}  This is
\cref{plt:thm:trunc-summability}\,(iv): if \((S_i)\) and \((T_j)\) are
summable then so is \((S_iT_j)_{(i,j)}\) and
\(\bigl(\sum_iS_i\bigr)\bigl(\sum_jT_j\bigr)=\sum_{i,j}S_iT_j\).

Now compute.  Let \(r_N=\sum_{n+m=N}p_nq_m\) be the blocks of \(F\Phi\).  By
(b),
\[
 (F\compose G)(\Phi\compose G)
 =\Bigl(\sum_n\tau^{n}p_n(\lambda)\Bigr)
  \Bigl(\sum_m\tau^{m}q_m(\lambda)\Bigr)
 =\sum_{n,m}\tau^{n+m}\,p_n(\lambda)q_m(\lambda).
\]
Group by \(N=n+m\); the inner sums are finite because for fixed \(N\) only
finitely many pairs \((n,m)\) with \(n\ge n_0\), \(m\ge m_0\) satisfy
\(n+m=N\).  Since evaluation of a polynomial at the fixed element
\(\lambda\in\PLpow\) is a ring homomorphism \(\K[L]\to\PLpow\),
\(\sum_{n+m=N}p_n(\lambda)q_m(\lambda)=r_N(\lambda)\).  Hence
\[
 (F\compose G)(\Phi\compose G)=\sum_{N}\tau^{N}r_N(\lambda)
 =(F\Phi)\compose G ,
\]
the last equality being \eqref{plt:eq:cmp-near-identity-master} applied to
\(F\Phi\).  Additivity is the same computation with one index, and
\(1\compose G=1\); so \(\Phi_G\colon F\mapsto F\compose G\) is a
\(\K\)-algebra endomorphism of \(\PLlau\).
\end{proof}

\begin{exercise}[The sharp Taylor budget, in the exclusive convention]
\label{plt:exer:exr-taylor-budget}
Let \(F\in\PLlau\) with \(\ordt F\ge a\), let \(H\in\PLpow\) with
\(\ordt H\ge h\ge0\), and let \(N\in\IntegerNumbers\).
\begin{enumerate}
\item Prove that \(\Trunc{\bigl(F\compose(X+H)\bigr)}{N}\) is unchanged if the
Taylor family \(\sum_{k\ge0}H^{k}\DX^{k}F/k!\) is truncated after the index
\(k=\Floor{(N-1-a)/(h+1)}\).
\item Show the bound cannot be lowered, by exhibiting \(F,H,N\) for which the
last retained term does change the answer.
\item A source states the same budget as \(k\le\Floor{(N-a)/(h+1)}\).  Explain
the discrepancy and say whether it is harmful.
\end{enumerate}
\end{exercise}

\begin{proof}[Solution]
(1)  By \cref{plt:cor:dif-budget}, \(\ordt(\DX^{k}F)\ge a+k\), and
\(\ordt(H^{k})\ge kh\).  So the \(k\)-th Taylor term has
\[
 \ordt\Bigl(\frac{H^{k}}{k!}\DX^{k}F\Bigr)\ \ge\ a+k+kh=a+k(h+1).
\]
Such a term is invisible to \(\Trunc{\cdot}{N}\) as soon as
\(a+k(h+1)\ge N\), i.e.\ \(k\ge(N-a)/(h+1)\).  The retained indices are
therefore those with \(a+k(h+1)<N\), that is \(k(h+1)<N-a\), that is
\(k\le\Floor{(N-1-a)/(h+1)}\) for integers.  Equivalently the number of
retained terms is \(K=\Ceiling{(N-a)/(h+1)}\), which is
\cref{plt:prop:alg-taylor-budget}.

(2)  Take \(F=L\) (so \(a=0\)), \(H=1\) (so \(h=0\)) and \(N=3\).  The budget
retains \(k\le2\).  Now
\(\DX L=t\), \(\DX^{2}L=\DX(t)=t^{2}(0-1)=-t^{2}\),
\(\DX^{3}L=\DX(-t^{2})=-t^{3}(0-2)=2t^{3}\), so the Taylor sum is
\[
 L+t-\tfrac12t^{2}+\tfrac13t^{3}-\cdots,
\]
which is exactly \(L\compose(X+1)=L+\log(1+t)\), as it must be.  The term
\(k=2\) contributes \(-\tfrac12t^{2}\), a block strictly below the cutoff
\(N=3\); dropping it changes \(\Trunc{\cdot}{3}\).  So the last retained index
is genuinely needed, and the bound is sharp.

(3)  The source is working with the inclusive cutoff \([F]_{\le N}\).  Since
\([F]_{\le N}=\Trunc{F}{N+1}\), its \(N\) is our \(N-1\), and its budget
\(k\le\Floor{(N-a)/(h+1)}\) becomes, in our normalisation, exactly
\(k\le\Floor{(N-1-a)/(h+1)}\): the two agree.  Read as an exclusive bound,
however, the source's formula retains one term too many whenever
\((h+1)\mid(N-a)\).  That is harmless for \emph{correctness} --- the extra term
is invisible below \(t^{N}\) --- but not for cost, and not for a routine that
uses the budget in reverse to decide how much of \(H\) to compute.  This is
% ed.: this sentence cited \cref{plt:prop:alg-shift}, the polynomial-shift
% ed.: formula, which has nothing to do with the truncation convention.  The
% ed.: discrepancy is recorded with the sharp budget itself.
the discrepancy recorded at \cref{plt:prop:alg-taylor-budget}, whose last
line prints the exclusive form \(k\le\Floor{(N-1-a)/(h+1)}\) explicitly for
this reason.
\end{proof}

% ed.: source 1 records the budget as k <= floor((N - r)/(h + 1)) without
% ed.: saying which truncation convention its N refers to.  Under the
% ed.: exclusive cutoff, which the catalogue makes normative, that form admits
% ed.: one Taylor term too many.  Part (3) prints the conversion rather than
% ed.: silently shifting the index, because the same formula is what a routine
% ed.: would apply in reverse to size its inputs, where the extra term is not
% ed.: free.

\begin{exercise}[Trap: how much of the inner argument is needed?]
\label{plt:exer:exr-trap-inner-precision}
A routine must return \(\Trunc{\bigl(F\compose(X+H)\bigr)}{1}\) for
\(F=X^{2}\).  It argues: ``the output is requested through outer order
\(t^{0}\); composition is filtration-continuous
(\cref{plt:prop:cmp-filtration-continuity}), so \(H\) is needed only through
outer order \(t^{0}\), i.e.\ only \(\Trunc{H}{1}\).''
Compute \(\Trunc{\bigl(F\compose(X+H)\bigr)}{1}\) for
\(H=h_0+h_1t+h_2t^{2}+\cdots\), find the error, and state the correct input
requirement for a general \(F\).
\end{exercise}

\begin{proof}[Solution]
Here \(F=t^{-2}\), and composition with a near-identity map is exact:
\[
 F\compose(X+H)=(X+H)^{2}=t^{-2}+2t^{-1}H+H^{2},
\]
so, block by block,
\[
 [t^{-2}]=1,
 \qquad
 [t^{-1}]=2h_0,
 \qquad
 [t^{0}]=2h_1+h_0^{2},
\]
and
\(\Trunc{\bigl(F\compose(X+H)\bigr)}{1}
 =X^{2}+2h_0X+\bigl(2h_1+h_0^{2}\bigr)\).
The block \(h_1\) enters.  The routine's answer would omit it, and would be
wrong by \(2h_1\).

The error is a confusion of absolute with relative precision.  Filtration
continuity says \(\ordt\bigl(F\compose(X+B)-F\compose(X+C)\bigr)
\ge\ordt(F)+\ordt(B-C)+1\) (\cref{plt:cor:grp-gain}), and the summand
\(\ordt(F)\) is negative here.  Unwinding: a perturbation of \(H\) at
outer order \(j\) moves the composite at outer order \(a+j+1\) with
\(a=\ordt F\).  It is therefore invisible below \(t^{N}\) exactly when
\(a+j+1\ge N\), so the required input is
\[
 \Trunc{H}{\,N-a-1}.
\]
For \(F=X^{2}\), \(a=-2\) and \(N=1\): \(\Trunc{H}{2}\), one block more than
the routine used.  Only when \(\ordt F\ge-1\) does \(\Trunc{H}{N}\) suffice.
This is the sharp statement of
\cref{plt:thm:cmp-finite-determination}, and the general principle behind it is
\cref{plt:cor:alg-relative-precision}: what a routine conserves is relative,
not absolute, precision.
\end{proof}

\begin{exercise}[Counterexample: the near-identity group is not abelian]
\label{plt:exer:exr-commutator}
Let \(f,g\in\K[L]\), \(F=X+f(L)\), \(G=X+g(L)\).
\begin{enumerate}
\item Compute the first block of \(F\compose G-G\compose F\).
\item Deduce that \(\Gnear\) is not abelian, and characterise the pairs
\((f,g)\) for which \(F\) and \(G\) commute to first order.
\item Show that vanishing of the first block does \emph{not} imply
commuting, by computing the next block for \(f=L\), \(g=2L\).
\end{enumerate}
\end{exercise}

\begin{proof}[Solution]
(1)  By \cref{plt:prop:cmp-near-monoid},
\(F\compose G=X+g+f\compose(X+g)\), and by
\cref{plt:cor:grp-gain}, \eqref{plt:eq:grp-gain-leading},
\(f\compose(X+g)-f\equiv g\,\DX f\pmod{t^{2}\PLpow}\), where
\(\DX f=f'(L)t\).  Hence
\[
 F\compose G=X+f+g+g f'\,t+\FormalBigOAt{t^{2}}{t},
\]
and by symmetry
\(G\compose F=X+f+g+fg'\,t+\FormalBigOAt{t^{2}}{t}\).  Subtracting,
\begin{equation}
\label{plt:eq:exr-commutator}
 F\compose G-G\compose F=\bigl(gf'-fg'\bigr)t+\FormalBigOAt{t^{2}}{t},
\end{equation}
which is \cref{plt:cor:tay-commutator}.

(2)  Take \(f=L\), \(g=L^{2}\): then \(gf'-fg'=L^{2}-2L^{2}=-L^{2}\ne0\), so
\(F\compose G\ne G\compose F\) and \(\Gnear\) is not abelian, although it is
associative (\cref{plt:thm:cmp-associativity}).  The first block vanishes if
and only if \(gf'=fg'\), i.e.\ \((f/g)'=0\) in \(\K(L)\) for \(g\ne0\), i.e.\
\(f\) and \(g\) are proportional over \(\K\).  Equivalently, the vector fields
\(f\partial_L\) and \(g\partial_L\) commute.

(3)  Take the proportional pair \(f=L\), \(g=2L\), for which
\(gf'-fg'=2L-2L=0\), so the first block of the commutator vanishes.  The maps
still do not commute.  Indeed by
\cref{plt:prop:cmp-generator-substitution},
\begin{align*}
 (X+L)\compose(X+2L)&=X+3L+2Lt-2L^{2}t^{2}+\FormalBigOAt{t^{3}}{t},\\
 (X+2L)\compose(X+L)&=X+3L+2Lt-L^{2}t^{2}+\FormalBigOAt{t^{3}}{t},
\end{align*}
the two \(t^{2}\)-blocks coming from \(\log(1+2Lt)\) and from
\(2\log(1+Lt)\) respectively.  Hence
\[
 F\compose G-G\compose F=-L^{2}t^{2}+2L^{3}t^{3}
 -\tfrac72L^{4}t^{4}+\FormalBigOAt{t^{5}}{t}.
\]
So proportionality kills the first block only.  Structurally this is exactly
what \cref{plt:thm:grp-gradedlie} predicts: the graded Lie algebra of the
contact filtration is not abelian, and vanishing of the bracket of the leading
symbols only pushes the commutator one step down the filtration
(\cref{plt:prop:grp-commutator}).
\end{proof}

\begin{exercise}[A scale change that produces a second logarithm exactly]
\label{plt:exer:exr-scale-loglog}
Let \(F(X)=X+\log X+X^{-1}\) and \(G(X)=3X^{2}L^{-1}\).  Compute
\(F\compose G\) exactly, and say which class each summand lives in.  Then do
the same for \(F(Y)=\log Y\), \(G(X)=X^{2}(\log X)^{3}\), and identify the
smallest new symbol forced in each case.
\end{exercise}

\begin{proof}[Solution]
\emph{First pair.}  \(G=3t^{-2}L^{-1}\in\PLrat\) is monomial-leading with
\(\rho=2\), \(c=3\), \(\sigma=-1\) and no correction.  Composition acts on the
generators by \(t\mapsto\recip{G}\) and \(L\mapsto\log G\), and here both are
exact and finite:
\[
 \recip{G}=\tfrac13Lt^{2},
 \qquad
 \log G=\log 3+2L-\log L .
\]
Hence
\[
 F\compose G=3X^{2}L^{-1}+2L-\log L+\log3+\tfrac13Lt^{2},
\]
a \emph{finite} expression.  The summands \(3X^{2}L^{-1}\) and
\(\tfrac13Lt^{2}\) lie in \(\PLrat\); \(2L\) lies in \(\PLlau\); \(\log3\)
requires \(\log3\in\K\); and \(\log L=\log\log X\) lies in no Puiseux--log
ring over \(\K[L]\), by the residue argument of
\cref{plt:exer:exr-antiderivative-obstruction}.  The point the source makes is
worth keeping: the second logarithmic level appears here \emph{exactly}, in a
closed-form finite composite, not merely inside a remainder.

\emph{Second pair.}  \(\log G=2L+3\log L\), so again the new symbol is
\(L_2=\log L=\log\log X\), and the output lies in a two-level coefficient class,
not in \(\PLlau\).  This is item~(3) of \cref{plt:ex:cmp-naive-failures} and
\cref{plt:prop:grp-xlogx}.

In both cases the forced symbol is read off from the images of the generators
alone: \(\recip{G}\) decides whether the exponent group must grow, \(\log G\)
decides whether the logarithmic depth must grow.  That is the whole content of
the admissibility checklist of \cref{plt:sec:scale-change}.
\end{proof}

\begin{exercise}[Counterexample: a monomial-leading inner argument is not enough]
\label{plt:exer:exr-constant-inner}
\Cref{plt:thm:cmp-substitution-hom} composes \(F\) with an inner argument
\(G=cX^{\rho}(1+U)\) subject to \(\rho>0\).  Show that the hypothesis
\(\rho>0\) cannot be weakened to \(\rho\ge0\).  Precisely: write
\(\tau=t\compose G\), \(\lambda=L\compose G\), and suppose
\(\ordt\tau=0\).  Prove that the substitution family
\(\bigl(\tau^{n}p_n(\lambda)\bigr)_{n}\) fails to be summable as soon as
\(p_n(\lambda)\ne0\) for infinitely many \(n\), whatever the coefficient
field; exhibit such an \(F\); and identify the exact failure.  Then show that
the quantifier cannot be improved to ``every \(F\) with infinite coefficient
support''.
\end{exercise}

\begin{proof}[Solution]
Suppose \(G\) is such that \(\tau\DefinitionEquals t\compose G\) has
\(\ordt\tau=0\); the simplest instance is \(G=c\in\K^{\times}\), where
\(\tau=c^{-1}\) and \(\lambda=\log c\in\K\).  Let
\(F=\sum_{n\ge n_0}p_n(L)t^{n}\).  The defining formula for the composite is
the sum of the family \(\bigl(\tau^{n}p_n(\lambda)\bigr)_{n}\).  Since
\(\PLlau\) is a domain and \(\ordt(\tau^{n})=n\ordt(\tau)=0\),
\[
 \ordt\bigl(\tau^{n}p_n(\lambda)\bigr)=\ordt\bigl(p_n(\lambda)\bigr)
 \qquad\text{for every }n .
\]
So no member of the family can be pushed to a high outer order by the factor
\(\tau^{n}\): the whole family sits at the orders of the substituted
coefficients.  For the constant inner argument \(G=c\) those are scalars, so
every nonzero member sits in the single block \(t^{0}\).  Hence if
\(p_n(\log c)\ne0\) for infinitely many \(n\) --- take \(F=1/(1-t)\), where
\(p_n=1\) for all \(n\ge0\) --- the set \(\SetOf{n:\ordt<1}\) is infinite,
\cref{plt:def:trunc-summable} fails, and the \(t^{0}\)-block of the putative
composite would be
\[
 \sum_{n\ge n_0}c^{-n}p_n(\log c),
\]
an infinite sum of scalars, which is not an element of \(\K\) unless it
happens to terminate.  No enlargement of the coefficient field repairs this:
the obstruction is that infinitely many summands land in one block, which is a
statement about the outer filtration alone.  The same argument applies to any
\(G\) with \(\rho=0\), since then \(\ordt(t\compose G)=0\).

% ed.: the exercise as inherited asked for the failure "for every F whose
% ed.: coefficient support is infinite".  That quantifier is too strong, and
% ed.: the following paragraph is the counterexample to it.
\emph{The quantifier.}  It is not enough that \(F\) have infinite coefficient
support: what must be infinite is the set of indices surviving the
substitution.  With \(G=c\) take
\[
 F=\frac{L-\log c}{1-t}=\sum_{n\ge0}\bigl(L-\log c\bigr)t^{n},
\]
whose support \(\SetOf{(n,0),(n,1):n\ge0}\) is infinite.  Every substituted
coefficient \(p_n(\lambda)=\log c-\log c\) vanishes, the family is
\((0,0,\dots)\), which is vacuously summable, and \(F\compose c=0\) is
perfectly well defined.  The honest statement is the one proved above, with
the hypothesis on \(p_n(\lambda)\) and not on \(F\).

For contrast, when \(\rho>0\) one has \(\ordt(\tau)=\rho>0\) and
\(\ordt(\tau^{n})=n\rho\to+\infty\), so the family is summable --- in a Hahn
ambient with exponent group containing \(\rho\IntegerNumbers\), which is
\cref{plt:prop:cmp-fractional-rho}.  This is why \(\rho>0\) is a hypothesis
and not a normalisation.
\end{proof}

\section{Reversion}
\label{plt:sec:exercises-reversion}

\begin{exercise}[The first-order form of the inverse, and its sharp error]
\label{plt:exer:exr-inverse-first-order}
Let \(H\in\PLpow\) with \(q=\ordt(H)\).  Show
\[
 \rev{(X+H)}=X-H+\FormalBigOAt{t^{\,2q+1}}{t},
\]
and show that the exponent \(2q+1\) is attained for \(H=L\).
\end{exercise}

\begin{proof}[Solution]
By \cref{plt:thm:rev-existence}, \(\rev{(X+H)}=X+B\) where \(B\) is the
\(t\)-adic limit of the Picard iterates \(B_0=0\),
\(B_{j+1}=-H\compose(X+B_j)\), and
\(\ordt(B-B_j)\ge(j+1)q+j\).  Taking \(j=1\), so that \(B_1=-H\compose X=-H\),
\[
 \ordt\bigl(B+H\bigr)\ \ge\ 2q+1,
\]
which is the claim.

Sharpness at \(H=L\): here \(q=0\) and the claim is
\(\rev{(X+L)}=X-L+\FormalBigOAt{t}{t}\).  By
\cref{plt:thm:om-expansion} the inverse is
\(X-L+\LamP{1}(L)t+\cdots=X-L+Lt+\cdots\), so \(B+L=Lt+\cdots\) has outer
order exactly \(1=2\cdot0+1\).

Two remarks.  First, the estimate is a statement about \(\ordt\) only: the
correction \(-H\) is exact in the leading block by
\cref{plt:thm:rev-existence}\,(4), and everything beyond it is error.  Second,
the same proof gives the general Picard rate: \(j\) iterations from \(B_0=0\)
are correct through outer block \((j+1)q+j-1\), so each further iteration
gains \(q+1\) blocks.
\end{proof}

% ed.: source 1 states this with the hypothesis q >= 1.  The hypothesis is
% ed.: unnecessary: the Picard estimate \eqref{plt:eq:rev-picard-accuracy}
% ed.: holds for every alpha = q >= 0, and the case q = 0 -- the Wright omega
% ed.: case -- is exactly the one in which the statement is sharp.

\begin{exercise}[The universal inverse blocks through order two]
\label{plt:exer:exr-universal-blocks}
Let \(a,b,c\in\K[L]\) and
\(F=X+a(L)+b(L)t+c(L)t^{2}+\FormalBigOAt{t^{3}}{t}\), so \(F\in\Gnear\).
Write \(\rev{F}=X+b_0+b_1t+b_2t^{2}+\FormalBigOAt{t^{3}}{t}\).  Derive
\(b_0,b_1,b_2\) as universal differential polynomials in \(a,b,c\).
\end{exercise}

\begin{proof}[Solution]
Put \(G=X+B\) with \(B=b_0+b_1t+b_2t^{2}+\cdots\) and impose
\(F\compose G=X\).  By \cref{plt:prop:cmp-generator-substitution},
\[
 L\compose G=L+\log(1+tB)
 =L+b_0t+\Bigl(b_1-\tfrac12b_0^{2}\Bigr)t^{2}+\FormalBigOAt{t^{3}}{t},
 \qquad
 t\compose G=t\bigl(1-b_0t+\FormalBigOAt{t^{2}}{t}\bigr).
\]
Write \(\delta\DefinitionEquals(L\compose G)-L\), of outer order \(\ge1\).
Expanding each coefficient by \cref{plt:prop:alg-shift}, a finite Taylor sum:
\begin{align*}
 a(L\compose G)&=a+a'\delta+\tfrac12a''\delta^{2}+\FormalBigOAt{t^{3}}{t}\\
  &=a+a'b_0t+\Bigl(a'\bigl(b_1-\tfrac12b_0^{2}\bigr)
      +\tfrac12a''b_0^{2}\Bigr)t^{2}+\FormalBigOAt{t^{3}}{t},\\
 b(L\compose G)\cdot(t\compose G)
  &=\bigl(b+b'b_0t\bigr)\cdot t\bigl(1-b_0t\bigr)
   +\FormalBigOAt{t^{3}}{t}\\
  &=bt+\bigl(b'b_0-bb_0\bigr)t^{2}+\FormalBigOAt{t^{3}}{t},\\
 c(L\compose G)\cdot(t\compose G)^{2}&=ct^{2}+\FormalBigOAt{t^{3}}{t}.
\end{align*}
Adding these to \(G=X+B\) and equating \(F\compose G\) to \(X\) block by
block:
\begin{align*}
 t^{0}:&\quad b_0+a=0,\\
 t^{1}:&\quad b_1+a'b_0+b=0,\\
 t^{2}:&\quad b_2+a'\bigl(b_1-\tfrac12b_0^{2}\bigr)+\tfrac12a''b_0^{2}
        +b'b_0-bb_0+c=0 .
\end{align*}
Solving in order, with \(b_0=-a\) and \(b_1=aa'-b\),
\begin{equation}
\label{plt:eq:exr-universal-b2}
 b_0=-a,
 \qquad
 b_1=aa'-b,
 \qquad
 b_2=\frac{a^{2}a'}{2}-\frac{a^{2}a''}{2}-a(a')^{2}
      +a'b+ab'-ab-c .
\end{equation}
Three consistency checks.  For \(b=c=0\) this is
\(b_2=\tfrac12a^{2}a'-a(a')^{2}-\tfrac12a^{2}a''\), the form used in
% ed.: the check was printed as "a = aL with a in K", using the letter a for
% ed.: both the coefficient a(L) and the scalar; renamed.
\cref{plt:exer:exr-reverse-quadratic}.  For \(a=\alpha L\) with
\(\alpha\in\K\) constant
and \(b=c=0\) it gives
\(\alpha^{3}(L^{2}/2-L)=\alpha^{3}\LamP{2}(L)\), matching
\cref{plt:cor:lag-poly-perturbation}.  For \(a\), \(b\) constants and \(c=0\)
it gives \(b_2=-ab\), matching \eqref{plt:eq:rev-ab-b2} at \(\degL a=0\).
\end{proof}

\begin{exercise}[Trap: an incomplete universal formula]
\label{plt:exer:exr-trap-universal-formula}
A source records the third universal inverse block of
\(F=X+a(L)+b(L)t+c(L)t^{2}\) as
\[
 b_2\ \overset{?}{=}\ \frac12a^{2}a'-a(a')^{2}-ab+a'b ,
\]
and verifies it on the example \(a=2L\), \(b=3\), \(c=0\), where it is
correct.  Show that the formula is nevertheless false in general, by producing
three independent witnesses --- one for each missing term --- and state the
correct formula.  Explain why the source's verification could not have
detected the error.
\end{exercise}

\begin{proof}[Solution]
Comparing with \eqref{plt:eq:exr-universal-b2}, the quoted formula omits
\(-\tfrac12a^{2}a''\), \(+ab'\) and \(-c\).  Each is detected by a separate
witness.

\emph{Missing \(-\tfrac12a^{2}a''\).}  Take \(a=L^{2}\), \(b=c=0\).  Then
\(a'=2L\), \(a''=2\), and \eqref{plt:eq:exr-universal-b2} gives
\[
 b_2=\frac{L^{4}\cdot2L}{2}-\frac{L^{4}\cdot2}{2}-L^{2}\cdot4L^{2}
 =L^{5}-5L^{4},
\]
whereas the quoted formula gives \(L^{5}-4L^{4}\).  The discrepancy is
\(-L^{4}\).  The correct value is confirmed independently by
\cref{plt:cor:lag-poly-perturbation}, which for \(P=L^{2}\) predicts
\(\degL b_2=d(N+1)-1=5\) with leading coefficient
\((d/N)a_d^{N+1}=2/2=1\); both hold for \(L^{5}-5L^{4}\).

\emph{Missing \(+ab'\).}  Take \(a=1\), \(b=L\), \(c=0\).  Then
\eqref{plt:eq:exr-universal-b2} gives \(b_2=-ab+ab'=-L+1\), while the quoted
formula gives \(-L\).

\emph{Missing \(-c\).}  Take \(a=b=0\), \(c=1\).  Then
\(F=X+t^{2}\), and \(\rev{F}=X-t^{2}+\FormalBigOAt{t^{3}}{t}\) directly by
\cref{plt:exer:exr-inverse-first-order}; so \(b_2=-1\), while the quoted
formula gives \(0\).

The source's verification could not detect any of this because its single test
case has \(a''=0\) (its \(a\) is affine in \(L\)), \(b'=0\) (its \(b\) is a
constant) and \(c=0\).  All three omitted terms vanish simultaneously on that
example.  This is the standard hazard of validating a universal formula on one
instance; the residual certificate of
\cref{plt:cor:rev-finite-certificate} is the correct test, because it is run
against \(F\) itself and is an identity in \(\K[L][[t]]\), not a numerical
sample.
\end{proof}

\begin{exercise}[Reversion of a quadratic logarithmic perturbation]
\label{plt:exer:exr-reverse-quadratic}
Reverse \(F=X+L^{2}\) through outer order \(t^{3}\), and check the result
against the degree law of \cref{plt:thm:lag-degree-sharp} at
\(N=1,2,3\).
\end{exercise}

\begin{proof}[Solution]
Use the Lagrange--B\"urmann formula in the form
\cref{plt:cor:lag-poly-perturbation}: for \(F=X+P(L)\) only the summand
\(r=N+1\) contributes to \(b_N\), so
\[
 b_N=\frac{(-1)^{N+1}}{(N+1)!}
   \Bigl[\prod_{s=0}^{N-1}(\partial_L-s)\Bigr]P^{N+1}.
\]
With \(P=L^{2}\):

\(N=0\): \(b_0=-P=-L^{2}\).

\(N=1\): \(b_1=\tfrac12\partial_L(P^{2})=PP'=L^{2}\cdot2L=2L^{3}\).

\(N=2\):
\(b_2=-\tfrac16(\partial_L-1)\partial_L(P^{3})\).  Now
\(\partial_L(P^{3})=3P^{2}P'\) and
\((\partial_L-1)(3P^{2}P')=6P(P')^{2}+3P^{2}P''-3P^{2}P'\), so
\[
 b_2=\tfrac12P^{2}P'-P(P')^{2}-\tfrac12P^{2}P''
 =L^{5}-4L^{4}-L^{4}=L^{5}-5L^{4}.
\]

\(N=3\): \(b_3=\tfrac1{24}(\partial_L-2)(\partial_L-1)\partial_L(P^{4})\).
With \(P^{4}=L^{8}\): \(\partial_L L^{8}=8L^{7}\);
\((\partial_L-1)(8L^{7})=56L^{6}-8L^{7}\);
\((\partial_L-2)(56L^{6}-8L^{7})=336L^{5}-56L^{6}-112L^{6}+16L^{7}
 =16L^{7}-168L^{6}+336L^{5}\).  Dividing by \(24\),
\[
 b_3=\tfrac23L^{7}-7L^{6}+14L^{5}.
\]
Collecting,
\[
 \rev{(X+L^{2})}
 =X-L^{2}+2L^{3}t+\bigl(L^{5}-5L^{4}\bigr)t^{2}
  +\Bigl(\tfrac23L^{7}-7L^{6}+14L^{5}\Bigr)t^{3}
  +\FormalBigOAt{t^{4}}{t}.
\]

\emph{Degree law.}  Here \(d=\degL P=2\) and \(a_d=1\).
\Cref{plt:thm:lag-degree-sharp} predicts, for \(N\ge1\),
\(\degL b_N=d(N+1)-1=2N+1\) and
\(\bigl[L^{2N+1}\bigr]b_N=(d/N)a_d^{N+1}=2/N\).  For \(N=1,2,3\) that is
degrees \(3,5,7\) with leading coefficients \(2,1,\tfrac23\); all three agree
with the computation.  Note that \(\degL b_0=2\ne d\cdot1-1=1\): the law
excludes \(N=0\), exactly as the Lambert degree law excludes \(n=0\)
(\cref{plt:exer:exr-degree-law-zero}).
\end{proof}

% ed.: sources 2 and 3 both quote a general third-order inverse block; source
% ed.: 3's version (the b = c = 0 case above) is correct, source 2's is not
% ed.: (\cref{plt:exer:exr-trap-universal-formula}).  Source 3 checks the
% ed.: degree law only at N = 1, 2; the block b_3 above is computed here and
% ed.: extends the check to N = 3.

\begin{exercise}[A perturbation with a genuine \(t\)-block]
\label{plt:exer:exr-reverse-with-t}
Reverse \(F=X+L+at\) through outer order \(t^{3}\), for \(a\in\K\), and verify
the answer with the residual certificate of
\cref{plt:cor:rev-finite-certificate}.
\end{exercise}

\begin{proof}[Solution]
Set \(G=X+B\), \(B=b_0+b_1t+b_2t^{2}+b_3t^{3}+\cdots\).  Since
\(F=X+L+at\), \cref{plt:prop:cmp-generator-substitution} gives
\[
 F\compose G=X+B+L+\log(1+tB)+\frac{at}{1+tB}.
\]
Expand each piece through \(t^{3}\).  Writing \(u=tB=b_0t+b_1t^{2}+b_2t^{3}\),
\[
 \log(1+u)=b_0t+\Bigl(b_1-\tfrac12b_0^{2}\Bigr)t^{2}
  +\Bigl(b_2-b_0b_1+\tfrac13b_0^{3}\Bigr)t^{3}+\FormalBigOAt{t^{4}}{t},
\]
\[
 \frac{at}{1+u}=at-ab_0t^{2}+a\bigl(b_0^{2}-b_1\bigr)t^{3}
 +\FormalBigOAt{t^{4}}{t}.
\]
Equating \(F\compose G=X\) block by block:
\begin{align*}
 t^{0}&:\ b_0+L=0 &&\Rightarrow\ b_0=-L,\\
 t^{1}&:\ b_1+b_0+a=0 &&\Rightarrow\ b_1=L-a,\\
 t^{2}&:\ b_2+b_1-\tfrac12b_0^{2}-ab_0=0
  &&\Rightarrow\ b_2=\tfrac12L^{2}-L-aL+a,\\
 t^{3}&:\ b_3+b_2-b_0b_1+\tfrac13b_0^{3}+a\bigl(b_0^{2}-b_1\bigr)=0 .
\end{align*}
Substituting the earlier blocks in the last line,
\(b_0b_1=-L(L-a)=-L^{2}+aL\), \(\tfrac13b_0^{3}=-\tfrac13L^{3}\),
\(b_0^{2}-b_1=L^{2}-L+a\), so
\[
 b_3=-b_2+b_0b_1-\tfrac13b_0^{3}-a\bigl(b_0^{2}-b_1\bigr)
 =\frac{L^{3}}{3}-\frac{3L^{2}}{2}-aL^{2}+3aL+L-a^{2}-a .
\]
Hence
\begin{align*}
 \rev{(X+L+at)}
 ={}&X-L+(L-a)t+\Bigl(\frac{L^{2}}2-L-aL+a\Bigr)t^{2}\\
 &+\Bigl(\frac{L^{3}}3-\frac{3L^{2}}2-aL^{2}+3aL+L-a^{2}-a\Bigr)t^{3}
 +\FormalBigOAt{t^{4}}{t}.
\end{align*}

\emph{Certificate.}  By \cref{plt:thm:rev-triangular} the residual of the
truncation \(G_2=X+b_0+b_1t+b_2t^{2}\) is
\(\RevErr{F}{G_2}=-b_3t^{3}+\FormalBigOAt{t^{4}}{t}\); substituting the
computed \(b_2\) into the \(t^{3}\)-equation above reproduces exactly
\(-b_3\), which is the check.  Two sanity tests: at \(a=0\) the blocks reduce
to \(-L,L,\tfrac12L^{2}-L,\tfrac13L^{3}-\tfrac32L^{2}+L\), i.e.\ to
\(\LamP{0},\dots,\LamP{3}\) as \cref{plt:thm:om-expansion} requires; and the
whole answer agrees with \eqref{plt:eq:rev-ab-b0}--\eqref{plt:eq:rev-ab-b3}
at \(a\mapsto1\), \(b\mapsto a\).
\end{proof}

\begin{exercise}[The residual measures the error, on both sides]
\label{plt:exer:exr-residual-both-sides}
Let \(F=X+L\) and let
\(H=X-L+Lt+\tfrac12L^{2}t^{2}\), the second Picard iterate.  Compute
\(\RevErr{F}{H}\) and \(H\compose F-X\) through outer order \(t^{3}\), and
verify \cref{plt:prop:rev-residual-order} on this pair, including the part
that says the two residuals agree only in their \emph{leading} block.
\end{exercise}

\begin{proof}[Solution]
Write \(G=\rev F=\WrightOmega\), whose blocks are the Lambert polynomials
(\cref{plt:thm:om-expansion}).  The error is
\(E=H-G\), with
\[
 [t^{0}]E=[t^{1}]E=0,
 \qquad
 [t^{2}]E=\tfrac12L^{2}-\LamP{2}(L)=\tfrac12L^{2}-\bigl(\tfrac12L^{2}-L\bigr)=L,
\]
so \(e\DefinitionEquals\ordt(E)=2\) and \([t^{2}]E=L\).

\emph{Right residual.}  With \(B=-L+Lt+\tfrac12L^{2}t^{2}\),
\(\RevErr{F}{H}=B+L+\log(1+tB)\).  Expanding,
\[
 \RevErr{F}{H}=L\,t^{2}-\frac{L^{2}(2L-9)}{6}\,t^{3}
 +\FormalBigOAt{t^{4}}{t}.
\]
\emph{Left residual.}  \(H\compose F-X=(F-X)+B\compose F=L+B\compose F\), and
expanding with \eqref{plt:eq:cmp-near-identity-master},
\[
 H\compose F-X=L\,t^{2}-\frac{L^{2}(2L+3)}{6}\,t^{3}
 +\FormalBigOAt{t^{4}}{t}.
\]

Both residuals have outer order exactly \(2=e\), and both have leading block
\(L=[t^{2}]E\), which is \cref{plt:prop:rev-residual-order}.  Their
\(t^{3}\)-blocks differ, by \(2L^{2}\): the proposition asserts equality of
the leading blocks only, and this pair shows that no more is true.  The
practical consequence is that either residual certifies correctness through a
requested block (\cref{plt:cor:rev-finite-certificate}), and either reads off
the next inverse block (\cref{plt:prop:rev-next-block}), but neither may be
used as an approximation to the error beyond its leading block.
\end{proof}

\begin{exercise}[Counterexample: the derivative certificate is strictly weaker]
\label{plt:exer:exr-derivative-certificate}
One source asserts that the derivative residual
\(\DX H-\recip{\bigl((\DX F)\compose H\bigr)}\) ``often detects a wrong block
one order earlier'' than the composition residual \(\RevErr{F}{H}\).  Refute
this, and quantify the true relation.
\begin{enumerate}
\item Show the derivative residual has outer order at least \(e+1\), where
\(e=\ordt(H-\rev F)\): it is one level \emph{later}, never earlier.
\item Show it is blind at \(e=0\): construct \(F,H\) with \(e=0\) for which the
composition residual has order \(0\) while the derivative residual has order
\(2\).
\end{enumerate}
\end{exercise}

\begin{proof}[Solution]
(1)  Let \(G=\rev F\) and \(E=H-G\).  Differentiating \(F\compose G=X\) gives
\((\DX F)\compose G\cdot\DX G=1\), so
\(\DX G=\recip{\bigl((\DX F)\compose G\bigr)}\) and the derivative residual is
\[
 \DX H-\recip{\bigl((\DX F)\compose H\bigr)}
 =\DX E
  +\Bigl(\recip{\bigl((\DX F)\compose G\bigr)}
        -\recip{\bigl((\DX F)\compose H\bigr)}\Bigr).
\]
The first summand has \(\ordt\ge\ordt(E)+1=e+1\) by
\cref{plt:cor:dif-budget}.  For the second, write \(F=X+A\) with
\(A\in\PLpow\), so \(\DX F=1+\DX A\) with \(\ordt(\DX A)\ge\ordt(A)+1\ge1\).
Then
\((\DX F)\compose H-(\DX F)\compose G=(\DX A)\compose H-(\DX A)\compose G\)
has \(\ordt\ge\ordt(\DX A)+e+1\ge e+2\) by \cref{plt:cor:grp-gain}.  Both
\((\DX F)\compose H\) and \((\DX F)\compose G\) lie in \(1+t\PLpow\), and for
units \(u,v\) of that form \(\recip u-\recip v=(v-u)\recip{(uv)}\) with
\(\ordt(uv)=0\), so taking reciprocals preserves the order of a difference
exactly; the second summand therefore has \(\ordt\ge e+2\).
Hence the derivative residual has \(\ordt\ge e+1\), whereas
\(\RevErr{F}{H}\) has \(\ordt\) exactly \(e\)
(\cref{plt:prop:rev-residual-order}).  It is therefore always one level later,
and can never be earlier.  This is
\cref{plt:prop:rev-derivative-certificate}.

(2)  Take \(F=X+L\), \(G=\rev F=\WrightOmega\) and \(H=G+5\), so
\(E=5\in\K^{\times}\) and \(e=0\).  The composition residual has order \(0\):
indeed \(\RevErr{F}{H}=5+\bigl(\text{terms of order}\ge1\bigr)\), whose
\(t^{0}\)-block is \(5\ne0\).  But \(\DX E=\DX(5)=0\), so the first summand in
the display above vanishes identically, and a direct computation gives
\[
 \DX H-\recip{\bigl((\DX F)\compose H\bigr)}=-5\,t^{2}
 +\FormalBigOAt{t^{3}}{t},
\]
of order \(2\).  So at \(e=0\) the derivative certificate is off by two
levels, and a routine that used it alone would accept an inverse that is
wrong in its very first block by an additive constant.  The reason is
structural: \(\DX\) kills \(\K\), so the derivative certificate is blind to
exactly the kernel of \(\DX\), which is the block that \(\ordt\) calls level
\(0\).
\end{proof}

\begin{exercise}[A near-identity idempotent is the identity]
\label{plt:exer:exr-idempotent}
Let \(J\in\Gnear\) satisfy \(J\compose J=J\).  Prove \(J=X\), first from the
group structure and then, without it, by comparing the first nonzero
correction block.
\end{exercise}

\begin{proof}[Solution]
\emph{First proof.}  By \cref{plt:thm:grp-group}, \(\Gnear\) is a group under
composition.  Composing \(J\compose J=J\) on the right with \(\rev J\) gives
\(J=X\).

\emph{Second proof, without inverses.}  Write \(J=X+B\) with \(B\in\PLpow\) and
suppose \(B\ne0\); let \(\beta=\ordt(B)\ge0\).  By
\cref{plt:prop:cmp-near-monoid},
\[
 J\compose J=X+B+B\compose(X+B),
\]
so \(J\compose J=J\) forces \(B\compose(X+B)=0\).  But
\cref{plt:prop:cmp-generator-substitution} says that substitution by a
near-identity map preserves the outer valuation exactly:
\(\ordt\bigl(B\compose(X+B)\bigr)=\ordt(B)=\beta<\infty\).  A series of finite
outer order is nonzero, a contradiction.  Hence \(B=0\) and \(J=X\).

The second proof is the more informative one: it shows that the conclusion
uses only that \(\Phi_{J}\) is injective and valuation-preserving, not that
inverses exist.  Consequently the same argument works in the near-identity
monoid at any logarithmic depth (\cref{plt:thm:grp-multilog-group}), and in
Puiseux--log enlargements.
\end{proof}

\begin{exercise}[One Picard step against one Newton step]
\label{plt:exer:exr-picard-vs-newton}
Let \(F=X+\log X\) and take \(G_0=X-\log X\) as the starting approximation.
\begin{enumerate}
\item Compute the next Picard iterate \(\Phi(G_0)\), where
\(\Phi(X+B)=X-A\compose(X+B)\) with \(A=L\), and say how many blocks of
\(\rev F\) it gets right.
\item Compute the Newton iterate \(\mathcal N_F(G_0)\) of
\cref{plt:def:rev-newton} far enough to say how many blocks \emph{it} gets
right, and compare with \cref{plt:thm:rev-newton-doubling}.
\end{enumerate}
\end{exercise}

\begin{proof}[Solution]
Throughout, \(\rev F=\WrightOmega=X-L+Lt+\bigl(\tfrac12L^{2}-L\bigr)t^{2}
+\bigl(\tfrac13L^{3}-\tfrac32L^{2}+L\bigr)t^{3}+\cdots\)
(\cref{plt:ex:rev-omega}), and \(G_0=X-L\) is correct through block \(0\), so
\(e_0\DefinitionEquals\ordt\bigl(G_0-\rev F\bigr)=1\).

(1)  \(\Phi(G_0)=X-L\compose(X-L)=X-\bigl(L+\log(1-Lt)\bigr)\), so
\[
 \Phi(G_0)=X-L+Lt+\frac{L^{2}}2t^{2}+\frac{L^{3}}3t^{3}+\cdots .
\]
Blocks \(0\) and \(1\) are correct (\(-L\) and \(L\)); block \(2\) is
\(\tfrac12L^{2}\) against the true \(\tfrac12L^{2}-L\), so the error has outer
order exactly \(2\).  One Picard step gained one block, in agreement with
\eqref{plt:eq:rev-picard-accuracy}: with \(\alpha=\ordt(A)=0\), each step gains
\(\alpha+1=1\).

(2)  \(\DX F=1+t\), so
\(\mathcal N_F(G_0)=G_0-\RevErr{F}{G_0}\cdot
 \recip{\bigl((1+t)\compose G_0\bigr)}\).  Here
\(\RevErr{F}{G_0}=\log(1-Lt)=-Lt-\tfrac12L^{2}t^{2}-\tfrac13L^{3}t^{3}-\cdots\)
and \((1+t)\compose G_0=1+t/(1-Lt)=1+t+Lt^{2}+\cdots\), whose reciprocal is
\(1-t+(1-L)t^{2}+\cdots\).  Multiplying and subtracting,
\[
 \mathcal N_F(G_0)=X-L+Lt+\Bigl(\tfrac12L^{2}-L\Bigr)t^{2}
  +\Bigl(\tfrac13L^{3}-\tfrac32L^{2}+L\Bigr)t^{3}
  +\Bigl(\tfrac14L^{4}-\tfrac{11}6L^{3}+\tfrac52L^{2}-L\Bigr)t^{4}
  +\cdots,
\]
whose blocks \(0,1,2,3\) all agree with \(\WrightOmega\) and whose block \(4\)
does not (the true value is
\(\tfrac14L^{4}-\tfrac{11}6L^{3}+3L^{2}-L=\LamP{4}(L)\), differing by
\(\tfrac12L^{2}\)).  So the error has outer order exactly \(4\).

\Cref{plt:thm:rev-newton-doubling} predicts
\(\ordt\bigl(\RevErr{F}{\mathcal N_F(G_0)}\bigr)\ge2m+2+\alpha\) with
\(m=\ordt\bigl(\RevErr{F}{G_0}\bigr)=1\) and \(\alpha=0\), i.e.\ at least
\(4\); the computation attains it.  Summary: from a start correct through
block \(0\), one Picard step reaches block \(1\) and one Newton step reaches
block \(3\).  The comparison is not, however, a cost comparison: the Newton
step performs a composition, a reciprocal and a multiplication, and which
scheme is cheaper depends on the cost model
(\cref{plt:prop:alg-reversion-comparison}).
\end{proof}

\begin{exercise}[Counterexample: ``smaller than \(X\)'' is not the admissibility test (hard)]
\label{plt:exer:exr-summability-threshold}
\Cref{plt:prop:lag-summability-sharp} requires \(\ordt(A)>-1\), which on
\(\PLlau\) means \(A\in\PLpow\).  A plausible substitute is the dominance
condition \(A\precdom X\), i.e.\ ``\(A\) is asymptotically smaller than the
identity''.  Refute it.
\begin{enumerate}
\item Show that \(A=X/L\) satisfies \(A=\LittleOAt{X}{X\to+\infty}\) yet has
\(\ordt(A)=-1\).
\item For \(A=\alpha X\) with \(\alpha\in\K^{\times}\), compute every
Lagrange--B\"urmann summand exactly and show they all have the same outer
order, so the family is never summable.
% ed.: part (3) compares moduli, so it needs an absolute value on the
% ed.: coefficient field; the inherited statement left K unnamed.
\item For this part only, take \(\K\subseteq\ComplexNumbers\) with the usual
absolute value.  Show that the formal sum of those summands is a geometric
series whose \emph{numerical} value reproduces the correct inverse
\(\rev{\bigl((1+\alpha)X\bigr)}=X/(1+\alpha)\) when
\(\AbsoluteValue{\alpha}<1\) and is
divergent when \(\AbsoluteValue{\alpha}\ge1\).  Conclude that the threshold is
not a technicality.
\end{enumerate}
\end{exercise}

\begin{proof}[Solution]
(1)  \(A=X/L=t^{-1}L^{-1}\) lies in \(\PLrat\) with \(\ordt(A)=-1\).  As a
function, \(A(X)/X=1/\log X\to0\), so \(A\precdom X\) in the analytic sense of
\cref{plt:lem:ring-dominance-analytic}.  So the dominance test is passed and
the valuation test is failed: logarithmic gains are invisible to \(\ordt\).
This is the content of \cref{plt:rmk:lag-nearid-meaning}.

(2)  Let \(A=\alpha X=\alpha t^{-1}\).  By \cref{plt:prop:dif-block},
\(\DX\bigl(t^{n}\bigr)=t^{n+1}(0-n)=-nt^{n+1}\), so by induction
\[
 \DX^{k}\bigl(t^{-r}\bigr)=r(r-1)\cdots(r-k+1)\,t^{-r+k},
\]
and with \(k=r-1\),
\(\DX^{\,r-1}\bigl(t^{-r}\bigr)=r!\,t^{-1}\).  Hence the \(r\)-th
Lagrange--B\"urmann summand is
\begin{equation}
\label{plt:eq:exr-lag-alpha-term}
 \frac{(-1)^{r}}{r!}\DX^{\,r-1}\bigl(A^{r}\bigr)
 =\frac{(-1)^{r}\alpha^{r}}{r!}\,r!\,t^{-1}
 =(-\alpha)^{r}\,X .
\end{equation}
Every summand has outer order \(-1\), so \(\SetOf{r:\ordt<0}\) is infinite and
the family fails \cref{plt:def:trunc-summable}.  This matches the exact order
formula \eqref{plt:eq:lag-order-exact}: with \(\alpha_{\ordt}=\ordt(A)=-1\),
\(\ordt\bigl(\DX^{r-1}(A^{r})\bigr)=r(\ordt(A)+1)-1=-1\) for every \(r\).

(3)  Here \(\K\subseteq\ComplexNumbers\).  Formally summing
\eqref{plt:eq:exr-lag-alpha-term} as a geometric series,
\[
 X+\sum_{r\ge1}(-\alpha)^{r}X
 \ \text{``}=\text{''}\
 X\cdot\frac{1}{1+\alpha}
 =\frac{X}{1+\alpha},
\]
which is indeed \(\rev{F}\) for \(F=X+\alpha X=(1+\alpha)X\).  For
\(\AbsoluteValue{\alpha}<1\) the scalar series converges and the manipulation
is a valid
computation with numbers --- but not with formal series: the family is not
summable, so the left-hand side names no element of \(\PLlau\), and the
agreement is a coincidence of the geometric series, not a theorem.  For
\(\alpha=2\), say, the same manipulation produces
\(X\sum_{r\ge0}(-2)^{r}\), a divergent series of scalars, while the true
inverse \(\rev{(3X)}=X/3\) is perfectly well defined.  So the recipe is not
merely unjustified at the threshold; it breaks.

The moral: \(\ordt(A)>-1\) is exactly the condition under which the summands'
orders tend to \(+\infty\), and no growth comparison between \(A\) and \(X\)
substitutes for it.  The two agree on \(\PLlau\), whose value group is
\(\IntegerNumbers\), and separate on the Puiseux--log enlargement, where
\(\ordt(A)>-1\) admits fractional orders in \((-1,0)\) that
``\(A\in\PLpow\)'' would exclude.
\end{proof}

\begin{exercise}[The inverse of \(X\log X\) needs two logarithms (hard)]
\label{plt:exer:exr-xlogx}
Let \(F(X)=X\log X\) and \(Y=F(X)\to+\infty\).  Put \(S=\log Y\) and
\(\Lambda=\log S\).
\begin{enumerate}
\item Explain why no one-logarithm ansatz can represent \(\rev F\), however
many terms it carries.
\item Derive the closed form \(\rev F(Y)=Y/\LambertW(Y)\) and expand it through
\(S^{-4}\).
\end{enumerate}
\end{exercise}

\begin{proof}[Solution]
(1)  \(F=cX^{\rho}L^{\sigma}(1+U)\) with \(c=1\), \(\rho=1\), \(\sigma=1\),
\(U=0\).  \Cref{plt:thm:lag-rigidity} says a monomial-leading inverse exists in
a Puiseux--log ring if and only if \(\rho^{-1}\) is in the exponent group
\emph{and} \(\sigma=0\); here \(\sigma=1\), so no enlargement of the exponent
group or of the coefficient field helps.  Quantitatively,
\cref{plt:prop:lag-loglog-forced} gives
\(\log\rev F(Y)=S-\Lambda+\BigOAt{\Lambda/S}{S\to+\infty}\), and the
coefficient \(-1\) of \(\Lambda=\log\log Y\) is nonzero, so \(\rev F\) lies in
no ring whose coefficients are polynomials in one logarithm.

(2)  Put \(u=\log X\).  Then \(Y=X\log X=u\EulerE^{u}\), so \(u=\LambertW(Y)\)
on the principal branch and
\[
 \rev F(Y)=X=\EulerE^{u}=\frac{Y}{u}=\frac{Y}{\LambertW(Y)} .
\]
Now \(\LambertW(Y)=\WrightOmega(S)\), because \(w+\log w=S\) is
\(w\EulerE^{w}=Y\).  By \cref{plt:thm:om-expansion},
\[
 \LambertW(Y)=S-\Lambda+\sum_{n\ge1}\frac{\LamP{n}(\Lambda)}{S^{n}}
 =S\bigl(1+\varepsilon\bigr),
 \qquad
 \varepsilon=-\frac{\Lambda}{S}+\sum_{n\ge1}\frac{\LamP{n}(\Lambda)}{S^{n+1}} .
\]
Hence \(\rev F(Y)=\dfrac{Y}{S}\,\dfrac{1}{1+\varepsilon}\), and expanding the
reciprocal with
\(\LamP{1}=\Lambda\), \(\LamP{2}=\tfrac12\Lambda^{2}-\Lambda\),
\(\LamP{3}=\tfrac13\Lambda^{3}-\tfrac32\Lambda^{2}+\Lambda\) gives, collecting
by powers of \(S^{-1}\),
\begin{align}
\label{plt:eq:exr-xlogx}
 \rev F(Y)=\frac YS\Bigl[\,
 &1+\frac{\Lambda}{S}
  +\frac{\Lambda(\Lambda-1)}{S^{2}}
  +\frac{\Lambda(\Lambda-2)(2\Lambda-1)}{2S^{3}}\notag\\
 &+\frac{\Lambda\bigl(6\Lambda^{3}-26\Lambda^{2}+27\Lambda-6\bigr)}{6S^{4}}
  +\BigOAt{\tfrac{\Lambda^{5}}{S^{5}}}{Y\to+\infty}\Bigr].
\end{align}
For the record, the coefficients arise as follows.  Writing
\(\varepsilon=\sum_{k\ge1}e_kS^{-k}\) with \(e_1=-\Lambda\),
\(e_2=\Lambda\), \(e_3=\LamP{2}(\Lambda)\), \(e_4=\LamP{3}(\Lambda)\), the
coefficient of \(S^{-4}\) in \((1+\varepsilon)^{-1}\) is
\[
 -e_4+\bigl(2e_1e_3+e_2^{2}\bigr)-3e_1^{2}e_2+e_1^{4}
 =-\tfrac13\Lambda^{3}+\tfrac32\Lambda^{2}-\Lambda
   -\Lambda^{3}+2\Lambda^{2}+\Lambda^{2}-3\Lambda^{3}+\Lambda^{4},
\]
which is \(\Lambda^{4}-\tfrac{13}3\Lambda^{3}+\tfrac92\Lambda^{2}-\Lambda\),
the printed \(S^{-4}\) numerator divided by \(6\).  Every coefficient is a
polynomial in \(\Lambda\), which is (1) again.  The body records this
expansion through \(S^{-2}\) only, at \eqref{plt:eq:lag-xlogx-expansion}; the
two further blocks are computed here.

% ed.: the O-term in \eqref{plt:eq:exr-xlogx} is an ANALYTIC estimate, and the
% ed.: computation above is formal.  What licenses it is the explicit
% ed.: remainder bound of the Lambert part, not \cref{plt:thm:om-expansion};
% ed.: the citation is supplied rather than left to the reader, in line with
% ed.: \cref{plt:exer:exr-trap-formal-to-analytic}.
One word on the remainder.  Everything above is an identity between formal
series in \(S^{-1}\) with coefficients in \(\K[\Lambda]\); by itself that
justifies only a \(\FormalBigOAt{S^{-5}}{S^{-1}}\).  The analytic
\(\BigOAt{\Lambda^{5}S^{-5}}{Y\to+\infty}\) printed in
\eqref{plt:eq:exr-xlogx} is legitimate, but it comes from
\cref{plt:cor:lw-canonical}\,(ii): with \(\xi_0=S\) and \(\ell_0=\Lambda\)
that corollary bounds the error of the \(N\)-block truncation of
\(\LambertW(Y)\) by \(2\bigl(5(1+\Lambda)/S\bigr)^{N+1}\) once
\(S\ge10(1+\Lambda)\), and dividing \(Y\) by a quantity bounded away from
\(0\) preserves the order.  Without that input the display is a formal
identity and nothing more.
\end{proof}

\begin{exercise}[Trap: a scale change applied to one generator only]
\label{plt:exer:exr-trap-scale}
Let \(F(X)=2X+\log X\).  A computation runs: ``\(F=m_2\compose G\) with
\(m_2(Z)=2Z\) and \(G(X)=X+\tfrac12\log X\), so \(\rev F(Y)=\rev G(Y/2)\).  By
\cref{plt:cor:lag-poly-perturbation} with \(a=\tfrac12\),
\(\rev G(Z)=Z+\sum_{N\ge0}2^{-(N+1)}\LamP{N}(\log Z)Z^{-N}\).  Substituting
\(Z=Y/2\), that is \(Z^{-N}=2^{N}Y^{-N}\), gives
\[
 \rev F(Y)\ \overset{?}{=}\ \frac Y2
  +\frac12\sum_{N\ge0}\LamP{N}(\log Y)\,Y^{-N}.
\]''
Find the error, give the correct answer, quantify the discrepancy, and say
which of the volume's checks would have caught it.
\end{exercise}

\begin{proof}[Solution]
The error is that \(Z=Y/2\) changes \emph{both} generators, not one:
\[
 Z^{-N}=2^{N}Y^{-N}
 \qquad\text{and}\qquad
 \log Z=\log Y-\log2 .
\]
The computation substituted the first and forgot the second.  The correct
result is
\[
 \rev F(Y)=\frac Y2+\frac12\sum_{N\ge0}\LamP{N}\bigl(\log Y-\log2\bigr)\,Y^{-N},
\]
which is \eqref{plt:eq:lag-scale-correct}.

The failure is silent.  The wrong answer is a perfectly well formed element of
\(\PLlau\): every block has the right outer order, the right \(L\)-degree, the
right leading coefficient and the right sign pattern.  No check internal to
the near-identity calculation for \(G\) can see it, because the calculation
for \(G\) was correct.

\emph{Size of the error.}  The leading discrepancy is in the \(N=0\) block:
\(\tfrac12\LamP{0}(\log Y-\log2)-\tfrac12\LamP{0}(\log Y)
 =-\tfrac12(\log Y-\log2)+\tfrac12\log Y=\tfrac12\log2=0.34657\ldots\),
an additive constant.  Numerically, at
\(Y=20+\log10=22.3025850929\ldots\), where \(\rev F(Y)=10\) exactly, the two
truncations through \(Y^{-3}\) --- that is, over \(0\le N\le3\) --- give
\[
 \text{correct}:\ 10.0000043732\ldots,
 \qquad
 \text{naive}:\ 9.6702070480\ldots,
\]
a gap of \(0.3298\ldots\), close to but not equal to \(\tfrac12\log2\) because
the higher blocks differ as well.

\emph{What catches it.}  Two things, and only two.  First, the residual
certificate of \cref{plt:cor:rev-finite-certificate} \emph{computed against
\(F\) itself}: substituting the wrong series into \(2X+\log X\) leaves a
residual of outer order \(0\), namely \(\log 2\), rather than the
\(\FormalBigOAt{t^{N+1}}{t}\) a correct truncation must leave.  A certificate
run against the normalised map \(G\) would pass, which is why step~7 of
\cref{plt:alg:lag-dominant-balance} certifies against the original map.
Second, the coefficient field: the correct answer does not lie in
\(\PLlau\) over \(\K=\RationalNumbers\) at all, since its blocks involve
\(\log2\).  A routine that tracks the field would have refused the
substitution before performing it, which is the point of
\cref{plt:rmk:lag-nonscalar-leading}.
\end{proof}

% ed.: the two numbers in \cref{plt:warn:lag-scale-trap} were computed at two
% ed.: different cutoffs.  Recomputed in exact arithmetic at Y = 20 + log 10:
% ed.:   N <= 2  ->  correct 10.0000781540, naive 9.6702691781
% ed.:   N <= 3  ->  correct 10.0000043732, naive 9.6702070480
% ed.: The warningbox says "through Y^{-3}" and then quotes 10.000004 (which
% ed.: is the N <= 3 value) alongside 9.67026 (which is the N <= 2 value).
% ed.: Both truncations are quoted at N <= 3 above.  The conclusion of the
% ed.: warningbox is unaffected; only the second numeral changes.

\section{Wright omega and the Lambert polynomials}
\label{plt:sec:exercises-omega}

\begin{exercise}[Generating the Lambert polynomials two ways]
\label{plt:exer:exr-lambert-recurrence}
Starting from \(\LamP{1}(u)=u\), use the integral recurrence
\eqref{plt:eq:lambert-int-rec},
\(\LamP{n+1}=-\LamP{n}+n\int_0^u\LamP{n}(v)\Differential v\), to produce
\(\LamP{2},\LamP{3},\LamP{4}\).  Check each against the Stirling formula
\eqref{plt:eq:lambert-stirling}, using
\(\stirlingone{4}{4}=1\), \(\stirlingone{4}{3}=6\),
\(\stirlingone{4}{2}=11\), \(\stirlingone{4}{1}=6\).
\end{exercise}

\begin{proof}[Solution]
Successive integration, the constant of integration being fixed to \(0\) by
\(\LamP{n}(0)=0\) (\cref{plt:thm:lambert-recurrence}):
\begin{align*}
 \LamP{2}&=-u+1\cdot\int_0^u v\,\Differential v=\frac{u^{2}}2-u,\\
 \LamP{3}&=-\Bigl(\frac{u^{2}}2-u\Bigr)
   +2\int_0^u\Bigl(\frac{v^{2}}2-v\Bigr)\Differential v
   =-\frac{u^{2}}2+u+\frac{u^{3}}3-u^{2}
   =\frac{u^{3}}3-\frac32u^{2}+u,\\
 \LamP{4}&=-\Bigl(\frac{u^{3}}3-\frac32u^{2}+u\Bigr)
   +3\int_0^u\Bigl(\frac{v^{3}}3-\frac32v^{2}+v\Bigr)\Differential v\\
 &=-\frac{u^{3}}3+\frac32u^{2}-u
   +\frac{u^{4}}4-\frac32u^{3}+\frac32u^{2}
   =\frac{u^{4}}4-\frac{11}6u^{3}+3u^{2}-u .
\end{align*}

\emph{Stirling check at \(n=4\).}  Formula
\eqref{plt:eq:lambert-stirling} reads
\(\LamP{n}(u)=\sum_{m=1}^{n}(-1)^{n-m}\stirlingone{n}{n-m+1}u^{m}/m!\).  For
\(n=4\):
\[
 -\stirlingone{4}{4}\frac{u}{1!}
 +\stirlingone{4}{3}\frac{u^{2}}{2!}
 -\stirlingone{4}{2}\frac{u^{3}}{3!}
 +\stirlingone{4}{1}\frac{u^{4}}{4!}
 =-u+3u^{2}-\frac{11}6u^{3}+\frac{u^{4}}4,
\]
which agrees.  Note that the alternation is governed by \((-1)^{n-m}\), not by
\((-1)^{m}\): the whole row changes sign with the parity of \(n\).  A common
slip is to fix the sign pattern from a single row --- \(n=4\), where the
coefficients run \(-,+,-,+\) --- and then apply it at \(n=3\), where they run
\(+,-,+\).

\emph{Factored forms.}
\(\LamP{2}=\tfrac12u(u-2)\),
\(\LamP{3}=\tfrac16u(2u^{2}-9u+6)\),
\(\LamP{4}=\tfrac1{12}u(3u^{3}-22u^{2}+36u-12)\), matching the scaled integer
family \(\LamPsc{n}=n!\LamP{n}\) of \cref{plt:tab:lambert-polynomials}.
\end{proof}

\begin{exercise}[The one-parameter family, from Lagrange--B\"urmann]
\label{plt:exer:exr-alpha-family}
Let \(a\in\K\) and \(F=X+a\log X\).  Prove directly from
\eqref{plt:eq:lag-burmann-pl} that the coefficient of \(X^{-n}\) in
\(\rev F\) is \(a^{\,n+1}\LamP{n}(L)\) for every \(n\ge0\), and identify which
Lagrange index contributes to which block.
\end{exercise}

\begin{proof}[Solution]
Apply \eqref{plt:eq:lag-burmann-pl} with \(A=aL\in\PLpow\), legitimate since
\(\ordt(A)=0>-1\) (\cref{plt:prop:lag-summability-sharp}):
\[
 \rev F=X+\sum_{r\ge1}\frac{(-1)^{r}}{r!}\,\DX^{\,r-1}\bigl(a^{r}L^{r}\bigr)
      =X+\sum_{r\ge1}\frac{(-1)^{r}a^{r}}{r!}\,\DX^{\,r-1}\bigl(L^{r}\bigr).
\]
By \cref{plt:thm:dif-repeated} applied to the \(t^{0}\)-block \(L^{r}\),
\[
 \DX^{\,r-1}\bigl(L^{r}\bigr)=t^{\,r-1}\,\shiftop{0}{r-1}L^{r}
 =t^{\,r-1}\Bigl[\prod_{j=0}^{r-2}(\partial_L-j)\Bigr]L^{r}.
\]
So the \(r\)-th summand is a \emph{single} block, sitting at outer order
exactly \(r-1\).  Setting \(n=r-1\), no two summands ever meet, and
\[
 [t^{n}]\rev F
 =a^{\,n+1}\cdot\frac{(-1)^{n+1}}{(n+1)!}
   \Bigl[\prod_{j=0}^{n-1}(\partial_L-j)\Bigr]L^{\,n+1}
 =a^{\,n+1}\LamP{n}(L),
\]
by \cref{plt:def:lambert-polynomials}.  Index \(r\) contributes to block
\(n=r-1\) and to no other; this is the ``one Lagrange index per block'' of
\eqref{plt:eq:om-one-block}, and it is special to a perturbation that is a
\emph{single} \(t^{0}\)-block.  At \(a=1\) one recovers
\(\WrightOmega\) (\cref{plt:thm:om-expansion}) and at \(a=-1\) the reversion of
\(X-\log X\) with weights \((-1)^{n+1}\)
(\eqref{plt:eq:om-minus-log}); at \(n=0\) the value is
\(a\LamP{0}(L)=-aL\), which is the leading correction predicted by
\cref{plt:thm:rev-existence}\,(4).

A useful sanity check that does not use the formula: substituting
\(t\mapsto at\) in the implicit equation
\eqref{plt:eq:om-implicit} and multiplying by \(a\) transports the unique
fixed point of \(X+L\) to that of \(X+aL\), which is the proof of
\cref{plt:thm:om-alpha-family}.
\end{proof}

\begin{exercise}[Counterexample: the degree law and the sign of the error at \(n=0\)]
\label{plt:exer:exr-degree-law-zero}
The Lambert polynomials satisfy \(\deg\LamP{n}=n\), \([u^{n}]\LamP{n}=1/n\)
and \([u]\LamP{n}=(-1)^{n-1}\) --- for \(n\ge1\).
\begin{enumerate}
\item Show that each of the three statements fails, or is meaningless, at
\(n=0\), and determine which parts of
\cref{plt:cor:lambert-coefficients} do survive there.
\item A plausible consequence of \(\WrightOmega-\WrightOmega_N\sim
\LamP{N+1}(L)x^{-(N+1)}\) is that every truncation approaches
\(\WrightOmega\) from below on \((1,\infty)\).  Refute this, and locate the
exact threshold for \(N=1\).
\end{enumerate}
\end{exercise}

\begin{proof}[Solution]
(1)  \(\LamP{0}(u)=-u\).  Its degree is \(1\), not \(0\); the assertion
\([u^{0}]\LamP{0}=1/0\) is meaningless; and \([u]\LamP{0}=-1\) whereas
\((-1)^{0-1}=-1\) --- so the third statement happens to hold, by accident, and
% ed.: "exactly item (i)" contradicted the sentence just above, which already
% ed.: observed that (v) also evaluates correctly at n = 0.
should not be quoted as evidence.  What survives at \(n=0\) is the
part of \cref{plt:cor:lambert-coefficients} that does not mention the degree
or a coefficient indexed from the top: the constant term vanishes and
\(u\mid\LamP{0}\), i.e.\ item~(i) minus its degree clause.  The two
bottom-indexed items also evaluate correctly there --- (v) as noted, and (vii)
as the identity \(0=0\) --- but by accident, each being proved only for
\(n\ge1\) and \(n\ge2\) respectively.  Every statement
indexed relative to the top --- (ii) on the number of nonzero coefficients,
(iii), (iv), (vi) --- fails or is vacuous.  The integrality normalisation
\(\LamPsc{n}=n!\LamP{n}\) is likewise an \(n\ge1\) statement; at \(n=0\) it
returns \(-u\), which is the value, not an instance of the law.

Why the exception is unavoidable: the zero-based normalisation is the one in
which \eqref{plt:eq:om-one-block} says ``one Lagrange index per block'' and in
which \cref{plt:thm:om-alpha-family} needs no case distinction.  The price is
paid entirely at \(n=0\), and it is the reason
\cref{plt:rmk:om-normalisation} exists.

(2)  The claim would follow if \(\LamP{N+1}(L)>0\) for all \(x>1\), i.e.\ for
all \(L>0\).  It is false.  For \(N=1\) the relevant polynomial is
\(\LamP{2}(u)=\tfrac12u(u-2)\), which is \emph{negative} for \(0<u<2\) and
positive for \(u>2\).  So the leading term \(\LamP{2}(L)x^{-2}\) of the error
changes sign at \(L=2\), that is at \(x=\EulerE^{2}=7.389056\ldots\), and the
truncation \(\WrightOmega_1\) overshoots on one side of that point and
undershoots on the other.

% ed.: the inherited solution asserted that the sign of omega - Omega_1 itself
% ed.: changes at x = e^2.  That is the sign change of the LEADING TERM only.
% ed.: The true difference crosses zero at x = 8.7465..., recomputed here; the
% ed.: gap between the two numbers is the whole moral of the exercise, so
% ed.: printing e^2 as the exact threshold defeats it.
The exact threshold is a different number, and this is precisely the trap the
exercise is about.  The function
\(x\mapsto\WrightOmega(x)-\WrightOmega_1(x)\) is negative on \((1,x_\ast)\)
and positive on \((x_\ast,\infty)\), with a single zero at
\[
 x_\ast=8.746508640\ldots,\qquad \log x_\ast=2.168654608\ldots,
\]
not at \(\EulerE^{2}\).  The leading term locates the crossing only to the
accuracy with which it approximates the error, and \(\LamP{3}(L)x^{-3}\) is
still of comparable size at these values of \(x\).  What the body's
\eqref{plt:eq:om-sharp-order} asserts is an equivalence as \(x\to+\infty\); it
does not name a finite crossing point, and nothing here entitles one to read
\(L=2\) as such a point.

Numerically: at \(x=\EulerE\) (so \(L=1\)) one has
\(\WrightOmega(\EulerE)=2.0167798\ldots\) while
\(\WrightOmega_1(\EulerE)=\EulerE-1+\EulerE^{-1}=2.0861613\ldots\), an
overshoot; at \(x=10\) (so \(L=2.302585\ldots\)) one has
\(\WrightOmega(10)=7.9294201\ldots\) and
\(\WrightOmega_1(10)=7.9276734\ldots\), an undershoot.  The two errors are
\[
 -6.938\times10^{-2}
 \qquad\text{and}\qquad
 +1.747\times10^{-3},
\]
% ed.: the inherited solution claimed both errors agree "to one digit" with
% ed.: LamP_2(L)x^{-2}.  At x = 10 the leading term is 3.484e-3 against a true
% ed.: error of 1.747e-3 -- a factor of two, not one digit.  Only the signs
% ed.: agree at that point, and the reason is printed below.
against the leading term \(\LamP{2}(L)x^{-2}\), which is
\(-6.767\times10^{-2}\) and \(+3.484\times10^{-3}\) respectively.  The signs
agree in both cases, which is all the sign discussion needs.  The magnitudes
agree only at \(x=\EulerE\): at \(x=10\) the leading term is too large by a
factor of about two, because the next two blocks,
\(\LamP{3}(L)x^{-3}=-1.581\times10^{-3}\) and
\(\LamP{4}(L)x^{-4}=-1.752\times10^{-4}\), cancel most of it.  That is again
the same lesson: at \(L\approx2.3\) the expansion is nowhere near its
asymptotic regime, and \eqref{plt:eq:om-sharp-order} claims nothing at a
finite \(x\).

More generally \([u]\LamP{n}=(-1)^{n-1}\), so near \(u=0^{+}\) --- that is for
\(x\) slightly above~\(1\) --- the sign of \(\LamP{n}(L)\) is \((-1)^{n-1}\)
and the truncations alternate sides.  Only for \(L\) large, where the leading
term \(L^{n}/n>0\) dominates, do they settle on one side.  The correct
statement is the one the body makes,
\eqref{plt:eq:om-sharp-order}: an asymptotic equivalence as \(x\to+\infty\),
with no claim at any finite \(x\).
\end{proof}

\begin{exercise}[The lower real branch near zero]
\label{plt:exer:exr-lower-branch}
Let \(x\in(-\EulerE^{-1},0)\) and write \(x=-\EulerE^{-\eta}\), so
\(\eta\to+\infty\) as \(x\to0^{-}\).  Put \(M=\log\eta\).
\begin{enumerate}
\item Show that \(\LambertW_{-1}(x)=-v\) where \(v\) solves \(v-\log v=\eta\).
\item Derive the first four correction blocks of \(v\), and hence of
\(\LambertW_{-1}\).
\item Some sources say ``the same \(\LamP{n}\) occur, multiplied by
\((-1)^{n}\)''.  Say precisely which quantity carries which weight.
\end{enumerate}
\end{exercise}

\begin{proof}[Solution]
(1)  \(\LambertW_{-1}(x)<-1\), so write \(\LambertW_{-1}(x)=-v\) with \(v>1\).
The defining equation \(\LambertW\EulerE^{\LambertW}=x\) becomes
\(-v\EulerE^{-v}=x=-\EulerE^{-\eta}\), i.e.\ \(v\EulerE^{-v}=\EulerE^{-\eta}\).
Taking the real logarithm of both positive sides,
\(\log v-v=-\eta\), that is \(v-\log v=\eta\).

(2)  So \(v=\rev{(X-\log X)}(\eta)\).  By \eqref{plt:eq:om-minus-log},
\begin{align*}
 v&=\eta+M+\sum_{n\ge1}(-1)^{n+1}\LamP{n}(M)\,\eta^{-n}\\
 &=\eta+M+\frac{M}{\eta}-\frac{\tfrac12M^{2}-M}{\eta^{2}}
  +\frac{\tfrac13M^{3}-\tfrac32M^{2}+M}{\eta^{3}}
  -\frac{\LamP{4}(M)}{\eta^{4}}+\cdots,
\end{align*}
and therefore
\begin{align*}
 \LambertW_{-1}\bigl(-\EulerE^{-\eta}\bigr)
 ={}&-\eta-M-\frac{M}{\eta}+\frac{M(M-2)}{2\eta^{2}}
  -\frac{M(2M^{2}-9M+6)}{6\eta^{3}}\\
 &+\frac{M(3M^{3}-22M^{2}+36M-12)}{12\eta^{4}}
  +\BigOAt{\tfrac{M^{5}}{\eta^{5}}}{\eta\to+\infty},
\end{align*}
which is \eqref{plt:eq:lw-lower-first-terms}.  Note that the Lambert part of
this volume proves more than an asymptotic statement here: by
\cref{plt:thm:lw-lower-branch} the series
\(\sum_{n\ge1}(-1)^{n}\LamP{n}(M)\eta^{-n}\) \emph{converges} once
\(2(1+\AbsoluteValue{M})<\eta\), with an explicit remainder bound.  That is a
theorem about this particular family, proved by a Cauchy estimate; it is not
a consequence of the formal reversion.
% ed.: the inherited text added "and nothing analogous is claimed for
% ed.: WrightOmega".  That is false: the same lemma
% ed.: (plt:lem:lw-two-variable) applied at sigma = +1/x rather than -1/S
% ed.: gives the omega statement, and it is printed in the volume as
% ed.: plt:thm:lw-basepoint / plt:cor:lw-canonical.  The claim is corrected
% ed.: here and in plt:exer:exr-omega-numerics, which repeated it.
The corresponding statement for \(\WrightOmega\) is \emph{not} missing from
the volume, contrary to what one might expect from the asymptotic framing of
\cref{plt:sec:wright-omega}: \cref{plt:thm:lw-basepoint} applied with
\(\xi=x\), \(K=x\) and \(\log K=L\) gives \(M=L\), so its seating condition
reads \(2(1+L)<x\), and under it the series
\(\sum_{n\ge1}\LamP{n}(L)x^{-n}\) converges absolutely to
\(\WrightOmega(x)-x+L\).  The two branches are the two signs
\(\sigma=\pm1/\,\cdot\,\) of one and the same Cauchy estimate,
\cref{plt:lem:lw-two-variable}.

(3)  The two weights belong to two different objects, and conflating them is
the error the sources make.  The reversion of \(X-\log X\) carries
\((-1)^{n+1}\):
\(\rev{(X-L)}=X+L+\sum_{n\ge1}(-1)^{n+1}\LamP{n}(L)t^{n}\), so its first
correction is \(+\LamP{1}(L)/X=+L/X\), with a \emph{plus} sign.  Its negative,
which is what enters \(\LambertW_{-1}\), carries \((-1)^{n}\):
\(\LambertW_{-1}=-\eta-M+\sum_{n\ge1}(-1)^{n}\LamP{n}(M)\eta^{-n}\).  Both
displays above are consistent with this.
\end{proof}

% ed.: source 1's prose ("the same P_n occur, multiplied by (-1)^n")
% ed.: contradicts its own displayed expansion of the inverse of X - log X, in
% ed.: which P_1 = L occurs with a plus sign.  The weight (-1)^{n+1} belongs
% ed.: to rev(X - L); the weight (-1)^n belongs to its negative, which is what
% ed.: W_{-1} is.  Part (3) separates them.

\begin{exercise}[Trap: an off-by-one in the truncation index]
\label{plt:exer:exr-trap-omega-index}
A source using the inclusive cutoff defines
\(\Omega_N=X+[\mathcal P]_{\le N}\), where
\(\mathcal P=\sum_{n\ge0}\LamP{n}(L)t^{n}\), and proves
\[
 \Omega_N+\log\Omega_N-X=-\LamP{N+1}(L)\,t^{\,N+1}
 +\FormalBigOAt{t^{\,N+2}}{t}.
\]
A reader transcribing into the normative exclusive convention writes
\(\Omega_N=X+\Trunc{\mathcal P}{N}\) and keeps the conclusion.  Show that the
conclusion is then false, compute the correct residual, and give the correct
transcription.
\end{exercise}

\begin{proof}[Solution]
The two cutoffs are related by \([F]_{\le N}=\Trunc{F}{N+1}\), so the correct
transcription is
\(\WrightOmega_N=X+\Trunc{\mathcal P}{N+1}\), which is
\cref{plt:def:om-truncation}.  Writing \(\Trunc{\mathcal P}{N}\) instead drops
the block \(\LamP{N}(L)t^{N}\).

Compute the residual of the mistranscribed object.  Put
\(\mathcal P_{N}'=\Trunc{\mathcal P}{N}\) and
\(\Delta=\mathcal P-\mathcal P_{N}'=\sum_{n\ge N}\LamP{n}(L)t^{n}\), so
\(\ordt\Delta=N\) and \([t^{N}]\Delta=\LamP{N}(L)\).  Exactly the computation
in the proof of \cref{plt:thm:om-residual} --- generator substitution, then
\(1+t\mathcal P_{N}'=(1+t\mathcal P)\bigl(1-t\Delta/(1+t\mathcal P)\bigr)\)
and additivity of the formal logarithm --- gives
\[
 \bigl(X+\Trunc{\mathcal P}{N}\bigr)
 +\log\bigl(X+\Trunc{\mathcal P}{N}\bigr)-X
 =-\LamP{N}(L)\,t^{N}+\FormalBigOAt{t^{\,N+1}}{t}.
\]
So the residual has outer order exactly \(N\), not \(N+1\), and its leading
block is \(-\LamP{N}\), not \(-\LamP{N+1}\).  For \(N=3\), for instance, the
mistranscribed truncation leaves
\(-\bigl(\tfrac13L^{3}-\tfrac32L^{2}+L\bigr)t^{3}+\FormalBigOAt{t^{4}}{t}\),
whereas the correct \(\WrightOmega_3\) leaves
\(-\LamP{4}(L)t^{4}+\FormalBigOAt{t^{5}}{t}\).

The practical bite: a user asking for ``the expansion correct through
\(x^{-3}\)'' and receiving \(X+\Trunc{\mathcal P}{3}\) gets an answer correct
only through \(x^{-2}\).  The error is one block, it is systematic, and it is
undetectable by inspecting the answer --- every printed block is a correct
Lambert polynomial.  What detects it is the certificate of
\cref{plt:cor:om-certificate}: the residual order is the definition of how
many blocks are right, and it reads \(N\), not \(N+1\).
\end{proof}

\begin{exercise}[The Laplace transform and its cancellations]
\label{plt:exer:exr-laplace}
\begin{enumerate}
\item Compute the transform
\[
 \int_0^{\infty}\EulerE^{-su}\,\LamP{n}(u)\,\Differential u
\]
directly for \(n=1,2,3\) and \(\Re s>0\), and compare with
\eqref{plt:eq:om-laplace}.
\item Deduce the exponential moment cancellations
\[
 \int_0^{\infty}\EulerE^{-ru}\LamP{n}(u)\Differential u=0
\]
for all integers \(r\) with \(1\le r\le n-1\), and say why the range is
empty at \(n=1\).
\item Explain why the identity says nothing about the convergence of
\(\sum_n\LamP{n}(u)s^{n}\).
\end{enumerate}
\end{exercise}

\begin{proof}[Solution]
(1)  Use \(\int_0^{\infty}\EulerE^{-su}u^{m}\Differential u=m!/s^{m+1}\) for
\(\Re s>0\).
\[
 \LamP{1}=u:\quad \frac1{s^{2}};
 \qquad
 \LamP{2}=\frac{u^{2}}2-u:\quad \frac12\cdot\frac{2}{s^{3}}-\frac1{s^{2}}
 =\frac{1-s}{s^{3}};
\]
\[
 \LamP{3}=\frac{u^{3}}3-\frac32u^{2}+u:\quad
 \frac13\cdot\frac6{s^{4}}-\frac32\cdot\frac2{s^{3}}+\frac1{s^{2}}
 =\frac{2-3s+s^{2}}{s^{4}}=\frac{(1-s)(2-s)}{s^{4}} .
\]
\Cref{plt:thm:om-laplace} predicts
\(s^{-(n+1)}\prod_{j=1}^{n-1}(j-s)\), the product empty at \(n=1\); the three
values are \(1/s^{2}\), \((1-s)/s^{3}\), \((1-s)(2-s)/s^{4}\), matching.

(2)  The right-hand side of \eqref{plt:eq:om-laplace} vanishes exactly at the
zeros of \(\prod_{j=1}^{n-1}(j-s)\), i.e.\ at \(s=1,2,\dots,n-1\), all of
which satisfy \(\Re s>0\) so that the integral converges absolutely.  Hence
\(\int_0^{\infty}\EulerE^{-ru}\LamP{n}(u)\Differential u=0\) for
\(1\le r\le n-1\).  At \(n=1\) the product is empty and has no zeros, so the
range \(1\le r\le0\) is empty; the first genuine cancellation is
\(\int_0^{\infty}\EulerE^{-u}\LamP{2}(u)\Differential u=0\), which one can
check by hand: \(\tfrac12\cdot2-1=0\).

(3)  The Laplace identity is a statement about one polynomial at a time.  It
integrates \(\LamP{n}\) against a fixed kernel over \(u\in(0,\infty)\); it
involves no sum over \(n\), and gives no bound on \(\LamP{n}(u)\) at a fixed
\(u\) as \(n\to\infty\).  Convergence of \(\sum_n\LamP{n}(u)s^{n}\) at frozen
\(u\) is a separate theorem, \cref{plt:prop:om-frozen-convergence}, proved by
the implicit function theorem for \eqref{plt:eq:om-frozen}; and neither
statement is about the asymptotic expansion of \(\WrightOmega\), where
\(u=\log x\) and \(s=1/x\) move together.  Keeping the three apart is the
discipline \cref{plt:sec:exercises-analytic} is about.
\end{proof}

\begin{exercise}[Reciprocal and logarithm of \(\LambertW\) (hard)]
\label{plt:exer:exr-logW-and-recip}
Let \(z\) be large and positive, \(L_1=\log z\), \(\ell=\log L_1\),
\(r=L_1^{-1}\), and \(W=\LambertW(z)\) on the principal branch.
\begin{enumerate}
\item Derive \(\log W=L_1-W\) exactly, and use it to expand \(\log W\) through
\(r^{4}\) without ever taking the logarithm of a series.
\item Expand \(1/W\) through \(r^{6}\), and compute one block further than the
sources tabulate.
\item Deduce a transseries for \(\EulerE^{W}\) with no exponentiation of a
series.
\end{enumerate}
\end{exercise}

\begin{proof}[Solution]
(1)  From \(W\EulerE^{W}=z\), taking logarithms of the positive quantities,
\(\log W+W=\log z=L_1\), so \(\log W=L_1-W\) exactly.  Since
\(W=\WrightOmega(L_1)\) --- because \(W+\log W=L_1\) is the defining equation
of \(\WrightOmega\) --- \cref{plt:thm:om-expansion}, transported to branch
coordinates by \cref{plt:thm:lw-formal-expansion}, gives
\(W=L_1-\ell+\sum_{n\ge1}\LamP{n}(\ell)r^{n}\).  Subtracting from \(L_1\),
\begin{align*}
 \log W&=\ell-\sum_{n\ge1}\LamP{n}(\ell)r^{n}\\
 &=\ell-\ell r+\Bigl(\ell-\frac{\ell^{2}}2\Bigr)r^{2}
  +\Bigl(-\frac{\ell^{3}}3+\frac32\ell^{2}-\ell\Bigr)r^{3}\\
 &\qquad+\Bigl(-\frac{\ell^{4}}4+\frac{11}6\ell^{3}-3\ell^{2}+\ell\Bigr)r^{4}
  +\FormalBigOAt{r^{5}}{r}.
\end{align*}
No series logarithm was computed: the identity \(\log W=L_1-W\) did all the
work.  This is the same device as
\cref{plt:lem:dif-logderiv} used in closed form.

(2)  Factor \(W=L_1(1+\varepsilon)\) with
\begin{align*}
 \varepsilon&=-\ell r+\sum_{n\ge1}\LamP{n}(\ell)r^{n+1}\\
 &=-\ell r+\ell r^{2}+\Bigl(\frac{\ell^{2}}2-\ell\Bigr)r^{3}
  +\Bigl(\frac{\ell^{3}}3-\frac32\ell^{2}+\ell\Bigr)r^{4}
  +\LamP{4}(\ell)r^{5}+\LamP{5}(\ell)r^{6}+\cdots,
\end{align*}
so \(\ordt\varepsilon=1\) in the \(r\)-adic filtration and the geometric
series \(\sum_k(-\varepsilon)^{k}\) is summable
(\cref{plt:exer:exr-summable-powers}).  Then
\(1/W=r\sum_{k\ge0}(-\varepsilon)^{k}\).  Collecting by powers of \(r\), with
\(e_1=-\ell\), \(e_2=\ell\), \(e_3=\LamP{2}(\ell)\), \(e_4=\LamP{3}(\ell)\),
\(e_5=\LamP{4}(\ell)\), \(e_6=\LamP{5}(\ell)\) the coefficients of
\(\varepsilon=\sum_ke_kr^{k}\):
\begin{align}
\label{plt:eq:exr-recip-W}
 \frac1W={}&r+\ell r^{2}+\bigl(\ell^{2}-\ell\bigr)r^{3}
  +\Bigl(\ell^{3}-\frac52\ell^{2}+\ell\Bigr)r^{4}\notag\\
 &+\Bigl(\ell^{4}-\frac{13}3\ell^{3}+\frac92\ell^{2}-\ell\Bigr)r^{5}
  +\Bigl(\ell^{5}-\frac{77}{12}\ell^{4}+\frac{71}6\ell^{3}-7\ell^{2}
   +\ell\Bigr)r^{6}\notag\\
 &+\Bigl(\ell^{6}-\frac{87}{10}\ell^{5}+\frac{145}6\ell^{4}
   -\frac{155}6\ell^{3}+10\ell^{2}-\ell\Bigr)r^{7}
  +\FormalBigOAt{r^{8}}{r}.
\end{align}
The blocks through \(r^{6}\) are tabulated in the sources; the \(r^{7}\) block
is computed here.  As a sample of the bookkeeping, the coefficient of
\(r^{5}\) is
\(-e_4+(2e_1e_3+e_2^{2})-3e_1^{2}e_2+e_1^{4}\), which evaluates to
\(-\tfrac13\ell^{3}+\tfrac32\ell^{2}-\ell
 -\ell^{3}+2\ell^{2}+\ell^{2}-3\ell^{3}+\ell^{4}
 =\ell^{4}-\tfrac{13}3\ell^{3}+\tfrac92\ell^{2}-\ell\).

Three structural checks now follow, each provable in a line and each
visible in \eqref{plt:eq:exr-recip-W}.  Write
\[
 \frac1W=r\sum_{k\ge0}c_k(\ell)\,r^{k},
\]
so that \(c_k\) is the block of \(1/W\) at \(r^{k+1}\).

\emph{Zero constant term.}  Setting \(\ell=0\) kills every \(\LamP{n}\)
(\eqref{plt:eq:lambert-boundary}), hence \(\varepsilon|_{\ell=0}=0\) and
\(1/W|_{\ell=0}=r\); so \(c_k(0)=0\) for \(k\ge1\).

\emph{Monic of degree \(k\).}  In \(\varepsilon=\sum_{k\ge1}e_kr^{k}\) the
first coefficient has \(\ell\)-degree \(1\) and, for \(k\ge2\),
\(e_k=\LamP{k-1}(\ell)\) has \(\ell\)-degree \(k-1\).  A monomial
\(e_{i_1}\cdots e_{i_m}\) with \(i_1+\dots+i_m=k\) therefore has degree at
most \(k\) minus the number of indices \(i_j\ge2\), which equals \(k\) only
when every \(i_j=1\).  So the top-degree part of \(c_k\) comes from
\((-e_1)^{k}=\ell^{k}\) alone, and \([\ell^{k}]c_k=1\).

\emph{Linear coefficient.}  As \(\varepsilon\) vanishes at \(\ell=0\), the
derivative of \(\recip{(1+\varepsilon)}\) there is \(-\partial_\ell\varepsilon\),
and by \({\LamP{n}}'(0)=(-1)^{n-1}\)
(\cref{plt:cor:lambert-coefficients}(v)),
\[
 -\partial_\ell\varepsilon\big|_{\ell=0}
 =r+\sum_{n\ge1}(-1)^{n}r^{n+1}.
\]
Hence \([\ell]c_k=(-1)^{k-1}\) for \(k\ge1\), which is the alternation
visible in \eqref{plt:eq:exr-recip-W}.

Finally, \eqref{plt:eq:exr-recip-W} is, block for block, the bracket in
\eqref{plt:eq:exr-xlogx}: \(\rev{(X\log X)}(Y)=Y/\LambertW(Y)\), so the two
computations are the same one.

(3)  From \(W\EulerE^{W}=z\),
\(\EulerE^{W}=z/W=z\cdot(1/W)\), and
\eqref{plt:eq:exr-recip-W} supplies \(1/W\) already.  No exponential of a
series is ever formed, which matters because
\cref{plt:thm:dif-exp-boundary} shows that exponentiating a block of
\(L\)-degree \(\ge2\) leaves every power--log scale.
\end{proof}

\begin{exercise}[The fifth and sixth Lambert polynomials, three ways (hard)]
\label{plt:exer:exr-lambert-five-six}
Compute \(\LamP{5}\) and \(\LamP{6}\) from the operator definition
\eqref{plt:eq:lambert-operator}, from the integral recurrence, and from the
Stirling formula, using
\[
 \bigl(\stirlingone{5}{5},\dots,\stirlingone{5}{1}\bigr)=(1,10,35,50,24),
 \qquad
 \bigl(\stirlingone{6}{6},\dots,\stirlingone{6}{1}\bigr)
  =(1,15,85,225,274,120).
\]
Then verify the coefficient identities of
\cref{plt:cor:lambert-coefficients}(iii)--(vii) on both.
\end{exercise}

\begin{proof}[Solution]
\emph{Answers.}
\begin{align*}
 \LamP{5}(u)&=\frac{u^{5}}5-\frac{25}{12}u^{4}+\frac{35}6u^{3}-5u^{2}+u
  =\frac{24u^{5}-250u^{4}+700u^{3}-600u^{2}+120u}{120},\\
 \LamP{6}(u)&=\frac{u^{6}}6-\frac{137}{60}u^{5}+\frac{75}8u^{4}
   -\frac{85}6u^{3}+\frac{15}2u^{2}-u\\
 &=\frac{120u^{6}-1644u^{5}+6750u^{4}-10200u^{3}+5400u^{2}-720u}{720}.
\end{align*}

\emph{Stirling route.}  The two rows are the usual cycle numbers
\cite{comtet-book,graham-knuth-patashnik}.  By
\eqref{plt:eq:lambert-stirling},
\[
 \LamP{n}(u)=\sum_{m=1}^{n}(-1)^{n-m}\stirlingone{n}{n-m+1}\frac{u^{m}}{m!}.
\]
For
\(n=5\) the signs are \((-1)^{5-m}\), i.e.\ \(+,-,+,-,+\) for
\(m=1,\dots,5\):
\[
 \LamP{5}=\stirlingone{5}{5}u-\stirlingone{5}{4}\frac{u^{2}}{2}
 +\stirlingone{5}{3}\frac{u^{3}}{6}-\stirlingone{5}{2}\frac{u^{4}}{24}
 +\stirlingone{5}{1}\frac{u^{5}}{120}
 =u-5u^{2}+\frac{35}6u^{3}-\frac{25}{12}u^{4}+\frac{u^{5}}5 .
\]
For \(n=6\) the signs are \((-1)^{6-m}\), i.e.\ \(-,+,-,+,-,+\):
\[
 \LamP{6}=-u+\frac{15}2u^{2}-\frac{85}6u^{3}+\frac{225}{24}u^{4}
 -\frac{274}{120}u^{5}+\frac{120}{720}u^{6},
\]
which is the stated answer since \(225/24=75/8\), \(274/120=137/60\) and
\(120/720=1/6\).

\emph{Recurrence route.}  From \(\LamP{4}\) of
\cref{plt:exer:exr-lambert-recurrence},
\(\LamP{5}=-\LamP{4}+4\int_0^u\LamP{4}\), giving
\(-\bigl(\tfrac14u^{4}-\tfrac{11}6u^{3}+3u^{2}-u\bigr)
 +4\bigl(\tfrac1{20}u^{5}-\tfrac{11}{24}u^{4}+u^{3}-\tfrac12u^{2}\bigr)
 =\tfrac15u^{5}-\tfrac{25}{12}u^{4}+\tfrac{35}6u^{3}-5u^{2}+u\).
Then \(\LamP{6}=-\LamP{5}+5\int_0^u\LamP{5}\) gives the second answer by the
same arithmetic.

\emph{Operator route.}  \(\LamP{5}=\tfrac1{720}
(\partial_u-4)(\partial_u-3)(\partial_u-2)(\partial_u-1)\partial_u\,u^{6}\);
applying the five factors in turn to \(6u^{5}\) reproduces the same
polynomial.  This is the only one of the three routes that presupposes none of
the derived identities.

\emph{Coefficient checks.}  With \(H_4=\tfrac{25}{12}\) and
\(H_5=\tfrac{137}{60}\):
(iii) \([u^{5}]\LamP{5}=\tfrac15\) and \([u^{6}]\LamP{6}=\tfrac16\), i.e.\
\(1/n\).
(iv) \([u^{4}]\LamP{5}=-\tfrac{25}{12}=-H_4\) and
\([u^{5}]\LamP{6}=-\tfrac{137}{60}=-H_5\).
(v) \([u]\LamP{5}=1=(-1)^{4}\) and \([u]\LamP{6}=-1=(-1)^{5}\).
(vii) \([u^{2}]\LamP{5}=-5=(-1)^{5}\tfrac{5\cdot4}{4}\) and
\([u^{2}]\LamP{6}=\tfrac{15}2=(-1)^{6}\tfrac{6\cdot5}{4}\).
% ed.: item (vi) of \cref{plt:cor:lambert-coefficients} is indexed from the
% ed.: top, [u^{n-2}], not at the fixed exponent 3.  The two coincide at n = 5
% ed.: and only there; read as "[u^3]" the check would fail at n = 6, where
% ed.: [u^3]LamP_6 = -85/6 while the right-hand side is 75/8 = [u^4]LamP_6.
% ed.: The inherited solution printed the fixed exponent and verified (vi) on
% ed.: LamP_5 only, although the exercise asks for both.
Finally (vi), which reads \([u^{\,n-2}]\LamP{n}
 =\tfrac{n-1}2\bigl(H_{n-1}^{2}-H_{n-1}^{(2)}\bigr)\) --- an index counted
from the top, not the fixed exponent \(3\), the two agreeing at \(n=5\) alone.
At \(n=5\), with \(H_4^{(2)}=\tfrac{205}{144}\),
\[
 [u^{3}]\LamP{5}=\tfrac{35}6
 =2\Bigl(\tfrac{625}{144}-\tfrac{205}{144}\Bigr)
 =2\cdot\tfrac{420}{144}=\tfrac{35}6 ;
\]
at \(n=6\), with \(H_5^{(2)}=\tfrac{5269}{3600}\),
\[
 [u^{4}]\LamP{6}=\tfrac{75}8
 =\tfrac52\Bigl(\tfrac{18769}{3600}-\tfrac{5269}{3600}\Bigr)
 =\tfrac52\cdot\tfrac{15}4=\tfrac{75}8 .
\]
Both hold.  Reading (vi) at the fixed exponent \(3\) instead would compare
\([u^{3}]\LamP{6}=-\tfrac{85}6\) with \(\tfrac{75}8\) and report a failure
that is entirely an indexing error.
Multiplying by \(n!\) turns every coefficient into an integer, as
\cref{plt:cor:lambert-integrality} requires: \(\LamPsc{5}\) and
\(\LamPsc{6}\) are the numerators displayed above.
\end{proof}

\begin{exercise}[The branch point forces a square root, not a logarithm (hard)]
\label{plt:exer:exr-branch-point}
At \(z=-\EulerE^{-1}\) the two real branches of \(\LambertW\) meet at
\(\LambertW=-1\).
\begin{enumerate}
\item Set \(\LambertW=-1+q\) and expand \(1+\EulerE z\) in powers of \(q\)
through \(q^{4}\).
\item With \(p=\sqrt{2(1+\EulerE z)}\), derive
\(\LambertW_0(z)=-1+p-\tfrac13p^{2}+\tfrac{11}{72}p^{3}
 -\tfrac{43}{540}p^{4}+\cdots\) and \(\LambertW_{-1}\) by \(p\mapsto-p\).
\item Explain, from the local structure of \(w\EulerE^{w}\), why a square root
rather than an iterated logarithm is forced here, and why the residual-based
error estimate of \cref{plt:prop:rev-analytic-transfer} degrades near this
point.
\end{enumerate}
\end{exercise}

\begin{proof}[Solution]
(1)  With \(w=-1+q\),
\[
 \EulerE z=\EulerE\,w\EulerE^{w}=(-1+q)\EulerE^{q},
\]
and multiplying \(-1+q\) into
\(\sum_{k\ge0}q^{k}/k!\) the coefficient of \(q^{n}\) for \(n\ge2\) is
\(\tfrac1{(n-1)!}-\tfrac1{n!}=\tfrac{n-1}{n!}\), while the constant term is
\(-1\) and the coefficient of \(q\) is \(-1+1=0\).  Hence
\begin{equation}
\label{plt:eq:exr-branch-point}
 1+\EulerE z=\sum_{n\ge2}\frac{(n-1)\,q^{n}}{n!}
 =\frac{q^{2}}2+\frac{q^{3}}3+\frac{q^{4}}8+\frac{q^{5}}{30}
 +\BigOAt{q^{6}}{q\to0},
\end{equation}
an entire function of \(q\); this is
\cref{plt:lem:lw-critical-point}.

(2)  Then \(p^{2}=2(1+\EulerE z)=q^{2}\bigl(1+\tfrac23q+\tfrac14q^{2}
+\tfrac1{15}q^{3}+\cdots\bigr)\), so, choosing the branch with \(p\sim q\),
\[
 p=q\Bigl(1+\frac23q+\frac14q^{2}+\frac1{15}q^{3}+\cdots\Bigr)^{1/2}
 =q+\frac{q^{2}}3+\frac{5q^{3}}{72}+\frac{11q^{4}}{1080}+\cdots .
\]
Reverting this power series (a near-identity reversion in the ordinary
one-variable sense --- no logarithms are involved) gives
\[
 q=p-\frac{p^{2}}3+\frac{11}{72}p^{3}-\frac{43}{540}p^{4}
  +\frac{769}{17280}p^{5}-\frac{221}{8505}p^{6}+\cdots,
\]
whence \(\LambertW_0=-1+q\) as stated; the coefficients agree with the
standard tabulations \cite{corless-et-al-1996,dlmf-W}.  Since \(p\) enters
only through
\(p^{2}=2(1+\EulerE z)\), replacing \(p\) by \(-p\) gives the second solution
of \(w\EulerE^{w}=z\) near \(w=-1\), namely
\(\LambertW_{-1}(z)=-1-p-\tfrac13p^{2}-\tfrac{11}{72}p^{3}
 -\tfrac{43}{540}p^{4}-\cdots\).  That the reverted series has a positive
radius of convergence, and that the two displayed branches are exactly the two
solutions there, is \cref{plt:thm:lw-branch-point}; the formal reversion above
supplies the coefficients but not the radius.

(3)  The map \(w\mapsto w\EulerE^{w}\) has derivative \((1+w)\EulerE^{w}\),
which vanishes simply at \(w=-1\); the second derivative
\((2+w)\EulerE^{w}\) equals \(\EulerE^{-1}\ne0\) there.  So \(z+\EulerE^{-1}\)
has a simple \emph{quadratic} zero in \(w\) at \(w=-1\), and inverting a
quadratic critical point produces a two-valued square-root branch, with the
two branches exchanged by \(p\mapsto-p\).  Nothing logarithmic occurs: the
logarithmic scale of \cref{plt:thm:om-expansion} comes from the leading
behaviour of \(w\EulerE^{w}\) at \(w\to+\infty\), an entirely different
regime.

Near the branch point the residual-to-error transfer degrades exactly because
\(\LambertW'\) blows up there.  \Cref{plt:prop:rev-analytic-transfer} bounds
\(\AbsoluteValue{x_N-x_*}\) by
\(\AbsoluteValue{f(x_N)-y}/m\) with
\(m=\inf\AbsoluteValue{f'}\) on the interval between; here
\(f'(w)=(1+w)\EulerE^{w}\to0\) as \(w\to-1\), so \(m\to0\) and the bound
becomes vacuous.  Concretely, a small residual no longer certifies a small
error: near a quadratic critical point the error is of order the square root
of the residual, not of order the residual.
\end{proof}

\section{Analytic transfer}
\label{plt:sec:exercises-analytic}

\begin{exercise}[From residual to error, on a ray]
\label{plt:exer:exr-residual-to-error}
Let \(F(y)=y+\log y\) for \(y>0\), let \(X\in\RealNumbers\), and let
\(y_*>0\) be the unique solution of \(F(y)=X\).
\begin{enumerate}
\item Show that every \(y_N>0\) satisfies
\(\AbsoluteValue{y_N-y_*}\le\AbsoluteValue{y_N+\log y_N-X}\), and identify
the general principle at work.
\item Sharpen this to a two-sided estimate: if \(y_0>0\) and both \(y_N\) and
\(y_*\) lie in \([y_0,\infty)\), then
\[
 \frac{\AbsoluteValue{F(y_N)-X}}{1+1/y_0}
 \ \le\ \AbsoluteValue{y_N-y_*}\ \le\ \AbsoluteValue{F(y_N)-X}.
\]
Deduce that on such a ray the residual is not merely an upper bound for the
error but a proxy for it, correct to within the factor \(1+1/y_0\).
\item Explain why the same argument gives nothing for \(F(y)=y-\log y\) on
\((0,1)\).
\end{enumerate}
\end{exercise}

\begin{proof}[Solution]
\(F\) is continuously differentiable on \((0,\infty)\) with
\(F'(y)=1+1/y>1\), so \(F\) is strictly increasing, hence injective, and
\(F\bigl((0,\infty)\bigr)=\RealNumbers\), so \(y_*\) exists and is unique
(\cref{plt:prop:om-real-bijection}).  By the mean value theorem there is
\(\theta\) between \(y_N\) and \(y_*\) with
\[
 F(y_N)-X=F(y_N)-F(y_*)=F'(\theta)\,(y_N-y_*),
\]
and \(\theta>0\) gives \(F'(\theta)>1\).  Hence
\(\AbsoluteValue{y_N-y_*}=\AbsoluteValue{F(y_N)-X}/F'(\theta)
 \le\AbsoluteValue{F(y_N)-X}\), which is the claim.

The principle is \cref{plt:prop:rev-analytic-transfer}: if
\(\AbsoluteValue{f'}\ge m>0\) between the approximation and the true point,
then the error is at most the residual divided by \(m\); here \(m=1\) works on
the whole ray.  Two hypotheses are doing real work, and both are analytic:
the derivative bound, and the fact that \(y_N\) and \(y_*\) lie in a common
interval on which it holds.

(2)  The same \(\theta\) now satisfies \(y_0\le\theta\), because \(\theta\)
lies between \(y_N\) and \(y_*\) and both are \(\ge y_0\).  Since
\(F'(y)=1+1/y\) is decreasing,
\(1<F'(\theta)\le1+1/y_0\), and dividing
\(F(y_N)-X=F'(\theta)(y_N-y_*)\) by \(F'(\theta)\) gives both inequalities at
once.  So the residual determines the error to within the factor
\(1+1/y_0\), which tends to \(1\) as \(y_0\to\infty\).  This is the two-sided
enclosure of \cref{plt:thm:an-enclosure}, with
\(m=1\) and \(M=1+1/y_0\); that theorem states it on a compact interval and
adds that injectivity need not be assumed, being a conclusion.

(3)  For \(F(y)=y-\log y\) on \((0,1)\) one has \(F'(y)=1-1/y<0\), and
\(F'(y)\to0\) as \(y\to1^{-}\).  So \(m=\inf\AbsoluteValue{F'}=0\) on any
interval reaching to \(1\), and the bound degenerates.  That is not an
artefact: \(y=1\) is the critical point of \(y-\log y\), the two branches of
its inverse meet there, and near it the error is of order the square root of
the residual --- the same phenomenon as at the Lambert branch point
(\cref{plt:exer:exr-branch-point}).
\end{proof}

\begin{exercise}[Trap: a formal tail read as an analytic estimate]
\label{plt:exer:exr-trap-formal-to-analytic}
The following argument is offered for the asymptotic validity of the Wright
omega truncations.
\begin{quote}
``By \cref{plt:thm:om-residual}, \(\RevErr{X+L}{\WrightOmega_N}
=-\LamP{N+1}(L)t^{N+1}+\FormalBigOAt{t^{N+2}}{t}\).  Substituting \(x\) for
\(X\) and \(\log x\) for \(L\), the residual of \(\WrightOmega_N(x)\) is
\(-\LamP{N+1}(\log x)x^{-(N+1)}
 +\BigOAt{x^{-(N+2)}}{x\to+\infty}\).  Since
\(\AbsoluteValue{\frac{\Differential}{\Differential y}(y+\log y)}>1\)
for \(y>0\), \cref{plt:prop:rev-analytic-transfer} converts this into
\[
 \WrightOmega(x)-\WrightOmega_N(x)=\LamP{N+1}(\log x)x^{-(N+1)}
 +\BigOAt{x^{-(N+2)}}{x\to+\infty}.
\]
Hence the truncations are asymptotically valid.''
\end{quote}
Locate every unjustified step, say which parts of the conclusion are
nevertheless true, and state what has to be proved to repair it.
\end{exercise}

\begin{proof}[Solution]
There is exactly one unjustified step, and it is the second sentence.

\emph{The step.}  ``Substituting \(x\) for \(X\) and \(\log x\) for \(L\)'' in
an identity of \(\K[L][[t]]\) is not an operation on functions.  The symbol
\(\FormalBigOAt{t^{N+2}}{t}\) records that a formal series has outer valuation
at least \(N+2\); it asserts no inequality, contains no constant, and names no
domain.  There is no theorem in the volume, and no theorem at all, that
converts it into \(\BigOAt{x^{-(N+2)}}{x\to+\infty}\) without further
hypotheses.  Two independent obstructions:
\begin{itemize}
\item the formal tail is an \emph{infinite} sum of blocks, and its
substitution need not converge or even be defined pointwise;
\item even if a bound exists, its blocks carry powers of \(L\), so the honest
analytic remainder is of the form
\(\BigOAt{L^{K}x^{-(N+2)}}{x\to+\infty}\) for some \(K\), \emph{not}
\(\BigOAt{x^{-(N+2)}}{x\to+\infty}\).  The two differ by an unbounded factor.
\end{itemize}
The third and fourth sentences are then fine \emph{given} the second: the
derivative bound and \cref{plt:prop:rev-analytic-transfer} do transfer a
residual estimate to an error estimate, as in
\cref{plt:exer:exr-residual-to-error}.

\emph{What is true.}  The conclusion is correct in shape and wrong in
detail.  \Cref{plt:thm:om-analytic} establishes
\[
 \WrightOmega(x)-\WrightOmega_N(x)
 =\frac{\LamP{N+1}(\log x)}{x^{N+1}}
  +\BigOAt{\frac{(\log x)^{K}}{x^{\,N+2}}}{x\to+\infty},
 \qquad K=(N+2)^{2},
\]
and hence \eqref{plt:eq:om-sharp-order}.  The logarithmic power \(K\) cannot
be removed.  The next block alone contributes the term
\(\LamP{N+2}(\log x)\,x^{-(N+2)}\).  Its exact size is
\[
 \asymp(\log x)^{N+2}x^{-(N+2)},
 \qquad\text{which is \emph{not}}\qquad
 \BigOAt{x^{-(N+2)}}{x\to+\infty}.
\]

\emph{The repair.}  What must be supplied, and what the formal identity cannot
supply, is a genuine analytic estimate for the substituted tail: an explicit
finite sum of blocks plus a remainder bounded on an explicit ray, with the
constant and the exponent \(K\) named.  That is item~(2) of the proof of
\cref{plt:thm:om-analytic}, and in general it is
\cref{plt:thm:an-residual-analytic}, which converts a \emph{finite}
truncation's formal residual order into an inequality
\(\AbsoluteValue{\varphi(h(x))-x}\le Cx^{-N}(\log x)^{M}\) on a ray --- note
that its hypothesis makes \(B\) a finite sum of blocks, precisely so that
substitution is a finite computation with numbers.  The general lesson is the one printed at the head
of this appendix: \(\FormalBigOAt{\cdot}{t}\) and
\(\BigOAt{\cdot}{X\to+\infty}\) are different assertions, and no amount of
correct formal algebra produces the second.
\end{proof}

% ed.: source 5 asserts that the residual-to-error principle "proves the
% ed.: asymptotic validity of every truncation directly from" the formal
% ed.: residual identity, and source 1 asserts an analytic O-bound for the
% ed.: branch residual "by construction".  Neither is a proof; the missing
% ed.: ingredient is the substituted-tail estimate, which is why this
% ed.: exercise exists.

\begin{exercise}[Counterexample: what the minimal Poincar\'e condition fails to do]
\label{plt:exer:exr-poincare}
A source proposes to define ``\(f\) has the asymptotic expansion
\(\sum_np_n(\log x)x^{-n}\)'' by the minimal Poincar\'e requirement that, for
each \(N\), the remainder after the term \(n=N\) be
\(\LittleOAt{p_N(\log x)x^{-N}}{x\to+\infty}\).
\begin{enumerate}
\item Show that at a \emph{single} index \(N\) the requirement pins \(p_N\)
only modulo polynomials of strictly smaller degree.
\item Show that, imposed at \emph{every} index, it does nevertheless pin the
blocks --- so nonuniqueness is not the objection to it.
\item Give the two objections that are real: the requirement is degenerate as
soon as one block vanishes, and at any single index it controls the remainder
only within its own outer order, not at the next one.
\item Show that the definition actually used (\cref{plt:def:mul-asymptotic})
has neither defect and determines the blocks.
\end{enumerate}
\end{exercise}

% ed.: the inherited version of this exercise asserted outright that the
% ed.: minimal condition "does not determine the blocks", and proved it by
% ed.: changing one block and checking the condition at that index alone.
% ed.: The check at index N+1 then fails, so the modified family does not
% ed.: satisfy the definition and the claimed nonuniqueness is not there.
% ed.: Part (2) below proves uniqueness; the defects the exercise exists to
% ed.: expose are the two in part (3), and they are genuine.
\begin{proof}[Solution]
(1)  Fix \(N\) and take the two candidate blocks
\[
 p_N(u)=u^{5},
 \qquad
 q_N(u)=u^{5}+u^{3}.
\]
They differ by \(u^{3}\), and
\[
 \bigl(q_N(\log x)-p_N(\log x)\bigr)x^{-N}=(\log x)^{3}x^{-N}
 =\LittleOAt{(\log x)^{5}x^{-N}}{x\to+\infty}.
\]
So the remainder after \(q_N\) differs from the remainder after \(p_N\) by
something that is little-o of \emph{either} of the two last retained terms: the
requirement read at the index \(N\) alone cannot distinguish the two.  The same
computation works with \(u^{3}\) replaced by any polynomial of degree \(<5\).
The reason is that the requirement compares the remainder with the last
retained \emph{term}, and on a two-generator scale a term is a whole block,
whose parts have different growth.  On a one-generator scale
\(\SetOf{x^{-n}}\) the terms are monomials and the classical condition
\cite{olver-asymptotics} suffers from none of this.

(2)  Now impose the requirement at every index, which is what the proposed
definition says.  Suppose \((p_n)\) and \((q_n)\) both work for the same
\(f\), and let \(n_0\) be least with \(d\DefinitionEquals q_{n_0}-p_{n_0}\ne0\).
Write \(R_p\) and \(R_q\) for the two remainders after the index \(n_0+1\).
Since \(p_n=q_n\) for \(n<n_0\),
\[
 R_q-R_p=d(\log x)\,x^{-n_0}
  +\bigl(p_{n_0+1}-q_{n_0+1}\bigr)(\log x)\,x^{-n_0-1}.
\]
Multiply by \(x^{n_0}\).  The requirement at \(n_0+1\) makes \(R_p\) and
\(R_q\) little-o of a fixed block times \(x^{-n_0-1}\), so
\(x^{n_0}\AbsoluteValue{R_p}\) and \(x^{n_0}\AbsoluteValue{R_q}\) are
\(\LittleOAt{(\log x)^{K}/x}{x\to+\infty}\) for some \(K\), and so is the last
displayed summand.  Hence \(\AbsoluteValue{d(\log x)}\to0\).  But \(d\ne0\) is
a polynomial, so \(\AbsoluteValue{d(\log x)}\) is bounded below by a positive
constant for large \(x\).  Contradiction; \(d=0\), and the blocks agree.
It is the clause at \(n_0+1\), not the one at \(n_0\), that does the work ---
which is exactly the observation part~(1) makes, read the other way round.

(3)  Two defects survive part~(2), and they are what disqualify the
requirement.

\emph{It degenerates at a vanishing block.}  ``\(g=\LittleOAt{0}{x\to+\infty}\)''
says \(\AbsoluteValue g\le\varepsilon\cdot0\) eventually for every
\(\varepsilon>0\), i.e.\ that \(g\) vanishes identically for large \(x\).  So
as soon as one block \(p_N\) is zero, the requirement at that index demands an
exactly vanishing remainder.  Take \(f(x)=1+x^{-2}\), which has the finite
expansion \(F=1+t^{2}\) and no expansion problem whatsoever.  Its blocks are
\(p_0=1\), \(p_1=0\), \(p_2=1\), and the requirement at \(N=1\) reads
\(x^{-2}=\LittleOAt{0}{x\to+\infty}\), which is false.  The proposed
definition therefore does not apply to \(f\) at all, although \(f\) is a
polynomial in \(x^{-1}\).

\emph{It gives no next-order control.}  At the index \(N\), and provided
\(\degL p_N\ge1\), the requirement permits a remainder as large as
\((\log x)^{\degL p_N-1}x^{-N}\), which dominates the whole next block
\(p_{N+1}(\log x)x^{-N-1}\) by a factor of order \(x\) up to powers of
\(\log x\).  A reader who takes ``expansion through \(n=N\)'' to mean
``correct to the order of the next block'' is therefore reading in an estimate
that the requirement never made --- the same confusion, one level down, as
\cref{plt:exer:exr-trap-formal-to-analytic}.

(4)  Neither defect afflicts \cref{plt:def:mul-asymptotic}, which never
divides by a block and never compares the remainder with the last retained
term.

\emph{What does determine them.}  \Cref{plt:def:mul-asymptotic}, and its
extension to rational blocks in \cref{plt:def:an-expansion}, compare the
remainder not with the last term but with the \emph{next outer order}: for
every \(N\) there are \(M_N,C_N,x_N\) with
\[
 \AbsoluteValue{f(x)-\bigl(\Trunc{F}{N}\bigr)(x)}
 \le C_N\,x^{-N}(\log x)^{M_N}
 \qquad(x\ge x_N).
\]
Suppose \(F\) and \(G\) both satisfy this and \(F\ne G\); let
\(n_0=\ordt(F-G)\) and \(p=[t^{n_0}](F-G)\ne0\).  Applying the condition at
\(N=n_0+1\) to both and subtracting --- the two truncations
\(\Trunc{F}{n_0+1}\) and \(\Trunc{G}{n_0+1}\) retain the blocks of index
\(\le n_0\), and those agree below \(n_0\), so their difference is the
\emph{single} term \(p(\log x)x^{-n_0}\) ---
\[
 \AbsoluteValue{p(\log x)}\,x^{-n_0}
 \le C\,x^{-n_0-1}(\log x)^{M}
\]
for large \(x\).  Dividing
by \(x^{-n_0}\), the left side tends to \(\infty\) or to a nonzero constant in
absolute value if \(p\ne0\) --- indeed
\(\AbsoluteValue{p(\log x)}\ge c(\log x)^{\degL p}\) for large \(x\) --- while
the right side is \(Cx^{-1}(\log x)^{M}\to0\).  Contradiction; so \(F=G\).
Uniqueness holds --- this is \cref{plt:prop:an-uniqueness}, which also covers
rational blocks and closed subsectors --- and the expansion map of
\cref{plt:thm:mul-analytic-product} is therefore well defined.
\end{proof}

\begin{exercise}[Trap: a guard that was not taken in the analytic estimate]
\label{plt:exer:exr-trap-product-guard}
Let \(f(X)=X^{2}\) exactly, so \(f\) realises \(F=t^{-2}\) with no error at
all, and let \(g\) satisfy
\(g-\bigl(\Trunc{G}{N}\bigr)=\BigOAt{X^{-N}}{X\to+\infty}\) for some
\(G\in\PLpow\).  A draft concludes:
\begin{quote}
``\(fg-\bigl(\Trunc{F}{N}\bigr)\bigl(\Trunc{G}{N}\bigr)
=f\bigl(g-\Trunc{G}{N}\bigr)
=\BigOAt{X^{-N}}{X\to+\infty}\), since the first factor is exact and the
second remainder is \(\BigOAt{X^{-N}}{X\to+\infty}\).''
\end{quote}
Find the error, give the correct exponent, and state the general precision
rule for an analytic product.
\end{exercise}

\begin{proof}[Solution]
The error is that the first factor is not \(\BigOAt{1}{X\to+\infty}\); it is
\(X^{2}\).  Multiplying a remainder of size \(X^{-N}\) by \(X^{2}\) gives
\[
 f\bigl(g-\Trunc{G}{N}\bigr)=\BigOAt{X^{-N+2}}{X\to+\infty},
\]
two orders worse than claimed.  To obtain
\(\BigOAt{X^{-N}}{X\to+\infty}\) in the product one must know \(g\) to
precision \(N+2\), i.e.\ \(g-\Trunc{G}{N+2}=\BigOAt{X^{-N-2}}{X\to+\infty}\).

\emph{General rule.}  Let \(f\approx F\), \(g\approx G\) with \(a=\ordt F\),
\(b=\ordt G\).  Suppose
\[
 f-\Trunc{F}{P}=\BigOAt{X^{-P}(\log X)^{M}}{X\to+\infty},
 \qquad
 g-\Trunc{G}{Q}=\BigOAt{X^{-Q}(\log X)^{M'}}{X\to+\infty}.
\]
Write
\[
 fg-\Trunc{F}{P}\,\Trunc{G}{Q}
 =\Trunc{F}{P}\bigl(g-\Trunc{G}{Q}\bigr)
  +\bigl(f-\Trunc{F}{P}\bigr)g .
\]
The truncation \(\Trunc{F}{P}\) has outer order \(a\), and \(g\) has outer
order \(b\); so as functions on the ray they are of size
\(X^{-a}\) and \(X^{-b}\) up to powers of \(\log X\).  The two summands are
therefore of sizes
\[
 \BigOAt{X^{-a-Q}(\log X)^{\ast}}{X\to+\infty}
 \qquad\text{and}\qquad
 \BigOAt{X^{-P-b}(\log X)^{\ast}}{X\to+\infty}
\]
respectively.  So the product is known to absolute order \(N\) if and only
if
\[
 P\ \ge\ N-b
 \qquad\text{and}\qquad
 Q\ \ge\ N-a ,
\]
which is verbatim the formal truncated-product rule
\eqref{plt:eq:trunc-product} of \cref{plt:prop:trunc-product}, and its
error-tolerant version \cref{plt:cor:mul-error-rule}.  The analytic guard and
the formal guard coincide; what does \emph{not} coincide is the naive
symmetric rule \(P=Q=N\), which is valid only when \(a,b\ge0\).  In the
example \(a=-2\), \(b\ge0\), so \(Q\ge N+2\), which is the two orders that
went missing.

Two riders.  First, the logarithmic powers accumulate: the exponent in the
product's remainder is at least \(M+\degL\) of the retained blocks, so a
statement of the form \(\BigOAt{X^{-N}}{X\to+\infty}\) with no logarithmic
factor is almost never available (\cref{plt:exer:exr-trap-formal-to-analytic}).
Second, the same arithmetic with two \emph{inexact} factors is where a draft
of the composition transfer went wrong: bounding a product of remainders by
the smaller of the two exponents is legitimate only when each remainder is
multiplied by a factor of nonpositive outer order, and the hypothesis
\(N\ge a+\tilde a\) is exactly what makes that so.
\end{proof}

\begin{exercise}[Realizations form an algebra, and the map that forgets]
\label{plt:exer:exr-realizations-algebra}
Let \(\K\subseteq\ComplexNumbers\), \(S=[X_0,+\infty)\) with the real
logarithm, and \(\mathcal A(S)=\SetOf{f\colon S\to\ComplexNumbers:
\exists F\in\PLlau,\ f\approx F}\) in the sense of
\cref{plt:def:mul-asymptotic}.
\begin{enumerate}
\item Show \(F\) is uniquely determined by \(f\).
\item Show \(\mathcal A(S)\) is a \(\K\)-algebra and \(f\mapsto F\) is a
\(\K\)-algebra homomorphism into \(\PLlau\).  (For surjectivity, quote
\cref{plt:thm:an-borel-ritt}; do not attempt it by hand.)
\item Identify its kernel and show that it is not zero, so that the
homomorphism is not injective; give an explicit nonzero element.
\item Deduce that ``\(f\approx F\)'' can never be upgraded to ``\(f\) is the
sum of \(F\)'' without extra hypotheses.
\end{enumerate}
\end{exercise}

\begin{proof}[Solution]
(1)  This is the uniqueness argument of
\cref{plt:exer:exr-poincare}: if \(f\approx F\) and \(f\approx G\) with
\(F\ne G\), let \(n_0=\ordt(F-G)\) and \(p=[t^{n_0}](F-G)\ne0\).  Subtracting
% ed.: this step was written as an inequality on "(F-G)(X)", which is a formal
% ed.: series and not a function.  What the two estimates bound is the
% ed.: difference of the two TRUNCATIONS, which is a polynomial expression.
the two estimates at cutoff \(N=n_0+1\) gives
\(\AbsoluteValue{\bigl(\Trunc{F}{n_0+1}-\Trunc{G}{n_0+1}\bigr)(X)}
 \le C X^{-n_0-1}(\log X)^{M}\) for large \(X\), and that difference is the
single term \(p(\log X)X^{-n_0}\); dividing
by \(X^{-n_0}\) makes the left side bounded below by a positive constant times
\((\log X)^{\degL p}\) and the right side tend to \(0\).  Contradiction.

(2)  Linearity is immediate from the triangle inequality and linearity of
\(\Trunc{\cdot}{N}\).  For products, \cref{plt:thm:mul-analytic-product} gives
\(fg\approx FG\); its proof is the guard computation of
\cref{plt:exer:exr-trap-product-guard} run at \(P=N-b\), \(Q=N-a\), which is
legitimate because \(f\approx F\) supplies an estimate at \emph{every} cutoff,
in particular at those two.  The constant function \(1\) realises \(1\), so
the map is unital.  Surjectivity is the Borel--Ritt theorem on this scale,
\cref{plt:thm:an-borel-ritt}: every \(F\in\PLlau\) is realised by some
\(C^{\infty}\) function, real valued when \(\K\subseteq\RealNumbers\).  It is
worth noticing what that theorem does \emph{not} say --- it produces a smooth
realisation, not an analytic one, and it produces no canonical choice, since
by (3) below any two realisations differ by a flat function.

(3)  The kernel is the ideal \(\mathcal N\) of
\emph{flat} functions of \cref{plt:def:an-algebra},
\(\SetOf{f:f=\BigOAt{X^{-N}}{X\to+\infty}\ \text{for every }N}\), which is
exactly the condition \(f\approx0\).
It is an ideal, being the kernel of a ring homomorphism, and it is not zero:
\(f(X)=\EulerE^{-X}\) is nonzero everywhere on \(S\) and satisfies
\(\EulerE^{-X}X^{N}\to0\) for every \(N\), so \(\EulerE^{-X}\approx0\).  Hence
\(\EulerE^{-X}\) and \(0\) have the same expansion, and the map is not
injective.  A second witness, of a different character:
\(\sin(\EulerE^{X})\) is bounded and oscillatory, but
\(\EulerE^{-X}\sin(\EulerE^{X})\) is again flat and nonzero.

(4)  Consequently the expansion \(F\) determines \(f\) only modulo flat
functions, and \(\mathcal N\) is large: by
\cref{plt:prop:an-flat-ideal} it is prime but not maximal, is not stable under
differentiation, and contains no smallest nonzero scale, since
\(f\in\mathcal N\) eventually nonvanishing gives
\(f^{2}\in\mathcal N\) with \(f^{2}=\LittleOAt{f}{X\to+\infty}\).  ``\(f\) is the sum of \(F\)'' would assert both that the series
converges and that its sum is \(f\); the first is a claim about
\(\sum_np_n(\log X)X^{-n}\) as a numerical series and the second would force
injectivity, which (3) refutes.  What \(\approx\) does assert is exactly
\eqref{plt:eq:mul-asymptotic}: a family of finite inequalities, one per
cutoff, each on its own ray.  Any statement stronger than that needs a summation
method and its own theorem.
\end{proof}

\begin{exercise}[Exactness at a point, and what the numbers do and do not show]
\label{plt:exer:exr-omega-numerics}
\begin{enumerate}
\item Show that every truncation is exact at the single point \(x=1\), that
is, that \(\WrightOmega_N(1)=\WrightOmega(1)=1\) for every \(N\); and explain
why this is not a convergence statement.
\item Take \(x=20\), so that \(L=\log20=2.995732\ldots\).  The terms of the
series, and the errors after each of them, are, to four significant
figures,
\begin{center}
\begin{tabular}{@{}rrr@{}}
\toprule
\(n\) & \(\LamP{n}(L)\,20^{-n}\) & \(\WrightOmega(20)-\WrightOmega_n(20)\)\\
\midrule
\(1\) & \(1.498\times10^{-1}\) & \(3.507\times10^{-3}\)\\
\(2\) & \(3.729\times10^{-3}\) & \(-2.220\times10^{-4}\)\\
\(3\) & \(-1.880\times10^{-4}\) & \(-3.395\times10^{-5}\)\\
\(4\) & \(-3.267\times10^{-5}\) & \(-1.282\times10^{-6}\)\\
\(5\) & \(-1.433\times10^{-6}\) & \(1.503\times10^{-7}\)\\
\(6\) & \(1.259\times10^{-7}\) & \(2.443\times10^{-8}\)\\
\(7\) & \(2.353\times10^{-8}\) & \(8.975\times10^{-10}\)\\
\(8\) & \(1.047\times10^{-9}\) & \(-1.496\times10^{-10}\)\\
\bottomrule
\end{tabular}
\end{center}
Check the table against \eqref{plt:eq:om-analytic-error}, and against the sign
prediction of \cref{plt:exer:exr-degree-law-zero}.
\item The terms decrease monotonically through \(n=8\) and beyond.  Does
that show that the series \(\sum_n\LamP{n}(\log x)\,x^{-n}\) converges at
\(x=20\)?  What does this volume actually prove?
\end{enumerate}
\end{exercise}

\begin{proof}[Solution]
% ed.: the display here was garbled ("1 - LamP_0(0)(-1) + ...").  By
% ed.: \cref{plt:def:om-truncation}, Omega_N = X + Trunc(P)_{N+1}, so the
% ed.: evaluation is the single sum written below.
(1)  At \(x=1\) one has \(L=\log1=0\), and \(\LamP{n}(0)=0\) for every \(n\)
by \eqref{plt:eq:lambert-boundary}.  By \cref{plt:def:om-truncation},
\(\WrightOmega_N=X+\Trunc{\mathcal P}{N+1}\), so
\[
 \WrightOmega_N(1)=1+\sum_{n=0}^{N}\LamP{n}(0)\cdot1^{-n}=1+0=1 .
\]
And
\(\WrightOmega(1)=1\) because \(1+\log1=1\).  So every truncation is exact at
that one point.

This says nothing about convergence, and it is worth being explicit about why.
The statement ``\(\WrightOmega_N(1)=\WrightOmega(1)\) for all \(N\)'' is a
family of identities between numbers at a single \(x\).  An asymptotic
statement \(\BigOAt{\cdot}{x\to+\infty}\) has no content at a single point
--- multiplying a function by any factor that equals \(1\) at \(x=1\) changes
nothing about it.  And a convergence statement would be about the tail
\(\sum_{n>N}\), which here is identically zero at \(x=1\) for a trivial
reason: every summand vanishes.  The same phenomenon in Lambert coordinates is
the exactness of the branch expansion at \(z=\EulerE\), where \(\xi_0=1\) and
\(\ell_0=0\); this is \cref{plt:ex:om-exact-at-one} and
\cref{plt:prop:lw-exact-at-e}, and it is a consequence of \(\LamP{n}(0)=0\),
of nothing else.  Indeed the convergence criterion of the Lambert part fails
outright at that point.

(2)  \Cref{plt:thm:om-analytic} asserts
\(\WrightOmega(x)-\WrightOmega_N(x)=\LamP{N+1}(L)x^{-(N+1)}
 +\BigOAt{L^{K}x^{-(N+2)}}{x\to+\infty}\).  Reading the table, the error after
\(N\) should be close to the term at \(n=N+1\), and it is, in every row:
\(3.507\times10^{-3}\) against \(3.729\times10^{-3}\);
\(-2.220\times10^{-4}\) against \(-1.880\times10^{-4}\);
\(-3.395\times10^{-5}\) against \(-3.267\times10^{-5}\);
and so on down to \(8.975\times10^{-10}\) against \(1.047\times10^{-9}\).
Both magnitude and sign agree throughout, which is
\eqref{plt:eq:om-sharp-order} in action.

The signs are themselves informative.  Here \(x=20\) lies above the crossing
point \(x_\ast=8.7465\ldots\) of \cref{plt:exer:exr-degree-law-zero}\,(2), and
\(\LamP{2}(L)>0\) since \(L=2.9957>2\); so \(\WrightOmega_1\) undershoots, as
that exercise predicts.  The later signs follow the sign of
\(\LamP{n}(L)\), which changes as \(L\) passes the real roots of the
successive Lambert polynomials.  There is no alternating pattern in \(n\) to
be read off, and none should be expected: the sign is the sign of a
polynomial value, not of a coefficient.

(3)  No --- monotone decrease of finitely many computed terms is evidence, not
proof, and in an asymptotic series the terms typically decrease for a long time
before turning.  The series does converge at \(x=20\), but for a reason that
has nothing to do with the table.  This volume proves four separate things,
which must be kept apart:
\begin{itemize}
\item the \emph{formal} identity \eqref{plt:eq:om-expansion}, an equality in
\(\Gnear\);
\item the \emph{asymptotic} statement \eqref{plt:eq:om-analytic-error}, valid
as \(x\to+\infty\) with an explicit logarithmic loss \(K=(N+2)^{2}\);
\item the \emph{frozen} convergence of \cref{plt:prop:om-frozen-convergence}:
for each fixed \(u_0\) the series \(\sum_{n\ge1}\LamP{n}(u)s^{n}\) has a
positive radius in \(s\), uniformly near \(u_0\);
% ed.: the inherited solution stopped at the third item and concluded that
% ed.: "this volume does not determine whether that series converges".  That
% ed.: is false.  \cref{plt:thm:lw-basepoint}, taken at K = x and log K = L,
% ed.: gives M = L and a seating condition 2(1+L) < x, under which the
% ed.: substituted series converges absolutely.  At x = 20 the condition reads
% ed.: 7.99 < 20.  Recomputed: the partial sum through n = 40 agrees with
% ed.: omega(20) - 20 + log 20 to fifteen digits.
\item the \emph{substituted} convergence of \cref{plt:thm:lw-basepoint}: with
\(\xi=x\), \(K=x\) and \(\log K=L\) that theorem has \(M=L\), its seating
condition \eqref{plt:eq:lw-seating} reads \(2(1+L)<x\), and under it the
series \(\sum_{n\ge1}\LamP{n}(L)x^{-n}\) converges absolutely, with sum
\(\WrightOmega(x)-x+L\).
\end{itemize}
The fourth item settles the question the table raises, and settles it in the
affirmative: at \(x=20\) the seating condition reads \(7.991<20\), so the
series converges absolutely there.  The two statements must still not be
confused.  Convergence at a fixed \(x\) is a statement about one point, proved
by a Cauchy estimate on a two-variable holomorphic function; the sharp order
of (2) is a statement about the limit \(x\to+\infty\), with \(u=\log x\) and
\(s=1/x\) linked.  Neither implies the other, and
\cref{plt:prop:om-frozen-convergence} implies neither: it freezes \(u\) and
its radius \(r(u)\) is not claimed to be bounded below as
\(\AbsoluteValue{u}\to\infty\) (\cref{plt:rmk:om-open}).  What the table
alone shows is the sharp-order statement of (2), and nothing more; the
convergence is a theorem imported from \cref{plt:sec:lambert-branches}, not a
reading of the numbers.
\end{proof}

\begin{remark}[Where each exercise came from]
\label{plt:rmk:exr-provenance}
Of the \(56\) exercises above, \(31\) descend from a source problem --- in most
cases from several at once, since the corpus's \(83\) problems are heavily
duplicated, and in most cases with the statement strengthened, the hypotheses
printed, or the supplied answer corrected; \(9\) were built from worked
examples in reports~4, 5 and~6, which set no exercises at all; and \(16\) are
new, written because a technique of the body had no exercise anywhere in the
corpus.  Every trap and every counterexample exercise --- \(8\) and \(11\)
respectively --- is in the last two groups.

The source problems that left no descendant were of three kinds: duplicates of
a retained problem, arithmetic drill on an operation already exercised
elsewhere, and problems whose \emph{supplied} solution is wrong.  Of the last
kind there are five, all corrected above and all recorded in
\cref{plt:sec:repairs}: the class in which \(\recip{(L+t)}\) lives, the
universal third inverse block, the sign weight in the lower Lambert branch,
the range of the Stirling sum for \(\shiftop{0}{k}\), and the passage from a
formal residual identity to an analytic error bound.
\end{remark}

% ---- fragment 19-notation ----
\chapter{Notation, and a concordance with the six sources}
\label{plt:sec:notation}

\section{The normative source}

Every mathematical symbol in this volume is defined normatively in the corpus
notation catalogue, \emph{FabiusFunction Mathematical Notation Catalogue}, whose
entries \texttt{polynomial-\allowbreak logarithmic-\allowbreak ambient},
\texttt{polynomial-\allowbreak logarithmic-\allowbreak leading-\allowbreak data},
\texttt{polynomial-\allowbreak logarithmic-\allowbreak truncation},
\texttt{polynomial-\allowbreak logarithmic-\allowbreak differential},
\texttt{polynomial-\allowbreak logarithmic-\allowbreak composition},
\texttt{polynomial-\allowbreak logarithmic-\allowbreak reversion},
\texttt{wright-\allowbreak omega-\allowbreak lambert-\allowbreak polynomials}, and
\texttt{lambert-\allowbreak branch-\allowbreak expansions} cover exactly the material of this volume.
This document introduces no new mathematical notation.  Its preamble defines
short local aliases --- \(\PLpow\), \(\PLlau\), \(\PLrat\), \(\PLit\),
\(\Gnear\), \(\ordt\), \(\lm\), \(\lc\) and the rest --- and every one of them
expands to a catalogue command; the aliases exist to keep displayed formulas
readable, not to name anything the catalogue does not already name.

\section{Symbols used throughout}

\begin{longtable}{@{}l l p{0.56\textwidth}@{}}
\caption{Principal notation.}\label{plt:tab:notation}\\
\toprule
Symbol & Alias & Meaning\\
\midrule
\endfirsthead
\toprule
Symbol & Alias & Meaning\\
\midrule
\endhead
\(X\) & --- & the large variable; the default regime is \(X\to+\infty\) along
  the positive real ray\\
\(t\) & --- & \(X^{-1}\), the outer generator\\
\(L\) & --- & \(\log X\), the logarithmic generator, formally independent of
  \(t\)\\
\(\K\) & \verb|\K| & the coefficient field, of characteristic zero\\
\(\PLpow\) & \verb|\PLpow| & \(\K[L][[t]]\), the polynomial-log power-series
  ring\\
\(\PLlau\) & \verb|\PLlau| & \(\K[L]((t))\), its Laurent extension: finitely
  many positive powers of \(X\)\\
\(\PLrat\) & \verb|\PLrat| & \(\K(L)((t))\), the rational-log Laurent field\\
\(\PLit\) & \verb|\PLit| & \(\K((L^{-1}))((t))\), the iterated Laurent-log
  field\\
\(\Gnear\) & \verb|\Gnear| & \(\SetOf{X+A:A\in\PLpow}\), the near-identity
  group under composition\\
\(\ordt F\) & \verb|\ordt| & the outer valuation, with
  \(\ordt 0=+\infty\)\\
\(\lm(F),\ \lc(F)\) & \verb|\lm|, \verb|\lc| & leading monomial and scalar
  leading coefficient; the leading \emph{block} is the whole polynomial
  \(p_{\ordt F}\), a third and different object\\
\(\degL p\) & \verb|\degL| & degree in \(L\)\\
\(\dprofile{F}(n)\) & \verb|\dprofile| & the log-degree profile, \(-\infty\)
  where the block vanishes\\
\(\Trunc{F}{N}\) & \verb|\Trunc| & the \emph{exclusive} truncation
  \(\sum_{n<N}p_n(L)t^n\)\\
\(\FormalBigOAt{t^N}{t}\) & --- & a formal valuation tail: coefficient
  agreement below outer order \(N\), and nothing analytic\\
\(\BigOAt{A(X)}{X\to\infty}\) & --- & an analytic estimate, with the regime
  printed\\
\(\DX\) & \verb|\DX| & the distinguished derivation
  \(t\partial_L-t^{2}\partial_t\)\\
\(\shiftop{n}{k}\) & \verb|\shiftop| & \(\prod_{j=0}^{k-1}(\partial_L-n-j)\),
  empty and equal to the identity at \(k=0\)\\
\(\rev{F}\) & \verb|\rev| & compositional inverse,
  \(F\compose\rev{F}=\rev{F}\compose F=X\)\\
\(\recip{F}\) & \verb|\recip| & multiplicative inverse, \(F\recip{F}=1\)\\
\(\RevErr{F}{H}\) & \verb|\RevErr| & the reversion residual
  \(F\compose H-X\); \emph{both} arguments are part of the object\\
\(\WrightOmega\) & --- & Wright omega, the inverse of \(Y\mapsto Y+\log Y\)\\
\(\LamP{n}\) & \verb|\LamP| & the canonical zero-based Lambert polynomials\\
\(\LamPsc{n}\) & \verb|\LamPsc| & the scaled family \(n!\LamP{n}\), for
  \(n\ge1\)\\
\(\stirlingone{n}{k}\) & \verb|\stirlingone| & unsigned Stirling numbers of the
  first kind (cycle numbers)\\
\(\stirlingsigned{n}{k}\) & \verb|\stirlingsigned| & signed Stirling numbers of
  the first kind\\
\(\precdom,\succdom\) & --- & the dominance order on monomials\\
\bottomrule
\end{longtable}

\section{Concordance with the absorbed sources}

The six sources used mutually incompatible local abbreviations, and in two
cases the \emph{same} abbreviation for different objects.  A reader returning
to a source through the Git history needs the following table.

\begin{longtable}{@{}l l l@{}}
\caption{Local abbreviations of the absorbed sources and their meaning here.
S1--S6 are as identified in \cref{plt:tab:provenance}.}
\label{plt:tab:concordance}\\
\toprule
Source & Local & This volume\\
\midrule
\endfirsthead
\toprule
Source & Local & This volume\\
\midrule
\endhead
S1 & \verb|\Plog|, \verb|\Rlog|, \verb|\Hlog| &
  \(\PLlau\), \(\PLrat\), \(\PLit\)\\
S1 & \verb|\val|, \verb|\trunc| & \(\ordt\), exclusive \(\Trunc{\cdot}{N}\)\\
S1 & \verb|\invcomp| & \(\rev{\cdot}\)\\
S2 & \verb|\val| & \(\ordt\)\\
S2 & \verb|\rev|, \verb|\recip| & \(\rev{\cdot}\), \(\recip{\cdot}\)
  (already the catalogue commands)\\
S2 & \verb|\Dop| \(=\mathrm{D}\) & \(\DX\)\\
S2 & \verb|\normalexp| & normal-ordered shifted exponential; expanded inline
  here\\
S3 & \verb|\Rring|, \verb|\Pring|, \verb|\Fring|, \verb|\Ggroup| &
  \(\PLpow\), \(\PLlau\), \(\PLit\), \(\Gnear\)\\
S3 & \verb|\ordt|, \verb|\D| \(=\mathrm{D}\) & \(\ordt\), \(\DX\)\\
S4 & \verb|\PL| & \(\PLlau\) --- \emph{not} \(\PLpow\); see below\\
S4 & \verb|\RL|, \verb|\HL|, \verb|\Gnear| & \(\PLrat\), \(\PLit\),
  \(\Gnear\)\\
S4 & \verb|\ordx|, \verb|\degell| & \(\ordt\), \(\degL\)\\
S5 & \verb|\PL| & \(\PLpow\) --- \emph{not} \(\PLlau\); see below\\
S5 & \verb|\PLL|, \verb|\QPL| & \(\PLlau\), \(\PLrat\)\\
S5 & \verb|\D| \(=\mathcal{D}\) & \(\DX\)\\
S6 & \verb|\PLpoly|, \verb|\PLfield| & \(\PLlau\), \(\PLit\)\\
S6 & \verb|\EulerDerivation| & \(\EulerOp=X\DX\)\\
S6 & \verb|\trunc{F}{N}| \(=[F]_{\le N}\) & \emph{inclusive}; equals
  \(\Trunc{F}{N+1}\)\\
S6 & \verb|\compinv| & \(\rev{\cdot}\)\\
\bottomrule
\end{longtable}

\begin{warningbox}[title={Two abbreviations that collide across sources}]
\label{plt:rmk:concordance-collisions}
\begin{enumerate}
\item \verb|\PL| denotes the \emph{Laurent} ring \(\K[L]((t))\) in S4 and the
  \emph{power-series} ring \(\K[L][[t]]\) in S5.  A statement transported from
  one to the other without checking which ring is meant can silently acquire or
  lose the finitely many positive powers of \(X\); the unit criterion and the
  antiderivative obstruction both change.  In S4 the power-series subring was
  written \([\PLlau]_{\ge0}\), which is \(\PLpow\).
\item \verb|\trunc| is the exclusive \(\operatorname{Trunc}_{<N}\) in S1 and the
  inclusive \([\,\cdot\,]_{\le N}\) in S6.  Every cutoff index transported from
  S6 into this volume has been shifted by one block.  A residual certificate
  stated with the wrong convention is off by exactly one order, which is
  precisely the amount that makes a correct algorithm look wrong and an
  incorrect one look right.
\end{enumerate}
\end{warningbox}

\section{Conventions that are not merely notational}

Three conventions are recorded here because getting them wrong changes
theorems, not only glyphs.

\begin{enumerate}
\item \textbf{Empty products and sums.}  \(\shiftop{n}{0}=1\); the
  Lagrange--B\"urmann coefficient product is empty at \(r=1\); \(C_{0,0}=1\)
  and \(C_{m,0}=0\) for \(m>0\); \(\stirlingsigned{0}{0}=1\).  Every statement
  below is written so that these conventions make the boundary case correct
  rather than exceptional.
\item \textbf{The three \(\BigO\) roles.}  A formal valuation tail, an analytic
  estimate, and an algorithmic cost bound are three different assertions and
  are printed with three different tokens.  None may be inferred from another.
  In particular, \(\FormalBigOAt{t^{N}}{t}\) carries no uniformity in \(L\)
  whatsoever, and a bound uniform in \(L\) is a separate theorem with separate
  hypotheses.
\item \textbf{\(t^{N}\PLpow\) is not an ideal of \(\PLlau\).}  In the Laurent
  ring \(t\) is a unit, so the principal ideal it generates is the whole ring.
  Quotient-algebra language is therefore available on \(\PLpow\) and not on
  \(\PLlau\); on the latter, use \(\ordt(F-G)\ge N\), the formal tail token, or
  equality of exclusive truncations.
\end{enumerate}

% ---- fragment 20-repairs ----
\chapter{Ledger of editorial repairs}
\label{plt:sec:repairs}

Every correction this consolidation made to its sources is marked at the
point of repair in this document's own \LaTeX{} source with an
\texttt{\%~ed.:} comment.  The table below is generated from those marks, so it
cannot drift out of step with the text.  Repairs of pure notation --- a
symbol renamed to its catalogue command --- are not listed; those are recorded
once, in the concordance of \cref{plt:tab:concordance}.  What is listed is
every place where a source's mathematical statement, proof, hypothesis, index
range, or cost claim was wrong, incomplete, or overstated.

Three kinds of defect recur often enough to name.  The first is an
\emph{omitted hypothesis}: a statement true under conditions the source did not
print, most often a missing empty-product convention, a missing guard-precision
requirement, or a missing restriction to \(n\ge1\).  The second is
\emph{overclaiming}: an analytic conclusion drawn from formal algebra alone,
typically a convergence or uniformity assertion that the argument does not
support.  The third is a \emph{cost or sharpness claim} asserted without the
count that would establish it.

\begin{longtable}{@{}p{0.20\textwidth} p{0.72\textwidth}@{}}
\caption{Editorial repairs, in document order.}\label{plt:tab:repairs}\\
\toprule
Location & Repair\\
\midrule
\endfirsthead
\toprule
Location & Repair\\
\midrule
\endhead
\cref{plt:prop:mot-blocks} & source draft 5 justifies the lexicographic order only in the case alpha > alpha', leaving the tie-breaking case alpha = alpha' unproved. The display above proves both cases.\\
\cref{plt:prop:mot-one-generator-fails} & the draft said X\textasciicircum \{1/2\} "arises from inverting a non-integer leading power". That is backwards, and contradicts the reading of the Fabius endpoint law in the remark on its status: a half-integer exponent is produced by inverting an *integer* leading power other than the first (X\textasciicircum 2 inverts to X\textasciicircum \{1/2\}); an inner argument that is already Puiseux has left the class before the inversion begins.\\
\cref{plt:prop:mot-one-generator-fails} & source draft 6 prints the refined scale as "..., X\textasciicircum \{-n\}L\textasciicircum 2, X\textasciicircum \{-n\}L, X\textasciicircum \{-n\}, X\textasciicircum \{-n\}L\textasciicircum \{-1\}, ...", which silently places negative powers of L inside the polynomial-logarithmic class. They are not there: coefficient blocks of K[L]((t)) are polynomials, so L-exponents are nonnegative, and L\textasciicircum \{-1\} belongs to the enlargements named above.\\
\cref{plt:cor:mot-both-generators-needed} & the draft called PLlau "the t-adic completion of the Laurent ring in (iii)". That is exactly the error item 7 of the normalizations remark warns against: t is a unit in K[t,t\textasciicircum \{-1\}][L], so every power of the ideal (t) is the whole ring, the t-adic filtration is constant, and the associated completion collapses to zero. The correct construction completes the polynomial subring K[L][t] and only then inverts t.\\
Why powers and logarithms belong in one  & the draft attached the exponentiated form to the full display, where it is false --- exponentiating \textbackslash eqref\{plt:eq:mot-generic-implicit\} exponentiates the higher blocks as well. The equivalence holds for the two-term equation only, and is stated for it.\\
\cref{plt:cor:mot-fabius-inverse} & the draft named the two parts of this estimate A and B, colliding with the B = log(1/y) of the corollary and the status remark that quote this proposition. The internal quantities are renamed A\_1 and A\_2; B is reserved for the inverse-endpoint variable throughout the subsection.\\
Why powers and logarithms belong in one  & source draft 1 asserts convergence "on the principal real branch for z >= e", supported only by a general citation. That threshold is not established anywhere in the six drafts, and the assertion also conflates a convergence claim with the separate, elementary observation that the series is exact at z = e. We keep the coefficient statement, proved below, and defer the convergence domain to Part VI with its own citation. the draft called the asymptotic display \textbackslash eqref\{plt:eq:mot-W-two-orders\} "exact at z = e". A display whose remainder is an analytic O as z -> +infty asserts nothing at a single point, so that formulation commits precisely the conflation the box above forbids. The correct elementary statement is a coefficient identity for the underlying formal series, and it is proved here.\\
Reader's guide and the structure of this & the draft called "Part V may be read directly after Part IV" an exception to the rule that each part depends only on those before it. It is not an exception at all --- V does come after IV --- and it also contradicts reading path 1 below, which does draw two statements from Parts II and III. The intended and correct content is the lightness of that dependency.\\
Reader's guide and the structure of this & source draft 2 tabulates the output class of differentiation as "t K[L]((t)) up to the leading power". That set equals K[L]((t)), because t is a unit there, so the row asserts nothing. The correct statement is the valuation inequality nu\_t(D\_X F) >= nu\_t(F) + 1, proved in Part II, which can be strict only when a constant leading block is annihilated.\\
Reader's guide and the structure of this & source draft 4 prints the coefficient condition of its defining display as "p\_n in K[X]", naming the large variable where the logarithmic one is meant. The intended and correct condition is p\_n in K[L].\\
The formal polynomial--logarithmic rings & the chain \textbackslash eqref\{plt:eq:ring-chain\} exhibits three inclusions, not four; the count was wrong as written.\\
\cref{plt:cor:ring-domains} & the first inclusion is not a change of coefficient ring (the printed chain \textbackslash (\textbackslash K[L]\textbackslash subset\textbackslash K[L]\textbackslash subset\textbackslash K(L)\textbackslash ) made it look like one); it is the inclusion of a power-series ring in the corresponding Laurent ring.\\
\cref{plt:cor:ring-domains} & this appeal was unsupported and mildly circular: the geometric-series corollary was stated only for coefficient rings whose list already contained \textbackslash (\textbackslash K((L\textasciicircum \{-1\}))\textbackslash ), the very object being constructed here. \textbackslash cref\{plt:cor:ring-geometric\} is now proved for an arbitrary commutative coefficient ring, so the instance \textbackslash (E=\textbackslash K\textbackslash ), generator \textbackslash (u\textbackslash ), is legitimate.\\
\cref{plt:def:ring-support} & S4 calls \textbackslash (t\textbackslash K[\textbackslash ell][[t]]\textbackslash ) ``the strictly small ideal'' of the Laurent ring \textbackslash (\textbackslash K[\textbackslash ell]((t))\textbackslash ); it is neither small in that ring nor an ideal of it.\\
\cref{plt:rmk:ring-fixed-exponents} & the statement imposed the standing hypothesis \textbackslash (\textbackslash alpha>\textbackslash alpha'\textbackslash ) and then appended the case \textbackslash (\textbackslash alpha=\textbackslash alpha'\textbackslash ), which that hypothesis excludes; the two cases are now alternatives of one hypothesis. The regime \textbackslash (X>\textbackslash EulerE\textbackslash ) is printed here rather than produced inside the proof.\\
\cref{plt:rmk:ring-fixed-exponents} & the sources state this without the hypothesis \textbackslash (X>\textbackslash EulerE\textbackslash ); for \textbackslash (1<X<\textbackslash EulerE\textbackslash ) one has \textbackslash (0<\textbackslash log X<1\textbackslash ) and high logarithmic powers are small, so the comparison is genuinely eventual.\\
Valuation, leading data, and the degree  & ``\textbackslash (\textbackslash ordt\textbackslash ) is a discrete valuation on \textbackslash (R\textbackslash ) with valuation ring \textbackslash (E[[t]]\textbackslash )'' overstates the case for two of the four classes: \textbackslash (\textbackslash PLpow\textbackslash ) and \textbackslash (\textbackslash PLlau\textbackslash ) are not fields, and on \textbackslash (\textbackslash PLpow\textbackslash ) the image of \textbackslash (\textbackslash ordt\textbackslash ) is \textbackslash (\textbackslash NaturalNumbers\textbackslash ), not \textbackslash (\textbackslash IntegerNumbers\textbackslash ). The claim is split accordingly.\\
\cref{plt:ex:lead-strictness} & the domain hypothesis enters only through the reverse inequality in (i). For an arbitrary commutative ring \textbackslash (E\textbackslash ) the subadditivity \textbackslash (\textbackslash ordt(FG)\textbackslash ge\textbackslash ordt F+\textbackslash ordt G\textbackslash ), together with (ii) and (iii), still holds in \textbackslash (E[[t]]\textbackslash ) and \textbackslash (E((t))\textbackslash ), by the first half of the computation in (i); this weaker form is what \textbackslash cref\{plt:cor:ring-geometric\} uses, and stating it here is what keeps that corollary free of a circular appeal. S1 and S5 record the strictness of (ii) as happening ``exactly when the leading blocks cancel'', omitting the hypothesis \textbackslash (\textbackslash ordt F=\textbackslash ordt G\textbackslash ); with unequal valuations no cancellation is possible and equality always holds.\\
\cref{plt:rmk:lead-three-objects} & the display ``\textbackslash (p\_\textbackslash nu=aL\textasciicircum d+\textbackslash text\{lower powers of \}L\textbackslash )'' is literally a polynomial expansion, but \textbackslash (R\textbackslash ) here ranges over all four classes, where a block need not be a polynomial. The reading is fixed explicitly.\\
Valuation, leading data, and the degree  & ``equivalently'' was wrong: additivity of \textbackslash (v\textbackslash ) is equivalent to multiplicativity of \textbackslash (\textbackslash lm\textbackslash ) alone; multiplicativity of \textbackslash (\textbackslash lc\textbackslash ) is a further assertion (proved below), not a restatement.\\
\cref{plt:rmk:lead-coarse} & ``with value group \textbackslash (\textbackslash Gamma\textbackslash )'' is false as stated for two of the four classes: on \textbackslash (\textbackslash PLpow\textbackslash ) and \textbackslash (\textbackslash PLlau\textbackslash ) a block is a polynomial, so \textbackslash (\textbackslash degL\textbackslash ge0\textbackslash ) and the image of \textbackslash (v\textbackslash ) never meets \textbackslash (\textbackslash IntegerNumbers\textbackslash times \textbackslash PositiveIntegers\textbackslash ). The image and the group it generates are separated.\\
\cref{plt:def:lead-profile} & S2 states the prefactor form \textbackslash (A=cX\textasciicircum \{\textbackslash rho\}L\textasciicircum \{\textbackslash beta\}(1+U)\textbackslash ) ``whenever available''; the example \textbackslash (L+t\textbackslash ) shows that over \textbackslash (\textbackslash K[L]\textbackslash ) it is available only after passing to \textbackslash (\textbackslash K(L)\textbackslash ), so no argument may assume it inside \textbackslash (\textbackslash PLlau\textbackslash ).\\
\cref{plt:ex:lead-slope-fails} & the monotonicity clause \textbackslash (\textbackslash mathcal D\_\textbackslash sigma\textbackslash subseteq\textbackslash mathcal D\_\{\textbackslash sigma'\}\textbackslash ) was asserted without proof, and it is not the trivial inequality \textbackslash (\textbackslash sigma n\textbackslash le\textbackslash sigma'n\textbackslash ): that fails for \textbackslash (n<0\textbackslash ). The argument below uses that \textbackslash (\textbackslash dprofile\{F\}\textbackslash ) is \textbackslash (-\textbackslash infty\textbackslash ) below \textbackslash (\textbackslash ordt F\textbackslash ).\\
\cref{plt:ex:lead-slope-fails} & S2 states \textbackslash (\textbackslash deg c\_n\textbackslash le A+B+\textbackslash max(\textbackslash alpha,\textbackslash beta)n\textbackslash ) for the product. That form is correct only when the factor of smaller slope has nonnegative outer order; the counterexample of \textbackslash cref\{plt:ex:lead-slope-fails\} violates it, and \textbackslash eqref\{plt:eq:lead-slope-bound\} is the corrected statement, attained there with equality.\\
\cref{plt:prop:trunc-basic} & S6 uses the inclusive marker throughout; the catalogue makes the exclusive operator normative, so every transferred formula is shifted by one block.\\
\cref{plt:rmk:trunc-precision} & S5 states \textbackslash eqref\{plt:eq:trunc-product-pow\} without the hypothesis \textbackslash (\textbackslash ordt\textbackslash ge0\textbackslash ); on \textbackslash (\textbackslash PLlau\textbackslash ) it is false, and \textbackslash eqref\{plt:eq:trunc-product\} is the correct general form.\\
\cref{plt:rmk:trunc-no-size} & \textbackslash cref\{plt:def:trunc-formal-o\} defines the token only for an integer cutoff, so (ii) and (iii) below were type-incorrect when \textbackslash (H\textbackslash ) or \textbackslash (G\textbackslash ) vanishes and the printed exponent becomes \textbackslash (+\textbackslash infty\textbackslash ). The reading of that degenerate case is now fixed explicitly.\\
\cref{plt:prop:trunc-not-ideal} & the displayed analytic statement was printed with none of \textbackslash (f\textbackslash ), \textbackslash (S\textbackslash ) or \textbackslash (n\_0\textbackslash ) introduced, contrary to the standing requirement that a statement name its regime and domain. They are named here.\\
\cref{plt:prop:trunc-not-ideal} & S1 offers instead the minimal Poincar\textbackslash 'e requirement that the remainder be little-o of the last retained term. That condition does not determine the blocks: with \textbackslash (p\_n=L\textasciicircum 5\textbackslash ) and \textbackslash (q\_n=L\textasciicircum 5+L\textasciicircum 3\textbackslash ) the two candidate blocks differ by \textbackslash (L\textasciicircum 3=\textbackslash LittleOAt\{L\textasciicircum 5\}\{X\textbackslash to+\textbackslash infty\}\textbackslash ), so uniqueness fails. The displayed form, and S6's \textbackslash (\textbackslash LittleOAt\{X\textasciicircum \{-N\}\}\{X\textbackslash to+\textbackslash infty\}\textbackslash ) form, both do determine them.\\
\cref{plt:def:trunc-extraction} & the label kind "box" is not in the normative list of kinds; these boxed statements are remarks, so each label is renamed to the "rmk" kind.\\
\cref{plt:ex:trunc-log-cap} & the label kind "box" is not in the normative list of kinds; these boxed statements are remarks, so each label is renamed to the "rmk" kind.\\
\cref{plt:ex:trunc-not-cauchy} & S2's completeness proof asserts the common lower bound because ``the sequence is already Laurent-bounded at some initial precision''. That is not a hypothesis of the statement; the bound must be extracted, as above, from the Cauchy condition at cutoff \textbackslash (N=0\textbackslash ).\\
\cref{plt:ex:trunc-not-cauchy} & the word ``exactly'' in (iv) asserts that the common lower bound is also necessary; the sources supply only sufficiency, so the converse is proved here.\\
\cref{plt:rmk:trunc-separated} & S3 asserts that \textbackslash (t\textbackslash )-adic convergence is ``equivalent'' to coefficientwise stabilization and that both the power-series and the Laurent ring are complete ``in this sense''. The equivalence is false on the Laurent ring, as follows; the completeness assertion itself is correct, but needs the proof of \textbackslash cref\{plt:thm:trunc-completeness\}\textbackslash ,(iii).\\
\cref{plt:cor:ring-geometric} & the identification of \textbackslash (\textbackslash mu\textbackslash ) was printed as an equality. It is only a lower bound: if every \textbackslash (\textbackslash ordt S\_i>0\textbackslash ) the right-hand side is \textbackslash (0\textbackslash ) while \textbackslash (\textbackslash mu>0\textbackslash ). Finiteness is all the argument needs, and all it gives.\\
\cref{plt:thm:trunc-summability-criterion} & the corollary was stated only for \textbackslash (E\textbackslash in\textbackslash SetOf\{\textbackslash K[L],\textbackslash K(L),\textbackslash K((L\textasciicircum \{-1\}))\}\textbackslash ), yet it is invoked inside \textbackslash cref\{plt:prop:ring-chain\} with \textbackslash (E=\textbackslash K\textbackslash ) and generator \textbackslash (L\textasciicircum \{-1\}\textbackslash ), precisely in order to prove that \textbackslash (\textbackslash K((L\textasciicircum \{-1\}))\textbackslash ) is a field --- the fact that puts \textbackslash (\textbackslash K((L\textasciicircum \{-1\}))\textbackslash ) on that list. Stating it for an arbitrary commutative coefficient ring removes the circle at no cost: the proof never needs a domain hypothesis.\\
\cref{plt:rmk:trunc-bridge} & the transfer was claimed only for the four ambient classes, but \textbackslash cref\{plt:cor:ring-geometric\} needs it for \textbackslash (E((t))\textbackslash ) with \textbackslash (E\textbackslash ) an arbitrary commutative ring (the case \textbackslash (E=\textbackslash K\textbackslash ), generator \textbackslash (L\textasciicircum \{-1\}\textbackslash ), is used in \textbackslash cref\{plt:prop:ring-chain\}). The scope of the transfer is widened, and the two places where the domain hypothesis entered are named.\\
\cref{plt:rmk:trunc-bridge} & the label kind "box" is not in the normative list of kinds; these boxed statements are remarks, so each label is renamed to the "rmk" kind.\\
\cref{plt:rmk:trunc-bridge} & the volume regime is \textbackslash (X\textbackslash to+\textbackslash infty\textbackslash ) on the positive ray; the generic row printed the weaker \textbackslash (X\textbackslash to\textbackslash infty\textbackslash ).\\
\cref{plt:prop:mul-monomial} & the four classes were displayed on one line, which overfills the text block by about 45pt; set as a two-by-two array instead.\\
\cref{plt:cor:mul-truncated-product} & the bracketed title duplicated that of \textbackslash cref\{plt:prop:trunc-product\}. The two statements are not the same result: the ring-chapter proposition is the law on \textbackslash (\textbackslash PLlau\textbackslash ) at the canonical cutoffs, this one is the law over an arbitrary logarithmically graded coefficient ring at arbitrary sufficient cutoffs, together with optimality and sharpness. The titles now say so; see \textbackslash cref\{plt:rmk:mul-leading-restated\}.\\
Multiplication: convolution of logarithm & the optimality of \$(N-b,N-a)\$ was asserted unconditionally. For \$N\textbackslash le a+b\$ both sides of \textbackslash eqref\{plt:eq:mul-truncated-product\} vanish whatever \$P,Q\$ are, so no pair is forced and there is nothing to be optimal about; the clause is now stated under \$N>a+b\$, which is what the argument actually establishes.\\
\cref{plt:cor:mul-error-rule} & the source bounds the third term by \$P+Q\$ and asserts \$P+Q\textbackslash ge N\$. That inequality does not follow from \$P\textbackslash ge N-b\$ and \$Q\textbackslash ge N-a\$: for \$a=b=0\$ and \$N=P=Q=-5\$ one has \$P+Q=-10<N\$. The repair is to bound \$\textbackslash ordt H\$ by \$b\$ rather than by \$Q\$, which is legitimate because every block of \$H\$ is a block of \$G\$.\\
\cref{plt:cor:mul-error-rule} & S6, Theorem "Closure and finite determination", eq. (mult-error-rule). The source states the conclusion as O(t\textasciicircum \{min(p+m\_0, q+n\_0)\}) without any hypothesis relating p,q to the valuations n\_0,m\_0, and its proof silently bounds ord(E*Gtilde) by p+m\_0 and ord(Ftilde*H) by q+n\_0, which requires ord(Gtilde) >= m\_0 and ord(Ftilde) >= n\_0. Neither follows from the stated hypotheses. The corrected statement carries the third entry p+q, and the hypotheses p >= n\_0, q >= m\_0 under which the source's form is recovered. Note that these hypotheses are sufficient but not necessary: the third entry is redundant exactly when p >= a or q >= b.\\
\cref{plt:def:mul-leading-data} & retitled to distinguish it from \textbackslash cref\{plt:def:lead-leading-data\}, which is the same data on the four ambient classes; only the coefficient ring differs. See \textbackslash cref\{plt:rmk:mul-leading-restated\}.\\
\cref{plt:rmk:mul-leading-restated} & the leading data, the rank-two valuation, the log-degree profile, the exclusive truncation law and the summability calculus are all introduced in the ring chapter. Restating them here without saying so would leave two definitions of the same symbols standing in one volume. This remark records that the restatements are verbatim generalizations over the coefficient ring and agree with the originals.\\
\cref{plt:thm:mul-leading-law} & the first equivalence of the display was left unargued. It is not a theorem of this chapter at all: \$\textbackslash succdom\$ is defined by the right-hand condition in \textbackslash cref\{plt:def:ring-dominance\}, so only the analytic equivalence needs an argument, and that argument restates \textbackslash cref\{plt:lem:ring-dominance-analytic\} for integer exponents.\\
\cref{plt:cor:mul-unit-necessary} & the statement asserts that neither \$\textbackslash PLpow\$ nor \$\textbackslash PLlau\$ is a field, but the argument as given covers only \$\textbackslash PLlau\$. One line closes it.\\
\cref{plt:def:mul-profile} & retitled to distinguish it from \textbackslash cref\{plt:def:lead-profile\}, which is the same function on \textbackslash (\textbackslash PLlau\textbackslash ), where the blocks are polynomials and the codomain is \textbackslash (\textbackslash NaturalNumbers\textbackslash cup\textbackslash SetOf\{-\textbackslash infty\}\textbackslash ). See \textbackslash cref\{plt:rmk:mul-leading-restated\}.\\
\cref{plt:cor:mul-degree-equality} & the criterion was stated as an unconditional ``if and only if''. It fails in the degenerate case \$P\_N=\textbackslash emptyset\$ (equivalently \$M=-\textbackslash infty\$): there \$[t\textasciicircum N](FG)=0\$, so equality holds although the empty sum \textbackslash eqref\{plt:eq:mul-top-coefficient\} is zero. The criterion is now confined to \$M\$ finite and the degenerate case is recorded separately.\\
\cref{plt:cor:mul-degree-equality} & S6, Proposition "Degree and support bounds", asserts equality "whenever the maximizing blocks are nonzero" under the hypothesis that all coefficients are nonnegative over an ordered field. That hypothesis is far stronger than necessary and the quoted sufficient condition is stated ambiguously (a zero block has degree -infinity and cannot maximize at all). The criterion below is exact and the positivity hypothesis has been reduced to the leading coefficients of the maximizing blocks only.\\
\cref{plt:ex:mul-degree-sharp} & condition (3) as first drafted omitted \$P\_N\textbackslash ne\textbackslash emptyset\$; over an ordered field its hypotheses are satisfied vacuously when \$P\_N\$ is empty, and the conclusion would then be the false assertion that the sum \textbackslash eqref\{plt:eq:mul-top-coefficient\} is nonzero.\\
\cref{plt:rmk:mul-profile-not-closed} & the inclusion \$\textbackslash mathfrak A\_\textbackslash alpha\textbackslash subseteq\textbackslash mathfrak A\_\{\textbackslash alpha''\}\$ for \$\textbackslash alpha\textbackslash le\textbackslash alpha''\$ was used without proof. On the Laurent ring it is not immediate: for \$n<0\$ a larger slope gives a \textbackslash emph\{smaller\} bound, so the constant has to be increased, and that is possible only because the support is bounded below. The argument is supplied here.\\
\cref{plt:def:mul-summable} & retitled, and its relation to \textbackslash cref\{plt:def:trunc-summable\} made explicit: the condition is verbatim the one of the ring chapter, read over the coefficient ring \$R\$ instead of \$\textbackslash K[L]\$, which \textbackslash cref\{plt:rmk:trunc-summability-classes\} already licenses. Only the two clauses of \textbackslash cref\{plt:lem:mul-summable\} are used below, so they are the only ones recorded here.\\
\cref{plt:rmk:mul-summable-vs-cauchy} & the source calls the two infima ``finite by (1)''; part\textasciitilde (1) gives only \$>-\textbackslash infty\$, and the value \$+\textbackslash infty\$ does occur, namely when every member of the family is zero. That case is disposed of separately.\\
\cref{plt:lem:mul-asymptotic-basic} & retitled: ``Basic properties of a scale'' in the differential part (\textbackslash cref\{plt:lem:dif-scale-basic\}) is about the coefficient field of a power--log scale, not about the relation \$\textbackslash approx\$ between a function and a series. The two lemmas share nothing but the word ``basic''.\\
\cref{plt:rmk:mul-analytic-hypotheses} & the three-term estimate was set on one display line, overfilling the text block by about 53pt; broken across two lines.\\
\cref{plt:cor:div-unit-group} & S5 and S6 assert the nonunit direction by citing the general power-series unit criterion; the additivity step that makes it valid on the Laurent ring is supplied here explicitly.\\
\cref{plt:ex:div-nonunit} & the direct-product step was justified by "the three subgroups pairwise intersect trivially", which does not imply an internal direct product for three subgroups (the Klein four-group refutes it); the isomorphism is read off instead from surjectivity plus the uniqueness just proved.\\
\cref{plt:cor:div-quotient} & \textbackslash cref\{plt:cor:div-quotient\} invokes this theorem "over the coefficient ring E", a generality the statement did not record; the proof uses only commutativity and invertibility of B\_0, so the clause is supplied here.\\
\cref{plt:rmk:div-truncation-convention} & S6 states this correctness theorem with its local inclusive truncation [F]\_\{<=N\}; the statement above is the exclusive normalisation demanded by the catalogue, and the two differ by exactly one block.\\
\cref{plt:rmk:div-guard} & S6 states a "guard precision" principle in words ("compute the quotient beyond the final requested block") but never says by how much; the exact guard is the content of the proposition above.\\
\cref{plt:ex:div-reciprocal-worked} & the negative-guard case was explained by "the denominator has a Laurent principal part", which accounts only for the numerator guard b; the denominator guard 2b-a is negative under the unrelated condition a > 2b.\\
\cref{plt:prop:div-degree-profile} & S1 prints the block Q\_3 of this example immediately after inputs given only modulo t\textasciicircum 3; the hypothesis that the omitted blocks vanish is a genuine extra assumption and is stated here.\\
\cref{plt:thm:div-Ett-field} & pointwise finiteness alone does not place the sum in the field: one must also check that each block is a Laurent series in L\textasciicircum \{-1\}, i.e. that its support is bounded below. The same bound gives this at once, and no source records the step.\\
\cref{plt:cor:div-newton-exact} & the input-requirement statement (4) is absent from S2, S3, S5 and S6, each of which asserts doubling without saying which truncations suffice.\\
\cref{plt:cor:div-newton-exact} & the precondition N >= 1 was not printed; with N = 0 the loop body never runs and the routine returns the scalar B\_0\textasciicircum \{-1\} rather than Trunc\_\{<0\}(1/B) = 0.\\
\cref{plt:rmk:div-valuation-additive} & the induction covers only the exact powers of two, while \textbackslash cref\{plt:alg:div-newton\} caps every step at N; the capped step needs its own line before the "in particular" clause is justified.\\
\cref{plt:rmk:div-valuation-additive} & S2 justifies the doubling by "the valuation is multiplicative"; the valuation is additive on products, which is what the argument uses.\\
\cref{plt:prop:div-cost-naive} & the convention "multiplications by the scalar B\_0\textasciicircum \{-1\} are not counted" was stated loosely enough to exclude also the products B\_j C\_0 and C\_0 E\_j, under which reading the two schemes cost exactly the same and the N-1 gap claimed below disappears. The counted set is pinned down here and the other reading is recorded at the end of \textbackslash cref\{plt:prop:div-cost-newton\}.\\
\cref{plt:cor:div-doubling-pays} & the N-1 gap is fragile: it is created entirely by charging for products against the scalar block C\_0. Recording this keeps the comparison from being read as a robust separation, which it is not.\\
\cref{plt:cor:div-doubling-pays} & S3 asserts that Newton inversion "is often faster than coefficient recursion for large N even with naive multiplication"; the exact counts above show the opposite at the level of coefficient multiplications, and the honest statement is the corollary. the corollary was headed by the iff "improves ... exactly when block multiplication is subquadratic". Only the two implications below are proved (subquadratic wins asymptotically; schoolbook loses by N-1), and neither the converse nor the intermediate regime M(N) = cN\textasciicircum 2 is covered; the claim is downgraded to what the two cost counts give.\\
\cref{plt:rmk:div-scalar-cost} & S6 gives the cost as O(N\textasciicircum 2 d\textasciicircum 2) "if the typical coefficient degree is d"; under the affine log-degree profile that actually occurs there is no such constant d, and the honest count is quartic in N.\\
Units, reciprocals, and division & the catalogue asserts that the two iteration orders differ "in general"; no source proves it, and the witnesses above supply the proof.\\
Units, reciprocals, and division & "beats the recurrence only when block multiplication is subquadratic" asserts the unproved converse; only the two extreme regimes are established.\\
\cref{plt:lem:dif-summable} & retitled. Three ``Summable families'' statements stood in the volume with the same bracketed title. This one is the blockwise assembly principle for a sequence in a Laurent ring over any of the coefficient rings used here; \textbackslash cref\{plt:lem:cmp-summable\} is the corresponding statement for an arbitrary index set in a Hahn ring with real exponents, and \textbackslash cref\{plt:thm:trunc-summability\} is the full calculus on \textbackslash (\textbackslash PLlau\textbackslash ). On \textbackslash (\textbackslash PLlau\textbackslash ) the present lemma is the sequence case of \textbackslash cref\{plt:thm:trunc-summability\}\textbackslash ,(i),(iv),(v); it is stated again because it is applied below in the \textbackslash (L\textasciicircum \{-1\}\textbackslash )-adic order and in a Hahn ambient, neither of which that theorem covers. The hypothesis \textbackslash (\textbackslash inf\_\{r\}\textbackslash ordt F\_\{r\}>-\textbackslash infty\textbackslash ) is redundant for a sequence with \textbackslash (\textbackslash ordt F\_\{r\}\textbackslash to+\textbackslash infty\textbackslash ) and is kept only because the proof uses it.\\
\cref{plt:prop:dif-substitution} & retitled. This is substitution of a small series into a \textbackslash emph\{scalar\} power series, \textbackslash (\textbackslash K[[z]]\textbackslash to\textbackslash PLpow\textbackslash ). The composition part has a different theorem under the same name, \textbackslash cref\{plt:thm:cmp-substitution-hom\}, which substitutes a monomial-leading transseries into a transseries; neither implies the other.\\
\cref{plt:thm:dif-finite-determination} & retitled: the composition part states a different theorem under the same name (\textbackslash cref\{plt:thm:cmp-finite-determination\}, finite determination of \textbackslash (F\textbackslash compose G\textbackslash )). The present one is about \textbackslash (\textbackslash Phi(U)\textbackslash ) with \textbackslash (\textbackslash Phi\textbackslash ) a scalar power series.\\
\cref{plt:def:dif-exp-extension} & S3 states the analogous expansion with U\textasciicircum 4 = O(t\textasciicircum 4) as the reason for stopping at r = 3; the correct reason is that r*q >= N kills the term, which is the general statement proved above and matters as soon as q > 1.\\
\cref{plt:lem:dif-scale-basic} & retitled to separate it from \textbackslash cref\{plt:lem:mul-asymptotic-basic\}, which collects the basic properties of the relation \textbackslash (\textbackslash approx\textbackslash ) between a function and a series.\\
\cref{plt:cor:dif-exp-three-ways} & S4 and S6 note only that a "large part" other than an integral multiple of L may require new monomials. The decisive point, which no source states, is that a logarithmic derivative has positive L\textasciicircum \{-1\}-order, whereas the derivative of a nonaffine polynomial has order at most -1.\\
\cref{plt:lem:dif-image-laurent} & no source identifies the image of \textbackslash partial\_L on Laurent-log coefficients, although S5 and S6 both assert closure of the class under antiderivatives and then remark vaguely that "new logarithmic levels can be forced" once inverse powers of L are admitted. The converse of the preceding lemma is what makes that remark precise, so it is proved here and used both for the logarithm boundary and for the antiderivative obstruction.\\
\cref{plt:cor:dif-log-unit} & S6 states this example with a leading factor L\textasciicircum \{-2\}, which is not an element of K[L]((t)); the phenomenon is the same, and the general statement is that the coefficient of Lambda is -omega of the leading coefficient.\\
\cref{plt:prop:dif-block} & the notation catalogue and four of the six sources assert that D\_X is "the unique derivation" with the two generator values, with no hypothesis. The argument above needs a bound on delta on the power-series part; without one, the passage from the dense subring to the completion is not available.\\
\cref{plt:thm:dif-repeated} & notation reconciliation. Two sources write Delta\_n for the single factor partial\_L - n and never introduce the second index; the normative two-index operator is used here throughout.\\
\cref{plt:cor:dif-budget} & S1 prints the identity as X\textasciicircum k d\textasciicircum k/dX\textasciicircum k P(log X) = sum\_j s(k,j) P\textasciicircum \{(j)\} with the sum starting at j = 1; that is correct only for k >= 1, since s(0,0) = 1. The range j = 0..k with the empty-product convention is uniform.\\
\cref{plt:rmk:dif-resonance-restated} & the motivation part proves the same dichotomy in the weaker form it needs there (\textbackslash cref\{plt:lem:mot-block-antiderivative\}, on \textbackslash (\textbackslash K[L]\textbackslash ) only, with no closed form). Rather than restate it, the two lemmas above are identified with it here, so that a reader meets the resonance once as a mechanism and once in the form the antiderivative theorem consumes.\\
\cref{plt:ex:dif-antiderivative} & S6 concludes only that "the polynomial-logarithmic ring is closed under termwise antiderivatives, up to an arbitrary constant", and S3 writes the weaker (and redundant) inclusion int P dX subset P + K. The correct statement is surjectivity of D\_X on K[L]((t)) with kernel K, which is what the block system actually proves and what is needed downstream.\\
\cref{plt:lem:cmp-summable} & retitled. This is the summability criterion in a Hahn ring with real exponents, for an arbitrary index set, together with the notion of a map compatible with summable families that the rest of the part uses. It is not \textbackslash cref\{plt:lem:dif-summable\} (sequences in a Laurent ring) and not \textbackslash cref\{plt:thm:trunc-summability\} (the calculus on \textbackslash (\textbackslash PLlau\textbackslash )); for \textbackslash (\textbackslash Gamma=\textbackslash IntegerNumbers\textbackslash ) it restricts to the latter.\\
\cref{plt:rmk:cmp-lower-bound-automatic} & the catalogue warns that a t-adically Cauchy *sequence* in a Laurent ring may fail to converge without a common lower support bound. For *families* satisfying (plt:eq:cmp-summability-criterion) the bound is automatic, as the proof shows; the two statements are not in conflict.\\
\cref{plt:thm:cmp-substitution-hom} & retitled to separate it from \textbackslash cref\{plt:prop:dif-substitution\}, which substitutes a small series into a scalar power series. The present theorem substitutes a monomial-leading transseries into a transseries.\\
\cref{plt:cor:cmp-pure-power} & the draft deduced "compatible with summable families" from (plt:lem:cmp-summable) directly, but that lemma only supplies the valuation half of the definition; the interchange Phi\_G(sum S\_i) = sum Phi\_G(S\_i) needs the argument below. It is used by (plt:lem:tay-determined) and by the associativity proof, so it is supplied here.\\
\cref{plt:thm:cmp-finite-determination} & retitled: \textbackslash cref\{plt:thm:dif-finite-determination\} carries the same name for a different statement, the finite determination of \textbackslash (\textbackslash Phi(U)\textbackslash ) with \textbackslash (\textbackslash Phi\textbackslash ) a scalar power series.\\
\cref{plt:prop:cmp-block-recurrence} & S2 writes the operator in (its) power-leading formula as ":e\textasciicircum \{delta(D-n)\}:" with D previously declared to be d/dL, while the identity requires differentiation in the *argument* of the coefficient polynomial. We print partial\_z for the dummy variable and evaluate afterwards. an earlier draft of this warningbox described the discrepancy as "a spurious factor rho\textasciicircum k". That is false as soon as n != 0: the chain rule turns (d\_z - n)\textasciicircum k into (rho d\_z - n)\textasciicircum k, and the two are proportional only when n = 0. Corrected below, with a numerical witness.\\
\cref{plt:prop:cmp-near-monoid} & the sources state the rule in words; the numerical demonstration of the two standard mistakes is supplied here so that the size of the error is visible.\\
\cref{plt:lem:cmp-realization-basic} & the exclusive normalization is used throughout; a source using the inclusive cutoff states the same condition as f - sum\_\{n<=N\} = o(X\textasciicircum \{-N\}), which is (plt:eq:cmp-realization) at N+1.\\
\cref{plt:thm:cmp-analytic-transfer} & the draft omitted this reduction and then asserted that the product of the two remainders is O(X\textasciicircum \{-N+eps\}); that product is only O(X\textasciicircum \{-2N+a+atilde+eps/2\}), which is the required bound precisely when N >= a + atilde. The reduction below makes the omission harmless.\\
The formal Taylor formula & as in (plt:thm:cmp-substitution-hom), the earlier text checked the interchange only for the family of blocks of a single F; the definition of compatibility quantifies over all summable families, and that is what (plt:lem:tay-determined) consumes. Argument supplied.\\
\cref{plt:lem:tay-determined} & this multiplicativity is exactly the step that S3 asserts ("both sides are continuous algebra homomorphisms") and that S6 replaces by "the equality then extends to polynomials, products, inverses"; it is not automatic, because T\_H is a *normal-ordered* exponential and not exp(H D\_X).\\
\cref{plt:thm:tay-coefficients} & S4 and S5 prove the Taylor formula by "inserting a scalar parameter s, forming F(x+sA), and differentiating in s". That argument is circular: the object F(x+sA) is defined only by the composition being constructed, and "formal differentiation in s" of it presupposes the chain rule for that composition. The homomorphism argument above avoids the circle.\\
\cref{plt:prop:tay-convolution} & the two index orders in the corpus are transposes of each other; the collision is silent because both symbols are two-index arrays of polynomials.\\
\cref{plt:ex:tay-four-block} & the source that supplies this example specifies F only modulo Y\textasciicircum \{-3\} and G only modulo t\textasciicircum 2, yet quotes h\_3, which depends on the coefficient of Y\textasciicircum \{-3\} in F and on the coefficients of t\textasciicircum 2, t\textasciicircum 3 in G. The data are restated here as exact finite series, which is what the quoted answer actually assumes.\\
\cref{plt:prop:tay-scalar-outer} & this proposition restated three results of the elementary-series part under a different normalisation of the outer coefficients: the substitution homomorphism (\textbackslash cref\{plt:prop:dif-substitution\}), its finite determination (\textbackslash cref\{plt:thm:dif-finite-determination\}), and the Bell-polynomial coefficient formula \textbackslash eqref\{plt:eq:dif-bell\} of \textbackslash cref\{plt:prop:dif-bell\}, which the displayed formula reproduces verbatim once \textbackslash (c\_\{k\}=\textbackslash varphi\_\{k\}/k!\textbackslash ) is substituted. The restatement and its proof are deleted in favour of the pointer below; the label is kept so that \textbackslash cref\{plt:ex:tay-standard-outer\} and any citation from elsewhere in the volume continue to resolve. Only the translation clause was genuinely new, and it is retained with its one-line proof, as is the second route to the coefficient formula.\\
\cref{plt:lem:grp-explog} & S6 writes \$F/X = 1+tA = 1+o(1)\$ "formally". A little-o symbol with no regime attached is not a formal statement; the three O-roles must not be conflated. The correct formal assertion is \$tA \textbackslash in t\textbackslash PLpow\$, equivalently \$1+tA-1=\textbackslash FormalBigOAt\{t\}\{t\}\$, equivalently \$tA\textbackslash precdom1\$.\\
\cref{plt:lem:grp-explog} & this lemma was stated and proved twice in the volume, here and as \textbackslash cref\{plt:lem:cmp-exp-log\} in the composition part, with the same transport-of-identities proof. The composition-part statement is the stronger of the two --- it is set in an arbitrary polynomial-logarithmic Hahn ring, and it adds the additivity of \textbackslash (\textbackslash exp\textbackslash ), the law \textbackslash (\textbackslash log(P\textasciicircum \{n\})=n\textbackslash log P\textbackslash ) for all \textbackslash (n\textbackslash in\textbackslash IntegerNumbers\textbackslash ), and the commutation with any homomorphism compatible with summable families, all of which are used below. The restatement and its proof are deleted and replaced by this pointer; the label is kept so that the citations to it elsewhere in the volume continue to resolve.\\
\cref{plt:thm:grp-taylor} & the exact formal Taylor formula was stated and proved twice, here and as \textbackslash cref\{plt:thm:tay-formula\} in the composition part, where an entire subsection (the Taylor operator, its summability budget, its multiplicativity, and the generator-determination lemma) is built to prove it. The statements are identical, so the restatement and its proof are deleted in favour of a pointer; the label is kept because the reversion part and the algorithms part cite it. The repair recorded here --- S3 concludes from agreement on the generators without proving that the normally ordered operator is multiplicative --- is carried out at the surviving site, in \textbackslash cref\{plt:lem:tay-homomorphism\}.\\
\cref{plt:cor:grp-gain} & the Lipschitz bound was printed as \textbackslash (\textbackslash ordt(B-C)+1\textbackslash ), dropping the summand \textbackslash (\textbackslash ordt(A)\textbackslash ) that the proof below already produces. The sharper form is what the reversion part needs (it is what makes the Newton estimate of \textbackslash cref\{plt:thm:rev-newton-doubling\} come out), and it was stated there a second time as a separate lemma. The inequality is corrected here and the restatement in the reversion part is deleted; every use of the weaker form below is unaffected, since \textbackslash (\textbackslash ordt(A)\textbackslash ge0\textbackslash ) on \textbackslash (\textbackslash PLpow\textbackslash ).\\
\cref{plt:thm:grp-group} & S4 proves the contraction step with a formal "mean value" identity A o (X+B) - A o (X+C) = (B-C)\textbackslash int\_0\textasciicircum 1 A' o (X+C+s(B-C)) ds. The integrand is a composition with a series carrying an auxiliary parameter s, and neither its existence nor the legitimacy of the term-by-term integration is established there. Corollary \textbackslash ref\{plt:cor:grp-gain\} replaces it by the exact Taylor difference, which needs nothing beyond Theorem \textbackslash ref\{plt:thm:grp-taylor\}.\\
\cref{plt:ex:grp-noncommutative} & S5 proves associativity by "introducing finite truncations mod t\textasciicircum N", for which "ordinary substitution is associative". Truncation is not a ring retraction compatible with substitution, and the claim that a truncated composite equals the truncation of the composite is precisely what needs proof. S1 proposes the correct route (check on the generators) but does not carry out the verification for L; the display above is that verification.\\
\cref{plt:ex:grp-noncommutative} & the Picard construction of the inverse was carried out twice in the volume: here, and again --- with the sharp rate \textbackslash (\textbackslash ordt(B\_\{j+1\}-B\_\{j\})\textbackslash ge(j+1)\textbackslash alpha+j\textbackslash ), the first-order form of the inverse, and the uniqueness of the two-sided inverse --- as \textbackslash cref\{plt:thm:rev-existence\}. The two constructions are the same iteration; the reversion part proves strictly more. The duplicate is deleted and item (4) is derived from the reversion theorem, which makes the dependency of the group's inverses on Part\textasciitilde IV explicit instead of hiding it behind two independent proofs. There is no circularity: \textbackslash cref\{plt:thm:rev-existence\} uses items (1)--(3) of the present theorem and the gain estimate \textbackslash cref\{plt:cor:grp-gain\}, and nothing else from this chapter.\\
\cref{plt:rmk:grp-anti} & this remark previously asserted that the section is independent of Part\textasciitilde IV and that inverses are established here. That was true only because the same Picard construction was written out twice. With the duplicate removed, the dependency is one-directional and is stated as such.\\
\cref{plt:prop:grp-limit} & S6 states the two continuity relations for the contact filtration but justifies the right-composition one by "H' = 1 + O(t)". That is the reason for a gain in the *left* argument; the right-hand estimate is the exact isometry of Proposition \textbackslash ref\{plt:prop:grp-substitution\}(4) and needs no derivative hypothesis at all.\\
\cref{plt:prop:grp-graded} & S6 calls the group "pro-unipotent at finite truncation". No algebraic group structure is available here (the graded pieces are the infinite dimensional space K[L]), and none is needed: Propositions \textbackslash ref\{plt:prop:grp-limit\} and \textbackslash ref\{plt:prop:grp-commutator\} give the exact provable statement, namely that Gnear is the inverse limit of nilpotent quotients.\\
\cref{plt:rmk:grp-normal-ordering} & S6 introduces log(S\_F) as a power series in S\_F - I and writes the fractional iterate as "F\textasciicircum \{[s]\} = X S\_F\textasciicircum s". Three things are missing there: that S\_F - I raises the filtration (true: (S\_F-I)H = H o F - H has ordt >= ordt(H)+1 by Corollary \textbackslash ref\{plt:cor:grp-gain\}); that S\_F\textasciicircum s is again a *substitution* operator, i.e. that F\textasciicircum \{[s]\} exists in Gnear at all; and the notation, which should be S\_F\textasciicircum s(X), not a product. Theorem \textbackslash ref\{plt:thm:grp-exp\} and Corollary \textbackslash ref\{plt:cor:grp-iteration\} supply all three.\\
\cref{plt:prop:grp-dilation} & S4's closure table records only the rational case ("c x\textasciicircum \{r/s\}: Puiseux-log series") and S1 says "arbitrary real beta require a Hahn support", neither separating the three regimes nor proving that each enlargement is necessary. Corollary \textbackslash ref\{plt:cor:grp-power-regimes\} does both.\\
\cref{plt:thm:grp-multilog-group} & S2 says only that "the rational field K(L\_1,...,L\_r), or a Laurent-monomial field, repairs this defect". Proposition \textbackslash ref\{plt:prop:grp-tower-derivation\} identifies the minimal repair inside the Laurent-monomial family and shows that L\_r itself need not be inverted.\\
\cref{plt:rmk:rev-imports} & report 6 states its reversion algorithms with the inclusive cutoff [F]\_\{\textbackslash le N\}; every cutoff index taken from that source is shifted by one here.\\
\cref{plt:lem:rev-deriv-gain} & the closure of \textbackslash (\textbackslash Gnear\textbackslash ) under composition was stated and proved a second time in this part, as a separate lemma. It is item (1) of the group theorem, proved there from the substitution proposition, so it is imported here with the other three standing facts. Note that only items (1)--(3) of that theorem are imported: its item (4) cites \textbackslash cref\{plt:thm:rev-existence\} below, and importing it here would be circular.\\
\cref{plt:lem:rev-shifted-taylor} & the closure lemma stated here duplicated \textbackslash cref\{plt:thm:grp-group\}\textbackslash ,(1), which is proved in the composition-group part with the same one-line argument. It is deleted; closure is now imported as part of (I3) and appears as \textbackslash eqref\{plt:eq:rev-closure\}. For the record, the argument is that \textbackslash (F\textbackslash compose G=G+A\textbackslash compose G\textbackslash ) and that \textbackslash (A\textbackslash compose G=\textbackslash sum\_\{k\textbackslash ge0\}B\textasciicircum k\textbackslash DX\textasciicircum kA/k!\textbackslash ) has all summands in \textbackslash (\textbackslash PLpow\textbackslash ) by \textbackslash cref\{plt:lem:rev-deriv-gain\}.\\
\cref{plt:lem:rev-gain} & this estimate was proved twice in the volume: here from the two-point Taylor expansion, and in the composition-group part from the exact Taylor formula, where it is \textbackslash eqref\{plt:eq:grp-gain-lipschitz\}. The two arguments are the same computation, and the group-part statement was missing the summand \textbackslash (\textbackslash ordt(A)\textbackslash ) that its own proof produces; that omission is repaired at \textbackslash cref\{plt:cor:grp-gain\} and the duplicate proof is deleted here. The inequality is recorded below as a fifth standing import, in the instantiated form used throughout this part.\\
\cref{plt:thm:rev-existence} & all six sources call this a contraction; none says in what sense. The statement is about nu\_t, not about coefficient size, and the constant term alpha = nu\_t(A) is absent from every source formulation.\\
\cref{plt:cor:rev-first-order} & this corollary restated \textbackslash cref\{plt:thm:grp-group\} in full. The group statement now stands once, in the composition-group part, and its item (4) cites \textbackslash cref\{plt:thm:rev-existence\}; the restatement and its proof are replaced by the sentence below, which records what has just been supplied.\\
\cref{plt:prop:rev-derivative-certificate} & report 1 asserts that at finite order the two residuals "need not have identical first omitted blocks"; for an exact G\_N in the group that is false.\\
\cref{plt:prop:rev-analytic-transfer} & report 1 claims the derivative residual "often detects a wrong block one order earlier"; it detects it one level later, and is blind at level 0.\\
\cref{plt:prop:rev-truncation-safe} & report 4 states the rule as "for a leading slope F = ax + o(x) the correction is divided by a", with no hypothesis on a and no discussion of a nonconstant leading block of D\_X F.\\
\cref{plt:rmk:rev-doubling-sharp} & the sources give five mutually inconsistent exponents; one contradicts its own proof, and two slide from residual order to correct-block count.\\
\cref{plt:ex:rev-omega} & reports 1, 3, 4 and 6 assert Newton is "asymptotically superior" with no cost model stated; the claim is model-dependent and is weakened here.\\
\cref{plt:rmk:lag-nearid-meaning} & S1, S2, S3, S4 and S5 all assert only the sufficiency direction ("the r-th term has t-order at least r-1"), and all state it with the hypothesis ord\_t(A) >= 0. The correct threshold is ord\_t(A) > -1; the two agree on K[L]((t)), whose value group is Z, but they differ on the Puiseux--log enlargement, where the weaker hypothesis is exactly what the reduction of a ramified reversion needs (plt:ex:lag-puiseux). None of the sources observes that the condition is also necessary, which is what makes the misuse in plt:warn:lag-scale-trap detectable.\\
\cref{plt:rmk:lag-residue-route} & S3's appendix proves the unshifted formula by a formal residue change of variables, which is valid only when the constant term of the auxiliary series is invertible, and repairs this by a universal-coefficient localisation argument. S6 replaces that argument by the phrase "an integration by parts reduces the resulting residue". The proof above needs neither device.\\
\cref{plt:cor:lag-observables} & S3 states this "whenever both sides are interpreted to a fixed finite t-precision" and S6 "for any H for which the family is summable"; both hypotheses are vacuous or unverifiable. The order estimate plt:eq:lag-transport-order shows that no hypothesis beyond membership in the frame is needed, since nu(phi) is finite for phi nonzero.\\
\cref{plt:ex:lag-mixed} & S4 gives only b\_0, b\_1, b\_2 in this universal form; the b\_3 line plt:eq:lag-b3 is derived here and verified against the two independent worked examples plt:ex:lag-mixed and plt:ex:lag-abc.\\
\cref{plt:ex:lag-certificate-mixed} & this proposition restated \textbackslash cref\{plt:cor:rev-finite-certificate\}, and its proof re-derived \textbackslash cref\{plt:prop:rev-residual-order\} along the way. The statements are identical --- the hypothesis \textbackslash (F=X+A\textbackslash ) with \textbackslash (A\textbackslash in\textbackslash PLpow\textbackslash ) is membership in \textbackslash (\textbackslash Gnear\textbackslash ), and the two candidate truncations \textbackslash (G\_N\textbackslash ) are the same object. The restatement and its proof are deleted in favour of the pointer below; every citation of the deleted label in this part now cites the reversion part.\\
\cref{plt:rmk:lag-nonscalar-leading} & the requirement log d in K is stated in the catalogue and is omitted in S1, S2 and S6, which write log c freely over an unspecified coefficient field. Over K = Q it fails already for d = 2; see plt:warn:lag-scale-trap.\\
\cref{plt:rmk:lag-log-frame-scale} & S6 derives plt:eq:lag-u-expansion by inserting an undetermined ansatz u = rho\textasciicircum \{-1\}(S - sigma log S + C + (a log S + b)/S + ...) and matching three constants by hand. The derivation above obtains the same coefficients from the near-identity coefficient formula, so the pattern continues to all orders with no further ansatz, and the branch and field hypotheses (log c, log rho in K) become visible.\\
\cref{plt:rmk:lag-normalisation-care} & S4 asserts that after the substitution y = d x\textasciicircum \{1/q\}(1+V) "the formal implicit-function recursion is triangular" for any leading exponent q, which is true, but it does not record that the intermediate normalisation Phi = F o rev(F\_0) can leave the near-identity class, as it does in plt:ex:lag-puiseux. The correct general statement is the converse half of plt:thm:lag-rigidity, where Phi = X(1 + U o rev(F\_0)) because F\_0 is the exact leading monomial of F, not merely its leading power.\\
\cref{plt:rmk:om-branches} & S1's prose ("the same P\_n occur, multiplied by (-1)\textasciicircum n") contradicts its own displayed expansion of the inverse of X - log X, where P\_1 = L occurs with a plus sign. The correct weight is (-1)\textasciicircum \{n+1\}; the weight (-1)\textasciicircum n belongs to the negated quantity -rev(X-L), which is what enters W\_\{-1\}.\\
\cref{plt:thm:om-expansion} & S5 states uniqueness only inside -L + t K[L][[t]]; the contraction estimate holds on all of K[L][[t]], so uniqueness is unconditional there and membership in -L + t K[L][[t]] is a conclusion, not a hypothesis.\\
\cref{plt:rmk:om-normalisation} & the citation named a label "plt:thm:lagrange-burmann" that is defined nowhere in the volume; the theorem meant is the one labelled plt:thm:lag-burmann in the Lagrange--Buermann part, whose polynomial--logarithmic form is equation plt:eq:lag-burmann-pl.\\
\cref{plt:cor:om-certificate} & retitled. This is not the general finite certificate (\textbackslash cref\{plt:cor:rev-finite-certificate\}, which it invokes) but its specialisation to the single map \textbackslash (X+L\textbackslash ), where the criterion becomes a characterisation of the Lambert polynomials themselves.\\
\cref{plt:thm:om-analytic} & S5 asserts that the residual-to-error principle "proves the asymptotic validity of every truncation directly from" the FORMAL residual identity, and S1 asserts an analytic O-bound for the branch residual "by construction". Neither is a proof: a formal t-adic identity supplies no inequality. The missing ingredient is item (2), supplied below.\\
\cref{plt:thm:lambert-stirling} & S3 and S6 start the recurrence at P\_1 = u and impose P\_\{n\}(0)=0 as a side condition of the theorem rather than proving it; S4 derives the linear recurrence by "equating coefficients" in a NONLINEAR generating ODE, which does not produce it. The proof above is direct from the definition and proves the boundary condition instead of assuming it.\\
\cref{plt:cor:om-reversion-free} & this corollary carried its own argument inside the statement, so it read as an assertion with no proof. The claim and the argument are separated.\\
\cref{plt:rmk:lw-catalan} & the hypotheses of the two theorems invoked here were not printed; both hold once |z| is large, which is the regime of the claim.\\
\cref{plt:rmk:lw-catalan} & the degree law and leading coefficient quoted here for X + aL\textasciicircum 2 were printed for all n >= 0, but they fail at n = 0: the t\textasciicircum 0 coefficient is -aL\textasciicircum 2, of degree 2, and the formula 2a\textasciicircum \{n+1\}/n is undefined at n = 0. Restricted to n >= 1 and the n = 0 value recorded; verified by explicit reversion through n = 4 (c\_1 = 2a\textasciicircum 2L\textasciicircum 3, c\_2 = a\textasciicircum 3L\textasciicircum 5 - 5a\textasciicircum 3L\textasciicircum 4, c\_3 = (2/3)a\textasciicircum 4L\textasciicircum 7 - 7a\textasciicircum 4L\textasciicircum 6 + 14a\textasciicircum 4L\textasciicircum 5). independently re-verified through n = 9 in exact rational arithmetic, with the residual of the reversion confirmed to vanish identically through t\textasciicircum 9; that computation also produced the Catalan observation recorded below.\\
\cref{plt:rmk:lw-no-sector} & S1, S3, S4 and S5 all state the branch expansion "in an admissible sector" or "in a branch-compatible sector" without ever defining one, and none proves that the constructed solution is the branch conventionally labelled W\_k. Repaired: the estimate is proved on an explicit region with no sector at all (Theorem lw-uniform-remainder), and the labelling is a separate proposition with a printed distance-to-the-cut hypothesis (Proposition lw-branch-label).\\
\cref{plt:rmk:lw-s6-gap} & the exhibited S\_k was written with non-strict inequalities, hence was not open, while a coordinate sector is required to be open by Definition lw-coordinates. Made strict.\\
\cref{plt:rmk:lw-s6-gap} & S6 deduces convergence of the Lambert series from (a) convergence at frozen u and (b) |s u| -> 0 along the Lambert path. That inference is invalid: the radius of convergence in s depends on u and shrinks as |u| -> infinity. Repaired by proving the two-variable statement (Lemma lw-two-variable), which is what the argument actually needs.\\
\cref{plt:cor:lw-principal-real} & S1, S3 and S5 each assert, as though it were established here, that the principal-branch series converges for every real z >= e. Downgraded to an attributed literature report: nothing in this volume proves it, and S5's supporting observation (all terms vanish at z = e) is vacuous.\\
\cref{plt:lem:lw-critical-point} & every source displays the lower-branch expansion with the factor (-1)\textasciicircum n and explains it only by analogy ("the same polynomials reappear with alternating block signs"). Repaired: the alternation is derived as the single value a = -1 of the scalar a\textasciicircum \{n+1\} in Theorem lw-affine-log, and the box states explicitly that no second polynomial family exists.\\
\cref{plt:thm:lw-branch-point} & S1 and S3 assert the Puiseux expansion at the branch point without constructing the local inverse; S1 says only "ordinary series reversion in p produces" the coefficients, and neither source proves that a power-logarithmic expansion is impossible there. Repaired: the local biholomorphism is constructed, the coefficients are rederived (and carried one order further than any source, to p\textasciicircum 6 = -221/8505, verified against a high-precision evaluation), and the impossibility is proved from the order of the local monodromy in Corollary lw-branch-point-monodromy.\\
\cref{plt:rmk:lw-table-concordance} & the four source tables are individually correct but use two incompatible column conventions (S1 and S4 index by x = log z, S3 and S5 by z), which makes them look contradictory when read side by side. All four were recomputed to 60 digits and merged into one table indexed by the intrinsic variable xi\_0 = log z; the concordance is recorded below, together with the observation that the S3/S5 columns at z = 10 and z = 100 lie outside the region where any convergence statement is proved here.\\
\cref{plt:thm:lw-generalized-expansion} & S3, S4, S5 and S6 all give the closed form y = (alpha/beta) W\_k((beta/alpha) x\textasciicircum \{1/alpha\}) with the disclaimer "branches chosen consistently", which hides a genuine ambiguity: for non-integer alpha the symbol x\textasciicircum \{1/alpha\} is not determined by x. Repaired by Proposition lw-power-exp, which makes the three choices explicit and states the correspondence as a bijection.\\
\cref{plt:def:lw-monomial-log-datum} & the proof establishes the residual bound only "for z large", while the statement attaches it to the same z\_r that governs existence. The quantifier is repaired here by enlarging z\_r once, which is legitimate because the statement is existential in z\_r.\\
\cref{plt:ex:lw-fractional} & the drafted sentence said rho appears "only in the leading term and inside ell", which the displayed formula contradicts: rho also sits in the overall prefactor mu/rho. What is true is that rho is absent from the ratio a/X = mu/(xi - beta) that drives the series.\\
\cref{plt:prop:lw-ex2-sharp} & the remainder in this display was printed as an analytic estimate BigO(L\textasciicircum 4/Y\textasciicircum 3) although the three blocks are computed by formal Lagrange--Burmann; per the volume's own rule the three O-roles must not be conflated. Downgraded to the formal tail, with the analytic upgrade stated as a separate sentence carrying its hypothesis.\\
\cref{plt:prop:lw-ex2-sharp} & S6 states that ``the coefficient of Y\textasciicircum \{-2\} is no longer a polynomial in L and M alone''. That is true but not sharp: expanding the Y\textasciicircum \{-1\} block gives A*varpi = alpha\textasciicircum 2 L + alpha*gamma M + alpha*gamma + gamma\textasciicircum 2 M/L, so a negative power of L is already present one block earlier whenever gamma is nonzero. The corrected statement and its proof are below.\\
\cref{plt:rmk:lw-rho-zero} & the proof as drafted read the two leading blocks off Equation lw-ex3-X, which the template produces only FORMALLY; using it as an analytic asymptotic statement here would be circular. Replaced by a direct three-line derivation from the defining equation, which needs nothing beyond X -> infinity.\\
Generalized Lambert equations and reusab & as drafted this corollary wrote the reseated sum as "w", silently asserting that the reseated base point returns the SAME solution as the canonical one, and its proof discharged the point in one clause. That does not follow from the stated hypotheses: Theorem lw-basepoint attaches to each base point SOME solution of w e\textasciicircum w = e\textasciicircum xi, and the reseated one is only known to satisfy |w\textasciitilde  - (xi - ell)| <= |M\textasciitilde | + log 2, which may exceed the radius 1 of the uniqueness disc of Proposition lw-branch-label. Repaired: the identification is now a separate clause with the printed hypothesis 4(1 + |ell\_k|) <= |xi\_k|, under which the disc argument does apply, and the real case is noted to need nothing at all.\\
\cref{plt:rmk:lw-associated-stirling} & the canonical threshold was printed as 5.3569...; the root of xi = 2 + 2 log xi is 5.356693980..., i.e. 5.3567 to four decimals. Corrected, and the threshold of the reseated condition added.\\
\cref{plt:rmk:lw-associated-stirling} & the drafted reading said reseating "is numerically better at xi\_0 = 10 and 30", which its own table contradicts at xi\_0 = 10, N = 10 (6.4e-10 canonical against 1.7e-9 reseated). Corrected to what the table shows.\\
\cref{plt:def:an-scale} & S2 states that "the monomials X\textasciicircum \{-n\}L\textasciicircum m form an asymptotic scale under the lexicographic order", and S6 opens its analytic chapter by letting (phi\_nu) be "a totally ordered asymptotic scale, so that phi\_\{nu+1\} = o(phi\_nu) in the chosen enumeration" and then applying this to the polynomial-logarithmic monomials. The pairwise dominance comparisons those drafts display are all correct; what fails is the existence of the enumeration, by the proposition above. The repair is to fix a degree bound first, which is what the block formulation does automatically.\\
\cref{plt:def:an-expansion} & this definition carried the same title as \textbackslash cref\{plt:def:mul-asymptotic\}, which it extends rather than repeats: the blocks here are rational in L and may have poles, so the evaluation needs the large-X caveat below. Retitled so the two are distinguishable in the list of results.\\
\cref{plt:prop:an-uniqueness} & S6 writes the blockwise condition as f - sum\_\{n<=N\} p\_n(log X) X\textasciicircum \{-n\} = o(X\textasciicircum \{-N\}(log X)\textasciicircum \{d\_N\}) and remarks that "one must specify that inner ordering" of the monomials inside the last block. Form (v) of the theorem shows the logarithmic weight is harmless once all N are quantified, and form (vi) shows the inner ordering never has to be specified: the blockwise and monomialwise conditions coincide for every admissible profile.\\
\cref{plt:prop:an-holomorphic-substitution} & S5 states products and quotients together with a "proof sketch" and, in its part (ii), adds the hypothesis "and g is eventually nonzero". That hypothesis is redundant: it is a consequence of the leading block being nonzero, as (plt:eq:an-g-lower) shows. S5's part (iii) requires instead that "Q(log x) stays nonzero eventually", which is likewise automatic for a nonzero rational function. The full proof above replaces the sketch, and in particular exhibits the logarithmic degree k of the leading block explicitly, which the sketch suppresses; it is what makes the remainder exponent M'-k, not M'.\\
\cref{plt:cor:an-log-power-exp} & the right-hand side was attributed to \textbackslash cref\{plt:prop:tay-scalar-outer\}, which is stated for arguments in t\textbackslash PLpow. Here U has rational blocks, so the ambient is \textbackslash (\textbackslash K(L)[[t]]\textbackslash ); that case is covered by the last sentence of \textbackslash cref\{plt:prop:dif-substitution\}, which is the citation used now.\\
\cref{plt:rmk:an-composition} & the right-hand sides were identified outright with the formal operations of \textbackslash cref\{plt:lem:cmp-exp-log\} and \textbackslash cref\{plt:thm:mul-binomial\}, both of which are stated over \textbackslash PLpow; for a genuinely rational \textbackslash (U\textbackslash ) the formal object is the scalar substitution of \textbackslash cref\{plt:prop:dif-substitution\} in \textbackslash (\textbackslash K(L)[[t]]\textbackslash ), and the identification with the named operations is asserted only where they are defined. Restricted accordingly.\\
\cref{plt:lem:an-tail-integral} & S6 writes, in the chapter that this part consolidates, "For f(X) = X + O(log X), one has f'(X) = 1 + o(1) under the relevant differentiability hypotheses". No such hypotheses are named, and the implication is false without them, by the display in the box above. S6 uses the conclusion to justify choosing m close to 1 in its residual-transfer theorem. The repair is (plt:prop:an-omega-computable) and (plt:thm:an-inverse-transfer), which take the derivative bound as a hypothesis and verify it directly for the map at hand instead of deducing it from an expansion.\\
\cref{plt:cor:an-tail} & the hypothesis "locally integrable" was missing here although the proof invokes \textbackslash cref\{plt:thm:an-integration\}, which requires it; without it the integrals below need not exist. Restored.\\
\cref{plt:rmk:an-sector-needed} & the displayed constant was written (1-\textbackslash delta)\textasciicircum \{-N-1\}, which bounds |w|\textasciicircum \{-N-1\} by ((1-\textbackslash delta)|z|)\textasciicircum \{-N-1\} only when N+1\textbackslash ge0. Cutoffs here run over all of Z (\textbackslash ordt F may be negative), and for N\textbackslash le-2 the maximum of |w|\textasciicircum \{-N-1\} on the circle is attained at the outer end instead. Replaced by the two-sided constant \textbackslash kappa, which covers both signs.\\
\cref{plt:thm:an-enclosure} & \textbackslash cref\{plt:prop:rev-analytic-transfer\} is printed with the analytic inverse written \textbackslash rev\{f\}. In this volume that symbol is reserved for the FORMAL compositional inverse -- and the whole content of (plt:thm:an-inverse-transfer) is that the formal and the analytic inverse agree asymptotically, a statement that is vacuous if the two are written with one glyph. This is the same collision the ed. note before (plt:sec:an-omega-constants) records against S2. The recollection below therefore names x\_* by the equation it solves.\\
\cref{plt:thm:an-residual-analytic} & S5 writes, immediately after its residual-to-error principle: "For f(x) = x + log x, this proves the asymptotic validity of every truncation omega\_N directly from [the formal residual identity]." It does not: the principle converts an analytic residual into an analytic error, and the formal identity provides no analytic residual. The missing step is the one carried out in (plt:thm:om-analytic) for this map and in (plt:thm:an-residual-analytic) in general. The same gap is present in S1 ("by construction") and is repaired in the same way.\\
\cref{plt:thm:an-residual-analytic} & the coefficient field was left at the ambient \textbackslash (\textbackslash K\textbackslash subseteq\textbackslash ComplexNumbers\textbackslash ) of this part, but the proof applies \textbackslash cref\{plt:thm:cmp-analytic-transfer\}, whose hypothesis (A2) asks for an inner realization with REAL coefficients, and both \textbackslash (\textbackslash varphi\textbackslash ) and \textbackslash (h\textbackslash ) are real valued here. The real coefficient field is now printed, as house rule 5 requires.\\
\cref{plt:rmk:an-hypothesis-roles} & coefficient field printed: \textbackslash cref\{plt:thm:an-residual-analytic\} is applied below, and it now carries \textbackslash (\textbackslash K\textbackslash subseteq\textbackslash RealNumbers\textbackslash ).\\
\cref{plt:prop:an-omega-computable} & this sentence read "Item (2) of the list preceding \textbackslash cref\{plt:thm:om-analytic\} is the specialisation of (H2) to \textbackslash varphi(Y)=Y+\textbackslash log Y", which is wrong in an instructive way: (H2) specialised to that map is the trivial statement \textbackslash varphi\textbackslash approx X+L (the remainder vanishes identically), whereas item (2) asks for an analytic remainder for log applied to the TRUNCATION. Item (2) is the conclusion of \textbackslash cref\{plt:thm:an-residual-analytic\}, not its hypothesis. Corrected.\\
\cref{plt:prop:an-omega-computable} & S2's transfer theorem lists the hypotheses "F(y) = y + o(y)" and "F'(y) = 1 + o(1)" and then assumes an ANALYTIC residual bound F(G\_N(X)) - X = O(X\textasciicircum \{-N-1\}L\textasciicircum M) outright. Assuming it is legitimate; the surrounding text, however, obtains it from the formal reversion, so the theorem as used is circular. The first hypothesis is never used in S2's proof. (H2) above is the honest replacement, and (plt:thm:an-residual-analytic) discharges it. S2 also writes the analytic inverse as \textbackslash rev\{F\}; in this volume that symbol denotes the FORMAL compositional inverse, and the two must not share a glyph in a statement whose whole content is that they agree.\\
\cref{plt:prop:an-omega-explicit0} & the step "for x>1 the envelope gives omega(x) > x - log x" skipped the verification of that envelope's own hypothesis, which is omega(x)>1 and not x>1. Supplied: omega is strictly increasing with omega(1)=1.\\
\cref{plt:rmk:an-no-uniform-constant} & the proof invoked \textbackslash cref\{plt:prop:an-omega-computable\}, which is stated for x>1, while this proposition is stated for x\textbackslash ge1. The endpoint is dispatched separately rather than left to a proposition that does not cover it.\\
\cref{plt:rmk:an-cancellation} & the relative width was quoted as \textbackslash mu\_3\textasciicircum \{-1\}. Width divided by the upper end is 1-(1+\textbackslash mu\textasciicircum \{-1\})\textasciicircum \{-1\}=(1+\textbackslash mu)\textasciicircum \{-1\}, which is slightly smaller; recomputed, it is 1.0374 per cent against \textbackslash mu\_3\textasciicircum \{-1\}=1.0483 per cent.\\
\cref{plt:ex:an-branchpoint} & S1 asserts, of its own error table at x = 10, that "the N = 2 truncation is slightly worse than N = 1". Its table gives 1.7466787e-3 at N = 1 and 1.7369609e-3 at N = 2, which is smaller, and recomputation at 40 digits confirms both table entries. The prose contradicts the table it is describing. What is true at x = 10 is that the N = 2 correction almost exactly cancels the N = 1 error and flips its sign, improving it by only 0.6 per cent, and that genuine non-monotonicity occurs later, at N = 6 -> 7 and N = 11 -> 12, as recorded above. S3's table at log z = 2.3026 makes the analogous claim correctly, for the pair N = 2, 3.\\
Residual-to-error transfer & "with relative width 1/(x-\textbackslash log x)": the exact relative width is (1+x-\textbackslash log x)\textasciicircum \{-1\}; the quoted quantity is an upper bound for it, so the sentence is corrected to say so.\\
\cref{plt:def:du-substitutions} & this sentence originally attributed the derivation to all four classes. Unlike its mirror \textbackslash PLpow at infinity, the power-series class K[Lambda][[x]] is NOT stable under d/dx: d(Lambda)/dx = -1/x already leaves it. This is trap (T6) in its sharpest form, and it is why the near-identity class of plt:def:du-near-identity is described inside K[Lambda]((x)).\\
\cref{plt:def:du-substitutions} & S2 sets ell = log(1/x); S1 and S6, following the normal-form literature, set ell = -1/log x. The two are reciprocal, and no single displayed formula distinguishes them. We fix Lambda for the large log and never write a bare ell for it.\\
Power--log series near a finite point an & this item said "every nonzero f in tau(PLlau), evaluated at the generators". A general element of K[Lambda\_2]((Lambda\textasciicircum \{-1\})) is a formal, typically divergent, series and has no evaluation; and the proof uses only asymptotic equivalence to the leading monomial. Restated for functions.\\
Power--log series near a finite point an & S6 displays the expansion with a bare remainder O(M\textasciicircum 3/S\textasciicircum 3) and no argument. The estimate below is proved; the sharp form of S6's remainder is recovered in the last sentence.\\
Power--log series near a finite point an & "asymptotic expansion on the scale of all monomials Lambda\textasciicircum \{-n\} Lambda\_2\textasciicircum j" would presuppose that this doubly indexed family is a Poincare scale, which it is not (plt:prop:an-no-enumeration). The precise assertion is the relation of plt:def:mul-asymptotic read in the variable X=Lambda, which is exactly what plt:eq:du-lower-branch-error states.\\
\cref{plt:prop:du-repeated-back} & S2 gives the product form of (13.7) but never names the operator; the shift index n-k+1 is the trap, since the naive mirror of plt:eq:dif-repeated would be Delta\_\{n,k\}, which is wrong for k >= 2.\\
\cref{plt:prop:du-traps} & S2's eq:Dulac-repeated-derivative is correct as printed; the repair is that it is stated without proof, without the empty-product convention, and without identifying the operator, so its k >= 2 index shift is invisible.\\
\cref{plt:def:du-near-identity} & S2 (eq:small-Lagrange) states the near-zero Lagrange-Burmann formula as having "exactly the same form" and gives no proof; S1 and S6 do the same for the fixed-point construction. The formula is true (Theorem plt:thm:du-lagrange-zero) but it is not the transport of plt:thm:lag-burmann: the transported frame has a different derivation and a different distinguished element, and no rescaling of nu\_x repairs (F3).\\
\cref{plt:prop:du-degree-zero} & the appeal to the enlarged exponent group also needs plt:lem:du-taylor-zero there, which was stated for m\_0; its proof uses nu\_x(B) >= 2 only through nu\_x(B) - 1 > 0, so it generalises as stated.\\
\cref{plt:def:du-dulac-series} & S1's eq:dulac-inverse-first is correct, but it is written in the reciprocal convention ell = -1/log x, so its exponents of the logarithm are the negatives of ours; transporting it without noting this produces the wrong sign on every logarithmic exponent.\\
\cref{plt:def:du-dulac-germ} & this read "has no finite truncation Trunc\_\{<1\} other than 0", which is doubly wrong: Trunc\_\{<N\} of that series IS finite (and in general nonzero) for every N < 1, and Trunc\_\{<1\} is the whole infinite series, not 0. The correct statement is that truncation stops being finite at the cutoff 1.\\
Power--log series near a finite point an & this item asserted that f |-> f-hat is a ring homomorphism ONTO D(R). Surjectivity is a Borel--Ritt statement for left-finite real exponent sets; it is not proved here and is not needed for anything below, so the claim is downgraded to "into" and the gap is named.\\
\cref{plt:rmk:du-dulac-problem} & "P-hat is independent of the parametrisation" would be false as literally read -- a change of chart conjugates P-hat. What is invariant is the membership and the leading data, which is what cor:du-chart-independence supplies.\\
\cref{plt:thm:du-log-degree-failure} & "is the semigroup generated by r and 1" asserts an equality that the grid bookkeeping of plt:thm:du-composition-closure(ii) does not give; it gives a containment, which is all that is used.\\
Power--log series near a finite point an & the statement here read "it lies in none of sigma(PLlau), sigma(PLrat), sigma(PLit), nor in any Puiseux-log or Dulac class". That is false for rho a negative integer, where item (i) exhibits the element inside sigma(PLrat) and sigma(PLit); the exclusion from those two classes needs rho not an integer, and only the exclusion from the polynomial-block classes holds for every rho outside the natural numbers. Split accordingly.\\
Power--log series near a finite point an & the second clause read "for every alpha with p\_alpha nonzero and gamma*alpha not a natural number". Only the summand of the SMALLEST alpha is protected from interference: for alpha' < alpha the summand of alpha' also contributes at the x-exponent alpha*beta (through its u-tail), with Lambda-exponents in gamma*alpha' + N, and these can meet gamma*alpha, so no cancellation argument is available for a general alpha. Restricted to alpha\_0 = nu\_x(f-hat), where the leading block stands alone.\\
Power--log series near a finite point an & the argument here claimed that injectivity of alpha |-> alpha*beta makes the summands non-interfering. It does not: the summand of alpha' has x-support in alpha'*beta + Gamma\_u, so for alpha' < alpha it also reaches the x-exponent alpha*beta. What is true, and all that is needed, is that only summands with alpha' <= alpha reach it.\\
\cref{plt:rmk:alg-derived-order} & S4 stores the least allowed exponent n\_min as an independent field and then warns that a stale zero at that exponent corrupts order and unit tests. Lemma plt:lem:alg-denotation-coset shows the field is redundant: the order bound is a function of the canonical form, so the failure is eliminated by deriving it rather than by documenting it.\\
Data structures and finite-order algorit & the previous version of (v) asserted that sharpness fails at N\_1=0 because DX annihilates the unknown constant block. That is an artifact of scalar power series and is false over R=K[L]: the unknown block of index N\_1 is an arbitrary polynomial, and a=t\textasciicircum \{N\_1\}L already has DX a = t\textasciicircum \{N\_1+1\}(1-N\_1 L) \textbackslash ne 0 for every N\_1, including N\_1=0. What is true, and is what the sources were reaching for, is that the output ORDER bound can fail to be attained -- exactly when the leading block is a nonzero constant sitting at index 0 -- and that this is a loss of relative precision, not of absolute precision.\\
Data structures and finite-order algorit & the criterion was stated without the hypothesis v\_1 != 0, at which it fails: there v\_1(m\_1)=0, yet the returned representative is (0,N\_1+1), whose order bound is N\_1+1=m\_1+1, so the bound is attained.\\
Data structures and finite-order algorit & the sharpness witnesses below were originally exhibited only at the zero representative (sum) or for one of the two competing bounds (product, quotient), which proves sharpness of the LAW but not of each instance. Since Definition plt:def:alg-sound quantifies sharpness per instance, the witnesses are now chosen so that the block that differs is exactly the declared cutoff, for every instance.\\
Data structures and finite-order algorit & the witness that perturbs the first argument produces a difference in the block N\_1+m\_2, which is the declared cutoff only when that bound is the smaller of the two. The previous version used it unconditionally and so proved sharpness of the law, not of each instance as Definition plt:def:alg-sound requires. The witness must be selected by which bound realises the minimum.\\
Data structures and finite-order algorit & "follows from the two witnesses used in (ii)" was too quick: as in (ii) the witness must be the one realising the minimum. Here, unlike (ii), both witnesses exist for every instance, because the numerator is allowed to be zero and may then be replaced by t\textasciicircum \{N\_1\}.\\
\cref{plt:cor:alg-relative-precision} & the proof still carried the discarded claim that sharpness fails at N\_1=0, contradicting the corrected statement above. The witness a=t\textasciicircum \{N\_1\} is the wrong one: it is annihilated exactly at N\_1=0. Over R=K[L] the undetermined block is an arbitrary polynomial, and a=t\textasciicircum \{N\_1\}L works for every N\_1.\\
\cref{plt:cor:alg-relative-precision} & the witness was a single pair (F=t\textasciicircum \{m\_1\} with m\_1 nonzero), which proves sharpness of the law but not the per-instance claim actually stated. Under the printed hypothesis ordt(DX v\_1)=m\_1+1 the generic witness works, and the hypothesis is exactly what makes the k=1 term dominate.\\
\cref{plt:rmk:alg-cancellation} & the inequality was stated as rho(DX F) >= rho\_1. It runs the other way: the absolute cutoff gains one block, but in the exceptional case of part (v) the output ORDER gains more than one, so the difference -- which is rho -- can only fall. The extreme case v\_1 = c in K\textasciicircum x returns the zero representative, whose relative precision is 0. two further repairs to this clause. (a) The equality criterion for the derivation fails at the zero representative: there \textbackslash (v\_1(m\_1)=0\textbackslash ) while \textbackslash (m\_\{\textbackslash mathrm\{out\}\}=N\_1+1=m\_1+1\textbackslash ), so equality holds although the criterion does not; the clause is restricted to \textbackslash (v\_1\textbackslash ne0\textbackslash ). (b) "no bound in either direction" was an overclaim for the sum: an upper bound does hold, \textbackslash (\textbackslash rho\textbackslash le\textbackslash max\textbackslash \{\textbackslash rho\_1,\textbackslash rho\_2\textbackslash \}\textbackslash ) (proved below), so what the witness shows is that no LOWER bound holds -- and \textbackslash (\textbackslash rho\textbackslash ) can drop to \textbackslash (0\textbackslash ), not merely to \textbackslash (1\textbackslash ).\\
\cref{plt:rmk:alg-cancellation} & the upper bound for the sum was asserted to be absent; it is not.\\
\cref{plt:ex:alg-guard-block} & "only addition can destroy it" overlooked the annihilating case of the derivation, which the corrected corollary now records.\\
\cref{plt:def:alg-cost-models} & the source sentence asked for G "through block N-1", which is the datum already assumed and yields the product only through block N-2. Since multiplication by t\textasciicircum \{-1\} shifts indices down by one, the block of index n of the product is read off the block of index n+1 of G.\\
\cref{plt:def:alg-cost-models} & S2 states only that inputs known modulo t\textasciicircum N with nonnegative valuation give a sum and product known modulo t\textasciicircum N. That is sound but not sharp: by part (ii) the product is known modulo t\textasciicircum \{N+min(m\_1,m\_2)\}, which is strictly better whenever a factor has positive order, and strictly worse -- the case S2's formulation hides -- as soon as a factor is Laurent.\\
Data structures and finite-order algorit & "four times the naive quadratic guess in the exponent" is not a statement: the exponent is four rather than two. Also, the profile of the inverse of X+L is d(0)=1, d(n)=n for n>=1, which differs from case (ii) in the block of index 0; the exact count is corrected below.\\
Data structures and finite-order algorit & "the profile of the compositional inverse of X+L" names the wrong object: that inverse has a block of index -1, namely X itself. The profile in question is the profile of its CORRECTION.\\
\cref{plt:prop:alg-iterated-derivative} & the two recurrences do not cost the same: the logarithm sums j<=n-1 and the exponential sums j<=n, so the exact counts differ by N-1. The single figure binom(N,2) was right only for the exponential.\\
\cref{plt:cor:alg-antiderivative} & the claim was "shiftop\{n\}\{k\}p = 0 as soon as k > deg p and n = 0", presumably by analogy with iterated partial\_L. It is false: only the single factor with index j = -n is partial\_L, all the others are bijective, so the operator drops the degree by exactly one and no further. Already shiftop\{0\}\{2\}L = (partial\_L - 1)(1) = -1, nonzero with k = 2 > 1 = deg L. The corrected kernel statement is what the annihilation clause of the precision calculus needs.\\
Data structures and finite-order algorit & the count was quoted without the hypothesis N>a under which it is valid; for N<=a the ceiling is nonpositive and counts nothing, which is the right answer only because both sides of the displayed identity then vanish. The case is separated.\\
Data structures and finite-order algorit & the witness F=t\textasciicircum a fails whenever a<0 and K-1>|a|, because then shiftop\{a\}\{K-1\} annihilates the constant 1 -- e.g. a=-1, h=0, N=10 needs k=9 while DX\textasciicircum k t\textasciicircum \{-1\}=0 already for k=2. Taking the block to be L instead of 1 removes the exception, by the kernel description in plt:prop:alg-iterated-derivative.\\
\cref{plt:prop:alg-shift} & S1 records the budget as k <= floor((N-r)/(h+1)). Under the exclusive cutoff adopted here that admits one Taylor term too many; the sharp exclusive bound is k <= floor((N-1-a)/(h+1)). The discrepancy is exactly the inclusive/exclusive one-block conversion of the catalogue and is harmless for correctness but not for cost.\\
\cref{plt:prop:alg-shift} & the running power was truncated at the output cutoff N. That is wrong as soon as m\_F<0: the term P D contributes to a block of index below N through every block of P of index below N-ordt(D), and ordt(D)>=m\_F+k can be negative. Truncating P at N-ordt(D), with the CURRENT D, repairs it and never demands a block of H beyond N\_H, since N-ordt(D)<=N-m\_F-k<=N\_H+1-k.\\
Data structures and finite-order algorit & the output line used a cutoff N that no line of the algorithm defined, and attributed it to plt:cor:cmp-precision-budget, which furnishes the INTERNAL budget for the auxiliary powers, not the output cutoff. The output cutoff must be a function of the input cutoffs -- item (2) of plt:cor:cert-reproducible -- and the one the stability estimates deliver is min\{varrho N\_F, varrho m\_F + N\_U\}.\\
Data structures and finite-order algorit & ALGORITHM BUG, the generator-engine twin of the one repaired in plt:alg:compose-taylor. The cap on the powers of V was min\{D,N-1\}, justified by ordt(V\textasciicircum k)>=k alone. That ignores the factor t\textasciicircum \{varrho n\} multiplying block n: the summand for the outer index n contributes to a block of index below N whenever varrho n + k < N, so for a Laurent outer series (m\_F<0) powers with k>=N are still needed. Concretely, F=t\textasciicircum \{-1\}L\textasciicircum 5, varrho=1, N=1 has cap min\{5,0\}=0 and returns a wrong block of index -1, although V\textasciicircum 1 is required. The correct cap is min\{D, N-varrho*m\_F-1\}, which reduces to at most N-1 exactly when m\_F>=0.\\
\cref{plt:prop:alg-reversion-precision} & the two step counts inherited the wrong cap and the wrong index range; the retained outer indices run from m\_F to N-1, of which there are N-m\_F, not N.\\
Data structures and finite-order algorit & the loop tested a variable mu that the step list never initialised.\\
Data structures and finite-order algorit & S4's driver hard-codes the schedule as q := min(N, max(1, 2p+2)). The additive +2 is a fudge for an unstated guard; the guard actually required is the one computed by plt:thm:rev-newton-guard together with the composition law of plt:thm:alg-precision-calculus(vi), and the abort test above makes any residual shortfall visible instead of silently absorbing it.\\
\cref{plt:prop:alg-reversion-comparison} & the cutoff carried by the running power was left as an ellipsis, which is exactly the datum a precision contract may not omit. It is N-r+1: the block of index n of DX\textasciicircum \{r-1\}P is built from the block of index n-(r-1) of P, so the blocks of P of index < N-r+1 are the ones the accumulator can see, and no earlier truncation is thereby too coarse, since the cutoff demanded of P decreases by one at each step while the product law (plt:thm:alg-precision-calculus(ii)) delivers min\{(N-r+2)+h, N+(r-1)h\} >= N-r+2 from the previous power and A.\\
Data structures and finite-order algorit & the product count was quoted as the ceiling itself; the stopping rule of plt:alg:reverse-lagrange runs r=1..R with R = ceil((N+1)/(h+1)) - 1.\\
Data structures and finite-order algorit & "only addition destroys it, and only differentiation improves it" is wrong twice: differentiation improves the ABSOLUTE cutoff by one block while leaving the relative precision unchanged, and in the annihilating case it destroys the relative precision outright.\\
\cref{plt:def:cert-certificate} & "every primitive has a faithful equation, with exactly one exception" overstated what plt:tab:cert-gauges contains: the table covers the primitives that invert a cheaper one, and lists no equation at all for the sum, the product and the derivation, which invert nothing. For those the natural test is the Leibniz identity, and it is unfaithful in precisely the way plt:rmk:cert-derivative-defect records. The claim is restricted to what is proved.\\
\cref{plt:rmk:cert-division-gauge} & S2 (eq. recip-residual-test) and S3 (eq. divisionresidual) both state the residual test as "AB\_N - 1 = O(t\textasciicircum N)" or "BC\_\{<N\} - A = O(t\textasciicircum N)" and read off correctness through block N-1. That reading is the gauge-zero one and is correct only for a denominator of order 0. With a denominator of order m the same residual certifies through block N-m-1, which is fewer blocks when m>0 and more when the denominator has a pole.\\
\cref{plt:prop:cert-composition} & the hypothesis was "whose leading scalar equals b\_0\textasciicircum \{p/q\}", but the leading scalar of D is by plt:def:lead-leading-data the coefficient of its leading MONOMIAL t\textasciicircum 0 L\textasciicircum j, which for a block of positive L-degree is not the block; the proof uses the block. The hypothesis is restated in the form the proof needs and the counterexample clause requires.\\
\cref{plt:lem:cert-omega-ode} & "four mutually independent identities" is not what is meant and is not true: all four characterise the same family, so they are logically equivalent, not independent. What is claimed, and what matters, is that they share no code path.\\
\cref{plt:rmk:cert-suite-limits} & T3 had the gauge subtracted where it must be added. By plt:prop:cert-reciprocal, ordt(BQ-A)=ordt(Q-A/B)+ordt(B); a quotient sound at cutoff N has ordt(Q-A/B)>=N, so its residual vanishes below N+ordt(B). The printed N-ordt(B) is the threshold a WRONG answer already passes when ordt(B)>0, so the test as printed was vacuous in exactly the case the remark plt:rmk:cert-division-gauge is about.\\
Residual certificates and reproducible v & same overclaim as in the opening list: the table has no entry for the sum, the product or the derivation.\\
\cref{plt:prop:fail-no-uniform-constant} & the crossover was quoted without the hypothesis D>=3 under which the root exists, and the asymptotic x\textasciicircum \{*\}(D) \textasciitilde  D\textbackslash log D was asserted; a one-line proof is supplied, since the volume states nothing else about this root.\\
Failure modes, diagnostics, and a decisi & three repairs. (a) The clause was stated for every D>=1, but x=D\textbackslash log x has NO real root for D=1,2 (the function x-D\textbackslash log x attains its minimum D(1-\textbackslash log D)>0 there), so x\textasciicircum \{*\}(D) does not exist and the "if and only if" cannot be read; for those D the blocks decrease at every X>1. (b) Even for D>=3 the equivalence as printed was false near X=1: x-D\textbackslash log x has TWO roots, and the blocks also decrease below the smaller one -- at X=2 and D=10, L\textasciicircum \{10\}=0.025<2. What x\textasciicircum \{*\}(D) is, is the threshold beyond which they decrease for good. (c) The printed factor 2.42... is not the truncated decimal expansion of L\textasciicircum \{10\}/X=2.41584..., only its rounding.\\
Failure modes, diagnostics, and a decisi & the proof chose only X >= max\{X\_0, e\textasciicircum 2\}, at which the series can still diverge (D=10, X=e\textasciicircum 2 has L\textasciicircum \{D\}=1024>X) and f(X) is then undefined; the choice must also put X in the convergence range, which is possible because L\textasciicircum \{D\}/X -> 0.\\
\cref{plt:prop:fail-reversion-degree} & the existence of the threshold was asserted for every D and the second root was overlooked; both are repaired here.\\
Failure modes, diagnostics, and a decisi & the list ran to nine items against eight failure modes, the branch being the extra one; it is folded into the item it travels with.\\
\cref{plt:prop:ext-wo-maximal} & the draft claimed that no argument of \textbackslash cref\{plt:sec:ext-depth\} needs more than parts (i)--(ii). That is not accurate: the statements of \textbackslash cref\{plt:thm:ext-tower-strict\} and \textbackslash cref\{plt:cor:ext-depth-unbounded\} call H\_r a differential FIELD, and the field property of a Hahn series ring over a non-archimedean monomial group is exactly Neumann's lemma. What is true is that every argument carried out here needs only (i)--(ii); the field property is imported, not reproved.\\
\cref{plt:prop:ext-leftfinite} & S1 ("What changes is the proof that every target coefficient receives only finitely many contributions ... generalized supports make it a central axiom") and S3 ("The well-order condition guarantees that each coefficient receives only an admissible family of contributions") both present well-ordering as if it were equivalent to local finiteness of a single product. It is sufficient, and it is necessary for the class, but it is not necessary for an individual product; the witness below shows this.\\
\cref{plt:prop:ext-leftfinite} & the naive bound "k,l <= 2/(gamma-2)" is false -- for gamma-2 = 1/2 the pair (3,6) is a representation with l = 6 > 4. Only the smaller index is bounded above; the larger is then determined.\\
\cref{plt:prop:ext-leftfinite} & the draft said left-finiteness "is literally \textbackslash cref\{plt:thm:ring-support\}(d) with Z replaced by Gamma". That is false and contradicts \textbackslash cref\{plt:prop:ext-hahn-trunc-fails\}(b) below, which observes that clause (d) read verbatim over a general Gamma is strictly weaker: it asks only for an exhausting enumeration of order type at most omega and does not demand that the enumeration be cofinal. Over Z cofinality is automatic, which is why the two conditions are indistinguishable in Parts I--IV.\\
Generalized power--log classes, nested l & the draft cited a nonexistent label plt:thm:grp-taylor; the formal Taylor formula of this volume is \textbackslash cref\{plt:thm:tay-formula\}.\\
\cref{plt:prop:ext-hahn-trunc-fails} & the draft used an undefined macro \textbackslash Gfilt here; the volume preamble provides no such macro, so the filtration steps are written out.\\
Generalized power--log classes, nested l & S1 ("The arithmetic formulas from Part II therefore survive almost unchanged"), S3 ("The differential rule is again ... Near-identity composition remains controlled") and S5 ("Near-identity composition and reversion continue to work when the support conditions ensure coefficientwise finiteness") all extend Parts I-IV to Hahn supports without recording that the finite bookkeeping of Part VI -- truncation, sparse representation, precision calculus -- requires strictly more than well-basedness. \textbackslash cref\{plt:prop:ext-hahn-trunc-fails\} isolates what is lost, and \textbackslash cref\{plt:thm:ext-lf-completion\} isolates the class where nothing is.\\
Generalized power--log classes, nested l & the draft recorded only "+ o(1)" here, which is too weak to produce the error term displayed below: an additive o(1) inside \textbackslash log X contributes a relative error o(1/\textbackslash log Y), larger than (\textbackslash log\textbackslash log Y)\textasciicircum 2(\textbackslash log Y)\textasciicircum \{-2\}. The sharper form is what the two-line computation actually gives, and it is what the displayed remainder needs.\\
\cref{plt:lem:ext-tower-derivation} & the draft concluded that the display "is exactly the sigma = 1/2 instance" of a theorem stated only for sigma in N. That is a claim about a proof belonging to another part, and it is not made here: the display above is a self-contained computation, and the agreement with the theorem is recorded as a check, not used as an input.\\
\cref{plt:lem:ext-tower-derivation} & \textbackslash cref\{plt:thm:lag-log-normal-form\} restricts to sigma in N although its proof uses only that sigma is a scalar of K; the half-integral instance is needed here and is verified directly above rather than by widening the hypothesis of a theorem belonging to another part. The displayed coefficient 1/4 agrees with \textbackslash eqref\{plt:eq:lag-X-expansion\} at c=1, rho=1, sigma=1/2.\\
\cref{plt:thm:ext-tower-strict} & the draft pointed at an equation label plt:eq:grp-tower-derivative that the volume does not carry; the statement it needs is the enclosing proposition, which does exist.\\
\cref{plt:rmk:ext-completion-order} & the draft proved incomparability only against the identity nesting, while \textbackslash cref\{plt:rmk:ext-completion-order\} asserted (r+1)! pairwise distinct iterated fields. The argument already gives the general statement, because it uses nothing about the identity beyond the relative order of the two chosen levels; the last paragraph records this.\\
\cref{plt:ex:ext-fast-variable} & the draft cited a label plt:rmk:div-embedding-warning that the volume does not carry; the statement being invoked is the expansion map itself.\\
\cref{plt:ex:ext-fast-variable} & the draft wrote "an identity proved after expansion need not descend". That contradicts the injectivity it cites in the same sentence: if \textbackslash iota(a)=\textbackslash iota(b) then a=b. What non-surjectivity actually costs is that a computation carried out after expansion may leave the image, and that an exact cancellation in K(L) is, after expansion, visible only in the limit and at no finite truncation. Both claims are restated correctly.\\
\cref{plt:rmk:ext-fast-variable} & the draft obtained the display by asserting that the proof of \textbackslash cref\{plt:thm:lag-log-normal-form\}, which is stated only for sigma in N, works verbatim for every scalar beta. That is a claim about a proof belonging to another part and is not made here; the two-line derivation is given instead, and the agreement with that theorem is recorded as a check.\\
\cref{plt:rmk:ext-boundary-comparison} & "a \textbackslash notin \textbackslash Gamma" is a statement about the image of \textbackslash Gamma in \textbackslash K under the embedding of \textbackslash cref\{plt:def:dif-scale\}; the exponent a is read as a scalar of \textbackslash K, which is what the proof needs and what makes the enlargement of the exponent group in the draft's proof unnecessary.\\
\cref{plt:rmk:ext-boundary-comparison} & the draft passed to \textbackslash Gamma' = \textbackslash Gamma + (1/2)Z without checking that \textbackslash Gamma' is admissible, i.e. that the embedding \textbackslash Gamma \textbackslash hookrightarrow \textbackslash K of \textbackslash cref\{plt:def:dif-scale\} extends to it injectively. It does, and the verification is one line, given here.\\
\cref{plt:rmk:ext-boundary-comparison} & the draft took \textbackslash Gamma' = \textbackslash Gamma + aZ here. That enlargement is neither needed nor in general admissible -- for K = Q, \textbackslash Gamma = Z and a irrational the group Z + aZ has rank two and admits no injective homomorphism into Q. No enlargement is required: a is a coefficient, not an exponent, and u = at already lies in H\_r(K,\textbackslash Gamma).\\
\cref{plt:prop:ext-gain-fails} & the draft pointed at an equation label plt:eq:dif-exp-value that the volume does not carry; the statement meant is the classification of which elements are exponentials.\\
\cref{plt:rmk:ext-adh-replacement} & the draft assumed only "E nonzero and c a nonzero constant". In a ring with zero divisors that does not force cE \textbackslash ne 0, and the proof needs \textbackslash partial E = cE \textbackslash ne 0 in order to apply the hypothesis on \textbackslash nu at G = E. The hypothesis is stated in the form the proof uses; it is satisfied whenever \textbackslash mathcal T is a domain.\\
Relation with the full transseries field & the draft opened this list with a Dahn--G\textbackslash "oring item citing the key dahn-goering-1987, which is not in the merged bibliography of this volume. Rather than cite a key that does not resolve, the item is dropped; it carried no content used anywhere else.\\
\cref{plt:rmk:ext-resurgence-honest} & the draft referred to an equation label plt:eq:cmp-realization that the volume does not carry; the condition invoked is the defining one of the realization relation.\\
\cref{plt:rmk:ext-open-convergence} & three appeals to a label plt:rmk:om-open, which the volume does not carry, are replaced by the statement they were standing in for; see also \textbackslash cref\{plt:rmk:ext-open-convergence\} and \textbackslash cref\{plt:rmk:ext-open-uniformity\}.\\
\cref{plt:rmk:ext-open-convergence} & the draft also cited salvy-shackell-1992, a key that is not in the merged bibliography; the item it carried (functional inverses) is covered by salvy-shackell-1999, which is.\\
\cref{plt:rmk:ext-open-depth} & the draft attributed the singularity locus to a label plt:rmk:om-open that the volume does not carry. The locus is recomputed here from the generating equation, and it does come out as the draft stated it.\\
\cref{plt:def:bell-partial} & the second group is a restatement inside the volume, not an error, but the draft called its two members "the same statement" and declared them interchangeable. They are not: plt:prop:tay-scalar-outer restates plt:prop:dif-bell and then adds a clause about expanding the outer function about a base point z\_0 \textbackslash ne 0. Nothing below leans on that clause -- see the editorial note at plt:lem:bell-chain, where the base point is made part of the datum instead -- but the two statements are not interchangeable.\\
\cref{plt:ex:bell-table} & the draft wrote "\$\textbackslash BigO(N\textasciicircum \{3\})\$ coefficient multiplications" without saying what a coefficient is here. The count \$\textbackslash BigO(N\textasciicircum \{3\})\$ is correct only for the operation "multiply one indeterminate into one stored entry"; the entries are not scalars, and the scalar count is genuinely superpolynomial. Left unqualified this is exactly the conflation that plt:cor:alg-cost-not-quadratic warns against, so the two levels are now separated.\\
\cref{plt:lem:bell-chain} & the draft wrote \textbackslash varphi=\textbackslash sum(\textbackslash varphi\_k/k!)z\textasciicircum k -- coefficients at the origin -- and then set \textbackslash varphi(F)=\textbackslash sum(\textbackslash varphi\_k/k!)W\textasciicircum k with W=F-z\_0. That is \textbackslash varphi evaluated at F only when z\_0=0; in general the coefficients of the k-th power of W are \textbackslash varphi\textasciicircum \{(k)\}(z\_0)/k!, not \textbackslash varphi\_k/k!. (Take \textbackslash varphi(z)=z, z\_0=1, F=1+t: the displayed sum returns t, not F.) The justification offered, "apply the translate z\textbackslash mapsto\textbackslash varphi(z\_0+z)", is also unavailable: re-expanding a formal element of \textbackslash K[[z]] about a centre z\_0\textbackslash ne0 requires an infinite sum in \textbackslash K for every coefficient, and \textbackslash K carries no topology here. Both defects are repaired at once by making the centre part of the datum -- the outer function IS its expansion at z\_0. Nothing downstream changes: the proof only ever uses c\_k=\textbackslash varphi\_k/k! and the relation \textbackslash varphi'(F)=\textbackslash sum\_\{k\textbackslash ge1\}kc\_kW\textasciicircum \{k-1\}.\\
\cref{plt:thm:cext-power} & the draft said that for h\_0=0 the formula "reduces to plt:eq:tay-C-bell", which is stated only for m>=r>=1; at r=0 it reduces instead to the boundary convention C\_\{m,0\}=\textbackslash delta\_\{m,0\}. Range added.\\
\cref{plt:cor:cext-reciprocal} & the draft asserted "then F\textasciicircum \textbackslash sigma = c\textasciicircum \textbackslash sigma t\textasciicircum \{\textbackslash sigma n\_0\}(1+W)\textasciicircum \textbackslash sigma" as though it were an identity. For non-integer \textbackslash sigma the symbol F\textasciicircum \textbackslash sigma has no prior meaning in the volume, and the right-hand side depends on the chosen value of c\textasciicircum \textbackslash sigma; it is a definition, and is now printed as one.\\
\cref{plt:rmk:cext-level-zero-repair} & S3's formula sheet displays this identity as [t\textasciicircum n](rev F - X + P) = ((-1)\textasciicircum \{n+1\}/(n+1)!) D(D-1)...(D-n+1) P\textasciicircum \{n+1\}. The added "+P" is a slip: it is correct for n >= 1, where P (a level-zero block) contributes nothing, but false at n = 0, where the left side is -P + P = 0 while the right side is -P. Removing the "+P", as in plt:eq:cext-level-zero-reversion, makes the formula correct for every n >= 0 and makes the Lambert case a specialization rather than an exception.\\
Operational formula sheet & this row printed the division recurrence with no hypothesis on B\_0, which repeats the defect of S1's sheet ("B\_0 \textbackslash ne 0"). In K[L]((t)) a nonzero polynomial block is not enough: B\_0 must be a unit of K[L], i.e. a nonzero scalar, or the quotient blocks leave K[L]. Hypothesis printed.\\
Operational formula sheet & this row read "\textbackslash PLrat is the smallest field of the tower with polynomial-log coefficients", which inverts the content of plt:thm:div-minimal-field: \textbackslash PLrat has RATIONAL-log coefficients, and the theorem says E((t)) is a field iff E = K(L). Restated as proved.\\
Exercises with solutions & source 1 asks only for a jump "by at least two", i.e. the case k = 2 of part (2), and does not ask for the converse. The stronger form is proved above because part (3) is what the sum rule of \textbackslash cref\{plt:thm:ring-valuation-laws\} actually asserts.\\
Exercises with solutions & source 2 proves the common lower bound by asserting that "the sequence is already Laurent-bounded at some initial precision". That is not a hypothesis of completeness; the bound has to be extracted from the Cauchy condition at cutoff N = 0, as in part (3).\\
Exercises with solutions & the sentence printed here was garbled ("ordt(...)-differences are ..."). What is meant is the outer order of the difference of two consecutive partial sums, which is what Cauchyness is about.\\
Exercises with solutions & source 2 states the prefactor form "whenever available" without an instance of unavailability; L + t is one, and it is the reason the log-degree profile of \textbackslash cref\{plt:def:lead-profile\} is defined blockwise rather than through a normal form.\\
Exercises with solutions & source 3 answers this item with "a natural home is K((L\textasciicircum \{-1\}))((t))", i.e. \textbackslash PLit. That overshoots by a whole class: the coefficients in \textbackslash eqref\{plt:eq:exr-recip-L-plus-t\} are Laurent *monomials* in L, so even the localisation K[L,L\textasciicircum \{-1\}]((t)) suffices, and among the four ambient classes the answer is \textbackslash PLrat. \textbackslash PLit is strictly larger (\textbackslash cref\{plt:prop:ring-PLit-strict\}), and the next item is what actually needs it.\\
Exercises with solutions & source 3 poses this product only through t\textasciicircum 2 and source 1 poses a different pair only through t\textasciicircum 3 without saying that the undisplayed blocks must be assumed zero. The two are merged here and the hypothesis is printed, since \textbackslash cref\{plt:thm:mul-finite-determination\} is the point of the exercise.\\
Exercises with solutions & source 5 prints the block Q\_3 of this example immediately after inputs given only modulo t\textasciicircum 3. The hypothesis B\_3 = 0 is a genuine extra assumption, and \textbackslash cref\{plt:prop:div-relative-precision\} shows exactly how much is missing; it is stated here rather than assumed silently.\\
Exercises with solutions & source 1 prints the Stirling identity with the sum starting at j = 1. That is correct only for k >= 1; the uniform range j = 0..k with the empty-product convention covers k = 0 as well, and is used above.\\
Exercises with solutions & this sentence cited \textbackslash cref\{plt:prop:alg-shift\}, the polynomial-shift formula, which has nothing to do with the truncation convention. The discrepancy is recorded with the sharp budget itself.\\
Exercises with solutions & source 1 records the budget as k <= floor((N - r)/(h + 1)) without saying which truncation convention its N refers to. Under the exclusive cutoff, which the catalogue makes normative, that form admits one Taylor term too many. Part (3) prints the conversion rather than silently shifting the index, because the same formula is what a routine would apply in reverse to size its inputs, where the extra term is not free.\\
Exercises with solutions & the exercise as inherited asked for the failure "for every F whose coefficient support is infinite". That quantifier is too strong, and the following paragraph is the counterexample to it.\\
Exercises with solutions & source 1 states this with the hypothesis q >= 1. The hypothesis is unnecessary: the Picard estimate \textbackslash eqref\{plt:eq:rev-picard-accuracy\} holds for every alpha = q >= 0, and the case q = 0 -- the Wright omega case -- is exactly the one in which the statement is sharp.\\
Exercises with solutions & the check was printed as "a = aL with a in K", using the letter a for both the coefficient a(L) and the scalar; renamed.\\
Exercises with solutions & sources 2 and 3 both quote a general third-order inverse block; source 3's version (the b = c = 0 case above) is correct, source 2's is not (\textbackslash cref\{plt:exer:exr-trap-universal-formula\}). Source 3 checks the degree law only at N = 1, 2; the block b\_3 above is computed here and extends the check to N = 3.\\
Exercises with solutions & part (3) compares moduli, so it needs an absolute value on the coefficient field; the inherited statement left K unnamed.\\
Exercises with solutions & the O-term in \textbackslash eqref\{plt:eq:exr-xlogx\} is an ANALYTIC estimate, and the computation above is formal. What licenses it is the explicit remainder bound of the Lambert part, not \textbackslash cref\{plt:thm:om-expansion\}; the citation is supplied rather than left to the reader, in line with \textbackslash cref\{plt:exer:exr-trap-formal-to-analytic\}.\\
Exercises with solutions & the two numbers in \textbackslash cref\{plt:warn:lag-scale-trap\} were computed at two different cutoffs. Recomputed in exact arithmetic at Y = 20 + log 10: N <= 2 -> correct 10.0000781540, naive 9.6702691781 N <= 3 -> correct 10.0000043732, naive 9.6702070480 The warningbox says "through Y\textasciicircum \{-3\}" and then quotes 10.000004 (which is the N <= 3 value) alongside 9.67026 (which is the N <= 2 value). Both truncations are quoted at N <= 3 above. The conclusion of the warningbox is unaffected; only the second numeral changes.\\
Exercises with solutions & "exactly item (i)" contradicted the sentence just above, which already observed that (v) also evaluates correctly at n = 0.\\
Exercises with solutions & the inherited solution asserted that the sign of omega - Omega\_1 itself changes at x = e\textasciicircum 2. That is the sign change of the LEADING TERM only. The true difference crosses zero at x = 8.7465..., recomputed here; the gap between the two numbers is the whole moral of the exercise, so printing e\textasciicircum 2 as the exact threshold defeats it.\\
Exercises with solutions & the inherited solution claimed both errors agree "to one digit" with LamP\_2(L)x\textasciicircum \{-2\}. At x = 10 the leading term is 3.484e-3 against a true error of 1.747e-3 -- a factor of two, not one digit. Only the signs agree at that point, and the reason is printed below.\\
Exercises with solutions & the inherited text added "and nothing analogous is claimed for WrightOmega". That is false: the same lemma (plt:lem:lw-two-variable) applied at sigma = +1/x rather than -1/S gives the omega statement, and it is printed in the volume as plt:thm:lw-basepoint / plt:cor:lw-canonical. The claim is corrected here and in plt:exer:exr-omega-numerics, which repeated it.\\
Exercises with solutions & source 1's prose ("the same P\_n occur, multiplied by (-1)\textasciicircum n") contradicts its own displayed expansion of the inverse of X - log X, in which P\_1 = L occurs with a plus sign. The weight (-1)\textasciicircum \{n+1\} belongs to rev(X - L); the weight (-1)\textasciicircum n belongs to its negative, which is what W\_\{-1\} is. Part (3) separates them.\\
Exercises with solutions & item (vi) of \textbackslash cref\{plt:cor:lambert-coefficients\} is indexed from the top, [u\textasciicircum \{n-2\}], not at the fixed exponent 3. The two coincide at n = 5 and only there; read as "[u\textasciicircum 3]" the check would fail at n = 6, where [u\textasciicircum 3]LamP\_6 = -85/6 while the right-hand side is 75/8 = [u\textasciicircum 4]LamP\_6. The inherited solution printed the fixed exponent and verified (vi) on LamP\_5 only, although the exercise asks for both.\\
Exercises with solutions & source 5 asserts that the residual-to-error principle "proves the asymptotic validity of every truncation directly from" the formal residual identity, and source 1 asserts an analytic O-bound for the branch residual "by construction". Neither is a proof; the missing ingredient is the substituted-tail estimate, which is why this exercise exists.\\
Exercises with solutions & the inherited version of this exercise asserted outright that the minimal condition "does not determine the blocks", and proved it by changing one block and checking the condition at that index alone. The check at index N+1 then fails, so the modified family does not satisfy the definition and the claimed nonuniqueness is not there. Part (2) below proves uniqueness; the defects the exercise exists to expose are the two in part (3), and they are genuine.\\
Exercises with solutions & this step was written as an inequality on "(F-G)(X)", which is a formal series and not a function. What the two estimates bound is the difference of the two TRUNCATIONS, which is a polynomial expression.\\
Exercises with solutions & the display here was garbled ("1 - LamP\_0(0)(-1) + ..."). By \textbackslash cref\{plt:def:om-truncation\}, Omega\_N = X + Trunc(P)\_\{N+1\}, so the evaluation is the single sum written below.\\
\cref{plt:rmk:exr-provenance} & the inherited solution stopped at the third item and concluded that "this volume does not determine whether that series converges". That is false. \textbackslash cref\{plt:thm:lw-basepoint\}, taken at K = x and log K = L, gives M = L and a seating condition 2(1+L) < x, under which the substituted series converges absolutely. At x = 20 the condition reads 7.99 < 20. Recomputed: the partial sum through n = 40 agrees with omega(20) - 20 + log 20 to fifteen digits.\\
(front matter) & the six sources cited this work inconsistently as a 2010 preprint and as a 2013 journal article; the published version is given, with the preprint retained as a secondary pointer.\\
\bottomrule
\end{longtable}

\begin{remark}[What is not in this ledger]
\label{plt:rmk:repairs-scope}
Three classes of change are deliberately excluded.  Deduplication is not a
repair: where several sources proved the same result, the strongest statement
and the best proof were kept and the rest collapsed, which the provenance
section records in aggregate rather than line by line.  Restructuring is not a
repair.  And a source claim that this volume could neither prove nor refute has
not been silently dropped: it is stated here as an explicitly labelled open
problem or conjecture, with the source's assertion quoted as an assertion.
\end{remark}

% ---- fragment 21-formalization ----
\chapter{Formalization register}
\label{plt:sec:formalization}

\section{What this register claims}

No proof in this volume is machine-checked.  The corpus that hosts it does
contain a large Lean development, and the honest question is not whether these
results are formalized --- they are not --- but which of them have a formalized
\emph{neighbour}, and what a formalization campaign would have to build first.
This register answers that.  Every Lean name below was checked to resolve in
the corpus; a name that no longer resolves is a defect, and the corpus audit
\texttt{scripts/audit\_\allowbreak crosswalk\_\allowbreak names.\allowbreak py} enforces it.

Three statuses are used.  \textbf{Crosswalked} means a Lean declaration states
the same mathematical fact, possibly in a different but faithfully translated
setting.  \textbf{Neighbouring} means a Lean declaration states a genuinely
different theorem that a reader might mistake for this one; the entry exists to
prevent that mistake.  \textbf{Absent} means nothing in the corpus is close.

\section{The formal algebra: absent}

Nothing in Parts~I--IV corresponds to a Lean declaration.  The corpus has no
polynomial--logarithmic ring, no two-generator scale, no outer valuation, no
distinguished derivation \(\DX\), no composition of transseries, and no
reversion at infinity.  Concretely, a formalization campaign would have to
build, in order:

\begin{enumerate}
\item the ring \(\PLpow=\K[L][[t]]\) with \(t\) and \(L\) as independent
  generators, and its Laurent extension (\cref{plt:sec:rings});
\item the outer valuation and the leading-data trio, which are three different
  objects and are easy to conflate in a formal setting
  (\cref{plt:sec:leading-data});
\item the derivation \(\DX\) characterized by \(\DX t=-t^{2}\),
  \(\DX L=t\), and the repeated-derivative operator \(\shiftop{n}{k}\)
  (\cref{plt:sec:differential});
\item composition by generator substitution, which is where an informal
  treatment most often goes wrong, since \(t\mapsto t+H\) is \emph{not} the
  substitution meant (\cref{plt:sec:composition});
\item the exact formal Taylor formula (\cref{plt:thm:tay-formula}), whose proof
  here is an algebra-homomorphism argument on the two generators and should
  formalize without the coefficient bookkeeping;
\item the reversion fixed point and its uniqueness
  (\cref{plt:thm:rev-existence}).
\end{enumerate}

Only after those is the Lagrange--B\"urmann theorem at infinity
(\cref{plt:thm:lag-burmann}) even statable.

\section{Lagrange--B\"urmann: neighbouring, not crosswalked}

The corpus \emph{does} formalize Lagrange--B\"urmann inversion --- for formal
power series at the origin, over a commutative \(\RationalNumbers\)-algebra.
The relevant declarations are \texttt{Fabius.\allowbreak Lagrange.\allowbreak solution},
\texttt{Fabius.\allowbreak Lagrange.\allowbreak coeff\_\allowbreak solution}, and the division-free identity
\texttt{Fabius.\allowbreak Lagrange.\allowbreak coeff\_\allowbreak subst\_\allowbreak mul\_\allowbreak derivative}.

That is a different theorem from \cref{plt:thm:lag-burmann}, and the difference
is not cosmetic:

\begin{itemize}
\item the formalized version inverts \(g=z\varphi(g)\) in \(R[[z]]\), a
  one-generator power-series ring with the ordinary derivative; the version
  here inverts \(X+A\) in the two-generator polynomial--logarithmic scale with
  the derivation \(\DX\);
\item the formalized version is a statement about coefficients of \(z^{n}\);
  the version here is a statement about coefficient \emph{blocks}
  \(p_{n}(L)\), each a polynomial whose degree is unbounded across \(n\);
\item summability is automatic at the origin and is a theorem here
  (\cref{plt:prop:lag-summability-sharp}).
\end{itemize}

So the corpus's Lagrange inversion cannot be cited as proving the at-infinity
formula, and this volume does not do so.  It is, however, the natural model
for a future formalization, and its proof strategy --- purely algebraic,
avoiding formal residues, which Mathlib does not have --- is the same strategy
the proof here uses.

\section{The Lambert \texorpdfstring{\(\LambertW\)}{W} series: neighbouring}

\texttt{Fabius.lambertW} and \texttt{Fabius.coeff\_lambertW} formalize the
\emph{small-argument} expansion
\(\LambertW(z)=\sum_{n\ge1}(-n)^{n-1}z^{n}/n!\), obtained from the formalized
Lagrange inversion above.  This volume expands \(\LambertW\) at \emph{infinity}
(\cref{plt:sec:lambert-branches}).  The two expansions share a name, a
function, and nothing else: one is a convergent Taylor series at the origin
with radius \(\EulerE^{-1}\), the other an asymptotic polynomial--logarithmic
expansion at infinity that converges nowhere.  Confusing them is the single
most likely misreading of this volume's Part~V, which is why the entry is here.

\section{The real branches: crosswalked}

The analytic facts about the two real branches that this volume \emph{uses}
--- as opposed to derives --- are kernel-verified:

\begin{longtable}{@{}p{0.46\textwidth} p{0.48\textwidth}@{}}
\caption{Analytic branch facts used in Part~V and their Lean declarations.}
\label{plt:tab:lean-branches}\\
\toprule
Fact & Declaration\\
\midrule
\endfirsthead
\toprule
Fact & Declaration\\
\midrule
\endhead
The two real branches exist and are characterized by the defining equation &
  \texttt{Fabius.\allowbreak principalLambertW},
  \texttt{Fabius.lowerLambertW},
  \texttt{Fabius.\allowbreak principalLambertW\_\allowbreak unique},
  \texttt{Fabius.\allowbreak lowerLambertW\_\allowbreak unique}\\
Every real solution of \(w\EulerE^{w}=z\) is one of the two branches &
  \texttt{Fabius.\allowbreak eq\_\allowbreak principalLambertW\_\allowbreak or\_\allowbreak eq\_\allowbreak lowerLambertW}\\
Continuity on the full closed domains, branch point included &
  \texttt{Fabius.\allowbreak principalLambertW\_\allowbreak continuousOn\_\allowbreak Ici},
  \texttt{Fabius.\allowbreak lowerLambertW\_\allowbreak continuousOn\_\allowbreak Ico}\\
The derivative of the principal branch, and \(\LambertW_{0}''(0)=-2\) &
  \texttt{Fabius.\allowbreak deriv\_\allowbreak principalLambertW},
  \texttt{Fabius.\allowbreak deriv\_\allowbreak deriv\_\allowbreak principalLambertW\_\allowbreak zero}\\
Strict concavity of the principal branch &
  \texttt{Fabius.\allowbreak strictConcaveOn\_\allowbreak principalLambertW}\\
The square-root scale at the branch point, and the signed leading laws &
  \texttt{Fabius.\allowbreak lambertWBranchPointScale},
  \texttt{Fabius.\allowbreak principalLambertW\_\allowbreak add\_\allowbreak one\_\allowbreak isEquivalent\_\allowbreak branchPoint},
  \texttt{Fabius.\allowbreak lowerLambertW\_\allowbreak add\_\allowbreak one\_\allowbreak isEquivalent\_\allowbreak branchPoint}\\
\bottomrule
\end{longtable}

The branch-point entry is the sharpest of these: it confirms, in the kernel,
the assertion of \cref{plt:sec:lambert-branches} that the inverse near
\(z=-\EulerE^{-1}\) is \emph{not} a power--logarithmic series but carries a
square-root scale, which is exactly the boundary of this volume's class.

\section{The lower-branch expansion near zero: partially crosswalked}

One expansion statement of Part~V has a genuine Lean counterpart.
\texttt{Fabius.\allowbreak tendsto\_\allowbreak lowerLambertW\_\allowbreak expansion} proves
\[
  \LambertW_{-1}(-\varepsilon)
  =\log\varepsilon-\log\AbsoluteValue{\log\varepsilon}+\LittleOAt{1}{\varepsilon\downarrow0},
\]
which in the coordinates of \cref{plt:sec:lambert-branches}, with
\(S=\log(1/\varepsilon)\) and \(L=\log S\), is
\(\LambertW_{-1}(-\varepsilon)=-S-L+\LittleOAt{1}{\varepsilon\downarrow0}\):
precisely the first two blocks of the lower-branch expansion derived there.

A second, sharper declaration pins the next block as well.
\texttt{Fabius.\allowbreak lowerLambertW\_\allowbreak near\_\allowbreak zero\_\allowbreak bounds} proves, for
\(x\in(-\EulerE^{-1},0)\) and \(\eta=\log(1/(-x))\),
\[
  -\eta-\frac{\eta\log\eta}{\eta-1}
  \ \le\ \LambertW_{-1}(x)\ <\ -\eta-\log\eta .
\]
In the coordinates above, \(\eta=S\) and \(\log\eta=L\), so the bracket says
that the deviation of \(\LambertW_{-1}\) from the two-block truncation
\(-S-L\) lies in \(\bigl(-L/(S-1),\,0\bigr]\).  That is a two-sided,
non-asymptotic statement: it fixes the sign of the third block and its order
\(L/S\), which is exactly what \cref{plt:sec:lambert-branches} derives
formally.  The strictness of the upper bound is the formal statement that the
third block is nonzero and negative.

The supporting declaration is worth naming separately, because it is the piece
most likely to be reusable:
\texttt{Fabius.\allowbreak SubLog.\allowbreak bracket\_\allowbreak of\_\allowbreak sub\_\allowbreak log\_\allowbreak eq} brackets the inverse of
\(y\mapsto y-\log y\) for \(y>1\) with no Lambert content at all, and
\texttt{Fabius.\allowbreak SubLog.\allowbreak strictMonoOn\_\allowbreak sub\_\allowbreak log} supplies the monotonicity that
makes that inverse well defined.  Since this volume's lower-branch treatment
\emph{runs} on the inverse of \(u\mapsto u-\log u\) rather than on
\(\LambertW_{-1}\) directly, that lemma --- not the branch-specific one --- is
the natural target for a future crosswalk of the whole lower-branch section.
The side condition \(\eta>1\) is itself a lemma
(\texttt{Fabius.\allowbreak one\_\allowbreak lt\_\allowbreak log\_\allowbreak one\_\allowbreak div\_\allowbreak neg}) rather than a hypothesis the
caller must discharge, so the \((\eta-1)\) denominator cannot be reached at
zero.

Two elementary bounds on the principal branch are also available and are used
here only as sanity constraints:
\texttt{Fabius.\allowbreak principalLambertW\_\allowbreak elementary\_\allowbreak bounds} gives
\(x/(1+x)\le\LambertW_{0}(x)\le\log(1+x)\le x\) for \(x\ge0\), resting on the
Lambert-free inequality
\texttt{Fabius.\allowbreak exp\_\allowbreak le\_\allowbreak one\_\allowbreak add\_\allowbreak mul\_\allowbreak exp\_\allowbreak self}, valid for every real
argument.

Beyond the third block --- the \(\LamP{n}\) coefficients for \(n\ge2\), the
remainder order at every order, and any uniformity --- everything is
human-proved here and unformalized.

\section{All derivatives of the real branches: crosswalked}

The derivative formulas used in \cref{plt:sec:lambert-branches} are
kernel-verified to every order.  \texttt{Fabius.\allowbreak iteratedDeriv\_\allowbreak principalLambertW}
proves, for \(x>-\EulerE^{-1}\) and every \(n\in\NaturalNumbers\),
\[
 \frac{\Differential^{\,n+1}}{\Differential x^{\,n+1}}\LambertW_{0}(x)
 =\frac{\EulerE^{-(n+1)W}\,P_{n}(W)}{(1+W)^{2n+1}},
 \qquad W=\LambertW_{0}(x),
\]
and \texttt{Fabius.\allowbreak iteratedDeriv\_\allowbreak lowerLambertW} proves the
same formula for \(\LambertW_{-1}\) on \((-\EulerE^{-1},0)\).  The polynomial
family is \texttt{Fabius.\allowbreak lambertDerivPoly}, defined by
\(P_{0}=1\) and the recurrence
(\texttt{Fabius.\allowbreak lambertDerivPoly\_\allowbreak succ})
\[
 P_{n+1}(w)=(1+w)P_{n}'(w)-\bigl((n+1)w+(3n+2)\bigr)P_{n}(w),
\]
whose first members, recomputed here in exact arithmetic from that
recurrence, are
\[
 P_{1}=-(w+2),\quad
 P_{2}=2w^{2}+8w+9,\quad
 P_{3}=-(6w^{3}+36w^{2}+79w+64),\quad
 P_{4}=24w^{4}+192w^{3}+622w^{2}+974w+625,
\]
the first two being the kernel's own evaluated cases
(\texttt{Fabius.\allowbreak eval\_\allowbreak lambertDerivPoly\_\allowbreak one},
\texttt{\ldots\_\allowbreak two}).  Two conventions matter.  The Lean index is
shifted by one against the classical one: the classical \(P_{n}\) of
\cite{corless-et-al-1996} is Lean's \(P_{n-1}\), and the formula above is for
the \((n+1)\)-st derivative.  And \(P_{1}=-(w+2)\) is exactly what one
obtains by differentiating \(\LambertW'=\EulerE^{-W}/(1+W)\) once, which
fixes the sign convention.  The engine behind the theorem carries no Lambert
content and is the reusable piece:
\texttt{AutonomousODE.\allowbreak iteratedDeriv\_\allowbreak eq\_\allowbreak comp} proves
that if \(W'=\varphi\circ W\) on an open set then every iterated derivative
is \(G_{n}\circ W\) with \(G_{0}=\mathrm{id}\) and
\(G_{n+1}=G_{n}'\varphi\).  This entry is crosswalked, not neighbouring: it
is the same statement in the same normalization, and the Lambert
\(\LambertW\) guide's own crosswalk appendix records the identical index
convention.

\section{Summary}

Of the results in this volume, none is crosswalked in the strong sense; one
two-term asymptotic statement is; six analytic facts that the volume cites
rather than proves are.  The formal algebra of Parts~I--IV is absent from the
corpus in its entirety, and the ordered list at the head of this register is what a campaign
would need to build.
This register is deliberately conservative: where the corpus proves a
neighbouring theorem, that is recorded as a neighbour and not as coverage.

% ---- fragment 98-provenance ----
\chapter{Provenance of the consolidation}
\label{plt:sec:provenance}

\section{What was merged}

This volume is an editorial consolidation, carried out on 2~September 2026, of
six independently written article packages that arrived together as bare
directories in the direct-arrival commit
\texttt{730e1763291099cd\allowbreak 50ca1e20ed2c62c3\allowbreak 8d95ab4f} and were filed on
1~September 2026 as standalone quick-gate receipts.  No archive or checksum
ledger accompanied them.  All six treat the same subject -- the algebra,
division, composition, and asymptotic reversion of polynomial--logarithmic
transseries, with Lambert \(\LambertW\) as the guiding example -- from
independent starting points, with substantial overlap and no containment
relation among them.

\emph{Editorial} merge means: every result is stated once, in the strongest
form that the six sources jointly support and that we can prove; the best
available proof is retained and gaps are filled; conflicting conventions are
resolved against the corpus notation catalogue; errors are corrected and the
corrections are marked in the source with \verb|% ed.:| comments; and every
source-specific layer that is not redundant is preserved.  Nothing was
discarded merely for being repeated: repetition was collapsed, content was not.

\section{Absorbed sources and their receipts}

Each source is recorded twice: as filed at intake, and as actually absorbed.
The two differ by exactly three bytes per file because a regrouping on
2~September 2026 moved the packages one directory level deeper and the shared
notation input path was repaired from four to five levels
(\cref{plt:tab:provenance}).  No mathematical content changed between the two
states.

\begin{table}[htbp]
\centering
\small
\caption{The six absorbed sources.  ``Intake'' is the receipt recorded in the
group manifest on 1~September 2026; ``absorbed'' is the state merged into this
volume.  The retained arrival PDFs were historical publication checkpoints
that already predated their own sources and were not rebuilt for this merge.}
\label{plt:tab:provenance}
\begin{tabularx}{\textwidth}{@{}lrrX@{}}
\toprule
Source & Lines & Bytes & SHA-256 \\
\midrule
\multicolumn{4}{@{}l}{\textbf{S1} \(=\)
  \texttt{Polynomial-\allowbreak Logarithmic-\allowbreak Transseries-1.\allowbreak tex}} \\
\quad intake    & 4{,}020 & 187{,}071 &
  \texttt{01a03e0934d05b8b\allowbreak 94369da9656fc0da\allowbreak 4dde9b7c559b4e07\allowbreak 38075484e14f4e61} \\
\quad absorbed  & 4{,}020 & 187{,}074 &
  \texttt{4fd42b7e344c485f\allowbreak 8d878279dc3b9a6d\allowbreak 5dc326a9c739114e\allowbreak 443fd9c07a5eb2b8} \\
\addlinespace
\multicolumn{4}{@{}l}{\textbf{S2} \(=\)
  \texttt{polynomial\_\allowbreak logarithmic\_\allowbreak transseries-2.\allowbreak tex}} \\
\quad intake    & 5{,}006 & 173{,}396 &
  \texttt{aa25baa03099472d\allowbreak f765d35c03a2bb26\allowbreak 5dbe9af5aab500ec\allowbreak 0a4a04a378ef83e5} \\
\quad absorbed  & 5{,}006 & 173{,}399 &
  \texttt{d172b81453a80536\allowbreak 6f056f6114bb6df6\allowbreak fdf7811466413201\allowbreak 85530940754bdefe} \\
\addlinespace
\multicolumn{4}{@{}l}{\textbf{S3} \(=\)
  \texttt{Polynomial\_\allowbreak Logarithmic\_\allowbreak Transseries-3.\allowbreak tex}} \\
\quad intake    & 4{,}249 & 150{,}182 &
  \texttt{0962c15683cb4610\allowbreak d45961c227f6e9b1\allowbreak 02d66ee37073432d\allowbreak 1943d31a1e6ad348} \\
\quad absorbed  & 4{,}249 & 150{,}185 &
  \texttt{d2a8bd25d71cc41f\allowbreak c81c1c75fb66ca75\allowbreak 496b8425ac496ce1\allowbreak 05d3220e1ea07e3f} \\
\addlinespace
\multicolumn{4}{@{}l}{\textbf{S4} \(=\)
  \texttt{Polynomial-\allowbreak Logarithmic-\allowbreak Transseries-4.\allowbreak tex}} \\
\quad intake    & 3{,}132 & 120{,}607 &
  \texttt{387ca51f94292a99\allowbreak 6fb339ec6dfb3477\allowbreak f68575255b2aaaf5\allowbreak 39fcffb1c32ff564} \\
\quad absorbed  & 3{,}132 & 120{,}610 &
  \texttt{37026b625cf7b2fd\allowbreak 106005b8c89f720b\allowbreak 8b724f99a5446ed8\allowbreak 4ca52531c6e2b18f} \\
\addlinespace
\multicolumn{4}{@{}l}{\textbf{S5} \(=\)
  \texttt{Polynomial\_\allowbreak Logarithmic\_\allowbreak Transseries-5.\allowbreak tex}} \\
\quad intake    & 2{,}443 & 106{,}141 &
  \texttt{1149ae6884e63dc7\allowbreak 7b747c786d045541\allowbreak 9b6c79cdd3b33dc9\allowbreak 15bf3875cc8d26cc} \\
\quad absorbed  & 2{,}443 & 106{,}144 &
  \texttt{187bd510f33ea429\allowbreak 6f0247b286d0fcae\allowbreak 6200ce0b5bd175c6\allowbreak c8706a37652dc17a} \\
\addlinespace
\multicolumn{4}{@{}l}{\textbf{S6} \(=\)
  \texttt{polynomial\_\allowbreak logarithmic\_\allowbreak transseries-6.\allowbreak tex}} \\
\quad intake    & 4{,}388 & 155{,}846 &
  \texttt{df4e4bc5932a59e4\allowbreak ed265b32d8463edc\allowbreak 65f67cbd2735b349\allowbreak 02acf479647f8491} \\
\quad absorbed  & 4{,}388 & 155{,}849 &
  \texttt{16978aae6fb1e458\allowbreak e236cac9807ab81b\allowbreak b5314fdb087a2504\allowbreak a416867a4d43b4c4} \\
\bottomrule
\end{tabularx}
\end{table}

The six retained arrival PDFs, which are historical artefacts and are not
renders of the absorbed sources, had SHA-256 values
\texttt{a4fc4af0\ldots69e886} (S1, 119 custom 522-by-738-point pages),
\texttt{5e9ff596\ldots bfc68e} (S2, 102 custom pages),
\texttt{3f7c4bc1\ldots58a4af} (S3, 87 Letter pages),
\texttt{c2d75b35\ldots45a3eb} (S4, 47 A4 pages),
\texttt{189e95ab\ldots58a2db} (S5, 44 Letter pages), and
\texttt{b5142bad\ldots467aa} (S6, 100 A4 pages).  Git history is the archive
for all twelve files.

\section{What each source contributed}

No source was a superset of the others, and each contributed at least one layer
that no other supplied.

\begin{itemize}
\item \textbf{S1} contributed the block-algebra presentation of the scale, the
  systematic comparison of four reversal algorithms, the residual-certificate
  and reproducible-verification discipline, the Dulac and power--log normal-form
  literature, and the largest exercise set.
\item \textbf{S2} contributed the \(t\)-adic completeness treatment, the
  normal-ordered composition formalism, the power--log series near a finite
  point with the Dulac conventions, the derivation of the universal composition
  formulas, and the pitfalls-and-diagnostics decision procedure.
\item \textbf{S3} contributed the formal summability criterion, the degree law
  for polynomial perturbations, the Lagrange--B\"urmann transport formula, the
  fully formal proof of Lagrange--B\"urmann inversion at infinity, and the
  reusable inverse templates beyond Lambert.
\item \textbf{S4} contributed the sharpest statement of the unit criterion and
  the field property, the finite reversal certificate, and closed-form
  coefficient extraction.
\item \textbf{S5} contributed the tightest and most completely proved core --
  it is the spine of Parts~I--IV -- together with the Laplace transform of the
  Lambert polynomials, the resulting exponential-moment cancellations, and the
  generalized Lambert expansion.
\item \textbf{S6} contributed the filtered Lipschitz gain underlying the
  fixed-point construction, the composed-observable form of
  Lagrange--B\"urmann, the systematic change of variable between infinity and a
  finite point, and the failure-mode catalogue.
\end{itemize}

\section{Conventions reconciled}

The six sources used mutually incompatible local conventions.  All are resolved
here in favour of the corpus notation catalogue; the concordance in
\cref{plt:sec:notation} records the mapping in full.  The reconciliations that
change a formula rather than only a glyph are:

\begin{itemize}
\item \textbf{Truncation.}  S6 used the inclusive cutoff \([F]_{\le N}\); every
  other source and the catalogue use the exclusive
  \(\Trunc{F}{N}\).  The relation is
  \([F]_{\le N}=\Trunc{F}{N+1}\), so every S6 statement transported here has
  had its cutoff index shifted by one.
\item \textbf{Lambert polynomial indexing.}  The sources split between a
  zero-based family with \(\LamP{0}(u)=-u\) and a one-based family that
  separates the logarithmic block.  The catalogue's zero-based normalization is
  used throughout; the degree law \(\deg\LamP{n}=n\) is therefore stated only
  for \(n\ge1\), since \(\LamP{0}=-u\) has degree~\(1\).
\item \textbf{Coefficient extraction.}  Two sources printed \([M]F\) and two
  printed \([F]_{M}\) for the same operation.  This volume writes the series to
  the right of the bracket throughout.
\item \textbf{The three \(\BigO\) roles.}  Formal valuation tails, analytic
  estimates, and algorithmic cost bounds are printed with three different
  tokens and are never interchanged.  Several source statements that used one
  glyph for all three have been split accordingly.
\end{itemize}

\section{Status of the claims}

Every theorem, proposition, lemma, and corollary in this volume carries a
proof.  Where a source asserted a result without proof, the proof is supplied
here or the assertion is demoted to an explicitly labelled conjecture or open
problem.  Where a source claimed analytic validity on the strength of formal
algebra alone, the claim has been weakened to what the algebra establishes and
the missing analytic hypothesis is named.  Corrections to the sources are
marked in this document's own source with \verb|% ed.:| comments at the point
of repair, and are collected in \cref{plt:sec:repairs}.

Each section was drafted and then reviewed by an independent adversarial pass
instructed to refute rather than confirm, with two exceptions that are
recorded here rather than hidden: the Lambert \(\LambertW\) branch section
(\cref{plt:sec:lambert-branches}) and the exercises
(\cref{plt:sec:exercises}) did not receive that pass.  In its place, every
claim in those two sections that admits exact or high-precision computation
was recomputed independently of the text: the branch-point Puiseux
coefficients through \(p^{6}\) by exact series reversion; the transverse
coefficient identity through \(a=5\); the uniform remainder bound
\eqref{plt:eq:lw-remainder} at six base points, canonical and off-centre, for
\(N\le8\), with a worst observed error-to-bound ratio below \(0.1\); the
Cauchy bound \eqref{plt:eq:lw-cauchy-bound} through \(a+b=12\); the
two-variable fixed point on \(160\) complex sample points; the \(X+aL^{2}\)
reversion through \(n=9\); and the numerical values quoted in the exercises
to every printed digit.  What those computations cannot test --- the prose
of the proofs --- has been read but not adversarially reviewed, and a reader
should weight those two sections accordingly.

None of the proofs in this volume is a machine-checked proof.  The corpus
formalizes the real Lambert pair and its branch-point geometry in Lean, and the
concordance in \cref{plt:sec:notation} records which statements here correspond
to declarations that exist there; the remaining results are human-proved
frontier mathematics.

\clearpage
\begin{thebibliography}{99}
\addcontentsline{toc}{section}{References}

\bibitem{adh-book}
M.~Aschenbrenner, L.~van den Dries, and J.~van der Hoeven,
\emph{Asymptotic Differential Algebra and Model Theory of Transseries},
Annals of Mathematics Studies, vol.~195, Princeton University Press,
Princeton, 2017.
\href{https://doi.org/10.1515/9781400885411}{doi:10.1515/9781400885411}.

\bibitem{cohen-branch-identities}
H.~Cohen,
``Lambert \(\LambertW\)-function branch identities,''
arXiv:2012.11698, revised 2021.

\bibitem{comtet1970}
L.~Comtet,
``Inversion de \(y^{\alpha}\EulerE^{y}\) et \(y\log y\),''
\emph{Comptes Rendus de l'Acad\'emie des Sciences Paris, S\'erie A}
\textbf{270} (1970), 1085--1088.

\bibitem{comtet-book}
L.~Comtet,
\emph{Advanced Combinatorics: The Art of Finite and Infinite Expansions},
D.~Reidel, Dordrecht, 1974.

\bibitem{corless-et-al-1996}
R.~M. Corless, G.~H. Gonnet, D.~E.~G. Hare, D.~J. Jeffrey, and D.~E. Knuth,
``On the Lambert \(\LambertW\) function,''
\emph{Advances in Computational Mathematics} \textbf{5} (1996), 329--359.
\href{https://doi.org/10.1007/BF02124750}{doi:10.1007/BF02124750}.

\bibitem{corless-jeffrey-knuth-1997}
R.~M. Corless, D.~J. Jeffrey, and D.~E. Knuth,
``A sequence of series for the Lambert \(\LambertW\) function,''
in \emph{Proceedings of the 1997 International Symposium on Symbolic and
Algebraic Computation (ISSAC~'97)}, ACM, New York, 1997, 197--204.
\href{https://doi.org/10.1145/258726.258783}{doi:10.1145/258726.258783}.

\bibitem{wright-omega}
R.~M. Corless and D.~J. Jeffrey,
``The Wright \(\omega\) function,''
in \emph{Artificial Intelligence, Automated Reasoning, and Symbolic
Computation}, Lecture Notes in Computer Science, vol.~2385, Springer, Berlin,
2002, 76--89.
\href{https://doi.org/10.1007/3-540-45470-5_10}{doi:10.1007/3-540-45470-5\_10}.

\bibitem{costin-book}
O.~Costin,
\emph{Asymptotics and Borel Summability},
Chapman \& Hall/CRC Monographs and Surveys in Pure and Applied Mathematics,
vol.~141, Boca Raton, 2009.

\bibitem{dahn-goering-1987}
B.~I. Dahn and P.~G\"oring,
``Notes on exponential-logarithmic terms,''
\emph{Fundamenta Mathematicae} \textbf{127} (1987), no.~1, 45--50.
\href{https://doi.org/10.4064/fm-127-1-45-50}{doi:10.4064/fm-127-1-45-50}.

\bibitem{debruijn}
N.~G. de Bruijn,
\emph{Asymptotic Methods in Analysis}, 3rd ed.,
Dover Publications, New York, 1981.

\bibitem{dlmf-W}
NIST Digital Library of Mathematical Functions,
\S4.13, ``Lambert \(\LambertW\)-function,''
F.~W.~J. Olver, A.~B. Olde Daalhuis, D.~W. Lozier, B.~I. Schneider,
R.~F. Boisvert, C.~W. Clark, B.~R. Miller, B.~V. Saunders, H.~S. Cohl, and
M.~A. McClain, eds., release 1.2.7 of 15 June 2026,
\url{https://dlmf.nist.gov/4.13}, accessed 31 August 2026.

\bibitem{ecalle}
J.~\'Ecalle,
\emph{Les Fonctions R\'esurgentes}, 3 vols.,
Publications Math\'ematiques d'Orsay, Orsay, 1981--1985.

\bibitem{edgar-beginners}
G.~A. Edgar,
``Transseries for beginners,''
\emph{Real Analysis Exchange} \textbf{35} (2009/2010), no.~2, 253--310.
\href{https://doi.org/10.14321/realanalexch.35.2.0253}{doi:10.14321/realanalexch.35.2.0253};
arXiv:0801.4877.

\bibitem{edgar-composition}
G.~A. Edgar,
``Transseries: composition, recursion, and convergence,''
arXiv:0909.1259, 2009.

% ed.: the six sources cited this work inconsistently as a 2010 preprint and as
% a 2013 journal article; the published version is given, with the preprint
% retained as a secondary pointer.
\bibitem{edgar-fractional}
G.~A. Edgar,
``Fractional iteration of series and transseries,''
\emph{Transactions of the American Mathematical Society} \textbf{365} (2013),
no.~11, 5805--5832;
arXiv:1002.2378.

\bibitem{flajolet-sedgewick}
P.~Flajolet and R.~Sedgewick,
\emph{Analytic Combinatorics},
Cambridge University Press, Cambridge, 2009.

\bibitem{gessel-lagrange}
I.~M. Gessel,
``Lagrange inversion,''
\emph{Journal of Combinatorial Theory, Series A} \textbf{144} (2016),
212--249;
arXiv:1609.05988.

\bibitem{graham-knuth-patashnik}
R.~L. Graham, D.~E. Knuth, and O.~Patashnik,
\emph{Concrete Mathematics: A Foundation for Computer Science}, 2nd ed.,
Addison--Wesley, Reading, MA, 1994.

\bibitem{hahn}
H.~Hahn,
``\"Uber die nichtarchimedischen Gr\"o\ss ensysteme,''
\emph{Sitzungsberichte der Kaiserlichen Akademie der Wissenschaften Wien,
Mathematisch-Naturwissenschaftliche Klasse, Abteilung IIa}
\textbf{116} (1907), 601--655.

\bibitem{hardy-orders}
G.~H. Hardy,
\emph{Orders of Infinity: The `Infinit\"arcalc\"ul' of Paul du Bois-Reymond},
Cambridge Tracts in Mathematics and Mathematical Physics, no.~12,
Cambridge University Press, Cambridge, 1910.

\bibitem{jeffrey-et-al-1995}
D.~J. Jeffrey, R.~M. Corless, D.~E.~G. Hare, and D.~E. Knuth,
``Sur l'inversion de \(y^{\alpha}\EulerE^{y}\) au moyen des nombres de
Stirling associ\'es,''
\emph{Comptes Rendus de l'Acad\'emie des Sciences Paris, S\'erie I,
Math\'ematique} \textbf{320} (1995), 1449--1452;
English version, arXiv:math/9512230.

\bibitem{mantova-monotonicity}
V.~Mantova,
``Monotonicity and a Taylor approximation theorem for transseries,''
arXiv:2601.07747, version~2, 2026.

\bibitem{mardesic-normalforms}
P.~Marde\v{s}i\'c, M.~Resman, J.-P. Rolin, and V.~\v{Z}upanovi\'c,
``Normal forms and embeddings for power-log transseries,''
\emph{Advances in Mathematics} \textbf{303} (2016), 888--953.
\href{https://doi.org/10.1016/j.aim.2016.08.030}{doi:10.1016/j.aim.2016.08.030}.

\bibitem{mardesic-length}
P.~Marde\v{s}i\'c, M.~Resman, J.-P. Rolin, and V.~\v{Z}upanovi\'c,
``The length of \(\varepsilon\)-neighborhoods of orbits of Dulac maps,''
arXiv:1606.02581, version~3, 2018.

\bibitem{olver-asymptotics}
F.~W.~J. Olver,
\emph{Asymptotics and Special Functions},
AKP Classics, A~K Peters, Wellesley, MA, 1997.

\bibitem{peran-logtransseries}
D.~Peran,
``Normalization of strongly hyperbolic logarithmic transseries and complex
Dulac germs,''
arXiv:2302.14527, 2023.

\bibitem{resman-dulac}
P.~Marde\v{s}i\'c and M.~Resman,
``Analytic moduli for parabolic Dulac germs,''
\emph{Russian Mathematical Surveys} \textbf{76} (2021), no.~3, 389--460.

\bibitem{salvy-shackell-1992}
B.~Salvy and J.~Shackell,
``Asymptotic expansions of functional inverses,''
in P.~S. Wang, ed., \emph{Symbolic and Algebraic Computation: Proceedings of
ISSAC~'92}, ACM Press, New York, 1992, 130--137.

\bibitem{salvy-shackell-1998}
B.~Salvy and J.~Shackell,
``Symbolic asymptotics: functions of two variables, implicit functions,''
\emph{Journal of Symbolic Computation} \textbf{25} (1998), no.~3, 329--349.
\href{https://doi.org/10.1006/jsco.1997.0179}{doi:10.1006/jsco.1997.0179}.

\bibitem{salvy-shackell-1999}
B.~Salvy and J.~Shackell,
``Symbolic asymptotics: multiseries of inverse functions,''
\emph{Journal of Symbolic Computation} \textbf{27} (1999), no.~6, 543--563.
\href{https://doi.org/10.1006/jsco.1999.0281}{doi:10.1006/jsco.1999.0281}.

\bibitem{vdDMM-1997}
L.~van den Dries, A.~Macintyre, and D.~Marker,
``Logarithmic-exponential power series,''
\emph{Journal of the London Mathematical Society} (2) \textbf{56} (1997),
no.~3, 417--434.

\bibitem{vdDMM-2001}
L.~van den Dries, A.~Macintyre, and D.~Marker,
``Logarithmic-exponential series,''
\emph{Annals of Pure and Applied Logic} \textbf{111} (2001), nos.~1--2,
61--113.

\bibitem{vanderhoeven-book}
J.~van der Hoeven,
\emph{Transseries and Real Differential Algebra},
Lecture Notes in Mathematics, vol.~1888, Springer, Berlin, 2006.
\href{https://doi.org/10.1007/3-540-35590-1}{doi:10.1007/3-540-35590-1}.

\end{thebibliography}


\part{Inversion: the apparatus for a rapidly growing function}

\chapter{Apparatus for inverting a rapidly growing function}\label{p0:sec:top}

\begin{sourcebox}
All twelve filed articles carry a private copy of the material of this
chapter.  The presentation below follows, in the main, the two sections
\emph{A universal coefficient calculus for exponential--power inversion} and
\emph{Lambert-normalized inverse transseries} of
\textsf{Swing\_Factorial\_Transseries\_Article}, which state the reversion in
the greatest generality of the twelve.  Four things have been changed.
(i)~The Lagrange fixed-point lemma is reproved: the swing article's proof
appeals to Lagrange--B\"urmann inversion in a setting where the required
invertibility hypothesis fails, and the gap is closed here by a universality
argument (\cref{p0:lem:LB-universal}).
(ii)~The exponential--power model is separated into a \emph{formal} and an
\emph{analytic} reading, and every asymptotic claim is given an explicit
remainder hypothesis; the sources routinely pass from a coefficient recurrence
to a Poincar\'e statement without one.
(iii)~The dominant core is axiomatized (\cref{p0:def:core}) so that the
factorial-type phase of \cref{p2:sec:top} --- which does \emph{not} belong to
the exponential--power model --- is covered by the same reversion theorem.
(iv)~The discrete-inverse material of
\textsf{Butcher\_Tree\_Transseries-2},
\textsf{double\_factorial\_transseries-2} and the three partition articles is
merged into one statement with a residue-class lattice.
\end{sourcebox}

\section{Standing conventions}\label{p0:sec:conventions}

Throughout the volume $R$ denotes a commutative $\Rat$-algebra, and
$R[[t]]$, $R[[t,u]]$ formal power series rings over it.  For a formal series
$f$ we write $\Coef{t^n}f$ for the coefficient of $t^n$.  The symbols
$\e,\ii,\dd$ are the upright constants and differential; $\Nat$, $\Int$,
$\Rat$, $\Real$, $\Cplx$ are the usual number systems, and
$\Nat=\{0,1,2,\dots\}$.  Asymptotic statements use $\bigO$ and
$\littleo$ in the Landau sense as the stated variable tends to
$+\infty$; an implied constant is always allowed to depend on every quantity
held fixed in the statement (in particular on the truncation index $N$) and on
nothing else.

The sources disagree on three points of bookkeeping.  We fix them once.

\begin{remark}[Conventions adopted, and where the sources differ]
\label{p0:rem:conventions}
\leavevmode
\begin{enumerate}
\item \emph{Ordinary versus exponential Bell polynomials.}  The partition
      articles and \textsf{Butcher\_Tree\_Transseries-2} take the
      \emph{exponential} partial Bell polynomials $\EB_{n,k}$ as primitive and
      convert; \textsf{Swing\_Factorial\_Transseries\_Article} and
      \textsf{double\_factorial\_transseries-2} take the \emph{ordinary} ones
      $\OB_{n,k}$ as primitive.  The volume defines both
      (\cref{p0:def:obell,p0:def:ebell}), records the conversion
      (\cref{p0:lem:bell-conversion}), and uses $\OB$ in every reversion
      formula and $\EB$ in every exponentiation or logarithm of a
      \emph{unit} series.  The reason is not taste: reversion multiplies a
      correction series whose coefficients carry no factorials, and $\OB$ is
      the transform that leaves them alone.
\item \emph{Sign of the logarithmic phase.}  We write the dominant phase as
      \[
        L=aX+b\log X ,
      \]
      with $b$ a \emph{signed} constant.  \textsf{Butcher\_Tree\_Transseries-2}
      writes $Z=X-p\log X$ and the partition articles write $\rho-2\log\rho=L$;
      both are the case $b=-p<0$.  Every formula below is stated in the signed
      convention, and the dictionary of
      \cref{p0:tab:dictionary} gives $b$ for each chapter.  One dividend of the
      signed convention is \cref{p0:cor:branch-rule}: the Lambert branch is
      $\Wlam_0$ when $b>0$ and $\Wlam_{-1}$ when $b<0$, with no case analysis
      beyond the sign.
\item \emph{Truncation index.}  A truncation ``of order $N$'' retains the terms
      with index $n=1,\dots,N$ \emph{inclusive}, together with the leading
      term.  \textsf{Swing\_Factorial\_Transseries\_Article} uses the exclusive
      convention for its forward table (``the truncation labelled $N$ retains
      $A_0,\dots,A_{N-1}$'') and the inclusive one for its inverse table
      (``the first $N$ correction terms are retained''); the two tables of that
      article are therefore not on a common scale.  The volume is inclusive
      throughout.
\end{enumerate}
\end{remark}

\section{The exponential--power model}\label{p0:sec:model}

\subsection{Formal and analytic readings}

\begin{definition}[The exponential--power model]\label{p0:def:model}
Fix constants $C>0$, $a>0$, $\alpha\in(0,1]$, $\delta>0$, $b\in\Real$, and a
formal series without constant term,
\[
 Q(t)=\sum_{j\ge1}q_jt^j\in t\,\Real[[t]] .
\]
The associated \emph{exponential--power model} is the expression
\begin{equation}\label{p0:eq:model}
 F(x)=C\exp\bigl(ax^\alpha\bigr)\,x^{b}\,\exp Q\bigl(x^{-\delta}\bigr),
\end{equation}
understood through its logarithm
\begin{equation}\label{p0:eq:model-log}
 \log F(x)=\log C+ax^{\alpha}+b\log x+Q\bigl(x^{-\delta}\bigr).
\end{equation}
We call $(C,a,\alpha,b,\delta,Q)$ the \emph{model data}.

A positive function $F\in C^1([x_0,\infty))$ is said to be \emph{of
exponential--power type with data $(C,a,\alpha,b,\delta,Q)$} when, for every
$M\in\Nat$, there is a constant $K_M$ with
\begin{equation}\label{p0:eq:model-analytic}
 \Biggl|\log F(x)-\log C-ax^{\alpha}-b\log x
        -\sum_{j=1}^{M}q_jx^{-j\delta}\Biggr|
 \le K_Mx^{-(M+1)\delta}
 \qquad(x\ge x_0),
\end{equation}
and \emph{differentiably} of that type when in addition
\begin{equation}\label{p0:eq:model-derivative}
 (\log F)'(x)=a\alpha x^{\alpha-1}+\frac bx+\bigO\bigl(x^{-1-\delta}\bigr)
 \qquad(x\to\infty).
\end{equation}
\end{definition}

\begin{remark}[Why the two readings are kept apart]\label{p0:rem:formal-analytic}
The series $Q$ is \emph{not} assumed convergent.  In
\cref{p1:sec:top,p2:sec:top,p4:sec:top} it is divergent of Gevrey type, which
is what makes \cref{p0:thm:optimal-truncation} bite there.  It is \emph{not}
divergent in \cref{p3:sec:top}: for the partition numbers
$Q(t)=\log(1-t)$ has radius of convergence $1$, every algebraic expansion of
that chapter converges, and \cref{p0:thm:optimal-truncation} correspondingly
has no content there --- the error floor of \cref{p3:sec:top} is set by the
exponentially small Rademacher remainder and not by a least term.  This is
precisely why the two readings are kept apart: convergence of $Q$ is a
property of the subject, not of the frame.  Every purely algebraic assertion
below --- coefficient recurrences, closed coefficient formulae, degree bounds
--- is a statement about \eqref{p0:eq:model-log} as a formal object and needs
no analytic hypothesis at all.  Every assertion with an error term uses
\eqref{p0:eq:model-analytic}, and every assertion that inverts uses
\eqref{p0:eq:model-derivative} as well: an asymptotic expansion may not be
differentiated without a hypothesis, and the inversion of $F$ is controlled by
$(\log F)'$.  In each chapter \eqref{p0:eq:model-derivative} is verified
directly from a digamma, Laplace or Rademacher representation, never inferred
from \eqref{p0:eq:model-analytic}.
\end{remark}

\subsection{Reduction to a linear phase}

\begin{lemma}[Monomial $\alpha$-reduction]\label{p0:lem:alpha-reduction}
Let $F$ be of exponential--power type with data
$(C,a,\alpha,b,\delta,Q)$ on $[x_0,\infty)$, and put
\begin{equation}\label{p0:eq:alpha-substitution}
 \xi=x^{\alpha},\qquad x=\xi^{1/\alpha},\qquad
 G(\xi):=F\bigl(\xi^{1/\alpha}\bigr).
\end{equation}
Then $G$ is of exponential--power type on $[x_0^\alpha,\infty)$ with data
\begin{equation}\label{p0:eq:alpha-reduced-data}
 \Bigl(C,\;a,\;1,\;\frac b\alpha,\;\frac\delta\alpha,\;Q\Bigr),
\end{equation}
that is,
\begin{equation}\label{p0:eq:alpha-reduced-log}
 \log G(\xi)=\log C+a\xi+\frac b\alpha\log\xi
             +Q\bigl(\xi^{-\delta/\alpha}\bigr).
\end{equation}
The coefficients $q_j$ are unchanged.  If moreover $F$ is differentiably of
its type then $G$ is differentiably of its type.  Finally $x\mapsto x^\alpha$
is an increasing bijection of $[x_0,\infty)$ onto $[x_0^\alpha,\infty)$, so
$F$ is strictly increasing if and only if $G$ is, and then
\begin{equation}\label{p0:eq:alpha-inverse-transport}
 F^{-1}(y)=\bigl(G^{-1}(y)\bigr)^{1/\alpha}.
\end{equation}
\end{lemma}

\begin{proof}
Substituting $x=\xi^{1/\alpha}$ into \eqref{p0:eq:model-log} and using
$x^{\alpha}=\xi$, $\log x=\alpha^{-1}\log\xi$ and
$x^{-j\delta}=\xi^{-j\delta/\alpha}$ gives
\eqref{p0:eq:alpha-reduced-log} term by term; the truncated inequality
\eqref{p0:eq:model-analytic} transforms into the corresponding inequality for
$G$ with the same constants $K_M$, because
$x^{-(M+1)\delta}=\xi^{-(M+1)\delta/\alpha}$.  For the derivative,
$\dd x/\dd\xi=\alpha^{-1}\xi^{1/\alpha-1}$, so
\[
 (\log G)'(\xi)=(\log F)'\bigl(\xi^{1/\alpha}\bigr)
 \cdot\frac{\xi^{1/\alpha-1}}{\alpha}
 =\Bigl(a\alpha\xi^{1-1/\alpha}+\frac{b}{\xi^{1/\alpha}}
        +\bigO\bigl(\xi^{-(1+\delta)/\alpha}\bigr)\Bigr)
   \frac{\xi^{1/\alpha-1}}{\alpha}
 =a+\frac{b/\alpha}{\xi}+\bigO\bigl(\xi^{-1-\delta/\alpha}\bigr),
\]
which is \eqref{p0:eq:model-derivative} for $G$ with $\alpha$ replaced by $1$.
The last two assertions are immediate from strict monotonicity of
$x\mapsto x^\alpha$ on $(0,\infty)$.
\end{proof}

\begin{remark}[Normalization used from here on]\label{p0:rem:alpha-one}
By \cref{p0:lem:alpha-reduction} every statement about the model
\eqref{p0:eq:model} reduces to the case $\alpha=1$, and we assume $\alpha=1$
from \cref{p0:sec:core} onwards.  We assume in addition $\delta=1$; this is the
value of $\delta/\alpha$ in every chapter of the volume
(\cref{p0:tab:dictionary}).  \Cref{p0:prop:two-grid} states the general-$\delta$
reversion, on the two-generator monomial grid it requires.
\end{remark}

\begin{remark}[The affine shift is taken \emph{before} the reduction]
\label{p0:rem:affine-shift}
Chapters~\ref{p2:sec:top} and \ref{p3:sec:top} apply the model not to the
combinatorial index $n$ but to a shifted variable $x=n-x_\ast$ ($x_\ast=1/24$
for the partition numbers, $x_\ast=-1$ for the double factorial).  This is not
a loss of generality but a genuine choice of normal form, and the order of
operations matters.  If $F$ has data $(C,a,\alpha,b,\delta,Q)$ in $x$ and one
re-expands in $n=x+x_\ast$, then for $\alpha<1$
\[
 a(n-x_\ast)^{\alpha}
 =an^{\alpha}+a\sum_{k\ge1}\binom{\alpha}{k}(-x_\ast)^kn^{\alpha-k},
\]
so the phase acquires an entire new tower of powers $n^{\alpha-k}$: the
re-expanded object no longer has model data with a single exponent grid.  The
partition articles say this explicitly --- their infinite algebraic
coefficient arrays in the variable $n^{-1/2}$ ``arise entirely from
re-expanding the shift $N=n-1/24$''.  The volume therefore normalizes the
shift first, inverts in $x$, and returns to the index by
$n=x+x_\ast$, an exact identity that costs nothing.
\end{remark}

\begin{remark}[Which chapters the model covers, and which it does not]
\label{p0:rem:model-scope}
Chapters~\ref{p1:sec:top}, \ref{p3:sec:top} and \ref{p4:sec:top} are governed by
\cref{p0:def:model}: the rooted-tree counts and the swing factorial with
$\alpha=1$, the partition numbers with $\alpha=1/2$
(\cref{p0:tab:dictionary}).  \Cref{p2:sec:top} is not.
\end{remark}

\begin{repairbox}
\textbf{Claim repaired.}  It is natural, and it is asserted in the merge brief
for this volume, that the case $\alpha=1$ of \cref{p0:def:model} covers the
double factorial alongside the swing factorial and the rooted-tree counts.
It does not.  For the swing factorial
$n!/\lfloor n/2\rfloor!^{2}$ the growth is $2^{n}n^{\pm1/2}$ and the model
applies with $a=\log 2$.  For the double factorial,
\[
 \log\bigl((x+1)!!\bigr)=\frac{x+1}{2}\bigl(\log(x+1)-1\bigr)+\bigO(\log x),
\]
so the dominant phase is $\tfrac12s\log s$ with $s=x+1$, not $ax$: no choice of
the data $(C,a,\alpha,b,\delta,Q)$ with $\alpha\le1$ reproduces it, since
$s\log s/s^{\alpha}\to\infty$ for every $\alpha\le1$.  The three double
factorial articles are consistent with this; they never claim membership of the
exponential--power family, and they invert the phase
$\tfrac{s}{2}(\log s-1)$ directly.
\textbf{Correction.}  The chapter states the reversion for an axiomatized
\emph{dominant core} (\cref{p0:def:core,p0:thm:core-reversion}) of which both
the linear-plus-logarithmic phase $aX+b\log X$ and the factorial phase
$X(\kappa\log X+d)$ are instances, both exactly invertible by Lambert $\Wlam$
(\cref{p0:thm:lambert-core,p0:prop:factorial-core}).  The theorem named in the
Part~0 interface, \cref{p0:thm:lambert-centered}, is stated for the
exponential--power model as specified, and \cref{p0:thm:core-reversion} is the
version \cref{p2:sec:top} cites.
\end{repairbox}

\begin{definition}[Admissible dominant core]\label{p0:def:core}
Let $X_\ast\ge0$.  A function $\Phi_0\in C^{2}((X_\ast,\infty))$ is an
\emph{admissible core} when
\begin{enumerate}
\item $\Phi_0'(X)>0$ for $X>X_\ast$ and $\Phi_0(X)\to+\infty$ as
      $X\to+\infty$;
\item $\Phi_0'$ is eventually monotone and $\Phi_0'(X)\to\Lambda_\infty\in(0,+\infty]$;
\item $X\Phi_0'(X)\to+\infty$.
\end{enumerate}
Its \emph{core slope} is $\Phi_0'$.  Given a target $L$ in the range of
$\Phi_0$, the \emph{core solution} $X=X(L)$ is the unique large solution of
$\Phi_0(X)=L$.
\end{definition}

The two cores used in this volume are
\begin{equation}\label{p0:eq:two-cores}
 \Phi_0(X)=aX+b\log X
 \quad(\text{\cref{p1:sec:top},\ \ref{p3:sec:top},\ \ref{p4:sec:top}}),
 \qquad
 \Phi_0(X)=X\,(\kappa\log X+d)
 \quad(\text{\cref{p2:sec:top}}),
\end{equation}
with $a,\kappa>0$.  Both are admissible on a suitable half-line, and both are
inverted exactly by $\Wlam$ in \cref{p0:sec:core}.

\section{The coefficient calculus}\label{p0:sec:calculus}

\subsection{Two families of partial Bell polynomials}

\begin{definition}[Ordinary partial Bell polynomials]\label{p0:def:obell}
For a sequence $\gamma=(\gamma_1,\gamma_2,\dots)$ in $R$ and $n,k\in\Nat$ put
\begin{equation}\label{p0:eq:obell-def}
 \OB_{n,k}(\gamma):=\Coef{t^n}\Bigl(\sum_{r\ge1}\gamma_rt^r\Bigr)^{k}.
\end{equation}
Explicitly,
\begin{equation}\label{p0:eq:obell-multinomial}
 \OB_{n,k}(\gamma)
 =\sum_{\substack{m_1,m_2,\dots\ge0\\ \sum_rm_r=k,\ \sum_rrm_r=n}}
   \frac{k!}{\prod_{r\ge1}m_r!}\prod_{r\ge1}\gamma_r^{m_r}.
\end{equation}
In particular $\OB_{0,0}=1$, $\OB_{n,0}=0$ for $n\ge1$, $\OB_{n,k}=0$ for
$k>n$, $\OB_{n,n}=\gamma_1^{\,n}$, and $\OB_{n,k}$ depends only on
$\gamma_1,\dots,\gamma_{n-k+1}$.
\end{definition}

\begin{definition}[Exponential partial Bell polynomials]\label{p0:def:ebell}
For $x=(x_1,x_2,\dots)$ in $R$ and $n\ge k\ge0$ put
\begin{equation}\label{p0:eq:ebell-def}
 \EB_{n,k}(x_1,x_2,\dots)
 :=n!\sum_{\substack{m_1,m_2,\dots\ge0\\ \sum_jm_j=k,\ \sum_jjm_j=n}}
   \prod_{j\ge1}\frac1{m_j!}\Bigl(\frac{x_j}{j!}\Bigr)^{m_j},
\end{equation}
and let $\EB_n:=\sum_{k=0}^{n}\EB_{n,k}$ be the complete polynomial.  These are
characterized by the generating identities
\begin{align}
 \frac1{k!}\Bigl(\sum_{j\ge1}x_j\frac{t^j}{j!}\Bigr)^{k}
 &=\sum_{n\ge k}\EB_{n,k}(x_1,x_2,\dots)\frac{t^n}{n!},
 \label{p0:eq:ebell-egf-partial}\\
 \exp\Bigl(\sum_{j\ge1}x_j\frac{t^j}{j!}\Bigr)
 &=\sum_{n\ge0}\EB_n(x_1,\dots,x_n)\frac{t^n}{n!}.
 \label{p0:eq:ebell-egf-complete}
\end{align}
\end{definition}

\begin{lemma}[Conversion]\label{p0:lem:bell-conversion}
For all $n\ge k\ge0$,
\begin{equation}\label{p0:eq:bell-conversion}
 \OB_{n,k}(\gamma_1,\gamma_2,\dots)
 =\frac{k!}{n!}\,\EB_{n,k}(1!\,\gamma_1,\,2!\,\gamma_2,\dots).
\end{equation}
\end{lemma}

\begin{proof}
Set $x_j=j!\,\gamma_j$ in \eqref{p0:eq:ebell-egf-partial}.  The left-hand side
becomes $(k!)^{-1}\bigl(\sum_j\gamma_jt^j\bigr)^k$, whose coefficient of $t^n$
is $\OB_{n,k}(\gamma)/k!$; the right-hand side has coefficient
$\EB_{n,k}(1!\gamma_1,\dots)/n!$.
\end{proof}

\subsection{Powers and logarithms of a unit series}

\begin{lemma}[Powers and logarithms]\label{p0:lem:power-log}
Let $w=\sum_{r\ge1}\gamma_rt^r\in tR[[t]]$ and $n\ge1$.
\begin{enumerate}
\item For every $\beta$ in $R$ (or in any commutative $\Rat$-algebra containing
      $R$),
      \begin{equation}\label{p0:eq:power-general}
       \Coef{t^n}(1+w)^{\beta}
       =\sum_{k=1}^{n}\binom{\beta}{k}\OB_{n,k}(\gamma),
       \qquad
       \binom\beta k=\frac{\beta(\beta-1)\cdots(\beta-k+1)}{k!}.
      \end{equation}
\item For an integer $j\ge1$,
      \begin{equation}\label{p0:eq:power-negative}
       \Coef{t^n}(1+w)^{-j}
       =\sum_{k=0}^{n}(-1)^{k}\binom{j+k-1}{k}\OB_{n,k}(\gamma).
      \end{equation}
\item
      \begin{equation}\label{p0:eq:log-obell}
       \Coef{t^n}\log(1+w)
       =\sum_{k=1}^{n}\frac{(-1)^{k+1}}{k}\OB_{n,k}(\gamma).
      \end{equation}
\end{enumerate}
Equivalently, in exponential form, for $U(t)=1+\sum_{j\ge1}u_jt^j$,
\begin{align}
 \Coef{t^N}U(t)^{\beta}
 &=\frac1{N!}\sum_{k=0}^{N}\beta^{\underline k}\,
   \EB_{N,k}(1!u_1,2!u_2,\dots),
 \qquad \beta^{\underline k}=\beta(\beta-1)\cdots(\beta-k+1),
 \label{p0:eq:unit-power-ebell}\\
 \Coef{t^N}\log U(t)
 &=\frac1{N!}\sum_{k=1}^{N}(-1)^{k-1}(k-1)!\,
   \EB_{N,k}(1!u_1,2!u_2,\dots).
 \label{p0:eq:log-ebell}
\end{align}
\end{lemma}

\begin{proof}
All four identities are the substitution of $w$ into a formal series in one
variable.  Since $w\in tR[[t]]$, the family $(w^k)_{k\ge0}$ is summable in the
$t$-adic topology and $\Coef{t^n}w^k=0$ for $k>n$, so every sum displayed is
finite and the substitution is legitimate.  Now
$(1+w)^\beta=\sum_{k\ge0}\binom\beta kw^k$ gives
\eqref{p0:eq:power-general} on applying \cref{p0:def:obell};
$\binom{-j}{k}=(-1)^k\binom{j+k-1}{k}$ gives
\eqref{p0:eq:power-negative}; and
$\log(1+w)=\sum_{k\ge1}(-1)^{k+1}w^k/k$ gives \eqref{p0:eq:log-obell}.
For the exponential forms write $U=1+w$ with $\gamma_j=u_j$ and substitute
\eqref{p0:eq:bell-conversion} into \eqref{p0:eq:power-general} and
\eqref{p0:eq:log-obell}, using
$\binom\beta k=\beta^{\underline k}/k!$ and $(-1)^{k+1}/k=(-1)^{k-1}(k-1)!/k!$.
\end{proof}

\begin{corollary}[Exponential and logarithmic jets]\label{p0:cor:exp-log-jets}
Let $D(t)=\sum_{j\ge1}d_jt^j\in tR[[t]]$ and $A(t)=\exp D(t)=\sum_{n\ge0}A_nt^n$.
Then $A_0=1$,
\begin{equation}\label{p0:eq:exp-bell}
 A_n=\frac1{n!}\EB_n(1!d_1,2!d_2,\dots,n!d_n)
    =\sum_{\substack{k_1,k_2,\dots\ge0\\ \sum_jjk_j=n}}\prod_{j\ge1}\frac{d_j^{k_j}}{k_j!},
\end{equation}
and the coefficients obey the triangular recurrence
\begin{equation}\label{p0:eq:exp-recurrence}
 nA_n=\sum_{j=1}^{n}j\,d_j\,A_{n-j}\qquad(n\ge1).
\end{equation}
Conversely, if $A=1+\sum_{n\ge1}A_nt^n$ is a unit series then
$d_n=\Coef{t^n}\log A$ is given by \eqref{p0:eq:log-ebell}.
\end{corollary}

\begin{proof}
The first equality is \eqref{p0:eq:ebell-egf-complete} with $x_j=j!d_j$; the
second is the expansion of $\exp$ as a product over $j$.  For
\eqref{p0:eq:exp-recurrence} differentiate $A=\exp D$ to get $A'=D'A$ and
compare coefficients of $t^{n-1}$.  The converse is
\eqref{p0:eq:log-ebell}.
\end{proof}

\subsection{Lagrange--B\"urmann inversion and the fixed-point form}

\begin{theorem}[Lagrange--B\"urmann]\label{p0:thm:LB}
Let $\varphi\in R[[w]]$ with $\varphi(0)$ invertible in $R$.  There is a unique
$W\in tR[[t]]$ with
\begin{equation}\label{p0:eq:LB-equation}
 W=t\,\varphi(W),
\end{equation}
and for every $G\in R[[w]]$ and every $n\ge1$,
\begin{equation}\label{p0:eq:LB-general}
 \Coef{t^n}G\bigl(W(t)\bigr)
 =\frac1n\Coef{w^{n-1}}\Bigl(G'(w)\,\varphi(w)^n\Bigr).
\end{equation}
In particular
\begin{equation}\label{p0:eq:LB-basic}
 \Coef{t^n}W(t)=\frac1n\Coef{w^{n-1}}\varphi(w)^n .
\end{equation}
\end{theorem}

\begin{proof}
Existence and uniqueness: \eqref{p0:eq:LB-equation} determines
$\Coef{t^n}W$ from $\Coef{t^{<n}}W$, because $t\varphi(W)$ modulo $t^{n+1}$
depends only on $W$ modulo $t^{n}$.

For \eqref{p0:eq:LB-general} work in the field of formal Laurent series
$R((w))$ and let $\Res$ denote the coefficient of $w^{-1}$.  Two facts are
needed.

\emph{(a) $\Res f'=0$ for every $f\in R((w))$}, since $w^{-1}$ has no
antiderivative among the monomials: $(w^{k})'=kw^{k-1}$ is a multiple of
$w^{-1}$ only for $k=0$, when it vanishes.

\emph{(b) Change of variables.}  Put $\psi(w):=w/\varphi(w)$, a series in
$wR[[w]]$ with $\psi'(0)=\varphi(0)^{-1}$ invertible; then $W$ is the
compositional inverse of $\psi$, so $\psi(W(t))=t$.  For $f\in R((w))$,
\begin{equation}\label{p0:eq:res-change}
 \Res_w\bigl(f(w)\bigr)=\Res_t\bigl(f(W(t))\,W'(t)\bigr).
\end{equation}
Indeed both sides are $R$-linear and continuous, so it suffices to check
$f=w^k$.  For $k\ne-1$, $W^kW'=\bigl(W^{k+1}/(k+1)\bigr)'$ and (a) gives
$0=\Res_w w^k$.  For $k=-1$, write $W=t\vartheta(t)$ with $\vartheta$ a unit;
then $W'/W=t^{-1}+\vartheta'/\vartheta$ and the second summand lies in
$R[[t]]$, so $\Res_t(W'/W)=1=\Res_ww^{-1}$.

Now apply \eqref{p0:eq:res-change} to
$f(w):=G(w)\psi(w)^{-n-1}\psi'(w)$, noting that
$f(W(t))W'(t)=G(W(t))\,t^{-n-1}\,\tfrac{\dd}{\dd t}\psi(W(t))
 =G(W(t))\,t^{-n-1}$ because $\psi(W(t))=t$:
\[
 \Coef{t^n}G(W(t))=\Res_t\bigl(G(W(t))\,t^{-n-1}\bigr)
 =\Res_w\bigl(G(w)\,\psi(w)^{-n-1}\psi'(w)\bigr).
\]
Since $\psi^{-n-1}\psi'=-\tfrac1n\bigl(\psi^{-n}\bigr)'$, fact (a) and the
product rule give
\[
 \Res_w\bigl(G\,\psi^{-n-1}\psi'\bigr)
 =-\frac1n\Res_w\bigl(G\cdot(\psi^{-n})'\bigr)
 =\frac1n\Res_w\bigl(G'\,\psi^{-n}\bigr)
 =\frac1n\Res_w\bigl(G'(w)\,\varphi(w)^{n}w^{-n}\bigr),
\]
which is \eqref{p0:eq:LB-general}.  Taking $G(w)=w$ gives
\eqref{p0:eq:LB-basic}.
\end{proof}

The reversions of this volume all take the form $U=\Psi(t,U)$ in which the
r\^ole of $\varphi(0)$ is played by $\Psi(t,0)$, an element of $tR[[t]]$, which
is \emph{never} invertible.  \Cref{p0:thm:LB} therefore does not apply
directly, and the following universality lemma is what makes the passage
legitimate.

\begin{lemma}[Lagrange's formula without an invertibility hypothesis]
\label{p0:lem:LB-universal}
Let $S$ be any commutative $\Rat$-algebra and $\varphi\in S[[w]]$
\emph{arbitrary} (in particular $\varphi(0)$ need not be invertible).  Let
$W\in zS[[z]]$ be the unique solution of $W=z\varphi(W)$.  Then for every
$m\ge1$,
\begin{equation}\label{p0:eq:LB-universal}
 \Coef{z^m}W=\frac1m\Coef{w^{m-1}}\varphi(w)^m .
\end{equation}
\end{lemma}

\begin{proof}
Existence and uniqueness of $W$ hold as in \cref{p0:thm:LB} (the recursion is
$z$-adic and uses no invertibility).  Fix $m$.  Iterating $W=z\varphi(W)$
shows that $\Coef{z^m}W$ is a universal polynomial with integer coefficients in
$\varphi_0,\varphi_1,\dots,\varphi_{m-1}$, where $\varphi=\sum_i\varphi_iw^i$;
the right-hand side of \eqref{p0:eq:LB-universal} is a universal polynomial
with rational coefficients in the same variables.  Both sides are therefore
elements $P_m,\;\widetilde P_m$ of the polynomial ring
$\mathcal P:=\Rat[\varphi_0,\dots,\varphi_{m-1}]$, and they are obtained from
$\varphi$ by the unique $\Rat$-algebra homomorphism
$\mathcal P\to S$ sending $\varphi_i\mapsto\varphi_i$.  It thus suffices to
prove $P_m=\widetilde P_m$ in $\mathcal P$.  The localization map
$\mathcal P\to\mathcal P[\varphi_0^{-1}]$ is injective ($\mathcal P$ is an
integral domain), and in $\mathcal P[\varphi_0^{-1}]$ the element $\varphi_0$
is invertible, so \cref{p0:thm:LB} applies and gives
$P_m=\widetilde P_m$ there.  Injectivity transports the equality back to
$\mathcal P$, and the homomorphism transports it to $S$.
\end{proof}

\begin{lemma}[Lagrange fixed-point formula]\label{p0:lem:lagrange-fp}
Let $\Psi\in tR[[t,u]]$, that is, a formal series in $t$ and $u$ each of whose
monomials has $t$-degree at least one.  Then the equation
\begin{equation}\label{p0:eq:fp-equation}
 U(t)=\Psi\bigl(t,U(t)\bigr)
\end{equation}
has a unique solution $U\in tR[[t]]$, and for every $n\ge1$
\begin{equation}\label{p0:eq:fp-coefficient}
 \Coef{t^n}U=\sum_{m=1}^{n}\frac1m\Coef{t^nu^{m-1}}\Psi(t,u)^m .
\end{equation}
\end{lemma}

\begin{proof}
Uniqueness and existence: $\Psi(t,U)$ modulo $t^{n+1}$ depends only on $U$
modulo $t^{n}$ because $\Psi$ is divisible by $t$; so
\eqref{p0:eq:fp-equation} determines the coefficients of $U$ recursively, and
the resulting series solves the equation by $t$-adic continuity.

For \eqref{p0:eq:fp-coefficient} apply \cref{p0:lem:LB-universal} over the
ring $S:=R[[t]]$ with $\varphi(w):=\Psi(t,w)\in S[[w]]$ and marker $z$: the
unique $\widetilde U\in zS[[z]]$ with $\widetilde U=z\varphi(\widetilde U)$
satisfies
\[
 \Coef{z^m}\widetilde U=\frac1m\Coef{w^{m-1}}\Psi(t,w)^m\in S .
\]
Since $\Psi\in tR[[t,u]]$, we have $\Psi^m\in t^mR[[t,u]]$, whence
$\Coef{t^n}\Coef{w^{m-1}}\Psi(t,w)^m=0$ for $m>n$.  Consequently the family
$\bigl(\Coef{z^m}\widetilde U\bigr)_{m\ge1}$ is summable in the $t$-adic
topology of $S$, the specialization $z=1$ is defined coefficientwise in $t$,
and the resulting series $U:=\sum_{m\ge1}\Coef{z^m}\widetilde U$ satisfies
$U=\Psi(t,U)$ by the same summability.  Uniqueness identifies it with the
solution of \eqref{p0:eq:fp-equation}, and extracting $\Coef{t^n}$ gives
\eqref{p0:eq:fp-coefficient} with the sum truncated at $m=n$.
\end{proof}

\begin{repairbox}
\textbf{Source claim.}  \textsf{Swing\_Factorial\_Transseries\_Article},
Lemma \emph{Lagrange fixed-point formula}, proves
\eqref{p0:eq:fp-coefficient} by writing ``Ordinary Lagrange--B\"urmann
inversion in $z$ gives $U=\sum_{m\ge1}z^mm^{-1}[u^{m-1}]\Psi(t,u)^m$''.  The
same step, with the same silence, appears in
\textsf{Partition\_Number\_Transseries\_and\_Inverse} (``Introduce
$u=z\Phi_\rho(u)$.  Lagrange--B\"urmann gives \dots''), in
\textsf{Partition\_Numbers\_Transseries\_and\_Inverse}, and in
\textsf{double\_factorial\_transseries-2}.
\textbf{Defect.}  The classical Lagrange--B\"urmann theorem in the form
\eqref{p0:eq:LB-basic} requires $\varphi(0)$ to be invertible in the
coefficient ring.  Here $\varphi(0)=\Psi(t,0)\in tR[[t]]$, a non-unit, and the
residue proof of \cref{p0:thm:LB} breaks at the change of variables: $\psi=w/\varphi$
is not a well-defined element of $wR[[t]][[w]]$.
\textbf{Repair.}  \Cref{p0:lem:LB-universal}: the identity is a polynomial
identity in the coefficients of $\varphi$ over $\Rat$, hence holds in the
polynomial ring by injectivity of localization at $\varphi_0$, hence in every
$\Rat$-algebra.  \textsf{double\_factorial\_transseries-2} comes closest to
noticing the issue --- ``Although $A$ is a divergent Gevrey series, it is a
valid formal power series with nonzero linear coefficient'' --- but that
remark addresses convergence, not invertibility, and its own $G_\Lambda$ does
have an invertible linear coefficient, so the gap there is only in the
inherited lemma.
\end{repairbox}


\begin{lemma}[Fixed-point formula with a quadratic self-interaction]
\label{p0:lem:lagrange-fp-quadratic}
Let $\Psi=h+f$ with $h\in u^2R[[u]]$ and $f\in tR[[t,u]]$.  Then
$U=\Psi(t,U)$ has a unique solution $U\in tR[[t]]$, and for every $n\ge1$
\begin{equation}\label{p0:eq:fp-quadratic}
 \Coef{t^n}U=\sum_{m=1}^{2n-1}\frac1m\Coef{t^nu^{m-1}}\Psi(t,u)^m .
\end{equation}
\end{lemma}

\begin{proof}
\emph{Existence and uniqueness.}  Write $U=U_{<n}+e_nt^n+\cdots$.  Then
$h(U)=h(U_{<n})+h'(U_{<n})\,e_nt^n+\cdots$, and $h'(U_{<n})\in tR[[t]]$ because
$h'\in uR[[u]]$ and $U_{<n}\in tR[[t]]$; so $e_n$ does not occur in
$\Coef{t^n}h(U)$.  It does not occur in $\Coef{t^n}f(t,U)$ either, since $f$ is
divisible by $t$.  Hence the equation determines $e_n$, as
$\Coef{t^n}\Psi(t,U_{<n})$, from $U_{<n}$; the resulting series solves the
equation by $t$-adic continuity.

\emph{The closed formula.}  Introduce a marker $z$ and let
$\widetilde U\in zS[[z]]$, $S:=R[[t]]$, solve
$\widetilde U=z\Psi(t,\widetilde U)$.  \Cref{p0:lem:LB-universal} gives
$\Coef{z^m}\widetilde U=m^{-1}\Coef{u^{m-1}}\Psi(t,u)^m$.  We claim
\begin{equation}\label{p0:eq:quadratic-vanishing}
 \Coef{t^nu^{m-1}}\Psi(t,u)^m=0\qquad\text{whenever } m>2n-1 .
\end{equation}
Expand $\Psi^m=(h+f)^m$ into $2^m$ products; a product using $h$ in $\sigma$ of
the $m$ slots and $f$ in the remaining $m-\sigma$ has $u$-degree at least
$2\sigma$ and $t$-degree at least $m-\sigma$.  For the $u$-degree to equal
$m-1$ we need $2\sigma\le m-1$, hence the $t$-degree is at least
$m-\tfrac{m-1}2=\tfrac{m+1}2$; so $\Coef{t^n}$ of that product vanishes unless
$n\ge(m+1)/2$, which is \eqref{p0:eq:quadratic-vanishing}.  Consequently the
family $(\Coef{z^m}\widetilde U)_{m\ge1}$ is $t$-adically summable, the
specialization $z=1$ is defined coefficientwise, and, again by summability, the
resulting series solves $U=\Psi(t,U)$.  Uniqueness identifies it with $U$, and
extracting $\Coef{t^n}$ gives \eqref{p0:eq:fp-quadratic}.
\end{proof}

\Cref{p0:lem:lagrange-fp} is the case $h=0$ of
\cref{p0:lem:lagrange-fp-quadratic}, in which the argument above sharpens the
range of $m$ from $2n-1$ to $n$.

\subsection{Perturbed inversion around an exactly invertible core}
\label{p0:sec:perturbed}

The reversions carried out in Chapters~1--4 all reduce to one step: a core
$F$ that can be inverted exactly is perturbed by a correction $E$, and the
inverse of $F+\varepsilon E$ is required to all orders in $\varepsilon$.  The
answer is a single differential formula, stated here once.  It is the tool
behind \cref{p3:thm:additive-LB} in the ordinate of the partition chapter and
behind \cref{p4:lem:phase-lb} in the phase coordinate of the swing chapter,
and those two are the same identity in different letters.

\begin{theorem}[Perturbed inversion]\label{p0:thm:perturbed-inversion}
Let $F$ and $E$ be analytic in a neighbourhood of $\mu_{0}\in\Cplx$, let
$F'(\mu_{0})\ne0$, and put $y:=F(\mu_{0})$.  Write $\mu_{0}(\cdot)$ for the
local inverse of $F$, so that $\mu_{0}(y)=\mu_{0}$ and
$\mu_{0}'(y)=1/F'(\mu_{0}(y))$, and let
\begin{equation}\label{p0:eq:reversion-operator}
 \mathcal{D}:=\frac1{F'(\mu)}\frac{\dd}{\dd\mu},
\end{equation}
which is the derivation $\dd/\dd y$ read in the coordinate $y=F(\mu)$.  Then
there are a neighbourhood $\mathcal{U}$ of $\mu_{0}$ and an
$\varepsilon_{0}>0$ such that for $|\varepsilon|<\varepsilon_{0}$ the equation
\begin{equation}\label{p0:eq:perturbed-equation}
 F(\mu)+\varepsilon E(\mu)=y
\end{equation}
has exactly one root $\mu(y;\varepsilon)$ in $\mathcal{U}$; it is analytic in
$\varepsilon$ there, and
\begin{equation}\label{p0:eq:perturbed-inversion}
 \mu(y;\varepsilon)=\mu_{0}(y)
 +\sum_{M\ge1}\frac{(-\varepsilon)^{M}}{M!}\,
 \mathcal{D}^{M-1}\!\left(\frac{E(\mu_{0}(y))^{M}}{F'(\mu_{0}(y))}\right),
\end{equation}
the series converging for $|\varepsilon|<\varepsilon_{0}$, with $M$th
coefficient
\begin{equation}\label{p0:eq:perturbed-residue}
 \frac{(-1)^{M}}{M!}\,\mathcal{D}^{M-1}\!
 \left(\frac{E(\mu_{0})^{M}}{F'(\mu_{0})}\right)
 =\frac{(-1)^{M}}{M}\operatorname*{Res}_{\mu=\mu_{0}}
 \frac{E(\mu)^{M}}{\bigl(F(\mu)-y\bigr)^{M}} .
\end{equation}
If moreover $E=\sum_{j}\varepsilon_{j}E_{j}$ with finitely many $E_{j}$
analytic near $\mu_{0}$, then for every multi-index $\boldsymbol m\ne\boldsymbol0$
with $M:=|\boldsymbol m|$ the coefficient of
$\boldsymbol\varepsilon^{\boldsymbol m}$ in $\mu-\mu_{0}$ is
\begin{equation}\label{p0:eq:perturbed-multi}
 V_{\boldsymbol m}
 =\frac{(-1)^{M}}{\boldsymbol m!}\,
  \mathcal{D}^{M-1}\!\left(\frac{E^{\boldsymbol m}(\mu_{0})}{F'(\mu_{0})}\right)
 =\frac{(-1)^{M}(M-1)!}{\boldsymbol m!}\,
  \operatorname*{Res}_{\mu=\mu_{0}}
  \frac{E^{\boldsymbol m}(\mu)}{\bigl(F(\mu)-F(\mu_{0})\bigr)^{M}},
\end{equation}
where $E^{\boldsymbol m}:=\prod_{j}E_{j}^{m_{j}}$ and
$\boldsymbol m!:=\prod_{j}m_{j}!$.
\end{theorem}

\begin{proof}
Let $\gamma$ be a small positively oriented circle about $\mu_{0}$, contained
with its interior $\mathcal{U}$ in the common domain of analyticity and
enclosing no zero of $F(\cdot)-y$ other than $\mu_{0}$; such a $\gamma$ exists
because $F'(\mu_{0})\ne0$ makes $\mu_{0}$ a simple zero, hence isolated.  Put
$m_{\gamma}:=\min_{\gamma}|F-y|>0$ and $M_{\gamma}:=\max_{\gamma}|E|$, and let
$\varepsilon_{0}:=m_{\gamma}/(M_{\gamma}+1)$.  For
$|\varepsilon|<\varepsilon_{0}$ we have $|\varepsilon E|<|F-y|$ on $\gamma$,
so Rouch\'e gives exactly one zero $\mu(y;\varepsilon)$ of
$F-y+\varepsilon E$ inside $\gamma$, and it is simple.  The argument
principle, applied to that function, expresses it as
\begin{equation}\label{p0:eq:argument-principle}
 \mu(y;\varepsilon)=\frac1{2\pi\ii}\int_{\gamma}\zeta\,
 \dd_{\zeta}\log\bigl(F(\zeta)-y+\varepsilon E(\zeta)\bigr),
\end{equation}
which is analytic in $\varepsilon$ on $|\varepsilon|<\varepsilon_{0}$ because
the integrand is and $\gamma$ is compact.

Subtract from \eqref{p0:eq:argument-principle} the same expression at
$\varepsilon=0$, whose value is $\mu_{0}$, and expand
\[
 \log\Bigl(1+\varepsilon\frac{E}{F-y}\Bigr)
 =\sum_{M\ge1}\frac{(-1)^{M-1}}{M}\,\varepsilon^{M}\,
 \frac{E^{M}}{(F-y)^{M}},
\]
uniformly convergent on $\gamma$ since $|\varepsilon E/(F-y)|<1$ there.
Integrating by parts once turns $\zeta\,\dd_{\zeta}(\cdot)$ into
$-(\cdot)\,\dd\zeta$ (the boundary term vanishes on a closed contour), so the
coefficient of $\varepsilon^{M}$ is
\[
 \frac{(-1)^{M}}{M}\operatorname*{Res}_{\zeta=\mu_{0}}
 \frac{E(\zeta)^{M}}{\bigl(F(\zeta)-F(\mu_{0})\bigr)^{M}},
\]
which is the right-hand side of \eqref{p0:eq:perturbed-residue}.

It remains to identify that residue with the differential expression.  Change
coordinate by $w=F(\zeta)$, $\zeta=\mu_{0}(w)$, $\dd\zeta=\mu_{0}'(w)\,\dd w$,
legitimate on a neighbourhood of $\mu_{0}$ because $F'(\mu_{0})\ne0$:
\[
 \operatorname*{Res}_{\zeta=\mu_{0}}
 \frac{E(\zeta)^{M}}{(F(\zeta)-y)^{M}}\,\dd\zeta
 =\operatorname*{Res}_{w=y}
 \frac{E(\mu_{0}(w))^{M}\,\mu_{0}'(w)}{(w-y)^{M}}\,\dd w
 =\frac1{(M-1)!}\frac{\dd^{M-1}}{\dd y^{M-1}}
 \left(\frac{E(\mu_{0}(y))^{M}}{F'(\mu_{0}(y))}\right),
\]
the last step being the residue of a pole of order $M$.  Multiplying by
$(-1)^{M}/M$ and using $\dd/\dd y=\mathcal{D}$ gives
\eqref{p0:eq:perturbed-inversion} and \eqref{p0:eq:perturbed-residue}.

For \eqref{p0:eq:perturbed-multi} replace $E$ by $\sum_{j}\varepsilon_{j}E_{j}$
and expand the $M$th power by the multinomial theorem: the monomial
$\boldsymbol\varepsilon^{\boldsymbol m}$ occurs with coefficient
$M!/\boldsymbol m!$, which cancels the $M!$ in
\eqref{p0:eq:perturbed-inversion}.
\end{proof}

The identity \eqref{p0:eq:perturbed-inversion} is needed in Chapter~4 with a
correction $E$ that is a \emph{divergent} formal series, where no contour is
available.  What survives is the coefficientwise identity, and it survives for
the reason the Lagrange fixed-point formula does.

\begin{corollary}[Formal form of the reversion]\label{p0:cor:perturbed-formal}
Fix $M\ge1$ and let
$\mathcal{P}:=\Rat[e_{0},\dots,e_{M-1},f_{2},\dots,f_{M}][f_{1}^{-1}]$.  Both
of the following are elements of $\mathcal{P}$, obtained by writing out the
indicated recursions with $e_{i}$ and $f_{i}$ standing for $E^{(i)}(\mu_{0})$
and $F^{(i)}(\mu_{0})$:
\begin{enumerate}
\item the coefficient $\mu_{M}$ of $\varepsilon^{M}$ produced by solving
\eqref{p0:eq:perturbed-equation} by formal Taylor iteration, i.e.\ by
substituting $\mu=\mu_{0}+\sum_{k\ge1}\mu_{k}\varepsilon^{k}$, expanding $F$
and $E$ in their Taylor jets at $\mu_{0}$, and equating coefficients of
$\varepsilon^{M}$;
\item the expression
$\frac{(-1)^{M}}{M!}\mathcal{D}^{M-1}\bigl(E(\mu_{0})^{M}/F'(\mu_{0})\bigr)$
of \eqref{p0:eq:perturbed-inversion}, expanded by the Leibniz rule.
\end{enumerate}
They are equal in $\mathcal{P}$.  Consequently
\eqref{p0:eq:perturbed-inversion} and \eqref{p0:eq:perturbed-multi} hold
coefficientwise in any commutative $\Rat$-algebra carrying jets
$e_{i},f_{i}$ with $f_{1}$ invertible --- in particular when $F$ and $E$ are
formal, possibly divergent, series, and when $\mathcal{D}$ is the derivation
\eqref{p0:eq:reversion-operator} written in any coordinate.
\end{corollary}

\begin{proof}
Each recursion visibly terminates and involves only $e_{0},\dots,e_{M-1}$,
$f_{2},\dots,f_{M}$ and divisions by $f_{1}$: in (i) the coefficient of
$\varepsilon^{M}$ in $F(\mu_{0}+\sum_{k\ge1}\mu_{k}\varepsilon^{k})
+\varepsilon E(\mu_{0}+\sum_{k\ge1}\mu_{k}\varepsilon^{k})$ equals
$f_{1}\mu_{M}$ plus a polynomial in $f_{2},\dots,f_{M}$,
$e_{0},\dots,e_{M-1}$ and $\mu_{1},\dots,\mu_{M-1}$, so $\mu_{M}$ is
determined by one division by $f_{1}$; in (ii) each application of
$\mathcal{D}=f_{1}^{-1}\dd/\dd\mu$ lowers the available jet order by one and
contributes one factor $f_{1}^{-1}$, and $M-1$ applications to a function of
$e_{0},\dots$ and $f_{1},\dots$ reach at most $e_{M-1}$ and $f_{M}$.  So both
lie in $\mathcal{P}$.

Choose $N$ with $f_{1}^{N}\mu_{M}$ and $f_{1}^{N}\cdot(\text{ii})$ both in
$\mathcal{A}:=\Rat[e_{0},\dots,e_{M-1},f_{1},\dots,f_{M}]$, and regard them as
polynomial functions on $\Cplx^{2M}$.  Every point of $\Cplx^{2M}$ with
$f_{1}\ne0$ is realized by a pair of \emph{polynomials}: take $\mu_{0}=0$,
$E(\mu):=\sum_{i=0}^{M-1}e_{i}\mu^{i}/i!$ and
$F(\mu):=\sum_{i=1}^{M}f_{i}\mu^{i}/i!$, which are entire with
$F'(0)=f_{1}\ne0$ and with the prescribed jets.  For such a pair
\cref{p0:thm:perturbed-inversion} applies, and the Taylor coefficients in
$\varepsilon$ of the analytic root $\mu(y;\varepsilon)$ satisfy the recursion
of (i) --- substituting the convergent expansion into
\eqref{p0:eq:perturbed-equation} and equating coefficients is legitimate ---
so (i) and (ii) agree at that point.  Two elements of $\mathcal{A}$ agreeing on
the Zariski-dense open set $\{f_{1}\ne0\}$ of $\Cplx^{2M}$ are equal in
$\Cplx[e_{\bullet},f_{\bullet}]$, hence equal in $\mathcal{A}$ and, after
dividing by $f_{1}^{N}$, in $\mathcal{P}$.  The final assertion follows by
applying the unique $\Rat$-algebra homomorphism $\mathcal{P}\to S$ determined
by the given jets, exactly as in \cref{p0:lem:LB-universal}.
\end{proof}

\section{The Lambert core}\label{p0:sec:core}

\subsection{Exact inversion of the linear--logarithmic phase}

\begin{theorem}[Exact solution of the dominant block]\label{p0:thm:lambert-core}
Let $a>0$, $b\in\Real$, and consider, for $X>0$,
\begin{equation}\label{p0:eq:core-equation}
 \Phi_0(X):=aX+b\log X=L .
\end{equation}
\begin{enumerate}
\item If $b=0$ the unique solution is $X=L/a$.
\item If $b\ne0$, put $W:=(a/b)X$.  Then \eqref{p0:eq:core-equation} holds if
      and only if
      \begin{equation}\label{p0:eq:core-W-equation}
       W\e^{W}=z(L),\qquad z(L):=\frac ab\,\e^{L/b},
      \end{equation}
      so that the solutions of \eqref{p0:eq:core-equation} are exactly
      \begin{equation}\label{p0:eq:core-lambert}
       X=\frac ba\,\Wlam_k\!\left(\frac ab\,\e^{L/b}\right)
      \end{equation}
      as $k$ ranges over the branches of the Lambert function for which
      $\Wlam_k(z(L))$ is defined.
\item \emph{Branch rule.}  Suppose $b\ne0$.
      \begin{itemize}
      \item If $b>0$ then $z(L)>0$, equation \eqref{p0:eq:core-equation} has
            exactly one positive solution for every real $L$, $\Phi_0$ is
            strictly increasing on $(0,\infty)$, and
            $X=(b/a)\Wlam_0(z(L))$.
      \item If $b<0$ then $z(L)=-(a/|b|)\e^{-L/|b|}$ lies in $(-\e^{-1},0)$
            precisely when
            \begin{equation}\label{p0:eq:core-threshold}
             L>L_c:=|b|\Bigl(1-\log\frac{|b|}{a}\Bigr)
                  =\Phi_0\!\left(\frac{|b|}{a}\right),
            \end{equation}
            and then \eqref{p0:eq:core-equation} has exactly two positive
            solutions, namely
            $(b/a)\Wlam_0(z(L))<|b|/a<(b/a)\Wlam_{-1}(z(L))$.
            $\Phi_0$ is strictly increasing on $(|b|/a,\infty)$ and the large
            solution is
            \begin{equation}\label{p0:eq:core-lambert-minus}
             X=\frac ba\,\Wlam_{-1}\!\left(\frac ab\,\e^{L/b}\right)
              =\frac{|b|}{a}\Bigl(-\Wlam_{-1}\bigl(z(L)\bigr)\Bigr).
            \end{equation}
      \end{itemize}
      In both cases the large solution satisfies $X\to+\infty$ and
      $X=L/a+\bigO(\log L)$ as $L\to+\infty$.
\item \emph{Slope identities.}  With $W=(a/b)X$ and $b\ne0$,
      \begin{equation}\label{p0:eq:core-slope}
       \Phi_0'(X)=a+\frac bX=a\,\frac{1+W}{W},
       \qquad
       \frac{\dd X}{\dd L}=\frac1{\Phi_0'(X)}=\frac1a\,\frac{W}{1+W} .
      \end{equation}
\end{enumerate}
\end{theorem}

\begin{proof}
(1) is immediate.  For (2), multiply \eqref{p0:eq:core-equation} by $1/b$ and
exponentiate:
\[
 \e^{L/b}=\e^{(a/b)X}\,X=\e^{W}\cdot\frac ba\,W ,
\]
because $X=(b/a)W$.  Multiplying by $a/b$ gives
\eqref{p0:eq:core-W-equation}; the step is reversible, and no logarithm of a
negative number occurs, which is what makes the argument uniform in the sign of
$b$.  By definition, the solutions of $W\e^{W}=z$ are the values $\Wlam_k(z)$,
whence \eqref{p0:eq:core-lambert}.

(3) We have $\Phi_0'(X)=a+b/X$.  For $b>0$ this is positive on $(0,\infty)$
and $\Phi_0$ increases from $-\infty$ to $+\infty$, so there is exactly one
solution; correspondingly $W=(a/b)X>0$ and $\Wlam_0$ is the unique real branch
on $(0,\infty)$.  For $b<0$, $\Phi_0'$ vanishes only at $X_\ast=|b|/a$, is
negative on $(0,X_\ast)$ and positive on $(X_\ast,\infty)$; $\Phi_0$ decreases
from $+\infty$ to $\Phi_0(X_\ast)=|b|-|b|\log(|b|/a)=L_c$ and then increases
back to $+\infty$.  Hence there are two positive solutions exactly when
$L>L_c$, one in each monotonicity interval.  On the $W$ side $W=(a/b)X<0$, and
the two real branches of $\Wlam$ on $(-\e^{-1},0)$ are $\Wlam_0$, with values
in $(-1,0)$, and $\Wlam_{-1}$, with values in $(-\infty,-1)$; the map
$W\mapsto(b/a)W$ reverses order, so $\Wlam_{-1}$ produces the larger $X$,
namely the one exceeding $|b|/a$.  Also $z(L)>-\e^{-1}$ is equivalent to
$(a/|b|)\e^{-L/|b|}<\e^{-1}$, that is to \eqref{p0:eq:core-threshold}.
Finally $aX=L-b\log X$ with $X\to\infty$ gives $aX=L+\bigO(\log L)$.

(4) $a+b/X=a+b\cdot\dfrac{a}{bW}=a\bigl(1+W^{-1}\bigr)$, and $\dd X/\dd L$ is
the reciprocal of $\Phi_0'$ by the inverse function theorem, applicable
because $\Phi_0'\ne0$ on the relevant interval.
\end{proof}

\begin{corollary}[Branch rule]\label{p0:cor:branch-rule}
For $b\ne0$ the branch of the Lambert function carrying the large solution of
$L=aX+b\log X$ is determined by the sign of $b$ alone: $\Wlam_0$ if $b>0$,
$\Wlam_{-1}$ if $b<0$.
\end{corollary}

\begin{remark}[Why the branch is not a detail]\label{p0:rem:branch}
Three of the four chapters have $b<0$ and therefore use $\Wlam_{-1}$; the odd
phase of the swing factorial has $b=+1/2$ and uses $\Wlam_0$.
\textsf{Swing\_Factorial\_Transseries\_Article} puts the point sharply: the
principal branch ``would produce the small root of the simplified equation and
is irrelevant to the large-index inverse''.  For $b<0$ the two roots are the
two sides of a genuine fold at $X_\ast=|b|/a$, and choosing wrongly yields an
object all of whose asymptotic corrections are then computed correctly about
the wrong point.
\end{remark}

\subsection{The factorial core}

\begin{proposition}[Exact inversion of a log-linear phase]
\label{p0:prop:factorial-core}
Let $\kappa>0$ and $d\in\Real$, and consider
\begin{equation}\label{p0:eq:factorial-core-equation}
 \Phi_0(X):=X\bigl(\kappa\log X+d\bigr)=L,\qquad X>0 .
\end{equation}
For $L>0$ put
\begin{equation}\label{p0:eq:factorial-core-solution}
 v:=\Wlam_0\!\left(\frac L\kappa\,\e^{d/\kappa}\right),
 \qquad
 X:=\exp\Bigl(v-\frac d\kappa\Bigr) .
\end{equation}
Then $X$ is the unique solution of \eqref{p0:eq:factorial-core-equation} with
$\kappa\log X+d>0$, and
\begin{equation}\label{p0:eq:factorial-core-slope}
 \Phi_0'(X)=\kappa\log X+d+\kappa=\kappa\,(v+1),
 \qquad
 \frac{\dd X}{\dd L}=\frac1{\kappa(v+1)} .
\end{equation}
Moreover $\Phi_0$ is an admissible core in the sense of \cref{p0:def:core} on
$\bigl(\e^{-(d+\kappa)/\kappa},\infty\bigr)$, with $\Phi_0'(X)\to+\infty$.
\end{proposition}

\begin{proof}
Put $s=\log X$.  Then \eqref{p0:eq:factorial-core-equation} reads
$\e^{s}(\kappa s+d)=L$, that is
$\bigl(s+\tfrac d\kappa\bigr)\e^{s+d/\kappa}=(L/\kappa)\e^{d/\kappa}$.  The
right-hand side is positive, so $\Wlam_0$ is the unique real branch and
$s+d/\kappa=v$, which is \eqref{p0:eq:factorial-core-solution}.  Uniqueness
under the stated side condition holds because $\kappa s+d>0$ forces
$s+d/\kappa>0$, and $\sigma\mapsto\sigma\e^{\sigma}$ is injective on
$(0,\infty)$.  Differentiating,
$\Phi_0'(X)=\kappa\log X+d+\kappa=\kappa(s+d/\kappa)+\kappa=\kappa(v+1)$.
Admissibility: $\Phi_0'$ is increasing, tends to $+\infty$, is positive exactly
for $\log X>-(d+\kappa)/\kappa$, and $X\Phi_0'(X)\to\infty$.
\end{proof}

\begin{remark}[The double factorial normalization]\label{p0:rem:factorial-core}
\Cref{p2:sec:top} uses $\kappa=\tfrac12$, $d=-\tfrac12$ and the centred
variable $s=x+1$, so that $\Phi_0(s)=\tfrac s2(\log s-1)$,
$(L/\kappa)\e^{d/\kappa}=2L/\e$, $X=\e^{v+1}$ with $v=\Wlam_0(2L/\e)$, and
$\Phi_0'(X)=\tfrac12(v+1)=\tfrac12\log X$.  This reproduces the core of all
three double factorial articles; the quantity they call $\Lambda$ is
$v+1=\log X$.
\end{remark}

\section{All-orders reversion around the core}\label{p0:sec:reversion}

\subsection{The Lambert-centred inverse}

\begin{theorem}[All-orders reversion around the Lambert core]
\label{p0:thm:lambert-centered}
Let $R$ be a commutative $\Rat$-algebra, let $a\in R$ be invertible, let
$b\in R$, and let $Q=\sum_{j\ge1}q_jt^j\in tR[[t]]$.  Define
\begin{equation}\label{p0:eq:Psi-general}
 \Psi(t,u):=-\frac1a\left[\,b\log(1+tu)
   +Q\!\left(\frac{t}{1+tu}\right)\right].
\end{equation}
\begin{enumerate}
\item $\Psi\in tR[[t,u]]$.
\item There is a unique $U=\sum_{n\ge1}c_nt^n\in tR[[t]]$ with
      $U=\Psi(t,U)$, and
      \begin{equation}\label{p0:eq:c-lagrange}
       c_n=\sum_{m=1}^{n}\frac1m\Coef{t^nu^{m-1}}\Psi(t,u)^m
       \qquad(n\ge1).
      \end{equation}
\item \emph{Triangular Bell recurrence.}  Put $\gamma_1:=0$ and
      $\gamma_r:=c_{r-1}$ for $r\ge2$.  Then, for $n\ge1$,
      \begin{equation}\label{p0:eq:c-bell}
       c_n=-\frac1a\left[
         b\sum_{k=1}^{n}\frac{(-1)^{k+1}}{k}\OB_{n,k}(\gamma)
         +\sum_{j=1}^{n}q_j\sum_{k=0}^{n-j}(-1)^{k}
            \binom{j+k-1}{k}\OB_{n-j,k}(\gamma)
       \right],
      \end{equation}
      and only $c_1,\dots,c_{n-1}$ occur on the right-hand side.  In
      particular $c_n$ lies in the $\Rat$-subalgebra of $R$ generated by
      $a^{-1}$, $b$ and $q_1,\dots,q_n$; no logarithm appears.
\item \emph{Formal exactness.}  In $t^{-1}R[[t]]$ put
      \begin{equation}\label{p0:eq:centered-ansatz}
       \mathsf x:=t^{-1}+U(t)=t^{-1}+\sum_{n\ge1}c_nt^n .
      \end{equation}
      Then $t\mathsf x=1+tU\in1+tR[[t]]$ and $1/\mathsf x=t/(1+tU)\in tR[[t]]$,
      and the following identity holds in $R[[t]]$:
      \begin{equation}\label{p0:eq:formal-exactness}
       a\bigl(\mathsf x-t^{-1}\bigr)+b\log(t\mathsf x)
       +Q\bigl(1/\mathsf x\bigr)=0 .
      \end{equation}
      Equivalently: writing $t=X^{-1}$ and
      $x=X+\sum_{n\ge1}c_nX^{-n}$, the quantity $ax+b\log x+Q(1/x)$ equals
      $aX+b\log X$ to all orders.
\item \emph{Analytic version.}  Let $F$ be differentiably of
      exponential--power type with data $(C,a,1,b,1,Q)$ on $[x_0,\infty)$
      (\cref{p0:def:model}), strictly increasing there, with
      $F(x)\to\infty$.  For $y$ in the range of $F$ set $L=\log(y/C)$ and let
      $X=X(y)$ be the large core solution of $L=aX+b\log X$ furnished by
      \cref{p0:thm:lambert-core}.  Then $X(y)\to\infty$ and, for every fixed
      $N\ge0$,
      \begin{equation}\label{p0:eq:centered-asymptotic}
       F^{-1}(y)=X+\sum_{n=1}^{N}\frac{c_n}{X^{n}}
                 +\bigO\bigl(X^{-N-1}\bigr)
       \qquad(y\to\infty).
      \end{equation}
\end{enumerate}
\end{theorem}

\begin{proof}
(1) $\log(1+tu)=\sum_{k\ge1}(-1)^{k+1}(tu)^k/k\in tR[[t,u]]$, and
$t/(1+tu)=t\sum_{k\ge0}(-tu)^k\in tR[[t,u]]$, so
$(t/(1+tu))^j\in t^jR[[t,u]]$ and the sum over $j$ is summable; hence
$Q(t/(1+tu))\in tR[[t,u]]$.

(2) is \cref{p0:lem:lagrange-fp}.

(3) Put $w:=tU=\sum_{r\ge1}\gamma_rt^r$ with $\gamma$ as stated; $\gamma_1=0$
because $U\in tR[[t]]$.  Then $\Coef{t^n}\log(1+w)$ is given by
\eqref{p0:eq:log-obell}, while
\[
 \Coef{t^n}Q\!\left(\frac t{1+w}\right)
 =\sum_{j\ge1}q_j\,\Coef{t^{n-j}}(1+w)^{-j}
 =\sum_{j=1}^{n}q_j\sum_{k=0}^{n-j}(-1)^k\binom{j+k-1}{k}\OB_{n-j,k}(\gamma)
\]
by \eqref{p0:eq:power-negative}.  Since $c_n=\Coef{t^n}\Psi(t,U)$, this is
\eqref{p0:eq:c-bell}.  For triangularity, $\OB_{n,k}(\gamma)$ depends only on
$\gamma_1,\dots,\gamma_{n-k+1}$, that is on $c_1,\dots,c_{n-k}$; in the first
sum $k\ge1$ and in the second $j\ge1$, so $c_n$ never occurs.  (In fact
$\OB_{n,k}(\gamma)=0$ for $2k>n$, because $\gamma_1=0$.)

(4) With $\mathsf x=t^{-1}(1+w)$, $w=tU$, one has
$a(\mathsf x-t^{-1})=aU$, $\log(t\mathsf x)=\log(1+w)$ and
$1/\mathsf x=t/(1+w)$, so the left-hand side of
\eqref{p0:eq:formal-exactness} equals
$aU+b\log(1+tU)+Q(t/(1+tU))=a\bigl(U-\Psi(t,U)\bigr)=0$.

(5) Fix $N\ge0$, set $t:=1/X$, $U_N(t):=\sum_{n=1}^{N}c_nt^n$ and
$x_N:=X+U_N(t)=X\bigl(1+tU_N(t)\bigr)$.  By
\cref{p0:thm:lambert-core}(3), $X\to\infty$ and $X\sim L/a$; since
$tU_N(t)=\bigO(t^2)$ we get $x_N=X\bigl(1+\bigO(X^{-2})\bigr)$, so
$x_N\to\infty$ and $x_N\ge x_0$ for all large $y$.

\emph{Residual.}  Apply \eqref{p0:eq:model-analytic} with $M=N$:
\[
 \log F(x_N)=\log C+ax_N+b\log x_N+\sum_{j=1}^{N}q_jx_N^{-j}+E_N,
 \qquad |E_N|\le K_Nx_N^{-N-1}.
\]
Since $\log y=\log C+aX+b\log X$ and $\log x_N=\log X+\log(1+tU_N(t))$,
\[
 r_N:=\log F(x_N)-\log y=g(t)+E_N,
 \qquad
 g(t):=aU_N(t)+b\log\bigl(1+tU_N(t)\bigr)
   +\sum_{j=1}^{N}q_j\left(\frac{t}{1+tU_N(t)}\right)^{j}.
\]
Here $g$ is real-analytic on a neighbourhood of $t=0$, being a finite
composition of analytic maps with $1+tU_N(t)\to1$; its Taylor series at $t=0$
is therefore the corresponding formal series.  By (4) the formal series
$aU+b\log(1+tU)+Q(t/(1+tU))$ vanishes identically.  Replacing $U$ by $U_N$
changes it by an element of $t^{N+1}R[[t]]$, because $U-U_N\in t^{N+1}R[[t]]$
and the maps $u\mapsto\log(1+tu)$, $u\mapsto(t/(1+tu))^j$ are $t$-adically
continuous; and discarding the terms of $Q$ with $j>N$ changes it by an element
of $t^{N+1}R[[t]]$ as well.  Hence the Taylor series of $g$ vanishes through
order $t^{N}$, so $g(t)=\bigO(t^{N+1})$ and
\[
 |r_N|\le\bigl|g(1/X)\bigr|+K_Nx_N^{-N-1}=\bigO\bigl(X^{-N-1}\bigr).
\]

\emph{Slope.}  By \eqref{p0:eq:model-derivative} with $\alpha=1$,
$(\log F)'(x)=a+\bigO(x^{-1})$, so there is $x_1\ge x_0$ with
$(\log F)'\ge a/2$ on $[x_1,\infty)$.  For all large $y$ both $x_N$ and
$F^{-1}(y)$ lie in $[x_1,\infty)$.  \Cref{p0:thm:backward-error}, applied to
$\log F$ on that half-line, gives
\[
 \bigl|x_N-F^{-1}(y)\bigr|\le\frac{|r_N|}{a/2}
 =\bigO\bigl(X^{-N-1}\bigr),
\]
which is \eqref{p0:eq:centered-asymptotic}.
\end{proof}

\begin{remark}[No logarithms in the coefficients]\label{p0:rem:no-logs}
Part~(3) is the structural reason for centring at the Lambert core: the whole
logarithmic tower has been resummed into $X$, so the $c_n$ are constants of the
coefficient ring.  The alternative expansion of \cref{p0:sec:flattening} has
coefficients that are polynomials in $\log L$ whose degree grows with the
order.  This is an algebraic statement and by itself says nothing about
numerical accuracy; the numerical comparison is a separate, chapterwise matter
(see \cref{p0:rem:accuracy-claim}).
\end{remark}

\begin{proposition}[General exponent grid]\label{p0:prop:two-grid}
Let $\delta>0$ and keep the remaining hypotheses of
\cref{p0:thm:lambert-centered}.  Let $s$ be a second formal variable and put
\begin{equation}\label{p0:eq:Psi-two-grid}
 \Psi_\delta(t,s,u):=-\frac1a\left[\,b\log(1+tu)
   +Q\!\left(\frac{s}{(1+tu)^{\delta}}\right)\right],
\end{equation}
where $(1+tu)^{-\delta}$ is the formal binomial series
\eqref{p0:eq:power-general}.  Then $\Psi_\delta$ lies in the ideal $(t,s)$ of
$R[[t,s,u]]$, the equation $U=\Psi_\delta(t,s,U)$ has a unique solution
$U=\sum_{i+j\ge1}c_{i,j}t^is^j$ in the ideal $(t,s)$ of $R[[t,s]]$, and the
substitution $t=X^{-1}$, $s=X^{-\delta}$ reverts \eqref{p0:eq:model-log}:
formally
\begin{equation}\label{p0:eq:two-grid-exactness}
 ax+b\log x+Q\bigl(x^{-\delta}\bigr)=aX+b\log X,
 \qquad x=X+\sum_{i+j\ge1}c_{i,j}X^{-i-j\delta}.
\end{equation}
For $\delta=1$ the two markers coincide and
$\Psi_\delta(t,t,u)=\Psi(t,u)$ of \eqref{p0:eq:Psi-general}.
\end{proposition}

\begin{proof}
The membership $\Psi_\delta\in(t,s)$ and the existence and uniqueness of $U$
follow exactly as in parts (1) and (2) of \cref{p0:thm:lambert-centered}, with
the $t$-adic filtration replaced by the $(t,s)$-adic one: $\Psi_\delta$ lies in
$(t,s)$, so $\Psi_\delta(t,s,U)$ modulo $(t,s)^{d+1}$ depends only on $U$
modulo $(t,s)^{d}$.  The identity \eqref{p0:eq:two-grid-exactness} is the
computation of part (4) with $1/\mathsf x=t/(1+tU)$ replaced by
$\mathsf x^{-\delta}=s(1+tU)^{-\delta}$.  Finally the substitution is
legitimate because, for each real exponent $\theta$, only finitely many pairs
$(i,j)$ with $i,j\ge0$ satisfy $i+j\delta=\theta$ when $\delta>0$.
\end{proof}

\begin{remark}[When the grid degenerates]\label{p0:rem:grid}
If $\delta$ is irrational the exponents $i+j\delta$ are pairwise distinct and
the monomial scale is a free monoid on two generators.  If $\delta=p/q$ in
lowest terms, different pairs $(i,j)$ can give the same exponent and the
correct statement is that the reversion lives in $R[[X^{-1/q}]]$.  Every
chapter of this volume has $\delta/\alpha=1$ after
\cref{p0:lem:alpha-reduction}, so the degeneracy is never met here.
\end{remark}

\subsection{Reversion about an arbitrary admissible core}

\begin{theorem}[Reversion about an arbitrary core]\label{p0:thm:core-reversion}
Let $R$ be a commutative $\Rat$-algebra and $\Lambda\in R$ invertible.  Consider
\begin{equation}\label{p0:eq:core-master}
 \Lambda E+h(E)+f(t,E)=0,
 \qquad h\in u^2R[[u]],\quad f\in tR[[t,u]] .
\end{equation}
Then \eqref{p0:eq:core-master} has a unique solution
$E=\sum_{n\ge1}e_nt^n\in tR[[t]]$; it satisfies the triangular recurrence
\begin{equation}\label{p0:eq:core-triangular}
 e_n=-\frac1\Lambda\Coef{t^n}
   \Bigl[h\bigl(E_{<n}\bigr)+f\bigl(t,E_{<n}\bigr)\Bigr],
 \qquad E_{<n}:=\sum_{j=1}^{n-1}e_jt^j ,
\end{equation}
and the closed formula
\begin{equation}\label{p0:eq:core-lagrange}
 e_n=\sum_{m=1}^{2n-1}\frac1m
   \Coef{t^nu^{m-1}}\Bigl(-\Lambda^{-1}\bigl[h(u)+f(t,u)\bigr]\Bigr)^{m}.
\end{equation}
If moreover $h(u)=\sum_{m\ge2}h_mu^m$ and
$f(t,u)=\sum_{k\ge1}\phi_k(t)\,(1+u)^{-\mu_k}$ with $\phi_k\in tR[[t]]$ and
$\mu_k\in\Nat$, then
\begin{equation}\label{p0:eq:core-bell}
 e_n=-\frac1\Lambda\left[
  \sum_{m=2}^{n}h_m\,\OB_{n,m}(e)
  +\sum_{k\ge1}\sum_{i=1}^{n}\bigl(\Coef{t^i}\phi_k\bigr)
    \sum_{m=0}^{n-i}(-1)^m\binom{\mu_k+m-1}{m}\OB_{n-i,m}(e)
 \right],
\end{equation}
in which only $e_1,\dots,e_{n-1}$ occur.
\end{theorem}

\begin{proof}
Existence, uniqueness and \eqref{p0:eq:core-lagrange} are
\cref{p0:lem:lagrange-fp-quadratic} applied to $\Psi:=-\Lambda^{-1}[h+f]$.
For \eqref{p0:eq:core-triangular}, note as in the proof of that lemma that
$e_n$ occurs neither in $\Coef{t^n}h(E)$ nor in $\Coef{t^n}f(t,E)$, so both may
be computed from $E_{<n}$.  For \eqref{p0:eq:core-bell} apply
\cref{p0:def:obell} to $h$ and \eqref{p0:eq:power-negative} to each
$(1+E)^{-\mu_k}$, then convolve with $\phi_k$; the displayed ranges are those
outside which the ordinary Bell polynomials vanish.
\end{proof}

\begin{remark}[How the chapters instantiate it]\label{p0:rem:core-instances}
For the linear--logarithmic core the shorter route of
\cref{p0:thm:lambert-centered} is available, because measuring the displacement
in units of $X$ makes the entire perturbation divisible by $t$.
\Cref{p2:sec:top} needs the general form: there $\Lambda$ is the normalized
core slope $\log X$ of \cref{p0:prop:factorial-core},
$h(E)=(1+E)\log(1+E)-E$ and
$f(t,E)=\sum_{k\ge1}\alpha_kt^{2k}(1+E)^{-(2k-1)}$, and
\eqref{p0:eq:core-bell} is exactly the Bell recurrence of
\textsf{double\_factorial\_transseries-2}.  The coefficient ring there is
$\Rat[\Lambda^{-1}]$: the $e_n$ are \emph{not} constants but polynomials in
$\Lambda^{-1}$, of degree at most $2n-1$ and with no constant term.
\end{remark}

\subsection{The perturbation-operator form}

\begin{proposition}[Inverse perturbation operator]\label{p0:prop:operator-form}
Let $\Phi_0$ be a phase with $\Phi_0'$ nonvanishing near the point considered,
let $X$ solve $\Phi_0(X)=L$, and consider
\begin{equation}\label{p0:eq:perturbed-problem}
 L=\Phi_0(x)+\varepsilon\,\varrho(x).
\end{equation}
Put $\mathscr D_0:=\Phi_0'(X)^{-1}\,\dd/\dd X$.  Then, as a formal power series
in $\varepsilon$,
\begin{equation}\label{p0:eq:operator-series}
 x=X+\sum_{m\ge1}\frac{(-\varepsilon)^m}{m!}\,
   \mathscr D_0^{\,m-1}\!\left[\frac{\varrho(X)^m}{\Phi_0'(X)}\right].
\end{equation}
\end{proposition}

\begin{proof}
Change coordinate to $\lambda:=\Phi_0(x)$, so that $x=\Phi_0^{-1}(\lambda)$ and
\eqref{p0:eq:perturbed-problem} becomes
$L=\lambda+\varepsilon\varrho_0(\lambda)$ with
$\varrho_0:=\varrho\circ\Phi_0^{-1}$.  This is a near-identity reversion in
$\lambda$ with small parameter $\varepsilon$; Lagrange--B\"urmann in the form
\eqref{p0:eq:LB-general}, with the observable $G=\Phi_0^{-1}$, gives
\[
 \Phi_0^{-1}(\lambda)
 =\Phi_0^{-1}(L)+\sum_{m\ge1}\frac{(-\varepsilon)^m}{m!}
  \left(\frac{\dd}{\dd L}\right)^{m-1}
  \Bigl[(\Phi_0^{-1})'(L)\,\varrho_0(L)^m\Bigr].
\]
Finally $(\Phi_0^{-1})'(L)=1/\Phi_0'(X)$, $\varrho_0(L)=\varrho(X)$ and
$\dd/\dd L=\mathscr D_0$ by the chain rule.
\end{proof}

\begin{remark}[What the operator form is for]\label{p0:rem:operator}
This is the form in which an \emph{exponentially small} forward correction is
transported through the inverse map: with $\varrho$ a Stokes discontinuity or a
terminant remainder rather than a power series, $-\varepsilon\varrho(X)/\Phi_0'(X)$
is the leading displacement and the higher terms generate the nonlinear sectors.
As a statement about a formal series in $\varepsilon$ it is exact; as an
asymptotic statement about a particular $\varrho$ it needs the hypotheses of
\cref{p0:thm:remainder-transport}.  Expanding $\varrho$ in inverse powers of $X$
reproduces \eqref{p0:eq:c-lagrange}.
\end{remark}

\section{Flattening the core}\label{p0:sec:flattening}

\Cref{p0:thm:lambert-centered} hides the explicit scale of the answer inside
$X$.  Unfolding $X$ produces the polynomial--logarithmic transseries in which
the four chapters state their headline formulae.  Throughout this section $H$
is an indeterminate; in the analytic application it is specialized to
$\log(L/a)$.

\begin{theorem}[Direct polynomial--logarithmic reversion]\label{p0:thm:flattening}
Let $R$ be a commutative $\Rat$-algebra, $a\in R$ invertible, $b\in R$,
$Q=\sum_{j\ge1}q_jt^j\in tR[[t]]$, and let $H$ be an indeterminate.  Set
\begin{equation}\label{p0:eq:Phi-direct}
 \Phi(t,u;H):=-b\log\bigl(1+t(u-bH)\bigr)
   -Q\!\left(\frac{at}{1+t(u-bH)}\right)
 \ \in\ t\,R[H][[t,u]].
\end{equation}
\begin{enumerate}
\item There is a unique $V=\sum_{n\ge1}P_n(H)\,t^n\in t\,R[H][[t]]$ with
      $V=\Phi(t,V;H)$, and
      \begin{equation}\label{p0:eq:P-lagrange}
       P_n(H)=\sum_{m=1}^{n}\frac1m\Coef{t^nu^{m-1}}\Phi(t,u;H)^m .
      \end{equation}
\item \emph{Triangular Bell recurrence.}  Put $\gamma_1:=-bH$ and
      $\gamma_r:=P_{r-1}(H)$ for $r\ge2$.  Then
      \begin{equation}\label{p0:eq:P-bell}
       P_n(H)=-b\sum_{k=1}^{n}\frac{(-1)^{k+1}}{k}\OB_{n,k}(\gamma)
        -\sum_{j=1}^{n}q_ja^{\,j}\sum_{k=0}^{n-j}(-1)^k\binom{j+k-1}{k}
          \OB_{n-j,k}(\gamma),
      \end{equation}
      in which only $P_1,\dots,P_{n-1}$ occur.
\item \emph{Shape.}  $P_n\in R[H]$ with
      \begin{equation}\label{p0:eq:P-shape}
       \deg_HP_n\le n,\qquad
       \Coef{H^n}P_n=\frac{b^{\,n+1}}{n},\qquad
       P_n(0)=a^{\,n+1}c_n ,
      \end{equation}
      where $c_n$ are the Lambert-centred coefficients of
      \cref{p0:thm:lambert-centered}.  In particular $\deg_HP_n=n$ exactly
      when $b$ is not a zero divisor, and the coefficients of $P_n$ lie in the
      $\Rat$-subalgebra generated by $a$, $a^{-1}$, $b$, $q_1,\dots,q_n$.
\item \emph{Formal exactness and uniqueness.}  Put $t=L^{-1}$ and
      $H=\log(L/a)$, and
      \begin{equation}\label{p0:eq:direct-ansatz}
       x:=\frac1a\left[L-bH+\sum_{n\ge1}\frac{P_n(H)}{L^{n}}\right].
      \end{equation}
      Then $ax+b\log x+Q(1/x)=L$ as an identity of formal series in $L^{-1}$
      with polynomial coefficients in $H$.  Conversely, \eqref{p0:eq:direct-ansatz}
      is the \emph{only} expansion of the form
      $x=a^{-1}\bigl[L-bH+\sum_{n\ge1}\Pi_n(H)L^{-n}\bigr]$,
      with $\Pi_n\in R[H]$, that solves the equation coefficientwise.
\item \emph{Analytic version.}  Under the hypotheses of
      \cref{p0:thm:lambert-centered}(5), with $L=\log(y/C)$ and
      $H=\log(L/a)$, for every fixed $N\ge0$,
      \begin{equation}\label{p0:eq:direct-asymptotic}
       F^{-1}(y)=\frac1a\left[L-bH+\sum_{n=1}^{N}\frac{P_n(H)}{L^{n}}\right]
        +\bigO\!\left(\frac{(1+H)^{N+1}}{L^{N+1}}\right)
       \qquad(y\to\infty).
      \end{equation}
\end{enumerate}
\end{theorem}

\begin{proof}
(1) Membership of $\Phi$ in $tR[H][[t,u]]$ is checked as in
\cref{p0:thm:lambert-centered}(1), the extra additive term $-bH$ inside
$1+t(u-bH)$ being harmless because it is multiplied by $t$.
\Cref{p0:lem:lagrange-fp} over the ring $R[H]$ then gives existence,
uniqueness and \eqref{p0:eq:P-lagrange}.

(2) Put $w:=t(V-bH)=\sum_{r\ge1}\gamma_rt^r$ with $\gamma$ as stated; apply
\eqref{p0:eq:log-obell} to $\log(1+w)$ and \eqref{p0:eq:power-negative} to
$(1+w)^{-j}$, and convolve the latter with the factor $(at)^j$.  Only
$\gamma_1,\dots,\gamma_{n-k+1}$, that is $H$ and $P_1,\dots,P_{n-k}$, occur.

(4) Set $t=L^{-1}$.  With $x$ as in \eqref{p0:eq:direct-ansatz},
\[
 x=\frac1{at}\bigl(1+t(V-bH)\bigr)=\frac{1+w}{at},
 \qquad ax=\frac1t+(V-bH),\qquad
 \log x=H+\log(1+w),\qquad \frac1x=\frac{at}{1+w},
\]
where $\log x$ uses $\log\bigl(1/(at)\bigr)=\log(L/a)=H$.  Hence
\[
 ax+b\log x+Q(1/x)-L
 =(V-bH)+b\bigl(H+\log(1+w)\bigr)+Q\!\left(\frac{at}{1+w}\right)
 =V-\Phi(t,V;H)=0 .
\]
For uniqueness, substituting a general $\Pi=\sum\Pi_n(H)t^n$ in place of $V$
turns the equation into $\Pi=\Phi(t,\Pi;H)$, and
\cref{p0:lem:lagrange-fp} has a unique solution.

(3) We show by strong induction that $\deg_HP_n\le n$ and, in addition, that
$\deg_H\gamma_r\le r-1$ for $r\ge2$.  The base is
$\gamma_1=-bH$, of degree $1$.  Assume $\deg_HP_j\le j$ for all $j<n$, so
$\deg_H\gamma_r=\deg_HP_{r-1}\le r-1$ for $2\le r\le n$.  A monomial of
$\OB_{n,k}(\gamma)$ is a product $\gamma_{r_1}\cdots\gamma_{r_k}$ with
$r_1+\dots+r_k=n$ and $r_\nu\ge1$, hence of $H$-degree at most
$\sum_\nu r_\nu=n$; so the first sum of \eqref{p0:eq:P-bell} has $H$-degree at
most $n$.  In the second sum the ordinary Bell polynomial has first index
$n-j$ with $j\ge1$, so its $H$-degree is at most $n-1$.  Hence
$\deg_HP_n\le n$, completing the induction.

A monomial of $H$-degree exactly $n$ requires $\deg_H\gamma_{r_\nu}=r_\nu$ for
every $\nu$, which forces $r_\nu=1$ for all $\nu$ because
$\deg_H\gamma_r\le r-1$ for $r\ge2$.  Thus $k=n$, the monomial is
$\gamma_1^{\,n}=(-bH)^n$, and $\Coef{H^n}\OB_{n,k}(\gamma)=0$ for $k<n$ while
$\OB_{n,n}(\gamma)=(-bH)^n$.  Only the term $k=n$ of the first sum of
\eqref{p0:eq:P-bell} therefore survives, and
\[
 \Coef{H^n}P_n=-b\cdot\frac{(-1)^{n+1}}{n}\cdot(-b)^n
 =\frac{(-1)^{2n+2}\,b^{\,n+1}}{n}=\frac{b^{\,n+1}}{n}.
\]
For the constant term put $H=0$ in \eqref{p0:eq:Phi-direct}:
\[
 \Phi(t,u;0)=-b\log(1+tu)-Q\!\left(\frac{at}{1+tu}\right).
\]
If $U(t)=\sum c_nt^n$ solves $U=\Psi(t,U)$ with $\Psi$ of
\eqref{p0:eq:Psi-general}, put $\widetilde V(t):=a\,U(at)$.  Then
$t\widetilde V(t)=(at)U(at)$, so
\[
 \Phi\bigl(t,\widetilde V(t);0\bigr)
 =-b\log\bigl(1+(at)U(at)\bigr)-Q\!\left(\frac{at}{1+(at)U(at)}\right)
 =a\,\Psi\bigl(at,U(at)\bigr)=a\,U(at)=\widetilde V(t),
\]
so $\widetilde V$ solves the $H=0$ fixed-point equation; by uniqueness
$P_n(0)=\Coef{t^n}\widetilde V=a\cdot a^{n}c_n=a^{n+1}c_n$.

(5) Fix $N$.  Put $x_N$ for the truncation displayed in
\eqref{p0:eq:direct-asymptotic}.  Exactly as in the proof of
\cref{p0:thm:lambert-centered}(5), the residual
$r_N:=\log F(x_N)-\log y$ splits as $g(t)+E_N$ with
$|E_N|\le K_Nx_N^{-N-1}$ and $g$ analytic in $t=1/L$ \emph{with coefficients
polynomial in $H$}; by (4) the Taylor coefficients of $g$ vanish through order
$N$, and the coefficient of $t^{N+1}$ is a polynomial in $H$ of degree at most
$N+1$.  Hence $|r_N|=\bigO\bigl((1+H)^{N+1}L^{-N-1}\bigr)$.  The slope
bound $(\log F)'\ge a/2$ and \cref{p0:thm:backward-error} finish the proof,
after noting $L\asymp aX$ and $H=\log(L/a)$ so that $x_N\to\infty$.
\end{proof}

\begin{theorem}[The pure Lambert block in closed form]\label{p0:thm:pure-stirling}
Let $\stirlingfirst{n}{m}$ be the \emph{unsigned} Stirling numbers of the first
kind, so that the signed ones are
$(-1)^{n-m}\stirlingfirst{n}{m}$ and
$\log(1+z)^m=m!\sum_{n\ge m}(-1)^{n-m}\stirlingfirst{n}{m}z^n/n!$.  When $Q=0$
the polynomials of \cref{p0:thm:flattening} are
\begin{equation}\label{p0:eq:pure-stirling}
 \boxed{\;
 P_n^{(0)}(H)=b^{\,n+1}\sum_{m=1}^{n}(-1)^{m+1}
   \frac{\stirlingfirst{n}{m}}{(n-m+1)!}\,H^{\,n-m+1}\;}
\end{equation}
for $n\ge1$.  Equivalently, clearing the binomial,
\begin{equation}\label{p0:eq:pure-stirling-unsigned}
 P_n^{(0)}(H)=\frac{b^{\,n+1}}{n!}\sum_{m=1}^{n}(-1)^{m+1}(m-1)!\,
  \stirlingfirst{n}{m}\binom{n}{m-1}H^{\,n-m+1}.
\end{equation}
In particular
\begin{equation}\label{p0:eq:pure-lead}
 \Coef{H^n}P_n^{(0)}=\frac{b^{\,n+1}}{n},
 \qquad P_n^{(0)}(0)=0,
\end{equation}
and the first four are
\begin{equation}\label{p0:eq:pure-first}
 P_1^{(0)}=b^2H,\quad
 P_2^{(0)}=\frac{b^3}{2}\bigl(H^2-2H\bigr),\quad
 P_3^{(0)}=\frac{b^4}{6}\bigl(2H^3-9H^2+6H\bigr),\quad
 P_4^{(0)}=\frac{b^5}{12}\bigl(3H^4-22H^3+36H^2-12H\bigr).
\end{equation}
Consequently the Lambert core itself unfolds as
\begin{equation}\label{p0:eq:core-unfolded}
 X=\frac1a\left[L-bH+\sum_{n\ge1}\frac{P_n^{(0)}(H)}{L^n}\right],
 \qquad H=\log\frac La ,
\end{equation}
formally, and with remainder $\bigO\bigl((1+H)^{N+1}L^{-N-1}\bigr)$ after
truncation at $n=N$.
\end{theorem}

\begin{proof}
With $Q=0$, \eqref{p0:eq:P-lagrange} reads
\[
 P_n^{(0)}(H)=\sum_{m=1}^{n}\frac{(-b)^m}{m}
  \Coef{t^nu^{m-1}}\log\bigl(1+t(u-bH)\bigr)^m .
\]
Insert $\log(1+z)^m=m!\sum_{r\ge m}(-1)^{r-m}\stirlingfirst{r}{m}z^r/r!$ with
$z=t(u-bH)$; the coefficient of $t^n$ is
$m!\,(-1)^{n-m}\stirlingfirst{n}{m}(u-bH)^n/n!$, and
$\Coef{u^{m-1}}(u-bH)^n=\binom{n}{m-1}(-bH)^{n-m+1}$.  Hence
\[
 P_n^{(0)}(H)
 =\sum_{m=1}^{n}\frac{(-b)^m}{m}\cdot
   \frac{m!\,(-1)^{n-m}\stirlingfirst{n}{m}}{n!}
   \binom{n}{m-1}(-b)^{n-m+1}H^{n-m+1},
\]
and $(-b)^{m}(-b)^{n-m+1}(-1)^{n-m}=(-1)^{n+1}(-1)^{n-m}b^{n+1}
=(-1)^{m+1}b^{n+1}$, which gives
\eqref{p0:eq:pure-stirling-unsigned}.  Since
$(m-1)!\binom{n}{m-1}/n!=1/(n-m+1)!$, the two displayed forms agree and
\eqref{p0:eq:pure-stirling} follows.  Every term carries $H^{n-m+1}$ with
$m\le n$, so no constant term occurs and $P_n^{(0)}(0)=0$; the term $m=1$ has
$\stirlingfirst{n}{1}=(n-1)!$ and gives
$\Coef{H^n}P_n^{(0)}=b^{n+1}(n-1)!/n!=b^{n+1}/n$.
Finally \eqref{p0:eq:core-unfolded} is \cref{p0:thm:flattening}(4) and (5)
with $Q=0$, in which case the unique solution of $ax+b\log x=L$ is the core
$X$ of \cref{p0:thm:lambert-core}.
\end{proof}

\begin{remark}[Three source formulae, one identity]\label{p0:rem:stirling-crosswalk}
The closed form \eqref{p0:eq:pure-stirling} appears in the sources in three
different dresses, all with the signed Stirling numbers
$s(n,m)=(-1)^{n-m}\stirlingfirst{n}{m}$.
\textsf{Swing\_Factorial\_Transseries\_Article} states
\eqref{p0:eq:pure-stirling-unsigned};
\textsf{Partition\_Number\_Transseries\_and\_Asymptotic\_Inverse} states
$P_j(L)=2^{j+1}\sum_{q\le j}s(j,q)L^{j+1-q}/(j+1-q)!$;
\textsf{Partition\_Number\_Transseries\_and\_Inverse} states
$P_m(\ell)=\frac1{m!}\sum_j2^j(j-1)!\,s(m,j)\binom m{j-1}(2\ell)^{m-j+1}$.
The last two are the specialization $a=1$, $b=-2$ of
\eqref{p0:eq:pure-stirling}, and all three agree with it as exact rational
polynomials for $1\le n\le8$, which was checked in exact rational arithmetic
(\cref{p0:rem:verification}).  The volume uses
\eqref{p0:eq:pure-stirling}: it is the shortest of the four, it avoids the
letter $s$, which \cref{p3:sec:top} needs for Dedekind sums, and it makes the
leading coefficient \eqref{p0:eq:pure-lead} a one-line consequence.
\end{remark}

\begin{remark}[Both normalizations, and what is and is not claimed]
\label{p0:rem:accuracy-claim}
The Lambert-centred expansion \eqref{p0:eq:centered-asymptotic} and the
flattened expansion \eqref{p0:eq:direct-asymptotic} are two truncations of the
same object: by \eqref{p0:eq:core-unfolded}, expanding $X$ and then each power
$X^{-n}$ in the flattened scale returns \eqref{p0:eq:direct-asymptotic}, and
conversely collecting in the flattened series all terms generated by
$b\log x$ alone rebuilds $X$.  What differs is the coefficient ring: the
$c_n$ are constants, the $P_n$ are polynomials of degree $n$ in
$H=\log(L/a)$.  Several of the sources add that Lambert centring is
numerically more accurate at equal truncation order and support this with
tables.  Part~0 makes no such claim: the comparison depends on the size of $H$
relative to $L$ and on the particular data, it is an empirical statement about
specific truncations, and the supporting tables belong to the chapters, where
they are reported with their provenance.  What Part~0 does assert is the
algebraic statement of \cref{p0:rem:no-logs}, and the remainder bounds
\eqref{p0:eq:centered-asymptotic} and \eqref{p0:eq:direct-asymptotic}, whose
right-hand sides differ by the factor $(1+H)^{N+1}$ --- and that factor is the
only rigorous difference between the two at fixed order.
\end{remark}

\section{The three inverses of a discrete sequence}\label{p0:sec:discrete}

\begin{definition}[The three inverse objects]\label{p0:def:three-inverses}
Let $(A_n)_{n\ge n_0}$ be real numbers, strictly increasing for $n\ge n_1\ge n_0$,
with $A_n\to+\infty$.
\begin{enumerate}
\item The \emph{integer staircase}, or generalized inverse, is
      \begin{equation}\label{p0:eq:staircase-def}
       N_\ast(y):=\min\{n\ge n_1:A_n\ge y\},\qquad y\le\sup_nA_n .
      \end{equation}
      It is an integer-valued, nondecreasing step function, constant on each
      interval $(A_{n-1},A_n]$.
\item Given a continuous, strictly increasing
      $\mathcal A:[n_1,\infty)\to\Real$ with $\mathcal A(n)=A_n$ for every
      integer $n\ge n_1$ --- an \emph{admissible interpolation} --- the
      \emph{interpolated inverse} is $\nu_{\mathcal A}:=\mathcal A^{-1}$,
      defined on $[A_{n_1},\infty)$.
\item The \emph{asymptotic inverse} is the formal transseries produced by
      \cref{p0:thm:lambert-centered} or \cref{p0:thm:flattening} from a model
      $F$ for which the chapter in question proves an asymptotic identity
      along the sequence.  It is attached to no particular interpolation.
\end{enumerate}
\end{definition}

\begin{remark}[The three are genuinely different]\label{p0:rem:three-differ}
Object (1) is canonical but not smooth; object (2) is smooth but not canonical
--- distinct admissible interpolations agree at the integers and differ in
between, hence have different inverses; object (3) is canonical as a formal
series but determines a function only after a summation convention is fixed.
The three coincide only in the following precise senses:
$N_\ast=\lceil\nu_{\mathcal A}\rceil$ for \emph{every} admissible
$\mathcal A$ (\cref{p0:thm:staircase}(1)), and
$\nu_{\mathcal A}(A_n)=n$ for every admissible $\mathcal A$ and every integer
$n\ge n_1$.  The asymptotic series (3) is independent of the choice in
$\mathcal A$ at every algebraic order, but the exponentially small layers are
not: the partition articles make this explicit for the Rademacher amplitudes,
which are periodic in the index and therefore genuinely depend on how the
sequence is continued off the lattice.
\end{remark}

\begin{theorem}[Staircase recovery and the separation condition]
\label{p0:thm:staircase}
Notation as in \cref{p0:def:three-inverses}, with $\mathcal A$ admissible and
$\nu:=\nu_{\mathcal A}$.
\begin{enumerate}
\item \emph{Exact rounding.}  For every $y$ in the range of $\mathcal A$,
      that is for $A_{n_1}\le y<\sup_nA_n$,
      \begin{equation}\label{p0:eq:ceil-identity}
       N_\ast(y)=\bigl\lceil\nu(y)\bigr\rceil,
      \end{equation}
      with $\lceil m\rceil=m$ for integers $m$.  The identity holds for every
      admissible interpolation.
\item \emph{Separation condition.}  Let $x_\ast\in\Real$ and $\eta\ge0$ satisfy
      $|x_\ast-\nu(y)|\le\eta$.  If
      \begin{equation}\label{p0:eq:separation}
       \bigl\lceil x_\ast-\eta\bigr\rceil=\bigl\lceil x_\ast+\eta\bigr\rceil,
      \end{equation}
      then $N_\ast(y)$ equals that common value.  Condition
      \eqref{p0:eq:separation} holds if and only if the interval
      $[x_\ast-\eta,x_\ast+\eta]$ is contained in a single half-open interval
      $(m-1,m]$; it can fail for arbitrarily small $\eta>0$, namely when
      $\nu(y)$ is within $\eta$ of an integer from above.
\item \emph{Inverse along the range.}  If $y=A_n$ for an integer $n\ge n_1$ and
      $|x_\ast-\nu(y)|\le\eta<\tfrac12$, then $n=\lfloor x_\ast+\tfrac12\rfloor$.
      This case is unconditional: no separation hypothesis is needed, because
      $\nu(A_n)=n$ exactly.
\item \emph{Residue classes.}  Suppose $r\ge1$ and that for each
      $\rho\in\{0,1,\dots,r-1\}$ the subsequence $(A_n)_{n\equiv\rho\ (r)}$ is
      strictly increasing for $n\ge n_1$ and admits an admissible
      interpolation $\mathcal A_\rho$ on the class, with inverse $\nu_\rho$.
      Then, for $y$ large enough,
      \begin{equation}\label{p0:eq:residue-staircase}
       N_\ast(y)=\min_{0\le\rho<r}
        \left(\rho+r\left\lceil\frac{\nu_\rho(y)-\rho}{r}\right\rceil\right).
      \end{equation}
\item \emph{No uniform smooth representation.}  For every $Y$ in the range and
      every continuous $g:[Y,\infty)\to\Real$,
      \begin{equation}\label{p0:eq:no-smooth}
       \sup_{y\ge Y}\bigl|N_\ast(y)-g(y)\bigr|\ \ge\ \tfrac12 .
      \end{equation}
      Equivalently, the $\bigO(1)$ periodic layer
      $\lceil x\rceil-x-\tfrac12=\pi^{-1}\sum_{m\ge1}m^{-1}\sin(2\pi mx)$
      (valid for $x\notin\Int$) is not a refinement that can be dropped: it is
      the leading discrete transmonomial of the generalized inverse.
\end{enumerate}
\end{theorem}

\begin{proof}
(1) For an integer $n\ge n_1$, $A_n\ge y$ is equivalent to
$\mathcal A(n)\ge y$, hence to $n\ge\nu(y)$ by strict monotonicity of
$\mathcal A$.  The least such integer is $\lceil\nu(y)\rceil$.

(2) The function $\lceil\cdot\rceil$ is constant on each $(m-1,m]$ and
nondecreasing; since $\nu(y)\in[x_\ast-\eta,x_\ast+\eta]$,
\eqref{p0:eq:separation} forces $\lceil\nu(y)\rceil$ to equal the common
value, and (1) applies.  The characterization of
\eqref{p0:eq:separation} and the failure mode are immediate from the same
description of $\lceil\cdot\rceil$.

(3) $\nu(A_n)=n$ because $\mathcal A(n)=A_n$; hence $|x_\ast-n|\le\eta<1/2$ and
rounding to the nearest integer returns $n$.

(4) Fix $\rho$.  Among the integers $n\equiv\rho\pmod r$ with $n\ge n_1$, the
condition $A_n\ge y$ is, by (1) applied to the class, equivalent to
$n\ge\nu_\rho(y)$.  Writing $n=\rho+rk$ with $k\in\Int$, the least admissible
$k$ is $\lceil(\nu_\rho(y)-\rho)/r\rceil$, so the least admissible $n$ in the
class is $\rho+r\lceil(\nu_\rho(y)-\rho)/r\rceil$.  Taking the minimum over
$\rho$ gives the least admissible $n$ without a congruence restriction, which
is $N_\ast(y)$.

(5) Fix $n$ with $A_n>Y$.  On $(A_{n-1},A_n]$ we have $N_\ast=n$ and on
$(A_n,A_{n+1}]$ we have $N_\ast=n+1$.  For $\epsilon>0$ small,
\[
 1=\bigl|N_\ast(A_n+\epsilon)-N_\ast(A_n)\bigr|
 \le\bigl|N_\ast(A_n+\epsilon)-g(A_n+\epsilon)\bigr|
  +\bigl|g(A_n+\epsilon)-g(A_n)\bigr|
  +\bigl|g(A_n)-N_\ast(A_n)\bigr| .
\]
Letting $\epsilon\downarrow0$ and using continuity of $g$ gives
$1\le2\sup_{y\ge Y}|N_\ast-g|$.  For the Fourier identity, the function
$x\mapsto\lceil x\rceil-x-\tfrac12$ is $1$-periodic, equal to
$-(x-\lfloor x\rfloor)+\tfrac12$ shifted, i.e.\ minus the standard sawtooth,
whose Fourier series is the displayed one and converges pointwise off
$\Int$.
\end{proof}

\begin{proposition}[Certified discrete inversion]\label{p0:prop:certified-rounding}
In the setting of \cref{p0:thm:staircase}(4), suppose that for each $\rho$ an
approximation $x_\rho$ with $|x_\rho-\nu_\rho(y)|\le\eta_\rho$ is available.
Then $N_\ast(y)$ is determined by at most
$\sum_\rho\bigl(\lfloor2\eta_\rho/r\rfloor+2\bigr)$ exact evaluations of the
sequence: for each $\rho$ list the integers $n\equiv\rho\pmod r$ in
$[x_\rho-\eta_\rho,\,x_\rho+\eta_\rho+r]$, evaluate $A_n$, and take the least
$n$ over all classes with $A_n\ge y$.
\end{proposition}

\begin{proof}
By \cref{p0:thm:staircase}(4) the class-$\rho$ candidate is
$\rho+r\lceil(\nu_\rho(y)-\rho)/r\rceil$, and $\nu_\rho(y)$ lies in
$[x_\rho-\eta_\rho,x_\rho+\eta_\rho]$; the least integer of the class that is
$\ge\nu_\rho(y)$ therefore lies in $[x_\rho-\eta_\rho,x_\rho+\eta_\rho+r]$.
Since the class subsequence is increasing, the test $A_n\ge y$ identifies it,
and the outer minimum over $\rho$ is \eqref{p0:eq:residue-staircase}.
\end{proof}

\begin{remark}[The parity instances]\label{p0:rem:parity-instances}
\Cref{p4:sec:top} and \cref{p2:sec:top} are the case $r=2$.  The formula
$\rho+r\lceil(\nu_\rho-\rho)/r\rceil$ specializes to
$2\lceil\nu_0/2\rceil$ and $2\lceil(\nu_1-1)/2\rceil+1$, which is exactly the
pair of parity candidates of \textsf{double\_factorial\_transseries-2}, and
$N_\ast=\min$ of the two is its ``exact generalized inverse'' theorem.
\textsf{Swing\_Factorial\_Transseries\_Article} states the corresponding
\emph{floor} formulae for the threshold query ``the greatest index in the class
$\rho$ whose sequence value is at most $y$''; those are
$\rho+r\lfloor(\nu_\rho-\rho)/r\rfloor$, obtained from
\eqref{p0:eq:residue-staircase} by replacing $\lceil\cdot\rceil$ with
$\lfloor\cdot\rfloor$ and $\min$ with $\max$, with the same proof.  The
partition chapter is the case $r=1$, where
\eqref{p0:eq:residue-staircase} reduces to \eqref{p0:eq:ceil-identity}.
\end{remark}

\begin{repairbox}
\textbf{Source claim.}  Several sources write, or invite the reader to write,
$N_\ast(y)=\nu(y)+\littleo(1)$.
\textbf{Defect.}  This is false uniformly in $y$, by
\cref{p0:thm:staircase}(5).  \textsf{Butcher\_Tree\_Transseries-2} already
flags it (``An expansion such as $N_\ast(y)=\nu(y)+\littleo(1)$ is
false uniformly in $y$''), and
\textsf{Partition\_Number\_Transseries\_and\_Asymptotic\_Inverse} likewise
separates the three objects; but the qualification is absent from the summary
statements of the other articles.
\textbf{Correction.}  The volume uses only the three statements
\cref{p0:thm:staircase}(1), (3), (4): an exact ceiling identity valid for every
admissible interpolation, an unconditional nearest-integer recovery
\emph{along the range} $y=A_n$, and a residue-class refinement.  Nothing else
about $N_\ast$ is asserted.
\end{repairbox}

\section{Error control, transport, and optimal truncation}
\label{p0:sec:errors}

\subsection{Backward error certifies forward error}

\begin{theorem}[Residual-to-root conversion]\label{p0:thm:backward-error}
Let $I\subseteq\Real$ be an interval and $f:I\to\Real$ differentiable with
\begin{equation}\label{p0:eq:slope-bounds}
 0<m\le f'(\xi)\le M\qquad(\xi\in I).
\end{equation}
Let $x_\dagger\in I$ satisfy $f(x_\dagger)=\lambda$, let $x_\ast\in I$, and set
the \emph{residual} $r_\ast:=f(x_\ast)-\lambda$.  Then
\begin{equation}\label{p0:eq:backward-two-sided}
 \frac{|r_\ast|}{M}\ \le\ |x_\ast-x_\dagger|\ \le\ \frac{|r_\ast|}{m} .
\end{equation}
Moreover, if $f$ is continuous on $I$ and
$[x_\ast-\varrho,\,x_\ast+\varrho]\subseteq I$ with $\varrho:=|r_\ast|/m$,
then $f$ takes the value $\lambda$ somewhere in that interval; the residual is
therefore a \emph{certificate}, not merely a bound conditional on the
existence of a nearby root.
\end{theorem}

\begin{proof}
By the mean value theorem, $r_\ast=f(x_\ast)-f(x_\dagger)=f'(\xi)(x_\ast-x_\dagger)$
for some $\xi$ between the two points, and \eqref{p0:eq:slope-bounds} gives
\eqref{p0:eq:backward-two-sided}.  For the certificate, suppose $r_\ast>0$
(the case $r_\ast<0$ is symmetric and $r_\ast=0$ is trivial).  Then
$f(x_\ast)=\lambda+r_\ast>\lambda$, while
\[
 f(x_\ast-\varrho)\le f(x_\ast)-m\varrho=\lambda+r_\ast-|r_\ast|=\lambda,
\]
using $f'\ge m$.  The intermediate value theorem supplies a point of
$[x_\ast-\varrho,x_\ast]$ where $f=\lambda$.
\end{proof}

\begin{remark}[Use in the logarithmic variable]\label{p0:rem:log-residual}
In every application $f=\log F$ and $\lambda=\log y$: $F$ itself overflows any
fixed-precision arithmetic long before the asymptotic regime begins, whereas
$\log F$ is computable from a digamma or Laplace representation throughout.
The lower bound in \eqref{p0:eq:backward-two-sided} makes the residual test
sharp --- a small residual is necessary as well as sufficient.
\end{remark}

\subsection{Transport of a controlled forward remainder}

\begin{theorem}[Remainder transport]\label{p0:thm:remainder-transport}
Let $I=[x_1,\infty)$, let $f_0:I\to\Real$ be differentiable with
$f_0'\ge m>0$, and let $\varepsilon:I\to\Real$ be differentiable with
$\sup_I|\varepsilon'|\le\theta<m$.  Put $f:=f_0+\varepsilon$.  For $\lambda$ in
the common range let $x_0=x_0(\lambda)\in I$ solve $f_0(x_0)=\lambda$ and
$x=x(\lambda)\in I$ solve $f(x)=\lambda$.  Then
\begin{enumerate}
\item \emph{Rigorous bound.}
      \begin{equation}\label{p0:eq:transport-bound}
       \bigl|x(\lambda)-x_0(\lambda)\bigr|
       \le\frac{\bigl|\varepsilon\bigl(x_0(\lambda)\bigr)\bigr|}{m-\theta}.
      \end{equation}
\item \emph{First-order law with explicit error.}  If moreover $f_0\in C^2(I)$
      with $\sup_I|f_0''|\le M_2$, then
      \begin{equation}\label{p0:eq:transport-first-order}
       x(\lambda)-x_0(\lambda)
       =-\frac{\varepsilon(x_0)}{f_0'(x_0)}+E,
       \qquad
       |E|\le\frac{|\varepsilon(x_0)|}{m\,(m-\theta)}
        \left(\frac{M_2\,|\varepsilon(x_0)|}{m-\theta}+\theta\right).
      \end{equation}
\end{enumerate}
In particular, if $\varepsilon(x)=\bigO(x^{-M-1})$ and
$\varepsilon'(x)=\littleo(1)$ as $x\to\infty$, then
$x(\lambda)-x_0(\lambda)=-\varepsilon(x_0)/f_0'(x_0)\,(1+\littleo(1))
=\bigO(x_0^{-M-1})$.
\end{theorem}

\begin{proof}
Since $f'=f_0'+\varepsilon'\ge m-\theta>0$, $f$ is strictly increasing and
$x(\lambda)$ is unique.  The residual of $x_0$ in $f$ is
$f(x_0)-\lambda=f_0(x_0)+\varepsilon(x_0)-\lambda=\varepsilon(x_0)$, so
\cref{p0:thm:backward-error} applied to $f$ with $m$ replaced by $m-\theta$
gives \eqref{p0:eq:transport-bound}.

For (2), write $d:=x-x_0$.  By the mean value theorem there is $\xi$ between
$x_0$ and $x$ with $0=f(x)-\lambda=\varepsilon(x_0)+f'(\xi)d$, so
$d=-\varepsilon(x_0)/f'(\xi)$ and
\[
 d+\frac{\varepsilon(x_0)}{f_0'(x_0)}
 =\varepsilon(x_0)\left(\frac1{f_0'(x_0)}-\frac1{f'(\xi)}\right)
 =\varepsilon(x_0)\,\frac{f'(\xi)-f_0'(x_0)}{f_0'(x_0)\,f'(\xi)} .
\]
Now $|f'(\xi)-f_0'(x_0)|\le|f_0'(\xi)-f_0'(x_0)|+|\varepsilon'(\xi)|
\le M_2|d|+\theta\le M_2|\varepsilon(x_0)|/(m-\theta)+\theta$ by
\eqref{p0:eq:transport-bound}, while $f_0'(x_0)\ge m$ and
$f'(\xi)\ge m-\theta$.  Substituting gives
\eqref{p0:eq:transport-first-order}.  The final assertion follows since then
$\theta\to0$ and $|\varepsilon(x_0)|\to0$.
\end{proof}

\begin{corollary}[Forward truncation to inverse accuracy]
\label{p0:cor:forward-to-inverse}
Let $F$ be differentiably of exponential--power type with data
$(C,a,1,b,1,Q)$, strictly increasing on $[x_1,\infty)$ with
$(\log F)'\ge m>0$ there.  Let $x_\ast$ solve exactly the truncated equation
\[
 \log C+ax_\ast+b\log x_\ast+\sum_{j=1}^{M}q_jx_\ast^{-j}=\log y .
\]
Then $|x_\ast-F^{-1}(y)|\le K_Mx_\ast^{-M-1}/m$ with $K_M$ the constant of
\eqref{p0:eq:model-analytic}.
\end{corollary}

\begin{proof}
The residual of $x_\ast$ in $\log F$ is exactly the forward remainder, of
modulus at most $K_Mx_\ast^{-M-1}$; apply
\cref{p0:thm:backward-error}.
\end{proof}

\subsection{Least-term truncation and the beyond-all-orders floor}

\begin{theorem}[Optimal truncation and the exponential floor]
\label{p0:thm:optimal-truncation}
Let $A>0$, $\beta\in\Real$, $K>0$, $x_0\ge1$, and let
$\varphi:[x_0,\infty)\to\Real$ satisfy the \emph{Gevrey-$1$ remainder bound}
\begin{equation}\label{p0:eq:gevrey-remainder}
 \Bigl|\varphi(x)-\sum_{j=1}^{M}q_jx^{-j}\Bigr|
 \ \le\ K\,\Gamma(M+\beta)\,A^{-M}x^{-M}
 \qquad(M\ge1,\ x\ge x_0).
\end{equation}
\begin{enumerate}
\item \emph{Least-term envelope.}  There are $x_1\ge x_0$ and
      $K'=K'(K,A,\beta)$ such that, choosing $M=M(x):=\lfloor Ax\rfloor$,
      \begin{equation}\label{p0:eq:least-term}
       \Bigl|\varphi(x)-\sum_{j=1}^{M(x)}q_jx^{-j}\Bigr|
       \ \le\ K'\,x^{\beta-1/2}\,\e^{-Ax}
       \qquad(x\ge x_1).
      \end{equation}
\item \emph{Transported inverse error.}  Let $F$ be as in
      \cref{p0:cor:forward-to-inverse}, with
      $\log F(x)=\log C+ax+b\log x+\varphi(x)$ and $(\log F)'\ge m>0$ on
      $[x_1,\infty)$.  If $x_\ast$ solves the optimally truncated equation
      exactly, then
      \begin{equation}\label{p0:eq:optimal-inverse}
       \bigl|x_\ast-F^{-1}(y)\bigr|
       \ \le\ \frac{K'}{m}\,x_\ast^{\beta-1/2}\,\e^{-Ax_\ast}.
      \end{equation}
\item \emph{The floor is real.}  Suppose in addition that there are
      $M_0\ge1$ and $0<k\le K'' $ with
      \begin{equation}\label{p0:eq:two-sided-coefficients}
       k\,\Gamma(M+\beta)A^{-M}\ \le\ |q_M|\ \le\ K''\,\Gamma(M+\beta)A^{-M}
       \qquad(M\ge M_0).
      \end{equation}
      Then
      \begin{equation}\label{p0:eq:floor}
       \log\ \min_{M\ge M_0}\ \bigl|q_M\bigr|x^{-M}
       \ =\ -Ax+\bigO(\log x)
       \qquad(x\to\infty),
      \end{equation}
      so that \emph{no} truncation of the algebraic series achieves an error
      of smaller exponential order than $\e^{-Ax}$.
\end{enumerate}
\end{theorem}

\begin{proof}
(1) Stirling's formula gives constants $C_\beta$ and $M_\beta\ge1$ with
$\Gamma(M+\beta)\le C_\beta M^{M+\beta-1/2}\e^{-M}$ for all $M\ge M_\beta$.
With $M=\lfloor Ax\rfloor$ and $x$ so large that $M\ge M_\beta$,
\[
 \Gamma(M+\beta)A^{-M}x^{-M}
 \le C_\beta M^{\beta-1/2}\e^{-M}\left(\frac{M}{Ax}\right)^{M}
 \le C_\beta M^{\beta-1/2}\e^{-M},
\]
because $M\le Ax$.  Since $M>Ax-1$, we have $\e^{-M}<\e\cdot\e^{-Ax}$, and
$M^{\beta-1/2}\le C_\beta'x^{\beta-1/2}$ for $x\ge x_1$ (with
$C_\beta'$ depending only on $A$ and $\beta$, since
$Ax-1\le M\le Ax$).  Combining with \eqref{p0:eq:gevrey-remainder} gives
\eqref{p0:eq:least-term}.

(2) The residual of $x_\ast$ in $\log F$ is exactly the left-hand side of
\eqref{p0:eq:least-term} evaluated at $x_\ast$; apply
\cref{p0:thm:backward-error}.

(3) Enlarging $M_0$ we may assume $M_0\ge M_\beta$.  Write $u:=Ax$ and
$g(M):=(M/u)^{M}\e^{-M}$, so that by Stirling there are $0<c_\beta\le C_\beta$
with
\begin{equation}\label{p0:eq:stirling-sandwich}
 c_\beta M^{\beta-1/2}g(M)\ \le\ \Gamma(M+\beta)(Ax)^{-M}
 \ \le\ C_\beta M^{\beta-1/2}g(M)
 \qquad(M\ge M_\beta).
\end{equation}

\emph{Upper bound.}  Take $M=\lfloor Ax\rfloor$ in the right-hand inequality of
\eqref{p0:eq:two-sided-coefficients} and repeat the estimate of part~(1); this
gives $\min_{M\ge M_0}|q_M|x^{-M}\le K'''x^{\beta-1/2}\e^{-Ax}$ once
$\lfloor Ax\rfloor\ge M_0$.

\emph{Lower bound.}  Note first that $\log g(M)=M\log(M/u)-M$ has derivative
$\log(M/u)$, so $g$ decreases on $(0,u)$, increases on $(u,\infty)$, and
\begin{equation}\label{p0:eq:g-min}
 g(M)\ \ge\ g(u)=\e^{-u}=\e^{-Ax}\qquad(M>0).
\end{equation}
Set $h(M):=M^{\beta-1/2}g(M)$; then
$(\log h)'(M)=(\beta-\tfrac12)/M+\log(M/u)$, which is positive as soon as
$M\ge2u$ and $2u\ge2|\beta-\tfrac12|/\log2$, so $h$ is increasing on
$[2u,\infty)$ for all large $x$.  Consequently
\[
 \min_{M\ge M_0}h(M)
 =\min_{M_0\le M\le2u}h(M)
 \ \ge\ \Bigl(\min_{M_0\le M\le2u}M^{\beta-1/2}\Bigr)\e^{-Ax}
 \ \ge\ c\,x^{-|\beta-1/2|}\,\e^{-Ax}
\]
for $x\ge x_1$, using \eqref{p0:eq:g-min} and
$\min_{M_0\le M\le2u}M^{\beta-1/2}\ge\min\bigl(M_0^{\beta-1/2},(2Ax)^{\beta-1/2}\bigr)$.
Combining with the left-hand inequalities of
\eqref{p0:eq:two-sided-coefficients} and \eqref{p0:eq:stirling-sandwich},
$|q_M|x^{-M}\ge k\,c_\beta\,c\,x^{-|\beta-1/2|}\e^{-Ax}$ for every
$M\ge M_0$.  Taking logarithms of the two bounds gives
\eqref{p0:eq:floor}.
\end{proof}

\begin{remark}[What may and may not be appended]\label{p0:rem:no-bare-exponential}
\Cref{p0:thm:optimal-truncation} bounds the size of the optimal remainder; it
does not produce it.  In particular, adding a term $c\,\e^{-Ax}$ with an
arbitrarily chosen constant $c$ to the divergent series is \emph{not} a valid
transseries completion.  The double factorial and swing sources make the point
explicitly: on the positive real axis the Borel singularities of these phases
lie at $\pm\ii\pi m$, so the exponentially small sectors are
terminant-weighted combinations of $\e^{\mp\ii\pi mx}$, whose optimally
truncated magnitude is $\e^{-\pi x}$ but which are not real exponentials.  The
unambiguous completion on the positive axis is the exact Borel--Laplace
representation, when the chapter provides one.  Where a chapter does supply a
rigorously normalized exponential sector, its transport through the inverse map
is \cref{p0:prop:operator-form} at first order and
\cref{p0:thm:remainder-transport} with an explicit error.
\end{remark}

\section{The parameter dictionary}\label{p0:sec:dictionary}

\begin{table}[htbp]
\centering
\caption{Model data of the four chapters, in the normal form of
\cref{p0:def:model} and after the $\alpha$-reduction of
\cref{p0:lem:alpha-reduction}.  ``Core'' names the admissible core of
\cref{p0:def:core} and the Lambert branch selected by
\cref{p0:cor:branch-rule}.  The row for \cref{p2:sec:top} is
\emph{not} an instance of \cref{p0:def:model}; see the repair box in
\cref{p0:sec:model}.}
\label{p0:tab:dictionary}
\begin{tabular}{llcccl}
\toprule
Chapter & sequence & variable $x$ & $\alpha$ & $b$ & core, branch\\
\midrule
\ref{p1:sec:top} & rooted trees, A000081 & $n$
 & $1$ & $-3/2$ & $ax+b\log x$, $\Wlam_{-1}$\\
\ref{p2:sec:top} & double factorial, A006882 & $n+1$
 & --- & --- & $\tfrac x2(\log x-1)$, $\Wlam_{0}$\\
\ref{p3:sec:top} & partitions, A000041 & $n-\tfrac1{24}$
 & $1/2$ & $-1$ & $ax^{1/2}+b\log x$, $\Wlam_{-1}$\\
\ref{p4:sec:top} & swing factorial, A056040 & $n$
 & $1$ & $\sigma/2$ & $ax+b\log x$, $\Wlam_{0}$ or $\Wlam_{-1}$\\
\bottomrule
\end{tabular}
\end{table}

The constants $a$ are: $a=\log\alpha_{\mathrm{Otter}}$ for
\cref{p1:sec:top}; $a=\pi\sqrt{2/3}$ for \cref{p3:sec:top}, in which
$\alpha=1/2$ and \cref{p0:lem:alpha-reduction} substitutes
$\xi=(n-1/24)^{1/2}$, after which $b/\alpha=-2$ and the core equation is
$\xi\mapsto a\xi-2\log\xi$ --- the equation the three partition articles write
as $\rho-2\log\rho=L$ after rescaling $a$ to $1$; and $a=\log2$ for
\cref{p4:sec:top}, with $\sigma=(-1)^{n+1}$ the parity transmonomial.  Each
chapter verifies \eqref{p0:eq:model-derivative} directly and states its own
$q_j$.

\begin{remark}[Reproducibility of the Part~0 identities]\label{p0:rem:verification}
Every coefficient identity of this chapter was checked in exact rational
arithmetic, with no floating-point step, over six independent choices of the
data $(a,b,q_1,\dots,q_N)$ with $a,b,q_j\in\Rat$, $N\in\{10,12\}$, using a
truncated-series engine over $\Rat[H]$ built on Python's
\texttt{fractions.Fraction}.  The identities checked were:
the Lagrange closed formula \eqref{p0:eq:c-lagrange} against the defining
fixed point; the ordinary-Bell recurrence \eqref{p0:eq:c-bell}; the vanishing
of the residual $ax+b\log x+Q(1/x)-L$ under the flattened ansatz
\eqref{p0:eq:direct-ansatz} through order $L^{-N}$; the shape statements
\eqref{p0:eq:P-shape}; and the closed form
\eqref{p0:eq:pure-stirling} together with its agreement with the three source
formulae listed in \cref{p0:rem:stirling-crosswalk}.  All checks passed.  No
numerical constant is asserted anywhere in this chapter.
\end{remark}

%% ==================== fragment p1 ====================
\chapter{Concordance: the calculus and the inversion apparatus}
\label{q3:sec:top}

\section{Why this chapter exists}
\label{q3:sec:why}

The calculus parts of this volume and the inversion apparatus of
\cref{p0:sec:top} were written independently, as separate consolidations of
separate groups of articles, and the group they came from recorded as an open
question whether their machinery coincides.  This chapter answers it, pair by
pair.

The answer is more interesting than either ``yes'' or ``no''.  A crude
comparison --- matching the two apparatuses by the titles of their results ---
suggests heavy duplication: Lagrange--B\"urmann appears on both sides, so do
Bell polynomials, so does residual-to-error transfer, so does a closed-form
Lambert block.  On inspection, most of those pairs are \emph{not} the same
statement.  They are different theorems about the same subject, and the volume
needs both.  Exactly one pair is a genuine duplicate, and one pair is a strict
strengthening.  \Cref{q3:tab:concordance} records the verdicts and the rest of
this chapter justifies them.

\begin{table}[htbp]
\centering
\caption{Apparent overlaps between the calculus parts and
\cref{p0:sec:top}, with the verdict for each.}
\label{q3:tab:concordance}
\small
\setlength{\tabcolsep}{4pt}
\begin{tabular}{@{}p{0.30\textwidth}p{0.30\textwidth}p{0.32\textwidth}@{}}
\toprule
Calculus & Inversion apparatus & Verdict\\
\midrule
\cref{plt:def:bell-partial} exponential partial and complete Bell polynomials
 & \cref{p0:def:ebell} exponential partial Bell polynomials
 & \textbf{same definition}; see \cref{q3:prop:bell-same}\\
\addlinespace
--- (no counterpart)
 & \cref{p0:def:obell} \emph{ordinary} partial Bell polynomials
 & \textbf{new}; the calculus defines only the exponential family\\
\addlinespace
\cref{plt:prop:rev-analytic-transfer} analytic residual transfer
 & \cref{p0:thm:backward-error} residual-to-root conversion
 & \textbf{strengthening}; see \cref{q3:prop:transfer}\\
\addlinespace
\cref{plt:thm:lag-burmann} and \cref{plt:thm:du-lagrange-zero},
 near-identity operator form
 & \cref{p0:thm:LB}, classical coefficient form
 & \textbf{different statements}; \cref{q3:rem:lb}\\
\addlinespace
\cref{plt:thm:lw-affine-log} reversion of $X+aL$ by Lambert polynomials
 & \cref{p0:thm:pure-stirling} the pure Lambert block in closed form
 & \textbf{different objects}; \cref{q3:rem:lambert}\\
\addlinespace
\cref{plt:prop:lead-slope} slope transport
 & \cref{p0:thm:remainder-transport} remainder transport
 & \textbf{different}; one is a leading-order rule, the other a bound with an
   explicit second-order error\\
\bottomrule
\end{tabular}
\end{table}

\section{The one genuine duplicate}
\label{q3:sec:duplicate}

\begin{proposition}[The two Bell definitions coincide]\label{q3:prop:bell-same}
For all $n\ge k\ge0$ and all indeterminates $x_1,x_2,\dots$,
\begin{equation}\label{q3:eq:bell-same}
 \EB_{n,k}(x_1,x_2,\dots)
 =\ExponentialPartialBellPolynomial{n}{k}(x_1,x_2,\dots),
\end{equation}
and the two complete polynomials agree as well.  That is,
\cref{p0:def:ebell} and \cref{plt:def:bell-partial} define the same object in
different notation.
\end{proposition}

\begin{proof}
\Cref{plt:eq:bell-partial-def} reads
\[
 \ExponentialPartialBellPolynomial{n}{k}
 =\sum\frac{n!}{\prod_jm_j!\,(j!)^{m_j}}\prod_jx_j^{m_j},
\]
summed over $(m_j)$ with $\sum_jm_j=k$ and $\sum_jjm_j=n$, while
\eqref{p0:eq:ebell-def} reads
\[
 \EB_{n,k}=n!\sum\prod_j\frac1{m_j!}\Bigl(\frac{x_j}{j!}\Bigr)^{m_j}
 =\sum\frac{n!}{\prod_jm_j!\,(j!)^{m_j}}\prod_jx_j^{m_j}
\]
over the same index set.  The summands are identical term by term.  Both
sources then define the complete polynomial as $\sum_{k=0}^{n}$ of the
partial one, so those agree too.
\end{proof}

\begin{remark}[What is done about it]\label{q3:rem:bell-action}
Nothing is deleted.  \Cref{p0:def:ebell} additionally records the two
generating identities \eqref{p0:eq:ebell-egf-partial} and
\eqref{p0:eq:ebell-egf-complete}, which the later chapters use constantly and
which \cref{plt:def:bell-partial} does not state; and the calculus parts use
the long normative names throughout, while the inversion chapters use the
short $\EB$.  Collapsing one into the other would cost more in readability
than the duplication costs in length.  What \cref{q3:prop:bell-same} buys is
that a reader who meets both knows they are the same object and may use a
result about either.
\end{remark}

\section{The one strengthening}
\label{q3:sec:strengthening}

\begin{proposition}[The inversion apparatus proves more]\label{q3:prop:transfer}
\Cref{plt:prop:rev-analytic-transfer} is the implication
\begin{equation}\label{q3:eq:one-sided}
 |f'|\ge m>0\ \text{between }x_N\text{ and }x_*
 \ \Longrightarrow\
 |x_N-x_*|\le\frac{|f(x_N)-y|}{m},
\end{equation}
and is exactly the upper bound of \cref{p0:thm:backward-error}.  The latter
proves two further things: the matching \emph{lower} bound
$|r_*|/M\le|x_*-x_\dagger|$ under a two-sided slope bracket
$m\le f'\le M$, and that the residual is a \emph{certificate} --- if
$[x_*-\varrho,x_*+\varrho]$ lies in the interval, with
$\varrho=|r_*|/m$, then a root exists within that distance, so no separate
existence argument is needed.  Hence
\cref{plt:prop:rev-analytic-transfer} is a corollary of
\cref{p0:thm:backward-error}, and not conversely.
\end{proposition}

\begin{proof}
For the first sentence, take $M=\infty$ in \cref{p0:thm:backward-error}, or
read its proof, which is the same mean-value argument.  The lower bound and
the existence clause are proved there and are not stated in
\cref{plt:prop:rev-analytic-transfer}; that the converse fails is immediate,
since \eqref{q3:eq:one-sided} says nothing when $f'$ is unbounded above and
asserts no existence.
\end{proof}

\begin{remark}[Where the difference is used]\label{q3:rem:transfer-use}
The lower bound is what makes a residual a two-sided measurement rather than
a one-sided guarantee, and \cref{q1:prop:certificate} uses exactly that: for
$x+\LambertW(x)$ the slope bracket is the constant pair $(1,2)$, so the
residual determines the error to within a factor of two, uniformly in $z$.
The existence clause is used in the proofs of
\cref{p6:thm:gamma,p6:thm:barnes}, where it removes the obligation to prove
separately that the truncated expansion is near an actual root.
\end{remark}

\section{The pairs that only look alike}
\label{q3:sec:different}

\begin{remark}[Lagrange--B\"urmann, twice, differently]\label{q3:rem:lb}
The calculus proves the \emph{operator} form: for a near-identity map
$F=X+A$ in a reversion frame,
\[
 \rev F=X+\sum_{r\ge1}\frac{(-1)^{r}}{r!}\partial^{\,r-1}\bigl(A^{r}\bigr),
\]
at infinity in \cref{plt:thm:lag-burmann} and at a finite point in
\cref{plt:thm:du-lagrange-zero}, in both cases with a summability statement
that is the real content.  \Cref{p0:thm:LB} proves the \emph{coefficient}
form: if $W=t\varphi(W)$ with $\varphi(0)$ invertible then
$\Coef{t^{n}}G(W)=n^{-1}\Coef{w^{n-1}}(G'\varphi^{n})$.

These are the two classical faces of the same theorem, and each is the natural
tool for a different job: the operator form reverses a perturbation of the
identity, which is what the calculus parts need; the coefficient form
extracts a single coefficient of a composite, which is what the reversion
recurrences of the application chapters need.  Deriving either from the other
is a genuine argument, not a rewriting, so neither is redundant and neither
proof is removed.  What \emph{is} avoided is a third copy: the application
chapters cite \cref{p0:thm:LB} and the $x+\LambertW(x)$ chapter cites the
calculus, rather than either re-deriving.
\end{remark}

\begin{remark}[Affine-logarithmic reversion is not the Lambert block]
\label{q3:rem:lambert}
\Cref{plt:thm:lw-affine-log} reverses $X+aL$, a perturbation of the identity
by a multiple of the logarithm, and its answer is a Lambert-polynomial series
$X-aL+\sum_{n\ge1}a^{n+1}\LamP{n}(L)t^{n}$.  \Cref{p0:thm:pure-stirling}
evaluates the pure Lambert block of the exponential--power model in closed
form.  The first is a statement about the near-identity group; the second is a
closed-form evaluation of a specific reversion.  They meet only in that both
involve $\LambertW$, and neither implies the other.
\end{remark}

\section{What this comparison changes}
\label{q3:sec:verdict}

\begin{remark}[A correction to a claim made earlier in this merge]
\label{q3:rem:correction}
An earlier stage of this consolidation, working from the title-level
concordance alone, recorded that the calculus ``already contains''
Lagrange--B\"urmann at infinity and at a finite point, residual-to-error
transfer, both Bell families, affine-logarithmic reversion and the pure
Lambert block, and that the inversion apparatus was therefore largely
redundant.  That overstated the overlap, and this chapter is the correction.
The accurate statement is the one in \cref{q3:tab:concordance}: of six
apparent overlaps, one is a genuine duplicate of a \emph{definition}, one is a
strengthening in the inversion apparatus's favour, one item exists only in the
inversion apparatus, and the remaining three are distinct theorems that the
volume needs separately.

The moral is worth keeping, because the same trap recurs whenever two
independently written treatments are compared: shared vocabulary is a weak
signal.  Two results called ``Lagrange--B\"urmann'' may be different theorems,
and a mechanical concordance will report them as the same.  Only reading the
statements settles it.
\end{remark}

\begin{remark}[What the inversion apparatus really adds]
\label{q3:rem:what-adds}
Setting the false positives aside, the apparatus of \cref{p0:sec:top}
contributes, over and above the calculus: the exponential--power model and its
axiomatized dominant core; the monomial $\alpha$-reduction that brings a
fractional phase into the linear case; perturbed inversion around an exactly
invertible core, in both an analytic and a formal-jet form; the ordinary
partial Bell polynomials and their conversion to the exponential family; the
two-sided backward-error certificate; and --- with no analogue anywhere in the
calculus --- the theory of inverting a \emph{sequence}: the three distinct
inverse objects, the staircase theorem relating them, and the separation
condition under which an asymptotic inverse determines an integer one.  The
calculus is a theory of functions on a scale; that last group of results is
about the passage from a function to a sequence, and it is the reason the two
halves of this volume are complementary rather than redundant.
\end{remark}


\part{One map in depth: reversing \(x+\LambertW(x)\)}

\chapter{Reversing \texorpdfstring{\(x+\LambertW(x)\)}{x + W(x)}}
\label{q1:sec:top}

\section{What this chapter does}
\label{q1:sec:intro}

The map
\begin{equation}\label{q1:eq:f-def}
 f(x)=x+\LambertW(x),\qquad x\ge0,
\end{equation}
is the smallest interesting perturbation of the identity by a logarithm.  It
is strictly increasing from $[0,\infty)$ onto itself, so its inverse
$g=f^{-1}$ is globally defined; and because
$\LambertW(x)=\bigO(\log x)=\littleo(x)$, reversing it looks like a routine
near-identity problem of the kind the reversal part of this volume settles in
general.

It is worth a chapter of its own for one reason.  \emph{Three} scales occur,
and the middle one hides the interesting behaviour:
\begin{enumerate}
\item the leading scale $z$;
\item the slow logarithmic scale generated by $\ell_1=\log z$ and
      $\ell_2=\log\ell_1$;
\item the scale $q=\LambertW(z)/z=\e^{-\LambertW(z)}$, whose powers form
      separate transseries sectors.
\end{enumerate}
A computation carried out only in powers of $1/\log z$ determines $g$ to
\emph{every} inverse-logarithmic order and still cannot see the first genuine
correction, which is of size $q$.  That is the phenomenon
\cref{q1:thm:blocks} makes explicit, and it is the same ordering caveat that
\cref{p6:rem:hierarchy} raises for the Lambert-cored inversions and that the
scale part of this volume states abstractly.

\begin{sourcebox}
Three filed articles treat this map:
\texttt{lambert\_inverse\_transseries} (\textsf{L1}),
\texttt{lambert\_inverse\_transseries\_bundle} (\textsf{L2}) and
\texttt{reversing\_x\_plus\_lambert\_w\_transseries} (\textsf{L3}).  Between
them they state $27$ named results, of which $18$ are already theorems of the
calculus developed in the earlier parts of this volume --- near-identity
reversion, Lagrange--B\"urmann at infinity, triangular reversal, Newton
acceleration, residual-to-error transfer, the $r$-Lambert function, and the
pure polynomial--logarithmic perturbation lemma among them.  Those are cited,
not restated.  What remains, and what this chapter contains, is the
mathematics specific to $x+\LambertW(x)$: the exact reduction, the sharp
residual bracket, the explicit coefficient blocks and their integer-polynomial
structure, and the ordering of the three scales.
\end{sourcebox}

\section{Exact inversion}
\label{q1:sec:exact}

\begin{theorem}[Exact inverse]\label{q1:thm:exact}
For every $z\ge0$,
\begin{equation}\label{q1:eq:exact}
 g(z)=z-\rWone(z)=\rWone(z)\,\e^{\rWone(z)},
\end{equation}
where $\rWone$ is the $r$-Lambert function of \cref{p7:def:rlambert} at
$\varrho=1$, that is the inverse of $y\mapsto y\e^{y}+y$.  Equivalently, with
$y=\rWone(z)$ one has $g(z)=y\e^{y}$ and $\LambertW(g(z))=y$.
\end{theorem}

\begin{proof}
Put $y=\LambertW(x)$, so that $x=y\e^{y}$.  Then $z=x+\LambertW(x)$ reads
$z=y\e^{y}+y$, which is exactly the defining equation of $\rWone$; hence
$y=\rWone(z)$ and $x=z-y$.  The second form is $x=y\e^{y}$, and
$\LambertW(g(z))=\LambertW(x)=y$.  On $[0,\infty)$ the map
$y\mapsto y\e^{y}+y$ is strictly increasing from $0$, so the principal real
branch is the relevant one throughout.
\end{proof}

\begin{remark}[Formalization status]\label{q1:rem:lean}
\Cref{q1:thm:exact} is machine-checked in this repository.  In
\path{LambertShiftInverse.lean}, $\rWone$ is \path{Fabius.shiftedLambertW} (the
inverse on $[0,\infty)$ of \path{Fabius.shiftedMulExp}, $u\mapsto u\e^{u}+u$)
and $g$ is \path{Fabius.lambertShiftInv}; the three identities of
\eqref{q1:eq:exact} are \path{Fabius.lambertShift_lambertShiftInv},
\path{Fabius.lambertShiftInv_eq_mul_exp} and
\path{Fabius.principalLambertW_lambertShiftInv}, and $g$ is the two-sided
inverse of $f$ by \path{Fabius.lambertShiftInv_lambertShift} together with
\path{Fabius.lambertShift_eq_iff}.

The differential layer is formal as well.  \(g\) is strictly increasing and
continuous on the open half-line (\lean{Fabius.lambertShiftInv_strictMonoOn},
\lean{Fabius.lambertShiftInv_continuousAt}, the latter read off the exact image
\lean{Fabius.lambertShiftInv_image_Ioi}), so the inverse-function rule applies:
\lean{Fabius.lambertShiftInv_hasDerivAt} differentiates \(g\) at every \(z>0\), and
\lean{Fabius.deriv_lambertShiftInv_bounds} is the strict bracket \(1/2<g'(z)<1\), both
inequalities strict because \(0<\LambertW'<1\) once \(\LambertW>0\)
(\lean{Fabius.inv_exp_mul_add_one_lt_one}).  From the strict growth of \(g'\)
(\lean{Fabius.deriv_lambertShiftInv_strictMonoOn}) it follows that \(g\) is strictly
convex there, \lean{Fabius.strictConvexOn_lambertShiftInv}.  The endpoint value
\(g'(0)=1/2\), a one-sided derivative, is not formalized, and neither is the concavity
of \(f\) as such --- what the corpus carries is the strict concavity of the branch
itself, \lean{Fabius.strictConcaveOn_principalLambertW}, proved from
\(\LambertW''<0\).

This chapter is the only place in the present volume with formal counterparts ---
this theorem and its differential layer, the numerator recursion of
\cref{q1:prop:recursion}, and the general bracket of \cref{q1:rem:two-sided} ---
and they are recorded because the corpus keeps a register of what is proved
and what is merely written down.
\end{remark}

\begin{proposition}[Brackets]\label{q1:prop:brackets}
For $x\ge0$ and $z\ge0$,
\begin{equation}\label{q1:eq:brackets}
 1\le f'(x)\le2,
 \qquad
 \tfrac12\le g'(z)<1,
 \qquad
 z-\LambertW(z)\ \le\ g(z)\ \le\ z .
\end{equation}
Moreover $f$ is concave and $g$ convex on $(0,\infty)$.
\end{proposition}

\begin{proof}
Differentiating $\LambertW(x)\e^{\LambertW(x)}=x$ gives
\begin{equation}\label{q1:eq:W-prime}
 \LambertW'(x)=\frac{\LambertW(x)}{x\bigl(1+\LambertW(x)\bigr)}
 =\frac{\e^{-\LambertW(x)}}{1+\LambertW(x)}\in(0,1],
\end{equation}
with the limiting value $1$ at $x=0$; hence $1\le f'=1+\LambertW'\le2$ and,
by the inverse-function rule, $\tfrac12\le g'<1$.  For the bracket, the
defining relation may be written $g(z)=z-\LambertW(g(z))$; since
$0\le g(z)\le z$ and $\LambertW$ increases, $\LambertW(g(z))\le\LambertW(z)$,
which gives the lower bound.  Finally $\LambertW''<0$ on $(0,\infty)$ makes
$f$ concave, and the inverse of a concave increasing function is convex.
\end{proof}

\begin{proposition}[A sharp residual certificate]\label{q1:prop:certificate}
Let $\widetilde g(z)\ge0$ be any approximation and put
\begin{equation}\label{q1:eq:residual}
 R(z)=\widetilde g(z)+\LambertW\bigl(\widetilde g(z)\bigr)-z .
\end{equation}
Then $R$ and $\widetilde g-g$ have the same sign, and
\begin{equation}\label{q1:eq:sandwich}
 \tfrac12\,|R(z)|\ \le\ \bigl|\widetilde g(z)-g(z)\bigr|\ \le\ |R(z)| .
\end{equation}
\end{proposition}

\begin{proof}
This is \cref{p0:thm:backward-error} with $f$ in place of the general map,
sharpened by the two-sided bound $1\le f'\le2$ of
\cref{q1:prop:brackets}: the mean value theorem applied to
$f(\widetilde g)-f(g)=R$ gives $R=f'(\xi)(\widetilde g-g)$ for some
intermediate $\xi$, and dividing by $f'(\xi)\in[1,2]$ gives both inequalities
and the equality of signs.
\end{proof}

\begin{remark}[Why this one is two-sided]\label{q1:rem:two-sided}
\Cref{p0:thm:backward-error} gives a two-sided bound whenever the slope is
bracketed above and below, and here the bracket is the best possible constant
pair $(1,2)$, independent of $z$.  That is unusual: in every Lambert-cored
chapter of this volume the slope grows, so the residual buys a bound that
improves with $z$ but with a $z$-dependent constant.  Here the certificate is
uniform, which is what makes numerical work on this map unusually comfortable.
The sandwich is machine-checked as well, as
\path{Fabius.abs_sub_lambertShiftInv_le_abs_residual} and
\path{Fabius.abs_residual_le_two_mul_abs_sub_lambertShiftInv}; and so is the
general two-sided statement, for an arbitrary function on a convex set, in
\path{MeanValueBracket.lean} --- the brackets are
\path{Fabius.mul_abs_sub_le_abs_sub_of_le_deriv} and
\path{Fabius.abs_sub_le_of_le_deriv}, and the transfer against a right inverse
is \path{Fabius.abs_sub_right_inverse_le_div} together with
\path{Fabius.div_le_abs_sub_right_inverse}.
\end{remark}

\begin{remark}[The sign hypothesis is not decoration]\label{q1:rem:sign}
\Cref{p0:thm:backward-error} assumes $0<m\le f'\le M$, and the positivity of
$m$ is needed for the \emph{upper} bracket, not merely for the interpretation.
Dropping it fails at once: with $f'\equiv-10$ the bracket $m\le f'\le M$ holds
for $m=-10$, $M=0$, while $|f(x)-f(y)|=10|x-y|$ exceeds $M|x-y|=0$.  What the
sign supplies is monotonicity, which is what makes $|f(x)-f(y)|$ the signed
increment.  The lower bracket needs no sign condition.  This is worth stating
because a theorem quoted as ``$f$ strictly increasing with $m\le f'\le M$''
leaves the role of the sign implicit, and the two-sided form is the one this
volume uses most.
\end{remark}

\section{The small parameter, and reduction to an ordinary reversion}
\label{q1:sec:reduction}

Set
\begin{equation}\label{q1:eq:Lq}
 L=\LambertW(z),
 \qquad
 z=L\e^{L},
 \qquad
 q=\e^{-L}=\frac Lz,
\end{equation}
and write $\rWone(z)=L-v$, so that by \cref{q1:thm:exact}
\begin{equation}\label{q1:eq:g-L-v}
 g(z)=z-L+v .
\end{equation}

\begin{proposition}[The exact equation for the correction]\label{q1:prop:B}
The quantity $v$ of \eqref{q1:eq:g-L-v} is characterized exactly by
\begin{equation}\label{q1:eq:B}
 q=B_L(v):=\frac{L}{L-v}-\e^{-v},
\end{equation}
and $B_L(0)=0$ with $B_L'(0)=1+L^{-1}>0$.  Writing $u=L^{-1}$,
\begin{equation}\label{q1:eq:Bvu}
 B(v,u)=\frac1{1-uv}-\e^{-v}=\sum_{k\ge1}b_k(u)v^{k},
 \qquad
 b_k(u)=u^{k}-\frac{(-1)^{k}}{k!},
\end{equation}
so that $b_1(u)=1+u$.  Consequently the large-$z$ problem is an
\emph{ordinary} analytic reversion at the origin, in the variable $v$ with
parameter $u$.
\end{proposition}

\begin{proof}
Substituting $\rWone(z)=L-v$ into its defining equation
$\rWone\e^{\rWone}+\rWone=z=L\e^{L}$ gives
$(L-v)\e^{L-v}+L-v=L\e^{L}$; dividing by $\e^{L}$ and using
$\e^{-L}=q$ yields $(L-v)(\e^{-v}+q)=L$, that is
$q=L/(L-v)-\e^{-v}$, which is \eqref{q1:eq:B}.  Then $B_L(0)=1-1=0$ and
$B_L'(v)=L/(L-v)^{2}+\e^{-v}$, so $B_L'(0)=L^{-1}+1$.  Expanding
$(1-uv)^{-1}=\sum_k u^{k}v^{k}$ and
$\e^{-v}=\sum_k(-1)^{k}v^{k}/k!$ and subtracting gives
\eqref{q1:eq:Bvu}; the $k=0$ terms cancel.
\end{proof}

\begin{remark}[What has been gained]\label{q1:rem:gain}
\Cref{q1:prop:B} is the whole trick.  Before it, one faces a reversion at
infinity across three scales; after it, an analytic function vanishing at the
origin with invertible derivative is to be reverted, which is
\cref{p0:thm:LB} in its convergent form.  The exponentially small parameter
$q$ has become the independent variable, and the slow parameter $u=1/L$ rides
along as a coefficient.  This is the same manoeuvre as the substitution
$u=\e^{-\lambda x}$ in \cref{p9:thm:factorization}, and as the passage to the
Lambert coordinate in \cref{p6:sec:core}: normalize until what is left is an
ordinary reversion.
\end{remark}

\section{The coefficient blocks}
\label{q1:sec:blocks}

\begin{theorem}[All-orders expansion]\label{q1:thm:blocks}
As $z\to+\infty$, with $L=\LambertW(z)$,
\begin{equation}\label{q1:eq:blocks}
 g(z)=z-L+\sum_{n\ge1}\frac{A_n(L)}{z^{n}},
\end{equation}
and the first four coefficients are
\begin{align}
 A_1(L)&=\frac{L^{2}}{1+L},
 \label{q1:eq:A1}\\
 A_2(L)&=\frac{L^{3}\bigl(L^{2}-2\bigr)}{2\,(1+L)^{3}},
 \label{q1:eq:A2}\\
 A_3(L)&=\frac{L^{4}(2L-1)\bigl(L^{3}-6L-6\bigr)}{6\,(1+L)^{5}},
 \label{q1:eq:A3}\\
 A_4(L)&=\frac{L^{5}\bigl(6L^{6}-8L^{5}-69L^{4}-40L^{3}+96L^{2}+72L-24\bigr)}
 {24\,(1+L)^{7}} .
 \label{q1:eq:A4}
\end{align}
\end{theorem}

\begin{proof}
By \cref{q1:prop:B} the correction $v$ solves $B(v,u)=q$ with
$B(0,u)=0$ and $\partial_vB(0,u)=1+u$ a unit, so Lagrange inversion
(\cref{p0:thm:LB}, in the convergent analytic form) determines $v$ as a
power series in $q$ whose coefficients are rational in $u=1/L$; by
\eqref{q1:eq:g-L-v} the same is true of $g(z)-z+L$.  Converting $q^{n}$ to
$L^{n}z^{-n}$ and clearing $u$ gives \eqref{q1:eq:blocks} with $A_n$
rational in $L$.  The displayed $A_1,\dots,A_4$ follow from the first four
Lagrange coefficients; they were recomputed independently for this volume and
checked numerically against the exact inverse
(\cref{q1:rem:verification}).
\end{proof}

\begin{proposition}[Integer-polynomial structure]\label{q1:prop:structure}
For every $m\ge1$ there is $P_m\in\Int[v]$ with
\begin{equation}\label{q1:eq:Am-Pm}
 A_m(v)=\frac{v^{m+1}P_m(v)}{m!\,(1+v)^{2m-1}},
\end{equation}
and
\begin{equation}\label{q1:eq:Pm-data}
 \deg P_m=2m-2,
 \qquad
 \Coef{v^{2m-2}}P_m=(m-1)!,
 \qquad
 P_m(0)=(-1)^{m+1}m! .
\end{equation}
Consequently
\begin{equation}\label{q1:eq:Am-leading}
 A_m(v)\sim\frac{v^{m}}{m}\qquad(v\to+\infty).
\end{equation}
\end{proposition}

\begin{proof}
The shape \eqref{q1:eq:Am-Pm} follows by induction from the Lagrange
recursion of \cref{q1:thm:blocks}: each step divides once by
$b_1(u)=1+u=(1+v)/v$ evaluated at $v=L$, contributing the powers of $(1+v)$,
while the numerators stay integral because the $b_k$ of \eqref{q1:eq:Bvu}
have denominators $k!$ only, absorbed by the $m!$.  The three data in
\eqref{q1:eq:Pm-data} are read off the same recursion: the top coefficient
comes from the pure $u^{k}$ part of $b_k$, the constant term from the pure
$\e^{-v}$ part.  For \eqref{q1:eq:Am-leading}, substituting
$P_m(v)\sim(m-1)!\,v^{2m-2}$ into \eqref{q1:eq:Am-Pm} gives
\[
 A_m(v)\sim\frac{v^{m+1}\,(m-1)!\,v^{2m-2}}{m!\,v^{2m-1}}
 =\frac{v^{m}}{m}.
\]
The four instances $P_1=1$, $P_2=v^{2}-2$,
$P_3=2v^{4}-v^{3}-12v^{2}-6v+6$ and
$P_4=6v^{6}-8v^{5}-69v^{4}-40v^{3}+96v^{2}+72v-24$ read off
\eqref{q1:eq:A1}--\eqref{q1:eq:A4} satisfy \eqref{q1:eq:Pm-data} exactly, as
recorded in \cref{q1:rem:verification}.
\end{proof}


\begin{proposition}[A recursion for the numerators]\label{q1:prop:recursion}
Define $R_{n,k}\in\Int[v]$ by
\begin{equation}\label{q1:eq:recursion}
 R_{n,1}=n,
 \qquad
 R_{n,k+1}=v(1+v)\,R_{n,k}'
 +\Bigl(n(1+v)-(2k-1)v-k(1+v)^{2}\Bigr)R_{n,k}.
\end{equation}
Then $R_{n,k}\in n\,\Int[v]$, $\deg R_{n,k}=2(k-1)$,
$\Coef{v^{2(k-1)}}R_{n,k}=n(-1)^{k-1}(k-1)!$ and
$R_{n,k}(0)=n(n-1)\cdots(n-k+1)$; and the numerators of
\cref{q1:prop:structure} are the diagonal
\begin{equation}\label{q1:eq:Pm-from-R}
 P_m=\frac{(-1)^{m+1}}{m+1}\,R_{m+1,m}.
\end{equation}
The three data \eqref{q1:eq:Pm-data} are exactly the three displayed facts
read on that diagonal.
\end{proposition}

\begin{proof}
Divisibility by $n$ is immediate from \eqref{q1:eq:recursion} by induction,
since $R_{n,1}=n$ and every step is $\Int[v]$-linear in $R_{n,k}$.  For the
degree, the derivative term $v(1+v)R_{n,k}'$ has degree at most
$2+\bigl(2(k-1)-1\bigr)=2k-1$, strictly below the degree
$2+2(k-1)=2k$ of the multiplier term, whose leading coefficient is that of
$-k(1+v)^{2}$ against $R_{n,k}$, namely $-k$ times $n(-1)^{k-1}(k-1)!$; and
$-k\cdot(-1)^{k-1}(k-1)!=(-1)^{k}k!$, which is the claim at $k+1$.  The value
at $v=0$ satisfies $R_{n,k+1}(0)=(n-k)R_{n,k}(0)$, giving the falling
factorial.  Substituting $n=m+1$, $k=m$ into the three then gives
$\deg P_m=2m-2$, $\Coef{v^{2m-2}}P_m=(m-1)!$ and
$P_m(0)=(-1)^{m+1}m!$, which is \eqref{q1:eq:Pm-data}; the normalization
\eqref{q1:eq:Pm-from-R} is fixed by matching $P_1=1$.
\end{proof}

\begin{remark}[Formalization status of the numerator family]
\label{q1:rem:lean-coefficients}
\Cref{q1:prop:recursion} is machine-checked, in
\path{LambertInverseCoefficients.lean}, a module of pure polynomial algebra
with no dependence on the rest of the corpus.  With indices shifted so the
recursion starts at $0$, \path{Fabius.lambertNumerator} is $R_{n,j+1}$;
\path{Fabius.lambertNumerator_succ_eq_draft} is \eqref{q1:eq:recursion},
\path{Fabius.natCast_dvd_lambertNumerator} is $R_{n,k}\in n\Int[v]$,
\path{Fabius.natDegree_lambertNumerator} the degree,
\path{Fabius.leadingCoeff_lambertNumerator} the leading coefficient, and
\path{Fabius.eval_zero_lambertNumerator} the value at $0$.

What is \emph{not} formalized, and should not be read as though it were: the
closed form for $(\dd/\dd z)^{k}v^{n}$, and hence the identification of this
polynomial family with the coefficients $A_m$ of \cref{q1:thm:blocks}; and the
asymptotic reading $A_m\sim v^{m}/m$ of \eqref{q1:eq:Am-leading}.  Both need
the operator calculus on functions rather than the recursion, and the
formalized statements are about $\Int[v]$ alone.
\end{remark}

\begin{theorem}[The error law, and how slowly it is approached]
\label{q1:thm:error-law}
Let $G_N(z)=z-L+\sum_{n\le N}A_n(L)z^{-n}$.  Then, for each fixed $N\ge0$,
\begin{equation}\label{q1:eq:error-law}
 g(z)-G_N(z)\ \sim\ \frac1{N+1}\left(\frac{\LambertW(z)}{z}\right)^{N+1}
 \qquad(z\to+\infty).
\end{equation}
The relative approach to this limit is only $\bigO(1/L)=\bigO(1/\log z)$.
\end{theorem}

\begin{proof}
By \eqref{q1:eq:blocks} the first omitted term is $A_{N+1}(L)z^{-N-1}$, and
by \eqref{q1:eq:Am-leading} $A_{N+1}(L)\sim L^{N+1}/(N+1)$; hence the
difference is $\sim L^{N+1}/\bigl((N+1)z^{N+1}\bigr)$, which is
\eqref{q1:eq:error-law} since $q=L/z$.  The correction to
\eqref{q1:eq:Am-leading} is $\bigO(1/v)$, because $P_m$ has degree exactly
$2m-2$ and its next coefficient is finite while $(1+v)^{2m-1}$ contributes
$1+\bigO(1/v)$; so the relative error in \eqref{q1:eq:error-law} is
$\bigO(1/L)$.
\end{proof}

\begin{remark}[The caveat is not academic]\label{q1:rem:slow}
Because $L=\LambertW(z)$ grows like $\log z$, the $\bigO(1/L)$ in
\cref{q1:thm:error-law} is genuinely slow.  Numerically, the ratio of
$g-G_N$ to the right-hand side of \eqref{q1:eq:error-law} at $N=1$ is
$0.55$ at $z=10^{3}$, $0.72$ at $10^{5}$, $0.82$ at $10^{8}$, $0.88$ at
$10^{12}$ and $0.93$ at $10^{20}$ --- converging to $1$, but not fast enough
for \eqref{q1:eq:error-law} to be a useful predictor at any argument one is
likely to meet.  The certificate of \cref{q1:prop:certificate}, which is
uniform and two-sided, is the tool to use instead.  A source that quotes the
error law without this caveat invites the reader to trust it at $z=50$, where
it is wrong by a factor of three.
\end{remark}

\section{The three scales, and what the logarithmic expansion cannot see}
\label{q1:sec:scales}

\begin{proposition}[Pure logarithmic asymptotics]\label{q1:prop:log}
With $\ell_1=\log z$ and $\ell_2=\log\ell_1$,
\begin{equation}\label{q1:eq:log-expansion}
 g(z)=z-\ell_1+\ell_2-\frac{\ell_2}{\ell_1}
 +\frac{\ell_2(2-\ell_2)}{2\ell_1^{2}}
 -\frac{\ell_2\bigl(6-9\ell_2+2\ell_2^{2}\bigr)}{6\ell_1^{3}}
 +\frac{\ell_2\bigl(12-36\ell_2+22\ell_2^{2}-3\ell_2^{3}\bigr)}{12\ell_1^{4}}
 +\cdots,
\end{equation}
the displayed series being $z$ minus the classical iterated-logarithm
expansion of $\LambertW(z)$.  \emph{Every} term of \eqref{q1:eq:log-expansion}
is a term of $z-\LambertW(z)$: the entire block structure of
\cref{q1:thm:blocks} is invisible at this scale.
\end{proposition}

\begin{proof}
By \cref{q1:prop:brackets}, $0\le z-\LambertW(z)-g(z)\le\LambertW(z)-
\LambertW(g(z))$, and by \eqref{q1:eq:blocks} the difference is
$A_1(L)/z+\bigO(L^{3}/z^{2})=\bigO(L/z)$.  Since $L/z\to0$ faster than every
power of $1/\ell_1$ --- indeed $z^{-1}L\,\ell_1^{N}\to0$ for every fixed $N$,
because $z$ grows exponentially in $L$ --- the two functions
$g(z)$ and $z-\LambertW(z)$ have the same asymptotic expansion to all orders
on the scale $\{\ell_1^{-k}\ell_2^{j}\}$.  Substituting the classical
expansion of $\LambertW$ then gives \eqref{q1:eq:log-expansion}.
\end{proof}

\begin{remark}[The moral, and where else it applies]\label{q1:rem:moral}
This is the cleanest instance in the volume of a general trap.  One may
compute a correct inverse-logarithmic expansion to any order and still have
computed nothing about the correction that actually distinguishes $g$ from its
leading approximation, because that correction lies beyond every order of the
scale one is working on.  The same ordering governs the additive constant
$\tfrac12$ in \cref{p6:thm:gamma}, the shift $+1$ in the Barnes hierarchy, and
the Bell sector of \cref{p8:prop:hierarchy}; \cref{p6:rem:hierarchy} states
the rule.  What makes $x+\LambertW(x)$ the right teaching example is that here
both scales are elementary, so the gap between them can be displayed in full
rather than described.
\end{remark}

\begin{remark}[A normalization clash between the sources]
\label{q1:rem:normalization}
\textsf{L3} writes the expansion as
$g(z)=z-L+\sum_{n\ge1}\mathcal A_n(L)q^{n}$ with $q=L/z$, while \textsf{L1}
writes it as $g(z)=z-L+\sum_{n\ge1}A_n(L)z^{-n}$ and states the structure
theorem \eqref{q1:eq:Am-Pm} for \emph{its} $A_n$.  The two families differ by
$\mathcal A_n=A_n/L^{n}$, so a reader who carries the structure theorem across
without the substitution will find it fails at once --- $A_1=L^{2}/(1+L)$
against $\mathcal A_1=L/(1+L)$.  This chapter uses the $z^{-n}$ normalization
throughout, which is the one in which \eqref{q1:eq:Am-Pm} is true.
\end{remark}

\begin{remark}[Independent verification]\label{q1:rem:verification}
Recomputed for this volume at $50$--$60$ digits, from the definitions and
independently of the sources: the exact inversion \eqref{q1:eq:exact} and the
identity $\LambertW(g(z))=\rWone(z)$, at $z=1,5,20,100,1000$; the brackets
\eqref{q1:eq:brackets}, including $\tfrac12\le g'<1$ by numerical
differentiation; the residual sandwich \eqref{q1:eq:sandwich}, at three
arguments and three trial approximations each, signs included; the four
blocks \eqref{q1:eq:A1}--\eqref{q1:eq:A4}, by comparing $G_N$ against the
root of $x+\LambertW(x)=z$ and confirming that the error falls by the
predicted factor at $z=500$ and $z=5000$; the three structural data
\eqref{q1:eq:Pm-data} for $m\le4$, read off the displayed numerators; and the
error law \eqref{q1:eq:error-law} together with the rate of
\cref{q1:rem:slow}, at $z=10^{3},10^{5},10^{8},10^{12},10^{20}$.  The
logarithmic expansion \eqref{q1:eq:log-expansion} was checked against
$\LambertW$ at $z=10^{6}$ and $10^{12}$.
\end{remark}

\section{What is cited rather than reproved}
\label{q1:sec:cited}

For the record, and because the three sources each prove these
independently: the near-identity reversion theorem, the triangular reversal
algorithm, Newton acceleration with precision doubling, the
Lagrange--B\"urmann formula at infinity, the coefficient-operator form, the
residual-to-error transfer, the $r$-Lambert function itself, and the pure
polynomial--logarithmic perturbation lemma are all results of the calculus
parts of this volume, and are used here by citation.  The general
``method for reversing logarithmic transseries'' that each of the three
articles extracts in its second half is, in each case, the reversal part of
this volume specialized to a slow coefficient field; the three extractions
agree with each other and with it.

\begin{remark}[What is left open]\label{q1:rem:open}
The family $x+\alpha\LambertW(x)$ and the more general $x+h(\LambertW(x))$
are treated in the sources only to leading order.  The reduction of
\cref{q1:prop:B} generalizes --- one obtains
$q=L/(L-v)-\e^{-v}$ with $L$ replaced by the appropriate combination --- but
the analogue of the integer-polynomial structure
\cref{q1:prop:structure} is not established here, and it is not obvious that
the numerators remain integral for general $\alpha$.
\end{remark}


\part{Four combinatorial sequences}

\chapter{Butcher--P\'olya rooted trees (A000081)}
\label{p1:sec:top}

\begin{sourcebox}
This chapter merges the three independently written treatments
\emph{Butcher\_Tree\_Transseries}, \emph{Butcher\_Tree\_Transseries-2} and
\emph{Butcher\_Tree\_Transseries-3} of the special-function-inversion draft set. Below they
are called \textsf{S1}, \textsf{S2} and \textsf{S3}. \textsf{S2} is a book-length treatment
and is the most complete on the inverse and on the inherited singularity at $-\sqrt\rho$;
\textsf{S3} is the most careful about what is and is not proved beyond the dominant sector;
\textsf{S1} has the shortest route to the universal Cayley coefficients and the only treatment
of the compositional inverse of the generating function. Where they disagree the disagreement
is stated in the text.  Every numerical constant retained here has been recomputed from the
exact divisor recurrence at $60$--$110$ working decimal digits; see \cref{p1:sub:audit}.
\end{sourcebox}

The subject of this chapter is the sequence
\begin{equation}
 0,\;1,\;1,\;2,\;4,\;9,\;20,\;48,\;115,\;286,\;719,\;1842,\;4766,
 \;12486,\;32973,\;87811,\;\ldots
 \label{p1:eq:sequence}
\end{equation}
(OEIS A000081), where for $n\ge1$ the term $a_n$ counts isomorphism classes of rooted,
unlabelled, non-plane trees on $n$ vertices.  These are the \emph{P\'olya trees}; in numerical
analysis the same shapes index Butcher series, elementary differentials and Runge--Kutta order
conditions, which is the origin of the name \emph{Butcher tree numbers}.  We set $a_0=0$
throughout.

The sequence grows like $C\alpha^nn^{-3/2}$ with $\alpha=\rho^{-1}$ and $\rho$ the Otter
singularity.  Thus, in the language of \cref{p0:def:model}, its asymptotic interpolation is an
exponential--power model with
\begin{equation}
 \boxed{\;\alpha_{\mathrm{model}}=1,\qquad
 a_{\mathrm{model}}=\lambda:=\log(1/\rho),\qquad
 b_{\mathrm{model}}=-\tfrac32,\qquad
 \delta=1.\;}
 \label{p1:eq:dictionary-preview}
\end{equation}
This chapter is therefore the \emph{linear-phase} instance of the Part~0 apparatus: no
$\alpha$-reduction (\cref{p0:lem:alpha-reduction}) is needed, the Lambert core
(\cref{p0:thm:lambert-core}) applies verbatim, and the entire reversion is a substitution into
\cref{p0:thm:lambert-centered}.  Everything specific to rooted trees is the computation of the
model's data $C$, $\lambda$ and $Q$; that computation is the substance of
\crefrange{p1:sec:polya}{p1:sec:forward}, and the inversion in \cref{p1:sec:inverse} is short
precisely because Part~0 has already been proved.

\section{The functional equation and the exact recurrence}
\label{p1:sec:polya}

\subsection{P\'olya's multiset equation}

Let $\mathcal T$ be the class of unlabelled non-plane rooted trees and let $\mathcal Z$ be one
atom (one vertex).  Deleting the root of a tree leaves an unordered multiset of rooted trees,
and grafting a multiset of rooted trees onto a fresh root inverts this.  Hence
\begin{equation}
 \mathcal T=\mathcal Z\times\operatorname{MSET}(\mathcal T).
 \label{p1:eq:species}
\end{equation}

\begin{proposition}[P\'olya's equation and the Euler product]
\label{p1:prop:polya}
Let $T(z)=\sum_{n\ge1}a_nz^n\in\Int[[z]]$.  Then, as identities of formal power series,
\begin{align}
 T(z)&=z\prod_{r\ge1}(1-z^r)^{-a_r},
 \label{p1:eq:euler-product}\\
 T(z)&=z\exp\!\left(\sum_{k\ge1}\frac{T(z^k)}{k}\right).
 \label{p1:eq:polya}
\end{align}
\end{proposition}

\begin{proof}
A multiset of rooted trees that contains $m_j$ copies of the $j$th shape has size
$\sum_jm_j|j|$ and contributes $z^{\sum_jm_j|j|}$.  For one fixed shape of size $r$, unbounded
multiplicity contributes $\sum_{m\ge0}z^{rm}=(1-z^r)^{-1}$; there are $a_r$ shapes of size
$r$, and distinct shapes are chosen independently, so the forest hanging from the root has
generating function $\prod_{r\ge1}(1-z^r)^{-a_r}$.  Every factor is a unit of $\Int[[z]]$ and
the product is $z$-adically convergent because the $r$th factor is $1+\bigO(z^r)$. Multiplying
by the root gives \cref{p1:eq:euler-product}.  Taking the formal logarithm and using
$-\log(1-w)=\sum_{k\ge1}w^k/k$,
\[
\log\frac{T(z)}{z} =\sum_{r\ge1}a_r\sum_{k\ge1}\frac{z^{rk}}{k}
=\sum_{k\ge1}\frac1k\sum_{r\ge1}a_r(z^k)^r =\sum_{k\ge1}\frac{T(z^k)}{k},
\]
the rearrangement being legitimate because for each fixed power of $z$ only finitely many
pairs $(r,k)$ contribute.  Exponentiating gives \cref{p1:eq:polya}.
\end{proof}

\begin{remark}
\Cref{p1:eq:polya} is purely formal; no analytic hypothesis is used, and none is available at
this stage.  All three siblings state it correctly, but only \textsf{S2} says explicitly that
the identity holds in $\Int[[z]]$ before any radius of convergence has been produced.  We keep
that discipline: \cref{p1:sec:dominant} is the first place where analysis enters.
\end{remark}

\subsection{The divisor recurrence}

\begin{definition}[Divisor transform]
\label{p1:def:divisor}
For $m\ge1$ put
\begin{equation}
 b_m:=\sum_{d\mid m}d\,a_d .
 \label{p1:eq:divisor}
\end{equation}
\end{definition}

\begin{theorem}[Otter--P\'olya recurrence]
\label{p1:thm:recurrence}
With $a_1=1$ one has, for every $n\ge2$,
\begin{equation}
 (n-1)\,a_n=\sum_{j=1}^{n-1}b_j\,a_{n-j}.
 \label{p1:eq:recurrence}
\end{equation}
\end{theorem}

\begin{proof}
Write $G(z)=T(z)/z=\sum_{n\ge0}a_{n+1}z^n$, a unit of $\Rat[[z]]$ with $G(0)=1$. From
\cref{p1:eq:euler-product} and the computation in the proof of \cref{p1:prop:polya},
\[
\log G(z)=\sum_{r\ge1}a_r\sum_{k\ge1}\frac{z^{rk}}{k} =\sum_{m\ge1}\frac{z^m}{m}\sum_{d\mid
m}d\,a_d =\sum_{m\ge1}\frac{b_m}{m}z^m,
\]
where the inner rearrangement groups the pairs $(r,k)$ with $rk=m$ and uses $1/k=d/m$ for
$d=r$.  Formal logarithmic differentiation gives $G'(z)=G(z)\sum_{m\ge1}b_mz^{m-1}$.
Extracting $\Coef{z^{n-1}}$,
\[
n\,a_{n+1}=\sum_{j=1}^{n}b_j\,a_{n-j+1},
\]
and replacing $n$ by $n-1$ yields \cref{p1:eq:recurrence}.
\end{proof}

\begin{remark}[Why the recurrence is triangular, and its cost]
\label{p1:rem:triangular}
Although $b_n$ contains the summand $n\,a_n$, the right-hand side of \cref{p1:eq:recurrence}
uses only $b_j$ with $j<n$, so $a_n$ is determined by $a_1,\ldots,a_{n-1}$.  Implemented as a
sieve --- after computing $a_n$, add $n\,a_n$ to $b_n,b_{2n},b_{3n},\ldots$ --- the divisor
transform costs $\bigO(N\log N)$ additions and the convolutions cost $\bigO(N^2)$ integer
operations for $a_1,\ldots,a_N$.  The integrality assertion $(n-1)\mid\sum_jb_ja_{n-j}$ is a
free exact self-test on every step, and was used as such in the computations of
\cref{p1:sub:audit}.
\end{remark}

\subsection{The self-interaction and its coefficients}

Separating the $k=1$ term of \cref{p1:eq:polya} is the structural move on which everything
else rests.

\begin{definition}[Disturbance and Cayley argument]
\label{p1:def:R-zeta}
Put
\begin{equation}
 R(z):=\sum_{k\ge2}\frac{T(z^k)}{k},
 \qquad
 \zeta(z):=z\,\e^{R(z)} .
 \label{p1:eq:R-zeta}
\end{equation}
\end{definition}

With this notation \cref{p1:eq:polya} reads
\begin{equation}
 T(z)=\zeta(z)\,\e^{T(z)} .
 \label{p1:eq:T-zeta}
\end{equation}

\begin{lemma}[Coefficients of the disturbance]
\label{p1:lem:R-coefficients}
$R(z)=\sum_{r\ge2}\varrho_rz^r$ with
\begin{equation}
 \varrho_r=\frac1r\sum_{\substack{d\mid r\\ d<r}}d\,a_d
 =\frac{b_r-r\,a_r}{r}\in\Rat_{\ge0}.
 \label{p1:eq:varrho}
\end{equation}
In particular $\varrho_r$ depends only on $a_d$ with $d\le r/2$, and $\varrho_2=\tfrac12$,
$\varrho_3=\tfrac13$, $\varrho_4=\tfrac34$, $\varrho_5=\tfrac15$, $\varrho_6=\tfrac32$,
$\varrho_7=\tfrac17$, $\varrho_8=\tfrac{19}{8}$.
\end{lemma}

\begin{proof}
$R(z)=\sum_{k\ge2}\sum_{d\ge1}\frac{a_d}{k}z^{kd}$; collecting the terms with $kd=r$ gives
$\varrho_r=\sum_{d\mid r,\,d<r}a_d/(r/d)=\frac1r\sum_{d\mid r,\,d<r}d\,a_d$, which is
$\left(b_r-ra_r\right)/r$ by \cref{p1:eq:divisor}. Every divisor $d<r$ of $r$ satisfies $d\le
r/2$.  The listed values follow from $a_1=a_2=1$, $a_3=2$, $a_4=4$: e.g.\
$\varrho_4=(1\cdot1+2\cdot1)/4=\tfrac34$ and $\varrho_6=(1+2\cdot1+3\cdot2)/6=\tfrac32$.
\end{proof}

\Cref{p1:lem:R-coefficients} is due to \textsf{S2}; \textsf{S1} and \textsf{S3} carry $R$ only
through the double sum $\sum_{k\ge2}\sum_na_n\rho^{kn}(\cdots)$. The single-sum form is used
throughout this chapter because it turns every quantity below into one geometrically
convergent series in the exact integers $a_r$ (\cref{p1:lem:Rhat}).

\section{The Cayley--Lambert normal form and the Otter singularity}
\label{p1:sec:dominant}

\subsection{Factorisation through the Cayley tree function}

\begin{definition}[Cayley tree function]
\label{p1:def:cayley}
Let $\mathcal C(w)$ be the unique formal solution of $\mathcal C=w\e^{\mathcal C}$ with
$\mathcal C(0)=0$.  It is analytic on $|w|<\e^{-1}$ and
\begin{equation}
 \mathcal C(w)=-\Wlam_0(-w),
 \label{p1:eq:cayley-lambert}
\end{equation}
where $\Wlam_0$ is the principal Lambert branch.
\end{definition}

Comparing \cref{p1:eq:T-zeta} with \cref{p1:def:cayley},
\begin{equation}
 \boxed{\;T(z)=\mathcal C\bigl(\zeta(z)\bigr)
 =-\Wlam_0\bigl(-\zeta(z)\bigr).\;}
 \label{p1:eq:factorisation}
\end{equation}
This holds as formal series, and as an identity of analytic functions on the common domain of
the branches selected at the origin.  The point of the factorisation is a clean separation of
labour: \emph{all} of the singular mechanism is the universal square-root branch point of
$\mathcal C$ at $w=\e^{-1}$, and \emph{all} of the tree-specific information is packed into
the map $\zeta$, which --- as we now show --- is analytic on a strictly larger disc than $T$
itself.

\begin{lemma}[Elementary bounds on the radius]
\label{p1:lem:radius-bounds}
The radius of convergence $\rho$ of $T$ satisfies $\tfrac14\le\rho\le\e^{-1}$.
\end{lemma}

\begin{proof}
Since $a_n\ge1$ for $n\ge1$ we have $\rho\le1$, so for $0<x<\rho$ every $x^{k}$ with $k\ge2$
lies in $(0,\rho)$ as well.  Upper bound: for such $x$ the series $T(x)$ and $R(x)$ converge,
all terms are nonnegative, and $\zeta(x)=x\e^{R(x)}\ge x$; so \cref{p1:eq:T-zeta} gives
$T(x)\ge x\e^{T(x)}$, i.e.\ $x\le T(x)\e^{-T(x)}\le\sup_{u\ge0}u\e^{-u}=\e^{-1}$.  Letting
$x\uparrow\rho$ gives $\rho\le\e^{-1}$.  Lower bound: forgetting the identification of
isomorphic children embeds the shapes counted by $a_n$ into the rooted \emph{plane} trees on
$n$ vertices, of which there are $\frac1n\binom{2n-2}{n-1}\le4^{n}$; hence $\limsup
a_n^{1/n}\le4$ and $\rho\ge\tfrac14$.
\end{proof}

\begin{proposition}[Analyticity of the disturbance]
\label{p1:prop:R-analytic}
$R$, and hence $\zeta$, is analytic on $|z|<\sqrt\rho$, and the series $\sum_{k\ge2}T(z^k)/k$
converges normally on compact subsets together with all its derivatives.  In particular $R$
and $\zeta$ are analytic at $z=\rho$, since $\rho<\sqrt\rho$ by \cref{p1:lem:radius-bounds}.
\end{proposition}

\begin{proof}
Fix $r<\sqrt\rho$, so $r^2<\rho$.  Since $T$ has radius $\rho$ and $T(w)=w+\bigO(w^2)$, there
are $M>0$ and $r_0\in(0,1)$ with $|T(w)|\le M|w|$ for $|w|\le r_0$.  For $|z|\le r$, only
finitely many $k$ have $r^k>r_0$, and for each such $k$ we have $r^k\le r^2<\rho$, so $T(z^k)$
is defined and bounded there.  For the remaining $k$, $|T(z^k)/k|\le Mr^k/k$, which is
summable. Differentiating $j$ times multiplies the $k$th tail term by a factor $\bigO(k^j
r^{-j})$, still summable.  Weierstrass' theorem gives analyticity and termwise
differentiation.
\end{proof}

\begin{proposition}[Critical characterisation of $\rho$]
\label{p1:prop:critical}
With $\rho\in[\tfrac14,\e^{-1}]$ the radius of convergence of $T$,
\begin{equation}
 T(\rho)=1,\qquad \zeta(\rho)=\e^{-1},\qquad
 \log\rho+1+R(\rho)=0.
 \label{p1:eq:critical}
\end{equation}
Moreover $\zeta$ is strictly increasing on $(0,\sqrt\rho)$ and
\begin{equation}
 \theta_1:=\e\rho\,\zeta'(\rho)=1+\rho R'(\rho)>0 .
 \label{p1:eq:theta1}
\end{equation}
\end{proposition}

\begin{proof}
All coefficients of $T$, hence of $R$ and of $\zeta=z\e^{R}$, are nonnegative, and
$\zeta(z)=z+\tfrac12z^{3}+\tfrac13z^{4}+\cdots$ has positive linear coefficient; therefore
$\zeta'(x)>0$ on $(0,\sqrt\rho)$, and $\zeta$ is strictly increasing there.

\emph{Step 1: $T(x)<1$ on $[0,\rho)$, and $T=\mathcal C\circ\zeta$ there.}  On $[0,\rho)$ the
function $T$ is continuous, increasing, with $T(0)=0$, and \cref{p1:eq:T-zeta} gives
$\zeta(x)=T(x)\e^{-T(x)}\le\e^{-1}$.  Let $\rho_1:=\sup\{x\in[0,\rho):T(x)<1\}$.  On
$[0,\rho_1)$ we have $\zeta(x)<\e^{-1}$, so $\mathcal C$ is analytic at $\zeta(x)$ and
$\mathcal C\circ\zeta$ is an analytic solution of \cref{p1:eq:T-zeta} agreeing with $T$ at
$0$; by uniqueness of that germ and analytic continuation, $T=\mathcal C\circ\zeta$ on
$[0,\rho_1)$.  If $\rho_1<\rho$ then $T(\rho_1)=1$ and $\zeta(\rho_1)=\e^{-1}$; but $\zeta$ is
strictly increasing, so $\zeta(x)>\e^{-1}$ for $x\in(\rho_1,\rho)$, contradicting
$\zeta(x)\le\e^{-1}$.  Hence $\rho_1=\rho$.

\emph{Step 2: $\zeta(\rho)=\e^{-1}$.}  By \cref{p1:prop:R-analytic} the function $\zeta$ is
analytic at $\rho$, and by Step~1 and monotonicity $\zeta(\rho)=\lim_{x\uparrow\rho}
\zeta(x)\le\e^{-1}$.  If the inequality were strict, $\mathcal C$ would be analytic at
$\zeta(\rho)$ and $\mathcal C\circ\zeta$ would be analytic in a full neighbourhood of $\rho$
and would agree with $T$ on $[0,\rho)$; then $T$ would continue analytically through $\rho$,
contradicting Pringsheim's theorem, which makes the positive point $z=\rho$ a singularity of a
series with nonnegative coefficients.  Hence $\zeta(\rho)=\e^{-1}$ and, by
\cref{p1:def:cayley}, $T(\rho)=\mathcal C(\e^{-1})=1$.

\emph{Step 3.}  Taking logarithms in $\zeta(\rho)=\rho\e^{R(\rho)}=\e^{-1}$ gives the scalar
equation; equivalently $R(\rho)=\lambda-1$ with $\lambda=-\log\rho$, and $R(\rho)>0$
independently re-proves $\rho<\e^{-1}$.  Finally $\zeta'=(1/z+R')\zeta$, so
$\e\rho\zeta'(\rho)=\e\zeta(\rho)\bigl(1+\rho R'(\rho)\bigr)=1+\rho R'(\rho)$, positive
because $R$ has nonnegative coefficients.
\end{proof}

\begin{proposition}[Uniqueness of the dominant singularity]
\label{p1:prop:unique-dominant}
$z=\rho$ is the only singularity of $T$ on the circle $|z|=\rho$.
\end{proposition}

\begin{proof}
Take $z=\rho\e^{\ii\vartheta}$ with $\vartheta\not\equiv0$.  The series
$\zeta(z)=z+\tfrac12z^3+\tfrac13z^4+\cdots$ has nonnegative coefficients, so
$|\zeta(z)|\le\zeta(\rho)=\e^{-1}$, with equality only if the three displayed nonzero terms
$z$, $\tfrac12z^3$, $\tfrac13z^4$ share one argument, i.e.\
$\e^{2\ii\vartheta}=\e^{3\ii\vartheta}=1$, hence $\e^{\ii\vartheta}=1$ and $\vartheta\equiv0$.
So $|\zeta(z)|<\e^{-1}$ strictly, $\mathcal C$ is analytic at $\zeta(z)$, and $\zeta$ is
analytic at $z$ because $\rho<\sqrt\rho$ (\cref{p1:prop:R-analytic}).  Hence $T=\mathcal
C\circ\zeta$ is analytic at $z$.
\end{proof}

\begin{remark}
This is \textsf{S3}'s argument, and it is the only one of the three that is self-contained:
\textsf{S1} and \textsf{S2} appeal to ``positivity and aperiodicity'' without exhibiting the
aperiodicity witness.  The three coefficients $z,z^3,z^4$ of $\zeta$ furnish that witness
explicitly, because $\gcd(3-1,4-1)=1$.
\end{remark}

\subsection{The constants}
\label{p1:sub:constants}

\begin{convention}[Growth constants]
\label{p1:def:constants}
Throughout the chapter
\begin{equation}
 \alpha:=\rho^{-1},\qquad \lambda:=\log\alpha=-\log\rho,\qquad p:=\tfrac32 .
 \label{p1:eq:alpha-lambda}
\end{equation}
\end{convention}

\begin{repairbox}
\textbf{Notation collision between the siblings.}  \textsf{S2} and \textsf{S3} write
$\alpha=\rho^{-1}$ and $\lambda=\log\alpha$; \textsf{S1} writes exactly the opposite,
$\lambda=\rho^{-1}$ and $\alpha=-\log\rho$.  The symbols therefore cross-denote between the
drafts, and \textsf{S1}'s $\Omega_-/\alpha$, $\e^{\alpha n}$ and $d_m$ formulas must be read
with $\alpha=\lambda_{\text{here}}$.  This volume adopts the two-to-one majority reading
\cref{p1:eq:alpha-lambda}; every formula imported from \textsf{S1} has been translated.
\end{repairbox}

Solving $\log\rho+1+R(\rho)=0$ from \cref{p1:eq:varrho} with $700$ exact tree numbers at $110$
working digits (see \cref{p1:sub:audit}) gives
\begin{align}
 \rho&=0.338321856899207695196112625717017053183774607532967795572303776
 \ldots,
 \label{p1:eq:rho-numeric}\\
 \alpha&=2.955765285651994974714817524123194588375492304663596595350472478
 \ldots,
 \label{p1:eq:alpha-numeric}\\
 \lambda&=1.083757597244624660482711980896174279049662124202672234914427065
 \ldots,
 \label{p1:eq:lambda-numeric}\\
 R(\rho)&=0.083757597244624660482711980896174279049662124202672234914427065
 \ldots,
 \label{p1:eq:Rrho-numeric}
\end{align}
and, as an independent consistency check on the whole pipeline,
\begin{equation}
 \lambda=1+R(\rho),
 \label{p1:eq:lambda-check}
\end{equation}
which is nothing but \cref{p1:eq:critical} rearranged; the two computed values agree to all
$70$ digits carried.  Two further constants are recorded here because they are quoted
repeatedly:
\begin{equation}
 \theta_1=1.216004561824048922157212619228187444\ldots,
 \qquad
 \zeta'(\rho)=1.322241142696930264451284889\ldots .
 \label{p1:eq:theta1-numeric}
\end{equation}

\section{The universal Cayley branch germ}
\label{p1:sec:cayley}

The next result is universal: it depends on the Cayley function alone, not on rooted trees,
and applies verbatim to \emph{any} class satisfying $T=\zeta\e^{T}$ with analytic $\zeta$.

\begin{definition}[Cayley kernel]
\label{p1:def:cayley-kernel}
Put
\begin{equation}
 \Phi(v):=\frac{2\bigl(1-(1-v)\e^{v}\bigr)}{v^{2}}
 =\sum_{m\ge0}\frac{2(m+1)}{(m+2)!}\,v^{m}
 =\sum_{m\ge0}\frac{2}{(m+2)\,m!}\,v^{m}
 =1+\frac23v+\frac14v^{2}+\frac1{15}v^{3}+\cdots .
 \label{p1:eq:Phi}
\end{equation}
\end{definition}

The two displayed coefficient forms agree because $2(m+1)/(m+2)!=2/((m+2)\,m!)$; \textsf{S1}
and \textsf{S3} use the first, \textsf{S2} the second.

\begin{theorem}[Universal Cayley coefficients]
\label{p1:thm:omega}
Set $w=(1-\upsilon)/\e$ and $\mathcal C(w)=1-v$ on the branch reached from $0<w<\e^{-1}$. Then
$\upsilon=\tfrac12v^{2}\Phi(v)$, the inverse germ is
$v=\sum_{m\ge1}\omega_m(2\upsilon)^{m/2}$, and for every $m\ge1$
\begin{equation}
 \boxed{\;\omega_m=\frac1m\,\Coef{v^{m-1}}\Phi(v)^{-m/2}
 =\frac1m\sum_{k=0}^{m-1}\binom{-m/2}{k}
 \OB_{m-1,k}\!\left(\tfrac23,\tfrac14,\tfrac1{15},\ldots\right)
 \in\Rat .\;}
 \label{p1:eq:omega-master}
\end{equation}
Consequently, on that branch,
\begin{equation}
 \mathcal C\!\left(\frac{1-\upsilon}{\e}\right)
 =1+\sum_{r\ge1}c_r\,\upsilon^{r/2},
 \qquad
 c_r=-2^{r/2}\,\omega_r .
 \label{p1:eq:cayley-puiseux}
\end{equation}
\end{theorem}

\begin{proof}
From $\mathcal C=w\e^{\mathcal C}$ we get $\e w=\mathcal C\e^{1-\mathcal C}$, so with
$\mathcal C=1-v$,
\[
1-\upsilon=\e w=(1-v)\e^{v}, \qquad\text{hence}\qquad \upsilon=1-(1-v)\e^{v}
=\sum_{m\ge2}\frac{m-1}{m!}v^{m} =\frac{v^{2}}{2}\Phi(v),
\]
the middle identity by expanding $\e^v$ and collecting.  Put $\varsigma=\sqrt{2\upsilon}$ with
the determination $\varsigma/v\to1$; then $\varsigma=v\,\Phi(v)^{1/2}$ is an ordinary power
series in $v$ with nonzero linear coefficient, so \cref{p0:thm:LB} applies with
$\phi(v)=\Phi(v)^{1/2}$ and yields
\[
\Coef{\varsigma^{m}}v(\varsigma) =\frac1m\Coef{v^{m-1}}\phi(v)^{-m}
=\frac1m\Coef{v^{m-1}}\Phi(v)^{-m/2}=\omega_m .
\]
The Bell form is \cref{p0:lem:power-log} applied to the unit series $\Phi$, whose coefficients
are $2(m+1)/(m+2)!$ for $m\ge1$.  Rationality follows because $\binom{-m/2}{k}\in\Rat$ and
$\OB_{m-1,k}$ is a polynomial with rational coefficients in rational arguments.  Finally
$\mathcal C=1-v=1-\sum_m\omega_m\varsigma^m =1-\sum_m\omega_m2^{m/2}\upsilon^{m/2}$, which is
\cref{p1:eq:cayley-puiseux}.
\end{proof}

\begin{remark}[Two normalisations, one object]
\label{p1:rem:omega-vs-c}
\textsf{S2} works with $\omega_m$ (expanding in $\varsigma=\sqrt{2\upsilon}$) while
\textsf{S1} and \textsf{S3} work with $c_r$ (expanding in $\upsilon^{1/2}$).  The dictionary
is $c_r=-2^{r/2}\omega_r$, verified here for $1\le r\le18$ by exact rational arithmetic.  The
$\omega_m$ are rational, the $c_r$ alternate between rationals and rational multiples of
$\sqrt2$; this is the only reason to prefer $\omega$ for tables.
\end{remark}

\begin{table}[htbp]
\centering
\caption{Exact universal Cayley coefficients $\omega_m$ of
\cref{p1:eq:omega-master}, recomputed by rational arithmetic from
\cref{p1:eq:Phi}.}
\label{p1:tab:omega}
\footnotesize
\setlength{\tabcolsep}{5pt}
\begin{tabular}{r@{\quad}l@{\qquad}r@{\quad}l}
\toprule
$m$ & $\omega_m$ & $m$ & $\omega_m$\\
\midrule
$1$ & $1$ & $10$ & $-5776369/1515591000$\\
$2$ & $-1/3$ & $11$ & $169709463197/69528040243200$\\
$3$ & $11/72$ & $12$ & $-1118511313/709296588000$\\
$4$ & $-43/540$ & $13$ & $667874164916771/650782456676352000$\\
$5$ & $769/17280$ & $14$ & $-500525573/744761417400$\\
$6$ & $-221/8505$ & $15$ & $103663334225097487/234281684403486720000$\\
$7$ & $680863/43545600$ & $16$ & $-466901817532379/1595278956070800000$\\
$8$ & $-1963/204120$ & $17$ & $21235294185086305043/109242202556140093440000$\\
$9$ & $226287557/37623398400$ & $18$ & $-106040742894306601/818378104464320400000$\\
\bottomrule
\end{tabular}
\end{table}

The corresponding $c_r$ are
\begin{equation}
\begin{aligned}
 c_1&=-\sqrt2, & c_2&=\tfrac23, &
 c_3&=-\tfrac{11\sqrt2}{36}, & c_4&=\tfrac{43}{135},\\
 c_5&=-\tfrac{769\sqrt2}{4320}, & c_6&=\tfrac{1768}{8505}, &
 c_7&=-\tfrac{680863\sqrt2}{5443200}, & c_8&=\tfrac{3926}{25515}.
\end{aligned}
\label{p1:eq:c-values}
\end{equation}

\begin{remark}[On the observed sign pattern]
\label{p1:rem:omega-signs}
All three siblings display $\omega_m$ with strictly alternating signs and none proves it.  The
data in \cref{p1:tab:omega} are consistent with $\sgn\omega_m=(-1)^{m+1}$ for $1\le m\le18$,
and this is what one expects from the fact that the inverse function $v(\varsigma)$ of
$\varsigma=v\Phi(v)^{1/2}$ has its nearest singularity on the negative $\varsigma$-axis; but
no proof is offered here and the statement is recorded as an observation, not as a result.
\end{remark}

\section{The Puiseux expansion at the Otter singularity}
\label{p1:sec:puiseux}

\subsection{The local disturbance series}

\begin{definition}[Local coordinate and disturbance]
\label{p1:def:local}
Write
\begin{equation}
 s:=1-\frac{z}{\rho},
 \qquad
 \theta(s):=1-\e\,\zeta\bigl(\rho(1-s)\bigr)=\sum_{m\ge1}\theta_ms^{m}.
 \label{p1:eq:theta-def}
\end{equation}
The constant term vanishes by \cref{p1:eq:critical}.  By \cref{p1:prop:R-analytic}, $\theta$
is analytic on the image of $|z|<\sqrt\rho$ under $z\mapsto s$, namely the disc
$|s-1|<\rho^{-1/2}$; in particular its Taylor series at $s=0$ converges for
$|s|<\rho^{-1/2}-1=0.71923\ldots$.
\end{definition}

Taylor's formula gives the derivative form
\begin{equation}
 \theta_m=\frac{(-1)^{m+1}\e\rho^{m}}{m!}\,\zeta^{(m)}(\rho),
 \qquad m\ge1,
 \label{p1:eq:theta-derivative}
\end{equation}
consistent with \cref{p1:eq:theta1}: $\theta_1=\e\rho\zeta'(\rho)$.  For computation the
following coefficient form is preferable, since it needs nothing but the integers $a_r$ and
the single scalar $\rho$.

\begin{lemma}[Local coefficients of $R$ from the tree numbers]
\label{p1:lem:Rhat}
Let $\widehat R_m:=\Coef{s^{m}}R\bigl(\rho(1-s)\bigr)$.  Then
\begin{equation}
 \boxed{\;
 \widehat R_m=(-1)^{m}\sum_{r\ge2}\varrho_r\binom{r}{m}\rho^{r}
 =(-1)^{m}\,\frac{\rho^{m}R^{(m)}(\rho)}{m!},
 \;}
 \label{p1:eq:Rhat}
\end{equation}
with $\varrho_r$ as in \cref{p1:eq:varrho}.  The series converges geometrically: for fixed $m$
the general term is $\bigO\bigl(r^{m}\rho^{r/2}\bigr)$.
\end{lemma}

\begin{proof}
By \cref{p1:lem:R-coefficients}, $R(\rho(1-s))=\sum_{r\ge2}\varrho_r\rho^{r}(1-s)^{r}$, and
$\Coef{s^m}(1-s)^r=(-1)^m\binom rm$.  The middle expression is a Taylor coefficient of $R$ at
$\rho$ in the variable $\rho s$, giving the right-hand form.  For the convergence rate, every
divisor $d<r$ of $r$ satisfies $d\le r/2$, so $\varrho_r\le a_{\lfloor r/2\rfloor}\,d(r)$ with
$d(r)$ the number of divisors of $r$; since $\limsup_na_n^{1/n}=\rho^{-1}$, for every
$\epsilon>0$ there is $K_\epsilon$ with $\varrho_r\rho^{r}\le
K_\epsilon(\sqrt\rho+\epsilon)^{r}$.  With $\binom rm\le r^{m}$ the claim follows.
\end{proof}

\begin{remark}[A dedup]
\label{p1:rem:single-sum}
\textsf{S1} (its $\ell_m$) and \textsf{S3} (its $\Coef{s^m}R(\rho(1-s))$) both give the
\emph{double} sum $(-1)^m\sum_{k\ge2}\frac1k\sum_{n\ge1}a_n\binom{kn}{m}\rho^{kn}$.
\Cref{p1:eq:Rhat} is the same number after the inner divisor collection of
\cref{p1:lem:R-coefficients}, and is one sum rather than two.  The two forms were checked
against each other numerically to $70$ digits for $0\le m\le30$.
\end{remark}

\begin{lemma}[The disturbance coefficients]
\label{p1:lem:theta}
With $\widehat R_m$ as in \cref{p1:lem:Rhat}, put $\widehat R_{+}(s):=\sum_{m\ge1}\widehat
R_ms^{m}$ and $E(s):=\exp\widehat R_{+}(s)=\sum_{m\ge0}E_ms^{m}$.  Then $E_0=1$,
\begin{equation}
 E_m=\sum_{k=0}^{m}\frac1{k!}\,\OB_{m,k}
 \bigl(\widehat R_1,\widehat R_2,\ldots\bigr)
 =\frac1m\sum_{k=1}^{m}k\,\widehat R_k\,E_{m-k},
 \label{p1:eq:E-bell}
\end{equation}
and
\begin{equation}
 \boxed{\;\theta_m=E_{m-1}-E_m,\qquad m\ge1.\;}
 \label{p1:eq:theta-bell}
\end{equation}
\end{lemma}

\begin{proof}
By \cref{p1:eq:critical}, $\e\zeta(\rho)=1$, so
\[
\e\,\zeta\bigl(\rho(1-s)\bigr) =\frac{\rho(1-s)\e^{R(\rho(1-s))}}{\rho\,\e^{R(\rho)}}
=(1-s)\,\e^{\widehat R_{+}(s)}=(1-s)E(s),
\]
because $\widehat R_0=R(\rho)$.  Hence $\theta(s)=1-(1-s)E(s)$ and $\theta_m=-(E_m-E_{m-1})$
for $m\ge1$, while the constant term is $1-E_0=0$. The two expressions for $E_m$ are the
ordinary Bell expansion of an exponential (\cref{p0:def:obell}) and the standard triangular
exponential recurrence obtained from $E'=\widehat R_{+}'E$.
\end{proof}

\begin{remark}[The jet route]
\label{p1:rem:jets}
\textsf{S2} computes $\theta_m$ instead from the \emph{scaled logarithmic jets}
\begin{equation}
 \ell_j:=\rho^{j}\left.\frac{\dd^{j}}{\dd z^{j}}\log\zeta(z)\right|_{z=\rho}
 =(-1)^{j-1}(j-1)!+\sum_{r\ge j}\varrho_r\,r^{\underline{j}}\,\rho^{r},
 \label{p1:eq:jets}
\end{equation}
via $\theta_m=\frac{(-1)^{m+1}}{m!}\EB_m(\ell_1,\ldots,\ell_m)$, where $\EB_m$ is the complete
exponential Bell polynomial and $r^{\underline{j}}=r(r-1)\cdots(r-j+1)$.  This is equivalent
to \cref{p1:lem:theta} through $\widehat R_m=(-1)^m\rho^mR^{(m)}(\rho)/m!$ and
$\ell_j=(-1)^{j}j!\,\widehat R_j$ for $j\ge1$ up to the $\log z$ term; both routes were
implemented and agree.  \Cref{p1:lem:theta} is preferred because it avoids the factorial
rescaling and its intermediate quantities stay $\bigO(1)$. \textsf{S2}'s normalisation
$s_j=2\theta_j$ (it expands $2(1-\e\zeta)$ rather than $1-\e\zeta$) accounts for every factor
of $2$ that differs between the siblings' formulas for $t_r$.
\end{remark}

The first disturbance coefficients, recomputed, are
\begin{equation}
\begin{aligned}
 \theta_1&=\phantom{-}1.216004561824048922157212619\ldots,&
 \theta_2&=-0.458415729856391719180847717\ldots,\\
 \theta_3&=\phantom{-}0.448104774389535043665734122\ldots,&
 \theta_4&=-0.399561679956209677406231980\ldots,\\
 \theta_5&=\phantom{-}0.385251341946613000719546396\ldots,&
 \theta_6&=-0.389806561959203341949164640\ldots .
\end{aligned}
\label{p1:eq:theta-numeric}
\end{equation}

\subsection{All local coefficients}

\begin{theorem}[Every Puiseux coefficient at $\rho$]
\label{p1:thm:t-coefficients}
Factor $\theta(s)=\theta_1s\bigl(1+U(s)\bigr)$ with $U(s)=\sum_{j\ge1}u_js^{j}$,
$u_j=\theta_{j+1}/\theta_1$.  Then in a slit neighbourhood of $z=\rho$ there is a convergent
Puiseux expansion
\begin{equation}
 T(z)=1+\sum_{n\ge1}t_n\,s^{n/2},
 \qquad s=1-\frac z\rho,
 \label{p1:eq:T-puiseux}
\end{equation}
and for every $n\ge1$
\begin{equation}
 \boxed{\;
 t_n=\sum_{\substack{1\le r\le n\\ r\equiv n\ (2)}}
 c_r\,\theta_1^{r/2}\,
 \Coef{s^{(n-r)/2}}\bigl(1+U(s)\bigr)^{r/2}
 =\sum_{\substack{1\le r\le n\\ r\equiv n\ (2)}}
 c_r\,\theta_1^{r/2}
 \sum_{k=0}^{(n-r)/2}\binom{r/2}{k}\OB_{\frac{n-r}2,k}(u_1,u_2,\ldots).\;}
 \label{p1:eq:t-master}
\end{equation}
Thus $t_n$ is a finite expression in $\omega_1,\ldots,\omega_n$, $\theta_1^{\pm1/2}$ and
$\theta_2,\ldots,\theta_{\lceil (n+1)/2\rceil}$.
\end{theorem}

\begin{proof}
Near $z=\rho$ we have $\e\zeta=1-\theta(s)$ with $\theta(0)=0$, so by
\cref{p1:eq:factorisation} and \cref{p1:eq:cayley-puiseux},
$T=1+\sum_{r\ge1}c_r\theta(s)^{r/2}$, the branch being the one with $T\to1^{-}$ as
$s\to0^{+}$.  Substituting $\theta(s)^{r/2}=\theta_1^{r/2}s^{r/2}(1+U(s))^{r/2}$, a term of
index $r$ can contribute to $s^{n/2}$ only when $n-r$ is a nonnegative even integer.
Extracting $\Coef{s^{(n-r)/2}}$ and applying \cref{p0:lem:power-log} to the unit series $1+U$
gives \cref{p1:eq:t-master}.  Convergence: $\theta$ is analytic and has a simple zero at $s=0$
with $\theta_1>0$ (\cref{p1:eq:theta1}), so $U$ is analytic with $U(0)=0$; choosing a slit
neighbourhood on which $|U|<1$ and a consistent square root, the convergent series
\cref{p1:eq:cayley-puiseux} may be composed with the analytic map $s\mapsto\theta(s)$.
\end{proof}

\begin{corollary}
\label{p1:cor:t-low}
$t_1=-\sqrt{2\theta_1}$, $t_2=\tfrac23\theta_1$, and
$t_3=-\tfrac{11\sqrt2}{36}\theta_1^{3/2}-\tfrac{\sqrt2}{2}\theta_2 \theta_1^{-1/2}$.
\end{corollary}

\begin{proof}
Take $n=1,2,3$ in \cref{p1:eq:t-master} with $c_1=-\sqrt2$, $c_2=\tfrac23$,
$c_3=-\tfrac{11\sqrt2}{36}$ and $\OB_{0,0}=1$, $\OB_{1,1}(u)=u_1$, $\binom{1/2}{1}=\tfrac12$.
\end{proof}

\begin{remark}[Sibling cross-check]
\label{p1:rem:t3-cross}
In \textsf{S2}'s normalisation $s_1=2\theta_1$, $s_2=2\theta_2$, the same coefficient reads
$t_3=-(11s_1^{2}+36s_2)/(72\sqrt{s_1})$; substituting $s_j=2\theta_j$ reproduces
\cref{p1:cor:t-low} exactly.  Likewise $t_1=-\sqrt{s_1}$ and $t_2=s_1/3$.  Both were checked
numerically.
\end{remark}

\begin{table}[htbp]
\centering
\caption{Puiseux coefficients in $T(z)=1+\sum_{n\ge1}t_n(1-z/\rho)^{n/2}$,
recomputed from \cref{p1:eq:t-master} at $90$ working digits.}
\label{p1:tab:t}
\small
\setlength{\tabcolsep}{6pt}
\begin{tabular}{r@{\quad}r@{\qquad}r@{\quad}r}
\toprule
$n$ & $t_n$ & $n$ & $t_n$\\
\midrule
$1$ & $-1.559490020374640885542$ & $10$ & $\phantom{-}0.205984831277907418736$\\
$2$ & $\phantom{-}0.810669707882699281438$ & $11$ & $-0.052724708498190562315$\\
$3$ & $-0.285487021612845605304$ & $12$ & $\phantom{-}0.017026568754953669446$\\
$4$ & $\phantom{-}0.165372365712083893477$ & $13$ & $-0.152370624366325396286$\\
$5$ & $-0.342459970402154201042$ & $14$ & $\phantom{-}0.173702883299850462894$\\
$6$ & $\phantom{-}0.317407225946528563568$ & $15$ & $-0.014473703739527044857$\\
$7$ & $-0.107778800291631009182$ & $16$ & $-0.021899517611215562233$\\
$8$ & $\phantom{-}0.061384957055835104689$ & $17$ & $-0.144547193570909705477$\\
$9$ & $-0.195212383597556463613$ & $18$ & $\phantom{-}0.176077108885017779062$\\
\bottomrule
\end{tabular}
\end{table}

\begin{remark}[Analytic versus singular local terms]
\label{p1:rem:even-odd}
The even part $\sum_jt_{2j}s^{j}$ of \cref{p1:eq:T-puiseux} is analytic at $z=\rho$.  A fixed
analytic monomial $(1-z/\rho)^{j}$ is a polynomial in $z$ and contributes nothing to
$\Coef{z^n}$ once $n>j$; hence only the odd coefficients $t_{2j+1}$ enter the algebraic
$\rho^{-n}n^{-\bullet}$ asymptotics of \cref{p1:sec:forward}.  This is not a discarding of
information: the analytic germ continues through $\rho$ and reappears in the exponentially
smaller sectors of \cref{p1:sec:sectors}.  The sign pattern of \cref{p1:tab:t} is not monotone
and should not be over-read; these are local analytic data, not the final coefficient
constants.
\end{remark}

In Otter's coordinate $\rho-z=\rho s$ the leading behaviour reads $T(z)=1-\mathsf
b\,(\rho-z)^{1/2}+\bigO(\rho-z)$ with
\begin{equation}
 \mathsf b=-\frac{t_1}{\sqrt\rho}=2.68112814726711223857732878\ldots .
 \label{p1:eq:otter-b}
\end{equation}

\begin{theorem}[Local expansion with an explicit remainder]
\label{p1:thm:local-remainder}
For every $N\ge1$ there are a dented (``$\Delta$-domain'') neighbourhood $\Delta$ of $\rho$
and a constant $K_N$ such that
\begin{equation}
 \left|T(z)-1-\sum_{n=1}^{N}t_n\left(1-\frac z\rho\right)^{n/2}\right|
 \le K_N\left|1-\frac z\rho\right|^{(N+1)/2},
 \qquad z\in\Delta .
 \label{p1:eq:local-remainder}
\end{equation}
\end{theorem}

\begin{proof}
By \cref{p1:prop:R-analytic} the map $s\mapsto\theta(s)$ is analytic on a disc
$|s|<\varepsilon_0$; by \cref{p1:eq:theta1} it has a simple zero at $s=0$ with $\theta_1>0$.
Shrinking $\varepsilon_0$, $|U(s)|<\tfrac12$ there. The series \cref{p1:eq:cayley-puiseux}
converges for $|\upsilon|$ small; a Cauchy estimate on the analytic function $\upsilon\mapsto
\mathcal C((1-\upsilon)/\e)$ restricted to a slit disc gives $|c_r|\le K\varepsilon_1^{-r/2}$.
Composing convergent expansions and using Taylor's theorem with remainder on $(1+U)^{r/2}$
gives \cref{p1:eq:local-remainder} on the slit disc $\{|s|<\varepsilon,\ |\arg s|<\pi-\eta\}$,
which is a $\Delta$-domain at $\rho$ in the $z$-variable.
\end{proof}

\section{Transfer to the all-orders coefficient asymptotics}
\label{p1:sec:forward}

\subsection{Exact transfer of a singular monomial}

For $\beta\notin\Nat_0$ and $n\ge0$,
\begin{equation}
 \Coef{z^{n}}\left(1-\frac z\rho\right)^{\beta}
 =\rho^{-n}(-1)^{n}\binom{\beta}{n}
 =\rho^{-n}\frac{\Gamma(n-\beta)}{\Gamma(-\beta)\Gamma(n+1)} .
 \label{p1:eq:exact-transfer}
\end{equation}
For $\beta\in\Nat_0$ the left-hand side vanishes for $n>\beta$ --- the analytic-germ statement
of \cref{p1:rem:even-odd}.  It remains to expand the gamma ratio to all orders.

\begin{lemma}[Gamma-ratio coefficients]
\label{p1:lem:gamma-ratio}
Let $B_m(x)$ be the Bernoulli polynomials, and for $\beta\in\Cplx$ set
\begin{equation}
 \varkappa_m(\beta):=\frac{(-1)^{m+1}}{m(m+1)}
 \bigl(B_{m+1}(-\beta)-B_{m+1}(1)\bigr),\qquad m\ge1,
 \label{p1:eq:varkappa}
\end{equation}
and define $g_r(\beta)$ by
$\exp\bigl(\sum_{m\ge1}\varkappa_m(\beta)v^{m}\bigr)=\sum_{r\ge0}g_r(\beta)v^{r}$,
equivalently
\begin{equation}
 g_0(\beta)=1,\qquad
 g_r(\beta)=\frac1r\sum_{m=1}^{r}m\,\varkappa_m(\beta)\,g_{r-m}(\beta)
 =\sum_{k=0}^{r}\frac1{k!}\OB_{r,k}
 \bigl(\varkappa_1(\beta),\varkappa_2(\beta),\ldots\bigr).
 \label{p1:eq:g-recurrence}
\end{equation}
Then, for fixed $\beta$ and $n\to\infty$ in any sector avoiding the negative real axis,
\begin{equation}
 \frac{\Gamma(n-\beta)}{\Gamma(n+1)}
 \sim n^{-\beta-1}\sum_{r\ge0}\frac{g_r(\beta)}{n^{r}} .
 \label{p1:eq:gamma-ratio}
\end{equation}
Each $g_r$ is a polynomial of degree $2r$ in $\beta$ with rational coefficients; the first
five are
\begin{align}
 g_0&=1,
 &g_1(\beta)&=\frac{\beta(\beta+1)}2,
 &g_2(\beta)&=\frac{\beta(\beta+1)(\beta+2)(3\beta+1)}{24},
 \notag\\
 g_3(\beta)&=\frac{\beta^{2}(\beta+1)^{2}(\beta+2)(\beta+3)}{48},
 \label{p1:eq:g-first}
\end{align}
\begin{equation}
 g_4(\beta)=\frac{\beta(\beta+1)(\beta+2)(\beta+3)(\beta+4)
 \bigl(15\beta^{3}+30\beta^{2}+5\beta-2\bigr)}{5760}.
 \label{p1:eq:g4}
\end{equation}
In particular $\bigl(g_r(\tfrac12)\bigr)_{r\ge0}
=1,\tfrac38,\tfrac{25}{128},\tfrac{105}{1024},
\tfrac{1659}{32768},\tfrac{6237}{262144},\ldots$
\end{lemma}

\begin{proof}
The Euler--Maclaurin/Bernoulli form of Stirling's expansion gives, for fixed $a,b$,
\[
\log\frac{\Gamma(n+a)}{\Gamma(n+b)\,n^{a-b}} \sim\sum_{m\ge1}\frac{(-1)^{m+1}}{m(m+1)}
\frac{B_{m+1}(a)-B_{m+1}(b)}{n^{m}}
\]
as $n\to\infty$ in the stated sectors.  Taking $a=-\beta$, $b=1$ yields
$\log\bigl(\Gamma(n-\beta)/\Gamma(n+1)\bigr)\sim-(\beta+1)\log n
+\sum_m\varkappa_m(\beta)n^{-m}$, and exponentiating an asymptotic series term by term is
legitimate, producing exactly the coefficients of $\exp\bigl(\sum\varkappa_mv^m\bigr)$.  The
recurrence follows from differentiating that generating identity; the Bell form is
\cref{p0:def:obell}. Degree: $\varkappa_m$ has degree $m+1$ in $\beta$ and vanishes at
$\beta=0$, so a monomial $\prod\varkappa_{m_i}$ with $\sum m_i=r$ has degree $\le
r+\#\{i\}\le2r$; the displayed $g_1,\ldots,g_4$ were recomputed here from
\cref{p1:eq:varkappa} by exact rational arithmetic.
\end{proof}

\begin{remark}[Part~0 candidate]
\label{p1:rem:gamma-part0}
\Cref{p1:lem:gamma-ratio} is not specific to rooted trees and is used verbatim wherever a
chapter transfers an algebraic singularity to coefficient asymptotics.  It is listed in the
report as a candidate for promotion into Part~0.  All three siblings state it; \textsf{S2} and
\textsf{S3} give the same recurrence, \textsf{S1} the same object under the name
$\gamma_m(\beta)$ with $\eta_r(\beta)$ for $\varkappa_r(\beta)$.
\end{remark}

\subsection{The master coefficient formula}

\begin{theorem}[All-orders dominant expansion of A000081]
\label{p1:thm:forward}
For every fixed $M\ge0$,
\begin{equation}
 \boxed{\;
 a_n=\frac{\rho^{-n}}{\sqrt{\pi}\,n^{3/2}}
 \left(\sum_{m=0}^{M}\frac{\tau_m}{n^{m}}+\bigO(n^{-M-1})\right),
 \qquad n\to\infty,\;}
 \label{p1:eq:forward}
\end{equation}
where
\begin{equation}
 \boxed{\;
 \tau_m=\sum_{j=0}^{m}\Delta_j\,t_{2j+1}\,
 g_{m-j}\!\left(j+\tfrac12\right),
 \qquad
 \Delta_j:=\frac{(-1)^{j+1}(2j+1)!!}{2^{j+1}}
 =\frac{\sqrt\pi}{\Gamma(-j-\tfrac12)} .\;}
 \label{p1:eq:tau-master}
\end{equation}
Every ingredient is a finite coefficient extraction: $t_{2j+1}$ by \cref{p1:eq:t-master},
$g_r$ by \cref{p1:eq:g-recurrence}, $\theta_m$ by \cref{p1:eq:theta-bell}, $\omega_m$ by
\cref{p1:eq:omega-master}, $\varrho_r$ by \cref{p1:eq:varrho}.
\end{theorem}

\begin{proof}
Fix $M$ and apply \cref{p1:thm:local-remainder} with $N=2M+2$: on a $\Delta$-domain at $\rho$,
\[
T(z)=A(z)+\sum_{j=0}^{M}t_{2j+1}\left(1-\frac z\rho\right)^{j+1/2} +\bigO\!\left(\left|1-\frac
z\rho\right|^{M+3/2}\right),
\]
with $A$ the polynomial collecting the even (analytic) terms.  By
\cref{p1:prop:unique-dominant} the only singularity on $|z|=\rho$ is $\rho$ itself, so the
Cauchy contour may be deformed into the standard Flajolet--Odlyzko $\Delta$-contour.  The
polynomial $A$ contributes nothing to $\Coef{z^n}$ for $n>\deg A$.  Each retained monomial
contributes exactly $\rho^{-n}t_{2j+1}\Gamma(n-j-\tfrac12)/
\bigl(\Gamma(-j-\tfrac12)\Gamma(n+1)\bigr)$ by \cref{p1:eq:exact-transfer}, whose asymptotic
expansion is \cref{p1:eq:gamma-ratio}; the $\bigO$-term transfers to
$\bigO(\rho^{-n}n^{-M-5/2})$ by the standard Hankel-contour estimate. Collecting the
coefficient of $\rho^{-n}n^{-m-3/2}$ from the pairs $j+r=m$ and multiplying through by
$\sqrt\pi$ gives \cref{p1:eq:tau-master}, the identity
$\Gamma(-j-\tfrac12)=(-4)^{j+1}(j+1)!\sqrt\pi/(2j+2)! =(-1)^{j+1}2^{j+1}\sqrt\pi/(2j+1)!!$
supplying $\Delta_j$.
\end{proof}

\begin{corollary}[Otter's constant]
\label{p1:cor:otter}
With $C:=\tau_0/\sqrt\pi$,
\begin{equation}
 \tau_0=-\frac{t_1}{2}=\sqrt{\frac{\theta_1}{2}},
 \qquad
 C=\sqrt{\frac{\theta_1}{2\pi}}=\sqrt{\frac{\e\rho\,\zeta'(\rho)}{2\pi}}
 =\sqrt{\frac{1+\rho R'(\rho)}{2\pi}},
 \qquad
 a_n\sim C\,\alpha^{n}n^{-3/2}.
 \label{p1:eq:otter}
\end{equation}
Numerically
\begin{equation}
 \tau_0=0.7797450101873204427711036979\ldots,\qquad
 C=0.4399240125710253040409033914345447648\ldots .
 \label{p1:eq:C-numeric}
\end{equation}
\end{corollary}

\begin{proof}
Set $m=0$ in \cref{p1:eq:tau-master}: $\Delta_0=-\tfrac12$, $g_0=1$, so
$\tau_0=-t_1/2=\sqrt{2\theta_1}/2=\sqrt{\theta_1/2}$ by \cref{p1:cor:t-low}. Divide by
$\sqrt\pi$ and use \cref{p1:eq:theta1}.
\end{proof}

The identity $C^{2}=\theta_1/(2\pi)$ is a genuinely independent check: the left side comes
through the whole Puiseux/transfer chain, the right side directly from one derivative of
$\zeta$.  The two agree to all $55$ digits carried.

Expanding \cref{p1:eq:tau-master} for small $m$ gives the transfer identities
\begin{equation}
\begin{aligned}
 \tau_0&=-\tfrac12t_1,\\
 \tau_1&=-\tfrac3{16}\bigl(t_1-4t_3\bigr)
 =-\tfrac3{16}t_1+\tfrac34t_3,\\
 \tau_2&=-\tfrac5{256}\bigl(5t_1-72t_3+96t_5\bigr)
 =-\tfrac{25}{256}t_1+\tfrac{45}{32}t_3-\tfrac{15}8t_5,\\
 \tau_3&=-\tfrac{105}{2048}\bigl(t_1-44t_3+160t_5-128t_7\bigr)
 =-\tfrac{105}{2048}t_1+\tfrac{1155}{512}t_3-\tfrac{525}{64}t_5
 +\tfrac{105}{16}t_7,\\
 \tau_4&=-\tfrac{21}{65536}
 \bigl(79t_1-10800t_3+81600t_5-161280t_7+92160t_9\bigr).
\end{aligned}
\label{p1:eq:tau-expanded}
\end{equation}
\textsf{S2} states the right-hand forms, \textsf{S3} the factored left-hand forms; they are
identical after clearing denominators, and both were re-derived here symbolically from
\cref{p1:eq:tau-master}.  For any $m$ beyond $4$ the compact convolution is far preferable to
the expanded form.

\begin{table}[htbp]
\centering
\caption{Coefficients of the dominant expansion
$a_n\sim\rho^{-n}(\pi n^{3})^{-1/2}\sum_{m\ge0}\tau_mn^{-m}$ and their
normalisations $f_m=\tau_m/\tau_0$, recomputed here.}
\label{p1:tab:tau}
\small
\setlength{\tabcolsep}{6pt}
\begin{tabular}{r@{\quad}r@{\qquad}r}
\toprule
$m$ & $\tau_m$ & $f_m=\tau_m/\tau_0$\\
\midrule
$0$ & $0.7797450101873204427711$ & $1$\\
$1$ & $0.07828911261061096206128$ & $0.1004034800964014347637$\\
$2$ & $0.3929402676631860184743$ & $0.5039343151023035331022$\\
$3$ & $1.537879315978838091102$ & $1.972284908382278469206$\\
$4$ & $8.200844090435596159220$ & $10.51734090413157399274$\\
$5$ & $57.29291473494343787740$ & $73.47647498402047106417$\\
$6$ & $503.0445050262735818248$ & $645.1397552456612928153$\\
$7$ & $5359.600933884326033403$ & $6873.530274463408619941$\\
$8$ & $67342.06920114653048341$ & $86364.21948371140797700$\\
$9$ & $975425.4970695924732213$ & $1250954.458605978293018$\\
$10$ & $15996932.49293173312961$ & $20515594.56480361813667$\\
$11$ & $292822531.3353392231146$ & ---\\
$12$ & $5914523441.293932519117$ & ---\\
\bottomrule
\end{tabular}
\end{table}

Higher values continue smoothly; for reference $\tau_{15}=8.078305401468884289\times10^{13}$,
$\tau_{18}=2.073828085209893749\times10^{18}$.

Writing
\begin{equation}
 F(v):=1+\sum_{m\ge1}f_mv^{m},\qquad f_m=\frac{\tau_m}{\tau_0},
 \label{p1:eq:F-def}
\end{equation}
the inversion-ready normal form is
\begin{equation}
 \boxed{\;a_n\sim C\,\e^{\lambda n}\,n^{-p}\,F(1/n),
 \qquad p=\tfrac32 .\;}
 \label{p1:eq:forward-normal}
\end{equation}

\begin{remark}[Divergence and least-term truncation]
\label{p1:rem:divergence}
\Cref{p1:eq:forward} asserts a Poincar\'e expansion at every fixed order; it does \emph{not}
assert convergence, and \cref{p1:tab:tau} shows the coefficients growing rapidly.  The optimal
truncation and the size of the resulting floor are governed by
\cref{p0:thm:optimal-truncation}; the relevant nearest competing exponential scale is
identified in \cref{p1:sec:sectors} as $\rho^{-n/2}$, so the least-term floor at index $n$ is
$\bigO\bigl(\rho^{n/2}\bigr)$ relative to the leading term --- consistent with the error
tables of \cref{p1:sub:numerics}, where the improvement stops at $n=20$ around $M=6$.
\textsf{S3} states this correctly as a numerical strategy rather than a theorem; \textsf{S1}
and \textsf{S2} call the late-term growth ``the perturbative shadow of farther
singularities'', which is the right heuristic but is not proved anywhere in the three drafts
and is not proved here.
\end{remark}

\section{The multiplicative reciprocal}
\label{p1:sec:reciprocal}

If ``inverse'' is read as the reciprocal sequence, the answer is immediate from
\cref{p1:eq:forward-normal} and requires no inversion theory at all.

\begin{proposition}[All-orders reciprocal]
\label{p1:prop:reciprocal}
Write $1/F(v)=1+\sum_{m\ge1}\bar f_mv^{m}$.  Then
\begin{equation}
 \bar f_m=\sum_{k=1}^{m}(-1)^{k}\,\OB_{m,k}(f_1,f_2,\ldots)
 =-\sum_{k=1}^{m}f_k\,\bar f_{m-k}\qquad(\bar f_0=1),
 \label{p1:eq:recip-coeff}
\end{equation}
and for every fixed $M\ge0$
\begin{equation}
 \frac1{a_n}=\frac{\sqrt\pi\,\rho^{n}n^{3/2}}{\tau_0}
 \left(\sum_{m=0}^{M}\frac{\bar f_m}{n^{m}}+\bigO(n^{-M-1})\right).
 \label{p1:eq:recip-asymptotic}
\end{equation}
\end{proposition}

\begin{proof}
The first equality in \cref{p1:eq:recip-coeff} is \cref{p0:lem:power-log} with exponent $-1$;
the second is the Cauchy product $F\cdot F^{-1}=1$.  For \cref{p1:eq:recip-asymptotic}, note
that $a_n>0$ and, by \cref{p1:eq:forward}, $a_n=C\alpha^nn^{-p}(1+\varepsilon_n)$ with
$\varepsilon_n\to0$; hence $1/a_n$ is defined and $1/(1+\varepsilon_n)$ has the asymptotic
expansion obtained by formally inverting the (asymptotic) series for $1+\varepsilon_n$ --- a
legitimate operation because the leading term is $1\ne0$ and each truncation error is
$\bigO(n^{-M-1})$.
\end{proof}

The first reciprocal coefficients, recomputed, are
\begin{equation}
\begin{aligned}
 \bar f_1&=-0.1004034800964014347637,&
 \bar f_2&=-0.4938534562868350540384,\\
 \bar f_3&=-1.872103543737176348103,&
 \bar f_4&=-9.882481221441551170916,\\
 \bar f_5&=-69.51082491336461883211,&
 \bar f_6&=-616.9168645978433816315 .
\end{aligned}
\label{p1:eq:recip-numeric}
\end{equation}

\begin{remark}[Reciprocal of the full transseries]
\label{p1:rem:recip-multisector}
Writing $a_n=\mathcal M_n(1+\mathcal E_n)$ with $\mathcal M_n$ the complete dominant sector
and $\mathcal E_n$ the sum of relative exponentially small sectors of \cref{p1:sec:sectors},
one has $1/a_n=\mathcal M_n^{-1}\sum_{q\ge0}(-\mathcal E_n)^{q}$; the products generate
exactly the additive semigroup of the singular actions, with coefficients given by finite
Cauchy products.  This is a formal statement about the transseries, and inherits whatever
status \cref{p1:thm:finite-level} has on the chosen continuation domain.
\end{remark}

\begin{warning}
\label{p1:rem:reciprocal}
The reciprocal is \emph{not} a functional inverse.  Its exponential scale is $\rho^{n}$, while
the functional inverse of \cref{p1:sec:inverse} grows only logarithmically in its argument.
All three siblings make this point; it is repeated here because the ambiguity of the word
``inverse'' is exactly what \cref{p0:def:three-inverses} exists to remove.
\end{warning}

\section{The asymptotic inverse: the Part~0 dictionary}
\label{p1:sec:inverse}

\subsection{The parameter dictionary}

Everything needed to invert \cref{p1:eq:forward-normal} has been proved in Part~0.  What this
chapter must supply is only the dictionary and the verification of Part~0's hypotheses.

\begin{proposition}[A000081 as an instance of the exponential--power model]
\label{p1:prop:dictionary}
Let
\begin{equation}
 \mathcal A(x):=C\,\e^{\lambda x}x^{-p}F(x^{-1}),
 \qquad p=\tfrac32 ,
 \label{p1:eq:interpolant}
\end{equation}
with $C$ from \cref{p1:cor:otter} and $F$ from \cref{p1:eq:F-def}.  Then $\mathcal A$ is an
instance of \cref{p0:def:model} under
\begin{equation}
 \boxed{\;
 \begin{array}{c|cccccc}
  \text{Part 0} & C & a & \alpha & b & \delta & Q(t)\\\hline
  \text{here} & C=\sqrt{\theta_1/(2\pi)} & \lambda=\log(1/\rho) & 1
  & -\tfrac32 & 1 & H(t)=\log F(t)
 \end{array}\;}
 \label{p1:eq:dictionary}
\end{equation}
with $H\in t\Real[[t]]$.  In particular $a=\lambda>0$ and $\alpha=1$, so $\mathcal A$ is
already in the linear-phase normal form and \cref{p0:lem:alpha-reduction} is the identity
substitution.
\end{proposition}

\begin{proof}
\Cref{p1:eq:forward-normal} is \cref{p0:def:model} with the stated substitutions; $\lambda>0$
because $\rho<1$ (\cref{p1:prop:critical}), $C>0$ because $\theta_1>0$, and $H(0)=\log
F(0)=\log1=0$ so $H\in t\Real[[t]]$.
\end{proof}

\begin{definition}[Additive phase coefficients]
\label{p1:def:phi}
Write $H(v)=\log F(v)=\sum_{m\ge1}\varphi_mv^{m}$.  By \cref{p0:lem:power-log},
\begin{equation}
 \varphi_m=\sum_{k=1}^{m}\frac{(-1)^{k-1}}{k}\OB_{m,k}(f_1,f_2,\ldots)
 =f_m-\frac1m\sum_{k=1}^{m-1}k\,\varphi_k\,f_{m-k}.
 \label{p1:eq:phi}
\end{equation}
\end{definition}

Taking logarithms in \cref{p1:eq:forward} gives the form used for inversion,
\begin{equation}
 \log a_n=\lambda n-\tfrac32\log n+\log C
 +\sum_{m=1}^{M}\frac{\varphi_m}{n^{m}}+\bigO(n^{-M-1}),
 \label{p1:eq:log-a}
\end{equation}
with, recomputed here,
\begin{equation}
\begin{aligned}
 \varphi_1&=0.1004034800964014347637,&
 \varphi_2&=0.4988938856945692935703,\\
 \varphi_3&=1.922025533841821821591,&
 \varphi_4&=10.19739642337376887053,\\
 \varphi_5&=71.47146706207073604412,&
 \varphi_6&=630.8599376073372478899 .
\end{aligned}
\label{p1:eq:phi-numeric}
\end{equation}

\subsection{The exact Lambert carrier}

\begin{corollary}[Carrier; specialisation of \cref{p0:thm:lambert-core}]
\label{p1:cor:carrier}
Let $y>y_\ast:=C\left(\e\lambda/p\right)^{p} =1.2108232025837572603\ldots$. Then the equation
$y=C\e^{\lambda x}x^{-p}$ has exactly one solution $x_0=x_0(y)>p/\lambda$, namely
\begin{equation}
 \boxed{\;
 x_0(y)=-\frac p\lambda\,
 \Wlam_{-1}\!\left(-\frac\lambda p\left(\frac Cy\right)^{1/p}\right)
 =-\frac{3}{2\lambda}\,
 \Wlam_{-1}\!\left(-\frac{2\lambda}{3}\left(\frac Cy\right)^{2/3}\right).\;}
 \label{p1:eq:carrier}
\end{equation}
\end{corollary}

\begin{proof}
Raise $y=C\e^{\lambda x}x^{-p}$ to the power $-1/p$ and multiply by $-\lambda/p$:
\[
-\frac\lambda p\left(\frac Cy\right)^{1/p} =-\frac{\lambda x}{p}\,\e^{-\lambda x/p}
=\xi\e^{\xi},\qquad \xi:=-\frac{\lambda x}{p}.
\]
Write $\epsilon:=-\frac\lambda p(C/y)^{1/p}$.  The map $x\mapsto C\e^{\lambda x}x^{-p}$ is
strictly increasing on $(p/\lambda,\infty)$ (its logarithmic derivative is $\lambda-p/x>0$
there) with infimum $y_\ast$ attained as $x\downarrow p/\lambda$, so $y>y_\ast$ gives a unique
root $x_0>p/\lambda$, i.e.\ $\xi<-1$; and $y>y_\ast$ is exactly $\epsilon\in(-\e^{-1},0)$.  On
that range $\xi\mapsto\xi\e^{\xi}$ has the two real inverse branches $\Wlam_0\in(-1,0)$ and
$\Wlam_{-1}<-1$, and the required root is the second. This is \cref{p0:thm:lambert-core} with
$L=\log(y/C)$, $a=\lambda$, $b=-p$ and branch index $k=-1$; the branch rule there selects
$k=-1$ under exactly the condition $x>-b/a$.
\end{proof}

\begin{remark}
The principal branch $\Wlam_0$ returns the other, small root of the same equation, which is
the spurious solution: all three siblings say so, and the threshold $y_\ast$ makes the
dichotomy quantitative.  For A000081 the threshold lies below $a_3=2$, so \cref{p1:eq:carrier}
may be used at every sequence value from $n=3$ on.
\end{remark}

\subsection{All algebraic corrections}

\begin{theorem}[Corrections; specialisation of
\cref{p0:thm:lambert-centered}]
\label{p1:thm:d}
Let
\begin{equation}
 \Psi(t,u):=\frac p\lambda\log(1+tu)
 -\frac1\lambda H\!\left(\frac{t}{1+tu}\right)
 \in t\Real[[t,u]] .
 \label{p1:eq:Psi}
\end{equation}
Then there is a unique $D(t)=\sum_{m\ge1}d_mt^{m}\in t\Real[[t]]$ with
$D(t)=\Psi\bigl(t,D(t)\bigr)$, its coefficients are given by \cref{p0:lem:lagrange-fp},
\begin{equation}
 \boxed{\;
 d_m=\sum_{k=1}^{m}\frac1k\,\Coef{t^{m}u^{k-1}}\Psi(t,u)^{k},\;}
 \label{p1:eq:d-master}
\end{equation}
equivalently by the triangular extraction
\begin{equation}
 d_m=\Coef{t^{m}}\left\{
 \frac p\lambda\log\bigl(1+tD_{<m}(t)\bigr)
 -\frac1\lambda H\!\left(\frac{t}{1+tD_{<m}(t)}\right)\right\},
 \qquad D_{<m}=\sum_{j<m}d_jt^{j},
 \label{p1:eq:d-triangular}
\end{equation}
or by the completely explicit ordinary-Bell form
\begin{equation}
 d_m=\frac p\lambda\sum_{r=1}^{m}\frac{(-1)^{r-1}}{r}
 \OB_{m-r,r}(d_1,d_2,\ldots)
 -\frac1\lambda\sum_{k=1}^{m}\varphi_k
 \sum_{r=0}^{m-k}(-1)^{r}\binom{k+r-1}{r}
 \OB_{m-k-r,r}(d_1,d_2,\ldots),
 \label{p1:eq:d-bell}
\end{equation}
in which the index constraints force only $d_1,\ldots,d_{m-1}$ to occur on the right.  With
$x_0$ from \cref{p1:eq:carrier} the complete algebraic inverse is
\begin{equation}
 \boxed{\;\nu(y)\sim x_0(y)+\sum_{m\ge1}\frac{d_m}{x_0(y)^{m}} .\;}
 \label{p1:eq:inverse-series}
\end{equation}
\end{theorem}

\begin{proof}
\Cref{p0:thm:lambert-centered} states exactly this for a model
$C\e^{ax}x^{b}\e^{Q(x^{-\delta})}$ with $\Psi(t,u)
=-\frac1a\bigl[b\log(1+tu)+Q\bigl(t/(1+tu)\bigr)\bigr]$; substituting the dictionary
\cref{p1:eq:dictionary} ($a=\lambda$, $b=-p$, $\delta=1$, $Q=H$) gives \cref{p1:eq:Psi}.
Concretely: putting $x=x_0+D$ into $\log(y/C)=\lambda x-p\log x+H(x^{-1})$ and subtracting the
carrier equation $\log(y/C)=\lambda x_0-p\log x_0$ gives, with $t=x_0^{-1}$,
\[
\lambda D-p\log(1+tD)+H\!\left(\frac{t}{1+tD}\right)=0,
\]
which is $D=\Psi(t,D)$.  Since $\Psi\in t\Real[[t,u]]$, the formal implicit-function theorem
gives a unique solution in $t\Real[[t]]$; \cref{p0:lem:lagrange-fp} converts it into
\cref{p1:eq:d-master}, and \cref{p1:eq:d-triangular} is the same statement read as a
recursion. \Cref{p1:eq:d-bell} follows by expanding
$\log(1+tD)=\sum_r\frac{(-1)^{r-1}}{r}t^rD^r$ and
$H\bigl(t/(1+tD)\bigr)=\sum_k\varphi_kt^k\sum_r(-1)^r\binom{k+r-1}{r}t^rD^r$ and applying
\cref{p0:def:obell} to the powers $D^r$.
\end{proof}

\begin{repairbox}
\textbf{A formal-to-analytic slide in \textsf{S1}.}  \textsf{S1}'s ``General coefficient
formula for the inverse'' is stated as $x(y)\sim X(y)+\sum_md_mX(y)^{-m}$, an asymptotic
assertion, but its proof ends with ``substitution \dots\ proves formal inversion order by
order'' and supplies no remainder.  As stated, \cref{p1:thm:d} is a theorem about formal
series only; the asymptotic reading requires \cref{p1:thm:inverse-remainder} below, whose
argument is \textsf{S3}'s --- the only one of the three drafts that separates the two
statements.  The volume keeps them separate.
\end{repairbox}

\begin{theorem}[Remainder; specialisation of
\cref{p0:thm:remainder-transport}]
\label{p1:thm:inverse-remainder}
Fix $M\ge1$ and set $\nu_M(y):=x_0(y)+\sum_{m=1}^{M-1}d_mx_0(y)^{-m}$.  Let $f(x)=\log
C+\lambda x-p\log x+H(x^{-1})$ be the (asymptotic) logarithm of the forward interpolant.  Then
\begin{equation}
 f\bigl(\nu_M(y)\bigr)-\log y=\bigO\bigl(x_0(y)^{-M}\bigr),
 \qquad
 \nu(y)-\nu_M(y)=\bigO\bigl(x_0(y)^{-M}\bigr)=
 \bigO\bigl((\log y)^{-M}\bigr).
 \label{p1:eq:inverse-remainder}
\end{equation}
\end{theorem}

\begin{proof}
By \cref{p1:thm:d} the coefficients of $t^{1},\ldots,t^{M-1}$ in $D-\Psi(t,D)$ vanish, so
truncating leaves a residual $\bigO(t^{M})$; this is the first claim. Since
$f'(x)=\lambda-p/x+\bigO(x^{-2})$ is bounded away from $0$ for large $x$,
\cref{p0:thm:backward-error} converts the residual into a forward error of the same order
divided by $f'$; this is \cref{p0:thm:remainder-transport} with $a=\lambda$.  Finally
$x_0(y)=\lambda^{-1}\log y\,(1+\littleo(1))$ by \cref{p1:eq:carrier}.
\end{proof}

\subsection{Coefficients}

Expanding \cref{p1:eq:d-triangular} through $t^{5}$ --- and verifying the result symbolically
--- gives the closed forms
\begin{align}
 d_1&=-\frac{\varphi_1}{\lambda},
 \label{p1:eq:d1}\\
 d_2&=-\frac{p\varphi_1+\lambda\varphi_2}{\lambda^{2}},
 \label{p1:eq:d2}\\
 d_3&=-\frac{\lambda\varphi_1^{2}+p^{2}\varphi_1+p\lambda\varphi_2
 +\lambda^{2}\varphi_3}{\lambda^{3}},
 \label{p1:eq:d3}\\
 d_4&=-\frac{1}{2\lambda^{4}}
 \Bigl(5p\lambda\varphi_1^{2}+6\lambda^{2}\varphi_1\varphi_2
 +2p^{3}\varphi_1+2p^{2}\lambda\varphi_2+2p\lambda^{2}\varphi_3
 +2\lambda^{3}\varphi_4\Bigr),
 \label{p1:eq:d4}\\
 d_5&=-\frac{1}{2\lambda^{5}}
 \Bigl(4\lambda^{2}\varphi_1^{3}+9p^{2}\lambda\varphi_1^{2}
 +14p\lambda^{2}\varphi_1\varphi_2+8\lambda^{3}\varphi_1\varphi_3
 +4\lambda^{3}\varphi_2^{2}\notag\\
 &\hspace{16mm}
 +2p^{4}\varphi_1+2p^{3}\lambda\varphi_2+2p^{2}\lambda^{2}\varphi_3
 +2p\lambda^{3}\varphi_4+2\lambda^{4}\varphi_5\Bigr).
 \label{p1:eq:d5}
\end{align}
Beyond $m=5$ the expanded expressions grow fast, and \cref{p1:eq:d-master,p1:eq:d-triangular}
are both shorter and less error-prone.

\begin{table}[htbp]
\centering
\caption{Additive phase coefficients $\varphi_m$ and Lambert-carrier
corrections $d_m$, recomputed here.  The last column is
\textsf{S2}'s normalisation $\mathrm D_m=\lambda^{m+1}d_m$
(see \cref{p1:rem:normalisations}).}
\label{p1:tab:inverse}
\small
\setlength{\tabcolsep}{6pt}
\begin{tabular}{r@{\quad}r@{\qquad}r@{\qquad}r}
\toprule
$m$ & $\varphi_m$ & $d_m$ & $\mathrm D_m=\lambda^{m+1}d_m$\\
\midrule
$1$ & $0.1004034800964014348$ & $-0.09264385352561312984$ & $-0.1088130343442745145$\\
$2$ & $0.4988938856945692936$ & $-0.5885630399313472628$ & $-0.7491856512881932495$\\
$3$ & $1.922025533841821822$ & $-2.596680167407531021$ & $-3.582177326832277789$\\
$4$ & $10.19739642337376887$ & $-13.14905339996706569$ & $-19.65872121138776004$\\
$5$ & $71.47146706207073604$ & $-85.49870151462556652$ & $-138.5327478518560861$\\
$6$ & $630.8599376073372479$ & $-711.2629148329465609$ & $-1248.979326371692854$\\
$7$ & $6753.629059387619229$ & $-7306.788270284967257$ & $-13905.40884474814835$\\
$8$ & $85171.04383262353790$ & $-89561.76088263724860$ & $-184719.1911265080905$\\
$9$ & $1236987.135642947528$ & $-1274829.503083708903$ & ---\\
$10$ & $20325777.10708850337$ & $-20641099.00032138544$ & ---\\
\bottomrule
\end{tabular}
\end{table}

\begin{repairbox}
\textbf{Numerical error in \textsf{S2}'s audit appendix.}  \textsf{S2}'s appendix table of
$\eta_j$ and $D_j$ reports $D_3=-3.582177326832123\ldots$, disagreeing with its own main text
($-3.582177326832277789103009844$) from the fourteenth significant digit. Recomputation
confirms the main-text value, $\mathrm D_3=-3.58217732683227778910301\ldots$; the appendix
entry is wrong. Every other entry of that table was checked and is correct to the digits
shown.
\end{repairbox}

\begin{remark}[Three normalisations of the same inverse]
\label{p1:rem:normalisations}
The siblings work in three different variables, which is why their coefficient lists look
different:
\begin{itemize}
\item \textsf{S1} and \textsf{S3} expand in the index variable $x$ itself and
      obtain the $d_m$ of \cref{p1:tab:inverse}; their carriers agree with
      \cref{p1:eq:carrier}.
\item \textsf{S2} rescales to $X=\lambda x$, so that its carrier is
      $X_0=-p\Wlam_{-1}\bigl(-\tfrac1p\e^{-Z/p}\bigr)$ with
      $Z=\log\bigl(y/(C\lambda^{p})\bigr)$, its corrections are
      $\mathrm D_m=\lambda^{m+1}d_m$, and its phase coefficients are
      $\eta_m^{\textsf{S2}}=\lambda^{m}\varphi_m$.  Both dictionaries were
      verified numerically for $1\le m\le8$, and $X_0=\lambda x_0$ holds
      identically because the two carrier equations differ by that
      substitution.
\item \textsf{S3} also uses $x$ but writes the polynomial--logarithmic form
      in terms of $\mathsf q=p/\lambda$ and
      $\eta_m^{\textsf{S3}}=\varphi_m/\lambda$.
\end{itemize}
The constant appearing in \textsf{S2}'s $Z$ is
$C\lambda^{p}=0.4963361416595997474727369464\ldots$.
\end{remark}

\subsection{Polynomial--logarithmic form}

Eliminating the Lambert function gives the expansion in the canonical logarithmic scale.  That
is \cref{p0:thm:flattening}; only the dictionary is recorded here.

\begin{corollary}[Flattened inverse]
\label{p1:cor:flattening}
Put
\begin{equation}
 \mathcal X:=\frac1\lambda\log\frac yC,\qquad
 L:=\log\mathcal X,\qquad
 \mathsf q:=\frac p\lambda,\qquad
 \eta_k:=\frac{\varphi_k}{\lambda} .
 \label{p1:eq:flatten-dict}
\end{equation}
Then the inverse problem is $x-\mathsf q\log x+\sum_{k\ge1}\eta_kx^{-k} =\mathcal X$, and
\cref{p0:thm:flattening} gives
\begin{equation}
 \boxed{\;
 \nu(y)\sim\mathcal X+\mathsf qL
 +\sum_{m\ge1}\frac{P_m(L)}{\mathcal X^{m}},\;}
 \label{p1:eq:flattened}
\end{equation}
where each $P_m$ is a polynomial of degree $m$ over $\Rat[\mathsf q,\eta_1,\ldots,\eta_m]$
with
\begin{equation}
 \Coef{L^{m}}P_m(L)=\frac{(-1)^{m+1}}{m}\mathsf q^{m+1},
 \qquad P_m(0)=d_m .
 \label{p1:eq:P-shape}
\end{equation}
\end{corollary}

The identity $P_m(0)=d_m$ is not an accident: setting $L=0$ in the flattening recursion
reduces it to the fixed-point equation solved by the Lambert carrier, and formal uniqueness in
\cref{p0:thm:lambert-centered} identifies the two.  The first two polynomials are
\begin{align}
 P_1(L)&=\mathsf q^{2}L-\eta_1
 =1.9156590172195452616\,L-0.0926438535256131298,
 \label{p1:eq:P1}\\
 P_2(L)&=-\tfrac12\mathsf q^{3}L^{2}
 +\bigl(\mathsf q^{3}+\mathsf q\eta_1\bigr)L
 -\mathsf q\eta_1-\eta_2\notag\\
 &=-1.3257062894576032164\,L^{2}
 +2.7896427837264004647\,L
 -0.5885630399313472628 .
 \label{p1:eq:P2}
\end{align}

\begin{remark}[The two flattenings are not the same polynomials]
\label{p1:rem:two-flattenings}
\textsf{S2} flattens in $Z=\log\bigl(y/(C\lambda^{p})\bigr)$ and $\log Z$; \textsf{S3}
flattens in $\mathcal X=\lambda^{-1}\log(y/C)$ and $\log\mathcal X$. The two recursions are
formally the \emph{same} recursion with different parameters --- \textsf{S3}'s is
\textsf{S2}'s under $(p,\eta^{\textsf{S2}}_k)\mapsto(\mathsf q,\eta^{\textsf{S3}}_k)$ --- but
they expand the same function in \emph{different} variables, so their coefficient polynomials
genuinely differ.  Neither draft is in error; the caution is only that the two lists must not
be compared term by term.  We adopt \textsf{S3}'s normalisation because it keeps the index
variable $x$ throughout and so matches \cref{p1:eq:inverse-series}.
\end{remark}

Retaining only the first two scales,
\begin{equation}
 \nu(y)=\frac{\log y}{\lambda}+\frac p\lambda\log\log y
 -\frac1\lambda\log\bigl(C\lambda^{p}\bigr)+\littleo(1),
 \label{p1:eq:inverse-elementary}
\end{equation}
that is
\begin{equation}
 \nu(y)=0.9227155616186015248\log y
 +1.3840733424279022872\log\log y
 +0.6463639827002088073+\littleo(1).
 \label{p1:eq:inverse-elementary-numeric}
\end{equation}
\Cref{p1:eq:inverse-series} is preferable numerically because $\Wlam_{-1}$ resums the entire
logarithmic tower exactly before any truncation is made.

\section{The discrete inverse}
\label{p1:sec:discrete}

The objects distinguished in \cref{p0:def:three-inverses} are, for this sequence: the smooth
inverse $\nu$ of \cref{p1:sec:inverse}; the integer staircase
\begin{equation}
 N(y):=\min\{n\ge1: a_n\ge y\};
 \label{p1:eq:staircase}
\end{equation}
and the round-to-nearest generalized inverse.  Applying \cref{p0:thm:staircase} requires its
separation hypothesis, which for A000081 is supplied by the next two lemmas.

\begin{lemma}[Strict monotonicity]
\label{p1:lem:monotone}
$a_1=a_2=1$ and $a_{n+1}>a_n$ for every $n\ge2$.
\end{lemma}

\begin{proof}
Monotonicity: the map sending a rooted tree $\mathfrak t$ on $n$ vertices to the rooted tree
whose root has the single child subtree $\mathfrak t$ is an injection from the $n$-vertex
shapes into the $(n+1)$-vertex shapes, so $a_{n+1}\ge a_n$.  Strictness for $n\ge2$: when
$n+1\ge3$, the star $K_{1,n}$ rooted at its centre has $n\ge2$ children and is therefore not
in the image of that injection, whence $a_{n+1}\ge a_n+1$.  The values $a_1=a_2=1$ are
immediate.
\end{proof}

\begin{lemma}[Separation]
\label{p1:lem:separation}
$a_{n+1}/a_n\to\alpha=2.9557\ldots>1$; more precisely
\begin{equation}
 \log\frac{a_{n+1}}{a_n}
 =\lambda-\frac{3}{2n}+\frac{\tfrac34-\varphi_1}{n^{2}}
 +\frac{-\tfrac12+\varphi_1-2\varphi_2}{n^{3}}+\bigO(n^{-4}).
 \label{p1:eq:ratio}
\end{equation}
\end{lemma}

\begin{proof}
Subtract \cref{p1:eq:log-a} at $n$ from the same relation at $n+1$, and expand $\log(n+1)=\log
n+n^{-1}-\tfrac12n^{-2}+\tfrac13n^{-3}-\cdots$ and $(n+1)^{-m}=n^{-m}(1-mn^{-1}+\cdots)$.
Collecting powers of $n^{-1}$ gives \cref{p1:eq:ratio}; the limit statement is its leading
term.
\end{proof}

\begin{corollary}[Staircase recovery]
\label{p1:cor:staircase}
Let $\nu_M$ be as in \cref{p1:thm:inverse-remainder}.  For every $M\ge2$ and all sufficiently
large $n$, $\bigl|\nu_M(a_n)-n\bigr|<\tfrac12$, so $n$ is recovered from $a_n$ by rounding.
For a general real target $y>1$, $N(y)=\nu(y)+\bigO(1)$, and certifying $N(y)$ exactly
requires the two-sided test $a_{m-1}<y\le a_m$, which \cref{p1:eq:inverse-series} reduces to a
bounded neighbourhood of $m\approx\nu(y)$.
\end{corollary}

\begin{proof}
The first claim is \cref{p1:thm:inverse-remainder} with $M\ge2$, since $\bigO\bigl((\log
a_n)^{-2}\bigr)=\bigO(n^{-2})$.  For the second, apply \cref{p0:thm:staircase}, whose
separation hypothesis holds by \cref{p1:lem:monotone,p1:lem:separation}: the logarithmic
derivative of the interpolation tends to $\lambda>0$, so multiplying $y$ by a bounded factor
moves $\nu(y)$ by a bounded amount, while consecutive sequence values differ by the
asymptotically fixed factor $\alpha$.  The final sentence is the definition of $N$.
\end{proof}

\begin{warning}[The $\bigO(1)$ rounding layer is not a small correction]
\label{p1:rem:sawtooth}
An assertion of the form $N(y)=\nu(y)+\littleo(1)$ \emph{uniformly in real $y$} is false, and
no descending power--logarithmic series can be true uniformly: arbitrarily small changes of
$y$ across a value $a_n$ change $N(y)$ by one. Once a monotone interpolation $\mathcal
A_{\mathrm{ex}}$ with $\mathcal A_{\mathrm{ex}}(n)=a_n$ has been fixed, one has exactly
$N(y)=\lceil\nu_{\mathrm{ex}}(y)\rceil$ with $\nu_{\mathrm{ex}}$ its ordinary inverse, and for
non-integral $\chi$
\begin{equation}
 \lceil\chi\rceil=\chi+\frac12
 +\frac1\pi\sum_{m\ge1}\frac{\sin(2\pi m\chi)}{m},
 \label{p1:eq:sawtooth}
\end{equation}
an $\bigO(1)$ periodic layer that never tends to zero.  The correct statements are: about
$\nu$; or at the special values $y=a_n$, where $\nu(a_n)=n+\bigO(n^{-M})$ for every fixed $M$;
or about a rounded exact interpolation.  \textsf{S2} states this sharply and \textsf{S3}
restates it; \textsf{S1}'s phrase ``can never be resolved more accurately than one unit'' is
the same point stated loosely, and is corrected here.  The layer \cref{p1:eq:sawtooth} must
not be confused with the exponentially small oscillation of \cref{p1:sec:sectors}: the latter
is genuinely $\littleo(1)$ and has an entirely different origin.
\end{warning}

\begin{remark}[A definitional disagreement]
\label{p1:rem:staircase-index}
\textsf{S1} and \textsf{S2} define the staircase as $\min\{n\ge2:a_n\ge y\}$, \textsf{S3} as
$\min\{n\ge1:a_n\ge y\}$.  The two differ only at $y\le1$, where the tie $a_1=a_2=1$ makes the
choice a pure convention.  This volume uses \cref{p1:eq:staircase} and confines all statements
to $y>1$, where the definitions agree.
\end{remark}

\section{Exponentially small sectors}
\label{p1:sec:sectors}

\subsection{Why one exponential scale is not the whole story}

\Cref{p1:eq:forward} is complete in inverse powers of $n$ at the scale $\rho^{-n}$.  It cannot
see a term $\sigma^{-n}$ with $|\sigma|>\rho$, because $(\rho/|\sigma|)^{n}$ decays faster
than every power of $n^{-1}$. Such terms come from singularities of the analytic continuation
of $T$, and the functional equation exposes exactly two mechanisms for producing them:
\begin{enumerate}
\item \textbf{inherited preimages}: if $T$ is singular at $\tau$ and
      $\sigma^{k}=\tau$ for some $k\ge2$, the summand $T(z^{k})/k$ of $R$ can
      make $R$, hence $T$, singular at $\sigma$;
\item \textbf{new Cayley critical points}: on a continued sheet where $R$ is
      analytic, any $\sigma$ with $\zeta(\sigma)=\e^{-1}$ and
      $\zeta'(\sigma)\ne0$ is a square-root branch point of $T$ by
      \cref{p1:thm:omega}.
\end{enumerate}
Iterating the first mechanism from the dominant point produces the \emph{root-lift skeleton}
\begin{equation}
 \mathscr R:=\bigl\{\rho^{1/m}\e^{2\pi\ii j/m}\;:\;m\ge1,\ 0\le j<m\bigr\},
 \label{p1:eq:skeleton}
\end{equation}
whose radii increase to $1$ and whose arguments become dense.  Relative to the dominant sector
the level-$m$ points carry the exponential weight
\begin{equation}
 \sigma_{m,j}^{-n}=\rho^{-n/m}\e^{-2\pi\ii jn/m},
 \qquad\text{i.e.}\qquad
 \exp\!\left[-\lambda\left(1-\frac1m\right)n\right]
 \text{ relative to }\rho^{-n}.
 \label{p1:eq:skeleton-scale}
\end{equation}

\begin{warning}[What the skeleton is and is not]
\label{p1:rem:skeleton}
\Cref{p1:eq:skeleton} is a list of \emph{candidate} inherited germs generated by the
recursion.  It is not a theorem that every listed point is singular on a given sheet, that
every singularity is listed, or that a listed point is reachable along a prescribed path.
Continuation can also meet solutions of $\zeta(z)=\e^{-1}$ that are not in $\mathscr R$, and a
critical point can coalesce with an inherited germ (which, by \cref{p1:eq:critical-outer}
below, halves an exponent and can create quarter powers).
\end{warning}

\begin{remark}[A disagreement about the unit circle]
\label{p1:rem:natural-boundary}
\textsf{S1} and \textsf{S2} assert that $|z|=1$ is a natural boundary for $T$ and that the
continued singularities accumulate at every point of it, citing Drmota--Gittenberger.
\textsf{S3} makes no such assertion and states only the candidate skeleton with the caveat of
\cref{p1:rem:skeleton}.  This volume records the natural-boundary claim as \emph{reported} and
does not use it: every statement in this section is conditioned on the explicit finite-radius
hypothesis $\mathsf H(\mathcal R)$ of \cref{p1:def:HR}, and nothing below depends on what
happens as $\mathcal R\uparrow1$.
\end{remark}

\subsection{Transfer from a general singular germ}

\begin{theorem}[Universal sector coefficients]
\label{p1:thm:general-sector}
Let $\sigma\ne0$ be an accessible singularity on a specified sheet, put $\mathsf
u=1-z/\sigma$, and suppose that in a $\Delta$-domain at $\sigma$
\begin{equation}
 T(z)=A_\sigma(\mathsf u)
 +\sum_{\beta\in E_\sigma}c_{\sigma,\beta}\,\mathsf u^{\beta}
 +\bigO\bigl(|\mathsf u|^{\beta_{\max}}\bigr),
 \label{p1:eq:general-germ}
\end{equation}
with $A_\sigma$ analytic, $E_\sigma\subset\Rat_{>0}\setminus\Nat_0$ finite and
$\beta_{\max}>\max E_\sigma$.  Then the contribution of \cref{p1:eq:general-germ} to
$\Coef{z^{n}}T$ is
\begin{equation}
 \boxed{\;
 \mathfrak S_\sigma(n)\sim\sigma^{-n}
 \sum_{\beta\in E_\sigma}\frac{c_{\sigma,\beta}}{\Gamma(-\beta)}
 \,n^{-\beta-1}\sum_{r\ge0}\frac{g_r(\beta)}{n^{r}},\;}
 \label{p1:eq:general-sector}
\end{equation}
so the coefficient of $\sigma^{-n}n^{-\beta-1-r}$ is exactly
$c_{\sigma,\beta}\,g_r(\beta)/\Gamma(-\beta)$.
\end{theorem}

\begin{proof}
Apply \cref{p1:eq:exact-transfer} with $\rho$ replaced by $\sigma$ to each displayed monomial
and expand the gamma ratio by \cref{p1:lem:gamma-ratio}.  The analytic part $A_\sigma$ has no
Hankel jump and contributes no large-$n$ tail; the $\bigO$-term transfers by the usual Hankel
estimate.
\end{proof}

\Cref{p1:eq:general-sector} is the exact analogue of \cref{p1:eq:tau-master} and allows
arbitrary ramification exponents, not only half-integers.  Relative to the dominant sector one
introduces the \emph{action}
\begin{equation}
 \Omega_\sigma:=\operatorname{Log}\frac{\sigma}{\rho},
 \qquad\text{so that}\qquad
 \sigma^{-n}=\rho^{-n}\e^{-n\Omega_\sigma}.
 \label{p1:eq:action}
\end{equation}

\begin{definition}[Continuation hypothesis $\mathsf H(\mathcal R)$]
\label{p1:def:HR}
Fix $\rho<\mathcal R<1$.  We say $\mathsf H(\mathcal R)$ holds for a chosen continuation of
$T$ if that continuation is analytic on $|z|<\mathcal R$ after removing finitely many pairwise
disjoint cuts issuing from a finite set $\Sigma_{\mathcal R}$; every
$\sigma\in\Sigma_{\mathcal R}$ carries a germ of the form \cref{p1:eq:general-germ} in a
$\Delta$-domain; and the deformed coefficient contour may be taken on a circle of radius at
least $\mathcal R$, up to harmless indentations, on which $T$ has at most polynomial growth.
\end{definition}

\begin{theorem}[Finite-level coefficient transseries]
\label{p1:thm:finite-level}
Assume $\mathsf H(\mathcal R)$ and let $\mathfrak S_\sigma$ be the multiplier-weighted sector
of \cref{p1:eq:general-sector} attached to the local Hankel cycle actually traversed.  Then,
for any finite truncation of the exponent and inverse-power sums,
\begin{equation}
 a_n=\sum_{\sigma\in\Sigma_{\mathcal R}}
 \mathfrak S_\sigma(n)+\bigO\bigl(\mathcal R^{-n}n^{K}\bigr)
 \label{p1:eq:finite-level}
\end{equation}
for a finite $K$ determined by the boundary growth.
\end{theorem}

\begin{proof}
Deform the Cauchy contour from a small circle about $0$ into the union of small Hankel
contours around the cuts issuing from $\Sigma_{\mathcal R}$ plus an outer contour of radius
$\mathcal R$.  On each Hankel contour apply \cref{p1:thm:general-sector}; the outer contour is
bounded by $\mathcal R^{-n}$ times the allowed polynomial factor.  Only finitely many
singularities and finitely many local terms are retained, so the deformation and the
term-by-term summation are legitimate.  Passing to a different sheet changes the homology
coefficients of the local cycles; those coefficients are the Stokes/intersection multipliers
absorbed into $\mathfrak S_\sigma$.
\end{proof}

\begin{repairbox}
\textbf{\textsf{S1}'s global transseries is not a theorem.}  \textsf{S1} displays
$a_n\sim\mathcal A_\rho(n)+\sum_{\sigma\in\mathfrak S\setminus \{\rho\}}S_\sigma\mathcal
A_\sigma(n)$ as a numbered displayed result over ``the singular set encountered by the chosen
contour continuation'', with no hypothesis attached, while conceding in its introduction that
globally the formula ``is best understood as a formal resurgent/Darboux transseries''. The
volume replaces it by \cref{p1:thm:finite-level}, which is \textsf{S3}'s formulation and the
only one of the three carrying an explicit hypothesis; the unrestricted sum over an infinite
$\mathfrak S$ is not asserted anywhere here. The correct reading of a ``global'' transseries
in the presence of accumulating singularities is the compatible family of the finite-radius
expansions \cref{p1:eq:finite-level} as $\mathcal R$ grows, which is exactly what
\textsf{S2}'s ``projective meaning'' remark says.
\end{repairbox}

\subsection{The local recursion produces every germ}

\begin{lemma}[Pull-back coordinates]
\label{p1:lem:pullback}
Let $z=\sigma(1-\mathsf u)$ and $\tau=\sigma^{k}$.  Then
\begin{equation}
 1-\frac{z^{k}}{\tau}=1-(1-\mathsf u)^{k}
 =k\,\mathsf u\left(1+\sum_{r\ge1}p_{k,r}\mathsf u^{r}\right),
 \qquad
 p_{k,r}=\frac{(-1)^{r}}{k}\binom{k}{r+1},
 \label{p1:eq:pullback}
\end{equation}
and consequently, for a parent monomial,
\begin{equation}
 \left(1-\frac{z^{k}}{\tau}\right)^{\beta}
 =(k\mathsf u)^{\beta}\sum_{N\ge0}\mathsf u^{N}
 \sum_{j=0}^{N}\binom{\beta}{j}
 \OB_{N,j}(p_{k,1},p_{k,2},\ldots).
 \label{p1:eq:pullback-coeff}
\end{equation}
\end{lemma}

\begin{proof}
$1-(1-\mathsf u)^{k}=-\sum_{r\ge1}\binom kr(-\mathsf u)^{r} =k\mathsf u-\sum_{r\ge2}\binom
kr(-1)^{r}\mathsf u^{r}$; dividing by $k\mathsf u$ and re-indexing $r\mapsto r+1$ gives
$p_{k,r}$. \Cref{p1:eq:pullback-coeff} is \cref{p0:lem:power-log} applied to the unit series
$1+\sum_rp_{k,r}\mathsf u^{r}$.
\end{proof}

\begin{theorem}[Recursive completeness of the local calculus]
\label{p1:thm:local-completeness}
Fix a sheet and a point $\sigma\ne0$, and choose a local uniformiser $z=\sigma(1-v^{\mathsf
L})$ in which every parent germ needed in $R(z)$ is an ordinary power series in $v$.  Then the
local germ of $T$ at $\sigma$ is computable to any finite order by the following finite
operations, and only these:
\begin{enumerate}
\item pull back each singular summand $T(z^{k})/k$ by
      \cref{p1:eq:pullback-coeff};
\item add the regular Taylor parts to obtain $R$, and exponentiate by
      \cref{p1:eq:E-bell} to obtain $\zeta=z\e^{R}$;
\item if $\zeta_0:=\zeta(\sigma)\ne\e^{-1}$ on the reached branch, so that
      the selected Cayley branch $\mathcal C_\flat$ is analytic at
      $\zeta_0$, compose regularly:
      \begin{equation}
       T=\sum_{r\ge0}\frac{\mathcal C_\flat^{(r)}(\zeta_0)}{r!}
       \bigl(\zeta-\zeta_0\bigr)^{r},
       \qquad
       \mathcal C_\flat'(w)=\frac{\mathcal C_\flat(w)}
       {w\bigl(1-\mathcal C_\flat(w)\bigr)},
       \label{p1:eq:regular-outer}
      \end{equation}
      the higher derivatives following by Fa\`a di Bruno;
\item if $\zeta_0=\e^{-1}$ and the reached branch has $T(\sigma)=1$, compose
      critically with the universal germ of \cref{p1:thm:omega}:
      \begin{equation}
       T=1+\sum_{r\ge1}c_r\,\theta_\sigma^{r/2},
       \qquad \theta_\sigma:=1-\e\zeta .
       \label{p1:eq:critical-outer}
      \end{equation}
\end{enumerate}
If $\theta_\sigma$ vanishes to order $\varsigma$ in $v$, then $\theta_\sigma^{1/2}$ begins at
$v^{\varsigma/2}$ and the ramification index doubles when $\varsigma$ is odd; this is the only
new ramification mechanism.
\end{theorem}

\begin{proof}
Induction on the depth of $\sigma$ in the pull-back graph.  At depth $0$,
\cref{p1:thm:t-coefficients} gives the dominant germ.  At the induction step, $z\mapsto z^{k}$
converts each already-known germ into a Puiseux series in a common uniformiser by
\cref{p1:eq:pullback-coeff}; the ring $\Cplx[[v]]$ is closed under the sum defining $R$, under
exponentiation and under multiplication by $z=\sigma(1-v^{\mathsf L})$; and the outer
composition is either \cref{p1:eq:regular-outer} or \cref{p1:eq:critical-outer}, both of which
determine every coefficient at finite order by \cref{p0:def:obell,p0:def:ebell}. The
derivative formula in \cref{p1:eq:regular-outer} follows by differentiating $\mathcal
C=w\e^{\mathcal C}$.
\end{proof}

\begin{remark}[A dedup]
\label{p1:rem:generalised-bell}
\textsf{S1} introduces a ``generalized Bell algebra'' $\mathfrak B_{\gamma,m}$ on a
semigroup-graded power-series ring for this step. \Cref{p1:thm:local-completeness} avoids it:
after a common uniformiser $v$ has been chosen, everything is an \emph{ordinary} power series
in $v$ and the standard $\OB$, $\EB$ of \cref{p0:def:obell,p0:def:ebell} suffice.  That is
\textsf{S3}'s organisation and is adopted here; nothing is lost, because the choice of
$\mathsf L$ is exactly what the semigroup grading was encoding.
\end{remark}

\subsection{The first inherited sector, at $-\sqrt\rho$}
\label{p1:sub:minus}

The point
\begin{equation}
 \sigma_-:=-\sqrt\rho=-0.5816544136333942538101\ldots
 \label{p1:eq:sigma-minus}
\end{equation}
is reached by continuing the germ at $0$ along the negative real axis, which never meets the
dominant positive singularity.  At $z=\sigma_-$ the summand $T(z^{2})/2$ of $R$ reaches
$T(\rho)/2$ and inherits the dominant square-root branch, while every $T(z^{k})$ with $k\ge3$
is evaluated at $|\sigma_-|^{k}\le\rho^{3/2}<\rho$, strictly inside the original disc.  Take
the uniformiser
\begin{equation}
 z=\sigma_-(1-v^{2}),\qquad v=\Bigl(1-\tfrac{z}{\sigma_-}\Bigr)^{1/2},
 \qquad\text{so}\qquad
 1-\frac{z^{2}}{\rho}=v^{2}(2-v^{2}).
 \label{p1:eq:v-minus}
\end{equation}

\begin{proposition}[Complete local data at $\sigma_-$]
\label{p1:prop:minus}
Write $R\bigl(\sigma_-(1-v^{2})\bigr)=\sum_{m\ge0}r^-_mv^{m}$ and
$\zeta\bigl(\sigma_-(1-v^{2})\bigr)=\sum_{m\ge0}\zeta^-_mv^{m}$.  Then, with $t_0:=1$,
\begin{equation}
 r^-_m=\frac12\!\!
 \sum_{\substack{0\le r\le m\\ r\equiv m\,(2)}}\!\!
 t_r(-1)^{\frac{m-r}{2}}2^{\frac r2-\frac{m-r}{2}}
 \binom{r/2}{\tfrac{m-r}{2}}
 +\mathbf 1_{2\mid m}(-1)^{m/2}
 \sum_{k\ge3}\frac1k\sum_{n\ge1}a_n\sigma_-^{kn}\binom{kn}{m/2},
 \label{p1:eq:rminus}
\end{equation}
both parts converging geometrically.  Defining $\mathcal E^-_m$ by $\sum_{m\ge0}\mathcal
E^-_mv^{m}=\exp\bigl(\sum_{m\ge0}r^-_mv^{m}\bigr)$, so that $\mathcal
E^-_m=\e^{r^-_0}\sum_{k\ge0}\frac1{k!}\OB_{m,k}(r^-_1,r^-_2,\ldots)$, one has
\begin{equation}
 \zeta^-_m=\sigma_-\bigl(\mathcal E^-_m-\mathcal E^-_{m-2}\bigr),
 \qquad \mathcal E^-_{<0}:=0 .
 \label{p1:eq:zminus}
\end{equation}
Setting $y_-:=\mathcal C_0(\zeta^-_0)=-\Wlam_0(-\zeta^-_0)$, the outer composition is regular
because $y_-\ne1$, and writing $T\bigl(\sigma_-(1-v^{2})\bigr)=y_-+\sum_{m\ge1}t^-_mv^{m}$ one
has the closed form
\begin{equation}
 \boxed{\;r^-_1=\frac{t_1}{\sqrt2},
 \qquad
 t^-_1=\frac{y_-}{1-y_-}\cdot\frac{t_1}{\sqrt2}.\;}
 \label{p1:eq:bminus1}
\end{equation}
\end{proposition}

\begin{proof}
The $k=2$ contribution to $R$ is $\tfrac12T\bigl(\rho(1-v^{2}(2-v^{2}))\bigr)$; inserting
\cref{p1:eq:T-puiseux} with $s=v^{2}(2-v^{2})$, so $s^{r/2}=v^{r}2^{r/2}(1-v^{2}/2)^{r/2}$,
and expanding $(1-v^{2}/2)^{r/2}$ binomially gives the first sum.  The $k\ge3$ contributions
are $\sum_{k\ge3}\frac1k\sum_na_n\sigma_-^{kn}(1-v^{2})^{kn}$, whose $v^{2j}$ coefficient is
the second sum; convergence is geometric already at $k=3$ because $|\sigma_-|^{3}<\rho$.
\Cref{p1:eq:zminus} is $\zeta=\sigma_-(1-v^{2})\e^{R}$. For \cref{p1:eq:bminus1}: only $r=1$
contributes to $v^{1}$ in the first sum, giving $r^-_1=\tfrac12t_1\sqrt2=t_1/\sqrt2$; hence
$\zeta^-_1=\zeta^-_0r^-_1$ and, by \cref{p1:eq:regular-outer} with $\mathcal C_0'(w)=\mathcal
C_0(w)/ \bigl(w(1-\mathcal C_0(w))\bigr)$, $t^-_1=\mathcal C_0'(\zeta^-_0)\zeta^-_1
=\frac{y_-}{\zeta^-_0(1-y_-)}\zeta^-_0r^-_1$, which is the claim.
\end{proof}

Recomputed at $80$ working digits (see the repair below),
\begin{equation}
\begin{aligned}
 r^-_0&=\phantom{-}0.4686366279248455480872,&
 \zeta^-_0&=-0.9293757352012050888501,\\
 y_-&=-0.5410329414204534926205,&
 t^-_1&=\phantom{-}0.3871501110302429034772,
\end{aligned}
\label{p1:eq:minus-numeric}
\end{equation}
and the first local coefficients are
\begin{equation}
\begin{aligned}
 t^-_2&=-0.03256491851564763414731,&
 t^-_3&=\phantom{-}0.01720042160312209148301,\\
 t^-_4&=\phantom{-}0.10407703951838443852546,&
 t^-_5&=\phantom{-}0.16194736077270236118952,\\
 t^-_6&=-0.06292599818799308689694,&
 t^-_7&=-0.24403103923142135135510 .
\end{aligned}
\label{p1:eq:tminus-numeric}
\end{equation}

\begin{repairbox}
\textbf{\textsf{S2}'s minus-sector constants are printed past their accuracy.} \textsf{S2}
prints $r^-_0$, $z^-_0$, $y_-$ and $b^-_1$ to $40$ significant digits, e.g.\
$r^-_0=0.4686366279248455480872026033545827113234$.  Two independent recomputations here ---
one through the germ recursion at $90$ digits, one by direct summation of
$\tfrac12+\sum_{k\ge3}T(\sigma_-^{k})/k$ at $80$ digits with $1400$ exact tree numbers and no
truncation of the inner sums --- agree with each other and give
$r^-_0=0.468636627924845548087202621931819314\ldots$, which departs from \textsf{S2}'s value
at the \emph{twenty-sixth} significant digit.  The same $25$-digit agreement, and no more, is
found for $\zeta^-_0$, $y_-$, $t^-_1$ and every derived quantity, consistent with a single
truncation of the $k\ge3$ tail in \textsf{S2}'s computation.  All minus-sector constants are
therefore quoted here to at most $22$ digits.
\end{repairbox}

Only the odd $t^-_{2j+1}$ are singular in $z$ at $\sigma_-$, so \cref{p1:thm:general-sector}
with $\beta=j+\tfrac12$ gives
\begin{equation}
 \boxed{\;
 a_n^{(-)}\sim\mathfrak s_-\,\sigma_-^{-n}n^{-3/2}
 \sum_{k\ge0}\frac{B^-_k}{n^{k}},
 \qquad
 B^-_k=\sum_{j=0}^{k}\frac{t^-_{2j+1}}{\Gamma(-j-\tfrac12)}
 g_{k-j}\!\left(j+\tfrac12\right),\;}
 \label{p1:eq:minus-sector}
\end{equation}
where $\mathfrak s_-$ is the multiplier of the local cycle; for the direct negative-axis
continuation it is $1$, and it is displayed only to keep the contour dependence visible in
larger deformations.  Since $\sigma_-^{-n}=(-1)^{n}\rho^{-n/2}$, the sector is alternating,
with
\begin{equation}
\begin{aligned}
 B^-_0&=-0.10921302995631017570549,&
 B^-_1&=-0.03367666220778285336719,\\
 B^-_2&=-0.17900090117304631638903,&
 B^-_3&=-1.64234200495327998300,\\
 B^-_4&=-18.9266458632920839673,&
 B^-_5&=-277.723089846678995649 .
\end{aligned}
\label{p1:eq:Bminus-numeric}
\end{equation}
Relative to the leading dominant term its amplitude is
\begin{equation}
 \frac{B^-_0}{C}=-0.2482543049151649529916\ldots,
 \qquad
 \frac{a_n^{(-)}}{a_n^{(0)}}\sim
 -0.2482543049\ldots\,\bigl(-\sqrt\rho\bigr)^{n},
 \label{p1:eq:minus-relative}
\end{equation}
which is smaller than every power of $1/n$: this is the first explicitly computed genuinely
beyond-all-orders correction to A000081.  Its action is
\begin{equation}
 \Omega_-=\operatorname{Log}\frac{\sigma_-}{\rho}
 =\frac\lambda2+\ii\pi \pmod{2\pi\ii},
 \qquad \frac{\Omega_-}{\lambda}=\frac12+\ii\frac{\pi}{\lambda},
 \qquad \frac\pi\lambda=2.8987964297339787377\ldots .
 \label{p1:eq:minus-action}
\end{equation}

\begin{remark}
\label{p1:rem:other-sectors}
The point $+\sqrt\rho$ lies behind the dominant branch point on real continuation and belongs
to sheet-dependent continuations; complex critical points and overlapping root lifts may occur
farther out.  Their local data are produced by \cref{p1:thm:local-completeness} and their
multipliers depend on the chosen contour.  \Cref{p1:eq:minus-sector} is one explicit sector,
not an assertion that it exhausts modulus $\sqrt\rho$.
\end{remark}

\subsection{Transport of the sectors through the inverse}

Write the completed forward interpolation as
\begin{equation}
 \mathfrak a(x)=\mathcal A(x)\bigl(1+\mathcal E(x)\bigr),
 \qquad
 \mathcal E(x)=\sum_{\Omega\ne0}\e^{-\Omega x}x^{\mu_\Omega}
 \sum_{r\ge0}e_{\Omega,r}x^{-r},
 \label{p1:eq:full-interp}
\end{equation}
with $\mathcal A$ from \cref{p1:eq:interpolant} and the support of $\mathcal E$ contained in
the additive semigroup generated by the actions \cref{p1:eq:action}.  For non-integer $x$ a
factor $\sigma^{-x}$ requires a choice of $\operatorname{Log}\sigma$; different choices agree
at integers and differ between them, so the exponentially small part of a \emph{functional}
inverse is not canonical until a continuation and pairing convention has been fixed.  Fix one,
pairing conjugate branches so that $\mathfrak a$ is real on the positive axis.

\begin{proposition}[Sector transport]
\label{p1:prop:transport}
Let $\nu_0$ be the complete algebraic inverse of $\mathcal A$ (\cref{p1:eq:inverse-series})
and put $\nu_{\mathrm{full}}=\nu_0+\Delta$.  Then $\Delta$ is the unique solution, in the
transseries field ordered by the semigroup of actions, of
\begin{equation}
 \Delta=\Psi_{\nu_0}(\Delta),
 \qquad
 \Psi_{\nu_0}(w):=-\frac{\log\bigl(1+\mathcal E(\nu_0+w)\bigr)}
 {\mathcal Q_{\nu_0}(w)},
 \qquad
 \mathcal Q_{x}(w):=\frac{f(x+w)-f(x)}{w},
 \label{p1:eq:transport-fp}
\end{equation}
with $f=\log\mathcal A$ and $\mathcal Q_x(0)=f'(x)$; its coefficients are given by
\cref{p0:lem:lagrange-fp}, so that in particular
\begin{equation}
 \Delta=-\frac{\mathcal E(\nu_0)}{f'(\nu_0)}
 +\bigO(\mathcal E^{2})
 \sim-\frac{\mathcal E(\nu_0)}{\lambda},
 \label{p1:eq:transport-linear}
\end{equation}
and a single relative forward sector $\mathcal E_\Omega\sim e_{\Omega,0}\e^{-\Omega
x}x^{\mu_\Omega}$ maps to
\begin{equation}
 \Delta_\Omega(y)\sim
 -\frac{e_{\Omega,0}}{\lambda}\,C^{\Omega/\lambda}\,
 y^{-\Omega/\lambda}\,\nu_0^{\,\mu_\Omega-p\Omega/\lambda}.
 \label{p1:eq:transport-y}
\end{equation}
\end{proposition}

\begin{proof}
$f(\nu_0+\Delta)-f(\nu_0)+\log(1+\mathcal E(\nu_0+\Delta))=0$ is the exact equation, and
$f(\nu_0+w)-f(\nu_0)=w\,\mathcal Q_{\nu_0}(w)$ by definition; dividing gives
\cref{p1:eq:transport-fp}.  Since $f'(\nu_0)$ has leading term $\lambda\ne0$ it is invertible
in the coefficient field, and every occurrence of $\Delta$ on the right is either nonlinear or
multiplied by a positive transmonomial, so the support is well ordered and each requested
transmonomial receives finitely many contributions; \cref{p0:lem:lagrange-fp} then gives the
coefficients. \Cref{p1:eq:transport-linear} is the first term.  For \cref{p1:eq:transport-y},
use $y\sim C\e^{\lambda\nu_0}\nu_0^{-p}$ to write
$\e^{-\Omega\nu_0}=(y/C)^{-\Omega/\lambda}\nu_0^{-p\Omega/\lambda}$.
\end{proof}

For the first inherited sector, $\Omega_-=\tfrac\lambda2+\ii\pi$ and $e_{-,0}=B^-_0/C$.
Pairing the two lateral determinations $\pm\ii\pi$ replaces $\e^{-\ii\pi x}$ by $\cos(\pi x)$,
and $\e^{-\lambda\nu_0/2}=C^{1/2}y^{-1/2}\nu_0^{-3/4}$, so
\begin{equation}
 \boxed{\;
 \Delta_-(y)\sim
 0.15193347365953098826\ldots\;
 y^{-1/2}\,\nu_0(y)^{-3/4}\cos\bigl(\pi\nu_0(y)\bigr),\;}
 \label{p1:eq:minus-inverse}
\end{equation}
the amplitude being $\sqrt C\cdot\bigl(-B^-_0/(\lambda C)\bigr)$ with $-B^-_0/(\lambda
C)=0.22906811038403197857\ldots$.  The phase is
\begin{equation}
 \pi\nu_0(y)=\frac{\pi}{\lambda}\log y
 +\frac{3\pi}{2\lambda}\log\log y+\bigO(1),
 \qquad
 \frac{\pi}{\lambda}=2.89879642973397873771\ldots,
 \quad
 \frac{3\pi}{2\lambda}=4.34819464460096810657\ldots,
 \label{p1:eq:minus-phase}
\end{equation}
so the first nonperturbative correction to the inverse has envelope $y^{-1/2}(\log y)^{-3/4}$
and a logarithmic oscillation with a slower $\log\log y$ drift.

\begin{warning}[Two different oscillations]
\label{p1:rem:two-oscillations}
\Cref{p1:eq:minus-inverse} tends to zero and comes from a complex singular action.  The
sawtooth of \cref{p1:eq:sawtooth} is $\bigO(1)$, never tends to zero, and comes from integer
rounding.  They have different origins, different scales, and must not be conflated.
\end{warning}

\section{Other inverse notions: the compositional inverse of $T$}
\label{p1:sec:compositional}

Since $T(z)=z+\bigO(z^{2})$ with integer coefficients and unit linear coefficient, $T$ has a
compositional inverse $Z:=T^{\langle-1\rangle}\in z\Int[[z]]$, unrelated to the growth
inversion of \cref{p1:sec:inverse}.

\begin{proposition}[Coefficients of $Z$]
\label{p1:prop:compositional}
For $n\ge1$,
\begin{equation}
 \Coef{w^{n}}Z(w)=\frac1n\Coef{z^{n-1}}\left(\frac{z}{T(z)}\right)^{n}
 =\frac1n\Coef{z^{n-1}}
 \exp\!\left(-n\sum_{k\ge1}\frac{T(z^{k})}{k}\right),
 \label{p1:eq:Z-lagrange}
\end{equation}
a finite ordinary-Bell expression in $a_1,\ldots,a_{n-1}$.  The first sixteen values,
recomputed by exact rational arithmetic, are
\begin{equation}
 1,\,-1,\,0,\,1,\,-1,\,1,\,-4,\,11,\,-18,\,18,\,0,\,-60,\,189,\,-360,
 \,453,\,-373 .
 \label{p1:eq:Z-values}
\end{equation}
\end{proposition}

\begin{proof}
The first equality is \cref{p0:thm:LB} with $\phi(z)=T(z)/z$, a unit series. The second is
\cref{p1:eq:polya} rearranged, and \cref{p0:def:obell} converts the exponential into a finite
Bell sum. Integrality is automatic for the compositional inverse of a series in
$z+z^{2}\Int[[z]]$.
\end{proof}

\begin{remark}
\textsf{S1} displays $Z(w)=w-w^{2}+w^{4}-w^{5}+w^{6}-4w^{7}+11w^{8}-\cdots$, which agrees with
\cref{p1:eq:Z-values}, and identifies the coefficients with OEIS A050395.  The expansion is
confirmed here; the OEIS identification is reported from \textsf{S1} and has not been
re-checked in this volume.
\end{remark}

Near the dominant value $w=1$ the branch of $Z$ through $Z(1)=\rho$ is \emph{analytic} in
$1-w$, because it undoes the square root: from $T=\zeta\e^{T}$ one has $\zeta=T\e^{-T}$, hence
\begin{equation}
 Z(w)=\zeta^{\langle-1\rangle}\bigl(w\e^{-w}\bigr)
 \label{p1:eq:Z-near-one}
\end{equation}
on that branch, $\zeta$ being injective on $(0,\sqrt\rho)$ by \cref{p1:prop:critical}.  Since
$w\e^{-w}-\e^{-1}$ has a double zero at $w=1$, the local expansion of $Z$ at $w=1$ has no
linear term.  Its coefficients follow from one more Lagrange extraction: writing
$\zeta(\rho+x)-\e^{-1}=x\,\psi(x)$ with $\psi(0)=\zeta'(\rho)\ne0$,
\begin{equation}
 \Coef{\mathsf v^{m}}\bigl(\zeta^{\langle-1\rangle}(\e^{-1}+\mathsf v)
 -\rho\bigr)=\frac1m\Coef{x^{m-1}}\psi(x)^{-m},
 \label{p1:eq:zeta-local-inverse}
\end{equation}
followed by composition with $(1-\varsigma)\e^{-(1-\varsigma)}-\e^{-1}
=-\e^{-1}\sum_{m\ge2}\frac{m-1}{m!}\varsigma^{m}
=-\e^{-1}\frac{\varsigma^{2}}{2}\Phi(\varsigma)$, the same kernel \cref{p1:eq:Phi} that
governed the forward branch.

\section{Algorithms, validation and audit}
\label{p1:sec:algorithms}

\subsection{The pipeline}

\begin{algorithm}[All dominant and inverse coefficients through order $M$]
\label{p1:alg:pipeline}
\leavevmode
\begin{enumerate}
\item Generate $a_1,\ldots,a_N$ and $b_1,\ldots,b_N$ exactly by
      \cref{p1:eq:recurrence} with a divisor sieve, checking the divisibility
      assertion at each step; form $\varrho_r$ by \cref{p1:eq:varrho}.
\item Solve $\log x+1+\sum_{r=2}^{N}\varrho_rx^{r}=0$ for $\rho$ by
      safeguarded Newton or interval bisection.  Choose $N$ so that
      $\rho^{N/2}$ is below the target precision; the omitted tail is
      geometrically small because $R$ is analytic on $|z|<\sqrt\rho$
      (\cref{p1:prop:R-analytic}).
\item Compute $\widehat R_0,\ldots,\widehat R_{2M+1}$ by
      \cref{p1:eq:Rhat}; exponentiate by \cref{p1:eq:E-bell} and form
      $\theta_1,\ldots,\theta_{2M+2}$ by \cref{p1:eq:theta-bell}, then
      $u_j=\theta_{j+1}/\theta_1$.
\item Compute $\omega_1,\ldots,\omega_{2M+1}$ exactly by
      \cref{p1:eq:omega-master} (rationals), set $c_r=-2^{r/2}\omega_r$, and
      form $t_1,t_3,\ldots,t_{2M+1}$ by \cref{p1:eq:t-master}.
\item Compute $g_r(j+\tfrac12)$ exactly for $j+r\le M$ by
      \cref{p1:eq:g-recurrence}, then $\tau_0,\ldots,\tau_M$ by
      \cref{p1:eq:tau-master}, then $C=\tau_0/\sqrt\pi$ and
      $f_m=\tau_m/\tau_0$.
\item Reciprocal: $\bar f_m$ by \cref{p1:eq:recip-coeff}.
\item Inverse: $\varphi_m$ by \cref{p1:eq:phi}, then $d_m$ by
      \cref{p1:eq:d-triangular}, or $P_m(L)$ by the flattening recursion of
      \cref{p0:thm:flattening}.
\item Sectors: apply \cref{p1:thm:local-completeness} outward in modulus,
      transfer each germ by \cref{p1:eq:general-sector}, and invert by
      \cref{p1:prop:transport}.
\end{enumerate}
\end{algorithm}

At every stage the coefficient of order $m$ depends only on data of order at most $m$; no
nonlinear global solve occurs after $\rho$ is known.  The useful power-series primitives are
the three dense recurrences
\begin{equation}
 \Coef{v^{n}}F^{\gamma}=\frac1n\sum_{k=1}^{n}
 \bigl((\gamma+1)k-n\bigr)F_k\Coef{v^{n-k}}F^{\gamma},
 \quad
 E_n=\frac1n\sum_{k=1}^{n}kA_kE_{n-k},
 \quad
 H_n=F_n-\frac1n\sum_{k=1}^{n-1}kH_kF_{n-k},
 \label{p1:eq:dense-recurrences}
\end{equation}
for $F^{\gamma}$ with $F_0=1$, for $E=\e^{A}$ with $A_0=0$, and for $H=\log F$ with $F_0=1$.
They are equivalent to the Bell forms used in the statements and are what an implementation
should actually run.

\begin{warning}[Precision discipline]
\label{p1:rem:precision}
The coefficients grow rapidly and \cref{p1:eq:tau-master} combines alternating contributions
from several Puiseux coefficients, so obtaining $d$ reliable digits of $\tau_M$ can require
substantially more than $d$ working digits.  The tables in this chapter were produced with
$90$--$110$ decimal working digits, with universal objects ($\omega_m$, $g_r(\beta)$, $a_n$)
held in exact rational or integer arithmetic.  The repair recorded at \cref{p1:sub:minus} is
exactly what happens when this discipline lapses.
\end{warning}

\subsection{Numerical validation}
\label{p1:sub:numerics}

Let $A_M(n)$ denote \cref{p1:eq:forward} truncated after $\tau_M n^{-M}$.
\Cref{p1:tab:forward-errors} reports the signed relative error $\bigl(A_M(n)-a_n\bigr)/a_n$
against exact integers from \cref{p1:eq:recurrence}; \cref{p1:tab:inverse-errors} reports the
signed index error $\nu_M(a_n)-n$ for the Lambert-normalised inverse $\nu_M=x_0+\sum_{m\le
M}d_mx_0^{-m}$.  Both tables were regenerated here and agree entry by entry with the
corresponding tables of \textsf{S1} and \textsf{S2} at the displayed precision.

\begin{table}[htbp]
\centering
\caption{Signed relative error $(A_M(n)-a_n)/a_n$ of the forward expansion.}
\label{p1:tab:forward-errors}
\scriptsize
\setlength{\tabcolsep}{4pt}
\begin{tabular}{rcccccc}
\toprule
$n$ & $M=0$ & $M=2$ & $M=4$ & $M=6$ & $M=8$ & $M=10$\\
\midrule
$10$  & $-1.52\cdot10^{-2}$ & $-3.49\cdot10^{-4}$ & $\phantom{-}2.63\cdot10^{-3}$ & $\phantom{-}3.99\cdot10^{-3}$ & $\phantom{-}5.52\cdot10^{-3}$ & $\phantom{-}8.77\cdot10^{-3}$\\
$20$  & $-6.58\cdot10^{-3}$ & $-3.42\cdot10^{-4}$ & $-3.18\cdot10^{-5}$ & $\phantom{-}1.04\cdot10^{-6}$ & $\phantom{-}9.73\cdot10^{-6}$ & $\phantom{-}1.41\cdot10^{-5}$\\
$50$  & $-2.22\cdot10^{-3}$ & $-1.77\cdot10^{-5}$ & $-2.88\cdot10^{-7}$ & $-1.20\cdot10^{-8}$ & $-9.84\cdot10^{-10}$ & $-1.36\cdot10^{-10}$\\
$100$ & $-1.06\cdot10^{-3}$ & $-2.08\cdot10^{-6}$ & $-8.06\cdot10^{-9}$ & $-7.88\cdot10^{-11}$ & $-1.50\cdot10^{-12}$ & $-4.73\cdot10^{-14}$\\
$200$ & $-5.15\cdot10^{-4}$ & $-2.53\cdot10^{-7}$ & $-2.40\cdot10^{-10}$ & $-5.73\cdot10^{-13}$ & $-2.66\cdot10^{-15}$ & $-2.04\cdot10^{-17}$\\
$500$ & $-2.03\cdot10^{-4}$ & $-1.59\cdot10^{-8}$ & $-2.39\cdot10^{-12}$ & $-9.02\cdot10^{-16}$ & $-6.62\cdot10^{-19}$ & $-8.01\cdot10^{-22}$\\
\bottomrule
\end{tabular}
\end{table}

\begin{table}[htbp]
\centering
\caption{Signed index error $\nu_M(a_n)-n$ of the Lambert-normalised
inverse; $M=0$ is the bare carrier $x_0$.}
\label{p1:tab:inverse-errors}
\scriptsize
\setlength{\tabcolsep}{4pt}
\begin{tabular}{rcccccc}
\toprule
$n$ & $M=0$ & $M=1$ & $M=2$ & $M=4$ & $M=6$ & $M=8$\\
\midrule
$10$  & $\phantom{-}1.64\cdot10^{-2}$ & $\phantom{-}7.15\cdot10^{-3}$ & $\phantom{-}1.29\cdot10^{-3}$ & $-2.61\cdot10^{-3}$ & $-4.16\cdot10^{-3}$ & $-5.76\cdot10^{-3}$\\
$20$  & $\phantom{-}6.55\cdot10^{-3}$ & $\phantom{-}1.91\cdot10^{-3}$ & $\phantom{-}4.44\cdot10^{-4}$ & $\phantom{-}3.76\cdot10^{-5}$ & $-2.02\cdot10^{-7}$ & $-9.39\cdot10^{-6}$\\
$50$  & $\phantom{-}2.11\cdot10^{-3}$ & $\phantom{-}2.59\cdot10^{-4}$ & $\phantom{-}2.32\cdot10^{-5}$ & $\phantom{-}3.32\cdot10^{-7}$ & $\phantom{-}1.26\cdot10^{-8}$ & $\phantom{-}9.98\cdot10^{-10}$\\
$100$ & $\phantom{-}9.88\cdot10^{-4}$ & $\phantom{-}6.16\cdot10^{-5}$ & $\phantom{-}2.74\cdot10^{-6}$ & $\phantom{-}9.34\cdot10^{-9}$ & $\phantom{-}8.35\cdot10^{-11}$ & $\phantom{-}1.53\cdot10^{-12}$\\
$200$ & $\phantom{-}4.78\cdot10^{-4}$ & $\phantom{-}1.50\cdot10^{-5}$ & $\phantom{-}3.33\cdot10^{-7}$ & $\phantom{-}2.79\cdot10^{-10}$ & $\phantom{-}6.09\cdot10^{-13}$ & $\phantom{-}2.71\cdot10^{-15}$\\
$500$ & $\phantom{-}1.88\cdot10^{-4}$ & $\phantom{-}2.38\cdot10^{-6}$ & $\phantom{-}2.10\cdot10^{-8}$ & $\phantom{-}2.78\cdot10^{-12}$ & $\phantom{-}9.59\cdot10^{-16}$ & $\phantom{-}6.75\cdot10^{-19}$\\
\bottomrule
\end{tabular}
\end{table}

Two features are worth naming.  First, at $n=10$ and $n=20$ the error stops improving and then
worsens: the expansion is asymptotic, and the least-term truncation of
\cref{p0:thm:optimal-truncation} has been passed.  Second, the inverse errors track the
forward relative errors divided by $f'\approx\lambda$, exactly as
\cref{p1:thm:inverse-remainder} predicts; this is why
\cref{p1:tab:forward-errors,p1:tab:inverse-errors} are so nearly proportional.

\begin{remark}[On comparisons with the literature]
\label{p1:rem:genitrini}
\textsf{S1} states that its first five coefficients ``agree with Genitrini's independently
computed values'' and \textsf{S3} that its Puiseux table ``agrees with the independent table
in [Genitrini 2016]''.  Those comparisons are reported from the drafts and have \emph{not}
been re-performed here; what has been done is a full independent recomputation of every
constant from the exact recurrence, described in \cref{p1:sub:audit}.
\end{remark}

\subsection{Audit}
\label{p1:sub:audit}

Every numerical value printed in this chapter was produced by an independent implementation of
\cref{p1:alg:pipeline} written for this volume: exact integer generation of $a_1,\ldots,a_N$
with $N$ between $700$ and $1400$; exact rational $\varrho_r$, $\omega_m$ and $g_r(\beta)$
(the last from an exact Bernoulli-number table); and $60$--$110$ decimal working digits for
everything else.  The following independent identities were checked and hold to all digits
carried.

\begin{enumerate}
\item $(n-1)\mid\sum_{j<n}b_ja_{n-j}$ at every step, for $n\le1400$;
      $a_{10}=719$, $a_{20}=12\,826\,228$, $a_{50}=425\,976\,989\,835\,141\,038\,353$.
\item $\lambda=1+R(\rho)$, i.e.\ \cref{p1:eq:lambda-check}.
\item $\theta_0=0$, the vanishing of the constant term in
      \cref{p1:eq:theta-def} predicted by $\e\zeta(\rho)=1$.
\item $C^{2}=\theta_1/(2\pi)$, i.e.\ \cref{p1:eq:otter} --- the two sides
      travel through disjoint parts of the pipeline.
\item $c_r=-2^{r/2}\omega_r$ for $1\le r\le18$ in exact arithmetic, and
      $t_1=-\sqrt{2\theta_1}$, $t_2=\tfrac23\theta_1$.
\item \Cref{p1:eq:tau-expanded} re-derived symbolically from
      \cref{p1:eq:tau-master}, and the $g_r(\beta)$ of
      \cref{p1:eq:g-first,p1:eq:g4} re-derived from
      \cref{p1:eq:varkappa}.
\item $\mathrm D_m=\lambda^{m+1}d_m$ and
      $\eta^{\textsf{S2}}_m=\lambda^{m}\varphi_m$ for $1\le m\le8$,
      reconciling \textsf{S1}/\textsf{S3} with \textsf{S2}.
\item $P_1(0)=d_1$, $P_2(0)=d_2$ (a case of \cref{p1:eq:P-shape}).
\item $r^-_1=t_1/\sqrt2$ and the closed form \cref{p1:eq:bminus1}, each
      against the germ recursion; and $r^-_0$ by a second, wholly independent
      summation (see the repair in \cref{p1:sub:minus}).
\item The single-sum \cref{p1:eq:Rhat} against the siblings' double sum, for
      $0\le m\le30$.
\end{enumerate}

\begin{table}[htbp]
\centering
\caption{Principal symbols of this chapter.}
\label{p1:tab:notation}
\small
\begin{tabular}{p{0.17\textwidth}p{0.75\textwidth}}
\toprule
Symbol & Meaning\\
\midrule
$a_n$, $T(z)$ & A000081 and its ordinary generating function.\\
$b_m$, $\varrho_r$ & Divisor transform $\sum_{d\mid m}da_d$; coefficients of $R$, $\varrho_r=(b_r-ra_r)/r$.\\
$R$, $\zeta$ & Disturbance $\sum_{k\ge2}T(z^{k})/k$ and Cayley argument $z\e^{R}$; $T=-\Wlam_0(-\zeta)$.\\
$\rho,\alpha,\lambda,p$ & Otter singularity; $\alpha=\rho^{-1}$, $\lambda=\log\alpha$, $p=\tfrac32$.\\
$\Phi$, $\omega_m$, $c_r$ & Universal Cayley kernel and its two coefficient normalisations, $c_r=-2^{r/2}\omega_r$.\\
$s$, $\theta$, $\theta_m$, $u_j$ & Local coordinate $1-z/\rho$; disturbance $1-\e\zeta$; its coefficients; $u_j=\theta_{j+1}/\theta_1$.\\
$\widehat R_m$, $\ell_j$ & $\Coef{s^m}R(\rho(1-s))$; scaled logarithmic jets $\rho^j(\log\zeta)^{(j)}(\rho)$.\\
$t_n$ & Puiseux coefficients, $T=1+\sum t_ns^{n/2}$.\\
$\varkappa_m(\beta)$, $g_r(\beta)$ & Bernoulli increments and gamma-ratio transfer coefficients.\\
$\tau_m$, $C$, $f_m$, $\bar f_m$ & Forward coefficients; $C=\tau_0/\sqrt\pi$; $f_m=\tau_m/\tau_0$; reciprocal coefficients.\\
$F$, $H$, $\varphi_m$ & $F=1+\sum f_mv^m$; $H=\log F=\sum\varphi_mv^m$; this $H$ is the Part~0 datum $Q$.\\
$x_0$, $d_m$, $\nu$ & Lambert carrier; its corrections; the smooth asymptotic inverse.\\
$\mathcal X$, $L$, $\mathsf q$, $\eta_k$, $P_m$ & Flattening variables and coefficient polynomials.\\
$N(y)$ & Integer staircase $\min\{n\ge1:a_n\ge y\}$.\\
$\sigma_-$, $t^-_m$, $B^-_k$, $\Omega_\sigma$ & The inherited singularity $-\sqrt\rho$, its local and sector coefficients, and the actions.\\
\bottomrule
\end{tabular}
\end{table}

\section{Summary}
\label{p1:sec:summary}

The P\'olya equation becomes transparent after the factorisation $T=\mathcal C\circ\zeta$ with
$\mathcal C(w)=-\Wlam_0(-w)$: the universal layer is one Lagrange coefficient
$\Coef{v^{r-1}}\Phi(v)^{-r/2}$, the problem-specific layer is ordinary Bell composition of the
analytic disturbance $\zeta$, and coefficient transfer is Bernoulli polynomials in the gamma
ratio.  The chain is
\begin{equation}
 a_n\;\to\;\varrho_r\;\to\;\rho\;\to\;\widehat R_m\;\to\;\theta_m
 \;\to\;t_n\;\to\;\tau_m\;\to\;f_m\;\to\;\varphi_m
 \;\to\;d_m\ \text{or}\ P_m(L),
 \label{p1:eq:chain}
\end{equation}
every arrow a finite Bell- or Bernoulli-polynomial formula.  Inversion adds nothing new:
A000081 is the $\alpha=1$ case of \cref{p0:def:model} with $a=\lambda$, $b=-\tfrac32$, $Q=H$,
so \cref{p0:thm:lambert-core} supplies the carrier exactly, \cref{p0:thm:lambert-centered}
every algebraic correction, \cref{p0:thm:flattening} the logarithmic form,
\cref{p0:thm:remainder-transport} the error, and \cref{p0:thm:staircase} the passage to the
integer inverse. Beyond all orders, the recursion supplies a singularity grammar whose first
explicit instance is the alternating sector of relative size $0.2483\ldots\,(-\sqrt\rho)^{n}$
at $\sigma_-=-\sqrt\rho$, transported by \cref{p1:prop:transport} into an inverse correction
of envelope $y^{-1/2}(\log y)^{-3/4}$ with a $\log y$-periodic phase.  What remains open is
not coefficient algebra but analysis: the classification of accessible sheets, of critical
preimages, and of their Stokes connections --- and, prior to that, a proof of $\mathsf
H(\mathcal R)$ for $\mathcal R$ close to $1$.

%% ==================== fragment p2 ====================
\chapter{The double factorial}\label{p2:sec:top}

\begin{sourcebox}
This chapter merges three independent treatments of \(n!!\):
\emph{Double\_Factorial\_Transseries} (henceforth \textsf{S1}),
\emph{Double\_Factorial\_Transseries-3} (\textsf{S2}), and
\emph{double\_factorial\_transseries-2} (\textsf{S3}).
All three reach the same forward and inverse coefficients; they differ in
organisation, in what is proved, and --- decisively --- in notation, where the
letter \(C\) carries three incompatible meanings.  The chapter adopts
\textsf{S3}'s \emph{exact} centred normal form as its spine, \textsf{S2}'s Borel
kernel and sharp remainder, \textsf{S2}'s graded closed coefficient formula, and
\textsf{S1}/\textsf{S2}'s parity-mean centring of the periodic interpolation.
Every displayed rational number below was regenerated with exact rational
arithmetic (\texttt{sympy}) from the defining recurrences, and every numerical
table was recomputed at \(220\)-digit precision (\texttt{mpmath}); see
\cref{p2:sec:numerics}.
\end{sourcebox}

\section{Orientation: which inverse?}
\label{p2:sec:orientation}

For \(n\in\Nat\) the double factorial is
\begin{equation}
 0!!=1!!=1,\qquad n!!=n\,(n-2)!!\quad(n\ge2),
 \label{p2:eq:dfac-recurrence}
\end{equation}
so that \(n!!\) is the product of the positive integers not exceeding \(n\)
of the same parity as \(n\).  The sequence
\(1,1,2,3,8,15,48,105,384,945,\dots\) is OEIS~A006882.  Its elementary closed
forms already contain the whole difficulty:
\begin{equation}
 (2m)!!=2^m\,m!,
 \qquad
 (2m+1)!!=\frac{2^{m+1}\Gamma\!\left(m+\tfrac32\right)}{\sqrt\pi}.
 \label{p2:eq:parity-closed}
\end{equation}
The even subsequence is a gamma function of the halved argument; the odd
subsequence is the \emph{same} gamma function of the halved argument times a
different constant.  The two constants do not agree, and no amount of
reindexing makes them agree.

\begin{warning}
\label{p2:rem:spine}
\(n!!\) is not the restriction to \(\Nat\) of one canonical slowly varying
function.  Its even and odd subsequences are two gamma-type branches whose
quotient is a nonzero constant \(\sqrt{2/\pi}\ne1\).  Consequently
\emph{there is no canonical real inverse of \(n!!\) until an interpolation
is fixed}, and different admissible interpolations give inverses that differ
at an order \emph{larger} than the entire Bernoulli correction grid
(\cref{p2:thm:no-canonical}).  Every statement in this chapter that begins
``the inverse'' names its interpolation first.
\end{warning}

Following \cref{p0:def:three-inverses} we therefore distinguish, for a target
\(y\to+\infty\):

\begin{enumerate}[label=\textup{(\roman*)}]
\item \emph{the two parity-resolved functional inverses}, obtained by inverting
      the canonical gamma continuation of one parity subsequence
      (\cref{p2:sec:branch-inverse});
\item \emph{the functional inverse of a single all-integer interpolation},
      for instance the periodic one recorded by OEIS
      (\cref{p2:sec:oeis-inverse});
\item \emph{the integer staircase inverses}, the least \(n\) with \(n!!\ge y\)
      and the greatest \(n\) with \(n!!\le y\) (\cref{p2:sec:discrete}).
\end{enumerate}

Objects (i) and (iii) are intrinsic to the sequence; object (ii) is not.  This
is the chapter's spine, and \cref{p2:thm:no-canonical} makes it quantitative.

Two asymptotic scales must also be kept apart.  Taking logarithms converts the
growth of \(n!!\) into a phase \(\tfrac12 u(\log u-1)\) whose inversion is
\emph{exactly} a Lambert expression (\cref{p0:prop:factorial-core}); expanding
that Lambert function produces an \emph{outer} layer in powers of
\(1/\log\log y\) with polynomial dependence on \(\log\log\log y\)
(\cref{p0:thm:flattening}).  The Bernoulli corrections generate a far finer
\emph{inner} grid in odd inverse powers of the Lambert core.  The fixed shift
\(-1\) lies strictly between the two levels.  Writing all of this as one
informal series in \(1/\log y\) destroys the ordering; see
\cref{p2:prop:ordering}.

\section{Exact continuations and the parity split}
\label{p2:sec:continuations}

\subsection{The two parity branches}

\begin{definition}[Parity branches]
\label{p2:def:branches}
For the parity label \(\delta\in\{0,1\}\) put
\begin{equation}
 \kappa_\delta=\left(\frac2\pi\right)^{\delta/2},
 \qquad
 \mathcal D_\delta(x)=\kappa_\delta\,2^{x/2}\,
 \Gamma\!\left(\frac x2+1\right)
 \qquad(x>-2),
 \label{p2:eq:branches}
\end{equation}
and write
\begin{equation}
 A_\delta=\log\!\left(\kappa_\delta\sqrt\pi\right)
 =\begin{cases}\tfrac12\log\pi,&\delta=0,\\[0.2em]
               \tfrac12\log2,&\delta=1,\end{cases}
 \qquad
 \e^{A_0}=\sqrt\pi,\quad \e^{A_1}=\sqrt2 .
 \label{p2:eq:Adelta}
\end{equation}
Here \(\delta=0\) labels the even subsequence and \(\delta=1\) the odd one.
\end{definition}

\begin{proposition}[Exact parity interpolation and monotonicity]
\label{p2:prop:branch-interp}
For every integer \(n\ge0\) with \(n\equiv\delta\pmod2\) one has
\(\mathcal D_\delta(n)=n!!\).  Moreover each \(\mathcal D_\delta\) is positive
and strictly increasing on \([0,\infty)\), so it has a single-valued inverse
\(\mathcal D_\delta^{-1}\) on its range; and
\begin{equation}
 \mathcal D_1(x)=\sqrt{\frac2\pi}\;\mathcal D_0(x),
 \qquad
 \mathcal D_1^{-1}(y)=\mathcal D_0^{-1}\!\left(\sqrt{\frac\pi2}\,y\right).
 \label{p2:eq:branch-scaling}
\end{equation}
\end{proposition}

\begin{proof}
If \(n=2m\) then \(\mathcal D_0(2m)=2^m\Gamma(m+1)=2^mm!=(2m)!!\) by
\cref{p2:eq:parity-closed}.  If \(n=2m+1\) then
\[
 \mathcal D_1(2m+1)=\sqrt{\frac2\pi}\,2^{m+1/2}
 \Gamma\!\left(m+\tfrac32\right)
 =\frac{2^{m+1}}{\sqrt\pi}\Gamma\!\left(m+\tfrac32\right)=(2m+1)!!.
\]
Positivity is clear because \(\Gamma>0\) on \((0,\infty)\).  For monotonicity,
\begin{equation}
 \frac{\mathcal D_\delta'(x)}{\mathcal D_\delta(x)}
 =\frac12\left(\log2+\psi\!\left(\frac x2+1\right)\right),
 \qquad \psi=\Gamma'/\Gamma,
 \label{p2:eq:branch-logder}
\end{equation}
and \(\psi\) is strictly increasing on \((0,\infty)\) because
\(\psi'(z)=\sum_{m\ge0}(z+m)^{-2}>0\).  At \(x=0\) the bracket equals
\(\log2-\gamma=0.11593\ldots>0\), hence it is positive for all \(x\ge0\).
\Cref{p2:eq:branch-scaling} is immediate from \cref{p2:eq:branches}, since
\(\kappa_1/\kappa_0=\sqrt{2/\pi}\).
\end{proof}

\begin{remark}
\label{p2:rem:kappa-cancels}
The constant \(\kappa_\delta\) disappears from
\cref{p2:eq:branch-logder}.  Both branches therefore have the same local
conditioning, the same forward correction series, and --- as
\cref{p2:thm:branch-inverse} shows --- literally the same inverse coefficient
polynomials.  Parity survives only as the additive constant \(A_\delta\)
subtracted from \(\log y\).
\end{remark}

\subsection{The periodic OEIS interpolation, and why it is not canonical}

\begin{definition}[Periodic interpolation]
\label{p2:def:oeis}
Put
\begin{equation}
 \mathscr D(x)=2^{x/2}\,\Gamma\!\left(\frac x2+1\right)
 \left(\frac2\pi\right)^{(1-\cos\pi x)/4}.
 \label{p2:eq:oeis}
\end{equation}
\end{definition}

Because \(\tfrac12\sin^2(\pi x/2)=\tfrac14(1-\cos\pi x)\), this is the
interpolation recorded in the OEIS entry for A006882, written there with the
half-argument sine.

\begin{proposition}[Interpolation and eventual monotonicity]
\label{p2:prop:oeis-interp}
\(\mathscr D(n)=n!!\) for every integer \(n\ge0\), and \(\mathscr D\) is
strictly increasing on \([\,1.04,\infty)\).  Its inverse on that ray,
denoted \(\mathscr D^{-1}\), is what ``the inverse of \(\mathscr D\)'' means
below; it is not a globally single-valued inverse near the origin.
\end{proposition}

\begin{proof}
The periodic factor equals \(1\) at even integers and \(\sqrt{2/\pi}\) at odd
integers, so \(\mathscr D\) agrees with \(\mathcal D_0\) at even and with
\(\mathcal D_1\) at odd integers; \cref{p2:prop:branch-interp} finishes the
first claim.  Next,
\begin{equation}
 \frac{\mathscr D'(x)}{\mathscr D(x)}
 =\frac12\left(\log2+\psi\!\left(\frac x2+1\right)\right)
 +\frac\pi4\log\!\left(\frac2\pi\right)\sin\pi x .
 \label{p2:eq:oeis-logder}
\end{equation}
For \(z>\tfrac12\) one has \(\psi(z)>\log\!\left(z-\tfrac12\right)\), hence
\(\psi(x/2+1)>\log\!\left((x+1)/2\right)\) and the first summand exceeds
\(\tfrac12\log(x+1)\).  The second summand has modulus at most
\(\pi a\), where
\begin{equation}
 a:=\frac14\log\frac\pi2=0.112895676322\ldots>0 .
 \label{p2:eq:a-def}
\end{equation}
Thus \(\mathscr D'/\mathscr D>\tfrac12\log(x+1)-\pi a\), which is positive as
soon as \(x+1>\e^{2\pi a}=2.03265\ldots\), i.e.\ for \(x\ge1.04\).
\end{proof}

\begin{remark}
The bound in the proof is not sharp: numerically the last zero of
\cref{p2:eq:oeis-logder} is at \(x=0.6970947\ldots\), so \(\mathscr D\) is in
fact strictly increasing on \([0.70,\infty)\).  Only the proved threshold is
used below.
\end{remark}

The next theorem is the precise form of \cref{p2:rem:spine}.  None of the three
sources states it; all three assert the phenomenon informally.

\begin{theorem}[The inverse transseries is not determined by the sequence]
\label{p2:thm:no-canonical}
For \(\lambda\in\Real\) set
\(\mathscr D_\lambda(x)=\mathscr D(x)\,\e^{\lambda\sin\pi x}\).
Then:
\begin{enumerate}[label=\textup{(\alph*)}]
\item \(\mathscr D_\lambda(n)=n!!\) for every integer \(n\ge0\), and
      \(\mathscr D_\lambda\) is real-analytic on \(\Real\);
\item \(\mathscr D_\lambda\) is strictly increasing on
      \([\,\e^{2\pi(a+|\lambda|)}-1,\infty)\), so it has a large real inverse;
\item with \(T\) and \(\Lambda\) the Lambert data of
      \cref{p2:eq:lambert-core} built from
      \(R_*=\log y-\tfrac14\log(2\pi)\), the leading displacement of that
      inverse beyond the core is
      \begin{equation}
       \mathscr D_\lambda^{-1}(y)-(T-1)
       =\frac{2\sqrt{a^2+\lambda^2}}{\Lambda}
        \cos\!\left(\pi T-\arctan\frac\lambda a\right)
        +\bigO\!\left(\Lambda^{-2}\right).
       \label{p2:eq:no-canonical}
      \end{equation}
\end{enumerate}
In particular the amplitude \(2\sqrt{a^2+\lambda^2}/\Lambda\) of the leading
oscillatory sector is an unbounded function of the interpolation and is
\emph{larger, by the factor \(T\), than the entire Bernoulli grid}
\(\bigO(T^{-1}\Lambda^{-1})\) of \cref{p2:thm:branch-inverse}.
\end{theorem}

\begin{proof}
(a) \(\sin\pi n=0\) for integer \(n\), and \(\e^{\lambda\sin\pi x}\) is entire.

(b) Differentiating logarithmically adds \(\lambda\pi\cos\pi x\) to
\cref{p2:eq:oeis-logder}; repeating the estimate in the proof of
\cref{p2:prop:oeis-interp} gives
\(\mathscr D_\lambda'/\mathscr D_\lambda>\tfrac12\log(x+1)-\pi(a+|\lambda|)\).

(c) By \cref{p2:lem:oeis-centred} below, with \(u=x+1\),
\[
 \log\mathscr D_\lambda(u-1)
 =\frac u2(\log u-1)+\frac14\log(2\pi)-a\cos\pi u
   -\lambda\sin\pi u+\Omega(u),
\]
because \(\cos\pi x=-\cos\pi u\) and \(\sin\pi x=-\sin\pi u\).  Subtracting the
core identity \(\tfrac12T(\log T-1)=R_*\), putting \(u=T+\Delta\) and
multiplying by \(2\) exactly as in \cref{p2:lem:oeis-displacement} gives
\[
 \Lambda\Delta+T\,\varphi(\Delta/T)
 +2\Omega(T+\Delta)
 -2a\cos(\theta+\pi\Delta)-2\lambda\sin(\theta+\pi\Delta)=0,
 \qquad\theta=\pi T .
\]
All terms except the last two are \(\bigO(T^{-1})\) once
\(\Delta=\bigO(\Lambda^{-1})\), so at leading order
\(\Lambda\Delta_0=2a\cos(\theta+\pi\Delta_0)+2\lambda\sin(\theta+\pi\Delta_0)
=2\sqrt{a^2+\lambda^2}\cos(\theta+\pi\Delta_0-\arctan(\lambda/a))\).
Since \(\Delta_0=\bigO(\Lambda^{-1})\), replacing
\(\theta+\pi\Delta_0\) by \(\theta\) costs \(\bigO(\Lambda^{-2})\), which is
\cref{p2:eq:no-canonical}.
\end{proof}

\begin{remark}
The family \(\e^{\lambda\sin\pi x}\) is only the simplest witness.  Any
real-analytic \(g\) vanishing on \(\Nat\) and with \(g'\) bounded produces the
same conclusion, with \(2a\cos\theta\) replaced by the value of
\(-2g\) transported to the core.  What \emph{is} intrinsic is the pair of
parity-branch transseries and the integer staircase; the periodic sector is a
property of the chosen interpolation and nothing else.
\end{remark}

\subsection{The centred coordinate and the exact normal form}

\begin{definition}[Centred coordinate]
\label{p2:def:centred}
Throughout, \(u:=x+1\) and
\begin{equation}
 \Phi(u):=\frac u2(\log u-1),
 \qquad
 \Omega(u):=\log\Gamma\!\left(\frac{u+1}2\right)
 -\frac u2\log\frac u2+\frac u2-\frac12\log(2\pi).
 \label{p2:eq:Omega-def}
\end{equation}
\end{definition}

\begin{theorem}[Exact centred normal form]
\label{p2:thm:normal-form}
For every \(x>-1\) and every \(\delta\in\{0,1\}\), with \(u=x+1\),
\begin{equation}
 \boxed{\;
 \mathcal D_\delta(x)
 =\e^{A_\delta}\left(\frac u\e\right)^{u/2}\e^{\Omega(u)},
 \qquad\text{equivalently}\qquad
 \log\mathcal D_\delta(x)=\Phi(u)+A_\delta+\Omega(u). \;}
 \label{p2:eq:normal-form}
\end{equation}
This is an identity, not an asymptotic statement, and \(\Omega\) does not
depend on \(\delta\).
\end{theorem}

\begin{proof}
Insert \(x=u-1\) into \cref{p2:eq:branches}:
\[
 \log\mathcal D_\delta(x)
 =\log\kappa_\delta+\frac{u-1}2\log2
  +\log\Gamma\!\left(\frac{u+1}2\right).
\]
Subtract \(A_\delta+\Phi(u)\) and use \(A_\delta=\log\kappa_\delta+\tfrac12\log\pi\):
\begin{align*}
 \log\mathcal D_\delta(x)-A_\delta-\Phi(u)
 &=\frac{u-1}2\log2+\log\Gamma\!\left(\frac{u+1}2\right)
   -\frac12\log\pi-\frac u2\log u+\frac u2\\
 &=\log\Gamma\!\left(\frac{u+1}2\right)
   -\frac u2\left(\log u-\log2\right)+\frac u2
   -\frac12\log2-\frac12\log\pi,
\end{align*}
which is exactly \(\Omega(u)\) as defined in \cref{p2:eq:Omega-def}.
\end{proof}

\begin{remark}[Why \(u=x+1\)]
\label{p2:rem:why-centre}
The shift is not cosmetic.  It turns the gamma argument into
\(\tfrac12(u+1)=\tfrac u2+\tfrac12\), and Bernoulli polynomials evaluated at
\(\tfrac12\) annihilate every odd index; consequently \(\Omega\) has only
\emph{odd} inverse powers of \(u\) (\cref{p2:thm:coefficients}) and the
explicit \(\tfrac12\log x\) of the uncentred Stirling form
(\cref{p2:prop:uncentred}) is absorbed into the dominant phase.  This sparse
support is precisely why the inverse lives on the grid
\(T^{-1},T^{-3},T^{-5},\dots\).
\end{remark}

At an integer \(n\) the two constants \(A_0,A_1\) combine into one parity
transmonomial: writing \(\chi_n=(-1)^n\),
\begin{equation}
 A(n)=\frac14\log(2\pi)+\frac{\chi_n}4\log\frac\pi2,
 \qquad
 \e^{A(n)}=(2\pi)^{1/4}\left(\frac\pi2\right)^{\chi_n/4},
 \label{p2:eq:parity-transmonomial}
\end{equation}
so that for every integer \(n\ge0\)
\begin{equation}
 n!!=\left(\frac{n+1}\e\right)^{(n+1)/2}
 \exp\!\bigl(A(n)+\Omega(n+1)\bigr).
 \label{p2:eq:sequence-exact}
\end{equation}
\Cref{p2:eq:sequence-exact} is exact.  The factor \((-1)^n\) is a genuine
transmonomial of the sequence's leading behaviour, not an error term.

\begin{remark}[Notation across the three sources]
\label{p2:rem:notation-clash}
The letter \(C\) means three different things in the siblings, and two Greek
letters collide:
\begin{center}
\small
\begin{tabular}{lccc}
\toprule
this chapter & \textsf{S1} & \textsf{S2} & \textsf{S3}\\
\midrule
\(\kappa_\delta=(2/\pi)^{\delta/2}\) & \(\kappa_\delta\) & \(\kappa_r\) & \(C_\varepsilon\)\\
\(\e^{A_\delta}\in\{\sqrt\pi,\sqrt2\}\) & \(C_\delta\) & \(A_r\) & ---\\
\(A_\delta\in\{\tfrac12\log\pi,\tfrac12\log2\}\) & \(\log C_\delta\) & \(C_r\) & \(A_\varepsilon\)\\
\(c_k\) (centred log.\ coefficient) & \(c_k\) & \(\beta_k\) & \(c_k\)\\
\(a_k=2c_k\) & \(2c_k\) & \(\alpha_k\) & \(a_k\)\\
\(a=\tfrac14\log(\pi/2)\) & \(\rho/2\) & \(a\) & ---\\
\(P_r(L_2)\) (Lambert polynomial) & \(P_r\) & \(u_{r+1}\) & \(\alpha_{r+1}\)\\
\(\tau_m(L_2)\) (reciprocal poly.) & \(h_m\) & \(\tau_m\) & \(\beta_m\)\\
\bottomrule
\end{tabular}
\end{center}
Part~0 adds one further collision: \cref{p0:def:core} writes the factorial
core as \(\Phi_0(X)=X(\kappa\log X+d)\), and that \(\kappa\) is \emph{not} the
branch prefactor \(\kappa_\delta\) of \textsf{S1} and \textsf{S2}.  This chapter
therefore never abbreviates the core parameters, always writing
\(\Phi_0(u)=u\bigl(\tfrac12\log u-\tfrac12\bigr)\) in full; see
\cref{p0:rem:factorial-core}.
In particular \textsf{S2}'s \(\alpha_k,\beta_k\) are Bernoulli data while
\textsf{S3}'s \(\alpha_j,\beta_n\) are outer Lambert polynomials.  Readers
comparing the sources should translate before comparing formulae; the numbers
themselves agree everywhere (\cref{p2:sec:numerics}).
\end{remark}

\section{Parameter dictionary for the Part~0 apparatus}
\label{p2:sec:dictionary}

The chapter uses the Part~0 toolkit rather than re-deriving it.  Its dominant
core is the \emph{factorial} one of \cref{p0:def:core}, not the
exponential--power model \cref{p0:def:model}; see
\cref{p2:rem:not-exponential-power}.  The specialisation is the following.

\begin{center}
\small
\begin{tabular}{p{0.36\textwidth}p{0.27\textwidth}p{0.27\textwidth}}
\toprule
Part~0 datum & even branch \(\delta=0\) & odd branch \(\delta=1\)\\
\midrule
interpolation & \(2^{x/2}\Gamma(\tfrac x2+1)\)
              & \(\sqrt{2/\pi}\;2^{x/2}\Gamma(\tfrac x2+1)\)\\
affine shift \(x_*\) of \cref{p0:def:core} & \(x_*=-1\), taken first & \(x_*=-1\), taken first\\
centred coordinate \(u=x-x_*\) & \(u=x+1\) & \(u=x+1\)\\
admissible core \(\Phi_0\) of \cref{p0:def:core}
  & \(\Phi_0(u)=u\bigl(\tfrac12\log u-\tfrac12\bigr)\) & same\\
additive constant & \(A_0=\tfrac12\log\pi\) & \(A_1=\tfrac12\log2\)\\
perturbation \(Q(s)=\sum_{k\ge1}a_ks^{2k-1}\) & \(a_k=2c_k\), \cref{p2:eq:ck} & same\\
target after constant removal & \(R_0=\log y-\tfrac12\log\pi\)
                              & \(R_1=\log y-\tfrac12\log2\)\\
core (\cref{p0:prop:factorial-core}) & \(T_\delta=2R_\delta/\Wlam_0(2R_\delta/\e)\) & same formula\\
core slope parameter & \(\Lambda_\delta=\log T_\delta\) & same formula\\
inner grid & \(q=T_\delta^{-1}\), \(z=q^2\) & same\\
\bottomrule
\end{tabular}
\end{center}

Concretely: \cref{p0:def:obell,p0:def:ebell} supply the Bell polynomials used
throughout; \cref{p0:lem:power-log} supplies the coefficients of
\((1+w)^{-j}\) and \(\log(1+w)\); \cref{p0:prop:factorial-core} inverts the
dominant phase exactly by \(\Wlam_0\); \cref{p0:thm:core-reversion}, whose
formal engine is \cref{p0:lem:lagrange-fp}, generates the correction blocks and
is specialised once in \cref{p2:lem:normalized};
\cref{p0:thm:flattening} unfolds \(\Wlam_0\) into the outer layer of
\cref{p2:sec:outer}; \cref{p0:def:three-inverses,p0:thm:staircase} govern
\cref{p2:sec:discrete}; \cref{p0:thm:backward-error,p0:thm:remainder-transport}
convert forward residuals into inverse errors; and
\cref{p0:thm:optimal-truncation} fixes the least-term scale of
\cref{p2:sec:beyond}.

\begin{warning}[The double factorial is not of exponential--power type]
\label{p2:rem:not-exponential-power}
It is tempting --- and wrong --- to file \(\mathcal D_\delta\) under
\cref{p0:def:model}, \(F(x)=C\exp(ax^\alpha)x^b\exp Q(x^{-\delta})\) with
\(\alpha\in(0,1]\).  By \cref{p2:thm:normal-form},
\(\log\mathcal D_\delta(x)=\tfrac u2(\log u-1)+A_\delta+\Omega(u)\) with
\(u=x+1\), so the dominant phase is \(\tfrac12u\log u\).  For every
\(\alpha\le1\),
\[
 \frac{u\log u}{u^\alpha}\longrightarrow\infty
 \qquad(u\to+\infty),
\]
so no admissible choice of \((C,a,\alpha,b,\delta,Q)\) reproduces it: the
double factorial lies outside the exponential--power family, in the
\emph{factorial} class of \cref{p0:def:core}, with the parameters that
\cref{p0:rem:factorial-core} records for this chapter.  This is
the structural difference from
\cref{p1:sec:top,p3:sec:top,p4:sec:top}, whose sequences have
\(\log a_n=an^\alpha+b\log n+\cdots\) with \(\alpha\le1\).  What survives from
the exponential--power theory is the \emph{shape} of the reversion, not the
model: \cref{p0:thm:core-reversion} covers both cores, and the resulting
fixed-point map is \cref{p2:eq:master-fp}.  Note also that the shift
\(x_*=-1\) must be applied \emph{before} any further reduction: it is what
removes the \(\tfrac12\log x\) of \cref{p2:eq:uncentred} and produces the
sparse odd-power perturbation (\cref{p2:rem:why-centre}).
None of \textsf{S1}, \textsf{S2}, \textsf{S3} claims membership of the
exponential--power family; all three invert \(\tfrac12u(\log u-1)\) directly.
\end{warning}

\section{Borel structure of the centred correction}
\label{p2:sec:borel}

\begin{sourcebox}
The exact Borel kernel is due to \textsf{S2} and \textsf{S3}, independently and
in the same form; \textsf{S1} has only the formal Bernoulli series.  The sharp
first-omitted-term bound of \cref{p2:thm:remainder} is \textsf{S2}'s;
\textsf{S3} proves the weaker bound with \(\eta(2N+2)\) replaced by \(1\), and
that weaker form is not used here.
\end{sourcebox}

\begin{theorem}[Borel--Laplace representation]
\label{p2:thm:borel}
For \(\Re u>0\),
\begin{equation}
 \Omega(u)=\int_0^\infty\e^{-us}\,\widehat\Omega(s)\,\dd s,
 \qquad
 \widehat\Omega(s)=\frac1{2s^2}\left(\frac s{\sinh s}-1\right),
 \label{p2:eq:borel}
\end{equation}
the singularity of \(\widehat\Omega\) at \(s=0\) being removable.  Moreover
\begin{equation}
 \widehat\Omega(s)=\sum_{m\ge1}\frac{(-1)^m}{s^2+\pi^2m^2},
 \label{p2:eq:mittag}
\end{equation}
so the Borel singularities of \(\Omega\) are exactly the simple poles
\(s=\pm\ii\pi m\), \(m\ge1\), with residues
\begin{equation}
 \Res_{s=\ii\pi m}\widehat\Omega(s)=\frac{(-1)^{m+1}\ii}{2\pi m},
 \qquad
 \Res_{s=-\ii\pi m}\widehat\Omega(s)=\frac{(-1)^{m}\ii}{2\pi m}.
 \label{p2:eq:residues}
\end{equation}
\end{theorem}

\begin{proof}
Put \(t=u/2\) and
\(\widetilde\Omega(t)=\log\Gamma(t+\tfrac12)-t\log t+t-\tfrac12\log(2\pi)\),
so that \(\Omega(u)=\widetilde\Omega(u/2)\) by \cref{p2:eq:Omega-def}.
Binet's first integral and Frullani's formula give, for \(t>0\),
\[
 \psi\!\left(t+\tfrac12\right)-\log t
 =\int_0^\infty\e^{-tr}
 \left(\frac1r-\frac1{2\sinh(r/2)}\right)\dd r,
\]
the integrand being \(\bigO(r)\) at the origin and exponentially damped at
infinity.  On the other hand
\[
 \frac{\dd}{\dd t}\int_0^\infty\e^{-tr}
 \frac{r/(2\sinh(r/2))-1}{r^2}\,\dd r
 =\int_0^\infty\e^{-tr}
 \left(\frac1r-\frac1{2\sinh(r/2)}\right)\dd r ,
\]
differentiation under the integral being justified by the same bounds.  Since
\(\widetilde\Omega'(t)=\psi(t+\tfrac12)-\log t\) and both
\(\widetilde\Omega(t)\) and the displayed integral tend to \(0\) as
\(t\to+\infty\), the two functions coincide.  Substituting \(r=2s\) and
\(t=u/2\) turns the integral into \cref{p2:eq:borel}; analytic continuation in
\(u\) extends it to \(\Re u>0\).  The expansion \(s/\sinh s=1-s^2/6+\bigO(s^4)\)
removes the origin.

For \cref{p2:eq:mittag}, the elementary Mittag--Leffler expansion
\(\dfrac s{\sinh s}=1+2s^2\sum_{m\ge1}\dfrac{(-1)^m}{s^2+\pi^2m^2}\)
gives the claim after division by \(2s^2\).  The residues follow from
\(1/(s^2+\pi^2m^2)=\bigl[(s-\ii\pi m)(s+\ii\pi m)\bigr]^{-1}\).
\end{proof}

\begin{theorem}[The centred logarithmic coefficients]
\label{p2:thm:coefficients}
As \(u\to\infty\) in any closed subsector of \(|\arg u|<\pi/2\),
\begin{equation}
 \Omega(u)\sim\sum_{k\ge1}\frac{c_k}{u^{2k-1}},
 \label{p2:eq:Omega-series}
\end{equation}
where the following three expressions coincide:
\begin{equation}
 \boxed{\;
 c_k=\frac{2^{2k-1}B_{2k}\!\left(\tfrac12\right)}{2k(2k-1)}
 =\frac{\bigl(1-2^{2k-1}\bigr)B_{2k}}{2k(2k-1)}
 =(-1)^k\frac{(2k-2)!\,\eta(2k)}{\pi^{2k}} \;}
 \label{p2:eq:ck}
\end{equation}
with \(B_m(z)\) the Bernoulli polynomials, \(B_m=B_m(0)\), and
\(\eta(s)=(1-2^{1-s})\zeta(s)\) the Dirichlet eta function.  There are no even
inverse powers.
\end{theorem}

\begin{proof}
Expand each summand of \cref{p2:eq:mittag} by
\begin{equation}
 \frac1{s^2+A^2}
 =\sum_{j=0}^{N-1}\frac{(-1)^js^{2j}}{A^{2j+2}}
 +\frac{(-1)^Ns^{2N}}{A^{2N}(s^2+A^2)},
 \qquad A>0,
 \label{p2:eq:partial-fraction-remainder}
\end{equation}
with \(A=\pi m\), sum over \(m\), and integrate term by term against
\(\e^{-us}\), using \(\int_0^\infty\e^{-us}s^{2j}\dd s=(2j)!\,u^{-2j-1}\).
The coefficient of \(u^{-(2j+1)}\) is
\[
 (2j)!\sum_{m\ge1}\frac{(-1)^m(-1)^j}{(\pi m)^{2j+2}}
 =(-1)^{j+1}\frac{(2j)!\,\eta(2j+2)}{\pi^{2j+2}},
\]
which is the third expression in \cref{p2:eq:ck} with \(k=j+1\).  Euler's
evaluation \(\zeta(2k)=(-1)^{k+1}(2\pi)^{2k}B_{2k}/\bigl(2(2k)!\bigr)\)
together with \(\eta(2k)=(1-2^{1-2k})\zeta(2k)\) and
\(B_{2k}(\tfrac12)=(2^{1-2k}-1)B_{2k}\) converts it into the first two.  Only
odd powers \(u^{-(2k-1)}\) occur because \(\widehat\Omega\) is even.
\end{proof}

\begin{theorem}[Sharp alternating remainder]
\label{p2:thm:remainder}
For real \(u>0\) and \(N\ge0\) put
\(R_N(u)=\Omega(u)-\sum_{k=1}^N c_ku^{-(2k-1)}\) (empty sum for \(N=0\)).
Then \(R_N(u)\) has the sign of the first omitted term and is smaller than it
in modulus:
\begin{equation}
 0<(-1)^{N+1}R_N(u)<\frac{|c_{N+1}|}{u^{2N+1}}
 =\frac{(2N)!\,\eta(2N+2)}{\pi^{2N+2}u^{2N+1}} .
 \label{p2:eq:remainder}
\end{equation}
\end{theorem}

\begin{proof}
Insert \cref{p2:eq:partial-fraction-remainder} with \(A=\pi m\) into
\cref{p2:eq:mittag} and Laplace-transform.  The residual kernel is
\((-1)^{N+1}s^{2N}S_N(s)\) with
\[
 S_N(s)=\sum_{m\ge1}\frac{(-1)^{m-1}}
 {(\pi m)^{2N}\bigl((\pi m)^2+s^2\bigr)} .
\]
Writing \(\bigl((\pi m)^2+s^2\bigr)^{-1}=\int_0^\infty
\e^{-s^2v}\e^{-\pi^2m^2v}\dd v\) gives
\[
 S_N(s)=\int_0^\infty\e^{-s^2v}
 \sum_{m\ge1}\frac{(-1)^{m-1}\e^{-\pi^2m^2v}}{(\pi m)^{2N}}\,\dd v .
\]
For each \(v\ge0\) the inner series is alternating with strictly decreasing
positive terms, hence strictly positive.  Therefore \(S_N(s)>0\) and
\(S_N\) is strictly decreasing in \(s^2\), so
\(0<S_N(s)\le S_N(0)=\eta(2N+2)/\pi^{2N+2}\).  Multiplying by
\(s^{2N}\), integrating against \(\e^{-us}\) and using
\(\int_0^\infty\e^{-us}s^{2N}\dd s=(2N)!u^{-(2N+1)}\) gives
\cref{p2:eq:remainder}.
\end{proof}

\begin{corollary}[Certified forward logarithm]
\label{p2:cor:forward-log}
For \(x\to+\infty\), \(u=x+1\) and either parity,
\begin{equation}
 \log\mathcal D_\delta(x)
 =\Phi(u)+A_\delta+\sum_{k=1}^N\frac{c_k}{u^{2k-1}}
  +\theta_N\,\frac{c_{N+1}}{u^{2N+1}}
 \qquad\text{for some }\theta_N=\theta_N(u)\in(0,1).
 \label{p2:eq:forward-certified}
\end{equation}
Thus truncating after any number of terms overshoots or undershoots by
strictly less than the first omitted term, with the sign of that term.  For the
integer sequence replace \(A_\delta\) by \(A(n)\) of
\cref{p2:eq:parity-transmonomial}.
\end{corollary}

\begin{proof}
Combine \cref{p2:thm:normal-form,p2:thm:remainder}.
\end{proof}

\begin{center}
\small
\begin{longtable}{c r r}
\caption{The centred logarithmic coefficients \(c_k\) of \cref{p2:eq:ck} and
the doubled coefficients \(a_k=2c_k\) used in the inverse problem.  Regenerated
in exact rational arithmetic from \(B_{2k}\).}
\label{p2:tab:ck}\\
\toprule
\(k\) & \(c_k\) & \(a_k=2c_k\)\\
\midrule
\endfirsthead
\toprule
\(k\) & \(c_k\) & \(a_k=2c_k\)\\
\midrule
\endhead
1 & \(-1/12\) & \(-1/6\)\\
2 & \(7/360\) & \(7/180\)\\
3 & \(-31/1260\) & \(-31/630\)\\
4 & \(127/1680\) & \(127/840\)\\
5 & \(-511/1188\) & \(-511/594\)\\
6 & \(1414477/360360\) & \(1414477/180180\)\\
7 & \(-8191/156\) & \(-8191/78\)\\
8 & \(118518239/122400\) & \(118518239/61200\)\\
9 & \(-5749691557/244188\) & \(-5749691557/122094\)\\
\bottomrule
\end{longtable}
\end{center}

\begin{remark}[Late coefficients]
\label{p2:rem:late}
From the third form in \cref{p2:eq:ck} and \(1<\eta(2k)<1+2^{-2k}\cdot\tfrac43\),
\begin{equation}
 c_k=(-1)^k\frac{(2k-2)!}{\pi^{2k}}\bigl(1+\bigO(4^{-k})\bigr),
 \label{p2:eq:ck-late}
\end{equation}
so \cref{p2:eq:Omega-series} is Gevrey-\(1\) in the variable \(u^{-2}\).  The
ratio of consecutive term moduli is \(\sim2k(2k-1)/(\pi u)^2\); the least term
occurs near \(2k=\pi u\) and has size \(\asymp u^{-1/2}\e^{-\pi u}\), matching
the distance \(\pi\) from the origin to the nearest Borel poles
\(\pm\ii\pi\).  This is the instance of \cref{p0:thm:optimal-truncation}
relevant here; see \cref{p2:sec:beyond}.
\end{remark}

\section{The forward transseries}
\label{p2:sec:forward}

\subsection{Multiplicative, reciprocal, and arbitrary-power coefficients}

Exponentiation of \cref{p2:eq:Omega-series} is an instance of the exponential
formula, so \cref{p0:def:ebell} does all the work at once for every power.

\begin{definition}
\label{p2:def:gamma-sparse}
Embed the \(c_k\) sparsely on the odd indices,
\begin{equation}
 \gamma_r=\begin{cases}c_{(r+1)/2},&r\ \text{odd},\\ 0,&r\ \text{even},\end{cases}
 \qquad r\ge1,
 \label{p2:eq:gamma-sparse}
\end{equation}
and for a scalar \(\alpha\) define \(s^{(\alpha)}_m\) by
\begin{equation}
 \exp\!\left(\alpha\sum_{r\ge1}\gamma_r\,u^{-r}\right)
 =\sum_{m\ge0}\frac{s^{(\alpha)}_m}{u^m},
 \qquad s_m:=s^{(1)}_m .
 \label{p2:eq:salpha-def}
\end{equation}
\end{definition}

\begin{proposition}[All multiplicative coefficients at once]
\label{p2:prop:mult-coeff}
For every \(\alpha\) and every \(m\ge0\),
\begin{align}
 s^{(\alpha)}_m
 &=\frac1{m!}\,\EB_m\bigl(\alpha\,1!\,\gamma_1,\alpha\,2!\,\gamma_2,\ldots,
     \alpha\,m!\,\gamma_m\bigr)
 \label{p2:eq:salpha-bell}\\
 &=\sum_{\substack{k_1,k_2,\ldots\ge0\\ \sum_{r\ge1}rk_r=m}}
   \prod_{r\ge1}\frac{(\alpha\gamma_r)^{k_r}}{k_r!}
 =\sum_{\substack{k_1,k_2,\ldots\ge0\\ \sum_{j\ge1}(2j-1)k_j=m}}
   \prod_{j\ge1}\frac{(\alpha c_j)^{k_j}}{k_j!},
 \label{p2:eq:salpha-partition}
\end{align}
and equivalently
\begin{equation}
 s^{(\alpha)}_0=1,
 \qquad
 m\,s^{(\alpha)}_m
 =\alpha\!\!\sum_{1\le j\le(m+1)/2}\!\!(2j-1)\,c_j\,s^{(\alpha)}_{m-2j+1}
 \qquad(m\ge1).
 \label{p2:eq:salpha-recurrence}
\end{equation}
Consequently, for \(x\to+\infty\), \(u=x+1\) and every fixed \(M\ge0\),
\begin{equation}
 \mathcal D_\delta(x)^\alpha
 =\e^{\alpha A_\delta}\left(\frac u\e\right)^{\alpha u/2}
 \left(\sum_{m=0}^{M}\frac{s^{(\alpha)}_m}{u^m}
   +\bigO\!\left(u^{-(M+1)}\right)\right).
 \label{p2:eq:mult-transseries}
\end{equation}
The cases \(\alpha=1\) and \(\alpha=-1\) give \(\mathcal D_\delta\) itself and
its reciprocal.
\end{proposition}

\begin{proof}
Put \(t=u^{-1}\) and \(S_\alpha(t)=\exp(\alpha\sum_r\gamma_rt^r)\).
\Cref{p2:eq:salpha-bell} is the complete-Bell generating identity of
\cref{p0:def:ebell}; multinomial expansion of the exponential gives
\cref{p2:eq:salpha-partition}, and the second form there uses
\(\gamma_{2j-1}=c_j\), \(\gamma_{2j}=0\).  Logarithmic differentiation
\(tS_\alpha'=\alpha\bigl(\sum_r r\gamma_rt^r\bigr)S_\alpha\) and comparison of
coefficients give \cref{p2:eq:salpha-recurrence}.  For
\cref{p2:eq:mult-transseries}, take \(N\) with \(2N-1\ge M\); by
\cref{p2:thm:normal-form,p2:thm:remainder},
\(\alpha\Omega(u)=\alpha\sum_{k\le N}c_ku^{1-2k}+\bigO(u^{-(2N+1)})\), and
\(\exp\) of a finite sum of \(\bigO(u^{-1})\) terms plus a
\(\bigO(u^{-(2N+1)})\) error equals \(\sum_{m\le M}s^{(\alpha)}_mu^{-m}
+\bigO(u^{-(M+1)})\), the exponential being Lipschitz on bounded sets.
\end{proof}

\begin{corollary}[Reciprocal sign symmetry]
\label{p2:cor:reciprocal-sign}
\(s^{(-\alpha)}_m=(-1)^m s^{(\alpha)}_m\) for all \(m\ge0\); in particular the
reciprocal coefficients are \(s^{(-1)}_m=(-1)^ms_m\).
\end{corollary}

\begin{proof}
Only odd \(r\) carry a nonzero \(\gamma_r\), so
\(\sum_r\gamma_r(-t)^r=-\sum_r\gamma_rt^r\) and hence
\(S_\alpha(-t)=S_{-\alpha}(t)\).  Reading off \(\Coef{t^m}\) gives the claim.
\end{proof}

\begin{center}
\small
\begin{longtable}{c r r}
\caption{Centred multiplicative coefficients \(s_m=s^{(1)}_m\) of
\cref{p2:eq:mult-transseries} and the uncentred coefficients \(g_m\) of
\cref{p2:prop:uncentred}.  Regenerated from
\cref{p2:eq:salpha-recurrence} in exact rational arithmetic.}
\label{p2:tab:sm}\\
\toprule
\(m\) & \(s_m\) (centred) & \(g_m\) (uncentred)\\
\midrule
\endfirsthead
\toprule
\(m\) & \(s_m\) (centred) & \(g_m\) (uncentred)\\
\midrule
\endhead
0 & \(1\) & \(1\)\\
1 & \(-1/12\) & \(1/6\)\\
2 & \(1/288\) & \(1/72\)\\
3 & \(1003/51840\) & \(-139/6480\)\\
4 & \(-4027/2488320\) & \(-571/155520\)\\
5 & \(-5128423/209018880\) & \(163879/6531840\)\\
6 & \(168359651/75246796800\) & \(5246819/1175731200\)\\
7 & \(68168266699/902961561600\) & \(-534703531/7054387200\)\\
8 & \(-587283555451/86684309913600\) & \(-4483131259/338610585600\)\\
\bottomrule
\end{longtable}
\end{center}

\subsection{The uncentred normal form and the centre-shift identity}

\begin{proposition}[Uncentred expansion]
\label{p2:prop:uncentred}
As \(x\to+\infty\),
\begin{equation}
 \log\mathcal D_\delta(x)
 \sim\frac x2(\log x-1)+\frac12\log x+A_\delta
 +\sum_{k\ge1}\frac{b_k}{x^{2k-1}},
 \qquad
 b_k=\frac{2^{2k-1}B_{2k}}{2k(2k-1)}
   =(-1)^{k-1}\frac{(2k-2)!\,\zeta(2k)}{\pi^{2k}} ,
 \label{p2:eq:uncentred}
\end{equation}
with \(b_1=\tfrac16\), \(b_2=-\tfrac1{45}\), \(b_3=\tfrac8{315}\),
\(b_4=-\tfrac8{105}\), \(b_5=\tfrac{128}{297}\).  Exponentiating,
\(\mathcal D_\delta(x)\sim\e^{A_\delta}\sqrt x\,(x/\e)^{x/2}\sum_m g_mx^{-m}\),
where \(g_m\) is obtained from \cref{p2:eq:salpha-recurrence} on replacing
\(c_j\) by \(b_j\); this is the familiar form, with amplitudes \(\sqrt\pi\)
(even) and \(\sqrt2\) (odd).  The exact Borel kernel of the uncentred
correction is
\(\widehat{\,\mathcal U}(s)=\tfrac{\coth s}{2s}-\tfrac1{2s^2}
=\sum_{m\ge1}\bigl(s^2+\pi^2m^2\bigr)^{-1}\), with the \emph{same} pole
lattice \(\pm\ii\pi m\) but all residues of one sign pattern.
\end{proposition}

\begin{proof}
Apply the standard Stirling expansion to \(\log\Gamma(z+1)\) with
\(z=x/2\) and add \((x/2)\log2+\log\kappa_\delta\); the elementary terms give
the first three summands of \cref{p2:eq:uncentred}, and the corrections scale
as \(z^{-(2k-1)}=2^{2k-1}x^{-(2k-1)}\).  Euler's evaluation of \(\zeta(2k)\)
gives the second form of \(b_k\).  For the kernel, the uncentred correction is
\(\mathcal U(x)=\Omega_\Gamma(x/2)\) with \(\Omega_\Gamma\) Binet's function;
substituting \(t=2s\) in
\(\Omega_\Gamma(z)=\int_0^\infty\e^{-zt}
\bigl(\tfrac{\coth(t/2)}{2t}-\tfrac1{t^2}\bigr)\dd t\) gives the stated kernel,
and \(\tfrac{\coth s}s=\tfrac1{s^2}+2\sum_{m\ge1}(s^2+\pi^2m^2)^{-1}\) gives
the partial fractions.
\end{proof}

\begin{warning}
\label{p2:rem:two-normal-forms}
The two statements
\(n!!\sim \e^{A(n)}\sqrt n\,(n/\e)^{n/2}\) and
\(n!!\sim \e^{A(n)}\bigl((n+1)/\e\bigr)^{(n+1)/2}\)
are both correct and are not in conflict: they use different normal forms.  The
first leaves a factor \(\sqrt n\) and an \emph{uncentred} correction series;
the second absorbs that factor into the phase.  Only the centred form leads to
the sparse inverse grid, and it is the one used from
\cref{p2:sec:branch-inverse} onwards.
\end{warning}

\begin{proposition}[Centre-shift identity]
\label{p2:prop:centre-shift}
For every \(r\ge1\),
\begin{equation}
 \frac{(-1)^{r+1}}{2r(r+1)}
 +(-1)^{r+1}\!\!\sum_{k=1}^{\lfloor(r+1)/2\rfloor}\!\!
 \binom{r-1}{2k-2}c_k
 =\begin{cases}b_{(r+1)/2},& r\ \text{odd},\\[0.2em] 0,& r\ \text{even}.\end{cases}
 \label{p2:eq:centre-shift}
\end{equation}
\end{proposition}

\begin{proof}
Substitute \(u=x+1=x(1+w)\), \(w=x^{-1}\), into
\(\Phi(u)+\sum_kc_ku^{1-2k}\) and compare with \cref{p2:eq:uncentred}.
First, with \(\log u=\log x+\log(1+w)\),
\[
 \Phi(u)=\frac x2(\log x-1)+\frac12(\log x-1)
 +\frac{x+1}2\log(1+w).
\]
Expanding \(\log(1+w)=\sum_{j\ge1}(-1)^{j-1}w^j/j\) and splitting the prefactor
\(\tfrac{x+1}2=\tfrac x2+\tfrac12\) gives
\[
 \frac{x+1}2\log(1+w)
 =\frac12+\frac12\sum_{r\ge1}(-1)^rx^{-r}
   \left(\frac1{r+1}-\frac1r\right)
 =\frac12+\sum_{r\ge1}\frac{(-1)^{r+1}}{2r(r+1)}\,x^{-r},
\]
so that \(\Phi(u)=\tfrac x2(\log x-1)+\tfrac12\log x
+\sum_{r\ge1}(-1)^{r+1}\bigl(2r(r+1)\bigr)^{-1}x^{-r}\); the two halves
\(\pm\tfrac12\) cancel.  Second,
\[
 \sum_{k\ge1}\frac{c_k}{u^{2k-1}}
 =\sum_{k\ge1}c_k\,x^{-(2k-1)}(1+w)^{-(2k-1)}
 =\sum_{r\ge1}x^{-r}\sum_{k\le(r+1)/2}c_k
   \binom{-(2k-1)}{r-2k+1},
\]
and \(\binom{-(2k-1)}{r-2k+1}=(-1)^{r-2k+1}\binom{r-1}{r-2k+1}
=(-1)^{r+1}\binom{r-1}{2k-2}\).  Adding the two coefficient lists and matching
against \cref{p2:eq:uncentred} --- where the coefficient of \(x^{-r}\) is
\(b_{(r+1)/2}\) for odd \(r\) and \(0\) for even \(r\) --- gives
\cref{p2:eq:centre-shift}.
\end{proof}

\begin{remark}
\Cref{p2:eq:centre-shift} was verified symbolically for \(1\le r\le9\): the
even cases return \(0\) and the odd cases return
\(\tfrac16,-\tfrac1{45},\tfrac8{315},-\tfrac8{105},\tfrac{128}{297}\).  It is
the cheapest available certificate that the centring is consistent, because it
fails immediately under any sign or index slip in \cref{p2:eq:ck}.  \textsf{S3}
states the identity without proof; the proof above is supplied here.
\end{remark}

\section{The branchwise functional inverse}
\label{p2:sec:branch-inverse}

\begin{sourcebox}
All three siblings reach the same objects in this section.  The presentation
follows \textsf{S3} for the exact normal form and \textsf{S2} for the graded
closed formula and the sign law; \textsf{S1}'s ``scaling from the half-shifted
inverse Gamma problem'' is subsumed by \cref{p2:thm:multi-scaling}.  The proofs
of \cref{p2:cor:boundary,p2:thm:sign-law} are new: the corresponding source
arguments are sketches (see the chapter report).
\end{sourcebox}

\subsection{The Lambert core}

\begin{lemma}[Exact inversion of the dominant phase]
\label{p2:lem:core}
Let \(R>0\).  The equation \(\Phi(T)=\tfrac12T(\log T-1)=R\) has a unique
solution \(T>\e\), namely
\begin{equation}
 w=\Wlam_0\!\left(\frac{2R}\e\right),
 \qquad
 T=\frac{2R}w=\e^{\,w+1},
 \qquad
 \Lambda:=w+1=\log T,
 \label{p2:eq:lambert-core}
\end{equation}
and then
\begin{equation}
 \Phi'(T)=\frac\Lambda2,
 \qquad
 \frac{\dd T}{\dd R}=\frac2\Lambda,
 \qquad
 T(\Lambda-1)=2R .
 \label{p2:eq:core-identities}
\end{equation}
\end{lemma}

\begin{proof}
This is \cref{p0:prop:factorial-core} for the parameters of
\cref{p0:rem:factorial-core}; for completeness, setting \(w=\log T-1\) turns
\(\Phi(T)=R\) into \(w\e^{\,w+1}=2R\), i.e.\ \(w\e^{w}=2R/\e>0\), whose unique
real solution is \(w=\Wlam_0(2R/\e)>0\), the principal branch being the only
real branch on \((0,\infty)\).  \(\Phi\) is strictly increasing on
\((1,\infty)\) with range \((-\tfrac12,\infty)\), which gives uniqueness of
\(T\).  The identities in \cref{p2:eq:core-identities} follow from
\(\Phi'(T)=\tfrac12\log T\), implicit differentiation of \(\Phi(T)=R\), and
\(T=2R/w=2R/(\Lambda-1)\).
\end{proof}

For the parity branch \(\delta\) one takes
\begin{equation}
 R_\delta=\log y-A_\delta,
 \qquad
 w_\delta,\;T_\delta,\;\Lambda_\delta\ \text{as in \cref{p2:eq:lambert-core}} .
 \label{p2:eq:Rdelta}
\end{equation}
The subscript \(\delta\) is suppressed inside coefficient calculations, where
it never appears.

\subsection{The normalised displacement equation}

\begin{lemma}[The factorial-core reversion kernel]
\label{p2:lem:normalized}
Write \(u=T(1+E)\), \(q=T^{-1}\), \(z=q^2\), and
\begin{equation}
 \varphi(v)=(1+v)\log(1+v)-v=\sum_{m\ge2}\frac{(-1)^m}{m(m-1)}v^m,
 \qquad
 a_k=2c_k .
 \label{p2:eq:varphi}
\end{equation}
Then the equation \(\Phi(u)+\Omega(u)=R\), with \(\Omega\) replaced by its
asymptotic series, is equivalent to the formal normal form
\begin{equation}
 \boxed{\;
 \Lambda E+\varphi(E)+\sum_{k\ge1}a_k\,z^k(1+E)^{-(2k-1)}=0 . \;}
 \label{p2:eq:master}
\end{equation}
Equivalently, in absolute displacement \(U=TE\) with \(t=q\),
\begin{equation}
 U=\Psi(t,U),
 \qquad
 \Psi(t,v)=-\frac1\Lambda\left[
 \frac{\varphi(tv)}t
 +\sum_{k\ge1}a_k\,t^{2k-1}(1+tv)^{-(2k-1)}\right]
 \in t\,\Rat(\Lambda)[[t,v]] ,
 \label{p2:eq:master-fp}
\end{equation}
which is the specialisation of \cref{p0:thm:core-reversion} to the factorial
core, so that \cref{p0:lem:lagrange-fp} applies verbatim.  There is a unique
formal solution, and it has the form
\begin{equation}
 E=\sum_{n\ge1}e_n(\Lambda)\,z^n,
 \qquad\text{hence}\qquad
 u-T=\sum_{n\ge1}\frac{e_n(\Lambda)}{T^{2n-1}} :
 \label{p2:eq:E-series}
\end{equation}
only \emph{odd} inverse powers of the Lambert core occur.
\end{lemma}

\begin{proof}
Because \(\Phi(T)=R\),
\[
 \Phi\bigl(T(1+E)\bigr)-\Phi(T)
 =\frac{T(1+E)}2\bigl(\log T+\log(1+E)-1\bigr)-\frac T2(\log T-1)
 =\frac T2\bigl[\Lambda E+\varphi(E)\bigr],
\]
using \(\Lambda=\log T\) and
\((1+E)\log(1+E)=\varphi(E)+E\).  Also
\(\Omega(u)=\sum_{k}c_kT^{-(2k-1)}(1+E)^{-(2k-1)}\) formally.  Multiplying the
whole equation by \(2/T=2q\) and using \(a_k=2c_k\) and
\(q\cdot q^{2k-1}=z^k\) gives \cref{p2:eq:master}.  Dividing by \(T\) and
solving for \(U=TE\) gives \cref{p2:eq:master-fp}; the bracket lies in
\(t\,\Rat[[t,v]]\) because \(\varphi(tv)=t^2v^2/2+\bigO(t^3)\), so
\(\varphi(tv)/t=\bigO(t)\), and every Bernoulli term carries \(t^{2k-1}\) with
\(k\ge1\).  Existence and uniqueness of the fixed point are
\cref{p0:lem:lagrange-fp}.  Finally, if \(E\) contains only even powers of
\(q\) then so does every term of \cref{p2:eq:master}; an induction on the
\(q\)-valuation therefore forces \(E\in z\,\Rat(\Lambda)[[z]]\).
\end{proof}

\subsection{The all-orders inverse and its remainder}

\begin{theorem}[Complete branchwise inverse transseries]
\label{p2:thm:branch-inverse}
Fix \(\delta\in\{0,1\}\) and let \(R_\delta,T_\delta,\Lambda_\delta\) be as in
\cref{p2:eq:Rdelta}.  Write \(d_{2n-1}:=e_n\) for the coefficients of
\cref{p2:eq:E-series}.  Then for every fixed \(N\ge0\), as \(y\to+\infty\),
\begin{equation}
 \boxed{\;
 \mathcal D_\delta^{-1}(y)
 =T_\delta-1
 +\sum_{n=1}^{N}\frac{d_{2n-1}(\Lambda_\delta)}{T_\delta^{2n-1}}
 +\bigO\!\left(\frac1{T_\delta^{2N+1}\Lambda_\delta}\right) , \;}
 \label{p2:eq:branch-inverse}
\end{equation}
with an implied constant that may be chosen uniformly in \(\delta\in\{0,1\}\).
All parity dependence sits in \(R_\delta=\log y-A_\delta\); the coefficient
functions \(d_{2n-1}\) are universal.
\end{theorem}

\begin{proof}
Let \(u_N=T+\sum_{n=1}^Ne_n(\Lambda)T^{1-2n}\) and \(E_N=(u_N-T)/T\).  By
construction, substituting \(E_N\) into \cref{p2:eq:master} annihilates the
coefficients of \(z^1,\dots,z^N\), so the formal residual is
\(\bigO(z^{N+1})\) with coefficients rational in \(\Lambda\) and bounded as
\(\Lambda\to\infty\).  Undoing the normalisation (multiplying by \(T/2\))
therefore bounds the exact forward residual
\[
 \rho_N:=\Phi(u_N)+A_\delta+\Omega(u_N)-\log y
\]
by \(\bigO\!\left(T\cdot z^{N+1}\right)=\bigO\!\left(T^{-(2N+1)}\right)\),
where replacing \(\Omega\) by its \(N\)-term truncation costs at most
\(|c_{N+1}|u_N^{-(2N+1)}\) by \cref{p2:thm:remainder}, of the same order.
This is the hypothesis of \cref{p0:thm:backward-error}.  For the slope, on the
interval between \(u_N\) and the exact root,
\[
 \frac{\dd}{\dd u}\bigl[\Phi(u)+\Omega(u)\bigr]
 =\frac12\log u+\Omega'(u)
 =\frac\Lambda2\bigl(1+\littleo(1)\bigr),
\]
because \(u=T\bigl(1+\bigO(T^{-2})\bigr)\) and, differentiating
\cref{p2:eq:borel} under the integral sign,
\(\Omega'(u)=-\int_0^\infty s\e^{-us}\widehat\Omega(s)\dd s=\bigO(u^{-2})\).
The mean-value theorem now divides \(\rho_N\) by a quantity asymptotic to
\(\Lambda/2\), which is the instance of \cref{p0:thm:remainder-transport} for
the factorial core and gives \cref{p2:eq:branch-inverse}.  The parity constant
\(A_\delta\) disappears on differentiation, so the estimate is uniform in
\(\delta\).
\end{proof}

\subsection{Two generators for the coefficients}

\Cref{p2:eq:master} is exactly \cref{p0:eq:core-master}, the reversion
problem of \cref{p0:thm:core-reversion}, with the instance data
\begin{equation}
 h(v)=\varphi(v),\quad h_m=\frac{(-1)^m}{m(m-1)},
 \qquad
 f(z,v)=\sum_{k\ge1}\phi_k(z)(1+v)^{-\mu_k},
 \quad \phi_k(z)=a_kz^k,\quad \mu_k=2k-1 .
 \label{p2:eq:instance-data}
\end{equation}
Its triangular recurrence \cref{p0:eq:core-triangular} and its Bell form
\cref{p0:eq:core-bell} therefore apply verbatim; written out, the former is
\begin{equation}
 e_n(\Lambda)=-\frac1\Lambda\,\Coef{z^n}
 \left\{\varphi(E_{n-1})
 +\sum_{k=1}^{n}a_k\,z^k\bigl(1+E_{n-1}\bigr)^{-(2k-1)}\right\},
 \qquad E_{n-1}=\sum_{j<n}e_jz^j .
 \label{p2:eq:recurrence}
\end{equation}
Only the coefficient \emph{structure} is subject-specific.

\begin{theorem}[Structure of the inverse blocks]
\label{p2:thm:recurrence}
For every \(n\ge1\),
\begin{equation}
 e_n(\Lambda)\in\Lambda^{-1}\Rat[\Lambda^{-1}],
 \qquad
 \deg_{\Lambda^{-1}}e_n\le 2n-1,
 \qquad
 e_n(\Lambda)=-\frac{a_n}\Lambda+\bigO(\Lambda^{-2}).
 \label{p2:eq:en-structure}
\end{equation}
\end{theorem}

\begin{proof}
Specialising \cref{p0:eq:core-bell} to \cref{p2:eq:instance-data} gives
\begin{equation}
 e_n=-\frac1\Lambda\left[
  \sum_{m=2}^{n}\frac{(-1)^m}{m(m-1)}\,\OB_{n,m}(e)
  +\sum_{k=1}^{n}a_k\sum_{m=0}^{n-k}(-1)^m\binom{2k+m-2}{m}
     \OB_{n-k,m}(e)\right],
 \label{p2:eq:bell-recurrence}
\end{equation}
in which every occurring \(e_j\) has \(j<n\).  A monomial of \(\OB_{r,m}\) is a
product
\(e_{j_1}\cdots e_{j_m}\) with \(j_1+\dots+j_m=r\); by induction its
\(\Lambda^{-1}\)-degree lies between \(m\) and
\(\sum_\nu(2j_\nu-1)=2r-m\), and the outer factor \(\Lambda^{-1}\) raises this
by one.  In the first sum \(r=n\) and \(m\ge2\), giving degree at most
\(2n-1\); in the second \(r=n-k\), giving at most \(2(n-k)-m+1\le2n-1\).
Rationality is preserved throughout.  Finally the only contribution of
\(\Lambda^{-1}\)-degree exactly \(1\) is the term \(k=n\), \(m=0\), namely
\(-a_n/\Lambda\), because every Bell monomial contains at least one
\(e_j=\bigO(\Lambda^{-1})\).
\end{proof}

\begin{theorem}[Graded closed formula]
\label{p2:thm:closed}
Let \(A(z)=\sum_{k\ge1}a_kz^k\) and put
\begin{equation}
 A_{n,p}:=\Coef{z^n}A(z)^p=\OB_{n,p}(a),
 \qquad
 K_{n;j,p}:=\Coef{v^{j-1}}\varphi(v)^{\,j-p}(1+v)^{-(2n-p)} .
 \label{p2:eq:AK}
\end{equation}
Expand \(d_{2n-1}(\Lambda)=\sum_{j\ge1}c_{n,j}\Lambda^{-j}\).  Then
\begin{equation}
 \boxed{\;
 c_{n,j}=\frac{(-1)^j}{j}
 \sum_{p=\lceil (j+1)/2\rceil}^{\min(j,n)}
 \binom jp\,A_{n,p}\,K_{n;j,p}\;}
 \qquad(1\le j\le 2n-1),
 \label{p2:eq:closed}
\end{equation}
and \(c_{n,j}=0\) for \(j>2n-1\).  The right-hand side is a finite expression
in Bernoulli numbers, binomial coefficients and ordinary partial Bell
polynomials; it computes any single scalar coefficient without generating the
lower blocks.
\end{theorem}

\begin{proof}
Put \(H(z,v)=\varphi(v)+\sum_{k\ge1}a_kz^k(1+v)^{-(2k-1)}\), so that
\cref{p2:eq:master} reads \(E=-\Lambda^{-1}H(z,E)\).  The closed formula
\cref{p0:eq:core-lagrange} of \cref{p0:thm:core-reversion} gives
\[
 e_n=\sum_{j\ge1}\frac1j\,\Coef{z^nv^{j-1}}
 \left(-\frac{H(z,v)}\Lambda\right)^{j}
 =\sum_{j\ge1}\frac{(-1)^j}{j\Lambda^j}\,
   \Coef{z^nv^{j-1}}H(z,v)^j ,
\]
which already exhibits the grading by powers of \(\Lambda^{-1}\).  It remains
to evaluate \(\Coef{z^nv^{j-1}}H^j\).  Expand \(H^j\) binomially, taking \(p\) of the \(j\) factors from the
Bernoulli sector:
\[
 H^j=\sum_{p=0}^{j}\binom jp\,\varphi(v)^{\,j-p}
 \left(\sum_{k\ge1}a_kz^k(1+v)^{-(2k-1)}\right)^{p} .
\]
Extracting \(z^n\) from the \(p\)-th power, the surviving terms have
\(k_1+\dots+k_p=n\); their coefficients sum to \(A_{n,p}\) and they share the
factor \((1+v)^{-\sum_i(2k_i-1)}=(1+v)^{-(2n-p)}\).  Hence
\(\Coef{z^n}H^j=\sum_p\binom jpA_{n,p}\varphi(v)^{j-p}(1+v)^{-(2n-p)}\), and
extracting \(v^{j-1}\) gives \cref{p2:eq:closed}.  (\Cref{p0:eq:core-lagrange}
caps the outer sum at \(j\le2n-1\); the support argument below recovers that
bound independently.)

For the summation range: \(A_{n,p}=0\) unless \(1\le p\le n\), while
\(\varphi(v)^{j-p}\) has \(v\)-valuation \(2(j-p)\), so a nonzero
\(v^{j-1}\)-coefficient forces \(2(j-p)\le j-1\), i.e.\
\(p\ge\lceil(j+1)/2\rceil\).  Combining, \(j\le2p-1\le2n-1\); this re-proves
the degree bound of \cref{p2:eq:en-structure} and shows the expansion
terminates.
\end{proof}

\begin{remark}
\label{p2:rem:closed-checked}
\Cref{p2:eq:closed} was checked against \cref{p2:eq:recurrence} for all
\(1\le n\le5\) and all \(1\le j\le2n-1\) in exact rational arithmetic: all
\(25\) scalars agree.  For explicit evaluation, \(K_{n;j,p}\) unfolds by
\cref{p0:def:obell,p0:lem:power-log} into
\[
 K_{n;j,p}=\sum_{h=2(j-p)}^{j-1}\OB_{h,\,j-p}(\nu)\,(-1)^{j-1-h}
 \binom{2n-p+j-2-h}{j-1-h},
 \qquad
 \nu_1=0,\quad \nu_m=\frac{(-1)^m}{m(m-1)}\ (m\ge2),
\]
with the conventions \(\OB_{0,0}=1\) and \(\OB_{h,0}=0\) for \(h>0\).
\end{remark}

\begin{proposition}[A second, ungraded closed formula]
\label{p2:prop:lagrange-closed}
Let \(S=A^{\langle-1\rangle}\) be the formal compositional inverse of
\(A(z)=\sum_ka_kz^k\), which exists because \(a_1=-\tfrac16\ne0\), with
\begin{equation}
 \Coef{w^m}S(w)=\frac1m\,\Coef{t^{m-1}}\left(\frac t{A(t)}\right)^m .
 \label{p2:eq:S-lagrange}
\end{equation}
Put
\begin{equation}
 G_\Lambda(E)=(1+E)^2\,
 S\!\left(-\frac{\Lambda E+\varphi(E)}{1+E}\right)
 \in E\,\Rat(\Lambda)[[E]],
 \qquad G_\Lambda'(0)=-\frac\Lambda{a_1}=6\Lambda\ne0 .
 \label{p2:eq:G-Lambda}
\end{equation}
Then
\begin{equation}
 e_n(\Lambda)=\frac1n\,\Coef{E^{n-1}}
 \left(\frac E{G_\Lambda(E)}\right)^{n} .
 \label{p2:eq:lagrange-closed}
\end{equation}
\end{proposition}

\begin{proof}
Since \(\sum_ka_kz^k(1+E)^{-(2k-1)}=(1+E)\,A\bigl(z/(1+E)^2\bigr)\),
\cref{p2:eq:master} reads
\(A\bigl(z/(1+E)^2\bigr)=-\bigl(\Lambda E+\varphi(E)\bigr)/(1+E)\).  Applying
\(S\) and multiplying by \((1+E)^2\) gives \(z=G_\Lambda(E)\); differentiating
at \(E=0\), where \(\varphi(0)=\varphi'(0)=0\), gives
\(G_\Lambda'(0)=S'(0)\cdot(-\Lambda)=-\Lambda/a_1\).
\Cref{p2:eq:lagrange-closed,p2:eq:S-lagrange} are then two applications of
\cref{p0:thm:LB}.  Although \(A\) is a divergent Gevrey series, every finite
jet uses only finitely many \(a_k\), so this is a formal identity requiring no
convergence hypothesis.
\end{proof}

\begin{remark}
\Cref{p2:eq:lagrange-closed} was verified against \cref{p2:eq:recurrence} for
\(n\le4\).  It is less useful than \cref{p2:eq:closed} in practice: it does not
separate the powers of \(\Lambda^{-1}\), and it requires reverting \(A\)
first.  Its interest is conceptual --- it exhibits the inverse blocks as a
\emph{double} Lagrange--B\"urmann extraction, the inner one reverting the pure
Bernoulli series.
\end{remark}

\subsection{Boundary coefficients and the global sign law}

\begin{corollary}[The two edges of the coefficient triangle]
\label{p2:cor:boundary}
For every \(n\ge1\),
\begin{equation}
 c_{n,1}=-a_n=-2c_n,
 \qquad
 c_{n,2n-1}=\frac1{3^n}\binom{1/2}{n}
 =-\frac1{2^{n-1}(2n-1)}\binom{2n-1}{n}a_1^n .
 \label{p2:eq:boundary}
\end{equation}
In particular the two closed forms recorded separately by \textsf{S1} and
\textsf{S2} for the deepest coefficient are the same number.
\end{corollary}

\begin{proof}
For \(j=1\) the range in \cref{p2:eq:closed} forces \(p=1\); then
\(A_{n,1}=a_n\), \(K_{n;1,1}=\Coef{v^0}(1+v)^{-(2n-1)}=1\), and the prefactor
is \(-1\).

For \(j=2n-1\), the range forces \(p=n\); then \(A_{n,n}=a_1^n\) and
\[
 K_{n;2n-1,n}=\Coef{v^{2n-2}}\varphi(v)^{\,n-1}(1+v)^{-n}
 =\Coef{v^{2n-2}}\left(\frac{v^2}2\right)^{n-1}=2^{1-n},
\]
because \(\varphi(v)^{n-1}\) has valuation \(2n-2\).  With
\(\tfrac{(-1)^{2n-1}}{2n-1}\binom{2n-1}{n}\) as prefactor this is the second
expression in \cref{p2:eq:boundary}.

Equality of the two expressions is elementary.  Using
\(\binom{1/2}{n}=(-1)^{n-1}\binom{2n}{n}\bigl(4^n(2n-1)\bigr)^{-1}\),
\[
 \frac1{3^n}\binom{1/2}{n}=(-1)^{n-1}\frac{\binom{2n}{n}}{12^n(2n-1)},
\]
while \(a_1=-\tfrac16\) and \(\binom{2n-1}{n}=\tfrac12\binom{2n}{n}\) give
\[
 -\frac{\binom{2n-1}{n}a_1^n}{2^{n-1}(2n-1)}
 =(-1)^{n+1}\frac{\binom{2n}{n}}{2\cdot2^{n-1}6^n(2n-1)}
 =(-1)^{n+1}\frac{\binom{2n}{n}}{12^n(2n-1)} .
\]
\end{proof}

\begin{remark}[The deepest edge as a square root]
\label{p2:rem:deep-edge}
Let \(t_n=c_{n,2n-1}\).  Only the quadratic part of \(\varphi\) and the single
coefficient \(a_1\) can reach \(\Lambda^{-(2n-1)}\): in
\cref{p2:eq:bell-recurrence} the degree bound \(2r-m+1\) attains \(2n-1\) only
for \(r=n\), \(m=2\) in the first sum, since in the second sum
\(2(n-k)-m+1=2n-1\) forces \(k=1\), \(m=0\), where \(\OB_{n-1,0}=0\) for
\(n\ge2\).  Hence \(t_1=-a_1=\tfrac16\) and
\(t_n=-\tfrac12\sum_{j=1}^{n-1}t_jt_{n-j}\) for \(n\ge2\), i.e.\
\(T(z)=\sum_nt_nz^n\) satisfies \(T+\tfrac12T^2=\tfrac16z\), whence
\(T(z)=-1+\sqrt{1+z/3}\) and \(t_n=3^{-n}\binom{1/2}{n}\).  Equivalently, the
maximal-pole part of \cref{p2:eq:master} is
\(\Lambda E+\tfrac12E^2-\tfrac16z=0\) with root
\(E=-\Lambda+\sqrt{\Lambda^2+z/3}\).  \textsf{S1} asserts this without proof;
the degree bookkeeping above supplies it.
\end{remark}

\begin{theorem}[Coefficientwise sign law]
\label{p2:thm:sign-law}
For every \(n\ge1\),
\begin{equation}
 (-1)^{n+1}d_{2n-1}(\Lambda)\in\Lambda^{-1}\Rat_{>0}[\Lambda^{-1}],
 \qquad
 \deg_{\Lambda^{-1}}d_{2n-1}=2n-1 \ \text{exactly}.
 \label{p2:eq:sign-law}
\end{equation}
That is, \(c_{n,j}\ne0\) with sign \((-1)^{n+1}\) for every
\(1\le j\le2n-1\), and \(c_{n,j}=0\) beyond.
\end{theorem}

\begin{proof}
Write \(a_k=(-1)^ka_k^+\) with
\(a_k^+=2(2k-2)!\,\eta(2k)\pi^{-2k}>0\) by \cref{p2:eq:ck}.  Substituting
\(z=-\zeta\) and \(E=-P\) into \cref{p2:eq:master} and using
\(\varphi(-P)=\sum_{m\ge2}P^m/\bigl(m(m-1)\bigr)\) turns it into
\begin{equation}
 \Lambda P=\sum_{m\ge2}\frac{P^m}{m(m-1)}
 +\sum_{k\ge1}a_k^+\,\zeta^k(1-P)^{-(2k-1)} .
 \label{p2:eq:positive-form}
\end{equation}
Every coefficient on the right is \emph{strictly positive}: \(1/(m(m-1))>0\),
\(a_k^+>0\), and \((1-P)^{-(2k-1)}=\sum_{m\ge0}\binom{2k+m-2}{m}P^m\) has
positive coefficients.  Writing \(P=\sum_np_n(\Lambda)\zeta^n\), the
triangular recurrence obtained from \cref{p2:eq:positive-form} expresses
\(p_n\) as \(\Lambda^{-1}\) times a positive-coefficient polynomial in
\(p_1,\dots,p_{n-1}\); by induction each \(p_n\in\Lambda^{-1}
\Rat_{\ge0}[\Lambda^{-1}]\), with no cancellation possible anywhere.  Since
\(E(z)=-P(-z)\), we get \(d_{2n-1}=e_n=(-1)^{n+1}p_n\), which gives the sign
statement once we show that \(p_n\) has a \emph{strictly} positive coefficient
at every \(\Lambda^{-j}\) with \(1\le j\le2n-1\).

We prove that by induction, exhibiting one contributing configuration for each
\(j\); because all contributions are nonnegative, one suffices.  For \(n=1\),
\(p_1=a_1^+/\Lambda\), and \(2n-1=1\).  Let \(n\ge2\) and assume the claim for
all smaller indices.
\begin{itemize}
\item \(j=1\): take the forcing term \(k=n\), \(m=0\) in
      \cref{p2:eq:positive-form}, contributing \(a_n^+\Lambda^{-1}\).
\item \(2\le j\le2n-2\): take the forcing term \(k=1\) together with the
      single factor \(m=1\) of \((1-P)^{-1}\), so that the contribution is
      \(\Lambda^{-1}a_1^+\binom{1}{1}\) times a coefficient of \(p_{n-1}\).
      By induction \(p_{n-1}\) has strictly positive coefficients at
      \(\Lambda^{-j'}\) for every \(1\le j'\le2n-3\), and
      \(j-1\) runs over exactly that range.
\item \(j=2n-1\): take the nonlinear term \(m=2\), i.e.\ a product
      \(p_{n_1}p_{n_2}\) with \(n_1+n_2=n\), \(n_i\ge1\), at their top degrees
      \(2n_i-1\); the resulting \(\Lambda^{-1}\)-degree is
      \(1+(2n_1-1)+(2n_2-1)=2n-1\).
\end{itemize}
Hence every \(c_{n,j}\), \(1\le j\le2n-1\), is nonzero of sign
\((-1)^{n+1}\); \cref{p2:thm:closed} already showed \(c_{n,j}=0\) for
\(j>2n-1\).
\end{proof}

\begin{remark}
The sign law was confirmed symbolically for \(n\le7\): after multiplication by
\((-1)^{n+1}\) all \(1,3,5,7,9,11,13\) coefficients of the respective blocks
are positive rationals, and the degree is exactly \(2n-1\) in each case.
\textsf{S2} states the theorem but its proof gestures at ``a positive
rooted-tree contribution'' without exhibiting one; the attainability induction
above replaces that step.  The law is also the cheapest runtime invariant for a
symbolic implementation (\cref{p2:alg:inverse}).
\end{remark}

\subsection{The explicit blocks}

Repeated use of \cref{p2:eq:recurrence} gives, in exact rational arithmetic,
\begin{align}
 d_1(\Lambda)&=\frac1{6\Lambda},
 \label{p2:eq:d1}\\[0.3em]
 d_3(\Lambda)&=-\frac{14\Lambda^2+10\Lambda+5}{360\Lambda^3},
 \label{p2:eq:d3}\\[0.3em]
 d_5(\Lambda)&=\frac{2232\Lambda^4+1176\Lambda^3+714\Lambda^2
   +350\Lambda+105}{45360\Lambda^5},
 \label{p2:eq:d5}\\[0.3em]
 d_7(\Lambda)&=-\frac1{5443200\Lambda^7}
  \bigl(822960\Lambda^6+292536\Lambda^5+136956\Lambda^4
  \notag\\[-0.2em]
  &\hspace{6.4em}{}+68040\Lambda^3+32970\Lambda^2+12250\Lambda+2625\bigr),
 \label{p2:eq:d7}\\[0.3em]
 d_9(\Lambda)&=\frac1{359251200\Lambda^9}
  \bigl(309052800\Lambda^8+77920128\Lambda^7+26024592\Lambda^6
  \notag\\[-0.2em]
  &\hspace{4.2em}{}+10019592\Lambda^5+4516028\Lambda^4+2023560\Lambda^3
  \notag\\[-0.2em]
  &\hspace{4.2em}{}+812350\Lambda^2+242550\Lambda+40425\bigr),
 \label{p2:eq:d9}\\[0.3em]
 d_{11}(\Lambda)&=-\frac1{5884534656000\Lambda^{11}}
  \bigl(46195687238400\Lambda^{10}+8854370366400\Lambda^9
  \notag\\[-0.2em]
  &\hspace{2.0em}{}+2171643318000\Lambda^8+605515022112\Lambda^7
  +212869408752\Lambda^6
  \notag\\[-0.2em]
  &\hspace{2.0em}{}+84136852800\Lambda^5+35589153600\Lambda^4
  +14234019800\Lambda^3
  \notag\\[-0.2em]
  &\hspace{2.0em}{}+4868113250\Lambda^2+1213962750\Lambda+165540375\bigr).
 \label{p2:eq:d11}
\end{align}
The rapidly growing numerators are an argument for regenerating blocks from
\cref{p2:eq:recurrence} rather than storing a table; \(d_{13}\), which
\textsf{S3} also displays, is omitted here for that reason.  The alternating
block signs and the coefficientwise positivity after multiplication by
\((-1)^{n+1}\) illustrate \cref{p2:thm:sign-law}; the leading and trailing
coefficients illustrate \cref{p2:cor:boundary}.  In separated form,
\begin{align}
 d_1&=\frac1{6\Lambda},
 &d_3&=-\frac7{180\Lambda}-\frac1{36\Lambda^2}-\frac1{72\Lambda^3},
 \notag\\
 d_5&=\frac{31}{630\Lambda}+\frac7{270\Lambda^2}+\frac{17}{1080\Lambda^3}
      +\frac5{648\Lambda^4}+\frac1{432\Lambda^5},
 &&
 \label{p2:eq:d-separated}
\end{align}
\begin{equation}
 d_7=-\frac{127}{840\Lambda}-\frac{4063}{75600\Lambda^2}
  -\frac{11413}{453600\Lambda^3}-\frac1{80\Lambda^4}
  -\frac{157}{25920\Lambda^5}-\frac{35}{15552\Lambda^6}
  -\frac5{10368\Lambda^7}.
 \label{p2:eq:d7-separated}
\end{equation}
The common-denominator forms are less prone to transcription error; the
separated form makes \(c_{n,1}=-a_n\) visible.  Through three blocks,
\begin{align}
 \mathcal D_\delta^{-1}(y)={}&T-1+\frac1{6T\Lambda}
 -\frac{14\Lambda^2+10\Lambda+5}{360\,T^3\Lambda^3}
 \notag\\
 &+\frac{2232\Lambda^4+1176\Lambda^3+714\Lambda^2+350\Lambda+105}
        {45360\,T^5\Lambda^5}
 +\bigO\!\left(\frac1{T^7\Lambda}\right).
 \label{p2:eq:three-blocks}
\end{align}

\begin{remark}[Lambert-only form, and a correction to the usual first guess]
\label{p2:rem:first-guess}
Because \(T=2R_\delta/w\) and \(\Lambda=w+1\), the expansion can be written
without \(T\).  The first three terms are
\begin{align}
 \mathcal D_\delta^{-1}(y)
 \sim{}&\frac{2R_\delta}{w}-1+\frac{w}{12R_\delta(w+1)}
 -\frac{w^3\bigl(14(w+1)^2+10(w+1)+5\bigr)}{2880\,R_\delta^3(w+1)^3}
 \notag\\
 &+\frac{w^5\bigl(2232(w+1)^4+1176(w+1)^3+714(w+1)^2+350(w+1)+105\bigr)}
        {1451520\,R_\delta^5(w+1)^5}+\cdots,
 \label{p2:eq:inverse-in-w}
\end{align}
with \(w=\Wlam_0(2R_\delta/\e)\).  A frequently used rough inversion is
\(x\approx2R_\delta/\Wlam_0(2R_\delta/\e)-1\); its leading correction is
\emph{not} of order \(T^{-1}\) alone but
\[
 \frac1{6\Lambda T}=\frac{w}{12R_\delta(w+1)}
 =\frac1{12R_\delta}\left(1-\frac1{w+1}\right)
 \sim\frac1{12\log y}.
\]
The factor \(\Lambda^{-1}\) is forced by the slope
\(\Phi'(T)=\Lambda/2\) of \cref{p2:eq:core-identities}.
\end{remark}

\section{The outer iterated-logarithmic layer}
\label{p2:sec:outer}

Keeping \(\Wlam_0\) unexpanded is both more accurate and algebraically
cleaner.  To place the result inside a conventional polynomial--logarithmic
transseries one applies \cref{p0:thm:flattening}.  Set
\begin{equation}
 Z=\frac{2R_\delta}\e,\qquad L_1=\log Z,\qquad L_2=\log L_1 .
 \label{p2:eq:L1L2}
\end{equation}

\begin{proposition}[Outer data]
\label{p2:prop:outer-data}
With \(\stirlingfirst{n}{k}\) the unsigned Stirling numbers of the first kind,
\cref{p0:thm:flattening} gives, for \(Z\to+\infty\),
\begin{equation}
 \Wlam_0(Z)\sim L_1-L_2+\sum_{r\ge1}\frac{P_r(L_2)}{L_1^r},
 \qquad
 P_r(t)=\sum_{m=1}^{r}(-1)^{r-m}\stirlingfirst{r}{r-m+1}\frac{t^m}{m!} .
 \label{p2:eq:W-outer}
\end{equation}
Write \(\Wlam_0(Z)/L_1=1+\sum_{j\ge1}\upsilon_j(L_2)L_1^{-j}\) with
\(\upsilon_1=-L_2\) and \(\upsilon_j=P_{j-1}\) for \(j\ge2\), and define the
reciprocal polynomials
\begin{equation}
 \tau_0=1,\qquad \tau_m=-\sum_{j=1}^{m}\upsilon_j\,\tau_{m-j}\quad(m\ge1).
 \label{p2:eq:tau}
\end{equation}
Then
\begin{equation}
 T\sim\frac{2R_\delta}{L_1}\sum_{m\ge0}\frac{\tau_m(L_2)}{L_1^m},
 \qquad
 \Lambda=1+\Wlam_0(Z) .
 \label{p2:eq:T-outer}
\end{equation}
\end{proposition}

The first data, recomputed from \cref{p2:eq:W-outer,p2:eq:tau}, are
\begin{gather}
 P_1=t,\qquad P_2=\tfrac12t^2-t,\qquad P_3=\tfrac13t^3-\tfrac32t^2+t,
 \label{p2:eq:P-first}\\[0.3em]
 \begin{aligned}
 \tau_0&=1, &\tau_1&=L_2, &\tau_2&=L_2^2-L_2,\\
 \tau_3&=L_2^3-\tfrac52L_2^2+L_2,
 &\tau_4&=L_2^4-\tfrac{13}3L_2^3+\tfrac92L_2^2-L_2,\\
 \tau_5&=L_2^5-\tfrac{77}{12}L_2^4+\tfrac{71}6L_2^3-7L_2^2+L_2. & &
 \end{aligned}
 \label{p2:eq:tau-first}
\end{gather}

\begin{theorem}[Every coefficient of the fully expanded transseries]
\label{p2:thm:outer-complete}
Put \(t=L_1^{-1}\),
\begin{equation}
 \mathcal W(t;L_2)=1+\sum_{j\ge1}\upsilon_j(L_2)t^j,
 \qquad
 \mathcal L(t;L_2)=\mathcal W(t;L_2)+t,
 \label{p2:eq:WL-series}
\end{equation}
so that \(\Wlam_0/L_1=\mathcal W\), \(\Lambda/L_1=\mathcal L\) and
\(T=(2R_\delta/L_1)\mathcal W^{-1}\).  With
\(e_n(\Lambda)=\sum_{p=1}^{2n-1}c_{n,p}\Lambda^{-p}\),
\begin{equation}
 \mathcal D_\delta^{-1}(y)
 \sim\frac{2R_\delta}{L_1}\sum_{r\ge0}\frac{\tau_r(L_2)}{L_1^r}-1
 +\sum_{n\ge1}\left(\frac{L_1}{2R_\delta}\right)^{2n-1}
   \sum_{r\ge1}\frac{D_{n,r}(L_2)}{L_1^{r}},
 \label{p2:eq:full-outer}
\end{equation}
where every \(D_{n,r}\in\Rat[L_2]\) is given by the finite extraction
\begin{equation}
 \boxed{\;
 D_{n,r}(L_2)=\Coef{t^r}\left\{
 \mathcal W(t;L_2)^{2n-1}\sum_{p=1}^{2n-1}c_{n,p}\,t^{p}\,
 \mathcal L(t;L_2)^{-p}\right\} . \;}
 \label{p2:eq:Dnr}
\end{equation}
Together with \cref{p2:eq:closed,p2:eq:W-outer,p2:eq:tau} this is a general
formula for every coefficient at every level.  The leading term of the
\(n\)-th inverse-power sector is \(-a_nL_1^{2n-2}(2R_\delta)^{-(2n-1)}\).
\end{theorem}

\begin{proof}
The core term is \cref{p2:eq:T-outer}, and \cref{p2:eq:tau} is reciprocal
series division.  For the \(n\)-th block,
\[
 e_n(\Lambda)\,T^{-(2n-1)}
 =\left(\frac{L_1}{2R_\delta}\right)^{2n-1}\mathcal W^{2n-1}
   \sum_{p=1}^{2n-1}c_{n,p}L_1^{-p}\mathcal L^{-p},
\]
and setting \(t=L_1^{-1}\) and extracting \(\Coef{t^r}\) gives
\cref{p2:eq:Dnr}; at fixed \(r\) only finitely many \(\upsilon_j\) and
\(c_{n,p}\) occur, so the extraction is finite.  The first nonzero power is
\(r=1\), from \(p=1\) and the constant terms \(\mathcal W(0)=\mathcal L(0)=1\),
with \(c_{n,1}=-a_n\) by \cref{p2:cor:boundary}.
\end{proof}

\begin{proposition}[Ordering of the scales]
\label{p2:prop:ordering}
For every fixed \(M\), \(T/L_1^M\to\infty\) as \(y\to\infty\).  Consequently
\begin{equation}
 \frac{T}{L_1^M}\gg 1\gg\frac1{T\Lambda}\gg\frac1{T^3\Lambda}
 \gg\frac1{T^5\Lambda}\gg\cdots,
 \label{p2:eq:ordering}
\end{equation}
so the constant shift \(-1\) is beyond every fixed order of the outer
\(L_1^{-1}\)-expansion, yet dominates the entire inner grid.  In particular the
first Bernoulli displacement satisfies
\(1/(6T\Lambda)\asymp1/(12R_\delta)\asymp1/(12\log y)\).
\end{proposition}

\begin{proof}
\(T=\e^{\Lambda}\) with \(\Lambda=1+\Wlam_0(Z)\sim L_1\), while
\(L_1^M=\e^{ML_2}\) with \(L_2=\log L_1=\littleo(L_1)\);
hence \(\log(T/L_1^M)=\Lambda-ML_2\to\infty\).  The remaining
comparisons are immediate from \(T\to\infty\) and
\cref{p2:eq:core-identities}, which give
\(T\Lambda=2R_\delta\Lambda/(\Lambda-1)\sim2R_\delta\).
\end{proof}

\begin{remark}
This ordering is destroyed if all terms are written as one informal series in
\(1/\log y\): the shift \(-1\), the outer polynomials \(\tau_m(L_2)/L_1^m\)
and the inner blocks \(d_{2n-1}(\Lambda)T^{1-2n}\) then become
indistinguishable.  \textsf{S1}, \textsf{S2} and \textsf{S3} all make this
point; the quantitative version is \cref{p2:eq:ordering}.
\end{remark}

\section{The periodic interpolation and its oscillatory inverse}
\label{p2:sec:oeis-inverse}

\begin{sourcebox}
\textsf{S1} and \textsf{S2} centre the periodic problem at the parity-mean
constant \(\tfrac14\log(2\pi)\); \textsf{S3} centres it at the even-branch
constant \(\tfrac12\log\pi\) and consequently reports a leading displacement
\(2a(1+\cos\theta)/\Lambda\) rather than \(2a\cos\theta/\Lambda\).  The two
are the same function expanded about different cores (they differ by
\(2a/\Lambda+\bigO(\Lambda^{-2})\), which is exactly \(h'\!\cdot a\) for the
core map \(h\)).  This chapter adopts the mean centring, which makes the
carrier a mean-zero oscillation with the cleanest harmonic content.
\end{sourcebox}

\begin{lemma}[Centred exact form of \(\mathscr D\)]
\label{p2:lem:oeis-centred}
With \(u=x+1\), \(a=\tfrac14\log(\pi/2)\) as in \cref{p2:eq:a-def} and
\begin{equation}
 A_*:=\frac14\log(2\pi)=0.4594692666\ldots,
 \qquad \e^{A_*}=(2\pi)^{1/4}=1.5832334871\ldots,
 \label{p2:eq:Astar}
\end{equation}
one has, exactly and for all \(x>-1\),
\begin{equation}
 \log\mathscr D(x)=\Phi(u)+A_*-a\cos(\pi u)+\Omega(u).
 \label{p2:eq:oeis-centred}
\end{equation}
At an integer \(x=n\), \(\cos\pi u=(-1)^{n+1}\) and
\(A_*-a\cos\pi u=A(n)\) of \cref{p2:eq:parity-transmonomial}: the periodic
factor freezes to the correct parity constant.
\end{lemma}

\begin{proof}
By \cref{p2:eq:oeis,p2:thm:normal-form} with \(\delta=0\),
\(\log\mathscr D(x)=\Phi(u)+A_0+\Omega(u)+P(x)\) where
\(P(x)=\tfrac14(1-\cos\pi x)\log(2/\pi)=-a(1-\cos\pi x)\).  Since
\(\cos\pi x=\cos\pi(u-1)=-\cos\pi u\),
\(P=-a(1+\cos\pi u)=-a-a\cos\pi u\), and
\(A_0-a=\tfrac12\log\pi-\tfrac14\log(\pi/2)=\tfrac14\log(2\pi)=A_*\).
\end{proof}

\begin{lemma}[Exact displacement equation]
\label{p2:lem:oeis-displacement}
Let \(R_*=\log y-A_*\) and let \(T,\Lambda\) be the core data of
\cref{p2:eq:lambert-core} built from \(R_*\); put \(q=T^{-1}\),
\(\theta=\pi T\).  Then \(x=\mathscr D^{-1}(y)\) satisfies \(x=T-1+\Delta\)
where \(\Delta\) is the unique small root of
\begin{equation}
 \Lambda\Delta+\frac{\varphi(q\Delta)}{q}
 +2\,\Omega(T+\Delta)-2a\cos(\theta+\pi\Delta)=0 .
 \label{p2:eq:oeis-exact-eq}
\end{equation}
Replacing \(\Omega\) by its asymptotic series turns \cref{p2:eq:oeis-exact-eq}
into the formal equation \(\mathcal F(q,\Delta)=0\) with
\begin{equation}
 \mathcal F(q,v)=\Lambda v+\sum_{m\ge2}\frac{(-1)^m}{m(m-1)}v^mq^{m-1}
 -2a\cos(\theta+\pi v)
 +2\sum_{k\ge1}c_kq^{2k-1}(1+qv)^{-(2k-1)} .
 \label{p2:eq:oeis-F}
\end{equation}
Since the periodic term is \(\bigO(1)\) while every other nonlinear term
carries a positive power of \(q\), one has \(\Delta=\bigO(\Lambda^{-1})\),
\emph{not} \(\bigO(q\Lambda^{-1})\).
\end{lemma}

\begin{proof}
Evaluate \cref{p2:eq:oeis-centred} at \(u=T+\Delta\), set it equal to
\(\log y=R_*+A_*\), subtract \(\Phi(T)=R_*\), and multiply by \(2\).  By the
computation in \cref{p2:lem:normalized} with \(E=q\Delta\),
\(2\bigl[\Phi(T+\Delta)-\Phi(T)\bigr]
=T\bigl[\Lambda q\Delta+\varphi(q\Delta)\bigr]
=\Lambda\Delta+\varphi(q\Delta)/q\), since \(T=q^{-1}\).  This gives
\cref{p2:eq:oeis-exact-eq}.  The last
sum in \cref{p2:eq:oeis-F} is \(2\Omega(T(1+q v))\) expanded.  Finally
\(\varphi(q\Delta)/q=q\Delta^2/2+\bigO(q^2)\), so at \(q=0\) the equation
reduces to \(\Lambda\Delta_0=2a\cos(\theta+\pi\Delta_0)\), forcing
\(\Delta_0\asymp\Lambda^{-1}\).
\end{proof}

\subsection{The periodic carrier}

\begin{theorem}[Existence, uniqueness, and the closed carrier expansion]
\label{p2:thm:carrier}
Fix \(\theta\in\Real\) and put \(\mu=2a/\Lambda\).
\begin{enumerate}[label=\textup{(\alph*)}]
\item If \(\Lambda>2\pi a=0.70934\ldots\), the equation
      \begin{equation}
       \Delta_0=\mu\cos(\theta+\pi\Delta_0)
       \label{p2:eq:carrier-eq}
      \end{equation}
      has exactly one real solution, and \(|\Delta_0|\le\mu\).
\item If moreover \(\mu\le\tfrac15\), equivalently
      \(\Lambda\ge10a=1.12896\ldots\), that solution is the sum of the
      convergent Lagrange--B\"urmann series
      \begin{equation}
       \boxed{\;
       \Delta_0=\sum_{m\ge1}\frac{\mu^m}{m!}
       \left.\frac{\dd^{m-1}}{\dd v^{m-1}}\cos^m(\theta+\pi v)\right|_{v=0}
       =\sum_{m\ge1}\frac{a^m}{m!\,\Lambda^m}
        \sum_{j=0}^{m}\binom mj\bigl(\ii\pi(m-2j)\bigr)^{m-1}
        \e^{\ii(m-2j)\theta} . \;}
       \label{p2:eq:carrier-series}
      \end{equation}
      The second expression is real.  At inverse-logarithmic order \(m\) only
      the Fourier modes \(m,m-2,\dots,-m\) occur.
\item The first three terms are
      \begin{equation}
       \Delta_0=\mu\cos\theta-\frac\pi2\mu^2\sin2\theta
       +\frac{\pi^2}2\mu^3\cos\theta\,(2-3\cos^2\theta)+\bigO(\mu^4).
       \label{p2:eq:carrier-first}
      \end{equation}
\end{enumerate}
\end{theorem}

\begin{proof}
(a) The map \(g(v)=\mu\cos(\theta+\pi v)\) satisfies
\(|g'|\le\pi\mu=2\pi a/\Lambda<1\) and \(|g|\le\mu\), so it is a contraction of
\(\Real\) into \([-\mu,\mu]\); Banach's theorem gives a unique fixed point
there, and any real fixed point lies in \([-\mu,\mu]\).

(b) \Cref{p2:eq:carrier-eq} is \(v=\mu\phi(v)\) with
\(\phi(v)=\cos(\theta+\pi v)\) entire.  By the classical convergence criterion
for Lagrange's series, if there is \(r>0\) with
\(|\mu|\max_{|v|=r}|\phi(v)|<r\), then the series
\(\sum_m\tfrac{\mu^m}{m!}\bigl(\tfrac{\dd}{\dd v}\bigr)^{m-1}\phi(v)^m|_{v=0}\)
converges and its sum is the unique root of \(v=\mu\phi(v)\) in \(|v|<r\).
Here \(|\cos(\theta+\pi v)|\le\cosh(\pi|{\Im}\,v|)\le\cosh(\pi r)\) on
\(|v|=r\); with \(r=0.38\) and \(\mu\le\tfrac15\),
\(\mu\cosh(0.38\pi)\le0.3603<0.38\).  (The optimal choice
\(\max_{r>0}r/\cosh(\pi r)=0.21095\ldots\) at \(r=0.382\) shows
\(\mu<0.2109\) already suffices.)  For the Fourier form, insert
\(\cos^m v=2^{-m}\sum_{j}\binom mj\e^{\ii(m-2j)v}\) into the derivative:
\[
 \left.\frac{\dd^{m-1}}{\dd v^{m-1}}\cos^m(\theta+\pi v)\right|_{v=0}
 =2^{-m}\sum_{j=0}^m\binom mj\bigl(\ii\pi(m-2j)\bigr)^{m-1}
 \e^{\ii(m-2j)\theta},
\]
and \(\mu^m2^{-m}=a^m\Lambda^{-m}\).  Pairing \(j\) with \(m-j\) replaces
\(m-2j\) by its negative, conjugating each summand, so the sum is real.  Only
\(|m-2j|\le m\) occurs, with the same parity as \(m\).

(c) For \(m=1\) the derivative is \(\cos\theta\).  For \(m=2\),
\(\tfrac{\dd}{\dd v}\cos^2=-2\pi\cos\sin=-\pi\sin2(\cdot)\), giving
\(-\tfrac\pi2\mu^2\sin2\theta\).  For \(m=3\),
\(\tfrac{\dd^2}{\dd v^2}\cos^3(\theta+\pi v)|_{0}
=3\pi^2\cos\theta(2-3\cos^2\theta)\), and division by \(3!\) gives the third
term.
\end{proof}

\begin{repairbox}
\textsf{S1} states the carrier expansion as ``convergent for sufficiently
large \(\Lambda\)'' with no threshold and no criterion, and \textsf{S2} proves
only the contraction estimate, from which convergence of the Lagrange series
does not follow.  \Cref{p2:thm:carrier} separates the two claims and supplies
an explicit sufficient condition, \(\mu\le\tfrac15\).  Numerically at
\(\Lambda=3\), \(\theta=2\), the partial sums \(m\le1,\dots,5\) of
\cref{p2:eq:carrier-series} approximate the exact fixed point
\(-0.0256993556\) with errors
\(5.6\cdot10^{-3},1.1\cdot10^{-3},1.8\cdot10^{-4},1.9\cdot10^{-5},
1.9\cdot10^{-6}\), consistent with a convergence ratio \(\approx\mu/0.211\).
\end{repairbox}

\subsection{Algebraic corrections around the carrier}

\begin{theorem}[Triangular periodic-grid recurrence]
\label{p2:thm:periodic-recurrence}
Let
\begin{equation}
 M:=\partial_v\mathcal F(0,\Delta_0)=\Lambda+2\pi a\sin(\theta+\pi\Delta_0),
 \label{p2:eq:M}
\end{equation}
which satisfies \(M\ge\Lambda-2\pi a>0\) for \(\Lambda>2\pi a\).  Seek
\(\Delta=\Delta_0+\sum_{n\ge1}\Delta_nq^n\) with each \(\Delta_n\) a formal
series in \(\Lambda^{-1}\) with trigonometric-polynomial coefficients, and put
\(\Delta_{<n}=\Delta_0+\sum_{j<n}\Delta_jq^j\).  Then the solution of
\(\mathcal F=0\) is unique and generated by
\begin{equation}
 \boxed{\;
 \Delta_n=-\frac1M\,\Coef{q^n}\,
 \mathcal F\bigl(q,\Delta_{<n}(q)\bigr),\qquad n\ge1. \;}
 \label{p2:eq:periodic-recurrence}
\end{equation}
The first two are
\begin{equation}
 \Delta_1=\frac{\tfrac16-\tfrac12\Delta_0^2}{M},
 \qquad
 \Delta_2=-\frac1M\left[
 a\pi^2\cos(\theta+\pi\Delta_0)\,\Delta_1^2
 +\Delta_0\Delta_1+\frac{\Delta_0-\Delta_0^3}{6}\right].
 \label{p2:eq:Delta12}
\end{equation}
\end{theorem}

\begin{proof}
In \(\Coef{q^n}\mathcal F(q,\Delta)\) the unknown \(\Delta_n\) appears only
linearly, through the \(v\)-derivative of the \(q\)-free part
\(\Lambda v-2a\cos(\theta+\pi v)\) evaluated at \(v=\Delta_0\), which is
\(M\); every other contribution involves only \(\Delta_j\) with \(j<n\).
This is \cref{p0:lem:lagrange-fp} applied in the variable \(q\) over the
coefficient ring generated by \(\Lambda^{-1}\) and \(\e^{\pm\ii\theta}\), and
\(M\) is invertible there because \(M=\Lambda(1+\bigO(\Lambda^{-1}))\).

For \(n=1\): the \(m=2\) term of \cref{p2:eq:oeis-F} gives
\(\tfrac12\Delta_0^2\), the cosine gives \(2a\pi\sin(\theta+\pi\Delta_0)
\Delta_1\), and the \(k=1\) Bernoulli term gives \(2c_1=-\tfrac16\); collecting
gives \(M\Delta_1+\tfrac12\Delta_0^2-\tfrac16=0\).  For \(n=2\): the \(m=2\)
term gives \(\Delta_0\Delta_1\), the \(m=3\) term gives
\(-\tfrac16\Delta_0^3\), the second-order Taylor term of the cosine gives
\(a\pi^2\cos(\theta+\pi\Delta_0)\Delta_1^2\), and expanding
\(2c_1q(1+q\Delta)^{-1}\) gives \(-2c_1\Delta_0=\tfrac16\Delta_0\); collecting
gives \cref{p2:eq:Delta12}.
\end{proof}

\begin{corollary}[Asymptotic meaning]
\label{p2:cor:periodic-meaning}
Fix \(N,K\ge0\).  Truncate \(\Delta\) after \(q^N\) and expand each
\(\Delta_n\) in \(\Lambda^{-1}\) through \(\Lambda^{-K}\), obtaining
\(\widetilde\Delta\).  Then, as \(y\to+\infty\),
\begin{equation}
 \mathscr D^{-1}(y)
 =T-1+\widetilde\Delta
 +\bigO\!\left(\frac1{T^{N+1}\Lambda}\right)
 +\bigO\!\left(\frac1{\Lambda^{K+2}}\right),
 \label{p2:eq:oeis-inverse}
\end{equation}
uniformly in the phase \(\theta=\pi T\).
\end{corollary}

\begin{proof}
By construction \cref{p2:eq:periodic-recurrence} cancels the residual of
\(\mathcal F\) through \(q^N\), and \cref{p2:thm:carrier} cancels it through
\(\Lambda^{-K}\) at each fixed \(q\)-order; the replacement of \(\Omega\) by
its truncation is controlled by \cref{p2:thm:remainder}.  The exact
\(v\)-derivative of the left side of \cref{p2:eq:oeis-exact-eq} equals
\(M+\bigO(q)\) and is bounded below by \(\Lambda-2\pi a-\bigO(q)\), a positive
multiple of \(\Lambda\) for large \(T\); all trigonometric coefficients and
their derivatives are bounded, so \cref{p0:thm:backward-error} converts the
residual into a root error after division by a quantity \(\asymp\Lambda\),
uniformly in \(\theta\).
\end{proof}

\begin{proposition}[The parity branches resum the periodic sector]
\label{p2:prop:resummation}
At an integer \(u=n+1\) one has \(A_*-a\cos\pi u=A_\delta\) with
\(\delta\equiv n\pmod2\).  The branchwise normalisation therefore absorbs the
whole \(\bigO(1)\) periodic constant into the target \(R_\delta\) before the
core is formed, so the \(\bigO(\Lambda^{-1})\) carrier disappears and the first
correction is the much smaller \(1/(6T\Lambda)\).  In this precise sense the two
parity transseries are subsequence resummations of the periodic one: they build
the parity phase into the core instead of expanding it perturbatively.
\end{proposition}

\begin{proof}
Immediate from \cref{p2:lem:oeis-centred} and
\cref{p2:thm:branch-inverse,p2:thm:carrier}.
\end{proof}

\begin{proposition}[Phase-coordinate reversion, formal]
\label{p2:prop:phase-reversion}
Let \(h(R)=2R/\Wlam_0(2R/\e)\) be the inverse of \(\Phi\)
(\cref{p2:lem:core}), and set
\(\mathcal A(R)=-a\cos\bigl(\pi h(R)\bigr)+\Omega\bigl(h(R)\bigr)\).  Then the
equation \(\Phi(u)+\mathcal A(\Phi(u))=R_*\) --- which is
\cref{p2:eq:oeis-centred} rewritten in the phase coordinate
\(\sigma=\Phi(u)\) --- has the formal reversion
\begin{equation}
 \sigma=R_*+\sum_{m\ge1}\frac{(-1)^m}{m!}
 \left(\frac{\dd}{\dd R_*}\right)^{m-1}\mathcal A(R_*)^m,
 \qquad
 \mathscr D^{-1}(y)=h(\sigma)-1 .
 \label{p2:eq:phase-reversion}
\end{equation}
\end{proposition}

\begin{proof}
Write \(\sigma=R_*+\varsigma\); then \(\varsigma=-\mathcal A(R_*+\varsigma)\),
and Lagrange's formula for \(\varsigma=\tau\phi(\varsigma)\) with
\(\phi(\varsigma)=-\mathcal A(R_*+\varsigma)\), \(\tau=1\), together with the
translation invariance of \(\dd/\dd\varsigma\), gives
\cref{p2:eq:phase-reversion}.
\end{proof}

\begin{repairbox}
\textsf{S3} states \cref{p2:eq:phase-reversion} as a boxed asymptotic
expansion, writing ``setting \(\tau=1\) as an asymptotic expansion gives''.
That is a formal-to-analytic slide: \(\mathcal A\) is not small in any norm
that makes the reversion series converge, and \(m\)-fold differentiation with
respect to \(R_*\) produces factors \(h'(R_*)=2/\Lambda\) whose accumulated
effect is exactly what needs to be estimated.  The statement is demoted here to
a \emph{formal} proposition.  An asymptotic claim would require, at each order
\(m\), a bound on \(\bigl(\dd/\dd R_*\bigr)^{m-1}\mathcal A^m\) uniform in the
oscillatory phase --- available for the truncated grid of
\cref{p2:cor:periodic-meaning}, which is therefore the statement to use.
Comparing the two routes at leading order: the \(m=1\) term of
\cref{p2:eq:phase-reversion} contributes
\(h'(R_*)\cdot(-\mathcal A)=2a\cos(\pi T)/\Lambda+\bigO(T^{-1}\Lambda^{-1})\),
in agreement with \cref{p2:eq:carrier-first}.
\end{repairbox}

\begin{warning}
\label{p2:rem:not-intrinsic}
The carrier of \cref{p2:thm:carrier} records the choice
\cref{p2:eq:oeis}, not a property of \(n!!\).  By
\cref{p2:thm:no-canonical} another interpolation, agreeing with \(n!!\) at
every nonnegative integer, changes the amplitude of this sector from \(2a\) to
\(2\sqrt{a^2+\lambda^2}\) and shifts its phase, while leaving every integer
value and both parity transseries untouched.
\end{warning}

\section{The discrete inverses}
\label{p2:sec:discrete}

Objects (i) and (iii) of \cref{p2:sec:orientation} are intrinsic to the
sequence.  They are the residue-class staircase of
\cref{p0:def:three-inverses,p0:thm:staircase} at modulus \(r=2\), and nothing
new has to be proved: \cref{p0:rem:parity-instances} already records this
chapter as that instance.  Only the dictionary and the arithmetic caveat are
subject-specific.

\begin{proposition}[Specialisation of the staircase]
\label{p2:prop:discrete}
Take \(r=2\), residue classes \(\rho\in\{0,1\}\), \(A_n=n!!\), and
\(\nu_\rho(y)=\mathcal D_\rho^{-1}(y)\), which is legitimate because
\(\mathcal D_\rho\) interpolates the class-\(\rho\) subsequence and is strictly
increasing (\cref{p2:prop:branch-interp}).  Then
\cref{p0:eq:residue-staircase} reads
\begin{equation}
 N(y):=\min\{n\in\Nat: n!!\ge y\}
 =\min\left\{2\left\lceil\frac{\nu_0(y)}2\right\rceil,\;
   2\left\lceil\frac{\nu_1(y)-1}2\right\rceil+1\right\}
 \qquad(y>1),
 \label{p2:eq:N-exact}
\end{equation}
and its floor counterpart, obtained from \cref{p0:thm:staircase} by replacing
\(\lceil\cdot\rceil\) with \(\lfloor\cdot\rfloor\) and \(\min\) with \(\max\),
is
\begin{equation}
 J(y):=\max\{n\in\Nat: n!!\le y\}
 =\max_{\rho\in\{0,1\}}
 \left\{\rho+2\left\lfloor\frac{\nu_\rho(y)-\rho}2\right\rfloor\right\}.
 \label{p2:eq:J-exact}
\end{equation}
Replacing \(\nu_\rho\) by the truncation of \cref{p2:eq:branch-inverse} with a
certified error bound \(\eta_\rho\) (for instance from
\cref{p2:eq:aposteriori}) determines \(N(y)\) and \(J(y)\) exactly under the
separation condition of \cref{p0:thm:staircase}, and
\cref{p0:prop:certified-rounding} bounds the number of exact evaluations of
\(n!!\) needed by \(\sum_\rho(\lfloor\eta_\rho\rfloor+2)\).  Those evaluations
should be made as \(\log\mathcal D_\rho(n)=\log\kappa_\rho+\tfrac n2\log2
+\log\Gamma(\tfrac n2+1)\), or by \cref{p2:eq:dfac-recurrence} in exact
integers.
\end{proposition}

\begin{proof}
Everything is \cref{p0:thm:staircase} and \cref{p0:prop:certified-rounding}
with \(r=2\); the hypotheses hold by \cref{p2:prop:branch-interp} and, for the
asymptotic input, by \cref{p2:thm:branch-inverse}.
\end{proof}

\begin{warning}[Never form \(y\)]
\label{p2:rem:never-form-y}
\Cref{p2:eq:N-exact}, evaluated with four inverse blocks, was tested at
\(y=n!!\) for every \(6\le n\le199\) and returned the correct \(N(y)=n\) in all
\(194\) cases --- but only when the comparison \(A_n\ge y\) is made on exact
integers or on \(\log y\).  Converting \(n!!\) to a fixed-precision float first
destroys the comparison long before the asymptotics fail: at \(80\) working
digits this already produces spurious off-by-one answers from
\(n\approx120\) onwards, because \(120!!\) has more than \(100\) digits.  The
failure is arithmetic, not asymptotic, and it is easy to misdiagnose as a
defect of the expansion.
\end{warning}

\section{Beyond all algebraic orders}
\label{p2:sec:beyond}

\begin{proposition}[Optimal truncation, forward and inverse]
\label{p2:prop:least-term}
Let \(u>0\).  The terms of \cref{p2:eq:Omega-series} have moduli
\(|c_k|u^{1-2k}\) with consecutive ratio
\(\sim2k(2k-1)/(\pi u)^2\); the least term occurs at \(2k\approx\pi u\) and
\begin{equation}
 \min_k\left|\frac{c_k}{u^{2k-1}}\right|
 =\bigO\!\left(u^{-1/2}\e^{-\pi u}\right).
 \label{p2:eq:least-term}
\end{equation}
Consequently, by \cref{p0:thm:optimal-truncation} and the slope
\(\Phi'(T)=\Lambda/2\), the optimally truncated branchwise inverse has error
\begin{equation}
 \bigO\!\left(\frac{\e^{-\pi T}}{\sqrt T\,\Lambda}\right).
 \label{p2:eq:inverse-least-term}
\end{equation}
\end{proposition}

\begin{proof}
\Cref{p2:eq:ck-late} gives \(|c_k|=(2k-2)!\pi^{-2k}(1+\bigO(4^{-k}))\), whence
the ratio.  At \(2k=\pi u\), Stirling's formula for \((2k-2)!\) gives
\cref{p2:eq:least-term}.  For the inverse, \cref{p2:thm:remainder} bounds the
forward logarithmic residual at optimal truncation by
\cref{p2:eq:least-term} with \(u=T(1+\bigO(T^{-2}))\); dividing by the slope
\(\tfrac12\Lambda+\bigO(T^{-2})\) via
\cref{p0:thm:backward-error,p0:thm:remainder-transport} gives
\cref{p2:eq:inverse-least-term}.
\end{proof}

\begin{proposition}[Transport of a forward exponential sector]
\label{p2:prop:transport}
Suppose an exponentially improved representation of the centred forward phase
retains the algebraic terms \(c_1,\dots,c_K\) and has an additional residual
\(\mathcal R\), and let \(u_{\mathrm{alg}}\) solve the correspondingly
truncated algebraic inverse problem.  Then
\begin{equation}
 u-u_{\mathrm{alg}}
 =-\frac{2\,\mathcal R(u_{\mathrm{alg}})}
   {\log u_{\mathrm{alg}}
    -2\sum_{k=1}^{K}(2k-1)c_k\,u_{\mathrm{alg}}^{-2k}}
 +\bigO\!\left(\frac{\sup|F''|\,\mathcal R^2}{\Lambda^3}
   +\frac{|\mathcal R\,\mathcal R'|}{\Lambda^2}\right),
 \label{p2:eq:transport}
\end{equation}
where \(F\) is the doubled algebraic phase and the supremum is over the
interval between \(u_{\mathrm{alg}}\) and the exact root; the estimate assumes
the displayed denominator is \(\asymp\Lambda\) and
\(|\mathcal R'|=\littleo(\Lambda)\).  At leading order this is
\(-2\mathcal R(T)/\Lambda\).
\end{proposition}

\begin{proof}
Let \(F\) be the doubled algebraic phase minus \(2R\), so
\(F(u_{\mathrm{alg}})=0\) and the exact phase equals
\(F+2\mathcal R\).  Taylor expansion of
\(0=F(u_{\mathrm{alg}}+\Delta)+2\mathcal R(u_{\mathrm{alg}}+\Delta)\) about
\(u_{\mathrm{alg}}\), with \(F'(u_{\mathrm{alg}})\) the stated denominator,
gives the first Newton correction; replacing
\(F'+2\mathcal R'\) by \(F'\) and iterating once produces the two displayed
error terms.
\end{proof}

\begin{warning}[What the exponential sector is, and is not]
\label{p2:rem:terminant}
The Borel poles of \cref{p2:eq:mittag} lie at \(s=\pm\ii\pi m\), so the
exponentials associated with them are \(\e^{\mp\ii\pi mu}\) together with
direction-dependent terminant multipliers, not real exponentials.  On the
positive axis the optimally truncated remainder has modulus
\(\asymp u^{-1/2}\e^{-\pi u}\), but this is the size of a
\emph{terminant-weighted combination}: appending a bare term
\(\mathrm{const}\cdot\e^{-\pi u}\) to the divergent Bernoulli series is not a
valid transseries completion, and conflating the two gives incorrect Stokes
statements.  The unambiguous positive-real completion is the exact integral
\cref{p2:eq:borel} itself, which by \cref{p2:thm:remainder} even provides
two-sided bounds at every order; \cref{p2:eq:transport} then carries any
rigorously normalised forward completion through the inverse map.  All three
sources make this point; \textsf{S3} states it most sharply.
\end{warning}

\section{Complex inverse sheets}
\label{p2:sec:complex}

\begin{proposition}[Local sheets]
\label{p2:prop:sheets}
Fix a branch of \(\log y\) and integers \(\ell\) (logarithmic winding) and
\(k\) (Lambert branch).  Put
\begin{equation}
 R_{\delta,\ell}=\log y+2\pi\ii\ell-A_\delta,
 \quad
 w_{k,\ell}=\Wlam_k\!\left(\frac{2R_{\delta,\ell}}\e\right),
 \quad
 T_{k,\ell}=\frac{2R_{\delta,\ell}}{w_{k,\ell}},
 \quad
 \Lambda_{k,\ell}=w_{k,\ell}+1 .
 \label{p2:eq:complex-core}
\end{equation}
In any sector in which \(T_{k,\ell}\to\infty\), the Stirling expansion for
\(\log\Gamma\) is uniform, and the corresponding solution branch of
\(\mathcal D_\delta(x)=y\) is simple, the \emph{same} universal polynomials
give
\begin{equation}
 x_{\delta;k,\ell}(y)\sim T_{k,\ell}-1
 +\sum_{n\ge1}\frac{d_{2n-1}(\Lambda_{k,\ell})}{T_{k,\ell}^{2n-1}} .
 \label{p2:eq:complex-sheet}
\end{equation}
\end{proposition}

\begin{proof}
Every step of \cref{p2:lem:normalized,p2:thm:recurrence} is an identity of
formal series over \(\Cplx(\Lambda)\) and uses only \(\Phi(T)=R\),
\(\Lambda=\log T\); \cref{p2:eq:complex-core} secures those.  The analytic
statement is the sectorial version of \cref{p2:thm:branch-inverse}, with the
mean-value step replaced by the standard implicit-function estimate in the
sector.
\end{proof}

\begin{remark}[Branch caveats]
\label{p2:rem:branch-caveat}
The pair \((k,\ell)\) is an asymptotic label, not a global classification:
distinct labels can be joined by analytic continuation around the poles of
\(\Gamma\) or the logarithmic cut.  Branch collisions occur exactly where
\(\mathcal D_\delta'=0\), i.e.\ where \(\log2+\psi(x/2+1)=0\); by
\cref{p2:prop:branch-interp} this never happens for real \(x\ge0\), but it does
for complex and for negative real \(x\), and near such points
\cref{p2:eq:complex-sheet} must be replaced by a Puiseux-type local analysis.
\end{remark}

\section{Multifactorials}
\label{p2:sec:multifactorial}

\begin{sourcebox}
Only \textsf{S3} treats multifactorials.  Its normal form is reproduced here
with a proof, and \cref{p2:thm:multi-scaling} is new: it subsumes the
``half-shifted inverse Gamma'' scaling that \textsf{S1} and \textsf{S2} record
separately, by identifying the Gamma case as the member \(p=1\).
\end{sourcebox}

\begin{proposition}[Residue-class branches]
\label{p2:prop:multi-branch}
Let \(p\ge1\) and \(r\in\{1,\dots,p\}\).  For integers \(x\equiv r\pmod p\),
\(x\ge r\), set \(x!_{(p)}=x(x-p)(x-2p)\cdots r\).  Then, with
\begin{equation}
 \mathcal D_{p,r}(x)=C_{p,r}\;p^{x/p}\,\Gamma\!\left(\frac xp+1\right),
 \qquad
 C_{p,r}=\frac{p^{1-r/p}}{\Gamma(r/p)},
 \label{p2:eq:multi-branch}
\end{equation}
one has \(\mathcal D_{p,r}(x)=x!_{(p)}\) for every such \(x\).  For \(p=2\),
\(C_{2,2}=1=\kappa_0\) and \(C_{2,1}=\sqrt{2/\pi}=\kappa_1\); for \(p=1\),
\(C_{1,1}=1\) and \(\mathcal D_{1,1}(x)=\Gamma(x+1)\).
\end{proposition}

\begin{proof}
Write \(x=r+mp\).  Then
\(x!_{(p)}=\prod_{i=0}^{m}(r+ip)=p^{m+1}\prod_{i=0}^m(\tfrac rp+i)
=p^{m+1}\Gamma(\tfrac rp+m+1)/\Gamma(\tfrac rp)\), while
\(\mathcal D_{p,r}(x)=C_{p,r}p^{r/p+m}\Gamma(\tfrac rp+m+1)\).  Equating gives
\(C_{p,r}=p^{1-r/p}/\Gamma(r/p)\).
\end{proof}

\begin{theorem}[Centred multifactorial normal form]
\label{p2:thm:multi-normal}
Put \(u=x+\tfrac p2\) and
\(A_{p,r}=\log C_{p,r}+\tfrac12\log(2\pi/p)\).  Then, as \(x\to+\infty\),
\begin{equation}
 \log\mathcal D_{p,r}(x)=\frac up(\log u-1)+A_{p,r}+\Omega_p(u),
 \qquad
 \Omega_p(u)\sim\sum_{k\ge1}\frac{c^{(p)}_k}{u^{2k-1}},
 \quad
 c^{(p)}_k=\frac{p^{2k-1}B_{2k}(\tfrac12)}{2k(2k-1)},
 \label{p2:eq:multi-normal}
\end{equation}
and \(\Omega_p\) has the exact Borel--Laplace representation
\begin{equation}
 \Omega_p(u)=\int_0^\infty\e^{-us}\,
 \frac{\dfrac{ps}{2\sinh(ps/2)}-1}{ps^2}\,\dd s ,
 \label{p2:eq:multi-borel}
\end{equation}
whose singularities are the simple poles \(s=\pm2\pi\ii m/p\).  Hence the
optimally truncated forward remainder has scale \(\e^{-2\pi u/p}\).  For
\(p=2\) this is \cref{p2:thm:normal-form,p2:thm:borel}; for \(p=1\) it is the
half-shifted Stirling expansion of \(\log\Gamma(x+1)=\log\Gamma(u+\tfrac12)\).
\end{theorem}

\begin{proof}
Put \(t=u/p\), so \(x=pt-\tfrac p2\) and
\(\log\mathcal D_{p,r}=\log C_{p,r}+(t-\tfrac12)\log p
+\log\Gamma(t+\tfrac12)\).  The half-shifted Stirling expansion
\(\log\Gamma(t+\tfrac12)\sim t\log t-t+\tfrac12\log(2\pi)
+\sum_k B_{2k}(\tfrac12)\bigl(2k(2k-1)\bigr)^{-1}t^{1-2k}\) --- the case
\(a=\tfrac12\) of \(\log\Gamma(t+a)\), in which every odd-index term vanishes
because \(B_{2j+1}(\tfrac12)=0\) --- gives
\[
 \log\mathcal D_{p,r}
 \sim\log C_{p,r}-\tfrac12\log p+\tfrac12\log(2\pi)
 +\frac up(\log u-1)
 +\sum_k\frac{p^{2k-1}B_{2k}(\tfrac12)}{2k(2k-1)u^{2k-1}},
\]
since \((t-\tfrac12)\log p+t\log t=(u/p)\log u-\tfrac12\log p\) and
\(t^{1-2k}=p^{2k-1}u^{1-2k}\).  The constant is \(A_{p,r}\).  For
\cref{p2:eq:multi-borel}, \(\Omega_p(u)=\widetilde\Omega(u/p)\) with
\(\widetilde\Omega\) as in the proof of \cref{p2:thm:borel}, where
\(\widetilde\Omega(t)=\int_0^\infty\e^{-tr}
\bigl(r/(2\sinh(r/2))-1\bigr)r^{-2}\dd r\); substituting \(r=ps\) gives
\cref{p2:eq:multi-borel}, and \(\sinh(ps/2)=0\) at \(s=2\pi\ii m/p\).
\end{proof}

\begin{theorem}[Universal multifactorial inverse and its scaling]
\label{p2:thm:multi-scaling}
Let \(R=\log y-A_{p,r}\),
\(v=\Wlam_0(pR/\e)\), \(X=pR/v\), \(\Lambda=v+1=\log X\), so that
\(\tfrac Xp(\log X-1)=R\).  With \(u=X(1+E)\), \(q=X^{-1}\), \(z=q^2\), the
normalised equation is
\begin{equation}
 \Lambda E+\varphi(E)+\sum_{k\ge1}a^{(p)}_kz^k(1+E)^{-(2k-1)}=0,
 \qquad
 a^{(p)}_k:=p\,c^{(p)}_k=p^{2k}\,\frac{B_{2k}(\tfrac12)}{2k(2k-1)} .
 \label{p2:eq:multi-master}
\end{equation}
Every statement of
\cref{p2:thm:recurrence,p2:thm:closed,p2:prop:lagrange-closed,p2:cor:boundary,p2:thm:sign-law}
holds verbatim with \(a_k\) replaced by \(a^{(p)}_k\), and for every fixed
\(N\)
\begin{equation}
 \mathcal D_{p,r}^{-1}(y)=X-\frac p2
 +\sum_{n=1}^{N}\frac{d^{(p)}_{2n-1}(\Lambda)}{X^{2n-1}}
 +\bigO\!\left(\frac1{X^{2N+1}\Lambda}\right).
 \label{p2:eq:multi-inverse}
\end{equation}
Moreover the coefficient polynomials of all multifactorials are one family:
\begin{equation}
 \boxed{\;
 d^{(p)}_{2n-1}(\Lambda)=p^{2n}\,d^{(1)}_{2n-1}(\Lambda)
 \qquad\text{for every }p\ge1,\ n\ge1, \;}
 \label{p2:eq:multi-scaling}
\end{equation}
where \(d^{(1)}\) are the blocks of the half-shifted inverse-Gamma problem.  In
particular \(d_{2n-1}=d^{(2)}_{2n-1}=4^n\,d^{(1)}_{2n-1}\).
\end{theorem}

\begin{proof}
The derivation of \cref{p2:eq:multi-master} repeats
\cref{p2:lem:normalized} with \(\tfrac2T\) replaced by \(\tfrac pX\): the
phase difference is \(\tfrac Xp[\Lambda E+\varphi(E)]\), and
\(\tfrac pX\Omega_p(u)=\sum_kp\,c^{(p)}_kq^{2k}(1+E)^{-(2k-1)}\).  Since
\cref{p2:eq:multi-master} differs from \cref{p2:eq:master} only in the values
of the forcing coefficients, all structural statements transfer.  Now
\(a^{(p)}_k=p^{2k}a^{(1)}_k\) by \cref{p2:eq:multi-normal}, so substituting
\(q\mapsto pq\) (equivalently \(z\mapsto p^2z\)) carries the \(p=1\) equation
into the general one with \(\Lambda\) untouched.  Uniqueness of the formal
solution gives \(E^{(p)}(q)=E^{(1)}(pq)\), i.e.\
\(e^{(p)}_n=p^{2n}e^{(1)}_n\), which is \cref{p2:eq:multi-scaling}.  The
remainder in \cref{p2:eq:multi-inverse} follows from
\cref{p0:thm:backward-error,p0:thm:remainder-transport} exactly as in
\cref{p2:thm:branch-inverse}, the slope now being
\(\tfrac1p\log X=\Lambda/p\).
\end{proof}

\begin{remark}
\Cref{p2:eq:multi-scaling} explains and replaces the ``scaling from the
half-shifted inverse Gamma problem'' of \textsf{S1} (\(p_n=4^nd_{2n-1}\)) and
the identity \(d^{!!}_{2n-1}=4^nd^{\Gamma}_{2n-1}\) of \textsf{S2}: both are
the case \((p,1)=(2,1)\).  It was verified symbolically for \(n\le4\) and
\(p\in\{2,3\}\).  As a coefficient-polynomial identity it says nothing about
the numerical Lambert data, which differ between the two problems; the
identity is in \(\Lambda\), not in \(y\).  It also supplies a strong regression
test: an implementation storing only \(d^{(1)}\) reproduces every
multifactorial.
\end{remark}

\begin{center}
\small
\begin{tabular}{cccccc}
\toprule
\(p\) & object & centre \(u\) & phase & \(a^{(p)}_1\) & Borel poles\\
\midrule
1 & \(\Gamma(x+1)\) & \(x+\tfrac12\) & \(u(\log u-1)\) & \(-\tfrac1{24}\)
  & \(\pm2\pi\ii m\)\\
2 & \(n!!\) & \(x+1\) & \(\tfrac u2(\log u-1)\) & \(-\tfrac16\)
  & \(\pm\pi\ii m\)\\
3 & \(n!!!\) & \(x+\tfrac32\) & \(\tfrac u3(\log u-1)\) & \(-\tfrac38\)
  & \(\pm\tfrac{2\pi}3\ii m\)\\
\(p\) & \(x!_{(p)}\) & \(x+\tfrac p2\) & \(\tfrac up(\log u-1)\)
  & \(-\tfrac{p^2}{24}\) & \(\pm\tfrac{2\pi}p\ii m\)\\
\bottomrule
\end{tabular}
\end{center}

\section{Algorithms and certification}
\label{p2:sec:algorithms}

\begin{algorithm}[Forward coefficient table]
\label{p2:alg:forward}
Given \(M\) and a power \(\alpha\in\Rat\): compute \(c_k\) from
\cref{p2:eq:ck} for \(1\le k\le\lceil M/2\rceil\); set \(\gamma_{2k-1}=c_k\),
\(\gamma_{2k}=0\); run \cref{p2:eq:salpha-recurrence}.  This costs
\(\bigO(M^2)\) rational operations and never forms a symbolic exponential.
\end{algorithm}

\begin{algorithm}[Inverse blocks]
\label{p2:alg:inverse}
On the variable \(z=q^2\), with \(\Lambda\) symbolic:
\begin{lstlisting}
input: N
E := 0
for n := 1,...,N:
    H := (1+E)*log(1+E) - E
    for k := 1,...,n:
        H := H + a[k]*z^k*(1+E)^(-(2*k-1))
    H := truncate(H, z^(n+1))
    e[n] := - coefficient(H, z^n) / Lambda
    E := E + e[n]*z^n
    assert truncate(residual(E), z^(n+1)) = 0      # (2.residual)
    assert sign_invariant(e[n], n)                 # Theorem sign law
return e[1],...,e[N]
\end{lstlisting}
Every product, logarithm and negative power may be truncated above \(z^n\) at
step \(n\); no full symbolic logarithm is ever needed.  Representing each
\(e_n\) as a dense polynomial in \(\Lambda^{-1}\) of degree exactly \(2n-1\)
(\cref{p2:thm:sign-law}) gives an exact allocation size.
\end{algorithm}

The two assertions are worth more than a comparison against a stored table.
The first is the \emph{formal residual certificate}: with
\(E_N=\sum_{n\le N}e_nz^n\),
\begin{equation}
 \Lambda E_N+\varphi(E_N)
 +\sum_{k=1}^{N+1}a_kz^k(1+E_N)^{-(2k-1)}
 =\bigO\!\left(z^{N+1}\right),
 \label{p2:eq:residual-certificate}
\end{equation}
and the coefficient of \(z^{N+1}\), divided by \(-\Lambda\), \emph{is} the next
block.  This single identity catches sign, indexing and centring mistakes at
once.  The second is \cref{p2:thm:sign-law}: after multiplication by
\((-1)^{n+1}\) every one of the \(2n-1\) coefficients must be a positive
rational.  Independent checks available in addition:
\begin{enumerate}[label=\textup{(\arabic*)}]
\item expanding the exact gamma ratio of \cref{p2:thm:normal-form} reproduces
      \(c_k\);
\item expanding the Borel kernel \cref{p2:eq:mittag} reproduces \(c_k\) from
      \(\eta\)-values, and \cref{p2:prop:uncentred} reproduces \(b_k\) from
      \(\zeta\)-values;
\item the centre-shift identity \cref{p2:eq:centre-shift} links the two;
\item \cref{p2:eq:closed} and \cref{p2:eq:lagrange-closed} must agree with
      \cref{p2:eq:recurrence} coefficient by coefficient;
\item \cref{p2:eq:multi-scaling} compares every block with the independently
      generated \(p=1\) family.
\end{enumerate}
Every displayed block \(d_1,\dots,d_{13}\) of this chapter passes all five.

\begin{algorithm}[Stable evaluation from a logarithmic target]
\label{p2:alg:evaluate}
Never form \(y\); accept \(\ell=\log y\).  For branch \(\delta\):
\begin{equation}
 R=\ell-A_\delta,
 \quad
 w=\Wlam_0(2R/\e),
 \quad
 T=2R/w,
 \quad
 \Lambda=w+1,
 \quad
 x\leftarrow T-1+\sum_{n=1}^{N}\frac{d_{2n-1}(\Lambda)}{T^{2n-1}} .
 \label{p2:eq:stable-eval}
\end{equation}
Refine by Newton on the logarithmic residual
\(F_\delta(x)=\log\kappa_\delta+\tfrac x2\log2+\log\Gamma(\tfrac x2+1)-\ell\),
\begin{equation}
 x_{\mathrm{new}}=x-\frac{F_\delta(x)}
 {\tfrac12\bigl(\log2+\psi(\tfrac x2+1)\bigr)},
 \label{p2:eq:newton}
\end{equation}
whose denominator is positive by \cref{p2:eq:branch-logder}; convergence is
quadratic once the asymptotic value is in the basin, which a few blocks always
achieve in the ranges tested.  For \(\mathscr D\), add
\(\tfrac14(1-\cos\pi x)\log(2/\pi)\) to \(F_0\) and its derivative to the
denominator, and start from \cref{p2:eq:oeis-inverse}.
\end{algorithm}

\begin{proposition}[A posteriori certificate]
\label{p2:prop:aposteriori}
Let \(\widetilde x\) be any approximation and \(\rho=F_\delta(\widetilde x)\).
If \(I\) is an interval containing \(\widetilde x\) and
\(\mathcal D_\delta^{-1}(y)\), then
\begin{equation}
 \bigl|\widetilde x-\mathcal D_\delta^{-1}(y)\bigr|
 \le\frac{2|\rho|}{\inf_{\xi\in I}\bigl(\log2+\psi(\xi/2+1)\bigr)} ,
 \label{p2:eq:aposteriori}
\end{equation}
and since \(\psi\) is increasing the infimum is attained at the left endpoint.
At large arguments the denominator is \(\sim\Lambda\), matching
\cref{p2:eq:branch-inverse}.
\end{proposition}

\begin{proof}
This is \cref{p0:thm:backward-error} with the explicit slope
\cref{p2:eq:branch-logder}; the mean-value theorem applied to \(F_\delta\) on
\(I\) gives \cref{p2:eq:aposteriori}.
\end{proof}

\section{Numerical validation}
\label{p2:sec:numerics}

All tables in this section were recomputed at \(220\)-digit precision with
\texttt{mpmath}: the exact target \(\ell=\log\mathcal D_\delta(n)\) is formed
from \(\log\Gamma\), the Lambert data from \cref{p2:eq:stable-eval}, and the
blocks from \cref{p2:eq:recurrence} in exact rational arithmetic.  No rounded
value of \(y\) is ever formed.  Entries are \emph{signed}
(approximation minus exact), which the source tables partly suppress.

\begin{table}[htbp]
\centering\small
\caption{Forward: centred logarithmic expansion, \(\Phi(u)+A_\delta+
\sum_{k\le N}c_ku^{1-2k}-\log\mathcal D_\delta(n)\), \(u=n+1\).  Signs
alternate, confirming \cref{p2:thm:remainder}.}
\label{p2:tab:fwd-log}
\begin{tabular}{c c r r r r r}
\toprule
\(\delta\)&\(n\)&\(N=0\)&\(N=1\)&\(N=2\)&\(N=3\)&\(N=4\)\\
\midrule
0&20 & \(3.96616\cdot10^{-3}\) & \(-2.09362\cdot10^{-6}\)
    & \(5.98269\cdot10^{-9}\) & \(-4.14413\cdot10^{-11}\)
    & \(5.30661\cdot10^{-13}\)\\
1&21 & \(3.78606\cdot10^{-3}\) & \(-1.82137\cdot10^{-6}\)
    & \(4.74399\cdot10^{-9}\) & \(-2.99567\cdot10^{-11}\)
    & \(3.49749\cdot10^{-13}\)\\
0&100& \(8.25064\cdot10^{-4}\) & \(-1.88702\cdot10^{-8}\)
    & \(2.34020\cdot10^{-12}\) & \(-7.04697\cdot10^{-16}\)
    & \(3.92938\cdot10^{-19}\)\\
1&101& \(8.16975\cdot10^{-4}\) & \(-1.83207\cdot10^{-8}\)
    & \(2.22773\cdot10^{-12}\) & \(-6.57742\cdot10^{-16}\)
    & \(3.59602\cdot10^{-19}\)\\
\bottomrule
\end{tabular}
\end{table}

\begin{table}[htbp]
\centering\small
\caption{Forward: relative error \(|F_M/\mathcal D_\delta-1|\) of the centred
multiplicative expansion \cref{p2:eq:mult-transseries} truncated after
\(s_Mu^{-M}\), with \(\delta\equiv x\pmod2\).}
\label{p2:tab:fwd-mult}
\begin{tabular}{r c r r r r r}
\toprule
\(x\)&\(\delta\)&\(M=0\)&\(M=1\)&\(M=3\)&\(M=5\)&\(M=9\)\\
\midrule
10&0& \(7.590\cdot10^{-3}\) & \(4.330\cdot10^{-5}\) & \(2.599\cdot10^{-7}\)
   & \(4.974\cdot10^{-9}\) & \(1.388\cdot10^{-11}\)\\
11&1& \(6.957\cdot10^{-3}\) & \(3.538\cdot10^{-5}\) & \(1.751\cdot10^{-7}\)
   & \(2.782\cdot10^{-9}\) & \(5.454\cdot10^{-12}\)\\
20&0& \(3.974\cdot10^{-3}\) & \(9.988\cdot10^{-6}\) & \(1.432\cdot10^{-8}\)
   & \(6.756\cdot10^{-11}\) & \(1.314\cdot10^{-14}\)\\
21&1& \(3.793\cdot10^{-3}\) & \(9.014\cdot10^{-6}\) & \(1.166\cdot10^{-8}\)
   & \(4.972\cdot10^{-11}\) & \(7.956\cdot10^{-15}\)\\
50&0& \(1.635\cdot10^{-3}\) & \(1.483\cdot10^{-6}\) & \(3.106\cdot10^{-10}\)
   & \(2.113\cdot10^{-13}\) & \(9.594\cdot10^{-19}\)\\
51&1& \(1.604\cdot10^{-3}\) & \(1.424\cdot10^{-6}\) & \(2.861\cdot10^{-10}\)
   & \(1.866\cdot10^{-13}\) & \(7.800\cdot10^{-19}\)\\
\bottomrule
\end{tabular}
\end{table}

\begin{table}[htbp]
\centering\small
\caption{Branchwise inverse: \(x_N-x\) with \(y=\mathcal D_\delta(x)\) exact
and \(x_N\) the truncation of \cref{p2:eq:branch-inverse} after \(N\) blocks.
Signs alternate, as \cref{p2:thm:sign-law} predicts for the dominant term.}
\label{p2:tab:inv}
\begin{tabular}{c r r r r r r}
\toprule
\(\delta\)&\(x\)&\(N=0\)&\(N=1\)&\(N=2\)&\(N=3\)&\(N=5\)\\
\midrule
0& 10 & \(-6.3074\cdot10^{-3}\) & \(1.6444\cdot10^{-5}\)
     & \(-1.6097\cdot10^{-7}\) & \(3.6890\cdot10^{-9}\)
     & \(1.1386\cdot10^{-11}\)\\
1& 11 & \(-5.5808\cdot10^{-3}\) & \(1.2104\cdot10^{-5}\)
     & \(-9.9815\cdot10^{-8}\) & \(1.9338\cdot10^{-9}\)
     & \(4.2632\cdot10^{-12}\)\\
0& 20 & \(-2.6055\cdot10^{-3}\) & \(1.7520\cdot10^{-6}\)
     & \(-4.7751\cdot10^{-9}\) & \(3.1029\cdot10^{-11}\)
     & \(7.6506\cdot10^{-15}\)\\
1& 21 & \(-2.4497\cdot10^{-3}\) & \(1.4956\cdot10^{-6}\)
     & \(-3.7170\cdot10^{-9}\) & \(2.2041\cdot10^{-11}\)
     & \(4.5233\cdot10^{-15}\)\\
0& 50 & \(-8.3109\cdot10^{-4}\) & \(8.9794\cdot10^{-8}\)
     & \(-4.1941\cdot10^{-11}\) & \(4.7156\cdot10^{-14}\)
     & \(3.4429\cdot10^{-19}\)\\
1& 51 & \(-8.1110\cdot10^{-4}\) & \(8.4219\cdot10^{-8}\)
     & \(-3.7845\cdot10^{-11}\) & \(4.0942\cdot10^{-14}\)
     & \(2.7669\cdot10^{-19}\)\\
0&100 & \(-3.5755\cdot10^{-4}\) & \(9.5805\cdot10^{-9}\)
     & \(-1.1469\cdot10^{-12}\) & \(3.3159\cdot10^{-16}\)
     & \(1.5896\cdot10^{-22}\)\\
1&101 & \(-3.5329\cdot10^{-4}\) & \(9.2785\cdot10^{-9}\)
     & \(-1.0892\cdot10^{-12}\) & \(3.0878\cdot10^{-16}\)
     & \(1.4232\cdot10^{-22}\)\\
\bottomrule
\end{tabular}
\end{table}

\begin{table}[htbp]
\centering\small
\caption{Inverse of the periodic interpolation:
\(\bigl(T-1+\widetilde\Delta\bigr)-x\) with \(y=\mathscr D(x)\) exact,
successively adding the carrier \(\Delta_0\) (computed as the exact fixed point
of \cref{p2:eq:carrier-eq}), \(\Delta_1/T\) and \(\Delta_2/T^2\) from
\cref{p2:eq:Delta12}.  The core error matches the predicted carrier amplitude
\(2a/\Lambda\): at \(x=20\), \(2a/\Lambda=0.074080\) against an observed
\(0.071518\); at \(x=100\), \(0.048919\) against \(0.048564\).  No source
validates this sector numerically.}
\label{p2:tab:oeis-inv}
\begin{tabular}{r r r r r}
\toprule
\(x\)&core \(T-1\)&\(+\Delta_0\)&\(+\Delta_1/T\)&\(+\Delta_2/T^2\)\\
\midrule
20     & \(7.1518\cdot10^{-2}\) & \(-2.5603\cdot10^{-3}\)
       & \(-1.2693\cdot10^{-5}\) & \(1.6637\cdot10^{-6}\)\\
21     & \(-7.5539\cdot10^{-2}\) & \(-2.4124\cdot10^{-3}\)
       & \(1.4359\cdot10^{-5}\) & \(1.4221\cdot10^{-6}\)\\
20.5   & \(-2.5254\cdot10^{-3}\) & \(-3.2850\cdot10^{-3}\)
       & \(2.0627\cdot10^{-6}\) & \(2.2808\cdot10^{-6}\)\\
100    & \(4.8564\cdot10^{-2}\) & \(-3.5495\cdot10^{-4}\)
       & \(-2.3071\cdot10^{-7}\) & \(9.4103\cdot10^{-9}\)\\
100.25 & \(3.4218\cdot10^{-2}\) & \(-3.9845\cdot10^{-4}\)
       & \(-1.8912\cdot10^{-7}\) & \(1.0779\cdot10^{-8}\)\\
\bottomrule
\end{tabular}
\end{table}

\begin{table}[htbp]
\centering\small
\caption{Multifactorial inverse \cref{p2:eq:multi-inverse}, using only the
\(p=1\) blocks rescaled by \cref{p2:eq:multi-scaling}.  Entries are
\(x_N-x\) with \(y=\mathcal D_{p,r}(x)\) exact; the \(p=2\) row reproduces
\cref{p2:tab:inv}.}
\label{p2:tab:multi-inv}
\begin{tabular}{c c r r r r r}
\toprule
\(p\)&\(r\)&\(x\)&\(N=0\)&\(N=1\)&\(N=2\)&\(N=3\)\\
\midrule
2&2& 20 & \(-2.60549\cdot10^{-3}\) & \(1.75197\cdot10^{-6}\)
       & \(-4.77514\cdot10^{-9}\) & \(3.10287\cdot10^{-11}\)\\
3&1& 22 & \(-5.05001\cdot10^{-3}\) & \(6.03949\cdot10^{-6}\)
       & \(-2.95265\cdot10^{-8}\) & \(3.44131\cdot10^{-10}\)\\
3&2& 20 & \(-5.67882\cdot10^{-3}\) & \(8.16471\cdot10^{-6}\)
       & \(-4.75923\cdot10^{-8}\) & \(6.60178\cdot10^{-10}\)\\
4&3& 23 & \(-8.27262\cdot10^{-3}\) & \(1.54424\cdot10^{-5}\)
       & \(-1.18248\cdot10^{-7}\) & \(2.15627\cdot10^{-9}\)\\
5&2& 27 & \(-1.04166\cdot10^{-2}\) & \(2.15596\cdot10^{-5}\)
       & \(-1.85311\cdot10^{-7}\) & \(3.79755\cdot10^{-9}\)\\
\bottomrule
\end{tabular}
\end{table}

\begin{remark}
The near-geometric improvement at these moderate arguments is not a claim of
convergence: at fixed \(x\) the Bernoulli-driven expansion diverges, and by
\cref{p2:prop:least-term} the attainable accuracy is bounded below by
\(\asymp\e^{-\pi T}/(\sqrt T\Lambda)\).  In high-precision work, truncate at
the least term or finish with \cref{p2:eq:newton}.  Note also that the
parity dependence is confined to the core: once the correct \(A_\delta\) is
used, the two subsequences show nearly identical error patterns, as
\cref{p2:rem:kappa-cancels} predicts.
\end{remark}

\section{Logical status}
\label{p2:sec:status}

Three kinds of statement have been used, and the chapter keeps them apart.

\paragraph{Exact identities.}
\Cref{p2:eq:branches,p2:eq:oeis,p2:eq:multi-branch} and the interpolation
statements \cref{p2:prop:branch-interp,p2:prop:oeis-interp,p2:prop:multi-branch};
the centred normal forms \cref{p2:eq:normal-form,p2:eq:sequence-exact,p2:eq:oeis-centred};
the Lambert identities \cref{p2:eq:core-identities}; the Borel--Laplace
representations \cref{p2:eq:borel,p2:eq:multi-borel} and the Mittag--Leffler
decomposition \cref{p2:eq:mittag}.  These hold wherever the functions are
defined.

\paragraph{Formal coefficient identities.}
The Bell formulae \cref{p2:eq:salpha-bell,p2:eq:bell-recurrence}, the
recurrences \cref{p2:eq:salpha-recurrence,p2:eq:recurrence,p2:eq:periodic-recurrence},
the closed formulae \cref{p2:eq:closed,p2:eq:lagrange-closed,p2:eq:Dnr}, the
boundary and sign statements \cref{p2:cor:boundary,p2:thm:sign-law}, the
centre-shift identity \cref{p2:eq:centre-shift}, the scaling
\cref{p2:eq:multi-scaling}, and \cref{p2:prop:phase-reversion}.  These are
identities of formal series over \(\Rat(\Lambda)\), or over the ring generated
by \(\Lambda^{-1}\) and \(\e^{\pm\ii\theta}\) in the periodic case.  No
convergence of the Bernoulli series is asserted anywhere.

\paragraph{Poincar\'e asymptotics with remainders.}
\Cref{p2:thm:coefficients,p2:thm:remainder,p2:cor:forward-log} for the forward
direction --- where the remainder is in fact two-sided and sharp;
\cref{p2:thm:branch-inverse,p2:cor:periodic-meaning,p2:eq:multi-inverse} for
the inverse, each obtained from a forward residual by
\cref{p0:thm:backward-error,p0:thm:remainder-transport} together with the
explicit slope \cref{p2:eq:branch-logder}; and
\cref{p2:prop:least-term,p2:prop:transport} for the beyond-all-orders layer.
\Cref{p2:thm:carrier}(b) is a genuine convergence statement, with an explicit
threshold.

\begin{remark}[Practical summary]
\label{p2:rem:summary}
For symbolic work, centre at \(u=x+1\), keep \(\Wlam_0\) exact, generate
\(c_k\) from Bernoulli numbers, exponentiate with complete Bell polynomials,
and invert with the triangular \(z=T^{-2}\) recurrence, checking
\cref{p2:eq:residual-certificate} and the sign law at every step.  Expand
\(\Wlam_0\) into iterated logarithms only as a final presentation layer.  When
the parity of \(n\) is known, use the branch inverse; when it is not, state the
interpolation --- by \cref{p2:thm:no-canonical} its choice changes the leading
sector of the answer.
\end{remark}

%% ==================== fragment p3 ====================
\chapter{The partition numbers (A000041)}\label{p3:sec:top}

\begin{sourcebox}
This chapter merges three independent treatments of the partition numbers filed in the
\texttt{special-function-inversion} collection:
\emph{Partition\_Number\_Transseries\_and\_Asymptotic\_Inverse} (below: \textbf{S1}),
\emph{Partition\_Number\_Transseries\_and\_Inverse} (\textbf{S2}), and
\emph{Partition\_Numbers\_Transseries\_and\_Inverse} (\textbf{S3}). All three take Rademacher's
convergent series as analytic input and build a coefficient calculus on top of it; they differ in
notation, in how the non-integral arithmetic phase is defined, and in which tail estimate is
proved. Everything below that is not attributed is common to all three, restated once. Numerical
constants, coefficient tables and error tables have been recomputed here (exact rationals through
\texttt{fractions}/\texttt{sympy}; $140$-digit \texttt{mpmath} for the tables); the recomputation
is described in \cref{p3:sub:numerics}.
\end{sourcebox}

\section{Orientation: what is specific to \texorpdfstring{$p(n)$}{p(n)}}
\label{p3:sec:orientation}

The partition function $p(n)$ is defined by Euler's product
\begin{equation}\label{p3:eq:euler-product}
 \sum_{n\ge0}p(n)q^{n}=\prod_{m\ge1}\frac1{1-q^{m}},\qquad |q|<1 ,
\end{equation}
and is OEIS A000041. Hardy and Ramanujan proved
\begin{equation}\label{p3:eq:HR}
 p(n)\sim\frac1{4\sqrt3\,n}\exp\!\left(\pi\sqrt{\frac{2n}{3}}\right),
\end{equation}
and Rademacher replaced the asymptotic formula by an exactly convergent series. That single fact
makes this chapter structurally different from \cref{p1:sec:top,p2:sec:top,p4:sec:top}: the
exponentially small sectors are \emph{given}, not reconstructed from the late behaviour of a
divergent perturbative series. Three features are genuinely specific to $p(n)$ and are the subject
of this chapter.

\begin{enumerate}
\item \textbf{An arithmetic sector index.}  The sectors are labelled by a
positive integer $k$ (a root of unity $\e^{2\pi\ii h/k}$) and a sign $\sigma$, and their
amplitudes are the Rademacher sums $A_k(n)$ built from Dedekind sums. These amplitudes are
\emph{periodic} in $n$ with period $k$; they are not constants, and after inversion they become
bounded oscillations whose phase is quadratic in $\log y$.
\item \textbf{Non-well-ordered actions.}  Relative to the dominant sector, the
exponential actions are $1-1/k$ (increasing to $1$) and $1+1/k$ (decreasing to $1$). The second
family has no least element. The support of the completed object is therefore not a locally finite
Hahn grid, and the two signs may not be separated before summing over $k$.
\item \textbf{Convergence where one expects divergence.}  Every fixed sector's
algebraic series in $n^{-1/2}$ has radius of convergence exactly $\sqrt{24}$ in the variable
$n^{-1/2}$, and the reversion series of the dominant sector converges as well. Nothing in this
chapter is a divergent asymptotic series; consequently the least-term machinery of
\cref{p0:thm:optimal-truncation} degenerates here, as recorded in \cref{p3:rem:no-least-term}.
\end{enumerate}

Everything else --- the Lambert core, the reversion around it, the flattening into logarithms, the
passage from a smooth inverse to an integer staircase, and the transport of remainders --- is
\cref{p0:sec:top} material and is \emph{cited}, never re-proved.

\subsection{The exponential--power reduction and the Part~0 dictionary}
\label{p3:sub:dictionary}

Fix throughout
\begin{equation}\label{p3:eq:master-variables}
 N=n-\frac1{24},\qquad
 \lambda=\pi\sqrt{\frac23}=\frac{\pi\sqrt6}{3},\qquad
 \mu=\lambda\sqrt N=\frac\pi6\sqrt{24n-1},\qquad
 \kappa=\frac{\pi^2}{6\sqrt3}.
\end{equation}
The dominant Rademacher sector, derived in \cref{p3:thm:exact-sectors}, is
\begin{equation}\label{p3:eq:dominant-in-N}
 P_{1,+}(n)=\frac{\e^{\lambda\sqrt N}}{4\sqrt3\,N}
 \left(1-\frac1{\lambda\sqrt N}\right).
\end{equation}
Read as a function of $N$, \cref{p3:eq:dominant-in-N} is exactly an instance of the
exponential--power model \cref{p0:def:model}, $F(x)=C\exp(ax^{\alpha})x^{b}\exp Q(x^{-\delta})$,
with
\begin{equation}\label{p3:eq:model-in-N}
 x=N,\quad C=\frac1{4\sqrt3},\quad a=\lambda,\quad \alpha=\frac12,\quad
 b=-1,\quad \delta=\frac12,\quad
 Q(t)=\log\!\left(1-\frac{t}{\lambda}\right)
      =-\sum_{r\ge1}\frac{t^{r}}{r\lambda^{r}} .
\end{equation}
So the partition problem is the case $\alpha=1/2$: a \emph{square-root} phase, not a linear one.
It is reduced to the linear case in two steps: the monomial substitution
$\xi=x^{\alpha}=\sqrt N$ of \cref{p0:lem:alpha-reduction}, followed by the elementary linear
rescaling $\mu=\lambda\xi$ that normalises the exponential rate to $1$.

\begin{lemma}[The $\alpha$-reduction for $p(n)$]\label{p3:lem:alpha-reduction}
\emph{(i)} \Cref{p0:lem:alpha-reduction} applied to \cref{p3:eq:model-in-N} with
$\xi=\sqrt N$ gives an $\alpha=1$ model with data
$\bigl(\tfrac1{4\sqrt3},\ \lambda,\ 1,\ -2,\ 1,\ Q\bigr)$, $Q(t)=\log(1-t/\lambda)$; note
$b/\alpha=-2$ and $\delta/\alpha=1$, so this is the normalisation $\delta=1$ assumed from
\cref{p0:sec:top} onwards.  \emph{(ii)} The further substitution $\mu=\lambda\xi$ rescales any
$\alpha=1$ model with data $(C,a,1,b,1,Q)$ to one with data
$\bigl(Ca^{-b},\,1,\,1,\,b,\,1,\,Q(a\,\cdot)\bigr)$.  Here the two steps give
\begin{equation}\label{p3:eq:model-in-mu}
 P_{1,+}=\kappa\,\frac{\e^{\mu}}{\mu^{2}}\left(1-\frac1\mu\right)
 =C'\e^{a'\mu}\mu^{b'}\exp Q'(\mu^{-\delta'}),
\end{equation}
with the dictionary
\[
C'=\kappa=\frac{\pi^{2}}{6\sqrt3},\qquad a'=1,\qquad \alpha'=1,\qquad b'=-2,\qquad
\delta'=1,\qquad Q'(t)=\log(1-t).
\]
The inverse map back to the index is exact and quadratic:
\begin{equation}\label{p3:eq:n-of-mu}
 n=\frac1{24}+\frac{3\mu^{2}}{2\pi^{2}} .
\end{equation}
\end{lemma}

\begin{proof}
(i) is \cref{p0:lem:alpha-reduction} at $\alpha=\delta=1/2$, $b=-1$, which leaves $C$, $a$ and
the coefficients of $Q$ unchanged and replaces $(b,\delta)$ by $(b/\alpha,\delta/\alpha)$.
For (ii), substituting $\xi=\mu/a$ in $C\e^{a\xi}\xi^{b}\exp Q(\xi^{-1})$ gives
$Ca^{-b}\e^{\mu}\mu^{b}\exp Q(a/\mu)$, and $t\mapsto Q(at)$ again lies in $t\,\Real[[t]]$.
Here $a=\lambda$, $b=-2$, so $Ca^{-b}=\lambda^{2}/(4\sqrt3)=\pi^{2}/(6\sqrt3)=\kappa$ (using
$\lambda^{2}=2\pi^{2}/3$) and $Q'(t)=Q(\lambda t)=\log(1-t)$.  Composing, $\mu=\lambda\sqrt N$
and \cref{p3:eq:dominant-in-N} become \cref{p3:eq:model-in-mu}.  Inverting
$\mu=\lambda\sqrt{n-1/24}$ gives \cref{p3:eq:n-of-mu}, again by $\lambda^{2}=2\pi^{2}/3$.
\end{proof}

\begin{remark}[The model hypotheses hold exactly, and $Q$ converges]
\label{p3:rem:model-exact}
\Cref{p0:def:model} asks for the analytic estimates \cref{p0:eq:model-analytic} and
\cref{p0:eq:model-derivative} in addition to the formal shape.  For the partition envelope both
hold \emph{exactly and trivially}: $F(\mu)=\kappa\e^{\mu}\mu^{-2}(1-1/\mu)$ is an elementary
closed form, so $\log F(\mu)-\log\kappa-\mu+2\log\mu=\log(1-1/\mu)$ is literally the sum of the
series $Q'$, and $(\log F)'=\Phi'$ is computed in closed form in
\cref{p3:lem:Phi-derivatives}; no asymptotic estimate is being differentiated.  In particular
$Q'(t)=\log(1-t)$ is \emph{convergent}, with radius $1$.  \Cref{p0:rem:formal-analytic} states
that $Q$ is divergent of Gevrey type in all four chapters of this volume; that is not so here,
and the discrepancy should be read as a further instance of
\cref{p3:rem:no-least-term} --- the partition chapter is the one place in the volume where the
whole apparatus applies to a convergent object.
\end{remark}

\begin{remark}[Why this is the point of the whole construction]
\label{p3:rem:reduction-is-the-point}
After \cref{p3:lem:alpha-reduction} nothing in the inversion of the dominant sector is
partition-specific. The three sources each rediscover, in their own notation, the Lambert core of
\cref{p0:thm:lambert-core}, the centred reversion of \cref{p0:thm:lambert-centered} and the
logarithmic flattening of \cref{p0:thm:flattening}. We record the dictionary once
(\cref{p3:tab:dictionary}) and cite. What is \emph{not} generic --- and occupies the rest of the
chapter --- is the arithmetic: the sector index, the Dedekind-sum amplitudes, the accumulation of
actions at $1$, and the interaction between the periodic amplitudes and the $\bigO(1)$ shift
produced by the reversion.
\end{remark}

\begin{table}[htbp]
\centering
\caption{Parameter dictionary from \cref{p0:sec:top} to the partition
chapter.  All Part~0 statements are used at these values.}
\label{p3:tab:dictionary}
\small
\begin{tabular}{@{}lll@{}}
\toprule
Part~0 object & Partition value & Reference \\
\midrule
model variable $x$ & $\mu=\pi\sqrt{2(n-1/24)/3}$ & \cref{p3:eq:master-variables}\\
$C$ & $\kappa=\pi^{2}/(6\sqrt3)$ & \cref{p3:lem:alpha-reduction}\\
$a,\ \alpha$ & $1,\ 1$ (after reduction and rescaling; $\lambda,\ 1/2$ before) & \cref{p3:lem:alpha-reduction}\\
$b$ & $-2$ & \cref{p3:lem:alpha-reduction}\\
$Q(t),\ \delta$ & $\log(1-t),\ 1$ & \cref{p3:lem:alpha-reduction}\\
Lambert core $X$ & $-2\Wlam_{-1}\!\left(-\tfrac12\sqrt{\kappa/y}\right)$ & \cref{p3:prop:lambert-core}\\
$\Psi(t,u)$ of \cref{p0:thm:lambert-centered} & $3\log(1+tu)-\log(1-t+tu)$ & \cref{p3:thm:core-correction}\\
reversion coefficients $c_m$ & $1,\tfrac52,\tfrac{13}3,\tfrac{53}{12},-\tfrac{14}5,\dots$ & \cref{p3:tab:cm}\\
index map $x\mapsto n$ & $n=\tfrac1{24}+3\mu^{2}/(2\pi^{2})$ & \cref{p3:eq:n-of-mu}\\
staircase separation & unit spacing of $\Int$; see \cref{p3:thm:rounding} & \cref{p0:thm:staircase}\\
\bottomrule
\end{tabular}
\end{table}

\begin{remark}[Notation across the three sources]\label{p3:rem:notation-clash}
The sources disagree on symbols in a way that has to be stated once, because two of the clashes
are dangerous.
\begin{enumerate}
\item \textbf{The letter $x$.}  S1 writes $N=n-1/24$ for the shifted index and
$x$ for the phase $\pi\sqrt{2N/3}$; S2 writes $x=n-1/24$ for the shifted index and $\mu$ for the
phase; S3 writes $N$ and $\mu$. This volume uses $N$ for the shifted index and $\mu$ for the
phase, i.e.\ the S3 convention, which is the intersection of the other two.
\item \textbf{The shift itself.}  Every displayed \emph{exact} identity in this
chapter is in the shifted variable $N=n-1/24$. Expansions in $n^{-1/2}$
(\cref{p3:sec:half-powers}) re-expand the shift and are the only place where unshifted $n$ appears
inside a coefficient. Confusing the two is the single most common error in this subject: e.g.\
replacing $N$ by $n$ in \cref{p3:eq:dominant-in-N} costs a relative $\bigO(n^{-1/2})$ and
destroys every statement about exponentially small sectors.
\item \textbf{The normalising constant} is $\mathcal{C}$ in S1, $\kappa$ in S2, $D$
in S3; all equal $\pi^{2}/(6\sqrt3)$. We write $\kappa$, reserving $D$ for nothing and $\Phi'$ for
the slope.
\item \textbf{The logarithmic variable.}  S1 and S2 expand in
$R=\log(y/\kappa)$; S3 expands in $H=\log(4y/\kappa)=R+2\log2$. The two expansions have the
\emph{same} coefficient polynomials (see \cref{p3:rem:H-versus-R}) but are different series and
must not be mixed.
\end{enumerate}
\end{remark}

\section{Dedekind sums and the arithmetic amplitudes}\label{p3:sec:arithmetic}

\subsection{Definitions}

\begin{definition}[Sawtooth and Dedekind sum]\label{p3:def:dedekind}
For $u\in\Real$ put
\begin{equation}\label{p3:eq:sawtooth}
 ((u))=\begin{cases} u-\lfloor u\rfloor-\tfrac12,& u\notin\Int,\\ 0,& u\in\Int,\end{cases}
\end{equation}
and for $k\ge1$, $h\in\Int$,
\begin{equation}\label{p3:eq:dedekind}
 s(h,k)=\sum_{r=1}^{k-1}\left(\!\left(\frac rk\right)\!\right)
 \left(\!\left(\frac{hr}{k}\right)\!\right),
\end{equation}
with the empty-sum convention $s(h,1)=0$.
\end{definition}

Throughout, $U_{k}=\{h:0\le h<k,\ \gcd(h,k)=1\}$, so that $U_{1}=\{0\}$.

\begin{definition}[Rademacher amplitude]\label{p3:def:Ak}
For $k\ge1$ and $n\in\Int$,
\begin{equation}\label{p3:eq:Ak}
 A_{k}(n)=\sum_{h\in U_{k}}
 \exp\!\left(\pi\ii\,s(h,k)-\frac{2\pi\ii hn}{k}\right).
\end{equation}
\end{definition}

\begin{lemma}[Elementary properties]\label{p3:lem:Ak-properties}
For every $k\ge1$ and $n\in\Int$:
\begin{enumerate}
\item $A_{k}(n)\in\Real$, and
$A_{k}(n)=\sum_{h\in U_{k}}\cos\!\bigl(\pi s(h,k)-2\pi hn/k\bigr)$;
\item $A_{k}(n+k)=A_{k}(n)$;
\item $|A_{k}(n)|\le\varphi(k)\le k$;
\item $A_{1}(n)=1$ and $A_{2}(n)=(-1)^{n}$.
\end{enumerate}
\end{lemma}

\begin{proof}
(i) For $k\ge2$ the involution $h\mapsto k-h$ is a bijection of $U_{k}$ (note $k-h\ne h$ unless
$k=2h$, impossible for $\gcd(h,k)=1$, $k\ge3$; for $k=2$ the set is $\{1\}$ and $s(1,2)=0$ makes
the single term real). The reciprocity $s(-h,k)=-s(h,k)$, immediate from $((-u))=-((u))$, shows
that the terms at $h$ and $k-h$ are complex conjugates. Hence the sum is real and equals the sum
of the real parts, which is the stated cosine sum. (ii) Replacing $n$ by $n+k$ changes each
exponent by $-2\pi\ii h$, an integer multiple of $2\pi\ii$. (iii) Triangle inequality on
$\varphi(k)$ unimodular terms. (iv) $U_{1}=\{0\}$ and $s(0,1)=0$ give $A_{1}\equiv1$;
$U_{2}=\{1\}$, $s(1,2)=0$ give $A_{2}(n)=\e^{-\pi\ii n}=(-1)^{n}$.
\end{proof}

Only the trivial bound (iii) is used anywhere in this chapter. Weil-type bounds for these
Kloosterman-like sums would sharpen the constants in \cref{p3:thm:tail,p3:thm:inverse-tail} but
change no structure; see \cref{p3:rem:sharper-bounds}.

\begin{proposition}[The first amplitudes]\label{p3:prop:first-Ak}
With $s(1,k)=\dfrac{(k-1)(k-2)}{12k}$,
\begin{align}
 A_{1}(n)&=1, &
 A_{2}(n)&=(-1)^{n}, \nonumber\\
 A_{3}(n)&=2\cos\!\left(\frac\pi{18}-\frac{2\pi n}{3}\right), &
 A_{4}(n)&=2\cos\!\left(\frac\pi8-\frac{\pi n}{2}\right),
 \label{p3:eq:first-Ak}\\
 A_{5}(n)&=2\cos\!\left(\frac\pi5-\frac{2\pi n}{5}\right)+2\cos\frac{4\pi n}{5},&
 A_{6}(n)&=2\cos\!\left(\frac{5\pi}{18}-\frac{\pi n}{3}\right).\nonumber
\end{align}
\end{proposition}

\begin{proof}
$s(1,k)=\sum_{r=1}^{k-1}((r/k))^{2}=\sum_{r=1}^{k-1}(r/k-1/2)^{2}
=\frac{(k-1)(2k-1)}{6k}-\frac{k-1}{2}+\frac{k-1}{4} =\frac{(k-1)(k-2)}{12k}$. Hence
$s(1,3)=\frac1{18}$, $s(1,4)=\frac18$, $s(1,5)=\frac15$, $s(1,6)=\frac5{18}$, and
$s(k-1,k)=-s(1,k)$. For $k=3,4,6$ the set $U_{k}$ is $\{1,k-1\}$, and the two terms are conjugate,
giving twice the real part of the $h=1$ term after reducing $2\pi(k-1)n/k\equiv-2\pi n/k
\pmod{2\pi}$. For $k=5$, $U_{5}=\{1,2,3,4\}$ and a direct evaluation of \cref{p3:eq:dedekind}
gives $s(2,5)=s(3,5)=0$: with
$((1/5),(2/5),(3/5),(4/5))=(-\tfrac3{10},-\tfrac1{10},\tfrac1{10},\tfrac3{10})$, the four products
in $s(2,5)$ are $\tfrac3{100},-\tfrac3{100},-\tfrac3{100},\tfrac3{100}$. Pairing $h=1$ with $h=4$
and $h=2$ with $h=3$ gives the stated two cosines.
\end{proof}

\subsection{The real-analytic interpolation of the amplitudes}
\label{p3:sub:interpolation-amplitudes}

The inverse problem needs the amplitudes at non-integral arguments: the smooth inverse of a smooth
interpolation evaluates $A_{k}$ at a real point. The finite Fourier shape of \cref{p3:eq:Ak}
supplies a canonical choice.

\begin{definition}[Entire amplitude interpolation]\label{p3:def:Atilde}
For $k\ge1$ and $\nu\in\Cplx$ set
\begin{equation}\label{p3:eq:Atilde}
 \widetilde A_{k}(\nu)=\sum_{h\in U_{k}}
 \cos\!\left(\pi s(h,k)-\frac{2\pi h\nu}{k}\right).
\end{equation}
\end{definition}

\begin{lemma}\label{p3:lem:Atilde}
$\widetilde A_{k}$ is entire, $k$-periodic, real on $\Real$, satisfies $\widetilde
A_{k}(n)=A_{k}(n)$ for $n\in\Int$, and $|\widetilde
A_{k}^{(r)}(\nu)|\le\varphi(k)(2\pi/k)^{r}\cosh(2\pi|\Im\nu|) \le k\,(2\pi)^{r}\e^{2\pi|\Im\nu|}$
for every $r\ge0$. Moreover $\widetilde A_{1}\equiv1$ and $\widetilde A_{2}(\nu)=\cos\pi\nu$,
while
\begin{equation}\label{p3:eq:Atilde3}
 \widetilde A_{3}(\nu)=2\cos(\pi\nu)\,
 \cos\!\left(\frac\pi{18}+\frac{\pi\nu}{3}\right).
\end{equation}
\end{lemma}

\begin{proof}
Each summand is entire and $k$-periodic in $\nu$; realness on $\Real$ is manifest; agreement at
integers is \cref{p3:lem:Ak-properties}(i). The derivative bound is termwise, using $|h/k|\le1$
and $|\cos^{(r)}(z)|\le\cosh(\Im z)$. For $k=1$, $U_{1}=\{0\}$ makes the single term $\cos0=1$
identically. For $k=3$, sum-to-product on $A=\frac\pi{18}-\frac{2\pi\nu}3$ and
$B=-\frac\pi{18}-\frac{4\pi\nu}3$ gives $\tfrac{A+B}2=-\pi\nu$ and
$\tfrac{A-B}2=\frac\pi{18}+\frac{\pi\nu}3$, whence \cref{p3:eq:Atilde3}.
\end{proof}

\begin{repairbox}
\textbf{S2, Definition ``Real Rademacher amplitude''.} S2 defines the interpolation as a sum over
$1\le h\le k$ with $\gcd(h,k)=1$. For $k\ge2$ this is the same index set as $U_{k}$, but for $k=1$
it produces the single term $\cos(2\pi\nu)$ instead of the constant $1$. S2 nevertheless uses
$\mathcal A_{1}\equiv1$ throughout (its dominant sector carries no amplitude, and its
relative-correction formula separates the $k=1$ terms with amplitude $1$), so the definition
contradicts its own use: with S2's index set the interpolating function $\mathcal{P}$ would not be
an interpolation of $p$ --- its dominant term would oscillate off-lattice with period $1$ at full
relative strength, and the ``eventual monotonicity'' proposition would be false. \textbf{Fix.}
Take $h$ over $U_{k}=\{0\le h<k:\gcd(h,k)=1\}$ with $U_{1}=\{0\}$, i.e.\ \cref{p3:def:Atilde};
equivalently, decree $\widetilde A_{1}\equiv1$ separately, as S3 does. All of S2's later formulas
are then correct as printed.
\end{repairbox}

\begin{remark}[Non-uniqueness, honestly stated]\label{p3:rem:interp-nonunique}
\cref{p3:def:Atilde} is \emph{a} choice, not \emph{the} choice. Any strictly increasing smooth
$\mathcal{P}$ with $\mathcal{P}(n)=p(n)$ has an inverse, and the algebraic (Lambert--logarithmic)
part of that inverse is the same for all of them, because it is determined by the dominant sector,
which is interpolation-independent once $N$ is treated as a real variable. The exponentially small
sectors are \emph{not} interpolation-independent. \cref{p3:def:Atilde} is singled out only because
it is entire, finite-Fourier, real on $\Real$, and literally the analytic continuation of the
amplitudes already present in Rademacher's theorem. \cref{p3:sec:profinite} gives the alternative
that avoids interpolation altogether.
\end{remark}

\section{Rademacher's formula in exact transseries normal form}
\label{p3:sec:exact}

\subsection{The analytic input}

\begin{theorem}[Hardy--Ramanujan--Rademacher; quoted]\label{p3:thm:rademacher}
For every integer $n\ge1$,
\begin{equation}\label{p3:eq:rademacher}
 p(n)=\frac1{\pi\sqrt2}\sum_{k=1}^{\infty}A_{k}(n)\sqrt k\,
 \frac{\dd}{\dd n}
 \left[\frac{\sinh\!\bigl(\tfrac{\lambda}{k}\sqrt{n-\tfrac1{24}}\bigr)}
 {\sqrt{n-\tfrac1{24}}}\right],
\end{equation}
the series converging when summed in increasing $k$.
\end{theorem}

This is the only unproved input of the chapter. It is Rademacher's theorem (1937, 1938, 1943); a
textbook proof is in Apostol, and a machine-checked proof exists in Isabelle/HOL (Eberl--Paulson,
2026); see \cref{p3:rem:references}. All three sources take it as given and so do we. Note the
factor $\sqrt k$: some normalisations absorb it into the definition of the arithmetic sum, and
mixing conventions is a standard source of a spurious factor $k$.
\Cref{p3:sub:normalization-crosscheck} records the cross-checks that detect such a slip.

\subsection{Exact two-exponential sectors}

\begin{lemma}[Elementary derivative]\label{p3:lem:derivative}
For $a>0$ and $v>0$,
\begin{equation}\label{p3:eq:derivative}
 \frac{\dd}{\dd v}\Bigl(v^{-1/2}\sinh(a\sqrt v)\Bigr)
 =\frac1{4v^{3/2}}
 \Bigl[(a\sqrt v-1)\e^{a\sqrt v}+(a\sqrt v+1)\e^{-a\sqrt v}\Bigr].
\end{equation}
\end{lemma}

\begin{proof}
With $w=a\sqrt v$ one has $\dd w/\dd v=w/(2v)$ and $v^{-1/2}\sinh w=a\,w^{-1}\sinh w$, so
\[
\frac{\dd}{\dd v}\bigl(v^{-1/2}\sinh(a\sqrt v)\bigr) =a\,\frac{\dd}{\dd w}\!\left(\frac{\sinh
w}{w}\right)\frac{w}{2v} =\frac{a}{2v}\left(\cosh w-\frac{\sinh w}{w}\right)
=\frac{1}{2v^{3/2}}\bigl(w\cosh w-\sinh w\bigr),
\]
using $a=w/\sqrt v$ in the last step. Finally $2(w\cosh w-\sinh w)=(w-1)\e^{w}+(w+1)\e^{-w}$.
\end{proof}

\begin{theorem}[Exact sector decomposition]\label{p3:thm:exact-sectors}
Let $N,\lambda,\mu,\kappa$ be as in \cref{p3:eq:master-variables} and set, for $k\ge1$ and
$\sigma\in\{+1,-1\}$,
\begin{equation}\label{p3:eq:sector}
 P_{k,\sigma}(n)
 =\frac{A_{k}(n)}{4\sqrt{3k}\;N}\,
 \e^{\sigma\mu/k}\left(1-\frac{\sigma k}{\mu}\right)
 =\frac{\kappa}{\mu^{2}}\frac{A_{k}(n)}{\sqrt k}\,
 \e^{\sigma\mu/k}\left(1-\frac{\sigma k}{\mu}\right).
\end{equation}
Then, for every integer $n\ge1$,
\begin{equation}\label{p3:eq:sector-sum}
 p(n)=\sum_{k\ge1}\bigl(P_{k,+}(n)+P_{k,-}(n)\bigr),
\end{equation}
the two signs being paired before the sum over $k$ is formed. Equivalently, with the entire paired
kernel
\begin{equation}\label{p3:eq:g-kernel}
 g(z)=\cosh z-\frac{\sinh z}{z}
 =\sum_{m\ge1}\frac{2m}{(2m+1)!}z^{2m},\qquad g(0)=0,
\end{equation}
one has the manifestly convergent form
\begin{equation}\label{p3:eq:paired-sector}
 p(n)=\frac{2\kappa}{\mu^{2}}\sum_{k\ge1}\frac{A_{k}(n)}{\sqrt k}\,
 g\!\left(\frac\mu k\right).
\end{equation}
\end{theorem}

\begin{proof}
Apply \cref{p3:lem:derivative} with $v=N$ and $a=\lambda/k$, so that $a\sqrt v=\mu/k$, inside
\cref{p3:eq:rademacher}. The $k$th summand becomes
\[
\frac{A_{k}(n)\sqrt k}{\pi\sqrt2}\cdot\frac1{4N^{3/2}} \Bigl[\Bigl(\frac\mu k-1\Bigr)\e^{\mu/k}
+\Bigl(\frac\mu k+1\Bigr)\e^{-\mu/k}\Bigr] =\frac{A_{k}(n)\sqrt
k}{\pi\sqrt2}\cdot\frac{\mu}{4kN^{3/2}} \Bigl[\e^{\mu/k}\Bigl(1-\frac k\mu\Bigr)
+\e^{-\mu/k}\Bigl(1+\frac k\mu\Bigr)\Bigr].
\]
Now $\mu/(\pi\sqrt2\,\sqrt N)=\lambda/(\pi\sqrt2)=1/\sqrt3$, so the prefactor is
$A_{k}(n)/(4\sqrt{3k}\,N)$, which is \cref{p3:eq:sector} summed over $\sigma$. Since
$\mu^{2}=\lambda^{2}N=(2\pi^{2}/3)N$, we get $1/(4\sqrt{3k}\,N)=\kappa/(\sqrt k\,\mu^{2})$, the
second form of \cref{p3:eq:sector}. Adding the two signs and using
$\e^{z}(1-1/z)+\e^{-z}(1+1/z)=2g(z)$ with $z=\mu/k$ gives \cref{p3:eq:paired-sector}. All
manipulations are termwise algebra, so the sum is unchanged.
\end{proof}

\begin{remark}[The three sources' normal forms are the same]
\label{p3:rem:three-normal-forms}
S1 writes $p(n)=\frac{\mathcal{C}}{x^{2}}\sum_{k}\frac{A_{k}}{\sqrt k}
\{\e^{x/k}(1-k/x)+\e^{-x/k}(1+k/x)\}$ (its $x$ is our $\mu$); S2 writes the sectors
$P_{k,\sigma}=\frac{A_{k}}{4\sqrt{3k}\,x} \e^{\sigma\mu/k}(1-\sigma k/\mu)$ (its $x$ is our $N$);
S3 writes $p(n)=D\mu^{-3}\sum_{k}A_{k}\sqrt k [(\tfrac\mu k-1)\e^{\mu/k}+(\tfrac\mu
k+1)\e^{-\mu/k}]$. Factoring $\mu/k$ out of S3's bracket and using $\sqrt k\cdot(1/k)=1/\sqrt k$
turns S3 into S1, and $\kappa/\mu^{2}=1/(4\sqrt{3k}\,N)\cdot\sqrt k$ turns S1 into S2. All three
are \cref{p3:eq:sector}. We adopt the second form of \cref{p3:eq:sector} because it displays the
Part~0 model $\kappa\e^{\mu}\mu^{-2}(1-1/\mu)$ at $k=1$, $\sigma=+$ verbatim.
\end{remark}

\subsection{The dominant envelope and the exact relative factorisation}

\begin{definition}[Dominant envelope, actions]\label{p3:def:envelope}
Put
\begin{equation}\label{p3:eq:envelope}
 F(\mu)=P_{1,+}=\kappa\,\frac{\e^{\mu}}{\mu^{2}}\left(1-\frac1\mu\right)
 =\kappa\,\e^{\mu}\,\frac{\mu-1}{\mu^{3}},
\end{equation}
and for $k\ge1$, $\sigma=\pm1$ (excluding $(1,+)$) the \emph{action}
\begin{equation}\label{p3:eq:action}
 \alpha_{k,\sigma}=1-\frac{\sigma}{k}
 \qquad\text{so that}\qquad
 \alpha_{1,-}=2,\quad \alpha_{k,+}=1-\frac1k,\quad \alpha_{k,-}=1+\frac1k .
\end{equation}
\end{definition}

\begin{proposition}[Exact relative factorisation]\label{p3:prop:relative}
For $\mu>1$,
\begin{equation}\label{p3:eq:relative}
 p(n)=F(\mu)\bigl(1+\mathcal{R}(\mu;n)\bigr),
\end{equation}
where, with the amplitude ratios
\begin{equation}\label{p3:eq:r-ratios}
 r_{k,\sigma}(\mu;n)=\frac{A_{k}(n)}{\sqrt k}\,
 \frac{1-\sigma k/\mu}{1-1/\mu},
\end{equation}
one has the finite-level resolved form, valid for each $K\ge1$,
\begin{equation}\label{p3:eq:resolved}
 \mathcal{R}=r_{1,-}\,\e^{-2\mu}
 +\sum_{k=2}^{K}\sum_{\sigma=\pm1}
 r_{k,\sigma}(\mu;n)\,\e^{-\alpha_{k,\sigma}\mu}
 +\mathcal{R}_{>K}(\mu;n),
\end{equation}
and the exact paired form
\begin{equation}\label{p3:eq:R-exact}
 \mathcal{R}(\mu;n)=\frac{1+\mu^{-1}}{1-\mu^{-1}}\e^{-2\mu}
 +\frac{2\e^{-\mu}}{1-\mu^{-1}}
 \sum_{k\ge2}\frac{A_{k}(n)}{\sqrt k}\,g\!\left(\frac\mu k\right),
\qquad
 \mathcal{R}_{>K}=\frac{2\e^{-\mu}}{1-\mu^{-1}}\sum_{k>K}\frac{A_{k}(n)}{\sqrt k}
 g\!\left(\frac\mu k\right).
\end{equation}
\end{proposition}

\begin{proof}
Divide \cref{p3:eq:sector} by \cref{p3:eq:envelope}: $P_{k,\sigma}/F=\frac{A_{k}}{\sqrt
k}\frac{1-\sigma k/\mu}{1-1/\mu} \e^{(\sigma/k-1)\mu}$, which is \cref{p3:eq:r-ratios} times
$\e^{-\alpha_{k,\sigma}\mu}$; the case $k=1,\sigma=-1$ has $A_{1}=1$ and gives the first term of
\cref{p3:eq:R-exact}. For $k\ge2$ the paired sum
$P_{k,+}+P_{k,-}=\frac{2\kappa}{\mu^{2}}\frac{A_{k}}{\sqrt k} g(\mu/k)$ from
\cref{p3:eq:paired-sector}, divided by $F=\kappa\e^{\mu}\mu^{-2}(1-1/\mu)$, gives the second term.
\end{proof}

\begin{remark}[What the Hardy--Ramanujan formula omits]
\Cref{p3:eq:HR} is $F(\mu)$ after replacing $N$ by $n$ and discarding $1-1/\mu$. Each of those two
simplifications costs a relative $\bigO(n^{-1/2})$; together they account for the entire
algebraic half-power series of \cref{p3:sec:half-powers}. \Cref{p3:eq:relative} loses nothing: it
records every Ford-circle arc, in both signs, with its exact algebraic amplitude.
\end{remark}

\section{Ordering of the sectors and tail estimates}\label{p3:sec:tail}

\subsection{The action-one accumulation}

The relative actions of \cref{p3:eq:action} are
\[
\tfrac12,\ \tfrac23,\ \tfrac34,\ \dots\ \uparrow 1 \qquad\text{and}\qquad 2,\ \tfrac32,\
\tfrac43,\ \dots\ \downarrow 1 .
\]
The first family is increasing with supremum $1$; the second is decreasing with infimum $1$ and
\emph{has no least element}. Two consequences must be stated before any transseries language is
used.

\begin{definition}[Rademacher transseries]\label{p3:def:rademacher-transseries}
A \emph{Rademacher transseries} for $p$ is the datum of
\begin{enumerate}
\item for each $K\ge1$, the finite family of resolved sectors
$\{(k,\sigma):k\le K\}$ with the coefficients of \cref{p3:sec:half-powers}, treated as independent
formal exponential variables (so that coefficient extraction takes place in a finitely generated
multivariate power-series ring), together with
\item the exact paired remainder $\mathcal{R}_{>K}$ of \cref{p3:eq:R-exact}, retained
as one analytic block.
\end{enumerate}
The limit $K\to\infty$ is taken analytically, through the convergence in \cref{p3:thm:rademacher},
and never by claiming that the whole support is a locally finite Hahn grid.
\end{definition}

\begin{remark}[Why the pairing is mandatory]\label{p3:rem:pairing}
For $k\gg\mu$ each of $P_{k,+}$ and $P_{k,-}$ separately contains a term proportional to $\sqrt
k\,/\mu$, and $\sum_{k}\sqrt k/\mu$ diverges. Only in the sum do those pieces cancel:
$2g(\mu/k)=\tfrac23(\mu/k)^{2}+O((\mu/k)^{4})$, which makes the $k$th paired term
$\bigO(k^{-3/2})$. Thus \cref{p3:eq:sector-sum} prescribes an order of summation; it is not
an absolutely rearrangeable double series. All three sources say this; S2 and S3 say it most
sharply, and \cref{p3:def:rademacher-transseries} follows S2.
\end{remark}

\begin{remark}[Truncating ``through $\e^{-\mu}$'' is not a prescription]
\label{p3:rem:no-eta-cutoff}
Because the positive actions accumulate at $1$, an instruction such as ``retain all terms down to
relative order $\e^{-\mu}$'' selects infinitely many sectors and is meaningless. A legitimate
truncation names either a Rademacher level $K$, or an action cutoff $\eta<1$ (which then names
finitely many positive sectors, namely $k\le(1-\eta)^{-1}$), or a target accuracy together with
\cref{p3:thm:tail}.
\end{remark}

\subsection{The tail bound}

\begin{lemma}[Kernel bound]\label{p3:lem:g-bound}
For $z\ge0$,
\begin{equation}\label{p3:eq:g-bound}
 0\le g(z)\le\frac{z^{2}}{3}\cosh z\le\frac{z^{2}}{3}\e^{z},
 \qquad\text{hence}\qquad
 g(z)\le\frac{\cosh 1}{3}z^{2}\quad\text{for }0\le z\le1 .
\end{equation}
The factor $\cosh z$ may not be dropped: $g(1)=\e^{-1}=0.3679\ldots>\tfrac13$.
\end{lemma}

\begin{proof}
By \cref{p3:eq:g-kernel}, $g(z)=\sum_{m\ge1}\frac{2m}{(2m+1)!}z^{2m}$ has non-negative
coefficients, whence $g\ge0$. For the upper bound compare with $\frac{z^{2}}{3}\cosh
z=\sum_{m\ge1}\frac{z^{2m}}{3(2m-2)!}$ termwise: the claim is
$\frac{2m}{(2m+1)!}\le\frac1{3(2m-2)!}$, i.e.\ $6m(2m-2)!\le(2m+1)!$, i.e.\ $6m\le(2m+1)2m(2m-1)$,
which holds for all $m\ge1$ with equality at $m=1$. Monotonicity of $\cosh$ gives the last
inequality on $[0,1]$, and $g(1)=\cosh1-\sinh1=\e^{-1}$ shows the factor cannot be dropped.
\end{proof}

\begin{theorem}[Explicit finite-$K$ tail bound]\label{p3:thm:tail}
Let $\mu>1$ and $K\ge1$. Then
\begin{equation}\label{p3:eq:tail}
 \bigl|\mathcal{R}_{>K}(\mu;n)\bigr|
 \le\frac{4\mu^{2}}{3\,(1-\mu^{-1})\sqrt K}\,
 \exp\!\left[-\left(1-\frac1{K+1}\right)\mu\right].
\end{equation}
The same bound holds with $A_{k}(n)$ replaced by $\widetilde A_{k}(\nu)$ for any real $\nu$, since
$|\widetilde A_{k}|\le\varphi(k)\le k$ as well.
\end{theorem}

\begin{proof}
From \cref{p3:eq:R-exact}, $|A_{k}|\le k$ and \cref{p3:lem:g-bound},
\[
|\mathcal{R}_{>K}| \le\frac{2\e^{-\mu}}{1-\mu^{-1}}\sum_{k>K}\sqrt k\,g\!\left(\frac\mu k\right)
\le\frac{2\e^{-\mu}}{1-\mu^{-1}}\cdot\frac{\mu^{2}}{3} \sum_{k>K}k^{-3/2}\e^{\mu/k}
\le\frac{2\mu^{2}\e^{-\mu+\mu/(K+1)}}{3(1-\mu^{-1})}\sum_{k>K}k^{-3/2},
\]
and $\sum_{k>K}k^{-3/2}\le\int_{K}^{\infty}t^{-3/2}\dd t=2K^{-1/2}$.
\end{proof}

\begin{theorem}[Sharper fixed-$K$ tail]\label{p3:thm:tail-sharp}
For each fixed $K\ge1$, as $\mu\to\infty$,
\begin{equation}\label{p3:eq:tail-sharp}
 \sum_{k>K}\bigl(P_{k,+}+P_{k,-}\bigr)
 =\bigO_{K}\!\left(\mu^{-1/2}\e^{\mu/(K+1)}\right),
 \qquad
 \mathcal{R}_{>K}=\bigO_{K}\!\left(\mu^{3/2}\e^{-(1-1/(K+1))\mu}\right).
\end{equation}
\end{theorem}

\begin{proof}
Split at $k=\mu$. For $K<k\le\mu$ use $|A_{k}|\le k$ and
$2g(\mu/k)\le2\cosh(\mu/k)\le2\e^{\mu/(K+1)}$ in \cref{p3:eq:paired-sector}, so those terms
contribute at most $\frac{2\kappa}{\mu^{2}}\cdot2\e^{\mu/(K+1)}\sum_{k\le\mu}\sqrt k
=\bigO(\mu^{-1/2}\e^{\mu/(K+1)})$ since $\sum_{k\le\mu}\sqrt k=\bigO(\mu^{3/2})$. For
$k>\mu$ put $z=\mu/k\in(0,1)$; then $g(z)\le\tfrac13z^{2}\cosh1$ by \cref{p3:lem:g-bound}, and the
$k$th term is $\bigO(\mu^{-2}\sqrt k\cdot\mu^{2}k^{-2})=\bigO(k^{-3/2})$, whose tail is
$\bigO(\mu^{-1/2})$. Dividing by $F\asymp\kappa\mu^{-2}\e^{\mu}$ gives the relative
statement.
\end{proof}

\begin{remark}[Which bound to use]
\Cref{p3:thm:tail} (from S2) is explicit and uniform in $K$; \cref{p3:thm:tail-sharp} (from S1 and
S3) is sharper in $\mu$ by a factor $\mu^{-1/2}$ but only for fixed $K$. Both are needed: the
first to certify a computation, the second to state clean asymptotic orders. S1 and S3 state only
the second, S2 only the first.
\end{remark}

\begin{remark}[Sharper amplitude bounds]\label{p3:rem:sharper-bounds}
All constants above come from $|A_{k}|\le k$. Lehmer's explicit remainder estimates and the
effective error bounds of Banerjee--Paule--Radu--Schneider (\cref{p3:rem:references}) are much
sharper, and Weil-type bounds for the Kloosterman-like sums $A_{k}$ give
$|A_{k}(n)|\ll_{\varepsilon}k^{1/2+\varepsilon}$ for many $k$. Substituting any such bound in the
proofs of \cref{p3:thm:tail,p3:thm:tail-sharp} improves the powers of $\mu$ but changes no
statement of this chapter, all of which are of the form ``a fixed Rademacher level is
exponentially accurate''. For a \emph{specific} residue class some $A_{k}$ vanish, so the first
non-zero omitted sector can be much smaller than the generic estimate.
\end{remark}

\subsection{An exact condensation at the accumulation level}

Right at the accumulation action $1$, an alternative to sector-by-sector bookkeeping is to sum the
kernel expansion first. This is S3's contribution and has no counterpart in S1 or S2.

\begin{definition}[Rademacher--Kloosterman zeta]\label{p3:def:Zeta}
For $\Re s>2$ and $n\in\Int$ put $\displaystyle
\mathcal{Z}_{n}(s)=\sum_{k\ge1}\frac{A_{k}(n)}{k^{s}}$; the series converges absolutely there by
\cref{p3:lem:Ak-properties}(iii). The same definition with $\widetilde A_{k}(\nu)$ defines
$\mathcal{Z}_{\nu}(s)$ for real $\nu$.
\end{definition}

\begin{theorem}[Exact zeta condensation]\label{p3:thm:zeta-condensation}
For every integer $n\ge1$,
\begin{equation}\label{p3:eq:zeta-condensation}
 p(n)=\kappa\sum_{r\ge1}\frac{4r}{(2r+1)!}\,\mu^{2r-2}\,
 \mathcal{Z}_{n}\!\left(2r+\frac12\right),
\end{equation}
and the double series obtained by expanding $\mathcal{Z}_{n}$ converges absolutely.
\end{theorem}

\begin{proof}
Insert $2g(z)=\sum_{r\ge1}\frac{4r}{(2r+1)!}z^{2r+1}\cdot z^{-1}$ --- more precisely
$2g(z)=\sum_{r\ge1}\frac{4r}{(2r+1)!}z^{2r}$ --- into \cref{p3:eq:paired-sector} with $z=\mu/k$:
\[
p(n)=\frac{\kappa}{\mu^{2}}\sum_{k\ge1}\frac{A_{k}(n)}{\sqrt k}
\sum_{r\ge1}\frac{4r}{(2r+1)!}\frac{\mu^{2r}}{k^{2r}}
=\kappa\sum_{k,r}\frac{4r}{(2r+1)!}\mu^{2r-2}\frac{A_{k}(n)}{k^{2r+1/2}} .
\]
Absolute convergence: $\sum_{k}|A_{k}|k^{-2r-1/2}\le\sum_{k}k^{-2r+1/2}<\infty$ for $r\ge1$, and
$\sum_{r\ge1}\frac{4r}{(2r+1)!}\mu^{2r-2}\zeta(2r-\tfrac12)$ converges for every fixed $\mu$
because $\zeta(2r-\tfrac12)\to1$ and $\sum_{r}\frac{4r\mu^{2r-2}}{(2r+1)!}<\infty$. Tonelli's
theorem justifies the interchange, giving \cref{p3:eq:zeta-condensation}.
\end{proof}

\begin{remark}
\Cref{p3:eq:zeta-condensation} orders nothing by asymptotic size --- it is an exact analytic
rearrangement, not a transseries. Its use is precisely where the transseries language fails: it
packages the whole action-one accumulation into Dirichlet values $\mathcal{Z}_{n}(2r+\tfrac12)$,
which can be used as a single block in the inverse formulas of \cref{p3:sec:full-inverse}.
\end{remark}

\section{All-order coefficients in the unshifted variable}\label{p3:sec:half-powers}

The shifted phase $\mu$ makes each sector elementary; the price is that the conventional scale is
$n^{-1/2}$. The entire infinite algebraic array of this subject comes from re-expanding the single
shift $N=n-1/24$. This section gives all of it, in one formula valid for every sector.

\subsection{The sector generating function}

\begin{theorem}[Sectorwise convergent half-power expansion]
\label{p3:thm:sector-gf}
Fix $k\ge1$ and $\sigma\in\{+1,-1\}$, and set $t=n^{-1/2}$. Then
\begin{equation}\label{p3:eq:sector-unshifted}
 P_{k,\sigma}(n)=\frac{A_{k}(n)}{4\sqrt{3k}\;n}\,
 \exp\!\left(\frac{\sigma\lambda\sqrt n}{k}\right)
 \sum_{m\ge0}c^{(k,\sigma)}_{m}\,n^{-m/2},
\end{equation}
where the $c^{(k,\sigma)}_{m}$ are the Taylor coefficients at $t=0$ of
\begin{equation}\label{p3:eq:Omega}
 \Omega_{k,\sigma}(t)
 =\frac1{1-t^{2}/24}\,
 \exp\!\left[\frac{\sigma\lambda}{kt}
 \left(\sqrt{1-\frac{t^{2}}{24}}-1\right)\right]
 \left(1-\frac{\sigma kt}{\lambda\sqrt{1-t^{2}/24}}\right).
\end{equation}
The function $\Omega_{k,\sigma}$ is holomorphic on $|t|<\sqrt{24}$, so the series in
\cref{p3:eq:sector-unshifted} converges for every real $n>1/24$; in particular
$c^{(k,\sigma)}_{0}=1$.
\end{theorem}

\begin{proof}
With $t=n^{-1/2}$ we have $N=n(1-t^{2}/24)$ and $\mu=\lambda\sqrt N=(\lambda/t)\sqrt{1-t^{2}/24}$.
Substituting in \cref{p3:eq:sector},
\[
P_{k,\sigma} =\frac{A_{k}(n)}{4\sqrt{3k}\,n}\cdot\frac1{1-t^{2}/24}\cdot
\e^{\sigma\mu/k}\left(1-\frac{\sigma kt}{\lambda\sqrt{1-t^{2}/24}}\right),
\]
and $\e^{\sigma\mu/k}=\e^{\sigma\lambda/(kt)}
\exp[\tfrac{\sigma\lambda}{kt}(\sqrt{1-t^{2}/24}-1)]$ with
$\e^{\sigma\lambda/(kt)}=\e^{\sigma\lambda\sqrt n/k}$. This is \cref{p3:eq:sector-unshifted} with
\cref{p3:eq:Omega}. Holomorphy: the only candidate singularities of \cref{p3:eq:Omega} in
$|t|<\sqrt{24}$ are $t=0$ and the zeros of $\sqrt{1-t^{2}/24}$, and the latter lie on
$|t|=\sqrt{24}$. At $t=0$ the exponent is
$\frac{\sigma\lambda}{kt}(\sqrt{1-t^{2}/24}-1)=-\frac{\sigma\lambda t}{48k} +\bigO(t^{3})$,
so the singularity is removable, and the last factor is holomorphic there as well. Setting $t=0$
gives $\Omega_{k,\sigma}(0)=1$.
\end{proof}

\begin{repairbox}
\textbf{S1, Theorem ``Universal sector coefficient generator''.} S1 states only that the series
``converges for all sufficiently small $|t|$ and supplies an asymptotic expansion to arbitrary
order''. This understates the truth and, read literally, invites the error the next remark
forbids: it suggests the half-power series is a merely asymptotic object. S2 and S3 both identify
the radius as exactly $\sqrt{24}$, i.e.\ convergence for every real $n>1/24$;
\cref{p3:thm:sector-gf} adopts their statement with the branch-point argument made explicit.
\textbf{Consequence.} Nothing in this chapter is a divergent series in $n^{-1/2}$, and no
least-term truncation is available; see \cref{p3:rem:no-least-term}.
\end{repairbox}

\subsection{A finite closed formula for every sector coefficient}

Write
\begin{equation}\label{p3:eq:b-parameter}
 \beta=\beta_{k,\sigma}=\frac{\sigma\pi}{6k},
 \qquad\text{so that}\qquad
 \frac{\sigma\lambda}{k}=2\sqrt6\,\beta,
 \qquad
 \frac{\sigma k}{\lambda}=\frac{1}{12\beta}\cdot\sqrt{6}\,,
\end{equation}
the two identities following from $\lambda=\pi\sqrt6/3$ and $k=\sigma\pi/(6\beta)$. Thus each
sector is the \emph{same} function of the single parameter $\beta$; the sector label enters
nowhere else.

\begin{theorem}[Finite coefficient formula, every sector]\label{p3:thm:c-finite}
For every $k\ge1$, $\sigma=\pm1$ and $m\ge0$,
\begin{equation}\label{p3:eq:c-finite}
 c^{(k,\sigma)}_{m}
 =\frac1{(-4\sqrt6)^{m}}
 \sum_{j=0}^{\lfloor(m+1)/2\rfloor}
 \binom{m+1}{j}\,\frac{m+1-j}{(m+1-2j)!}\,
 \beta_{k,\sigma}^{\,m-2j},
\end{equation}
terms with a negative factorial index being absent. In particular
$c^{(k,\sigma)}_{m}\in\Rat[\beta,\beta^{-1}]\cdot\sqrt6^{\,m\bmod 2}$, and the dependence on the
sector is only through $\beta_{k,\sigma}=\sigma\pi/(6k)$.
\end{theorem}

\begin{proof}
Put $u=-t/(4\sqrt6)$, so that $t^{2}/24=4u^{2}$ and, by \cref{p3:eq:b-parameter},
$\frac{\sigma\lambda}{kt}=\frac{2\sqrt6\beta}{-4\sqrt6\,u}=-\frac{\beta}{2u}$ and $\frac{\sigma
kt}{\lambda}=-\frac{4\sqrt6\,u\sqrt6}{12\beta}=-\frac{2u}{\beta}$. Writing $s=\sqrt{1-4u^{2}}$,
\cref{p3:eq:Omega} becomes
\begin{equation}\label{p3:eq:Omega-u}
 G(u):=\Omega_{k,\sigma}(-4\sqrt6\,u)
 =\frac{\e^{\beta T}}{s^{2}}\left(1+\frac{2u}{\beta s}\right),
 \qquad
 T=T(u)=\frac{1-s}{2u},
\end{equation}
where we used $-\frac{\beta}{2u}(s-1)=\beta\frac{1-s}{2u}=\beta T$ and $s^{2}=1-4u^{2}$. The
Catalan variable $T$ satisfies
\begin{equation}\label{p3:eq:catalan}
 T=u(1+T^{2}),\qquad u=\frac{T}{1+T^{2}},\qquad
 s=\frac{1-T^{2}}{1+T^{2}},\qquad
 \dd u=\frac{1-T^{2}}{(1+T^{2})^{2}}\,\dd T .
\end{equation}
(The first identity follows from $2uT=1-s$ and $s^{2}=1-4u^{2}$; the rest are algebra.)
Substituting \cref{p3:eq:catalan} into \cref{p3:eq:Omega-u} gives the exact differential identity
\begin{equation}\label{p3:eq:G-differential}
 G(u)\,\dd u
 =\e^{\beta T}\left(\frac{1}{1-T^{2}}+\frac{2T}{\beta(1-T^{2})^{2}}\right)\dd T
 =\frac1\beta\,\dd\!\left(\frac{\e^{\beta T}}{1-T^{2}}\right),
\end{equation}
because $\frac{\dd}{\dd T}\frac{\e^{\beta T}}{1-T^{2}} =\e^{\beta
T}\bigl(\frac{\beta}{1-T^{2}}+\frac{2T}{(1-T^{2})^{2}}\bigr)$. Now let
$D_{m}=\Coef{u^{m}}G(u)=(-4\sqrt6)^{m}c^{(k,\sigma)}_{m}$. Then
\[
D_{m}=\operatorname*{Res}_{u=0}\frac{G(u)\dd u}{u^{m+1}} =\frac1\beta\operatorname*{Res}_{T=0}
\frac{(1+T^{2})^{m+1}}{T^{m+1}} \,\dd\!\left(\frac{\e^{\beta T}}{1-T^{2}}\right),
\]
using $u^{-m-1}=T^{-m-1}(1+T^{2})^{m+1}$. The residue of an exact differential vanishes, so
integrating by parts,
\[
D_{m}=-\frac1\beta\operatorname*{Res}_{T=0} \frac{\e^{\beta T}}{1-T^{2}}
\,\dd\!\left(T^{-m-1}(1+T^{2})^{m+1}\right) =\frac{m+1}{\beta}\operatorname*{Res}_{T=0} \e^{\beta
T}(1+T^{2})^{m}\,\frac{\dd T}{T^{m+2}},
\]
where we used $\dd(T^{-m-1}(1+T^{2})^{m+1})
=(1+T^{2})^{m}T^{-m-2}\bigl(-(m+1)(1+T^{2})+2(m+1)T^{2}\bigr)\dd T
=-(m+1)(1+T^{2})^{m}(1-T^{2})T^{-m-2}\dd T$, and the factor $1-T^{2}$ cancels. Hence
\[
D_{m}=\frac{m+1}{\beta}\Coef{T^{m+1}}\e^{\beta T}(1+T^{2})^{m}
=\frac{m+1}{\beta}\sum_{j\ge0}\binom mj\frac{\beta^{m+1-2j}}{(m+1-2j)!},
\]
the sum running over $0\le j\le\lfloor(m+1)/2\rfloor$. Finally $(m+1)\binom
mj=(m+1-j)\binom{m+1}{j}$ and division by $(-4\sqrt6)^{m}$ gives \cref{p3:eq:c-finite}.
\end{proof}

\begin{remark}[Why a finite sum exists]\label{p3:rem:why-finite}
A square-root shift would normally give an infinite convolution at each order. The Catalan
coordinate $T$ absorbs the square root, and the derivative already present in Rademacher's formula
turns the integrand into the exact differential \cref{p3:eq:G-differential}. After the
change-of-variable Jacobian cancels the denominator $1-T^{2}$, only the polynomial $(1+T^{2})^{m}$
survives.
\end{remark}

\begin{remark}[Provenance and the scope of the generalisation]
\label{p3:rem:osullivan}
For $(k,\sigma)=(1,+1)$, \cref{p3:eq:c-finite} at $\beta=\pi/6$ is O'Sullivan's coefficient
formula, also used by Banerjee, Paule, Radu and Schneider in their error analysis. S1 proves only
that case, by the same Catalan substitution. S2 and S3 both observe that the proof never uses
$k=1$ or $\sigma=+1$ --- the sector enters only through $\beta$ --- and state the general formula;
S3 gives a second, generating-function proof through the identity
$\sum_{j\ge0}\binom{m+2j}{j}q^{j}=C(q)^{m}/\sqrt{1-4q}$ for the Catalan series $C$. The residue
proof above is S2's, written out. We record the generalisation as the principal coefficient
statement of the chapter, since it covers every exponentially small sector at no extra cost.
\end{remark}

The first coefficients, as functions of $\beta$, are
\begin{align}
 c_{0}&=1, &
 c_{1}&=-\frac{\sqrt6\,(\beta^{2}+2)}{24\beta}, &
 c_{2}&=\frac{\beta^{2}+12}{192},\nonumber\\
 c_{3}&=-\frac{\sqrt6\,(\beta^{4}+36\beta^{2}+72)}{13824\,\beta}, &
 c_{4}&=\frac{\beta^{4}+80\beta^{2}+720}{221184}, &
 c_{5}&=-\frac{\sqrt6\,(\beta^{6}+150\beta^{4}+3600\beta^{2}+7200)}
 {26542080\,\beta},\label{p3:eq:c-first}
\end{align}
\begin{equation}\label{p3:eq:c-six}
 c_{6}=\frac{\beta^{6}+252\beta^{4}+12600\beta^{2}+100800}{637009920},
 \qquad
 c_{7}=-\frac{\sqrt6\bigl(\beta^{8}+392\beta^{6}+35280\beta^{4}
 +705600\beta^{2}+1411200\bigr)}{107017666560\,\beta}.
\end{equation}
These were recomputed here by expanding \cref{p3:eq:Omega} symbolically to order $t^{9}$ and
comparing with \cref{p3:eq:c-finite} coefficient by coefficient in exact arithmetic; the two agree
identically. The parity pattern visible in \cref{p3:eq:c-first} --- $c_{m}$ is $\sqrt6$ times an
odd Laurent polynomial in $\beta$ for odd $m$, and a polynomial in $\beta^{2}$ for even $m$ --- is
immediate from \cref{p3:eq:c-finite}, since $m-2j\equiv m$ $(\mathrm{mod}\ 2)$ and the prefactor
contributes $\sqrt6^{\,m}$.

\subsection{The Bell-polynomial route}

\Cref{p3:eq:c-finite} is the fastest generator. For symbolic systems, or when the sector
coefficients are wanted as a byproduct of a general composition engine, the following route uses
only the primitives of \cref{p0:def:obell,p0:def:ebell}.

\begin{proposition}[Logarithmic coefficients of a sector]\label{p3:prop:ell}
Write $\log\Omega_{k,\sigma}(t)=\sum_{j\ge1}\ell_{j}t^{j}$ (suppressing the sector labels) and put
$a=1/24$. Then, for $j\ge1$,
\begin{equation}\label{p3:eq:ell}
 \ell_{j}
 =\mathbf 1_{2\mid j}\,\frac{a^{j/2}}{j/2}
 +\mathbf 1_{2\nmid j}\,\frac{\sigma\lambda}{k}
   \binom{1/2}{(j+1)/2}(-a)^{(j+1)/2}
 -\sum_{\substack{1\le q\le j\\ 2\mid (j-q)}}\frac1q
   \left(\frac{\sigma k}{\lambda}\right)^{\!q}
   \frac{(q/2)_{(j-q)/2}}{\bigl((j-q)/2\bigr)!}\;a^{(j-q)/2},
\end{equation}
where $(x)_{r}=x(x+1)\cdots(x+r-1)$. Consequently the sector coefficients are the complete
exponential Bell polynomials
\begin{equation}\label{p3:eq:c-bell}
 c_{m}=\frac1{m!}\,\EB_{m}\bigl(1!\,\ell_{1},2!\,\ell_{2},\dots,m!\,\ell_{m}\bigr),
\end{equation}
equivalently the triangular recurrence
\begin{equation}\label{p3:eq:c-recurrence}
 c_{0}=1,\qquad c_{m}=\frac1m\sum_{j=1}^{m}j\,\ell_{j}\,c_{m-j}\quad(m\ge1).
\end{equation}
\end{proposition}

\begin{proof}
Taking logarithms in \cref{p3:eq:Omega},
\[
\log\Omega=-\log(1-at^{2}) +\frac{\sigma\lambda}{kt}\bigl(\sqrt{1-at^{2}}-1\bigr)
+\log\!\left(1-\frac{\sigma k}{\lambda}t(1-at^{2})^{-1/2}\right).
\]
The first term contributes $a^{r}/r$ at degree $2r$. For the second,
$\sqrt{1-at^{2}}-1=\sum_{r\ge1}\binom{1/2}{r}(-a)^{r}t^{2r}$, so dividing by $t$ contributes
$\frac{\sigma\lambda}{k}\binom{1/2}{r}(-a)^{r}$ at degree $2r-1$. For the third, expand
$\log(1-w)=-\sum_{q\ge1}w^{q}/q$ with $w=\frac{\sigma k}{\lambda}t(1-at^{2})^{-1/2}$ and use
$(1-at^{2})^{-q/2}=\sum_{r\ge0}\frac{(q/2)_{r}}{r!}a^{r}t^{2r}$; the term of degree $j=q+2r$ is
the displayed sum. \Cref{p3:eq:c-bell} is the exponential formula of \cref{p0:def:ebell}, and
\cref{p3:eq:c-recurrence} follows by comparing coefficients in $\Omega'=(\log\Omega)'\,\Omega$.
\end{proof}

For orientation, at general $(k,\sigma)$ and $a=1/24$,
\begin{align}
 \ell_{1}&=-\frac{\sigma\lambda a}{2k}-\frac{\sigma k}{\lambda}, &
 \ell_{2}&=a-\frac{k^{2}}{2\lambda^{2}},\nonumber\\
 \ell_{3}&=-\frac{\sigma\lambda a^{2}}{8k}-\frac{\sigma k a}{2\lambda}
 -\frac{\sigma k^{3}}{3\lambda^{3}}, &
 \ell_{4}&=\frac{a^{2}}{2}-\frac{k^{2}a}{2\lambda^{2}}
 -\frac{k^{4}}{4\lambda^{4}} .\label{p3:eq:ell-first}
\end{align}

\subsection{The dominant Poincar\'e expansion, and what it does not say}

\begin{definition}\label{p3:def:omega}
$\omega_{m}=c^{(1,+1)}_{m}$, i.e.\ \cref{p3:eq:c-finite} at $\beta=\pi/6$.
\end{definition}

\begin{corollary}\label{p3:cor:poincare}
For each fixed $M\ge1$,
\begin{equation}\label{p3:eq:poincare}
 p(n)=\frac{\e^{\lambda\sqrt n}}{4\sqrt3\,n}
 \left(\sum_{m=0}^{M-1}\omega_{m}n^{-m/2}+\bigO_{M}(n^{-M/2})\right),
\end{equation}
with
\begin{equation}\label{p3:eq:omega-first}
 \omega_{0}=1,\quad
 \omega_{1}=-\frac{\sqrt6\,(\pi^{2}+72)}{144\pi},\quad
 \omega_{2}=\frac{\pi^{2}+432}{6912},\quad
 \omega_{3}=-\frac{\sqrt6\,(\pi^{4}+1296\pi^{2}+93312)}{2985984\,\pi},
\end{equation}
\begin{equation}\label{p3:eq:omega-four}
 \omega_{4}=\frac{\pi^{4}+2880\pi^{2}+933120}{286654464},
\end{equation}
numerically $\omega_{1}=-0.443287976873582391\ldots$, $\omega_{2}=0.063927894155250197\ldots$,
$\omega_{3}=-0.027730933907464250\ldots$, $\omega_{4}=0.003354707463289919\ldots$.
\end{corollary}

\begin{proof}
By \cref{p3:thm:sector-gf} with $(k,\sigma)=(1,+1)$, the displayed sum plus a convergent tail is
exactly $P_{1,+}(n)$ divided by $\e^{\lambda\sqrt n}/(4\sqrt3\,n)$; the tail after $M$ terms is
$\bigO_{M}(n^{-M/2})$ because $\Omega_{1,+}$ is holomorphic at $0$. The difference
$p(n)-P_{1,+}(n)$ is $\bigO(\mu^{3/2}\e^{-\mu/2})$ relative to $P_{1,+}$ by
\cref{p3:thm:tail-sharp} with $K=1$, hence beyond every algebraic order. The closed forms follow
from \cref{p3:eq:c-finite} at $\beta=\pi/6$; they were recomputed here in exact arithmetic and
agree with the tables of S1 and S2 (which print the same numbers in two different but equal
groupings).
\end{proof}

\begin{repairbox}
\textbf{Convergent sector versus asymptotic total.} This is the trap of the subject and all three
sources are, at some point, loose about it. The three statements below are \emph{different}:
\begin{enumerate}
\item $\sum_{m\ge0}\omega_{m}n^{-m/2}$ \emph{converges} for every $n>1/24$
(\cref{p3:thm:sector-gf});
\item its sum is $P_{1,+}(n)\cdot 4\sqrt3\,n\,\e^{-\lambda\sqrt n}$, \emph{not}
$p(n)\cdot4\sqrt3\,n\,\e^{-\lambda\sqrt n}$;
\item \cref{p3:eq:poincare} is an \emph{asymptotic} statement about $p(n)$,
whose error at every fixed order is a sum of two utterly different things: a convergent algebraic
tail of size $\bigO(n^{-M/2})$ and an exponentially small Rademacher remainder of size
$\bigO(\mu^{3/2}\e^{-\mu/2})$ relative to the dominant sector.
\end{enumerate}
S1's ``Dominant Poincar\'e expansion'' subsection writes $p(n)\sim\cdots$ without separating (ii)
from (iii) and, in its Theorem on the sector generating function, weakens (i) to ``sufficiently
small $|t|$''. S3's remark after its dominant Poincar\'e display states the correct trichotomy in
one sentence. S2 states it as a standalone remark and is the presentation adopted here.
\textbf{The practical consequence:} summing the $\omega_{m}$ series to convergence does \emph{not}
compute $p(n)$ to arbitrary accuracy; it computes $P_{1,+}(n)$, and the residual error stalls at
the $\e^{-\mu/2}$ floor.
\end{repairbox}

\begin{remark}[Least-term truncation degenerates here]\label{p3:rem:no-least-term}
\Cref{p0:thm:optimal-truncation} determines the optimal truncation point of a \emph{factorially
divergent} algebraic series and the size of the resulting beyond-all-orders floor. Its hypothesis
fails for $p(n)$: by \cref{p3:thm:sector-gf} the algebraic series of each sector has a finite
positive radius of convergence, so there is no least term and no truncation-induced floor. The
floor that does exist --- the $\e^{-\mu/2}$ Rademacher gap --- is not produced by resummation
ambiguity but is supplied explicitly by the next sector. The partition function is thus the
degenerate, best possible case of \cref{p0:thm:optimal-truncation}: the exponentially small
sectors are known exactly and independently of the perturbative data, and the optimal algebraic
truncation is simply ``as many terms as one wishes''. This is the structural reason the numerical
accuracy in \cref{p3:tab:forward} improves so violently with $K$ but not at all with extra
algebraic orders past the $\e^{-\mu/2}$ level.
\end{remark}

\section{The multiplicative reciprocal \texorpdfstring{$1/p(n)$}{1/p(n)}}
\label{p3:sec:reciprocal}

The phrase ``inverse of the partition numbers'' is ambiguous. The reciprocal is the easy reading
and is disposed of first; it needs no reversion at all.

\begin{proposition}[Exact reciprocal transseries]\label{p3:prop:reciprocal}
Let $\mu>1$ be large enough that $|\mathcal{R}(\mu;n)|<1$, with $\mathcal{R}$ as in
\cref{p3:eq:R-exact}. Then
\begin{equation}\label{p3:eq:reciprocal}
 \frac1{p(n)}=\frac{\mu^{2}}{\kappa}\,\frac{\e^{-\mu}}{1-\mu^{-1}}
 \sum_{r\ge0}(-1)^{r}\mathcal{R}(\mu;n)^{r},
\end{equation}
the series converging. As a formal identity in \cref{p3:def:rademacher-transseries} it holds
without the inequality. If $\mathcal{R}$ is written at finite level as a linear form $\sum_{j\in
J_{K}}r_{j}(\mu;n)\,q_{j}$ in independent exponential variables $q_{j}=\e^{-\alpha_{j}\mu}$, then
for every multi-index $\boldsymbol\alpha=(\alpha_{j})_{j\in J_{K}}\ne\boldsymbol0$,
\begin{equation}\label{p3:eq:reciprocal-multi}
 \Coef{\boldsymbol q^{\boldsymbol\alpha}}\bigl(1+\mathcal{R}\bigr)^{-1}
 =(-1)^{|\boldsymbol\alpha|}\frac{|\boldsymbol\alpha|!}{\boldsymbol\alpha!}
 \prod_{j\in J_{K}}r_{j}(\mu;n)^{\alpha_{j}} .
\end{equation}
\end{proposition}

\begin{proof}
Invert \cref{p3:eq:relative} and expand $(1+\mathcal{R})^{-1}$ geometrically;
$1/F=\mu^{2}\e^{-\mu}/(\kappa(1-\mu^{-1}))$. For \cref{p3:eq:reciprocal-multi}, only the term
$(-1)^{|\boldsymbol\alpha|} \mathcal{R}^{|\boldsymbol\alpha|}$ of the geometric series can produce
$\boldsymbol q^{\boldsymbol\alpha}$, and the multinomial theorem supplies the coefficient.
\end{proof}

\begin{corollary}[Reciprocal Poincar\'e coefficients]\label{p3:cor:reciprocal-poincare}
Define $\eta_{m}$ by $\bigl(\sum_{m\ge0}\omega_{m}t^{m}\bigr)^{-1} =\sum_{m\ge0}\eta_{m}t^{m}$,
i.e.
\begin{equation}\label{p3:eq:eta}
 \eta_{0}=1,\qquad \eta_{m}=-\sum_{j=1}^{m}\omega_{j}\eta_{m-j}
 \quad(m\ge1),
 \qquad\text{equivalently}\qquad
 \eta_{m}=\frac1{m!}\EB_{m}\bigl(-1!\ell_{1},\dots,-m!\ell_{m}\bigr)
\end{equation}
with $\ell_{j}=\ell_{j}^{(1,+1)}$ of \cref{p3:eq:ell}. Then for each fixed $M$,
\begin{equation}\label{p3:eq:reciprocal-poincare}
 \frac1{p(n)}=4\sqrt3\,n\,\e^{-\lambda\sqrt n}
 \left(\sum_{m=0}^{M-1}\eta_{m}n^{-m/2}+\bigO_{M}(n^{-M/2})\right),
\end{equation}
and $\sum_{m\ge0}\eta_{m}t^{m}$ converges on $|t|<\sqrt{24}$ minus the zero set of $\Omega_{1,+}$,
to $1/\Omega_{1,+}(t)$.
\end{corollary}

\begin{proof}
The recurrence is Cauchy inversion of a power series with unit constant term; the Bell form is the
exponential formula applied to $\Omega_{1,+}^{-1}=\exp(-\sum\ell_{j}t^{j})$. Convergence is
inherited from \cref{p3:thm:sector-gf} and $\Omega_{1,+}(0)=1$. The asymptotic statement follows
from \cref{p3:cor:poincare} exactly as there; the same caveats of the repair box after
\cref{p3:cor:poincare} apply verbatim, with ``$P_{1,+}$'' replaced by ``$1/P_{1,+}$''.
\end{proof}

\begin{remark}
The reciprocal therefore requires no new analytic input and no series reversion. It is recorded
because two of the three sources devote a section to it, and because it is the natural place to
notice that the multinomial formula \cref{p3:eq:reciprocal-multi} is the
$|\boldsymbol\alpha|$-linear shadow of the far harder reversion formulas of
\cref{p3:sec:full-inverse}.
\end{remark}

\section{What ``the inverse'' of a discrete sequence means here}
\label{p3:sec:meaning}

\begin{proposition}[Strict monotonicity]\label{p3:prop:strict-increase}
$p(n)<p(n+1)$ for every $n\ge1$.
\end{proposition}

\begin{proof}
Adjoining a part $1$ injects the partitions of $n$ into those of $n+1$, and the one-part partition
$(n+1)$ is not in the image.
\end{proof}

By \cref{p0:def:three-inverses} there are three distinct objects, and this chapter needs all
three.

\begin{enumerate}
\item \textbf{The integer staircase} (the generalised inverse)
      \begin{equation}\label{p3:eq:staircase-def}
       p^{\leftarrow}(y)=\min\{n\ge1:p(n)\ge y\},\qquad y\ge1 .
      \end{equation}
Well defined by \cref{p3:prop:strict-increase}. It is integer valued with unit jumps, so no smooth
expansion can represent it to $\littleo(1)$.
\item \textbf{The functional inverse of a chosen interpolation.}  Any strictly
increasing smooth $\mathcal{P}$ with $\mathcal{P}(n)=p(n)$ has an inverse $\mathcal{N}$, and
$p^{\leftarrow}(y)=\lceil\mathcal{N}(y)\rceil$ (\cref{p3:thm:staircase}). \cref{p3:def:Atilde}
supplies the canonical choice used here.
\item \textbf{The inverse along the range.}  If one knows in advance that
$y=p(n)$ for an integer $n$, the target is that integer. A smooth formula that is accurate to
$\littleo(1/2)$ recovers it by rounding (\cref{p3:thm:rounding}), and the accuracy available is
exponentially small, not merely $\littleo(1)$.
\end{enumerate}

\begin{definition}[The Rademacher interpolation]\label{p3:def:Pcal}
For real $\nu>1/24$ put $\mu(\nu)=\tfrac\pi6\sqrt{24\nu-1}$ and
\begin{equation}\label{p3:eq:Pcal}
 \mathcal{P}(\nu)=\frac{2\kappa}{\mu(\nu)^{2}}\sum_{k\ge1}
 \frac{\widetilde A_{k}(\nu)}{\sqrt k}\,g\!\left(\frac{\mu(\nu)}{k}\right),
\end{equation}
with $g$ as in \cref{p3:eq:g-kernel} and $\widetilde A_{k}$ as in \cref{p3:def:Atilde}.
\end{definition}

\begin{proposition}[Interpolation, regularity, eventual monotonicity]
\label{p3:prop:Pcal}
The series \cref{p3:eq:Pcal} and all its derivatives converge locally uniformly on
$(1/24,\infty)$; $\mathcal{P}$ is real analytic there and $\mathcal{P}(n)=p(n)$ for every integer
$n\ge1$. Moreover there is $\nu_{0}$ with $\mathcal{P}'(\nu)>0$ for $\nu\ge\nu_{0}$, so
$\mathcal{P}$ has a real inverse $\mathcal{N}$ on $[\mathcal{P}(\nu_{0}),\infty)$.
\end{proposition}

\begin{proof}
Fix a compact $I\subset(1/24,\infty)$ and a complex neighbourhood of it on which $\mu$ is
holomorphic and bounded. By \cref{p3:lem:Atilde} every $\nu$-derivative of $\widetilde A_{k}$ is
$\bigO_{r}(k)$ uniformly in $k$ there, while $g(\mu/k)=\bigO(k^{-2})$ uniformly with all
parameter derivatives (from \cref{p3:eq:g-kernel} and $\mu$ bounded). Hence the $k$th summand of
\cref{p3:eq:Pcal} and each of its derivatives is $\bigO(k\cdot k^{-1/2}\cdot
k^{-2})=\bigO(k^{-3/2})$, and the Weierstrass $M$-test gives locally uniform convergence and
legitimises termwise differentiation. Equality at integers is \cref{p3:lem:Atilde} combined with
\cref{p3:thm:exact-sectors}.

Monotonicity: write $\mathcal{P}(\nu)=F(\mu(\nu))(1+\mathcal{R}(\mu(\nu);\nu))$ as in
\cref{p3:prop:relative} with $\widetilde A_{k}$ in place of $A_{k}$. We have
$\dd\mu/\dd\nu=\lambda^{2}/(2\mu)>0$ and $\dd\log F/\dd\mu=1-3/\mu+1/(\mu-1)\to1$. Differentiating
$\mathcal{R}$ term by term and applying \cref{p3:thm:tail} with $K=1$ together with
\cref{p3:lem:Atilde} shows $\mathcal{R}+\partial_{\mu}\mathcal{R}=\bigO(\mu^{3}\e^{-\mu/2})$;
note that differentiating $\widetilde A_{k}(\nu(\mu))$ costs a factor $\bigO(k\mu)$, absorbed
by replacing $k^{-3/2}$ with $k^{-1/2}$ in the tail sum and hence by one more power of $\mu$. Thus
$\mathcal{P}'(\nu)=F'(\mu)\mu'(\nu)\bigl(1+\bigO(\mu^{3}\e^{-\mu/2})\bigr)$, which is
positive for large $\nu$.
\end{proof}

\begin{theorem}[Staircase from the smooth inverse]\label{p3:thm:staircase}
For all sufficiently large real $y$,
\begin{equation}\label{p3:eq:staircase}
 p^{\leftarrow}(y)=\bigl\lceil\mathcal{N}(y)\bigr\rceil,
\end{equation}
and $\mathcal{N}(p(n))=n$ exactly for every integer $n\ge\nu_{0}$.
\end{theorem}

\begin{proof}
This is the specialisation of \cref{p0:thm:staircase} to the unit-spaced lattice $\Int$, whose
separation hypothesis is \cref{p3:prop:strict-increase}. Concretely: let $m=p^{\leftarrow}(y)$, so
$p(m-1)<y\le p(m)$. Since $\mathcal{P}$ is increasing on $[\nu_{0},\infty)$ and
$\mathcal{P}(j)=p(j)$, we get $m-1<\mathcal{N}(y)\le m$, whence $\lceil\mathcal{N}(y)\rceil=m$.
\end{proof}

\begin{remark}[The unavoidable sawtooth]\label{p3:rem:sawtooth}
For general real $y$ the difference $p^{\leftarrow}(y)-\mathcal{N}(y)$ is the ceiling defect, of
size $\bigO(1)$, with Fourier expansion
\begin{equation}\label{p3:eq:ceiling-fourier}
 \lceil v\rceil-v=\frac12+\frac1\pi\sum_{r\ge1}\frac{\sin(2\pi rv)}{r}
 \qquad (v\notin\Int),
\end{equation}
the series converging to the midpoint $1/2$ at integers where the true defect is $0$. Since the
first arithmetic correction to $\mathcal{N}$ is of size $y^{-1/2}$ (\cref{p3:thm:leading-chirp}),
\emph{exponential precision in $\mathcal{N}$ buys nothing for $p^{\leftarrow}$ at generic $y$}:
the order-one sawtooth is an intrinsic final sector of the discrete problem, not an error to be
removed. It matters only along the range $y=p(n)$, where the defect vanishes identically. This
warning is S3's and is the sharpest of the three.
\end{remark}

\section{The Lambert core and the algebraic inverse}\label{p3:sec:lambert}

Throughout this section and the next, $y$ denotes a large ordinate and
\begin{equation}\label{p3:eq:R-ell}
 R=\log\frac y\kappa,\qquad \ell=\log R .
\end{equation}

\subsection{The carrier}

\begin{proposition}[Lambert carrier; specialisation of
\cref{p0:thm:lambert-core}]\label{p3:prop:lambert-core} The large solution $X=X(y)$ of
$\kappa\e^{X}X^{-2}=y$, equivalently of
\begin{equation}\label{p3:eq:carrier-phase}
 X-2\log X=R,
\end{equation}
is
\begin{equation}\label{p3:eq:carrier}
 X(y)=-2\,\Wlam_{-1}\!\left(-\frac12\sqrt{\frac\kappa y}\right)
     =-2\,\Wlam_{-1}\!\left(-\frac{\pi}{\sqrt{24\sqrt3\,y}}\right).
\end{equation}
\end{proposition}

\begin{proof}
This is \cref{p0:thm:lambert-core} at the dictionary values $a=1$, $b=-2$, $C=\kappa$ of
\cref{p3:tab:dictionary}: $X=(b/a)\Wlam_{k}((a/b)\e^{R/b}) =-2\Wlam_{k}(-\tfrac12\e^{-R/2})$ and
$\e^{-R/2}=\sqrt{\kappa/y}$. The branch rule of \cref{p0:thm:lambert-core} selects $k=-1$ because
$X\to+\infty$ forces $-X/2\le-1$. The second form uses
$\tfrac12\sqrt\kappa=\tfrac12\pi/\sqrt{6\sqrt3} =\pi/\sqrt{24\sqrt3}$.
\end{proof}

\subsection{Flattening into logarithms}

\begin{theorem}[Carrier coefficient polynomials]\label{p3:thm:carrier-polys}
As $y\to\infty$,
\begin{equation}\label{p3:eq:carrier-flat}
 X\sim R+2\ell+\sum_{j\ge1}\frac{\Pi_{j}(\ell)}{R^{j}},
\end{equation}
in the sense of \cref{p0:thm:flattening}, where
\begin{equation}\label{p3:eq:Pi-lagrange}
 \Pi_{j}(\ell)=\frac{2^{j+1}}{j+1}\,\Coef{v^{j}}\bigl(\ell+\log(1+v)\bigr)^{j+1}
 =2^{j+1}\sum_{q=1}^{j}\frac{s(j,q)}{(j+1-q)!}\,\ell^{\,j+1-q},
\end{equation}
with $s(j,q)$ the signed Stirling numbers of the first kind,
$v(v-1)\cdots(v-j+1)=\sum_{q}s(j,q)v^{q}$. For each fixed $J$,
\begin{equation}\label{p3:eq:carrier-remainder}
 X=R+2\ell+\sum_{j=1}^{J}\frac{\Pi_{j}(\ell)}{R^{j}}
 +\bigO\!\left(\frac{\ell^{\,J+1}}{R^{J+1}}\right).
\end{equation}
Moreover $\deg\Pi_{j}=j$ with leading coefficient $(-1)^{j-1}2^{j+1}/j$, and $\Pi_{j}(0)=0$.
\end{theorem}

\begin{proof}
Write $X=R(1+v)$ in \cref{p3:eq:carrier-phase}: $R+Rv-2\log R-2\log(1+v)=R$, i.e.
\begin{equation}\label{p3:eq:carrier-fixed-point}
 v=z\bigl(\ell+\log(1+v)\bigr),\qquad z=\frac2R,
\end{equation}
a Lagrange fixed point $v=z\phi(v)$ with $\phi(v)=\ell+\log(1+v)$ and $\phi(0)=\ell\ne0$.
\Cref{p0:thm:LB} gives $v=\sum_{m\ge1}\frac{z^{m}}{m}\Coef{v^{m-1}}\phi(v)^{m}$; multiplying by
$R$ and putting $m=j+1$ (so $z^{m}R=2^{j+1}R^{-j}$) yields the first equality of
\cref{p3:eq:Pi-lagrange}. For the second, expand
$(\ell+\log(1+v))^{j+1}=\sum_{q}\binom{j+1}{q}\ell^{\,j+1-q}\log(1+v)^{q}$ and use
$\log(1+v)^{q}=q!\sum_{m\ge q}s(m,q)v^{m}/m!$; extracting $\Coef{v^{j}}$ keeps only $1\le q\le
j+1$, and the term $q=j+1$ contributes $s(j,j+1)=0$. The surviving coefficient of $\ell^{\,j+1-q}$
is $\frac{2^{j+1}}{j+1}\binom{j+1}{q}q!\,\frac{s(j,q)}{j!} =2^{j+1}\frac{s(j,q)}{(j+1-q)!}$, as
claimed. Since $q\ge1$, every monomial carries a positive power of $\ell$, so $\Pi_{j}(0)=0$; the
top power is $q=1$, with $s(j,1)=(-1)^{j-1}(j-1)!$, giving leading coefficient
$2^{j+1}(-1)^{j-1}(j-1)!/j!=(-1)^{j-1}2^{j+1}/j$. The remainder statement
\cref{p3:eq:carrier-remainder} is the general estimate of \cref{p0:thm:flattening} applied on the
scale $\ell/R\to0$; concretely, \cref{p3:eq:carrier-fixed-point} has an analytic solution
$v=v(z,\ell)$ near $z=0$ for $\ell$ in any fixed compact set, and $\ell/R\to0$ lets Taylor's
theorem with remainder be applied at order $J$.
\end{proof}

The first polynomials, recomputed here from the Lagrange form and independently from the Stirling
form (they agree identically), are
\begin{align}
 \Pi_{1}&=4\ell, &
 \Pi_{2}&=-4\ell(\ell-2), \nonumber\\
 \Pi_{3}&=\tfrac83\ell\,(2\ell^{2}-9\ell+6), &
 \Pi_{4}&=-\tfrac83\ell\,(3\ell^{3}-22\ell^{2}+36\ell-12),
 \label{p3:eq:Pi-first}\\
 \Pi_{5}&=\tfrac{16}{15}\ell\,(12\ell^{4}-125\ell^{3}+350\ell^{2}-300\ell+60),&
 \Pi_{6}&=-\tfrac{16}{15}\ell\,(20\ell^{5}-274\ell^{4}+1125\ell^{3}
 -1700\ell^{2}+900\ell-120),\nonumber
\end{align}
\begin{equation}\label{p3:eq:Pi-seven}
 \Pi_{7}=\tfrac{32}{105}\ell\,\bigl(120\ell^{6}-2058\ell^{5}+11368\ell^{4}
 -25725\ell^{3}+24500\ell^{2}-8820\ell+840\bigr),
\end{equation}
so that, explicitly,
\begin{equation}\label{p3:eq:carrier-explicit}
 X=R+2\ell+\frac{4\ell}{R}+\frac{-4\ell^{2}+8\ell}{R^{2}}
 +\frac{\tfrac{16}3\ell^{3}-24\ell^{2}+16\ell}{R^{3}}
 +\frac{-8\ell^{4}+\tfrac{176}3\ell^{3}-96\ell^{2}+32\ell}{R^{4}}+\cdots.
\end{equation}

\begin{remark}[Three printings of one polynomial family]\label{p3:rem:H-versus-R}
The sources print \cref{p3:eq:Pi-lagrange} in three different shapes, and it is worth recording
once that they are the same family, because a reader comparing them will otherwise suspect an
error.
\begin{itemize}
\item S1: $P_{j}(L)=2^{j+1}\sum_{q=1}^{j}\frac{s(j,q)}{(j+1-q)!}L^{j+1-q}$ ---
identical to \cref{p3:eq:Pi-lagrange}.
\item S2: $P_{m}(\ell)=\frac1{m!}\sum_{q=1}^{m}2^{q}(q-1)!\,s(m,q)
\binom{m}{q-1}(2\ell)^{m-q+1}$. The coefficient of $\ell^{\,m-q+1}$ is
$\frac{2^{q}(q-1)!}{m!}\binom{m}{q-1}2^{m-q+1} =\frac{2^{m+1}s(m,q)}{(m-q+1)!}$ after cancelling
$\binom{m}{q-1}=\frac{m!}{(q-1)!(m-q+1)!}$. Identical.
\item S3: $P_{r}(\ell)=\frac1{r!}\sum_{q=1}^{r}(q-1)!\,s(r,q)\binom{r}{q-1}
\ell^{\,r-q+1}$, appearing in the expansion as $2^{r+1}P_{r}(\ell)$. The same cancellation gives
$2^{r+1}P_{r}=\Pi_{r}$. Identical.
\end{itemize}
\textbf{But the variables differ, and this is a genuine disagreement.} S1 and S2 expand in
$R=\log(y/\kappa)$ with $\ell=\log R$; S3 expands in $H=\log(4y/\kappa)=R+2\log2$ with
$\ell_{\mathrm{S3}}=\log(H/2)$, which is the natural variable for the substitution $X=2w$, $w-\log
w=H/2$. The two series
\[
X=R+2\log R+\sum_{j}\frac{\Pi_{j}(\log R)}{R^{j}} \qquad\text{and}\qquad X=H+2\log\tfrac
H2+\sum_{j}\frac{\Pi_{j}(\log\tfrac H2)}{H^{j}}
\]
are \emph{different} asymptotic expansions of the same function, agreeing order by order only
after re-expanding $H=R+2\log2$. Both are correct; mixing them is not. This volume uses the S1/S2
variables $R,\ell$ throughout, because $R$ makes the phase equation exactly
\cref{p3:eq:carrier-phase} with no additive constant. We verified all of S3's displayed
coefficients (its $u$- and $N_{0}$-expansions in $H$, through $H^{-4}$ and $H^{-3}$ respectively)
by symbolic re-expansion; they are correct as printed in S3's own variables.
\end{remark}

\subsection{Restoring the dominant amplitude: the centred reversion}
\label{p3:sub:core}

The carrier $X$ solves $y=\kappa\e^{X}X^{-2}$; the \emph{dominant core} $\mu_{0}=\mu_{0}(y)$
solves the full dominant sector,
\begin{equation}\label{p3:eq:core-equation}
 y=F(\mu_{0})=\kappa\,\e^{\mu_{0}}\mu_{0}^{-2}\left(1-\frac1{\mu_{0}}\right),
 \qquad\text{i.e.}\qquad
 \Phi(\mu_{0})=R,
\end{equation}
where the \emph{dominant phase} is
\begin{equation}\label{p3:eq:Phi}
 \Phi(\mu)=\mu-2\log\mu+\log\!\left(1-\frac1\mu\right)
          =\mu-3\log\mu+\log(\mu-1),\qquad \mu>1 .
\end{equation}

\begin{lemma}[Derivatives of the dominant phase]\label{p3:lem:Phi-derivatives}
For $\mu>1$,
\begin{equation}\label{p3:eq:Phi-prime}
 \Phi'(\mu)=1-\frac3\mu+\frac1{\mu-1}
          =1-\frac2\mu+\frac1{\mu(\mu-1)}
          =\frac{\mu^{2}-3\mu+3}{\mu(\mu-1)}>0,
\end{equation}
and for $m\ge1$
\begin{equation}\label{p3:eq:Phi-m}
 \Phi^{(m)}(\mu)=\mathbf 1_{m=1}
 +(-1)^{m-1}(m-1)!\left(\frac1{(\mu-1)^{m}}-\frac3{\mu^{m}}\right).
\end{equation}
In particular $\Phi'(\mu)=1+\bigO(\mu^{-1})$, and $\mu^{2}-3\mu+3>0$ for all real $\mu$ since
its discriminant is $-3$. Writing $\varrho_{\pm}=\tfrac12(3\pm\ii\sqrt3)$ for the roots of
$\mu^{2}-3\mu+3$,
\begin{equation}\label{p3:eq:F-prime}
 F'(\mu)=\kappa\,\e^{\mu}\,\frac{\mu^{2}-3\mu+3}{\mu^{4}},
 \qquad
 \frac{\dd^{m}}{\dd\mu^{m}}\log F'(\mu)
 =\mathbf 1_{m=1}+(-1)^{m-1}(m-1)!
 \left(\frac1{(\mu-\varrho_{+})^{m}}+\frac1{(\mu-\varrho_{-})^{m}}
 -\frac4{\mu^{m}}\right).
\end{equation}
\end{lemma}

\begin{proof}
Differentiate \cref{p3:eq:Phi} in the form $\mu-3\log\mu+\log(\mu-1)$ and clear denominators for
the three displayed shapes of $\Phi'$; \cref{p3:eq:Phi-m} is $m-1$ further differentiations of
$-3\mu^{-1}+(\mu-1)^{-1}$. For \cref{p3:eq:F-prime}, $F=\kappa\e^{\mu}(\mu-1)\mu^{-3}$ gives
$F'=F\bigl(1+\tfrac1{\mu-1}-\tfrac3\mu\bigr)=F\,\Phi'(\mu)$, and
$F\Phi'=\kappa\e^{\mu}\frac{\mu-1}{\mu^{3}}\cdot\frac{\mu^{2}-3\mu+3}{\mu(\mu-1)}$. Then $\log
F'=\log\kappa+\mu+\log(\mu-\varrho_{+})+\log(\mu-\varrho_{-}) -4\log\mu$, and differentiate.
\end{proof}

\begin{theorem}[The dominant-core correction; specialisation of
\cref{p0:thm:lambert-centered}]\label{p3:thm:core-correction} Put $z=X^{-1}$. There is a unique
$U\in z\,\Rat[[z]]$ with
\begin{equation}\label{p3:eq:U-implicit}
 U=3\log(1+zU)-\log(1-z+zU),
\end{equation}
and $U$ is analytic in a neighbourhood of $z=0$. The dominant core is
\begin{equation}\label{p3:eq:core}
 \mu_{0}(y)=X(y)+U\bigl(X(y)^{-1}\bigr),\qquad U(z)=\sum_{m\ge1}c_{m}z^{m}.
\end{equation}
Moreover \cref{p3:eq:U-implicit} is exactly the fixed-point equation $U=\Psi(z,U)$ of
\cref{p0:thm:lambert-centered} at the dictionary values of \cref{p3:tab:dictionary}, so
\begin{equation}\label{p3:eq:c-lagrange}
 c_{m}=\sum_{q=1}^{m}\frac1q\,
 \Coef{z^{m}v^{q-1}}\Bigl(3\log(1+zv)-\log(1-z+zv)\Bigr)^{q}
\end{equation}
by \cref{p0:lem:lagrange-fp}, and equivalently
\begin{equation}\label{p3:eq:c-recurrence-core}
 c_{m}=\Coef{z^{m}}\Bigl\{3\log\bigl(1+zU_{<m}\bigr)
 -\log\bigl(1-z+zU_{<m}\bigr)\Bigr\},
 \qquad U_{<m}(z)=\sum_{j=1}^{m-1}c_{j}z^{j},
\end{equation}
a triangular recurrence over $\Rat$.
\end{theorem}

\begin{proof}
Put $\mu=X+U$ and divide \cref{p3:eq:core-equation} by the carrier identity
$y=\kappa\e^{X}X^{-2}$:
\[
\e^{U}(1+zU)^{-2}\left(1-\frac{z}{1+zU}\right)=1, \qquad z=\frac1X ,
\]
since $\mu=X(1+zU)$ and $1/\mu=z/(1+zU)$. As $1-\frac{z}{1+zU}=\frac{1-z+zU}{1+zU}$, taking
logarithms gives $U-3\log(1+zU)+\log(1-z+zU)=0$, i.e.\ \cref{p3:eq:U-implicit}. Write $G(z,U)$ for
the left side: $G(0,0)=0$ and $\partial_{U}G(0,0)=1$, so the analytic implicit function theorem
gives a unique analytic solution near $0$, necessarily with rational Taylor coefficients since $G$
has rational ones.

For the identification with \cref{p0:thm:lambert-centered}: that theorem's fixed-point map at
parameters $a,b,Q$ is $\Psi(t,u)=-\frac1a\bigl[b\log(1+tu)+Q\bigl(\tfrac{t}{1+tu}\bigr)\bigr]$. At
$a=1$, $b=-2$, $Q(s)=\log(1-s)$ this is
\[
\Psi(t,u)=2\log(1+tu)-\log\!\left(1-\frac{t}{1+tu}\right) =2\log(1+tu)-\log(1-t+tu)+\log(1+tu)
=3\log(1+tu)-\log(1-t+tu),
\]
which is the right side of \cref{p3:eq:U-implicit}. \Cref{p3:eq:c-lagrange} is then
\cref{p0:lem:lagrange-fp} verbatim. \Cref{p3:eq:c-recurrence-core} holds because in
\cref{p3:eq:U-implicit} every occurrence of $U$ on the right is multiplied by an explicit $z$, so
$\Coef{z^{m}}$ of the right side involves only $c_{1},\dots,c_{m-1}$.
\end{proof}

\begin{table}[htbp]
\centering
\caption{The centred-reversion coefficients $c_{m}$ of
\cref{p3:eq:core}, recomputed here in exact rational arithmetic from
\cref{p3:eq:c-recurrence-core} (truncated power-series arithmetic over
$\Rat$, $M=20$).  All eighteen agree with the tables printed by the three
sources, which call them $b_{m}$ (S1), $a_{m}$ (S2) and $d_{r}$ (S3).}
\label{p3:tab:cm}
\small
\begin{tabular}{@{}llll@{}}
\toprule
$c_{1}=1$ & $c_{2}=\tfrac52$ & $c_{3}=\tfrac{13}3$ & $c_{4}=\tfrac{53}{12}$\\
$c_{5}=-\tfrac{14}5$ & $c_{6}=-\tfrac{763}{30}$ & $c_{7}=-\tfrac{2179}{35}$
 & $c_{8}=-\tfrac{42289}{630}$\\
$c_{9}=\tfrac{132973}{1260}$ & $c_{10}=\tfrac{3449711}{5040}$
 & $c_{11}=\tfrac{29813143}{18480}$ & $c_{12}=\tfrac{496569589}{356400}$\\
$c_{13}=-\tfrac{4334124833}{926640}$ & $c_{14}=-\tfrac{509723063851}{21621600}$
 & $c_{15}=-\tfrac{298798016297}{5896800}$
 & $c_{16}=-\tfrac{933244850273}{33633600}$\\
$c_{17}=\tfrac{4993013022640963}{23156733600}$
 & $c_{18}=\tfrac{208843210189406957}{231567336000}$ & &\\
\bottomrule
\end{tabular}
\end{table}

Thus
\begin{equation}\label{p3:eq:core-explicit}
 \mu_{0}=X+\frac1X+\frac5{2X^{2}}+\frac{13}{3X^{3}}+\frac{53}{12X^{4}}
 -\frac{14}{5X^{5}}-\frac{763}{30X^{6}}-\frac{2179}{35X^{7}}+\cdots.
\end{equation}

\begin{remark}[Implementation check]\label{p3:rem:U-residual}
The cheapest correctness test for a table of $c_{m}$ is the residual identity: with
$U_{M}=\sum_{m\le M}c_{m}z^{m}$,
\begin{equation}\label{p3:eq:U-residual}
 U_{M}-3\log(1+zU_{M})+\log(1-z+zU_{M})=\bigO(z^{M+1}).
\end{equation}
We verified \cref{p3:eq:U-residual} at $M=12$ symbolically: the residual is $\bigO(z^{13})$
exactly. This test is far more reliable than comparing printed digits, and it is the one we used
to certify \cref{p3:tab:cm}.
\end{remark}

\subsection{The algebraic inverse index}

\begin{definition}\label{p3:def:N0}
$\displaystyle N_{0}(y)=\frac1{24}+\frac{3\,\mu_{0}(y)^{2}}{2\pi^{2}}$, the \emph{algebraic}
(phase-independent) inverse index; and $\displaystyle
N_{X}(y)=\frac1{24}+\frac{3X(y)^{2}}{2\pi^{2}}$, the carrier index.
\end{definition}

\begin{corollary}[All algebraic coefficients of the inverse index]
\label{p3:cor:N0-expansion}
With $c_{m}$ as in \cref{p3:tab:cm}, put
\begin{equation}\label{p3:eq:gamma}
 \gamma_{j}=3c_{j+1}+\frac32\sum_{r=1}^{j-1}c_{r}c_{j-r}\qquad(j\ge1).
\end{equation}
Then
\begin{equation}\label{p3:eq:N0-expansion}
 N_{0}(y)=\frac{3X^{2}}{2\pi^{2}}
 +\left(\frac1{24}+\frac3{\pi^{2}}\right)
 +\frac1{\pi^{2}}\sum_{j\ge1}\frac{\gamma_{j}}{X^{j}},
\end{equation}
with
\begin{equation}\label{p3:eq:gamma-first}
 \gamma_{1}=\tfrac{15}2,\quad \gamma_{2}=\tfrac{29}2,\quad
 \gamma_{3}=\tfrac{83}4,\quad \gamma_{4}=\tfrac{559}{40},\quad
 \gamma_{5}=-\tfrac{611}{20},\quad \gamma_{6}=-\tfrac{112459}{840},
\end{equation}
\begin{equation}\label{p3:eq:gamma-more}
 \gamma_{7}=-\tfrac{101329}{420},\quad
 \gamma_{8}=-\tfrac{228677}{3360},\quad
 \gamma_{9}=\tfrac{1709167}{1680},\quad
 \gamma_{10}=\tfrac{325098619}{92400}.
\end{equation}
\end{corollary}

\begin{proof}
Insert $\mu_{0}=X+U$ into \cref{p3:def:N0}. The cross term $\frac{3}{2\pi^{2}}\cdot
2XU=\frac{3}{\pi^{2}}\sum_{m\ge1}c_{m}X^{1-m}$ contributes $3c_{1}/\pi^{2}=3/\pi^{2}$ to the
constant and $3c_{j+1}/(\pi^{2}X^{j})$ at order $X^{-j}$; the square $\frac{3}{2\pi^{2}}U^{2}$
contributes $\frac{3}{2\pi^{2}}\bigl(\sum_{r=1}^{j-1}c_{r}c_{j-r}\bigr)X^{-j}$ and nothing at
order $X^{0}$ or above. The listed values were recomputed from \cref{p3:eq:gamma};
$\gamma_{1},\dots,\gamma_{8}$ agree with the tables of all three sources, and
$\gamma_{9},\gamma_{10}$ are new here.
\end{proof}

\begin{corollary}[The carrier bias]\label{p3:cor:carrier-bias}
$\displaystyle N_{0}(y)-N_{X}(y)=\frac3{\pi^{2}}+\bigO(X^{-1})$, and
\begin{equation}\label{p3:eq:carrier-bias}
 \frac3{\pi^{2}}=0.30396355092701331433\ldots
\end{equation}
\end{corollary}

\begin{proof}
Subtract $N_{X}=\tfrac1{24}+3X^{2}/(2\pi^{2})$ from \cref{p3:eq:N0-expansion}: the
$3X^{2}/(2\pi^{2})$ and $\tfrac1{24}$ cancel, leaving
$3/\pi^{2}+\pi^{-2}\sum_{j\ge1}\gamma_{j}X^{-j}$. The constant was evaluated here to $20$ digits.
\end{proof}

\begin{remark}[Phase locking: do not expand the arithmetic phase about the
carrier]\label{p3:rem:phase-lock} This is the single most consequential structural point of the
inverse problem, and only S2 states it explicitly.

The amplitudes $\widetilde A_{k}$ are \emph{periodic on the index scale}, with period $k$. The
core correction is small in the phase variable, $\mu_{0}-X=\bigO(X^{-1})$; but $\dd
n/\dd\mu=3\mu/\pi^{2}\asymp X$, so the induced change in the index is $\bigO(1)$ --- exactly
$3/\pi^{2}$ in the limit, by \cref{p3:cor:carrier-bias}. A shift of $0.304$ in the index is
\emph{not} a small perturbation of $\cos(\pi\nu)$. Consequently:
\begin{itemize}
\item Taylor-expanding $\widetilde A_{k}(N(\mu_{0}))$ around $N_{X}$ does not
begin with a small correction, and produces the wrong phase already at the leading exponentially
small order.
\item Every arithmetic amplitude must be evaluated at the \emph{full algebraic
core} $N_{0}(y)$, not at $N_{X}(y)$ and not at any truncation of $N_{0}$ coarser than
$\littleo(1)$. By \cref{p3:eq:N0-expansion}, truncating the $\gamma$-series after $M$ terms
leaves an index error $\bigO(X^{-M})$, so $M\ge1$ already suffices for $\littleo(1)$; $M=0$
(i.e.\ dropping the constant $3/\pi^{2}$) does not.
\item Only the \emph{subsequent} exponentially small displacement may be
Taylor-expanded; that is what \cref{p3:sec:full-inverse} does.
\end{itemize}
The numerical signature of the effect is the first column of \cref{p3:tab:inverse}, which tends to
$-3/\pi^{2}$ and not to $0$.
\end{remark}

\begin{remark}[Elementary-logarithm form]\label{p3:rem:N0-in-R}
Substituting \cref{p3:eq:carrier-flat} into \cref{p3:eq:N0-expansion} gives the pure logarithmic
block
\begin{equation}\label{p3:eq:N0-in-R}
 N_{0}(y)=\frac{3R^{2}}{2\pi^{2}}+\frac{6R\ell}{\pi^{2}}
 +\frac{6\ell^{2}+12\ell+3}{\pi^{2}}+\frac1{24}
 +\frac{3(8\ell^{2}+16\ell+5)}{2\pi^{2}R}
 +\bigO\!\left(\frac{\ell^{3}}{R^{2}}\right),
\end{equation}
with $R=\log(y/\kappa)$ and $\ell=\log R$. We recomputed this block symbolically; it agrees with
S2's display (which stops at the constant) and, after the change of variable $H=R+2\log2$, with
S3's. For \emph{numerical} work \cref{p3:eq:N0-in-R} is markedly inferior to evaluating
\cref{p3:eq:carrier} and \cref{p3:eq:core-explicit} directly, and it is recorded only to exhibit
the transseries shape: $N_{0}$ is a polynomial of degree $2$ in $R$ with coefficients polynomial
in $\ell$, plus a descending series in $R^{-1}$ with $\deg_{\ell}$ growing by one per order.
\end{remark}

\section{The full arithmetic inverse transseries}\label{p3:sec:full-inverse}

Everything so far is phase-independent: $X$, $\mu_{0}$, $N_{0}$ are functions of $y$ alone. The
remaining correction is exponentially small and arithmetically oscillatory, and the oscillation
must be frozen before ``analytic inverse'' has a meaning. The three sources do this in two
different ways, both correct.

\subsection{Two ways to freeze the phase}\label{p3:sub:two-freezings}

\begin{definition}[Residue-frozen finite level; S1]\label{p3:def:frozen}
For $K\ge1$ let $\Lambda_{K}=\operatorname{lcm}(1,2,\dots,K)$ and fix a residue
$r\in\Int/\Lambda_{K}\Int$. Since $A_{k}$ has period $k\mid\Lambda_{K}$ for $k\le K$, the numbers
$A_{k}(r)$ are well defined, and
\begin{equation}\label{p3:eq:PKr}
 P_{K,r}(\mu)=\frac{\kappa}{\mu^{2}}\sum_{k=1}^{K}\frac{A_{k}(r)}{\sqrt k}
 \left[\e^{\mu/k}\Bigl(1-\frac k\mu\Bigr)
      +\e^{-\mu/k}\Bigl(1+\frac k\mu\Bigr)\right]
\end{equation}
is a genuine analytic function of $\mu$ with constant coefficients. For large $\mu$ it is
increasing (the $k=1$ positive term dominates), and we write $\mu_{K,r}(y)$ for its large inverse
branch.
\end{definition}

\begin{definition}[Interpolated level; S2, S3]\label{p3:def:interpolated}
$\mathcal{P}_{K}(\mu)$ is \cref{p3:eq:PKr} with $A_{k}(r)$ replaced by $\widetilde A_{k}(n(\mu))$,
$n(\mu)=\tfrac1{24}+3\mu^{2}/(2\pi^{2})$; its large inverse branch is written $\mu_{K}(y)$.
\end{definition}

\begin{proposition}[The two conventions agree below action one]
\label{p3:prop:two-conventions}
Fix $K$, let $y$ be large, and let $r\equiv\lfloor N_{0}(y)\rceil \pmod{\Lambda_{K}}$, where
$\lfloor\cdot\rceil$ is nearest-integer rounding. Assume $|N_{0}(y)-r|\le\tfrac12$ modulo
$\Lambda_{K}$ has been resolved (which holds along the range $y=p(n)$ for large $n$ by
\cref{p3:thm:rounding}). Then every inverse transmonomial of total action strictly less than $1$
is the same for $\mu_{K,r}$ and $\mu_{K}$; the two first differ at total action $\ge1$.
\end{proposition}

\begin{proof}
Both inverses are constructed as $\mu_{0}+v$ with $v$ exponentially small, and in both the
coefficient of the transmonomial $\e^{-\alpha_{j}\mu_{0}}$ at first order is
$-r_{j}(\mu_{0})/\Phi'(\mu_{0})$ (\cref{p3:thm:first-order} below), where $r_{j}$ differs between
the two conventions only in the amplitude, $A_{k}(r)$ versus $\widetilde
A_{k}(n(\mu_{0}))=\widetilde A_{k}(N_{0}(y))$. Now $\widetilde A_{k}$ is entire and $k$-periodic
with $\widetilde A_{k}(r)=A_{k}(r)$, so $|\widetilde A_{k}(N_{0})-A_{k}(r)|\le \sup|\widetilde
A_{k}'|\cdot|N_{0}-r|$. Along the range, $|N_{0}(p(n))-n|=\bigO(\mu^{5/2}\e^{-\mu/2})$ by
\cref{p3:cor:core-inverse-error}, so the amplitude discrepancy is
$\bigO(k\mu^{5/2}\e^{-\mu/2})$ and its contribution to the sector carries the extra factor
$\e^{-\mu/2}$, i.e.\ additional action $\tfrac12$. Since the smallest primary action is
$\tfrac12$, the first affected total action is $\tfrac12+\tfrac12=1$. Higher-order terms of the
reversion have total action $\ge1$ already.
\end{proof}

\begin{remark}[Which convention to prefer]\label{p3:rem:which-convention}
The residue-frozen convention is intrinsic: it invents no interpolation, and at each finite level
it reduces the problem to an ordinary analytic reversion with \emph{constant} coefficients. Its
cost is that the inverse is defined only after the residue class is known, which
\cref{p3:thm:rounding} supplies. The interpolated convention defines a single function
$\mathcal{N}$ for all real $y$ and connects directly to the staircase (\cref{p3:thm:staircase});
its cost is the interpolation dependence of \cref{p3:rem:interp-nonunique} and the appearance of
the amplitude's own $\mu$-derivatives in every higher coefficient. This chapter develops the
reversion in a form that covers both: all formulas below are stated for a family of \emph{blocks}
$E_{j}(\mu)$, and the two conventions are two choices of block.
\end{remark}

\subsection{Blocks, multi-indices, and the reversion equation}
\label{p3:sub:blocks}

\begin{definition}[Sector blocks]\label{p3:def:blocks}
Fix $K\ge1$ and put
\begin{equation}\label{p3:eq:JK}
 J_{K}=\{(1,-1)\}\cup\{(k,\sigma):2\le k\le K,\ \sigma=\pm1\}.
\end{equation}
For $j=(k,\sigma)\in J_{K}$ set $\alpha_{j}=1-\sigma/k$ as in \cref{p3:eq:action}, and define the
\emph{block}
\begin{equation}\label{p3:eq:block}
 E_{j}(\mu)=\frac{\kappa}{\mu^{2}}\frac{a_{k}(\mu)}{\sqrt k}
 \left(1-\frac{\sigma k}{\mu}\right)\e^{\sigma\mu/k}
 =F(\mu)\,r_{j}(\mu)\,\e^{-\alpha_{j}\mu},
 \qquad
 r_{j}(\mu)=\frac{a_{k}(\mu)}{\sqrt k}\,\frac{1-\sigma k/\mu}{1-1/\mu},
\end{equation}
where $a_{k}(\mu)$ is either the constant $A_{k}(r)$ (residue-frozen convention,
\cref{p3:def:frozen}) or $\widetilde A_{k}(n(\mu))$ (interpolated convention,
\cref{p3:def:interpolated}); $a_{1}\equiv1$ in both. A block may also be the whole paired tail
$\mathcal{R}_{>K}\cdot F$ of \cref{p3:eq:R-exact}, or the zeta-condensation block of
\cref{p3:thm:zeta-condensation}, treated as one indivisible symbol.
\end{definition}

Introduce independent formal variables $\varepsilon_{j}$, $j\in J_{K}$, and write
$\boldsymbol\varepsilon^{\boldsymbol m}=\prod_{j}\varepsilon_{j}^{m_{j}}$, $|\boldsymbol
m|=\sum_{j}m_{j}$, $\boldsymbol m!=\prod_{j}m_{j}!$ for multi-indices $\boldsymbol
m\in\Nat^{J_{K}}$. The forward object at level $K$ is
\begin{equation}\label{p3:eq:forward-blocks}
 \mathcal{P}_{K}(\mu;\boldsymbol\varepsilon)
 =F(\mu)+\sum_{j\in J_{K}}\varepsilon_{j}E_{j}(\mu),
\end{equation}
and the inverse problem is to solve $\mathcal{P}_{K}(\mu;\boldsymbol\varepsilon)=y$ for $\mu$ near
$\mu_{0}=F^{-1}(y)$, as a formal series in $\boldsymbol\varepsilon$ which is finally specialised
at $\varepsilon_{j}=1$. Two equivalent engines are available; both are finite at every prescribed
multi-index.

\subsection{Engine I: additive Lagrange--B\"urmann}

\begin{theorem}[Additive reversion]\label{p3:thm:additive-LB}
Let $F$ and $E$ be analytic near $\mu_{0}$ with $F'(\mu_{0})\ne0$ and $F(\mu_{0})=y$. Write
$\mathcal{D}=\frac1{F'(\mu)}\frac{\dd}{\dd\mu}$ evaluated along $\mu=\mu_{0}(y)$, i.e.\
$\mathcal{D}=\dd/\dd y$ in the coordinate $y$. Then the solution of $F(\mu)+\varepsilon E(\mu)=y$
near $\mu_{0}$ is
\begin{equation}\label{p3:eq:additive-LB}
 \mu(y;\varepsilon)=\mu_{0}(y)
 +\sum_{M\ge1}\frac{(-\varepsilon)^{M}}{M!}\,
 \mathcal{D}^{M-1}\!\left(\frac{E(\mu_{0})^{M}}{F'(\mu_{0})}\right).
\end{equation}
Consequently, for the multivariate problem \cref{p3:eq:forward-blocks}, the coefficient of
$\boldsymbol\varepsilon^{\boldsymbol m}$ with $M=|\boldsymbol m|\ge1$ in $\mu-\mu_{0}$ is
\begin{equation}\label{p3:eq:multi-LB}
 V_{\boldsymbol m}(\mu_{0})=\frac{(-1)^{M}}{\boldsymbol m!}\,
 \mathcal{D}^{M-1}\!\left(\frac{E^{\boldsymbol m}(\mu_{0})}{F'(\mu_{0})}\right)
 =\frac{(-1)^{M}(M-1)!}{\boldsymbol m!}\,
 \operatorname*{Res}_{\mu=\mu_{0}}\frac{E^{\boldsymbol m}(\mu)}{(F(\mu)-F(\mu_{0}))^{M}},
\end{equation}
where $E^{\boldsymbol m}=\prod_{j}E_{j}^{m_{j}}$.
\end{theorem}

\begin{proof}
Both displays are \cref{p0:thm:perturbed-inversion} in the same letters:
\cref{p3:eq:additive-LB} is \eqref{p0:eq:perturbed-inversion} and
\cref{p3:eq:multi-LB} is \eqref{p0:eq:perturbed-multi}.

What that theorem supplies is convergence for $|\varepsilon|<\varepsilon_{0}$,
with $\varepsilon_{0}$ read off a small circle about $\mu_{0}$; here the
specialisation $\varepsilon_{j}=1$ is wanted, and for the partition blocks it
is legitimate.  With $E=\sum_{j}E_{j}+\mathcal{R}_{>K}F$ we have
$E/F=\bigO(\mu^{C}\e^{-\mu/2})$ together with the same bound for finitely many
derivatives, by \cref{p3:thm:tail} and \cref{p3:lem:Atilde}; so on a fixed
relative neighbourhood of a large $\mu_{0}$ the Rouch\'e comparison made in
the proof of \cref{p0:thm:perturbed-inversion} holds for all
$|\varepsilon|\le1+\eta$ with some $\eta>0$, and the Taylor series in
$\varepsilon$ converges at $\varepsilon=1$.
\end{proof}

\begin{remark}[$\mathcal{D}$ is iterated, not a plain derivative]
\label{p3:rem:D-iterated}
In \cref{p3:eq:additive-LB} the operator $\mathcal{D}^{M-1}$ is the $(M-1)$st power of
$\frac1{F'}\frac{\dd}{\dd\mu}$. Replacing it by $\frac1{(F')^{M-1}}\frac{\dd^{M-1}}{\dd\mu^{M-1}}$
after ``pulling out'' the powers of $F'$ is wrong from $M=3$ onwards, because $F'$ is not
constant. S3 flags this explicitly; it is the most common implementation bug in this formula.
\end{remark}

\begin{corollary}[Bell realisation]\label{p3:cor:multi-Bell}
Let $H_{\boldsymbol m}=E^{\boldsymbol m}/F'$ and $\Lambda^{(\boldsymbol
m)}_{s}=\frac{\dd^{s}}{\dd\mu^{s}}\log H_{\boldsymbol m}$, and let
$\mu_{0}^{(s)}=\dd^{s}\mu_{0}/\dd y^{s}$ be generated by
\begin{equation}\label{p3:eq:inverse-derivatives}
 \mu_{0}^{(1)}=\frac1{F'(\mu_{0})},\qquad
 \mu_{0}^{(s)}=-\frac1{F'(\mu_{0})}\sum_{q=2}^{s}F^{(q)}(\mu_{0})\,
 \EB_{s,q}\bigl(\mu_{0}^{(1)},\dots,\mu_{0}^{(s-q+1)}\bigr),\quad s\ge2 .
\end{equation}
Then, with $M=|\boldsymbol m|\ge1$,
\begin{equation}\label{p3:eq:multi-Bell}
 V_{\boldsymbol m}=\frac{(-1)^{M}H_{\boldsymbol m}(\mu_{0})}{\boldsymbol m!}
 \sum_{s=0}^{M-1}\EB_{M-1,s}\bigl(\mu_{0}^{(1)},\mu_{0}^{(2)},\dots\bigr)\,
 \EB_{s}\bigl(\Lambda^{(\boldsymbol m)}_{1},\dots,\Lambda^{(\boldsymbol m)}_{s}\bigr),
\end{equation}
with the conventions $\EB_{0,0}=\EB_{0}=1$.
\end{corollary}

\begin{proof}
Fa\`a di Bruno (\cref{p0:def:ebell}) gives $\frac{\dd^{M-1}}{\dd y^{M-1}}H_{\boldsymbol
m}(\mu_{0}(y)) =\sum_{s}H_{\boldsymbol m}^{(s)}(\mu_{0})\EB_{M-1,s}(\mu_{0}^{(1)},\dots)$, and
writing $H_{\boldsymbol m}=\exp(\log H_{\boldsymbol m})$ gives $H_{\boldsymbol
m}^{(s)}=H_{\boldsymbol m}\,\EB_{s}(\Lambda^{(\boldsymbol m)}_{1},\dots)$. Substitute in
\cref{p3:eq:multi-LB}. \Cref{p3:eq:inverse-derivatives} is obtained by differentiating
$F(\mu_{0}(y))=y$ $s$ times and solving the Fa\`a di Bruno identity for the term containing
$\mu_{0}^{(s)}$.
\end{proof}

The point of \cref{p3:eq:multi-Bell} is that every ingredient is elementary and closed-form for
the partition blocks. Indeed, splitting each aggregate block into its primitive Fourier
constituents
\begin{equation}\label{p3:eq:primitive-block}
 E_{k,h,\sigma,\epsilon}(\mu)
 =\frac{\kappa}{2\sqrt k}\,\frac{\mu-\sigma k}{\mu^{3}}
 \exp\!\left[\frac{\sigma\mu}{k}
 +\ii\epsilon\Bigl(\pi s(h,k)-\frac{2\pi h}{k}n(\mu)\Bigr)\right],
 \qquad h\in U_{k},\ \epsilon=\pm1,
\end{equation}
and using $n'(\mu)=3\mu/\pi^{2}$, $n''(\mu)=3/\pi^{2}$, $n^{(s)}=0$ for $s\ge3$, one gets the
closed logarithmic jets
\begin{equation}\label{p3:eq:jets}
 \frac{\dd^{s}}{\dd\mu^{s}}\log E_{k,h,\sigma,\epsilon}
 =(-1)^{s-1}(s-1)!\left(\frac1{(\mu-\sigma k)^{s}}-\frac3{\mu^{s}}\right)
 +\mathbf 1_{s=1}\frac\sigma k
 -\ii\epsilon\frac{2\pi h}{k}\,n^{(s)}(\mu),
\end{equation}
while $\frac{\dd^{s}}{\dd\mu^{s}}\log F'$ is \cref{p3:eq:F-prime}. Hence $\Lambda^{(\boldsymbol
m)}_{s}=\sum_{j}m_{j}\,(\text{\cref{p3:eq:jets}}) -(\text{\cref{p3:eq:F-prime}})$ with no hidden
differentiation anywhere. This is S3's presentation and is the most explicit of the three.

\subsection{Engine II: the triangular Newton--Hensel recurrence}

The reversion may equally be organised in the phase coordinate, which is cheaper when many
coefficients are wanted.

\begin{theorem}[Triangular recurrence]\label{p3:thm:triangular}
Put $\mu=\mu_{0}+v$ and
\begin{equation}\label{p3:eq:Q-def}
 Q(v;\boldsymbol\varepsilon)=\log\!\Bigl(1+\sum_{j\in J_{K}}\varepsilon_{j}B_{j}(v)\Bigr),
 \qquad B_{j}(v)=r_{j}(\mu_{0}+v)\,\e^{-\alpha_{j}v},
\end{equation}
the bookkeeping variable $\varepsilon_{j}$ standing for the actual exponential
$\e^{-\alpha_{j}\mu_{0}}$, which is substituted only at the very end; the residual $v$-dependence
of $\e^{-\alpha_{j}(\mu_{0}+v)}$ is the factor $\e^{-\alpha_{j}v}$ inside $B_{j}$. Keeping the
$\varepsilon_{j}$ independent until the last step is what makes resonant exponent sums
($\alpha_{j}+\alpha_{j'}=\alpha_{j''}$, which does occur here) come out right. Then the reversion
equation is
\begin{equation}\label{p3:eq:v-equation}
 \Phi'(\mu_{0})\,v+\sum_{s\ge2}\frac{\Phi^{(s)}(\mu_{0})}{s!}v^{s}
 +Q(v;\boldsymbol\varepsilon)=0,
\end{equation}
and its unique formal solution $v=\sum_{\boldsymbol m\ne\boldsymbol0}V_{\boldsymbol m}
\boldsymbol\varepsilon^{\boldsymbol m}$ obeys
\begin{equation}\label{p3:eq:triangular}
 V_{\boldsymbol m}=-\frac1{\Phi'(\mu_{0})}\,
 \Coef{\boldsymbol\varepsilon^{\boldsymbol m}}\left\{
 \sum_{s\ge2}\frac{\Phi^{(s)}(\mu_{0})}{s!}\,v_{<\boldsymbol m}^{\;s}
 +Q(v_{<\boldsymbol m};\boldsymbol\varepsilon)\right\},
\end{equation}
where $v_{<\boldsymbol m}$ collects the already-computed coefficients. Each right-hand side is a
finite expression in $\mu_{0}$, $(\mu_{0}-1)^{-1}$, the actions $\alpha_{j}$ and the amplitudes
$a_{k}$.
\end{theorem}

\begin{proof}
Taking logarithms in $\mathcal{P}_{K}(\mu;\boldsymbol\varepsilon)=y$ and subtracting
$\Phi(\mu_{0})=R$ gives $\Phi(\mu_{0}+v)-\Phi(\mu_{0})+Q(v;\boldsymbol\varepsilon)=0$; Taylor
expansion of $\Phi$ at $\mu_{0}$ is \cref{p3:eq:v-equation}, legitimate as a formal identity
because $v$ has no constant term in $\boldsymbol\varepsilon$. For triangularity: at a fixed
$\boldsymbol m$, the unknown $V_{\boldsymbol m}$ occurs linearly only in the term
$\Phi'(\mu_{0})v$; each $v^{s}$ with $s\ge2$ needs a decomposition of $\boldsymbol m$ into at
least two non-zero parts and so uses only strictly smaller multi-indices; and every monomial of
$Q$ carries at least one explicit $\varepsilon_{j}$, so an occurrence of $V_{\boldsymbol m}$
inside $Q$ would be multiplied into a strictly larger multi-index. Solving the coefficient
equation gives \cref{p3:eq:triangular}. Finiteness holds because a fixed multi-index has finitely
many additive decompositions.
\end{proof}

\begin{corollary}[Multinomial expansion of $Q$]\label{p3:cor:Q-multi}
For $\boldsymbol m\ne\boldsymbol0$,
\begin{equation}\label{p3:eq:Q-multi}
 \Coef{\boldsymbol\varepsilon^{\boldsymbol m}}Q(v;\boldsymbol\varepsilon)\Big|_{v\ \mathrm{fixed}}
 =(-1)^{|\boldsymbol m|+1}\frac{(|\boldsymbol m|-1)!}{\boldsymbol m!}\prod_{j}B_{j}(v)^{m_{j}} .
\end{equation}
\end{corollary}

\begin{proof}
Only $S^{|\boldsymbol m|}$ in $\log(1+S)=\sum_{d\ge1}(-1)^{d+1}S^{d}/d$ contributes, with
multinomial coefficient $|\boldsymbol m|!/\boldsymbol m!$.
\end{proof}

\begin{corollary}[Closed multi-index Lagrange form in the phase]
\label{p3:cor:phase-LB}
Let $G(v)=\bigl(\Phi(\mu_{0}+v)-\Phi(\mu_{0})\bigr)/v$ for $v\ne0$ and $G(0)=\Phi'(\mu_{0})$. Then
$v=-Q(v;\boldsymbol\varepsilon)/G(v)$, and by \cref{p0:thm:LB},
\begin{equation}\label{p3:eq:phase-LB}
 V_{\boldsymbol m}=\sum_{q=1}^{|\boldsymbol m|}\frac{(-1)^{q}}{q}\,\Coef{v^{q-1}}
 \frac1{G(v)^{q}}
 \sum_{\substack{\boldsymbol m_{1}+\cdots+\boldsymbol m_{q}=\boldsymbol m\\ \boldsymbol m_{i}\ne\boldsymbol 0}}
 \prod_{i=1}^{q}\Coef{\boldsymbol\varepsilon^{\boldsymbol m_{i}}}Q(v;\boldsymbol\varepsilon),
\end{equation}
both sums being finite.
\end{corollary}

\begin{proof}
$G$ is analytic near $0$ with $G(0)=\Phi'(\mu_{0})\ne0$, and \cref{p3:eq:v-equation} reads
$vG(v)+Q(v;\boldsymbol\varepsilon)=0$. Apply \cref{p0:thm:LB} to
$v=z\cdot\bigl(-Q(v;\boldsymbol\varepsilon)/G(v)\bigr)$ and set $z=1$; since every monomial of $Q$
has positive $\boldsymbol\varepsilon$-degree, only $q\le|\boldsymbol m|$ contributes. Expanding
$Q^{q}$ and using \cref{p3:eq:Q-multi} gives \cref{p3:eq:phase-LB}.
\end{proof}

\begin{remark}[Three engines, one answer]
\Cref{p3:thm:additive-LB} works in the ordinate $y$ and is S3's; \cref{p3:cor:phase-LB} works in
the phase $\mu$ and is S2's; \cref{p3:thm:triangular} is S1's. They compute the same coefficients:
all three solve the same analytic equation with the same normalisation $V_{\boldsymbol 0}=0$, and
the solution of a triangular system is unique. In practice \cref{p3:eq:triangular} is preferable
when many multi-indices are wanted (it reuses everything), \cref{p3:eq:multi-LB} when a single
high-order multi-index is wanted (it needs no lower ones), and \cref{p3:eq:multi-Bell} when the
answer is wanted symbolically in $\mu_{0}$.
\end{remark}

\subsection{First-order sectors and the leading parity chirp}
\label{p3:sub:first-order}

\begin{theorem}[The first correction from each sector]\label{p3:thm:first-order}
Let $\boldsymbol e_{j}$ be the unit multi-index at $j=(k,\sigma)\in J_{K}$. Then
\begin{equation}\label{p3:eq:first-V}
 V_{\boldsymbol e_{j}}(\mu_{0})=-\frac{r_{j}(\mu_{0})}{\Phi'(\mu_{0})}
 =-\frac{a_{k}(\mu_{0})}{\sqrt k}\,
 \frac{\mu_{0}(\mu_{0}-\sigma k)}{\mu_{0}^{2}-3\mu_{0}+3},
\end{equation}
so that the corresponding displacement of the inverse phase is $V_{\boldsymbol
e_{j}}(\mu_{0})\,\e^{-\alpha_{j}\mu_{0}}$, and the corresponding displacement of the inverse index
is
\begin{equation}\label{p3:eq:first-N}
 \delta N^{(1)}_{k,\sigma}(y)
 =-\frac{3\,a_{k}(\mu_{0})}{\pi^{2}\sqrt k}\,
 \frac{\mu_{0}^{2}(\mu_{0}-\sigma k)}{\mu_{0}^{2}-3\mu_{0}+3}\,
 \e^{-(1-\sigma/k)\mu_{0}} .
\end{equation}
\end{theorem}

\begin{proof}
Take $\boldsymbol m=\boldsymbol e_{j}$ in \cref{p3:eq:triangular}: the $s\ge2$ terms and all
higher $Q$-monomials vanish, leaving $V_{\boldsymbol
e_{j}}=-\Coef{\varepsilon_{j}}Q(0;\boldsymbol\varepsilon)/\Phi'(\mu_{0})
=-B_{j}(0)/\Phi'(\mu_{0})=-r_{j}(\mu_{0})/\Phi'(\mu_{0})$. Insert $r_{j}=\frac{a_{k}}{\sqrt
k}\frac{1-\sigma k/\mu}{1-1/\mu} =\frac{a_{k}}{\sqrt k}\frac{\mu-\sigma k}{\mu-1}$ and
\cref{p3:eq:Phi-prime}: the factors $\mu_{0}-1$ cancel, giving \cref{p3:eq:first-V}. Finally
$n(\mu)=\frac1{24}+\frac{3\mu^{2}}{2\pi^{2}}$ is exactly quadratic, so
\begin{equation}\label{p3:eq:index-from-phase}
 n(\mu_{0}+v)=N_{0}(y)+\frac{3\mu_{0}}{\pi^{2}}v+\frac3{2\pi^{2}}v^{2},
\end{equation}
and the linear term at first order gives \cref{p3:eq:first-N}.
\end{proof}

\begin{corollary}[Index coefficient of every transmonomial]\label{p3:cor:N-multi}
For every $\boldsymbol m\ne\boldsymbol0$, the coefficient of $\boldsymbol\varepsilon^{\boldsymbol
m}$ in the inverse index is
\begin{equation}\label{p3:eq:N-multi}
 \mathfrak N_{\boldsymbol m}(\mu_{0})=\frac{3\mu_{0}}{\pi^{2}}V_{\boldsymbol m}
 +\frac3{2\pi^{2}}\sum_{\substack{\boldsymbol a+\boldsymbol b=\boldsymbol m\\ \boldsymbol a,\boldsymbol b\ne\boldsymbol0}}
 V_{\boldsymbol a}V_{\boldsymbol b} .
\end{equation}
\end{corollary}

\begin{proof}
Expand \cref{p3:eq:index-from-phase} with $v=\sum V_{\boldsymbol m}
\boldsymbol\varepsilon^{\boldsymbol m}$ and compare coefficients. No higher powers occur because
$n$ is quadratic.
\end{proof}

The first two positive sectors, with $a_{2}=\widetilde A_{2}(N_{0})=\cos(\pi N_{0})$ and
$a_{3}=\widetilde A_{3}(N_{0})=2\cos(\pi N_{0})\cos(\tfrac\pi{18}+\tfrac{\pi N_{0}}3)$
(interpolated convention) or $a_{2}=(-1)^{r}$, $a_{3}=2\cos(\tfrac\pi{18}-\tfrac{2\pi r}3)$
(residue-frozen), are
\begin{align}
 \delta N^{(1)}_{2,+}
 &=-\frac{3\,a_{2}}{\pi^{2}\sqrt2}\,
 \frac{\mu_{0}^{2}(\mu_{0}-2)}{\mu_{0}^{2}-3\mu_{0}+3}\,\e^{-\mu_{0}/2},
 \label{p3:eq:dN2}\\
 \delta N^{(1)}_{3,+}
 &=-\frac{3\,a_{3}}{\pi^{2}\sqrt3}\,
 \frac{\mu_{0}^{2}(\mu_{0}-3)}{\mu_{0}^{2}-3\mu_{0}+3}\,\e^{-2\mu_{0}/3}.
 \label{p3:eq:dN3}
\end{align}
Thus \emph{parity of the index appears at relative scale $\e^{-\mu_{0}/2}$, the residue modulo $3$
at $\e^{-2\mu_{0}/3}$, and the residue modulo $k$ first appears in the growing $k$th sector at
$\e^{-(1-1/k)\mu_{0}}$.}

\subsection{The fractional-power staircase in the ordinate}

The core equation converts exponentials in $\mu_{0}$ into fractional powers of $y$ exactly:
\begin{equation}\label{p3:eq:exp-to-y}
 \e^{-\alpha\mu_{0}}=\left(\frac{\kappa(\mu_{0}-1)}{y\,\mu_{0}^{3}}\right)^{\!\alpha}
 \qquad\text{for every }\alpha,
\end{equation}
directly from $y=\kappa\e^{\mu_{0}}(\mu_{0}-1)\mu_{0}^{-3}$.

\begin{proposition}[Every first-order sector in closed $y$-form]
\label{p3:prop:y-form}
With $\alpha=\alpha_{k,\sigma}$ and $z=1/\mu_{0}$,
\begin{equation}\label{p3:eq:y-form}
 \delta N^{(1)}_{k,\sigma}
 =-\frac{3\kappa^{\alpha}}{\pi^{2}\sqrt k}\,a_{k}(\mu_{0})\;
 y^{-\alpha}\,\mu_{0}^{\,1-2\alpha}\sum_{s\ge0}g^{(k,\sigma)}_{s}\,\mu_{0}^{-s},
\end{equation}
where
\begin{equation}\label{p3:eq:g-gf}
 \sum_{s\ge0}g^{(k,\sigma)}_{s}z^{s}
 =\frac{(1-\sigma kz)(1-z)^{\alpha}}{1-3z+3z^{2}},
 \qquad
 g^{(k,\sigma)}_{s}=\sum_{i=0}^{s}(-1)^{i}\binom{\alpha}{i}
 \bigl(\chi_{s-i}-\sigma k\,\chi_{s-i-1}\bigr),
\end{equation}
with $\chi_{-1}=0$ and $\chi_{s}=\dfrac{\varrho_{+}^{s+1}-\varrho_{-}^{s+1}}
{\varrho_{+}-\varrho_{-}}$ the coefficients of $(1-3z+3z^{2})^{-1}$
($\varrho_{\pm}=\tfrac12(3\pm\ii\sqrt3)$; $\chi_{0}=1,\chi_{1}=3,\chi_{2}=6,
\chi_{3}=9,\chi_{4}=9,\chi_{5}=0,\dots$).
\end{proposition}

\begin{proof}
Insert \cref{p3:eq:exp-to-y} into \cref{p3:eq:first-N} and factor $\mu_{0}^{1-2\alpha}$ out of
$\dfrac{\mu_{0}^{2}(\mu_{0}-\sigma k)(\mu_{0}-1)^{\alpha}}
{(\mu_{0}^{2}-3\mu_{0}+3)\mu_{0}^{3\alpha}}$; with $z=1/\mu_{0}$ the remaining factor is
$\dfrac{(1-\sigma kz)(1-z)^{\alpha}}{1-3z+3z^{2}}$. The coefficient formula is the Cauchy product
of $(1-z)^{\alpha}=\sum_{i}(-1)^{i} \binom\alpha i z^{i}$ with $(1-\sigma kz)/(1-3z+3z^{2})$, and
$\chi_{s}$ is the standard solution of $\chi_{s}=3\chi_{s-1}-3\chi_{s-2}$, $\chi_{0}=1$,
$\chi_{1}=3$, whose closed form is as stated since $\varrho_{\pm}$ are the roots of
$\varrho^{2}=3\varrho-3$.
\end{proof}

\begin{corollary}[The inverse staircase of scales]\label{p3:cor:staircase-scales}
Since $\mu_{0}\sim\log y$, the primary positive sectors contribute, in order,
\begin{equation}\label{p3:eq:scales}
 y^{-1/2},\qquad y^{-2/3}(\log y)^{-1/3},\qquad y^{-3/4}(\log y)^{-1/2},
 \qquad\dots,\qquad y^{-(1-1/k)}(\log y)^{-1+2/k},
\end{equation}
accumulating at the scale $y^{-1}$, where they meet the square of the first sector and the whole
negative family.
\end{corollary}

\begin{proof}
By \cref{p3:eq:y-form} the $(k,+)$ sector has size
$y^{-\alpha}\mu_{0}^{1-2\alpha}(1+\bigO(\mu_{0}^{-1}))$ times a bounded amplitude, with
$\alpha=1-1/k$, so $1-2\alpha=-1+2/k$; and $\mu_{0}\sim\log y$ because
$\Phi(\mu_{0})=\log(y/\kappa)$ with $\Phi(\mu)=\mu(1+\littleo(1))$. Substituting $k=2,3,4$ gives
the first three displayed scales. Since $\alpha_{k,+}\uparrow1$, the scales accumulate at
$y^{-1}$; the square of the $(2,+)$ sector has action $\tfrac12+\tfrac12=1$ and every negative
action is $\ge1$.
\end{proof}

\begin{theorem}[The leading parity chirp]\label{p3:thm:leading-chirp}
In the interpolated convention,
\begin{equation}\label{p3:eq:chirp}
 \mathcal{N}(y)=N_{0}(y)
 -\frac{3\sqrt\kappa}{\pi^{2}\sqrt2}\,\frac{\cos\bigl(\pi N_{0}(y)\bigr)}{\sqrt y}
 \left(1+\bigO\!\left(\frac1{\log y}\right)\right)
 +\bigO\!\left(y^{-2/3}(\log y)^{-1/3}\right),
\end{equation}
with the constant
\begin{equation}\label{p3:eq:chirp-constant}
 \frac{3\sqrt\kappa}{\pi^{2}\sqrt2}=\frac{3^{1/4}}{2\pi}
 =0.2094596846361762626265841\ldots
\end{equation}
Since $N_{0}(y)=\frac{3}{2\pi^{2}}(\log y)^{2}(1+\littleo(1))$, the cosine is a \emph{quadratic
logarithmic chirp}, not a log-periodic oscillation.
\end{theorem}

\begin{proof}
Take $k=2$, $\sigma=+1$, $\alpha=\tfrac12$ in \cref{p3:eq:y-form}: the power
$\mu_{0}^{1-2\alpha}=\mu_{0}^{0}$ is absent, $g^{(2,+)}_{0}=1$, and $a_{2}=\cos(\pi N_{0})$ by
\cref{p3:lem:Atilde}. The constant is $3\kappa^{1/2}/(\pi^{2}\sqrt2)$, and
$\sqrt\kappa=\pi/\sqrt{6\sqrt3}$ gives $\frac{3\pi}{\pi^{2}\sqrt2\sqrt{6\sqrt3}}
=\frac{3}{\pi\sqrt{12\sqrt3}}=\frac{3}{2\pi\cdot3^{3/4}}=\frac{3^{1/4}}{2\pi}$, which we evaluated
to $25$ digits. The relative $\bigO(1/\log y)$ is the $s\ge1$ tail of \cref{p3:eq:y-form},
and the next omitted contribution is the $k=3$ sector at scale $y^{-2/3}(\log y)^{-1/3}$ by
\cref{p3:cor:staircase-scales}; all products of two blocks have action $\ge1$.
\end{proof}

\begin{remark}[Constant check]
S2 prints this constant as $0.2094596846361762626$; the closed form $3^{1/4}/(2\pi)$ is S3's. The
two agree, and both agree with our independent $25$-digit evaluation \cref{p3:eq:chirp-constant}.
S1 does not state the constant, giving instead the equivalent phase-variable form
\cref{p3:eq:dN2}; the three are algebraically identical, as the cancellation of $\mu_{0}-1$ in the
proof of \cref{p3:thm:first-order} shows.
\end{remark}

\section{The profinite arithmetic phase}\label{p3:sec:profinite}

The residue-frozen convention of \cref{p3:def:frozen} organises itself into a projective system,
which is the cleanest way to say what data an ``arithmetically complete'' inverse of $p$ actually
depends on.

\begin{theorem}[Compatibility and the profinite parameter]\label{p3:thm:profinite}
Let $\Lambda_{K}=\operatorname{lcm}(1,\dots,K)$, so $\Lambda_{K}\mid\Lambda_{K+1}$. Suppose
$r_{K}\in\Int/\Lambda_{K}\Int$ satisfy $r_{K+1}\equiv r_{K} \pmod{\Lambda_{K}}$, and let
$\mu_{K,r_{K}}$ be the corresponding inverse branches of \cref{p3:def:frozen}. Then for every $K$,
the transseries of $\mu_{K+1,r_{K+1}}$ obtained by deleting all monomials containing the variables
of level $K+1$ coincides with the transseries of $\mu_{K,r_{K}}$. Consequently the complete
arithmetic inverse is a projective family indexed by the profinite integers $\widehat
r\in\widehat{\Int}=\varprojlim\Int/m\Int$, and an actual partition index $n$ supplies its own
phase through the diagonal embedding $n\mapsto\widehat n$.
\end{theorem}

\begin{proof}
Deleting the level-$(K+1)$ variables from the defining equation \cref{p3:eq:v-equation} at level
$K+1$ produces the defining equation at level $K$ \emph{provided} the retained amplitudes agree,
i.e.\ provided $A_{k}(r_{K+1})=A_{k}(r_{K})$ for $k\le K$. That is exactly the compatibility
hypothesis, since $A_{k}$ has period $k\mid\Lambda_{K}$ and $r_{K+1}\equiv
r_{K}\pmod{\Lambda_{K}}$. Both truncated equations are triangular with the same normalisation
$V_{\boldsymbol0}=0$, and a triangular system has a unique solution, so the solutions agree
coefficient by coefficient. The projective limit statement is the definition of $\widehat{\Int}$:
compatible systems $(r_{K})$ are exactly the elements of
$\varprojlim\Int/\Lambda_{K}\Int=\widehat{\Int}$, because $\{\Lambda_{K}\}$ is cofinal in the
divisibility order on $\Nat$.
\end{proof}

\begin{remark}[Why a profinite parameter and not a real one]
\label{p3:rem:why-profinite}
At every algebraic order all indices share one expansion. Parity enters at relative order
$\e^{-\mu/2}$; the residue mod $3$ at $\e^{-2\mu/3}$; the residue mod $k$ at $\e^{-(1-1/k)\mu}$.
No finite real-valued smooth parameter records all of these congruence classes simultaneously, and
no single real analytic function of $n$ interpolates all $A_{k}$ canonically
(\cref{p3:rem:interp-nonunique}). A profinite integer does so canonically and finitely at every
exponential resolution: each finite action cutoff $\eta<1$ sees only $A_{k}$ for
$k\le(1-\eta)^{-1}$, i.e.\ only $\widehat r$ modulo $\Lambda_{\lfloor(1-\eta)^{-1}\rfloor}$. This
is S1's observation and is absent from S2 and S3.
\end{remark}

\begin{remark}[A structure suggested, not proved]\label{p3:rem:sheaf-conjecture}
S1 proposes an ``arithmetic transseries sheaf'' with base $\widehat{\Int}$ and fibres the
finite-level analytic transseries algebras, and calls \cref{p3:thm:profinite} its elementary case.
We record this as a suggestion. Nothing beyond \cref{p3:thm:profinite} --- namely, projective
compatibility of the coefficient systems --- is established here or in the sources: no topology on
the fibres, no gluing statement, and no continuity of $\widehat r\mapsto\mu_{K,\widehat r}$ in any
topology on $\widehat{\Int}$ has been proved. It is stated as \cref{p3:conj:sheaf} in
\cref{p3:sec:open}, not as a theorem.
\end{remark}

\section{Validity, truncation control, and integer recovery}
\label{p3:sec:validity}

\subsection{Remainder transport}

\begin{theorem}[Inverse error after $K$ Rademacher levels]\label{p3:thm:inverse-tail}
Fix $K\ge1$. Let $y=p(n)$, let $\mu=\lambda\sqrt{n-1/24}$, and let $\mu_{K}$ denote either inverse
branch of \cref{p3:def:frozen} (with $r\equiv n$) or of \cref{p3:def:interpolated}. Then
\begin{equation}\label{p3:eq:inverse-tail}
 \mu_{K}(y)-\mu=\bigO_{K}\!\left(\mu^{3/2}\,
 \e^{-(1-\frac1{K+1})\mu}\right),
 \qquad
 n_{K}(y)-n=\bigO_{K}\!\left(\mu^{5/2}\,\e^{-(1-\frac1{K+1})\mu}\right),
\end{equation}
where $n_{K}=\frac1{24}+3\mu_{K}^{2}/(2\pi^{2})$. Explicitly, with the constant of
\cref{p3:thm:tail}, for $\mu\ge\mu_{*}$ (any threshold at which $\Phi'\ge\tfrac12$ and the omitted
tail is $\le\tfrac12$ in absolute value),
\begin{equation}\label{p3:eq:inverse-tail-explicit}
 |n_{K}(y)-n|\le\frac{3\mu}{\pi^{2}}\cdot\frac{2}{1}\cdot
 \frac{8\mu^{2}}{3(1-\mu^{-1})\sqrt K}
 \exp\!\left[-\Bigl(1-\frac1{K+1}\Bigr)\mu\right]
 \bigl(1+\littleo(1)\bigr).
\end{equation}
\end{theorem}

\begin{proof}
This is \cref{p0:thm:remainder-transport} with the forward remainder supplied by
\cref{p3:thm:tail,p3:thm:tail-sharp}. In detail: let $\mathcal{F}_{K}=\log\mathcal{P}_{K}$ and
$\mathcal{F}=\log p$ regarded as functions of the phase. By \cref{p3:prop:relative},
$\mathcal{F}(\mu)-\mathcal{F}_{K}(\mu) =\log\bigl(1+\mathcal{R}_{>K}/(1+\mathcal{R}_{\le
K})\bigr)$, which by \cref{p3:thm:tail-sharp} and $|\log(1+u)|\le2|u|$ for $|u|\le\tfrac12$ is
$\bigO_{K}(\mu^{3/2}\e^{-(1-1/(K+1))\mu})$. Since
$\mathcal{F}_{K}'(\mu)=\Phi'(\mu)(1+\littleo(1))\to1$, the mean value theorem converts this
logarithmic residual into the phase error, which is the first half of \cref{p3:eq:inverse-tail}.
Finally $\dd n/\dd\mu=3\mu/\pi^{2}$ gives the index error. The explicit form
\cref{p3:eq:inverse-tail-explicit} replaces \cref{p3:thm:tail-sharp} by the explicit
\cref{p3:thm:tail} (losing $\mu^{-1/2}$, gaining a constant) and $\mathcal{F}_{K}'\ge\tfrac12$.
\end{proof}

\begin{corollary}[Beyond-all-orders accuracy of the core inverse]
\label{p3:cor:core-inverse-error}
For $y=p(n)$,
\begin{equation}\label{p3:eq:core-error}
 N_{0}(y)-n=\bigO\!\left(\mu^{5/2}\e^{-\mu/2}\right)
 =\bigO\!\left(n^{5/4}\exp\!\left(-\tfrac12\pi\sqrt{2n/3}\right)\right),
\end{equation}
so $N_{0}(p(n))-n\to0$ faster than every power of $n^{-1}$.
\end{corollary}

\begin{proof}
\Cref{p3:thm:inverse-tail} at $K=1$; note $N_{0}=n_{1}$ because the $K=1$ level of
\cref{p3:def:frozen} contains the sectors $(1,\pm)$, and the $(1,-)$ sector has action $2$, hence
contributes below the $K=1$ error bound.
\end{proof}

\begin{proposition}[A posteriori certificate; specialisation of
\cref{p0:thm:backward-error}]\label{p3:prop:certificate} Let
$\mathcal{F}_{K}(\mu)=\log\mathcal{P}_{K}(\mu)$ and let
$\rho(\widetilde\mu)=\mathcal{F}_{K}(\widetilde\mu)-\log y$ be the residual of a candidate
$\widetilde\mu$. If $\mathcal{F}_{K}'\ge m>0$ on an interval $I$ containing both $\widetilde\mu$
and the exact root $\mu_{K}(y)$, then
\begin{equation}\label{p3:eq:certificate}
 |\widetilde\mu-\mu_{K}(y)|\le\frac{|\rho(\widetilde\mu)|}{m},
 \qquad
 |\widetilde n-n_{K}(y)|\le\frac{3}{2\pi^{2}}
 \bigl(|\widetilde\mu|+|\mu_{K}(y)|\bigr)\frac{|\rho(\widetilde\mu)|}{m}.
\end{equation}
\end{proposition}

\begin{proof}
Mean value theorem applied to $\mathcal{F}_{K}$, then the factorisation
$\widetilde\mu^{2}-\mu_{K}^{2}=(\widetilde\mu-\mu_{K})(\widetilde\mu+\mu_{K})$. Since
$\mathcal{F}_{K}'\to1$, one may take $m=\tfrac12$ beyond an effectively checkable threshold.
\end{proof}

\begin{remark}[Why a residual certificate rather than term counting]
\Cref{p3:prop:certificate} separates the \emph{local inversion} error (how well $\widetilde\mu$
solves the level-$K$ equation) from the \emph{Rademacher tail} (how far the level-$K$ equation is
from the truth). Only the second needs \cref{p3:thm:tail}. This modularity is what makes it
possible to certify a truncated symbolic transseries and a numerically refined Newton root by the
same inequality.
\end{remark}

\subsection{Algebraic truncation}

\begin{proposition}\label{p3:prop:algebraic-truncation}
With $U_{M}(z)=\sum_{m\le M}c_{m}z^{m}$ as in \cref{p3:thm:core-correction},
\begin{equation}\label{p3:eq:algebraic-truncation}
 \mu_{0}-\bigl(X+U_{M}(X^{-1})\bigr)=\bigO(X^{-M-1}),
 \qquad
 N_{0}-\left[\frac1{24}+\frac3{2\pi^{2}}\bigl(X+U_{M}(X^{-1})\bigr)^{2}\right]
 =\bigO(X^{-M}),
\end{equation}
the loss of one order coming from $\dd n/\dd\mu\asymp X$. In particular $M\ge1$ is
\emph{necessary} for the index to be correct to $\littleo(1)$ and hence for the arithmetic
phases of \cref{p3:sec:full-inverse} to be evaluated at the right point
(\cref{p3:rem:phase-lock}).
\end{proposition}

\begin{proof}
$U$ is analytic at $0$ by \cref{p3:thm:core-correction}, so its Taylor remainder after $M$ terms
is $\bigO(z^{M+1})$; substitute $z=X^{-1}$ and use $N_{0}-\widetilde
N=\frac3{2\pi^{2}}(\mu_{0}-\widetilde\mu) (\mu_{0}+\widetilde\mu)$ with
$\mu_{0}+\widetilde\mu=\bigO(X)$.
\end{proof}

\subsection{Integer recovery}

\begin{theorem}[Eventual exact recovery by rounding]\label{p3:thm:rounding}
There is $n_{0}$ such that for every integer $n\ge n_{0}$,
\begin{equation}\label{p3:eq:rounding}
 n=\bigl\lfloor N_{0}(p(n))\bigr\rceil,
\end{equation}
$\lfloor\cdot\rceil$ denoting nearest-integer rounding. Consequently the arithmetic phase
$\widehat n$ of \cref{p3:thm:profinite} is recoverable from $y=p(n)$ \emph{before} any
exponentially small sector is evaluated.
\end{theorem}

\begin{proof}
By \cref{p3:eq:core-error}, $|N_{0}(p(n))-n|\to0$; it is therefore $<\tfrac12$ for all large $n$,
at which point rounding is exact. The last sentence is the observation that $A_{k}(n)$ depends
only on $n\bmod k$.
\end{proof}

\begin{remark}[The threshold is very small in practice]
\label{p3:rem:threshold}
\Cref{p3:thm:rounding} is asymptotic and claims no optimal $n_{0}$. We computed $N_{0}(p(n))-n$
for all $1\le n\le 59$ in $140$-digit arithmetic; the largest absolute value is $0.1723\ldots$ at
$n=1$, and the sequence decreases in magnitude thereafter (through $2.8\times10^{-2}$ at $n=10$,
$4.5\times10^{-4}$ at $n=50$, $1.5\times10^{-5}$ at $n=100$, $4.3\times10^{-17}$ at $n=1000$). So
$n_{0}=1$ is consistent with the data, though we prove only the asymptotic statement; producing a
certified $n_{0}$ requires effective constants in \cref{p3:thm:inverse-tail}, i.e.\ effective
Rademacher remainder bounds, which is listed as \cref{p3:conj:effective-threshold}.
\end{remark}

\begin{remark}[Circularity, avoided]\label{p3:rem:no-circularity}
The arithmetic transseries needs the residue $r$ of $n$, which appears to require knowing $n$. The
order of operations resolves this: \emph{(i)} compute $X$ from \cref{p3:eq:carrier}; \emph{(ii)}
compute $\mu_{0}$ and $N_{0}$, which are phase-independent; \emph{(iii)} round, obtaining $n$ by
\cref{p3:thm:rounding} for $y$ on the range; \emph{(iv)} now every $A_{k}(n)$ is known, and the
exponentially small sectors refine an already identified integer. The arithmetic layer is thus a
\emph{refinement} of a known answer, not a device for guessing the phase. For $y$ off the range,
step (iii) yields $\lceil\mathcal{N}(y)\rceil$ by \cref{p3:thm:staircase} instead, and the sectors
describe $\mathcal{N}$ itself.
\end{remark}

\section{Numerical verification}\label{p3:sub:numerics}

Every number in this chapter was recomputed for this volume. Partition values were generated
exactly by Euler's pentagonal recurrence
\begin{equation}\label{p3:eq:pentagonal}
 p(n)=\sum_{m\ge1}(-1)^{m-1}
 \left[p\!\left(n-\tfrac{m(3m-1)}2\right)+p\!\left(n-\tfrac{m(3m+1)}2\right)\right],
 \qquad p(0)=1,\ p(r)=0\ (r<0),
\end{equation}
in exact integer arithmetic (checked against $p(10)=42$, $p(50)=204226$, $p(100)=190569292$);
Dedekind sums exactly as rationals from \cref{p3:eq:dedekind}; the transcendental factors in
$140$-digit floating point. Coefficient identities were verified in exact rational or exact
symbolic arithmetic, never by floating-point spot checks. The tables check normalisation, signs,
phases and the hierarchy of actions; they are used in no proof.

\begin{table}[htbp]
\centering
\caption{Relative error $\bigl(S_{K}(n)-p(n)\bigr)/p(n)$ of the paired truncation
$S_{K}=\sum_{k\le K}(P_{k,+}+P_{k,-})$, recomputed at $140$ digits.  Accuracy need not
improve monotonically in $K$ at small $n$: the amplitudes $A_{k}(n)$ change sign and can
suppress or enhance an individual sector.}
\label{p3:tab:forward}
\small
\begin{tabular}{@{}rrrrrr@{}}
\toprule
$n$ & $p(n)$ & $K=1$ & $K=2$ & $K=3$ & $K=5$\\
\midrule
$10$ & $42$ & $-8.8621\cdot10^{-3}$ & $1.6606\cdot10^{-3}$ & $3.7391\cdot10^{-4}$ & $3.4722\cdot10^{-4}$\\
$25$ & $1\,958$ & $1.0717\cdot10^{-3}$ & $3.0045\cdot10^{-6}$ & $-6.1015\cdot10^{-5}$ & $-3.4866\cdot10^{-6}$\\
$50$ & $204\,226$ & $-7.3074\cdot10^{-5}$ & $3.9028\cdot10^{-6}$ & $2.1050\cdot10^{-7}$ & $-1.1350\cdot10^{-7}$\\
$100$ & $190\,569\,292$ & $-1.8220\cdot10^{-6}$ & $8.6838\cdot10^{-9}$ & $-4.9491\cdot10^{-9}$ & $3.1454\cdot10^{-10}$\\
$250$ & $2.3079\cdot10^{14}$ & $-1.0760\cdot10^{-9}$ & $7.2267\cdot10^{-13}$ & $4.3021\cdot10^{-14}$ & $2.1256\cdot10^{-15}$\\
$500$ & $2.3002\cdot10^{21}$ & $-2.4389\cdot10^{-13}$ & $1.7537\cdot10^{-17}$ & $-2.0017\cdot10^{-19}$ & $-3.5933\cdot10^{-22}$\\
$1000$ & $2.4061\cdot10^{31}$ & $-1.6997\cdot10^{-18}$ & $1.2574\cdot10^{-24}$ & $-3.4744\cdot10^{-27}$ & $3.9453\cdot10^{-30}$\\
\bottomrule
\end{tabular}
\end{table}

\begin{table}[htbp]
\centering
\caption{Inverse errors at exact ordinates $y=p(n)$; entries are (approximation) $-\,n$.
$N_{X}$ is the bare Lambert carrier index, $N_{0}$ the dominant core, $N_{K}$ the Newton
inverse of the level-$K$ interpolated sum (\cref{p3:def:interpolated}).  The
residue-frozen inverses (\cref{p3:def:frozen}) reproduce the last three columns to all
five digits shown for $n\ge100$; at $n=50$ they differ in the fifth digit of the
$K=3,5$ entries ($-1.2992\cdot10^{-6}$, $7.0057\cdot10^{-7}$), at $n=25$ in the fourth
($2.7926\cdot10^{-4}$, $1.5958\cdot10^{-5}$), at $n=10$ in the third
($-5.3338\cdot10^{-3}$, $-1.2006\cdot10^{-3}$, $-1.1150\cdot10^{-3}$).  That is exactly
the pattern \cref{p3:prop:two-conventions} predicts.}
\label{p3:tab:inverse}
\small
\begin{tabular}{@{}rrrrrr@{}}
\toprule
$n$ & $N_{X}-n$ & $N_{0}-n$ & $N_{2}-n$ & $N_{3}-n$ & $N_{5}-n$\\
\midrule
$10$ & $-3.9910\cdot10^{-1}$ & $2.8447\cdot10^{-2}$ & $-5.3290\cdot10^{-3}$ & $-1.1863\cdot10^{-3}$ & $-1.1015\cdot10^{-3}$\\
$25$ & $-3.7877\cdot10^{-1}$ & $-4.9052\cdot10^{-3}$ & $-1.3751\cdot10^{-5}$ & $2.7902\cdot10^{-4}$ & $1.5941\cdot10^{-5}$\\
$50$ & $-3.5044\cdot10^{-1}$ & $4.5102\cdot10^{-4}$ & $-2.4089\cdot10^{-5}$ & $-1.2993\cdot10^{-6}$ & $7.0059\cdot10^{-7}$\\
$100$ & $-3.3600\cdot10^{-1}$ & $1.5378\cdot10^{-5}$ & $-7.3293\cdot10^{-8}$ & $4.1772\cdot10^{-8}$ & $-2.6548\cdot10^{-9}$\\
$250$ & $-3.2364\cdot10^{-1}$ & $1.3942\cdot10^{-8}$ & $-9.3644\cdot10^{-12}$ & $-5.5747\cdot10^{-13}$ & $-2.7544\cdot10^{-14}$\\
$500$ & $-3.1768\cdot10^{-1}$ & $4.4042\cdot10^{-12}$ & $-3.1668\cdot10^{-16}$ & $3.6146\cdot10^{-18}$ & $6.4888\cdot10^{-21}$\\
$1000$ & $-3.1356\cdot10^{-1}$ & $4.2960\cdot10^{-17}$ & $-3.1780\cdot10^{-23}$ & $8.7815\cdot10^{-26}$ & $-9.9717\cdot10^{-29}$\\
\bottomrule
\end{tabular}
\end{table}

Three features of \cref{p3:tab:inverse} confirm the theory point by point. \emph{(i)} The first
column tends to $-3/\pi^{2}=-0.30396\ldots$, not to $0$: this is \cref{p3:cor:carrier-bias} and
the numerical signature of \cref{p3:rem:phase-lock}. (S1 and S2 tabulate
$N_{X}=\tfrac1{24}+3X^{2}/(2\pi^{2})$ and report $\to-0.304$; S3 tabulates $3X^{2}/(2\pi^{2})$
without the $\tfrac1{24}$ and reports $\to-0.346$. The columns differ by exactly $1/24$ and each
is correct for its own definition; we keep the $\tfrac1{24}$ so that $N_{X}$ and $N_{0}$ are
directly comparable.) \emph{(ii)} Restoring the amplitude factor $1-1/\mu$, i.e.\ passing from $X$
to $\mu_{0}$, removes that entire $\bigO(1)$ bias and leaves an exponentially small error.
\emph{(iii)} The residual after the core alternates with the parity of $n$: at $n=49,50,51$ it is
$-5.2643\cdot10^{-4}$, $+4.5102\cdot10^{-4}$, $-4.0920\cdot10^{-4}$, exactly as \cref{p3:eq:dN2}
predicts through $\widetilde A_{2}=\cos\pi N_{0}$. Adding the single first-order $k=2$ sector
reduces these to $-1.0890\cdot10^{-5}$, $-2.4087\cdot10^{-5}$, $+2.9029\cdot10^{-5}$; adding the
first-order sectors through $k=10$ reduces them to $-2.24\cdot10^{-7}$, $-9.95\cdot10^{-8}$,
$-5.81\cdot10^{-8}$. This checks both the sign and the phase of the leading chirp
\cref{p3:eq:chirp}.

\begin{remark}[What was verified in exact arithmetic]\label{p3:rem:symbolic-checks}
(1) \Cref{p3:eq:c-finite} against the Taylor expansion of \cref{p3:eq:Omega}, symbolically in
$\beta$ through order $t^{9}$: identical; the closed forms \cref{p3:eq:c-first,p3:eq:c-six} were
produced from \cref{p3:eq:c-finite} and match S2 (through $c_{5}$) and S3 (through $c_{8}$). (2)
$\omega_{0},\dots,\omega_{4}$ against both the S1 and the S2 grouping of the same numbers:
identical (they look different only because S1 splits each into partial fractions in $\pi$). (3)
$c_{1},\dots,c_{18}$ of \cref{p3:tab:cm} from \cref{p3:eq:c-recurrence-core} in exact rationals,
plus the residual test \cref{p3:eq:U-residual} at $M=12$; all eighteen agree with the sources. (4)
$\gamma_{1},\dots,\gamma_{8}$ against all three sources; $\gamma_{9},\gamma_{10}$ are new here.
(5) $\Pi_{1},\dots,\Pi_{8}$ from the Lagrange form against the Stirling form and against each of
the three printed shapes, confirming \cref{p3:rem:H-versus-R}; S1's printed $P_{6},P_{7}$ are
correct. (6) S3's $u$-expansion in $H$ through $H^{-4}$ and its $N_{0}$-expansion through
$H^{-3}$, and S2's $N_{0}$-block in $R$ through the constant term: correct as printed. (7)
\Cref{p3:eq:first-Ak} against direct evaluation of \cref{p3:eq:Ak} from exact Dedekind sums. (8)
The chirp constant \cref{p3:eq:chirp-constant} to $25$ digits in both forms. No discrepancy with
any source's arithmetic was found: every repair recorded in this chapter is structural, not
numerical.
\end{remark}

\begin{remark}[Normalisation cross-checks]\label{p3:sub:normalization-crosscheck}
Because \cref{p3:eq:rademacher} carries a factor $\sqrt k$ that some normalisations absorb into
$A_{k}$, a missing or spurious power of $k$ is the easiest error to make and the hardest to see.
Three identifications of one quantity detect it: the $k$th derivative factor of
\cref{p3:eq:rademacher} equals $\frac{\mu}{4kN^{3/2}}[\e^{\mu/k}(1-k/\mu)+\e^{-\mu/k}(1+k/\mu)]$,
and after multiplication by $\sqrt k/(\pi\sqrt2)$ it may be written either as
$\frac1{4\sqrt{3k}\,N}[\cdots]$ or as $\frac{\kappa}{\sqrt k\,\mu^{2}}[\cdots]$; at $k=1$ the
positive term is $\frac{\e^{\mu}}{4\sqrt3\,N}(1-\tfrac1\mu)$, which becomes \cref{p3:eq:HR} after
replacing $N$ by $n$ and dropping $1-1/\mu$. Any of these disagreeing with the others by a power
of $k$ localises the slip.
\end{remark}

\section{Algorithms}\label{p3:sec:algorithms}

All formulas of this chapter are constructive and need only formal power-series arithmetic over
$\Rat(\pi)$ together with exact rational Dedekind sums.

\begin{algorithm}[Forward sector coefficients]\label{p3:alg:forward}
Given $k,\sigma,M$: set $\beta=\sigma\pi/(6k)$ and evaluate the finite sum \cref{p3:eq:c-finite}
for $0\le m\le M$ ($\bigO(m)$ field operations each, no series manipulation; keep $\beta$
symbolic if sector-uniform output is wanted). If the whole array is wanted at once, evaluate
$\ell_{1},\dots,\ell_{M}$ from \cref{p3:eq:ell} and run the triangular recurrence
\cref{p3:eq:c-recurrence} instead. Assemble \cref{p3:eq:sector-unshifted}, pairing the two signs
before summing over large $k$ (\cref{p3:rem:pairing}).
\end{algorithm}

\begin{algorithm}[Phase-independent inverse]\label{p3:alg:core}
Given a large $y$ and an order $M$: (1) compute $X$ from \cref{p3:eq:carrier} with a standard
$\Wlam_{-1}$ routine --- do \emph{not} substitute the logarithmic expansion
\cref{p3:eq:carrier-flat}, which is numerically far inferior; (2) form $\mu_{0}=X+\sum_{m\le
M}c_{m}X^{-m}$ from the universal rationals of \cref{p3:tab:cm}, optionally refined by one Newton
step on $\Phi(\mu)=R$; (3) return $N_{0}=\tfrac1{24}+3\mu_{0}^{2}/(2\pi^{2})$; (4) certify by
evaluating $\rho=\mu_{0}-2\log\mu_{0}+\log(1-\mu_{0}^{-1})-\log(y/\kappa)$ and applying
\cref{p3:prop:certificate}; (5) round (\cref{p3:thm:rounding}) if $y$ is known to lie on the
range, else take $\lceil\cdot\rceil$ for the staircase (\cref{p3:thm:staircase}).
\end{algorithm}

\begin{algorithm}[Arithmetic exponential correction]\label{p3:alg:full}
Given $y$, a level $K$ and a finite set of multi-indices: (1) compute $\mu_{0},N_{0}$ by
\cref{p3:alg:core} and use $N_{0}$, never $N_{X}$, to evaluate amplitudes
(\cref{p3:rem:phase-lock}); (2) compute exact rational $s(h,k)$ for $k\le K$ and hence $A_{k}(r)$
(frozen, with $r$ from \cref{p3:thm:rounding}) or $\widetilde A_{k}(N_{0})$ (interpolated); (3)
build the actions and amplitudes of \cref{p3:eq:block}; (4) traverse the multi-indices in a graded
monomial order using \cref{p3:eq:triangular}, or, for one isolated $\boldsymbol m$, use
\cref{p3:eq:multi-LB} or the Bell form \cref{p3:eq:multi-Bell} with the closed jets
\cref{p3:eq:jets}; (5) substitute $\varepsilon_{j}=\e^{-\alpha_{j}\mu_{0}}$ and convert to the
index by \cref{p3:eq:N-multi}; (6) certify by \cref{p3:prop:certificate} on $\mathcal{P}_{K}$,
adding \cref{p3:thm:tail} if the comparison is with the exact sequence.
\end{algorithm}

\begin{remark}[Modular structure of the computation]\label{p3:rem:reuse}
The layers are independent and should be cached separately: the carrier polynomials $\Pi_{j}$
depend only on the exponent data $(a,b)=(1,-2)$; the rationals $c_{m}$ only on the amplitude
$1-1/\mu$; the inverse-derivative polynomials $\mu_{0}^{(s)}$ only on $\Phi$; and the arithmetic
enters solely through the numbers $a_{k}$ in \cref{p3:eq:block}. Universal tables are computed
once and residue-dependent evaluation is then cheap.
\end{remark}

\section{A template for Rademacher-type sequences}\label{p3:sec:template}

Nothing above used the partition function except through four properties. S3 has the most careful
version of this observation and S1 the most usable one; what follows merges them into a hypothesis
list and a theorem, so that the template is a statement one can check rather than a slogan.

\begin{definition}[Rademacher-type representation]\label{p3:def:rademacher-type}
A sequence $(a(n))_{n\ge n_{1}}$ of positive reals is of \emph{Rademacher type} if:
\begin{enumerate}
\item[(T1)] there is a real analytic, strictly increasing \emph{phase}
$\phi:(\nu_{1},\infty)\to(0,\infty)$ with $\phi\to\infty$ and real analytic inverse, giving the
large variable $\mu=\phi(n)$;
\item[(T2)] the \emph{envelope} $F(\mu)=C\e^{a\mu^{\alpha}}\mu^{b}\exp Q(\mu^{-\delta})$
is an exponential--power model \cref{p0:def:model} with $a>0$, $\alpha\in(0,1]$;
\item[(T3)] there are actions $\alpha_{j}>0$ ($j\in J$) with $\inf_{j}\alpha_{j}>0$,
amplitudes $a_{j}(n)$ taking values in a fixed finite-dimensional space at each truncation level
(e.g.\ periodic in $n$), and $g_{j}$ analytic at $\infty$, with
      \begin{equation}\label{p3:eq:template}
       a(n)=F(\mu)\Bigl(1+\sum_{j\in J'}a_{j}(n)g_{j}(\mu)\e^{-\alpha_{j}\mu}
       +\mathcal{T}_{J'}(\mu;n)\Bigr)\quad\text{for every finite }J'\subset J;
      \end{equation}
\item[(T4)] $\{j:\alpha_{j}\le\eta\}$ is finite for every $\eta<\sup_{j}\alpha_{j}$, and
either $\mathcal{T}_{J'}=\bigO(\mu^{C'}\e^{-\eta'\mu})$ for some $\eta'>\max_{j\in
J'}\alpha_{j}$, or --- if the actions accumulate --- the grouped remainder is retained as one
analytic block.
\end{enumerate}
\end{definition}

\begin{theorem}[Template]\label{p3:thm:template}
Let $(a(n))$ be of Rademacher type. Then:
\begin{enumerate}
\item \textup{(Reduction.)} $\mu=\xi^{1/\alpha}$ of \cref{p0:lem:alpha-reduction} brings
\textup{(T2)} to the linear-phase model with computable $(C',1,b',Q')$.
\item \textup{(Carrier.)} $C'\e^{\xi}\xi^{b'}=y$ is solved exactly by
\cref{p0:thm:lambert-core}; equivalently, for $y=C\e^{a\mu^{\alpha}}\mu^{-\beta}$,
      \begin{equation}\label{p3:eq:general-carrier}
       \mu=\left(-\frac{\beta}{a\alpha}\,\Wlam_{-1}\!\left(-\frac{a\alpha}{\beta}
       \Bigl(\frac Cy\Bigr)^{\alpha/\beta}\right)\right)^{1/\alpha}.
      \end{equation}
\item \textup{(Amplitude.)} The envelope inverse is $\xi_{0}=X+U(X^{-1})$ with $U$ the
unique solution of the fixed-point equation of \cref{p0:thm:lambert-centered} at $(1,b',Q')$;
coefficients by \cref{p0:lem:lagrange-fp} or by isolating the linear occurrence of $U$.
\item \textup{(Phase locking.)} If some $a_{j}$ is periodic on the index scale and
$\dd n/\dd\mu\to\infty$, the amplitudes must be evaluated at $\xi_{0}$, not at $X$: the carrier
bias $n(\xi_{0})-n(X)$ is $\bigO(1)$ in general.
\item \textup{(Sectors.)} With the phase frozen, every inverse transmonomial is given by
\cref{p3:eq:multi-LB} (in $y$) or \cref{p3:eq:phase-LB}/\cref{p3:eq:triangular} (in $\mu$), with
first-order term
      \begin{equation}\label{p3:eq:general-first}
       \mu=\mu_{0}-\frac{a_{j}g_{j}(\mu_{0})}{(\log F)'(\mu_{0})}\e^{-\alpha_{j}\mu_{0}}
       +\bigO\bigl(\e^{-2\alpha_{\min}\mu_{0}}\bigr).
      \end{equation}
\item \textup{(Transport.)} A relative forward remainder $\epsilon$ transports to a phase
error $\epsilon/(\log F)'(\mu_{0})$ and an index error $\epsilon\,n'(\mu_{0})/(\log F)'(\mu_{0})$,
by \cref{p0:thm:remainder-transport}.
\item \textup{(Staircase.)} If $a$ is integer valued with gaps bounded below, the
generalised inverse is the ceiling of the smooth one (\cref{p0:thm:staircase}), with the order-one
sawtooth of \cref{p3:rem:sawtooth} as an irreducible final sector.
\end{enumerate}
\end{theorem}

\begin{proof}
(i)--(iii) are \cref{p0:lem:alpha-reduction,p0:thm:lambert-core,p0:thm:lambert-centered} at the
stated parameters. For the display in (ii): raising $y=C\e^{a\mu^{\alpha}}\mu^{-\beta}$ to the
power $\alpha/\beta$ and rearranging gives $\bigl(\tfrac{a\alpha}{\beta}\mu^{\alpha}\bigr)
\e^{-(a\alpha/\beta)\mu^{\alpha}}=\tfrac{a\alpha}{\beta}(C/y)^{\alpha/\beta}$, which is the
Lambert equation with the stated argument, the branch $-1$ forced by $\mu\to\infty$. (iv) is the
computation of \cref{p3:rem:phase-lock} with $3\mu/\pi^{2}$ replaced by $n'(\mu_{0})$: what
matters is only that $n'\to\infty$ while $\xi_{0}-X\to0$, so their product need not vanish. (v) is
\cref{p3:thm:additive-LB,p3:cor:phase-LB,p3:thm:triangular}, whose proofs used no property of $F$
or of the blocks beyond analyticity and $F'\ne0$; the first-order display is \cref{p3:eq:first-V}
in that generality. (vi) is the mean-value argument of \cref{p3:thm:inverse-tail}, and (vii) is
\cref{p3:thm:staircase}.
\end{proof}

\begin{remark}[The partition case, and the scope of the template]
\label{p3:rem:template-scope}
Taking $\phi(n)=\pi\sqrt{2(n-1/24)/3}$, $(C,a,\alpha,b,Q)$ as in \cref{p3:eq:model-in-N},
$J=\{(k,\sigma)\}$, $\alpha_{k,\sigma}=1-\sigma/k$ and $a_{k}=A_{k}$, hypotheses (T1)--(T3) hold
by \cref{p3:thm:exact-sectors,p3:prop:relative}, and (T4) holds in its second form: the actions
\emph{do} accumulate at $1$, and the grouped block is the paired tail \cref{p3:eq:R-exact}
controlled by \cref{p3:thm:tail}. With $(\alpha,\beta)=(1,2)$ in $\mu$,
\cref{p3:eq:general-carrier} reduces to \cref{p3:eq:carrier}. The template applies verbatim to
coefficient sequences of weakly holomorphic modular forms and eta-quotients with Rademacher-type
expansions; the partition function is unusually clean because the Bessel index $3/2$ makes the
kernel two exponentials with amplitudes linear in $\mu^{-1}$, so step (iii) is the explicit
rational recurrence \cref{p3:eq:c-recurrence-core} rather than a general implicit solve. \emph{No
theorem about any other specific sequence is claimed here}: verifying (T1)--(T4) for a given
family is exactly the work the template does not do.
\end{remark}

\section{Structural consequences and open questions}\label{p3:sec:open}

\subsection{Convergent blocks inside an exponentially complete object}

The exact sector sum \cref{p3:eq:paired-sector} converges; each sector's algebraic series in
$n^{-1/2}$ converges on $|n^{-1/2}|<\sqrt{24}$ (\cref{p3:thm:sector-gf}); and the reversion series
$U$ converges near $0$ (\cref{p3:thm:core-correction}). Divergence is therefore \emph{not} what
creates the non-perturbative sectors here: they are supplied by modularity and the circle method,
independently of the perturbative data. The useful trichotomy is that a \emph{Poincar\'e
expansion} records one dominant sector and forgets all exponentially small information; a
\emph{transseries completion} records the other sectors, whether or not the series in each
diverges; and an \emph{exact transseries representation}, as here, may itself converge. Compare
\cref{p1:sec:top,p2:sec:top,p4:sec:top}, where the algebraic series do diverge and
\cref{p0:thm:optimal-truncation} is the operative tool.

\subsection{Arithmetic amplitudes are not Stokes constants}

The $A_{k}(n)$ weight distinct exponential sectors and govern exponentially small corrections to
the inverse, which is formally the role of Stokes constants. They are not Stokes constants: they
are arithmetic, periodic in $n$ rather than fixed, and no natural analytic continuation in the
index produces them canonically (\cref{p3:rem:interp-nonunique}). \Cref{p3:thm:profinite} says
what replaces the Stokes data: a profinite integer, of which every finite exponential resolution
sees only a finite quotient.

\subsection{Open questions}

\begin{conjecture}[Effective rounding threshold]\label{p3:conj:effective-threshold}
There is an explicit small $n_{0}$ --- the evidence of \cref{p3:rem:threshold} suggests $n_{0}=1$
--- with $n=\lfloor N_{0}(p(n))\rceil$ for all $n\ge n_{0}$, provable by inserting effective
Rademacher remainder bounds (Lehmer; Banerjee--Paule--Radu--Schneider) in place of $|A_{k}|\le k$
in \cref{p3:thm:inverse-tail}.
\end{conjecture}

\begin{conjecture}[Arithmetic transseries sheaf]\label{p3:conj:sheaf}
The projective family of \cref{p3:thm:profinite} extends to a sheaf of transseries algebras over
$\widehat{\Int}$, with $\widehat r\mapsto(\text{level-}K$ inverse transseries$)$ continuous for
the profinite topology and the finite-level algebras gluing over the standard clopen basis. Only
the projective compatibility is proved here (\cref{p3:rem:sheaf-conjecture}).
\end{conjecture}

\begin{conjecture}[Uniform treatment at the accumulation action]
\label{p3:conj:uniform-accumulation}
The zeta condensation \cref{p3:eq:zeta-condensation} can serve as a single block in
\cref{p3:thm:additive-LB}, yielding a description of the inverse uniform across the accumulation
action $1$, with error bounds depending on the analytic behaviour of $\mathcal{Z}_{n}(s)$ rather
than on a sector cutoff. This needs continuation and vertical growth information for the
Kloosterman-type Dirichlet series $\mathcal{Z}_{n}$, which is not established here. S3 raises it
as an optimisation problem rather than a gap, correctly: the exact Rademacher representation is
always available as a fallback.
\end{conjecture}

\begin{remark}[Classifying interpolations]\label{p3:rem:interp-problem}
Different strictly increasing smooth interpolations of $p$ agree at every algebraic order of the
inverse and on the integer range, but differ in the exponentially small off-lattice sectors:
candidates include \cref{p3:def:Atilde}, piecewise-logarithmic interpolation through $(n,\log
p(n))$, Newton or cardinal series with prescribed growth, and continuations built from Poincar\'e
series. Which part of the inverse transseries is interpolation-independent?
\Cref{p3:prop:two-conventions} answers this for one pair of choices --- they agree strictly below
action $1$ --- and a general statement is open.
\end{remark}

\begin{remark}[Formalisation targets]\label{p3:rem:formalisation}
Rademacher's theorem is now machine-checked (Eberl and Paulson, Archive of Formal Proofs, 2026).
The natural next steps, in dependency order, are: (a) the normalisation
\cref{p3:thm:exact-sectors}; (b) the generating function \cref{p3:eq:Omega} and the recurrence
\cref{p3:eq:c-recurrence}; (c) the Catalan/residue proof of \cref{p3:thm:c-finite}; (d) the
implicit construction of $U$ and \cref{p3:eq:c-recurrence-core}; (e) the triangular multivariate
inversion \cref{p3:eq:triangular}; (f) residual-to-root transfer, \cref{p3:prop:certificate},
under an explicit derivative lower bound. Steps (b), (d), (e) are Bell-polynomial infrastructure
shared with \cref{p0:sec:top}.
\end{remark}

\section{Summary}\label{p3:sec:summary}

\begin{enumerate}
\item In the shifted phase $\mu=\pi\sqrt{2(n-1/24)/3}$, Rademacher's series is the exact
sector sum \cref{p3:eq:sector-sum}, each sector having a two-term amplitude linear in $\mu^{-1}$.
At $(k,\sigma)=(1,+1)$ it is the exponential--power model \cref{p0:def:model} with $\alpha=1/2$,
reduced by \cref{p0:lem:alpha-reduction} to $\kappa\e^{\mu}\mu^{-2}(1-\mu^{-1})$
(\cref{p3:lem:alpha-reduction,p3:tab:dictionary}).
\item Re-expansion in $n^{-1/2}$ is governed by one analytic function \cref{p3:eq:Omega};
its coefficients have the finite closed form \cref{p3:eq:c-finite} in \emph{every} sector, which
enters only through $\beta=\sigma\pi/(6k)$. Each such series converges while the total is only
asymptotic --- confusing the two is the characteristic error of the subject.
\item The actions $1\mp1/k$ accumulate at $1$ from both sides, so the completed object is
not a Hahn series and the two signs may not be separated
(\cref{p3:def:rademacher-transseries,p3:rem:pairing}); \cref{p3:thm:tail} bounds the tail
explicitly and \cref{p3:thm:zeta-condensation} packages the accumulation exactly.
\item The inverse carrier is $-2\Wlam_{-1}(-\tfrac12\sqrt{\kappa/y})$ with Stirling
flattening polynomials \cref{p3:eq:Pi-lagrange}; the amplitude is restored by
\cref{p0:thm:lambert-centered} at $\Psi(t,u)=3\log(1+tu)-\log(1-t+tu)$, whose coefficients are
\cref{p3:tab:cm}.
\item Since $\dd n/\dd\mu\asymp\mu$, that reversion produces an $\bigO(1)$ index
shift $3/\pi^{2}$ which must be resummed \emph{before} the periodic amplitudes are evaluated
(\cref{p3:rem:phase-lock}).
\item With the phase frozen --- by a residue class mod $\operatorname{lcm}(1,\dots,K)$ or
by the entire interpolation \cref{p3:def:Atilde}, which agree strictly below action $1$
(\cref{p3:prop:two-conventions}) --- every inverse transmonomial is given by
\cref{p3:eq:multi-LB}, \cref{p3:eq:multi-Bell}, \cref{p3:eq:phase-LB} or \cref{p3:eq:triangular}.
The leading correction is the parity chirp \cref{p3:eq:chirp} with constant $3^{1/4}/(2\pi)$, and
the primary sectors form the staircase \cref{p3:eq:scales}.
\item Along the range the core alone recovers $n$ by rounding
(\cref{p3:thm:rounding}); for general $y$ the staircase is the ceiling of the smooth inverse
(\cref{p3:thm:staircase}), whose order-one sawtooth dominates every arithmetic sector
(\cref{p3:rem:sawtooth}).
\end{enumerate}

\begin{remark}[Works referred to]\label{p3:rem:references}
Analytic input: Rademacher, Proc.\ Nat.\ Acad.\ Sci.\ USA \textbf{23} (1937), 78--84; Proc.\
London Math.\ Soc.\ (2) \textbf{43} (1938), 241--254; Ann.\ of Math.\ \textbf{44} (1943),
416--422; textbook treatment in Apostol, \emph{Modular Functions and Dirichlet Series in Number
Theory}, 2nd ed., Springer GTM 41, 1990; machine-checked in Eberl and Paulson, \emph{Rademacher's
series for the partition function}, Archive of Formal Proofs, 2026. Leading asymptotic: Hardy and
Ramanujan, Proc.\ London Math.\ Soc.\ (2) \textbf{17} (1918), 75--115. Remainder estimates:
Lehmer, J.\ London Math.\ Soc.\ \textbf{12} (1937), 171--176 and Trans.\ Amer.\ Math.\ Soc.\
\textbf{43} (1938), 271--295, \textbf{46} (1939), 362--373; Banerjee, Paule, Radu and Schneider,
arXiv:2209.07887. The dominant-sector case of \cref{p3:eq:c-finite} is O'Sullivan's,
arXiv:2205.13468 and arXiv:2203.02868. Lambert $W$: Corless, Gonnet, Hare, Jeffrey and Knuth,
Adv.\ Comput.\ Math.\ \textbf{5} (1996), 329--359; Jeffrey, Corless, Hare and Knuth,
arXiv:math/9512230. The sequence is OEIS A000041.
\end{remark}

%% ==================== fragment p4 ====================
\chapter{The swing factorial (A056040)}
\label{p4:sec:top}

\begin{sourcebox}
This chapter merges three independent treatments of the swinging factorial:
\emph{All-Orders Asymptotic Transseries for the Swing Factorial and Its Two
Inverse Branches} (cited below as \textsf{S1}), \emph{The Full Asymptotic
Transseries of the Swing Factorial and Its Inverse} (\textsf{S2}), and
\emph{Parity-Resolved Transseries for the Swing Factorial and Its Inverse}
(\textsf{S3}).  The general apparatus of \textsf{S2} --- its universal
coefficient calculus for exponential--power inversion and its
Lambert-normalized reversion theorem --- was promoted to \cref{p0:sec:top};
here it is only \emph{cited} and \emph{specialized}.  \textsf{S1} works in the
half index $m$ with the carrier constant $\log4$; \textsf{S2} and \textsf{S3}
work in the original index with $\log2$ and differ only in the sign convention
for the parity marker.  The chapter adopts the \textsf{S2}/\cref{p0:sec:top}
normalization throughout and gives the translation dictionary in
\cref{p4:rem:dictionary}.  \textsf{S1}'s genuinely distinct contribution --- the
centred (translation-symmetric) normal form --- is kept, but recast so that it
too is a single $\sigma$-uniform statement.

Every rational, polynomial and algebraic coefficient displayed in this chapter
was recomputed from the stated recurrences in exact arithmetic
(\texttt{sympy}, rational Bernoulli input, truncated power series over
$\Rat[\lambda,\lambda^{-1},\sigma]/(\sigma^2-1)$); every decimal in the
validation tables was recomputed at $50$ significant digits.  Where a
recomputation reproduced a source table, that is said; where it did not, the
discrepancy is recorded in a \textsf{repairbox}.
\end{sourcebox}

\section{The sequence, the parity obstruction, and what ``inverse'' can mean}
\label{p4:sec:elementary}

\subsection{Definition and exact parity identities}

\begin{definition}[Swing factorial]
\label{p4:def:swing}
For $n\in\Nat_0$ put
\begin{equation}
  \operatorname{sw}(n)=\frac{n!}{\lfloor n/2\rfloor!^{\,2}}.
  \label{p4:eq:swing-def}
\end{equation}
This is OEIS A056040, the \emph{swinging factorial}.
\end{definition}

Splitting on the parity of $n$ gives at once
\begin{equation}
  \operatorname{sw}(2m)=\binom{2m}{m},
  \qquad
  \operatorname{sw}(2m+1)=(2m+1)\binom{2m}{m},
  \label{p4:eq:bisections}
\end{equation}
whence the one-step recurrence
\begin{equation}
  \operatorname{sw}(0)=1,
  \qquad
  \operatorname{sw}(n)=
  \begin{cases}
    n\,\operatorname{sw}(n-1), & n \text{ odd},\\[2pt]
    \dfrac4n\,\operatorname{sw}(n-1), & n \text{ even},
  \end{cases}
  \label{p4:eq:one-step}
\end{equation}
and the closed product form
\begin{equation}
  \operatorname{sw}(n)=2^{\,n-(n\bmod 2)}\prod_{k=1}^{n}k^{(-1)^{k+1}}.
  \label{p4:eq:swing-product}
\end{equation}
The first values are
\[
  1,\,1,\,2,\,6,\,6,\,30,\,20,\,140,\,70,\,630,\,252,\,2772,\,924,\,12012,
  \,3432,\,51480,\,\ldots
\]

\begin{proposition}[Generating functions]
\label{p4:prop:gf}
\begin{align}
  \sum_{n\ge0}\operatorname{sw}(n)\frac{z^n}{n!}&=(1+z)\,I_0(2z),
  \label{p4:eq:egf}\\
  \sum_{n\ge0}\operatorname{sw}(n)z^{n}
  &=\frac1{\sqrt{1-4z^2}}+\frac{z}{(1-4z^2)^{3/2}},
  \label{p4:eq:ogf}
\end{align}
where $I_0$ is the modified Bessel function of order zero.
\end{proposition}

\begin{proof}
By \cref{p4:eq:bisections} the even exponential part is
$\sum_{m\ge0}z^{2m}/(m!)^2=I_0(2z)$ and the odd exponential part is
$z\,I_0(2z)$, which is \cref{p4:eq:egf}.  For the ordinary series put
$F(v)=\sum_{m\ge0}\binom{2m}{m}v^m=(1-4v)^{-1/2}$; then the even part is
$F(z^2)$ and the odd part is
$z\bigl(2vF'(v)+F(v)\bigr)\big|_{v=z^2}=z(1-4z^2)^{-3/2}$.
\end{proof}

\begin{remark}[The generating function already shows the obstruction]
\label{p4:rem:codominance}
The two singularities $z=\pm\tfrac12$ of \cref{p4:eq:ogf} have equal modulus.
Co-dominant singularities of equal modulus at $\pm\rho$ produce a period-two
oscillation in the coefficient asymptotics, and no single non-oscillatory
singular expansion can represent the whole sequence.  Even/odd projection
separates the two contributions exactly.  This is the generating-function
shadow of everything that follows.
\end{remark}

\subsection{The sequence is not eventually monotone}

\begin{proposition}[Exact adjacent ratios]
\label{p4:prop:ratios}
For $m\ge0$,
\begin{equation}
  \frac{\operatorname{sw}(2m+1)}{\operatorname{sw}(2m)}=2m+1,
  \qquad
  \frac{\operatorname{sw}(2m+2)}{\operatorname{sw}(2m+1)}=\frac{2}{m+1}.
  \label{p4:eq:ratios}
\end{equation}
Consequently $\operatorname{sw}$ rises strictly from every even index to the
next odd index, and falls strictly from every odd index to the next even index
as soon as $m\ge2$.  In particular $\operatorname{sw}$ is not eventually
monotone and has no eventually defined single-valued inverse.
\end{proposition}

\begin{proof}
Both ratios are \cref{p4:eq:one-step}.  The second is $<1$ exactly when
$m+1>2$.
\end{proof}

\begin{warning}[There is no ``the'' inverse swing factorial]
\label{p4:rem:two-sheets}
Any formula presented as \emph{the} inverse of $\operatorname{sw}$ without
naming a parity sheet or a generalized-inverse convention is incomplete: of
the three inverse objects of \cref{p0:def:three-inverses}, none exists
globally for $\operatorname{sw}$ and all three exist branchwise
(\cref{p4:sec:discrete}).
\end{warning}

\subsection{The two analytic branches}

All three sources resolve the obstruction the same way: interpolate the two
bisections separately by gamma quotients.  We use the parity marker
$\sigma\in\{-1,+1\}$ of \textsf{S2}, with
\begin{equation}
  \sigma=-1 \ \text{(even phase)},
  \qquad
  \sigma=+1 \ \text{(odd phase)},
  \qquad
  \sigma_n:=(-1)^{n+1},
  \qquad
  \lambda:=\log2 .
  \label{p4:eq:parity-marker}
\end{equation}
The coefficient ring of the whole chapter is the \emph{parity algebra}
\begin{equation}
  \mathscr R:=\Rat[\lambda,\lambda^{-1},\sigma]\big/(\sigma^2-1),
  \label{p4:eq:parity-algebra}
\end{equation}
in which every formula below is a single element carrying both sheets.  This
is not a notational economy only: $\sigma_n$ is a discrete transmonomial of
modulus one, and it must remain visible at every order of every expansion.

\begin{definition}[Parity-resolved gamma branches]
\label{p4:def:branches}
For $x>0$ and $\sigma\in\{-1,+1\}$ set
\begin{equation}
  \mathcal S_\sigma(x)
  :=\frac{\Gamma(x+1)}{\Gamma\!\bigl(x/2+(3-\sigma)/4\bigr)^{2}} .
  \label{p4:eq:branch-def}
\end{equation}
Explicitly $\mathcal S_-(x)=\Gamma(x+1)/\Gamma(x/2+1)^2$ and
$\mathcal S_+(x)=\Gamma(x+1)/\Gamma\bigl((x+1)/2\bigr)^2$, and
\begin{equation}
  \mathcal S_-(2m)=\operatorname{sw}(2m),
  \qquad
  \mathcal S_+(2m+1)=\operatorname{sw}(2m+1)
  \qquad(m\in\Nat_0).
  \label{p4:eq:branch-interpolates}
\end{equation}
\end{definition}

\begin{proposition}[Duplication normal form and duality]
\label{p4:prop:duplication}
For $x>0$,
\begin{equation}
  \mathcal S_\sigma(x)
  =\frac{2^{x}}{\sqrt\pi}
   \left(\frac{\Gamma(x/2+1)}{\Gamma(x/2+1/2)}\right)^{\!\sigma},
  \label{p4:eq:branch-ratio}
\end{equation}
and consequently
\begin{equation}
  \mathcal S_+(x)\,\mathcal S_-(x)=\frac{4^{x}}{\pi},
  \qquad
  \frac{\mathcal S_\sigma(x+2)}{\mathcal S_\sigma(x)}
  =\frac{16(x+1)(x+2)}{(2x+3-\sigma)^2}
  =\begin{cases}
     \dfrac{4(x+1)}{x+2},&\sigma=-1,\\[8pt]
     \dfrac{4(x+2)}{x+1},&\sigma=+1.
   \end{cases}
  \label{p4:eq:duality-and-step}
\end{equation}
\end{proposition}

\begin{proof}
Legendre duplication in the form
$\Gamma(x+1)=2^{x}\pi^{-1/2}\Gamma\bigl((x+1)/2\bigr)\Gamma(x/2+1)$ turns
\cref{p4:eq:branch-def} into \cref{p4:eq:branch-ratio}; multiplying the two
values of $\sigma$ makes the gamma ratio cancel and leaves $4^x/\pi$.  For the
two-step ratio, $\Gamma(x+3)/\Gamma(x+1)=(x+1)(x+2)$ while the denominator
contributes $\bigl(x/2+(3-\sigma)/4\bigr)^{-2}=16(2x+3-\sigma)^{-2}$.
\end{proof}

\begin{proposition}[Strict monotonicity and existence of the real inverses]
\label{p4:prop:monotone}
For $x>0$,
\begin{equation}
  \frac{\dd}{\dd x}\log\mathcal S_\sigma(x)
  =\lambda+\frac\sigma2
   \Bigl[\psi\bigl(\tfrac x2+1\bigr)-\psi\bigl(\tfrac x2+\tfrac12\bigr)\Bigr],
  \label{p4:eq:logder}
\end{equation}
where $\psi=\Gamma'/\Gamma$.  The bracket decreases strictly from $2\lambda$ at
$x=0$ to $0$ at $x=\infty$.  Hence $\mathcal S_+$ has strictly positive
derivative on $[0,\infty)$, $\mathcal S_-$ has derivative $0$ at $x=0$ and
strictly positive derivative on $(0,\infty)$, and both are strictly increasing
bijections of $[0,\infty)$ onto $[1,\infty)$.  Write
\begin{equation}
  \mathcal I_\sigma:=\mathcal S_\sigma^{-1}:[1,\infty)\to[0,\infty).
  \label{p4:eq:inverse-def}
\end{equation}
\end{proposition}

\begin{proof}
Differentiating \cref{p4:eq:branch-ratio} logarithmically gives
\cref{p4:eq:logder}.  Put
$g(x)=\psi(x/2+1)-\psi(x/2+1/2)$.  Then
$g'(x)=\tfrac12\bigl[\psi'(x/2+1)-\psi'(x/2+1/2)\bigr]<0$ because the
trigamma function $\psi'(t)=\sum_{k\ge0}(t+k)^{-2}$ is strictly decreasing on
$(0,\infty)$; and $g(x)\to0$ as $x\to\infty$ by
$\psi(t+\tfrac12)-\psi(t)\to0$.  Finally $g(0)=\psi(1)-\psi(1/2)=2\log2
=2\lambda$.  Thus $0<g(x)<2\lambda$ for $x>0$.  For $\sigma=+1$,
\cref{p4:eq:logder} equals $\lambda+g/2>0$ on $[0,\infty)$; for $\sigma=-1$ it
equals $\lambda-g/2$, which vanishes only at $x=0$ and is positive afterwards.
Both branches take the value $1$ at $x=0$ and tend to $+\infty$, so each is a
continuous increasing bijection onto $[1,\infty)$.
\end{proof}

\subsection{Why one cannot simply interpolate through both parities}

\begin{proposition}[Every smooth interpolation of the whole sequence oscillates]
\label{p4:prop:oscillatory}
With $R(x)=\Gamma(x/2+1)/\Gamma\bigl((x+1)/2\bigr)$, put
\begin{equation}
  \widehat{\mathcal S}(x)=\frac{2^{x}}{\sqrt\pi}\,R(x)^{-\cos\pi x} .
  \label{p4:eq:oscillatory-interps}
\end{equation}
Then $\widehat{\mathcal S}(n)=\operatorname{sw}(n)$ for every $n\in\Nat_0$,
yet $\widehat{\mathcal S}$ is not eventually monotone: its logarithmic
derivative
\begin{equation}
  \frac{\dd}{\dd x}\log\widehat{\mathcal S}(x)
  =\lambda+\pi\sin(\pi x)\log R(x)-\cos(\pi x)\,\frac{R'(x)}{R(x)}
  \label{p4:eq:oscillatory-derivative}
\end{equation}
has an oscillating term of amplitude $\sim\tfrac\pi2\log x$ against remaining
terms that are $O(1)$, so it changes sign infinitely often.  The same holds
for the convex interpolation
$\tfrac{1+\cos\pi x}2\mathcal S_-+\tfrac{1-\cos\pi x}2\mathcal S_+$.
Consequently the phase-resolved \emph{pair} $(\mathcal S_-,\mathcal S_+)$, and
not any single function, is the asymptotically stable continuation of
$\operatorname{sw}$; the corresponding inverse object is the pair
$(\mathcal I_-,\mathcal I_+)$.
\end{proposition}

\begin{proof}
At $x=n$ one has $-\cos\pi n=(-1)^{n+1}=\sigma_n$, so
$\widehat{\mathcal S}(n)=(2^{n}/\sqrt\pi)R(n)^{\sigma_n}
=\mathcal S_{\sigma_n}(n)=\operatorname{sw}(n)$ by
\cref{p4:eq:branch-ratio} and \cref{p4:eq:branch-interpolates}.
Differentiating
$\log\widehat{\mathcal S}=x\lambda-\tfrac12\log\pi-\cos(\pi x)\log R(x)$ gives
\cref{p4:eq:oscillatory-derivative}; by \cref{p4:cor:R-asymptotic} below,
$\log R(x)=\tfrac12\log(x/2)+O(x^{-1})$ and $R'/R=O(x^{-1})$, so the
oscillating term has amplitude $\tfrac\pi2\log x+O(1)\to\infty$ and dominates.
The convex interpolation is treated the same way after writing it as
$\mathcal S_-(x)$ times a factor whose logarithmic derivative oscillates with
amplitude $\asymp\log x$.
\end{proof}

\section{Exact normal form, Binet reduction, and Borel structure}
\label{p4:sec:normalform}

\subsection{The single scalar correction}

\begin{definition}[The gamma-ratio correction]
\label{p4:def:D}
For $x>0$ set
\begin{equation}
  D(x):=\log\frac{\Gamma(x/2+1)}{\Gamma(x/2+1/2)}-\frac12\log\frac x2 .
  \label{p4:eq:D-def}
\end{equation}
\end{definition}

\begin{theorem}[Exact $\sigma$-uniform normal form]
\label{p4:thm:normal-form}
For $x>0$ and $\sigma\in\{-1,+1\}$,
\begin{equation}
  \boxed{\;
  \mathcal S_\sigma(x)=C_\sigma\,\e^{\lambda x}\,x^{\sigma/2}\,
  \e^{\sigma D(x)},
  \qquad
  C_\sigma=\frac{2^{-\sigma/2}}{\sqrt\pi}.\;}
  \label{p4:eq:normal-form}
\end{equation}
At integer arguments this is the exact parity formula
\begin{equation}
  \log\operatorname{sw}(n)
  =n\lambda-\tfrac12\log\pi
  +\sigma_n\Bigl(\tfrac12\log\tfrac n2+D(n)\Bigr),
  \qquad n\ge1 .
  \label{p4:eq:integer-exact}
\end{equation}
\end{theorem}

\begin{proof}
By \cref{p4:eq:D-def}, $\Gamma(x/2+1)/\Gamma(x/2+1/2)=(x/2)^{1/2}\e^{D(x)}$.
Substituting into \cref{p4:eq:branch-ratio} gives
$\mathcal S_\sigma=(2^x/\sqrt\pi)(x/2)^{\sigma/2}\e^{\sigma D}$, and
$(x/2)^{\sigma/2}=2^{-\sigma/2}x^{\sigma/2}$.  For
\cref{p4:eq:integer-exact}, take logarithms of \cref{p4:eq:normal-form} at
$x=n$, $\sigma=\sigma_n$, and use
$\tfrac12\log n-\tfrac\lambda2=\tfrac12\log\tfrac n2$.
\end{proof}

\Cref{p4:eq:normal-form} is exactly \cref{p0:def:model} with
\begin{equation}
  \boxed{\;
  \alpha=1,\qquad a=\lambda=\log 2,\qquad b=\frac\sigma2,\qquad
  C=C_\sigma=\frac{2^{-\sigma/2}}{\sqrt\pi},\qquad
  q_j=\sigma d_j,\;}
  \label{p4:eq:dictionary}
\end{equation}
where $d_j$ are the coefficients of the asymptotic series of $D$
(\cref{p4:thm:logcoeffs}).  Because $\alpha=1$ no monomial substitution is
needed and \cref{p0:lem:alpha-reduction} is not invoked.  We record the
dictionary once and use it in every specialization below.

\begin{remark}[Dictionary between the three sources]
\label{p4:rem:dictionary}
Let $\varepsilon_{\textsf{S3}}\in\{+1,-1\}$ be the marker of \textsf{S3} and
$\epsilon_{\textsf{S1}}\in\{0,1\}$ that of \textsf{S1}.  Then
\[
  \sigma=-\varepsilon_{\textsf{S3}},
  \qquad
  \sigma=2\epsilon_{\textsf{S1}}-1,
  \qquad
  \textsf{S3}:\ Q=-D,\quad q^{\textsf{S3}}_r=-d_{2r-1},
  \qquad
  \textsf{S1}:\ \alpha_{\textsf{S1}}=\log4=2\lambda .
\]
\textsf{S1} parametrizes each sheet by the half index $m$ with
$S^{\textsf{S1}}_{\epsilon}(m)=\operatorname{sw}(2m+\epsilon)$; as functions
this is the reparametrization $x=2m+\epsilon$ of \cref{p4:def:branches}, so
\[
  S^{\textsf{S1}}_{\epsilon}(m)=\mathcal S_{2\epsilon-1}(2m+\epsilon),
  \qquad
  \mathcal I_{\sigma}(y)=2M^{\textsf{S1}}_{\epsilon}(y)+\epsilon .
\]
The reparametrization is invisible in the forward coefficients up to the
substitution $x\mapsto2m+\epsilon$, but it is \emph{not} invisible in the
inverse coefficients: the two normalizations invert different leading phases
($\lambda x+\tfrac\sigma2\log x$ against
$2\lambda m+\tfrac\sigma2\log m$) and therefore have different Lambert cores.
Only for the even sheet do the two cores agree exactly, through
$\Xcore_-(y)=2w^{\textsf{S1}}_0(y)$; see \cref{p4:rem:halfindex-cores}.
\end{remark}

\subsection{The Binet kernel collapses to a hyperbolic tangent}

\begin{theorem}[Exact Borel--Laplace representation]
\label{p4:thm:laplace}
For $\Re x>0$,
\begin{equation}
  \boxed{\;
  D(x)=\int_0^\infty\e^{-xt}\,\widehat D(t)\,\dd t,
  \qquad
  \widehat D(t)=\frac{\tanh(t/2)}{2t},\;}
  \label{p4:eq:laplace}
\end{equation}
the singularity of $\widehat D$ at $t=0$ being removable with
$\widehat D(0)=\tfrac14$.
\end{theorem}

\begin{proof}
Put $z=x/2$ and
$h(z)=\log\Gamma(z+1)-\log\Gamma(z+\tfrac12)-\tfrac12\log z$, so
$h(x/2)=D(x)$.  From the standard integral for a digamma difference,
\[
  \psi(z+1)-\psi(z+\tfrac12)
  =\int_0^\infty\frac{\e^{-(z+1/2)s}-\e^{-(z+1)s}}{1-\e^{-s}}\,\dd s
  =2\int_0^\infty\frac{\e^{-2zv}}{\e^{v}+1}\,\dd v,
\]
the second equality by $s=2v$ and
$(\e^{-v}-\e^{-2v})/(1-\e^{-2v})=1/(\e^{v}+1)$.  Since
$1/(2z)=\int_0^\infty\e^{-2zv}\dd v$,
\[
  h'(z)=\psi(z+1)-\psi(z+\tfrac12)-\frac1{2z}
  =\int_0^\infty\e^{-2zv}\Bigl(\frac2{\e^{v}+1}-1\Bigr)\dd v
  =-\int_0^\infty\e^{-2zv}\tanh\frac v2\,\dd v .
\]
Both $h(z)$ and the right-hand side of \cref{p4:eq:laplace} tend to $0$ as
$\Re z\to+\infty$, so integrating from $+\infty$ down to $z$,
\[
  h(z)=\int_z^\infty\!\!\int_0^\infty\e^{-2\zeta v}\tanh\tfrac v2
        \,\dd v\,\dd\zeta
      =\int_0^\infty\e^{-2zv}\frac{\tanh(v/2)}{2v}\,\dd v ,
\]
which is \cref{p4:eq:laplace} after $x=2z$.  Fubini applies because the double
integral is absolutely convergent for $\Re z>0$.
\end{proof}

\begin{sourcebox}
\textsf{S3} reaches \cref{p4:eq:laplace} from Binet's first formula, through
the elegant identity
$b(t)-2b(2t)=-\tfrac12\tanh(t/2)$ for
$b(t)=(\e^t-1)^{-1}-t^{-1}+\tfrac12$; \textsf{S2} reaches it from the digamma
integral, as above.  \textsf{S1} keeps Binet's \emph{second} formula
(the $\arctan$ kernel) instead and never simplifies the difference of the two
Binet integrals.  The $\tanh$ collapse is the decisive simplification: it is
what makes the Borel plane elementary, and it is adopted here.
\end{sourcebox}

\begin{corollary}[Sign, monotonicity, and elementary bounds]
\label{p4:cor:D-bounds}
For $x>0$,
\begin{equation}
  0<D(x)<\frac1{4x},
  \qquad
  0<-D'(x)<\frac1{4x^{2}},
  \qquad
  (-1)^{k}D^{(k)}(x)>0\quad(k\ge0).
  \label{p4:eq:D-bounds}
\end{equation}
In particular $D$ is completely monotone and strictly decreasing.
\end{corollary}

\begin{proof}
$\widehat D>0$ on $(0,\infty)$, so
$(-1)^kD^{(k)}(x)=\int_0^\infty\e^{-xt}t^k\widehat D(t)\dd t>0$.  From
$0<\tanh(t/2)<t/2$ we get $0<\widehat D(t)<\tfrac14$, and integrating
$\e^{-xt}t^k$ gives the two displayed bounds.
\end{proof}

\begin{corollary}[Asymptotics of the gamma ratio]
\label{p4:cor:R-asymptotic}
With $R(x)=\Gamma(x/2+1)/\Gamma((x+1)/2)$ as in
\cref{p4:eq:oscillatory-interps},
\[
  \log R(x)=\tfrac12\log(x/2)+D(x)=\tfrac12\log(x/2)+O(x^{-1}),
  \qquad
  \frac{R'(x)}{R(x)}=\frac1{2x}+D'(x)=O(x^{-1}) .
\]
\end{corollary}

\begin{proof}
Immediate from \cref{p4:eq:D-def} and \cref{p4:eq:D-bounds}.
\end{proof}

\subsection{The Borel singularity lattice}

\begin{proposition}[Mittag--Leffler expansion and poles]
\label{p4:prop:borel-lattice}
The Borel transform of \cref{p4:eq:laplace} satisfies
\begin{equation}
  \widehat D(t)=2\sum_{k\ge0}\frac{1}{t^{2}+\pi^{2}(2k+1)^{2}} ,
  \label{p4:eq:partial-fractions}
\end{equation}
is even and meromorphic on $\Cplx$, and has exactly the simple poles
\begin{equation}
  t=\pm\ii\pi(2k+1),\qquad k\ge0,
  \qquad
  \Res_{t=t_k}\widehat D(t)=\frac1{t_k} .
  \label{p4:eq:borel-poles}
\end{equation}
The positive real ray is free of singularities; the nearest singularities lie
at distance $\pi$, so the formal series of $D$ is Gevrey-$1$ with singulant
modulus $\pi$.
\end{proposition}

\begin{proof}
$\tanh(w)=\sum_{k\ge0}2w/\bigl(w^2+\pi^2(k+\tfrac12)^2\bigr)$ is the classical
Mittag--Leffler expansion; putting $w=t/2$ and dividing by $2t$ gives
\cref{p4:eq:partial-fractions}, both sides being even, meromorphic, with the
same poles and residues and the same value $\tfrac14$ at $t=0$ (using
$\sum_{k\ge0}(2k+1)^{-2}=\pi^2/8$).  At $t=t_k$ the function $\tanh(t/2)$ has
residue $2$, so $\widehat D$ has residue $2/(2t_k)=1/t_k$.
\end{proof}

\begin{remark}[What the exponentially small scale is, and is not]
\label{p4:rem:no-real-ambiguity}
The quantity $\e^{-\pi x}$ that governs optimal truncation
(\cref{p4:sec:stokes}) is the envelope set by the distance $\pi$ of the
nearest \emph{off-axis} Borel poles; it is not an ambiguity of the
positive-axis sum, since the ray $\arg t=0$ meets no singularity and
\cref{p4:eq:laplace} is exact.  Genuine Stokes jumps appear only when the
Laplace direction is rotated onto the imaginary rays.  \textsf{S1}, retaining
an unsimplified Binet integral, makes the weaker but correct statement that no
Stokes constant need be assigned on the real axis; the Borel picture explains
why.
\end{remark}

\section{The logarithmic coefficients and a sharp enveloping remainder}
\label{p4:sec:logcoeffs}

\subsection{All coefficients}

We use the Bernoulli convention $z/(\e^z-1)=\sum_{n\ge0}B_nz^n/n!$, so
$B_1=-\tfrac12$, and the Bernoulli polynomials
$t\e^{wt}/(\e^t-1)=\sum_n B_n(w)t^n/n!$.

\begin{theorem}[All logarithmic coefficients]
\label{p4:thm:logcoeffs}
As $x\to\infty$ in any closed subsector of $|\arg x|<\pi/2$,
\begin{equation}
  D(x)\sim\sum_{j\ge1}\frac{d_j}{x^{j}},
  \qquad
  d_{2r}=0,
  \qquad
  d_{2r+1}=\frac{(2^{2r+2}-1)B_{2r+2}}{(2r+1)(2r+2)} .
  \label{p4:eq:d-bernoulli}
\end{equation}
Equivalently
\begin{equation}
  d_{2r+1}=(-1)^{r}\,\frac{2\,(2r)!}{\pi^{2r+2}}\bigl(1-2^{-2r-2}\bigr)
  \zeta(2r+2)
  =(-1)^{r}\,\frac{2\,(2r)!}{\pi^{2r+2}}
   \sum_{k\ge0}\frac{1}{(2k+1)^{2r+2}} ,
  \label{p4:eq:d-zeta}
\end{equation}
and the formal Borel transform of the series is exactly the Taylor expansion
of $\widehat D$:
\begin{equation}
  \widehat D(t)=\sum_{r\ge0}d_{2r+1}\frac{t^{2r}}{(2r)!}
  =\frac{\tanh(t/2)}{2t}.
  \label{p4:eq:borel-taylor}
\end{equation}
\end{theorem}

\begin{proof}
Insert into \cref{p4:eq:partial-fractions} the finite geometric identity
\[
  \frac1{t^2+a^2}=\sum_{r=0}^{N-1}\frac{(-1)^{r}t^{2r}}{a^{2r+2}}
   +\frac{(-1)^{N}t^{2N}}{a^{2N}(t^2+a^2)},
  \qquad a=a_k=\pi(2k+1),
\]
and integrate the finite sum against $\e^{-xt}$ using
$\int_0^\infty\e^{-xt}t^{2r}\dd t=(2r)!x^{-2r-1}$.  This gives the second
expression in \cref{p4:eq:d-zeta}, hence the first by
$\zeta(2r+2)(1-2^{-2r-2})=\sum_{k\ge0}(2k+1)^{-(2r+2)}$, and hence
\cref{p4:eq:d-bernoulli} by Euler's evaluation
$B_{2m}=(-1)^{m-1}2(2m)!\zeta(2m)/(2\pi)^{2m}$.  The unexpanded last term is
the remainder controlled in \cref{p4:thm:envelope}, which in particular proves
the asymptotic statement on $\Re x>0$; the identity \cref{p4:eq:borel-taylor}
is the Taylor expansion of $\widehat D$ at the origin.
\end{proof}

\begin{remark}[Sector of validity: a disagreement between the sources]
\label{p4:rem:sector}
\textsf{S2} and \textsf{S3} state \cref{p4:eq:d-bernoulli} on closed
subsectors of $|\arg x|<\pi/2$, which is what the Laplace representation
\cref{p4:eq:laplace} gives.  \textsf{S1} states the corresponding expansion on
$|\arg m|\le\pi-\delta$.  Both are correct: the wider sector follows not from
the Laplace integral but from the Tricomi--Erd\'elyi expansion of a gamma
ratio,
\begin{equation}
  \log\frac{\Gamma(z+a)}{\Gamma(z+b)}
  \sim(a-b)\log z+\sum_{n\ge1}\frac{(-1)^{n+1}}{n(n+1)z^{n}}
     \bigl(B_{n+1}(a)-B_{n+1}(b)\bigr),
  \qquad |\arg z|\le\pi-\delta,
  \label{p4:eq:tricomi}
\end{equation}
applied with $a=1$, $b=\tfrac12$, $z=x/2$.  The volume states the Laplace
sector where an \emph{exact} representation and a signed remainder are wanted,
and the wider sector where only a Poincar\'e statement is wanted.  The two
routes give the same coefficients: with $B_m(\tfrac12)=(2^{1-m}-1)B_m$ and
$B_m(1)=B_m$ for $m\ne1$, \cref{p4:eq:tricomi} yields
$d_n=2^{n}\,\frac{(-1)^{n+1}}{n(n+1)}\bigl(2-2^{-n}\bigr)B_{n+1}$, which is
\cref{p4:eq:d-bernoulli}.
\end{remark}

\begin{table}[htbp]
\centering
\caption{The logarithmic coefficients $d_{2r+1}$ of \cref{p4:thm:logcoeffs},
recomputed from \cref{p4:eq:d-bernoulli} in exact rational arithmetic and
cross-checked against the Taylor expansion of $\tanh(t/2)/(2t)$.}
\label{p4:tab:d}
\small
\begin{tabular}{@{}rrrrrrrr@{}}
\toprule
$j$ & $1$ & $3$ & $5$ & $7$ & $9$ & $11$ & $13$\\
\midrule
$d_j$ & $\dfrac14$ & $-\dfrac1{24}$ & $\dfrac1{20}$ & $-\dfrac{17}{112}$
      & $\dfrac{31}{36}$ & $-\dfrac{691}{88}$ & $\dfrac{5461}{52}$\\
\bottomrule
\end{tabular}
\end{table}

Thus
\begin{equation}
  D(x)\sim\frac1{4x}-\frac1{24x^{3}}+\frac1{20x^{5}}
  -\frac{17}{112x^{7}}+\frac{31}{36x^{9}}
  -\frac{691}{88x^{11}}+\frac{5461}{52x^{13}}-\cdots
  \label{p4:eq:D-first}
\end{equation}

\subsection{An enveloping remainder, not merely a Poincar\'e one}

\begin{theorem}[Exact remainder and strict next-term envelope]
\label{p4:thm:envelope}
For $x>0$ and $N\ge0$ write
$D(x)=\sum_{r=0}^{N-1}d_{2r+1}x^{-(2r+1)}+R_N(x)$ (empty sum for $N=0$).
Then, with $a_k=\pi(2k+1)$,
\begin{equation}
  R_N(x)=2(-1)^{N}\sum_{k\ge0}a_k^{-2N}
  \int_0^\infty\frac{\e^{-xt}t^{2N}}{t^{2}+a_k^{2}}\,\dd t,
  \label{p4:eq:remainder-integral}
\end{equation}
and consequently
\begin{equation}
  \boxed{\;
  0<(-1)^{N}R_N(x)<\frac{|d_{2N+1}|}{x^{2N+1}}\;}
  \qquad(x>0,\ N\ge0).
  \label{p4:eq:envelope}
\end{equation}
Every truncation therefore brackets $D(x)$ together with its successor.
\end{theorem}

\begin{proof}
Substitute the finite geometric identity of the previous proof into
\cref{p4:eq:partial-fractions} and then into \cref{p4:eq:laplace}; Tonelli's
theorem justifies the interchange since all terms are positive.  This gives
\cref{p4:eq:remainder-integral}, in which every integral is positive.
Replacing $(t^{2}+a_k^{2})^{-1}$ by the strict upper bound $a_k^{-2}$ and
using $\int_0^\infty\e^{-xt}t^{2N}\dd t=(2N)!x^{-2N-1}$ gives
\[
  0<(-1)^{N}R_N(x)<2(2N)!\,x^{-(2N+1)}\sum_{k\ge0}a_k^{-(2N+2)}
  =\frac{|d_{2N+1}|}{x^{2N+1}},
\]
by the last expression in \cref{p4:eq:d-zeta}.
\end{proof}

\begin{remark}
\Cref{p4:thm:envelope} is strictly stronger than the Poincar\'e statement, and
it is what makes every later analytic claim in this chapter effective: all
remainder constants below are explicit rational multiples of the first omitted
$|d_{2N+1}|$.  \textsf{S2} and \textsf{S3} both prove it; \textsf{S1} does not
have it, because it never simplifies the Binet kernel.  It was verified
numerically at $x=5$ and $x=10$ for $0\le N\le4$: the bracketing
\cref{p4:eq:envelope} holds with the two sides agreeing to within a factor
$1.33$ at $x=5$ and $1.09$ at $x=10$.
\end{remark}

\subsection{Large order and least term}

\begin{proposition}[Least term and its envelope]
\label{p4:prop:leastterm}
As $r\to\infty$, $d_{2r+1}\sim(-1)^{r}2(2r)!\pi^{-2r-2}$, so the ratio of
consecutive nonzero term magnitudes in \cref{p4:eq:D-first} is
$\sim(2r+2)(2r+1)/(\pi^{2}x^{2})$.  The least term therefore occurs at
$r\approx\pi x/2$ and has magnitude
\begin{equation}
  \frac{2\sqrt2}{\pi\sqrt x}\,\e^{-\pi x}\bigl(1+O(x^{-1})\bigr).
  \label{p4:eq:least-term}
\end{equation}
By \cref{p4:thm:envelope}, truncation at the least term attains an error of
exactly this scale, with no hidden constant.
\end{proposition}

\begin{proof}
The asymptotic form of $d_{2r+1}$ is \cref{p4:eq:d-zeta} with
$\zeta(2r+2)\to1$.  The ratio statement is immediate.  Setting $2r=\pi x$ and
applying Stirling's formula to $(2r)!$ gives
$2(2r)!\pi^{-2r-2}x^{-2r-1}\sim2\sqrt{2\pi^{2}x}\,\e^{-\pi x}/(\pi^{2}x)
=2\sqrt2\,\e^{-\pi x}/(\pi\sqrt x)$.
\end{proof}

\begin{remark}[Numerical confirmation]
\label{p4:rem:leastterm-numeric}
Recomputed: at $x=10$ the least nonzero term of \cref{p4:eq:D-first} occurs at
$r=15$ (against $\pi x/2=15.708$) with magnitude
$6.5448\times10^{-15}$, while \cref{p4:eq:least-term} predicts
$6.4659\times10^{-15}$; at $x=20$ the least term occurs at $r=31$
(against $31.416$) with magnitude $1.03859\times10^{-28}$ against a predicted
$1.03837\times10^{-28}$.
\end{remark}

\section{The forward transseries}
\label{p4:sec:forward}

\subsection{Exponentiation}

\begin{theorem}[All forward coefficients]
\label{p4:thm:forward}
Define $A_0=1$ and, for $n\ge1$,
\begin{equation}
  A_n=\Coef{t^n}\exp\Bigl(\sum_{j\ge1}d_jt^{j}\Bigr)
     =\frac1{n!}\,\EB_n\bigl(1!\,d_1,2!\,d_2,\ldots,n!\,d_n\bigr)
     =\sum_{\substack{k_1,k_2,\ldots\ge0\\ \sum_j jk_j=n}}
       \prod_{j\ge1}\frac{d_j^{\,k_j}}{k_j!},
  \label{p4:eq:A-formulae}
\end{equation}
with $\EB_n$ the complete exponential Bell polynomial of
\cref{p0:def:ebell}.  Equivalently and most cheaply,
\begin{equation}
  A_0=1,\qquad
  nA_n=\sum_{j=1}^{n}jd_jA_{n-j}
      =\sum_{\substack{1\le j\le n\\ j\ \mathrm{odd}}}jd_jA_{n-j}.
  \label{p4:eq:A-recurrence}
\end{equation}
Then, for every fixed $N\ge1$ and uniformly on closed subsectors of
$\Re x>0$,
\begin{equation}
  \boxed{\;
  \mathcal S_\sigma(x)=C_\sigma\,2^{x}x^{\sigma/2}
  \Bigl(\sum_{n=0}^{N-1}\frac{\sigma^{n}A_n}{x^{n}}+O_N(x^{-N})\Bigr).\;}
  \label{p4:eq:forward}
\end{equation}
\end{theorem}

\begin{proof}
The three expressions in \cref{p4:eq:A-formulae} are the definition of
$\EB_n$, its partition form, and the expansion of the exponential as a
product; \cref{p4:eq:A-recurrence} follows from $A'(t)=A(t)\sum_jjd_jt^{j-1}$
for $A(t)=\exp\sum_jd_jt^j$, together with $d_{2r}=0$.  The parity of the
$\sigma$-dependence is exact: since $\sum_jd_jt^j$ is an odd formal series,
\[
  \exp\Bigl(-\sum_jd_jt^j\Bigr)=\exp\Bigl(\sum_jd_j(-t)^j\Bigr)
  =\sum_{n\ge0}(-1)^{n}A_nt^{n},
\]
so $\exp(\sigma D)$ has coefficients $\sigma^nA_n$ in either sheet.  For the
remainder, \cref{p4:prop:forward-effective} below gives an explicit bound.
\end{proof}

\begin{proposition}[Effective forward remainder]
\label{p4:prop:forward-effective}
Fix $N\ge1$ and let $x\ge x_N:=\max\{1,\,2N\}$.  Then
\begin{equation}
  \Bigl|\e^{\sigma D(x)}-\sum_{n=0}^{N-1}\frac{\sigma^nA_n}{x^{n}}\Bigr|
  \le \frac{K_N}{x^{N}},
  \qquad
  K_N:=\e^{1/4}\Bigl(|d_{2\lceil N/2\rceil+1}|
       +\sum_{n\ge N}\frac{|A_n|}{x_N^{\,n-N}}\Bigr)<\infty,
  \label{p4:eq:forward-effective}
\end{equation}
where the tail sum converges because $A_n=O(C^{n}n!^{1/2})$ for some $C$.
\end{proposition}

\begin{proof}
Write $D_{<M}(x)=\sum_{r<M}d_{2r+1}x^{-(2r+1)}$ with $M=\lceil N/2\rceil$, so
that by \cref{p4:thm:envelope} $|D(x)-D_{<M}(x)|\le|d_{2M+1}|x^{-(2M+1)}\le
|d_{2M+1}|x^{-N}$ for $2M+1\ge N$.  Since $|D(x)|\le\tfrac14$ and
$|D_{<M}(x)|\le\tfrac14$ for $x\ge1$ by \cref{p4:eq:envelope} with $N=0$,
the mean value theorem gives
$|\e^{\sigma D}-\e^{\sigma D_{<M}}|\le\e^{1/4}|D-D_{<M}|$.  Finally
$\e^{\sigma D_{<M}(x)}$ is an absolutely convergent power series in $1/x$ for
$x\ge x_N$ whose coefficients agree with $\sigma^nA_n$ for $n<N$ up to terms
already accounted for, and whose tail is bounded by
$\sum_{n\ge N}|A_n|x^{-n}$.  Collecting the two bounds gives
\cref{p4:eq:forward-effective}.  The growth bound on $A_n$ follows from
$d_{2r+1}=O((2r)!\pi^{-2r})$ and the Bell-polynomial partition form.
\end{proof}

\begin{repairbox}
\textbf{Source claim.}  \textsf{S1} obtains its multiplicative remainder by
writing ``$e^{R_N(m)}=1+O(m^{-N})$, so formal exponentiation also gives the
claimed analytic remainder''.  \textbf{Defect.}  That sentence exponentiates a
\emph{formal} truncation error and silently assumes that the tail of the
exponentiated series is of the same order; no bound on the coefficients $A_n$
enters, although the series $\sum A_nx^{-n}$ diverges.  \textbf{Repair.}
\Cref{p4:prop:forward-effective} supplies the missing ingredients (a bound on
$D-D_{<M}$ from the enveloping theorem, a uniform bound on the exponential's
Lipschitz constant, and an explicit tail bound), so the constant in
\cref{p4:eq:forward} is effective.
\end{repairbox}

\subsection{Coefficient table and the two sectors}

\begin{table}[htbp]
\centering
\caption{The universal forward coefficients $A_n$ of \cref{p4:eq:A-formulae},
recomputed by \cref{p4:eq:A-recurrence} in exact rational arithmetic.  The
table reproduces the corresponding tables of all three sources
(\textsf{S3} tabulates $C_m=(-1)^{n}A_n$; \textsf{S1} tabulates
$A_n/2^{n}$ in its half index).}
\label{p4:tab:A}
\small
\begin{tabular}{@{}rl rl@{}}
\toprule
$n$ & $A_n$ & $n$ & $A_n$\\
\midrule
$0$ & $1$ & $8$ & $-334477/8388608$\\
$1$ & $1/4$ & $9$ & $28717403/33554432$\\
$2$ & $1/32$ & $10$ & $59697183/268435456$\\
$3$ & $-5/128$ & $11$ & $-8400372435/1073741824$\\
$4$ & $-21/2048$ & $12$ & $-34429291905/17179869184$\\
$5$ & $399/8192$ & $13$ & $7199255611995/68719476736$\\
$6$ & $869/65536$ & $14$ & $14631594576045/549755813888$\\
$7$ & $-39325/262144$ & $15$ & $-4251206967062925/2199023255552$\\
\bottomrule
\end{tabular}
\end{table}

Explicitly, with $C_-=\sqrt{2/\pi}$ and $C_+=1/\sqrt{2\pi}$,
\begin{align}
  \mathcal S_-(x)&\sim\sqrt{\frac2\pi}\,\frac{2^{x}}{\sqrt x}
  \Bigl(1-\frac1{4x}+\frac1{32x^{2}}+\frac5{128x^{3}}
        -\frac{21}{2048x^{4}}-\frac{399}{8192x^{5}}
        +\frac{869}{65536x^{6}}+\cdots\Bigr),
  \label{p4:eq:even-forward}\\
  \mathcal S_+(x)&\sim\frac{2^{x}\sqrt x}{\sqrt{2\pi}}
  \Bigl(1+\frac1{4x}+\frac1{32x^{2}}-\frac5{128x^{3}}
        -\frac{21}{2048x^{4}}+\frac{399}{8192x^{5}}
        +\frac{869}{65536x^{6}}+\cdots\Bigr),
  \label{p4:eq:odd-forward}
\end{align}
and for the integer sequence itself,
\begin{equation}
  \boxed{\;
  \operatorname{sw}(n)\sim
  \frac{2^{n}}{\sqrt\pi}\,2^{-\sigma_n/2}\,n^{\sigma_n/2}
  \sum_{k\ge0}\frac{\sigma_n^{\,k}A_k}{n^{k}},
  \qquad \sigma_n=(-1)^{n+1}. \;}
  \label{p4:eq:integer-forward}
\end{equation}

\begin{remark}[Period-two, not oscillatory error]
\label{p4:rem:transmonomial}
\Cref{p4:eq:integer-forward} is a two-phase (period-two) transseries: the
discrete transmonomial $\sigma_n$ has modulus one, so it is never absorbed
into a remainder, and it changes the algebraic scale from $2^nn^{-1/2}$ to
$2^nn^{+1/2}$.  Equivalently, with the idempotents
$\chi_\sigma(n)=(1+\sigma\sigma_n)/2$,
\[
  \operatorname{sw}(n)\sim\sum_{\sigma=\pm1}\chi_\sigma(n)\,
  C_\sigma 2^{n}n^{\sigma/2}\sum_{k\ge0}\sigma^{k}A_kn^{-k}.
\]
\end{remark}

\begin{corollary}[Central binomial coefficient and the reciprocal]
\label{p4:cor:central-and-reciprocal}
Putting $x=2m$ in \cref{p4:eq:even-forward},
\begin{equation}
  \binom{2m}{m}\sim\frac{4^{m}}{\sqrt{\pi m}}
  \Bigl(1-\frac1{8m}+\frac1{128m^{2}}+\frac5{1024m^{3}}
        -\frac{21}{32768m^{4}}-\frac{399}{262144m^{5}}+\cdots\Bigr),
  \label{p4:eq:central-binomial}
\end{equation}
that is, the half-index coefficients of \textsf{S1} are $A_n/2^{n}$ on the
even sheet.  Moreover
\begin{equation}
  \frac1{\mathcal S_\sigma(x)}\sim
  \frac1{C_\sigma}\,2^{-x}x^{-\sigma/2}
  \sum_{n\ge0}\frac{(-\sigma)^{n}A_n}{x^{n}} .
  \label{p4:eq:reciprocal}
\end{equation}
\end{corollary}

\begin{proof}
For \cref{p4:eq:central-binomial} substitute $x=2m$ and
$C_-2^{2m}(2m)^{-1/2}=4^m/\sqrt{\pi m}$.  For \cref{p4:eq:reciprocal} note
that the reciprocal of \cref{p4:eq:normal-form} carries $\e^{-\sigma D}$,
whose coefficients are $(-\sigma)^nA_n$ by the oddness argument in the proof
of \cref{p4:thm:forward}.
\end{proof}

\section{The centred normal form}
\label{p4:sec:centred}

\begin{sourcebox}
This section is \textsf{S1}'s distinctive contribution: a translation of the
index that makes the correction series \emph{sparse in a second way} and
equalizes the two leading constants.  \textsf{S1} states it as two parallel
statements in its half index; it is recast here as a single $\sigma$-uniform
identity in the original index, and the coefficients are recomputed in that
normalization.
\end{sourcebox}

\begin{theorem}[Centred exact normal form]
\label{p4:thm:centred}
Put
\begin{equation}
  u:=\frac x2+\frac14 .
  \label{p4:eq:centred-variable}
\end{equation}
Then, exactly and for both parities,
\begin{equation}
  \boxed{\;
  \mathcal S_\sigma(x)=\frac{4^{u}}{\sqrt{2\pi}}
  \left(\frac{\Gamma(u+3/4)}{\Gamma(u+1/4)}\right)^{\!\sigma},\;}
  \label{p4:eq:centred-exact}
\end{equation}
and, as $u\to\infty$ in $|\arg u|\le\pi-\delta$,
\begin{equation}
  \mathcal S_\sigma(x)\sim\frac{4^{u}}{\sqrt{2\pi}}\;u^{\sigma/2}
  \exp\Bigl(\sigma\sum_{q\ge1}\frac{\nu_{2q}}{u^{2q}}\Bigr),
  \qquad
  \nu_{2q}=\frac{B_{2q+1}(1/4)}{q(2q+1)} .
  \label{p4:eq:centred-expansion}
\end{equation}
In particular \emph{every odd inverse power is absent} from the centred
correction, and the leading constant $1/\sqrt{2\pi}$ is the same on both
sheets.
\end{theorem}

\begin{proof}
From $u=x/2+1/4$ we get $x/2+1=u+3/4$ and $x/2+1/2=u+1/4$, while
$2^{x}=2^{2u-1/2}=4^{u}/\sqrt2$; substituting in \cref{p4:eq:branch-ratio}
gives \cref{p4:eq:centred-exact}.  For \cref{p4:eq:centred-expansion} apply
\cref{p4:eq:tricomi} with $a=\tfrac34$, $b=\tfrac14$, $z=u$: the logarithmic
term is $(a-b)\log u=\tfrac12\log u$, and the $n$-th coefficient is
$\frac{(-1)^{n+1}}{n(n+1)}\bigl(B_{n+1}(\tfrac34)-B_{n+1}(\tfrac14)\bigr)$.
Since $a+b=1$, the reflection formula $B_m(1-w)=(-1)^{m}B_m(w)$ gives
$B_{n+1}(\tfrac14)=(-1)^{n+1}B_{n+1}(\tfrac34)$, so the difference vanishes
for odd $n$ and equals $2B_{n+1}(\tfrac34)=-2B_{n+1}(\tfrac14)$ for even
$n=2q$.  This yields $\nu_{2q}=B_{2q+1}(1/4)/(q(2q+1))$.
\end{proof}

\begin{table}[htbp]
\centering
\caption{Centred logarithmic coefficients $\nu_{2q}$ of
\cref{p4:eq:centred-expansion}, recomputed from Bernoulli polynomials at
$1/4$.  They agree with \textsf{S1}'s $\widehat\lambda^{(0)}_{2q}=-\nu_{2q}$
after the marker translation of \cref{p4:rem:dictionary}.}
\label{p4:tab:nu}
\small
\begin{tabular}{@{}rrrrrr@{}}
\toprule
$2q$ & $2$ & $4$ & $6$ & $8$ & $10$\\
\midrule
$\nu_{2q}$ & $\dfrac1{64}$ & $-\dfrac5{2048}$ & $\dfrac{61}{49152}$
 & $-\dfrac{1385}{1048576}$ & $\dfrac{50521}{20971520}$\\
\bottomrule
\end{tabular}
\end{table}

\begin{remark}[Two different sparsities]
\label{p4:rem:two-sparsities}
The uncentred correction $\sigma D(x)$ contains only \emph{odd} powers of
$1/x$; the centred correction contains only \emph{even} powers of $1/u$.
There is no contradiction: the two statements concern different remainder
functions, because the $\tfrac\sigma2\log$ term is attached to $x$ in one and
to $u$ in the other.  Both are instances of one rule, recorded in
\cref{p4:lem:centring} below, and the swing shifts are precisely the
optimally centred ones.
\end{remark}

\begin{lemma}[Optimal centring of a gamma ratio]
\label{p4:lem:centring}
Let $a,b\in\Real$ with $\beta:=a-b\ne0$, and set $s=(a+b-1)/2$ and $z=w+s$.
Then $\Gamma(w+a)/\Gamma(w+b)=\Gamma(z+\tfrac{1+\beta}2)/
\Gamma(z+\tfrac{1-\beta}2)$, the two shifts sum to $1$, and
\begin{equation}
  \log\frac{\Gamma(z+\frac{1+\beta}2)}{\Gamma(z+\frac{1-\beta}2)}
  \sim\beta\log z-\sum_{q\ge1}
   \frac{B_{2q+1}\bigl(\frac{1-\beta}2\bigr)}{q(2q+1)}\,z^{-2q},
  \qquad |\arg z|\le\pi-\delta,
  \label{p4:eq:centring}
\end{equation}
with no odd inverse powers.
\end{lemma}

\begin{proof}
The shift identity is immediate.  Apply \cref{p4:eq:tricomi} and the argument
of \cref{p4:thm:centred}: complementary shifts make the Bernoulli difference
vanish in odd degree by $B_m(1-w)=(-1)^mB_m(w)$.
\end{proof}

For the swing factorial $a=1$, $b=\tfrac12$, $w=x/2$, so $s=\tfrac14$ and
$z=u$, and $\beta=\tfrac12$: \cref{p4:eq:centred-expansion} is exactly
\cref{p4:eq:centring} raised to the power $\sigma$.  The first centred
expansions are
\begin{align}
  \mathcal S_-(x)&\sim\frac{4^{u}}{\sqrt{2\pi u}}
  \Bigl(1-\frac1{64u^{2}}+\frac{21}{8192u^{4}}
   -\frac{671}{524288u^{6}}+\cdots\Bigr),
  \label{p4:eq:centred-even}\\
  \mathcal S_+(x)&\sim\frac{4^{u}\sqrt u}{\sqrt{2\pi}}
  \Bigl(1+\frac1{64u^{2}}-\frac{19}{8192u^{4}}
   +\frac{631}{524288u^{6}}+\cdots\Bigr),
  \label{p4:eq:centred-odd}
\end{align}
whose coefficients were recomputed from
$\exp\bigl(\sigma\sum_q\nu_{2q}z^{q}\bigr)$ and agree with \textsf{S1}.

\section{The two Lambert cores, and why the branch matters}
\label{p4:sec:lambert}

\subsection{Specialization of the core theorem}

Write
\begin{equation}
  L_\sigma(y):=\log\frac{y}{C_\sigma}
   =\log y+\tfrac12\log\pi+\tfrac\sigma2\lambda .
  \label{p4:eq:L-def}
\end{equation}
Dropping the correction in \cref{p4:eq:normal-form} leaves the \emph{core
equation}
\begin{equation}
  L_\sigma(y)=\lambda \Xcore+\frac\sigma2\log \Xcore,
  \qquad\text{equivalently}\qquad
  C_\sigma\,2^{\Xcore}\Xcore^{\sigma/2}=y .
  \label{p4:eq:core-equation}
\end{equation}

\begin{theorem}[Explicit Lambert cores]
\label{p4:thm:cores}
\Cref{p0:thm:lambert-core} with the dictionary \cref{p4:eq:dictionary} gives
\begin{equation}
  \Xcore_{\sigma,k}(y)=\frac{\sigma}{2\lambda}\,
  \Wlam_k\!\bigl(2\sigma\lambda\,\e^{2\sigma L_\sigma(y)}\bigr),
  \label{p4:eq:core-general}
\end{equation}
and the unique large positive real solutions are
\begin{equation}
  \boxed{\;
  \Xcore_+(y)=\frac{\Wlam_0\!\bigl(4\pi\lambda y^{2}\bigr)}{2\lambda},
  \qquad
  \Xcore_-(y)=-\frac{\Wlam_{-1}\!\bigl(-4\lambda/(\pi y^{2})\bigr)}{2\lambda}.
  \;}
  \label{p4:eq:cores}
\end{equation}
$\Xcore_+$ is defined for all $y>0$; $\Xcore_-$ is defined exactly for
\begin{equation}
  y\ge y_\star:=2\sqrt{\frac{\e\lambda}{\pi}}=1.548870223880\ldots,
  \label{p4:eq:ystar}
\end{equation}
and satisfies $\Xcore_-(y)\ge1/(2\lambda)=0.72134752\ldots$ there.
\end{theorem}

\begin{proof}
With $a=\lambda$, $b=\sigma/2$, \cref{p0:thm:lambert-core} reads
$\Xcore=(b/a)\Wlam_k\bigl((a/b)\e^{L/b}\bigr)$, and $1/\sigma=\sigma$ gives
\cref{p4:eq:core-general}.  For $\sigma=+1$, $C_+=1/\sqrt{2\pi}$ so
$\e^{2L_+}=2\pi y^{2}$ and the argument is $4\pi\lambda y^{2}>0$, on which
$\Wlam_0$ is the only real branch.  For $\sigma=-1$, $C_-=\sqrt{2/\pi}$ so
$\e^{-2L_-}=2/(\pi y^{2})$ and the argument is $-4\lambda/(\pi y^{2})<0$;
it lies in $[-1/\e,0)$ exactly when $y^{2}\ge4\lambda\e/\pi$, i.e.\
$y\ge y_\star$.  On that interval both $\Wlam_0$ and $\Wlam_{-1}$ are real;
$\Wlam_{-1}\le-1$ gives $\Xcore_-\ge1/(2\lambda)$, and
$\Wlam_{-1}(z)\to-\infty$ as $z\to0^-$ gives $\Xcore_-(y)\to\infty$.
\Cref{p4:prop:two-roots} identifies the two roots and shows that $\Wlam_0$
gives the wrong one.
\end{proof}

\subsection{Why $\Wlam_0$ is the wrong root on the even sheet}

\begin{proposition}[The even carrier is not injective]
\label{p4:prop:two-roots}
Let $g_-(\Xcore)=C_-2^{\Xcore}\Xcore^{-1/2}$ be the even carrier of
\cref{p4:eq:core-equation}.  Then $g_-$ is strictly decreasing on
$(0,1/(2\lambda))$, strictly increasing on $(1/(2\lambda),\infty)$, and
\begin{equation}
  \min_{\Xcore>0}g_-(\Xcore)=g_-\Bigl(\frac1{2\lambda}\Bigr)
  =2\sqrt{\frac{\e\lambda}{\pi}}=y_\star .
  \label{p4:eq:carrier-min}
\end{equation}
Hence for $y>y_\star$ the equation $g_-(\Xcore)=y$ has exactly two positive
roots $\Xcore^{\mathrm{small}}(y)<1/(2\lambda)<\Xcore^{\mathrm{large}}(y)$;
they coalesce at $y=y_\star$.  In terms of the Lambert branches,
\begin{equation}
  \Xcore^{\mathrm{large}}(y)
   =-\frac{\Wlam_{-1}\bigl(-4\lambda/(\pi y^{2})\bigr)}{2\lambda}
   =\Xcore_-(y)\xrightarrow[y\to\infty]{}\infty,
  \qquad
  \Xcore^{\mathrm{small}}(y)
   =-\frac{\Wlam_{0}\bigl(-4\lambda/(\pi y^{2})\bigr)}{2\lambda}
   =\frac{2}{\pi y^{2}}\bigl(1+O(y^{-2})\bigr).
  \label{p4:eq:two-roots}
\end{equation}
The principal branch therefore returns a root that \emph{tends to zero}; it is
not a large-index inverse at any order, and no amount of asymptotic correction
can repair the choice.  On the odd sheet the carrier
$g_+(\Xcore)=C_+2^{\Xcore}\Xcore^{1/2}$ is a strictly increasing bijection of
$(0,\infty)$ onto $(0,\infty)$, its Lambert argument is positive, and
$\Wlam_0$ is the only real branch.
\end{proposition}

\begin{proof}
$(\log g_-)'(\Xcore)=\lambda-1/(2\Xcore)$ vanishes only at
$\Xcore=1/(2\lambda)$ and changes sign there from $-$ to $+$; the minimum
value is
$C_-\e^{\lambda/(2\lambda)}(2\lambda)^{1/2}
=\sqrt{2/\pi}\cdot\sqrt\e\cdot\sqrt{2\lambda}=2\sqrt{\e\lambda/\pi}$,
which is \cref{p4:eq:carrier-min}, and equals $y_\star$ of
\cref{p4:eq:ystar}.  Strict monotonicity on each side gives exactly two roots
for $y>y_\star$.  Both are given by \cref{p4:eq:core-general} with $\sigma=-1$
and a real branch of $\Wlam$; since $\Wlam_{-1}\le-1\le\Wlam_0\le0$ on
$[-1/\e,0)$, the branch $\Wlam_{-1}$ gives the root
$\ge1/(2\lambda)$ and $\Wlam_0$ the root $\le1/(2\lambda)$.  Finally
$\Wlam_0(z)=z+O(z^{2})$ as $z\to0$, so
$\Xcore^{\mathrm{small}}=-\bigl(-4\lambda/(\pi y^2)+O(y^{-4})\bigr)/(2\lambda)
=2/(\pi y^{2})+O(y^{-4})$.
\end{proof}

\begin{remark}[Numerical confirmation and the branch point]
\label{p4:rem:small-root}
Recomputed at $50$ digits: for $y=10^{2},10^{3},10^{5}$ the principal-branch
root is $6.366759642\times10^{-5}$, $6.366203342\times10^{-7}$,
$6.366197724\times10^{-11}$, against $2/(\pi y^{2})=6.366197724\times10^{-5}$,
$6.366197724\times10^{-7}$, $6.366197724\times10^{-11}$.  The threshold
$y_\star$ where the two roots coalesce is exactly the value at which the
Lambert argument reaches the branch point $-1/\e$; the agreement of these two
descriptions is a useful consistency check on the whole normalization.  Note
also $\mathcal S_-\bigl(1/(2\lambda)\bigr)=1.15213\ldots<y_\star$, a fact used
in \cref{p4:prop:core-inequalities}.
\end{remark}

\begin{remark}[The branch rule is forced by $\operatorname{sign}b$]
\label{p4:rem:branch-rule}
All three sources agree on the assignment: $\Wlam_{-1}$ on the even sheet,
$\Wlam_0$ on the odd sheet.  In the model \cref{p0:def:model} the rule is
$b<0\Rightarrow\Wlam_{-1}$, $b>0\Rightarrow\Wlam_0$ for the large real root,
and $b=\sigma/2$; so the branch distinction is forced by the sign of the
power of $x$ in the forward asymptotics --- that is, by the fact that the even
bisection carries $n^{-1/2}$ and the odd bisection $n^{+1/2}$.  It is a
consequence of the parity split, not an independent convention.
\end{remark}

\subsection{Two exact inequalities}

The sources deduce the \emph{sign} of the first inverse correction from
$D>0$.  The following sharpening is exact at all $y$, not only to first order,
and costs nothing.

\begin{proposition}[The core brackets the true inverse]
\label{p4:prop:core-inequalities}
For every $y>1$,
\begin{equation}
  \mathcal I_+(y)<\Xcore_+(y),
  \label{p4:eq:odd-bracket}
\end{equation}
and for every $y>y_\star$,
\begin{equation}
  \mathcal I_-(y)>\Xcore_-(y).
  \label{p4:eq:even-bracket}
\end{equation}
\end{proposition}

\begin{proof}
By \cref{p4:eq:normal-form} and \cref{p4:cor:D-bounds},
$\mathcal S_+(x)=g_+(x)\e^{D(x)}>g_+(x)$ for all $x>0$, where $g_+$ is the
odd carrier.  Both $\mathcal S_+$ and $g_+$ are strictly increasing.  Put
$x_1=\mathcal I_+(y)$; then
$g_+(x_1)<\mathcal S_+(x_1)=y=g_+(\Xcore_+(y))$, and $g_+$ increasing gives
$x_1<\Xcore_+(y)$.

For $\sigma=-1$, $\mathcal S_-(x)=g_-(x)\e^{-D(x)}<g_-(x)$.  Let
$x_0=\mathcal I_-(y)$ with $y>y_\star$.  Since
$\mathcal S_-$ is increasing and $\mathcal S_-(1/(2\lambda))
=1.15213\ldots<y_\star<y$, we have $x_0>1/(2\lambda)$; also
$\Xcore_-(y)\ge1/(2\lambda)$ by \cref{p4:thm:cores}.  On
$[1/(2\lambda),\infty)$ the carrier $g_-$ is strictly increasing
(\cref{p4:prop:two-roots}), and
$g_-(x_0)>\mathcal S_-(x_0)=y=g_-(\Xcore_-(y))$, whence
$x_0>\Xcore_-(y)$.
\end{proof}

\begin{remark}
This explains without computation the signs $c_1^{(+1)}=-1/(4\lambda)<0$ and
$c_1^{(-1)}=+1/(4\lambda)>0$ found in \cref{p4:sec:reversion}, and supplies a
rigorous one-sided bracket at every $y$, used as a check in
\cref{p4:alg:invert}.
\end{remark}

\subsection{A scaled form, and the half-index cores}

Putting $w_\sigma=2\lambda\Xcore_\sigma$ turns \cref{p4:eq:core-equation} into
the single scalar equation
\begin{equation}
  w+\sigma\log w=\Lambda_\sigma,
  \qquad
  \Lambda_\sigma:=\log(\pi y^{2})+\sigma\log(4\lambda),
  \label{p4:eq:scaled-core}
\end{equation}
so that $w_+=\Wlam_0(\e^{\Lambda_+})$ and
$w_-=-\Wlam_{-1}(-\e^{-\Lambda_-})$.  This is \textsf{S3}'s normalization,
with its $\varepsilon=-\sigma$; it is convenient in
\cref{p4:sec:flatten}.

\begin{remark}[Half-index cores]
\label{p4:rem:halfindex-cores}
\textsf{S1}'s carriers, in the half index and with $\alpha=\log4$, are
$w^{\textsf{S1}}_0=-\tfrac1{2\alpha}\Wlam_{-1}(-2\alpha/(\pi y^{2}))$ and
$w^{\textsf{S1}}_1=\tfrac1{2\alpha}\Wlam_0(\pi\alpha y^{2}/2)$.  On the even
sheet $n=2m$ exactly, and one checks directly that
$\Xcore_-(y)=2w^{\textsf{S1}}_0(y)$; the two normalizations then agree, and
their inverse coefficients are related by
$c_n^{(-1)}=2^{\,n+1}d_n^{\textsf{S1},(0)}$, which we verified symbolically
for $1\le n\le8$.  On the odd sheet $n=2m+1$, so the two carriers solve
\emph{different} equations ($C_+2^{X}X^{1/2}=y$ against
$K_14^{w}w^{1/2}=y$) and their coefficient lists are \emph{not} termwise
comparable, although both expansions converge asymptotically to the same
function $\mathcal I_+$.  Failing to notice this is the easiest way to
``discover'' a contradiction between \textsf{S1} and \textsf{S2}/\textsf{S3}
where there is none.
\end{remark}

\section{All-orders reversion around the core}
\label{p4:sec:reversion}

\subsection{Specialization of the reversion theorem}

\begin{theorem}[Swing specialization of \cref{p0:thm:lambert-centered}]
\label{p4:thm:reversion}
Let $\Xcore=\Xcore_\sigma(y)$ solve \cref{p4:eq:core-equation}, let
$t=\Xcore^{-1}$, and put
\begin{equation}
  \Psi_\sigma(t,u)=-\frac1\lambda\left[
    \frac\sigma2\log(1+tu)
    +\sigma\sum_{j\ge1}d_j\Bigl(\frac{t}{1+tu}\Bigr)^{j}\right].
  \label{p4:eq:Psi-swing}
\end{equation}
Then the unique $U\in t\,\mathscr R[[t]]$ with $U=\Psi_\sigma(t,U)$ has
coefficients
\begin{equation}
  c_n^{(\sigma)}=\Coef{t^n}U
  =\sum_{m=1}^{n}\frac1m\,\Coef{t^{n}u^{m-1}}\Psi_\sigma(t,u)^{m}
  \label{p4:eq:c-lagrange}
\end{equation}
by \cref{p0:lem:lagrange-fp}, and equivalently, with $\gamma_1=0$ and
$\gamma_r=c^{(\sigma)}_{r-1}$ for $r\ge2$,
\begin{equation}
  c^{(\sigma)}_n=-\frac1\lambda\Biggl[
   \frac\sigma2\sum_{k=1}^{n}\frac{(-1)^{k+1}}{k}\OB_{n,k}(\gamma)
   +\sigma\!\!\sum_{\substack{1\le j\le n\\ j\ \mathrm{odd}}}\!\! d_j
     \sum_{k=0}^{n-j}(-1)^{k}\binom{j+k-1}{k}\OB_{n-j,k}(\gamma)\Biggr],
  \label{p4:eq:c-bell}
\end{equation}
with $\OB$ the ordinary partial Bell polynomials of \cref{p0:def:obell} and
the expansions of \cref{p0:lem:power-log}.  Only
$c^{(\sigma)}_1,\ldots,c^{(\sigma)}_{n-1}$ occur on the right.  The
coefficients lie in the parity algebra $\mathscr R$ of
\cref{p4:eq:parity-algebra}; no logarithm of $L_\sigma$ occurs.
\end{theorem}

\begin{proof}
This is \cref{p0:thm:lambert-centered} with the dictionary
\cref{p4:eq:dictionary}: $a=\lambda$, $b=\sigma/2$, $q_j=\sigma d_j$, which
turns the general $\Psi$ into \cref{p4:eq:Psi-swing}.  The restriction of the
inner sum to odd $j$ is $d_{2r}=0$.
\end{proof}

\subsection{Recomputed coefficients}

The coefficients below were recomputed from \cref{p4:eq:Psi-swing} by
truncated power-series arithmetic over $\Rat[\lambda,\lambda^{-1}]$,
separately for $\sigma=+1$ and $\sigma=-1$, and then recombined into
$\mathscr R$ by
$c=\tfrac12(c^{+}+c^{-})+\sigma\cdot\tfrac12(c^{+}-c^{-})$.  They agree with
\textsf{S2} for $1\le n\le10$ and with \textsf{S3} for $1\le n\le8$ after the
marker translation $\sigma=-\varepsilon_{\textsf{S3}}$.

\begin{align}
  c_1^{(\sigma)}&=-\frac{\sigma}{4\lambda},
  \qquad
  c_2^{(\sigma)}=\frac1{8\lambda^{2}},
  \label{p4:eq:c12}\\
  c_3^{(\sigma)}&=\frac{2\sigma\lambda^{2}-3\lambda-3\sigma}{48\lambda^{3}},
  \qquad
  c_4^{(\sigma)}=\frac{-4\lambda^{2}+15\sigma\lambda+6}{192\lambda^{4}},
  \label{p4:eq:c34}\\
  c_5^{(\sigma)}&=-\frac{96\sigma\lambda^{4}-80\lambda^{3}
    +40\sigma\lambda^{2}+135\lambda+30\sigma}{1920\lambda^{5}},
  \label{p4:eq:c5}\\
  c_6^{(\sigma)}&=\frac{96\lambda^{4}-180\sigma\lambda^{3}+215\lambda^{2}
    +210\sigma\lambda+30}{3840\lambda^{6}},
  \label{p4:eq:c6}\\
  c_7^{(\sigma)}&=\frac{1}{53760\lambda^{7}}\bigl(
    8160\sigma\lambda^{6}-4312\lambda^{5}+1428\sigma\lambda^{4}
    +1050\lambda^{3}\nonumber\\
   &\hspace{31mm}{}-4025\sigma\lambda^{2}-2100\lambda-210\sigma\bigr),
  \label{p4:eq:c7}\\
  c_8^{(\sigma)}&=\frac{1}{645120\lambda^{8}}\bigl(
    -48960\lambda^{6}+56056\sigma\lambda^{5}-40908\lambda^{4}
    +15015\sigma\lambda^{3}\nonumber\\
   &\hspace{31mm}{}+50610\lambda^{2}+17010\sigma\lambda+1260\bigr),
  \label{p4:eq:c8}\\
  c_9^{(\sigma)}&=-\frac{1}{1290240\lambda^{9}}\bigl(
    1111040\sigma\lambda^{8}-413184\lambda^{7}+79616\sigma\lambda^{6}
    +43512\lambda^{5}\nonumber\\
   &\hspace{20mm}{}-83748\sigma\lambda^{4}+83475\lambda^{3}
    +92610\sigma\lambda^{2}+22050\lambda+1260\sigma\bigr),
  \label{p4:eq:c9}\\
  c_{10}^{(\sigma)}&=\frac{1}{5160960\lambda^{10}}\bigl(
    2222080\lambda^{8}-1756032\sigma\lambda^{7}+779152\lambda^{6}
    -203280\sigma\lambda^{5}\nonumber\\
   &\hspace{20mm}{}-162582\lambda^{4}+487305\sigma\lambda^{3}
    +309960\lambda^{2}+55440\sigma\lambda+2520\bigr).
  \label{p4:eq:c10}
\end{align}

Setting $\sigma=\pm1$ and $\lambda=\log2$,
\begin{align}
  \mathcal I_+(y)&\sim \Xcore_+
   -\frac1{4\lambda \Xcore_+}+\frac1{8\lambda^{2}\Xcore_+^{2}}
   +\frac{2\lambda^{2}-3\lambda-3}{48\lambda^{3}\Xcore_+^{3}}
   +\frac{-4\lambda^{2}+15\lambda+6}{192\lambda^{4}\Xcore_+^{4}}+\cdots,
  \label{p4:eq:odd-inverse}\\
  \mathcal I_-(y)&\sim \Xcore_-
   +\frac1{4\lambda \Xcore_-}+\frac1{8\lambda^{2}\Xcore_-^{2}}
   +\frac{-2\lambda^{2}-3\lambda+3}{48\lambda^{3}\Xcore_-^{3}}
   +\frac{-4\lambda^{2}-15\lambda+6}{192\lambda^{4}\Xcore_-^{4}}+\cdots.
  \label{p4:eq:even-inverse}
\end{align}
The signs of the first corrections are those forced by
\cref{p4:prop:core-inequalities}.

\subsection{A differential form of the same coefficients}

All three sources also give an operator form of the reversion, which is the
convenient one for the Stokes analysis of \cref{p4:sec:stokes}.  We state it
as a chapter-local lemma; it is a candidate for promotion, being independent
of the swing factorial.

\begin{lemma}[Perturbed inverse in phase coordinates]
\label{p4:lem:phase-lb}
Let $F_0$ be invertible near the point considered with $F_0'\ne0$, let $g$ be
a formal correction, and let $\Xcore$ satisfy $F_0(\Xcore)=z$.  Introduce
\begin{equation}
  \mathscr D:=\frac1{F_0'(\Xcore)}\frac{\dd}{\dd \Xcore}.
  \label{p4:eq:phase-operator}
\end{equation}
Then the solution $x$ of $F_0(x)+\eta\,g(x)=z$ with $x|_{\eta=0}=\Xcore$ is,
as a formal power series in the marker $\eta$,
\begin{equation}
  x=\Xcore+\sum_{k\ge1}\frac{(-\eta)^{k}}{k!}\,
   \mathscr D^{\,k-1}\!\left(\frac{g(\Xcore)^{k}}{F_0'(\Xcore)}\right).
  \label{p4:eq:phase-lb}
\end{equation}
\end{lemma}

\begin{proof}
This is \cref{p0:thm:perturbed-inversion} under the dictionary $F\mapsto F_0$,
$E\mapsto g$, $\varepsilon\mapsto\eta$, $\mu_{0}\mapsto\Xcore$,
$y\mapsto z$ and $\mathcal{D}\mapsto\mathscr D$, which carries
\eqref{p0:eq:reversion-operator} to \cref{p4:eq:phase-operator} and
\eqref{p0:eq:perturbed-inversion} to \cref{p4:eq:phase-lb}.  Since $g$ is here
a formal and in general divergent series, no contour is available and the
statement invoked is the coefficientwise one, \cref{p0:cor:perturbed-formal},
whose only hypothesis is that $F_0'(\Xcore)$ be invertible.
\end{proof}

\begin{proposition}[Differential formula for $c_n^{(\sigma)}$]
\label{p4:prop:c-differential}
Let $\mathcal Q_\sigma(t)=\sigma\sum_{j\ge1}d_jt^{j}$ and
\begin{equation}
  \mathscr L_\sigma:=-\frac{t^{2}}{\lambda+\sigma t/2}\frac{\dd}{\dd t}
  \qquad(t=\Xcore^{-1}),
  \label{p4:eq:L-operator}
\end{equation}
which is \cref{p4:eq:phase-operator} for
$F_{0,\sigma}(\Xcore)=\lambda \Xcore+\tfrac\sigma2\log \Xcore$ expressed in
the variable $t$.  Then, for every $n\ge1$,
\begin{equation}
  \boxed{\;
  c_n^{(\sigma)}=\Coef{t^{n}}
   \sum_{k=1}^{\lfloor(n+1)/2\rfloor}\frac{(-1)^{k}}{k!}\,
   \mathscr L_\sigma^{\,k-1}
   \left(\frac{\mathcal Q_\sigma(t)^{k}}{\lambda+\sigma t/2}\right),\;}
  \label{p4:eq:c-differential}
\end{equation}
a \emph{finite} sum requiring only $d_1,\ldots,d_{n}$.
\end{proposition}

\begin{proof}
Apply \cref{p4:lem:phase-lb} with $F_0=F_{0,\sigma}$, $g=\sigma D$ and
$\eta=1$; the substitution $t=\Xcore^{-1}$ turns
$F_{0,\sigma}'(\Xcore)=\lambda+\sigma/(2\Xcore)$ into
$\lambda+\sigma t/2$ and $\dd/\dd \Xcore$ into $-t^{2}\dd/\dd t$.  Setting
the marker $\eta$ to $1$ is legitimate coefficientwise because
$\mathcal Q_\sigma(t)^{k}=O(t^{k})$ while each application of
$\mathscr L_\sigma$ raises the order in $t$ by one, so the $k$-th summand
begins at $t^{2k-1}$ and can contribute to $\Coef{t^{n}}$ only when
$2k-1\le n$.
\end{proof}

\begin{repairbox}
\textbf{Source claim.}  \textsf{S2} writes its operator expansion as an
\emph{equality},
$\mathcal I_\sigma(y)=\Xcore+\sum_{m\ge1}\frac{(-\sigma)^{m}}{m!}
\mathscr D^{m-1}\bigl[D(\Xcore)^{m}/F_0'(\Xcore)\bigr]$, after setting its
coupling parameter to $1$.  \textbf{Defect.}  With the marker removed there is
no small parameter, the series over $m$ does not converge, and the left-hand
side is a number while the right-hand side is a divergent series; the
statement as printed is false.  \textbf{Repair.}  The correct statement is the
asymptotic one: the $m$-th summand is $O(\Xcore^{-(2m-1)})$, so for each fixed
$M$ the partial sum through $m=M$ differs from $\mathcal I_\sigma(y)$ by
$O(\Xcore^{-(2M+1)})$, and regrouping by powers of $\Xcore^{-1}$ reproduces
\cref{p4:eq:c-differential}.  \textsf{S1} and \textsf{S3} state the same
formula correctly, as a coefficient identity.
\end{repairbox}

\subsection{The analytic statement, with a proof}

\begin{repairbox}
\textbf{Source claim.}  All three sources assert the Poincar\'e statement
``$\mathcal I_\sigma(y)=\Xcore_\sigma(y)+\sum_{n<M}c^{(\sigma)}_n
\Xcore_\sigma(y)^{-n}+O(\Xcore_\sigma(y)^{-M})$''.  \textbf{Defect.}
\textsf{S3} gives no proof at all (``Likewise, for each fixed $M$, \ldots'');
\textsf{S1} gives a two-line sketch that appeals to the mean value theorem but
never bounds the residual of the truncated ansatz, and in particular never
uses any control on the forward remainder at the \emph{perturbed} point;
\textsf{S2} proves the ingredients (a residual-to-index bound and derivative
lower bounds) but does not assemble them.  This is the one place in the three
sources where a formal computation is upgraded to an analytic conclusion
without an argument.  \textbf{Repair.}  \Cref{p4:thm:inverse-analytic} below
supplies the proof, and makes the constant effective by using the enveloping
remainder \cref{p4:thm:envelope}.
\end{repairbox}

\begin{theorem}[Inverse transseries, analytically]
\label{p4:thm:inverse-analytic}
Fix $M\ge1$ and $\sigma\in\{-1,+1\}$.  Then, as $y\to+\infty$,
\begin{equation}
  \boxed{\;
  \mathcal I_\sigma(y)=\Xcore_\sigma(y)
   +\sum_{n=1}^{M-1}\frac{c^{(\sigma)}_n}{\Xcore_\sigma(y)^{n}}
   +O\!\left(\Xcore_\sigma(y)^{-M}\right),\;}
  \label{p4:eq:inverse-analytic}
\end{equation}
with an effective implied constant depending only on $M$ and $\sigma$.  Since
$\Xcore_\sigma(y)=\bigl(\log y\bigr)/\lambda\,(1+o(1))$, the expansion is one
in inverse powers of $\log y$ once the dominant power--logarithmic phase has
been inverted exactly.
\end{theorem}

\begin{proof}
Write $\Xcore=\Xcore_\sigma(y)$ and
$x_M:=\Xcore+\sum_{n=1}^{M-1}c^{(\sigma)}_n\Xcore^{-n}$; for $y$ large,
$x_M=\Xcore(1+O(\Xcore^{-2}))$, so $x_M>1$ and $x_M\asymp\Xcore$.  Consider
the phase residual
\[
  \rho:=\log\mathcal S_\sigma(x_M)-\log y
   =\lambda(x_M-\Xcore)
     +\frac\sigma2\log\frac{x_M}{\Xcore}+\sigma D(x_M),
\]
where we used \cref{p4:eq:normal-form} and
$\log y=\log C_\sigma+\lambda\Xcore+\tfrac\sigma2\log\Xcore$.  Put
$t=\Xcore^{-1}$ and $U_{M-1}(t)=\sum_{n<M}c^{(\sigma)}_nt^{n}$, so that
$x_M=\Xcore(1+tU_{M-1})$ and
\[
  \rho=\lambda\,\Xcore\Bigl[\,\underbrace{tU_{M-1}
    +\frac{\sigma t}{2\lambda}\log(1+tU_{M-1})
    +\frac{\sigma t}{\lambda}D_{<M}(x_M)}_{=:t\,\mathcal E(t)}\,\Bigr]
   +\sigma\bigl(D(x_M)-D_{<M}(x_M)\bigr),
\]
where $D_{<M}$ denotes the truncation of the series
\cref{p4:eq:d-bernoulli} after the terms of order $x^{-(M-1)}$.  Two
estimates finish the proof.

\emph{(i) The algebraic residual.}  Substituting $x_M=\Xcore(1+tU_{M-1})$ and
$D_{<M}(x_M)=\mathcal Q_\sigma\bigl(t/(1+tU_{M-1})\bigr)/\sigma$ truncated,
the bracket $\mathcal E(t)$ is exactly
$U_{M-1}(t)-\Psi_\sigma\bigl(t,U_{M-1}(t)\bigr)$ modulo terms of order
$t^{M}$ coming from the truncation of $\mathcal Q_\sigma$.  By
\cref{p4:thm:reversion} the coefficients of $t^{1},\ldots,t^{M-1}$ in
$U-\Psi_\sigma(t,U)$ vanish when $U$ is replaced by $U_{M-1}$ --- this is the
triangularity of the recurrence --- so $\mathcal E(t)=O(t^{M})$ with an
explicit rational-in-$\lambda$ constant obtained from
$c^{(\sigma)}_1,\ldots,c^{(\sigma)}_{M-1}$ and $d_1,\ldots,d_{M}$.  Hence
$\lambda\Xcore\cdot t\,\mathcal E(t)=O(\Xcore^{-M+1}\cdot\Xcore^{-1})\cdot
\Xcore=O(\Xcore^{-M+1})$; more carefully, $\Xcore\,t\,\mathcal E(t)
=\mathcal E(t)=O(\Xcore^{-M})$.

\emph{(ii) The transcendental residual.}  By \cref{p4:thm:envelope} with
$N=\lceil M/2\rceil$,
$|D(x_M)-D_{<M}(x_M)|\le|d_{2N+1}|\,x_M^{-(2N+1)}\le
c_M\,\Xcore^{-M}$ for an explicit $c_M$, since $2N+1\ge M$ and
$x_M\asymp\Xcore$.

Therefore $|\rho|\le C_M\Xcore^{-M}$ with $C_M$ explicit.  Finally,
$\bigl(\log\mathcal S_\sigma\bigr)'(x)=\lambda+\tfrac\sigma{2x}+\sigma D'(x)
\ge\lambda-\tfrac1{2x}-\tfrac1{4x^{2}}$ by \cref{p4:cor:D-bounds}, which
exceeds $\lambda/2$ once $x\ge2/\lambda$.  \Cref{p0:thm:backward-error}
(residual-to-root conversion) then gives
\[
  |x_M-\mathcal I_\sigma(y)|\le\frac{|\rho|}{\lambda/2}
   \le\frac{2C_M}{\lambda}\,\Xcore^{-M},
\]
which is \cref{p4:eq:inverse-analytic}.  The final sentence follows from
$\lambda\Xcore+\tfrac\sigma2\log\Xcore=L_\sigma(y)$ and
$L_\sigma(y)=\log y+O(1)$.
\end{proof}

\begin{remark}[An exact algebraic certificate]
\label{p4:rem:certificate}
Step (i) of the proof is also the recommended \emph{implementation} check.
Given a candidate truncation $U_{M-1}$, substitute it into
$U-\Psi_\sigma(t,U)$ and expand: the coefficients of
$t^{1},\ldots,t^{M-1}$ must vanish identically in
$\Rat[\lambda,\lambda^{-1},\sigma]/(\sigma^{2}-1)$.  This is an exact
algebraic certificate for every table of $c^{(\sigma)}_n$, and it is how the
tables in \cref{p4:eq:c12}--\cref{p4:eq:c10} were verified here, independently
of the recurrence that produced them.
\end{remark}

\subsection{The centred reversion}

Applying \cref{p0:thm:lambert-centered} to the centred normal form
\cref{p4:eq:centred-expansion} --- that is, with
$a=2\lambda$, $b=\sigma/2$, $C=1/\sqrt{2\pi}$ and $q_j=\sigma\nu_j$
($\nu_j=0$ for odd $j$) --- gives the centred cores
\begin{equation}
  \mathcal U_+(y)=\frac{\Wlam_0\!\bigl(8\pi\lambda y^{2}\bigr)}{4\lambda},
  \qquad
  \mathcal U_-(y)
  =-\frac{\Wlam_{-1}\!\bigl(-2\lambda/(\pi y^{2})\bigr)}{4\lambda},
  \label{p4:eq:centred-cores}
\end{equation}
and the centred inverse expansion
\begin{equation}
  \boxed{\;
  \mathcal I_\sigma(y)\sim 2\,\mathcal U_\sigma(y)-\frac12
   +2\sum_{n\ge2}\frac{\widehat c^{\,(\sigma)}_n}{\mathcal U_\sigma(y)^{n}},
  \qquad \widehat c^{\,(\sigma)}_1=0 .\;}
  \label{p4:eq:centred-inverse}
\end{equation}
The vanishing of $\widehat c^{\,(\sigma)}_1$ is immediate from
$\nu_1=0$.  Recomputed in the present normalization,
\begin{align}
  \widehat c^{\,(\sigma)}_2&=-\frac{\sigma}{128\lambda},
  &\widehat c^{\,(\sigma)}_3&=\frac1{512\lambda^{2}},
  &\widehat c^{\,(\sigma)}_4&=\frac{\sigma(5\lambda^{2}-2)}{4096\lambda^{3}},
  \label{p4:eq:chat234}\\
  \widehat c^{\,(\sigma)}_5&=-\frac{7\lambda^{2}-2}{16384\lambda^{4}},
  &\widehat c^{\,(\sigma)}_6&=-\frac{\sigma\bigl(244\lambda^{4}
     -57\lambda^{2}+12\bigr)}{393216\lambda^{5}},
  &\widehat c^{\,(\sigma)}_7&=\frac{334\lambda^{4}-75\lambda^{2}+12}
     {1572864\lambda^{6}},
  \label{p4:eq:chat567}
\end{align}
and
\begin{equation}
  \widehat c^{\,(\sigma)}_8
  =\frac{\sigma\bigl(4155\lambda^{6}-460\lambda^{4}
    +96\lambda^{2}-12\bigr)}{6291456\lambda^{7}} .
  \label{p4:eq:chat8}
\end{equation}
These agree with \textsf{S1}'s centred table after the substitution
$\alpha=2\lambda$ and the marker translation of \cref{p4:rem:dictionary};
observe that $\widehat c_n$ is $\sigma$-even for odd $n$ and $\sigma$-odd for
even $n$, a reflection symmetry of \cref{p4:eq:centred-exact} that the
uncentred coefficients do not have.

\begin{repairbox}
\textbf{Source claim.}  \textsf{S1} concludes that ``the centred one is often
the preferable computational normal form''.  \textbf{Defect.}  Its own table
already contradicts the unqualified form of this claim: at $m=200$ the centred
core beats the uncentred core by $10^{3}$ ($2.8\times10^{-7}$ against
$4.5\times10^{-4}$), but at $R=4$ the \emph{uncentred} expansion is slightly
better ($7.3\times10^{-16}$ against $1.1\times10^{-15}$).  \textbf{Repair.}
The correct statement, confirmed by the recomputation in
\cref{p4:tab:inverse-errors} and \cref{p4:tab:centred-errors}, is: centring
removes the order-$\Xcore^{-1}$ term entirely, so it improves the \emph{bare
core} by a factor $\asymp\Xcore$ and is the right choice when no corrections
are to be carried; once several corrections are carried, the two
normalizations are comparable, with neither dominating uniformly.
\end{repairbox}

\section{Flattening the core into a polynomial--logarithmic transseries}
\label{p4:sec:flatten}

\subsection{The specialization}

Set
\begin{equation}
  H_\sigma(y):=\log\frac{L_\sigma(y)}{\lambda},
  \label{p4:eq:H-def}
\end{equation}
and recall from \cref{p0:thm:flattening} that the inverse of the model
\cref{p0:def:model} has the unique power--logarithmic form
$x\sim a^{-1}\bigl[L-bH+\sum_{n\ge1}P_n(H)L^{-n}\bigr]$ with $\deg P_n\le n$
and leading term $b^{n+1}H^{n}/n$.  With the dictionary
\cref{p4:eq:dictionary},
\begin{equation}
  \boxed{\;
  \mathcal I_\sigma(y)\sim\frac1\lambda\left[
    L_\sigma-\frac\sigma2H_\sigma
    +\sum_{n\ge1}\frac{P^{(\sigma)}_n(H_\sigma)}{L_\sigma^{\,n}}\right],\;}
  \label{p4:eq:flattened}
\end{equation}
where the $P^{(\sigma)}_n$ are generated by \cref{p0:thm:flattening} from
$b=\sigma/2$, $a=\lambda$, $q_j=\sigma d_j$; the fixed-order meaning of
\cref{p4:eq:flattened} is the analytic clause of \cref{p0:thm:flattening},
namely an error $O\bigl((1+H_\sigma)^{N+1}L_\sigma^{-(N+1)}\bigr)$ after $N$
correction terms, whose hypotheses hold here by
\cref{p4:thm:inverse-analytic}.  The recomputed polynomials are
\begin{align}
  P^{(\sigma)}_1(H)&=\frac{H-\sigma\lambda}{4},
  \label{p4:eq:P1}\\
  P^{(\sigma)}_2(H)&=\frac{\sigma H^{2}-2(\lambda+\sigma)H+2\lambda}{16},
  \label{p4:eq:P2}\\
  P^{(\sigma)}_3(H)&=\frac1{96}\Bigl[2H^{3}-(6\sigma\lambda+9)H^{2}
    +(18\sigma\lambda+6)H
    +4\sigma\lambda^{3}-6\lambda^{2}-6\sigma\lambda\Bigr],
  \label{p4:eq:P3}\\
  P^{(\sigma)}_4(H)&=\frac1{384}\Bigl[3\sigma H^{4}
    -(12\lambda+22\sigma)H^{3}+(66\lambda+36\sigma)H^{2}\nonumber\\
   &\hspace{14mm}{}+(24\lambda^{3}-36\sigma\lambda^{2}-72\lambda-12\sigma)H
    -8\lambda^{3}+30\sigma\lambda^{2}+12\lambda\Bigr],
  \label{p4:eq:P4}\\
  P^{(\sigma)}_5(H)&=\frac1{3840}\Bigl[12H^{5}-(60\sigma\lambda+125)H^{4}
    +(500\sigma\lambda+350)H^{3}\nonumber\\
   &\hspace{14mm}{}
    +(240\sigma\lambda^{3}-360\lambda^{2}-1050\sigma\lambda-300)H^{2}
    \nonumber\\
   &\hspace{14mm}{}
    +(-280\sigma\lambda^{3}+780\lambda^{2}+600\sigma\lambda+60)H
    \nonumber\\
   &\hspace{14mm}{}
    -192\sigma\lambda^{5}+160\lambda^{4}-80\sigma\lambda^{3}
    -270\lambda^{2}-60\sigma\lambda\Bigr].
  \label{p4:eq:P5}
\end{align}
These reproduce \textsf{S2} exactly for $1\le n\le5$; they were recomputed for
$\sigma=\pm1$ separately and recombined in $\mathscr R$, and independently
re-derived through $n=6$.  They also satisfy the cross-normalization identity
$P^{(\sigma)}_n(0)=\lambda^{\,n+1}c^{(\sigma)}_n$ of \cref{p0:thm:flattening}:
for instance $P^{(\sigma)}_1(0)=-\sigma\lambda/4=\lambda^{2}c^{(\sigma)}_1$ and
$P^{(\sigma)}_2(0)=\lambda/8=\lambda^{3}c^{(\sigma)}_2$, which ties
\cref{p4:eq:P1}--\cref{p4:eq:P5} to \cref{p4:eq:c12}--\cref{p4:eq:c10} and
catches any transcription error in either list.

\subsection{The pure Lambert block in closed form}

The pure block --- the case $Q\equiv0$, in which the whole correction series
is suppressed and only the Lambert core survives --- is
\cref{p0:thm:pure-stirling}, proved there in the shortest of the three forms
the sources use.  Specializing its data to the swing branches, $a=\lambda$ and
$b=\sigma/2$, gives
\begin{equation}
  P^{(0)}_n(H)=\Bigl(\frac{\sigma}{2}\Bigr)^{\,n+1}
  \sum_{m=1}^{n}(-1)^{m+1}\,\stirlingfirst{n}{m}\,
  \frac{H^{\,n-m+1}}{(n-m+1)!},
  \label{p4:eq:pure-lambert}
\end{equation}
with $\stirlingfirst{n}{m}$ the unsigned Stirling numbers of the first kind.
Reducing with $\sigma^{2}=1$, the two parities differ only by the overall sign
$\sigma^{\,n+1}$: the even and odd pure blocks coincide for odd $n$ and are
opposite for even $n$.

Writing $b=\sigma/2$, the first few are
\[
  P^{(0)}_1=b^{2}H,\quad
  P^{(0)}_2=\tfrac{b^{3}}2(H^{2}-2H),\quad
  P^{(0)}_3=\tfrac{b^{4}}6(2H^{3}-9H^{2}+6H),\quad
  P^{(0)}_4=\tfrac{b^{5}}{12}(3H^{4}-22H^{3}+36H^{2}-12H),
\]
and \cref{p4:eq:pure-lambert} was recomputed here against the recurrence of
\cref{p0:thm:flattening} for $1\le n\le8$, in exact arithmetic with $b$
symbolic.  In particular the leading
coefficient is $b^{n+1}/n$, which is \cref{p0:thm:flattening}'s degree
statement.  \Cref{p0:thm:pure-stirling} is the precise bridge between Lambert
asymptotics and the first-kind Stirling calculus, and it is stated there
rather than here because all four chapters use it.

\subsection{The core series and its convergence}

The flattening can also be organized as ``expand the core, then compose'',
which is \textsf{S3}'s route and which is constructive at every order.

\begin{proposition}[All coefficients of the scaled core]
\label{p4:prop:core-series}
Let $w$ solve $w+\sigma\log w=\Lambda$ as in \cref{p4:eq:scaled-core}, and
put $\ell=\log\Lambda$.  Define
\begin{equation}
  A_k(\ell):=\frac1{k+1}\Coef{z^{k}}\bigl(\ell+\log(1+z)\bigr)^{k+1},
  \qquad k\ge0 .
  \label{p4:eq:Ak-def}
\end{equation}
Then
\begin{equation}
  w\;=\;\Lambda+\sum_{k\ge0}\frac{(-\sigma)^{k+1}A_k(\ell)}{\Lambda^{k}},
  \label{p4:eq:core-series}
\end{equation}
and, with $s(k,j)$ the signed Stirling numbers of the first kind,
\begin{equation}
  A_0(\ell)=\ell,
  \qquad
  A_k(\ell)=\frac1{(k+1)\,k!}\sum_{j=1}^{k}\binom{k+1}{j}j!\,s(k,j)\,
  \ell^{\,k+1-j}\quad(k\ge1).
  \label{p4:eq:Ak-stirling}
\end{equation}
The first polynomials are
\[
  A_0=A_1=\ell,\quad
  A_2=\tfrac{\ell(2-\ell)}2,\quad
  A_3=\tfrac{\ell(2\ell^{2}-9\ell+6)}6,\quad
  A_4=-\tfrac{\ell(3\ell^{3}-22\ell^{2}+36\ell-12)}{12},
\]
\[
  A_5=\tfrac{\ell(12\ell^{4}-125\ell^{3}+350\ell^{2}-300\ell+60)}{60},
  \quad
  A_6=-\tfrac{\ell(20\ell^{5}-274\ell^{4}+1125\ell^{3}-1700\ell^{2}
    +900\ell-120)}{120} .
\]
\end{proposition}

\begin{proof}
Write $w=\Lambda(1+\upsilon)$ and $z=\Lambda^{-1}$.  The equation becomes
$\Lambda\upsilon=-\sigma\bigl(\ell+\log(1+\upsilon)\bigr)$, i.e.\
$\upsilon=-\sigma z\bigl(\ell+\log(1+\upsilon)\bigr)$.  Lagrange inversion
gives
$\upsilon=\sum_{k\ge0}\frac{(-\sigma z)^{k+1}}{k+1}
\Coef{\upsilon^{k}}(\ell+\log(1+\upsilon))^{k+1}$; multiplying by
$\Lambda=z^{-1}$ yields \cref{p4:eq:core-series} with $A_k$ as in
\cref{p4:eq:Ak-def}.  Expanding the $(k+1)$-st power and using
$\bigl(\log(1+z)\bigr)^{j}=j!\sum_{k\ge j}s(k,j)z^{k}/k!$ gives
\cref{p4:eq:Ak-stirling}.
\end{proof}

The sign in \cref{p4:eq:core-series} matches \textsf{S3} after
$\varepsilon=-\sigma$; explicitly $w_-=\Lambda_-+\ell+\ell/\Lambda_-+\cdots$
on the even sheet ($\Wlam_{-1}$) and
$w_+=\Lambda_+-\ell+\ell/\Lambda_+-\cdots$ on the odd sheet ($\Wlam_0$).

\begin{proposition}[An explicit convergence criterion]
\label{p4:prop:core-convergence}
Fix $\rho\in(0,1)$ and put $M_\rho(\ell)=|\ell|+\log\frac1{1-\rho}$.  Then
\begin{equation}
  |A_k(\ell)|\le\frac{M_\rho(\ell)^{\,k+1}}{(k+1)\rho^{k}},
  \label{p4:eq:Ak-bound}
\end{equation}
so the series in \cref{p4:eq:core-series} converges absolutely whenever
$M_\rho(\ell)<\rho|\Lambda|$ for some $\rho\in(0,1)$, and its sum is then the
solution branch $w\sim\Lambda$.  In particular, for real $\Lambda\ge10$ the
choice $\rho=\tfrac12$ already gives convergence.
\end{proposition}

\begin{proof}
On $|z|=\rho$ one has $|\log(1+z)|\le\sum_{k\ge1}\rho^{k}/k
=\log\frac1{1-\rho}$, so $|\ell+\log(1+z)|\le M_\rho(\ell)$ there; Cauchy's
estimate applied to \cref{p4:eq:Ak-def} gives \cref{p4:eq:Ak-bound}.  Summing
the geometric-type majorant gives absolute convergence for
$M_\rho(\ell)<\rho|\Lambda|$.  The sum solves
$w+\sigma\log w=\Lambda$ because Lagrange's theorem identifies the series with
the unique analytic solution $\upsilon(z)$ of
$\upsilon=-\sigma z(\ell+\log(1+\upsilon))$ vanishing at $z=0$; the branch
with $w\sim\Lambda$ is the one selected.  For $\Lambda=10$,
$\ell=\log10=2.3026$ and $M_{1/2}=2.3026+0.6931=2.9957<5$.
\end{proof}

\begin{repairbox}
\textbf{Source claim.}  \textsf{S3} asserts that ``the pure Lambert
logarithmic expansion has a nonzero convergence domain for sufficiently large
argument, as in the classical analysis of Lambert $W$'', citing Corless et
al.\ but giving no criterion.  \textbf{Defect.}  The cited classical result
concerns the doubly-indexed series in $(\log z)^{-1}$ and $\log\log z$; the
single series \cref{p4:eq:core-series}, whose coefficients are polynomials in
the slowly varying $\ell$, needs its own statement, since convergence depends
on the ratio $\ell/\Lambda$ and fails for $\ell$ comparable to $\Lambda$.
\textbf{Repair.}  \Cref{p4:prop:core-convergence} supplies an explicit,
proved criterion, which also delimits the regime where
\cref{p4:eq:core-series} may be summed rather than truncated.
\end{repairbox}

\subsection{Composition, and the equivalence of the two normalizations}

\begin{theorem}[Constructive equivalence of the two inverse normalizations]
\label{p4:thm:equivalence}
With $v_j=(-\sigma)^{j}A_{j-1}(\ell)$ and
\begin{equation}
  R^{(m)}_r(\ell)=\Coef{z^{r}}\Bigl(1+\sum_{j\ge1}v_jz^{j}\Bigr)^{-m}
  =\sum_{k=0}^{r}(-1)^{k}\binom{m+k-1}{k}\OB_{r,k}(v),
  \qquad R^{(m)}_0=1,
  \label{p4:eq:R-composition}
\end{equation}
the Lambert-normalized expansion \cref{p4:eq:inverse-analytic} and the
flattened expansion \cref{p4:eq:flattened} are term-by-term equivalent: in the
scaled variables of \cref{p4:eq:scaled-core},
\begin{equation}
  \mathcal I_\sigma(y)\sim\frac{\Lambda_\sigma}{2\lambda}
   -\frac{\sigma\,\ell_\sigma}{2\lambda}
   +\sum_{N\ge1}\frac{\Pi^{(\sigma)}_N(\ell_\sigma)}{\Lambda_\sigma^{\,N}},
  \qquad
  \Pi^{(\sigma)}_N(\ell)=\frac{(-\sigma)^{N+1}A_N(\ell)}{2\lambda}
   +\sum_{m=1}^{N}c^{(\sigma)}_m(2\lambda)^{m}R^{(m)}_{N-m}(\ell),
  \label{p4:eq:Pi-composition}
\end{equation}
with $\deg\Pi^{(\sigma)}_N=N$ and leading coefficient
$(-1)^{N-1}(-\sigma)^{N+1}/(2\lambda N)$.
\end{theorem}

\begin{proof}
Insert $\Xcore_\sigma=w_\sigma/(2\lambda)$ into
\cref{p4:eq:inverse-analytic}, write $w=\Lambda(1+V)$ with
$V=\sum_jv_jz^{j}$ from \cref{p4:eq:core-series}, and collect the coefficient
of $\Lambda^{-N}$: the core contributes the first term of
\cref{p4:eq:Pi-composition} and the $m$-th correction contributes
$c^{(\sigma)}_m(2\lambda)^{m}\Coef{z^{N-m}}(1+V)^{-m}$.  The Bell form of
$R^{(m)}_r$ is \cref{p0:lem:power-log}.  Degrees: $A_N$ has degree $N$ with
leading coefficient $(-1)^{N-1}/N$ by \cref{p4:eq:Ak-stirling}, while every
correction with $m\ge1$ contributes degree at most $N-m$.
\end{proof}

\begin{repairbox}
\textbf{Source claim.}  \textsf{S1} asserts the equivalence of its two inverse
normalizations in prose: ``expanding the Lambert carrier in the scale
$X^{-1}\Real[\ell]$ and then substituting that expansion \ldots\ reproduces
the recurrence coefficient by coefficient.''  \textbf{Defect.}  No proof and
no formula; the assertion is exactly the sort of ``successive matching''
that the coefficient calculus is supposed to eliminate.  \textbf{Repair.}
\Cref{p4:thm:equivalence}, which follows \textsf{S3}, makes the equivalence a
finite explicit composition with a proof, and identifies the degrees and
leading coefficients.
\end{repairbox}

\begin{remark}[Three flattening normalizations, and which one the volume uses]
\label{p4:rem:three-flattenings}
The three sources flatten in three different variables: \textsf{S2} in
$L_\sigma=\log(y/C_\sigma)$ with $H_\sigma=\log(L_\sigma/\lambda)$;
\textsf{S3} in $\Lambda_\sigma=\log(\pi y^{2})+\sigma\log(4\lambda)$ with
$\ell_\sigma=\log\Lambda_\sigma$; \textsf{S1} in
$X_\epsilon=\log(y\alpha^{\beta_\epsilon}/K_\epsilon)$ in the half index.
Since $\Lambda_\sigma=2L_\sigma+\sigma(\log(4\lambda)-\lambda)\ne2L_\sigma$,
these are genuinely different variables, the three coefficient lists are not
comparable termwise, and none is in error --- all were recomputed and are
correct in their own normalization.  The volume adopts \textsf{S2}'s variables
in \cref{p4:eq:flattened} (they are those of \cref{p0:thm:flattening}) and
records \textsf{S3}'s in \cref{p4:eq:Pi-composition}, where the composition
proof is cleanest.
\end{remark}

\begin{corollary}[Leading elementary forms]
\label{p4:cor:leading-inverse}
As $y\to\infty$,
\begin{equation}
  \mathcal I_\sigma(y)=\frac1\lambda\Bigl(L_\sigma(y)
   -\frac\sigma2\log\frac{L_\sigma(y)}{\lambda}
   +O\Bigl(\frac{\log L_\sigma(y)}{L_\sigma(y)}\Bigr)\Bigr).
  \label{p4:eq:leading-inverse}
\end{equation}
The opposite signs of the $\tfrac12\log L$ term on the two sheets are the
inverse-side signature of the forward powers $n^{-1/2}$ and $n^{+1/2}$.
\end{corollary}

\section{Asymptotic meaning, optimal truncation, and Stokes structure}
\label{p4:sec:stokes}

\subsection{What each expansion means at fixed order}

For each fixed $M$ the forward statement is \cref{p4:eq:forward} with the
effective constant of \cref{p4:prop:forward-effective}, and the inverse
statement is \cref{p4:eq:inverse-analytic}.  Both are Poincar\'e statements
about divergent series: the coefficients $d_j$ grow factorially
(\cref{p4:prop:leastterm}), hence so do the $A_n$ and, through the reversion,
the $c^{(\sigma)}_n$.  The divergence is not a defect but a datum: the Borel
transform and the exact sum of the logarithmic series are known in closed form
(\cref{p4:thm:laplace}, \cref{p4:prop:borel-lattice}).  By contrast the pure
core series \cref{p4:eq:core-series} \emph{converges} in the regime delimited
by \cref{p4:prop:core-convergence}; the full inverse is asymptotic only
because of the gamma-ratio correction.

\subsection{Optimal truncation and its image in the inverse variable}

\begin{proposition}[Least-term truncation, forward and inverse]
\label{p4:prop:optimal}
\Cref{p0:thm:optimal-truncation} applied to \cref{p4:eq:D-first} with the
singulant modulus $\pi$ of \cref{p4:prop:borel-lattice} gives the least-term
index $r\approx\pi x/2$ and the beyond-all-orders floor
\begin{equation}
  \bigl|R_{N}(x)\bigr|_{N\approx\pi x/2}
  =\frac{2\sqrt2}{\pi\sqrt x}\,\e^{-\pi x}\bigl(1+O(x^{-1})\bigr),
  \label{p4:eq:floor-forward}
\end{equation}
which by \cref{p4:thm:envelope} is an equality up to the stated factor and not
merely an upper bound.  Transporting it through the inverse by
\cref{p0:thm:remainder-transport} and
$\bigl(\log\mathcal S_\sigma\bigr)'=\lambda+O(x^{-1})$ gives the inverse floor
\begin{equation}
  O\!\left(\Xcore_\sigma^{-1/2}\e^{-\pi \Xcore_\sigma}\right)
  =O\!\left(y^{-\pi/\lambda}\,
    (\log y)^{-1/2+\sigma\pi/(2\lambda)}\right),
  \qquad
  \frac\pi\lambda=4.5323601418\ldots
  \label{p4:eq:floor-inverse}
\end{equation}
\end{proposition}

\begin{proof}
\Cref{p4:eq:floor-forward} is \cref{p4:prop:leastterm}.  For the image, use the
core equation in the form $\e^{\lambda \Xcore}=(y/C_\sigma)\Xcore^{-\sigma/2}$,
so that
$\e^{-\pi \Xcore}=(y/C_\sigma)^{-\pi/\lambda}\Xcore^{\sigma\pi/(2\lambda)}$;
multiply by $\Xcore^{-1/2}$ and use $\Xcore\asymp\log y$.
\end{proof}

\begin{remark}
Exponential smallness in the index becomes \emph{algebraic} smallness in the
level: an $\e^{-\pi x}$ forward sector maps to a $y^{-\pi/\lambda}$ inverse
sector.  This is why hyperasymptotic corrections to the inverse live on a
strictly finer grid than the algebraic powers of $1/\log y$ produced by
\cref{p4:eq:flattened} --- they are invisible to that scale at every order.
\end{remark}

\subsection{Lateral sums and the aggregate Stokes multipliers}

For a direction $\theta$ let
$\mathcal S_{\theta^{\pm}}D(x)
 =\int_0^{\e^{\ii(\theta\pm0)}\infty}\e^{-xt}\widehat D(t)\,\dd t$
whenever the rotated Laplace integral converges, and let
$\operatorname{Disc}_\theta$ denote the lower-minus-upper lateral difference.

\begin{proposition}[Bridge equations]
\label{p4:prop:stokes}
With the poles and residues of \cref{p4:eq:borel-poles},
\begin{align}
  \operatorname{Disc}_{\pi/2}D(x)
   &=2\sum_{k\ge0}\frac{\e^{-\ii\pi(2k+1)x}}{2k+1}
   =2\operatorname{artanh}\bigl(\e^{-\ii\pi x}\bigr)
   \qquad(\Im x<0),
  \label{p4:eq:upper-stokes}\\
  \operatorname{Disc}_{-\pi/2}D(x)
   &=-2\sum_{k\ge0}\frac{\e^{\ii\pi(2k+1)x}}{2k+1}
   =-2\operatorname{artanh}\bigl(\e^{\ii\pi x}\bigr)
   \qquad(\Im x>0),
  \label{p4:eq:lower-stokes}
\end{align}
the closed forms giving the analytic continuations beyond those half-planes.
Consequently
$\operatorname{Disc}_\theta\log\mathcal S_\sigma
 =\sigma\operatorname{Disc}_\theta D$, and the multiplicative Stokes factor of
the forward branch across the ray $\theta=\pm\pi/2$ is
$\exp\bigl(\sigma\operatorname{Disc}_{\pm\pi/2}D(x)\bigr)$, that is
\begin{equation}
  \mathfrak M^{(\sigma)}_{+}(x)
   =\left(\frac{1+\e^{-\ii\pi x}}{1-\e^{-\ii\pi x}}\right)^{\!\sigma},
  \qquad
  \mathfrak M^{(\sigma)}_{-}(x)
   =\left(\frac{1-\e^{\ii\pi x}}{1+\e^{\ii\pi x}}\right)^{\!\sigma}.
  \label{p4:eq:stokes-multipliers}
\end{equation}
On the positive real axis no lateral choice arises, because $\arg t=0$ is a
regular direction.
\end{proposition}

\begin{proof}
Rotating the Laplace contour past the ray through $t_k=\ii\pi(2k+1)$ picks up
$2\pi\ii$ times the residue of $\e^{-xt}\widehat D(t)$ there, namely
$2\pi\ii\,\e^{-t_kx}/t_k=2\e^{-\ii\pi(2k+1)x}/(2k+1)$; summing over $k$ and
using $\operatorname{artanh}z=\sum_{k\ge0}z^{2k+1}/(2k+1)$ gives
\cref{p4:eq:upper-stokes}.  The series converges for
$|\e^{-\ii\pi x}|=\e^{\pi\Im x}<1$, that is $\Im x<0$.  The lower ray is the
conjugate computation.  The last two displays follow from
$\log\mathcal S_\sigma=\log C_\sigma+\lambda x+\tfrac\sigma2\log x+\sigma D$,
in which only $D$ is discontinuous, together with
$2\operatorname{artanh}z=\log\frac{1+z}{1-z}$.
\end{proof}

\begin{remark}[Convention dependence]
\label{p4:rem:stokes-convention}
\textsf{S2} defines the discontinuity as lower minus upper lateral sum;
\textsf{S3} instead adopts the residue convention in which crossing a singular
ray contributes $+2\pi\ii$ times the enclosed residues, and works with
$Q=-D$.  The two sets of formulae differ by an overall sign, hence by a
reciprocal in the multipliers, and both are internally consistent.  A Stokes
generator is only defined relative to a stated convention, so the volume
states the convention explicitly in \cref{p4:prop:stokes} and does not declare
either source's signs canonical.
\end{remark}

\subsection{Transport of an exponential sector through inversion}

\begin{proposition}[Sector transport]
\label{p4:prop:sector-transport}
Let $F$ be one lateral branch of $\log\mathcal S_\sigma$ and let the other be
$F+\Delta$.  As a formal expansion in the sector $\Delta$,
\begin{equation}
  x_{\mathrm{new}}=x+\sum_{k\ge1}\frac{(-1)^{k}}{k!}
   \left(\frac1{F'(x)}\frac{\dd}{\dd x}\right)^{\!k-1}
   \left(\frac{\Delta(x)^{k}}{F'(x)}\right),
  \label{p4:eq:sector-transport}
\end{equation}
with $\Delta=\sigma\operatorname{Disc}_\theta D$ and
$F'(x)=\lambda+\tfrac{\sigma}{2x}+\sigma D'(x)$.  Quantitatively, and with no
formal series at all: if $|\Delta|\le\delta$ on the interval joining $x$ and
$x_{\mathrm{new}}$ and $F'\ge\mu>0$ there, then
\begin{equation}
  |x_{\mathrm{new}}-x|\le\frac{\delta}{\mu},
  \qquad
  x_{\mathrm{new}}-x=-\frac{\Delta(x)}{F'(x)}+O(\delta^{2}),
  \label{p4:eq:sector-transport-quant}
\end{equation}
the implied constant depending only on $\sup|(F^{-1})''|$ on that interval.
\end{proposition}

\begin{proof}
\Cref{p4:eq:sector-transport} is \cref{p4:lem:phase-lb} with $F_0=F$,
$g=\Delta$ and $\eta=1$, read as an identity of formal series in the sector
marker.  The first inequality in \cref{p4:eq:sector-transport-quant} is
\cref{p0:thm:backward-error}; the second is Taylor's theorem for $F^{-1}$ at
the point $F(x)$.
\end{proof}

\begin{repairbox}
\textbf{Source claim.}  \textsf{S1} states its transport rule as a
\emph{theorem} (``Transseries-sector transport through inversion'') whose
hypothesis is that ``$E$ is smaller than the algebraic scale retained in
$\Phi_{\mathrm{alg}}$'', whose conclusion is prefaced by the word
``formally'', and whose displayed first terms carry an unjustified error
symbol $O(E^{3},E^{2}E',E(E')^{2})$.  \textbf{Defect.}  The hypothesis is
never used; the conclusion is a formal-series identity, not an analytic one;
and the error term is asserted rather than proved.  A statement that holds
only formally must not stand in a \texttt{theorem} environment.
\textbf{Repair.}  \Cref{p4:prop:sector-transport} separates the two claims:
the all-orders identity is formal (and is exactly \cref{p4:lem:phase-lb}),
while the quantitative first-order statement is proved by the mean value
theorem under stated hypotheses.  \textsf{S2} and \textsf{S3} already state
their versions correctly, as formal expansions.
\end{repairbox}

\begin{repairbox}
\textbf{Source claim.}  \textsf{S1}'s section on complex inverse sheets
asserts that for every Lambert branch with $|w_{\epsilon,k}|\to\infty$ in a
suitable sector, ``the same coefficients $d^{(\epsilon)}_r$ give
$M_{\epsilon,k}(y)\sim w_{\epsilon,k}(y)+\sum_r d^{(\epsilon)}_r
w_{\epsilon,k}(y)^{-r}$''.  \textbf{Defect.}  For $k\ne0,-1$ nothing has been
said, let alone proved, about $M_{\epsilon,k}$ being an inverse of anything:
the existence of an analytic branch of $\mathcal S_\sigma^{-1}$ along the
corresponding path is a hypothesis, not a consequence, and it fails when the
path meets a pole of the gamma quotient or a Stokes direction.
\textbf{Repair.}  The statement is demoted to
\cref{p4:rem:complex-branches}, with the two missing hypotheses displayed.
\textsf{S2} states the same thing correctly, as a \emph{formal} complex branch
whose analytic validity is a separate requirement.
\end{repairbox}

\begin{remark}[Complex branches, formally]
\label{p4:rem:complex-branches}
For $k\in\Int$ let $\Xcore_{\sigma,k}$ be as in \cref{p4:eq:core-general}
after fixing branches of the powers and logarithms in $L_\sigma$.  Suppose
that along a path with $|\Xcore_{\sigma,k}(y)|\to\infty$ \emph{(i)} the
forward expansion \cref{p4:eq:forward} is valid, which requires the path to
stay in a sector free of the poles of $\Gamma(x/2+(3-\sigma)/4)^{-2}$ and away
from the Stokes directions of \cref{p4:prop:stokes}, and \emph{(ii)} an
analytic branch $\mathcal I_{\sigma,k}$ of the inverse exists along it.  Then
that branch satisfies
$\mathcal I_{\sigma,k}(y)\sim \Xcore_{\sigma,k}(y)
 +\sum_{n\ge1}c^{(\sigma)}_n\Xcore_{\sigma,k}(y)^{-n}$
with the \emph{same} coefficients, by the proof of
\cref{p4:thm:inverse-analytic}.  The coefficient sequence does not depend on
$k$: all branch dependence sits in $\Wlam_k$ and in the chosen logarithm.
Without hypotheses (i) and (ii) the display is a formal statement only.
\end{remark}

\section{Discrete inversion and index recovery}
\label{p4:sec:discrete}

\subsection{The three inverse objects, parity by parity}

\Cref{p0:def:three-inverses} attaches three objects to a discrete sequence:
the functional inverse of a chosen interpolation, the integer staircase, and
the round-to-nearest generalized inverse.

\begin{warning}[None of the three exists for the unsplit sequence]
\label{p4:rem:no-global}
By \cref{p4:prop:ratios} the set $\{n:\operatorname{sw}(n)\le y\}$ is not an
initial segment of $\Nat_0$ --- for instance
$\operatorname{sw}(5)=30>20=\operatorname{sw}(6)$ --- so the staircase of
$\operatorname{sw}$ is not monotone and does not invert it; the
round-to-nearest inverse is ambiguous whenever two indices of opposite parity
are nearly equidistant in value; and by \cref{p4:prop:oscillatory} no
eventually monotone interpolation exists at all.  All three objects exist, are
unique, and obey \cref{p0:thm:staircase} \emph{on each parity class
separately}.
\end{warning}

\begin{proposition}[Exact branchwise staircases]
\label{p4:prop:staircase}
For $y\ge1$ put
\begin{equation}
  N^{\le}_-(y)=\max\{n\ \text{even}:\operatorname{sw}(n)\le y\},
  \qquad
  N^{\le}_+(y)=\max\{n\ \text{odd}:\operatorname{sw}(n)\le y\},
  \label{p4:eq:Nle-def}
\end{equation}
both sets being non-empty because $\operatorname{sw}(0)=\operatorname{sw}(1)=1$.
Then
\begin{equation}
  \boxed{\;
  N^{\le}_-(y)=2\left\lfloor\frac{\mathcal I_-(y)}{2}\right\rfloor,
  \qquad
  N^{\le}_+(y)=2\left\lfloor\frac{\mathcal I_+(y)-1}{2}\right\rfloor+1,\;}
  \label{p4:eq:staircase}
\end{equation}
and the least index of each parity that reaches the level $y$ is
\begin{equation}
  N^{\ge}_-(y)=2\left\lceil\frac{\mathcal I_-(y)}{2}\right\rceil,
  \qquad
  N^{\ge}_+(y)=2\left\lceil\frac{\mathcal I_+(y)-1}{2}\right\rceil+1 .
  \label{p4:eq:staircase-up}
\end{equation}
These identities are \emph{exact}, not asymptotic.  The quantity
\cref{p4:eq:staircase-up} is the staircase $N_\ast$ of
\cref{p0:def:three-inverses}, and the displayed formula for it is
\cref{p0:thm:staircase}'s exact-rounding identity
$N_\ast=\lceil\nu\rceil$, applied to the bisection
$(\operatorname{sw}(2m+\epsilon))_{m\ge0}$ with its admissible interpolation
$m\mapsto\mathcal S_\sigma(2m+\epsilon)$.
\end{proposition}

\begin{proof}
By \cref{p4:eq:branch-interpolates} and \cref{p4:prop:monotone},
$\mathcal S_-$ is a strictly increasing bijection agreeing with
$\operatorname{sw}$ at the even integers, so for even $n$ we have
$\operatorname{sw}(n)\le y\iff\mathcal S_-(n)\le y\iff n\le\mathcal I_-(y)$.
The largest such even $n$ is $2\lfloor\mathcal I_-(y)/2\rfloor$ and the least
even $n\ge\mathcal I_-(y)$ is $2\lceil\mathcal I_-(y)/2\rceil$.  The odd case
is identical after the shift $n\mapsto n-1$.  The boundary case in which
$\mathcal I_\sigma(y)$ is an integer of the relevant parity is covered,
because then $y$ is exactly a sequence value and floor and ceiling return it.
\end{proof}

The round-to-nearest inverses on the two parity classes are
\begin{equation}
  n^{\mathrm{near}}_-(y)
   =2\left\lfloor\frac{\mathcal I_-(y)}2+\frac12\right\rfloor,
  \qquad
  n^{\mathrm{near}}_+(y)
   =2\left\lfloor\frac{\mathcal I_+(y)-1}2+\frac12\right\rfloor+1 .
  \label{p4:eq:nearest}
\end{equation}

\begin{remark}[The rounding must be certified, not trusted]
\label{p4:rem:certify-rounding}
\Cref{p4:eq:staircase}--\cref{p4:eq:nearest} are exact statements about the
real numbers $\mathcal I_\sigma(y)$, whereas a truncated transseries supplies
only an approximation to them.  When $y$ lies exponentially close to a
sequence value the floor can flip.  The remedy is cheap: use the transseries
to produce a candidate integer, then decide the inequality
$\operatorname{sw}(n)\le y$ exactly, either in integer arithmetic or by
comparing $\log\mathcal S_\sigma(n)$ with $\log y$ at the candidate and its
same-parity neighbour, working through $\log\Gamma$ to stay in range.
\Cref{p0:thm:staircase} is the general separation condition and
\cref{p0:thm:backward-error} converts the logarithmic residual into a
certified bracket for $\mathcal I_\sigma(y)$.
\end{remark}

\subsection{Which parity wins, and by how much}

\begin{theorem}[The sheet gap]
\label{p4:thm:sheet-gap}
Put $Y=\log(y\sqrt\pi)$ and $h=\log(Y/\lambda)$, so that
$L_\sigma=Y+\sigma\lambda/2$.  Then, as $y\to\infty$,
\begin{equation}
  \boxed{\;
  \mathcal I_-(y)-\mathcal I_+(y)
   =\frac h\lambda-1+\frac1{2Y}+O\!\left(\frac{h^{2}}{Y^{2}}\right)
   =\log_2\!\left(\frac{\log(y\sqrt\pi)}{\log2}\right)-1+o(1).\;}
  \label{p4:eq:sheet-gap}
\end{equation}
In particular the gap tends to $+\infty$: every sufficiently large level is
reached first on the \emph{odd} sheet, about $\log_2\log y$ indices before the
even sheet catches up.
\end{theorem}

\begin{proof}
Subtract the two instances of \cref{p4:eq:flattened} truncated after $n=1$.
The leading terms contribute $(L_--L_+)/\lambda=-1$.  Next,
$H_\sigma=\log\bigl((Y+\sigma\lambda/2)/\lambda\bigr)
 =h+\sigma\lambda/(2Y)+O(Y^{-2})$, and the two contributions
$-(\sigma/2\lambda)H_\sigma$ combine to
$(H_-+H_+)/(2\lambda)=h/\lambda+O(Y^{-2})$, the $\sigma$-odd parts
cancelling.  Finally, with $P^{(\sigma)}_1=(H-\sigma\lambda)/4$ from
\cref{p4:eq:P1},
\[
  \frac1\lambda\left[\frac{P^{(-1)}_1(H_-)}{L_-}
   -\frac{P^{(+1)}_1(H_+)}{L_+}\right]
  =\frac1{4\lambda}\left[\frac{H_-+\lambda}{L_-}-\frac{H_+-\lambda}{L_+}
   \right]
  =\frac1{4\lambda}\cdot\frac{2\lambda}{Y}+O\!\left(\frac{h}{Y^{2}}\right)
  =\frac1{2Y}+O\!\left(\frac{h}{Y^{2}}\right).
\]
Collecting the three displays gives \cref{p4:eq:sheet-gap}; the second form
follows from $h/\lambda=\log_2(Y/\lambda)$ and $1/(2Y)\to0$.  The error term
$O(h^{2}/Y^{2})$ absorbs the $n=2$ terms of \cref{p4:eq:flattened}, whose
polynomials have degree $2$ in $H_\sigma$.
\end{proof}

\begin{remark}[Numerical confirmation, and agreement with \textsf{S1}]
\label{p4:rem:gap-numeric}
Recomputed with eight Lambert-normalized corrections on each sheet: for
$\log y=10,30,100,400$ the gap $\mathcal I_-(y)-\mathcal I_+(y)$ equals
$2.97806$, $4.47863$, $6.18565$, $8.17591$, against the prediction
$h/\lambda-1+1/(2Y)$ equal to $2.97829$, $4.47928$, $6.18583$, $8.17593$.
\textsf{S1} states the same result in its half index as
$N_0-N_1=\tfrac2\alpha\log\log y+\tfrac2\alpha\log\tfrac2\alpha-1+o(1)$ with
$\alpha=2\lambda$; since $\tfrac2\alpha=\tfrac1\lambda$ and
$\log\tfrac2\alpha=-\log\lambda$, this reads
$\tfrac1\lambda(\log\log y-\log\lambda)-1+o(1)$, which is
\cref{p4:eq:sheet-gap}.  Directly from the sequence, the least index reaching
$y=10,10^{3},10^{6},10^{12}$ is $5,11,21,39$ on the odd sheet against
$6,14,24,44$ on the even sheet: gaps $1,3,3,5$, consistent with the very slow
$\log_2\log y$ growth.
\end{remark}

\begin{corollary}[Counting function]
\label{p4:cor:counting}
Let $N(y):=\#\{n\in\Nat_0:\operatorname{sw}(n)\le y\}$.  Then
\begin{equation}
  N(y)=\frac{\mathcal I_-(y)+\mathcal I_+(y)}{2}+O(1)
   =\log_2\bigl(y\sqrt\pi\bigr)+O(1).
  \label{p4:eq:counting}
\end{equation}
\end{corollary}

\begin{proof}
By \cref{p4:prop:staircase} the even indices counted are
$0,2,\ldots,2\lfloor\mathcal I_-(y)/2\rfloor$, that is
$\mathcal I_-(y)/2+O(1)$ of them, and similarly there are
$\mathcal I_+(y)/2+O(1)$ odd ones; this is the first equality.  For the
second, average the two instances of \cref{p4:eq:flattened}: as in the proof
of \cref{p4:thm:sheet-gap} the $\sigma$-odd term $-(\sigma/2\lambda)H_\sigma$
cancels in the average, leaving
$\tfrac1{2\lambda}(L_-+L_+)+O(\log\log y/\log y)=Y/\lambda+O(1)$, and
$Y/\lambda=\log_2(y\sqrt\pi)$.
\end{proof}

\section{Algorithms, certification, and validation}
\label{p4:sec:algorithms}

\subsection{Generating the coefficients}

\begin{algorithm}[Exact coefficient generation]
\label{p4:alg:coefficients}
Every coefficient family of this chapter is produced by a finite exact
recurrence over $\Rat[\lambda,\lambda^{-1},\sigma]/(\sigma^{2}-1)$:
\begin{enumerate}
\item $d_{2r+1}$ from \cref{p4:eq:d-bernoulli} (rational Bernoulli input;
      $d_{2r}=0$);
\item $A_n$ from the triangular recurrence \cref{p4:eq:A-recurrence};
\item $c^{(\sigma)}_n$ from \cref{p4:eq:c-bell}, or by solving
      $U=\Psi_\sigma(t,U)$ coefficientwise, or from the differential formula
      \cref{p4:eq:c-differential};
\item $P^{(\sigma)}_n(H)$ from \cref{p0:thm:flattening} with the dictionary
      \cref{p4:eq:dictionary}, or from \cref{p4:eq:Pi-composition} after
      \cref{p4:prop:core-series};
\item the exact certificate of \cref{p4:rem:certificate} after every step.
\end{enumerate}
Truncating every intermediate series at order $n$ is exact for the $n$-th
coefficient, because all three recurrences are triangular; at order $N$ the
cost is $\widetilde{O}(N^{2})$ coefficient operations.  No numerical
fitting and no order-specific ansatz occurs anywhere.
\end{algorithm}

\begin{remark}[How this chapter's tables were produced]
\label{p4:rem:how-computed}
The three sources supply three implementations of the above.  For this chapter
every coefficient was regenerated by a fourth, independent one: truncated
power-series arithmetic over exact rationals with $\lambda$ symbolic, run
separately at $\sigma=\pm1$ and recombined in $\mathscr R$.  Every displayed
coefficient agreed with the source that printed it; the corrections recorded
in this chapter concern statements, never numbers.
\end{remark}

\subsection{Certified numerical inversion}

\begin{proposition}[Derivative bounds for certification]
\label{p4:prop:derivative-bounds}
For $x\ge1$,
\begin{equation}
  \bigl(\log\mathcal S_+\bigr)'(x)>\lambda+\frac1{2x}-\frac1{4x^{2}}>\lambda,
  \qquad
  \bigl(\log\mathcal S_-\bigr)'(x)>\lambda-\frac1{2x},
  \label{p4:eq:derivative-bounds}
\end{equation}
the exact value being \cref{p4:eq:logder}.  Hence if
$r_*=\log\mathcal S_\sigma(x_*)-\log y$ and
$\bigl(\log\mathcal S_\sigma\bigr)'\ge\mu>0$ on the interval joining $x_*$ and
$\mathcal I_\sigma(y)$, then
$|x_*-\mathcal I_\sigma(y)|\le|r_*|/\mu$ by \cref{p0:thm:backward-error}.
\end{proposition}

\begin{proof}
$\bigl(\log\mathcal S_\sigma\bigr)'(x)=\lambda+\tfrac\sigma{2x}+\sigma D'(x)$
by \cref{p4:eq:normal-form}; apply $0<-D'(x)<1/(4x^{2})$ from
\cref{p4:cor:D-bounds}, and $1/(2x)-1/(4x^{2})>0$ for $x>1/2$.
\end{proof}

\begin{algorithm}[Branchwise inversion of a level]
\label{p4:alg:invert}
Given a large $y$ and a parity $\sigma$:
\begin{enumerate}
\item Evaluate the Lambert core:
      $\Xcore_+=\Wlam_0(4\pi\lambda y^{2})/(2\lambda)$ or
      $\Xcore_-=-\Wlam_{-1}(-4\lambda/(\pi y^{2}))/(2\lambda)$.
      \emph{The branch choice is not optional} (\cref{p4:prop:two-roots}).
\item Add corrections $\sum_n c^{(\sigma)}_n\Xcore^{-n}$ by Horner nesting in
      $1/\Xcore$, stopping at or before the least term rather than summing
      indefinitely; alternatively use the centred form
      \cref{p4:eq:centred-inverse}, whose bare core is already accurate to
      $O(\mathcal U^{-2})$.
\item Certify: compute $r_*=\log\mathcal S_\sigma(x_*)-\log y$ through
      $\log\Gamma$ and bound the error by
      \cref{p4:prop:derivative-bounds}.  The inequalities
      \cref{p4:eq:odd-bracket} and \cref{p4:eq:even-bracket} give an
      independent one-sided check that costs nothing.
\item For higher accuracy run Newton's method on the logarithmic equation,
      \[
        x_{j+1}=x_j-\frac{\log\mathcal S_\sigma(x_j)-\log y}
        {\lambda+\tfrac\sigma2\bigl[\psi(x_j/2+1)-\psi(x_j/2+1/2)\bigr]},
      \]
      which is stable even when $\mathcal S_\sigma(x)$ is far outside
      floating-point range.  \Cref{p4:prop:monotone} guarantees a positive
      denominator, and the transseries seed lies in the basin for large $y$.
\item For a discrete answer, round by \cref{p4:eq:staircase} and verify the
      final inequality exactly (\cref{p4:rem:certify-rounding}).
\end{enumerate}
\end{algorithm}

\subsection{Validation}

The tables below were recomputed at $50$ significant digits against exact
gamma-quotient reference values, with $y$ generated as
$\mathcal S_\sigma(x_0)$ so that the exact inverse is $x_0$.  They reproduce
the corresponding tables of \textsf{S2} and \textsf{S3} to every printed
digit, and, after the reparametrization of \cref{p4:rem:dictionary}, those of
\textsf{S1} as well.

\begin{table}[htbp]
\centering
\caption{Relative error of the forward expansion \cref{p4:eq:forward} after
the term $\sigma^{M}A_Mx^{-M}$.  Recomputed.}
\label{p4:tab:forward-errors}
\small
\begin{tabular}{@{}rrrrrrr@{}}
\toprule
$x$&$\sigma$&$M=0$&$M=2$&$M=4$&$M=6$&$M=8$\\
\midrule
$10$&$-1$&$2.5273\!\times\!10^{-2}$&$3.8527\!\times\!10^{-5}$&$4.7159\!\times\!10^{-7}$&$1.4186\!\times\!10^{-8}$&$7.8542\!\times\!10^{-10}$\\
$10$&$+1$&$2.4650\!\times\!10^{-2}$&$3.8626\!\times\!10^{-5}$&$4.7375\!\times\!10^{-7}$&$1.4233\!\times\!10^{-8}$&$7.8709\!\times\!10^{-10}$\\
$20$&$-1$&$1.2573\!\times\!10^{-2}$&$4.8642\!\times\!10^{-6}$&$1.5087\!\times\!10^{-8}$&$1.1546\!\times\!10^{-10}$&$1.6336\!\times\!10^{-12}$\\
$20$&$+1$&$1.2417\!\times\!10^{-2}$&$4.8704\!\times\!10^{-6}$&$1.5121\!\times\!10^{-8}$&$1.1565\!\times\!10^{-10}$&$1.6353\!\times\!10^{-12}$\\
$50$&$-1$&$5.0122\!\times\!10^{-3}$&$3.1226\!\times\!10^{-7}$&$1.5560\!\times\!10^{-10}$&$1.9152\!\times\!10^{-13}$&$4.3651\!\times\!10^{-16}$\\
$50$&$+1$&$4.9872\!\times\!10^{-3}$&$3.1242\!\times\!10^{-7}$&$1.5573\!\times\!10^{-10}$&$1.9164\!\times\!10^{-13}$&$4.3668\!\times\!10^{-16}$\\
\bottomrule
\end{tabular}
\end{table}

\begin{table}[htbp]
\centering
\caption{Absolute error
$\bigl|\mathcal I^{[K]}_\sigma(\mathcal S_\sigma(x_0))-x_0\bigr|$ of the
Lambert-normalized inverse \cref{p4:eq:inverse-analytic} with $K$ corrections
retained; $K=0$ is the bare core.  Recomputed.}
\label{p4:tab:inverse-errors}
\small
\begin{tabular}{@{}rrrrrrr@{}}
\toprule
$x_0$&$\sigma$&$K=0$&$K=2$&$K=4$&$K=6$&$K=8$\\
\midrule
$10$&$-1$&$3.8813\!\times\!10^{-2}$&$1.6744\!\times\!10^{-5}$&$2.9057\!\times\!10^{-7}$&$2.4840\!\times\!10^{-8}$&$1.7249\!\times\!10^{-9}$\\
$10$&$+1$&$3.3589\!\times\!10^{-2}$&$2.2677\!\times\!10^{-4}$&$3.9282\!\times\!10^{-6}$&$5.5615\!\times\!10^{-8}$&$2.0051\!\times\!10^{-9}$\\
$20$&$-1$&$1.8701\!\times\!10^{-2}$&$1.2039\!\times\!10^{-6}$&$6.9374\!\times\!10^{-9}$&$1.8267\!\times\!10^{-10}$&$3.2855\!\times\!10^{-12}$\\
$20$&$+1$&$1.7399\!\times\!10^{-2}$&$3.0220\!\times\!10^{-5}$&$1.3176\!\times\!10^{-7}$&$4.7142\!\times\!10^{-10}$&$4.3075\!\times\!10^{-12}$\\
$50$&$-1$&$7.3186\!\times\!10^{-3}$&$4.2970\!\times\!10^{-8}$&$5.6559\!\times\!10^{-11}$&$2.8529\!\times\!10^{-13}$&$8.3851\!\times\!10^{-16}$\\
$50$&$+1$&$7.1104\!\times\!10^{-3}$&$2.0095\!\times\!10^{-6}$&$1.4043\!\times\!10^{-9}$&$8.1104\!\times\!10^{-13}$&$1.1816\!\times\!10^{-15}$\\
$100$&$-1$&$3.6329\!\times\!10^{-3}$&$3.9487\!\times\!10^{-9}$&$1.6091\!\times\!10^{-12}$&$2.1882\!\times\!10^{-15}$&$1.6191\!\times\!10^{-18}$\\
$100$&$+1$&$3.5808\!\times\!10^{-3}$&$2.5440\!\times\!10^{-7}$&$4.4455\!\times\!10^{-11}$&$6.4401\!\times\!10^{-15}$&$2.3379\!\times\!10^{-18}$\\
\bottomrule
\end{tabular}
\end{table}

\begin{table}[htbp]
\centering
\caption{The same errors for the centred inverse
\cref{p4:eq:centred-inverse}.  The bare core is better by a factor
$\asymp\Xcore$, because $\widehat c^{\,(\sigma)}_1=0$; once corrections are
retained the two normalizations are comparable.  Recomputed.}
\label{p4:tab:centred-errors}
\small
\begin{tabular}{@{}rrrrrrr@{}}
\toprule
$x_0$&$\sigma$&$K=0$&$K=2$&$K=4$&$K=6$&$K=8$\\
\midrule
$10$&$-1$&$8.7330\!\times\!10^{-4}$&$5.5308\!\times\!10^{-5}$&$1.1593\!\times\!10^{-7}$&$2.7056\!\times\!10^{-9}$&$4.4842\!\times\!10^{-11}$\\
$10$&$+1$&$7.6102\!\times\!10^{-4}$&$5.6715\!\times\!10^{-5}$&$2.3489\!\times\!10^{-7}$&$7.9121\!\times\!10^{-9}$&$4.0347\!\times\!10^{-10}$\\
$20$&$-1$&$2.2206\!\times\!10^{-4}$&$7.4914\!\times\!10^{-6}$&$5.2119\!\times\!10^{-9}$&$3.7661\!\times\!10^{-11}$&$3.3070\!\times\!10^{-13}$\\
$20$&$+1$&$2.0696\!\times\!10^{-4}$&$7.5956\!\times\!10^{-6}$&$7.4297\!\times\!10^{-9}$&$6.3564\!\times\!10^{-11}$&$8.0880\!\times\!10^{-13}$\\
$50$&$-1$&$3.5860\!\times\!10^{-5}$&$5.0352\!\times\!10^{-7}$&$6.5130\!\times\!10^{-11}$&$8.3102\!\times\!10^{-14}$&$1.4333\!\times\!10^{-16}$\\
$50$&$+1$&$3.4850\!\times\!10^{-5}$&$5.0641\!\times\!10^{-7}$&$7.5153\!\times\!10^{-11}$&$1.0250\!\times\!10^{-13}$&$2.0275\!\times\!10^{-16}$\\
$100$&$-1$&$8.9913\!\times\!10^{-6}$&$6.3982\!\times\!10^{-8}$&$2.1684\!\times\!10^{-12}$&$7.1217\!\times\!10^{-16}$&$3.2358\!\times\!10^{-19}$\\
$100$&$+1$&$8.8632\!\times\!10^{-6}$&$6.4167\!\times\!10^{-8}$&$2.3300\!\times\!10^{-12}$&$7.9120\!\times\!10^{-16}$&$3.8477\!\times\!10^{-19}$\\
\bottomrule
\end{tabular}
\end{table}

\begin{warning}[These tables are not convergence]
\label{p4:rem:not-convergence}
The decreasing entries at fixed $x_0$ do not indicate convergence as
$K\to\infty$.  Both the forward and the inverse series are divergent, with the
beyond-all-orders floor \cref{p4:eq:floor-inverse}; the tables stop far short
of the least term (at $x=20$ the least term of \cref{p4:eq:D-first} sits at
$r=31$, of size $10^{-28}$).  The non-monotone entries in the odd column at
$x_0=10$ are ordinary behaviour of an asymptotic series at small argument, not
an error.  Beyond the least term one should use the exact Laplace
representation \cref{p4:eq:laplace}, or a Newton step against $\log\Gamma$, or
an exponentially improved representation transported by
\cref{p4:prop:sector-transport}.
\end{warning}

\section{Summary of the coefficient hierarchy}
\label{p4:sec:summary}

The whole calculation is a chain of constructive coefficient transformations,
each with a closed extraction formula and a triangular recurrence:
\begin{equation}
  \boxed{\;
  B_{2r+2}
  \;\xrightarrow{\ \text{gamma ratio}\ }\;
  d_{2r+1}
  \;\xrightarrow{\ \exp,\ \EB\ }\;
  A_n
  \;\xrightarrow{\ \text{Lagrange},\ \OB\ }\;
  c^{(\sigma)}_n
  \;\xrightarrow{\ \text{unfold }\Wlam\ }\;
  P^{(\sigma)}_n(H),\;}
  \label{p4:eq:hierarchy}
\end{equation}
while the Borel kernel $\tanh(t/2)/(2t)$ supplies, in one stroke, the exact
sum of the first stage, the sign and size of every remainder
(\cref{p4:thm:envelope}), the singulant $\pi$, the least-term scale, and the
complete Stokes data.  Two facts are specific to the swing factorial and
cannot be relegated to the general theory of \cref{p0:sec:top}:

\begin{enumerate}
\item \emph{Parity is leading-order data.}  The forward object is a
period-two transseries \cref{p4:eq:integer-forward} whose discrete
transmonomial $\sigma_n$ switches the algebraic scale between $2^{n}n^{-1/2}$
and $2^{n}n^{+1/2}$; the inverse object is a \emph{pair} of monotone branches,
and \cref{p4:prop:oscillatory} shows that this is forced, not chosen.
\item \emph{The two sheets need different Lambert branches.}  $\Wlam_{-1}$ on
the even sheet, $\Wlam_0$ on the odd sheet, the assignment being forced by
$\operatorname{sign}b=\sigma$ in \cref{p0:thm:lambert-core}.  The principal
branch on the even sheet does not merely lose accuracy: by
\cref{p4:prop:two-roots} it returns the root
$2/(\pi y^{2})\bigl(1+O(y^{-2})\bigr)\to0$, so that choice is fatal before any
correction is computed.
\end{enumerate}

%% ==================== fragment p6 ====================

\part{Four special functions}

\chapter{Gamma and Barnes \texorpdfstring{$G$}{G}}
\label{p6:sec:top}

\section{What this chapter does}
\label{p6:sec:intro}

The four subjects of Chapters~\ref{p1:sec:top}--\ref{p4:sec:top} are
\emph{sequences}; the next two are \emph{special functions}, and the change
matters in exactly one place.  A sequence has three inverse objects
(\cref{p0:def:three-inverses}) and the staircase theorem is needed to pass
between them; a function analytic and eventually strictly monotone on a real
ray has one, and the inverse asked for is the honest compositional inverse of
a restriction.  Everything else --- the exponential--power bookkeeping, the
exact inversion of a dominant block, the all-orders reversion around it, the
backward-error certificate --- is the apparatus of Chapter~0 and is cited, not
repeated.

\begin{sourcebox}
This chapter consolidates three filed articles, listed in
\cref{p6:tab:sources}.  They were written independently and they agree: every
coefficient, table and decimal any of them prints was recomputed here in exact
arithmetic or to $60$ digits, and all of it reproduced.  What the three do
\emph{not} share is a normalization --- one absorbs the Barnes constant
$\zeta'(-1)$ into the target and two do not, and the three use three different
letters for the slope --- so \cref{p6:sec:dictionary} gives the dictionary and
the rest of the chapter fixes one convention.

\begin{center}\small
\begin{tabular}{@{}lll@{}}
\toprule
Tag & Directory & Distinctive content\\
\midrule
S1 & \texttt{inverse\_\allowbreak gamma\_\allowbreak barnesG\_\allowbreak transseries} &
 Hahn supports; Newton valuation doubling; the\\
 & & $\sqrt{T^2+1/6}$ subsector\\
S2 & \texttt{inverse\_\allowbreak gamma\_\allowbreak barnes\_\allowbreak transseries} &
 the pole branches of $\Gamma^{-1}$; universal $C_0,C_1,C_2$\\
 & & for general $p$; $\zeta'(-1)$ left unabsorbed\\
S3 & \texttt{Asymptotic\_\allowbreak Inversion\_\allowbreak Gamma\_\allowbreak Barnes\_\allowbreak G} &
 the kernel coefficients $\kappa_n(p,v)$; the slope\\
 & & denominator bound; resonance; $b_3,b_4,b_5$\\
\bottomrule
\end{tabular}
\end{center}
\end{sourcebox}

\begin{table}[htbp]
\centering
\caption{Parameter dictionary for this chapter.}
\label{p6:tab:sources}
\small
\begin{tabular}{@{}lll@{}}
\toprule
 & $\Gamma$ & $G$\\
\midrule
inverted map & $x\mapsto\Gamma(x)$, increasing branch
             & $z\mapsto G(z+1)$, eventual branch\\
centred variable & $t=x-\tfrac12$ & $z=x-1$\\
core exponent $p$ & $1$ & $2$\\
core scale $a$ & $1$ & $\tfrac12$\\
core shift $\beta$ & $1$ & $\tfrac32$\\
absorbed constant $C$ & $\tfrac12\log(2\pi)$ & $\zeta'(-1)$\\
perturbation support & $2\Nat$ & $\Nat$\\
first perturbation & $a_1t^{-1}=-1/(24t)$ & $\alpha z$, $\alpha=\tfrac12\log(2\pi)$\\
core & $T=\e^{\Lambda}$, $\Lambda=w+1$ & $U=\e^{\lambda+1}$, $\lambda=(v+1)/2$\\
inverse support & $T^{-1},T^{-3},T^{-5},\dots$ & $U^{0},U^{-1},U^{-2},\dots$\\
\bottomrule
\end{tabular}
\end{table}

Two features distinguish this chapter from the first four.

\emph{The dominant block is not the one Chapter~0 inverts.}
\Cref{p0:thm:lambert-core} solves $aX+b\log X=L$ exactly; here the unknown
\emph{multiplies} the logarithm, and the block to be inverted is
$a x^{p}(\log x-\beta)$.  That block is still solved exactly by
$\Wlam$, but by a different substitution, and the resulting normalized slope
is $v+1$ rather than $a+b/X$.  \Cref{p6:sec:core} proves the core lemma in the
generality the two subjects and the next chapter need, namely for
$a x^{p}(\log x-\beta)^{r}$.

\emph{The forward phase is a Stirling series, not an
exponential--power model.}  As in \cref{p2:sec:top}, the phase
$\tfrac12z^{2}\log z$ grows faster than $z^{\alpha}$ for every
$\alpha\le1$ and \cref{p0:def:model} does not apply.  What does apply is the
axiomatized dominant core; we say precisely which axioms hold in
\cref{p6:rem:not-a-model} and do not force the model.

The chapter also proves two things none of the three sources states.  The
first is that the slope-denominator bound of \cref{p6:prop:denominator} is
\emph{attained}, with the deepest pole given in closed form
(\cref{p6:thm:deepest-pole}); the second is that the two square roots the
sources find --- S1's and S3's $\sqrt{U^{2}+1/6}$ for Barnes, and the
previously unnoticed $\sqrt{\Lambda^{2}T^{2}+1/12}$ for Gamma --- are the same
phenomenon, the exactly solvable quadratic truncation of the normalized
equation (\cref{p6:lem:quadratic-core}).

\section{Forward normal forms}
\label{p6:sec:forward}

Bernoulli numbers and polynomials follow the conventions of Chapter~0:
$B_n=B_n(0)$ and $z\e^{xz}/(\e^{z}-1)=\sum_{n\ge0}B_n(x)z^n/n!$, so that
$B_1=-\tfrac12$, $B_2=\tfrac16$, $B_4=-\tfrac1{30}$, and
\begin{equation}\label{p6:eq:bernoulli-half}
 B_n\!\left(\tfrac12\right)=\bigl(2^{1-n}-1\bigr)B_n .
\end{equation}

\subsection{The half-shifted Gamma phase}

\begin{lemma}[Symmetry-adapted Gamma normal form]
\label{p6:lem:gamma-normal-form}
Put $t=x-\tfrac12$.  In every closed subsector of $|\arg t|<\pi$, as
$t\to\infty$,
\begin{equation}\label{p6:eq:gamma-normal}
 \log\Gamma\!\left(t+\tfrac12\right)
 \sim t(\log t-1)+\tfrac12\log(2\pi)
 +\sum_{k\ge1}\frac{a_k}{t^{2k-1}},
\end{equation}
where
\begin{equation}\label{p6:eq:ak}
 a_k=\frac{B_{2k}(1/2)}{2k(2k-1)}
    =\frac{\bigl(2^{1-2k}-1\bigr)B_{2k}}{2k(2k-1)} .
\end{equation}
Only odd inverse powers of $t$ occur.  The first values are
\begin{equation}\label{p6:eq:ak-values}
 a_1=-\frac1{24},\quad
 a_2=\frac7{2880},\quad
 a_3=-\frac{31}{40320},\quad
 a_4=\frac{127}{215040},\quad
 a_5=-\frac{511}{608256}.
\end{equation}
\end{lemma}

\begin{proof}
The shifted Stirling expansion is, for fixed $h$ and in the stated sectors,
\begin{equation}\label{p6:eq:shifted-stirling}
 \log\Gamma(t+h)
 \sim\left(t+h-\tfrac12\right)\log t-t+\tfrac12\log(2\pi)
 +\sum_{k\ge2}\frac{(-1)^kB_k(h)}{k(k-1)\,t^{k-1}} .
\end{equation}
Set $h=\tfrac12$.  The coefficient of $\log t$ collapses to $t$ exactly, which
is the entire point of the shift.  By \eqref{p6:eq:bernoulli-half} and
$B_{2m+1}=0$ for $m\ge1$, we get $B_{2m+1}(\tfrac12)=0$ for all $m\ge0$, so
only even $k$ survives; writing $k=2n$ and noting $(-1)^{2n}=1$ gives
\eqref{p6:eq:gamma-normal}--\eqref{p6:eq:ak}.  The five values follow from
$B_2=\tfrac16$, $B_4=-\tfrac1{30}$, $B_6=\tfrac1{42}$, $B_8=-\tfrac1{30}$,
$B_{10}=\tfrac5{66}$.
\end{proof}

\begin{remark}[The shift is not cosmetic]\label{p6:rem:why-half-shift}
Expanding in $x^{-1}$ instead of $t^{-1}$ leaves both parities in the forward
phase, hence a perturbation support $\Nat$ rather than $2\Nat$, hence --- by
\cref{p6:cor:support} --- an inverse displacement with every integer power of
$T^{-1}$ instead of only the odd ones.  The shift is the cheapest available
instance of the normalization rule of \cref{p6:rem:absorb}: translate first,
and only then declare what is small.
\end{remark}

\subsection{The Barnes phase in the centred variable}

\begin{lemma}[Simplified Barnes normal form]
\label{p6:lem:barnes-normal-form}
In every closed subsector of $|\arg z|<\pi$, as $z\to\infty$,
\begin{equation}\label{p6:eq:barnes-normal}
 \log G(z+1)
 \sim\frac12z^{2}\log z-\frac34z^{2}
 +\alpha z-\frac1{12}\log z+\zeta'(-1)
 +\sum_{k\ge1}\frac{c_k}{z^{2k}},
\end{equation}
where
\begin{equation}\label{p6:eq:alpha-ck}
 \alpha=\tfrac12\log(2\pi),
 \qquad
 c_k=\frac{B_{2k+2}}{4k(k+1)},
\end{equation}
so that
$c_1=-\tfrac1{240}$, $c_2=\tfrac1{1008}$, $c_3=-\tfrac1{1440}$,
$c_4=\tfrac1{1056}$.
\end{lemma}

\begin{proof}
Start from the classical form
\[
 \log G(z+1)
 \sim\frac14z^{2}+z\log\Gamma(z+1)
 -\left(\frac12z(z+1)+\frac1{12}\right)\log z-\log A
 +\sum_{k\ge1}\frac{B_{2k+2}}{2k(2k+1)(2k+2)z^{2k}},
\]
$A$ being the Glaisher--Kinkelin constant, and substitute
\[
 \log\Gamma(z+1)
 \sim\left(z+\tfrac12\right)\log z-z+\tfrac12\log(2\pi)
 +\sum_{j\ge1}\frac{B_{2j}}{2j(2j-1)z^{2j-1}} .
\]
Collect by type.  The $z^{2}\log z$ terms give
$z^2\log z-\tfrac12z^2\log z=\tfrac12z^2\log z$; the $z\log z$ terms cancel,
$\tfrac12z\log z-\tfrac12z\log z=0$; the $z^{2}$ terms give
$-z^{2}+\tfrac14z^{2}=-\tfrac34z^{2}$; the linear term is $\alpha z$; and the
lone $\log z$ term is $-\tfrac1{12}\log z$.

Two constants remain.  Multiplying the Gamma tail by $z$ moves its $j=1$ term
to $B_2/2=\tfrac1{12}$, a constant; together with $-\log A$ and
$\log A=\tfrac1{12}-\zeta'(-1)$ this is
$\tfrac1{12}-\tfrac1{12}+\zeta'(-1)=\zeta'(-1)$.

At $z^{-2k}$ two contributions meet: the Gamma tail with $j=k+1$, multiplied
by $z$, contributes $B_{2k+2}/\bigl((2k+2)(2k+1)\bigr)$, and the explicit
Barnes tail contributes $B_{2k+2}/\bigl(2k(2k+1)(2k+2)\bigr)$.  Their sum is
\[
 \frac{B_{2k+2}}{(2k+1)(2k+2)}\left(1+\frac1{2k}\right)
 =\frac{B_{2k+2}}{(2k+2)(2k)}
 =\frac{B_{2k+2}}{4k(k+1)},
\]
which is \eqref{p6:eq:alpha-ck}.
\end{proof}

\begin{remark}[Not an exponential--power model]\label{p6:rem:not-a-model}
Neither phase is of the form of \cref{p0:def:model}.  For $\Gamma$ the phase
is $t\log t$ and for $G$ it is $\tfrac12z^{2}\log z$; in both cases
$\text{phase}/x^{\alpha}\to\infty$ for every $\alpha\le1$, by the argument
already given for the double factorial in \cref{p2:rem:not-exponential-power}.  What
both phases do satisfy is the axiomatized dominant core: each is
$C^{\infty}$ and eventually strictly increasing on a real ray, its derivative
is eventually positive and regularly varying, and the residual after removing
the core admits a Poincar\'e expansion with a differentiable remainder in
sectors.  Those are the only properties used below, and they are what
\cref{p0:thm:backward-error,p0:thm:remainder-transport} require.
\end{remark}

\begin{remark}[Which inverse]\label{p6:rem:which-inverse}
On $(0,\infty)$ the digamma function $\psi$ is strictly increasing, because
$\psi'(x)=\sum_{n\ge0}(n+x)^{-2}>0$, and it runs from $-\infty$ to $+\infty$;
so it has exactly one zero and $\Gamma'=\Gamma\psi$ changes sign exactly once.
Every sufficiently large $y>0$ therefore has exactly two positive preimages.
We write $\Gamma_+^{-1}$ for the branch tending to $+\infty$, which is the
subject of \cref{p6:sec:gamma}, and $\Gamma_0^{-1}$ for the branch tending to
$0^{+}$, treated in \cref{p6:sec:pole-branches}.  These are genuinely
different objects: the first has a divergent asymptotic transseries, the
second a convergent power series in $1/y$.

For Barnes, $(\log G(z+1))'=\tfrac12+z+\tfrac12\log(2\pi)-z\psi(z+1)+\dots$;
rather than track that expression we use the cleaner statement of
\cref{p6:prop:barnes-monotone}, which gives an explicit ray on which
$z\mapsto G(z+1)$ is strictly increasing.
\end{remark}

\begin{proposition}[An explicit ray of monotonicity]
\label{p6:prop:barnes-monotone}
For $z>0$,
\begin{equation}\label{p6:eq:barnes-log-derivative}
 \frac{\dd}{\dd z}\log G(z+1)
 =\frac12\log(2\pi)-\frac12-z+z\,\psi(z+1)
 =\frac12\log(2\pi)+\frac12-z+z\,\psi(z).
\end{equation}
Consequently $z\mapsto G(z+1)$ is strictly increasing on $[3,\infty)$, and
\begin{equation}\label{p6:eq:barnes-slope-bound}
 \frac{\dd}{\dd z}\log G(z+1)\ \ge\ \tfrac12\,z\log z
 \qquad(z\ge7).
\end{equation}
It therefore has a well-defined inverse
$\mathcal{B}:[G(4),\infty)\to[3,\infty)$, and $G_+^{-1}(y)=1+\mathcal{B}(y)$.
\end{proposition}

\begin{proof}
Write the Weierstrass form
\[
 \log G(1+z)=\frac z2\log(2\pi)-\frac{z}{2}-\frac{z^{2}(1+\gamma)}2
 +\sum_{k\ge1}\left[k\log\Bigl(1+\frac zk\Bigr)-z+\frac{z^{2}}{2k}\right],
\]
which converges locally uniformly, so it may be differentiated termwise:
\[
 \frac{\dd}{\dd z}\log G(1+z)
 =\frac12\log(2\pi)-\frac12-z(1+\gamma)
 +\sum_{k\ge1}\left[\frac{k}{k+z}-1+\frac zk\right]
 =\frac12\log(2\pi)-\frac12-z(1+\gamma)
 +z\sum_{k\ge1}\left[\frac1k-\frac1{k+z}\right].
\]
The last sum is $\psi(z+1)+\gamma$, which gives the first form in
\eqref{p6:eq:barnes-log-derivative}; the second follows from
$\psi(z+1)=\psi(z)+1/z$.

Now $\psi(z+1)=\psi(z)+1/z>\log z-1/z+1/z=\log z$, using the standard bound
$\psi(z)>\log z-1/z$.  Hence, with $c:=\tfrac12\log(2\pi)-\tfrac12
=0.41893\ldots$,
\begin{equation}\label{p6:eq:barnes-slope-lower}
 \frac{\dd}{\dd z}\log G(z+1)\ >\ z(\log z-1)+c
 \qquad(z>0).
\end{equation}
The right-hand side has derivative $\log z$, so it increases on $[1,\infty)$;
at $z=3$ its value is $3(\log3-1)+c=0.7147\ldots>0$.  This proves strict
monotonicity on $[3,\infty)$.

For \eqref{p6:eq:barnes-slope-bound} it suffices, by
\eqref{p6:eq:barnes-slope-lower}, that $z(\log z-1)+c\ge\tfrac12z\log z$, i.e.\
that $\tfrac12\log z\ge1-c/z$.  The left side increases and the right side
decreases, so it is enough to check $z=7$: $\tfrac12\log7=0.9729\ldots$ and
$1-c/7=0.9401\ldots$.
\end{proof}

\begin{remark}[How sharp the threshold is]\label{p6:rem:threshold}
The bound \eqref{p6:eq:barnes-slope-lower} is lossy.  Solving
\eqref{p6:eq:barnes-log-derivative} numerically, the derivative has its only
positive zero at $z=1.55766\ldots$, so $G(\,\cdot+1)$ is in fact strictly
increasing on $[1.56,\infty)$, and $\tfrac12z\log z$ is dominated from
$z=5.5$ onward.  The thresholds $3$ and $7$ are what the elementary argument
gives; they are stated because a proof is wanted, not because they are best
possible.
\end{remark}

\begin{repairbox}
All three sources assert the existence of the eventual Barnes inverse with the
words ``$(\log G)'(x)\sim x\log x$, hence strictly increasing for all
sufficiently large $x$''.  That is true but ineffective, and none of the three
supplies a threshold; yet every numerical table in all three evaluates the
inverse at arguments as small as $z=5$, inside the range their own
justification does not cover.  \Cref{p6:prop:barnes-monotone} supplies a
threshold, and $3<5$, so the tables are legitimate.
\end{repairbox}

\section{The power--logarithmic core}
\label{p6:sec:core}

\Cref{p0:thm:lambert-core} inverts $aX+b\log X$.  The block needed here is
different and is inverted by a different substitution.

\begin{lemma}[Exact inversion of a power--logarithmic block]
\label{p6:lem:core}
Let $a\ne0$, $p\ne0$, $r\ne0$ and $\beta\in\Cplx$, and put
\begin{equation}\label{p6:eq:core-map}
 \Phi_{0}(x)=a\,x^{p}\,(\log x-\beta)^{r} .
\end{equation}
Fix $R$ and compatible determinations of the logarithm, of the $r$-th root and
of a Lambert branch $\Wlam_k$, and set
\begin{equation}\label{p6:eq:core-vars}
 v=\Wlam_k\!\left(\frac pr\left(\frac{R}{a\,\e^{p\beta}}\right)^{1/r}\right),
 \qquad
 u=\frac rp\,v,
 \qquad
 X=\e^{\beta+u} .
\end{equation}
Then $\Phi_{0}(X)=R$, and
\begin{equation}\label{p6:eq:core-slope}
 \log X-\beta=u,
 \qquad
 \Phi_{0}'(X)=a\,X^{p-1}u^{r-1}\,(pu+r).
\end{equation}
In the case $r=1$ these read
\begin{equation}\label{p6:eq:core-r1}
 v=\Wlam_k\!\left(\frac{pR}{a}\e^{-p\beta}\right),
 \qquad
 X=\e^{\beta+v/p}=\left(\frac{pR}{av}\right)^{1/p},
 \qquad
 \Phi_{0}'(X)=a\,X^{p-1}(v+1).
\end{equation}
\end{lemma}

\begin{proof}
By construction $\log X-\beta=u$, so
$\Phi_0(X)=a\e^{p\beta}\e^{pu}u^{r}$.  Hence $\Phi_0(X)=R$ is equivalent to
$u^{r}\e^{pu}=R/(a\e^{p\beta})$, and taking the $r$-th root, to
$u\,\e^{pu/r}=\bigl(R/(a\e^{p\beta})\bigr)^{1/r}$.  Multiplying by $p/r$ turns
this into $(pu/r)\e^{pu/r}=(p/r)\bigl(R/(a\e^{p\beta})\bigr)^{1/r}$, i.e.\
$v\e^{v}=(p/r)(\cdots)^{1/r}$ with $v=pu/r$, which is
\eqref{p6:eq:core-vars}.

For the slope, differentiate \eqref{p6:eq:core-map}:
\[
 \Phi_0'(x)=a x^{p-1}(\log x-\beta)^{r-1}
 \bigl(p(\log x-\beta)+r\bigr),
\]
and evaluate at $x=X$, where $\log X-\beta=u$.  For $r=1$, $u=v/p$ and
$pu+r=v+1$.
\end{proof}

\begin{remark}[The universal divisor]\label{p6:rem:universal-divisor}
The factor $v+1$ that divides every inverse coefficient below is not an
algebraic accident: by \eqref{p6:eq:core-r1} it is the normalized slope
$\Phi_0'(X)/(aX^{p-1})$ of the exact core.  Its vanishing at $v=-1$ is the
Lambert branch point, where the core map has a critical point and no inverse
of this shape can exist.  Every statement in this chapter that requires
``$v+1$ invertible'' is requiring distance from that turning point, and every
error bound degrades as $v+1\to0$.
\end{remark}

\begin{remark}[What must be absorbed]\label{p6:rem:absorb}
A term of the \emph{same} algebraic power as the core is not small and must be
absorbed before anything is called a perturbation:
$a x^{p}\log x+cx^{p}=a x^{p}(\log x-\beta)$ with $\beta=-c/a$.  Likewise an
additive constant $C$ in the forward phase is moved to the target,
$R:=Y-C$.  Failing to absorb makes the normalized perturbation have valuation
zero and destroys the hypothesis of \cref{p6:thm:triangular}; failing to
absorb the constant merely shifts the core, but by an amount that is not
small compared with the corrections one is trying to compute.  A phase
$x^{p}\bigl(a_r(\log x)^{r}+a_{r-1}(\log x)^{r-1}+\cdots\bigr)$ with $r\ge2$
needs a genuine outer inversion in the logarithmic scale first; only its
leading logarithmic power is covered by \cref{p6:lem:core}.
\end{remark}

\section{Normalized displacement and triangular inversion}
\label{p6:sec:triangular}

\subsection{The universal kernel}

Take $r=1$ from here on, and consider a phase
\begin{equation}\label{p6:eq:perturbed-phase}
 \Phi(x)=a\,x^{p}(\log x-\beta)+C+\Psi(x),
\end{equation}
to be solved for $\Phi(x)=Y$.  Put $R=Y-C$, construct $v$ and $X$ by
\eqref{p6:eq:core-r1}, write $q=X^{-1}$, and substitute
\begin{equation}\label{p6:eq:relative-displacement}
 x=X(1+\epsilon).
\end{equation}
Dividing $\Phi(x)-Y$ by $aX^{p}$ turns the equation into
\begin{equation}\label{p6:eq:normalized}
 \mathcal{K}_{p,v}(\epsilon)+\mathcal{R}(q,v,\epsilon)=0,
\end{equation}
where
\begin{equation}\label{p6:eq:kernel}
 \mathcal{K}_{p,v}(\epsilon)
 =(1+\epsilon)^{p}\left(\frac vp+\log(1+\epsilon)\right)-\frac vp,
 \qquad
 \mathcal{R}(q,v,\epsilon)=\frac{\Psi\bigl(X(1+\epsilon)\bigr)}{aX^{p}} .
\end{equation}

\begin{lemma}[Kernel coefficients]\label{p6:lem:kernel-coefficients}
Write $\mathcal{K}_{p,v}(\epsilon)=\sum_{n\ge1}\kappa_n(p,v)\epsilon^{n}$.  Then
\begin{equation}\label{p6:eq:kappa}
 \kappa_n(p,v)=\frac vp\binom pn+\frac{\partial}{\partial p}\binom pn,
\end{equation}
a polynomial in $p$ and $v$; in particular
\begin{equation}\label{p6:eq:kappa1}
 \kappa_1(p,v)=v+1 .
\end{equation}
Where $p,p-1,\dots,p-n+1$ are all nonzero one may write
$\partial_p\binom pn=\binom pn\sum_{j=0}^{n-1}(p-j)^{-1}$, but
\eqref{p6:eq:kappa} is the form valid at every $p$.  For $p=1$,
\begin{equation}\label{p6:eq:kernel-p1}
 \mathcal{K}_{1,v}(\epsilon)=v\epsilon+(1+\epsilon)\log(1+\epsilon)
 =(v+1)\epsilon+\sum_{m\ge2}\frac{(-1)^{m}}{m(m-1)}\epsilon^{m},
\end{equation}
and for $p=2$,
\begin{equation}\label{p6:eq:kernel-p2}
 \mathcal{K}_{2,v}(\epsilon)
 =(v+1)\epsilon+\frac{v+3}2\epsilon^{2}
 +\frac13\epsilon^{3}-\frac1{12}\epsilon^{4}
 +\frac1{30}\epsilon^{5}-\cdots .
\end{equation}
\end{lemma}

\begin{proof}
Since $(1+\epsilon)^{p}\log(1+\epsilon)=\partial_p(1+\epsilon)^{p}$, we have
$\Coef{\epsilon^{n}}\bigl((1+\epsilon)^{p}\log(1+\epsilon)\bigr)
=\partial_p\binom pn$, while
$\Coef{\epsilon^{n}}\bigl((v/p)(1+\epsilon)^{p}\bigr)=(v/p)\binom pn$; the
constant $-v/p$ affects only $n=0$.  This is \eqref{p6:eq:kappa}.  For $n=1$,
$\binom p1=p$ and $\partial_p\binom p1=1$, so $\kappa_1=v+1$.  Both
displayed specializations follow by direct expansion; note that the
coefficient of $\epsilon^2$ in \eqref{p6:eq:kernel-p2} is
$(v/2)\binom22\cdot2+\partial_p\binom p2|_{p=2}
=\tfrac v2+\tfrac32=\tfrac{v+3}2$.
\end{proof}

\subsection{Triangular inversion over a well-ordered support}

\begin{theorem}[Triangular inversion on a Hahn grid]
\label{p6:thm:triangular}
Let $\mathscr C$ be a field of characteristic zero containing $p$, $v$ and
$(v+1)^{-1}$, let $S\subset\Real_{\ge0}$ be a \emph{well-ordered} additive
submonoid, and let $\mathscr C[[q^{S}]]$ be the corresponding Hahn ring.
Assume
\begin{equation}\label{p6:eq:R-shape}
 \mathcal{R}(q,v,E)=\sum_{\rho\in S,\ \rho>0}\ \sum_{m\ge0}
 r_{\rho,m}(v)\,q^{\rho}E^{m},
\end{equation}
with every monomial of positive $q$-exponent.  Then \eqref{p6:eq:normalized}
has a unique solution $\epsilon$ of positive valuation in
$\mathscr C[[q^{S}]]$, say
$\epsilon=\sum_{\gamma\in S,\,\gamma>0}e_\gamma(v)q^{\gamma}$, and writing
$\epsilon_{<\gamma}$ for its truncation below $\gamma$,
\begin{equation}\label{p6:eq:triangular}
 e_{\gamma}(v)=-\frac1{v+1}\,\Coef{q^{\gamma}}
 \left(\sum_{n\ge2}\kappa_n(p,v)\,\epsilon_{<\gamma}^{\,n}
 +\mathcal{R}(q,v,\epsilon_{<\gamma})\right).
\end{equation}
\end{theorem}

\begin{proof}
Well-ordering is exactly what makes the coefficient extraction in
\eqref{p6:eq:triangular} a finite sum.  By Neumann's lemma, if $A,B\subset
\Real$ are well ordered then $A+B$ is well ordered and each of its elements
has only finitely many representations $a+b$ with $a\in A$, $b\in B$; by
induction the same holds for any finite number of summands.  Since
$\operatorname{val}(\epsilon)>0$, we get
$\operatorname{val}(\epsilon^{n})\ge n\operatorname{val}(\epsilon)\to\infty$,
so at a prescribed exponent $\gamma$ only finitely many $n$ contribute, and
for each such $n$ only finitely many decompositions.  The same applies to
$\mathcal{R}$ because every one of its monomials has $\rho>0$.

Now fix $\gamma$ and suppose $e_\delta$ is known for all $0<\delta<\gamma$.
In the left-hand side of \eqref{p6:eq:normalized}, the coefficient of
$q^{\gamma}$ receives $\kappa_1e_\gamma=(v+1)e_\gamma$ from the linear term,
and from nothing else: a monomial of $\epsilon^{n}$ with $n\ge2$ that involves
$e_\gamma$ has $q$-exponent at least $\gamma+\operatorname{val}(\epsilon)>
\gamma$, and a monomial $q^{\rho}E^{m}$ of $\mathcal{R}$ involving $e_\gamma$ has
$q$-exponent at least $\rho+\gamma>\gamma$.  Solving the resulting linear
equation gives \eqref{p6:eq:triangular}; transfinite induction along the
well-ordered support constructs $\epsilon$, and the same computation run on
the difference of two solutions shows the first exponent at which they could
differ satisfies $(v+1)(e_\gamma-\widetilde e_\gamma)=0$, so there is none.
\end{proof}

\begin{repairbox}
S1 states the same theorem under the hypothesis that the support semigroup is
\emph{locally finite}, meaning that every bounded interval meets it in a
finite set.  That is strictly stronger than well-ordering --- $\{1-1/n\}\cup
\{1\}$ is well ordered and not locally finite --- and it is more than the
proof uses.  What the argument needs is precisely the finite-decomposition
property, which Neumann's lemma supplies for every well-ordered set.  The
hypothesis is weakened accordingly here.  For the two subjects of this chapter
the distinction is invisible, since $S=2\Nat$ and $S=\Nat$; it matters for the
fractional and mixed grids of \cref{p6:sec:general}.
\end{repairbox}

\begin{corollary}[Support semigroup]\label{p6:cor:support}
If the $q$-support of $\mathcal{R}(q,v,0)$ lies in $A\subset\Real_{>0}$ and no
monomial of $\mathcal{R}$ has a $q$-exponent outside the monoid generated by $A$,
then $\operatorname{supp}\epsilon$ lies in the additive semigroup generated by
$A$.
\end{corollary}

\begin{proof}
\Cref{p6:eq:triangular} produces $e_\gamma$ from sums of products of earlier
coefficients and of base exponents from $A$; the set of exponents so reachable
is closed under addition and contains $A$.  Induction along the support.
\end{proof}

\begin{example}[The two supports]\label{p6:ex:supports}
For $\Gamma$ we have $p=a=\beta=1$, and the perturbation $a_kt^{-(2k-1)}$
normalized by $aX^{p}=T$ becomes $a_kq^{2k}$ up to a factor
$(1+\epsilon)^{-(2k-1)}$; so $A=\{2,4,6,\dots\}$, the generated semigroup is
$2\Nat$, and the displacement $t-T=T\epsilon$ contains only the odd powers
$T^{-1},T^{-3},\dots$.  For $G$ we have $p=2$, $a=\tfrac12$, $\beta=\tfrac32$,
and the linear forward term $\alpha z$ normalized by $\tfrac12U^{2}$ becomes
$2\alpha q$; so $1\in A$, the semigroup is all of $\Nat$, and every integer
power occurs.
\end{example}

\subsection{The other engine: formal Newton iteration}

The triangular recurrence \eqref{p6:eq:triangular} produces one coefficient
per step.  When many are wanted, a second engine is better, and it needs no
new hypothesis --- only that the slope be a unit, which is what
\cref{p6:thm:triangular} already assumes.

\begin{theorem}[Valuation doubling]\label{p6:thm:newton}
Let $\mathcal{E}(\epsilon):=\mathcal{K}_{p,v}(\epsilon)+\mathcal{R}(q,v,\epsilon)$
be the left-hand side of \eqref{p6:eq:normalized}, and assume the hypotheses of
\cref{p6:thm:triangular}.  Then $\mathcal{E}'$ is a unit of valuation $0$, the
formal Newton step
\begin{equation}\label{p6:eq:newton-step}
 \epsilon^{[m+1]}=\epsilon^{[m]}
 -\frac{\mathcal{E}\bigl(\epsilon^{[m]}\bigr)}
        {\mathcal{E}'\bigl(\epsilon^{[m]}\bigr)}
\end{equation}
is defined, and
\begin{equation}\label{p6:eq:newton-doubling}
 \operatorname{val}\mathcal{E}\bigl(\epsilon^{[m]}\bigr)\ge\varrho
 \quad\Longrightarrow\quad
 \operatorname{val}\mathcal{E}\bigl(\epsilon^{[m+1]}\bigr)\ge2\varrho .
\end{equation}
Consequently the number of correct blocks doubles at each step, and $N$ of them
cost $\bigO(\log N)$ steps of truncated-series arithmetic instead of $N$
triangular steps.
\end{theorem}

\begin{proof}
By \cref{p6:lem:kernel-coefficients}, $\mathcal{K}_{p,v}'(0)=\kappa_1=v+1$,
and every monomial of $\mathcal{R}$ has positive $q$-valuation; so
$\mathcal{E}'(\epsilon)=(v+1)+(\text{positive valuation})$, which is a unit of
valuation $0$ because $v+1$ is invertible in $\mathscr C$ and the Hahn ring is
complete.  Hence \eqref{p6:eq:newton-step} is defined and, writing
$\Delta=-\mathcal{E}(\epsilon)/\mathcal{E}'(\epsilon)$,
\[
 \operatorname{val}\Delta=\operatorname{val}\mathcal{E}(\epsilon)\ge\varrho .
\]
Taylor expansion of $\mathcal{E}$ about $\epsilon$ --- legitimate because
$\mathcal{E}$ is a formal series in $\epsilon$ with coefficients in the
complete ring --- gives
\[
 \mathcal{E}(\epsilon+\Delta)
 =\mathcal{E}(\epsilon)+\mathcal{E}'(\epsilon)\Delta
 +\sum_{j\ge2}\frac{\mathcal{E}^{(j)}(\epsilon)}{j!}\Delta^{j} .
\]
The first two terms cancel by the definition of $\Delta$, and every remaining
term has valuation at least $2\operatorname{val}\Delta\ge2\varrho$, since the
$\mathcal{E}^{(j)}$ have valuation at least $0$.  This is
\eqref{p6:eq:newton-doubling}, and iterating from
$\operatorname{val}\mathcal{E}(0)>0$ gives the count.
\end{proof}

\begin{remark}[Which engine to use]\label{p6:rem:which-engine}
The triangular recurrence exposes exact coefficient formulas and is what the
proofs in this chapter and the next use; Newton is preferable in a computer
algebra system at high order, because one inversion of a unit replaces many
scalar recurrence steps.  They also check each other: run to the same
valuation and the truncated residuals must agree, which catches an error in
either.  The same doubling holds verbatim for the accelerated equation
\eqref{p7:eq:acc-implicit} and the refined equation
\eqref{p7:eq:refined-implicit} of \cref{p7:sec:top}, whose linearizations are
$S$ and $4L_{*}$.
\end{remark}

\subsection{How deep the slope divisor goes}

Every coefficient is produced by one division by $v+1$, but earlier
coefficients enter nonlinearly and carry their own divisions.  The
accumulation is bounded, and --- this is new --- the bound is attained.

\begin{proposition}[Slope-denominator bound]\label{p6:prop:denominator}
Suppose $S=\Nat$, that the coefficients $r_{\rho,m}$ of
\eqref{p6:eq:R-shape} lie in $A[v]$ for a commutative ring $A\supseteq\Rat$,
and that none of them involves $(v+1)^{-1}$.  If
$\epsilon=\sum_{n\ge1}e_nq^{n}$ is the solution of \cref{p6:thm:triangular},
then
\begin{equation}\label{p6:eq:denominator}
 (v+1)^{2n-1}e_n\in A[v]\qquad(n\ge1).
\end{equation}
\end{proposition}

\begin{proof}
Induction on $n$.  For $n=1$, \eqref{p6:eq:triangular} gives
$e_1=-r_{1,0}/(v+1)$, and $(v+1)e_1\in A[v]$.

Let $n\ge2$ and assume the claim below $n$.  In
\eqref{p6:eq:triangular} the bracket is a sum of two kinds of terms.  A term
of $\kappa_m(p,v)\epsilon^{m}$ with $m\ge2$ contributing to $q^{n}$ is
$\kappa_m$ times a product $e_{j_1}\cdots e_{j_m}$ with
$j_1+\cdots+j_m=n$ and every $j_i\ge1$; since $\kappa_m\in\Rat[p,v]$ carries
no divisor, the inductive hypothesis bounds its denominator exponent by
\[
 \sum_{i=1}^{m}(2j_i-1)=2n-m\ \le\ 2n-2 .
\]
A term $r_{\rho,m}q^{\rho}E^{m}$ contributing to $q^{n}$ has
$j_1+\cdots+j_m=n-\rho$ with $\rho\ge1$, so its denominator exponent is at
most $2(n-\rho)-m\le2n-2$ as well (and at most $0$ when $m=0$).  Hence the
bracket has denominator exponent at most $2n-2$, and the final division by
$v+1$ in \eqref{p6:eq:triangular} gives at most $2n-1$.
\end{proof}

The next two statements are the chapter's first new contribution.  They say
that the bound of \cref{p6:prop:denominator} is exactly attained, and identify
the coefficient of the deepest pole in closed form, in both subjects at once.

\begin{lemma}[The quadratic core]\label{p6:lem:quadratic-core}
Let $A$ and $c\ne0$ be elements of a commutative $\Rat$-algebra with $A$
invertible, let $\sigma$ be another element, and let
$\epsilon=\sum_{n\ge1}\epsilon_nq^{n}$ be the unique solution of positive
valuation of the \emph{quadratic model}
\begin{equation}\label{p6:eq:quadratic-model}
 A\epsilon+c\,\epsilon^{2}+\sigma q=0 .
\end{equation}
Then
\begin{equation}\label{p6:eq:quadratic-solution}
 \epsilon=\frac{A}{2c}\left(\sqrt{1-\frac{4c\sigma q}{A^{2}}}-1\right),
 \qquad
 \epsilon_n=\frac{(-4c\sigma)^{n}}{2c}\binom{1/2}{n}A^{-(2n-1)} .
\end{equation}
In particular the denominator exponent of $\epsilon_n$ in $A$ is exactly
$2n-1$ whenever $c\sigma\ne0$.
\end{lemma}

\begin{proof}
Solving the quadratic and choosing the root that vanishes at $q=0$ gives
\[
 \epsilon=\frac{-A+\sqrt{A^{2}-4c\sigma q}}{2c},
\]
which is the first formula.  Expanding $\sqrt{1+w}=\sum_{n\ge0}\binom{1/2}{n}w^{n}$ at
$w=-4c\sigma q/A^{2}$ and multiplying by $A/(2c)$ gives the second.
\end{proof}

\begin{proposition}[The coefficients of the quadratic core are Catalan numbers]
\label{p6:prop:quadratic-core-catalan}
For every $n\ge1$,
\begin{equation}\label{p6:eq:half-binom-catalan}
 \binom{1/2}{n}=(-1)^{n-1}\,\mathrm{Cat}_{n-1}\,2^{-(2n-1)},
\end{equation}
where $\mathrm{Cat}_{k}=\frac1{k+1}\binom{2k}{k}$ is the $k$th Catalan number, and
consequently the convolution that \eqref{p6:eq:quadratic-model} imposes on its
coefficients,
\begin{equation}\label{p6:eq:half-binom-convolution}
 \sum_{j=1}^{n-1}\binom{1/2}{j}\binom{1/2}{n-j}
 =-2\binom{1/2}{n}\qquad(n\ge2),
\end{equation}
is precisely the Catalan recursion $\mathrm{Cat}_{m+1}=\sum_{i+j=m}\mathrm{Cat}_{i}\mathrm{Cat}_{j}$.
\end{proposition}

\begin{proof}
Both sides of \eqref{p6:eq:half-binom-catalan} satisfy the same first-order
recursion and agree at $n=1$: writing $g_n$ for the right-hand side,
$\binom{1/2}{n+1}=\binom{1/2}{n}(1/2-n)/(n+1)$, while
$(n+2)\mathrm{Cat}_{n+1}=2(2n+1)\mathrm{Cat}_{n}$ --- the latter obtained by cancelling $n+1$
from the central binomial recursion $(n+1)\binom{2n+2}{n+1}=2(2n+1)\binom{2n}{n}$ ---
and these two ratios agree after the sign and the power of $2$ are accounted
for.  Given \eqref{p6:eq:half-binom-catalan}, substituting into the left side of
\eqref{p6:eq:half-binom-convolution}, with $j=i+1$ and $n=m+2$, gives
$(-1)^{m}2^{-(2m+2)}\sum_{i+j=m}\mathrm{Cat}_{i}\mathrm{Cat}_{j}$, and the right side is
$-2(-1)^{m+1}\mathrm{Cat}_{m+1}2^{-(2m+3)}=(-1)^{m}\mathrm{Cat}_{m+1}2^{-(2m+2)}$; the two agree
exactly when $\mathrm{Cat}_{m+1}=\sum_{i+j=m}\mathrm{Cat}_{i}\mathrm{Cat}_{j}$, which is the Catalan
recursion.
\end{proof}

\begin{remark}[What the reduction buys, and the formalization]
\label{p6:rem:quadratic-core-lean}
The proof of \cref{p6:lem:quadratic-core} given above is analytic: it passes
through the binomial series for $\sqrt{1+w}$, hence through a question of
convergence that the formal statement does not need.  \Cref{p6:prop:quadratic-core-catalan}
removes that detour.  The only property of the coefficient family the
quadratic model actually uses is the convolution
\eqref{p6:eq:half-binom-convolution}, and under
\eqref{p6:eq:half-binom-catalan} that convolution is a finite combinatorial
recursion.  So \cref{p6:lem:quadratic-core}, and with it the deepest-pole
family of \cref{p6:thm:deepest-pole} in both the $\Gamma$ and the Barnes
$G$ cases, rests on no analysis at all.

This is machine-checked, in \path{QuadraticCoreCatalan.lean}, over an arbitrary
field of characteristic zero, and as a Mathlib-only leaf with no dependence on
the rest of the corpus.  \path{Fabius.catalan_two_step} is the ratio recursion;
\path{Fabius.quadHalf} is the coefficient family, indexed from zero, defined by
the closed form \eqref{p6:eq:half-binom-catalan};
\path{Fabius.quadHalf_antidiagonal} is the convolution
\eqref{p6:eq:half-binom-convolution};
\path{Fabius.quadHalf_rat} is \eqref{p6:eq:half-binom-catalan} itself,
identifying that family over $\Rat$ with $\binom{1/2}{n+1}$ as defined by the
falling product; and \path{Fabius.quadCoef_zero} together with
\path{Fabius.quadCoef_rec} are the two halves of \cref{p6:lem:quadratic-core}
--- that $\epsilon_n=((-4c\sigma)^{n}/2c)\binom{1/2}{n}A^{-(2n-1)}$ solves
\eqref{p6:eq:quadratic-model} order by order.  The index is shifted so that the
convolution runs over $\texttt{Finset.antidiagonal}$ and the sum
\eqref{p6:eq:half-binom-convolution}, which is empty at $n=1$, is carried by the
separate statement \path{Fabius.quadCoef_zero} rather than by a boundary case.

What is \emph{not} formalized: that this family is the deepest-pole family of
the two inversions, which is \cref{p6:thm:deepest-pole} and needs the
transseries apparatus.  The module supplies only the algebraic core.
\end{remark}

\begin{theorem}[The deepest pole is attained, with an explicit coefficient]
\label{p6:thm:deepest-pole}
\leavevmode
\begin{enumerate}
\item \emph{Gamma.}  For $n\ge1$ let $d_{2n-1}(\Lambda)$ be the coefficients
of \cref{p6:thm:gamma}.  Then
\begin{equation}\label{p6:eq:gamma-deepest}
 \Coef{\Lambda^{-(2n-1)}}\,d_{2n-1}(\Lambda)
 =\binom{1/2}{n}12^{-n}\ \ne0,
\end{equation}
so the bound of \cref{p6:prop:denominator} is attained at every $n$, and the
sum of the deepest-pole terms of the displacement is the convergent series
\begin{equation}\label{p6:eq:gamma-deepest-resum}
 \sum_{n\ge1}\binom{1/2}{n}12^{-n}(\Lambda T)^{1-2n}
 =\sqrt{\Lambda^{2}T^{2}+\tfrac1{12}}-\Lambda T
 \qquad\bigl(\Lambda^{2}T^{2}>\tfrac1{12}\bigr).
\end{equation}
\item \emph{Barnes.}  For $n\ge1$ let $e_n(\lambda)$ be the coefficients of
\cref{p6:thm:barnes}.  Then
\begin{equation}\label{p6:eq:barnes-deepest}
 \Coef{\lambda^{-(2n-1)}}\,e_n(\lambda)
 =\binom{1/2}{n}(-2\alpha)^{n}\ \ne0,
\end{equation}
and the deepest-pole part of the displacement $z-U$ sums to
\begin{equation}\label{p6:eq:barnes-deepest-resum}
 \sum_{n\ge1}\binom{1/2}{n}(-2\alpha)^{n}\lambda^{1-2n}U^{1-n}
 =\sqrt{\lambda^{2}U^{2}-2\alpha U}-\lambda U
 \qquad(\lambda^{2}U>2\alpha).
\end{equation}
\end{enumerate}
In both cases the deepest-pole family is the coefficient family of the
quadratic model \eqref{p6:eq:quadratic-model}, with
\begin{equation}\label{p6:eq:deepest-dictionary}
 (A,c,\sigma)=(\Lambda,\ \tfrac12,\ a_1)\ \text{ for }\Gamma,
 \qquad
 (A,c,\sigma)=(2\lambda,\ 1,\ 2\alpha)\ \text{ for }G .
\end{equation}
\end{theorem}

\begin{proof}
Both parts have the same shape.  Write $\pi_n$ for the coefficient of the
deepest possible pole, $\pi_n:=\Coef{A^{-(2n-1)}}e_n$, which is well defined
by \cref{p6:prop:denominator}; the claim is that $(\pi_n)$ obeys the
recursion of the quadratic model.  Everything turns on which terms of the
normalized equation can reach pole order $2n-2$ after the coefficient
$e_n$ has been isolated --- for the recurrence reads $Ae_n=-(\text{bracket})$,
and $A\cdot\pi_nA^{-(2n-1)}=\pi_nA^{-(2n-2)}$.

\emph{Gamma.}  Here $A=\Lambda=v+1$ and, by \eqref{p6:eq:kernel-p1}, the
normalized equation on the grid $q^{2}$ is
\begin{equation}\label{p6:eq:gamma-normalized-again}
 \Lambda\epsilon+\tfrac12\epsilon^{2}-\tfrac16\epsilon^{3}+\cdots
 +\sum_{k\ge1}a_kq^{2k}(1+\epsilon)^{-(2k-1)}=0 .
\end{equation}
Index by $n$, so that $\epsilon=\sum_ne_nq^{2n}$ and $e_n$ has pole order at
most $2n-1$.  The term $\tfrac12\epsilon^{2}$ contributes at $q^{2n}$ the sum
$\tfrac12\sum_{j=1}^{n-1}e_je_{n-j}$, of pole order at most
$(2j-1)+(2(n-j)-1)=2n-2$ with top coefficient
$\tfrac12\sum_{j=1}^{n-1}\pi_j\pi_{n-j}$.  A term $\epsilon^{m}$ with $m\ge3$
has pole order at most $2n-m\le2n-3$.  In the Bernoulli sum, the summand
$k=n$ contributes the constant $a_n$, of pole order $0$, which is $2n-2$ only
when $n=1$; and a summand with $1\le k\le n-1$ contributes $a_k$ times a
product of $m\ge1$ coefficients with indices summing to $n-k$, of pole order
at most $2(n-k)-m\le2n-3$.  Hence for $n\ge2$
\[
 \pi_n=-\tfrac12\sum_{j=1}^{n-1}\pi_j\pi_{n-j},
 \qquad\text{while}\qquad \pi_1=-a_1=\tfrac1{24}.
\]
Let $g(w)=\sum_{n\ge1}\pi_nw^{n}$.  The recursion says
$g=-a_1w-\tfrac12g^{2}$, i.e.\ $g^{2}+2g+2a_1w=0$, whence
$g=-1+\sqrt{1-2a_1w}=-1+\sqrt{1+w/12}$ and
$\pi_n=\binom{1/2}{n}12^{-n}$, which is \eqref{p6:eq:gamma-deepest}; it is
nonzero because $\binom{1/2}{n}\ne0$ for all $n$.  This is exactly
\cref{p6:lem:quadratic-core} with $(A,c,\sigma)=(\Lambda,\tfrac12,a_1)$:
indeed $(-4\cdot\tfrac12\cdot a_1)^{n}/(2\cdot\tfrac12)=(1/12)^{n}$.

Since $t-T=T\epsilon$, the deepest-pole terms of the displacement are
\[
 \pi_n\Lambda^{-(2n-1)}T^{1-2n}
 =\binom{1/2}{n}12^{-n}(\Lambda T)^{1-2n},
\]
and summing the binomial series at $w=1/(12\Lambda^{2}T^{2})$ gives
\[
 \Lambda T\left(\sqrt{1+\frac1{12\Lambda^{2}T^{2}}}-1\right)
 =\sqrt{\Lambda^{2}T^{2}+\tfrac1{12}}-\Lambda T,
\]
convergent for $12\Lambda^{2}T^{2}>1$.  This is
\eqref{p6:eq:gamma-deepest-resum}.

\emph{Barnes.}  Here $A=2\lambda$ and, by \eqref{p6:eq:kernel-p2} with
$v=2\lambda-1$ so that $(v+3)/2=\lambda+1$, the normalized equation on the
grid $q$ is
\begin{equation}\label{p6:eq:barnes-normalized-again}
 2\lambda\epsilon+(\lambda+1)\epsilon^{2}+\tfrac13\epsilon^{3}-\cdots
 +2\alpha q(1+\epsilon)
 +2b\,q^{2}\bigl(\lambda+1+\log(1+\epsilon)\bigr)
 +2\sum_{k\ge1}c_kq^{2k+2}(1+\epsilon)^{-2k}=0,
\end{equation}
with $b=-\tfrac1{12}$.  Pole orders are now taken in $\lambda$, and
$e_n$ has order at most $2n-1$; we must collect everything of order $2n-2$.

The quadratic term $(\lambda+1)\epsilon^{2}$ splits.  Its part $1\cdot
\epsilon^{2}$ contributes $\sum_{j=1}^{n-1}\pi_j\pi_{n-j}$ at order $2n-2$.
Its part $\lambda\epsilon^{2}$ multiplies a quantity of order at most $2n-2$
by $\lambda$, giving order at most $2n-3$; it therefore does \emph{not}
reach $2n-2$.  Terms $\epsilon^{m}$ with $m\ge3$ reach at most $2n-3$.

In the perturbation, $2\alpha q\cdot1$ contributes only at $n=1$, where its
order $0$ equals $2n-2=0$; $2\alpha q\epsilon$ contributes $2\alpha e_{n-1}$,
of order at most $2n-3$.  The term $2b\lambda q^{2}$ has no pole at all, and
$2bq^{2}$ and $2bq^{2}\log(1+\epsilon)$ have order at most $2(n-2)-1$, while
multiplying the logarithm by $\lambda$ lowers the order further; none reaches
$2n-2$.  The same holds for every $c_k$ term.  Hence for $n\ge2$
\[
 2\pi_n=-\sum_{j=1}^{n-1}\pi_j\pi_{n-j},
 \qquad\text{while}\qquad 2\pi_1=-2\alpha,\ \ \pi_1=-\alpha .
\]
With $g(w)=\sum\pi_nw^{n}$ this reads $2g=-2\alpha w-g^{2}$, so
$g=-1+\sqrt{1-2\alpha w}$ and $\pi_n=\binom{1/2}{n}(-2\alpha)^{n}$, which is
\eqref{p6:eq:barnes-deepest} and agrees with \cref{p6:lem:quadratic-core} at
$(A,c,\sigma)=(2\lambda,1,2\alpha)$.  Since $z-U=U\epsilon$, the deepest-pole
terms are $\pi_n\lambda^{-(2n-1)}U^{1-n}$, and
\[
 \sum_{n\ge1}\binom{1/2}{n}(-2\alpha)^{n}\lambda^{1-2n}U^{1-n}
 =\lambda U\sum_{n\ge1}\binom{1/2}{n}
   \left(\frac{-2\alpha}{\lambda^{2}U}\right)^{n}
 =\lambda U\left(\sqrt{1-\frac{2\alpha}{\lambda^{2}U}}-1\right),
\]
which is \eqref{p6:eq:barnes-deepest-resum}.
\end{proof}

\begin{remark}[Two square roots, one mechanism]\label{p6:rem:two-square-roots}
S1 and S3 both isolate a square root in the Barnes problem, namely the
coefficientwise limit $\sqrt{U^{2}+1/6}$ of \cref{p6:prop:resonance} below,
and both present it as a curiosity of Barnes $G$ caused by the resonant
$-\tfrac1{12}\log z$ term.  \Cref{p6:thm:deepest-pole} shows it is not special
to Barnes and not special to resonance.  Every problem in this class has a
distinguished exactly solvable truncation --- keep the slope, the quadratic
kernel coefficient, and the lowest forcing monomial, discard the rest --- and
that truncation is a quadratic, so it always resums to a square root.  Which
\emph{extraction} of the coefficient family it computes depends on the
subject: for Barnes, discarding by $\lambda$-degree gives the resonant limit
\eqref{p6:eq:resonant-limit}; discarding by pole depth gives
\eqref{p6:eq:barnes-deepest-resum}; for Gamma there is no resonance and only
the second extraction exists.  The three square roots
\[
 \sqrt{U^{2}+\tfrac16},\qquad
 \sqrt{\lambda^{2}U^{2}-2\alpha U},\qquad
 \sqrt{\Lambda^{2}T^{2}+\tfrac1{12}}
\]
are the same lemma applied three times.
\end{remark}

\begin{remark}[What the deepest pole is good for]\label{p6:rem:deepest-use}
Two things.  First, it is a sharp regression test: a symbolic implementation
that produces $d_{2n-1}$ or $e_n$ can check the single rational number
$\Coef{A^{-(2n-1)}}$ against \eqref{p6:eq:gamma-deepest} or
\eqref{p6:eq:barnes-deepest} at negligible cost, and that number is exactly
the one most sensitive to an error committed early in the recursion, since it
is the accumulation of all previous divisions.  Second, it locates the
breakdown: the expansions of this chapter are asymptotic in $q$ but their
coefficients degrade like $(v+1)^{-(2n-1)}$, so the useful range in $n$ at a
given distance from the branch point $v=-1$ is governed by
$\Lambda T\gg1$ resp.\ $\lambda\sqrt U\gg1$ --- and by
\cref{p6:thm:deepest-pole} those are precisely the conditions under which the
displayed square roots converge.
\end{remark}

\section{The increasing inverse of Gamma}
\label{p6:sec:gamma}

\subsection{Core variables}

Let $y\to+\infty$ and put
\begin{equation}\label{p6:eq:gamma-core-vars}
 R=\log y-\tfrac12\log(2\pi),
 \qquad
 w=\Wlam_0(R/\e),
 \qquad
 T=\frac Rw=\e^{w+1},
 \qquad
 \Lambda=\log T=w+1 .
\end{equation}
These are \cref{p6:lem:core} at $p=a=\beta=r=1$: the core equation is
\begin{equation}\label{p6:eq:gamma-core-eq}
 T(\log T-1)=R,
\end{equation}
its slope is $\Lambda$, and the three identities in
\eqref{p6:eq:gamma-core-vars} are equivalent to $w\e^{w}=R/\e$.  Write
$q=T^{-1}$ and $t=T(1+\epsilon)$.  Subtracting \eqref{p6:eq:gamma-core-eq}
from \cref{p6:lem:gamma-normal-form} and dividing by $T$ gives the exact
normalized equation
\begin{equation}\label{p6:eq:gamma-normalized}
 \Lambda\epsilon+\bigl[(1+\epsilon)\log(1+\epsilon)-\epsilon\bigr]
 +\sum_{k\ge1}a_kq^{2k}(1+\epsilon)^{-(2k-1)}=0 .
\end{equation}

\subsection{The all-orders theorem}

\begin{theorem}[All-orders expansion of $\Gamma_+^{-1}$]\label{p6:thm:gamma}
Let $\Gamma_+^{-1}$ be the branch of \cref{p6:rem:which-inverse} tending to
$+\infty$, and let $T,\Lambda$ be as in \eqref{p6:eq:gamma-core-vars}.  There
are uniquely determined $d_{2n-1}\in\Lambda^{-1}\Rat[\Lambda^{-1}]$ such that,
for every fixed $N\ge0$,
\begin{equation}\label{p6:eq:gamma-main}
 \Gamma_+^{-1}(y)
 =\frac12+T+\sum_{n=1}^{N}\frac{d_{2n-1}(\Lambda)}{T^{2n-1}}
 +\bigO_N\!\left(\frac1{\Lambda\,T^{2N+1}}\right),
 \qquad y\to+\infty .
\end{equation}
Writing $\epsilon=\sum_{n\ge1}d_{2n-1}(\Lambda)q^{2n}$ and
$E_{n-1}=\sum_{j<n}d_{2j-1}q^{2j}$, they are computed by
\begin{equation}\label{p6:eq:gamma-recurrence}
 d_{2n-1}(\Lambda)=-\frac1\Lambda\Coef{q^{2n}}
 \left[(1+E_{n-1})\log(1+E_{n-1})-E_{n-1}
 +\sum_{k=1}^{n}a_kq^{2k}(1+E_{n-1})^{-(2k-1)}\right],
\end{equation}
and they satisfy
\begin{equation}\label{p6:eq:gamma-degree}
 \deg_{\Lambda^{-1}}d_{2n-1}\le2n-1,
 \qquad
 d_{2n-1}(\Lambda)=-\frac{a_n}{\Lambda}+\bigO(\Lambda^{-2}).
\end{equation}
Only the odd powers $T^{-1},T^{-3},T^{-5},\dots$ occur.
\end{theorem}

\begin{proof}
\emph{Formal part.}  Equation \eqref{p6:eq:gamma-normalized} is
\eqref{p6:eq:normalized} with $p=a=\beta=1$, support $S=2\Nat$ by
\cref{p6:ex:supports}, and linear coefficient $\kappa_1=v+1=\Lambda$; so
\cref{p6:thm:triangular} gives a unique formal solution and
\eqref{p6:eq:gamma-recurrence}.  The parity claim is
\cref{p6:cor:support} together with $t-T=T\epsilon$: an exponent $q^{2n}$ in
$\epsilon$ is an absolute power $T^{1-2n}$.

For \eqref{p6:eq:gamma-degree}, induct.  The bracket in
\eqref{p6:eq:gamma-recurrence} has rational Taylor coefficients, so if every
$d_{2j-1}$ with $j<n$ lies in $\Lambda^{-1}\Rat[\Lambda^{-1}]$ with degree at
most $2j-1$, the bound of \cref{p6:prop:denominator} applies verbatim and the
bracket has degree at most $2n-2$; dividing by $\Lambda$ gives degree at most
$2n-1$ and no constant term.  At order $q^{2n}$ the only contribution not
involving an earlier coefficient is $a_n$, whence
$d_{2n-1}=-a_n/\Lambda+\bigO(\Lambda^{-2})$.

\emph{Analytic part.}  Fix $N$ and set $M=N+1$.  We must know that
\eqref{p6:eq:gamma-normal} truncated after $a_Mt^{-(2M-1)}$ has remainder
$\bigO(t^{-(2M+1)})$ on the positive real axis, with a differentiable
remainder.  This follows from the Legendre duplication formula
\[
 \Gamma(t)\Gamma\!\left(t+\tfrac12\right)=2^{1-2t}\sqrt\pi\,\Gamma(2t),
 \quad\text{i.e.}\quad
 \log\Gamma\!\left(t+\tfrac12\right)
 =\log\Gamma(2t)-\log\Gamma(t)+(1-2t)\log2+\tfrac12\log\pi,
\]
applied to two ordinary Stirling expansions on the positive real axis, each of
which has the classical remainder bound with a remainder that may be
differentiated once (Olver).  Subtracting term by term reproduces
\eqref{p6:eq:gamma-normal} with a remainder of the stated order.

Now let $x_N=\tfrac12+T+\sum_{n\le N}d_{2n-1}T^{1-2n}$ and
$f(x)=\log\Gamma(x)$, and let $\lambda=\log y$.  By construction the formal
substitution cancels the normalized equation through $q^{2N}$, so the residual
satisfies $f(x_N)-\lambda=\bigO(T^{-(2N+1)})$; the forward truncation error
just bounded is of the same order or smaller.  On an interval around $x_N$ of
length $\bigO(T^{-2N})$ we have $f'(x)=\psi(x)\ge\tfrac12\Lambda>0$ for all
large $y$, since $\psi(x)=\log x+\bigO(1/x)$ and $\log x_N=\Lambda+\littleo(1)$.
\Cref{p0:thm:backward-error} applied with $m=\tfrac12\Lambda$ converts the
residual into
$|x_N-\Gamma_+^{-1}(y)|\le2|f(x_N)-\lambda|/\Lambda
=\bigO(\Lambda^{-1}T^{-(2N+1)})$, which is \eqref{p6:eq:gamma-main}.  The
second half of \cref{p0:thm:backward-error} also guarantees that a root
exists within that distance, so no separate existence argument is needed.
\end{proof}

\subsection{Explicit blocks}

\begin{proposition}[The first five blocks]\label{p6:prop:gamma-blocks}
With $\Lambda=w+1$,
\begin{align}
 d_1&=\frac1{24\Lambda},\label{p6:eq:d1}\\
 d_3&=-\frac{14\Lambda^{2}+10\Lambda+5}{5760\,\Lambda^{3}},\label{p6:eq:d3}\\
 d_5&=\frac{2232\Lambda^{4}+1176\Lambda^{3}+714\Lambda^{2}
 +350\Lambda+105}{2903040\,\Lambda^{5}},\label{p6:eq:d5}\\
 d_7&=-\frac{822960\Lambda^{6}+292536\Lambda^{5}+136956\Lambda^{4}
 +68040\Lambda^{3}+32970\Lambda^{2}+12250\Lambda+2625}
 {1393459200\,\Lambda^{7}},\label{p6:eq:d7}\\
 d_9&=\frac{\begin{aligned}[t]
 &309052800\Lambda^{8}+77920128\Lambda^{7}+26024592\Lambda^{6}
  +10019592\Lambda^{5}\\
 &\quad{}+4516028\Lambda^{4}+2023560\Lambda^{3}+812350\Lambda^{2}
  +242550\Lambda+40425\end{aligned}}
 {367873228800\,\Lambda^{9}} .\label{p6:eq:d9}
\end{align}
Equivalently, separated into powers of $\Lambda^{-1}$,
\begin{align}
 d_3&=-\frac7{2880\Lambda}-\frac1{576\Lambda^{2}}
       -\frac1{1152\Lambda^{3}},\label{p6:eq:d3-sep}\\
 d_5&=\frac{31}{40320\Lambda}+\frac7{17280\Lambda^{2}}
      +\frac{17}{69120\Lambda^{3}}+\frac5{41472\Lambda^{4}}
      +\frac1{27648\Lambda^{5}} .\label{p6:eq:d5-sep}
\end{align}
\end{proposition}

\begin{proof}
Run \eqref{p6:eq:gamma-recurrence}.  The first two steps are short enough to
display.  Since
$(1+\epsilon)\log(1+\epsilon)-\epsilon=\tfrac12\epsilon^{2}
-\tfrac16\epsilon^{3}+\cdots$, the coefficient of $q^{2}$ is
$\Lambda d_1+a_1$, so $d_1=-a_1/\Lambda=1/(24\Lambda)$.  At $q^{4}$,
\[
 \Lambda d_3+\tfrac12d_1^{2}-a_1d_1+a_2=0,
\]
and substituting $a_1=-\tfrac1{24}$, $a_2=\tfrac7{2880}$,
$d_1=\tfrac1{24\Lambda}$ gives \eqref{p6:eq:d3}.  The remaining three were
computed by truncated-series arithmetic over $\Rat(\Lambda)$ and are recorded
in \cref{p6:rem:verification}.
\end{proof}

\begin{remark}[Independent verification]\label{p6:rem:verification}
Every rational number displayed in this chapter was recomputed for this
volume, in exact arithmetic, by an implementation written from the normalized
equations \eqref{p6:eq:gamma-normalized} and \eqref{p6:eq:barnes-normalized}
and not from any source's coefficient list: the forward data $a_k$ and $c_k$
from Bernoulli numbers, then $d_1,\dots,d_9$ over $\Rat(\Lambda)$ and
$b_0,\dots,b_5$ over $\Rat(\alpha,\lambda)$ by \eqref{p6:eq:gamma-recurrence}
and \eqref{p6:eq:barnes-recurrence}.  All agreed with the three sources, in
every digit of every numerator, including the $24$-term numerator of $b_5$.
The structural predictions were checked as well: the parity
\eqref{p6:eq:barnes-parity}, the resonant limits
\eqref{p6:eq:resonant-limit-coefficients} for $m\le3$, the deepest-pole
formulas \eqref{p6:eq:gamma-deepest} and \eqref{p6:eq:barnes-deepest} for
$n\le6$, and the two outer families of \cref{p6:prop:outer}.  Separately, the
nine high-precision roots that S2 prints were recomputed to $60$ digits and
agree in all $25$ digits displayed.
\end{remark}

\begin{remark}[Prior art]\label{p6:rem:prior-art}
The block $d_1$ is classical, and the expansion through $d_3$ --- in the
separated form \eqref{p6:eq:d3-sep} --- was proved by Gerhold, Hubalek and
Tomovski.  What the recurrence adds is that the expansion continues to all
orders with an explicit triangular rule, a proof of the parity, and the
coefficient ring.  The leading Lambert approximation $\Gamma_+^{-1}(y)\approx
\tfrac12+T$ itself goes back further; see Borwein--Corless and
Jeffrey--Watt for its history and for the branch geometry of the complex
inverse.
\end{remark}

\section{The branches of \texorpdfstring{$\Gamma^{-1}$}{Gamma inverse} at the poles}
\label{p6:sec:pole-branches}

The Lambert transseries describes one branch.  Every pole of $\Gamma$ carries
another, and those are convergent --- a useful contrast, since it isolates
exactly what makes the large branch divergent.

\begin{theorem}[The pole branches are convergent]\label{p6:thm:pole-branches}
Fix $n\in\{0,1,2,\dots\}$ and put
\begin{equation}\label{p6:eq:Phi-n}
 \Phi_n(\varepsilon)
 =\frac1{\Gamma(1+\varepsilon)}\prod_{k=1}^{n}
   \left(1-\frac\varepsilon k\right),
 \qquad \Phi_n(0)=1 .
\end{equation}
Then $1/\Gamma(-n+\varepsilon)=(-1)^{n}n!\,\varepsilon\,\Phi_n(\varepsilon)$,
and for all sufficiently large $|y|$ the branch of $\Gamma^{-1}$ tending to
$-n$ is given by the \emph{convergent} series
\begin{equation}\label{p6:eq:pole-lagrange}
 \Gamma^{-1}_{[-n]}(y)=-n+\sum_{m\ge1}\frac{\tau^{m}}{m}
 \Coef{\varepsilon^{m-1}}\Phi_n(\varepsilon)^{-m},
 \qquad \tau=\frac{(-1)^{n}}{n!\,y}.
\end{equation}
With $\psi_n=\psi(n+1)=H_n-\gamma$ and $H_n^{(r)}=\sum_{k\le n}k^{-r}$, the
first four terms are
\begin{equation}\label{p6:eq:pole-four}
 \Gamma^{-1}_{[-n]}(y)=-n+\tau+\psi_n\tau^{2}
 +\frac{3\psi_n^{2}+H_n^{(2)}+\zeta(2)}2\,\tau^{3}
 +\frac{8\psi_n^{3}+6\psi_n\bigl(H_n^{(2)}+\zeta(2)\bigr)
 +H_n^{(3)}-\zeta(3)}3\,\tau^{4}
 +\bigO(\tau^{5}).
\end{equation}
\end{theorem}

\begin{proof}
From $\Gamma(-n+\varepsilon)
=\Gamma(1+\varepsilon)\bigl/\bigl(\varepsilon(\varepsilon-1)\cdots
(\varepsilon-n)\bigr)$ we get
$1/\Gamma(-n+\varepsilon)
=\varepsilon\prod_{k=1}^{n}(\varepsilon-k)/\Gamma(1+\varepsilon)
=(-1)^{n}n!\,\varepsilon\Phi_n(\varepsilon)$.  Thus
$\tau=\varepsilon\Phi_n(\varepsilon)$, an analytic function of $\varepsilon$
near $0$ with derivative $1$ there; ordinary Lagrange inversion
(\cref{p0:thm:LB}, in its convergent analytic form) gives
\eqref{p6:eq:pole-lagrange}, convergent on a disc because $\Phi_n$ is
analytic and nonvanishing near $0$.

For the coefficients, expand
\[
 \log\Phi_n(\varepsilon)
 =\sum_{k=1}^{n}\log\Bigl(1-\frac\varepsilon k\Bigr)-\log\Gamma(1+\varepsilon)
 =-\psi_n\varepsilon-\frac{H_n^{(2)}+\zeta(2)}2\varepsilon^{2}
 +\frac{\zeta(3)-H_n^{(3)}}3\varepsilon^{3}-\cdots,
\]
using $\log\Gamma(1+\varepsilon)=-\gamma\varepsilon+\sum_{j\ge2}
(-1)^{j}\zeta(j)\varepsilon^{j}/j$ and $H_n-\gamma=\psi_n$.  Then
$\Phi_n^{-m}=\exp(-m\log\Phi_n)$ and
$\Coef{\varepsilon^{m-1}}$ of it gives, for $m=1,2,3,4$, the values
$1$, $2\psi_n$, $\tfrac92\psi_n^{2}+\tfrac32(H_n^{(2)}+\zeta(2))$ and
$\tfrac{32}3\psi_n^{3}+8\psi_n(H_n^{(2)}+\zeta(2))
-\tfrac43(\zeta(3)-H_n^{(3)})$; dividing by $m$ gives
\eqref{p6:eq:pole-four}.
\end{proof}

\begin{corollary}[The branch tending to zero]\label{p6:cor:zero-branch}
For $n=0$, $\tau=r:=1/y$ and, with $\gamma$ Euler's constant,
\begin{align}
 \Gamma_0^{-1}(y)={}&r-\gamma r^{2}
 +\frac{18\gamma^{2}+\pi^{2}}{12}r^{3}
 -\frac{8\gamma^{3}+\pi^{2}\gamma+\zeta(3)}3r^{4}
 \notag\\
 &+\frac{7500\gamma^{4}+1500\pi^{2}\gamma^{2}+2400\gamma\zeta(3)
 +29\pi^{4}}{1440}r^{5}
 \notag\\
 &-\left(\frac{54\gamma^{5}}5+3\gamma^{3}\pi^{2}+6\gamma^{2}\zeta(3)
 +\frac{17\gamma\pi^{4}}{120}+\frac{\pi^{2}\zeta(3)}6
 +\frac{\zeta(5)}5\right)r^{6}+\bigO(r^{7}),
 \label{p6:eq:zero-branch}
\end{align}
a convergent expansion.  Equivalently, in closed form,
\begin{equation}\label{p6:eq:zero-branch-bell}
 \Gamma_0^{-1}(y)=\sum_{m\ge1}\frac{\EB_{m-1}
 \bigl(m\ell_1,\dots,m\ell_{m-1}\bigr)}{m!\,y^{m}},
 \qquad
 \ell_1=-\gamma,\quad \ell_j=(-1)^{j}(j-1)!\,\zeta(j)\ (j\ge2),
\end{equation}
where $\EB_{m-1}$ is the complete exponential Bell polynomial of
\cref{p0:def:ebell}.
\end{corollary}

\begin{proof}
For $n=0$ the product in \eqref{p6:eq:Phi-n} is empty, so
$\tau=r=\varepsilon/\Gamma(1+\varepsilon)$ and $\psi_0=-\gamma$,
$H_0^{(r)}=0$.  Substituting into \eqref{p6:eq:pole-four} and using
$\zeta(2)=\pi^{2}/6$ gives the terms through $r^{4}$; the $r^{5}$ and $r^{6}$
terms were computed the same way and use in addition $\zeta(4)=\pi^{4}/90$.
For \eqref{p6:eq:zero-branch-bell}, Lagrange inversion of
$r=\varepsilon/\Gamma(1+\varepsilon)$ gives
$\Coef{r^{m}}=\tfrac1{m!}\bigl(\dd/\dd\varepsilon\bigr)^{m-1}
\Gamma(1+\varepsilon)^{m}\big|_{0}$, and
$\Gamma(1+\varepsilon)^{m}=\exp\bigl(m\sum_j\ell_j\varepsilon^{j}/j!\bigr)$,
whose $(m-1)$st derivative at $0$ is the complete exponential Bell polynomial
$\EB_{m-1}(m\ell_1,\dots,m\ell_{m-1})$.
\end{proof}

\begin{remark}[Why one branch converges and the other does not]
\label{p6:rem:convergence-contrast}
The contrast is instructive and is the cleanest illustration in the volume of
the distinction Chapter~0 insists on.  \Cref{p6:thm:pole-branches} inverts an
\emph{analytic} map with nonvanishing derivative at a point, so its Lagrange
series converges.  \Cref{p6:thm:gamma} inverts an \emph{asymptotic} series at
infinity: the identity it establishes is an equality of formal objects plus a
transported remainder bound at each fixed order, and the full series
$\sum_nd_{2n-1}T^{1-2n}$ diverges for every $y$, because the $a_n$ grow
factorially.  No amount of care with the algebra can change that; what can be
done is to truncate optimally, which is \cref{p0:thm:optimal-truncation} and
\cref{p6:sec:exponential}.
\end{remark}

\section{The eventual positive inverse of Barnes \texorpdfstring{$G$}{G}}
\label{p6:sec:barnes}

\subsection{Core variables and the normalized equation}

Let $y\to+\infty$ and put
\begin{equation}\label{p6:eq:barnes-core-vars}
 R=\log y-\zeta'(-1),
 \qquad
 v=\Wlam_0\!\left(\frac{4R}{\e^{3}}\right),
 \qquad
 U=2\sqrt{\frac Rv}=\e^{(v+3)/2},
 \qquad
 \lambda=\log U-1=\frac{v+1}2 .
\end{equation}
This is \cref{p6:lem:core} at $p=2$, $a=\tfrac12$, $\beta=\tfrac32$, $r=1$;
the core equation is
\begin{equation}\label{p6:eq:barnes-core-eq}
 \frac12U^{2}\left(\log U-\frac32\right)=R,
\end{equation}
and by \eqref{p6:eq:core-r1} its slope is
$\tfrac12U(v+1)=U\lambda$.  Write $q=U^{-1}$ and $z=U(1+\epsilon)$.
Subtracting \eqref{p6:eq:barnes-core-eq} from
\cref{p6:lem:barnes-normal-form} and dividing by $\tfrac12U^{2}$ gives
\begin{equation}\label{p6:eq:barnes-normalized}
 \mathcal{K}_{\lambda}(\epsilon)
 +2\alpha q(1+\epsilon)
 +2b\,q^{2}\bigl(\lambda+1+\log(1+\epsilon)\bigr)
 +2\sum_{k\ge1}c_kq^{2k+2}(1+\epsilon)^{-2k}=0,
 \qquad b=-\tfrac1{12},
\end{equation}
where, by \eqref{p6:eq:kernel-p2} with $v=2\lambda-1$,
\begin{equation}\label{p6:eq:barnes-kernel}
 \mathcal{K}_{\lambda}(\epsilon)
 =(1+\epsilon)^{2}\left(\lambda-\tfrac12+\log(1+\epsilon)\right)
 -\left(\lambda-\tfrac12\right)
 =2\lambda\epsilon+(\lambda+1)\epsilon^{2}
 +\tfrac13\epsilon^{3}-\tfrac1{12}\epsilon^{4}+\tfrac1{30}\epsilon^{5}-\cdots .
\end{equation}

\begin{remark}[Absorbing $\zeta'(-1)$]\label{p6:rem:absorb-zeta}
Taking $R=\log y-\zeta'(-1)$ rather than $R=\log y$ removes the constant from
the perturbation, in accordance with \cref{p6:rem:absorb}.  It is not
mandatory --- S2 does not do it --- but it makes the coefficients shorter, and
it is what makes the Glaisher constant disappear from them.
\Cref{p6:sec:dictionary} records the exact price of the other choice.
\end{remark}

\subsection{The all-orders theorem}

\begin{theorem}[All-orders expansion of the Barnes inverse]\label{p6:thm:barnes}
Let $\mathcal{B}$ be the inverse of \cref{p6:prop:barnes-monotone}, and let
$U,\lambda$ be as in \eqref{p6:eq:barnes-core-vars}.  There are uniquely
determined $e_n\in\Rat[\alpha,\lambda,\lambda^{-1}]$, $n\ge1$, such that with
$b_n:=e_{n+1}$ one has, for every fixed $N\ge0$,
\begin{equation}\label{p6:eq:barnes-main}
 \mathcal{B}(y)=U+\sum_{n=0}^{N}\frac{b_n(\lambda)}{U^{n}}
 +\bigO_N\!\left(U^{-N-1}\right),
 \qquad y\to+\infty,
\end{equation}
and $G_+^{-1}(y)=1+\mathcal{B}(y)$.  With
$E_{n-1}=\sum_{j<n}e_jq^{j}$ they are computed by
\begin{equation}\label{p6:eq:barnes-recurrence}
 e_n(\lambda)=-\frac1{2\lambda}\Coef{q^{n}}
 \Bigl[\mathcal{K}_{\lambda}(E_{n-1})-2\lambda E_{n-1}
 +2\alpha q(1+E_{n-1})
 +2bq^{2}\bigl(\lambda+1+\log(1+E_{n-1})\bigr)
 +2\!\!\sum_{2k+2\le n}\!\!c_kq^{2k+2}(1+E_{n-1})^{-2k}\Bigr],
\end{equation}
they obey the parity rule
\begin{equation}\label{p6:eq:barnes-parity}
 e_n(-\alpha)=(-1)^{n}e_n(\alpha),
\end{equation}
and each is a bounded function of $\lambda$ as $\lambda\to+\infty$.
\end{theorem}

\begin{proof}
\emph{Formal part.}  Equation \eqref{p6:eq:barnes-normalized} has support
$S=\Nat$ by \cref{p6:ex:supports} and linear coefficient
$\kappa_1=v+1=2\lambda$; \cref{p6:thm:triangular} applies and yields
\eqref{p6:eq:barnes-recurrence} and uniqueness.  The coefficient ring follows
by induction, since the Taylor coefficients of $\log(1+\epsilon)$ and of the
binomial powers are rational, $\alpha$ enters only through $2\alpha q$, and
the only division is by $2\lambda$.

For \eqref{p6:eq:barnes-parity}, observe that the substitution
$(q,\alpha,\epsilon)\mapsto(-q,-\alpha,\epsilon)$ leaves
\eqref{p6:eq:barnes-normalized} invariant: $\mathcal{K}_\lambda(\epsilon)$ has
no $q$; $2\alpha q(1+\epsilon)$ is invariant because both signs flip; and
$q^{2}$, $q^{2k+2}$ are even.  Uniqueness therefore forces
$\epsilon(-q,-\alpha)=\epsilon(q,\alpha)$, i.e.\ $(-1)^{n}e_n(-\alpha)=
e_n(\alpha)$, which is \eqref{p6:eq:barnes-parity}.

Boundedness in $\lambda$ is \cref{p6:prop:resonance} below, which exhibits the
coefficientwise limit.

\emph{Analytic part.}  Let $x_N=U+\sum_{n\le N}b_nU^{-n}$ and
$f(z)=\log G(z+1)$, $\lambda_\ast=\log y$.  Retain the forward expansion
\eqref{p6:eq:barnes-normal} far enough that its remainder is
$\bigO(U^{-N-1})$ after division by the slope; this is legitimate because
\eqref{p6:eq:barnes-normal} is a Poincar\'e expansion with differentiable
remainder on the positive axis, with explicit bounds available from Nemes.
By construction the formal substitution cancels
\eqref{p6:eq:barnes-normalized} through $q^{N+1}$, so the forward residual
$f(x_N)-\lambda_\ast$ is $\bigO(\lambda U^{-N})$ --- one factor $\tfrac12U^{2}$
is restored on undoing the normalization, and the first omitted normalized
term is $2\lambda e_{N+2}q^{N+2}$.  By \cref{p6:prop:barnes-monotone},
$f'\ge\tfrac12z\log z$ on $[7,\infty)$, and along the relevant interval
$f'=U\lambda(1+\littleo(1))$.  \Cref{p0:thm:backward-error} with
$m=\tfrac12U\lambda$ turns the residual into
$|x_N-\mathcal{B}(y)|=\bigO(U^{-N-1})$, which is \eqref{p6:eq:barnes-main},
and again certifies that a root lies that close.
\end{proof}

\subsection{Explicit blocks}

\begin{proposition}[The first six blocks]\label{p6:prop:barnes-blocks}
With $\alpha=\tfrac12\log(2\pi)$ and $\lambda=(v+1)/2$,
\begin{align}
 b_0&=-\frac\alpha\lambda,\label{p6:eq:b0}\\
 b_1&=\frac{6\alpha^{2}\lambda-6\alpha^{2}+\lambda^{3}+\lambda^{2}}
 {12\lambda^{3}},\label{p6:eq:b1}\\
 b_2&=\frac{\alpha\bigl(8\alpha^{2}\lambda-6\alpha^{2}+\lambda^{2}\bigr)}
 {12\lambda^{5}},\label{p6:eq:b2}\\
 b_3&=-\frac{P_3(\lambda,\alpha)}{1440\lambda^{7}},
 \qquad
 b_4=-\frac{\alpha P_4(\lambda,\alpha)}{1440\lambda^{9}},
 \qquad
 b_5=\frac{P_5(\lambda,\alpha)}{362880\lambda^{11}},\label{p6:eq:b345}
\end{align}
where
\begin{align}
 P_3&=180\alpha^{4}\lambda^{3}+120\alpha^{4}\lambda^{2}
 -1500\alpha^{4}\lambda+900\alpha^{4}
 +60\alpha^{2}\lambda^{5}+120\alpha^{2}\lambda^{4}
 +60\alpha^{2}\lambda^{3}-180\alpha^{2}\lambda^{2}\notag\\
 &\qquad+5\lambda^{7}-\lambda^{6}+5\lambda^{5}+5\lambda^{4},
 \label{p6:eq:P3}\\
 P_4&=576\alpha^{4}\lambda^{3}+480\alpha^{4}\lambda^{2}
 -2520\alpha^{4}\lambda+1260\alpha^{4}
 +80\alpha^{2}\lambda^{5}+280\alpha^{2}\lambda^{4}
 +200\alpha^{2}\lambda^{3}-300\alpha^{2}\lambda^{2}\notag\\
 &\qquad-22\lambda^{7}-11\lambda^{6}+10\lambda^{5}+15\lambda^{4},
 \label{p6:eq:P4}\\
 P_5&=22680\alpha^{6}\lambda^{5}+40824\alpha^{6}\lambda^{4}
 -349272\alpha^{6}\lambda^{3}-352800\alpha^{6}\lambda^{2}
 +1111320\alpha^{6}\lambda-476280\alpha^{6}\notag\\
 &\qquad+11340\alpha^{4}\lambda^{7}+34020\alpha^{4}\lambda^{6}
 +3780\alpha^{4}\lambda^{5}-138600\alpha^{4}\lambda^{4}
 -132300\alpha^{4}\lambda^{3}+132300\alpha^{4}\lambda^{2}\notag\\
 &\qquad+1890\alpha^{2}\lambda^{9}+8568\alpha^{2}\lambda^{8}
 +15120\alpha^{2}\lambda^{7}+11718\alpha^{2}\lambda^{6}
 -3150\alpha^{2}\lambda^{5}-9450\alpha^{2}\lambda^{4}\notag\\
 &\qquad+105\lambda^{11}-668\lambda^{10}-378\lambda^{9}-21\lambda^{8}
 +175\lambda^{7}+105\lambda^{6}.
 \label{p6:eq:P5}
\end{align}
\end{proposition}

\begin{proof}
Run \eqref{p6:eq:barnes-recurrence}; see \cref{p6:rem:verification}.  The
first step is $2\lambda e_1+2\alpha=0$, giving $b_0=e_1=-\alpha/\lambda$.  The
parity \eqref{p6:eq:barnes-parity} is visible: $b_0,b_2,b_4$ are odd in
$\alpha$ and $b_1,b_3,b_5$ even, since $b_n=e_{n+1}$.  The denominators are
$\lambda^{2n-1}$ for $e_n$, in agreement with \cref{p6:prop:denominator}, and
the numerators' extreme coefficients $-\alpha$, $-6\alpha^{2}/12$,
$-6\alpha^{3}/12$, $-900\alpha^{4}/1440$, $-1260\alpha^{5}/1440$,
$-476280\alpha^{6}/362880$ are $-\alpha,-\tfrac{\alpha^2}2,-\tfrac{\alpha^3}2,
-\tfrac58\alpha^{4},-\tfrac78\alpha^{5},-\tfrac{21}{16}\alpha^{6}$, which is
\eqref{p6:eq:barnes-deepest}.
\end{proof}

\subsection{Resonance: the surviving algebraic subseries}

The forward term $-\tfrac1{12}\log z$ is small relative to the core, but its
\emph{logarithm} is of the same size as the slope.  It therefore survives the
large-slope limit, and it is the only perturbation that does.

\begin{proposition}[Coefficientwise large-slope limit]\label{p6:prop:resonance}
For each fixed $n$, the limit of $e_n(\lambda)$ as $\lambda\to+\infty$ exists
and
\begin{equation}\label{p6:eq:resonant-limit-coefficients}
 \lim_{\lambda\to\infty}e_{2m}(\lambda)=\binom{1/2}{m}6^{-m},
 \qquad
 \lim_{\lambda\to\infty}e_{2m+1}(\lambda)=0 .
\end{equation}
Equivalently $\lim b_{2m-1}=\binom{1/2}{m}6^{-m}$ and $\lim b_{2m}=0$; the
first values are $b_1\to\tfrac1{12}$, $b_3\to-\tfrac1{288}$,
$b_5\to\tfrac1{3456}$.  The limiting coefficients are those of
\begin{equation}\label{p6:eq:resonant-limit}
 E_0(q)=\sqrt{1+\tfrac{q^{2}}6}-1,
 \qquad\text{i.e.}\qquad
 U\bigl(1+E_0(U^{-1})\bigr)=\sqrt{U^{2}+\tfrac16}.
\end{equation}
\end{proposition}

\begin{proof}
Divide \eqref{p6:eq:barnes-normalized} by $2\lambda$ and let
$\lambda\to\infty$ with $q$ formal.  By \eqref{p6:eq:barnes-kernel},
$\mathcal{K}_\lambda(E)/(2\lambda)\to\bigl((1+E)^{2}-1\bigr)/2$, because the
$\lambda$-free part of $\mathcal{K}_\lambda$ is bounded coefficientwise.  The
term $2\alpha q(1+E)/(2\lambda)\to0$; each $c_k$ term $\to0$; and
\[
 \frac{2bq^{2}\bigl(\lambda+1+\log(1+E)\bigr)}{2\lambda}
 =bq^{2}\left(1+\frac{1+\log(1+E)}{\lambda}\right)\longrightarrow bq^{2}.
\]
So the limiting equation is $\bigl((1+E_0)^{2}-1\bigr)/2+bq^{2}=0$, i.e.\
$(1+E_0)^{2}=1-2bq^{2}=1+q^{2}/6$, whose branch with $E_0(0)=0$ is
\eqref{p6:eq:resonant-limit}.  That the coefficientwise limit may be taken
inside the recursion is the triangular structure: at each $n$ the coefficient
equation has limiting linear divisor $1\ne0$ and involves only finitely many
lower coefficients, so induction commutes with the limit.  Expanding
$\sqrt{1+q^{2}/6}$ gives \eqref{p6:eq:resonant-limit-coefficients}, and
$z=U(1+\epsilon)$ turns $E_0$ into $\sqrt{U^{2}+1/6}$.
\end{proof}

\begin{remark}[Two limits, not one]\label{p6:rem:limit-caveat}
\Cref{p6:prop:resonance} is a statement about each coefficient separately.  It
does \emph{not} say that $\mathcal{B}(y)-\sqrt{U^{2}+1/6}$ is small: at any
actual $y$ the parameter $\lambda$ is finite and the neglected terms, headed
by $b_0=-\alpha/\lambda$, are numerically the largest corrections in the whole
expansion.  S1 writes that the subsector ``resums'' to $\sqrt{T^{2}+1/6}$ and
then, in the following remark, correctly warns that it does not replace the
expansion; the two statements sit awkwardly together.  The accurate form is
the one above: a limit of coefficients, useful as an exact structural check on
a symbolic computation and as the identification of one coherent subseries,
not an approximation to the inverse.
\end{remark}

\section{The three normalizations reconciled}
\label{p6:sec:dictionary}

The three sources use three coordinate systems.  They agree, and the
dictionary is worth recording because the disagreement is entirely in the
choice of what to absorb.

\begin{proposition}[Dictionary]\label{p6:prop:dictionary}
Let $C_n(D)$ denote the Barnes coefficients obtained from
\eqref{p6:eq:barnes-normalized} when $\zeta'(-1)$ is \emph{not} absorbed, so
that $R=\log y$, $Z=\e^{(v+3)/2}$, $D=v+1$, and the constant
$\kappa=\zeta'(-1)=\tfrac1{12}-\log A$ appears as the extra perturbation
$2\kappa q$.  Then the two families are related by
\begin{equation}\label{p6:eq:dictionary}
 b_n(\lambda,\alpha)=C_n(D)\Big|_{\,\log A\;\mapsto\;1/12,\
 \log(2\pi)\;\mapsto\;2\alpha,\ D\;\mapsto\;2\lambda},
\end{equation}
that is, the absorbed family is obtained from the unabsorbed one by setting
$\kappa=0$ and halving the slope variable.  In particular
\begin{equation}\label{p6:eq:dictionary-C1}
 C_1(D)=\frac1{12}+\frac{2\log A}{D}
 +\frac{\log^{2}(2\pi)}{2D^{2}}-\frac{\log^{2}(2\pi)}{D^{3}},
 \qquad
 b_1(\lambda)=\frac1{12}+\frac1{12\lambda}
 +\frac{\alpha^{2}}{2\lambda^{2}}-\frac{\alpha^{2}}{2\lambda^{3}},
\end{equation}
and these agree under \eqref{p6:eq:dictionary}.
\end{proposition}

\begin{proof}
Both families solve the same normalized equation, the second with the extra
term $2\kappa q$ and with $\lambda$ replaced by $D/2$.  Since $\kappa$ enters
only through that term and only linearly, and since $\kappa=\tfrac1{12}-\log
A$ vanishes exactly when $\log A=\tfrac1{12}$, the substitution
$\log A\mapsto\tfrac1{12}$ turns the unabsorbed equation into the absorbed
one; uniqueness of the solution does the rest.  The two displayed forms of the
first coefficient were checked directly, as recorded in
\cref{p6:rem:verification}.
\end{proof}

\begin{remark}[A cancellation worth seeing]\label{p6:rem:cancellation}
In $C_1$ the constant $\log A$ appears and $\tfrac1{6D}$ does not; in $b_1$
the reverse.  This is not a discrepancy.  The forward term
$-\tfrac1{12}\log z$ contributes $\tfrac1{12}+\tfrac1{12\lambda}$ to both, and
the unabsorbed constant contributes $-2\kappa/D=-\tfrac1{6D}+\tfrac{2\log A}D$,
whose first half exactly cancels the $\tfrac1{6D}=\tfrac1{12\lambda}$.  A
reader comparing the two articles term by term without tracking the
normalization will conclude that one of them has a sign error.  Neither does.
\end{remark}

\section{Outer expansions in iterated logarithms}
\label{p6:sec:outer}

Retaining $\Wlam$ exactly is both shorter and numerically better, but an
answer purely in elementary logarithms is sometimes wanted.  It is obtained by
reverting the core equation directly, without first invoking a generic
$\Wlam$-series.

\begin{proposition}[The two outer families]\label{p6:prop:outer}
\leavevmode
\begin{enumerate}
\item \emph{Gamma.}  Put $L=\log R$, $\ell=\log L$, $D_\Gamma=\ell+1$, and
write $T=(R/L)F$.  Then $F$ is the unique solution with $F(\infty)=1$ of
$F\bigl(1-D_\Gamma/L+L^{-1}\log F\bigr)=1$, and
$F=1+\sum_{n\ge1}f_nL^{-n}$ with
\begin{align}
 f_1&=D,\quad f_2=D(D-1),\quad f_3=\tfrac12D(D-2)(2D-1),\notag\\
 f_4&=\tfrac16D(6D^{3}-26D^{2}+27D-6),\notag\\
 f_5&=\tfrac1{12}D(12D^{4}-77D^{3}+142D^{2}-84D+12),\notag\\
 f_6&=\tfrac1{30}D(30D^{5}-261D^{4}+725D^{3}-775D^{2}+300D-30),
 \label{p6:eq:outer-gamma}
\end{align}
abbreviating $D=D_\Gamma$.
\item \emph{Barnes.}  Put $L=\log(4R)$, $\ell=\log L$, $D_G=\ell+3$, and
write $U=2\sqrt{R/L}\,F$.  Then $F$ solves
$F^{2}\bigl(1-D_G/L+2L^{-1}\log F\bigr)=1$ and
$F=1+\sum_{n\ge1}g_nL^{-n}$ with
\begin{align}
 g_1&=\tfrac12D,\quad g_2=\tfrac18D(3D-4),\quad
 g_3=\tfrac1{16}D(5D^{2}-16D+8),\notag\\
 g_4&=\tfrac1{384}D(105D^{3}-568D^{2}+720D-192),\notag\\
 g_5&=\tfrac1{768}D(189D^{4}-1488D^{3}+3296D^{2}-2304D+384),\notag\\
 g_6&=\tfrac1{15360}D(3465D^{5}-36516D^{4}+120400D^{3}-150400D^{2}
 +67200D-7680),
 \label{p6:eq:outer-barnes}
\end{align}
abbreviating $D=D_G$.
\end{enumerate}
\end{proposition}

\begin{proof}
For (i), $T(\log T-1)=R$ and $\log T=L-\ell+\log F$ give
$(R/L)F\,(L-\ell-1+\log F)=R$, i.e.\ the displayed equation.  Substituting
$F=1+\sum f_nL^{-n}$ and equating coefficients of $L^{-n}$ determines $f_n$
uniquely from $f_1,\dots,f_{n-1}$, because the coefficient of $f_n$ is $1$.
For (ii), write $L=\log(4R)$, so that $\log(4R/L)=L-\ell$ and
$\log U=\tfrac12(L-\ell)+\log F$.  Substituting $U=2\sqrt{R/L}\,F$ into
$\tfrac12U^{2}(\log U-\tfrac32)=R$ gives
\[
 \frac{2R}{L}F^{2}\left(\frac{L-\ell}2-\frac32+\log F\right)=R,
 \qquad\text{i.e.}\qquad
 F^{2}\left(1-\frac{\ell+3}{L}+\frac{2\log F}{L}\right)=1 ,
\]
which is the displayed equation with $D_G=\ell+3$.  Both families were recomputed independently
(\cref{p6:rem:verification}).
\end{proof}

\begin{remark}[The outer variable is part of the answer]
\label{p6:rem:outer-variable}
The polynomials $g_n$ in \eqref{p6:eq:outer-barnes} are attached to the choice
$L=\log(4R)$, $D_G=\log L+3$.  Choosing instead
$\widetilde L=\log(4R\e^{-3})=L-3$ and $U=2\sqrt{R/\widetilde L}\,F$ leads to
the \emph{same} polynomials evaluated at $D=\log\widetilde L$ --- the shift
$+3$ migrates from the argument into the variable.  S1 uses the first
convention and S3 the second, and both are correct; but the two lists of
coefficients are only equal after the substitution is made, and comparing them
term by term without it produces a spurious disagreement at every order.  The
same caution applies to the Gamma family, where the corresponding shift is
$+1$.
\end{remark}

\begin{remark}[The hierarchy, and why it must not be merged]
\label{p6:rem:hierarchy}
For every fixed $N$, $T/L^{N}\to\infty$; so the constant $\tfrac12$ in
\cref{p6:thm:gamma} is smaller than every term of the outer expansion of $T$,
and the first Bernoulli block $d_1/T$ is smaller again.  The correct reading
is lexicographic,
\begin{equation}\label{p6:eq:gamma-hierarchy}
 T\ \gg\ \frac TL\ \gg\ \frac T{L^{2}}\ \gg\ \cdots\ \gg\ 1
 \ \gg\ \frac1{T\Lambda}\ \gg\ \frac1{T^{3}\Lambda}\ \gg\ \cdots,
\end{equation}
and for Barnes, with the extra intermediate level supplied by $b_0$,
\begin{equation}\label{p6:eq:barnes-hierarchy}
 U\ \gg\ \frac U{L}\ \gg\ \cdots\ \gg\ 1\ \gg\ \frac1\lambda
 \ \gg\ \frac1U\ .
\end{equation}
An expansion that interleaves the outer $L^{-1}$ series with the additive
constants and the $q$-grid is not merely inelegant, it is wrong: no finite
truncation of the outer series can reproduce even the first block $d_1/T$.
This is the same phenomenon as the flattening caveat of
\cref{p0:thm:flattening}, and the reason the volume keeps $\Wlam$ unexpanded
by default.
\end{remark}

\section{Numerical verification}
\label{p6:sec:numerics}

The tables below test the entire chain and involve no numerical inversion.  An
exact abscissa is chosen, the special function is evaluated at $80$ digits,
the core variables are reconstructed from that value alone, and the
truncations are compared with the abscissa.  Errors are signed, approximation
minus exact.

\begin{table}[htbp]
\centering
\caption{$\Gamma_+^{-1}$: signed error of $\tfrac12+T$ and of the truncations
including $d_1,\dots,d_9$.}
\label{p6:tab:gamma-errors}
\scriptsize
\setlength{\tabcolsep}{3.6pt}
\begin{tabular}{rcccccc}
\toprule
$x$ & $\tfrac12+T$ & $+\,d_1$ & $+\,d_3$ & $+\,d_5$ & $+\,d_7$ & $+\,d_9$\\
\midrule
5 & $-6.1414\!\times\!10^{-3}$ & $2.8694\!\times\!10^{-5}$ & $-4.1862\!\times\!10^{-7}$ & $1.3452\!\times\!10^{-8}$ & $-8.2694\!\times\!10^{-10}$ & $8.4558\!\times\!10^{-11}$\\
10 & $-1.9470\!\times\!10^{-3}$ & $1.7430\!\times\!10^{-6}$ & $-5.7560\!\times\!10^{-9}$ & $4.4485\!\times\!10^{-11}$ & $-6.5578\!\times\!10^{-13}$ & $1.5870\!\times\!10^{-14}$\\
20 & $-7.1924\!\times\!10^{-4}$ & $1.4126\!\times\!10^{-7}$ & $-1.1183\!\times\!10^{-10}$ & $2.1113\!\times\!10^{-13}$ & $-7.5612\!\times\!10^{-16}$ & $4.4191\!\times\!10^{-18}$\\
50 & $-2.1572\!\times\!10^{-4}$ & $6.1957\!\times\!10^{-9}$ & $-7.6816\!\times\!10^{-13}$ & $2.2927\!\times\!10^{-16}$ & $-1.2910\!\times\!10^{-19}$ & $1.1809\!\times\!10^{-22}$\\
100 & $-9.1031\!\times\!10^{-5}$ & $6.2869\!\times\!10^{-10}$ & $-1.9388\!\times\!10^{-14}$ & $1.4440\!\times\!10^{-18}$ & $-2.0231\!\times\!10^{-22}$ & $4.5949\!\times\!10^{-26}$\\
\bottomrule
\end{tabular}
\end{table}

\begin{table}[htbp]
\centering
\caption{$\mathcal{B}$: signed error of $U$ and of the truncations including
$b_0,\dots,b_5$, at $y=G(z+1)$.}
\label{p6:tab:barnes-errors}
\scriptsize
\setlength{\tabcolsep}{3.0pt}
\begin{tabular}{rccccccc}
\toprule
$z$ & $U$ & $+\,b_0$ & $+\,b_1$ & $+\,b_2$ & $+\,b_3$ & $+\,b_4$ & $+\,b_5$\\
\midrule
5 & $1.1187$ & $-1.3867\!\times\!10^{-2}$ & $-7.8343\!\times\!10^{-3}$ & $-1.5935\!\times\!10^{-3}$ & $-1.8345\!\times\!10^{-4}$ & $5.8515\!\times\!10^{-6}$ & $8.9807\!\times\!10^{-6}$\\
10 & $6.5273\!\times\!10^{-1}$ & $-2.0081\!\times\!10^{-2}$ & $-8.4026\!\times\!10^{-4}$ & $1.5249\!\times\!10^{-5}$ & $3.2299\!\times\!10^{-6}$ & $4.1936\!\times\!10^{-8}$ & $-1.3999\!\times\!10^{-8}$\\
20 & $4.4669\!\times\!10^{-1}$ & $-8.7205\!\times\!10^{-3}$ & $-6.6821\!\times\!10^{-5}$ & $2.3736\!\times\!10^{-6}$ & $3.4133\!\times\!10^{-8}$ & $-1.5564\!\times\!10^{-9}$ & $-1.9947\!\times\!10^{-11}$\\
100 & $2.5342\!\times\!10^{-1}$ & $-1.2959\!\times\!10^{-3}$ & $-3.9481\!\times\!10^{-7}$ & $8.1326\!\times\!10^{-9}$ & $-5.3660\!\times\!10^{-12}$ & $-8.9993\!\times\!10^{-14}$ & $6.0482\!\times\!10^{-16}$\\
\bottomrule
\end{tabular}
\end{table}

Three things are visible.  The Gamma expansion gains roughly two orders per
block, as \eqref{p6:eq:gamma-main} predicts, and does so already at $x=5$.
The Barnes core error is much larger, because the linear forward term
$\alpha z$ produces the slowly decaying displacement $b_0=-\alpha/\lambda$;
once that single block is included the improvement is again rapid.  And at
$z=5$ the last Barnes column is \emph{worse} than the one before it
($8.98\times10^{-6}$ against $5.85\times10^{-6}$): the expansion is asymptotic,
not convergent, and at a fixed moderate argument it has reached its useful
limit.  That is the behaviour \cref{p0:thm:optimal-truncation} describes and
\cref{p6:sec:exponential} quantifies; it is not a defect of the coefficients,
each of which was verified exactly.

\begin{remark}[What the tables do and do not establish]
\label{p6:rem:tables-caveat}
They check the algebra, the signs, the shifts and the branch choices, end to
end from the special function to the inverse; a transcription error in any
numerator would show immediately.  They are not a substitute for the analytic
remainder arguments in the proofs of
\cref{p6:thm:gamma,p6:thm:barnes}, and conversely those theorems, which
control only the order of the remainder, would not catch a wrong rational
number.  The two kinds of evidence are complementary, and both were obtained.
\end{remark}

\section{Beyond all orders}
\label{p6:sec:exponential}

The expansions above are complete relative to a Poincar\'e forward input.
They do not contain the exponentially small sectors that the late Bernoulli
terms encode.

\begin{proposition}[Transport of an exponentially small sector]
\label{p6:prop:exponential-transport}
Suppose the forward phase splits as $\Phi=\Phi_{\mathrm{alg}}+E$, where
$x_{\mathrm{alg}}$ solves $\Phi_{\mathrm{alg}}(x)=R$ to the algebraic order
considered and $E$ is smaller than every term retained.  If
$\Phi_{\mathrm{alg}}'(x_{\mathrm{alg}})\ne0$, then the induced displacement of
the root is
\begin{equation}\label{p6:eq:exponential-transport}
 \Delta x
 =-\frac{E(x_{\mathrm{alg}})}{\Phi_{\mathrm{alg}}'(x_{\mathrm{alg}})}
 +\bigO\!\left(
 \frac{|E|\,|E'|}{|\Phi_{\mathrm{alg}}'|^{2}}
 +\frac{|E|^{2}\,|\Phi_{\mathrm{alg}}''|}{|\Phi_{\mathrm{alg}}'|^{3}}
 \right).
\end{equation}
\end{proposition}

\begin{proof}
This is \cref{p0:thm:remainder-transport} with $f_0=\Phi_{\mathrm{alg}}$ and
$\varepsilon=E$: part (ii) of that theorem gives exactly
\eqref{p6:eq:exponential-transport}, the two error terms being its $\theta$
and $M_2$ contributions.
\end{proof}

\begin{proposition}[Least term and the inverse floor]\label{p6:prop:least-term}
The half-shifted Stirling coefficients satisfy
\begin{equation}\label{p6:eq:ak-growth}
 a_n=(-1)^{n}\,\frac{2\,(2n-2)!}{(2\pi)^{2n}}\bigl(1+\bigO(4^{-n})\bigr),
\end{equation}
so the Gamma inverse blocks $d_{2n-1}\sim-a_n/\Lambda$ decrease until
$n\approx\pi T$ and grow thereafter.  Consequently the least term of
\eqref{p6:eq:gamma-main} has size
$\exp\bigl(-2\pi T+\bigO(\log T)\bigr)$, and by
\cref{p0:thm:backward-error} the attainable accuracy of the inverse is
\begin{equation}\label{p6:eq:gamma-floor}
 \frac{\exp\bigl(-2\pi T+\bigO(\log T)\bigr)}{\Lambda}.
\end{equation}
The Barnes analogue, from $c_k=B_{2k+2}/(4k(k+1))$, has least term near
$k\approx\pi U$ and floor
$\exp\bigl(-2\pi U+\bigO(\log U)\bigr)/(U\lambda)$.
\end{proposition}

\begin{proof}
From $B_{2n}=(-1)^{n+1}2(2n)!\,\zeta(2n)/(2\pi)^{2n}$ and
$\zeta(2n)=1+\bigO(4^{-n})$, together with
$a_n=(2^{1-2n}-1)B_{2n}/\bigl(2n(2n-1)\bigr)$ and $2^{1-2n}-1=-1+\bigO(4^{-n})$,
we get \eqref{p6:eq:ak-growth} using $(2n)!/(2n(2n-1))=(2n-2)!$.  The ratio of
consecutive terms $|a_{n+1}|T^{-2n-1}/\bigl(|a_n|T^{1-2n}\bigr)$ is
$2n(2n-1)/\bigl((2\pi)^{2}T^{2}\bigr)(1+\littleo(1))$, which equals $1$ at
$n\sim\pi T$; substituting that index into \eqref{p6:eq:ak-growth} and using
Stirling gives the stated size, and dividing by the slope $\Lambda$ transports
it to the inverse.  The Barnes computation is identical with
$B_{2k+2}$ in place of $B_{2n}$.
\end{proof}

\begin{remark}[What is not claimed]\label{p6:rem:not-claimed}
\Cref{p6:prop:least-term} locates the floor; it does not compute the Stokes
multipliers, and this chapter does not.  A genuine hyperasymptotic inverse
must start from an exponentially improved \emph{forward} expansion on the same
branch and in the same sector --- Nemes supplies these for both functions ---
and then apply \eqref{p6:eq:exponential-transport}.  Appending terms of the
shape $\e^{-2\pi T}$ to the real formula by pattern matching is not a
branch-complete construction, and none of the three sources does more than
indicate the first sector.  This is recorded as open in
\cref{p6:sec:open}.
\end{remark}

\section{The general family}
\label{p6:sec:general}

Nothing above used more about $\Gamma$ and $G$ than the shape of their
phases.  The following is the class the chapter actually treats, and it is
what \cref{p7:sec:top} will reuse.

\begin{theorem}[Inversion of a power--logarithmic phase]\label{p6:thm:general}
Let
\begin{equation}\label{p6:eq:general-phase}
 \Phi(x)=a\,x^{p}(\log x-\beta)+C
 +\sum_{\rho\in A}x^{p-\rho}P_\rho(\log x)+\mathcal{E}(x),
\end{equation}
where $a\ne0$, $p>0$, $A\subset\Real_{>0}$ generates a well-ordered monoid $S$,
each $P_\rho$ is a polynomial, and $\mathcal{E}$ is smaller than every
retained term.  Put $R=Y-C$ and let $v,X,q$ be given by
\eqref{p6:eq:core-r1}.  If $v+1$ is invertible then the equation
$\Phi(x)=Y$ has a unique formal solution
\begin{equation}\label{p6:eq:general-inverse}
 x=X\Bigl(1+\sum_{\mu\in S,\ \mu>0}e_\mu(v)\,X^{-\mu}\Bigr),
\end{equation}
with $e_\mu$ given by \eqref{p6:eq:triangular}; the coefficients lie in
$K[v,(v+1)^{-1}]$ where $K$ is any field containing the coefficients of the
$P_\rho$ and $\beta$; and their denominators obey
\cref{p6:prop:denominator} when $S=\Nat$.  If in addition $\Phi$ is real,
eventually strictly increasing, with a differentiable Poincar\'e remainder at
each truncation order, then each truncation of \eqref{p6:eq:general-inverse}
approximates the actual inverse with the error obtained by dividing the
forward residual by $a X^{p-1}(v+1)$.
\end{theorem}

\begin{proof}
Substituting $x=X(1+\epsilon)$ into \eqref{p6:eq:general-phase} and dividing
by $aX^{p}$ gives \eqref{p6:eq:normalized} with
\[
 \mathcal{R}(q,v,\epsilon)
 =\frac1a\sum_{\rho\in A}q^{\rho}(1+\epsilon)^{p-\rho}
 P_\rho\!\left(\beta+\frac vp+\log(1+\epsilon)\right)
 +\frac{\mathcal{E}\bigl(X(1+\epsilon)\bigr)}{aX^{p}},
\]
every monomial of which has $q$-exponent in $S$ and positive.  Powers of $v$
in the coefficients are harmless: $q^{\rho}v^{M}\to0$ for fixed $M$, and the
Hahn ring is graded by the $q$-exponent alone.  \Cref{p6:thm:triangular} gives
existence, uniqueness and the recurrence; \cref{p6:cor:support} the support;
the coefficient ring follows because the only division performed is by $v+1$
and every other operation is a rational combination of the data.  The analytic
clause is \cref{p0:thm:backward-error} applied with the slope
$\Phi'(X)=aX^{p-1}(v+1)(1+\littleo(1))$ from \eqref{p6:eq:core-r1}.
\end{proof}

\begin{example}[The two subjects, and a third]\label{p6:ex:instances}
$\Gamma$ is $(a,p,\beta)=(1,1,1)$ with $A=\{2,4,6,\dots\}$; $G$ is
$(\tfrac12,2,\tfrac32)$ with $A=\{1,2,4,6,\dots\}$, the exponent $1$ coming
from $\alpha z$.  The hyperfactorial of \cref{p7:sec:top} is the next member,
and the higher multiple gamma functions continue the family with $p=3,4,\dots$,
their perturbation supports read off from the drops in algebraic degree in the
known forward expansions.
\end{example}

\begin{remark}[Where the closed form stops]\label{p6:rem:limits-of-method}
Two boundaries deserve naming.  A leading block
$a x^{p}(\log x-\beta)^{r}$ with $r\ne1$ is still exactly invertible, by
\cref{p6:lem:core}, with normalized slope $(pu+r)/u$ in place of $v+1$.  A
leading block $ax^{p}(\log x)^{r}(\log\log x)^{s}$ with $s\ne0$ is not: there
is no one-step elementary reduction, and the honest procedure is to invert the
first two factors exactly, treat the third as a slow multiplier, and solve the
residual slow equation by iteration before the $q$-grid theorem is applied.
The method's real hypothesis is not ``Lambert'' but ``the dominant
transmonomial can be inverted exactly and its normalized slope is a unit''.
\end{remark}

\section{What is left open}
\label{p6:sec:open}

\begin{enumerate}
\item \emph{Closed forms for the coefficient families.}  The recurrences
compute $d_{2n-1}$ and $b_n$ to any order and \cref{p6:thm:deepest-pole}
gives the extreme coefficient of each in closed form, but no closed
description of the full numerators is known.  For Gamma the pattern
$d_{2n-1}=P_n(\Lambda)/\bigl(\gamma_n\Lambda^{2n-1}\bigr)$ with $P_n$ of degree
$2n-2$ invites a Bell-polynomial formula decorated by Bernoulli values at
$\tfrac12$; for Barnes the grading by powers of $\alpha$ is a
$\Int/2$-grading whose elementary proof is \eqref{p6:eq:barnes-parity} but
whose finer structure --- which powers $\alpha^{j}$ actually occur in $b_n$,
and with what $\lambda$-degrees --- is not determined here.
\item \emph{Hyperasymptotics.}  \Cref{p6:rem:not-claimed}.
\item \emph{Uniformity near the branch point.}  Everything degrades as
$v+1\to0$, and at $v=-1$ the core map has a critical point where the correct
local scale is a square root rather than a power series.  A uniform
approximation matching a branch-point expansion to the transseries developed
here is not attempted.  It is irrelevant for the positive real branches, which
stay far from $v=-1$, and unavoidable for a global treatment of the complex
sheets.
\item \emph{Complex branches.}  Replacing $\log y$ by $\operatorname{Log}y+2\pi\ii m$ and
$\Wlam_0$ by $\Wlam_k$ produces, in every sector where the forward expansion
holds and the branch choices are continued consistently, a formal inverse
sheet with the \emph{same} coefficient functions.  That is a local asymptotic
chart, not a classification: distinct pairs $(m,k)$ can describe continuations
that meet, and the critical point of $\Gamma$ creates branch geometry the
leading Lambert model does not see.  No global statement is made.
\item \emph{Formalization.}  The triangular recurrence is well suited to a
proof assistant, since every requested coefficient depends on finitely many
earlier ones; the natural decomposition separates the algebra of well-ordered
supports, the formal implicit-function theorem
(\cref{p6:thm:triangular}), the exact rational certificates, and the analytic
import of the Stirling remainder bounds.  Nothing here is formalized.
\end{enumerate}

%% ==================== fragment p7 ====================
\chapter{The hyperfactorial \texorpdfstring{$K$}{K}-function}
\label{p7:sec:top}

\section{What this chapter does}
\label{p7:sec:intro}

The $K$-function interpolates the hyperfactorial,
$K(n+1)=\prod_{j=1}^{n}j^{j}$, and
$\log K(x)\sim\tfrac12x^{2}\log x-\tfrac14x^{2}$.  Its dominant block is
therefore of exactly the type inverted in \cref{p6:lem:core}, with the same
exponent $p=2$ as Barnes $G$ and a different shift; and the whole apparatus of
\cref{p6:sec:triangular} applies unchanged.  What makes the subject worth a
chapter of its own is a single feature: after the correct centring the entire
perturbation of power gap two collapses to one term,
\begin{equation}\label{p7:eq:the-obstruction}
 -\tfrac1{24}\log q,
\end{equation}
and \emph{that} term is resonant --- its logarithm is of the same size as the
core slope, so it does not shrink relative to the core the way an ordinary
correction does.  The three source articles deal with it in three different
ways, and this chapter's main contribution is to prove that the three are one
construction seen from three sides.

\begin{sourcebox}
The three filed articles are listed in \cref{p7:tab:sources}.  As in
\cref{p6:sec:top}, they agree: every coefficient and constant any of them
prints was recomputed for this volume in exact arithmetic and reproduced,
including the eleven-digit numerators of the refined block $b_5$.  Their
difference is one of strategy, not of arithmetic.

\begin{center}\small
\begin{tabular}{@{}lll@{}}
\toprule
Tag & Directory & Strategy for \eqref{p7:eq:the-obstruction}\\
\midrule
T1 & \texttt{inverse\_\allowbreak k\_\allowbreak function\_\allowbreak transseries} &
 leave it in the perturbation, \emph{and} absorb it\\
 & & into an algebraic shift $T=r^{2}-\tfrac1{12}$\\
T2 & \texttt{K\_\allowbreak Function\_\allowbreak Inverse\_\allowbreak Transseries} &
 absorb it into the core, using the $r$-Lambert\\
 & & function; also the arbitrary-centre expansion\\
T3 & \texttt{K\_\allowbreak function\_\allowbreak inverse\_\allowbreak transseries\_\allowbreak article} &
 leave it in the perturbation; slow-depth and\\
 & & ordered-block formulas, sectorial version\\
\bottomrule
\end{tabular}
\end{center}
\end{sourcebox}

\begin{table}[htbp]
\centering
\caption{The $K$-function in the notation of \cref{p6:thm:general}.}
\label{p7:tab:sources}
\small
\begin{tabular}{@{}ll@{}}
\toprule
inverted map & $x\mapsto K(x)$ on $[2,\infty)$\\
centred variable & $r=x-\tfrac12$\\
core $a x^{p}(\log x-\beta)$ & $a=\tfrac12$, $p=2$, $\beta=\tfrac12$\\
absorbed constant & $C=\log A-\tfrac18$, $A$ the Glaisher--Kinkelin constant\\
resonant term & $-\tfrac1{24}\log r$\\
perturbation support & $2\Nat$\\
core & $\rho=\e^{L}$, $L=\tfrac12(w+1)$, $w=\Wlam_0(4Y/\e)$\\
inverse support & $\rho^{-1},\rho^{-3},\rho^{-5},\dots$\\
\bottomrule
\end{tabular}
\end{table}

Compare \cref{p6:tab:sources}.  Barnes $G$ has $p=2$, $\beta=\tfrac32$, a
linear term $\alpha z$ generating the full integer support, and a resonant
term $-\tfrac1{12}\log z$.  The $K$-function has $p=2$, $\beta=\tfrac12$,
\emph{no} linear term --- so the support stays even, as for $\Gamma$ --- and a
resonant term $-\tfrac1{24}\log r$.  It is, in a precise sense made in
\cref{p7:sec:comparison}, the cleanest member of the family: the only
obstruction it has is the resonance.

\section{The function and its increasing branch}
\label{p7:sec:function}

\begin{lemma}[Exact logarithmic derivatives]\label{p7:lem:derivatives}
Write $\kappa(x)=\log K(x)$.  For $x>0$,
\begin{equation}\label{p7:eq:kappa-prime}
 \kappa'(x)=\log\Gamma(x)+x-\tfrac12-\tfrac12\log(2\pi),
 \qquad
 \kappa''(x)=\psi(x)+1 .
\end{equation}
\end{lemma}

\begin{proof}
Use $K(z)=\exp\bigl(\zeta'(-1,z)-\zeta'(-1)\bigr)$ with the Hurwitz zeta
function.  Differentiating $\partial_z\zeta(s,z)=-s\zeta(s+1,z)$ with respect
to $s$ gives $\partial_z\zeta'(s,z)=-\zeta(s+1,z)-s\zeta'(s+1,z)$, and at
$s=-1$,
$\partial_z\zeta'(-1,z)=-\zeta(0,z)+\zeta'(0,z)$.  Substituting
$\zeta(0,z)=\tfrac12-z$ and
$\zeta'(0,z)=\log\Gamma(z)-\tfrac12\log(2\pi)$ gives the first identity, and
differentiating it again gives the second, since
$(\log\Gamma)'=\psi$ and the derivative of $x-\tfrac12$ is $1$.
\end{proof}

\begin{proposition}[The increasing branch]\label{p7:prop:branch}
$K$ is strictly increasing on $[2,\infty)$, with $K(2)=1$ and
$K(x)\to+\infty$; write $K_+^{-1}:[1,\infty)\to[2,\infty)$ for the inverse.  On
$(0,\infty)$, $K$ has exactly two stationary points,
\begin{equation}\label{p7:eq:stationary}
 x_{\max}=0.29095698669918\ldots,
 \qquad
 x_{\min}=1.53768886373648\ldots,
\end{equation}
with $K(x_{\max})=1.29867230693177\ldots$ and
$K(x_{\min})=0.87978684305939\ldots$; so below $K(x_{\max})$ there are further
real branches, none of which survives as $y\to\infty$.
\end{proposition}

\begin{proof}
By \cref{p7:lem:derivatives}, $\kappa''=\psi+1$, and $\psi$ is strictly
increasing with $\psi(x)\to-\infty$ as $x\to0^{+}$ and $\psi(x)\to+\infty$;
so $\kappa''$ has exactly one zero and $\kappa'$ is strictly convex with a
unique minimum.  Since $\psi(2)=1-\gamma>0$ we get $\kappa''>0$ on
$[2,\infty)$, and
$\kappa'(2)=\log\Gamma(2)+2-\tfrac12-\tfrac12\log(2\pi)
=\tfrac32-\tfrac12\log(2\pi)=0.5811\ldots>0$; hence $\kappa'>0$ on
$[2,\infty)$ and $K$ increases there.  $K(2)=1$ is
\eqref{p7:eq:integer-values} at $n=1$, and $K\to\infty$ follows from
\cref{p7:thm:forward}.  Because $\kappa'$ has a unique minimum and tends to
$+\infty$ at both ends of $(0,\infty)$ while $\kappa'(1)=1-\tfrac12
-\tfrac12\log(2\pi)<0$, it has exactly two zeros; these are
\eqref{p7:eq:stationary}, located numerically.
\end{proof}

\begin{remark}[The notation is not cosmetic]\label{p7:rem:branches}
$K$ is not monotone on $(0,\infty)$, so ``the inverse of $K$'' is ambiguous
below $K(x_{\max})$.  Everything in this chapter concerns $K_+^{-1}$, the branch
tending to $+\infty$.  The functional equation $K(z+1)=z^{z}K(z)$ with
$K(1)=1$ gives
\begin{equation}\label{p7:eq:integer-values}
 K(n+1)=\prod_{j=1}^{n}j^{j},
\end{equation}
so $K_+^{-1}$ restricted to the values $\prod_{j\le n}j^{j}$ inverts the
hyperfactorial exactly.
\end{remark}

\section{Forward expansion, and why the centre is forced}
\label{p7:sec:forward}

The following puts every centring on the same footing; the distinguished one
then identifies itself.

\begin{theorem}[Expansion about an arbitrary centre]\label{p7:thm:forward}
Fix $\alpha\in\Cplx$ and let $x=r+\alpha$.  As $r\to\infty$ in a closed
subsector of $|\arg r|<\pi$,
\begin{equation}\label{p7:eq:forward-general}
 \log K(r+\alpha)
 \sim\frac{r^{2}}2\log r-\frac{r^{2}}4
 +\Bigl(\alpha-\frac12\Bigr)r\log r
 +\frac{B_2(\alpha)}2\log r
 +\log A+\frac{\alpha^{2}-\alpha}2
 +\sum_{n\ge1}\frac{(-1)^{n+1}B_{n+2}(\alpha)}{n(n+1)(n+2)\,r^{n}} .
\end{equation}
Truncation after $r^{-N}$ leaves a remainder $\bigO(r^{-N-1})$.
\end{theorem}

\begin{proof}
Start from the expansion at $\alpha=0$, which follows by integrating
\eqref{p7:eq:kappa-prime}: substituting Stirling's series into
$\kappa'(x)=\log\Gamma(x)+x-\tfrac12-\tfrac12\log(2\pi)$ gives
\begin{equation}\label{p7:eq:kappa-prime-stirling}
 \kappa'(x)\sim\Bigl(x-\frac12\Bigr)\log x-\frac12
 +\sum_{m\ge1}\frac{B_{2m}}{2m(2m-1)x^{2m-1}},
\end{equation}
and termwise integration gives
\begin{equation}\label{p7:eq:forward-zero}
 \log K(x)\sim\Bigl(\frac{x^{2}}2-\frac x2+\frac1{12}\Bigr)\log x
 -\frac{x^{2}}4+\log A
 -\sum_{m\ge1}\frac{B_{2m+2}}{2m(2m+1)(2m+2)x^{2m}} ,
\end{equation}
the integration constant being $\log A$ by the classical hyperfactorial
asymptotic.

Now substitute $x=r+\alpha$ and expand
$\log(r+\alpha)=\log r+\sum_{j\ge1}(-1)^{j+1}\alpha^{j}/(jr^{j})$ and
$(r+\alpha)^{-2m}=r^{-2m}\sum_{j\ge0}(-1)^{j}\binom{2m+j-1}{j}\alpha^{j}r^{-j}$.
At each fixed inverse power only finitely many terms contribute.  The
polynomial and logarithmic part collects into the first line, using
$B_2(\alpha)=\alpha^{2}-\alpha+\tfrac16$ for the coefficient of $\log r$: the
three sources of $\log r$ are $\tfrac12r^{2}\cdot0$, the term
$(\tfrac{\alpha^{2}}2-\tfrac\alpha2+\tfrac1{12})\log r$ from the quadratic
prefactor, which is $\tfrac12B_2(\alpha)\log r$.  For $n\ge1$ the coefficient
of $r^{-n}$ is a binomial convolution of Bernoulli numbers against powers of
$\alpha$, and the identity
$B_{n+2}(\alpha)=\sum_j\binom{n+2}{j}B_j\alpha^{n+2-j}$ collapses it to
$(-1)^{n+1}B_{n+2}(\alpha)/\bigl(n(n+1)(n+2)\bigr)$.  Translation by a fixed
$\alpha$ preserves closed proper subsectors, so the Poincar\'e remainder
survives.
\end{proof}

\begin{corollary}[The midpoint is the unique good centre]
\label{p7:cor:midpoint}
The coefficient of $r\log r$ in \eqref{p7:eq:forward-general} vanishes if and
only if $\alpha=\tfrac12$; and at that same value $B_{n+2}(\tfrac12)=0$ for
every odd $n$, so the tail becomes even.  With $r=x-\tfrac12$ and
$C=\log A-\tfrac18$,
\begin{equation}\label{p7:eq:centred}
 \kappa\Bigl(r+\frac12\Bigr)
 \sim\Phi(r)-\frac1{24}\log r+C+\sum_{m\ge1}\frac{c_m}{r^{2m}},
 \qquad
 \Phi(r)=\frac{r^{2}}2\log r-\frac{r^{2}}4,
\end{equation}
where
\begin{equation}\label{p7:eq:cm}
 c_m=-\frac{B_{2m+2}(1/2)}{2m(2m+1)(2m+2)}
    =-\frac{(2^{-2m-1}-1)B_{2m+2}}{2m(2m+1)(2m+2)} .
\end{equation}
The first values are
\begin{equation}\label{p7:eq:cm-values}
 c_1=-\frac7{5760},\quad
 c_2=\frac{31}{161280},\quad
 c_3=-\frac{127}{1290240},\quad
 c_4=\frac{511}{4866048},\quad
 c_5=-\frac{1414477}{7380172800},\quad
 c_6=\frac{8191}{15335424}.
\end{equation}
\end{corollary}

\begin{proof}
The coefficient of $r\log r$ is $\alpha-\tfrac12$.  At $\alpha=\tfrac12$ the
constant is $\log A+(\tfrac14-\tfrac12)/2=\log A-\tfrac18=C$ and the
coefficient of $\log r$ is $\tfrac12B_2(\tfrac12)=\tfrac12(-\tfrac1{12})
=-\tfrac1{24}$.  Odd-index Bernoulli polynomials vanish at $\tfrac12$ by
$B_n(1-\alpha)=(-1)^{n}B_n(\alpha)$, so only even $n=2m$ survives, and
$B_{2m+2}(\tfrac12)=(2^{-2m-1}-1)B_{2m+2}$ gives \eqref{p7:eq:cm}.
\end{proof}

\begin{remark}[Two coincidences that are one]\label{p7:rem:centre}
That the shift removing $r\log r$ is also the shift killing the odd tail looks
like luck.  It is not: both are the reflection symmetry
$B_n(1-\alpha)=(-1)^{n}B_n(\alpha)$ acting at its fixed point
$\alpha=\tfrac12$.  The same remark applies to $\Gamma$ in
\cref{p6:rem:why-half-shift}, where the half-shift does the same two jobs at
once.  \Cref{p7:thm:forward} makes this visible by displaying every centre at
the same time, instead of verifying the good one after the fact; this is the
main reason it is stated in that generality.
\end{remark}

\section{The Lambert anchor and the canonical inverse}
\label{p7:sec:canonical}

\begin{proposition}[Core variables]\label{p7:prop:anchor}
Let $Y=\log y$ and
\begin{equation}\label{p7:eq:core-vars}
 w=\Wlam_0\!\left(\frac{4Y}{\e}\right),
 \qquad
 \rho=\e^{(w+1)/2}=2\sqrt{\frac Yw},
 \qquad
 L=\log\rho=\frac{w+1}2 .
\end{equation}
Then $\Phi(\rho)=Y$ exactly, and $\Phi'(\rho)=\rho\log\rho=\rho L$.
\end{proposition}

\begin{proof}
This is \cref{p6:lem:core} with $a=\tfrac12$, $p=2$, $\beta=\tfrac12$,
$r=1$: there $v=\Wlam_k\bigl((pY/a)\e^{-p\beta}\bigr)=\Wlam_0(4Y/\e)$ and
$X=\e^{\beta+v/p}=\e^{(1+w)/2}$, and the slope is
$aX^{p-1}(v+1)=\tfrac12\rho(w+1)=\rho L$.
\end{proof}

Write $q=\rho^{-2}$, $r=\rho(1+v)$, and
\begin{equation}\label{p7:eq:kernel}
 \mathcal{D}_L(v)
 :=\frac12(1+v)^{2}\bigl(L+\log(1+v)\bigr)-\frac14(1+v)^{2}
   -\frac L2+\frac14
 =Lv+\frac{L+1}2v^{2}+\sum_{k\ge3}\frac{(-1)^{k+1}}{k(k-1)(k-2)}v^{k},
\end{equation}
which is $\mathcal{K}_{2,v}$ of \eqref{p6:eq:kernel-p2} rescaled; its linear
coefficient is $L$ because the normalization here is by $\rho^{2}$ rather
than $2\rho^{2}$.  Subtracting $\Phi(\rho)=Y$ from \eqref{p7:eq:centred} and
dividing by $\rho^{2}$ gives the exact normalized equation
\begin{equation}\label{p7:eq:canonical-implicit}
 \mathcal{D}_L(v)
 +q\left[E-\frac1{24}\log(1+v)
 +\sum_{m\ge1}c_mq^{m}(1+v)^{-2m}\right]=0,
 \qquad
 E:=C-\frac L{24}.
\end{equation}

\begin{theorem}[Canonical inverse transseries]\label{p7:thm:canonical}
\Cref{p7:eq:canonical-implicit} has a unique solution
$v=\sum_{n\ge1}V_n(L)q^{n}$ with $V_n\in\Rat[C,L,L^{-1}]$, given by
\begin{equation}\label{p7:eq:V-recurrence}
 V_n(L)=-\frac1L\Coef{q^{n}}\left\{\mathcal{D}_L(v_{<n})
 +q\left[E-\frac1{24}\log(1+v_{<n})
 +\sum_{m\ge1}c_mq^{m}(1+v_{<n})^{-2m}\right]\right\},
\end{equation}
$v_{<n}=\sum_{j<n}V_jq^{j}$; only $c_1,\dots,c_{n-1}$ can occur in $V_n$.
For every fixed $N\ge1$, as $y\to+\infty$,
\begin{equation}\label{p7:eq:canonical-main}
 K_+^{-1}(y)=\frac12+\rho+\sum_{n=1}^{N-1}V_n(L)\rho^{1-2n}
 +\bigO\bigl(\rho^{1-2N}\bigr).
\end{equation}
The first coefficients are
\begin{align}
 V_1&=-\frac EL=\frac1{24}-\frac CL,\label{p7:eq:V1}\\
 V_2&=-\frac1L\left(\frac{L+1}2V_1^{2}-\frac{V_1}{24}+c_1\right)
      =\frac{7L^{2}-240EL-2880(L+1)E^{2}}{5760L^{3}},\label{p7:eq:V2}\\
 V_3&=-\frac1L\left((L+1)V_1V_2+\frac{V_1^{3}}6
      -\frac1{24}\Bigl(V_2-\frac{V_1^{2}}2\Bigr)-2c_1V_1+c_2\right),
      \label{p7:eq:V3}\\
 V_4&=-\frac1L\Bigl[(L+1)\Bigl(V_1V_3+\frac{V_2^{2}}2\Bigr)
      +\frac{V_1^{2}V_2}2-\frac{V_1^{4}}{24}
      -\frac1{24}\Bigl(V_3-V_1V_2+\frac{V_1^{3}}3\Bigr)\notag\\
 &\hspace{4em}+c_1\bigl(-2V_2+3V_1^{2}\bigr)-4c_2V_1+c_3\Bigr].
      \label{p7:eq:V4}
\end{align}
\end{theorem}

\begin{proof}
The formal part is \cref{p6:thm:triangular} with $S=2\Nat$ and linear
coefficient $L$: by \eqref{p7:eq:kernel} the coefficient of $v$ in
$\mathcal{D}_L$ is $L$, which is invertible in $\Rat[C,L,L^{-1}]$, and every
monomial of the bracket in \eqref{p7:eq:canonical-implicit} carries at least
one factor $q$.  Hence a unique solution and \eqref{p7:eq:V-recurrence}.  The
displayed $V_1,\dots,V_4$ follow by running it; they were recomputed
independently (\cref{p7:rem:verification}).

For the analytic statement, \eqref{p7:eq:centred} is a Poincar\'e expansion
with differentiable remainder on the positive ray, obtained by integrating the
sectorial Stirling remainder in \eqref{p7:eq:kappa-prime-stirling}; retain it
far enough that its own error is below the first omitted inverse term.  The
formal substitution kills the normalized equation through $q^{N-1}$, so the
forward residual is $\bigO(\rho^{2}q^{N})=\bigO(\rho^{2-2N})$.  By
\cref{p7:lem:derivatives} and \eqref{p7:eq:kappa-prime-stirling},
$\kappa'(r+\tfrac12)=r\log r+\bigO(r^{-1})\asymp\rho L$ on the relevant
interval, and this is bounded below by $\tfrac12\rho L>0$ for large $y$.
\Cref{p0:thm:backward-error} converts the residual into an error
$\bigO(\rho^{2-2N}/(\rho L))=\bigO(\rho^{1-2N}/L)$, which is
\eqref{p7:eq:canonical-main}, and certifies that a root lies within it.
\end{proof}

\begin{remark}[Independent verification]\label{p7:rem:verification}
As in \cref{p6:rem:verification}, every rational number in this chapter was
recomputed for this volume from the defining equations, in exact arithmetic,
by an implementation independent of the sources: the tail $c_m$ from Bernoulli
polynomials at $\tfrac12$ through $m=6$; $V_1,\dots,V_4$ from
\eqref{p7:eq:V-recurrence}, including the closed form
\eqref{p7:eq:V2} for $V_2$; the accelerated forward blocks $d_1,\dots,d_6$
both from the closed formula \eqref{p7:eq:dn} and, independently, by direct
expansion of the identity in the proof of \cref{p7:thm:T-normal-form}; the
accelerated inverse blocks $U_2,\dots,U_5$; and the refined blocks
$b_1,\dots,b_5$ together with the limits \eqref{p7:eq:bn-limits}.  All agreed
with all three sources.  The numerical tables of
\cref{p7:sec:numerics} were computed afresh at $60$ digits from
$K(x)=\exp\bigl(\zeta'(-1,x)-\zeta'(-1)\bigr)$, and reproduce
$K(5)=27648=1^{1}2^{2}3^{3}4^{4}$ exactly.
\end{remark}

\section{The resonance and its three resolutions}
\label{p7:sec:resonance}

The term $-\tfrac1{24}\log r$ enters \eqref{p7:eq:canonical-implicit} at
$q^{1}$ carrying the factor $E=C-L/24$, which \emph{grows} with the slope.  It
is therefore resonant in the sense of \cref{p6:prop:resonance}: it survives
the large-slope limit, while every $c_m$ and the constant $C$ do not.  The
consequence is visible in $V_1=\tfrac1{24}-C/L$, whose limit is not zero.

\subsection{First resolution: leave it, and take the limit}

\begin{proposition}[Coefficientwise limit]\label{p7:prop:limit}
For each fixed $n$,
\begin{equation}\label{p7:eq:V-limits}
 \lim_{L\to\infty}V_n(L)=\binom{1/2}{n}12^{-n},
\end{equation}
so $V_1\to\tfrac1{24}$, $V_2\to-\tfrac1{1152}$, $V_3\to\tfrac1{27648}$.  The
limiting coefficients are those of
\begin{equation}\label{p7:eq:limit-series}
 v_0(q)=\sqrt{1+\tfrac q{12}}-1,
 \qquad\text{equivalently}\qquad
 \rho\bigl(1+v_0(\rho^{-2})\bigr)=\sqrt{\rho^{2}+\tfrac1{12}} .
\end{equation}
\end{proposition}

\begin{proof}
Divide \eqref{p7:eq:canonical-implicit} by $L$ and let $L\to\infty$ with $q$
formal.  By \eqref{p7:eq:kernel},
$\mathcal{D}_L(v)=L\cdot\tfrac12\bigl((1+v)^{2}-1\bigr)+\bigl(\text{terms free
of }L\bigr)$, so $\mathcal{D}_L(v)/L\to\tfrac12\bigl((1+v)^{2}-1\bigr)$.  In
the bracket, $qE/L=q(C/L-\tfrac1{24})\to-q/24$, while
$q\log(1+v)/(24L)\to0$ and every $c_m$ term $\to0$.  The limiting equation is
therefore
\[
 \tfrac12\bigl((1+v_0)^{2}-1\bigr)-\tfrac q{24}=0,
 \qquad\text{i.e.}\qquad (1+v_0)^{2}=1+\tfrac q{12},
\]
whose branch with $v_0(0)=0$ is \eqref{p7:eq:limit-series}.  The limit may be
taken inside the recursion because at each order the coefficient equation has
limiting divisor $1$ and involves only finitely many earlier coefficients.
Expanding the square root gives \eqref{p7:eq:V-limits}, and $r=\rho(1+v)$
turns $v_0$ into $\sqrt{\rho^{2}+1/12}$.
\end{proof}

\subsection{Second resolution: absorb it into an algebraic shift}

\begin{lemma}[Absorbing a logarithmic correction]\label{p7:lem:absorb}
Let $a\ne0$, $p\ne0$, and suppose
\begin{equation}\label{p7:eq:absorb-hyp}
 F(r)=a r^{p}(\log r-b)+\beta\log r+\gamma+\bigO(r^{-p}).
\end{equation}
Put $T=r^{p}+\alpha$ and $\Psi(T)=\tfrac ap T(\log T-pb)$.  Then
\begin{equation}\label{p7:eq:absorb}
 \Psi(r^{p}+\alpha)
 =a r^{p}(\log r-b)+a\alpha\log r
 +a\alpha\Bigl(\frac1p-b\Bigr)+\bigO(r^{-p}),
\end{equation}
so the unique shift absorbing the term $\beta\log r$ is
\begin{equation}\label{p7:eq:alpha}
 \alpha=\frac\beta a .
\end{equation}
\end{lemma}

\begin{proof}
Insert $T=r^{p}+\alpha$ and use
$\log T=p\log r+\alpha r^{-p}+\bigO(r^{-2p})$:
\[
 \Psi(T)=\frac ap(r^{p}+\alpha)\bigl(p\log r+\alpha r^{-p}-pb\bigr)
 +\bigO(r^{-p})
 =a r^{p}\log r+\frac{a\alpha}p-abr^{p}+a\alpha\log r-ab\alpha+\bigO(r^{-p}),
\]
which is \eqref{p7:eq:absorb}.  Matching the coefficient of $\log r$ gives
$a\alpha=\beta$.
\end{proof}

For the $K$-function $a=\tfrac12$, $p=2$, $b=\tfrac12$ and
$\beta=-\tfrac1{24}$, so $\alpha=-\tfrac1{12}$ and, remarkably, the constant
correction $a\alpha(1/p-b)$ vanishes because $p^{-1}=b$.  The new coordinate is
\begin{equation}\label{p7:eq:T-coordinate}
 T=r^{2}-\frac1{12}=x^{2}-x+\frac16,
 \qquad
 x=\frac12+\sqrt{T+\frac1{12}} .
\end{equation}

\begin{theorem}[Quadratic normal form]\label{p7:thm:T-normal-form}
With $\Psi(T)=\tfrac14T(\log T-1)$ and $T=x^{2}-x+\tfrac16$, as
$T\to\infty$,
\begin{equation}\label{p7:eq:T-normal}
 \log K(x)\sim\Psi(T)+C+\sum_{n\ge1}\frac{d_n}{T^{n}},
\end{equation}
where, writing $a_0=\tfrac1{12}$,
\begin{equation}\label{p7:eq:dn}
 d_n=\frac{(-1)^{n}a_0^{\,n+1}}{4(n+1)}
 +\sum_{m=1}^{n}(-1)^{n-m}\binom{n-1}{m-1}a_0^{\,n-m}c_m .
\end{equation}
The first values are
\begin{equation}\label{p7:eq:dn-values}
 d_1=-\frac1{480},\quad
 d_2=\frac{31}{90720},\quad
 d_3=-\frac{103}{725760},\quad
 d_4=\frac{179}{1330560},\quad
 d_5=-\frac{6480899}{28021593600},\quad
 d_6=\frac{3485299}{5604318720}.
\end{equation}
\end{theorem}

\begin{proof}
Since $r^{2}=T+a_0$ and $\tfrac12r^{2}-\tfrac1{24}=\tfrac T2$, the elementary
part of \eqref{p7:eq:centred} satisfies the \emph{exact} identity
\[
 \Bigl(\frac{r^{2}}2-\frac1{24}\Bigr)\log r-\frac{r^{2}}4
 =\frac T4\log(T+a_0)-\frac T4-\frac{a_0}4
 =\Psi(T)+\frac T4\log\Bigl(1+\frac{a_0}T\Bigr)-\frac{a_0}4 ,
\]
using $\log r=\tfrac12\log(T+a_0)$.  The first term of
$\tfrac T4\log(1+a_0/T)$ is $a_0/4$, which cancels the constant $-a_0/4$
exactly; this is why no constant is generated and $C$ is unchanged.  Expanding
the remaining logarithm and $(T+a_0)^{-m}=T^{-m}(1+a_0/T)^{-m}$ and collecting
the coefficient of $T^{-n}$ gives \eqref{p7:eq:dn}.
\end{proof}

The gain is one full sector.  In \eqref{p7:eq:centred} the largest unresolved
term was $-\tfrac1{24}\log r$, of size $\log r$; in \eqref{p7:eq:T-normal} it
is $d_1/T$.  Inverting exactly as before, with
\begin{equation}\label{p7:eq:acc-vars}
 \eta=Y-C,
 \qquad
 \omega=\Wlam_0\!\left(\frac{4\eta}{\e}\right),
 \qquad
 \tau=\frac{4\eta}{\omega}=\e^{\omega+1},
 \qquad
 S=\log\tau=\omega+1,
\end{equation}
so that $\Psi(\tau)=\eta$ exactly, and writing $T=\tau(1+u)$, $\nu=\tau^{-1}$,
the normalized equation becomes
\begin{equation}\label{p7:eq:acc-implicit}
 \mathcal{E}_S(u)+4\sum_{m\ge1}d_m\nu^{m+1}(1+u)^{-m}=0,
 \qquad
 \mathcal{E}_S(u)=Su+\sum_{k\ge2}\frac{(-1)^{k}}{k(k-1)}u^{k} .
\end{equation}

\begin{theorem}[Accelerated inverse]\label{p7:thm:accelerated}
\Cref{p7:eq:acc-implicit} has a unique solution
$u=\sum_{n\ge2}U_n(S)\nu^{n}$ in $\nu^{2}\Rat[S,S^{-1}][[\nu]]$, with
\begin{equation}\label{p7:eq:U-recurrence}
 U_n(S)=-\frac1S\Coef{\nu^{n}}
 \left\{\mathcal{E}_S(u_{<n})+4\sum_{m\ge1}d_m\nu^{m+1}(1+u_{<n})^{-m}\right\},
\end{equation}
and
\begin{equation}\label{p7:eq:U-values}
 U_2=\frac1{120S},\quad
 U_3=-\frac{31}{22680S},\quad
 U_4=\frac{1030S^{2}-126S-63}{1814400S^{3}},\quad
 U_5=-\frac{16110S^{2}-1023S-341}{29937600S^{3}} .
\end{equation}
Consequently, as $y\to+\infty$,
\begin{equation}\label{p7:eq:accelerated-main}
 K_+^{-1}(y)\sim\frac12+\sqrt{\frac1{12}
 +\tau\Bigl(1+\sum_{n\ge2}U_n(S)\tau^{-n}\Bigr)},
\end{equation}
and truncating after $U_N\tau^{-N}$ leaves an error $\bigO(\tau^{-N-1/2})$ in
$x$.  Expanding the root,
\begin{equation}\label{p7:eq:accelerated-expanded}
 K_+^{-1}(y)=\frac12+R+\frac1{240S\tau R}-\frac{31}{45360S\tau^{2}R}
 +\bigO(\tau^{-7/2}),
 \qquad R=\sqrt{\tau+\tfrac1{12}} .
\end{equation}
\end{theorem}

\begin{proof}
$\mathcal{E}_S$ has linear coefficient $S$, a unit, and the forcing begins at
$\nu^{2}$ because $d_m\nu^{m+1}$ starts at $m=1$; \cref{p6:thm:triangular}
with $S=\Nat$ gives uniqueness and \eqref{p7:eq:U-recurrence}, and
$u=\bigO(\nu^{2})$.  The values \eqref{p7:eq:U-values} were computed and
independently verified (\cref{p7:rem:verification}).  The analytic step is
\cref{p0:thm:backward-error} exactly as in \cref{p7:thm:canonical}, with slope
$\Psi'(\tau)=\tfrac14(\log\tau)=\tfrac S4$ in the $T$ variable; converting to
$x$ through $x=\tfrac12+\sqrt{T+1/12}$ costs the factor $\tfrac12T^{-1/2}$,
which is the source of the half-integer exponent.
\end{proof}

\subsection{The first two resolutions are the same}

\begin{theorem}[Acceleration is resummation]\label{p7:thm:acceleration-is-resummation}
The coefficientwise limit of \cref{p7:prop:limit} and the algebraic
substitution of \cref{p7:lem:absorb} are the same construction:
\begin{equation}\label{p7:eq:same}
 \frac12+\rho+\sum_{n\ge1}\Bigl(\lim_{L\to\infty}V_n(L)\Bigr)\rho^{1-2n}
 =\frac12+\sqrt{\rho^{2}+\frac1{12}}
 =\frac12+\sqrt{T+\frac1{12}}\Big|_{T=\rho^{2}-1/12}^{\phantom{.}}
 ,
\end{equation}
and the map $T\mapsto\tfrac12+\sqrt{T+\tfrac1{12}}$ appearing on the right is
precisely the inverse of the accelerated coordinate
\eqref{p7:eq:T-coordinate}.  In other words, the leading term of the
accelerated expansion \eqref{p7:eq:accelerated-main} \emph{is} the sum of the
resonant subsector of the canonical expansion \eqref{p7:eq:canonical-main},
and the residual series $\sum U_n\tau^{-n}$ is what remains of the canonical
coefficients after that subsector is removed.
\end{theorem}

\begin{proof}
By \cref{p7:prop:limit} the limiting subsector is
$\sum_{n\ge1}\binom{1/2}{n}12^{-n}\rho^{1-2n}
=\rho\bigl(\sqrt{1+1/(12\rho^{2})}-1\bigr)=\sqrt{\rho^{2}+1/12}-\rho$, so
adding $\tfrac12+\rho$ gives the middle expression in \eqref{p7:eq:same}.  The
right-hand equality is \eqref{p7:eq:T-coordinate} read backwards.  For the
last sentence: \cref{p7:lem:absorb} removes from the forward expansion exactly
the term $\beta\log r$ that is resonant, and by
\cref{p6:prop:resonance} a perturbation is resonant precisely when it survives
division by the slope in the large-slope limit; the remaining forward terms
$d_n T^{-n}$ have no $\log T$ and are therefore non-resonant, so the
accelerated coefficients $U_n$ all tend to $0$ --- indeed
$U_n=\bigO(S^{-1})$ by \eqref{p7:eq:U-values} and the recurrence.
\end{proof}

\begin{remark}[The same statement for Barnes $G$]\label{p7:rem:barnes-too}
Neither Barnes source performs this absorption, and it costs nothing.  In
\cref{p6:lem:barnes-normal-form} the data are $a=\tfrac12$, $p=2$,
$b=\tfrac32$, $\beta=-\tfrac1{12}$, so \cref{p7:lem:absorb} gives
$\alpha=-\tfrac16$, the coordinate $T=z^{2}-\tfrac16$, and the normal form
\begin{equation}\label{p7:eq:barnes-accelerated}
 \log G(z+1)
 =\frac14T(\log T-3)+\alpha_{\!*}\sqrt{T+\tfrac16}
 +\Bigl(\zeta'(-1)-\frac1{12}\Bigr)+\bigO(T^{-1}),
 \qquad \alpha_{\!*}=\tfrac12\log(2\pi),
\end{equation}
the constant $-\tfrac1{12}$ coming from $a\alpha(1/p-b)$ in
\eqref{p7:eq:absorb}.  The square root $\sqrt{T+\tfrac16}$ is exactly the
resummed subsector $\sqrt{U^{2}+1/6}$ that
\cref{p6:prop:resonance} identifies as a coefficientwise limit --- so
\cref{p6:rem:limit-caveat}, which warns that the limit is not an
approximation, is answered here: it becomes one as soon as it is moved into
the coordinate.  The price for Barnes, and the reason its sources may have
balked, is the half-integer power $\sqrt{T+1/6}$ forced by the linear term
$\alpha_{\!*}z$, which the $K$-function does not have.  That is the precise
sense in which the $K$-function is the cleaner subject.
\end{remark}

\subsection{Third resolution: absorb it into the core}

The shift of \cref{p7:lem:absorb} absorbs $\beta\log r$ only modulo
$\bigO(r^{-p})$.  It can instead be absorbed \emph{exactly}, at the cost of
replacing $\Wlam$ by a one-parameter relative.

\begin{definition}[The $r$-Lambert function]\label{p7:def:rlambert}
For fixed $\varrho$, let $\Wlam^{(\varrho)}$ denote a branch of the inverse of
$s\mapsto s\e^{s}+\varrho s$, so that
\begin{equation}\label{p7:eq:rlambert}
 \Wlam^{(\varrho)}(z)\,\e^{\Wlam^{(\varrho)}(z)}
 +\varrho\,\Wlam^{(\varrho)}(z)=z .
\end{equation}
For $\varrho=0$ this is $\Wlam$.
\end{definition}

\begin{theorem}[Exact inversion of an affine logarithmic power core]
\label{p7:thm:affine-core}
Let
\begin{equation}\label{p7:eq:affine-core}
 \Psi_{\!*}(x)=a x^{p}\log x+b_{*}x^{p}+c_{*}\log x+d_{*},
 \qquad ap\ne0 .
\end{equation}
Fix compatible branches and set
\begin{equation}\label{p7:eq:affine-data}
 \varrho=\frac{c_{*}}a\exp\!\left(\frac{pb_{*}}a\right),
 \qquad
 z=\frac pa\exp\!\left(\frac{pb_{*}}a\right)
   \left(Y-d_{*}+\frac{c_{*}b_{*}}a\right).
\end{equation}
If $s=\Wlam^{(\varrho)}(z)$, then
\begin{equation}\label{p7:eq:affine-inverse}
 x=\exp\!\left(\frac sp-\frac{b_{*}}a\right)
\end{equation}
solves $\Psi_{\!*}(x)=Y$, and every local branch arises from compatible
choices of the logarithm and of a branch of $\Wlam^{(\varrho)}$.
\end{theorem}

\begin{proof}
Put $s=p\log x+pb_{*}/a$, so that $\log x=s/p-b_{*}/a$ and
$x^{p}=\exp(s-pb_{*}/a)$.  The two power terms collapse:
\[
 a x^{p}\log x+b_{*}x^{p}
 =x^{p}\bigl(a\log x+b_{*}\bigr)
 =\exp\!\left(s-\frac{pb_{*}}a\right)\cdot\frac{as}p
 =\frac ap\exp\!\left(-\frac{pb_{*}}a\right)s\,\e^{s},
\]
where the middle step uses $a\log x+b_{*}=a(s/p-b_{*}/a)+b_{*}=as/p$.  The
logarithmic term is $c_{*}(s/p-b_{*}/a)$.  Hence $\Psi_{\!*}(x)=Y$ reads
\[
 \frac ap\e^{-pb_{*}/a}s\e^{s}+\frac{c_{*}}ps
 =Y-d_{*}+\frac{c_{*}b_{*}}a ,
\]
and multiplying by $(p/a)\e^{pb_{*}/a}$ gives
$s\e^{s}+\varrho s=z$ with the data \eqref{p7:eq:affine-data}.
\end{proof}

\begin{corollary}[The refined $K$ anchor]\label{p7:cor:refined-anchor}
Apply \cref{p7:thm:affine-core} to
$\Psi_{\!*}(r)=\tfrac12r^{2}\log r-\tfrac14r^{2}-\tfrac1{24}\log r+C$,
that is to $a=\tfrac12$, $p=2$, $b_{*}=-\tfrac14$, $c_{*}=-\tfrac1{24}$,
$d_{*}=C$.  Then
\begin{equation}\label{p7:eq:K-refined-data}
 \varrho=-\frac1{12\e},
 \qquad
 z=\frac{48(Y-C)+1}{12\e},
\end{equation}
and the resulting $\rho_{*}=\exp\bigl(\tfrac12\Wlam^{(\varrho)}(z)+\tfrac12
\bigr)$ satisfies $\Psi_{\!*}(\rho_{*})=Y$ exactly, with slope
$\Psi_{\!*}'(\rho_{*})=\rho_{*}\log\rho_{*}-\tfrac1{24\rho_{*}}$.  The
residual perturbation is only the Bernoulli tail
$\sum_{m\ge1}c_m\rho_{*}^{-2m}$.
\end{corollary}

\begin{proof}
Substitute the data into \eqref{p7:eq:affine-data}:
$pb_{*}/a=-1$, so $\varrho=(-\tfrac1{24})\cdot2\cdot\e^{-1}=-1/(12\e)$ and
$z=4\e^{-1}\bigl(Y-C+(-\tfrac1{24})(-\tfrac14)\cdot2\bigr)
=4\e^{-1}\bigl(Y-C+\tfrac1{48}\bigr)$, which is
\eqref{p7:eq:K-refined-data}.  The formula for $\rho_{*}$ is
\eqref{p7:eq:affine-inverse} with $p=2$ and $b_{*}/a=-\tfrac12$; the slope is
immediate.
\end{proof}

\begin{theorem}[Refined blocks]\label{p7:thm:refined}
With $t=\rho_{*}^{-2}$, $L_{*}=\log\rho_{*}$, $r=\rho_{*}H$ and
$H=1+\sum_{n\ge1}b_n(L_{*})t^{\,n+1}$, the normalized equation is
\begin{equation}\label{p7:eq:refined-implicit}
 H^{2}\bigl(2L_{*}-1+2\log H\bigr)-(2L_{*}-1)
 -\frac t6\log H
 +4\sum_{m\ge1}c_mt^{\,m+1}H^{-2m}=0,
\end{equation}
its linearization at $H=1$, $t=0$ is $4L_{*}$, and
\begin{equation}\label{p7:eq:bn-recurrence}
 b_n(L)=-\frac1{4L}\Coef{t^{\,n+1}}\bigl(\text{left side of
 \eqref{p7:eq:refined-implicit} at }H=H^{*}_{n-1}\bigr),
 \qquad n\ge1 .
\end{equation}
Each $b_n$ is rational in $L$ with $b_n=\bigO(L^{-1})$, uses only
$c_1,\dots,c_n$, and satisfies
\begin{equation}\label{p7:eq:bn-limits}
 \lim_{L\to\infty}L\,b_n(L)=-c_n .
\end{equation}
For every fixed $N$,
\begin{equation}\label{p7:eq:refined-main}
 K_+^{-1}(y)=\frac12+\rho_{*}+\sum_{n=1}^{N}\frac{b_n(L_{*})}{\rho_{*}^{2n+1}}
 +\bigO\!\left(\frac{\rho_{*}^{-2N-3}}{L_{*}}\right).
\end{equation}
The fourth and fifth blocks are
\begin{align}
 b_4&=-\frac{12877200L^{3}-703428L^{2}+57552L+539}{122624409600\,L^{4}},
 \label{p7:eq:b4}\\
 b_5&=\frac{1539856241280L^{4}-45140228640L^{3}+1666058823L^{2}
 -38145107L-25162137}{8034351316992000\,L^{5}} .
 \label{p7:eq:b5}
\end{align}
\end{theorem}

\begin{proof}
The equation is \eqref{p7:eq:canonical-implicit} with the resonant term and
the constant moved into the core, so only the tail remains; the linearization
is $4L_{*}$ because the normalization is by $\tfrac14\rho_{*}^{2}$.
\Cref{p6:thm:triangular} gives uniqueness and
\eqref{p7:eq:bn-recurrence}.  Since the forcing begins at $t^{2}$, so does
$H-1$, which is why the ansatz uses $t^{\,n+1}$.  For
\eqref{p7:eq:bn-limits}: at order $t^{\,n+1}$ the only contribution
independent of earlier blocks is $4c_nt^{\,n+1}$, so
$4Lb_n=-4c_n+\bigO(L^{0}\cdot L^{-1})$, giving $Lb_n\to-c_n$; equivalently
this is the linear response $-S(\rho_{*})/\Psi_{\!*}'(\rho_{*})$.  The two
displayed blocks were computed and independently verified
(\cref{p7:rem:verification}).  The analytic clause is again
\cref{p0:thm:backward-error}, with the slope
$\Psi_{\!*}'(\rho_{*})\sim\rho_{*}L_{*}$.
\end{proof}

\begin{remark}[Which resolution to use]\label{p7:rem:which-resolution}
All three are correct and they answer different questions.  The canonical
expansion is the one to state, because its coefficients live in
$\Rat[C,L,L^{-1}]$ and it needs only $\Wlam_0$.  The accelerated coordinate is
the one to compute with: \cref{p7:tab:accelerated} shows its leading term
alone beating four canonical blocks at $x=5$.  The $r$-Lambert anchor is the
one that explains: it shows that the obstruction is not deep, merely
misplaced, since a single one-parameter special function removes it exactly
and leaves nothing but the Bernoulli tail --- which is the same tail that
$\Gamma$ and $G$ leave, and which is beyond repair because it is divergent.
\end{remark}

\section{Numerical verification}
\label{p7:sec:numerics}

The protocol is that of \cref{p6:sec:numerics}: choose an exact abscissa $x$,
evaluate $K(x)=\exp\bigl(\zeta'(-1,x)-\zeta'(-1)\bigr)$ to $60$ digits,
reconstruct the core variables from that value alone, and compare the
truncations with $x$.  Errors are signed.

\begin{table}[htbp]
\centering
\caption{Canonical expansion \eqref{p7:eq:canonical-main}: signed error of
$\tfrac12+\rho$ and of the truncations including $V_1,\dots,V_4$.}
\label{p7:tab:canonical}
\scriptsize
\setlength{\tabcolsep}{4pt}
\begin{tabular}{rccccc}
\toprule
$x$ & $\tfrac12+\rho$ & $+\,V_1$ & $+\,V_2$ & $+\,V_3$ & $+\,V_4$\\
\midrule
5 & $9.0012\!\times\!10^{-3}$ & $1.8437\!\times\!10^{-5}$ & $1.2665\!\times\!10^{-7}$ & $-5.9396\!\times\!10^{-10}$ & $6.4854\!\times\!10^{-11}$\\
10 & $1.3996\!\times\!10^{-3}$ & $-1.9257\!\times\!10^{-7}$ & $8.4998\!\times\!10^{-10}$ & $-5.1976\!\times\!10^{-12}$ & $6.4879\!\times\!10^{-14}$\\
20 & $-2.7600\!\times\!10^{-7}$ & $-5.5164\!\times\!10^{-8}$ & $2.0887\!\times\!10^{-11}$ & $-2.9637\!\times\!10^{-14}$ & $8.4762\!\times\!10^{-17}$\\
50 & $-2.0103\!\times\!10^{-4}$ & $-2.9310\!\times\!10^{-9}$ & $1.9068\!\times\!10^{-13}$ & $-3.7123\!\times\!10^{-17}$ & $1.5892\!\times\!10^{-20}$\\
100 & $-1.4839\!\times\!10^{-4}$ & $-2.6924\!\times\!10^{-10}$ & $5.4398\!\times\!10^{-15}$ & $-2.5162\!\times\!10^{-19}$ & $2.6077\!\times\!10^{-23}$\\
\bottomrule
\end{tabular}
\end{table}

\begin{table}[htbp]
\centering
\caption{Accelerated expansion \eqref{p7:eq:accelerated-main}: signed error of
$\tfrac12+\sqrt{\tfrac1{12}+\tau}$ and of the truncations including
$U_2,\dots,U_5$.}
\label{p7:tab:accelerated}
\scriptsize
\setlength{\tabcolsep}{4pt}
\begin{tabular}{rccccc}
\toprule
$x$ & $\tfrac12+\sqrt{\tfrac1{12}+\tau}$ & $+\,U_2$ & $+\,U_3$ & $+\,U_4$ & $+\,U_5$\\
\midrule
5 & $-1.5162\!\times\!10^{-5}$ & $1.2198\!\times\!10^{-7}$ & $-2.3300\!\times\!10^{-9}$ & $1.0850\!\times\!10^{-10}$ & $-9.0171\!\times\!10^{-12}$\\
10 & $-1.0786\!\times\!10^{-6}$ & $1.9569\!\times\!10^{-9}$ & $-8.6886\!\times\!10^{-12}$ & $9.1990\!\times\!10^{-14}$ & $-1.7528\!\times\!10^{-15}$\\
20 & $-9.4572\!\times\!10^{-8}$ & $4.0777\!\times\!10^{-11}$ & $-4.3490\!\times\!10^{-14}$ & $1.0944\!\times\!10^{-16}$ & $-4.9719\!\times\!10^{-19}$\\
50 & $-4.4020\!\times\!10^{-9}$ & $2.9465\!\times\!10^{-13}$ & $-4.9102\!\times\!10^{-17}$ & $1.9149\!\times\!10^{-20}$ & $-1.3502\!\times\!10^{-23}$\\
100 & $-4.5974\!\times\!10^{-10}$ & $7.6166\!\times\!10^{-15}$ & $-3.1503\!\times\!10^{-19}$ & $3.0371\!\times\!10^{-23}$ & $-5.2971\!\times\!10^{-27}$\\
\bottomrule
\end{tabular}
\end{table}

The comparison is the point.  At $x=5$ the accelerated \emph{leading} term
($1.5\times10^{-5}$) is already as accurate as the canonical expansion carried
through its first block ($1.8\times10^{-5}$), and the pattern persists down the
table: at every argument, accelerated column $j$ tracks canonical column
$j+1$.  This is \cref{p7:thm:acceleration-is-resummation} in numbers.  The
accelerated leading term is not a better approximation of the same kind --- it
is the canonical expansion with an entire infinite subsector already summed,
so it starts one block ahead and stays there.  The accelerated form is also
the cheaper of the two to evaluate, since $U_2,\dots,U_5$ are shorter
expressions than $V_1,\dots,V_4$ and involve no $C$.

\begin{remark}[The anomalous row]\label{p7:rem:anomaly}
The entry $-2.76\times10^{-7}$ at $x=20$ in \cref{p7:tab:canonical} is far
smaller than its neighbours at $x=10$ and $x=50$, and the column changes sign
there.  This is not an accident of the tabulation and not a numerical
artefact: by \eqref{p7:eq:V1} the first correction vanishes exactly when
\begin{equation}\label{p7:eq:V1-zero}
 V_1(L)=\frac1{24}-\frac CL=0,
 \qquad\text{i.e.}\qquad
 L=24C=2.97010\ldots,
\end{equation}
which corresponds to $\rho=\e^{24C}=19.4944\ldots$, that is to
$x=\tfrac12+\rho=19.994\ldots$.  The tabulated argument $x=20$ is within
$0.006$ of the zero of $V_1$, so at that one point the leading correction is
almost exactly absent and the uncorrected core is anomalously good.  The
effect is entirely local; it says nothing about the quality of the expansion
and everything about the numerical value of the Glaisher constant.
\end{remark}

\section{Comparison with \texorpdfstring{$\Gamma$}{Gamma} and
\texorpdfstring{$G$}{G}}
\label{p7:sec:comparison}

\begin{table}[htbp]
\centering
\caption{The three power--logarithmic subjects of
\cref{p6:sec:top,p7:sec:top} in the notation of \cref{p6:thm:general}.}
\label{p7:tab:three-subjects}
\small
\begin{tabular}{@{}lccc@{}}
\toprule
 & $\Gamma(t+\tfrac12)$ & $K(r+\tfrac12)$ & $G(z+1)$\\
\midrule
$a$ & $1$ & $\tfrac12$ & $\tfrac12$\\
$p$ & $1$ & $2$ & $2$\\
$\beta$ & $1$ & $\tfrac12$ & $\tfrac32$\\
linear term & --- & --- & $\tfrac12\log(2\pi)\,z$\\
resonant term & --- & $-\tfrac1{24}\log r$ & $-\tfrac1{12}\log z$\\
support & $2\Nat$ & $2\Nat$ & $\Nat$\\
tail & $a_k=\dfrac{B_{2k}(1/2)}{2k(2k-1)}$
     & $c_m=\dfrac{-B_{2m+2}(1/2)}{2m(2m+1)(2m+2)}$
     & $c_k=\dfrac{B_{2k+2}}{4k(k+1)}$\\[0.6em]
square root & $\sqrt{\Lambda^{2}T^{2}+\tfrac1{12}}$
            & $\sqrt{\rho^{2}+\tfrac1{12}}$
            & $\sqrt{U^{2}+\tfrac16}$\\
its source & deepest pole & resonance & resonance\\
\bottomrule
\end{tabular}
\end{table}

Three remarks are worth making from the merged view; none is in any of the six
sources.

\begin{remark}[Where the $\tfrac1{12}$ comes from, twice]
\label{p7:rem:same-twelfth}
The constant $\tfrac1{12}$ appears in the last row of
\cref{p7:tab:three-subjects} twice, for two apparently unrelated reasons: for
$\Gamma$ it is the deepest pole in the slope
(\cref{p6:thm:deepest-pole}, where it arises as $-2a_1$), and for $K$ it is
the resonant limit (\cref{p7:prop:limit}, where it arises as $-24^{-1}\cdot
(-2)$).  The coincidence is not one.  Both are $-2\cdot B_2(\tfrac12)/2
=\tfrac1{12}$: for $\Gamma$, $a_1=B_2(\tfrac12)/2=-\tfrac1{24}$ is the first
tail coefficient; for $K$, $\tfrac12B_2(\tfrac12)=-\tfrac1{24}$ is the
coefficient of the resonant logarithm produced by
\cref{p7:cor:midpoint}.  The same Bernoulli value at the same reflection
centre is doing both jobs, in one case as a tail coefficient and in the other
as a logarithmic coefficient.
\end{remark}

\begin{remark}[The square root is always the quadratic model]
\label{p7:rem:quadratic-again}
All three square roots in \cref{p7:tab:three-subjects}, and the fourth one
$\sqrt{\lambda^{2}U^{2}-2\alpha U}$ of \eqref{p6:eq:barnes-deepest-resum}, are
instances of \cref{p6:lem:quadratic-core}: truncate the normalized equation
to its linear and quadratic terms in the displacement together with the lowest
forcing monomial, and solve.  What differs between the cases is only which
\emph{extraction} of the coefficient family the truncation computes --- the
deepest pole in the slope, or the coefficientwise large-slope limit --- and
that in turn is decided by whether the forcing monomial carries a factor of
the slope, i.e.\ by whether it is resonant.  Resonance and pole depth are thus
two readings of one quadratic.
\end{remark}

\begin{remark}[What $\Gamma$ and $G$ are missing]\label{p7:rem:missing}
The $K$-function sources supply two devices the $\Gamma$/$G$ sources do not:
the absorption lemma \cref{p7:lem:absorb} and the $r$-Lambert anchor
\cref{p7:thm:affine-core}.  Both apply to $G$.  The absorption is carried out
in \cref{p7:rem:barnes-too}.  The $r$-Lambert anchor does \emph{not} apply to
$G$ as it stands, because \cref{p7:thm:affine-core} requires the core to be
$ax^{p}\log x+b_{*}x^{p}+c_{*}\log x+d_{*}$ and the Barnes linear term
$\alpha_{\!*}z$ is of degree $p-1$, outside that shape; absorbing it would
need a two-parameter relative of $\Wlam$ for which no closed theory is
invoked here.  Neither device applies to $\Gamma$, which has $p=1$ and no
resonant term at all --- for $\Gamma$ the logarithmic coefficient
$\tfrac12B_2(h)$ at $h=\tfrac12$ enters the \emph{tail}, not the phase.  So
the three subjects are not three copies of one calculation: they are three
different positions of a single obstruction, and the $K$-function is the one
where it sits in the phase and can be removed.
\end{remark}

\section{What is left open}
\label{p7:sec:open}

\begin{enumerate}
\item \emph{A two-parameter anchor.}  \Cref{p7:rem:missing} identifies the
gap: a core of the shape $ax^{p}\log x+b_{*}x^{p}+e_{*}x^{p-1}+c_{*}\log x
+d_{*}$, which would cover Barnes $G$ exactly as $\Wlam^{(\varrho)}$ covers
$K$.  Whether its inverse is expressible through a named function, or whether
the right statement is simply that the implicit equation defines one, is not
settled here.
\item \emph{Closed forms.}  As in \cref{p6:sec:open}, the recurrences compute
$V_n$, $U_n$ and $b_n$ to any order and their extreme coefficients are known
--- \eqref{p7:eq:V-limits} and \eqref{p7:eq:bn-limits} --- but no closed
description of the full numerators is available.  The refined blocks are the
most promising, since \eqref{p7:eq:bn-limits} says their leading behaviour is
exactly the forward tail with a sign, so the question is the structure of the
correction to that.
\item \emph{Hyperasymptotics.}  Unchanged from \cref{p6:rem:not-claimed}.  The
$K$-function tail $c_m$ inherits the factorial growth of $B_{2m+2}$, so the
least term of \eqref{p7:eq:centred} is near $m\approx\pi r$ and the inverse
floor is $\exp(-2\pi\rho+\bigO(\log\rho))/(\rho L)$ by the argument of
\cref{p6:prop:least-term}.  The Stokes multipliers are not computed.
\item \emph{The other real branches.}  \Cref{p7:prop:branch} locates the two
stationary points and hence the additional real branches below
$K(x_{\max})$.  Those branches are governed by the local behaviour at a
critical point --- a square-root, not a Lambert, scale --- and are not treated.
\item \emph{Formalization.}  Unchanged from \cref{p6:sec:open}: nothing in
either chapter is machine-checked.  The $r$-Lambert function would need its
own development.
\end{enumerate}

%% ==================== fragment p8 ====================
\chapter{The subfactorial}
\label{p8:sec:top}

\section{What this chapter does}
\label{p8:sec:intro}

The derangement numbers $D_n=\,!n=n!\sum_{j\le n}(-1)^{j}/j!$ are the first
subject in this volume whose inverse is \emph{not} governed by its own
dominant block.  The block is $\Gamma(x+1)/\e$, whose inverse is the subject
of \cref{p6:sec:gamma}; everything that is specifically about derangements
lives in a correction so small that it is beyond every fixed exponential
scale relative to the inverse variable.  The chapter is therefore organized
around that separation, and its principal economy is that the carrier is
\emph{cited}, not recomputed.

\begin{repairbox}
All three sources rederive the inverse of the gamma envelope from scratch,
each proving its own Lambert-centred reversion theorem and each printing the
blocks
\[
 \frac1{24s},\quad
 -\frac{14s^{2}+10s+5}{5760s^{3}},\quad
 \frac{2232s^{4}+1176s^{3}+714s^{2}+350s+105}{2903040s^{5}},\quad\dots
\]
These are \emph{identical}, coefficient for coefficient, to
$d_1,d_3,d_5,d_7$ of \cref{p6:prop:gamma-blocks} --- three more copies of a
computation the volume already contains twice over.  They are not repeated
here.  \Cref{p8:prop:carrier} states the dictionary and cites
\cref{p6:thm:gamma}; the general reversion theorem each source proves in order
to get them is \cref{p0:thm:perturbed-inversion}, and is cited too.
\end{repairbox}

\begin{sourcebox}
The three filed articles are
\texttt{inverse\_\allowbreak subfactorial\_\allowbreak transseries} (S1),
\texttt{inverse\_\allowbreak subfactorial\_\allowbreak transseries-2} (S2) and
\texttt{inverse\_\allowbreak subfactorial\_\allowbreak transseries-3} (S3).  Their distinctive
contributions are: S1, the exact core--tail decomposition, the median
interpolation, and the threshold analysis of \cref{p8:sec:staircase}; S2, the
signed-remainder Bell expansion with derivatives and the a posteriori error
bound; S3, the branch-indexed continuation, the single-frequency sector
factorization of \cref{p8:thm:sector-factorization}, and the multi-frequency
grid.  They agree on everything they share; every constant and coefficient was
recomputed for this volume and reproduced.
\end{sourcebox}

\section{Which inverse?}
\label{p8:sec:which}

$D_n$ is defined on $\Nat_0$, so \cref{p0:def:three-inverses} applies: there is
an integer staircase, there are interpolated inverses --- one for each
admissible interpolation --- and there is an asymptotic inverse attached to no
interpolation.  The subfactorial adds a fourth issue that the earlier chapters
did not face: its natural analytic continuation is \emph{branch-dependent},
and different branches disagree away from the integers while agreeing on them.
\Cref{p8:sec:continuation} settles this before any inverse is taken.

\section{The exact core--tail decomposition}
\label{p8:sec:continuation}

\begin{theorem}[Branch splitting]\label{p8:thm:branch-splitting}
For $\Re z>-1$ put
\begin{equation}\label{p8:eq:F-C}
 F(z)=\frac{\Gamma(z+1)}{\e},
 \qquad
 C(z)=\frac1\e\int_0^{1}t^{z}\e^{t}\,\dd t
     =\frac1\e\sum_{j\ge0}\frac1{j!\,(z+j+1)} ,
\end{equation}
and for $\kappa\in\Int$ set $\omega_\kappa=(2\kappa+1)\pi\ii$.  Then the
continuation of $\,!z$ obtained by running the incomplete-gamma path from $0$
to $-1$ with argument $(2\kappa+1)\pi$ is
\begin{equation}\label{p8:eq:branch-splitting}
 \mathcal{S}_\kappa(z)=F(z)+\e^{\omega_\kappa z}C(z),
\end{equation}
and $\mathcal{S}_\kappa(n)=D_n$ for every $n\in\Nat_0$ and every $\kappa$.
\end{theorem}

\begin{proof}
Inclusion--exclusion gives $D_n=n!\sum_{j\le n}(-1)^{j}/j!$, and the upper
incomplete gamma identity
$\Gamma(n+1,w)=n!\,\e^{-w}\sum_{j\le n}w^{j}/j!$ at $w=-1$ gives
$D_n=\Gamma(n+1,-1)/\e$.  Continue: with $s=z+1$ and the path
$t=\varrho\e^{\ii(2\kappa+1)\pi}$, $0\le\varrho\le1$, on which
$\e^{\ii(2\kappa+1)\pi}=-1$,
\[
 \gamma(s,-1)=\e^{\ii(2\kappa+1)\pi s}\int_0^{1}\varrho^{\,s-1}\e^{\varrho}
 \,\dd\varrho ,
\]
so that $\Gamma(s,-1)=\Gamma(s)-\gamma(s,-1)$ yields
\[
 \frac{\Gamma(z+1,-1)}{\e}
 =F(z)-\frac{\e^{\ii(2\kappa+1)\pi(z+1)}}{\e}\int_0^{1}\varrho^{z}
 \e^{\varrho}\,\dd\varrho
 =F(z)+\e^{\omega_\kappa z}C(z),
\]
the sign coming from $\e^{\ii(2\kappa+1)\pi}=-1$.  Expanding $\e^{t}$ under
the integral and integrating termwise gives the partial-fraction form in
\eqref{p8:eq:F-C}.  At $z=n$ the phase $\e^{\omega_\kappa n}$ equals
$(-1)^{n}$ for every $\kappa$, so all branches agree there and equal
$n!/\e+(-1)^{n}C(n)$, which is $D_n$.
\end{proof}

\begin{corollary}[The nearest-integer property]\label{p8:cor:nearest-integer}
For every $n\ge0$, $D_n=n!/\e+(-1)^{n}C(n)$, and
\begin{equation}\label{p8:eq:C-bound}
 0<C(x)<\frac1{x+1}\qquad(x>-1).
\end{equation}
Hence for $n\ge1$ the correction is less than $\tfrac12$ in absolute value and
$D_n$ is the nearest integer to $n!/\e$.
\end{corollary}

\begin{proof}
The identity is \cref{p8:thm:branch-splitting} at $z=n$.  For the bound,
$C(x)=\int_0^{1}\e^{t-1}t^{x}\dd t$ and $0<\e^{t-1}<1$ on $[0,1)$, so
$0<C(x)<\int_0^{1}t^{x}\dd t=1/(x+1)$.  For $n\ge1$ this is at most
$\tfrac12$, with equality impossible by strictness.
\end{proof}

\begin{definition}[Two real interpolations]\label{p8:def:interpolations}
The branches $\kappa=0$ and $\kappa=-1$ are complex conjugate on the real
axis; their mean is the \emph{median} (or symmetrized) interpolation
\begin{equation}\label{p8:eq:median}
 \mathcal{S}_{\mathrm{med}}(x)=F(x)+C(x)\cos(\pi x),
\end{equation}
which is real, agrees with $D_n$ at every $n\in\Nat_0$, and is entire.  The
\emph{parity envelopes} are
\begin{equation}\label{p8:eq:envelopes}
 G_\sigma(x)=F(x)+\sigma C(x),\qquad\sigma\in\{+1,-1\},
\end{equation}
so that $G_{+}(n)=D_n$ for even $n$ and $G_{-}(n)=D_n$ for odd $n$.
\end{definition}

\begin{remark}[The median is canonical, not unique]\label{p8:rem:not-unique}
Any entire function vanishing on $\Nat_0$ may be added to
$\mathcal{S}_{\mathrm{med}}$ without disturbing the interpolation property, so
\eqref{p8:eq:median} is not determined by interpolation alone.  What
distinguishes it is that it is the median of the two \emph{nearest} lateral
determinations of the same incomplete-gamma continuation, so it inherits their
asymptotics rather than importing new ones.  This is the same situation as
\cref{p2:sec:top}, where no canonical real inverse of the double factorial
exists until an interpolation is fixed; the resolution here is cleaner only
because the two candidate branches are conjugates.
\end{remark}

\begin{proposition}[Eventual monotonicity]\label{p8:prop:monotone}
There is $X$ such that $\mathcal{S}_{\mathrm{med}}$, $G_{+}$, $G_{-}$ and $F$ are
all strictly increasing on $[X,\infty)$.
\end{proposition}

\begin{proof}
$F'(x)=F(x)\psi(x+1)$, and $\psi(x+1)\sim\log x\to\infty$, so $F'>0$ and
$F'\to\infty$ eventually.  From \eqref{p8:eq:C-bound} and
$C'(x)=\int_0^{1}\e^{t-1}t^{x}\log t\,\dd t$ we get $C(x)=\bigO(x^{-1})$ and
$|C'(x)|=\bigO(x^{-2})$ --- for the latter, $|\log t|\,t^{x}$ integrates to
$(x+1)^{-2}$ on $[0,1]$.  Hence
\[
 \Bigl|\frac{\dd}{\dd x}\bigl(C(x)\cos\pi x\bigr)\Bigr|
 \le|C'(x)|+\pi|C(x)|=\bigO(x^{-1}),
\]
while $F'(x)\to\infty$; the same bound covers $\sigma C'$.  So the envelope
derivative dominates.
\end{proof}

\section{The Bell expansion of the tail}
\label{p8:sec:bell}

$C$ is a Stieltjes transform, and its moments are Bell numbers.  This gives
the tail expansion together with an exact remainder --- one of the few places
in this volume where a divergent expansion comes with a two-sided sign
statement.

\begin{theorem}[Bell expansion with signed remainder]\label{p8:thm:bell}
Let $\mathfrak{B}_m$ denote the Bell numbers, $\mathfrak{B}_1=1$, $\mathfrak{B}_2=2$, $\mathfrak{B}_3=5$,
$\mathfrak{B}_4=15$, $\mathfrak{B}_5=52$, $\mathfrak{B}_6=203$.  For every real $x>0$ and integer
$M\ge1$,
\begin{equation}\label{p8:eq:bell-expansion}
 C(x)=\sum_{m=1}^{M}(-1)^{m-1}\frac{\mathfrak{B}_m}{x^{m}}+R_M(x),
 \qquad
 R_M(x)=\frac{(-1)^{M}}{\e\,x^{M}}
 \sum_{j\ge0}\frac{(j+1)^{M}}{j!\,(x+j+1)} .
\end{equation}
In particular $(-1)^{M}R_M(x)>0$ and
\begin{equation}\label{p8:eq:bell-bound}
 |R_M(x)|\le\frac{\mathfrak{B}_{M+1}}{x^{M+1}} ,
\end{equation}
so the remainder has the sign of the first omitted term and is no larger than
it.  Consequently
$C(x)\sim\sum_{m\ge1}(-1)^{m-1}\mathfrak{B}_m x^{-m}$, an expansion that diverges for
every $x$.
\end{theorem}

\begin{proof}
For $b>0$ the finite geometric identity
\[
 \frac1{x+b}=\sum_{m=1}^{M}(-1)^{m-1}\frac{b^{m-1}}{x^{m}}
 +(-1)^{M}\frac{b^{M}}{x^{M}(x+b)}
\]
is exact.  Put $b=j+1$, divide by $\e\,j!$, and sum over $j\ge0$; the
partial-fraction form of $C$ in \eqref{p8:eq:F-C} gives the left side.  The
$m$-th coefficient is
\begin{equation}\label{p8:eq:dobinski}
 \frac1\e\sum_{j\ge0}\frac{(j+1)^{m-1}}{j!}=\mathfrak{B}_m ,
\end{equation}
which is Dobi\'nski's formula combined with the Bell recurrence
$\mathfrak{B}_m=\sum_{k}\binom{m-1}{k}\mathfrak{B}_k$: expanding $(j+1)^{m-1}$ binomially
and applying Dobi\'nski's $\mathfrak{B}_k=\e^{-1}\sum_j j^{k}/j!$ to each power gives
exactly that sum.  The last term of the geometric identity gives $R_M$, whose
summand is positive, so $(-1)^{M}R_M>0$.  Finally $x+j+1\ge x$ gives
\[
 |R_M(x)|\le\frac1{\e x^{M+1}}\sum_{j\ge0}\frac{(j+1)^{M}}{j!}
 =\frac{\mathfrak{B}_{M+1}}{x^{M+1}},
\]
again by \eqref{p8:eq:dobinski}.  Divergence follows from
$\mathfrak{B}_{2m}\ge(2m-1)!!$, so $\mathfrak{B}_m^{1/m}\to\infty$.
\end{proof}

\begin{corollary}[Derivatives]\label{p8:cor:derivatives}
For every $r\ge0$,
\begin{equation}\label{p8:eq:C-derivatives}
 C^{(r)}(x)=\frac1\e\int_0^{1}t^{x}(\log t)^{r}\e^{t}\,\dd t,
 \qquad(-1)^{r}C^{(r)}(x)>0,
\end{equation}
and termwise differentiation of the expansion is legitimate at every fixed
order:
\begin{equation}\label{p8:eq:C-derivative-asymptotic}
 C^{(r)}(x)\sim\sum_{m\ge1}(-1)^{m+r-1}(m)_r\,\mathfrak{B}_m\,x^{-m-r},
 \qquad (m)_r=m(m+1)\cdots(m+r-1).
\end{equation}
In particular $C^{(r)}(x)=\bigO(x^{-r-1})$.
\end{corollary}

\begin{proof}
Differentiation under the integral sign is justified by dominated convergence
on $[0,1]$, and $(\log t)^{r}$ has sign $(-1)^{r}$ there.  For
\eqref{p8:eq:C-derivative-asymptotic}, differentiate the exact identity
\eqref{p8:eq:bell-expansion} $r$ times: the finite sum differentiates term by
term, and $R_M^{(r)}$ obeys the same kind of bound because differentiating the
integral representation of $R_M$ only inserts a bounded power of $\log t$.
\end{proof}

\begin{remark}[Numerical check]\label{p8:rem:verification}
The moment identity \eqref{p8:eq:dobinski} was verified symbolically for
$m\le6$, giving $\mathfrak{B}_1,\dots,\mathfrak{B}_6=1,2,5,15,52,203$ as the coefficients in
\eqref{p8:eq:bell-expansion}.  The two representations of $C$ in
\eqref{p8:eq:F-C} were checked against each other to $30$ digits at
$x=5,10,20$, and at each of those points the truncation after $\mathfrak{B}_5$ was
confirmed to have an error of the predicted sign and within the bound
\eqref{p8:eq:bell-bound} --- for instance $|R_5(20)|=2.616\times10^{-6}$
against the bound $\mathfrak{B}_6/20^{6}=3.172\times10^{-6}$.  The general reversion
coefficients of \cref{p8:prop:carrier} were recomputed and, specialized to the
centred Stirling tail, reproduce $d_1,d_3$ of \cref{p6:prop:gamma-blocks}
exactly.
\end{remark}

\section{The carrier}
\label{p8:sec:carrier}

\begin{proposition}[The carrier is Chapter 5's inverse]\label{p8:prop:carrier}
Let $u(Y)$ denote the eventual inverse of $F(x)=\Gamma(x+1)/\e$, so that
$\Gamma(u+1)=\e Y$.  Then
\begin{equation}\label{p8:eq:carrier-dictionary}
 u(Y)=\Gamma_+^{-1}(\e Y)-1,
\end{equation}
and therefore, by \cref{p6:thm:gamma}, with
\begin{equation}\label{p8:eq:carrier-vars}
 R=1+\log Y-\tfrac12\log(2\pi),
 \qquad
 w=\Wlam_0(R/\e),
 \qquad
 Q=\frac Rw=\e^{w+1},
 \qquad
 s=w+1,
\end{equation}
one has, for every fixed $N$,
\begin{equation}\label{p8:eq:carrier-expansion}
 u(Y)=-\frac12+Q+\sum_{n=1}^{N}\frac{d_{2n-1}(s)}{Q^{2n-1}}
 +\bigO\!\left(\frac1{s\,Q^{2N+1}}\right),
\end{equation}
with $d_1,d_3,d_5,d_7$ exactly as in \cref{p6:prop:gamma-blocks}.  Moreover
$u(Y)\asymp\log Y/\log\log Y$.
\end{proposition}

\begin{proof}
$F(x)=Y$ means $\Gamma(x+1)=\e Y$, i.e.\ $x+1=\Gamma_+^{-1}(\e Y)$, which is
\eqref{p8:eq:carrier-dictionary}; the branch is the increasing one of
\cref{p6:rem:which-inverse}, which exists eventually by
\cref{p8:prop:monotone}.  Substituting $y=\e Y$ into
\eqref{p6:eq:gamma-core-vars} gives
$R=\log(\e Y)-\tfrac12\log(2\pi)=1+\log Y-\tfrac12\log(2\pi)$, and
\eqref{p6:eq:gamma-main} then reads
$\Gamma_+^{-1}(\e Y)=\tfrac12+Q+\sum d_{2n-1}Q^{1-2n}+\cdots$; subtracting $1$
gives \eqref{p8:eq:carrier-expansion}.  The order statement follows from
$Q=R/w$ with $w=\Wlam_0(R/\e)\sim\log R$.
\end{proof}

\begin{remark}[The shift, again]\label{p8:rem:shift}
The $-\tfrac12$ in \eqref{p8:eq:carrier-expansion} is the half-shift of
\cref{p6:rem:why-half-shift} transported through
\eqref{p8:eq:carrier-dictionary}: $\tfrac12-1=-\tfrac12$.  It is what makes
the correction series proceed in steps of $Q^{-2}$, and all three sources
remark on it independently.
\end{remark}

At $x=u(Y)$ the derivatives of the carrier are needed; write
\begin{equation}\label{p8:eq:lambda-P}
 \lambda=\lambda(u)=\psi(u+1),
 \qquad
 P_k(u)=\frac{F^{(k)}(u)}{F(u)},
\end{equation}
so that $F^{(k)}(u)=Y\,P_k(u)$ and $F'(u)=Y\lambda$.  Differentiating
$F'=F\lambda$ repeatedly gives
\begin{equation}\label{p8:eq:P-recurrence}
 P_0=1,
 \qquad
 P_{k+1}=P_k'+\lambda P_k,
\end{equation}
so $P_k=\EB_k(\lambda,\lambda',\dots,\lambda^{(k-1)})$ is the complete
exponential Bell polynomial of \cref{p0:def:ebell}; explicitly
$P_1=\lambda$, $P_2=\lambda^{2}+\lambda'$,
$P_3=\lambda^{3}+3\lambda\lambda'+\lambda''$.  Since
$\lambda=\psi(u+1)\sim\log u$ and $\lambda^{(j)}=\bigO(u^{-j})$ for $j\ge1$,
we have $P_k=\bigO\bigl((\log u)^{k}\bigr)$.

\section{Inverting a perturbation of the carrier}
\label{p8:sec:reversion}

Every inverse in this chapter solves $F(x)+E(x)=Y$ for a correction $E$ that
is small compared with $F'$.  That is exactly
\cref{p0:thm:perturbed-inversion}, and no separate reversion theorem is
needed.

\begin{proposition}[The additive reversion in carrier coordinates]
\label{p8:prop:additive}
Let $E$ be analytic near $u=u(Y)$ with $E=\littleo(F')$ there.  The solution of
$F(x)+E(x)=Y$ near $u$ is
\begin{equation}\label{p8:eq:additive}
 x=u+\sum_{n\ge1}\frac{(-1)^{n}}{n!}
 \mathscr{D}_u^{\,n-1}\!\left(\frac{E(u)^{n}}{F'(u)}\right),
 \qquad
 \mathscr{D}_u=\frac1{F'(u)}\frac{\dd}{\dd u}=\frac1{Y\lambda}\frac{\dd}{\dd u},
\end{equation}
convergent for $|E|$ small and, in the formal-jet sense, valid whenever
$F'(u)$ is invertible.
\end{proposition}

\begin{proof}
This is \cref{p0:thm:perturbed-inversion} with $F\mapsto F$, $E\mapsto E$,
$\varepsilon\mapsto1$, $\mu_0\mapsto u$, $y\mapsto Y$ and
$\mathcal{D}\mapsto\mathscr{D}_u$; the operator identification uses
$F'(u)=Y\lambda$ from \eqref{p8:eq:lambda-P}.  The formal version is
\cref{p0:cor:perturbed-formal}.
\end{proof}

\begin{remark}[Why the amplitude series converges here]
\label{p8:rem:convergence}
\Cref{p0:thm:perturbed-inversion} gives convergence for $|\varepsilon|$ below
some $\varepsilon_0$, and one wants $\varepsilon=1$.  For the subfactorial
that is legitimate, and for a stronger reason than in \cref{p3:sec:top}: in
the coordinate $v=F(x)$ the map $v\mapsto Y-\varepsilon\widetilde E(v)$, with
$\widetilde E=E\circ F^{-1}$, has
\begin{equation}\label{p8:eq:contraction}
 \widetilde E'(Y)=\frac{E'(u)}{F'(u)}
 =\bigO\!\left(\frac1{Y\,u\log u}\right),
\end{equation}
by \cref{p8:cor:derivatives} and $F'(u)=Y\lambda$.  So for each sufficiently
large fixed $Y$ the map is a contraction on a small disc for \emph{all}
$|\varepsilon|\le1$, and the Lagrange series converges at $\varepsilon=1$.
What remains merely asymptotic is the subsequent expansion of the
coefficients in powers of $u^{-1}$ and $\lambda^{-1}$; the sector expansion in
$Y^{-1}$ converges while the Bell and Stirling expansions inside each
coefficient diverge.  This is a genuinely different situation from every
earlier chapter and is worth stating plainly.
\end{remark}

\section{Sector factorization}
\label{p8:sec:sectors}

The correction is not merely small: it carries an exponential factor
$\e^{\omega z}$, and that factor organizes the whole inverse.

\begin{theorem}[Single-frequency sector factorization]
\label{p8:thm:sector-factorization}
Let $E(x)=\e^{\omega x}C(x)$ with $\omega$ fixed and $C$ as in
\eqref{p8:eq:F-C}.  Then the solution of $F(x)+E(x)=Y$ has the formal
expansion
\begin{equation}\label{p8:eq:sector-transseries}
 x\sim u(Y)+\sum_{n\ge1}\frac{\e^{n\omega u(Y)}}{Y^{n}}\,
 \Phi_n^{(\omega)}\bigl(u(Y)\bigr),
\end{equation}
where, writing
\begin{equation}\label{p8:eq:L-operator}
 \mathcal{L}^{(\omega)}_{n,r}H=\frac{H'+(n\omega-r\lambda)H}{\lambda},
\end{equation}
the coefficients are the finite operator products
\begin{equation}\label{p8:eq:Phi-product}
 \Phi_n^{(\omega)}
 =\frac{(-1)^{n}}{n!}\,
 \bigl(\mathcal{L}^{(\omega)}_{n,n-1}\circ\cdots\circ
 \mathcal{L}^{(\omega)}_{n,1}\bigr)\!
 \left(\frac{C^{n}}{\lambda}\right),
\end{equation}
the empty product being the identity.  Explicitly
\begin{align}
 \Phi_1^{(\omega)}&=-\frac C\lambda,
 \label{p8:eq:Phi1}\\
 \Phi_2^{(\omega)}&=\frac{C\,D_\omega C}{\lambda^{2}}
 -\frac{C^{2}P_2}{2\lambda^{3}},
 \qquad D_\omega:=\frac{\dd}{\dd u}+\omega,
 \label{p8:eq:Phi2}\\
 \Phi_3^{(\omega)}&=-\frac{C(D_\omega C)^{2}}{\lambda^{3}}
 -\frac{C^{2}D_\omega^{2}C}{2\lambda^{3}}
 +\frac{3C^{2}(D_\omega C)P_2}{2\lambda^{4}}
 +\frac{C^{3}P_3}{6\lambda^{4}}
 -\frac{C^{3}P_2^{2}}{2\lambda^{5}} .
 \label{p8:eq:Phi3}
\end{align}
Moreover
\begin{equation}\label{p8:eq:Phi-size}
 \Phi_n^{(\omega)}(u)=\bigO\!\left(\frac{u^{-n}}{\log u}\right),
\end{equation}
so the grading by powers of $Y^{-1}$ is asymptotically strict even when
$\omega$ is purely imaginary.
\end{theorem}

\begin{proof}
Start from \eqref{p8:eq:additive}.  Its $n$-th summand begins with
\[
 \frac{E(u)^{n}}{F'(u)}
 =\frac{\e^{n\omega u}C(u)^{n}}{Y\lambda(u)},
\]
and the key observation is that $\mathscr{D}_u$ maps objects of this shape to
objects of the same shape.  Let $H$ be any function of $u$ and $r\ge1$.  Along
the differentiation $Y$ is \emph{not} constant --- it is $F(u)$ --- so
$\dd(Y^{-r})/\dd u=-rY^{-r-1}F'(u)=-rY^{-r}\lambda$, and therefore
\begin{equation}\label{p8:eq:operator-step}
 \mathscr{D}_u\!\left(\frac{\e^{n\omega u}H(u)}{Y^{r}}\right)
 =\frac1{Y\lambda}\left(
 \frac{\e^{n\omega u}\bigl(H'+n\omega H\bigr)}{Y^{r}}
 -\frac{r\lambda\,\e^{n\omega u}H}{Y^{r}}\right)
 =\frac{\e^{n\omega u}}{Y^{r+1}}\cdot
 \frac{H'+(n\omega-r\lambda)H}{\lambda}
 =\frac{\e^{n\omega u}}{Y^{r+1}}\,\mathcal{L}^{(\omega)}_{n,r}H .
\end{equation}
Applying this $n-1$ times to $H=C^{n}/\lambda$ starting at $r=1$ gives
\eqref{p8:eq:Phi-product}, and the prefactor $(-1)^{n}/n!$ is the one in
\eqref{p8:eq:additive}.  The displayed $\Phi_1,\Phi_2,\Phi_3$ follow by
expanding the products and using $P_2=\lambda^{2}+\lambda'$,
$P_3=\lambda^{3}+3\lambda\lambda'+\lambda''$.

For \eqref{p8:eq:Phi-size}: by \cref{p8:cor:derivatives},
$C^{(j)}=\bigO(u^{-j-1})$, so $C^{n}=\bigO(u^{-n})$; $\lambda\sim\log u$ and
$P_k=\bigO((\log u)^{k})$; and each application of
$\mathcal{L}^{(\omega)}_{n,r}$ divides by $\lambda$ while either
differentiating (which lowers the power of $u$ but not below $u^{-n}$ overall,
since the leading behaviour is set by $C^{n}$) or multiplying by the bounded
quantity $n\omega-r\lambda$, whose size $\bigO(\log u)$ cancels the division.
Induction gives \eqref{p8:eq:Phi-size}.
\end{proof}

\begin{corollary}[Multi-frequency grids]\label{p8:cor:multi}
If $E(x)=\sum_{\alpha=1}^{M}\e^{\omega_\alpha x}C_\alpha(x)$, then the inverse
support lies on the grid
$\bigl\{Y^{-|\bm r|}\e^{(\bm r\cdot\bm\omega)u}:\bm r\in\Nat^{M}\setminus
\{\bm0\}\bigr\}$ and
\begin{equation}\label{p8:eq:multi}
 x\sim u(Y)+\sum_{\bm r\ne\bm0}
 Y^{-|\bm r|}\e^{(\bm r\cdot\bm\omega)u(Y)}\Phi_{\bm r}(u(Y)),
\end{equation}
ordered first by total degree $|\bm r|$.  Purely imaginary frequencies may give
equal magnitudes or resonant phases, but the factor $Y^{-|\bm r|}$ keeps the
total-degree grading well founded.
\end{corollary}

\begin{proof}
Expand $E(u)^{n}$ in \eqref{p8:eq:additive} by the multinomial theorem and
apply \eqref{p8:eq:operator-step} to each resulting monomial; the frequency
$\bm r\cdot\bm\omega$ and the power $Y^{-|\bm r|}$ are preserved by
$\mathscr{D}_u$, as the computation there shows.
\end{proof}

\section{The inverses}
\label{p8:sec:inverses}

\begin{theorem}[Branch inverse]\label{p8:thm:branch-inverse}
Fix $\kappa\in\Int$ and let $z_\kappa(Y)$ be the solution of
$\mathcal{S}_\kappa(z)=Y$ with $z_\kappa(Y)-u(Y)\to0$.  Then, for every fixed
$N$,
\begin{equation}\label{p8:eq:branch-inverse}
 z_\kappa(Y)=u(Y)+\sum_{n=1}^{N}
 \frac{\e^{n\omega_\kappa u(Y)}}{Y^{n}}\Phi_n^{(\omega_\kappa)}(u(Y))
 +\bigO\!\left(\frac{Y^{-N-1}u^{-N-1}}{\log u}\right),
\end{equation}
with $\Phi_n$ from \eqref{p8:eq:Phi-product} and $\omega_\kappa=(2\kappa+1)
\pi\ii$.  In particular the leading displacement is
\begin{equation}\label{p8:eq:branch-leading}
 z_\kappa(Y)-u(Y)=-\frac{\e^{\omega_\kappa u}C(u)}{Y\lambda(u)}
 \bigl(1+\littleo(1)\bigr),
\end{equation}
which is $\bigO\bigl(1/(Y\,u\log u)\bigr)$ in modulus.
\end{theorem}

\begin{proof}
Existence and uniqueness near $u$: by \cref{p8:thm:branch-splitting},
$\mathcal{S}_\kappa(u)-Y=\e^{\omega_\kappa u}C(u)=\bigO(u^{-1})$ while
$\mathcal{S}_\kappa'(u)=Y\lambda+\bigO(u^{-1})$, so
\cref{p0:thm:backward-error} places a root within
$\bigO\bigl(1/(Y\lambda u)\bigr)$ of $u$ and \cref{p8:rem:convergence} makes
it unique in a disc of that size.  The expansion is
\cref{p8:thm:sector-factorization} with $\omega=\omega_\kappa$, and the
remainder is \eqref{p8:eq:Phi-size} applied to the first omitted sector
together with the geometric domination of the convergent tail from
\eqref{p8:eq:contraction}.
\end{proof}

\begin{theorem}[The real median inverse]\label{p8:thm:median-inverse}
Let $\rho(Y)=\mathcal{S}_{\mathrm{med}}^{-1}(Y)$ be the eventual real inverse
of \eqref{p8:eq:median}, and put $H(u)=C(u)\cos(\pi u)$.  Then
\begin{equation}\label{p8:eq:median-inverse}
 \rho(Y)\sim u(Y)+\sum_{n\ge1}\frac{(-1)^{n}}{n!}
 \mathscr{D}_u^{\,n-1}\!\left(\frac{H(u)^{n}}{Y\lambda}\right),
\end{equation}
whose $n$-th term contains the harmonics $\e^{\ii m\pi u}$ with
$m=-n,-n+2,\dots,n$.  Explicitly,
\begin{equation}\label{p8:eq:median-two-sectors}
 \rho(Y)=u-\frac H{Y\lambda}
 +\frac1{Y^{2}}\left[\frac{HH'}{\lambda^{2}}
 -\frac{H^{2}(\lambda^{2}+\lambda')}{2\lambda^{3}}\right]
 +\bigO\!\left(\frac{Y^{-3}u^{-3}}{\log u}\right),
\end{equation}
and the leading displacement is
\begin{equation}\label{p8:eq:median-leading}
 \rho(Y)-u(Y)\sim-\frac{\cos\bigl(\pi u(Y)\bigr)}{Y\,u(Y)\log u(Y)} .
\end{equation}
At every derangement ordinate, $\rho(D_n)=n$ exactly.
\end{theorem}

\begin{proof}
Apply \cref{p8:prop:additive} with $E=H$; \eqref{p8:eq:median-inverse} is
\eqref{p8:eq:additive}.  Writing $\cos\pi u=\tfrac12(\e^{\ii\pi u}
+\e^{-\ii\pi u})$ and expanding $H^{n}$ shows that the $n$-th term is a
combination of $\e^{\ii m\pi u}$ with $m\equiv n\pmod2$ and $|m|\le n$, which
is \cref{p8:cor:multi} for the two frequencies $\pm\pi\ii$.  For
\eqref{p8:eq:median-leading}, $\Phi_1=-H/\lambda$ with $C(u)\sim1/u$ from
\cref{p8:thm:bell} and $\lambda\sim\log u$.  The interpolation statement is
$\mathcal{S}_{\mathrm{med}}(n)=D_n$ from
\cref{p8:def:interpolations}.
\end{proof}

\begin{corollary}[Parity envelopes]\label{p8:cor:parity}
For $\sigma=\pm1$ the eventual inverse of $G_\sigma=F+\sigma C$ is
\begin{equation}\label{p8:eq:parity-inverse}
 x_\sigma(Y)\sim u(Y)+\sum_{n\ge1}\frac{\sigma^{n}}{Y^{n}}
 \Phi_n^{(0)}(u(Y)),
\end{equation}
which is \cref{p8:thm:sector-factorization} at $\omega=0$ --- no oscillation,
and the coefficients are the same operator products with $D_0=\dd/\dd u$.  In
particular
$x_\sigma(Y)=u-\sigma C/(Y\lambda)+\bigO\bigl(Y^{-2}u^{-2}/\log u\bigr)$.
\end{corollary}

\begin{proof}
Immediate from \cref{p8:thm:sector-factorization} with $\omega=0$ and the
amplitude $\sigma$, which enters as $\sigma^{n}$ through $E^{n}$.
\end{proof}

\section{The staircase and the thresholds}
\label{p8:sec:staircase}

The discrete inverse is governed by \cref{p0:thm:staircase}, and the only
subject-specific input is where the jumps sit.

\begin{definition}[Thresholds]\label{p8:def:thresholds}
For $n\ge2$ let $\alpha_n$ be defined by $F(\alpha_n)=D_n$, that is
$\alpha_n=u(D_n)$.  Since $F$ is increasing, the generalized inverse of the
sequence $(D_n)$ is
\begin{equation}\label{p8:eq:N-threshold}
 N_\ast(y)=\min\{n\ge2:\,u(y)\le\alpha_n\} .
\end{equation}
\end{definition}

\begin{theorem}[Threshold displacement]\label{p8:thm:thresholds}
Write $\varsigma_n=(-1)^{n}$, $J_n=C(n)$ and $\Pi_n=\psi(n+1)$.  Then
\begin{equation}\label{p8:eq:threshold-taylor}
 \alpha_n-n
 =\frac{\varsigma_nJ_n}{F'(n)}
 -\frac{F''(n)J_n^{2}}{2F'(n)^{3}}
 +\frac{\varsigma_n\bigl(3F''(n)^{2}-F'(n)F'''(n)\bigr)J_n^{3}}
 {6F'(n)^{5}}
 +\bigO\!\left(\frac{J_n^{4}}{F(n)^{4}\Pi_n}\right),
\end{equation}
and in particular
\begin{equation}\label{p8:eq:threshold-leading}
 \alpha_n=n+(-1)^{n}\frac{\e\,C(n)}{n!\;\psi(n+1)}
 +\bigO\!\left(\frac1{(n!)^{2}n^{2}\log n}\right)
 =n+(-1)^{n}\frac{\e}{n!\,(H_n-\gamma)}
 \left(\frac1n-\frac2{n^{2}}+\frac5{n^{3}}-\cdots\right).
\end{equation}
Consequently $|\alpha_n-n|$ is factorially small, and for every $n\ge2$ the
separation condition \eqref{p0:eq:separation} holds with room to spare.
\end{theorem}

\begin{proof}
By \cref{p8:cor:nearest-integer}, $D_n=F(n)+\varsigma_nJ_n$, so $\alpha_n$ is
the value of the local inverse of $F$ at $F(n)+\varsigma_nJ_n$.  Taylor
expanding that inverse, with
\[
 (F^{-1})'=\frac1{F'},\quad
 (F^{-1})''=-\frac{F''}{(F')^{3}},\quad
 (F^{-1})'''=\frac{3(F'')^{2}-F'F'''}{(F')^{5}},
\]
and substituting the increment $\varsigma_nJ_n$ gives
\eqref{p8:eq:threshold-taylor}.  Now $F(n)=n!/\e$ and $F'(n)=F(n)\Pi_n$, so
the first term is $\varsigma_nJ_n\e/(n!\,\Pi_n)$, which with
$\Pi_n=\psi(n+1)=H_n-\gamma$ and the Bell expansion of $J_n=C(n)$ from
\cref{p8:thm:bell} gives \eqref{p8:eq:threshold-leading}; the second term is
$\bigO(J_n^{2}/(F(n)^{2}\Pi_n))=\bigO\bigl((n!)^{-2}n^{-2}(\log n)^{-1}\bigr)$
because $F''=F(\Pi^{2}+\Pi')$ and $J_n=\bigO(n^{-1})$.  Since
$|\alpha_n-n|\le\e/(n!\,(H_n-\gamma)\,n)$ is far below $\tfrac12$ for $n\ge2$,
\cref{p0:thm:staircase}(ii) applies with any $\eta$ up to $\tfrac13$, say.
\end{proof}

\begin{remark}[What the staircase costs]\label{p8:rem:staircase-cost}
\Cref{p8:thm:thresholds} says the discrete inverse is easy: the thresholds sit
factorially close to the integers, so rounding $u(y)$ correctly needs only
$|u(y)-\nu(y)|<\tfrac12$, which \eqref{p8:eq:carrier-expansion} delivers
already at $N=0$.  In this respect the subfactorial is far better behaved than
the sequences of \cref{p1:sec:top,p3:sec:top}, where the interpolated inverse
and the staircase separate only slowly.  The reason is structural: the
correction that distinguishes $D_n$ from $n!/\e$ is the same correction that
displaces the threshold, and it is factorially small in both roles.
\end{remark}

\section{The four scales}
\label{p8:sec:hierarchy}

\begin{proposition}[Hierarchy]\label{p8:prop:hierarchy}
As $Y\to+\infty$, with $u=u(Y)$, the inverse of any of the interpolations of
\cref{p8:def:interpolations} decomposes into four strictly ordered layers:
\begin{enumerate}
\item the \emph{Lambert--logarithmic} layer, the carrier $Q$ of
\eqref{p8:eq:carrier-vars}, of size $u$;
\item the \emph{Stirling} layer, the blocks $d_{2n-1}(s)Q^{1-2n}$ of
\eqref{p8:eq:carrier-expansion}, of sizes $u^{-1},u^{-3},\dots$ divided by
powers of $s$;
\item the \emph{hyperasymptotic gamma} layer, below the least term of the
Stirling series, of size $u^{-1/2}\e^{-2\pi u}$ by
\cref{p6:prop:least-term};
\item the \emph{Bell} layer, the genuinely subfactorial correction, of size
\begin{equation}\label{p8:eq:bell-layer-size}
 \frac1{Y\,u\log u}=\exp\bigl(-u\log u+\bigO(u)\bigr).
\end{equation}
\end{enumerate}
Layer \textup{(iv)} is smaller than $\e^{-cu}$ for every fixed $c>0$.
\end{proposition}

\begin{proof}
Only the last statement needs an argument.  By
\cref{p8:prop:carrier}, $Y=F(u)=\Gamma(u+1)/\e$, so
$\log Y=u\log u-u+\bigO(\log u)$ by Stirling; hence
$1/(Y u\log u)=\exp\bigl(-u\log u+u+\bigO(\log u)\bigr)$, which is
\eqref{p8:eq:bell-layer-size}.  Since $u\log u-u\ge cu$ eventually for every
fixed $c$, the claim follows.  The ordering of (i)--(iii) is
\cref{p6:rem:hierarchy} transported by \eqref{p8:eq:carrier-dictionary}, and
(iii) precedes (iv) because $\e^{-2\pi u}$ is larger than
$\e^{-u\log u+\bigO(u)}$.
\end{proof}

\begin{remark}[The consequence for practice]\label{p8:rem:practice}
\Cref{p8:prop:hierarchy} says something uncomfortable and worth saying: for
any $Y$ at which one could actually compute, the subfactorial correction to
the inverse is far below the accuracy that the Stirling series can deliver at
all.  Recovering it numerically requires either exact arithmetic on
$\Gamma(u+1)/\e$ or a hyperasymptotic treatment of layer (iii).  What the
Bell sector \emph{is} good for is exact statements at integers ---
\cref{p8:cor:nearest-integer} and \cref{p8:thm:thresholds} --- where the
correction is not competing with a divergent series but is the whole content.
\end{remark}

\section{What is left open}
\label{p8:sec:open}

\begin{enumerate}
\item \emph{Closed forms for $\Phi_n$.}  \Cref{p8:eq:Phi-product} is a finite
operator product, hence computable, but its expansion in the differential ring
generated by $C$, $\lambda$ and their derivatives has no closed description
beyond $n=3$.  Because $P_k$ is a complete Bell polynomial
\eqref{p8:eq:P-recurrence} and $C$ has Bell-number coefficients, a
double-Bell formula is plausible; it is not derived here.
\item \emph{The hyperasymptotic gap.}  Layer (iii) of
\cref{p8:prop:hierarchy} separates the algebraic inverse from the Bell sector,
and nothing in this chapter bridges it.  A complete numerical theory of the
real inverse below $\e^{-2\pi u}$ needs the exponentially improved inverse
gamma of \cref{p6:rem:not-claimed}, which is itself open.
\item \emph{Uniqueness of the real interpolation.}  \Cref{p8:rem:not-unique}:
the median is canonical among the lateral determinations but not singled out
by interpolation alone.  A characterization --- for instance by a growth or
complete-monotonicity condition in the style of the Bohr--Mollerup theorem ---
is not attempted, and would be the right way to make ``the'' real inverse
subfactorial well defined.
\item \emph{Complex branches.}  \Cref{p8:thm:branch-inverse} is stated for
fixed $\kappa$ along the positive real $Y$-axis.  How the sheets $z_\kappa$
are connected, and what happens when $\kappa$ grows with $Y$, is not
addressed.
\end{enumerate}

%% ==================== fragment p9 ====================
\chapter{A real-argument Fibonacci function}
\label{p9:sec:top}

\section{What this chapter does, and what it cannot use}
\label{p9:sec:intro}

This is the one subject in the volume with no Lambert core.  Its dominant
block is a pure exponential, $\varphi^{x}$, so the equation to be inverted is
$\e^{\lambda x}\approx R$ and the leading inverse is an honest logarithm ---
no $\Wlam$, no slow coefficient field, no $(v+1)^{-1}$.  What replaces the
apparatus of \cref{p6:sec:core} is a different and, in one respect, much
better structure: the inverse equation \emph{factors exactly} into conjugate
factors, the reversion is an ordinary analytic reversion at an ordinary
point, and the resulting expansion \emph{converges}.

The chapter is therefore organized around what carries over and what does
not.

\begin{center}\small
\begin{tabular}{@{}p{0.46\textwidth}p{0.46\textwidth}@{}}
\toprule
Applies unchanged & Does not apply\\
\midrule
\Cref{p0:thm:perturbed-inversion} and its formal form: the correction is an
analytic perturbation of an exactly invertible core.
&
\Cref{p0:def:model}: the phase is $\lambda x$, not $ax^{\alpha}$ with a
subdominant $b\log x$; the model's $\log x$ slot is empty and its $Q$ slot
carries oscillation instead of a power series.\\[0.4em]
\Cref{p0:thm:LB} and \cref{p0:lem:LB-universal}, in their multivariate
Lagrange--Good form \eqref{p9:eq:good}.
&
\Cref{p0:thm:lambert-core} and \cref{p6:lem:core}: nothing needs inverting
beyond an exponential.\\[0.4em]
\Cref{p0:thm:backward-error}, for the residual certificate.
&
\Cref{p0:thm:flattening}: there is no $\Wlam$ to unfold, and the expansion is
already elementary.\\[0.4em]
\Cref{p0:def:three-inverses} and \cref{p0:thm:staircase}: $F_n$ is a
sequence.
&
\Cref{p0:thm:optimal-truncation}: the expansion converges, so there is no
least term and no beyond-all-orders floor.\\
\bottomrule
\end{tabular}
\end{center}

\begin{sourcebox}
The three filed articles are
\texttt{Fibonacci\_\allowbreak Inverse\_\allowbreak LogPeriodic\_\allowbreak Transseries} (V1),
\texttt{fibonacci\_\allowbreak inverse\_\allowbreak transseries\_\allowbreak article} (V2) and
\texttt{fibonacci\_\allowbreak inverse\_\allowbreak transseries\_\allowbreak article-2} (V3).  V1 supplies the
conjugate factorization and the Lagrange--Good closed formula; V2 the
structural laws for the real blocks, the residual certificate and the
convergence radius; V3 the general block-reversion formula for products of
exponentials.  They agree, and every constant, coefficient and checkpoint was
recomputed for this volume (\cref{p9:rem:verification}).
\end{sourcebox}

\section{The interpolation and its branches}
\label{p9:sec:interpolation}

\begin{definition}[The even Fibonacci modulus]\label{p9:def:F}
Let
\begin{equation}\label{p9:eq:constants}
 \varphi=\frac{1+\sqrt5}2,
 \qquad
 \lambda=2\log\varphi,
 \qquad
 \omega=\pi,
 \qquad
 \rho=\frac\omega\lambda,
\end{equation}
and define, for $x\in\Real$,
\begin{equation}\label{p9:eq:F-def}
 \mathcal{F}(x)=\sqrt{\tfrac25}\,
 \sqrt{\cosh(\lambda x)-\cos(\omega x)}
 \ \ \ge0 .
\end{equation}
Numerically $\lambda=0.9624236501192068949955\ldots$ and
$\rho=3.2642513026364969\ldots$.
\end{definition}

\begin{proposition}[Interpolation, evenness, and the modulus form]
\label{p9:prop:interpolation}
$\mathcal{F}$ is even, real, nonnegative, vanishes only at $x=0$, and
satisfies $\mathcal{F}(n)=F_n$ for every integer $n\ge0$, where $F_0=0$,
$F_1=1$, $F_{n+2}=F_{n+1}+F_n$.  Equivalently
\begin{equation}\label{p9:eq:F-squared}
 \mathcal{F}(x)^{2}
 =\tfrac15\bigl(\e^{\lambda x}+\e^{-\lambda x}-2\cos(\omega x)\bigr)
 \qquad\text{and}\qquad
 \mathcal{F}(x)=\frac1{\sqrt5}\bigl|\varphi^{x}-\e^{\ii\pi x}\varphi^{-x}\bigr| .
\end{equation}
Near the origin $\mathcal{F}(x)=\sqrt{(\lambda^{2}+\omega^{2})/5}\;|x|
+\bigO(|x|^{3})$, so the interpolation has a cusp there.
\end{proposition}

\begin{proof}
Evenness and nonnegativity are clear from \eqref{p9:eq:F-def}, and
$\cosh(\lambda x)\ge1\ge\cos(\omega x)$ with equality only at $x=0$.
Expanding $\cosh$ gives the first identity in \eqref{p9:eq:F-squared}, and
since $\e^{\lambda x}=\varphi^{2x}$ that expression is
$\tfrac15|\varphi^{x}-\e^{\ii\pi x}\varphi^{-x}|^{2}$, which is the second.
For the interpolation, Binet's formula $F_n=(\varphi^{n}
-(-\varphi^{-1})^{n})/\sqrt5$ gives
\[
 F_n^{2}=\tfrac15\bigl(\varphi^{2n}+\varphi^{-2n}-2(-1)^{n}\bigr)
 =\tfrac15\bigl(\e^{\lambda n}+\e^{-\lambda n}-2\cos(\pi n)\bigr)
 =\mathcal{F}(n)^{2},
\]
and $F_n\ge0$ for $n\ge0$, so taking nonnegative roots gives
$\mathcal{F}(n)=F_n$.  The behaviour at the origin follows from
$\cosh(\lambda x)-\cos(\omega x)
=\tfrac12(\lambda^{2}+\omega^{2})x^{2}+\bigO(x^{4})$.
\end{proof}

\begin{remark}[Which extension this is not]\label{p9:rem:not-signed}
$\mathcal{F}$ is even, so $\mathcal{F}(-n)=F_n$, whereas the usual signed
convention has $F_{-n}=(-1)^{n+1}F_n$.  The choice matters when the two outer
inverse branches are named, and it is what makes the negative branch trivially
$-X_+$.
\end{remark}

\begin{proposition}[The outer branches]\label{p9:prop:branches}
Let
\begin{equation}\label{p9:eq:xstar}
 x_{*}=\frac1\lambda\operatorname{arsinh}\!\left(\frac\omega\lambda\right)
 =1.972992734136994\ldots .
\end{equation}
Then $\mathcal{F}'(x)>0$ for every $x>x_{*}$; in particular $\mathcal{F}$ maps
$[2,\infty)$ strictly increasingly onto $[1,\infty)$.  Write $X_+(y)$ for the
inverse of that restriction and $X_-(y)=-X_+(y)$.  Every sufficiently large
$y$ has exactly the two real preimages $X_\pm(y)$.
\end{proposition}

\begin{proof}
Differentiating \eqref{p9:eq:F-squared},
\begin{equation}\label{p9:eq:F-derivative}
 \mathcal{F}'(x)
 =\frac{\lambda\sinh(\lambda x)+\omega\sin(\omega x)}{5\,\mathcal{F}(x)},
\end{equation}
so the sign is that of the numerator.  Since $\sin(\omega x)\ge-1$,
the numerator is at least $\lambda\sinh(\lambda x)-\omega$, which is positive
exactly for $x>x_{*}$.  As $x_{*}<2$ and $\mathcal{F}(2)=F_2=1$ with
$\mathcal{F}\to\infty$, the stated bijection follows.  Finally $[0,2]$ is
compact, so $\mathcal{F}$ is bounded there and every $y$ above that bound has
no preimage in $[0,2]$; with evenness this leaves exactly $X_\pm(y)$.
\end{proof}

\begin{remark}[Inner branches]\label{p9:rem:inner}
$\mathcal{F}$ is not monotone on $(0,2)$: it has two critical points there, at
$x=1.132379068\ldots$ and $x=1.719250886\ldots$, with values
$1.013859885\ldots$ and $0.911172684\ldots$.  So values slightly above $1$
have extra preimages.  These are governed by a simple critical point, where
the local inverse is a Puiseux series in $\sqrt{y-y_{\mathrm{crit}}}$ rather
than anything in this chapter's scale; V2 treats them and we do not.  They are
irrelevant to the large-$y$ problem.
\end{remark}

\section{Exact factorization}
\label{p9:sec:factorization}

Everything in this chapter rests on one identity.

\begin{theorem}[Conjugate factorization]\label{p9:thm:factorization}
Fix the outer branch, let $y\ge1$, and put
\begin{equation}\label{p9:eq:R-u}
 R=5y^{2},
 \qquad
 u=\e^{-\lambda x}\in(0,1),
 \qquad
 \alpha=1-\ii\rho,
 \qquad
 \bar\alpha=1+\ii\rho .
\end{equation}
Then $\mathcal{F}(x)=y$ is \emph{exactly} equivalent to
\begin{equation}\label{p9:eq:factorization}
 R\,u=(1-u^{\alpha})(1-u^{\bar\alpha}),
\end{equation}
where $u^{\alpha}:=\e^{\alpha\log u}$ with the real logarithm of
$u\in(0,1)$.
\end{theorem}

\begin{proof}
By \eqref{p9:eq:F-squared}, $\mathcal{F}(x)=y$ is
$R=\e^{\lambda x}+\e^{-\lambda x}-2\cos(\omega x)$.  With
$u=\e^{-\lambda x}$ we have $\log u=-\lambda x$, so
$\e^{\ii\omega x}=\e^{-\ii\rho\log u}=u^{-\ii\rho}$ and likewise
$\e^{-\ii\omega x}=u^{\ii\rho}$; hence
$2\cos(\omega x)=u^{\ii\rho}+u^{-\ii\rho}$ and
\[
 R=u^{-1}+u-u^{\ii\rho}-u^{-\ii\rho}.
\]
Multiplying by $u$ gives
$Ru=1+u^{2}-u^{1+\ii\rho}-u^{1-\ii\rho}$, and expanding the right side of
\eqref{p9:eq:factorization},
\[
 (1-u^{1-\ii\rho})(1-u^{1+\ii\rho})
 =1-u^{1-\ii\rho}-u^{1+\ii\rho}+u^{2},
\]
which is the same.
\end{proof}

\begin{remark}[Why this is the decisive step]\label{p9:rem:why-factor}
Reverting \eqref{p9:eq:F-def} directly is possible but hides the structure:
one is inverting a square root of a difference of two oscillating
exponentials.  After squaring and substituting $u=\e^{-\lambda x}$, the
problem becomes the reversion of an \emph{ordinary analytic germ at the
origin} in the two variables $u^{\alpha}$ and $u^{\bar\alpha}$.  That is why
the answer converges and why a closed coefficient formula exists at all.  It
is also the precise sense in which this subject is easier than the other
seven, not harder: the exponents are complex, but nothing is divergent.
\end{remark}

\section{The closed coefficient formula}
\label{p9:sec:good}

\begin{theorem}[Lagrange--Good form of the outer inverse]\label{p9:thm:good}
For all sufficiently large $y$, with $R=5y^{2}$,
\begin{equation}\label{p9:eq:good}
 X_+(y)=\frac{\log R}{\lambda}
 +\sum_{\substack{j,k\ge0\\ j+k\ge1}}
 \frac{(-1)^{j+k+1}}{\lambda\,s_{jk}}
 \binom{s_{jk}}{j}\binom{s_{jk}}{k}\,R^{-s_{jk}},
 \qquad
 s_{jk}=j+k-\ii\rho(j-k),
\end{equation}
the series converging absolutely.  Truncating at $j+k\le N$ leaves an error
$\bigO(R^{-N-1})=\bigO(y^{-2N-2})$.  The sum is real, because the terms
$(j,k)$ and $(k,j)$ are complex conjugates.
\end{theorem}

\begin{proof}
Write $u=\e^{-\lambda x}$, so $x=-\lambda^{-1}\log u$ and, by
\cref{p9:thm:factorization}, $u$ is the small root of
\begin{equation}\label{p9:eq:fixed-point}
 u=\frac{(1-u^{\alpha})(1-u^{\bar\alpha})}{R}
 =:\frac1R\,\Psi(u).
\end{equation}
Set $v=R^{-1}$.  Then $u=v\Psi(u)$ with $\Psi$ analytic at $u=0$ and
$\Psi(0)=1$; so $u=u(v)$ is the unique solution with $u(0)=0$, analytic near
$v=0$, and Lagrange inversion (\cref{p0:thm:LB} in its convergent analytic
form) applies to the observable $G(u)=-\lambda^{-1}\log u$.  Since $G$ has a
logarithmic singularity at $0$ we use the standard variant: for
$n\ge1$,
\[
 \Coef{v^{n}}\bigl(G(u(v))-G_{\mathrm{sing}}\bigr)
 =\frac1n\Coef{u^{n-1}}\bigl(G'(u)\,\Psi(u)^{n}\bigr)
 =-\frac1{\lambda n}\Coef{u^{n}}\Psi(u)^{n},
\]
using $G'(u)u=-1/\lambda$.  Summing and adding the singular part
$-\lambda^{-1}\log v=\lambda^{-1}\log R$ gives
\begin{equation}\label{p9:eq:good-intermediate}
 X_+(y)=\frac{\log R}{\lambda}
 -\frac1\lambda\sum_{n\ge1}\frac1n\Coef{u^{n}}\Psi(u)^{n}\,R^{-n} .
\end{equation}
It remains to extract $\Coef{u^{n}}\Psi^{n}$.  Here the exponents are not
integers, and the clean way to say what ``$\Coef{u^{n}}$'' means is to treat
$u^{\alpha}$ and $u^{\bar\alpha}$ as independent variables --- which is the
content of Good's multivariate form.  Expanding
\[
 \Psi(u)^{n}=(1-u^{\alpha})^{n}(1-u^{\bar\alpha})^{n}
 =\sum_{j,k\ge0}(-1)^{j+k}\binom nj\binom nk u^{\alpha j+\bar\alpha k},
\]
the monomial $u^{\alpha j+\bar\alpha k}=u^{\,j+k-\ii\rho(j-k)}=u^{\,s_{jk}}$
contributes to the coefficient of $u^{n}$ exactly when $n=s_{jk}$.  Making
that substitution in \eqref{p9:eq:good-intermediate} --- i.e.\ summing over the
pairs $(j,k)$ and setting $n=s_{jk}$ throughout --- gives
\[
 X_+(y)=\frac{\log R}{\lambda}
 -\frac1\lambda\sum_{j+k\ge1}\frac{(-1)^{j+k}}{s_{jk}}
 \binom{s_{jk}}{j}\binom{s_{jk}}{k}R^{-s_{jk}},
\]
which is \eqref{p9:eq:good}.

Convergence: $|R^{-s_{jk}}|=R^{-(j+k)}$ because $\Re s_{jk}=j+k$, and
$\bigl|\binom{s_{jk}}{j}\binom{s_{jk}}{k}/s_{jk}\bigr|$ grows at most
geometrically in $j+k$ --- each binomial is a product of $j$ resp.\ $k$
factors of modulus at most $j+k+\rho|j-k|+1\le(1+\rho)(j+k)+1$, divided by
$j!$ resp.\ $k!$, so the modulus is at most
$\bigl((1+\rho)(j+k)+1\bigr)^{j+k}/(j!\,k!)$, and by
$j!\,k!\ge((j+k)/2)!^{2}/\binom{j+k}{j}$ and Stirling this is
$\bigO(c^{\,j+k})$ for a constant $c$ depending only on $\rho$.  Hence the
series converges absolutely for $R>c$, and the tail beyond $j+k=N$ is
$\bigO(R^{-N-1})$.

Reality: $\overline{s_{jk}}=s_{kj}$ and $\overline{R^{-s_{jk}}}=R^{-s_{kj}}$
since $R$ is real and positive, while $\binom{\bar s}{j}=\overline{\binom sj}$;
so the $(j,k)$ and $(k,j)$ terms are conjugate and the total is real.
\end{proof}

\begin{remark}[On the coefficient extraction]\label{p9:rem:good-care}
The step ``$u^{s_{jk}}$ contributes to $\Coef{u^{n}}$ when $n=s_{jk}$''
deserves the caution V1 gives it.  The function $\Psi$ is not a power series
in $u$ with integer exponents; it is a series in the two monomials
$u^{\alpha},u^{\bar\alpha}$, whose real exponents are $j+k$ but whose
imaginary parts are not zero.  The correct formal setting is a Hahn series in
the group generated by $\alpha$ and $\bar\alpha$, ordered by real part, and
Lagrange inversion in that setting is Good's theorem applied to the two
variables $u^{\alpha}$, $u^{\bar\alpha}$ with the diagonal specialization.
The reordering used to reach \eqref{p9:eq:good} is legitimate because the
series converges absolutely, which the proof establishes independently.
\end{remark}

\section{The real form: log-periodic sectors}
\label{p9:sec:real}

\begin{theorem}[Real sectors]\label{p9:thm:real-sectors}
Put
\begin{equation}\label{p9:eq:xi-theta}
 \xi=\frac{\log R}\lambda=\log_\varphi\bigl(\sqrt5\,y\bigr),
 \qquad
 \theta=\rho\log R=\pi\xi,
 \qquad
 q=R^{-1}=\frac1{5y^{2}} .
\end{equation}
Then
\begin{equation}\label{p9:eq:real-form}
 X_+(y)=\xi+\sum_{n\ge1}D_n(\theta)\,q^{n},
\end{equation}
where each $D_n$ is a real trigonometric polynomial in $\theta$ whose
harmonics all have the same parity as $n$ and whose degree is at most $n$.
The first three are
\begin{align}
 D_1(\theta)&=\frac2\lambda\cos\theta,
 \label{p9:eq:D1}\\
 D_2(\theta)&=-\frac1\lambda\bigl(2+\cos2\theta+2\rho\sin2\theta\bigr),
 \label{p9:eq:D2}\\
 D_3(\theta)&=\frac1\lambda\Bigl[(6-\rho^{2})\cos\theta+5\rho\sin\theta
 +\Bigl(\tfrac23-3\rho^{2}\Bigr)\cos3\theta+3\rho\sin3\theta\Bigr].
 \label{p9:eq:D3}
\end{align}
In particular the leading correction is
\begin{equation}\label{p9:eq:leading}
 X_+(y)=\log_\varphi\bigl(\sqrt5\,y\bigr)
 +\frac{\cos\bigl(\pi\log_\varphi(\sqrt5\,y)\bigr)}{5\,(\log\varphi)\,y^{2}}
 +\bigO(y^{-4}).
\end{equation}
\end{theorem}

\begin{proof}
Group the terms of \eqref{p9:eq:good} by $n=j+k$.  Since
$R^{-s_{jk}}=R^{-n}\e^{\ii\rho(j-k)\log R}=q^{n}\e^{\ii(j-k)\theta}$, the
$n$-th group is $q^{n}$ times a finite combination of $\e^{\ii h\theta}$ with
$h=j-k$ running over $-n,-n+2,\dots,n$; so the harmonics have the parity of
$n$ and degree at most $n$, and by the conjugation of
\cref{p9:thm:good} the combination is real.  This is
\eqref{p9:eq:real-form}.

For $n=1$ the pairs are $(1,0)$ and $(0,1)$ with $s=1\mp\ii\rho$; since
$\binom s1=s$ and $\binom s0=1$, each term is
$(-1)^{2}\lambda^{-1}s^{-1}\cdot s\cdot R^{-s}=\lambda^{-1}R^{-s}$, and the
two add to
$\lambda^{-1}q\bigl(\e^{\ii\theta}+\e^{-\ii\theta}\bigr)
=2\lambda^{-1}q\cos\theta$, which is \eqref{p9:eq:D1}.  Substituting
$q=1/(5y^{2})$, $\lambda=2\log\varphi$ and $\theta=\pi\xi$ gives
\eqref{p9:eq:leading}.  The coefficients \eqref{p9:eq:D2} and
\eqref{p9:eq:D3} follow from the pairs with $j+k=2$ and $j+k=3$ by the same
computation; they were also obtained independently and checked against
\eqref{p9:eq:good} numerically (\cref{p9:rem:verification}).
\end{proof}

\begin{remark}[No limiting normalized coefficient]\label{p9:rem:no-limit}
Every earlier chapter had a distinguished large-parameter limit of its
coefficient family --- the resonant limit of \cref{p6:prop:resonance} and
\cref{p7:prop:limit}, or the deepest pole of \cref{p6:thm:deepest-pole} ---
and in each case that limit summed to a square root.  Here there is no such
limit.  The coefficients $D_n(\theta)$ do not converge as $y\to\infty$:
$\theta=\pi\xi$ increases without bound and $D_n$ is a nonconstant periodic
function of it, so $D_n(\theta)$ oscillates forever between fixed bounds.
This is not a defect of the expansion --- \eqref{p9:eq:real-form} converges,
which none of the others do --- but it does mean that the resummation
technique of \cref{p6:lem:quadratic-core} has no analogue here, and that
statements like ``the $n$-th coefficient tends to \dots'' must be replaced by
statements about the mean or the extremes of $D_n$ over a period.
\end{remark}

\section{Convergence, and a residual certificate}
\label{p9:sec:convergence}

\begin{proposition}[Residual certificate]\label{p9:prop:certificate}
Let $\widehat x\ge2$ and $y\ge1$.  Then
\begin{equation}\label{p9:eq:certificate}
 \bigl|X_+(y)-\widehat x\bigr|
 \le\frac{\bigl|\mathcal{F}(\widehat x)-y\bigr|}
 {\displaystyle\min_{[\,\widehat x\wedge X_+,\ \widehat x\vee X_+]}\mathcal{F}'}
 ,
 \qquad
 \mathcal{F}'(x)\ \ge\
 \frac{\lambda\sinh(\lambda x)-\omega}{5\,\mathcal{F}(x)}>0
 \quad(x>x_{*}).
\end{equation}
In particular, for $x\ge2$ one may use
$\mathcal{F}'(x)\ge\lambda\bigl(\mathcal{F}(x)-\tfrac1{\mathcal{F}(x)}\bigr)/2$
for $\mathcal{F}(x)$ large, so the inverse is well conditioned: a residual of
size $\varepsilon$ certifies an error $\bigO(\varepsilon/(\lambda y))$.
\end{proposition}

\begin{proof}
The first display is \cref{p0:thm:backward-error} applied to
$f=\mathcal{F}$ on the interval joining $\widehat x$ to the root, together
with the lower bound on $\mathcal{F}'$ from \eqref{p9:eq:F-derivative} and
$\sin\ge-1$, valid past $x_{*}$ by \cref{p9:prop:branches}.  For the
conditioning, $5\mathcal{F}^{2}=\e^{\lambda x}+\e^{-\lambda x}-2\cos\omega x$
gives $\e^{\lambda x}=5\mathcal{F}^{2}/2+\bigO(1)$, hence
$\lambda\sinh(\lambda x)=\lambda\e^{\lambda x}/2+\bigO(1)
=\bigO(\lambda\mathcal{F}^{2})$ and
$\mathcal{F}'\asymp\lambda\mathcal{F}$; dividing gives the stated error scale.
\end{proof}

\begin{remark}[What is different here]\label{p9:rem:convergent}
\Cref{p9:thm:good} is a \emph{convergent} expansion of an analytic function,
not a Poincar\'e asymptotic statement.  There is no least term, no
beyond-all-orders floor, and \cref{p0:thm:optimal-truncation} has nothing to
say: one simply takes more terms.  Every other chapter in this volume inverts
a divergent Stirling- or Bernoulli-type series and must stop somewhere.  The
reason for the difference is visible in \cref{p9:thm:factorization}: the
forward function is, after squaring, an \emph{entire} combination of
exponentials, with no asymptotic series anywhere in it.
\end{remark}

\section{Is the log-periodicity intrinsic?}
\label{p9:sec:intrinsic}

The synthesis chapter asks this of every subject with an unexpected feature,
and here it has a clean answer in two halves.

\begin{proposition}[The oscillation is forced; its shape is not]
\label{p9:prop:intrinsic}
Let $\mathcal{G}$ be any function that interpolates the Fibonacci numbers,
$\mathcal{G}(n)=F_n$ for $n\ge0$, and is eventually strictly increasing with
an inverse $X_{\mathcal G}$.  Then
\begin{equation}\label{p9:eq:forced}
 X_{\mathcal G}(F_n)=n
 =\log_\varphi\bigl(\sqrt5\,F_n\bigr)
 +\frac{(-1)^{n}}{\sqrt5\,\varphi^{n}\log\varphi\,F_n}
 \bigl(1+\bigO(\varphi^{-2n})\bigr),
\end{equation}
so along the sequence the displacement from the pure logarithm is
\emph{necessarily} of size $F_n^{-2}$ and \emph{necessarily} alternates in
sign.  By contrast, the values of $X_{\mathcal G}$ between the $F_n$, and
hence the harmonic content of the correction as a function of $\log y$,
depend on $\mathcal G$ and are not determined by the sequence.
\end{proposition}

\begin{proof}
Binet gives $\sqrt5F_n=\varphi^{n}-(-1)^{n}\varphi^{-n}$, so
\[
 \log_\varphi(\sqrt5F_n)
 =n+\log_\varphi\bigl(1-(-1)^{n}\varphi^{-2n}\bigr)
 =n-\frac{(-1)^{n}\varphi^{-2n}}{\log\varphi}
 +\bigO(\varphi^{-4n}).
\]
Rearranging and using $\varphi^{-2n}=\bigl(\sqrt5F_n\varphi^{n}\bigr)^{-1}
\bigl(1+\bigO(\varphi^{-2n})\bigr)$ gives \eqref{p9:eq:forced}.  Since
$\mathcal{G}(n)=F_n$ forces $X_{\mathcal G}(F_n)=n$ for every interpolation,
the displayed relation is a statement about the \emph{sequence}, independent of
$\mathcal G$.  For the second half, let $\chi$ be any nonzero entire function
vanishing on $\Nat_0$ and small enough that $\mathcal{G}+\chi$ is still
eventually increasing; it interpolates the same numbers and its inverse
differs from $X_{\mathcal G}$ between the integers.
\end{proof}

\begin{remark}[Reading the answer]\label{p9:rem:intrinsic-reading}
So the log-periodicity is intrinsic in the only sense that a statement about a
\emph{sequence} can be: the second Binet term $(-1)^{n}\varphi^{-n}$ is
$\varphi^{-2n}$ relative to the first, which is $R^{-1}$, and it alternates.
Every interpolation must reproduce that.  What $\mathcal{F}$ contributes is
the specific choice $\cos(\pi x)$ --- the minimal-growth real interpolation of
$(-1)^{n}$ --- and \emph{that} is what makes the correction a pure cosine at
first order rather than some other function with the same values at integers.
An article that says ``the inverse Fibonacci function is log-periodic'' is
therefore making a statement half about Fibonacci and half about its own
choice of continuation; the volume records both halves.
\end{remark}

\section{The general mechanism: inverting a product of exponentials}
\label{p9:sec:general}

The factorization of \cref{p9:thm:factorization} is an instance of a general
shape, and the closed formula of \cref{p9:thm:good} generalizes with it.

\begin{theorem}[Inverse of a finite exponential product]\label{p9:thm:product}
Let $\lambda>0$, let $\alpha_1,\dots,\alpha_m$ have positive real part, and
let $c_1,\dots,c_m$ and $\beta_1,\dots,\beta_m$ be complex.  Suppose
\begin{equation}\label{p9:eq:product-forward}
 R=\e^{\lambda x}\prod_{r=1}^{m}
 \bigl(1-c_r\e^{-\lambda\alpha_r x}\bigr)^{\beta_r},
\end{equation}
with the principal branch of each factor near $1$.  Then for $R$ large the
solution with $x\to\infty$ is
\begin{equation}\label{p9:eq:product-inverse}
 x=\frac{\log R}\lambda
 -\frac1\lambda\sum_{\bm\jmath\in\Nat^{m}\setminus\{\bm0\}}
 \frac1{s_{\bm\jmath}}
 \prod_{r=1}^{m}(-c_r)^{j_r}\binom{s_{\bm\jmath}\beta_r}{j_r}\,
 R^{-s_{\bm\jmath}},
 \qquad
 s_{\bm\jmath}=\sum_{r=1}^{m}j_r\alpha_r,
\end{equation}
the series converging absolutely for $R$ large.
\end{theorem}

\begin{proof}
Put $u=\e^{-\lambda x}$, so that \eqref{p9:eq:product-forward} reads
$Ru=\Psi(u):=\prod_r(1-c_ru^{\alpha_r})^{\beta_r}$, i.e.\ $u=v\Psi(u)$ with
$v=R^{-1}$ and $\Psi(0)=1$.  As in \cref{p9:thm:good}, Lagrange inversion for
the observable $-\lambda^{-1}\log u$ gives
\[
 x=\frac{\log R}\lambda-\frac1\lambda\sum_{n\ge1}\frac1n
 \Coef{u^{n}}\Psi(u)^{n}\,R^{-n},
\]
and the generalized binomial theorem gives
\[
 \Psi(u)^{n}=\prod_{r}(1-c_ru^{\alpha_r})^{n\beta_r}
 =\sum_{\bm\jmath\in\Nat^{m}}
 \prod_{r}(-c_r)^{j_r}\binom{n\beta_r}{j_r}\;
 u^{\,\sum_rj_r\alpha_r}.
\]
The monomial $u^{s_{\bm\jmath}}$ contributes at $n=s_{\bm\jmath}$, which gives
\eqref{p9:eq:product-inverse}.  Absolute convergence follows as before:
$|R^{-s_{\bm\jmath}}|=R^{-\Re s_{\bm\jmath}}$ with
$\Re s_{\bm\jmath}\ge\delta|\bm\jmath|$ for
$\delta=\min_r\Re\alpha_r>0$, while the binomial products grow at most
geometrically in $|\bm\jmath|$.
\end{proof}

\begin{corollary}[The Fibonacci case]\label{p9:cor:fib-instance}
Taking $m=2$, $\alpha_1=\alpha=1-\ii\rho$, $\alpha_2=\bar\alpha$,
$c_1=c_2=\beta_1=\beta_2=1$ in \cref{p9:thm:product} returns
\eqref{p9:eq:good}, with $s_{(j,k)}=j\alpha+k\bar\alpha=j+k-\ii\rho(j-k)
=s_{jk}$ and $\prod_r(-c_r)^{j_r}=(-1)^{j+k}$.
\end{corollary}

\begin{proof}
Substitute those values into \eqref{p9:eq:product-inverse}.  The exponent is
$s_{(j,k)}=j\alpha+k\bar\alpha=j+k-\ii\rho(j-k)=s_{jk}$; each $\binom{s\beta_r}{j_r}$
becomes $\binom{s_{jk}}{j}$ resp.\ $\binom{s_{jk}}{k}$ because $\beta_r=1$;
and $\prod_r(-c_r)^{j_r}=(-1)^{j+k}$, so the overall sign is
$-(-1)^{j+k}=(-1)^{j+k+1}$, which is the sign in \eqref{p9:eq:good}.
\end{proof}

\begin{remark}[Where the shape comes from, and where it stops]
\label{p9:rem:shape}
A forward function of the form \eqref{p9:eq:product-forward} arises whenever a
dominant exponential is multiplied by finitely many exponentially small
corrections --- which is exactly what a Binet-type formula with finitely many
characteristic roots gives.  Conjugate pairs $\alpha,\bar\alpha$ produce the
log-periodicity, since $R^{-s}$ then carries $\e^{\ii\Im(s)\log R}$.  Two
boundaries are worth naming.  If the correction is an \emph{infinite} product
or a general analytic germ, the closed formula is replaced by a
Lagrange--Good expansion whose coefficients are no longer products of two
binomials; V1 and V3 both develop that case.  And if the carrier is mixed,
$\e^{\lambda x}x^{\mu}$, then $\Wlam$ reappears and one is back in the world of
\cref{p6:sec:core} --- which is the precise point at which this chapter's
subject rejoins the rest of the volume.
\end{remark}

\section{Numerical verification}
\label{p9:sec:numerics}

\begin{remark}[What was checked]\label{p9:rem:verification}
All of the following were recomputed for this volume at $50$ digits, from the
definitions and independently of the sources.  The constants
$\lambda=0.9624236501192068949955\ldots$,
$\rho=3.2642513026364969\ldots$ and
$x_{*}=1.972992734136994\ldots$ of
\eqref{p9:eq:constants} and \eqref{p9:eq:xstar}.  The interpolation
$\mathcal{F}(n)=F_n$ for $0\le n\le10$, to $40$ digits.  The integer
checkpoints $X_+(2)=3$, $X_+(5)=5$ and $X_+(55)=10$, each to $40$ digits ---
these test the branch convention as well as the formula.  The exact
factorization \eqref{p9:eq:factorization}, at $y=3,10,100$, to $35$ digits.
The closed formula \eqref{p9:eq:good}: its imaginary part vanishes to $54$
digits or better, and its truncations converge at the predicted rate.  And the
agreement between \eqref{p9:eq:good} truncated at $j+k\le3$ and the
trigonometric sectors \eqref{p9:eq:D1}--\eqref{p9:eq:D3}, which matched to the
full working precision at $y=100$ and $y=1000$.
\end{remark}

\begin{table}[htbp]
\centering
\caption{Absolute error of \eqref{p9:eq:real-form} truncated after
$D_N$, against the root of $\mathcal{F}(x)=y$ computed to $50$ digits.}
\label{p9:tab:errors}
\small
\begin{tabular}{rcccc}
\toprule
$y$ & $N=0$ (i.e.\ $\xi$) & $N=1$ & $N=2$ & $N=3$\\
\midrule
$10$ & $5.4520\!\times\!10^{-4}$ & $1.1350\!\times\!10^{-5}$
 & $1.5430\!\times\!10^{-7}$ & $2.1470\!\times\!10^{-9}$\\
$100$ & $3.0100\!\times\!10^{-5}$ & $3.5620\!\times\!10^{-9}$
 & $2.9720\!\times\!10^{-13}$ & $5.7170\!\times\!10^{-18}$\\
$1000$ & $4.1410\!\times\!10^{-7}$ & $1.7020\!\times\!10^{-13}$
 & $2.5800\!\times\!10^{-19}$ & $4.2100\!\times\!10^{-25}$\\
\bottomrule
\end{tabular}
\end{table}

Each additional block gains a factor $R^{-1}=1/(5y^{2})$, as
\cref{p9:thm:good} predicts: from $y=100$ to $y=1000$ the $N=3$ error drops by
about six orders of magnitude per block, and $R^{N+1}$ times the error stays
bounded across the table.  Unlike every other error table in this volume, the
columns would continue to improve indefinitely.

\section{Comparison, and what is left open}
\label{p9:sec:open}

\begin{remark}[Where this subject sits]\label{p9:rem:comparison}
Against the seven others: it is the only one whose inverse expansion
converges; the only one with no Lambert core and no slow coefficient field;
the only one whose coefficient family has no large-parameter limit
(\cref{p9:rem:no-limit}); and the only one where the interesting structure ---
oscillation --- comes from a \emph{pair of complex conjugate exponents} rather
than from a logarithm.  What it shares with the rest is the organizing
principle: invert a dominant block exactly, and revert around it.  That is
\cref{p0:thm:perturbed-inversion}, and it is the only piece of Chapter~0 this
chapter genuinely needs.
\end{remark}

\begin{enumerate}
\item \emph{Closed forms for $D_n$.}  \Cref{p9:thm:real-sectors} gives the
first three and a rule for the harmonics, but no closed formula for the
coefficients of $\cos h\theta$ and $\sin h\theta$ inside $D_n$.  They are
finite explicit sums by \eqref{p9:eq:good}; a generating function in
$\theta$ is not derived.
\item \emph{The exact radius.}  The proof of \cref{p9:thm:good} gives
convergence for $R$ larger than some constant depending on $\rho$, obtained
from a crude bound on the binomials.  The sharp radius --- equivalently, the
nearest singularity of $v\mapsto u(v)$, which is where
$\Psi'(u)v=1$ --- is not computed here; V1 states one and we did not verify
it, so no value is quoted.
\item \emph{Inner branches.}  \Cref{p9:rem:inner}: the preimages coming from
the two critical points in $(1,2)$ are Puiseux, not log-periodic, and are not
treated.
\item \emph{Complex branches.}  Replacing $\log R$ by $\log R+2\pi\ii m$ gives
formal sheets, but $\mathcal{F}$ is built from a modulus and is not analytic,
so the continuation question is genuinely different from that of
\cref{p6:sec:open}(iv) and is not addressed.
\item \emph{Other interpolations.}  \Cref{p9:prop:intrinsic} shows the answer
depends on the continuation.  Which interpolation is canonical --- by minimal
growth, by complete monotonicity of some transform, or by agreement with the
signed convention $F_{-n}=(-1)^{n+1}F_n$ rather than with evenness --- is
open, and is the same question left open for the subfactorial in
\cref{p8:sec:open}(iii) and for the double factorial in \cref{p2:sec:top}.
\end{enumerate}

%% ==================== fragment p5 ====================

\part{Two classical sequences by their generating function}

\chapter{The Bell numbers: a Lambert saddle}
\label{q2:sec:bell}

\section{What this chapter does}
\label{q2:sec:bell-intro}

The Bell numbers $\mathfrak{B}_n$, counting set partitions, have exponential
generating function $\exp(\e^{z}-1)$ and a saddle point at
\begin{equation}\label{q2:eq:saddle}
 r\e^{r}=n,\qquad r=\Wlam_0(n).
\end{equation}
That the saddle is a Lambert point is what places this subject in the present
volume: $r$ is the same object the Lambert parts study, and the natural
transmonomial turns out to be $\e^{-r}=r/n$, which is the reciprocal of the
Lambert scale.

The leading Moser--Wyman--de Bruijn approximation is classical.  What the two
filed sources supply, and what is kept here, is the data that turns it into a
usable all-orders formula: a named asymptotic scale, a finite rule computing
every coefficient, the algebraic structure of those coefficients, and a
remainder theorem.

\begin{sourcebox}
Two articles: \texttt{Bell\_Number\_Asymptotic\_Transseries} (\textsf{B1}) and
\texttt{Bell\_Number\_Transseries\_Article} (\textsf{B2}).  They agree; every
coefficient either prints was recomputed for this volume against exact Bell
numbers.  \textsf{B1} carries the finite partition rule and the explicit
$c_1,c_2,c_3$; \textsf{B2} the Stirling interface and the singulant ordering.
Both flag their nonprincipal-saddle sections as formal, and that labelling is
preserved.
\end{sourcebox}

\section{The saddle expansion to all orders}
\label{q2:sec:bell-allorders}

\begin{definition}[Touchard cumulants]\label{q2:def:touchard}
For $j\ge1$ put
\begin{equation}\label{q2:eq:touchard}
 T_j(r)=\e^{-r}(r\partial_r)^{j}\e^{r}
       =\sum_{k=1}^{j}\stirlingsecond{j}{k}r^{k},
 \qquad
 q_j(r)=\frac{T_j(r)}{r\,(1+r)^{j/2}} .
\end{equation}
\end{definition}

\begin{lemma}[Central and complementary arcs]\label{q2:lem:localization}
Write $\Phi_r(\theta)$ for the exponent of the Cauchy integrand on the circle
$|z|=r$, normalized so that $\Phi_r(0)=0$.  There are absolute constants
$\delta,c_1,c_2>0$ such that, for all large $n$ and $r=\Wlam_0(n)$,
\begin{align}
 \Re\Phi_r(\theta)&\le-c_1n(1+r)\theta^{2},
 && |\theta|\le\delta/r,
 \label{q2:eq:central}\\
 \Re\Phi_r(\theta)&\le-c_2\,n/r^{2},
 && \delta/r\le|\theta|\le\pi .
 \label{q2:eq:outer}
\end{align}
Consequently the complementary arc contributes
\begin{equation}\label{q2:eq:outer-small}
 \int_{\delta/r\le|\theta|\le\pi}\e^{\Phi_r(\theta)}\dd\theta
 =\bigO\!\bigl(\e^{-c_2n/r^{2}}\bigr),
\end{equation}
which is smaller than $(r/n)^{M}$ for every fixed $M$.
\end{lemma}

\begin{proof}
On $|z|=r$ the exponent is $\e^{r\cos\theta}\cos(r\sin\theta)-\e^{r}
-\ii n\theta$ up to the terms that cancel at $\theta=0$; its real part is
$\e^{r\cos\theta}\cos(r\sin\theta)-\e^{r}$.  For $|\theta|\le\delta/r$ with
$\delta$ small, expanding to second order gives
$\Re\Phi_r(\theta)=-\tfrac12\e^{r}r(1+r)\theta^{2}\bigl(1+\bigO(\delta)\bigr)$,
and $\e^{r}r=n$, which is \eqref{q2:eq:central}.  For
$\delta/r\le|\theta|\le\pi$ the factor $\cos(r\sin\theta)$ is bounded away
from $1$ except near $\theta=0$, while $\e^{r\cos\theta}\le\e^{r}$; a short
computation --- splitting at $|\theta|=\pi/2$ and using
$\cos\theta\le1-2\theta^{2}/\pi^{2}$ on $[0,\pi/2]$ --- gives
$\Re\Phi_r\le-c_2\e^{r}/r=-c_2n/r^{2}$, which is \eqref{q2:eq:outer}.
Integrating \eqref{q2:eq:outer} over an arc of length at most $2\pi$ gives
\eqref{q2:eq:outer-small}, and $\e^{-c_2n/r^{2}}$ beats every power of
$r/n=\e^{-r}$ because $n/r^{2}=\e^{r}/r$ grows faster than any multiple
of $r$.
\end{proof}

\begin{theorem}[All-orders saddle expansion]\label{q2:thm:bell}
Let $r=\Wlam_0(n)$.  For every fixed $M\ge1$, as $n\to\infty$,
\begin{equation}\label{q2:eq:bell-main}
 \mathfrak{B}_n
 =\frac{n!\,\exp(\e^{r}-1)}{r^{n}\sqrt{2\pi n(1+r)}}
 \left(\sum_{m=0}^{M-1}\frac{c_m(r)}{n^{m}}
 +\bigO_M\!\bigl((r/n)^{M}\bigr)\right),
\end{equation}
where $c_0=1$ and, for $m\ge1$,
\begin{equation}\label{q2:eq:cm-partition}
 c_m(r)=\!\!\sum_{\substack{k_1,\dots,k_{2m}\ge0\\ \sum_{s}s\,k_s=2m}}\!\!
 (-1)^{m+K}(2m+2K-1)!!
 \prod_{s=1}^{2m}\frac{q_{s+2}(r)^{k_s}}{\bigl((s+2)!\bigr)^{k_s}k_s!},
 \qquad K=\sum_{s}k_s .
\end{equation}
Equivalently, $c_m$ is the Gaussian average of a complete exponential Bell
polynomial in the $q_j$.
\end{theorem}

\begin{proof}
The Cauchy integral
$\mathfrak{B}_n/n!=(2\pi\ii)^{-1}\oint\exp(\e^{z}-1)z^{-n-1}\dd z$ is put on
the circle $|z|=r$ with $r$ the saddle \eqref{q2:eq:saddle}; writing
$z=r\e^{\ii\theta/\sqrt{n(1+r)}}$ normalizes the quadratic term, and the
higher Taylor coefficients of the exponent are exactly the $q_j$ of
\eqref{q2:eq:touchard}.  Expanding $\exp(\text{cubic and higher})$ and
integrating the resulting Gaussian moments term by term gives
\eqref{q2:eq:cm-partition}, the double factorial being the Gaussian moment
$\mathbb{E}Z^{2m+2K}=(2m+2K-1)!!$.

Two steps make this an asymptotic statement rather than a formal one, and
both are \cref{q2:lem:localization}.  First, \eqref{q2:eq:outer-small} discards
the complementary arc at a cost smaller than every retained term, so the
integral may be replaced by the central one.  Second, \eqref{q2:eq:central}
bounds the central integrand by a Gaussian, which justifies integrating the
expansion term by term and bounds the tail after $M$ terms by the first
omitted Gaussian moment; that is the $\bigO_M\bigl((r/n)^{M}\bigr)$ of
\eqref{q2:eq:bell-main}.  What \eqref{q2:eq:cm-partition} adds to the classical
mechanism is that every coefficient is a \emph{finite} explicit sum, which is
what \cref{q2:prop:bell-structure} then exploits.
\end{proof}

\begin{proposition}[Structure of the coefficients]\label{q2:prop:bell-structure}
For every $m\ge1$,
\begin{equation}\label{q2:eq:bell-structure}
 c_m(r)=\frac{P_m(r)}{(1+r)^{3m}},
 \qquad P_m\in\Rat[r],
 \qquad \deg P_m\le4m,
 \qquad c_m(r)=\bigO(r^{m})\ \ (r\to\infty).
\end{equation}
The first three are
\begin{align}
 c_1(r)&=-\frac{2r^{4}+9r^{3}+16r^{2}+6r+2}{24\,(1+r)^{3}},
 \label{q2:eq:c1}\\
 c_2(r)&=\frac{4r^{8}-12r^{7}-143r^{6}-384r^{5}-588r^{4}-636r^{3}
 +100r^{2}+24r+4}{1152\,(1+r)^{6}},
 \label{q2:eq:c2}\\
 c_3(r)&=\frac{\begin{aligned}[t]
 &1112r^{12}+12708r^{11}+56082r^{10}+118683r^{9}+97176r^{8}-74394r^{7}\\
 &\quad-259046r^{6}-456444r^{5}-635424r^{4}+115308r^{3}+39432r^{2}
 +10008r+1112\end{aligned}}
 {414720\,(1+r)^{9}} .
 \label{q2:eq:c3}
\end{align}
\end{proposition}

\begin{proof}
By \eqref{q2:eq:touchard}, $T_j=\sum_{k=1}^{j}\stirlingsecond jk r^{k}$ is a
polynomial of degree $j$ with $T_j(0)=0$; write $T_j=r\,\widetilde T_j$ with
$\widetilde T_j\in\Int[r]$ of degree $j-1$.  Fix a term of
\eqref{q2:eq:cm-partition}, indexed by $(k_1,\dots,k_{2m})$ with
$\sum_s sk_s=2m$ and $K=\sum_sk_s$.  Then
\[
 \prod_{s}q_{s+2}^{k_s}
 =\frac{\prod_s T_{s+2}^{k_s}}
        {r^{K}\,(1+r)^{\sum_sk_s(s+2)/2}}
 =\frac{\prod_s\widetilde T_{s+2}^{\,k_s}}{(1+r)^{m+K}},
\]
because $\prod_sT_{s+2}^{k_s}=r^{K}\prod_s\widetilde T_{s+2}^{k_s}$ --- the
factor $r^{K}$ cancels exactly --- and
$\sum_sk_s(s+2)/2=\bigl(\sum_ssk_s+2K\bigr)/2=m+K$.  The surviving numerator
has degree $\sum_sk_s(s+1)=2m+K$.

Now $s\ge1$ forces $K\le\sum_ssk_s=2m$, so $(1+r)^{m+K}$ divides
$(1+r)^{3m}$; and on putting the term over the common denominator
$(1+r)^{3m}$ its numerator acquires the factor $(1+r)^{2m-K}$, reaching
degree exactly $(2m+K)+(2m-K)=4m$, independently of the index.  Summing over
the finitely many indices gives \eqref{q2:eq:bell-structure} with
$\deg P_m\le4m$.  Since the denominator has degree $3m$, this also gives
$c_m=\bigO(r^{4m-3m})=\bigO(r^{m})$.  The three displayed coefficients follow
from \eqref{q2:eq:cm-partition} at $m=1,2,3$; they were recomputed and checked
against exact Bell numbers (\cref{q2:rem:verification}).
\end{proof}

\begin{remark}[Why the Poincar\'e form hides the right scale]
\label{q2:rem:bell-scale}
Because $c_m(r)=\bigO(r^{m})$, the $m$-th term of \eqref{q2:eq:bell-main} is
of size $(r/n)^{m}$, not $n^{-m}$.  The genuine transmonomial is therefore
\begin{equation}\label{q2:eq:bell-transmonomial}
 \frac rn=\e^{-r},
\end{equation}
and rewriting the expansion in powers of $\e^{-r}$ --- which requires
absorbing the Stirling series for $n!$ --- turns every coefficient into a
\emph{bounded} rational function of $r$.  Then
\begin{equation}\label{q2:eq:bell-log}
 \log\mathfrak{B}_n
 \sim n\Bigl(r-1+\frac1r\Bigr)-1-\frac12\log(1+r)
 +\sum_{m\ge1}\mathcal{D}_m(r)\,\e^{-mr},
\end{equation}
with each $\mathcal{D}_m$ bounded and given by finite coefficient operations.
This is the same distinction the volume draws repeatedly: an expansion may be
correct and still be written on the wrong scale, in which case its
coefficients grow and its ordering is misleading.
\end{remark}

\begin{remark}[Three ranked sectors]\label{q2:rem:bell-sectors}
Substituting the iterated-logarithm expansion of $r=\Wlam_0(n)$, with
$L=\log n$ and $\lambda=\log L$, splits $\log\mathfrak{B}_n$ into
\[
 nL\,\Rat[\lambda][[L^{-1}]],
 \qquad
 \log L+\Rat[\lambda][[L^{-1}]],
 \qquad
 \sum_{m\ge1}n^{-m}L^{m}\Rat[\lambda][[L^{-1}]],
\]
the third family lying beyond every power of $L^{-1}$ relative to the first
two.  Nonprincipal complex saddles lie beyond all three.  Both sources label
their treatment of those saddles \emph{formal}, because identifying the Stokes
multipliers of the original Cauchy contour is a separate global
steepest-descent problem; that labelling is kept here, and the sectors are not
claimed as asymptotics.
\end{remark}

\chapter{The Fubini numbers: an exact pole lattice}
\label{q2:sec:fubini}

\section{What this chapter does}
\label{q2:sec:fubini-intro}

The Fubini (ordered Bell) numbers
$\mathfrak{F}_n=\sum_j j!\stirlingsecond{n}{j}$ count weak orderings.  Their
exponential generating function is
\begin{equation}\label{q2:eq:fubini-egf}
 \sum_{n\ge0}\mathfrak{F}_n\frac{z^{n}}{n!}=\frac1{2-\e^{z}},
\end{equation}
which is meromorphic with \emph{simple} poles at every solution of
$\e^{z}=2$, that is at
\begin{equation}\label{q2:eq:lattice}
 \zeta_k=\rho+2\pi\ii k,\qquad \rho=\log2,\qquad k\in\Int .
\end{equation}
This is the one subject in the volume whose transseries is exact and
convergent rather than asymptotic --- as is, for a different reason, the
Fibonacci inverse of \cref{p9:sec:top}.  The contrast with the Bell numbers,
whose saddle expansion diverges, is the point of putting the two chapters
side by side, and \cref{q2:sec:compare} draws it.

\begin{sourcebox}
Three articles: \texttt{Fubini\_Number\_Full\_Transseries},
\texttt{Fubini\_Number\_Transseries} and
\texttt{Fubini\_Number\_Transseries\_Article}.  They agree.  Their common
content is the pole-lattice formula and the factorial normalization; the first
adds the meromorphic-generating-function principle and the higher-order-pole
analysis, which is kept because it is what explains the shape of the answer.
\end{sourcebox}

\section{The exact formula}
\label{q2:sec:fubini-exact}

\begin{theorem}[Exact pole-lattice transseries]\label{q2:thm:fubini}
For every integer $n\ge1$,
\begin{equation}\label{q2:eq:fubini-exact}
 \mathfrak{F}_n
 =\frac{n!}2\sum_{k\in\Int}\zeta_k^{-n-1}
 =\frac{n!}{2\rho^{\,n+1}}
 \left[1+2\sum_{k\ge1}\left(\frac{\rho}{|\zeta_k|}\right)^{n+1}
 \cos\bigl((n+1)\arg\zeta_k\bigr)\right],
\end{equation}
both series converging absolutely.  This is an identity, not an asymptotic
statement.
\end{theorem}

\begin{proof}
Every zero of $2-\e^{z}$ is simple, since $(2-\e^{z})'=-\e^{z}=-2\ne0$ there,
and the zeros are exactly the $\zeta_k$ of \eqref{q2:eq:lattice}.  The residue
of $(2-\e^{z})^{-1}$ at $\zeta_k$ is $1/(-\e^{\zeta_k})=-\tfrac12$.  The
Mittag--Leffler expansion of the meromorphic function
$(2-\e^{z})^{-1}$, which converges because the poles have spacing $2\pi$ in
the imaginary direction and the function is bounded on the intervening
horizontal strips, therefore gives
\[
 \frac1{2-\e^{z}}=-\frac12\sum_{k\in\Int}\frac1{z-\zeta_k}+\text{entire},
\]
and the entire part vanishes by the decay of both sides as
$\Re z\to-\infty$.  Extracting $\Coef{z^{n}/n!}$ from
$-\tfrac12(z-\zeta_k)^{-1}=\tfrac12\sum_{m\ge0}z^{m}\zeta_k^{-m-1}$ gives
$\mathfrak{F}_n/n!=\tfrac12\sum_k\zeta_k^{-n-1}$, which is the first equality;
absolute convergence holds for $n\ge1$ because $|\zeta_k|\sim2\pi|k|$ and
$n+1\ge2$.  Pairing $k$ with $-k$, and using
$\zeta_{-k}=\overline{\zeta_k}$, turns the sum into the real cosine form.
\end{proof}

\begin{remark}[What a simple pole does and does not contribute]
\label{q2:rem:simple-pole}
A simple pole contributes \emph{one} exact sector with a \emph{constant}
amplitude.  It does not generate an independent $1/n$ fluctuation series.  In
the elementary normalization the inverse powers of $n$ appear solely because
$n!$ has been replaced by Stirling's expansion; keep $n!$ exact and there are
none.  This is worth stating because the opposite is easy to assume: a
singularity-analysis habit expects each singularity to carry its own algebraic
series.  Higher-order poles are the case where that expectation is correct,
and there the amplitudes are polynomial in $n$, governed by generalized
Bernoulli numbers.
\end{remark}

\begin{corollary}[Elementary-scale form]\label{q2:cor:fubini-elementary}
Let $\Lambda_k=\operatorname{Log}(\zeta_k/\rho)$ with the principal logarithm,
and let $c_j$ be defined by the Stirling series of the scaled gamma factor,
\begin{equation}\label{q2:eq:stirling-cj}
 \exp\left(\sum_{r\ge1}\frac{B_{r+1}}{r(r+1)}t^{r}\right)=\sum_{j\ge0}c_jt^{j},
 \qquad
 c_0=1,\ c_1=\tfrac1{12},\ c_2=\tfrac1{288},\ c_3=-\tfrac{139}{51840}.
\end{equation}
Then
\begin{equation}\label{q2:eq:fubini-elementary}
 \mathfrak{F}_n\ \sim\
 \frac{\sqrt{2\pi n}}{2\rho}\left(\frac n{\e\rho}\right)^{n}
 \sum_{k\in\Int}\e^{-n\Lambda_k}\sum_{j\ge0}\frac{C_{j,k}}{n^{j}},
 \qquad
 C_{j,k}=c_j\,\frac{\rho}{\zeta_k},
\end{equation}
where the $k$-sum converges for every $n\ge1$ and the $j$-sum is the Poincar\'e
expansion of the gamma factor.
\end{corollary}

\begin{proof}
Substitute $n!=\sqrt{2\pi n}\,(n/\e)^{n}\sum_j c_jn^{-j}$, which is
\eqref{q2:eq:stirling-cj} exponentiated, into \eqref{q2:eq:fubini-exact}, and
write $\zeta_k^{-n-1}=\rho^{-n-1}\e^{-n\Lambda_k}\cdot\rho/\zeta_k$.  The
amplitude $C_{j,k}$ is then the product of the $j$-th Stirling coefficient and
the constant $\rho/\zeta_k$ of \cref{q2:rem:simple-pole}; no $k$-dependence
enters the $j$-sum, which is why the same series multiplies every sector.
\end{proof}

\section{Certified truncation}
\label{q2:sec:fubini-truncation}

The lattice sum converges, so truncating it is a numerical choice rather than
an asymptotic one, and the cost can be bounded exactly.  Write
\begin{equation}\label{q2:eq:Tn}
 T_n=\sum_{k\in\Int}u_k^{\,n+1},
 \qquad
 u_k=\frac{\rho}{\zeta_k},
 \qquad
 T_{n,K}=\sum_{|k|\le K}u_k^{\,n+1},
\end{equation}
so that $\mathfrak{F}_n=\tfrac12 n!\,\rho^{-n-1}T_n$ by
\cref{q2:thm:fubini}.

\begin{theorem}[Pole-tail bounds]\label{q2:thm:pole-tail}
For $n\ge1$ and $K\ge0$,
\begin{equation}\label{q2:eq:tail}
 |T_n-T_{n,K}|\ \le\ 2\sum_{k>K}|u_k|^{\,n+1}
 \ \le\ 2\Bigl(\frac{\rho}{2\pi}\Bigr)^{n+1}\zeta(n+1,K+1),
\end{equation}
$\zeta(s,a)$ being the Hurwitz zeta function, and for $K\ge1$ also
\begin{equation}\label{q2:eq:tail-integral}
 |T_n-T_{n,K}|\ \le\ \frac2n\Bigl(\frac{\rho}{2\pi}\Bigr)^{n+1}K^{-n}.
\end{equation}
For the leading approximation $K=0$ one has the uniform estimate
\begin{equation}\label{q2:eq:uniform}
 |T_n-1|\ \le\ \frac{\rho^{2}}{12}=0.04003775115985011\ldots
 \qquad(n\ge1).
\end{equation}
\end{theorem}

\begin{proof}
$|\zeta_k|^{2}=\rho^{2}+4\pi^{2}k^{2}\ge4\pi^{2}k^{2}$, so
$|u_k|\le\rho/(2\pi|k|)$, and pairing $k$ with $-k$ gives the first inequality
in \eqref{q2:eq:tail}; summing $|u_k|^{n+1}\le(\rho/2\pi)^{n+1}k^{-n-1}$ over
$k>K$ gives the second, by the definition of the Hurwitz zeta function.  For
\eqref{q2:eq:tail-integral}, compare that sum with
$\int_K^{\infty}t^{-n-1}\dd t=K^{-n}/n$.  For \eqref{q2:eq:uniform}, take
$K=0$ in \eqref{q2:eq:tail}: $|T_n-1|\le2(\rho/2\pi)^{n+1}\zeta(n+1)$, which
decreases in $n$, so its value at $n=1$ bounds it; that value is
$2(\rho/2\pi)^{2}\zeta(2)=2\cdot\rho^{2}/(4\pi^{2})\cdot\pi^{2}/6
=\rho^{2}/12$.
\end{proof}

\begin{proposition}[A sufficient linear pole budget]\label{q2:prop:budget}
Fix $\delta>0$ and choose $K=K_n$ with
\begin{equation}\label{q2:eq:budget}
 K_n+1\ \ge\ \frac{(1+\delta)\,n}{2\pi\e}.
\end{equation}
Then the bound \eqref{q2:eq:tail-integral}, transported to $\mathfrak{F}_n$, is
$\bigO_\delta\bigl(n^{-1/2}(1+\delta)^{-n}\bigr)$ relative to
$\mathfrak{F}_n$.  In particular a number of poles \emph{linear} in $n$
suffices to determine $\mathfrak{F}_n$ to within less than $\tfrac12$, so the
integer is recovered by rounding.
\end{proposition}

\begin{proof}
By \eqref{q2:eq:tail-integral} the relative error is at most
$\tfrac2n(\rho/2\pi K_n)^{n}\rho/(2\pi)$ divided by $T_n\ge1-\rho^{2}/12$,
which by \eqref{q2:eq:budget} is at most a constant times
$n^{-1}\bigl(\e/(1+\delta)\bigr)^{n}\rho/(2\pi)\cdot(\rho/2\pi)^{n}
\cdot(2\pi)^{n}$; collecting the powers and using
$\rho\e<2\pi$ leaves $\bigO_\delta((1+\delta)^{-n})$, and the elementary
Stirling upper bound $n!\le C\sqrt n\,(n/\e)^{n}$ contributes the remaining
$n^{-1/2}$.
\end{proof}

\section{Weighted Fubini polynomials}
\label{q2:sec:weighted}

Everything above is the case $x=1$ of a one-parameter family, and the
parameter costs nothing.

\begin{theorem}[Full weighted pole formula]\label{q2:thm:weighted}
For $x>0$ let $\mathfrak{F}_n(x)=\sum_j j!\stirlingsecond nj x^{j}$, whose
exponential generating function is $\bigl(1-x(\e^{z}-1)\bigr)^{-1}$.  Put
\begin{equation}\label{q2:eq:rhox}
 \rho_x=\log\Bigl(1+\frac1x\Bigr),
 \qquad
 \zeta_k(x)=\rho_x+2\pi\ii k .
\end{equation}
Then, for every $n\ge1$,
\begin{equation}\label{q2:eq:weighted}
 \mathfrak{F}_n(x)=\frac{n!}{1+x}\sum_{k\in\Int}\zeta_k(x)^{-n-1},
\end{equation}
an exact convergent identity, which at $x=1$ is \cref{q2:thm:fubini}.
\end{theorem}

\begin{proof}
The poles of $\bigl(1-x(\e^{z}-1)\bigr)^{-1}$ are the solutions of
$\e^{z}=1+1/x$, that is the $\zeta_k(x)$ of \eqref{q2:eq:rhox}, and each is
simple because the derivative of the denominator is $-x\e^{z}=-(1+x)\ne0$
there; the residue is therefore $-1/(1+x)$, constant in $k$.  The
Mittag--Leffler argument of \cref{q2:thm:fubini} applies verbatim, with
$\tfrac12$ replaced by $1/(1+x)$, and gives \eqref{q2:eq:weighted}.  At $x=1$,
$\rho_1=\log2=\rho$ and $1/(1+x)=\tfrac12$.
\end{proof}

\begin{remark}[What the parameter shows]\label{q2:rem:weighted}
The residue is constant in $k$ for every $x$, which is
\cref{q2:rem:simple-pole} again: a simple pole contributes one
constant-amplitude sector whatever the weight.  What the weight moves is the
\emph{location} of the lattice, through $\rho_x=\log(1+1/x)$; as $x\to\infty$
the lattice spacing stays $2\pi$ but $\rho_x\to0$, so the sectors become
comparable in size and the leading-pole approximation degrades.  The
uniformity of \eqref{q2:eq:uniform} in $n$ is therefore \emph{not} uniform in
$x$, and a statement about $\mathfrak{F}_n(x)$ on a compact $x$-interval needs
the bound of \cref{q2:thm:pole-tail} with $\rho$ replaced by
$\min_x\rho_x$.
\end{remark}

\section{Two sequences, two kinds of expansion}
\label{q2:sec:compare}

\begin{remark}[The contrast]\label{q2:rem:contrast}
The Bell and Fubini numbers are both counted by set partitions and both have
an entire or meromorphic exponential generating function, yet their
expansions are of opposite type.  $\exp(\e^{z}-1)$ is entire, so there is no
singularity to expand about and the answer comes from a \emph{saddle}; saddle
expansions are asymptotic and the Bell series diverges.  $(2-\e^{z})^{-1}$ is
meromorphic, so the answer comes from its \emph{poles}; the pole sum is exact
and converges.  The dividing line is not the combinatorics but the location of
the singularities of the generating function --- which is the same principle
that makes \cref{p9:sec:top} the one convergent inversion in this volume.
\end{remark}

\begin{proposition}[One coefficient family serves both]
\label{q2:prop:same-coefficients}
Let $\gamma_m=\lim_{r\to\infty}c_m(r)/r^{m}$ be the leading large-saddle
coefficients of the Bell expansion \eqref{q2:eq:bell-main}, and let $c_j$ be
the Stirling coefficients \eqref{q2:eq:stirling-cj} that multiply every
Fubini sector in \eqref{q2:eq:fubini-elementary}.  Then, for $m=1,2,3$,
\begin{equation}\label{q2:eq:gamma-c}
 \gamma_m=(-1)^{m}c_m,
\end{equation}
that is
$\gamma_1=-\tfrac1{12}$, $\gamma_2=\tfrac1{288}$,
$\gamma_3=\tfrac{139}{51840}$ against
$c_1=\tfrac1{12}$, $c_2=\tfrac1{288}$, $c_3=-\tfrac{139}{51840}$.
\end{proposition}

\begin{proof}
Take $r\to\infty$ in \eqref{q2:eq:c1}--\eqref{q2:eq:c3}: the numerator's
leading term over $(1+r)^{3m}$ gives
$\gamma_1=-2/24=-\tfrac1{12}$, $\gamma_2=4/1152=\tfrac1{288}$ and
$\gamma_3=1112/414720=\tfrac{139}{51840}$.  Comparison with the expansion of
\eqref{q2:eq:stirling-cj} gives \eqref{q2:eq:gamma-c}.
\end{proof}

\begin{remark}[Why that is not a coincidence, and what is not proved]
\label{q2:rem:gamma-c}
Both families come from the same place: the gamma factor.  In the Bell
expansion the factor $n!$ sits in front of \eqref{q2:eq:bell-main} and its
Stirling series is what \cref{q2:rem:bell-scale} absorbs when passing to the
$\e^{-r}$ scale; in the Fubini expansion the same $n!$ is what
\cref{q2:cor:fubini-elementary} replaces.  The sign alternation is the
inversion of the scale --- $r/n$ against $1/n$.  Neither source states
\eqref{q2:eq:gamma-c}, because neither treats both sequences.  What is
established here is the identity for $m\le3$, verified in exact arithmetic;
the general statement $\gamma_m=(-1)^{m}c_m$ is \emph{not} proved, and is
recorded as a conjecture in \cref{q2:rem:open}.
\end{remark}

\begin{remark}[Independent verification]\label{q2:rem:verification}
Recomputed for this volume, independently of the sources.  Fubini: the
sequence $1,1,3,13,75,541,4683,47293,545835,\dots$ from
$\sum_jj!\stirlingsecond nj$; the exact formula \eqref{q2:eq:fubini-exact}
against those integers at $n=2,3,5,8,12,20$, agreeing to relative error
$5\times10^{-15}$ at $n=2$ and below $10^{-40}$ from $n=8$ on, the residual
being the truncation of the lattice sum and not an error in the identity --- at
$n=1$ the tail of the $k$-sum is only $\bigO(K^{-1})$ and a truncation at
$K=4000$ leaves $6\times10^{-6}$, exactly the observed discrepancy; the cosine
form against the same integers; and the coefficients $c_0=1$, $c_1=1/12$,
$c_2=1/288$, $c_3=-139/51840$ from \eqref{q2:eq:stirling-cj}.  Bell: the
weighted formula \eqref{q2:eq:weighted} against
$\sum_jj!\stirlingsecond njx^{j}$ at $x=\tfrac12,1,2,3$ and $n=3,6,10$, exact
to $10^{-17}$ or better in every one of the twelve cases; the constant
$\rho^{2}/12=0.04003775115985011\ldots$ of \eqref{q2:eq:uniform}; and the
expansion \eqref{q2:eq:bell-main} against exact Bell numbers at
$n=20,50,200,1000$, the relative error falling from $1.5\times10^{-2}$ to
$7.6\times10^{-7}$ at $n=20$ and from $5.6\times10^{-4}$ to
$1.5\times10^{-12}$ at $n=1000$ as $c_1,c_2,c_3$ are included; the structural
claims \eqref{q2:eq:bell-structure} for $m\le3$, degree and denominator
exactly as stated; and \eqref{q2:eq:gamma-c}.
\end{remark}

\begin{remark}[What is left open]\label{q2:rem:open}
\begin{enumerate}
\item \emph{The identity $\gamma_m=(-1)^{m}c_m$ for all $m$.}  Verified for
$m\le3$ (\cref{q2:prop:same-coefficients}).  A proof should come from tracking
the gamma factor through \eqref{q2:eq:cm-partition} rather than from computing
both sides.
\item \emph{Nonprincipal Bell saddles.}  Formal, as both sources say; the
Stokes multipliers of the Cauchy contour are not determined.
\item \emph{Higher-order Fubini poles.}  The meromorphic principle of
\cref{q2:rem:simple-pole} predicts polynomial-in-$n$ amplitudes governed by
generalized Bernoulli numbers; the general statement is quoted from the source
and not reproved here.
\end{enumerate}
\end{remark}


\part{Synthesis}

\chapter{Synthesis: what the eight subjects share, and where they part}
\label{p5:sec:top}

\begin{sourcebox}
This chapter belongs to the consolidation rather than to any source.  It
compares the eight inversions carried out in
\cref{p1:sec:top,p2:sec:top,p3:sec:top,p4:sec:top,p6:sec:top,p7:sec:top,%
p8:sec:top,p9:sec:top}, states precisely how much of each is an instance of
\cref{p0:sec:top}, and records what the twenty-four articles left open.  Nothing here is asserted beyond what the earlier chapters
prove.
\end{sourcebox}

\section{One inversion, four dictionaries}\label{p5:sec:one-inversion}

The four subjects of this volume were written about independently, and their
\emph{forward} analyses have almost nothing in common: a singularity analysis
of a Pólya functional equation, a pair of gamma-quotient branches, a
convergent Rademacher series over root-of-unity sectors, and a second, quite
different pair of gamma quotients.  Their \emph{inverse} analyses are, to a
degree that the individual articles could not see, the same computation.

The reason is \cref{p0:def:core}.  In each subject one first isolates a
\emph{dominant core} $\Phi_0$ --- an increasing function whose inverse is
available in closed form --- and writes the exact equation for the index $x$
in terms of the observed value $y$ as
\[
  \Phi_0(x)+\text{(a correction that is small relative to $\Phi_0'$)}=L(y).
\]
The core is then inverted exactly by the Lambert $\Wlam$ function, and the
correction is removed to all orders by one triangular recurrence,
\cref{p0:thm:core-reversion} (in the exponential--power case,
\cref{p0:thm:lambert-centered}).  The recurrence does not know which sequence
it came from: it sees only the core data and the coefficients $q_j$ of the
correction.

\Cref{p0:tab:dictionary} records the resulting dictionary.  Three points in it
deserve emphasis, because each was invisible from inside a single subject.

\paragraph{The partition numbers are the same problem after one substitution.}
Of the four, the partition numbers look least like the others: the
Hardy--Ramanujan phase is $\pi\sqrt{2N/3}$, a \emph{square root} of the
variable, where the other three carry a phase linear in the index or in
$x\log x$.  The three partition articles accordingly build their inversion by
hand around the equation they write as $\rho-2\log\rho=L$.  By
\cref{p0:lem:alpha-reduction} that equation \emph{is} the linear-phase core
$a\xi+b\log\xi$ in the variable $\xi=\sqrt{N}$, with $b/\alpha=-2$; the
substitution costs nothing and removes the apparent exception.  What the
substitution does \emph{not} do is commute with the arithmetic shift, and
\cref{p0:rem:affine-shift} shows why the order matters: re-expanding
$N=n-\tfrac1{24}$ in powers of $n$ after the reduction manufactures an entire
second tower of exponents $n^{\alpha-k}$, which is exactly the infinite
algebraic array that one source reports and attributes, correctly, to the
shift alone.

\paragraph{The two ``factorials'' are not the same kind of object.}
It is tempting --- and this volume's own brief made the mistake --- to expect
the double factorial and the swing factorial to sit side by side.  They do
not.  The swing factorial $n!/\lfloor n/2\rfloor!^{2}$ has its factorials
cancel to leave growth $2^{n}n^{\pm1/2}$, so it is an exponential--power model
with $a=\log2$ (\cref{p4:sec:top}).  The double factorial does not: its
logarithm is $\tfrac{s}{2}(\log s-1)+\bigO(\log s)$ with $s=n+1$, a phase of
factorial type, and since $s\log s/s^{\alpha}\to\infty$ for every
$\alpha\le1$, no choice of model data reproduces it
(\cref{p2:sec:top}).  What survives the difference is exactly the axiomatized
core: both $aX+b\log X$ and $X(\kappa\log X+d)$ are admissible in the sense of
\cref{p0:def:core}, both are inverted exactly by $\Wlam$
(\cref{p0:thm:lambert-core,p0:prop:factorial-core}), and the reversion above
them is one theorem, \cref{p0:thm:core-reversion}.  The right statement is
therefore not ``four instances of one model'' but ``four instances of one
\emph{core calculus}, of which three share a model''.

\paragraph{The Lambert branch is determined by a sign.}
Each of the twelve articles verifies by hand, for its own case, which branch
of $\Wlam$ produces the large root; several remark that the appearance of
$\Wlam_{-1}$ is essential and surprising.  With the signed convention of
\cref{p0:rem:conventions} it is neither: by \cref{p0:cor:branch-rule} the
branch is $\Wlam_0$ when $b>0$ and $\Wlam_{-1}$ when $b<0$, with the fold at
$X_\ast=\lvert b\rvert/a$ and the threshold
$L_c=\lvert b\rvert(1-\log(\lvert b\rvert/a))$.  Reading
\cref{p0:tab:dictionary} down the $b$ column now predicts the branch in every
chapter, including the two branches of \cref{p4:sec:top}, which differ only in
the sign of $\sigma$.

\section{Where the common frame stops}\label{p5:sec:where-it-stops}

The shared apparatus covers the passage from an exact forward description to
an all-orders inverse.  It covers neither end of that passage, and it is worth
being explicit about which parts of the four chapters are irreducibly their
own.

\begin{enumerate}[label=(\alph*),leftmargin=*]
\item \emph{Getting the forward description at all.}  The four forward
      analyses share no method.  \Cref{p1:sec:top} needs a singularity
      analysis of $T(z)=z\exp\sum_{k\ge1}T(z^k)/k$, the location and nature of
      the Otter singularity, and a transfer theorem;
      \cref{p2:sec:top,p4:sec:top} need Binet's representation and Stirling's
      series for gamma quotients; \cref{p3:sec:top} needs the circle method,
      or rather its exact endpoint, Rademacher's convergent series.  Nothing
      in \cref{p0:sec:top} produces any of these, and nothing in one of them
      helps with another.
\item \emph{The arithmetic layer.}  The sectors of \cref{p3:sec:top} are
      indexed by a denominator $k$ and carry Kloosterman-type amplitudes built
      from Dedekind sums.  This structure has no analogue in the other three
      subjects, whose exponentially small corrections come from a competing
      singularity (\cref{p1:sec:top}) or from Borel geometry
      (\cref{p2:sec:top,p4:sec:top}).  The frame transports each sector
      through the inversion once the sector is given; it does not explain
      where the sectors come from.
\item \emph{Parity.}  Two of the four sequences are not the restriction of any
      single smooth interpolation, and in both cases this is forced, not
      accidental: the even and odd subsequences have different exponents
      ($n^{\pm1/2}$ for the swing factorial) or different gamma
      representations (the double factorial).  \Cref{p0:def:three-inverses}
      and \cref{p0:thm:staircase}(4) handle the consequence uniformly --- the
      integer inverse is a minimum over residue classes --- but the
      \emph{fact} of the split is a property of the sequence.
\item \emph{Which interpolation.}  For a discrete sequence the asymptotic
      inverse is well defined only after an interpolation is chosen, and the
      choice is not innocent.  \Cref{p0:rem:three-differ} separates the three
      inverse objects; \cref{p2:sec:top} shows what changes when the OEIS
      periodic interpolation is used instead of the branchwise gamma
      continuation, namely a log-periodic oscillation in the inverse that is
      absent from either branch.  One partition source records the sharper
      version of the same phenomenon: the algebraic part of the inverse is
      interpolation-independent, the nonperturbative part is not.
\end{enumerate}

\section{The three inverse objects, subject by subject}
\label{p5:sec:three-objects}

\Cref{p0:def:three-inverses} distinguishes the staircase inverse $N_\ast$, the
inverse of a chosen interpolation, and the asymptotic inverse.  The four
chapters do not stand in the same relation to them, and the differences are
not cosmetic.

\begin{center}
\small
\begin{tabular}{@{}lccc@{}}
\toprule
Chapter & interpolation available & parity split & exact staircase recovery \\
\midrule
\ref{p1:sec:top} rooted trees & from the singular expansion & no & along the range \\
\ref{p2:sec:top} double factorial & two branches, or OEIS & yes & per residue class \\
\ref{p3:sec:top} partitions & from Rademacher & no & along the range \\
\ref{p4:sec:top} swing factorial & two branches & yes & per residue class \\
\bottomrule
\end{tabular}
\end{center}

Two cautions apply to the whole table.  First, by \cref{p0:thm:staircase}(5)
the staircase $N_\ast$ is \emph{not} uniformly approximable by any continuous
function: $\sup_{y\ge Y}\lvert N_\ast(y)-g(y)\rvert\ge\tfrac12$ for every
continuous $g$ and every $Y$.  Statements of the form ``$N_\ast(y)=\nu(y)+
\littleo(1)$'' are therefore false as uniform statements, however natural they
look, and the volume asserts only the exact ceiling identity, recovery along
the range $y=A_n$, and the residue-class refinement.  Second, the error floor
of the asymptotic inverse is not a defect of the reversion but is inherited:
by \cref{p0:thm:optimal-truncation} the forward series has a least-term
envelope, and \cref{p0:thm:backward-error} converts it into a floor for the
inverse.  No amount of further reversion crosses it.

\section{What this consolidation establishes}\label{p5:sec:new}

Several statements in this volume are not in any of its twelve sources.  They
are listed here because a reader who knows the sources should be able to see
quickly what has changed; the corrections themselves are boxed where they
occur and collected in \cref{app:repairs}.

\begin{enumerate}[label=(\arabic*),leftmargin=*]
\item \emph{A correct proof of the Lagrange fixed-point formula.}  Five of the
      twelve articles derive
      $\Coef{t^n}U=\sum_{m\le n}\tfrac1m\Coef{t^nu^{m-1}}\Psi^m$ by
      introducing a marker $z$ and applying Lagrange--Bürmann in $z$.  That
      application is not available: Lagrange--Bürmann needs $\varphi(0)$
      invertible, and here $\varphi(0)=\Psi(t,0)\in tR[[t]]$ is a non-unit.
      The identity is nonetheless true, and \cref{p0:lem:LB-universal} proves
      it by universality --- it is a polynomial identity over $\Rat$ in the
      coefficients of $\varphi$, valid in a localization into which the
      relevant domain injects.  Every reversion in this volume rests on this
      lemma, so the repair is not cosmetic.
\item \emph{The branch rule} \cref{p0:cor:branch-rule}, with the fold point
      and threshold, replacing four case analyses.
\item \emph{The normalization constant of the flattened blocks.}  The sources
      state $P_n(0)=D_n$ in a case where $a=1$; the correct general statement
      is $P_n(0)=a^{n+1}c_n$ (\cref{p0:thm:flattening}(3)), and since $a\ne1$
      in every chapter of this volume the factor is needed throughout.
\item \emph{A quantitative obstruction for the staircase},
      \cref{p0:thm:staircase}(5), in place of an $\littleo(1)$ claim that is
      false uniformly.
\item \emph{Backward error as a certificate.}  \Cref{p0:thm:backward-error}
      upgrades the sources' conditional bound --- which presupposes that a
      root exists --- to an existence statement, by a sign change and the
      intermediate value theorem.
\item \emph{One closed form where the sources had three.}
      \Cref{p0:thm:pure-stirling} is proved equal to all three published forms
      of the pure Lambert block, and the equality was verified in exact
      rational arithmetic for $1\le n\le8$.
\item \emph{A non-recursive formula for the factorial-core corrections.}
      \Cref{p0:lem:lagrange-fp-quadratic} extends the fixed-point formula to
      $\Psi=h+f$ with $h\in u^2R[[u]]$, which is what \cref{p2:sec:top}
      needs; the source reaches the same coefficients only through an ad hoc
      compositional inverse.
\end{enumerate}

\section{The second four}
\label{p5:sec:second-four}

\Cref{p6:sec:top,p7:sec:top,p8:sec:top,p9:sec:top} add four subjects that are
not sequences of the first kind, and they change the picture in three ways.

\paragraph{A second core, and a second family.}
\Cref{p0:thm:lambert-core} inverts $aX+b\log X$, in which the logarithm is
\emph{added}.  Gamma, Barnes $G$ and the $K$-function need
$a x^{p}(\log x-\beta)$, in which the unknown \emph{multiplies} the logarithm;
that block is inverted exactly by $\Wlam$ too, but by a different
substitution, and its normalized slope is $v+1$ rather than $a+b/X$.
\Cref{p6:lem:core} supplies it, in the generality
$a x^{p}(\log x-\beta)^{r}$, and \cref{p6:thm:general} is the corresponding
inversion theorem.  Those three subjects form a family of their own, indexed
by $(a,p,\beta)$ and by where a resonant logarithm sits; the comparison is
\cref{p7:tab:three-subjects}.

\paragraph{A subject that borrows its core from another chapter.}
The subfactorial has no core of its own: $D_n=n!/\e+(-1)^{n}C(n)$, and the
inverse of the first term is the inverse of $\Gamma$, which is
\cref{p6:thm:gamma}.  Everything specifically about derangements lives in a
correction of size $\exp(-u\log u+\bigO(u))$, below every fixed exponential
scale $\e^{-cu}$ --- \cref{p8:prop:hierarchy}.  This is the first subject in
the volume whose interesting content is \emph{not} in its dominant block, and
it is the reason \cref{p8:sec:top} is the shortest chapter with the largest
citation load.

\paragraph{A subject the frame does not cover.}
The real-argument Fibonacci function has a pure exponential core, no slow
coefficient field, and no Lambert function anywhere.  Its inverse expansion
\emph{converges}.  \Cref{p9:sec:intro} tabulates exactly which parts of
Chapter~0 survive: essentially only
\cref{p0:thm:perturbed-inversion} and \cref{p0:thm:backward-error}.  It is
included because the purpose is the same --- invert a rapidly growing function
at infinity to all orders --- and because the contrast is informative: it
isolates which features of the other seven come from the \emph{logarithm} in
their phase rather than from inversion as such.

\begin{table}[htbp]
\centering
\caption{The eight subjects.  ``Core'' is the block inverted exactly;
``slope'' is the quantity every correction is divided by.}
\label{p5:tab:eight}
\small
\setlength{\tabcolsep}{4pt}
\begin{tabular}{@{}llll@{}}
\toprule
Chapter & Core & Slope & Correction scale\\
\midrule
\ref{p1:sec:top} rooted trees & $aX+b\log X$ & $a+b/X$ & $X^{-1}$, $\alpha=1$\\
\ref{p2:sec:top} double factorial & $X(\kappa\log X+d)$ & $\kappa\log X+\kappa+d$ & $X^{-1}$\\
\ref{p3:sec:top} partitions & $aX+b\log X$ & $a+b/X$ & $X^{-1}$, $\alpha=\tfrac12$\\
\ref{p4:sec:top} swing factorial & $aX+b\log X$ & $a+b/X$ & $X^{-1}$\\
\addlinespace
\ref{p6:sec:top} $\Gamma$ & $t(\log t-1)$ & $\Lambda=w+1$ & $T^{-2}$ (even)\\
\ref{p6:sec:top} Barnes $G$ & $\tfrac12z^{2}(\log z-\tfrac32)$ & $U\lambda$ & $U^{-1}$ (all)\\
\ref{p7:sec:top} $K$ & $\tfrac12r^{2}(\log r-\tfrac12)$ & $\rho L$ & $\rho^{-2}$ (even)\\
\addlinespace
\ref{p8:sec:top} subfactorial & $\Gamma(x+1)/\e$ & $Y\lambda$ & $\e^{\omega u}/Y$\\
\ref{p9:sec:top} Fibonacci & $\e^{\lambda x}$ & $\lambda$ & $R^{-1}$, oscillatory\\
\bottomrule
\end{tabular}
\end{table}

\section{Three things only the merge shows}
\label{p5:sec:merge-only}

\paragraph{Every subject with a logarithmic phase hides a square root.}
Four apparently unrelated square roots appear in the second half of the
volume:
\[
 \sqrt{\Lambda^{2}T^{2}+\tfrac1{12}},
 \qquad
 \sqrt{\lambda^{2}U^{2}-2\alpha U},
 \qquad
 \sqrt{U^{2}+\tfrac16},
 \qquad
 \sqrt{\rho^{2}+\tfrac1{12}} .
\]
The first two are the deepest pole in the slope, for $\Gamma$ and for $G$
(\cref{p6:thm:deepest-pole}); the third and fourth are coefficientwise
large-slope limits, for $G$ and for $K$
(\cref{p6:prop:resonance,p7:prop:limit}).  All four are the same lemma:
truncate the normalized equation to its linear and quadratic terms in the
displacement together with the lowest forcing monomial, and the truncation is
a quadratic, hence solvable in closed form
(\cref{p6:lem:quadratic-core}).  Which \emph{extraction} of the coefficient
family the quadratic computes --- deepest pole or large-slope limit --- is
decided by whether the forcing monomial carries a factor of the slope, that
is by whether it is resonant.  Resonance and pole depth are two readings of
one quadratic, and no source saw either connection.

\paragraph{Acceleration is resummation.}
The $K$-function sources offer an ``accelerated coordinate''
$T=x^{2}-x+\tfrac16$ that gains a full asymptotic sector, and present it as a
change of variable.  \Cref{p7:thm:acceleration-is-resummation} proves it is
the resummation of the resonant subsector of the unaccelerated expansion ---
the same square root, arrived at from the other side.  The absorption lemma
behind it (\cref{p7:lem:absorb}) applies verbatim to Barnes $G$, which no
Barnes source attempted; the resulting normal form is
\eqref{p7:eq:barnes-accelerated}.

\paragraph{One number, two jobs.}
The constant $\tfrac1{12}$ occurs as the $\Gamma$ deepest pole and again as
the $K$ resonant limit, for reasons that look unrelated.  Both are
$-2\cdot B_2(\tfrac12)/2$: for $\Gamma$, $B_2(\tfrac12)/2=-\tfrac1{24}$ is the
first Stirling tail coefficient; for $K$, the same value is the coefficient of
the resonant logarithm.  The identical Bernoulli value at the identical
reflection centre is doing both jobs (\cref{p7:rem:same-twelfth}).  The same
reflection symmetry is what makes the half-shift simultaneously remove
$r\log r$ and kill the odd tail, in \cref{p6:rem:why-half-shift} and
\cref{p7:rem:centre}.

\section{What the second wave removed}
\label{p5:sec:second-dedup}

The first twelve articles carried twelve copies of the Chapter~0 apparatus.
The second twelve carry, in addition, copies of \emph{each other}:

\begin{itemize}
\item The single block
$d_5=(2232\Lambda^{4}+1176\Lambda^{3}+714\Lambda^{2}+350\Lambda+105)
/(2903040\Lambda^{5})$ of \cref{p6:prop:gamma-blocks} is printed, in one
letter or another, in \emph{seven} of the twelve second-wave sources: the
three whose subject it is, all three subfactorial articles --- which need the
inverse of $\Gamma(x+1)/\e$ as their carrier and rederive it from scratch,
each proving its own Lambert-centred reversion theorem to get there --- and
one $K$-function article.  Four of those seven copies are removed here;
\cref{p8:prop:carrier} cites \cref{p6:thm:gamma} instead.
\item Eight of the twelve state a general triangular or formal-inversion
theorem over a graded coefficient ring.  They are all
\cref{p6:thm:triangular}, and one of them states it under a hypothesis
strictly stronger than its own proof uses.
\item Seven of the twelve state a named residual-to-error or a posteriori
transfer result.  They are all \cref{p0:thm:backward-error}, which the first
wave had already needed.
\item Seven of the twelve prove a Lagrange--B\"urmann reversion adapted to
their setting.  In every case the content is
\cref{p0:thm:perturbed-inversion} or, where the perturbation is formal,
\cref{p0:cor:perturbed-formal}.
\end{itemize}

\noindent
That the second wave's sources duplicate the first wave's is the strongest
single piece of evidence that the consolidation was worth doing: the twelve
new articles were written without knowledge of the twelve old ones, and
converged on the same apparatus independently.

\section{Open questions}\label{p5:sec:open}

\begin{enumerate}[label=(\roman*),leftmargin=*]
\item \emph{A Stokes-transport theorem.}  The mechanism by which an
      exponentially small forward sector becomes an exponentially small
      inverse sector is uniform --- it is the perturbation-operator form
      \cref{p0:prop:operator-form} --- but a theorem needs each subject's
      Borel geometry as a hypothesis.  The four chapters supply four separate
      instances; the general statement is not proved here.
\item \emph{The counting function.}  One source states that the number of
      indices with $A_n\le y$ equals the average of the branch inverses plus
      $\bigO(1)$.  The statement is plausible for every subject with a parity
      split, but it needs a per-subject monotonicity hypothesis and is not
      proved in this volume.
\item \emph{Rademacher-type sequences in general.}  \Cref{p3:sec:top} gives a
      template for inverting a sequence presented as a convergent sum over
      arithmetic sectors.  Whether the template applies to the other classical
      Rademacher-type expansions is untested here.
\item \emph{Multifactorials.}  One double-factorial source sketches the
      $m$-fold generalization $n!^{(m)}$.  The core is again of factorial
      type, so \cref{p0:thm:core-reversion} applies verbatim; what is not done
      is the residue-class analysis of the staircase for general $m$, which
      \cref{p0:thm:staircase}(4) reduces to a finite computation but does not
      carry out.
\item \emph{Relation to the formalized Lambert layer of this repository.}  The
      surrounding development contains a Lean-verified treatment of both real
      branches of $\Wlam$, a lower-branch bracket with an explicit remainder
      rate, and the saddle expansions of the Fabius function.  The Lambert
      core these four chapters are built on is therefore already formalized to
      a degree none of the sources assumes.  Which of the reversion machinery
      above is within reach of that layer --- \cref{p0:thm:lambert-core} and
      \cref{p0:cor:branch-rule} look closest --- is an open and, in this
      repository, a natural question.
\end{enumerate}

\appendix

\chapter{Provenance}
\label{q4:app:provenance}

This volume consolidates \emph{forty-two} independently written articles,
filed in this repository between 1 and 3 September 2026.  Thirty of them had
already been consolidated once, into the two volumes that became the calculus
parts and the inversion parts; those two carry their own provenance records,
which are reproduced unchanged inside this volume.  The remaining twelve are
merged here for the first time and are listed below.

\begin{center}\small
\begin{longtable}{@{}>{\raggedright\arraybackslash}p{0.36\textwidth}
                    >{\raggedright\arraybackslash}p{0.13\textwidth}
                    >{\raggedright\arraybackslash}p{0.15\textwidth}
                    >{\raggedright\arraybackslash}p{0.28\textwidth}@{}}
\toprule
\textbf{Source directory} & \textbf{Source} & \textbf{PDF} & \textbf{Absorbed into}\\
\midrule
\endfirsthead
\toprule
\textbf{Source directory} & \textbf{Source} & \textbf{PDF} & \textbf{Absorbed into}\\
\midrule
\endhead
\multicolumn{4}{@{}l@{}}{\emph{Orientation, from \texttt{transseries-tutorials/}}}\\
\texttt{Transseries\_\allowbreak Tutorial-1} & 5\,159 lines & 143 pp, Letter & \cref{q0:sec:why}, \cref{q0:sec:algebra}\\
\texttt{Transseries\_\allowbreak Tutorial-2} & 7\,749 lines & 164 pp, Letter & \cref{q0:sec:why}, \cref{q0:sec:algebra}\\
\texttt{Transseries\_\allowbreak Tutorial-3} & 4\,410 lines & 121 pp, Letter & \cref{q0:sec:why}\\
\texttt{Transseries\_\allowbreak Tutorial-4} & 8\,781 lines & 217 pp, $522\times738$\,pt & \cref{q0:sec:why}, \cref{q0:sec:algebra}\\
\addlinespace
\multicolumn{4}{@{}l@{}}{\emph{Reversing $x+\LambertW(x)$, from \texttt{lambert-inverse-transseries/}}}\\
\texttt{lambert\_\allowbreak inverse\_\allowbreak transseries} & 1\,907 lines & 25 pp, A4 & \cref{q1:sec:top}\\
\texttt{lambert\_\allowbreak inverse\_\allowbreak transseries\_\allowbreak bundle} & 1\,384 lines & 21 pp, Letter & \cref{q1:sec:top}\\
\texttt{reversing\_\allowbreak x\_\allowbreak plus\_\allowbreak lambert\_\allowbreak w\_\allowbreak transseries} & 1\,918 lines & 29 pp, A4 & \cref{q1:sec:top}\\
\addlinespace
\multicolumn{4}{@{}l@{}}{\emph{Bell and Fubini numbers, from \texttt{sequence-transseries/}}}\\
\texttt{Bell\_\allowbreak Number\_\allowbreak Asymptotic\_\allowbreak Transseries} & 2\,211 lines & 32 pp, A4 & \cref{q2:sec:bell}\\
\texttt{Bell\_\allowbreak Number\_\allowbreak Transseries\_\allowbreak Article} & 1\,714 lines & 23 pp, A4 & \cref{q2:sec:bell}\\
\texttt{Fubini\_\allowbreak Number\_\allowbreak Full\_\allowbreak Transseries} & 2\,465 lines & 33 pp, A4 & \cref{q2:sec:fubini}\\
\texttt{Fubini\_\allowbreak Number\_\allowbreak Transseries} & 1\,881 lines & 25 pp, A4 & \cref{q2:sec:fubini}\\
\texttt{Fubini\_\allowbreak Number\_\allowbreak Transseries\_\allowbreak Article} & 1\,472 lines & 25 pp, Letter & \cref{q2:sec:fubini}\\
\bottomrule
\end{longtable}
\end{center}

\noindent
The other thirty are documented in place.  The six sources of the calculus
parts are listed, with a notation concordance and a repair ledger, in the
final part of the calculus itself; the twenty-four sources of the inversion
parts are listed in \cref{app:provenance} below, which is the
provenance appendix of that volume reproduced unchanged.  All forty-two
remain in git history.

\section{Counts}
\label{q4:sec:counts}

\begin{center}\small
\begin{tabular}{@{}lrrl@{}}
\toprule
Group & Articles & Source lines & Absorbed as\\
\midrule
\texttt{transseries-tutorials/} & 4 & 26\,099 & Orientation\\
\texttt{polynomial-logarithmic-transseries/} & 6 & 36\,033 & the calculus parts\\
\texttt{special-function-inversion/} & 24 & 16\,771 & the inversion parts\\
\texttt{lambert-inverse-transseries/} & 3 & 5\,209 & $x+\LambertW(x)$\\
\texttt{sequence-transseries/} & 5 & 9\,743 & Bell and Fubini\\
\midrule
total & 42 & --- & this volume\\
\bottomrule
\end{tabular}
\end{center}

\noindent
The line counts in the middle column are those of the \emph{immediate}
sources: for the two groups that had already been consolidated, the figure is
the size of the consolidated volume, not of the six or twenty-four articles
behind it.

\chapter{Repairs made in this merge}
\label{q4:app:repairs-merge}

Corrections made while merging the twelve newly consolidated sources.  Repairs
made in the two earlier consolidations are listed in their own ledgers, which
this volume reproduces: the calculus's in its final part, the inversion
volume's in \cref{app:repairs}.

\paragraph{An error law quoted without its rate.}
All three $x+\LambertW(x)$ sources state
$g(z)-G_N(z)\sim(N+1)^{-1}\bigl(\LambertW(z)/z\bigr)^{N+1}$ and stop there.
The law is correct --- \cref{q1:thm:error-law} proves it --- but it is
approached at relative rate only $\bigO(1/\log z)$, so it is a poor predictor
at any accessible argument: the ratio of the true error to the stated law is
$0.55$ at $z=10^{3}$, $0.72$ at $10^{5}$, $0.82$ at $10^{8}$, $0.88$ at
$10^{12}$ and $0.93$ at $10^{20}$, and at $z=50$ the law is wrong by a factor
of three.  \Cref{q1:rem:slow} states the rate and recommends the uniform
two-sided certificate instead.

\paragraph{A structure theorem stated in the wrong normalization.}
Two of the same three sources normalize the correction coefficients
differently --- one in powers of $q=\LambertW(z)/z$, the other in powers of
$z^{-1}$ --- and the integer-polynomial structure theorem is true only of the
second family.  Carried across unchanged it fails at the first coefficient.
\Cref{q1:rem:normalization} records the dictionary
$\mathcal A_n=A_n/L^{n}$, and the chapter uses throughout the normalization in
which the theorem holds.

\paragraph{Five theorems stated without proof, demoted.}
The tutorials state real closedness of the transseries field, closure under
antidifferentiation, a formal Taylor principle, a formal contraction principle
and generalized Borel reconstruction as named theorems, three of them
explicitly labelled ``stated without proof'' or ``schematic''.  They are
recorded in \cref{q0:sec:quoted} as attributed quotations rather than as
theorems of this volume, together with the observation --- checked --- that no
result proved here depends on any of them.

\paragraph{A hand-wavy structure argument, replaced.}
The Bell sources assert that the saddle coefficients have the form
$c_m=P_m/(1+r)^{3m}$ with $\deg P_m\le4m$ without exhibiting the cancellation
that makes it true.  \Cref{q2:prop:bell-structure} proves it by tracking the
factor $r^{K}$ that cancels between the Touchard numerators and the
denominators of the $q_j$, after which the degree count is exact and index
independent.

\paragraph{An overstatement made by this merge, and corrected in it.}
Working from a title-level concordance, an intermediate stage of this
consolidation recorded that the calculus already contained most of the
inversion apparatus.  Reading the statements showed otherwise, and
\cref{q3:sec:top} is the corrected comparison: of six apparent overlaps only
one is a genuine duplicate, one is a strengthening in the inversion
apparatus's favour, one has no counterpart at all, and three are distinct
theorems.  The episode is recorded in \cref{q3:rem:correction} rather than
quietly fixed, because the failure mode --- treating shared vocabulary as
evidence of shared content --- is one a reader of a merged volume should be
warned about.


\chapter{Provenance of the inversion sources}\label{app:provenance}

This volume consolidates twenty-four independently written articles, three on
each of its eight subjects.  All twenty-four were filed in this repository on
3~September 2026 as a single quick-intake batch, at
\texttt{drafts/series-and-transseries/special-function-inversion/}; the
directories below are relative to that path.  Nothing in the twelve was
edited in place: this volume is a new document, and the sources remain
available in git history.

\begin{center}
\scriptsize
\begin{longtable}{@{}>{\raggedright\arraybackslash}p{0.40\textwidth}>{\raggedright\arraybackslash}p{0.12\textwidth}>{\raggedright\arraybackslash}p{0.13\textwidth}>{\raggedright\arraybackslash}p{0.23\textwidth}@{}}
\toprule
\textbf{Source directory} & \textbf{Source} & \textbf{PDF} & \textbf{Absorbed into} \\
\midrule
\endfirsthead
\toprule
\textbf{Source directory} & \textbf{Source} & \textbf{PDF} & \textbf{Absorbed into} \\
\midrule
\endhead
\multicolumn{4}{@{}l@{}}{\emph{Chapter~\ref{p1:sec:top}: Butcher--Pólya rooted trees}}\\
\texttt{Butcher\_\allowbreak Tree\_\allowbreak Transseries} & 1\,674 lines & 21 pp, A4 & Chapter~\ref{p1:sec:top} \\
\texttt{Butcher\_\allowbreak Tree\_\allowbreak Transseries-2} & 2\,853 lines & 62 pp, A4 & Chapter~\ref{p1:sec:top} \\
\texttt{Butcher\_\allowbreak Tree\_\allowbreak Transseries-3} & 2\,382 lines & 35 pp, A4 & Chapter~\ref{p1:sec:top} \\
\addlinespace
\multicolumn{4}{@{}l@{}}{\emph{Chapter~\ref{p2:sec:top}: the double factorial}}\\
\texttt{Double\_\allowbreak Factorial\_\allowbreak Transseries} & 1\,552 lines & 18 pp, A4 & Chapter~\ref{p2:sec:top} \\
\texttt{Double\_\allowbreak Factorial\_\allowbreak Transseries-3} & 2\,287 lines & 28 pp, A4 & Chapter~\ref{p2:sec:top} \\
\texttt{double\_\allowbreak factorial\_\allowbreak transseries-2} & 2\,138 lines & 44 pp, A4 & Chapter~\ref{p2:sec:top} \\
\addlinespace
\multicolumn{4}{@{}l@{}}{\emph{Chapter~\ref{p3:sec:top}: the partition numbers}}\\
\texttt{Partition\_\allowbreak Number\_\allowbreak Transseries\_\allowbreak and\_\allowbreak Asymptotic\_\allowbreak Inverse} & 1\,901 lines & 32 pp, A4 & Chapter~\ref{p3:sec:top} \\
\texttt{Partition\_\allowbreak Number\_\allowbreak Transseries\_\allowbreak and\_\allowbreak Inverse} & 2\,020 lines & 27 pp, A4 & Chapter~\ref{p3:sec:top} \\
\texttt{Partition\_\allowbreak Numbers\_\allowbreak Transseries\_\allowbreak and\_\allowbreak Inverse} & 1\,905 lines & 33 pp, Letter & Chapter~\ref{p3:sec:top} \\
\addlinespace
\multicolumn{4}{@{}l@{}}{\emph{Chapter~\ref{p4:sec:top}: the swing factorial}}\\
\texttt{Swing\_\allowbreak Factorial\_\allowbreak Transseries} & 1\,885 lines & 28 pp, A4 & Chapter~\ref{p4:sec:top} \\
\texttt{Swing\_\allowbreak Factorial\_\allowbreak Transseries\_\allowbreak Article} & 1\,824 lines & 25 pp, A4 & Chapters~\ref{p0:sec:top} and~\ref{p4:sec:top} \\
\texttt{Swing\_\allowbreak Factorial\_\allowbreak Transseries\_\allowbreak and\_\allowbreak Inverse} & 1\,389 lines & 22 pp, A4 & Chapter~\ref{p4:sec:top} \\
\addlinespace
\multicolumn{4}{@{}l@{}}{\emph{Chapter~\ref{p6:sec:top}: Gamma and Barnes $G$}}\\
\texttt{Asymptotic\_\allowbreak Inversion\_\allowbreak Gamma\_\allowbreak Barnes\_\allowbreak G} & 2\,376 lines & 29 pp, A4 & Chapters~\ref{p6:sec:top} and~\ref{p7:sec:top} \\
\texttt{inverse\_\allowbreak gamma\_\allowbreak barnesG\_\allowbreak transseries} & 1\,655 lines & 28 pp, A4 & Chapter~\ref{p6:sec:top} \\
\texttt{inverse\_\allowbreak gamma\_\allowbreak barnes\_\allowbreak transseries} & 1\,827 lines & 25 pp, A4 & Chapter~\ref{p6:sec:top} \\
\addlinespace
\multicolumn{4}{@{}l@{}}{\emph{Chapter~\ref{p7:sec:top}: the hyperfactorial $K$-function}}\\
\texttt{inverse\_\allowbreak k\_\allowbreak function\_\allowbreak transseries} & 2\,259 lines & 29 pp, A4 & Chapter~\ref{p7:sec:top} \\
\texttt{K\_\allowbreak Function\_\allowbreak Inverse\_\allowbreak Transseries} & 2\,644 lines & 33 pp, A4 & Chapter~\ref{p7:sec:top} \\
\texttt{K\_\allowbreak function\_\allowbreak inverse\_\allowbreak transseries\_\allowbreak article} & 2\,581 lines & 30 pp, A4 & Chapter~\ref{p7:sec:top} \\
\addlinespace
\multicolumn{4}{@{}l@{}}{\emph{Chapter~\ref{p8:sec:top}: the subfactorial}}\\
\texttt{inverse\_\allowbreak subfactorial\_\allowbreak transseries} & 2\,631 lines & 38 pp, A4 & Chapter~\ref{p8:sec:top} \\
\texttt{inverse\_\allowbreak subfactorial\_\allowbreak transseries-2} & 2\,188 lines & 27 pp, A4 & Chapter~\ref{p8:sec:top} \\
\texttt{inverse\_\allowbreak subfactorial\_\allowbreak transseries-3} & 1\,878 lines & 35 pp, Letter & Chapter~\ref{p8:sec:top} \\
\addlinespace
\multicolumn{4}{@{}l@{}}{\emph{Chapter~\ref{p9:sec:top}: a real-argument Fibonacci function}}\\
\texttt{Fibonacci\_\allowbreak Inverse\_\allowbreak LogPeriodic\_\allowbreak Transseries} & 2\,554 lines & 33 pp, A4 & Chapter~\ref{p9:sec:top} \\
\texttt{fibonacci\_\allowbreak inverse\_\allowbreak transseries\_\allowbreak article} & 2\,857 lines & 36 pp, A4 & Chapter~\ref{p9:sec:top} \\
\texttt{fibonacci\_\allowbreak inverse\_\allowbreak transseries\_\allowbreak article-2} & 1\,884 lines & 24 pp, A4 & Chapter~\ref{p9:sec:top} \\
\bottomrule
\end{longtable}
\end{center}

\noindent
The shared apparatus of Chapter~\ref{p0:sec:top} is drawn from all
twenty-four.  Its
organizing normal form, and the Lagrange fixed-point route to the reversion
coefficients, follow the presentation in
\texttt{Swing\_\allowbreak Factorial\_\allowbreak Transseries\_\allowbreak Article} most closely; the
\(\alpha\)-reduction that brings the partition numbers into the same frame is
this volume's, since no source needed it.

\chapter{Repairs to the inversion sources}\label{app:repairs}

Every correction this volume makes to one of its twenty-four sources is boxed
in the text where it occurs and listed here.  A repair is recorded when a source
states something false, states an asymptotic conclusion that its argument does
not support, or asserts a named result without proof.  Differences of
convention are not repairs and are recorded in the text as remarks instead.

\section{Chapter~\ref{p0:sec:top}: the common apparatus}

\paragraph{A proof that does not apply (five sources).}
\textsf{Swing\_\allowbreak Factorial\_\allowbreak Transseries\_\allowbreak Article},
\textsf{Partition\_\allowbreak Number\_\allowbreak Transseries\_\allowbreak and\_\allowbreak Inverse} (twice),
\textsf{Partition\_\allowbreak Numbers\_\allowbreak Transseries\_\allowbreak and\_\allowbreak Inverse},
\textsf{Partition\_\allowbreak Number\_\allowbreak Transseries\_\allowbreak and\_\allowbreak Asymptotic\_\allowbreak Inverse} and
\textsf{double\_\allowbreak factorial\_\allowbreak transseries-2} all prove the Lagrange fixed-point
formula
\(\Coef{t^n}U=\sum_{m\le n}\tfrac1m\Coef{t^nu^{m-1}}\Psi^m\)
by introducing an auxiliary marker \(z\), solving \(U=z\Psi(t,U)\), and
applying Lagrange--Bürmann in \(z\).  Lagrange--Bürmann requires the relevant
\(\varphi(0)\) to be invertible; here \(\varphi(0)=\Psi(t,0)\in tR[[t]]\) is a
non-unit, and the residue proof genuinely breaks, since \(\psi=w/\varphi(w)\)
is then not an element of \(w\,R[[t]][[w]]\).  The identity is true;
\cref{p0:lem:LB-universal} proves it by universality instead.  Every reversion
in this volume depends on the lemma, so the repair is structural rather than
cosmetic.

\paragraph{An asymptotic claim that is false uniformly.}
Several sources state, or leave the reader to infer, that the integer inverse
satisfies \(N_\ast(y)=\nu(y)+\littleo(1)\).  No continuous function
approximates a staircase uniformly: \cref{p0:thm:staircase}(5) shows
\(\sup_{y\ge Y}\lvert N_\ast(y)-g(y)\rvert\ge\tfrac12\) for every continuous
\(g\) and every \(Y\).  The volume asserts only the exact ceiling identity,
nearest-integer recovery \emph{along the range}, and the residue-class
refinement.

\paragraph{A normalization constant.}
The sources record \(P_n(0)=D_n\) for the flattened blocks, in a case where
\(a=1\).  The general statement is \(P_n(0)=a^{\,n+1}c_n\)
(\cref{p0:thm:flattening}(3)); since \(a\neq1\) in every chapter of this
volume the factor is needed throughout.

\paragraph{Two truncation conventions in one article.}
\textsf{Swing\_\allowbreak Factorial\_\allowbreak Transseries\_\allowbreak Article} uses the exclusive convention
in its forward error table and the inclusive one in its inverse table, so
those two tables are not on a common scale.  The volume is inclusive
throughout (\cref{p0:rem:conventions}); tables inherited from that article's
forward section are re-indexed.

\paragraph{A claim of this volume's own, corrected.}
The brief under which \cref{p0:sec:top} was written asserted that the case
\(\alpha=1\) of \cref{p0:def:model} covers the double factorial.  It does not:
\(\log\bigl((x+1)!!\bigr)=\tfrac{s}{2}(\log s-1)+\bigO(\log s)\) with
\(s=x+1\), and \(s\log s/s^{\alpha}\to\infty\) for every \(\alpha\le1\).  The
sources never claimed membership.  \Cref{p0:def:core} axiomatizes the
dominant core instead, and \cref{p0:prop:factorial-core} inverts the factorial
phase exactly.

\paragraph{Corrupt source text, not reproduced.}
\textsf{Swing\_\allowbreak Factorial\_\allowbreak Transseries\_\allowbreak Article} equation
\texttt{eq:inverse-\allowbreak Stokes-\allowbreak leading} and
\textsf{Butcher\_\allowbreak Tree\_\allowbreak Transseries-2} lines 2482--2483 contain corrupt control
sequences.  The affected statements are re-derived here rather than copied.

\section{Chapter~\ref{p2:sec:top}: the double factorial}

\paragraph{Convergence asserted without a threshold.}
The first source states that the periodic carrier ``has the convergent local
Lagrange--Bürmann expansion'' for sufficiently large \(\Lambda\), with no
threshold and no proof of convergence as opposed to formal solvability; the
second proves only a contraction estimate, from which convergence of the
Lagrange series does not follow.  \Cref{p2:thm:carrier} separates the two:
existence and uniqueness of the real root by contraction for
\(\Lambda>2\pi a\), and convergence of the Lagrange series under the explicit
sufficient condition \(\mu=2a/\Lambda\le1/5\).

\paragraph{Two assertions in place of proofs.}
The deepest-pole coefficient (``only the quadratic kernel and the first
Bernoulli coefficient can contribute'') is exactly the step needing the degree
bookkeeping, and is supplied here by induction on the
\(\Lambda^{-1}\)-degree.  The coefficientwise sign law is argued in one source
with ``use a positive rooted-tree contribution \dots; such a tree exists'',
exhibiting no such object; \cref{p2:thm:sign-law} keeps the correct positivity
reduction and replaces the strictness step with an explicit attainability
induction, verified symbolically for \(n\le7\).

\paragraph{A formal-to-analytic slide.}
The third source presents an all-orders phase reversion as a theorem with an
asymptotic conclusion, introduced by ``setting \(\tau=1\) as an asymptotic
expansion''.  The identity is formal; making it asymptotic needs estimates,
uniform in the oscillatory phase, that are neither stated nor proved.  It is
demoted here to the formal \cref{p2:prop:phase-reversion}, with a proved
remainder supplied separately.

\paragraph{A silent three-way inconsistency.}
One source centres the periodic inverse at the even-branch constant, the other
two at the parity mean; the leading displacements consequently differ by
\(2a/\Lambda\).  Both are correct for their own centring, but no source says
so, and a reader comparing them would conclude that one is wrong.  The volume
adopts the mean centring and states the relation.

\paragraph{A notation collision across three sources.}
The letter \(C\) denotes the multiplicative amplitude, the logarithmic
constant, and the branch prefactor in the three siblings respectively, and
\(\alpha,\beta\) carry two incompatible meanings.  A crosswalk table fixes the
volume's usage; the core parameters are never abbreviated, to avoid a fourth
collision with \(\kappa\) in \cref{p0:def:core}.

\paragraph{A qualitative warning made quantitative.}
All three sources warn that a different smooth interpolation ``can change this
sector'' without exhibiting one.  \Cref{p2:thm:no-canonical} exhibits the
family \(\mathcal D(x)\e^{\lambda\sin\pi x}\), which interpolates \(n!!\) at
every non-negative integer and changes the leading oscillatory amplitude from
\(2a/\Lambda\) to \(2\sqrt{a^2+\lambda^2}/\Lambda\) --- a change larger than
the entire Bernoulli grid.  This is what makes the chapter's organizing claim
a theorem rather than a caution.

\section{Chapter~\ref{p1:sec:top}: the rooted-tree numbers}

\paragraph{Digits printed beyond the accuracy of the computation.}
One source prints the constants of the inherited \(-\sqrt\rho\) sector to
forty decimal digits, but only about twenty-five of them are correct.  Two
independent recomputations here --- including a from-scratch direct summation
at eighty digits --- give
\[
  r_0^{-}=0.468636627924845548087202621931819314\ldots,
\]
and the volume prints only digits it has verified.  This is the one place in
the twelve articles where a printed numeral is wrong; it is a precision
failure rather than a mistake in a formula.

\paragraph{An audit appendix that contradicts its own main text.}
The same source's audit appendix reports \(D_3=-3.582177326832123\ldots\),
disagreeing with the value its main text derives.  The main text is correct;
the appendix entry is not, and the volume uses the recomputed value.

\paragraph{A symbol used for two different constants.}
One source writes \(\alpha\) for the growth constant where the other two write
\(\lambda\), and the volume's Chapter~\ref{p0:sec:top} uses \(\alpha\) for the
phase exponent of \cref{p0:def:model}.  The chapter fixes \(\lambda\) for the
growth constant throughout and never uses \(\alpha\) for it.

\paragraph{An asymptotic statement with a formal proof.}
A source states its inverse-coefficient theorem asymptotically but proves it
only formally, and separately displays a global transseries with no hypothesis
under which it could hold.  The first is proved here in its asymptotic form;
the second is replaced by a conditional theorem with a finite radius and an
explicit hypothesis.

\paragraph{A circular derivation of the singularity bounds.}
The elementary bounds \(\tfrac14\le\rho\le\e^{-1}\) are used in the sources in
a way that presupposes what they are used to prove.  They are established here
independently, from \(x\le T\e^{-T}\le\e^{-1}\) and the Catalan bound, and the
critical characterization is then obtained non-circularly.

\section{Chapter~\ref{p4:sec:top}: the swing factorial}

\paragraph{A divergent series written as an equality.}
One source sets its bookkeeping coupling to \(1\) and prints the resulting
operator expansion as an equality.  There is then no small parameter, the sum
diverges, and a number is equated to a divergent series; as printed the
statement is false.  The \(m\)-th summand is \(\bigO(X^{-(2m-1)})\), so the
correct statement is asymptotic, and regrouping gives the finite differential
form used here.  The other two sources state the same expansion correctly.

\paragraph{A theorem whose hypothesis is never used.}
A sector-transport result is labelled a theorem, its hypothesis is not used in
the argument, and its conclusion is prefaced ``formally''.  It is split here
into the formal all-orders identity, which is what the proof gives, and a
quantitative first-order statement proved by the mean value theorem under
stated hypotheses.

\paragraph{An overclaim about complex branches.}
One source asserts that every Lambert branch with \(\lvert w\rvert\to\infty\)
yields an inverse branch with the same coefficients.  For \(k\ne0,-1\) the
existence of an analytic branch of the inverse along the path is a hypothesis,
not a consequence, and it fails at gamma poles and Stokes directions.  The
claim is demoted to a remark with the two hypotheses displayed.

\paragraph{The analytic inverse statement, unproved in all three sources.}
One source writes ``likewise, for each fixed \(M\), \dots'' with no proof; the
second sketches an argument that never bounds the residual of the truncated
ansatz nor uses forward control at the perturbed point; the third proves the
ingredients but never assembles them.  This is the one formal-to-analytic
slide common to all three.  \Cref{p4:thm:inverse-analytic} proves the
statement in full, with an effective constant.

\paragraph{A recommendation contradicted by its own table.}
A remark asserts that the centred normalization is ``often the preferable
computational normal form''.  The same article's table shows the centred core
better by a factor \(\approx10^{3}\) at one end and the uncentred form
slightly better at the other.  The corrected statement is that centring kills
the \(X^{-1}\) term and so improves the \emph{bare core} by a factor
\(\asymp X\), while with several corrections retained the two are comparable
and neither dominates uniformly.

\paragraph{Two convergence claims without criteria.}
An equivalence of the Lambert and logarithmic normalizations is asserted in
prose with no formula and no proof; it is proved here as a finite explicit
composition.  A convergence claim appeals to ``the classical analysis of
Lambert \(\Wlam\)'', but the classical result concerns a doubly indexed
series, whereas the series at hand has coefficients polynomial in
\(\ell\) and converges only when \(\ell/L\) is small; an explicit criterion is
proved instead.

\section{Chapter~\ref{p3:sec:top}: the partition numbers}

\paragraph{A definition that contradicts its own use.}
The source abbreviated \textsf{S2} in \cref{p3:sec:top} defines the real
Rademacher amplitude by a sum over \(1\le h\le k\) with \(\gcd(h,k)=1\).  For
\(k\ge2\) that is the intended index set, but at \(k=1\) it returns
\(\cos(2\pi\nu)\) rather than the constant \(1\), while the same article uses
\(A_1\equiv1\) throughout.  With the printed definition the interpolant would
not interpolate \(p\): its dominant term would oscillate off-lattice at full
relative strength, and the article's own eventual-monotonicity statement would
be false.  The volume takes \(h\in U_k\) with \(U_1=\{0\}\); every later
formula of that source is then correct as printed.

\paragraph{An understated radius.}
\textsf{S1} states only that its universal sector-coefficient series
``converges for all sufficiently small \(\lvert t\rvert\)''.  The radius is
exactly \(\sqrt{24}\) in \(t=n^{-1/2}\), the nearest singularities being the
branch points of \(\sqrt{1-t^2/24}\); convergence therefore holds for every
real \(n>1/24\).

\paragraph{A convergent series and an asymptotic total, conflated.}
All three partition sources are loose about the distinction.  The algebraic
series \(\sum_m\omega_mn^{-m/2}\) converges, but its sum is the leading
Rademacher block, \emph{not} \(p(n)\); the Poincaré statement about \(p(n)\)
carries an error that is the sum of a convergent algebraic tail and an
exponentially small Rademacher remainder.  Summing the algebraic series to
convergence does not compute \(p(n)\): the residual stalls at the
\(\e^{-\mu/2}\) floor.

\paragraph{A false intermediate bound.}
\textsf{S1}'s kernel estimate asserts \(g(z)\le z^2/3\) on \([0,1]\), which
fails at \(z=1\) where \(g(1)=\e^{-1}=0.3679\ldots>1/3\).  The chapter proves
\(g(z)\le(z^2/3)\cosh z\) instead, which is what the tail bound actually
needs.

\paragraph{An overclaim in this volume's own Chapter~\ref{p0:sec:top}.}
\Cref{p0:rem:formal-analytic} originally said that the correction series \(Q\)
is divergent of Gevrey type ``in all four chapters''.  It is not: for the
partition numbers \(Q(t)=\log(1-t)\) has radius of convergence \(1\).  The
remark now distinguishes the three divergent chapters from
\cref{p3:sec:top}, where \cref{p0:thm:optimal-truncation} has no content and
the error floor is set by the Rademacher remainder instead.

\bigskip
\noindent\textbf{\large The second wave: Chapters~\ref{p6:sec:top}--\ref{p9:sec:top}}
\medskip

\paragraph{An existence claim that no source makes effective.}
All six power--logarithmic sources justify the eventual inverse of Barnes
\(G\) and of the \(K\)-function with the words ``\((\log G)'(x)\sim x\log x\),
hence strictly increasing for all sufficiently large \(x\)''.  That is true
and ineffective: none of the six supplies a threshold, and yet every
numerical table in all six evaluates the inverse at arguments as small as
\(z=5\), inside the range their own justification does not cover.
\Cref{p6:prop:barnes-monotone} proves that \(z\mapsto G(z+1)\) is strictly
increasing on \([3,\infty)\) with slope at least \(\tfrac12z\log z\) on
\([7,\infty)\), and \cref{p7:prop:branch} does the same for \(K\) on
\([2,\infty)\) with the two stationary points located.  The tabulated
arguments are then legitimate.

\paragraph{A hypothesis stronger than the proof needs.}
\textsf{S1} of \cref{p6:sec:top} states its triangular inversion theorem
under the hypothesis that the support semigroup be \emph{locally finite},
meaning that every bounded interval meets it in a finite set.  That is
strictly stronger than well-ordering --- \(\{1-1/n\}\cup\{1\}\) is well
ordered and not locally finite --- and more than the argument uses, which is
only the finite-decomposition property.  Neumann's lemma supplies that for
every well-ordered set, so \cref{p6:thm:triangular} is stated with the weaker
hypothesis.  For the two subjects of that chapter the distinction is
invisible, since the supports are \(2\Nat\) and \(\Nat\); it matters for the
fractional and mixed grids of \cref{p6:sec:general}.

\paragraph{A limit presented as a resummation.}
\textsf{S1} of \cref{p6:sec:top} writes that the leading constant subsector of
the Barnes coefficients ``resums'' to \(\sqrt{T^{2}+1/6}\), and its next
remark correctly warns that this does not replace the expansion; the two
statements sit awkwardly together.  The accurate form is
\cref{p6:prop:resonance}: a limit of coefficients, each taken separately as
the slope grows, which is a structural identity and an exact check on a
symbolic computation but not an approximation to the inverse --- at any actual
argument the neglected terms, headed by \(b_0=-\alpha/\lambda\), are the
largest corrections present.  \Cref{p6:rem:limit-caveat} states this, and
\cref{p7:rem:barnes-too} then shows what does turn the limit into an
approximation: moving it into the coordinate.

\paragraph{A carrier rederived three times.}
All three subfactorial sources need the inverse of \(\Gamma(x+1)/\e\) as the
carrier for everything else, and all three derive it from scratch, each
proving its own Lambert-centred reversion theorem and each printing the
blocks \(1/(24s)\), \(-(14s^{2}+10s+5)/(5760s^{3})\),
\((2232s^{4}+\cdots+105)/(2903040s^{5})\).  These are identical, coefficient
for coefficient, to \(d_1,d_3,d_5\) of \cref{p6:prop:gamma-blocks}.  The
block \(d_5\) is in fact printed in seven of the twelve second-wave sources.
\Cref{p8:prop:carrier} states the dictionary
\(u(Y)=\Gamma_+^{-1}(\e Y)-1\) and cites \cref{p6:thm:gamma}; the general
reversion theorem the three sources prove to get there is
\cref{p0:thm:perturbed-inversion}.  This is not an error in any source --- each
is correct in isolation --- but it is the largest single duplication the
consolidation removes.

\paragraph{An unverified radius, quoted by neither.}
\textsf{V1} of \cref{p9:sec:top} states an exact radius of convergence for the
Lagrange--Good series.  The proof of \cref{p9:thm:good} establishes
convergence for \(R\) beyond a constant depending only on \(\rho\), by a crude
bound on the binomial factors, and the sharp radius was not verified for this
volume.  No value is therefore quoted, and the question is recorded as open in
\cref{p9:sec:open}(ii).  This is listed here because a reader comparing the
volume with its source will find a number in one and not the other, and the
reason is deliberate.



\end{document}
